% Punto de entrada aditivo del sucesor; el main.tex rector queda inmutable.
\RequirePackage{fix-cm}
\documentclass[11pt,a4paper,twoside,openany]{book}

\usepackage[T1]{fontenc}
\usepackage[utf8]{inputenc}
\DeclareUnicodeCharacter{1D49}{\textsuperscript{e}}
\DeclareUnicodeCharacter{02B3}{\textsuperscript{r}}
\DeclareUnicodeCharacter{03BC}{\ensuremath{\mu}}
\DeclareUnicodeCharacter{03C4}{\ensuremath{\tau}}
\DeclareUnicodeCharacter{03C3}{\ensuremath{\sigma}}
\DeclareUnicodeCharacter{03C0}{\ensuremath{\pi}}
\DeclareUnicodeCharacter{03B3}{\ensuremath{\gamma}}
\DeclareUnicodeCharacter{03A9}{\ensuremath{\Omega}}
\DeclareUnicodeCharacter{03C1}{\ensuremath{\rho}}
\DeclareUnicodeCharacter{03C6}{\ensuremath{\varphi}}
\DeclareUnicodeCharacter{03C8}{\ensuremath{\psi}}
\DeclareUnicodeCharacter{03A5}{\ensuremath{\Upsilon}}
\DeclareUnicodeCharacter{03B7}{\ensuremath{\eta}}
\DeclareUnicodeCharacter{03B1}{\ensuremath{\alpha}}
\DeclareUnicodeCharacter{03B5}{\ensuremath{\varepsilon}}
\DeclareUnicodeCharacter{03B6}{\ensuremath{\zeta}}
\DeclareUnicodeCharacter{03B8}{\ensuremath{\theta}}
\DeclareUnicodeCharacter{03BB}{\ensuremath{\lambda}}
\DeclareUnicodeCharacter{00B1}{\ensuremath{\pm}}
\DeclareUnicodeCharacter{2192}{\ensuremath{\longrightarrow}}
\DeclareUnicodeCharacter{2194}{\ensuremath{\leftrightarrow}}
\DeclareUnicodeCharacter{21A6}{\ensuremath{\longmapsto}}
\DeclareUnicodeCharacter{21D2}{\ensuremath{\Rightarrow}}
\DeclareUnicodeCharacter{2202}{\ensuremath{\partial}}
% lcdf-typetools glyphtounicode.tex, Version 2.95
% Contents: Glyph mapping information for pdftex, used for PDF searching
% Generated from:
% - glyphlist.txt, Version 2.0
% - texglyphlist.txt, Version 2.95
% - texglyphlist-g2u.txt, Version 2.95
\pdfglyphtounicode{A}{0041}
\pdfglyphtounicode{AE}{00C6}
\pdfglyphtounicode{AEacute}{01FC}
\pdfglyphtounicode{AEmacron}{01E2}
\pdfglyphtounicode{AEsmall}{00E6}
\pdfglyphtounicode{Aacute}{00C1}
\pdfglyphtounicode{Aacutesmall}{00E1}
\pdfglyphtounicode{Abreve}{0102}
\pdfglyphtounicode{Abreveacute}{1EAE}
\pdfglyphtounicode{Abrevecyrillic}{04D0}
\pdfglyphtounicode{Abrevedotbelow}{1EB6}
\pdfglyphtounicode{Abrevegrave}{1EB0}
\pdfglyphtounicode{Abrevehookabove}{1EB2}
\pdfglyphtounicode{Abrevetilde}{1EB4}
\pdfglyphtounicode{Acaron}{01CD}
\pdfglyphtounicode{Acircle}{24B6}
\pdfglyphtounicode{Acircumflex}{00C2}
\pdfglyphtounicode{Acircumflexacute}{1EA4}
\pdfglyphtounicode{Acircumflexdotbelow}{1EAC}
\pdfglyphtounicode{Acircumflexgrave}{1EA6}
\pdfglyphtounicode{Acircumflexhookabove}{1EA8}
\pdfglyphtounicode{Acircumflexsmall}{00E2}
\pdfglyphtounicode{Acircumflextilde}{1EAA}
\pdfglyphtounicode{Acute}{00B4}
\pdfglyphtounicode{Acutesmall}{00B4}
\pdfglyphtounicode{Acyrillic}{0410}
\pdfglyphtounicode{Adblgrave}{0200}
\pdfglyphtounicode{Adieresis}{00C4}
\pdfglyphtounicode{Adieresiscyrillic}{04D2}
\pdfglyphtounicode{Adieresismacron}{01DE}
\pdfglyphtounicode{Adieresissmall}{00E4}
\pdfglyphtounicode{Adotbelow}{1EA0}
\pdfglyphtounicode{Adotmacron}{01E0}
\pdfglyphtounicode{Agrave}{00C0}
\pdfglyphtounicode{Agravesmall}{00E0}
\pdfglyphtounicode{Ahookabove}{1EA2}
\pdfglyphtounicode{Aiecyrillic}{04D4}
\pdfglyphtounicode{Ainvertedbreve}{0202}
\pdfglyphtounicode{Alpha}{0391}
\pdfglyphtounicode{Alphatonos}{0386}
\pdfglyphtounicode{Amacron}{0100}
\pdfglyphtounicode{Amonospace}{FF21}
\pdfglyphtounicode{Aogonek}{0104}
\pdfglyphtounicode{Aring}{00C5}
\pdfglyphtounicode{Aringacute}{01FA}
\pdfglyphtounicode{Aringbelow}{1E00}
\pdfglyphtounicode{Aringsmall}{00E5}
\pdfglyphtounicode{Asmall}{0061}
\pdfglyphtounicode{Atilde}{00C3}
\pdfglyphtounicode{Atildesmall}{00E3}
\pdfglyphtounicode{Aybarmenian}{0531}
\pdfglyphtounicode{B}{0042}
\pdfglyphtounicode{Bcircle}{24B7}
\pdfglyphtounicode{Bdotaccent}{1E02}
\pdfglyphtounicode{Bdotbelow}{1E04}
\pdfglyphtounicode{Becyrillic}{0411}
\pdfglyphtounicode{Benarmenian}{0532}
\pdfglyphtounicode{Beta}{0392}
\pdfglyphtounicode{Bhook}{0181}
\pdfglyphtounicode{Blinebelow}{1E06}
\pdfglyphtounicode{Bmonospace}{FF22}
\pdfglyphtounicode{Brevesmall}{02D8}
\pdfglyphtounicode{Bsmall}{0062}
\pdfglyphtounicode{Btopbar}{0182}
\pdfglyphtounicode{C}{0043}
\pdfglyphtounicode{Caarmenian}{053E}
\pdfglyphtounicode{Cacute}{0106}
\pdfglyphtounicode{Caron}{02C7}
\pdfglyphtounicode{Caronsmall}{02C7}
\pdfglyphtounicode{Ccaron}{010C}
\pdfglyphtounicode{Ccedilla}{00C7}
\pdfglyphtounicode{Ccedillaacute}{1E08}
\pdfglyphtounicode{Ccedillasmall}{00E7}
\pdfglyphtounicode{Ccircle}{24B8}
\pdfglyphtounicode{Ccircumflex}{0108}
\pdfglyphtounicode{Cdot}{010A}
\pdfglyphtounicode{Cdotaccent}{010A}
\pdfglyphtounicode{Cedillasmall}{00B8}
\pdfglyphtounicode{Chaarmenian}{0549}
\pdfglyphtounicode{Cheabkhasiancyrillic}{04BC}
\pdfglyphtounicode{Checyrillic}{0427}
\pdfglyphtounicode{Chedescenderabkhasiancyrillic}{04BE}
\pdfglyphtounicode{Chedescendercyrillic}{04B6}
\pdfglyphtounicode{Chedieresiscyrillic}{04F4}
\pdfglyphtounicode{Cheharmenian}{0543}
\pdfglyphtounicode{Chekhakassiancyrillic}{04CB}
\pdfglyphtounicode{Cheverticalstrokecyrillic}{04B8}
\pdfglyphtounicode{Chi}{03A7}
\pdfglyphtounicode{Chook}{0187}
\pdfglyphtounicode{Circumflexsmall}{02C6}
\pdfglyphtounicode{Cmonospace}{FF23}
\pdfglyphtounicode{Coarmenian}{0551}
\pdfglyphtounicode{Csmall}{0063}
\pdfglyphtounicode{D}{0044}
\pdfglyphtounicode{DZ}{01F1}
\pdfglyphtounicode{DZcaron}{01C4}
\pdfglyphtounicode{Daarmenian}{0534}
\pdfglyphtounicode{Dafrican}{0189}
\pdfglyphtounicode{Dbar}{0110}
\pdfglyphtounicode{Dcaron}{010E}
\pdfglyphtounicode{Dcedilla}{1E10}
\pdfglyphtounicode{Dcircle}{24B9}
\pdfglyphtounicode{Dcircumflexbelow}{1E12}
\pdfglyphtounicode{Dcroat}{0110}
\pdfglyphtounicode{Ddotaccent}{1E0A}
\pdfglyphtounicode{Ddotbelow}{1E0C}
\pdfglyphtounicode{Decyrillic}{0414}
\pdfglyphtounicode{Deicoptic}{03EE}
\pdfglyphtounicode{Delta}{2206}
\pdfglyphtounicode{Deltagreek}{0394}
\pdfglyphtounicode{Dhook}{018A}
\pdfglyphtounicode{Dieresis}{00A8}
\pdfglyphtounicode{DieresisAcute}{F6CC}
\pdfglyphtounicode{DieresisGrave}{F6CD}
\pdfglyphtounicode{Dieresissmall}{00A8}
\pdfglyphtounicode{Digamma}{D875 DFCB}
\pdfglyphtounicode{Digammagreek}{03DC}
\pdfglyphtounicode{Djecyrillic}{0402}
\pdfglyphtounicode{Dlinebelow}{1E0E}
\pdfglyphtounicode{Dmonospace}{FF24}
\pdfglyphtounicode{Dotaccentsmall}{02D9}
\pdfglyphtounicode{Dslash}{0110}
\pdfglyphtounicode{Dsmall}{0064}
\pdfglyphtounicode{Dtopbar}{018B}
\pdfglyphtounicode{Dz}{01F2}
\pdfglyphtounicode{Dzcaron}{01C5}
\pdfglyphtounicode{Dzeabkhasiancyrillic}{04E0}
\pdfglyphtounicode{Dzecyrillic}{0405}
\pdfglyphtounicode{Dzhecyrillic}{040F}
\pdfglyphtounicode{E}{0045}
\pdfglyphtounicode{Eacute}{00C9}
\pdfglyphtounicode{Eacutesmall}{00E9}
\pdfglyphtounicode{Ebreve}{0114}
\pdfglyphtounicode{Ecaron}{011A}
\pdfglyphtounicode{Ecedillabreve}{1E1C}
\pdfglyphtounicode{Echarmenian}{0535}
\pdfglyphtounicode{Ecircle}{24BA}
\pdfglyphtounicode{Ecircumflex}{00CA}
\pdfglyphtounicode{Ecircumflexacute}{1EBE}
\pdfglyphtounicode{Ecircumflexbelow}{1E18}
\pdfglyphtounicode{Ecircumflexdotbelow}{1EC6}
\pdfglyphtounicode{Ecircumflexgrave}{1EC0}
\pdfglyphtounicode{Ecircumflexhookabove}{1EC2}
\pdfglyphtounicode{Ecircumflexsmall}{00EA}
\pdfglyphtounicode{Ecircumflextilde}{1EC4}
\pdfglyphtounicode{Ecyrillic}{0404}
\pdfglyphtounicode{Edblgrave}{0204}
\pdfglyphtounicode{Edieresis}{00CB}
\pdfglyphtounicode{Edieresissmall}{00EB}
\pdfglyphtounicode{Edot}{0116}
\pdfglyphtounicode{Edotaccent}{0116}
\pdfglyphtounicode{Edotbelow}{1EB8}
\pdfglyphtounicode{Efcyrillic}{0424}
\pdfglyphtounicode{Egrave}{00C8}
\pdfglyphtounicode{Egravesmall}{00E8}
\pdfglyphtounicode{Eharmenian}{0537}
\pdfglyphtounicode{Ehookabove}{1EBA}
\pdfglyphtounicode{Eightroman}{2167}
\pdfglyphtounicode{Einvertedbreve}{0206}
\pdfglyphtounicode{Eiotifiedcyrillic}{0464}
\pdfglyphtounicode{Elcyrillic}{041B}
\pdfglyphtounicode{Elevenroman}{216A}
\pdfglyphtounicode{Emacron}{0112}
\pdfglyphtounicode{Emacronacute}{1E16}
\pdfglyphtounicode{Emacrongrave}{1E14}
\pdfglyphtounicode{Emcyrillic}{041C}
\pdfglyphtounicode{Emonospace}{FF25}
\pdfglyphtounicode{Encyrillic}{041D}
\pdfglyphtounicode{Endescendercyrillic}{04A2}
\pdfglyphtounicode{Eng}{014A}
\pdfglyphtounicode{Enghecyrillic}{04A4}
\pdfglyphtounicode{Enhookcyrillic}{04C7}
\pdfglyphtounicode{Eogonek}{0118}
\pdfglyphtounicode{Eopen}{0190}
\pdfglyphtounicode{Epsilon}{0395}
\pdfglyphtounicode{Epsilontonos}{0388}
\pdfglyphtounicode{Ercyrillic}{0420}
\pdfglyphtounicode{Ereversed}{018E}
\pdfglyphtounicode{Ereversedcyrillic}{042D}
\pdfglyphtounicode{Escyrillic}{0421}
\pdfglyphtounicode{Esdescendercyrillic}{04AA}
\pdfglyphtounicode{Esh}{01A9}
\pdfglyphtounicode{Esmall}{0065}
\pdfglyphtounicode{Eta}{0397}
\pdfglyphtounicode{Etarmenian}{0538}
\pdfglyphtounicode{Etatonos}{0389}
\pdfglyphtounicode{Eth}{00D0}
\pdfglyphtounicode{Ethsmall}{00F0}
\pdfglyphtounicode{Etilde}{1EBC}
\pdfglyphtounicode{Etildebelow}{1E1A}
\pdfglyphtounicode{Euro}{20AC}
\pdfglyphtounicode{Ezh}{01B7}
\pdfglyphtounicode{Ezhcaron}{01EE}
\pdfglyphtounicode{Ezhreversed}{01B8}
\pdfglyphtounicode{F}{0046}
\pdfglyphtounicode{FFIsmall}{0066 0066 0069}
\pdfglyphtounicode{FFLsmall}{0066 0066 006C}
\pdfglyphtounicode{FFsmall}{0066 0066}
\pdfglyphtounicode{FIsmall}{0066 0069}
\pdfglyphtounicode{FLsmall}{0066 006C}
\pdfglyphtounicode{Fcircle}{24BB}
\pdfglyphtounicode{Fdotaccent}{1E1E}
\pdfglyphtounicode{Feharmenian}{0556}
\pdfglyphtounicode{Feicoptic}{03E4}
\pdfglyphtounicode{Fhook}{0191}
\pdfglyphtounicode{Finv}{2132}
\pdfglyphtounicode{Fitacyrillic}{0472}
\pdfglyphtounicode{Fiveroman}{2164}
\pdfglyphtounicode{Fmonospace}{FF26}
\pdfglyphtounicode{Fourroman}{2163}
\pdfglyphtounicode{Fsmall}{0066}
\pdfglyphtounicode{G}{0047}
\pdfglyphtounicode{GBsquare}{3387}
\pdfglyphtounicode{Gacute}{01F4}
\pdfglyphtounicode{Gamma}{0393}
\pdfglyphtounicode{Gammaafrican}{0194}
\pdfglyphtounicode{Gangiacoptic}{03EA}
\pdfglyphtounicode{Gbreve}{011E}
\pdfglyphtounicode{Gcaron}{01E6}
\pdfglyphtounicode{Gcedilla}{0122}
\pdfglyphtounicode{Gcircle}{24BC}
\pdfglyphtounicode{Gcircumflex}{011C}
\pdfglyphtounicode{Gcommaaccent}{0122}
\pdfglyphtounicode{Gdot}{0120}
\pdfglyphtounicode{Gdotaccent}{0120}
\pdfglyphtounicode{Gecyrillic}{0413}
\pdfglyphtounicode{Germandbls}{0053 0053}
\pdfglyphtounicode{Germandblssmall}{0073 0073}
\pdfglyphtounicode{Ghadarmenian}{0542}
\pdfglyphtounicode{Ghemiddlehookcyrillic}{0494}
\pdfglyphtounicode{Ghestrokecyrillic}{0492}
\pdfglyphtounicode{Gheupturncyrillic}{0490}
\pdfglyphtounicode{Ghook}{0193}
\pdfglyphtounicode{Gimarmenian}{0533}
\pdfglyphtounicode{Gjecyrillic}{0403}
\pdfglyphtounicode{Gmacron}{1E20}
\pdfglyphtounicode{Gmir}{2141}
\pdfglyphtounicode{Gmonospace}{FF27}
\pdfglyphtounicode{Grave}{0060}
\pdfglyphtounicode{Gravesmall}{0060}
\pdfglyphtounicode{Gsmall}{0067}
\pdfglyphtounicode{Gsmallhook}{029B}
\pdfglyphtounicode{Gstroke}{01E4}
\pdfglyphtounicode{H}{0048}
\pdfglyphtounicode{H18533}{25CF}
\pdfglyphtounicode{H18543}{25AA}
\pdfglyphtounicode{H18551}{25AB}
\pdfglyphtounicode{H22073}{25A1}
\pdfglyphtounicode{HPsquare}{33CB}
\pdfglyphtounicode{Haabkhasiancyrillic}{04A8}
\pdfglyphtounicode{Hadescendercyrillic}{04B2}
\pdfglyphtounicode{Hardsigncyrillic}{042A}
\pdfglyphtounicode{Hbar}{0126}
\pdfglyphtounicode{Hbrevebelow}{1E2A}
\pdfglyphtounicode{Hcedilla}{1E28}
\pdfglyphtounicode{Hcircle}{24BD}
\pdfglyphtounicode{Hcircumflex}{0124}
\pdfglyphtounicode{Hdieresis}{1E26}
\pdfglyphtounicode{Hdotaccent}{1E22}
\pdfglyphtounicode{Hdotbelow}{1E24}
\pdfglyphtounicode{Hmonospace}{FF28}
\pdfglyphtounicode{Hoarmenian}{0540}
\pdfglyphtounicode{Horicoptic}{03E8}
\pdfglyphtounicode{Hsmall}{0068}
\pdfglyphtounicode{Hungarumlaut}{02DD}
\pdfglyphtounicode{Hungarumlautsmall}{02DD}
\pdfglyphtounicode{Hzsquare}{3390}
\pdfglyphtounicode{I}{0049}
\pdfglyphtounicode{IAcyrillic}{042F}
\pdfglyphtounicode{IJ}{0132}
\pdfglyphtounicode{IUcyrillic}{042E}
\pdfglyphtounicode{Iacute}{00CD}
\pdfglyphtounicode{Iacutesmall}{00ED}
\pdfglyphtounicode{Ibreve}{012C}
\pdfglyphtounicode{Icaron}{01CF}
\pdfglyphtounicode{Icircle}{24BE}
\pdfglyphtounicode{Icircumflex}{00CE}
\pdfglyphtounicode{Icircumflexsmall}{00EE}
\pdfglyphtounicode{Icyrillic}{0406}
\pdfglyphtounicode{Idblgrave}{0208}
\pdfglyphtounicode{Idieresis}{00CF}
\pdfglyphtounicode{Idieresisacute}{1E2E}
\pdfglyphtounicode{Idieresiscyrillic}{04E4}
\pdfglyphtounicode{Idieresissmall}{00EF}
\pdfglyphtounicode{Idot}{0130}
\pdfglyphtounicode{Idotaccent}{0130}
\pdfglyphtounicode{Idotbelow}{1ECA}
\pdfglyphtounicode{Iebrevecyrillic}{04D6}
\pdfglyphtounicode{Iecyrillic}{0415}
\pdfglyphtounicode{Ifractur}{2111}
\pdfglyphtounicode{Ifraktur}{2111}
\pdfglyphtounicode{Igrave}{00CC}
\pdfglyphtounicode{Igravesmall}{00EC}
\pdfglyphtounicode{Ihookabove}{1EC8}
\pdfglyphtounicode{Iicyrillic}{0418}
\pdfglyphtounicode{Iinvertedbreve}{020A}
\pdfglyphtounicode{Iishortcyrillic}{0419}
\pdfglyphtounicode{Imacron}{012A}
\pdfglyphtounicode{Imacroncyrillic}{04E2}
\pdfglyphtounicode{Imonospace}{FF29}
\pdfglyphtounicode{Iniarmenian}{053B}
\pdfglyphtounicode{Iocyrillic}{0401}
\pdfglyphtounicode{Iogonek}{012E}
\pdfglyphtounicode{Iota}{0399}
\pdfglyphtounicode{Iotaafrican}{0196}
\pdfglyphtounicode{Iotadieresis}{03AA}
\pdfglyphtounicode{Iotatonos}{038A}
\pdfglyphtounicode{Ismall}{0069}
\pdfglyphtounicode{Istroke}{0197}
\pdfglyphtounicode{Itilde}{0128}
\pdfglyphtounicode{Itildebelow}{1E2C}
\pdfglyphtounicode{Izhitsacyrillic}{0474}
\pdfglyphtounicode{Izhitsadblgravecyrillic}{0476}
\pdfglyphtounicode{J}{004A}
\pdfglyphtounicode{Jaarmenian}{0541}
\pdfglyphtounicode{Jcircle}{24BF}
\pdfglyphtounicode{Jcircumflex}{0134}
\pdfglyphtounicode{Jecyrillic}{0408}
\pdfglyphtounicode{Jheharmenian}{054B}
\pdfglyphtounicode{Jmonospace}{FF2A}
\pdfglyphtounicode{Jsmall}{006A}
\pdfglyphtounicode{K}{004B}
\pdfglyphtounicode{KBsquare}{3385}
\pdfglyphtounicode{KKsquare}{33CD}
\pdfglyphtounicode{Kabashkircyrillic}{04A0}
\pdfglyphtounicode{Kacute}{1E30}
\pdfglyphtounicode{Kacyrillic}{041A}
\pdfglyphtounicode{Kadescendercyrillic}{049A}
\pdfglyphtounicode{Kahookcyrillic}{04C3}
\pdfglyphtounicode{Kappa}{039A}
\pdfglyphtounicode{Kastrokecyrillic}{049E}
\pdfglyphtounicode{Kaverticalstrokecyrillic}{049C}
\pdfglyphtounicode{Kcaron}{01E8}
\pdfglyphtounicode{Kcedilla}{0136}
\pdfglyphtounicode{Kcircle}{24C0}
\pdfglyphtounicode{Kcommaaccent}{0136}
\pdfglyphtounicode{Kdotbelow}{1E32}
\pdfglyphtounicode{Keharmenian}{0554}
\pdfglyphtounicode{Kenarmenian}{053F}
\pdfglyphtounicode{Khacyrillic}{0425}
\pdfglyphtounicode{Kheicoptic}{03E6}
\pdfglyphtounicode{Khook}{0198}
\pdfglyphtounicode{Kjecyrillic}{040C}
\pdfglyphtounicode{Klinebelow}{1E34}
\pdfglyphtounicode{Kmonospace}{FF2B}
\pdfglyphtounicode{Koppacyrillic}{0480}
\pdfglyphtounicode{Koppagreek}{03DE}
\pdfglyphtounicode{Ksicyrillic}{046E}
\pdfglyphtounicode{Ksmall}{006B}
\pdfglyphtounicode{L}{004C}
\pdfglyphtounicode{LJ}{01C7}
\pdfglyphtounicode{LL}{004C 004C}
\pdfglyphtounicode{Lacute}{0139}
\pdfglyphtounicode{Lambda}{039B}
\pdfglyphtounicode{Lcaron}{013D}
\pdfglyphtounicode{Lcedilla}{013B}
\pdfglyphtounicode{Lcircle}{24C1}
\pdfglyphtounicode{Lcircumflexbelow}{1E3C}
\pdfglyphtounicode{Lcommaaccent}{013B}
\pdfglyphtounicode{Ldot}{013F}
\pdfglyphtounicode{Ldotaccent}{013F}
\pdfglyphtounicode{Ldotbelow}{1E36}
\pdfglyphtounicode{Ldotbelowmacron}{1E38}
\pdfglyphtounicode{Liwnarmenian}{053C}
\pdfglyphtounicode{Lj}{01C8}
\pdfglyphtounicode{Ljecyrillic}{0409}
\pdfglyphtounicode{Llinebelow}{1E3A}
\pdfglyphtounicode{Lmonospace}{FF2C}
\pdfglyphtounicode{Lslash}{0141}
\pdfglyphtounicode{Lslashsmall}{0142}
\pdfglyphtounicode{Lsmall}{006C}
\pdfglyphtounicode{M}{004D}
\pdfglyphtounicode{MBsquare}{3386}
\pdfglyphtounicode{Macron}{00AF}
\pdfglyphtounicode{Macronsmall}{00AF}
\pdfglyphtounicode{Macute}{1E3E}
\pdfglyphtounicode{Mcircle}{24C2}
\pdfglyphtounicode{Mdotaccent}{1E40}
\pdfglyphtounicode{Mdotbelow}{1E42}
\pdfglyphtounicode{Menarmenian}{0544}
\pdfglyphtounicode{Mmonospace}{FF2D}
\pdfglyphtounicode{Msmall}{006D}
\pdfglyphtounicode{Mturned}{019C}
\pdfglyphtounicode{Mu}{039C}
\pdfglyphtounicode{N}{004E}
\pdfglyphtounicode{NJ}{01CA}
\pdfglyphtounicode{Nacute}{0143}
\pdfglyphtounicode{Ncaron}{0147}
\pdfglyphtounicode{Ncedilla}{0145}
\pdfglyphtounicode{Ncircle}{24C3}
\pdfglyphtounicode{Ncircumflexbelow}{1E4A}
\pdfglyphtounicode{Ncommaaccent}{0145}
\pdfglyphtounicode{Ndotaccent}{1E44}
\pdfglyphtounicode{Ndotbelow}{1E46}
\pdfglyphtounicode{Ng}{014A}
\pdfglyphtounicode{Nhookleft}{019D}
\pdfglyphtounicode{Nineroman}{2168}
\pdfglyphtounicode{Nj}{01CB}
\pdfglyphtounicode{Njecyrillic}{040A}
\pdfglyphtounicode{Nlinebelow}{1E48}
\pdfglyphtounicode{Nmonospace}{FF2E}
\pdfglyphtounicode{Nowarmenian}{0546}
\pdfglyphtounicode{Nsmall}{006E}
\pdfglyphtounicode{Ntilde}{00D1}
\pdfglyphtounicode{Ntildesmall}{00F1}
\pdfglyphtounicode{Nu}{039D}
\pdfglyphtounicode{O}{004F}
\pdfglyphtounicode{OE}{0152}
\pdfglyphtounicode{OEsmall}{0153}
\pdfglyphtounicode{Oacute}{00D3}
\pdfglyphtounicode{Oacutesmall}{00F3}
\pdfglyphtounicode{Obarredcyrillic}{04E8}
\pdfglyphtounicode{Obarreddieresiscyrillic}{04EA}
\pdfglyphtounicode{Obreve}{014E}
\pdfglyphtounicode{Ocaron}{01D1}
\pdfglyphtounicode{Ocenteredtilde}{019F}
\pdfglyphtounicode{Ocircle}{24C4}
\pdfglyphtounicode{Ocircumflex}{00D4}
\pdfglyphtounicode{Ocircumflexacute}{1ED0}
\pdfglyphtounicode{Ocircumflexdotbelow}{1ED8}
\pdfglyphtounicode{Ocircumflexgrave}{1ED2}
\pdfglyphtounicode{Ocircumflexhookabove}{1ED4}
\pdfglyphtounicode{Ocircumflexsmall}{00F4}
\pdfglyphtounicode{Ocircumflextilde}{1ED6}
\pdfglyphtounicode{Ocyrillic}{041E}
\pdfglyphtounicode{Odblacute}{0150}
\pdfglyphtounicode{Odblgrave}{020C}
\pdfglyphtounicode{Odieresis}{00D6}
\pdfglyphtounicode{Odieresiscyrillic}{04E6}
\pdfglyphtounicode{Odieresissmall}{00F6}
\pdfglyphtounicode{Odotbelow}{1ECC}
\pdfglyphtounicode{Ogoneksmall}{02DB}
\pdfglyphtounicode{Ograve}{00D2}
\pdfglyphtounicode{Ogravesmall}{00F2}
\pdfglyphtounicode{Oharmenian}{0555}
\pdfglyphtounicode{Ohm}{2126}
\pdfglyphtounicode{Ohookabove}{1ECE}
\pdfglyphtounicode{Ohorn}{01A0}
\pdfglyphtounicode{Ohornacute}{1EDA}
\pdfglyphtounicode{Ohorndotbelow}{1EE2}
\pdfglyphtounicode{Ohorngrave}{1EDC}
\pdfglyphtounicode{Ohornhookabove}{1EDE}
\pdfglyphtounicode{Ohorntilde}{1EE0}
\pdfglyphtounicode{Ohungarumlaut}{0150}
\pdfglyphtounicode{Oi}{01A2}
\pdfglyphtounicode{Oinvertedbreve}{020E}
\pdfglyphtounicode{Omacron}{014C}
\pdfglyphtounicode{Omacronacute}{1E52}
\pdfglyphtounicode{Omacrongrave}{1E50}
\pdfglyphtounicode{Omega}{2126}
\pdfglyphtounicode{Omegacyrillic}{0460}
\pdfglyphtounicode{Omegagreek}{03A9}
\pdfglyphtounicode{Omegainv}{2127}
\pdfglyphtounicode{Omegaroundcyrillic}{047A}
\pdfglyphtounicode{Omegatitlocyrillic}{047C}
\pdfglyphtounicode{Omegatonos}{038F}
\pdfglyphtounicode{Omicron}{039F}
\pdfglyphtounicode{Omicrontonos}{038C}
\pdfglyphtounicode{Omonospace}{FF2F}
\pdfglyphtounicode{Oneroman}{2160}
\pdfglyphtounicode{Oogonek}{01EA}
\pdfglyphtounicode{Oogonekmacron}{01EC}
\pdfglyphtounicode{Oopen}{0186}
\pdfglyphtounicode{Oslash}{00D8}
\pdfglyphtounicode{Oslashacute}{01FE}
\pdfglyphtounicode{Oslashsmall}{00F8}
\pdfglyphtounicode{Osmall}{006F}
\pdfglyphtounicode{Ostrokeacute}{01FE}
\pdfglyphtounicode{Otcyrillic}{047E}
\pdfglyphtounicode{Otilde}{00D5}
\pdfglyphtounicode{Otildeacute}{1E4C}
\pdfglyphtounicode{Otildedieresis}{1E4E}
\pdfglyphtounicode{Otildesmall}{00F5}
\pdfglyphtounicode{P}{0050}
\pdfglyphtounicode{Pacute}{1E54}
\pdfglyphtounicode{Pcircle}{24C5}
\pdfglyphtounicode{Pdotaccent}{1E56}
\pdfglyphtounicode{Pecyrillic}{041F}
\pdfglyphtounicode{Peharmenian}{054A}
\pdfglyphtounicode{Pemiddlehookcyrillic}{04A6}
\pdfglyphtounicode{Phi}{03A6}
\pdfglyphtounicode{Phook}{01A4}
\pdfglyphtounicode{Pi}{03A0}
\pdfglyphtounicode{Piwrarmenian}{0553}
\pdfglyphtounicode{Pmonospace}{FF30}
\pdfglyphtounicode{Psi}{03A8}
\pdfglyphtounicode{Psicyrillic}{0470}
\pdfglyphtounicode{Psmall}{0070}
\pdfglyphtounicode{Q}{0051}
\pdfglyphtounicode{Qcircle}{24C6}
\pdfglyphtounicode{Qmonospace}{FF31}
\pdfglyphtounicode{Qsmall}{0071}
\pdfglyphtounicode{R}{0052}
\pdfglyphtounicode{Raarmenian}{054C}
\pdfglyphtounicode{Racute}{0154}
\pdfglyphtounicode{Rcaron}{0158}
\pdfglyphtounicode{Rcedilla}{0156}
\pdfglyphtounicode{Rcircle}{24C7}
\pdfglyphtounicode{Rcommaaccent}{0156}
\pdfglyphtounicode{Rdblgrave}{0210}
\pdfglyphtounicode{Rdotaccent}{1E58}
\pdfglyphtounicode{Rdotbelow}{1E5A}
\pdfglyphtounicode{Rdotbelowmacron}{1E5C}
\pdfglyphtounicode{Reharmenian}{0550}
\pdfglyphtounicode{Rfractur}{211C}
\pdfglyphtounicode{Rfraktur}{211C}
\pdfglyphtounicode{Rho}{03A1}
\pdfglyphtounicode{Ringsmall}{02DA}
\pdfglyphtounicode{Rinvertedbreve}{0212}
\pdfglyphtounicode{Rlinebelow}{1E5E}
\pdfglyphtounicode{Rmonospace}{FF32}
\pdfglyphtounicode{Rsmall}{0072}
\pdfglyphtounicode{Rsmallinverted}{0281}
\pdfglyphtounicode{Rsmallinvertedsuperior}{02B6}
\pdfglyphtounicode{S}{0053}
\pdfglyphtounicode{SF010000}{250C}
\pdfglyphtounicode{SF020000}{2514}
\pdfglyphtounicode{SF030000}{2510}
\pdfglyphtounicode{SF040000}{2518}
\pdfglyphtounicode{SF050000}{253C}
\pdfglyphtounicode{SF060000}{252C}
\pdfglyphtounicode{SF070000}{2534}
\pdfglyphtounicode{SF080000}{251C}
\pdfglyphtounicode{SF090000}{2524}
\pdfglyphtounicode{SF100000}{2500}
\pdfglyphtounicode{SF110000}{2502}
\pdfglyphtounicode{SF190000}{2561}
\pdfglyphtounicode{SF200000}{2562}
\pdfglyphtounicode{SF210000}{2556}
\pdfglyphtounicode{SF220000}{2555}
\pdfglyphtounicode{SF230000}{2563}
\pdfglyphtounicode{SF240000}{2551}
\pdfglyphtounicode{SF250000}{2557}
\pdfglyphtounicode{SF260000}{255D}
\pdfglyphtounicode{SF270000}{255C}
\pdfglyphtounicode{SF280000}{255B}
\pdfglyphtounicode{SF360000}{255E}
\pdfglyphtounicode{SF370000}{255F}
\pdfglyphtounicode{SF380000}{255A}
\pdfglyphtounicode{SF390000}{2554}
\pdfglyphtounicode{SF400000}{2569}
\pdfglyphtounicode{SF410000}{2566}
\pdfglyphtounicode{SF420000}{2560}
\pdfglyphtounicode{SF430000}{2550}
\pdfglyphtounicode{SF440000}{256C}
\pdfglyphtounicode{SF450000}{2567}
\pdfglyphtounicode{SF460000}{2568}
\pdfglyphtounicode{SF470000}{2564}
\pdfglyphtounicode{SF480000}{2565}
\pdfglyphtounicode{SF490000}{2559}
\pdfglyphtounicode{SF500000}{2558}
\pdfglyphtounicode{SF510000}{2552}
\pdfglyphtounicode{SF520000}{2553}
\pdfglyphtounicode{SF530000}{256B}
\pdfglyphtounicode{SF540000}{256A}
\pdfglyphtounicode{SS}{0053 0053}
\pdfglyphtounicode{SSsmall}{0073 0073}
\pdfglyphtounicode{Sacute}{015A}
\pdfglyphtounicode{Sacutedotaccent}{1E64}
\pdfglyphtounicode{Sampigreek}{03E0}
\pdfglyphtounicode{Scaron}{0160}
\pdfglyphtounicode{Scarondotaccent}{1E66}
\pdfglyphtounicode{Scaronsmall}{0161}
\pdfglyphtounicode{Scedilla}{015E}
\pdfglyphtounicode{Schwa}{018F}
\pdfglyphtounicode{Schwacyrillic}{04D8}
\pdfglyphtounicode{Schwadieresiscyrillic}{04DA}
\pdfglyphtounicode{Scircle}{24C8}
\pdfglyphtounicode{Scircumflex}{015C}
\pdfglyphtounicode{Scommaaccent}{0218}
\pdfglyphtounicode{Sdotaccent}{1E60}
\pdfglyphtounicode{Sdotbelow}{1E62}
\pdfglyphtounicode{Sdotbelowdotaccent}{1E68}
\pdfglyphtounicode{Seharmenian}{054D}
\pdfglyphtounicode{Sevenroman}{2166}
\pdfglyphtounicode{Shaarmenian}{0547}
\pdfglyphtounicode{Shacyrillic}{0428}
\pdfglyphtounicode{Shchacyrillic}{0429}
\pdfglyphtounicode{Sheicoptic}{03E2}
\pdfglyphtounicode{Shhacyrillic}{04BA}
\pdfglyphtounicode{Shimacoptic}{03EC}
\pdfglyphtounicode{Sigma}{03A3}
\pdfglyphtounicode{Sixroman}{2165}
\pdfglyphtounicode{Smonospace}{FF33}
\pdfglyphtounicode{Softsigncyrillic}{042C}
\pdfglyphtounicode{Ssmall}{0073}
\pdfglyphtounicode{Stigmagreek}{03DA}
\pdfglyphtounicode{T}{0054}
\pdfglyphtounicode{Tau}{03A4}
\pdfglyphtounicode{Tbar}{0166}
\pdfglyphtounicode{Tcaron}{0164}
\pdfglyphtounicode{Tcedilla}{0162}
\pdfglyphtounicode{Tcircle}{24C9}
\pdfglyphtounicode{Tcircumflexbelow}{1E70}
\pdfglyphtounicode{Tcommaaccent}{0162}
\pdfglyphtounicode{Tdotaccent}{1E6A}
\pdfglyphtounicode{Tdotbelow}{1E6C}
\pdfglyphtounicode{Tecyrillic}{0422}
\pdfglyphtounicode{Tedescendercyrillic}{04AC}
\pdfglyphtounicode{Tenroman}{2169}
\pdfglyphtounicode{Tetsecyrillic}{04B4}
\pdfglyphtounicode{Theta}{0398}
\pdfglyphtounicode{Thook}{01AC}
\pdfglyphtounicode{Thorn}{00DE}
\pdfglyphtounicode{Thornsmall}{00FE}
\pdfglyphtounicode{Threeroman}{2162}
\pdfglyphtounicode{Tildesmall}{02DC}
\pdfglyphtounicode{Tiwnarmenian}{054F}
\pdfglyphtounicode{Tlinebelow}{1E6E}
\pdfglyphtounicode{Tmonospace}{FF34}
\pdfglyphtounicode{Toarmenian}{0539}
\pdfglyphtounicode{Tonefive}{01BC}
\pdfglyphtounicode{Tonesix}{0184}
\pdfglyphtounicode{Tonetwo}{01A7}
\pdfglyphtounicode{Tretroflexhook}{01AE}
\pdfglyphtounicode{Tsecyrillic}{0426}
\pdfglyphtounicode{Tshecyrillic}{040B}
\pdfglyphtounicode{Tsmall}{0074}
\pdfglyphtounicode{Twelveroman}{216B}
\pdfglyphtounicode{Tworoman}{2161}
\pdfglyphtounicode{U}{0055}
\pdfglyphtounicode{Uacute}{00DA}
\pdfglyphtounicode{Uacutesmall}{00FA}
\pdfglyphtounicode{Ubreve}{016C}
\pdfglyphtounicode{Ucaron}{01D3}
\pdfglyphtounicode{Ucircle}{24CA}
\pdfglyphtounicode{Ucircumflex}{00DB}
\pdfglyphtounicode{Ucircumflexbelow}{1E76}
\pdfglyphtounicode{Ucircumflexsmall}{00FB}
\pdfglyphtounicode{Ucyrillic}{0423}
\pdfglyphtounicode{Udblacute}{0170}
\pdfglyphtounicode{Udblgrave}{0214}
\pdfglyphtounicode{Udieresis}{00DC}
\pdfglyphtounicode{Udieresisacute}{01D7}
\pdfglyphtounicode{Udieresisbelow}{1E72}
\pdfglyphtounicode{Udieresiscaron}{01D9}
\pdfglyphtounicode{Udieresiscyrillic}{04F0}
\pdfglyphtounicode{Udieresisgrave}{01DB}
\pdfglyphtounicode{Udieresismacron}{01D5}
\pdfglyphtounicode{Udieresissmall}{00FC}
\pdfglyphtounicode{Udotbelow}{1EE4}
\pdfglyphtounicode{Ugrave}{00D9}
\pdfglyphtounicode{Ugravesmall}{00F9}
\pdfglyphtounicode{Uhookabove}{1EE6}
\pdfglyphtounicode{Uhorn}{01AF}
\pdfglyphtounicode{Uhornacute}{1EE8}
\pdfglyphtounicode{Uhorndotbelow}{1EF0}
\pdfglyphtounicode{Uhorngrave}{1EEA}
\pdfglyphtounicode{Uhornhookabove}{1EEC}
\pdfglyphtounicode{Uhorntilde}{1EEE}
\pdfglyphtounicode{Uhungarumlaut}{0170}
\pdfglyphtounicode{Uhungarumlautcyrillic}{04F2}
\pdfglyphtounicode{Uinvertedbreve}{0216}
\pdfglyphtounicode{Ukcyrillic}{0478}
\pdfglyphtounicode{Umacron}{016A}
\pdfglyphtounicode{Umacroncyrillic}{04EE}
\pdfglyphtounicode{Umacrondieresis}{1E7A}
\pdfglyphtounicode{Umonospace}{FF35}
\pdfglyphtounicode{Uogonek}{0172}
\pdfglyphtounicode{Upsilon}{03A5}
\pdfglyphtounicode{Upsilon1}{03D2}
\pdfglyphtounicode{Upsilonacutehooksymbolgreek}{03D3}
\pdfglyphtounicode{Upsilonafrican}{01B1}
\pdfglyphtounicode{Upsilondieresis}{03AB}
\pdfglyphtounicode{Upsilondieresishooksymbolgreek}{03D4}
\pdfglyphtounicode{Upsilonhooksymbol}{03D2}
\pdfglyphtounicode{Upsilontonos}{038E}
\pdfglyphtounicode{Uring}{016E}
\pdfglyphtounicode{Ushortcyrillic}{040E}
\pdfglyphtounicode{Usmall}{0075}
\pdfglyphtounicode{Ustraightcyrillic}{04AE}
\pdfglyphtounicode{Ustraightstrokecyrillic}{04B0}
\pdfglyphtounicode{Utilde}{0168}
\pdfglyphtounicode{Utildeacute}{1E78}
\pdfglyphtounicode{Utildebelow}{1E74}
\pdfglyphtounicode{V}{0056}
\pdfglyphtounicode{Vcircle}{24CB}
\pdfglyphtounicode{Vdotbelow}{1E7E}
\pdfglyphtounicode{Vecyrillic}{0412}
\pdfglyphtounicode{Vewarmenian}{054E}
\pdfglyphtounicode{Vhook}{01B2}
\pdfglyphtounicode{Vmonospace}{FF36}
\pdfglyphtounicode{Voarmenian}{0548}
\pdfglyphtounicode{Vsmall}{0076}
\pdfglyphtounicode{Vtilde}{1E7C}
\pdfglyphtounicode{W}{0057}
\pdfglyphtounicode{Wacute}{1E82}
\pdfglyphtounicode{Wcircle}{24CC}
\pdfglyphtounicode{Wcircumflex}{0174}
\pdfglyphtounicode{Wdieresis}{1E84}
\pdfglyphtounicode{Wdotaccent}{1E86}
\pdfglyphtounicode{Wdotbelow}{1E88}
\pdfglyphtounicode{Wgrave}{1E80}
\pdfglyphtounicode{Wmonospace}{FF37}
\pdfglyphtounicode{Wsmall}{0077}
\pdfglyphtounicode{X}{0058}
\pdfglyphtounicode{Xcircle}{24CD}
\pdfglyphtounicode{Xdieresis}{1E8C}
\pdfglyphtounicode{Xdotaccent}{1E8A}
\pdfglyphtounicode{Xeharmenian}{053D}
\pdfglyphtounicode{Xi}{039E}
\pdfglyphtounicode{Xmonospace}{FF38}
\pdfglyphtounicode{Xsmall}{0078}
\pdfglyphtounicode{Y}{0059}
\pdfglyphtounicode{Yacute}{00DD}
\pdfglyphtounicode{Yacutesmall}{00FD}
\pdfglyphtounicode{Yatcyrillic}{0462}
\pdfglyphtounicode{Ycircle}{24CE}
\pdfglyphtounicode{Ycircumflex}{0176}
\pdfglyphtounicode{Ydieresis}{0178}
\pdfglyphtounicode{Ydieresissmall}{00FF}
\pdfglyphtounicode{Ydotaccent}{1E8E}
\pdfglyphtounicode{Ydotbelow}{1EF4}
\pdfglyphtounicode{Yen}{00A5}
\pdfglyphtounicode{Yericyrillic}{042B}
\pdfglyphtounicode{Yerudieresiscyrillic}{04F8}
\pdfglyphtounicode{Ygrave}{1EF2}
\pdfglyphtounicode{Yhook}{01B3}
\pdfglyphtounicode{Yhookabove}{1EF6}
\pdfglyphtounicode{Yiarmenian}{0545}
\pdfglyphtounicode{Yicyrillic}{0407}
\pdfglyphtounicode{Yiwnarmenian}{0552}
\pdfglyphtounicode{Ymonospace}{FF39}
\pdfglyphtounicode{Ysmall}{0079}
\pdfglyphtounicode{Ytilde}{1EF8}
\pdfglyphtounicode{Yusbigcyrillic}{046A}
\pdfglyphtounicode{Yusbigiotifiedcyrillic}{046C}
\pdfglyphtounicode{Yuslittlecyrillic}{0466}
\pdfglyphtounicode{Yuslittleiotifiedcyrillic}{0468}
\pdfglyphtounicode{Z}{005A}
\pdfglyphtounicode{Zaarmenian}{0536}
\pdfglyphtounicode{Zacute}{0179}
\pdfglyphtounicode{Zcaron}{017D}
\pdfglyphtounicode{Zcaronsmall}{017E}
\pdfglyphtounicode{Zcircle}{24CF}
\pdfglyphtounicode{Zcircumflex}{1E90}
\pdfglyphtounicode{Zdot}{017B}
\pdfglyphtounicode{Zdotaccent}{017B}
\pdfglyphtounicode{Zdotbelow}{1E92}
\pdfglyphtounicode{Zecyrillic}{0417}
\pdfglyphtounicode{Zedescendercyrillic}{0498}
\pdfglyphtounicode{Zedieresiscyrillic}{04DE}
\pdfglyphtounicode{Zeta}{0396}
\pdfglyphtounicode{Zhearmenian}{053A}
\pdfglyphtounicode{Zhebrevecyrillic}{04C1}
\pdfglyphtounicode{Zhecyrillic}{0416}
\pdfglyphtounicode{Zhedescendercyrillic}{0496}
\pdfglyphtounicode{Zhedieresiscyrillic}{04DC}
\pdfglyphtounicode{Zlinebelow}{1E94}
\pdfglyphtounicode{Zmonospace}{FF3A}
\pdfglyphtounicode{Zsmall}{007A}
\pdfglyphtounicode{Zstroke}{01B5}
\pdfglyphtounicode{a}{0061}
\pdfglyphtounicode{aabengali}{0986}
\pdfglyphtounicode{aacute}{00E1}
\pdfglyphtounicode{aadeva}{0906}
\pdfglyphtounicode{aagujarati}{0A86}
\pdfglyphtounicode{aagurmukhi}{0A06}
\pdfglyphtounicode{aamatragurmukhi}{0A3E}
\pdfglyphtounicode{aarusquare}{3303}
\pdfglyphtounicode{aavowelsignbengali}{09BE}
\pdfglyphtounicode{aavowelsigndeva}{093E}
\pdfglyphtounicode{aavowelsigngujarati}{0ABE}
\pdfglyphtounicode{abbreviationmarkarmenian}{055F}
\pdfglyphtounicode{abbreviationsigndeva}{0970}
\pdfglyphtounicode{abengali}{0985}
\pdfglyphtounicode{abopomofo}{311A}
\pdfglyphtounicode{abreve}{0103}
\pdfglyphtounicode{abreveacute}{1EAF}
\pdfglyphtounicode{abrevecyrillic}{04D1}
\pdfglyphtounicode{abrevedotbelow}{1EB7}
\pdfglyphtounicode{abrevegrave}{1EB1}
\pdfglyphtounicode{abrevehookabove}{1EB3}
\pdfglyphtounicode{abrevetilde}{1EB5}
\pdfglyphtounicode{acaron}{01CE}
\pdfglyphtounicode{acircle}{24D0}
\pdfglyphtounicode{acircumflex}{00E2}
\pdfglyphtounicode{acircumflexacute}{1EA5}
\pdfglyphtounicode{acircumflexdotbelow}{1EAD}
\pdfglyphtounicode{acircumflexgrave}{1EA7}
\pdfglyphtounicode{acircumflexhookabove}{1EA9}
\pdfglyphtounicode{acircumflextilde}{1EAB}
\pdfglyphtounicode{acute}{00B4}
\pdfglyphtounicode{acutebelowcmb}{0317}
\pdfglyphtounicode{acutecmb}{0301}
\pdfglyphtounicode{acutecomb}{0301}
\pdfglyphtounicode{acutedeva}{0954}
\pdfglyphtounicode{acutelowmod}{02CF}
\pdfglyphtounicode{acutetonecmb}{0341}
\pdfglyphtounicode{acyrillic}{0430}
\pdfglyphtounicode{adblgrave}{0201}
\pdfglyphtounicode{addakgurmukhi}{0A71}
\pdfglyphtounicode{adeva}{0905}
\pdfglyphtounicode{adieresis}{00E4}
\pdfglyphtounicode{adieresiscyrillic}{04D3}
\pdfglyphtounicode{adieresismacron}{01DF}
\pdfglyphtounicode{adotbelow}{1EA1}
\pdfglyphtounicode{adotmacron}{01E1}
\pdfglyphtounicode{ae}{00E6}
\pdfglyphtounicode{aeacute}{01FD}
\pdfglyphtounicode{aekorean}{3150}
\pdfglyphtounicode{aemacron}{01E3}
\pdfglyphtounicode{afii00208}{2015}
\pdfglyphtounicode{afii08941}{20A4}
\pdfglyphtounicode{afii10017}{0410}
\pdfglyphtounicode{afii10018}{0411}
\pdfglyphtounicode{afii10019}{0412}
\pdfglyphtounicode{afii10020}{0413}
\pdfglyphtounicode{afii10021}{0414}
\pdfglyphtounicode{afii10022}{0415}
\pdfglyphtounicode{afii10023}{0401}
\pdfglyphtounicode{afii10024}{0416}
\pdfglyphtounicode{afii10025}{0417}
\pdfglyphtounicode{afii10026}{0418}
\pdfglyphtounicode{afii10027}{0419}
\pdfglyphtounicode{afii10028}{041A}
\pdfglyphtounicode{afii10029}{041B}
\pdfglyphtounicode{afii10030}{041C}
\pdfglyphtounicode{afii10031}{041D}
\pdfglyphtounicode{afii10032}{041E}
\pdfglyphtounicode{afii10033}{041F}
\pdfglyphtounicode{afii10034}{0420}
\pdfglyphtounicode{afii10035}{0421}
\pdfglyphtounicode{afii10036}{0422}
\pdfglyphtounicode{afii10037}{0423}
\pdfglyphtounicode{afii10038}{0424}
\pdfglyphtounicode{afii10039}{0425}
\pdfglyphtounicode{afii10040}{0426}
\pdfglyphtounicode{afii10041}{0427}
\pdfglyphtounicode{afii10042}{0428}
\pdfglyphtounicode{afii10043}{0429}
\pdfglyphtounicode{afii10044}{042A}
\pdfglyphtounicode{afii10045}{042B}
\pdfglyphtounicode{afii10046}{042C}
\pdfglyphtounicode{afii10047}{042D}
\pdfglyphtounicode{afii10048}{042E}
\pdfglyphtounicode{afii10049}{042F}
\pdfglyphtounicode{afii10050}{0490}
\pdfglyphtounicode{afii10051}{0402}
\pdfglyphtounicode{afii10052}{0403}
\pdfglyphtounicode{afii10053}{0404}
\pdfglyphtounicode{afii10054}{0405}
\pdfglyphtounicode{afii10055}{0406}
\pdfglyphtounicode{afii10056}{0407}
\pdfglyphtounicode{afii10057}{0408}
\pdfglyphtounicode{afii10058}{0409}
\pdfglyphtounicode{afii10059}{040A}
\pdfglyphtounicode{afii10060}{040B}
\pdfglyphtounicode{afii10061}{040C}
\pdfglyphtounicode{afii10062}{040E}
\pdfglyphtounicode{afii10063}{F6C4}
\pdfglyphtounicode{afii10064}{F6C5}
\pdfglyphtounicode{afii10065}{0430}
\pdfglyphtounicode{afii10066}{0431}
\pdfglyphtounicode{afii10067}{0432}
\pdfglyphtounicode{afii10068}{0433}
\pdfglyphtounicode{afii10069}{0434}
\pdfglyphtounicode{afii10070}{0435}
\pdfglyphtounicode{afii10071}{0451}
\pdfglyphtounicode{afii10072}{0436}
\pdfglyphtounicode{afii10073}{0437}
\pdfglyphtounicode{afii10074}{0438}
\pdfglyphtounicode{afii10075}{0439}
\pdfglyphtounicode{afii10076}{043A}
\pdfglyphtounicode{afii10077}{043B}
\pdfglyphtounicode{afii10078}{043C}
\pdfglyphtounicode{afii10079}{043D}
\pdfglyphtounicode{afii10080}{043E}
\pdfglyphtounicode{afii10081}{043F}
\pdfglyphtounicode{afii10082}{0440}
\pdfglyphtounicode{afii10083}{0441}
\pdfglyphtounicode{afii10084}{0442}
\pdfglyphtounicode{afii10085}{0443}
\pdfglyphtounicode{afii10086}{0444}
\pdfglyphtounicode{afii10087}{0445}
\pdfglyphtounicode{afii10088}{0446}
\pdfglyphtounicode{afii10089}{0447}
\pdfglyphtounicode{afii10090}{0448}
\pdfglyphtounicode{afii10091}{0449}
\pdfglyphtounicode{afii10092}{044A}
\pdfglyphtounicode{afii10093}{044B}
\pdfglyphtounicode{afii10094}{044C}
\pdfglyphtounicode{afii10095}{044D}
\pdfglyphtounicode{afii10096}{044E}
\pdfglyphtounicode{afii10097}{044F}
\pdfglyphtounicode{afii10098}{0491}
\pdfglyphtounicode{afii10099}{0452}
\pdfglyphtounicode{afii10100}{0453}
\pdfglyphtounicode{afii10101}{0454}
\pdfglyphtounicode{afii10102}{0455}
\pdfglyphtounicode{afii10103}{0456}
\pdfglyphtounicode{afii10104}{0457}
\pdfglyphtounicode{afii10105}{0458}
\pdfglyphtounicode{afii10106}{0459}
\pdfglyphtounicode{afii10107}{045A}
\pdfglyphtounicode{afii10108}{045B}
\pdfglyphtounicode{afii10109}{045C}
\pdfglyphtounicode{afii10110}{045E}
\pdfglyphtounicode{afii10145}{040F}
\pdfglyphtounicode{afii10146}{0462}
\pdfglyphtounicode{afii10147}{0472}
\pdfglyphtounicode{afii10148}{0474}
\pdfglyphtounicode{afii10192}{F6C6}
\pdfglyphtounicode{afii10193}{045F}
\pdfglyphtounicode{afii10194}{0463}
\pdfglyphtounicode{afii10195}{0473}
\pdfglyphtounicode{afii10196}{0475}
\pdfglyphtounicode{afii10831}{F6C7}
\pdfglyphtounicode{afii10832}{F6C8}
\pdfglyphtounicode{afii10846}{04D9}
\pdfglyphtounicode{afii299}{200E}
\pdfglyphtounicode{afii300}{200F}
\pdfglyphtounicode{afii301}{200D}
\pdfglyphtounicode{afii57381}{066A}
\pdfglyphtounicode{afii57388}{060C}
\pdfglyphtounicode{afii57392}{0660}
\pdfglyphtounicode{afii57393}{0661}
\pdfglyphtounicode{afii57394}{0662}
\pdfglyphtounicode{afii57395}{0663}
\pdfglyphtounicode{afii57396}{0664}
\pdfglyphtounicode{afii57397}{0665}
\pdfglyphtounicode{afii57398}{0666}
\pdfglyphtounicode{afii57399}{0667}
\pdfglyphtounicode{afii57400}{0668}
\pdfglyphtounicode{afii57401}{0669}
\pdfglyphtounicode{afii57403}{061B}
\pdfglyphtounicode{afii57407}{061F}
\pdfglyphtounicode{afii57409}{0621}
\pdfglyphtounicode{afii57410}{0622}
\pdfglyphtounicode{afii57411}{0623}
\pdfglyphtounicode{afii57412}{0624}
\pdfglyphtounicode{afii57413}{0625}
\pdfglyphtounicode{afii57414}{0626}
\pdfglyphtounicode{afii57415}{0627}
\pdfglyphtounicode{afii57416}{0628}
\pdfglyphtounicode{afii57417}{0629}
\pdfglyphtounicode{afii57418}{062A}
\pdfglyphtounicode{afii57419}{062B}
\pdfglyphtounicode{afii57420}{062C}
\pdfglyphtounicode{afii57421}{062D}
\pdfglyphtounicode{afii57422}{062E}
\pdfglyphtounicode{afii57423}{062F}
\pdfglyphtounicode{afii57424}{0630}
\pdfglyphtounicode{afii57425}{0631}
\pdfglyphtounicode{afii57426}{0632}
\pdfglyphtounicode{afii57427}{0633}
\pdfglyphtounicode{afii57428}{0634}
\pdfglyphtounicode{afii57429}{0635}
\pdfglyphtounicode{afii57430}{0636}
\pdfglyphtounicode{afii57431}{0637}
\pdfglyphtounicode{afii57432}{0638}
\pdfglyphtounicode{afii57433}{0639}
\pdfglyphtounicode{afii57434}{063A}
\pdfglyphtounicode{afii57440}{0640}
\pdfglyphtounicode{afii57441}{0641}
\pdfglyphtounicode{afii57442}{0642}
\pdfglyphtounicode{afii57443}{0643}
\pdfglyphtounicode{afii57444}{0644}
\pdfglyphtounicode{afii57445}{0645}
\pdfglyphtounicode{afii57446}{0646}
\pdfglyphtounicode{afii57448}{0648}
\pdfglyphtounicode{afii57449}{0649}
\pdfglyphtounicode{afii57450}{064A}
\pdfglyphtounicode{afii57451}{064B}
\pdfglyphtounicode{afii57452}{064C}
\pdfglyphtounicode{afii57453}{064D}
\pdfglyphtounicode{afii57454}{064E}
\pdfglyphtounicode{afii57455}{064F}
\pdfglyphtounicode{afii57456}{0650}
\pdfglyphtounicode{afii57457}{0651}
\pdfglyphtounicode{afii57458}{0652}
\pdfglyphtounicode{afii57470}{0647}
\pdfglyphtounicode{afii57505}{06A4}
\pdfglyphtounicode{afii57506}{067E}
\pdfglyphtounicode{afii57507}{0686}
\pdfglyphtounicode{afii57508}{0698}
\pdfglyphtounicode{afii57509}{06AF}
\pdfglyphtounicode{afii57511}{0679}
\pdfglyphtounicode{afii57512}{0688}
\pdfglyphtounicode{afii57513}{0691}
\pdfglyphtounicode{afii57514}{06BA}
\pdfglyphtounicode{afii57519}{06D2}
\pdfglyphtounicode{afii57534}{06D5}
\pdfglyphtounicode{afii57636}{20AA}
\pdfglyphtounicode{afii57645}{05BE}
\pdfglyphtounicode{afii57658}{05C3}
\pdfglyphtounicode{afii57664}{05D0}
\pdfglyphtounicode{afii57665}{05D1}
\pdfglyphtounicode{afii57666}{05D2}
\pdfglyphtounicode{afii57667}{05D3}
\pdfglyphtounicode{afii57668}{05D4}
\pdfglyphtounicode{afii57669}{05D5}
\pdfglyphtounicode{afii57670}{05D6}
\pdfglyphtounicode{afii57671}{05D7}
\pdfglyphtounicode{afii57672}{05D8}
\pdfglyphtounicode{afii57673}{05D9}
\pdfglyphtounicode{afii57674}{05DA}
\pdfglyphtounicode{afii57675}{05DB}
\pdfglyphtounicode{afii57676}{05DC}
\pdfglyphtounicode{afii57677}{05DD}
\pdfglyphtounicode{afii57678}{05DE}
\pdfglyphtounicode{afii57679}{05DF}
\pdfglyphtounicode{afii57680}{05E0}
\pdfglyphtounicode{afii57681}{05E1}
\pdfglyphtounicode{afii57682}{05E2}
\pdfglyphtounicode{afii57683}{05E3}
\pdfglyphtounicode{afii57684}{05E4}
\pdfglyphtounicode{afii57685}{05E5}
\pdfglyphtounicode{afii57686}{05E6}
\pdfglyphtounicode{afii57687}{05E7}
\pdfglyphtounicode{afii57688}{05E8}
\pdfglyphtounicode{afii57689}{05E9}
\pdfglyphtounicode{afii57690}{05EA}
\pdfglyphtounicode{afii57694}{FB2A}
\pdfglyphtounicode{afii57695}{FB2B}
\pdfglyphtounicode{afii57700}{FB4B}
\pdfglyphtounicode{afii57705}{FB1F}
\pdfglyphtounicode{afii57716}{05F0}
\pdfglyphtounicode{afii57717}{05F1}
\pdfglyphtounicode{afii57718}{05F2}
\pdfglyphtounicode{afii57723}{FB35}
\pdfglyphtounicode{afii57793}{05B4}
\pdfglyphtounicode{afii57794}{05B5}
\pdfglyphtounicode{afii57795}{05B6}
\pdfglyphtounicode{afii57796}{05BB}
\pdfglyphtounicode{afii57797}{05B8}
\pdfglyphtounicode{afii57798}{05B7}
\pdfglyphtounicode{afii57799}{05B0}
\pdfglyphtounicode{afii57800}{05B2}
\pdfglyphtounicode{afii57801}{05B1}
\pdfglyphtounicode{afii57802}{05B3}
\pdfglyphtounicode{afii57803}{05C2}
\pdfglyphtounicode{afii57804}{05C1}
\pdfglyphtounicode{afii57806}{05B9}
\pdfglyphtounicode{afii57807}{05BC}
\pdfglyphtounicode{afii57839}{05BD}
\pdfglyphtounicode{afii57841}{05BF}
\pdfglyphtounicode{afii57842}{05C0}
\pdfglyphtounicode{afii57929}{02BC}
\pdfglyphtounicode{afii61248}{2105}
\pdfglyphtounicode{afii61289}{2113}
\pdfglyphtounicode{afii61352}{2116}
\pdfglyphtounicode{afii61573}{202C}
\pdfglyphtounicode{afii61574}{202D}
\pdfglyphtounicode{afii61575}{202E}
\pdfglyphtounicode{afii61664}{200C}
\pdfglyphtounicode{afii63167}{066D}
\pdfglyphtounicode{afii64937}{02BD}
\pdfglyphtounicode{agrave}{00E0}
\pdfglyphtounicode{agujarati}{0A85}
\pdfglyphtounicode{agurmukhi}{0A05}
\pdfglyphtounicode{ahiragana}{3042}
\pdfglyphtounicode{ahookabove}{1EA3}
\pdfglyphtounicode{aibengali}{0990}
\pdfglyphtounicode{aibopomofo}{311E}
\pdfglyphtounicode{aideva}{0910}
\pdfglyphtounicode{aiecyrillic}{04D5}
\pdfglyphtounicode{aigujarati}{0A90}
\pdfglyphtounicode{aigurmukhi}{0A10}
\pdfglyphtounicode{aimatragurmukhi}{0A48}
\pdfglyphtounicode{ainarabic}{0639}
\pdfglyphtounicode{ainfinalarabic}{FECA}
\pdfglyphtounicode{aininitialarabic}{FECB}
\pdfglyphtounicode{ainmedialarabic}{FECC}
\pdfglyphtounicode{ainvertedbreve}{0203}
\pdfglyphtounicode{aivowelsignbengali}{09C8}
\pdfglyphtounicode{aivowelsigndeva}{0948}
\pdfglyphtounicode{aivowelsigngujarati}{0AC8}
\pdfglyphtounicode{akatakana}{30A2}
\pdfglyphtounicode{akatakanahalfwidth}{FF71}
\pdfglyphtounicode{akorean}{314F}
\pdfglyphtounicode{alef}{05D0}
\pdfglyphtounicode{alefarabic}{0627}
\pdfglyphtounicode{alefdageshhebrew}{FB30}
\pdfglyphtounicode{aleffinalarabic}{FE8E}
\pdfglyphtounicode{alefhamzaabovearabic}{0623}
\pdfglyphtounicode{alefhamzaabovefinalarabic}{FE84}
\pdfglyphtounicode{alefhamzabelowarabic}{0625}
\pdfglyphtounicode{alefhamzabelowfinalarabic}{FE88}
\pdfglyphtounicode{alefhebrew}{05D0}
\pdfglyphtounicode{aleflamedhebrew}{FB4F}
\pdfglyphtounicode{alefmaddaabovearabic}{0622}
\pdfglyphtounicode{alefmaddaabovefinalarabic}{FE82}
\pdfglyphtounicode{alefmaksuraarabic}{0649}
\pdfglyphtounicode{alefmaksurafinalarabic}{FEF0}
\pdfglyphtounicode{alefmaksurainitialarabic}{FEF3}
\pdfglyphtounicode{alefmaksuramedialarabic}{FEF4}
\pdfglyphtounicode{alefpatahhebrew}{FB2E}
\pdfglyphtounicode{alefqamatshebrew}{FB2F}
\pdfglyphtounicode{aleph}{2135}
\pdfglyphtounicode{allequal}{224C}
\pdfglyphtounicode{alpha}{03B1}
\pdfglyphtounicode{alphatonos}{03AC}
\pdfglyphtounicode{amacron}{0101}
\pdfglyphtounicode{amonospace}{FF41}
\pdfglyphtounicode{ampersand}{0026}
\pdfglyphtounicode{ampersandmonospace}{FF06}
\pdfglyphtounicode{ampersandsmall}{0026}
\pdfglyphtounicode{amsquare}{33C2}
\pdfglyphtounicode{anbopomofo}{3122}
\pdfglyphtounicode{angbopomofo}{3124}
\pdfglyphtounicode{angbracketleft}{27E8}
\pdfglyphtounicode{angbracketright}{27E9}
\pdfglyphtounicode{angkhankhuthai}{0E5A}
\pdfglyphtounicode{angle}{2220}
\pdfglyphtounicode{anglebracketleft}{3008}
\pdfglyphtounicode{anglebracketleftvertical}{FE3F}
\pdfglyphtounicode{anglebracketright}{3009}
\pdfglyphtounicode{anglebracketrightvertical}{FE40}
\pdfglyphtounicode{angleleft}{2329}
\pdfglyphtounicode{angleright}{232A}
\pdfglyphtounicode{angstrom}{212B}
\pdfglyphtounicode{anoteleia}{0387}
\pdfglyphtounicode{anticlockwise}{27F2}
\pdfglyphtounicode{anudattadeva}{0952}
\pdfglyphtounicode{anusvarabengali}{0982}
\pdfglyphtounicode{anusvaradeva}{0902}
\pdfglyphtounicode{anusvaragujarati}{0A82}
\pdfglyphtounicode{aogonek}{0105}
\pdfglyphtounicode{apaatosquare}{3300}
\pdfglyphtounicode{aparen}{249C}
\pdfglyphtounicode{apostrophearmenian}{055A}
\pdfglyphtounicode{apostrophemod}{02BC}
\pdfglyphtounicode{apple}{F8FF}
\pdfglyphtounicode{approaches}{2250}
\pdfglyphtounicode{approxequal}{2248}
\pdfglyphtounicode{approxequalorimage}{2252}
\pdfglyphtounicode{approximatelyequal}{2245}
\pdfglyphtounicode{approxorequal}{224A}
\pdfglyphtounicode{araeaekorean}{318E}
\pdfglyphtounicode{araeakorean}{318D}
\pdfglyphtounicode{arc}{2312}
\pdfglyphtounicode{archleftdown}{21B6}
\pdfglyphtounicode{archrightdown}{21B7}
\pdfglyphtounicode{arighthalfring}{1E9A}
\pdfglyphtounicode{aring}{00E5}
\pdfglyphtounicode{aringacute}{01FB}
\pdfglyphtounicode{aringbelow}{1E01}
\pdfglyphtounicode{arrowboth}{2194}
\pdfglyphtounicode{arrowbothv}{2195}
\pdfglyphtounicode{arrowdashdown}{21E3}
\pdfglyphtounicode{arrowdashleft}{21E0}
\pdfglyphtounicode{arrowdashright}{21E2}
\pdfglyphtounicode{arrowdashup}{21E1}
\pdfglyphtounicode{arrowdblboth}{21D4}
\pdfglyphtounicode{arrowdblbothv}{21D5}
\pdfglyphtounicode{arrowdbldown}{21D3}
\pdfglyphtounicode{arrowdblleft}{21D0}
\pdfglyphtounicode{arrowdblright}{21D2}
\pdfglyphtounicode{arrowdblup}{21D1}
\pdfglyphtounicode{arrowdown}{2193}
\pdfglyphtounicode{arrowdownleft}{2199}
\pdfglyphtounicode{arrowdownright}{2198}
\pdfglyphtounicode{arrowdownwhite}{21E9}
\pdfglyphtounicode{arrowheaddownmod}{02C5}
\pdfglyphtounicode{arrowheadleftmod}{02C2}
\pdfglyphtounicode{arrowheadrightmod}{02C3}
\pdfglyphtounicode{arrowheadupmod}{02C4}
\pdfglyphtounicode{arrowhorizex}{F8E7}
\pdfglyphtounicode{arrowleft}{2190}
\pdfglyphtounicode{arrowleftbothalf}{21BD}
\pdfglyphtounicode{arrowleftdbl}{21D0}
\pdfglyphtounicode{arrowleftdblstroke}{21CD}
\pdfglyphtounicode{arrowleftoverright}{21C6}
\pdfglyphtounicode{arrowlefttophalf}{21BC}
\pdfglyphtounicode{arrowleftwhite}{21E6}
\pdfglyphtounicode{arrownortheast}{2197}
\pdfglyphtounicode{arrownorthwest}{2196}
\pdfglyphtounicode{arrowparrleftright}{21C6}
\pdfglyphtounicode{arrowparrrightleft}{21C4}
\pdfglyphtounicode{arrowright}{2192}
\pdfglyphtounicode{arrowrightbothalf}{21C1}
\pdfglyphtounicode{arrowrightdblstroke}{21CF}
\pdfglyphtounicode{arrowrightheavy}{279E}
\pdfglyphtounicode{arrowrightoverleft}{21C4}
\pdfglyphtounicode{arrowrighttophalf}{21C0}
\pdfglyphtounicode{arrowrightwhite}{21E8}
\pdfglyphtounicode{arrowsoutheast}{2198}
\pdfglyphtounicode{arrowsouthwest}{2199}
\pdfglyphtounicode{arrowtableft}{21E4}
\pdfglyphtounicode{arrowtabright}{21E5}
\pdfglyphtounicode{arrowtailleft}{21A2}
\pdfglyphtounicode{arrowtailright}{21A3}
\pdfglyphtounicode{arrowtripleleft}{21DA}
\pdfglyphtounicode{arrowtripleright}{21DB}
\pdfglyphtounicode{arrowup}{2191}
\pdfglyphtounicode{arrowupdn}{2195}
\pdfglyphtounicode{arrowupdnbse}{21A8}
\pdfglyphtounicode{arrowupdownbase}{21A8}
\pdfglyphtounicode{arrowupleft}{2196}
\pdfglyphtounicode{arrowupleftofdown}{21C5}
\pdfglyphtounicode{arrowupright}{2197}
\pdfglyphtounicode{arrowupwhite}{21E7}
\pdfglyphtounicode{arrowvertex}{F8E6}
\pdfglyphtounicode{asciicircum}{005E}
\pdfglyphtounicode{asciicircummonospace}{FF3E}
\pdfglyphtounicode{asciitilde}{007E}
\pdfglyphtounicode{asciitildemonospace}{FF5E}
\pdfglyphtounicode{ascript}{0251}
\pdfglyphtounicode{ascriptturned}{0252}
\pdfglyphtounicode{asmallhiragana}{3041}
\pdfglyphtounicode{asmallkatakana}{30A1}
\pdfglyphtounicode{asmallkatakanahalfwidth}{FF67}
\pdfglyphtounicode{asterisk}{002A}
\pdfglyphtounicode{asteriskaltonearabic}{066D}
\pdfglyphtounicode{asteriskarabic}{066D}
\pdfglyphtounicode{asteriskcentered}{2217}
\pdfglyphtounicode{asteriskmath}{2217}
\pdfglyphtounicode{asteriskmonospace}{FF0A}
\pdfglyphtounicode{asterisksmall}{FE61}
\pdfglyphtounicode{asterism}{2042}
\pdfglyphtounicode{asuperior}{0061}
\pdfglyphtounicode{asymptoticallyequal}{2243}
\pdfglyphtounicode{at}{0040}
\pdfglyphtounicode{atilde}{00E3}
\pdfglyphtounicode{atmonospace}{FF20}
\pdfglyphtounicode{atsmall}{FE6B}
\pdfglyphtounicode{aturned}{0250}
\pdfglyphtounicode{aubengali}{0994}
\pdfglyphtounicode{aubopomofo}{3120}
\pdfglyphtounicode{audeva}{0914}
\pdfglyphtounicode{augujarati}{0A94}
\pdfglyphtounicode{augurmukhi}{0A14}
\pdfglyphtounicode{aulengthmarkbengali}{09D7}
\pdfglyphtounicode{aumatragurmukhi}{0A4C}
\pdfglyphtounicode{auvowelsignbengali}{09CC}
\pdfglyphtounicode{auvowelsigndeva}{094C}
\pdfglyphtounicode{auvowelsigngujarati}{0ACC}
\pdfglyphtounicode{avagrahadeva}{093D}
\pdfglyphtounicode{aybarmenian}{0561}
\pdfglyphtounicode{ayin}{05E2}
\pdfglyphtounicode{ayinaltonehebrew}{FB20}
\pdfglyphtounicode{ayinhebrew}{05E2}
\pdfglyphtounicode{b}{0062}
\pdfglyphtounicode{babengali}{09AC}
\pdfglyphtounicode{backslash}{005C}
\pdfglyphtounicode{backslashmonospace}{FF3C}
\pdfglyphtounicode{badeva}{092C}
\pdfglyphtounicode{bagujarati}{0AAC}
\pdfglyphtounicode{bagurmukhi}{0A2C}
\pdfglyphtounicode{bahiragana}{3070}
\pdfglyphtounicode{bahtthai}{0E3F}
\pdfglyphtounicode{bakatakana}{30D0}
\pdfglyphtounicode{bar}{007C}
\pdfglyphtounicode{bardbl}{2225}
\pdfglyphtounicode{barmonospace}{FF5C}
\pdfglyphtounicode{bbopomofo}{3105}
\pdfglyphtounicode{bcircle}{24D1}
\pdfglyphtounicode{bdotaccent}{1E03}
\pdfglyphtounicode{bdotbelow}{1E05}
\pdfglyphtounicode{beamedsixteenthnotes}{266C}
\pdfglyphtounicode{because}{2235}
\pdfglyphtounicode{becyrillic}{0431}
\pdfglyphtounicode{beharabic}{0628}
\pdfglyphtounicode{behfinalarabic}{FE90}
\pdfglyphtounicode{behinitialarabic}{FE91}
\pdfglyphtounicode{behiragana}{3079}
\pdfglyphtounicode{behmedialarabic}{FE92}
\pdfglyphtounicode{behmeeminitialarabic}{FC9F}
\pdfglyphtounicode{behmeemisolatedarabic}{FC08}
\pdfglyphtounicode{behnoonfinalarabic}{FC6D}
\pdfglyphtounicode{bekatakana}{30D9}
\pdfglyphtounicode{benarmenian}{0562}
\pdfglyphtounicode{bet}{05D1}
\pdfglyphtounicode{beta}{03B2}
\pdfglyphtounicode{betasymbolgreek}{03D0}
\pdfglyphtounicode{betdagesh}{FB31}
\pdfglyphtounicode{betdageshhebrew}{FB31}
\pdfglyphtounicode{beth}{2136}
\pdfglyphtounicode{bethebrew}{05D1}
\pdfglyphtounicode{betrafehebrew}{FB4C}
\pdfglyphtounicode{between}{226C}
\pdfglyphtounicode{bhabengali}{09AD}
\pdfglyphtounicode{bhadeva}{092D}
\pdfglyphtounicode{bhagujarati}{0AAD}
\pdfglyphtounicode{bhagurmukhi}{0A2D}
\pdfglyphtounicode{bhook}{0253}
\pdfglyphtounicode{bihiragana}{3073}
\pdfglyphtounicode{bikatakana}{30D3}
\pdfglyphtounicode{bilabialclick}{0298}
\pdfglyphtounicode{bindigurmukhi}{0A02}
\pdfglyphtounicode{birusquare}{3331}
\pdfglyphtounicode{blackcircle}{25CF}
\pdfglyphtounicode{blackdiamond}{25C6}
\pdfglyphtounicode{blackdownpointingtriangle}{25BC}
\pdfglyphtounicode{blackleftpointingpointer}{25C4}
\pdfglyphtounicode{blackleftpointingtriangle}{25C0}
\pdfglyphtounicode{blacklenticularbracketleft}{3010}
\pdfglyphtounicode{blacklenticularbracketleftvertical}{FE3B}
\pdfglyphtounicode{blacklenticularbracketright}{3011}
\pdfglyphtounicode{blacklenticularbracketrightvertical}{FE3C}
\pdfglyphtounicode{blacklowerlefttriangle}{25E3}
\pdfglyphtounicode{blacklowerrighttriangle}{25E2}
\pdfglyphtounicode{blackrectangle}{25AC}
\pdfglyphtounicode{blackrightpointingpointer}{25BA}
\pdfglyphtounicode{blackrightpointingtriangle}{25B6}
\pdfglyphtounicode{blacksmallsquare}{25AA}
\pdfglyphtounicode{blacksmilingface}{263B}
\pdfglyphtounicode{blacksquare}{25A0}
\pdfglyphtounicode{blackstar}{2605}
\pdfglyphtounicode{blackupperlefttriangle}{25E4}
\pdfglyphtounicode{blackupperrighttriangle}{25E5}
\pdfglyphtounicode{blackuppointingsmalltriangle}{25B4}
\pdfglyphtounicode{blackuppointingtriangle}{25B2}
\pdfglyphtounicode{blank}{2423}
\pdfglyphtounicode{blinebelow}{1E07}
\pdfglyphtounicode{block}{2588}
\pdfglyphtounicode{bmonospace}{FF42}
\pdfglyphtounicode{bobaimaithai}{0E1A}
\pdfglyphtounicode{bohiragana}{307C}
\pdfglyphtounicode{bokatakana}{30DC}
\pdfglyphtounicode{bparen}{249D}
\pdfglyphtounicode{bqsquare}{33C3}
\pdfglyphtounicode{braceex}{F8F4}
\pdfglyphtounicode{braceleft}{007B}
\pdfglyphtounicode{braceleftbt}{F8F3}
\pdfglyphtounicode{braceleftmid}{F8F2}
\pdfglyphtounicode{braceleftmonospace}{FF5B}
\pdfglyphtounicode{braceleftsmall}{FE5B}
\pdfglyphtounicode{bracelefttp}{F8F1}
\pdfglyphtounicode{braceleftvertical}{FE37}
\pdfglyphtounicode{braceright}{007D}
\pdfglyphtounicode{bracerightbt}{F8FE}
\pdfglyphtounicode{bracerightmid}{F8FD}
\pdfglyphtounicode{bracerightmonospace}{FF5D}
\pdfglyphtounicode{bracerightsmall}{FE5C}
\pdfglyphtounicode{bracerighttp}{F8FC}
\pdfglyphtounicode{bracerightvertical}{FE38}
\pdfglyphtounicode{bracketleft}{005B}
\pdfglyphtounicode{bracketleftbt}{F8F0}
\pdfglyphtounicode{bracketleftex}{F8EF}
\pdfglyphtounicode{bracketleftmonospace}{FF3B}
\pdfglyphtounicode{bracketlefttp}{F8EE}
\pdfglyphtounicode{bracketright}{005D}
\pdfglyphtounicode{bracketrightbt}{F8FB}
\pdfglyphtounicode{bracketrightex}{F8FA}
\pdfglyphtounicode{bracketrightmonospace}{FF3D}
\pdfglyphtounicode{bracketrighttp}{F8F9}
\pdfglyphtounicode{breve}{02D8}
\pdfglyphtounicode{brevebelowcmb}{032E}
\pdfglyphtounicode{brevecmb}{0306}
\pdfglyphtounicode{breveinvertedbelowcmb}{032F}
\pdfglyphtounicode{breveinvertedcmb}{0311}
\pdfglyphtounicode{breveinverteddoublecmb}{0361}
\pdfglyphtounicode{bridgebelowcmb}{032A}
\pdfglyphtounicode{bridgeinvertedbelowcmb}{033A}
\pdfglyphtounicode{brokenbar}{00A6}
\pdfglyphtounicode{bstroke}{0180}
\pdfglyphtounicode{bsuperior}{0062}
\pdfglyphtounicode{btopbar}{0183}
\pdfglyphtounicode{buhiragana}{3076}
\pdfglyphtounicode{bukatakana}{30D6}
\pdfglyphtounicode{bullet}{2022}
\pdfglyphtounicode{bulletinverse}{25D8}
\pdfglyphtounicode{bulletoperator}{2219}
\pdfglyphtounicode{bullseye}{25CE}
\pdfglyphtounicode{c}{0063}
\pdfglyphtounicode{caarmenian}{056E}
\pdfglyphtounicode{cabengali}{099A}
\pdfglyphtounicode{cacute}{0107}
\pdfglyphtounicode{cadeva}{091A}
\pdfglyphtounicode{cagujarati}{0A9A}
\pdfglyphtounicode{cagurmukhi}{0A1A}
\pdfglyphtounicode{calsquare}{3388}
\pdfglyphtounicode{candrabindubengali}{0981}
\pdfglyphtounicode{candrabinducmb}{0310}
\pdfglyphtounicode{candrabindudeva}{0901}
\pdfglyphtounicode{candrabindugujarati}{0A81}
\pdfglyphtounicode{capslock}{21EA}
\pdfglyphtounicode{careof}{2105}
\pdfglyphtounicode{caron}{02C7}
\pdfglyphtounicode{caronbelowcmb}{032C}
\pdfglyphtounicode{caroncmb}{030C}
\pdfglyphtounicode{carriagereturn}{21B5}
\pdfglyphtounicode{cbopomofo}{3118}
\pdfglyphtounicode{ccaron}{010D}
\pdfglyphtounicode{ccedilla}{00E7}
\pdfglyphtounicode{ccedillaacute}{1E09}
\pdfglyphtounicode{ccircle}{24D2}
\pdfglyphtounicode{ccircumflex}{0109}
\pdfglyphtounicode{ccurl}{0255}
\pdfglyphtounicode{cdot}{010B}
\pdfglyphtounicode{cdotaccent}{010B}
\pdfglyphtounicode{cdsquare}{33C5}
\pdfglyphtounicode{cedilla}{00B8}
\pdfglyphtounicode{cedillacmb}{0327}
\pdfglyphtounicode{ceilingleft}{2308}
\pdfglyphtounicode{ceilingright}{2309}
\pdfglyphtounicode{cent}{00A2}
\pdfglyphtounicode{centigrade}{2103}
\pdfglyphtounicode{centinferior}{00A2}
\pdfglyphtounicode{centmonospace}{FFE0}
\pdfglyphtounicode{centoldstyle}{00A2}
\pdfglyphtounicode{centsuperior}{00A2}
\pdfglyphtounicode{chaarmenian}{0579}
\pdfglyphtounicode{chabengali}{099B}
\pdfglyphtounicode{chadeva}{091B}
\pdfglyphtounicode{chagujarati}{0A9B}
\pdfglyphtounicode{chagurmukhi}{0A1B}
\pdfglyphtounicode{chbopomofo}{3114}
\pdfglyphtounicode{cheabkhasiancyrillic}{04BD}
\pdfglyphtounicode{check}{2713}
\pdfglyphtounicode{checkmark}{2713}
\pdfglyphtounicode{checyrillic}{0447}
\pdfglyphtounicode{chedescenderabkhasiancyrillic}{04BF}
\pdfglyphtounicode{chedescendercyrillic}{04B7}
\pdfglyphtounicode{chedieresiscyrillic}{04F5}
\pdfglyphtounicode{cheharmenian}{0573}
\pdfglyphtounicode{chekhakassiancyrillic}{04CC}
\pdfglyphtounicode{cheverticalstrokecyrillic}{04B9}
\pdfglyphtounicode{chi}{03C7}
\pdfglyphtounicode{chieuchacirclekorean}{3277}
\pdfglyphtounicode{chieuchaparenkorean}{3217}
\pdfglyphtounicode{chieuchcirclekorean}{3269}
\pdfglyphtounicode{chieuchkorean}{314A}
\pdfglyphtounicode{chieuchparenkorean}{3209}
\pdfglyphtounicode{chochangthai}{0E0A}
\pdfglyphtounicode{chochanthai}{0E08}
\pdfglyphtounicode{chochingthai}{0E09}
\pdfglyphtounicode{chochoethai}{0E0C}
\pdfglyphtounicode{chook}{0188}
\pdfglyphtounicode{cieucacirclekorean}{3276}
\pdfglyphtounicode{cieucaparenkorean}{3216}
\pdfglyphtounicode{cieuccirclekorean}{3268}
\pdfglyphtounicode{cieuckorean}{3148}
\pdfglyphtounicode{cieucparenkorean}{3208}
\pdfglyphtounicode{cieucuparenkorean}{321C}
\pdfglyphtounicode{circle}{25CB}
\pdfglyphtounicode{circleR}{00AE}
\pdfglyphtounicode{circleS}{24C8}
\pdfglyphtounicode{circleasterisk}{229B}
\pdfglyphtounicode{circlecopyrt}{20DD}
\pdfglyphtounicode{circledivide}{2298}
\pdfglyphtounicode{circledot}{2299}
\pdfglyphtounicode{circleequal}{229C}
\pdfglyphtounicode{circleminus}{2296}
\pdfglyphtounicode{circlemultiply}{2297}
\pdfglyphtounicode{circleot}{2299}
\pdfglyphtounicode{circleplus}{2295}
\pdfglyphtounicode{circlepostalmark}{3036}
\pdfglyphtounicode{circlering}{229A}
\pdfglyphtounicode{circlewithlefthalfblack}{25D0}
\pdfglyphtounicode{circlewithrighthalfblack}{25D1}
\pdfglyphtounicode{circumflex}{02C6}
\pdfglyphtounicode{circumflexbelowcmb}{032D}
\pdfglyphtounicode{circumflexcmb}{0302}
\pdfglyphtounicode{clear}{2327}
\pdfglyphtounicode{clickalveolar}{01C2}
\pdfglyphtounicode{clickdental}{01C0}
\pdfglyphtounicode{clicklateral}{01C1}
\pdfglyphtounicode{clickretroflex}{01C3}
\pdfglyphtounicode{clockwise}{27F3}
\pdfglyphtounicode{club}{2663}
\pdfglyphtounicode{clubsuitblack}{2663}
\pdfglyphtounicode{clubsuitwhite}{2667}
\pdfglyphtounicode{cmcubedsquare}{33A4}
\pdfglyphtounicode{cmonospace}{FF43}
\pdfglyphtounicode{cmsquaredsquare}{33A0}
\pdfglyphtounicode{coarmenian}{0581}
\pdfglyphtounicode{colon}{003A}
\pdfglyphtounicode{colonmonetary}{20A1}
\pdfglyphtounicode{colonmonospace}{FF1A}
\pdfglyphtounicode{colonsign}{20A1}
\pdfglyphtounicode{colonsmall}{FE55}
\pdfglyphtounicode{colontriangularhalfmod}{02D1}
\pdfglyphtounicode{colontriangularmod}{02D0}
\pdfglyphtounicode{comma}{002C}
\pdfglyphtounicode{commaabovecmb}{0313}
\pdfglyphtounicode{commaaboverightcmb}{0315}
\pdfglyphtounicode{commaaccent}{F6C3}
\pdfglyphtounicode{commaarabic}{060C}
\pdfglyphtounicode{commaarmenian}{055D}
\pdfglyphtounicode{commainferior}{002C}
\pdfglyphtounicode{commamonospace}{FF0C}
\pdfglyphtounicode{commareversedabovecmb}{0314}
\pdfglyphtounicode{commareversedmod}{02BD}
\pdfglyphtounicode{commasmall}{FE50}
\pdfglyphtounicode{commasuperior}{002C}
\pdfglyphtounicode{commaturnedabovecmb}{0312}
\pdfglyphtounicode{commaturnedmod}{02BB}
\pdfglyphtounicode{compass}{263C}
\pdfglyphtounicode{complement}{2201}
\pdfglyphtounicode{compwordmark}{200C}
\pdfglyphtounicode{congruent}{2245}
\pdfglyphtounicode{contourintegral}{222E}
\pdfglyphtounicode{control}{2303}
\pdfglyphtounicode{controlACK}{0006}
\pdfglyphtounicode{controlBEL}{0007}
\pdfglyphtounicode{controlBS}{0008}
\pdfglyphtounicode{controlCAN}{0018}
\pdfglyphtounicode{controlCR}{000D}
\pdfglyphtounicode{controlDC1}{0011}
\pdfglyphtounicode{controlDC2}{0012}
\pdfglyphtounicode{controlDC3}{0013}
\pdfglyphtounicode{controlDC4}{0014}
\pdfglyphtounicode{controlDEL}{007F}
\pdfglyphtounicode{controlDLE}{0010}
\pdfglyphtounicode{controlEM}{0019}
\pdfglyphtounicode{controlENQ}{0005}
\pdfglyphtounicode{controlEOT}{0004}
\pdfglyphtounicode{controlESC}{001B}
\pdfglyphtounicode{controlETB}{0017}
\pdfglyphtounicode{controlETX}{0003}
\pdfglyphtounicode{controlFF}{000C}
\pdfglyphtounicode{controlFS}{001C}
\pdfglyphtounicode{controlGS}{001D}
\pdfglyphtounicode{controlHT}{0009}
\pdfglyphtounicode{controlLF}{000A}
\pdfglyphtounicode{controlNAK}{0015}
\pdfglyphtounicode{controlRS}{001E}
\pdfglyphtounicode{controlSI}{000F}
\pdfglyphtounicode{controlSO}{000E}
\pdfglyphtounicode{controlSOT}{0002}
\pdfglyphtounicode{controlSTX}{0001}
\pdfglyphtounicode{controlSUB}{001A}
\pdfglyphtounicode{controlSYN}{0016}
\pdfglyphtounicode{controlUS}{001F}
\pdfglyphtounicode{controlVT}{000B}
\pdfglyphtounicode{coproduct}{2A3F}
\pdfglyphtounicode{copyright}{00A9}
\pdfglyphtounicode{copyrightsans}{00A9}
\pdfglyphtounicode{copyrightserif}{00A9}
\pdfglyphtounicode{cornerbracketleft}{300C}
\pdfglyphtounicode{cornerbracketlefthalfwidth}{FF62}
\pdfglyphtounicode{cornerbracketleftvertical}{FE41}
\pdfglyphtounicode{cornerbracketright}{300D}
\pdfglyphtounicode{cornerbracketrighthalfwidth}{FF63}
\pdfglyphtounicode{cornerbracketrightvertical}{FE42}
\pdfglyphtounicode{corporationsquare}{337F}
\pdfglyphtounicode{cosquare}{33C7}
\pdfglyphtounicode{coverkgsquare}{33C6}
\pdfglyphtounicode{cparen}{249E}
\pdfglyphtounicode{cruzeiro}{20A2}
\pdfglyphtounicode{cstretched}{0297}
\pdfglyphtounicode{ct}{0063 0074}
\pdfglyphtounicode{curlyand}{22CF}
\pdfglyphtounicode{curlyleft}{21AB}
\pdfglyphtounicode{curlyor}{22CE}
\pdfglyphtounicode{curlyright}{21AC}
\pdfglyphtounicode{currency}{00A4}
\pdfglyphtounicode{cwm}{200C}
\pdfglyphtounicode{cyrBreve}{02D8}
\pdfglyphtounicode{cyrFlex}{00A0 0311}
\pdfglyphtounicode{cyrbreve}{02D8}
\pdfglyphtounicode{cyrflex}{00A0 0311}
\pdfglyphtounicode{d}{0064}
\pdfglyphtounicode{daarmenian}{0564}
\pdfglyphtounicode{dabengali}{09A6}
\pdfglyphtounicode{dadarabic}{0636}
\pdfglyphtounicode{dadeva}{0926}
\pdfglyphtounicode{dadfinalarabic}{FEBE}
\pdfglyphtounicode{dadinitialarabic}{FEBF}
\pdfglyphtounicode{dadmedialarabic}{FEC0}
\pdfglyphtounicode{dagesh}{05BC}
\pdfglyphtounicode{dageshhebrew}{05BC}
\pdfglyphtounicode{dagger}{2020}
\pdfglyphtounicode{daggerdbl}{2021}
\pdfglyphtounicode{dagujarati}{0AA6}
\pdfglyphtounicode{dagurmukhi}{0A26}
\pdfglyphtounicode{dahiragana}{3060}
\pdfglyphtounicode{dakatakana}{30C0}
\pdfglyphtounicode{dalarabic}{062F}
\pdfglyphtounicode{dalet}{05D3}
\pdfglyphtounicode{daletdagesh}{FB33}
\pdfglyphtounicode{daletdageshhebrew}{FB33}
\pdfglyphtounicode{daleth}{2138}
\pdfglyphtounicode{dalethatafpatah}{05D3 05B2}
\pdfglyphtounicode{dalethatafpatahhebrew}{05D3 05B2}
\pdfglyphtounicode{dalethatafsegol}{05D3 05B1}
\pdfglyphtounicode{dalethatafsegolhebrew}{05D3 05B1}
\pdfglyphtounicode{dalethebrew}{05D3}
\pdfglyphtounicode{dalethiriq}{05D3 05B4}
\pdfglyphtounicode{dalethiriqhebrew}{05D3 05B4}
\pdfglyphtounicode{daletholam}{05D3 05B9}
\pdfglyphtounicode{daletholamhebrew}{05D3 05B9}
\pdfglyphtounicode{daletpatah}{05D3 05B7}
\pdfglyphtounicode{daletpatahhebrew}{05D3 05B7}
\pdfglyphtounicode{daletqamats}{05D3 05B8}
\pdfglyphtounicode{daletqamatshebrew}{05D3 05B8}
\pdfglyphtounicode{daletqubuts}{05D3 05BB}
\pdfglyphtounicode{daletqubutshebrew}{05D3 05BB}
\pdfglyphtounicode{daletsegol}{05D3 05B6}
\pdfglyphtounicode{daletsegolhebrew}{05D3 05B6}
\pdfglyphtounicode{daletsheva}{05D3 05B0}
\pdfglyphtounicode{daletshevahebrew}{05D3 05B0}
\pdfglyphtounicode{dalettsere}{05D3 05B5}
\pdfglyphtounicode{dalettserehebrew}{05D3 05B5}
\pdfglyphtounicode{dalfinalarabic}{FEAA}
\pdfglyphtounicode{dammaarabic}{064F}
\pdfglyphtounicode{dammalowarabic}{064F}
\pdfglyphtounicode{dammatanaltonearabic}{064C}
\pdfglyphtounicode{dammatanarabic}{064C}
\pdfglyphtounicode{danda}{0964}
\pdfglyphtounicode{dargahebrew}{05A7}
\pdfglyphtounicode{dargalefthebrew}{05A7}
\pdfglyphtounicode{dasiapneumatacyrilliccmb}{0485}
\pdfglyphtounicode{dbar}{0111}
\pdfglyphtounicode{dblGrave}{00A0 030F}
\pdfglyphtounicode{dblanglebracketleft}{300A}
\pdfglyphtounicode{dblanglebracketleftvertical}{FE3D}
\pdfglyphtounicode{dblanglebracketright}{300B}
\pdfglyphtounicode{dblanglebracketrightvertical}{FE3E}
\pdfglyphtounicode{dblarchinvertedbelowcmb}{032B}
\pdfglyphtounicode{dblarrowdwn}{21CA}
\pdfglyphtounicode{dblarrowheadleft}{219E}
\pdfglyphtounicode{dblarrowheadright}{21A0}
\pdfglyphtounicode{dblarrowleft}{21D4}
\pdfglyphtounicode{dblarrowright}{21D2}
\pdfglyphtounicode{dblarrowup}{21C8}
\pdfglyphtounicode{dblbracketleft}{27E6}
\pdfglyphtounicode{dblbracketright}{27E7}
\pdfglyphtounicode{dbldanda}{0965}
\pdfglyphtounicode{dblgrave}{00A0 030F}
\pdfglyphtounicode{dblgravecmb}{030F}
\pdfglyphtounicode{dblintegral}{222C}
\pdfglyphtounicode{dbllowline}{2017}
\pdfglyphtounicode{dbllowlinecmb}{0333}
\pdfglyphtounicode{dbloverlinecmb}{033F}
\pdfglyphtounicode{dblprimemod}{02BA}
\pdfglyphtounicode{dblverticalbar}{2016}
\pdfglyphtounicode{dblverticallineabovecmb}{030E}
\pdfglyphtounicode{dbopomofo}{3109}
\pdfglyphtounicode{dbsquare}{33C8}
\pdfglyphtounicode{dcaron}{010F}
\pdfglyphtounicode{dcedilla}{1E11}
\pdfglyphtounicode{dcircle}{24D3}
\pdfglyphtounicode{dcircumflexbelow}{1E13}
\pdfglyphtounicode{dcroat}{0111}
\pdfglyphtounicode{ddabengali}{09A1}
\pdfglyphtounicode{ddadeva}{0921}
\pdfglyphtounicode{ddagujarati}{0AA1}
\pdfglyphtounicode{ddagurmukhi}{0A21}
\pdfglyphtounicode{ddalarabic}{0688}
\pdfglyphtounicode{ddalfinalarabic}{FB89}
\pdfglyphtounicode{dddhadeva}{095C}
\pdfglyphtounicode{ddhabengali}{09A2}
\pdfglyphtounicode{ddhadeva}{0922}
\pdfglyphtounicode{ddhagujarati}{0AA2}
\pdfglyphtounicode{ddhagurmukhi}{0A22}
\pdfglyphtounicode{ddotaccent}{1E0B}
\pdfglyphtounicode{ddotbelow}{1E0D}
\pdfglyphtounicode{decimalseparatorarabic}{066B}
\pdfglyphtounicode{decimalseparatorpersian}{066B}
\pdfglyphtounicode{decyrillic}{0434}
\pdfglyphtounicode{defines}{225C}
\pdfglyphtounicode{degree}{00B0}
\pdfglyphtounicode{dehihebrew}{05AD}
\pdfglyphtounicode{dehiragana}{3067}
\pdfglyphtounicode{deicoptic}{03EF}
\pdfglyphtounicode{dekatakana}{30C7}
\pdfglyphtounicode{deleteleft}{232B}
\pdfglyphtounicode{deleteright}{2326}
\pdfglyphtounicode{delta}{03B4}
\pdfglyphtounicode{deltaturned}{018D}
\pdfglyphtounicode{denominatorminusonenumeratorbengali}{09F8}
\pdfglyphtounicode{dezh}{02A4}
\pdfglyphtounicode{dhabengali}{09A7}
\pdfglyphtounicode{dhadeva}{0927}
\pdfglyphtounicode{dhagujarati}{0AA7}
\pdfglyphtounicode{dhagurmukhi}{0A27}
\pdfglyphtounicode{dhook}{0257}
\pdfglyphtounicode{dialytikatonos}{0385}
\pdfglyphtounicode{dialytikatonoscmb}{0344}
\pdfglyphtounicode{diamond}{2666}
\pdfglyphtounicode{diamondmath}{22C4}
\pdfglyphtounicode{diamondsolid}{2666}
\pdfglyphtounicode{diamondsuitwhite}{2662}
\pdfglyphtounicode{dieresis}{00A8}
\pdfglyphtounicode{dieresisacute}{00A0 0308 0301}
\pdfglyphtounicode{dieresisbelowcmb}{0324}
\pdfglyphtounicode{dieresiscmb}{0308}
\pdfglyphtounicode{dieresisgrave}{00A0 0308 0300}
\pdfglyphtounicode{dieresistonos}{0385}
\pdfglyphtounicode{difference}{224F}
\pdfglyphtounicode{dihiragana}{3062}
\pdfglyphtounicode{dikatakana}{30C2}
\pdfglyphtounicode{dittomark}{3003}
\pdfglyphtounicode{divide}{00F7}
\pdfglyphtounicode{dividemultiply}{22C7}
\pdfglyphtounicode{divides}{2223}
\pdfglyphtounicode{divisionslash}{2215}
\pdfglyphtounicode{djecyrillic}{0452}
\pdfglyphtounicode{dkshade}{2593}
\pdfglyphtounicode{dlinebelow}{1E0F}
\pdfglyphtounicode{dlsquare}{3397}
\pdfglyphtounicode{dmacron}{0111}
\pdfglyphtounicode{dmonospace}{FF44}
\pdfglyphtounicode{dnblock}{2584}
\pdfglyphtounicode{dochadathai}{0E0E}
\pdfglyphtounicode{dodekthai}{0E14}
\pdfglyphtounicode{dohiragana}{3069}
\pdfglyphtounicode{dokatakana}{30C9}
\pdfglyphtounicode{dollar}{0024}
\pdfglyphtounicode{dollarinferior}{0024}
\pdfglyphtounicode{dollarmonospace}{FF04}
\pdfglyphtounicode{dollaroldstyle}{0024}
\pdfglyphtounicode{dollarsmall}{FE69}
\pdfglyphtounicode{dollarsuperior}{0024}
\pdfglyphtounicode{dong}{20AB}
\pdfglyphtounicode{dorusquare}{3326}
\pdfglyphtounicode{dotaccent}{02D9}
\pdfglyphtounicode{dotaccentcmb}{0307}
\pdfglyphtounicode{dotbelowcmb}{0323}
\pdfglyphtounicode{dotbelowcomb}{0323}
\pdfglyphtounicode{dotkatakana}{30FB}
\pdfglyphtounicode{dotlessi}{0131}
\pdfglyphtounicode{dotlessj}{0237}
\pdfglyphtounicode{dotlessjstrokehook}{0284}
\pdfglyphtounicode{dotmath}{22C5}
\pdfglyphtounicode{dotplus}{2214}
\pdfglyphtounicode{dottedcircle}{25CC}
\pdfglyphtounicode{doubleyodpatah}{FB1F}
\pdfglyphtounicode{doubleyodpatahhebrew}{FB1F}
\pdfglyphtounicode{downfall}{22CE}
\pdfglyphtounicode{downslope}{29F9}
\pdfglyphtounicode{downtackbelowcmb}{031E}
\pdfglyphtounicode{downtackmod}{02D5}
\pdfglyphtounicode{dparen}{249F}
\pdfglyphtounicode{dsuperior}{0064}
\pdfglyphtounicode{dtail}{0256}
\pdfglyphtounicode{dtopbar}{018C}
\pdfglyphtounicode{duhiragana}{3065}
\pdfglyphtounicode{dukatakana}{30C5}
\pdfglyphtounicode{dz}{01F3}
\pdfglyphtounicode{dzaltone}{02A3}
\pdfglyphtounicode{dzcaron}{01C6}
\pdfglyphtounicode{dzcurl}{02A5}
\pdfglyphtounicode{dzeabkhasiancyrillic}{04E1}
\pdfglyphtounicode{dzecyrillic}{0455}
\pdfglyphtounicode{dzhecyrillic}{045F}
\pdfglyphtounicode{e}{0065}
\pdfglyphtounicode{eacute}{00E9}
\pdfglyphtounicode{earth}{2641}
\pdfglyphtounicode{ebengali}{098F}
\pdfglyphtounicode{ebopomofo}{311C}
\pdfglyphtounicode{ebreve}{0115}
\pdfglyphtounicode{ecandradeva}{090D}
\pdfglyphtounicode{ecandragujarati}{0A8D}
\pdfglyphtounicode{ecandravowelsigndeva}{0945}
\pdfglyphtounicode{ecandravowelsigngujarati}{0AC5}
\pdfglyphtounicode{ecaron}{011B}
\pdfglyphtounicode{ecedillabreve}{1E1D}
\pdfglyphtounicode{echarmenian}{0565}
\pdfglyphtounicode{echyiwnarmenian}{0587}
\pdfglyphtounicode{ecircle}{24D4}
\pdfglyphtounicode{ecircumflex}{00EA}
\pdfglyphtounicode{ecircumflexacute}{1EBF}
\pdfglyphtounicode{ecircumflexbelow}{1E19}
\pdfglyphtounicode{ecircumflexdotbelow}{1EC7}
\pdfglyphtounicode{ecircumflexgrave}{1EC1}
\pdfglyphtounicode{ecircumflexhookabove}{1EC3}
\pdfglyphtounicode{ecircumflextilde}{1EC5}
\pdfglyphtounicode{ecyrillic}{0454}
\pdfglyphtounicode{edblgrave}{0205}
\pdfglyphtounicode{edeva}{090F}
\pdfglyphtounicode{edieresis}{00EB}
\pdfglyphtounicode{edot}{0117}
\pdfglyphtounicode{edotaccent}{0117}
\pdfglyphtounicode{edotbelow}{1EB9}
\pdfglyphtounicode{eegurmukhi}{0A0F}
\pdfglyphtounicode{eematragurmukhi}{0A47}
\pdfglyphtounicode{efcyrillic}{0444}
\pdfglyphtounicode{egrave}{00E8}
\pdfglyphtounicode{egujarati}{0A8F}
\pdfglyphtounicode{eharmenian}{0567}
\pdfglyphtounicode{ehbopomofo}{311D}
\pdfglyphtounicode{ehiragana}{3048}
\pdfglyphtounicode{ehookabove}{1EBB}
\pdfglyphtounicode{eibopomofo}{311F}
\pdfglyphtounicode{eight}{0038}
\pdfglyphtounicode{eightarabic}{0668}
\pdfglyphtounicode{eightbengali}{09EE}
\pdfglyphtounicode{eightcircle}{2467}
\pdfglyphtounicode{eightcircleinversesansserif}{2791}
\pdfglyphtounicode{eightdeva}{096E}
\pdfglyphtounicode{eighteencircle}{2471}
\pdfglyphtounicode{eighteenparen}{2485}
\pdfglyphtounicode{eighteenperiod}{2499}
\pdfglyphtounicode{eightgujarati}{0AEE}
\pdfglyphtounicode{eightgurmukhi}{0A6E}
\pdfglyphtounicode{eighthackarabic}{0668}
\pdfglyphtounicode{eighthangzhou}{3028}
\pdfglyphtounicode{eighthnotebeamed}{266B}
\pdfglyphtounicode{eightideographicparen}{3227}
\pdfglyphtounicode{eightinferior}{2088}
\pdfglyphtounicode{eightmonospace}{FF18}
\pdfglyphtounicode{eightoldstyle}{0038}
\pdfglyphtounicode{eightparen}{247B}
\pdfglyphtounicode{eightperiod}{248F}
\pdfglyphtounicode{eightpersian}{06F8}
\pdfglyphtounicode{eightroman}{2177}
\pdfglyphtounicode{eightsuperior}{2078}
\pdfglyphtounicode{eightthai}{0E58}
\pdfglyphtounicode{einvertedbreve}{0207}
\pdfglyphtounicode{eiotifiedcyrillic}{0465}
\pdfglyphtounicode{ekatakana}{30A8}
\pdfglyphtounicode{ekatakanahalfwidth}{FF74}
\pdfglyphtounicode{ekonkargurmukhi}{0A74}
\pdfglyphtounicode{ekorean}{3154}
\pdfglyphtounicode{elcyrillic}{043B}
\pdfglyphtounicode{element}{2208}
\pdfglyphtounicode{elevencircle}{246A}
\pdfglyphtounicode{elevenparen}{247E}
\pdfglyphtounicode{elevenperiod}{2492}
\pdfglyphtounicode{elevenroman}{217A}
\pdfglyphtounicode{ellipsis}{2026}
\pdfglyphtounicode{ellipsisvertical}{22EE}
\pdfglyphtounicode{emacron}{0113}
\pdfglyphtounicode{emacronacute}{1E17}
\pdfglyphtounicode{emacrongrave}{1E15}
\pdfglyphtounicode{emcyrillic}{043C}
\pdfglyphtounicode{emdash}{2014}
\pdfglyphtounicode{emdashvertical}{FE31}
\pdfglyphtounicode{emonospace}{FF45}
\pdfglyphtounicode{emphasismarkarmenian}{055B}
\pdfglyphtounicode{emptyset}{2205}
\pdfglyphtounicode{enbopomofo}{3123}
\pdfglyphtounicode{encyrillic}{043D}
\pdfglyphtounicode{endash}{2013}
\pdfglyphtounicode{endashvertical}{FE32}
\pdfglyphtounicode{endescendercyrillic}{04A3}
\pdfglyphtounicode{eng}{014B}
\pdfglyphtounicode{engbopomofo}{3125}
\pdfglyphtounicode{enghecyrillic}{04A5}
\pdfglyphtounicode{enhookcyrillic}{04C8}
\pdfglyphtounicode{enspace}{2002}
\pdfglyphtounicode{eogonek}{0119}
\pdfglyphtounicode{eokorean}{3153}
\pdfglyphtounicode{eopen}{025B}
\pdfglyphtounicode{eopenclosed}{029A}
\pdfglyphtounicode{eopenreversed}{025C}
\pdfglyphtounicode{eopenreversedclosed}{025E}
\pdfglyphtounicode{eopenreversedhook}{025D}
\pdfglyphtounicode{eparen}{24A0}
\pdfglyphtounicode{epsilon}{03B5}
\pdfglyphtounicode{epsilon1}{03F5}
\pdfglyphtounicode{epsiloninv}{03F6}
\pdfglyphtounicode{epsilontonos}{03AD}
\pdfglyphtounicode{equal}{003D}
\pdfglyphtounicode{equaldotleftright}{2252}
\pdfglyphtounicode{equaldotrightleft}{2253}
\pdfglyphtounicode{equalmonospace}{FF1D}
\pdfglyphtounicode{equalorfollows}{22DF}
\pdfglyphtounicode{equalorgreater}{2A96}
\pdfglyphtounicode{equalorless}{2A95}
\pdfglyphtounicode{equalorprecedes}{22DE}
\pdfglyphtounicode{equalorsimilar}{2242}
\pdfglyphtounicode{equalsdots}{2251}
\pdfglyphtounicode{equalsmall}{FE66}
\pdfglyphtounicode{equalsuperior}{207C}
\pdfglyphtounicode{equivalence}{2261}
\pdfglyphtounicode{equivasymptotic}{224D}
\pdfglyphtounicode{erbopomofo}{3126}
\pdfglyphtounicode{ercyrillic}{0440}
\pdfglyphtounicode{ereversed}{0258}
\pdfglyphtounicode{ereversedcyrillic}{044D}
\pdfglyphtounicode{escyrillic}{0441}
\pdfglyphtounicode{esdescendercyrillic}{04AB}
\pdfglyphtounicode{esh}{0283}
\pdfglyphtounicode{eshcurl}{0286}
\pdfglyphtounicode{eshortdeva}{090E}
\pdfglyphtounicode{eshortvowelsigndeva}{0946}
\pdfglyphtounicode{eshreversedloop}{01AA}
\pdfglyphtounicode{eshsquatreversed}{0285}
\pdfglyphtounicode{esmallhiragana}{3047}
\pdfglyphtounicode{esmallkatakana}{30A7}
\pdfglyphtounicode{esmallkatakanahalfwidth}{FF6A}
\pdfglyphtounicode{estimated}{212E}
\pdfglyphtounicode{esuperior}{0065}
\pdfglyphtounicode{eta}{03B7}
\pdfglyphtounicode{etarmenian}{0568}
\pdfglyphtounicode{etatonos}{03AE}
\pdfglyphtounicode{eth}{00F0}
\pdfglyphtounicode{etilde}{1EBD}
\pdfglyphtounicode{etildebelow}{1E1B}
\pdfglyphtounicode{etnahtafoukhhebrew}{0591}
\pdfglyphtounicode{etnahtafoukhlefthebrew}{0591}
\pdfglyphtounicode{etnahtahebrew}{0591}
\pdfglyphtounicode{etnahtalefthebrew}{0591}
\pdfglyphtounicode{eturned}{01DD}
\pdfglyphtounicode{eukorean}{3161}
\pdfglyphtounicode{euro}{20AC}
\pdfglyphtounicode{evowelsignbengali}{09C7}
\pdfglyphtounicode{evowelsigndeva}{0947}
\pdfglyphtounicode{evowelsigngujarati}{0AC7}
\pdfglyphtounicode{exclam}{0021}
\pdfglyphtounicode{exclamarmenian}{055C}
\pdfglyphtounicode{exclamdbl}{203C}
\pdfglyphtounicode{exclamdown}{00A1}
\pdfglyphtounicode{exclamdownsmall}{00A1}
\pdfglyphtounicode{exclammonospace}{FF01}
\pdfglyphtounicode{exclamsmall}{0021}
\pdfglyphtounicode{existential}{2203}
\pdfglyphtounicode{ezh}{0292}
\pdfglyphtounicode{ezhcaron}{01EF}
\pdfglyphtounicode{ezhcurl}{0293}
\pdfglyphtounicode{ezhreversed}{01B9}
\pdfglyphtounicode{ezhtail}{01BA}
\pdfglyphtounicode{f}{0066}
\pdfglyphtounicode{fadeva}{095E}
\pdfglyphtounicode{fagurmukhi}{0A5E}
\pdfglyphtounicode{fahrenheit}{2109}
\pdfglyphtounicode{fathaarabic}{064E}
\pdfglyphtounicode{fathalowarabic}{064E}
\pdfglyphtounicode{fathatanarabic}{064B}
\pdfglyphtounicode{fbopomofo}{3108}
\pdfglyphtounicode{fcircle}{24D5}
\pdfglyphtounicode{fdotaccent}{1E1F}
\pdfglyphtounicode{feharabic}{0641}
\pdfglyphtounicode{feharmenian}{0586}
\pdfglyphtounicode{fehfinalarabic}{FED2}
\pdfglyphtounicode{fehinitialarabic}{FED3}
\pdfglyphtounicode{fehmedialarabic}{FED4}
\pdfglyphtounicode{feicoptic}{03E5}
\pdfglyphtounicode{female}{2640}
\pdfglyphtounicode{ff}{0066 0066}
\pdfglyphtounicode{ffi}{0066 0066 0069}
\pdfglyphtounicode{ffl}{0066 0066 006C}
\pdfglyphtounicode{fi}{0066 0069}
\pdfglyphtounicode{fifteencircle}{246E}
\pdfglyphtounicode{fifteenparen}{2482}
\pdfglyphtounicode{fifteenperiod}{2496}
\pdfglyphtounicode{figuredash}{2012}
\pdfglyphtounicode{filledbox}{25A0}
\pdfglyphtounicode{filledrect}{25AC}
\pdfglyphtounicode{finalkaf}{05DA}
\pdfglyphtounicode{finalkafdagesh}{FB3A}
\pdfglyphtounicode{finalkafdageshhebrew}{FB3A}
\pdfglyphtounicode{finalkafhebrew}{05DA}
\pdfglyphtounicode{finalkafqamats}{05DA 05B8}
\pdfglyphtounicode{finalkafqamatshebrew}{05DA 05B8}
\pdfglyphtounicode{finalkafsheva}{05DA 05B0}
\pdfglyphtounicode{finalkafshevahebrew}{05DA 05B0}
\pdfglyphtounicode{finalmem}{05DD}
\pdfglyphtounicode{finalmemhebrew}{05DD}
\pdfglyphtounicode{finalnun}{05DF}
\pdfglyphtounicode{finalnunhebrew}{05DF}
\pdfglyphtounicode{finalpe}{05E3}
\pdfglyphtounicode{finalpehebrew}{05E3}
\pdfglyphtounicode{finaltsadi}{05E5}
\pdfglyphtounicode{finaltsadihebrew}{05E5}
\pdfglyphtounicode{firsttonechinese}{02C9}
\pdfglyphtounicode{fisheye}{25C9}
\pdfglyphtounicode{fitacyrillic}{0473}
\pdfglyphtounicode{five}{0035}
\pdfglyphtounicode{fivearabic}{0665}
\pdfglyphtounicode{fivebengali}{09EB}
\pdfglyphtounicode{fivecircle}{2464}
\pdfglyphtounicode{fivecircleinversesansserif}{278E}
\pdfglyphtounicode{fivedeva}{096B}
\pdfglyphtounicode{fiveeighths}{215D}
\pdfglyphtounicode{fivegujarati}{0AEB}
\pdfglyphtounicode{fivegurmukhi}{0A6B}
\pdfglyphtounicode{fivehackarabic}{0665}
\pdfglyphtounicode{fivehangzhou}{3025}
\pdfglyphtounicode{fiveideographicparen}{3224}
\pdfglyphtounicode{fiveinferior}{2085}
\pdfglyphtounicode{fivemonospace}{FF15}
\pdfglyphtounicode{fiveoldstyle}{0035}
\pdfglyphtounicode{fiveparen}{2478}
\pdfglyphtounicode{fiveperiod}{248C}
\pdfglyphtounicode{fivepersian}{06F5}
\pdfglyphtounicode{fiveroman}{2174}
\pdfglyphtounicode{fivesuperior}{2075}
\pdfglyphtounicode{fivethai}{0E55}
\pdfglyphtounicode{fl}{0066 006C}
\pdfglyphtounicode{flat}{266D}
\pdfglyphtounicode{floorleft}{230A}
\pdfglyphtounicode{floorright}{230B}
\pdfglyphtounicode{florin}{0192}
\pdfglyphtounicode{fmonospace}{FF46}
\pdfglyphtounicode{fmsquare}{3399}
\pdfglyphtounicode{fofanthai}{0E1F}
\pdfglyphtounicode{fofathai}{0E1D}
\pdfglyphtounicode{follownotdbleqv}{2ABA}
\pdfglyphtounicode{follownotslnteql}{2AB6}
\pdfglyphtounicode{followornoteqvlnt}{22E9}
\pdfglyphtounicode{follows}{227B}
\pdfglyphtounicode{followsequal}{2AB0}
\pdfglyphtounicode{followsorcurly}{227D}
\pdfglyphtounicode{followsorequal}{227F}
\pdfglyphtounicode{fongmanthai}{0E4F}
\pdfglyphtounicode{forall}{2200}
\pdfglyphtounicode{forces}{22A9}
\pdfglyphtounicode{forcesbar}{22AA}
\pdfglyphtounicode{fork}{22D4}
\pdfglyphtounicode{four}{0034}
\pdfglyphtounicode{fourarabic}{0664}
\pdfglyphtounicode{fourbengali}{09EA}
\pdfglyphtounicode{fourcircle}{2463}
\pdfglyphtounicode{fourcircleinversesansserif}{278D}
\pdfglyphtounicode{fourdeva}{096A}
\pdfglyphtounicode{fourgujarati}{0AEA}
\pdfglyphtounicode{fourgurmukhi}{0A6A}
\pdfglyphtounicode{fourhackarabic}{0664}
\pdfglyphtounicode{fourhangzhou}{3024}
\pdfglyphtounicode{fourideographicparen}{3223}
\pdfglyphtounicode{fourinferior}{2084}
\pdfglyphtounicode{fourmonospace}{FF14}
\pdfglyphtounicode{fournumeratorbengali}{09F7}
\pdfglyphtounicode{fouroldstyle}{0034}
\pdfglyphtounicode{fourparen}{2477}
\pdfglyphtounicode{fourperiod}{248B}
\pdfglyphtounicode{fourpersian}{06F4}
\pdfglyphtounicode{fourroman}{2173}
\pdfglyphtounicode{foursuperior}{2074}
\pdfglyphtounicode{fourteencircle}{246D}
\pdfglyphtounicode{fourteenparen}{2481}
\pdfglyphtounicode{fourteenperiod}{2495}
\pdfglyphtounicode{fourthai}{0E54}
\pdfglyphtounicode{fourthtonechinese}{02CB}
\pdfglyphtounicode{fparen}{24A1}
\pdfglyphtounicode{fraction}{2044}
\pdfglyphtounicode{franc}{20A3}
\pdfglyphtounicode{frown}{2322}
\pdfglyphtounicode{g}{0067}
\pdfglyphtounicode{gabengali}{0997}
\pdfglyphtounicode{gacute}{01F5}
\pdfglyphtounicode{gadeva}{0917}
\pdfglyphtounicode{gafarabic}{06AF}
\pdfglyphtounicode{gaffinalarabic}{FB93}
\pdfglyphtounicode{gafinitialarabic}{FB94}
\pdfglyphtounicode{gafmedialarabic}{FB95}
\pdfglyphtounicode{gagujarati}{0A97}
\pdfglyphtounicode{gagurmukhi}{0A17}
\pdfglyphtounicode{gahiragana}{304C}
\pdfglyphtounicode{gakatakana}{30AC}
\pdfglyphtounicode{gamma}{03B3}
\pdfglyphtounicode{gammalatinsmall}{0263}
\pdfglyphtounicode{gammasuperior}{02E0}
\pdfglyphtounicode{gangiacoptic}{03EB}
\pdfglyphtounicode{gbopomofo}{310D}
\pdfglyphtounicode{gbreve}{011F}
\pdfglyphtounicode{gcaron}{01E7}
\pdfglyphtounicode{gcedilla}{0123}
\pdfglyphtounicode{gcircle}{24D6}
\pdfglyphtounicode{gcircumflex}{011D}
\pdfglyphtounicode{gcommaaccent}{0123}
\pdfglyphtounicode{gdot}{0121}
\pdfglyphtounicode{gdotaccent}{0121}
\pdfglyphtounicode{gecyrillic}{0433}
\pdfglyphtounicode{gehiragana}{3052}
\pdfglyphtounicode{gekatakana}{30B2}
\pdfglyphtounicode{geomequivalent}{224E}
\pdfglyphtounicode{geometricallyequal}{2251}
\pdfglyphtounicode{gereshaccenthebrew}{059C}
\pdfglyphtounicode{gereshhebrew}{05F3}
\pdfglyphtounicode{gereshmuqdamhebrew}{059D}
\pdfglyphtounicode{germandbls}{00DF}
\pdfglyphtounicode{gershayimaccenthebrew}{059E}
\pdfglyphtounicode{gershayimhebrew}{05F4}
\pdfglyphtounicode{getamark}{3013}
\pdfglyphtounicode{ghabengali}{0998}
\pdfglyphtounicode{ghadarmenian}{0572}
\pdfglyphtounicode{ghadeva}{0918}
\pdfglyphtounicode{ghagujarati}{0A98}
\pdfglyphtounicode{ghagurmukhi}{0A18}
\pdfglyphtounicode{ghainarabic}{063A}
\pdfglyphtounicode{ghainfinalarabic}{FECE}
\pdfglyphtounicode{ghaininitialarabic}{FECF}
\pdfglyphtounicode{ghainmedialarabic}{FED0}
\pdfglyphtounicode{ghemiddlehookcyrillic}{0495}
\pdfglyphtounicode{ghestrokecyrillic}{0493}
\pdfglyphtounicode{gheupturncyrillic}{0491}
\pdfglyphtounicode{ghhadeva}{095A}
\pdfglyphtounicode{ghhagurmukhi}{0A5A}
\pdfglyphtounicode{ghook}{0260}
\pdfglyphtounicode{ghzsquare}{3393}
\pdfglyphtounicode{gihiragana}{304E}
\pdfglyphtounicode{gikatakana}{30AE}
\pdfglyphtounicode{gimarmenian}{0563}
\pdfglyphtounicode{gimel}{05D2}
\pdfglyphtounicode{gimeldagesh}{FB32}
\pdfglyphtounicode{gimeldageshhebrew}{FB32}
\pdfglyphtounicode{gimelhebrew}{05D2}
\pdfglyphtounicode{gjecyrillic}{0453}
\pdfglyphtounicode{glottalinvertedstroke}{01BE}
\pdfglyphtounicode{glottalstop}{0294}
\pdfglyphtounicode{glottalstopinverted}{0296}
\pdfglyphtounicode{glottalstopmod}{02C0}
\pdfglyphtounicode{glottalstopreversed}{0295}
\pdfglyphtounicode{glottalstopreversedmod}{02C1}
\pdfglyphtounicode{glottalstopreversedsuperior}{02E4}
\pdfglyphtounicode{glottalstopstroke}{02A1}
\pdfglyphtounicode{glottalstopstrokereversed}{02A2}
\pdfglyphtounicode{gmacron}{1E21}
\pdfglyphtounicode{gmonospace}{FF47}
\pdfglyphtounicode{gohiragana}{3054}
\pdfglyphtounicode{gokatakana}{30B4}
\pdfglyphtounicode{gparen}{24A2}
\pdfglyphtounicode{gpasquare}{33AC}
\pdfglyphtounicode{gradient}{2207}
\pdfglyphtounicode{grave}{0060}
\pdfglyphtounicode{gravebelowcmb}{0316}
\pdfglyphtounicode{gravecmb}{0300}
\pdfglyphtounicode{gravecomb}{0300}
\pdfglyphtounicode{gravedeva}{0953}
\pdfglyphtounicode{gravelowmod}{02CE}
\pdfglyphtounicode{gravemonospace}{FF40}
\pdfglyphtounicode{gravetonecmb}{0340}
\pdfglyphtounicode{greater}{003E}
\pdfglyphtounicode{greaterdbleqlless}{2A8C}
\pdfglyphtounicode{greaterdblequal}{2267}
\pdfglyphtounicode{greaterdot}{22D7}
\pdfglyphtounicode{greaterequal}{2265}
\pdfglyphtounicode{greaterequalorless}{22DB}
\pdfglyphtounicode{greaterlessequal}{22DB}
\pdfglyphtounicode{greatermonospace}{FF1E}
\pdfglyphtounicode{greatermuch}{226B}
\pdfglyphtounicode{greaternotdblequal}{2A8A}
\pdfglyphtounicode{greaternotequal}{2A88}
\pdfglyphtounicode{greaterorapproxeql}{2A86}
\pdfglyphtounicode{greaterorequalslant}{2A7E}
\pdfglyphtounicode{greaterorequivalent}{2273}
\pdfglyphtounicode{greaterorless}{2277}
\pdfglyphtounicode{greaterornotdbleql}{2269}
\pdfglyphtounicode{greaterornotequal}{2269}
\pdfglyphtounicode{greaterorsimilar}{2273}
\pdfglyphtounicode{greateroverequal}{2267}
\pdfglyphtounicode{greatersmall}{FE65}
\pdfglyphtounicode{gscript}{0261}
\pdfglyphtounicode{gstroke}{01E5}
\pdfglyphtounicode{guhiragana}{3050}
\pdfglyphtounicode{guillemotleft}{00AB}
\pdfglyphtounicode{guillemotright}{00BB}
\pdfglyphtounicode{guilsinglleft}{2039}
\pdfglyphtounicode{guilsinglright}{203A}
\pdfglyphtounicode{gukatakana}{30B0}
\pdfglyphtounicode{guramusquare}{3318}
\pdfglyphtounicode{gysquare}{33C9}
\pdfglyphtounicode{h}{0068}
\pdfglyphtounicode{haabkhasiancyrillic}{04A9}
\pdfglyphtounicode{haaltonearabic}{06C1}
\pdfglyphtounicode{habengali}{09B9}
\pdfglyphtounicode{hadescendercyrillic}{04B3}
\pdfglyphtounicode{hadeva}{0939}
\pdfglyphtounicode{hagujarati}{0AB9}
\pdfglyphtounicode{hagurmukhi}{0A39}
\pdfglyphtounicode{haharabic}{062D}
\pdfglyphtounicode{hahfinalarabic}{FEA2}
\pdfglyphtounicode{hahinitialarabic}{FEA3}
\pdfglyphtounicode{hahiragana}{306F}
\pdfglyphtounicode{hahmedialarabic}{FEA4}
\pdfglyphtounicode{haitusquare}{332A}
\pdfglyphtounicode{hakatakana}{30CF}
\pdfglyphtounicode{hakatakanahalfwidth}{FF8A}
\pdfglyphtounicode{halantgurmukhi}{0A4D}
\pdfglyphtounicode{hamzaarabic}{0621}
\pdfglyphtounicode{hamzadammaarabic}{0621 064F}
\pdfglyphtounicode{hamzadammatanarabic}{0621 064C}
\pdfglyphtounicode{hamzafathaarabic}{0621 064E}
\pdfglyphtounicode{hamzafathatanarabic}{0621 064B}
\pdfglyphtounicode{hamzalowarabic}{0621}
\pdfglyphtounicode{hamzalowkasraarabic}{0621 0650}
\pdfglyphtounicode{hamzalowkasratanarabic}{0621 064D}
\pdfglyphtounicode{hamzasukunarabic}{0621 0652}
\pdfglyphtounicode{hangulfiller}{3164}
\pdfglyphtounicode{hardsigncyrillic}{044A}
\pdfglyphtounicode{harpoondownleft}{21C3}
\pdfglyphtounicode{harpoondownright}{21C2}
\pdfglyphtounicode{harpoonleftbarbup}{21BC}
\pdfglyphtounicode{harpoonleftright}{21CC}
\pdfglyphtounicode{harpoonrightbarbup}{21C0}
\pdfglyphtounicode{harpoonrightleft}{21CB}
\pdfglyphtounicode{harpoonupleft}{21BF}
\pdfglyphtounicode{harpoonupright}{21BE}
\pdfglyphtounicode{hasquare}{33CA}
\pdfglyphtounicode{hatafpatah}{05B2}
\pdfglyphtounicode{hatafpatah16}{05B2}
\pdfglyphtounicode{hatafpatah23}{05B2}
\pdfglyphtounicode{hatafpatah2f}{05B2}
\pdfglyphtounicode{hatafpatahhebrew}{05B2}
\pdfglyphtounicode{hatafpatahnarrowhebrew}{05B2}
\pdfglyphtounicode{hatafpatahquarterhebrew}{05B2}
\pdfglyphtounicode{hatafpatahwidehebrew}{05B2}
\pdfglyphtounicode{hatafqamats}{05B3}
\pdfglyphtounicode{hatafqamats1b}{05B3}
\pdfglyphtounicode{hatafqamats28}{05B3}
\pdfglyphtounicode{hatafqamats34}{05B3}
\pdfglyphtounicode{hatafqamatshebrew}{05B3}
\pdfglyphtounicode{hatafqamatsnarrowhebrew}{05B3}
\pdfglyphtounicode{hatafqamatsquarterhebrew}{05B3}
\pdfglyphtounicode{hatafqamatswidehebrew}{05B3}
\pdfglyphtounicode{hatafsegol}{05B1}
\pdfglyphtounicode{hatafsegol17}{05B1}
\pdfglyphtounicode{hatafsegol24}{05B1}
\pdfglyphtounicode{hatafsegol30}{05B1}
\pdfglyphtounicode{hatafsegolhebrew}{05B1}
\pdfglyphtounicode{hatafsegolnarrowhebrew}{05B1}
\pdfglyphtounicode{hatafsegolquarterhebrew}{05B1}
\pdfglyphtounicode{hatafsegolwidehebrew}{05B1}
\pdfglyphtounicode{hbar}{0127}
\pdfglyphtounicode{hbopomofo}{310F}
\pdfglyphtounicode{hbrevebelow}{1E2B}
\pdfglyphtounicode{hcedilla}{1E29}
\pdfglyphtounicode{hcircle}{24D7}
\pdfglyphtounicode{hcircumflex}{0125}
\pdfglyphtounicode{hdieresis}{1E27}
\pdfglyphtounicode{hdotaccent}{1E23}
\pdfglyphtounicode{hdotbelow}{1E25}
\pdfglyphtounicode{he}{05D4}
\pdfglyphtounicode{heart}{2665}
\pdfglyphtounicode{heartsuitblack}{2665}
\pdfglyphtounicode{heartsuitwhite}{2661}
\pdfglyphtounicode{hedagesh}{FB34}
\pdfglyphtounicode{hedageshhebrew}{FB34}
\pdfglyphtounicode{hehaltonearabic}{06C1}
\pdfglyphtounicode{heharabic}{0647}
\pdfglyphtounicode{hehebrew}{05D4}
\pdfglyphtounicode{hehfinalaltonearabic}{FBA7}
\pdfglyphtounicode{hehfinalalttwoarabic}{FEEA}
\pdfglyphtounicode{hehfinalarabic}{FEEA}
\pdfglyphtounicode{hehhamzaabovefinalarabic}{FBA5}
\pdfglyphtounicode{hehhamzaaboveisolatedarabic}{FBA4}
\pdfglyphtounicode{hehinitialaltonearabic}{FBA8}
\pdfglyphtounicode{hehinitialarabic}{FEEB}
\pdfglyphtounicode{hehiragana}{3078}
\pdfglyphtounicode{hehmedialaltonearabic}{FBA9}
\pdfglyphtounicode{hehmedialarabic}{FEEC}
\pdfglyphtounicode{heiseierasquare}{337B}
\pdfglyphtounicode{hekatakana}{30D8}
\pdfglyphtounicode{hekatakanahalfwidth}{FF8D}
\pdfglyphtounicode{hekutaarusquare}{3336}
\pdfglyphtounicode{henghook}{0267}
\pdfglyphtounicode{herutusquare}{3339}
\pdfglyphtounicode{het}{05D7}
\pdfglyphtounicode{hethebrew}{05D7}
\pdfglyphtounicode{hhook}{0266}
\pdfglyphtounicode{hhooksuperior}{02B1}
\pdfglyphtounicode{hieuhacirclekorean}{327B}
\pdfglyphtounicode{hieuhaparenkorean}{321B}
\pdfglyphtounicode{hieuhcirclekorean}{326D}
\pdfglyphtounicode{hieuhkorean}{314E}
\pdfglyphtounicode{hieuhparenkorean}{320D}
\pdfglyphtounicode{hihiragana}{3072}
\pdfglyphtounicode{hikatakana}{30D2}
\pdfglyphtounicode{hikatakanahalfwidth}{FF8B}
\pdfglyphtounicode{hiriq}{05B4}
\pdfglyphtounicode{hiriq14}{05B4}
\pdfglyphtounicode{hiriq21}{05B4}
\pdfglyphtounicode{hiriq2d}{05B4}
\pdfglyphtounicode{hiriqhebrew}{05B4}
\pdfglyphtounicode{hiriqnarrowhebrew}{05B4}
\pdfglyphtounicode{hiriqquarterhebrew}{05B4}
\pdfglyphtounicode{hiriqwidehebrew}{05B4}
\pdfglyphtounicode{hlinebelow}{1E96}
\pdfglyphtounicode{hmonospace}{FF48}
\pdfglyphtounicode{hoarmenian}{0570}
\pdfglyphtounicode{hohipthai}{0E2B}
\pdfglyphtounicode{hohiragana}{307B}
\pdfglyphtounicode{hokatakana}{30DB}
\pdfglyphtounicode{hokatakanahalfwidth}{FF8E}
\pdfglyphtounicode{holam}{05B9}
\pdfglyphtounicode{holam19}{05B9}
\pdfglyphtounicode{holam26}{05B9}
\pdfglyphtounicode{holam32}{05B9}
\pdfglyphtounicode{holamhebrew}{05B9}
\pdfglyphtounicode{holamnarrowhebrew}{05B9}
\pdfglyphtounicode{holamquarterhebrew}{05B9}
\pdfglyphtounicode{holamwidehebrew}{05B9}
\pdfglyphtounicode{honokhukthai}{0E2E}
\pdfglyphtounicode{hookabovecomb}{0309}
\pdfglyphtounicode{hookcmb}{0309}
\pdfglyphtounicode{hookpalatalizedbelowcmb}{0321}
\pdfglyphtounicode{hookretroflexbelowcmb}{0322}
\pdfglyphtounicode{hoonsquare}{3342}
\pdfglyphtounicode{horicoptic}{03E9}
\pdfglyphtounicode{horizontalbar}{2015}
\pdfglyphtounicode{horncmb}{031B}
\pdfglyphtounicode{hotsprings}{2668}
\pdfglyphtounicode{house}{2302}
\pdfglyphtounicode{hparen}{24A3}
\pdfglyphtounicode{hsuperior}{02B0}
\pdfglyphtounicode{hturned}{0265}
\pdfglyphtounicode{huhiragana}{3075}
\pdfglyphtounicode{huiitosquare}{3333}
\pdfglyphtounicode{hukatakana}{30D5}
\pdfglyphtounicode{hukatakanahalfwidth}{FF8C}
\pdfglyphtounicode{hungarumlaut}{02DD}
\pdfglyphtounicode{hungarumlautcmb}{030B}
\pdfglyphtounicode{hv}{0195}
\pdfglyphtounicode{hyphen}{002D}
\pdfglyphtounicode{hyphenchar}{002D}
\pdfglyphtounicode{hypheninferior}{002D}
\pdfglyphtounicode{hyphenmonospace}{FF0D}
\pdfglyphtounicode{hyphensmall}{FE63}
\pdfglyphtounicode{hyphensuperior}{002D}
\pdfglyphtounicode{hyphentwo}{2010}
\pdfglyphtounicode{i}{0069}
\pdfglyphtounicode{iacute}{00ED}
\pdfglyphtounicode{iacyrillic}{044F}
\pdfglyphtounicode{ibengali}{0987}
\pdfglyphtounicode{ibopomofo}{3127}
\pdfglyphtounicode{ibreve}{012D}
\pdfglyphtounicode{icaron}{01D0}
\pdfglyphtounicode{icircle}{24D8}
\pdfglyphtounicode{icircumflex}{00EE}
\pdfglyphtounicode{icyrillic}{0456}
\pdfglyphtounicode{idblgrave}{0209}
\pdfglyphtounicode{ideographearthcircle}{328F}
\pdfglyphtounicode{ideographfirecircle}{328B}
\pdfglyphtounicode{ideographicallianceparen}{323F}
\pdfglyphtounicode{ideographiccallparen}{323A}
\pdfglyphtounicode{ideographiccentrecircle}{32A5}
\pdfglyphtounicode{ideographicclose}{3006}
\pdfglyphtounicode{ideographiccomma}{3001}
\pdfglyphtounicode{ideographiccommaleft}{FF64}
\pdfglyphtounicode{ideographiccongratulationparen}{3237}
\pdfglyphtounicode{ideographiccorrectcircle}{32A3}
\pdfglyphtounicode{ideographicearthparen}{322F}
\pdfglyphtounicode{ideographicenterpriseparen}{323D}
\pdfglyphtounicode{ideographicexcellentcircle}{329D}
\pdfglyphtounicode{ideographicfestivalparen}{3240}
\pdfglyphtounicode{ideographicfinancialcircle}{3296}
\pdfglyphtounicode{ideographicfinancialparen}{3236}
\pdfglyphtounicode{ideographicfireparen}{322B}
\pdfglyphtounicode{ideographichaveparen}{3232}
\pdfglyphtounicode{ideographichighcircle}{32A4}
\pdfglyphtounicode{ideographiciterationmark}{3005}
\pdfglyphtounicode{ideographiclaborcircle}{3298}
\pdfglyphtounicode{ideographiclaborparen}{3238}
\pdfglyphtounicode{ideographicleftcircle}{32A7}
\pdfglyphtounicode{ideographiclowcircle}{32A6}
\pdfglyphtounicode{ideographicmedicinecircle}{32A9}
\pdfglyphtounicode{ideographicmetalparen}{322E}
\pdfglyphtounicode{ideographicmoonparen}{322A}
\pdfglyphtounicode{ideographicnameparen}{3234}
\pdfglyphtounicode{ideographicperiod}{3002}
\pdfglyphtounicode{ideographicprintcircle}{329E}
\pdfglyphtounicode{ideographicreachparen}{3243}
\pdfglyphtounicode{ideographicrepresentparen}{3239}
\pdfglyphtounicode{ideographicresourceparen}{323E}
\pdfglyphtounicode{ideographicrightcircle}{32A8}
\pdfglyphtounicode{ideographicsecretcircle}{3299}
\pdfglyphtounicode{ideographicselfparen}{3242}
\pdfglyphtounicode{ideographicsocietyparen}{3233}
\pdfglyphtounicode{ideographicspace}{3000}
\pdfglyphtounicode{ideographicspecialparen}{3235}
\pdfglyphtounicode{ideographicstockparen}{3231}
\pdfglyphtounicode{ideographicstudyparen}{323B}
\pdfglyphtounicode{ideographicsunparen}{3230}
\pdfglyphtounicode{ideographicsuperviseparen}{323C}
\pdfglyphtounicode{ideographicwaterparen}{322C}
\pdfglyphtounicode{ideographicwoodparen}{322D}
\pdfglyphtounicode{ideographiczero}{3007}
\pdfglyphtounicode{ideographmetalcircle}{328E}
\pdfglyphtounicode{ideographmooncircle}{328A}
\pdfglyphtounicode{ideographnamecircle}{3294}
\pdfglyphtounicode{ideographsuncircle}{3290}
\pdfglyphtounicode{ideographwatercircle}{328C}
\pdfglyphtounicode{ideographwoodcircle}{328D}
\pdfglyphtounicode{ideva}{0907}
\pdfglyphtounicode{idieresis}{00EF}
\pdfglyphtounicode{idieresisacute}{1E2F}
\pdfglyphtounicode{idieresiscyrillic}{04E5}
\pdfglyphtounicode{idotbelow}{1ECB}
\pdfglyphtounicode{iebrevecyrillic}{04D7}
\pdfglyphtounicode{iecyrillic}{0435}
\pdfglyphtounicode{ieungacirclekorean}{3275}
\pdfglyphtounicode{ieungaparenkorean}{3215}
\pdfglyphtounicode{ieungcirclekorean}{3267}
\pdfglyphtounicode{ieungkorean}{3147}
\pdfglyphtounicode{ieungparenkorean}{3207}
\pdfglyphtounicode{igrave}{00EC}
\pdfglyphtounicode{igujarati}{0A87}
\pdfglyphtounicode{igurmukhi}{0A07}
\pdfglyphtounicode{ihiragana}{3044}
\pdfglyphtounicode{ihookabove}{1EC9}
\pdfglyphtounicode{iibengali}{0988}
\pdfglyphtounicode{iicyrillic}{0438}
\pdfglyphtounicode{iideva}{0908}
\pdfglyphtounicode{iigujarati}{0A88}
\pdfglyphtounicode{iigurmukhi}{0A08}
\pdfglyphtounicode{iimatragurmukhi}{0A40}
\pdfglyphtounicode{iinvertedbreve}{020B}
\pdfglyphtounicode{iishortcyrillic}{0439}
\pdfglyphtounicode{iivowelsignbengali}{09C0}
\pdfglyphtounicode{iivowelsigndeva}{0940}
\pdfglyphtounicode{iivowelsigngujarati}{0AC0}
\pdfglyphtounicode{ij}{0133}
\pdfglyphtounicode{ikatakana}{30A4}
\pdfglyphtounicode{ikatakanahalfwidth}{FF72}
\pdfglyphtounicode{ikorean}{3163}
\pdfglyphtounicode{ilde}{02DC}
\pdfglyphtounicode{iluyhebrew}{05AC}
\pdfglyphtounicode{imacron}{012B}
\pdfglyphtounicode{imacroncyrillic}{04E3}
\pdfglyphtounicode{imageorapproximatelyequal}{2253}
\pdfglyphtounicode{imatragurmukhi}{0A3F}
\pdfglyphtounicode{imonospace}{FF49}
\pdfglyphtounicode{increment}{2206}
\pdfglyphtounicode{infinity}{221E}
\pdfglyphtounicode{iniarmenian}{056B}
\pdfglyphtounicode{integerdivide}{2216}
\pdfglyphtounicode{integral}{222B}
\pdfglyphtounicode{integralbottom}{2321}
\pdfglyphtounicode{integralbt}{2321}
\pdfglyphtounicode{integralex}{F8F5}
\pdfglyphtounicode{integraltop}{2320}
\pdfglyphtounicode{integraltp}{2320}
\pdfglyphtounicode{intercal}{22BA}
\pdfglyphtounicode{interrobang}{203D}
\pdfglyphtounicode{interrobangdown}{2E18}
\pdfglyphtounicode{intersection}{2229}
\pdfglyphtounicode{intersectiondbl}{22D2}
\pdfglyphtounicode{intersectionsq}{2293}
\pdfglyphtounicode{intisquare}{3305}
\pdfglyphtounicode{invbullet}{25D8}
\pdfglyphtounicode{invcircle}{25D9}
\pdfglyphtounicode{invsmileface}{263B}
\pdfglyphtounicode{iocyrillic}{0451}
\pdfglyphtounicode{iogonek}{012F}
\pdfglyphtounicode{iota}{03B9}
\pdfglyphtounicode{iotadieresis}{03CA}
\pdfglyphtounicode{iotadieresistonos}{0390}
\pdfglyphtounicode{iotalatin}{0269}
\pdfglyphtounicode{iotatonos}{03AF}
\pdfglyphtounicode{iparen}{24A4}
\pdfglyphtounicode{irigurmukhi}{0A72}
\pdfglyphtounicode{ismallhiragana}{3043}
\pdfglyphtounicode{ismallkatakana}{30A3}
\pdfglyphtounicode{ismallkatakanahalfwidth}{FF68}
\pdfglyphtounicode{issharbengali}{09FA}
\pdfglyphtounicode{istroke}{0268}
\pdfglyphtounicode{isuperior}{0069}
\pdfglyphtounicode{iterationhiragana}{309D}
\pdfglyphtounicode{iterationkatakana}{30FD}
\pdfglyphtounicode{itilde}{0129}
\pdfglyphtounicode{itildebelow}{1E2D}
\pdfglyphtounicode{iubopomofo}{3129}
\pdfglyphtounicode{iucyrillic}{044E}
\pdfglyphtounicode{ivowelsignbengali}{09BF}
\pdfglyphtounicode{ivowelsigndeva}{093F}
\pdfglyphtounicode{ivowelsigngujarati}{0ABF}
\pdfglyphtounicode{izhitsacyrillic}{0475}
\pdfglyphtounicode{izhitsadblgravecyrillic}{0477}
\pdfglyphtounicode{j}{006A}
\pdfglyphtounicode{jaarmenian}{0571}
\pdfglyphtounicode{jabengali}{099C}
\pdfglyphtounicode{jadeva}{091C}
\pdfglyphtounicode{jagujarati}{0A9C}
\pdfglyphtounicode{jagurmukhi}{0A1C}
\pdfglyphtounicode{jbopomofo}{3110}
\pdfglyphtounicode{jcaron}{01F0}
\pdfglyphtounicode{jcircle}{24D9}
\pdfglyphtounicode{jcircumflex}{0135}
\pdfglyphtounicode{jcrossedtail}{029D}
\pdfglyphtounicode{jdotlessstroke}{025F}
\pdfglyphtounicode{jecyrillic}{0458}
\pdfglyphtounicode{jeemarabic}{062C}
\pdfglyphtounicode{jeemfinalarabic}{FE9E}
\pdfglyphtounicode{jeeminitialarabic}{FE9F}
\pdfglyphtounicode{jeemmedialarabic}{FEA0}
\pdfglyphtounicode{jeharabic}{0698}
\pdfglyphtounicode{jehfinalarabic}{FB8B}
\pdfglyphtounicode{jhabengali}{099D}
\pdfglyphtounicode{jhadeva}{091D}
\pdfglyphtounicode{jhagujarati}{0A9D}
\pdfglyphtounicode{jhagurmukhi}{0A1D}
\pdfglyphtounicode{jheharmenian}{057B}
\pdfglyphtounicode{jis}{3004}
\pdfglyphtounicode{jmonospace}{FF4A}
\pdfglyphtounicode{jparen}{24A5}
\pdfglyphtounicode{jsuperior}{02B2}
\pdfglyphtounicode{k}{006B}
\pdfglyphtounicode{kabashkircyrillic}{04A1}
\pdfglyphtounicode{kabengali}{0995}
\pdfglyphtounicode{kacute}{1E31}
\pdfglyphtounicode{kacyrillic}{043A}
\pdfglyphtounicode{kadescendercyrillic}{049B}
\pdfglyphtounicode{kadeva}{0915}
\pdfglyphtounicode{kaf}{05DB}
\pdfglyphtounicode{kafarabic}{0643}
\pdfglyphtounicode{kafdagesh}{FB3B}
\pdfglyphtounicode{kafdageshhebrew}{FB3B}
\pdfglyphtounicode{kaffinalarabic}{FEDA}
\pdfglyphtounicode{kafhebrew}{05DB}
\pdfglyphtounicode{kafinitialarabic}{FEDB}
\pdfglyphtounicode{kafmedialarabic}{FEDC}
\pdfglyphtounicode{kafrafehebrew}{FB4D}
\pdfglyphtounicode{kagujarati}{0A95}
\pdfglyphtounicode{kagurmukhi}{0A15}
\pdfglyphtounicode{kahiragana}{304B}
\pdfglyphtounicode{kahookcyrillic}{04C4}
\pdfglyphtounicode{kakatakana}{30AB}
\pdfglyphtounicode{kakatakanahalfwidth}{FF76}
\pdfglyphtounicode{kappa}{03BA}
\pdfglyphtounicode{kappasymbolgreek}{03F0}
\pdfglyphtounicode{kapyeounmieumkorean}{3171}
\pdfglyphtounicode{kapyeounphieuphkorean}{3184}
\pdfglyphtounicode{kapyeounpieupkorean}{3178}
\pdfglyphtounicode{kapyeounssangpieupkorean}{3179}
\pdfglyphtounicode{karoriisquare}{330D}
\pdfglyphtounicode{kashidaautoarabic}{0640}
\pdfglyphtounicode{kashidaautonosidebearingarabic}{0640}
\pdfglyphtounicode{kasmallkatakana}{30F5}
\pdfglyphtounicode{kasquare}{3384}
\pdfglyphtounicode{kasraarabic}{0650}
\pdfglyphtounicode{kasratanarabic}{064D}
\pdfglyphtounicode{kastrokecyrillic}{049F}
\pdfglyphtounicode{katahiraprolongmarkhalfwidth}{FF70}
\pdfglyphtounicode{kaverticalstrokecyrillic}{049D}
\pdfglyphtounicode{kbopomofo}{310E}
\pdfglyphtounicode{kcalsquare}{3389}
\pdfglyphtounicode{kcaron}{01E9}
\pdfglyphtounicode{kcedilla}{0137}
\pdfglyphtounicode{kcircle}{24DA}
\pdfglyphtounicode{kcommaaccent}{0137}
\pdfglyphtounicode{kdotbelow}{1E33}
\pdfglyphtounicode{keharmenian}{0584}
\pdfglyphtounicode{kehiragana}{3051}
\pdfglyphtounicode{kekatakana}{30B1}
\pdfglyphtounicode{kekatakanahalfwidth}{FF79}
\pdfglyphtounicode{kenarmenian}{056F}
\pdfglyphtounicode{kesmallkatakana}{30F6}
\pdfglyphtounicode{kgreenlandic}{0138}
\pdfglyphtounicode{khabengali}{0996}
\pdfglyphtounicode{khacyrillic}{0445}
\pdfglyphtounicode{khadeva}{0916}
\pdfglyphtounicode{khagujarati}{0A96}
\pdfglyphtounicode{khagurmukhi}{0A16}
\pdfglyphtounicode{khaharabic}{062E}
\pdfglyphtounicode{khahfinalarabic}{FEA6}
\pdfglyphtounicode{khahinitialarabic}{FEA7}
\pdfglyphtounicode{khahmedialarabic}{FEA8}
\pdfglyphtounicode{kheicoptic}{03E7}
\pdfglyphtounicode{khhadeva}{0959}
\pdfglyphtounicode{khhagurmukhi}{0A59}
\pdfglyphtounicode{khieukhacirclekorean}{3278}
\pdfglyphtounicode{khieukhaparenkorean}{3218}
\pdfglyphtounicode{khieukhcirclekorean}{326A}
\pdfglyphtounicode{khieukhkorean}{314B}
\pdfglyphtounicode{khieukhparenkorean}{320A}
\pdfglyphtounicode{khokhaithai}{0E02}
\pdfglyphtounicode{khokhonthai}{0E05}
\pdfglyphtounicode{khokhuatthai}{0E03}
\pdfglyphtounicode{khokhwaithai}{0E04}
\pdfglyphtounicode{khomutthai}{0E5B}
\pdfglyphtounicode{khook}{0199}
\pdfglyphtounicode{khorakhangthai}{0E06}
\pdfglyphtounicode{khzsquare}{3391}
\pdfglyphtounicode{kihiragana}{304D}
\pdfglyphtounicode{kikatakana}{30AD}
\pdfglyphtounicode{kikatakanahalfwidth}{FF77}
\pdfglyphtounicode{kiroguramusquare}{3315}
\pdfglyphtounicode{kiromeetorusquare}{3316}
\pdfglyphtounicode{kirosquare}{3314}
\pdfglyphtounicode{kiyeokacirclekorean}{326E}
\pdfglyphtounicode{kiyeokaparenkorean}{320E}
\pdfglyphtounicode{kiyeokcirclekorean}{3260}
\pdfglyphtounicode{kiyeokkorean}{3131}
\pdfglyphtounicode{kiyeokparenkorean}{3200}
\pdfglyphtounicode{kiyeoksioskorean}{3133}
\pdfglyphtounicode{kjecyrillic}{045C}
\pdfglyphtounicode{klinebelow}{1E35}
\pdfglyphtounicode{klsquare}{3398}
\pdfglyphtounicode{kmcubedsquare}{33A6}
\pdfglyphtounicode{kmonospace}{FF4B}
\pdfglyphtounicode{kmsquaredsquare}{33A2}
\pdfglyphtounicode{kohiragana}{3053}
\pdfglyphtounicode{kohmsquare}{33C0}
\pdfglyphtounicode{kokaithai}{0E01}
\pdfglyphtounicode{kokatakana}{30B3}
\pdfglyphtounicode{kokatakanahalfwidth}{FF7A}
\pdfglyphtounicode{kooposquare}{331E}
\pdfglyphtounicode{koppacyrillic}{0481}
\pdfglyphtounicode{koreanstandardsymbol}{327F}
\pdfglyphtounicode{koroniscmb}{0343}
\pdfglyphtounicode{kparen}{24A6}
\pdfglyphtounicode{kpasquare}{33AA}
\pdfglyphtounicode{ksicyrillic}{046F}
\pdfglyphtounicode{ktsquare}{33CF}
\pdfglyphtounicode{kturned}{029E}
\pdfglyphtounicode{kuhiragana}{304F}
\pdfglyphtounicode{kukatakana}{30AF}
\pdfglyphtounicode{kukatakanahalfwidth}{FF78}
\pdfglyphtounicode{kvsquare}{33B8}
\pdfglyphtounicode{kwsquare}{33BE}
\pdfglyphtounicode{l}{006C}
\pdfglyphtounicode{labengali}{09B2}
\pdfglyphtounicode{lacute}{013A}
\pdfglyphtounicode{ladeva}{0932}
\pdfglyphtounicode{lagujarati}{0AB2}
\pdfglyphtounicode{lagurmukhi}{0A32}
\pdfglyphtounicode{lakkhangyaothai}{0E45}
\pdfglyphtounicode{lamaleffinalarabic}{FEFC}
\pdfglyphtounicode{lamalefhamzaabovefinalarabic}{FEF8}
\pdfglyphtounicode{lamalefhamzaaboveisolatedarabic}{FEF7}
\pdfglyphtounicode{lamalefhamzabelowfinalarabic}{FEFA}
\pdfglyphtounicode{lamalefhamzabelowisolatedarabic}{FEF9}
\pdfglyphtounicode{lamalefisolatedarabic}{FEFB}
\pdfglyphtounicode{lamalefmaddaabovefinalarabic}{FEF6}
\pdfglyphtounicode{lamalefmaddaaboveisolatedarabic}{FEF5}
\pdfglyphtounicode{lamarabic}{0644}
\pdfglyphtounicode{lambda}{03BB}
\pdfglyphtounicode{lambdastroke}{019B}
\pdfglyphtounicode{lamed}{05DC}
\pdfglyphtounicode{lameddagesh}{FB3C}
\pdfglyphtounicode{lameddageshhebrew}{FB3C}
\pdfglyphtounicode{lamedhebrew}{05DC}
\pdfglyphtounicode{lamedholam}{05DC 05B9}
\pdfglyphtounicode{lamedholamdagesh}{05DC 05B9 05BC}
\pdfglyphtounicode{lamedholamdageshhebrew}{05DC 05B9 05BC}
\pdfglyphtounicode{lamedholamhebrew}{05DC 05B9}
\pdfglyphtounicode{lamfinalarabic}{FEDE}
\pdfglyphtounicode{lamhahinitialarabic}{FCCA}
\pdfglyphtounicode{laminitialarabic}{FEDF}
\pdfglyphtounicode{lamjeeminitialarabic}{FCC9}
\pdfglyphtounicode{lamkhahinitialarabic}{FCCB}
\pdfglyphtounicode{lamlamhehisolatedarabic}{FDF2}
\pdfglyphtounicode{lammedialarabic}{FEE0}
\pdfglyphtounicode{lammeemhahinitialarabic}{FD88}
\pdfglyphtounicode{lammeeminitialarabic}{FCCC}
\pdfglyphtounicode{lammeemjeeminitialarabic}{FEDF FEE4 FEA0}
\pdfglyphtounicode{lammeemkhahinitialarabic}{FEDF FEE4 FEA8}
\pdfglyphtounicode{largecircle}{25EF}
\pdfglyphtounicode{latticetop}{22A4}
\pdfglyphtounicode{lbar}{019A}
\pdfglyphtounicode{lbelt}{026C}
\pdfglyphtounicode{lbopomofo}{310C}
\pdfglyphtounicode{lcaron}{013E}
\pdfglyphtounicode{lcedilla}{013C}
\pdfglyphtounicode{lcircle}{24DB}
\pdfglyphtounicode{lcircumflexbelow}{1E3D}
\pdfglyphtounicode{lcommaaccent}{013C}
\pdfglyphtounicode{ldot}{0140}
\pdfglyphtounicode{ldotaccent}{0140}
\pdfglyphtounicode{ldotbelow}{1E37}
\pdfglyphtounicode{ldotbelowmacron}{1E39}
\pdfglyphtounicode{leftangleabovecmb}{031A}
\pdfglyphtounicode{lefttackbelowcmb}{0318}
\pdfglyphtounicode{less}{003C}
\pdfglyphtounicode{lessdbleqlgreater}{2A8B}
\pdfglyphtounicode{lessdblequal}{2266}
\pdfglyphtounicode{lessdot}{22D6}
\pdfglyphtounicode{lessequal}{2264}
\pdfglyphtounicode{lessequalgreater}{22DA}
\pdfglyphtounicode{lessequalorgreater}{22DA}
\pdfglyphtounicode{lessmonospace}{FF1C}
\pdfglyphtounicode{lessmuch}{226A}
\pdfglyphtounicode{lessnotdblequal}{2A89}
\pdfglyphtounicode{lessnotequal}{2A87}
\pdfglyphtounicode{lessorapproxeql}{2A85}
\pdfglyphtounicode{lessorequalslant}{2A7D}
\pdfglyphtounicode{lessorequivalent}{2272}
\pdfglyphtounicode{lessorgreater}{2276}
\pdfglyphtounicode{lessornotdbleql}{2268}
\pdfglyphtounicode{lessornotequal}{2268}
\pdfglyphtounicode{lessorsimilar}{2272}
\pdfglyphtounicode{lessoverequal}{2266}
\pdfglyphtounicode{lesssmall}{FE64}
\pdfglyphtounicode{lezh}{026E}
\pdfglyphtounicode{lfblock}{258C}
\pdfglyphtounicode{lhookretroflex}{026D}
\pdfglyphtounicode{lira}{20A4}
\pdfglyphtounicode{liwnarmenian}{056C}
\pdfglyphtounicode{lj}{01C9}
\pdfglyphtounicode{ljecyrillic}{0459}
\pdfglyphtounicode{ll}{006C 006C}
\pdfglyphtounicode{lladeva}{0933}
\pdfglyphtounicode{llagujarati}{0AB3}
\pdfglyphtounicode{llinebelow}{1E3B}
\pdfglyphtounicode{llladeva}{0934}
\pdfglyphtounicode{llvocalicbengali}{09E1}
\pdfglyphtounicode{llvocalicdeva}{0961}
\pdfglyphtounicode{llvocalicvowelsignbengali}{09E3}
\pdfglyphtounicode{llvocalicvowelsigndeva}{0963}
\pdfglyphtounicode{lmiddletilde}{026B}
\pdfglyphtounicode{lmonospace}{FF4C}
\pdfglyphtounicode{lmsquare}{33D0}
\pdfglyphtounicode{lochulathai}{0E2C}
\pdfglyphtounicode{logicaland}{2227}
\pdfglyphtounicode{logicalnot}{00AC}
\pdfglyphtounicode{logicalnotreversed}{2310}
\pdfglyphtounicode{logicalor}{2228}
\pdfglyphtounicode{lolingthai}{0E25}
\pdfglyphtounicode{longdbls}{017F 017F}
\pdfglyphtounicode{longs}{017F}
\pdfglyphtounicode{longsh}{017F 0068}
\pdfglyphtounicode{longsi}{017F 0069}
\pdfglyphtounicode{longsl}{017F 006C}
\pdfglyphtounicode{longst}{017F 0074}
\pdfglyphtounicode{lowlinecenterline}{FE4E}
\pdfglyphtounicode{lowlinecmb}{0332}
\pdfglyphtounicode{lowlinedashed}{FE4D}
\pdfglyphtounicode{lozenge}{25CA}
\pdfglyphtounicode{lparen}{24A7}
\pdfglyphtounicode{lscript}{2113}
\pdfglyphtounicode{lslash}{0142}
\pdfglyphtounicode{lsquare}{2113}
\pdfglyphtounicode{lsuperior}{006C}
\pdfglyphtounicode{ltshade}{2591}
\pdfglyphtounicode{luthai}{0E26}
\pdfglyphtounicode{lvocalicbengali}{098C}
\pdfglyphtounicode{lvocalicdeva}{090C}
\pdfglyphtounicode{lvocalicvowelsignbengali}{09E2}
\pdfglyphtounicode{lvocalicvowelsigndeva}{0962}
\pdfglyphtounicode{lxsquare}{33D3}
\pdfglyphtounicode{m}{006D}
\pdfglyphtounicode{mabengali}{09AE}
\pdfglyphtounicode{macron}{00AF}
\pdfglyphtounicode{macronbelowcmb}{0331}
\pdfglyphtounicode{macroncmb}{0304}
\pdfglyphtounicode{macronlowmod}{02CD}
\pdfglyphtounicode{macronmonospace}{FFE3}
\pdfglyphtounicode{macute}{1E3F}
\pdfglyphtounicode{madeva}{092E}
\pdfglyphtounicode{magujarati}{0AAE}
\pdfglyphtounicode{magurmukhi}{0A2E}
\pdfglyphtounicode{mahapakhhebrew}{05A4}
\pdfglyphtounicode{mahapakhlefthebrew}{05A4}
\pdfglyphtounicode{mahiragana}{307E}
\pdfglyphtounicode{maichattawalowleftthai}{F895}
\pdfglyphtounicode{maichattawalowrightthai}{F894}
\pdfglyphtounicode{maichattawathai}{0E4B}
\pdfglyphtounicode{maichattawaupperleftthai}{F893}
\pdfglyphtounicode{maieklowleftthai}{F88C}
\pdfglyphtounicode{maieklowrightthai}{F88B}
\pdfglyphtounicode{maiekthai}{0E48}
\pdfglyphtounicode{maiekupperleftthai}{F88A}
\pdfglyphtounicode{maihanakatleftthai}{F884}
\pdfglyphtounicode{maihanakatthai}{0E31}
\pdfglyphtounicode{maitaikhuleftthai}{F889}
\pdfglyphtounicode{maitaikhuthai}{0E47}
\pdfglyphtounicode{maitholowleftthai}{F88F}
\pdfglyphtounicode{maitholowrightthai}{F88E}
\pdfglyphtounicode{maithothai}{0E49}
\pdfglyphtounicode{maithoupperleftthai}{F88D}
\pdfglyphtounicode{maitrilowleftthai}{F892}
\pdfglyphtounicode{maitrilowrightthai}{F891}
\pdfglyphtounicode{maitrithai}{0E4A}
\pdfglyphtounicode{maitriupperleftthai}{F890}
\pdfglyphtounicode{maiyamokthai}{0E46}
\pdfglyphtounicode{makatakana}{30DE}
\pdfglyphtounicode{makatakanahalfwidth}{FF8F}
\pdfglyphtounicode{male}{2642}
\pdfglyphtounicode{maltesecross}{2720}
\pdfglyphtounicode{mansyonsquare}{3347}
\pdfglyphtounicode{maqafhebrew}{05BE}
\pdfglyphtounicode{mars}{2642}
\pdfglyphtounicode{masoracirclehebrew}{05AF}
\pdfglyphtounicode{masquare}{3383}
\pdfglyphtounicode{mbopomofo}{3107}
\pdfglyphtounicode{mbsquare}{33D4}
\pdfglyphtounicode{mcircle}{24DC}
\pdfglyphtounicode{mcubedsquare}{33A5}
\pdfglyphtounicode{mdotaccent}{1E41}
\pdfglyphtounicode{mdotbelow}{1E43}
\pdfglyphtounicode{measuredangle}{2221}
\pdfglyphtounicode{meemarabic}{0645}
\pdfglyphtounicode{meemfinalarabic}{FEE2}
\pdfglyphtounicode{meeminitialarabic}{FEE3}
\pdfglyphtounicode{meemmedialarabic}{FEE4}
\pdfglyphtounicode{meemmeeminitialarabic}{FCD1}
\pdfglyphtounicode{meemmeemisolatedarabic}{FC48}
\pdfglyphtounicode{meetorusquare}{334D}
\pdfglyphtounicode{mehiragana}{3081}
\pdfglyphtounicode{meizierasquare}{337E}
\pdfglyphtounicode{mekatakana}{30E1}
\pdfglyphtounicode{mekatakanahalfwidth}{FF92}
\pdfglyphtounicode{mem}{05DE}
\pdfglyphtounicode{memdagesh}{FB3E}
\pdfglyphtounicode{memdageshhebrew}{FB3E}
\pdfglyphtounicode{memhebrew}{05DE}
\pdfglyphtounicode{menarmenian}{0574}
\pdfglyphtounicode{merkhahebrew}{05A5}
\pdfglyphtounicode{merkhakefulahebrew}{05A6}
\pdfglyphtounicode{merkhakefulalefthebrew}{05A6}
\pdfglyphtounicode{merkhalefthebrew}{05A5}
\pdfglyphtounicode{mhook}{0271}
\pdfglyphtounicode{mhzsquare}{3392}
\pdfglyphtounicode{middledotkatakanahalfwidth}{FF65}
\pdfglyphtounicode{middot}{00B7}
\pdfglyphtounicode{mieumacirclekorean}{3272}
\pdfglyphtounicode{mieumaparenkorean}{3212}
\pdfglyphtounicode{mieumcirclekorean}{3264}
\pdfglyphtounicode{mieumkorean}{3141}
\pdfglyphtounicode{mieumpansioskorean}{3170}
\pdfglyphtounicode{mieumparenkorean}{3204}
\pdfglyphtounicode{mieumpieupkorean}{316E}
\pdfglyphtounicode{mieumsioskorean}{316F}
\pdfglyphtounicode{mihiragana}{307F}
\pdfglyphtounicode{mikatakana}{30DF}
\pdfglyphtounicode{mikatakanahalfwidth}{FF90}
\pdfglyphtounicode{minus}{2212}
\pdfglyphtounicode{minusbelowcmb}{0320}
\pdfglyphtounicode{minuscircle}{2296}
\pdfglyphtounicode{minusmod}{02D7}
\pdfglyphtounicode{minusplus}{2213}
\pdfglyphtounicode{minute}{2032}
\pdfglyphtounicode{miribaarusquare}{334A}
\pdfglyphtounicode{mirisquare}{3349}
\pdfglyphtounicode{mlonglegturned}{0270}
\pdfglyphtounicode{mlsquare}{3396}
\pdfglyphtounicode{mmcubedsquare}{33A3}
\pdfglyphtounicode{mmonospace}{FF4D}
\pdfglyphtounicode{mmsquaredsquare}{339F}
\pdfglyphtounicode{mohiragana}{3082}
\pdfglyphtounicode{mohmsquare}{33C1}
\pdfglyphtounicode{mokatakana}{30E2}
\pdfglyphtounicode{mokatakanahalfwidth}{FF93}
\pdfglyphtounicode{molsquare}{33D6}
\pdfglyphtounicode{momathai}{0E21}
\pdfglyphtounicode{moverssquare}{33A7}
\pdfglyphtounicode{moverssquaredsquare}{33A8}
\pdfglyphtounicode{mparen}{24A8}
\pdfglyphtounicode{mpasquare}{33AB}
\pdfglyphtounicode{mssquare}{33B3}
\pdfglyphtounicode{msuperior}{006D}
\pdfglyphtounicode{mturned}{026F}
\pdfglyphtounicode{mu}{00B5}
\pdfglyphtounicode{mu1}{00B5}
\pdfglyphtounicode{muasquare}{3382}
\pdfglyphtounicode{muchgreater}{226B}
\pdfglyphtounicode{muchless}{226A}
\pdfglyphtounicode{mufsquare}{338C}
\pdfglyphtounicode{mugreek}{03BC}
\pdfglyphtounicode{mugsquare}{338D}
\pdfglyphtounicode{muhiragana}{3080}
\pdfglyphtounicode{mukatakana}{30E0}
\pdfglyphtounicode{mukatakanahalfwidth}{FF91}
\pdfglyphtounicode{mulsquare}{3395}
\pdfglyphtounicode{multicloseleft}{22C9}
\pdfglyphtounicode{multicloseright}{22CA}
\pdfglyphtounicode{multimap}{22B8}
\pdfglyphtounicode{multiopenleft}{22CB}
\pdfglyphtounicode{multiopenright}{22CC}
\pdfglyphtounicode{multiply}{00D7}
\pdfglyphtounicode{mumsquare}{339B}
\pdfglyphtounicode{munahhebrew}{05A3}
\pdfglyphtounicode{munahlefthebrew}{05A3}
\pdfglyphtounicode{musicalnote}{266A}
\pdfglyphtounicode{musicalnotedbl}{266B}
\pdfglyphtounicode{musicflatsign}{266D}
\pdfglyphtounicode{musicsharpsign}{266F}
\pdfglyphtounicode{mussquare}{33B2}
\pdfglyphtounicode{muvsquare}{33B6}
\pdfglyphtounicode{muwsquare}{33BC}
\pdfglyphtounicode{mvmegasquare}{33B9}
\pdfglyphtounicode{mvsquare}{33B7}
\pdfglyphtounicode{mwmegasquare}{33BF}
\pdfglyphtounicode{mwsquare}{33BD}
\pdfglyphtounicode{n}{006E}
\pdfglyphtounicode{nabengali}{09A8}
\pdfglyphtounicode{nabla}{2207}
\pdfglyphtounicode{nacute}{0144}
\pdfglyphtounicode{nadeva}{0928}
\pdfglyphtounicode{nagujarati}{0AA8}
\pdfglyphtounicode{nagurmukhi}{0A28}
\pdfglyphtounicode{nahiragana}{306A}
\pdfglyphtounicode{nakatakana}{30CA}
\pdfglyphtounicode{nakatakanahalfwidth}{FF85}
\pdfglyphtounicode{nand}{22BC}
\pdfglyphtounicode{napostrophe}{0149}
\pdfglyphtounicode{nasquare}{3381}
\pdfglyphtounicode{natural}{266E}
\pdfglyphtounicode{nbopomofo}{310B}
\pdfglyphtounicode{nbspace}{00A0}
\pdfglyphtounicode{ncaron}{0148}
\pdfglyphtounicode{ncedilla}{0146}
\pdfglyphtounicode{ncircle}{24DD}
\pdfglyphtounicode{ncircumflexbelow}{1E4B}
\pdfglyphtounicode{ncommaaccent}{0146}
\pdfglyphtounicode{ndotaccent}{1E45}
\pdfglyphtounicode{ndotbelow}{1E47}
\pdfglyphtounicode{negationslash}{0338}
\pdfglyphtounicode{nehiragana}{306D}
\pdfglyphtounicode{nekatakana}{30CD}
\pdfglyphtounicode{nekatakanahalfwidth}{FF88}
\pdfglyphtounicode{newsheqelsign}{20AA}
\pdfglyphtounicode{nfsquare}{338B}
\pdfglyphtounicode{ng}{014B}
\pdfglyphtounicode{ngabengali}{0999}
\pdfglyphtounicode{ngadeva}{0919}
\pdfglyphtounicode{ngagujarati}{0A99}
\pdfglyphtounicode{ngagurmukhi}{0A19}
\pdfglyphtounicode{ngonguthai}{0E07}
\pdfglyphtounicode{nhiragana}{3093}
\pdfglyphtounicode{nhookleft}{0272}
\pdfglyphtounicode{nhookretroflex}{0273}
\pdfglyphtounicode{nieunacirclekorean}{326F}
\pdfglyphtounicode{nieunaparenkorean}{320F}
\pdfglyphtounicode{nieuncieuckorean}{3135}
\pdfglyphtounicode{nieuncirclekorean}{3261}
\pdfglyphtounicode{nieunhieuhkorean}{3136}
\pdfglyphtounicode{nieunkorean}{3134}
\pdfglyphtounicode{nieunpansioskorean}{3168}
\pdfglyphtounicode{nieunparenkorean}{3201}
\pdfglyphtounicode{nieunsioskorean}{3167}
\pdfglyphtounicode{nieuntikeutkorean}{3166}
\pdfglyphtounicode{nihiragana}{306B}
\pdfglyphtounicode{nikatakana}{30CB}
\pdfglyphtounicode{nikatakanahalfwidth}{FF86}
\pdfglyphtounicode{nikhahitleftthai}{F899}
\pdfglyphtounicode{nikhahitthai}{0E4D}
\pdfglyphtounicode{nine}{0039}
\pdfglyphtounicode{ninearabic}{0669}
\pdfglyphtounicode{ninebengali}{09EF}
\pdfglyphtounicode{ninecircle}{2468}
\pdfglyphtounicode{ninecircleinversesansserif}{2792}
\pdfglyphtounicode{ninedeva}{096F}
\pdfglyphtounicode{ninegujarati}{0AEF}
\pdfglyphtounicode{ninegurmukhi}{0A6F}
\pdfglyphtounicode{ninehackarabic}{0669}
\pdfglyphtounicode{ninehangzhou}{3029}
\pdfglyphtounicode{nineideographicparen}{3228}
\pdfglyphtounicode{nineinferior}{2089}
\pdfglyphtounicode{ninemonospace}{FF19}
\pdfglyphtounicode{nineoldstyle}{0039}
\pdfglyphtounicode{nineparen}{247C}
\pdfglyphtounicode{nineperiod}{2490}
\pdfglyphtounicode{ninepersian}{06F9}
\pdfglyphtounicode{nineroman}{2178}
\pdfglyphtounicode{ninesuperior}{2079}
\pdfglyphtounicode{nineteencircle}{2472}
\pdfglyphtounicode{nineteenparen}{2486}
\pdfglyphtounicode{nineteenperiod}{249A}
\pdfglyphtounicode{ninethai}{0E59}
\pdfglyphtounicode{nj}{01CC}
\pdfglyphtounicode{njecyrillic}{045A}
\pdfglyphtounicode{nkatakana}{30F3}
\pdfglyphtounicode{nkatakanahalfwidth}{FF9D}
\pdfglyphtounicode{nlegrightlong}{019E}
\pdfglyphtounicode{nlinebelow}{1E49}
\pdfglyphtounicode{nmonospace}{FF4E}
\pdfglyphtounicode{nmsquare}{339A}
\pdfglyphtounicode{nnabengali}{09A3}
\pdfglyphtounicode{nnadeva}{0923}
\pdfglyphtounicode{nnagujarati}{0AA3}
\pdfglyphtounicode{nnagurmukhi}{0A23}
\pdfglyphtounicode{nnnadeva}{0929}
\pdfglyphtounicode{nohiragana}{306E}
\pdfglyphtounicode{nokatakana}{30CE}
\pdfglyphtounicode{nokatakanahalfwidth}{FF89}
\pdfglyphtounicode{nonbreakingspace}{00A0}
\pdfglyphtounicode{nonenthai}{0E13}
\pdfglyphtounicode{nonuthai}{0E19}
\pdfglyphtounicode{noonarabic}{0646}
\pdfglyphtounicode{noonfinalarabic}{FEE6}
\pdfglyphtounicode{noonghunnaarabic}{06BA}
\pdfglyphtounicode{noonghunnafinalarabic}{FB9F}
\pdfglyphtounicode{noonhehinitialarabic}{FEE7 FEEC}
\pdfglyphtounicode{nooninitialarabic}{FEE7}
\pdfglyphtounicode{noonjeeminitialarabic}{FCD2}
\pdfglyphtounicode{noonjeemisolatedarabic}{FC4B}
\pdfglyphtounicode{noonmedialarabic}{FEE8}
\pdfglyphtounicode{noonmeeminitialarabic}{FCD5}
\pdfglyphtounicode{noonmeemisolatedarabic}{FC4E}
\pdfglyphtounicode{noonnoonfinalarabic}{FC8D}
\pdfglyphtounicode{notapproxequal}{2247}
\pdfglyphtounicode{notarrowboth}{21AE}
\pdfglyphtounicode{notarrowleft}{219A}
\pdfglyphtounicode{notarrowright}{219B}
\pdfglyphtounicode{notbar}{2224}
\pdfglyphtounicode{notcontains}{220C}
\pdfglyphtounicode{notdblarrowboth}{21CE}
\pdfglyphtounicode{notdblarrowleft}{21CD}
\pdfglyphtounicode{notdblarrowright}{21CF}
\pdfglyphtounicode{notelement}{2209}
\pdfglyphtounicode{notelementof}{2209}
\pdfglyphtounicode{notequal}{2260}
\pdfglyphtounicode{notexistential}{2204}
\pdfglyphtounicode{notfollows}{2281}
\pdfglyphtounicode{notfollowsoreql}{2AB0 0338}
\pdfglyphtounicode{notforces}{22AE}
\pdfglyphtounicode{notforcesextra}{22AF}
\pdfglyphtounicode{notgreater}{226F}
\pdfglyphtounicode{notgreaterdblequal}{2267 0338}
\pdfglyphtounicode{notgreaterequal}{2271}
\pdfglyphtounicode{notgreaternorequal}{2271}
\pdfglyphtounicode{notgreaternorless}{2279}
\pdfglyphtounicode{notgreaterorslnteql}{2A7E 0338}
\pdfglyphtounicode{notidentical}{2262}
\pdfglyphtounicode{notless}{226E}
\pdfglyphtounicode{notlessdblequal}{2266 0338}
\pdfglyphtounicode{notlessequal}{2270}
\pdfglyphtounicode{notlessnorequal}{2270}
\pdfglyphtounicode{notlessorslnteql}{2A7D 0338}
\pdfglyphtounicode{notparallel}{2226}
\pdfglyphtounicode{notprecedes}{2280}
\pdfglyphtounicode{notprecedesoreql}{2AAF 0338}
\pdfglyphtounicode{notsatisfies}{22AD}
\pdfglyphtounicode{notsimilar}{2241}
\pdfglyphtounicode{notsubset}{2284}
\pdfglyphtounicode{notsubseteql}{2288}
\pdfglyphtounicode{notsubsetordbleql}{2AC5 0338}
\pdfglyphtounicode{notsubsetoreql}{228A}
\pdfglyphtounicode{notsucceeds}{2281}
\pdfglyphtounicode{notsuperset}{2285}
\pdfglyphtounicode{notsuperseteql}{2289}
\pdfglyphtounicode{notsupersetordbleql}{2AC6 0338}
\pdfglyphtounicode{notsupersetoreql}{228B}
\pdfglyphtounicode{nottriangeqlleft}{22EC}
\pdfglyphtounicode{nottriangeqlright}{22ED}
\pdfglyphtounicode{nottriangleleft}{22EA}
\pdfglyphtounicode{nottriangleright}{22EB}
\pdfglyphtounicode{notturnstile}{22AC}
\pdfglyphtounicode{nowarmenian}{0576}
\pdfglyphtounicode{nparen}{24A9}
\pdfglyphtounicode{nssquare}{33B1}
\pdfglyphtounicode{nsuperior}{207F}
\pdfglyphtounicode{ntilde}{00F1}
\pdfglyphtounicode{nu}{03BD}
\pdfglyphtounicode{nuhiragana}{306C}
\pdfglyphtounicode{nukatakana}{30CC}
\pdfglyphtounicode{nukatakanahalfwidth}{FF87}
\pdfglyphtounicode{nuktabengali}{09BC}
\pdfglyphtounicode{nuktadeva}{093C}
\pdfglyphtounicode{nuktagujarati}{0ABC}
\pdfglyphtounicode{nuktagurmukhi}{0A3C}
\pdfglyphtounicode{numbersign}{0023}
\pdfglyphtounicode{numbersignmonospace}{FF03}
\pdfglyphtounicode{numbersignsmall}{FE5F}
\pdfglyphtounicode{numeralsigngreek}{0374}
\pdfglyphtounicode{numeralsignlowergreek}{0375}
\pdfglyphtounicode{numero}{2116}
\pdfglyphtounicode{nun}{05E0}
\pdfglyphtounicode{nundagesh}{FB40}
\pdfglyphtounicode{nundageshhebrew}{FB40}
\pdfglyphtounicode{nunhebrew}{05E0}
\pdfglyphtounicode{nvsquare}{33B5}
\pdfglyphtounicode{nwsquare}{33BB}
\pdfglyphtounicode{nyabengali}{099E}
\pdfglyphtounicode{nyadeva}{091E}
\pdfglyphtounicode{nyagujarati}{0A9E}
\pdfglyphtounicode{nyagurmukhi}{0A1E}
\pdfglyphtounicode{o}{006F}
\pdfglyphtounicode{oacute}{00F3}
\pdfglyphtounicode{oangthai}{0E2D}
\pdfglyphtounicode{obarred}{0275}
\pdfglyphtounicode{obarredcyrillic}{04E9}
\pdfglyphtounicode{obarreddieresiscyrillic}{04EB}
\pdfglyphtounicode{obengali}{0993}
\pdfglyphtounicode{obopomofo}{311B}
\pdfglyphtounicode{obreve}{014F}
\pdfglyphtounicode{ocandradeva}{0911}
\pdfglyphtounicode{ocandragujarati}{0A91}
\pdfglyphtounicode{ocandravowelsigndeva}{0949}
\pdfglyphtounicode{ocandravowelsigngujarati}{0AC9}
\pdfglyphtounicode{ocaron}{01D2}
\pdfglyphtounicode{ocircle}{24DE}
\pdfglyphtounicode{ocircumflex}{00F4}
\pdfglyphtounicode{ocircumflexacute}{1ED1}
\pdfglyphtounicode{ocircumflexdotbelow}{1ED9}
\pdfglyphtounicode{ocircumflexgrave}{1ED3}
\pdfglyphtounicode{ocircumflexhookabove}{1ED5}
\pdfglyphtounicode{ocircumflextilde}{1ED7}
\pdfglyphtounicode{ocyrillic}{043E}
\pdfglyphtounicode{odblacute}{0151}
\pdfglyphtounicode{odblgrave}{020D}
\pdfglyphtounicode{odeva}{0913}
\pdfglyphtounicode{odieresis}{00F6}
\pdfglyphtounicode{odieresiscyrillic}{04E7}
\pdfglyphtounicode{odotbelow}{1ECD}
\pdfglyphtounicode{oe}{0153}
\pdfglyphtounicode{oekorean}{315A}
\pdfglyphtounicode{ogonek}{02DB}
\pdfglyphtounicode{ogonekcmb}{0328}
\pdfglyphtounicode{ograve}{00F2}
\pdfglyphtounicode{ogujarati}{0A93}
\pdfglyphtounicode{oharmenian}{0585}
\pdfglyphtounicode{ohiragana}{304A}
\pdfglyphtounicode{ohookabove}{1ECF}
\pdfglyphtounicode{ohorn}{01A1}
\pdfglyphtounicode{ohornacute}{1EDB}
\pdfglyphtounicode{ohorndotbelow}{1EE3}
\pdfglyphtounicode{ohorngrave}{1EDD}
\pdfglyphtounicode{ohornhookabove}{1EDF}
\pdfglyphtounicode{ohorntilde}{1EE1}
\pdfglyphtounicode{ohungarumlaut}{0151}
\pdfglyphtounicode{oi}{01A3}
\pdfglyphtounicode{oinvertedbreve}{020F}
\pdfglyphtounicode{okatakana}{30AA}
\pdfglyphtounicode{okatakanahalfwidth}{FF75}
\pdfglyphtounicode{okorean}{3157}
\pdfglyphtounicode{olehebrew}{05AB}
\pdfglyphtounicode{omacron}{014D}
\pdfglyphtounicode{omacronacute}{1E53}
\pdfglyphtounicode{omacrongrave}{1E51}
\pdfglyphtounicode{omdeva}{0950}
\pdfglyphtounicode{omega}{03C9}
\pdfglyphtounicode{omega1}{03D6}
\pdfglyphtounicode{omegacyrillic}{0461}
\pdfglyphtounicode{omegalatinclosed}{0277}
\pdfglyphtounicode{omegaroundcyrillic}{047B}
\pdfglyphtounicode{omegatitlocyrillic}{047D}
\pdfglyphtounicode{omegatonos}{03CE}
\pdfglyphtounicode{omgujarati}{0AD0}
\pdfglyphtounicode{omicron}{03BF}
\pdfglyphtounicode{omicrontonos}{03CC}
\pdfglyphtounicode{omonospace}{FF4F}
\pdfglyphtounicode{one}{0031}
\pdfglyphtounicode{onearabic}{0661}
\pdfglyphtounicode{onebengali}{09E7}
\pdfglyphtounicode{onecircle}{2460}
\pdfglyphtounicode{onecircleinversesansserif}{278A}
\pdfglyphtounicode{onedeva}{0967}
\pdfglyphtounicode{onedotenleader}{2024}
\pdfglyphtounicode{oneeighth}{215B}
\pdfglyphtounicode{onefitted}{0031}
\pdfglyphtounicode{onegujarati}{0AE7}
\pdfglyphtounicode{onegurmukhi}{0A67}
\pdfglyphtounicode{onehackarabic}{0661}
\pdfglyphtounicode{onehalf}{00BD}
\pdfglyphtounicode{onehangzhou}{3021}
\pdfglyphtounicode{oneideographicparen}{3220}
\pdfglyphtounicode{oneinferior}{2081}
\pdfglyphtounicode{onemonospace}{FF11}
\pdfglyphtounicode{onenumeratorbengali}{09F4}
\pdfglyphtounicode{oneoldstyle}{0031}
\pdfglyphtounicode{oneparen}{2474}
\pdfglyphtounicode{oneperiod}{2488}
\pdfglyphtounicode{onepersian}{06F1}
\pdfglyphtounicode{onequarter}{00BC}
\pdfglyphtounicode{oneroman}{2170}
\pdfglyphtounicode{onesuperior}{00B9}
\pdfglyphtounicode{onethai}{0E51}
\pdfglyphtounicode{onethird}{2153}
\pdfglyphtounicode{oogonek}{01EB}
\pdfglyphtounicode{oogonekmacron}{01ED}
\pdfglyphtounicode{oogurmukhi}{0A13}
\pdfglyphtounicode{oomatragurmukhi}{0A4B}
\pdfglyphtounicode{oopen}{0254}
\pdfglyphtounicode{oparen}{24AA}
\pdfglyphtounicode{openbullet}{25E6}
\pdfglyphtounicode{option}{2325}
\pdfglyphtounicode{ordfeminine}{00AA}
\pdfglyphtounicode{ordmasculine}{00BA}
\pdfglyphtounicode{orthogonal}{221F}
\pdfglyphtounicode{orunderscore}{22BB}
\pdfglyphtounicode{oshortdeva}{0912}
\pdfglyphtounicode{oshortvowelsigndeva}{094A}
\pdfglyphtounicode{oslash}{00F8}
\pdfglyphtounicode{oslashacute}{01FF}
\pdfglyphtounicode{osmallhiragana}{3049}
\pdfglyphtounicode{osmallkatakana}{30A9}
\pdfglyphtounicode{osmallkatakanahalfwidth}{FF6B}
\pdfglyphtounicode{ostrokeacute}{01FF}
\pdfglyphtounicode{osuperior}{006F}
\pdfglyphtounicode{otcyrillic}{047F}
\pdfglyphtounicode{otilde}{00F5}
\pdfglyphtounicode{otildeacute}{1E4D}
\pdfglyphtounicode{otildedieresis}{1E4F}
\pdfglyphtounicode{oubopomofo}{3121}
\pdfglyphtounicode{overline}{203E}
\pdfglyphtounicode{overlinecenterline}{FE4A}
\pdfglyphtounicode{overlinecmb}{0305}
\pdfglyphtounicode{overlinedashed}{FE49}
\pdfglyphtounicode{overlinedblwavy}{FE4C}
\pdfglyphtounicode{overlinewavy}{FE4B}
\pdfglyphtounicode{overscore}{00AF}
\pdfglyphtounicode{ovowelsignbengali}{09CB}
\pdfglyphtounicode{ovowelsigndeva}{094B}
\pdfglyphtounicode{ovowelsigngujarati}{0ACB}
\pdfglyphtounicode{owner}{220B}
\pdfglyphtounicode{p}{0070}
\pdfglyphtounicode{paampssquare}{3380}
\pdfglyphtounicode{paasentosquare}{332B}
\pdfglyphtounicode{pabengali}{09AA}
\pdfglyphtounicode{pacute}{1E55}
\pdfglyphtounicode{padeva}{092A}
\pdfglyphtounicode{pagedown}{21DF}
\pdfglyphtounicode{pageup}{21DE}
\pdfglyphtounicode{pagujarati}{0AAA}
\pdfglyphtounicode{pagurmukhi}{0A2A}
\pdfglyphtounicode{pahiragana}{3071}
\pdfglyphtounicode{paiyannoithai}{0E2F}
\pdfglyphtounicode{pakatakana}{30D1}
\pdfglyphtounicode{palatalizationcyrilliccmb}{0484}
\pdfglyphtounicode{palochkacyrillic}{04C0}
\pdfglyphtounicode{pansioskorean}{317F}
\pdfglyphtounicode{paragraph}{00B6}
\pdfglyphtounicode{parallel}{2225}
\pdfglyphtounicode{parenleft}{0028}
\pdfglyphtounicode{parenleftaltonearabic}{FD3E}
\pdfglyphtounicode{parenleftbt}{F8ED}
\pdfglyphtounicode{parenleftex}{F8EC}
\pdfglyphtounicode{parenleftinferior}{208D}
\pdfglyphtounicode{parenleftmonospace}{FF08}
\pdfglyphtounicode{parenleftsmall}{FE59}
\pdfglyphtounicode{parenleftsuperior}{207D}
\pdfglyphtounicode{parenlefttp}{F8EB}
\pdfglyphtounicode{parenleftvertical}{FE35}
\pdfglyphtounicode{parenright}{0029}
\pdfglyphtounicode{parenrightaltonearabic}{FD3F}
\pdfglyphtounicode{parenrightbt}{F8F8}
\pdfglyphtounicode{parenrightex}{F8F7}
\pdfglyphtounicode{parenrightinferior}{208E}
\pdfglyphtounicode{parenrightmonospace}{FF09}
\pdfglyphtounicode{parenrightsmall}{FE5A}
\pdfglyphtounicode{parenrightsuperior}{207E}
\pdfglyphtounicode{parenrighttp}{F8F6}
\pdfglyphtounicode{parenrightvertical}{FE36}
\pdfglyphtounicode{partialdiff}{2202}
\pdfglyphtounicode{paseqhebrew}{05C0}
\pdfglyphtounicode{pashtahebrew}{0599}
\pdfglyphtounicode{pasquare}{33A9}
\pdfglyphtounicode{patah}{05B7}
\pdfglyphtounicode{patah11}{05B7}
\pdfglyphtounicode{patah1d}{05B7}
\pdfglyphtounicode{patah2a}{05B7}
\pdfglyphtounicode{patahhebrew}{05B7}
\pdfglyphtounicode{patahnarrowhebrew}{05B7}
\pdfglyphtounicode{patahquarterhebrew}{05B7}
\pdfglyphtounicode{patahwidehebrew}{05B7}
\pdfglyphtounicode{pazerhebrew}{05A1}
\pdfglyphtounicode{pbopomofo}{3106}
\pdfglyphtounicode{pcircle}{24DF}
\pdfglyphtounicode{pdotaccent}{1E57}
\pdfglyphtounicode{pe}{05E4}
\pdfglyphtounicode{pecyrillic}{043F}
\pdfglyphtounicode{pedagesh}{FB44}
\pdfglyphtounicode{pedageshhebrew}{FB44}
\pdfglyphtounicode{peezisquare}{333B}
\pdfglyphtounicode{pefinaldageshhebrew}{FB43}
\pdfglyphtounicode{peharabic}{067E}
\pdfglyphtounicode{peharmenian}{057A}
\pdfglyphtounicode{pehebrew}{05E4}
\pdfglyphtounicode{pehfinalarabic}{FB57}
\pdfglyphtounicode{pehinitialarabic}{FB58}
\pdfglyphtounicode{pehiragana}{307A}
\pdfglyphtounicode{pehmedialarabic}{FB59}
\pdfglyphtounicode{pekatakana}{30DA}
\pdfglyphtounicode{pemiddlehookcyrillic}{04A7}
\pdfglyphtounicode{perafehebrew}{FB4E}
\pdfglyphtounicode{percent}{0025}
\pdfglyphtounicode{percentarabic}{066A}
\pdfglyphtounicode{percentmonospace}{FF05}
\pdfglyphtounicode{percentsmall}{FE6A}
\pdfglyphtounicode{period}{002E}
\pdfglyphtounicode{periodarmenian}{0589}
\pdfglyphtounicode{periodcentered}{00B7}
\pdfglyphtounicode{periodhalfwidth}{FF61}
\pdfglyphtounicode{periodinferior}{002E}
\pdfglyphtounicode{periodmonospace}{FF0E}
\pdfglyphtounicode{periodsmall}{FE52}
\pdfglyphtounicode{periodsuperior}{002E}
\pdfglyphtounicode{perispomenigreekcmb}{0342}
\pdfglyphtounicode{perpcorrespond}{2A5E}
\pdfglyphtounicode{perpendicular}{22A5}
\pdfglyphtounicode{pertenthousand}{2031}
\pdfglyphtounicode{perthousand}{2030}
\pdfglyphtounicode{peseta}{20A7}
\pdfglyphtounicode{pfsquare}{338A}
\pdfglyphtounicode{phabengali}{09AB}
\pdfglyphtounicode{phadeva}{092B}
\pdfglyphtounicode{phagujarati}{0AAB}
\pdfglyphtounicode{phagurmukhi}{0A2B}
\pdfglyphtounicode{phi}{03C6}
\pdfglyphtounicode{phi1}{03D5}
\pdfglyphtounicode{phieuphacirclekorean}{327A}
\pdfglyphtounicode{phieuphaparenkorean}{321A}
\pdfglyphtounicode{phieuphcirclekorean}{326C}
\pdfglyphtounicode{phieuphkorean}{314D}
\pdfglyphtounicode{phieuphparenkorean}{320C}
\pdfglyphtounicode{philatin}{0278}
\pdfglyphtounicode{phinthuthai}{0E3A}
\pdfglyphtounicode{phisymbolgreek}{03D5}
\pdfglyphtounicode{phook}{01A5}
\pdfglyphtounicode{phophanthai}{0E1E}
\pdfglyphtounicode{phophungthai}{0E1C}
\pdfglyphtounicode{phosamphaothai}{0E20}
\pdfglyphtounicode{pi}{03C0}
\pdfglyphtounicode{pi1}{03D6}
\pdfglyphtounicode{pieupacirclekorean}{3273}
\pdfglyphtounicode{pieupaparenkorean}{3213}
\pdfglyphtounicode{pieupcieuckorean}{3176}
\pdfglyphtounicode{pieupcirclekorean}{3265}
\pdfglyphtounicode{pieupkiyeokkorean}{3172}
\pdfglyphtounicode{pieupkorean}{3142}
\pdfglyphtounicode{pieupparenkorean}{3205}
\pdfglyphtounicode{pieupsioskiyeokkorean}{3174}
\pdfglyphtounicode{pieupsioskorean}{3144}
\pdfglyphtounicode{pieupsiostikeutkorean}{3175}
\pdfglyphtounicode{pieupthieuthkorean}{3177}
\pdfglyphtounicode{pieuptikeutkorean}{3173}
\pdfglyphtounicode{pihiragana}{3074}
\pdfglyphtounicode{pikatakana}{30D4}
\pdfglyphtounicode{pisymbolgreek}{03D6}
\pdfglyphtounicode{piwrarmenian}{0583}
\pdfglyphtounicode{planckover2pi}{210F}
\pdfglyphtounicode{planckover2pi1}{210F}
\pdfglyphtounicode{plus}{002B}
\pdfglyphtounicode{plusbelowcmb}{031F}
\pdfglyphtounicode{pluscircle}{2295}
\pdfglyphtounicode{plusminus}{00B1}
\pdfglyphtounicode{plusmod}{02D6}
\pdfglyphtounicode{plusmonospace}{FF0B}
\pdfglyphtounicode{plussmall}{FE62}
\pdfglyphtounicode{plussuperior}{207A}
\pdfglyphtounicode{pmonospace}{FF50}
\pdfglyphtounicode{pmsquare}{33D8}
\pdfglyphtounicode{pohiragana}{307D}
\pdfglyphtounicode{pointingindexdownwhite}{261F}
\pdfglyphtounicode{pointingindexleftwhite}{261C}
\pdfglyphtounicode{pointingindexrightwhite}{261E}
\pdfglyphtounicode{pointingindexupwhite}{261D}
\pdfglyphtounicode{pokatakana}{30DD}
\pdfglyphtounicode{poplathai}{0E1B}
\pdfglyphtounicode{postalmark}{3012}
\pdfglyphtounicode{postalmarkface}{3020}
\pdfglyphtounicode{pparen}{24AB}
\pdfglyphtounicode{precedenotdbleqv}{2AB9}
\pdfglyphtounicode{precedenotslnteql}{2AB5}
\pdfglyphtounicode{precedeornoteqvlnt}{22E8}
\pdfglyphtounicode{precedes}{227A}
\pdfglyphtounicode{precedesequal}{2AAF}
\pdfglyphtounicode{precedesorcurly}{227C}
\pdfglyphtounicode{precedesorequal}{227E}
\pdfglyphtounicode{prescription}{211E}
\pdfglyphtounicode{prime}{2032}
\pdfglyphtounicode{primemod}{02B9}
\pdfglyphtounicode{primereverse}{2035}
\pdfglyphtounicode{primereversed}{2035}
\pdfglyphtounicode{product}{220F}
\pdfglyphtounicode{projective}{2305}
\pdfglyphtounicode{prolongedkana}{30FC}
\pdfglyphtounicode{propellor}{2318}
\pdfglyphtounicode{propersubset}{2282}
\pdfglyphtounicode{propersuperset}{2283}
\pdfglyphtounicode{proportion}{2237}
\pdfglyphtounicode{proportional}{221D}
\pdfglyphtounicode{psi}{03C8}
\pdfglyphtounicode{psicyrillic}{0471}
\pdfglyphtounicode{psilipneumatacyrilliccmb}{0486}
\pdfglyphtounicode{pssquare}{33B0}
\pdfglyphtounicode{puhiragana}{3077}
\pdfglyphtounicode{pukatakana}{30D7}
\pdfglyphtounicode{punctdash}{2014}
\pdfglyphtounicode{pvsquare}{33B4}
\pdfglyphtounicode{pwsquare}{33BA}
\pdfglyphtounicode{q}{0071}
\pdfglyphtounicode{qadeva}{0958}
\pdfglyphtounicode{qadmahebrew}{05A8}
\pdfglyphtounicode{qafarabic}{0642}
\pdfglyphtounicode{qaffinalarabic}{FED6}
\pdfglyphtounicode{qafinitialarabic}{FED7}
\pdfglyphtounicode{qafmedialarabic}{FED8}
\pdfglyphtounicode{qamats}{05B8}
\pdfglyphtounicode{qamats10}{05B8}
\pdfglyphtounicode{qamats1a}{05B8}
\pdfglyphtounicode{qamats1c}{05B8}
\pdfglyphtounicode{qamats27}{05B8}
\pdfglyphtounicode{qamats29}{05B8}
\pdfglyphtounicode{qamats33}{05B8}
\pdfglyphtounicode{qamatsde}{05B8}
\pdfglyphtounicode{qamatshebrew}{05B8}
\pdfglyphtounicode{qamatsnarrowhebrew}{05B8}
\pdfglyphtounicode{qamatsqatanhebrew}{05B8}
\pdfglyphtounicode{qamatsqatannarrowhebrew}{05B8}
\pdfglyphtounicode{qamatsqatanquarterhebrew}{05B8}
\pdfglyphtounicode{qamatsqatanwidehebrew}{05B8}
\pdfglyphtounicode{qamatsquarterhebrew}{05B8}
\pdfglyphtounicode{qamatswidehebrew}{05B8}
\pdfglyphtounicode{qarneyparahebrew}{059F}
\pdfglyphtounicode{qbopomofo}{3111}
\pdfglyphtounicode{qcircle}{24E0}
\pdfglyphtounicode{qhook}{02A0}
\pdfglyphtounicode{qmonospace}{FF51}
\pdfglyphtounicode{qof}{05E7}
\pdfglyphtounicode{qofdagesh}{FB47}
\pdfglyphtounicode{qofdageshhebrew}{FB47}
\pdfglyphtounicode{qofhatafpatah}{05E7 05B2}
\pdfglyphtounicode{qofhatafpatahhebrew}{05E7 05B2}
\pdfglyphtounicode{qofhatafsegol}{05E7 05B1}
\pdfglyphtounicode{qofhatafsegolhebrew}{05E7 05B1}
\pdfglyphtounicode{qofhebrew}{05E7}
\pdfglyphtounicode{qofhiriq}{05E7 05B4}
\pdfglyphtounicode{qofhiriqhebrew}{05E7 05B4}
\pdfglyphtounicode{qofholam}{05E7 05B9}
\pdfglyphtounicode{qofholamhebrew}{05E7 05B9}
\pdfglyphtounicode{qofpatah}{05E7 05B7}
\pdfglyphtounicode{qofpatahhebrew}{05E7 05B7}
\pdfglyphtounicode{qofqamats}{05E7 05B8}
\pdfglyphtounicode{qofqamatshebrew}{05E7 05B8}
\pdfglyphtounicode{qofqubuts}{05E7 05BB}
\pdfglyphtounicode{qofqubutshebrew}{05E7 05BB}
\pdfglyphtounicode{qofsegol}{05E7 05B6}
\pdfglyphtounicode{qofsegolhebrew}{05E7 05B6}
\pdfglyphtounicode{qofsheva}{05E7 05B0}
\pdfglyphtounicode{qofshevahebrew}{05E7 05B0}
\pdfglyphtounicode{qoftsere}{05E7 05B5}
\pdfglyphtounicode{qoftserehebrew}{05E7 05B5}
\pdfglyphtounicode{qparen}{24AC}
\pdfglyphtounicode{quarternote}{2669}
\pdfglyphtounicode{qubuts}{05BB}
\pdfglyphtounicode{qubuts18}{05BB}
\pdfglyphtounicode{qubuts25}{05BB}
\pdfglyphtounicode{qubuts31}{05BB}
\pdfglyphtounicode{qubutshebrew}{05BB}
\pdfglyphtounicode{qubutsnarrowhebrew}{05BB}
\pdfglyphtounicode{qubutsquarterhebrew}{05BB}
\pdfglyphtounicode{qubutswidehebrew}{05BB}
\pdfglyphtounicode{question}{003F}
\pdfglyphtounicode{questionarabic}{061F}
\pdfglyphtounicode{questionarmenian}{055E}
\pdfglyphtounicode{questiondown}{00BF}
\pdfglyphtounicode{questiondownsmall}{00BF}
\pdfglyphtounicode{questiongreek}{037E}
\pdfglyphtounicode{questionmonospace}{FF1F}
\pdfglyphtounicode{questionsmall}{003F}
\pdfglyphtounicode{quotedbl}{0022}
\pdfglyphtounicode{quotedblbase}{201E}
\pdfglyphtounicode{quotedblleft}{201C}
\pdfglyphtounicode{quotedblmonospace}{FF02}
\pdfglyphtounicode{quotedblprime}{301E}
\pdfglyphtounicode{quotedblprimereversed}{301D}
\pdfglyphtounicode{quotedblright}{201D}
\pdfglyphtounicode{quoteleft}{2018}
\pdfglyphtounicode{quoteleftreversed}{201B}
\pdfglyphtounicode{quotereversed}{201B}
\pdfglyphtounicode{quoteright}{2019}
\pdfglyphtounicode{quoterightn}{0149}
\pdfglyphtounicode{quotesinglbase}{201A}
\pdfglyphtounicode{quotesingle}{0027}
\pdfglyphtounicode{quotesinglemonospace}{FF07}
\pdfglyphtounicode{r}{0072}
\pdfglyphtounicode{raarmenian}{057C}
\pdfglyphtounicode{rabengali}{09B0}
\pdfglyphtounicode{racute}{0155}
\pdfglyphtounicode{radeva}{0930}
\pdfglyphtounicode{radical}{221A}
\pdfglyphtounicode{radicalex}{F8E5}
\pdfglyphtounicode{radoverssquare}{33AE}
\pdfglyphtounicode{radoverssquaredsquare}{33AF}
\pdfglyphtounicode{radsquare}{33AD}
\pdfglyphtounicode{rafe}{05BF}
\pdfglyphtounicode{rafehebrew}{05BF}
\pdfglyphtounicode{ragujarati}{0AB0}
\pdfglyphtounicode{ragurmukhi}{0A30}
\pdfglyphtounicode{rahiragana}{3089}
\pdfglyphtounicode{rakatakana}{30E9}
\pdfglyphtounicode{rakatakanahalfwidth}{FF97}
\pdfglyphtounicode{ralowerdiagonalbengali}{09F1}
\pdfglyphtounicode{ramiddlediagonalbengali}{09F0}
\pdfglyphtounicode{ramshorn}{0264}
\pdfglyphtounicode{rangedash}{2013}
\pdfglyphtounicode{ratio}{2236}
\pdfglyphtounicode{rbopomofo}{3116}
\pdfglyphtounicode{rcaron}{0159}
\pdfglyphtounicode{rcedilla}{0157}
\pdfglyphtounicode{rcircle}{24E1}
\pdfglyphtounicode{rcommaaccent}{0157}
\pdfglyphtounicode{rdblgrave}{0211}
\pdfglyphtounicode{rdotaccent}{1E59}
\pdfglyphtounicode{rdotbelow}{1E5B}
\pdfglyphtounicode{rdotbelowmacron}{1E5D}
\pdfglyphtounicode{referencemark}{203B}
\pdfglyphtounicode{reflexsubset}{2286}
\pdfglyphtounicode{reflexsuperset}{2287}
\pdfglyphtounicode{registered}{00AE}
\pdfglyphtounicode{registersans}{00AE}
\pdfglyphtounicode{registerserif}{00AE}
\pdfglyphtounicode{reharabic}{0631}
\pdfglyphtounicode{reharmenian}{0580}
\pdfglyphtounicode{rehfinalarabic}{FEAE}
\pdfglyphtounicode{rehiragana}{308C}
\pdfglyphtounicode{rehyehaleflamarabic}{0631 FEF3 FE8E 0644}
\pdfglyphtounicode{rekatakana}{30EC}
\pdfglyphtounicode{rekatakanahalfwidth}{FF9A}
\pdfglyphtounicode{resh}{05E8}
\pdfglyphtounicode{reshdageshhebrew}{FB48}
\pdfglyphtounicode{reshhatafpatah}{05E8 05B2}
\pdfglyphtounicode{reshhatafpatahhebrew}{05E8 05B2}
\pdfglyphtounicode{reshhatafsegol}{05E8 05B1}
\pdfglyphtounicode{reshhatafsegolhebrew}{05E8 05B1}
\pdfglyphtounicode{reshhebrew}{05E8}
\pdfglyphtounicode{reshhiriq}{05E8 05B4}
\pdfglyphtounicode{reshhiriqhebrew}{05E8 05B4}
\pdfglyphtounicode{reshholam}{05E8 05B9}
\pdfglyphtounicode{reshholamhebrew}{05E8 05B9}
\pdfglyphtounicode{reshpatah}{05E8 05B7}
\pdfglyphtounicode{reshpatahhebrew}{05E8 05B7}
\pdfglyphtounicode{reshqamats}{05E8 05B8}
\pdfglyphtounicode{reshqamatshebrew}{05E8 05B8}
\pdfglyphtounicode{reshqubuts}{05E8 05BB}
\pdfglyphtounicode{reshqubutshebrew}{05E8 05BB}
\pdfglyphtounicode{reshsegol}{05E8 05B6}
\pdfglyphtounicode{reshsegolhebrew}{05E8 05B6}
\pdfglyphtounicode{reshsheva}{05E8 05B0}
\pdfglyphtounicode{reshshevahebrew}{05E8 05B0}
\pdfglyphtounicode{reshtsere}{05E8 05B5}
\pdfglyphtounicode{reshtserehebrew}{05E8 05B5}
\pdfglyphtounicode{revasymptequal}{22CD}
\pdfglyphtounicode{reversedtilde}{223D}
\pdfglyphtounicode{reviahebrew}{0597}
\pdfglyphtounicode{reviamugrashhebrew}{0597}
\pdfglyphtounicode{revlogicalnot}{2310}
\pdfglyphtounicode{revsimilar}{223D}
\pdfglyphtounicode{rfishhook}{027E}
\pdfglyphtounicode{rfishhookreversed}{027F}
\pdfglyphtounicode{rhabengali}{09DD}
\pdfglyphtounicode{rhadeva}{095D}
\pdfglyphtounicode{rho}{03C1}
\pdfglyphtounicode{rho1}{03F1}
\pdfglyphtounicode{rhook}{027D}
\pdfglyphtounicode{rhookturned}{027B}
\pdfglyphtounicode{rhookturnedsuperior}{02B5}
\pdfglyphtounicode{rhosymbolgreek}{03F1}
\pdfglyphtounicode{rhotichookmod}{02DE}
\pdfglyphtounicode{rieulacirclekorean}{3271}
\pdfglyphtounicode{rieulaparenkorean}{3211}
\pdfglyphtounicode{rieulcirclekorean}{3263}
\pdfglyphtounicode{rieulhieuhkorean}{3140}
\pdfglyphtounicode{rieulkiyeokkorean}{313A}
\pdfglyphtounicode{rieulkiyeoksioskorean}{3169}
\pdfglyphtounicode{rieulkorean}{3139}
\pdfglyphtounicode{rieulmieumkorean}{313B}
\pdfglyphtounicode{rieulpansioskorean}{316C}
\pdfglyphtounicode{rieulparenkorean}{3203}
\pdfglyphtounicode{rieulphieuphkorean}{313F}
\pdfglyphtounicode{rieulpieupkorean}{313C}
\pdfglyphtounicode{rieulpieupsioskorean}{316B}
\pdfglyphtounicode{rieulsioskorean}{313D}
\pdfglyphtounicode{rieulthieuthkorean}{313E}
\pdfglyphtounicode{rieultikeutkorean}{316A}
\pdfglyphtounicode{rieulyeorinhieuhkorean}{316D}
\pdfglyphtounicode{rightangle}{221F}
\pdfglyphtounicode{rightanglene}{231D}
\pdfglyphtounicode{rightanglenw}{231C}
\pdfglyphtounicode{rightanglese}{231F}
\pdfglyphtounicode{rightanglesw}{231E}
\pdfglyphtounicode{righttackbelowcmb}{0319}
\pdfglyphtounicode{righttriangle}{22BF}
\pdfglyphtounicode{rihiragana}{308A}
\pdfglyphtounicode{rikatakana}{30EA}
\pdfglyphtounicode{rikatakanahalfwidth}{FF98}
\pdfglyphtounicode{ring}{02DA}
\pdfglyphtounicode{ringbelowcmb}{0325}
\pdfglyphtounicode{ringcmb}{030A}
\pdfglyphtounicode{ringhalfleft}{02BF}
\pdfglyphtounicode{ringhalfleftarmenian}{0559}
\pdfglyphtounicode{ringhalfleftbelowcmb}{031C}
\pdfglyphtounicode{ringhalfleftcentered}{02D3}
\pdfglyphtounicode{ringhalfright}{02BE}
\pdfglyphtounicode{ringhalfrightbelowcmb}{0339}
\pdfglyphtounicode{ringhalfrightcentered}{02D2}
\pdfglyphtounicode{ringinequal}{2256}
\pdfglyphtounicode{rinvertedbreve}{0213}
\pdfglyphtounicode{rittorusquare}{3351}
\pdfglyphtounicode{rlinebelow}{1E5F}
\pdfglyphtounicode{rlongleg}{027C}
\pdfglyphtounicode{rlonglegturned}{027A}
\pdfglyphtounicode{rmonospace}{FF52}
\pdfglyphtounicode{rohiragana}{308D}
\pdfglyphtounicode{rokatakana}{30ED}
\pdfglyphtounicode{rokatakanahalfwidth}{FF9B}
\pdfglyphtounicode{roruathai}{0E23}
\pdfglyphtounicode{rparen}{24AD}
\pdfglyphtounicode{rrabengali}{09DC}
\pdfglyphtounicode{rradeva}{0931}
\pdfglyphtounicode{rragurmukhi}{0A5C}
\pdfglyphtounicode{rreharabic}{0691}
\pdfglyphtounicode{rrehfinalarabic}{FB8D}
\pdfglyphtounicode{rrvocalicbengali}{09E0}
\pdfglyphtounicode{rrvocalicdeva}{0960}
\pdfglyphtounicode{rrvocalicgujarati}{0AE0}
\pdfglyphtounicode{rrvocalicvowelsignbengali}{09C4}
\pdfglyphtounicode{rrvocalicvowelsigndeva}{0944}
\pdfglyphtounicode{rrvocalicvowelsigngujarati}{0AC4}
\pdfglyphtounicode{rsuperior}{0072}
\pdfglyphtounicode{rtblock}{2590}
\pdfglyphtounicode{rturned}{0279}
\pdfglyphtounicode{rturnedsuperior}{02B4}
\pdfglyphtounicode{ruhiragana}{308B}
\pdfglyphtounicode{rukatakana}{30EB}
\pdfglyphtounicode{rukatakanahalfwidth}{FF99}
\pdfglyphtounicode{rupeemarkbengali}{09F2}
\pdfglyphtounicode{rupeesignbengali}{09F3}
\pdfglyphtounicode{rupiah}{20A8}
\pdfglyphtounicode{ruthai}{0E24}
\pdfglyphtounicode{rvocalicbengali}{098B}
\pdfglyphtounicode{rvocalicdeva}{090B}
\pdfglyphtounicode{rvocalicgujarati}{0A8B}
\pdfglyphtounicode{rvocalicvowelsignbengali}{09C3}
\pdfglyphtounicode{rvocalicvowelsigndeva}{0943}
\pdfglyphtounicode{rvocalicvowelsigngujarati}{0AC3}
\pdfglyphtounicode{s}{0073}
\pdfglyphtounicode{sabengali}{09B8}
\pdfglyphtounicode{sacute}{015B}
\pdfglyphtounicode{sacutedotaccent}{1E65}
\pdfglyphtounicode{sadarabic}{0635}
\pdfglyphtounicode{sadeva}{0938}
\pdfglyphtounicode{sadfinalarabic}{FEBA}
\pdfglyphtounicode{sadinitialarabic}{FEBB}
\pdfglyphtounicode{sadmedialarabic}{FEBC}
\pdfglyphtounicode{sagujarati}{0AB8}
\pdfglyphtounicode{sagurmukhi}{0A38}
\pdfglyphtounicode{sahiragana}{3055}
\pdfglyphtounicode{sakatakana}{30B5}
\pdfglyphtounicode{sakatakanahalfwidth}{FF7B}
\pdfglyphtounicode{sallallahoualayhewasallamarabic}{FDFA}
\pdfglyphtounicode{samekh}{05E1}
\pdfglyphtounicode{samekhdagesh}{FB41}
\pdfglyphtounicode{samekhdageshhebrew}{FB41}
\pdfglyphtounicode{samekhhebrew}{05E1}
\pdfglyphtounicode{saraaathai}{0E32}
\pdfglyphtounicode{saraaethai}{0E41}
\pdfglyphtounicode{saraaimaimalaithai}{0E44}
\pdfglyphtounicode{saraaimaimuanthai}{0E43}
\pdfglyphtounicode{saraamthai}{0E33}
\pdfglyphtounicode{saraathai}{0E30}
\pdfglyphtounicode{saraethai}{0E40}
\pdfglyphtounicode{saraiileftthai}{F886}
\pdfglyphtounicode{saraiithai}{0E35}
\pdfglyphtounicode{saraileftthai}{F885}
\pdfglyphtounicode{saraithai}{0E34}
\pdfglyphtounicode{saraothai}{0E42}
\pdfglyphtounicode{saraueeleftthai}{F888}
\pdfglyphtounicode{saraueethai}{0E37}
\pdfglyphtounicode{saraueleftthai}{F887}
\pdfglyphtounicode{sarauethai}{0E36}
\pdfglyphtounicode{sarauthai}{0E38}
\pdfglyphtounicode{sarauuthai}{0E39}
\pdfglyphtounicode{satisfies}{22A8}
\pdfglyphtounicode{sbopomofo}{3119}
\pdfglyphtounicode{scaron}{0161}
\pdfglyphtounicode{scarondotaccent}{1E67}
\pdfglyphtounicode{scedilla}{015F}
\pdfglyphtounicode{schwa}{0259}
\pdfglyphtounicode{schwacyrillic}{04D9}
\pdfglyphtounicode{schwadieresiscyrillic}{04DB}
\pdfglyphtounicode{schwahook}{025A}
\pdfglyphtounicode{scircle}{24E2}
\pdfglyphtounicode{scircumflex}{015D}
\pdfglyphtounicode{scommaaccent}{0219}
\pdfglyphtounicode{sdotaccent}{1E61}
\pdfglyphtounicode{sdotbelow}{1E63}
\pdfglyphtounicode{sdotbelowdotaccent}{1E69}
\pdfglyphtounicode{seagullbelowcmb}{033C}
\pdfglyphtounicode{second}{2033}
\pdfglyphtounicode{secondtonechinese}{02CA}
\pdfglyphtounicode{section}{00A7}
\pdfglyphtounicode{seenarabic}{0633}
\pdfglyphtounicode{seenfinalarabic}{FEB2}
\pdfglyphtounicode{seeninitialarabic}{FEB3}
\pdfglyphtounicode{seenmedialarabic}{FEB4}
\pdfglyphtounicode{segol}{05B6}
\pdfglyphtounicode{segol13}{05B6}
\pdfglyphtounicode{segol1f}{05B6}
\pdfglyphtounicode{segol2c}{05B6}
\pdfglyphtounicode{segolhebrew}{05B6}
\pdfglyphtounicode{segolnarrowhebrew}{05B6}
\pdfglyphtounicode{segolquarterhebrew}{05B6}
\pdfglyphtounicode{segoltahebrew}{0592}
\pdfglyphtounicode{segolwidehebrew}{05B6}
\pdfglyphtounicode{seharmenian}{057D}
\pdfglyphtounicode{sehiragana}{305B}
\pdfglyphtounicode{sekatakana}{30BB}
\pdfglyphtounicode{sekatakanahalfwidth}{FF7E}
\pdfglyphtounicode{semicolon}{003B}
\pdfglyphtounicode{semicolonarabic}{061B}
\pdfglyphtounicode{semicolonmonospace}{FF1B}
\pdfglyphtounicode{semicolonsmall}{FE54}
\pdfglyphtounicode{semivoicedmarkkana}{309C}
\pdfglyphtounicode{semivoicedmarkkanahalfwidth}{FF9F}
\pdfglyphtounicode{sentisquare}{3322}
\pdfglyphtounicode{sentosquare}{3323}
\pdfglyphtounicode{seven}{0037}
\pdfglyphtounicode{sevenarabic}{0667}
\pdfglyphtounicode{sevenbengali}{09ED}
\pdfglyphtounicode{sevencircle}{2466}
\pdfglyphtounicode{sevencircleinversesansserif}{2790}
\pdfglyphtounicode{sevendeva}{096D}
\pdfglyphtounicode{seveneighths}{215E}
\pdfglyphtounicode{sevengujarati}{0AED}
\pdfglyphtounicode{sevengurmukhi}{0A6D}
\pdfglyphtounicode{sevenhackarabic}{0667}
\pdfglyphtounicode{sevenhangzhou}{3027}
\pdfglyphtounicode{sevenideographicparen}{3226}
\pdfglyphtounicode{seveninferior}{2087}
\pdfglyphtounicode{sevenmonospace}{FF17}
\pdfglyphtounicode{sevenoldstyle}{0037}
\pdfglyphtounicode{sevenparen}{247A}
\pdfglyphtounicode{sevenperiod}{248E}
\pdfglyphtounicode{sevenpersian}{06F7}
\pdfglyphtounicode{sevenroman}{2176}
\pdfglyphtounicode{sevensuperior}{2077}
\pdfglyphtounicode{seventeencircle}{2470}
\pdfglyphtounicode{seventeenparen}{2484}
\pdfglyphtounicode{seventeenperiod}{2498}
\pdfglyphtounicode{seventhai}{0E57}
\pdfglyphtounicode{sfthyphen}{00AD}
\pdfglyphtounicode{shaarmenian}{0577}
\pdfglyphtounicode{shabengali}{09B6}
\pdfglyphtounicode{shacyrillic}{0448}
\pdfglyphtounicode{shaddaarabic}{0651}
\pdfglyphtounicode{shaddadammaarabic}{FC61}
\pdfglyphtounicode{shaddadammatanarabic}{FC5E}
\pdfglyphtounicode{shaddafathaarabic}{FC60}
\pdfglyphtounicode{shaddafathatanarabic}{0651 064B}
\pdfglyphtounicode{shaddakasraarabic}{FC62}
\pdfglyphtounicode{shaddakasratanarabic}{FC5F}
\pdfglyphtounicode{shade}{2592}
\pdfglyphtounicode{shadedark}{2593}
\pdfglyphtounicode{shadelight}{2591}
\pdfglyphtounicode{shademedium}{2592}
\pdfglyphtounicode{shadeva}{0936}
\pdfglyphtounicode{shagujarati}{0AB6}
\pdfglyphtounicode{shagurmukhi}{0A36}
\pdfglyphtounicode{shalshelethebrew}{0593}
\pdfglyphtounicode{sharp}{266F}
\pdfglyphtounicode{shbopomofo}{3115}
\pdfglyphtounicode{shchacyrillic}{0449}
\pdfglyphtounicode{sheenarabic}{0634}
\pdfglyphtounicode{sheenfinalarabic}{FEB6}
\pdfglyphtounicode{sheeninitialarabic}{FEB7}
\pdfglyphtounicode{sheenmedialarabic}{FEB8}
\pdfglyphtounicode{sheicoptic}{03E3}
\pdfglyphtounicode{sheqel}{20AA}
\pdfglyphtounicode{sheqelhebrew}{20AA}
\pdfglyphtounicode{sheva}{05B0}
\pdfglyphtounicode{sheva115}{05B0}
\pdfglyphtounicode{sheva15}{05B0}
\pdfglyphtounicode{sheva22}{05B0}
\pdfglyphtounicode{sheva2e}{05B0}
\pdfglyphtounicode{shevahebrew}{05B0}
\pdfglyphtounicode{shevanarrowhebrew}{05B0}
\pdfglyphtounicode{shevaquarterhebrew}{05B0}
\pdfglyphtounicode{shevawidehebrew}{05B0}
\pdfglyphtounicode{shhacyrillic}{04BB}
\pdfglyphtounicode{shiftleft}{21B0}
\pdfglyphtounicode{shiftright}{21B1}
\pdfglyphtounicode{shimacoptic}{03ED}
\pdfglyphtounicode{shin}{05E9}
\pdfglyphtounicode{shindagesh}{FB49}
\pdfglyphtounicode{shindageshhebrew}{FB49}
\pdfglyphtounicode{shindageshshindot}{FB2C}
\pdfglyphtounicode{shindageshshindothebrew}{FB2C}
\pdfglyphtounicode{shindageshsindot}{FB2D}
\pdfglyphtounicode{shindageshsindothebrew}{FB2D}
\pdfglyphtounicode{shindothebrew}{05C1}
\pdfglyphtounicode{shinhebrew}{05E9}
\pdfglyphtounicode{shinshindot}{FB2A}
\pdfglyphtounicode{shinshindothebrew}{FB2A}
\pdfglyphtounicode{shinsindot}{FB2B}
\pdfglyphtounicode{shinsindothebrew}{FB2B}
\pdfglyphtounicode{shook}{0282}
\pdfglyphtounicode{sigma}{03C3}
\pdfglyphtounicode{sigma1}{03C2}
\pdfglyphtounicode{sigmafinal}{03C2}
\pdfglyphtounicode{sigmalunatesymbolgreek}{03F2}
\pdfglyphtounicode{sihiragana}{3057}
\pdfglyphtounicode{sikatakana}{30B7}
\pdfglyphtounicode{sikatakanahalfwidth}{FF7C}
\pdfglyphtounicode{siluqhebrew}{05BD}
\pdfglyphtounicode{siluqlefthebrew}{05BD}
\pdfglyphtounicode{similar}{223C}
\pdfglyphtounicode{similarequal}{2243}
\pdfglyphtounicode{sindothebrew}{05C2}
\pdfglyphtounicode{siosacirclekorean}{3274}
\pdfglyphtounicode{siosaparenkorean}{3214}
\pdfglyphtounicode{sioscieuckorean}{317E}
\pdfglyphtounicode{sioscirclekorean}{3266}
\pdfglyphtounicode{sioskiyeokkorean}{317A}
\pdfglyphtounicode{sioskorean}{3145}
\pdfglyphtounicode{siosnieunkorean}{317B}
\pdfglyphtounicode{siosparenkorean}{3206}
\pdfglyphtounicode{siospieupkorean}{317D}
\pdfglyphtounicode{siostikeutkorean}{317C}
\pdfglyphtounicode{six}{0036}
\pdfglyphtounicode{sixarabic}{0666}
\pdfglyphtounicode{sixbengali}{09EC}
\pdfglyphtounicode{sixcircle}{2465}
\pdfglyphtounicode{sixcircleinversesansserif}{278F}
\pdfglyphtounicode{sixdeva}{096C}
\pdfglyphtounicode{sixgujarati}{0AEC}
\pdfglyphtounicode{sixgurmukhi}{0A6C}
\pdfglyphtounicode{sixhackarabic}{0666}
\pdfglyphtounicode{sixhangzhou}{3026}
\pdfglyphtounicode{sixideographicparen}{3225}
\pdfglyphtounicode{sixinferior}{2086}
\pdfglyphtounicode{sixmonospace}{FF16}
\pdfglyphtounicode{sixoldstyle}{0036}
\pdfglyphtounicode{sixparen}{2479}
\pdfglyphtounicode{sixperiod}{248D}
\pdfglyphtounicode{sixpersian}{06F6}
\pdfglyphtounicode{sixroman}{2175}
\pdfglyphtounicode{sixsuperior}{2076}
\pdfglyphtounicode{sixteencircle}{246F}
\pdfglyphtounicode{sixteencurrencydenominatorbengali}{09F9}
\pdfglyphtounicode{sixteenparen}{2483}
\pdfglyphtounicode{sixteenperiod}{2497}
\pdfglyphtounicode{sixthai}{0E56}
\pdfglyphtounicode{slash}{002F}
\pdfglyphtounicode{slashmonospace}{FF0F}
\pdfglyphtounicode{slong}{017F}
\pdfglyphtounicode{slongdotaccent}{1E9B}
\pdfglyphtounicode{slurabove}{2322}
\pdfglyphtounicode{slurbelow}{2323}
\pdfglyphtounicode{smile}{2323}
\pdfglyphtounicode{smileface}{263A}
\pdfglyphtounicode{smonospace}{FF53}
\pdfglyphtounicode{sofpasuqhebrew}{05C3}
\pdfglyphtounicode{softhyphen}{00AD}
\pdfglyphtounicode{softsigncyrillic}{044C}
\pdfglyphtounicode{sohiragana}{305D}
\pdfglyphtounicode{sokatakana}{30BD}
\pdfglyphtounicode{sokatakanahalfwidth}{FF7F}
\pdfglyphtounicode{soliduslongoverlaycmb}{0338}
\pdfglyphtounicode{solidusshortoverlaycmb}{0337}
\pdfglyphtounicode{sorusithai}{0E29}
\pdfglyphtounicode{sosalathai}{0E28}
\pdfglyphtounicode{sosothai}{0E0B}
\pdfglyphtounicode{sosuathai}{0E2A}
\pdfglyphtounicode{space}{0020}
\pdfglyphtounicode{spacehackarabic}{0020}
\pdfglyphtounicode{spade}{2660}
\pdfglyphtounicode{spadesuitblack}{2660}
\pdfglyphtounicode{spadesuitwhite}{2664}
\pdfglyphtounicode{sparen}{24AE}
\pdfglyphtounicode{sphericalangle}{2222}
\pdfglyphtounicode{square}{25A1}
\pdfglyphtounicode{squarebelowcmb}{033B}
\pdfglyphtounicode{squarecc}{33C4}
\pdfglyphtounicode{squarecm}{339D}
\pdfglyphtounicode{squarediagonalcrosshatchfill}{25A9}
\pdfglyphtounicode{squaredot}{22A1}
\pdfglyphtounicode{squarehorizontalfill}{25A4}
\pdfglyphtounicode{squareimage}{228F}
\pdfglyphtounicode{squarekg}{338F}
\pdfglyphtounicode{squarekm}{339E}
\pdfglyphtounicode{squarekmcapital}{33CE}
\pdfglyphtounicode{squareln}{33D1}
\pdfglyphtounicode{squarelog}{33D2}
\pdfglyphtounicode{squaremg}{338E}
\pdfglyphtounicode{squaremil}{33D5}
\pdfglyphtounicode{squareminus}{229F}
\pdfglyphtounicode{squaremm}{339C}
\pdfglyphtounicode{squaremsquared}{33A1}
\pdfglyphtounicode{squaremultiply}{22A0}
\pdfglyphtounicode{squareoriginal}{2290}
\pdfglyphtounicode{squareorthogonalcrosshatchfill}{25A6}
\pdfglyphtounicode{squareplus}{229E}
\pdfglyphtounicode{squaresolid}{25A0}
\pdfglyphtounicode{squareupperlefttolowerrightfill}{25A7}
\pdfglyphtounicode{squareupperrighttolowerleftfill}{25A8}
\pdfglyphtounicode{squareverticalfill}{25A5}
\pdfglyphtounicode{squarewhitewithsmallblack}{25A3}
\pdfglyphtounicode{squiggleleftright}{21AD}
\pdfglyphtounicode{squiggleright}{21DD}
\pdfglyphtounicode{srsquare}{33DB}
\pdfglyphtounicode{ssabengali}{09B7}
\pdfglyphtounicode{ssadeva}{0937}
\pdfglyphtounicode{ssagujarati}{0AB7}
\pdfglyphtounicode{ssangcieuckorean}{3149}
\pdfglyphtounicode{ssanghieuhkorean}{3185}
\pdfglyphtounicode{ssangieungkorean}{3180}
\pdfglyphtounicode{ssangkiyeokkorean}{3132}
\pdfglyphtounicode{ssangnieunkorean}{3165}
\pdfglyphtounicode{ssangpieupkorean}{3143}
\pdfglyphtounicode{ssangsioskorean}{3146}
\pdfglyphtounicode{ssangtikeutkorean}{3138}
\pdfglyphtounicode{ssuperior}{0073}
\pdfglyphtounicode{st}{0073 0074}
\pdfglyphtounicode{star}{22C6}
\pdfglyphtounicode{sterling}{00A3}
\pdfglyphtounicode{sterlingmonospace}{FFE1}
\pdfglyphtounicode{strokelongoverlaycmb}{0336}
\pdfglyphtounicode{strokeshortoverlaycmb}{0335}
\pdfglyphtounicode{subset}{2282}
\pdfglyphtounicode{subsetdbl}{22D0}
\pdfglyphtounicode{subsetdblequal}{2AC5}
\pdfglyphtounicode{subsetnoteql}{228A}
\pdfglyphtounicode{subsetnotequal}{228A}
\pdfglyphtounicode{subsetorequal}{2286}
\pdfglyphtounicode{subsetornotdbleql}{2ACB}
\pdfglyphtounicode{subsetsqequal}{2291}
\pdfglyphtounicode{succeeds}{227B}
\pdfglyphtounicode{suchthat}{220B}
\pdfglyphtounicode{suhiragana}{3059}
\pdfglyphtounicode{sukatakana}{30B9}
\pdfglyphtounicode{sukatakanahalfwidth}{FF7D}
\pdfglyphtounicode{sukunarabic}{0652}
\pdfglyphtounicode{summation}{2211}
\pdfglyphtounicode{sun}{263C}
\pdfglyphtounicode{superset}{2283}
\pdfglyphtounicode{supersetdbl}{22D1}
\pdfglyphtounicode{supersetdblequal}{2AC6}
\pdfglyphtounicode{supersetnoteql}{228B}
\pdfglyphtounicode{supersetnotequal}{228B}
\pdfglyphtounicode{supersetorequal}{2287}
\pdfglyphtounicode{supersetornotdbleql}{2ACC}
\pdfglyphtounicode{supersetsqequal}{2292}
\pdfglyphtounicode{svsquare}{33DC}
\pdfglyphtounicode{syouwaerasquare}{337C}
\pdfglyphtounicode{t}{0074}
\pdfglyphtounicode{tabengali}{09A4}
\pdfglyphtounicode{tackdown}{22A4}
\pdfglyphtounicode{tackleft}{22A3}
\pdfglyphtounicode{tadeva}{0924}
\pdfglyphtounicode{tagujarati}{0AA4}
\pdfglyphtounicode{tagurmukhi}{0A24}
\pdfglyphtounicode{taharabic}{0637}
\pdfglyphtounicode{tahfinalarabic}{FEC2}
\pdfglyphtounicode{tahinitialarabic}{FEC3}
\pdfglyphtounicode{tahiragana}{305F}
\pdfglyphtounicode{tahmedialarabic}{FEC4}
\pdfglyphtounicode{taisyouerasquare}{337D}
\pdfglyphtounicode{takatakana}{30BF}
\pdfglyphtounicode{takatakanahalfwidth}{FF80}
\pdfglyphtounicode{tatweelarabic}{0640}
\pdfglyphtounicode{tau}{03C4}
\pdfglyphtounicode{tav}{05EA}
\pdfglyphtounicode{tavdages}{FB4A}
\pdfglyphtounicode{tavdagesh}{FB4A}
\pdfglyphtounicode{tavdageshhebrew}{FB4A}
\pdfglyphtounicode{tavhebrew}{05EA}
\pdfglyphtounicode{tbar}{0167}
\pdfglyphtounicode{tbopomofo}{310A}
\pdfglyphtounicode{tcaron}{0165}
\pdfglyphtounicode{tccurl}{02A8}
\pdfglyphtounicode{tcedilla}{0163}
\pdfglyphtounicode{tcheharabic}{0686}
\pdfglyphtounicode{tchehfinalarabic}{FB7B}
\pdfglyphtounicode{tchehinitialarabic}{FB7C}
\pdfglyphtounicode{tchehmedialarabic}{FB7D}
\pdfglyphtounicode{tchehmeeminitialarabic}{FB7C FEE4}
\pdfglyphtounicode{tcircle}{24E3}
\pdfglyphtounicode{tcircumflexbelow}{1E71}
\pdfglyphtounicode{tcommaaccent}{0163}
\pdfglyphtounicode{tdieresis}{1E97}
\pdfglyphtounicode{tdotaccent}{1E6B}
\pdfglyphtounicode{tdotbelow}{1E6D}
\pdfglyphtounicode{tecyrillic}{0442}
\pdfglyphtounicode{tedescendercyrillic}{04AD}
\pdfglyphtounicode{teharabic}{062A}
\pdfglyphtounicode{tehfinalarabic}{FE96}
\pdfglyphtounicode{tehhahinitialarabic}{FCA2}
\pdfglyphtounicode{tehhahisolatedarabic}{FC0C}
\pdfglyphtounicode{tehinitialarabic}{FE97}
\pdfglyphtounicode{tehiragana}{3066}
\pdfglyphtounicode{tehjeeminitialarabic}{FCA1}
\pdfglyphtounicode{tehjeemisolatedarabic}{FC0B}
\pdfglyphtounicode{tehmarbutaarabic}{0629}
\pdfglyphtounicode{tehmarbutafinalarabic}{FE94}
\pdfglyphtounicode{tehmedialarabic}{FE98}
\pdfglyphtounicode{tehmeeminitialarabic}{FCA4}
\pdfglyphtounicode{tehmeemisolatedarabic}{FC0E}
\pdfglyphtounicode{tehnoonfinalarabic}{FC73}
\pdfglyphtounicode{tekatakana}{30C6}
\pdfglyphtounicode{tekatakanahalfwidth}{FF83}
\pdfglyphtounicode{telephone}{2121}
\pdfglyphtounicode{telephoneblack}{260E}
\pdfglyphtounicode{telishagedolahebrew}{05A0}
\pdfglyphtounicode{telishaqetanahebrew}{05A9}
\pdfglyphtounicode{tencircle}{2469}
\pdfglyphtounicode{tenideographicparen}{3229}
\pdfglyphtounicode{tenparen}{247D}
\pdfglyphtounicode{tenperiod}{2491}
\pdfglyphtounicode{tenroman}{2179}
\pdfglyphtounicode{tesh}{02A7}
\pdfglyphtounicode{tet}{05D8}
\pdfglyphtounicode{tetdagesh}{FB38}
\pdfglyphtounicode{tetdageshhebrew}{FB38}
\pdfglyphtounicode{tethebrew}{05D8}
\pdfglyphtounicode{tetsecyrillic}{04B5}
\pdfglyphtounicode{tevirhebrew}{059B}
\pdfglyphtounicode{tevirlefthebrew}{059B}
\pdfglyphtounicode{tfm:cmbsy10/diamond}{2662}
\pdfglyphtounicode{tfm:cmbsy10/heart}{2661}
\pdfglyphtounicode{tfm:cmbsy5/diamond}{2662}
\pdfglyphtounicode{tfm:cmbsy5/heart}{2661}
\pdfglyphtounicode{tfm:cmbsy6/diamond}{2662}
\pdfglyphtounicode{tfm:cmbsy6/heart}{2661}
\pdfglyphtounicode{tfm:cmbsy7/diamond}{2662}
\pdfglyphtounicode{tfm:cmbsy7/heart}{2661}
\pdfglyphtounicode{tfm:cmbsy8/diamond}{2662}
\pdfglyphtounicode{tfm:cmbsy8/heart}{2661}
\pdfglyphtounicode{tfm:cmbsy9/diamond}{2662}
\pdfglyphtounicode{tfm:cmbsy9/heart}{2661}
\pdfglyphtounicode{tfm:cmmi10/phi}{03D5}
\pdfglyphtounicode{tfm:cmmi10/phi1}{03C6}
\pdfglyphtounicode{tfm:cmmi12/phi}{03D5}
\pdfglyphtounicode{tfm:cmmi12/phi1}{03C6}
\pdfglyphtounicode{tfm:cmmi5/phi}{03D5}
\pdfglyphtounicode{tfm:cmmi5/phi1}{03C6}
\pdfglyphtounicode{tfm:cmmi6/phi}{03D5}
\pdfglyphtounicode{tfm:cmmi6/phi1}{03C6}
\pdfglyphtounicode{tfm:cmmi7/phi}{03D5}
\pdfglyphtounicode{tfm:cmmi7/phi1}{03C6}
\pdfglyphtounicode{tfm:cmmi8/phi}{03D5}
\pdfglyphtounicode{tfm:cmmi8/phi1}{03C6}
\pdfglyphtounicode{tfm:cmmi9/phi}{03D5}
\pdfglyphtounicode{tfm:cmmi9/phi1}{03C6}
\pdfglyphtounicode{tfm:cmmib10/phi}{03D5}
\pdfglyphtounicode{tfm:cmmib10/phi1}{03C6}
\pdfglyphtounicode{tfm:cmmib5/phi}{03D5}
\pdfglyphtounicode{tfm:cmmib5/phi1}{03C6}
\pdfglyphtounicode{tfm:cmmib6/phi}{03D5}
\pdfglyphtounicode{tfm:cmmib6/phi1}{03C6}
\pdfglyphtounicode{tfm:cmmib7/phi}{03D5}
\pdfglyphtounicode{tfm:cmmib7/phi1}{03C6}
\pdfglyphtounicode{tfm:cmmib8/phi}{03D5}
\pdfglyphtounicode{tfm:cmmib8/phi1}{03C6}
\pdfglyphtounicode{tfm:cmmib9/phi}{03D5}
\pdfglyphtounicode{tfm:cmmib9/phi1}{03C6}
\pdfglyphtounicode{tfm:cmsy10/diamond}{2662}
\pdfglyphtounicode{tfm:cmsy10/heart}{2661}
\pdfglyphtounicode{tfm:cmsy5/heart}{2661}
\pdfglyphtounicode{tfm:cmsy6/diamond}{2662}
\pdfglyphtounicode{tfm:cmsy6/heart}{2661}
\pdfglyphtounicode{tfm:cmsy7/diamond}{2662}
\pdfglyphtounicode{tfm:cmsy7/heart}{2661}
\pdfglyphtounicode{tfm:cmsy8/diamond}{2662}
\pdfglyphtounicode{tfm:cmsy8/heart}{2661}
\pdfglyphtounicode{tfm:cmsy9/diamond}{2662}
\pdfglyphtounicode{tfm:cmsy9/heart}{2661}
\pdfglyphtounicode{tfm:eurb10/phi}{03D5}
\pdfglyphtounicode{tfm:eurb10/phi1}{03C6}
\pdfglyphtounicode{tfm:eurb5/phi}{03D5}
\pdfglyphtounicode{tfm:eurb5/phi1}{03C6}
\pdfglyphtounicode{tfm:eurb6/phi}{03D5}
\pdfglyphtounicode{tfm:eurb6/phi1}{03C6}
\pdfglyphtounicode{tfm:eurb7/phi}{03D5}
\pdfglyphtounicode{tfm:eurb7/phi1}{03C6}
\pdfglyphtounicode{tfm:eurb8/phi}{03D5}
\pdfglyphtounicode{tfm:eurb8/phi1}{03C6}
\pdfglyphtounicode{tfm:eurb9/phi}{03D5}
\pdfglyphtounicode{tfm:eurb9/phi1}{03C6}
\pdfglyphtounicode{tfm:eurm10/phi}{03D5}
\pdfglyphtounicode{tfm:eurm10/phi1}{03C6}
\pdfglyphtounicode{tfm:eurm5/phi}{03D5}
\pdfglyphtounicode{tfm:eurm5/phi1}{03C6}
\pdfglyphtounicode{tfm:eurm6/phi}{03D5}
\pdfglyphtounicode{tfm:eurm6/phi1}{03C6}
\pdfglyphtounicode{tfm:eurm7/phi}{03D5}
\pdfglyphtounicode{tfm:eurm7/phi1}{03C6}
\pdfglyphtounicode{tfm:eurm8/phi}{03D5}
\pdfglyphtounicode{tfm:eurm8/phi1}{03C6}
\pdfglyphtounicode{tfm:eurm9/phi}{03D5}
\pdfglyphtounicode{tfm:eurm9/phi1}{03C6}
\pdfglyphtounicode{tfm:fplmbi/phi}{03D5}
\pdfglyphtounicode{tfm:fplmbi/phi1}{03C6}
\pdfglyphtounicode{tfm:fplmri/phi}{03D5}
\pdfglyphtounicode{tfm:fplmri/phi1}{03C6}
\pdfglyphtounicode{tfm:lmbsy10/diamond}{2662}
\pdfglyphtounicode{tfm:lmbsy10/heart}{2661}
\pdfglyphtounicode{tfm:lmbsy5/diamond}{2662}
\pdfglyphtounicode{tfm:lmbsy5/heart}{2661}
\pdfglyphtounicode{tfm:lmbsy7/diamond}{2662}
\pdfglyphtounicode{tfm:lmbsy7/heart}{2661}
\pdfglyphtounicode{tfm:lmmi10/phi}{03D5}
\pdfglyphtounicode{tfm:lmmi10/phi1}{03C6}
\pdfglyphtounicode{tfm:lmmi12/phi}{03D5}
\pdfglyphtounicode{tfm:lmmi12/phi1}{03C6}
\pdfglyphtounicode{tfm:lmmi5/phi}{03D5}
\pdfglyphtounicode{tfm:lmmi5/phi1}{03C6}
\pdfglyphtounicode{tfm:lmmi6/phi}{03D5}
\pdfglyphtounicode{tfm:lmmi6/phi1}{03C6}
\pdfglyphtounicode{tfm:lmmi7/phi}{03D5}
\pdfglyphtounicode{tfm:lmmi7/phi1}{03C6}
\pdfglyphtounicode{tfm:lmmi8/phi}{03D5}
\pdfglyphtounicode{tfm:lmmi8/phi1}{03C6}
\pdfglyphtounicode{tfm:lmmi9/phi}{03D5}
\pdfglyphtounicode{tfm:lmmi9/phi1}{03C6}
\pdfglyphtounicode{tfm:lmmib10/phi}{03D5}
\pdfglyphtounicode{tfm:lmmib10/phi1}{03C6}
\pdfglyphtounicode{tfm:lmmib5/phi}{03D5}
\pdfglyphtounicode{tfm:lmmib5/phi1}{03C6}
\pdfglyphtounicode{tfm:lmmib7/phi}{03D5}
\pdfglyphtounicode{tfm:lmmib7/phi1}{03C6}
\pdfglyphtounicode{tfm:lmsy10/diamond}{2662}
\pdfglyphtounicode{tfm:lmsy10/heart}{2661}
\pdfglyphtounicode{tfm:lmsy5/diamond}{2662}
\pdfglyphtounicode{tfm:lmsy5/heart}{2661}
\pdfglyphtounicode{tfm:lmsy6/diamond}{2662}
\pdfglyphtounicode{tfm:lmsy6/heart}{2661}
\pdfglyphtounicode{tfm:lmsy7/diamond}{2662}
\pdfglyphtounicode{tfm:lmsy7/heart}{2661}
\pdfglyphtounicode{tfm:lmsy8/diamond}{2662}
\pdfglyphtounicode{tfm:lmsy8/heart}{2661}
\pdfglyphtounicode{tfm:lmsy9/diamond}{2662}
\pdfglyphtounicode{tfm:lmsy9/heart}{2661}
\pdfglyphtounicode{tfm:msam10/diamond}{2662}
\pdfglyphtounicode{tfm:msam5/diamond}{2662}
\pdfglyphtounicode{tfm:msam6/diamond}{2662}
\pdfglyphtounicode{tfm:msam7/diamond}{2662}
\pdfglyphtounicode{tfm:msam8/diamond}{2662}
\pdfglyphtounicode{tfm:msam9/diamond}{2662}
\pdfglyphtounicode{tfm:pxbmia/phi}{03D5}
\pdfglyphtounicode{tfm:pxbmia/phi1}{03C6}
\pdfglyphtounicode{tfm:pxbsy/diamond}{2662}
\pdfglyphtounicode{tfm:pxbsy/heart}{2661}
\pdfglyphtounicode{tfm:pxbsya/diamond}{2662}
\pdfglyphtounicode{tfm:pxmia/phi}{03D5}
\pdfglyphtounicode{tfm:pxmia/phi1}{03C6}
\pdfglyphtounicode{tfm:pxsy/diamond}{2662}
\pdfglyphtounicode{tfm:pxsy/heart}{2661}
\pdfglyphtounicode{tfm:pxsya/diamond}{2662}
\pdfglyphtounicode{tfm:pzdr/a1}{2701}
\pdfglyphtounicode{tfm:pzdr/a10}{2721}
\pdfglyphtounicode{tfm:pzdr/a100}{275E}
\pdfglyphtounicode{tfm:pzdr/a101}{2761}
\pdfglyphtounicode{tfm:pzdr/a102}{2762}
\pdfglyphtounicode{tfm:pzdr/a103}{2763}
\pdfglyphtounicode{tfm:pzdr/a104}{2764}
\pdfglyphtounicode{tfm:pzdr/a105}{2710}
\pdfglyphtounicode{tfm:pzdr/a106}{2765}
\pdfglyphtounicode{tfm:pzdr/a107}{2766}
\pdfglyphtounicode{tfm:pzdr/a108}{2767}
\pdfglyphtounicode{tfm:pzdr/a109}{2660}
\pdfglyphtounicode{tfm:pzdr/a11}{261B}
\pdfglyphtounicode{tfm:pzdr/a110}{2665}
\pdfglyphtounicode{tfm:pzdr/a111}{2666}
\pdfglyphtounicode{tfm:pzdr/a112}{2663}
\pdfglyphtounicode{tfm:pzdr/a117}{2709}
\pdfglyphtounicode{tfm:pzdr/a118}{2708}
\pdfglyphtounicode{tfm:pzdr/a119}{2707}
\pdfglyphtounicode{tfm:pzdr/a12}{261E}
\pdfglyphtounicode{tfm:pzdr/a120}{2460}
\pdfglyphtounicode{tfm:pzdr/a121}{2461}
\pdfglyphtounicode{tfm:pzdr/a122}{2462}
\pdfglyphtounicode{tfm:pzdr/a123}{2463}
\pdfglyphtounicode{tfm:pzdr/a124}{2464}
\pdfglyphtounicode{tfm:pzdr/a125}{2465}
\pdfglyphtounicode{tfm:pzdr/a126}{2466}
\pdfglyphtounicode{tfm:pzdr/a127}{2467}
\pdfglyphtounicode{tfm:pzdr/a128}{2468}
\pdfglyphtounicode{tfm:pzdr/a129}{2469}
\pdfglyphtounicode{tfm:pzdr/a13}{270C}
\pdfglyphtounicode{tfm:pzdr/a130}{2776}
\pdfglyphtounicode{tfm:pzdr/a131}{2777}
\pdfglyphtounicode{tfm:pzdr/a132}{2778}
\pdfglyphtounicode{tfm:pzdr/a133}{2779}
\pdfglyphtounicode{tfm:pzdr/a134}{277A}
\pdfglyphtounicode{tfm:pzdr/a135}{277B}
\pdfglyphtounicode{tfm:pzdr/a136}{277C}
\pdfglyphtounicode{tfm:pzdr/a137}{277D}
\pdfglyphtounicode{tfm:pzdr/a138}{277E}
\pdfglyphtounicode{tfm:pzdr/a139}{277F}
\pdfglyphtounicode{tfm:pzdr/a14}{270D}
\pdfglyphtounicode{tfm:pzdr/a140}{2780}
\pdfglyphtounicode{tfm:pzdr/a141}{2781}
\pdfglyphtounicode{tfm:pzdr/a142}{2782}
\pdfglyphtounicode{tfm:pzdr/a143}{2783}
\pdfglyphtounicode{tfm:pzdr/a144}{2784}
\pdfglyphtounicode{tfm:pzdr/a145}{2785}
\pdfglyphtounicode{tfm:pzdr/a146}{2786}
\pdfglyphtounicode{tfm:pzdr/a147}{2787}
\pdfglyphtounicode{tfm:pzdr/a148}{2788}
\pdfglyphtounicode{tfm:pzdr/a149}{2789}
\pdfglyphtounicode{tfm:pzdr/a15}{270E}
\pdfglyphtounicode{tfm:pzdr/a150}{278A}
\pdfglyphtounicode{tfm:pzdr/a151}{278B}
\pdfglyphtounicode{tfm:pzdr/a152}{278C}
\pdfglyphtounicode{tfm:pzdr/a153}{278D}
\pdfglyphtounicode{tfm:pzdr/a154}{278E}
\pdfglyphtounicode{tfm:pzdr/a155}{278F}
\pdfglyphtounicode{tfm:pzdr/a156}{2790}
\pdfglyphtounicode{tfm:pzdr/a157}{2791}
\pdfglyphtounicode{tfm:pzdr/a158}{2792}
\pdfglyphtounicode{tfm:pzdr/a159}{2793}
\pdfglyphtounicode{tfm:pzdr/a16}{270F}
\pdfglyphtounicode{tfm:pzdr/a160}{2794}
\pdfglyphtounicode{tfm:pzdr/a161}{2192}
\pdfglyphtounicode{tfm:pzdr/a162}{27A3}
\pdfglyphtounicode{tfm:pzdr/a163}{2194}
\pdfglyphtounicode{tfm:pzdr/a164}{2195}
\pdfglyphtounicode{tfm:pzdr/a165}{2799}
\pdfglyphtounicode{tfm:pzdr/a166}{279B}
\pdfglyphtounicode{tfm:pzdr/a167}{279C}
\pdfglyphtounicode{tfm:pzdr/a168}{279D}
\pdfglyphtounicode{tfm:pzdr/a169}{279E}
\pdfglyphtounicode{tfm:pzdr/a17}{2711}
\pdfglyphtounicode{tfm:pzdr/a170}{279F}
\pdfglyphtounicode{tfm:pzdr/a171}{27A0}
\pdfglyphtounicode{tfm:pzdr/a172}{27A1}
\pdfglyphtounicode{tfm:pzdr/a173}{27A2}
\pdfglyphtounicode{tfm:pzdr/a174}{27A4}
\pdfglyphtounicode{tfm:pzdr/a175}{27A5}
\pdfglyphtounicode{tfm:pzdr/a176}{27A6}
\pdfglyphtounicode{tfm:pzdr/a177}{27A7}
\pdfglyphtounicode{tfm:pzdr/a178}{27A8}
\pdfglyphtounicode{tfm:pzdr/a179}{27A9}
\pdfglyphtounicode{tfm:pzdr/a18}{2712}
\pdfglyphtounicode{tfm:pzdr/a180}{27AB}
\pdfglyphtounicode{tfm:pzdr/a181}{27AD}
\pdfglyphtounicode{tfm:pzdr/a182}{27AF}
\pdfglyphtounicode{tfm:pzdr/a183}{27B2}
\pdfglyphtounicode{tfm:pzdr/a184}{27B3}
\pdfglyphtounicode{tfm:pzdr/a185}{27B5}
\pdfglyphtounicode{tfm:pzdr/a186}{27B8}
\pdfglyphtounicode{tfm:pzdr/a187}{27BA}
\pdfglyphtounicode{tfm:pzdr/a188}{27BB}
\pdfglyphtounicode{tfm:pzdr/a189}{27BC}
\pdfglyphtounicode{tfm:pzdr/a19}{2713}
\pdfglyphtounicode{tfm:pzdr/a190}{27BD}
\pdfglyphtounicode{tfm:pzdr/a191}{27BE}
\pdfglyphtounicode{tfm:pzdr/a192}{279A}
\pdfglyphtounicode{tfm:pzdr/a193}{27AA}
\pdfglyphtounicode{tfm:pzdr/a194}{27B6}
\pdfglyphtounicode{tfm:pzdr/a195}{27B9}
\pdfglyphtounicode{tfm:pzdr/a196}{2798}
\pdfglyphtounicode{tfm:pzdr/a197}{27B4}
\pdfglyphtounicode{tfm:pzdr/a198}{27B7}
\pdfglyphtounicode{tfm:pzdr/a199}{27AC}
\pdfglyphtounicode{tfm:pzdr/a2}{2702}
\pdfglyphtounicode{tfm:pzdr/a20}{2714}
\pdfglyphtounicode{tfm:pzdr/a200}{27AE}
\pdfglyphtounicode{tfm:pzdr/a201}{27B1}
\pdfglyphtounicode{tfm:pzdr/a202}{2703}
\pdfglyphtounicode{tfm:pzdr/a203}{2750}
\pdfglyphtounicode{tfm:pzdr/a204}{2752}
\pdfglyphtounicode{tfm:pzdr/a205}{276E}
\pdfglyphtounicode{tfm:pzdr/a206}{2770}
\pdfglyphtounicode{tfm:pzdr/a21}{2715}
\pdfglyphtounicode{tfm:pzdr/a22}{2716}
\pdfglyphtounicode{tfm:pzdr/a23}{2717}
\pdfglyphtounicode{tfm:pzdr/a24}{2718}
\pdfglyphtounicode{tfm:pzdr/a25}{2719}
\pdfglyphtounicode{tfm:pzdr/a26}{271A}
\pdfglyphtounicode{tfm:pzdr/a27}{271B}
\pdfglyphtounicode{tfm:pzdr/a28}{271C}
\pdfglyphtounicode{tfm:pzdr/a29}{2722}
\pdfglyphtounicode{tfm:pzdr/a3}{2704}
\pdfglyphtounicode{tfm:pzdr/a30}{2723}
\pdfglyphtounicode{tfm:pzdr/a31}{2724}
\pdfglyphtounicode{tfm:pzdr/a32}{2725}
\pdfglyphtounicode{tfm:pzdr/a33}{2726}
\pdfglyphtounicode{tfm:pzdr/a34}{2727}
\pdfglyphtounicode{tfm:pzdr/a35}{2605}
\pdfglyphtounicode{tfm:pzdr/a36}{2729}
\pdfglyphtounicode{tfm:pzdr/a37}{272A}
\pdfglyphtounicode{tfm:pzdr/a38}{272B}
\pdfglyphtounicode{tfm:pzdr/a39}{272C}
\pdfglyphtounicode{tfm:pzdr/a4}{260E}
\pdfglyphtounicode{tfm:pzdr/a40}{272D}
\pdfglyphtounicode{tfm:pzdr/a41}{272E}
\pdfglyphtounicode{tfm:pzdr/a42}{272F}
\pdfglyphtounicode{tfm:pzdr/a43}{2730}
\pdfglyphtounicode{tfm:pzdr/a44}{2731}
\pdfglyphtounicode{tfm:pzdr/a45}{2732}
\pdfglyphtounicode{tfm:pzdr/a46}{2733}
\pdfglyphtounicode{tfm:pzdr/a47}{2734}
\pdfglyphtounicode{tfm:pzdr/a48}{2735}
\pdfglyphtounicode{tfm:pzdr/a49}{2736}
\pdfglyphtounicode{tfm:pzdr/a5}{2706}
\pdfglyphtounicode{tfm:pzdr/a50}{2737}
\pdfglyphtounicode{tfm:pzdr/a51}{2738}
\pdfglyphtounicode{tfm:pzdr/a52}{2739}
\pdfglyphtounicode{tfm:pzdr/a53}{273A}
\pdfglyphtounicode{tfm:pzdr/a54}{273B}
\pdfglyphtounicode{tfm:pzdr/a55}{273C}
\pdfglyphtounicode{tfm:pzdr/a56}{273D}
\pdfglyphtounicode{tfm:pzdr/a57}{273E}
\pdfglyphtounicode{tfm:pzdr/a58}{273F}
\pdfglyphtounicode{tfm:pzdr/a59}{2740}
\pdfglyphtounicode{tfm:pzdr/a6}{271D}
\pdfglyphtounicode{tfm:pzdr/a60}{2741}
\pdfglyphtounicode{tfm:pzdr/a61}{2742}
\pdfglyphtounicode{tfm:pzdr/a62}{2743}
\pdfglyphtounicode{tfm:pzdr/a63}{2744}
\pdfglyphtounicode{tfm:pzdr/a64}{2745}
\pdfglyphtounicode{tfm:pzdr/a65}{2746}
\pdfglyphtounicode{tfm:pzdr/a66}{2747}
\pdfglyphtounicode{tfm:pzdr/a67}{2748}
\pdfglyphtounicode{tfm:pzdr/a68}{2749}
\pdfglyphtounicode{tfm:pzdr/a69}{274A}
\pdfglyphtounicode{tfm:pzdr/a7}{271E}
\pdfglyphtounicode{tfm:pzdr/a70}{274B}
\pdfglyphtounicode{tfm:pzdr/a71}{25CF}
\pdfglyphtounicode{tfm:pzdr/a72}{274D}
\pdfglyphtounicode{tfm:pzdr/a73}{25A0}
\pdfglyphtounicode{tfm:pzdr/a74}{274F}
\pdfglyphtounicode{tfm:pzdr/a75}{2751}
\pdfglyphtounicode{tfm:pzdr/a76}{25B2}
\pdfglyphtounicode{tfm:pzdr/a77}{25BC}
\pdfglyphtounicode{tfm:pzdr/a78}{25C6}
\pdfglyphtounicode{tfm:pzdr/a79}{2756}
\pdfglyphtounicode{tfm:pzdr/a8}{271F}
\pdfglyphtounicode{tfm:pzdr/a81}{25D7}
\pdfglyphtounicode{tfm:pzdr/a82}{2758}
\pdfglyphtounicode{tfm:pzdr/a83}{2759}
\pdfglyphtounicode{tfm:pzdr/a84}{275A}
\pdfglyphtounicode{tfm:pzdr/a85}{276F}
\pdfglyphtounicode{tfm:pzdr/a86}{2771}
\pdfglyphtounicode{tfm:pzdr/a87}{2772}
\pdfglyphtounicode{tfm:pzdr/a88}{2773}
\pdfglyphtounicode{tfm:pzdr/a89}{2768}
\pdfglyphtounicode{tfm:pzdr/a9}{2720}
\pdfglyphtounicode{tfm:pzdr/a90}{2769}
\pdfglyphtounicode{tfm:pzdr/a91}{276C}
\pdfglyphtounicode{tfm:pzdr/a92}{276D}
\pdfglyphtounicode{tfm:pzdr/a93}{276A}
\pdfglyphtounicode{tfm:pzdr/a94}{276B}
\pdfglyphtounicode{tfm:pzdr/a95}{2774}
\pdfglyphtounicode{tfm:pzdr/a96}{2775}
\pdfglyphtounicode{tfm:pzdr/a97}{275B}
\pdfglyphtounicode{tfm:pzdr/a98}{275C}
\pdfglyphtounicode{tfm:pzdr/a99}{275D}
\pdfglyphtounicode{tfm:rpxbmi/phi}{03D5}
\pdfglyphtounicode{tfm:rpxbmi/phi1}{03C6}
\pdfglyphtounicode{tfm:rpxmi/phi}{03D5}
\pdfglyphtounicode{tfm:rpxmi/phi1}{03C6}
\pdfglyphtounicode{tfm:rpzdr/a1}{2701}
\pdfglyphtounicode{tfm:rpzdr/a10}{2721}
\pdfglyphtounicode{tfm:rpzdr/a100}{275E}
\pdfglyphtounicode{tfm:rpzdr/a101}{2761}
\pdfglyphtounicode{tfm:rpzdr/a102}{2762}
\pdfglyphtounicode{tfm:rpzdr/a103}{2763}
\pdfglyphtounicode{tfm:rpzdr/a104}{2764}
\pdfglyphtounicode{tfm:rpzdr/a105}{2710}
\pdfglyphtounicode{tfm:rpzdr/a106}{2765}
\pdfglyphtounicode{tfm:rpzdr/a107}{2766}
\pdfglyphtounicode{tfm:rpzdr/a108}{2767}
\pdfglyphtounicode{tfm:rpzdr/a109}{2660}
\pdfglyphtounicode{tfm:rpzdr/a11}{261B}
\pdfglyphtounicode{tfm:rpzdr/a110}{2665}
\pdfglyphtounicode{tfm:rpzdr/a111}{2666}
\pdfglyphtounicode{tfm:rpzdr/a112}{2663}
\pdfglyphtounicode{tfm:rpzdr/a117}{2709}
\pdfglyphtounicode{tfm:rpzdr/a118}{2708}
\pdfglyphtounicode{tfm:rpzdr/a119}{2707}
\pdfglyphtounicode{tfm:rpzdr/a12}{261E}
\pdfglyphtounicode{tfm:rpzdr/a120}{2460}
\pdfglyphtounicode{tfm:rpzdr/a121}{2461}
\pdfglyphtounicode{tfm:rpzdr/a122}{2462}
\pdfglyphtounicode{tfm:rpzdr/a123}{2463}
\pdfglyphtounicode{tfm:rpzdr/a124}{2464}
\pdfglyphtounicode{tfm:rpzdr/a125}{2465}
\pdfglyphtounicode{tfm:rpzdr/a126}{2466}
\pdfglyphtounicode{tfm:rpzdr/a127}{2467}
\pdfglyphtounicode{tfm:rpzdr/a128}{2468}
\pdfglyphtounicode{tfm:rpzdr/a129}{2469}
\pdfglyphtounicode{tfm:rpzdr/a13}{270C}
\pdfglyphtounicode{tfm:rpzdr/a130}{2776}
\pdfglyphtounicode{tfm:rpzdr/a131}{2777}
\pdfglyphtounicode{tfm:rpzdr/a132}{2778}
\pdfglyphtounicode{tfm:rpzdr/a133}{2779}
\pdfglyphtounicode{tfm:rpzdr/a134}{277A}
\pdfglyphtounicode{tfm:rpzdr/a135}{277B}
\pdfglyphtounicode{tfm:rpzdr/a136}{277C}
\pdfglyphtounicode{tfm:rpzdr/a137}{277D}
\pdfglyphtounicode{tfm:rpzdr/a138}{277E}
\pdfglyphtounicode{tfm:rpzdr/a139}{277F}
\pdfglyphtounicode{tfm:rpzdr/a14}{270D}
\pdfglyphtounicode{tfm:rpzdr/a140}{2780}
\pdfglyphtounicode{tfm:rpzdr/a141}{2781}
\pdfglyphtounicode{tfm:rpzdr/a142}{2782}
\pdfglyphtounicode{tfm:rpzdr/a143}{2783}
\pdfglyphtounicode{tfm:rpzdr/a144}{2784}
\pdfglyphtounicode{tfm:rpzdr/a145}{2785}
\pdfglyphtounicode{tfm:rpzdr/a146}{2786}
\pdfglyphtounicode{tfm:rpzdr/a147}{2787}
\pdfglyphtounicode{tfm:rpzdr/a148}{2788}
\pdfglyphtounicode{tfm:rpzdr/a149}{2789}
\pdfglyphtounicode{tfm:rpzdr/a15}{270E}
\pdfglyphtounicode{tfm:rpzdr/a150}{278A}
\pdfglyphtounicode{tfm:rpzdr/a151}{278B}
\pdfglyphtounicode{tfm:rpzdr/a152}{278C}
\pdfglyphtounicode{tfm:rpzdr/a153}{278D}
\pdfglyphtounicode{tfm:rpzdr/a154}{278E}
\pdfglyphtounicode{tfm:rpzdr/a155}{278F}
\pdfglyphtounicode{tfm:rpzdr/a156}{2790}
\pdfglyphtounicode{tfm:rpzdr/a157}{2791}
\pdfglyphtounicode{tfm:rpzdr/a158}{2792}
\pdfglyphtounicode{tfm:rpzdr/a159}{2793}
\pdfglyphtounicode{tfm:rpzdr/a16}{270F}
\pdfglyphtounicode{tfm:rpzdr/a160}{2794}
\pdfglyphtounicode{tfm:rpzdr/a161}{2192}
\pdfglyphtounicode{tfm:rpzdr/a162}{27A3}
\pdfglyphtounicode{tfm:rpzdr/a163}{2194}
\pdfglyphtounicode{tfm:rpzdr/a164}{2195}
\pdfglyphtounicode{tfm:rpzdr/a165}{2799}
\pdfglyphtounicode{tfm:rpzdr/a166}{279B}
\pdfglyphtounicode{tfm:rpzdr/a167}{279C}
\pdfglyphtounicode{tfm:rpzdr/a168}{279D}
\pdfglyphtounicode{tfm:rpzdr/a169}{279E}
\pdfglyphtounicode{tfm:rpzdr/a17}{2711}
\pdfglyphtounicode{tfm:rpzdr/a170}{279F}
\pdfglyphtounicode{tfm:rpzdr/a171}{27A0}
\pdfglyphtounicode{tfm:rpzdr/a172}{27A1}
\pdfglyphtounicode{tfm:rpzdr/a173}{27A2}
\pdfglyphtounicode{tfm:rpzdr/a174}{27A4}
\pdfglyphtounicode{tfm:rpzdr/a175}{27A5}
\pdfglyphtounicode{tfm:rpzdr/a176}{27A6}
\pdfglyphtounicode{tfm:rpzdr/a177}{27A7}
\pdfglyphtounicode{tfm:rpzdr/a178}{27A8}
\pdfglyphtounicode{tfm:rpzdr/a179}{27A9}
\pdfglyphtounicode{tfm:rpzdr/a18}{2712}
\pdfglyphtounicode{tfm:rpzdr/a180}{27AB}
\pdfglyphtounicode{tfm:rpzdr/a181}{27AD}
\pdfglyphtounicode{tfm:rpzdr/a182}{27AF}
\pdfglyphtounicode{tfm:rpzdr/a183}{27B2}
\pdfglyphtounicode{tfm:rpzdr/a184}{27B3}
\pdfglyphtounicode{tfm:rpzdr/a185}{27B5}
\pdfglyphtounicode{tfm:rpzdr/a186}{27B8}
\pdfglyphtounicode{tfm:rpzdr/a187}{27BA}
\pdfglyphtounicode{tfm:rpzdr/a188}{27BB}
\pdfglyphtounicode{tfm:rpzdr/a189}{27BC}
\pdfglyphtounicode{tfm:rpzdr/a19}{2713}
\pdfglyphtounicode{tfm:rpzdr/a190}{27BD}
\pdfglyphtounicode{tfm:rpzdr/a191}{27BE}
\pdfglyphtounicode{tfm:rpzdr/a192}{279A}
\pdfglyphtounicode{tfm:rpzdr/a193}{27AA}
\pdfglyphtounicode{tfm:rpzdr/a194}{27B6}
\pdfglyphtounicode{tfm:rpzdr/a195}{27B9}
\pdfglyphtounicode{tfm:rpzdr/a196}{2798}
\pdfglyphtounicode{tfm:rpzdr/a197}{27B4}
\pdfglyphtounicode{tfm:rpzdr/a198}{27B7}
\pdfglyphtounicode{tfm:rpzdr/a199}{27AC}
\pdfglyphtounicode{tfm:rpzdr/a2}{2702}
\pdfglyphtounicode{tfm:rpzdr/a20}{2714}
\pdfglyphtounicode{tfm:rpzdr/a200}{27AE}
\pdfglyphtounicode{tfm:rpzdr/a201}{27B1}
\pdfglyphtounicode{tfm:rpzdr/a202}{2703}
\pdfglyphtounicode{tfm:rpzdr/a203}{2750}
\pdfglyphtounicode{tfm:rpzdr/a204}{2752}
\pdfglyphtounicode{tfm:rpzdr/a205}{276E}
\pdfglyphtounicode{tfm:rpzdr/a206}{2770}
\pdfglyphtounicode{tfm:rpzdr/a21}{2715}
\pdfglyphtounicode{tfm:rpzdr/a22}{2716}
\pdfglyphtounicode{tfm:rpzdr/a23}{2717}
\pdfglyphtounicode{tfm:rpzdr/a24}{2718}
\pdfglyphtounicode{tfm:rpzdr/a25}{2719}
\pdfglyphtounicode{tfm:rpzdr/a26}{271A}
\pdfglyphtounicode{tfm:rpzdr/a27}{271B}
\pdfglyphtounicode{tfm:rpzdr/a28}{271C}
\pdfglyphtounicode{tfm:rpzdr/a29}{2722}
\pdfglyphtounicode{tfm:rpzdr/a3}{2704}
\pdfglyphtounicode{tfm:rpzdr/a30}{2723}
\pdfglyphtounicode{tfm:rpzdr/a31}{2724}
\pdfglyphtounicode{tfm:rpzdr/a32}{2725}
\pdfglyphtounicode{tfm:rpzdr/a33}{2726}
\pdfglyphtounicode{tfm:rpzdr/a34}{2727}
\pdfglyphtounicode{tfm:rpzdr/a35}{2605}
\pdfglyphtounicode{tfm:rpzdr/a36}{2729}
\pdfglyphtounicode{tfm:rpzdr/a37}{272A}
\pdfglyphtounicode{tfm:rpzdr/a38}{272B}
\pdfglyphtounicode{tfm:rpzdr/a39}{272C}
\pdfglyphtounicode{tfm:rpzdr/a4}{260E}
\pdfglyphtounicode{tfm:rpzdr/a40}{272D}
\pdfglyphtounicode{tfm:rpzdr/a41}{272E}
\pdfglyphtounicode{tfm:rpzdr/a42}{272F}
\pdfglyphtounicode{tfm:rpzdr/a43}{2730}
\pdfglyphtounicode{tfm:rpzdr/a44}{2731}
\pdfglyphtounicode{tfm:rpzdr/a45}{2732}
\pdfglyphtounicode{tfm:rpzdr/a46}{2733}
\pdfglyphtounicode{tfm:rpzdr/a47}{2734}
\pdfglyphtounicode{tfm:rpzdr/a48}{2735}
\pdfglyphtounicode{tfm:rpzdr/a49}{2736}
\pdfglyphtounicode{tfm:rpzdr/a5}{2706}
\pdfglyphtounicode{tfm:rpzdr/a50}{2737}
\pdfglyphtounicode{tfm:rpzdr/a51}{2738}
\pdfglyphtounicode{tfm:rpzdr/a52}{2739}
\pdfglyphtounicode{tfm:rpzdr/a53}{273A}
\pdfglyphtounicode{tfm:rpzdr/a54}{273B}
\pdfglyphtounicode{tfm:rpzdr/a55}{273C}
\pdfglyphtounicode{tfm:rpzdr/a56}{273D}
\pdfglyphtounicode{tfm:rpzdr/a57}{273E}
\pdfglyphtounicode{tfm:rpzdr/a58}{273F}
\pdfglyphtounicode{tfm:rpzdr/a59}{2740}
\pdfglyphtounicode{tfm:rpzdr/a6}{271D}
\pdfglyphtounicode{tfm:rpzdr/a60}{2741}
\pdfglyphtounicode{tfm:rpzdr/a61}{2742}
\pdfglyphtounicode{tfm:rpzdr/a62}{2743}
\pdfglyphtounicode{tfm:rpzdr/a63}{2744}
\pdfglyphtounicode{tfm:rpzdr/a64}{2745}
\pdfglyphtounicode{tfm:rpzdr/a65}{2746}
\pdfglyphtounicode{tfm:rpzdr/a66}{2747}
\pdfglyphtounicode{tfm:rpzdr/a67}{2748}
\pdfglyphtounicode{tfm:rpzdr/a68}{2749}
\pdfglyphtounicode{tfm:rpzdr/a69}{274A}
\pdfglyphtounicode{tfm:rpzdr/a7}{271E}
\pdfglyphtounicode{tfm:rpzdr/a70}{274B}
\pdfglyphtounicode{tfm:rpzdr/a71}{25CF}
\pdfglyphtounicode{tfm:rpzdr/a72}{274D}
\pdfglyphtounicode{tfm:rpzdr/a73}{25A0}
\pdfglyphtounicode{tfm:rpzdr/a74}{274F}
\pdfglyphtounicode{tfm:rpzdr/a75}{2751}
\pdfglyphtounicode{tfm:rpzdr/a76}{25B2}
\pdfglyphtounicode{tfm:rpzdr/a77}{25BC}
\pdfglyphtounicode{tfm:rpzdr/a78}{25C6}
\pdfglyphtounicode{tfm:rpzdr/a79}{2756}
\pdfglyphtounicode{tfm:rpzdr/a8}{271F}
\pdfglyphtounicode{tfm:rpzdr/a81}{25D7}
\pdfglyphtounicode{tfm:rpzdr/a82}{2758}
\pdfglyphtounicode{tfm:rpzdr/a83}{2759}
\pdfglyphtounicode{tfm:rpzdr/a84}{275A}
\pdfglyphtounicode{tfm:rpzdr/a85}{276F}
\pdfglyphtounicode{tfm:rpzdr/a86}{2771}
\pdfglyphtounicode{tfm:rpzdr/a87}{2772}
\pdfglyphtounicode{tfm:rpzdr/a88}{2773}
\pdfglyphtounicode{tfm:rpzdr/a89}{2768}
\pdfglyphtounicode{tfm:rpzdr/a9}{2720}
\pdfglyphtounicode{tfm:rpzdr/a90}{2769}
\pdfglyphtounicode{tfm:rpzdr/a91}{276C}
\pdfglyphtounicode{tfm:rpzdr/a92}{276D}
\pdfglyphtounicode{tfm:rpzdr/a93}{276A}
\pdfglyphtounicode{tfm:rpzdr/a94}{276B}
\pdfglyphtounicode{tfm:rpzdr/a95}{2774}
\pdfglyphtounicode{tfm:rpzdr/a96}{2775}
\pdfglyphtounicode{tfm:rpzdr/a97}{275B}
\pdfglyphtounicode{tfm:rpzdr/a98}{275C}
\pdfglyphtounicode{tfm:rpzdr/a99}{275D}
\pdfglyphtounicode{tfm:rtxbmi/phi}{03D5}
\pdfglyphtounicode{tfm:rtxbmi/phi1}{03C6}
\pdfglyphtounicode{tfm:rtxmi/phi}{03D5}
\pdfglyphtounicode{tfm:rtxmi/phi1}{03C6}
\pdfglyphtounicode{tfm:txbmia/phi}{03D5}
\pdfglyphtounicode{tfm:txbmia/phi1}{03C6}
\pdfglyphtounicode{tfm:txbsy/diamond}{2662}
\pdfglyphtounicode{tfm:txbsy/heart}{2661}
\pdfglyphtounicode{tfm:txbsya/diamond}{2662}
\pdfglyphtounicode{tfm:txmia/phi}{03D5}
\pdfglyphtounicode{tfm:txmia/phi1}{03C6}
\pdfglyphtounicode{tfm:txsy/diamond}{2662}
\pdfglyphtounicode{tfm:txsy/heart}{2661}
\pdfglyphtounicode{tfm:txsya/diamond}{2662}
\pdfglyphtounicode{tfm:zd/a1}{2701}
\pdfglyphtounicode{tfm:zd/a10}{2721}
\pdfglyphtounicode{tfm:zd/a100}{275E}
\pdfglyphtounicode{tfm:zd/a101}{2761}
\pdfglyphtounicode{tfm:zd/a102}{2762}
\pdfglyphtounicode{tfm:zd/a103}{2763}
\pdfglyphtounicode{tfm:zd/a104}{2764}
\pdfglyphtounicode{tfm:zd/a105}{2710}
\pdfglyphtounicode{tfm:zd/a106}{2765}
\pdfglyphtounicode{tfm:zd/a107}{2766}
\pdfglyphtounicode{tfm:zd/a108}{2767}
\pdfglyphtounicode{tfm:zd/a109}{2660}
\pdfglyphtounicode{tfm:zd/a11}{261B}
\pdfglyphtounicode{tfm:zd/a110}{2665}
\pdfglyphtounicode{tfm:zd/a111}{2666}
\pdfglyphtounicode{tfm:zd/a112}{2663}
\pdfglyphtounicode{tfm:zd/a117}{2709}
\pdfglyphtounicode{tfm:zd/a118}{2708}
\pdfglyphtounicode{tfm:zd/a119}{2707}
\pdfglyphtounicode{tfm:zd/a12}{261E}
\pdfglyphtounicode{tfm:zd/a120}{2460}
\pdfglyphtounicode{tfm:zd/a121}{2461}
\pdfglyphtounicode{tfm:zd/a122}{2462}
\pdfglyphtounicode{tfm:zd/a123}{2463}
\pdfglyphtounicode{tfm:zd/a124}{2464}
\pdfglyphtounicode{tfm:zd/a125}{2465}
\pdfglyphtounicode{tfm:zd/a126}{2466}
\pdfglyphtounicode{tfm:zd/a127}{2467}
\pdfglyphtounicode{tfm:zd/a128}{2468}
\pdfglyphtounicode{tfm:zd/a129}{2469}
\pdfglyphtounicode{tfm:zd/a13}{270C}
\pdfglyphtounicode{tfm:zd/a130}{2776}
\pdfglyphtounicode{tfm:zd/a131}{2777}
\pdfglyphtounicode{tfm:zd/a132}{2778}
\pdfglyphtounicode{tfm:zd/a133}{2779}
\pdfglyphtounicode{tfm:zd/a134}{277A}
\pdfglyphtounicode{tfm:zd/a135}{277B}
\pdfglyphtounicode{tfm:zd/a136}{277C}
\pdfglyphtounicode{tfm:zd/a137}{277D}
\pdfglyphtounicode{tfm:zd/a138}{277E}
\pdfglyphtounicode{tfm:zd/a139}{277F}
\pdfglyphtounicode{tfm:zd/a14}{270D}
\pdfglyphtounicode{tfm:zd/a140}{2780}
\pdfglyphtounicode{tfm:zd/a141}{2781}
\pdfglyphtounicode{tfm:zd/a142}{2782}
\pdfglyphtounicode{tfm:zd/a143}{2783}
\pdfglyphtounicode{tfm:zd/a144}{2784}
\pdfglyphtounicode{tfm:zd/a145}{2785}
\pdfglyphtounicode{tfm:zd/a146}{2786}
\pdfglyphtounicode{tfm:zd/a147}{2787}
\pdfglyphtounicode{tfm:zd/a148}{2788}
\pdfglyphtounicode{tfm:zd/a149}{2789}
\pdfglyphtounicode{tfm:zd/a15}{270E}
\pdfglyphtounicode{tfm:zd/a150}{278A}
\pdfglyphtounicode{tfm:zd/a151}{278B}
\pdfglyphtounicode{tfm:zd/a152}{278C}
\pdfglyphtounicode{tfm:zd/a153}{278D}
\pdfglyphtounicode{tfm:zd/a154}{278E}
\pdfglyphtounicode{tfm:zd/a155}{278F}
\pdfglyphtounicode{tfm:zd/a156}{2790}
\pdfglyphtounicode{tfm:zd/a157}{2791}
\pdfglyphtounicode{tfm:zd/a158}{2792}
\pdfglyphtounicode{tfm:zd/a159}{2793}
\pdfglyphtounicode{tfm:zd/a16}{270F}
\pdfglyphtounicode{tfm:zd/a160}{2794}
\pdfglyphtounicode{tfm:zd/a161}{2192}
\pdfglyphtounicode{tfm:zd/a162}{27A3}
\pdfglyphtounicode{tfm:zd/a163}{2194}
\pdfglyphtounicode{tfm:zd/a164}{2195}
\pdfglyphtounicode{tfm:zd/a165}{2799}
\pdfglyphtounicode{tfm:zd/a166}{279B}
\pdfglyphtounicode{tfm:zd/a167}{279C}
\pdfglyphtounicode{tfm:zd/a168}{279D}
\pdfglyphtounicode{tfm:zd/a169}{279E}
\pdfglyphtounicode{tfm:zd/a17}{2711}
\pdfglyphtounicode{tfm:zd/a170}{279F}
\pdfglyphtounicode{tfm:zd/a171}{27A0}
\pdfglyphtounicode{tfm:zd/a172}{27A1}
\pdfglyphtounicode{tfm:zd/a173}{27A2}
\pdfglyphtounicode{tfm:zd/a174}{27A4}
\pdfglyphtounicode{tfm:zd/a175}{27A5}
\pdfglyphtounicode{tfm:zd/a176}{27A6}
\pdfglyphtounicode{tfm:zd/a177}{27A7}
\pdfglyphtounicode{tfm:zd/a178}{27A8}
\pdfglyphtounicode{tfm:zd/a179}{27A9}
\pdfglyphtounicode{tfm:zd/a18}{2712}
\pdfglyphtounicode{tfm:zd/a180}{27AB}
\pdfglyphtounicode{tfm:zd/a181}{27AD}
\pdfglyphtounicode{tfm:zd/a182}{27AF}
\pdfglyphtounicode{tfm:zd/a183}{27B2}
\pdfglyphtounicode{tfm:zd/a184}{27B3}
\pdfglyphtounicode{tfm:zd/a185}{27B5}
\pdfglyphtounicode{tfm:zd/a186}{27B8}
\pdfglyphtounicode{tfm:zd/a187}{27BA}
\pdfglyphtounicode{tfm:zd/a188}{27BB}
\pdfglyphtounicode{tfm:zd/a189}{27BC}
\pdfglyphtounicode{tfm:zd/a19}{2713}
\pdfglyphtounicode{tfm:zd/a190}{27BD}
\pdfglyphtounicode{tfm:zd/a191}{27BE}
\pdfglyphtounicode{tfm:zd/a192}{279A}
\pdfglyphtounicode{tfm:zd/a193}{27AA}
\pdfglyphtounicode{tfm:zd/a194}{27B6}
\pdfglyphtounicode{tfm:zd/a195}{27B9}
\pdfglyphtounicode{tfm:zd/a196}{2798}
\pdfglyphtounicode{tfm:zd/a197}{27B4}
\pdfglyphtounicode{tfm:zd/a198}{27B7}
\pdfglyphtounicode{tfm:zd/a199}{27AC}
\pdfglyphtounicode{tfm:zd/a2}{2702}
\pdfglyphtounicode{tfm:zd/a20}{2714}
\pdfglyphtounicode{tfm:zd/a200}{27AE}
\pdfglyphtounicode{tfm:zd/a201}{27B1}
\pdfglyphtounicode{tfm:zd/a202}{2703}
\pdfglyphtounicode{tfm:zd/a203}{2750}
\pdfglyphtounicode{tfm:zd/a204}{2752}
\pdfglyphtounicode{tfm:zd/a205}{276E}
\pdfglyphtounicode{tfm:zd/a206}{2770}
\pdfglyphtounicode{tfm:zd/a21}{2715}
\pdfglyphtounicode{tfm:zd/a22}{2716}
\pdfglyphtounicode{tfm:zd/a23}{2717}
\pdfglyphtounicode{tfm:zd/a24}{2718}
\pdfglyphtounicode{tfm:zd/a25}{2719}
\pdfglyphtounicode{tfm:zd/a26}{271A}
\pdfglyphtounicode{tfm:zd/a27}{271B}
\pdfglyphtounicode{tfm:zd/a28}{271C}
\pdfglyphtounicode{tfm:zd/a29}{2722}
\pdfglyphtounicode{tfm:zd/a3}{2704}
\pdfglyphtounicode{tfm:zd/a30}{2723}
\pdfglyphtounicode{tfm:zd/a31}{2724}
\pdfglyphtounicode{tfm:zd/a32}{2725}
\pdfglyphtounicode{tfm:zd/a33}{2726}
\pdfglyphtounicode{tfm:zd/a34}{2727}
\pdfglyphtounicode{tfm:zd/a35}{2605}
\pdfglyphtounicode{tfm:zd/a36}{2729}
\pdfglyphtounicode{tfm:zd/a37}{272A}
\pdfglyphtounicode{tfm:zd/a38}{272B}
\pdfglyphtounicode{tfm:zd/a39}{272C}
\pdfglyphtounicode{tfm:zd/a4}{260E}
\pdfglyphtounicode{tfm:zd/a40}{272D}
\pdfglyphtounicode{tfm:zd/a41}{272E}
\pdfglyphtounicode{tfm:zd/a42}{272F}
\pdfglyphtounicode{tfm:zd/a43}{2730}
\pdfglyphtounicode{tfm:zd/a44}{2731}
\pdfglyphtounicode{tfm:zd/a45}{2732}
\pdfglyphtounicode{tfm:zd/a46}{2733}
\pdfglyphtounicode{tfm:zd/a47}{2734}
\pdfglyphtounicode{tfm:zd/a48}{2735}
\pdfglyphtounicode{tfm:zd/a49}{2736}
\pdfglyphtounicode{tfm:zd/a5}{2706}
\pdfglyphtounicode{tfm:zd/a50}{2737}
\pdfglyphtounicode{tfm:zd/a51}{2738}
\pdfglyphtounicode{tfm:zd/a52}{2739}
\pdfglyphtounicode{tfm:zd/a53}{273A}
\pdfglyphtounicode{tfm:zd/a54}{273B}
\pdfglyphtounicode{tfm:zd/a55}{273C}
\pdfglyphtounicode{tfm:zd/a56}{273D}
\pdfglyphtounicode{tfm:zd/a57}{273E}
\pdfglyphtounicode{tfm:zd/a58}{273F}
\pdfglyphtounicode{tfm:zd/a59}{2740}
\pdfglyphtounicode{tfm:zd/a6}{271D}
\pdfglyphtounicode{tfm:zd/a60}{2741}
\pdfglyphtounicode{tfm:zd/a61}{2742}
\pdfglyphtounicode{tfm:zd/a62}{2743}
\pdfglyphtounicode{tfm:zd/a63}{2744}
\pdfglyphtounicode{tfm:zd/a64}{2745}
\pdfglyphtounicode{tfm:zd/a65}{2746}
\pdfglyphtounicode{tfm:zd/a66}{2747}
\pdfglyphtounicode{tfm:zd/a67}{2748}
\pdfglyphtounicode{tfm:zd/a68}{2749}
\pdfglyphtounicode{tfm:zd/a69}{274A}
\pdfglyphtounicode{tfm:zd/a7}{271E}
\pdfglyphtounicode{tfm:zd/a70}{274B}
\pdfglyphtounicode{tfm:zd/a71}{25CF}
\pdfglyphtounicode{tfm:zd/a72}{274D}
\pdfglyphtounicode{tfm:zd/a73}{25A0}
\pdfglyphtounicode{tfm:zd/a74}{274F}
\pdfglyphtounicode{tfm:zd/a75}{2751}
\pdfglyphtounicode{tfm:zd/a76}{25B2}
\pdfglyphtounicode{tfm:zd/a77}{25BC}
\pdfglyphtounicode{tfm:zd/a78}{25C6}
\pdfglyphtounicode{tfm:zd/a79}{2756}
\pdfglyphtounicode{tfm:zd/a8}{271F}
\pdfglyphtounicode{tfm:zd/a81}{25D7}
\pdfglyphtounicode{tfm:zd/a82}{2758}
\pdfglyphtounicode{tfm:zd/a83}{2759}
\pdfglyphtounicode{tfm:zd/a84}{275A}
\pdfglyphtounicode{tfm:zd/a85}{276F}
\pdfglyphtounicode{tfm:zd/a86}{2771}
\pdfglyphtounicode{tfm:zd/a87}{2772}
\pdfglyphtounicode{tfm:zd/a88}{2773}
\pdfglyphtounicode{tfm:zd/a89}{2768}
\pdfglyphtounicode{tfm:zd/a9}{2720}
\pdfglyphtounicode{tfm:zd/a90}{2769}
\pdfglyphtounicode{tfm:zd/a91}{276C}
\pdfglyphtounicode{tfm:zd/a92}{276D}
\pdfglyphtounicode{tfm:zd/a93}{276A}
\pdfglyphtounicode{tfm:zd/a94}{276B}
\pdfglyphtounicode{tfm:zd/a95}{2774}
\pdfglyphtounicode{tfm:zd/a96}{2775}
\pdfglyphtounicode{tfm:zd/a97}{275B}
\pdfglyphtounicode{tfm:zd/a98}{275C}
\pdfglyphtounicode{tfm:zd/a99}{275D}
\pdfglyphtounicode{tfm:zpzdr-reversed/a1}{2701}
\pdfglyphtounicode{tfm:zpzdr-reversed/a10}{2721}
\pdfglyphtounicode{tfm:zpzdr-reversed/a100}{275E}
\pdfglyphtounicode{tfm:zpzdr-reversed/a101}{2761}
\pdfglyphtounicode{tfm:zpzdr-reversed/a102}{2762}
\pdfglyphtounicode{tfm:zpzdr-reversed/a103}{2763}
\pdfglyphtounicode{tfm:zpzdr-reversed/a104}{2764}
\pdfglyphtounicode{tfm:zpzdr-reversed/a105}{2710}
\pdfglyphtounicode{tfm:zpzdr-reversed/a106}{2765}
\pdfglyphtounicode{tfm:zpzdr-reversed/a107}{2766}
\pdfglyphtounicode{tfm:zpzdr-reversed/a108}{2767}
\pdfglyphtounicode{tfm:zpzdr-reversed/a109}{2660}
\pdfglyphtounicode{tfm:zpzdr-reversed/a11}{261B}
\pdfglyphtounicode{tfm:zpzdr-reversed/a110}{2665}
\pdfglyphtounicode{tfm:zpzdr-reversed/a111}{2666}
\pdfglyphtounicode{tfm:zpzdr-reversed/a112}{2663}
\pdfglyphtounicode{tfm:zpzdr-reversed/a117}{2709}
\pdfglyphtounicode{tfm:zpzdr-reversed/a118}{2708}
\pdfglyphtounicode{tfm:zpzdr-reversed/a119}{2707}
\pdfglyphtounicode{tfm:zpzdr-reversed/a12}{261E}
\pdfglyphtounicode{tfm:zpzdr-reversed/a120}{2460}
\pdfglyphtounicode{tfm:zpzdr-reversed/a121}{2461}
\pdfglyphtounicode{tfm:zpzdr-reversed/a122}{2462}
\pdfglyphtounicode{tfm:zpzdr-reversed/a123}{2463}
\pdfglyphtounicode{tfm:zpzdr-reversed/a124}{2464}
\pdfglyphtounicode{tfm:zpzdr-reversed/a125}{2465}
\pdfglyphtounicode{tfm:zpzdr-reversed/a126}{2466}
\pdfglyphtounicode{tfm:zpzdr-reversed/a127}{2467}
\pdfglyphtounicode{tfm:zpzdr-reversed/a128}{2468}
\pdfglyphtounicode{tfm:zpzdr-reversed/a129}{2469}
\pdfglyphtounicode{tfm:zpzdr-reversed/a13}{270C}
\pdfglyphtounicode{tfm:zpzdr-reversed/a130}{2776}
\pdfglyphtounicode{tfm:zpzdr-reversed/a131}{2777}
\pdfglyphtounicode{tfm:zpzdr-reversed/a132}{2778}
\pdfglyphtounicode{tfm:zpzdr-reversed/a133}{2779}
\pdfglyphtounicode{tfm:zpzdr-reversed/a134}{277A}
\pdfglyphtounicode{tfm:zpzdr-reversed/a135}{277B}
\pdfglyphtounicode{tfm:zpzdr-reversed/a136}{277C}
\pdfglyphtounicode{tfm:zpzdr-reversed/a137}{277D}
\pdfglyphtounicode{tfm:zpzdr-reversed/a138}{277E}
\pdfglyphtounicode{tfm:zpzdr-reversed/a139}{277F}
\pdfglyphtounicode{tfm:zpzdr-reversed/a14}{270D}
\pdfglyphtounicode{tfm:zpzdr-reversed/a140}{2780}
\pdfglyphtounicode{tfm:zpzdr-reversed/a141}{2781}
\pdfglyphtounicode{tfm:zpzdr-reversed/a142}{2782}
\pdfglyphtounicode{tfm:zpzdr-reversed/a143}{2783}
\pdfglyphtounicode{tfm:zpzdr-reversed/a144}{2784}
\pdfglyphtounicode{tfm:zpzdr-reversed/a145}{2785}
\pdfglyphtounicode{tfm:zpzdr-reversed/a146}{2786}
\pdfglyphtounicode{tfm:zpzdr-reversed/a147}{2787}
\pdfglyphtounicode{tfm:zpzdr-reversed/a148}{2788}
\pdfglyphtounicode{tfm:zpzdr-reversed/a149}{2789}
\pdfglyphtounicode{tfm:zpzdr-reversed/a15}{270E}
\pdfglyphtounicode{tfm:zpzdr-reversed/a150}{278A}
\pdfglyphtounicode{tfm:zpzdr-reversed/a151}{278B}
\pdfglyphtounicode{tfm:zpzdr-reversed/a152}{278C}
\pdfglyphtounicode{tfm:zpzdr-reversed/a153}{278D}
\pdfglyphtounicode{tfm:zpzdr-reversed/a154}{278E}
\pdfglyphtounicode{tfm:zpzdr-reversed/a155}{278F}
\pdfglyphtounicode{tfm:zpzdr-reversed/a156}{2790}
\pdfglyphtounicode{tfm:zpzdr-reversed/a157}{2791}
\pdfglyphtounicode{tfm:zpzdr-reversed/a158}{2792}
\pdfglyphtounicode{tfm:zpzdr-reversed/a159}{2793}
\pdfglyphtounicode{tfm:zpzdr-reversed/a16}{270F}
\pdfglyphtounicode{tfm:zpzdr-reversed/a160}{2794}
\pdfglyphtounicode{tfm:zpzdr-reversed/a161}{2192}
\pdfglyphtounicode{tfm:zpzdr-reversed/a162}{27A3}
\pdfglyphtounicode{tfm:zpzdr-reversed/a163}{2194}
\pdfglyphtounicode{tfm:zpzdr-reversed/a164}{2195}
\pdfglyphtounicode{tfm:zpzdr-reversed/a165}{2799}
\pdfglyphtounicode{tfm:zpzdr-reversed/a166}{279B}
\pdfglyphtounicode{tfm:zpzdr-reversed/a167}{279C}
\pdfglyphtounicode{tfm:zpzdr-reversed/a168}{279D}
\pdfglyphtounicode{tfm:zpzdr-reversed/a169}{279E}
\pdfglyphtounicode{tfm:zpzdr-reversed/a17}{2711}
\pdfglyphtounicode{tfm:zpzdr-reversed/a170}{279F}
\pdfglyphtounicode{tfm:zpzdr-reversed/a171}{27A0}
\pdfglyphtounicode{tfm:zpzdr-reversed/a172}{27A1}
\pdfglyphtounicode{tfm:zpzdr-reversed/a173}{27A2}
\pdfglyphtounicode{tfm:zpzdr-reversed/a174}{27A4}
\pdfglyphtounicode{tfm:zpzdr-reversed/a175}{27A5}
\pdfglyphtounicode{tfm:zpzdr-reversed/a176}{27A6}
\pdfglyphtounicode{tfm:zpzdr-reversed/a177}{27A7}
\pdfglyphtounicode{tfm:zpzdr-reversed/a178}{27A8}
\pdfglyphtounicode{tfm:zpzdr-reversed/a179}{27A9}
\pdfglyphtounicode{tfm:zpzdr-reversed/a18}{2712}
\pdfglyphtounicode{tfm:zpzdr-reversed/a180}{27AB}
\pdfglyphtounicode{tfm:zpzdr-reversed/a181}{27AD}
\pdfglyphtounicode{tfm:zpzdr-reversed/a182}{27AF}
\pdfglyphtounicode{tfm:zpzdr-reversed/a183}{27B2}
\pdfglyphtounicode{tfm:zpzdr-reversed/a184}{27B3}
\pdfglyphtounicode{tfm:zpzdr-reversed/a185}{27B5}
\pdfglyphtounicode{tfm:zpzdr-reversed/a186}{27B8}
\pdfglyphtounicode{tfm:zpzdr-reversed/a187}{27BA}
\pdfglyphtounicode{tfm:zpzdr-reversed/a188}{27BB}
\pdfglyphtounicode{tfm:zpzdr-reversed/a189}{27BC}
\pdfglyphtounicode{tfm:zpzdr-reversed/a19}{2713}
\pdfglyphtounicode{tfm:zpzdr-reversed/a190}{27BD}
\pdfglyphtounicode{tfm:zpzdr-reversed/a191}{27BE}
\pdfglyphtounicode{tfm:zpzdr-reversed/a192}{279A}
\pdfglyphtounicode{tfm:zpzdr-reversed/a193}{27AA}
\pdfglyphtounicode{tfm:zpzdr-reversed/a194}{27B6}
\pdfglyphtounicode{tfm:zpzdr-reversed/a195}{27B9}
\pdfglyphtounicode{tfm:zpzdr-reversed/a196}{2798}
\pdfglyphtounicode{tfm:zpzdr-reversed/a197}{27B4}
\pdfglyphtounicode{tfm:zpzdr-reversed/a198}{27B7}
\pdfglyphtounicode{tfm:zpzdr-reversed/a199}{27AC}
\pdfglyphtounicode{tfm:zpzdr-reversed/a2}{2702}
\pdfglyphtounicode{tfm:zpzdr-reversed/a20}{2714}
\pdfglyphtounicode{tfm:zpzdr-reversed/a200}{27AE}
\pdfglyphtounicode{tfm:zpzdr-reversed/a201}{27B1}
\pdfglyphtounicode{tfm:zpzdr-reversed/a202}{2703}
\pdfglyphtounicode{tfm:zpzdr-reversed/a203}{2750}
\pdfglyphtounicode{tfm:zpzdr-reversed/a204}{2752}
\pdfglyphtounicode{tfm:zpzdr-reversed/a205}{276E}
\pdfglyphtounicode{tfm:zpzdr-reversed/a206}{2770}
\pdfglyphtounicode{tfm:zpzdr-reversed/a21}{2715}
\pdfglyphtounicode{tfm:zpzdr-reversed/a22}{2716}
\pdfglyphtounicode{tfm:zpzdr-reversed/a23}{2717}
\pdfglyphtounicode{tfm:zpzdr-reversed/a24}{2718}
\pdfglyphtounicode{tfm:zpzdr-reversed/a25}{2719}
\pdfglyphtounicode{tfm:zpzdr-reversed/a26}{271A}
\pdfglyphtounicode{tfm:zpzdr-reversed/a27}{271B}
\pdfglyphtounicode{tfm:zpzdr-reversed/a28}{271C}
\pdfglyphtounicode{tfm:zpzdr-reversed/a29}{2722}
\pdfglyphtounicode{tfm:zpzdr-reversed/a3}{2704}
\pdfglyphtounicode{tfm:zpzdr-reversed/a30}{2723}
\pdfglyphtounicode{tfm:zpzdr-reversed/a31}{2724}
\pdfglyphtounicode{tfm:zpzdr-reversed/a32}{2725}
\pdfglyphtounicode{tfm:zpzdr-reversed/a33}{2726}
\pdfglyphtounicode{tfm:zpzdr-reversed/a34}{2727}
\pdfglyphtounicode{tfm:zpzdr-reversed/a35}{2605}
\pdfglyphtounicode{tfm:zpzdr-reversed/a36}{2729}
\pdfglyphtounicode{tfm:zpzdr-reversed/a37}{272A}
\pdfglyphtounicode{tfm:zpzdr-reversed/a38}{272B}
\pdfglyphtounicode{tfm:zpzdr-reversed/a39}{272C}
\pdfglyphtounicode{tfm:zpzdr-reversed/a4}{260E}
\pdfglyphtounicode{tfm:zpzdr-reversed/a40}{272D}
\pdfglyphtounicode{tfm:zpzdr-reversed/a41}{272E}
\pdfglyphtounicode{tfm:zpzdr-reversed/a42}{272F}
\pdfglyphtounicode{tfm:zpzdr-reversed/a43}{2730}
\pdfglyphtounicode{tfm:zpzdr-reversed/a44}{2731}
\pdfglyphtounicode{tfm:zpzdr-reversed/a45}{2732}
\pdfglyphtounicode{tfm:zpzdr-reversed/a46}{2733}
\pdfglyphtounicode{tfm:zpzdr-reversed/a47}{2734}
\pdfglyphtounicode{tfm:zpzdr-reversed/a48}{2735}
\pdfglyphtounicode{tfm:zpzdr-reversed/a49}{2736}
\pdfglyphtounicode{tfm:zpzdr-reversed/a5}{2706}
\pdfglyphtounicode{tfm:zpzdr-reversed/a50}{2737}
\pdfglyphtounicode{tfm:zpzdr-reversed/a51}{2738}
\pdfglyphtounicode{tfm:zpzdr-reversed/a52}{2739}
\pdfglyphtounicode{tfm:zpzdr-reversed/a53}{273A}
\pdfglyphtounicode{tfm:zpzdr-reversed/a54}{273B}
\pdfglyphtounicode{tfm:zpzdr-reversed/a55}{273C}
\pdfglyphtounicode{tfm:zpzdr-reversed/a56}{273D}
\pdfglyphtounicode{tfm:zpzdr-reversed/a57}{273E}
\pdfglyphtounicode{tfm:zpzdr-reversed/a58}{273F}
\pdfglyphtounicode{tfm:zpzdr-reversed/a59}{2740}
\pdfglyphtounicode{tfm:zpzdr-reversed/a6}{271D}
\pdfglyphtounicode{tfm:zpzdr-reversed/a60}{2741}
\pdfglyphtounicode{tfm:zpzdr-reversed/a61}{2742}
\pdfglyphtounicode{tfm:zpzdr-reversed/a62}{2743}
\pdfglyphtounicode{tfm:zpzdr-reversed/a63}{2744}
\pdfglyphtounicode{tfm:zpzdr-reversed/a64}{2745}
\pdfglyphtounicode{tfm:zpzdr-reversed/a65}{2746}
\pdfglyphtounicode{tfm:zpzdr-reversed/a66}{2747}
\pdfglyphtounicode{tfm:zpzdr-reversed/a67}{2748}
\pdfglyphtounicode{tfm:zpzdr-reversed/a68}{2749}
\pdfglyphtounicode{tfm:zpzdr-reversed/a69}{274A}
\pdfglyphtounicode{tfm:zpzdr-reversed/a7}{271E}
\pdfglyphtounicode{tfm:zpzdr-reversed/a70}{274B}
\pdfglyphtounicode{tfm:zpzdr-reversed/a71}{25CF}
\pdfglyphtounicode{tfm:zpzdr-reversed/a72}{274D}
\pdfglyphtounicode{tfm:zpzdr-reversed/a73}{25A0}
\pdfglyphtounicode{tfm:zpzdr-reversed/a74}{274F}
\pdfglyphtounicode{tfm:zpzdr-reversed/a75}{2751}
\pdfglyphtounicode{tfm:zpzdr-reversed/a76}{25B2}
\pdfglyphtounicode{tfm:zpzdr-reversed/a77}{25BC}
\pdfglyphtounicode{tfm:zpzdr-reversed/a78}{25C6}
\pdfglyphtounicode{tfm:zpzdr-reversed/a79}{2756}
\pdfglyphtounicode{tfm:zpzdr-reversed/a8}{271F}
\pdfglyphtounicode{tfm:zpzdr-reversed/a81}{25D7}
\pdfglyphtounicode{tfm:zpzdr-reversed/a82}{2758}
\pdfglyphtounicode{tfm:zpzdr-reversed/a83}{2759}
\pdfglyphtounicode{tfm:zpzdr-reversed/a84}{275A}
\pdfglyphtounicode{tfm:zpzdr-reversed/a85}{276F}
\pdfglyphtounicode{tfm:zpzdr-reversed/a86}{2771}
\pdfglyphtounicode{tfm:zpzdr-reversed/a87}{2772}
\pdfglyphtounicode{tfm:zpzdr-reversed/a88}{2773}
\pdfglyphtounicode{tfm:zpzdr-reversed/a89}{2768}
\pdfglyphtounicode{tfm:zpzdr-reversed/a9}{2720}
\pdfglyphtounicode{tfm:zpzdr-reversed/a90}{2769}
\pdfglyphtounicode{tfm:zpzdr-reversed/a91}{276C}
\pdfglyphtounicode{tfm:zpzdr-reversed/a92}{276D}
\pdfglyphtounicode{tfm:zpzdr-reversed/a93}{276A}
\pdfglyphtounicode{tfm:zpzdr-reversed/a94}{276B}
\pdfglyphtounicode{tfm:zpzdr-reversed/a95}{2774}
\pdfglyphtounicode{tfm:zpzdr-reversed/a96}{2775}
\pdfglyphtounicode{tfm:zpzdr-reversed/a97}{275B}
\pdfglyphtounicode{tfm:zpzdr-reversed/a98}{275C}
\pdfglyphtounicode{tfm:zpzdr-reversed/a99}{275D}
\pdfglyphtounicode{thabengali}{09A5}
\pdfglyphtounicode{thadeva}{0925}
\pdfglyphtounicode{thagujarati}{0AA5}
\pdfglyphtounicode{thagurmukhi}{0A25}
\pdfglyphtounicode{thalarabic}{0630}
\pdfglyphtounicode{thalfinalarabic}{FEAC}
\pdfglyphtounicode{thanthakhatlowleftthai}{F898}
\pdfglyphtounicode{thanthakhatlowrightthai}{F897}
\pdfglyphtounicode{thanthakhatthai}{0E4C}
\pdfglyphtounicode{thanthakhatupperleftthai}{F896}
\pdfglyphtounicode{theharabic}{062B}
\pdfglyphtounicode{thehfinalarabic}{FE9A}
\pdfglyphtounicode{thehinitialarabic}{FE9B}
\pdfglyphtounicode{thehmedialarabic}{FE9C}
\pdfglyphtounicode{thereexists}{2203}
\pdfglyphtounicode{therefore}{2234}
\pdfglyphtounicode{theta}{03B8}
\pdfglyphtounicode{theta1}{03D1}
\pdfglyphtounicode{thetasymbolgreek}{03D1}
\pdfglyphtounicode{thieuthacirclekorean}{3279}
\pdfglyphtounicode{thieuthaparenkorean}{3219}
\pdfglyphtounicode{thieuthcirclekorean}{326B}
\pdfglyphtounicode{thieuthkorean}{314C}
\pdfglyphtounicode{thieuthparenkorean}{320B}
\pdfglyphtounicode{thirteencircle}{246C}
\pdfglyphtounicode{thirteenparen}{2480}
\pdfglyphtounicode{thirteenperiod}{2494}
\pdfglyphtounicode{thonangmonthothai}{0E11}
\pdfglyphtounicode{thook}{01AD}
\pdfglyphtounicode{thophuthaothai}{0E12}
\pdfglyphtounicode{thorn}{00FE}
\pdfglyphtounicode{thothahanthai}{0E17}
\pdfglyphtounicode{thothanthai}{0E10}
\pdfglyphtounicode{thothongthai}{0E18}
\pdfglyphtounicode{thothungthai}{0E16}
\pdfglyphtounicode{thousandcyrillic}{0482}
\pdfglyphtounicode{thousandsseparatorarabic}{066C}
\pdfglyphtounicode{thousandsseparatorpersian}{066C}
\pdfglyphtounicode{three}{0033}
\pdfglyphtounicode{threearabic}{0663}
\pdfglyphtounicode{threebengali}{09E9}
\pdfglyphtounicode{threecircle}{2462}
\pdfglyphtounicode{threecircleinversesansserif}{278C}
\pdfglyphtounicode{threedeva}{0969}
\pdfglyphtounicode{threeeighths}{215C}
\pdfglyphtounicode{threegujarati}{0AE9}
\pdfglyphtounicode{threegurmukhi}{0A69}
\pdfglyphtounicode{threehackarabic}{0663}
\pdfglyphtounicode{threehangzhou}{3023}
\pdfglyphtounicode{threeideographicparen}{3222}
\pdfglyphtounicode{threeinferior}{2083}
\pdfglyphtounicode{threemonospace}{FF13}
\pdfglyphtounicode{threenumeratorbengali}{09F6}
\pdfglyphtounicode{threeoldstyle}{0033}
\pdfglyphtounicode{threeparen}{2476}
\pdfglyphtounicode{threeperiod}{248A}
\pdfglyphtounicode{threepersian}{06F3}
\pdfglyphtounicode{threequarters}{00BE}
\pdfglyphtounicode{threequartersemdash}{F6DE}
\pdfglyphtounicode{threeroman}{2172}
\pdfglyphtounicode{threesuperior}{00B3}
\pdfglyphtounicode{threethai}{0E53}
\pdfglyphtounicode{thzsquare}{3394}
\pdfglyphtounicode{tihiragana}{3061}
\pdfglyphtounicode{tikatakana}{30C1}
\pdfglyphtounicode{tikatakanahalfwidth}{FF81}
\pdfglyphtounicode{tikeutacirclekorean}{3270}
\pdfglyphtounicode{tikeutaparenkorean}{3210}
\pdfglyphtounicode{tikeutcirclekorean}{3262}
\pdfglyphtounicode{tikeutkorean}{3137}
\pdfglyphtounicode{tikeutparenkorean}{3202}
\pdfglyphtounicode{tilde}{02DC}
\pdfglyphtounicode{tildebelowcmb}{0330}
\pdfglyphtounicode{tildecmb}{0303}
\pdfglyphtounicode{tildecomb}{0303}
\pdfglyphtounicode{tildedoublecmb}{0360}
\pdfglyphtounicode{tildeoperator}{223C}
\pdfglyphtounicode{tildeoverlaycmb}{0334}
\pdfglyphtounicode{tildeverticalcmb}{033E}
\pdfglyphtounicode{timescircle}{2297}
\pdfglyphtounicode{tipehahebrew}{0596}
\pdfglyphtounicode{tipehalefthebrew}{0596}
\pdfglyphtounicode{tippigurmukhi}{0A70}
\pdfglyphtounicode{titlocyrilliccmb}{0483}
\pdfglyphtounicode{tiwnarmenian}{057F}
\pdfglyphtounicode{tlinebelow}{1E6F}
\pdfglyphtounicode{tmonospace}{FF54}
\pdfglyphtounicode{toarmenian}{0569}
\pdfglyphtounicode{tohiragana}{3068}
\pdfglyphtounicode{tokatakana}{30C8}
\pdfglyphtounicode{tokatakanahalfwidth}{FF84}
\pdfglyphtounicode{tonebarextrahighmod}{02E5}
\pdfglyphtounicode{tonebarextralowmod}{02E9}
\pdfglyphtounicode{tonebarhighmod}{02E6}
\pdfglyphtounicode{tonebarlowmod}{02E8}
\pdfglyphtounicode{tonebarmidmod}{02E7}
\pdfglyphtounicode{tonefive}{01BD}
\pdfglyphtounicode{tonesix}{0185}
\pdfglyphtounicode{tonetwo}{01A8}
\pdfglyphtounicode{tonos}{0384}
\pdfglyphtounicode{tonsquare}{3327}
\pdfglyphtounicode{topatakthai}{0E0F}
\pdfglyphtounicode{tortoiseshellbracketleft}{3014}
\pdfglyphtounicode{tortoiseshellbracketleftsmall}{FE5D}
\pdfglyphtounicode{tortoiseshellbracketleftvertical}{FE39}
\pdfglyphtounicode{tortoiseshellbracketright}{3015}
\pdfglyphtounicode{tortoiseshellbracketrightsmall}{FE5E}
\pdfglyphtounicode{tortoiseshellbracketrightvertical}{FE3A}
\pdfglyphtounicode{totaothai}{0E15}
\pdfglyphtounicode{tpalatalhook}{01AB}
\pdfglyphtounicode{tparen}{24AF}
\pdfglyphtounicode{trademark}{2122}
\pdfglyphtounicode{trademarksans}{2122}
\pdfglyphtounicode{trademarkserif}{2122}
\pdfglyphtounicode{tretroflexhook}{0288}
\pdfglyphtounicode{triagdn}{25BC}
\pdfglyphtounicode{triaglf}{25C4}
\pdfglyphtounicode{triagrt}{25BA}
\pdfglyphtounicode{triagup}{25B2}
\pdfglyphtounicode{triangle}{25B3}
\pdfglyphtounicode{triangledownsld}{25BC}
\pdfglyphtounicode{triangleinv}{25BD}
\pdfglyphtounicode{triangleleft}{25C1}
\pdfglyphtounicode{triangleleftequal}{22B4}
\pdfglyphtounicode{triangleleftsld}{25C0}
\pdfglyphtounicode{triangleright}{25B7}
\pdfglyphtounicode{trianglerightequal}{22B5}
\pdfglyphtounicode{trianglerightsld}{25B6}
\pdfglyphtounicode{trianglesolid}{25B2}
\pdfglyphtounicode{ts}{02A6}
\pdfglyphtounicode{tsadi}{05E6}
\pdfglyphtounicode{tsadidagesh}{FB46}
\pdfglyphtounicode{tsadidageshhebrew}{FB46}
\pdfglyphtounicode{tsadihebrew}{05E6}
\pdfglyphtounicode{tsecyrillic}{0446}
\pdfglyphtounicode{tsere}{05B5}
\pdfglyphtounicode{tsere12}{05B5}
\pdfglyphtounicode{tsere1e}{05B5}
\pdfglyphtounicode{tsere2b}{05B5}
\pdfglyphtounicode{tserehebrew}{05B5}
\pdfglyphtounicode{tserenarrowhebrew}{05B5}
\pdfglyphtounicode{tserequarterhebrew}{05B5}
\pdfglyphtounicode{tserewidehebrew}{05B5}
\pdfglyphtounicode{tshecyrillic}{045B}
\pdfglyphtounicode{tsuperior}{0074}
\pdfglyphtounicode{ttabengali}{099F}
\pdfglyphtounicode{ttadeva}{091F}
\pdfglyphtounicode{ttagujarati}{0A9F}
\pdfglyphtounicode{ttagurmukhi}{0A1F}
\pdfglyphtounicode{tteharabic}{0679}
\pdfglyphtounicode{ttehfinalarabic}{FB67}
\pdfglyphtounicode{ttehinitialarabic}{FB68}
\pdfglyphtounicode{ttehmedialarabic}{FB69}
\pdfglyphtounicode{tthabengali}{09A0}
\pdfglyphtounicode{tthadeva}{0920}
\pdfglyphtounicode{tthagujarati}{0AA0}
\pdfglyphtounicode{tthagurmukhi}{0A20}
\pdfglyphtounicode{tturned}{0287}
\pdfglyphtounicode{tuhiragana}{3064}
\pdfglyphtounicode{tukatakana}{30C4}
\pdfglyphtounicode{tukatakanahalfwidth}{FF82}
\pdfglyphtounicode{turnstileleft}{22A2}
\pdfglyphtounicode{turnstileright}{22A3}
\pdfglyphtounicode{tusmallhiragana}{3063}
\pdfglyphtounicode{tusmallkatakana}{30C3}
\pdfglyphtounicode{tusmallkatakanahalfwidth}{FF6F}
\pdfglyphtounicode{twelvecircle}{246B}
\pdfglyphtounicode{twelveparen}{247F}
\pdfglyphtounicode{twelveperiod}{2493}
\pdfglyphtounicode{twelveroman}{217B}
\pdfglyphtounicode{twentycircle}{2473}
\pdfglyphtounicode{twentyhangzhou}{5344}
\pdfglyphtounicode{twentyparen}{2487}
\pdfglyphtounicode{twentyperiod}{249B}
\pdfglyphtounicode{two}{0032}
\pdfglyphtounicode{twoarabic}{0662}
\pdfglyphtounicode{twobengali}{09E8}
\pdfglyphtounicode{twocircle}{2461}
\pdfglyphtounicode{twocircleinversesansserif}{278B}
\pdfglyphtounicode{twodeva}{0968}
\pdfglyphtounicode{twodotenleader}{2025}
\pdfglyphtounicode{twodotleader}{2025}
\pdfglyphtounicode{twodotleadervertical}{FE30}
\pdfglyphtounicode{twogujarati}{0AE8}
\pdfglyphtounicode{twogurmukhi}{0A68}
\pdfglyphtounicode{twohackarabic}{0662}
\pdfglyphtounicode{twohangzhou}{3022}
\pdfglyphtounicode{twoideographicparen}{3221}
\pdfglyphtounicode{twoinferior}{2082}
\pdfglyphtounicode{twomonospace}{FF12}
\pdfglyphtounicode{twonumeratorbengali}{09F5}
\pdfglyphtounicode{twooldstyle}{0032}
\pdfglyphtounicode{twoparen}{2475}
\pdfglyphtounicode{twoperiod}{2489}
\pdfglyphtounicode{twopersian}{06F2}
\pdfglyphtounicode{tworoman}{2171}
\pdfglyphtounicode{twostroke}{01BB}
\pdfglyphtounicode{twosuperior}{00B2}
\pdfglyphtounicode{twothai}{0E52}
\pdfglyphtounicode{twothirds}{2154}
\pdfglyphtounicode{u}{0075}
\pdfglyphtounicode{uacute}{00FA}
\pdfglyphtounicode{ubar}{0289}
\pdfglyphtounicode{ubengali}{0989}
\pdfglyphtounicode{ubopomofo}{3128}
\pdfglyphtounicode{ubreve}{016D}
\pdfglyphtounicode{ucaron}{01D4}
\pdfglyphtounicode{ucircle}{24E4}
\pdfglyphtounicode{ucircumflex}{00FB}
\pdfglyphtounicode{ucircumflexbelow}{1E77}
\pdfglyphtounicode{ucyrillic}{0443}
\pdfglyphtounicode{udattadeva}{0951}
\pdfglyphtounicode{udblacute}{0171}
\pdfglyphtounicode{udblgrave}{0215}
\pdfglyphtounicode{udeva}{0909}
\pdfglyphtounicode{udieresis}{00FC}
\pdfglyphtounicode{udieresisacute}{01D8}
\pdfglyphtounicode{udieresisbelow}{1E73}
\pdfglyphtounicode{udieresiscaron}{01DA}
\pdfglyphtounicode{udieresiscyrillic}{04F1}
\pdfglyphtounicode{udieresisgrave}{01DC}
\pdfglyphtounicode{udieresismacron}{01D6}
\pdfglyphtounicode{udotbelow}{1EE5}
\pdfglyphtounicode{ugrave}{00F9}
\pdfglyphtounicode{ugujarati}{0A89}
\pdfglyphtounicode{ugurmukhi}{0A09}
\pdfglyphtounicode{uhiragana}{3046}
\pdfglyphtounicode{uhookabove}{1EE7}
\pdfglyphtounicode{uhorn}{01B0}
\pdfglyphtounicode{uhornacute}{1EE9}
\pdfglyphtounicode{uhorndotbelow}{1EF1}
\pdfglyphtounicode{uhorngrave}{1EEB}
\pdfglyphtounicode{uhornhookabove}{1EED}
\pdfglyphtounicode{uhorntilde}{1EEF}
\pdfglyphtounicode{uhungarumlaut}{0171}
\pdfglyphtounicode{uhungarumlautcyrillic}{04F3}
\pdfglyphtounicode{uinvertedbreve}{0217}
\pdfglyphtounicode{ukatakana}{30A6}
\pdfglyphtounicode{ukatakanahalfwidth}{FF73}
\pdfglyphtounicode{ukcyrillic}{0479}
\pdfglyphtounicode{ukorean}{315C}
\pdfglyphtounicode{umacron}{016B}
\pdfglyphtounicode{umacroncyrillic}{04EF}
\pdfglyphtounicode{umacrondieresis}{1E7B}
\pdfglyphtounicode{umatragurmukhi}{0A41}
\pdfglyphtounicode{umonospace}{FF55}
\pdfglyphtounicode{underscore}{005F}
\pdfglyphtounicode{underscoredbl}{2017}
\pdfglyphtounicode{underscoremonospace}{FF3F}
\pdfglyphtounicode{underscorevertical}{FE33}
\pdfglyphtounicode{underscorewavy}{FE4F}
\pdfglyphtounicode{union}{222A}
\pdfglyphtounicode{uniondbl}{22D3}
\pdfglyphtounicode{unionmulti}{228E}
\pdfglyphtounicode{unionsq}{2294}
\pdfglyphtounicode{universal}{2200}
\pdfglyphtounicode{uogonek}{0173}
\pdfglyphtounicode{uparen}{24B0}
\pdfglyphtounicode{upblock}{2580}
\pdfglyphtounicode{upperdothebrew}{05C4}
\pdfglyphtounicode{uprise}{22CF}
\pdfglyphtounicode{upsilon}{03C5}
\pdfglyphtounicode{upsilondieresis}{03CB}
\pdfglyphtounicode{upsilondieresistonos}{03B0}
\pdfglyphtounicode{upsilonlatin}{028A}
\pdfglyphtounicode{upsilontonos}{03CD}
\pdfglyphtounicode{upslope}{29F8}
\pdfglyphtounicode{uptackbelowcmb}{031D}
\pdfglyphtounicode{uptackmod}{02D4}
\pdfglyphtounicode{uragurmukhi}{0A73}
\pdfglyphtounicode{uring}{016F}
\pdfglyphtounicode{ushortcyrillic}{045E}
\pdfglyphtounicode{usmallhiragana}{3045}
\pdfglyphtounicode{usmallkatakana}{30A5}
\pdfglyphtounicode{usmallkatakanahalfwidth}{FF69}
\pdfglyphtounicode{ustraightcyrillic}{04AF}
\pdfglyphtounicode{ustraightstrokecyrillic}{04B1}
\pdfglyphtounicode{utilde}{0169}
\pdfglyphtounicode{utildeacute}{1E79}
\pdfglyphtounicode{utildebelow}{1E75}
\pdfglyphtounicode{uubengali}{098A}
\pdfglyphtounicode{uudeva}{090A}
\pdfglyphtounicode{uugujarati}{0A8A}
\pdfglyphtounicode{uugurmukhi}{0A0A}
\pdfglyphtounicode{uumatragurmukhi}{0A42}
\pdfglyphtounicode{uuvowelsignbengali}{09C2}
\pdfglyphtounicode{uuvowelsigndeva}{0942}
\pdfglyphtounicode{uuvowelsigngujarati}{0AC2}
\pdfglyphtounicode{uvowelsignbengali}{09C1}
\pdfglyphtounicode{uvowelsigndeva}{0941}
\pdfglyphtounicode{uvowelsigngujarati}{0AC1}
\pdfglyphtounicode{v}{0076}
\pdfglyphtounicode{vadeva}{0935}
\pdfglyphtounicode{vagujarati}{0AB5}
\pdfglyphtounicode{vagurmukhi}{0A35}
\pdfglyphtounicode{vakatakana}{30F7}
\pdfglyphtounicode{vav}{05D5}
\pdfglyphtounicode{vavdagesh}{FB35}
\pdfglyphtounicode{vavdagesh65}{FB35}
\pdfglyphtounicode{vavdageshhebrew}{FB35}
\pdfglyphtounicode{vavhebrew}{05D5}
\pdfglyphtounicode{vavholam}{FB4B}
\pdfglyphtounicode{vavholamhebrew}{FB4B}
\pdfglyphtounicode{vavvavhebrew}{05F0}
\pdfglyphtounicode{vavyodhebrew}{05F1}
\pdfglyphtounicode{vcircle}{24E5}
\pdfglyphtounicode{vdotbelow}{1E7F}
\pdfglyphtounicode{vector}{20D7}
\pdfglyphtounicode{vecyrillic}{0432}
\pdfglyphtounicode{veharabic}{06A4}
\pdfglyphtounicode{vehfinalarabic}{FB6B}
\pdfglyphtounicode{vehinitialarabic}{FB6C}
\pdfglyphtounicode{vehmedialarabic}{FB6D}
\pdfglyphtounicode{vekatakana}{30F9}
\pdfglyphtounicode{venus}{2640}
\pdfglyphtounicode{verticalbar}{007C}
\pdfglyphtounicode{verticallineabovecmb}{030D}
\pdfglyphtounicode{verticallinebelowcmb}{0329}
\pdfglyphtounicode{verticallinelowmod}{02CC}
\pdfglyphtounicode{verticallinemod}{02C8}
\pdfglyphtounicode{vewarmenian}{057E}
\pdfglyphtounicode{vhook}{028B}
\pdfglyphtounicode{vikatakana}{30F8}
\pdfglyphtounicode{viramabengali}{09CD}
\pdfglyphtounicode{viramadeva}{094D}
\pdfglyphtounicode{viramagujarati}{0ACD}
\pdfglyphtounicode{visargabengali}{0983}
\pdfglyphtounicode{visargadeva}{0903}
\pdfglyphtounicode{visargagujarati}{0A83}
\pdfglyphtounicode{visiblespace}{2423}
\pdfglyphtounicode{visualspace}{2423}
\pdfglyphtounicode{vmonospace}{FF56}
\pdfglyphtounicode{voarmenian}{0578}
\pdfglyphtounicode{voicediterationhiragana}{309E}
\pdfglyphtounicode{voicediterationkatakana}{30FE}
\pdfglyphtounicode{voicedmarkkana}{309B}
\pdfglyphtounicode{voicedmarkkanahalfwidth}{FF9E}
\pdfglyphtounicode{vokatakana}{30FA}
\pdfglyphtounicode{vparen}{24B1}
\pdfglyphtounicode{vtilde}{1E7D}
\pdfglyphtounicode{vturned}{028C}
\pdfglyphtounicode{vuhiragana}{3094}
\pdfglyphtounicode{vukatakana}{30F4}
\pdfglyphtounicode{w}{0077}
\pdfglyphtounicode{wacute}{1E83}
\pdfglyphtounicode{waekorean}{3159}
\pdfglyphtounicode{wahiragana}{308F}
\pdfglyphtounicode{wakatakana}{30EF}
\pdfglyphtounicode{wakatakanahalfwidth}{FF9C}
\pdfglyphtounicode{wakorean}{3158}
\pdfglyphtounicode{wasmallhiragana}{308E}
\pdfglyphtounicode{wasmallkatakana}{30EE}
\pdfglyphtounicode{wattosquare}{3357}
\pdfglyphtounicode{wavedash}{301C}
\pdfglyphtounicode{wavyunderscorevertical}{FE34}
\pdfglyphtounicode{wawarabic}{0648}
\pdfglyphtounicode{wawfinalarabic}{FEEE}
\pdfglyphtounicode{wawhamzaabovearabic}{0624}
\pdfglyphtounicode{wawhamzaabovefinalarabic}{FE86}
\pdfglyphtounicode{wbsquare}{33DD}
\pdfglyphtounicode{wcircle}{24E6}
\pdfglyphtounicode{wcircumflex}{0175}
\pdfglyphtounicode{wdieresis}{1E85}
\pdfglyphtounicode{wdotaccent}{1E87}
\pdfglyphtounicode{wdotbelow}{1E89}
\pdfglyphtounicode{wehiragana}{3091}
\pdfglyphtounicode{weierstrass}{2118}
\pdfglyphtounicode{wekatakana}{30F1}
\pdfglyphtounicode{wekorean}{315E}
\pdfglyphtounicode{weokorean}{315D}
\pdfglyphtounicode{wgrave}{1E81}
\pdfglyphtounicode{whitebullet}{25E6}
\pdfglyphtounicode{whitecircle}{25CB}
\pdfglyphtounicode{whitecircleinverse}{25D9}
\pdfglyphtounicode{whitecornerbracketleft}{300E}
\pdfglyphtounicode{whitecornerbracketleftvertical}{FE43}
\pdfglyphtounicode{whitecornerbracketright}{300F}
\pdfglyphtounicode{whitecornerbracketrightvertical}{FE44}
\pdfglyphtounicode{whitediamond}{25C7}
\pdfglyphtounicode{whitediamondcontainingblacksmalldiamond}{25C8}
\pdfglyphtounicode{whitedownpointingsmalltriangle}{25BF}
\pdfglyphtounicode{whitedownpointingtriangle}{25BD}
\pdfglyphtounicode{whiteleftpointingsmalltriangle}{25C3}
\pdfglyphtounicode{whiteleftpointingtriangle}{25C1}
\pdfglyphtounicode{whitelenticularbracketleft}{3016}
\pdfglyphtounicode{whitelenticularbracketright}{3017}
\pdfglyphtounicode{whiterightpointingsmalltriangle}{25B9}
\pdfglyphtounicode{whiterightpointingtriangle}{25B7}
\pdfglyphtounicode{whitesmallsquare}{25AB}
\pdfglyphtounicode{whitesmilingface}{263A}
\pdfglyphtounicode{whitesquare}{25A1}
\pdfglyphtounicode{whitestar}{2606}
\pdfglyphtounicode{whitetelephone}{260F}
\pdfglyphtounicode{whitetortoiseshellbracketleft}{3018}
\pdfglyphtounicode{whitetortoiseshellbracketright}{3019}
\pdfglyphtounicode{whiteuppointingsmalltriangle}{25B5}
\pdfglyphtounicode{whiteuppointingtriangle}{25B3}
\pdfglyphtounicode{wihiragana}{3090}
\pdfglyphtounicode{wikatakana}{30F0}
\pdfglyphtounicode{wikorean}{315F}
\pdfglyphtounicode{wmonospace}{FF57}
\pdfglyphtounicode{wohiragana}{3092}
\pdfglyphtounicode{wokatakana}{30F2}
\pdfglyphtounicode{wokatakanahalfwidth}{FF66}
\pdfglyphtounicode{won}{20A9}
\pdfglyphtounicode{wonmonospace}{FFE6}
\pdfglyphtounicode{wowaenthai}{0E27}
\pdfglyphtounicode{wparen}{24B2}
\pdfglyphtounicode{wreathproduct}{2240}
\pdfglyphtounicode{wring}{1E98}
\pdfglyphtounicode{wsuperior}{02B7}
\pdfglyphtounicode{wturned}{028D}
\pdfglyphtounicode{wynn}{01BF}
\pdfglyphtounicode{x}{0078}
\pdfglyphtounicode{xabovecmb}{033D}
\pdfglyphtounicode{xbopomofo}{3112}
\pdfglyphtounicode{xcircle}{24E7}
\pdfglyphtounicode{xdieresis}{1E8D}
\pdfglyphtounicode{xdotaccent}{1E8B}
\pdfglyphtounicode{xeharmenian}{056D}
\pdfglyphtounicode{xi}{03BE}
\pdfglyphtounicode{xmonospace}{FF58}
\pdfglyphtounicode{xparen}{24B3}
\pdfglyphtounicode{xsuperior}{02E3}
\pdfglyphtounicode{y}{0079}
\pdfglyphtounicode{yaadosquare}{334E}
\pdfglyphtounicode{yabengali}{09AF}
\pdfglyphtounicode{yacute}{00FD}
\pdfglyphtounicode{yadeva}{092F}
\pdfglyphtounicode{yaekorean}{3152}
\pdfglyphtounicode{yagujarati}{0AAF}
\pdfglyphtounicode{yagurmukhi}{0A2F}
\pdfglyphtounicode{yahiragana}{3084}
\pdfglyphtounicode{yakatakana}{30E4}
\pdfglyphtounicode{yakatakanahalfwidth}{FF94}
\pdfglyphtounicode{yakorean}{3151}
\pdfglyphtounicode{yamakkanthai}{0E4E}
\pdfglyphtounicode{yasmallhiragana}{3083}
\pdfglyphtounicode{yasmallkatakana}{30E3}
\pdfglyphtounicode{yasmallkatakanahalfwidth}{FF6C}
\pdfglyphtounicode{yatcyrillic}{0463}
\pdfglyphtounicode{ycircle}{24E8}
\pdfglyphtounicode{ycircumflex}{0177}
\pdfglyphtounicode{ydieresis}{00FF}
\pdfglyphtounicode{ydotaccent}{1E8F}
\pdfglyphtounicode{ydotbelow}{1EF5}
\pdfglyphtounicode{yeharabic}{064A}
\pdfglyphtounicode{yehbarreearabic}{06D2}
\pdfglyphtounicode{yehbarreefinalarabic}{FBAF}
\pdfglyphtounicode{yehfinalarabic}{FEF2}
\pdfglyphtounicode{yehhamzaabovearabic}{0626}
\pdfglyphtounicode{yehhamzaabovefinalarabic}{FE8A}
\pdfglyphtounicode{yehhamzaaboveinitialarabic}{FE8B}
\pdfglyphtounicode{yehhamzaabovemedialarabic}{FE8C}
\pdfglyphtounicode{yehinitialarabic}{FEF3}
\pdfglyphtounicode{yehmedialarabic}{FEF4}
\pdfglyphtounicode{yehmeeminitialarabic}{FCDD}
\pdfglyphtounicode{yehmeemisolatedarabic}{FC58}
\pdfglyphtounicode{yehnoonfinalarabic}{FC94}
\pdfglyphtounicode{yehthreedotsbelowarabic}{06D1}
\pdfglyphtounicode{yekorean}{3156}
\pdfglyphtounicode{yen}{00A5}
\pdfglyphtounicode{yenmonospace}{FFE5}
\pdfglyphtounicode{yeokorean}{3155}
\pdfglyphtounicode{yeorinhieuhkorean}{3186}
\pdfglyphtounicode{yerahbenyomohebrew}{05AA}
\pdfglyphtounicode{yerahbenyomolefthebrew}{05AA}
\pdfglyphtounicode{yericyrillic}{044B}
\pdfglyphtounicode{yerudieresiscyrillic}{04F9}
\pdfglyphtounicode{yesieungkorean}{3181}
\pdfglyphtounicode{yesieungpansioskorean}{3183}
\pdfglyphtounicode{yesieungsioskorean}{3182}
\pdfglyphtounicode{yetivhebrew}{059A}
\pdfglyphtounicode{ygrave}{1EF3}
\pdfglyphtounicode{yhook}{01B4}
\pdfglyphtounicode{yhookabove}{1EF7}
\pdfglyphtounicode{yiarmenian}{0575}
\pdfglyphtounicode{yicyrillic}{0457}
\pdfglyphtounicode{yikorean}{3162}
\pdfglyphtounicode{yinyang}{262F}
\pdfglyphtounicode{yiwnarmenian}{0582}
\pdfglyphtounicode{ymonospace}{FF59}
\pdfglyphtounicode{yod}{05D9}
\pdfglyphtounicode{yoddagesh}{FB39}
\pdfglyphtounicode{yoddageshhebrew}{FB39}
\pdfglyphtounicode{yodhebrew}{05D9}
\pdfglyphtounicode{yodyodhebrew}{05F2}
\pdfglyphtounicode{yodyodpatahhebrew}{FB1F}
\pdfglyphtounicode{yohiragana}{3088}
\pdfglyphtounicode{yoikorean}{3189}
\pdfglyphtounicode{yokatakana}{30E8}
\pdfglyphtounicode{yokatakanahalfwidth}{FF96}
\pdfglyphtounicode{yokorean}{315B}
\pdfglyphtounicode{yosmallhiragana}{3087}
\pdfglyphtounicode{yosmallkatakana}{30E7}
\pdfglyphtounicode{yosmallkatakanahalfwidth}{FF6E}
\pdfglyphtounicode{yotgreek}{03F3}
\pdfglyphtounicode{yoyaekorean}{3188}
\pdfglyphtounicode{yoyakorean}{3187}
\pdfglyphtounicode{yoyakthai}{0E22}
\pdfglyphtounicode{yoyingthai}{0E0D}
\pdfglyphtounicode{yparen}{24B4}
\pdfglyphtounicode{ypogegrammeni}{037A}
\pdfglyphtounicode{ypogegrammenigreekcmb}{0345}
\pdfglyphtounicode{yr}{01A6}
\pdfglyphtounicode{yring}{1E99}
\pdfglyphtounicode{ysuperior}{02B8}
\pdfglyphtounicode{ytilde}{1EF9}
\pdfglyphtounicode{yturned}{028E}
\pdfglyphtounicode{yuhiragana}{3086}
\pdfglyphtounicode{yuikorean}{318C}
\pdfglyphtounicode{yukatakana}{30E6}
\pdfglyphtounicode{yukatakanahalfwidth}{FF95}
\pdfglyphtounicode{yukorean}{3160}
\pdfglyphtounicode{yusbigcyrillic}{046B}
\pdfglyphtounicode{yusbigiotifiedcyrillic}{046D}
\pdfglyphtounicode{yuslittlecyrillic}{0467}
\pdfglyphtounicode{yuslittleiotifiedcyrillic}{0469}
\pdfglyphtounicode{yusmallhiragana}{3085}
\pdfglyphtounicode{yusmallkatakana}{30E5}
\pdfglyphtounicode{yusmallkatakanahalfwidth}{FF6D}
\pdfglyphtounicode{yuyekorean}{318B}
\pdfglyphtounicode{yuyeokorean}{318A}
\pdfglyphtounicode{yyabengali}{09DF}
\pdfglyphtounicode{yyadeva}{095F}
\pdfglyphtounicode{z}{007A}
\pdfglyphtounicode{zaarmenian}{0566}
\pdfglyphtounicode{zacute}{017A}
\pdfglyphtounicode{zadeva}{095B}
\pdfglyphtounicode{zagurmukhi}{0A5B}
\pdfglyphtounicode{zaharabic}{0638}
\pdfglyphtounicode{zahfinalarabic}{FEC6}
\pdfglyphtounicode{zahinitialarabic}{FEC7}
\pdfglyphtounicode{zahiragana}{3056}
\pdfglyphtounicode{zahmedialarabic}{FEC8}
\pdfglyphtounicode{zainarabic}{0632}
\pdfglyphtounicode{zainfinalarabic}{FEB0}
\pdfglyphtounicode{zakatakana}{30B6}
\pdfglyphtounicode{zaqefgadolhebrew}{0595}
\pdfglyphtounicode{zaqefqatanhebrew}{0594}
\pdfglyphtounicode{zarqahebrew}{0598}
\pdfglyphtounicode{zayin}{05D6}
\pdfglyphtounicode{zayindagesh}{FB36}
\pdfglyphtounicode{zayindageshhebrew}{FB36}
\pdfglyphtounicode{zayinhebrew}{05D6}
\pdfglyphtounicode{zbopomofo}{3117}
\pdfglyphtounicode{zcaron}{017E}
\pdfglyphtounicode{zcircle}{24E9}
\pdfglyphtounicode{zcircumflex}{1E91}
\pdfglyphtounicode{zcurl}{0291}
\pdfglyphtounicode{zdot}{017C}
\pdfglyphtounicode{zdotaccent}{017C}
\pdfglyphtounicode{zdotbelow}{1E93}
\pdfglyphtounicode{zecyrillic}{0437}
\pdfglyphtounicode{zedescendercyrillic}{0499}
\pdfglyphtounicode{zedieresiscyrillic}{04DF}
\pdfglyphtounicode{zehiragana}{305C}
\pdfglyphtounicode{zekatakana}{30BC}
\pdfglyphtounicode{zero}{0030}
\pdfglyphtounicode{zeroarabic}{0660}
\pdfglyphtounicode{zerobengali}{09E6}
\pdfglyphtounicode{zerodeva}{0966}
\pdfglyphtounicode{zerogujarati}{0AE6}
\pdfglyphtounicode{zerogurmukhi}{0A66}
\pdfglyphtounicode{zerohackarabic}{0660}
\pdfglyphtounicode{zeroinferior}{2080}
\pdfglyphtounicode{zeromonospace}{FF10}
\pdfglyphtounicode{zerooldstyle}{0030}
\pdfglyphtounicode{zeropersian}{06F0}
\pdfglyphtounicode{zerosuperior}{2070}
\pdfglyphtounicode{zerothai}{0E50}
\pdfglyphtounicode{zerowidthjoiner}{FEFF}
\pdfglyphtounicode{zerowidthnonjoiner}{200C}
\pdfglyphtounicode{zerowidthspace}{200B}
\pdfglyphtounicode{zeta}{03B6}
\pdfglyphtounicode{zhbopomofo}{3113}
\pdfglyphtounicode{zhearmenian}{056A}
\pdfglyphtounicode{zhebrevecyrillic}{04C2}
\pdfglyphtounicode{zhecyrillic}{0436}
\pdfglyphtounicode{zhedescendercyrillic}{0497}
\pdfglyphtounicode{zhedieresiscyrillic}{04DD}
\pdfglyphtounicode{zihiragana}{3058}
\pdfglyphtounicode{zikatakana}{30B8}
\pdfglyphtounicode{zinorhebrew}{05AE}
\pdfglyphtounicode{zlinebelow}{1E95}
\pdfglyphtounicode{zmonospace}{FF5A}
\pdfglyphtounicode{zohiragana}{305E}
\pdfglyphtounicode{zokatakana}{30BE}
\pdfglyphtounicode{zparen}{24B5}
\pdfglyphtounicode{zretroflexhook}{0290}
\pdfglyphtounicode{zstroke}{01B6}
\pdfglyphtounicode{zuhiragana}{305A}
\pdfglyphtounicode{zukatakana}{30BA}
 \pdfgentounicode=1
\usepackage[english,french,shorthands=off]{babel}
\usepackage{newtxtext}
\usepackage[protrusion=true,expansion=true,tracking=true]{microtype}
\usepackage[a4paper,inner=2.72cm,outer=2.34cm,top=2.42cm,bottom=2.68cm,
  headheight=31pt,headsep=18pt,footskip=30pt]{geometry}
\usepackage{amsmath,amsthm,amscd,mathtools,mathrsfs,bm}
\usepackage{newtxmath}
% Las jerarquías de títulos, el índice y las tablas emplean escalas
% intermedias. Se mantiene la familia RSFS y se habilita su escalado continuo,
% de modo que la notación matemática conserve el cuerpo tipográfico exacto
% sin sustituciones silenciosas de tamaño.
\DeclareFontFamily{U}{rsfs}{\skewchar\font127}
\DeclareFontShape{U}{rsfs}{m}{n}{%
  <-6> s*[1.0] rsfs5
  <6-8> s*[1.0] rsfs7
  <8-> s*[1.0] rsfs10
}{}
\usepackage{array,booktabs,longtable,tabularx,multirow}
\usepackage{enumitem,graphicx,float,pdfpages,pdflscape}
\usepackage{subcaption}
\usepackage[table]{xcolor}
\pdfinclusioncopyfonts=1
\usepackage{tikz}
\usepackage{tikz-cd}
\usetikzlibrary{arrows.meta,positioning,fit,calc,matrix,shapes.geometric,
  decorations.pathreplacing,decorations.markings,backgrounds,patterns}
\usepackage[most]{tcolorbox}
\usepackage{setspace,csquotes,xurl,fancyhdr,fancyvrb,listings,tocloft,titlesec,needspace,etoolbox}
\usepackage{xstring}
\usepackage{hyperref}
% Los marcadores PDF son parte de la jerarquía científica de la obra.  Esta
% tabla conserva en texto Unicode la notación que aparece en los títulos, en
% lugar de dejar que hyperref suprima silenciosamente flechas y constantes.
\pdfstringdefDisableCommands{%
  \def\({}\def\){}\def\[{}\def\]{}%
  \def\ensuremath#1{#1}%
  \def\mathrm#1{#1}\def\mathbf#1{#1}\def\mathsf#1{#1}%
  \def\mathbb#1{#1}\def\mathcal#1{#1}\def\mathfrak#1{#1}%
  \def\operatorname#1{#1}\def\text#1{#1}\def\bm#1{#1}%
  \def\allowbreak{}%
  \def\frac#1#2{#1/#2}\def\sqrt#1{racine(#1)}%
  \def\to{→}\def\longrightarrow{→}\def\mapsto{↦}%
  \def\leftrightarrow{↔}\def\Rightarrow{⇒}\def\Longrightarrow{⇒}%
  \def\pi{π}\def\varphi{φ}\def\phi{φ}\def\alpha{α}%
  \def\varepsilon{ε}\def\epsilon{ε}\def\mu{μ}\def\lambda{λ}%
  \def\Omega{Ω}\def\zeta{ζ}\def\partial{∂}%
  \def\ast{*}\def\cdot{·}\def\natural{natural}\def\log{log}%
}
\usepackage[nameinlink,noabbrev]{cleveref}
\crefname{appendix}{annexe}{annexes}
\Crefname{appendix}{Annexe}{Annexes}
% Referencias descriptivas a resultados del tratado que se transducen en esta
% monografía bajo una numeración propia.  La sustitución evita remitir al lector
% a la numeración editorial de otra obra sin duplicar los resultados.
\makeatletter
\newcommand{\HMTDeclareDescriptor}[2]{%
  \expandafter\gdef\csname hmt@xref@#1\endcsname{#2}}
\newcommand{\HMTIfDescriptorTF}[3]{%
  \ifcsname hmt@xref@#1\endcsname #2\else #3\fi}

\HMTDeclareDescriptor{app:catalogo-semillas-app-rev9}{chapitre consacré au domaine des germes APP et au catalogue canonique des émissions}
\HMTDeclareDescriptor{chap:83-08}{chapitre consacré à la décomposition généalogique des cinq régions orientées de \(\pi\)}
\HMTDeclareDescriptor{chap:alpha}{chapitre consacré à la clôture dodécaphasique de \(\alpha\)}
\HMTDeclareDescriptor{chap:banda-particula-virtual}{chapitre consacré à la bande de particule, à l’antisection et à la virtualité}
\HMTDeclareDescriptor{chap:bandera-excepcional-acumulativa}{chapitre consacré à l’incidence dodécaphasique et au corridor Witt--Leech--Monstre}
\HMTDeclareDescriptor{chap:biografia-pi-autonoma}{chapitre consacré à la construction du caractère de clôture et à la biographie structurelle de \(\pi\)}
\HMTDeclareDescriptor{chap:cinco-ciclos-pi}{chapitre consacré à la minimalité des clôtures des cinq régions de \(\pi\)}
\HMTDeclareDescriptor{chap:funtor-dimensional-hmt}{chapitre consacré aux droites de grandeur et à la réalisation dimensionnelle du noyau}
\HMTDeclareDescriptor{chap:integral-genealogica}{chapitre consacré à l’algèbre généalogique des chemins et à sa représentation intégrale}
\HMTDeclareDescriptor{chap:numero-medicion}{chapitre consacré aux systèmes cylindriques, à la survie et à la construction du nombre}
\HMTDeclareDescriptor{chap:orden-determinantal-accion}{chapitre consacré à l’ordre déterminantal de la droite d’action}
\HMTDeclareDescriptor{chap:pantallas-dualidad}{chapitre consacré aux sous-espaces, aux quotients, aux réductions de rang et à la dualité}
\HMTDeclareDescriptor{chap:semillas-emisor-54}{chapitre consacré à l’émetteur d’indice \(54\) et à ses germes orientés}
\HMTDeclareDescriptor{chap:sistema-inverso-continuidad-genealogica}{chapitre consacré au système inverse et à la continuité généalogique}
\HMTDeclareDescriptor{chap:tpk-definicion}{chapitre consacré au système de transport et à l’évolution de phase}
\HMTDeclareDescriptor{def:banda-dodecafasica-ruta}{définition de la bande dodécaphasique de route}
\HMTDeclareDescriptor{def:odometro-9-12}{définition de l’odomètre nonadique--dodécaphasique}
\HMTDeclareDescriptor{eq:G-interno}{identité constitutive interne de la gravitation}
\HMTDeclareDescriptor{eq:consistencia-proyectiva-medidas-u024}{identité de cohérence projective des mesures cylindriques}
\HMTDeclareDescriptor{eq:continuo-lector-arquimediano}{identité de couverture du lecteur archimédien}
\HMTDeclareDescriptor{eq:descomposiciones-cinco-regiones-pi}{décomposition des cinq régions de \(\pi\)}
\HMTDeclareDescriptor{eq:frontal-doble-proyeccion}{diagramme frontal de la double projection}
\HMTDeclareDescriptor{eq:modular-finite-first-law}{première identité modulaire finie avec terme de défaut}
\HMTDeclareDescriptor{eq:modular-inverse-channels}{identité de reconstruction conjointe des canaux d’aire}
\HMTDeclareDescriptor{eq:realizacion-circular-ct108}{identité de réalisation circulaire du calendrier à \(108\) phases}
\HMTDeclareDescriptor{espiral:chap:realizacion-em}{chapitre consacré à la réalisation électromagnétique restreinte}
\HMTDeclareDescriptor{espiral:eq:cociclos-radiales}{identité des cocycles radiaux du relèvement tritique}
\HMTDeclareDescriptor{hap:def:entropia-bicapa-rev6}{définition de l’entropie bicouche orientée}
\HMTDeclareDescriptor{hmtctx:p12:chap:orbitas-elevacion-r36}{chapitre consacré aux orbites et au relèvement au calibre \(R_{36}\)}
\HMTDeclareDescriptor{mdg:chap:torres}{chapitre consacré à la hiérarchie globale des tours dimensionnelles}
\HMTDeclareDescriptor{mdg:sec:cierre-electron-two-way-rev11}{section consacrée à la clôture électronique bidirectionnelle}
\HMTDeclareDescriptor{mdg:sec:ley-general-masa-accion-autonoma}{section consacrée à l’opérateur universel d’action}
\HMTDeclareDescriptor{mdg:sec:selector-gravitatorio-rev11}{section consacrée au sélecteur gravitationnel naturel}
\HMTDeclareDescriptor{p12-83-c19-s1:hmtch:eq-celula-trit-completa}{identité de la cellule tritique complète}
\HMTDeclareDescriptor{prop:cierre-antiseccion-ct108}{proposition de clôture des formes du calendrier à \(108\) phases}
\HMTDeclareDescriptor{prop:proyeccion-fase-longitud}{proposition de projection phase--longueur}
\HMTDeclareDescriptor{rm:ch:cosmo}{chapitre consacré aux échelles cosmologiques du calendrier à \(108\) phases}
\HMTDeclareDescriptor{rm:ch:metrology}{chapitre consacré aux constantes quantiques--électriques, à l’échelle thermique et au système de Planck}
\HMTDeclareDescriptor{sec:tres-generadores-concretos-trit}{section consacrée aux trois générateurs concrets des régimes tritiques}
\HMTDeclareDescriptor{sec:w12-w24-elevacion}{section consacrée au plan résiduel dodécaphasique et au relèvement des hexades}
\HMTDeclareDescriptor{thm:censo-atlas-81}{théorème de classification des familles et générations de l’atlas nonadique}
\HMTDeclareDescriptor{thm:corredor-hmt-moonshine}{théorème du corridor de Leech à l’algèbre de Moonshine}
\HMTDeclareDescriptor{thm:dos-realizaciones-fuente-cuadratica-rev3}{théorème des deux réalisations d’une présentation quadratique entière}
\HMTDeclareDescriptor{thm:generacion-monodromica-conjunta-rev7}{théorème de filtration conjointe à profondeur finie}
\HMTDeclareDescriptor{thm:homogeneidad-pch}{théorème d’homogénéité de l’action de Palatini--Cartan--Holst}
\HMTDeclareDescriptor{thm:principio-conjugacion-masa-banda}{théorème de conjugaison de masse et de caractère de bande}
\HMTDeclareDescriptor{thm:tpk-cociente-predictivo-36}{théorème du quotient minimal de mémoire prédictive}

% Copias semánticas, no simples alias: \cref y \Cref son órdenes robustas
% gestionadas por xparse y un \let conservaría el despachador mutable,
% provocando recursión después de \RenewDocumentCommand.
\NewCommandCopy{\HMTOriginalCref}{\cref}
\NewCommandCopy{\HMTOriginalCrefCapital}{\Cref}
\NewCommandCopy{\HMTOriginalEqref}{\eqref}
% Los capítulos transducidos conservan centenares de referencias semánticas
% cuya numeración pertenecía al tratado integral.  Una referencia que ya ha
% recibido etiqueta local conserva su numeración propia; una referencia que
% sólo existe en el volumen rector se expresa por su clase matemática.  De
% este modo el texto sigue siendo legible y autosuficiente, sin signos ``??''
% ni remisiones ficticias a una numeración ausente.
\newcommand{\HMTFallbackReference}[1]{%
  \IfSubStr{#1}{eq:}{l’identité correspondante}{%
  \IfSubStr{#1}{thm:}{le théorème correspondant}{%
  \IfSubStr{#1}{theorem:}{le théorème correspondant}{%
  \IfSubStr{#1}{prop:}{la proposition correspondante}{%
  \IfSubStr{#1}{lem:}{le lemme correspondant}{%
  \IfSubStr{#1}{cor:}{le corollaire correspondant}{%
  \IfSubStr{#1}{def:}{la définition correspondante}{%
  \IfSubStr{#1}{fig:}{la figure correspondante}{%
  \IfSubStr{#1}{tab:}{le tableau correspondant}{%
  \IfSubStr{#1}{chap:}{le chapitre correspondant}{%
  \IfSubStr{#1}{ch:}{le chapitre correspondant}{%
  \IfSubStr{#1}{sec:}{la section correspondante}{%
  \IfSubStr{#1}{crit:}{le critère correspondant}{%
  le résultat correspondant}}}}}}}}}}}}}}
\newcommand{\HMTFallbackReferenceCapital}[1]{%
  \IfSubStr{#1}{eq:}{L’identité correspondante}{%
  \IfSubStr{#1}{thm:}{Le théorème correspondant}{%
  \IfSubStr{#1}{theorem:}{Le théorème correspondant}{%
  \IfSubStr{#1}{prop:}{La proposition correspondante}{%
  \IfSubStr{#1}{lem:}{Le lemme correspondant}{%
  \IfSubStr{#1}{cor:}{Le corollaire correspondant}{%
  \IfSubStr{#1}{def:}{La définition correspondante}{%
  \IfSubStr{#1}{fig:}{La figure correspondante}{%
  \IfSubStr{#1}{tab:}{Le tableau correspondant}{%
  \IfSubStr{#1}{chap:}{Le chapitre correspondant}{%
  \IfSubStr{#1}{ch:}{Le chapitre correspondant}{%
  \IfSubStr{#1}{sec:}{La section correspondante}{%
  \IfSubStr{#1}{crit:}{Le critère correspondant}{%
  Le résultat correspondant}}}}}}}}}}}}}}
\newcommand{\HMTIfLocalLabelTF}[3]{%
  \ifcsname r@#1\endcsname #2\else #3\fi}
\RenewDocumentCommand{\cref}{s m}{%
  \HMTIfDescriptorTF{#2}{\csname hmt@xref@#2\endcsname}{%
    \HMTIfLocalLabelTF{#2}{%
      \IfBooleanTF{#1}{\HMTOriginalCref*{#2}}{\HMTOriginalCref{#2}}%
    }{\HMTFallbackReference{#2}}}}
\RenewDocumentCommand{\Cref}{s m}{%
  \HMTIfDescriptorTF{#2}{\csname hmt@xref@#2\endcsname}{%
    \HMTIfLocalLabelTF{#2}{%
      \IfBooleanTF{#1}{\HMTOriginalCrefCapital*{#2}}{\HMTOriginalCrefCapital{#2}}%
    }{\HMTFallbackReferenceCapital{#2}}}}
\RenewDocumentCommand{\eqref}{m}{%
  \HMTIfDescriptorTF{#1}{\textup{(\csname hmt@xref@#1\endcsname)}}{%
    \HMTIfLocalLabelTF{#1}{\HMTOriginalEqref{#1}}{\textup{(identité correspondante)}}}}
\makeatother
 \usepackage[noautomatic,quiet]{imakeidx}
\makeindex
\allowdisplaybreaks

\definecolor{hmtblue}{RGB}{24,57,89}
\definecolor{hmtgold}{RGB}{145,101,26}
\definecolor{hmtgreen}{RGB}{29,104,78}
\definecolor{hmtred}{RGB}{151,54,45}
\definecolor{hmtviolet}{RGB}{92,64,125}
\definecolor{hmtgray}{RGB}{87,95,102}
\definecolor{hmtpaper}{RGB}{248,247,243}
\definecolor{hmtlightblue}{RGB}{232,240,247}
\definecolor{hmtlightgold}{RGB}{249,242,225}
\definecolor{hmtlightgreen}{RGB}{232,244,238}
\definecolor{hmtlightred}{RGB}{250,235,233}
\definecolor{hmtlightviolet}{RGB}{241,236,248}

% Paleta pública común para los cuadros, el atlas y los inventarios
% integrales de Mecánica Dimensional. Los nombres en versales son alias
% tipográficos; no introducen una segunda identidad visual.
\colorlet{HMTNavy}{hmtblue}
\colorlet{HMTBlue}{hmtblue}
\colorlet{HMTGold}{hmtgold}
\colorlet{HMTLight}{hmtlightblue}
\colorlet{HMTGray}{hmtgray}
\colorlet{HMTLine}{hmtlightblue}
\providecommand{\HMTtablehead}[1]{\color{white}\sffamily\bfseries #1}
\providecommand{\RaggedRight}{\raggedright}

\newcommand{\HMT}{\ensuremath{\mathrm{HMT}}}
\providecommand{\gammaH}{\gamma_{\HMT}}
\providecommand{\ii}{\mathrm i}
\providecommand{\C}{\mathbb C}
\providecommand{\Z}{\mathbb Z}
\providecommand{\U}{\mathcal U}
\providecommand{\arsinh}{\operatorname{arsinh}}
\providecommand{\arcosh}{\operatorname{arcosh}}
\providecommand{\HHM}{K_{\mathrm{HHM}}}
\providecommand{\Dmod}{\Delta_{4}^{\mathrm{mod}}}
\providecommand{\Val}{\operatorname{Val}}
\providecommand{\Ext}{\operatorname{Ext}}
\providecommand{\ev}{\operatorname{ev}}
\providecommand{\End}{\operatorname{End}}
\providecommand{\diam}{\operatorname{diam}}
\providecommand{\im}{\operatorname{im}}
\providecommand{\supp}{\operatorname{supp}}
\providecommand{\tr}{\operatorname{tr}}
\providecommand{\Gnine}{\mathcal G_9}
\providecommand{\NTHW}{\mathfrak N_{\mathrm{HW}}}
\providecommand{\readerquestion}[1]{%
  \begin{quote}\small\itshape\color{hmtblue}#1\end{quote}}
\newcommand{\TRIT}{\ensuremath{\widehat{\mathbb T}}}
\newcommand{\AAPP}{\mathcal A_9}
\newcommand{\AAPPlift}{\mathcal A_9^{\#}}
\newcommand{\GraphAPP}{\Gamma_{\mathrm{APP},9}}
\newcommand{\MonNine}{\mathscr M_9}
\newcommand{\XTPK}{\mathcal X_{\mathrm{TPK}}}
\newcommand{\CellRQ}{\mathcal C_{9}^{\mathrm{rq}}}
\newcommand{\WordMonoid}{\mathfrak W_{9}^{\times}}
\newcommand{\Cdoce}{C_{12}}
\newcommand{\Kseal}{K}
\newcommand{\Usgn}{U_{12}^{\mathrm{sgn}}}
\newcommand{\Etrans}{\mathcal E_{3,4}}
\newcommand{\Rpi}{\mathscr R_{\pi}}
\newcommand{\Reu}{\mathscr R_{e}}
\newcommand{\Rphi}{\mathscr R_{\varphi}}
\newcommand{\RR}{\mathbb R}
\newcommand{\ZZ}{\mathbb Z}
\newcommand{\QQ}{\mathbb Q}
\newcommand{\FF}{\mathbb F}
\newcommand{\CC}{\mathbb C}
\newcommand{\dd}{\,\mathrm d}
\newcommand{\ee}{\mathrm e}
\newcommand{\Aut}{\operatorname{Aut}}
\newcommand{\Hol}{\operatorname{Hol}}
\newcommand{\dr}{\operatorname{dr}_9}
\newcommand{\Surv}{\operatorname{Surv}}
\newcommand{\val}{\operatorname{val}}
\newcommand{\Dnine}{D_9}
\newcommand{\Wnine}{\mathscr W_9}
\newcommand{\Krad}{\mathcal K_{\#}^{+}}
\newcommand{\Hplus}{\mathcal H_{+}}
\newcommand{\Dom}{\operatorname{Dom}}
\newcommand{\Tr}{\operatorname{Tr}}
\newcommand{\Spec}{\operatorname{Spec}}
\providecommand{\Sha}{\operatorname{Sha}}
\newcommand{\Irr}{\mathfrak I}
\newcommand{\NN}{\mathbb N}
\newcommand{\N}{\mathbb N}
\newcommand{\R}{\mathbb R}
\providecommand{\Adeg}{A_{\rm deg}}
\providecommand{\Arad}{A_{\rm rad}}
\providecommand{\Cdeg}{C^{*}_{\rm deg}}
\providecommand{\Crad}{C^{*}_{\rm rad}}
\providecommand{\etaRetrad}{\eta_{\rm ret}^{(\rm rad)}}
\providecommand{\Gtr}{\ensuremath{\Gamma_{\rm tr}}}
\providecommand{\HTr}{\ensuremath{\mathsf{Tr}_{q}}}
\providecommand{\HAdd}{\ensuremath{\mathsf{Add}}}
\providecommand{\HDef}{\ensuremath{\mathsf{Def}}}
\providecommand{\HRes}{\ensuremath{\mathsf{Res}}}
\providecommand{\HLoc}{\ensuremath{\mathsf{Loc}}}

\theoremstyle{plain}
\newtheorem{theorem}{Théorème}[chapter]
\newtheorem{lemma}[theorem]{Lemme}
\newtheorem{proposition}[theorem]{Proposition}
\newtheorem{corollary}[theorem]{Corollaire}
\theoremstyle{definition}
\newtheorem{definition}[theorem]{Définition}
\newtheorem{construction}[theorem]{Construction}
\newtheorem{algorithm}[theorem]{Algorithme}
\newtheorem{ablation}[theorem]{Proposition de nécessité structurelle}
\newtheorem{criterion}[theorem]{Critère de comparaison}
\theoremstyle{remark}
\newtheorem{remark}[theorem]{Remarque}

\tcbset{enhanced,breakable,boxrule=.45pt,arc=0mm,left=2.5mm,right=2.5mm,top=1.8mm,bottom=1.8mm}
\newtcolorbox{exactbox}[1][]{colback=white,colframe=black,title={Résultat exact},fonttitle=\bfseries,#1}

\newcolumntype{Y}{>{\raggedright\arraybackslash}X}
\newcolumntype{R}{>{\raggedleft\arraybackslash}X}
\newcolumntype{L}[1]{>{\raggedright\arraybackslash}p{#1}}
\newcolumntype{C}{>{\centering\arraybackslash}p{1.65cm}}

\setlength{\emergencystretch}{3em}
 % Emisión de afirmaciones vivas durante la compilación de REV.9.
% Una marca sólo cuenta si la rama TeX que la contiene se ejecuta.
\newwrite\HMTClaimsStream
\immediate\openout\HMTClaimsStream=\jobname.hmtclaims
\newcommand{\HMTClaim}[1]{%
  \immediate\write\HMTClaimsStream{HMT-CLAIM:#1}%
}
\AtEndDocument{\immediate\closeout\HMTClaimsStream}
 % Compatibilidad tipada de las tres aportaciones incorporadas en la edición
% sucesora de 2026-08-12. Este archivo no redefine objetos de REV.3: aporta
% únicamente nombres, entornos y operadores que los cuadernos autónomos
% necesitan para residir en el grafo único del tratado.

\RequirePackage{xspace}

\definecolor{hmtteal}{RGB}{10,107,112}

\providecommand{\Rhyper}{\mathcal R_{\mathrm{hyp}}}
\providecommand{\OmegaInf}{\Omega_{\infty}}
\providecommand{\Clop}{\operatorname{Clop}}
\providecommand{\Closed}{\operatorname{Closed}}
\providecommand{\Proj}{\operatorname{Proj}}
\providecommand{\id}{\operatorname{id}}
\providecommand{\coker}{\operatorname{coker}}
\providecommand{\ord}{\operatorname{ord}}
\providecommand{\HALT}{\mathsf{HALT}}
\providecommand{\Prob}{\mathcal P}
\providecommand{\status}[1]{\emph{#1}}
\providecommand{\citep}{\cite}
\providecommand{\QMS}{\ensuremath{\mathrm{QMS}}\xspace}
\providecommand{\QLS}{\ensuremath{\mathrm{QLS}}\xspace}
\providecommand{\QPM}{\ensuremath{\mathrm{QPM}}\xspace}
\providecommand{\UCP}{\ensuremath{\mathrm{UCP}}\xspace}
\providecommand{\MUB}{\ensuremath{\mathrm{MUB}}\xspace}
\providecommand{\PVM}{\ensuremath{\mathrm{PVM}}\xspace}
\providecommand{\POVM}{\ensuremath{\mathrm{POVM}}\xspace}
\providecommand{\Mat}{\operatorname{Mat}}
\providecommand{\Her}{\operatorname{Her}}
\providecommand{\mconv}{\operatorname{mconv}}
\providecommand{\Span}{\operatorname{span}}
\providecommand{\Ad}{\operatorname{Ad}}
\providecommand{\diag}{\operatorname{diag}}

\theoremstyle{definition}
\newtheorem{example}[theorem]{Exemple}

\newtcolorbox{resultbox}[1][]{%
  colback=hmtlightgreen,colframe=hmtgreen,
  title={Résultat},fonttitle=\bfseries,#1}
\newtcolorbox{controlbox}[1][]{%
  colback=white,colframe=hmtgold!75!black,
  title={Obstruction structurelle},fonttitle=\bfseries,#1}
\newtcolorbox{sourcebox}[1][]{%
  colback=white,colframe=hmtblue!55,
  title={Prémisses et dérivation},fonttitle=\bfseries,#1}
\newtcolorbox{intuitionbox}[1][]{%
  colback=hmtlightgold,colframe=hmtgold,
  title={Lecture scientifique},fonttitle=\bfseries,#1}
\newtcolorbox{workbox}[1][]{%
  colback=white,colframe=hmtgray,
  title={Extensions mathématiques},fonttitle=\bfseries,#1}
\newtcolorbox{importbox}[1][]{%
  colback=white,colframe=hmtblue!55,
  title={Théorème antécédent},fonttitle=\bfseries,#1}
\newtcolorbox{bridgebox}[1][]{%
  colback=hmtlightgold,colframe=hmtgold,
  title={Résultat du pont matriciel},fonttitle=\bfseries,#1}
 % Estado material V3 comentario-only: conserva los 219 resultados de V5 y
% registra las dos correcciones editoriales autorizadas del puente matricial.
% Los sidecars 218/219 previos quedan como testigos inactivos e inmutables.
% HMT--MD REV.4: sidecar comentario-only de continuidad material V3/219.
% Encadena V2, conserva 219 IDs y no publica teoremas, títulos ni contenido paginado.
% Las correcciones registrales V3 no alteran enunciados ni fuerzas probatorias.
% HMT_RESULT_BEGIN:R3-001:residence
% {"material_control":"REV3_BASE_GRANULAR+PASS_NO_PERDIDA_REV3","material_state":"CONSERVADO_POR_BASE_GRANULAR_REV3","residences":[{"path":"/Users/ruben/Documents/HMT2/CONTINUIDAD_EDITORIAL/INTEGRACION_REV3_VON_NEUMANN_MATRICIAL_REALIZACION_2026-08-12/REV3_BASE_GRANULAR.json","sha256":"d89983feac6e064f0ce77d6506d60b94dd1eaf00ed174352d3e32bfc5a9b88d4"}],"result_id":"R3-001","statement_sha256":"ec22d0aed84b1a5ccd6d8675f494a2d54bb746ae05ca67325087c42e7697061d"}
% HMT_RESULT_END:R3-001:residence
% HMT_RESULT_BEGIN:R3-002:residence
% {"material_control":"REV3_BASE_GRANULAR+PASS_NO_PERDIDA_REV3","material_state":"CONSERVADO_POR_BASE_GRANULAR_REV3","residences":[{"path":"/Users/ruben/Documents/HMT2/CONTINUIDAD_EDITORIAL/INTEGRACION_REV3_VON_NEUMANN_MATRICIAL_REALIZACION_2026-08-12/REV3_BASE_GRANULAR.json","sha256":"d89983feac6e064f0ce77d6506d60b94dd1eaf00ed174352d3e32bfc5a9b88d4"}],"result_id":"R3-002","statement_sha256":"94b5e013944e5d31216adcca582bf4469c088ecd478d1b65dccc1a3343942490"}
% HMT_RESULT_END:R3-002:residence
% HMT_RESULT_BEGIN:R3-003:residence
% {"material_control":"REV3_BASE_GRANULAR+PASS_NO_PERDIDA_REV3","material_state":"CONSERVADO_POR_BASE_GRANULAR_REV3","residences":[{"path":"/Users/ruben/Documents/HMT2/CONTINUIDAD_EDITORIAL/INTEGRACION_REV3_VON_NEUMANN_MATRICIAL_REALIZACION_2026-08-12/REV3_BASE_GRANULAR.json","sha256":"d89983feac6e064f0ce77d6506d60b94dd1eaf00ed174352d3e32bfc5a9b88d4"}],"result_id":"R3-003","statement_sha256":"1c2e8dc0a9e7ba35279bfc5f48df699c614e4d96af09dcdaf1b76a665374c7cc"}
% HMT_RESULT_END:R3-003:residence
% HMT_RESULT_BEGIN:R3-004:residence
% {"material_control":"REV3_BASE_GRANULAR+PASS_NO_PERDIDA_REV3","material_state":"CONSERVADO_POR_BASE_GRANULAR_REV3","residences":[{"path":"/Users/ruben/Documents/HMT2/CONTINUIDAD_EDITORIAL/INTEGRACION_REV3_VON_NEUMANN_MATRICIAL_REALIZACION_2026-08-12/REV3_BASE_GRANULAR.json","sha256":"d89983feac6e064f0ce77d6506d60b94dd1eaf00ed174352d3e32bfc5a9b88d4"}],"result_id":"R3-004","statement_sha256":"067caf5027a3ca629eaa3f299767eebe36ab8da36d4a594f0490a6fcba2959a3"}
% HMT_RESULT_END:R3-004:residence
% HMT_RESULT_BEGIN:R3-005:residence
% {"material_control":"REV3_BASE_GRANULAR+PASS_NO_PERDIDA_REV3","material_state":"CONSERVADO_POR_BASE_GRANULAR_REV3","residences":[{"path":"/Users/ruben/Documents/HMT2/CONTINUIDAD_EDITORIAL/INTEGRACION_REV3_VON_NEUMANN_MATRICIAL_REALIZACION_2026-08-12/REV3_BASE_GRANULAR.json","sha256":"d89983feac6e064f0ce77d6506d60b94dd1eaf00ed174352d3e32bfc5a9b88d4"}],"result_id":"R3-005","statement_sha256":"7f015f268d0e7772da1a33238cd19903733ee7cfa8947c6372241210ebcd0f01"}
% HMT_RESULT_END:R3-005:residence
% HMT_RESULT_BEGIN:R3-006:residence
% {"material_control":"REV3_BASE_GRANULAR+PASS_NO_PERDIDA_REV3","material_state":"CONSERVADO_POR_BASE_GRANULAR_REV3","residences":[{"path":"/Users/ruben/Documents/HMT2/CONTINUIDAD_EDITORIAL/INTEGRACION_REV3_VON_NEUMANN_MATRICIAL_REALIZACION_2026-08-12/REV3_BASE_GRANULAR.json","sha256":"d89983feac6e064f0ce77d6506d60b94dd1eaf00ed174352d3e32bfc5a9b88d4"}],"result_id":"R3-006","statement_sha256":"9e3f5af6f313c9c4c256d7e6a1f22a0bc0d112d95a7909a4b61b35638a89ea73"}
% HMT_RESULT_END:R3-006:residence
% HMT_RESULT_BEGIN:R3-007:residence
% {"material_control":"REV3_BASE_GRANULAR+PASS_NO_PERDIDA_REV3","material_state":"CONSERVADO_POR_BASE_GRANULAR_REV3","residences":[{"path":"/Users/ruben/Documents/HMT2/CONTINUIDAD_EDITORIAL/INTEGRACION_REV3_VON_NEUMANN_MATRICIAL_REALIZACION_2026-08-12/REV3_BASE_GRANULAR.json","sha256":"d89983feac6e064f0ce77d6506d60b94dd1eaf00ed174352d3e32bfc5a9b88d4"}],"result_id":"R3-007","statement_sha256":"afe6a7f442b4b9b73f09e57dcf993019f691bf05fabf3cf8a866c0cb30e0b20e"}
% HMT_RESULT_END:R3-007:residence
% HMT_RESULT_BEGIN:R3-008:residence
% {"material_control":"REV3_BASE_GRANULAR+PASS_NO_PERDIDA_REV3","material_state":"CONSERVADO_POR_BASE_GRANULAR_REV3","residences":[{"path":"/Users/ruben/Documents/HMT2/CONTINUIDAD_EDITORIAL/INTEGRACION_REV3_VON_NEUMANN_MATRICIAL_REALIZACION_2026-08-12/REV3_BASE_GRANULAR.json","sha256":"d89983feac6e064f0ce77d6506d60b94dd1eaf00ed174352d3e32bfc5a9b88d4"}],"result_id":"R3-008","statement_sha256":"3f0aaf6fa196a66818755e0fe8d0a900ec8170dc90de754fc9b133303d4627b7"}
% HMT_RESULT_END:R3-008:residence
% HMT_RESULT_BEGIN:R3-009:residence
% {"material_control":"REV3_BASE_GRANULAR+PASS_NO_PERDIDA_REV3","material_state":"CONSERVADO_POR_BASE_GRANULAR_REV3","residences":[{"path":"/Users/ruben/Documents/HMT2/CONTINUIDAD_EDITORIAL/INTEGRACION_REV3_VON_NEUMANN_MATRICIAL_REALIZACION_2026-08-12/REV3_BASE_GRANULAR.json","sha256":"d89983feac6e064f0ce77d6506d60b94dd1eaf00ed174352d3e32bfc5a9b88d4"}],"result_id":"R3-009","statement_sha256":"84c9171db979cfcda487379afcdf640f4d78458cf4f7d5697a3ec0b30a2f3096"}
% HMT_RESULT_END:R3-009:residence
% HMT_RESULT_BEGIN:R3-010:residence
% {"material_control":"REV3_BASE_GRANULAR+PASS_NO_PERDIDA_REV3","material_state":"CONSERVADO_POR_BASE_GRANULAR_REV3","residences":[{"path":"/Users/ruben/Documents/HMT2/CONTINUIDAD_EDITORIAL/INTEGRACION_REV3_VON_NEUMANN_MATRICIAL_REALIZACION_2026-08-12/REV3_BASE_GRANULAR.json","sha256":"d89983feac6e064f0ce77d6506d60b94dd1eaf00ed174352d3e32bfc5a9b88d4"}],"result_id":"R3-010","statement_sha256":"3ce9bef7730788d82ad408181bd8942d9cc48d02b7e2b99b36b5f763b3a20968"}
% HMT_RESULT_END:R3-010:residence
% HMT_RESULT_BEGIN:R3-011:residence
% {"material_control":"REV3_BASE_GRANULAR+PASS_NO_PERDIDA_REV3","material_state":"CONSERVADO_POR_BASE_GRANULAR_REV3","residences":[{"path":"/Users/ruben/Documents/HMT2/CONTINUIDAD_EDITORIAL/INTEGRACION_REV3_VON_NEUMANN_MATRICIAL_REALIZACION_2026-08-12/REV3_BASE_GRANULAR.json","sha256":"d89983feac6e064f0ce77d6506d60b94dd1eaf00ed174352d3e32bfc5a9b88d4"}],"result_id":"R3-011","statement_sha256":"5de78d3622665d2d9585dcb93d9f3d762ca07c9959ab042c4c085b8ef5943597"}
% HMT_RESULT_END:R3-011:residence
% HMT_RESULT_BEGIN:R3-012:residence
% {"material_control":"REV3_BASE_GRANULAR+PASS_NO_PERDIDA_REV3","material_state":"CONSERVADO_POR_BASE_GRANULAR_REV3","residences":[{"path":"/Users/ruben/Documents/HMT2/CONTINUIDAD_EDITORIAL/INTEGRACION_REV3_VON_NEUMANN_MATRICIAL_REALIZACION_2026-08-12/REV3_BASE_GRANULAR.json","sha256":"d89983feac6e064f0ce77d6506d60b94dd1eaf00ed174352d3e32bfc5a9b88d4"}],"result_id":"R3-012","statement_sha256":"a873e4e36ef58361dfd662e6f6b08f0484e978fb07f397dab890d355b8ae0df8"}
% HMT_RESULT_END:R3-012:residence
% HMT_RESULT_BEGIN:R3-013:residence
% {"material_control":"REV3_BASE_GRANULAR+PASS_NO_PERDIDA_REV3","material_state":"CONSERVADO_POR_BASE_GRANULAR_REV3","residences":[{"path":"/Users/ruben/Documents/HMT2/CONTINUIDAD_EDITORIAL/INTEGRACION_REV3_VON_NEUMANN_MATRICIAL_REALIZACION_2026-08-12/REV3_BASE_GRANULAR.json","sha256":"d89983feac6e064f0ce77d6506d60b94dd1eaf00ed174352d3e32bfc5a9b88d4"}],"result_id":"R3-013","statement_sha256":"892af66958bb697bdbc9d87aa439e0dcebdecff82281fa756e35aac1088a3c74"}
% HMT_RESULT_END:R3-013:residence
% HMT_RESULT_BEGIN:R3-014:residence
% {"material_control":"REV3_BASE_GRANULAR+PASS_NO_PERDIDA_REV3","material_state":"CONSERVADO_POR_BASE_GRANULAR_REV3","residences":[{"path":"/Users/ruben/Documents/HMT2/CONTINUIDAD_EDITORIAL/INTEGRACION_REV3_VON_NEUMANN_MATRICIAL_REALIZACION_2026-08-12/REV3_BASE_GRANULAR.json","sha256":"d89983feac6e064f0ce77d6506d60b94dd1eaf00ed174352d3e32bfc5a9b88d4"}],"result_id":"R3-014","statement_sha256":"3b135fbe4cbe5d0e292ffcee032f81cef6c993d833f85dde740957bd50df821f"}
% HMT_RESULT_END:R3-014:residence
% HMT_RESULT_BEGIN:R3-015:residence
% {"material_control":"REV3_BASE_GRANULAR+PASS_NO_PERDIDA_REV3","material_state":"CONSERVADO_POR_BASE_GRANULAR_REV3","residences":[{"path":"/Users/ruben/Documents/HMT2/CONTINUIDAD_EDITORIAL/INTEGRACION_REV3_VON_NEUMANN_MATRICIAL_REALIZACION_2026-08-12/REV3_BASE_GRANULAR.json","sha256":"d89983feac6e064f0ce77d6506d60b94dd1eaf00ed174352d3e32bfc5a9b88d4"}],"result_id":"R3-015","statement_sha256":"f94da8f805b5238f940f359590ae265b58be35f6bf4257fac8fa16cf8933c5f8"}
% HMT_RESULT_END:R3-015:residence
% HMT_RESULT_BEGIN:R3-016:residence
% {"material_control":"REV3_BASE_GRANULAR+PASS_NO_PERDIDA_REV3","material_state":"CONSERVADO_POR_BASE_GRANULAR_REV3","residences":[{"path":"/Users/ruben/Documents/HMT2/CONTINUIDAD_EDITORIAL/INTEGRACION_REV3_VON_NEUMANN_MATRICIAL_REALIZACION_2026-08-12/REV3_BASE_GRANULAR.json","sha256":"d89983feac6e064f0ce77d6506d60b94dd1eaf00ed174352d3e32bfc5a9b88d4"}],"result_id":"R3-016","statement_sha256":"660c7b55c80036946d6df58772d5bebbfded71d302316af3ab877a9e3cd1f3b7"}
% HMT_RESULT_END:R3-016:residence
% HMT_RESULT_BEGIN:R3-017:residence
% {"material_control":"REV3_BASE_GRANULAR+PASS_NO_PERDIDA_REV3","material_state":"CONSERVADO_POR_BASE_GRANULAR_REV3","residences":[{"path":"/Users/ruben/Documents/HMT2/CONTINUIDAD_EDITORIAL/INTEGRACION_REV3_VON_NEUMANN_MATRICIAL_REALIZACION_2026-08-12/REV3_BASE_GRANULAR.json","sha256":"d89983feac6e064f0ce77d6506d60b94dd1eaf00ed174352d3e32bfc5a9b88d4"}],"result_id":"R3-017","statement_sha256":"0834862c1f2fd11986bcab78a6a9ad03cd89152459ab8f8a6b83e9d1d4a0d0a8"}
% HMT_RESULT_END:R3-017:residence
% HMT_RESULT_BEGIN:R3-018:residence
% {"material_control":"REV3_BASE_GRANULAR+PASS_NO_PERDIDA_REV3","material_state":"CONSERVADO_POR_BASE_GRANULAR_REV3","residences":[{"path":"/Users/ruben/Documents/HMT2/CONTINUIDAD_EDITORIAL/INTEGRACION_REV3_VON_NEUMANN_MATRICIAL_REALIZACION_2026-08-12/REV3_BASE_GRANULAR.json","sha256":"d89983feac6e064f0ce77d6506d60b94dd1eaf00ed174352d3e32bfc5a9b88d4"}],"result_id":"R3-018","statement_sha256":"7a6ecf5a52d2b3bb92ca8b82bb5e3852c1a1f729a12ca8d39ad7353a61b2259b"}
% HMT_RESULT_END:R3-018:residence
% HMT_RESULT_BEGIN:R3-019:residence
% {"material_control":"REV3_BASE_GRANULAR+PASS_NO_PERDIDA_REV3","material_state":"CONSERVADO_POR_BASE_GRANULAR_REV3","residences":[{"path":"/Users/ruben/Documents/HMT2/CONTINUIDAD_EDITORIAL/INTEGRACION_REV3_VON_NEUMANN_MATRICIAL_REALIZACION_2026-08-12/REV3_BASE_GRANULAR.json","sha256":"d89983feac6e064f0ce77d6506d60b94dd1eaf00ed174352d3e32bfc5a9b88d4"}],"result_id":"R3-019","statement_sha256":"949668449b4331519738227ea1e782aefbfbe4067e371147c0821de5c4e2314c"}
% HMT_RESULT_END:R3-019:residence
% HMT_RESULT_BEGIN:R3-020:residence
% {"material_control":"REV3_BASE_GRANULAR+PASS_NO_PERDIDA_REV3","material_state":"CONSERVADO_POR_BASE_GRANULAR_REV3","residences":[{"path":"/Users/ruben/Documents/HMT2/CONTINUIDAD_EDITORIAL/INTEGRACION_REV3_VON_NEUMANN_MATRICIAL_REALIZACION_2026-08-12/REV3_BASE_GRANULAR.json","sha256":"d89983feac6e064f0ce77d6506d60b94dd1eaf00ed174352d3e32bfc5a9b88d4"}],"result_id":"R3-020","statement_sha256":"e6097d0d922f2e8e30dd23b38e77066a77f46464e19a875d95d0e9c21b9acd23"}
% HMT_RESULT_END:R3-020:residence
% HMT_RESULT_BEGIN:C-001:residence
% {"material_control":"PASS_MATRIZ_ATOMICA_AUTOCONTENIDA_REV3","material_state":"CONTROL_ACTIVO_EN_PLAN_Y_FUENTE","residences":[{"path":"/Users/ruben/Documents/HMT2/CONTINUIDAD_EDITORIAL/INTEGRACION_REV3_VON_NEUMANN_MATRICIAL_REALIZACION_2026-08-12/verificar_matriz_atomica_autocontenida.py","sha256":"a6598c0e1da8e8a8c61acc4150033a66348c039fb2ebf241791b0517dc458472"}],"result_id":"C-001","statement_sha256":"beae79d27d25d820609da41c6347a8561c6d41d85b1ec0525092421ad8721cf3"}
% HMT_RESULT_END:C-001:residence
% HMT_RESULT_BEGIN:C-002:residence
% {"material_control":"PASS_MATRIZ_ATOMICA_AUTOCONTENIDA_REV3","material_state":"CONTROL_ACTIVO_EN_PLAN_Y_FUENTE","residences":[{"path":"/Users/ruben/Documents/HMT2/CONTINUIDAD_EDITORIAL/INTEGRACION_REV3_VON_NEUMANN_MATRICIAL_REALIZACION_2026-08-12/verificar_matriz_atomica_autocontenida.py","sha256":"a6598c0e1da8e8a8c61acc4150033a66348c039fb2ebf241791b0517dc458472"}],"result_id":"C-002","statement_sha256":"5cf2ae9f5be634f470caa6670cf89d72015dc631a63870dc323d018d9f03d0e0"}
% HMT_RESULT_END:C-002:residence
% HMT_RESULT_BEGIN:C-003:residence
% {"material_control":"PASS_MATRIZ_ATOMICA_AUTOCONTENIDA_REV3","material_state":"CONTROL_ACTIVO_EN_PLAN_Y_FUENTE","residences":[{"path":"/Users/ruben/Documents/HMT2/CONTINUIDAD_EDITORIAL/INTEGRACION_REV3_VON_NEUMANN_MATRICIAL_REALIZACION_2026-08-12/verificar_matriz_atomica_autocontenida.py","sha256":"a6598c0e1da8e8a8c61acc4150033a66348c039fb2ebf241791b0517dc458472"}],"result_id":"C-003","statement_sha256":"2cabb42c3d350ac50184c7358426ca161c539e6dfd45b760168861d7461247ad"}
% HMT_RESULT_END:C-003:residence
% HMT_RESULT_BEGIN:C-004:residence
% {"material_control":"PASS_MATRIZ_ATOMICA_AUTOCONTENIDA_REV3","material_state":"CONTROL_ACTIVO_EN_PLAN_Y_FUENTE","residences":[{"path":"/Users/ruben/Documents/HMT2/CONTINUIDAD_EDITORIAL/INTEGRACION_REV3_VON_NEUMANN_MATRICIAL_REALIZACION_2026-08-12/verificar_matriz_atomica_autocontenida.py","sha256":"a6598c0e1da8e8a8c61acc4150033a66348c039fb2ebf241791b0517dc458472"}],"result_id":"C-004","statement_sha256":"9ff96e268873d82ee4ac8616df2a572ca23330c63f928a3e7a67ea10e2d46084"}
% HMT_RESULT_END:C-004:residence
% HMT_RESULT_BEGIN:C-005:residence
% {"material_control":"PASS_MATRIZ_ATOMICA_AUTOCONTENIDA_REV3","material_state":"CONTROL_ACTIVO_EN_PLAN_Y_FUENTE","residences":[{"path":"/Users/ruben/Documents/HMT2/CONTINUIDAD_EDITORIAL/INTEGRACION_REV3_VON_NEUMANN_MATRICIAL_REALIZACION_2026-08-12/verificar_matriz_atomica_autocontenida.py","sha256":"a6598c0e1da8e8a8c61acc4150033a66348c039fb2ebf241791b0517dc458472"}],"result_id":"C-005","statement_sha256":"f28286ac6df6ee110b8bdd8f2de9acd882a13dafee92dd7bc0645accfe4384d5"}
% HMT_RESULT_END:C-005:residence
% HMT_RESULT_BEGIN:C-006:residence
% {"material_control":"PASS_MATRIZ_ATOMICA_AUTOCONTENIDA_REV3","material_state":"CONTROL_ACTIVO_EN_PLAN_Y_FUENTE","residences":[{"path":"/Users/ruben/Documents/HMT2/CONTINUIDAD_EDITORIAL/INTEGRACION_REV3_VON_NEUMANN_MATRICIAL_REALIZACION_2026-08-12/verificar_matriz_atomica_autocontenida.py","sha256":"a6598c0e1da8e8a8c61acc4150033a66348c039fb2ebf241791b0517dc458472"}],"result_id":"C-006","statement_sha256":"61914b1116661f7db384d4e45bba72ec85e7ff331800fc50353e8e93f850dd76"}
% HMT_RESULT_END:C-006:residence
% HMT_RESULT_BEGIN:C-007:residence
% {"material_control":"PASS_MATRIZ_ATOMICA_AUTOCONTENIDA_REV3","material_state":"CONTROL_ACTIVO_EN_PLAN_Y_FUENTE","residences":[{"path":"/Users/ruben/Documents/HMT2/CONTINUIDAD_EDITORIAL/INTEGRACION_REV3_VON_NEUMANN_MATRICIAL_REALIZACION_2026-08-12/verificar_matriz_atomica_autocontenida.py","sha256":"a6598c0e1da8e8a8c61acc4150033a66348c039fb2ebf241791b0517dc458472"}],"result_id":"C-007","statement_sha256":"12e1672ff9d4988000a45bcd1a7b89ef0b97692c2ca5f6756ca06ad9670ede54"}
% HMT_RESULT_END:C-007:residence
% HMT_RESULT_BEGIN:MAS-RECTOR-001:residence
% {"material_control":"PASS_PAQUETE_RECTOR_MASAS_HMT_MD_2026_08_06","material_state":"CONSERVADO_CON_APTO_FOCAL","residences":[{"path":"/Users/ruben/Documents/New project/PUBLICACION_HMT/PAQUETE_RECTOR_LEY_GENERAL_MASAS_HMT_MD_2026-08-06/MATRIZ_RECTORA.tsv","sha256":"d4d58d9e5394804b83f48a2996ef884b275669b0f9468a28c9ed7c0a58affdc5"}],"result_id":"MAS-RECTOR-001","statement_sha256":"0a768984b144ee78bd8e25d36a2945dd57a6f83a5071c4fcfc7ba1a0b68f8d7f"}
% HMT_RESULT_END:MAS-RECTOR-001:residence
% HMT_RESULT_BEGIN:MAS-RECTOR-002:residence
% {"material_control":"PASS_PAQUETE_RECTOR_MASAS_HMT_MD_2026_08_06","material_state":"CONSERVADO_CON_APTO_FOCAL","residences":[{"path":"/Users/ruben/Documents/New project/PUBLICACION_HMT/PAQUETE_RECTOR_LEY_GENERAL_MASAS_HMT_MD_2026-08-06/MATRIZ_RECTORA.tsv","sha256":"d4d58d9e5394804b83f48a2996ef884b275669b0f9468a28c9ed7c0a58affdc5"}],"result_id":"MAS-RECTOR-002","statement_sha256":"89fa9d9d6538d06602255e139245ecc085a236981ab5f973f21ada8ca0f83cd3"}
% HMT_RESULT_END:MAS-RECTOR-002:residence
% HMT_RESULT_BEGIN:MAS-RECTOR-003:residence
% {"material_control":"PASS_PAQUETE_RECTOR_MASAS_HMT_MD_2026_08_06","material_state":"CONSERVADO_CON_APTO_FOCAL","residences":[{"path":"/Users/ruben/Documents/New project/PUBLICACION_HMT/PAQUETE_RECTOR_LEY_GENERAL_MASAS_HMT_MD_2026-08-06/MATRIZ_RECTORA.tsv","sha256":"d4d58d9e5394804b83f48a2996ef884b275669b0f9468a28c9ed7c0a58affdc5"}],"result_id":"MAS-RECTOR-003","statement_sha256":"a859bf6639b8355be0ff5ebcf8ba32630249b5d4842f80ebc5272733f0715109"}
% HMT_RESULT_END:MAS-RECTOR-003:residence
% HMT_RESULT_BEGIN:MAS-RECTOR-004:residence
% {"material_control":"PASS_PAQUETE_RECTOR_MASAS_HMT_MD_2026_08_06","material_state":"CONSERVADO_CON_APTO_FOCAL","residences":[{"path":"/Users/ruben/Documents/New project/PUBLICACION_HMT/PAQUETE_RECTOR_LEY_GENERAL_MASAS_HMT_MD_2026-08-06/MATRIZ_RECTORA.tsv","sha256":"d4d58d9e5394804b83f48a2996ef884b275669b0f9468a28c9ed7c0a58affdc5"}],"result_id":"MAS-RECTOR-004","statement_sha256":"5478fc19a7db005a9c0a02175258a722aa9289c8aae2b40c0f8be1a5539331e6"}
% HMT_RESULT_END:MAS-RECTOR-004:residence
% HMT_RESULT_BEGIN:MAS-RECTOR-005:residence
% {"material_control":"PASS_PAQUETE_RECTOR_MASAS_HMT_MD_2026_08_06","material_state":"CONSERVADO_CON_APTO_FOCAL","residences":[{"path":"/Users/ruben/Documents/New project/PUBLICACION_HMT/PAQUETE_RECTOR_LEY_GENERAL_MASAS_HMT_MD_2026-08-06/MATRIZ_RECTORA.tsv","sha256":"d4d58d9e5394804b83f48a2996ef884b275669b0f9468a28c9ed7c0a58affdc5"}],"result_id":"MAS-RECTOR-005","statement_sha256":"9ac8216130e6ddb43e8032b3eb72c4dd18ab87c05f238ee8fa65f06fdea782d2"}
% HMT_RESULT_END:MAS-RECTOR-005:residence
% HMT_RESULT_BEGIN:MAS-RECTOR-006:residence
% {"material_control":"PASS_PAQUETE_RECTOR_MASAS_HMT_MD_2026_08_06","material_state":"CONSERVADO_CON_APTO_FOCAL","residences":[{"path":"/Users/ruben/Documents/New project/PUBLICACION_HMT/PAQUETE_RECTOR_LEY_GENERAL_MASAS_HMT_MD_2026-08-06/MATRIZ_RECTORA.tsv","sha256":"d4d58d9e5394804b83f48a2996ef884b275669b0f9468a28c9ed7c0a58affdc5"}],"result_id":"MAS-RECTOR-006","statement_sha256":"f64d301daf63969b4ae52cb1259e81411210bb18d9c5e297f815a8f7ea6698e3"}
% HMT_RESULT_END:MAS-RECTOR-006:residence
% HMT_RESULT_BEGIN:MAS-RECTOR-007:residence
% {"material_control":"PASS_PAQUETE_RECTOR_MASAS_HMT_MD_2026_08_06","material_state":"CONSERVADO_CON_APTO_FOCAL","residences":[{"path":"/Users/ruben/Documents/New project/PUBLICACION_HMT/PAQUETE_RECTOR_LEY_GENERAL_MASAS_HMT_MD_2026-08-06/MATRIZ_RECTORA.tsv","sha256":"d4d58d9e5394804b83f48a2996ef884b275669b0f9468a28c9ed7c0a58affdc5"}],"result_id":"MAS-RECTOR-007","statement_sha256":"537017eed6d0ecd896c8a40aed8e3197aeef62bb6895d356b95a7bbc2948643e"}
% HMT_RESULT_END:MAS-RECTOR-007:residence
% HMT_RESULT_BEGIN:MAS-RECTOR-008:residence
% {"material_control":"PASS_PAQUETE_RECTOR_MASAS_HMT_MD_2026_08_06","material_state":"CONSERVADO_CON_APTO_FOCAL","residences":[{"path":"/Users/ruben/Documents/New project/PUBLICACION_HMT/PAQUETE_RECTOR_LEY_GENERAL_MASAS_HMT_MD_2026-08-06/MATRIZ_RECTORA.tsv","sha256":"d4d58d9e5394804b83f48a2996ef884b275669b0f9468a28c9ed7c0a58affdc5"}],"result_id":"MAS-RECTOR-008","statement_sha256":"9563efbb323a0406d27d86e100b82f2b9ca50dd08d2197ca5d02418e4656c614"}
% HMT_RESULT_END:MAS-RECTOR-008:residence
% HMT_RESULT_BEGIN:MAS-RECTOR-009:residence
% {"material_control":"PASS_PAQUETE_RECTOR_MASAS_HMT_MD_2026_08_06","material_state":"CONSERVADO_CON_APTO_FOCAL","residences":[{"path":"/Users/ruben/Documents/New project/PUBLICACION_HMT/PAQUETE_RECTOR_LEY_GENERAL_MASAS_HMT_MD_2026-08-06/MATRIZ_RECTORA.tsv","sha256":"d4d58d9e5394804b83f48a2996ef884b275669b0f9468a28c9ed7c0a58affdc5"}],"result_id":"MAS-RECTOR-009","statement_sha256":"cda0251b6312a403ff360c8c40e29d102a3395cef7e4b3f6ea6b2b72ef982e76"}
% HMT_RESULT_END:MAS-RECTOR-009:residence
% HMT_RESULT_BEGIN:MAS-RECTOR-010:residence
% {"material_control":"PASS_PAQUETE_RECTOR_MASAS_HMT_MD_2026_08_06","material_state":"CONSERVADO_CON_APTO_FOCAL","residences":[{"path":"/Users/ruben/Documents/New project/PUBLICACION_HMT/PAQUETE_RECTOR_LEY_GENERAL_MASAS_HMT_MD_2026-08-06/MATRIZ_RECTORA.tsv","sha256":"d4d58d9e5394804b83f48a2996ef884b275669b0f9468a28c9ed7c0a58affdc5"}],"result_id":"MAS-RECTOR-010","statement_sha256":"5913a73e04b4c752216af319b77e125d31914639200b39b45d1a6b28b536508b"}
% HMT_RESULT_END:MAS-RECTOR-010:residence
% HMT_RESULT_BEGIN:MAS-RECTOR-011:residence
% {"material_control":"PASS_PAQUETE_RECTOR_MASAS_HMT_MD_2026_08_06","material_state":"CONSERVADO_CON_APTO_FOCAL","residences":[{"path":"/Users/ruben/Documents/New project/PUBLICACION_HMT/PAQUETE_RECTOR_LEY_GENERAL_MASAS_HMT_MD_2026-08-06/MATRIZ_RECTORA.tsv","sha256":"d4d58d9e5394804b83f48a2996ef884b275669b0f9468a28c9ed7c0a58affdc5"}],"result_id":"MAS-RECTOR-011","statement_sha256":"b56b1d18f5b984455801b89eace1fbf9143fa965b645cd443b7bbef759b09251"}
% HMT_RESULT_END:MAS-RECTOR-011:residence
% HMT_RESULT_BEGIN:MAS-RECTOR-012:residence
% {"material_control":"PASS_PAQUETE_RECTOR_MASAS_HMT_MD_2026_08_06","material_state":"CONSERVADO_CON_APTO_FOCAL","residences":[{"path":"/Users/ruben/Documents/New project/PUBLICACION_HMT/PAQUETE_RECTOR_LEY_GENERAL_MASAS_HMT_MD_2026-08-06/MATRIZ_RECTORA.tsv","sha256":"d4d58d9e5394804b83f48a2996ef884b275669b0f9468a28c9ed7c0a58affdc5"}],"result_id":"MAS-RECTOR-012","statement_sha256":"aa319b29fb1a3ad47b2d521716fea82da5660479225223eac504798c36b42f24"}
% HMT_RESULT_END:MAS-RECTOR-012:residence
% HMT_RESULT_BEGIN:MAS-RECTOR-013:residence
% {"material_control":"PASS_PAQUETE_RECTOR_MASAS_HMT_MD_2026_08_06","material_state":"CONSERVADO_CON_APTO_FOCAL","residences":[{"path":"/Users/ruben/Documents/New project/PUBLICACION_HMT/PAQUETE_RECTOR_LEY_GENERAL_MASAS_HMT_MD_2026-08-06/MATRIZ_RECTORA.tsv","sha256":"d4d58d9e5394804b83f48a2996ef884b275669b0f9468a28c9ed7c0a58affdc5"}],"result_id":"MAS-RECTOR-013","statement_sha256":"c9f012d4f1b9fa1d9a642670582f65c700f433930615ddc0500e4beef203d5e9"}
% HMT_RESULT_END:MAS-RECTOR-013:residence
% HMT_RESULT_BEGIN:MAS-RECTOR-014:residence
% {"material_control":"PASS_PAQUETE_RECTOR_MASAS_HMT_MD_2026_08_06","material_state":"CONSERVADO_CON_APTO_FOCAL","residences":[{"path":"/Users/ruben/Documents/New project/PUBLICACION_HMT/PAQUETE_RECTOR_LEY_GENERAL_MASAS_HMT_MD_2026-08-06/MATRIZ_RECTORA.tsv","sha256":"d4d58d9e5394804b83f48a2996ef884b275669b0f9468a28c9ed7c0a58affdc5"}],"result_id":"MAS-RECTOR-014","statement_sha256":"12aadab1bbdd24c3cf826e52a333c34f70ced2a37f655b2919eaaaf3281fef04"}
% HMT_RESULT_END:MAS-RECTOR-014:residence
% HMT_RESULT_BEGIN:MAS-RECTOR-015:residence
% {"material_control":"PASS_PAQUETE_RECTOR_MASAS_HMT_MD_2026_08_06","material_state":"CONSERVADO_CON_APTO_FOCAL","residences":[{"path":"/Users/ruben/Documents/New project/PUBLICACION_HMT/PAQUETE_RECTOR_LEY_GENERAL_MASAS_HMT_MD_2026-08-06/MATRIZ_RECTORA.tsv","sha256":"d4d58d9e5394804b83f48a2996ef884b275669b0f9468a28c9ed7c0a58affdc5"}],"result_id":"MAS-RECTOR-015","statement_sha256":"ae01c9c511c63968eed7516b99a2cf7e2a65d0d1c6f56e756638a3d524bb5f9f"}
% HMT_RESULT_END:MAS-RECTOR-015:residence
% HMT_RESULT_BEGIN:MAS-RECTOR-016:residence
% {"material_control":"PASS_PAQUETE_RECTOR_MASAS_HMT_MD_2026_08_06","material_state":"CONSERVADO_CON_APTO_FOCAL","residences":[{"path":"/Users/ruben/Documents/New project/PUBLICACION_HMT/PAQUETE_RECTOR_LEY_GENERAL_MASAS_HMT_MD_2026-08-06/MATRIZ_RECTORA.tsv","sha256":"d4d58d9e5394804b83f48a2996ef884b275669b0f9468a28c9ed7c0a58affdc5"}],"result_id":"MAS-RECTOR-016","statement_sha256":"685df6bdb27ac04026246ca3f09b212f9aabebf22929fb63ddfe1222da2a6e76"}
% HMT_RESULT_END:MAS-RECTOR-016:residence
% HMT_RESULT_BEGIN:MAS-RECTOR-017:residence
% {"material_control":"PASS_PAQUETE_RECTOR_MASAS_HMT_MD_2026_08_06","material_state":"CONSERVADO_CON_APTO_FOCAL","residences":[{"path":"/Users/ruben/Documents/New project/PUBLICACION_HMT/PAQUETE_RECTOR_LEY_GENERAL_MASAS_HMT_MD_2026-08-06/MATRIZ_RECTORA.tsv","sha256":"d4d58d9e5394804b83f48a2996ef884b275669b0f9468a28c9ed7c0a58affdc5"}],"result_id":"MAS-RECTOR-017","statement_sha256":"2d6f3806c75ac84503e6b8ce4c6fe49dbc650e525915991cfef4ecc29e12bb55"}
% HMT_RESULT_END:MAS-RECTOR-017:residence
% HMT_RESULT_BEGIN:MAS-RECTOR-018:residence
% {"material_control":"PASS_PAQUETE_RECTOR_MASAS_HMT_MD_2026_08_06","material_state":"CONSERVADO_CON_APTO_FOCAL","residences":[{"path":"/Users/ruben/Documents/New project/PUBLICACION_HMT/PAQUETE_RECTOR_LEY_GENERAL_MASAS_HMT_MD_2026-08-06/MATRIZ_RECTORA.tsv","sha256":"d4d58d9e5394804b83f48a2996ef884b275669b0f9468a28c9ed7c0a58affdc5"}],"result_id":"MAS-RECTOR-018","statement_sha256":"1fc9dcfc0f77f0e3dcf950b1b3ffb6ee7ee7314e9e8407bb22f75fa7efd5300c"}
% HMT_RESULT_END:MAS-RECTOR-018:residence
% HMT_RESULT_BEGIN:MAS-RECTOR-019:residence
% {"material_control":"PASS_PAQUETE_RECTOR_MASAS_HMT_MD_2026_08_06","material_state":"CONSERVADO_CON_APTO_FOCAL","residences":[{"path":"/Users/ruben/Documents/New project/PUBLICACION_HMT/PAQUETE_RECTOR_LEY_GENERAL_MASAS_HMT_MD_2026-08-06/MATRIZ_RECTORA.tsv","sha256":"d4d58d9e5394804b83f48a2996ef884b275669b0f9468a28c9ed7c0a58affdc5"}],"result_id":"MAS-RECTOR-019","statement_sha256":"ebbe49ae2f4afbc32a86801a2996b2c2d2bf71c0ff4f86d9dc3e58cf6366ae60"}
% HMT_RESULT_END:MAS-RECTOR-019:residence
% HMT_RESULT_BEGIN:MAS-RECTOR-020:residence
% {"material_control":"PASS_PAQUETE_RECTOR_MASAS_HMT_MD_2026_08_06","material_state":"CONSERVADO_CON_APTO_FOCAL","residences":[{"path":"/Users/ruben/Documents/New project/PUBLICACION_HMT/PAQUETE_RECTOR_LEY_GENERAL_MASAS_HMT_MD_2026-08-06/MATRIZ_RECTORA.tsv","sha256":"d4d58d9e5394804b83f48a2996ef884b275669b0f9468a28c9ed7c0a58affdc5"}],"result_id":"MAS-RECTOR-020","statement_sha256":"ea7797a9dc2c4c52d16299bf176fa0b0178158922903c442a8cb1aa85af17879"}
% HMT_RESULT_END:MAS-RECTOR-020:residence
% HMT_RESULT_BEGIN:MAS-RECTOR-021:residence
% {"material_control":"PASS_PAQUETE_RECTOR_MASAS_HMT_MD_2026_08_06","material_state":"CONSERVADO_CON_APTO_FOCAL","residences":[{"path":"/Users/ruben/Documents/New project/PUBLICACION_HMT/PAQUETE_RECTOR_LEY_GENERAL_MASAS_HMT_MD_2026-08-06/MATRIZ_RECTORA.tsv","sha256":"d4d58d9e5394804b83f48a2996ef884b275669b0f9468a28c9ed7c0a58affdc5"}],"result_id":"MAS-RECTOR-021","statement_sha256":"7f875aaae09118d70a3022037a232a45a3e2fe74a8c92aae01e5486aa2ef7f95"}
% HMT_RESULT_END:MAS-RECTOR-021:residence
% HMT_RESULT_BEGIN:MAS-RECTOR-022:residence
% {"material_control":"PASS_PAQUETE_RECTOR_MASAS_HMT_MD_2026_08_06","material_state":"CONSERVADO_CON_APTO_FOCAL","residences":[{"path":"/Users/ruben/Documents/New project/PUBLICACION_HMT/PAQUETE_RECTOR_LEY_GENERAL_MASAS_HMT_MD_2026-08-06/MATRIZ_RECTORA.tsv","sha256":"d4d58d9e5394804b83f48a2996ef884b275669b0f9468a28c9ed7c0a58affdc5"}],"result_id":"MAS-RECTOR-022","statement_sha256":"c6e9560103ad03e63b2460672442d06e18c09c1f866c7d34ee4946f543255d51"}
% HMT_RESULT_END:MAS-RECTOR-022:residence
% HMT_RESULT_BEGIN:MAS-RECTOR-023:residence
% {"material_control":"PASS_PAQUETE_RECTOR_MASAS_HMT_MD_2026_08_06","material_state":"CONSERVADO_CON_APTO_FOCAL","residences":[{"path":"/Users/ruben/Documents/New project/PUBLICACION_HMT/PAQUETE_RECTOR_LEY_GENERAL_MASAS_HMT_MD_2026-08-06/MATRIZ_RECTORA.tsv","sha256":"d4d58d9e5394804b83f48a2996ef884b275669b0f9468a28c9ed7c0a58affdc5"}],"result_id":"MAS-RECTOR-023","statement_sha256":"dd2590f85503d95b6ed6280bb6540f982aadade680f17a9ca0298c3abb7fe3c6"}
% HMT_RESULT_END:MAS-RECTOR-023:residence
% HMT_RESULT_BEGIN:MAS-RECTOR-024:residence
% {"material_control":"PASS_PAQUETE_RECTOR_MASAS_HMT_MD_2026_08_06","material_state":"CONSERVADO_CON_APTO_FOCAL","residences":[{"path":"/Users/ruben/Documents/New project/PUBLICACION_HMT/PAQUETE_RECTOR_LEY_GENERAL_MASAS_HMT_MD_2026-08-06/MATRIZ_RECTORA.tsv","sha256":"d4d58d9e5394804b83f48a2996ef884b275669b0f9468a28c9ed7c0a58affdc5"}],"result_id":"MAS-RECTOR-024","statement_sha256":"4d3e6032d279cb845f8bb3fe2e3bcd153b9585996d56fd93c0e721793c358eae"}
% HMT_RESULT_END:MAS-RECTOR-024:residence
% HMT_RESULT_BEGIN:MAS-RECTOR-025:residence
% {"material_control":"PASS_PAQUETE_RECTOR_MASAS_HMT_MD_2026_08_06","material_state":"CONSERVADO_CON_APTO_FOCAL","residences":[{"path":"/Users/ruben/Documents/New project/PUBLICACION_HMT/PAQUETE_RECTOR_LEY_GENERAL_MASAS_HMT_MD_2026-08-06/MATRIZ_RECTORA.tsv","sha256":"d4d58d9e5394804b83f48a2996ef884b275669b0f9468a28c9ed7c0a58affdc5"}],"result_id":"MAS-RECTOR-025","statement_sha256":"6e3bb65bdf6f87c71551bf92d9f6b8e7c60ba9785857316070c5d88c9a59fb9a"}
% HMT_RESULT_END:MAS-RECTOR-025:residence
% HMT_RESULT_BEGIN:MAS-RECTOR-026:residence
% {"material_control":"PASS_PAQUETE_RECTOR_MASAS_HMT_MD_2026_08_06","material_state":"CONSERVADO_CON_APTO_FOCAL","residences":[{"path":"/Users/ruben/Documents/New project/PUBLICACION_HMT/PAQUETE_RECTOR_LEY_GENERAL_MASAS_HMT_MD_2026-08-06/MATRIZ_RECTORA.tsv","sha256":"d4d58d9e5394804b83f48a2996ef884b275669b0f9468a28c9ed7c0a58affdc5"}],"result_id":"MAS-RECTOR-026","statement_sha256":"33fe0ca572be6ea55887ec4b42e78250c112e42699de77e13eaf5ddfdf608e12"}
% HMT_RESULT_END:MAS-RECTOR-026:residence
% HMT_RESULT_BEGIN:MAS-RECTOR-027:residence
% {"material_control":"PASS_PAQUETE_RECTOR_MASAS_HMT_MD_2026_08_06","material_state":"CONSERVADO_CON_APTO_FOCAL","residences":[{"path":"/Users/ruben/Documents/New project/PUBLICACION_HMT/PAQUETE_RECTOR_LEY_GENERAL_MASAS_HMT_MD_2026-08-06/MATRIZ_RECTORA.tsv","sha256":"d4d58d9e5394804b83f48a2996ef884b275669b0f9468a28c9ed7c0a58affdc5"}],"result_id":"MAS-RECTOR-027","statement_sha256":"6168ae29b3e71bf7885c166f21fc815528ad39b20fa81a628be34700e902949c"}
% HMT_RESULT_END:MAS-RECTOR-027:residence
% HMT_RESULT_BEGIN:MAS-RECTOR-028:residence
% {"material_control":"PASS_PAQUETE_RECTOR_MASAS_HMT_MD_2026_08_06","material_state":"CONSERVADO_CON_APTO_FOCAL","residences":[{"path":"/Users/ruben/Documents/New project/PUBLICACION_HMT/PAQUETE_RECTOR_LEY_GENERAL_MASAS_HMT_MD_2026-08-06/MATRIZ_RECTORA.tsv","sha256":"d4d58d9e5394804b83f48a2996ef884b275669b0f9468a28c9ed7c0a58affdc5"}],"result_id":"MAS-RECTOR-028","statement_sha256":"52f3251b7aec29d108b1b6c06d3cc064189cfe5cad74e31c311b1eab0be4b3d1"}
% HMT_RESULT_END:MAS-RECTOR-028:residence
% HMT_RESULT_BEGIN:MAS-PUBLIC-023:residence
% {"material_control":"PASS_PAQUETE_RECTOR_MASAS_HMT_MD_2026_08_06","material_state":"CONSERVADO_CON_APTO_FOCAL","residences":[{"path":"/Users/ruben/Documents/New project/PUBLICACION_HMT/PAQUETE_RECTOR_LEY_GENERAL_MASAS_HMT_MD_2026-08-06/MATRIZ_RECTORA.tsv","sha256":"d4d58d9e5394804b83f48a2996ef884b275669b0f9468a28c9ed7c0a58affdc5"}],"result_id":"MAS-PUBLIC-023","statement_sha256":"2490d5f9a177b2e5c480be8158dc8b1bccbb5f0785cb1e610be8e85f3a6b8918"}
% HMT_RESULT_END:MAS-PUBLIC-023:residence
% HMT_RESULT_BEGIN:MAS-COND-017:residence
% {"material_control":"PASS_PAQUETE_RECTOR_MASAS_HMT_MD_2026_08_06","material_state":"CONSERVADO_CON_APTO_FOCAL","residences":[{"path":"/Users/ruben/Documents/New project/PUBLICACION_HMT/PAQUETE_RECTOR_LEY_GENERAL_MASAS_HMT_MD_2026-08-06/MATRIZ_RECTORA.tsv","sha256":"d4d58d9e5394804b83f48a2996ef884b275669b0f9468a28c9ed7c0a58affdc5"}],"result_id":"MAS-COND-017","statement_sha256":"d9b0c82973bb442a515af30c9bfe351d9b9ab183eb81db8250f32c615d5f50ba"}
% HMT_RESULT_END:MAS-COND-017:residence
% HMT_RESULT_BEGIN:INF-001:residence
% {"material_control":"GATE_MATERIAL_64+PASS_PUBLICACION_INTEGRAL","material_state":"INTEGRADO_O_REMITIDO_SIN_PERDIDA","residences":[{"path":"manuscrito/sections/hmt/05_tpk_numero_medicion.tex","sha256":"51deffafc13a6d83bc07e15f77892d1ec32ba05f54cc2b01d6417ea13e3599a2"}],"result_id":"INF-001","statement_sha256":"d2d4fc4dc0a639778129fc6046b7f974c7a2cc39dae5b6aede5bd2ce96d91b1d"}
% HMT_RESULT_END:INF-001:residence
% HMT_RESULT_BEGIN:INF-002:residence
% {"material_control":"GATE_MATERIAL_64+PASS_PUBLICACION_INTEGRAL","material_state":"INTEGRADO_O_REMITIDO_SIN_PERDIDA","residences":[{"path":"manuscrito/sections/hmt/05_tpk_numero_medicion.tex","sha256":"51deffafc13a6d83bc07e15f77892d1ec32ba05f54cc2b01d6417ea13e3599a2"}],"result_id":"INF-002","statement_sha256":"8e32b4ab88538fa1ca2b52a628278c2f084dc54f423cf78093a3270700d157a5"}
% HMT_RESULT_END:INF-002:residence
% HMT_RESULT_BEGIN:INF-003:residence
% {"material_control":"GATE_MATERIAL_64+PASS_PUBLICACION_INTEGRAL","material_state":"INTEGRADO_O_REMITIDO_SIN_PERDIDA","residences":[{"path":"manuscrito/sections/hmt/05_tpk_numero_medicion.tex","sha256":"51deffafc13a6d83bc07e15f77892d1ec32ba05f54cc2b01d6417ea13e3599a2"}],"result_id":"INF-003","statement_sha256":"667c50d6578b21e1abcff2bc241db42bbe6e270c7268fef0d64172e29d096836"}
% HMT_RESULT_END:INF-003:residence
% HMT_RESULT_BEGIN:N9-001:residence
% {"material_control":"GATE_MATERIAL_64+PASS_PUBLICACION_INTEGRAL","material_state":"INTEGRADO_O_REMITIDO_SIN_PERDIDA","residences":[{"path":"manuscrito/sections/hmt/06aa_genealogia_filtracion_g9_autonoma.tex","sha256":"ba4816ac627ae2d7eb4141421a19a240bf66022f8728f4ce8bf873bd439efb39"}],"result_id":"N9-001","statement_sha256":"f38e038d91253c1364af870bc92f1e5a243d0a5e5c69c5f85bc4d05cc2ef6fb9"}
% HMT_RESULT_END:N9-001:residence
% HMT_RESULT_BEGIN:N9-002:residence
% {"material_control":"GATE_MATERIAL_64+PASS_PUBLICACION_INTEGRAL","material_state":"INTEGRADO_O_REMITIDO_SIN_PERDIDA","residences":[{"path":"manuscrito/sections/hmt/06aa_genealogia_filtracion_g9_autonoma.tex","sha256":"ba4816ac627ae2d7eb4141421a19a240bf66022f8728f4ce8bf873bd439efb39"}],"result_id":"N9-002","statement_sha256":"e4c29231470c3fe69fe80b927a6249121b8a79287a407ca4318e1889a8ed9fb6"}
% HMT_RESULT_END:N9-002:residence
% HMT_RESULT_BEGIN:N9-003:residence
% {"material_control":"GATE_MATERIAL_64+PASS_PUBLICACION_INTEGRAL","material_state":"INTEGRADO_O_REMITIDO_SIN_PERDIDA","residences":[{"path":"manuscrito/sections/hmt/06b_cinco_ciclos_pi.tex","sha256":"83f8a2be69639ac01847a0e25720d3e58b193ae689abfe5fd567a552fb301012"},{"path":"manuscrito/sections/hmt/06d_biografias_numeros_rev6.tex","sha256":"a01607a1a6f5ee4fadbd482daa384b09b192d315da15048ff661d2baf7d197f6"},{"path":"manuscrito/sections/hmt/06e_certificado_generacion_monodromica_rev7.tex","sha256":"2a0cf8028e12605a97923c53726fcb62c1798c458af8e9f78f12c1129b0e1b94"}],"result_id":"N9-003","statement_sha256":"a3b1e710d9bbef1129d8ddb4da929c36fe0b2eaf7ea7e20bbc91312eff72cb57"}
% HMT_RESULT_END:N9-003:residence
% HMT_RESULT_BEGIN:N9-004:residence
% {"material_control":"GATE_MATERIAL_64+PASS_PUBLICACION_INTEGRAL","material_state":"INTEGRADO_O_REMITIDO_SIN_PERDIDA","residences":[{"path":"manuscrito/sections/hmt/06e_certificado_generacion_monodromica_rev7.tex","sha256":"2a0cf8028e12605a97923c53726fcb62c1798c458af8e9f78f12c1129b0e1b94"}],"result_id":"N9-004","statement_sha256":"5288a67a9b48b60a9a1896b482933707612c1ae1e4b4862ba2b23f9ff393e77f"}
% HMT_RESULT_END:N9-004:residence
% HMT_RESULT_BEGIN:HEL-001:residence
% {"material_control":"GATE_MATERIAL_64+PASS_PUBLICACION_INTEGRAL","material_state":"INTEGRADO_O_REMITIDO_SIN_PERDIDA","residences":[{"path":"manuscrito/sections/hmt/03l_elevacion_cohomologica_dimension_rev2.tex","sha256":"1b0f6b86dc275a3947b64e319b0b19789d62c83bcc3714d5c0b8a2cea2243a8c"}],"result_id":"HEL-001","statement_sha256":"3ed30cd71d3b16fcdaf015298295033782765571ae3a2c7d60f65aa8b0bee52d"}
% HMT_RESULT_END:HEL-001:residence
% HMT_RESULT_BEGIN:OPS-001:residence
% {"material_control":"GATE_MATERIAL_64+PASS_PUBLICACION_INTEGRAL","material_state":"INTEGRADO_O_REMITIDO_SIN_PERDIDA","residences":[{"path":"manuscrito/sections/hmt/03k_genealogia_operaciones_rev2.tex","sha256":"fae7aca94c6fb20754431ed79a5aa2f005a7cf2f5681736fc5826da4c694fb98"}],"result_id":"OPS-001","statement_sha256":"61ed21a147f32626c42baccde2ffe7594b14cf6c1c83912b96d5f3f8776405e5"}
% HMT_RESULT_END:OPS-001:residence
% HMT_RESULT_BEGIN:ZF-001:residence
% {"material_control":"GATE_MATERIAL_64+PASS_PUBLICACION_INTEGRAL","material_state":"INTEGRADO_O_REMITIDO_SIN_PERDIDA","residences":[{"path":"manuscrito/sections/integracion_20260812/von_neumann_64/04_zfc_fundacion_ordinales.tex","sha256":"5d81fd31bf31a950f9d068f278e76ef0ebfb37cb8cf6ee5483575ab87f5fddb0"}],"result_id":"ZF-001","statement_sha256":"1d6985e3c4c3690eb43465daa6bcc3d7fb39cb3d0ae1fe31a3cd6c70db621314"}
% HMT_RESULT_END:ZF-001:residence
% HMT_RESULT_BEGIN:ZF-002:residence
% {"material_control":"GATE_MATERIAL_64+PASS_PUBLICACION_INTEGRAL","material_state":"INTEGRADO_O_REMITIDO_SIN_PERDIDA","residences":[{"path":"manuscrito/sections/integracion_20260812/von_neumann_64/04_zfc_fundacion_ordinales.tex","sha256":"5d81fd31bf31a950f9d068f278e76ef0ebfb37cb8cf6ee5483575ab87f5fddb0"}],"result_id":"ZF-002","statement_sha256":"e78afa0785eb7978d61364af148ab43c40092242ad1edf80db0e5c19f11915cd"}
% HMT_RESULT_END:ZF-002:residence
% HMT_RESULT_BEGIN:EXT-001:residence
% {"material_control":"GATE_MATERIAL_64+PASS_PUBLICACION_INTEGRAL","material_state":"INTEGRADO_O_REMITIDO_SIN_PERDIDA","residences":[{"path":"manuscrito/sections/hmt/14d_ontologia_numeros_rev11.tex","sha256":"c96b3c94af4668f15c6a061260d1c1900c52b69b8a4a5d2442268f6ff1d0abcc"}],"result_id":"EXT-001","statement_sha256":"664eeafa4a1a0cfad63aa9414daaccdf8f0ea8a15f92851a674ad7cd5d26c9a8"}
% HMT_RESULT_END:EXT-001:residence
% HMT_RESULT_BEGIN:FND-001:residence
% {"material_control":"GATE_MATERIAL_64+PASS_PUBLICACION_INTEGRAL","material_state":"INTEGRADO_O_REMITIDO_SIN_PERDIDA","residences":[{"path":"manuscrito/sections/integracion_20260812/von_neumann_64/04_zfc_fundacion_ordinales.tex","sha256":"5d81fd31bf31a950f9d068f278e76ef0ebfb37cb8cf6ee5483575ab87f5fddb0"},{"path":"manuscrito/sections/hmt/00a_continuo_numero_medida_hmt_revfinal.tex","sha256":"1be7648180a45eb337db43d6d18a63624876b9f568ea29998c80b00d13b1a819"}],"result_id":"FND-001","statement_sha256":"6b0867dcae5d6924d26cbfdfb829827b9b60d7461148fab933762bd8ab425ac7"}
% HMT_RESULT_END:FND-001:residence
% HMT_RESULT_BEGIN:ORD-001:residence
% {"material_control":"GATE_MATERIAL_64+PASS_PUBLICACION_INTEGRAL","material_state":"INTEGRADO_O_REMITIDO_SIN_PERDIDA","residences":[{"path":"manuscrito/sections/integracion_20260812/von_neumann_64/04_zfc_fundacion_ordinales.tex","sha256":"5d81fd31bf31a950f9d068f278e76ef0ebfb37cb8cf6ee5483575ab87f5fddb0"},{"path":"manuscrito/sections/hmt/14f_cierre_genealogico_numero.tex","sha256":"4522230d8297a1420ccb586d0d6c1fefcb08b5e8e2deeda85b11cc73afd8ded5"}],"result_id":"ORD-001","statement_sha256":"1c7c37b7877465ee0663efb162f12aa0614e44999754b59066edcfa489f7f6ae"}
% HMT_RESULT_END:ORD-001:residence
% HMT_RESULT_BEGIN:ORD-002:residence
% {"material_control":"GATE_MATERIAL_64+PASS_PUBLICACION_INTEGRAL","material_state":"INTEGRADO_O_REMITIDO_SIN_PERDIDA","residences":[{"path":"manuscrito/sections/integracion_20260812/von_neumann_64/04_zfc_fundacion_ordinales.tex","sha256":"5d81fd31bf31a950f9d068f278e76ef0ebfb37cb8cf6ee5483575ab87f5fddb0"}],"result_id":"ORD-002","statement_sha256":"c8188c7c87197d0287f9d7bde9efd04e7246384d5271159053cefc1b62278bf1"}
% HMT_RESULT_END:ORD-002:residence
% HMT_RESULT_BEGIN:PWR-001:residence
% {"material_control":"GATE_MATERIAL_64+PASS_PUBLICACION_INTEGRAL","material_state":"INTEGRADO_O_REMITIDO_SIN_PERDIDA","residences":[{"path":"manuscrito/sections/integracion_20260812/von_neumann_64/05_potencia_eleccion.tex","sha256":"cb412bdee19cc8e65181f75c0666a966dc20111f1dafcaee563bf1ee0e3d2684"}],"result_id":"PWR-001","statement_sha256":"15e7dd2441fdb8f717dbe8bbfb825a15a1165399b2158782c52216b71f1cb8e4"}
% HMT_RESULT_END:PWR-001:residence
% HMT_RESULT_BEGIN:PWR-002:residence
% {"material_control":"GATE_MATERIAL_64+PASS_PUBLICACION_INTEGRAL","material_state":"INTEGRADO_O_REMITIDO_SIN_PERDIDA","residences":[{"path":"manuscrito/sections/integracion_20260812/von_neumann_64/05_potencia_eleccion.tex","sha256":"cb412bdee19cc8e65181f75c0666a966dc20111f1dafcaee563bf1ee0e3d2684"}],"result_id":"PWR-002","statement_sha256":"667fc09531d0de7134ded48a16fb499d56212ba6d7734d3cbc1e04ec1144ad52"}
% HMT_RESULT_END:PWR-002:residence
% HMT_RESULT_BEGIN:CHO-001:residence
% {"material_control":"GATE_MATERIAL_64+PASS_PUBLICACION_INTEGRAL","material_state":"INTEGRADO_O_REMITIDO_SIN_PERDIDA","residences":[{"path":"manuscrito/sections/integracion_20260812/von_neumann_64/05_potencia_eleccion.tex","sha256":"cb412bdee19cc8e65181f75c0666a966dc20111f1dafcaee563bf1ee0e3d2684"},{"path":"manuscrito/sections/hmt/14f_cierre_genealogico_numero.tex","sha256":"4522230d8297a1420ccb586d0d6c1fefcb08b5e8e2deeda85b11cc73afd8ded5"}],"result_id":"CHO-001","statement_sha256":"8b2d6fc14cefc2d7b601d295b4441c23d96827e08dd639c47877a67cb5b36929"}
% HMT_RESULT_END:CHO-001:residence
% HMT_RESULT_BEGIN:CHO-002:residence
% {"material_control":"GATE_MATERIAL_64+PASS_PUBLICACION_INTEGRAL","material_state":"INTEGRADO_O_REMITIDO_SIN_PERDIDA","residences":[{"path":"manuscrito/sections/integracion_20260812/von_neumann_64/05_potencia_eleccion.tex","sha256":"cb412bdee19cc8e65181f75c0666a966dc20111f1dafcaee563bf1ee0e3d2684"}],"result_id":"CHO-002","statement_sha256":"26fee3b2fa8f6cd1bc1073d172a9682b085c58a85cae5e2a6864e89d45f7ddea"}
% HMT_RESULT_END:CHO-002:residence
% HMT_RESULT_BEGIN:CAT-001:residence
% {"material_control":"GATE_MATERIAL_64+PASS_PUBLICACION_INTEGRAL","material_state":"INTEGRADO_O_REMITIDO_SIN_PERDIDA","residences":[{"path":"manuscrito/sections/integracion_20260812/von_neumann_64/04_zfc_fundacion_ordinales.tex","sha256":"5d81fd31bf31a950f9d068f278e76ef0ebfb37cb8cf6ee5483575ab87f5fddb0"},{"path":"manuscrito/sections/hmt/00b_complejo_discreto_medida_rev2.tex","sha256":"c4bd51e17232daf45c4fe95b5863bfa0f945f2a926d22ec6c9c8eccdfa5f1dd8"},{"path":"manuscrito/sections/hmt/03e_funtor_fase_longitud.tex","sha256":"b1a874c3ecb9a434f01fa3ed05c0cef286ac57a85a9269945a994c08df70e78b"},{"path":"manuscrito/sections/hmt/14f_cierre_genealogico_numero.tex","sha256":"4522230d8297a1420ccb586d0d6c1fefcb08b5e8e2deeda85b11cc73afd8ded5"}],"result_id":"CAT-001","statement_sha256":"cd69268cc61050d66556a6a33bd6fb31ab1fc444f678229ff48d09dfeeb7ec1b"}
% HMT_RESULT_END:CAT-001:residence
% HMT_RESULT_BEGIN:TOP-001:residence
% {"material_control":"GATE_MATERIAL_64+PASS_PUBLICACION_INTEGRAL","material_state":"INTEGRADO_O_REMITIDO_SIN_PERDIDA","residences":[{"path":"manuscrito/sections/integracion_20260812/von_neumann_64/04_zfc_fundacion_ordinales.tex","sha256":"5d81fd31bf31a950f9d068f278e76ef0ebfb37cb8cf6ee5483575ab87f5fddb0"},{"path":"manuscrito/sections/hmt/05_tpk_numero_medicion.tex","sha256":"51deffafc13a6d83bc07e15f77892d1ec32ba05f54cc2b01d6417ea13e3599a2"}],"result_id":"TOP-001","statement_sha256":"cf2c0f501167aac25e8c173e2c9c6cca8aa4cc123f668b06c5a7ca8c42f21eb3"}
% HMT_RESULT_END:TOP-001:residence
% HMT_RESULT_BEGIN:HIL-001:residence
% {"material_control":"GATE_MATERIAL_64+PASS_PUBLICACION_INTEGRAL","material_state":"INTEGRADO_O_REMITIDO_SIN_PERDIDA","residences":[{"path":"manuscrito/sections/integracion_20260812/von_neumann_64/06_hilbert_heisenberg.tex","sha256":"6f35e2e9b56af3cbf33d20c28ae1b5bfac2b92f4a0827f444be39774046e130b"},{"path":"manuscrito/sections/hmt/03l_elevacion_cohomologica_dimension_rev2.tex","sha256":"1b0f6b86dc275a3947b64e319b0b19789d62c83bcc3714d5c0b8a2cea2243a8c"}],"result_id":"HIL-001","statement_sha256":"7529a7dd47eb815a976287f21bd184c40d9cc3dfb0dcb1d2c3f1a2b4b4eeaaa7"}
% HMT_RESULT_END:HIL-001:residence
% HMT_RESULT_BEGIN:UHF-001:residence
% {"material_control":"GATE_MATERIAL_64+PASS_PUBLICACION_INTEGRAL","material_state":"INTEGRADO_O_REMITIDO_SIN_PERDIDA","residences":[{"path":"manuscrito/sections/integracion_20260812/von_neumann_64/07_uhf_factor_ii1.tex","sha256":"a3f413168ac1514605dc7e9bf8a7866e73420380c5a93390fbd5f5d55fa25765"},{"path":"manuscrito/sections/hmt/03l_elevacion_cohomologica_dimension_rev2.tex","sha256":"1b0f6b86dc275a3947b64e319b0b19789d62c83bcc3714d5c0b8a2cea2243a8c"}],"result_id":"UHF-001","statement_sha256":"0cf650bcfd502563a4ff8cab59c8c8c4ae997565ebed1f058034f5cce7e5b378"}
% HMT_RESULT_END:UHF-001:residence
% HMT_RESULT_BEGIN:II1-001:residence
% {"material_control":"GATE_MATERIAL_64+PASS_PUBLICACION_INTEGRAL","material_state":"INTEGRADO_O_REMITIDO_SIN_PERDIDA","residences":[{"path":"manuscrito/sections/integracion_20260812/von_neumann_64/07_uhf_factor_ii1.tex","sha256":"a3f413168ac1514605dc7e9bf8a7866e73420380c5a93390fbd5f5d55fa25765"},{"path":"manuscrito/sections/hmt/03l_elevacion_cohomologica_dimension_rev2.tex","sha256":"1b0f6b86dc275a3947b64e319b0b19789d62c83bcc3714d5c0b8a2cea2243a8c"}],"result_id":"II1-001","statement_sha256":"dc01c84423637eb3704db747bb7a96106332b3cc6f8f36258e1af6cf02c8ead6"}
% HMT_RESULT_END:II1-001:residence
% HMT_RESULT_BEGIN:DIM-001:residence
% {"material_control":"GATE_MATERIAL_64+PASS_PUBLICACION_INTEGRAL","material_state":"INTEGRADO_O_REMITIDO_SIN_PERDIDA","residences":[{"path":"manuscrito/sections/integracion_20260812/von_neumann_64/07_uhf_factor_ii1.tex","sha256":"a3f413168ac1514605dc7e9bf8a7866e73420380c5a93390fbd5f5d55fa25765"},{"path":"manuscrito/sections/hmt/03l_elevacion_cohomologica_dimension_rev2.tex","sha256":"1b0f6b86dc275a3947b64e319b0b19789d62c83bcc3714d5c0b8a2cea2243a8c"}],"result_id":"DIM-001","statement_sha256":"9b8936bc39190cf70a463a94ca6fb311c2d8d8e66262fb2a3d44d831c5653c4b"}
% HMT_RESULT_END:DIM-001:residence
% HMT_RESULT_BEGIN:CLK-001:residence
% {"material_control":"GATE_MATERIAL_64+PASS_PUBLICACION_INTEGRAL","material_state":"INTEGRADO_O_REMITIDO_SIN_PERDIDA","residences":[{"path":"manuscrito/sections/integracion_20260812/von_neumann_64/08_reloj_logica_ergodicidad.tex","sha256":"aeda0eee8e4376ae44e5f6505cf4414ba767a64705a111c66ef2917d7411a0d9"},{"path":"manuscrito/sections/hmt/19_rotaciones_vacancias_criticidad.tex","sha256":"2d46e077e597811a4b0545b2a6ded98b5d56d0d66c65d71f416a1d06c282c363"}],"result_id":"CLK-001","statement_sha256":"b89435d53375a0c540e471c24439fbb5b83ca598038eaa6e60ad19138fe1d106"}
% HMT_RESULT_END:CLK-001:residence
% HMT_RESULT_BEGIN:INT-001:residence
% {"material_control":"GATE_MATERIAL_64+PASS_PUBLICACION_INTEGRAL","material_state":"INTEGRADO_O_REMITIDO_SIN_PERDIDA","residences":[{"path":"manuscrito/sections/integracion_20260812/von_neumann_64/08_reloj_logica_ergodicidad.tex","sha256":"aeda0eee8e4376ae44e5f6505cf4414ba767a64705a111c66ef2917d7411a0d9"}],"result_id":"INT-001","statement_sha256":"aa1ffba490fcc7437280d3a98fe5f9a543b85b037100711af707edb8fdc41527"}
% HMT_RESULT_END:INT-001:residence
% HMT_RESULT_BEGIN:QLG-001:residence
% {"material_control":"GATE_MATERIAL_64+PASS_PUBLICACION_INTEGRAL","material_state":"INTEGRADO_O_REMITIDO_SIN_PERDIDA","residences":[{"path":"manuscrito/sections/integracion_20260812/von_neumann_64/08_reloj_logica_ergodicidad.tex","sha256":"aeda0eee8e4376ae44e5f6505cf4414ba767a64705a111c66ef2917d7411a0d9"},{"path":"manuscrito/sections/integracion_20260812/von_neumann_64/06_hilbert_heisenberg.tex","sha256":"6f35e2e9b56af3cbf33d20c28ae1b5bfac2b92f4a0827f444be39774046e130b"},{"path":"manuscrito/sections/hmt/03_app_ortograma.tex","sha256":"aaf0fd1dec2c270d544cd19562b4461c57ebb009437e051d973be43395428390"}],"result_id":"QLG-001","statement_sha256":"5af886b5ee8e9ef572b0db2f8860b44c9a5685f03b042ad7934083b6157af5d3"}
% HMT_RESULT_END:QLG-001:residence
% HMT_RESULT_BEGIN:ERG-001:residence
% {"material_control":"GATE_MATERIAL_64+PASS_PUBLICACION_INTEGRAL","material_state":"INTEGRADO_O_REMITIDO_SIN_PERDIDA","residences":[{"path":"manuscrito/sections/integracion_20260812/von_neumann_64/08_reloj_logica_ergodicidad.tex","sha256":"aeda0eee8e4376ae44e5f6505cf4414ba767a64705a111c66ef2917d7411a0d9"},{"path":"manuscrito/sections/hmt/19_rotaciones_vacancias_criticidad.tex","sha256":"2d46e077e597811a4b0545b2a6ded98b5d56d0d66c65d71f416a1d06c282c363"}],"result_id":"ERG-001","statement_sha256":"0e5e1bd65b68a7b812b621df955227e4bd7880b25b43edc0e1a9db24e3a80b8e"}
% HMT_RESULT_END:ERG-001:residence
% HMT_RESULT_BEGIN:DEN-001:residence
% {"material_control":"GATE_MATERIAL_64+PASS_PUBLICACION_INTEGRAL","material_state":"INTEGRADO_O_REMITIDO_SIN_PERDIDA","residences":[{"path":"manuscrito/sections/integracion_20260812/von_neumann_64/08_reloj_logica_ergodicidad.tex","sha256":"aeda0eee8e4376ae44e5f6505cf4414ba767a64705a111c66ef2917d7411a0d9"},{"path":"manuscrito/sections/md/04ab_postulados_cuanticos_genealogicos_rev2.tex","sha256":"1895f84bd75b104e873a3d73c5b0e8ae1b5df606bd6feefeb4f9666a92966804"}],"result_id":"DEN-001","statement_sha256":"d3f54d28f9a40e95ef70680ed00ece354f13bef3fa9616bf65771ff984899f55"}
% HMT_RESULT_END:DEN-001:residence
% HMT_RESULT_BEGIN:LAN-001:residence
% {"material_control":"GATE_MATERIAL_64+PASS_PUBLICACION_INTEGRAL","material_state":"INTEGRADO_O_REMITIDO_SIN_PERDIDA","residences":[{"path":"manuscrito/sections/md/04g_informacion_reversible_landauer_rev7.tex","sha256":"3f60a3da983f28f197fe487896eda3667d44247d0ad30f6f3d1a598b368f0e8e"}],"result_id":"LAN-001","statement_sha256":"197b4b23802d24d8e911bcd3787ba07daf3a5694505b103e27acf6a540021b23"}
% HMT_RESULT_END:LAN-001:residence
% HMT_RESULT_BEGIN:LAN-002:residence
% {"material_control":"GATE_MATERIAL_64+PASS_PUBLICACION_INTEGRAL","material_state":"INTEGRADO_O_REMITIDO_SIN_PERDIDA","residences":[{"path":"manuscrito/sections/integracion_20260812/bloque_64_landauer.tex","sha256":"ef8b27ff2b888ef3c5aa3228093c84ed139d63a8e8a504a4b6da8b56f256cbe6"}],"result_id":"LAN-002","statement_sha256":"aa495d4ee76ccc4ec19bbfee8fca847aa39dfb3e7785b0654264e883afa8e9f3"}
% HMT_RESULT_END:LAN-002:residence
% HMT_RESULT_BEGIN:GDL-001:residence
% {"material_control":"GATE_MATERIAL_64+PASS_PUBLICACION_INTEGRAL","material_state":"INTEGRADO_O_REMITIDO_SIN_PERDIDA","residences":[{"path":"manuscrito/sections/integracion_20260812/von_neumann_64/10_godel_turing.tex","sha256":"c17b938966df53455f1927258ae4e4fd2603927ea1827a08620dfdbb386e7863"}],"result_id":"GDL-001","statement_sha256":"e65a071d6f27cf044d0d665fa43c1c77e8b1c800787b0879ee63f693a954b4eb"}
% HMT_RESULT_END:GDL-001:residence
% HMT_RESULT_BEGIN:GDL-002:residence
% {"material_control":"GATE_MATERIAL_64+PASS_PUBLICACION_INTEGRAL","material_state":"INTEGRADO_O_REMITIDO_SIN_PERDIDA","residences":[{"path":"manuscrito/sections/integracion_20260812/von_neumann_64/10_godel_turing.tex","sha256":"c17b938966df53455f1927258ae4e4fd2603927ea1827a08620dfdbb386e7863"}],"result_id":"GDL-002","statement_sha256":"30af7308ca4cbed2f47ea0f0aaf875baaf611eda5be904a33538de7e5ee41864"}
% HMT_RESULT_END:GDL-002:residence
% HMT_RESULT_BEGIN:IND-001:residence
% {"material_control":"GATE_MATERIAL_64+PASS_PUBLICACION_INTEGRAL","material_state":"INTEGRADO_O_REMITIDO_SIN_PERDIDA","residences":[{"path":"manuscrito/sections/integracion_20260812/von_neumann_64/11_forcing_borel_minimax.tex","sha256":"c4b14714d216d63c6db3b930602a3cd7953d40dd61d50a093bd50d840ccfc33c"}],"result_id":"IND-001","statement_sha256":"5e22be0c9938a9dceab95188d043eebbbffab47eced16a62dc2ded6048b7891d"}
% HMT_RESULT_END:IND-001:residence
% HMT_RESULT_BEGIN:FOR-001:residence
% {"material_control":"GATE_MATERIAL_64+PASS_PUBLICACION_INTEGRAL","material_state":"INTEGRADO_O_REMITIDO_SIN_PERDIDA","residences":[{"path":"manuscrito/sections/integracion_20260812/von_neumann_64/11_forcing_borel_minimax.tex","sha256":"c4b14714d216d63c6db3b930602a3cd7953d40dd61d50a093bd50d840ccfc33c"},{"path":"manuscrito/sections/hmt/05_tpk_numero_medicion.tex","sha256":"51deffafc13a6d83bc07e15f77892d1ec32ba05f54cc2b01d6417ea13e3599a2"}],"result_id":"FOR-001","statement_sha256":"1a8d72a91afdb3a1411a9d0431a099066dffdd6a4603d5d3bc5cb570330fa919"}
% HMT_RESULT_END:FOR-001:residence
% HMT_RESULT_BEGIN:FOR-002:residence
% {"material_control":"GATE_MATERIAL_64+PASS_PUBLICACION_INTEGRAL","material_state":"INTEGRADO_O_REMITIDO_SIN_PERDIDA","residences":[{"path":"manuscrito/sections/integracion_20260812/von_neumann_64/11_forcing_borel_minimax.tex","sha256":"c4b14714d216d63c6db3b930602a3cd7953d40dd61d50a093bd50d840ccfc33c"}],"result_id":"FOR-002","statement_sha256":"59f4f58f84a7735d60c6fd2454ba7965e566bd09c58b3d45392f878cb3dec369"}
% HMT_RESULT_END:FOR-002:residence
% HMT_RESULT_BEGIN:BOR-001:residence
% {"material_control":"GATE_MATERIAL_64+PASS_PUBLICACION_INTEGRAL","material_state":"INTEGRADO_O_REMITIDO_SIN_PERDIDA","residences":[{"path":"manuscrito/sections/integracion_20260812/von_neumann_64/11_forcing_borel_minimax.tex","sha256":"c4b14714d216d63c6db3b930602a3cd7953d40dd61d50a093bd50d840ccfc33c"}],"result_id":"BOR-001","statement_sha256":"ac8e973306e90a8e1806a33402b6ef6e087583edf2b49e475fcf64fc7bffc827"}
% HMT_RESULT_END:BOR-001:residence
% HMT_RESULT_BEGIN:MIN-001:residence
% {"material_control":"GATE_MATERIAL_64+PASS_PUBLICACION_INTEGRAL","material_state":"INTEGRADO_O_REMITIDO_SIN_PERDIDA","residences":[{"path":"manuscrito/sections/integracion_20260812/von_neumann_64/11_forcing_borel_minimax.tex","sha256":"c4b14714d216d63c6db3b930602a3cd7953d40dd61d50a093bd50d840ccfc33c"}],"result_id":"MIN-001","statement_sha256":"ed4856dacc07059c530b3404157bcc15541d967d4e205eaba615ed2e7f06bd2e"}
% HMT_RESULT_END:MIN-001:residence
% HMT_RESULT_BEGIN:III-001:residence
% {"material_control":"GATE_MATERIAL_64+PASS_PUBLICACION_INTEGRAL","material_state":"INTEGRADO_O_REMITIDO_SIN_PERDIDA","residences":[{"path":"manuscrito/sections/integracion_20260812/von_neumann_64/08_reloj_logica_ergodicidad.tex","sha256":"aeda0eee8e4376ae44e5f6505cf4414ba767a64705a111c66ef2917d7411a0d9"},{"path":"manuscrito/sections/md/04e_termodinamica_rutas.tex","sha256":"ab1429f9145eee89e5287095b94f327291fa4aa247ae0abcc95e3f5611de87c1"}],"result_id":"III-001","statement_sha256":"ed27269e518d94d2dd09835c4f6d5c6be473aa08906e8ed2d633e2686f5bb4a7"}
% HMT_RESULT_END:III-001:residence
% HMT_RESULT_BEGIN:TRI-001:residence
% {"material_control":"GATE_MATERIAL_64+PASS_PUBLICACION_INTEGRAL","material_state":"INTEGRADO_O_REMITIDO_SIN_PERDIDA","residences":[{"path":"manuscrito/sections/integracion_20260812/von_neumann_64/12_cierre_triple_doble_proyeccion.tex","sha256":"9773623f121e57b194d95202c9b01d24af38a7f1e391fd219beb76ce9dbf2a86"},{"path":"manuscrito/sections/md/10c_puente_retorno_holonomia_cartan.tex","sha256":"b206dc7147f5eb4295dace2d1960e61cb5f6c4b40a0f6828535d1ba4024a7ba9"}],"result_id":"TRI-001","statement_sha256":"7717c12ab00203b9ef7e04dfea02661344e850d6889ac9e04b4906e296f8b383"}
% HMT_RESULT_END:TRI-001:residence
% HMT_RESULT_BEGIN:EXC-001:residence
% {"material_control":"GATE_MATERIAL_64+PASS_PUBLICACION_INTEGRAL","material_state":"INTEGRADO_O_REMITIDO_SIN_PERDIDA","residences":[{"path":"manuscrito/sections/integracion_20260812/von_neumann_64/12_cierre_triple_doble_proyeccion.tex","sha256":"9773623f121e57b194d95202c9b01d24af38a7f1e391fd219beb76ce9dbf2a86"},{"path":"manuscrito/sections/hmt/09_doble_proyeccion.tex","sha256":"ff538add4f3afedca7f4cbd5cea553e3b8132e05117117c40be30e55d33dcdba"},{"path":"manuscrito/sections/hmt/15_bandera_y_rama_excepcional.tex","sha256":"6c42a7012881b036232a5e12c6e2ce382d48964ef9ed12cc0a4d5256126154e5"},{"path":"manuscrito/sections/hmt/15c_cierre_micro_macro_nivel_tres.tex","sha256":"0a5cf2476d9825b9005494bfab1b937deddbba26d5c69f76a7822077be3797c5"},{"path":"manuscrito/sections/hmt/15d_voa_moonshine_orden_dos_rev11.tex","sha256":"6bc88137eb614a0bba834212c9f8f95be09e48e8f570f6c8126ac33dbe007e86"}],"result_id":"EXC-001","statement_sha256":"b9c1e5194c3ddd0cd113377935663265ca76fce3648e3ea5576c0e8b41aa6747"}
% HMT_RESULT_END:EXC-001:residence
% HMT_RESULT_BEGIN:MD-001:residence
% {"material_control":"GATE_MATERIAL_64+PASS_PUBLICACION_INTEGRAL","material_state":"INTEGRADO_O_REMITIDO_SIN_PERDIDA","residences":[{"path":"manuscrito/sections/integracion_20260812/von_neumann_64/12_cierre_triple_doble_proyeccion.tex","sha256":"9773623f121e57b194d95202c9b01d24af38a7f1e391fd219beb76ce9dbf2a86"},{"path":"manuscrito/sections/md/05f_operadores_masa_two_way_rev11.tex","sha256":"e0c576ff54f405c41adea7093ca6b117134e9322749c969f1fc4cca6e8615d35"},{"path":"manuscrito/sections/md/06h_atlas_familias_rev6.tex","sha256":"f6cf1c90f4d42d24598380ca0bff5c9ee66c3269aad819c08d85d32a31066367"}],"result_id":"MD-001","statement_sha256":"3def61d573286af3af73231496757ed80209391421e2f5b759252b1e4f75a582"}
% HMT_RESULT_END:MD-001:residence
% HMT_RESULT_BEGIN:PSC-001:residence
% {"material_control":"GATE_MATERIAL_64+PASS_PUBLICACION_INTEGRAL","material_state":"INTEGRADO_O_REMITIDO_SIN_PERDIDA","residences":[{"path":"manuscrito/sections/integracion_20260812/von_neumann_64/11_forcing_borel_minimax.tex","sha256":"c4b14714d216d63c6db3b930602a3cd7953d40dd61d50a093bd50d840ccfc33c"}],"result_id":"PSC-001","statement_sha256":"aa11f37512474aa438b0ea7a0c6719d2337b671cab88d78f0bfbda69d2d029d4"}
% HMT_RESULT_END:PSC-001:residence
% HMT_RESULT_BEGIN:COV-001:residence
% {"material_control":"GATE_MATERIAL_64+PASS_PUBLICACION_INTEGRAL","material_state":"INTEGRADO_O_REMITIDO_SIN_PERDIDA","residences":[{"path":"/Users/ruben/Documents/HMT2/CONTINUIDAD_EDITORIAL/INTEGRACION_REV3_VON_NEUMANN_MATRICIAL_REALIZACION_2026-08-12/MATRIZ_ATOMICA_AUTOCONTENIDA_REV3_V4.tsv","sha256":"a89bbb749d42ca8e123b0549913409dadda6196187891b4bf7f83b993d2ffa9b"}],"result_id":"COV-001","statement_sha256":"2d2a3767b8f662d7a68bd271a690c4b557d2bc3f3a256077632ca30e518bd4d7"}
% HMT_RESULT_END:COV-001:residence
% HMT_RESULT_BEGIN:IINF-001:residence
% {"material_control":"GATE_MATERIAL_64+PASS_PUBLICACION_INTEGRAL","material_state":"INTEGRADO_O_REMITIDO_SIN_PERDIDA","residences":[{"path":"manuscrito/sections/integracion_20260812/von_neumann_64/07_uhf_factor_ii1.tex","sha256":"a3f413168ac1514605dc7e9bf8a7866e73420380c5a93390fbd5f5d55fa25765"}],"result_id":"IINF-001","statement_sha256":"217587fdc70a6738c237837099cbbba8d41f6d59c40619c866bf48c95e0198d7"}
% HMT_RESULT_END:IINF-001:residence
% HMT_RESULT_BEGIN:DEN-003:residence
% {"material_control":"GATE_MATERIAL_64+PASS_PUBLICACION_INTEGRAL","material_state":"INTEGRADO_O_REMITIDO_SIN_PERDIDA","residences":[{"path":"manuscrito/sections/integracion_20260812/von_neumann_64/08_reloj_logica_ergodicidad.tex","sha256":"aeda0eee8e4376ae44e5f6505cf4414ba767a64705a111c66ef2917d7411a0d9"}],"result_id":"DEN-003","statement_sha256":"921dc3fb0d30a2764747bddc32233eb05b1589bbcbb17f4a98f3aae0ea462a13"}
% HMT_RESULT_END:DEN-003:residence
% HMT_RESULT_BEGIN:CARD-001:residence
% {"material_control":"GATE_MATERIAL_64+PASS_PUBLICACION_INTEGRAL","material_state":"INTEGRADO_O_REMITIDO_SIN_PERDIDA","residences":[{"path":"manuscrito/sections/integracion_20260812/von_neumann_64/04_zfc_fundacion_ordinales.tex","sha256":"5d81fd31bf31a950f9d068f278e76ef0ebfb37cb8cf6ee5483575ab87f5fddb0"},{"path":"manuscrito/sections/hmt/00_construccion_app_gnomon_rev11.tex","sha256":"f027dd8cb64274e9f984a180e01da910cdb1fc6235d856ad6dc736f7477381b9"}],"result_id":"CARD-001","statement_sha256":"657e91ba07505f78793eeac979c673b8bc6ec88f1adc9067345b9c06ef806c89"}
% HMT_RESULT_END:CARD-001:residence
% HMT_RESULT_BEGIN:PWR-003:residence
% {"material_control":"GATE_MATERIAL_64+PASS_PUBLICACION_INTEGRAL","material_state":"INTEGRADO_O_REMITIDO_SIN_PERDIDA","residences":[{"path":"manuscrito/sections/integracion_20260812/von_neumann_64/05_potencia_eleccion.tex","sha256":"cb412bdee19cc8e65181f75c0666a966dc20111f1dafcaee563bf1ee0e3d2684"},{"path":"manuscrito/sections/hmt/09_doble_proyeccion.tex","sha256":"ff538add4f3afedca7f4cbd5cea553e3b8132e05117117c40be30e55d33dcdba"}],"result_id":"PWR-003","statement_sha256":"eaefebdd49a0763208853c604924034d7d4c705c6371b2990f7fec1b29b0032a"}
% HMT_RESULT_END:PWR-003:residence
% HMT_RESULT_BEGIN:WEYL-001:residence
% {"material_control":"GATE_MATERIAL_64+PASS_PUBLICACION_INTEGRAL","material_state":"INTEGRADO_O_REMITIDO_SIN_PERDIDA","residences":[{"path":"manuscrito/sections/integracion_20260812/von_neumann_64/06_hilbert_heisenberg.tex","sha256":"6f35e2e9b56af3cbf33d20c28ae1b5bfac2b92f4a0827f444be39774046e130b"},{"path":"manuscrito/sections/hmt/14_numero_heisenberg_y_d108.tex","sha256":"36c7812d22940e488c5a5434ffb85399d9430a12d32f6a376138c58474b3247b"}],"result_id":"WEYL-001","statement_sha256":"3c9c006467a020b44e1f152943e542f7efd63f1fbafa62c1459d799fb5a7fb8d"}
% HMT_RESULT_END:WEYL-001:residence
% HMT_RESULT_BEGIN:DSC-001:residence
% {"material_control":"GATE_MATERIAL_64+PASS_PUBLICACION_INTEGRAL","material_state":"INTEGRADO_O_REMITIDO_SIN_PERDIDA","residences":[{"path":"manuscrito/sections/md/02_pantalla_accion.tex","sha256":"40796ae0bdab2ff0e8c2cc2f07181c0f5efc77c14aa4cca2d5334e899dc0f1ef"}],"result_id":"DSC-001","statement_sha256":"67da6399420386a7c52ecce6942df44aef3e17b49f2b6c2ea0005b3f95cae4ca"}
% HMT_RESULT_END:DSC-001:residence
% HMT_RESULT_BEGIN:REC-001:residence
% {"material_control":"GATE_MATERIAL_64+PASS_PUBLICACION_INTEGRAL","material_state":"INTEGRADO_O_REMITIDO_SIN_PERDIDA","residences":[{"path":"manuscrito/sections/md/04h_entropias_tipadas_rev2.tex","sha256":"f47cec558bb3f375af069dfde73f82b69bf48d9f168a635630f0e9a25bab3d91"}],"result_id":"REC-001","statement_sha256":"0bd633094af48b759b62a734d6cefa9feb253f605c7052145c392284c4b6a32f"}
% HMT_RESULT_END:REC-001:residence
% HMT_RESULT_BEGIN:LAN-003:residence
% {"material_control":"GATE_MATERIAL_64+PASS_PUBLICACION_INTEGRAL","material_state":"INTEGRADO_O_REMITIDO_SIN_PERDIDA","residences":[{"path":"manuscrito/sections/md/04g_informacion_reversible_landauer_rev7.tex","sha256":"3f60a3da983f28f197fe487896eda3667d44247d0ad30f6f3d1a598b368f0e8e"}],"result_id":"LAN-003","statement_sha256":"2a5b4ef14486d6da04b11ef915d58b57331c845e9d106cf402e883204d561861"}
% HMT_RESULT_END:LAN-003:residence
% HMT_RESULT_BEGIN:FIN-001:residence
% {"material_control":"GATE_MATERIAL_64+PASS_PUBLICACION_INTEGRAL","material_state":"INTEGRADO_O_REMITIDO_SIN_PERDIDA","residences":[{"path":"manuscrito/sections/integracion_20260812/von_neumann_64/10_godel_turing.tex","sha256":"c17b938966df53455f1927258ae4e4fd2603927ea1827a08620dfdbb386e7863"}],"result_id":"FIN-001","statement_sha256":"a73d6ef20cd13e14f35306811ecee9eb3057af5c97ddf492c7ba1f86dd04f1ce"}
% HMT_RESULT_END:FIN-001:residence
% HMT_RESULT_BEGIN:DEN-002:residence
% {"material_control":"GATE_MATERIAL_64+PASS_PUBLICACION_INTEGRAL","material_state":"INTEGRADO_O_REMITIDO_SIN_PERDIDA","residences":[{"path":"manuscrito/sections/integracion_20260812/von_neumann_64/08_reloj_logica_ergodicidad.tex","sha256":"aeda0eee8e4376ae44e5f6505cf4414ba767a64705a111c66ef2917d7411a0d9"}],"result_id":"DEN-002","statement_sha256":"35b2ba3c1b4020b9163f1ef3d443b6bcb565b09c8a8622d543cf8e4f57289896"}
% HMT_RESULT_END:DEN-002:residence
% HMT_RESULT_BEGIN:UNI-001:residence
% {"material_control":"GATE_MATERIAL_64+PASS_PUBLICACION_INTEGRAL","material_state":"INTEGRADO_O_REMITIDO_SIN_PERDIDA","residences":[{"path":"manuscrito/sections/integracion_20260812/von_neumann_64/12_cierre_triple_doble_proyeccion.tex","sha256":"9773623f121e57b194d95202c9b01d24af38a7f1e391fd219beb76ce9dbf2a86"}],"result_id":"UNI-001","statement_sha256":"c937c36037e7f572961ab452004513b8c1c29a3ba0d578a7b0da6590a849a422"}
% HMT_RESULT_END:UNI-001:residence
% HMT_RESULT_BEGIN:UNI-002:residence
% {"material_control":"GATE_MATERIAL_64+PASS_PUBLICACION_INTEGRAL","material_state":"INTEGRADO_O_REMITIDO_SIN_PERDIDA","residences":[{"path":"manuscrito/sections/integracion_20260812/von_neumann_64/12_cierre_triple_doble_proyeccion.tex","sha256":"9773623f121e57b194d95202c9b01d24af38a7f1e391fd219beb76ce9dbf2a86"}],"result_id":"UNI-002","statement_sha256":"6c0f67e5980d3e7df62bb7613076467f3e8c1e71d7bd85ae9b8e5a5a6e6eefac"}
% HMT_RESULT_END:UNI-002:residence
% HMT_RESULT_BEGIN:P01:residence
% {"material_control":"PASS_PUENTE_GENEALOGICO_MATRICIAL","material_state":"INTEGRADO_EN_GRAFO_ACTIVO","residences":[{"path":"manuscrito/sections/integracion_20260812/matricial_52/02_hmt_en_profundidad.tex","sha256":"fe6c23eac09e147199090a091fe2564e5587b7b1155ea85518d2bf47aade4cfc"}],"result_id":"P01","statement_sha256":"53d8dc07ede69f64f0d718b277459101f4824f60da865637ded87662edab5245"}
% HMT_RESULT_END:P01:residence
% HMT_RESULT_BEGIN:P02:residence
% {"material_control":"PASS_PUENTE_GENEALOGICO_MATRICIAL","material_state":"INTEGRADO_EN_GRAFO_ACTIVO","residences":[{"path":"manuscrito/sections/integracion_20260812/matricial_52/02_hmt_en_profundidad.tex","sha256":"fe6c23eac09e147199090a091fe2564e5587b7b1155ea85518d2bf47aade4cfc"}],"result_id":"P02","statement_sha256":"69c00b394120ef47751db68d34eb4b918078692555e4fa5e663b2f0297a8548c"}
% HMT_RESULT_END:P02:residence
% HMT_RESULT_BEGIN:P03:residence
% {"material_control":"PASS_PUENTE_GENEALOGICO_MATRICIAL","material_state":"INTEGRADO_EN_GRAFO_ACTIVO","residences":[{"path":"manuscrito/sections/integracion_20260812/matricial_52/02_hmt_en_profundidad.tex","sha256":"fe6c23eac09e147199090a091fe2564e5587b7b1155ea85518d2bf47aade4cfc"}],"result_id":"P03","statement_sha256":"2fe6041370eb6e955a3d0ba0f2e07811bc0aebd09c6793150785e46a510b9359"}
% HMT_RESULT_END:P03:residence
% HMT_RESULT_BEGIN:P04:residence
% {"material_control":"PASS_PUENTE_GENEALOGICO_MATRICIAL","material_state":"INTEGRADO_EN_GRAFO_ACTIVO","residences":[{"path":"manuscrito/sections/integracion_20260812/matricial_52/02_hmt_en_profundidad.tex","sha256":"fe6c23eac09e147199090a091fe2564e5587b7b1155ea85518d2bf47aade4cfc"}],"result_id":"P04","statement_sha256":"c9c80bd3c7cc1697f5e94b0d7401b0649ab52d4f48fd42943fe4bd715aa71c05"}
% HMT_RESULT_END:P04:residence
% HMT_RESULT_BEGIN:P05:residence
% {"material_control":"PASS_PUENTE_GENEALOGICO_MATRICIAL","material_state":"INTEGRADO_EN_GRAFO_ACTIVO","residences":[{"path":"manuscrito/sections/integracion_20260812/matricial_52/02_hmt_en_profundidad.tex","sha256":"fe6c23eac09e147199090a091fe2564e5587b7b1155ea85518d2bf47aade4cfc"}],"result_id":"P05","statement_sha256":"87558450ea6ada90154488b513914bbcbe56a44e45cb40411baa8054fa0a9b11"}
% HMT_RESULT_END:P05:residence
% HMT_RESULT_BEGIN:P06:residence
% {"material_control":"PASS_PUENTE_GENEALOGICO_MATRICIAL","material_state":"INTEGRADO_EN_GRAFO_ACTIVO","residences":[{"path":"manuscrito/sections/integracion_20260812/matricial_52/02_hmt_en_profundidad.tex","sha256":"fe6c23eac09e147199090a091fe2564e5587b7b1155ea85518d2bf47aade4cfc"}],"result_id":"P06","statement_sha256":"b3609e5998fdaf08a9020b1ff80a65442e473516d2ecf24b5b68474c1f6c94c0"}
% HMT_RESULT_END:P06:residence
% HMT_RESULT_BEGIN:P07:residence
% {"material_control":"PASS_PUENTE_GENEALOGICO_MATRICIAL","material_state":"INTEGRADO_EN_GRAFO_ACTIVO","residences":[{"path":"manuscrito/sections/integracion_20260812/matricial_52/02_hmt_en_profundidad.tex","sha256":"fe6c23eac09e147199090a091fe2564e5587b7b1155ea85518d2bf47aade4cfc"},{"path":"manuscrito/sections/md/04ab_postulados_cuanticos_genealogicos_rev2.tex","sha256":"1895f84bd75b104e873a3d73c5b0e8ae1b5df606bd6feefeb4f9666a92966804"}],"result_id":"P07","statement_sha256":"f8147d53ff5923fc1c43835f30a9f07bf581e5e699eb63dd74ac1414533a629e"}
% HMT_RESULT_END:P07:residence
% HMT_RESULT_BEGIN:P08:residence
% {"material_control":"PASS_PUENTE_GENEALOGICO_MATRICIAL","material_state":"INTEGRADO_EN_GRAFO_ACTIVO","residences":[{"path":"manuscrito/sections/integracion_20260812/matricial_52/02_hmt_en_profundidad.tex","sha256":"fe6c23eac09e147199090a091fe2564e5587b7b1155ea85518d2bf47aade4cfc"}],"result_id":"P08","statement_sha256":"f9a8d4c5ec3b2afef25f422ad348af58e28287d1d33e38fd8b67b7529c84166d"}
% HMT_RESULT_END:P08:residence
% HMT_RESULT_BEGIN:P09:residence
% {"material_control":"PASS_PUENTE_GENEALOGICO_MATRICIAL","material_state":"INTEGRADO_EN_GRAFO_ACTIVO","residences":[{"path":"manuscrito/sections/integracion_20260812/matricial_52/01_objeto_y_diccionario.tex","sha256":"95d254edbb6dd7184d3a5db5ad83006aee1adbbdd4ee0a5f631e89d6300648da"}],"result_id":"P09","statement_sha256":"cb5502d8e34b02758b04dda42755bf0a91a83e8ba03dbc45f8e4453942b6ce77"}
% HMT_RESULT_END:P09:residence
% HMT_RESULT_BEGIN:P10:residence
% {"material_control":"PASS_PUENTE_GENEALOGICO_MATRICIAL","material_state":"INTEGRADO_EN_GRAFO_ACTIVO","residences":[{"path":"manuscrito/sections/integracion_20260812/matricial_52/01_objeto_y_diccionario.tex","sha256":"95d254edbb6dd7184d3a5db5ad83006aee1adbbdd4ee0a5f631e89d6300648da"}],"result_id":"P10","statement_sha256":"42f48b2ff4254875c1d787ae3c1d25e5f5e086049b6b5f8e963f462ceab4b4d9"}
% HMT_RESULT_END:P10:residence
% HMT_RESULT_BEGIN:P11:residence
% {"material_control":"PASS_PUENTE_GENEALOGICO_MATRICIAL","material_state":"INTEGRADO_EN_GRAFO_ACTIVO","residences":[{"path":"manuscrito/sections/integracion_20260812/matricial_52/01_objeto_y_diccionario.tex","sha256":"95d254edbb6dd7184d3a5db5ad83006aee1adbbdd4ee0a5f631e89d6300648da"}],"result_id":"P11","statement_sha256":"bcd1e3fbe8b15b11e4f83caac1d5b8879ef6f7e0ff3881fdd1ffba891a6f32ce"}
% HMT_RESULT_END:P11:residence
% HMT_RESULT_BEGIN:P12:residence
% {"material_control":"PASS_PUENTE_GENEALOGICO_MATRICIAL","material_state":"INTEGRADO_EN_GRAFO_ACTIVO","residences":[{"path":"manuscrito/sections/integracion_20260812/matricial_52/04_cuadrados_cuanticos.tex","sha256":"7131342598b20a77a08f6d9d98e12d41ec47e331372ea6c19754df21a58ab585"}],"result_id":"P12","statement_sha256":"3eca57321c7a0f5f618f0ffad0376d39eb8c8ab5968847fb959fdbb09fc734ab"}
% HMT_RESULT_END:P12:residence
% HMT_RESULT_BEGIN:P13:residence
% {"material_control":"PASS_PUENTE_GENEALOGICO_MATRICIAL","material_state":"INTEGRADO_EN_GRAFO_ACTIVO","residences":[{"path":"manuscrito/sections/integracion_20260812/matricial_52/04_cuadrados_cuanticos.tex","sha256":"7131342598b20a77a08f6d9d98e12d41ec47e331372ea6c19754df21a58ab585"}],"result_id":"P13","statement_sha256":"55db6d54634c091c6187c36369faa76bdf9f508b14a4a3b18f5a348d2c9cf011"}
% HMT_RESULT_END:P13:residence
% HMT_RESULT_BEGIN:P14:residence
% {"material_control":"PASS_PUENTE_GENEALOGICO_MATRICIAL","material_state":"INTEGRADO_EN_GRAFO_ACTIVO","residences":[{"path":"manuscrito/sections/integracion_20260812/matricial_52/03_bigraduacion_ucp.tex","sha256":"c054ec7cb105bc2fda9c3228bfaee808f01db65e38e126a8865a3420ed37864d"}],"result_id":"P14","statement_sha256":"57d9a1d7ea29cb176674b4f73e0033f34f82419195809fa7e8ff1f235d95e410"}
% HMT_RESULT_END:P14:residence
% HMT_RESULT_BEGIN:P15:residence
% {"material_control":"PASS_PUENTE_GENEALOGICO_MATRICIAL","material_state":"INTEGRADO_EN_GRAFO_ACTIVO","residences":[{"path":"manuscrito/sections/integracion_20260812/matricial_52/03_bigraduacion_ucp.tex","sha256":"c054ec7cb105bc2fda9c3228bfaee808f01db65e38e126a8865a3420ed37864d"}],"result_id":"P15","statement_sha256":"83ae0271d5db4482e0f2f1c4ec1fc0a1b3e5bbde57a18fc5a75a65cf03960e6f"}
% HMT_RESULT_END:P15:residence
% HMT_RESULT_BEGIN:P16:residence
% {"material_control":"PASS_PUENTE_GENEALOGICO_MATRICIAL","material_state":"INTEGRADO_EN_GRAFO_ACTIVO","residences":[{"path":"manuscrito/sections/integracion_20260812/matricial_52/03_bigraduacion_ucp.tex","sha256":"c054ec7cb105bc2fda9c3228bfaee808f01db65e38e126a8865a3420ed37864d"}],"result_id":"P16","statement_sha256":"bdf8eff661630a7ba8edf744bf36bded26cbd2417e7ef7e7259b9c6c7505a21e"}
% HMT_RESULT_END:P16:residence
% HMT_RESULT_BEGIN:P17:residence
% {"material_control":"PASS_PUENTE_GENEALOGICO_MATRICIAL","material_state":"INTEGRADO_EN_GRAFO_ACTIVO","residences":[{"path":"manuscrito/sections/integracion_20260812/matricial_52/03_bigraduacion_ucp.tex","sha256":"c054ec7cb105bc2fda9c3228bfaee808f01db65e38e126a8865a3420ed37864d"}],"result_id":"P17","statement_sha256":"b0cce98edf7fa782d1064e6b7ac175d850f4ad264c1b6d8e62288f2568ad45f5"}
% HMT_RESULT_END:P17:residence
% HMT_RESULT_BEGIN:P18:residence
% {"material_control":"PASS_PUENTE_GENEALOGICO_MATRICIAL","material_state":"INTEGRADO_EN_GRAFO_ACTIVO","residences":[{"path":"manuscrito/sections/integracion_20260812/matricial_52/04_cuadrados_cuanticos.tex","sha256":"7131342598b20a77a08f6d9d98e12d41ec47e331372ea6c19754df21a58ab585"}],"result_id":"P18","statement_sha256":"1473f4abe925a8e28298eb8bf1dfd39fb9af561154966420712347199f61b2ce"}
% HMT_RESULT_END:P18:residence
% HMT_RESULT_BEGIN:P19:residence
% {"material_control":"PASS_PUENTE_GENEALOGICO_MATRICIAL","material_state":"INTEGRADO_EN_GRAFO_ACTIVO","residences":[{"path":"manuscrito/sections/integracion_20260812/matricial_52/04_cuadrados_cuanticos.tex","sha256":"7131342598b20a77a08f6d9d98e12d41ec47e331372ea6c19754df21a58ab585"}],"result_id":"P19","statement_sha256":"fd3e9565106ad07b0ffeff64653cdd8f4e25c01bd8125c56b954e4ad58015bce"}
% HMT_RESULT_END:P19:residence
% HMT_RESULT_BEGIN:P20:residence
% {"material_control":"PASS_PUENTE_GENEALOGICO_MATRICIAL","material_state":"INTEGRADO_EN_GRAFO_ACTIVO","residences":[{"path":"manuscrito/sections/integracion_20260812/matricial_52/04_cuadrados_cuanticos.tex","sha256":"7131342598b20a77a08f6d9d98e12d41ec47e331372ea6c19754df21a58ab585"}],"result_id":"P20","statement_sha256":"f0301b81ceca519949c82bef3e2c9c9d0415d1943aadd55e64bafd83dc369e25"}
% HMT_RESULT_END:P20:residence
% HMT_RESULT_BEGIN:P21:residence
% {"material_control":"PASS_PUENTE_GENEALOGICO_MATRICIAL","material_state":"INTEGRADO_EN_GRAFO_ACTIVO","residences":[{"path":"manuscrito/sections/integracion_20260812/matricial_52/08_sistemas_conicos_y_universalidad.tex","sha256":"2b94acee5d5df46c4c6e2c4b9048962de8961703fc585df45d08a42f9f4ae0d2"}],"result_id":"P21","statement_sha256":"9677a2892d35d3020f115f04121dc2d8f34f6bcda52cdb28ac4ad708ca6fcfe2"}
% HMT_RESULT_END:P21:residence
% HMT_RESULT_BEGIN:P22:residence
% {"material_control":"PASS_PUENTE_GENEALOGICO_MATRICIAL","material_state":"INTEGRADO_EN_GRAFO_ACTIVO","residences":[{"path":"manuscrito/sections/integracion_20260812/matricial_52/05_geometria_simplectica_y_lectores.tex","sha256":"087482ebc4a880dba1096a9364dab1a05b0f0bddd3495b6f5b967fb9ccb79ccd"}],"result_id":"P22","statement_sha256":"f79d8ca5659786ae2dfdcb69f4a77431e62f85c8e159a0cb789c16dd900f0e08"}
% HMT_RESULT_END:P22:residence
% HMT_RESULT_BEGIN:P23:residence
% {"material_control":"PASS_PUENTE_GENEALOGICO_MATRICIAL","material_state":"INTEGRADO_EN_GRAFO_ACTIVO","residences":[{"path":"manuscrito/sections/integracion_20260812/matricial_52/05_geometria_simplectica_y_lectores.tex","sha256":"087482ebc4a880dba1096a9364dab1a05b0f0bddd3495b6f5b967fb9ccb79ccd"}],"result_id":"P23","statement_sha256":"4ebed4b0f0e733875915bad691b12c053780dfaa50e1f0f4898837314d3edf9a"}
% HMT_RESULT_END:P23:residence
% HMT_RESULT_BEGIN:P24:residence
% {"material_control":"PASS_PUENTE_GENEALOGICO_MATRICIAL","material_state":"INTEGRADO_EN_GRAFO_ACTIVO","residences":[{"path":"manuscrito/sections/integracion_20260812/matricial_52/06_operatorios_doble_proyeccion.tex","sha256":"0f183ed734248c42a405ba89f43a6ac6085d705a188f911c44ed76e8eb4bf745"}],"result_id":"P24","statement_sha256":"6196ff6b2fa346db4f4fd64dc27219eba3734ae04559393e65a5ad14e3c297e0"}
% HMT_RESULT_END:P24:residence
% HMT_RESULT_BEGIN:P25:residence
% {"material_control":"PASS_PUENTE_GENEALOGICO_MATRICIAL","material_state":"INTEGRADO_EN_GRAFO_ACTIVO","residences":[{"path":"manuscrito/sections/integracion_20260812/matricial_52/06_operatorios_doble_proyeccion.tex","sha256":"0f183ed734248c42a405ba89f43a6ac6085d705a188f911c44ed76e8eb4bf745"}],"result_id":"P25","statement_sha256":"339d4b25f9c31d5aa693c11a2377a405ec7ba754affe3143c351475a2d50885d"}
% HMT_RESULT_END:P25:residence
% HMT_RESULT_BEGIN:P26:residence
% {"material_control":"PASS_PUENTE_GENEALOGICO_MATRICIAL","material_state":"INTEGRADO_EN_GRAFO_ACTIVO","residences":[{"path":"manuscrito/sections/integracion_20260812/matricial_52/06_operatorios_doble_proyeccion.tex","sha256":"0f183ed734248c42a405ba89f43a6ac6085d705a188f911c44ed76e8eb4bf745"}],"result_id":"P26","statement_sha256":"1f0955a242b45acf69644be639ccca7ceac39f2d83009c4aafcdb1172802a601"}
% HMT_RESULT_END:P26:residence
% HMT_RESULT_BEGIN:P27:residence
% {"material_control":"PASS_PUENTE_GENEALOGICO_MATRICIAL","material_state":"INTEGRADO_EN_GRAFO_ACTIVO","residences":[{"path":"manuscrito/sections/integracion_20260812/matricial_52/06_operatorios_doble_proyeccion.tex","sha256":"0f183ed734248c42a405ba89f43a6ac6085d705a188f911c44ed76e8eb4bf745"}],"result_id":"P27","statement_sha256":"60ff5f6df5943b80338a53b54928755aa4cae85b90764d63a153d6ad50dfed24"}
% HMT_RESULT_END:P27:residence
% HMT_RESULT_BEGIN:P28:residence
% {"material_control":"PASS_PUENTE_GENEALOGICO_MATRICIAL","material_state":"INTEGRADO_EN_GRAFO_ACTIVO","residences":[{"path":"manuscrito/sections/integracion_20260812/matricial_52/07_continuo_tensores_y_memoria.tex","sha256":"3df0b8a366a6458f7acf2199859a25f5b40379b3f9039ccaa0b959e2f25e36a5"}],"result_id":"P28","statement_sha256":"de45d3b80b483962a79b544ce428c4298718a36362562637d064386d09251da4"}
% HMT_RESULT_END:P28:residence
% HMT_RESULT_BEGIN:P29:residence
% {"material_control":"PASS_PUENTE_GENEALOGICO_MATRICIAL","material_state":"INTEGRADO_EN_GRAFO_ACTIVO","residences":[{"path":"manuscrito/sections/integracion_20260812/matricial_52/07_continuo_tensores_y_memoria.tex","sha256":"3df0b8a366a6458f7acf2199859a25f5b40379b3f9039ccaa0b959e2f25e36a5"}],"result_id":"P29","statement_sha256":"eed86be2eeabf3a9de1c3150b3fb89bdee7fc9a0e51ea1e9052cfab634322586"}
% HMT_RESULT_END:P29:residence
% HMT_RESULT_BEGIN:P30:residence
% {"material_control":"PASS_PUENTE_GENEALOGICO_MATRICIAL","material_state":"INTEGRADO_EN_GRAFO_ACTIVO","residences":[{"path":"manuscrito/sections/integracion_20260812/matricial_52/09_metrologia_y_programa.tex","sha256":"d5cf28b96e5cde3577601d038089754378d0eb0e6118b134ed2b27e98004a1e7"}],"result_id":"P30","statement_sha256":"16dae690fc342fd1d4bbb3ee5c7067af92476cfdcc61842908df31b55a844797"}
% HMT_RESULT_END:P30:residence
% HMT_RESULT_BEGIN:P31:residence
% {"material_control":"PASS_PUENTE_GENEALOGICO_MATRICIAL","material_state":"INTEGRADO_EN_GRAFO_ACTIVO","residences":[{"path":"manuscrito/sections/integracion_20260812/matricial_52/08_sistemas_conicos_y_universalidad.tex","sha256":"2b94acee5d5df46c4c6e2c4b9048962de8961703fc585df45d08a42f9f4ae0d2"}],"result_id":"P31","statement_sha256":"cabf80d0abd369646b4013a8ac2aad89b7f4c6cd7e0b0092d5f4e07d76e74471"}
% HMT_RESULT_END:P31:residence
% HMT_RESULT_BEGIN:P32:residence
% {"material_control":"PASS_PUENTE_GENEALOGICO_MATRICIAL","material_state":"INTEGRADO_EN_GRAFO_ACTIVO","residences":[{"path":"manuscrito/sections/integracion_20260812/matricial_52/08_sistemas_conicos_y_universalidad.tex","sha256":"2b94acee5d5df46c4c6e2c4b9048962de8961703fc585df45d08a42f9f4ae0d2"}],"result_id":"P32","statement_sha256":"d808a7bd125136593c39a42df20427716bf90ad4b1ac4a033227a3c4e370c64a"}
% HMT_RESULT_END:P32:residence
% HMT_RESULT_BEGIN:QPM-001:residence
% {"material_control":"PASS_PUENTE_GENEALOGICO_MATRICIAL","material_state":"INTEGRADO_EN_GRAFO_ACTIVO","residences":[{"path":"manuscrito/sections/integracion_20260812/matricial_52/01_objeto_y_diccionario.tex","sha256":"95d254edbb6dd7184d3a5db5ad83006aee1adbbdd4ee0a5f631e89d6300648da"}],"result_id":"QPM-001","statement_sha256":"876638c74d1df52e66cf8b186c110f2cbb879465ca4aed51526e0ca3b191405f"}
% HMT_RESULT_END:QPM-001:residence
% HMT_RESULT_BEGIN:QLS-001:residence
% {"material_control":"PASS_PUENTE_GENEALOGICO_MATRICIAL","material_state":"INTEGRADO_EN_GRAFO_ACTIVO","residences":[{"path":"manuscrito/sections/integracion_20260812/matricial_52/01_objeto_y_diccionario.tex","sha256":"95d254edbb6dd7184d3a5db5ad83006aee1adbbdd4ee0a5f631e89d6300648da"}],"result_id":"QLS-001","statement_sha256":"60abad49376bb9a2df1df2851bfc068d7fd0e59b90f09c5812f8ac93138d5308"}
% HMT_RESULT_END:QLS-001:residence
% HMT_RESULT_BEGIN:QLSI-001:residence
% {"material_control":"PASS_PUENTE_GENEALOGICO_MATRICIAL","material_state":"INTEGRADO_EN_GRAFO_ACTIVO","residences":[{"path":"manuscrito/sections/integracion_20260812/matricial_52/01_objeto_y_diccionario.tex","sha256":"95d254edbb6dd7184d3a5db5ad83006aee1adbbdd4ee0a5f631e89d6300648da"}],"result_id":"QLSI-001","statement_sha256":"3cc12a3dcfe1865a34b892c489519c5c02e989e8075c94e9782f1e01e67b2c1d"}
% HMT_RESULT_END:QLSI-001:residence
% HMT_RESULT_BEGIN:QMSI-001:residence
% {"material_control":"PASS_PUENTE_GENEALOGICO_MATRICIAL","material_state":"INTEGRADO_EN_GRAFO_ACTIVO","residences":[{"path":"manuscrito/sections/integracion_20260812/matricial_52/01_objeto_y_diccionario.tex","sha256":"95d254edbb6dd7184d3a5db5ad83006aee1adbbdd4ee0a5f631e89d6300648da"}],"result_id":"QMSI-001","statement_sha256":"641653814fe7d90a81405d016aa63d2840bdf02215f16c828afe4ece7989c37b"}
% HMT_RESULT_END:QMSI-001:residence
% HMT_RESULT_BEGIN:QMSI-002:residence
% {"material_control":"PASS_PUENTE_GENEALOGICO_MATRICIAL","material_state":"INTEGRADO_EN_GRAFO_ACTIVO","residences":[{"path":"manuscrito/sections/integracion_20260812/matricial_52/01_objeto_y_diccionario.tex","sha256":"95d254edbb6dd7184d3a5db5ad83006aee1adbbdd4ee0a5f631e89d6300648da"}],"result_id":"QMSI-002","statement_sha256":"79afee550d7e55837c8ed6a9067de90d22cfee54bbff7c98636ded289843016f"}
% HMT_RESULT_END:QMSI-002:residence
% HMT_RESULT_BEGIN:QMSI-003:residence
% {"material_control":"PASS_PUENTE_GENEALOGICO_MATRICIAL","material_state":"INTEGRADO_EN_GRAFO_ACTIVO","residences":[{"path":"manuscrito/sections/integracion_20260812/matricial_52/01_objeto_y_diccionario.tex","sha256":"95d254edbb6dd7184d3a5db5ad83006aee1adbbdd4ee0a5f631e89d6300648da"}],"result_id":"QMSI-003","statement_sha256":"1bbe978d6a9ba0ec29aa916bb4b2a72df5bf3936b89bf14a66b2916e1d4fd090"}
% HMT_RESULT_END:QMSI-003:residence
% HMT_RESULT_BEGIN:QBR-001:residence
% {"material_control":"PASS_PUENTE_GENEALOGICO_MATRICIAL","material_state":"INTEGRADO_EN_GRAFO_ACTIVO","residences":[{"path":"manuscrito/sections/integracion_20260812/matricial_52/02_hmt_en_profundidad.tex","sha256":"fe6c23eac09e147199090a091fe2564e5587b7b1155ea85518d2bf47aade4cfc"},{"path":"manuscrito/sections/md/04ab_postulados_cuanticos_genealogicos_rev2.tex","sha256":"1895f84bd75b104e873a3d73c5b0e8ae1b5df606bd6feefeb4f9666a92966804"}],"result_id":"QBR-001","statement_sha256":"38ad3305cd9c643e9067d133a541df231165278ab3b01e906f494ec45be28104"}
% HMT_RESULT_END:QBR-001:residence
% HMT_RESULT_BEGIN:QDY-001:residence
% {"material_control":"PASS_PUENTE_GENEALOGICO_MATRICIAL","material_state":"INTEGRADO_EN_GRAFO_ACTIVO","residences":[{"path":"manuscrito/sections/integracion_20260812/matricial_52/02_hmt_en_profundidad.tex","sha256":"fe6c23eac09e147199090a091fe2564e5587b7b1155ea85518d2bf47aade4cfc"},{"path":"manuscrito/sections/md/04ab_postulados_cuanticos_genealogicos_rev2.tex","sha256":"1895f84bd75b104e873a3d73c5b0e8ae1b5df606bd6feefeb4f9666a92966804"}],"result_id":"QDY-001","statement_sha256":"bffa8460b3d4625a598b73816f5f9e8a9218458a12b69fd3cc36c8ba18d69a89"}
% HMT_RESULT_END:QDY-001:residence
% HMT_RESULT_BEGIN:QKR-001:residence
% {"material_control":"PASS_PUENTE_GENEALOGICO_MATRICIAL","material_state":"INTEGRADO_EN_GRAFO_ACTIVO","residences":[{"path":"manuscrito/sections/integracion_20260812/matricial_52/02_hmt_en_profundidad.tex","sha256":"fe6c23eac09e147199090a091fe2564e5587b7b1155ea85518d2bf47aade4cfc"},{"path":"manuscrito/sections/md/01a_medicion_cuantica_genealogica_rev2.tex","sha256":"725bbfb625e4bd74c3f5845dc2d3ad117b173fac63f11d300f7f6b8344fb505a"}],"result_id":"QKR-001","statement_sha256":"e458e06f5341341d7cbbeee4dbcc1fa333c7c1d293266aec038128fe9c334858"}
% HMT_RESULT_END:QKR-001:residence
% HMT_RESULT_BEGIN:QMO-001:residence
% {"material_control":"PASS_PUENTE_GENEALOGICO_MATRICIAL","material_state":"INTEGRADO_EN_GRAFO_ACTIVO","residences":[{"path":"manuscrito/sections/integracion_20260812/matricial_52/02_hmt_en_profundidad.tex","sha256":"fe6c23eac09e147199090a091fe2564e5587b7b1155ea85518d2bf47aade4cfc"},{"path":"manuscrito/sections/md/04ab_postulados_cuanticos_genealogicos_rev2.tex","sha256":"1895f84bd75b104e873a3d73c5b0e8ae1b5df606bd6feefeb4f9666a92966804"}],"result_id":"QMO-001","statement_sha256":"14d863637e09c435b127f3b5fc5dccc9f9d873d837809884a8efc4b2da2953be"}
% HMT_RESULT_END:QMO-001:residence
% HMT_RESULT_BEGIN:DIC-001:residence
% {"material_control":"PASS_PUENTE_GENEALOGICO_MATRICIAL","material_state":"INTEGRADO_EN_GRAFO_ACTIVO","residences":[{"path":"manuscrito/sections/integracion_20260812/matricial_52/01_objeto_y_diccionario.tex","sha256":"95d254edbb6dd7184d3a5db5ad83006aee1adbbdd4ee0a5f631e89d6300648da"}],"result_id":"DIC-001","statement_sha256":"51d602fff6bbd77d754ba0ebe0732f3998ac22c81a82514440408ff0f6545d30"}
% HMT_RESULT_END:DIC-001:residence
% HMT_RESULT_BEGIN:QMC-001:residence
% {"material_control":"PASS_PUENTE_GENEALOGICO_MATRICIAL","material_state":"INTEGRADO_EN_GRAFO_ACTIVO","residences":[{"path":"manuscrito/sections/integracion_20260812/matricial_52/04_cuadrados_cuanticos.tex","sha256":"7131342598b20a77a08f6d9d98e12d41ec47e331372ea6c19754df21a58ab585"}],"result_id":"QMC-001","statement_sha256":"900c10fc899cb62a7793541419f3c6ce103296ab9690d816954062e980960504"}
% HMT_RESULT_END:QMC-001:residence
% HMT_RESULT_BEGIN:QDT-001:residence
% {"material_control":"PASS_PUENTE_GENEALOGICO_MATRICIAL","material_state":"INTEGRADO_EN_GRAFO_ACTIVO","residences":[{"path":"manuscrito/sections/integracion_20260812/matricial_52/05_geometria_simplectica_y_lectores.tex","sha256":"087482ebc4a880dba1096a9364dab1a05b0f0bddd3495b6f5b967fb9ccb79ccd"}],"result_id":"QDT-001","statement_sha256":"a42d2bae56ce5410c45794269f3181c68a38ea3fb2d5861cebb013294580287d"}
% HMT_RESULT_END:QDT-001:residence
% HMT_RESULT_BEGIN:ENV-001:residence
% {"material_control":"PASS_PUENTE_GENEALOGICO_MATRICIAL","material_state":"INTEGRADO_EN_GRAFO_ACTIVO","residences":[{"path":"manuscrito/sections/integracion_20260812/matricial_52/03_bigraduacion_ucp.tex","sha256":"c054ec7cb105bc2fda9c3228bfaee808f01db65e38e126a8865a3420ed37864d"}],"result_id":"ENV-001","statement_sha256":"a2c7c6251ce55a5d5d21e2825a8eecad64f3da571d4130d8119999126371726c"}
% HMT_RESULT_END:ENV-001:residence
% HMT_RESULT_BEGIN:RNG-001:residence
% {"material_control":"PASS_PUENTE_GENEALOGICO_MATRICIAL","material_state":"INTEGRADO_EN_GRAFO_ACTIVO","residences":[{"path":"manuscrito/sections/integracion_20260812/matricial_52/06_operatorios_doble_proyeccion.tex","sha256":"0f183ed734248c42a405ba89f43a6ac6085d705a188f911c44ed76e8eb4bf745"}],"result_id":"RNG-001","statement_sha256":"6f564d18bb3910cde9dd86c3bcc20cf5be7d388abb45c8acebdc31553910900c"}
% HMT_RESULT_END:RNG-001:residence
% HMT_RESULT_BEGIN:DYN-001:residence
% {"material_control":"PASS_PUENTE_GENEALOGICO_MATRICIAL","material_state":"INTEGRADO_EN_GRAFO_ACTIVO","residences":[{"path":"manuscrito/sections/integracion_20260812/matricial_52/07_continuo_tensores_y_memoria.tex","sha256":"3df0b8a366a6458f7acf2199859a25f5b40379b3f9039ccaa0b959e2f25e36a5"}],"result_id":"DYN-001","statement_sha256":"931716fef6fe78eab6afb720f6b8621e7cf08908dfd2b425ac48605bb361219a"}
% HMT_RESULT_END:DYN-001:residence
% HMT_RESULT_BEGIN:TEN-001:residence
% {"material_control":"PASS_PUENTE_GENEALOGICO_MATRICIAL","material_state":"INTEGRADO_EN_GRAFO_ACTIVO","residences":[{"path":"manuscrito/sections/integracion_20260812/matricial_52/07_continuo_tensores_y_memoria.tex","sha256":"3df0b8a366a6458f7acf2199859a25f5b40379b3f9039ccaa0b959e2f25e36a5"}],"result_id":"TEN-001","statement_sha256":"458096f924295967751d528c1f5b73f3ae28b5e1cec81ab6ac3036eda48c54fd"}
% HMT_RESULT_END:TEN-001:residence
% HMT_RESULT_BEGIN:CHP-001:residence
% {"material_control":"PASS_PUENTE_GENEALOGICO_MATRICIAL","material_state":"INTEGRADO_EN_GRAFO_ACTIVO","residences":[{"path":"manuscrito/sections/integracion_20260812/matricial_52/08_sistemas_conicos_y_universalidad.tex","sha256":"2b94acee5d5df46c4c6e2c4b9048962de8961703fc585df45d08a42f9f4ae0d2"}],"result_id":"CHP-001","statement_sha256":"8fce51f0c85e0120694a493a3cd73f2edb220a7f02de5bc0d820e4697456ff1f"}
% HMT_RESULT_END:CHP-001:residence
% HMT_RESULT_BEGIN:SDP-001:residence
% {"material_control":"PASS_PUENTE_GENEALOGICO_MATRICIAL","material_state":"INTEGRADO_EN_GRAFO_ACTIVO","residences":[{"path":"manuscrito/sections/integracion_20260812/matricial_52/08_sistemas_conicos_y_universalidad.tex","sha256":"2b94acee5d5df46c4c6e2c4b9048962de8961703fc585df45d08a42f9f4ae0d2"}],"result_id":"SDP-001","statement_sha256":"0fa05be883df7a0e3a59ceb8d9386d4775fc4456de540e0b4fa29e7d7716aa46"}
% HMT_RESULT_END:SDP-001:residence
% HMT_RESULT_BEGIN:EQV-001:residence
% {"material_control":"PASS_PUENTE_GENEALOGICO_MATRICIAL","material_state":"INTEGRADO_EN_GRAFO_ACTIVO","residences":[{"path":"manuscrito/sections/integracion_20260812/matricial_52/05_geometria_simplectica_y_lectores.tex","sha256":"087482ebc4a880dba1096a9364dab1a05b0f0bddd3495b6f5b967fb9ccb79ccd"}],"result_id":"EQV-001","statement_sha256":"0ef0dde22fd08d35555aa0a88a1eb4fa608e3b298286a47559fb223af824f79f"}
% HMT_RESULT_END:EQV-001:residence
% HMT_RESULT_BEGIN:QMT-001:residence
% {"material_control":"PASS_PUENTE_GENEALOGICO_MATRICIAL","material_state":"INTEGRADO_EN_GRAFO_ACTIVO","residences":[{"path":"manuscrito/sections/integracion_20260812/matricial_52/04_cuadrados_cuanticos.tex","sha256":"7131342598b20a77a08f6d9d98e12d41ec47e331372ea6c19754df21a58ab585"}],"result_id":"QMT-001","statement_sha256":"e40aff576b435e19d95904fbf6c160fedf5d8f5bcbe4634eab9591b8c6960630"}
% HMT_RESULT_END:QMT-001:residence
% HMT_RESULT_BEGIN:TRN-001:residence
% {"material_control":"PASS_PUENTE_GENEALOGICO_MATRICIAL","material_state":"INTEGRADO_EN_GRAFO_ACTIVO","residences":[{"path":"manuscrito/sections/integracion_20260812/matricial_52/09_metrologia_y_programa.tex","sha256":"d5cf28b96e5cde3577601d038089754378d0eb0e6118b134ed2b27e98004a1e7"}],"result_id":"TRN-001","statement_sha256":"076c0e63737e5420b981d0d4223a441dd2be4af3e62fe56af2608779d7c60bb0"}
% HMT_RESULT_END:TRN-001:residence
% HMT_RESULT_BEGIN:MPSH-001:residence
% {"material_control":"PASS_PUENTE_GENEALOGICO_MATRICIAL","material_state":"INTEGRADO_EN_GRAFO_ACTIVO","residences":[{"path":"manuscrito/sections/integracion_20260812/matricial_52/07_continuo_tensores_y_memoria.tex","sha256":"3df0b8a366a6458f7acf2199859a25f5b40379b3f9039ccaa0b959e2f25e36a5"}],"result_id":"MPSH-001","statement_sha256":"e8ee9870050a98d998dcfa91b8f581a5626e54e41b516743f9666140829b1f9e"}
% HMT_RESULT_END:MPSH-001:residence
% HMT_RESULT_BEGIN:UNIIMP-001:residence
% {"material_control":"PASS_PUENTE_GENEALOGICO_MATRICIAL","material_state":"INTEGRADO_EN_GRAFO_ACTIVO","residences":[{"path":"manuscrito/sections/integracion_20260812/matricial_52/08_sistemas_conicos_y_universalidad.tex","sha256":"2b94acee5d5df46c4c6e2c4b9048962de8961703fc585df45d08a42f9f4ae0d2"}],"result_id":"UNIIMP-001","statement_sha256":"342e5b5e49f7e76a154e4376612d523325fc967212ee3b2f4156d578ed7818be"}
% HMT_RESULT_END:UNIIMP-001:residence
% HMT_RESULT_BEGIN:PRG-001:residence
% {"material_control":"PASS_PUENTE_GENEALOGICO_MATRICIAL","material_state":"INTEGRADO_EN_GRAFO_ACTIVO","residences":[{"path":"manuscrito/sections/integracion_20260812/matricial_52/09_metrologia_y_programa.tex","sha256":"d5cf28b96e5cde3577601d038089754378d0eb0e6118b134ed2b27e98004a1e7"}],"result_id":"PRG-001","statement_sha256":"42475f38a74d5f5c373f586a656ec2ec049a0f9fd5b8ed896dfb8bd11bf3aec5"}
% HMT_RESULT_END:PRG-001:residence
% HMT_RESULT_BEGIN:QCT-001:residence
% {"material_control":"HASH_REGENERACION+NON_RESULTS_6+SCOPE_LIMITS_5","material_state":"CONTROL_FOCAL_ACTIVO","residences":[{"path":"manuscrito/sections/integracion_20260812/matricial_52/A_pruebas_y_certificados.tex","sha256":"ba316b48d94c537b548e76a95774f361aab946f725468bb6acc0700fbe9b0cdb"},{"path":"pruebas/integracion_20260812/verificar_puente_genealogico_matricial.py","sha256":"4661187f03768de0b87fd454620287e9bb407a9402fa6d7e43e6d2969dd2dc69"},{"path":"certificados/integracion_20260812/puente_genealogico_matricial.json","sha256":"fd45cc4e1e1d12225874e214e27fb648e2cea586a5162ba49afdabb981216683"}],"result_id":"QCT-001","statement_sha256":"6a04ea7ae52a0c3f2571d9787272c10ddc18456acf94ddfef7adef2764393d8c"}
% HMT_RESULT_END:QCT-001:residence
% HMT_RESULT_BEGIN:PRV-001:residence
% {"material_control":"HASH_REGENERACION+NON_RESULTS_6+SCOPE_LIMITS_5","material_state":"CONTROL_FOCAL_ACTIVO","residences":[{"path":"manuscrito/sections/integracion_20260812/matricial_52/C_fuentes_y_bibliografia.tex","sha256":"ba54c19f8729fa9e16033a263c3409a80a6fbb6788dfddcb037a657c14f598ff"},{"path":"pruebas/integracion_20260812/verificar_puente_genealogico_matricial.py","sha256":"4661187f03768de0b87fd454620287e9bb407a9402fa6d7e43e6d2969dd2dc69"},{"path":"certificados/integracion_20260812/puente_genealogico_matricial.json","sha256":"fd45cc4e1e1d12225874e214e27fb648e2cea586a5162ba49afdabb981216683"}],"result_id":"PRV-001","statement_sha256":"3898b9fa7cd9ca405abd8eb49767227624e28b2704db43c64ed70e7a0061bf68"}
% HMT_RESULT_END:PRV-001:residence
% HMT_RESULT_BEGIN:MOR-001:residence
% {"material_control":"GATE_MATERIAL_68_22_22+AUDITORIA_ESTATICA","material_state":"INTEGRADO_O_REMITIDO_SIN_PERDIDA","residences":[{"path":"manuscrito/sections/hmt/14f_cierre_genealogico_numero.tex","sha256":"4522230d8297a1420ccb586d0d6c1fefcb08b5e8e2deeda85b11cc73afd8ded5"}],"result_id":"MOR-001","statement_sha256":"81ec51bdba281285df6bfdc1ea9bb82c7f59300a8a07d463fa438037c958c27a"}
% HMT_RESULT_END:MOR-001:residence
% HMT_RESULT_BEGIN:MOR-002:residence
% {"material_control":"GATE_MATERIAL_68_22_22+AUDITORIA_ESTATICA","material_state":"INTEGRADO_O_REMITIDO_SIN_PERDIDA","residences":[{"path":"manuscrito/sections/hmt/03_app_ortograma.tex","sha256":"aaf0fd1dec2c270d544cd19562b4461c57ebb009437e051d973be43395428390"}],"result_id":"MOR-002","statement_sha256":"4d2f31444c015ef52d258636bae153f77c798e12f034cb624fdf8d1721eda040"}
% HMT_RESULT_END:MOR-002:residence
% HMT_RESULT_BEGIN:MOR-003:residence
% {"material_control":"GATE_MATERIAL_68_22_22+AUDITORIA_ESTATICA","material_state":"INTEGRADO_O_REMITIDO_SIN_PERDIDA","residences":[{"path":"manuscrito/sections/hmt/02_trit_geometrias.tex","sha256":"aefb5c1ec6732534c810d0021834eff92a662f59072b97c35c175a6ce7952fa6"}],"result_id":"MOR-003","statement_sha256":"56c13700ad07f59cb91c4278006ecbd0864a2c64b2bc08c14219693a1d5c1961"}
% HMT_RESULT_END:MOR-003:residence
% HMT_RESULT_BEGIN:MOR-004:residence
% {"material_control":"GATE_MATERIAL_68_22_22+AUDITORIA_ESTATICA","material_state":"INTEGRADO_O_REMITIDO_SIN_PERDIDA","residences":[{"path":"manuscrito/sections/hmt/03k_genealogia_operaciones_rev2.tex","sha256":"fae7aca94c6fb20754431ed79a5aa2f005a7cf2f5681736fc5826da4c694fb98"}],"result_id":"MOR-004","statement_sha256":"e709d75fe6d60011165caa4daee14d1c7b07aa4b57f2395a2ef27267672e3448"}
% HMT_RESULT_END:MOR-004:residence
% HMT_RESULT_BEGIN:MOR-005:residence
% {"material_control":"GATE_MATERIAL_68_22_22+AUDITORIA_ESTATICA","material_state":"INTEGRADO_O_REMITIDO_SIN_PERDIDA","residences":[{"path":"manuscrito/sections/integracion_20260812/realizacion_68/01_cilindros_kronecker_dirac.tex","sha256":"782450c6fc1e432641df8b3cc605425c604109f1db5830873b0b778c59ae6946"}],"result_id":"MOR-005","statement_sha256":"39a6a2f053bee151e082142609deae0baa2067c95df7e2cc57eabe0e0cdb0b35"}
% HMT_RESULT_END:MOR-005:residence
% HMT_RESULT_BEGIN:MOR-006:residence
% {"material_control":"GATE_MATERIAL_68_22_22+AUDITORIA_ESTATICA","material_state":"INTEGRADO_O_REMITIDO_SIN_PERDIDA","residences":[{"path":"manuscrito/sections/integracion_20260812/realizacion_68/01_cilindros_kronecker_dirac.tex","sha256":"782450c6fc1e432641df8b3cc605425c604109f1db5830873b0b778c59ae6946"}],"result_id":"MOR-006","statement_sha256":"969223c89123c8150a4a16dde30b45f11ccb2c9ba9a773533e000cbb936e1d0c"}
% HMT_RESULT_END:MOR-006:residence
% HMT_RESULT_BEGIN:MOR-007:residence
% {"material_control":"GATE_MATERIAL_68_22_22+AUDITORIA_ESTATICA","material_state":"INTEGRADO_O_REMITIDO_SIN_PERDIDA","residences":[{"path":"manuscrito/sections/integracion_20260812/bloque_64_operatorial.tex","sha256":"46e39285e0cee5f56142e4c31f8aece125754649847a8864878db03c39927c0f"}],"result_id":"MOR-007","statement_sha256":"cd14cec597ee88698599f74945cfe83dee11ed410ee4f5f9c55eea289b84c15d"}
% HMT_RESULT_END:MOR-007:residence
% HMT_RESULT_BEGIN:MOR-008:residence
% {"material_control":"GATE_MATERIAL_68_22_22+AUDITORIA_ESTATICA","material_state":"INTEGRADO_O_REMITIDO_SIN_PERDIDA","residences":[{"path":"manuscrito/sections/md/01a_medicion_cuantica_genealogica_rev2.tex","sha256":"725bbfb625e4bd74c3f5845dc2d3ad117b173fac63f11d300f7f6b8344fb505a"},{"path":"manuscrito/sections/md/04ab_postulados_cuanticos_genealogicos_rev2.tex","sha256":"1895f84bd75b104e873a3d73c5b0e8ae1b5df606bd6feefeb4f9666a92966804"}],"result_id":"MOR-008","statement_sha256":"a220ba7c55fda4a44314abe4540421d4c8d4b4e75f7ffa6f6bc273734217fc40"}
% HMT_RESULT_END:MOR-008:residence
% HMT_RESULT_BEGIN:MOR-009:residence
% {"material_control":"GATE_MATERIAL_68_22_22+AUDITORIA_ESTATICA","material_state":"INTEGRADO_O_REMITIDO_SIN_PERDIDA","residences":[{"path":"manuscrito/sections/md/04ab_postulados_cuanticos_genealogicos_rev2.tex","sha256":"1895f84bd75b104e873a3d73c5b0e8ae1b5df606bd6feefeb4f9666a92966804"},{"path":"manuscrito/sections/md/00b_realizacion_espectral_numero_particula.tex","sha256":"741b8542e90d89c9973f78712ccfb0625689f72a5af27ecedcc81d193f38d8f6"}],"result_id":"MOR-009","statement_sha256":"3b85ec28f733e941f275615780570ab9c2a3d81fc638a62c0e6728f5ae1f5925"}
% HMT_RESULT_END:MOR-009:residence
% HMT_RESULT_BEGIN:MOR-010:residence
% {"material_control":"GATE_MATERIAL_68_22_22+AUDITORIA_ESTATICA","material_state":"INTEGRADO_O_REMITIDO_SIN_PERDIDA","residences":[{"path":"manuscrito/sections/integracion_20260812/realizacion_68/02_realizacion_fisica_discreta_corregida.tex","sha256":"219a1a44dc7a0979b3728cebdc4308038cfbef74342bb0fcba50fe754c30292e"}],"result_id":"MOR-010","statement_sha256":"f2c89228291df78e200ef2085a4f4d478ac1be2e6684f9a715f658123de31342"}
% HMT_RESULT_END:MOR-010:residence
% HMT_RESULT_BEGIN:MOR-011:residence
% {"material_control":"GATE_MATERIAL_68_22_22+AUDITORIA_ESTATICA","material_state":"INTEGRADO_O_REMITIDO_SIN_PERDIDA","residences":[{"path":"manuscrito/sections/integracion_20260812/realizacion_68/05_involucion_constitutiva_bidireccional.tex","sha256":"c5cc64e9d5b8014a44b0c30ef8b2885e64b0379a107a073c6026b15381f76a7c"}],"result_id":"MOR-011","statement_sha256":"1ff2ade4f9f1721656ba26a95ce77ff15b5def7ccc20e9c135257ef0a7b97f0b"}
% HMT_RESULT_END:MOR-011:residence
% HMT_RESULT_BEGIN:MOR-012:residence
% {"material_control":"GATE_MATERIAL_68_22_22+AUDITORIA_ESTATICA","material_state":"INTEGRADO_O_REMITIDO_SIN_PERDIDA","residences":[{"path":"manuscrito/sections/integracion_20260812/realizacion_68/05_involucion_constitutiva_bidireccional.tex","sha256":"c5cc64e9d5b8014a44b0c30ef8b2885e64b0379a107a073c6026b15381f76a7c"}],"result_id":"MOR-012","statement_sha256":"c71dfcd2e0330c46692b8833f1bc5d41d7878bfc9db5a82f4cbfc84d7549ab7c"}
% HMT_RESULT_END:MOR-012:residence
% HMT_RESULT_BEGIN:MOR-013:residence
% {"material_control":"GATE_MATERIAL_68_22_22+AUDITORIA_ESTATICA","material_state":"INTEGRADO_O_REMITIDO_SIN_PERDIDA","residences":[{"path":"manuscrito/sections/md/04c_realizaciones_nula_material.tex","sha256":"d4bc1ef3e11a544359a35157497479312fce619dec624b4e01d33b6df16a84d9"}],"result_id":"MOR-013","statement_sha256":"ffe00c7b9b7d0edf7f35a609adfde95e7d3498e2638a94dfd78b707d810b9890"}
% HMT_RESULT_END:MOR-013:residence
% HMT_RESULT_BEGIN:MOR-014:residence
% {"material_control":"GATE_MATERIAL_68_22_22+AUDITORIA_ESTATICA","material_state":"INTEGRADO_O_REMITIDO_SIN_PERDIDA","residences":[{"path":"manuscrito/sections/md/04c_realizaciones_nula_material.tex","sha256":"d4bc1ef3e11a544359a35157497479312fce619dec624b4e01d33b6df16a84d9"},{"path":"manuscrito/sections/md/06ka_clausura_nulo_material_electron_rev3.tex","sha256":"4e83aee8d1ef22ee3bc6e71b10989c62a7a01c4f41153c30efc054909337c201"}],"result_id":"MOR-014","statement_sha256":"892c56b47c56b0f46d3c0955bfbb8bd5ef3f1c07286a797739e61c565e0dc2a2"}
% HMT_RESULT_END:MOR-014:residence
% HMT_RESULT_BEGIN:MOR-015:residence
% {"material_control":"GATE_MATERIAL_68_22_22+AUDITORIA_ESTATICA","material_state":"INTEGRADO_O_REMITIDO_SIN_PERDIDA","residences":[{"path":"manuscrito/sections/integracion_20260812/realizacion_68/03_hopf_villarceau_memoria_5_2.tex","sha256":"5ff5c6ec5855f74338db3440d9a7ddeaa0d93fec54b1ca5bb5094d6a14bd0fff"}],"result_id":"MOR-015","statement_sha256":"49efabbeafb0ec3ca4300a5684772f9b39fbc9836d768664adc1d443c2be762b"}
% HMT_RESULT_END:MOR-015:residence
% HMT_RESULT_BEGIN:MOR-016:residence
% {"material_control":"GATE_MATERIAL_68_22_22+AUDITORIA_ESTATICA","material_state":"INTEGRADO_O_REMITIDO_SIN_PERDIDA","residences":[{"path":"manuscrito/sections/md/04b_banda_particula_virtual.tex","sha256":"d9b8c3db3bc26d37b5a78251a5fbcf67076eaacf35c194b87bf6552d1edfac08"},{"path":"manuscrito/sections/md/21_persistencia_ruta_mobius_familias_actualizada.tex","sha256":"375c7e897d21777f05f9a506bb2429e5de76351b55f10e15f275913d9c306999"}],"result_id":"MOR-016","statement_sha256":"d549af55e36e03afb5ec1ca58bf41f7c6c0c5bd39e17a4a0320aff3b421fec3b"}
% HMT_RESULT_END:MOR-016:residence
% HMT_RESULT_BEGIN:MOR-017:residence
% {"material_control":"GATE_MATERIAL_68_22_22+AUDITORIA_ESTATICA","material_state":"INTEGRADO_O_REMITIDO_SIN_PERDIDA","residences":[{"path":"manuscrito/sections/integracion_20260812/realizacion_68/04_control_cglmp_exacto_68.tex","sha256":"278ced97b864e5c8a06650fca5a102c1a213716af14d3148e8206da54f91fdd7"}],"result_id":"MOR-017","statement_sha256":"8aceaf43834801dfe28d27e802dea64160a4e53dd788c2579d1ea1828dd1aba1"}
% HMT_RESULT_END:MOR-017:residence
% HMT_RESULT_BEGIN:MOR-018:residence
% {"material_control":"GATE_MATERIAL_68_22_22+AUDITORIA_ESTATICA","material_state":"INTEGRADO_O_REMITIDO_SIN_PERDIDA","residences":[{"path":"manuscrito/sections/md/10_constitutivas_si.tex","sha256":"852439fe31f4004a04c10217e89f3414c049bc8803c5b86e3d5b2097c04da0a1"}],"result_id":"MOR-018","statement_sha256":"1c4bd3e069ce1c87d0a969b52dbb80e067cb34adc73c877aa272e4018e841ed0"}
% HMT_RESULT_END:MOR-018:residence
% HMT_RESULT_BEGIN:MOR-019:residence
% {"material_control":"GATE_MATERIAL_68_22_22+AUDITORIA_ESTATICA","material_state":"INTEGRADO_O_REMITIDO_SIN_PERDIDA","residences":[{"path":"manuscrito/sections/md/03b_escala_accion_determinantal.tex","sha256":"73e012418f7e982e277e3e978a0fb03d76d0aebacc8fe085341c70121b20d832"},{"path":"manuscrito/sections/md/04e_termodinamica_rutas.tex","sha256":"ab1429f9145eee89e5287095b94f327291fa4aa247ae0abcc95e3f5611de87c1"}],"result_id":"MOR-019","statement_sha256":"12f6cc8a7f2131682e716fb00c3f2546dbb1f63249994d5d8c1c14d83510c994"}
% HMT_RESULT_END:MOR-019:residence
% HMT_RESULT_BEGIN:MOR-020:residence
% {"material_control":"GATE_MATERIAL_68_22_22+AUDITORIA_ESTATICA","material_state":"INTEGRADO_O_REMITIDO_SIN_PERDIDA","residences":[{"path":"manuscrito/sections/md/05aa_ley_general_masa_accion_autonoma.tex","sha256":"43dbef95faacb5db35585fa447b16406038f05b5296852f259687292bc0313a2"},{"path":"manuscrito/sections/md/05f_operadores_masa_two_way_rev11.tex","sha256":"e0c576ff54f405c41adea7093ca6b117134e9322749c969f1fc4cca6e8615d35"},{"path":"manuscrito/sections/md/06m_tabla_completa_masas_hmt_sm_revfinal.tex","sha256":"1dd8ca6908a9b94af9d5533d57785a2222d455f88e7ac70aa9acf34667ed35b1"}],"result_id":"MOR-020","statement_sha256":"ef9647c7358991a2b5046f1b253d9497443b3f0c31a95b7f6a9906100d45378e"}
% HMT_RESULT_END:MOR-020:residence
% HMT_RESULT_BEGIN:MOR-021:residence
% {"material_control":"GATE_MATERIAL_68_22_22+AUDITORIA_ESTATICA","material_state":"INTEGRADO_O_REMITIDO_SIN_PERDIDA","residences":[{"path":"manuscrito/sections/md/06k_electron_two_way_rev11.tex","sha256":"e9857301411fd54a3f9d9ad99f56466fc2d7aa72b59d7a4a12f14076fc151e0b"},{"path":"manuscrito/sections/md/06ka_clausura_nulo_material_electron_rev3.tex","sha256":"4e83aee8d1ef22ee3bc6e71b10989c62a7a01c4f41153c30efc054909337c201"}],"result_id":"MOR-021","statement_sha256":"79bbb961eec13af6b1fb10be55c0f3a17f871661f753bc37ab8c5e6ce19c0a6a"}
% HMT_RESULT_END:MOR-021:residence
% HMT_RESULT_BEGIN:MOR-022:residence
% {"material_control":"GATE_MATERIAL_68_22_22+AUDITORIA_ESTATICA","material_state":"INTEGRADO_O_REMITIDO_SIN_PERDIDA","residences":[{"path":"manuscrito/sections/md/03d_identidad_autodual_planck_ct108_rev3.tex","sha256":"522a1ce5c537fb26513227fc0766eadefd1b6fd898382137129677aaa69cfaed"}],"result_id":"MOR-022","statement_sha256":"e741dca3aef92e05581285a7ecd9802c495fc2b37445a4604229e67b619c98da"}
% HMT_RESULT_END:MOR-022:residence
% HMT_RESULT_BEGIN:MOR-023:residence
% {"material_control":"GATE_MATERIAL_68_22_22+AUDITORIA_ESTATICA","material_state":"INTEGRADO_O_REMITIDO_SIN_PERDIDA","residences":[{"path":"manuscrito/sections/integracion_20260812/realizacion_68/02_realizacion_fisica_discreta_corregida.tex","sha256":"219a1a44dc7a0979b3728cebdc4308038cfbef74342bb0fcba50fe754c30292e"}],"result_id":"MOR-023","statement_sha256":"bfaa0b7d864dd992d466327275ed93894d90946d897887db756d6d92268920b5"}
% HMT_RESULT_END:MOR-023:residence
% HMT_RESULT_BEGIN:MOR-024:residence
% {"material_control":"GATE_MATERIAL_68_22_22+AUDITORIA_ESTATICA","material_state":"INTEGRADO_O_REMITIDO_SIN_PERDIDA","residences":[{"path":"manuscrito/sections/md/11c_gravedad_holonomia_rev6.tex","sha256":"87d99cb98c92f5e7a771a9ab750690274e370d49a630e461b9d524e02cf6cf71"},{"path":"manuscrito/sections/md/11n_selector_gravitatorio_rev11.tex","sha256":"14b45e691dfd01916515b4a15b910eba1ff698031f794e68ba13d72da9973c9d"}],"result_id":"MOR-024","statement_sha256":"7736b062aa565d44a6b82fd9985d0d4dd7dd7bf8810791d818212081641834d4"}
% HMT_RESULT_END:MOR-024:residence
% HMT_RESULT_BEGIN:MOR-025:residence
% {"material_control":"GATE_MATERIAL_68_22_22+AUDITORIA_ESTATICA","material_state":"INTEGRADO_O_REMITIDO_SIN_PERDIDA","residences":[{"path":"manuscrito/sections/integracion_20260812/realizacion_68/06_programas_mach_cosmologia_68.tex","sha256":"5cca10348cd22a33e86ac0edaaffbd486e433e7e3d80bdf4c97a9368ab868e89"}],"result_id":"MOR-025","statement_sha256":"560fed3a547431a0a0bb145c63be86a48bd71bbf9d20c9ce42e8660b1aa64996"}
% HMT_RESULT_END:MOR-025:residence
% HMT_RESULT_BEGIN:MOR-026:residence
% {"material_control":"GATE_MATERIAL_68_22_22+AUDITORIA_ESTATICA","material_state":"INTEGRADO_O_REMITIDO_SIN_PERDIDA","residences":[{"path":"manuscrito/sections/hmt/14f_cierre_genealogico_numero.tex","sha256":"4522230d8297a1420ccb586d0d6c1fefcb08b5e8e2deeda85b11cc73afd8ded5"},{"path":"manuscrito/sections/hmt/19b_zeta_vacancias_primos_rev11.tex","sha256":"37b1f8fff4663d03826091250843b7a33c8a83e3fd08659c5fa6613a4c346534"}],"result_id":"MOR-026","statement_sha256":"e27ece362c8e14e5d61ce86444197f785ee77f3c6b2a79b2057a895ebb30db26"}
% HMT_RESULT_END:MOR-026:residence
% HMT_RESULT_BEGIN:MOR-027:residence
% {"material_control":"GATE_MATERIAL_68_22_22+AUDITORIA_ESTATICA","material_state":"INTEGRADO_O_REMITIDO_SIN_PERDIDA","residences":[{"path":"manuscrito/sections/hmt/14d_ontologia_numeros_rev11.tex","sha256":"c96b3c94af4668f15c6a061260d1c1900c52b69b8a4a5d2442268f6ff1d0abcc"},{"path":"manuscrito/sections/hmt/14f_cierre_genealogico_numero.tex","sha256":"4522230d8297a1420ccb586d0d6c1fefcb08b5e8e2deeda85b11cc73afd8ded5"}],"result_id":"MOR-027","statement_sha256":"5b813e0066c4fb64f0e94a22e81f2741e55dc4dec333c9aea2d7d02f31523a5d"}
% HMT_RESULT_END:MOR-027:residence
% HMT_RESULT_BEGIN:MOR-028:residence
% {"material_control":"GATE_MATERIAL_68_22_22+AUDITORIA_ESTATICA","material_state":"INTEGRADO_O_REMITIDO_SIN_PERDIDA","residences":[{"path":"manuscrito/sections/hmt/03f_puente_central_app_emrh_rev2.tex","sha256":"f4088242d198ece37ba686e3314432777582aa93e3b7960e39986a5e9a54358f"}],"result_id":"MOR-028","statement_sha256":"3efd7a28586cf53974d753f4d22e4e7b42a73517f5454c21aefb37550f6b1eae"}
% HMT_RESULT_END:MOR-028:residence
% HMT_RESULT_BEGIN:MOR-029:residence
% {"material_control":"GATE_MATERIAL_68_22_22+AUDITORIA_ESTATICA","material_state":"INTEGRADO_O_REMITIDO_SIN_PERDIDA","residences":[{"path":"manuscrito/sections/hmt/06e_certificado_generacion_monodromica_rev7.tex","sha256":"2a0cf8028e12605a97923c53726fcb62c1798c458af8e9f78f12c1129b0e1b94"}],"result_id":"MOR-029","statement_sha256":"91fbdaa9566d579a9f7f0fc50c56eae1187bb6ce719915dc9f11c1031024480f"}
% HMT_RESULT_END:MOR-029:residence
% HMT_RESULT_BEGIN:MOR-030:residence
% {"material_control":"GATE_MATERIAL_68_22_22+AUDITORIA_ESTATICA","material_state":"INTEGRADO_O_REMITIDO_SIN_PERDIDA","residences":[{"path":"manuscrito/sections/integracion_20260812/realizacion_68/08_interfaces_falsadores_realizacion_68.tex","sha256":"323f27caa4ce2d4e3e96795394eb2843c9549f0c7137bde4a8702bf79561e05e"}],"result_id":"MOR-030","statement_sha256":"ae3db0712fc04e6f7197110ed9ac53af8a4e6d1ce41056c0f1849cbe71519bc2"}
% HMT_RESULT_END:MOR-030:residence
% HMT_RESULT_BEGIN:PLS-001:residence
% {"material_control":"GATE_MATERIAL_68_22_22+AUDITORIA_ESTATICA","material_state":"INTEGRADO_O_REMITIDO_SIN_PERDIDA","residences":[{"path":"manuscrito/sections/md/03d_identidad_autodual_planck_ct108_rev3.tex","sha256":"522a1ce5c537fb26513227fc0766eadefd1b6fd898382137129677aaa69cfaed"}],"result_id":"PLS-001","statement_sha256":"dd9b9e630ef17a042a9d2eb03dfb5bfb4f12afef20d488d3eab6c9d3236bc69a"}
% HMT_RESULT_END:PLS-001:residence
% HMT_RESULT_BEGIN:MAD-001:residence
% {"material_control":"GATE_MATERIAL_68_22_22+AUDITORIA_ESTATICA","material_state":"INTEGRADO_O_REMITIDO_SIN_PERDIDA","residences":[{"path":"manuscrito/sections/md/03d_identidad_autodual_planck_ct108_rev3.tex","sha256":"522a1ce5c537fb26513227fc0766eadefd1b6fd898382137129677aaa69cfaed"}],"result_id":"MAD-001","statement_sha256":"5bce3c37e8568b7428605807933327d2bddbd9668451a4eac13bc22f055af271"}
% HMT_RESULT_END:MAD-001:residence
% HMT_RESULT_BEGIN:AGR-001:residence
% {"material_control":"GATE_MATERIAL_68_22_22+AUDITORIA_ESTATICA","material_state":"INTEGRADO_O_REMITIDO_SIN_PERDIDA","residences":[{"path":"manuscrito/sections/md/03d_identidad_autodual_planck_ct108_rev3.tex","sha256":"522a1ce5c537fb26513227fc0766eadefd1b6fd898382137129677aaa69cfaed"}],"result_id":"AGR-001","statement_sha256":"4fee4acf8c84a7dbbaf710eff510d2cb37e76574c593f4912526198dafaa1747"}
% HMT_RESULT_END:AGR-001:residence
% HMT_RESULT_BEGIN:PLE-001:residence
% {"material_control":"GATE_MATERIAL_68_22_22+AUDITORIA_ESTATICA","material_state":"INTEGRADO_O_REMITIDO_SIN_PERDIDA","residences":[{"path":"manuscrito/sections/md/03d_identidad_autodual_planck_ct108_rev3.tex","sha256":"522a1ce5c537fb26513227fc0766eadefd1b6fd898382137129677aaa69cfaed"},{"path":"manuscrito/sections/md/06ka_clausura_nulo_material_electron_rev3.tex","sha256":"4e83aee8d1ef22ee3bc6e71b10989c62a7a01c4f41153c30efc054909337c201"}],"result_id":"PLE-001","statement_sha256":"354993392866e30255c6503239114f91f02cba03ce48f1bd36427071e5939fc8"}
% HMT_RESULT_END:PLE-001:residence
% HMT_RESULT_BEGIN:ANG-001:residence
% {"material_control":"GATE_MATERIAL_68_22_22+AUDITORIA_ESTATICA","material_state":"INTEGRADO_O_REMITIDO_SIN_PERDIDA","residences":[{"path":"manuscrito/sections/integracion_20260812/realizacion_68/03_hopf_villarceau_memoria_5_2.tex","sha256":"5ff5c6ec5855f74338db3440d9a7ddeaa0d93fec54b1ca5bb5094d6a14bd0fff"}],"result_id":"ANG-001","statement_sha256":"f8f0e5ef3962f4ee1a30115f54f35751e6d96da9a4b54f83017c293e67683683"}
% HMT_RESULT_END:ANG-001:residence
% HMT_RESULT_BEGIN:BDR-001:residence
% {"material_control":"GATE_MATERIAL_68_22_22+AUDITORIA_ESTATICA","material_state":"INTEGRADO_O_REMITIDO_SIN_PERDIDA","residences":[{"path":"manuscrito/sections/integracion_20260812/realizacion_68/05_involucion_constitutiva_bidireccional.tex","sha256":"c5cc64e9d5b8014a44b0c30ef8b2885e64b0379a107a073c6026b15381f76a7c"}],"result_id":"BDR-001","statement_sha256":"6a0929cf24c9401330b5cc9ca4baa6fbfe43f7d9efa70a75883aa7f5168a4682"}
% HMT_RESULT_END:BDR-001:residence
% HMT_RESULT_BEGIN:OCC-001:residence
% {"material_control":"GATE_MATERIAL_68_22_22+AUDITORIA_ESTATICA","material_state":"INTEGRADO_O_REMITIDO_SIN_PERDIDA","residences":[{"path":"manuscrito/sections/md/04e_termodinamica_rutas.tex","sha256":"ab1429f9145eee89e5287095b94f327291fa4aa247ae0abcc95e3f5611de87c1"}],"result_id":"OCC-001","statement_sha256":"ee25ce148dd36c42aeaace51d97662de4acb35aa1ed8f59a27a40f937515c612"}
% HMT_RESULT_END:OCC-001:residence
% HMT_RESULT_BEGIN:BOL-001:residence
% {"material_control":"GATE_MATERIAL_68_22_22+AUDITORIA_ESTATICA","material_state":"INTEGRADO_O_REMITIDO_SIN_PERDIDA","residences":[{"path":"manuscrito/sections/md/04e_termodinamica_rutas.tex","sha256":"ab1429f9145eee89e5287095b94f327291fa4aa247ae0abcc95e3f5611de87c1"}],"result_id":"BOL-001","statement_sha256":"b336e9229e7eae82bc749151dce330d09f3ec616c75b361fac0728e7c87f6c18"}
% HMT_RESULT_END:BOL-001:residence
% HMT_RESULT_BEGIN:COS-001:residence
% {"material_control":"GATE_MATERIAL_68_22_22+AUDITORIA_ESTATICA","material_state":"INTEGRADO_O_REMITIDO_SIN_PERDIDA","residences":[{"path":"manuscrito/sections/integracion_20260812/realizacion_68/06_programas_mach_cosmologia_68.tex","sha256":"5cca10348cd22a33e86ac0edaaffbd486e433e7e3d80bdf4c97a9368ab868e89"}],"result_id":"COS-001","statement_sha256":"bd40c86eaf9fa174e90b2113b1e1001dc7db1a73fb69eb230c228131867c4969"}
% HMT_RESULT_END:COS-001:residence
% HMT_RESULT_BEGIN:GWV-001:residence
% {"material_control":"GATE_MATERIAL_68_22_22+AUDITORIA_ESTATICA","material_state":"INTEGRADO_O_REMITIDO_SIN_PERDIDA","residences":[{"path":"manuscrito/sections/integracion_20260812/realizacion_68/06_programas_mach_cosmologia_68.tex","sha256":"5cca10348cd22a33e86ac0edaaffbd486e433e7e3d80bdf4c97a9368ab868e89"}],"result_id":"GWV-001","statement_sha256":"3cf172b69a92d50f0b499502ca1a16ea845a1c3adb4ce9e5c07cfda16f9c38fa"}
% HMT_RESULT_END:GWV-001:residence
% HMT_RESULT_BEGIN:DRK-001:residence
% {"material_control":"GATE_MATERIAL_68_22_22+AUDITORIA_ESTATICA","material_state":"INTEGRADO_O_REMITIDO_SIN_PERDIDA","residences":[{"path":"manuscrito/sections/integracion_20260812/realizacion_68/06_programas_mach_cosmologia_68.tex","sha256":"5cca10348cd22a33e86ac0edaaffbd486e433e7e3d80bdf4c97a9368ab868e89"}],"result_id":"DRK-001","statement_sha256":"cfbb50fa8ff9e04eb1292280e0f1cac17fb57e837209733564e0026a2be31639"}
% HMT_RESULT_END:DRK-001:residence
% HMT_RESULT_BEGIN:CMK-001:residence
% {"material_control":"GATE_MATERIAL_68_22_22+AUDITORIA_ESTATICA","material_state":"INTEGRADO_O_REMITIDO_SIN_PERDIDA","residences":[{"path":"manuscrito/sections/integracion_20260812/realizacion_68/07_cmk_coestructura_dinamica.tex","sha256":"66c38ef3af1739a961696c1d7c23c719df6530be5afced7b53e2c3de4b79d1fd"}],"result_id":"CMK-001","statement_sha256":"d549072e45a1534a427abeaf657bff1597774b49ed6b4c93ac49b3a4a6c836a6"}
% HMT_RESULT_END:CMK-001:residence
% HMT_RESULT_BEGIN:MRC-001:residence
% {"material_control":"GATE_MATERIAL_68_22_22+AUDITORIA_ESTATICA","material_state":"INTEGRADO_O_REMITIDO_SIN_PERDIDA","residences":[{"path":"manuscrito/sections/integracion_20260812/realizacion_68/09_certificado_y_disposicion_68.tex","sha256":"87dc8390235dd0f1c606ca180c5b56639642fa4818d7400f4ce633c8df913d0c"}],"result_id":"MRC-001","statement_sha256":"80a7a36147cfc1163735bc538f698ff947ea95661565393b3489da4e24a85dd7"}
% HMT_RESULT_END:MRC-001:residence
% HMT_RESULT_BEGIN:PRG68-001:residence
% {"material_control":"GATE_MATERIAL_68_22_22+AUDITORIA_ESTATICA","material_state":"INTEGRADO_O_REMITIDO_SIN_PERDIDA","residences":[{"path":"manuscrito/sections/integracion_20260812/realizacion_68/08_interfaces_falsadores_realizacion_68.tex","sha256":"323f27caa4ce2d4e3e96795394eb2843c9549f0c7137bde4a8702bf79561e05e"}],"result_id":"PRG68-001","statement_sha256":"cf2a64c0deb4ed89400145c71e75b98276573e3fd427b0a93e0ecf04d4114f26"}
% HMT_RESULT_END:PRG68-001:residence
% HMT_RESULT_BEGIN:GTR-001:residence
% {"material_control":"GATE_MATERIAL_68_22_22+AUDITORIA_ESTATICA","material_state":"INTEGRADO_O_REMITIDO_SIN_PERDIDA","residences":[{"path":"manuscrito/sections/md/11c_gravedad_holonomia_rev6.tex","sha256":"87d99cb98c92f5e7a771a9ab750690274e370d49a630e461b9d524e02cf6cf71"},{"path":"manuscrito/sections/md/11n_selector_gravitatorio_rev11.tex","sha256":"14b45e691dfd01916515b4a15b910eba1ff698031f794e68ba13d72da9973c9d"}],"result_id":"GTR-001","statement_sha256":"ddfba4182cba6421426c961c29f8e8fc1053105316f55429bf5db6218eae8a25"}
% HMT_RESULT_END:GTR-001:residence
% HMT_RESULT_BEGIN:AMB-001:residence
% {"material_control":"PASS_AMBIVALENCE_AREAL_VOLUMETRIC+GATE_MATERIAL_V5_219","material_state":"CONSERVADO_Y_ATOMIZADO_EN_DELTA_SUCESOR","residences":[{"path":"manuscrito/sections/hmt/05b_extensiones_dinamica_cociente.tex","sha256":"2c562fa974ec9637a3fe09b0594b57cf0652f1f6813e6eefcae2a14929248882"}],"result_id":"AMB-001","statement_sha256":"a33b07c534544852bd3e4e64ece4fe0c7c8af500d5eb8b8addbb97736170835e"}
% HMT_RESULT_END:AMB-001:residence
 % Sidecar operativo sucesor: conserva los 219 resultados atómicos de REV.4 y
% añade los 29 resultados CHI, separado del estado material testigo.
% Sidecar atómico de la edición integral reforzada: 261 IDs estables.
% PREPARED: no es recibo y no autoriza por sí solo la compilación.
% HMT_RESULT_BEGIN:R3-001:residence
% HMT_ATOMIC_RESULT {"material_state":"CONSERVADO_POR_BASE_GRANULAR_REV3","proof_strength":"EXACTO_INTERNO","provenance":"RESULTADO_RECUPERADO","reinforcement_state":"DEPLOYED_WITHOUT_NEW_RESULT_ID","residence":"portada, resumen, introducción, síntesis","result_id":"R3-001","statement":"Título y subtítulo como tesis: cierre holográfico, estructura discreta del continuo y origen de constantes","status":"PRESENT_IN_SUCCESSOR_SOURCE"}
% HMT_RESULT_END:R3-001:residence
% HMT_RESULT_BEGIN:R3-002:residence
% HMT_ATOMIC_RESULT {"material_state":"CONSERVADO_POR_BASE_GRANULAR_REV3","proof_strength":"EXACTO_INTERNO","provenance":"RESULTADO_RECUPERADO","reinforcement_state":"DEPLOYED_WITHOUT_NEW_RESULT_ID","residence":"`hmt/00_construccion_app_gnomon_rev11`, `03_app_ortograma`, `03f`","result_id":"R3-002","statement":"APP completo: suma, producto, residuo, cociente, gnomón, fibras, acarreo y bloque central","status":"PRESENT_IN_SUCCESSOR_SOURCE"}
% HMT_RESULT_END:R3-002:residence
% HMT_RESULT_BEGIN:R3-003:residence
% HMT_ATOMIC_RESULT {"material_state":"CONSERVADO_POR_BASE_GRANULAR_REV3","proof_strength":"EXACTO_INTERNO","provenance":"RESULTADO_RECUPERADO","reinforcement_state":"DEPLOYED_WITHOUT_NEW_RESULT_ID","residence":"`hmt/02_trit_geometrias`, `03g`","result_id":"R3-003","statement":"TRIT orientado y temporal, con presente torsional y memoria","status":"PRESENT_IN_SUCCESSOR_SOURCE"}
% HMT_RESULT_END:R3-003:residence
% HMT_RESULT_BEGIN:R3-004:residence
% HMT_ATOMIC_RESULT {"material_state":"CONSERVADO_POR_BASE_GRANULAR_REV3","proof_strength":"EXACTO_INTERNO","provenance":"RESULTADO_RECUPERADO","reinforcement_state":"DEPLOYED_WITHOUT_NEW_RESULT_ID","residence":"/Users/ruben/Documents/HMT2/EDICION_AUTONOMA_HMT_MD_2026-08-09_REV_3_EN_TRABAJO/01_FUENTE_REVISADA/manuscrito/sections/hmt/03h_arbol_operatorio_tpk_rev11.tex; /Users/ruben/Documents/HMT2/EDICION_AUTONOMA_HMT_MD_2026-08-09_REV_3_EN_TRABAJO/01_FUENTE_REVISADA/manuscrito/sections/hmt/03i_inventario_operatorio_tpk_rev11.tex; /Users/ruben/Documents/HMT2/EDICION_AUTONOMA_HMT_MD_2026-08-09_REV_3_EN_TRABAJO/01_FUENTE_REVISADA/manuscrito/sections/hmt/03k_genealogia_operaciones_rev2.tex; /Users/ruben/Documents/HMT2/EDICION_AUTONOMA_HMT_MD_2026-08-09_REV_3_EN_TRABAJO/01_FUENTE_REVISADA/manuscrito/sections/hmt/03l_elevacion_cohomologica_dimension_rev2.tex; /Users/ruben/Documents/HMT2/EDICION_AUTONOMA_HMT_MD_2026-08-09_REV_3_EN_TRABAJO/01_FUENTE_REVISADA/manuscrito/sections/hmt/20_arquitectura_operatoria_tpk_actualizada.tex","result_id":"R3-004","statement":"TPK-full: objetos, rutas, memoria, frontera, prolongación, operadores y cálculo genealógico","status":"PRESENT_IN_SUCCESSOR_SOURCE"}
% HMT_RESULT_END:R3-004:residence
% HMT_RESULT_BEGIN:R3-005:residence
% HMT_ATOMIC_RESULT {"material_state":"CONSERVADO_POR_BASE_GRANULAR_REV3","proof_strength":"EXACTO_INTERNO","provenance":"RESULTADO_RECUPERADO","reinforcement_state":"DEPLOYED_WITHOUT_NEW_RESULT_ID","residence":"`hmt/00a`, `00b`, `14f`","result_id":"R3-005","statement":"Cuenta→medida, continuo como secciones compatibles y recta como evaluación","status":"PRESENT_IN_SUCCESSOR_SOURCE"}
% HMT_RESULT_END:R3-005:residence
% HMT_RESULT_BEGIN:R3-006:residence
% HMT_ATOMIC_RESULT {"material_state":"CONSERVADO_POR_BASE_GRANULAR_REV3","proof_strength":"EXACTO_INTERNO","provenance":"RESULTADO_RECUPERADO","reinforcement_state":"DEPLOYED_WITHOUT_NEW_RESULT_ID","residence":"`hmt/06aa`, `06b--06f`","result_id":"R3-006","statement":"Monodromía \\(9\\to1\\) con memoria, generación coinductiva de \\(\\pi,e,\\varphi\\), 396/729/43/387 y cinco regiones de \\(\\pi\\)","status":"PRESENT_IN_SUCCESSOR_SOURCE"}
% HMT_RESULT_END:R3-006:residence
% HMT_RESULT_BEGIN:R3-007:residence
% HMT_ATOMIC_RESULT {"material_state":"CONSERVADO_POR_BASE_GRANULAR_REV3","proof_strength":"EXACTO_INTERNO","provenance":"RESULTADO_RECUPERADO","reinforcement_state":"DEPLOYED_WITHOUT_NEW_RESULT_ID","residence":"capítulos 8, Planck, coordenadas conjugadas y Barbero","result_id":"R3-007","statement":"\\(\\alpha_{\\rm HMT}\\), Planck, Barbero y constantes en su orden causal","status":"PRESENT_IN_SUCCESSOR_SOURCE"}
% HMT_RESULT_END:R3-007:residence
% HMT_RESULT_BEGIN:R3-008:residence
% HMT_ATOMIC_RESULT {"material_state":"CONSERVADO_POR_BASE_GRANULAR_REV3","proof_strength":"EXACTO_INTERNO","provenance":"RESULTADO_RECUPERADO","reinforcement_state":"DEPLOYED_WITHOUT_NEW_RESULT_ID","residence":"`hmt/09`, `09g`, `09e`","result_id":"R3-008","statement":"Doble proyección desde el mismo estado holonómico integral","status":"PRESENT_IN_SUCCESSOR_SOURCE"}
% HMT_RESULT_END:R3-008:residence
% HMT_RESULT_BEGIN:R3-009:residence
% HMT_ATOMIC_RESULT {"material_state":"CONSERVADO_POR_BASE_GRANULAR_REV3","proof_strength":"EXACTO_INTERNO","provenance":"RESULTADO_RECUPERADO","reinforcement_state":"DEPLOYED_WITHOUT_NEW_RESULT_ID","residence":"`hmt/09c`, `15`, `15c`, `15d`","result_id":"R3-009","statement":"Corredor Paley--Witt--Golay--Leech--\\(V^\\natural\\)--Monstruo--Moonshine","status":"PRESENT_IN_SUCCESSOR_SOURCE"}
% HMT_RESULT_END:R3-009:residence
% HMT_RESULT_BEGIN:R3-010:residence
% HMT_ATOMIC_RESULT {"material_state":"CONSERVADO_POR_BASE_GRANULAR_REV3","proof_strength":"EXACTO_INTERNO","provenance":"RESULTADO_RECUPERADO","reinforcement_state":"DEPLOYED_WITHOUT_NEW_RESULT_ID","residence":"`md/00b_realizacion_espectral_numero_particula`","result_id":"R3-010","statement":"Número→partícula: sección, ruta, espectro, persistencia y frontera MD","status":"PRESENT_IN_SUCCESSOR_SOURCE"}
% HMT_RESULT_END:R3-010:residence
% HMT_RESULT_BEGIN:R3-011:residence
% HMT_ATOMIC_RESULT {"material_state":"CONSERVADO_POR_BASE_GRANULAR_REV3","proof_strength":"EXACTO_INTERNO","provenance":"RESULTADO_RECUPERADO","reinforcement_state":"DEPLOYED_WITHOUT_NEW_RESULT_ID","residence":"`md/04ab_postulados_cuanticos_genealogicos_rev2`","result_id":"R3-011","statement":"Postulados cuánticos HMT: Hilbert, rayos/densidades, observables, PVM, Born, Stone--Schrödinger y tensor","status":"PRESENT_IN_SUCCESSOR_SOURCE"}
% HMT_RESULT_END:R3-011:residence
% HMT_RESULT_BEGIN:R3-012:residence
% HMT_ATOMIC_RESULT {"material_state":"CONSERVADO_POR_BASE_GRANULAR_REV3","proof_strength":"EXACTO_INTERNO","provenance":"RESULTADO_RECUPERADO","reinforcement_state":"DEPLOYED_WITHOUT_NEW_RESULT_ID","residence":"`md/01a_medicion_cuantica_genealogica_rev2`","result_id":"R3-012","statement":"Medición: fibra, Lüders, Kraus, CPTP, Stinespring y PVM común","status":"PRESENT_IN_SUCCESSOR_SOURCE"}
% HMT_RESULT_END:R3-012:residence
% HMT_RESULT_BEGIN:R3-013:residence
% HMT_ATOMIC_RESULT {"material_state":"CONSERVADO_POR_BASE_GRANULAR_REV3","proof_strength":"EXACTO_INTERNO","provenance":"RESULTADO_RECUPERADO","reinforcement_state":"DEPLOYED_WITHOUT_NEW_RESULT_ID","residence":"`md/04aa`, capítulos cuánticos","result_id":"R3-013","statement":"Caminos, operador unitario \\(27\\times27\\), límite cilíndrico y reconstrucción cuántica","status":"PRESENT_IN_SUCCESSOR_SOURCE"}
% HMT_RESULT_END:R3-013:residence
% HMT_RESULT_BEGIN:R3-014:residence
% HMT_ATOMIC_RESULT {"material_state":"CONSERVADO_POR_BASE_GRANULAR_REV3","proof_strength":"EXACTO_INTERNO","provenance":"RESULTADO_RECUPERADO","reinforcement_state":"DEPLOYED_WITHOUT_NEW_RESULT_ID","residence":"`md/04d`, `04e`, `04g`, `04h`","result_id":"R3-014","statement":"Landauer finito/informacional y estadística cuántica tipada","status":"PRESENT_IN_SUCCESSOR_SOURCE"}
% HMT_RESULT_END:R3-014:residence
% HMT_RESULT_BEGIN:R3-015:residence
% HMT_ATOMIC_RESULT {"material_state":"CONSERVADO_POR_BASE_GRANULAR_REV3","proof_strength":"EXACTO_INTERNO","provenance":"RESULTADO_RECUPERADO","reinforcement_state":"DEPLOYED_WITHOUT_NEW_RESULT_ID","residence":"`md/03d`, `06k`, `06ka`","result_id":"R3-015","statement":"Planck--CT108--Compton, recta autodual y clausura electrónica de rango uno","status":"PRESENT_IN_SUCCESSOR_SOURCE"}
% HMT_RESULT_END:R3-015:residence
% HMT_RESULT_BEGIN:R3-016:residence
% HMT_ATOMIC_RESULT {"material_state":"CONSERVADO_POR_BASE_GRANULAR_REV3","proof_strength":"EXACTO_INTERNO","provenance":"RESULTADO_RECUPERADO","reinforcement_state":"DEPLOYED_WITHOUT_NEW_RESULT_ID","residence":"`md/05aa`, `05f`, `06m`, `06i`","result_id":"R3-016","statement":"Ley operatorial de masas: \\(10^{-34}\\), \\(K_D/K_\\Omega\\), 81/56/13/324, 23+23 y 17 condicionadas","status":"PRESENT_IN_SUCCESSOR_SOURCE"}
% HMT_RESULT_END:R3-016:residence
% HMT_RESULT_BEGIN:R3-017:residence
% HMT_ATOMIC_RESULT {"material_state":"CONSERVADO_POR_BASE_GRANULAR_REV3","proof_strength":"EXACTO_INTERNO","provenance":"RESULTADO_RECUPERADO","reinforcement_state":"DEPLOYED_WITHOUT_NEW_RESULT_ID","residence":"`md/10c`, `11`, `11c`, `11m`, `11p`","result_id":"R3-017","statement":"Cartan--Holst, gravitación, autoinercia y sus fronteras declaradas","status":"PRESENT_IN_SUCCESSOR_SOURCE"}
% HMT_RESULT_END:R3-017:residence
% HMT_RESULT_BEGIN:R3-018:residence
% HMT_ATOMIC_RESULT {"material_state":"CONSERVADO_POR_BASE_GRANULAR_REV3","proof_strength":"EXACTO_INTERNO","provenance":"RESULTADO_RECUPERADO","reinforcement_state":"DEPLOYED_WITHOUT_NEW_RESULT_ID","residence":"`hmt/17_sintesis`, `md/12_sintesis`","result_id":"R3-018","statement":"Síntesis causal APP→TRIT→TPK→MD, incluyendo defectos 4/2/6/54","status":"PRESENT_IN_SUCCESSOR_SOURCE"}
% HMT_RESULT_END:R3-018:residence
% HMT_RESULT_BEGIN:R3-019:residence
% HMT_ATOMIC_RESULT {"material_state":"CONSERVADO_POR_BASE_GRANULAR_REV3","proof_strength":"EXACTO_INTERNO","provenance":"RESULTADO_RECUPERADO","reinforcement_state":"DEPLOYED_WITHOUT_NEW_RESULT_ID","residence":"activos de REV.3","result_id":"R3-019","statement":"Portada, referencias, tablas y QA ya aptos","status":"PRESENT_IN_SUCCESSOR_SOURCE"}
% HMT_RESULT_END:R3-019:residence
% HMT_RESULT_BEGIN:R3-020:residence
% HMT_ATOMIC_RESULT {"material_state":"CONSERVADO_POR_BASE_GRANULAR_REV3","proof_strength":"EXACTO_INTERNO","provenance":"RESULTADO_RECUPERADO","reinforcement_state":"DEPLOYED_WITHOUT_NEW_RESULT_ID","residence":"grafo completo","result_id":"R3-020","statement":"Todo archivo y toda afirmación de REV.3 no afectados por el delta","status":"PRESENT_IN_SUCCESSOR_SOURCE"}
% HMT_RESULT_END:R3-020:residence
% HMT_RESULT_BEGIN:C-001:residence
% HMT_ATOMIC_RESULT {"material_state":"CONTROL_ACTIVO_EN_PLAN_Y_FUENTE","proof_strength":"EXACTO_INTERNO","provenance":"FORMALIZACION_NUEVA","reinforcement_state":"DEPLOYED_WITHOUT_NEW_RESULT_ID","residence":"gobierno transversal","result_id":"C-001","statement":"La monografía de 68 pp abrevia KMS como \\(F_{A,B}(t+i\\beta)=F_{B,A}(t)\\) bajo una definición que no lo permite","status":"CONTROLLED"}
% HMT_RESULT_END:C-001:residence
% HMT_RESULT_BEGIN:C-002:residence
% HMT_ATOMIC_RESULT {"material_state":"CONTROL_ACTIVO_EN_PLAN_Y_FUENTE","proof_strength":"EXACTO_INTERNO","provenance":"FORMALIZACION_NUEVA","reinforcement_state":"DEPLOYED_WITHOUT_NEW_RESULT_ID","residence":"gobierno transversal","result_id":"C-002","statement":"MOR-010 enumera siete piezas, mientras teorema y certificado tienen cuatro componentes","status":"CONTROLLED"}
% HMT_RESULT_END:C-002:residence
% HMT_RESULT_BEGIN:C-003:residence
% HMT_ATOMIC_RESULT {"material_state":"CONTROL_ACTIVO_EN_PLAN_Y_FUENTE","proof_strength":"EXACTO_INTERNO","provenance":"FORMALIZACION_NUEVA","reinforcement_state":"DEPLOYED_WITHOUT_NEW_RESULT_ID","residence":"gobierno transversal","result_id":"C-003","statement":"El verificador de 68 pp declara parte de la funtorialidad mediante `True`","status":"CONTROLLED"}
% HMT_RESULT_END:C-003:residence
% HMT_RESULT_BEGIN:C-004:residence
% HMT_ATOMIC_RESULT {"material_state":"CONTROL_ACTIVO_EN_PLAN_Y_FUENTE","proof_strength":"EXACTO_INTERNO","provenance":"FORMALIZACION_NUEVA","reinforcement_state":"DEPLOYED_WITHOUT_NEW_RESULT_ID","residence":"gobierno transversal","result_id":"C-004","statement":"QMS no conmutativos del puente son semiclásicos","status":"CONTROLLED"}
% HMT_RESULT_END:C-004:residence
% HMT_RESULT_BEGIN:C-005:residence
% HMT_ATOMIC_RESULT {"material_state":"CONTROL_ACTIVO_EN_PLAN_Y_FUENTE","proof_strength":"EXACTO_INTERNO","provenance":"FORMALIZACION_NUEVA","reinforcement_state":"DEPLOYED_WITHOUT_NEW_RESULT_ID","residence":"gobierno transversal","result_id":"C-005","statement":"El límite inverso HMT y el continuo MPS son objetos diferentes","status":"CONTROLLED"}
% HMT_RESULT_END:C-005:residence
% HMT_RESULT_BEGIN:C-006:residence
% HMT_ATOMIC_RESULT {"material_state":"CONTROL_ACTIVO_EN_PLAN_Y_FUENTE","proof_strength":"EXACTO_INTERNO","provenance":"FORMALIZACION_NUEVA","reinforcement_state":"DEPLOYED_WITHOUT_NEW_RESULT_ID","residence":"gobierno transversal","result_id":"C-006","statement":"`06_teorema_morfismo.tex` presupone un dominio común y hace preservar inversas a magnitudes positivas; además sugerir que la realización positiva se recupera de la orientada pierde variación total","status":"CONTROLLED"}
% HMT_RESULT_END:C-006:residence
% HMT_RESULT_BEGIN:C-007:residence
% HMT_ATOMIC_RESULT {"material_state":"CONTROL_ACTIVO_EN_PLAN_Y_FUENTE","proof_strength":"EXACTO_INTERNO","provenance":"FORMALIZACION_NUEVA","reinforcement_state":"DEPLOYED_WITHOUT_NEW_RESULT_ID","residence":"gobierno transversal","result_id":"C-007","statement":"El apéndice del puente matricial afirma que sus seis no-certificaciones se incluyen en el JSON, pero `scope_limits` contiene cinco controles distintos","status":"CONTROLLED"}
% HMT_RESULT_END:C-007:residence
% HMT_RESULT_BEGIN:MAS-RECTOR-001:residence
% HMT_ATOMIC_RESULT {"material_state":"CONSERVADO_CON_APTO_FOCAL","proof_strength":"EXACTO_INTERNO","provenance":"RESULTADO_RECUPERADO","reinforcement_state":"DEPLOYED_WITHOUT_NEW_RESULT_ID","residence":"/Users/ruben/Documents/HMT2/EDICION_AUTONOMA_HMT_MD_2026-08-09_REV_3_EN_TRABAJO/01_FUENTE_REVISADA/manuscrito/sections/md/06h_atlas_familias_rev6.tex; /Users/ruben/Documents/HMT2/EDICION_AUTONOMA_HMT_MD_2026-08-09_REV_3_EN_TRABAJO/01_FUENTE_REVISADA/manuscrito/sections/md/06i_inventario_humano_23_sectores_rev8.tex","result_id":"MAS-RECTOR-001","statement":"81 celdas; 56 familias; 13 multisecciones; 324 rutas orientadas; sector nulo separado","status":"PRESENT_IN_SUCCESSOR_SOURCE"}
% HMT_RESULT_END:MAS-RECTOR-001:residence
% HMT_RESULT_BEGIN:MAS-RECTOR-002:residence
% HMT_ATOMIC_RESULT {"material_state":"CONSERVADO_CON_APTO_FOCAL","proof_strength":"EXACTO_INTERNO","provenance":"ARQUITECTURA_AUTORAL_PREEXISTENTE","reinforcement_state":"DEPLOYED_WITHOUT_NEW_RESULT_ID","residence":"/Users/ruben/Documents/HMT2/EDICION_AUTONOMA_HMT_MD_2026-08-09_REV_3_EN_TRABAJO/01_FUENTE_REVISADA/manuscrito/sections/md/00b_realizacion_espectral_numero_particula.tex; /Users/ruben/Documents/HMT2/EDICION_AUTONOMA_HMT_MD_2026-08-09_REV_3_EN_TRABAJO/01_FUENTE_REVISADA/manuscrito/sections/md/05aa_ley_general_masa_accion_autonoma.tex","result_id":"MAS-RECTOR-002","statement":"extensión superviviente→ruta→ledger→firma Z5→carácter→torres→recta de masa","status":"PRESENT_IN_SUCCESSOR_SOURCE"}
% HMT_RESULT_END:MAS-RECTOR-002:residence
% HMT_RESULT_BEGIN:MAS-RECTOR-003:residence
% HMT_ATOMIC_RESULT {"material_state":"CONSERVADO_CON_APTO_FOCAL","proof_strength":"EXACTO_INTERNO","provenance":"RESULTADO_RECUPERADO","reinforcement_state":"DEPLOYED_WITHOUT_NEW_RESULT_ID","residence":"/Users/ruben/Documents/HMT2/EDICION_AUTONOMA_HMT_MD_2026-08-09_REV_3_EN_TRABAJO/01_FUENTE_REVISADA/manuscrito/sections/hmt/03_app_ortograma.tex","result_id":"MAS-RECTOR-003","statement":"Bc=9P_++5P_-; salto primordial=4; defecto global posterior 2→6→54","status":"PRESENT_IN_SUCCESSOR_SOURCE"}
% HMT_RESULT_END:MAS-RECTOR-003:residence
% HMT_RESULT_BEGIN:MAS-RECTOR-004:residence
% HMT_ATOMIC_RESULT {"material_state":"CONSERVADO_CON_APTO_FOCAL","proof_strength":"EXACTO_INTERNO","provenance":"RESULTADO_RECUPERADO","reinforcement_state":"DEPLOYED_WITHOUT_NEW_RESULT_ID","residence":"/Users/ruben/Documents/HMT2/EDICION_AUTONOMA_HMT_MD_2026-08-09_REV_3_EN_TRABAJO/01_FUENTE_REVISADA/manuscrito/sections/md/03b_escala_accion_determinantal.tex","result_id":"MAS-RECTOR-004","statement":"SNF(H4)=(1,2,2,4); |coker|=16; Sigma_H=phi alpha^16; ord10=-34","status":"PRESENT_IN_SUCCESSOR_SOURCE"}
% HMT_RESULT_END:MAS-RECTOR-004:residence
% HMT_RESULT_BEGIN:MAS-RECTOR-005:residence
% HMT_ATOMIC_RESULT {"material_state":"CONSERVADO_CON_APTO_FOCAL","proof_strength":"DEMOSTRADO_CON_ESTRUCTURA_DE_PARTIDA_EXPLICITA","provenance":"RESULTADO_RECUPERADO","reinforcement_state":"DEPLOYED_WITHOUT_NEW_RESULT_ID","residence":"/Users/ruben/Documents/HMT2/EDICION_AUTONOMA_HMT_MD_2026-08-09_REV_3_EN_TRABAJO/01_FUENTE_REVISADA/manuscrito/sections/md/03b_escala_accion_determinantal.tex; /Users/ruben/Documents/HMT2/EDICION_AUTONOMA_HMT_MD_2026-08-09_REV_3_EN_TRABAJO/01_FUENTE_REVISADA/manuscrito/sections/md/03c_elipse_accion_holonomia_rev8.tex","result_id":"MAS-RECTOR-005","statement":"delta_2w=H5-etaRet=(90/pi)Cpi(alpha); R_act=H5/etaRet","status":"PRESENT_IN_SUCCESSOR_SOURCE"}
% HMT_RESULT_END:MAS-RECTOR-005:residence
% HMT_RESULT_BEGIN:MAS-RECTOR-006:residence
% HMT_ATOMIC_RESULT {"material_state":"CONSERVADO_CON_APTO_FOCAL","proof_strength":"DEMOSTRADO_CON_ESTRUCTURA_DE_PARTIDA_EXPLICITA","provenance":"FORMALIZACION_NUEVA","reinforcement_state":"DEPLOYED_WITHOUT_NEW_RESULT_ID","residence":"/Users/ruben/Documents/HMT2/EDICION_AUTONOMA_HMT_MD_2026-08-09_REV_3_EN_TRABAJO/01_FUENTE_REVISADA/manuscrito/sections/md/03d_identidad_autodual_planck_ct108_rev3.tex; /Users/ruben/Documents/HMT2/EDICION_AUTONOMA_HMT_MD_2026-08-09_REV_3_EN_TRABAJO/01_FUENTE_REVISADA/manuscrito/sections/md/05aa_ley_general_masa_accion_autonoma.tex","result_id":"MAS-RECTOR-006","statement":"hbar_ret=etaRet 10^-34 U_S; omega0=2pi/(108t0); m_clk=hbar_ret omega0 c^-2","status":"PRESENT_IN_SUCCESSOR_SOURCE"}
% HMT_RESULT_END:MAS-RECTOR-006:residence
% HMT_RESULT_BEGIN:MAS-RECTOR-007:residence
% HMT_ATOMIC_RESULT {"material_state":"CONSERVADO_CON_APTO_FOCAL","proof_strength":"EXACTO_INTERNO","provenance":"RESULTADO_RECUPERADO","reinforcement_state":"DEPLOYED_WITHOUT_NEW_RESULT_ID","residence":"/Users/ruben/Documents/HMT2/EDICION_AUTONOMA_HMT_MD_2026-08-09_REV_3_EN_TRABAJO/01_FUENTE_REVISADA/manuscrito/sections/md/05f_operadores_masa_two_way_rev11.tex","result_id":"MAS-RECTOR-007","statement":"K_D=(D27+D36,D1,(D4-D3)/2); 14 perfiles","status":"PRESENT_IN_SUCCESSOR_SOURCE"}
% HMT_RESULT_END:MAS-RECTOR-007:residence
% HMT_RESULT_BEGIN:MAS-RECTOR-008:residence
% HMT_ATOMIC_RESULT {"material_state":"CONSERVADO_CON_APTO_FOCAL","proof_strength":"EXACTO_INTERNO","provenance":"RESULTADO_RECUPERADO","reinforcement_state":"DEPLOYED_WITHOUT_NEW_RESULT_ID","residence":"/Users/ruben/Documents/HMT2/EDICION_AUTONOMA_HMT_MD_2026-08-09_REV_3_EN_TRABAJO/01_FUENTE_REVISADA/manuscrito/sections/md/05f_operadores_masa_two_way_rev11.tex","result_id":"MAS-RECTOR-008","statement":"K_Omega=(Omega,54,P6); 9 perfiles","status":"PRESENT_IN_SUCCESSOR_SOURCE"}
% HMT_RESULT_END:MAS-RECTOR-008:residence
% HMT_RESULT_BEGIN:MAS-RECTOR-009:residence
% HMT_ATOMIC_RESULT {"material_state":"CONSERVADO_CON_APTO_FOCAL","proof_strength":"EXACTO_INTERNO","provenance":"CERTIFICADO_NUEVO","reinforcement_state":"DEPLOYED_WITHOUT_NEW_RESULT_ID","residence":"/Users/ruben/Documents/HMT2/EDICION_AUTONOMA_HMT_MD_2026-08-09_REV_3_EN_TRABAJO/01_FUENTE_REVISADA/manuscrito/sections/md/05f_operadores_masa_two_way_rev11.tex","result_id":"MAS-RECTOR-009","statement":"40 pares; 18 coincidencias; valor común único (80,54,6)","status":"PRESENT_IN_SUCCESSOR_SOURCE"}
% HMT_RESULT_END:MAS-RECTOR-009:residence
% HMT_RESULT_BEGIN:MAS-RECTOR-010:residence
% HMT_ATOMIC_RESULT {"material_state":"CONSERVADO_CON_APTO_FOCAL","proof_strength":"EXACTO_INTERNO","provenance":"RESULTADO_RECUPERADO","reinforcement_state":"DEPLOYED_WITHOUT_NEW_RESULT_ID","residence":"/Users/ruben/Documents/HMT2/EDICION_AUTONOMA_HMT_MD_2026-08-09_REV_3_EN_TRABAJO/01_FUENTE_REVISADA/manuscrito/sections/md/06f_biografia_electron_rev6.tex","result_id":"MAS-RECTOR-010","statement":"(5,5,++); tau12=+++---+++---; D=(54,56,68,48,32)","status":"PRESENT_IN_SUCCESSOR_SOURCE"}
% HMT_RESULT_END:MAS-RECTOR-010:residence
% HMT_RESULT_BEGIN:MAS-RECTOR-011:residence
% HMT_ATOMIC_RESULT {"material_state":"CONSERVADO_CON_APTO_FOCAL","proof_strength":"DEMOSTRADO_CON_ESTRUCTURA_DE_PARTIDA_EXPLICITA","provenance":"RESULTADO_RECUPERADO","reinforcement_state":"DEPLOYED_WITHOUT_NEW_RESULT_ID","residence":"/Users/ruben/Documents/HMT2/EDICION_AUTONOMA_HMT_MD_2026-08-09_REV_3_EN_TRABAJO/01_FUENTE_REVISADA/manuscrito/sections/md/06f_biografia_electron_rev6.tex; /Users/ruben/Documents/HMT2/EDICION_AUTONOMA_HMT_MD_2026-08-09_REV_3_EN_TRABAJO/01_FUENTE_REVISADA/manuscrito/sections/md/06ka_clausura_nulo_material_electron_rev3.tex","result_id":"MAS-RECTOR-011","statement":"e0=(sqrt3/4)[exp(phi/pi^2)-22alpha^3](1+15Delta4)=0.51099892248806607250987087719134943763; etapa basal, no salida final","status":"PRESENT_IN_SUCCESSOR_SOURCE"}
% HMT_RESULT_END:MAS-RECTOR-011:residence
% HMT_RESULT_BEGIN:MAS-RECTOR-012:residence
% HMT_ATOMIC_RESULT {"material_state":"CONSERVADO_CON_APTO_FOCAL","proof_strength":"EXACTO_INTERNO","provenance":"RESULTADO_RECUPERADO","reinforcement_state":"DEPLOYED_WITHOUT_NEW_RESULT_ID","residence":"/Users/ruben/Documents/HMT2/EDICION_AUTONOMA_HMT_MD_2026-08-09_REV_3_EN_TRABAJO/01_FUENTE_REVISADA/manuscrito/sections/md/06k_electron_two_way_rev11.tex; /Users/ruben/Documents/HMT2/EDICION_AUTONOMA_HMT_MD_2026-08-09_REV_3_EN_TRABAJO/01_FUENTE_REVISADA/manuscrito/sections/md/06ka_clausura_nulo_material_electron_rev3.tex","result_id":"MAS-RECTOR-012","statement":"kappa_e=80-54alpha+6Delta4; m_e^beta=e0 R_act^kappa_e=0.51099895069000080860825716483792339562463486771513350353936201542059113716425185","status":"PRESENT_IN_SUCCESSOR_SOURCE"}
% HMT_RESULT_END:MAS-RECTOR-012:residence
% HMT_RESULT_BEGIN:MAS-RECTOR-013:residence
% HMT_ATOMIC_RESULT {"material_state":"CONSERVADO_CON_APTO_FOCAL","proof_strength":"EXACTO_INTERNO","provenance":"RESULTADO_RECUPERADO","reinforcement_state":"DEPLOYED_WITHOUT_NEW_RESULT_ID","residence":"/Users/ruben/Documents/HMT2/EDICION_AUTONOMA_HMT_MD_2026-08-09_REV_3_EN_TRABAJO/01_FUENTE_REVISADA/manuscrito/sections/md/06f_biografia_electron_rev6.tex; /Users/ruben/Documents/HMT2/EDICION_AUTONOMA_HMT_MD_2026-08-09_REV_3_EN_TRABAJO/01_FUENTE_REVISADA/manuscrito/sections/md/06k_electron_two_way_rev11.tex","result_id":"MAS-RECTOR-013","statement":"firma cruda; firma par CT108; eventos; v_e=0 no se fusionan","status":"PRESENT_IN_SUCCESSOR_SOURCE"}
% HMT_RESULT_END:MAS-RECTOR-013:residence
% HMT_RESULT_BEGIN:MAS-RECTOR-014:residence
% HMT_ATOMIC_RESULT {"material_state":"CONSERVADO_CON_APTO_FOCAL","proof_strength":"EXACTO_INTERNO","provenance":"RESULTADO_RECUPERADO","reinforcement_state":"DEPLOYED_WITHOUT_NEW_RESULT_ID","residence":"/Users/ruben/Documents/HMT2/EDICION_AUTONOMA_HMT_MD_2026-08-09_REV_3_EN_TRABAJO/01_FUENTE_REVISADA/manuscrito/sections/md/05a_jerarquia_espectral_infinita_rev7.tex","result_id":"MAS-RECTOR-014","statement":"s_n=q_-^n-q_+^n; c_n=q_-^n+q_+^n; S=<90,120>={90,120}∪{30r:r≥6}","status":"PRESENT_IN_SUCCESSOR_SOURCE"}
% HMT_RESULT_END:MAS-RECTOR-014:residence
% HMT_RESULT_BEGIN:MAS-RECTOR-015:residence
% HMT_ATOMIC_RESULT {"material_state":"CONSERVADO_CON_APTO_FOCAL","proof_strength":"EXACTO_INTERNO","provenance":"FORMALIZACION_NUEVA","reinforcement_state":"DEPLOYED_WITHOUT_NEW_RESULT_ID","residence":"/Users/ruben/Documents/HMT2/EDICION_AUTONOMA_HMT_MD_2026-08-09_REV_3_EN_TRABAJO/01_FUENTE_REVISADA/manuscrito/sections/md/05a_jerarquia_espectral_infinita_rev7.tex","result_id":"MAS-RECTOR-015","statement":"k[S0]=k[X90,Y120]/(X90^4-Y120^3); dos genealogías a 360 antes del cociente","status":"PRESENT_IN_SUCCESSOR_SOURCE"}
% HMT_RESULT_END:MAS-RECTOR-015:residence
% HMT_RESULT_BEGIN:MAS-RECTOR-016:residence
% HMT_ATOMIC_RESULT {"material_state":"CONSERVADO_CON_APTO_FOCAL","proof_strength":"EXACTO_INTERNO","provenance":"RESULTADO_RECUPERADO","reinforcement_state":"DEPLOYED_WITHOUT_NEW_RESULT_ID","residence":"/Users/ruben/Documents/HMT2/EDICION_AUTONOMA_HMT_MD_2026-08-09_REV_3_EN_TRABAJO/01_FUENTE_REVISADA/manuscrito/sections/md/05f_operadores_masa_two_way_rev11.tex","result_id":"MAS-RECTOR-016","statement":"B_M^D y B_M^Omega diagonales exactos con espectros completos","status":"PRESENT_IN_SUCCESSOR_SOURCE"}
% HMT_RESULT_END:MAS-RECTOR-016:residence
% HMT_RESULT_BEGIN:MAS-RECTOR-017:residence
% HMT_ATOMIC_RESULT {"material_state":"CONSERVADO_CON_APTO_FOCAL","proof_strength":"EXACTO_INTERNO","provenance":"FORMALIZACION_NUEVA","reinforcement_state":"DEPLOYED_WITHOUT_NEW_RESULT_ID","residence":"/Users/ruben/Documents/HMT2/EDICION_AUTONOMA_HMT_MD_2026-08-09_REV_3_EN_TRABAJO/01_FUENTE_REVISADA/manuscrito/sections/md/05f_operadores_masa_two_way_rev11.tex","result_id":"MAS-RECTOR-017","statement":"M_M^{beta,r}=R_act^{B_M^r/2} M_M^(0) R_act^{B_M^r/2}","status":"PRESENT_IN_SUCCESSOR_SOURCE"}
% HMT_RESULT_END:MAS-RECTOR-017:residence
% HMT_RESULT_BEGIN:MAS-RECTOR-018:residence
% HMT_ATOMIC_RESULT {"material_state":"CONSERVADO_CON_APTO_FOCAL","proof_strength":"EXACTO_INTERNO","provenance":"CERTIFICADO_NUEVO","reinforcement_state":"DEPLOYED_WITHOUT_NEW_RESULT_ID","residence":"/Users/ruben/Documents/HMT2/EDICION_AUTONOMA_HMT_MD_2026-08-09_REV_3_EN_TRABAJO/01_FUENTE_REVISADA/manuscrito/sections/md/06m_tabla_completa_masas_hmt_sm_revfinal.tex","result_id":"MAS-RECTOR-018","statement":"Toda clase conserva residencia, operador/salida y estatuto; electrón beta y nulos cerrados; la publicación mantiene 23 filas de salida interna y 23 filas de comparabilidad física, una por clase","status":"PRESENT_IN_SUCCESSOR_SOURCE"}
% HMT_RESULT_END:MAS-RECTOR-018:residence
% HMT_RESULT_BEGIN:MAS-RECTOR-019:residence
% HMT_ATOMIC_RESULT {"material_state":"CONSERVADO_CON_APTO_FOCAL","proof_strength":"EXACTO_INTERNO","provenance":"RESULTADO_RECUPERADO","reinforcement_state":"DEPLOYED_WITHOUT_NEW_RESULT_ID","residence":"/Users/ruben/Documents/HMT2/EDICION_AUTONOMA_HMT_MD_2026-08-09_REV_3_EN_TRABAJO/01_FUENTE_REVISADA/manuscrito/sections/md/06m_tabla_completa_masas_hmt_sm_revfinal.tex","result_id":"MAS-RECTOR-019","statement":"m=0 en carta nula; kappa no aplicable","status":"PRESENT_IN_SUCCESSOR_SOURCE"}
% HMT_RESULT_END:MAS-RECTOR-019:residence
% HMT_RESULT_BEGIN:MAS-RECTOR-020:residence
% HMT_ATOMIC_RESULT {"material_state":"CONSERVADO_CON_APTO_FOCAL","proof_strength":"INTERFAZ_FISICA","provenance":"RESULTADO_RECUPERADO","reinforcement_state":"DEPLOYED_WITHOUT_NEW_RESULT_ID","residence":"/Users/ruben/Documents/HMT2/EDICION_AUTONOMA_HMT_MD_2026-08-09_REV_3_EN_TRABAJO/01_FUENTE_REVISADA/manuscrito/sections/md/06i_inventario_humano_23_sectores_rev8.tex; /Users/ruben/Documents/HMT2/EDICION_AUTONOMA_HMT_MD_2026-08-09_REV_3_EN_TRABAJO/01_FUENTE_REVISADA/manuscrito/sections/md/06m_tabla_completa_masas_hmt_sm_revfinal.tex","result_id":"MAS-RECTOR-020","statement":"Rama up/down, generación, color y carácter fino; scalarización exige join individual y esquema/escala","status":"PRESENT_IN_SUCCESSOR_SOURCE"}
% HMT_RESULT_END:MAS-RECTOR-020:residence
% HMT_RESULT_BEGIN:MAS-RECTOR-021:residence
% HMT_ATOMIC_RESULT {"material_state":"CONSERVADO_CON_APTO_FOCAL","proof_strength":"INTERFAZ_FISICA","provenance":"RESULTADO_RECUPERADO","reinforcement_state":"DEPLOYED_WITHOUT_NEW_RESULT_ID","residence":"/Users/ruben/Documents/HMT2/EDICION_AUTONOMA_HMT_MD_2026-08-09_REV_3_EN_TRABAJO/01_FUENTE_REVISADA/manuscrito/sections/md/06i_inventario_humano_23_sectores_rev8.tex; /Users/ruben/Documents/HMT2/EDICION_AUTONOMA_HMT_MD_2026-08-09_REV_3_EN_TRABAJO/01_FUENTE_REVISADA/manuscrito/sections/md/08b_transporte_ckm_pmns_rev11.tex; /Users/ruben/Documents/HMT2/EDICION_AUTONOMA_HMT_MD_2026-08-09_REV_3_EN_TRABAJO/01_FUENTE_REVISADA/manuscrito/sections/md/05a1_generaciones_quarks_neutrinos_rev2.tex","result_id":"MAS-RECTOR-021","statement":"Tres clases de sabor neutrínico sobre el par operatorial común B_F^D/B_F^Omega; los tres sabores comparten H_nu=directsum_{j=1}^4 H_{nu,G_j} de dimensión 20; f_e,f_mu,f_tau es base de interacción y G_j profundidad estructural; PMNS transporta bases y no crea autovalores, escala ni masa de sabor","status":"PRESENT_IN_SUCCESSOR_SOURCE"}
% HMT_RESULT_END:MAS-RECTOR-021:residence
% HMT_RESULT_BEGIN:MAS-RECTOR-022:residence
% HMT_ATOMIC_RESULT {"material_state":"CONSERVADO_CON_APTO_FOCAL","proof_strength":"RECONOCIMIENTO_EXTERNO","provenance":"RESULTADO_RECUPERADO","reinforcement_state":"DEPLOYED_WITHOUT_NEW_RESULT_ID","residence":"/Users/ruben/Documents/HMT2/EDICION_AUTONOMA_HMT_MD_2026-08-09_REV_3_EN_TRABAJO/01_FUENTE_REVISADA/manuscrito/sections/md/08_ckm_cp.tex; /Users/ruben/Documents/HMT2/EDICION_AUTONOMA_HMT_MD_2026-08-09_REV_3_EN_TRABAJO/01_FUENTE_REVISADA/manuscrito/sections/md/08b_transporte_ckm_pmns_rev11.tex","result_id":"MAS-RECTOR-022","statement":"Cuatro fórmulas angulares, matriz unitaria y J; contraste posterior","status":"PRESENT_IN_SUCCESSOR_SOURCE"}
% HMT_RESULT_END:MAS-RECTOR-022:residence
% HMT_RESULT_BEGIN:MAS-RECTOR-023:residence
% HMT_ATOMIC_RESULT {"material_state":"CONSERVADO_CON_APTO_FOCAL","proof_strength":"PROGRAMA_ABIERTO","provenance":"RESULTADO_RECUPERADO","reinforcement_state":"DEPLOYED_WITHOUT_NEW_RESULT_ID","residence":"/Users/ruben/Documents/HMT2/EDICION_AUTONOMA_HMT_MD_2026-08-09_REV_3_EN_TRABAJO/01_FUENTE_REVISADA/manuscrito/sections/md/06g_composicion_especies_rev6.tex","result_id":"MAS-RECTOR-023","statement":"Rutas compuestas con frontera/enlace; 432 y 1728 expresiones no son especies","status":"PRESENT_IN_SUCCESSOR_SOURCE"}
% HMT_RESULT_END:MAS-RECTOR-023:residence
% HMT_RESULT_BEGIN:MAS-RECTOR-024:residence
% HMT_ATOMIC_RESULT {"material_state":"CONSERVADO_CON_APTO_FOCAL","proof_strength":"EXACTO_INTERNO","provenance":"RESULTADO_RECUPERADO","reinforcement_state":"DEPLOYED_WITHOUT_NEW_RESULT_ID","residence":"/Users/ruben/Documents/HMT2/EDICION_AUTONOMA_HMT_MD_2026-08-09_REV_3_EN_TRABAJO/01_FUENTE_REVISADA/manuscrito/sections/md/06c_sectores_topologicos_no_omision.tex","result_id":"MAS-RECTOR-024","statement":"Anillos, dimensiones y trenzado son codominios propios, no masas escalares automáticas","status":"PRESENT_IN_SUCCESSOR_SOURCE"}
% HMT_RESULT_END:MAS-RECTOR-024:residence
% HMT_RESULT_BEGIN:MAS-RECTOR-025:residence
% HMT_ATOMIC_RESULT {"material_state":"CONSERVADO_CON_APTO_FOCAL","proof_strength":"EXACTO_INTERNO","provenance":"RESULTADO_RECUPERADO","reinforcement_state":"DEPLOYED_WITHOUT_NEW_RESULT_ID","residence":"/Users/ruben/Documents/HMT2/EDICION_AUTONOMA_HMT_MD_2026-08-09_REV_3_EN_TRABAJO/01_FUENTE_REVISADA/manuscrito/sections/md/04b_banda_particula_virtual.tex","result_id":"MAS-RECTOR-025","statement":"Sección de prolongación finita con ledger hasta L; no especie estable ni masa imaginaria","status":"PRESENT_IN_SUCCESSOR_SOURCE"}
% HMT_RESULT_END:MAS-RECTOR-025:residence
% HMT_RESULT_BEGIN:MAS-RECTOR-026:residence
% HMT_ATOMIC_RESULT {"material_state":"CONSERVADO_CON_APTO_FOCAL","proof_strength":"EXACTO_INTERNO","provenance":"CERTIFICADO_NUEVO","reinforcement_state":"DEPLOYED_WITHOUT_NEW_RESULT_ID","residence":"/Users/ruben/Documents/HMT2/EDICION_AUTONOMA_HMT_MD_2026-08-09_REV_3_EN_TRABAJO/01_FUENTE_REVISADA/manuscrito/sections/md/06m_tabla_completa_masas_hmt_sm_revfinal.tex","result_id":"MAS-RECTOR-026","statement":"Salida interna congelada antes de abrir el catálogo externo","status":"PRESENT_IN_SUCCESSOR_SOURCE"}
% HMT_RESULT_END:MAS-RECTOR-026:residence
% HMT_RESULT_BEGIN:MAS-RECTOR-027:residence
% HMT_ATOMIC_RESULT {"material_state":"CONSERVADO_CON_APTO_FOCAL","proof_strength":"EXACTO_INTERNO","provenance":"RESULTADO_RECUPERADO","reinforcement_state":"DEPLOYED_WITHOUT_NEW_RESULT_ID","residence":"/Users/ruben/Documents/HMT2/EDICION_AUTONOMA_HMT_MD_2026-08-09_REV_3_EN_TRABAJO/01_FUENTE_REVISADA/manuscrito/sections/md/06m_tabla_completa_masas_hmt_sm_revfinal.tex","result_id":"MAS-RECTOR-027","statement":"Sólo trazabilidad histórica; no residencia narrativa pública ni dominio","status":"PRESENT_IN_SUCCESSOR_SOURCE"}
% HMT_RESULT_END:MAS-RECTOR-027:residence
% HMT_RESULT_BEGIN:MAS-RECTOR-028:residence
% HMT_ATOMIC_RESULT {"material_state":"CONSERVADO_CON_APTO_FOCAL","proof_strength":"DEMOSTRADO_CON_ESTRUCTURA_DE_PARTIDA_EXPLICITA","provenance":"RESULTADO_RECUPERADO","reinforcement_state":"DEPLOYED_WITHOUT_NEW_RESULT_ID","residence":"/Users/ruben/Documents/HMT2/EDICION_AUTONOMA_HMT_MD_2026-08-09_REV_3_EN_TRABAJO/01_FUENTE_REVISADA/manuscrito/sections/md/06l_vacancias_corriente_rev11.tex","result_id":"MAS-RECTOR-028","statement":"Matriz [[21,22],[6,7]] con det=15 acompaña (-22,+15); uso local, no global","status":"PRESENT_IN_SUCCESSOR_SOURCE"}
% HMT_RESULT_END:MAS-RECTOR-028:residence
% HMT_RESULT_BEGIN:MAS-PUBLIC-023:residence
% HMT_ATOMIC_RESULT {"material_state":"CONSERVADO_CON_APTO_FOCAL","proof_strength":"EXACTO_INTERNO","provenance":"CERTIFICADO_NUEVO","reinforcement_state":"DEPLOYED_WITHOUT_NEW_RESULT_ID","residence":"/Users/ruben/Documents/HMT2/EDICION_AUTONOMA_HMT_MD_2026-08-09_REV_3_EN_TRABAJO/01_FUENTE_REVISADA/manuscrito/sections/md/06m_tabla_completa_masas_hmt_sm_revfinal.tex","result_id":"MAS-PUBLIC-023","statement":"las 23 claves internas reaparecen una vez y en el mismo orden en la capa pública de comparabilidad; las seis clases de codominio no escalar se comparan en su codominio nativo y no reciben una masa universal inventada","status":"PRESENT_IN_SUCCESSOR_SOURCE"}
% HMT_RESULT_END:MAS-PUBLIC-023:residence
% HMT_RESULT_BEGIN:MAS-COND-017:residence
% HMT_ATOMIC_RESULT {"material_state":"CONSERVADO_CON_APTO_FOCAL","proof_strength":"INTERFAZ_FISICA","provenance":"RESULTADO_RECUPERADO","reinforcement_state":"DEPLOYED_WITHOUT_NEW_RESULT_ID","residence":"/Users/ruben/Documents/HMT2/EDICION_AUTONOMA_HMT_MD_2026-08-09_REV_3_EN_TRABAJO/01_FUENTE_REVISADA/manuscrito/sections/md/06m_tabla_completa_masas_hmt_sm_revfinal.tex","result_id":"MAS-COND-017","statement":"17 evaluaciones exactas, 11 elementales y 6 compuestas, residen en una tabla separada y rotulada como no prospectiva; no son salidas cerradas, no seleccionan ruta y sólo admiten residual cuando los escalares, unidad, esquema y escala son comparables","status":"PRESENT_IN_SUCCESSOR_SOURCE"}
% HMT_RESULT_END:MAS-COND-017:residence
% HMT_RESULT_BEGIN:INF-001:residence
% HMT_ATOMIC_RESULT {"material_state":"INTEGRADO_O_REMITIDO_SIN_PERDIDA","proof_strength":"EXACTO_INTERNO","provenance":"RESULTADO_RECUPERADO","reinforcement_state":"DEPLOYED_WITHOUT_NEW_RESULT_ID","residence":"apertura del capítulo del continuo, tras `hmt/00b`","result_id":"INF-001","statement":"El continuo HMT no se forma con valores desnudos: cada casilla transporta bloque, firma, acarreo, fase, etiqueta y orientación. \\(\\xi_m=(B_m,\\sigma_m,c_m,\\omega_m)\\).","status":"PRESENT_IN_SUCCESSOR_SOURCE"}
% HMT_RESULT_END:INF-001:residence
% HMT_RESULT_BEGIN:INF-002:residence
% HMT_ATOMIC_RESULT {"material_state":"INTEGRADO_O_REMITIDO_SIN_PERDIDA","proof_strength":"EXACTO_INTERNO","provenance":"RESULTADO_RECUPERADO","reinforcement_state":"DEPLOYED_WITHOUT_NEW_RESULT_ID","residence":"bloque fundacional, § Estado, sección y valor, tras INF-001","result_id":"INF-002","statement":"Un número HMT es una sección compatible, no su expansión ni un conjunto acabado sin genealogía. \\(\\Omega_\\infty=\\varprojlim_m\\mathcal S_m\\).","status":"PRESENT_IN_SUCCESSOR_SOURCE"}
% HMT_RESULT_END:INF-002:residence
% HMT_RESULT_BEGIN:INF-003:residence
% HMT_ATOMIC_RESULT {"material_state":"INTEGRADO_O_REMITIDO_SIN_PERDIDA","proof_strength":"DEMOSTRADO_CON_ESTRUCTURA_DE_PARTIDA_EXPLICITA","provenance":"RESULTADO_RECUPERADO","reinforcement_state":"DEPLOYED_WITHOUT_NEW_RESULT_ID","residence":"bloque fundacional, § Existencia y unicidad coinductivas, tras INF-002","result_id":"INF-003","statement":"Finitud, no vaciedad y sobreyectividad dan existencia por König; la unicidad exige además fibras unitarias compatibles a toda profundidad.","status":"PRESENT_IN_SUCCESSOR_SOURCE"}
% HMT_RESULT_END:INF-003:residence
% HMT_RESULT_BEGIN:N9-001:residence
% HMT_ATOMIC_RESULT {"material_state":"INTEGRADO_O_REMITIDO_SIN_PERDIDA","proof_strength":"DEMOSTRADO_CON_ESTRUCTURA_DE_PARTIDA_EXPLICITA","provenance":"ARQUITECTURA_AUTORAL_PREEXISTENTE","reinforcement_state":"DEPLOYED_WITHOUT_NEW_RESULT_ID","residence":"`hmt/06aa` y transición inicial","result_id":"N9-001","statement":"La cadena APP--TRIT--TPK abre la frontera coinductiva; el retorno de fase no es retorno de estado.","status":"PRESENT_IN_SUCCESSOR_SOURCE"}
% HMT_RESULT_END:N9-001:residence
% HMT_RESULT_BEGIN:N9-002:residence
% HMT_ATOMIC_RESULT {"material_state":"INTEGRADO_O_REMITIDO_SIN_PERDIDA","proof_strength":"EXACTO_INTERNO","provenance":"RESULTADO_RECUPERADO","reinforcement_state":"DEPLOYED_WITHOUT_NEW_RESULT_ID","residence":"`hmt/06aa`","result_id":"N9-002","statement":"Las tres evaluaciones son lecturas coordinadas de un mismo estado; \\(\\varphi\\) fija el eje y \\(\\pi/e\\) constituyen la fibra especular orientada.","status":"PRESENT_IN_SUCCESSOR_SOURCE"}
% HMT_RESULT_END:N9-002:residence
% HMT_RESULT_BEGIN:N9-003:residence
% HMT_ATOMIC_RESULT {"material_state":"INTEGRADO_O_REMITIDO_SIN_PERDIDA","proof_strength":"EXACTO_INTERNO","provenance":"RESULTADO_RECUPERADO","reinforcement_state":"DEPLOYED_WITHOUT_NEW_RESULT_ID","residence":"`hmt/06b--06e`","result_id":"N9-003","statement":"Una salida visible conserva cinco genealogías regionales. \\(5=3_{++}+1_{--}+1_\\perp\\), \\(432+4\\cdot144=1008\\).","status":"PRESENT_IN_SUCCESSOR_SOURCE"}
% HMT_RESULT_END:N9-003:residence
% HMT_RESULT_BEGIN:N9-004:residence
% HMT_ATOMIC_RESULT {"material_state":"INTEGRADO_O_REMITIDO_SIN_PERDIDA","proof_strength":"EXACTO_INTERNO","provenance":"CERTIFICADO_NUEVO","reinforcement_state":"DEPLOYED_WITHOUT_NEW_RESULT_ID","residence":"`hmt/06e--06f`","result_id":"N9-004","statement":"Las 1.000 cifras son una restricción reproducible, no el techo ni por sí solas la prueba lógica del límite.","status":"PRESENT_IN_SUCCESSOR_SOURCE"}
% HMT_RESULT_END:N9-004:residence
% HMT_RESULT_BEGIN:HEL-001:residence
% HMT_ATOMIC_RESULT {"material_state":"INTEGRADO_O_REMITIDO_SIN_PERDIDA","proof_strength":"EXACTO_INTERNO","provenance":"RESULTADO_RECUPERADO","reinforcement_state":"DEPLOYED_WITHOUT_NEW_RESULT_ID","residence":"`hmt/03l`","result_id":"HEL-001","statement":"«Diez vuelve a uno» significa misma fase y memoria incrementada.","status":"PRESENT_IN_SUCCESSOR_SOURCE"}
% HMT_RESULT_END:HEL-001:residence
% HMT_RESULT_BEGIN:OPS-001:residence
% HMT_ATOMIC_RESULT {"material_state":"INTEGRADO_O_REMITIDO_SIN_PERDIDA","proof_strength":"EXACTO_INTERNO","provenance":"RESULTADO_RECUPERADO","reinforcement_state":"DEPLOYED_WITHOUT_NEW_RESULT_ID","residence":"`hmt/03k`","result_id":"OPS-001","statement":"Suma, producto, potencia, raíz, logaritmo, derivación e integración deben elevarse con memoria de fibra.","status":"PRESENT_IN_SUCCESSOR_SOURCE"}
% HMT_RESULT_END:OPS-001:residence
% HMT_RESULT_BEGIN:ZF-001:residence
% HMT_ATOMIC_RESULT {"material_state":"INTEGRADO_O_REMITIDO_SIN_PERDIDA","proof_strength":"EXACTO_INTERNO","provenance":"RESULTADO_RECUPERADO","reinforcement_state":"DEPLOYED_WITHOUT_NEW_RESULT_ID","residence":"nuevo bloque «Genealogía, extensionalidad y teoría de conjuntos» después de completar APP--TRIT--TPK y `03l`","result_id":"ZF-001","statement":"HMT no refuta ZFC: construye estructuras genealógicas codificables en ZFC y muestra qué información pierde la decategorización conjuntista.","status":"PRESENT_IN_SUCCESSOR_SOURCE"}
% HMT_RESULT_END:ZF-001:residence
% HMT_RESULT_BEGIN:ZF-002:residence
% HMT_ATOMIC_RESULT {"material_state":"INTEGRADO_O_REMITIDO_SIN_PERDIDA","proof_strength":"PROGRAMA_ABIERTO","provenance":"RESULTADO_NUEVO","reinforcement_state":"DEPLOYED_WITHOUT_NEW_RESULT_ID","residence":"bloque lógico, § Alcance axiomático, tras ZF-001","result_id":"ZF-002","statement":"No se afirmará que HMT sea ya un modelo interno completo de ZF, IZF o NBG: Separación, Reemplazo, Potencia plena e Infinito transfinito no se han verificado conjuntamente.","status":"PRESENT_IN_SUCCESSOR_SOURCE"}
% HMT_RESULT_END:ZF-002:residence
% HMT_RESULT_BEGIN:EXT-001:residence
% HMT_ATOMIC_RESULT {"material_state":"INTEGRADO_O_REMITIDO_SIN_PERDIDA","proof_strength":"EXACTO_INTERNO","provenance":"RESULTADO_RECUPERADO","reinforcement_state":"DEPLOYED_WITHOUT_NEW_RESULT_ID","residence":"bloque fundacional, § Extensionalidad tipada, tras DSC-001","result_id":"EXT-001","statement":"Dos secciones son idénticas si coinciden a toda profundidad; igualdad de valor es un cociente salvo lector fiel.","status":"PRESENT_IN_SUCCESSOR_SOURCE"}
% HMT_RESULT_END:EXT-001:residence
% HMT_RESULT_BEGIN:FND-001:residence
% HMT_ATOMIC_RESULT {"material_state":"INTEGRADO_O_REMITIDO_SIN_PERDIDA","proof_strength":"EXACTO_INTERNO","provenance":"RESULTADO_RECUPERADO","reinforcement_state":"DEPLOYED_WITHOUT_NEW_RESULT_ID","residence":"bloque lógico, § Fundación y coinducción, tras TOP-001","result_id":"FND-001","statement":"Los prefijos son bien fundados por rango natural; una sección infinita es coinductiva y no introduce ciclos de pertenencia.","status":"PRESENT_IN_SUCCESSOR_SOURCE"}
% HMT_RESULT_END:FND-001:residence
% HMT_RESULT_BEGIN:ORD-001:residence
% HMT_ATOMIC_RESULT {"material_state":"INTEGRADO_O_REMITIDO_SIN_PERDIDA","proof_strength":"EXACTO_INTERNO","provenance":"FORMALIZACION_NUEVA","reinforcement_state":"DEPLOYED_WITHOUT_NEW_RESULT_ID","residence":"subsección «Ordinal, cardinal y ruta»","result_id":"ORD-001","statement":"El ordinal de von Neumann se eleva conservando el estado recorrido; la proyección recupera el ordinal desnudo.","status":"PRESENT_IN_SUCCESSOR_SOURCE"}
% HMT_RESULT_END:ORD-001:residence
% HMT_RESULT_BEGIN:ORD-002:residence
% HMT_ATOMIC_RESULT {"material_state":"INTEGRADO_O_REMITIDO_SIN_PERDIDA","proof_strength":"PROGRAMA_ABIERTO","provenance":"FORMALIZACION_NUEVA","reinforcement_state":"DEPLOYED_WITHOUT_NEW_RESULT_ID","residence":"bloque lógico, § Jerarquía transfinita, tras ORD-001","result_id":"ORD-002","statement":"El corpus cierra profundidades finitas y secciones de longitud \\(\\omega\\); la extensión a todo ordinal es programa.","status":"PRESENT_IN_SUCCESSOR_SOURCE"}
% HMT_RESULT_END:ORD-002:residence
% HMT_RESULT_BEGIN:PWR-001:residence
% HMT_ATOMIC_RESULT {"material_state":"INTEGRADO_O_REMITIDO_SIN_PERDIDA","proof_strength":"EXACTO_INTERNO","provenance":"FORMALIZACION_NUEVA","reinforcement_state":"DEPLOYED_WITHOUT_NEW_RESULT_ID","residence":"subsección «Potencia genealógica»","result_id":"PWR-001","statement":"Los predicados decidibles a profundidad finita forman los clopen del límite.","status":"PRESENT_IN_SUCCESSOR_SOURCE"}
% HMT_RESULT_END:PWR-001:residence
% HMT_RESULT_BEGIN:PWR-002:residence
% HMT_ATOMIC_RESULT {"material_state":"INTEGRADO_O_REMITIDO_SIN_PERDIDA","proof_strength":"EXACTO_INTERNO","provenance":"FORMALIZACION_NUEVA","reinforcement_state":"DEPLOYED_WITHOUT_NEW_RESULT_ID","residence":"bloque lógico, § Potencia cerrada, tras PWR-001","result_id":"PWR-002","statement":"Familias compatibles de subconjuntos finitos reconstruyen cerrados, no todos los subconjuntos.","status":"PRESENT_IN_SUCCESSOR_SOURCE"}
% HMT_RESULT_END:PWR-002:residence
% HMT_RESULT_BEGIN:CHO-001:residence
% HMT_ATOMIC_RESULT {"material_state":"INTEGRADO_O_REMITIDO_SIN_PERDIDA","proof_strength":"EXACTO_INTERNO","provenance":"RESULTADO_RECUPERADO","reinforcement_state":"DEPLOYED_WITHOUT_NEW_RESULT_ID","residence":"subsección «Elección tipada»","result_id":"CHO-001","statement":"Se separan elección conjuntista, natural, algebraica y efectiva. Toda sección natural, algebraica o efectiva es en particular una sección conjuntista, pero naturalidad, algebraicidad y efectividad son requisitos adicionales generalmente incomparables entre sí.","status":"PRESENT_IN_SUCCESSOR_SOURCE"}
% HMT_RESULT_END:CHO-001:residence
% HMT_RESULT_BEGIN:CHO-002:residence
% HMT_ATOMIC_RESULT {"material_state":"INTEGRADO_O_REMITIDO_SIN_PERDIDA","proof_strength":"EXACTO_INTERNO","provenance":"RESULTADO_RECUPERADO","reinforcement_state":"DEPLOYED_WITHOUT_NEW_RESULT_ID","residence":"bloque lógico, § Elección tipada y obstrucción homomorfa, tras CHO-001","result_id":"CHO-002","statement":"Hay representantes conjuntistas pero no sección homomorfa; eso mide memoria estructural, no fracaso de AC.","status":"PRESENT_IN_SUCCESSOR_SOURCE"}
% HMT_RESULT_END:CHO-002:residence
% HMT_RESULT_BEGIN:CAT-001:residence
% HMT_ATOMIC_RESULT {"material_state":"INTEGRADO_O_REMITIDO_SIN_PERDIDA","proof_strength":"EXACTO_INTERNO","provenance":"RESULTADO_RECUPERADO","reinforcement_state":"DEPLOYED_WITHOUT_NEW_RESULT_ID","residence":"cierre categorial del nuevo bloque","result_id":"CAT-001","statement":"Objetos como torres tipadas; morfismos conmutan restricciones; las secciones son el funtor límite.","status":"PRESENT_IN_SUCCESSOR_SOURCE"}
% HMT_RESULT_END:CAT-001:residence
% HMT_RESULT_BEGIN:TOP-001:residence
% HMT_ATOMIC_RESULT {"material_state":"INTEGRADO_O_REMITIDO_SIN_PERDIDA","proof_strength":"EXACTO_INTERNO","provenance":"FORMALIZACION_NUEVA","reinforcement_state":"DEPLOYED_WITHOUT_NEW_RESULT_ID","residence":"nota técnica subordinada a CAT-001","result_id":"TOP-001","statement":"La categoría pequeña de caminos TPK admite una envolvente de presheaves; no se le atribuyen elección o lógica clásica sin prueba.","status":"PRESENT_IN_SUCCESSOR_SOURCE"}
% HMT_RESULT_END:TOP-001:residence
% HMT_RESULT_BEGIN:HIL-001:residence
% HMT_ATOMIC_RESULT {"material_state":"INTEGRADO_O_REMITIDO_SIN_PERDIDA","proof_strength":"EXACTO_INTERNO","provenance":"RESULTADO_RECUPERADO","reinforcement_state":"DEPLOYED_WITHOUT_NEW_RESULT_ID","residence":"`hmt/03l` y `md/04ab`","result_id":"HIL-001","statement":"Cada profundidad posee una realización de Hilbert y una álgebra matricial.","status":"PRESENT_IN_SUCCESSOR_SOURCE"}
% HMT_RESULT_END:HIL-001:residence
% HMT_RESULT_BEGIN:UHF-001:residence
% HMT_ATOMIC_RESULT {"material_state":"INTEGRADO_O_REMITIDO_SIN_PERDIDA","proof_strength":"EXACTO_INTERNO","provenance":"FORMALIZACION_NUEVA","reinforcement_state":"DEPLOYED_WITHOUT_NEW_RESULT_ID","residence":"nuevo bloque «Torre nonádica de observables» tras `hmt/03l`","result_id":"UHF-001","statement":"Conservar simultáneamente las nueve fibras preserva unidad y traza y genera la rama UHF--II_1; una esquina marcada no lo hace y su corredor de tipo I_infty se conserva separadamente en IINF-001.","status":"PRESENT_IN_SUCCESSOR_SOURCE"}
% HMT_RESULT_END:UHF-001:residence
% HMT_RESULT_BEGIN:II1-001:residence
% HMT_ATOMIC_RESULT {"material_state":"INTEGRADO_O_REMITIDO_SIN_PERDIDA","proof_strength":"EXACTO_INTERNO","provenance":"FORMALIZACION_NUEVA","reinforcement_state":"DEPLOYED_WITHOUT_NEW_RESULT_ID","residence":"torre operatorial, § Cierre tracial hiperinfinito II_1, tras IINF-001","result_id":"II1-001","statement":"El límite unital de todas las puertas produce UHF nonádico y su cierre tracial es el factor hiperinfinito \\(II_1\\).","status":"PRESENT_IN_SUCCESSOR_SOURCE"}
% HMT_RESULT_END:II1-001:residence
% HMT_RESULT_BEGIN:DIM-001:residence
% HMT_ATOMIC_RESULT {"material_state":"INTEGRADO_O_REMITIDO_SIN_PERDIDA","proof_strength":"EXACTO_INTERNO","provenance":"ARQUITECTURA_AUTORAL_PREEXISTENTE","reinforcement_state":"DEPLOYED_WITHOUT_NEW_RESULT_ID","residence":"torre operatorial, § Dimensión continua de proyectores, tras II1-001","result_id":"DIM-001","statement":"Para todo t∈[0,1], m_d=floor(9^d t) permite elegir proyectores anidados cuyo límite fuerte P_t satisface tau(P_t)=t; en el factor finito, P~Q iff tau(P)=tau(Q).","status":"PRESENT_IN_SUCCESSOR_SOURCE"}
% HMT_RESULT_END:DIM-001:residence
% HMT_RESULT_BEGIN:CLK-001:residence
% HMT_ATOMIC_RESULT {"material_state":"INTEGRADO_O_REMITIDO_SIN_PERDIDA","proof_strength":"EXACTO_INTERNO","provenance":"RESULTADO_RECUPERADO","reinforcement_state":"DEPLOYED_WITHOUT_NEW_RESULT_ID","residence":"subsección «Realización por reloj profinito»","result_id":"CLK-001","statement":"Una segunda ruta llega abstractamente al mismo factor mediante producto cruzado.","status":"PRESENT_IN_SUCCESSOR_SOURCE"}
% HMT_RESULT_END:CLK-001:residence
% HMT_RESULT_BEGIN:INT-001:residence
% HMT_ATOMIC_RESULT {"material_state":"INTEGRADO_O_REMITIDO_SIN_PERDIDA","proof_strength":"PROGRAMA_ABIERTO","provenance":"FORMALIZACION_NUEVA","reinforcement_state":"DEPLOYED_WITHOUT_NEW_RESULT_ID","residence":"torre operatorial, § Entrelazador APP--reloj, tras CLK-001","result_id":"INT-001","statement":"La isomorfía abstracta de factores no identifica genealogías ni los antecedentes \\(C^*\\).","status":"PRESENT_IN_SUCCESSOR_SOURCE"}
% HMT_RESULT_END:INT-001:residence
% HMT_RESULT_BEGIN:QLG-001:residence
% HMT_ATOMIC_RESULT {"material_state":"INTEGRADO_O_REMITIDO_SIN_PERDIDA","proof_strength":"EXACTO_INTERNO","provenance":"RESULTADO_RECUPERADO","reinforcement_state":"DEPLOYED_WITHOUT_NEW_RESULT_ID","residence":"`md/04ab` después de PVM/contextos","result_id":"QLG-001","statement":"\\(\\operatorname{Proj}(R_{\\rm hyp})\\) es un retículo ortomodular completo, continuo y no distributivo; cada contexto conmutativo es booleano.","status":"PRESENT_IN_SUCCESSOR_SOURCE"}
% HMT_RESULT_END:QLG-001:residence
% HMT_RESULT_BEGIN:ERG-001:residence
% HMT_ATOMIC_RESULT {"material_state":"INTEGRADO_O_REMITIDO_SIN_PERDIDA","proof_strength":"EXACTO_INTERNO","provenance":"RESULTADO_RECUPERADO","reinforcement_state":"DEPLOYED_WITHOUT_NEW_RESULT_ID","residence":"capítulo de torre/reloj","result_id":"ERG-001","statement":"Los promedios temporales proyectan sobre invariantes; para el reloj ergódico, sobre la media de Haar.","status":"PRESENT_IN_SUCCESSOR_SOURCE"}
% HMT_RESULT_END:ERG-001:residence
% HMT_RESULT_BEGIN:DEN-001:residence
% HMT_ATOMIC_RESULT {"material_state":"INTEGRADO_O_REMITIDO_SIN_PERDIDA","proof_strength":"EXACTO_INTERNO","provenance":"RESULTADO_RECUPERADO","reinforcement_state":"DEPLOYED_WITHOUT_NEW_RESULT_ID","residence":"`md/04ab` y torre \\(II_1\\)","result_id":"DEN-001","statement":"La incertidumbre sobre rutas se representa por estado positivo; no obliga a borrar la genealogía subyacente.","status":"PRESENT_IN_SUCCESSOR_SOURCE"}
% HMT_RESULT_END:DEN-001:residence
% HMT_RESULT_BEGIN:LAN-001:residence
% HMT_ATOMIC_RESULT {"material_state":"INTEGRADO_O_REMITIDO_SIN_PERDIDA","proof_strength":"EXACTO_INTERNO","provenance":"RESULTADO_RECUPERADO","reinforcement_state":"DEPLOYED_WITHOUT_NEW_RESULT_ID","residence":"`md/04g`","result_id":"LAN-001","statement":"Una actualización biyectiva U del estado completo satisface H(X|U(X))=0 y, para el término ideal de borrado lógico, Q_L^{min}[U]=0.","status":"PRESENT_IN_SUCCESSOR_SOURCE"}
% HMT_RESULT_END:LAN-001:residence
% HMT_RESULT_BEGIN:LAN-002:residence
% HMT_ATOMIC_RESULT {"material_state":"INTEGRADO_O_REMITIDO_SIN_PERDIDA","proof_strength":"EXACTO_INTERNO","provenance":"FORMALIZACION_NUEVA","reinforcement_state":"DEPLOYED_WITHOUT_NEW_RESULT_ID","residence":"`md/04g`, después del caso finito","result_id":"LAN-002","statement":"Automorfismos traza-preservantes conservan entropía; una esperanza condicional incrementa la entropía exactamente por entropía relativa.","status":"PRESENT_IN_SUCCESSOR_SOURCE"}
% HMT_RESULT_END:LAN-002:residence
% HMT_RESULT_BEGIN:GDL-001:residence
% HMT_ATOMIC_RESULT {"material_state":"INTEGRADO_O_REMITIDO_SIN_PERDIDA","proof_strength":"EXACTO_INTERNO","provenance":"RESULTADO_RECUPERADO","reinforcement_state":"DEPLOYED_WITHOUT_NEW_RESULT_ID","residence":"nuevo bloque «Límite efectivo de la clausura»","result_id":"GDL-001","statement":"Cada nivel puede calcularse mientras la unicidad de la fibra global codifica HALT.","status":"PRESENT_IN_SUCCESSOR_SOURCE"}
% HMT_RESULT_END:GDL-001:residence
% HMT_RESULT_BEGIN:GDL-002:residence
% HMT_ATOMIC_RESULT {"material_state":"INTEGRADO_O_REMITIDO_SIN_PERDIDA","proof_strength":"EXACTO_INTERNO","provenance":"RESULTADO_RECUPERADO","reinforcement_state":"DEPLOYED_WITHOUT_NEW_RESULT_ID","residence":"bloque metalógico, § Incompletitud semántica Gödel--Turing, tras GDL-001","result_id":"GDL-002","statement":"Toda teoría recursivamente axiomatizable y correcta para la familia \\(U_e\\) deja alguna sentencia sin decidir; no implica inconsistencia de ZFC.","status":"PRESENT_IN_SUCCESSOR_SOURCE"}
% HMT_RESULT_END:GDL-002:residence
% HMT_RESULT_BEGIN:IND-001:residence
% HMT_ATOMIC_RESULT {"material_state":"INTEGRADO_O_REMITIDO_SIN_PERDIDA","proof_strength":"EXACTO_INTERNO","provenance":"FORMALIZACION_NUEVA","reinforcement_state":"DEPLOYED_WITHOUT_NEW_RESULT_ID","residence":"bloque metalógico, § Frontera Gödel--Cohen, tras GDL-002","result_id":"IND-001","statement":"Bajo las hipótesis de consistencia correspondientes, CH es independiente de ZFC y AC es independiente de ZF; HMT enriquece objetos, no decide por sí sola esos enunciados.","status":"PRESENT_IN_SUCCESSOR_SOURCE"}
% HMT_RESULT_END:IND-001:residence
% HMT_RESULT_BEGIN:FOR-001:residence
% HMT_ATOMIC_RESULT {"material_state":"INTEGRADO_O_REMITIDO_SIN_PERDIDA","proof_strength":"EXACTO_INTERNO","provenance":"FORMALIZACION_NUEVA","reinforcement_state":"DEPLOYED_WITHOUT_NEW_RESULT_ID","residence":"subsección de filtros de prefijos","result_id":"FOR-001","statement":"Una sección determina un filtro que atraviesa todos los densos de profundidad; eso no basta para genericidad sobre un modelo.","status":"PRESENT_IN_SUCCESSOR_SOURCE"}
% HMT_RESULT_END:FOR-001:residence
% HMT_RESULT_BEGIN:FOR-002:residence
% HMT_ATOMIC_RESULT {"material_state":"INTEGRADO_O_REMITIDO_SIN_PERDIDA","proof_strength":"INTERFAZ_FISICA","provenance":"FORMALIZACION_NUEVA","reinforcement_state":"DEPLOYED_WITHOUT_NEW_RESULT_ID","residence":"bloque metalógico, § Forcing cilíndrico y filtro genérico, tras FOR-001","result_id":"FOR-002","statement":"Árbol perfecto, numerable y sin átomos da forcing cilíndrico Cohen-equivalente; rama única da forcing trivial.","status":"PRESENT_IN_SUCCESSOR_SOURCE"}
% HMT_RESULT_END:FOR-002:residence
% HMT_RESULT_BEGIN:BOR-001:residence
% HMT_ATOMIC_RESULT {"material_state":"INTEGRADO_O_REMITIDO_SIN_PERDIDA","proof_strength":"EXACTO_INTERNO","provenance":"FORMALIZACION_NUEVA","reinforcement_state":"DEPLOYED_WITHOUT_NEW_RESULT_ID","residence":"subsección final del bloque","result_id":"BOR-001","statement":"Los cilindros generan una clase Borel estable por complemento, unión numerable e imágenes inversas continuas; sus juegos usan determinación boreliana clásica.","status":"PRESENT_IN_SUCCESSOR_SOURCE"}
% HMT_RESULT_END:BOR-001:residence
% HMT_RESULT_BEGIN:MIN-001:residence
% HMT_ATOMIC_RESULT {"material_state":"INTEGRADO_O_REMITIDO_SIN_PERDIDA","proof_strength":"INTERFAZ_FISICA","provenance":"FORMALIZACION_NUEVA","reinforcement_state":"DEPLOYED_WITHOUT_NEW_RESULT_ID","residence":"bloque metalógico, § Minimax finito y proyectivo, tras BOR-001","result_id":"MIN-001","statement":"En cada profundidad finita, von Neumann da valor mixto; un límite global requiere compactitud, continuidad y consistencia proyectiva, y no selecciona por sí solo rama pura.","status":"PRESENT_IN_SUCCESSOR_SOURCE"}
% HMT_RESULT_END:MIN-001:residence
% HMT_RESULT_BEGIN:III-001:residence
% HMT_ATOMIC_RESULT {"material_state":"INTEGRADO_O_REMITIDO_SIN_PERDIDA","proof_strength":"PROGRAMA_ABIERTO","provenance":"RESULTADO_NUEVO","reinforcement_state":"DEPLOYED_WITHOUT_NEW_RESULT_ID","residence":"frontera tras \\(II_1\\)","result_id":"III-001","statement":"La construcción actual produce naturalmente \\(II_1\\); tipo \\(III\\) exige GNS no tracial/KMS o cociclo de Radon--Nikodym no trivial.","status":"PRESENT_IN_SUCCESSOR_SOURCE"}
% HMT_RESULT_END:III-001:residence
% HMT_RESULT_BEGIN:TRI-001:residence
% HMT_ATOMIC_RESULT {"material_state":"INTEGRADO_O_REMITIDO_SIN_PERDIDA","proof_strength":"INTERFAZ_FISICA","provenance":"FORMALIZACION_NUEVA","reinforcement_state":"DEPLOYED_WITHOUT_NEW_RESULT_ID","residence":"después de exponer íntegramente la doble proyección y completar el corredor Paley--Witt--Golay--Leech--\\(V^\\natural\\)--Monstruo--Moonshine","result_id":"TRI-001","statement":"Punto, dimensión y genealogía requieren tres capas coordinadas.","status":"PRESENT_IN_SUCCESSOR_SOURCE"}
% HMT_RESULT_END:TRI-001:residence
% HMT_RESULT_BEGIN:EXC-001:residence
% HMT_ATOMIC_RESULT {"material_state":"INTEGRADO_O_REMITIDO_SIN_PERDIDA","proof_strength":"DEMOSTRADO_CON_ESTRUCTURA_DE_PARTIDA_EXPLICITA","provenance":"RESULTADO_RECUPERADO","reinforcement_state":"DEPLOYED_WITHOUT_NEW_RESULT_ID","residence":"`hmt/09g`, `15`","result_id":"EXC-001","statement":"El mismo estado superviviente alimenta lectura real e incidencia excepcional hasta Leech--\\(V^\\natural\\)--Monstruo/Moonshine; el Monstruo no actúa literalmente sobre dígitos.","status":"PRESENT_IN_SUCCESSOR_SOURCE"}
% HMT_RESULT_END:EXC-001:residence
% HMT_RESULT_BEGIN:MD-001:residence
% HMT_ATOMIC_RESULT {"material_state":"INTEGRADO_O_REMITIDO_SIN_PERDIDA","proof_strength":"PROGRAMA_ABIERTO","provenance":"RESULTADO_RECUPERADO","reinforcement_state":"DEPLOYED_WITHOUT_NEW_RESULT_ID","residence":"`md/00b`, antes de Planck/electrón","result_id":"MD-001","statement":"Una sección compatible no es todavía una partícula: la ruta es su historia observable y la partícula una sección persistente enriquecida antes del morfismo de realización física.","status":"PRESENT_IN_SUCCESSOR_SOURCE"}
% HMT_RESULT_END:MD-001:residence
% HMT_RESULT_BEGIN:PSC-001:residence
% HMT_ATOMIC_RESULT {"material_state":"INTEGRADO_O_REMITIDO_SIN_PERDIDA","proof_strength":"EXACTO_INTERNO","provenance":"ARQUITECTURA_AUTORAL_PREEXISTENTE","reinforcement_state":"DEPLOYED_WITHOUT_NEW_RESULT_ID","residence":"apéndice crítico, fuera del tronco probatorio","result_id":"PSC-001","statement":"Las ramas AD/MM, CH y «cierre fuerte» conservan prioridad intuitiva, no fuerza probatoria vigente.","status":"PRESENT_IN_SUCCESSOR_SOURCE"}
% HMT_RESULT_END:PSC-001:residence
% HMT_RESULT_BEGIN:COV-001:residence
% HMT_ATOMIC_RESULT {"material_state":"INTEGRADO_O_REMITIDO_SIN_PERDIDA","proof_strength":"EXACTO_INTERNO","provenance":"RESULTADO_RECUPERADO","reinforcement_state":"DEPLOYED_WITHOUT_NEW_RESULT_ID","residence":"gate de cobertura final","result_id":"COV-001","statement":"Constantes, electrón, masas, Barbero, CKM/CP, gravedad y corredor excepcional se preservan por remisión, sin versión degradada.","status":"PRESENT_IN_SUCCESSOR_SOURCE"}
% HMT_RESULT_END:COV-001:residence
% HMT_RESULT_BEGIN:IINF-001:residence
% HMT_ATOMIC_RESULT {"material_state":"INTEGRADO_O_REMITIDO_SIN_PERDIDA","proof_strength":"EXACTO_INTERNO","provenance":"RESULTADO_RECUPERADO","reinforcement_state":"DEPLOYED_WITHOUT_NEW_RESULT_ID","residence":"torre operatorial, § Rama marcada de tipo I_infty, tras UHF-001","result_id":"IINF-001","statement":"Las esquinas j_{d,j}(a)=a⊗p_j, implementadas por V_{d,j}aV_{d,j}^*, generan en el límite los operadores de rango finito; su clausura normativa es K(H_infty) y su clausura débil B(H_infty), factor de tipo I_infty.","status":"PRESENT_IN_SUCCESSOR_SOURCE"}
% HMT_RESULT_END:IINF-001:residence
% HMT_RESULT_BEGIN:DEN-003:residence
% HMT_ATOMIC_RESULT {"material_state":"INTEGRADO_O_REMITIDO_SIN_PERDIDA","proof_strength":"EXACTO_INTERNO","provenance":"RESULTADO_RECUPERADO","reinforcement_state":"DEPLOYED_WITHOUT_NEW_RESULT_ID","residence":"torre \\(II_1\\), inmediatamente después de DEN-002","result_id":"DEN-003","statement":"Si h=9^d rho, entonces log h=d log9 I+log rho y S_tau(h)=S_vN(rho)-d log9; bajo rho_{d+1}=rho_d⊗I9/9, S_vN aumenta en log9 y S_tau permanece invariante.","status":"PRESENT_IN_SUCCESSOR_SOURCE"}
% HMT_RESULT_END:DEN-003:residence
% HMT_RESULT_BEGIN:CARD-001:residence
% HMT_ATOMIC_RESULT {"material_state":"INTEGRADO_O_REMITIDO_SIN_PERDIDA","proof_strength":"EXACTO_INTERNO","provenance":"RESULTADO_RECUPERADO","reinforcement_state":"DEPLOYED_WITHOUT_NEW_RESULT_ID","residence":"bloque lógico, § Cardinalidad como funtor no fiel, tras ORD-002","result_id":"CARD-001","statement":"Contar es decategorizar: el perfil (|X_m|)_m no determina restricciones, acarreo, orientación o memoria.","status":"PRESENT_IN_SUCCESSOR_SOURCE"}
% HMT_RESULT_END:CARD-001:residence
% HMT_RESULT_BEGIN:PWR-003:residence
% HMT_ATOMIC_RESULT {"material_state":"INTEGRADO_O_REMITIDO_SIN_PERDIDA","proof_strength":"EXACTO_INTERNO","provenance":"FORMALIZACION_NUEVA","reinforcement_state":"DEPLOYED_WITHOUT_NEW_RESULT_ID","residence":"bloque lógico, § Tres escalas de potencia, tras PWR-002","result_id":"PWR-003","statement":"En un Cantor HMT, clopen, cerrados y potencia plena tienen tamaños distintos.","status":"PRESENT_IN_SUCCESSOR_SOURCE"}
% HMT_RESULT_END:PWR-003:residence
% HMT_RESULT_BEGIN:WEYL-001:residence
% HMT_ATOMIC_RESULT {"material_state":"INTEGRADO_O_REMITIDO_SIN_PERDIDA","proof_strength":"EXACTO_INTERNO","provenance":"RESULTADO_RECUPERADO","reinforcement_state":"DEPLOYED_WITHOUT_NEW_RESULT_ID","residence":"`hmt/14_numero_heisenberg_y_d108.tex`, `md/04ab`","result_id":"WEYL-001","statement":"Fase y desplazamiento forman un par no conmutativo; fijado el carácter central primitivo, la representación irreducible es única hasta unitario.","status":"PRESENT_IN_SUCCESSOR_SOURCE"}
% HMT_RESULT_END:WEYL-001:residence
% HMT_RESULT_BEGIN:DSC-001:residence
% HMT_ATOMIC_RESULT {"material_state":"INTEGRADO_O_REMITIDO_SIN_PERDIDA","proof_strength":"EXACTO_INTERNO","provenance":"RESULTADO_RECUPERADO","reinforcement_state":"DEPLOYED_WITHOUT_NEW_RESULT_ID","residence":"bloque fundacional, § Presentación y descenso por fibras, tras INF-003","result_id":"DSC-001","statement":"Existe un único operador descendido O_bar:V→Y con O=O_bar q si y sólo si O es constante en cada fibra de q.","status":"PRESENT_IN_SUCCESSOR_SOURCE"}
% HMT_RESULT_END:DSC-001:residence
% HMT_RESULT_BEGIN:REC-001:residence
% HMT_ATOMIC_RESULT {"material_state":"INTEGRADO_O_REMITIDO_SIN_PERDIDA","proof_strength":"EXACTO_INTERNO","provenance":"RESULTADO_RECUPERADO","reinforcement_state":"DEPLOYED_WITHOUT_NEW_RESULT_ID","residence":"`md/04g`, antes del teorema de Landauer lógico","result_id":"REC-001","statement":"Si q_hat=(q,M) es inyectivo sobre supp(X), entonces H(X|q(X),M(X))=0; H(X|q(X)) mide la fibra olvidada por el lector visible.","status":"PRESENT_IN_SUCCESSOR_SOURCE"}
% HMT_RESULT_END:REC-001:residence
% HMT_RESULT_BEGIN:LAN-003:residence
% HMT_ATOMIC_RESULT {"material_state":"INTEGRADO_O_REMITIDO_SIN_PERDIDA","proof_strength":"EXACTO_INTERNO","provenance":"RESULTADO_RECUPERADO","reinforcement_state":"DEPLOYED_WITHOUT_NEW_RESULT_ID","residence":"`md/04g`, después de LAN-001","result_id":"LAN-003","statement":"Para un reset no inyectivo de tres estados equiprobables, H(X)=log 3 y el término ideal satisface Q_reset≥k_B Θ log 3.","status":"PRESENT_IN_SUCCESSOR_SOURCE"}
% HMT_RESULT_END:LAN-003:residence
% HMT_RESULT_BEGIN:FIN-001:residence
% HMT_ATOMIC_RESULT {"material_state":"INTEGRADO_O_REMITIDO_SIN_PERDIDA","proof_strength":"EXACTO_INTERNO","provenance":"FORMALIZACION_NUEVA","reinforcement_state":"DEPLOYED_WITHOUT_NEW_RESULT_ID","residence":"nuevo bloque «Límite efectivo de la clausura», antes de GDL-001","result_id":"FIN-001","statement":"Fijado N, con dominios finitos y observables atómicos computables, toda fórmula del lenguaje finito L_N es semánticamente decidible.","status":"PRESENT_IN_SUCCESSOR_SOURCE"}
% HMT_RESULT_END:FIN-001:residence
% HMT_RESULT_BEGIN:DEN-002:residence
% HMT_ATOMIC_RESULT {"material_state":"INTEGRADO_O_REMITIDO_SIN_PERDIDA","proof_strength":"EXACTO_INTERNO","provenance":"RESULTADO_RECUPERADO","reinforcement_state":"DEPLOYED_WITHOUT_NEW_RESULT_ID","residence":"torre \\(II_1\\), después de DEN-001","result_id":"DEN-002","statement":"Para h=9^d rho, S_tau(h)=S_vN(rho)-d log9; la extensión rho_{d+1}=rho_d⊗I9/9 aumenta S_vN en log9 y deja invariante S_tau.","status":"PRESENT_IN_SUCCESSOR_SOURCE"}
% HMT_RESULT_END:DEN-002:residence
% HMT_RESULT_BEGIN:UNI-001:residence
% HMT_ATOMIC_RESULT {"material_state":"INTEGRADO_O_REMITIDO_SIN_PERDIDA","proof_strength":"EXACTO_INTERNO","provenance":"RESULTADO_RECUPERADO","reinforcement_state":"DEPLOYED_WITHOUT_NEW_RESULT_ID","residence":"cierre APP--EMRH","result_id":"UNI-001","statement":"La frontera distingue 999=729+243+27 de 1000=729+243+27+1=729+271; el acarreo c_{m+1}=1 conserva la unidad visible y registra una vuelta adicional.","status":"PRESENT_IN_SUCCESSOR_SOURCE"}
% HMT_RESULT_END:UNI-001:residence
% HMT_RESULT_BEGIN:UNI-002:residence
% HMT_ATOMIC_RESULT {"material_state":"INTEGRADO_O_REMITIDO_SIN_PERDIDA","proof_strength":"PROGRAMA_ABIERTO","provenance":"ARQUITECTURA_AUTORAL_PREEXISTENTE","reinforcement_state":"DEPLOYED_WITHOUT_NEW_RESULT_ID","residence":"transición APP--EMRH hacia monodromía e inclusión unital nonádica","result_id":"UNI-002","statement":"La compatibilidad entre c_{m+1}=1, la monodromía 9→1 y a↦a⊗I9 es una obligación de entrelazamiento: debe construirse un diagrama que transporte la misma unidad con memoria entre los tres dominios.","status":"PRESENT_IN_SUCCESSOR_SOURCE"}
% HMT_RESULT_END:UNI-002:residence
% HMT_RESULT_BEGIN:P01:residence
% HMT_ATOMIC_RESULT {"material_state":"INTEGRADO_EN_GRAFO_ACTIVO","proof_strength":"EXACTO_INTERNO","provenance":"FORMALIZACION_NUEVA","reinforcement_state":"DEPLOYED_WITHOUT_NEW_RESULT_ID","residence":"APP vigente","result_id":"P01","statement":"APP tiene soporte $X_9$, 81 vértices y todos los caminos finitos.","status":"PRESENT_IN_SUCCESSOR_SOURCE"}
% HMT_RESULT_END:P01:residence
% HMT_RESULT_BEGIN:P02:residence
% HMT_ATOMIC_RESULT {"material_state":"INTEGRADO_EN_GRAFO_ACTIVO","proof_strength":"EXACTO_INTERNO","provenance":"FORMALIZACION_NUEVA","reinforcement_state":"DEPLOYED_WITHOUT_NEW_RESULT_ID","residence":"APP vigente","result_id":"P02","statement":"$\\Sigma$ es latina; $\\Pi$ sólo permuta en unidades.","status":"PRESENT_IN_SUCCESSOR_SOURCE"}
% HMT_RESULT_END:P02:residence
% HMT_RESULT_BEGIN:P03:residence
% HMT_ATOMIC_RESULT {"material_state":"INTEGRADO_EN_GRAFO_ACTIVO","proof_strength":"EXACTO_INTERNO","provenance":"FORMALIZACION_NUEVA","reinforcement_state":"DEPLOYED_WITHOUT_NEW_RESULT_ID","residence":"APP/TRIT","result_id":"P03","statement":"La elevación residuo--cociente recompone $1215=54+9\\cdot129$.","status":"PRESENT_IN_SUCCESSOR_SOURCE"}
% HMT_RESULT_END:P03:residence
% HMT_RESULT_BEGIN:P04:residence
% HMT_ATOMIC_RESULT {"material_state":"INTEGRADO_EN_GRAFO_ACTIVO","proof_strength":"EXACTO_INTERNO","provenance":"FORMALIZACION_NUEVA","reinforcement_state":"DEPLOYED_WITHOUT_NEW_RESULT_ID","residence":"`hmt/03f`","result_id":"P04","statement":"$B_c=9P_++5P_-$ y el salto es cuatro.","status":"PRESENT_IN_SUCCESSOR_SOURCE"}
% HMT_RESULT_END:P04:residence
% HMT_RESULT_BEGIN:P05:residence
% HMT_ATOMIC_RESULT {"material_state":"INTEGRADO_EN_GRAFO_ACTIVO","proof_strength":"EXACTO_INTERNO","provenance":"FORMALIZACION_NUEVA","reinforcement_state":"DEPLOYED_WITHOUT_NEW_RESULT_ID","residence":"TRIT","result_id":"P05","statement":"TRIT visible no determina el estado levantado.","status":"PRESENT_IN_SUCCESSOR_SOURCE"}
% HMT_RESULT_END:P05:residence
% HMT_RESULT_BEGIN:P06:residence
% HMT_ATOMIC_RESULT {"material_state":"INTEGRADO_EN_GRAFO_ACTIVO","proof_strength":"EXACTO_INTERNO","provenance":"FORMALIZACION_NUEVA","reinforcement_state":"DEPLOYED_WITHOUT_NEW_RESULT_ID","residence":"TPK","result_id":"P06","statement":"TPK compone rutas con hoja, acarreo y memoria.","status":"PRESENT_IN_SUCCESSOR_SOURCE"}
% HMT_RESULT_END:P06:residence
% HMT_RESULT_BEGIN:P07:residence
% HMT_ATOMIC_RESULT {"material_state":"INTEGRADO_EN_GRAFO_ACTIVO","proof_strength":"EXACTO_INTERNO","provenance":"RESULTADO_RECUPERADO","reinforcement_state":"DEPLOYED_WITHOUT_NEW_RESULT_ID","residence":"`md/04ab`","result_id":"P07","statement":"Los operadores de traslación X_u y fase Z_v satisfacen Z_vX_u=omega^{u·v}X_uZ_v, generan M_{9^d}(C) y, fijado el carácter central primitivo, dan la representación irreducible finita única salvo equivalencia unitaria.","status":"PRESENT_IN_SUCCESSOR_SOURCE"}
% HMT_RESULT_END:P07:residence
% HMT_RESULT_BEGIN:P08:residence
% HMT_ATOMIC_RESULT {"material_state":"INTEGRADO_EN_GRAFO_ACTIVO","proof_strength":"EXACTO_INTERNO","provenance":"RESULTADO_RECUPERADO","reinforcement_state":"DEPLOYED_WITHOUT_NEW_RESULT_ID","residence":"bloque UHF--\\(II_1\\)","result_id":"P08","statement":"El límite inductivo de a↦a⊗I_9 es UHF(9^infty), canónicamente isomorfa a UHF(3^infty); la representación GNS de su traza única cierra en el factor hiperinfinito II_1.","status":"PRESENT_IN_SUCCESSOR_SOURCE"}
% HMT_RESULT_END:P08:residence
% HMT_RESULT_BEGIN:P09:residence
% HMT_ATOMIC_RESULT {"material_state":"INTEGRADO_EN_GRAFO_ACTIVO","proof_strength":"EXACTO_INTERNO","provenance":"FORMALIZACION_NUEVA","reinforcement_state":"DEPLOYED_WITHOUT_NEW_RESULT_ID","residence":"apertura del corredor matricial","result_id":"P09","statement":"Un QMS es una cuadrícula positiva con sumas $I_s$.","status":"PRESENT_IN_SUCCESSOR_SOURCE"}
% HMT_RESULT_END:P09:residence
% HMT_RESULT_BEGIN:P10:residence
% HMT_ATOMIC_RESULT {"material_state":"INTEGRADO_EN_GRAFO_ACTIVO","proof_strength":"EXACTO_INTERNO","provenance":"RESULTADO_RECUPERADO","reinforcement_state":"DEPLOYED_WITHOUT_NEW_RESULT_ID","residence":"corredor matricial, § Semiclasicidad como factorización conmutativa, tras P09/QPM-001","result_id":"P10","statement":"A es semiclásico si y sólo si existe un mapa positivo unital Phi:C(S_n)→M_s tal que A_ij=Phi(f_ij), donde f_ij(sigma)=1_{sigma(i)=j}; por ser conmutativo el dominio, Phi es completamente positivo.","status":"PRESENT_IN_SUCCESSOR_SOURCE"}
% HMT_RESULT_END:P10:residence
% HMT_RESULT_BEGIN:P11:residence
% HMT_ATOMIC_RESULT {"material_state":"INTEGRADO_EN_GRAFO_ACTIVO","proof_strength":"EXACTO_INTERNO","provenance":"RESULTADO_RECUPERADO","reinforcement_state":"DEPLOYED_WITHOUT_NEW_RESULT_ID","residence":"corredor matricial, § QLS fáciles y semiclásicos, tras QLS-001","result_id":"P11","statement":"Un cuadrado latino cuántico es semiclásico si y sólo si es fácil, es decir, procede de un cuadrado latino clásico y una base ortonormal fija.","status":"PRESENT_IN_SUCCESSOR_SOURCE"}
% HMT_RESULT_END:P11:residence
% HMT_RESULT_BEGIN:P12:residence
% HMT_ATOMIC_RESULT {"material_state":"INTEGRADO_EN_GRAFO_ACTIVO","proof_strength":"EXACTO_INTERNO","provenance":"FORMALIZACION_NUEVA","reinforcement_state":"DEPLOYED_WITHOUT_NEW_RESULT_ID","residence":"capítulo qutrit","result_id":"P12","statement":"Cuatro contextos HMT producen cuatro QLS fáciles.","status":"PRESENT_IN_SUCCESSOR_SOURCE"}
% HMT_RESULT_END:P12:residence
% HMT_RESULT_BEGIN:P13:residence
% HMT_ATOMIC_RESULT {"material_state":"INTEGRADO_EN_GRAFO_ACTIVO","proof_strength":"EXACTO_INTERNO","provenance":"RESULTADO_NUEVO","reinforcement_state":"DEPLOYED_WITHOUT_NEW_RESULT_ID","residence":"cuadrados operatoriales, § No-go de QLS cíclicos, tras P12","result_id":"P13","statement":"Dos filas MUB no pueden apilarse en una QLS.","status":"PRESENT_IN_SUCCESSOR_SOURCE"}
% HMT_RESULT_END:P13:residence
% HMT_RESULT_BEGIN:P14:residence
% HMT_ATOMIC_RESULT {"material_state":"INTEGRADO_EN_GRAFO_ACTIVO","proof_strength":"EXACTO_INTERNO","provenance":"FORMALIZACION_NUEVA","reinforcement_state":"DEPLOYED_WITHOUT_NEW_RESULT_ID","residence":"tras TPK y torre UHF","result_id":"P14","statement":"Para n fijo, profundidad genealógica m y nivel matricial s definen una bigraduación; si n varía, la familia es trigraduada por (m,n,s).","status":"PRESENT_IN_SUCCESSOR_SOURCE"}
% HMT_RESULT_END:P14:residence
% HMT_RESULT_BEGIN:P15:residence
% HMT_ATOMIC_RESULT {"material_state":"INTEGRADO_EN_GRAFO_ACTIVO","proof_strength":"EXACTO_INTERNO","provenance":"FORMALIZACION_NUEVA","reinforcement_state":"DEPLOYED_WITHOUT_NEW_RESULT_ID","residence":"sistema inverso UCP, § Compresiones y truncación, tras P14","result_id":"P15","statement":"Si Phi_r∈K_{m+1,s_r} y V_r:C^s→C^{s_r} satisfacen sum_r V_r^*V_r=I_s, entonces la columna V=(V_1,\\ldots,V_k)^T es isométrica y la restricción genealógica conmuta con la combinación matricial: rho_m^s(sum_r V_r^*Phi_r V_r)=sum_r V_r^*(rho_m^{s_r}Phi_r)V_r.","status":"PRESENT_IN_SUCCESSOR_SOURCE"}
% HMT_RESULT_END:P15:residence
% HMT_RESULT_BEGIN:P16:residence
% HMT_ATOMIC_RESULT {"material_state":"INTEGRADO_EN_GRAFO_ACTIVO","proof_strength":"EXACTO_INTERNO","provenance":"FORMALIZACION_NUEVA","reinforcement_state":"DEPLOYED_WITHOUT_NEW_RESULT_ID","residence":"sistema inverso UCP, § Dilatación cilíndrica compatible y PVM común, tras P15","result_id":"P16","statement":"Existe una única Phi_infty:C(X_infty)→M_s que extiende todas las cartas cilíndricas y una dilatación común Phi_infty(f)=V*∫f dP V, con E_x^(m)=V*P([x]_m)V; la dilatación mínima es única hasta unitario y, fijada pi, la PVM P es única.","status":"PRESENT_IN_SUCCESSOR_SOURCE"}
% HMT_RESULT_END:P16:residence
% HMT_RESULT_BEGIN:P17:residence
% HMT_ATOMIC_RESULT {"material_state":"INTEGRADO_EN_GRAFO_ACTIVO","proof_strength":"EXACTO_INTERNO","provenance":"FORMALIZACION_NUEVA","reinforcement_state":"DEPLOYED_WITHOUT_NEW_RESULT_ID","residence":"continuación UHF","result_id":"P17","statement":"Si Psi_m=Psi_{m+1}∘iota_m, existe una única UCP Psi_infty:A_infty→M_s y una dilatación común Psi_m(a)=V* pi(iota_{m,infty}(a))V; pi puede hacerse fiel.","status":"PRESENT_IN_SUCCESSOR_SOURCE"}
% HMT_RESULT_END:P17:residence
% HMT_RESULT_BEGIN:P18:residence
% HMT_ATOMIC_RESULT {"material_state":"INTEGRADO_EN_GRAFO_ACTIVO","proof_strength":"EXACTO_INTERNO","provenance":"FORMALIZACION_NUEVA","reinforcement_state":"DEPLOYED_WITHOUT_NEW_RESULT_ID","residence":"capítulo «Cuadrados cuánticos de APP»","result_id":"P18","statement":"Existe QMS aditivo $(9,3)$.","status":"PRESENT_IN_SUCCESSOR_SOURCE"}
% HMT_RESULT_END:P18:residence
% HMT_RESULT_BEGIN:P19:residence
% HMT_ATOMIC_RESULT {"material_state":"INTEGRADO_EN_GRAFO_ACTIVO","proof_strength":"EXACTO_INTERNO","provenance":"RESULTADO_NUEVO","reinforcement_state":"DEPLOYED_WITHOUT_NEW_RESULT_ID","residence":"cuadrados APP, § Construcción A^Sigma, tras P18","result_id":"P19","statement":"Tres contextos producen QMS $(9,3)$ no conmutativo y semiclásico.","status":"PRESENT_IN_SUCCESSOR_SOURCE"}
% HMT_RESULT_END:P19:residence
% HMT_RESULT_BEGIN:P20:residence
% HMT_ATOMIC_RESULT {"material_state":"INTEGRADO_EN_GRAFO_ACTIVO","proof_strength":"EXACTO_INTERNO","provenance":"RESULTADO_NUEVO","reinforcement_state":"DEPLOYED_WITHOUT_NEW_RESULT_ID","residence":"cuadrados APP, § Construcción A^(9), tras P19","result_id":"P20","statement":"Las doce direcciones producen QMS $(12,3)$ semiclásico.","status":"PRESENT_IN_SUCCESSOR_SOURCE"}
% HMT_RESULT_END:P20:residence
% HMT_RESULT_BEGIN:P21:residence
% HMT_ATOMIC_RESULT {"material_state":"INTEGRADO_EN_GRAFO_ACTIVO","proof_strength":"EXACTO_INTERNO","provenance":"RESULTADO_NUEVO","reinforcement_state":"DEPLOYED_WITHOUT_NEW_RESULT_ID","residence":"cuadrados APP, § Construcción B(12,3) y cubo de tres índices, tras P20","result_id":"P21","statement":"Las mismas \\POVM producen cubos de órdenes 9 y 12.","status":"PRESENT_IN_SUCCESSOR_SOURCE"}
% HMT_RESULT_END:P21:residence
% HMT_RESULT_BEGIN:P22:residence
% HMT_ATOMIC_RESULT {"material_state":"INTEGRADO_EN_GRAFO_ACTIVO","proof_strength":"EXACTO_INTERNO","provenance":"FORMALIZACION_NUEVA","reinforcement_state":"DEPLOYED_WITHOUT_NEW_RESULT_ID","residence":"subsección de fibración finita","result_id":"P22","statement":"P^1(Z/9Z) tiene doce puntos y r tiene cuatro fibras de tres; la fibración es canónica, mientras el orden de cada fibra y la bijección Xi hacia resultados qutrit constituyen un calibre.","status":"PRESENT_IN_SUCCESSOR_SOURCE"}
% HMT_RESULT_END:P22:residence
% HMT_RESULT_BEGIN:P23:residence
% HMT_ATOMIC_RESULT {"material_state":"INTEGRADO_EN_GRAFO_ACTIVO","proof_strength":"EXACTO_INTERNO","provenance":"FORMALIZACION_NUEVA","reinforcement_state":"DEPLOYED_WITHOUT_NEW_RESULT_ID","residence":"geometría qutrit, § Fibración 12→4×3 y POVM calibrada, tras P22","result_id":"P23","statement":"F_k^(pi)=(3/pi)P_k para k=0,1,2 y F_r^(pi)=((1-3/pi)/6)I_3 para r=3,…,8 forman una POVM; para un proyector Q de otro contexto MUB, Tr(QF_k^(pi))=(3/pi)(1/3)=1/pi","status":"PRESENT_IN_SUCCESSOR_SOURCE"}
% HMT_RESULT_END:P23:residence
% HMT_RESULT_BEGIN:P24:residence
% HMT_ATOMIC_RESULT {"material_state":"INTEGRADO_EN_GRAFO_ACTIVO","proof_strength":"EXACTO_INTERNO","provenance":"RESULTADO_NUEVO","reinforcement_state":"DEPLOYED_WITHOUT_NEW_RESULT_ID","residence":"después de separación espectral APP","result_id":"P24","statement":"El polinomio de CC* es t^5(t-1)(t-5/7)(t-1/3)^2 y la descomposición de Halmos da C*(P_{Sigma,2},P_{Pi,2})≅C^4⊕M_2(C)⊕M_2(C), de dimensión 12.","status":"PRESENT_IN_SUCCESSOR_SOURCE"}
% HMT_RESULT_END:P24:residence
% HMT_RESULT_BEGIN:P25:residence
% HMT_ATOMIC_RESULT {"material_state":"INTEGRADO_EN_GRAFO_ACTIVO","proof_strength":"EXACTO_INTERNO","provenance":"RESULTADO_NUEVO","reinforcement_state":"DEPLOYED_WITHOUT_NEW_RESULT_ID","residence":"doble proyección operatorial, § Intervalo libre del bloque B_c, tras P24","result_id":"P25","statement":"Para todo nivel s, W_s(B_c)={Y=Y*:5I_s≼Y≼9I_s}; 7 es centro, 2 radio operatorial y 4 diámetro espectral.","status":"PRESENT_IN_SUCCESSOR_SOURCE"}
% HMT_RESULT_END:P25:residence
% HMT_RESULT_BEGIN:P26:residence
% HMT_ATOMIC_RESULT {"material_state":"INTEGRADO_EN_GRAFO_ACTIVO","proof_strength":"EXACTO_INTERNO","provenance":"FORMALIZACION_NUEVA","reinforcement_state":"DEPLOYED_WITHOUT_NEW_RESULT_ID","residence":"doble proyección, antes de TRI-001","result_id":"P26","statement":"Sus PVM diagonales conmutan y admiten el refinamiento común P_{a,e}=1_{R_m=a}1_{Q_m=e}; si los lectores respetan truncaciones, el refinamiento pasa a una PVM cilíndrica común.","status":"PRESENT_IN_SUCCESSOR_SOURCE"}
% HMT_RESULT_END:P26:residence
% HMT_RESULT_BEGIN:P27:residence
% HMT_ATOMIC_RESULT {"material_state":"INTEGRADO_EN_GRAFO_ACTIVO","proof_strength":"EXACTO_INTERNO","provenance":"FORMALIZACION_NUEVA","reinforcement_state":"DEPLOYED_WITHOUT_NEW_RESULT_ID","residence":"doble proyección operatorial, § Límite de dilatación MUB, tras P26","result_id":"P27","statement":"No admiten una PVM conjunta aguda: sus proyectores tienen solapamiento 1/3 y no conmutan; una realización conjunta requiere ruido, POVM o un ambiente mayor.","status":"PRESENT_IN_SUCCESSOR_SOURCE"}
% HMT_RESULT_END:P27:residence
% HMT_RESULT_BEGIN:P28:residence
% HMT_ATOMIC_RESULT {"material_state":"INTEGRADO_EN_GRAFO_ACTIVO","proof_strength":"EXACTO_INTERNO","provenance":"RESULTADO_RECUPERADO","reinforcement_state":"DEPLOYED_WITHOUT_NEW_RESULT_ID","residence":"frontera del capítulo de continuo operatorial","result_id":"P28","statement":"El teorema importado afirma que la MPS posee límite continuo si y sólo si E_A es infinitamente divisible, equivalentemente E_A=P exp(L), P^2=P y PLP=PL; una potencia E_A^r basta para el límite continuo grueso.","status":"PRESENT_IN_SUCCESSOR_SOURCE"}
% HMT_RESULT_END:P28:residence
% HMT_RESULT_BEGIN:P29:residence
% HMT_ATOMIC_RESULT {"material_state":"INTEGRADO_EN_GRAFO_ACTIVO","proof_strength":"EXACTO_INTERNO","provenance":"FORMALIZACION_NUEVA","reinforcement_state":"DEPLOYED_WITHOUT_NEW_RESULT_ID","residence":"medición/expectativa tracial","result_id":"P29","statement":"E_m es UCP, satisface E_m(a⊗I_9)=a y posee Kraus W_a psi=(1/3)psi⊗e_a y una dilatación cuyo ancilla registra el dígito promediado; en la diagonal, la retracción de rama kappa_{m,a}(f)(x)=f(x,a) es multiplicativa y no equivale a promediar. Esa memoria no se identifica aún con TPK.","status":"PRESENT_IN_SUCCESSOR_SOURCE"}
% HMT_RESULT_END:P29:residence
% HMT_RESULT_BEGIN:P30:residence
% HMT_ATOMIC_RESULT {"material_state":"INTEGRADO_EN_GRAFO_ACTIVO","proof_strength":"EXACTO_INTERNO","provenance":"FORMALIZACION_NUEVA","reinforcement_state":"DEPLOYED_WITHOUT_NEW_RESULT_ID","residence":"APP matricial","result_id":"P30","statement":"El doble centrado inserta una señal APP en un QMS.","status":"PRESENT_IN_SUCCESSOR_SOURCE"}
% HMT_RESULT_END:P30:residence
% HMT_RESULT_BEGIN:P31:residence
% HMT_ATOMIC_RESULT {"material_state":"INTEGRADO_EN_GRAFO_ACTIVO","proof_strength":"PROGRAMA_ABIERTO","provenance":"ARQUITECTURA_AUTORAL_PREEXISTENTE","reinforcement_state":"DEPLOYED_WITHOUT_NEW_RESULT_ID","residence":"perspectivas operatoriales","result_id":"P31","statement":"El programa asigna a cada tipo c un cono de lectores positivos A(c), con transformaciones naturales por niveles; requiere funtor de conos, objeto dualizante, evaluaciones, coevaluaciones y compatibilidad de la composición TPK. La profundidad de separación es el menor m donde los sistemas mínimo y máximo difieren; el nivel matricial de detección es el menor s que exhibe esa diferencia; y el defecto de memoria es la mínima información de ruta necesaria para levantar un punto del sistema máximo al mínimo.","status":"PRESENT_IN_SUCCESSOR_SOURCE"}
% HMT_RESULT_END:P31:residence
% HMT_RESULT_BEGIN:P32:residence
% HMT_ATOMIC_RESULT {"material_state":"INTEGRADO_EN_GRAFO_ACTIVO","proof_strength":"PROGRAMA_ABIERTO","provenance":"ARQUITECTURA_AUTORAL_PREEXISTENTE","reinforcement_state":"DEPLOYED_WITHOUT_NEW_RESULT_ID","residence":"perspectivas, § Universalidad genealógicamente fiel, tras UNIIMP-001","result_id":"P32","statement":"La versión HMT exige un funtor entre categorías de rutas que preserve composición, retorno y memoria, además de los datos del marco target--context importado.","status":"PRESENT_IN_SUCCESSOR_SOURCE"}
% HMT_RESULT_END:P32:residence
% HMT_RESULT_BEGIN:QPM-001:residence
% HMT_ATOMIC_RESULT {"material_state":"INTEGRADO_EN_GRAFO_ACTIVO","proof_strength":"EXACTO_INTERNO","provenance":"FORMALIZACION_NUEVA","reinforcement_state":"DEPLOYED_WITHOUT_NEW_RESULT_ID","residence":"apertura del corredor matricial, tras P09","result_id":"QPM-001","statement":"Una QPM es un QMS cuyas celdas son proyecciones; es conmutativa sólo si todas esas proyecciones conmutan conjuntamente. La definición no se obtiene dilatando cada fila por separado.","status":"PRESENT_IN_SUCCESSOR_SOURCE"}
% HMT_RESULT_END:QPM-001:residence
% HMT_RESULT_BEGIN:QLS-001:residence
% HMT_ATOMIC_RESULT {"material_state":"INTEGRADO_EN_GRAFO_ACTIVO","proof_strength":"EXACTO_INTERNO","provenance":"FORMALIZACION_NUEVA","reinforcement_state":"DEPLOYED_WITHOUT_NEW_RESULT_ID","residence":"apertura del corredor matricial, antes de QLSI-001","result_id":"QLS-001","statement":"Un QLS es una cuadrícula de vectores unitarios v_ij en C^n cuyas filas y columnas son bases ortonormales; las proyecciones P_ij=|v_ij><v_ij| forman un QMS proyectivo. Un QLS fácil se obtiene de un cuadrado latino clásico y una base ortonormal fija.","status":"PRESENT_IN_SUCCESSOR_SOURCE"}
% HMT_RESULT_END:QLS-001:residence
% HMT_RESULT_BEGIN:QLSI-001:residence
% HMT_ATOMIC_RESULT {"material_state":"INTEGRADO_EN_GRAFO_ACTIVO","proof_strength":"EXACTO_INTERNO","provenance":"RESULTADO_RECUPERADO","reinforcement_state":"DEPLOYED_WITHOUT_NEW_RESULT_ID","residence":"apertura del corredor matricial, tras P11","result_id":"QLSI-001","statement":"La envolvente matricial convexa de los QLS contiene estrictamente más que la clase semiclásica.","status":"PRESENT_IN_SUCCESSOR_SOURCE"}
% HMT_RESULT_END:QLSI-001:residence
% HMT_RESULT_BEGIN:QMSI-001:residence
% HMT_ATOMIC_RESULT {"material_state":"INTEGRADO_EN_GRAFO_ACTIVO","proof_strength":"EXACTO_INTERNO","provenance":"RESULTADO_RECUPERADO","reinforcement_state":"DEPLOYED_WITHOUT_NEW_RESULT_ID","residence":"apertura del corredor matricial, tras P10","result_id":"QMSI-001","statement":"Un QMS es semiclásico si y sólo si admite una dilatación simultánea a una QPM conmutativa mediante una única compresión común.","status":"PRESENT_IN_SUCCESSOR_SOURCE"}
% HMT_RESULT_END:QMSI-001:residence
% HMT_RESULT_BEGIN:QMSI-002:residence
% HMT_ATOMIC_RESULT {"material_state":"INTEGRADO_EN_GRAFO_ACTIVO","proof_strength":"EXACTO_INTERNO","provenance":"RESULTADO_RECUPERADO","reinforcement_state":"DEPLOYED_WITHOUT_NEW_RESULT_ID","residence":"apertura matricial, § Existencia de QMS no semiclásicos, tras QMSI-001","result_id":"QMSI-002","statement":"Para n≥3 y s≥2 existen QMS no semiclásicos.","status":"PRESENT_IN_SUCCESSOR_SOURCE"}
% HMT_RESULT_END:QMSI-002:residence
% HMT_RESULT_BEGIN:QMSI-003:residence
% HMT_ATOMIC_RESULT {"material_state":"INTEGRADO_EN_GRAFO_ACTIVO","proof_strength":"EXACTO_INTERNO","provenance":"RESULTADO_RECUPERADO","reinforcement_state":"DEPLOYED_WITHOUT_NEW_RESULT_ID","residence":"apertura del corredor matricial, tras QMSI-002","result_id":"QMSI-003","statement":"Para todo n≥3 existen QMS sobre M_2(C) que no admiten dilatación a ninguna QPM.","status":"PRESENT_IN_SUCCESSOR_SOURCE"}
% HMT_RESULT_END:QMSI-003:residence
% HMT_RESULT_BEGIN:QBR-001:residence
% HMT_ATOMIC_RESULT {"material_state":"INTEGRADO_EN_GRAFO_ACTIVO","proof_strength":"EXACTO_INTERNO","provenance":"RESULTADO_RECUPERADO","reinforcement_state":"DEPLOYED_WITHOUT_NEW_RESULT_ID","residence":"`md/04ab`, después de la realización hilbertiana finita","result_id":"QBR-001","statement":"Los observables se representan por operadores autoadjuntos y sus PVM; para un suceso cilíndrico P y estado normalizado Psi, la probabilidad es ||P Psi||^2.","status":"PRESENT_IN_SUCCESSOR_SOURCE"}
% HMT_RESULT_END:QBR-001:residence
% HMT_RESULT_BEGIN:QDY-001:residence
% HMT_ATOMIC_RESULT {"material_state":"INTEGRADO_EN_GRAFO_ACTIVO","proof_strength":"EXACTO_INTERNO","provenance":"RESULTADO_RECUPERADO","reinforcement_state":"DEPLOYED_WITHOUT_NEW_RESULT_ID","residence":"`md/04ab`, dinámica cuántica","result_id":"QDY-001","statement":"Un grupo unitario fuertemente continuo posee un generador autoadjunto H y satisface la ecuación de Schrödinger en su dominio.","status":"PRESENT_IN_SUCCESSOR_SOURCE"}
% HMT_RESULT_END:QDY-001:residence
% HMT_RESULT_BEGIN:QKR-001:residence
% HMT_ATOMIC_RESULT {"material_state":"INTEGRADO_EN_GRAFO_ACTIVO","proof_strength":"EXACTO_INTERNO","provenance":"RESULTADO_RECUPERADO","reinforcement_state":"DEPLOYED_WITHOUT_NEW_RESULT_ID","residence":"`md/01a`, instrumentos y canales","result_id":"QKR-001","statement":"Una familia de operadores de Kraus con suma K_a^*K_a=I define un canal completamente positivo y traza-preservante; cada resultado conserva su efecto e instrumento tipados.","status":"PRESENT_IN_SUCCESSOR_SOURCE"}
% HMT_RESULT_END:QKR-001:residence
% HMT_RESULT_BEGIN:QMO-001:residence
% HMT_ATOMIC_RESULT {"material_state":"INTEGRADO_EN_GRAFO_ACTIVO","proof_strength":"EXACTO_INTERNO","provenance":"RESULTADO_RECUPERADO","reinforcement_state":"DEPLOYED_WITHOUT_NEW_RESULT_ID","residence":"`md/04ab`, composición de registros","result_id":"QMO-001","statement":"Una realización monoidal fuerte transporta la composición independiente de registros al producto tensorial, con asociatividad, unidad y coherencias declaradas.","status":"PRESENT_IN_SUCCESSOR_SOURCE"}
% HMT_RESULT_END:QMO-001:residence
% HMT_RESULT_BEGIN:DIC-001:residence
% HMT_ATOMIC_RESULT {"material_state":"INTEGRADO_EN_GRAFO_ACTIVO","proof_strength":"EXACTO_INTERNO","provenance":"FORMALIZACION_NUEVA","reinforcement_state":"DEPLOYED_WITHOUT_NEW_RESULT_ID","residence":"apertura del corredor matricial","result_id":"DIC-001","statement":"El diccionario conserva exactamente: tamaño exterior n↔cuadrícula publicada; tamaño interior s↔carta operatorial; profundidad m↔genealogía; POVM↔familia de lectores positivos; PVM↔contexto espectral; UCP↔lector matricial; compresión↔sombra; dilatación↔antecedente operatorio; semiclasicidad↔POVM primaria común reordenada; no semiclasicidad↔inexistencia de una POVM primaria común que reproduzca todas las celdas; sistema de operadores↔espacio de lectores; rango matricial↔sombras a todos los s; límite inductivo↔cierre de refinamientos operatorios; y límite inverso↔sección infinita compatible. La composición genuinamente bidimensional es la interpretación HMT propuesta, no la definición importada.","status":"PRESENT_IN_SUCCESSOR_SOURCE"}
% HMT_RESULT_END:DIC-001:residence
% HMT_RESULT_BEGIN:QMC-001:residence
% HMT_ATOMIC_RESULT {"material_state":"INTEGRADO_EN_GRAFO_ACTIVO","proof_strength":"EXACTO_INTERNO","provenance":"RESULTADO_RECUPERADO","reinforcement_state":"DEPLOYED_WITHOUT_NEW_RESULT_ID","residence":"capítulo de cuadrados cuánticos, antes de las instancias \\((9,3)\\)","result_id":"QMC-001","statement":"A es un QMS (n,s), admite la descomposición semiclásica A=sum_t P_{sigma_t}⊗E_t y toda dilatación Naimark común E_t=V*Q_tV induce un QPM conmutativo simultáneo.","status":"PRESENT_IN_SUCCESSOR_SOURCE"}
% HMT_RESULT_END:QMC-001:residence
% HMT_RESULT_BEGIN:QDT-001:residence
% HMT_ATOMIC_RESULT {"material_state":"INTEGRADO_EN_GRAFO_ACTIVO","proof_strength":"EXACTO_INTERNO","provenance":"FORMALIZACION_NUEVA","reinforcement_state":"DEPLOYED_WITHOUT_NEW_RESULT_ID","residence":"geometría qutrit, antes de la fibración \\(12\\to4\\times3\\)","result_id":"QDT-001","statement":"M_3(C)=CI⊕direct_sum_{a∈P^1(F_3)}D_a^0 ortogonalmente para Hilbert--Schmidt y dim=9=1+4·2.","status":"PRESENT_IN_SUCCESSOR_SOURCE"}
% HMT_RESULT_END:QDT-001:residence
% HMT_RESULT_BEGIN:ENV-001:residence
% HMT_ATOMIC_RESULT {"material_state":"INTEGRADO_EN_GRAFO_ACTIVO","proof_strength":"EXACTO_INTERNO","provenance":"FORMALIZACION_NUEVA","reinforcement_state":"DEPLOYED_WITHOUT_NEW_RESULT_ID","residence":"dilatación cilíndrica, después de P16","result_id":"ENV-001","statement":"Sumar una representación fiel pi_f a una dilatación conserva todas las compresiones finitas y hace fiel la representación ambiental; esto no prueba todavía fidelidad dinámica TPK.","status":"PRESENT_IN_SUCCESSOR_SOURCE"}
% HMT_RESULT_END:ENV-001:residence
% HMT_RESULT_BEGIN:RNG-001:residence
% HMT_ATOMIC_RESULT {"material_state":"INTEGRADO_EN_GRAFO_ACTIVO","proof_strength":"EXACTO_INTERNO","provenance":"RESULTADO_NUEVO","reinforcement_state":"DEPLOYED_WITHOUT_NEW_RESULT_ID","residence":"rango matricial APP, después de P24","result_id":"RNG-001","statement":"(X,Y) pertenece a W_s(APP_2) exactamente cuando existen cuatro efectos escalares E_ab y dos mapas CP Psi_{5/7},Psi_{1/3}:M_2→M_s que satisfacen las ecuaciones completas de normalización y coordenadas.","status":"PRESENT_IN_SUCCESSOR_SOURCE"}
% HMT_RESULT_END:RNG-001:residence
% HMT_RESULT_BEGIN:DYN-001:residence
% HMT_ATOMIC_RESULT {"material_state":"INTEGRADO_EN_GRAFO_ACTIVO","proof_strength":"EXACTO_INTERNO","provenance":"FORMALIZACION_NUEVA","reinforcement_state":"DEPLOYED_WITHOUT_NEW_RESULT_ID","residence":"continuo operatorial, después de Stinespring","result_id":"DYN-001","statement":"La compatibilidad UCP a toda profundidad no implica ancilla finito uniforme; la fidelidad dinámica exige J_{m+1}T_m^TPK=tilde T_m J_m y naturalidad de las compresiones.","status":"PRESENT_IN_SUCCESSOR_SOURCE"}
% HMT_RESULT_END:DYN-001:residence
% HMT_RESULT_BEGIN:TEN-001:residence
% HMT_ATOMIC_RESULT {"material_state":"INTEGRADO_EN_GRAFO_ACTIVO","proof_strength":"PROGRAMA_ABIERTO","provenance":"ARQUITECTURA_AUTORAL_PREEXISTENTE","reinforcement_state":"DEPLOYED_WITHOUT_NEW_RESULT_ID","residence":"perspectivas del continuo tensorial","result_id":"TEN-001","statement":"El programa exige declarar tensor local, grupo de simetría, régimen exacto/aproximado/asintótico y estabilidad de rango; propone una torre coherente de raíces E_{kℓ}^ℓ=E_k equipada con registros TPK.","status":"PRESENT_IN_SUCCESSOR_SOURCE"}
% HMT_RESULT_END:TEN-001:residence
% HMT_RESULT_BEGIN:CHP-001:residence
% HMT_ATOMIC_RESULT {"material_state":"INTEGRADO_EN_GRAFO_ACTIVO","proof_strength":"PROGRAMA_ABIERTO","provenance":"FORMALIZACION_NUEVA","reinforcement_state":"DEPLOYED_WITHOUT_NEW_RESULT_ID","residence":"sistemas cónicos, antes de universalidad","result_id":"CHP-001","statement":"Las condiciones sum_j J(Phi_ij)=J(id) y sum_i J(Phi_ij)=J(id), con J(Phi_ij) positivo, linealizan la normalización mágica; la fidelidad de hoja y ruta es una obligación adicional.","status":"PRESENT_IN_SUCCESSOR_SOURCE"}
% HMT_RESULT_END:CHP-001:residence
% HMT_RESULT_BEGIN:SDP-001:residence
% HMT_ATOMIC_RESULT {"material_state":"INTEGRADO_EN_GRAFO_ACTIVO","proof_strength":"EXACTO_INTERNO","provenance":"FORMALIZACION_NUEVA","reinforcement_state":"DEPLOYED_WITHOUT_NEW_RESULT_ID","residence":"cierre del capítulo qutrit","result_id":"SDP-001","statement":"Se buscan Q_sigma positivos con sum_sigma Q_sigma=I y A_ij=sum_{sigma: sigma(i)=j}Q_sigma; factibilidad certifica semiclasicidad e inviabilidad requiere testigo dual.","status":"PRESENT_IN_SUCCESSOR_SOURCE"}
% HMT_RESULT_END:SDP-001:residence
% HMT_RESULT_BEGIN:EQV-001:residence
% HMT_ATOMIC_RESULT {"material_state":"INTEGRADO_EN_GRAFO_ACTIVO","proof_strength":"EXACTO_INTERNO","provenance":"FORMALIZACION_NUEVA","reinforcement_state":"DEPLOYED_WITHOUT_NEW_RESULT_ID","residence":"geometría simpléctica y lectores, después del calibre \\(3/\\pi\\)","result_id":"EQV-001","statement":"El calibre preliminar falla un test antiunitario y no es aún entrelazador; la promoción exige Xi(g·ell)=U_g Xi(ell)U_g* en todas las fibras, formulada como naturalidad con origen y destino cuando g es una flecha.","status":"PRESENT_IN_SUCCESSOR_SOURCE"}
% HMT_RESULT_END:EQV-001:residence
% HMT_RESULT_BEGIN:QMT-001:residence
% HMT_ATOMIC_RESULT {"material_state":"INTEGRADO_EN_GRAFO_ACTIVO","proof_strength":"PROGRAMA_ABIERTO","provenance":"FORMALIZACION_NUEVA","reinforcement_state":"DEPLOYED_WITHOUT_NEW_RESULT_ID","residence":"cierre del capítulo de cuadrados cuánticos APP","result_id":"QMT-001","statement":"La tabla conserva tres filas distintas: A^Sigma (9,3), que cierra hoja suma, positividad, normalización y semiclasicidad pero no hoja producto, cociente ni rutas; A^(9) (9,3), que cierra tres contextos, no conmutatividad y dilatación común pero no cuarto contexto, incidencia Hensel, gnomon ni memoria TPK; y B (12,3), que cierra el atlas 12→4×3 y su dilatación común pero no equivarianza, dinámica ni monodromía.","status":"PRESENT_IN_SUCCESSOR_SOURCE"}
% HMT_RESULT_END:QMT-001:residence
% HMT_RESULT_BEGIN:TRN-001:residence
% HMT_ATOMIC_RESULT {"material_state":"INTEGRADO_EN_GRAFO_ACTIVO","proof_strength":"PROGRAMA_ABIERTO","provenance":"ARQUITECTURA_AUTORAL_PREEXISTENTE","reinforcement_state":"DEPLOYED_WITHOUT_NEW_RESULT_ID","residence":"cierre metrológico del corredor matricial","result_id":"TRN-001","statement":"Conserva ocho filas completas: (1) monodromía 9→1→retorno del observable con estado ambiental distinto→clasificar dilataciones que preservan holonomía; (2) pi,e,phi,alpha→cuatro lectores compatibles de una genealogía→existencia y minimalidad de un sistema UCP común; (3) doble proyección→dos marginales de un antecedente cilíndrico→antecedente no conmutativo y separación conjunta de fibras; (4) B_c y salto 4→intervalo matricial libre [5I,9I]→extremos y transporte de la jerarquía 4;2;6;54; (5) Planck y CT108→punto autodual y lectores conjugados→realización como simetría de un sistema cónico con memoria; (6) ley de masas→familia positiva sobre fibras y rutas→rango matricial, extremos y realizaciones por sector; (7) CKM/CP→transporte de rango tres y fase no trivial→canal covariante que preserve espectro y orientación; (8) simetría excepcional→segundo codominio de la doble proyección→entrelazador entre atlas nonádico, código y representación.","status":"PRESENT_IN_SUCCESSOR_SOURCE"}
% HMT_RESULT_END:TRN-001:residence
% HMT_RESULT_BEGIN:MPSH-001:residence
% HMT_ATOMIC_RESULT {"material_state":"INTEGRADO_EN_GRAFO_ACTIVO","proof_strength":"PROGRAMA_ABIERTO","provenance":"ARQUITECTURA_AUTORAL_PREEXISTENTE","reinforcement_state":"DEPLOYED_WITHOUT_NEW_RESULT_ID","residence":"inmediatamente después de P28","result_id":"MPSH-001","statement":"Aplicar el teorema MPS a HMT exige construir una celda local uniforme, tensor A, canal E_A, forma irreducible, divisibilidad infinita propia o de una potencia y transporte al límite de hoja, acarreo y monodromía.","status":"PRESENT_IN_SUCCESSOR_SOURCE"}
% HMT_RESULT_END:MPSH-001:residence
% HMT_RESULT_BEGIN:UNIIMP-001:residence
% HMT_ATOMIC_RESULT {"material_state":"INTEGRADO_EN_GRAFO_ACTIVO","proof_strength":"EXACTO_INTERNO","provenance":"RESULTADO_RECUPERADO","reinforcement_state":"DEPLOYED_WITHOUT_NEW_RESULT_ID","residence":"perspectivas operatoriales, antes de P32","result_id":"UNIIMP-001","statement":"Un simulador s:P⊗C→T⊗C es universal cuando existe la reducción laxa prescrita al simulador trivial; un monótono compatible cuyo supremo sobre la imagen funcional quede por debajo de algún target constituye un no-go.","status":"PRESENT_IN_SUCCESSOR_SOURCE"}
% HMT_RESULT_END:UNIIMP-001:residence
% HMT_RESULT_BEGIN:PRG-001:residence
% HMT_ATOMIC_RESULT {"material_state":"INTEGRADO_EN_GRAFO_ACTIVO","proof_strength":"PROGRAMA_ABIERTO","provenance":"ARQUITECTURA_AUTORAL_PREEXISTENTE","reinforcement_state":"DEPLOYED_WITHOUT_NEW_RESULT_ID","residence":"cierre del corredor matricial","result_id":"PRG-001","statement":"Los cierres son: sistema APP filtrado, QMS APP--TPK fiel, frontera semiclásica, canal de continuo, sistema cónico intrínseco y realización metrológica UCP conjunta; ninguno se promueve sin dominio, prueba/certificado, positividad, compatibilidad, falsadores y trazabilidad.","status":"PRESENT_IN_SUCCESSOR_SOURCE"}
% HMT_RESULT_END:PRG-001:residence
% HMT_RESULT_BEGIN:QCT-001:residence
% HMT_ATOMIC_RESULT {"material_state":"CONTROL_FOCAL_ACTIVO","proof_strength":"EXACTO_INTERNO","provenance":"CERTIFICADO_NUEVO","reinforcement_state":"DEPLOYED_WITHOUT_NEW_RESULT_ID","residence":"apéndice digital y gate del puente","result_id":"QCT-001","statement":"Registra exactamente dieciocho controles: PVM qutrit; idempotencia; MUB; no-go QLS; QMS aditivo (9,3); QMS de tres contextos; no conmutatividad; QMS atlas (12,3); descomposición semiclásica; cubos de órdenes 9 y 12; P^1(Z/9)→P^1(F3); truncación--compresión; lector 1/pi; calibre 3/pi; polinomio de Gram; intersecciones de Halmos; y C*-álgebra APP d=2.","status":"CONTROLLED"}
% HMT_RESULT_END:QCT-001:residence
% HMT_RESULT_BEGIN:PRV-001:residence
% HMT_ATOMIC_RESULT {"material_state":"CONTROL_FOCAL_ACTIVO","proof_strength":"EXACTO_INTERNO","provenance":"FORMALIZACION_NUEVA","reinforcement_state":"DEPLOYED_WITHOUT_NEW_RESULT_ID","residence":"registro bibliográfico del corredor matricial","result_id":"PRV-001","statement":"Conserva autores, firmas e identificadores exactamente: De las Cuevas--Drescher--Netzer, Quantum magic squares: dilations and their limitations—DOI 10.1063/5.0022344—QMS/QPM/dilatación; Magic squares: Latin, Semiclassical and Quantum, De las Cuevas--Netzer--Valentiner-Branth—DOI 10.1063/5.0127393—QLS/envolventes; Quantum information theory and free semialgebraic geometry, De las Cuevas--Netzer—arXiv:2102.04240—torres/Choi; Beyond Operator Systems, De les Coves--Mirte van der Eyden--Netzer—DOI 10.1016/j.jmaa.2025.130102—sistemas cónicos; Continuum limits of Matrix Product States, De las Cuevas--Schuch--Pérez-García--Cirac—DOI 10.1103/PhysRevB.98.174303—divisibilidad; Tensor decompositions on simplicial complexes with invariance, De las Cuevas--Hoogsteder Riera--Netzer—DOI 10.1016/j.jsc.2024.102299—simetría; Approximate tensor decompositions, De las Cuevas--Klingler--Netzer—DOI 10.1063/5.0033876—aproximación; A Framework for Universality in Physics, Computer Science, and Beyond, Gonda--Reinhart--Stengele--De les Coves—DOI 10.46298/compositionality-6-3—target--context. Separa literatura simpléctica vecina y fija De las Cuevas/De les Coves.","status":"CONTROLLED"}
% HMT_RESULT_END:PRV-001:residence
% HMT_RESULT_BEGIN:MOR-001:residence
% HMT_ATOMIC_RESULT {"material_state":"INTEGRADO_O_REMITIDO_SIN_PERDIDA","proof_strength":"EXACTO_INTERNO","provenance":"ARQUITECTURA_AUTORAL_PREEXISTENTE","reinforcement_state":"DEPLOYED_WITHOUT_NEW_RESULT_ID","residence":"apertura continuo y `md/00b`","result_id":"MOR-001","statement":"Una cifra es una coordenada finita publicada; un valor es la imagen arquimediana, espectral o física de una sección; un número HMT es la sección completa con profundidad, hoja, orientación, acarreo, ruta, frontera, memoria y supervivencia; una simetría HMT conserva la identidad enriquecida y transporta su memoria.","status":"PRESENT_IN_SUCCESSOR_SOURCE"}
% HMT_RESULT_END:MOR-001:residence
% HMT_RESULT_BEGIN:MOR-002:residence
% HMT_ATOMIC_RESULT {"material_state":"INTEGRADO_O_REMITIDO_SIN_PERDIDA","proof_strength":"EXACTO_INTERNO","provenance":"RESULTADO_RECUPERADO","reinforcement_state":"DEPLOYED_WITHOUT_NEW_RESULT_ID","residence":"APP","result_id":"MOR-002","statement":"suma/producto y memoria de fibras","status":"PRESENT_IN_SUCCESSOR_SOURCE"}
% HMT_RESULT_END:MOR-002:residence
% HMT_RESULT_BEGIN:MOR-003:residence
% HMT_ATOMIC_RESULT {"material_state":"INTEGRADO_O_REMITIDO_SIN_PERDIDA","proof_strength":"EXACTO_INTERNO","provenance":"ARQUITECTURA_AUTORAL_PREEXISTENTE","reinforcement_state":"DEPLOYED_WITHOUT_NEW_RESULT_ID","residence":"TRIT","result_id":"MOR-003","statement":"orientación temporal y acarreo","status":"PRESENT_IN_SUCCESSOR_SOURCE"}
% HMT_RESULT_END:MOR-003:residence
% HMT_RESULT_BEGIN:MOR-004:residence
% HMT_ATOMIC_RESULT {"material_state":"INTEGRADO_O_REMITIDO_SIN_PERDIDA","proof_strength":"EXACTO_INTERNO","provenance":"FORMALIZACION_NUEVA","reinforcement_state":"DEPLOYED_WITHOUT_NEW_RESULT_ID","residence":"TPK","result_id":"MOR-004","statement":"composición con memoria","status":"PRESENT_IN_SUCCESSOR_SOURCE"}
% HMT_RESULT_END:MOR-004:residence
% HMT_RESULT_BEGIN:MOR-005:residence
% HMT_ATOMIC_RESULT {"material_state":"INTEGRADO_O_REMITIDO_SIN_PERDIDA","proof_strength":"EXACTO_INTERNO","provenance":"RESULTADO_RECUPERADO","reinforcement_state":"DEPLOYED_WITHOUT_NEW_RESULT_ID","residence":"continuo, tras `hmt/00b`","result_id":"MOR-005","statement":"9^n 1_C mu -> delta_x","status":"PRESENT_IN_SUCCESSOR_SOURCE"}
% HMT_RESULT_END:MOR-005:residence
% HMT_RESULT_BEGIN:MOR-006:residence
% HMT_ATOMIC_RESULT {"material_state":"INTEGRADO_O_REMITIDO_SIN_PERDIDA","proof_strength":"EXACTO_INTERNO","provenance":"FORMALIZACION_NUEVA","reinforcement_state":"DEPLOYED_WITHOUT_NEW_RESULT_ID","residence":"bisagra continuo--cuántica, § Kronecker--cilindro--delta de Dirac, tras MOR-005","result_id":"MOR-006","statement":"proyectores condicionales y núcleo","status":"PRESENT_IN_SUCCESSOR_SOURCE"}
% HMT_RESULT_END:MOR-006:residence
% HMT_RESULT_BEGIN:MOR-007:residence
% HMT_ATOMIC_RESULT {"material_state":"INTEGRADO_O_REMITIDO_SIN_PERDIDA","proof_strength":"EXACTO_INTERNO","provenance":"RESULTADO_RECUPERADO","reinforcement_state":"DEPLOYED_WITHOUT_NEW_RESULT_ID","residence":"UHF--\\(II_1\\) y `md/04ab`","result_id":"MOR-007","statement":"La inclusión unital nonádica genera UHF(9^infty), canónicamente isomorfa a UHF(3^infty); la traza GNS conduce al factor hiperinfinito II_1. El cuaderno de 64 páginas conserva la prueba propietaria completa.","status":"PRESENT_IN_SUCCESSOR_SOURCE"}
% HMT_RESULT_END:MOR-007:residence
% HMT_RESULT_BEGIN:MOR-008:residence
% HMT_ATOMIC_RESULT {"material_state":"INTEGRADO_O_REMITIDO_SIN_PERDIDA","proof_strength":"EXACTO_INTERNO","provenance":"RESULTADO_RECUPERADO","reinforcement_state":"DEPLOYED_WITHOUT_NEW_RESULT_ID","residence":"`md/01a`, `04ab`","result_id":"MOR-008","statement":"Born--Lüders--Kraus","status":"PRESENT_IN_SUCCESSOR_SOURCE"}
% HMT_RESULT_END:MOR-008:residence
% HMT_RESULT_BEGIN:MOR-009:residence
% HMT_ATOMIC_RESULT {"material_state":"INTEGRADO_O_REMITIDO_SIN_PERDIDA","proof_strength":"INTERFAZ_FISICA","provenance":"RESULTADO_RECUPERADO","reinforcement_state":"DEPLOYED_WITHOUT_NEW_RESULT_ID","residence":"`md/04ab`, `00b`","result_id":"MOR-009","statement":"Q:P_TPK→Hilb_{R,J}; la extensión a flechas inversas se afirma sólo en el subdominio reversible donde el transporte unitario fuente teoremática la demuestra","status":"PRESENT_IN_SUCCESSOR_SOURCE"}
% HMT_RESULT_END:MOR-009:residence
% HMT_RESULT_BEGIN:MOR-010:residence
% HMT_ATOMIC_RESULT {"material_state":"INTEGRADO_O_REMITIDO_SIN_PERDIDA","proof_strength":"PROGRAMA_ABIERTO","provenance":"FORMALIZACION_NUEVA","reinforcement_state":"DEPLOYED_WITHOUT_NEW_RESULT_ID","residence":"después de `md/00b`, antes de la ley dimensional/Planck","result_id":"MOR-010","statement":"R_MD^{disc,+} es el candidato de producto con blanco Hilbert general y sólo queda construido tras demostrar identidad y composición componente a componente; U_reg existe sólo en el subdominio unitario; R_MD^{disc,or} existe únicamente cuando espectro e Hilbert son invertibles, las relaciones localizadas se respetan y la dimensional toma valores en Gr(M_D^+); positiva y orientada son lectoras hermanas; no se presupone ni exige una factorización horizontal entre ellas. El monoide de coeficientes satisface (n,a,b,lambda)(n',a',b',lambda')=(n+n',a+a',b+b',lambda lambda') y sólo después se aplica la completación de Grothendieck.","status":"PRESENT_IN_SUCCESSOR_SOURCE"}
% HMT_RESULT_END:MOR-010:residence
% HMT_RESULT_BEGIN:MOR-011:residence
% HMT_ATOMIC_RESULT {"material_state":"INTEGRADO_O_REMITIDO_SIN_PERDIDA","proof_strength":"EXACTO_INTERNO","provenance":"RESULTADO_RECUPERADO","reinforcement_state":"DEPLOYED_WITHOUT_NEW_RESULT_ID","residence":"`md/11m`","result_id":"MOR-011","statement":"C U(gamma) C^-1=U(gamma^-1)","status":"PRESENT_IN_SUCCESSOR_SOURCE"}
% HMT_RESULT_END:MOR-011:residence
% HMT_RESULT_BEGIN:MOR-012:residence
% HMT_ATOMIC_RESULT {"material_state":"INTEGRADO_O_REMITIDO_SIN_PERDIDA","proof_strength":"INTERFAZ_FISICA","provenance":"ARQUITECTURA_AUTORAL_PREEXISTENTE","reinforcement_state":"DEPLOYED_WITHOUT_NEW_RESULT_ID","residence":"`md/11m`","result_id":"MOR-012","statement":"retardo/adelanto y borde","status":"PRESENT_IN_SUCCESSOR_SOURCE"}
% HMT_RESULT_END:MOR-012:residence
% HMT_RESULT_BEGIN:MOR-013:residence
% HMT_ATOMIC_RESULT {"material_state":"INTEGRADO_O_REMITIDO_SIN_PERDIDA","proof_strength":"EXACTO_INTERNO","provenance":"RESULTADO_RECUPERADO","reinforcement_state":"DEPLOYED_WITHOUT_NEW_RESULT_ID","residence":"`md/04c`","result_id":"MOR-013","statement":"N_sp(uP_++vP_-)=uv","status":"PRESENT_IN_SUCCESSOR_SOURCE"}
% HMT_RESULT_END:MOR-013:residence
% HMT_RESULT_BEGIN:MOR-014:residence
% HMT_ATOMIC_RESULT {"material_state":"INTEGRADO_O_REMITIDO_SIN_PERDIDA","proof_strength":"INTERFAZ_FISICA","provenance":"RESULTADO_RECUPERADO","reinforcement_state":"DEPLOYED_WITHOUT_NEW_RESULT_ID","residence":"`md/04c`, `06ka`","result_id":"MOR-014","statement":"producto positivo/interior","status":"PRESENT_IN_SUCCESSOR_SOURCE"}
% HMT_RESULT_END:MOR-014:residence
% HMT_RESULT_BEGIN:MOR-015:residence
% HMT_ATOMIC_RESULT {"material_state":"INTEGRADO_O_REMITIDO_SIN_PERDIDA","proof_strength":"PROGRAMA_ABIERTO","provenance":"RESULTADO_NUEVO","reinforcement_state":"DEPLOYED_WITHOUT_NEW_RESULT_ID","residence":"`md/04ac` y persistencia","result_id":"MOR-015","statement":"gamma_+ es fibra de Hopf y gamma_- se proyecta con grado dos; sus clases (1,1) y (1,-1) dan las familias de Villarceau. La conjugación compleja total preserva cada familia no orientada y revierte su orientación: no intercambia ambas familias.","status":"PRESENT_IN_SUCCESSOR_SOURCE"}
% HMT_RESULT_END:MOR-015:residence
% HMT_RESULT_BEGIN:MOR-016:residence
% HMT_ATOMIC_RESULT {"material_state":"INTEGRADO_O_REMITIDO_SIN_PERDIDA","proof_strength":"PROGRAMA_ABIERTO","provenance":"RESULTADO_RECUPERADO","reinforcement_state":"DEPLOYED_WITHOUT_NEW_RESULT_ID","residence":"/Users/ruben/Documents/HMT2/EDICION_AUTONOMA_HMT_MD_2026-08-09_REV_3_EN_TRABAJO/01_FUENTE_REVISADA/manuscrito/sections/md/21_persistencia_ruta_mobius_familias_actualizada.tex","result_id":"MOR-016","statement":"El doble retorno 2pi/4pi y el fibrado real de línea L_chi con chi(1)=-1 son exactos; un círculo aislado no es Möbius. Ascender el signo C_M a la holonomía del lazo Villarceau requiere un único entrelazador físico.","status":"PRESENT_IN_SUCCESSOR_SOURCE"}
% HMT_RESULT_END:MOR-016:residence
% HMT_RESULT_BEGIN:MOR-017:residence
% HMT_ATOMIC_RESULT {"material_state":"INTEGRADO_O_REMITIDO_SIN_PERDIDA","proof_strength":"EXACTO_INTERNO","provenance":"FORMALIZACION_NUEVA","reinforcement_state":"DEPLOYED_WITHOUT_NEW_RESULT_ID","residence":"`md/04fa`","result_id":"MOR-017","statement":"En el escenario qutrit declarado, las POVM completas producen una distribución no señalizante y la cota local es I_3≤2. Para el estado máximamente entrelazado, I_3=4(3+2sqrt(3))/9=2.872934…; para |Psi_gamma>=(|00>+gamma|11>+|22>)/sqrt(2+gamma^2), gamma=(sqrt(11)-sqrt(3))/2, se alcanza I_3=1+sqrt(11/3)=2.914854…. Por tanto, máxima entropía no equivale a máxima violación. Derivar el estado y los instrumentos íntegramente desde el registro HMT sigue siendo una obligación separada.","status":"PRESENT_IN_SUCCESSOR_SOURCE"}
% HMT_RESULT_END:MOR-017:residence
% HMT_RESULT_BEGIN:MOR-018:residence
% HMT_ATOMIC_RESULT {"material_state":"INTEGRADO_O_REMITIDO_SIN_PERDIDA","proof_strength":"INTERFAZ_FISICA","provenance":"RESULTADO_RECUPERADO","reinforcement_state":"DEPLOYED_WITHOUT_NEW_RESULT_ID","residence":"`md/10_constitutivas_si`","result_id":"MOR-018","statement":"mu,epsilon,Z,c","status":"PRESENT_IN_SUCCESSOR_SOURCE"}
% HMT_RESULT_END:MOR-018:residence
% HMT_RESULT_BEGIN:MOR-019:residence
% HMT_ATOMIC_RESULT {"material_state":"INTEGRADO_O_REMITIDO_SIN_PERDIDA","proof_strength":"INTERFAZ_FISICA","provenance":"RESULTADO_RECUPERADO","reinforcement_state":"DEPLOYED_WITHOUT_NEW_RESULT_ID","residence":"Planck y paquete de masas","result_id":"MOR-019","statement":"Sigma_H y 10^-34","status":"PRESENT_IN_SUCCESSOR_SOURCE"}
% HMT_RESULT_END:MOR-019:residence
% HMT_RESULT_BEGIN:MOR-020:residence
% HMT_ATOMIC_RESULT {"material_state":"INTEGRADO_O_REMITIDO_SIN_PERDIDA","proof_strength":"INTERFAZ_FISICA","provenance":"RESULTADO_RECUPERADO","reinforcement_state":"DEPLOYED_WITHOUT_NEW_RESULT_ID","residence":"`md/05aa`, `05f`, `06m`","result_id":"MOR-020","statement":"operadores K_D/K_Omega y torres","status":"PRESENT_IN_SUCCESSOR_SOURCE"}
% HMT_RESULT_END:MOR-020:residence
% HMT_RESULT_BEGIN:MOR-021:residence
% HMT_ATOMIC_RESULT {"material_state":"INTEGRADO_O_REMITIDO_SIN_PERDIDA","proof_strength":"INTERFAZ_FISICA","provenance":"RESULTADO_RECUPERADO","reinforcement_state":"DEPLOYED_WITHOUT_NEW_RESULT_ID","residence":"`md/06k`, `06ka`, `06m`","result_id":"MOR-021","statement":"basal->beta (80,54,6)","status":"PRESENT_IN_SUCCESSOR_SOURCE"}
% HMT_RESULT_END:MOR-021:residence
% HMT_RESULT_BEGIN:MOR-022:residence
% HMT_ATOMIC_RESULT {"material_state":"INTEGRADO_O_REMITIDO_SIN_PERDIDA","proof_strength":"EXACTO_INTERNO","provenance":"RESULTADO_RECUPERADO","reinforcement_state":"DEPLOYED_WITHOUT_NEW_RESULT_ID","residence":"`md/03d`, con remisiones a `06k` y `06ka`","result_id":"MOR-022","statement":"Con T_108=108t_0, omega_108=2pi/T_108, R_108=c/omega_108 y m_108=hbar omega_108/c^2, se cumple hbar/(m_108 c)=R_108 y h/(m_108 c)=2pi R_108=C_108.","status":"PRESENT_IN_SUCCESSOR_SOURCE"}
% HMT_RESULT_END:MOR-022:residence
% HMT_RESULT_BEGIN:MOR-023:residence
% HMT_ATOMIC_RESULT {"material_state":"INTEGRADO_O_REMITIDO_SIN_PERDIDA","proof_strength":"INTERFAZ_FISICA","provenance":"RESULTADO_RECUPERADO","reinforcement_state":"DEPLOYED_WITHOUT_NEW_RESULT_ID","residence":"nuevo teorema producto + `md/10c`","result_id":"MOR-023","statement":"holonomía, coframe, conexión","status":"PRESENT_IN_SUCCESSOR_SOURCE"}
% HMT_RESULT_END:MOR-023:residence
% HMT_RESULT_BEGIN:MOR-024:residence
% HMT_ATOMIC_RESULT {"material_state":"INTEGRADO_O_REMITIDO_SIN_PERDIDA","proof_strength":"INTERFAZ_FISICA","provenance":"RESULTADO_RECUPERADO","reinforcement_state":"DEPLOYED_WITHOUT_NEW_RESULT_ID","residence":"`md/11`, `11c`, `11p`","result_id":"MOR-024","statement":"A=(omega,e), F=(R,T)","status":"PRESENT_IN_SUCCESSOR_SOURCE"}
% HMT_RESULT_END:MOR-024:residence
% HMT_RESULT_BEGIN:MOR-025:residence
% HMT_ATOMIC_RESULT {"material_state":"INTEGRADO_O_REMITIDO_SIN_PERDIDA","proof_strength":"INTERFAZ_FISICA","provenance":"RESULTADO_RECUPERADO","reinforcement_state":"DEPLOYED_WITHOUT_NEW_RESULT_ID","residence":"`md/11m`","result_id":"MOR-025","statement":"Si b es el lector de base, todo observable constante en las fibras de b factoriza de manera única a través de b. Una semilla s_0 se transporta consistentemente alrededor de un lazo exactamente cuando s_0 pertenece a Fix(Hol). El dato de fibra asciende al estado TPK completo sólo si la realización empleada es inyectiva.","status":"PRESENT_IN_SUCCESSOR_SOURCE"}
% HMT_RESULT_END:MOR-025:residence
% HMT_RESULT_BEGIN:MOR-026:residence
% HMT_ATOMIC_RESULT {"material_state":"INTEGRADO_O_REMITIDO_SIN_PERDIDA","proof_strength":"EXACTO_INTERNO","provenance":"RESULTADO_RECUPERADO","reinforcement_state":"DEPLOYED_WITHOUT_NEW_RESULT_ID","residence":"Euler/09b/14f","result_id":"MOR-026","statement":"J^2=-I y z->iz","status":"PRESENT_IN_SUCCESSOR_SOURCE"}
% HMT_RESULT_END:MOR-026:residence
% HMT_RESULT_BEGIN:MOR-027:residence
% HMT_ATOMIC_RESULT {"material_state":"INTEGRADO_O_REMITIDO_SIN_PERDIDA","proof_strength":"EXACTO_INTERNO","provenance":"FORMALIZACION_NUEVA","reinforcement_state":"DEPLOYED_WITHOUT_NEW_RESULT_ID","residence":"/Users/ruben/Documents/HMT2/EDICION_AUTONOMA_HMT_MD_2026-08-09_REV_3_EN_TRABAJO/01_FUENTE_REVISADA/manuscrito/sections/hmt/14d_ontologia_numeros_rev11.tex","result_id":"MOR-027","statement":"valor regularizado de zeta(-1)","status":"PRESENT_IN_SUCCESSOR_SOURCE"}
% HMT_RESULT_END:MOR-027:residence
% HMT_RESULT_BEGIN:MOR-028:residence
% HMT_ATOMIC_RESULT {"material_state":"INTEGRADO_O_REMITIDO_SIN_PERDIDA","proof_strength":"EXACTO_INTERNO","provenance":"ARQUITECTURA_AUTORAL_PREEXISTENTE","reinforcement_state":"DEPLOYED_WITHOUT_NEW_RESULT_ID","residence":"`hmt/03f`, `04_cubo_emrh`","result_id":"MOR-028","statement":"doble evaluación del antecedente","status":"PRESENT_IN_SUCCESSOR_SOURCE"}
% HMT_RESULT_END:MOR-028:residence
% HMT_RESULT_BEGIN:MOR-029:residence
% HMT_ATOMIC_RESULT {"material_state":"INTEGRADO_O_REMITIDO_SIN_PERDIDA","proof_strength":"EXACTO_INTERNO","provenance":"RESULTADO_RECUPERADO","reinforcement_state":"DEPLOYED_WITHOUT_NEW_RESULT_ID","residence":"`hmt/06aa--06f`","result_id":"MOR-029","statement":"Para todo K, la misma regla particiona exhaustivamente las 729 prolongaciones, selecciona una única superviviente y conmuta con las restricciones; los cilindros compatibles anidados determinan una sección coinductiva única.","status":"PRESENT_IN_SUCCESSOR_SOURCE"}
% HMT_RESULT_END:MOR-029:residence
% HMT_RESULT_BEGIN:MOR-030:residence
% HMT_ATOMIC_RESULT {"material_state":"INTEGRADO_O_REMITIDO_SIN_PERDIDA","proof_strength":"EXACTO_INTERNO","provenance":"FORMALIZACION_NUEVA","reinforcement_state":"DEPLOYED_WITHOUT_NEW_RESULT_ID","residence":"junto al teorema producto y síntesis final","result_id":"MOR-030","statement":"Los falsadores exactos son: continuo, pérdida de masa unitaria o de convergencia débil; Hilbert, ausencia de PVM compatible o dinámica no unitaria; morfismo discreto, fallo de composición; conjugación, C^2≠I, CJC^{-1}≠-J o CHC^{-1}≠H; absorbedor, falta de operador de onda o de condiciones causales; constitutivas, inversión que no intercambia mu/epsilon o cambia c sin extensión; masa, selección por PDG o promoción condicionada; electrón, imposibilidad de obtener carga, espín o estabilidad del mismo objeto; Planck, otra solución positiva autodual; Cartan, dependencia de subdivisión o confusión con Delta_4; Bell, señalización o instrumentos escogidos tras el blanco; filtración nonádica, consulta de prefijos objetivo o pérdida de una región de pi; Euler/zeta, regularización tratada como suma ordinaria; cosmología, incapacidad de reproducir simultáneamente expansión, lentes y estructura.","status":"PRESENT_IN_SUCCESSOR_SOURCE"}
% HMT_RESULT_END:MOR-030:residence
% HMT_RESULT_BEGIN:PLS-001:residence
% HMT_ATOMIC_RESULT {"material_state":"INTEGRADO_O_REMITIDO_SIN_PERDIDA","proof_strength":"EXACTO_INTERNO","provenance":"RESULTADO_RECUPERADO","reinforcement_state":"DEPLOYED_WITHOUT_NEW_RESULT_ID","residence":"`md/03d`, subsección CT108--Compton","result_id":"PLS-001","statement":"Para G(theta)=theta c_int^5 t_0^2/hbar_int, la clausura circular selecciona de manera única theta_*=(54/pi)^2 y hace coincidir el radio CT108 con las escalas de Compton, gravitatoria y de Planck en la convención declarada. Se distinguen r_g=Gm/c^2 y r_s=2r_g.","status":"PRESENT_IN_SUCCESSOR_SOURCE"}
% HMT_RESULT_END:PLS-001:residence
% HMT_RESULT_BEGIN:MAD-001:residence
% HMT_ATOMIC_RESULT {"material_state":"INTEGRADO_O_REMITIDO_SIN_PERDIDA","proof_strength":"EXACTO_INTERNO","provenance":"RESULTADO_RECUPERADO","reinforcement_state":"DEPLOYED_WITHOUT_NEW_RESULT_ID","residence":"`md/03d`, recta autodual de masas","result_id":"MAD-001","statement":"Con P_±=(I±j)/2, Z(m)=(m_P/m)P_++(m/m_P)P_- satisface N_sp(Z(m))=1; la conjugación implementa m↦m_P^2/m, el punto fijo positivo es m=m_P y x=log(m/m_P) es la rapidez de la recta.","status":"PRESENT_IN_SUCCESSOR_SOURCE"}
% HMT_RESULT_END:MAD-001:residence
% HMT_RESULT_BEGIN:AGR-001:residence
% HMT_ATOMIC_RESULT {"material_state":"INTEGRADO_O_REMITIDO_SIN_PERDIDA","proof_strength":"DEMOSTRADO_CON_ESTRUCTURA_DE_PARTIDA_EXPLICITA","provenance":"RESULTADO_RECUPERADO","reinforcement_state":"DEPLOYED_WITHOUT_NEW_RESULT_ID","residence":"`md/03d`, covariancia pre/retorno","result_id":"AGR-001","statement":"Las ramas pre y retorno transforman recíprocamente acción y acoplamiento gravitatorio, mantienen invariante su producto geométrico y conservan la longitud característica construida.","status":"PRESENT_IN_SUCCESSOR_SOURCE"}
% HMT_RESULT_END:AGR-001:residence
% HMT_RESULT_BEGIN:PLE-001:residence
% HMT_ATOMIC_RESULT {"material_state":"INTEGRADO_O_REMITIDO_SIN_PERDIDA","proof_strength":"INTERFAZ_FISICA","provenance":"RESULTADO_RECUPERADO","reinforcement_state":"DEPLOYED_WITHOUT_NEW_RESULT_ID","residence":"`md/03d`, `06k`, `06ka`","result_id":"PLE-001","statement":"El punto autodual de Planck fija el origen de la coordenada; la genealogía electrónica determina un desplazamiento posterior y la clausura escalar de rango uno. Este cierre discreto no demuestra por sí solo carga, espín, estabilidad, factor de forma, momento magnético ni un Hamiltoniano electromagnético ligado.","status":"PRESENT_IN_SUCCESSOR_SOURCE"}
% HMT_RESULT_END:PLE-001:residence
% HMT_RESULT_BEGIN:ANG-001:residence
% HMT_ATOMIC_RESULT {"material_state":"INTEGRADO_O_REMITIDO_SIN_PERDIDA","proof_strength":"INTERFAZ_FISICA","provenance":"ARQUITECTURA_AUTORAL_PREEXISTENTE","reinforcement_state":"DEPLOYED_WITHOUT_NEW_RESULT_ID","residence":"Hopf/Villarceau, después de MOR-016","result_id":"ANG-001","statement":"Para epsilon=+1, Spec(B_mem^+)={7,3}; resolver R+r=7ell y R-r=3ell da exactamente (R,r)=(5ell,2ell) y sin(beta)=2/5. La lectura como bordes de toro es interfaz declarada; el bloque total conserva por separado la razón 7:2.","status":"PRESENT_IN_SUCCESSOR_SOURCE"}
% HMT_RESULT_END:ANG-001:residence
% HMT_RESULT_BEGIN:BDR-001:residence
% HMT_ATOMIC_RESULT {"material_state":"INTEGRADO_O_REMITIDO_SIN_PERDIDA","proof_strength":"INTERFAZ_FISICA","provenance":"RESULTADO_RECUPERADO","reinforcement_state":"DEPLOYED_WITHOUT_NEW_RESULT_ID","residence":"constitutivas, con remisión desde bidireccionalidad","result_id":"BDR-001","statement":"La inversión conserva c_hat=exp(-E), invierte Z_hat=exp(B) e intercambia mu_hat y epsilon_hat; la realización vigente tiene velocidad común en ambos sentidos e impedancia orientada. Velocidades distintas exigen una extensión anisótropa adicional.","status":"PRESENT_IN_SUCCESSOR_SOURCE"}
% HMT_RESULT_END:BDR-001:residence
% HMT_RESULT_BEGIN:OCC-001:residence
% HMT_ATOMIC_RESULT {"material_state":"INTEGRADO_O_REMITIDO_SIN_PERDIDA","proof_strength":"EXACTO_INTERNO","provenance":"RESULTADO_RECUPERADO","reinforcement_state":"DEPLOYED_WITHOUT_NEW_RESULT_ID","residence":"Gibbs--KMS y ocupación","result_id":"OCC-001","statement":"Gobierna la identidad tipada de REV.3 F_{A,B}(t+i beta)=Tr(rho alpha_t(B)A), no la abreviación errónea del PDF de 68 pp; para 0<C<I, K_1=log((I-C)C^{-1}) y C=(I+e^{K_1})^{-1}, con tipos CAR/CCR separados.","status":"PRESENT_IN_SUCCESSOR_SOURCE"}
% HMT_RESULT_END:OCC-001:residence
% HMT_RESULT_BEGIN:BOL-001:residence
% HMT_ATOMIC_RESULT {"material_state":"INTEGRADO_O_REMITIDO_SIN_PERDIDA","proof_strength":"EXACTO_INTERNO","provenance":"RESULTADO_RECUPERADO","reinforcement_state":"DEPLOYED_WITHOUT_NEW_RESULT_ID","residence":"`md/04e`, `04g` y síntesis, con remisión CT108","result_id":"BOL-001","statement":"Boltzmann actúa como k_B:L_Theta→L_E y, para el reloj CT108, satisface k_B Theta_clk log(3)=hbar_int omega_0 con omega_0=2pi/(108t_0). La identidad fija el producto acción--información; insertar la coordenada SI declarada de k_B sólo evalúa Theta_clk y no constituye una derivación autónoma de ese número SI.","status":"PRESENT_IN_SUCCESSOR_SOURCE"}
% HMT_RESULT_END:BOL-001:residence
% HMT_RESULT_BEGIN:COS-001:residence
% HMT_ATOMIC_RESULT {"material_state":"INTEGRADO_O_REMITIDO_SIN_PERDIDA","proof_strength":"PROGRAMA_ABIERTO","provenance":"ARQUITECTURA_AUTORAL_PREEXISTENTE","reinforcement_state":"DEPLOYED_WITHOUT_NEW_RESULT_ID","residence":"perspectivas cosmológicas","result_id":"COS-001","statement":"La coexistencia de cartas permite un modelo estratificado, pero la aceleración sólo queda explicada tras derivar una ecuación efectiva con signo observado y límites de curvatura.","status":"PRESENT_IN_SUCCESSOR_SOURCE"}
% HMT_RESULT_END:COS-001:residence
% HMT_RESULT_BEGIN:GWV-001:residence
% HMT_ATOMIC_RESULT {"material_state":"INTEGRADO_O_REMITIDO_SIN_PERDIDA","proof_strength":"PROGRAMA_ABIERTO","provenance":"RESULTADO_RECUPERADO","reinforcement_state":"DEPLOYED_WITHOUT_NEW_RESULT_ID","residence":"gravedad sin gravitón y espectro","result_id":"GWV-001","statement":"Linealizar y reducir por calibre puede producir dos grados transversos de una onda gravitatoria colectiva; esto no obliga ni excluye un gravitón elemental, cuya existencia es una pregunta espectral adicional.","status":"PRESENT_IN_SUCCESSOR_SOURCE"}
% HMT_RESULT_END:GWV-001:residence
% HMT_RESULT_BEGIN:DRK-001:residence
% HMT_ATOMIC_RESULT {"material_state":"INTEGRADO_O_REMITIDO_SIN_PERDIDA","proof_strength":"PROGRAMA_ABIERTO","provenance":"ARQUITECTURA_AUTORAL_PREEXISTENTE","reinforcement_state":"DEPLOYED_WITHOUT_NEW_RESULT_ID","residence":"perspectivas cosmológicas","result_id":"DRK-001","statement":"Una alternativa estructural debe reproducir curvas de rotación, lentes, crecimiento, CMB y expansión con los mismos parámetros; Landauer nulo reversible no genera por sí solo un ciclo, que requiere dinámica de frontera acción--información.","status":"PRESENT_IN_SUCCESSOR_SOURCE"}
% HMT_RESULT_END:DRK-001:residence
% HMT_RESULT_BEGIN:CMK-001:residence
% HMT_ATOMIC_RESULT {"material_state":"INTEGRADO_O_REMITIDO_SIN_PERDIDA","proof_strength":"PROGRAMA_ABIERTO","provenance":"ARQUITECTURA_AUTORAL_PREEXISTENTE","reinforcement_state":"DEPLOYED_WITHOUT_NEW_RESULT_ID","residence":"cierre Euler/TPK","result_id":"CMK-001","statement":"El programa debe construir un diagrama de módulos y transformaciones naturales que tipen pares adjuntos o involutivos; una enumeración de operadores no constituye la coestructura.","status":"PRESENT_IN_SUCCESSOR_SOURCE"}
% HMT_RESULT_END:CMK-001:residence
% HMT_RESULT_BEGIN:MRC-001:residence
% HMT_ATOMIC_RESULT {"material_state":"INTEGRADO_O_REMITIDO_SIN_PERDIDA","proof_strength":"EXACTO_INTERNO","provenance":"CERTIFICADO_NUEVO","reinforcement_state":"DEPLOYED_WITHOUT_NEW_RESULT_ID","residence":"apéndice digital y gate de la monografía de realización","result_id":"MRC-001","statement":"Conserva diez familias: cilindros 0--1--infinito; espectros APP/memoria; identidad Perron--Fibonacci; inversión antiunitaria; producto semidirecto de Poincaré; constitutivas; norma partida y clausura material; Planck/autodualidad; CGLMP; e invariantes de filtración nonádica y masas.","status":"PRESENT_IN_SUCCESSOR_SOURCE"}
% HMT_RESULT_END:MRC-001:residence
% HMT_RESULT_BEGIN:PRG68-001:residence
% HMT_ATOMIC_RESULT {"material_state":"INTEGRADO_O_REMITIDO_SIN_PERDIDA","proof_strength":"PROGRAMA_ABIERTO","provenance":"ARQUITECTURA_AUTORAL_PREEXISTENTE","reinforcement_state":"DEPLOYED_WITHOUT_NEW_RESULT_ID","residence":"cierre del capítulo de realización física","result_id":"PRG68-001","statement":"Los seis cierres distintos son: límite suave Poincaré--Cartan independiente de subdivisión; un único operador C global que separe C, P y T; electrón ligado con carga, espín, estabilidad, factor de forma y momento magnético; mapa desde Delta_4 y la clausura electrónica hasta S_matter y la corriente de espín; estado e instrumentos HMT que deriven Bell/CGLMP sin usar el blanco; y selector general theta_G del transductor gravitatorio sin introducir la coordenada externa.","status":"PRESENT_IN_SUCCESSOR_SOURCE"}
% HMT_RESULT_END:PRG68-001:residence
% HMT_RESULT_BEGIN:GTR-001:residence
% HMT_ATOMIC_RESULT {"material_state":"INTEGRADO_O_REMITIDO_SIN_PERDIDA","proof_strength":"EXACTO_INTERNO","provenance":"RESULTADO_RECUPERADO","reinforcement_state":"DEPLOYED_WITHOUT_NEW_RESULT_ID","residence":"`md/11c`, `11n` y síntesis","result_id":"GTR-001","statement":"El mapa G_int(theta_G)=theta_G c_int^5 t_0^2/hbar_int tiene tipo dimensional exacto. En la clausura circular CT108 se especializa con theta_*=(54/pi)^2; una carta declarada evalúa 6.67430×10^-11. El selector general theta_G sigue siendo una obligación distinta y no se obtiene introduciendo esa coordenada externa.","status":"PRESENT_IN_SUCCESSOR_SOURCE"}
% HMT_RESULT_END:GTR-001:residence
% HMT_RESULT_BEGIN:AMB-001:residence
% HMT_ATOMIC_RESULT {"material_state":"CONSERVADO_Y_ATOMIZADO_EN_DELTA_SUCESOR","proof_strength":"DEMOSTRADO_CON_ESTRUCTURA_DE_PARTIDA_EXPLICITA","provenance":"RESULTADO_RECUPERADO","reinforcement_state":"DEPLOYED_WITHOUT_NEW_RESULT_ID","residence":"manuscrito/sections/hmt/05b_extensiones_dinamica_cociente.tex","result_id":"AMB-001","statement":"El cociente de modos del TPK induce la incidencia primitiva F_av; la medida de Parry del envolvente de memoria uno satisface mu_ar/mu_vol=phi_HMT^-2. El lector coinductivo marca una sola vez el retorno de frontera y produce pi_HMT/phi_HMT^2. Para R(x,y)=x/y^2 y J(x,y)=(y,x), las orientaciones conjugadas son pi_HMT/phi_HMT^2 y phi_HMT/pi_HMT^2; son funcionales coordinados y no intercambiables del mismo estado holonómico integral con memoria, sin factorización horizontal.","status":"PRESENT_IN_SUCCESSOR_SOURCE"}
% HMT_RESULT_END:AMB-001:residence
% HMT_RESULT_BEGIN:CHI-II-HDG-01:residence
% HMT_ATOMIC_RESULT {"material_state":"AUDITADO_PENDIENTE_INTEGRACION","proof_strength":"DEMOSTRADO_CON_ESTRUCTURA_DE_PARTIDA_EXPLICITA","provenance":"FORMALIZACION_NUEVA","reinforcement_state":"DEPLOYED_WITHOUT_NEW_RESULT_ID","residence":"manuscrito/sections/sucesora/02_estructura_discreta_continuo_20260817.tex:Hodge_macroobjetos","result_id":"CHI-II-HDG-01","statement":"La realización tipada R_Hdg:Xχ→Perf conserva la historia enriquecida del carácter χ y publica cada estado de Xχ como complejo perfecto, sin reconstruir el antecedente desde su sombra cohomológica.","status":"PRESENT_IN_SUCCESSOR_SOURCE"}
% HMT_RESULT_END:CHI-II-HDG-01:residence
% HMT_RESULT_BEGIN:CHI-II-HDG-02:residence
% HMT_ATOMIC_RESULT {"material_state":"AUDITADO_PENDIENTE_INTEGRACION","proof_strength":"DEMOSTRADO_CON_ESTRUCTURA_DE_PARTIDA_EXPLICITA","provenance":"FORMALIZACION_NUEVA","reinforcement_state":"DEPLOYED_WITHOUT_NEW_RESULT_ID","residence":"manuscrito/sections/sucesora/02_estructura_discreta_continuo_20260817.tex:Hodge_identidad_realizacion","result_id":"CHI-II-HDG-02","statement":"La identidad cl∘chloc∘R_Hdg=ev∞ conmuta exactamente en Xχ: la realización perfecta, su carácter de Chern localizado y la clase de ciclo publican la misma evaluación terminal.","status":"PRESENT_IN_SUCCESSOR_SOURCE"}
% HMT_RESULT_END:CHI-II-HDG-02:residence
% HMT_RESULT_BEGIN:CHI-II-HDG-03:residence
% HMT_ATOMIC_RESULT {"material_state":"INTEGRADO_EN_SUCESORA_PENDIENTE_DE_SELLO","proof_strength":"DEMOSTRADO_CON_ESTRUCTURA_DE_PARTIDA_EXPLICITA","provenance":"FORMALIZACION_NUEVA","reinforcement_state":"DEPLOYED_WITHOUT_NEW_RESULT_ID","residence":"manuscrito/sections/sucesora/02_estructura_discreta_continuo_20260817.tex:eq:continuo-hodge-completitud; manuscrito/sections/sucesora/04_hodge_bsd_fragment_20260817.tex:thm:hodge-suc-generacion","result_id":"CHI-II-HDG-03","statement":"Xχ cubre el dominio Hdg y, para toda clase α de ese dominio, existe ξ_α∈Xχ con β_α=chloc(R_Hdg(ξ_α)); por la identidad de realización, cl(β_α)=α.","status":"PRESENT_IN_SUCCESSOR_SOURCE"}
% HMT_RESULT_END:CHI-II-HDG-03:residence
% HMT_RESULT_BEGIN:CHI-IV-HDG-01:residence
% HMT_ATOMIC_RESULT {"material_state":"AUDITADO_PENDIENTE_INTEGRACION","proof_strength":"DEMOSTRADO_CON_ESTRUCTURA_DE_PARTIDA_EXPLICITA","provenance":"RESULTADO_NUEVO","reinforcement_state":"DEPLOYED_WITHOUT_NEW_RESULT_ID","residence":"manuscrito/sections/sucesora/04_cinco_realizaciones_fundamentales_20260817.tex:Hodge_APP_MV_coend","result_id":"CHI-IV-HDG-01","statement":"La celda APP–Mayer–Vietoris κ materializa suma, producto, diferencia y carry en Perf; el coend sobre historias compatibles totaliza las nueve fases coordinadas y conserva como terminal Cone(I−U_Γ9), con memoria visible.","status":"PRESENT_IN_SUCCESSOR_SOURCE"}
% HMT_RESULT_END:CHI-IV-HDG-01:residence
% HMT_RESULT_BEGIN:CHI-IV-HDG-C01:residence
% HMT_ATOMIC_RESULT {"material_state":"AUDITADO_PENDIENTE_INTEGRACION","proof_strength":"EXACTO_INTERNO","provenance":"RESULTADO_NUEVO","reinforcement_state":"DEPLOYED_WITHOUT_NEW_RESULT_ID","residence":"manuscrito/sections/sucesora/04_cinco_realizaciones_fundamentales_20260817.tex:Hodge_control_no_go_eliptico","result_id":"CHI-IV-HDG-C01","statement":"En una curva elíptica compleja E no existe un frame racional finito que cubra H²(E,Q(1)) y sea invariante bajo todas las traslaciones, ni un selector natural no basado de la clase de un punto; la obstrucción obliga a conservar la rigidificación y su descenso.","status":"PRESENT_IN_SUCCESSOR_SOURCE"}
% HMT_RESULT_END:CHI-IV-HDG-C01:residence
% HMT_RESULT_BEGIN:CHI-II-BSD-01:residence
% HMT_ATOMIC_RESULT {"material_state":"AUDITADO_PENDIENTE_INTEGRACION","proof_strength":"DEMOSTRADO_CON_ESTRUCTURA_DE_PARTIDA_EXPLICITA","provenance":"FORMALIZACION_NUEVA","reinforcement_state":"DEPLOYED_WITHOUT_NEW_RESULT_ID","residence":"manuscrito/sections/sucesora/02_estructura_discreta_continuo_20260817.tex:BSD_Gmn_AW","result_id":"CHI-II-BSD-01","statement":"El complejo G_mn de cadenas normalizadas del nervio de historias TPK supervivientes porta la diagonal Alexander–Whitney; ésta conserva la unidad genealógica y hace group-like cada vértice de estado completo.","status":"PRESENT_IN_SUCCESSOR_SOURCE"}
% HMT_RESULT_END:CHI-II-BSD-01:residence
% HMT_RESULT_BEGIN:CHI-II-BSD-02:residence
% HMT_ATOMIC_RESULT {"material_state":"AUDITADO_PENDIENTE_INTEGRACION","proof_strength":"DEMOSTRADO_CON_ESTRUCTURA_DE_PARTIDA_EXPLICITA","provenance":"RESULTADO_NUEVO","reinforcement_state":"DEPLOYED_WITHOUT_NEW_RESULT_ID","residence":"manuscrito/sections/sucesora/02_estructura_discreta_continuo_20260817.tex:BSD_Manin_producto_fibrado","result_id":"CHI-II-BSD-02","statement":"El blanco conservativo BSD es Manin×_A G: el producto fibrado acopla el complejo Manin basado con G sobre el aumento A y conserva simultáneamente la proyección modular y la historia profunda.","status":"PRESENT_IN_SUCCESSOR_SOURCE"}
% HMT_RESULT_END:CHI-II-BSD-02:residence
% HMT_RESULT_BEGIN:CHI-II-BSD-03:residence
% HMT_ATOMIC_RESULT {"material_state":"AUDITADO_PENDIENTE_INTEGRACION","proof_strength":"DEMOSTRADO_CON_ESTRUCTURA_DE_PARTIDA_EXPLICITA","provenance":"RESULTADO_NUEVO","reinforcement_state":"DEPLOYED_WITHOUT_NEW_RESULT_ID","residence":"manuscrito/sections/sucesora/02_estructura_discreta_continuo_20260817.tex:BSD_publicacion_acoplada_K_PT","result_id":"CHI-II-BSD-03","statement":"Las publicaciones Kato y Poitou–Tate se construyen sobre Manin×_A G y publican la misma unidad genealógica de G; la comparación conserva raíz y complemento sin factorizar por el mero aumento.","status":"PRESENT_IN_SUCCESSOR_SOURCE"}
% HMT_RESULT_END:CHI-II-BSD-03:residence
% HMT_RESULT_BEGIN:CHI-IV-BSD-01:residence
% HMT_ATOMIC_RESULT {"material_state":"AUDITADO_PENDIENTE_INTEGRACION","proof_strength":"DEMOSTRADO_CON_ESTRUCTURA_DE_PARTIDA_EXPLICITA","provenance":"RESULTADO_NUEVO","reinforcement_state":"DEPLOYED_WITHOUT_NEW_RESULT_ID","residence":"manuscrito/sections/sucesora/04_cinco_realizaciones_fundamentales_20260817.tex:BSD_raiz_dualidad_HFS","result_id":"CHI-IV-BSD-01","statement":"La reducción ortogonal y la dualidad chain-level fijan J_red; la orientación λ elimina el twist y fuerza p_root zK=zPT, mientras H_FS construye h* para contraer el complemento finita–singular sin alterar la raíz.","status":"PRESENT_IN_SUCCESSOR_SOURCE"}
% HMT_RESULT_END:CHI-IV-BSD-01:residence
% HMT_RESULT_BEGIN:CHI-IV-BSD-02:residence
% HMT_ATOMIC_RESULT {"material_state":"INTEGRADO_EN_SUCESORA_PENDIENTE_DE_SELLO","proof_strength":"DEMOSTRADO_CON_ESTRUCTURA_DE_PARTIDA_EXPLICITA","provenance":"FORMALIZACION_NUEVA","reinforcement_state":"DEPLOYED_WITHOUT_NEW_RESULT_ID","residence":"manuscrito/sections/sucesora/04_hodge_bsd_fragment_20260817.tex:thm:bsd-suc-bockstein,thm:bsd-suc-local-global,eq:bsd-suc-formula-final","result_id":"CHI-IV-BSD-02","statement":"El regulador derivado R_Boc^der coincide como elemento basado con el término inicial LT_T(S); unido a la igualdad raíz Kato/PT, la altura, Sha, Tamagawa, periodo y torsión ya tipados, publica la fórmula BSD final sin definir un lado por el otro.","status":"PRESENT_IN_SUCCESSOR_SOURCE"}
% HMT_RESULT_END:CHI-IV-BSD-02:residence
% HMT_RESULT_BEGIN:CHI-IV-BSD-C01:residence
% HMT_ATOMIC_RESULT {"material_state":"AUDITADO_PENDIENTE_INTEGRACION","proof_strength":"EXACTO_INTERNO","provenance":"RESULTADO_NUEVO","reinforcement_state":"DEPLOYED_WITHOUT_NEW_RESULT_ID","residence":"manuscrito/sections/sucesora/04_cinco_realizaciones_fundamentales_20260817.tex:BSD_control_twist4","result_id":"CHI-IV-BSD-C01","statement":"La ablación del pullback genealógico o la inserción del twist raíz 4 conserva sombras modulares parciales pero destruye la unidad basada: zK^(4)=4r+3e+b pasa módulo 3 y falla desde módulo 9, por lo que no puede satisfacer p_root zK=zPT.","status":"PRESENT_IN_SUCCESSOR_SOURCE"}
% HMT_RESULT_END:CHI-IV-BSD-C01:residence
% HMT_RESULT_BEGIN:CHI-II-RH-01:residence
% HMT_ATOMIC_RESULT {"material_state":"AUDITADO_PENDIENTE_INTEGRACION","proof_strength":"EXACTO_INTERNO","provenance":"FORMALIZACION_NUEVA","reinforcement_state":"DEPLOYED_WITHOUT_NEW_RESULT_ID","residence":"manuscrito/sections/sucesora/02_estructura_discreta_continuo_20260817.tex:RH_macroobjeto","result_id":"CHI-II-RH-01","statement":"La Parte II construye M_RH=(Xhat_infinito,Gamma9,Rhat=(R,K),Carr_R,S0,M,O9,Xi,P,D_R_cof): Gamma9 conserva la fase y eleva la profundidad; Carr_R compone como cociclo; S0=sum_{q=0}^{N-1}M^{*q}O9^*O9M^q+M^{*N}S0M^N conserva el terminal; y B_W=Xi^*(I-P)Xi publica el complemento de nivel/fibra sin borrar memoria.","status":"PRESENT_IN_SUCCESSOR_SOURCE"}
% HMT_RESULT_END:CHI-II-RH-01:residence
% HMT_RESULT_BEGIN:CHI-IV-RH-01:residence
% HMT_ATOMIC_RESULT {"material_state":"AUDITADO_PENDIENTE_INTEGRACION","proof_strength":"DEMOSTRADO_CON_ESTRUCTURA_DE_PARTIDA_EXPLICITA","provenance":"RESULTADO_RECUPERADO","reinforcement_state":"DEPLOYED_WITHOUT_NEW_RESULT_ID","residence":"manuscrito/sections/sucesora/04_cinco_realizaciones_fundamentales_20260817.tex:RH_teorema","result_id":"CHI-IV-RH-01","statement":"Con q=5/9, M9=Phi A8···A0=qV9 y la identidad finita de memoria dan B_W,R=L_R^*L_R≥0 al extinguirse su terminal; la realización común produce B_W=Xi^*(I-P)Xi≥0 y, sobre D_R_cof con los canales completos, el criterio reversible publica Re(rho)=1/2 para todo cero no trivial.","status":"PRESENT_IN_SUCCESSOR_SOURCE"}
% HMT_RESULT_END:CHI-IV-RH-01:residence
% HMT_RESULT_BEGIN:CHI-IV-RH-C01:residence
% HMT_ATOMIC_RESULT {"material_state":"AUDITADO_PENDIENTE_INTEGRACION","proof_strength":"EXACTO_INTERNO","provenance":"RESULTADO_RECUPERADO","reinforcement_state":"DEPLOYED_WITHOUT_NEW_RESULT_ID","residence":"manuscrito/sections/sucesora/04_cinco_realizaciones_fundamentales_20260817.tex:RH_control_ablacion","result_id":"CHI-IV-RH-C01","statement":"El control retira por separado ortogonalidad, antecedente común, canal primo/gamma/polar, cofinalidad, carry, memoria o terminal; cada ablación rompe una igualdad anterior al criterio y no puede presentarse como una variante del teorema.","status":"PRESENT_IN_SUCCESSOR_SOURCE"}
% HMT_RESULT_END:CHI-IV-RH-C01:residence
% HMT_RESULT_BEGIN:CHI-II-NS-01:residence
% HMT_ATOMIC_RESULT {"material_state":"AUDITADO_PENDIENTE_INTEGRACION","proof_strength":"EXACTO_INTERNO","provenance":"FORMALIZACION_NUEVA","reinforcement_state":"DEPLOYED_WITHOUT_NEW_RESULT_ID","residence":"manuscrito/sections/sucesora/02_estructura_discreta_continuo_20260817.tex:NS_macroobjeto","result_id":"CHI-II-NS-01","statement":"La Parte II construye M_NS,m=(Xi_NS,m,C_{m←n},M_m^+,M_m^-,D_m,H_m,E9,J9,P9,Q_micro,R_m,E_m,ell_m): C_{m←n}=P_mB(u_n,u_n)-B_m(P_m u_n,P_m u_n) conserva el carry; E9J9=I, P9=J9E9 y Q_micro=I-P9 tipan la fibra; E9||Xi_{m+1}||²+E9 ell_m=||Xi_m||² conserva el terminal y archiva memoria positiva.","status":"PRESENT_IN_SUCCESSOR_SOURCE"}
% HMT_RESULT_END:CHI-II-NS-01:residence
% HMT_RESULT_BEGIN:CHI-IV-NS-01:residence
% HMT_ATOMIC_RESULT {"material_state":"AUDITADO_PENDIENTE_INTEGRACION","proof_strength":"EXACTO_INTERNO","provenance":"RESULTADO_RECUPERADO","reinforcement_state":"DEPLOYED_WITHOUT_NEW_RESULT_ID","residence":"manuscrito/sections/sucesora/04_cinco_realizaciones_fundamentales_20260817.tex:NS_teorema","result_id":"CHI-IV-NS-01","statement":"El puerto del macroobjeto satisface [B_{s,m}+(1-eta)nu D_{s,m}]^-≤C_NS ell_m+b_{parcial,m}; al componerlo con la telescopía y la energía se obtiene sup_m||u_m||_{L∞H^s}²+sup_m||u_m||_{L²H^{s+1}}²≤C_T(u0,f,Xi0) para s>5/2, y se publican existencia global, unicidad y suavidad para t>0 en el dominio declarado.","status":"PRESENT_IN_SUCCESSOR_SOURCE"}
% HMT_RESULT_END:CHI-IV-NS-01:residence
% HMT_RESULT_BEGIN:CHI-IV-NS-C01:residence
% HMT_ATOMIC_RESULT {"material_state":"AUDITADO_PENDIENTE_INTEGRACION","proof_strength":"EXACTO_INTERNO","provenance":"RESULTADO_RECUPERADO","reinforcement_state":"DEPLOYED_WITHOUT_NEW_RESULT_ID","residence":"manuscrito/sections/sucesora/04_cinco_realizaciones_fundamentales_20260817.tex:NS_control_ablacion","result_id":"CHI-IV-NS-C01","statement":"El control elimina por separado retorno conjunto, carry de Galerkin, memoria positiva, Q_micro, terminal o uniformidad del puerto; una rama de norma 3^J, C_NS dependiente del corte o b_parcial,m no sumable impiden la telescopía y rechazan la publicación global.","status":"PRESENT_IN_SUCCESSOR_SOURCE"}
% HMT_RESULT_END:CHI-IV-NS-C01:residence
% HMT_RESULT_BEGIN:CHI-II-YM-01:residence
% HMT_ATOMIC_RESULT {"material_state":"AUDITADO_PENDIENTE_INTEGRACION","proof_strength":"EXACTO_INTERNO","provenance":"FORMALIZACION_NUEVA","reinforcement_state":"DEPLOYED_WITHOUT_NEW_RESULT_ID","residence":"manuscrito/sections/sucesora/02_estructura_discreta_continuo_20260817.tex:YM_macroobjeto","result_id":"CHI-II-YM-01","statement":"La Parte II construye M_YM,m=(Lambda_m,mu_m,T_m,J_m,Theta_{nu,m},K_t,D_m^YM,P_Omega_m,Q_m^phys): Lambda_m=a_m Z^4 y K_{a_{m+1}}^9=K_{a_m}; (T_{m+1})^9J_m=J_mT_m y Theta_{nu,m+1}J_m=J_mTheta_{nu,m} conservan nivel/fibra, carry de fusión y memoria; P_Omega_m es el terminal raíz.","status":"PRESENT_IN_SUCCESSOR_SOURCE"}
% HMT_RESULT_END:CHI-II-YM-01:residence
% HMT_RESULT_BEGIN:CHI-IV-YM-01:residence
% HMT_ATOMIC_RESULT {"material_state":"AUDITADO_PENDIENTE_INTEGRACION","proof_strength":"EXACTO_INTERNO","provenance":"RESULTADO_RECUPERADO","reinforcement_state":"DEPLOYED_WITHOUT_NEW_RESULT_ID","residence":"manuscrito/sections/sucesora/04_cinco_realizaciones_fundamentales_20260817.tex:YM_teorema","result_id":"CHI-IV-YM-01","statement":"Sobre el macroobjeto común, omega_m(overline{F(Theta_nu,m U)}F(U))=||B_nu,mF||²≥0 para nu=1,…,4; el generador D_m^YM satisface ker D_m^YM=C Omega_m y D_m^YM≥3kappa_av(I-P_Omega_m); q_{m+1}^9=q_m y -a_m^{-1}log q_m=3kappa_av publican una teoría local, vacío simple y brecha Delta_YM^HMT≥3kappa_av>0 en el sector colectivo gauge-invariante.","status":"PRESENT_IN_SUCCESSOR_SOURCE"}
% HMT_RESULT_END:CHI-IV-YM-01:residence
% HMT_RESULT_BEGIN:CHI-IV-YM-C01:residence
% HMT_ATOMIC_RESULT {"material_state":"AUDITADO_PENDIENTE_INTEGRACION","proof_strength":"EXACTO_INTERNO","provenance":"RESULTADO_RECUPERADO","reinforcement_state":"DEPLOYED_WITHOUT_NEW_RESULT_ID","residence":"manuscrito/sections/sucesora/04_cinco_realizaciones_fundamentales_20260817.tex:YM_control_ablacion","result_id":"CHI-IV-YM-C01","statement":"El control elimina por separado semigrupo común, medida compartida, localidad uniforme, entrelazamiento de refinamiento, carry, memoria, vacío terminal o escala física; cada ablación rompe una identidad previa a la brecha y no puede repararse usando Wilson como selector retrospectivo.","status":"PRESENT_IN_SUCCESSOR_SOURCE"}
% HMT_RESULT_END:CHI-IV-YM-C01:residence
% HMT_RESULT_BEGIN:CHI-I-TPK-01:residence
% HMT_ATOMIC_RESULT {"material_state":"INTEGRADO_EN_LIENZO_SUCESOR_NO_ACTIVADO","proof_strength":"EXACTO_INTERNO","provenance":"FORMALIZACION_NUEVA","reinforcement_state":"DEPLOYED_WITHOUT_NEW_RESULT_ID","residence":"manuscrito/sections/sucesora/01_nucleo_formal_hmt_20260817.tex:conexion_nonadica_conjunta","result_id":"CHI-I-TPK-01","statement":"La holonomía única Γ9=Hol_C9(∇_TPK) conserva el retorno de fase con evolución del estado; en una realización isométrica satisface U_Γ9 J=JT+η, J*η=0 e I−T*T=η*η, y la iteración anti-especular conserva la suma de memoria junto con el terminal finito positivo.","status":"PRESENT_IN_SUCCESSOR_SOURCE"}
% HMT_RESULT_END:CHI-I-TPK-01:residence
% HMT_RESULT_BEGIN:CHI-I-EMRH-01:residence
% HMT_ATOMIC_RESULT {"material_state":"INTEGRADO_EN_LIENZO_SUCESOR_NO_ACTIVADO","proof_strength":"EXACTO_INTERNO","provenance":"RESULTADO_RECUPERADO","reinforcement_state":"DEPLOYED_WITHOUT_NEW_RESULT_ID","residence":"manuscrito/sections/sucesora/01_nucleo_formal_hmt_20260817.tex:completacion_cubica_conservacion_proyectiva","result_id":"CHI-I-EMRH-01","statement":"Para C̄9^(d)=(E9⊔{⋆})^d, el estrato con r coordenadas de frontera tiene cardinal B_r(d)=binom(d,r)9^(d−r), la proyección que olvida la última coordenada satisface (π_d)_*μ_d=μ_{d−1} y la masa total permanece uno; en dimensión tres, 1000=729+243+27+1=729+271.","status":"PRESENT_IN_SUCCESSOR_SOURCE"}
% HMT_RESULT_END:CHI-I-EMRH-01:residence
% HMT_RESULT_BEGIN:CHI-I-FAV-01:residence
% HMT_ATOMIC_RESULT {"material_state":"INTEGRADO_EN_LIENZO_SUCESOR_NO_ACTIVADO","proof_strength":"DEMOSTRADO_CON_ESTRUCTURA_DE_PARTIDA_EXPLICITA","provenance":"FORMALIZACION_NUEVA","reinforcement_state":"DEPLOYED_WITHOUT_NEW_RESULT_ID","residence":"manuscrito/sections/sucesora/01_nucleo_formal_hmt_20260817.tex:cociente_areal_volumetrico","result_id":"CHI-I-FAV-01","statement":"El bloque primitivo (+1,−1,0) induce la órbita de modos (Σ,Π,Π); el lector tipado de hojas transporta su incidencia a F_av=[[0,1],[1,1]], y la repetición del calendario produce N9=3F_av, N27=9F_av y N108=36F_av.","status":"PRESENT_IN_SUCCESSOR_SOURCE"}
% HMT_RESULT_END:CHI-I-FAV-01:residence
% HMT_RESULT_BEGIN:CHI-III-VAC-01:residence
% HMT_ATOMIC_RESULT {"material_state":"AUDITADO_PENDIENTE_INTEGRACION","proof_strength":"EXACTO_INTERNO","provenance":"FORMALIZACION_NUEVA","reinforcement_state":"DEPLOYED_WITHOUT_NEW_RESULT_ID","residence":"manuscrito/sections/sucesora/03_genealogia_constantes_20260817.tex:descomposicion_Fitting_vacancias","result_id":"CHI-III-VAC-01","statement":"En T9=(Z/9Z)^2, M=[[3,4],[6,7]] satisface M²−M=3I; E=M³ es idempotente y N=M−E cumple N²=0 y EN=NE, de modo que T9=im(E)⊕ker(E) y la filtración de imágenes tiene cardinales 81→27→9→9.","status":"PRESENT_IN_SUCCESSOR_SOURCE"}
% HMT_RESULT_END:CHI-III-VAC-01:residence
% HMT_RESULT_BEGIN:CHI-III-MASS-01:residence
% HMT_ATOMIC_RESULT {"material_state":"AUDITADO_PENDIENTE_INTEGRACION","proof_strength":"DEMOSTRADO_CON_ESTRUCTURA_DE_PARTIDA_EXPLICITA","provenance":"FORMALIZACION_NUEVA","reinforcement_state":"DEPLOYED_WITHOUT_NEW_RESULT_ID","residence":"manuscrito/sections/sucesora/03_genealogia_constantes_20260817.tex:dominio_completo_ley_general_masas","result_id":"CHI-III-MASS-01","statement":"El dominio interno vinculante de la Ley General de Masas comprende 81 celdas, 56 familias de ruta, 13 multisecciones, 324 rutas orientadas y un sector nulo separado; cada extensión superviviente publica ruta, registro, firma, carácter y operador de fibra antes de toda carta metrológica. Las antiguas ventanas de diez o doce masas son sólo testigos históricos.","status":"PRESENT_IN_SUCCESSOR_SOURCE"}
% HMT_RESULT_END:CHI-III-MASS-01:residence
% HMT_RESULT_BEGIN:CHI-III-ELECTRON-01:residence
% HMT_ATOMIC_RESULT {"material_state":"AUDITADO_PENDIENTE_INTEGRACION","proof_strength":"DEMOSTRADO_CON_ESTRUCTURA_DE_PARTIDA_EXPLICITA","provenance":"FORMALIZACION_NUEVA","reinforcement_state":"DEPLOYED_WITHOUT_NEW_RESULT_ID","residence":"manuscrito/sections/sucesora/03_genealogia_constantes_20260817.tex:electron_central_two_way","result_id":"CHI-III-ELECTRON-01","statement":"La fibra electrónica ocupa el único centro (5,5) con ruta Gamma_e=(5,5,++); sus lectores independientes coinciden en K_D(Gamma_e)=K_Omega(Gamma_e)=(80,54,6), fijan kappa_e=80−54alpha_HMT+6Delta_4 y, después del basal y de la recta de acción, publican m_e^beta=0.5109989506900008086082571648… MeV. El basal permanece como etapa, no como salida rectora.","status":"PRESENT_IN_SUCCESSOR_SOURCE"}
% HMT_RESULT_END:CHI-III-ELECTRON-01:residence
% HMT_RESULT_BEGIN:CHI-III-BOLTZ-01:residence
% HMT_ATOMIC_RESULT {"material_state":"AUDITADO_PENDIENTE_INTEGRACION","proof_strength":"EXACTO_INTERNO","provenance":"FORMALIZACION_NUEVA","reinforcement_state":"DEPLOYED_WITHOUT_NEW_RESULT_ID","residence":"manuscrito/sections/sucesora/03_genealogia_constantes_20260817.tex:boltzmann_accion_informacion","result_id":"CHI-III-BOLTZ-01","statement":"La constante de Boltzmann se tipa como el isomorfismo positivo k_B:L_Theta→L_E; una decisión TRIT uniforme aporta I_TRIT=log 3 y la periodicidad de orden 108 fija k_B Theta_clk log 3=hbar_int 2pi/(108 t0). La construcción determina internamente la energía por unidad de información; la carta kelvin–joule publica después k_B=1.380649×10^−23 J K^−1.","status":"PRESENT_IN_SUCCESSOR_SOURCE"}
% HMT_RESULT_END:CHI-III-BOLTZ-01:residence
% HMT_RESULT_BEGIN:CHI-III-PLANCK108-01:residence
% HMT_ATOMIC_RESULT {"material_state":"AUDITADO_PENDIENTE_INTEGRACION","proof_strength":"DEMOSTRADO_CON_ESTRUCTURA_DE_PARTIDA_EXPLICITA","provenance":"FORMALIZACION_NUEVA","reinforcement_state":"DEPLOYED_WITHOUT_NEW_RESULT_ID","residence":"manuscrito/sections/sucesora/03_genealogia_constantes_20260817.tex:planck_periodicidad_108","result_id":"CHI-III-PLANCK108-01","statement":"Las cartas de acción hbar_pre y hbar_ret admiten una realización circular de período T_108=108t0, con omega_108=2pi/T_108, R_108=c_int/omega_108 y C_108=108 ell0; para m_108,a=hbar_a omega_108/c_int² se cumplen barlambda_C=R_108 y lambda_C=C_108. La condición de radio gravitatorio reducido selecciona theta_108=(54/pi)² y coordina las unidades internas de Planck en ambas cartas.","status":"PRESENT_IN_SUCCESSOR_SOURCE"}
% HMT_RESULT_END:CHI-III-PLANCK108-01:residence
% HMT_RESULT_BEGIN:CHI-VI-GRAV-ALG-01:residence
% HMT_ATOMIC_RESULT {"material_state":"INTEGRADO_EN_LIENZO_SUCESOR_NO_ACTIVADO","proof_strength":"EXACTO_INTERNO","provenance":"FORMALIZACION_NUEVA","reinforcement_state":"DEPLOYED_WITHOUT_NEW_RESULT_ID","residence":"manuscrito/sections/sucesora/06_mecanica_dimensional_20260817.tex:sector_tensorial_icosaedrico_y_proyeccion_TT","result_id":"CHI-VI-GRAV-ALG-01","statement":"La representación icosaédrica de A5 sobre V12 posee el proyector ortogonal P5 de rango cinco; la igualdad de caracteres identifica V5=im(P5) con Sym²_0(R³) mediante una isometría A5-equivariante única hasta signo, y el proyector transversal sin traza Lambda(n) tiene rango dos. La cadena exacta es 12→5→5→2 y no se confunde con una realización gravitatoria dinámica.","status":"PRESENT_IN_SUCCESSOR_SOURCE"}
% HMT_RESULT_END:CHI-VI-GRAV-ALG-01:residence
% HMT_RESULT_BEGIN:AMP-I-RECTA-01:residence
% HMT_ATOMIC_RESULT {"material_state":"AMPLIFICACION_AUTORIZADA_2026_08_17","proof_strength":"EXACTO_INTERNO","provenance":"FORMALIZACION_NUEVA","reinforcement_state":"DEPLOYED_WITHOUT_NEW_RESULT_ID","residence":"manuscrito/sections/ampliacion_20260817/00_recta_horizonte_completacion.tex","result_id":"AMP-I-RECTA-01","statement":"La sucesión de cartas [0,10]→[0,1000]→[0,10^m] construye un sistema compatible de marcas, intervalos y fronteras cuya evaluación final vive en R; el quinto postulado euclídeo se trata como antecedente histórico de cierre geométrico y la completación se formula mediante mapas y límites, no como metáfora.","status":"PRESENT_IN_AMPLIFIED_SOURCE"}
% HMT_RESULT_END:AMP-I-RECTA-01:residence
% HMT_RESULT_BEGIN:AMP-I-UNIDAD-01:residence
% HMT_ATOMIC_RESULT {"material_state":"AMPLIFICACION_AUTORIZADA_2026_08_17","proof_strength":"EXACTO_INTERNO","provenance":"FORMALIZACION_NUEVA","reinforcement_state":"DEPLOYED_WITHOUT_NEW_RESULT_ID","residence":"manuscrito/sections/ampliacion_20260817/00_recta_horizonte_completacion.tex","result_id":"AMP-I-UNIDAD-01","statement":"La unidad normalizada se conserva bajo refinamiento, elevación dimensional y doble publicación; en dimensión d la estratificación binomial de (E9⊔{★})^d suma 10^d y los push-forward proyectivos conservan masa uno.","status":"PRESENT_IN_AMPLIFIED_SOURCE"}
% HMT_RESULT_END:AMP-I-UNIDAD-01:residence
% HMT_RESULT_BEGIN:AMP-II-EXC-01:residence
% HMT_ATOMIC_RESULT {"material_state":"AMPLIFICACION_AUTORIZADA_2026_08_17","proof_strength":"DEMOSTRADO_CON_ESTRUCTURA_DE_PARTIDA_EXPLICITA","provenance":"RESULTADO_RECUPERADO","reinforcement_state":"DEPLOYED_WITHOUT_NEW_RESULT_ID","residence":"manuscrito/sections/sucesora/02_estructura_discreta_continuo_20260817.tex; manuscrito/sections/sucesora/06_mecanica_dimensional_20260817.tex","result_id":"AMP-II-EXC-01","statement":"La rama excepcional se bifurca tipadamente: C_W→W12→W24 por diseños, y C_W→N(A2^12)→Λ24→V_Λ→V^natural→Monster→Moonshine por retículo y orbifold; no existe una serialización Witt→Golay→A2^12.","status":"PRESENT_IN_AMPLIFIED_SOURCE"}
% HMT_RESULT_END:AMP-II-EXC-01:residence
% HMT_RESULT_BEGIN:AMP-III-TIPOS-01:residence
% HMT_ATOMIC_RESULT {"material_state":"AMPLIFICACION_AUTORIZADA_2026_08_17","proof_strength":"DEMOSTRADO_CON_ESTRUCTURA_DE_PARTIDA_EXPLICITA","provenance":"RESULTADO_RECUPERADO","reinforcement_state":"DEPLOYED_WITHOUT_NEW_RESULT_ID","residence":"manuscrito/sections/sucesora/03_genealogia_constantes_20260817.tex; manuscrito/sections/sucesora/06_mecanica_dimensional_20260817.tex","result_id":"AMP-III-TIPOS-01","statement":"G se expone en tres capas —transductor exacto, selector circular de orden 108 y criterio holonómico general—; Boltzmann deriva internamente k_B Θ_clk log3 y sólo después abre la carta SI; raíz de -1 permanece en Parte II y -1/12 en Parte III.","status":"PRESENT_IN_AMPLIFIED_SOURCE"}
% HMT_RESULT_END:AMP-III-TIPOS-01:residence
% HMT_RESULT_BEGIN:AMP-IV-REAL-01:residence
% HMT_ATOMIC_RESULT {"material_state":"AMPLIFICACION_AUTORIZADA_2026_08_17","proof_strength":"DEMOSTRADO_CON_ESTRUCTURA_DE_PARTIDA_EXPLICITA","provenance":"FORMALIZACION_NUEVA","reinforcement_state":"DEPLOYED_WITHOUT_NEW_RESULT_ID","residence":"manuscrito/sections/ampliacion_20260817/04_principio_realizacion_y_supersesion.tex","result_id":"AMP-IV-REAL-01","statement":"Una familia F_N:S_N^surv→C_N compatible con truncación, orientación, transporte, acarreo, unidad y terminal induce F_∞ por la propiedad universal del límite inverso y transporta conjuntamente observable y memoria.","status":"PRESENT_IN_AMPLIFIED_SOURCE"}
% HMT_RESULT_END:AMP-IV-REAL-01:residence
% HMT_RESULT_BEGIN:AMP-IV-RHNS-01:residence
% HMT_ATOMIC_RESULT {"material_state":"AMPLIFICACION_AUTORIZADA_2026_08_17","proof_strength":"DEMOSTRADO_CON_ESTRUCTURA_DE_PARTIDA_EXPLICITA","provenance":"RESULTADO_RECUPERADO","reinforcement_state":"DEPLOYED_WITHOUT_NEW_RESULT_ID","residence":"manuscrito/sections/ampliacion_20260817/04a_refuerzo_riemann_weil.tex; manuscrito/sections/ampliacion_20260817/04b_refuerzo_navier_stokes.tex","result_id":"AMP-IV-RHNS-01","statement":"RH incorpora el residual finito 729=81+648 y R_m,R=0 antes del criterio; NS incorpora cota conjunta, ∂_t u en L2_loc H^{s-1}, pegado por unicidad y presión normalizada, sin sustituir sus demostraciones vigentes.","status":"PRESENT_IN_AMPLIFIED_SOURCE"}
% HMT_RESULT_END:AMP-IV-RHNS-01:residence
% HMT_RESULT_BEGIN:AMP-IV-YM-01:residence
% HMT_ATOMIC_RESULT {"material_state":"AMPLIFICACION_AUTORIZADA_2026_08_17","proof_strength":"DEMOSTRADO_CON_ESTRUCTURA_DE_PARTIDA_EXPLICITA","provenance":"RESULTADO_RECUPERADO","reinforcement_state":"DEPLOYED_WITHOUT_NEW_RESULT_ID","residence":"manuscrito/sections/ampliacion_20260817/04c_refuerzo_yang_mills.tex","result_id":"AMP-IV-YM-01","statement":"La familia proyectiva común produce cuatro formas OS, un Hamiltoniano común, acción fuerte de E(4), funcionales de Schwinger, publicación en H^{-s}_loc(R4), s>2, lectores clover y F², vacío simple y brecha exacta 3κ_av.","status":"PRESENT_IN_AMPLIFIED_SOURCE"}
% HMT_RESULT_END:AMP-IV-YM-01:residence
% HMT_RESULT_BEGIN:AMP-IV-HDG-01:residence
% HMT_ATOMIC_RESULT {"material_state":"AMPLIFICACION_AUTORIZADA_2026_08_17","proof_strength":"DEMOSTRADO_CON_ESTRUCTURA_DE_PARTIDA_EXPLICITA","provenance":"FORMALIZACION_NUEVA","reinforcement_state":"DEPLOYED_WITHOUT_NEW_RESULT_ID","residence":"manuscrito/sections/ampliacion_20260817/04d_hodge_completitud_desplegada.tex","result_id":"AMP-IV-HDG-01","statement":"La completitud de Xχ(X,p) se despliega como teorema coinductivo de realizaciones finitas compatibles; produce primero ξ_α y sólo después β_α=ch_p^loc(R_Hdg(ξ_α)), preservando el cuadrado de publicación.","status":"PRESENT_IN_AMPLIFIED_SOURCE"}
% HMT_RESULT_END:AMP-IV-HDG-01:residence
% HMT_RESULT_BEGIN:AMP-IV-BSD-01:residence
% HMT_ATOMIC_RESULT {"material_state":"AMPLIFICACION_AUTORIZADA_2026_08_17","proof_strength":"DEMOSTRADO_CON_ESTRUCTURA_DE_PARTIDA_EXPLICITA","provenance":"FORMALIZACION_NUEVA","reinforcement_state":"DEPLOYED_WITHOUT_NEW_RESULT_ID","residence":"manuscrito/sections/ampliacion_20260817/04e_bsd_comparadores_desplegados.tex","result_id":"AMP-IV-BSD-01","statement":"Los comparadores Kato–FERS/Iwasawa, Poitou–Tate/Poincaré–Néron–Tate, Smith–Tamagawa–Shafarevich, torsión y período se construyen y componen sobre una única línea orientada; la igualdad raíz elimina la unidad relativa y publica rango, finitud de Sha y fórmula final.","status":"PRESENT_IN_AMPLIFIED_SOURCE"}
% HMT_RESULT_END:AMP-IV-BSD-01:residence
% HMT_RESULT_BEGIN:AMP-IV-SUP-01:residence
% HMT_ATOMIC_RESULT {"material_state":"AMPLIFICACION_AUTORIZADA_2026_08_17","proof_strength":"EXACTO_INTERNO","provenance":"RESULTADO_RECUPERADO","reinforcement_state":"DEPLOYED_WITHOUT_NEW_RESULT_ID","residence":"manuscrito/sections/ampliacion_20260817/04_principio_realizacion_y_supersesion.tex; manuscrito/appendices_acumulativa_revfinal.tex","result_id":"AMP-IV-SUP-01","statement":"Los rótulos históricos que declaraban abiertos RH, NS, YM, Hodge y BSD quedan supersedidos en esta edición por los cinco teoremas de Parte IV; se conservan únicamente como trazabilidad de formulaciones anteriores y controles tipados, sin autoridad sobre el corte vigente.","status":"PRESENT_IN_AMPLIFIED_SOURCE"}
% HMT_RESULT_END:AMP-IV-SUP-01:residence
% HMT_RESULT_BEGIN:AMP-V-GRAV5-01:residence
% HMT_ATOMIC_RESULT {"material_state":"AMPLIFICACION_AUTORIZADA_2026_08_17","proof_strength":"EXACTO_INTERNO","provenance":"FORMALIZACION_NUEVA","reinforcement_state":"DEPLOYED_WITHOUT_NEW_RESULT_ID","residence":"manuscrito/sections/ampliacion_20260817/05_sector_tensorial_gravitatorio_cinco_componentes.tex","result_id":"AMP-V-GRAV5-01","statement":"La acción icosaédrica dodecafásica induce V12≅V1⊕V3⊕V3'⊕V5; el proyector central P5, el entrelazador explícito Γ5 hacia Sym0²(R3), la base STF y la reducción TT construyen la cadena 12→5→5→2 y separan portador quíntuple, espín dos masivo, GR y clases E(2).","status":"PRESENT_IN_AMPLIFIED_SOURCE"}
% HMT_RESULT_END:AMP-V-GRAV5-01:residence
% HMT_RESULT_BEGIN:AMP-V-GRAVPHY-01:residence
% HMT_ATOMIC_RESULT {"material_state":"AMPLIFICACION_AUTORIZADA_2026_08_17","proof_strength":"INTERFAZ_FISICA","provenance":"FORMALIZACION_NUEVA","reinforcement_state":"DEPLOYED_WITHOUT_NEW_RESULT_ID","residence":"manuscrito/sections/ampliacion_20260817/05_sector_tensorial_gravitatorio_cinco_componentes.tex","result_id":"AMP-V-GRAVPHY-01","statement":"El transductor G, la realización circular de orden 108 y el sector tensorial algebraico conservan sus conclusiones exactas; la realización PCH se expone como cadena tipada que determina el espacio físico mediante ecuaciones, restricciones y calibre, sin atribuir cinco polarizaciones a Einstein, Feynman ni GR.","status":"PRESENT_IN_AMPLIFIED_SOURCE"}
% HMT_RESULT_END:AMP-V-GRAVPHY-01:residence
% HMT_RESULT_BEGIN:AMP-EPI-R-01:residence
% HMT_ATOMIC_RESULT {"material_state":"AMPLIFICACION_AUTORIZADA_2026_08_17","proof_strength":"DEMOSTRADO_CON_ESTRUCTURA_DE_PARTIDA_EXPLICITA","provenance":"FORMALIZACION_NUEVA","reinforcement_state":"DEPLOYED_WITHOUT_NEW_RESULT_ID","residence":"manuscrito/sections/ampliacion_20260817/99_epilogo_retorno_recta_real.tex","result_id":"AMP-EPI-R-01","statement":"El epílogo retorna a R y demuestra que la recta real es la publicación arquimediana de una sección compatible, mientras la publicación excepcional conserva incidencia; ambas proceden del mismo estado y ninguna reconstruye por sí sola su genealogía.","status":"PRESENT_IN_AMPLIFIED_SOURCE"}
% HMT_RESULT_END:AMP-EPI-R-01:residence
 % Sidecar sucesor de la reparación semántica aprobada: 32 residencias
% rectoras, comentario-only y enlazadas al nuevo recibo de precompilación.
% Sidecar probatorio de la reparación semántica 2026-08-19.
% No produce salida tipográfica; cada intervalo liga una residencia rectora.
% HMT_REPAIR_RESULT_BEGIN::HMT-CORE-APP-TRIT-TPK
% RESULT_ID: HMT-CORE-APP-TRIT-TPK
% STATUS: DEMONSTRATED
% RESIDENCE: manuscrito/sections/sucesora/01_nucleo_formal_hmt_20260817.tex#eq:nucleo-app-dos-hojas
% STATEMENT_SHA256: 6657284d9bd2918464c9d964efdaae3608a94c55cd238f3de467fe85c00ad0ac
% STATEMENT: APP construye X_9; TRIT tipa sus dos hojas; después se forman los dominios de semilla, las emisiones y las palabras visibles, y sólo entonces se produce el estado holonómico integral que conserva hoja, orientación, residuo, cociente, acarreo, firma, cilindro, frontera, ruta, memoria, multisección y supervivencia.
% HMT_REPAIR_RESULT_END::HMT-CORE-APP-TRIT-TPK
% HMT_REPAIR_RESULT_BEGIN::HMT-CORE-NONADIC-JOINT-CONTINUITY
% RESULT_ID: HMT-CORE-NONADIC-JOINT-CONTINUITY
% STATUS: DEMONSTRATED
% RESIDENCE: manuscrito/sections/sucesora/01_nucleo_formal_hmt_20260817.tex#thm:nucleo-continuidad-conexion-nonadica-conjunta
% STATEMENT_SHA256: bea2a7fd24c115638fea7eb3682c82f9b684160dc0a4e855abb04024c2e484c4
% STATEMENT: Γ_9 devuelve la fase tras nueve pasos sin reinicializar el estado; conserva memoria, acarreo y terminal, y sus potencias son composiciones relacionales compatibles con truncamientos.
% HMT_REPAIR_RESULT_END::HMT-CORE-NONADIC-JOINT-CONTINUITY
% HMT_REPAIR_RESULT_BEGIN::HMT-RECTOR-CONTINUUM-CLOSURE-CONSTANTS
% RESULT_ID: HMT-RECTOR-CONTINUUM-CLOSURE-CONSTANTS
% STATUS: DEMONSTRATED
% RESIDENCE: manuscrito/sections/sucesora/02_estructura_discreta_continuo_20260817.tex#eq:continuo-teorema-fundamental-discreto
% STATEMENT_SHA256: 34c506d5e3311116e4a5f5d4baf1febf44fb1fcbd22e17a6540197f17479fbd7
% STATEMENT: HMT cierra el problema del continuo en el dominio genealógico: las publicaciones continuas son lectores compatibles de secciones discretas completas; las constantes emergen del mismo antecedente, y los cinco problemas del Milenio son especializaciones concretas y subordinadas del cierre rector.
% HMT_REPAIR_RESULT_END::HMT-RECTOR-CONTINUUM-CLOSURE-CONSTANTS
% HMT_REPAIR_RESULT_BEGIN::HMT-CORE-NO-ABLATION
% RESULT_ID: HMT-CORE-NO-ABLATION
% STATUS: DEMONSTRATED
% RESIDENCE: gestion/reparacion_semantica_20260818/MAPA_RECONSTRUCCION_CAUSAL_FUENTE_20260818.md §7.1
% STATEMENT_SHA256: 6a8421d9f55a806faa959c19578d710ae27e286555775c9a1f34551268c7e592
% STATEMENT: El mapa de olvido U retiene una sombra visible, conjuntista, ordinal, escalar o convencional y elimina genealogía, memoria, orientación, cociente/residuo, acarreo, frontera/ruta, supervivencia o multisección; ninguna objeción a esa sombra se transfiere al objeto completo sin probar que U conserva y refleja la propiedad examinada.
% HMT_REPAIR_RESULT_END::HMT-CORE-NO-ABLATION
% HMT_REPAIR_RESULT_BEGIN::HMT-ZFC-VON-NEUMANN-NONFAITHFUL-CH
% RESULT_ID: HMT-ZFC-VON-NEUMANN-NONFAITHFUL-CH
% STATUS: DEMONSTRATED
% RESIDENCE: manuscrito/sections/integracion_20260812/von_neumann_64/04_zfc_fundacion_ordinales.tex#thm:zfc-perdida-genealogia-olvido
% STATEMENT_SHA256: 9fcf0f78fc08d6066f7bf7f9f8aeb86d0a5e87348f2efb1affef488e061e0b9b
% STATEMENT: Tres resultados distintos quedan tipados: U_60 colapsa dos objetos no isomorfos con la misma Sig_60 y fibras Ext_81 de cardinales 0 y 12, por lo que no refleja isomorfías ni clasifica; V no es fiel porque id_ρ y Γ_9 tienen la misma sombra unipuntual pero grados de acarreo/memoria (0,0) y (1,1); los ordinales finitos con memorias k y k+1 proyectan al mismo ordinal de von Neumann. CH es un blanco distinto.
% HMT_REPAIR_RESULT_END::HMT-ZFC-VON-NEUMANN-NONFAITHFUL-CH
% HMT_REPAIR_RESULT_BEGIN::HMT-CORE-DOUBLE-PUBLICATION
% RESULT_ID: HMT-CORE-DOUBLE-PUBLICATION
% STATUS: DEMONSTRATED
% RESIDENCE: manuscrito/sections/sucesora/02_estructura_discreta_continuo_20260817.tex#eq:continuo-doble-publicacion
% STATEMENT_SHA256: 3051702377a6942a887f59b9323a35081f9d1ba17e9b3e8a72326b7039ab0e43
% STATEMENT: El lector arquimediano contrae cilindros compatibles y el lector incidencial conserva la genealogía de frontera; sus codominios son distintos. La rama excepcional es bifurcada: Hadamard organiza U_12^sgn en paralelo y Paley produce C_W, desde donde parten las ramas Witt–Golay–Mathieu y Niemeier–Leech–V^natural–Monstruo–Moonshine.
% HMT_REPAIR_RESULT_END::HMT-CORE-DOUBLE-PUBLICATION
% HMT_REPAIR_RESULT_BEGIN::HMT-MILLENNIUM-RH-WEIL-CLOSED
% RESULT_ID: HMT-MILLENNIUM-RH-WEIL-CLOSED
% STATUS: DEMONSTRATED
% RESIDENCE: manuscrito/sections/sucesora/04_cinco_realizaciones_fundamentales_20260817.tex#thm:realizaciones-rh
% STATEMENT_SHA256: a4f088786afdbee666309a8fc5034b3722fade65c9f6a259d05da2b5732626b5
% STATEMENT: En cada región R se fija q=5/9 y se construye E_R=j_-Q_RXi_R con AE_R=qE_R y E_R*E_R=B_{W,R}; la telescopía finita conserva el terminal q^(2N)B_{W,R}, cuya extinción sólo se toma al pasar N→∞, y L_R satisface L_R*L_R=B_{W,R}⪰0. El descenso cofinal publica la forma de Weil global como cuadrado positivo y el criterio de Weil reconoce la hipótesis de Riemann.
% HMT_REPAIR_RESULT_END::HMT-MILLENNIUM-RH-WEIL-CLOSED
% HMT_REPAIR_RESULT_BEGIN::HMT-MILLENNIUM-NS-CLOSED
% RESULT_ID: HMT-MILLENNIUM-NS-CLOSED
% STATUS: DEMONSTRATED
% RESIDENCE: manuscrito/sections/sucesora/04_cinco_realizaciones_fundamentales_20260817.tex#thm:realizaciones-ns
% STATEMENT_SHA256: f8e920750843d201de3d006e4d249eb58a18849ab02f1db10372401a049235db
% STATEMENT: La rama propietaria PC–Fock–KYP construye prospectivamente el Gram del terminal y las pérdidas futuras, deriva la identidad Stein y la telescopía con terminal, factoriza directamente el bloque KYP, establece balance, observabilidad y registro de torre, y produce el presupuesto raíz uniforme y la cota H^s global. El comparador G_{M,j}=P_aux*U^term_{M,j}P_aux es sólo un control no gobernante.
% HMT_REPAIR_RESULT_END::HMT-MILLENNIUM-NS-CLOSED
% HMT_REPAIR_RESULT_BEGIN::HMT-MILLENNIUM-YM-CLOSED
% RESULT_ID: HMT-MILLENNIUM-YM-CLOSED
% STATUS: DEMONSTRATED
% RESIDENCE: manuscrito/sections/sucesora/04_cinco_realizaciones_fundamentales_20260817.tex#thm:realizaciones-ym
% STATEMENT_SHA256: 7bae1f46e683874449e24e2e5e76323c00b6953291777330a09a1c2324fc5167
% STATEMENT: La misma familia proyectiva construye medida, red local, generador D_m^YM, semigrupo K_{m,t}=e^(−tD_m^YM), vacío, cuatro formas OS, acción fuerte E(4) y brecha exacta. Se cumplen K_{m,t}P_{Ω_m}=P_{Ω_m} y ||K_{m,t}(I−P_{Ω_m})||≤e^(−3κ_av t); el testigo W_c alcanza la cota. El lector clover/F² reconoce después la curvatura del mismo proceso.
% HMT_REPAIR_RESULT_END::HMT-MILLENNIUM-YM-CLOSED
% HMT_REPAIR_RESULT_BEGIN::HMT-MILLENNIUM-HODGE-CLOSED
% RESULT_ID: HMT-MILLENNIUM-HODGE-CLOSED
% STATUS: DEMONSTRATED
% RESIDENCE: manuscrito/sections/sucesora/04_hodge_bsd_fragment_20260817.tex#thm:hodge-suc-generacion
% STATEMENT_SHA256: 871872c44580d0e280556a6d45a7a549870a5b324223da5310a7a61a0ee0f72b
% STATEMENT: La realización gobernante parte de X_χ=lim_K Surv_K(X,p), conmuta con el carácter local y la publicación de clases y, para cada α, construye ξ_α y β_α con cl(β_α)=α. Los núcleos relativos y la presentación coend son controles de autonomía expositiva, no la cadena gobernante.
% HMT_REPAIR_RESULT_END::HMT-MILLENNIUM-HODGE-CLOSED
% HMT_REPAIR_RESULT_BEGIN::HMT-MILLENNIUM-BSD-CLOSED
% RESULT_ID: HMT-MILLENNIUM-BSD-CLOSED
% STATUS: DEMONSTRATED
% RESIDENCE: manuscrito/sections/sucesora/04_hodge_bsd_fragment_20260817.tex#eq:bsd-suc-formula-final
% STATEMENT_SHA256: 58f17d0b3189c35e80a258c376c281f796aee2740fea34308df90316cbe709c7
% STATEMENT: La coálgebra Alexander–Whitney y su counidad construyen G_{m,n} y el producto fibrado P_E^gen; una misma unidad genera las publicaciones Kato/PT, la igualdad de raíz, el relleno h_*, la bandera de Bockstein, los comparadores y la igualdad basada que cierra BSD. La unipotencia generador a generador es sólo control de despliegue.
% HMT_REPAIR_RESULT_END::HMT-MILLENNIUM-BSD-CLOSED
% HMT_REPAIR_RESULT_BEGIN::HMT-EXCEPTIONAL-CORRIDOR
% RESULT_ID: HMT-EXCEPTIONAL-CORRIDOR
% STATUS: DEMONSTRATED
% RESIDENCE: manuscrito/sections/hmt/09c_genealogia_excepcional_completa_autonoma.tex#sec:genealogia-excepcional-completa-autonoma
% STATEMENT_SHA256: 9e1f0885079d7e94e3bbf4142ac1fef08314e4d1d537ffc3403e23bd7d3bb08b
% STATEMENT: Hadamard organiza U_12^sgn sin convertirse en antecedente serial de Paley; la cadena C_P→A_W→C_W bifurca desde C_W hacia Witt–Golay–Mathieu y hacia Niemeier–Leech–V_Λ–V^natural; después Aut(V^natural)=Monstruo y J=T_{1A} precede a la familia {T_g} de Moonshine.
% HMT_REPAIR_RESULT_END::HMT-EXCEPTIONAL-CORRIDOR
% HMT_REPAIR_RESULT_BEGIN::HMT-INTERNAL-I-QUARTER-TURN
% RESULT_ID: HMT-INTERNAL-I-QUARTER-TURN
% STATUS: DEMONSTRATED
% RESIDENCE: manuscrito/sections/hmt/09a_paley_codigo_witt.tex#eq:AW-raiz-menos-uno
% STATEMENT_SHA256: 89af0b10c8bc5da185fb92a05117661adee2496d45d963a0922f2c33d9cf247e
% STATEMENT: La matriz incidencial A_W satisface A_W²=−I en el sector Paley–Witt; ésta es la raíz finita interna. La estructura J_E²=−I y su lector exponencial publican después la realización compleja arquimediana, sin identificar ambos operadores por nombre.
% HMT_REPAIR_RESULT_END::HMT-INTERNAL-I-QUARTER-TURN
% HMT_REPAIR_RESULT_BEGIN::HMT-INTERNAL-MINUS-ONE-TWELFTH
% RESULT_ID: HMT-INTERNAL-MINUS-ONE-TWELFTH
% STATUS: DEMONSTRATED
% RESIDENCE: manuscrito/sections/sucesora/03_genealogia_constantes_20260817.tex#eq:constantes-menos-un-doce-proyectores
% STATEMENT_SHA256: d04beccf9787ff880f26f73a4ff543adb5aab378144657d59b5d156e5fd65284
% STATEMENT: El valor −1/12 emerge en seis dominios tipados: incidencia 90/120, diferencia de periodos, Gram positivo con pseudoinversa, extractor de escala, exponente eta y traza de proyectores del corrimiento cíclico S_12:Q^12→Q^12, con S_12^12=I, donde D_r=I−S_12^r y los proyectores sobre ker D_3 y ker D_4 tienen rangos 3 y 4. La igualdad ζ(−1)=−1/12 es un reconocimiento analítico separado, no el origen de esas realizaciones.
% HMT_REPAIR_RESULT_END::HMT-INTERNAL-MINUS-ONE-TWELFTH
% HMT_REPAIR_RESULT_BEGIN::HMT-LORENTZ-MINKOWSKI
% RESULT_ID: HMT-LORENTZ-MINKOWSKI
% STATUS: DEMONSTRATED
% RESIDENCE: manuscrito/sections/hmt/03_app_ortograma.tex#prop:app-trit-espectral-markov-lorentz-rev8
% STATEMENT_SHA256: 7da5585d1b7752e0cd76fc580631d618b0aedef2c4bc51e326b4cc8bae1649e5
% STATEMENT: El bloque central B_c de APP produce M_c=B_c/√45∈SO^+(1,1) y satisface M_c^Tη_{1,1}M_c=η_{1,1}; ésta es una realización exacta de Lorentz y del plano de Minkowski 1+1. Toda elevación 3+1 requiere un mapa dimensional tipado y no demueve el resultado 1+1.
% HMT_REPAIR_RESULT_END::HMT-LORENTZ-MINKOWSKI
% HMT_REPAIR_RESULT_BEGIN::HMT-HILBERT-REALIZATION
% RESULT_ID: HMT-HILBERT-REALIZATION
% STATUS: DEMONSTRATED
% RESIDENCE: manuscrito/sections/integracion_20260812/von_neumann_64/06_hilbert_heisenberg.tex
% STATEMENT_SHA256: 5323eb3e082071f29482590ea4ede54b91c92314e2f7d6f865ddf5cde5fea80d
% STATEMENT: El espacio de Hilbert matricial, el límite UHF y su cierre GNS, la dinámica compleja y Born se construyen por ramas tipadas que no se implican automáticamente: C_9^(d) produce H_d,A_d,UHF(9^∞), su traza τ produce la representación GNS (π_τ,H_τ,Ω_τ) y π_τ(A_(9^∞))ʺ=R_hyp; J_E,H∈B(H_τ), con J_E²=−I y H=H*, producen U(t); el lector Born GNS usa P=P*=P²∈A_UHF actuando por π_τ(P) y p_τ(P)=<Ω_τ,π_τ(P)Ω_τ>; una segunda realización finita conserva Pr_ρ(a)=Tr(ρE_a) desde una PVM y un estado normalizado.
% HMT_REPAIR_RESULT_END::HMT-HILBERT-REALIZATION
% HMT_REPAIR_RESULT_BEGIN::HMT-T-DUALITY-M-SCREENS
% RESULT_ID: HMT-T-DUALITY-M-SCREENS
% STATUS: DEMONSTRATED
% RESIDENCE: manuscrito/sections/hmt/16b_pantallas_dualidad_teoria_m.tex#thm:dualidad-visible-memoria
% STATEMENT_SHA256: a6c3111b2029fd7c5e93e357ea6b65a0cd583f22589de276bcbf2a3215fc5d31
% STATEMENT: En D_dual=Z²×R_(>0), T(m,w,R_0)=(w,m,R_0^(−1)) satisface T²=id y preserva E_(R_0)(m,w)=m²/R_0²+w²R_0² mediante E_(R_0)(m,w)=E_(R_0^(−1))(w,m); su realización unitaria convive con las reducciones 11→10 y 26→24. La interfaz física separada R_M:X_HMT^enr→X_11^phys debe producir y compatibilizar g_(MN)∈Γ(Sym² T*X_11), C_(MNP)∈Ω³(X_11), ψ_M∈Γ(S_11⊗T*X_11), su acción y transformaciones, sin contaminar el teorema interno.
% HMT_REPAIR_RESULT_END::HMT-T-DUALITY-M-SCREENS
% HMT_REPAIR_RESULT_BEGIN::HMT-CONSTANT-PI
% RESULT_ID: HMT-CONSTANT-PI
% STATUS: DEMONSTRATED
% RESIDENCE: manuscrito/sections/ampliacion_20260818/02_generacion_infinita_pi_e_phi.tex#eq:generacion-infinita-gaussiana-pi-20260818
% STATEMENT_SHA256: 26b11a14165eacab3aa7d06fcb90f31ff0f2eb0dd328b3f7c11cf183fee1293f
% STATEMENT: La órbita y el carácter 5/4/239 se seleccionan antes de abrir intervalos reales; la identidad de Machin reconoce posteriormente el lector ya construido.
% HMT_REPAIR_RESULT_END::HMT-CONSTANT-PI
% HMT_REPAIR_RESULT_BEGIN::HMT-CONSTANT-E
% RESULT_ID: HMT-CONSTANT-E
% STATUS: DEMONSTRATED
% RESIDENCE: manuscrito/sections/ampliacion_20260818/02_generacion_infinita_pi_e_phi.tex#eq:generacion-infinita-recurrencia-e-20260818
% STATEMENT_SHA256: 2c21447e726bae153a60dc20b39c536b01f8b4dba8836c1176eeb252af969ec0
% STATEMENT: La ley multiplicativa del canal fuerza la recurrencia a_n=1/n! antes de reconocer la serie exponencial y su valor e.
% HMT_REPAIR_RESULT_END::HMT-CONSTANT-E
% HMT_REPAIR_RESULT_BEGIN::HMT-CONSTANT-PHI
% RESULT_ID: HMT-CONSTANT-PHI
% STATUS: DEMONSTRATED
% RESIDENCE: capítulo autónomo de la constante de autoescala, ecuación de prolongación 20260818
% STATEMENT_SHA256: 8151c9f0c4164daffee20672e7f6751ed9452b96c233d2ecd5e81aedca948338
% STATEMENT: La incidencia areal–volumétrica induce F_av y su rayo de Perron fuerza x²−x−1=0; la raíz positiva se reconoce después como φ.
% HMT_REPAIR_RESULT_END::HMT-CONSTANT-PHI
% HMT_REPAIR_RESULT_BEGIN::HMT-CONSTANT-ALPHA
% RESULT_ID: HMT-CONSTANT-ALPHA
% STATUS: DEMONSTRATED
% RESIDENCE: manuscrito/sections/sucesora/03_genealogia_constantes_20260817.tex#thm:constantes-alpha-independencia-principal
% STATEMENT_SHA256: 8eb80629b88f75eaf8f624a0e6992e43812a576d8f34b530137934ff6bf0d7a6
% STATEMENT: En el perfil HMT completo, B_term y U_12^sgn son proyecciones paralelas del mismo estado terminal; R_sgn extrae U, H_12 invertible da K, y el acarreo y la recurrencia producen α sin α, CODATA ni masa como entrada.
% HMT_REPAIR_RESULT_END::HMT-CONSTANT-ALPHA
% HMT_REPAIR_RESULT_BEGIN::HMT-EULER-HOLOGRAPHIC
% RESULT_ID: HMT-EULER-HOLOGRAPHIC
% STATUS: DEMONSTRATED
% RESIDENCE: manuscrito/sections/ampliacion_20260818/03_extension_holografica_euler.tex#thm:euler-holografica-tres-cierres-20260818
% STATEMENT_SHA256: f54720a289461a9177197a7184922777c9bb186d4617980d744cae299400c025
% STATEMENT: A_W²=−I proporciona la raíz finita incidencial; en una rama distinta, J_E²=−I y el lector exponencial producen la identidad arquimediana de Euler, extendida por J_τ a tres ramas. El contraángulo no nace de Euler: sólo después de α_HMT, la microfibra 169, D_A, H_5 y la recurrencia regional generan C*, cuya identidad inversa se reconoce al final.
% HMT_REPAIR_RESULT_END::HMT-EULER-HOLOGRAPHIC
% HMT_REPAIR_RESULT_BEGIN::HMT-EMRH-CUBIC-UNIT
% RESULT_ID: HMT-EMRH-CUBIC-UNIT
% STATUS: DEMONSTRATED
% RESIDENCE: manuscrito/sections/ampliacion_20260818/01_emrh_desarrollo_integral.tex#eq:emrh-cierre-apice-20260818
% STATEMENT_SHA256: 6f630c70e8d9bf2ba536e9a19b4f59ce7f4f90709e8c1f3410e2cf49cfab5724
% STATEMENT: La completación correcta es 10^3=729+271, con 271 como corona E y 729 como bloque alfa; además (1/271+1/729)^−1=271·729/1000 expresa exactamente la razón inversa asociada a la completación volumétrica.
% HMT_REPAIR_RESULT_END::HMT-EMRH-CUBIC-UNIT
% HMT_REPAIR_RESULT_BEGIN::HMT-CONSTANTS-FOUR-STRATA
% RESULT_ID: HMT-CONSTANTS-FOUR-STRATA
% STATUS: DEMONSTRATED
% RESIDENCE: manuscrito/sections/ampliacion_20260818/05_constantes_fisicas_genealogia.tex#thm:amp18-constantes-cuatro-estratos
% STATEMENT_SHA256: f67dfab70770a2b7aed78a5925e7f56252e00166de0b0e1ce805a0be3520e083
% STATEMENT: Toda constante se presenta en el orden constructor discreto/genealógico, caracterización analítica, transducción dimensional equivariante y comparación metrológica; la carta SI no construye el invariante.
% HMT_REPAIR_RESULT_END::HMT-CONSTANTS-FOUR-STRATA
% HMT_REPAIR_RESULT_BEGIN::HMT-CONSTANT-PLANCK
% RESULT_ID: HMT-CONSTANT-PLANCK
% STATUS: DEMONSTRATED
% RESIDENCE: manuscrito/sections/md/03b_escala_accion_determinantal.tex#thm:orden-accion
% STATEMENT_SHA256: 0bc899b4d06add106a366bf912d47203626d8fde1c6911f1073da88cc17eeef0
% STATEMENT: El lector determinantal Σ_H=φ_HMT α_HMT^16 y su cokernel local de orden 16 fijan la recta de acción y la década −34; la mantisa y la carta J·s se distinguen, y cualquier residual permanece visible.
% HMT_REPAIR_RESULT_END::HMT-CONSTANT-PLANCK
% HMT_REPAIR_RESULT_BEGIN::HMT-CONSTANT-BOLTZMANN
% RESULT_ID: HMT-CONSTANT-BOLTZMANN
% STATUS: DEMONSTRATED
% RESIDENCE: manuscrito/sections/ampliacion_20260818/05_constantes_fisicas_genealogia.tex#prop:amp18-constantes-covariancia-termica
% STATEMENT_SHA256: c61e51204132bcfb86f480ea2592cef625c3aa9bdbbce4632ba8c5995cc47828
% STATEMENT: El invariante interno fija k_B Θ_clk log 3=ħ_int 2π/(108t_0); la covariancia Θ→λΘ, k_B→k_B/λ demuestra que el kelvin fija la carta, no el objeto.
% HMT_REPAIR_RESULT_END::HMT-CONSTANT-BOLTZMANN
% HMT_REPAIR_RESULT_BEGIN::HMT-CONSTITUTIVES-MU0-EPSILON0
% RESULT_ID: HMT-CONSTITUTIVES-MU0-EPSILON0
% STATUS: DEMONSTRATED
% RESIDENCE: manuscrito/sections/ampliacion_20260818/05_constantes_fisicas_genealogia.tex#eq:amp18-constantes-identidades-vacio
% STATEMENT_SHA256: 5a67e3f2c22019847d22853c8038fa4258f414e3c63a0d1568a2e08908c456fd
% STATEMENT: A partir de α_HMT y una única carta global de e,h,c se construye la impedancia Z_0 y, mediante Z_0=μ_0c y Z_0c=1/ε_0, se publican μ_0 y ε_0; sus valores SI son comparación posterior.
% HMT_REPAIR_RESULT_END::HMT-CONSTITUTIVES-MU0-EPSILON0
% HMT_REPAIR_RESULT_BEGIN::HMT-CONSTANT-G
% RESULT_ID: HMT-CONSTANT-G
% STATUS: DEMONSTRATED
% RESIDENCE: manuscrito/sections/ampliacion_20260818/05_constantes_fisicas_genealogia.tex#eq:amp18-constantes-selector-108
% STATEMENT_SHA256: 4e8377ac898d655608bc820b018d8760487159e3f6349a41d5aef00743e4fae5
% STATEMENT: El selector circular fija θ_108=(54/π_HMT)^2 y el transductor exacto define G_int(θ_108). La vía holonómica independiente da la plantilla homogénea G_hol=(c_int^3/ħ_int)(δ_θΛ_27)^2, cuya cifra exige construir y normalizar Λ_27; la carta decimal SI se mantiene separada y no construye ninguna de las dos vías.
% HMT_REPAIR_RESULT_END::HMT-CONSTANT-G
% HMT_REPAIR_RESULT_BEGIN::HMT-MASS-LAW-ATLAS
% RESULT_ID: HMT-MASS-LAW-ATLAS
% STATUS: DEMONSTRATED
% RESIDENCE: manuscrito/sections/ampliacion_20260818/04_electron_y_ley_general_masas.tex#eq:amp18-masas-ley-operatorial
% STATEMENT_SHA256: 4666f47d67c5e3e1b827ca0cb3802e083fbee5770e215ceea03a3ab43eefdc84
% STATEMENT: La partícula es sección persistente y la masa es carácter positivo de esa persistencia. El electrón central único (5,5) construye primero el basal adimensional \widetilde{\mathfrak e}_0=(√3/4)[exp(φ_HMT/π_HMT²)−22α_HMT³](1+15Δ_4), después \widetilde{\mathfrak e}_e^β=\widetilde{\mathfrak e}_0R_act^(80−54α_HMT+6Δ_4), y sólo entonces la carta común T_E:=\varsigma_{u_E}:R_{>0}^{adim}→R_{>0}u_E, T_E(x)=xu_E, u_E=1 MeV publica \mathfrak e_e^β=T_E(\widetilde{\mathfrak e}_e^β)=m_e^βc²=0.5109989506900008086082571648… MeV y m_e^β=\mathfrak e_e^β/c²=0.5109989506900008086082571648… MeV/c². La misma carta común gobierna el atlas 81/56/13/324; los 23 sectores son sólo un inventario humano separado.
% HMT_REPAIR_RESULT_END::HMT-MASS-LAW-ATLAS
% HMT_REPAIR_RESULT_BEGIN::HMT-BARBERO-IMMIRZI
% RESULT_ID: HMT-BARBERO-IMMIRZI
% STATUS: DEMONSTRATED
% RESIDENCE: manuscrito/sections/md/07_hexada_octada_barbero.tex#chap:barbero
% STATEMENT_SHA256: 360316071f48bb431d8e0b9a737194c74511243d354eafeaf5e398f932b048e9
% STATEMENT: Las incidencias 90/120, las series de los dos canales, el residuo 1/24 y la minimalidad diofántica fijan los pesos (12,1) y el funcional Γ_6→8; la identificación γ=Γ ocurre después y no define retroactivamente A o C*.
% HMT_REPAIR_RESULT_END::HMT-BARBERO-IMMIRZI
% HMT_REPAIR_RESULT_BEGIN::HMT-CKM-CP
% RESULT_ID: HMT-CKM-CP
% STATUS: DEMONSTRATED
% RESIDENCE: manuscrito/sections/ampliacion_20260818/05_constantes_fisicas_genealogia.tex#eq:amp18-constantes-ckm-cuatro-ecuaciones
% STATEMENT_SHA256: a245142e9d5486d59a0cd73afac1263a4916dd3b49d3437964457a148a0d3878
% STATEMENT: El registro de rutas y el único par angular construyen cuatro ángulos, la matriz unitaria CKM y un invariante de Jarlskog no nulo; la comparación PDG es posterior y CP no se confunde con CPT.
% HMT_REPAIR_RESULT_END::HMT-CKM-CP
% HMT_REPAIR_RESULT_BEGIN::HMT-GOVERNANCE-NO-REGRESSION
% RESULT_ID: HMT-GOVERNANCE-NO-REGRESSION
% STATUS: DEMONSTRATED
% RESIDENCE: gestion/reparacion_semantica_20260818/MAPA_RECONSTRUCCION_CAUSAL_FUENTE_20260818.md §7
% STATEMENT_SHA256: c59a160a4fe22784718235a339c7a3a3b5df6305829c5ef7e443d93c2ca38b7a
% STATEMENT: Todo artefacto obsoleto que invierta APP→TRIT→estado del TPK, haga externas las semillas, convierta los cinco cierres en tareas pendientes o subordine el continuo al marco clásico se retira de la gobernanza o queda como testigo inactivo; no se arrastra a la edición sucesora.
% HMT_REPAIR_RESULT_END::HMT-GOVERNANCE-NO-REGRESSION
 % Entorno editorial para incorporar cuerpos científicos completos sin
% crear capitulos nuevos. El grupo conserva intactos los archivos de origen:
% solo traduce su jerarquia visible un nivel hacia abajo durante el input.
% La procedencia material de cada uso se registra en comentarios del destino
% y en gestion/RESTAURACION_16_CLAVES_SIN_DESTINO_20260824.tsv.
\newenvironment{HMTScientificOwnerSubordinated}{%
  \begingroup
  \let\HMTScientificOwnerSection\section
  \let\HMTScientificOwnerSubsection\subsection
  \let\HMTScientificOwnerSubsubsection\subsubsection
  \let\HMTScientificOwnerParagraph\paragraph
  \let\HMTScientificOwnerSubparagraph\subparagraph
  \let\chapter\HMTScientificOwnerSection
  \let\section\HMTScientificOwnerSubsection
  \let\subsection\HMTScientificOwnerSubsubsection
  \let\subsubsection\HMTScientificOwnerParagraph
  \let\paragraph\HMTScientificOwnerSubparagraph
}{%
  \par
  \endgroup
}
 \colorlet{hmtlight}{hmtpaper}
\theoremstyle{definition}
\newtheorem{principle}[theorem]{Principe}

% Jerarquía tipográfica de monografía científica. Los títulos de capítulo
% comparten línea con su numeración, se componen con cuerpo moderado y dejan
% espacio suficiente para el primer párrafo o la primera ecuación.
\titleformat{\part}[display]
  {\centering\normalfont\bfseries\color{hmtblue}}
  {\Large\scshape\partname\ \thepart}
  {.45ex}
  {\Huge}
\titlespacing*{\part}{0pt}{1.25cm}{1.0cm}
\titleformat{\chapter}[hang]
  {\normalfont\bfseries\color{hmtblue}\LARGE}
  {\normalsize\scshape\chaptertitlename\ \thechapter}
  {1.0em}
  {\raggedright}
\titlespacing*{\chapter}{0pt}{-0.35cm}{0.75cm}
\titleformat{name=\chapter,numberless}[block]
  {\normalfont\bfseries\color{hmtblue}\LARGE}
  {}{0pt}{\raggedright}
\titlespacing*{name=\chapter,numberless}{0pt}{-0.35cm}{.75cm}
\titleformat{\section}
  {\normalfont\bfseries\Large\color{hmtblue}}
  {\thesection}{.65em}{}
\titlespacing*{\section}{0pt}{2.0ex plus .45ex minus .2ex}{.65ex}
\titleformat{\subsection}
  {\normalfont\bfseries\large}
  {\thesubsection}{.65em}{}
\titlespacing*{\subsection}{0pt}{1.65ex plus .4ex minus .2ex}{.55ex}
\titleformat{\subsubsection}
  {\normalfont\itshape\normalsize}
  {\thesubsubsection}{.65em}{}
\titlespacing*{\subsubsection}{0pt}{1.35ex plus .3ex minus .2ex}{.4ex}

% Jerarquía editorial de la edición revisada. Los archivos heredados
% conservan sus rótulos y etiquetas, pero sus capítulos se presentan como
% secciones dentro de capítulos macroscópicos.
\let\HMTBookChapter\chapter
\let\HMTBookSection\section
\let\HMTBookSubsection\subsection
\let\HMTBookSubsubsection\subsubsection
\let\HMTBookParagraph\paragraph
\let\HMTBookSubparagraph\subparagraph
\newcommand{\HMTDemoteHeadings}{%
  \let\chapter\HMTBookSection
  \let\section\HMTBookSubsection
  \let\subsection\HMTBookSubsubsection
  \let\subsubsection\HMTBookParagraph
  \let\paragraph\HMTBookSubparagraph}
\newcommand{\HMTRestoreHeadings}{%
  \let\chapter\HMTBookChapter
  \let\section\HMTBookSection
  \let\subsection\HMTBookSubsection
  \let\subsubsection\HMTBookSubsubsection
  \let\paragraph\HMTBookParagraph
  \let\subparagraph\HMTBookSubparagraph}

% Notación de la frontera con Mecánica Dimensional.
\providecommand{\APP}{\ensuremath{\mathrm{APP}}}
\providecommand{\TPK}{\ensuremath{\mathrm{TPK}}}
\providecommand{\R}{\mathbb R}
\providecommand{\Z}{\mathbb Z}
\providecommand{\N}{\mathbb N}
\providecommand{\F}{\mathbb F}
\providecommand{\M}{\mathbb M}
\providecommand{\one}{\mathbf 1}
\providecommand{\Halg}{\mathfrak H}
\providecommand{\cA}{\mathcal A}
\providecommand{\cC}{\mathcal C}
\providecommand{\cR}{\mathcal R}
\providecommand{\cX}{\mathcal X}
\providecommand{\dd}{\mathrm d}
\providecommand{\e}{\mathrm e}
\providecommand{\ii}{\mathrm i}
\providecommand{\Znine}{\mathbb Z/9\mathbb Z}
\providecommand{\Cstar}{C^{*}}
\providecommand{\Cstarrad}{C^{*}_{\rm rad}}
\providecommand{\hbarHMT}{\hbar_{\mathrm{HMT}}}
\providecommand{\etaRet}{\eta_{\mathrm{ret}}^{(\mathrm{deg})}}
\providecommand{\etaRetRad}{\eta_{\mathrm{ret}}^{(\mathrm{rad})}}
\providecommand{\alphaAngle}{\alpha_{\angle,\mathrm{deg}}}
\providecommand{\alphaHMT}{\alpha}
\providecommand{\piHMT}{\pi}
\providecommand{\phiHMT}{\varphi}
\providecommand{\eexp}{\mathrm e}
\providecommand{\Dfour}{\Delta_{4}}
\providecommand{\zCorr}{\zeta_{270}^{\mathrm{corr}}}
\providecommand{\sSpec}{s_{270}^{\mathrm{spec}}}
\providecommand{\Kfield}{\mathbb K}
\providecommand{\rank}{\operatorname{rank}}
\providecommand{\nint}{\operatorname{nint}}
\providecommand{\MD}{\ensuremath{\mathrm{MD}}}
\providecommand{\CT}{\ensuremath{\mathrm{CT}_{108}}}
\providecommand{\RR}{\mathbb R}
\providecommand{\ZZ}{\mathbb Z}
\providecommand{\PP}{\mathbb P}
\providecommand{\hbarpre}{\hbar_{\mathrm{pre}}}
\providecommand{\hbarret}{\hbar_{\mathrm{ret}}}
\providecommand{\Ract}{R_{\mathrm{act}}}
\providecommand{\ellP}{\ell_{\mathrm P}}
\providecommand{\mP}{m_{\mathrm P}}
\providecommand{\tP}{t_{\mathrm P}}
\providecommand{\EP}{E_{\mathrm P}}
\providecommand{\rg}{r_{\mathrm g}}
\providecommand{\rs}{r_{\mathrm s}}
\providecommand{\lambdabarC}{\bar\lambda_{\mathrm C}}
\providecommand{\lambdaC}{\lambda_{\mathrm C}}
\providecommand{\Nsp}{N_{\mathrm{sp}}}
\providecommand{\thetaStar}{\theta_{\!*}}
\providecommand{\Gammae}{\Gamma_{e}}
\providecommand{\mebeta}{m_{e}^{\beta}}
% En esta monografía las letras universales no designan duplicados de las
% constantes: son las mismas coordenadas producidas por la construcción y
% reconocidas después en la realización arquimediana.
\renewcommand{\alphaHMT}{\alpha}
\renewcommand{\piHMT}{\pi}
\renewcommand{\phiHMT}{\varphi}
\providecommand{\epigraphtext}[2]{%
  \begin{flushright}%
  \begin{minipage}{0.76\textwidth}\itshape
  #1\par\smallskip
  \raggedleft\normalfont--- #2
  \end{minipage}%
  \end{flushright}%
}

\newtcolorbox{historybox}[1][]{%
  colback=hmtlightblue,colframe=hmtblue,
  title={Fait historique documenté},fonttitle=\bfseries,#1}
\newtcolorbox{readingbox}[1][]{%
  colback=hmtlightgold,colframe=hmtgold,
  title={Lectura HMT},fonttitle=\bfseries,#1}
\newtcolorbox{statusbox}[1][]{%
  colback=white,colframe=hmtblue!55,
  title={Hypothèses, domaine et conclusion},fonttitle=\bfseries,#1}
\newtcolorbox{warningbox}[1][]{%
  colback=white,colframe=hmtgold!75!black,
  title={Condition de compatibilité},fonttitle=\bfseries,#1}
\newtcolorbox{thesisbox}[1][]{%
  colback=white,colframe=hmtviolet,
  title={Tesis rectora},fonttitle=\bfseries,#1}
\newtcolorbox{proofstatusbox}[1][]{%
  colback=white,colframe=hmtgreen,
  title={Estatuto probatorio},fonttitle=\bfseries,#1}
\newtcolorbox{falsifierbox}[1][]{%
  colback=white,colframe=hmtred,
  title={Critère de réfutation},fonttitle=\bfseries,#1}
\newtcolorbox{falsadorBox}[1][]{%
  colback=hmtlightred,colframe=hmtred,
  title={Control negativo},fonttitle=\bfseries,#1}
\newtcolorbox{fronteraBox}[1][]{%
  colback=hmtlightgold,colframe=hmtgold,
  title={Frontière de validité},fonttitle=\bfseries,#1}
\newtcolorbox{programbox}[1][]{%
  colback=hmtlightgold,colframe=hmtgold,
  title={Hypothèses et programme de démonstration},fonttitle=\bfseries,#1}
\newtcolorbox{takeawaybox}[1][]{%
  colback=white,colframe=hmtblue,
  title={Résultat de la section},fonttitle=\bfseries,#1}
% Entornos académicos conservados por los fuentes teoremáticas científicos extensos.
% Se tipan aquí con la jerarquía gráfica única del tratado para evitar que cada
% incorporación restablezca una tipografía o un preámbulo autónomos.
\newtcolorbox{resultado}[1][]{%
  colback=white,colframe=hmtgreen,
  title={Résultat},fonttitle=\bfseries,#1}
\newtcolorbox{definicionnarrativa}[1][]{%
  colback=hmtlightblue,colframe=hmtblue,
  title={Définition},fonttitle=\bfseries,#1}
\newtcolorbox{lecturafisica}[1][]{%
  colback=hmtlightgold,colframe=hmtgold,
  title={Réalisation physique},fonttitle=\bfseries,#1}
\newtcolorbox{frontera}[1][]{%
  colback=hmtlightgold,colframe=hmtgold,
  title={Frontière de validité},fonttitle=\bfseries,#1}
\newtheorem{falsifier}[theorem]{Critère de réfutation}
\newtheorem{test}[theorem]{Programme de comparaison}
\newenvironment{keyresult}
  {\begin{quote}\color{hmtblue}\bfseries}
  {\end{quote}}

\lstdefinestyle{hmtcode}{
  basicstyle=\ttfamily\scriptsize,
  backgroundcolor=\color{hmtpaper},
  frame=single,
  rulecolor=\color{hmtgray!45},
  breaklines=true,
  showstringspaces=false,
  columns=fullflexible,
  keepspaces=true
}

\hypersetup{
  colorlinks=true,
  bookmarksdepth=3,
  pdflang={es-ES},
  linkcolor=hmtblue,
  citecolor=hmtblue,
  urlcolor=hmtgold,
  pdftitle={Double projection holographique d'un cristal temporel apériodique},
  pdfsubject={Construction du continu, génération de constantes et stabilisateurs exceptionnels de l'incidence de frontière},
  pdfauthor={Oumar Haidara Fall; Rubén Ramos Balsa},
  pdfkeywords={arithmétique nonadique, orientation ternaire, holonomie,
    cristal temporel apériodique, continu, constantes fondamentales,
    symétries exceptionnelles, action, information, gravité}
}

\setstretch{1.035}
\setlength{\parskip}{0.16em}
\setlength{\parindent}{1.25em}
% El índice científico de la edición definitiva hace visibles capítulos y
% secciones rectoras. Los niveles inferiores permanecen en el cuerpo para
% conservar la demostración sin convertir el índice en un inventario.
\setcounter{tocdepth}{3}
\setcounter{secnumdepth}{3}
\renewcommand{\contentsname}{Table des matières}
\renewcommand{\cftchappagefont}{\fontsize{11.5}{13.4}\selectfont\bfseries}
\renewcommand{\cftchapfont}{\fontsize{11.5}{13.4}\selectfont\bfseries}
\renewcommand{\cftsecfont}{\fontsize{10.75}{12.8}\selectfont\bfseries}
\renewcommand{\cftsecpagefont}{\fontsize{10.75}{12.8}\selectfont\bfseries}
\renewcommand{\cftsubsecfont}{\fontsize{10.25}{12.2}\selectfont}
\renewcommand{\cftsubsecpagefont}{\fontsize{10.25}{12.2}\selectfont}
\renewcommand{\cftsubsubsecfont}{\fontsize{9.5}{11.4}\selectfont\itshape}
\renewcommand{\cftsubsubsecpagefont}{\fontsize{9.5}{11.4}\selectfont\itshape}
\setlength{\cftbeforechapskip}{.38em}
\setlength{\cftbeforesecskip}{.05em}
\setlength{\cftchapnumwidth}{3.25em}
\setlength{\cftsecindent}{3.25em}
\setlength{\cftsecnumwidth}{4.15em}
\setlength{\cftsubsecnumwidth}{5.4em}
\setlist[itemize]{leftmargin=1.65em,itemsep=.25em,topsep=.35em}
\setlist[enumerate]{leftmargin=2em,itemsep=.25em,topsep=.35em}
\emergencystretch=3em
\widowpenalty=10000
\clubpenalty=10000
\displaywidowpenalty=10000
\brokenpenalty=10000
\predisplaypenalty=500
\postdisplaypenalty=500
\raggedbottom
% Cada epígrafe debe permanecer unido a un desarrollo mínimo. Se protegen
% tanto los mandos públicos como sus copias editoriales, pues los capítulos
% heredados se degradan localmente mediante \HMTDemoteHeadings.
\pretocmd{\section}{\Needspace{9\baselineskip}}{}{}
\pretocmd{\subsection}{\Needspace{7\baselineskip}}{}{}
\pretocmd{\subsubsection}{\Needspace{5\baselineskip}}{}{}
\pretocmd{\HMTBookSection}{\Needspace{9\baselineskip}}{}{}
\pretocmd{\HMTBookSubsection}{\Needspace{7\baselineskip}}{}{}
\pretocmd{\HMTBookSubsubsection}{\Needspace{5\baselineskip}}{}{}
\pretocmd{\HMTBookParagraph}{\Needspace{9\baselineskip}}{}{}

\crefname{chapter}{chapitre}{chapitres}
\Crefname{chapter}{Chapitre}{Chapitres}
\crefname{section}{section}{sections}
\Crefname{section}{Section}{Sections}
\crefname{subsection}{section}{sections}
\Crefname{subsection}{Section}{Sections}
\crefname{subsubsection}{section}{sections}
\Crefname{subsubsection}{Section}{Sections}
\crefname{figure}{figure}{figures}
\Crefname{figure}{Figure}{Figures}
\crefname{table}{tableau}{tableaux}
\Crefname{table}{Tableau}{Tableaux}
\crefname{equation}{équation}{équations}
\Crefname{equation}{Équation}{Équations}
\crefname{theorem}{théorème}{théorèmes}
\Crefname{theorem}{Théorème}{Théorèmes}
\crefname{lemma}{lemme}{lemmes}
\Crefname{lemma}{Lemme}{Lemmes}
\crefname{proposition}{proposition}{propositions}
\Crefname{proposition}{Proposition}{Propositions}
\crefname{corollary}{corollaire}{corollaires}
\Crefname{corollary}{Corollaire}{Corollaires}
\crefname{definition}{définition}{définitions}
\Crefname{definition}{Définition}{Définitions}
\crefname{construction}{construction}{constructions}
\Crefname{construction}{Construction}{Constructions}
\crefname{algorithm}{algorithme}{algorithmes}
\Crefname{algorithm}{Algorithme}{Algorithmes}
\crefname{ablation}{proposition de nécessité structurelle}
  {propositions de nécessité structurelle}
\Crefname{ablation}{Proposition de nécessité structurelle}
  {Propositions de nécessité structurelle}
\crefname{remark}{remarque}{remarques}
\Crefname{remark}{Remarque}{Remarques}
\crefname{principle}{principe}{principes}
\Crefname{principle}{Principe}{Principes}
\crefname{criterion}{critère de comparaison}{critères de comparaison}
\Crefname{criterion}{Critère de comparaison}{Critères de comparaison}
\AtBeginDocument{%
  \renewcommand{\crefrangeconjunction}{ à\nobreakspace}%
  \renewcommand{\crefpairconjunction}{ et\nobreakspace}%
  \renewcommand{\crefmiddleconjunction}{, }%
  \renewcommand{\creflastconjunction}{ et\nobreakspace}%
  \renewcommand{\crefpairgroupconjunction}{ et\nobreakspace}%
  \renewcommand{\crefmiddlegroupconjunction}{, }%
  \renewcommand{\creflastgroupconjunction}{ et\nobreakspace}%
}

\pagestyle{fancy}
\fancyhf{}
\fancyhead[LE,RO]{\small\thepage}
% Las cajas de marca se reservan dentro del 84 % central del encabezado. Así
% los títulos extensos pueden ocupar dos líneas sin invadir el folio exterior.
\fancyhead[LO]{\parbox[b]{.84\headwidth}{\raggedright\footnotesize\nouppercase{\leftmark}}}
\fancyhead[RE]{\parbox[b]{.84\headwidth}{\raggedleft\footnotesize\nouppercase{\leftmark}}}
\renewcommand{\headrulewidth}{0.25pt}
\fancypagestyle{plain}{%
  \fancyhf{}\fancyfoot[C]{\thepage}\renewcommand{\headrulewidth}{0pt}}

% Jerarquía óptica moderada aprobada para el sucesor y mapa declarativo de
% rótulos públicos. El adaptador del radión aporta sólo sus macros locales;
% no importa el estilo autónomo ni modifica la tipografía general.
% Escala tipográfica aprobada por Rubén el 28 de agosto de 2026.
% Este archivo sustituye exclusivamente la escala de encabezamientos; no
% altera los cuerpos científicos ni sus relaciones de numeración.
\titleformat{\part}[display]
  {\centering\normalfont\color{hmtblue}}
  {\fontsize{12.4}{14.2}\selectfont\mdseries\scshape\partname\ \thepart}
  {2.5pt}
  {\fontsize{16.2}{18.4}\selectfont\bfseries}
\titlespacing*{\part}{0pt}{1.1cm}{.82cm}

\titleformat{\chapter}[display]
  {\normalfont\color{hmtblue}}
  {\fontsize{9.4}{11.2}\selectfont\mdseries\scshape\chaptertitlename\ \thechapter}
  {2.4pt}
  {\fontsize{13.7}{15.7}\selectfont\bfseries\raggedright}
\titlespacing*{\chapter}{0pt}{-.22cm}{.62cm}

\titleformat{name=\chapter,numberless}[block]
  {\normalfont\fontsize{13.7}{15.7}\selectfont\bfseries\color{hmtblue}}
  {}{0pt}{\raggedright}
\titlespacing*{name=\chapter,numberless}{0pt}{-.22cm}{.62cm}

\titleformat{\section}
  {\normalfont\fontsize{11.4}{13.4}\selectfont\bfseries\color{hmtblue}}
  {\thesection}{.65em}{}
\titlespacing*{\section}{0pt}{13pt plus 2pt minus 1pt}{3.5pt}

\titleformat{\subsection}
  {\normalfont\fontsize{10.4}{12.3}\selectfont\bfseries}
  {\thesubsection}{.65em}{}
\titlespacing*{\subsection}{0pt}{8pt plus 1.5pt minus 1pt}{3pt}

\titleformat{\subsubsection}
  {\normalfont\fontsize{9.7}{11.6}\selectfont\itshape}
  {\thesubsubsection}{.65em}{}
\titlespacing*{\subsubsection}{0pt}{6pt plus 1.2pt minus .8pt}{2.5pt}

% El índice reproduce la misma jerarquía descendente y evita que un nivel
% subordinado aparezca ópticamente mayor que su antecedente.
\renewcommand{\cftchappagefont}{\fontsize{10.5}{12.4}\selectfont\bfseries}
\renewcommand{\cftchapfont}{\fontsize{10.5}{12.4}\selectfont\bfseries}
\renewcommand{\cftsecfont}{\fontsize{9.9}{11.8}\selectfont\bfseries}
\renewcommand{\cftsecpagefont}{\fontsize{9.9}{11.8}\selectfont\bfseries}
\renewcommand{\cftsubsecfont}{\fontsize{9.35}{11.2}\selectfont}
\renewcommand{\cftsubsecpagefont}{\fontsize{9.35}{11.2}\selectfont}
\renewcommand{\cftsubsubsecfont}{\fontsize{8.9}{10.8}\selectfont\itshape}
\renewcommand{\cftsubsubsecpagefont}{\fontsize{8.9}{10.8}\selectfont\itshape}
% Los folios de cuatro cifras requieren una caja propia: el margen derecho
% evita que el último término del rótulo colisione con la paginación.
\cftsetpnumwidth{3em}
\cftsetrmarg{3.8em}
 % Mapa declarativo de retitulado aprobado. No altera los cuerpos científicos.
% Cada entrada procede de una coincidencia estática uno-a-uno y conserva el nivel.
\makeatletter
\DeclareUnicodeCharacter{00B2}{\ensuremath{^{2}}}
\DeclareUnicodeCharacter{00B3}{\ensuremath{^{3}}}
\DeclareUnicodeCharacter{00B7}{\ensuremath{\cdot}}
\DeclareUnicodeCharacter{00D7}{\ensuremath{\times}}
\DeclareUnicodeCharacter{0393}{\ensuremath{\Gamma}}
\DeclareUnicodeCharacter{03B1}{\ensuremath{\alpha}}
\DeclareUnicodeCharacter{03B6}{\ensuremath{\zeta}}
\DeclareUnicodeCharacter{03BE}{\ensuremath{\xi}}
\DeclareUnicodeCharacter{2013}{--}
\DeclareUnicodeCharacter{2074}{\ensuremath{^{4}}}
\DeclareUnicodeCharacter{2076}{\ensuremath{^{6}}}
\DeclareUnicodeCharacter{207B}{\ensuremath{^{-}}}
\DeclareUnicodeCharacter{2080}{\ensuremath{_{0}}}
\DeclareUnicodeCharacter{2081}{\ensuremath{_{1}}}
\DeclareUnicodeCharacter{2082}{\ensuremath{_{2}}}
\DeclareUnicodeCharacter{2083}{\ensuremath{_{3}}}
\DeclareUnicodeCharacter{2085}{\ensuremath{_{5}}}
\DeclareUnicodeCharacter{2087}{\ensuremath{_{7}}}
\DeclareUnicodeCharacter{2088}{\ensuremath{_{8}}}
\DeclareUnicodeCharacter{2089}{\ensuremath{_{9}}}
\DeclareUnicodeCharacter{210F}{\ensuremath{\hbar}}
\DeclareUnicodeCharacter{2192}{\ensuremath{\to}}
\DeclareUnicodeCharacter{21A6}{\ensuremath{\mapsto}}
\DeclareUnicodeCharacter{2212}{\ensuremath{-}}
\DeclareUnicodeCharacter{2218}{\ensuremath{\circ}}
\DeclareUnicodeCharacter{221A}{\ensuremath{\surd}}
\DeclareUnicodeCharacter{2260}{\ensuremath{\ne}}
\DeclareUnicodeCharacter{2261}{\ensuremath{\equiv}}
\DeclareUnicodeCharacter{27E8}{\ensuremath{\langle}}
\DeclareUnicodeCharacter{27E9}{\ensuremath{\rangle}}
\newcommand{\HMTTitleMapEntry}[2]{%
  \edef\HMTTitleProbe{\detokenize{#1}}%
  \ifx\HMTTitleKey\HMTTitleProbe\def\HMTResolvedTitle{#2}\fi}
\newcommand{\HMTResolvePublicTitle}[1]{%
  \def\HMTResolvedTitle{#1}%
  \edef\HMTTitleKey{\detokenize{#1}}%
  % El prefacio fue aprobado como texto autoral literal; su subtítulo no se
  % retitula ni se normaliza.
  \HMTTitleMapEntry{Interaction complexe entre deux machines simples}{Interaction complexe entre deux machines simples}% ADP102-0001-LITERAL
  \HMTTitleMapEntry{Introduction générale}{Architecture causale du problème du continu et portée mathématico-physique de HMT–MD}% ADP102-0002
  \HMTTitleMapEntry{Objet initial et ordre causal : échelle finie, horizon d'observation et publication archimédienne}{Horizon projectif : dépendance entre échelle, bord et mesure}% ADP102-0003
  \HMTTitleMapEntry{Intervalle orienté, marques, cellules et arête de clôture}{Construction de la droite finie par l'ordre, l'échelle et la conservation de l'unité}% ADP102-0004
  \HMTTitleMapEntry{Raffinement, survie et seuils de prolongement compatible}{Relation de limite entre infinitésimal, prolongement infini et seuil projectif}% ADP102-0005
  \HMTTitleMapEntry{Évaluation archimédienne des sections compatibles et reconstruction de \(\mathbb R\)}{Limite compatible d'horizons finis et réalisation de la droite réelle}% ADP102-0006
  \HMTTitleMapEntry{Noyau formel de l'holographie modulaire triadique}{Construction et architecture opératorielle du noyau formel APP–TRIT–TPK : support arithmétique, orientation ternaire et dynamique généalogique}% ADP102-0007
  \HMTTitleMapEntry{Relèvements et cocycle de retenue}{Extensions résidu–quotient et cocycle nonadique de retenue}% ADP102-0008
  \HMTTitleMapEntry{Bloc central de \(2\) positions}{Bloc central bivarié d'APP et décomposition en 2 modes spectraux}% ADP102-0009
  \HMTTitleMapEntry{Courbure mixte des \(2\) lois}{Courbure discrète de l'interaction entre les feuillets additif et multiplicatif}% ADP102-0010
  \HMTTitleMapEntry{Correspondance entre \(2\) cartes d’une même classe résiduelle}{Conjugaison entre les cartes positive et résiduelle d'une même classe modulo 9}% ADP102-0011
  \HMTTitleMapEntry{Noyau intérieur et complément gnomonique}{Décomposition de l'orthogramme APP en noyau intérieur et complément gnomonique}% ADP102-0012
  \HMTTitleMapEntry{Segments marqués, interstices et cellule nonadique}{Construction de la cellule nonadique à partir de segments marqués et d'interstices}% ADP102-0013
  \HMTTitleMapEntry{Carte positionnelle finie et projection résiduelle décimale}{Projection résiduelle décimale de la carte positionnelle finie d'APP}% ADP102-0014
  \HMTTitleMapEntry{Extraction de la cellule de \(9\) phases}{Construction de la cellule nonadique à 9 phases par projection résiduelle}% ADP102-0015
  \HMTTitleMapEntry{Position, hauteur et division euclidienne}{Coordonnées de position et de hauteur induites par la division euclidienne}% ADP102-0016
  \HMTTitleMapEntry{Domaine inaugural de la cellule nonadique et exclusion causale de l’APP, du TRIT, du TPK et des constantes analytiques}{Domaine arithmétique autonome de la cellule nonadique et précédence causale par rapport à APP–TRIT–TPK et aux constantes analytiques}% ADP102-0017
  \HMTTitleMapEntry{Loi du défaut à module fixe et longueur variable}{Invariance du défaut cyclotomique à module fixé sous variation de longueur}% ADP102-0018
  \HMTTitleMapEntry{Opérateur cyclotomique primaire}{Opérateur cyclotomique d'APP et action sur les 12 fibres de type A₂}% ADP102-0019
  \HMTTitleMapEntry{Localisation de l’agrégat \texorpdfstring{\(54\)}{54}}{Dérivation du défaut somme–produit 459−405=54 dans le bloc central d'APP}% ADP102-0020
  \HMTTitleMapEntry{Résonance de l’opérateur d’incidence en rang \(6\)}{Rang 6 de l'opérateur d'incidence et multiplicité de son mode résonant}% ADP102-0021
  \HMTTitleMapEntry{TRIT comme unité minimale d’information topologique : orientation ternaire, mémoire de retenue et ordre temporel}{TRIT comme unité ternaire orientée : géométries internes, ordre temporel et conditions dynamiques de cristal temporel}% ADP102-0022
  \HMTTitleMapEntry{Composition et cocycle de retenue}{Cocycle de retenue de la composition ternaire orientée}% ADP102-0023
  \HMTTitleMapEntry{Famille universelle d’algèbres quadratiques}{Classification unifiée des 3 algèbres quadratiques du TRIT}% ADP102-0024
  \HMTTitleMapEntry{Topological Phase Kernel : base observable et état holonomique intégral}{TPK : base observable, état holonomique intégral et séparation typée de ses projections}% ADP102-0025
  \HMTTitleMapEntry{La base double des chemins APP}{Catégorie bivariée des chemins APP comme base de l'état holonomique intégral TPK}% ADP102-0026
  \HMTTitleMapEntry{Phase nonadique, secteur dodécaphasique et calendrier}{Extension non scindée C₁₂→C₁₀₈→C₉ et registre dodécaphasique de la phase nonadique}% ADP102-0027
  \HMTTitleMapEntry{Cylindre et frontière}{Compatibilité cylindrique et signature de frontière de l'état holonomique intégral}% ADP102-0028
  \HMTTitleMapEntry{Transport des fibres}{Opérateur cardinal \(T_d\) de transport entre fibres enrichies avec conservation de la voie et de la mémoire}% ADP102-0029
  \HMTTitleMapEntry{Relations d’extension}{Relation typée \(\operatorname{Ext}_r\) de prolongement entre fibres enrichies}% ADP102-0030
  \HMTTitleMapEntry{Projections et pertes typées}{Quotients observables de l'état holonomique intégral et champs généalogiques conservés}% ADP102-0031
  \HMTTitleMapEntry{Arbre opératoriel du conteneur}{Insertion des 6 familles fibrées de l'état holonomique intégral dans la hiérarchie opératoire du TPK}% ADP102-0032
  \HMTTitleMapEntry{Algèbre d'opérateurs du Topological Phase Kernel et dynamique causale d'actualisation de l'état holonomique intégral}{Hiérarchie typée des 34 constructions canoniques hétérogènes et dynamique causale d'actualisation du TPK}% ADP102-0033
  \HMTTitleMapEntry{Classification de l'état holonomique intégral en \(8\) classes structurelles invariantes}{Hiérarchie typée L₀–L₇ des 34 constructions canoniques hétérogènes du TPK}% ADP102-0034
  \HMTTitleMapEntry{Décomposition fonctionnelle en \(14\) regroupements compatibles avec les opérateurs APP–TRIT–TPK}{Regroupement fonctionnel des 34 constructions canoniques en 14 familles de composition APP–TRIT–TPK}% ADP102-0035
  \HMTTitleMapEntry{Domaine et typage des \(34\) entrées de l’opérateur d’actualisation}{Classes mathématiques, domaines, codomaines et dépendances causales des 34 constructions canoniques}% ADP102-0036
  \HMTTitleMapEntry{Composition causale de l'actualisation de l'état holonomique intégral}{Actualisation de l'état holonomique intégral au moyen de \(U_t=\operatorname{Upd}_t\circ\operatorname{Tra}_t\circ\operatorname{Sel}_t\)}% ADP102-0037
  \HMTTitleMapEntry{Séparation typée entre opérateurs de transport, de mémoire, de calendrier, d’extraction et de retour}{Séparation typée entre opérateurs, états, projections et publications du TPK}% ADP102-0038
  \HMTTitleMapEntry{Décomposition opératorielle des canaux d’incidence \texorpdfstring{\(90/120\)}{90/120} : sélecteur temporel, générateur nilpotent et composante angulaire}{Incidence tritique 90/120 et séparation typée des opérateurs temporel, nilpotent et angulaire}% ADP102-0039
  \HMTTitleMapEntry{Distribution des indices \(90/120\) entre les \(6\) objets du système d’incidence}{Emplacements typés des indices 90/120 dans 6 objets opératoriels non équivalents}% ADP102-0040
  \HMTTitleMapEntry{Défaut normalisé de l’incidence}{Défaut normalisé 1−90/120=1/4 de l'incidence tritique}% ADP102-0041
  \HMTTitleMapEntry{Invariants de typage}{Invariants de domaine, d'orientation et de mémoire des réalisations 90/120}% ADP102-0042
  \HMTTitleMapEntry{État local et rétroaction de feuillet}{Opérateur local de transition orientée et actualisation du feuillet APP}% ADP102-0043
  \HMTTitleMapEntry{Cylindres des \(2\) lecteurs}{Systèmes cylindriques des 2 lecteurs arithmétiques et compatibilité de troncature}% ADP102-0044
  \HMTTitleMapEntry{Décomposition du Topological Phase Kernel en opérateurs d’avancée, de rotation, de retenue et de mémoire}{Factorisation de la transition TPK sous la forme \(U_t=\operatorname{Upd}_t\circ\operatorname{Tra}_t\circ\operatorname{Sel}_t\)}% ADP102-0045
  \HMTTitleMapEntry{Support ordonné et coordonnée cyclique}{Coordonnée cyclique sur le support temporel orienté à 12 secteurs}% ADP102-0046
  \HMTTitleMapEntry{Extension cyclique non scindée}{Extension non scindée 0→C₁₂→C₁₀₈→C₉→0 et cocycle de retenue}% ADP102-0047
  \HMTTitleMapEntry{Extracteur signé d’événements}{Extraction de l'état dodécaphasique signé à partir du registre temporel enrichi}% ADP102-0048
  \HMTTitleMapEntry{\(2\) mots dodécaphasiques de provenances différentes}{Distinction généalogique entre 2 mots dodécaphasiques de même période}% ADP102-0049
  \HMTTitleMapEntry{Raffinement minimal par microinstants}{Subdivision temporelle minimale compatible avec la semi-conjugaison}% ADP102-0050
  \HMTTitleMapEntry{\(3\) opérations quaternaires non interchangeables}{Commutateurs de 3 opérations quaternaires et séparation de leurs actions}% ADP102-0051
  \HMTTitleMapEntry{Lecteur tritique de l’état orienté : ordre temporel et transport de retenue}{Lecteur tritique orienté avec transport de quotient et de retenue}% ADP102-0052
  \HMTTitleMapEntry{Quotient minimal de mémoire prédictive}{Quotient de mémoire d'ordre 1 et suffisance prédictive de la projection modale}% ADP102-0053
  \HMTTitleMapEntry{Mémoire verticale non bornée}{Croissance non bornée de la profondeur de mémoire sous les retours de phase}% ADP102-0054
  \HMTTitleMapEntry{Réalisations géométriques du TRIT : régimes elliptique, parabolique et hyperbolique et double orientation temporelle}{TRIT comme état ternaire orienté : mémoire de retenue, ordre temporel et 3 régimes quadratiques}% ADP102-0055
  \HMTTitleMapEntry{APP, TRIT et TPK : support, orientation et transport}{Composition APP–TRIT–TPK : support arithmétique, orientation ternaire et transport avec mémoire}% ADP102-0056
  \HMTTitleMapEntry{TPK comme système de transport}{TPK comme dynamique de sélection, transport, mémoire, retour et déploiement}% ADP102-0057
  \HMTTitleMapEntry{Périodicité temporelle de (108) phases et lois de retour à (27), (54) et (108) itérations}{Périodicité temporelle de 108 phases et retours à 27, 54 et 108 itérations}% ADP102-0058
  \HMTTitleMapEntry{APP, TRIT et TPK : définitions directrices}{Composition générative APP→TRIT→TPK et construction de l'état holonomique intégral}% ADP102-0059
  \HMTTitleMapEntry{Conséquences et publications du continu construit}{Connexion nonadique conjointe Γ₉ comme correspondance compatible du système inverse}% ADP102-0060
  \HMTTitleMapEntry{Conditionnement comme élagage des futurs}{Conditionnement de la mesure par restriction de l'arbre des prolongements futurs}% ADP102-0061
  \HMTTitleMapEntry{Holonomie nonadique \texorpdfstring{\(\Gamma_9\)}{Γ₉} du Topological Phase Kernel : fibre \texorpdfstring{\(\mathbb F_3^6\)}{F₃⁶}, prolongement coinductif et retour de phase avec mémoire d'état}{Holonomie nonadique Γ₉ du Topological Phase Kernel : transport de l'état holonomique intégral, prolongation coinductive et retour 9→1 avec avance de mémoire}% ADP102-0062
  \HMTTitleMapEntry{Connexion nonadique conjointe, compression et mémoire terminale}{Holonomie conjointe Γ₉ : compression observable T, composante orthogonale η et conservation terminale de mémoire}% ADP102-0063
  \HMTTitleMapEntry{Relations locales et holonomie}{Relations locales de transition et construction de l'holonomie nonadique Γ₉}% ADP102-0064
  \HMTTitleMapEntry{Obstructions de l'holonomie \(\Gamma_9\) : retour \(9\to1\), indices de composition et conservation de la phase, de la retenue, de la frontière, de la route et de la mémoire}{Obstructions de l'holonomie Γ₉ : retour 9→1, indices de composition et conservation de la phase, de la retenue, du bord, de la voie et de la mémoire}% ADP102-0065
  \HMTTitleMapEntry{Décomposition de la microfibre de clôture de \texorpdfstring{\(\pi\)}{π} en \texorpdfstring{\(5\)}{5} régions orientées : multiplicité \texorpdfstring{\(432+4\cdot144=1008\)}{432+4·144=1008}, axe \texorpdfstring{\(\varphi\)}{φ} et involution spéculaire \texorpdfstring{\(\pi/e\)}{π/e}}{Classification orientée de la microfibre de clôture de π : 5 régions, multiplicité 432+4·144=1008, axe φ et involution π/e}% ADP102-0066
  \HMTTitleMapEntry{Génération coinductive de \texorpdfstring{\(\pi\)}{π}, \texorpdfstring{\(e\)}{e} et \texorpdfstring{\(\varphi\)}{φ} par raffinement nonadique compatible}{Génération coinductive de π, e et φ par raffinement nonadique compatible}% ADP102-0067
  \HMTTitleMapEntry{Relèvement visible et chaîne de blocs}{Prolongement coinductif de l'émission exponentielle au moyen de blocs compatibles}% ADP102-0068
  \HMTTitleMapEntry{Signature décimale et carré positionnel}{Construction de la signature décimale au moyen du carré positionnel ternaire–décimal}% ADP102-0069
  \HMTTitleMapEntry{Retour visible et rupture de mémoire}{Retour de phase avec actualisation non triviale de la mémoire de l'état}% ADP102-0070
  \HMTTitleMapEntry{Certificat adversarial de la génération de \(e\) : unicité orbitale, récurrence factorielle, encadrements rationnels et prolongation cylindrique}{Unicité orbitale, récurrence factorielle et convergence cylindrique de la génération de e}% ADP102-0071
  \HMTTitleMapEntry{Relèvement par blocs et 1ʳᵉ frontière coinductive}{Prolongement par blocs et construction du 1er bord coinductif}% ADP102-0072
  \HMTTitleMapEntry{Du calendrier tritique au multigraphe d’incidence}{Construction du multigraphe d'incidence à partir du calendrier tritique}% ADP102-0073
  \HMTTitleMapEntry{Caractérisation, encadrements et prolongement}{Caractérisation de l'auto-échelle au moyen d'encadrements rationnels et du prolongement coinductif}% ADP102-0075
  \HMTTitleMapEntry{Sélection diagonale minimale}{Sélection diagonale minimale de cylindres compatibles et convergence à précision arbitraire}% ADP102-0076
  \HMTTitleMapEntry{Transfert de Markov compatible avec toute la limite projective}{Dynamique de Markov comme quotient modal de mémoire d'ordre 1 de l'holonomie Γ₉, compatible avec la limite inverse}% ADP102-0077
  \HMTTitleMapEntry{2e système de coordonnées : différences de pas \(3\) et \(4\)}{Système différentiel de coordonnées induit par les opérateurs de pas 3 et 4}% ADP102-0078
  \HMTTitleMapEntry{Obstructions de la prolongation affine de \(\alpha\) : télescopage, inversion de \(K\), frontière \(662/663\) et isolement des données cibles}{Obstructions du prolongement affine de α : télescopie, inversion de K, bord 662/663 et isolement des données cibles}% ADP102-0079
  \HMTTitleMapEntry{Construction de \texorpdfstring{\(\sqrt{-1}\)}{√(-1)} au moyen de \texorpdfstring{\(J^2=-I\)}{J²=-I} et coordination des géométries elliptique, parabolique et hyperbolique}{Construction de √(−1) et famille exponentielle tritique d'Euler : opérateurs J_τ et coordination des régimes elliptique, parabolique et hyperbolique}% ADP102-0080
  \HMTTitleMapEntry{Le noyau de transduction de Hensel--Witt}{Noyau de transduction Hensel–Witt : domaine, opération et préservation généalogique}% ADP102-0081
  \HMTTitleMapEntry{Noyau de transduction de Hensel–Witt : microrégions de \(\pi\), holonomie \(\Gamma_9\), registre dodécaphasique et transport hexade–octade}{Noyau de transduction Hensel–Witt : microrégions de π, holonomie Γ₉, registre dodécaphasique et transport hexade–octade}% ADP102-0082
  \HMTTitleMapEntry{Du code ternaire au système de Witt}{Transduction du code ternaire vers le système de Witt}% ADP102-0083
  \HMTTitleMapEntry{Hexade, octade et corridor \texorpdfstring{\(90/120\)}{90/120}}{Inclusion hexade–octade et transport d'incidence par les canaux 90/120}% ADP102-0084
  \HMTTitleMapEntry{Extracteur d'échelle triadique}{Opérateur triadique d'extraction d'échelle et son image rationnelle}% ADP102-0085
  \HMTTitleMapEntry{Défaut des puissances de nombres premiers et opérateur nombre}{Représentation du défaut des puissances premières au moyen de l'opérateur nombre}% ADP102-0086
  \HMTTitleMapEntry{Ablations du sélecteur et de la continuation}{Conditions d'unicité du sélecteur régional et du prolongement analytique}% ADP102-0087
  \HMTTitleMapEntry{Chaîne de génération \(\mathcal F_\pi\to b_\pi\to\rho_\pi\to D_A\to C_*^{\rm HMT}\to\Gamma_{6\to8}\) : sélection régionale, prolongation déterminantale et coordonnée orientée}{Chaîne de génération \(\mathcal F_\pi\to b_\pi\to\rho_\pi\to D_A\to C_*^{\mathrm{HMT}}\to\Gamma_{6\to8}\) : sélection régionale, prolongation déterminantale et coordonnée orientée}% ADP102-0088
  \HMTTitleMapEntry{Décomposition d'incidence hexade--octade dans les secteurs \texorpdfstring{\(q_{90}\)}{q₉₀} et \texorpdfstring{\(q_{120}\)}{q₁₂₀} : opérateur d'entrelacement \texorpdfstring{\(6\to8\)}{6→8} et structure précurseure du spectre d'aire}{Canaux d'incidence 90/120 du Topological Phase Kernel : demi-couronnes tritiques, factorisation des feuillets et réalisation hexade–octade du morphisme 6→8}% ADP102-0089
  \HMTTitleMapEntry{Géométrie dodécaphasique de mélange : antécédent angulaire CKM--CP et lecteur précurseur de l'électron}{Opérateur de mélange dodécaphasique : phase CKM–CP et construction du lecteur prédimensionnel de l'électron}% ADP102-0090
  \HMTTitleMapEntry{Morphismes archimédien et exceptionnel à partir d'un domaine enrichi commun}{Principe d'Ambivalence : involution des lecteurs surfacique et volumique sur l'état holonomique intégral}% ADP102-0091
  \HMTTitleMapEntry{L'opération locale sur les 729 mots}{Domaine enrichi commun et codomaines non équivalents des 2 projections}% ADP102-0092
  \HMTTitleMapEntry{Naturalité et compatibilité à toute profondeur}{Séparation conjointe des histoires par les projections archimédienne et exceptionnelle}% ADP102-0093
  \HMTTitleMapEntry{Couplage postgénératif avec la clôture opératorielle d'Euler}{Invariant d'incidence commun entre l'état holonomique intégral et le système de Witt}% ADP102-0094
  \HMTTitleMapEntry{Globalisation naturelle des projections archimédienne et exceptionnelle sur un système projectif commun}{Décomposition de la microfibre de π en 5 régions orientées : 4 composantes périphériques et 1 axe fixe}% ADP102-0095
  \HMTTitleMapEntry{Systèmes compatibles à profondeur finie}{Prolongement en réseau de l'incidence et construction d'écrans dimensionnels}% ADP102-0096
  \HMTTitleMapEntry{Jauge de la publication exceptionnelle}{Factorisation du lecteur d'incidence à travers l'état holonomique intégral}% ADP102-0097
  \HMTTitleMapEntry{Naturalité conjointe des publications de l'état généalogique}{Double publication de l'état holonomique intégral : évaluation archimédienne et conservation de l'incidence}% ADP102-0098
  \HMTTitleMapEntry{Séparation conjointe de la valeur et de l'incidence exceptionnelle}{Unicité du calibre initial de la projection exceptionnelle}% ADP102-0099
  \HMTTitleMapEntry{Antécédent non commutatif commun de la double projection}{Algèbre non commutative commune et factorisation des 2 projections}% ADP102-0100
  \HMTTitleMapEntry{De l'espace des histoires à la droite réelle}{Quotient de l'espace des histoires compatibles et réalisation de la droite réelle}% ADP102-0101
  \HMTTitleMapEntry{Moyenne archimédienne et double projection du même état}{Compatibilité de l'évaluation archimédienne moyenne avec la double projection de l'état holonomique intégral}% ADP102-0102
  \HMTTitleMapEntry{Calendrier unitaire, mémoire modulaire et mode stable}{Mode stable du calendrier unitaire sous actualisation de mémoire modulaire}% ADP102-0103
  \HMTTitleMapEntry{Publications analytiques d'histoires HMT}{Réalisations analytiques d'histoires HMT au moyen des cartes archimédienne et bidirectionnelle}% ADP102-0104
  \HMTTitleMapEntry{Conséquence pour une action positive}{Positivité de l'action induite par l'opérateur d'entrelacement dodécaphasique–Witt}% ADP102-0105
  \HMTTitleMapEntry{Du saut spectral \texorpdfstring{\(4\)}{4} au réseau de Leech : déploiement APP \texorpdfstring{\(4\to2\to6\to54\)}{4→2→6→54} et corridor Paley--Witt--Mathieu--Golay--Niemeier}{Du saut spectral 4 au réseau de Leech : déploiement APP 4→2→6→54 et corridor Paley–Witt–Mathieu–Golay–Niemeier}% ADP102-0106
  \HMTTitleMapEntry{Nature de la structure obtenue}{Classification algébrique de la structure ternaire autoduale obtenue}% ADP102-0107
  \HMTTitleMapEntry{Loi d'incidence \texorpdfstring{\(4+1\)}{4+1} sur une tétrade}{Décomposition 4+1 de l'incidence d'une tétrade en composantes périphériques et axiale}% ADP102-0108
  \HMTTitleMapEntry{Module et forme discriminants}{Décomposition de l'incidence d'une tétrade en 4 composantes périphériques et 1 composante axiale}% ADP102-0109
  \HMTTitleMapEntry{\(12\) facteurs et application de recollement}{Correspondance entre strates d'incidence et orbite libre du triplet positif}% ADP102-0110
  \HMTTitleMapEntry{Caractère et noyau de congruence}{Recollement discriminant de 12 facteurs A₂ et construction de l'extension entière}% ADP102-0111
  \HMTTitleMapEntry{Courbures mixtes d’APP}{Système ordonné d'incidences et orientation de la chambre radiale}% ADP102-0112
  \HMTTitleMapEntry{Représentation de Weyl--Heisenberg et phase centrale}{Construction du caractère radial et détermination de son noyau de congruence}% ADP102-0113
  \HMTTitleMapEntry{Statut exact de la hiérarchie \(4\to2\to6\to54\) et de l'identification réticulaire.}{Statut exact de la hiérarchie 4→2→6→54 et de l'identification de réseau.}% ADP102-0114
  \HMTTitleMapEntry{Trace graduée du générateur ternaire}{Trace graduée de l'action du générateur ternaire sur \(V^{\natural}\)}% ADP102-0115
  \HMTTitleMapEntry{Mot de support total et isométrie sans points fixes}{Isométrie sans points fixes induite par le mot ternaire de support total}% ADP102-0116
  \HMTTitleMapEntry{Construction de l'espace de Hilbert à partir de la connexion nonadique et de ses opérateurs de phase}{Réalisation projective de Bloch à partir de la structure complexe interne}% ADP102-0117
  \HMTTitleMapEntry{La branche marquée change de type}{Immersion par coins d'une phase marquée : idéal compact \(\mathcal K(\mathcal H_\infty)\) et facteur \(\mathcal B(\mathcal H_\infty)\) de type \(\mathrm I_\infty\)}% ADP102-0118
  \HMTTitleMapEntry{Reste visible et quotient de frontière}{Extension opératorielle d'APP : modes visibles et modes de mémoire}% ADP102-0119
  \HMTTitleMapEntry{Involution de dualité T}{Complétion 1000=729+243+27+1 et stratification de l'unité}% ADP102-0120
  \HMTTitleMapEntry{Complétion stratifiée de l'unité}{Décomposition 243=1+22+90+120+10 et typage de ses composantes}% ADP102-0121
  \HMTTitleMapEntry{Le bloc \texorpdfstring{\(243\)}{243} comme frontière et boule ternaire parfaite}{Dérivation des canaux de constantes à partir de la décomposition du bloc 243}% ADP102-0122
  \HMTTitleMapEntry{Décomposition interne du bloc \texorpdfstring{\(243\)}{243}}{Construction de la dimension 11 au moyen de lecteurs résiduels compatibles}% ADP102-0123
  \HMTTitleMapEntry{Dimension \texorpdfstring{\(11\)}{11} et lectures compatibles}{Clôture du registre généalogique sous l'involution de dualité T}% ADP102-0124
  \HMTTitleMapEntry{Identités vérifiables de la construction}{Identités exactes d'involution, de modularité et de conservation résiduelle}% ADP102-0125
  \HMTTitleMapEntry{Niveau, échange linéaire de masses et modularité : \(3\) contrôles de non-fusion}{Séparation par niveau, échange linéaire des masses et modularité des 3 réalisations}% ADP102-0126
  \HMTTitleMapEntry{\(2\) descendants du code ternaire étendu}{Réduction dimensionnelle 26→24→11→10 et coordination avec les écrans de théorie M}% ADP102-0127
  \HMTTitleMapEntry{La représentation icosaédrique de dimension \(12\)}{Coordination des écrans exacts avec la réalisation physique de théorie M}% ADP102-0128
  \HMTTitleMapEntry{Réduction isotrope de rang \(26\to24\)}{Représentation opératorielle d'APP et triplet spectral supersymétrique en dimension 11}% ADP102-0129
  \HMTTitleMapEntry{Coordination du corridor Moonshine avec les écrans de théorie M : graduation, réseau lorentzien et double réalisation de l'état HMT}{Théorème de compatibilité entre le corridor Moonshine et les écrans de théorie M par graduation et transport de réseau}% ADP102-0130
  \HMTTitleMapEntry{Antécédent commun des lecteurs : raffinement spectral et limite de dilatation conjointe}{Raffinement spectral commun et limite de la dilatation conjointe des lecteurs matriciels}% ADP102-0131
  \HMTTitleMapEntry{Sous-structures, objet global et projections}{Construction de 4 carrés latins quantiques cycliques dans une classe opératorielle élémentaire}% ADP102-0132
  \HMTTitleMapEntry{Correspondances directes et inverses entre les formalismes}{Réalisation de 3 contextes opératoriels dans un carré quantique (9,3)}% ADP102-0133
  \HMTTitleMapEntry{Extension de la bigraduation à 3 indices}{Représentation de l'atlas généalogique complet dans un carré quantique (12,3)}% ADP102-0134
  \HMTTitleMapEntry{4 carrés latins quantiques cycliques de classe opératorielle facile}{Classification cyclique de 4 carrés latins quantiques dans une classe opératorielle élémentaire}% ADP102-0135
  \HMTTitleMapEntry{Lemme cyclique des carrés magiques quantiques associés à une mesure opératorielle positive}{Bigraduation de la famille généalogico–matricielle}% ADP102-0136
  \HMTTitleMapEntry{Cônes matriciels, images et extensions complètement positives}{Représentation matricielle de la composition généalogique APP–TRIT–TPK}% ADP102-0137
  \HMTTitleMapEntry{Construction explicite de cubes magiques quantiques et compatibilité de leurs marges}{TPK comme catégorie de voies avec sélection, transport, mémoire et retour}% ADP102-0138
  \HMTTitleMapEntry{Combinaisons magiques d'applications CP}{Compositions complètement positives compatibles avec la bigraduation généalogico–matricielle}% ADP102-0139
  \HMTTitleMapEntry{Hiérarchie spectrale \(4\to2\to6\to54\) du bloc central d'APP}{Réalisation métrologique ultérieure de la double projection de l'état holonomique intégral}% ADP102-0140
  \HMTTitleMapEntry{Raffinement nonadique par espérance conditionnelle}{Décomposition symplectique 12→4×3 des coordonnées de phase}% ADP102-0141
  \HMTTitleMapEntry{Taxonomie fonctionnelle de \(114\) expressions numériques et de \(106\) lectures opératorielles}{Classification opératorielle de 114 expressions numériques par domaine, générateur, invariant et codomaine}% ADP102-0142
  \HMTTitleMapEntry{Transfert de phase sur le multigraphe des blocs complets}{Microfibre de clôture de π : 5 régions orientées, multiplicités \(432+4\cdot144=1008\) et projection commune \(w_\pi\)}% ADP102-0143
  \HMTTitleMapEntry{\(5\) cycles orientés de clôture de la microfibre}{Publication de mémoire réduite d'un bloc nonadique complet au moyen du multigraphe de transfert de phase}% ADP102-0144
  \HMTTitleMapEntry{Opérateur d’incidence et demi-tour orienté}{Holonomies orientées de clôture des 5 régions de π}% ADP102-0145
  \HMTTitleMapEntry{Composition rationnelle et compensateur \texorpdfstring{\(239\)}{239}}{Détermination de la pente 1/5 par la multiplicité de la microfibre régionale de π}% ADP102-0146
  \HMTTitleMapEntry{Identité gaussienne et \(1^{\mathrm a}\) demi-tour orienté}{Composition rationnelle du demi-tour et détermination du terme compensateur 239}% ADP102-0147
  \HMTTitleMapEntry{Prolongement coinductif de la section de \(\pi\) par le système inverse des histoires survivantes}{Filtration de F₃⁶ par admissibilité régionale et sélection de sous-cylindres compatibles}% ADP102-0148
  \HMTTitleMapEntry{Système inverse des histoires compatibles}{Unicité de la section diagonale minimale dans le système d'encadrements rationnels}% ADP102-0149
  \HMTTitleMapEntry{Signature positionnelle décimale et carré de compatibilité}{Prolongement de la section exponentielle au moyen de la chaîne de blocs visibles}% ADP102-0150
  \HMTTitleMapEntry{Retour de phase avec changement d’état et conservation de la mémoire parcourue}{Construction de la signature décimale et du carré positionnel de la section exponentielle}% ADP102-0151
  \HMTTitleMapEntry{Semi-groupe factoriel du caractère exponentiel : récurrence \(a_n=1/n!\)}{Retour de phase avec inégalité des états holonomiques intégraux : \(g(K+9)=g(K)\) et \(B_{K+9}\ne B_K\)}% ADP102-0152
  \HMTTitleMapEntry{Certificat interne d’unicité orbitale, de récurrence factorielle, de convergence et de prolongement cylindrique}{Prolongement coinductif de e au moyen de l'holonomie nonadique Γ₉ du TPK}% ADP102-0153
  \HMTTitleMapEntry{Demi-droite positive invariante et unicité de l’auto-échelle}{Construction du multigraphe d'incidence modale à partir du calendrier tritique}% ADP102-0154
  \HMTTitleMapEntry{Certificat interne d’unicité, de compatibilité et de convergence de l’auto-échelle de Perron}{Caractérisation de φ par le rayon de Perron, l'emboîtement des encadrements et le prolongement coinductif}% ADP102-0156
  \HMTTitleMapEntry{Dérivations internes de la constante de structure fine \(\alpha\) : clôture entière dodécaphasique réversible et racine analytique unique du noyau \(E_{21/22}\)}{Clôture dodécaphasique de \(\alpha\) et caractérisation analytique au moyen de \(E_{21/22}\) : génération réversible, racine simple et prolongement nonadique d'ordre 9}% ADP102-0157
  \HMTTitleMapEntry{État terminal commun et 2 lectures internes coordonnées}{Évaluations terminale et dodécaphasique de l'état holonomique intégral commun : blocs \(B_{\mathrm{term}}\) et carte signée \(U_{12}^{\mathrm{sgn}}\)}% ADP102-0158
  \HMTTitleMapEntry{Voie A : clôture entière dodécaphasique}{Génération réversible de \(\alpha\) au moyen de la clôture dodécaphasique entière}% ADP102-0159
  \HMTTitleMapEntry{Domaine Hensel--Witt, applications de Teichmüller et cible rationnelle \(-1/12\)}{Opérateur Hensel–Witt de transduction entre microrégions de π, holonomie Γ₉, registre dodécaphasique et incidence hexade–octade}% ADP102-0160
  \HMTTitleMapEntry{Application de Teichmüller du code ternaire à l’anneau de Witt}{Construction du système de Witt à partir du codage ternaire compatible}% ADP102-0161
  \HMTTitleMapEntry{Incidence hexade--octade et canaux modulaires \(90/120\)}{Réalisation hexade–octade des canaux 90/120 précédemment dérivés de l'incidence tritique du TPK}% ADP102-0162
  \HMTTitleMapEntry{4 routes compatibles vers \texorpdfstring{\(-1/12\)}{-1/12}}{4 réalisations opératorielles de −1/12 et séparation de leurs domaines}% ADP102-0163
  \HMTTitleMapEntry{Dérivation déterminantale de l’échelle d’action : forme normale de Smith, conoyau d’ordre \(16\) et valuation \(10^{-34}\)}{Dérivation déterminantale de l'échelle absolue d'action : forme normale de Smith, conoyau d'ordre 16 et valuation 10⁻³⁴}% ADP102-0164
  \HMTTitleMapEntry{Mesure de quotient et caractère positif de l’action : hilbertisation canonique, défaut régional--dodécaphasique et loi fonctorielle sur les chemins}{Construction de la mesure positive d'action sur le quotient des émissions : hilbertisation, défaut régional–dodécaphasique et fonctorialité des chemins}% ADP102-0165
  \HMTTitleMapEntry{Dérivation autoduale de la constante de Planck : périodicité temporelle de \(108\) phases, action et longueur de Compton}{Dérivation autoduale de la constante de Planck et de sa droite d'action : périodicité temporelle de 108 phases, fermeture de Compton et covariance bidirectionnelle}% ADP102-0166
  \HMTTitleMapEntry{Équivalence des normalisations sur l’identité autoduale}{Commutativité des normalisations angulaire, temporelle et gravitationnelle de l'identité autoduale}% ADP102-0167
  \HMTTitleMapEntry{Opérateur angulaire prédimensionnel : décomposition en \(27\) secteurs, rayon spectral et composition effective de la double projection}{Opérateur de phase antérieur à la transduction dimensionnelle : décomposition en 27 secteurs, rayon spectral et factorisation de la double projection}% ADP102-0168
  \HMTTitleMapEntry{Caractères quadratiques et constantes spectrales : Catalan, Apéry et tour d’Euler--Stieltjes}{Caractères quadratiques et constantes spectrales : Catalan, Apéry, Euler–Mascheroni et hiérarchie de Stieltjes}% ADP102-0169
  \HMTTitleMapEntry{Criticité unimodale de la famille dodécaphasique : profil analytique exact, approximants de Feigenbaum et réduction de l’universalité au certificat hyperbolique}{Criticité unimodale de la famille dodécaphasique : profil analytique exact, approximants de Feigenbaum et certificat hyperbolique de renormalisation}% ADP102-0170
  \HMTTitleMapEntry{Arithmétique du mode \(30\) : identités de Rogers--Ramanujan, normalisateur \(899\) et relèvement nonadique de Pell}{Arithmétique de la famille modulaire dodécaphasique de période 30 : identités de Rogers–Ramanujan, normalisateur 899 et relèvement nonadique de Pell}% ADP102-0171
  \HMTTitleMapEntry{Théorème directeur du réseau radical et de la loi puissance--lacune}{Classification du réseau radical 3→6→9 et dérivation de la loi puissance–lacune}% ADP102-0172
  \HMTTitleMapEntry{Orbite cyclique \(369\to693\to936\to369\)}{Orbite cyclique 369→693→936→369}% ADP102-0173
  \HMTTitleMapEntry{Noyau \(7227\), décomposition spectrale et lien avec la croix radicale}{Détermination du noyau arithmétique 7227 par incidence avec la croix radicale d'APP}% ADP102-0174
  \HMTTitleMapEntry{Encodage mécanique des lacunes par les puissances nonadiques}{Codage sturmien des lacunes induit par les puissances nonadiques}% ADP102-0175
  \HMTTitleMapEntry{Racine interne de \(-1\), structure exponentielle et géométries elliptique, parabolique et hyperbolique}{Famille exponentielle triadique \(J_\tau^2=-\tau I\) : réalisations elliptique, parabolique et hyperbolique de l'identité d'Euler}% ADP102-0176
  \HMTTitleMapEntry{Réalisation continue du type quadratique induit}{Extension continue de l'opérateur quadratique commun et conservation de son type spectral}% ADP102-0177
  \HMTTitleMapEntry{Support opératoriel commun de \(3\) lectures géométriques}{Conjugaison des réalisations arithmétique, angulaire et torsionnelle au moyen de l'opérateur quadratique commun}% ADP102-0178
  \HMTTitleMapEntry{Chaîne constitutive du secteur gravitationnel : échelle chronologique \(108\), réduction tensorielle, holonomie et évaluation métrologique de \(G\)}{Composition causale des invariants de phase, d'action et de longueur dans la dérivation de G}% ADP102-0179
  \HMTTitleMapEntry{Cardinalités d’incidence 90 et 120}{Dérivation tritique des cardinalités 90/120 et réalisation ultérieure dans l'incidence hexade–octade}% ADP102-0180
  \HMTTitleMapEntry{Séparation causale entre le spectre d’aire, la mémoire modulaire et la réalisation Cartan--Holst}{Distinction typée entre Barbero–aire, aire–hamiltonien modulaire et réalisation Cartan–Holst}% ADP102-0182
  \HMTTitleMapEntry{Transducteur de Hensel non filtré et décalage complet sur les fibres ternaires}{Transducteur entier de Hensel comme opérateur de décalage complet}% ADP102-0183
  \HMTTitleMapEntry{Surjectivité de la carte archimédienne non filtrée sur \(\mathbb R\)}{Couverture exacte de la carte réelle antécédente au moyen de sections compatibles}% ADP102-0184
  \HMTTitleMapEntry{2 axes de classification : valeur visible et généalogie de transport}{Classification bifactorielle du nombre par valeur archimédienne et généalogie de transport}% ADP102-0185
  \HMTTitleMapEntry{Position de \texorpdfstring{\(\pi,e,\varphi\) et \(\alpha\)}{pi, e, phi et alpha} dans la classification}{Position de π,e,φ et α dans la classification}% ADP102-0186
  \HMTTitleMapEntry{Carte de Hensel non filtrée, survie et frontière globale}{Carte antécédente, survie coinductive et construction du bord global}% ADP102-0187
  \HMTTitleMapEntry{La famille opératorielle d'Euler comme grammaire de clôture}{Factorisation de la fermeture généalogique du nombre au moyen du système opératoriel d'Euler}% ADP102-0188
  \HMTTitleMapEntry{Dépendances mathématiques.}{Dépendance causale entre section inverse, évaluation archimédienne et conservation exceptionnelle de l'incidence}% ADP102-0189
  \HMTTitleMapEntry{\(4\) exigences de choix dans le système projectif}{Conditions de choix nécessaires au prolongement global des sections compatibles}% ADP102-0190
  \HMTTitleMapEntry{Décidabilité effective de chaque fenêtre finie de multiplicités projectives}{Décidabilité algorithmique de chaque niveau fini du système de préfixes}% ADP102-0191
  \HMTTitleMapEntry{Relation exacte avec les \(9\) phases du Topological Phase Kernel}{Restriction de l'obstruction uniforme au système temporel nonadique à 9 phases}% ADP102-0192
  \HMTTitleMapEntry{De la cohérence constructible à l'indépendance par forcing}{Cohérence relative dans L et indépendance par forcing}% ADP102-0193
  \HMTTitleMapEntry{Optimisation minimax sur des fenêtres finies de préfixes compatibles}{Théorème minimax à chaque niveau fini de l'ensemble partiellement ordonné des préfixes compatibles}% ADP102-0194
  \HMTTitleMapEntry{La chaîne que le pont reçoit démontrée}{Antécédents théorématiques du foncteur de transduction généalogico–métrologique}% ADP102-0195
  \HMTTitleMapEntry{Dérivation bidirectionnelle de l'exposant électronique à partir du parcours enrichi}{Bande de réalisation de l'électron et séparation entre masse basale et clôture bêta}% ADP102-0196
  \HMTTitleMapEntry{\(2\) branches causales et leur confluence}{Confluence des branches déterminantale et dodécaphasique de l'échelle d'action}% ADP102-0197
  \HMTTitleMapEntry{Carte métrologique ultérieure et évaluation dimensionnelle des invariants engendrés}{Systèmes de transduction entre nombre généalogique et grandeur physique}% ADP102-0198
  \HMTTitleMapEntry{Principe d'ambivalence : involution des lecteurs surfacique et volumétrique, équivalence locale et orientation temporelle conjuguée}{Principe d'ambivalence en mécanique dimensionnelle : involution surfacique–volumétrique, systèmes auto-inertiels et orientation temporelle conjuguée}% ADP102-0199
  \HMTTitleMapEntry{Conditions d'une évaluation prospective}{Conditions d'évaluation prospective du principe d'Ambivalence sur l'état holonomique intégral}% ADP102-0200
  \HMTTitleMapEntry{Cristal temporel fini HMT : forçage, sous-harmonique et stabilité spectrale}{Cristal temporel fini HMT : hamiltonien périodiquement forcé, réponse sous-harmonique et stabilité spectrale}% ADP102-0201
  \HMTTitleMapEntry{Forçage construit à partir de l’horloge à 108 états}{Construction du forçage périodique à partir du système temporel à 108 états}% ADP102-0202
  \HMTTitleMapEntry{Stabilité spectrale finie}{Jet nilpotent de l'évolution triadique et prolongement compatible}% ADP102-0203
  \HMTTitleMapEntry{Théorème de confluence HMT–MD : forme électrogravitationnelle, système d’Einstein–Schrödinger et coordination de l’action, de la matière, des champs, de la géométrie et de l’information}{Théorème de confluence HMT–MD : système Einstein–Schrödinger et coordination électrogravitationnelle de l'action, de la matière, des champs, de la géométrie et de l'information}% ADP102-0204
  \HMTTitleMapEntry{Indépendance horizontale et compatibilité verticale}{Naturalité des réalisations et conservation verticale des invariants de l'état holonomique intégral}% ADP102-0205
  \HMTTitleMapEntry{1re confluence : couplage, vide et secteur électrofaible}{Couplage constitutif entre le vide et le secteur électrofaible}% ADP102-0206
  \HMTTitleMapEntry{2e confluence : action, gravité et forme électrogravitationnelle}{Réciprocité entre action et gravité dans la forme électrogravitationnelle}% ADP102-0207
  \HMTTitleMapEntry{3e confluence : géométrie, information et mélange}{Confluence entre géométrie, information et mélange}% ADP102-0208
  \HMTTitleMapEntry{4e confluence : écran, torsion et mémoire hélicoïdale}{Transduction de la mémoire hélicoïdale dans la réalisation torsionnelle}% ADP102-0209
  \HMTTitleMapEntry{Du déterminant électrofaible à l’opérateur constitutif positif du vide et à sa complétion spectrale \(3\to4\)}{Du déterminant électrofaible à l'opérateur constitutif positif du vide et à sa complétion spectrale 3→4}% ADP102-0210
  \HMTTitleMapEntry{Valeur publiée et mémoire de la fibre}{Désintégration spectrale de la mesure des trajectoires et conditionnement des futurs}% ADP102-0211
  \HMTTitleMapEntry{Fondements algébriques de la couture de phase et de l'automesure}{État conjoint observateur–observable, instrument réversible et automesure}% ADP102-0212
  \HMTTitleMapEntry{Information paire et transport géométrique impair dans la bicouche}{Décomposition paire–impaire de l'état bicouche : information interfeuillet et transport géométrique orienté}% ADP102-0214
  \HMTTitleMapEntry{\texorpdfstring{\(3\)}{3} réalisations de la mémoire généalogique : transport dodécaphasique, fonction de Green non locale avec générateur local et expansion cosmologique}{Réalisations dynamique, propagative et cosmologique de la mémoire généalogique : transport dodécaphasique, fonction de Green et expansion spatiale}% ADP102-0215
  \HMTTitleMapEntry{Chaîne constitutive du secteur gravitationnel : échelle chronologique \texorpdfstring{\(108\)}{108}, réduction tensorielle, holonomie et évaluation métrologique de \texorpdfstring{\(G\)}{G}}{Descente icosaédrique A₅/C₅ et sélection du porteur gravitationnel}% ADP102-0216
  \HMTTitleMapEntry{Entrelaceur avec les tenseurs symétriques sans trace et réduction transverse \texorpdfstring{\(12\to5\to5\to2\)}{12→5→5→2}}{Opérateur d'entrelacement avec les tenseurs symétriques sans trace et réduction transverse 12→5→5→2}% ADP102-0217
  \HMTTitleMapEntry{Critères structurels du sélecteur gravitationnel : transducteur dimensionnel, holonomie enrichie et séparation entre sélecteurs torsionnel et circulaire}{Sélecteur holonomique gravitationnel : transduction dimensionnelle, naturalité, raffinement et séparation des régimes torsionnel et circulaire}% ADP102-0218
  \HMTTitleMapEntry{Carte décimale déclarée et comparaison métrologique}{Reconnaissance métrologique ultérieure de G et indépendance causale de la dérivation}% ADP102-0219
  \HMTTitleMapEntry{Tétrade, connexion et \(2\) défauts}{Tétrade, connexion de spin et torsion de la réalisation Cartan–Holst}% ADP102-0220
  \HMTTitleMapEntry{Action du 1er ordre}{Action de Palatini–Holst et équations variationnelles}% ADP102-0221
  \HMTTitleMapEntry{Confluence horizon–aire–information–entropie–action}{Théorème de confluence entre énergie d'horizon, aire, information, entropie et action}% ADP102-0222
  \HMTTitleMapEntry{\(2\) orientations d’une même propagation}{Propagateurs avancé et retardé sur des voies orientées}% ADP102-0223
  \HMTTitleMapEntry{Orbite elliptique, action et \(3\) invariants J}{Invariants de l'involution temporelle et conjugaison particule–antiparticule}% ADP102-0224
  \HMTTitleMapEntry{Cosmologie de périodicité \texorpdfstring{\(108\)}{108} : inversion de feuillet, horizon, Big Bang, rebond torsionnel et secteur sombre}{Cosmologie torsionnelle de périodicité 108 : expansion, retour avec mémoire et rebond d'Einstein–Cartan}% ADP102-0225
  \HMTTitleMapEntry{Séparation typologique des \(2\) réalisations cosmologiques}{Séparation entre accélération de de Sitter et rebond torsionnel d'Einstein–Cartan}% ADP102-0226
  \HMTTitleMapEntry{\(3\) feuilletages exacts de de Sitter}{Feuilletages cosmologiques compatibles avec l'inversion de feuillet}% ADP102-0227
  \HMTTitleMapEntry{Univers à frontière}{État de bord de l'horizon cosmologique et mémoire de retour}% ADP102-0228
  \HMTTitleMapEntry{Données spectrales d'une réalisation supersymétrique : \((\mathcal H_{\rm phys},H,\mathcal Q)\), supercharge et compactification en dimension \(11\)}{Réalisation de théorie M par compactification, cordes, branes et dimension 11}% ADP102-0229
  \HMTTitleMapEntry{Loi générale des masses sur l'atlas \texorpdfstring{\(81/56/13/324\)}{81/56/13/324} : trajectoires, registre, signature, caractère, tours \texorpdfstring{\(90/120\)}{90/120} et secteur nul}{Loi générale des masses sur l'atlas 81/56/13/324 : trajectoires, registre, signature, caractère, semi-groupe harmonique ⟨90,120⟩ et secteur nul}% ADP102-0230
  \HMTTitleMapEntry{Algorithme universel de réalisation matérielle et égalisateur spectral des lecteurs \(K_D\) et \(K_\Omega\)}{Construction fonctorielle des familles de masse sur l'atlas}% ADP102-0231
  \HMTTitleMapEntry{Bifurcation nulle et positivité}{Bifurcation du secteur nul antérieure à la branche positive de masse}% ADP102-0232
  \HMTTitleMapEntry{Accès au domaine exhaustif}{Construction exhaustive de l'atlas 81/56/13/324 et conservation du secteur nul}% ADP102-0233
  \HMTTitleMapEntry{Domaine interne et problème de réalisation}{Théorème d'exhaustivité du domaine matériel et stabilité du secteur nul sous raffinement}% ADP102-0234
  \HMTTitleMapEntry{Produit fibré entre descripteurs et routes}{Produit fibré entre invariants d'espèce et voies généalogiques}% ADP102-0235
  \HMTTitleMapEntry{Lecteurs bidirectionnels \texorpdfstring{\(K_D\)}{K D} et \texorpdfstring{\(K_\Omega\)}{K Omega} : profils, coïncidences, tours \texorpdfstring{\(90/120\)}{90/120} et opérateurs par multisection}{Lecteurs bidirectionnels \(K_D\) et \(K_\Omega\) : profils, coïncidences, semi-groupe harmonique \(\langle90,120\rangle\) et opérateurs par multisection}% ADP102-0236
  \HMTTitleMapEntry{Lecteurs \(K_D\) et \(K_\Omega\) : profils dodécaphasiques, domaine des parcours et codomaine spectral}{Lecteurs \(K_D\) et \(K_\Omega\) : profils dodécaphasiques, domaine des parcours et codomaine spectral}% ADP102-0237
  \HMTTitleMapEntry{Réalisations spectrales, scalaires et nulles de la Loi générale des masses : familles matérielles, transport CKM--PMNS et secteurs photonique et gluonique}{Réalisations spectrales, scalaires et nulles de la Loi générale des masses : familles matérielles, mélange quark CKM, mélange neutrinique PMNS et secteurs nuls photonique et gluonique}% ADP102-0238
  \HMTTitleMapEntry{Couleur, saveurs de quarks et mélange CKM}{Opérateurs de masse par fibre, orbite et multisection de l'atlas}% ADP102-0239
  \HMTTitleMapEntry{États composés et loi de liaison}{Correspondance par domaine et codomaine entre l'inventaire matériel et les réalisations de l'atlas}% ADP102-0240
  \HMTTitleMapEntry{Classificateur des espèces matérielles par fibre, représentation, charge, spin, saveur et génération}{Classification des espèces matérielles par fibre, représentation, charge, spin, saveur et génération}% ADP102-0241
  \HMTTitleMapEntry{Typage probatoire de la persistance, de l'atlas et des familles.}{Conditions de persistance et correspondance entre atlas, familles et représentations}% ADP102-0242
  \HMTTitleMapEntry{Ontologie matérielle, charge, spin, statistique, couleur et saveur}{Définitions structurelles de particule, masse, charge, spin, statistique, couleur, saveur et génération}% ADP102-0243
  \HMTTitleMapEntry{Théorème de clôture typée de la réalisation massique}{Théorème de clôture de la réalisation matérielle au moyen de sections persistantes et de représentations finies}% ADP102-0244
  \HMTTitleMapEntry{Confrontation métrologique et archive exhaustive des réalisations}{Reconnaissance métrologique ultérieure et inventaire exhaustif des réalisations matérielles}% ADP102-0245
  \HMTTitleMapEntry{Architecture naturelle commune des réalisations terminales du continu généalogique}{Théorème de factorisation conjointe des \(5\) réalisations terminales à travers la structure discrète du continu}% ADP102-0246
  \HMTTitleMapEntry{Publication terminale du continu unique par \(5\) lecteurs non permutables}{Publication terminale au moyen de 5 lecteurs non permutables et facteur commun de l'état holonomique intégral}% ADP102-0247
  \HMTTitleMapEntry{Fermetures dynamiques du continu : transport solénoïdal et quotient de jauge}{Dynamique et courbure du continu généalogique antérieures aux réalisations classiques}% ADP102-0248
  \HMTTitleMapEntry{Dynamique et courbure du continu généalogique}{Factorisation conjointe des réalisations terminales à partir du même état du continu}% ADP102-0249
  \HMTTitleMapEntry{Fermetures algébriques du continu et limites de la reconstruction extensionnelle}{Factorisation commune, cohérence extensionnelle et limites de la reconstruction terminale}% ADP102-0250
  \HMTTitleMapEntry{L’exemple des ordinaux de von Neumann}{Représentation ordinale de von Neumann comme modèle de cohérence extensionnelle}% ADP102-0251
  \HMTTitleMapEntry{Domaine, applications et compatibilités de l’objet terminal commun}{Vue interne \(G_T^{\mathrm{int}}\) du continu et application \(\operatorname{App}_{T,d_T}\) : conservation de la fibre, de la retenue, de la mémoire et du terminal}% ADP102-0252
  \HMTTitleMapEntry{Déploiement interne des réalisations corrélatives avant leur identification aux problèmes classiques}{Émergence corrélative des 5 actions terminales sur le même état du continu}% ADP102-0253
  \HMTTitleMapEntry{\(5\) lecteurs terminaux non permutables du continu généalogique}{Théorème de non-permutabilité des 5 lecteurs terminaux et classification de leurs codomaines}% ADP102-0254
  \HMTTitleMapEntry{Obstructions à la projection ablative du continu : perte de phase, d'orientation ou de mémoire généalogique}{Projecteurs positifs de la forme de Guinand–Weil et localisation spectrale}% ADP102-0255
  \HMTTitleMapEntry{Fixation de la droite critique par l’involution fonctionnelle \(s\mapsto1-\bar s\)}{Fixation de la droite critique par l'involution fonctionnelle \texorpdfstring{\(s\mapsto 1-\bar{s}\)}{s -> 1 - conjugué(s)}}% ADP102-0256
  \HMTTitleMapEntry{Distinction typée entre les zéros de \(\zeta\), de \(\xi\) et du déterminant spectral}{Séparation typée des zéros de ζ, ξ et du déterminant spectral}% ADP102-0257
  \HMTTitleMapEntry{Représentant positif de la classe résiduelle nulle : \(n\equiv0\pmod 9\) et \(\rho_9(n)=9\)}{Représentant positif de la classe résiduelle nulle : \(n\equiv0\pmod 9\) et \(\rho_9(n)=9\)}% ADP102-0258
  \HMTTitleMapEntry{Isométrie de transport tordu \(\widetilde U_k\) sur les états cylindriques, les poids projectifs et la mémoire de Weil}{Isométrie de transport tordu \(\widetilde U_k\) sur les états cylindriques, les poids projectifs et la mémoire de Weil}% ADP102-0259
  \HMTTitleMapEntry{Invariant gradué du corridor spectral \(4\to2\to6\to54\)}{Invariant gradué du corridor spectral 4→2→6→54}% ADP102-0260
  \HMTTitleMapEntry{Trace couplée du spectre de Weil et de la graduation Moonshine}{Construction indépendante du champ spectral sur le double cercle}% ADP102-0261
  \HMTTitleMapEntry{Domaine hydrodynamique et publication de Galerkin}{Domaine solénoïdal et conservation de la divergence nulle}% ADP102-0262
  \HMTTitleMapEntry{Donnée unique de départ et noyaux denses}{Données initiales, compatibilité de bord et propagation régulière}% ADP102-0263
  \HMTTitleMapEntry{Réalisation finie de l'inégalité de Kalman–Yakubovich–Popov et fonction auxiliaire dans la clôture énergétique}{Inégalité KYP et contrôle coercif de l'évolution solénoïdale}% ADP102-0264
  \HMTTitleMapEntry{Extraction matérielle de la borne de Sobolev}{Estimation de Sobolev et exclusion de la concentration singulière}% ADP102-0265
  \HMTTitleMapEntry{La fibre d'incidence est \(1\) coordonnée physique de l'état complet}{Courbure de fibre et coercivité de l'action de jauge}% ADP102-0266
  \HMTTitleMapEntry{Vide simple et témoin qui atteint l'écart spectral}{Unicité du vide invariant de jauge et saturation de l'écart \(3\kappa_{\mathrm{av}}\) par le vecteur propre \(W_c\)}% ADP102-0267
  \HMTTitleMapEntry{\(4\) reconstructions d'Osterwalder–Schrader d'une dynamique commune}{Positivité d'Osterwalder–Schrader et reconstruction quantique}% ADP102-0268
  \HMTTitleMapEntry{Théorie de jauge euclidienne et gap de Yang--Mills}{Écart spectral du secteur de jauge et limite projective}% ADP102-0269
  \HMTTitleMapEntry{Présentation générateur par générateur de la nouvelle carte}{Générateurs cohomologiques et projecteurs compatibles}% ADP102-0270
  \HMTTitleMapEntry{Fibre initiale, Mittag--Leffler et publication affirmative}{Lefschetz–Murre et descente algébrique des classes de Hodge}% ADP102-0271
  \HMTTitleMapEntry{Bocksteins de toutes les pages}{Homomorphismes de Bockstein de la suite spectrale : compatibilité entre pages et détermination de l'ordre en T du complexe universel}% ADP102-0272
  \HMTTitleMapEntry{Table complète Kato--FERS avant le déterminant}{Déterminant de Kato–FERS, rang et terme principal de BSD}% ADP102-0273
  \HMTTitleMapEntry{Distinction entre structure et application}{Vue interne \(G_T\) et lecteur terminal \(P_T\) : indépendance logique entre existence structurelle et réalisation classique}% ADP102-0274
  \HMTTitleMapEntry{Obstructions communes aux \(5\) expressions terminales : perte de domaine, de naturalité, de certificat ou de limite compatible}{Obstructions communes des lecteurs terminaux : perte de domaine, de naturalité ou de limite compatible}% ADP102-0275
  \HMTTitleMapEntry{Formalisation assistée du noyau algébrique}{Formalisation vérifiable du noyau algébrique et certification d'identités finies}% ADP102-0276
  \HMTTitleMapEntry{Correspondance entre énoncés, preuves, certificats et sources théorématiques scientifiques}{Correspondance entre énoncés, démonstrations, certificats et sources scientifiques primaires}% ADP102-0277
  \HMTTitleMapEntry{Conventions de typage et de représentation des données}{Conventions mathématiques et représentation des données}% ADP102-0278
  \HMTTitleMapEntry{Hétérogénéité typée et couverture du catalogue}{Hétérogénéité des domaines physiques et exhaustivité du catalogue}% ADP102-0279
  \HMTTitleMapEntry{Fiches individuelles des chemins orientés}{Déploiement intégral des 324 voies orientées : correspondance entre 53 champs canoniques, 57 champs enrichis et 49 invariants publiés par section}% ADP102-0280
  \HMTTitleMapEntry{Inventaire universel typé de 471 masses et 384 largeurs}{Inventaire exhaustif de 471 masses et 384 largeurs : identification, provenance et correspondance avec l'atlas}% ADP102-0281
  \HMTTitleMapEntry{Source, déduplication et limites de type}{Provenance des données, unicité des enregistrements et domaines de validité}% ADP102-0282
  \HMTTitleMapEntry{Code d’objet et lecteur HMT--MD}{Identification des objets et application du lecteur HMT--MD}% ADP102-0283
  % Correspondencias espirales: sustituciones científicas unívocas aprobadas
  % por cotejo focal. Se modifican sólo los rótulos; los cuerpos permanecen
  % íntegros en sus residencias distribuidas.
  \HMTTitleMapEntry{TPK : sélection, transport, actualisation et retour}{Composition dynamique du Topological Phase Kernel sur l'état holonomique intégral : sélection tritique, transport cardinal, actualisation de la mémoire et retour typé}% ADP102-0284
  \HMTTitleMapEntry{La transition typée}{Définition et factorisation de l'opérateur de transition enrichie \(U_t\)}% ADP102-0285
  \HMTTitleMapEntry{Composition de la retenue}{Loi de cocycle bimodulaire de la retenue sous concaténation des voies}% ADP102-0286
  \HMTTitleMapEntry{Calendrier dodécaphasique et non-scindement}{Extension cyclique non scindée \(0\to C_{12}\to C_{108}\to C_9\to0\) et chronologie interne du TPK}% ADP102-0287
  \HMTTitleMapEntry{Hélice nonadique de phase et de profondeur}{Revêtement hélicoïdal de la phase nonadique et coordonnée entière de profondeur de mémoire}% ADP102-0288
  \HMTTitleMapEntry{Immersion géométrique normalisée}{Immersion \(H_9\) du quotient--résidu nonadique dans \(\mathbb R^3\) et calcul des invariants de Frenet}% ADP102-0289
  \HMTTitleMapEntry{Cercles emboîtés et cylindres}{Sections circulaires et feuilletage cylindrique du revêtement phase--mémoire}% ADP102-0290
  \HMTTitleMapEntry{Feuillets orbitaux et couture de l'espace de suspension}{Orbites du produit croisé et relation d'identification de l'espace de suspension}% ADP102-0291
  \HMTTitleMapEntry{Cas stationnaire affine}{Itération de l'endomorphisme affine stationnaire et classification de ses points fixes}% ADP102-0292
  \HMTTitleMapEntry{Inversion des voies}{Anti-involution du sous-groupoïde des voies réversibles et conjugaison de l'holonomie affine}% ADP102-0293
  \HMTTitleMapEntry{Miroir des constantes et axe fixe}{Involution \(\pi/e\) de la fibre des constantes, invariant \(u^\pi+u^e\) et axe fixe \(\varphi\)}% ADP102-0294
  \HMTTitleMapEntry{La géométrie visible comme publication de l'état holonomique intégral}{Évaluation archimédienne de l'état holonomique intégral et défaut de fidélité généalogique}% ADP102-0295
  \HMTTitleMapEntry{L'inversion de l'ordre explicatif}{Factorisation causale de l'histoire enrichie à l'évaluation archimédienne}% ADP102-0296
  \HMTTitleMapEntry{État, histoire et publication}{Espace des états holonomiques intégraux, catégorie des histoires TPK et évaluation archimédienne}% ADP102-0297
  \HMTTitleMapEntry{Cercle intrinsèque et cercle projeté}{Distinction entre orbite radiale intrinsèque et projection circulaire d'une trajectoire hélicoïdale}% ADP102-0298
  \HMTTitleMapEntry{La cellule nonadique et la naissance bidimensionnelle}{Construction du support nonadique bidimensionnel \(X_9=(\mathbb Z/9\mathbb Z)^2\)}% ADP102-0299
  \HMTTitleMapEntry{Antécédents APP de la classification radiale}{Invariants APP qui induisent le caractère radial accumulé}% ADP102-0300
  \HMTTitleMapEntry{État ternaire relevé}{État TRIT enrichi par résidu et retenue : multiplicité des représentants du seuil}% ADP102-0301
  \HMTTitleMapEntry{Algèbres quadratiques}{Famille \(\mathcal A_\tau=\mathbb R[J]/(J^2+\tau I)\) et classification elliptique, parabolique et hyperbolique}% ADP102-0302
  \HMTTitleMapEntry{Double projection}{Factorisation conjointe des publications archimédienne et exceptionnelle de l'état holonomique intégral}% ADP102-0303
  \HMTTitleMapEntry{Représentation locale d'une transition}{Représentation affine \(e\mapsto h_e\) des générateurs de la catégorie des voies}% ADP102-0304
  \HMTTitleMapEntry{Holonomies préfixes}{Produits affines ordonnés \texorpdfstring{\(H_j=h_{e_j}\cdots h_{e_1}\)}{H_j = h_ej ... h_e1} et conservation des préfixes}% ADP102-0305
  \HMTTitleMapEntry{Catégorie relevée et caractères élémentaires}{Catégorie des voies enrichies par 2 orientations TRIT et définition des caractères \(p\) et \(a\)}% ADP102-0306
  \HMTTitleMapEntry{Décomposition équilibrée}{Décomposition unique \(p=r_3(p)+3c_3(p)\) en résidu équilibré et retenue}% ADP102-0307
  \HMTTitleMapEntry{Involution}{Action de \((\tau_+,\tau_-)\mapsto(-\tau_-,-\tau_+)\) sur les caractères \(p\) et \(a\)}% ADP102-0308
  \HMTTitleMapEntry{Définition et domaine}{Définition de \(L_p\) sur \(\mathbb R_{>0}\) et domaine de l'inverse \(\exp_p\)}% ADP102-0309
  \HMTTitleMapEntry{Addition déformée}{Loi de composition \(L_p(xy)=L_p(x)+L_p(y)+pL_p(x)L_p(y)\)}% ADP102-0310
  \HMTTitleMapEntry{Non-exhaustivité et portée exacte}{Domaine classificatoire de \(p\in\{2,1,0,-1,-2\}\) et extension affine générale}% ADP102-0311
  \HMTTitleMapEntry{Obstruction au lecteur réduit}{Défaut de naturalité du lecteur réduit \((p,a,\Lambda,C,m,r,\theta)\)}% ADP102-0312
  \HMTTitleMapEntry{Définition du lecteur complet}{Lecteur radial enrichi \(\mathcal P^{\mathrm{enr}}_{\mathrm{rad},N}\) : domaine, codomaine et conservation des préfixes}% ADP102-0313
  \HMTTitleMapEntry{Conservation comme traçabilité}{Conservation dynamique de la retenue, du bord et de la mémoire par le lecteur radial}% ADP102-0314
  \HMTTitleMapEntry{Subdivision quadruple}{Naturalité du lecteur sous le raffinement quadruple \(\operatorname{sd}_4\)}% ADP102-0315
  \HMTTitleMapEntry{Classification généalogique et non-fidélité de la publication}{Classification par la signature radiale pleine et défaut de fidélité de l'évaluation archimédienne}% ADP102-0316
  \HMTTitleMapEntry{Classification par 3 incréments visibles}{Classification cinématique au moyen des incréments \((\omega,\beta,\mu)\)}% ADP102-0317
  \HMTTitleMapEntry{3 fibres d'une même évaluation visible}{\(3\) classes généalogiques ayant une même image archimédienne}% ADP102-0318
  \HMTTitleMapEntry{Loi de vis}{Identité spectrale \(V^{-1}\mathscr S V=(2+\sqrt3)\mathrm i\mathscr S\) et équation de la suspension logarithmique}% ADP102-0319
  \HMTTitleMapEntry{Mode tritique commun}{Mode singulier commun des feuillets APP dans la carte canonique \(1,\ldots,9\)}% ADP102-0320
  \HMTTitleMapEntry{Ledger material}{Registre généalogique de la voie matérielle : extension survivante, signature entière et caractère adimensionnel}% ADP102-0321
  \HMTTitleMapEntry{Publication dimensionnelle ultérieure}{Réalisation de l'opérateur constitutif du vide dans les champs \((E_\sigma,B_\sigma)\)}% ADP102-0322
  \HMTTitleMapEntry{Courbe plane polaire}{Courbure et longueur d'arc d'une courbe polaire \(r=r(\theta)\)}% ADP102-0323
  \HMTTitleMapEntry{Spécialisations}{Évaluation de \(\kappa_{\mathrm F}\) dans les familles logarithmique, archimédienne, de Fermat, hyperbolique et lituus}% ADP102-0324
  \HMTTitleMapEntry{Suspension spatiale}{Immersion spatiale \(\Gamma(\theta)\) et critères différentiels de Frenet--Lancret}% ADP102-0325
  \HMTTitleMapEntry{Trichromie géométrique}{Coordination des régimes \(J^2=-I,0,+I\) avec la géométrie de Frenet}% ADP102-0326
  \HMTTitleMapEntry{Conclusion}{Théorème de fermeture des correspondances spirales généalogiques}% ADP102-0327
  \HMTTitleMapEntry{Produit semi-direct affine et composition des voies}{Monoïde affine \(\operatorname{End}(U_{\mathrm{APP}})\ltimes U_{\mathrm{APP}}\) et composition ordonnée des voies}% ADP102-0328
  \HMTTitleMapEntry{Tableau exhaustif de la carte locale à 2 orientations}{Énumération exhaustive de la carte locale \((\tau_+,\tau_-)\mapsto(p,a,r_3,c_3)\)}% ADP102-0329
  \HMTTitleMapEntry{5 régions de \texorpdfstring{\(\pi\)}{pi} et non-fidélité de l'évaluation}{Décomposition de la microfibre de \(\pi\) en 5 régions et multiplicité de l'évaluation archimédienne}% ADP102-0330
  \HMTTitleMapEntry{Algèbre puissance--logarithmique et 5 publications canoniques}{Algèbre de \(L_p\) et classification des 5 familles canoniques \(p\in\{2,1,0,-1,-2\}\)}% ADP102-0331
  \HMTTitleMapEntry{Dérivation des 5 courbes canoniques}{Évaluation de \(\exp_p(u_0+\beta m)\) dans les 5 familles canoniques}% ADP102-0332
  \HMTTitleMapEntry{Inversion des feuillets}{Conjugaison \(p\leftrightarrow-p\) des feuillets puissance--logarithmiques}% ADP102-0333
  \HMTTitleMapEntry{Secteur affine général}{Classification du secteur affine par la paire \((p,[\Lambda,C])\)}% ADP102-0334
  \HMTTitleMapEntry{Naturalité, équivariance et non-fidélité}{Naturalité sous troncature, équivariance nonadique et défaut de fidélité archimédienne}% ADP102-0335
  \HMTTitleMapEntry{Le lecteur de préfixes}{Lecteur de préfixes affines et conservation du registre enrichi}% ADP102-0336
  \HMTTitleMapEntry{Raffinements non triviaux}{Critère de compatibilité pour les raffinements généraux de voies}% ADP102-0337
  \HMTTitleMapEntry{Équivariance, non invariance}{Équivariance du cocycle radial sous l'holonomie \(\Gamma_9\)}% ADP102-0338
  \HMTTitleMapEntry{3 démonstrations de non-fidélité}{\(3\) obstructions explicites à la fidélité de l'évaluation archimédienne}% ADP102-0339
  \HMTTitleMapEntry{Même extrémité, ordre affine différent}{Coïncidence du produit final et dépendance de l'ordre des facteurs affines}% ADP102-0340
  \HMTTitleMapEntry{Même courbe, région généalogique différente}{Coïncidence de l'image archimédienne entre régions généalogiques distinctes}% ADP102-0341
  \HMTTitleMapEntry{Classification par niveau mathématique}{Stratification par domaine : histoire, lecteur, évaluation archimédienne et réalisation physique}% ADP102-0342
  \HMTTitleMapEntry{Classification par rotation, rayon et mémoire}{Classification par les incréments \((\omega,\beta,\mu)\)}% ADP102-0343
  \HMTTitleMapEntry{Réalisations ultérieures et contrôles de type}{Réalisations électromagnétique et matérielle avec séparation des domaines et codomaines}% ADP102-0344
  \HMTTitleMapEntry{Centre électronique et lecteur ultérieur}{Compatibilité du lecteur radial avec le point fixe électronique \((5,5)\) et l'atlas \(81/56/13/324\)}% ADP102-0345
  \HMTTitleMapEntry{Matrice fermée de contrôles négatifs}{Matrice des conditions de validité par niveau causal}% ADP102-0346
  \HMTTitleMapEntry{Première prolongation et frontière coinductive}{Prolongement coinductif initial de \(w_6\) jusqu'au bord de compatibilité d'ordre \(1\)}% ADP102-0347
  \HMTTitleMapEntry{Voie B : noyau de vacances \texorpdfstring{\(E_{21/22}\)}{E21/22}}{Caractérisation de \(\alpha\) comme racine simple du noyau analytique \(E_{21/22}\)}% ADP102-0348
  % Purismo residual: rótulos activos localizados en el índice del sucesor.
  % Las sustituciones siguientes no alteran el nivel jerárquico ni el cuerpo.
  \HMTTitleMapEntry{Chronologie observable des deux curseurs}{Évolution chronologique de la paire de curseurs et détermination de l'état observable}% ADP102-0349
  \HMTTitleMapEntry{Émetteur fini à deux curseurs}{Émetteur fini associé à la paire de curseurs et génération d'états admissibles}% ADP102-0350
  \HMTTitleMapEntry{D'une cellule à deux coordonnées : naissance d'APP}{Construction d'APP par décomposition de 1 cellule en 2 coordonnées}% ADP102-0351
  \HMTTitleMapEntry{La branche de \(20\) blocs}{Prolongement de la section survivante au moyen de 20 blocs compatibles}% ADP102-0352
  \HMTTitleMapEntry{Du préfixe à un point réel}{Évaluation archimédienne du préfixe compatible et détermination du point réel}% ADP102-0353
  \HMTTitleMapEntry{Une émission reproductible de douze triades}{Émission reproductible de 12 triades et reconstruction de leur état holonomique intégral}% ADP102-0354
  \HMTTitleMapEntry{Construction du caractère de clôture et des trajectoires généalogiques structurelles de \(\pi\)}{Construction du caractère de clôture et des trajectoires généalogiques de \(\pi\)}% ADP102-0355
  \HMTTitleMapEntry{La pente 1/5 imposée par les \(5\) régions}{Détermination de la pente 1/5 par le bilan des 5 régions orientées}% ADP102-0356
  \HMTTitleMapEntry{Fenêtre dodécaphasique de l'état signé et clôture réversible du cycle}{Restriction finie de l'état dodécaphasique signé et clôture réversible du cycle}% ADP102-0357
  \HMTTitleMapEntry{Fenêtre dodécaphasique finie}{Domaine fini de la clôture dodécaphasique}% ADP102-0358
  \HMTTitleMapEntry{Prolongement distant \texorpdfstring{\(\mathsf T_{\alpha,\infty}\)}{T alpha infini} de la voie entière}{Prolongement coinductif \(\mathsf T_{\alpha,\infty}\) de la clôture dodécaphasique entière}% ADP102-0359
  \HMTTitleMapEntry{Théorème des quatre canaux entiers}{Théorème de classification des 4 canaux entiers}% ADP102-0360
  \HMTTitleMapEntry{Premier encodeur intrinsèque et sa limite}{Codeur intrinsèque initial et existence de sa limite}% ADP102-0361
  \HMTTitleMapEntry{Hypothèses de typage et contrôles négatifs}{Conditions de typage et critères d'incompatibilité de la théorie des vacances nonadiques}% ADP102-0362
  \HMTTitleMapEntry{Voie chronologique de longueur, d'action et de spectre}{Composition chronologique de longueur, action et spectre dans la détermination de \(G\)}% ADP102-0363
  \HMTTitleMapEntry{Opérateur tordu et contrôle négatif déterministe}{Opérateur avec torsion et obstruction déterministe à la continuation multiplicative}% ADP102-0364
  \HMTTitleMapEntry{Stratification mathématique de l’objet de Majorana en quatre niveaux}{Stratification mathématique de l'objet Majorana en 4 niveaux}% ADP102-0365
  \HMTTitleMapEntry{Six caractères et quatre torsions de tresse}{Classification de 6 caractères par 4 classes de torsion de tresse}% ADP102-0366
  \HMTTitleMapEntry{obstructions et conditions spécifiques de non-dégénérescence}{Obstructions spécifiques et conditions de non-dégénérescence}% ADP102-0367
  \HMTTitleMapEntry{Quark haut et antiquark haut}{Secteurs quark \(u\) et antiquark \(\bar u\)}% ADP102-0368
  \HMTTitleMapEntry{Quark bas et antiquark bas}{Secteurs quark \(d\) et antiquark \(\bar d\)}% ADP102-0369
  \HMTTitleMapEntry{Contrôles de réfutation.}{Conditions de validité et d'incompatibilité de la réalisation équivariante des espèces}% ADP102-0370
  \HMTTitleMapEntry{Le carré local de \texorpdfstring{\(e\)}{e} et la rupture de mémoire}{Carré positionnel local de \(e\) et différenciation généalogique d'états ayant le même bloc visible}% ADP102-0371
  \HMTTitleMapEntry{La loi générale avant ses réalisations par espèce}{Antécédence causale de la Loi Générale des Masses par rapport aux réalisations par espèce}% ADP102-0372
  \HMTTitleMapEntry{Fondement unifié de la masse}{Factorisation extension–route et conservation du registre généalogique}% ADP102-0373
  \HMTTitleMapEntry{La bande de parcours}{Définition de la bande de route par l'orbite d'une section persistante}% ADP102-0374
  \HMTTitleMapEntry{Le canal d'incidence est une courbure locale, non un facteur auxiliaire}{Interprétation du canal d'incidence comme courbure locale du système de jauge}% ADP102-0375
  % Publicaciones terminales del capítulo 95: se conserva cada desarrollo
  % completo y se sustituye la enumeración narrativa por objeto, operación y resultado.
  \HMTTitleMapEntry{1re résolution : Riemann et la positivité comme complément construit}{Publication spectrale–polaire du continu par le complément positif de Guinand–Weil}% ADP102-0376
  \HMTTitleMapEntry{Existence, régularité et unicité globales pour Navier–Stokes par borne uniforme et compacité}{Publication solénoïdale du continu : borne coercive, compacité et régularité globale}% ADP102-0377
  \HMTTitleMapEntry{3e résolution : Yang–Mills et le saut du secteur de courbure}{Publication de jauge du continu : positivité par réflexion et saut du secteur de courbure}% ADP102-0378
  \HMTTitleMapEntry{4e résolution : Hodge et construction effective de la préimage algébrique}{Publication cohomologique du continu : relèvement compatible et construction effective de la préimage algébrique}% ADP102-0379
  \HMTTitleMapEntry{5e résolution : BSD et coïncidence des expressions analytique et arithmétique}{Publication déterminantielle du continu : comparaison Kato–Poitou–Tate et identité du terme principal de Birch–Swinnerton–Dyer}% ADP102-0380
  \HMTTitleMapEntry{L’identité de diffusion qui conserve la puissance croisée}{Identité de dispersion et conservation de la puissance croisée}% ADP102-0381
  \HMTTitleMapEntry{Persistance de l’objet terminal sous le passage à la limite}{Persistance de l'état terminal solénoïdal lors du passage à la limite}% ADP102-0382
  \HMTTitleMapEntry{Existence, régularité et unicité globale de la solution de Navier–Stokes obtenue par borne uniforme et compacité}{Théorème de clôture solénoïdale globale par borne uniforme et compacité}% ADP102-0383
  \HMTTitleMapEntry{L’opérateur universel d’extension}{Opérateur universel d'extension des classes cohomologiques compatibles}% ADP102-0384
  \HMTTitleMapEntry{Factorisation des entrelaceurs terminaux RH–NS–YM–Hodge–BSD à travers le continu généalogique commun}{Pertes d'information spécifiques aux 5 publications classiques et conservation dans l'état holonomique intégral}% ADP102-0385
  \HMTTitleMapEntry{Factorisation de RH–NS–YM–Hodge–BSD comme expressions terminales du continu généalogique commun}{Théorème de factorisation terminale de RH–NS–YM–Hodge–BSD à travers le continu généalogique commun}% ADP102-0386
  \HMTTitleMapEntry{Émergence corrélative de \texorpdfstring{\(5\)}{5} actions intrinsèques non exhaustives}{Décomposition expositive de l'action corrélative en 5 vues intrinsèques du même état}% ADP102-0387
  \HMTTitleMapEntry{Un même bord : flux intérieur et tension de jauge}{Factorisation commune de l'opérateur de frontière dans les réalisations solénoïdale et de jauge}% ADP102-0388
  \HMTTitleMapEntry{L'état commun de bord}{État holonomique intégral de frontière et décomposition paritaire du transport}% ADP102-0389
  \HMTTitleMapEntry{L'opérateur minimal et son spectre}{Complexe de bord minimal et équivalence spectrale des laplaciens de degrés 0 et 1}% ADP102-0390
  \HMTTitleMapEntry{Lecture hydrodynamique}{Réalisation solénoïdale de l'opérateur de frontière : contraction du flot et conservation de l'énergie}% ADP102-0391
  \HMTTitleMapEntry{Lecture Yang--Mills}{Réalisation de jauge de l'opérateur de frontière : coercivité, quotient de jauge et saut spectral}% ADP102-0392
  % Promoción jerárquica de los certificados causales de \(\Gamma_9\) y de
  % la generación trítica de los canales enteros.
  \HMTTitleMapEntry{Filtration monodromique et calcul à profondeur prescrite}{Certification finie de Γ₉ : 396 niveaux, 2.376 trits, 43 tours complets, 387 paires \((K,K+9)\) par canal et fibre ambiante de 729 éléments}% ADP102-0393
  \HMTTitleMapEntry{Recensement fini des transitions et génération décimale}{Recensement exhaustif de 396 transitions, 2.376 trits, 43 tours et 387 paires de même phase}% ADP102-0394
  \HMTTitleMapEntry{Structure entière des lecteurs de constantes}{Classification arithmétique des canaux entiers et compatibilité avec la demi-couronne tritique}% ADP102-0395
}
\let\HMTNativeChapter\chapter
\let\HMTNativeSection\section
\let\HMTNativeSubsection\subsection
\let\HMTNativeSubsubsection\subsubsection
\let\HMTNativeParagraph\paragraph
\let\HMTNativeSubparagraph\subparagraph
\RenewDocumentCommand{\chapter}{s o m}{%
  \HMTResolvePublicTitle{#3}%
  \IfBooleanTF{#1}{\HMTNativeChapter*{\HMTResolvedTitle}}{%
    \IfNoValueTF{#2}{\HMTNativeChapter{\HMTResolvedTitle}}{\HMTNativeChapter[\HMTResolvedTitle]{\HMTResolvedTitle}}}}
\RenewDocumentCommand{\section}{s o m}{%
  \HMTResolvePublicTitle{#3}%
  \IfBooleanTF{#1}{\HMTNativeSection*{\HMTResolvedTitle}}{%
    \IfNoValueTF{#2}{\HMTNativeSection{\HMTResolvedTitle}}{\HMTNativeSection[\HMTResolvedTitle]{\HMTResolvedTitle}}}}
\RenewDocumentCommand{\subsection}{s o m}{%
  \HMTResolvePublicTitle{#3}%
  \IfBooleanTF{#1}{\HMTNativeSubsection*{\HMTResolvedTitle}}{%
    \IfNoValueTF{#2}{\HMTNativeSubsection{\HMTResolvedTitle}}{\HMTNativeSubsection[\HMTResolvedTitle]{\HMTResolvedTitle}}}}
\RenewDocumentCommand{\subsubsection}{s o m}{%
  \HMTResolvePublicTitle{#3}%
  \IfBooleanTF{#1}{\HMTNativeSubsubsection*{\HMTResolvedTitle}}{%
    \IfNoValueTF{#2}{\HMTNativeSubsubsection{\HMTResolvedTitle}}{\HMTNativeSubsubsection[\HMTResolvedTitle]{\HMTResolvedTitle}}}}
\RenewDocumentCommand{\paragraph}{s o m}{%
  \HMTResolvePublicTitle{#3}%
  \IfBooleanTF{#1}{\HMTNativeParagraph*{\HMTResolvedTitle}}{%
    \IfNoValueTF{#2}{\HMTNativeParagraph{\HMTResolvedTitle}}{\HMTNativeParagraph[\HMTResolvedTitle]{\HMTResolvedTitle}}}}
\RenewDocumentCommand{\subparagraph}{s o m}{%
  \HMTResolvePublicTitle{#3}%
  \IfBooleanTF{#1}{\HMTNativeSubparagraph*{\HMTResolvedTitle}}{%
    \IfNoValueTF{#2}{\HMTNativeSubparagraph{\HMTResolvedTitle}}{\HMTNativeSubparagraph[\HMTResolvedTitle]{\HMTResolvedTitle}}}}
% Las utilidades de degradación posteriores deben conservar el mapa de títulos.
\let\HMTBookChapter\chapter
\let\HMTBookSection\section
\let\HMTBookSubsection\subsection
\let\HMTBookSubsubsection\subsubsection
\let\HMTBookParagraph\paragraph
\let\HMTBookSubparagraph\subparagraph
\makeatother
 % Adaptador de compatibilidad del fuente teoremática íntegro del radión.
%
% El tratado rector ya carga tipografía, geometría, hipervínculos, teoremas,
% tablas, TikZ y tcolorbox.  Este archivo aporta únicamente los nombres y
% entornos exclusivos del fuente teoremática, sin importar su preámbulo autónomo ni
% alterar la jerarquía tipográfica general.

\makeatletter
\@ifundefined{color@RadNavy}{\definecolor{RadNavy}{HTML}{102A43}}{}
\@ifundefined{color@RadBlue}{\definecolor{RadBlue}{HTML}{1F5E7A}}{}
\@ifundefined{color@RadTeal}{\definecolor{RadTeal}{HTML}{168C8C}}{}
\@ifundefined{color@RadCopper}{\definecolor{RadCopper}{HTML}{B66A3C}}{}
\@ifundefined{color@RadGold}{\definecolor{RadGold}{HTML}{D4A73A}}{}
\@ifundefined{color@RadInk}{\definecolor{RadInk}{HTML}{1A202C}}{}
\@ifundefined{color@RadMist}{\definecolor{RadMist}{HTML}{EDF4F7}}{}
\@ifundefined{color@RadCream}{\definecolor{RadCream}{HTML}{FAF7F0}}{}
\@ifundefined{color@RadRose}{\definecolor{RadRose}{HTML}{8D3F55}}{}
\@ifundefined{color@RadGreen}{\definecolor{RadGreen}{HTML}{2F7D5D}}{}
\@ifundefined{color@RadGray}{\definecolor{RadGray}{HTML}{66788A}}{}

\providecommand{\TRIT}{\ensuremath{\mathrm{TRIT}}}
\providecommand{\HMT}{\ensuremath{\mathrm{HMT}}}
\providecommand{\Rstate}{\mathcal R_{12}}
\providecommand{\epsR}{\varepsilon_R}
\providecommand{\epsone}{\varepsilon_R^{(1)}}
\providecommand{\pih}{\pi}
\providecommand{\alphah}{\alpha}
\providecommand{\hret}{\hbar_{\mathrm{ret}}}
\providecommand{\hfull}{h_{\mathrm{ret}}}
\providecommand{\diff}{\mathop{}\!\mathrm d}
\providecommand{\e}{\mathrm e}
\providecommand{\R}{\mathbb R}
\providecommand{\C}{\mathbb C}
\providecommand{\F}{\mathbb F}
\providecommand{\Z}{\mathbb Z}
\providecommand{\dd}{\,\mathrm d}
\providecommand{\status}[1]{\textsf{\bfseries #1}}
\providecommand{\sourcepath}[1]{\nolinkurl{#1}}

\@ifundefined{typebox}{%
  \newtcolorbox{typebox}[1][]{enhanced,breakable,colback=white,
    colframe=RadBlue,boxrule=.7pt,arc=1.5mm,left=4mm,right=4mm,
    top=3mm,bottom=3mm,title={Reçu de types},fonttitle=\bfseries,
    coltitle=RadNavy,#1}%
}{}
\@ifundefined{narrativebox}{%
  \newtcolorbox{narrativebox}[1][]{enhanced,breakable,colback=white,
    colframe=RadGold,borderline west={2.4pt}{0pt}{RadGold},
    boxrule=.25pt,arc=0mm,left=5mm,right=4mm,top=3mm,bottom=3mm,#1}%
}{}

% Las once portadillas internas del fuente teoremática autónomo se convierten en
% divisiones científicas ordinarias dentro del capítulo 52.  Se evita así
% introducir partes autónomas, gigantismo tipográfico o páginas vacías.
\providecommand{\partopener}[3]{\section{#2}\noindent\emph{#3}\par\medskip}

\providecommand{\BigRadionEquation}{%
  \begin{equation*}
  \boxed{\displaystyle
  \frac{180^{\circ}}{\pih}
  =50^{\circ}+A^{\circ}-\frac{20}{81}\,\Delta_4^{\circ}+\epsR^{\circ}}
  \end{equation*}}
\makeatother
 % Adaptador mínimo para las correspondencias espirales genealógicas.
% El preámbulo autónomo no se importa: sólo se registran símbolos y colores
% que los cuerpos científicos usan de manera efectiva.
\makeatletter
\@ifundefined{color@HMTNavy}{\definecolor{HMTNavy}{HTML}{17324D}}{}
\@ifundefined{color@HMTBlue}{\definecolor{HMTBlue}{HTML}{315E7D}}{}
\@ifundefined{color@HMTRed}{\definecolor{HMTRed}{HTML}{8A3D36}}{}
\@ifundefined{color@HMTGold}{\definecolor{HMTGold}{HTML}{9B7B3E}}{}
\@ifundefined{color@HMTGray}{\definecolor{HMTGray}{HTML}{5F6870}}{}
\@ifundefined{color@HMTLight}{\definecolor{HMTLight}{HTML}{F2F4F5}}{}
\@ifundefined{color@HMTLine}{\definecolor{HMTLine}{HTML}{C9CFD3}}{}
\providecommand{\enr}{\mathrm{enr}}
\providecommand{\AffEnd}{\operatorname{AffEnd}}
\providecommand{\Hist}{\operatorname{Hist}}
\providecommand{\Led}{\operatorname{Led}}
\providecommand{\Carr}{\operatorname{Carr}}
\providecommand{\ev}{\operatorname{ev}}
\providecommand{\sd}{\operatorname{sd}}
\providecommand{\id}{\operatorname{id}}
\providecommand{\im}{\operatorname{im}}
\providecommand{\Spec}{\operatorname{Spec}}
\providecommand{\arcosh}{\operatorname{arcosh}}
\providecommand{\N}{\mathbb N}
\providecommand{\T}{\mathbb T}
\providecommand{\ii}{\mathrm i}
\providecommand{\twoway}{\textit{bidirectionnel}}
\providecommand{\ph}{\mathrm{ph}}
\providecommand{\mem}{\mathrm{mem}}
\providecommand{\rad}{\mathrm{rad}}
\providecommand{\arq}{\mathrm{arq}}
\providecommand{\e}{\mathrm e}
\providecommand{\F}{\mathbb F}
\providecommand{\Z}{\mathbb Z}
\providecommand{\R}{\mathbb R}
\providecommand{\dd}{\,\mathrm d}

% Los fuentes teoremáticas de correspondencias espirales se distribuyen por
% residencia causal. Sus etiquetas conservan un espacio de nombres único y
% sus divisiones autónomas se subordinan a la jerarquía del capítulo receptor.
\newenvironment{HMTSpiralIntegratedBody}{%
  \begingroup
  % Las copias públicas ya llevan el prefijo material ``espiral:'' en cada
  % etiqueta y referencia. Esto incluye los labels diferidos de amsmath, que
  % no podían asegurarse mediante una redefinición local de \label.
  \let\part\HMTBookSection
  \let\chapter\HMTBookSubsection
  \let\section\HMTBookSubsubsection
  \let\subsection\HMTBookParagraph
  \let\subsubsection\HMTBookSubparagraph
}{\endgroup}

\newenvironment{HMTSpiralChapterBody}{%
  \begingroup
  \let\chapter\HMTBookSection
  \let\section\HMTBookSubsection
  \let\subsection\HMTBookSubsubsection
  \let\subsubsection\HMTBookParagraph
}{\endgroup}

\newenvironment{HMTSpiralAppendixBody}{%
  \begingroup
  \let\part\HMTBookChapter
  \let\chapter\HMTBookSection
  \let\section\HMTBookSubsection
  \let\subsection\HMTBookSubsubsection
  \let\subsubsection\HMTBookParagraph
}{\endgroup}
\makeatother
 
% Identificador estable para builds reproducibles.
\ifdefined\pdfinfoomitdate
  \pdfinfoomitdate=1
\fi
\ifdefined\pdfsuppressptexinfo
  \pdfsuppressptexinfo=15
\fi
\ifdefined\pdftrailerid
  \pdftrailerid{<554E49444144484F4C4F47524146494341>
                <44494D454E53494F4E414C484D54323026>}
\fi

\makeatletter\begin{document}
\frontmatter
\hypersetup{pageanchor=false}
\begin{titlepage}
\thispagestyle{empty}
\centering
\vspace*{1.8cm}
{\color{hmtgold}\rule{0.68\textwidth}{0.8pt}\par}
\vspace{1.35cm}
{\fontsize{31}{37}\selectfont\color{hmtblue}\bfseries
Double projection holographique\par
d'un cristal temporel apériodique\par}
\vspace{1.0cm}
{\Large\bfseries
Construction du continu comme limite projective de cylindres compatibles,\par
génération coinductive de \(\pi,\varphi,e,\alpha\) et stabilisateurs\par
exceptionnels de l'incidence de frontière\par}
\vspace{1.1cm}
{\color{hmtgold}\rule{0.68\textwidth}{0.8pt}\par}
\vfill
{\large Oumar Haidara Fall\par}
{\large Rubén Ramos Balsa\par}
\vspace*{1.8cm}
\end{titlepage}
\hypersetup{pageanchor=true}

\chapter*{Résumé}
\addcontentsline{toc}{chapter}{Résumé}

On construit une arithmétique d'histoires orientées dans laquelle l'unité est
conservée sous raffinement tandis qu'augmentent la résolution, la hauteur
holonomique et la mémoire. Le support fini de tous les chemins possibles
maintient simultanément les évaluations additive et multiplicative, avec
leurs résidus, quotients et retenues ; une fibre tritique distingue les régimes
elliptique, parabolique et hyperbolique ; et un noyau de phase topologique sélectionne,
transporte et actualise chaque état sans confondre le retour de phase avec le
retour de l'état. La composition de ces trois structures détermine une
algèbre associative non commutative, graduée par la mémoire, et un système de
raffinements compatibles dont l'espace d'histoires est une limite projective.
Les observables cylindriques se prolongent de manière covariante et constituent
la limite inductive duale. Cette double variance fournit la définition
opératoire de l'holographie employée dans l'article.

La dynamique combine une phase nonadique finie, un codage mécanique de
pente $\delta=\log_{10}(10/9)$ et un cocycle de retenue. Le résultat est
un ordre topologique-temporel apériodique de complexité linéaire et d'entropie
topologique nulle, muni d'une holonomie qui élève la mémoire à chaque retour
nonadique. Ses relations de Weyl séparent rigoureusement la périodicité locale
de l'apériodicité globale. Sur ce support, la prolongation
coinductive génère corrélativement $\pi$, $\varphi$, $e$ et $\alpha$ : les
trois premières coordonnées réalisent, respectivement, la clôture elliptique,
l'autoéchelle hyperbolique et le transport exponentiel ; la constante de
structure fine est l'invariant de compatibilité qui réunit, par deux
dérivations internes, la clôture dodécaphasique et la prolongation apériodique.

La même structure produit une vacance unitaire de retenue multiplicative
dans une fibre additive invariante. Son algèbre locale est
$\mathrm{Cl}_{1,1}\cong M_2(\R)$ et contient simultanément des générateurs
elliptiques, hyperboliques et des directions nulles de Witt. Cette section détermine
le noyau électronique et prolonge la généalogie vers l'action, les échelles
de Planck et de gravité, le paramètre de Barbero--Immirzi, le spectre d'aire et
l'information de frontière. La seconde projection conserve l'incidence et
conduit, au moyen des configurations de Paley et de Witt, du code de Golay, du
réseau de Leech et de l'algèbre de vertex $V^\natural$, aux stabilisateurs
exceptionnels et au groupe Monstre. Ainsi, la réalisation archimédienne et la
réalisation exceptionnelle ne constituent pas des sources indépendantes : ce sont deux
représentations d'une unique structure discrète du continu.

Le lecteur dodécaphasique uniforme détermine également une suite de premières racines de retour critique. On démontre que le rapport de leurs écarts converge vers la constante de Feigenbaum, en conservant la sélection initiale, l’orientation et la mémoire de la construction.

\vspace{4mm}
\noindent\textbf{Mots-clés :} holographie modulaire triadique ; limite
projective ; holonomie nonadique ; mot sturmien ; cristal temporel apériodique ;
algèbre de chemins ; constante de structure fine ; Clifford déployé ;
Barbero--Immirzi ; réseau de Leech ; groupe Monstre.

\clearpage

\chapter*{Résumé}
\addcontentsline{toc}{chapter}{Résumé}

Nous construisons une arithmétique d'histoires orientées dans laquelle l'unité est
conservée sous raffinement tandis que la résolution, la hauteur holonomique et la mémoire
augmentent. Le support fini de tous les chemins possibles conserve simultanément les évaluations additive et
multiplicative, avec leurs restes,
quotients et retenues ; une fibre tritique distingue les régimes elliptique, parabolique
et hyperbolique ; et un noyau de phase topologique sélectionne, transporte
et actualise chaque état sans identifier le retour de phase au retour de l'état.
Leur composition détermine une algèbre graduée par la mémoire, associative et
non commutative ainsi qu'un système de raffinements compatibles dont l'espace
d'histoires est une limite projective. Les observables cylindriques se prolongent de manière covariante et
forment la limite inductive duale. Cette dualité de variance est la notion opératoire
d'holographie utilisée tout au long de l'article.

La dynamique combine une phase nonadique finie, un codage mécanique de pente
$\delta=\log_{10}(10/9)$ et un cocycle de retenue. Elle produit un ordre
topologique--temporel apériodique de complexité linéaire et d'entropie topologique
nulle, muni d'une holonomie qui élève la mémoire à chaque retour nonadique.
Sur ce support, la prolongation coinductive génère $\pi$, $\varphi$,
$e$ et $\alpha$ comme sorties corrélées. La même structure produit une
vacance unitaire de retenue multiplicative sur une fibre additive invariante, dont
l'algèbre locale est $\mathrm{Cl}_{1,1}\cong M_2(\R)$. On montre ainsi que la réalisation
archimédienne et la réalisation exceptionnelle de l'incidence de frontière sont
deux représentations d'une même structure discrète du continu.

Le lecteur dodécaphasique uniforme détermine également une suite de premières racines de retour critique. On démontre que le rapport de leurs écarts converge vers la constante de Feigenbaum, en conservant la sélection initiale, l’orientation et la mémoire de la construction.

\vspace{4mm}
\noindent\textbf{Mots-clés :} holographie modulaire triadique ; limite projective ;
holonomie nonadique ; mot sturmien ; cristal temporel apériodique ; algèbre de chemins ;
constante de structure fine ; algèbre de Clifford déployée ; paramètre de Barbero--Immirzi ;
réseau de Leech ; groupe Monstre.
\clearpage
\tableofcontents
\clearpage

\let\cleardoublepage\clearpage
\mainmatter
\part{Arithmétique bifeuilletée, orientation trigéométrique et dynamique de phase}

\chapter{L'unité comme section compatible d'une arithmétique d'histoires}
\begingroup
\renewcommand{\chapter}[2][]{}%
\chapter{Thèse et architecture mathématique}
\label{sec:tesis-inaugural-helicoidal-hmt-md}

Le lecteur dodécaphasique uniforme détermine également une suite de premières racines de retour critique. On démontre que le rapport de leurs écarts converge vers la constante de Feigenbaum, en conservant la sélection initiale, l’orientation et la mémoire de la construction.


\section{L'unité comme invariant de raffinement}

La question centrale ne consiste pas à obtenir une collection de valeurs
scalaires à partir d'une table finie. Elle consiste à déterminer la structure
minimale qui conserve une unité à travers des raffinements illimités,
distingue trois classes de courbure, transporte la mémoire de ses parcours et
génère, au moyen de représentations typées, nombres, constantes, sections,
champs et symétries. La droite réelle n'est pas présupposée comme domaine universel de
cette construction. Elle apparaît au terme d'une contraction compatible de
cylindres ; elle conserve la valeur scalaire de l'histoire, mais pas nécessairement
la totalité de sa généalogie.

Pour chaque profondeur $n\geq0$, soit $X_n$ l'ensemble fini des états
admissibles. Un état n'est pas un chiffre isolé, mais une fibre structurée

\begin{equation}
x=\bigl(a,b;\Sigma,\Pi;R^+,q^+,R^\times,q^\times;
\kappa;\tau,o;\theta_9,\theta_{12};h;\partial;\mu;\rho;s\bigr),
\label{eq:estado-integral}
\end{equation}

où $(a,b)$ localise la cellule ; $\Sigma$ et $\Pi$ sont les feuillets additif et
multiplicatif ; $(R^+,q^+)$ et $(R^\times,q^\times)$ conservent le résidu et le
quotient ; $\kappa$ est la retenue ; $\tau\in\{+1,0,-1\}$ type la courbure ;
$o$ fixe l'orientation ; $\theta_9$ et $\theta_{12}$ sont les phases nonadique et
dodécaphasique ; $h$ mesure la hauteur hélicoïdale ; $\partial$ enregistre l'incidence
de frontière ; $\mu$ est la mémoire ; $\rho$ conserve le chemin et $s$ la condition
de survie. Ces coordonnées ne forment pas un produit libre : elles satisfont
les relations de l'arithmétique des chemins, le typage trigéométrique et la
dynamique de transport qui seront construits dans les chapitres suivants.

Les applications

\begin{equation}
r_{m,n}:X_m\longrightarrow X_n,\qquad m>n,
\label{eq:restriccion}
\end{equation}

oublient la dernière couche de résolution, mais préservent toute relation visible à
la profondeur $n$. L'espace des histoires infiniment raffinables est

\begin{equation}
X_\infty=\varprojlim_n X_n.
\label{eq:limite-inverso}
\end{equation}

Si $e_x$ désigne l'idempotent caractéristique du cylindre déterminé par
$x\in X_n$, la prolongation des observables satisfait

\begin{equation}
r_{m,n}^{*}e_x=\sum_{y:\,r_{m,n}(y)=x}e_y.
\label{eq:elevacion-idempotente}
\end{equation}

Par conséquent,

\begin{equation}
r_{m,n}^{*}\!\left(\sum_{x\in X_n}e_x\right)
=\sum_{y\in X_m}e_y,
\label{eq:conservacion-unidad}
\end{equation}

de sorte que l'unité est conservée même lorsque le nombre d'extensions augmente. L'évolution
ne conserve pas une configuration immobile : elle conserve l'augmentation, les
relations et la possibilité de reconstruire chaque antécédent à partir de l'histoire
complète.

\section{Double variance et double projection}

La profondeur des états et l'extension des observables avancent dans des
sens catégoriques opposés. Aux restrictions de
\eqref{eq:restriccion} correspondent des morphismes de prolongation

\begin{equation}
r_{m,n}^{*}:\Halg_n\longrightarrow\Halg_m,
\qquad f\longmapsto f\circ r_{m,n},
\end{equation}

et l'algèbre des observables cylindriques s'obtient comme limite inductive :

\begin{equation}
\Halg_{\mathrm{cyl}}=\varinjlim_n(\Halg_n,r_{n+1,n}^{*}).
\label{eq:limite-directo}
\end{equation}

Le couple formé par \eqref{eq:limite-inverso} et
\eqref{eq:limite-directo} fournit la structure holographique fondamentale :
l'état acquiert de la résolution par compatibilité contravariante et le langage
d'observation acquiert de la portée par prolongation covariante. Il ne s'agit pas d'une
ressemblance figurative entre échelles, mais de la naturalité de
l'évaluation

\begin{equation}
\langle r_{m,n}^{*}a,y\rangle_m
=\langle a,r_{m,n}y\rangle_n.
\label{eq:naturalidad-evaluacion}
\end{equation}

La double projection du titre désigne deux représentations de ce même
couple. La première contracte des cylindres compatibles vers une réalisation
archimédienne ; la seconde retient l'incidence de frontière et la prolonge
vers des structures exceptionnelles :

\begin{equation}
\rho_{\mathrm{arq}}:\Halg_{\mathrm{cyl}}\longrightarrow\cA_{\R},
\qquad
\rho_{\mathrm{exc}}:\Halg_{\mathrm{cyl}}\longrightarrow\cA_{\mathrm{exc}}.
\label{eq:doble-proyeccion}
\end{equation}

Les deux possèdent la même source et conservent des invariants différents. La première
retient la coordonnée scalaire d'une section compatible ; la seconde retient
les motifs d'incidence qui déterminent ses stabilisateurs. Le mot
\emph{double} ne duplique pas l'holographie : il distingue deux codomaines non
équivalents d'un unique mécanisme de raffinement et de lecture.

\section{Énoncé central}

\begin{theorem}[Évolution conservative et production de symétries]
Il existe une algèbre graduée d'histoires orientées
\begin{equation}
\Halg=\left[
\varinjlim_n
\frac{\R\langle e_x,J_x,T_u\rangle}
{\cR_{\mathrm{APP\text{-}TRIT\text{-}TPK}}}
\right]\rtimes_{\Gamma_9}\N,
\label{eq:algebra-rectora}
\end{equation}
qui agit sur $X_\infty$ et satisfait les propriétés suivantes :
\begin{enumerate}[label=\textup{(\roman*)},leftmargin=9mm]
\item l'unité et l'augmentation sont conservées sous tout raffinement compatible ;
\item le retour nonadique ferme la phase et élève la mémoire ;
\item les régimes elliptique, parabolique et hyperbolique appartiennent à un unique
calcul fonctionnel trigéométrique ;
\item le facteur symbolique possède une complexité linéaire, une entropie topologique nulle
et un ordre apériodique à longue portée ;
\item la structure discrète du continu émerge conjointement de l'espace des
histoires et de l'algèbre des observables ;
\item $\pi$, $\varphi$, $e$ et $\alpha$ sont des sorties corrélées de la
même prolongation coinductive ;
\item une symétrie d'une réalisation est le stabilisateur des relations et
de l'incidence de frontière que cette réalisation conserve.
\end{enumerate}
\end{theorem}

La démonstration sera répartie tout au long de l'article. On n'utilisera pas de
constante connue pour sélectionner sa propre généalogie, ni de groupe de
symétrie pour construire rétrospectivement l'objet sur lequel il agit. L'ordre
logique sera toujours le même : arithmétique bifeuilletée, typage trigéométrique,
transport holonomique, raffinement compatible, continu, constantes,
sections physiques et stabilisateurs de frontière.

\section{Irréductibilité du noyau générateur}

La structure possède trois composantes fonctionnellement irréductibles. Si l'on
identifie les deux feuillets arithmétiques, le potentiel
$\Delta_{\Sigma\Pi}$ et le commutateur de leurs projecteurs disparaissent. Si l'on identifie
les trois régimes tritiques, le spectre disparaît

\begin{equation}
J_+^2=-I,\qquad J_0^2=0,\qquad J_-^2=I.
\end{equation}

Si l'on trivialise le cocycle de transport, l'incrément

\begin{equation}
[D_M,\Gamma_9]=\Gamma_9
\label{eq:conmutador-memoria-intro}
\end{equation}

s'annule et le tour de phase devient indiscernable du retour de l'état.
Chaque quotient d'oubli détruit un invariant de type différent que les deux autres
facteurs ne reconstruisent pas. Cette irréductibilité explique pourquoi la
structure recherchée ne peut se réduire à une table arithmétique, à un mot
apériodique, à une algèbre de Clifford ou à une matrice isolée.

La conséquence conceptuelle est précise. Le nombre est une coordonnée scalaire
d'une représentation ; la constante est un invariant stable de retour ou
d'autoéchelle ; la particule est une section persistante ; le champ transporte
des sections ; la symétrie stabilise les relations de sa frontière. Ces cinq
classes procèdent de la même algèbre, mais ne partagent pas de codomaine et ne peuvent pas
être confondues entre elles.
 \endgroup

\chapter{Support arithmétique de tous les chemins possibles}
\label{chap:83-01}

% RESULTADO_RECUPERADO. Owner integro: sections/hmt/03_app_ortograma.tex
% Destino causal: capitulo 1; la jerarquia se subordina sin alterar el owner.
% Los dos bloques científicos del owner se activan en una única residencia.
\begingroup
\def\HMTAPPFormal{1}
\def\HMTAPPDevelopments{1}
\begin{HMTScientificOwnerSubordinated}
\ifdefined\HMTAPPFormal
\chapter{APP : support, opérations et chemins}
\label{chap:app}
\index{APP}

\noindent
Le sigle APP provient de \emph{All Possible Paths}. Il désigne une unique
structure de chemins sur un support arithmétique fini : le graphe
sous-jacent possède quatre-vingt-un sommets, tandis que l'ensemble des chemins
finis, réuni sur toutes les longueurs, est dénombrable. Sa matrice
arithmétique principale est la table multiplicative nonadique. L'addition est la loi
interne dont l'accumulation engendre chaque ligne et chaque colonne de cette table et agit
aussi comme observable de comparaison sur le même support. APP réunit ainsi
addition, produit et
composition de trajectoires avant d'effectuer la projection ternaire ou
l'évaluation archimédienne. La définition suivante fixe de manière exhaustive
l'univers de chemins considéré.

\section{Domaine, générateurs et relations préformelles du support APP}

Le support d'APP est un tore arithmétique fini : chacune de ses deux
coordonnées parcourt les neuf classes résiduelles et revient à son origine après
neuf translations unitaires. Ses quatre-vingt-un points sont donc des
paires de classes modulo neuf. Le graphe de Cayley associé conserve cette
structure cyclique : un sommet représente une position, une arête orientée
représente un déplacement cardinal unitaire et une suite finie
d'arêtes représente un chemin. Ainsi, \emph{tous les chemins
possibles} désigne l'ensemble des mots cardinaux finis ayant leur point
initial dans le tore.

Chaque sommet reçoit simultanément deux évaluations arithmétiques. Le
\emph{feuillet additif} publie la somme de ses deux coordonnées et le
\emph{feuillet multiplicatif} publie leur produit. Le terme \emph{feuillet}
désigne ici une évaluation du même point et du même chemin ; APP conserve
un unique support et deux fonctions sur celui-ci.

La différence opératoire entre les deux feuillets est exacte. Dans une ligne ou une colonne,
l'addition agit par translation et accumule de manière répétée le même incrément.
La multiplication par une unité permute également les classes ; la
multiplication par un diviseur de zéro replie plusieurs classes sur un même
résidu. Pour conserver l'information que ce repliement masque, APP élève
chaque évaluation à ses coordonnées résidu--quotient :
\[
 i+j=R^+_{ij}+9q^+_{ij},
 \qquad
 ij=R^\times_{ij}+9q^\times_{ij}.
\]
Les paires \((R^+,q^+)\) et \((R^\times,q^\times)\) récupèrent les valeurs
entières antérieures à la réduction. Lors de la concaténation de blocs, la retenue additive
\(\kappa_9\) enregistre l'écart entre additionner des représentants entiers et
additionner d'abord leurs résidus ; avec les deux quotients, elle rend la
récupération compositionnelle. Ainsi, l'accumulation additive construit les trajectoires visibles,
tandis que le relèvement résidu--quotient déploie les fibres créées par le
repliement multiplicatif.

Pour un chemin \(\gamma=(x_0,\ldots,x_r)\), la généalogie visible est la
suite des positions et des paires d'évaluations
\[
 \mathcal G_{\rm vis}(\gamma)
 =\bigl(\gamma;(\Sigma(x_k),\Pi(x_k))_{k=0}^{r}\bigr).
\]
Sa généalogie enrichie ajoute, à chaque pas, les quotients, la retenue et
l'orientation transportée. La projection enrichie--visible conserve le
chemin et les résidus et oublie précisément ces coordonnées de mémoire. Cette
distinction sera l'entrée de l'état complet du TRIT et du Topological Phase Kernel.
Le domaine complet des (104\,976) conditions initiales à double curseur et
les (468) fibres d'émission qu'elles produisent sont publiés intégralement dans
\cref{app:catalogo-semillas-app-rev9} ; chaque germe y est distingué de
l'émission à laquelle il appartient et la multiplicité de la fibre est conservée.

\section{Support arithmétique, translations et observables}

Soit
\[
\dr(n)=1+((n-1)\bmod9)\in\{1,\ldots,9\}.
\]
Cette application choisit le représentant \(9\) pour la classe nulle. Nous définissons
\(X_9=(\ZZ/9\ZZ)^2\), sur lequel agit le groupe de translations
\((\ZZ/9\ZZ)^2\), et fixons l'ensemble symétrique de générateurs
\[
 \mathscr D=\{(1,0),(0,1),(-1,0),(0,-1)\}.
\]
Le graphe de Cayley
\[
 \GraphAPP=\operatorname{Cay}(X_9,\mathscr D)
\]
possède \(81\) sommets et quatre arêtes orientées ayant leur origine à chaque sommet.
Un chemin orienté de longueur \(r\) est une suite
\(\gamma=(x_0,\ldots,x_r)\) dans \(X_9\) telle que
\(x_{k+1}-x_k\in\mathscr D\) pour tout \(k<r\). Nous notons
\(\operatorname{Path}(\GraphAPP)\) l'ensemble de tous les chemins finis
ainsi définis.

\begin{definition}[Structure APP]
La structure arithmétique des chemins possibles est le tuple
\[
 \AAPP=(X_9,\GraphAPP,\Sigma,\Pi,
        \operatorname{Path}(\GraphAPP)),
\]
où les deux observables de sommet sont
\[
\Sigma(i,j)=\dr(i+j),
\qquad
\Pi(i,j)=\dr(ij),
\qquad1\le i,j\le9.
\]
\end{definition}

\subsection*{Cadre cardinal et mots de trajectoire}

Pour rendre explicite l'orientation du tore, nous utilisons les coordonnées
\((\text{ligne},\text{colonne})\), avec les lignes croissantes vers le sud et les
colonnes croissantes vers l'est, et écrivons
\[
 \delta(N)=(-1,0),\qquad \delta(E)=(0,1),\qquad
 \delta(S)=(1,0),\qquad \delta(O)=(0,-1).
\]
Soit \(\mathscr D_{\rm card}=\{E,N,O,S\}\), où \(O\) désigne l'ouest. Étant donnés un point initial
\(x_0\in X_9\) et un mot \(w=d_1\cdots d_r\in
\mathscr D_{\rm card}^{r}\), sa trajectoire est
\[
 \gamma(x_0;w)_k
 =x_0+\sum_{j=1}^{k}\delta(d_j)\pmod 9,
 \qquad 0\leq k\leq r.
\]
La correspondance \((x_0,w)\mapsto\gamma(x_0;w)\) est bijective entre
\(X_9\times\mathscr D_{\rm card}^{r}\) et les chemins orientés de
longueur \(r\). En particulier, il existe exactement \(81\cdot4^r\) chemins
de cette longueur lorsque les retours et les reculs sont autorisés.

Le mot \(w\) se ferme sur le tore précisément lorsque
\[
\#S(w)-\#N(w)\equiv0\pmod9,
\qquad
\#E(w)-\#O(w)\equiv0\pmod9.
\]
Pour un mot fermé, le déplacement entier calculé avant réduction
modulo neuf définit le vecteur d'enroulement
\[
 \operatorname{wind}(w)=\frac1{9}
 \bigl(\#S(w)-\#N(w),\#E(w)-\#O(w)\bigr)\in\ZZ^2.
\]
L'inversion de l'orientation inverse l'ordre du mot, échange
\(E\leftrightarrow O\) et \(N\leftrightarrow S\), et transforme
\(\operatorname{wind}(w)\) en son opposé. Ce cadre cardinal est celui qui est
transporté ultérieurement dans le noyau de phase et fait partie de la définition
de chaque trajectoire.

Les expressions sont écrites ici dans la carte des représentants
\(\{1,\ldots,9\}\). La formule
\(\widetilde\Sigma(i,j)=\dr(i+j-1)\) est une carte translatée. Elle représente le
même tore uniquement lorsque la translation est également appliquée aux indices, aux conditions
initiales, aux trajectoires positionnelles et aux étiquettes. Toutes les démonstrations suivantes emploient
\(\Sigma(i,j)=\dr(i+j)\) et n'alternent pas entre les deux cartes.

\begin{figure}[H]
\centering
\resizebox{\textwidth}{!}{\begin{tikzpicture}[font=\scriptsize,cell/.style={minimum width=4.2mm,minimum height=4.2mm,inner sep=0pt,draw=hmtgray!35}]
  \begin{scope}[xshift=0cm]
    \node[font=\bfseries\scriptsize,hmtblue,align=center] at (1.72,.92)
      {matrice arithmétique APP\\[-.5mm]\(\Pi(i,j)=\dr(ij)\)};
    \foreach \i in {1,...,9}{
      \node[font=\bfseries] at (-.35,{-(\i-1)*.43}) {\i};
      \foreach \j in {1,...,9}{
        \pgfmathtruncatemacro{\v}{mod(\i*\j-1,9)+1}
        \pgfmathtruncatemacro{\border}{(\i==9)||(\j==9)}
        \pgfmathtruncatemacro{\rad}{(mod(\i,3)==0)||(mod(\j,3)==0)}
        \ifnum\border=1
          \node[cell,fill=hmtlightblue,text=hmtblue,font=\bfseries] at ({(\j-1)*.43},{-(\i-1)*.43}) {\v};
        \else
          \ifnum\rad=1
            \node[cell,fill=hmtlightviolet] at ({(\j-1)*.43},{-(\i-1)*.43}) {\v};
          \else
            \node[cell,fill=white] at ({(\j-1)*.43},{-(\i-1)*.43}) {\v};
          \fi
        \fi
      }
    }
    \foreach \j in {1,...,9}{\node[font=\bfseries] at ({(\j-1)*.43},.38) {\j};}
    \draw[hmtgold,very thick,rounded corners=1pt] (1.075,-1.075) rectangle (1.935,-1.935);
    \draw[hmtblue,very thick] (3.72,.22)--(3.72,-3.72)--(-.22,-3.72);
    \node[hmtblue,anchor=north] at (1.75,-4.02) {subconjunto coordenado de la clase nula};
  \end{scope}

  \begin{scope}[xshift=5.25cm]
    \node[draw=hmtblue,rounded corners=2pt,fill=hmtlightblue,
      text width=6.25cm,align=left,inner sep=8pt] at (2.9,-1.35) {%
      \textbf{Ley interna de acumulación}\par\smallskip
      \[
        \Pi(i,j)=\dr(\underbrace{i+\cdots+i}_{j\ {\rm sumandos}})
        =\dr(ij).
      \]
      La somme engendre les lignes multiplicatives et, dans la direction
      transversale, les colonnes. L'observable
      \(\Sigma(i,j)=\dr(i+j)\) conserve cette loi élémentaire pour comparer
      les deux lectures sur le même point.};
    \node[draw=hmtgold,rounded corners=2pt,fill=hmtlightgold,
      text width=6.25cm,align=center,inner sep=8pt] at (2.9,-3.65) {%
      \(\displaystyle
      \AAPP=(X_9,\GraphAPP,\Pi;\Sigma,
      \operatorname{Path}(\GraphAPP))\)\par
      une seule APP, une matrice principale, une loi additive interne et quatre
      directions orientées \(N,E,S,O\).};
  \end{scope}
\end{tikzpicture}
 }
\caption{L'unique matrice APP et sa loi interne. La table multiplicative
nonadique occupe le corps principal ; chacune de ses lignes est engendrée par
accumulation additive. L'ombrage violet indique les lignes ou colonnes non
unitaires ; la ligne bleue identifie le sous-ensemble déterminé par le
représentant \(9\) de la classe nulle ; le cadre doré souligne le bloc
central \(2\times2\).}
\label{fig:app-ortograma}
\end{figure}

Chaque point de \(X_9\) reçoit simultanément les valeurs du produit et de la somme,
mais n'appartient pas pour autant à deux matrices APP distinctes. Les chemins de
\(\operatorname{Path}(\GraphAPP)\) parcourent un unique support ; \(\Pi\)
publie sa table principale et \(\Sigma\) conserve l'opération interne qui
engendre les accumulations et permet de mesurer leur incompatibilité.

\begin{proposition}[Structure des lignes et des colonnes]
\label{prop:app-latina-unidades}
La matrice de \(\Sigma\) est un carré latin d'ordre neuf. La matrice de
\(\Pi\) ne l'est pas ; cependant, pour chaque
\(u\in(\ZZ/9\ZZ)^\times=\{1,2,4,5,7,8\}\), la ligne et la colonne
correspondant à \(u\) sont des permutations des neuf classes résiduelles.
\end{proposition}

\begin{proof}
Une fois \(i\) fixé, l'application \(j\mapsto i+j\) est une translation de
\(\ZZ/9\ZZ\) ; il en va de même en fixant \(j\). Ainsi, chaque représentant
apparaît exactement une fois dans chaque ligne et chaque colonne de \(\Sigma\). Pour
\(\Pi\), la multiplication par \(u\) permute les classes si et seulement si \(u\)
est une unité. Les lignes correspondant à \(3,6,9\) présentent des répétitions,
de sorte que la propriété latine ne s'étend pas à toute la matrice
multiplicative.
\end{proof}

\section{Structure multiplicative de l'anneau local}

Chaque ligne de la table multiplicative est une addition répétée. La ligne \(i\)
accumule \(i\) le long d'une direction ; la colonne \(j\) accumule \(j\)
le long de la direction perpendiculaire. La cellule \((i,j)\) est l'intersection
des deux accumulations et renvoie \(ij=ji\). La table représente ainsi le produit
et la commutativité au moyen de deux accumulations transversales sur la même
grille.

La deuxième ligne est
\[
2,4,6,8,1,3,5,7,9.
\]
Avant le retour modulaire apparaissent les nombres pairs, et ensuite les impairs. La troisième
ligne est
\[
3,6,9,3,6,9,3,6,9,
\]
et rend le radical visible. Soit
\[
 R:=\ZZ/9\ZZ,
 \qquad
 \mathfrak m:=(3)=\{\overline0,\overline3,\overline6\}.
\]
La multiplication par trois définit
\[
 M_3:R\longrightarrow R,
 \qquad x\longmapsto3x,
 \qquad
 \ker M_3=\operatorname{im}M_3=\mathfrak m.
\]
Par le premier théorème d'isomorphisme, elle induit une identification linéaire
\begin{equation}
 \overline M_3:R/\mathfrak m\xrightarrow{\ \sim\ }\mathfrak m,
 \qquad \overline x\longmapsto3x,
 \label{eq:multiplicador-tres-radical-rev3}
\end{equation}
d'espaces vectoriels sur \(\mathbb F_3\). Ce n'est pas un isomorphisme
d'anneaux : \(\mathfrak m^2=0\), tandis que le quotient est un corps.

En représentants positifs, les trois phases de
\eqref{eq:multiplicador-tres-radical-rev3} sont
\begin{equation}
 369\longmapsto693\longmapsto936\longmapsto369,
 \label{eq:orbita-radical-369-rev3}
\end{equation}
produites par le déplacement de phase \(j\mapsto j+1\). La ligne
correspondant à \(6\equiv-3\pmod9\) inverse l'orientation du même
cycle. Cette orbite radicale ne doit pas être confondue avec l'entier
\(369\,693\,100\), qui apparaîtra ensuite comme nombre de chiffres de
\(9^{9^9}\).

En réduction directe modulo trois,
\[
3,6,9\longmapsto0,0,0.
\]
Si nous extrayons d'abord le facteur trois et conservons sa coordonnée interne,
\[
\eta_{\mathfrak m}(3a)=a\bmod3,
\qquad
3,6,9\longmapsto1,2,0.
\]
La seconde application conserve la coordonnée interne de l'idéal engendré par
trois. Elle ne doit pas être confondue avec la projection ternaire orientée qui sera définie
dans la section consacrée au TRIT\@.

\section{Sous-ensemble de coordonnées de la classe nulle}

Dans le graphe de l'observable multiplicative,
\[
\Pi(9,j)=\Pi(i,9)=9.
\]
La dernière ligne et la dernière colonne forment une union de dix-sept positions,
toutes de valeur neuf. Elles ne constituent pas une frontière intrinsèque du tore
arithmétique : c'est le sous-ensemble de la carte dans lequel au moins une coordonnée
représente la classe nulle par \(9\). Il existe en outre quatre représentants nuls intérieurs en
\[
(3,3),(3,6),(6,3),(6,6).
\]
La figure distingue ainsi ce sous-ensemble et les zéros internes de l'idéal
non unitaire. Le symbole supplémentaire employé plus loin dans la complétion
cubique appartient à un espace étendu et ne s'identifie à aucun
d'entre eux.

\section{Extension résidu--quotient}

La réduction modulo neuf conserve le résidu et élimine le quotient de la
division euclidienne. Pour étudier les deux composantes, nous définissons, en chaque point,
\[
i+j=9q^+_{ij}+R^+_{ij},
\qquad
ij=9q^\times_{ij}+R^\times_{ij},
\]
où \(R^+=\Sigma\), \(R^\times=\Pi\) et les \(q\) conservent le quotient.
L'extension résidu--quotient d'APP est notée
\[
\AAPPlift=(R^+,q^+;R^\times,q^\times).
\]

Soit \(s_9:\ZZ/9\ZZ\to\{1,\ldots,9\}\) la section de représentants fixée
par \(\dr\). La retenue additive associée à cette section est
\[
 \kappa_9(a,b)
 =\frac{s_9(a)+s_9(b)-s_9(a+b)}9\in\ZZ.
\]

\begin{lemma}[Identité cocyclique de la retenue]
\label{lem:app-cociclo-acarreo}
Pour tous \(a,b,c\in\ZZ/9\ZZ\),
\[
 \kappa_9(a,b)+\kappa_9(a+b,c)
 =\kappa_9(b,c)+\kappa_9(a,b+c).
\]
\end{lemma}

\begin{proof}
Les deux membres sont égaux à
\[
 \frac{s_9(a)+s_9(b)+s_9(c)-s_9(a+b+c)}9.\qedhere
\]
\end{proof}

Ce cocycle enregistre l'écart entre additionner des représentants entiers et
additionner d'abord dans le quotient. C'est la forme compositionnelle de la coordonnée de
quotient qui réapparaîtra dans la concaténation en base \(1000\).

Les totaux bruts sont
\[
\sum_{i,j}(i+j)
=9\sum_i i+9\sum_j j
=2\cdot9\cdot45
=810,
\]
et
\[
\sum_{i,j}ij
=\left(\sum_i i\right)\left(\sum_j j\right)
=45^2
=2025.
\]
Les sommes visibles s'obtiennent à partir des tables :
\[
\sum R^+=405,
\qquad
\sum R^\times=459.
\]
Par conséquent, les sommes totales des coordonnées de quotient sont
\[
\sum q^+=\frac{810-405}{9}=45,
\qquad
\sum q^\times=\frac{2025-459}{9}=174.
\]

\begin{exactbox}[title={Décomposition exacte de la différence intégrée}]
\[
\boxed{
2025-810
=1215
=(459-405)+9(174-45)
=54+9\cdot129.}
\]
La différence entre les sommes de résidus est \(54\), tandis que la
différence entre les sommes de quotients est \(129\). L'égalité recompose la
différence totale entre les observables entières et démontre que la composante
de quotient contient une information arithmétique non déterminée par la classe
résiduelle.
\end{exactbox}

\fi

\ifdefined\HMTAPPDevelopments
\chapter{Opérateurs arithmétiques et géométrie locale de l'espace des chemins}
\label{chap:app-operadores}

\section{Localisation de l'excès dans le radical nilpotent}

Le groupe des unités de \(\ZZ/9\ZZ\) est
\[
(\ZZ/9\ZZ)^\times=\{1,2,4,5,7,8\},
\]
et l'idéal maximal est
\[
\mathfrak m=(3)=\{0,3,6\}=\{9,3,6\}_{\rm visual},
\qquad\mathfrak m^2=0.
\]
Chaque ligne unitaire de \(\Pi\) est une permutation de \(1,\ldots,9\) et a pour somme
\(45\). Les lignes \(3\) et \(6\) ont pour somme \(54\) ; la ligne \(9\), \(81\).
Donc
\[
459=6\cdot45+2\cdot54+81
\]
et
\[
459-405=2(54-45)+(81-45)=54.
\]
Tout l'excès se localise dans le secteur non unitaire.

L'union des lignes et des colonnes \(3,6,9\) contient \(45\) cellules et a pour somme
\(297\). Son carré central \(3\times3\) a pour somme \(81=3\cdot27\) ; les bras
extérieurs, \(216=3\cdot72\). Ainsi,
\[
297=216+81=3(72+27)=3\cdot99.
\]
L'interprétation de \(72\) et de \(27\) s'obtient plus loin à partir du
bloc central, sélectionné de manière unique.

\section{Espérance conditionnelle par rapport à la projection \(\mathbb Z/9\mathbb Z\to\mathbb Z/3\mathbb Z\)}

La réduction canonique
\(\ZZ/9\ZZ\to\ZZ/3\ZZ\) détermine la partition
\[
 C_1=\{1,4,7\},\qquad
 C_2=\{2,5,8\},\qquad
 C_0=\{3,6,9\}.
\]
L'espérance conditionnelle uniforme moyenne chaque observable sur les neuf points de chaque bloc
\(C_a\times C_b\). Les matrices agrégées sont
\[
 M_{\Pi}=
 \begin{pmatrix}4&5&6\\5&4&6\\6&6&9\end{pmatrix},
 \qquad
 M_{\Sigma}=
 \begin{pmatrix}5&6&4\\6&4&5\\4&5&6\end{pmatrix}.
\]
Puisque chaque bloc contient neuf positions,
\[
 9\sum_{a,b}(M_\Pi)_{ab}=459,\qquad
 9\sum_{a,b}(M_\Sigma)_{ab}=405.
\]
Ainsi, le moyennage conserve les totaux des deux observables et localise de
nouveau leur différence :
\[
 \sum M_\Pi-\sum M_\Sigma=51-45=6,\qquad 9\cdot6=54.
\]

\begin{theorem}[Spectre de l'opérateur multiplicatif agrégé]
\label{thm:app-espectro-agregado}
L'opérateur matriciel \(M_\Pi\) satisfait
\[
 \det M_\Pi=-9
\]
et son spectre réel est
\[
 \operatorname{Spec}(M_\Pi)
 =\{-1,\,9-6\sqrt2,\,9+6\sqrt2\}.
\]
En particulier, la forme quadratique associée possède la signature \((2,1)\).
\end{theorem}

\begin{proof}
Le développement du déterminant donne \(-9\). Le polynôme caractéristique
se factorise sous la forme
\[
 (\lambda+1)\bigl((\lambda-9)^2-72\bigr).
\]
Ses racines sont celles indiquées et \(9-6\sqrt2>0\) ; il existe donc deux
directions positives et une négative.
\end{proof}

La décomposition spectrale admet une forme entière sur un sous-réseau
invariant. Pour
\[
 v_-=(1,-1,0),\qquad v_+=(1,1,0),\qquad v_0=(0,0,1),
\]
on a \(M_\Pi v_-=-v_-\), tandis que la restriction de \(M_\Pi\) au
sous-module engendré par \(v_+,v_0\) est représentée par
\[
 \begin{pmatrix}9&6\\12&9\end{pmatrix}
 =3H,\qquad
 H=\begin{pmatrix}3&2\\4&3\end{pmatrix}\in SL(2,\ZZ).
\]
Les valeurs propres de \(H\) sont
\[
 3\pm2\sqrt2=(1\pm\sqrt2)^2.
\]
Les trois vecteurs \(v_-,v_+,v_0\) engendrent un sous-réseau d'indice deux dans
\(\ZZ^3\), car le déterminant de la matrice qui les a pour colonnes vaut
\(-2\). On n'affirme donc pas ici une somme directe entière de tout
\(\ZZ^3\) : sur ce sous-réseau invariant, l'opérateur agrégé sépare une
direction de valeur propre \(-1\) et un sous-module hyperbolique régi par
l'unité quadratique \(3+2\sqrt2\) ; sur \(\RR^3\), la décomposition spectrale
est complète.

\section{Commutateur des projecteurs arithmétiques et séparation de leurs images}

La différence entre \(\Sigma\) et \(\Pi\) ne se réduit pas à leurs sommes globales.
Pour \(d\ge2\), soit \(\mathscr A_d=\{1,\ldots,9\}^d\) et définissons
\[
 \Sigma_d(x)=\dr(x_1+\cdots+x_d),\qquad
 \Pi_d(x)=\dr(x_1\cdots x_d).
\]
Dans l'espace de Hilbert \(\ell^2(\mathscr A_d)\), muni de la mesure
uniforme, notons \(P_{\Sigma,d}\) et \(P_{\Pi,d}\) les espérances
conditionnelles orthogonales sur les fonctions mesurables par rapport à
\(\Sigma_d\) et à \(\Pi_d\), respectivement. La table conjointe normalisée est
\[
 C^{(d)}_{ab}
 =\frac{N^{(d)}_{ab}}{\sqrt{n_a m_b}},
\quad
 N^{(d)}_{ab}
 =\#\{x:\Sigma_d(x)=a,\ \Pi_d(x)=b\},
\]
où \(n_a=\sum_bN^{(d)}_{ab}\) et \(m_b=\sum_aN^{(d)}_{ab}\).

Sur le support bidimensionnel original, la corrélation conjointe exacte est
\[
N^{(2)}=
\begin{pmatrix}
0&0&2&0&0&2&3&0&2\\
3&0&2&0&0&2&0&0&2\\
0&2&0&0&2&0&0&2&3\\
0&0&2&3&0&2&0&0&2\\
0&0&2&3&0&2&0&0&2\\
0&2&0&0&2&0&0&2&3\\
3&0&2&0&0&2&0&0&2\\
0&0&2&0&0&2&3&0&2\\
0&2&0&0&2&0&0&2&3
\end{pmatrix}.
\]
Les lignes, ordonnées selon la valeur de \(\Sigma\), ont pour somme neuf ; les colonnes,
ordonnées selon la valeur de \(\Pi\), ont pour somme
\((6,6,12,6,6,12,6,6,21)\). La corrélation conjointe des deux observables est ainsi déterminée par énumération complète
avant l'introduction de toute
application d'évaluation ultérieure.

\begin{proposition}[Commutateurs en dimensions deux et trois]
\label{prop:app-no-conmutatividad}
Les deux paires de projecteurs satisfont
\[
 \|[P_{\Sigma,2},P_{\Pi,2}]\|=\frac{\sqrt2}{3},
 \qquad
 \|[P_{\Sigma,3},P_{\Pi,3}]\|=\frac{\sqrt3}{4}.
\]
En particulier, les observables additive et multiplicative n'induisent pas de
décompositions simultanément compatibles.
\end{proposition}

\begin{proof}
Les valeurs singulières \(s\) de \(C^{(d)}\) sont les cosinus des angles
principaux entre les deux sous-espaces. Dans le bloc bidimensionnel
correspondant, la norme du commutateur est \(s\sqrt{1-s^2}\). L'énumération
exacte des \(9^d\) mots produit, pour \(d=2\), les valeurs propres non
triviales
\[
 s^2\in\left\{\frac57,\frac13,\frac13\right\}.
\]
Le maximum de \(s^2(1-s^2)\) est \(2/9\), d'où résulte
\(\sqrt2/3\). Pour \(d=3\), la même construction donne le maximum
\(s^2(1-s^2)=3/16\), atteint en \(s^2=1/4\), et la norme est donc
\(\sqrt3/4\). Dans les deux cas, le commutateur est non nul.
\end{proof}

\begin{proposition}[Intersection des sous-espaces observables]
\label{prop:app-interseccion-proyectores}
L'énumération exacte des mots de \(\mathscr A_d\), pour
\(2\leq d\leq15\), vérifie
\[
 \operatorname{Ran}P_{\Sigma,d}\cap
 \operatorname{Ran}P_{\Pi,d}
 =\CC\mathbf 1.
\]
Dans cet intervalle de dimensions, les fonctions constantes constituent
l'unique sous-espace commun aux deux décompositions arithmétiques.
\end{proposition}

\begin{proof}
Pour deux projections orthogonales, la dimension de l'intersection de leurs
images coïncide avec la multiplicité de la valeur singulière \(1\) dans la matrice
de recouvrement \(C^{(d)}\). Le calcul entier des tables
\(N^{(d)}\), suivi du calcul exact de cette multiplicité, donne la valeur un
pour chaque \(d=2,\ldots,15\). Le vecteur constant appartient aux deux images ;
il engendre donc toute l'intersection.
\end{proof}

La suite des normes ne s'extrapole pas à partir des deux premiers cas.
En particulier, pour \(d=4\), le même calcul donne
\[
 \|[P_{\Sigma,4},P_{\Pi,4}]\|
 =\sqrt{\frac{2618939}{22240656}}
 \simeq0.3431538652,
\]
ce qui ne coïncide pas avec la prolongation élémentaire \(2/(d+1)\). Ce cas
constitue un contre-exemple à toute formule inférée uniquement de
\(d=2,3\).

\section{Bloc central et décomposition spectrale}

La caractérisation du bloc central s'obtient au moyen d'une condition de
congruence. Pour deux représentants consécutifs \(k,k+1\), considérons
\[
B_k=
\begin{pmatrix}
k^2&k(k+1)\\
k(k+1)&(k+1)^2
\end{pmatrix}\pmod9.
\]

\begin{theorem}[Unicité du bloc diagonal adjacent]
\label{thm:app-bloque-central}
Il existe un unique bloc \(B_k\) dont les entrées diagonales coïncident. Il est
déterminé par \(k=4\) et vaut
\[
B_c=
\begin{pmatrix}7&2\\2&7\end{pmatrix}.
\]
\end{theorem}

\begin{proof}
L'égalité des diagonales équivaut à
\[
k^2\equiv(k+1)^2\pmod9,
\]
c'est-à-dire \(2k+1\equiv0\pmod9\). Comme \(2\) est une unité de
\(\ZZ/9\ZZ\), la congruence a pour unique solution \(k\equiv4\pmod9\).
\end{proof}

Par conséquent, les positions \(4,5\) de l'observable multiplicative forment
\[
B_c=
\begin{pmatrix}7&2\\2&7\end{pmatrix}
=7I+2R,
\qquad
R=\begin{pmatrix}0&1\\1&0\end{pmatrix},
\qquad R^2=I.
\]
Avec \(P_\pm=(I\pm R)/2\),
\[
B_c=9P_++5P_-.
\]
\begin{theorem}[Hiérarchie typée de la séparation somme--produit]
\label{thm:app-jerarquia-cuatro-dos-seis-cincuentaycuatro}
Le bloc central a un écart spectral local (9-5=4). Son coefficient
hors diagonale est (2). La compression par les trois classes modulo trois
produit la différence agrégée (51-45=6), et le déploiement sur les neuf
positions produit le défaut global
\[
 9\cdot6=459-405=54.
\]
Les entiers (4), (2), (6) et (54) appartiennent, respectivement, à
l'écart spectral, au couplage local, à la compression modulaire et à
l'intégration bidimensionnelle.
\end{theorem}

\begin{proof}
La décomposition \(B_c=7I+2R=9P_++5P_-\) donne l'écart et le
coefficient local. Les matrices agrégées \(M_\Pi,M_\Sigma\) construites
ci-dessus satisfont
\(\sum M_\Pi=51\) et \(\sum M_\Sigma=45\). Chaque entrée agrégée représente
neuf positions, de sorte que le déploiement multiplie la différence par
neuf et récupère les totaux (459) et (405).
\end{proof}
Leurs concaténations ordonnées produisent les blocs \(72\) et \(27\) :
\[
72+27=99,
\qquad
72-27=45=\det B_c.
\]
La concaténation des entrées diagonales et antidiagonales produit,
respectivement, \(77\) et \(22\). Par conséquent,
\[
 72+27=77+22=99,\qquad
 77-22=55=54+1.
\]
Ces égalités expriment deux partitions exactes du même bloc ; elles
n'identifient pas à elles seules une constante externe.

\begin{proposition}[Relation locale--globale du secteur radical]
\label{prop:app-local-global-radical}
La partition par lignes du bloc central reproduit, avec un facteur d'échelle
trois, la partition du secteur radical :
\[
 3\cdot72=216,\qquad
 3\cdot27=81,\qquad
 3\cdot99=297.
\]
\end{proposition}

\begin{proof}
Le bloc \(C_0\times C_0\) contient neuf représentants nuls et a pour somme
\(81\). Les quatre bras restants des lignes et des colonnes de \(C_0\)
ont pour somme \(216\). Les identités découlent de \(72+27=99\).
\end{proof}

De plus,
\[
M_c:=\frac{B_c}{\sqrt{45}}
=\exp\left(\frac12\log\frac95\,R\right)
\in SO^+(1,1).
\]
Si \(\eta=\operatorname{diag}(1,-1)\), la matrice normalisée \(M_c\) satisfait
\(M_c^{\mathsf T}\eta M_c=\eta\), \(\det M_c=1\) et appartient à
\(SO^+(1,1)\). Cette réalisation finie se limite au bloc \(2\times2\)
construit ; son interprétation spatio-temporelle requiert le pont dimensionnel
développé dans le traité de mécanique dimensionnelle.

\begin{proposition}[Famille spectrale APP--TRIT et lecture
Markov--Lorentz locale]
\label{prop:app-trit-espectral-markov-lorentz-rev8}
Pour \(\tau\in\{-1,0,+1\}\), soit
\[
 B_\tau=(4+3\tau)I+(5-3\tau)R.
\]
Alors
\[
 \operatorname{Spec}(B_\tau)=\{9,\,6\tau-1\},
 \qquad
 B_\tau=9P_++(6\tau-1)P_-,
\]
et le mode différentiel récupère l'orientation au moyen de
\[
 \tau=\frac{\lambda_-+1}{6}.
\]
En particulier, \(B_{+1}=B_c\). Sa normalisation stochastique
\[
 P_B:=\frac{1}{9}B_c=P_++\frac59P_-
\]
satisfait, pour tout \(n\geq0\),
\[
 P_B^n=P_++\left(\frac59\right)^n P_-.
\]
Par conséquent, la contraction du mode antisymétrique est
\(\kappa_B=5/9\) et l'écart spectral de cette chaîne locale est \(4/9\).
Si
\[
 \eta_c:=\frac12\log\frac95,
\]
la normalisation lorentzienne du même bloc satisfait
\[
 \frac{B_c}{\sqrt{45}}=\exp(\eta_c R),
 \qquad
 \boxed{\frac59=e^{-2\eta_c}}.
\]
\end{proposition}

\begin{proof}
Comme \(R\) agit avec les valeurs propres \(+1\) et \(-1\) sur
\(\operatorname{im}P_+\) et \(\operatorname{im}P_-\), respectivement, les
deux valeurs propres de \(B_\tau\) sont
\[
 (4+3\tau)+(5-3\tau)=9,\qquad
 (4+3\tau)-(5-3\tau)=6\tau-1.
\]
La formule des puissances découle de
\(P_\pm^2=P_\pm\) et de \(P_+P_-=0\). Enfin,
\(e^{-2\eta_c}=e^{-\log(9/5)}=5/9\).
\end{proof}

Cette coïncidence identifie deux lectures exactes du bloc local : la
mémoire différentielle décroît avec le même quotient que celui qui encode la rapidité
dans la réalisation déployée. Elle ne transforme pas \(B_\tau\) en le noyau de
Markov global du TPK, ni \(5/9\) en une constante physique. Elle n'étend pas non plus
à elle seule la réalisation \(SO^+(1,1)\) à un espace de Minkowski
\(3+1\) : ces étapes requièrent des applications dimensionnelles supplémentaires.

\section{Courbure discrète induite par l'interaction des feuillets somme et produit et son défaut de commutation}

Sur le tore, nous écrivons, maintenant en résidus,
\[
S(x,y)=x+y,
\qquad
P(x,y)=xy\pmod9.
\]
Pour toute fonction \(T:X_9\to\ZZ/9\ZZ\), nous fixons les différences progressives
\[
 \Delta_xT(x,y)=T(x+1,y)-T(x,y),\qquad
 \Delta_yT(x,y)=T(x,y+1)-T(x,y),
\]
où les coordonnées sont prises modulo neuf. Cette convention fixe également
l'orientation positive de chaque plaquette : on parcourt d'abord la direction
\(x\), puis la direction \(y\), et l'on revient par les arêtes opposées.
Il est important de typer correctement ces deux objets. \(S\) et \(P\) sont
des potentiels de sommet ; les variables de lien sont leurs différences
discrètes
\[
A_x^T=\Delta_xT,
\qquad
A_y^T=\Delta_yT.
\]
Pour un mélange ordonné de potentiels \((T_x,T_y)\), nous définissons la
courbure de plaquette par
\[
F(T_x,T_y)
=\Delta_x A_y^{T_y}-\Delta_y A_x^{T_x}
=\Delta_x\Delta_yT_y-\Delta_y\Delta_xT_x.
\]
Comme
\[
\Delta_xS=\Delta_yS=1,
\qquad
\Delta_xP=y,
\qquad
\Delta_yP=x,
\]
on obtient sur les \(81\) plaquettes :
\[
F(S,S)=F(P,P)=0,
\qquad
F(S,P)=+1,
\qquad
F(P,S)=-1\pmod9.
\]
Les connexions associées à une unique observable sont plates ; la composition
mixte des deux observables produit une courbure de signe opposé selon
l'ordre. Pour un rectangle de coordonnées \(R_{a,b}\) contenu
dans la carte, nous définissons le Wilson additif de bord comme la somme orientée des
variables de lien. L'annulation des arêtes intérieures donne
\[
 W_{\partial R_{a,b}}(S,P)=ab,\qquad
 W_{\partial R_{a,b}}(P,S)=-ab\pmod9,
\]
qui est l'identité de Stokes finie pour cette convention.

La distinction entre potentiel et lien est algébriquement nécessaire. Si l'on remplaçait
directement \(A_x,A_y\) par \(S,P\), des termes dépendant de
\(x,y\) apparaîtraient et les valeurs constantes précédentes ne seraient pas obtenues. La formulation
typée est celle qui produit le cocycle mixte utilisé plus loin dans la fermeture
de Heisenberg finie.

\begin{remark}
La différence intégrée \(54\) et la courbure mixte \(\pm1\) sont des invariants
distincts. La première compare des sommes globales ; la seconde dépend de l'ordre des
deux potentiels et se localise sur chaque plaquette. La formulation au moyen
d'opérateurs autoadjoints et de commutateurs sera développée après la construction de
l'espace d'états correspondant.
\end{remark}

\section{Développement de l'opérateur de transition par chemins}

La définition d'APP comprend l'espace
\(\operatorname{Path}(\GraphAPP)\). Une fois fixé un opérateur de
transition \(U\) sur les sommets de \(\GraphAPP\), nous supposons que
\(U_{yx}=0\) dès que \(x\to y\) n'est pas une arête du graphe. Sous cette
hypothèse de support, ses puissances admettent le développement
\begin{proposition}[Développement fini par chemins]
\label{prop:app-suma-caminos}
Pour chaque \(R\ge1\) et chaque paire de sommets \(i,f\),
\[
(U^R)_{fi}
=\sum_{\substack{\gamma:i\to f\\|\gamma|=R}}
\prod_{k=0}^{R-1}U_{x_{k+1},x_k},
\qquad
\gamma=(x_0=i,x_1,\ldots,x_R=f).
\]
\end{proposition}

\begin{proof}
Pour \(R=1\), l'identité est la définition des éléments de matrice de
\(U\). Si elle est vraie pour \(R\), alors
\[
 (U^{R+1})_{fi}
 =\sum_j U_{fj}(U^R)_{ji}.
\]
Substituer l'hypothèse de récurrence concatène à chaque chemin de longueur \(R\)
de \(i\) à \(j\) la dernière arête \(j\to f\) ; ainsi, chaque chemin de
longueur~\(R+1\) apparaît exactement une fois.\qedhere
\end{proof}

Par conséquent, APP réunit en une seule structure le tore arithmétique fini,
les observables de somme et de produit, la décomposition résidu--quotient, le
radical nilpotent et une courbure mixte orientée. Le représentant \(9\) de
la classe nulle fixe une carte visible du quotient ; le remplacer par \(0\) sans
transporter simultanément la carte modifie les registres utilisés par la
dynamique ultérieure.
\fi
 \end{HMTScientificOwnerSubordinated}
\endgroup

% Unidad U003. Redacción autónoma; tablas y censos reconstruidos desde 1063/live.

\section{Matrice arithmétique de l’APP et structure de l’anneau local}
\label{p12-83-c1-s1:chap:app-hojas-algebra-local}

La \cref{p12-83-c1-s2:chap:app-completacion-gnomonica} a construit un unique support
\(X_9\). L’APP possède une unique matrice arithmétique principale : la multiplication
nonadique publiée par racine numérique positive. La somme ne constitue pas une
deuxième table APP ; elle est la loi interne d’accumulation qui engendre les lignes et
les colonnes de la matrice multiplicative et fournit une observable auxiliaire de
comparaison. Cette distinction sera maintenue tout au long de l’ouvrage.

\subsection{Matrice multiplicative nonadique et racine numérique}

Pour \(a,b\in\mathcal D_9=\{1,\ldots,9\}\), l’entrée de l’APP s’obtient
directement par
\begin{equation}
 \boxed{\Pi(a,b)=\operatorname{dr}_9(ab)
 =1+((ab-1)\bmod9).}
 \label{p12-83-c1-s1:eq:app-matriz-multiplicativa-directa}
\end{equation}
L’unique matrice APP est donc
\[
\begin{array}{c|ccccccccc}
\Pi&1&2&3&4&5&6&7&8&9\\\hline
1&1&2&3&4&5&6&7&8&9\\
2&2&4&6&8&1&3&5&7&9\\
3&3&6&9&3&6&9&3&6&9\\
4&4&8&3&7&2&6&1&5&9\\
5&5&1&6&2&7&3&8&4&9\\
6&6&3&9&6&3&9&6&3&9\\
7&7&5&3&1&8&6&4&2&9\\
8&8&7&6&5&4&3&2&1&9\\
9&9&9&9&9&9&9&9&9&9
\end{array}
\]
Chaque ligne peut également être reconstruite par addition répétée de son indice,
mais cette récurrence décrit la loi interne de la matrice et non un deuxième
feuillet. Les lignes et colonnes d’indices \(1,2,4,5,7,8\) parcourent les neuf
classes ; celles d’indices \(3,6,9\) sont exactement celles qui contractent le support.

Dans la carte résiduelle, l’ensemble des unités est
\[
 R_9^\times=(\mathbb Z/9\mathbb Z)^\times
   =\{[1],[2],[4],[5],[7],[8]\}.
\]
L’idéal maximal est
\[
 \mathfrak m=(3)=\{[0],[3],[6]\}.
\]

\begin{theorem}[Classification des lignes multiplicatives]
\label{p12-83-c1-s1:thm:app-clasificacion-filas-producto}
Pour \(a\in\mathcal D_9\), l'application
\[
 m_a:\mathcal D_9\longrightarrow\mathcal D_9,
 \qquad b\longmapsto\Pi(a,b)
\]
est bijective si et seulement si \([a]_9\in R_9^\times\). Si
\([a]_9\in\mathfrak m\setminus\{0\}\), son image a trois éléments. Si
\([a]_9=[0]_9\), son image a un seul élément.
\end{theorem}

\begin{proof}
La réduction résiduelle de \(m_a\) est la multiplication par \([a]_9\). Dans un
anneau fini, cette multiplication est bijective exactement lorsque \([a]_9\)
est une unité. Si \([a]_9=[3]\) ou \([6]\), son image est l’idéal
\(\mathfrak m\), de cardinal trois. Si \([a]_9=[0]\), toute entrée est envoyée
sur la classe nulle, publiée comme \(9\).
\end{proof}

\begin{proposition}[Latinité de la loi additive interne]
\label{p12-83-c1-s1:prop:app-latinidad-aditiva}
L’observable auxiliaire
\[
 \Sigma(a,b)=\operatorname{dr}_9(a+b)
\]
agit par translations : pour chaque \(a,c\in\mathcal D_9\), il existe un unique
\(b\in\mathcal D_9\) tel que \(\Sigma(a,b)=c\), et de même pour l’autre
coordonnée.
\end{proposition}

\begin{proof}
Dans la carte résiduelle, \([a+b]_9=[c]_9\) a l’unique solution
\([b]_9=[c-a]_9\). La section positive transporte cette classe vers un unique
élément de \(\mathcal D_9\).
\end{proof}

La différence de géométrie des fibres est ainsi typée sans dupliquer la matrice
APP : la loi additive interne ne fait que translater, tandis que la matrice multiplicative
permute sur les unités et contracte sur le radical.

\begin{proposition}[Structure linéaire du radical]
\label{p12-83-c1-s1:prop:app-radical-lineal}
L'application
\[
 M_3:R_9\longrightarrow R_9,
 \qquad x\longmapsto3x
\]
satisfait
\[
 \ker M_3=\operatorname{im}M_3=\mathfrak m.
\]
L’application induite
\[
 R_9/\mathfrak m\longrightarrow\mathfrak m,
 \qquad [x]\longmapsto3x
\]
est un isomorphisme d’espaces vectoriels sur \(\mathbb F_3\), mais non un
isomorphisme d’anneaux.
\end{proposition}

\begin{proof}
La congruence \(3x\equiv0\pmod9\) équivaut à \(x\equiv0\pmod3\), d’où
le noyau \(\mathfrak m\). L’image est
\(\{[0],[3],[6]\}=\mathfrak m\). Le premier théorème d’isomorphisme donne
l’isomorphisme linéaire. Ce n’est pas un isomorphisme d’anneaux parce que
\(\mathfrak m^2=0\) dans \(R_9\), tandis que le quotient
\(R_9/\mathfrak m\cong\mathbb F_3\) a une unité non nulle.
\end{proof}

La section positive publie le cycle radical non trivial comme
\[
 369\longmapsto693\longmapsto936\longmapsto369,
\]
selon la position à partir de laquelle on lit le triplet. Cette notation conserve
l’ordre des coordonnées ; elle ne transforme pas l’idéal en trois nombres entiers
indépendants.

\subsection{Relèvements et cocycle de retenue}

Soit \(s_0:R_9\to\{0,\ldots,8\}\) la section résiduelle normalisée. La
retenue binaire de la somme est
\[
 \kappa_0(a,b)
 =\frac{s_0(a)+s_0(b)-s_0(a+b)}9.
\]

\begin{lemma}[Identité cocyclique]
\label{p12-83-c1-s1:lem:app-identidad-cociclo-acarreo}
Pour \(a,b,c\in R_9\),
\[
 \kappa_0(a,b)+\kappa_0(a+b,c)
 =\kappa_0(b,c)+\kappa_0(a,b+c).
\]
\end{lemma}

\begin{proof}
Multiplier par neuf et développer l’un quelconque des deux membres produit
\[
 s_0(a)+s_0(b)+s_0(c)-s_0(a+b+c).
\]
Par conséquent, ils sont égaux.
\end{proof}

Nous définissons \(\chi_0(a)=1\) si \(a=[0]_9\) et \(0\) sinon. Comme
\[
 s_+(a)=s_0(a)+9\chi_0(a),
\]
les retenues des deux sections sont reliées par
\[
 \kappa_+(a,b)=
 \kappa_0(a,b)+\chi_0(a)+\chi_0(b)-\chi_0(a+b).
\]
Le changement de représentant modifie le cocycle par un cobord. C’est la
manière exacte de transporter la mémoire lorsque la classe nulle est publiée comme
\(9\) au lieu de \(0\).

\subsection{Recensements intégrés des résidus et des quotients}

En sommant sur toutes les paires \((i,j)\in\mathcal D_9^2\), on obtient
\[
 \sum_{i,j}(i+j)=810,
 \qquad
 \sum_{i,j}ij=2025.
\]
Les résidus positifs satisfont
\[
 \sum_{i,j}R^+_{i,j}=405,
 \qquad
 \sum_{i,j}R^\times_{i,j}=459.
\]
Par les identités relevées,
\[
 \sum_{i,j}Q^+_{i,j}=\frac{810-405}{9}=45,
 \qquad
 \sum_{i,j}Q^\times_{i,j}=\frac{2025-459}{9}=174.
\]

\begin{theorem}[Séparation exacte de l’excès somme--produit]
\label{p12-83-c1-s1:thm:app-separacion-exceso}
La différence intégrée entre produit et somme se décompose de manière unique
en une partie publiée et une partie de quotient :
\[
 \boxed{
  2025-810
  =(459-405)+9(174-45)
  =54+9\cdot129.}
\]
\end{theorem}

\begin{proof}
Nous soustrayons, terme à terme, les identités
\(ij=R^\times_{i,j}+9Q^\times_{i,j}\) et
\(i+j=R^+_{i,j}+9Q^+_{i,j}\), puis sommons sur les quatre-vingt-une paires. Les
recensements précédents donnent l’égalité. L’unicité procède de la division
euclidienne dans chaque paire.
\end{proof}

La valeur \(54\) n’est pas la différence complète entre deux matrices supposées.
C’est le défaut visible entre la matrice multiplicative et l’accumulateur additif
dans la carte positive. Le terme \(9\cdot129\) est la contribution
qui demeure dans les quotients. L’omettre reviendrait à confondre publication et
généalogie.

\subsection{Bloc central de \(2\) positions}

Pour \(k\in\mathbb Z/9\mathbb Z\), nous considérons le bloc symétrique
\[
 B_k=
 \begin{pmatrix}
  k^2&k(k+1)\\
  k(k+1)&(k+1)^2
 \end{pmatrix}\pmod9.
\]
Les diagonales coïncident exactement lorsque
\[
 k^2\equiv(k+1)^2\pmod9
 \quad\Longleftrightarrow\quad
 2k+1\equiv0\pmod9.
\]
Comme \(2^{-1}=5\pmod9\), l’unique solution est \(k=4\). Ainsi apparaît le
bloc
\[
 B_c=
 \begin{pmatrix}7&2\\2&7\end{pmatrix}.
\]
Soit
\[
 R=\begin{pmatrix}0&1\\1&0\end{pmatrix},
 \qquad
 P_\pm=\frac12(I\pm R).
\]
Alors
\[
 B_c=7I+2R=9P_++5P_-.
\]

\begin{theorem}[Unicité et spectre du bloc central]
\label{p12-83-c1-s1:thm:app-bloque-central}
Le bloc \(B_c\) est l’unique bloc adjacent de la famille \(B_k\) ayant
des diagonales égales. Ses valeurs propres dans les sous-espaces symétrique et
antisymétrique sont, respectivement, \(9\) et \(5\).
\end{theorem}

\begin{proof}
La congruence précédente prouve l’unicité. Puisque \(R\) agit par
\(+1\) sur \(\operatorname{im}P_+\) et par \(-1\) sur
\(\operatorname{im}P_-\), la matrice \(7I+2R\) agit par \(9\) et \(5\),
respectivement.
\end{proof}

Les deux lectures ordonnées de ses chiffres donnent
\[
 72+27=99,
 \qquad
 72-27=45=\det B_c,
\]
et les deux diagonales donnent
\[
 77+22=99,
 \qquad
 77-22=55=54+1.
\]
Ces identités sont des contrôles internes du bloc ; elles ne définissent pas à elles seules
le défaut global.

\subsection{Courbure mixte des \(2\) lois}

Dans la carte résiduelle, nous écrivons
\[
 S(x,y)=x+y,
 \qquad
 P(x,y)=xy.
\]
Pour une fonction \(T:X_9\to R_9\), soient
\[
 \Delta_xT(x,y)=T(x+1,y)-T(x,y),
 \qquad
 \Delta_yT(x,y)=T(x,y+1)-T(x,y).
\]
Nous définissons la courbure ordonnée d’une paire de lois par
\[
 \mathcal F(T_x,T_y)
 :=\Delta_x\Delta_yT_y-\Delta_y\Delta_xT_x.
\]

\begin{proposition}[Orientations de la composition mixte]
\label{p12-83-c1-s1:prop:app-curvatura-mixta}
On a
\[
 \mathcal F(S,S)=0,
 \qquad
 \mathcal F(P,P)=0,
\]
et, pour les compositions mixtes,
\[
 \mathcal F(S,P)=+1,
 \qquad
 \mathcal F(P,S)=-1
 \quad\text{dans }R_9.
\]
\end{proposition}

\begin{proof}
Les différences secondes mixtes de la loi linéaire \(S\) s’annulent. Pour
\(P(x,y)=xy\), les deux différences mixtes pures valent un et s’annulent lorsqu’elles sont
comparées dans le même ordre. En combinant une loi linéaire et une loi bilinéaire,
il reste une unique unité dont le signe change lorsque l’ordre des feuillets est inversé.
\end{proof}

La courbure ne crée pas une autre matrice APP. Elle enregistre que l’ordre
\(\Sigma\to\Pi\) et l’ordre \(\Pi\to\Sigma\) sont des orientations opposées,
tandis que les secteurs purs restent plats. Cette triade
\(-1,0,+1\) sera typée dans l’unité TRIT.

\subsection{Obstructions de la matrice arithmétique de l’APP : translations, radical local, bloc central, courbure mixte et conservation de la retenue}

\begin{criterion}[Critères de réfutation de la matrice APP]
L’unité est réfutée si :
\begin{enumerate}
 \item la loi additive interne cesse d’agir par translations ;
 \item une ligne multiplicative non unitaire est déclarée bijective ;
 \item on identifie \(R_9/\mathfrak m\) à \(\mathfrak m\) comme anneau ;
 \item l’excès \(1215\) est remplacé par \(54\) sans conserver les quotients ;
 \item le bloc central est choisi sans résoudre \(2k+1=0\pmod9\) ;
 \item l’inversion de l’ordre mixte ne change pas le signe de la courbure ;
 \item la retenue est éliminée lors du passage entre les sections \(s_0\) et
       \(s_+\).
\end{enumerate}
\end{criterion}

L’APP est maintenant construite comme support, catégorie de trajectoires, matrice
multiplicative nonadique, loi additive interne, radical local, cocycle de
retenue et orientation mixte. L’unité
suivante ne résumera pas ces structures : elle prendra la triade d’orientations
qui vient d’émerger et construira l’état complet du TRIT avec ses trois régimes
quadratiques.


% Unidad U002. Redacción autónoma a partir del generador integral 1063/live.

\section{Complétion gnomonique et support arithmétique de l’APP}
\label{p12-83-c1-s2:chap:app-completacion-gnomonica}

La présente unité construit le support de l’APP à partir de la cellule
nonadique de la \cref{p12-83-c1-s4:chap:segmentos-marcas-intersticios}. Aucun
plan continu, réseau infini ni constante numérique n’est présupposé. L’objet de
départ est fini : deux copies ordonnées de la carte positive de
\(\mathbb Z/9\mathbb Z\). La complétion de leurs interstices intérieurs
produit exactement quatre-vingt-un états, et la loi de translation résiduelle
transforme ce support carré en un tore discret.

Dans toute l’unité, APP désigne une structure mathématique et non une image.
Le sigle APP nomme désormais cette structure ; les propriétés qui
suivent dépendent de ses objets et de ses morphismes, non d’une explicitation verbale du
sigle.

\subsection{Correspondance entre \(2\) cartes d’une même classe résiduelle}

On définit
\[
 \mathcal D_9:=\{1,2,\ldots,9\},
 \qquad
 \mathcal D_8:=\{1,2,\ldots,8\}.
\]
L’application de publication résiduelle est
\[
 \lambda_9:\mathcal D_9\longrightarrow\mathbb Z/9\mathbb Z,
 \qquad
 \lambda_9(a)=[a]_9.
\]
Par conséquent, \(\lambda_9(9)=[0]_9\). Son inverse comme application
d’ensembles est la section positive
\[
 s_+([a]_9)=
 \begin{cases}
  a,&[a]_9\neq[0]_9,\quad 1\leq a\leq8,\\
  9,&[a]_9=[0]_9.
 \end{cases}
\]
Les deux expressions \(9\) et \([0]_9\) ne sont pas deux phases : ce sont deux
représentants du même élément dans des cartes différentes.

Pour deux coordonnées, nous écrirons
\[
 \lambda_9^{\times2}:\mathcal D_9^2
       \longrightarrow X_9:=(\mathbb Z/9\mathbb Z)^2,
 \qquad
 (a,b)\longmapsto([a]_9,[b]_9).
\]
Cette application est une bijection d’ensembles. La carte
\(\mathcal D_9^2\) rend visibles les marques positives ; le groupe \(X_9\)
rend visibles la somme résiduelle et les translations.

\begin{lemma}[Changement de carte sans perte d’état]
\label{p12-83-c1-s2:lem:app-cambio-carta-sin-perdida}
Les applications \(\lambda_9^{\times2}\) et
\(s_+^{\times2}\) sont inverses. En particulier, l’écriture positive et
l’écriture résiduelle contiennent exactement les mêmes quatre-vingt-un états.
\end{lemma}

\begin{proof}
Dans chaque coordonnée, on a
\(\lambda_9\circ s_+=\operatorname{id}_{\mathbb Z/9\mathbb Z}\) et
\(s_+\circ\lambda_9=\operatorname{id}_{\mathcal D_9}\). Le produit
cartésien de ces identités fournit les deux identités requises.
\end{proof}

\subsection{Noyau intérieur et complément gnomonique}

Le noyau intérieur est
\[
 \mathcal N_8:=\mathcal D_8\times\mathcal D_8,
 \qquad |\mathcal N_8|=8^2=64.
\]
Dans le plan des indices, \(\mathcal N_8\) est le carré qui n’utilise la
marque de clôture dans aucune coordonnée. Pour publier également la classe nulle,
il faut compléter la neuvième ligne et la neuvième colonne. Nous définissons le
complément gnomonique disjoint par
\[
 \mathcal G_9:=
  (\{9\}\times\mathcal D_9)
  \sqcup
  (\mathcal D_8\times\{9\}).
\]
Le coin \((9,9)\) appartient à la première pièce ; c’est pourquoi la seconde utilise
\(\mathcal D_8\) et non \(\mathcal D_9\). En conséquence,
\[
 |\mathcal G_9|=9+8=17,
 \qquad
 \mathcal D_9^2=\mathcal N_8\sqcup\mathcal G_9,
 \qquad
 81=64+17.
\]

\begin{proposition}[Minimalité de la complétion]
\label{p12-83-c1-s2:prop:app-minimalidad-completacion}
Soit \(Y\) un sous-ensemble de \(\mathcal D_9^2\) qui contient
\(\mathcal N_8\) et dont la projection sur chaque coordonnée est tout
\(\mathcal D_9\). Si, de plus, \(Y\) contient toutes les paires nécessaires
pour translater une coordonnée tandis que l’autre reste arbitraire,
alors \(\mathcal G_9\subseteq Y\). Par conséquent, l’ajout de dix-sept
états est la complétion minimale compatible avec les deux translations.
\end{proposition}

\begin{proof}
Pour translater la première coordonnée à travers la classe de clôture tandis que
la seconde reste à une valeur arbitraire \(b\in\mathcal D_9\), toutes les paires
\((9,b)\) sont nécessaires. On obtient neuf états. Pour
translater la seconde coordonnée à travers la clôture tandis que la première prend
une valeur arbitraire \(a\in\mathcal D_8\), les huit paires
\((a,9)\) restantes sont nécessaires. Le coin \((9,9)\) a déjà été compté dans la première
famille. Au total, \(9+8=17\) états sont nécessaires, exactement ceux de
\(\mathcal G_9\).
\end{proof}

L’hypothèse de compatibilité avec les deux translations est essentielle. Il suffirait
d’ajouter une seule marque par coordonnée pour que les deux projections soient
surjectives, mais cet ensemble ne contiendrait pas les états de frontière
nécessaires pour déplacer une coordonnée en gardant l’autre fixe. L’APP exige le
produit complet, non deux axes déconnectés.

\subsection{Réalisation exacte de l’intervalle interstitiel par des cellules APP et des coordonnées de frontière}

Soit \(P_9\) le chemin ordonné à neuf marques et huit arêtes intérieures
\[
 E(P_9)=\{e_1,\ldots,e_8\}.
\]
La clôture des extrémités incorpore une arête \(e_9\) et produit le cycle
\(C_9\) :
\[
 E(C_9)=E(P_9)\sqcup\{e_9\}.
\]
Les paires d’interstices intérieurs forment \(E(P_9)^2\). Les paires
d’interstices du cycle complet forment \(E(C_9)^2\). La différence est
\begin{align}
 E(C_9)^2\setminus E(P_9)^2
   &=(\{e_9\}\times E(C_9))
     \sqcup(E(P_9)\times\{e_9\}).
\label{p12-83-c1-s2:eq:app-complemento-intersticial}
\end{align}
La décomposition est disjointe parce que le coin \((e_9,e_9)\) est réservé à
la première pièce.

\begin{theorem}[Équivalence du modèle de marques et du modèle
d’interstices]
\label{p12-83-c1-s2:thm:app-equivalencia-modelos}
L'application
\[
 \beta:E(C_9)^2\longrightarrow\mathcal D_9^2,
 \qquad
 \beta(e_i,e_j)=(i,j),
\]
est bijective. Ses restrictions satisfont
\[
 \beta(E(P_9)^2)=\mathcal N_8,
 \qquad
 \beta(E(C_9)^2\setminus E(P_9)^2)=\mathcal G_9.
\]
Par conséquent, l’identité \(81=64+17\) compte simultanément des paires de
marques publiées et des paires d’interstices fermés.
\end{theorem}

\begin{proof}
Chaque arête de \(C_9\) a un indice unique dans \(\mathcal D_9\).
L’application \((i,j)\mapsto(e_i,e_j)\) est l’inverse explicite de \(\beta\).
Si \(i,j\leq8\), la paire appartient à \(E(P_9)^2\) et son image appartient à
\(\mathcal D_8^2\). Si au moins l’un des indices est neuf, la paire appartient
au complément de \cref{p12-83-c1-s2:eq:app-complemento-intersticial} ; séparer le cas
\(i=9\) du cas \(i\leq8,j=9\) produit exactement les deux pièces de
\(\mathcal G_9\).
\end{proof}

\subsection{Clôture toroïdale et translations}

Soient
\[
 \mathbf e_1=([1]_9,[0]_9),
 \qquad
 \mathbf e_2=([0]_9,[1]_9).
\]
L’ensemble des directions élémentaires est
\[
 \mathscr D_{\rm APP}:=
 \{+\mathbf e_1,-\mathbf e_1,+\mathbf e_2,-\mathbf e_2\}.
\]
Nous définissons le graphe de Cayley orienté
\[
 \Gamma_{\rm APP}:=\operatorname{Cay}(X_9,\mathscr D_{\rm APP}).
\]
Ses sommets sont les quatre-vingt-un éléments de \(X_9\). Pour chaque
\(x\in X_9\) et \(d\in\mathscr D_{\rm APP}\), il existe une arête orientée
\(x\xrightarrow{d}x+d\). L’arête inverse porte l’étiquette \(-d\).

Dans la carte positive, la translation unitaire s’écrit au moyen du successeur
publié et de sa retenue :
\[
 u_+(a)=
 \begin{cases}
  a+1,&a<9,\\
  1,&a=9,
 \end{cases}
 \qquad
 \kappa_+(a)=
 \begin{cases}
  0,&a<9,\\
  1,&a=9.
 \end{cases}
\]
L’identité relevée est
\begin{equation}
 a+1=u_+(a)+9\kappa_+(a).
\label{p12-83-c1-s2:eq:app-sucesor-publicado-acarreo}
\end{equation}
En appliquant \(\lambda_9\), le terme de retenue s’annule et il reste
\[
 \lambda_9(u_+(a))=\lambda_9(a)+[1]_9.
\]

\begin{lemma}[Compatibilité entre clôture, résidu et relèvement]
\label{p12-83-c1-s2:lem:app-compatibilidad-cierre-residuo}
L’ajout de \(e_9\), l’identification \(9\sim[0]_9\) et la translation
\(x\mapsto x+\mathbf e_i\) décrivent trois opérations compatibles dans
des domaines distincts. Le quotient résiduel publie le retour ; la retenue de
\cref{p12-83-c1-s2:eq:app-sucesor-publicado-acarreo} conserve le relèvement que le
quotient omet.
\end{lemma}

\begin{proof}
La clôture topologique incorpore l’arête qui relie les extrémités du chemin. La
carte résiduelle identifie l’extrémité publiée à la classe nulle. Si
\(a<9\), le successeur a une retenue nulle ; si \(a=9\), l’identité relevée
est \(10=1+9\). Les deux cas descendent à la même somme résiduelle, mais la
valeur de \(\kappa_+\) les distingue.
\end{proof}

\begin{corollary}[Support toroïdal]
\label{p12-83-c1-s2:cor:app-soporte-toroidal}
Identifier les côtés opposés de la carte \(\mathcal D_9^2\) produit une
réalisation géométrique du groupe \(X_9\). Chaque translation élémentaire est
définie en tout sommet et possède un inverse.
\end{corollary}

\subsection{Catégorie des trajectoires APP : objets, morphismes de prolongement et composition compatible}

\begin{definition}[Trajectoire orientée]
Une trajectoire de longueur \(n\geq0\) est une paire
\[
 \gamma=(x_0;d_1,\ldots,d_n),
 \qquad
 x_0\in X_9,\qquad d_k\in\mathscr D_{\rm APP}.
\]
Ses sommets sont
\[
 x_k=x_0+\sum_{r=1}^k d_r,
 \qquad 0\leq k\leq n.
\]
La source est \(s(\gamma)=x_0\) et la cible est
\(t(\gamma)=x_n\). Pour \(n=0\), on obtient la trajectoire identité en
\(x_0\).
\end{definition}

Si \(t(\gamma)=s(\delta)\), leur concaténation est
\[
 \delta\circ\gamma
  =(x_0;d_1,\ldots,d_n,e_1,\ldots,e_m),
\]
où \(e_1,\ldots,e_m\) sont les directions de \(\delta\). La
concaténation est associative et les trajectoires de longueur zéro sont des
identités. Nous obtenons ainsi la catégorie libre
\[
 \mathsf{Path}(\Gamma_{\rm APP}).
\]

\begin{proposition}[Finitude du support et dénombrabilité des trajectoires]
\label{p12-83-c1-s2:prop:app-soporte-finito-trayectorias-numerables}
L’ensemble des états de l’APP est fini, de cardinal \(81\). Pour une source
fixée, il existe \(4^n\) mots de directions de longueur \(n\). L’union
de toutes les trajectoires finies est dénombrable.
\end{proposition}

\begin{proof}
La première affirmation procède de \(|X_9|=9^2\). À chacun des \(n\)
pas, il y a quatre choix indépendants, d’où \(4^n\). Enfin,
\(\bigcup_{n\geq0}X_9\times\mathscr D_{\rm APP}^n\) est une union dénombrable
d’ensembles finis.
\end{proof}

Cette proposition sépare deux niveaux qui ne seront pas identifiés : l’APP possède un
support fini, mais admet des trajectoires de toute longueur. Le
prolongement d’une route ne crée pas de nouveaux sommets du support ; il crée une nouvelle
généalogie sur ceux-ci.

\subsection{Matrice multiplicative et loi additive interne sur un unique support}

Pour \(n\in\mathbb Z_{>0}\), la racine numérique positive est
\[
 \operatorname{dr}_9(n):=1+((n-1)\bmod9).
\]
Sur la carte positive, nous définissons la matrice principale et la loi interne
\[
 \Sigma(a,b):=\operatorname{dr}_9(a+b),
 \qquad
 \Pi(a,b):=\operatorname{dr}_9(ab).
\]
Les deux applications ont le même domaine \(\mathcal D_9^2\) et le même
codomaine \(\mathcal D_9\), mais elles n’ont pas le même statut : \(\Pi\) publie
l’unique matrice APP et \(\Sigma\) engendre ses accumulations. Leurs versions résiduelles
sont
\[
 \bar\Sigma([a]_9,[b]_9)=[a+b]_9,
 \qquad
 \bar\Pi([a]_9,[b]_9)=[ab]_9.
\]
Le changement de carte est compatible avec les deux opérations :
\begin{align*}
 \lambda_9\circ\Sigma
   &=\bar\Sigma\circ\lambda_9^{\times2},\\
 \lambda_9\circ\Pi
   &=\bar\Pi\circ\lambda_9^{\times2}.
\end{align*}

Pour ne pas perdre le quotient entier, on introduit les relèvements
\begin{align}
 a+b &= R^+_{a,b}+9Q^+_{a,b},
 &1\leq R^+_{a,b}\leq9,
\label{p12-83-c1-s2:eq:app-lift-suma}\\
 ab &= R^\times_{a,b}+9Q^\times_{a,b},
 &1\leq R^\times_{a,b}\leq9.
\label{p12-83-c1-s2:eq:app-lift-producto}
\end{align}
Ici
\[
 R^+_{a,b}=\Sigma(a,b),
 \qquad
 R^\times_{a,b}=\Pi(a,b),
\]
et les quotients sont déterminés de manière univoque par
\[
 Q^+_{a,b}=\frac{a+b-R^+_{a,b}}9,
 \qquad
 Q^\times_{a,b}=\frac{ab-R^\times_{a,b}}9.
\]

\begin{definition}[APP généalogique]
\label{p12-83-c1-s2:def:app-genealogica-autonoma}
La structure APP à cette étape est le tuple
\[
 \operatorname{APP}:=
 \bigl(
  X_9,
  \Gamma_{\rm APP},
  \mathsf{Path}(\Gamma_{\rm APP}),
  \Sigma,\Pi,
  R^+,Q^+,R^\times,Q^\times
 \bigr).
\]
La généalogie d’une trajectoire conserve, pour chaque sommet visité, la
position, la direction d’entrée, l’entrée multiplicative, l’accumulation
additive et leurs résidus et quotients respectifs.
\end{definition}

\begin{theorem}[Construction causale du support APP]
\label{p12-83-c1-s2:thm:app-cadena-constructiva-autonoma}
La chaîne
\[
 \mathcal D_9
 \longrightarrow
 \mathcal D_8^2
 \longrightarrow
 \mathcal D_8^2\sqcup\mathcal G_9
 \longrightarrow
 \mathcal D_9^2
 \xrightarrow{\lambda_9^{\times2}}
 X_9
 \longrightarrow
 \Gamma_{\rm APP}
 \longrightarrow
 \mathsf{Path}(\Gamma_{\rm APP})
 \longrightarrow
 (\Sigma,\Pi)
\]
construit, sans perte de cardinalité, le support, le mouvement, les trajectoires,
la matrice multiplicative de l’APP et sa loi additive interne.
\end{theorem}

\begin{proof}
La \cref{p12-83-c1-s2:prop:app-minimalidad-completacion} construit la carte complète à
partir du noyau. Le \cref{p12-83-c1-s2:thm:app-equivalencia-modelos} identifie son
modèle interstitiel. Le \cref{p12-83-c1-s2:lem:app-cambio-carta-sin-perdida} transporte
les quatre-vingt-un états vers \(X_9\). Le graphe de Cayley incorpore les quatre
translations réversibles sans changer les sommets. La catégorie libre
incorpore toutes les concaténations finies. Enfin, \(\Sigma\) et
\(\Pi\) sont des applications totales distinctes sur le même domaine ; les
équations \cref{p12-83-c1-s2:eq:app-lift-suma,p12-83-c1-s2:eq:app-lift-producto} conservent la partie
que la réduction résiduelle ne publie pas.
\end{proof}

\subsection{Décomposition de l’orthogramme APP en noyau \(8\times8\), gnomon de \(17\) états et carte \(9\times9\) à côtés opposés identifiés}

La \cref{fig:app-ortograma} doit être lue en trois couches : le carré intérieur
de soixante-quatre états, le gnomon de dix-sept états et
l’identification des côtés opposés. La carte de neuf sur neuf complète ainsi le
noyau de huit sur huit par le gnomon de clôture. La matrice
multiplicative et la loi additive interne sont évaluées en chaque sommet du même
état, conservent des relèvements résidu--quotient distincts et ne constituent pas
deux carrés APP.

\subsection{Obstructions à la fidélité du support APP : cardinalités \(81/17\), translations réversibles et séparation des feuillets somme et produit}

\begin{criterion}[Critères de réfutation du support APP]
La construction est réfutée si l’un quelconque des faits
suivants se produit :
\begin{enumerate}
 \item le support complet a un cardinal différent de \(81\) ;
 \item le complément gnomonique n’a pas pour cardinal \(17\) ;
 \item un état de \(X_9\) est dépourvu de l’une des quatre translations ;
 \item on identifie \(\Sigma\) et \(\Pi\) parce qu’ils coïncident en une valeur ;
 \item on élimine \(Q^+\) ou \(Q^\times\) lors du retour par une frontière ;
 \item les trajectoires de longueurs différentes sont présentées comme de nouveaux
       sommets du support ;
 \item on interprète \(9\) et \([0]_9\) comme des phases indépendantes.
\end{enumerate}
\end{criterion}

L’unité a construit l’APP jusqu’à son évaluation généalogique complète. Elle n’a pas encore
orienté la différence entre la matrice et son accumulateur, elle n’a pas
introduit le TRIT
et elle n’a pas défini le Topological Phase Kernel. L’unité suivante étudiera
la structure algébrique des lignes et des colonnes, le secteur des unités et le
cocycle de retenue qui prépare cette orientation.


% Unidad U005. Redacción autónoma desde 03c/03h/20 de la autoridad 1063/live.

\long\def\HMTDeferredTPKBase{%
\section{Topological Phase Kernel : base observable et état holonomique intégral}
\label{p12-83-c1-s3:chap:tpk-categoria-estado-enriquecido}

L’APP a fourni un support, une matrice multiplicative et une loi additive
interne ; TRIT a
fourni une orientation relevée. Le Topological Phase Kernel, ci-après
TPK, spécifie comment ces données sont transportées conjointement, quelles
coordonnées sont conservées lors de la clôture d’une route et quels lecteurs peuvent être publiés
sans identifier la publication à l’état. Son objet mathématique est une
catégorie sur une base de chemins, munie d’une famille explicite de
fibres et de relations de prolongement.

\subsection{La base double des chemins APP}

Soit \(\mathsf P_{\rm APP}\) la catégorie libre construite dans la
\cref{p12-83-c1-s2:chap:app-completacion-gnomonica}. Nous définissons
\[
 \mathsf P_{\rm APP}^{(2)}
 :=\mathsf P_{\rm APP}\times\mathsf P_{\rm APP}.
\]
Un objet est une paire de positions
\((p^\Sigma,p^\Pi)\in X_9^2\). Un morphisme est une paire de trajectoires ; l’une
d’elles peut être l’identité. La première composante est évaluée par
\(\Sigma\), la seconde par \(\Pi\). Ce ne sont pas deux réseaux : ce sont deux curseurs
typés sur le même support APP.

L’ensemble des directions cardinales est
\[
 \mathscr D_{\rm card}:=\{N,E,S,O\},
\]
avec décalages
\[
 \delta(N)=(-1,0),\quad
 \delta(E)=(0,1),\quad
 \delta(S)=(1,0),\quad
 \delta(O)=(0,-1)
 \quad\text{dans }X_9.
\]
L’involution d’orientation est
\[
 \overline N=S,
 \quad \overline S=N,
 \quad \overline E=O,
 \quad \overline O=E.
\]

\subsection{Phase nonadique, secteur dodécaphasique et calendrier}

L’horloge complète de cette étape est
\[
 \Theta_{108}:=\mathbb Z/108\mathbb Z.
\]
Pour \(\theta\in\Theta_{108}\), soit
\(\widetilde\theta\in\{0,\ldots,107\}\) son représentant normalisé.
Nous définissons trois coordonnées distinctes :
\begin{align*}
 g_0(\theta)&=[\widetilde\theta+1]_9\in C_9,\\
 r(\theta)&=1+(\widetilde\theta\bmod9)\in\mathcal D_9,\\
 \beta(\theta)&=
 \left[\left\lfloor\frac{\widetilde\theta}{9}\right\rfloor\right]_{12}
 \in C_{12}.
\end{align*}
La première est la classe de phase, la deuxième son étiquette positive et la troisième
le secteur dodécaphasique. L’application \(\beta\) est une application par blocs ;
ce n’est pas la réduction canonique \(\theta\bmod12\).

La signature et le feuillet actif sont
\[
 \tau(\theta):=\varepsilon_9(r(\theta))\in\mathsf{TRIT},
\]
et
\[
 \mu(\theta)=
 \begin{cases}
  \Sigma,&\tau(\theta)=+1,\\
  \Pi,&\tau(\theta)=-1,\\
  \varnothing,&\tau(\theta)=0.
 \end{cases}
\]
Le symbole \(\varnothing\) désigne un pas de seuil où aucun curseur
n’est translaté ; il ne désigne pas une absence d’état.

\begin{definition}[Configuration observable]
\label{p12-83-c1-s3:def:tpk-configuracion-observable}
La base observable est
\[
 \mathcal Q_{\rm obs}
 :=X_9^2\times\mathscr D_{\rm card}^2\times\Theta_{108}.
\]
Un élément s’écrit
\[
 q=(p^\Sigma,p^\Pi,d^\Sigma,d^\Pi,\theta).
\]
Il contient deux positions, deux orientations cardinales et un indice de
calendrier.
\end{definition}

\subsection{Sélecteur de feuillet \(M_{\rm ph}:\mathcal D_9\to\mathsf{TRIT}\) et évolution conjointe de la position, de l’orientation et du calendrier}

La fonction
\[
 M_{\rm ph}:\mathcal D_9\longrightarrow\mathsf{TRIT},
 \qquad M_{\rm ph}(r)=\varepsilon_9(r)
\]
est le sélecteur interne de feuillet. Il sélectionne l’évaluation qui agit ; il ne déplace
par lui-même aucune position.

Nous définissons l’indicateur d’inversion de l’orientation par
\[
 \eta_{27}(\theta)=
 \begin{cases}
  -1,&\widetilde\theta+1\in\{27,54,81,108\},\\
  +1,&\text{sinon}.
 \end{cases}
\]
Les nouvelles directions sont
\[
 d^{\nu,+}=
 \begin{cases}
  \overline{d^\nu},&\eta_{27}(\theta)=-1,\\
  d^\nu,&\eta_{27}(\theta)=+1,
 \end{cases}
 \qquad \nu\in\{\Sigma,\Pi\}.
\]
Les positions évoluent par
\[
 (p^{\Sigma,+},p^{\Pi,+})=
 \begin{cases}
 (p^\Sigma+\delta(d^\Sigma),p^\Pi),&\tau(\theta)=+1,\\
 (p^\Sigma,p^\Pi+\delta(d^\Pi)),&\tau(\theta)=-1,\\
 (p^\Sigma,p^\Pi),&\tau(\theta)=0.
 \end{cases}
\]

\begin{definition}[Évolution observable]
\label{p12-83-c1-s3:def:tpk-evolucion-observable}
L'application
\[
 F_{\rm obs}:\mathcal Q_{\rm obs}\longrightarrow\mathcal Q_{\rm obs}
\]
est défini par
\[
 F_{\rm obs}(q)=
 (p^{\Sigma,+},p^{\Pi,+},d^{\Sigma,+},d^{\Pi,+},[\theta+1]_{108}).
\]
La catégorie libre du graphe \(q\to F_{\rm obs}(q)\) est notée
\(\mathsf P_{\rm obs}\).
\end{definition}

Chaque flèche de \(\mathsf P_{\rm obs}\) induit une paire de trajectoires dans
l’APP, de sorte qu’il existe un foncteur
\[
 B:\mathsf P_{\rm obs}\longrightarrow
 \mathsf P_{\rm APP}^{(2)}.
\]
Si \(\tau=+1\), la deuxième composante de \(B\) est l’identité ; si
\(\tau=-1\), c’est la première ; si \(\tau=0\), les deux le sont.

\begin{proposition}[Séparation entre sélection, mouvement et observation]
\label{p12-83-c1-s3:prop:tpk-separacion-seleccion-movimiento}
Le sélecteur \(M_{\rm ph}\), les translations
\(T_d(p)=p+\delta(d)\) et tout échantillonnage temporel possèdent des domaines et des
codomaines distincts. Ce ne sont pas des opérateurs interchangeables.
\end{proposition}

\begin{proof}
\(M_{\rm ph}\) prend une étiquette de phase et renvoie une orientation tritique.
\(T_d\) est un automorphisme de \(X_9\). Un échantillonnage agit sur une
suite ou une catégorie de flèches. Identifier deux d’entre eux violerait le
typage. \(F_{\rm obs}\) les compose dans un ordre déclaré sans effacer cette
distinction.
\end{proof}

\subsection{Décomposition de l'état holonomique intégral en \(6\) familles fibrées}

Sur chaque \(q\in\mathcal Q_{\rm obs}\), nous fixons la fibre totale
\[
 \mathcal F_q=
 \mathsf O_q\times
 \mathsf H_q\times
 \mathsf C_q\times
 \mathsf S_q\times
 \mathsf B_q\times
 \Lambda_q.
\]
Chaque facteur a un type et une fonction propres.

\subsubsection{Détermination conjointe de l'orientation et du feuillet de l'état holonomique intégral}

\[
 \mathsf O_q\subseteq
 \mathscr D_{\rm card}^2\times\{+1,-1\}.
\]
Les deux directions coïncident avec celles stockées dans \(q\) ; la dernière
coordonnée enregistre l’orientation du feuillet lorsque la réalisation l’exige.
L’étiquette de feuillet ne remplace pas les directions cardinales.

\subsubsection{Décomposition euclidienne de l’état en résidu et quotient compatibles}

La première fibre de relèvement est
\[
 \mathsf H^{(1)}=\mathcal D_9\times\mathbb Z,
 \qquad
 \operatorname{rec}_9(r,u)=r+9u.
\]
Pour chaque entier \(n\), il existe d’uniques
\[
 r=\operatorname{dr}_9(n),
 \qquad
 u=\frac{n-r}{9}.
\]
Une profondeur supérieure consiste en une tour déclarée de ces fibres et de leurs
applications de troncature. Le terme « Hensel » ne remplace pas la définition de
cette tour.

\subsubsection{Opérateur de retenue et actualisation conjointe du résidu, du quotient et de la mémoire}

\(\mathsf C_q\) est un groupe abélien. Dans les lecteurs de blocs décimaux,
on prend \(\mathsf C_q=\mathbb Z\). L’accroissement n’est pas une coordonnée
indépendante, mais une cochaîne sur les flèches :
\[
 \kappa(r)\in\mathsf C_{q'}
 \quad\text{pour }r:q\to q'.
\]

\subsubsection{Construction de la signature entière à partir de l’orientation, de la retenue et de la frontière}

La signature est dérivée par une application typée
\[
 \operatorname{Sig}_q:
 \mathsf H_q\times\mathsf C_q\times\Lambda_q
 \longrightarrow\mathsf S_q.
\]
Dans la réalisation de frontière nonadique, on conserve individuellement
\[
 q_\partial,a_\partial,c_\partial\in\mathbb F_3^6,
 \qquad
 \operatorname{colw}\in\mathbb Z_{\geq0}^{6},
 \qquad
 \rho_W,r_\partial\in\mathbb F_3^3.
\]
Ces coordonnées sont la marge, la combinaison axiale, la condition cubique, les poids de
colonne, la classe résiduelle de Witt et le retour ternaire. Elles ne sont pas condensées en un
unique mot si une preuve ultérieure utilise l’une d’elles.

\subsubsection{Cylindre et frontière}

\[
 \mathsf B_q:=\mathsf{Cyl}_q\times\mathsf{Fr}_q.
\]
Le cylindre localise une famille de préfixes ; la frontière détermine quelles
extensions sont compatibles. La survie sera une relation entre
des états munis des deux données, non un nombre ajouté a posteriori.

\subsubsection{Registre généalogique de transport}

\(\Lambda_q\) contient les registres de transport dont l’extrémité est
\(q\). Un prolongement \(r:q\to q'\) agit par concaténation :
\[
 \lambda\longmapsto r\circ\lambda.
\]
Le registre généalogique conserve la route ; son extrémité observable ne la détermine pas.

\subsection{Transport des fibres}

Pour chaque flèche \(r:q\to q'\), on déclare des applications
\[
 \mathsf H(r):\mathsf H_q\to\mathsf H_{q'},
 \quad
 \mathsf C(r):\mathsf C_q\to\mathsf C_{q'},
 \quad
 \Lambda(r):\Lambda_q\to\Lambda_{q'}.
\]
Elles respectent les identités et la composition, et \(\mathsf C(r)\) est un homomorphisme.
La retenue satisfait
\begin{equation}
 \kappa(r_2\circ r_1)
 =\kappa(r_2)+\mathsf C(r_2)\bigl(\kappa(r_1)\bigr).
\label{p12-83-c1-s3:eq:tpk-cocadena-transporte}
\end{equation}
Par conséquent, le transport des coordonnées déterminées est
\begin{align*}
 h'&=\mathsf H(r)(h),\\
 c'&=\mathsf C(r)(c)+\kappa(r),\\
 \lambda'&=\Lambda(r)(\lambda),\\
 \sigma'&=\operatorname{Sig}_{q'}(h',c',\lambda').
\end{align*}

\begin{lemma}[Associativité de la retenue transportée]
\label{p12-83-c1-s3:lem:tpk-asociatividad-acarreo}
L’équation \cref{p12-83-c1-s3:eq:tpk-cocadena-transporte} garantit que transporter un
état par \(r_1\) puis par \(r_2\) produit la même retenue que
le transporter par \(r_2\circ r_1\).
\end{lemma}

\begin{proof}
Le transport successif donne
\[
 \mathsf C(r_2)\bigl(\mathsf C(r_1)c+\kappa(r_1)\bigr)+\kappa(r_2).
\]
Par fonctorialité de \(\mathsf C\), le premier terme est
\(\mathsf C(r_2r_1)c\) ; la somme des deux autres est
\(\kappa(r_2r_1)\) d’après \cref{p12-83-c1-s3:eq:tpk-cocadena-transporte}.
\end{proof}

\subsection{Relations d’extension}

L’orientation et la frontière ne sont pas forcées par une fonction universelle.
Pour chaque flèche \(r:q\to q'\), on déclare une relation
\[
 \operatorname{Ext}_r\subseteq\mathcal F_q\times\mathcal F_{q'}.
\]
On exige les deux conditions
\[
 \Delta_q\subseteq\operatorname{Ext}_{\operatorname{id}_q},
 \qquad
 \operatorname{Ext}_{r_2}\circ\operatorname{Ext}_{r_1}
 \subseteq\operatorname{Ext}_{r_2\circ r_1}.
\]
La première conserve les identités ; la seconde conserve la composition. Une
relation peut admettre plusieurs extensions locales et n’en laisser qu’une après un
horizon plus grand.

\begin{definition}[Topological Phase Kernel]
\label{p12-83-c1-s3:def:tpk-categoria-autonoma}
Un objet du TPK est un tuple
\[
 x=(q,o,h,c,\sigma,C,b,\lambda)
\]
tel que
\[
 q\in\mathcal Q_{\rm obs},\quad
 o\in\mathsf O_q,\quad
 h\in\mathsf H_q,\quad
 c\in\mathsf C_q,\quad
 \lambda\in\Lambda_q,
\]
\[
 \sigma=\operatorname{Sig}_q(h,c,\lambda),
 \qquad
 (C,b)\in\mathsf{Cyl}_q\times\mathsf{Fr}_q.
\]
Nous associons à chaque objet sa composante fibrale
\[
 \mathbf f(x):=(o,h,c,\sigma,(C,b),\lambda)\in\mathcal F_q.
\]
Un morphisme \(x\to x'\) est un triplet \((r,x,x')\) dans lequel
\(r:q\to q'\) appartient à \(\mathsf P_{\rm obs}\), les coordonnées
transportables satisfont les lois précédentes et
\((\mathbf f(x),\mathbf f(x'))\in
\operatorname{Ext}_r\). La composition est
\[
 (r_2,x',x'')\circ(r_1,x,x')=(r_2r_1,x,x''),
\]
et l’identité de \(x\) est \((\operatorname{id}_q,x,x)\).
\end{definition}

\begin{theorem}[Structure catégorielle du TPK]
\label{p12-83-c1-s3:thm:tpk-es-categoria}
La \cref{p12-83-c1-s3:def:tpk-categoria-autonoma} définit une catégorie et l’application
\[
 p:\mathsf{TPK}\longrightarrow\mathsf P_{\rm obs},
 \qquad p(x)=q,\qquad p(r,x,x')=r,
\]
est un foncteur.
\end{theorem}

\begin{proof}
La composition des routes est associative. Le
\cref{p12-83-c1-s3:lem:tpk-asociatividad-acarreo} prouve la compatibilité de la
coordonnée de retenue ; la fonctorialité déclarée de \(\mathsf H\),
\(\mathsf C\) et \(\Lambda\) prouve celle des autres coordonnées
déterminées. La stabilité de \(\operatorname{Ext}\) garantit que le
composé reste admissible, et sa condition diagonale fournit les
identités. L’application \(p\) conserve les identités et la composition par
construction.
\end{proof}

\subsection{Projections et pertes typées}

Les projections élémentaires sont
\[
 \phi(x)=g_0(\theta),
 \qquad
 \tau(x)=\varepsilon_9(r(\theta)),
 \qquad
 \mu(x)=\mu(\theta).
\]
La projection vers l’état TRIT local est
\[
 \pi_{\rm TRIT}(x)
 =\bigl(\tau,\mu,o,h,c,\sigma,\lambda,\phi\bigr).
\]
Elle oublie le cylindre, la frontière et une partie de \(q\) ; elle ne définit pas un autre conteneur.
Deux états peuvent avoir la même image par \(\pi_{\rm TRIT}\) et différer
par l’information oubliée.

\subsection{Critère de survie des routes par conservation de la frontière projective}

Soit \(\mathcal X_{\rm TPK}^{(n)}\) la famille des états à la profondeur
\(n\). Pour \(N\geq n\), nous définissons
\[
 \operatorname{Surv}_n^{(N)}=
 \left\{x_n:
 \begin{array}{l}
 \text{il existe }x_{n+1},\ldots,x_N\text{ avec}\\
 (x_j,x_{j+1})\in\operatorname{Ext}_{r_j}
 \text{ pour }n\leq j<N
 \end{array}\right\}.
\]

\begin{proposition}[Filtration de survie]
\label{p12-83-c1-s3:prop:tpk-filtracion-supervivencia}
Pour \(N\geq n\),
\[
 \operatorname{Surv}_n^{(N+1)}
 \subseteq
 \operatorname{Surv}_n^{(N)}.
\]
Si le schéma est défini à toute profondeur, les états
indéfiniment prolongeables sont
\[
 \operatorname{Surv}_n
 =\bigcap_{N\geq n}\operatorname{Surv}_n^{(N)}.
\]
\end{proposition}

\begin{proof}
Un prolongement jusqu’à \(N+1\) se tronque en un prolongement jusqu’à \(N\),
ce qui prouve l’inclusion. L’intersection exprime exactement la
prolongeabilité à tout horizon fini.
\end{proof}

\subsection{Arbre opératoriel du conteneur}

\begin{figure}[htbp]
 \centering
 \begin{tikzpicture}[
  node distance=7mm and 9mm,
  every node/.style={font=\small,align=center},
  root/.style={draw=black,very thick,rounded corners,inner sep=5pt},
  op/.style={draw=black,rounded corners,inner sep=4pt},
  edge/.style={-{Latex[length=2mm]},semithick}
 ]
  \node[root] (tpk) {TPK\\$\mathcal X_{\rm TPK}^{\rm enr}$};
  \node[op,below left=of tpk] (sel) {sélection\\$M_{\rm ph}$};
  \node[op,below=of tpk] (mov) {transport\\$T_d,F_{\rm obs}$};
  \node[op,below right=of tpk] (fib) {fibres\\$\mathsf H,\mathsf C,\mathsf S,
    \mathsf B,\Lambda$};
  \node[op,below=17mm of mov] (cal) {calendrier\\$C_{12}\hookrightarrow
    C_{108}\twoheadrightarrow C_9$};
  \node[op,below=of cal] (ext) {extension, cylindres et survie};
  \draw[edge] (tpk) -- (sel);
  \draw[edge] (tpk) -- (mov);
  \draw[edge] (tpk) -- (fib);
  \draw[edge] (sel) |- (cal);
  \draw[edge] (mov) -- (cal);
  \draw[edge] (fib) |- (cal);
  \draw[edge] (cal) -- (ext);
 \end{tikzpicture}
 \caption{Inclusion opératoire du Topological Phase Kernel. Le sélecteur,
 le transport, les lecteurs et les fibres ne sont pas des conteneurs parallèles : ce sont
 des applications ou des coordonnées du même état holonomique intégral.}
 \label{p12-83-c1-s3:fig:tpk-arbol-operatorio-autonomo}
\end{figure}

\subsection{Obstructions de l'état holonomique intégral du TPK : typage du sélecteur, cochaîne de retenue, extension compatible et conservation généalogique}

\begin{criterion}[Critères de réfutation de l'état holonomique intégral]
La définition est réfutée si :
\begin{enumerate}
 \item le sélecteur \(M_{\rm ph}\) est utilisé comme translation ;
 \item on identifie \(\theta\), \(g_0(\theta)\), \(r(\theta)\) et
       \(\beta(\theta)\) ;
 \item on permet de choisir une signature indépendamment de
       \((h,c,\lambda)\) ;
 \item la composition des retenues ne satisfait pas
       \cref{p12-83-c1-s3:eq:tpk-cocadena-transporte} ;
 \item on remplace \(\operatorname{Ext}_r\) par un choix fonctionnel non
       démontré ;
 \item on appelle état complet une projection qui a oublié le cylindre,
       la frontière ou le registre généalogique ;
 \item le retour de phase efface l’une des coordonnées de fibre.
\end{enumerate}
\end{criterion}

Cette unité a défini le conteneur. La suivante énumérera ses lecteurs,
opérateurs, odomètre et calendriers, et démontrera pourquoi les retours de
neuf, vingt-sept et cent huit pas ne sont pas des réinitialisations de l’état.
}%


% Unidad U001. Redacción autónoma; no procede del borrador condensado.

\section{Segments marqués, interstices et cellule nonadique}
\label{p12-83-c1-s4:chap:segmentos-marcas-intersticios}

La construction commence avant toute constante et avant tout choix
d’unité physique. La donnée inaugurale n’est pas un nombre réel déjà constitué, mais un
segment ordonné, une famille de marques et la relation d’adjacence entre
marques consécutives. Cette séparation est nécessaire parce qu’une marque, un
intervalle, la longueur de cet intervalle et le numéral qui publie cette
longueur sont des objets de types différents.

\subsection{Carte positionnelle finie et projection résiduelle décimale}

\begin{definition}[Segment marqué d’ordre dix]
Soit
\[
 \mathsf M_{10}:=\{m_0,m_1,\ldots,m_{10}\}
\]
un ensemble totalement ordonné par
\(m_0<m_1<\cdots<m_{10}\). Son graphe d’adjacence est le chemin
\[
 \mathsf P_{10}:\quad
 m_0\longrightarrow m_1\longrightarrow\cdots\longrightarrow m_{10}.
\]
Nous noterons
\[
 \mathsf I_k:=(m_k,m_{k+1}),\qquad 0\leq k\leq9,
\]
l’interstice abstrait compris entre deux marques consécutives. La
famille des interstices est
\(\mathsf E_{10}:=\{\mathsf I_0,\ldots,\mathsf I_9\}\).
\end{definition}

La notation ouverte \((m_k,m_{k+1})\) ne présuppose pas encore la droite réelle. Elle
indique seulement la paire ordonnée d’extrémités et l’arête qui les relie. Lorsque l’on choisit la
carte archimédienne
\[
 \chi_{10}(m_k)=k,
\]
l’interstice abstrait \(\mathsf I_k\) est publié par l’intervalle
semi-ouvert \([k,k+1)\). Par conséquent, l’indice \(k\), l’arête
\(\mathsf I_k\) et l’intervalle \([k,k+1)\) ne seront pas identifiés dans le texte :
le premier est une étiquette, le deuxième un élément combinatoire et le
troisième sa réalisation dans une carte.

\begin{lemma}[Recensement des marques et des interstices]
\label{p12-83-c1-s4:lem:censo-marcas-intersticios}
Le segment \(\mathsf P_{10}\) contient onze marques et dix interstices. Plus
généralement, un chemin à \(n+1\) marques possède exactement \(n\) arêtes.
\end{lemma}

\begin{proof}
L'application
\[
 \{0,\ldots,n-1\}\longrightarrow E(\mathsf P_n),
 \qquad k\longmapsto(m_k,m_{k+1}),
\]
est bijective. L’affirmation particulière s’obtient avec \(n=10\).
\end{proof}

L’extrémité \(m_0\) fixe l’origine de la carte et \(m_{10}\) fixe son seuil.
Aucune des deux ne constitue par elle-même un interstice supplémentaire. Cette
observation élémentaire empêche trois confusions ultérieures : compter les extrémités
comme des cellules, confondre la clôture avec un nouvel état de phase et traiter le zéro
comme une dixième phase indépendante du cycle nonadique.

\subsection{Extraction de la cellule de \(9\) phases}

La carte décimale distingue l’ancrage \([0,1)\) de la cellule périodique. Les
neuf marques positives
\[
 \mathsf D_9:=\{1,2,3,4,5,6,7,8,9\}
\]
seront utilisées comme représentants publics de
\(\mathbb Z/9\mathbb Z\), en faisant correspondre la marque \(9\) à la classe
résiduelle \(\overline 0\). Nous définissons
\[
 \lambda_9:\mathsf D_9\longrightarrow\mathbb Z/9\mathbb Z,
 \qquad \lambda_9(a)=\overline a.
\]

\begin{proposition}[Carte positive des classes nonadiques]
\label{p12-83-c1-s4:prop:carta-positiva-nonadica}
L’application \(\lambda_9\) est bijective. Son inverse est la section positive
\[
 s_+(\overline a)=
 \begin{cases}
 a,&1\leq a\leq8,\\
 9,&\overline a=\overline0.
 \end{cases}
\]
En particulier, l’écriture publique \(9\) et l’écriture résiduelle \(0\)
représentent la même classe, mais appartiennent à des cartes distinctes.
\end{proposition}

\begin{proof}
Les neuf classes de \(\mathbb Z/9\mathbb Z\) possèdent des représentants uniques dans
\(\{0,1,\ldots,8\}\). Remplacer le représentant \(0\) par \(9\) produit
une autre transversale complète, de sorte que \(\lambda_9\) est bijective et \(s_+\)
est son inverse.
\end{proof}

La cellule de phase ne s’obtient pas en effaçant arbitrairement l’origine. Le tronçon
\([0,1)\) fixe la carte décimale depuis laquelle on observe le cycle ; les neuf
phases sont les neuf classes publiées comme \(1,\ldots,9\). Dans la carte
relevée, le pas qui ferme la cellule est représenté par \([9,10)\) ; dans la
carte résiduelle, ce même pas retourne de \(\overline0\) à \(\overline1\).
La clôture et le changement de représentant ne sont pas la même opération.

\subsection{Position, hauteur et division euclidienne}

\begin{definition}[Coordonnées nonadiques]
Pour chaque entier \(n\geq1\), nous définissons
\[
 \rho_9(n):=1+((n-1)\bmod9),
 \qquad
 q_9(n):=\left\lfloor\frac{n-1}{9}\right\rfloor.
\]
La première coordonnée est appelée position de phase ; la seconde, hauteur de
tour ou quotient de mémoire.
\end{definition}

\begin{theorem}[Décomposition nonadique unique]
\label{p12-83-c1-s4:thm:descomposicion-nonadica-unica}
Tout entier \(n\geq1\) admet une unique écriture
\[
 n=9q+\rho,
 \qquad q\in\mathbb Z_{\geq0},\quad \rho\in\mathsf D_9.
\]
Dans cette écriture, \(q=q_9(n)\) et \(\rho=\rho_9(n)\).
\end{theorem}

\begin{proof}
Nous appliquons la division euclidienne à \(n-1\) : il existe d’uniques
\(q\geq0\) et \(r\in\{0,\ldots,8\}\) tels que
\(n-1=9q+r\). En prenant \(\rho=r+1\), on obtient
\(n=9q+\rho\), avec \(1\leq\rho\leq9\). L’unicité de \(q,r\) implique celle
de \(q,\rho\).
\end{proof}

\begin{corollary}[Retour de phase avec accroissement de mémoire]
\label{p12-83-c1-s4:cor:retorno-fase-incremento-memoria}
Pour tout \(n\geq1\),
\[
 \rho_9(n+9)=\rho_9(n),
 \qquad
 q_9(n+9)=q_9(n)+1.
\]
Par conséquent, un tour complet conserve la phase visible et change
l’état relevé.
\end{corollary}

\begin{proof}
Si \(n=9q+\rho\), alors \(n+9=9(q+1)+\rho\) ; l’unicité du
\cref{p12-83-c1-s4:thm:descomposicion-nonadica-unica} donne les deux identités.
\end{proof}

L’évolution d’un pas s’écrit sans ambiguïté comme
\[
 (q,\rho)\longmapsto
 \begin{cases}
 (q,\rho+1),&1\leq\rho<9,\\
 (q+1,1),&\rho=9.
 \end{cases}
\]
La deuxième ligne n’est pas une réinitialisation. Elle retrouve l’étiquette initiale de phase et
augmente le quotient qui enregistre combien de tours ont été effectués.

\subsection{Périodicité de la projection circulaire et injectivité du relèvement hélicoïdal nonadique}

Soit
\[
 \vartheta_9(\rho)=\frac{2\pi(\rho-1)}9.
\]
La projection de phase sur le cercle unité est
\[
 \mathcal R_9(n)=
 \bigl(\cos\vartheta_9(\rho_9(n)),
       \sin\vartheta_9(\rho_9(n))\bigr).
\]
Pour conserver simultanément la position et la hauteur, nous introduisons
l’immersion hélicoïdale
\[
 \mathcal H_9(n)=
 \left(
  \cos\vartheta_9(\rho_9(n)),
  \sin\vartheta_9(\rho_9(n)),
  q_9(n)+\frac{\rho_9(n)-1}{9}
 \right).
\]

\begin{proposition}[La projection circulaire ne détermine pas l’état]
\label{p12-83-c1-s4:prop:rueda-no-determina-estado}
L’application \(\mathcal R_9\) est périodique de période neuf. L’application
\(\mathcal H_9\) est injective. En particulier,
\[
 \mathcal R_9(n+9)=\mathcal R_9(n),
 \qquad
 \mathcal H_9(n+9)\neq\mathcal H_9(n).
\]
\end{proposition}

\begin{proof}
La première égalité découle du
\cref{p12-83-c1-s4:cor:retorno-fase-incremento-memoria}. La troisième coordonnée de
\(\mathcal H_9(n+9)\) dépasse d’une unité celle de \(\mathcal H_9(n)\),
de sorte que les deux points sont distincts. Enfin, à partir de la troisième coordonnée,
on retrouve \(n\) en multipliant par neuf et en ajoutant un à la partie de phase ;
par conséquent, \(\mathcal H_9\) est injective.
\end{proof}

Cette proposition est le modèle minimal du principe qui traversera tout
l’ouvrage :
\[
 \boxed{\text{retour de phase}\neq\text{retour de l’état}.}
\]
La connexion nonadique du Topological Phase Kernel conservera davantage de champs que
l’hélice élémentaire ---feuillet, orientation, retenue, frontière, terminal et
mémoire de route---, mais elle ne modifiera pas cette distinction inaugurale.

\subsection{Correspondance entre cardinalité discrète et mesure projective}

\begin{definition}[Dénombrement]
Le dénombrement d’un ensemble fini \(X\) est son cardinal \(|X|\). Il n’exige
ni orientation, ni échelle, ni lecteur supplémentaire.
\end{definition}

\begin{definition}[Donnée de mesure]
Une donnée de mesure discrète est un triplet
\[
 (x,\mathcal L,\mathcal C),
\]
où \(x\) est un état, \(\mathcal L\) est un lecteur typé et
\(\mathcal C\) est une carte de publication compatible avec le codomaine de
\(\mathcal L\). Sa sortie visible est
\[
 \mathcal C(\mathcal L(x)).
\]
\end{definition}

Le même ensemble peut avoir un unique cardinal et admettre plusieurs mesures,
parce que la cochaîne locale, l’orientation d’intégration ou la
carte de publication peuvent varier. Inversement, deux états distincts peuvent publier la
même coordonnée. C’est pourquoi l’ouvrage conservera toujours deux niveaux :
\[
 \text{état avec généalogie}
 \quad\longrightarrow\quad
 \text{coordonnée publiée}.
\]

\begin{criterion}[Critères de réfutation de l’unité inaugurale]
La construction est réfutée si l’un quelconque des faits suivants
se produit :
\begin{enumerate}
 \item on compte dix phases indépendantes dans la cellule résiduelle ;
 \item on identifie la marque \(9\) à l’entier \(0\) sans déclarer le
 changement de carte ;
 \item le pas \(9\to1\) force \(q\mapsto0\) ;
 \item deux entiers distincts possèdent la même paire \((q_9,\rho_9)\) ;
 \item une unité physique est attribuée à la cellule avant la définition du
 transducteur dimensionnel.
\end{enumerate}
\end{criterion}

\subsection{Domaine inaugural de la cellule nonadique et exclusion causale de l’APP, du TRIT, du TPK et des constantes analytiques}

L’unité construite jusqu’ici contient uniquement :
\[
 \mathsf P_{10},\quad
 \mathsf D_9,\quad
 \lambda_9,\quad
 (q_9,\rho_9),\quad
 \mathcal R_9,\quad
 \mathcal H_9.
\]
Elle ne contient pas encore l’APP, le TRIT, le Topological Phase Kernel ni une constante
analytique. L’unité suivante prendra deux copies de la cellule nonadique,
construira leur complétion gnomonique et démontrera comment apparaît le support
bidimensionnel sur lequel agissent la matrice multiplicative de l’APP et sa loi
additive interne.


% Unidad U003B. Esqueleto ciclotómico de APP y completación binomial.
% Redacción autónoma a partir de los objetos, no de la arquitectura del PDF abreviado.

\section{Squelette cyclotomique de l’APP et complétion binomiale de l’unité}
\label{p12-83-c1-s5:chap:app-esqueleto-completacion}

Le défaut entre la matrice APP et son accumulateur interne n’est pas un nombre
isolé. Sa loi en
longueur cartésienne, sa décomposition en orbites d’ordre trois et la
complétion de la cellule nonadique produisent trois objets différents qui doivent
rester typés :
\[
 \text{défaut scalaire}\quad D_d,
 \qquad
 \text{module cyclotomique}\quad W_{\rm APP},
 \qquad
 \text{complétion}\quad\overline{\mathcal C}_9.
\]
Le premier compte ; le deuxième conserve l’orientation et le rang ; le troisième ajoute
des strates de frontière. Les égalités entre leurs cardinalités n’autorisent pas à
identifier leurs éléments.

\subsection{Loi du défaut à module fixe et longueur variable}

Pour chaque entier \(d\geq1\), soit
\[
 B_d:=\{1,\ldots,9\}^d.
\]
La racine numérique nonadique est représentée par
\[
 \operatorname{dr}_9(n)=
 \begin{cases}
 n\bmod9,&n\not\equiv0\pmod9,\\
 9,&n\equiv0\pmod9.
 \end{cases}
\]
Nous définissons l’accumulation additive et l’évaluation matricielle multiplicative
sur le même domaine :
\[
 A_d(x_1,\ldots,x_d)
 =\operatorname{dr}_9\!\left(\sum_{j=1}^d x_j\right),
\]
\[
 M_d(x_1,\ldots,x_d)
 =\operatorname{dr}_9\!\left(\prod_{j=1}^d x_j\right).
\]
Leurs masses visibles et leur défaut sont
\[
 S_\Sigma(d):=\sum_{x\in B_d}A_d(x),
 \qquad
 S_\Pi(d):=\sum_{x\in B_d}M_d(x),
\]
\[
 D_d:=S_\Pi(d)-S_\Sigma(d).
\]
Ici, \(d\) compte les positions cartésiennes. Le module \(9\) ne varie pas et
aucune dimension physique n’a encore été introduite.

\begin{lemma}[Masse de l’accumulateur additif]
\label{p12-83-c1-s5:lem:masa-hoja-aditiva-longitud}
Pour tout \(d\geq1\),
\[
 S_\Sigma(d)=45\,9^{d-1}.
\]
\end{lemma}

\begin{proof}
Les \(d-1\) premières coordonnées et un résidu publié
\(r\in\{1,\ldots,9\}\) étant fixés, il existe une unique dernière coordonnée qui donne à la
somme le résidu \(r\). Chaque valeur visible possède ainsi \(9^{d-1}\)
antécédents. En sommant les poids \(1+\cdots+9=45\), on obtient la formule.
\end{proof}

\begin{lemma}[Recensement de la matrice multiplicative]
\label{p12-83-c1-s5:lem:censo-hoja-multiplicativa-longitud}
Pour chacun des six résidus unitaires, il y a \(6^{d-1}\) antécédents
par valeur fixée du produit. Les résidus \(3\) et \(6\) ont
\(d6^{d-1}\) antécédents chacun. Le résidu nul, publié comme \(9\),
a
\[
 9^d-6^d-2d6^{d-1}
\]
antécédents. En conséquence,
\[
 S_\Pi(d)=9^{d+1}-9(d+3)6^{d-1}.
\]
\end{lemma}

\begin{proof}
Le groupe des unités de \(\mathbb Z/9\mathbb Z\) a six éléments. En
fixant \(d-1\) unités, la dernière est déterminée par le produit
unitaire prescrit, ce qui donne \(6^{d-1}\) solutions pour chacune des
six classes.

Pour obtenir un produit de valuation ternaire exactement égale à un, on choisit laquelle
des \(d\) coordonnées appartient à l’une des deux classes \(3,6\), tandis que
les autres sont des unités. Le produit \(3\) ou \(6\) étant fixé, la dernière unité
est déterminée ; d’où \(d6^{d-1}\) pour chaque classe. Tous les tuples
restants ont un produit nul modulo \(9\). Soustraire de \(9^d\) les deux
recensements précédents donne leur multiplicité. La somme pondérée par les
représentants visibles produit la dernière identité.
\end{proof}

\begin{theorem}[Loi dimensionnelle interne de l’APP]
\label{p12-83-c1-s5:thm:ley-dimensional-app-autonoma}
Pour tout \(d\geq1\),
\[
 \boxed{D_d=4\,9^d-9(d+3)6^{d-1}.}
\]
En particulier,
\[
 D_1=0,
 \qquad
 D_2=54,
\]
et
\[
 \sum_{d\geq1}D_dz^d
 =\frac{54z^2(1-3z)}{(1-9z)(1-6z)^2}.
\]
De plus,
\[
 \frac{D_d}{9^d}
 =4-(d+3)\left(\frac23\right)^{d-1}
 \longrightarrow4.
\]
\end{theorem}

\begin{proof}
On soustrait les formules des deux lemmes. Pour la fonction génératrice, on utilise
\[
 \sum_{d\geq1}9^dz^d=\frac{9z}{1-9z},
 \qquad
 \sum_{d\geq1}(d+3)6^{d-1}z^d
 =\frac{z(4-18z)}{(1-6z)^2},
\]
et l’on réduit le numérateur commun. La formule normalisée résulte de la division
par \(9^d\) ; le terme exponentiel multiplié par \(d+3\) converge vers
zéro.
\end{proof}

L’égalité \(D_2=54\) reproduit le recensement \(459-405\), mais la formule
précédente et la loi à module variable
\(\Delta_{3^m}^{\rm pl}=m3^{2m-1}\) sont des énoncés distincts : la première
fait varier \(d\) en maintenant le module \(9\) ; la seconde fait varier le module.

\subsection{Opérateur cyclotomique primaire}

Soit
\[
 G_{A_2}=\begin{pmatrix}2&-1\\-1&2\end{pmatrix},
 \qquad
 C=\begin{pmatrix}0&-1\\1&-1\end{pmatrix}.
\]

\begin{proposition}[Coxeter de type \(A_2\)]
\label{p12-83-c1-s5:prop:coxeter-a2-autonomo}
On vérifie
\[
 C^3=I_2,
 \qquad
 C^{\mathsf T}G_{A_2}C=G_{A_2},
 \qquad
 \det(XI_2-C)=X^2+X+1.
\]
Par conséquent, \(C\) préserve la forme \(A_2\), est d’ordre exactement trois et
ne possède aucun vecteur fixe non nul.
\end{proposition}

\begin{proof}
La multiplication donne
\[
 C^2=\begin{pmatrix}-1&1\\-1&0\end{pmatrix},
 \qquad C^3=I_2.
\]
La substitution directe dans \(C^{\mathsf T}G_{A_2}C\) renvoie
\(G_{A_2}\). Enfin, \(1\) n’est pas racine de \(X^2+X+1\), donc
\(\ker(C-I_2)=0\).
\end{proof}

\subsection{Décomposition de l’APP en \(12\) fibres de type \(A_2\)}

Considérons l’action
\[
 a:\mathbb Z/9\mathbb Z\longrightarrow\mathbb Z/9\mathbb Z,
 \qquad a(r)=4r.
\]
Comme \(4^3\equiv1\pmod9\), sa restriction aux unités a deux orbites
libres :
\[
 \mathcal O_1=(1,4,7),
 \qquad
 \mathcal O_2=(2,8,5).
\]
Pour une orbite \(\mathcal O=(r,4r,7r)\), nous définissons le module d’augmentation
\[
 A(\mathcal O):=
 \ker\!\left(\mathbb Z[\mathcal O]
 \xrightarrow{\,\sum\,}\mathbb Z\right).
\]
Dans la base
\[
 b_1=e_r-e_{4r},
 \qquad b_2=e_{4r}-e_{7r},
\]
la forme euclidienne restreinte a pour matrice de Gram \(G_{A_2}\), et le cycle
\(r\mapsto4r\mapsto7r\mapsto r\) agit par \(C\).

Pour \(x=(x_1,x_2,x_3)\in(\mathbb Z/9\mathbb Z)^3\), les trois paires non
ordonnées sont
\[
 \{1,2\},\quad\{2,3\},\quad\{3,1\}.
\]
Sur chaque paire, on évalue deux lois typées :
\[
 \ell_{ij}^{\Sigma}(x)=x_i+x_j,
 \qquad
 \ell_{ij}^{\Pi}(x)=x_ix_j.
\]
Comme
\[
 \ell_{ij}^{\Sigma}(ax)=a\ell_{ij}^{\Sigma}(x),
 \qquad
 \ell_{ij}^{\Pi}(ax)=a^{-1}\ell_{ij}^{\Pi}(x),
\]
la somme porte l’orientation \(C\) et le produit l’orientation inverse
\(C^{-1}\).

\begin{definition}[Module cyclotomique de l’APP]
\label{p12-83-c1-s5:def:modulo-ciclotomico-app-autonomo}
On définit
\[
 \boxed{
 W_{\rm APP}:=
 \bigoplus_{\{i,j\}}
 \bigoplus_{\mu\in\{\Sigma,\Pi\}}
 \bigoplus_{\mathcal O\in\{\mathcal O_1,\mathcal O_2\}}
 A(\mathcal O).}
\]
Il y a \(3\cdot2\cdot2=12\) termes, chacun de rang deux ; par conséquent
\[
 W_{\rm APP}\cong A_2^{12},
 \qquad \operatorname{rank}W_{\rm APP}=24.
\]
Les étiquettes paire--loi--orbite font partie de l’objet.
\end{definition}

Pour \(u\in\mathbb Z/9\mathbb Z\), nous écrivons
\[
 \beta_{\mathcal O}(u)=
 \begin{cases}
 (1,0),&u=r,\\
 (0,1),&u=4r,\\
 (-1,-1),&u=7r,\\
 (0,0),&u\notin\mathcal O,
 \end{cases}
\]
dans la base \((b_1,b_2)\). L’application d’état est
\[
 \Psi(x)_{ij,\mu,\mathcal O}
 :=\beta_{\mathcal O}(\ell_{ij}^{\mu}(x)).
\]

\begin{theorem}[Équivariance et rang du squelette d’état]
\label{p12-83-c1-s5:thm:esqueleto-estatal-app-autonomo}
Soit \(\rho(a)\) l’action qui utilise \(C\) dans les six fibres additives et
\(C^{-1}\) dans les six fibres multiplicatives. Alors
\[
 \boxed{\Psi(ax)=\rho(a)\Psi(x)}
\]
pour les \(729\) états de \((\mathbb Z/9\mathbb Z)^3\). Le sous-module
rationnel engendré par les images a pour rang \(24\).
\end{theorem}

\begin{proof}
Dans la loi additive, multiplier toutes les coordonnées par \(4\) avance d’une
position dans chaque orbite ; dans la loi multiplicative, cela multiplie la valeur par
\(4^2\equiv7\equiv4^{-1}\pmod9\) et inverse donc le cycle. Cela prouve
l’équivariance composante par composante.

Pour que le rang ne reste pas caché dans une affirmation informatique, nous fixons
l’ordre des coordonnées
\[
 (12,\Sigma,\mathcal O_1),(12,\Sigma,\mathcal O_2),
 (12,\Pi,\mathcal O_1),(12,\Pi,\mathcal O_2),
\]
suivi du même ordre pour les paires \(23\) et \(31\), et, à l’intérieur de chaque
fibre, les deux coordonnées \((b_1,b_2)\). Considérons les états
\[
\begin{gathered}
001,002,004,005,010,011,012,013,014,015,016,020,\\
022,023,040,050,101,102,104,105,110,120,140,150.
\end{gathered}
\]
La matrice entière \(M_\Psi\) dont les lignes sont leurs vingt-quatre images a ses
entrées dans ({-1,0,1}). L’élimination sans fractions de
Bareiss produit comme mineurs principaux successifs
\[
1,1,1,-1,-1,-1,-1,1,-1,1,1,1,1,-2,-1,1,1,1,2,1,4,-6,12,9.
\]
En particulier,
\[
 \det M_\Psi=9\neq0.
\]
Ces vingt-quatre images sont linéairement indépendantes. Comme le
codomaine a déjà pour rang \(24\), elles engendrent \(W_{\rm APP}\otimes\mathbb Q\).
\end{proof}

La preuve présente un certificat fini reproductible ; elle ne déduit pas le rang plein
de la seule surjectivité de chaque projection composante.

\subsection{Localisation de l’agrégat \texorpdfstring{\(54\)}{54}}

Soit
\[
 \mathcal M:=\mathbb Z[\mathbb Z/9\mathbb Z]
\]
et, pour \(\mu\in\{\Sigma,\Pi\}\),
\[
 \nu_\mu:=\sum_{x,y\in\mathbb Z/9\mathbb Z}
 e_{\ell^\mu(x,y)}.
\]
Le recensement des deux tables donne
\[
 \delta_{\Pi\Sigma}:=\nu_\Pi-\nu_\Sigma
 =12e_0+3e_3+3e_6
 -3\sum_{u\in(\mathbb Z/9\mathbb Z)^\times}e_u.
\]
Ses coefficients sont constants sur les orbites de \(a\) ; par conséquent
\[
 (I-a)\delta_{\Pi\Sigma}=0.
\]
Nous définissons la projection d’augmentation
\[
 T:\mathcal M\longrightarrow
 A(\mathcal O_1)\oplus A(\mathcal O_2),
\]
\[
 T(\nu)=
 \left(((I-a)\nu)|_{\mathcal O_1},
       ((I-a)\nu)|_{\mathcal O_2}\right).
\]
Alors
\[
 T(\delta_{\Pi\Sigma})=0.
\]
Cependant, le lecteur visible
\[
 w(e_0)=9,
 \qquad w(e_r)=r\quad(1\leq r\leq8)
\]
satisfait
\[
 w(\delta_{\Pi\Sigma})=54.
\]

\begin{corollary}[Obstruction isotypique]
\label{p12-83-c1-s5:cor:obstruccion-isotipica-app-autonoma}
Le scalaire \(54\) appartient au secteur \(C_3\)-invariant et ne peut
être retrouvé par une fonctionnelle linéaire qui se factorise exclusivement par
les modules d’augmentation. En particulier,
\[
 \sum_x\Psi(x)=0
\]
et, pour toute fonctionnelle linéaire \(L\) sur \(W_{\rm APP}\otimes\mathbb Q\),
\[
 L\!\left(\sum_x\Psi(x)\right)=0\neq54.
\]
\end{corollary}

Cette obstruction empêche de présenter \(54\) comme une somme linéaire rétrospective
des douze fibres. Le couplage APP--Fock ultérieur conserve
l’orientation et utilise un objet supplémentaire ; il ne contredit pas le corollaire.

\subsection{Complétion cubique de la cellule nonadique}

Soit \(E_9\) l’ensemble sous-jacent à \(\mathbb Z/9\mathbb Z\), et soit
\(\star\notin E_9\) un état de frontière. Nous définissons
\[
 \mathcal C_9:=E_9^3,
 \qquad
 \overline{\mathcal C}_9:=(E_9\sqcup\{\star\})^3.
\]
Le symbole \(\star\) n’est pas la classe zéro : il indique qu’une coordonnée a
atteint la strate de clôture.

\begin{theorem}[Stratification binomiale de l’unité cubique]
\label{p12-83-c1-s5:thm:emrh-binomial-autonoma}
La strate ayant exactement \(r\) coordonnées égales à \(\star\) contient
\[
 \binom3r9^{3-r}
\]
éléments. Par conséquent,
\[
 \boxed{1000=729+243+27+1=729+271.}
\]
En retirant seulement le sommet \((\star,\star,\star)\), il reste
\[
 999=729+243+27=729+270.
\]
\end{theorem}

\begin{proof}
On choisit les \(r\) coordonnées de frontière de \(\binom3r\) manières ; chaque
coordonnée restante admet neuf valeurs. Les quatre strates sont disjointes
et épuisent le produit cartésien. La somme est le développement de
\((9+1)^3\).
\end{proof}

\begin{figure}[htbp]
\centering
\resizebox{0.96\linewidth}{!}{\begin{tikzpicture}[font=\small]
  \begin{scope}[x={(.95cm,0cm)},y={(.45cm,.33cm)},z={(0cm,.95cm)}]
  \coordinate (O) at (0,0,0);
  \coordinate (X) at (3.4,0,0);
  \coordinate (Y) at (0,3.4,0);
  \coordinate (XY) at (3.4,3.4,0);
  \coordinate (Z) at (0,0,3.4);
  \coordinate (XZ) at (3.4,0,3.4);
  \coordinate (YZ) at (0,3.4,3.4);
  \coordinate (XYZ) at (3.4,3.4,3.4);
  \draw[hmtgray,dashed,semithick] (O)--(Y)--(XY);
  \draw[hmtgray,dashed,semithick] (Y)--(YZ);
  \path[fill=hmtlightblue,fill opacity=.78]
    (O)--(X)--(XZ)--(Z)--cycle;
  \path[fill=hmtlightgold,fill opacity=.62]
    (X)--(XY)--(XYZ)--(XZ)--cycle;
  \path[fill=hmtlightgreen,fill opacity=.72]
    (Z)--(XZ)--(XYZ)--(YZ)--cycle;
  \draw[hmtblue,thick] (O)--(X)--(XZ)--(Z)--cycle;
  \draw[hmtgold,thick] (X)--(XY)--(XYZ)--(XZ);
  \draw[hmtgreen,thick] (Z)--(YZ)--(XYZ);
  \node[font=\bfseries\Large,hmtgreen,fill=white,fill opacity=.78,
    text opacity=1,inner sep=2pt] at (1.7,.55,1.72) {\(9^3=729\)};
  \node[font=\scriptsize,hmtgray,fill=white,fill opacity=.78,
    text opacity=1,inner sep=1.5pt] at (1.7,.55,1.22)
    {intérieur nonadique};
  \draw[-{Latex[length=2mm]},hmtblue,semithick]
    (O) -- (1.18,0,0)
    node[pos=.78,below=5pt,font=\scriptsize] {\(x=(x_0,x_1)\)};
  \draw[-{Latex[length=2mm]},hmtgreen,semithick]
    (O) -- (0,1.18,0)
    node[pos=1,right=5pt,font=\scriptsize] {\(y=(y_0,y_1)\)};
  \draw[-{Latex[length=2mm]},hmtblue,semithick]
    (O) -- (0,0,1.18)
    node[pos=.78,left=2pt,font=\scriptsize] {\(z=(z_0,z_1)\)};
  \end{scope}
  \node[align=center,text width=5.2cm,font=\scriptsize,hmtgray,
    anchor=north] at (2.55,-.72)
    {coordonnées ternaires appariées dans trois directions indépendantes};
  \draw[decorate,decoration={brace,amplitude=5pt},hmtgold,thick]
    (5.20,.65)--(5.20,4.80);
  \node[anchor=west,align=left,text width=1.25cm,font=\scriptsize,hmtgold]
    at (5.40,3.82) {du côté \(9\) à \(10\)\\couronne \(271\)};
  \begin{scope}[xshift=7cm,yshift=.45cm]
    \node[anchor=west,font=\bfseries\large,hmtblue] at (0,4.05)
      {Complétion cubique};
    \draw[fill=hmtlightgreen,draw=hmtgreen] (0,3.15) rectangle (6.1,3.75);
    \node at (3.05,3.45) {intérieur : \(729\)};
    \draw[fill=hmtlightblue,draw=hmtblue] (0,2.35) rectangle (2.9,2.95);
    \node[align=center,text width=2.75cm,font=\scriptsize] at (1.45,2.65)
      {caras nuevas\\\(243=3\cdot9^2\)};
    \draw[fill=hmtlightgold,draw=hmtgold] (2.9,2.35) rectangle (4.9,2.95);
    \node[align=center,text width=1.8cm,font=\scriptsize] at (3.9,2.65)
      {arêtes\\\(27=3\cdot9\)};
    \draw[fill=hmtlightred,draw=hmtred] (4.9,2.35) rectangle (6.1,2.95);
    \node[align=center,font=\scriptsize] at (5.5,2.65)
      {sommet\\\(1\)};
    \node[font=\bfseries\Large,hmtgold] at (3.05,1.65)
      {\(1000=729+243+27+1=729+271\)};
    \node[align=center,hmtgray] at (3.05,.65)
      {frontière interne du cube de côté \(9\) : \(386\) points\\
       couronne de la complétion \(9^3\subset10^3\) : \(271\) états};
  \end{scope}
\end{tikzpicture}
 }
\caption{Complétion cubique de la cellule nonadique. L’intérieur de
cardinal \(729\) et les strates de frontière de cardinalités \(243\), \(27\) et
\(1\) sont des parties disjointes d’un unique objet.}
\label{p12-83-c1-s5:fig:cubo-emrh-autonomo}
\end{figure}

Cette stratification reçoit dans HMT le nom d’extrême et moyenne raison
holographique. Mathématiquement, c’est la complétion binomiale typée de la cellule
nonadique ; le nom ne remplace pas sa définition.

Pour tout \(d\geq1\), soit
\[
 \overline{\mathcal C}^{(d)}_9=(E_9\sqcup\{\star\})^d.
\]
Attribuons le poids \(1/10\) à chaque symbole d’une coordonnée et prenons la mesure
produit \(\mu_d\).

\begin{theorem}[Compatibilité projective de la mesure]
\label{p12-83-c1-s5:thm:medida-proyectiva-completacion}
La projection qui oublie la dernière coordonnée satisfait
\[
 (\pi_d)_*\mu_d=\mu_{d-1}.
\]
De plus,
\[
 10^d=\sum_{r=0}^{d}\binom dr9^{d-r},
 \qquad
 \mu_d(\overline{\mathcal C}^{(d)}_9)=1.
\]
Le nombre d’états augmente avec \(d\), tandis que la masse totale normalisée
est conservée.
\end{theorem}

\begin{proof}
Pour tout cylindre \(A\times(E_9\sqcup\{\star\})\),
\[
 \mu_d(A\times(E_9\sqcup\{\star\}))
 =\mu_{d-1}(A)\sum_{u\in E_9\sqcup\{\star\}}\frac1{10}
 =\mu_{d-1}(A).
\]
Les cylindres engendrent l’algèbre finie des sous-ensembles, ce qui détermine la
mesure image. L’identité de cardinalités est à nouveau le binôme de Newton.
\end{proof}

\subsection{Morphisme du bloc central de l’APP vers sa réalisation euclidienne}

Le bloc central déjà construit est
\[
 B_c=\begin{pmatrix}7&2\\2&7\end{pmatrix}.
\]
La lecture par lignes et la lecture diagonale donnent
\[
 72+27=77+22=99,
\]
\[
 72-27=45=\det B_c,
 \qquad
 77-22=55=54+1.
\]
La relation avec la complétion ne s’obtient pas en concaténant des symboles, mais
par la division euclidienne
\[
 729=10\cdot72+9,
 \qquad
 271=10\cdot27+1.
\]

\begin{theorem}[Unicité du relèvement décimal]
\label{p12-83-c1-s5:thm:levantamiento-decimal-emrh-autonomo}
Pour \(L_d(n)=10n+d\), avec \(d\in\{0,\ldots,9\}\), il existe une unique paire
\((d,d')\) telle que
\[
 L_d(72)=729,
 \qquad
 L_d(72)+L_{d'}(27)=1000,
 \qquad
 d+d'=10.
\]
C’est la paire
\[
 (d,d')=(9,1).
\]
\end{theorem}

\begin{proof}
La première égalité est \(720+d=729\), donc \(d=9\). La deuxième devient
\(729+270+d'=1000\), d’où \(d'=1\). La troisième est vérifiée et
l’unicité du quotient et du reste exclut une autre solution.
\end{proof}

On obtient ainsi le diagramme causal
\[
 B_c\longrightarrow(72,27)
 \longleftarrow((72,9),(27,1))
 \longleftarrow(729,271).
\]
Le bloc et sa lecture appartiennent à l’APP ; les dividendes et les restes appartiennent
à la complétion. Le diagramme partage des quotients, il n’inverse pas la
causalité.

\subsection{Incidence du cube et séparation d’avec les hexades et les octades}

Soit \(M_d\) la matrice d’incidence des \(2d\) faces de codimension un
d’un hypercube sur ses \(2^d\) sommets. Chaque face contient
\(2^{d-1}\) sommets ; deux faces opposées sont disjointes et deux faces
transversales partagent \(2^{d-2}\) sommets.

\begin{theorem}[Spectre d’incidence hypercubique]
\label{p12-83-c1-s5:thm:espectro-incidencia-cubo-autonomo}
\[
 \operatorname{Spec}(M_dM_d^{\mathsf T})
 =\left\{
 d2^{d-1}\,^{(1)},
 2^{d-1}\,^{(d)},
 0^{(d-1)}
 \right\}.
\]
Para \(d=3\),
\[
 \operatorname{Spec}(M_3M_3^{\mathsf T})
 =\{12,4,4,4,0,0\}.
\]
\end{theorem}

\begin{proof}
L’espace des lignes se décompose en : la droite constante ; les \(d\)
différences de chaque paire de faces opposées ; et les (d-1) combinaisons de
sommes de paires dont la somme totale est nulle. Les intersections décrites précédemment
font que la matrice de Gram agit par \(d2^{d-1}\), \(2^{d-1}\) et \(0\),
respectivement.
\end{proof}

Le cube tridimensionnel possède six faces, huit sommets et douze arêtes. Ces
cardinalités anticipent un motif d’incidence (6/8), mais une face n’est pas
une hexade et l’ensemble des sommets n’est pas une octade. L’identification
exceptionnelle exige un transporteur ultérieur qui conserve les supports et la
jauge.

\subsection{Résonance de l’opérateur d’incidence en rang \(6\)}

Définissons les strates qui conservent exactement \(k\) coordonnées
nonadiques en rang \(d\) :
\[
 \mathcal I_k(d):=\binom dk9^k.
\]
Les défauts de l’APP sont
\[
 D_{\rm vis}=54,
 \qquad
 D_{\rm tot}=1215,
 \qquad
 D_{\rm coc}=129,
 \qquad
 1215=54+9\cdot129.
\]

\begin{theorem}[Unicité de la résonance binomiale]
\label{p12-83-c1-s5:thm:resonancia-binomial-autonoma}
L’unique entier positif \(d\) qui satisfait simultanément
\[
 \mathcal I_1(d)=54,
 \qquad
 \mathcal I_2(d)=1215
\]
est \(d=6\). Dans ce rang,
\[
 \frac{\mathcal I_2(6)-\mathcal I_1(6)}9=129.
\]
\end{theorem}

\begin{proof}
La première égalité est \(9d=54\), qui force \(d=6\). Alors
\[
 \mathcal I_2(6)=81\binom62=81\cdot15=1215,
\]
et
\[
 \frac{1215-54}{9}=129.
\]
Il ne reste aucun autre rang possible parce que la première égalité détermine déjà \(d\).
\end{proof}

\subsection{Obstructions du squelette cyclotomique de l’APP : fibres \(A_2^{12}\), complétion \(729+271\), résonance de rang \(6\) et incidence cubique}

\begin{criterion}[Réfutation typée de l’unité]
La construction est réfutée si l’un quelconque des faits
suivants se produit :
\begin{enumerate}
 \item la loi de \(D_d\) ne reproduit pas \(D_1=0\) et \(D_2=54\) ;
 \item la variation de longueur est fusionnée avec la variation de module ;
 \item \(C\) cesse de préserver \(G_{A_2}\) ou acquiert un vecteur fixe ;
 \item les étiquettes paire--feuillet--orbite ne produisent pas douze fibres ;
 \item le mineur certifié de \(\Psi\) a un déterminant nul ;
 \item une fonctionnelle qui se factorise par le module d’augmentation retrouve \(54\) ;
 \item le symbole de frontière \(\star\) est identifié au résidu zéro ;
 \item la somme des quatre strates n’est pas égale à \(1000\) ;
 \item une face cubique est identifiée à une hexade sans un transporteur
d’incidence.
\end{enumerate}
\end{criterion}

L’unité a construit, sans les reconnaître rétrospectivement, le module
\(A_2^{12}\), la complétion \(729+271\), la mesure projective et la
résonance de rang six. Ces structures alimenteront de manière différente
la génération des constantes et les branches exceptionnelles.


% Unidad U006F. Inventario operatorio completo del Topological Phase Kernel.

\long\def\HMTDeferredTPKDevelopment{%
\section{Algèbre d'opérateurs du Topological Phase Kernel et dynamique causale d'actualisation de l'état holonomique intégral}
\label{p12-83-c1-s6:chap:inventario-operatorio-completo-tpk}

Le Topological Phase Kernel n'est ni une réduction modulo trois ni une horloge de
cent huit positions. C'est une catégorie d'états holonomiques intégraux munie
d'opérateurs de sélection, de transport, de lecture, de mémoire, de périodicité,
de prolongation et de publication. Pour empêcher qu'une abréviation ne remplace
l'objet, chaque opérateur est enregistré par son type et par l'information qu'il
conserve.

\subsection{Schéma typé des 34 constructions canoniques du TPK : domaine, codomaine, action, invariants et mémoire}

\begin{definition}[Fiche opératoire]
\label{p12-83-c1-s6:def:u006f-ficha-operatoria}
Une réalisation canonique du TPK est un sextuplet
\begin{equation}
 \mathcal O=
 (D_{\mathcal O},C_{\mathcal O},f_{\mathcal O},
 I_{\mathcal O},L_{\mathcal O},M_{\mathcal O}),
 \label{p12-83-c1-s6:eq:u006f-ficha}
\end{equation}
où :
\begin{enumerate}
 \item \(f_{\mathcal O}:D_{\mathcal O}\to C_{\mathcal O}\) est sa règle ;
 \item \(I_{\mathcal O}\) est la famille des relations conservées ;
 \item \(L_{\mathcal O}\) déclare l’information publiée ou perdue ;
 \item \(M_{\mathcal O}\) énumère la mémoire qui demeure dans l'état
 holonomique intégral.
\end{enumerate}
\end{definition}

Deux réalisations ne représentent le même opérateur que s’il existe une
équivalence typée qui entrelace les règles et les invariants. Partager une période,
une cardinalité ou une sortie visible ne suffit pas.

\subsection{Classification de l'état holonomique intégral en \(8\) classes structurelles invariantes}

L’inventaire est partitionné en :
\begin{enumerate}
 \item[\textbf{C1.}] supports et états ;
 \item[\textbf{C2.}] sélection interne ;
 \item[\textbf{C3.}] transporte;
 \item[\textbf{C4.}] lecteurs et quotients ;
 \item[\textbf{C5.}] mémoire et frontière ;
 \item[\textbf{C6.}] périodicité et retour ;
 \item[\textbf{C7.}] prolongement ;
 \item[\textbf{C8.}] salidas.
\end{enumerate}
La classe est déterminée par le domaine, le codomaine et la fonction, non par le nom
d’un chiffre associé.

\Needspace{.92\textheight}
\subsection{Décomposition fonctionnelle en \(14\) regroupements compatibles avec les opérateurs APP–TRIT–TPK}

\begingroup
\small
\renewcommand{\arraystretch}{1.12}
\begin{longtable}{@{}>{\raggedright\arraybackslash}p{.20\textwidth}>{\raggedright\arraybackslash}p{.29\textwidth}>{\raggedright\arraybackslash}p{.43\textwidth}@{}}
\toprule
Regroupement&Opérateurs ou espaces&Fonction et invariants\\
\midrule
\endfirsthead
\toprule
Regroupement&Opérateurs ou espaces&Fonction et invariants\\
\midrule
\endhead
État et mémoire de route&
\(\mathcal X_{\rm TPK}^{\rm enr}\), \(\mathcal F_q\)&
Ils conservent la position, le feuillet, l’orientation, le résidu, le quotient, la retenue, la signature,
la frontière, le cylindre, le préfixe, la provenance et la mémoire.\\

Sélection interne&
\(M_{\rm ph}:\{1,\ldots,9\}\to\{-1,0,+1\}\)&
Active l’évaluation additive, multiplicative ou le pas de seuil ; ne déplace pas
l’état.\\

Transporte espacial&
\(T_N,T_E,T_S,T_O,F_{\rm loc},F_{\rm obs}\)&
Réalise des translations cardinales, une rétroaction de feuillet et une chronologie à deux
curseurs.\\

Lectores modulares&
\(\rho_9,q_9,\rho_3^{\rm or},c_3,q_{1000},r_{1000}\)&
Publient la phase, l’orientation et le bloc décimal en conservant les quotients de
relèvement.\\

Bloc et mémoire&
\(\mathcal K_{27}=F_{\rm obs}^{27}\), \(\mathcal N_{27}\)&
Le premier compose vingt-sept transports ; le second filtre des événements de
mémoire. Un indice identique n’implique pas un opérateur identique.\\

État dodécaphasique&
\(\mathcal R_{12},C_{12},\Theta_{108},\Gamma_{108}\)&
Distingue le registre enrichi, le secteur, la loi cyclique et le support ordonné.\\

Périodicité et profondeurs&
\(C_{12}\hookrightarrow C_{108}\twoheadrightarrow C_9\),
\(w_{108},w_{1080}\)&
Conserve l’extension non scindée, les mots de longueurs différentes et
la mémoire verticale.\\

Canaux \(90/120\)&
\(\mathcal F_6,\mathcal F_8,\mathcal I_{6\to8},
S_{90},S_{120}\)&
Sépare les incidences finies, les queues analytiques et leurs normalisations.\\

Bidirectionnalité&
\(P_\pm,\operatorname{rev},\operatorname{Ext}_\pm,
N_\pm,\beta,C_M\)&
Réalise l’échange, l’inversion de route, le prolongement futur et passé,
la valence, le bilan et la conjugaison.\\

Émission et coinduction&
\(L_0,L_1,R_{36},G_9,\varprojlim\operatorname{Cyl}_m\)&
Relève le germe, ouvre la frontière et organise des prolongements finis et
compatibles.\\

Lecteurs numériques&
\(\mathscr C_\pi,\mathscr P_e,F_{\rm av}\)&
Construisent des encadrements rationnels, une propagation factorielle et une auto-échelle de
Fibonacci.\\

Projection exceptionnelle&
\(\Gamma_{A_W}\), poids, support projectif&
Transporte la même émission vers une incidence ternaire sans sélectionner de chiffres
archimédiens.\\

Atlas et réalisation&
signature de route, atlas \(81/56/13/324\), \(C_M\)&
Classe les histoires et les caractères avant d’appliquer une carte dimensionnelle.\\

Validation finie&
recensements, preuves de nécessité structurelle, identités d’intégrité et
contrôles d’indépendance&
Vérifie la couverture à la profondeur exécutée et sépare la construction,
la vérification indépendante et les données de confrontation.\\
\bottomrule
\end{longtable}
\endgroup

Les huit classes décrivent la nature mathématique ; les quatorze regroupements,
la fonction d’exécution. Ce sont deux gradations du même inventaire.

\Needspace{.92\textheight}
\subsection{Domaine et typage des \(34\) entrées de l’opérateur d’actualisation}

\begin{equation}
 \mathfrak O_{\rm TPK}:=(\mathfrak o_1,\ldots,\mathfrak o_{34}).
 \label{p12-83-c1-s6:eq:u006f-registro-operadores}
\end{equation}

\begingroup
\scriptsize
\renewcommand{\arraystretch}{1.12}
\begin{longtable}{@{}r p{.19\textwidth}p{.28\textwidth}p{.42\textwidth}@{}}
\toprule
\#&Symbole&Domaine et codomaine&Règle ou fonction\\
\midrule
\endfirsthead
\toprule
\#&Symbole&Domaine et codomaine&Règle ou fonction\\
\midrule
\endhead
1&\(\mathcal A_9\)&\((\mathbb Z/9)^2\), matrice et loi interne&
Cellule arithmétique additive--multiplicative.\\
2&\((\rho_9,q_9)\)&\(\mathbb Z_{>0}\to
\{1,\ldots,9\}\times\mathbb Z_{\geq0}\)&
Division nonadique avec représentant positif et quotient.\\
3&\(\rho_3^{\rm or},c_3\)&\(\mathbb Z\to
\{-1,0,+1\}\times\mathbb Z\)&
Relèvement ternaire orienté \(n=\rho_3^{\rm or}(n)+3c_3(n)\).\\
4&\(\iota_{1000}:=(q_{1000},r_{1000})\)&\(\mathbb Z\to
\mathbb Z\times\{0,\ldots,999\}\)&
Division euclidienne \(N=1000q_{1000}+r_{1000}\).\\
5&\(T_d\)&\(X_9\to X_9\), \(d\in\{N,E,S,O\}\)&
Translation cardinale orientée.\\
6&\(F_{\rm loc}\)&\(\mathcal Q_{\rm loc}\times
\mathscr D_{\rm card}\to\mathcal Q_{\rm loc}\)&
Rétroaction locale de position et de feuillet.\\
7&\(T_{\min}\)&\(\Omega_{\rm seed}\to\mathcal U_6\)&
Émetteur minimal des deux curseurs, y compris la signature coordonnée.\\
8&\(G\)&\(\operatorname{im}s_\Sigma\times
\operatorname{im}s_\Pi\to\{0,\ldots,999\}^6\)&
Couplage décimal coordonnée par coordonnée.\\
9&\(\Psi_6\)&\(\mathcal U_6\to\mathbb F_3^6\)&
Réduction ternaire du catalogue.\\
10&\(\mathcal R_3\)&\(\mathcal Q_{\rm loc}\times
\mathscr D_{\rm card}^3\to\{-1,0,+1\}\)&
Facteur observable local de trois pas.\\
11&\(\mathcal K_{27}\)&\(\mathcal Q_{\rm obs}\to\mathcal Q_{\rm obs}\)&
Composition \(F_{\rm obs}^{27}\).\\
12&\(\mathcal X_{\rm TPK}^{\rm enr}\)&catégorie d'états holonomiques intégraux&
Conserve toutes les coordonnées typées du noyau.\\
13&\(\mathcal H\)&relation sur les états holonomiques intégraux&
Transition de Hensel qui retient le résidu et le quotient.\\
14&\(\mathcal C_{3,1000}\)&états croisés
\((N,L,D,K,A,B)\)&
Actualisation exacte des cylindres ternaires--décimaux.\\
15&\(\Sigma_{\rm B}\)&mémoires de frontière \(\to\) signature conjointe&
Publie l’orientation de bloc en conservant son registre généalogique.\\
16&\({\rm EF}\text{--}G_9\)&relation d’extension&
Filtre les extensions compatibles et organise la généalogie nonadique.\\
17&\(\Gamma_9\)&histoires enrichies \(\longrightarrow\) histoires&
Holonomie d’un tour de phase avec mémoire de frontière.\\
18&\(X_\chi\)&\(\varprojlim_m\operatorname{Surv}^{(\chi)}_m\)&
Section globale lorsque la non-vacuité, la compatibilité et l’unicité sont
démontrées à toute profondeur.\\
19&\(\mathcal P_{\rm APP}^{(2)}\)&catégorie bivariée de chemins&
Transporte simultanément l’évaluation multiplicative et l’accumulation
additive.\\
20&\(F_{\rm obs}\)&\(\mathcal Q_{\rm obs}\to\mathcal Q_{\rm obs}\)&
Chronologie observable déterministe à deux curseurs.\\
21&\(\mathcal A_{42}\)&alphabet de quarante-deux triades&
Support décimal élargi de l’émission finie.\\
22&\(\operatorname{Ext}_r\)&
\(\operatorname{Ext}_r\subseteq\mathcal F_q\times\mathcal F_{q'}\)&
Relation typée de prolongement des fibres.\\
23&\(\mathcal K_{\rm ph}\)&\(\mathcal U_6\times C_9\) sur lui-même&
Transfert de continuations compatibles avec phase explicite.\\
24&\((P,H,Q)\)&\(\mathcal X_\times\to\mathcal X_\times\)&
Avancée, demi-tour et quart de tour des germes.\\
25&\((c_{\rm mode},c_{\rm or})\)&caminos \(\to\mathbb Z^2\)&
Cocycles de changement de mode et d’orientation.\\
26&\(E_{20}\)&états finis \(\to\) vingt blocs&
Prolongement triadique fini avec livre de provenance.\\
27&\(\mathcal R_{12}\)&histoire enrichie \(\to\) registre de douze
secteurs&
Système temporel orienté avec route, signature, retenue et registre généalogique.\\
28&\(R_{\rm sgn}\)&état terminal \(\to U_{12}^{\rm sgn}
\subseteq\mathbb Z^{12}\)&
Lecteur de l’observable harmonique signée.\\
29&\(\mathcal H_{12}\)&\(\mathbb Z^{12}\to\mathbb Z^{12}\)&
Transformation par trois orbites de pas trois ; inverse
\(\mathcal H_{12}/4\) sur l’image pertinente.\\
30&\(\mathcal E_{\rm dod}=(D_3,D_4,Q)\)&
\(\mathbb Z^{12}\to\mathbb Z^{25}\)&
Extracteur transversal conjoint d’arités trois et quatre.\\
31&\(\mathcal C_\alpha\)&blocs et retenues \(\to\) blocs&
Récurrence dodécaphasique de clôture en base mille.\\
32&\((L_{\rm tick},\chi_{\rm mass})\)&histoire enrichie \(\to\)
caractère temporel&
Interface dimensionnelle ultérieure ; elle ne participe pas au constructeur numérique de
cet ouvrage.\\
33&\(\mathcal W_\pm\)&réseau bidirectionnel d’indice douze&
Lecture dans les deux sens en conservant l’orientation et la provenance.\\
34&\(\mathcal A_{\rm species}\)&familles d’extension \(\to\) atlas&
Interface de réalisation ultérieure, non un quatrième primitif.\\
\bottomrule
\end{longtable}
\endgroup

\subsection{Composition causale de l'actualisation de l'état holonomique intégral}

\begin{definition}[État canonique complet]
\label{p12-83-c1-s6:def:u006f-esquema-estado}
Soit
\[
 \mathcal E:=
 \mathbb Z_{\rm quo}\times\mathbb Z_{\rm car}\times
 \mathsf{Mem}\times\mathsf{Term}\times\mathsf{Cyl}\times
 \mathsf{Fr}\times\mathsf{Pref}\times\Lambda .
\]
L'état holonomique intégral canonique à l'instant \(t\) est
\begin{equation}
 x_t=(p_t,m_t,\sigma_t,b_t,r_t;
 q_t,c_t,\eta_t,s_t,C_t,\partial_t,N_t,\lambda_t)
 \in
 X_9\times\{\Sigma,\Pi\}\times\mathsf{TRIT}\times
 C_{12}\times C_9\times\mathcal E.
 \label{p12-83-c1-s6:eq:u006f-esquema-estado}
\end{equation}
Les abréviations d’état utilisées dans d’autres unités sont des projections de
ce tuple. En particulier, ni le quotient \(q_t\), ni la retenue
\(c_t\), ni le terminal \(s_t\), ni le cylindre \(C_t\), ni la frontière
\(\partial_t\), ni le préfixe \(N_t\), ni le livre \(\lambda_t\) ne sont omis.
\end{definition}

Soit \(\rho_t\) la flèche observable déterminée par \(x_t\). Les lois de
transport de U005 et la transition \(\mathcal C_{3,1000}\) déterminent un
endomorphisme de mémoire
\begin{align}
 \mathbf E_t:\mathcal E&\longrightarrow\mathcal E,\notag\\
 (q,c,\eta,s,C,\partial,N,\lambda)&\longmapsto
 \left(
 \begin{aligned}
 &\mathsf H(\rho_t)q,
  \mathsf C(\rho_t)c+\kappa(\rho_t),
  \mathsf M(\rho_t)\eta,\\
 &\mathsf T(\rho_t)s,
  \mathsf{Cyl}(\rho_t)C,
  \mathsf{Fr}(\rho_t)\partial,\\
 &\mathsf{Pref}(\rho_t)N,
  \Lambda(\rho_t)\lambda
 \end{aligned}
 \right).
 \label{p12-83-c1-s6:eq:u006f-actualizacion-memoria}
\end{align}
Chaque composante publiée dans \eqref{p12-83-c1-s6:eq:u006f-actualizacion-memoria} possède le
domaine indiqué par sa fibre ; \(\kappa(\rho_t)\) est la cochaîne de retenue.

Nous définissons trois endomorphismes du produit canonique. Si
\(\sigma'=M_{\rm ph}(g(t))\) et \(m'=\mu(\sigma',m)\), alors
\begin{align}
 \mathsf{Sel}_t(p,m,\sigma,b,r;e)
  &:=(p,m',\sigma',b,r;e),\label{p12-83-c1-s6:eq:u006f-selector-tipado}\\
 \mathsf{Tra}_t(p,m,\sigma,b,r;e)
  &:=(T_{d(t)}p,m,\sigma,b,r;e),\label{p12-83-c1-s6:eq:u006f-transporte-tipado}\\
 \mathsf{Upd}_t(p,m,\sigma,b,r;e)
  &:=(p,m,\sigma,b',r';\mathbf E_t(e)),
 \label{p12-83-c1-s6:eq:u006f-memoria-tipada}
\end{align}
où \((b',r')\) est la destination de \(\rho_t:(b,r)\to(b',r')\). L’actualisation
causale est la composition d’applications, non la composition d’une application avec une valeur
tritique :
\begin{equation}
 \boxed{\mathcal U_t:=
 \mathsf{Upd}_t\circ\mathsf{Tra}_t\circ\mathsf{Sel}_t.}
 \label{p12-83-c1-s6:eq:u006f-tuberia}
\end{equation}

\begin{theorem}[Clôture de l’actualisation causale]
\label{p12-83-c1-s6:thm:u006f-cierre-tuberia}
L’application \(\mathcal U_t\) est un endomorphisme de l’état canonique
complet. Les lecteurs, prolongements et projections sont appliqués après
\(\mathcal U_t\) uniquement lorsque leur domaine a été produit.
\end{theorem}

\begin{proof}
\(\mathsf{Sel}_t\) ne change que le mode et le TRIT ;
\(\mathsf{Tra}_t\) ne change que la position APP ; et
\(\mathsf{Upd}_t\) applique le produit typé
\eqref{p12-83-c1-s6:eq:u006f-actualizacion-memoria} et la flèche de phase. Les trois
codomaines sont le même produit de la
\cref{p12-83-c1-s6:def:u006f-esquema-estado} ; leur composition est donc définie et
conserve toutes les coordonnées de l’état.
\end{proof}

\subsection{Séparation typée entre opérateurs de transport, de mémoire, de calendrier, d’extraction et de retour}

L’inventaire impose
\begin{equation}
 \mathcal S_4\neq D_4\neq\operatorname{sd}_4,
 \qquad
 \mathcal R_{12}\neq C_{12}\neq\Theta_{108}\neq\Gamma_{108},
 \qquad
 F_{\rm obs}\neq\mathcal K_{\rm ph},
 \label{p12-83-c1-s6:eq:u006f-no-identificacion-uno}
\end{equation}
et
\begin{equation}
 U_{12}^{\rm cat}\neq U_{12}^{\rm sgn},
 \qquad
 \mathcal E_{\rm dod}\neq D_{\rm raw}^{90,120}
 \neq\mathcal K_{6\to8}.
 \label{p12-83-c1-s6:eq:u006f-no-identificacion-dos}
\end{equation}
Ce sont des incompatibilités de type, non de simples inégalités de valeur.

\begin{criterion}[Vérification de complétude opératoire]
\label{p12-83-c1-s6:crit:u006f-completitud}
Une exécution n'est opératoirement complète que si :
\begin{enumerate}
 \item elle déclare les entrées de \(\mathfrak O_{\rm TPK}\) employées ;
 \item elle enregistre le domaine et le codomaine de chaque composition ;
 \item elle conserve l’ordre causal de \eqref{p12-83-c1-s6:eq:u006f-tuberia} ;
 \item elle distingue dynamique, transfert, lecture, prolongement et sortie ;
 \item elle permet de reconstruire à partir de \(\lambda_t\) chaque changement de bloc et
 chaque retenue ;
 \item elle ne promeut pas une vérification finie en une section infinie sans
 prouver la non-vacuité, la compatibilité et l’unicité à toute profondeur.
\end{enumerate}
Tout symbole utilisé sans type ou toute fusion interdite par
\eqref{p12-83-c1-s6:eq:u006f-no-identificacion-uno} et
\eqref{p12-83-c1-s6:eq:u006f-no-identificacion-dos} constitue un échec vérifiable.
\end{criterion}


% Unidad U006E. Separación tipada de los canales 90/120 y de los símbolos
% epsilon. La realización de incidencia se construye después de W12 y W24.

\section{Décomposition opératorielle des canaux d’incidence \texorpdfstring{\(90/120\)}{90/120} : sélecteur temporel, générateur nilpotent et composante angulaire}
\label{p12-83-c1-s7:chap:canales-90-120-epsilon}

Les entiers \(90\) et \(120\) apparaissent dans plusieurs strates de la construction. Leur
coïncidence numérique n’établit pas une identité d’objets. De même,
la lettre epsilon désigne trois constructions aux domaines incompatibles. La
finalité de cette unité est de fixer leurs types avant de les employer et d’indiquer
avec précision quelles réalisations exigent des structures qui n’ont pas encore été
construites.

\subsection{Opérateurs temporel, nilpotent et angulaire : domaines et codomaines}

\subsubsection{Action du sélecteur temporel tritique sur \(\{-1,0,+1\}\) : ouverture, seuil et retour avec mémoire}

Le sélecteur interne du calendrier est
\begin{equation}
 \varepsilon_9:\{1,\ldots,9\}
 \longrightarrow\{-1,0,+1\},
 \qquad
 \varepsilon_9(r)=
 \begin{cases}
  +1,&r\in\{1,4,7\},\\
  -1,&r\in\{2,5,8\},\\
  0,&r\in\{3,6,9\}.
 \end{cases}
 \label{p12-83-c1-s7:eq:u006e-epsilon-temporal}
\end{equation}
Son codomaine est le TRIT orienté. Il décide quel feuillet agit et ne déplace
par lui-même aucune position.

\subsubsection{Générateur parabolique nilpotent \(J_0^2=0\) et évolution au seuil temporel}

La classe
\begin{equation}
 \varepsilon_{\rm nil}
 \in
 \mathcal Q_{\rm n}:=
 \mathbb F_3[\varepsilon_{\rm nil}]
 /(\varepsilon_{\rm nil}^{\,2})
 \label{p12-83-c1-s7:eq:u006e-epsilon-nilpotente}
\end{equation}
satisfait
\begin{equation}
 \varepsilon_{\rm nil}\neq0,
 \qquad
 \varepsilon_{\rm nil}^{\,2}=0.
 \label{p12-83-c1-s7:eq:u006e-epsilon-nilpotente-relacion}
\end{equation}
C’est un élément d’une algèbre quadratique dégénérée ; ce n’est pas un sélecteur
temporel.

\subsubsection{Résidu angulaire induit après la sélection temporelle et conservation de l’orientation}

Le symbole
\begin{equation}
 \varepsilon_{\rm res}^{\circ}\in\mathbb R
 \label{p12-83-c1-s7:eq:u006e-epsilon-residual-tipo}
\end{equation}
désigne, lorsque ses paramètres géométriques seront introduits, le résidu d’une
identité angulaire. À cette étape, seul son type est réservé. Il n’est obtenu
ni à partir de \(\varepsilon_9\) ni à partir de
\(\varepsilon_{\rm nil}\), et n’intervient pas dans la construction de
\(\pi,e,\varphi\) ou de \(\alpha\).

\begin{proposition}[Non-identification des symboles epsilon]
\label{p12-83-c1-s7:prop:u006e-no-identificacion-epsilon}
Il n’existe pas d’égalité bien typée entre
\[
 \varepsilon_9,\qquad
 \varepsilon_{\rm nil},\qquad
 \varepsilon_{\rm res}^{\circ}.
\]
\end{proposition}

\begin{proof}
Le premier objet est une fonction entre ensembles finis ; le deuxième, un
élément nilpotent d’une algèbre sur \(\mathbb F_3\) ; le troisième, un
scalaire réel avec unité angulaire. Leurs domaines, codomaines et lois de
composition sont distincts.
\end{proof}

\subsection{Distribution des indices \(90/120\) entre les \(6\) objets du système d’incidence}

\begin{definition}[Extracteur dodécaphasique]
\label{p12-83-c1-s7:def:u006e-extractor-dodecafasico}
L’extracteur entier
\begin{equation}
 \mathcal E_{\rm dod}:=(D_3,D_4,Q):
 \mathbb Z^{12}\longrightarrow\mathbb Z^{25}
 \label{p12-83-c1-s7:eq:u006e-extractor-dodecafasico}
\end{equation}
lit des différences d’arités trois et quatre et une coordonnée de somme. Les
indices ne sont pas des longueurs angulaires.
\end{definition}

\begin{definition}[Différence brute de séries]
\label{p12-83-c1-s7:def:u006e-diferencia-cruda}
Pour \(|q|<1\), on définit
\begin{equation}
 S_m^{\rm raw}(q):=\frac{q^m}{1-q^{3m}},
 \qquad
 D_{\rm raw}^{90,120}(q)
 :=S_{90}^{\rm raw}(q)-S_{120}^{\rm raw}(q).
 \label{p12-83-c1-s7:eq:u006e-series-crudas}
\end{equation}
Cet objet est une fonction analytique de \(q\) ; ce n’est pas un recensement fini.
\end{definition}

\begin{definition}[Familles marquées d’incidence]
\label{p12-83-c1-s7:def:u006e-familias-incidencia}
Une fois construites une hexade \(H\), son image
\(\ell(H)\) dans une dodécade binaire \(D\), et l’octade relevée
\(O_H\), on définira
\begin{align}
 \mathcal F_6(H)
 &:=\ell(H)\times\binom{\ell(H)}2,
 \label{p12-83-c1-s7:eq:u006e-familia-seis}\\
 \mathcal F_8(H)
 &:=O_H\times\binom{\ell(H)}2.
 \label{p12-83-c1-s7:eq:u006e-familia-ocho}
\end{align}
Par construction,
\begin{equation}
 |\mathcal F_6(H)|
 =6\binom62=90,
 \qquad
 |\mathcal F_8(H)|
 =8\binom62=120.
 \label{p12-83-c1-s7:eq:u006e-cardinales-incidencia}
\end{equation}
L’inclusion \(\ell(H)\hookrightarrow O_H\) induira
\begin{equation}
 \mathcal I_{6\to8}:
 \mathcal F_6(H)\hookrightarrow\mathcal F_8(H).
 \label{p12-83-c1-s7:eq:u006e-inclusion-incidencia}
\end{equation}
\end{definition}

Les \cref{p12-83-c1-s7:eq:u006e-familia-seis,p12-83-c1-s7:eq:u006e-familia-ocho} fixent dès maintenant le
type. Leur existence effective sera démontrée après la construction de
\(C_W\), de \(W_{12}\), de la donnée binaire \(G_{24},D,\ell\) et de \(W_{24}\).
Cette branche n’est pas anticipée ici.

\begin{definition}[Combinaison normalisée de queues]
\label{p12-83-c1-s7:def:u006e-combinacion-colas}
Pour des paramètres \(q_-,q_+\in(0,1)\), soient
\begin{equation}
 S_m(q):=\frac{q^m}{1-q^{3m}},
 \qquad
 \Delta S_m:=S_m(q_-)-S_m(q_+).
 \label{p12-83-c1-s7:eq:u006e-colas-normalizadas}
\end{equation}
La fonctionnelle
\begin{equation}
 \mathcal K_{6\to8}
 :=12\,\Delta S_{90}-\Delta S_{120}
 \label{p12-83-c1-s7:eq:u006e-funcional-seis-ocho}
\end{equation}
est une combinaison analytique de queues. Son coefficient douze appartient à sa
normalisation et n’en fait pas le registre
\(\mathcal R_{12}\).
\end{definition}

\begin{definition}[Semi-groupe des degrés et des moments]
\label{p12-83-c1-s7:def:u006e-semigrupo-momentos}
Le semi-groupe additif engendré par les deux degrés est
\begin{equation}
 \mathfrak S_{90,120}
 :=\langle90,120\rangle
 =\{90a+120b:a,b\in\mathbb N_0\}.
 \label{p12-83-c1-s7:eq:u006e-semigrupo}
\end{equation}
Pour \(n\in\mathfrak S_{90,120}\), le moment orienté est
\begin{equation}
 s_n:=q_-^n-q_+^n.
 \label{p12-83-c1-s7:eq:u006e-momento}
\end{equation}
\end{definition}

Les moments sont des monômes différentiels indexés par un semi-groupe ; les
séries de \eqref{p12-83-c1-s7:eq:u006e-colas-normalizadas} sont des sommes infinies de ces
moments sur des progressions spécifiques.

\subsection{Défaut normalisé de l’incidence}

Lorsque l’inclusion \(\mathcal I_{6\to8}\) aura été construite, sa
différence de cardinalités sera
\[
 |\mathcal F_6(H)|-|\mathcal F_8(H)|=90-120=-30.
\]
Si l’on déclare la normalisation
\[
 24\binom62=360,
\]
on obtient
\begin{equation}
 \frac{90-120}{24\binom62}
 =\frac{6-8}{24}
 =-\frac1{12}.
 \label{p12-83-c1-s7:eq:u006e-defecto-normalizado}
\end{equation}
L’inclusion détermine \(-30\) ; le rationnel \(-1/12\) dépend en outre de
la normalisation déclarée. Il ne doit pas être présenté comme si la cardinalité
fixait à elle seule le dénominateur.

\subsection{Invariants de typage}

\begin{theorem}[Séparation des six objets]
\label{p12-83-c1-s7:thm:u006e-separacion-seis-objetos}
Les objets
\begin{equation}
 \mathcal E_{\rm dod},\quad
 D_{\rm raw}^{90,120},\quad
 \mathcal I_{6\to8},\quad
 \mathcal K_{6\to8},\quad
 \mathfrak S_{90,120},\quad
 \varepsilon_{\rm res}^{\circ}
 \label{p12-83-c1-s7:eq:u006e-seis-objetos}
\end{equation}
sont mutuellement distincts par leur type : ce sont, respectivement, un lecteur entier,
une fonction analytique, un morphisme fini, une combinaison de séries, un
semi-groupe d’indices et un scalaire réel angulaire.
\end{theorem}

\begin{proof}
La conclusion s’obtient en comparant leurs domaines et codomaines :
\[
 \mathbb Z^{12}\to\mathbb Z^{25},\quad
 \mathbb D\to\mathbb C,\quad
 \mathcal F_6(H)\to\mathcal F_8(H),\quad
 (0,1)^2\to\mathbb R,\quad
 \mathbb N_0^2\to\mathbb N_0,\quad
 \{\ast\}\to\mathbb R.
\]
Il n’y a pas deux signatures égales.
\end{proof}

\begin{criterion}[Critères de réfutation]
\label{p12-83-c1-s7:crit:u006e-refutacion}
La classification est réfutée si l’un de ses domaines ou
codomaines est altéré, si le recensement de
\eqref{p12-83-c1-s7:eq:u006e-cardinales-incidencia} ne reproduit pas \(90\) et \(120\), si
l’on attribue \(-1/12\) à l’inclusion sans déclarer la normalisation, ou si
l’on présente le résidu angulaire comme une sortie du générateur de constantes de
cette monographie.
\end{criterion}


% Unidad U006A. Factor local, observacion estroboscopica y transicion exacta
% de cilindros.

\section{Facteur local de transport et transition ternaire--décimale}
\label{p12-83-c1-s8:chap:tpk-factor-local-transicion-cilindros}

Cette unité sépare trois opérations qui interviennent dans le transport du
Topological Phase Kernel et qui ne partagent pas de domaine : la réalisation cardinale
d’un trit, l’échantillonnage stroboscopique d’une suite et l’actualisation
exacte d’une paire de cylindres ternaires--décimaux. La première est une dynamique
finie ; la deuxième est un lecteur temporel ; la troisième conserve les entiers
qui permettent de décider de la compatibilité sans arithmétique à virgule flottante.

\subsection{État local et rétroaction de feuillet}

Soit
\[
 X_9=(\mathbb Z/9\mathbb Z)^2,
 \qquad
 \mathscr D_{\rm card}=\{N,E,S,O\},
\]
avec
\[
 \delta(N)=(-1,0),\quad \delta(E)=(0,1),\quad
 \delta(S)=(1,0),\quad \delta(O)=(0,-1).
\]
Les évaluations visibles de l’APP sont
\[
 L_\Sigma(x,y)=\rho_9((x+1)+(y+1)),
 \qquad
 L_\Pi(x,y)=\rho_9((x+1)(y+1)),
\]
où \(\rho_9(n)=1+((n-1)\bmod9)\). Nous définissons en outre
\[
 \rho_3^{\rm or}([0])=0,
 \qquad
 \rho_3^{\rm or}([1])=+1,
 \qquad
 \rho_3^{\rm or}([2])=-1.
\]

\begin{definition}[Facteur local de l’APP--TRIT]
\label{p12-83-c1-s8:def:u006a-factor-local}
L’espace des états locaux est
\[
 \mathcal Q_{\rm loc}=X_9\times\{\Sigma,\Pi\},
 \qquad |\mathcal Q_{\rm loc}|=162.
\]
Pour \(q=(p,m)\), \(d\in\mathscr D_{\rm card}\), nous posons
\[
 p_d=p+\delta(d),
 \qquad
 a_d(q)=\rho_3^{\rm or}([L_m(p_d)]_3),
\]
et
\[
 \mu(a,m)=
 \begin{cases}
  \Sigma,&a=+1,\\
  \Pi,&a=-1,\\
  m,&a=0.
 \end{cases}
\]
La transition locale est
\[
 F_d^{\rm loc}(p,m)=(p_d,\mu(a_d(p,m),m)).
\]
\end{definition}

La position est modifiée avant l’évaluation du feuillet. Le trit nul conserve le
dernier mode actif ; il n’élimine pas le feuillet et ne réinitialise pas l’histoire.

Pour un mot cardinal \(u=d_1d_2d_3\), soit
\[
 q_j=F_{d_j}^{\rm loc}(q_{j-1}),\qquad q_0=q,
\]
et notons
\[
 R_3(q,u)\in\{-1,0,+1\}
\]
le trit produit au troisième déplacement.

\begin{definition}[Observation de pas trois]
\label{p12-83-c1-s8:def:u006a-obs3}
Sur une suite orientée \(a=(a_t)_{t\geq1}\), nous définissons
\[
 \operatorname{Obs}_3(a)=(a_3,a_6,a_9,\ldots).
\]
L’application \(R_3\) prend un état et un mot cardinal de longueur trois ;
\(\operatorname{Obs}_3\) prend une suite et publie une sous-suite. Ce sont
des opérateurs de types distincts.
\end{definition}

\begin{theorem}[Surjectivité exacte du facteur échantillonné]
\label{p12-83-c1-s8:thm:u006a-sobreyectividad-r3}
Pour tout \(q\in\mathcal Q_{\rm loc}\) et tout
\(b\in\{-1,0,+1\}\), la fibre
\[
 \mathcal R_3(q,b)
 =\{u\in\mathscr D_{\rm card}^3:R_3(q,u)=b\}
\]
est non vide et satisfait
\[
 \boxed{8\leq |\mathcal R_3(q,b)|\leq40.}
\]
En particulier, les \(162\cdot3=486\) fibres état--sortie sont non vides.
\end{theorem}

\begin{proof}
À partir de chacun des \(162\) états, on énumère les \(4^3=64\) mots
cardinaux et l’on applique littéralement la récurrence de la
\cref{p12-83-c1-s8:def:u006a-factor-local}. Le recensement exhaustif des \(486\) paires
\((q,b)\) a pour minimum huit, pour maximum quarante et aucune valeur nulle.
L’énoncé est donc une identité finie sur le système défini, non
une hypothèse de mélange.
\end{proof}

\begin{corollary}[Réalisation d’une suite orientée arbitraire]
\label{p12-83-c1-s8:cor:u006a-realizacion-sucesion}
Pour toute suite
\((b_k)_{k\geq1}\in\{-1,0,+1\}^{\mathbb N}\), il existe une histoire cardinale
dont l’observation satisfait
\[
 (a_3,a_6,a_9,\ldots)=(b_1,b_2,b_3,\ldots).
\]
\end{corollary}

\begin{proof}
Un ordre total de \(\mathscr D_{\rm card}^3\) étant fixé, à l’étape \(k\), on
prend le premier élément de \(\mathcal R_3(q_{k-1},b_k)\). L’état terminal
du bloc appartient de nouveau à \(\mathcal Q_{\rm loc}\), de sorte que le
choix peut être itéré indéfiniment.
\end{proof}

Le corollaire démontre la capacité de réalisation du facteur échantillonné. Il ne
sélectionne pas les sections de \(\pi,e,\varphi\), ne fixe pas les deux trits
intermédiaires et ne remplace pas la compatibilité de l'état holonomique intégral.

\subsection{Cylindres des \(2\) lecteurs}

Un mot ternaire de longueur \(L\), interprété comme l’entier \(N\),
détermine
\[
 I_3(N,L)=\left[\frac{N}{3^L},\frac{N+1}{3^L}\right).
\]
Un préfixe de \(K\) triades décimales, interprété comme l’entier \(D\),
détermine
\[
 I_{1000}(D,K)=
 \left[\frac{D}{1000^K},\frac{D+1}{1000^K}\right).
\]
Nous introduisons les résidus croisés
\[
 A=N1000^K-D3^L,
 \qquad
 B=(N+1)1000^K-D3^L=A+1000^K.
\]

\begin{lemma}[Critère entier de compatibilité]
\label{p12-83-c1-s8:lem:u006a-criterio-entero}
Les cylindres \(I_3(N,L)\) et \(I_{1000}(D,K)\) se rencontrent si et seulement si
\[
 \boxed{A<3^L,\qquad B>0.}
\]
De manière équivalente,
\[
 N1000^K<(D+1)3^L,
 \qquad
 (N+1)1000^K>D3^L.
\]
\end{lemma}

\begin{proof}
Deux intervalles semi-ouverts \([a,b)\) et \([c,d)\) se rencontrent
exactement lorsque \(a<d\) et \(c<b\). La multiplication par
\(3^L1000^K>0\) produit les deux inégalités. Soustraire \(D3^L\) et
\(N1000^K\), respectivement, donne la forme en \(A,B\).
\end{proof}

\subsection{Opérateur de complétion \(C_{3,1000}\) : extension ternaire et conservation de l’unité en \(10^3\) positions}

Un nouveau bloc de six trits est représenté par
\[
 u\in\{0,\ldots,3^6-1\}=\{0,\ldots,728\}.
\]
L’actualisation ternaire est
\[
 N'=729N+u,
 \qquad L'=L+6.
\]
La publication décimale présente deux cas. Si le cylindre raffiné ne détermine pas
encore une nouvelle triade, alors \(K'=K\), \(D'=D\). S’il détermine
\(\delta\in\{0,\ldots,999\}\), alors
\[
 K'=K+1,
 \qquad D'=1000D+\delta.
\]

\begin{definition}[État croisé et transition exacte]
\label{p12-83-c1-s8:def:u006a-c31000}
L’état des cylindres est le sextuplet
\[
 \mathfrak c=(N,L,D,K,A,B)
\]
soumise à
\[
 A=N1000^K-D3^L,
 \qquad B=A+1000^K.
\]
La transition \(\mathcal C_{3,1000}\) adjoint \(u\) et, le cas échéant,
\(\delta\), et actualise
\[
 (N,L,D,K,A,B)\longmapsto
 (N',L',D',K',A',B')
\]
par
\[
 A'=
 \begin{cases}
  729A+u1000^K,&K'=K,\\[1mm]
  729000A+u1000^{K+1}-\delta\,729\,3^L,&K'=K+1,
 \end{cases}
 \qquad
 B'=A'+1000^{K'}.
\]
Le prolongement est admissible exactement lorsque
\[
 A'<3^{L'},\qquad B'>0.
\]
\end{definition}

\begin{theorem}[Exactitude et suffisance de l’état croisé]
\label{p12-83-c1-s8:thm:u006a-exactitud-c31000}
La transition de la \cref{p12-83-c1-s8:def:u006a-c31000} conserve exactement les
identités
\[
 A'=N'1000^{K'}-D'3^{L'},
 \qquad
 B'=(N'+1)1000^{K'}-D'3^{L'}.
\]
Par conséquent, \((A',B',L',K')\) suffit à décider de la compatibilité de la
nouvelle paire de cylindres, tandis que \((N',D')\) conserve sa provenance
reconstructive.
\end{theorem}

\begin{proof}
Si aucune triade décimale n’apparaît,
\[
 N'1000^K-D3^{L+6}
 =(729N+u)1000^K-729D3^L
 =729A+u1000^K.
\]
Si aparece \(\delta\),
\begin{align*}
 N'1000^{K+1}-D'3^{L+6}
 &=(729N+u)1000^{K+1}-(1000D+\delta)7293^L\\
 &=729000A+u1000^{K+1}-\delta\,729\,3^L.
\end{align*}
Dans les deux cas, augmenter \(N'\) d’une unité ajoute \(1000^{K'}\), ce qui
donne \(B'\). Le critère de la \cref{p12-83-c1-s8:lem:u006a-criterio-entero} conclut.
\end{proof}

\begin{proposition}[Compatibilité avec la concaténation]
\label{p12-83-c1-s8:prop:u006a-concatenacion-c31000}
Appliquer successivement \(\mathcal C_{3,1000}\) à deux blocs produit la
même paire \((N'',D'')\) et les mêmes résidus croisés qu’adjoindre d’un seul
coup les deux mots ternaires et les triades décimales publiées dans le
même ordre.
\end{proposition}

\begin{proof}
Les récurrences \(N\mapsto729N+u\) et
\(D\mapsto1000D+\delta\) sont les lois de concaténation positionnelle en
bases \(729\) et \(1000\). Le
\cref{p12-83-c1-s8:thm:u006a-exactitud-c31000} identifie \(A,B\) à des fonctions de ces
quatre entiers ; leur actualisation par étapes coïncide donc avec
l’évaluation postérieure à la concaténation.
\end{proof}

\subsection{Décomposition du Topological Phase Kernel en opérateurs d’avancée, de rotation, de retenue et de mémoire}

La chaîne correcte est
\[
 (q,u)\xmapsto{\ R_3\ }b,
 \qquad
 (a_t)_{t\geq1}
 \xmapsto{\ \operatorname{Obs}_3\ }
 (a_{3k})_{k\geq1},
 \qquad
 (\mathfrak c,u,\delta)
 \xmapsto{\ \mathcal C_{3,1000}\ }
 \mathfrak c'.
\]
\(R_3\) prouve la réalisabilité locale ; \(\operatorname{Obs}_3\) publie une
sous-suite ; \(\mathcal C_{3,1000}\) transporte la compatibilité entre deux
bases. Aucun des trois ne peut remplacer la relation de survie
du TPK complet.

\begin{criterion}[Critères de réfutation]
\label{p12-83-c1-s8:crit:u006a-refutacion}
La construction est réfutée dans son domaine si l’un quelconque des
faits suivants se produit :
\begin{enumerate}
 \item l’une des \(486\) fibres \(\mathcal R_3(q,b)\) est vide ou se trouve
       en dehors de l’intervalle \([8,40]\) ;
 \item le trit nul change le mode actif ;
 \item on identifie \(R_3\) à \(\operatorname{Obs}_3\) ;
 \item la récurrence de \(A'\) ne reproduit pas sa définition directe ;
 \item un prolongement déclaré compatible viole \(A'<3^{L'}\) ou
       \(B'>0\) ;
 \item la décision dépend d’une tolérance numérique absente des
       inégalités entières.
\end{enumerate}
\end{criterion}


% Unidad U006D. Registro dodecafásico integral y descomposición 1+5+6.

\section{Registre dodécaphasique entier, extension cyclique et décomposition \texorpdfstring{\(1+5+6\)}{1+5+6}}
\label{p12-83-c1-s9:chap:registro-dodecafasico-integral}

La période observable de cent huit positions contient douze secteurs de
neuf positions. Cette égalité de cardinalités ne suffit pas à décrire
l’état temporel : il faut distinguer le support ordonné, la loi cyclique,
l’indice sectoriel et le registre enrichi des transports, orientations,
retenues et incidences. Cette unité construit les quatre objets et démontre
la reconstruction entière des douze coordonnées de mémoire.

\subsection{Support ordonné et coordonnée cyclique}

Pour \(m\in\{1,\ldots,12\}\), nous définissons
\begin{equation}
 E_m:=\{9(m-1)+1,\ldots,9m\}.
 \label{p12-83-c1-s9:eq:u006d-bloques-nonadicos}
\end{equation}
Le support ordonné et le groupe cyclique sont
\begin{equation}
 \Gamma_{108}:=
 E_1\sqcup\cdots\sqcup E_{12},
 \qquad
 \Theta_{108}:=\mathbb Z/108\mathbb Z.
 \label{p12-83-c1-s9:eq:u006d-gamma-theta}
\end{equation}
\(\Gamma_{108}\) conserve la position, l’ordre et l’appartenance à un bloc ;
\(\Theta_{108}\) conserve la loi de translation. Ce ne sont pas le même objet.

Pour \(\theta\in\Theta_{108}\), soit
\(\widetilde\theta\in\{0,\ldots,107\}\) son représentant. Nous définissons
\begin{equation}
 \beta(\theta):=
 \left[
 \left\lfloor\frac{\widetilde\theta}{9}\right\rfloor
 \right]_{12},
 \qquad
 r(\theta):=1+(\widetilde\theta\bmod9).
 \label{p12-83-c1-s9:eq:u006d-coordenadas-bloque-residuo}
\end{equation}
La première coordonnée localise le secteur dodécaphasique ; la seconde publie
la position positive \(1,\ldots,9\) à l’intérieur du secteur.

\subsection{Extension cyclique non scindée}

Soient \(C_n=\mathbb Z/n\mathbb Z\). Les applications
\begin{equation}
 \iota:C_{12}\longrightarrow C_{108},
 \quad \iota([b]_{12})=[9b]_{108},
 \qquad
 p:C_{108}\longrightarrow C_9,
 \quad p([\theta]_{108})=[\theta]_9
 \label{p12-83-c1-s9:eq:u006d-mapas-extension}
\end{equation}
forment la suite
\begin{equation}
 0\longrightarrow C_{12}
 \xrightarrow{\;\iota\;}C_{108}
 \xrightarrow{\;p\;}C_9
 \longrightarrow0.
 \label{p12-83-c1-s9:eq:u006d-extension}
\end{equation}
Pour la section ensembliste
\[
 s([r]_9):=[\widetilde r]_{108},
 \qquad 0\leq\widetilde r<9,
\]
le défaut d’additivité est
\begin{equation}
 c(r,r'):=
 \left[
 \left\lfloor
 \frac{\widetilde r+\widetilde r'}9
 \right\rfloor
 \right]_{12},
 \qquad
 s(r)+s(r')
 =s(r+r')+\iota(c(r,r')).
 \label{p12-83-c1-s9:eq:u006d-cociclo-extension}
\end{equation}

\begin{theorem}[Exactitude, non-scindement et classe cohomologique]
\label{p12-83-c1-s9:thm:u006d-extension-no-escindida}
La suite \eqref{p12-83-c1-s9:eq:u006d-extension} est exacte et ne se scinde pas. Pour
l’action triviale,
\begin{equation}
 H^2(C_9,C_{12})
 \simeq C_{12}/9C_{12}
 \simeq C_3,
 \label{p12-83-c1-s9:eq:u006d-cohomologia-extension}
\end{equation}
et le cocycle \eqref{p12-83-c1-s9:eq:u006d-cociclo-extension} représente sa classe
génératrice.
\end{theorem}

\begin{proof}
L’image de \(\iota\) est le sous-groupe des multiples de neuf, qui coïncide
avec \(\ker p\) ; \(p\) est surjective. Cela prouve l’exactitude. Si
l’extension se scindait, \(C_{108}\) serait isomorphe à
\(C_{12}\times C_9\), mais
\[
 \exp(C_{108})=108,
 \qquad
 \exp(C_{12}\times C_9)=\operatorname{mcm}(12,9)=36.
\]
L’identité de cocycle est la division euclidienne de
\(\widetilde r+\widetilde r'\) par neuf. Pour un groupe cyclique à action
triviale,
\(H^2(C_9,C_{12})\simeq C_{12}/9C_{12}\) ; la retenue unitaire se projette sur
\([1]\), qui engendre le quotient d’ordre trois.
\end{proof}

Le non-scindement type le retour :
\[
 g_0(\theta+9)=g_0(\theta),
 \qquad
 \beta(\theta+9)=\beta(\theta)+1\pmod{12}.
\]
Neuf itérations ferment la phase visible et déplacent un secteur ; cent
huit ferment les deux coordonnées de l'horloge. L'état holonomique intégral conserve
en outre la mémoire verticale, le cylindre, la frontière et la généalogie.

\subsection{Registre temporel enrichi : encodage conjoint de la phase, du feuillet, de la frontière et de la mémoire}

\begin{definition}[Registre dodécaphasique orienté]
\label{p12-83-c1-s9:def:u006d-registro-dodecafasico}
Pour une histoire enrichie, soit \(\tau_m\) la sous-route ordonnée sur
\(E_m\), \(\sigma_m\) son orientation, \(\kappa_m\) la retenue qui franchit la
frontière du bloc et \(\Lambda_m\) son livre local d’incidences. Le
registre complet est
\begin{equation}
 \mathcal R_{12}:=
 \bigl(
 (E_m,\tau_m,\sigma_m,\kappa_m,\Lambda_m)
 \bigr)_{m=1}^{12}.
 \label{p12-83-c1-s9:eq:u006d-registro-r12}
\end{equation}
\end{definition}

On distingue ainsi quatre niveaux :
\begin{equation}
 \boxed{
 \mathcal R_{12}
 \neq C_{12}
 \neq\Theta_{108}
 \neq\Gamma_{108}.}
 \label{p12-83-c1-s9:eq:u006d-cuatro-objetos}
\end{equation}
Le premier est un registre enrichi ; le deuxième, un indice de secteur ; le
troisième, un groupe de translations ; le quatrième, un support ordonné. Il existe
des projections entre eux, mais pas d’identifications.

Soit \(q_{\rm mem}\in\mathbb Z_{\geq0}\) le nombre non borné de tours
nonadiques. La coordonnée dodécaphasique n’est que son résidu
\[
 \beta=[q_{\rm mem}]_{12}.
\]
Après cent huit pas,
\[
 \beta\longmapsto\beta,
 \qquad
 q_{\rm mem}\longmapsto q_{\rm mem}+12.
\]
Le calendrier revient ; la mémoire verticale n’est pas réinitialisée.

\subsection{Extracteur signé d’événements}

Soit \((d_t)_{t=1}^{108}\) le registre de chiffres sur
\(\Gamma_{108}\). Nous définissons
\begin{equation}
 \mathcal D_{\rm evt}:=\{2,4,6,8\},
 \qquad
 \Sigma_{12}:=(+,+,+,-,-,-,+,+,+,-,-,-).
 \label{p12-83-c1-s9:eq:u006d-datos-eventos}
\end{equation}
Pour \(m=0,\ldots,11\), soit
\begin{equation}
 e_m:=
 \Sigma_{12}(m+1)\,
 \#\{t\in E_{m+1}:d_t\in\mathcal D_{\rm evt}\}
 \in\mathbb Z.
 \label{p12-83-c1-s9:eq:u006d-eventos-firmados}
\end{equation}
L’application conserve le recensement local et l’orientation sectorielle comme
facteurs séparés.

Nous apparions les deux demi-couronnes par
\begin{equation}
 p_i:=e_i+e_{i+6},
 \qquad
 a_i:=e_i-e_{i+6},
 \qquad 0\leq i\leq5.
 \label{p12-83-c1-s9:eq:u006d-pares-antipodales}
\end{equation}
Nous définissons le recensement total, les cinq différences par rapport à la sixième paire et
les six antisymétries :
\begin{equation}
 s_C:=\sum_{i=0}^{5}p_i,
 \qquad
 M_i:=p_i-p_5\quad(0\leq i\leq4),
 \qquad
 A_6:=(a_0,\ldots,a_5).
 \label{p12-83-c1-s9:eq:u006d-coordenadas-uno-cinco-seis}
\end{equation}

\begin{definition}[Lecteur entier \(1+5+6\)]
\label{p12-83-c1-s9:def:u006d-lector-integral}
Le lecteur linéaire est
\begin{equation}
 \mathcal L_{12}:\mathbb Z^{12}\longrightarrow\mathbb Z^{12},
 \qquad
 \mathcal L_{12}(e)
 :=
 (s_C,M_0,\ldots,M_4,a_0,\ldots,a_5).
 \label{p12-83-c1-s9:eq:u006d-lector-l12}
\end{equation}
La distribution des coordonnées est
\begin{equation}
 \underbrace{1}_{s_C}
 +\underbrace{5}_{M_0,\ldots,M_4}
 +\underbrace{6}_{a_0,\ldots,a_5}
 =12.
 \label{p12-83-c1-s9:eq:u006d-uno-cinco-seis}
\end{equation}
\end{definition}

\begin{theorem}[Reconstruction exacte et image entière]
\label{p12-83-c1-s9:thm:u006d-inversa-l12}
Le lecteur \(\mathcal L_{12}\) est injectif. Sur son image,
\begin{align}
 p_5&=
 \frac{s_C-\sum_{i=0}^{4}M_i}{6},
 &
 p_i&=p_5+M_i\quad(0\leq i\leq4),
 \label{p12-83-c1-s9:eq:u006d-inversa-p}\\
 e_i&=\frac{p_i+a_i}{2},
 &
 e_{i+6}&=\frac{p_i-a_i}{2}
 \quad(0\leq i\leq5).
 \label{p12-83-c1-s9:eq:u006d-inversa-e}
\end{align}
Un tuple de sortie appartient à \(\operatorname{im}\mathcal L_{12}\) si et
seulement si
\begin{equation}
 s_C-\sum_{i=0}^{4}M_i\equiv0\pmod6
 \label{p12-83-c1-s9:eq:u006d-condicion-divisibilidad}
\end{equation}
et, pour les \(p_i\) reconstruits,
\begin{equation}
 p_i\equiv a_i\pmod2
 \qquad(0\leq i\leq5).
 \label{p12-83-c1-s9:eq:u006d-condicion-paridad}
\end{equation}
\end{theorem}

\begin{proof}
De \(M_i=p_i-p_5\), on obtient
\[
 s_C=\sum_{i=0}^{4}(p_5+M_i)+p_5
 =6p_5+\sum_{i=0}^{4}M_i,
\]
ce qui donne la première équation de \eqref{p12-83-c1-s9:eq:u006d-inversa-p} ; les autres sont
immédiates. Additionner et soustraire
\(p_i=e_i+e_{i+6}\) et \(a_i=e_i-e_{i+6}\) produit
\eqref{p12-83-c1-s9:eq:u006d-inversa-e}. La divisibilité par six et les six
congruences de parité sont précisément les conditions nécessaires et
suffisantes pour que toutes les coordonnées reconstruites soient entières.
\end{proof}

La division par six apparaît après le recensement orienté, non pendant celui-ci. Les
conditions d’intégralité de cette unité sont exactement
\eqref{p12-83-c1-s9:eq:u006d-condicion-divisibilidad} et
\eqref{p12-83-c1-s9:eq:u006d-condicion-paridad} ; on n’attribue pas ici de forme normale de
Smith à une matrice qui n’aurait pas été préalablement définie.

\subsection{\(2\) mots dodécaphasiques de provenances différentes}

La distinction doit être conservée
\begin{equation}
 U_{12}^{\rm cat}:=U_6\Vert U_6,
 \qquad
 U_{12}^{\rm sgn}:=\mathcal H_{12}(K).
 \label{p12-83-c1-s9:eq:u006d-dos-palabras-doce}
\end{equation}
La première concatène deux émissions du catalogue de longueur six. La seconde
est une coordonnée signée reliée réversiblement à \(K\) par l’opérateur
\(\mathcal H_{12}\). Leur longueur commune n’implique pas une même provenance,
sémantique ou loi de transport.

\begin{criterion}[Critères de réfutation]
\label{p12-83-c1-s9:crit:u006d-refutacion}
La construction est réfutée si la suite
\eqref{p12-83-c1-s9:eq:u006d-extension} n’est pas exacte, si sa classe cohomologique est
triviale, si un tuple de l’image viole
\eqref{p12-83-c1-s9:eq:u006d-condicion-divisibilidad} ou
\eqref{p12-83-c1-s9:eq:u006d-condicion-paridad}, si les formules inverses ne retrouvent pas
les douze entiers initiaux, ou si le retour de \(\beta\) est interprété comme
une réinitialisation de \(q_{\rm mem}\).
\end{criterion}


% Unidad U006B. Funtor ciclotomico y elevacion modular de caminos.

\section{Foncteur cyclotomique de phase--longueur et relèvement modulaire des chemins}
\label{p12-83-c1-s10:chap:funtor-ciclotomico-elevacion-caminos}

Le quotient odométrique \(\mathcal O_{9,12}\) conserve simultanément la
position nonadique, le registre dodécaphasique et la longueur entière de chaque
flèche. Dans cette unité, sa phase est réalisée sur des fibres \(A_2\) et l’on prouve
le relèvement cardinal qui remplace chaque arête par quatre arêtes
orientées. La représentation cyclotomique et la longueur absolue
resteront des coordonnées distinctes.

\subsection{Représentation de Coxeter d’ordre \(3\) sur les fibres dodécaphasiques de type \(A_2\)}

Soit
\begin{equation}
 G_{A_2}:=
 \begin{pmatrix}2&-1\\-1&2\end{pmatrix},
 \qquad
 C:=
 \begin{pmatrix}0&-1\\1&-1\end{pmatrix}.
 \label{p12-83-c1-s10:eq:u006b-dato-coxeter}
\end{equation}
Par le \cref{p12-83-c1-s5:prop:coxeter-a2-autonomo},
\begin{equation}
 C^3=I_2,
 \qquad
 C^{\mathsf T}G_{A_2}C=G_{A_2},
 \qquad
 \det(XI_2-C)=X^2+X+1.
 \label{p12-83-c1-s10:eq:u006b-identidades-coxeter}
\end{equation}
Ainsi, \(C\) réalise le caractère cyclotomique de \(C_3\) sur \(A_2\).

Pour chaque \(b\in\mathbb Z/12\mathbb Z\), soit \(A_{2,b}\) une copie de
\(A_2\), soit
\[
 \iota_b:A_{2,b}\xrightarrow{\;\sim\;}A_2
\]
une identification entière et soit
\(P_b\in O(G_{A_2})\) un repère qui dépend de \(b\), mais non du résidu
\(r\). Nous définissons
\begin{equation}
 \rho_b:=
 \iota_b^{-1}P_bCP_b^{-1}\iota_b,
 \qquad
 \rho_b^3=\operatorname{id}_{A_{2,b}}.
 \label{p12-83-c1-s10:eq:u006b-rho-fibra}
\end{equation}
L’indépendance par rapport à \(r\) lie le caractère à la fibre
dodécaphasique ; \(r\) conserve la position interne du tour nonadique.

\subsection{Catégorie odométrique du TPK : cylindres, troncatures et morphismes de succession}

Rappelons que \(\mathcal O_{9,12}\) a pour objets
\[
 (b,r)\in\mathbb Z/12\mathbb Z\times\{0,\ldots,8\}.
\]
Pour \(n\in\mathbb N\), l’unique flèche de longueur \(n\) est
\begin{equation}
 \gamma_{(b,r),n}:
 (b,r)\longrightarrow
 \left(
 b+\kappa_r(n),[r+n]_9
 \right),
 \qquad
 \kappa_r(n):=
 \left\lfloor\frac{r+n}{9}\right\rfloor.
 \label{p12-83-c1-s10:eq:u006b-flecha-odometro}
\end{equation}
La composition s’appuie sur
\begin{equation}
 \kappa_r(n+m)
 =\kappa_r(n)+\kappa_{[r+n]_9}(m).
 \label{p12-83-c1-s10:eq:u006b-cociclo-odometro}
\end{equation}
La longueur \(\ell(\gamma_{(b,r),n})=n\) est additive.

\begin{definition}[Foncteur cyclotomique de phase--longueur]
\label{p12-83-c1-s10:def:u006b-funtor-ciclotomico}
Soit \(\mathbf{Rep}_{\mathbb C}(C_3)\) la catégorie des représentations
complexes de \(C_3\), et soit \(\mathbf B\mathbb N\) la catégorie monoïdale
à un seul objet avec endomorphismes \((\mathbb N,+)\). Nous définissons
\begin{equation}
 \mathfrak F_{\rm cyc}:
 \mathcal O_{9,12}
 \longrightarrow
 \mathbf{Rep}_{\mathbb C}(C_3)\times\mathbf B\mathbb N
 \label{p12-83-c1-s10:eq:u006b-signatura-funtor}
\end{equation}
par
\[
 \mathfrak F_{\rm cyc}(b,r)
 =(A_{2,b}\otimes_{\mathbb Z}\mathbb C,\rho_b)
\]
et
\begin{equation}
 \mathfrak F_{\rm cyc}(\gamma)
 :=
 \left(
 \iota_{b'}^{-1}P_{b'}C^{\ell(\gamma)}
 P_b^{-1}\iota_b,
 \ell(\gamma)
 \right)
 \label{p12-83-c1-s10:eq:u006b-flecha-funtor}
\end{equation}
pour \(\gamma:(b,r)\to(b',r')\).
\end{definition}

\begin{theorem}[Fonctorialité et séparation des registres]
\label{p12-83-c1-s10:thm:u006b-functorialidad}
L’affectation de la \cref{p12-83-c1-s10:def:u006b-funtor-ciclotomico} est un foncteur. Sa
première composante conserve uniquement la classe de longueur modulo trois ; la
seconde conserve l’entier complet et, par réduction, détermine cette classe.
La composante cyclotomique ne permet pas de reconstruire la longueur absolue.
\end{theorem}

\begin{proof}
Si \(\gamma_1:(b_0,r_0)\to(b_1,r_1)\) et
\(\gamma_2:(b_1,r_1)\to(b_2,r_2)\) ont pour longueurs \(n_1,n_2\), la
trivialisation intermédiaire s’annule :
\begin{align}
 &\bigl(
 \iota_{b_2}^{-1}P_{b_2}C^{n_2}P_{b_1}^{-1}\iota_{b_1}
 \bigr)
 \bigl(
 \iota_{b_1}^{-1}P_{b_1}C^{n_1}P_{b_0}^{-1}\iota_{b_0}
 \bigr)\notag\\
 &\qquad=
 \iota_{b_2}^{-1}P_{b_2}C^{n_1+n_2}
 P_{b_0}^{-1}\iota_{b_0}.
 \label{p12-83-c1-s10:eq:u006b-composicion-funtor}
\end{align}
La seconde composante se compose par addition. Comme \(C^n\) ne dépend que de
\([n]_3\), mais que \(\ell=n\) conserve l’entier, les deux coordonnées ne sont pas
identifiées.
\end{proof}

Si
\[
 Q:\mathsf{TPK}\longrightarrow\mathcal O_{9,12}
\]
est le foncteur odométrique de U006, la réalisation typée est
\begin{equation}
 \mathfrak F_{\rm cyc}\circ Q:
 \mathsf{TPK}\longrightarrow
 \mathbf{Rep}_{\mathbb C}(C_3)\times\mathbf B\mathbb N.
 \label{p12-83-c1-s10:eq:u006b-compuesto-tpk-ciclotomico}
\end{equation}

\subsection{Relèvement modulaire des chemins cardinaux}

Soit \(\mathcal A_{N,d}\) la catégorie libre des chemins cardinaux de
\((\mathbb Z/N\mathbb Z)^d\). Soient
\[
 N\ge2,\quad u\ge1,\quad\gcd(u,N)=1,
 \quad k:=\operatorname{ord}_N(u),
 \quad q:=\frac{u^k-1}{N}.
\]

\begin{definition}[Substitution cardinale par \(u\)]
\label{p12-83-c1-s10:def:u006b-sustitucion-cardinal}
Le foncteur
\[
 \mathcal S_u:\mathcal A_{N,d}\longrightarrow\mathcal A_{N,d}
\]
envoie un objet \(p\) sur \(up\) et chaque arête orientée sur la concaténation
ordonnée de \(u\) copies de même direction.
\end{definition}

\begin{theorem}[Relèvement modulaire en endomorphisme de chemins]
\label{p12-83-c1-s10:thm:u006b-elevacion-modular}
Pour tout objet \(p\), chemin \(\gamma\) et
\[
 \kappa_{N,r}(n):=
 \left\lfloor\frac{r+n}{N}\right\rfloor,
\]
on a
\begin{align}
 \mathcal S_u^k(p)&=p,
 \label{p12-83-c1-s10:eq:u006b-retorno-objetos}\\
 \ell(\mathcal S_u^k\gamma)&=u^k\ell(\gamma),
 \label{p12-83-c1-s10:eq:u006b-longitud-sustituida}\\
 \kappa_{N,r}(u^kn)-\kappa_{N,r}(n)&=qn.
 \label{p12-83-c1-s10:eq:u006b-diferencia-vueltas}
\end{align}
\end{theorem}

\begin{proof}
La catégorie est libre, de sorte que les images des objets et des générateurs
déterminent un unique foncteur. Chaque application multiplie la longueur par
\(u\). Comme \(u^k=1+qN\), les objets reviennent modulo \(N\), et
\[
 \left\lfloor
 \frac{r+n+qNn}{N}
 \right\rfloor
 =
 \left\lfloor\frac{r+n}{N}\right\rfloor+qn.
\]
\end{proof}

Dans le cas nonadique,
\begin{equation}
 (N,u,k,q)=(9,4,3,7),
 \qquad
 4^3=64=1+7\cdot9.
 \label{p12-83-c1-s10:eq:u006b-parametros-nonadicos}
\end{equation}
Par conséquent,
\begin{equation}
 \boxed{
 \mathcal S_4^3(e)
 \text{ conserve les extrémités et contient sept tours nonadiques
 orientés par arête initiale}.}
 \label{p12-83-c1-s10:eq:u006b-siete-vueltas}
\end{equation}
Les extrémités et la position résiduelle reviennent en trois applications ;
l'état holonomique intégral ne revient pas nécessairement. Sans preuve
d'inversibilité, \(\mathcal S_4^3\) est un endomorphisme vertical de chemins,
non une monodromie.

\begin{remark}[Séparation d’avec la dilatation de l’horloge]
\label{p12-83-c1-s10:rem:u006b-s4-no-d4}
\(\mathcal S_4\) substitue les objets et les arêtes dans une catégorie de chemins.
L’application \(D_4\) de U006C multipliera par quatre la coordonnée
\(\Theta_{108}\). Elles partagent un facteur entier, mais ni le domaine, ni le codomaine, ni la
loi de composition.
\end{remark}

\begin{criterion}[Critères de réfutation]
\label{p12-83-c1-s10:crit:u006b-refutacion}
Le bloc est réfuté si l’une quelconque des identités
\eqref{p12-83-c1-s10:eq:u006b-identidades-coxeter},
\eqref{p12-83-c1-s10:eq:u006b-cociclo-odometro},
\eqref{p12-83-c1-s10:eq:u006b-composicion-funtor} ou
\eqref{p12-83-c1-s10:eq:u006b-parametros-nonadicos} échoue ; il est également mal typé si l’on
identifie la phase \(C^n\), la longueur \(n\), la substitution
\(\mathcal S_4\) et la dilatation \(D_4\).
\end{criterion}
}%


% Unidad U004. Redacción autónoma desde los generadors TRIT de 1063/live.

\section{TRIT comme unité minimale d’information topologique : orientation ternaire, mémoire de retenue et ordre temporel}
\label{p12-83-c1-s11:chap:trit-levantado-regimenes}

La courbure mixte de l’APP a produit trois valeurs orientées : zéro dans les feuillets
purs et des signes opposés lorsque l’ordre somme--produit est interverti. Le TRIT
type cette triade et conserve, outre le résidu visible, le quotient nécessaire
pour reconstruire l’entier relevé. Cette unité sépare expressément quatre
niveaux : classe ternaire, représentant équilibré, retenue et régime
quadratique.

\begin{definition}[Unité d’information topologique TRIT]
Le \emph{TRIT} est l’unité élémentaire orientée dont l’espace des états est
\(\mathsf{TRIT}=\{-1,0,+1\}\). Sa capacité logarithmique native est
\(\log 3\) et son expression en unités binaires est \(\log_2 3\). À la différence
d’un bit, ses trois états ne se réduisent pas à un encodage cardinal : ils typent,
respectivement, l’ouverture, le seuil et le retour, et conservent le quotient de
retenue requis pour reconstruire la transition qui a produit le résidu
visible.
\end{definition}

\subsection{Système ternaire orienté et représentation équilibrée des états TRIT}

Nous définissons l’ensemble orienté
\[
 \mathsf{TRIT}:=\{-1,0,+1\}.
\]
La carte équilibrée est la bijection d’ensembles
\[
 \operatorname{bal}:\mathbb F_3\longrightarrow\mathsf{TRIT},
 \qquad
 [0]_3\mapsto0,\qquad[1]_3\mapsto+1,\qquad[2]_3\mapsto-1.
\]
Il ne s’agit pas d’une égalité de types. \(\mathbb F_3\) porte la somme et le produit
d’un corps ; \(\mathsf{TRIT}\) porte l’ordre et l’orientation. Toute opération
transportée entre les deux doit indiquer la carte utilisée.

L’orientation des neuf phases est
\[
 \varepsilon_9(r)=
 \begin{cases}
 +1,&r\in\{1,4,7\},\\
 -1,&r\in\{2,5,8\},\\
 0,&r\in\{3,6,9\}.
 \end{cases}
\]
Sa table complète est
\[
\begin{array}{c|ccccccccc}
r&1&2&3&4&5&6&7&8&9\\\hline
\varepsilon_9(r)&+1&-1&0&+1&-1&0&+1&-1&0
\end{array}.
\]
La périodicité de trois n’élimine pas la phase nonadique : les phases
\(1,4,7\), par exemple, ont la même orientation et des positions distinctes
dans le tour.

\subsection{Algorithme de division ternaire équilibrée et unicité de la paire quotient–résidu}

\begin{theorem}[Relèvement ternaire local]
\label{p12-83-c1-s11:thm:trit-levantamiento-local}
Pour chaque \(n\in\{-4,-3,\ldots,4\}\), il existe une unique paire
\((r,c)\in\mathsf{TRIT}^2\) telle que
\[
 n=r+3c.
\]
\end{theorem}

\begin{proof}
L’existence se vérifie dans la table
\[
\begin{array}{c|rrrrrrrrr}
n&-4&-3&-2&-1&0&1&2&3&4\\\hline
r&-1&0&+1&-1&0&+1&-1&0&+1\\
c&-1&-1&-1&0&0&0&+1&+1&+1
\end{array}.
\]
Si \(r+3c=r'+3c'\), alors \(r-r'=3(c'-c)\). Comme
\(r-r'\in\{-2,-1,0,1,2\}\), le seul multiple de trois possible est zéro ;
par conséquent, \(r=r'\) et \(c=c'\).
\end{proof}

La projection visible \((r,c)\mapsto r\) possède, au-dessus de la valeur zéro, trois
états locaux :
\[
 0^-:=(0,-1),
 \qquad
 0^0:=(0,0),
 \qquad
 0^+:=(0,+1).
\]
Ils publient le même résidu, mais reconstruisent \(-3,0,+3\). Cette fibre triple
est la forme minimale d'une distinction qui réapparaîtra dans l'état holonomique intégral
du Topological Phase Kernel.

\begin{corollary}[Itération équilibrée]
\label{p12-83-c1-s11:cor:trit-iteracion-balanceada}
Tout entier \(N\) admet un développement ternaire équilibré fini
\[
 N=r_0+3r_1+\cdots+3^mr_m,
 \qquad r_j\in\mathsf{TRIT}.
\]
Il s’obtient en itérant la division \(N_k=r_k+3N_{k+1}\), où \(r_k\) est le
représentant équilibré de \(N_k\bmod3\).
\end{corollary}

\begin{proof}
La division euclidienne modulo trois donne un résidu dans \(\{0,1,2\}\) ; remplacer
\(2\) par \(-1\) accroît le quotient d’une unité et produit
\(r_k\in\mathsf{TRIT}\). La valeur absolue du quotient diminue en dehors
d’un ensemble fini, de sorte que le processus se termine. L’unicité s’obtient
en comparant le premier chiffre non coïncident et en réduisant modulo trois.
\end{proof}

\subsection{Composition et cocycle de retenue}

Pour \(r,s\in\mathsf{TRIT}\), nous définissons \(r\oplus s\in\mathsf{TRIT}\) et
\(\kappa_3(r,s)\in\mathbb Z\) par l’unique égalité
\[
 r+s=(r\oplus s)+3\kappa_3(r,s).
\]
La table est
\[
\begin{array}{c|ccc}
\oplus&-1&0&+1\\\hline
-1&+1&-1&0\\
0&-1&0&+1\\
+1&0&+1&-1
\end{array},
\qquad
\begin{array}{c|ccc}
\kappa_3&-1&0&+1\\\hline
-1&-1&0&0\\
0&0&0&0\\
+1&0&0&+1
\end{array}.
\]

\begin{proposition}[Identité cocyclique ternaire]
\label{p12-83-c1-s11:prop:trit-cociclo}
Pour \(r,s,t\in\mathsf{TRIT}\),
\[
 \kappa_3(r,s)+\kappa_3(r\oplus s,t)
 =\kappa_3(s,t)+\kappa_3(r,s\oplus t).
\]
\end{proposition}

\begin{proof}
Les deux membres sont le quotient total qui apparaît lorsque l’on écrit
\(r+s+t\) comme un représentant équilibré plus un multiple de trois.
L’unicité du relèvement prouve l’égalité.
\end{proof}

Si \((r,c)\) et \((s,d)\) sont des états relevés, leur somme est
\[
 (r,c)\boxplus(s,d)
 :=\bigl(r\oplus s,\ c+d+\kappa_3(r,s)\bigr).
\]
La \cref{p12-83-c1-s11:prop:trit-cociclo} prouve que \(\boxplus\) est associative et que la
reconstruction \(\operatorname{rec}_3(r,c)=r+3c\) satisfait
\[
 \operatorname{rec}_3(x\boxplus y)
 =\operatorname{rec}_3(x)+\operatorname{rec}_3(y).
\]

\subsection{Involution d’orientation et ordre temporel}

Pour un mot tritique \(\tau=(\tau_1,\ldots,\tau_n)\), nous définissons
\[
 \mathcal C_{\rm rev}(\tau_1,\ldots,\tau_n)
 :=(-\tau_n,\ldots,-\tau_1).
\]

\begin{lemma}[Renversement orienté]
\label{p12-83-c1-s11:lem:trit-reversion}
L’application \(\mathcal C_{\rm rev}\) est une involution. Elle échange les
secteurs \(+1\) et \(-1\) et fixe le secteur \(0\).
\end{lemma}

\begin{proof}
L’appliquer deux fois inverse deux fois l’ordre et change deux fois chaque
signe. On retrouve ainsi le mot initial. Les autres affirmations sont
immédiates dans chaque composante.
\end{proof}

Lorsqu’un paramètre temporel orienté est fixé, la convention du TPK est
\[
 +1:\text{retour avec mémoire},
 \qquad
 0:\text{seuil présent},
 \qquad
 -1:\text{ouverture de prolongement}.
\]
Cette affectation n’ajoute pas un quatrième état temporel. La direction est conservée
dans le signe et la position exacte est conservée dans la phase nonadique.

\subsection{Famille universelle d’algèbres quadratiques}

Pour \(\tau\in\mathsf{TRIT}\), nous définissons
\[
 \mathbb A_\tau:=\mathbb R[J_\tau]/(J_\tau^2+\tau I).
\]
Les trois cas sont :
\[
\begin{array}{c|c|c|c}
\tau&J_\tau^2&\text{régime}&\exp(tJ_\tau)\\\hline
+1&-I&\text{elliptique}&\cos t\,I+\sin t\,J_{+1}\\
0&0&\text{parabolique}&I+tJ_0\\
-1&+I&\text{hyperbolique}&\cosh t\,I+\sinh t\,J_{-1}
\end{array}.
\]

\begin{theorem}[Exponentielle quadratique uniforme]
\label{p12-83-c1-s11:thm:trit-exponencial-uniforme}
Si \(J_\tau^2=-\tau I\), alors
\[
 \exp(tJ_\tau)
 =C_\tau(t)I+S_\tau(t)J_\tau,
\]
où
\[
\begin{array}{c|ccc}
\tau&+1&0&-1\\\hline
C_\tau(t)&\cos t&1&\cosh t\\
S_\tau(t)&\sin t&t&\sinh t
\end{array}.
\]
\end{theorem}

\begin{proof}
Nous séparons la série exponentielle en puissances paires et impaires. Les relations
\(J_{+1}^{2m}=(-1)^mI\), \(J_0^2=0\) et
\(J_{-1}^{2m}=I\) produisent, respectivement, les séries du cosinus et du sinus,
la troncature linéaire et les séries hyperboliques.
\end{proof}

Sur \(\mathbb F_3\) apparaissent les analogues
\[
 \mathbb F_9=\mathbb F_3[\iota]/(\iota^2+1),
 \qquad
 \mathbb D_3=\mathbb F_3[\epsilon]/(\epsilon^2),
 \qquad
 \mathbb S_3=\mathbb F_3[j]/(j^2-1).
\]
Le premier est un corps à neuf éléments ; le deuxième contient un nilpotent
non nul ; le troisième se décompose au moyen des idempotents
\((1\pm j)/2\) et est isomorphe à \(\mathbb F_3\times\mathbb F_3\). La
coïncidence de cardinalités n’identifie pas ces trois algèbres.

La source entière du régime elliptique est
\[
 \mathbb Q_{\mathbb Z}:=\mathbb Z[u]/(u^2+1).
\]
La réduction modulo trois envoie \(u\) sur \(\iota\) et produit
\(\mathbb F_9\). L’extension des scalaires à \(\mathbb R\) et une
représentation \(u\mapsto J_{+1}\) produisent le réalisateur réel elliptique.
Il n’existe donc pas d’application non typée \(\mathbb F_9\to\mathbb R\) : les deux
réalisations procèdent d’une source entière commune.

\subsection{Représentations matricielles du TRIT sélectionnées par le signe de \(J_\tau^2\)}

Soit
\[
 C=\begin{pmatrix}0&-1\\1&-1\end{pmatrix},
 \qquad C^3=I,
\]
et définissons
\[
 J_e:=\frac{2C+I}{\sqrt3}
 =\frac1{\sqrt3}\begin{pmatrix}1&-2\\2&-1\end{pmatrix}.
\]
Une multiplication directe donne
\[
 (2C+I)^2=-3I,
 \qquad J_e^2=-I.
\]
Le réalisateur parabolique est
\[
 J_p:=N=\begin{pmatrix}0&1\\0&0\end{pmatrix},
 \qquad N^2=0.
\]
Pour le régime hyperbolique, nous prenons
\[
 F_{\rm av}=\begin{pmatrix}0&1\\1&1\end{pmatrix},
 \qquad
 A:=F_{\rm av}^2=\begin{pmatrix}1&1\\1&2\end{pmatrix},
\]
et
\[
 J_h:=\frac{2A-3I}{\sqrt5}
 =\frac1{\sqrt5}\begin{pmatrix}-1&2\\2&1\end{pmatrix}.
\]
Alors
\[
 (2A-3I)^2=5I,
 \qquad J_h^2=I.
\]

La forme bilinéaire
\[
 G=\begin{pmatrix}2&1\\1&-2\end{pmatrix}
\]
satisfait
\[
 A^TGA=G.
\]
De plus,
\[
 \cosh(2\log\varphi)=\frac32,
 \qquad
 \sinh(2\log\varphi)=\frac{\sqrt5}{2},
\]
de sorte que
\[
 A=\exp(2\log\varphi\,J_h)
  =\cosh(2\log\varphi)I+\sinh(2\log\varphi)J_h.
\]
Cette identité sera retrouvée causalement dans la biographie de
\(\varphi\) ; ici, seule la représentation matricielle est vérifiée.

Les trois contrôles exponentiels sont
\[
 \exp(\pi J_e)+I=0,
 \qquad
 \exp(J_p)=I+J_p,
 \qquad
 \exp(2\log\varphi\,J_h)=F_{\rm av}^2.
\]
La présence de \(\pi\) et de \(\varphi\) dans ces identités de vérification
n’autorise pas à les utiliser comme constructeurs de ces constantes. Leurs
constructeurs apparaîtront ensuite à partir de la généalogie nonadique.

\subsection{Couplage local APP--TRIT}

Avec l’échange
\[
 R=\begin{pmatrix}0&1\\1&0\end{pmatrix}
\]
nous définissons, pour \(\tau\in\mathsf{TRIT}\),
\[
 B_\tau=(4+3\tau)I+(5-3\tau)R.
\]
Puisque \(P_\pm=(I\pm R)/2\),
\[
 B_\tau=9P_++(6\tau-1)P_-.
\]

\begin{theorem}[Récupération spectrale de l’orientation]
\label{p12-83-c1-s11:thm:app-trit-recuperacion-espectral}
Le spectre de \(B_\tau\) est
\[
 \operatorname{Spec}(B_\tau)=\{9,6\tau-1\}.
\]
La valeur propre différentielle \(\lambda_-\) détermine univoquement le TRIT :
\[
 \tau=\frac{\lambda_-+1}{6}.
\]
De plus, le bloc central de l’APP est \(B_{+1}\).
\end{theorem}

\begin{proof}
Les projecteurs \(P_+\) et \(P_-\) diagonalisent \(R\) avec les valeurs propres
\(+1\) et \(-1\). La formule spectrale s’obtient par substitution. La deuxième
identité isole \(\tau\). Pour \(\tau=+1\),
\(B_{+1}=7I+2R=\begin{psmallmatrix}7&2\\2&7\end{psmallmatrix}\), qui est
l’unique bloc de la \cref{p12-83-c1-s1:thm:app-bloque-central}.
\end{proof}

Par conséquent, le TRIT n’est pas ajouté comme étiquette rétrospective : il peut être retrouvé
à partir du mode différentiel d’une famille construite sur l’échange interne
de l’APP.

\subsection{Obstructions du TRIT : séparation résidu–représentant–retenue, involutivité et récupération spectrale de \(\tau\)}

\begin{criterion}[Critères de réfutation du TRIT]
L’unité est réfutée si :
\begin{enumerate}
 \item on identifie la classe résiduelle, le représentant équilibré et la retenue ;
 \item on affirme que \(0^-\), \(0^0\) et \(0^+\) sont le même état ;
 \item la composition \(\boxplus\) cesse de reconstruire la somme entière ;
 \item \(\mathcal C_{\rm rev}\) n’est pas involutive ;
 \item on identifie les trois algèbres parce qu’elles ont le même cardinal ;
 \item un réalisateur ne satisfait pas \(J_\tau^2=-\tau I\) ;
 \item la valeur de \(\tau\) est choisie après observation du résultat et n’est pas
       retrouvée à partir du spectre de \(B_\tau\).
\end{enumerate}
\end{criterion}

L’APP fournit le support, la matrice multiplicative et la loi additive
interne ; TRIT fournit l’orientation,
le relèvement et le type quadratique. L’unité suivante construira le
Topological Phase Kernel comme catégorie de transport de ces données, sans
confondre le sélecteur interne avec l’état complet.

\HMTDeferredTPKBase
\HMTDeferredTPKDevelopment


% Unidad U006C. Dilatación aritmética, cociclo de acarreo y
% semiconjugación cuaternaria.

\section{Dilatation quaternaire du registre, obstruction observable et semi-conjugaison raffinée}
\label{p12-83-c1-s12:chap:dilatacion-cuaternaria-semiconjugacion}

La substitution cardinale \(\mathcal S_4\) construite dans l’unité précédente
agit sur les chemins. Dans cette unité, on définit une deuxième opération, de
nature arithmétique, sur le registre
\(\mathbb Z/12\mathbb Z\times\mathbb Z/9\mathbb Z\), et une troisième
opération, catégorielle, qui subdivise chaque transition observable en quatre
microtransitions. L’objectif n’est pas de les fusionner, mais de démontrer la relation
exacte entre leurs réalisations.

\subsection{Dilatation sur le registre produit}

Soit
\[
 \mathcal Q_{12,9}:=
 \mathbb Z/12\mathbb Z\times\mathbb Z/9\mathbb Z.
\]
Pour \(r\in\mathbb Z/9\mathbb Z\), nous notons
\(\widetilde r\in\{0,\ldots,8\}\) son représentant normalisé.

\begin{definition}[Dilatation quaternaire du registre]
\label{p12-83-c1-s12:def:u006c-dilatacion-registro}
Nous définissons
\begin{equation}
 D_4(b,r):=
 \left(
 \left[
  4b+\left\lfloor\frac{4\widetilde r}{9}\right\rfloor
 \right]_{12},
 [4r]_9
 \right).
 \label{p12-83-c1-s12:eq:u006c-d4-producto}
\end{equation}
La coordonnée totale est
\begin{equation}
 \Theta:\mathcal Q_{12,9}\xrightarrow{\;\sim\;}
 \mathbb Z/108\mathbb Z,
 \qquad
 \Theta(b,r):=[9b+\widetilde r]_{108}.
 \label{p12-83-c1-s12:eq:u006c-coordenada-total}
\end{equation}
\end{definition}

\begin{proposition}[Réalisation par multiplication modulaire]
\label{p12-83-c1-s12:prop:u006c-d4-total}
Dans la coordonnée totale,
\begin{equation}
 \boxed{
 \Theta\circ D_4\circ\Theta^{-1}(\theta)=[4\theta]_{108}.}
 \label{p12-83-c1-s12:eq:u006c-d4-total}
\end{equation}
En particulier,
\begin{equation}
 \operatorname{im}(D_4)
 =4\mathbb Z/108\mathbb Z
 \simeq\mathbb Z/27\mathbb Z,
 \label{p12-83-c1-s12:eq:u006c-imagen-d4}
\end{equation}
et la restriction induite sur l’image est d’ordre neuf.
\end{proposition}

\begin{proof}
Pour \(\theta=9b+\widetilde r\),
\[
 4\theta
 =9\left(
  4b+\left\lfloor\frac{4\widetilde r}{9}\right\rfloor
 \right)
 +(4\widetilde r\bmod9).
\]
C’est exactement la décomposition quotient--résidu de
\eqref{p12-83-c1-s12:eq:u006c-d4-producto}. L’image est formée des classes
multiples de quatre. Après division par quatre, on obtient
\(\mathbb Z/27\mathbb Z\), et
\[
\begin{aligned}
 4^1&\equiv4,& 4^2&\equiv16,& 4^3&\equiv10,\\
 4^4&\equiv13,& 4^5&\equiv25,& 4^6&\equiv19,\\
 4^7&\equiv22,& 4^8&\equiv7,& 4^9&\equiv1\pmod{27}.
\end{aligned}
\]
Aucune puissance précédente n’est égale à un ; par conséquent,
\(\operatorname{ord}_{27}(4)=9\).
\end{proof}

\subsection{Identité exacte de retenue}

Pour \(r\in\{0,\ldots,8\}\) et \(n\in\mathbb N\), nous définissons
\[
 \kappa_r(n):=\left\lfloor\frac{r+n}{9}\right\rfloor,
 \qquad
 r':=[r+n]_9.
\]

\begin{theorem}[Cocycle quaternaire de retenue]
\label{p12-83-c1-s12:thm:u006c-cociclo-d4}
On a
\begin{equation}
 \boxed{
 \kappa_{[4r]_9}(4n)
 =4\kappa_r(n)
 +\left\lfloor\frac{4\widetilde r'}{9}\right\rfloor
 -\left\lfloor\frac{4\widetilde r}{9}\right\rfloor.}
 \label{p12-83-c1-s12:eq:u006c-cociclo-d4}
\end{equation}
La dernière différence est le cobord produit par le choix des
représentants \(r\mapsto\widetilde r\).
\end{theorem}

\begin{proof}
Écrivons
\[
 r+n=9\kappa_r(n)+\widetilde r'.
\]
En multipliant par quatre,
\[
 4r+4n=36\kappa_r(n)+4\widetilde r'.
\]
Séparer le quotient et le résidu modulo neuf et soustraire le quotient déjà contenu dans
\(4\widetilde r\) produit
\eqref{p12-83-c1-s12:eq:u006c-cociclo-d4}.
\end{proof}

\begin{corollary}[Troisième itération]
\label{p12-83-c1-s12:cor:u006c-tercera-iteracion}
Pour tout \((b,r)\in\mathcal Q_{12,9}\),
\begin{equation}
 \boxed{D_4^3(b,r)=([4b+7\widetilde r]_{12},r).}
 \label{p12-83-c1-s12:eq:u006c-d4-cubo}
\end{equation}
\end{corollary}

\begin{proof}
Par la \cref{p12-83-c1-s12:prop:u006c-d4-total},
\[
 \Theta(D_4^3(b,r))
 =[4^3(9b+\widetilde r)]_{108}
 =[64(9b+\widetilde r)]_{108}.
\]
Comme \(64=1+7\cdot9\), le résidu modulo neuf redevient \(r\), tandis que
la coordonnée de bloc se transforme en
\([4b+7\widetilde r]_{12}\).
\end{proof}

Le retour de \(r\) n’est pas le retour du registre complet. La différence
\(7\widetilde r\) conserve la mémoire du représentant interne traversé.

\subsection{Obstruction dans le quotient observable}

Dans la configuration observable de
\cref{p12-83-c1-s3:def:tpk-configuracion-observable}, nous écrivons
\[
 q=(p^\Sigma,p^\Pi,d^\Sigma,d^\Pi,\theta)
 \in\mathcal Q_{\rm obs}.
\]
La projection positionnelle est
\begin{equation}
 \pi_{\rm pos}(q):=(p^\Sigma,p^\Pi)\in X_9^2.
 \label{p12-83-c1-s12:eq:u006c-proyeccion-posicional}
\end{equation}
La signature temporelle est
\[
 \tau(q):=\varepsilon_9(r(\theta))\in\{-1,0,+1\}.
\]
Nous définissons le déplacement observable orienté
\begin{equation}
 a(q):=
 \begin{cases}
  (\delta(d^\Sigma),0),&\tau(q)=+1,\\
  (0,\delta(d^\Pi)),&\tau(q)=-1,\\
  (0,0),&\tau(q)=0.
 \end{cases}
 \label{p12-83-c1-s12:eq:u006c-desplazamiento-observable}
\end{equation}
Alors
\[
 \pi_{\rm pos}(F_{\rm obs}q)=\pi_{\rm pos}(q)+a(q),
\]
avec l’actualisation des directions et de l’horloge définie dans U005.

\begin{theorem}[Obstruction à l’endofoncteur observable ordinaire]
\label{p12-83-c1-s12:thm:u006c-obstruccion}
Il n’existe pas d’endofoncteur de l’espace observable ordinaire qui satisfasse
simultanément les conditions suivantes :
\begin{enumerate}
 \item multiplier la position géométrique par quatre ;
 \item remplacer chaque transition par quatre microtransitions véritables ;
 \item conserver le feuillet et l’orientation à chaque microtransition ;
 \item envoyer une fibre dodécaphasique complète vers une seule fibre
 dodécaphasique.
\end{enumerate}
L’obstruction persiste à un pas de seuil avec \(a(q)=0\).
\end{theorem}

\begin{proof}
Soit
\[
 \mathcal B_b:=\{(b,r):0\leq r\leq8\}.
\]
Remplacer une itération élémentaire par quatre force
\[
 h(\theta+1)=h(\theta)+4
\]
et, par récurrence,
\[
 h(\theta)=4\theta+c\pmod{108}.
\]
Prenons \(0\leq\widetilde c<108\) et écrivons
\(\widetilde c=9c_b+c_0\), avec \(0\leq c_0<9\). Si toute la fibre
\(\mathcal B_b\) était envoyée vers une seule fibre, l’indice de bloc de
destination devrait être indépendant de \(r\) ; il contient cependant
\begin{equation}
 4b+c_b+
 \left\lfloor\frac{4r+c_0}{9}\right\rfloor
 \pmod{12},
 \label{p12-83-c1-s12:eq:u006c-bloque-obstruccion}
\end{equation}
qui n’est pas constant en \(r\). De plus, l’opérateur de transition neutre
satisfait \(P_0^{\rm TPK}=I\), différent des trivialisations
cyclotomiques \(P_b\) ; le calendrier ordinaire \((+1,-1,0)\) fait que
quatre itérations élémentaires visitent les deux modes arithmétiques et le mode
neutre.
\end{proof}

\subsection{Raffinement minimal par microinstants}

\begin{definition}[État observable subdivisé]
\label{p12-83-c1-s12:def:u006c-estado-subdividido}
Soit
\begin{equation}
 \widetilde{\mathcal Q}_4
 :=\mathcal Q_{\rm obs}\times\{0,1,2,3\},
 \qquad
 \widetilde\pi(q,j)
 :=4\pi_{\rm pos}(q)+j\,a(q).
 \label{p12-83-c1-s12:eq:u006c-estado-refinado}
\end{equation}
On définit
\begin{equation}
 \widetilde F(q,j):=
 \begin{cases}
  (q,j+1),&0\leq j<3,\\
  (F_{\rm obs}q,0),&j=3,
 \end{cases}
 \qquad
 \iota_4(q):=(q,0).
 \label{p12-83-c1-s12:eq:u006c-dinamica-refinada}
\end{equation}
\end{definition}

\begin{theorem}[Semi-conjugaison à quatre pas]
\label{p12-83-c1-s12:thm:u006c-semiconjugacion}
L’inclusion \(\iota_4\) satisfait
\begin{equation}
 \boxed{
 \widetilde F^4\circ\iota_4
 =\iota_4\circ F_{\rm obs}.}
 \label{p12-83-c1-s12:eq:u006c-semiconjugacion}
\end{equation}
Les positions intermédiaires sont exactement
\[
 4\pi_{\rm pos}(q)+j\,a(q),
 \qquad j=0,1,2,3.
\]
\end{theorem}

\begin{proof}
À partir de \((q,0)\), les trois premières applications n’accroissent que l’indice
de microinstant. La quatrième applique \(F_{\rm obs}\) et ramène l’indice à
zéro. Cela donne
\[
 (q,0)\mapsto(q,1)\mapsto(q,2)\mapsto(q,3)
 \mapsto(F_{\rm obs}q,0).
\]
La formule des positions est la définition de
\(\widetilde\pi\). Le raffinement rend visible une subdivision déterminée ;
il n’ajoute pas de décision dynamique.
\end{proof}

\subsection{Catégorie subdivisée et relèvement cyclotomique}

Soit \(\operatorname{sd}_4\mathsf P_{\rm obs}\) la catégorie obtenue en
remplaçant chaque générateur
\(\gamma:q\to F_{\rm obs}q\) par la chaîne
\begin{equation}
 (q,0)\to(q,1)\to(q,2)\to(q,3)\to(F_{\rm obs}q,0).
 \label{p12-83-c1-s12:eq:u006c-cadena-subdivision}
\end{equation}
Le foncteur
\[
 \operatorname{sd}_4:
 \mathsf P_{\rm obs}\longrightarrow
 \operatorname{sd}_4\mathsf P_{\rm obs}
\]
envoie chaque générateur sur cette chaîne. Le foncteur de contraction
\[
 \operatorname{col}_4:
 \operatorname{sd}_4\mathsf P_{\rm obs}
 \longrightarrow\mathsf P_{\rm obs}
\]
recompose les quatre microflèches et satisfait
\begin{equation}
 \operatorname{col}_4\circ\operatorname{sd}_4
 =\operatorname{id}_{\mathsf P_{\rm obs}}.
 \label{p12-83-c1-s12:eq:u006c-colapso}
\end{equation}

Dans une chaîne qui reste dans la fibre \(b\), les trois premières
microflèches agissent par \(\rho_b\) ; la quatrième incorpore le transport
interfibre s’il existe un franchissement dodécaphasique. Comme \(\rho_b^3=I\), le produit
ordonné des quatre actions coïncide avec le morphisme cyclotomique de la
flèche originale.

\begin{proposition}[Invariance cyclotomique par subdivision]
\label{p12-83-c1-s12:prop:u006c-invariancia-ciclotomica}
Il existe un relèvement typé
\[
 \widetilde{\mathfrak F}_{\rm cyc}:
 \operatorname{sd}_4\mathsf P_{\rm obs}
 \longrightarrow
 \mathbf{Rep}_{\mathbb C}(C_3)\times\mathbf B\mathbb N
\]
tel que, pour toute flèche \(\gamma\) de
\(\mathsf P_{\rm obs}\),
\begin{equation}
 \boxed{
 \widetilde{\mathfrak F}_{\rm cyc}
 \bigl(\operatorname{sd}_4\gamma\bigr)
 =
 \bigl(\mathfrak F_{\rm cyc}\circ Q\bigr)(\gamma).}
 \label{p12-83-c1-s12:eq:u006c-invariancia-ciclotomica}
\end{equation}
\end{proposition}

\begin{proof}
On attribue aux microflèches internes les puissances successives de
\(\rho_b\), et à la quatrième le transport entre les repères des fibres
d’origine et de destination. La composition des trois premiers facteurs est
l’identité parce que \(\rho_b^3=I\). Le facteur restant est exactement le
morphisme que \(\mathfrak F_{\rm cyc}\circ Q\) attribue à la flèche
contractée. La catégorie subdivisée est libre sur ces microflèches, de sorte que
l’affectation s’étend de manière unique.
\end{proof}

\subsection{\(3\) opérations quaternaires non interchangeables}

L’architecture contient trois opérations distinctes :
\begin{enumerate}
 \item \(\mathcal S_4\), substitution cardinale d’objets et de chemins ;
 \item \(D_4\), multiplication par quatre dans
 \(\Theta_{108}\) ;
 \item \(\operatorname{sd}_4\), subdivision fonctorielle d’une transition
 observable.
\end{enumerate}
Leurs compatibilités sont données, respectivement, par
\[
 4^3=1+7\cdot9,\qquad
 \Theta D_4=4\Theta,\qquad
 \widetilde F^4\iota_4=\iota_4F_{\rm obs}.
\]
Partager l’entier quatre ne transforme aucune d’elles en une autre.

\begin{criterion}[Critères de réfutation]
\label{p12-83-c1-s12:crit:u006c-refutacion}
La construction est réfutée si l’une quelconque des identités
\eqref{p12-83-c1-s12:eq:u006c-d4-total},
\eqref{p12-83-c1-s12:eq:u006c-cociclo-d4},
\eqref{p12-83-c1-s12:eq:u006c-d4-cubo},
\eqref{p12-83-c1-s12:eq:u006c-semiconjugacion} ou
\eqref{p12-83-c1-s12:eq:u006c-invariancia-ciclotomica} échoue ; également si une réalisation
observable satisfait les quatre conditions du
\cref{p12-83-c1-s12:thm:u006c-obstruccion} sans conserver un registre équivalent au
microinstant.
\end{criterion}


% Unidad U006. Lectores, odómetro y calendario; redacción autónoma.

\section{Lecteurs modulaires, odomètre et calendrier dodécaphasique}
\label{p12-83-c1-s13:chap:tpk-lectores-odometro-calendario}

L'état holonomique intégral du TPK ne s'identifie à aucune de ses lectures.
Cette unité définit les lecteurs modulo neuf, modulo trois et modulo mille,
l'odomètre phase--secteur, le quotient prédictif de trente-six états et
l'extension cyclique de cent huit positions. Chaque construction déclare ce qu'elle
publie et quelle mémoire elle conserve.

\subsection{Schéma typé des lecteurs modulaires du TPK : domaine, codomaine, publication et mémoire conservée}

\begin{definition}[Fiche opératoire]
\label{p12-83-c1-s13:def:tpk-ficha-operatoria}
Un opérateur canonique du TPK est enregistré au moyen d’un sextuplet
\[
 \mathcal O=(D_{\mathcal O},C_{\mathcal O},f_{\mathcal O},
 I_{\mathcal O},L_{\mathcal O},M_{\mathcal O}),
\]
où \(D_{\mathcal O}\) est le domaine, \(C_{\mathcal O}\) le codomaine,
\(f_{\mathcal O}\) la règle, \(I_{\mathcal O}\) les invariants préservés,
\(L_{\mathcal O}\) l'information publiée ou oubliée et
\(M_{\mathcal O}\) la mémoire qui demeure dans l'état holonomique intégral.
\end{definition}

Deux symboles de même cardinalité ou de même période ne représentent pas le même
opérateur, à moins qu’il existe une équivalence entrelaçant leurs domaines,
règles et invariants. Cette convention empêchera de confondre le quotient
prédictif \(C_{36}^{\rm pred}\), la frontière \(R_{36}\) et un défaut
cardinal de valeur trente-six.

\subsection{Lecteur nonadique de l'état holonomique intégral : phase, transition et retour avec mémoire}

Pour \(n\geq1\), nous définissons
\[
 \rho_9(n)=1+((n-1)\bmod9),
 \qquad
 q_9(n)=\left\lfloor\frac{n-1}{9}\right\rfloor.
\]
Alors
\[
 n=9q_9(n)+\rho_9(n).
\]
La fiche de \(\rho_9\) est
\[
 \bigl(
 \mathbb Z_{>0},\mathcal D_9,\rho_9,
 n\equiv\rho_9(n)\pmod9,
 \text{phase positive},
 q_9(n)
 \bigr).
\]
La racine numérique positive coïncide numériquement avec \(\rho_9\), mais sa
sémantique est distincte : l’une se présente comme section du quotient ; l’autre,
comme somme itérée de chiffres.

\subsection{Lecteur tritique de l’état orienté : ordre temporel et transport de retenue}

La division équilibrée s’écrit
\[
 n=r_3(n)+3c_3(n),
 \qquad r_3(n)\in\mathsf{TRIT},\quad c_3(n)\in\mathbb Z.
\]
Sa fiche est
\[
 \bigl(
 \mathbb Z,\mathsf{TRIT}\times\mathbb Z,
 n\mapsto(r_3(n),c_3(n)),
 \operatorname{rec}_3(r,c)=r+3c,
 r_3,
 c_3
 \bigr).
\]
La publication peut ne montrer que \(r_3\), mais le TPK conserve
\(c_3\) si une composition ultérieure exige une reconstruction.

\subsection{Lecteur positionnel par blocs en base \(10^3\)}

La division euclidienne définit
\[
 N=1000q_{1000}(N)+r_{1000}(N),
 \qquad 0\leq r_{1000}(N)<1000.
\]
Sa fiche est
\[
 \bigl(
 \mathbb Z,\mathbb Z\times\{0,\ldots,999\},
 N\mapsto(q_{1000},r_{1000}),
 \operatorname{rec}_{1000}(q,r)=1000q+r,
 r_{1000},
 q_{1000}
 \bigr).
\]

\begin{proposition}[Concaténation et lecture résiduelle]
\label{p12-83-c1-s13:prop:tpk-concatenacion-base-mil}
Si \(u_1,\ldots,u_m\in\{0,\ldots,999\}\) et
\[
 D_m=1000D_{m-1}+u_m,
 \qquad D_0=0,
\]
alors
\[
 D_m=\sum_{j=1}^m u_j1000^{m-j}.
\]
De plus,
\[
 D_m\equiv\sum_{j=1}^m u_j\pmod9
\]
parce que \(1000\equiv1\pmod9\).
\end{proposition}

\begin{proof}
La première égalité s’obtient par récurrence sur \(m\). La seconde procède
de sa réduction modulo neuf.
\end{proof}

La congruence ne reconstruit pas l’ordre des blocs. Le registre généalogique conserve la
séquence \((u_1,\ldots,u_m)\) ; le lecteur résiduel ne publie que leur somme
modulo neuf.

\subsection{Catégorie odométrique du TPK : objets cylindriques et composition par troncature}

\begin{definition}[Odomètre nonadique--dodécaphasique]
\label{p12-83-c1-s13:def:tpk-odometro-9-12}
La catégorie \(\mathcal O_{9,12}\) a pour objets
\[
 (b,r)\in C_{12}\times\{0,\ldots,8\}.
\]
Une flèche de longueur \(n\in\mathbb N\) est
\[
 (b,r)\xrightarrow{\ n\ }
 \left(
  b+\left\lfloor\frac{r+n}{9}\right\rfloor,
  r+n\bmod9
 \right).
\]
La composition additionne les longueurs.
\end{definition}

Soit
\[
 \kappa_r(n)=\left\lfloor\frac{r+n}{9}\right\rfloor.
\]

\begin{lemma}[Cocycle de l’odomètre]
\label{p12-83-c1-s13:lem:tpk-cociclo-odometro}
Si \(r_1=r+n_1\bmod9\), alors
\[
 \kappa_r(n_1+n_2)
 =\kappa_r(n_1)+\kappa_{r_1}(n_2).
\]
\end{lemma}

\begin{proof}
Nous écrivons \(r+n_1=9\kappa_r(n_1)+r_1\), avec
\(0\leq r_1<9\). En ajoutant \(n_2\) et en divisant de nouveau par neuf,
on obtient l’identité.
\end{proof}

Pour un état \(x\) dont la base contient \(\theta\), nous définissons
\[
 Q_0(x)=(\beta(\theta),\widetilde\theta\bmod9).
\]
Une flèche TPK a une longueur égale au nombre de générateurs observables qui
la composent.

\begin{theorem}[Foncteur phase--longueur]
\label{p12-83-c1-s13:thm:tpk-funtor-fase-longitud}
L’affectation qui envoie un état \(x\) sur \(Q_0(x)\) et une flèche de
longueur \(n\) sur la flèche odométrique de longueur \(n\) définit un foncteur
\[
 Q:\mathsf{TPK}\longrightarrow\mathcal O_{9,12}.
\]
\end{theorem}

\begin{proof}
Chaque générateur augmente \(\theta\) d’une unité. Après \(n\) pas, le
résidu de phase augmente de \(n\) modulo neuf et le secteur augmente de
\(\lfloor(r+n)/9\rfloor\). L’identité a une longueur nulle et le
\cref{p12-83-c1-s13:lem:tpk-cociclo-odometro} prouve la compatibilité avec la composition.
\end{proof}

Après neuf pas,
\[
 (b,r)\longmapsto(b+1,r).
\]
Après cent huit,
\[
 (b,r)\longmapsto(b,r).
\]
Ces identités appartiennent au quotient odométrique. Elles n’affirment pas que
le cylindre, la frontière, le registre généalogique ou la mémoire verticale soient revenus.

\subsection{Quotient minimal de mémoire prédictive}

Soit
\[
 X_{9,12}=\{(a,m):a\in\mathbb Z/9\mathbb Z, 0\leq m<12\}
\]
avec successeur
\[
 S(a,m)=
 \begin{cases}
  (a,m+1),&m<11,\\
  (a+1\bmod9,0),&m=11.
 \end{cases}
\]
Nous définissons les observables
\[
 \beta_{\rm pred}(a,m)=4a\pmod{12},
 \qquad
 d_{\rm pred}(a,m)=1+(m+\beta_{\rm pred}(a,m)\pmod{12}).
\]
Pour une observable \(q\), le mot futur inclusif d’horizon \(h\)
est
\[
 Q_h^q(x)=(q(x),q(Sx),\ldots,q(S^hx)).
\]

\begin{theorem}[Quotient prédictif minimal]
\label{p12-83-c1-s13:thm:tpk-cociente-predictivo-36}
Les observables \(\beta_{\rm pred}\), \(d_{\rm pred}\) et la paire
\((\beta_{\rm pred},d_{\rm pred})\) produisent le même quotient stable
\[
 \pi_{36}(a,m)=(a\bmod3,m).
\]
Il a trente-six états et des fibres de cardinal trois. Les horizons
minimaux sont \(11\), \(8\) et \(0\), respectivement.
\end{theorem}

\begin{proof}
Pour un système fini \((X,T)\), soit
\[
 E_h=\operatorname{Eq}(Q_h^q).
\]
Alors
\[
 E_{h+1}=E_h\cap(T\times T)^{-1}(E_h).
\]
La chaîne décroissante se stabilise et sa limite est la plus grande équivalence
\(T\)-compatible contenue dans \(\operatorname{Eq}(q)\). Dans le système
déclaré, l’énumération exacte produit le nombre de classes suivant :
\[
\begin{array}{c|rrrrrrrrrrrr}
h&0&1&2&3&4&5&6&7&8&9&10&11\\\hline
\#(X/E_h^{\beta})&3&6&9&12&15&18&21&24&27&30&33&36
\end{array},
\]
et
\[
\begin{array}{c|rrrrrrrrr}
h&0&1&2&3&4&5&6&7&8\\\hline
\#(X/E_h^{d})&12&15&18&21&24&27&30&33&36.
\end{array}
\]
La paire distingue déjà les trente-six classes à l’horizon zéro. Dans tous
les cas, les classes stables conservent \(m\) et \(a\bmod3\) ; les trois
possibilités de \(a\) ayant le même résidu modulo trois forment chaque fibre.
\end{proof}

Nous notons cet objet \(C_{36}^{\rm pred}\). Ce n’est pas la frontière
\(R_{36}\), qui contiendra trois canaux ordonnés de six trits, ni un recensement
de trente-six éléments d’une autre construction.

\subsection{Extension centrale du calendrier}

On définit
\[
 \iota:C_{12}\longrightarrow C_{108},
 \qquad \iota([b]_{12})=[9b]_{108},
\]
et
\[
 p:C_{108}\longrightarrow C_9,
 \qquad p([\theta]_{108})=[\theta]_9.
\]
On obtient la suite exacte
\begin{equation}
 0\longrightarrow C_{12}
 \xrightarrow{\ \iota\ }C_{108}
 \xrightarrow{\ p\ }C_9
 \longrightarrow0.
\label{p12-83-c1-s13:eq:tpk-extension-12-108-9}
\end{equation}
Une section ensembliste \(s([r]_9)=[r]_{108}\),
\(0\leq r<9\), induit le cocycle
\[
 c(r,r')=
 \left[\left\lfloor\frac{r+r'}{9}\right\rfloor\right]_{12}.
\]

\begin{proposition}[L’extension ne se scinde pas]
\label{p12-83-c1-s13:prop:tpk-extension-no-escindida}
La suite \cref{p12-83-c1-s13:eq:tpk-extension-12-108-9} n’est pas isomorphe à l’extension
produit \(C_{12}\times C_9\).
\end{proposition}

\begin{proof}
L’exposant de \(C_{108}\) est \(108\). L’exposant de
\(C_{12}\times C_9\) est
\(\operatorname{mcm}(12,9)=36\). Deux groupes finis isomorphes ont le
même exposant ; ils ne sont donc pas isomorphes.
\end{proof}

\begin{theorem}[Nécessité et période exacte de cent huit positions]
\label{p12-83-c1-s13:thm:tpk-periodo-exacto-108}
La translation \(\theta\mapsto\theta+1\) sur \(C_{108}\) est d’ordre
exactement \(108\). En conséquence, une itération \(n\) fait revenir le
calendrier complet si et seulement si \(108\mid n\).
\end{theorem}

\begin{proof}
La \(n\)-ième puissance envoie \([\theta]\) sur \([\theta+n]_{108}\). C’est
l’identité pour tout \(\theta\) exactement lorsque \([n]_{108}=0\), c’est-à-
dire lorsque \(108\mid n\).
\end{proof}

En particulier, \(1008\) n’est pas un autre retour du calendrier :
\[
 1008\equiv36\pmod{108}.
\]
Après \(1008\) pas, la phase modulo neuf revient, mais la coordonnée de
\(C_{108}\) s’est déplacée de trente-six positions. Naturellement,
\(216,324,\ldots\) font également revenir le calendrier, mais ce sont des répétitions du
tour fondamental de longueur \(108\).

\begingroup\fontsize{10.2}{11.7}\selectfont
\setlength{\parskip}{2pt}\setlength{\abovedisplayskip}{5pt}\setlength{\belowdisplayskip}{5pt}
\subsection{Mémoire verticale non bornée}

Soit \(q_{\rm mem}\in\mathbb Z_{\geq0}\) le nombre non borné de tours
nonadiques et soit
\[
 \beta=[q_{\rm mem}]_{12}.
\]
Alors un tour de neuf pas satisfait
\[
 q_{\rm mem}\longmapsto q_{\rm mem}+1,
 \qquad
 \beta\longmapsto\beta+1\pmod{12}.
\]
Après cent huit pas,
\[
 \beta\longmapsto\beta,
 \qquad
 q_{\rm mem}\longmapsto q_{\rm mem}+12.
\]
Le calendrier revient ; la mémoire verticale, non.

\begin{corollary}[Hiérarchie des retours]
\label{p12-83-c1-s13:cor:tpk-jerarquia-retornos}
Dans le quotient de phase,
\(9\) pas constituent un tour. Les morphismes composés
\[
 \mathcal K_{27}=F_{\rm obs}^{27},
 \qquad
 F_{\rm obs}^{54},
 \qquad
 F_{\rm obs}^{108}
\]
contiennent, respectivement, trois, six et douze tours nonadiques. Aucun
n’efface par définition le registre généalogique ou la frontière.
\end{corollary}

\subsection{Système minimal de types de l'état holonomique intégral et compatibilités entre ses lecteurs}

La table suivante répond explicitement à la question des objets internes qui
existent déjà dans le TPK avant l’introduction des germes :
\[
\begin{array}{p{.23\textwidth}p{.27\textwidth}p{.38\textwidth}}
\toprule
Objeto u operador&Dominio y codominio&Función\\\midrule
\(X_9\)&\((\mathbb Z/9)^2\)&soporte APP\\
\(T_N,T_E,T_S,T_O\)&\(X_9\to X_9\)&traslaciones cardinales\\
\(M_{\rm ph}\)&\(\mathcal D_9\to\mathsf{TRIT}\)&selección de hoja\\
\(F_{\rm obs}\)&\(\mathcal Q_{\rm obs}\to\mathcal Q_{\rm obs}\)&sucesor observable\\
\((\rho_9,q_9)\)&\(\mathbb Z_{>0}\to\mathcal D_9\times\mathbb Z_{\ge0}\)&fase y altura\\
\((r_3,c_3)\)&\(\mathbb Z\to\mathsf{TRIT}\times\mathbb Z\)&orientación y carry\\
\((q_{1000},r_{1000})\)&\(\mathbb Z\to\mathbb Z\times\{0,\ldots,999\}\)&bloque decimal y cociente\\
\(Q\)&\(\mathsf{TPK}\to\mathcal O_{9,12}\)&fase y longitud\\
\(C_{36}^{\rm pred}\)&cociente de \(X_{9,12}\)&memoria predictiva mínima\\
\(C_{108}\)&calendario cíclico&extensión no escindida\\
\(q_{\rm mem}\)&\(\mathbb Z_{\ge0}\)&memoria vertical no acotada\\
\(\operatorname{Ext}_r\)&
\(\operatorname{Ext}_r\subseteq\mathcal F_q\times\mathcal F_{q'}\)&
prolongaciones compatibles\\
\bottomrule
\end{array}
\]

\subsection{Obstructions des lecteurs et du calendrier du TPK : quotients euclidiens, \(C_{36}^{\mathrm{pred}}\), extension \(C_{108}\) et mémoire verticale}

\begin{criterion}[Critères de réfutation des lecteurs et du calendrier]
L’unité est réfutée si :
\begin{enumerate}
 \item un lecteur résiduel est présenté comme une transition complète ;
 \item on omet le quotient d’une division euclidienne utilisée ensuite ;
 \item on identifie \(C_{36}^{\rm pred}\) à \(R_{36}\) ;
 \item on affirme que l’extension \(C_{108}\) est le produit
       \(C_{12}\times C_9\) ;
 \item on déclare un retour complet pour un nombre non divisible par \(108\) ;
 \item le retour de \(\beta\) force \(q_{\rm mem}=0\) ;
 \item \(\mathcal K_{27}\) est présenté comme une horloge distincte et non comme
       une puissance du transport.
\end{enumerate}
\end{criterion}

\section{Structure opératorielle complète de l’APP et hiérarchie arithmétique--géométrique}

Le développement opératoriel de l’APP fixe la hiérarchie
\(4\to2\to6\to54\), les projecteurs arithmétiques, le bloc central et la
géométrie locale des chemins. Ces conséquences restent subordonnées au
noyau formel déjà construit et n’introduisent pas une seconde origine de l’APP.

L’unité suivante construira la subdivision quadruple, le foncteur
cyclotomique et les canaux internes de quatre-vingt-dix et cent vingt positions.
Ce n’est qu’ensuite que l’espace exhaustif des germes sera introduit.
 

\par\endgroup
\chapter{Architecture opératoire de la sélection, du transport et de l'actualisation}
\section{Structure catégorielle du TPK : états holonomiques intégraux, morphismes de
transport et classification opératoire}
\label{sec:arbol-operatorio-tpk-rev11}
\index{TPK!classification opératoire}
\index{sélecteur de phase!sous-opérateur du TPK}

Le TPK est le conteneur catégorique qui transporte sur APP les états
orientés de TRIT. Sa définition précède l'inventaire des réalisations :
on fixe d'abord les objets, les morphismes, la composition et les classes opératoires ; ce n'est
qu'ensuite que l'on introduit les opérateurs qui réalisent chaque classe.

\HMTClaim{HMT-R-TPK-ENRICHED-STATE}
\begin{definition}[TPK comme catégorie sur la base observable avec famille de fibres]
\label{def:estado-enriquecido-tpk-rev11}
Soit \(\mathsf P_{\rm obs}\) la catégorie des routes observables sur le graphe
de Cayley d'APP. Pour chaque configuration \(q\), on fixe la fibre
\[
 \mathcal F_q=
 \mathsf O_q\times\mathsf H_q\times\mathsf C_q\times
 \mathsf S_q\times\mathsf B_q\times\mathsf\Lambda_q,
\]
dont les facteurs conservent, respectivement, l'orientation et le feuillet ; le résidu et le
quotient ; la retenue ; la signature ; le cylindre et la frontière ; et le registre de transport. Un
objet du TPK est une paire \(x=(q,f)\), avec \(f\in\mathcal F_q\). Un morphisme
\(\Gamma:x\to y\) est une route \(\gamma:q\to q'\) de
\(\mathsf P_{\rm obs}\) accompagnée du transport de fibre qui préserve les
relations d'extension déclarées entre \(f\) et \(f'\). L'identité est la
route vide avec transport identité et la composition est la concaténation des
routes avec composition des transports de fibre. La projection
\[
 p:\mathcal X_{\rm TPK}^{\rm enr}\longrightarrow\mathsf P_{\rm obs}
\]
envoie chaque état sur sa configuration observable et chaque morphisme sur la route qu'il
transporte. Cette définition fixe une catégorie sur la base et une famille de
fibres ; elle ne présuppose pas l'existence de relèvements cartésiens pour toute
flèche de \(\mathsf P_{\rm obs}\).
\end{definition}

\begin{definition}[Classification ontologico--opératoire du TPK]
\label{def:categorias-operatorias-tpk-rev11}
La structure interne du TPK s'organise en huit classes différenciées par le type de
domaine, de codomaine et de fonction :
\begin{enumerate}
 \item \textbf{Support et état.} La base est \(\mathsf P_{\rm obs}\) et le
 porteur est \(\mathcal X_{\rm TPK}^{\rm enr}\to\mathsf P_{\rm obs}\). Cette
 classe contient des objets ; les sept suivantes contiennent des applications entre objets
 ou des lecteurs de ces objets.
 \item \textbf{Sélection interne.} Le sélecteur
 \[
 M_{\rm ph}:\{1,\ldots,9\}\longrightarrow\{+1,-1,0\}
 \]
 active l'évaluation additive, la multiplicative ou le passage de seuil. Il agit
 sur la coordonnée de feuillet et ne modifie pas à lui seul la position dans APP.
 \item \textbf{Transport.} Pour chaque direction cardinale \(d\),
 \(T_d:X_9\to X_9\), \(p\mapsto p+\delta(d)\), transporte la position.
 L'évolution \(F_{\rm obs}:\mathcal Q_{\rm obs}\to\mathcal Q_{\rm obs}\)
 compose sélection, translation, orientation et avancement de phase.
 \item \textbf{Lecteurs et quotients.} \(\rho_9\) publie la classe nonadique ;
 \((r_3,c_3)\) conserve le résidu tritique et le quotient ; et
 \((q_{1000},r_{1000})\) sépare le bloc décimal et la retenue. Ce sont des applications de
 lecture, non des transitions de l'état.
 \item \textbf{Mémoire et frontière.} Les transports de
 \(\mathsf H_q,\mathsf C_q,\mathsf S_q,\mathsf B_q\) et
 \(\mathsf\Lambda_q\) actualisent quotient, retenue, signature, cylindre, frontière
 et registre. Leur projection visible peut coïncider même si les états
 holonomiques intégraux diffèrent.
 \item \textbf{Périodicité et retour.} L'extension
 \[
 0\longrightarrow C_{12}\longrightarrow C_{108}
 \longrightarrow C_9\longrightarrow0
 \]
 type le secteur, le calendrier et la phase visible. Le retour nonadique restaure la
 phase ; la périodicité de cent huit itérations restaure le calendrier ;
 aucun n'efface
 par définition la mémoire de frontière. Le lacet
 \(\mathcal K_{27}=F_{\rm obs}^{27}\) est un morphisme composé, non une autre
 horloge.
 \item \textbf{Prolongation.} Les relations
 \(\operatorname{Ext}_r\subseteq\mathcal F_q\times\mathcal F_{q'}\)
 sélectionnent les continuations compatibles. Les cylindres et leur système inverse
 publient la survie finie et, sous les hypothèses de compatibilité
 déclarées, une section globale.
 \item \textbf{Sorties.} Les foncteurs d'évaluation contractent le même état
 holonomique intégral vers la sortie archimédienne, la sortie d'incidence
 exceptionnelle ou les classificateurs de l'atlas. Ces sorties sont des codomaines de
 lecture ; elles ne constituent pas de nouveaux conteneurs aux côtés du TPK.
\end{enumerate}
\end{definition}

Chaque opérateur qui apparaît ci-dessous est déterminé par le sextuplet
\((D,C,f,I,L,M)\) : domaine \(D\), codomaine \(C\), règle \(f\), invariants
préservés \(I\), information publiée ou perdue \(L\) et mémoire conservée
\(M\). L'égalité des cardinaux, des périodes ou des noms historiques n'identifie pas
deux opérateurs ayant des sextuplets distincts. Une fois cette
classification fixée, le diagramme suivant résume cinq regroupements fonctionnels
de réalisations qui raffinent les huit classes opératoires du TPK :
\[
\begin{tikzpicture}[
  node distance=7mm and 9mm,
  every node/.style={font=\small,align=center},
  root/.style={draw=hmtblue,very thick,rounded corners,fill=hmtlightblue,
    inner sep=5pt},
  op/.style={draw=hmtblue,rounded corners,fill=white,inner sep=4pt},
  sub/.style={draw=hmtgray,rounded corners,fill=hmtpaper,inner sep=3pt},
  edge/.style={-{Latex[length=2mm]},semithick,hmtblue}
]
\node[root] (tpk) {TPK\\$\mathcal X_{\rm TPK}^{\rm enr}$};
\node[op,below left=of tpk] (sel) {selección\\de lectura};
\node[op,below=of tpk] (trans) {transporte\\y retorno};
\node[op,below right=of tpk] (lift) {elevación\\y memoria};
\node[sub,below=of sel] (mov) {$M_{\rm ph}$};
\node[sub,below=of trans] (card) {$T_d,\ \mathcal K_{27},\ F_{\rm obs}$};
\node[sub,below=of lift] (fib) {$\mathsf H,\mathsf C,\mathsf S,
  \mathsf B,\mathsf\Lambda$};
\node[op,below=17mm of trans] (clock) {calendario y registro\\
  $C_{12}\hookrightarrow C_{108}\twoheadrightarrow C_9$, $\Theta_{108}$};
\node[op,below=of clock] (ext) {prolongación, cilindros, frontera\\
  y filtración de supervivencia};
\draw[edge] (tpk) -- (sel);
\draw[edge] (tpk) -- (trans);
\draw[edge] (tpk) -- (lift);
\draw[edge] (sel) -- (mov);
\draw[edge] (trans) -- (card);
\draw[edge] (lift) -- (fib);
\draw[edge] (mov) |- (clock);
\draw[edge] (card) -- (clock);
\draw[edge] (fib) |- (clock);
\draw[edge] (clock) -- (ext);
\end{tikzpicture}
\]
Le diagramme fixe une inclusion, non une succession de théories. En particulier,
le sélecteur interne, les translations cardinales, l'observation
stroboscopique et les lecteurs modulaires conservent les types fixés dans la
\cref{def:categorias-operatorias-tpk-rev11}. Le TPK coordonne leurs compositions
et conserve leur provenance dans un même état holonomique intégral.

\HMTClaim{HMT-R-TPK-OPERATOR-SEPARATION}
\begin{proposition}[Séparation de la sélection, du mouvement et de l'observation]
\label{prop:tpk-separacion-operadores-rev11}
Le sélecteur \(M_{\rm ph}\), le transport cardinal \(T_d\) et l'échantillonnage
\(\operatorname{Obs}_3\) sont des morphismes non interchangeables : le premier choisit
l'évaluation active ; le second modifie la position et l'enroulement ; le
troisième publie une sous-suite temporelle. Leur composition au sein du TPK
détermine un chemin observable, tandis que la fibre \(\mathcal F_q\) conserve
les données nécessaires pour distinguer des prolongations de même extrémité.
\end{proposition}

\begin{proof}
Les trois morphismes ont des codomaines différents : \(M_{\rm ph}\) prend ses valeurs
dans l'ensemble ternaire d'activation, \(T_d\) est un automorphisme du tore et
\(\operatorname{Obs}_3\) agit sur des suites. Une identification entre deux
d'entre eux violerait le typage du domaine ou du codomaine. L'extension par fibres
intègre en outre résidu--quotient, retenue, frontière et registre, de sorte
que la projection sur l'extrémité du chemin est une application d'oubli explicite.
\end{proof}

\begin{statusbox}[title={Nomenclature canonique}]
La dénomination publique du conteneur est \emph{TPK}. Lorsqu'il sera nécessaire
de nommer la donnée transportée, on emploiera \emph{état holonomique intégral du TPK} ou
\(\mathcal X_{\rm TPK}^{\rm enr}\). Le morphisme \(M_{\rm ph}\) désigne
exclusivement la sélection interne de phase.
\end{statusbox}
 \section{Opérateurs fondamentaux du TPK : domaines, actions, invariants et
mémoire}
\label{sec:inventario-operatorio-tpk-rev11}
\index{TPK!inventaire opératoire}

\cref{def:categorias-operatorias-tpk-rev11} a d'abord fixé les huit
classes ontologico--opératoires : support et état ; sélection ; transport ;
lecteurs et quotients ; mémoire et frontière ; périodicité et retour ;
prolongation ; et sorties. Cette section définit leurs réalisations. Pour chaque
opérateur, on déclare le sextuplet \((D,C,f,I,L,M)\) du domaine, du codomaine, de la règle,
des invariants, de l'information publiée ou perdue et de la mémoire conservée. Le recensement
nominal ultérieur regroupe des symboles déjà définis par ces six données et ne
remplace ni leurs domaines ni leurs règles.

\begin{definition}[Fiche typée d'un opérateur du TPK]
Un opérateur canonique est un sextuplet
\[
 \mathcal O=(D_{\mathcal O},C_{\mathcal O},f_{\mathcal O},
 I_{\mathcal O},L_{\mathcal O},M_{\mathcal O}),
\]
où \(f_{\mathcal O}:D_{\mathcal O}\to C_{\mathcal O}\) est la règle
d'action, \(I_{\mathcal O}\) est la famille de relations qu'elle conserve,
\(L_{\mathcal O}\) déclare les coordonnées publiées et la relation
d'équivalence induite par l'oubli, et \(M_{\mathcal O}\) énumère les données
qui demeurent dans l'état holonomique intégral après la publication. Deux
réalisations ne représentent le même opérateur que lorsqu'il existe une
équivalence typée qui entrelace leurs règles et leurs invariants ; partager
une période, un cardinal ou une sortie visible ne suffit pas.
\end{definition}

\subsection{Lecteurs modulaires construits}

La réduction nonadique est fixée par le représentant publié
\begin{equation}
 \rho_9(n)=1+((n-1)\bmod 9)
 =n-9\left\lfloor\frac{n-1}{9}\right\rfloor,
 \label{eq:rho9-construida-tpk-rev11}
\end{equation}
de sorte que la classe nulle s'écrit \(9\). Pour \(n\geq1\), la racine
numérique satisfait
\[
 \operatorname{dr}_9(n)=\rho_9(n),
\]
mais les deux expressions conservent des sémantiques distinctes : la première est une
section du quotient ; la seconde publie la somme itérée des chiffres. Le relèvement
tritique écrit de manière unique
\begin{equation}
 n=r_3(n)+3c_3(n),
 \qquad r_3(n)\in\{-1,0,+1\},\quad c_3(n)\in\mathbb Z,
 \label{eq:lift-mod3-tpk-rev11}
\end{equation}
et conserve simultanément le résidu orienté et le quotient de retenue.
Enfin, le lecteur de blocs décimaux est la division euclidienne
\begin{equation}
 N=1000q_{1000}(N)+r_{1000}(N),
 \qquad 0\leq r_{1000}(N)<1000,
 \label{eq:lector-mod1000-tpk-rev11}
\end{equation}
de sorte que la base \(1000\), le résidu modulo \(1000\) et la retenue \(q_{1000}\)
sont trois données liées et non un même objet.

\subsection{Classification des réalisations opératoires par domaine et fonction}

Les définitions précédentes régissent la lecture de la table. Ses quatorze
regroupements fonctionnels raffinent les huit classes opératoires et résument les
réalisations qui partagent une fonction ; une classe peut donc occuper plus
d'une ligne. La table ne sert pas de définition de remplacement. Les symboles qui
partagent un cardinal restent séparés par leur domaine et leur codomaine.

\begingroup
\footnotesize
\renewcommand{\arraystretch}{1.06}
\begin{longtable}{@{}p{.21\textwidth}p{.31\textwidth}p{.40\textwidth}@{}}
\toprule
Regroupement fonctionnel & Opérateurs et espaces & Domaine, action et invariants \\
\midrule
\endfirsthead
\toprule
Regroupement fonctionnel & Opérateurs et espaces & Domaine, action et invariants \\
\midrule
\endhead
État et mémoire de route &
\(\mathcal X_{\rm TPK}^{\rm enr}\), fibre \(\mathcal F_q\), mémoire de route &
Ils conservent position, feuillet, orientation, résidu, quotient, retenue, signature,
frontière, cylindre, préfixe, généalogie et mémoire. \\

Sélection de feuillet &
\(M_{\rm ph}:\{1,\ldots,9\}\to\{-1,0,+1\}\) &
Active l'évaluation additive, multiplicative ou le passage du seuil ; ne
déplace pas à lui seul l'état. \\

Transporte espacial &
\(T_N,T_E,T_S,T_O\), moteur réversible à deux curseurs,
\(F_{\rm obs}\) &
Réalise les translations du graphe de Cayley, met à jour le feuillet et l'orientation et
conserve chaque transition de la route. \\

Lectores modulares &
\(\rho_9\), \((r_3,c_3)\), \((q_{1000},r_{1000})\) &
Publient la phase nonadique, l'orientation tritique et le bloc décimal sans effacer les
quotients de relèvement. \\

Bloc et mémoire &
\(\mathcal K_{27}=F_{\rm obs}^{27}\), filtre de mémoire
\(\mathcal N_{27}\) &
Le premier compose vingt-sept transports ; le second sélectionne les
événements de mémoire. Un même cardinal n'implique pas un même opérateur. \\

État dodécaphasique &
\(\mathcal R_{12}\), centres \(\theta_m=30m-10\) degrés,
raffinement \(\mathcal R_{120}\) &
Organise les douze phases, leur opposition et le raffinement angulaire par incréments de
trois degrés. \\

Périodicité et profondeurs &
\(C_{12}\hookrightarrow C_{108}\twoheadrightarrow C_9\),
\(\Pi_{108}^{\rm cur}\), \(w_{108}\), \(w_{1080}\),
\(\Gamma_0(1080)\) &
Distingue la périodicité totale d'ordre \(108\), la projection des curseurs, les
mots de 108 et 1080 trits et le niveau modulaire 1080. \\

Canaux (90/120) &
masque \(C_m\), \(S_{90},S_{120}\), grammaire orbitale,
supertrace et invariants \(\nu_{120},\nu_{270}\) &
Sépare les deux généalogies, leurs queues géométriques, leurs poids et les
coordonnées qu'elles extraient de l'état. \\

Bidirectionnalité &
\(P_\pm\), inversion de route, \(\operatorname{Ext}_\pm\),
\(N_\pm\), \(\beta\), \(C_M=-\operatorname{rev}\) &
Réalise, respectivement, l'échange spectral, l'aller et le retour spatial,
la prolongation future/passée, la valence, le bilan et la conjugaison du mot. \\

Émission et coinduction &
\(L_0,L_1\), \(R_{36}\), \(G_9\), partition \(3^6=729\),
\(\varprojlim\operatorname{Cyl}_m\) &
Élève le germe, ouvre la frontière et énumère exhaustivement les
prolongations finies. La limite numérique et le relèvement enrichi
appartiennent à des types distincts. \\

Lecteurs numériques &
\(\mathscr C_\pi\), \(\mathscr P_e\), \(F_{\rm av}\) &
Construisent, par encadrements rationnels, récurrence factorielle et auto-échelle de
Fibonacci, les sections publiées de \(\pi,e,\varphi\). \\

Projection exceptionnelle &
relèvement de Witt, énumérateur de poids et ombre non sélectrice &
Transporte le même état vers l'incidence ternaire sans intervenir dans la
sélection archimédienne des chiffres. \\

Atlas et réalisation &
signature de route, atlas (81/56/13/324), involution \(C_M\), centre
électronique et opérateurs de fibre &
Transforme l'état de route en classificateurs, caractères et opérateurs avant
d'appliquer toute carte dimensionnelle. \\

Validation finie &
recensements exhaustifs, ablations et indépendance causale &
Contrôle la couverture et l'unicité à la profondeur calculée, et vérifie
qu'aucun préfixe cible n'intervient dans l'application génératrice. \\
\bottomrule
\end{longtable}
\endgroup

\Needspace{46\baselineskip}
\subsection{Index fonctionnel des opérateurs canoniques}

La classification précédente organise les réalisations par fonction. La table suivante
réunit les opérateurs après la définition catégorielle et les lecteurs
construits. Chaque ligne déclare le domaine, le codomaine, la règle et les invariants ; les
symboles permettent de reconnaître une même définition à travers ses
différentes réalisations sans identifier des opérateurs de types différents.

\begingroup
\footnotesize
\renewcommand{\arraystretch}{1.06}
\enlargethispage{2\baselineskip}
\begin{longtable}{@{}r >{\raggedright\arraybackslash}p{.42\textwidth}
  >{\raggedright\arraybackslash}p{.22\textwidth}
  >{\raggedright\arraybackslash}p{.22\textwidth}@{}}
\toprule
\# & Opérateur canonique & Symbole & Capa tipada \\
\midrule
\endfirsthead
\toprule
\# & Opérateur canonique & Symbole & Capa tipada \\
\midrule
\endhead
1 & Cellule arithmétique additive--multiplicative d'ordre neuf
  & \(\mathcal A_9\) & support arithmétique \\
2 & Décomposition résiduelle et quotient nonadiques
  & \((\rho_9,q_9)\) & réduction et élévation \\
3 & Réduction ternaire orientée
  & \(\rho_3^{\rm or}\) & réduction et élévation \\
4 & Lecteur positionnel ternaire--décimal en base mille
  & \(\iota_{1000}\) & réduction et élévation \\
5 & Action cardinale orientée sur le tore arithmétique
  & \(T_d,\ d\in\{N,E,S,O\}\) & transporte cardinal \\
6 & Rétroaction locale additive--multiplicative
  & \(F_{\rm loc}\) & transporte cardinal \\
7 & Émetteur minimal à deux curseurs
  & \(T_{\min}\) & émission finie \\
8 & Coupleur décimal de signatures arithmétiques
  & \(G\) & émission finie \\
9 & Projection ternaire du catalogue
  & \(\Psi_6\) & émission finie \\
10 & Surjectivité locale du transport cardinal
  & \(\mathcal R_3\) & transporte cardinal \\
11 & Lacet de transport de vingt-sept itérations
  & \(\mathcal K_{27}\) & transporte cardinal \\
12 & État holonomique intégral du noyau topologique de phase
  & \(\mathcal X_{\rm TPK}^{\rm enr}\) & état holonomique intégral \\
13 & Transition d'état de Hensel
  & \(\mathcal H\) & état holonomique intégral \\
14 & Transition exacte de cylindres ternaires--décimaux
  & \(\mathcal C_{3,1000}\) & état holonomique intégral \\
15 & Signature conjointe de frontière
  & \(\Sigma_{\rm B}\) & état holonomique intégral \\
16 & Relation d'extension et de survie nonadique
  & \({\rm EF}\text{--}G_9\) & prolongation et monodromie \\
17 & Monodromie nonadique avec holonomie de frontière
  & \(\mathcal M_9\) & prolongation et monodromie \\
18 & Section globale du système inverse des survivants
  & \(X_\chi=\varprojlim_m{\rm Surv}^{(\chi)}_m\)
  & prolongation et monodromie \\
19 & Catégorie bivariée de chemins arithmétiques
  & \(\mathcal P_{\rm APP}^{(2)}\) & transporte cardinal \\
20 & Chronologie observable à deux curseurs
  & \(F_{\rm obs}\) & transporte cardinal \\
21 & Alphabet décimal de quarante-deux triades
  & \(\mathcal A_{42}\) & émission finie \\
22 & Relation typée d'extension de fibres
  & \({\rm Ext}_r\) & état holonomique intégral \\
23 & Transfert nonadique de fibres compatibles
  & \(K_{\rm ph}\) & prolongation et monodromie \\
24 & Raffinement des germes par avancement et rotation
  & \((P,H,Q)\) & transporte cardinal \\
25 & Cocycles de mode et d'orientation
  & \((c_{\rm mode},c_{\rm or})\) & état holonomique intégral \\
26 & Prolongation triadique de vingt blocs
  & \(E_{20}\) & prolongation et monodromie \\
27 & Système dodécaphasique orienté
  & \(\mathcal R_{12}\) & lecture dodécaphasique \\
28 & Observable harmonique dodécaphasique signée
  & \(U_{12}^{\rm sgn}\) & lecture dodécaphasique \\
29 & Transformation harmonique par orbites de pas trois
  & \(\mathcal H_{12}\) & lecture dodécaphasique \\
30 & Extracteur transversal de différences cycliques
  & \(\mathcal E_{3,4}\) & lecture dodécaphasique \\
31 & Récurrence dodécaphasique de fermeture en base mille
  & \(\mathcal C_\alpha\) & fermeture en base mille \\
32 & Transducteur d'histoire enrichie en caractère angulaire
  & \((L_{\rm tick},\chi_{\rm mass})\) & interfaz dimensional \\
33 & Réseau bidirectionnel d'indice douze
  & \(\mathcal W_{\pm}\) & interfaz dimensional \\*
34 & Atlas nonadique des familles d'extension
  & \(\mathcal A_{\rm species}\) & interfaz dimensional \\
\bottomrule
\end{longtable}
\endgroup

Le symbole \(L_{\rm tick}\) désigne le lecteur d'itération et de retour. Son
codomaine est la mémoire temporelle de l'état ; il ne définit pas d'unité temporelle
physique indépendante.

Le sélecteur interne \(M_{\rm ph}\) appartient au TPK\@. Le transport cardinal
\(T_d\) occupe un autre facteur de la transition \(\mathcal U_t\) et réalise le
mouvement sur APP\@. Cette séparation conserve un unique conteneur,
distingue sélection de phase et transport, et préserve l'architecture complète
par types.

\Needspace{9\baselineskip}
\begin{resultbox}[title={Profondeur finie et limite enrichie}]
L'énumération, les recensements et la contraction des cylindres numériques sont des
résultats exacts dans leurs domaines. L'existence et l'unicité d'une section
enrichie soumise conjointement à tous les prédicats de survie
s'obtiennent lorsque les fibres de relèvement sont non vides, unitaires et
compatibles à toute profondeur. Un calcul sur une fenêtre finie
vérifie les fibres de cette fenêtre, mais n'implique pas à lui seul le
quantificateur sur toutes les profondeurs.
\end{resultbox}

\begin{theorem}[Composition de la dynamique interne]
\label{thm:composicion-dinamica-tpk-rev11}
Une transition observable du TPK est la composition typée
\[
 \mathcal U_t
 =\operatorname{Rec}_t\circ\operatorname{Mem}_t
 \circ T_{d(t)}\circ M_{\rm ph}(g(t))
\]
sur \(\mathcal X_{\rm TPK}^{\rm enr}\), suivie, le cas échéant, d'un
lecteur modulaire, d'une extension cylindrique ou de l'une des deux projections.
Chaque facteur agit sur une composante distincte de l'état ; ainsi, le
sélecteur, le mouvement, l'horloge, la mémoire et la publication ne peuvent pas
se réduire à une seule réduction modulaire.
\end{theorem}

La table précédente énumère les familles opératoires. Les équivalences entre
réalisations locales doivent entrelacer leurs domaines, leurs règles et leurs invariants.
Les développements ultérieurs sélectionnent des compositions de ces opérations ;
ils n'introduisent pas un quatrième objet primitif aux côtés d'APP, TRIT et TPK.
 \chapter{Transport, sélection et prolongement dans le noyau topologique de phase}
\label{sec:arquitectura-tpk}

Le sigle TPK désigne le \emph{Topological Phase Kernel}, ou noyau
topologique de phase.\footnote{La dénomination conserve également la
dédicace de l'auteur à Tan Peng Kian.} Il ne désigne pas un unique algorithme de 108
pas, un tableau de résultats ni le catalogue de 468 émissions. C'est le
système d'opérateurs qui fait coexister, sur la catégorie des chemins
d'APP, les lectures additive et multiplicative sans identifier leurs généalogies.

Le mot \emph{topologique} a ici un contenu précis. Le support
d'APP est le graphe de Cayley du groupe quotient
\[
 X_9=(\mathbb Z/9\mathbb Z)^2
\]
avec ses arêtes cardinales. Une voie fermée admet une classe d'enroulement
après relèvement au revêtement \(\mathbb Z^2\). TPK transporte cette classe,
l'orientation de la voie et l'holonomie de phase. Deux chemins qui se terminent
dans la même cellule peuvent donc représenter des états distincts. Le noyau
de phase est précisément l'information qui reste invisible lorsque l'on ne publie
que la cellule finale ou son reste.

La base dynamique est la catégorie libre \(\mathsf P_{\rm APP}\) du graphe
cardinal d'APP\@. Puisque les deux évaluations arithmétiques n'induisent pas la
même partition, le transport vit sur la catégorie bivariée
\[
 \mathsf P_{\rm APP}^{(2)}
 =\mathsf P_{\rm APP}\times\mathsf P_{\rm APP}.
\]
La première composante est évaluée au moyen du feuillet additif \(\Sigma\) et la
seconde au moyen du feuillet multiplicatif \(\Pi\). Cette séparation est
constitutive, mais ne duplique pas APP : elle enregistre deux lectures ordonnées d'une
même matrice de chemins. La composante multiplicative est le tableau principal ;
l'additive conserve la loi d'accumulation qui engendre ses lignes et permet de
mesurer l'information éliminée par une projection visible.

\ifdefined\HMTTPKRepeatAPP
\section{Carte canonique et architecture somme--produit}

Le support s'écrit en coordonnées de curseur
\[
 I_9=\{0,\ldots,8\},\qquad X_9=I_9^2,
\]
mais ses étiquettes arithmétiques appartiennent à \(\{1,\ldots,9\}\). La carte
qui relie les deux niveaux est
\[
 \chi:I_9\longrightarrow\{1,\ldots,9\},\qquad \chi(i)=i+1.
\]
Soit
\[
 \rho_9(n)=1+((n-1)\bmod9),\qquad
 q_9(n)=\frac{n-\rho_9(n)}9.
\]
En particulier, \(9\) est le représentant visible de la classe nulle. Les deux
observables d'un même sommet sont
\[
 \Sigma(i,j)=\rho_9\bigl(\chi(i)+\chi(j)\bigr),\qquad
 \Pi(i,j)=\rho_9\bigl(\chi(i)\chi(j)\bigr).
\]
Cette convention empêche de mélanger la coordonnée de curseur zéro avec l'étiquette
visible neuf. APP est le tuple
\[
 \mathcal A_9=
 \bigl(X_9,\operatorname{Cay}(X_9),\Sigma,\Pi,
       \operatorname{Path}(X_9)\bigr),
\]
et non la juxtaposition de deux tableaux : chaque germe \((i,j)\) porte à la fois
les deux évaluations et chaque voie conserve l'ordre dans lequel elle les consulte.

Le relèvement qui retient le reste et le quotient est
\[
 \mathcal A_9^\#=(R^+,Q^+;R^\times,Q^\times),
\]
avec, pour \(a=\chi(i)\) et \(b=\chi(j)\),
\[
 a+b=R^+_{ij}+9Q^+_{ij},\qquad
 ab=R^\times_{ij}+9Q^\times_{ij},
\]
\[
 R^+_{ij}=\Sigma(i,j),\qquad R^\times_{ij}=\Pi(i,j).
\]
Les totaux exacts sont
\[
 \sum R^+=405,\quad \sum R^\times=459,\quad
 \sum Q^+=45,\quad \sum Q^\times=174.
\]
Par conséquent,
\[
 2025-810=1215=(459-405)+9(174-45)=54+9\cdot129.
\]
L'égalité ne crée pas TPK par simple comptage : elle identifie la donnée que le
transport doit conserver. Le reste visible diffère de \(54\) et la mémoire
du quotient de \(129\).

La composition du feuillet additif possède le cocycle de retenue
\[
 \kappa_9(a,b)=
 \frac{\rho_9(a)+\rho_9(b)-\rho_9(a+b)}9,
\]
qui satisfait
\[
 \kappa_9(a,b)+\kappa_9(a+b,c)
 =\kappa_9(b,c)+\kappa_9(a,b+c).
\]
Le secteur multiplicatif possède en outre le radical
\[
 \mathfrak m=(3)=\{3,6,9\},\qquad \mathfrak m^2=0
 \quad\text{dans }\mathbb Z/9\mathbb Z.
\]
Les six lignes unitaires ont pour somme \(45\), les lignes \(3,6\) ont pour somme \(54\) et
la ligne \(9\) a pour somme \(81\) ; ainsi
\[
 459=6\cdot45+2\cdot54+81,
\]
et tout l'excès \(459-405=54\) se situe dans le secteur non unitaire. Cette
localisation, le cocycle de retenue et la différence des quotients sont trois
données distinctes de la même architecture somme--produit.
\fi

\section{Réductions résiduelles et représentation décimale}

L'architecture utilise trois réductions, chacune ayant un rôle distinct :
\[
 \mathbb Z\xrightarrow{(\rho_9,q_9)}
 \{1,\ldots,9\}\times\mathbb Z,
 \qquad
 \rho_3^{\rm or}:\mathbb Z/3\mathbb Z\longrightarrow
 \{-1,0,+1\},
\]
\[
 (d_1,d_2,\ldots)\in\{0,\ldots,999\}^{\mathbb N}
 \xrightarrow{\iota_{1000}}
 \sum_{m\geq1}d_m1000^{-m}.
\]
L'orientation ternaire est
\[
 \rho_3^{\rm or}([0])=0,\qquad
 \rho_3^{\rm or}([1])=+1,\qquad
 \rho_3^{\rm or}([2])=-1.
\]
La première sépare le reste visible et le quotient nonadique ; la deuxième attribue le
type orienté de la transition ; la troisième regroupe la lecture archimédienne
en triades décimales, dont le registre conserve la retenue. Aucune ne peut
remplacer une autre.
En particulier, la réduction modulo trois ne reconstruit pas le quotient modulo
neuf, et une triade décimale n'est pas un mot de six trits.
La somme accumule les déplacements ; le produit replie les classes et contracte les
lignes non inversibles. La nécessité de TPK naît du fait que les deux lectures
partagent le support, mais ne partagent ni la partition, ni la période, ni le quotient.

\Needspace{36\baselineskip}
\section{Correspondance entre symboles du TPK, domaines et compositions}
\label{sec:diccionario-interno-tpk}

Les noms historiques sont fixés ici par objet et non par ressemblance
numérique. Ce tableau fait partie de la définition publique de TPK.

\begin{longtable}{@{}>{\raggedright\arraybackslash}p{.19\textwidth}
>{\raggedright\arraybackslash}p{.30\textwidth}
>{\raggedright\arraybackslash}p{.43\textwidth}@{}}
\toprule
Nom&Objet&Fonction au sein du noyau\\\midrule
\endhead
coordonnée nonadique&
\(\rho_9(n)\) avec \(q_9(n)\)&
Sépare le représentant visible modulo \(9\) du quotient qui conserve la
mémoire du franchissement de frontière. Ce n'est pas une racine numérique.\\
sélecteur temporel ternaire&
\(M_{\rm ph}(t)=\varepsilon(t)\in\{+1,-1,0\}\)&
En \(1,4,7\), il active l'accumulation additive ; en \(2,5,8\), il active
l'accumulation multiplicative ; en \(3,6,9\), il maintient les deux curseurs et enregistre le seuil.
Il ne désigne pas les directions spatiales.\\
transporte cardinal&
\(T_N,T_E,T_S,T_O\)&
Déplace la position sur le tore APP et conserve le mot de voie et son
enroulement. Il est indépendant du sélecteur temporel.\\
lecteur ternaire&
\(\rho_3^{\rm or}\)&
Convertit les trois classes résiduelles en orientation
\(0,+1,-1\). Il ne reconstruit pas à lui seul le quotient nonadique.\\
lecteur en base mille&
\(\iota_{1000}\)&
Regroupe la publication archimédienne en triades décimales et transporte ses
retenues. C'est un lecteur positionnel, non une congruence ``modulo \(1000\)''.\\
retour de vingt-sept&
\(\mathcal K_{27}=F_{\rm obs}^{27}\)&
Compose vingt-sept mises à jour positionnelles et orientées. Son quatrième
itéré restaure la phase observable, tandis que le registre complet conserve
la mémoire accumulée.\\
registre dodécaphasique&
\(\Lambda\in\mathcal R_{12}\)&
Conserve six positions, cinq différences et un total. C'est le support du
transport réversible \(1+5+6=12\).\\
lecture bidirectionnelle&
\(q_+,q_-\) et ses poids conjugués&
Conserve simultanément les orientations de l'aller et du retour. En mécanique
dimensionnelle, une signature \((n_A,s_C)\) produit d'abord
\(W_\pm=6n_A\pm s_C\) ; ce n'est qu'ensuite qu'apparaît
\(s_C/6\) dans le caractère angulaire.\\
\bottomrule
\end{longtable}

Le sélecteur temporel, le transport cardinal et l'orientation ternaire sont
donc trois applications distinctes. Leur séparation conserve la différence entre
\emph{quand} une lecture est activée, \emph{vers où} la voie avance et
\emph{avec quelle orientation} le pas est enregistré. TPK contient les trois et
conserve leurs relations.

\ifdefined\HMTTPKWordArithmetic
\section{De la cellule locale au monoïde des mots nonadiques}

L'architecture somme--produit ne s'arrête pas au tableau local. Pour prouver que
le quotient transporté engendre une arithmétique globale, il faut construire
l'opération sur les mots sans introduire comme définition la multiplication
entière que l'on entend retrouver. Soit
\[
 \mathcal D_9=\{1,\ldots,9\},\qquad
 \mathcal W_9=\{\varnothing\}\sqcup
 \bigsqcup_{\ell\geq1}\mathcal D_9^\ell .
\]
Les mots s'écrivent à partir de la position la moins significative et s'évaluent
par
\[
 \operatorname{ev}_9(u_0,\ldots,u_{\ell-1})
 =\sum_{j=0}^{\ell-1}u_j9^j,\qquad
 \operatorname{ev}_9(\varnothing)=0.
\]
Le symbole \(9\) conserve le franchissement de frontière ; il n'est pas remplacé par un zéro.

\begin{theorem}[Forme normale nonadique]
L'évaluation \(\operatorname{ev}_9\) est une bijection entre
\(\mathcal W_9\) et \(\mathbb N_0\).
\end{theorem}
\begin{proof}
Les mots de longueur \(\ell\) couvrent exactement l'intervalle
\[
 \frac{9^\ell-1}{8}
 \leq\operatorname{ev}_9(u)
 \leq9\frac{9^\ell-1}{8}.
\]
Le minimum de l'intervalle suivant dépasse d'une unité le maximum du
précédent. Pour \(n>0\), la division de \(n-1\) par neuf détermine de manière
unique
\[
 d=1+((n-1)\bmod9)\in\mathcal D_9
\]
et un quotient strictement inférieur, auquel la même règle s'applique
récursivement. Le processus se termine et est réversible.
\end{proof}
Par exemple,
\[
 9\leftrightarrow(9),\quad
 10\leftrightarrow(1,1),\quad
 81\leftrightarrow(9,8),\quad
 1000\leftrightarrow(1,3,3,1).
\]
La dernière identité exhibe dans une seule forme normale les deux partitions
fondamentales
\[
 1000=729+271=729+243+27+1;
\]
leur interprétation par canaux est introduite ensuite, non dans la définition
de la numération.

\section{Somme, produit de Horner et conservation du quotient}

La somme native \(u\boxplus_\#v\) parcourt les mots par positions. À chaque
position, elle applique la cellule additive aux deux chiffres présents et, s'il existe
une retenue entrante, applique une seconde cellule au reste provisoire. Le
reste est publié et la somme des quotients passe à la position suivante.
L'invariant partiel est
\[
 \sum_{j<k}(u_j+v_j)9^j
 =\sum_{j<k}w_j9^j+c_k9^k,
\]
d'où l'on obtient
\[
 \operatorname{ev}_9(u\boxplus_\#v)
 =\operatorname{ev}_9(u)+\operatorname{ev}_9(v).
\]

Pour \(d\in\mathcal D_9\), le transducteur scalaire
\(d\odot_\#u\) utilise d'abord
\[
 u_jd=9q_j+r_j,\qquad
 (r_j,q_j)=\bigl(R^\times(u_j,d),Q^\times(u_j,d)\bigr),
\]
et normalise la collision entre \(r_j\) et le quotient entrant au moyen de la
cellule additive. Ainsi,
\[
 \operatorname{ev}_9(d\odot_\#u)
 =d\operatorname{ev}_9(u).
\]
Si \(v=(v_0,\ldots,v_{m-1})\), on définit, de droite à gauche,
\[
 a_m=\varnothing,\qquad
 a_j=(9\odot_\#a_{j+1})\boxplus_\#(v_j\odot_\#u),
\]
et
\[
 u\boxtimes_\#v:=a_0.
\]
C'est une construction de Horner réalisée uniquement avec les deux cellules
locales et leurs quotients.

Cette description fixe la dépendance opératoire, mais ne constitue pas un
second emplacement du résultat. Le produit global, sa preuve de Horner et
l'identité
\(operatorname{ev}_9(u\boxtimes_\#v)
=\operatorname{ev}_9(u)\operatorname{ev}_9(v)\)
sont établis une seule fois dans le
\cref{eq:ev-multiplicative}, au sein du chapitre spécialisé sur le produit
natif.
L'exemple de frontière le plus court est
\[
 (9)\boxtimes_\#(9)=(9,8),\qquad 9+8\cdot9=81.
\]
Ne conserver que \(R^\times\) détruirait ce théorème : \(2\cdot5=10\)
serait réduit au reste visible \(1\). Le quotient n'est pas une annotation
auxiliaire, mais la mémoire nécessaire pour que le produit survive au passage de
la frontière.

\section{Coproduit et irréductibilité}

L'ouverture multiplicative est définie de manière interne par
\[
 \Delta_\#e_w
 =\sum_{u\boxtimes_\#v=w}e_u\otimes e_v.
\]
Sa version réduite élimine les facteurs \((1)\). Si
\(\widetilde\Delta_\#\) désigne cette version et
\[
 A_\#=
 \bigl(\overline{\widetilde\Delta_\#}\bigr)^*
 \overline{\widetilde\Delta_\#},
\]
alors \(A_\#e_w=d_\#(w)e_w\), où \(d_\#(w)\) compte les paires ordonnées
non triviales qui produisent \(w\). Le sélecteur
\[
 P_{\mathrm{Irr},\#}
 =(I-P_{(1)})\,\mathbf1_{\{0\}}(A_\#)
\]
choisit exactement les mots distincts de l'unité sans ouverture native non
triviale. L'irréductibilité se lit donc dans la fibre même du produit ;
elle n'est pas introduite au moyen d'une table externe de nombres premiers.

Ce résultat est exact dans le monoïde des formes normales. Son extension aux
histoires complètes exige une compatibilité supplémentaire. Si \(P_n\) projette
sur des histoires survivantes de profondeur \(n\), une famille
d'injections \(\iota_n\) doit satisfaire
\[
 P_n\iota_n\mu_\#
 =P_n\mu_n(\iota_n\otimes\iota_n).
\]
Les théorèmes précédents construisent \(\mu_\#\) et fixent cette équation ; ils
n'affirment pas à eux seuls l'existence ou l'unicité de la famille
\((\iota_n)_n\). Cette séparation empêche de confondre l'arithmétique globale des
mots, déjà démontrée, avec son prolongement monoïdal à l'espace profond de
mémoire TPK\@.

\fi

\section{Action cardinale et rétroaction arithmétique}

Soient
\[
 \nu_N=(-1,0),\quad \nu_E=(0,1),\quad
 \nu_S=(1,0),\quad \nu_O=(0,-1)
\]
les quatre générateurs cardinaux. Pour \(d\in\{N,E,S,O\}\),
\[
 T_d(i,j)=(i,j)+\nu_d\pmod9.
\]
La lettre \(O\) signifie ouest. Ces directions sont des déplacements sur
APP et non des noms des canaux \(\pi,e,\varphi\).

Un mot \(w=d_1\cdots d_r\) détermine
\[
 \gamma(x_0;w)_k=x_0+\sum_{j=1}^{k}\nu_{d_j}\pmod9.
\]
Lorsqu'il se ferme, son relèvement conserve
\[
 \operatorname{wind}(w)=\frac1{9}
 \bigl(\#S-\#N,\#E-\#O\bigr)\in\mathbb Z^2.
\]
La lecture bidirectionnelle conserve le mot, l'orientation et le feuillet. Son
renversement formel inverse l'ordre et échange
\(N\leftrightarrow S\), \(E\leftrightarrow O\). Son éventuelle lecture comme
avance--retard est une interface physique ultérieure, non une partie de cette
définition de voie.

L'état local observable est
\[
 s_t=(i_t,j_t,m_t),
 \qquad m_t\in\{\Sigma,\Pi\}.
\]
On évalue le feuillet actif,
\[
 a_t=\rho_3^{\rm or}\bigl(m_t(i_t,j_t)\bigr),
\]
et le mode est mis à jour au moyen de
\[
 m_{t+1}=
 \begin{cases}
 \Sigma,&a_t=+1,\\
 \Pi,&a_t=-1,\\
 m_t,&a_t=0.
 \end{cases}
\]
Ainsi, chaque pas cardinal modifie à la fois la position et la lecture
arithmétique future. La voie ne se superpose pas à un mot déjà construit : c'est
l'histoire positionnelle qui le produit.

\begin{theorem}[Surjectivité du facteur échantillonné de pas trois]
Pour chacun des \(162=9^2\cdot2\) états \((i,j,m)\) et chaque trit
orienté cible \(b\in\{-1,0,+1\}\), il existe une voie
\(u\in\{N,E,S,O\}^3\) dont le troisième pas émet \(b\). Dans les \(486\) cas
état--cible, le nombre de voies admissibles est compris entre
\(8\) et \(40\). Par conséquent, pour toute suite cible
\((b_k)_{k\geq1}\), il existe une histoire cardinale dont la sortie satisfait
\[
 (a_3,a_6,a_9,\ldots)=(b_1,b_2,b_3,\ldots).
\]
Les deux trits intermédiaires \(a_{3k-2}\) et \(a_{3k-1}\) de chaque bloc ne
sont pas prescrits par cet énoncé.
\end{theorem}
\begin{proof}
On énumère les \(4^3=64\) voies de trois pas à partir de chaque état. Le
certificat exhaustif inclus enregistre, pour chacun des
\(162\cdot3\) cas, la voie exécutée, l'état final et le nombre total de
témoins ; leurs indices \(0,1,2\) sont transportés par
\(\rho_3^{\rm or}\) vers \(0,+1,-1\). Tous les décomptes sont positifs ; le
minimum est \(8\) et le maximum est \(40\). Après chaque bloc, le même
résultat s'applique à l'état atteint.
Le choix, par exemple, du premier témoin dans l'ordre lexicographique
produit récursivement une histoire pour toute suite cible.
\end{proof}

La surjectivité concerne le facteur observé tous les trois pas, non le
mot émis au temps unitaire. Ce théorème établit cette capacité
échantillonnée du support des voies. La sélection des sections
\(\pi,e,\varphi\) appartient à la relation enrichie EF--G9 définie
plus loin : la capacité de réalisation et la sélection généalogique sont deux
opérateurs distincts.

\section{Chronologie observable des deux curseurs}

Le calendrier se formule sur
\[
 \Theta_{108}=\mathbb Z/108\mathbb Z,\qquad
 \beta(\theta)=
 \left[\left\lfloor\frac{\widetilde\theta}{9}\right\rfloor\right]_{12},
 \qquad
 r(\theta)=1+(\theta\bmod9),
\]
où \(\widetilde\theta\in\{0,\ldots,107\}\). La signature active est
\[
 \varepsilon(r)=
 \begin{cases}
 +1,&r\in\{1,4,7\},\\
 -1,&r\in\{2,5,8\},\\
 0,&r\in\{3,6,9\}.
 \end{cases}
\]
L'espace observable à deux curseurs est
\[
 \mathcal Q_{\rm obs}
 =X_9^2\times\{N,E,S,O\}^2\times\Theta_{108}.
\]
Si \(\varepsilon=+1\), le curseur additif avance ; si
\(\varepsilon=-1\), le multiplicatif ; si \(\varepsilon=0\), aucun.
Les deux orientations s'inversent après les pas
\(27,54,81,108\). Cette règle définit une permutation
\[
 F_{\rm obs}:\mathcal Q_{\rm obs}\longrightarrow\mathcal Q_{\rm obs}.
\]
Les positions et orientations sont restaurées après \(54\) pas et la phase
complète après \(108\) :
\[
 F_{\rm obs}^{108}=I.
\]
La chronologie d'une seule voie est donc réversible et ne sélectionne pas à elle
seule une distribution sur les continuations ; la sélection de tous les
prolongements compatibles appartient à l'opérateur de transfert des fibres
défini ci-dessous.

\section{Émetteur fini à deux curseurs}

Un germe de curseur est
\[
 \mathcal S=\mathbb T_9^2\times\{N,E,S,O\},
 \qquad |\mathcal S|=324.
\]
TPK minimal reçoit un curseur additif \(P\in\mathcal S\) et un curseur
multiplicatif \(M\in\mathcal S\). Leurs signatures de six fenêtres sont
\(s(P)\) et \(e(M)\). À chaque coordonnée, on applique
\[
 G(s,e)=100s+10\bigl((s+e)\bmod10\bigr)
       +\bigl((3s+5e+1)\bmod10\bigr).
\]
Le chiffre des centaines retrouve \(s\) ; \(s\) étant connue, celui des dizaines retrouve
\(e\). Ainsi \(G\) est injective sur les signatures effectives. L'émetteur
se factorise comme
\[
 \mathcal S^2\xrightarrow{s\times e}
 \operatorname{im}s\times\operatorname{im}e
 \xrightarrow{G}\mathcal U_6,
\]
avec
\[
 |\mathcal S^2|=104\,976,\qquad
 |\operatorname{im}s|=18,\qquad
 |\operatorname{im}e|=26,\qquad
 |\mathcal U_6|=18\cdot26=468.
\]
La factorisation \(468=18\cdot26\) est donc une décomposition de
l'objet : chaque émission conserve une signature additive et une multiplicative.

La projection
\[
 \Psi_6:\mathcal U_6\longrightarrow\mathbb F_3^6,
 \qquad \Psi_6(U)=-U\pmod3,
\]
possède une image de \(243\) mots. La chaîne typée est
\[
 104\,976\ \text{germes}
 \longrightarrow468\ \text{émissions décimales }U_6
 \longrightarrow243\ \text{mots ternaires }w_6.
\]
L'espace ambiant \(\mathbb F_3^6\) contient \(729\) mots ; il ne s'identifie pas
à l'image effective de \(243\). \(U_6\) et \(w_6\) ne sont pas non plus de même
type.

L'émetteur minimal répète, dans la carte dodécaphasique du catalogue,
\[
 U_{12}^{\rm cat}=U_6\Vert U_6.
\]
Sa période visible six explique pourquoi il ne peut pas remplacer l'état signé
\(U_{12}^{\rm sgn}\). Cette différence n'est pas une insuffisance de TPK : elle sépare
la projection de catalogue de l'histoire enrichie.

Chaque bloc décimal émis appartient à un alphabet exact de
\(7\cdot6=42\) triades. Si \(\overline{d_1d_2d_3}\) est l'une d'elles,
\[
 \boxed{d_3\equiv5d_2-2d_1+1\pmod{10}.}
\]
Le balayage des \(104\,976\) germes ne produit aucune triade en dehors de cette
loi. Il convient donc de distinguer la chaîne complète des quotients :
\[
 \mathcal S^2
 \longrightarrow\Omega_{\rm seed}
 \longrightarrow\mathcal U_6
 \xrightarrow{\Psi_6}\operatorname{im}\Psi_6\subset\mathbb F_3^6.
\]
Les présentations microscopiques, les \(468\) émissions décimales, les
\(243\) mots ternaires et l'espace ambiant de \(729\) mots sont quatre
objets distincts.

\paragraph{Commutativité orbitale des morphismes typés.}
Soit
\(W_{243}:=\operatorname{im}\Psi_6\subset\mathbb F_3^6\). L'entonnoir
complet et son quotient diédrique se représentent sans transformer les inclusions en
sélections :
\[
 \underbrace{\mathcal S^2}_{104\,976\ \mathrm{semillas}}
 \xrightarrow[\text{recensement des fibres}]{\text{émission}}
 \underbrace{\mathcal U_6}_{468\ U_6}
 \xrightarrow[\text{projection visible}]{\Psi_6}
 \underbrace{W_{243}}_{243\ w_6}
 \hookrightarrow
 \underbrace{\mathbb F_3^6}_{729\ \mathrm{palabras}},
 \qquad
 W_{243}\xrightarrow{\,/D_3^{\rm cat}\,}
 \underbrace{\mathcal O_{43}}_{43\ \text{orbites}}.
 \tag{TPK-embudo}\label{eq:embudo-orbital-tipado-rev8}
\]
La première flèche regroupe les conditions initiales selon l'émission qu'elles produisent ;
la deuxième oublie la présentation décimale ; la troisième est une inclusion dans
l'espace ambiant algébrique ; la dernière est un quotient par l'action
\(D_3^{\rm cat}\) sur \(W_{243}\). En particulier, il n'y a pas de quotient
\(729\to43\) implicite. Le recensement exact est
\[
 243=38\cdot6+5\cdot3,
\]
c'est-à-dire trente-huit orbites de taille six et cinq de taille trois.
Ces \(43\) classes orbitales ne sont pas non plus les \(43\) tours nonadiques
complets par canal du témoin G9 de mille chiffres : le premier \(43\) appartient
au quotient fini \(W_{243}/D_3^{\rm cat}\) ; le second est un compteur de
profondeur d'une exécution certifiée.

\section{Configuration de l'état holonomique intégral du TPK : feuillet, orientation, résidu, quotient, retenue, frontière, chemin et mémoire}

L'unité dynamique complète est une histoire, non un mot isolé. À la
profondeur \(m\), nous écrivons
\[
 v_m=
 \bigl(
 B_m,\boldsymbol\delta_m,\boldsymbol x_m,
 \boldsymbol{\mathcal C}_m,\Sigma_m,
 \omega_m,\ell_m,g(m),\Lambda_m
 \bigr),
\]
où
\[
 B_m=(u_m^\pi,u_m^e,u_m^\varphi)\in(\mathbb F_3^6)^3.
\]
Ses composantes ont des fonctions distinctes :
\begin{enumerate}
 \item \(B_m\) est le bloc visible conjoint des trois canaux ;
 \item \(\boldsymbol\delta_m\) conserve les nouvelles triades décimales et la
       retenue associée ;
 \item \(\boldsymbol x_m\) est l'état de Hensel profond ;
 \item \(\boldsymbol{\mathcal C}_m\) contient les cylindres ternaires et
       décimaux et leurs restes entiers ;
 \item \(\Sigma_m\) est la signature conjointe de frontière ;
 \item \(\omega_m\) et \(\ell_m\) conservent l'orientation et le feuillet bidirectionnel ;
 \item \(g(m)\in\mathbb Z/9\mathbb Z\) est la phase nonadique ;
 \item \(\Lambda_m\) est le registre dodécaphasique orienté.
\end{enumerate}

\begin{definition}[Présentation opératoire et catégorie TPK]
L'inventaire des opérateurs de l'état holonomique intégral du TPK est
\[
 \mathfrak T_{\rm HMT}=
 \bigl(
 \mathcal A_9,\TRIT,T_{N,E,S,O},T_{\min},
 H,\mathcal C_{3,1000},\Sigma_B,EF,\mathcal G_9,
 \mathcal R_{12}
 \bigr)
\]
qui transporte les histoires \((v_0,\ldots,v_m)\) et leurs troncatures. Son
porteur formel est la catégorie \(\mathsf{TPK}\) sur la base observable :
ses objets sont des états holonomiques intégraux et ses morphismes sont les prolongations
\(\operatorname{Ext}_r\) stables par identité et concaténation, définies
ci-dessous. L'inventaire ne remplace pas cette catégorie. La projection
visible élimine des composantes de \(v_m\) ; elle ne définit pas d'équivalence avec
l'état complet.
\end{definition}

Cette définition contient exactement ce qui disparaît lorsque les curseurs sont réinitialisés
ou que seul le reste est conservé : le chemin, le feuillet, le quotient, la
retenue, l'orientation et la frontière. TPK est donc le point où
l'arithmétique modulaire cesse d'être une opération sans histoire.

\section{Cylindres ternaires-décimaux et état d'extension}

Soit \(N\) l'entier déterminé par un préfixe ternaire de longueur \(L\), et
soit \(D\) l'entier déterminé par \(K\) triades décimales. Nous définissons
\[
 A=N1000^K-D3^L,
 \qquad
 B=(N+1)1000^K-D3^L.
\]
Les cylindres se coupent exactement lorsque
\[
 \boxed{A<3^L,\qquad B>0.}
\]
En incorporant un bloc ternaire \(u\in\{0,\ldots,728\}\),
\[
 N'=729N+u.
\]
S'il n'apparaît pas de nouvelle triade décimale,
\[
 A'=729A+u1000^K.
\]
Si aparece \(\delta\in\{0,\ldots,999\}\),
\[
 D'=1000D+\delta,
\]
\[
 A'=729000A+u1000^{K+1}-\delta\,729\,3^L.
\]
Dans les deux cas,
\[
 B'=A'+1000^{K+\Delta K}.
\]
Ces récurrences entières transportent la donnée décimale sans en faire un
objectif extérieur de la mise à jour.

L'état de Hensel conserve l'information qui n'apparaît pas dans le reste. La
matrice active est
\[
A_H=\begin{pmatrix}
2&2&2&1&2&1\\2&1&2&2&1&1\\1&0&0&1&0&2\\
0&0&1&1&1&0\\2&2&2&0&2&1\\2&1&2&1&0&0
\end{pmatrix},\qquad \det A_H=1,
\]
et la division reste--quotient est
\[
 y=xA_H,\qquad u=y\bmod3,\qquad x^+=\frac{y-u}{3}.
\]
Pour les trois canaux, on forme en outre les cocycles conjoints
\[
 q=(\pi+e+\varphi)\bmod3,
 \qquad
 c=\frac{\pi+e+\varphi-q}{3},
\]
\[
 a=(\pi+e-\varphi)\bmod3,
 \qquad
 h=\frac{\pi+e-\varphi-a}{3}.
\]
La mise à jour du second ordre corrige l'état suivant avec \(c\) et
\(h\). C'est pourquoi deux branches ayant le même bloc visible ne sont pas nécessairement le
même état : elles peuvent différer par leur extension de Hensel et par leurs cylindres.
Comme \(\det A_H=1\), la réduction de \(A_H\) appartient à
\(\operatorname{GL}_6(\mathbb F_3)\). Ainsi, avant d'imposer les cylindres,
les frontières ou la survie, la carte résiduelle de Hensel couvre exactement
\(\mathbb F_3^6\). Cette couverture brute n'implique pas la couverture du système
filtré.

Formellement, sur chaque configuration observable \(q\), on fixe une fibre
\[
 \mathcal F_q=
 \mathsf O_q\times\mathsf H_q\times\mathsf C_q\times
 \mathsf S_q\times\mathsf B_q\times\mathsf\Lambda_q,
\]
correspondant, respectivement, à l'orientation et au feuillet, au reste et au quotient,
à la retenue, à la signature, au cylindre et à la frontière, et au registre de transport. Pour une
flèche cardinale \(r:q\to q'\), le prolongement admissible ne se réduit pas à une
fonction sur les restes ; c'est la relation
\[
 \operatorname{Ext}_r\subseteq\mathcal F_q\times\mathcal F_{q'},
 \qquad
 \operatorname{Ext}_{r_2}\circ\operatorname{Ext}_{r_1}
 \subseteq\operatorname{Ext}_{r_2\circ r_1}.
\]
L'incrément de retenue est un cocycle :
\[
 \kappa(r_2\circ r_1)
 =\kappa(r_2)+\mathsf C(r_2)\bigl(\kappa(r_1)\bigr).
\]
L'application d'oubli est ainsi typée : projeter une extension sur sa voie
observable élimine des composantes de \(\mathcal F_q\), mais n'autorise pas à
identifier les deux objets.

\section{Retours composés et calendriers de transport}

La longueur d'une trace ne détermine pas son type. La présentation académique
sépare le retour composé, le calendrier, la trace et l'atlas selon leurs
états conservés respectifs.

Le symbole \(\mathcal K_{27}\) est réservé à la composition de
vingt-sept mises à jour d'un lacet de transport positionnel et orienté.
L'entier \(27\) est sa longueur de composition ; il n'en fait pas une
trace de 108 itérations ni un émetteur infini. Une autre construction historique
regroupe quatre lacets dont les périodes conjointes atteignent 108 ; ce rotor conjoint
n'est pas désigné par \(\mathcal K_{27}\) isolé.

\begin{longtable}{@{}>{\raggedright\arraybackslash}p{.24\textwidth}>{\raggedright\arraybackslash}p{.31\textwidth}>{\raggedright\arraybackslash}p{.37\textwidth}@{}}
\toprule
Objet&Estado conservado&Significado de 108\\\midrule
\endhead
Contador simplificado&position et vitesse&sa période réelle est 18 ; c'est un
contrôle historique.\\
Rotor conjoint de quatre lacets&quatre voies, lecture, rotation et mouvement&
les périodes individuelles sont \(54,54,36,36\) ; l'état conjoint a une
période de 108.\\
Trace brute cadencée&quatre positions, orientations et phase&le contenu
dynamique se répète à 54.\\
Trace de \(T_{\min}\)&deux positions et deux directions&l'état revient à
54 et la sortie est \(U_6\Vert U_6\).\\
Pseudocode réinitialisé&curseurs réinitialisés à chaque fenêtre&se réduit à
\((222,222,222,333,333,333)^2\) ; cela ne représente pas l'état holonomique intégral du TPK.\\
Rutas condicionadas&chemin et mot prescrit&108 est un horizon, non une période.\\
Odomètre sectoriel&\(\mathbb Z/9\mathbb Z\times\{0,\ldots,11\}\)&l'état
complet a une période de 108 ; l'observable sectoriel, de 36.\\
Projection brute&quotient par oubli de \(X_{\rm TPK}^{\rm enr}\)&ne possède pas
à elle seule la lecture signée.\\
Traza enriquecida&tous les champs de \(v_m\) pendant 108 mises à jour&est
la restriction temporelle de l'état holonomique intégral du TPK.\\
Atlas des voies&81 cellules/56 signatures ou 324 germes orientés&c'est une taxonomie,
non une orbite de 108 états.\\
\bottomrule
\end{longtable}

Le lacet \(\mathcal K_{27}\) est l'opérateur de transport fini de
vingt-sept mises à jour. Les fenêtres, le rotor conjoint et la projection
calendaire reçoivent des symboles distincts, de sorte que la longueur partagée
ne remplace ni le domaine ni la composition de chaque objet.

\section{Transfert des fibres et mémoire de rotation}

La factorisation \(\mathcal U_6\simeq\mathcal S_+\times\mathcal S_\times\)
permet de définir, par rapport au comptage uniforme, les espérances conditionnelles
\[
 (E_+f)(s,e)=\frac1{18}\sum_{s'\in\mathcal S_+}f(s',e),
 \qquad
 (E_\times f)(s,e)=\frac1{26}
 \sum_{e'\in\mathcal S_\times}f(s,e').
\]
\(E_+\) renouvelle la fibre additive et conserve la multiplicative ;
\(E_\times\) réalise l'opération duale. Ce sont les seuls projecteurs de Markov
qui marginalisent complètement la coordonnée active, conservent l'inactive et
préservent le comptage uniforme.

Soit \(\tau=(+1,-1,0)^3\) le mot de phase sur \(C_9\), et
\(P_{+1}=E_+\), \(P_{-1}=E_\times\), \(P_0=I\). Le transfert nonadique est
\[
 \mathcal K_{\rm ph}\bigl((u,r),(v,r+1)\bigr)
 =P_{\tau(r)}(u,v).
\]
Il est doublement stochastique, irréductible, de période \(9\), et possède une unique
distribution stationnaire \(1/(9\cdot468)\). De plus,
\[
 \mathcal K_{\rm ph}^{\,9}=E_+E_\times,\qquad
 (E_+E_\times)(u,v)=\frac1{468}.
\]
Cette égalité exprime le renouvellement simultané de toutes les continuations
compatibles, non le successeur déterministe d'une seule voie.

Trois opérations géométriques sur les germes conservent une information
supplémentaire : \(P\), avancée dans l'orientation active ; \(H\), demi-tour ; et
\(Q\), quart de tour. La paire \((P,H)\) raffine les \(26\) signatures produit
en \(28\) classes. En particulier, la signature \(555555\) se sépare en trois feuillets
de huit représentants, et
\[
 18\cdot28=504=468+36.
\]
Le terme \(36\) compte les feuillets supplémentaires du raffinement et n'est pas
l'objet \(R_{36}\). La paire \((P,Q)\) produit \(162\) classes de deux
éléments, permutées transitivement par
\[
 J(i,j,d)=\bigl(7-i,\,7-j,\,\overline d\bigr)\pmod9.
\]

Dans une trace de \(108\) pas apparaissent \(36\) signes de chaque type,
\(72\) changements effectifs entre les deux feuillets et quatre inversions
d'orientation ; le cocycle signé total est nul. Ces décomptes sont des invariants
du registre d'évolution et relient la dynamique nonadique au registre dodécaphasique.

\section{Prolongement ternaire et frontière conjointe}

Relativement aux germes, frontières et lecteurs nommés du profil
publié, les mots visibles initiaux sont
\[
 w_6^\pi=010211,\qquad w_6^e=201101,\qquad
 w_6^\varphi=121200.
\]
Après cinq blocs, les préfixes sont
\[
\begin{aligned}
w_{30}^{\pi}&=010211|012222|010211|002111|110221,\\
w_{30}^{e}&=201101|121221|102011|012222|102011,\\
w_{30}^{\varphi}&=121200|112202|121020|010210|010200.
\end{aligned}
\]
La première frontière profonde est
\[
 R_{36}=
 \begin{pmatrix}222220\\021222\\102011\end{pmatrix}.
\]
Pour chaque canal, on définit le préfixe
\(w_{36}^\chi:=w_{30}^\chi|u_5^\chi\) ; \(R_{36}\) est le triplet ordonné des
trois sixièmes blocs. Aucun n'est un émetteur ni un état terminal.

Pour \(B_k=(u_k^\pi,u_k^e,u_k^\varphi)\), les signatures
\[
 q_k=u_k^\pi+u_k^e+u_k^\varphi,\qquad
 a_k=u_k^\pi+u_k^e-u_k^\varphi
\]
déterminent
\[
 u_k^\varphi=2(q_k-a_k),\qquad
 u_k^\pi+u_k^e=q_k-u_k^\varphi
 \quad\text{dans }\mathbb F_3^6.
\]
Ainsi, \(\varphi\) est l'axe fixé et la liberté résiduelle est la répartition
interne \(\pi/e\). Dans la neuvième phase,
\[
 B_9=(100100|020112|010122)
\]
admet trois sorties locales ayant le même axe \(\varphi\) et le même agrégat
\(\pi+e\) ; les trois survivent à un horizon et seule
\[
 B_{10}=(101222|112221|112002)
\]
survit au second. C'est la descente exécutable \(3\to1\) d'EF--G9.
La branche publiée se poursuit jusqu'à vingt blocs par canal : \(w_{30}\) et
\(R_{36}\) sont préfixe et frontière, non clôture.

\section{Relation d'extension et monodromie nonadique}
\label{sec:extension-monodromia-tpk}

Soit \(\mathcal H_m\) l'ensemble des histoires enrichies admissibles de
profondeur \(m\), et soit
\[
 p_{n,m}:\mathcal H_n\longrightarrow\mathcal H_m
\]
la troncature. Pour \(N\geq m\),
\[
 \operatorname{Surv}_m^{(N)}=p_{N,m}(\mathcal H_N),
 \qquad
 \operatorname{Surv}_m=\bigcap_{N\geq m}\operatorname{Surv}_m^{(N)}.
\]
La relation EF--G9 conserve les extensions qui appartiennent à ce noyau
stable. La phase satisfait
\[
 g(m+9)=g(m),
\]
mais l'état transporté satisfait, en général,
\[
 v_{m+9}\ne v_m.
\]
Le retour agit sur la fibre des états d'une même phase ; il ne réinitialise pas
l'histoire.

\begin{definition}[Nombre HMT]
\label{def:numero-hmt-tpk-completo}
Un nombre HMT du canal \(\chi\) est uniquement une section globale
compatible
\[
 x=(x_m)_{m\geq0}\in
 X_\chi:=\varprojlim_m\operatorname{Surv}_m^{(\chi)}.
\]
Le lecteur et la valeur publiée ne font pas partie de la section.
\end{definition}

Une présentation ou mesure du nombre est une paire \((x,L)\), où
\(L\) est un lecteur déclaré dont le domaine contient \(x\). Lorsque \(L\) est le
lecteur archimédien des cylindres, \(L(x)\) est la valeur publiée ; \(x\)
conserve en outre l'orientation, le quotient, la retenue et la frontière.

\begin{theorem}[Existence d'un prolongement global]
\label{thm:prolongacion-global-tpk-completo}
Soit \(h\in\mathcal H_m\). Si l'arbre des prolongements de \(h\) est
à branchement fini et possède au moins un nœud à toute profondeur, alors
il existe une histoire infinie dont les troncatures prolongent \(h\). Si, de plus,
chaque niveau stable est un singleton et que les états uniques sont compatibles, la
section globale est unique.
\end{theorem}
\begin{proof}
Les niveaux finis non vides forment un arbre à branchement fini. Le
lemme de König fournit une branche infinie. La compatibilité des
troncatures identifie cette branche à un élément de la limite inverse. Si
chaque niveau contient un seul état, il n'existe pas de seconde branche compatible.
\end{proof}

Le théorème précédent est abstrait et conditionnel : il n'affirme pas à lui seul la
non-vacuité ni l'unicité stable d'un canal concret. Pour
\(\pi,e,\varphi\), ces hypothèses sont établies ensuite au moyen de la filtration
rationnelle exécutable du
\cref{thm:generacion-monodromica-conjunta-rev7}. Cet emplacement ne conserve
ici que le principe de limite inverse et n'anticipe pas son application.

\section{Composition des opérateurs fondamentaux}

Les opérateurs précédents conduisent à l'état holonomique terminal sans
aplatir ses coordonnées :
\[
 \begin{aligned}
 \mathrm{APP}&\longrightarrow\mathrm{TRIT}
 \longrightarrow\text{route orientée}
 \longrightarrow\text{état holonomique intégral}\\
 &\longrightarrow\text{cylindres et survie}
 \longrightarrow x_{\rm term}.
 \end{aligned}
\]
La fermeture dodécaphasique reçoit alors deux évaluations coordonnées du même
état :
\[
 x_{\rm term}
 \xrightarrow{(\pi_{\rm term},R_{\rm sgn})}
 \bigl(B_{\rm term},U_{12}^{\rm sgn}\bigr),
 \qquad
 U_{12}^{\rm sgn}\xleftrightarrow{\mathcal H_{12}}K,
\]
\[
 (P,E,\Phi,K,c_{13}=0)\longleftrightarrow\alpha_{12}.
\]
Il n'apparaît pas de flèche \(B_{\rm term}\to U_{12}^{\rm sgn}\) : les deux
coordonnées proviennent de \(x_{\rm term}\). Le passage ultérieur à la mécanique
dimensionnelle ne crée pas non plus une autre machine ; il prolonge la généalogie conservée
au moyen du registre chronologique enrichi, de la signature, du caractère et des tours.

\ifdefined\HMTInternalTraceability
\section{Registre académique et traçabilité}

Le fichier \path{datos/registro_operadores_tpk.json} fait partie du paquet.
Pour chaque opérateur, il déclare le nom académique, l'alias historique, le domaine,
le codomaine, la transition, les invariants, le statut et la source locale. Il inclut également
la séparation des dix constructions historiques associées au nombre
108. La porte de publication vérifie ce registre et bloque une sortie
qui réduirait de nouveau TPK à son catalogue ou à une trace brute.

Dans le reste de l'article, quatre abréviations, déjà typées, seront utilisées :
\[
 w_6\in\mathbb F_3^6,\qquad
 w_{30}\in\mathbb F_3^{30},\qquad
 R_{36}\in(\mathbb F_3^6)^3,\qquad
 U_{12}^{\rm sgn}\in\mathbb Z^{12}.
\]
La première est un mot germe ; la deuxième, un préfixe de cinq blocs ;
 la troisième, la frontière conjointe d'entrée dans EF--G9 ; la quatrième, la
coordonnée signée déclarée de l'état dodécaphasique enrichi. Aucune égalité de cardinalités n'autorise à
les identifier.
\fi
  \part{Algèbre holonomique des chemins et structure électronique locale}

\chapter{Présentation universelle de l'algèbre holonomique}
\section{Données génératrices}

Une profondeur \(n\) étant fixée, soit \(X_n\) l'ensemble fini des états
APP--TRIT--TPK qui conservent, au niveau correspondant, feuillet, régime,
orientation, résidu, quotient, retenue, frontière, hauteur, chemin, mémoire et
survie. Soit \(\mathcal C_n\) la petite catégorie engendrée par les pas

\[
U_t={\rm Upd}_t\circ{\rm Tra}_t\circ{\rm Sel}_t.
\]

Les objets de \(\mathcal C_n\) sont les états de \(X_n\) ; ses flèches sont les
histoires admissibles. La composition existe lorsque l'état terminal de la
première histoire coïncide avec l'état initial de la seconde.

À chaque \(x\in X_n\) est associée la fibre tritique

\[
A_x=A_{\tau(x)}
=\mathbb R[J_x]/(J_x^2+\tau(x)e_x),
\qquad \tau(x)\in\{+1,0,-1\}.
\]

Le symbole \(e_x\) désigne l'identité locale de la fibre. Pour chaque flèche
\(u:x\to y\), le transport TPK induit une action

\[
\rho_u:A_x\longrightarrow A_y
\]

sur les coefficients qui survivent au pas. Lorsqu'une transition change
de régime, \(\rho_u\) s'interprète comme le bimodule typé de cette
transition, non comme une identification des deux algèbres quadratiques.

Le registre de mémoire de concaténation définit un multiplicateur

\[
\Omega(v,u)\in Z(A_z)^\times,
\qquad
u:x\to y,\qquad v:y\to z,
\]

qui enregistre la retenue, la torsion et la contribution de frontière produites
lors de la composition. Sa loi de compatibilité est

\[
\rho_w(\Omega(v,u))\Omega(w,v\circ u)
=\Omega(w,v)\Omega(w\circ v,u).
\tag{1.1}
\]

L'équation (1.1) est l'identité de cocycle qui empêche la mémoire d'une
histoire de dépendre du parenthésage de sa composition.

\section{Présentation par générateurs et relations}

Soit \(\mathcal F_n\) l'algèbre associative libre engendrée par

\[
\{e_x,J_x:x\in X_n\}\cup\{T_u:u\in{\rm Mor}(\mathcal C_n)\}.
\]

Nous définissons l'idéal \(\mathcal I_{\rm HMT}\) au moyen des relations

\[
e_xe_y=\delta_{xy}e_x,
\qquad
\sum_{x\in X_n}e_x=1_n,
\tag{2.1}
\]

\[
J_x=e_xJ_xe_x,
\qquad
J_x^2=-\tau(x)e_x,
\tag{2.2}
\]

\[
T_u=e_{t(u)}T_ue_{s(u)},
\tag{2.3}
\]

\[
T_ua=\rho_u(a)T_u,
\qquad a\in A_{s(u)},
\tag{2.4}
\]

et

\[
T_vT_u=\Omega(v,u)T_{v\circ u}
\tag{2.5}
\]

lorsque \(v\circ u\) existe ; le produit est nul lorsque les flèches ne sont pas
composables. L'algèbre holonomique finie est

\[
\boxed{
\mathfrak H_n=\mathcal F_n/\mathcal I_{\rm HMT}.}
\tag{2.6}
\]

La somme de \(\mathfrak H_n\) superpose des chemins possibles. Le produit concatène
des histoires. Le facteur \(\Omega\) empêche la concaténation d'effacer la retenue.
Ainsi, somme et produit ont acquis un sens simultanément
arithmétique et dynamique.

\section{Théorème d'associativité}

\textbf{Théorème 3.1.} Si \(\rho\) respecte les identités locales et si \(\Omega\)
satisfait (1.1), le produit défini par (2.5) est associatif.

\textbf{Démonstration.} Soient \(u:x\to y\), \(v:y\to z\) et \(w:z\to t\). Pour
des coefficients \(a,b,c\) situés dans les fibres terminales pertinentes,

\[
\begin{aligned}
((cT_w)(bT_v))(aT_u)
={}&c\rho_w(b)\Omega(w,v)
\rho_{w\circ v}(a)\Omega(w\circ v,u)T_{wvu},\\
(cT_w)((bT_v)(aT_u))
={}&c\rho_w(b\rho_v(a)\Omega(v,u))
\Omega(w,v\circ u)T_{wvu}.
\end{aligned}
\]

La covariance de \(\rho\) égalise les facteurs contenant \(a\) et \(b\).
L'identité (1.1) égalise les deux produits de \(\Omega\). L'associativité
sur les monômes s'ensuit, et la bilinéarité l'étend à toute
\(\mathfrak H_n\). \(\square\)

La démonstration précise la fonction de la mémoire. Sans \(\Omega\), la
concaténation conserverait le chemin nu, mais non l'écart accumulé
entre ses relèvements. Avec \(\Omega\), deux manières de regrouper la même
histoire produisent le même état complet.

\section{Propriété universelle de représentation}

\textbf{Théorème 4.1.} Soit \(\{E_x\}_{x\in X_n}\) une famille de modules. Supposons
que l'on se donne des représentations locales

\[
\pi_x:A_x\to\operatorname{End}(E_x)
\]

et des opérateurs \(V_u:E_x\to E_y\) satisfaisant

\[
V_u\pi_x(a)=\pi_y(\rho_u(a))V_u,
\tag{4.1}
\]

\[
V_vV_u=\pi_z(\Omega(v,u))V_{v\circ u}.
\tag{4.2}
\]

Il existe une unique représentation

\[
\Pi:\mathfrak H_n\to
\operatorname{End}\!\left(\bigoplus_{x\in X_n}E_x\right)
\]

qui réalise \(e_x,J_x,T_u\) comme les projections, générateurs et transports
donnés.

\textbf{Démonstration.} L'affectation des générateurs définit un homomorphisme
à partir de l'algèbre libre \(\mathcal F_n\). Les équations (4.1)--(4.2) et les
relations locales impliquent que \(\mathcal I_{\rm HMT}\) appartient à son
noyau. L'homomorphisme se factorise par le quotient (2.6). L'unicité découle
du fait que les classes de \(e_x,J_x,T_u\) engendrent le quotient. \(\square\)

La propriété universelle transforme la précédence causale du générateur en une affirmation
algébrique. Un lecteur ultérieur est légitime lorsqu'il représente les relations
du générateur ; il n'a pas besoin de les reconstruire à partir de son codomaine. Le continu,
l'électron, la gravité et la branche exceptionnelle sont des réalisations distinctes
d'une même présentation.

\section{Le cocycle nonadique élémentaire}

La section canonique \(s:C_9\to\{0,\ldots,8\}\subset C_{108}\) produit

\[
c_9(a,b)=\left\lfloor\frac{a+b}{9}\right\rfloor\in C_{12}.
\tag{5.1}
\]

L’identité

\[
c_9(a,b)+c_9(a+b\bmod9,c)
=c_9(b,c)+c_9(a,b+c\bmod9)
\tag{5.2}
\]

est exacte. Si \(X\) est l'opérateur de relèvement et \(\lambda\) l'image du
générateur de mémoire, alors

\[
X^9=\lambda I
\]

et

\[
X^aX^b
=\lambda^{c_9(a,b)}X^{a+b\bmod9}.
\tag{5.3}
\]

L'équation (5.3) est la forme minimale de l'hélice. L'exposant visible
revient modulo neuf ; le facteur \(\lambda\) conserve le tour. Si
\(\lambda\ne1\), retour de phase et retour de l'état sont des opérations
distinctes.

\section{Raffinement des états et prolongation des lecteurs}

Les ensembles finis d'états forment un système projectif

\[
\cdots\xrightarrow{r_{n+2,n+1}}X_{n+1}
\xrightarrow{r_{n+1,n}}X_n.
\]

L'état profond est

\[
X_\infty=\varprojlim_nX_n.
\tag{6.1}
\]

Les fonctions et les idempotents de cylindre se déplacent en sens opposé :

\[
r_{n+1,n}^*:f\longmapsto f\circ r_{n+1,n}.
\tag{6.2}
\]

Si \(e_x\) est l'indicatrice du cylindre \(x\in X_n\),

\[
r^*e_x=\sum_{y:r(y)=x}e_y.
\tag{6.3}
\]

La famille d'algèbres forme donc un système inductif :

\[
\mathfrak H_{\rm cyl}=\varinjlim_n\mathfrak H_n.
\tag{6.4}
\]

\textbf{Théorème 6.1.} Les inclusions (6.2) sont unitaires ; les idempotents de
cylindre restent idempotents ; et tout lecteur compatible à profondeur
finie détermine un unique lecteur sur \(\mathfrak H_{\rm cyl}\).

\textbf{Démonstration.} L'image réciproque d'une partition est une partition. Par conséquent,
\(r^*(fg)=r^*f\,r^*g\), \(r^*1_n=1_{n+1}\) et la somme (6.3) est orthogonale.
La propriété universelle de la limite inductive prolonge toute famille compatible
de lecteurs. \(\square\)

Ce théorème apporte le sens exact de «fractaliser vers l'intérieur et
déployer vers l'extérieur». Les états se complètent par restriction ; les
lecteurs se complètent par prolongation. La conservation de l'unité est le
caractère unitaire de ces inclusions.

\section{Holonomie et graduation de mémoire}

Soit \(d_M\) le degré de mémoire. L'holonomie satisfait

\[
d_M(\Gamma_9x)=d_M(x)+1.
\tag{7.1}
\]

Dans le domaine réversible, \(\Gamma_9\) induit un automorphisme de
\(\mathfrak H_{\rm cyl}\) et l'on peut former

\[
\mathfrak H_{\rm rev}\rtimes_{\Gamma_9}\mathbb Z.
\tag{7.2}
\]

Dans le domaine général, l'avance produit un endomorphisme et l'objet est

\[
\mathfrak H_{\rm cyl}\rtimes_{\Gamma_9}\mathbb N.
\tag{7.3}
\]

La distinction entre (7.2) et (7.3) empêche d'introduire une inversion là où
seule une avance a été construite. L'inversion hélicoïdale déjà typée agit dans
(7.2) ; le miroir \(\pi/e\) agit sur les coordonnées du lecteur et constitue
une autre involution.

\section{Algèbre tritique totale}

Soit

\[
\mathbb A_{\rm tri}=A_+\oplus A_0\oplus A_-
\]

avec les idempotents centraux \(p_+,p_0,p_-\). Les générateurs satisfont

\[
J_+^2=-p_+,
\qquad
J_0^2=0,
\qquad
J_-^2=p_-.
\tag{8.1}
\]

Soit

\[
\mathbb J=p_+J_++p_0J_0+p_-J_-.
\]

L'exponentielle définie par sa série interne satisfait

\[
\boxed{
\begin{aligned}
\operatorname{Exp}_\star(t\mathbb J)
={}&p_+(\cos t+J_+\sin t)\\
&+p_0(1+tJ_0)\\
&+p_-(\cosh t+J_-\sinh t).
\end{aligned}}
\tag{8.2}
\]

\textbf{Démonstration.} Dans la fibre elliptique, les puissances paires et impaires de
\(J_+\) alternent les signes. Dans la fibre parabolique, toutes les
puissances d'ordre au moins deux disparaissent. Dans la fibre hyperbolique, les puissances paires
sont l'identité et les impaires sont \(J_-\). Regrouper les trois séries prouve
(8.2). \(\square\)

L'identité conserve simultanément les trois géométries. Aucune d'elles
ne s'obtient par prolongement informel d'une autre ; les trois sont des termes
orthogonaux du même objet.

\section{Euler comme projection elliptique}

Soit \(\pi{}\) la clôture elliptique déjà produite par la connexion
nonadique. Alors

\[
\cos\pi{}=-1,
\qquad
\sin\pi{}=0
\]

dans la réalisation circulaire, et (8.2) donne

\[
p_+\operatorname{Exp}_\star(\pi{}\mathbb J)+p_+=0.
\tag{9.1}
\]

L'équation (9.1) retrouve l'identité d'Euler sur la face elliptique. L'élément
total conserve en outre

\[
p_0\operatorname{Exp}_\star(\pi{}\mathbb J)
=p_0(1+\pi{}J_0),
\]

et

\[
p_-\operatorname{Exp}_\star(\pi{}\mathbb J)
=p_-(\cosh\pi{}+J_-\sinh\pi{}).
\]

L'information omise par la représentation complexe est précisément le jet
de seuil et la dilatation. L'identité tritique ne remplace pas l'identité
d'Euler : elle identifie quelle projection la produit et quelle mémoire reste hors de cette
projection.

\section{L'autoéchelle hyperbolique et \(\varphi\)}

Soit

\[
F_{\rm av}=\begin{pmatrix}0&1\\1&1\end{pmatrix}.
\]

Son carré a pour déterminant un et pour trace trois. Nous définissons

\[
H_\varphi=\frac{2F_{\rm av}^2-3I}{\sqrt5}.
\tag{10.1}
\]

Un calcul direct donne

\[
H_\varphi^2=I.
\tag{10.2}
\]

Comme

\[
\cosh(2\log\varphi)=\frac32,
\qquad
\sinh(2\log\varphi)=\frac{\sqrt5}{2},
\]

on obtient

\[
\boxed{
F_{\rm av}^2
=\operatorname{Exp}(2\log\varphi{}\,H_\varphi).}
\tag{10.3}
\]

Le régime hyperbolique ne se limite pas à admettre une unité quadratique : la
matrice d'incidence aire--volume réalise exactement cette unité. Le rapport
\(\varphi^{-2}\) est la direction stable d'une évolution unimodulaire.

\section{La fonction structurelle de \(e\)}

La série \(\operatorname{Exp}_\star\) transforme le générateur additif en
transport multiplicatif :

\[
\operatorname{Exp}_\star(a+b)
=\operatorname{Exp}_\star(a)\operatorname{Exp}_\star(b)
\]

lorsque \(a\) et \(b\) commutent. Sa réalisation scalaire sur l'identité
définit

\[
\operatorname{Exp}_\star(1)=e{}I.
\tag{11.1}
\]

L'équation (11.1) ne remplace pas la génération décimale nonadique de \(e\).
Elle identifie sa fonction algébrique : \(e\) est l'échelle qui convertit la somme
infinitésimale en une loi multiplicative de transport. Tel est le pont
intrinsèque entre les deux feuillets APP dans la représentation exponentielle.

\section{Alpha comme fonctionnelle de transgression dodécaphasique}

Le cocycle (5.1) représente l'impossibilité de choisir une section
homomorphe \(C_9\to C_{108}\). Le tour visible produit un élément du
noyau dodécaphasique. La clôture de \(\alpha\) reçoit cet élément avec le
mot signé et le détermine sur une droite scalaire.

Il convient de distinguer le cocycle de son évaluation. Soit

\[
\mathcal T_\alpha:
Z^2(C_9;C_{12})\times X_{12}^{\rm sgn}
\longrightarrow\mathcal L_\alpha
\tag{12.1}
\]

la fonctionnelle qui réalise la concaténation avec mémoire. Alors

\[
\alpha{}=\mathcal T_\alpha(c_9,K).
\tag{12.2}
\]

L'équation (12.2) n'identifie pas \(\alpha\) à une classe de cohomologie.
Elle affirme que \(\alpha\) est la réalisation scalaire de cette classe après
l'intégration de l'état dodécaphasique. La seconde voie détermine la même section
au moyen de

\[
P_9(\alpha{})=\delta,
\qquad
\delta=\log_{10}\frac{10}{9},
\tag{12.3}
\]

et les deux se réunissent dans le produit fibré de résolution.

Ce type algébrique explique la singularité de \(\alpha\). \(\pi,\varphi,e\)
sont des coordonnées de clôture, d'échelle et de transport. \(\alpha\) est la fonctionnelle
qui mesure comment ces coordonnées survivent au relèvement non scindé.

\section{Matrice compacte du système corrélé}

Les quatre fonctions se réunissent sans se confondre dans

\[
\boxed{
\begin{pmatrix}
\operatorname{Exp}_\star(\pi{}J_+) & 0&0&0\\
0&\operatorname{Exp}_\star(I)&0&0\\
0&0&\operatorname{Exp}_\star(2\log\varphi{}H_\varphi)&0\\
0&0&0&\operatorname{Exp}_\star(P_9(\alpha{})J_0)
\end{pmatrix}
=
\begin{pmatrix}
-I&0&0&0\\
0&e{}I&0&0\\
0&0&F_{\rm av}^2&0\\
0&0&0&I+\delta J_0
\end{pmatrix}.}
\tag{13.1}
\]

La matrice (13.1) est une identité de fonctions, non une table de valeurs. Ses
quatre blocs signifient clôture elliptique, multiplication exponentielle,
autoéchelle hyperbolique et couture parabolique. La compatibilité dodécaphasique
de la dernière ligne conserve la généalogie conjointe.

\section{Le coin électronique}

Soit

\[
P=\begin{pmatrix}1&0\\0&0\end{pmatrix},
\qquad
Q=\begin{pmatrix}3/4&\sqrt3/4\\\sqrt3/4&1/4\end{pmatrix}.
\]

On définit

\[
S=2P-I,
\qquad
J=\frac4{\sqrt3}[P,Q].
\]

Alors

\[
S^2=I,
\qquad
J^2=-I,
\qquad
SJ=-JS.
\tag{14.1}
\]

Par conséquent,

\[
\boxed{
e_{\rm vac}\mathfrak H_ne_{\rm vac}\twoheadrightarrow
{\rm Cl}_{1,1}\cong M_2(\mathbb R).}
\tag{14.2}
\]

La flèche (14.2) exprime que le bloc électronique est un quotient ou une
représentation d'un coin local de l'arithmétique holonomique. La
non-commutativité n'est pas ajoutée depuis la mécanique quantique : elle réside déjà dans
l'ordre des feuillets APP. La racine interne de \(-1\) est le générateur elliptique ;
les involutions déployées conservent la vacance et la dilatation.

\section{Lecteurs et types de sortie}

Soit

\[
\operatorname{Read}(\mathfrak H)
=\coprod_B\operatorname{Rep}_{\rm comp}(\mathfrak H,B).
\tag{15.1}
\]

Une constante s'obtient comme

\[
c=\sigma_B(\rho(x_\infty)),
\tag{15.2}
\]

où \(\rho\) est une représentation compatible et \(\sigma_B\) un invariant
scalaire propre à son codomaine. Cette définition inclut sans confondre :

\[
\begin{array}{c|c|c}
\text{sortie}&B&\sigma_B\\ \hline
\pi&\text{module elliptique}&\text{temps minimal de clôture}\\
\varphi&\text{module hyperbolique}&\text{rayon spectral stable}\\
e&\text{semi-groupe exponentiel}&\text{évaluation unitaire}\\
\alpha&\text{produit fibré dodécaphasique}&\text{recollement scalaire}\\
G_{\rm Cat}&\text{module de caractère}&\operatorname{Tr}(N^{-2}Q_4)\\
k_B&\text{droites thermique et énergétique}&\text{transduction d’unité}\\
\gamma_{\rm BI}&\text{module aire--volume}&(12,-1)\text{ sur }90/120\\
G&\text{module gravitationnel}&\text{caractère contragrédient}
\end{array}
\]

Une particule n'est pas un chiffre de (15.2) : c'est une section persistante d'un
module. Un champ est le transport de ces sections. Une symétrie est le
groupe d'automorphismes qui préserve les relations et le bord du module.

\section{Théorème de conservation et de symétrie émergente}

Soit \(1_n=\sum_{x\in X_n}e_x\). Les inclusions de raffinement satisfont

\[
r^*1_n=1_{n+1}.
\tag{16.1}
\]

Une histoire individuelle définit en outre une famille compatible
d'idempotents de cylindre. La première est l'unité globale de l'algèbre ; la
seconde est l'unité généalogique d'une section. Les distinguer évite de
confondre conservation de l'espace des possibilités et conservation d'un
chemin particulier.

Pour une représentation \(\rho\) et une frontière \(b_\partial\), nous définissons

\[
\operatorname{Sym}(\rho,b_\partial)
=\{g\in\operatorname{Aut}(\rho(\mathfrak H)):
g\rho(\mathcal R_{\rm HMT})=\rho(\mathcal R_{\rm HMT}),
\;gb_\partial=b_\partial\}.
\tag{16.2}
\]

La symétrie est, par définition, le stabilisateur d'une conservation déjà
construite. Dans la branche exceptionnelle, (16.2) conduit aux stabilisateurs
d'incidence Witt--Golay--Leech--\(V^\natural\)--Monstre. Dans la branche physique,
elle conduit aux groupes qui préservent les modules électronique, de jauge ou
tensoriel. On n'impose pas une symétrie pour sélectionner l'état ; on obtient
la symétrie qui laisse sa projection invariante.

\section{Information, entropie et quotients}

Si \(U\) est bijectif sur l'état complet,

\[
H(X\mid U(X))=0.
\tag{17.1}
\]

Soit \(q\) un quotient qui oublie une coordonnée. Alors

\[
H(X)=H(qX)+H(X\mid qX).
\tag{17.2}
\]

Un trit uniforme oublié apporte

\[
H(X\mid qX)=\log3.
\]

Le transducteur thermique détermine

\[
Q_{\rm reset}\ge k_B\Theta\log3.
\tag{17.3}
\]

Les équations (17.1)--(17.3) montrent que l'entropie ne mesure pas l'avance de
l'holonomie. Elle mesure l'information écartée par un lecteur. La complexité
sturmienne \(p(m)=m+1\) permet une mémoire infinie avec une entropie topologique nulle ;
le coût thermique n'apparaît que lorsqu'on rend une projection irréversible.

\section{Théorème central}

\textbf{Théorème 18.1.} L'objet

\[
\boxed{
\mathfrak H_{\rm HMT}^{\rm alg}
=\left[\varinjlim_n
\bigl(\mathbb R\langle e_x,J_x,T_u\rangle/
\mathcal R_{\rm HMT}\bigr)\right]
\rtimes_{\Gamma_9}\mathbb N}
\tag{18.1}
\]

est une algèbre holonomique trigéométrique, graduée par la phase et la mémoire, avec un
produit associatif tordu par la retenue. Son espace d'états est la
limite projective \(X_\infty\). Toute réalisation compatible d'APP, de TRIT et de TPK
se factorise par (18.1). Ses représentations typées déterminent la structure
discrète conjointe du continu, le système corrélé, l'électron, l'action,
la gravité, les constantes et les symétries sans transformer leurs valeurs
terminales en générateurs.

La notation de (18.1) est délibérément algébrique : elle n'impose pas de norme
avant le choix du codomaine. Si
\(\rho:\mathfrak H_{\rm HMT}^{\rm alg}\to B(\mathcal K)\) est une
représentation compatible et bornée, sa réalisation complétée est

\[
\mathfrak H_{\rm HMT}^{\rho}
=\overline{\rho(\mathfrak H_{\rm HMT}^{\rm alg})}^{\|\cdot\|}.
\tag{18.2}
\]

Ainsi, chaque lecteur apporte la topologie qu'il conserve effectivement ;
aucune complétion conventionnelle n'est introduite comme partie du générateur.

\textbf{Démonstration.} L'associativité procède du Théorème 3.1. L'universalité,
du Théorème 4.1. La conservation de l'unité et la dualité de raffinement,
du Théorème 6.1. L'élévation de mémoire procède du cocycle nonadique (5.3).
La coexistence des courbures procède de (8.1). Les réalisations s'obtiennent
par la propriété universelle et les lecteurs (15.1)--(15.2). \(\square\)

\section{Certification finie}

La vérification exhaustive sur les domaines finis démontre :

\begin{enumerate}
\item les \(9^3\) identités du cocycle de retenue ;
\item la représentation matricielle \(X^9=2I\) et sa loi tordue ;
\item les trois relations \(J_+^2=-I\), \(J_0^2=0\), \(J_-^2=I\) ;
\item les trois branches de l'exponentielle ;
\item \(H_\varphi^2=I\) y \(\exp(2\log\varphi H_\varphi)=F_{\rm av}^2\);
\item le coin \({\rm Cl}_{1,1}\) ;
\item le caractère unitaire contravariant/covariant du raffinement ;
\item la synchronie \(A_2\) avec avance diagonale \(-16\) ;
\item la différence entre une permutation réversible et un quotient tritique ;
\item le couple nul de Witt contenu dans le coin électronique ;
\item les relations exactes du couple d'observables de phase et le commutateur de mémoire ;
\item l'irréductibilité d'APP, de TRIT et de TPK par quotients d'oubli, ainsi que l'inversion du potentiel dans le lecteur de séquences.
\end{enumerate}

Ces identités constituent une vérification finie reproductible du système algébrique.

\section{Le couple nul de Witt au sein du coin électronique}

Le coin électronique apporte une réalisation locale encore plus complète de
la trigéométrie. Soient

\[
S=2P-I,\qquad
J=\frac{[P,Q]}{\sqrt3/4}.
\tag{20.1}
\]

Les relations déjà démontrées sont

\[
S^2=I,\qquad J^2=-I,\qquad SJ=-JS.
\tag{20.2}\label{eq:relaciones-cl11}
\]

Définissons maintenant

\[
n_+=\frac{S+J}{2},\qquad
n_-=\frac{S-J}{2}.
\tag{20.3}
\]

Por (20.2),

\[
n_+^2=n_-^2=0,\qquad
n_+n_-+n_-n_+=I.
\tag{20.4}\label{eq:par-nulo-witt}
\]

La démonstration est immédiate mais structurellement décisive :
\((S\pm J)^2=S^2+J^2\pm(SJ+JS)=0\), tandis que la somme des
deux produits croisés vaut
\(\tfrac12(S^2-J^2)=I\). Ainsi, le même bloc
\({\rm Cl}_{1,1}\cong M_2(\mathbb R)\) contient :

\[
\begin{array}{c|c|c}
\text{générateur local}&\text{carré}&\text{régime}\\ \hline
J&-I&\text{elliptique}\\
n_\pm&0&\text{parabolique}\\
S&I&\text{hyperbolique}.
\end{array}
\tag{20.5}
\]

C'est une décomposition de Witt locale exacte. La vacance électronique
acquiert ainsi une seconde caractérisation équivalente dans son coin : outre
le fait d'être la perte minimale d'une unité de retenue multiplicative sous les
conservations APP, elle est une direction nulle parabolique entre les deux
orientations \(n_+\) et \(n_-\). TRIT type globalement ces trois classes et TPK
les transporte avec feuillet, phase et mémoire. La structure d'incidence
Paley--Witt ultérieure conserve son propre domaine combinatoire ; la
décomposition (20.3)--(20.5) apporte le modèle local à partir duquel on peut
formuler sa réalisation sans égaler les deux objets.

\section{Produit croisé des observables de phase nonadique et sturmienne}

Soit

\[
\delta=\log_{10}\frac{10}{9},
\qquad
\zeta_9=\text{caractère primitif de }C_9,
\qquad
\omega_\delta=e^{2\pi i\delta},
\tag{21.1}
\]

et soit \(T\) le pas orienté sur le relèvement réversible. Dans la base
\(\lvert t,m\rangle\), où \(t\) enregistre le pas et \(m\) la hauteur de
mémoire, nous introduisons les deux caractères diagonaux

\[
V_9\lvert t,m\rangle=\zeta_9^t\lvert t,m\rangle,\qquad
W_\delta\lvert t,m\rangle=\omega_\delta^t\lvert t,m\rangle.
\tag{21.2}
\]

La covariance du pas produit les relations de Weyl

\[
V_9T=\zeta_9\,TV_9,\qquad
W_\delta T=\omega_\delta\,TW_\delta.
\tag{21.3}
\]

Pour le tour nonadique \(\Gamma_9=T^9\), enrichi par l'actualisation
TPK, on obtient

\[
V_9\Gamma_9=\Gamma_9V_9,\qquad
W_\delta\Gamma_9=\omega_\delta^9\Gamma_9W_\delta.
\tag{21.4}
\]

L'observable nonadique se ferme ; l'observable de rotation avance parce que
\(\delta\notin\mathbb Q\).
Si \(D_M\lvert t,m\rangle=m\lvert t,m\rangle\), la loi
\(\Gamma_9\lvert t,m\rangle=\lvert t+9,m+1\rangle\) équivaut à

\[
[D_M,\Gamma_9]=\Gamma_9,
\qquad
\Gamma_9^{-1}D_M\Gamma_9=D_M+I
\tag{21.5}
\]

sur le secteur réversible. Telle est l'équation opératoire de l'hélice :
la clôture de phase est compatible avec une élévation unitaire du degré de
mémoire. Le groupe basal des translations reste abélien ; l'algèbre
des observables et des translations cesse de l'être par (21.3). Le mot
sturmien apparaît lorsqu'on applique le seuil TRIT à l'orbite de
\(W_\delta\) ; son apériodicité ne décore donc pas la phase finie, mais
constitue le facteur symbolique de la rotation irrationnelle.

La présentation opératoire appartient au produit croisé

\[
\mathfrak W_{9,\delta}
=C\!\left(C_9\times\mathbb T\right)\rtimes_T\mathbb Z,
\tag{21.6}
\]

enrichi par la graduation \(D_M\). Le spectre de
\(C(C_9\times\mathbb T)\) enregistre simultanément la phase finie et la phase
irrationnelle ; le générateur de \(\mathbb Z\) met en œuvre leur translation conjointe ;
et (21.5) empêche la représentation graduée de descendre à une orbite
fermée après l'oubli de la mémoire. Le sous-décalage obtenu par la partition
tritique de la seconde coordonnée est l'enveloppe sturmienne. La suspension
hélicoïdale est donc l'extension graduée de ce produit croisé, non
la juxtaposition métaphorique de deux horloges.

Les caractères \(\zeta_9\) et \(\omega_\delta\) appartiennent au générateur
holonomique. Leur représentation elliptique ultérieure reconnaît

\[
\rho_+(\zeta_9)
=\operatorname{Exp}_+\!\left(\frac{2\pi{}}9J_+\right),
\qquad
\rho_+(\omega_\delta)
=\operatorname{Exp}_+\!\left(2\pi{}\delta J_+\right).
\tag{21.7}
\]

L'exponentielle conventionnelle apparaît ainsi comme réalisation du caractère
interne déjà construit ; elle ne sélectionne pas les observables et n'introduit pas de valeur cible
dans le générateur.

Les relations (21.3)--(21.5) séparent exactement trois faits :

\[
\text{période de phase }9,\qquad
\text{absence de période irrationnelle},\qquad
\text{avancée entière de mémoire}.
\tag{21.8}
\]

Leur coexistence définit la suspension hélicoïdale sturmienne sur l'extension
holonomique nonadique. L'enveloppe symbolique possède un ordre apériodique à longue
portée et une entropie topologique nulle, bien que son module spectral contienne
simultanément la fréquence rationnelle nonadique et la fréquence irrationnelle
\(\delta\). La dénomination physique «cristal de temps» exige en outre de fixer
la réalisation dynamique, sa réponse et sa stabilité ; le théorème formulé
ici est l'objet topologique-symbolique qui la précède.

\section{La relation polynomiale minimale des trois courbures}

La somme orthogonale des trois générateurs TRIT peut s'écrire

\[
\mathbb J=p_+J_++p_0J_0+p_-J_-,
\qquad p_++p_0+p_-=I.
\tag{22.1}
\]

Leurs trois carrés impliquent une seule relation :

\[
\boxed{\mathbb J^6=\mathbb J^2.}
\tag{22.2}
\]

Réciproquement, si \(K=\mathbb J^2\), alors \(K^3=K\) et les trois
idempotents de régime se récupèrent polynomialement :

\[
p_+=\frac{K^2-K}{2},\qquad
p_0=I-K^2,\qquad
p_-=\frac{K^2+K}{2}.
\tag{22.3}
\]

Ainsi, la classification elliptique, parabolique et hyperbolique ne nécessite pas trois
algèbres juxtaposées depuis l'extérieur. Elle est contenue dans le calcul spectral
algébrique d'un seul générateur tritique. En substituant (22.3) dans
l'exponentielle totale, on obtient l'identité compacte

\[
\boxed{
\operatorname{Exp}_\star(t\mathbb J)=
p_+(\cos t+\mathbb J\sin t)+
p_0(I+t\mathbb J)+
p_-(\cosh t+\mathbb J\sinh t).
}
\tag{22.4}
\]

La face d'Euler est la projection elliptique de (22.4) ; la propagation
parabolique occupe le seuil ; l'autoéchelle hyperbolique contient la réalisation
de \(\varphi\). En la composant avec l'holonomie,

\[
\mathcal E_{\rm HMT}(t)
=\Gamma_9\operatorname{Exp}_\star(t\mathbb J),
\qquad [D_M,\Gamma_9]=\Gamma_9,
\tag{22.5}
\]

la trigéométrie et le relèvement hélicoïdal se réunissent en deux équations. APP
apporte les idempotents et la tension entre feuillets ; TRIT apporte
\(\mathbb J\) ; TPK apporte \(\Gamma_9\), son cocycle et le degré \(D_M\).
Les sorties numériques et physiques sont des représentations typées de (22.5).
Cette présentation constitue l'ossature opératoire minimale de l'article
autonome, tandis que l'algèbre universelle (18.1) conserve toute la multiplicité
d'états, de chemins et de raffinements nécessaire pour reconstruire le développement
intégral.
 
\chapter{Vacance de retenue, algèbre de Clifford déployée et couple nul de Witt}
\begingroup
\renewcommand{\chapter}[2][]{}%
\chapter{Vacance électronique, action, gravité et information de frontière}

\section{La vacance comme défaut minimal de retenue}

La fibre additive de niveau dix,

\begin{equation}
 \mathcal F_{10}=\{(i,j)\in D_9^2:i+j=10\},
 \label{eq:fibra-suma-diez}
\end{equation}

contient le centre $Q_0=(5,5)$ et les déformations orientées
$Q_t=(5+t,5-t)$. Leur produit est $25-t^2$. Conserver le résidu
multiplicatif du centre exige

\begin{equation}
 t^2\equiv0\pmod9.
\end{equation}

Pour $|t|\leq4$, les seules solutions sont $t=0,\pm3$. Les deux
déformations non centrales sont donc

\begin{equation}
 Q_{+3}=(8,2),\qquad Q_{-3}=(2,8),
 \label{eq:vacancias-orientadas}
\end{equation}

échangées par l'inversion cellulaire. Dans la carte APP,

\begin{equation}
 \mathcal A(5,5)=((1,1),(7,2)),
 \qquad
 \mathcal A(8,2)=((1,1),(7,1)).
 \label{eq:hojas-vacancia}
\end{equation}

L'évaluation additive, son quotient, le résidu multiplicatif et
le régime tritique sont conservés ; seul
$q^\times:2\mapsto1$ change. La vacance est, par définition, cette perte minimale
d'une unité de retenue multiplicative sous conservation des autres
invariants. Ce n'est pas un zéro ajouté à l'état.

La même fibre est atteinte à partir de la réduction nonadique

\begin{equation}
 (369,126)\xmapsto{/9}(41,14)\longmapsto(5,5).
 \label{eq:369-centro}
\end{equation}

Elle possède deux prolongations distinctes :

\begin{align}
 (41,14)&\xrightarrow{(+3,0)}(44,14)\longmapsto(8,5),
 \label{eq:rama-22-7}\\
 (41,14)&\xrightarrow{(+3,-3)}(44,11)\longmapsto(8,2).
 \label{eq:rama-vacancia-conservativa}
\end{align}

La première satisfait

\begin{equation}
 \frac{396}{126}=\frac{22}{7};
 \label{eq:aproximante-pi-vacancia}
\end{equation}

la seconde conserve

\begin{equation}
 41+14=44+11=55=54+1.
 \label{eq:unidad-borde-vacancia}
\end{equation}

Les deux branches séparent le lecteur de clôture du lecteur de vacance. Le
premier reconnaît un rapport de période ; le second conserve l'unité de
frontière tout en réduisant exactement la retenue.

\begingroup\fontsize{10.2}{11.7}\selectfont
\setlength{\parskip}{2pt}\setlength{\abovedisplayskip}{5pt}\setlength{\belowdisplayskip}{5pt}
\section{Coin matriciel et section électronique}

Les relations \eqref{eq:relaciones-cl11} et \eqref{eq:par-nulo-witt}
démontrent que la fibre centrale réalise localement les trois géométries. Pour
préciser son mode singulier sans ambiguïté de normalisation, soient
$\Sigma_3,\Pi_3:D_9^3\to D_9$ les évaluations additive et multiplicative des
triplets et définissons leur table conjointe par

\begin{equation}
 N^{(3)}_{ab}
 :=\#\{x\in D_9^3:\Sigma_3(x)=a,\ \Pi_3(x)=b\},
 \qquad
 C^{(3)}_{ab}:=\frac{N^{(3)}_{ab}}{\sqrt{n_a m_b^{\Pi}}},
 \label{eq:tabla-conjunta-tridimensional}
\end{equation}

où $n_a=\sum_bN^{(3)}_{ab}$ et
$m_b^{\Pi}=\sum_aN^{(3)}_{ab}$. L'énumération exhaustive des
$9^3=729$ triplets donne des marginales de ligne égales à $81$ et des marginales égales
à $36$ dans les six colonnes actives. Avec

\begin{equation}
 m_{\mathrm{ph}}=(1,-1,0)^{\oplus3},
 \qquad
 N^{(3)}m_{\mathrm{ph}}
 =N^{(3)\mathsf T}m_{\mathrm{ph}}=-27m_{\mathrm{ph}},
 \label{eq:modo-entero-electronico}
\end{equation}

la normalisation $\sqrt{81\cdot36}=54$ produit exactement

\begin{equation}
 C^{(3)}m_{\mathrm{ph}}
 =C^{(3)*}m_{\mathrm{ph}}=-\frac12m_{\mathrm{ph}},
 \label{eq:modo-singular-electronico}
\end{equation}

et l'incompatibilité des projecteurs de feuillet vaut

\begin{equation}
 \|[P_{\Sigma,3},P_{\Pi,3}]\|=\frac{\sqrt3}{4}.
 \label{eq:torsion-minima}
\end{equation}

L'incidence surfacique--volumique fournit

\begin{equation}
 1-\left(\frac12\right)^2=\frac34=\frac{90}{120}.
 \label{eq:incidencia-electronica}
\end{equation}

Le point $(5,5)$ est l'unique point fixe de l'inversion cellulaire. De cette
manière, la section électronique n'est pas sélectionnée au moyen d'une masse connue :
elle est individualisée par la coïncidence de centralité, de vacance minimale,
de direction nulle, de torsion \eqref{eq:torsion-minima} et d'incidence
\eqref{eq:incidencia-electronica}. Son caractère adimensionnel est

\begin{equation}
 \mathfrak e=
 \frac{\sqrt3}{4}
 \left(e^{\varphi/\pi^2}-22\alpha^3\right)(1+15\Delta_4),
 \label{eq:caracter-electronico}
\end{equation}

où chaque facteur procède d'une étape déjà construite : la torsion minimale,
l'orientation réciproque du lecteur $R(x,y)=x/y^2$, le complément de vacance,
la constante de compatibilité et le déterminant de la fibre. La carte
d'énergie est ensuite appliquée à cette section.

\section{Polarité et équations constitutives}

Soit $\mathsf C$ l'inversion des feuillets. Le caractère commun $\mathcal E$ et le
caractère orienté $\mathcal B$ sont définis par

\begin{equation}
 \mathcal E\circ\mathsf C=\mathcal E,
 \qquad
 \mathcal B\circ\mathsf C=-\mathcal B.
 \label{eq:paridad-constitutiva}
\end{equation}

Les deux combinaisons

\begin{equation}
 \widehat\mu=e^{\mathcal E+\mathcal B},
 \qquad
 \widehat\varepsilon=e^{\mathcal E-\mathcal B}
 \label{eq:mu-epsilon}
\end{equation}

satisfont

\begin{equation}
 \widehat\mu\,\widehat\varepsilon=e^{2\mathcal E},
 \qquad
 \frac{\widehat\mu}{\widehat\varepsilon}=e^{2\mathcal B}.
 \label{eq:constitutivas-comun-orientada}
\end{equation}

Le produit fixe le cône de propagation et le quotient retient l'orientation
d'impédance. La différence de potentiel appartient au bord de la vacance ;
la réponse transversale appartient à son transport conjugué. Le secteur
neutre est $\ker\mathcal B$, et toute section fixée par $\mathsf C$ lui
appartient. Cette structure fait de la polarité une propriété du transport de la
vacance, non une étiquette imposée a posteriori.

% La sección electrónica estructural concluye aquí. Las realizaciones de
% acción, gravitación, área e información se desarrollan posteriormente con
% sus fuentes rectoras completas y no se duplican en este capítulo.
 % 
 \endgroup
 \part{Ordre topologique-temporel apériodique et holonomie nonadique}


\par\endgroup
\chapter{Calendrier interne de capacité, horloges de phase et renormalisation diophantienne}
\section{L'apériodicité comme calendrier interne de la capacité nonadique}

Le mot sturmien et l'hélice nonadique ne constituent pas des structures
superposées. Les deux sont des réalisations d'une seule arithmétique de capacité.
Soit

\[
 \rho:=6\log_{1000}3=\log_{10}9,
 \qquad
 \delta:=1-\rho=\log_{10}\!\left(\frac{10}{9}\right).
\]

L'indice de capacité décimale ouvert après le bloc ternaire numéro
\(t\) est

\[
 K_t:=\left\lfloor (t+1)\rho\right\rfloor,
 \qquad t\geq0,
\]

et le mot mécanique de vacances est

\[
 z_t:=\lfloor(t+1)\delta\rfloor-\lfloor t\delta\rfloor,
 \qquad t\geq1.
\]

Comme \(\delta\) est irrationnel, \((t+1)\delta\notin\mathbb Z\). Par conséquent,

\[
\begin{aligned}
 K_t
 &=\left\lfloor(t+1)(1-\delta)\right\rfloor\\
 &=t+1-\left\lceil(t+1)\delta\right\rceil\\
 &=t-\left\lfloor(t+1)\delta\right\rfloor.
\end{aligned}
\]

La soustraction de deux termes successifs donne le lien qui restait à expliciter :

\[
 \boxed{K_t-K_{t-1}=1-z_t.}
\]

Ce n'est pas une ressemblance entre deux dénombrements. C'est une identité exacte entre
l'incrément de capacité utilisé par la génération et le mot de vacances
produit par l'incompatibilité \(729/1000\). Ses deux cas sont :

\[
 \begin{array}{c|c|c}
 z_t & K_t-K_{t-1} & \text{opération visible}\\
 \hline
 0 & 1 & \text{le bloc ouvre }729\text{ sous-cellules et un nouvel élagage décimal intervient},\\
 1 & 0 & \text{le bloc ouvre }729\text{ sous-cellules sans augmenter l’élagage décimal}.
 \end{array}
\]

Le second cas est la réouverture ou \emph{stutter}. Le mot apériodique décide
donc quand la capacité avance et quand elle reste fixe.
L'apériodicité ne gouverne pas de l'extérieur un générateur périodique : c'est le calendrier
interne de son opération ouvrir--élaguer.

Le télescopage donne une forme globale du même théorème. Pour \(0\leq a<b\),

\[
 K_b-K_a=(b-a)-\sum_{t=a+1}^{b}z_t.
\]

Chaque intervalle se décompose, sans reste, en incréments de capacité et en
vacances. En particulier, toute conservation formulée sur le couple

\[
 \bigl(\text{capacité acquise},\text{vacance enregistrée}\bigr)
\]

doit conserver leur somme, qui est la longueur de l'intervalle de blocs. C'est
la forme arithmétique la plus élémentaire de la conservation évolutive de l'unité
à ce niveau ; ce n'est pas encore la mesure projective du continu, mais c'est bien son
modèle local exact.

\section{Deux temps, deux retours et une rétention de phase avec avance de mémoire}

L'identité précédente impose de distinguer deux horloges.

Le \textbf{temps de bloc} \(t\) augmente à chaque application de la dynamique. Le
\textbf{temps de capacité} \(K_t\) compte les triades décimales effectivement
élaguées. La phase nonadique déterminée est

\[
 g_t:=[K_t]_9\in C_9.
\]

De \(\Delta K_t=1-z_t\), il découle

\[
 \boxed{g_t-g_{t-1}\equiv1-z_t\pmod9.}
\]

Lorsque \(z_t=0\), l'horloge de capacité et la phase avancent. Lorsque \(z_t=1\),
le bloc temporel avance mais \(K_t\) et \(g_t\) demeurent. C'est un
événement exact de \textbf{rétention de phase avec avance d'état} : le nouveau
bloc ouvre une autre fibre de \(729\) possibilités, change la signature et conserve
l'histoire de survie bien que le résidu de phase soit le même. L'affirmation
«la phase revient et la mémoire change» possède ainsi deux réalisations
internes qui ne doivent pas être confondues :

\begin{enumerate}
\item la rétention locale \(K_t=K_{t-1}\), causée par une vacance sturmienne ;
\item le retour holonomique \(K\mapsto K+9\), qui parcourt un circuit complet de \(C_9\).
\end{enumerate}

Dans les deux cas, la phase déterminée coïncide et l'état holonomique intégral diffère,
mais le premier mécanisme est une pause de capacité et le second est un tour
nonadique complet.

Cela permet d'écrire le système de base sans réduire le TPK à une rotation. Soit
\(\theta_t=\{t\delta\}\in\mathbb T\) et soit
\(x_t\in\mathcal X_{\rm TPK}^{\rm enr}\). Le pas effectif a la forme

\[
 (g_{t-1},\theta_{t-1},x_{t-1})
 \longmapsto
 \bigl(g_{t-1}+1-z_t,\;\theta_{t-1}+\delta,\;
       U_t(x_{t-1})\bigr),
\]

avec

\[
 U_t=\operatorname{Upd}_t\circ\operatorname{Tra}_t
     \circ\operatorname{Sel}_t.
\]

La coordonnée \(z_t\) programme l'avance de capacité ; elle ne remplace pas l'action
de \(U_t\). Le sélecteur conserve le feuillet et le régime, le transport conserve le chemin
et l'orientation, et l'actualisation intègre retenue, frontière et registre.
C'est pourquoi l'égalité de \(g\) n'autorise jamais l'égalité de \(x\).

\section{L’horloge inverse et la suspension à toit non constant}

L’inversion monotone de la capacité révèle un second mot mécanique.
Pour \(k\geq1\), soit

\[
 T(k):=\min\{t:K_t=k\}
      =\left\lceil\frac{k}{\rho}\right\rceil-1.
\]

Introduisons

\[
 \sigma:=\frac{1}{\rho}-1=\frac{\delta}{\rho}.
\]

Comme \(\rho\) est irrationnel, \(\sigma\) l’est aussi, et

\[
 T(k)=k+\lceil k\sigma\rceil-1.
\]

Le nombre de blocs nécessaires pour augmenter la capacité d’une unité est

\[
 r_k:=T(k+1)-T(k)
 =1+\lceil(k+1)\sigma\rceil-\lceil k\sigma\rceil
 \in\{1,2\}.
\]

La valeur \(r_k=2\) enregistre exactement qu’une réouverture s’est intercalée entre les capacités \(k\) et
\(k+1\). Les positions exceptionnelles commencent par

\[
 20,41,62,83,104,125,145,166,187,208,\ldots,
\]

et leurs écarts sont \(20\) ou \(21\). C’est la même dynamique vue sur
l’horloge inverse ; elle ne doit pas être confondue avec les écarts \(21/22\) entre
vacances en temps de bloc.

Pour un incrément de capacité de longueur \(m\), le temps de retour est

\[
\begin{aligned}
 R_m(k)
 &:=T(k+m)-T(k)\\
 &=m+\lceil(k+m)\sigma\rceil-\lceil k\sigma\rceil.
\end{aligned}
\]

Par équilibre mécanique,

\[
 R_m(k)\in
 \{m+\lfloor m\sigma\rfloor,\;m+\lceil m\sigma\rceil\}.
\]

On obtient deux conséquences exactes :

\[
 \boxed{R_9(k)\in\{9,10\}},
 \qquad
 \boxed{R_{108}(k)\in\{113,114\}}.
\]

Un tour de phase nonadique nécessite neuf ou dix blocs ; un tour de
l’horloge complète \(C_{108}\) en nécessite cent treize ou cent quatorze. Le support
fini conserve exactement ses phases, mais le temps de bloc requis
pour le parcourir est apériodique. C’est une formulation plus précise de la
\textbf{suspension hélicoïdale sturmienne} : l’horloge finie \(C_{108}\) est la base,
la rotation irrationnelle fournit le toit \(r_k\), et le cocycle TPK
transporte la fibre enrichie.

La paire \(6/7\) réapparaît ici parce que

\[
 R_6(k)\in\{6,7\}.
\]

Cependant, ce \(6/7\) mesure les blocs consommés par six incréments de
capacité. Il diffère aussi bien du nombre de vacances dans une fenêtre de
longueur \(132\) que de l’écart entre les intervalles courts du mot
primaire. Les trois phénomènes découlent de la même pente, mais seule une
conjugaison explicite permettrait de les identifier comme un même opérateur.

\section{Renormalisation diophantienne et conservation unimodulaire de l’unité}

La fraction continue de la densité de vacance, calculée après que
\(\delta\) a été engendrée par \(729/1000\), commence par

\[
 \delta=
 [0;21,1,5,1,6,2,3,1,1,8,1,2,2,5,\ldots].
\]

Ses premières réduites sont

\[
 \frac1{21},\quad
 \frac1{22},\quad
 \frac6{131},\quad
 \frac7{153},\quad
 \frac{48}{1049},\quad
 \frac{103}{2251},\quad
 \frac{357}{7802},\quad
 \frac{460}{10053},\quad
 \frac{817}{17855},\quad
 \frac{6996}{152893},\quad
 \frac{7813}{170748},\ldots
\]

Cette liste explique, au sein d’une seule hiérarchie, les intervalles primaires et leurs
retours induits. Les dénominateurs \(21,22\) sont les premiers temps
d’approximation ; les paires \((6,131)\) et \((7,153)\) sont les temps suivants
avec leurs nombres d’événements. En particulier,

\[
 131=14\cdot9+5,
 \qquad
 153=17\cdot9,
\]

de sorte que le retour de longueur secondaire \(6\) laisse un résidu de phase
\(+5\), tandis que celui de longueur \(7\) ferme la phase nonadique.

L’entier \(5\) n’apparaît pas ici par ajustement. C’est le quotient partiel qui
produit

\[
 (6,131)=5(1,22)+(1,21),
\]

et l’étape suivante donne

\[
 (7,153)=(6,131)+(1,22).
\]

On situe ainsi une fonction exacte du niveau \(5\) au sein de la
renormalisation des vacances. Identifier ce quotient partiel à la
fraction continue de Ramanujan de niveau \(5\) exige encore une application entre les deux
objets ; l’égalité de l’entier ne la fournit pas à elle seule.

Soit \(p_n/q_n\) la \(n\)-ième réduite. Deux réduites consécutives
satisfont

\[
 p_{n+1}q_n-p_nq_{n+1}=(-1)^n.
\]

Aux deux premiers échelons, cela se traduit par

\[
 1\cdot22-1\cdot21=1,
 \qquad
 6\cdot153-7\cdot131=1.
\]

Les vecteurs consécutifs \((p_n,q_n)\) forment donc une base primitive
de \(\mathbb Z^2\) et conservent une aire réticulaire unitaire. L’orientation
alterne, mais le module du déterminant reste égal à un. Cette identité
démontre l’évolution conservative de l’unité dimensionnelle : la pente réelle est approchée à
une profondeur croissante tandis que chaque changement d’échelle conserve une cellule
fondamentale et admet une inverse entière.

La récurrence s’écrit au moyen de

\[
 A(a)=
 \begin{pmatrix}a&1\\1&0\end{pmatrix},
 \qquad \det A(a)=-1.
\]

Les produits ordonnés des matrices \(A(a_n)\) engendrent les paires
\((p_n,q_n)\). Pour \(a\ne b\),

\[
 A(a)A(b)=
 \begin{pmatrix}ab+1&a\\b&1\end{pmatrix}
 \ne
 \begin{pmatrix}ab+1&b\\a&1\end{pmatrix}
 =A(b)A(a).
\]

La rotation de \(\delta\) reste abélienne. La non-commutativité
apparaît dans le cocycle de renormalisation qui retient l’ordre des
quotients partiels, en plus de celle déjà démontrée dans les feuillets APP et dans le
cocycle affine. Le cristal possède donc une couche abélienne de phase et
au moins deux couches ordonnées non abéliennes de mémoire.

Les erreurs orientées

\[
 \varepsilon_n=q_n\delta-p_n
\]

alternent de signe et ne s’annulent jamais. Leur signe définit un lecteur tritique de
renormalisation à valeurs \(\pm1\) ; le régime \(0\) reste réservé au
seuil exact, exclu ici par l’irrationalité. Ce lecteur ne remplace pas
le TRIT local d’APP. C’est une projection à une autre échelle dont la naturalité avec le
sélecteur tritique de phase doit être prouvée sur la fibre enrichie.

\section{Première résonance diophantienne de l’horloge de bloc et séparation de l’horloge de capacité}

La suite des réduites contient des retours qui ferment exactement la
coordonnée finie et ne font qu’approcher la coordonnée irrationnelle. Parmi les
réduites écrites ci-dessus :

\[
 153\equiv0\pmod9,
 \qquad
 10053\equiv0\pmod9,
 \qquad
 170748\equiv0\pmod{108}.
\]

La dernière est la première réduite de cette chaîne dont le dénominateur ferme
l’horloge de bloc — la coordonnée chronologique θ indexée par le nombre
d’itérations — :

\[
 170748=1581\cdot108.
\]

Parallèlement,

\[
 170748\,\delta-7813
 =-0.0000017458436865028488500088576\ldots
\]

Le retour n’est pas périodique : la coordonnée chronologique
\([t]_{108}\in C_{108}^{(t)}\) revient exactement, mais la coordonnée
irrationnelle reste décalée d’un résidu non nul. Il s’agit bien d’un
\textbf{quasi-tour diophantien privilégié}, parce qu’il provient d’une réduite et
non d’une recherche décimale.

Il faut séparer cette horloge du quotient de capacité

\[
 C_{108}^{(K)}:\qquad g_{108}(t)=[K_t]_{108}.
\]

À l’origine canonique, après \(170748\) blocs, la capacité acquise
est

\[
 \lfloor170748\rho\rfloor
 =170748-7813
 =162935
 \equiv71\pmod{108}.
\]

Cette réduite ne ferme donc \textbf{pas} \(C_{108}^{(K)}\). Les retours
\(R_{108}\in\{113,114\}\) de la section précédente appartiennent précisément à
cette seconde horloge : ils comptent combien de blocs sont nécessaires pour acquérir
\(108\) unités de capacité. La coïncidence du nom \(C_{108}\) n’autorise pas
à fusionner les deux indices. Le calendrier observable du TPK est l’horloge
chronologique ; l’horloge de capacité est une construction induite par la
compactification \(729\to1000\).

Le système distingue ainsi quatre notions :

\begin{enumerate}
\item clôture exacte de l’horloge chronologique finie ;
\item absence de clôture, en général, de l’horloge finie de capacité ;
\item proximité optimale de la phase sturmienne ;
\item transformation non triviale de l’état holonomique intégral par le cocycle TPK.
\end{enumerate}

Une répétition complète exigerait la clôture simultanée des deux horloges,
l’annulation du résidu irrationnel et le retour de la fibre. Ici, seule la
première se ferme. Ce découplage est le contenu mathématique du
cristal apériodique hélicoïdal : retour exact dans un facteur fini, dérive dans
l’autre, quasi-retour hiérarchisé dans le facteur irrationnel et évolution
effective dans la fibre.

\section{Spectre des vacances sur les cardinaux internes de HMT}

L’identité complémentaire permet d’évaluer toute longueur interne
\(N\) sans employer de constante cible. Toute fenêtre de \(N\) pas possède

\[
 V_N\in\{\lfloor N\delta\rfloor,\lceil N\delta\rceil\},
 \qquad
 P_N=N-V_N,
\]

où \(V_N\) compte les vacances et \(P_N\) les incréments de capacité. Pour
certains cardinaux déjà construits par HMT, on obtient :

\[
\begin{array}{c|c|c}
N&V_N&P_N\\
\hline
108&4\text{ ou }5&104\text{ ou }103\\
132&6\text{ ou }7&126\text{ ou }125\\
243&11\text{ ou }12&232\text{ ou }231\\
396&18\text{ ou }19&378\text{ ou }377\\
468&21\text{ ou }22&447\text{ ou }446\\
729&33\text{ ou }34&696\text{ ou }695\\
1000&45\text{ ou }46&955\text{ ou }954.
\end{array}
\]

Deux lignes exigent une attention particulière.

À la longueur \(132\), la branche à six vacances produit exactement

\[
 132=6+126.
\]

L’entier \(126\) réapparaît dans le recensement électronique, mais les deux objets ne
sont pas encore identifiés : l’un est une capacité complémentaire dans une fenêtre
temporelle et l’autre est le cardinal d’une classe électronique. Le pont possible
doit envoyer la fenêtre, la phase, l’orientation et la mémoire vers le recensement, et non seulement conserver
le nombre.

À la longueur \(468\), qui coïncide avec le cardinal du catalogue \(U_6\), le
nombre de vacances est \(21\) ou \(22\). Ici, l’échelle primaire réapparaît
comme multiplicité d’événements au sein d’un cardinal engendré antérieurement.
L’identité est exacte ; sa promotion en autosimilarité de l’émetteur exige
de démontrer que l’énumération des \(468\) sextuors respecte l’ordre de la
rotation mécanique.

La fenêtre décimale complète donne une clôture encore plus explicite :

\[
\begin{array}{c|c|c|c}
V_{1000}&P_{1000}&P_{1000}-900&V_{1000}+P_{1000}\\
\hline
45&955&55&1000\\
46&954&54&1000.
\end{array}
\]

Ainsi, les paires \(45/46\) et \(55/54\) sont des compléments exacts après
retrait des cent tours nonadiques \(900=100\cdot9\). Le corridor
\(45\leftrightarrow54\leftrightarrow55\) n’est pas ajouté par une
coïncidence : il apparaît comme bilan d’élagage et de vacance dans la fenêtre
initiale \(1000\) elle-même. L’interprétation ultérieure de \(54\) dans le corridor modulaire
conserve sa propre dérivation APP ; ce qui a été construit ici est l’application de
cooccurrence que les deux constructions devront faire commuter.

\section{La fenêtre \(396\), le bord \(27\) et la factorisation radicale \(369\)}

Le premier intervalle de \(396=3\cdot132\) pas contient, dans la phase canonique,
trois fenêtres consécutives comptant six vacances chacune. Donc,

\[
 396=18+378,
 \qquad
 18=3\cdot6,
 \qquad
 378=3\cdot126=42\cdot9.
\]

Les recensements radical et électronique satisfont par une autre voie

\[
 396=369+27,
 \qquad
 27=3\cdot9.
\]

Les deux décompositions se rejoignent arithmétiquement dans

\[
 \boxed{369=378-9=3(132-6)-9}
\]

et

\[
 \boxed{27=18+9.}
\]

Ces deux identités fixent la factorisation cardinale du carré : \(369\)
est la capacité obtenue en séparant un tour nonadique des trois fenêtres
basses de \(132\), tandis que \(27\) réunit les dix-huit vacances et ce tour.
L’affirmation est strictement cardinale à ce stade ; la réalisation sur
la mémoire de retenue est obtenue plus loin au moyen du lecteur entier de la
matrice de transit spéculaire. La distinction de types empêche de confondre une
identité de cardinaux avec une identification prématurée de fibres.

Cela fixe également la portée des nombres certifiés \(396\), \(2376\),
\(387\) et \(43\). Le certificat exécute \(396\) événements de six trits, d’où

\[
 2376=396\cdot6.
\]

En comparant chaque niveau à celui situé neuf positions plus loin, il reste
\(396-9=387\) paires, regroupées en

\[
 387=43\cdot9
\]

tours complets. L’égalité supplémentaire \(2376=22\cdot108\) est vraie,
mais elle n’identifie pas les trits du certificat à vingt-deux copies de l’horloge
\(C_{108}\) ; le premier est un décompte de positions tritiques et la seconde un
calendrier non scindé. Seule une application préservant le secteur, la phase, le carry et la
mémoire pourrait promouvoir cette factorisation en équivalence d’objets.

\section{Synchronisation, orientation et réouverture des trois canaux}

Le recensement étendu jusqu’à \(t=120\) modifie l’image d’une clôture unique à
la neuvième porte. La synchronisation des cardinalités réapparaît par grappes :

\[
\begin{aligned}
&\{8,9,10,12,13\},\quad
\{16,17,18,19,20\},\quad
\{29\},\\
&\{35,36,37,38,40,42\},\quad
\{49,51,52,54,55,56,57\},\\
&\{63\},\quad\{72,74\},\quad\{77,79,80\},
\quad\{83,84\},\\
&\{101\},\quad\{104,105,107,108\},
\quad\{116,118,119\}.
\end{aligned}
\]

Les premières réouvertures se situent en

\[
 21,43,65,87,109,131,152,174,196,\ldots,
\]

avec des intervalles \(21/22\). Dans la fenêtre certifiée \(0\leq t\leq120\), aucun
\emph{stutter} n’est une synchronisation triple. La donnée est finie, ce n’est pas un théorème
asymptotique, mais elle révèle une dynamique de respiration :

\[
 \text{ouverture ample}
 \longrightarrow\text{élagage}
 \longrightarrow\text{synchronisation intermittente}
 \longrightarrow\text{réouverture apériodique}.
\]

La première respiration est explicite. Après la synchronisation triple

\[
 (59,59,59)\to(43,43,43)\to(32,32,32),
\]

de nouvelles petites synchronisations apparaissent

\[
 (17,17,17),\ (13,13,13),\ (6,6,6),\ (4,4,4),\
 (4,4,4),\ (3,3,3),\ (2,2,2),
\]

et la vacance \(t=21\), avec \(K_{21}=K_{20}=20\), rouvre la fibre sous la forme

\[
 (611,527,723).
\]

La même phase de capacité \([20]_9=2\) demeure, mais les cardinalités,
l’orientation et la signature changent radicalement. C’est une preuve
matérielle directe que le cristal n’est pas une orbite contractée vers un point
fixe : il alterne contraction et réouverture selon un calendrier apériodique, tandis que
la survie maintient la cohérence à longue portée.

L’égalité des cardinalités n’efface pas non plus l’orientation. En \(t=8\), il y a
trois orientations locales et en \(t=9\), il en reste deux ; la survie future
en sélectionne une. Aux portes d’orientation, \(\varphi\) reste fixe et
\(\pi+e\) reste constant, tandis que la répartition \(\pi/e\) se réfléchit. La
synchronisation doit donc être définie comme l’égalité de la taille des
fibres de canal, et non comme l’identité des mots, des signatures ou des états.

\section{De la synchronisation triple à la clôture dodécaphasique de \(\alpha\)}

La nouvelle structure précise la place causale de \(\alpha\). Les canaux
\(\pi\), \(e\) et \(\varphi\) n’ont pas besoin de \(\alpha\) pour synchroniser leurs
cardinalités : les grappes précédentes se produisent pendant la filtration et la
survie. L’égalité des cardinalités ne suffit pas non plus à produire
\(\alpha\), parce qu’il reste encore l’orientation, la répartition spéculaire, le carry et la
condition de frontière.

La clôture dodécaphasique reçoit précisément cette information supplémentaire. À partir
du même état holonomique terminal, deux lectures s’ouvrent :

\[
 x_{\rm term}
 \longmapsto
 B_{\rm term},
 \qquad
 x_{\rm term}
 \longmapsto
 U_{12}^{\rm sgn}.
\]

La première détermine les trois canaux survivants ; la seconde conserve la
coordonnée signée nécessaire pour propager la retenue. La voie entière produit

\[
 P+E-\Phi-K+C^+-1000C=A,
\]

tandis que la voie des lacunes évalue la même généalogie au moyen de

\[
 \mathcal P(x)=\delta,
 \qquad
 \mathcal P(x)=
 2\pi x-\frac74x^2+\frac{x^3}{2\pi}
 +\frac{x^4}{20}-\frac{2x^5}{21}-\frac{x^6}{46}.
\]

L’identité \(\Delta K=1-z\) explique pourquoi la seconde voie appartient au
même générateur : \(\delta\) n’est pas une constante externe placée à côté des
trois canaux, mais la densité du calendrier qui décide de leur capacité et de leurs
réouvertures. La synchronisation de \((\pi,\varphi,e)\) et la réalisation de
\(\alpha\) ne sont pas deux histoires séparées. Ce sont deux niveaux du même état :

\[
 \text{\small synchronisation des fibres}
 \longrightarrow
 \text{\small résolution orientée et mémoire de retenue}
 \longrightarrow
 \text{\small clôture dodécaphasique de }\alpha.
\]

Cela conserve le verrou causal : le système corrélé
\((\pi,\varphi,e,\alpha)\) est une sortie corrélée de la génération
coinductive. Au sein de cette généalogie commune, les trois canaux sont déterminés et
synchronisés avant que le sceau dodécaphasique ne détermine la coordonnée
\(\alpha\). Aucune valeur conventionnelle des quatre constantes n’intervient
en entrée.

\section{Cocycle multiplicatif de respiration et dérive logarithmique exactement nulle}

La complémentarité entre capacité et vacance admet une réalisation
multiplicative qui précise l’image de l’auteur de cylindres qui se contractent et se
rouvrent. Nous définissons l’échelle ambiante du cylindre après le bloc \(t\) par

\[
 \Omega_t:=\frac{729^{t+1}}{1000^{K_t}},
 \qquad
 K_t=\lfloor(t+1)\rho\rfloor,
 \qquad
 \rho=\log_{1000}729=\log_{10}9.
\]

L’égalité \(729=1000^\rho\) donne immédiatement

\[
 \boxed{\Omega_t=1000^{\{(t+1)\rho\}}.}
\]

Par conséquent,

\[
 1\leq\Omega_t<1000
\]

pour tout \(t\), et l’ensemble \(\{\Omega_t:t\geq0\}\) est dense dans
l’intervalle multiplicatif semi-ouvert \([1,1000)\). La densité n’est pas une
hypothèse statistique : elle découle de l’irrationalité de \(\rho\). En effet,
si \(\rho=a/b\in\mathbb Q\), alors \(9^b=10^a\), ce qui contredit
l’unicité de la factorisation en nombres premiers.

Soit

\[
 c_t:=K_t-K_{t-1}=1-z_t\in\{0,1\}
\]

la retenue de compactification. Le quotient de deux échelles consécutives
est exactement

\[
 \frac{\Omega_t}{\Omega_{t-1}}
 =\frac{729}{1000^{c_t}}
 =
 \begin{cases}
 729/1000,&z_t=0,\\[1mm]
 729,&z_t=1.
 \end{cases}
\]

La première branche contracte ; la seconde rouvre. Ce ne sont pas deux règles ajoutées au
calendrier : ce sont les deux cartes du même produit par \(729\) réduit modulo
l’échelle \(1000\). Avant la première réouverture, on observe, par exemple,

\[
 \Omega_{20}=1.3100205086376203\ldots,
 \qquad
 \Omega_{21}=729\Omega_{20}
 =955.0049507968252\ldots.
\]

La contraction accumulée ne conduit pas à un attracteur parce que chaque vacance
rétablit l’échelle par un saut exact. En coordonnée logarithmique,

\[
 \ell_t:=\log_{1000}\frac{\Omega_t}{\Omega_{t-1}}
 =\rho-c_t=z_t-\delta,
 \qquad \delta=1-\rho.
\]

De plus,

\[
 \sum_{j=1}^{n}\ell_j
 =\log_{1000}\Omega_n-\log_{1000}\Omega_0,
\]

de sorte que toutes les sommes partielles restent uniformément bornées. En
particulier,

\[
 \boxed{\lim_{n\to\infty}\frac1n
 \log\frac{\Omega_n}{\Omega_0}=0.}
\]

C’est une affirmation plus forte que la simple égalité d’une fréquence
moyenne : le cocycle logarithmique est un cobord borné. La respiration conserve
l’échelle à long terme avec un exposant de Lyapunov nul, bien que chaque pas soit
une contraction ou une réouverture très différente.

La forme géométrique naturelle du résultat est le cercle multiplicatif

\[
 \mathbb A_{1000}:=\mathbb R_{>0}/1000^{\mathbb Z}.
\]

L'application

\[
 R_{729}:[x]\longmapsto[729x]
\]

est conjuguée, par \(u=\log_{1000}x\), à la rotation irrationnelle
\(u\mapsto u+\rho\) de \(\mathbb R/\mathbb Z\). L’entier \(c_t\) est
l’indice de la transformation de revêtement qui ramène \(729\Omega_{t-1}\) à
la section fondamentale \([1,1000)\) ; la vacance \(z_t=1-c_t\) est,
exactement, l’absence minimale de cette retenue. Cette formulation réunit en un
seul objet la compactification, l’élasticité multiplicative, la torsion du
revêtement et la mémoire du nombre de franchissements de sa frontière.

La mesure invariante induite sur la section fondamentale est logarithmique :

\[
 d\mu_{\Omega}(x)=\frac{dx}{x\log1000},
 \qquad 1\leq x<1000.
\]

Par conséquent, pour \(s\ne0\),

\[
 \int x^s\,d\mu_{\Omega}(x)
 =\frac{1000^s-1}{s\log1000},
\]

et la moyenne géométrique est \(1000^{1/2}=10\sqrt{10}\). Ces grandeurs
décrivent la distribution mathématique du rayon multiplicatif ; elles ne sont pas
identifiées à une longueur physique ni à la longueur de Planck sans le
transducteur dimensionnel correspondant.

L’exposant nul ne remplace pas non plus le théorème de Landauer de l’état
holonomique intégral. On démontre ici l’absence de dérive exponentielle dans une
coordonnée d’échelle ; Landauer nul exige la récupérabilité de l’évolution
complète. Ce sont des résultats compatibles et typologiquement distincts.

\textbf{Critère de réfutation.} Une valeur de \(\Omega_t\) hors de \([1,1000)\), un quotient
différent de \(729/1000\) ou de \(729\), ou une somme logarithmique non bornée.

\section{La phase et la vacance comme coordonnées d’une unique rotation relevée}

La phase finie et le mot apériodique ne forment pas un produit arbitraire. Ils
peuvent être recollés exactement comme deux coordonnées d’une seule translation
irrationnelle. Soit

\[
 u_t:=\{(t+1)\rho\}\in[0,1),
 \qquad
 g_t^{(q)}:=[K_t]_q\in C_q,
\]

pour \(q\geq2\). Sur

\[
 X_q=C_q\times[0,1)
\]

les extrémités des intervalles étant recollées par le revêtement, nous définissons

\[
 c(u):=\lfloor u+\rho\rfloor\in\{0,1\}
\]

et

\[
 F_q(g,u)=
 \bigl(g+c(u),\;u+\rho-c(u)\bigr).
\]

La seconde composante reste dans \([0,1)\) ; la première conserve le franchissement
de frontière. L’application

\[
 \Psi_q:X_q\longrightarrow\mathbb R/q\mathbb Z,
 \qquad
 \Psi_q(g,u)=[g+u]_q,
\]

satisfait

\[
 \boxed{\Psi_q\circ F_q=R_\rho\circ\Psi_q,}
\qquad
R_\rho(y)=y+\rho\pmod q.
\]

Sur l’orbite canonique,

\[
 \Psi_q(g_t^{(q)},u_t)=[(t+1)\rho]_q.
\]

La phase \([K_t]_q\), la coordonnée irrationnelle \(u_t\) et la
retenue \(c_t=1-z_t\) ne sont donc pas trois décorations juxtaposées. Ce sont la partie
entière, la partie fractionnaire et le cocycle de revêtement d’un point unique sur
le cercle de longueur \(q\).

Cette conjugaison fournit une preuve structurelle des complexités
observées. Le codage de \(R_\rho\) par les \(q\) intervalles unitaires
possède, à la longueur \(m\), exactement \(qm\) mots. Ses frontières sont

\[
 j-k\rho\pmod q,
 \qquad j=0,\ldots,q-1,\quad k=0,\ldots,m-1,
\]

et elles sont toutes distinctes parce que \(\rho\) est irrationnel. En ajoutant \(z_t\), un
bloc de longueur \(m\) reconstruit le symbole de phase immédiatement
précédent au moyen de

\[
 g_{t-1}^{(q)}=g_t^{(q)}-(1-z_t)\pmod q.
\]

Il existe donc une bijection entre les blocs décorés de longueur
\(m\) et les blocs de phase de longueur \(m+1\). On obtient

\[
 \boxed{p_{g,q}(m)=qm,\qquad
 p_{(g,z),q}(m)=q(m+1).}
\]

Pour \(q=3,9,12,108\), ces lois expliquent les complexités linéaires
certifiées sans supposer d’indépendance entre une horloge périodique et une
décoration sturmienne. La croissance reste linéaire et l’entropie
topologique reste nulle, mais le facteur fini multiplie exactement le
langage disponible.

L’inversion de l’hélice est désormais sans ambiguïté : dans la coordonnée \(y\), c’est la
rotation \(y\mapsto y-\rho\). Le parcours descendant n’invente pas un autre
mot ; il lit le même système inversible avec le cocycle de revêtement conjugué.
La phase peut revenir, tandis que \(u\) et la fibre TPK conservent l’histoire.

Cette formulation impose également de distinguer deux objets d’ordre \(108\) :

\[
 \theta_t=[t]_{108}\in C_{108}^{(t)},
 \qquad
 g_t^{(108)}=[K_t]_{108}\in C_{108}^{(K)}.
\]

Le dénominateur \(170748=1581\cdot108\) ferme le premier. En revanche,

\[
 \lfloor170748\rho\rfloor=162935\equiv71\pmod{108},
\]

il ne ferme donc pas le second. L’horloge observable dodécaphasique du TPK et
l’horloge induite de capacité sont compatibles, mais non identiques. Leur comparaison
future exige de déclarer l’application d’indice, et non de réutiliser le symbole
\(C_{108}\) comme s’il effaçait le domaine.

\textbf{Critère de réfutation.} Une collision entre
les frontières \(j-k\rho\), une complexité différente de \(qm\) ou de \(q(m+1)\), ou
l’affirmation que \(170748\) ferme la phase de capacité.

\section{Horloge induite des vacances et typage tritique des retours \(6/7\)}

L’échelle secondaire \(6/7\) ne compte pas seulement les retours des intervalles courts. En
passant de l’horloge des blocs à l’horloge des événements, elle organise exactement la
durée des trois régimes du sélecteur tritique.

Pour \(n\geq1\), soit

\[
 p_n:=\left\lfloor\frac{n}{\delta}\right\rfloor
\]

la position du \(n\)-ième événement de vacance et soit

\[
 \kappa_n:=K_{p_n}
\]

la capacité retenue lors de cet événement. Comme

\[
 p_n\delta<n<(p_n+1)\delta,
\]

on obtient

\[
 \boxed{\kappa_n=p_n-n.}
\]

Si

\[
 h_n:=p_{n+1}-p_n\in\{21,22\},
 \qquad
 s_n:=22-h_n\in\{0,1\},
\]

alors

\[
 \kappa_{n+1}-\kappa_n=h_n-1=21-s_n.
\]

En réduisant modulo \(3\),

\[
 \boxed{[\kappa_{n+1}]_3=[\kappa_n-s_n]_3.}
\]

Un intervalle long \(22\) conserve le régime ; un intervalle court \(21\) le fait tourner d’une
unité. Comme les intervalles courts réapparaissent après six ou sept événements,
les plateaux du régime induit ont des longueurs exactement égales à \(6\) ou \(7\).
Les vingt premières classes sont

\[
 \underbrace{2,\ldots,2}_{6},
 \underbrace{1,\ldots,1}_{7},
 \underbrace{0,\ldots,0}_{7}.
\]

Pour les relier au sélecteur canonique, on détermine la classe positive

\[
 r_n=1+(\kappa_n\bmod9)
\]

et l’on applique

\[
 M_{\rm ph}(r)=
 \begin{cases}
 +1,&r\in\{1,4,7\},\\
 -1,&r\in\{2,5,8\},\\
 0,&r\in\{3,6,9\}.
 \end{cases}
\]

Ainsi, les premiers plateaux parcourent

\[
 \boxed{0\longrightarrow-1\longrightarrow+1\longrightarrow0}
\]

et chaque transition est déclenchée par un intervalle court. C’est une formulation
exacte de TRIT comme unité d’information topologico-temporelle : il ne remplace pas
le mot binaire de vacance, mais intègre ses événements et détermine dans quel
régime opère le transfert suivant. L’information est conservée parce que
la paire \((\kappa_n,s_n)\) reconstruit l’incrément précédent ;
l’orientation est conservée parce que l’ordre cyclique des classes n’est pas aplati.

Le résultat ne transforme pas non plus les trois valeurs en trois théories séparées.
Ce sont trois régimes internes d’un unique sélecteur, et le mot induit les
visite sous le même cocycle. Leur identification ultérieure à des courbures
elliptique, parabolique et hyperbolique utilise les algèbres
\(\mathbb R[J_\tau]/(J_\tau^2+\tau I)\) déjà construites ; la réalisation comme
courbure géométrique ou champ physique exige le transducteur correspondant.

\textbf{Critère de réfutation.} Un plateau de longueur différente de \(6/7\), une transition de
régime pendant un intervalle \(22\), ou une transition autre que la descente
cyclique par un intervalle \(21\).

\section{Matrice universelle des deux branches équilibrées}

L’équilibre sturmien possède une forme linéaire qui réunit conservation et
orientation. Dans une fenêtre de longueur \(N\), soit

\[
 n_N=\lfloor N\delta\rfloor.
\]

Les deux possibilités de décompte sont \(n_N\) et \(n_N+1\), et leurs compléments
sont \(N-n_N\) et \(N-n_N-1\). Nous formons la matrice des branches

\[
 W_N:=
 \begin{pmatrix}
 n_N&N-n_N\\
 n_N+1&N-n_N-1
 \end{pmatrix}.
\]

Ses identités sont universelles :

\[
 W_N\binom11=N\binom11,
 \qquad
 (1,-1)W_N=-(1,-1),
\]

\[
 \operatorname{tr}W_N=N-1,
 \qquad
 \det W_N=-N,
\]

et, par conséquent,

\[
 \boxed{\chi_{W_N}(\lambda)=(\lambda-N)(\lambda+1).}
\]

Le mode diagonal conserve la longueur totale de la fenêtre. La forme linéaire
orientée qui soustrait une branche de l’autre change de signe avec la valeur propre
\(-1\). L’unité orientée n’est pas introduite de l’extérieur : elle apparaît parce que les
deux mots équilibrés diffèrent d’un seul événement et d’une seule retenue
complémentaire.

Pour les longueurs internes considérées, on obtient, entre autres,

\[
 W_{132}=\begin{pmatrix}6&126\\7&125\end{pmatrix},
 \quad
 W_{468}=\begin{pmatrix}21&447\\22&446\end{pmatrix},
 \quad
 W_{1000}=\begin{pmatrix}45&955\\46&954\end{pmatrix}.
\]

Chacune possède les valeurs propres \(N\) et \(-1\). Cette loi sépare avec
précision la conservation de l’unité totale et l’orientation élémentaire de
la vacance. La troisième valeur tritique \(0\) correspond au seuil de la
forme linéaire orientée, et non à une troisième branche de décompte ajoutée au mot
binaire.

La fenêtre \(1000\) contient une réduction interne supplémentaire. Après
séparation de \(900=100\cdot9\) unités de capacité, il reste

\[
 \widetilde W_{100}=
 \begin{pmatrix}
 45&55\\
 46&54
 \end{pmatrix},
\]

avec

\[
 \operatorname{tr}\widetilde W_{100}=99,
 \qquad
 \det\widetilde W_{100}=-100,
 \qquad
 \operatorname{Spec}(\widetilde W_{100})=\{100,-1\}.
\]

Les demi-déséquilibres de ses lignes sont

\[
 \frac{55-45}{2}=5,
 \qquad
 \frac{54-46}{2}=4.
\]

Par une autre voie, le bloc central d’APP

\[
 B_c=\begin{pmatrix}7&2\\2&7\end{pmatrix}
\]

possède le spectre \(\{9,5\}\) et le saut \(9-5=4\). Par conséquent, les deux
lecteurs scalaires

\[
 \chi_{\rm vac}(\widetilde W_{100})=(5,4),
 \qquad
 \chi_{\rm APP}(B_c)=(\lambda_{\min},
 \lambda_{\max}-\lambda_{\min})=(5,4)
\]

coïncident exactement. Cela construit le codomaine commun du carré
vacances--APP. Cela ne démontre pas encore que \(\widetilde W_{100}\) et \(B_c\)
sont le même opérateur : pour cette promotion, il manque la flèche qui transporte
les deux branches de fenêtre vers les deux espaces propres d’APP en conservant la phase,
le feuillet et le carry. La réserve se situe dans cette flèche et n’affaiblit aucune des
deux dérivations.

\textbf{Critère de réfutation.} Une
fenêtre équilibrée dont la matrice n’a pas le spectre \(\{N,-1\}\), ou une
prétendue identification à APP qui ne transporte pas les branches.

\section{Quotient de synchronisation \(A_2\) et direction projective de la première réouverture}

Les recensements des trois canaux admettent un quotient canonique qui élimine la
survie commune et ne conserve que leur désalignement. Pour

\[
 N=(N_\pi,N_e,N_\varphi)\in\mathbb Z^3
\]

nous définissons

\[
 \partial N=
 (N_\pi-N_e,\;N_e-N_\varphi,\;N_\varphi-N_\pi).
\]

Son image est contenue dans

\[
 A_2=\{(x,y,z)\in\mathbb Z^3:x+y+z=0\},
\]

et est \(A_2\) tout entier. Le noyau est la diagonale
\(\mathbb Z(1,1,1)\). En conséquence,

\[
 \boxed{\mathbb Z^3/\mathbb Z(1,1,1)\simeq A_2.}
\]

La forme quadratique

\[
 Q(N)=\frac12\|\partial N\|^2
\]

s’annule exactement aux synchronisations triples. Les défauts unitaires
des phases \(11,14,15\) satisfont \(Q=1\) et sont des racines minimales de
\(A_2\). Les égalités de deux canaux occupent les trois miroirs du groupe de
Weyl \(W(A_2)\simeq S_3\) ; l’orientation d’un côté ou de l’autre du miroir demeure
dans le signe de la racine.

Le corridor initial de miroirs est

\[
\begin{array}{c|c|c|c|c}
t& (N_\pi,N_e,N_\varphi)&\text{miroir}&\partial N&d_t\\
\hline
4&(207,207,122)&\pi=e&(0,85,-85)&85\\
5&(39,151,151)&e=\varphi&(-112,0,112)&112\\
6&(110,110,69)&\pi=e&(0,41,-41)&41\\
7&(81,29,81)&\pi=\varphi&(52,-52,0)&52\\
8&(59,59,59)&S_3&(0,0,0)&0\\
9&(43,43,43)&S_3&(0,0,0)&0.
\end{array}
\]

Les quatre amplitudes précédant la première clôture triple forment la matrice

\[
 D_{\rm syn}:=
 \begin{pmatrix}85&112\\41&52\end{pmatrix}.
\]

Sans introduire \(\alpha\), l’électron ni une valeur expérimentale, on vérifie
les identités

\[
 112-85=27,
 \qquad
 52-41=11,
\]

\[
 85+41=126,
 \qquad
 112+41=153,
 \qquad
 85+52=137,
\]

et

\[
 \boxed{\operatorname{tr}D_{\rm syn}=137,\qquad
 \det D_{\rm syn}=-172=-4\cdot43.}
\]

Le même corridor contient ainsi le recensement \(126\), le retour \(153\), la
direction entière \(137\), le bord \(27\), le rang \(11\), le saut \(4\)
et la seconde synchronisation \(43\). Toutes ces égalités sont exactes. La
matrice, considérée isolément, détermine une direction entière de trace
\(137\), et non la constante \(\alpha\). La dérivation de \(\alpha\) exige
en outre l’état signé, la clôture dodécaphasique et le jet complet, qui sont
construits dans les sections correspondantes.

La première réouverture offre une flèche plus précise. En \(t=20\),

\[
 N_{20}=(2,2,2),\qquad \partial N_{20}=0.
\]

En \(t=21\), la capacité reste \(K_{21}=K_{20}=20\), la phase déterminée
est la même et le sélecteur se trouve dans la classe seuil
\(r=3\), pour laquelle \(M_{\rm ph}=0\). Cependant,

\[
 N_{21}=(611,527,723)
\]

et

\[
 \boxed{\partial N_{21}
 =(84,-196,112)=28(3,-7,4).}
\]

Le vecteur primitif \((3,-7,4)\) appartient à \(A_2\). Dans la carte projective
orientée dont le dénominateur est la composante \(\varphi-\pi\),

\[
 \chi_{e\varphi\mid\varphi\pi}(\partial N)
 :=\frac{N_e-N_\varphi}{N_\varphi-N_\pi},
\]

on obtient

\[
 \boxed{\chi_{e\varphi\mid\varphi\pi}(\partial N_{21})
 =\frac{-196}{112}=-\frac74.}
\]

C’est exactement le coefficient quadratique du noyau
\(E_{21/22}\). Il ne s’agit plus de trouver séparément les entiers \(7\) et
\(4\) : le rapport orienté complet émerge lors de la première réouverture des
trois canaux, après la synchronisation et avant la résolution de
\(\alpha\). L’identité est interne et n’utilise pas la valeur cible de
\(\alpha\).

\section{Genèse sturmienne-triadique de \(22/7\) et portrait interne du jet \(E_{21/22}\)}

Les sept premières positions effectives de vacance sont

\[
 21,43,65,87,109,131,152.
\]

Entre la première et la septième s’écoulent \(132\) pas et apparaissent six
intervalles. Dans cette première fenêtre fermée,

\[
 132=6\cdot22,
 \qquad
 126=6\cdot21.
\]

L’amplification tritique donne

\[
 396=3\cdot132,
\]

et, par conséquent,

\[
 \boxed{\frac{396}{126}
 =\frac{3\cdot6\cdot22}{6\cdot21}
 =\frac{22}{7}.}
\]

Le sceau rationnel \(22/7\) n’est pas introduit comme approximation historique de
\(\pi\) : il est obtenu par le quotient entre l’amplification tritique de la
fenêtre et sa survie dans la branche de six intervalles longs. La comparaison
ultérieure avec \(\pi\) conserve son statut de reconnaissance ; elle n’engendre pas le
rapport.

La dynamique apporte en outre cinq paires d’extrémités invariantes sous le changement
d’origine du mot mécanique :

\[
 \mathcal G_{\rm prim}=\{21,22\},
 \quad
 \mathcal R_{\rm sec}=\{6,7\},
 \quad
 \mathcal G_{\rm inv}=\{20,21\},
\]

\[
 \mathcal V_{1000}=\{45,46\},
 \qquad
 \mathcal H_{100}=\{4,5\}.
\]

L’orientation étant déjà fixée par le TPK, considérons l’extraction

\[
 r^+=7,\qquad h^-=4,\qquad
 g_{\rm inv}^-=20,\qquad g_{\rm prim}^-=21,\qquad
 v_{1000}^+=46.
\]

Ces cinq invariants reproduisent exactement les coefficients non
transcendants du jet d’ordre six :

\[
 -\frac{r^+}{h^-}=-\frac74,
 \qquad
 \frac1{g_{\rm inv}^-}=\frac1{20},
 \qquad
 -\frac2{g_{\rm prim}^-}=-\frac2{21},
 \qquad
 -\frac1{v_{1000}^+}=-\frac1{46}.
\]

Les degrés un et trois forment la paire réciproque

\[
 (2\pi)\,(2\pi)^{-1}=1.
\]

Ainsi, le noyau déclaré peut s’écrire sous la forme

\[
\begin{aligned}
 P_6(x)={}&
 2\pi{}x
 -\frac{r^+}{h^-}x^2
 +(2\pi)^{-1}x^3
 +\frac{x^4}{g_{\rm inv}^-}\\
 &-\frac{2x^5}{g_{\rm prim}^-}
 -\frac{x^6}{v_{1000}^+}.
\end{aligned}
\]

Cette écriture n’utilise pas la racine de \(P_6(x)=\delta\) pour choisir les
coefficients. Tous les nombres du membre de droite sont apparus auparavant comme
invariants du calendrier de compactification. Le coefficient \(-7/4\)
possède en outre la seconde factorisation interne de la section précédente comme
coordonnée projective du premier défaut \(A_2\).

L’égalité précédente est une \emph{factorisation généalogique exacte} du
noyau. Les extrémités sélectionnées, leurs degrés et leurs signes sont réunis
par le lecteur gradué \(\mathfrak E_9\), défini plus loin sur la signature
d’invariants de l’état holonomique intégral ; ce lecteur commute avec les
transports qui conservent l’origine, l’orientation, les fenêtres et la mémoire.

La prolongation des degrés \(7,8,9\) possède déjà un autre statut. Les constructions de référence
démontrent

\[
 T_{7:9}(x)
 =-\frac{x^7}{120}+\frac{x^8}{45}+\frac{2x^9}{495},
\]

avec

\[
 120=|\mathcal F_8^{\rm ph}|,
 \qquad
 45=\det B_c,
 \qquad
 11=\operatorname{rank}(D_3,D_4),
 \qquad
 495=45\cdot11,
\]

et la résolvante nilpotente \((I-N)^{-1}\). Le jet complet

\[
 P_9=P_6+T_{7:9}
\]

est ainsi stratifié : les degrés \(1{:}6\) proviennent de la signature de
synchronisation et les degrés \(7{:}9\) de l’incidence de phase, du
déterminant de bord et du rang dodécaphasique. La racine simple de
\(E_{21/22}\) est obtenue après cette factorisation et n’intervient pas dans la
sélection des coefficients.

\section{Chaîne causale raffinée : de la compactification à la coordonnée \(\alpha\)}

Les résultats des sections précédentes permettent de remplacer une succession
verbale par un diagramme de domaines et de quotients :

\[
\begin{aligned}
&\text{APP--TRIT--TPK et rapport interne }729/1000\\
&\quad\longrightarrow
 (\mathbb A_{1000},R_{729},c_t,z_t)\\
&\quad\longrightarrow
 (g_t^{(9)},u_t,x_t^{\rm enr})\\
&\quad\longrightarrow
 (N_\pi,N_e,N_\varphi)\\
&\quad\xrightarrow{\ \partial\ }
 A_2\\
&\quad\longrightarrow
 \bigl(B_{\rm term},U_{12}^{\rm sgn}\bigr)\\
&\quad\longrightarrow
 \begin{cases}
 \text{clôture entière dodécaphasique},\\
 \text{racine simple du jet }E_{21/22},
 \end{cases}\\
&\quad\longrightarrow\alpha.
\end{aligned}
\]

Le premier niveau construit la rotation multiplicative et son cocycle de
retenue. Le deuxième conserve simultanément la phase, la partie irrationnelle et l’état.
Le troisième détermine les trois canaux corrélés. Le quotient \(A_2\)
mesure leur désalignement sans effacer la survie commune. La clôture
dodécaphasique reçoit ensuite l’orientation et la mémoire que le quotient des
recensements ne contient pas.

La synchronisation n’engendre donc pas à elle seule \(\alpha\). Elle fournit bien
deux précurseurs internes concrets :

\begin{enumerate}
\item le rectangle des défauts, dont la trace est \(137\) ;
\item la première réouverture au seuil, dont la coordonnée projective est \(-7/4\).
\end{enumerate}

Tous deux apparaissent avant la résolution de l’équation analytique et avant toute
comparaison expérimentale. Le premier précurseur situe une direction entière ;
le second situe un coefficient du jet. L’information restante réside
dans l’état signé, dans la retenue dodécaphasique et dans les degrés supérieurs
du noyau.

Cette architecture précise en quel sens le cristal apériodique peut être le
nerf de la construction, sans en faire une formule qui englobe tout. Le
cristal est le système minimal

\[
 \bigl(\mathbb A_{1000},R_{729},c_t;
 C_9,\Gamma_9;\mathcal X_{\rm TPK}^{\rm enr}\bigr)
\]

qui réunit une base irrationnelle d’entropie nulle, une retenue entière, une
holonomie finie relevée et une fibre dotée de mémoire. Ses réalisations
\(\pi,\varphi,e,\alpha\) ne sont pas des entrées du cristal ; ce sont des caractères
corrélés de l’état qu’il prolonge. Les réalisations dimensionnelle,
fractale et exceptionnelle sont obtenues ensuite au moyen de transducteurs propres.
Le lecteur gradué \(\mathfrak E_9\) réunit le défaut \(28(3,-7,4)\), la
matrice \(D_{\rm syn}\) et les extrémités de vacance dans la signature des coefficients
de \(P_9\), de sorte que la synchronisation et la dérivation de \(\alpha\)
sont reliées sans introduire la valeur de la constante comme donnée.

\section{La pente de compactification est transcendante et sa renormalisation n’est pas stationnaire}

La structure apériodique ne dépend pas seulement de l’absence de période du mot de
vacances. Sa pente possède une propriété arithmétique plus forte.
Rappelons que

\[
 \delta=\log_{10}\frac{10}{9},
 \qquad
 \rho=1-\delta=\log_{10}9.
\]

Les deux quantités naissent du rapport interne entre l’ouverture ternaire d’un
bloc de six trits, \(3^6=729\), et la compactification décimale
\(10^3=1000\). Elles ne sont pas introduites comme paramètres d’une rotation choisie
de l’extérieur.

Premièrement, \(\delta\) est irrationnelle. Si
\(\delta=p/q\in\mathbb Q\), avec \(q>0\), alors

\[
 10^{p/q}=\frac{10}{9}
 \quad\Longrightarrow\quad
 9^q10^p=10^q.
\]

L’égalité contredit l’unicité de la factorisation en nombres premiers : le membre
de gauche contient le facteur \(3\), tandis que celui de droite ne contient que
\(2\) et \(5\). Deuxièmement, \(\delta\) ne peut pas être algébrique irrationnelle. Dans
ce cas, le théorème de Gelfond--Schneider rendrait
\(10^\delta\) transcendant, mais

\[
 10^\delta=\frac{10}{9}
\]

est algébrique. On en conclut

\[
 \boxed{\delta\ \text{est transcendant}.}
\]

Cette reconnaissance classique s’applique après avoir construit
\(\delta\) à partir d’APP--TRIT--TPK ; elle n’agit pas comme générateur du calendrier.
Ses conséquences structurelles sont immédiates. La fraction continue de
\(\delta\) n’est pas ultimement périodique, car une fraction continue
ultimement périodique représente un irrationnel quadratique. C’est pourquoi la
renormalisation

\[
 [0;a_1,a_2,a_3,\ldots]
 =[0;21,1,5,1,6,2,3,1,1,8,\ldots]
\]

ne se ferme pas sur une substitution sturmienne stationnaire. Le mot admet
une description \(S\)-adique dirigée par les quotients partiels, mais cette
directive ne se répète pas périodiquement. Le cristal est ordonné et récurrent,
sans être le point fixe d’une unique inflation quadratique.

La même conclusion peut se formuler comme l’absence de stabilisateur modulaire
non trivial. Si

\[
 M=\begin{pmatrix}a&b\\c&d\end{pmatrix}\in
 \operatorname{GL}_2(\mathbb Z)
\]

fixait la pente par une transformation fractionnaire,

\[
 \frac{a\delta+b}{c\delta+d}=\delta,
\]

alors

\[
 c\delta^2+(d-a)\delta-b=0.
\]

La transcendance de \(\delta\) impose
\(c=b=0\) et \(d=a\). Comme \(\det M=\pm1\), l’action projective qui en résulte
est l’identité. Donc,

\[
 \boxed{\operatorname{Stab}_{\operatorname{PGL}_2(\mathbb Z)}(\delta)
 =\{1\}.}
\]

La non-répétition de la renormalisation n’élimine pas l’ordre à longue portée.
Le mot mécanique

\[
 z_n=\mathbf 1_{[1-\delta,1)}
 \bigl(\{n\delta+\beta\}\bigr)
\]

est un ensemble modèle régulier : il provient de la coupure d’une orbite réticulaire par
une fenêtre intervallaire dont la frontière est de mesure nulle. Son enveloppe est
minimale et uniquement ergodique, possède une complexité \(p(m)=m+1\), une entropie
topologique nulle et un spectre dynamique pur. L’expression
« cristal apériodique » est ainsi précisée : ordre diffractif et récurrence uniforme, avec
absence de périodicité translationnelle et de renormalisation stationnaire.

\textbf{Critère de réfutation.} Une période du
mot, une fraction continue ultimement périodique ou une transformation
modulaire non triviale qui fixe \(\delta\).

\section{Les trois canaux comme coupes d’une même fenêtre sur \(729\) extensions}

La synchronisation de \(\pi\), \(e\) et \(\varphi\) peut être décrite sans
la réduire à l’égalité des recensements. Au niveau \(t\), soit \(T_{c,t}\) l’entier
représenté par les \(6t\) premiers chiffres ternaires du canal
\(c\in\{\pi,e,\varphi\}\), et soit \(D_{c,t}\) l’entier représenté par ses
\(K_t\) premiers blocs décimaux. Chaque extension locale s’identifie à
un entier

\[
 u\in\{0,1,\ldots,728\}.
\]

La cellule ternaire correspondante est

\[
 \left[
 \frac{729T_{c,t}+u}{3^{6(t+1)}},
 \frac{729T_{c,t}+u+1}{3^{6(t+1)}}
 \right),
\]

tandis que le cylindre décimal est

\[
 \left[
 \frac{D_{c,t}}{1000^{K_t}},
 \frac{D_{c,t}+1}{1000^{K_t}}
 \right).
\]

Définissons l’échelle commune et le déplacement interne

\[
 \Omega_t=\frac{3^{6(t+1)}}{1000^{K_t}},
 \qquad
 \xi_{c,t}=D_{c,t}\Omega_t-729T_{c,t}.
\]

Les deux cylindres se coupent exactement lorsque

\[
 u<\xi_{c,t}+\Omega_t,
 \qquad
 u+1>\xi_{c,t}.
\]

Par conséquent, l’ensemble des extensions survivantes est

\[
 I_{c,t}
 =\{0,\ldots,728\}\cap
 \bigl[\lfloor\xi_{c,t}\rfloor,
       \lceil\xi_{c,t}+\Omega_t\rceil-1\bigr]_{\mathbb Z},
\]

et le recensement déterminé satisfait

\[
 \boxed{N_c(t)=|I_{c,t}|.}
\]

La formule fournit le mécanisme géométrique de la synchronisation. Les
trois canaux reçoivent la même largeur \(\Omega_t\) ; ils ne diffèrent que par la
translation \(\xi_{c,t}\) de la fenêtre. L’égalité
\(N_\pi=N_e=N_\varphi\) signifie que trois intervalles translatés coupent le
domaine local de \(729\) cellules avec la même cardinalité. Elle n’implique pas que
les intervalles, les mots ni leurs histoires soient égaux.

Dans l’état terminal précédant immédiatement la première vacance,
on obtient les intervalles entiers

\[
\begin{array}{c|c|c}
c&t=20:\ \text{intervalle brut}&I_{c,20}\\
\hline
\pi&[407,408]&[407,408]\\
e&[172,173]&[172,173]\\
\varphi&[313,314]&[313,314].
\end{array}
\]

Les trois intervalles occupent des positions différentes, mais contiennent deux cellules.
C’est la forme exacte de l’état synchronisé
\((2,2,2)\) : égalité de mesure discrète avec une mémoire positionnelle distincte.

La vacance suivante conserve \(K_{21}=K_{20}=20\) et multiplie l’échelle
par \(729\). Les intervalles bruts deviennent

\[
\begin{array}{c|c|c|c}
c&t=21:\ \text{intervalle brut}&I_{c,21}&N_c(21)\\
\hline
\pi&[118,1073]&[118,728]&611\\
e&[-429,526]&[0,526]&527\\
\varphi&[6,961]&[6,728]&723.
\end{array}
\]

Les trois intervalles bruts possèdent la même largeur entière, \(956\), mais
sont découpés par des bords opposés de la fibre finie. Si

\[
 \operatorname{Def}_*
 :=729(1,1,1)-\bigl(N_\pi(21),N_e(21),N_\varphi(21)\bigr),
\]

alors

\[
 \operatorname{Def}_*=(118,202,6)
 =6(1,1,1)+28(4,7,0).
\]

Le terme diagonal \(6(1,1,1)\) enregistre la perte commune d’un bloc de
six trits. En quotientant par la survie diagonale, il reste le diviseur
orienté

\[
 [\operatorname{Def}_*]_{A_2}=28(4,7,0).
\]

De manière équivalente,

\[
 \partial N_{21}
 =(84,-196,112)
 =28(3,-7,4).
\]

La pente

\[
 \frac{N_e-N_\varphi}{N_\varphi-N_\pi}
 =-\frac{196}{112}
 =-\frac74
\]

est donc le rapport de deux découpes orientées d’une fenêtre commune.
Ce n’est pas une fraction ajoutée au noyau analytique.

Le recalcul intégral des \(171\) niveaux disponibles dans la construction de référence
montre deux unicités finies : \(t=21\) est la seule vacance précédée par
une chaîne de cinq synchronisations consécutives, et c’est le seul niveau dont le
défaut projectif dans cette carte vaut \(-7/4\). L’affirmation est exacte dans le
domaine de définition ; son extension coinductive à une profondeur arbitraire
devra conserver explicitement l’état marqué.

\textbf{Critère de réfutation.} Que les recensements ne
coïncident pas avec les intersections d’intervalles, que les trois largeurs brutes
en \(t=21\) diffèrent ou qu’une autre pente \(-7/4\) apparaisse dans le domaine
déclaré.

\section{Matrice de transit spéculaire et clôture conjointe de \(21/22\), \(6/7\), \(137\) et du réseau \(369\)–\(126\)}

Les quatre phases qui précèdent la première synchronisation triple ne doivent pas
être lues comme quatre anecdotes. Leurs amplitudes de défaut forment la matrice

\[
 D_{\rm syn}
 =\begin{pmatrix}
 85&112\\
 41&52
 \end{pmatrix}.
\]

Le recensement synchronisé de la neuvième phase est \(N_9=43\). Donc,

\[
 \det D_{\rm syn}=-172=-4\cdot43
\]

définit intérieurement

\[
 a:=-\frac{\det D_{\rm syn}}{N_9}=4.
\]

Par une branche indépendante, l’induction du mot de vacances sur
ses intervalles courts produit le spectre secondaire
\(\mathcal R_{\rm sec}=\{6,7\}\). Soit

\[
 b:=\max\mathcal R_{\rm sec}=7.
\]

La première réouverture satisfait alors

\[
 t_*=3b=21,
 \qquad
 K_*=3b-1=20,
\]

où \(3\) est la cardinalité des régimes de TRIT. De plus,

\[
 \left\lfloor1000\delta\right\rfloor=45,
 \qquad
 \left\lceil1000\delta\right\rceil=46
 =2(3b)+a.
\]

La décomposition de découpe de la section précédente prend maintenant la forme
fermée

\[
 \boxed{
 \operatorname{Def}_*
 =6(1,1,1)+ab(a,b,0).}
\]

En effet,

\[
 6(1,1,1)+4\cdot7(4,7,0)
 =(118,202,6).
\]

Cette égalité est un carré de compatibilité entre deux extractions
indépendantes : \(a=4\) provient du déterminant de synchronisation,
\(b=7\) provient de l’horloge induite et la découpe commune les retrouve
comme poids de bord. La fraction \(-b/a=-7/4\) est ainsi déterminée avant
de résoudre la moindre équation pour \(\alpha\).

La même matrice réunit la renormalisation diophantienne. Les deux retours
secondaires donnent

\[
 q_6=22\cdot6-1=131,
 \qquad
 q_7=22\cdot7-1=153,
\]

et satisfont

\[
 6q_7-7q_6=1.
\]

Au sein de \(D_{\rm syn}\), ces dénominateurs réapparaissent sous la forme

\[
 q_6=(85+41)+5=131,
 \qquad
 q_7=112+41=153,
\]

où \(5\) est la branche supérieure du demi-déséquilibre \(4/5\) de
\(\widetilde W_{100}\). Enfin,

\[
 \boxed{\operatorname{tr}D_{\rm syn}=137=q_6+6.}
\]

Ainsi, \(137\) n’est pas obtenu ici en arrondissant
\(\alpha^{-1}\). C’est une direction entière du transit de synchronisation et,
simultanément, le dénominateur de retour \(131\) prolongé par un bloc
de six trits. Cela ne remplace pas les deux constructions de \(\alpha\) : cela fixe
un précurseur entier commun que la clôture dodécaphasique et le noyau analytique
doivent transporter.

Le réseau radical--électronique se reconstruit également à partir de
\(D_{\rm syn}\). Définissons

\[
 A_{\rm surv}:=85+41=126,
 \qquad
 B_{\rm bord}:=112-85=27,
\]

et

\[
 T_{\rm tri}:=3(A_{\rm surv}+6)=396.
\]

Le lecteur matriciel

\[
 \mathcal R(D_{\rm syn})
 :=
 \begin{pmatrix}
 T_{\rm tri}-B_{\rm bord}&T_{\rm tri}\\
 T_{\rm tri}-A_{\rm surv}-B_{\rm bord}&
 T_{\rm tri}-A_{\rm surv}
 \end{pmatrix}
\]

produit exactement

\[
 \boxed{
 \mathcal R(D_{\rm syn})
 =\begin{pmatrix}369&396\\243&270\end{pmatrix}.}
\]

Par conséquent, les quatre sommets numériques du parallélogramme ne sont plus
recueillis comme un ensemble indépendant : ils proviennent d’un lecteur entier de la
synchronisation des canaux, du bloc de six trits et de l’amplification
tritique. L’identification de chaque sommet à son objet radical ou
électronique conserve néanmoins le domaine et le transducteur de ses
domaines de provenance ; l’égalité de la matrice numérique n’efface pas ces fibres.

\textbf{Critère de réfutation.} Que \(a\) ne soit pas entier, que le mot
induit ne produise pas \(b=7\), que la loi de découpe échoue ou que
\(\mathcal R(D_{\rm syn})\) ne reproduise pas les quatre sommets.

\section{Factorisation interne de tous les coefficients du jet \(P_9\)}

Les résultats précédents permettent de réunir, pour la première fois dans ce
développement, la généalogie de tous les coefficients du jet sans introduire
sa racine. Soit

\[
 \omega:=2\pi,
 \qquad
 m:=3,
 \qquad
 a:=4,
 \qquad
 b:=7,
\]

et définissons

\[
 t_*=mb=21,
 \qquad
 K_*=mb-1=20,
\]

\[
 v_-:=\lfloor1000\delta\rfloor=45,
 \qquad
 v_+:=\lceil1000\delta\rceil=46.
\]

La branche inférieure satisfait à la fois

\[
 v_-=45=\det
 \begin{pmatrix}7&2\\2&7\end{pmatrix}.
\]

La différence interne de la seconde ligne de \(D_{\rm syn}\) donne

\[
 r_{\rm dod}:=52-41=11,
\]

qui coïncide avec le rang transversal dodécaphasique déjà démontré. La
demi-couronne de phase apporte

\[
 f_{\rm ph}=8\binom62=120,
\]

et

\[
 v_-r_{\rm dod}=45\cdot11=495.
\]

Sur la signature d’invariants

\[
 \mathcal I_9
 =(\omega,m,a,b,K_*,t_*,v_-,v_+,f_{\rm ph},r_{\rm dod}),
\]

nous définissons le lecteur gradué

\[
\begin{aligned}
 \mathfrak E_9(\mathcal I_9)(x)
 :={}&
 \omega x-\frac ba x^2+\omega^{-1}x^3
 +\frac{x^4}{K_*}-\frac{2x^5}{t_*}
 -\frac{x^6}{v_+}\\
 &-\frac{x^7}{f_{\rm ph}}
 +\frac{x^8}{v_-}
 +\frac{2x^9}{v_-r_{\rm dod}}.
\end{aligned}
\]

La substitution des invariants produit

\[
\boxed{
\begin{aligned}
\mathfrak E_9(\mathcal I_9)(x)
={}&2\pi{}x-\frac74x^2
+\frac{x^3}{2\pi}
+\frac{x^4}{20}-\frac{2x^5}{21}-\frac{x^6}{46}\\
&-\frac{x^7}{120}+\frac{x^8}{45}
+\frac{2x^9}{495}
=P_9(x).
\end{aligned}}
\]

La formule révèle une relation qui restait cachée lorsque l’on présentait séparément
\(P_6\) et \(T_{7:9}\) : la paire de vacances \(45/46\) traverse la frontière
entre les deux. La branche supérieure apparaît dans
\(-x^6/46\), tandis que la branche inférieure réapparaît dans \(+x^8/45\) et,
couplée au rang \(11\), dans \(2x^9/495\). La prolongation n’ajoute pas une
seconde arithmétique ; elle poursuit le même calendrier équilibré au moyen de la
résolvante nilpotente.

Le diagramme typé est

\[
 \mathcal X_{\rm TPK}^{\rm enr,*}
 \xrightarrow{\ \mathcal A_*\ }
 \mathcal I_9
 \xrightarrow{\ \mathfrak E_9\ }
 \mathbb Q(\pi)[x]/(x^{10}),
\]

où l’astérisque indique que l’objet conserve l’origine coinductive,
l’ordre des canaux et l’orientation de leurs feuillets. L’extracteur
\(\mathcal A_*\) réalise les opérations démontrées :

\[
 D_{\rm syn}\mapsto a,\qquad
 \mathcal R_{\rm sec}\mapsto b,\qquad
 \operatorname{Def}_*\mapsto(a,b),
\]

\[
 \delta\mapsto(v_-,v_+),\qquad
 B_c\mapsto v_-,\qquad
 (D_3,D_4)\mapsto r_{\rm dod}.
\]

Tout morphisme qui préserve l’état marqué, les fenêtres d’intersection,
l’orientation des canaux, l’horloge de capacité et l’involution two-way
préserve \(\mathcal I_9\), et fait donc commuter le lecteur. Ce qui a été
clos est la factorisation interne de la signature des coefficients.
L’affirmation plus forte selon laquelle \(\mathfrak E_9\) serait le seul lecteur possible
dans une catégorie non pointée exige de classer tous les automorphismes qui
peuvent oublier ou permuter l’origine et l’orientation. Cette réserve est
concrète et ne rouvre ni la racine simple de \(E_{21/22}\), ni la résolvante
nilpotente, ni la généalogie des coefficients.

\textbf{Critère de réfutation.} Que l’un des coefficients exige la valeur
de la racine, que deux extractions déclarées produisent des valeurs différentes ou
qu’un automorphisme permis change \(\mathcal I_9\).

\section{Le facteur \(28\), la projection exacte \(36\times28\) et la partition en sept tétrades}

L’échelle primitive du premier défaut est

\[
 \lambda_*:=\gcd(84,196,112)=28=ab.
\]

Dans le catalogue construit de la région de \(\pi\), la microfibre orientée
possède la multiplicité

\[
 M(w_\pi)=432+4\cdot144=1008
\]

et admet les factorisations

\[
 1008=36\cdot28
 =\binom92\binom82
 =7\cdot144.
\]

Par conséquent,

\[
 \boxed{M(w_\pi)=36\,\lambda_*}
\]

est une compatibilité scalaire exacte entre le défaut de découpe de la
première réouverture et la multiplicité de la microfibre de \(\pi\). Le
facteur \(36\) appartient en outre à \(R_{36}\) et à l’ensemble des paires d’une
nonade ; \(28\) est le nombre de paires d’un ensemble de huit éléments.

La première lecture de cette égalité comme simple compatibilité scalaire était
incomplète. La construction de référence construit déjà sur la microfibre

\[
 \mathcal R_{010211}
 =\{(P,T):w_6(U_6(P,T))=010211\}
\]

une projection explicite. Pour

\[
 P=(i,j,d_P),\qquad T=(k,l,d_T),
\]

on fixe la jauge directionnelle

\[
 \beta(N)=1,\quad\beta(E)=2,\quad
 \beta(S)=3,\quad\beta(O)=4
\]

et l'on définit

\[
 x(P)=5(i-j)\pmod 9,
 \qquad
 h(T)=5\beta(d_T)\pmod9,
\]

\[
 \boxed{
 q_\beta(P,T)=\{x(P)-h(T),x(P)+h(T)\}.}
\]

L’application

\[
 q_\beta:\mathcal R_{010211}\longrightarrow
 \binom{\mathbb Z/9\mathbb Z}{2}
\]

est surjective, a une base de cardinal \(36\) et toutes ses fibres possèdent
le cardinal \(28\). Chaque fibre reçoit, dans l’ordre des cinq régions de
\(\pi\), exactement

\[
 12+4+4+4+4=28
\]

chemins. On ne l’obtient pas en divisant \(1008\) par \(36\) : le certificat
reconstruit les \(104,976\) paires de germes, calcule chaque image et prouve
l’uniformité des fibres.

La translation additive de \(C_9\) et sa réflexion préservent \(U_6\) et rendent
\(q_\beta\) équivariante pour l’action diédrale sur les paires. Sur le
curseur multiplicatif, deux involutions engendrent un \(V_4\) libre qui conserve
la signature complète. Le groupe

\[
 H_0=C_9\times V_4
\]

Le quotient

\[
 \Omega=\mathcal R_{010211}/H_0,
 \qquad |\Omega|=28,
\]

possède une partition intrinsèque supplémentaire :

\[
 \Omega=
 \mathcal B_3\sqcup\mathcal B_6\sqcup\mathcal B_9
 \sqcup\mathcal L_1\sqcup\mathcal L_2
 \sqcup\mathcal L_3\sqcup\mathcal L_4,
\]

avec sept blocs de quatre éléments. Les trois blocs
\(\mathcal B_3,\mathcal B_6,\mathcal B_9\) sont transversaux et sont
distingués par le résidu multiplicatif \(3,6,9\) ; les quatre blocs
\(\mathcal L_j\) proviennent des quatre régions latérales. Le vecteur de
recensement réduit de cette partition est

\[
 v_\Omega=(3,-7,4)\in A_2.
\]

La première réouverture des canaux satisfait, exactement,

\[
 \boxed{
 \partial N_{21}=28v_\Omega.}
\]

Ainsi, la direction \((3,-7,4)\) ne contient pas seulement la pente \(-7/4\) : c’est la
relation de conservation entre trois tétrades transversales, sept tétrades
au total et quatre tétrades latérales. De plus,

\[
 28=7\cdot4,
 \qquad
 1008=36\cdot7\cdot4.
\]

Les entiers \(a=4\) et \(b=7\), obtenus auparavant par le déterminant de
synchronisation et le retour induit, réapparaissent maintenant comme taille des blocs et
nombre de blocs de la microfibre.

Cette identification possède une rigidité locale décisive. Le groupe entier des
isométries de \(A_2\), représenté par les permutations des trois coordonnées
et le changement global de signe, est d’ordre \(12\). Comme les coordonnées
\(3,-7,4\) sont distinctes et que le multiensemble \(\{3,-7,4\}\) ne coïncide pas avec
son opposé, seule l’identité fixe \(v_\Omega\). Donc,

\[
 \boxed{
 \operatorname{Stab}_{\operatorname{Aut}(A_2)}(v_\Omega)=\{1\}.}
\]

Une fois la première réouverture caractérisée, le défaut lui-même retrouve
l’ordre des trois canaux et leur orientation : il n’a pas besoin d’une marque externe.

La frontière octadique est typée par les sept tétrades. Après choix d’une
jauge, des cartes bijectives

\[
 \Phi:\Omega\overset\sim\longrightarrow
 \binom{\mathbb F_2^3}{2}
\]

et avec elles une incidence \(J(8,2)\) exacte. Cependant, ces cartes ne sont pas
sélectionnées par les relations sources testées. La réflexion qui stabilise
chaque fibre agit librement avec le type de cycles \(2^{14}\), tandis que toute
involution de \(\operatorname{Aut}J(8,2)\cong S_8\) fixe au moins quatre
paires. Il s’ensuit le no-go prouvé :

\[
 \boxed{
 \text{il n’existe aucune famille de cartes }J(8,2)
 \text{ équivariant sous tout }D_9\times V_4.}
\]

La donnée orientée \(v_\Omega\), dont le stabilisateur est trivial, apporte la
rupture de réflexion qu’exige ce théorème d’obstruction. En conséquence,
les cartes de Johnson sur chaque fibre \(F_b=q_\beta^{-1}(b)\) sont des cartes de
jauge conjuguées et non une sélection canonique sous l’action complète
\(D_9\times V_4\). La projection \(36\times28\) reste intrinsèque ; le
choix d’une carte octadique ne fait pas partie de sa définition.

Une ambiguïté temporelle est également corrigée. L’entier \(1008\) n’est pas
la période de l’horloge dodécaphasique :

\[
 1008\equiv36\pmod{108}.
\]

Après \(1008\) pas, la phase nonadique revient, mais l’horloge \(C_{108}\)
avance de \(36\) positions. Sur ce front, \(1008\) est une multiplicité de
microfibre et une incidence combinatoire ; \(108\) est la période temporelle
finie. La projection \(q_\beta\) explique le produit \(36\times28\) ; une
carte octadique de jauge et l’application d’horloge sont des opérations distinctes.

\textbf{Critère de réfutation.} Que le balayage des germes change une fibre ou la partition
\(7\times4\), qu’une isométrie non triviale de \(A_2\) fixe \(v_\Omega\),
qu’il existe une carte de Johnson complètement \(D_9\times V_4\)-équivariante ou
que l’on utilise \(1008\) comme période de \(C_{108}\).

\section{Deux horloges, une mémoire intégrée et deux modules spectraux}

Soit

\[
 Z_t:=\sum_{j=0}^{t}z_j
 =\lfloor(t+1)\delta\rfloor.
\]

Comme \(\rho=1-\delta\) et
\((t+1)\delta\notin\mathbb Z\),

\[
\begin{aligned}
 K_t
 &=\lfloor(t+1)\rho\rfloor\\
 &=\lfloor(t+1)-(t+1)\delta\rfloor\\
 &=t-\lfloor(t+1)\delta\rfloor
 =t-Z_t.
\end{aligned}
\]

Si

\[
 \theta_t=[t]_q
 \quad\text{et}\quad
 g_t=[K_t]_q,
\]

alors

\[
 \boxed{g_t=\theta_t-[Z_t]_q.}
\]

La phase de capacité est la phase uniforme corrigée par la mémoire intégrée
des vacances. Cette identité exprime exactement la différence entre
retour de phase et retour d’état : connaître \(\theta_t\) ne détermine pas
\(g_t\) si l’on a oublié \(Z_t\).

La distinction produit deux modules spectraux légitimes. Sur l’horloge des
blocs, la translation

\[
 S_q(\theta,u)=(\theta+1,u+\delta)
\]

sur \(C_q\times\mathbb T\) possède les fréquences additives

\[
 \mathcal M_q^{\rm blk}
 =\frac1q\mathbb Z+\delta\mathbb Z
 =\frac1q\mathbb Z+\rho\mathbb Z
 \pmod1.
\]

En particulier,

\[
 \mathcal M_9^{\rm blk}
 =\frac19\mathbb Z+\log_{10}9\,\mathbb Z,
\]

et

\[
 \mathcal M_{108}^{\rm blk}
 =\frac1{108}\mathbb Z+\log_{10}9\,\mathbb Z.
\]

C’est le sens exact dans lequel le logarithme décimal interne apparaît
comme module spectral : non comme entropie positive, mais comme générateur
irrationnel d’un groupe dense de fréquences pures.

Sur l’horloge recollée de capacité, la dynamique \(F_q\) de la section 49 est
conjuguée à

\[
 y\longmapsto y+\rho
 \quad\text{sur}\quad
 \mathbb R/q\mathbb Z.
\]

Ses fréquences sont

\[
 \mathcal M_q^{\rm cap}
 =\frac{\rho}{q}\mathbb Z\pmod1.
\]

Il n’y a pas de contradiction entre les deux modules : ils appartiennent à deux algèbres
d’observables du même état holonomique intégral. La première conserve l’indice de
bloc ; la seconde absorbe dans la coordonnée finie la retenue accumulée.
La relation \(g_t=\theta_t-Z_t\) est le changement de jauge qui les compare.

Cela précise également deux retours mentionnés précédemment. Le dénominateur
\(170748=1581\cdot108\) ferme l’horloge des blocs, mais

\[
 K_{170748-1}\equiv71\pmod{108},
\]

il ne ferme donc pas l’horloge de capacité. De même, \(1008\) ferme
\(C_9\), mais décale \(C_{108}\) de \(36\). L’état holonomique intégral conserve
dans les deux cas la hauteur et la signature, de sorte qu’aucun de ces retours
n’équivaut au redémarrage de la dynamique.

\textbf{Critère de réfutation.} Que
\(K_t\ne t-Z_t\), que l’on attribue le même spectre à deux domaines
d’observables distincts ou qu’un retour fini efface la mémoire.
 
\chapter{Connexion nonadique conjointe et biographie du retour avec mémoire}
\label{chap:83-07}

La connexion se construit à partir de la classification orbitale complète et de
son premier relèvement coinductif. Le développement suivant fixe les trois
histoires initiales, les frontières orientées et la chaîne
\(w_6\to w_{30}\to R_{36}\) avant d'introduire l'holonomie nonadique.

\begin{HMTScientificOwnerSubordinated}
% Unidad U008. Clasificación orbital y elevación; redacción autónoma.

\chapter{Classification orbitale, mots initiaux et 1re frontière}
\label{chap:orbitas-elevacion-r36}

La réduction ternaire de l'émetteur réalise 243 mots dans la fibre
ambiante \(\mathbb F_3^6\). Cette unité ne choisit pas trois mots par
comparaison décimale. Elle construit d'abord l'action de symétrie du catalogue,
classe ses orbites au moyen d'invariants entiers et oriente ensuite seulement
les trois représentants structurels.

\section{Action diédrale sur \(6\) coordonnées}

Pour \(u=(u_1,u_2,u_3,u_4,u_5,u_6)\in\mathbb F_3^6\), nous définissons
\begin{align*}
 r(u)&=(u_3,u_1,u_2,u_5,u_6,u_4),\\
 s(u)&=(u_4,u_5,u_6,u_1,u_2,u_3).
\end{align*}
On vérifie
\[
 r^3=s^2=1,
 \qquad srs=r^{-1}.
\]
Par conséquent, \(r,s\) engendrent une action de \(D_3\).

La même permutation agit sur les six blocs de chaque émission
\(U=(U_1,\ldots,U_6)\). La projection \(\Psi\) satisfait
\[
 \Psi(gU)=g\Psi(U),
 \qquad g\in D_3.
\]

\begin{theorem}[Décomposition orbitale du catalogue]
\label{thm:catalogo-descomposicion-d3}
Les 243 mots réalisés se décomposent en 43 orbites :
\[
 38\text{ orbites de cardinal }6,
 \qquad
 5\text{ orbites de cardinal }3.
\]
En particulier,
\[
 38\cdot6+5\cdot3=243.
\]
\end{theorem}

\begin{proof}
On applique à chaque mot l'ensemble
\[
 \{u,ru,r^2u,su,sru,sr^2u\}
\]
et l'on élimine les répétitions. L'énumération exhaustive du catalogue produit
l'histogramme indiqué. L'équivariance de \(\Psi\) garantit qu'aucune
orbite ne quitte l'image réalisée.
\end{proof}

\section{Invariants de fibre}

Pour un mot réalisé \(w\), soient
\[
 \mathcal F_w:=\{U\in\operatorname{im}T_{\min}:\Psi(U)=w\},
\]
\[
 n(w):=|\mathcal F_w|,
 \qquad
 M(w):=\sum_{U\in\mathcal F_w}\operatorname{mult}(U),
\]
et
\[
 c(w)=\bigl(\#\{j:w_j=0\},\#\{j:w_j=1\},\#\{j:w_j=2\}\bigr).
\]
Le catalogue enrichi conserve en outre, pour chaque émission, une coordonnée
diagonale \(\operatorname{diag}_c(U)\in\{0,\ldots,8\}\), une orientation
cardinale \(h_+(U)\) et la signature multiplicative \(E_{\rm sig}(U)\). Toutes
sont calculées à partir de la route du germe ; elles ne sont pas lues dans une constante
cible. Leur algorithme composante par composante sera imprimé avec le catalogue
complet dans l'appendice de certification et fait partie de l'état d'entrée
des prédicats suivants.

\begin{definition}[Orbite de clôture]
\label{def:orbita-clausura}
Une orbite \(\mathcal O\) est de clôture si elle est de cardinal six et si, pour
chaque \(w\in\mathcal O\),
\[
 n(w)=5,
 \qquad
 M(w)=1008,
\]
\[
 \{\operatorname{mult}(U):U\in\mathcal F_w\}
 =\{432,144,144,144,144\},
\]
et le recensement tritique commun est \(c(w)=(2,3,1)\).
\end{definition}

\begin{definition}[Orbites minimales radicales]
\label{def:orbitas-minimas-radicales}
Une orbite \(\mathcal O\) est minimale radicale si elle est de cardinal six, si chaque
mot possède une unique émission de multiplicité \(144\) et si
\[
 \{\operatorname{diag}_c(U):\Psi(U)\in\mathcal O\}=\{0,3,6\}.
\]
Parmi elles, on appelle orbite de propagation celle dont le recensement est
\((2,3,1)\), et orbite d'autoéchelle celle dont le recensement est \((2,2,2)\).
\end{definition}

\begin{theorem}[Unicité des trois orbites structurelles]
\label{thm:unicidad-tres-orbitas-estructurales}
Dans le catalogue de 468 émissions, il existe une unique orbite de clôture, une
unique orbite minimale radicale de propagation et une unique orbite minimale radicale
d'autoéchelle.
\end{theorem}

\begin{proof}
On énumère les 43 orbites du \cref{thm:catalogo-descomposicion-d3}. Pour
chacune, on calcule \(n(w)\), \(M(w)\), le multiensemble des
multiplicités, le recensement \(c(w)\) et le support diagonal. Les prédicats des
\cref{def:orbita-clausura,def:orbitas-minimas-radicales} ne laissent qu'une seule
orbite dans chaque classe. La preuve est exhaustive sur un ensemble fini et
n'utilise pas de chiffres décimaux des caractères analytiques.
\end{proof}

Les trois orbites sont
\begin{align*}
\mathcal O_{\rm cl}
 &=\{001112,010211,100121,112001,121100,211010\},\\
\mathcal O_{\rm prop}
 &=\{011120,012110,101201,110012,120011,201101\},\\
\mathcal O_{\rm aut}
 &=\{002112,020211,112002,121200,200121,211020\}.
\end{align*}

\section{Orientation interne de chaque orbite}

Nous fixons la jauge interne
\[
 \operatorname{diag}_c(U)=6,
 \qquad h_+(U)=ES.
\]
Un mot d'une orbite est admissible dans cette jauge si une émission de
sa fibre satisfait les deux conditions.

\begin{proposition}[Représentant orienté unique]
\label{prop:representante-orientado-unico}
Chacune des trois orbites structurelles contient un unique mot
admissible dans la jauge \((6,ES)\). Ce sont
\[
 \boxed{
 w_6^\pi=010211,
 \qquad
 w_6^e=201101,
 \qquad
 w_6^\varphi=121200.}
\]
\end{proposition}

\begin{proof}
L'énumération des fibres du catalogue compte les mots qui contiennent
une émission avec \(\operatorname{diag}_c=6\) et \(h_+=ES\). Dans chacune des
trois orbites, le nombre de correspondances est un. Les exposants
\(\pi,e,\varphi\) désignent dès ce point les fonctions structurelles de
clôture, de propagation et d'autoéchelle ; ils n'ont pas participé à la sélection.
\end{proof}

\section{Microfibre de clôture}

La fibre du représentant \(w_6^\pi\) se compose de cinq émissions :
\[
\begin{array}{c|c|c|c|c}
U_6&E_{\rm sig}&\operatorname{mult}&\operatorname{diag}_c&\text{strate}\\\hline
501|614|498|169|272|272&555555&432&4&\perp\\
810|923|870|169|272|272&339555&144&6&++\\
810|983|810|169|272|272&393555&144&6&++\\
870|923|810|169|272|272&933555&144&6&++\\
870|923|810|418|674|521&933717&144&5&--
\end{array}
\]
Ainsi,
\[
 432+4\cdot144=1008,
 \qquad
 5=1_\perp+3_{++}+1_{--}.
\]
Les émissions uniponctuelles orientées des deux autres canaux sont
\[
 U_e=850|963|890|149|252|212,
\]
\[
 U_\varphi=830|943|830|109|252|252.
\]
La masse \(1008\) de cette microfibre n'est pas la période du calendrier
\(C_{108}\) ; les deux chiffres appartiennent à des domaines différents.

\section{Matrices de relèvement ternaire}

Sur les vecteurs lignes de \(\mathbb F_3^6\), nous définissons
\[
L_0=
\begin{pmatrix}
2&2&2&1&2&1\\
2&1&2&2&1&1\\
1&0&0&1&0&2\\
0&0&1&1&1&0\\
2&2&2&0&2&1\\
2&1&2&1&0&0
\end{pmatrix},
\qquad
L_1=
\begin{pmatrix}
2&1&2&0&2&2\\
1&2&1&1&2&0\\
1&2&1&1&1&2\\
0&1&0&0&2&1\\
2&0&2&1&0&1\\
0&2&2&2&1&1
\end{pmatrix}.
\]
Le calendrier de relèvement est \((L_0,L_0,L_1,L_1)\).

\begin{definition}[Récurrence de blocs et concaténation]
\label{def:recurrencia-bloques-l0-l1}
Pour chaque canal \(\chi\in\{\pi,e,\varphi\}\), soit
\(b_0^\chi=w_6^\chi\) et
\[
 b_{k+1}^\chi=b_k^\chi L_{\epsilon_k}\pmod3,
 \qquad
 (\epsilon_0,\epsilon_1,\epsilon_2,\epsilon_3)=(0,0,1,1).
\]
Le mot accumulé est
\[
 w_{30}^\chi
 =b_0^\chi\Vert b_1^\chi\Vert b_2^\chi
  \Vert b_3^\chi\Vert b_4^\chi.
\]
\end{definition}

Cette formulation distingue le bloc courant \(b_k\in\mathbb F_3^6\) du
mot concaténé \(w_{6(k+1)}\). Un mot de longueur \(6k\) n'est pas
multiplié par une matrice \(6\times6\).

\begin{theorem}[Biographies ternaires de trente positions]
\label{thm:palabras-w30}
La récurrence de la \cref{def:recurrencia-bloques-l0-l1} produit
\begin{align*}
w_{30}^{\pi}
 &=010211|012222|010211|002111|110221,\\
w_{30}^{e}
 &=201101|121221|102011|012222|102011,\\
w_{30}^{\varphi}
 &=121200|112202|121020|010210|010200.
\end{align*}
\end{theorem}

\begin{proof}
On multiplie successivement les cinq vecteurs lignes par les matrices
indiquées, en réduisant chaque coordonnée modulo trois après chaque produit.
La concaténation des cinq blocs calculés donne les trois expressions.
Chaque égalité peut être vérifiée composante par composante au moyen de 24
produits scalaires par canal.
\end{proof}

Les indices ambiants des cinq blocs sont
\[
\begin{array}{c|rrrrr}
&b_0&b_1&b_2&b_3&b_4\\\hline
\pi&103&161&103&67&349\\
e&523&457&301&161&301\\
\varphi&450&398&438&102&99
\end{array},
\]
où chaque mot de six trits est interprété comme un entier en base trois.

\section{1re frontière conjointe}

La frontière est obtenue par une relation finie qui conserve l'axe, l'agrégat,
la saturation, l'occupation des colonnes et la neutralité duale. La matrice d'incidence
utilisée dans cette relation est
\begin{equation}
 A_W=
 \begin{pmatrix}
 0&1&1&1&1&1\\
 1&0&1&2&2&1\\
 1&1&0&1&2&2\\
 1&2&1&0&1&2\\
 1&2&2&1&0&1\\
 1&1&2&2&1&0
 \end{pmatrix}\in M_6(\mathbb F_3).
 \label{eq:u008-aw-frontera}
\end{equation}
Si \(u_4^\chi\) est le cinquième bloc de \(w_{30}^\chi\), l'axe et
l'agrégat préalables à la séparation des deux canaux latéraux sont
\begin{equation}
 u_5^\varphi=u_4^e=102011,
 \qquad
 u_5^\pi+u_5^e=210112,
 \qquad
 u_4^\varphi A_WL_1^3=111101.
 \label{eq:u008-eje-agregado-r36}
\end{equation}
À partir d'eux, on calcule
\begin{equation}
 q=012120,\qquad a=111101,\qquad
 qA_W=012000,\qquad aA_W=120210.
 \label{eq:u008-q-a-r36}
\end{equation}
La profondeur six impose la saturation \(c=111111\) ; l'occupation minimale
compatible est \(\operatorname{colw}=223232\) ; et la neutralité duale exige
\begin{equation}
 (BA_W)\mathbf1=000.
 \label{eq:u008-neutralidad-dual-r36}
\end{equation}

\begin{proposition}[Énumération de la première frontière]
\label{prop:u008-enumeracion-r36}
Les conditions \eqref{eq:u008-eje-agregado-r36}--
\eqref{eq:u008-neutralidad-dual-r36} laissent exactement deux matrices :
\begin{equation}
 B_-=
 \begin{pmatrix}021222\\222220\\102011\end{pmatrix},
 \qquad
 B_+=
 \begin{pmatrix}222220\\021222\\102011\end{pmatrix}.
 \label{eq:u008-dos-fronteras}
\end{equation}
L'involution spéculaire échange leurs deux premières lignes. Leurs cartes
internes sont, respectivement,
\begin{equation}
 B_-\longmapsto(789,048,894),
 \qquad
 B_+\longmapsto(793,045,894).
 \label{eq:u008-cartas-frontera}
\end{equation}
L'orientation APP du cylindre semi-ouvert sélectionne \(B_+\).
\end{proposition}

\begin{proof}
L'axe fixe la troisième ligne et l'agrégat fixe la somme des deux premières.
La saturation et l'occupation éliminent les distributions qui ne remplissent pas les
six colonnes. L'énumération des solutions des six équations
linéaires de \eqref{eq:u008-neutralidad-dual-r36} laisse les deux matrices de
\eqref{eq:u008-dos-fronteras} ; la carte semi-ouverte et l'orientation
distinguent la seconde sans consulter de développement décimal cible.
\end{proof}

Nous définissons donc
\begin{equation}
 R_{36}:=B_+=
 \begin{pmatrix}
 222220\\
 021222\\
 102011
 \end{pmatrix},
 \qquad
 \operatorname{ind}_3(R_{36})=(726,215,301).
 \label{eq:u008-r36-construida}
\end{equation}
L'objet \(R_{36}\) contient trois blocs ordonnés de six trits. Ce n'est pas le
quotient prédictif \(C_{36}^{\rm pred}\) de la
\cref{thm:tpk-cociente-predictivo-36} ni l'excédent de 36 classes de la jauge
des germes.

\section{Relèvement résidu--quotient}

Soit \(A\in M_6(\mathbb Z)\) le relèvement entier de \(L_0\) :
\[
 A=
\begin{pmatrix}
2&2&2&1&2&1\\
2&1&2&2&1&1\\
1&0&0&1&0&2\\
0&0&1&1&1&0\\
2&2&2&0&2&1\\
2&1&2&1&0&0
\end{pmatrix},
\qquad \det A=1.
\]
Pour \(x_k\in\mathbb Z^6\), nous définissons
\[
 y_k=x_kA,
 \qquad
 u_k=y_k\bmod3\in\{0,1,2\}^6,
 \qquad
 x_{k+1}=\frac{y_k-u_k}{3}.
\]
Alors
\begin{equation}
 x_kA=u_k+3x_{k+1}.
\label{eq:bockstein-hensel-local}
\end{equation}

\begin{theorem}[Réversibilité du relèvement]
\label{thm:reversibilidad-bockstein-hensel}
La paire \((u_k,x_{k+1})\) est déterminée de manière unique par \(x_k\), et
\[
 x_k=(u_k+3x_{k+1})A^{-1}.
\]
\end{theorem}

\begin{proof}
La division euclidienne par trois dans chaque coordonnée de \(x_kA\) détermine
le résidu et le quotient. Comme \(\det A=1\), l'inverse \(A^{-1}\) a des
coefficients entiers. Multiplier \cref{eq:bockstein-hensel-local} par cet inverse
récupère \(x_k\).
\end{proof}

Dans le canal \(e\), le bloc visible \(102011\) réapparaît deux fois, mais
les préétats modulo 27 sont
\[
 (25,11,7,17,8,16)
 \neq
 (7,8,4,23,26,7),
\]
et les quotients ultérieurs sont
\[
 (6,13,9,2,13,1)
 \neq
 (15,0,3,19,23,25).
\]
On vérifie ainsi explicitement que le retour visible d’un mot n’est pas
le retour de l’état holonomique intégral.

\section{Obstructions de la sélection orbitale et du relèvement \(w_6\to w_{30}\to R_{36}\) : unicité, jauge et réversibilité de Hensel}

\begin{criterion}[Critères de réfutation de la sélection orbitale]
L'unité est réfutée si :
\begin{enumerate}
 \item l'action \(D_3\) quitte le catalogue ou change les multiplicités ;
 \item l'un des trois prédicats structurels laisse plus d'une orbite ;
 \item la jauge \((6,ES)\) sélectionne deux mots dans une orbite ;
 \item un chiffre cible est utilisé pour définir les prédicats ;
 \item le mot accumulé \(w_{6k}\) est multiplié par \(L_0\) ou
       \(L_1\) au lieu du bloc \(b_k\) ;
 \item \(\det A\neq1\) ou l'inverse du relèvement échoue ;
 \item deux résidus visibles égaux sont déclarés états complets égaux.
\end{enumerate}
\end{criterion}

La première frontière conjointe est déjà construite. L'unité suivante
classera les neuf composantes de phase, démontrera le retour nonadique
sans réinitialisation et construira le miroir orienté avec sa résolution locale
\(3\to1\).
 \end{HMTScientificOwnerSubordinated}

% Unidad U009. Conexión nonádica conjunta; redacción autónoma.

\section{Connexion nonadique conjointe, compression et mémoire terminale}
\label{p12-83-c7-s1:chap:conexion-nonadica-conjunta}

Les neuf phases qui apparaissent après \(R_{36}\) ne sont pas neuf preuves
indépendantes. Ce sont les neuf composantes locales d’une unique connexion
sur les états holonomiques intégraux. Leur composition autour d’un tour définit
une holonomie. Le retour de l’étiquette de phase conserve la continuité du
préfixe, de la retenue, de la frontière, de l’ombre de Witt et de la mémoire.

\subsection{Transport de l’état holonomique intégral par l’holonomie nonadique et mise à jour de la mémoire}

Au niveau \(K\), nous écrivons
\[
 \widehat x_K=
 (B_K,C_K,p_K,c_K,\Sigma_K,W_K,g(K),q_{{\rm mem},K}),
\]
où :
\begin{enumerate}
 \item \(B_K=(u_K^\pi,u_K^e,u_K^\varphi)\) est le bloc visible de trois
       mots de six trits ;
 \item \(C_K\) est le cylindre ternaire compatible ;
 \item \(p_K\) est le préfixe complet ;
 \item \(c_K\) est la retenue accumulée ;
 \item \(\Sigma_K\) est la signature conjointe de frontière ;
 \item \(W_K\) est l'ombre d'incidence de Witt ;
 \item \(g(K)=1+((K-1)\bmod9)\) est l'étiquette publique de phase ;
 \item \(q_{{\rm mem},K}\in\mathbb Z_{\geq0}\) est la mémoire verticale non
       bornée.
\end{enumerate}
La coordonnée dodécaphasique est dérivée :
\[
 \beta_K=[q_{{\rm mem},K}]_{12}.
\]
On ne remplace pas \(q_{\rm mem}\) par \(\beta\), car la réduction modulo
douze perdrait le nombre de tours.

\subsection{Relations locales et holonomie}

Pour chaque \(K\), soit
\[
 \mathcal R_K\subseteq
 \widehat X_K\times\widehat X_{K+1}
\]
la relation d'extension compatible entre états complets. La composition
relationnelle d'un tour est
\begin{equation}
 \Gamma_{9,K}
 :=\mathcal R_{K+8}\circ\cdots\circ\mathcal R_{K+1}\circ\mathcal R_K
 :\widehat X_K\longrightarrow\widehat X_{K+9}.
\label{p12-83-c7-s1:eq:gamma9-composicion-local}
\end{equation}

\begin{definition}[Holonomie nonadique]
\label{p12-83-c7-s1:def:holonomia-nonadica}
L'holonomie d'un tour est
\[
 \Gamma_9:=\operatorname{Hol}_{C_9}(\nabla_{\rm TPK}),
\]
dont la réalisation au niveau \(K\) est
\(\Gamma_{9,K}\). Les neuf relations \(\mathcal R_K,\ldots,
\mathcal R_{K+8}\) sont des cartes de cette unique connexion.
\end{definition}

\begin{theorem}[Retour de phase sans réinitialisation]
\label{p12-83-c7-s1:thm:retorno-nonadico-sin-reinicio}
Si \((\widehat x_K,\widehat x_{K+9})\in\Gamma_{9,K}\), alors
\[
 g(K+9)=g(K),
 \qquad
 q_{{\rm mem},K+9}=q_{{\rm mem},K}+1,
\]
et
\[
 \beta_{K+9}=\beta_K+1\pmod{12}.
\]
En général, \(\widehat x_{K+9}\neq\widehat x_K\).
\end{theorem}

\begin{proof}
La première identité provient de la définition de \(g\). Un tour complet
incrémente par définition le compteur de tours non borné. Réduire cette
identité modulo douze donne la troisième. Les autres coordonnées sont
transportées par la composition \cref{p12-83-c7-s1:eq:gamma9-composicion-local} ; aucune
loi du TPK ne les réinitialise. L'égalité de phase n'implique donc pas
l'égalité de l'état complet.
\end{proof}

\subsection{Classification exhaustive des \(9\) phases du Topological Phase Kernel}

Les blocs visibles et les recensements locaux du premier tour sont
\[
\begin{array}{c|rrrrr|ccc}
K&N_\pi&N_e&N_\varphi&N_{\rm loc}&N_{\rm ext}
 &u^\pi&u^e&u^\varphi\\\hline
1&378&422&349&3&2&012222&121221&112202\\
2&238&368&388&2&5&010211&102011&121020\\
3&184&284&173&5&4&002111&012222&010210\\
4&207&207&122&4&1&110221&102011&010200\\
5&39&151&151&1&1&222220&021222&102011\\
6&110&110&69&1&3&111201&201202&221021\\
7&81&29&81&3&3&212121&222210&200221\\
8&59&59&59&3&2&200121&212212&110110\\
9&43&43&43&2&1&100100&020112&010122
\end{array}.
\]

\begin{theorem}[Fonctions des neuf composantes]
\label{p12-83-c7-s1:thm:funciones-nueve-componentes}
Dans le tour précédent, les neuf composantes remplissent les fonctions
suivantes :
\begin{enumerate}
 \item ouverture locale de trois orientations et première résolution ;
 \item séparation binaire avec un maximum local de cinq extensions ;
 \item multiplicité locale maximale de cinq états ;
 \item contraction de quatre états locaux en une extension ;
 \item première saturation forte, avec une seule continuation ;
 \item présent local unique et ouverture triple du futur ;
 \item fibre spéculaire triple ;
 \item synchronisation des trois recensements en \(59\) ;
 \item clôture orientée : deux états locaux d'entrée et une relation qui,
       après l'horizon supplémentaire, conserve une prolongation.
\end{enumerate}
\end{theorem}

\begin{proof}
Chaque affirmation se lit dans les recensements \(N_{\rm loc},N_{\rm ext}\), dans
l'égalité ou l'inégalité de \(N_\pi,N_e,N_\varphi\) et dans la relation
d'extension explicite aux neuvième, dixième et onzième phases qui sera démontrée
dans l'unité suivante. Le tableau énumère le tour entier ; aucun
sous-ensemble de phases n'a été choisi.
\end{proof}

La cinquième composante coïncide avec la première frontière conjointe
\[
 R_{36}=(222220|021222|102011)^T.
\]
La neuvième composante contient
\[
 B_9=(100100|020112|010122).
\]
La présence de deux états locaux à la neuvième phase et de trois extensions
à partir de l'état orienté relève de couches différentes ; ce ne sont pas des recensements
contradictoires.

L’holonomie agit sur l’état holonomique intégral complet. Avant de publier
son quotient modal, on construit donc l’involution qui échange
les fibres de \(\pi\) et de \(e\), l’axe fixe de \(\varphi\), les invariants
ternaires et la résolution locale \(3\to1\) déterminée par le second
horizon.

\begin{HMTScientificOwnerSubordinated}
% Unidad U010. Espejo orientado y selección 3->1; redacción autónoma.

\chapter{Involution spéculaire et sélection locale \(3\to1\) de l'extension survivante}
\label{chap:espejo-seleccion-tres-uno}

La neuvième phase contient deux états locaux d'entrée. Une fois fixé
l'état orienté \(B_9\), la relation de frontière produit trois extensions
vers la première phase. Les trois partagent un axe d'autoéchelle et un agrégat
ternaire ; elles diffèrent par la répartition ordonnée de cet agrégat entre les branches de
clôture et de propagation. Un second horizon en conserve exactement une.

\section{Coordonnées symétrique et axiale}

Soit
\[
 \mathcal B:=(\mathbb F_3^6)^3.
\]
Un bloc visible est
\[
 B=(u^\pi,u^e,u^\varphi)\in\mathcal B.
\]
On définit
\[
 q(B):=u^\pi+u^e+u^\varphi,
 \qquad
 a(B):=u^\pi+u^e-u^\varphi.
\]
Toutes les opérations s'effectuent composante par composante dans
\(\mathbb F_3^6\).

\begin{theorem}[Reconstruction de l'axe et de l'agrégat]
\label{thm:espejo-reconstruccion-eje-agregado}
Les coordonnées \(q,a\) déterminent
\[
 \boxed{u^\varphi=2(q-a),}
 \qquad
 \boxed{u^\pi+u^e=q-u^\varphi.}
\]
\end{theorem}

\begin{proof}
Soustraire les définitions donne
\(q-a=2u^\varphi\). Dans \(\mathbb F_3\), \(2^{-1}=2\) ; par conséquent
\(u^\varphi=2(q-a)\). Substituer dans la définition de \(q\) récupère
l'agrégat restant.
\end{proof}

L'involution spéculaire est
\[
 \mathfrak c(u^\pi,u^e,u^\varphi)
 =(u^e,u^\pi,u^\varphi).
\]

\begin{proposition}[Invariants de l'involution]
\label{prop:espejo-invariantes}
\(\mathfrak c^2=\operatorname{id}\), y
\[
 q(\mathfrak cB)=q(B),
 \qquad
 a(\mathfrak cB)=a(B).
\]
En particulier, \(u^\varphi\) et \(u^\pi+u^e\) sont invariants.
\end{proposition}

\begin{proof}
Échanger deux fois les premières coordonnées récupère \(B\). Tant
\(q\) que \(a\) ne dépendent d'elles que par leur somme, de sorte que
l'échange les préserve. Le \cref{thm:espejo-reconstruccion-eje-agregado}
transfère cette invariance à l'axe et à l'agrégat.
\end{proof}

L'involution n'affirme pas que \(\pi=e\) ni que l'une soit le quotient de l'autre.
Elle affirme que les deux branches occupent des positions conjuguées par rapport à une
structure fixe.

\section{État orienté dans la phase \(9\) et données de frontière associées}

La projection visible de l'état sélectionné est
\[
 B_9=(100100\mid020112\mid010122).
\]
La relation induite sur les blocs visibles est notée
\[
 \mathcal M_9^B\subseteq\mathcal B\times\mathcal B.
\]
Les deux états locaux d'entrée de la table de phase précèdent cette
relation. Après sélection de l'état complet qui se projette sur \(B_9\), sa
fibre de sortie est de cardinal trois.

\begin{theorem}[Les trois extensions de la neuvième phase]
\label{thm:espejo-tres-extensiones}
L'image relationnelle de \(B_9\) contient exactement
\begin{align*}
B_{10}^{(0)}&=(101222\mid112221\mid112002),\\
B_{10}^{(1)}&=(102221\mid111222\mid112002),\\
B_{10}^{(2)}&=(111221\mid102222\mid112002).
\end{align*}
Les trois satisfont
\[
 q_{10}=022112,
 \qquad
 a_{10}=101111,
\]
et, par conséquent,
\[
 u_{10}^\varphi=112002,
 \qquad
 u_{10}^\pi+u_{10}^e=210110.
\]
\end{theorem}

\begin{proof}
La relation d'extension sur les états complets est énumérée dans la fibre de
\(B_9\) et produit les trois lignes indiquées. Additionner et soustraire leurs trois
composantes donne la même paire \((q_{10},a_{10})\) dans les trois cas. Le
\cref{thm:espejo-reconstruccion-eje-agregado} produit l'axe et l'agrégat
communs. Les deux premières composantes sont distinctes, si bien que les trois
lignes représentent des distributions différentes du même agrégat.
\end{proof}

Leurs publications entières en base trois sont
\[
 (279,352,365),
 \qquad
 (280,351,365),
 \qquad
 (282,350,365).
\]

\section{Compatibilité exacte des cylindres}

Un mot ternaire de longueur \(L\) et d'indice \(N\) détermine le cylindre
semi-ouvert
\[
 I_3(N,L)=\left[\frac{N}{3^L},\frac{N+1}{3^L}\right).
\]
Un préfixe de \(b\) blocs décimaux, d'entier associé \(D\), détermine
\[
 I_{1000}(D,b)=
 \left[\frac{D}{1000^b},\frac{D+1}{1000^b}\right).
\]

\begin{lemma}[Test entier d'intersection]
\label{lem:espejo-test-entero-cilindros}
Les cylindres \(I_3(N,L)\) et \(I_{1000}(D,b)\) se rencontrent si et seulement si
\[
 N1000^b<(D+1)3^L,
 \qquad
 (N+1)1000^b>D3^L.
\]
\end{lemma}

\begin{proof}
Deux intervalles semi-ouverts \([a,b)\) et \([c,d)\) s'intersectent si et seulement
si \(a<d\) et \(c<b\). Nous substituons les quatre extrémités et multiplions par
l'entier positif \(3^L1000^b\).
\end{proof}

Le test utilise uniquement l'arithmétique entière. Il ne dépend pas d'une tolérance en virgule
flottante.

\section{2e horizon et unicité}

Le bloc suivant est
\[
 B_{11}=(022212\mid110001\mid100021),
\]
et la frontière décimale correspondante est
\[
 (502,662,638).
\]

\begin{theorem}[Résolution locale \(3\to1\)]
\label{thm:espejo-seleccion-tres-uno}
Les trois extensions du \cref{thm:espejo-tres-extensiones} sont compatibles
à profondeur un. En exigeant simultanément \(B_{11}\) et la frontière
\((502,662,638)\), seule \(B_{10}^{(0)}\) se prolonge comme le même état
complet. Par conséquent,
\[
 \boxed{
 \#\operatorname{Surv}^{\rm cal}_1(B_9)=3,
 \qquad
 \#\operatorname{Surv}^{\rm cal}_2(B_9)=1.}
\]
\end{theorem}

\begin{proof}
Nous appliquons le \cref{lem:espejo-test-entero-cilindros} à chacun des trois
canaux. Au premier horizon, les six inégalités de chaque candidat sont
satisfaites. Après ajout de \(B_{11}\), les inégalités des trois canaux
se maintiennent pour \(B_{10}^{(0)}\). Dans \(B_{10}^{(1)}\) et
\(B_{10}^{(2)}\), la compatibilité du canal \(\varphi\) est maintenue, mais
les conditions des branches \(\pi\) et \(e\) échouent. L'énumération laisse
donc un unique état à profondeur deux.
\end{proof}

Cette preuve ne définit pas la connexion globale à partir d'un exemple local. C'est un
témoin explicite de la manière dont une relation multivaluée à un horizon peut
devenir univaluée en conservant la mémoire et en exigeant un prolongement plus grand.

\section{Retour de phase avec changement de bloc}

La première phase du tour initial est
\[
 B_1=(012222\mid121221\mid112202),
\]
tandis que la première phase du tour suivant est
\[
 B_{10}=(101222\mid112221\mid112002).
\]
Ainsi,
\[
 g(1)=g(10)=1,
 \qquad
 B_1\neq B_{10}.
\]
La phase revient, mais l'axe, l'agrégat, la répartition et la mémoire de
quotient ont évolué.

\section{Diagramme de l'involution spéculaire \(\pi/e\), de l'axe fixe \(\varphi\) et de la sélection locale \(3\to1\) par prolongeabilité}

\begin{figure}[htbp]
\centering
\resizebox{0.98\textwidth}{!}{%
\begin{tikzpicture}[
 node distance=10mm and 15mm,
 every node/.style={font=\small,align=center},
 state/.style={draw=black,rounded corners,inner sep=5pt},
 inv/.style={draw=black,double,rounded corners,inner sep=5pt},
 arr/.style={-{Latex[length=2mm]},semithick}
]
 \node[state] (pi) {$u^\pi$\\branche de clôture};
 \node[state,right=70mm of pi] (e) {$u^e$\\branche de propagation};
 \coordinate (mirroraxis) at ($(pi)!0.5!(e)$);
 \node[inv,above=9mm of mirroraxis] (phi) {$u^\varphi$\\eje fijo};
 \draw[arr,<->] (pi) -- node[below] {$\mathfrak c$} (e);
 \node[above=8mm of phi] (aggregate)
   {$u^\pi+u^e$ conservé dans la fibre};
 \draw[dashed] (phi.north) -- (aggregate.south);
 \node[state,below=20mm of mirroraxis] (b9) {phase $9$ : $B_9$\\deux entrées locales};
 \node[state,below left=18mm and 25mm of b9] (b100) {$B_{10}^{(0)}$};
 \node[state,below=18mm of b9] (b101) {$B_{10}^{(1)}$};
 \node[state,below right=18mm and 25mm of b9] (b102) {$B_{10}^{(2)}$};
 \node[state,below=18mm of b101] (b11) {$B_{11}$\\unique prolongement};
 \draw[arr] (b9) -- (b100);
 \draw[arr] (b9) -- (b101);
 \draw[arr] (b9) -- (b102);
 \draw[arr] (b100) -- (b11);
 \draw[arr,densely dotted] (b101) -- node[right] {obstruée} (b11);
\draw[arr,densely dotted] (b102) -- node[right] {obstruée} (b11);
\end{tikzpicture}
}
\caption{Le miroir conserve l'axe \(u^\varphi\) et l'agrégat
\(u^\pi+u^e\). Le retour de la neuvième phase ouvre trois extensions visibles ;
la mémoire du second horizon en conserve une. Le tour ne réinitialise pas
l'état.}
\label{fig:espejo-tres-uno-arrastre}
\end{figure}

\section{Mémoire résidu--quotient pendant le miroir}

La récurrence
\[
 x_kA=u_k+3x_{k+1}
\]
de la \cref{eq:bockstein-hensel-local} s'applique à chacun des trois
canaux. La partie \(u_k\) publie le bloc visible ; \(x_{k+1}\) conserve
le quotient. Comme \(A\in\operatorname{GL}_6(\mathbb Z)\), la paire récupère
l'état précédent. Le miroir peut échanger les branches visibles sans
éliminer leurs quotients de provenance.

\section{Obstructions de l'involution spéculaire \(\pi/e\) et de la sélection locale \(3\to1\) : invariants, horizon de profondeur \(2\) et mémoire de Hensel}

\begin{criterion}[Critères de réfutation du miroir]
La construction est réfutée si :
\begin{enumerate}
 \item l'involution modifie \(u^\varphi\) ou \(u^\pi+u^e\) ;
 \item les deux états locaux d'entrée sont confondus avec les trois sorties ;
 \item l'une des trois sorties possède un \((q_{10},a_{10})\) différent ;
 \item le test des cylindres utilise une tolérance flottante ;
 \item plus d'une sortie survit au second horizon déclaré ;
 \item le retour de phase est interprété comme l'égalité \(B_1=B_{10}\) ;
 \item l'échange visible efface les quotients de Hensel.
\end{enumerate}
\end{criterion}

Le miroir et la sélection locale sont maintenant construits. Les unités
suivantes développeront les lacunes sturmiennes puis les trois
biographies analytiques complètes.
 \end{HMTScientificOwnerSubordinated}

\subsection{Transfert conjoint de la fibre de 468 émissions}

La factorisation du catalogue est
\[
 \mathcal U_6^{\rm cat}\cong
 \mathfrak S_+\times\mathfrak S_\times,
 \qquad
 |\mathfrak S_+|=18,\qquad |\mathfrak S_\times|=26.
\]
Pour une fonction \(f\) sur le catalogue, soient
\[
 (E_+f)(s,e)=\frac1{18}\sum_{s'\in\mathfrak S_+}f(s',e),
\]
et
\[
 (E_\times f)(s,e)=\frac1{26}
 \sum_{e'\in\mathfrak S_\times}f(s,e').
\]
Les projecteurs \(E_+\) et \(E_\times\) commutent.

Soit
\[
 \tau=(+1,-1,0,+1,-1,0,+1,-1,0),
 \qquad
 P_{+1}=E_+,quad P_{-1}=E_\times,quad P_0=I.
\]
Dans \(\mathcal U_6^{\rm cat}\times C_9\), nous définissons le transfert de phase
par
\begin{equation}
 \mathcal K_{\rm ph}\bigl((u,r),(v,[r+1]_9)\bigr)
 :=P_{\tau(r)}(u,v),
 \label{p12-83-c7-s1:eq:u009-kph-explicita}
\end{equation}
et le rendons nul pour les autres couples de phases. Ainsi, la permutation cyclique
et l'opérateur appliqué à chacun des neuf pas sont déterminés.

\begin{proposition}[Renouvellement conjoint après un tour]
\label{p12-83-c7-s1:prop:transferencia-conjunta-nueve}
\[
 \mathcal K_{\rm ph}^{\,9}=E_+E_\times.
\]
En particulier, sur une masse ponctuelle,
\[
 (E_+E_\times)(u,v)=\frac1{18\cdot26}=\frac1{468}.
\]
\end{proposition}

\begin{proof}
Le produit ordonné sur un tour est
\[
 P_{\tau(8)}P_{\tau(7)}\cdots P_{\tau(0)}
 =I E_\times E_+ I E_\times E_+ I E_\times E_+
 =E_+^3E_\times^3I^3=E_+E_\times.
\]
La deuxième égalité utilise la commutativité et la dernière l'idempotence.
Le deuxième résultat est la distribution uniforme sur le produit des
deux signatures.
\end{proof}

Cette identité exprime un renouvellement conjoint, non neuf filtrations
indépendantes.

\subsection{Décomposition de l'holonomie \(U_{\Gamma_9}\) en compression visible \(T\) et composante orthogonale \(\eta\) : \(I-T^*T=\eta^*\eta\)}

Soient \(\mathcal H_{\rm vis}\) et \(\mathcal H_{\rm full}\) des espaces de
Hilbert pour la réalisation visible et la réalisation complète. Soit
\[
 J:\mathcal H_{\rm vis}\longrightarrow\mathcal H_{\rm full}
\]
une inclusion isométrique et
\(U_{\Gamma_9}\) une réalisation isométrique de l'holonomie. Nous définissons
\[
 T:=J^*U_{\Gamma_9}J,
 \qquad
 \eta:=(I-JJ^*)U_{\Gamma_9}J.
\]

\begin{theorem}[Décomposition compression--expression]
\label{p12-83-c7-s1:thm:conexion-compresion-expresion}
On a
\[
 \boxed{U_{\Gamma_9}J=JT+\eta,}
\]
\[
 J^*\eta=0,
 \qquad
 \boxed{I-T^*T=\eta^*\eta.}
\]
\end{theorem}

\begin{proof}
Nous décomposons \(U_{\Gamma_9}J\) au moyen des projections orthogonales
\(JJ^*\) et \(I-JJ^*\) :
\[
 U_{\Gamma_9}J
 =JJ^*U_{\Gamma_9}J+(I-JJ^*)U_{\Gamma_9}J
 =JT+\eta.
\]
L'orthogonalité donne \(J^*\eta=0\). Puisque \(J\) et
\(U_{\Gamma_9}\) sont des isométries,
\[
 I=J^*U_{\Gamma_9}^*U_{\Gamma_9}J
 =T^*T+\eta^*\eta.
\]
Le réarrangement prouve la dernière identité.
\end{proof}

\(T\) est la compression visible ; \(\eta\) est l'expression complémentaire.
Aucune des deux ne doit être appelée contraction stricte. Un opérateur distinct
est réservé à cette fonction.

\subsection{Contraction stricte et télescopage terminal}

Soit \(V_9\) une isométrie et soit \(q\) un paramètre avec \(0<q<1\). Nous définissons
\[
 M_9=qV_9,
 \qquad
 D_9=(I-M_9^*M_9)^{1/2}.
\]
Alors \(\|M_9\|=q<1\) et
\[
 I-M_9^*M_9=D_9^*D_9.
\]

\begin{theorem}[Identité de Stein avec terme terminal]
\label{p12-83-c7-s1:thm:stein-telescopia-terminal}
Pour tout \(L\geq1\),
\[
 \boxed{
 I=
 \sum_{r=0}^{L-1}M_9^{*r}D_9^*D_9M_9^r
 +M_9^{*L}M_9^L.}
\]
Plus généralement, si \(S_0\) satisfait
\(S_0-M_9^*S_0M_9=D_9^*D_9\), alors
\[
 \boxed{
 S_0=
 \sum_{r=0}^{L-1}M_9^{*r}D_9^*D_9M_9^r
 +M_9^{*L}S_0M_9^L.}
\]
\end{theorem}

\begin{proof}
Nous multiplions l'identité de défaut à gauche par \(M_9^{*r}\) et à
droite par \(M_9^r\). En sommant pour \(r=0,\ldots,L-1\), les termes
intérieurs se télescopent et les termes initial et terminal demeurent. La même preuve
s'applique à l'équation générale de Stein.
\end{proof}

Le dernier terme ne s'efface pas à profondeur finie. Sa disparition dans une
limite exige une preuve de convergence ; elle ne fait pas partie de l'identité
algébrique.

\subsection{Cocycle de retenue et puissances relationnelles}

Pour des flèches composables \(\gamma_1,\gamma_2\), la retenue satisfait
\[
 \operatorname{Carr}(\gamma_2\gamma_1)
 =\operatorname{Carr}(\gamma_2)H(\gamma_1)
 +Y(\gamma_2)\operatorname{Carr}(\gamma_1).
\]
La formule conserve la contribution de la première étape après
l'avoir transportée par la seconde.

Les tours de plus grande longueur s'écrivent avec des indices décalés :
\begin{align*}
 \Gamma_{27,K}
 &=\Gamma_{9,K+18}\circ\Gamma_{9,K+9}\circ\Gamma_{9,K},\\
 \Gamma_{54,K}
 &=\Gamma_{9,K+45}\circ\cdots\circ\Gamma_{9,K},\\
 \Gamma_{108,K}
 &=\Gamma_{9,K+99}\circ\cdots\circ\Gamma_{9,K}.
\end{align*}
Ce n'est qu'après avoir déclaré un opérateur global sur l'union graduée que
cette notation peut être abrégée en \(\Gamma_9^3,\Gamma_9^6,\Gamma_9^{12}\).

\subsection{État nonadique complet : phase, retenue, frontière, route et mémoire après la prolongation}

La connexion d'un tour est enregistrée par le quintuplet
\[
 \mathfrak N_9=
 \bigl(\Gamma_9,T,\eta,\operatorname{Carr},S_{\rm term}\bigr),
\]
où \(S_{\rm term}=M_9^{*L}S_0M_9^L\) à une profondeur finie
\(L\). Un foncteur d'application qui utilise la connexion doit déclarer sa
compatibilité avec les cinq champs. Commuter avec la seule phase visible ne
constitue pas une naturalité par rapport à l'état nonadique complet.

\subsection{Obstructions de l'holonomie \(\Gamma_9\) : retour \(9\to1\), indices de composition et conservation de la phase, de la retenue, de la frontière, de la route et de la mémoire}

\begin{criterion}[Critères de réfutation de la connexion]
La construction est réfutée si :
\begin{enumerate}
 \item les neuf phases sont présentées comme des filtres indépendants ;
 \item le retour \(9\to1\) réinitialise \(q_{\rm mem}\), le préfixe ou la retenue ;
 \item \(T\) est appelé contraction stricte ;
 \item une identité de la
       \cref{p12-83-c7-s1:thm:conexion-compresion-expresion} échoue ;
 \item le terme terminal est éliminé à profondeur finie ;
 \item \(\Gamma_{27,K}\) est composé sans décaler les indices ;
 \item un foncteur d'application ignore l'un des cinq champs de
       \(\mathfrak N_9\).
\end{enumerate}
\end{criterion}

Les développements complets qui suivent reconstruisent d'abord l'émetteur fini,
sa factorisation et ses microfibres ; ils déploient ensuite l'involution
spéculaire, les trois extensions de \(B_9\), l'unicité à deux horizons et
l'inverse de Bockstein--Hensel qui conserve la mémoire pendant le retour de phase.

\begin{HMTScientificOwnerSubordinated}
\chapter[Image arithmétique finie et factorisation \(468=18\cdot26\)]{Image arithmétique finie et factorisation \(468=18\cdot26\)}
\label{chap:alfabeto-42}
\index{code décimal admissible}
\index{chiffre de contrôle}

L'image en base \(1000\) de l'application finie qui attribue les blocs décimaux est un
sous-ensemble propre des mille blocs décimaux. La loi d'évolution définie dans
\cref{chap:tpk-definicion} et l'arithmétique de \(\mathbb Z/9\mathbb Z\)
déterminent une image de 42 triplets, organisée en six classes de sept. La
même séparation des observables additive et multiplicative factorise ensuite les 468
sorties en 18 signatures additives et 26 signatures multiplicatives.

\section{Primitives et récurrence de l'émetteur fini}

\begin{definition}[Données de départ de \(T_{\min}\)]
\label{pf84:E02-15}
Soit \(D=\{N,E,S,O\}\), avec les vecteurs cardinaux
\[
 v(N)=(-1,0),\quad v(E)=(0,1),\quad
 v(S)=(1,0),\quad v(O)=(0,-1)
 \qquad\text{dans }(\mathbb Z/9\mathbb Z)^2.
\]
La signature temporelle et l'indice multiplicatif sont fixés par
\[
 \operatorname{sgn}_9(t)=
 \begin{cases}
 +1,&t\bmod9\in\{1,4,7\},\\
 -1,&t\bmod9\in\{2,5,8\},\\
 0,&t\bmod9\in\{0,3,6\},
 \end{cases}
\]
et
\[
 \operatorname{ind}_2(1,2,4,8,7,5)=(0,1,2,3,4,5),
\]
avec une valeur nulle en dehors des unités. Une entrée de \(T_{\min}\) est une paire
de germes
\[
 I_9:=\{0,1,\ldots,8\},\qquad
 s_+=(i_+,j_+,d_+),\qquad s_\times=(i_\times,j_\times,d_\times)
 \in I_9^2\times D.
\]
Ces tables, les lecteurs \(L_\Sigma,L_\Pi\), les six blocs de neuf
instants et l'inversion des deux vecteurs lorsque \(27\mid t\) constituent toutes les
données de départ. On n'introduit ni \(\pi\), ni \(e\), ni \(\varphi\), ni aucune autre
valeur cible.
\end{definition}

\begin{algorithm}[Récurrence de 54 pas]
\label{pf84:E02-16}
Au départ, les deux curseurs occupent les positions de leurs germes et portent
les vecteurs \(v(d_+)\), \(v(d_\times)\). Pour
\(t=1,\ldots,54\), dans cet ordre :
\begin{enumerate}
  \item si \(\operatorname{sgn}_9(t)=+1\), on accumule
  \(L_\Sigma(i_+(t),j_+(t))\) et l'on déplace le curseur additif ;
  \item si \(\operatorname{sgn}_9(t)=-1\), on accumule
  \(\operatorname{ind}_2(L_\Pi(i_\times(t),j_\times(t)))\) et l'on déplace le
  curseur multiplicatif ;
  \item si \(\operatorname{sgn}_9(t)=0\), les deux curseurs restent immobiles ;
  \item après cette mise à jour, les deux vecteurs sont inversés si
  \(27\mid t\).
\end{enumerate}
Les accumulateurs sont réinitialisés au début de chaque bloc
\(9(r-1)+1\leq t\leq9r\). Si \(S_r,E_r\) sont leurs valeurs finales, on émet
\[
 U_r=100\bar S_r+10(\bar S_r+\bar E_r)_{10}
       +(3S_r+5E_r+1)_{10},
 \qquad
 \bar S_r=S_r\bmod10,\quad \bar E_r=E_r\bmod10,
\]
et \(w_r=(-U_r)\bmod3\). On obtient ainsi
\[
 T_{\min}(s_+,s_\times)=(U_1,\ldots,U_6)
 \in\mathcal U_6^{\rm cat},
 \qquad
 \Psi(U_6)=(-U_6)\bmod3=:w_6\in\mathbb F_3^6.
\]
\end{algorithm}

\begin{criterion}[Contrôle dual indépendant de la valeur cible]
\label{pf84:E02-17}
La récurrence admet deux formulations indépendantes : factorisation des
canaux et simulation conjointe. Le parcours exhaustif des
\(104\,976\) entrées produit le même sextuplet par les deux voies pour chaque
germe, sans consulter de constante cible. Un écart pour un seul
germe réfuterait l'équivalence.
\end{criterion}

\section{Dynamique bivariée et observables par blocs}

Nous travaillons sur le tore \(\mathbb Z_9^2\). Avec la convention précédente, un
germe complet est
\[
\omega=(p_+,p_\times,h_+,h_\times)
\in(\mathbb Z_9^2)^2\times\{N,E,S,O\}^2,
\]
de sorte que
\[
|\Omega|=9^4\cdot4^2=104\,976.
\]
Dans chaque fenêtre de neuf pas, les temps \(1,4,7\) évaluent l'observable additive,
les temps \(2,5,8\) évaluent l'observable multiplicative et les temps \(3,6,9\)
incrémentent le compteur neutre. Si \(S_m,E_m,T_m\) sont les accumulateurs
de la fenêtre \(m\), l'observable décimale est
\[
d_1\equiv S_m,\qquad
d_2\equiv S_m+E_m,\qquad
d_3\equiv3S_m+5E_m+7T_m\pmod{10},
\]
et le triplet décimal est \(U_m=100d_1+10d_2+d_3\).

\begin{lemma}[Compteur neutre]
Pour toute fenêtre, \(T_m=3\).
\end{lemma}
\begin{proof}
Chaque bloc de neuf temps contient exactement les trois temps neutres
\(3,6,9\). Le compteur est donc incrémenté trois fois.
\end{proof}

\section{Image de l'accumulateur additif : \(\operatorname{Im}S_m=\{6,9,12,15,18,21,24\}\)}

Soit \(s=i+j\pmod9\). L'observable additive visible est
\[
f(s)=\operatorname{dr}_9(s+2)=(2,3,4,5,6,7,8,9,1).
\]
Une translation cardinale modifie \(s\) de \(+1\) ou de \(-1\). Dans une fenêtre, on
lit trois termes consécutifs ; inverser l'orientation inverse seulement leur
ordre. Les neuf sommes cycliques sont
\[
9,12,15,18,21,24,18,12,6.
\]

\begin{lemma}[Image de l'accumulateur additif]
\[
S_m\in\{6,9,12,15,18,21,24\}.
\]
En particulier, \(d_1=S_m\bmod10\) parcourt exactement
\(\{1,2,4,5,6,8,9\}\).
\end{lemma}
\nopagebreak[4]
\begin{proof}
La liste précédente énumère toutes les positions initiales de la trajectoire et les
deux orientations. Leurs résidus décimaux sont au nombre de sept et distincts.\qedhere
\end{proof}

\section{Image de l'accumulateur multiplicatif : \(\operatorname{Im}E_m=\{0,1,3,5,7,9\}\)}

Le groupe des unités \((\mathbb Z/9\mathbb Z)^\times\) est cyclique d'ordre
six. Nous fixons la fonction de coordonnée multiplicative
\[
\ell_\times(1,2,3,4,5,6,7,8,9)=(0,1,0,2,5,0,4,3,0).
\]
Si la ligne ou la colonne parcourue possède un facteur divisible par trois, chaque
évaluation appartient à \(\{3,6,9\}\) et apporte zéro. Si le facteur fixe est une
unité, les sommes de trois termes consécutifs de
\(b\mapsto\ell_\times(\operatorname{dr}_9(ab))\) sont :
\[
\begin{array}{c|c}
a&\text{somme possible}\\
\midrule
1&\{1,3,7,9\}\\
2&\{3,5,9\}\\
4&\{1,5,7\}\\
5&\{1,5,7\}\\
7&\{3,5,9\}\\
8&\{1,3,7,9\}.
\end{array}
\]

\begin{lemma}[Image de l'accumulateur multiplicatif]
\[
E_m\in\{0,1,3,5,7,9\}.
\]
\end{lemma}
\begin{proof}
Le cas non inversible apporte \(0\). La table finie des six unités
contient toutes les possibilités du cas inversible et leur union est
\(\{1,3,5,7,9\}\).
\end{proof}

\section{Congruence décimale de l'alphabet ternaire : unicité du 3e chiffre}

\begin{theorem}[Congruence de contrôle]
Tout triplet \(\overline{d_1d_2d_3}\) émis par le système satisfait
\[
\boxed{d_3\equiv5d_2-2d_1+1\pmod{10}.}
\]
\end{theorem}
\begin{proof}
De \(d_2\equiv S_m+E_m\), on obtient
\(E_m\equiv d_2-d_1\pmod{10}\). Comme \(T_m=3\),
\[
d_3\equiv3d_1+5(d_2-d_1)+21
\equiv5d_2-2d_1+1\pmod{10}.\qedhere
\]
\end{proof}

\begin{theorem}[Image exacte de l'application décimale]
Le système produit exactement \(7\cdot6=42\) triplets décimaux. Regroupés par
\(E=d_2-d_1\pmod{10}\), ce sont :
\[
\begin{array}{c|rrrrrrr}
E&\multicolumn{7}{c}{\text{triplets décimaux}}\\
\midrule
0&114&227&443&556&669&885&998\\
1&129&232&458&561&674&890&903\\
3&149&252&478&581&694&810&923\\
5&169&272&498&501&614&830&943\\
7&189&292&418&521&634&850&963\\
9&109&212&438&541&654&870&983.
\end{array}
\]
\end{theorem}
\begin{proof}
Il existe sept possibilités pour \(d_1\) et six pour \(E\). Pour chaque paire,
\(d_2\equiv d_1+E\) et la congruence de contrôle fixe un unique \(d_3\),
de sorte qu'il existe au plus 42 triplets. La table applique ces deux formules
aux 42 paires et ne contient pas de répétitions. L'énumération des
104\,976 germes retrouve tous les triplets de la table et aucun en dehors
de celle-ci ; la borne est atteinte.
\end{proof}

\section{Factorisation des \(2\) observables}

Soit \(\mathcal S=\mathbb Z_9^2\times\{N,E,S,O\}\), avec
\(|\mathcal S|=324\). À chaque germe additif \(P\in\mathcal S\), on
associe la signature de six fenêtres
\[
s(P)=(S_1,\ldots,S_6)\pmod{10},
\]
et à chaque germe multiplicatif \(T\in\mathcal S\),
\[
e(T)=(E_1,\ldots,E_6)\pmod{10}.
\]
Les deux signatures se combinent coordonnée par coordonnée au moyen de
\[
G(s,e)_k=
100s_k+10(s_k+e_k)_{10}+(3s_k+5e_k+1)_{10}.
\]
L'application \(G\) est injective : à partir d'un triplet, on récupère
\(s_k=d_1\) et \(e_k=d_2-d_1\pmod{10}\), tandis que le troisième chiffre vérifie
l'appartenance à l'image.

\begin{lemma}[Classification exacte du canal additif]
\label{pf84:E02-18}
\[
 |\operatorname{im}s|=18,
\]
et chaque signature additive possède exactement \(18\) préimages.
\end{lemma}
\begin{proof}
La signature dépend uniquement de
\((i+j\bmod9,\varepsilon)\), où \(\varepsilon=+1\) pour \(E,S\) et
\(-1\) pour \(N,O\). Les \(9\cdot2\) paires produisent des signatures distinctes ; chacune
possède neuf positions sur sa diagonale et deux directions, c'est-à-dire
\(18\) préimages. Un nombre différent de signatures ou une classe de taille
différente de \(18\) réfuterait la classification.
\end{proof}

\begin{lemma}[Classification exacte du canal multiplicatif]
\label{pf84:E02-19}
\[
 |\operatorname{im}e|=26.
\]
Il existe \(24\) classes de taille \(8\), une de taille \(24\) et une de taille
\(108\).
\end{lemma}
\begin{proof}
L'énumération des \(9^2\cdot4=324\) germes au moyen de la table explicite
de \(\ell_\times\) donne \(26\) sextuplets
\((E_1,\ldots,E_6)\) et l'histogramme
\[
 24\cdot8+24+108=324.
\]
Une classe supplémentaire, absente ou de taille différente réfuterait le recensement.
\end{proof}

\begin{theorem}[Factorisation exacte et injective du catalogue]
\label{pf84:E02-20}
\[
|\operatorname{im}s|=18,\qquad
|\operatorname{im}e|=26,\qquad
|\operatorname{im}G|=468=18\cdot26.
\]
Chaque signature additive possède 18 préimages. Les signatures multiplicatives présentent
l'histogramme
\[
24\text{ classes de taille }8,\qquad
1\text{ de taille }24,\qquad
1\text{ de taille }108.
\]
\end{theorem}
\begin{proof}
Les deux lemmes précédents donnent les cardinaux des signatures. Dans chaque composante
émise, le chiffre des centaines est \(\bar S_r\), tandis que celui des dizaines est
\(\bar S_r+\bar E_r\bmod10\). Par conséquent,
\[
 \bar S_r=\left\lfloor\frac{U_r}{100}\right\rfloor,\qquad
 \bar E_r=
 \left(\left\lfloor\frac{U_r}{10}\right\rfloor-\bar S_r\right)\bmod10.
\]
Le chiffre des unités vérifie \(3S_r+5E_r+1\bmod10\) ; il n'introduit pas d'autre
identification. Ainsi, \(G\) est injective et son image possède
\(18\cdot26=468\) éléments. Les \(24\) classes multiplicatives de taille
\(8\), combinées aux \(18\) signatures additives, produisent \(432\) sextuplets
de multiplicité \(144\) ; les deux autres classes produisent \(18\) sextuplets de
multiplicité \(432\) et \(18\) de multiplicité \(1944\). En particulier,
\[
 432\cdot144+18\cdot432+18\cdot1944=104\,976.
\]
Deux paires de signatures ayant le même \(U_6\), ou un histogramme qui ne restitue pas le
domaine complet, réfuteraient le théorème.
\end{proof}

La composition complète se factorise donc sous la forme
\[
\mathcal S^2\xrightarrow{s\times e}
\operatorname{im}s\times\operatorname{im}e
\xrightarrow{G}\mathcal U_6
\xrightarrow{\rho_3}\mathcal W_6\subset\mathbb F_3^6.
\]
La dernière projection possède 243 valeurs. La première descente conserve la
séparation somme--produit ; la deuxième conserve 468 états décimaux ; la
troisième élimine la coordonnée de fibre et retient un mot ternaire.

\begin{corollary}[Recensement de la projection observable]
\label{cor:censo-proyeccion-observable}
\label{pf84:E02-21}
L'énumération complète détermine la chaîne de cardinaux
\[
 104\,976\longrightarrow468\longrightarrow243\subset729.
\]
Les \(468\) éléments intermédiaires sont des sextuplets d'entiers compris entre
zéro et \(999\) ; seule leur image finale appartient à \(\mathbb F_3^6\). Les
multiplicités des fibres de la première application satisfont
\[
 432\cdot144+18\cdot432+18\cdot1944=104\,976.
\]
\end{corollary}

\begin{proof}
Le théorème de factorisation donne \(18\cdot26=468\), et l'énumération
exhaustive des \(104\,976\) germes produit exactement les trois
multiplicités indiquées. La réduction composante par composante parcourt les
\(243\) mots enregistrés dans le catalogue ternaire. Les catalogues de
signatures et de routes joints reproduisent les deux énumérations.
\end{proof}

\begin{proposition}[Microfibre quintuple du mot \(010211\)]
\label{pf84:E02-22}
Pour la réduction restreinte
\[
 \Psi:\mathcal U_6^{\rm cat}\longrightarrow\mathbb F_3^6,
 \qquad \Psi(U_6)=(-U_6)\bmod3,
\]
la fibre de \(010211\) se compose exactement des cinq sextuplets suivants :
\begingroup
\scriptsize
\[
\begin{array}{@{}cclc@{}}
\toprule
U_6&\sigma_+&\sigma_\times&\text{multiplicité}\\
\midrule
(501,614,498,169,272,272)&(5,6,4,1,2,2)&(5,5,5,5,5,5)&432\\
(810,923,870,169,272,272)&(8,9,8,1,2,2)&(3,3,9,5,5,5)&144\\
(810,983,810,169,272,272)&(8,9,8,1,2,2)&(3,9,3,5,5,5)&144\\
(870,923,810,169,272,272)&(8,9,8,1,2,2)&(9,3,3,5,5,5)&144\\
(870,923,810,418,674,521)&(8,9,8,4,6,5)&(9,3,3,7,1,7)&144\\
\bottomrule
\end{array}
\]
\endgroup
Ce sont cinq préimages du catalogue \(U_6\), non cinq histoires complètes
TPK. Elles conservent des signatures et des multiplicités distinctes même si elles partagent le même
mot visible.
\end{proposition}

\begin{proof}
Le balayage exhaustif du catalogue de \(468\) sextuplets sélectionne exactement
les cinq lignes de la table. La fonction de récupération des signatures reproduit
les deux sextuplets de chaque ligne et la somme des multiplicités est
\(432+4\cdot144=1008\). L'énumération vérifie en outre que chaque cellule est le
produit cartésien de ses germes de canal. Un cardinal différent de cinq
ou l'identification de deux lignes malgré leurs données enrichies réfuterait
l'énoncé.
\end{proof}
 \end{HMTScientificOwnerSubordinated}

\begin{HMTScientificOwnerSubordinated}
\chapter{Systèmes inverses, représentation de Weyl--Heisenberg et quotient prédictif}
\label{chap:numero-heisenberg-d108}
\index{nombre HMT}
\index{Heisenberg!groupe fini}

La construction distingue trois niveaux : le catalogue fini des germes,
l'espace des états qui conserve la retenue et l'orientation, et la limite
inverse des prolongements compatibles. Cette stratification permet de
formaliser la contraction des intervalles cylindriques jusqu'à un point réel,
la persistance des données de trajectoire dans une même fibre d'évaluation et
la représentation finie de Weyl--Heisenberg associée à la forme alternée
du tore nonadique qui porte les observables additif et multiplicatif.

Les symboles \(V_t\), \(X_t\) et \(\rho_{t+1,t}\) utilisés ici sont les
mêmes ensembles d'états locaux, préfixes admissibles et troncatures
introduits dans le \cref{chap:numero-medicion}, avec le changement purement
notationnel \(t=m\) et \(\rho_{n,m}=p_{n,m}\). Cette section explicite leurs
coordonnées et n'introduit pas de second système inverse.

\section{Projection observable et espace complet des états du TPK}

La projection observable du système de transport de phase (TPK) est
l'application finie de codage
\[
104\,976\longrightarrow468\longrightarrow243.
\]
qui classe les germes, signatures et mots \(w_6\). L'espace complet des
états du TPK conserve en outre
\[
\begin{gathered}
\text{orientation bidirectionnelle},\quad
\text{état de Hensel},\quad
\text{cylindre et retenue},\\
\text{signature conjointe de frontière},\quad
\text{phase nonadique visible }g_{\rm vis}(m).
\end{gathered}
\]
Deux candidats ayant une même signature locale peuvent appartenir à des classes distinctes
lorsque la frontière suivante est conservée. Nous n'introduisons pas de second tuple
d'état. Pour
\[
x_m=(q_m,o_m,h_m,c_m,\sigma_m,C_m,b_m,\lambda_m)\in\XTPK^{(m)}
\]
nous utilisons la carte de coordonnées
\[
\chi_m(x_m)=
\bigl(
B_m,h_m,(C_m,b_m),\sigma_m,
o_m,\mu(q_m),g_{\rm vis}(m),\lambda_m
\bigr),
\]
avec
\[
B_m=(u_m^\pi,u_m^e,u_m^\varphi)\in(\mathbb F_3^6)^3.
\]
Ici \(h_m\) est l'extension de Hensel ; \((C_m,b_m)\), le cylindre et sa
frontière ; \(o_m,\mu(q_m)\), l'orientation et l'observable bidirectionnel ; et
\[
g_{\rm vis}(m)=1+((m-1)\bmod9)\in\{1,\ldots,9\},
\]
l'étiquette visible de phase nonadique. La coordonnée de classe correspondante
est \(g_0(\theta_m)\in\ZZ/9\ZZ\), avec \(\theta_m=[m-1]_{108}\). La signature conjointe
est la coordonnée déjà typée par le noyau de phase,
\[
\sigma_m=\operatorname{Sig}_{q_m}(h_m,c_m,\lambda_m)
=(q_\partial,a_\partial,c_\partial,
\operatorname{colw},\rho_W,r_\partial).
\]
La signature conjointe ne se réduit pas à trois signatures de canal indépendantes, car la
sélection associée à l'involution \(\pi/e\) et à l'axe \(\varphi\) dépend de
données corrélées. Deux nœuds ayant le même \(B_m\) représentent des états
distincts lorsqu'ils diffèrent par l'extension, le cylindre, la retenue, la signature, l'orientation
ou le registre de transport.

\subsection{Compatibilité exacte des cylindres sous troncature et prolongement nonadique}

Soit \(N\) l'entier d'un préfixe ternaire de longueur \(L\), et soit \(D\)
l'entier obtenu en concaténant ses \(b\) triades décimales. Nous définissons
\[
A=N1000^b-D3^L,
\qquad
B=(N+1)1000^b-D3^L.
\]
Les cylindres se coupent si et seulement si
\[
\boxed{A<3^L,\qquad B>0.}
\]
En ajoutant un bloc ternaire \(u\in\{0,\ldots,728\}\), on a
\[
N'=729N+u.
\]
S'il n'apparaît pas de nouvelle triade décimale,
\[
A'=729A+u1000^b.
\]
Si \(\delta\) apparaît, de sorte que \(D'=1000D+\delta\), alors
\[
A'=729000A+u1000^{b+1}-\delta\,729\,3^L.
\]
Si \(\Delta b\in\{0,1\}\) enregistre l'apparition de cette nouvelle triade, dans
les deux cas
\[
B'=A'+1000^{b+\Delta b}.
\]
La transition, la paire admissible \((u,\delta)\) étant fixée, ne consulte pas de valeur
cible. La filtration conjointe par intersection de cylindres et retour
nonadique agit à un niveau différent : elle sélectionne les paires de frontière qui
admettent un prolongement compatible.

\subsection{Registre cyclique de 108 étapes et sa projection réduite}

Nous appelons \emph{registre cyclique de 108 étapes} la suite de
\(108\) pas regroupés en douze secteurs de neuf pas, avec la convention
d'indice fixée dans le \cref{chap:tpk-definicion}.

Soit \(\XTPK^{[108]}\subseteq\XTPK\) le domaine des états complets dont le
registre \(\lambda\) contient les \(108\) pas déclarés. Notons
\(\mathcal C_{108}^{\rm red}\) l'ensemble des registres qui résultent de la
conservation des seuls champs réduits, et \(\mathcal U_{12}^{\rm sgn}\)
l'ensemble des dodécaphases signées admissibles. L'instance distingue deux applications
déclarées :
\[
F:\XTPK^{[108]}\longrightarrow\mathcal C_{108}^{\rm red},
\qquad
R:\XTPK^{[108]}\longrightarrow\mathcal U_{12}^{\rm sgn}.
\]
L'application \(F\) ne retient que les champs du registre réduit, tandis que
\(R\) évalue la dodécaphase signée. L'existence d'une application \(\Xi\)
telle que \(R=\Xi\circ F\) équivaut à la factorisation de \(R\) par une
projection qui peut éliminer la coordonnée de l'observable, l'orientation,
la retenue ou la frontière. Cette factorisation constitue une proposition
supplémentaire et ne fait pas implicitement partie de la reconstruction sur l'espace
complet des états.

Le registre réduit est, par définition, une projection de l'état complet,
non l'état de départ. Les résultats ultérieurs utilisent \(R\) sur
l'état complet et ne présupposent pas de factorisation par \(F\).

\section{Systèmes inverses de survivants}

La transition TPK n'est pas nécessairement une fonction. Pour typer correctement
ce fait, soit \(V_t\) l'ensemble fini des états complets possibles à
la profondeur \(t\), et soit
\[
\mathcal R_t\subseteq V_t\times V_{t+1}
\]
la relation de transition admissible, tous les champs de l’état
holonomique intégral étant conservés. Nous définissons alors \(X_t\), non comme l’ensemble des
états terminaux isolés, mais comme l’ensemble des préfixes admissibles, c’est-à-dire
des histoires admissibles complètes :
\[
X_t=
\Bigl\{
(v_0,\ldots,v_t)\in\prod_{j=0}^{t}V_j:
(v_j,v_{j+1})\in\mathcal R_j\ (0\le j<t),
\ \text{et les clôtures sont satisfaites jusqu’à }t
\Bigr\}.
\]
Sur ces objets, la troncature est bien une application canonique :
\[
\begin{aligned}
\rho_{t+1,t}:X_{t+1}&\longrightarrow X_t,\\
(v_0,\ldots,v_t,v_{t+1})&\longmapsto(v_0,\ldots,v_t).
\end{aligned}
\]
Elle est bien définie précisément parce que tout préfixe d'une histoire
admissible reste une histoire admissible pour les fermetures déjà imposées.
Si un filtre utilisait une information future et pouvait invalider
rétroactivement un préfixe, il faudrait d'abord remplacer \(X_t\) par le
sous-ensemble des préfixes prolongeables ; avec cette convention, la même troncature
redevient bien définie.

Une histoire HMT globale est une famille compatible de préfixes
\[
x=(x_t)_{t\ge0},
\qquad
x_t\in X_t,
\qquad
\rho_{t+1,t}(x_{t+1})=x_t.
\]
Elle équivaut à une suite infinie \((v_0,v_1,\ldots)\) dont les paires
consécutives appartiennent aux relations \(\mathcal R_t\) et dont les préfixes
satisfont toutes les fermetures. L'espace des histoires est la limite inverse
\[
X_\infty=\varprojlim X_t.
\]
Cette formulation évite d'imposer une transition locale univoque
\(B_t\mapsto B_{t+1}\). La transition peut être une relation finie, plusieurs
branches peuvent partager le même état terminal visible et l'unicité peut
n'apparaître qu'après avoir imposé la compatibilité future. En conservant le
préfixe complet, on ne perd pas prématurément l'information permettant
de les distinguer.

\begin{theorem}[Section coinductive unique]
Supposons que, pour tout \(t\), l'ensemble des survivants \(S_t\subseteq
X_t\) soit non vide, que les restrictions
\(\rho_{t+1,t}:S_{t+1}\to S_t\) soient compatibles et que chaque \(S_t\) contienne un
unique préfixe. Alors \(\varprojlim S_t\) contient une unique histoire
globale.
\end{theorem}

\begin{proof}
Soit \(x_t\) l'unique préfixe de \(S_t\). La compatibilité impose
\(\rho_{t+1,t}(x_{t+1})=x_t\), de sorte que \((x_t)_t\) appartient à la limite
inverse. Toute autre histoire aurait à chaque niveau un membre de \(S_t\)
et, par unicité, coïnciderait préfixe par préfixe avec \((x_t)_t\).
\end{proof}

Le théorème formalise le passage de la compatibilité niveau par niveau à une histoire
globale. Il ne transforme pas une fenêtre finie en une conclusion coinductive : son
application à une famille concrète exige de vérifier à chaque profondeur la non-
vacuité, la compatibilité des restrictions et l'unicité des préfixes
survivants.

\section{Connexion nonadique de survie}
\label{sec:conexion-nonadica}

La contraction des cylindres ne procède pas bloc par bloc comme une substitution
déterministe. À certaines frontières, il existe plusieurs répartitions locales d'une
même donnée agrégée ; la prolongeabilité décide laquelle appartient à l'histoire
publiée. La structure résultante est une connexion sur les neuf classes de
phase : à l'achèvement d'un tour, la classe de phase revient, mais l'état de la
fibre a changé. Cette section construit la fenêtre finie complète dans laquelle
ce retour est observé.

Le domaine du calcul est le registre de transport calibré joint. Ses états initiaux,
sa matrice entière de propagation et ses relations de frontière sont des prémisses
publiées, et non les sorties d’une énumération autonome de TPK à profondeur
arbitraire. Toutes les affirmations de cette section sont exactes relativement à
ces prémisses. L’extension à toute profondeur est régie par le théorème
coinductif précédent et exige de vérifier ses hypothèses à chaque niveau.

\subsection{Préfixe de \(30\) trits et 1re frontière}

Nous écrivons \(u_k^\chi\in\FF_3^6\), avec
\(\chi\in\{\pi,e,\varphi\}\), pour le bloc visible du canal \(\chi\) à la
profondeur \(k\). Les cinq premiers blocs forment trois préfixes de trente
trits :
\begin{align*}
w_{30}^{\pi}&=010211|012222|010211|002111|110221,\\
w_{30}^{e}&=201101|121221|102011|012222|102011,\\
w_{30}^{\varphi}&=121200|112202|121020|010210|010200.
\end{align*}
Le sixième bloc est la première frontière postérieure à ces préfixes :
\[
R_{36}=
\begin{pmatrix}
u_5^\pi\\ u_5^e\\ u_5^\varphi
\end{pmatrix}
=
\begin{pmatrix}
222220\\021222\\102011
\end{pmatrix}.
\]
La notation \(R_{36}\) indique que trente-six trits par
canal ont été exposés. Elle ne désigne pas un état terminal : le registre reconstruit ci-dessous
contient vingt blocs, et la récurrence admet des continuations ultérieures
lorsque les états de Hensel correspondants sont fournis.

\subsection{Axe d'auto-échelle et fibre résiduelle}

\ifdefined\HMTAxialTheoremDefined
La décomposition axiale de la signature et l'identification de l'axe
\(\varphi\) face à la fibre résiduelle \(\pi/e\) ont déjà été démontrées dans le
\cref{thm:eje-espejo-nonadico}. On conserve ici leur conséquence pour la
fenêtre nonadique sans dupliquer le théorème ni la preuve.
\else
Pour
\[
B_k=(u_k^\pi,u_k^e,u_k^\varphi)\in(\FF_3^6)^3
\]
nous définissons, coordonnée par coordonnée,
\[
q_k=u_k^\pi+u_k^e+u_k^\varphi,
\qquad
a_k=u_k^\pi+u_k^e-u_k^\varphi.
\]

\begin{theorem}[Décomposition axiale de la signature de frontière]
\label{thm:eje-espejo-nonadico}
La paire \((q_k,a_k)\) détermine de manière univoque le canal axial et l'agrégat des
deux autres canaux :
\[
u_k^\varphi=2(q_k-a_k),
\qquad
u_k^\pi+u_k^e=q_k-u_k^\varphi.
\]
Par conséquent, une fibre de l'application
\[
q_{\rm ax}(u^\pi,u^e,u^\varphi)
=(u^\pi+u^e,u^\varphi)
\]
ne peut conserver d'ambiguïté que dans la répartition interne entre \(u^\pi\) et
\(u^e\).
\end{theorem}

\begin{proof}
La soustraction des deux équations de définition donne
\(q_k-a_k=2u_k^\varphi\). L'inverse de \(2\) dans \(\FF_3\) est encore
\(2\), d'où
\(u_k^\varphi=2(q_k-a_k)\). Soustraire ce vecteur de \(q_k\) produit
\(u_k^\pi+u_k^e\). L'opération ne retrouve pas chaque terme séparément ;
l'unique fibre résiduelle est donc l'ensemble des répartitions de l'agrégat fixé.
\end{proof}

Dans la terminologie géométrique du modèle, \(\varphi\) constitue l'axe
fixé par la signature et la paire \(\pi/e\) la fibre résiduelle. L'expression ne
postule pas de symétrie analytique des constantes réelles : elle nomme la
décomposition linéaire exacte qui vient d'être démontrée dans \(\FF_3^6\).
\fi

\subsection{Orbite complète des \(9\) phases consécutives du Topological Phase Kernel}

Nous fixons la phase absolue au moyen de
\[
g(k)=
\begin{cases}
k\bmod9,&k\not\equiv0\pmod9,\\
9,&k\equiv0\pmod9.
\end{cases}
\]
Comme \(w_{30}\) occupe \(k=0,\ldots,4\), le premier tour complet postérieur
au préfixe parcourt \(k=5,\ldots,13\), avec les phases
\(5,6,7,8,9,1,2,3,4\). L'énumération des \(729=3^6\) candidats de
signature agrégée à chaque niveau produit la fenêtre suivante :
\begin{center}
\small
\begin{tabular}{@{}ccrrlll@{}}
\toprule
\(k\)&\(g(k)\)&locaux&sélectionnés&\(u_k^\pi\)&\(u_k^e\)&\(u_k^\varphi\)\\
\midrule
5&5&1&1&222220&021222&102011\\
6&6&1&1&111201&201202&221021\\
7&7&3&1&212121&222210&200221\\
8&8&3&1&200121&212212&110110\\
9&9&2&1&100100&020112&010122\\
10&1&1&1&101222&112221&112002\\
11&2&1&1&022212&110001&100021\\
12&3&1&1&012012&202222&000120\\
13&4&1&1&111210&112102&211122\\
\bottomrule
\end{tabular}
\end{center}
La colonne « locaux » compte les états d'entrée admissibles au niveau
indiqué ; elle ne compte pas les sorties de la relation de frontière. Cette distinction
est essentielle dans \(k=9\) : les deux états locaux d'entrée précèdent une
fibre de trois sorties de l'état sélectionné.

\subsection{Retour 9→1 de la phase nonadique avec changement d'état et avancée de la mémoire}

Soit
\[
 \pi_B:V_k\longrightarrow\mathcal B:=(\FF_3^6)^3
\]
la projection sur le triplet visible. La relation de transport sur les états
complets induit, une fois fixées les données de fibre du registre, une relation
visible \(\MonNine^B\). L'état sélectionné qui entre dans la neuvième phase est
\[
B_9=(100100\mid020112\mid010122).
\]
Son image relationnelle se compose exactement de
\[
\begin{aligned}
B_{10}^{(0)}&=(101222\mid112221\mid112002),\\
B_{10}^{(1)}&=(102221\mid111222\mid112002),\\
B_{10}^{(2)}&=(111221\mid102222\mid112002).
\end{aligned}
\]
Les trois sorties ont
\[
q_{10}=022112,
\qquad a_{10}=101111,
\]
et, par le \cref{thm:eje-espejo-nonadico}, partagent
\[
u_{10}^\varphi=112002,
\qquad u_{10}^\pi+u_{10}^e=210110.
\]
Elles diffèrent exclusivement par les trois répartitions publiées de cet agrégat
entre \(\pi\) et \(e\).

\begin{theorem}[Sélection par deux niveaux de prolongeabilité]
\label{thm:seleccion-novena-fase}
Dans la fenêtre calibrée, les trois sorties précédentes admettent un prolongement
de profondeur un. En imposant également le bloc suivant
\[
B_{11}=(022212\mid110001\mid100021)
\]
et la frontière décimale \((502,662,638)\), seule
\(B_{10}^{(0)}\) admet un prolongement du même état complet. Par conséquent
\[
\#\operatorname{Surv}^{\rm cal}_{1}(B_9)=3,
\qquad
\#\operatorname{Surv}^{\rm cal}_{2}(B_9)=1.
\]
\end{theorem}

\begin{proof}
Les triplets décimaux associés aux trois sorties sont, respectivement,
\[
(279,352,365),\quad(280,351,365),\quad(282,350,365).
\]
Pour un mot ternaire de longueur \(L\), d'entier associé \(N\), et un
préfixe de \(b\) triplets, d'entier associé \(D\), la compatibilité des
deux cylindres équivaut aux deux inégalités entières
\[
N1000^b<(D+1)3^L,
\qquad
(N+1)1000^b>D3^L.
\]
Les six inégalités de chaque sortie sont satisfaites au niveau actuel. En
annexant \(B_{11}\) et \((502,662,638)\), elles sont à nouveau satisfaites pour les trois
canaux de \(B_{10}^{(0)}\). Pour \(B_{10}^{(1)}\) et \(B_{10}^{(2)}\),
l'inégalité du canal \(\varphi\) est conservée, mais celles de \(\pi\) et
\(e\) échouent. L'évaluation s'effectue avec des entiers dans le certificat fini correspondant ; aucune
tolérance numérique n'intervient. Par conséquent, l'ensemble des candidats
se réduit de trois à un lors du passage de la profondeur un à la profondeur deux.
\end{proof}

Les deux autres mots ne peuvent précéder une même queue ternaire que si
l'extension profonde de Hensel est modifiée. Ils ne sont donc pas des prolongements du
même état complet du registre : le reste visible coïncide en partie, mais
la généalogie de la fibre est distincte.

\subsection{Retour de phase avec changement d'état}

La phase satisfait \(g(k+9)=g(k)\). L'état visible, en revanche, ne se
réinitialise pas. En particulier,
\[
B_1=(012222\mid121221\mid112202),
\qquad
B_{10}=(101222\mid112221\mid112002),
\]
de sorte que \(g(1)=g(10)=1\) mais \(B_1\ne B_{10}\). La vérification des
neuf paires \(k\mapsto k+9\), pour \(k=1,\ldots,9\), donne en outre
\[
u_{k+9}^\varphi\ne u_k^\varphi,
\qquad
u_{k+9}^\pi+u_{k+9}^e\ne u_k^\pi+u_k^e.
\]
Par conséquent, la fenêtre satisfait exactement
\[
\boxed{g(k+9)=g(k),\qquad B_{k+9}\ne B_k
\quad(1\le k\le9).}
\]
Il s'agit d'un retour de phase avec holonomie visible non triviale. La dénomination
\emph{monodromie nonadique} désigne cette action de retour sur la
fibre enrichie ; elle n'introduit pas de dixième phase. Pour élever l'observation
finie à une monodromie globale, il faut construire l'action à toute profondeur
et démontrer la compatibilité coinductive exigée au début de la section.

\subsection{La branche de \(20\) blocs}

Le retour précédent ne clôt pas l'émission. Les vingt blocs que le
certificat reconstruit à partir des trois états entiers publiés sont
\begin{align*}
\pi:\;&010211|012222|010211|002111|110221|222220|111201|\\
&212121|200121|100100|101222|022212|012012|111210|\\
&121011|200220|120210|000101|022010|020111,\\[2pt]
e:\;&201101|121221|102011|012222|102011|021222|201202|\\
&222210|212212|020112|112221|110001|202222|112102|\\
&102010|022020|222002|012010|001112|022212,\\[2pt]
\varphi:\;&121200|112202|121020|010210|010200|102011|221021|\\
&200221|110110|010122|112002|100021|000120|211122|\\
&202200|202110|221000|100102|111122|020010.
\end{align*}
Ainsi, \(w_{30}\) est le préfixe et \(R_{36}\) la première frontière d’une branche
explicitement continuée ; aucun des deux objets ne remplace l’état
holonomique intégral qui transporte cette continuation.

\subsection{Décodage exact de \(12\) triplets décimaux au moyen du lecteur nonadique}

Soit \(N_{\chi,k}\) l'entier représenté en base trois par les \(k\) premiers
blocs du canal \(\chi\). Avec \(k=13\), le mot a \(78\) trits et
définit
\[
D_{\chi,12}
=\left\lfloor\frac{1000^{12}N_{\chi,13}}{3^{78}}\right\rfloor.
\]
L'écriture de \(D_{\chi,12}\) en douze blocs de trois chiffres est
\[
\begin{array}{c|rrrrrrrrrrrr}
\pi&141&592&653&589&793&238&462&643&383&279&502&884\\
e&718&281&828&459&045&235&360&287&471&352&662&497\\
\varphi&618&033&988&749&894&848&204&586&834&365&638&117.
\end{array}
\]
L'identité est entière : elle équivaut à
\[
D_{\chi,12}3^{78}
\le 1000^{12}N_{\chi,13}
<(D_{\chi,12}+1)3^{78}.
\]
Les deux cylindres se coupent donc dans chaque canal et l’opération de
lecture produit exactement les trente-six chiffres publiés. Ce
résultat est une conséquence du registre calibré. Les états initiaux de
la carte jointe ont été construits historiquement avec ces constantes comme
données de référence ; cette identité n’établit donc pas la provenance
indépendante des états initiaux à partir des primitives d’APP\@.

\subsection{Calendrier de réouverture des cylindres et périodicité du prolongement compatible}

La conversion d'un bloc de six trits en triplets décimaux a pour pente
moyenne
\[
\lambda=6\log_{1000}3\in(0,1),
\qquad k(t)=\lfloor\lambda(t+1)\rfloor.
\]
Un indice \(t\) est une réouverture sans nouveau triplet lorsque
\(k(t+1)=k(t)\). Soit \(\delta=1-\lambda\).

\begin{theorem}[Calendrier de Beatty de la lecture ternaire--décimale]
\label{thm:calendario-beatty}
L'ensemble des réouvertures est
\[
\mathcal T=
\left\{\left\lceil\frac{m}{\delta}\right\rceil-2:m\ge1\right\}.
\]
Ses premiers éléments sont
\[
20,42,64,86,108,130,151,173,195,217,\ldots,
\]
et deux réouvertures consécutives sont distantes de \(21\) ou \(22\) blocs.
\end{theorem}

\begin{proof}
La quantité \(\lambda=\log_{1000}729\) est irrationnelle : une égalité
\(\lambda=p/q\) impliquerait \(729^q=1000^p\), ce qui est impossible par factorisation
en nombres premiers. \(\delta\) est également irrationnelle, d'où
\[
k(t)=(t+1)-\lceil\delta(t+1)\rceil.
\]
L'égalité \(k(t+1)=k(t)\) se produit exactement lorsque la fonction plafond du
second terme augmente d'une unité. La \(m\)-ième occurrence s'obtient
en résolvant
\(\delta(t+1)<m<\delta(t+2)\), ce qui donne
\(t=\lceil m/\delta\rceil-2\). Les différences d'une suite de Beatty
appartiennent à
\[
\left\{\left\lfloor\delta^{-1}\right\rfloor,
\left\lceil\delta^{-1}\right\rceil\right\}=\{21,22\}.\qedhere
\]
\end{proof}

\subsection{Récurrence de Bockstein--Hensel du reste et du quotient}

La mémoire transportée par la branche provient de la suite exacte
\[
0\longrightarrow\ZZ^6\xrightarrow{\,\times3\,}\ZZ^6
\longrightarrow\FF_3^6\longrightarrow0
\]
et de la matrice unimodulaire
\[
A=
\begin{pmatrix}
2&2&2&1&2&1\\
2&1&2&2&1&1\\
1&0&0&1&0&2\\
0&0&1&1&1&0\\
2&2&2&0&2&1\\
2&1&2&1&0&0
\end{pmatrix},
\qquad\det A=1.
\]
Pour un état ligne \(x_k^\chi\in\ZZ^6\), on définit
\[
y_k^\chi=x_k^\chi A,
\qquad
u_k^\chi=y_k^\chi\bmod3\in\{0,1,2\}^6,
\qquad
x_{k+1}^\chi=\frac{y_k^\chi-u_k^\chi}{3}.
\]

\begin{proposition}[Conservation exacte de la mémoire du quotient]
\label{prop:memoria-bockstein-hensel}
\(x_k^\chi\) et \(A\) étant fixés, la paire
\((u_k^\chi,x_{k+1}^\chi)\) est unique et satisfait
\[
x_k^\chi A=u_k^\chi+3x_{k+1}^\chi.
\]
Réciproquement, comme \(A\in\operatorname{GL}_6(\ZZ)\), la paire détermine
de manière univoque \(x_k^\chi\).
\end{proposition}

\begin{proof}
La division euclidienne de chaque coordonnée de \(x_k^\chi A\) par trois produit
un reste unique dans \(\{0,1,2\}\) et un quotient entier unique. Cela prouve
le premier énoncé et l'identité. Puisque \(\det A=1\), l'inverse
\(A^{-1}\) a des entrées entières ; par conséquent
\[
x_k^\chi=(u_k^\chi+3x_{k+1}^\chi)A^{-1}
\]
retrouve l'état précédent sans ambiguïté.
\end{proof}

Le reste \(u_k^\chi\) est la partie visible et le quotient
\(x_{k+1}^\chi\) la mémoire 3-adique transportée. Le certificat reproduit
soixante pas coordonnés ---vingt par canal---, vérifie \(AA^{-1}=I\) par
arithmétique entière et confronte tous les blocs au registre. La dénomination
Bockstein--Hensel fait ici référence à cette séparation reste--quotient et à son
itération 3-adique ; elle ne s'identifie pas sans données supplémentaires à un opérateur de
Bockstein cohomologique.

\section{Du préfixe à un point réel}

Associons à chaque préfixe compatible \(x_t\in X_t\) un intervalle fermé
\(I_t(x_t)\subset\mathbb R\). L'interprétation géométrique est celle d'un
cylindre : tous les futurs qui partagent le même préfixe se projettent à l'intérieur
de \(I_t\).

\begin{theorem}[Convergence archimédienne d'intervalles emboîtés]
Soit \((x_t)_{t\ge0}\) une histoire compatible. Si
\[
I_{t+1}(x_{t+1})\subseteq I_t(x_t)
\]
et
\[
\operatorname{diam} I_t(x_t)\longrightarrow0,
\]
alors il existe un unique \(r\in\mathbb R\) tel que
\[
\bigcap_{t\ge0}I_t(x_t)=\{r\}.
\]
\end{theorem}

\begin{proof}
C'est le principe des intervalles emboîtés. Les extrémités gauches forment une
suite croissante bornée et les droites une suite décroissante bornée. Leurs limites
existent ; leur différence est bornée par le diamètre, qui tend vers
zéro. Les deux limites coïncident en un point unique contenu dans tous les
intervalles.
\end{proof}

Soit \(X_\infty^{\rm Arch}\subseteq X_\infty\) le domaine des histoires pour
lesquelles les cylindres sont emboîtés et leurs diamètres tendent vers zéro. Si ces
deux propriétés font partie de l'admissibilité globale, alors
\(X_\infty^{\rm Arch}=X_\infty\) ; sinon, cette restriction de domaine est
nécessaire. Nous définissons l'évaluation
\[
\operatorname{ev}_{\mathbb R}:X_\infty^{\rm Arch}\longrightarrow\mathbb R,
\qquad
\operatorname{ev}_{\mathbb R}(x)=
\text{l’unique point de }\bigcap_{t\ge0}I_t(x_t).
\]
L'unicité de la valeur n'impose pas l'unicité de l'histoire. Deux généalogies
\(x\ne y\) peuvent satisfaire
\[
\operatorname{ev}_{\mathbb R}(x)
=\operatorname{ev}_{\mathbb R}(y).
\]
Nous définissons l'équivalence de lecture
\[
x\sim_{\rm ev}y
\quad\Longleftrightarrow\quad
\operatorname{ev}_{\mathbb R}(x)
=\operatorname{ev}_{\mathbb R}(y).
\]

\begin{theorem}[Quotient de l'évaluation archimédienne]
L'application
\[
\begin{aligned}
\overline{\operatorname{ev}}_{\mathbb R}:
X_\infty^{\rm Arch}/\!\sim_{\rm ev}
&\longrightarrow
\operatorname{Im}(\operatorname{ev}_{\mathbb R})\subseteq\mathbb R,\\
[x]&\longmapsto\operatorname{ev}_{\mathbb R}(x)
\end{aligned}
\]
est une bijection. En particulier, le quotient s'identifie à tout
\(\mathbb R\) si et seulement si l'évaluation
\(\operatorname{ev}_{\mathbb R}\) est surjective.
\end{theorem}

\begin{proof}
La définition de \(\sim_{\rm ev}\) rend la valeur indépendante du
représentant, de sorte que l'application est bien définie. Elle est injective car
deux classes ayant la même valeur sont, par définition, une seule classe ; elle est
surjective sur l'image par la définition même de l'image. Le dernier
énoncé équivaut à
\(\operatorname{Im}(\operatorname{ev}_{\mathbb R})=\mathbb R\).
\end{proof}

Le théorème de couverture du lecteur archimédien construit précisément ces
préfixes : pour tout intervalle compact \(K\subset\mathbb R\),
\(\operatorname{ev}_{\mathbb R}(\Omega_{\infty,K})=K\), et les intervalles
\([-m,m]\) forment une famille dirigée dont l'union est \(\mathbb R\). Par
conséquent, chaque réel possède au moins une histoire HMT compatible et
l'énoncé démontré est
\[
X_\infty^{\rm Arch}/\!\sim_{\rm ev}
\cong\operatorname{Im}(\operatorname{ev}_{\mathbb R})
=\mathbb R.
\]
La valeur archimédienne est donc le quotient qui identifie les fibres de
l'évaluation ; le nombre HMT enrichi conserve en outre la classe de
trajectoire, la retenue et l'orientation. L'identification ne découle pas de la
simple formation abstraite du quotient : elle découle du théorème de couverture déjà
démontré dans la construction discrète du continu.

La descente algébrique est un autre problème et doit être résolue opération par
opération. Supposons que deux opérations candidates aient été construites,
\(\oplus\) et \(\odot\), sur \(X_\infty^{\rm Arch}\), internes à ce
domaine. La somme passe au quotient si et seulement si
\[
x\sim_{\rm ev}x',\ y\sim_{\rm ev}y'
\Longrightarrow
x\oplus y\sim_{\rm ev}x'\oplus y',
\]
tandis que le produit passe au quotient, de manière indépendante, si et seulement si
\[
x\sim_{\rm ev}x',\ y\sim_{\rm ev}y'
\Longrightarrow
x\odot y\sim_{\rm ev}x'\odot y'.
\]
La première congruence n'implique pas la seconde. En outre, pour que les opérations
induites soient précisément la somme et le produit ordinaires du quotient
archimédien, il faut prouver séparément
\[
\operatorname{ev}_{\mathbb R}(x\oplus y)
=\operatorname{ev}_{\mathbb R}(x)+\operatorname{ev}_{\mathbb R}(y),
\]
\[
\operatorname{ev}_{\mathbb R}(x\odot y)
=\operatorname{ev}_{\mathbb R}(x)\operatorname{ev}_{\mathbb R}(y),
\]
et que l'image de l'évaluation soit stable par l'opération correspondante.
Ce n'est qu'après ces certificats que le quotient hérite de la structure de
sous-anneau ---ou de sous-corps, si les opposés et les inverses sont également construits---
de \(\mathbb R\). La bijection ensembliste du théorème précédent est formelle ;
le contenu spécifique de HMT consiste à construire l'évaluation, à caractériser ses
fibres, à démontrer la couverture et à établir chaque loi de descente.

\section{Représentation de Weyl--Heisenberg associée au tore nonadique}

La forme d'aire orientée du tore nonadique est
\[
\sigma((a,b),(c,d))=ad-bc\pmod9.
\]
Cette forme alternée est celle que l'on obtient, avec l'orientation choisie, en
intégrant la courbure mixte d'APP sur des parallélogrammes du tore. Soit
\(\omega=\exp(2\pi i/9)\) et
\(\mathcal H_9=\mathbb C^9\), de base \(\{|n\rangle:n\in\mathbb Z_9\}\).
Nous définissons
\[
X|n\rangle=|n+1\rangle,
\qquad
Z|n\rangle=\omega^n|n\rangle.
\]
Alors
\[
ZX=\omega XZ,
\]
et, plus généralement,
\[
(X^aZ^b)(X^cZ^d)
=\omega^{bc}X^{a+c}Z^{b+d}.
\]
Par conséquent
\[
 c\bigl((a,b),(c,d)\bigr)=bc\pmod9
\]
est le \(2\)-cocycle polarisé qui apparaît dans cette section de l'extension
centrale. Nous définissons explicitement
\[
 \operatorname{Heis}_9
 =\{\omega^kX^aZ^b:a,b,k\in\ZZ/9\ZZ\}.
\]
C'est un groupe d'ordre \(9^3=729\). Il s'agit d'une extension centrale associée
au groupe des translations du support \(X_9\) d'APP, non d'un sous-groupe du
graphe ni de la catégorie des chemins d'APP\@.
Le commutateur des deux transports est
\[
(X^aZ^b)(X^cZ^d)(X^aZ^b)^{-1}(X^cZ^d)^{-1}
=\omega^{bc-ad}I.
\]
Par conséquent, son exposant est \(-\sigma((a,b),(c,d))\). Le cocycle
polarisé \(c\) et la forme \(\sigma\) ne sont pas le même objet : la seconde est,
avec cette convention, l'opposée de l'antisymétrisation
\(c(u,v)-c(v,u)\). La phase du commutateur est ainsi l'inverse de la phase de
Wilson pour l'orientation fixée.

Pour construire des observables, nous prenons
\[
Q_{\rm H}=\operatorname{diag}(0,1,\ldots,8),
\qquad
F_{jk}=\frac13\omega^{jk},
\qquad
P_{\rm H}=F^\dagger Q_{\rm H}F.
\]
Les opérateurs \(Q_{\rm H},P_{\rm H}\) sont autoadjoints. Leurs bases propres sont mutuellement
non biaisées car \(|F_{jk}|^2=1/9\). Pour tout vecteur unitaire
\(\psi\in\mathcal H_9\), l'inégalité de Robertson donne
\[
\operatorname{Var}_\psi(Q_{\rm H})\,
\operatorname{Var}_\psi(P_{\rm H})
\ge
\frac14\left|\langle\psi,[Q_{\rm H},P_{\rm H}]\psi\rangle\right|^2,
\]
et la borne entropique pour les bases mutuellement non biaisées produit
\[
H_{Q_{\rm H}}(\psi)+H_{P_{\rm H}}(\psi)\ge\log9.
\]
Ici \(H\) est l'entropie de Shannon, avec logarithme naturel, des
distributions spectrales correspondantes. Les relations construites se
résument, sans leur attribuer de causalité physique supplémentaire, par
\[
\begin{gathered}
\text{forme alternée }(X_9,\sigma)
\longrightarrow
\operatorname{Heis}_9,\\
\operatorname{Heis}_9\longrightarrow(X,Z)
\longrightarrow(Q_{\rm H},P_{\rm H})
\longrightarrow\text{incertitude finie}.
\end{gathered}
\]
Il s'agit d'un modèle algébrique fini d'opérateurs ; il ne spécifie pas à lui seul
une dynamique physique ni des unités de mesure.

\section{Dynamique cyclique sur \texorpdfstring{\(\mathbb Z_9\times\mathbb Z_{12}\)}{Z9 x Z12} et quotient prédictif minimal}

Le système dynamique fini \(\mathfrak C_{9,12}\) permet de déterminer la mémoire
minimale nécessaire pour que l'évolution observée soit déterministe. Soit
\[
X=\{(b,m):b\in\mathbb Z/9\mathbb Z, 0\le m<12\}
\]
avec successeur
\[
S(b,m)=
\begin{cases}
(b,m+1),&m<11,\\
(b+1\bmod9,0),&m=11.
\end{cases}
\]
Les étiquettes marginales sont
\[
\beta(b,m)=4b\pmod{12},
\qquad
d(b,m)=1+(m+\beta(b,m)\pmod{12}).
\]
Pour une application d'observation \(q\), le mot futur inclusif d'horizon \(h\) est
\[
Q_h^q(x)=\bigl(q(x),q(Sx),\ldots,q(S^hx)\bigr).
\]

\begin{theorem}[Quotient prédictif minimal]
Soient \(X\) fini, \(T:X\to X\) et \(q:X\to V\). Définissons
\[
\operatorname{Eq}(f):=\{(x,y)\in X^2:f(x)=f(y)\},
\qquad
E_h=\operatorname{Eq}(Q_h^q),
\qquad
E_\infty=\bigcap_{k\ge0}\operatorname{Eq}(q\circ T^k).
\]
Alors
\(E_{h+1}=E_h\cap(T\times T)^{-1}(E_h)\), la suite se stabilise en temps
fini et \(X/E_\infty\) est le quotient ayant le plus petit nombre d'états qui
raffine \(q\) et admet une dynamique déterministe induite.
\end{theorem}

\begin{proof}
\(E_h\) exige l'égalité des observations aux temps \(0,\ldots,h\), tandis que
\((T\times T)^{-1}(E_h)\) exige l'égalité en \(1,\ldots,h+1\). Leur intersection est
\(E_{h+1}\). Comme \(X\) est fini, le nombre de classes ne peut pas croître
indéfiniment. Au premier niveau stable, \(T\) respecte les classes et
passe au quotient. Toute équivalence \(T\)-compatible contenue dans
\(\operatorname{Eq}(q)\)
n'identifie que des états ayant le même futur complet, de sorte qu'elle est
contenue dans \(E_\infty\) ; d'où la propriété minimale.
\end{proof}

L'évaluation exhaustive de \(\mathfrak C_{9,12}\) donne :
\begin{center}
\small
\begin{tabular}{@{}lrrrr@{}}
\toprule
observable & valeurs instantanées & indice \(h\) & remarques & quotient stable\\
\midrule
\(\beta\) & 3 & 11 & 12 & \(C_{36}\), fibra 3\\
\(d\) & 12 & 8 & 9 & \(C_{36}\), fibra 3\\
\((\beta,d)\) & 36 & 0 & 1 & \(C_{36}\), fibra 3\\
\bottomrule
\end{tabular}
\end{center}
Dans les deux observables marginaux,
\[
\operatorname{Eq}(Q_{11}^{\beta})
=\operatorname{Eq}(Q_8^d)
=\operatorname{Eq}(\pi_{36}),
\qquad
\pi_{36}(b,m)=(b\bmod3,m).
\]
Le facteur \(36\) n'est pas introduit comme objectif. Il apparaît comme le quotient
prédictif minimal de chaque marginal. La phase \(m\) est une coordonnée d'état
nécessaire pour prédire ces futurs ; cet énoncé mathématique ne présuppose pas
d'interprétation causale physique. La transformation \(b\mapsto b+3\) est une
symétrie invisible pour ces observables à tout instant.

En tant que système dynamique abstrait, le quotient est un cycle de trente-six
états. Écrire \(C_{36}\) ne choisit toutefois pas d'origine distinguée :
pour obtenir une coordonnée \(\mathbb Z/36\mathbb Z\), il faut la fixer. Cela ne
transforme pas non plus automatiquement ces états en les trente-six paires d'une
nonade. En effet, \(\operatorname{Aut}J(9,2)\cong S_9\). Si les deux points d'une
paire sont dans des cycles distincts de longueurs \(r,s\) d'une permutation,
l'orbite de la paire est de longueur \(\operatorname{mcm}(r,s)\), avec \(r+s\leq9\),
et donc de longueur au plus \(20\). S'ils sont dans un même cycle de longueur
\(r\leq9\), l'orbite est de longueur \(r\), sauf dans le cas antipodal, où elle
est de longueur \(r/2\). Aucun automorphisme n'induit donc d'orbite de longueur
\(36\) sur les paires. Par conséquent, le passage
\[
C_{36}\longleftrightarrow \binom{[9]}2
\]
reste une identification calibrée, non une identification équivariante naturelle.
Même un ordre hamiltonien des paires prouverait une conservation locale
de l'incidence pas à pas, non un automorphisme global ni la canonicité du
dictionnaire.

\section{Nombre HMT comme histoire compatible dans \(X_\infty\) : identité généalogique, structure de lecture et évaluation archimédienne}

Les résultats précédents déterminent l'instance suivante de la définition
directrice. Dans la notation de cette section, un nombre HMT est une histoire
compatible
\[
x\in X_\infty=\varprojlim X_t.
\]
Un lecteur \(L\) est une structure supplémentaire : la paire \((x,L)\) constitue une
présentation lue ou une mesure, mais ne modifie pas le type d'identité de
\(x\). Lorsque \(x\) appartient à \(X_\infty^{\rm Arch}\), sa valeur
archimédienne est l'intersection ponctuelle des cylindres ; l'ensemble des
valeurs ainsi construit est
\(\operatorname{Im}(\operatorname{ev}_{\mathbb R})\subseteq\mathbb R\), et
son identification à toute la droite exige le théorème supplémentaire de couverture.
Sa mémoire minimale relative à un observable est le quotient prédictif des
futurs ; et la forme alternée du support somme--produit admet la
représentation finie de Weyl--Heisenberg construite ci-dessus. La valeur, la classe de trajectoire et l'incertitude sont
définies par trois opérations sur des domaines spécifiés d'une même
architecture de transport.
 \end{HMTScientificOwnerSubordinated}
  \part{Structure discrète du continu et double projection holographique}

\chapter{Système projectif de trajectoires compatibles}
\label{chap:83-03}

\section{Construction généalogique du continu}

\subsection{APP, TRIT et TPK : définitions directrices}

La construction part de trois objets irréductibles et causalement
ordonnés. L'objet \emph{All Possible Paths} (ci-après, \APP)
construit le support nonadique commun, son graphe de tous les chemins finis et
les deux feuillets d'évaluation somme--produit. L'unité tritique (ci-après,
\TRIT) est l'unité locale minimale de relation topologique, géométrique,
informationnelle et, une fois ordonnée, temporelle : sa signature visible possède trois
états et son relèvement conserve la retenue qu'une lecture binaire ou
résiduelle perdrait. Le \emph{Topological Phase Kernel} (ci-après, \TPK)
est la catégorie de transport qui compose ces unités en routes et conserve,
au sein du même état, feuillet, orientation, frontière, registre généalogique, mémoire, phase et
survie.

\begin{center}
\begin{tabularx}{.94\textwidth}{@{}p{.10\textwidth}Y Y@{}}
\toprule
Objet & Construction & Invariantes transportados\\
\midrule
APP & carte positive, tore fini, graphe de Cayley, chemins et feuillets
\(\Sigma,\Pi\) & marque, cellule, voisinage, somme, produit, quotient et clôture\\
TRIT & signature ternaire visible et relèvement topologique--temporel local &
résidu, retenue, feuillet, orientation, seuil et réversion\\
TPK & catégorie de routes, fibres et extensions compatibles & cylindre,
frontière, registre généalogique, mémoire, phase, holonomie et survie\\
\bottomrule
\end{tabularx}
\end{center}

La construction exacte d'APP est établie sur un segment marqué et une
carte finie ; la droite réelle apparaît ensuite comme publication limite. L'intervalle
de référence, ses marques et le chemin orienté qui les parcourt sont :
\begin{equation}
I_{10}=[0,10],\qquad
\mathcal M_{10}=I_{10}\cap\mathbb Z=\{0,1,\ldots,10\},
\qquad
\mathsf P_{10}:0\longrightarrow1\longrightarrow\cdots
\longrightarrow9\longrightarrow10,
\label{p12-83-c3-s1:eq:app-marked-segment}
\end{equation}
Les extrémités fixent le bord et les neuf marques intérieures
\(I_{10}^{\circ}\cap\mathbb Z\) forment l'alphabet positif. L'objet
mathématique exact de départ est la carte
\begin{equation}
\mathcal D_9=\{1,\ldots,9\},
\qquad \mathcal D_8=\{1,\ldots,8\},
\qquad
\lambda_9:\mathcal D_9\xrightarrow{\sim}\mathbb Z/9\mathbb Z,
\qquad \lambda_9(9)=[0]_9.
\label{p12-83-c3-s1:eq:app-positive-chart}
\end{equation}
La marque \(9\) est simultanément le dernier représentant positif et la
publication de la classe nulle. Pour tout \(n\geq1\), la racine numérique positive et
le quotient sont séparés au moyen de
\begin{equation}
\operatorname{dr}_9(n)=1+((n-1)\bmod9),
\qquad q_9(n)=\left\lfloor\frac{n-1}{9}\right\rfloor,
\qquad n=9q_9(n)+\operatorname{dr}_9(n).
\label{p12-83-c3-s1:eq:positive-digital-root-carry}
\end{equation}
Au chemin \(1\to2\to\cdots\to9\) s'ajoute
l'unique arête de clôture \(9\to1\). Le successeur et sa mémoire sont séparés :
\begin{equation}
u_+(a)=\begin{cases}a+1,&a<9,\\1,&a=9,\end{cases}
\qquad
\kappa_+(a)=\begin{cases}0,&a<9,\\1,&a=9,\end{cases}
\qquad a+1=u_+(a)+9\kappa_+(a).
\label{p12-83-c3-s1:eq:app-successor-carry}
\end{equation}
C'est pourquoi \(9\to1\) clôt la phase, augmente la mémoire et prolonge l'état
généalogique.

Le carré de la carte produit le tore fini et son graphe de mouvements :
\begin{equation}
X_9=(\mathbb Z/9\mathbb Z)^2,
\qquad
\Gamma_{\rm APP}=\operatorname{Cay}
\bigl(X_9,\{\pm e_1,\pm e_2\}\bigr),
\qquad
\operatorname{Path}(\Gamma_{\rm APP}).
\label{p12-83-c3-s1:eq:app-cayley-torus}
\end{equation}
L'identification des côtés opposés maintient le cardinal de la carte et ajoute
la loi de continuation. Le noyau de \(8^2=64\) interstices est complété par le gnomon
minimal qui publie la classe nulle dans les deux coordonnées :
\begin{equation}
\mathcal D_9^2=\mathcal D_8^2\sqcup\mathcal G_9,
\qquad
81=64+(9+8)=64+17.
\label{p12-83-c3-s1:eq:app-gnomon-public}
\end{equation}
La même loi se prolonge à des cartes de plus grande profondeur ; l'échelle change,
mais la clôture, le bord partagé et la mémoire de tour conservent leur type.

La complétion volumétrique associée s'exprime par l'identité de
strates
\begin{equation}
10^3=(9+1)^3=729+243+27+1.
\label{p12-83-c3-s1:eq:completion-unit}
\end{equation}
\begin{figure}[tbp]
\centering
\begin{minipage}[t]{.53\textwidth}
\centering
\resizebox{.96\linewidth}{!}{%
\begin{tikzpicture}[
  font=\scriptsize,
  cell/.style={minimum width=4.35mm,minimum height=4.35mm,
    inner sep=0pt,draw=black!28}]
  \node[font=\bfseries] at (1.72,.88)
    {Feuillet multiplicatif de la carte APP : $\Pi(i,j)=\operatorname{dr}_9(ij)$};
  \foreach \i in {1,...,9}{
    \node[font=\bfseries] at (-.35,{-(\i-1)*.44}) {\i};
    \foreach \j in {1,...,9}{
      \pgfmathtruncatemacro{\v}{mod(\i*\j-1,9)+1}
      \pgfmathtruncatemacro{\border}{(\i==9)||(\j==9)}
      \ifnum\border=1
        \node[cell,fill=black!16,font=\bfseries]
          at ({(\j-1)*.44},{-(\i-1)*.44}) {\v};
      \else
        \node[cell,fill=white]
          at ({(\j-1)*.44},{-(\i-1)*.44}) {\v};
      \fi
    }
  }
  \foreach \j in {1,...,9}{
    \node[font=\bfseries] at ({(\j-1)*.44},.38) {\j};
  }
  \draw[very thick] (3.74,.22)--(3.74,-3.74)--(-.22,-3.74);
  \draw[densely dashed] (-.22,.22) rectangle (3.30,-3.30);
  \node[align=center,anchor=north,font=\scriptsize] at (1.76,-3.98)
    {$8^2=64$ cellules intérieures ; gnomon de clôture
     $9+8=17$\\couronne positive de neufs : la marque $9$ publie $[0]_9$};
\end{tikzpicture}}
\end{minipage}\hfill
\begin{minipage}[t]{.44\textwidth}
\centering
\resizebox{.98\linewidth}{!}{%
\begin{tikzpicture}[font=\scriptsize]
  \begin{scope}[x={(.78cm,0cm)},y={(.36cm,.26cm)},z={(0cm,.78cm)}]
    \coordinate (O) at (0,0,0);
    \coordinate (X) at (3.1,0,0);
    \coordinate (Y) at (0,3.1,0);
    \coordinate (XY) at (3.1,3.1,0);
    \coordinate (Z) at (0,0,3.1);
    \coordinate (XZ) at (3.1,0,3.1);
    \coordinate (YZ) at (0,3.1,3.1);
    \coordinate (XYZ) at (3.1,3.1,3.1);
    \path[fill=black!5] (O)--(X)--(XZ)--(Z)--cycle;
    \path[fill=black!12] (X)--(XY)--(XYZ)--(XZ)--cycle;
    \path[fill=black!20] (Z)--(XZ)--(XYZ)--(YZ)--cycle;
    \draw[densely dashed] (O)--(Y)--(XY) (Y)--(YZ);
    \draw[thick] (O)--(X)--(XZ)--(Z)--cycle
      (X)--(XY)--(XYZ)--(XZ) (Z)--(YZ)--(XYZ);
    \node[font=\bfseries\large,fill=white,inner sep=1.5pt]
      at (1.55,.52,1.55) {$9^3=729$};
  \end{scope}
  \node[font=\bfseries] at (3.62,3.68) {Cube d'extrême et moyenne raison holographique};
  \draw[fill=black!8] (0.2,-.78) rectangle (3.15,-.05);
  \node[align=center,font=\tiny] at (1.68,-.415)
    {intérieur\\$9^3=729$};
  \draw[fill=black!16] (3.15,-.78) rectangle (5.15,-.05);
  \node[align=center,font=\tiny] at (4.15,-.415)
    {$3$ caras nuevas\\$3\cdot9^2=243$};
  \draw[fill=black!24] (5.15,-.78) rectangle (6.35,-.05);
  \node[align=center,font=\tiny] at (5.75,-.415)
    {$3$ arêtes\\$3\cdot9=27$};
  \draw[fill=black!34] (6.35,-.78) rectangle (7.02,-.05);
  \node[align=center,font=\tiny] at (6.685,-.415)
    {$1$ sommet};
  \node[font=\bfseries] at (3.61,-1.30)
    {$10^3=729+243+27+1=729+271$};
  \node[align=center] at (3.61,-1.88)
    {couronne de croissance $9^3\subset10^3$ : $271$ états\\
     frontière interne du cube de côté $9$ : $386$ points};
\end{tikzpicture}}
\end{minipage}
\caption{Deux publications finies du support. À gauche : la table de
produit avec racine numérique positive et son gnomon de clôture. À droite : la
complétion cubique EMRH, qui distingue intérieur, faces, arêtes et sommet.
La couronne de croissance $271$ et la frontière interne $386$ sont deux
comptages typés : le premier mesure l'extension $9^3\subset10^3$ et le
second la frontière du cube de côté $9$.}
\label{fig:app-emrh-mini}
\end{figure}
 Le volume intérieur, les faces, les arêtes et la clôture ponctuelle sont ainsi
typés comme des couches distinctes d’une même unité complétée. Sur les 81
cellules de la carte positive agissent les évaluations
\[
\Sigma,\Pi:\mathcal D_9^2\longrightarrow\mathcal D_9,
\qquad
\Sigma(i,j)=\operatorname{dr}_9(i+j),
\qquad
\Pi(i,j)=\operatorname{dr}_9(ij),
\]
Leurs résidus et quotients relevés sont définis par
\[
\begin{aligned}
R^+(i,j)&:=\Sigma(i,j),&
Q^+(i,j)&:=\frac{i+j-R^+(i,j)}9,\\
R^\times(i,j)&:=\Pi(i,j),&
Q^\times(i,j)&:=\frac{ij-R^\times(i,j)}9,
\end{aligned}
\]
et, par conséquent,
\[
\begin{aligned}
\widehat\Sigma(i,j)&:=\bigl(R^+(i,j),Q^+(i,j)\bigr),\\
\widehat\Pi(i,j)&:=\bigl(R^\times(i,j),Q^\times(i,j)\bigr),
\end{aligned}
\qquad
\widehat\Sigma,\widehat\Pi:
\mathcal D_9^2\longrightarrow\mathcal D_9\times\mathbb Z_{\geq0}.
\]
La bijection \(\lambda_9^{\times2}:\mathcal D_9^2\to X_9\) transporte
ces opérations et leurs relèvements sur le tore. APP est donc le tuple
\begin{equation}
\APP=\bigl(\mathcal D_9,\lambda_9,X_9,\Gamma_{\rm APP},
\operatorname{Path}(\Gamma_{\rm APP}),
\Sigma,\Pi,\widehat\Sigma,\widehat\Pi\bigr),
\label{p12-83-c3-s1:eq:app-full-definition-public}
\end{equation}
Le tableau constitue une représentation bidimensionnelle de ce tuple
complet. Le feuillet additif est latin ; le feuillet multiplicatif conserve
le radical ternaire de l'anneau \(\mathbb Z/9\mathbb Z\)
\begin{equation}
R_9:=\mathbb Z/9\mathbb Z,
\qquad
\mathfrak m:=([3]_9)=\{[0]_9,[3]_9,[6]_9\}\triangleleft R_9,
\qquad \mathfrak m^2=(0).
\label{p12-83-c3-s1:eq:app-radical-369}
\end{equation}
La rotation interne est l'automorphisme
\[
\rho:\mathfrak m^3\longrightarrow\mathfrak m^3,
\qquad \rho(x_1,x_2,x_3)=(x_2,x_3,x_1),
\]
dont l'orbite distinguée est
\[
([3]_9,[6]_9,[0]_9)\xrightarrow{\rho}
([6]_9,[0]_9,[3]_9)\xrightarrow{\rho}
([0]_9,[3]_9,[6]_9)\xrightarrow{\rho}
([3]_9,[6]_9,[0]_9).
\]
La réduction directe envoie \(3,6,9\) sur \(0,0,0\). La coordonnée interne
de l'idéal conserve, en revanche, le facteur extrait :
\begin{equation}
\eta_{\mathfrak m}:(\mathfrak m,+)
\xrightarrow{\;\sim\;}(\mathbb Z/3\mathbb Z,+),
\qquad
\eta_{\mathfrak m}([3a]_9):=[a]_3,
\qquad
([3]_9,[6]_9,[0]_9)
\xmapsto{\;\eta_{\mathfrak m}^{\times3}\;}
([1]_3,[2]_3,[0]_3),
\label{p12-83-c3-s1:eq:app-369-120}
\end{equation}
de sorte que le secteur \(3,6,9\) demeure entier et que sa coordonnée interne se
lit comme \(1,2,0\). Cette application appartient au radical d'APP ; la
projection temporelle de TRIT constitue le lecteur calendaire ultérieur. Les
trois classes calendaires
sont séparées comme
\[
\mathcal C_+=\{1,4,7\},\qquad
\mathcal C_-=\{2,5,8\},\qquad
\mathcal C_0=\{3,6,9\}.
\]
L'incompatibilité entre les deux feuillets devient un invariant mesurable. Sur
\(\mathscr A_3=\mathcal D_9^3\), muni de la mesure uniforme
\(\nu_3(x)=9^{-3}\), soit
\(\mathcal H_3=L^2(\mathscr A_3,\nu_3)\). Nous définissons
\[
\Sigma_3,\Pi_3:\mathscr A_3\longrightarrow\mathcal D_9,
\qquad
\Sigma_3(x)=\operatorname{dr}_9(x_1+x_2+x_3),
\qquad
\Pi_3(x)=\operatorname{dr}_9(x_1x_2x_3),
\]
et les projecteurs orthogonaux
\[
P_{\Sigma,3}:=\mathbb E_{\nu_3}[\,\cdot\mid\sigma(\Sigma_3)],
\qquad
P_{\Pi,3}:=\mathbb E_{\nu_3}[\,\cdot\mid\sigma(\Pi_3)]
\in\mathcal B(\mathcal H_3).
\]
L'énumération conjointe des \(9^3\) mots détermine alors
\[
\bigl\|[P_{\Sigma,3},P_{\Pi,3}]\bigr\|=\frac{\sqrt3}{4},
\qquad
B_c=7I+2R=9P_++5P_-,
\qquad
459-405=54.
\]
Le premier nombre réapparaît comme préfacteur électronique ; le deuxième organise
le bloc central ; le troisième enregistre le défaut global somme--produit.

\subsection{TRIT : unité locale topologique et unité de relation temporelle}

HMT fixe d'abord la signature visible
\begin{equation}
\mathrm{TRIT}_{\rm vis}=\{-1,0,+1\},
\qquad \tau\in\mathrm{TRIT}_{\rm vis}.
\label{p12-83-c3-s1:eq:trit-visible-signature}
\end{equation}
Face à l'unité binaire, le TRIT conserve trois régimes locaux. Dans la
distribution uniforme, son information et son entropie se distinguent comme
\begin{equation}
I_{\rm TRIT}=\log3\ \text{nats}=\log_2 3\ \text{bits},
\qquad S_{\rm TRIT}=k_BI_{\rm TRIT}.
\label{p12-83-c3-s1:eq:trit-information-unit}
\end{equation}
Le bit demeure l'unité binaire ; le TRIT est l'unité minimale capable de
distinguer simultanément retour, seuil et ouverture. L'égalité
\eqref{p12-83-c3-s1:eq:trit-information-unit} quantifie l'information du triplet ; le
coût \(k_B\Theta\log3\) appartient au lecteur thermique de réinitialisation ultérieur.
La mécanique dimensionnelle conserve en outre l'antécédent arithmétique de cette signature.
Pour \((i,j)\in\mathcal D_9^2\), la cellule arithmétique relevée associée est
\begin{equation}
c_{\rm MD}(i,j;\tau)=
\bigl(i,j;R^+,Q^+,R^\times,Q^\times;\tau\bigr).
\label{p12-83-c3-s1:eq:trit-complete-local-unit}
\end{equation}
Les résidus \(R^+,R^\times\) sont la publication visible des deux feuillets ;
les quotients \(Q^+,Q^\times\) conservent la mémoire exacte de la division
euclidienne ; et \(\tau\) oriente la transition. La signature visible occupe le premier
niveau ; la cellule qui la porte ajoute les antécédents arithmétiques, et l’état
holonomique intégral complète ensuite la frontière, le chemin, la phase et la survie.

Il convient de séparer quatre strates. La signature visible est \(\tau\) ; le couple
résidu--retenue est son relèvement local ; \(c_{\rm MD}\) est la cellule
arithmétique d'APP munie de cette signature ; et l'état TPK ajoute extension,
frontière, registre généalogique, phase et survie. Cette hiérarchie conserve l'identité
propre de la coordonnée observée, de son relèvement, de la cellule porteuse et de
l'histoire enrichie qui la transporte.

Dans la carte ternaire équilibrée, l'entier local relevé
\(n\in\{-4,-3,\ldots,3,4\}\) se décompose de manière unique comme
\begin{equation}
\widehat n=(r,\kappa),
\qquad n=r+3\kappa,
\qquad r,\kappa\in\{-1,0,+1\}.
\label{p12-83-c3-s1:eq:trit-balanced-lift}
\end{equation}
La signature visible est \(\tau=r\) ; \(\kappa\) conserve la retenue que la projection
\(n\mapsto r\) oublierait.
La fibre au-dessus du zéro visible contient trois états distincts,
\[
0^+=(0,+1),\qquad 0^0=(0,0),\qquad 0^-=(0,-1),
\]
car \(n=-3,0,3\) possèdent le même résidu visible et des retenues distinctes. C'est
le relèvement topologique minimal de MD : une transition orientée d'APP
est munie de l'information nécessaire pour l'inverser et pour décider quelles
prolongations demeurent compatibles.

La signature possède en outre une règle de mémoire de feuillet. Si \(m_k\) est le mode
arithmétique actif, alors
\begin{equation}
m_{k+1}=
\begin{cases}
\Sigma,&\tau_k=+1,\\
\Pi,&\tau_k=-1,\\
m_k,&\tau_k=0.
\end{cases}
\label{p12-83-c3-s1:eq:trit-zero-memory}
\end{equation}
Le zéro code la rétention du feuillet précédent. C'est pourquoi le mot
\((+1,-1,0)\) publie la séquence \((\Sigma,\Pi,\Pi)\) : le troisième état
conserve le mode multiplicatif actif.

La relation temporelle s'obtient en paramétrant ces cellules de manière ordonnée.
Chaque signature porte une algèbre quadratique
\begin{equation}
\mathbb A_\tau=\mathbb R[X]/(X^2+\tau),
\qquad J_\tau=[X],
\label{p12-83-c3-s1:eq:trit-quadratic-algebra-local}
\end{equation}
de sorte que
\[
\begin{array}{c@{\qquad}c@{\qquad}l}
\toprule
\tau & J_\tau^2 & \text{relation temporelle ordonnée}\\
\midrule
+1 & -I & \text{retour elliptique, mémoire et passé},\\
0  & 0  & \text{seuil parabolique et présent},\\
-1 & +I & \text{ouverture hyperbolique bidirectionnelle et futur}.\\
\bottomrule
\end{array}
\]
Le TRIT précède le temps extérieur : il classe d'abord le type algébrique et
topologique du transport ; ensuite, un paramétrage ordonné le publie
comme relation temporelle. La double direction est typée par l'involution
exacte
\begin{equation}
\mathcal C_{\rm rev}(\tau_1,\ldots,\tau_n)
=(-\tau_n,\ldots,-\tau_1),
\qquad \mathcal C_{\rm rev}^2=I,
\label{p12-83-c3-s1:eq:trit-time-reversal}
\end{equation}
qui échange retour et ouverture et conserve le seuil. Dans le feuillet
hyperbolique, \(P_\pm=(I\pm J_{-1})/2\) séparent les deux directions propres et
la réversion échange leurs orientations. Topologie et temps sont ainsi les
deux lectures coordonnées du TRIT.

L'involution combinatoire \(\mathcal C_{\rm rev}\) et une conjugaison
opératorielle \(C\) appartiennent cependant à des domaines distincts. Dans une
réalisation dynamique séparée, munie de \(H=H^*\),
\(CJ_EC^{-1}=-J_E\) et \(CHC^{-1}=H\), nous fixons une échelle d'action interne
\(\hbar_{\rm int}>0\) et écrivons
\(h_{\rm int}:=2\pi\hbar_{\rm int}\). Alors
\begin{equation}
U(t)=\exp(-J_EHt/\hbar_{\rm int}),
\qquad CU(t)C^{-1}=U(-t).
\label{p12-83-c3-s1:eq:trit-bidirectional-propagator}
\end{equation}
L'entrelaceur qui transporte les mots TRIT vers la réalisation dynamique fixe
la correspondance entre \(\mathcal C_{\rm rev}\) et \(C\). La double
orientation locale fournit le domaine de cette construction.

\subsection{TPK : catégorie de transport et évolution de phase}

Le \emph{Topological Phase Kernel} est la catégorie de transport sur la
base observable d’APP qui compose les chiffres, les mots, les opérateurs et les relations
de prolongation dans un même état holonomique intégral. Son architecture interne se répartit en
huit classes opératoires, différenciées par leur domaine, leur codomaine et leur fonction :
\begin{enumerate}[label=\textup{(\roman*)},leftmargin=2.25em]
\item support observable et état holonomique intégral ;
\item sélecteur interne de feuillet \(M_{\rm ph}:\{1,\ldots,9\}\to\{+1,-1,0\}\) ;
\item translations cardinales et évolution visible \(F_{\rm obs}\) ;
\item lecteurs et quotients nonadique, tritique et de bloc \(1000\) ;
\item mémoire, retenue, cylindre, frontière et registre généalogique ;
\item périodicité et retour
\(C_{12}\hookrightarrow C_{108}\twoheadrightarrow C_9\),
\(\mathcal K_{27}\) et canaux \(90/120\) ;
\item prolongations \(\operatorname{Ext}_r\), cylindres, survie et
extension naturelle bidirectionnelle ;
\item foncteurs de sortie archimédienne, exceptionnelle étalonnée et dimensionnelle.
\end{enumerate}
Le retour de 27 pas est le morphisme composé
\begin{equation}
\mathcal K_{27}:=F_{\rm obs}^{27},
\qquad \mathcal K_{27}^{4}=F_{\rm obs}^{108}.
\label{p12-83-c3-s1:eq:tpk-kernel-27}
\end{equation}
Son quatrième itéré restaure la phase observable du calendrier, tandis que la
mémoire et l'état TPK complet poursuivent leur trajectoire ;
\(\mathcal K_{27}\) est le morphisme de retour interne de \(C_{108}\).
Le calendrier \(C_{108}\), la connexion nonadique, l'extension bilatérale et
les lecteurs occupent des places typées au sein de cette architecture. Soit
\(\mathsf P_{\rm APP}\) la catégorie
libre du graphe dirigé d'APP et
\(\mathsf P_{\rm APP}^{(2)}=\mathsf P_{\rm APP}\times
\mathsf P_{\rm APP}\) sa lecture simultanée somme--produit. Soit en outre
\[
\mathcal Q_{\rm obs}
=X_9^2\times\{N,E,S,O\}^2\times\Theta_{108},
\qquad \Theta_{108}=\mathbb Z/108\mathbb Z,
\]
où, pour le représentant
\(\widetilde\theta\in\{0,\ldots,107\}\),
\[
r(\theta)=1+(\widetilde\theta\bmod9),
\qquad g_0(\theta)=[\widetilde\theta]_9,
\]
\[
\tau(\theta)=
\begin{cases}
+1,&r(\theta)\in\{1,4,7\},\\
-1,&r(\theta)\in\{2,5,8\},\\
0,&r(\theta)\in\{3,6,9\},
\end{cases}
\qquad
\mu(\theta)=
\begin{cases}
\Sigma,&\tau(\theta)=+1,\\
\Pi,&\tau(\theta)=-1,\\
\varnothing,&\tau(\theta)=0.
\end{cases}
\]
L'étiquette calendaire \(\mu=\varnothing\) indique le seuil, et la mémoire
dynamique \(m_k\) de \eqref{p12-83-c3-s1:eq:trit-zero-memory} retient le dernier feuillet
actif. Une configuration
\(q=(p^\Sigma,p^\Pi,d^\Sigma,d^\Pi,\theta)\) conserve les deux positions,
les deux directions cardinales et l'indice du calendrier. L'évolution
visible \(F_{\rm obs}\) active le feuillet indiqué par TRIT, transporte la
position correspondante et fait avancer \(\theta\). Notons
\(\mathsf P_{\rm obs}\) la catégorie libre engendrée par les flèches
\(q\to F_{\rm obs}(q)\). Chaque flèche induit un couple de chemins d'APP et
définit le foncteur de base
\begin{equation}
B:\mathsf P_{\rm obs}\longrightarrow\mathsf P_{\rm APP}^{(2)}.
\label{p12-83-c3-s1:eq:tpk-bivariant-base-functor}
\end{equation}

Au-dessus de chaque \(q\) est disposée la fibre typée
\begin{equation}
\mathcal F_q=
\mathsf O_q\times\mathsf H_q\times\mathsf C_q\times\mathsf S_q
\times\mathsf B_q\times\mathsf\Lambda_q,
\label{p12-83-c3-s1:eq:tpk-fiber}
\end{equation}
Ses six facteurs conservent respectivement orientation et feuillet ; résidu et
quotient ; retenue ; signature ; cylindre et frontière ; et registre ordonné de
transport. Mémoire et survie sont des propriétés de la composition et des
relations de prolongation.

Pour une route \(r:q\to q'\), la relation
\begin{equation}
\operatorname{Ext}_r\subseteq\mathcal F_q\times\mathcal F_{q'}
\label{p12-83-c3-s1:eq:tpk-extension-relation}
\end{equation}
contient toutes les continuations admissibles et satisfait
\[
\Delta_q\subseteq\operatorname{Ext}_{\operatorname{id}_q},
\qquad
\operatorname{Ext}_{r_2}\circ\operatorname{Ext}_{r_1}
\subseteq\operatorname{Ext}_{r_2\circ r_1}.
\]
Les transports respectent les identités et la composition :
\begin{equation}
\begin{gathered}
\mathsf H(\operatorname{id}_q)=I,
\quad \mathsf C(\operatorname{id}_q)=I,
\quad \mathsf\Lambda(\operatorname{id}_q)=I,
\quad \kappa(\operatorname{id}_q)=0,\\
\mathsf H(r_2\!\circ r_1)=\mathsf H(r_2)\mathsf H(r_1),
\quad
\mathsf C(r_2\!\circ r_1)=\mathsf C(r_2)\mathsf C(r_1),
\quad
\mathsf\Lambda(r_2\!\circ r_1)=
\mathsf\Lambda(r_2)\mathsf\Lambda(r_1).
\end{gathered}
\label{p12-83-c3-s1:eq:tpk-functorial-transport}
\end{equation}
Outre ces lois fonctorielles, les coordonnées transportées obéissent à
\begin{equation}
\begin{aligned}
h'&=\mathsf H(r)(h),\\
c'&=\mathsf C(r)(c)+\kappa(r),\\
\lambda'&=\mathsf\Lambda(r)(\lambda),\\
\sigma'&=\operatorname{Sig}_{q'}(h',c',\lambda'),
\end{aligned}
\qquad
\kappa(r_2\circ r_1)
=\kappa(r_2)+\mathsf C(r_2)\kappa(r_1).
\label{p12-83-c3-s1:eq:tpk-typed-transport}
\end{equation}
La dernière identité est le cocycle de retenue qui rend associative la
composition avec mémoire.

Le TPK est alors la catégorie \(\mathsf{TPK}\) dont les objets sont
\begin{equation}
x=(q,o,h,c,\sigma,C,b,\lambda),
\qquad (o,h,c,\sigma,(C,b),\lambda)\in\mathcal F_q,
\label{p12-83-c3-s1:eq:tpk-enriched-object}
\end{equation}
et dont les morphismes \(x\to x'\) sont les triplets \((r,x,x')\) tels que
les coordonnées satisfont \eqref{p12-83-c3-s1:eq:tpk-typed-transport} et
\((f_x,f_{x'})\in\operatorname{Ext}_r\), où
\(f_x\in\mathcal F_q\) et \(f_{x'}\in\mathcal F_{q'}\). La composition est
\begin{equation}
(r_2,x',x'')\circ(r_1,x,x')
=(r_2\circ r_1,x,x''),
\label{p12-83-c3-s1:eq:tpk-morphism-composition}
\end{equation}
\Needspace{7\baselineskip}
l'identité est \((\operatorname{id}_q,x,x)\), et le foncteur
\begin{equation}
p:\mathsf{TPK}\longrightarrow\mathsf P_{\rm obs},
\qquad p(x)=q,
\qquad p(r,x,x')=r,
\label{p12-83-c3-s1:eq:tpk-base-functor}
\end{equation}
publie la configuration visible sans identifier les états d'une même fibre.
La définition fixe une catégorie au-dessus d'une base ; elle ne présuppose pas qu'il s'agisse d'une
fibration catégorique ni que toute flèche admette un relèvement cartésien.
Les projections locales sont séparées par type :
\[
\pi_{\rm TRIT}:\operatorname{Ob}(\mathsf{TPK})
\longrightarrow\mathrm{TRIT}_{\rm vis},
\qquad
\pi_{\rm TRIT}(x)=\tau(\theta),
\]
\[
\pi_{\rm loc}(x)=
\bigl(\pi_{\rm TRIT}(x),\mu(\theta),o,h,c,
\sigma,\lambda,g_0(\theta)\bigr).
\]
La première ne retrouve que la signature TRIT visible. La seconde conserve
en outre l’étiquette calendaire du feuillet, l’orientation, la paire
résidu--quotient, la retenue, la signature de transport, le registre généalogique et la phase,
mais oublie encore le cylindre, son étiquette de bord et d’autres coordonnées de
\(q\). C’est pourquoi ni le TRIT visible ni son registre local
ne remplacent l’état holonomique intégral du TPK.

Le nom réunit trois fonctions inséparables. \emph{Topological} renvoie à
la catégorie de chemins, à l'incidence, au bord et à la prolongation par
cylindres ; \emph{Phase} renvoie à l'indice nonadique, au calendrier orienté et
à l'holonomie ; \emph{Kernel} désigne le noyau persistant de transport qui
demeure avant de contracter l'état au moyen d'un lecteur. Formellement,
l'objet est la catégorie précédente : une topologie supplémentaire sur les objets ou
les morphismes, ou son identification au noyau linéaire d'une application, exige
des données indépendantes.

En conséquence, deux routes ayant la même extrémité visible peuvent demeurer
des objets distincts par leur feuillet, quotient, retenue, frontière, cylindre, phase ou
registre. Le TPK ne choisit pas une sortie scalaire : il conserve l'antécédent commun
dont partent les publications archimédienne, exceptionnelle et dimensionnelle.

% Estas consecuencias se publican después de la construcción correlativa y
% del paso al límite de los capítulos 4 y 5. Se definen aquí para preservar
% íntegramente su procedencia, pero no se ejecutan antes de que exista el
% continuo al que se aplican.
\long\def\HMTDeferredContinuumConsequences{%
\section{Conséquences et publications du continu construit}

\subsection{Théorème d'assemblage terminal des constructions corrélatives engendrées par l'état commun}

Chaque profondeur \(k\) possède un état holonomique intégral \(\mathcal E_k\) avec
toutes ses émissions survivantes et un unique constructeur corrélatif
\begin{equation}
\mathcal O_k:\mathcal E_k\longrightarrow\mathcal Z_k,
\qquad
\widehat x_k\longmapsto
\mathcal O\bigl(\{\mathfrak e_\alpha:
\alpha\in\operatorname{Ch}_k(\widehat x_k)\}\bigr).
\label{p12-83-c3-s1:eq:five-features-single-constructor}
\end{equation}
La valeur de \(\mathcal O_k\) conserve simultanément les coordonnées
spectrale, solénoïdale, de jauge, cohérente et déterminantale,
\begin{equation}
(\mathrm{sp},\mathrm{sol},\mathrm{gau},\mathrm{coh},\mathrm{det}),
\label{p12-83-c3-s1:eq:five-features-coordinates}
\end{equation}
comme cinq caractéristiques du même état et du même registre généalogique. Les
troncatures
\[
u^{\mathfrak e}_{\ell,k}:\mathcal E_\ell\to\mathcal E_k,
\qquad
u^{\mathcal O}_{\ell,k}:\mathcal Z_\ell\to\mathcal Z_k
\]
satisfont l'unique naturalité
\begin{equation}
u^{\mathcal O}_{\ell,k}\circ\mathcal O_\ell
=\mathcal O_k\circ u^{\mathfrak e}_{\ell,k},
\qquad \ell\ge k.
\label{p12-83-c3-s1:eq:five-features-naturality}
\end{equation}
C'est pourquoi l'action corrélative passe à la limite sous la forme
\begin{equation}
\mathcal O_\infty:
\mathcal C_{\rm cont}^{\rm disc}:=\varprojlim_k\mathcal E_k
\longrightarrow
\mathcal Z_\infty:=\varprojlim_k\mathcal Z_k.
\label{p12-83-c3-s1:eq:five-features-infinite-operator}
\end{equation}
Les cinq applications terminales sont des vues intrinsèques de cette unique valeur :
\begin{equation}
\boxed{
\mathfrak G_T^{\rm int}:=
\operatorname{View}_T\!\left(
\mathcal O_\infty(\mathcal C_{\rm cont}^{\rm disc})\right),
\quad
T\in\{\mathrm{sp},\mathrm{sol},\mathrm{gau},
\mathrm{coh},\mathrm{det}\}.}
\label{p12-83-c3-s1:eq:five-features-views}
\end{equation}
La simultanéité précède les résolutions spécialisées : chaque vue expose
la coordonnée requise et l'état complet conserve les quatre autres.

\subsection{Survie et limite projective}

Soit \(\operatorname{Surv}^{(\chi)}_N\) l'ensemble des histoires de
profondeur \(N\) qui satisfont simultanément les lois actives du canal
\(\chi\) et admettent une prolongation. Nous définissons les troncatures compatibles
\[
\tau_{N+1,N}:\operatorname{Surv}^{(\chi)}_{N+1}
\longrightarrow\operatorname{Surv}^{(\chi)}_N
\]
Celles-ci conservent le préfixe enrichi. Le continu généalogique est la limite
\begin{equation}
X_\chi=\varprojlim_N
\bigl(\operatorname{Surv}^{(\chi)}_N,\tau_{N+1,N}\bigr),
\qquad
\tau_{N+1,N}(x_{N+1})=x_N.
\label{p12-83-c3-s1:eq:inverse-survival}
\end{equation}
L'espace global des sections survivantes et l'une de ses sections
enrichies sont notés respectivement
\begin{equation}
X_\infty:=\bigsqcup_\chi X_\chi,
\qquad
\widehat\Omega_\infty\in X_\infty.
\label{p12-83-c3-s1:eq:global-survival-space}
\end{equation}
Le lecteur archimédien global est
\begin{equation}
P_{\rm arq}:=\operatorname{ev}_{\mathbb R}:
X_\infty\longrightarrow\mathbb R;
\label{p12-83-c3-s1:eq:archimedean-reader}
\end{equation}
sa restriction à chaque canal évalue l'intersection réduite à un point des cylindres
compatibles de la section correspondante.
Lorsque le lecteur exceptionnel intervient, \(X_\infty^{\rm cal}\subseteq
X_\infty\) désigne le sous-espace dans lequel le type d'étalonnage
de \eqref{p12-83-c3-s1:eq:exceptional-calibration} a été fixé.
La survie détermine la section par compatibilité interne : chaque
histoire doit pouvoir être prolongée dans le système complet. Le lecteur numérique
intervient ensuite pour certifier et publier l'intersection du
descendant déjà compatible.

\subsection{Connexion nonadique conjointe, holonomie et mémoire verticale}

À la profondeur \(K\), l'holonomie du retour nonadique est initialement la
correspondance typée
\begin{equation}
\Gamma_{9,K}=\operatorname{Hol}_{C_9}(\nabla_{\rm TPK}):
\widehat{\mathcal X}_K\rightrightarrows
\widehat{\mathcal X}_{K+9}.
\label{p12-83-c3-s1:eq:gamma9}
\end{equation}
Ses itérés sont des compositions relationnelles. Ce n'est que lorsque chaque état possède
une unique prolongation compatible qu'elle se restreint à une application, et ce n'est que lorsque cette
application est inversible qu'elle constitue une monodromie. Cette distinction conserve toutes
les extensions survivantes avant de sélectionner une section.
Après neuf phases, la phase locale revient, tandis que l'état complet avance :
\[
g(K+9)=g(K),\qquad q_{{\rm mem},K+9}=q_{{\rm mem},K}+1,
\qquad \widehat x_{K+9}\ne\widehat x_K\ \text{en général}.
\]
Bloc, cylindre, préfixe, retenue et mémoire continuent d'évoluer. Telle
est la différence entre périodicité visible et retour avec biographie.

L'état transporté au niveau (K) peut être présenté comme
\begin{equation}
\widehat x_K=
(B_K,C_K,p_K,c_K,\Sigma_K,W_K,g(K),q_{{\rm mem},K},s_K),
\label{p12-83-c3-s1:eq:enriched-state-k}
\end{equation}
où chaque coordonnée possède une loi de transport propre. Le retour de phase
est représenté sur le catalogue
\[
\mathcal U_6^{\rm cat}\cong\mathfrak S_+\times\mathfrak S_\times,
\qquad |\mathfrak S_+|=18,
\qquad |\mathfrak S_\times|=26.
\]
Les moyennes conditionnelles
\[
(E_+f)(s,e)=\frac1{18}\sum_{s'\in\mathfrak S_+}f(s',e),
\qquad
(E_\times f)(s,e)=\frac1{26}\sum_{e'\in\mathfrak S_\times}f(s,e')
\]
commutent. Le mot de phase
\(\tau=(+1,-1,0,+1,-1,0,+1,-1,0)\) applique
\(E_+,E_\times,I\) dans cet ordre, et un tour se clôt au moyen de
\begin{equation}
\mathcal K_{\rm ph}^{\,9}
=(E_+E_\times)\otimes I_{\ell^2(C_9)},
\qquad
\bigl\langle\delta_u,E_+E_\times\delta_v\bigr\rangle
=\frac1{468}\quad(u,v\in\mathcal U_6^{\rm cat}).
\label{p12-83-c3-s1:eq:joint-nine-holonomy}
\end{equation}
Les feuillets somme et produit n'agissent pas comme des filtres indépendants : ce sont deux
composantes d'une unique holonomie qui ne se clôt qu'après avoir parcouru le
cycle complet.

La colligation compression--expression sépare ce qui revient de ce qui demeure
comme mémoire. Si (J) plonge l'espace des histoires dans l'espace de
phase, si (T) est la compression survivante et \(\eta\) le canal de défaut,
alors
\begin{equation}
\boxed{
U_{\Gamma_9}J=JT+\eta,
\qquad J^*\eta=0,
\qquad I-T^*T=\eta^*\eta.}
\label{p12-83-c3-s1:eq:compression-expression}
\end{equation}
La première égalité décompose le transport en retour comprimé et
expression de mémoire ; la deuxième prouve leur orthogonalité ; la troisième mesure
exactement l'information qui ne tient pas dans l'ombre périodique. L'affirmation
centrale est ainsi formalisée : c'est la phase qui revient, non l'état complet.

La mémoire verticale s'organise au moyen de l'extension cyclique orientée
\begin{equation}
0\longrightarrow C_{12}\xrightarrow{\,\iota\,}C_{108}
\xrightarrow{\,p\,}C_9\longrightarrow0,
\qquad
c(r,r')=\left[\left\lfloor\frac{r+r'}9\right\rfloor\right]_{12}.
\label{p12-83-c3-s1:eq:c108-extension}
\end{equation}
Ici \(\iota([b]_{12})=[9b]_{108}\) et
\(p([\theta]_{108})=[\theta]_9\). L'extension n'est pas scindée :
\(C_{108}\) est d'exposant 108, tandis que
\(C_{12}\times C_9\) est d'exposant 36.
Le cocycle enregistre la retenue de chaque tour de la base \(C_9\) dans la
mémoire dodécaphasique \(C_{12}\). \(C_{108}\) est donc un calendrier
orienté de phase et de mémoire. Ce n'est qu'après la carte dimensionnelle que cette
extension se réalise comme une horloge quantique finie de 108 pas.

Le relèvement local de MD qui porte la signature TRIT est l'unité
topologique--temporelle de ce cristal généalogique.
Sur
\(\mathcal H_{\rm clk}=\mathbb C^{108}\), de base
\(\{|j\rangle:j\in\mathbb Z/108\mathbb Z\}\), soient
\[
\zeta=e^{2\pi\ii/108},\qquad
|\widetilde k\rangle=\frac1{\sqrt{108}}
\sum_{j=0}^{107}\zeta^{jk}|j\rangle,
\qquad U_{\rm clk}|j\rangle=|j+1\rangle.
\]
\Needspace{12\baselineskip}
Avec \(\omega_0=2\pi/(108t_0)\), le hamiltonien
\begin{equation}
H_{\rm clk}=\hbar_{\rm int}\omega_0
\sum_{k=0}^{107}k\,|\widetilde k\rangle\langle\widetilde k|
\label{p12-83-c3-s1:eq:trit-time-crystal-hamiltonian}
\end{equation}
engendre exactement
\begin{equation}
U_{\rm step}=\exp\!\left(
-\frac{\ii H_{\rm clk}t_0}{\hbar_{\rm int}}\right)=U_{\rm clk},
\qquad U_{\rm step}^{108}=I,
\label{p12-83-c3-s1:eq:trit-time-crystal-period}
\end{equation}
Si \(Z_{\rm clk}|j\rangle=\zeta^j|j\rangle\), alors
\begin{equation}
\langle j_0|U_{\rm step}^{-n}Z_{\rm clk}U_{\rm step}^{n}|j_0\rangle
=\zeta^{j_0+n},
\label{p12-83-c3-s1:eq:trit-time-crystal-observable}
\end{equation}
dont l’orbite possède la période primitive 108. Dans le vocabulaire HMT, c’est le
cristal temporel généalogique interne : le calendrier observable se ferme,
\(C_{12}\) enregistre le tour et l’état holonomique intégral ne doit pas nécessairement
revenir. La réalisation ultérieure comme phase physique macroscopique exige,
en outre, un protocole de Floquet, une réponse sous-harmonique primitive, de la robustesse,
de la persistance et un mécanisme de protection ; cette carte ne modifie pas le théorème
fini de périodicité unitaire.

\subsection{Publication archimédienne, nombre généalogique et mesure}

La publication archimédienne utilise des cylindres ternaires semi-ouverts
\begin{equation}
I_{N,L}=\left[\frac{N}{3^L},\frac{N+1}{3^L}\right),
\qquad \operatorname{diam}I_{N,L}=3^{-L}\longrightarrow0.
\label{p12-83-c3-s1:eq:ternary-cylinders}
\end{equation}
Une section compatible détermine des cylindres emboîtés et, par complétude, un
unique point réel. Le nombre HMT est la section elle-même,
\begin{equation}
x=(x_N)_N\in X_\chi
=\varprojlim_N\operatorname{Surv}^{(\chi)}_N,
\label{p12-83-c3-s1:eq:hmt-number-as-section}
\end{equation}
non le tuple formé par elle et ses évaluations. La valeur apparaît lorsqu'un
lecteur oublie une partie de la généalogie.
La précision est donc profondeur : mesurer signifie descendre à un
cylindre plus fin tandis que le préfixe causal demeure dans l'antécédent.
La droite réelle est reconstruite comme quotient archimédien d'un espace
d'histoires beaucoup plus riche.

\subsection{Réalisation dimensionnelle de l'état HMT : application d'écran, grandeur observable et carte spectrale}

La mécanique dimensionnelle commence là où l'état généalogique se couple à une
frontière de réalisation. Un corps \(K_n\) étant fixé, soit \(M_n(x)\) le module
engendré par les canaux de prolongation d'un état \(x\), soit
\(\mathcal R_n(x)\) le sous-module de ses relations et soit \(\beta_n(x)\) la
donnée de bord conservée. L'écran dimensionnel est
\begin{equation}
\Sigma_n(x)=\bigl(K_n,M_n(x),\beta_n(x),\mathcal R_n(x)\bigr),
\qquad
d_{\Sigma_n}(x)=\dim_{K_n}\frac{M_n(x)}{\mathcal R_n(x)}.
\label{p12-83-c3-s1:eq:md-screen}
\end{equation}
Si \(M_n(x)\simeq K_n^{g_n}\) et si les colonnes de \(R_n(x)\) engendrent les
relations, alors
\begin{equation}
\boxed{d_{\Sigma_n}(x)=g_n-\operatorname{rank}_{K_n}R_n(x).}
\label{p12-83-c3-s1:eq:md-rank}
\end{equation}
La dimension n'est pas adjointe comme étiquette : c'est le nombre de canaux
indépendants qui survivent après l'imposition des relations de frontière.

Une grandeur physique est une section d'une droite dimensionnelle \(L_D\) ; une
unité est une carte \(\chi_u:L_D\to\mathbb R\) qui publie sa coordonnée.
C'est pourquoi
\begin{equation}
\text{état TPK}\longrightarrow\Sigma_n(x)
\longrightarrow s_D\in\Gamma(L_D)
\xrightarrow{\;\chi_u\;}[s_D]_u\in\mathbb R
\label{p12-83-c3-s1:eq:md-measurement-chain}
\end{equation}
sépare exactement structure, type dimensionnel et nombre mesuré. Mesurer
consiste à coupler une section à un appareil réversible, à imposer la
compatibilité de bord et à appliquer le lecteur ; changer d'unités change
\(\chi_u\), non la section ni sa généalogie.

\subsection{Double projection mathématique et réalisation dimensionnelle}

La branche exceptionnelle conserve des données de jauge dont la publication réelle n'a pas
besoin. Soit \(\mathscr C_{\rm exc}\) l'espace des étalonnages et soit
\begin{equation}
\mathfrak E:
X_\infty^{\rm cal}\times\mathscr C_{\rm exc}
\longrightarrow\mathbf{Exc}
\label{p12-83-c3-s1:eq:exceptional-calibrated-reader}
\end{equation}
le lecteur d'incidence. Nous fixons une fois
\begin{equation}
\begin{aligned}
\chi_{\rm exc}=\bigl(&\text{ordre projectif},\ \text{calibre de Paley},
 \ \text{dodécade},\\
 &\text{alignement d’incidence},\ \text{chambre réticulaire}\bigr)
 \in\mathscr C_{\rm exc}
\end{aligned}
\label{p12-83-c3-s1:eq:exceptional-calibration}
\end{equation}
et abrégeons
\(\mathfrak E_{\chi_{\rm exc}}(x):=\mathfrak E(x,\chi_{\rm exc})\).
La jauge initiale est transportée par l'histoire ; elle n'est pas choisie de nouveau à
chaque profondeur.

La même limite étalonnée admet alors deux lecteurs de codomaines
distincts :
\begin{equation}
\mathbb R
\xleftarrow{\;\operatorname{ev}_{\mathbb R}\;}
X_\infty^{\rm cal}
\xrightarrow{\;\mathfrak E_{\chi_{\rm exc}}\;}
\mathbf{Exc}.
\label{p12-83-c3-s1:eq:double-projection}
\end{equation}
La première branche contracte les cylindres jusqu'à une valeur continue. La seconde
conserve incidences, orientations, jauge et frontière. Toutes deux descendent en
parallèle du même état source. Cette bifurcation explique conjointement la
continuité métrique et l'émergence de symétries exceptionnelles. La seconde
branche est une reconstruction étalonnée : elle est exacte une fois
\(\chi_{\rm exc}\) fixé, mais l'état nu ne sélectionne pas rétrospectivement
les cinq données qui composent cette jauge.

La réalisation dimensionnelle ne sérialise pas ces deux branches. C'est un troisième lecteur
du même antécédent :
\begin{equation}
P_{\rm dim}:X_\infty\longrightarrow
\mathcal R_{\rm MD},
\qquad
\mathcal R_{\rm MD}=
\{\text{action, masse, champ, géométrie, information}\}.
\label{p12-83-c3-s1:eq:dimensional-publication}
\end{equation}
De cette manière, un nombre, une symétrie et une grandeur physique peuvent partager
une généalogie sans se confondre ni se réduire les uns aux autres.

La thèse du continu est ainsi formulée avec précision : la réalité
archimédienne est l'ombre de rayon nul d'histoires discrètes infiniment
raffinables ; ce que l'ombre omet réapparaît comme mémoire, incidence,
géométrie et structure physique.

% RESULTADO_RECUPERADO. Owner íntegro de la dinámica del cociente, el retorno
% observable y los cociclos de modo y orientación. Se subordina a la única
% construcción APP--TRIT--TPK sin abrir un origen alternativo.
\begin{HMTScientificOwnerSubordinated}
\chapter{Extensions survivantes, monodromie et dynamique du quotient}
\label{chap:extensiones-dinamica-cociente}
\index{survie!extension globale}
\index{monodromie!retour nonadique}
\index{mesure!catalogue d’émissions}

Le langage des chemins nécessite une précision décisive. Un chemin
d'APP, une émission visible, une histoire enrichie de TPK et une section
globale ne sont pas quatre noms pour le même objet. Ils constituent quatre niveaux
consécutifs d'une même construction. Le chemin conserve l'ordre
positionnel ; l'émission conserve une lecture finie ; l'histoire ajoute la signature,
la retenue, l'orientation, la frontière et la mémoire ; la section globale exige que tous
ses préfixes survivent simultanément à toute profondeur.

Cette hiérarchie permet de séparer deux questions qui étaient restées
entremêlées. La première est logique : sous quelles hypothèses un préfixe fini
appartient réellement à une extension infinie compatible. La seconde est
dynamique : ce que signifie transporter une seule histoire et ce que signifie
transférer simultanément toutes ses continuations compatibles. La
distinction est essentielle dans APP\@. Une trajectoire déterminée mène un état à
un autre ; un transfert de chemins mène une distribution d'extensions à
une autre distribution.

Cette section résout le problème fini de profondeur six. La
chronologie déterministe publiée possède un retour visible et ne peut pas sélectionner
une probabilité unique. En revanche, le transport biaisé par la phase nonadique,
lorsque chaque canal actif est interprété comme le renouvellement normalisé de toute
sa fibre de continuations compatibles, possède un retour de neuf pas
uniforme sur les (468) émissions. Les poids (1/3) cessent alors d'être
un choix extérieur : ils proviennent des trois apparitions de chaque signature
(+1,-1,0) dans le cycle nonadique orienté. La normalisation au sein de chaque fibre reste
une prémisse mathématique précise --- le comptage uniforme de tous les
chemins compatibles --- et ne doit pas être confondue avec le successeur d'une seule
trajectoire. Le passage à la limite inverse de toutes les profondeurs conserve un
énoncé propre, formulé à la fin de la section.

\section{\(5\) niveaux de la notion de chemin}

Soit \(\Gamma_{\rm APP}\) la catégorie des chemins cardinaux du tore APP\@.
Pour chaque profondeur \(m\), soit \(\mathcal H_m\) l'ensemble fini des
histoires enrichies admissibles
\[
 h_m=(v_0,e_1,v_1,\ldots,e_m,v_m),
\]
où chaque flèche \(e_j\) transporte tous les champs déclarés par TPK\@. Si
\(n\geq m\), l'application
\[
 p_{n,m}:\mathcal H_n\longrightarrow\mathcal H_m
\]
élimine la queue postérieure à la profondeur \(m\). L'ombre positionnelle
\[
 \operatorname{sh}_m:\mathcal H_m\longrightarrow\Gamma_{\rm APP}
\]
oublie la fibre de mémoire et conserve le chemin APP sous-jacent.

\begin{definition}[Hiérarchie des extensions]
\label{def:jerarquia-extensiones}
Nous distinguons les objets suivants.
\begin{enumerate}
  \item Un \emph{chemin APP} est un morphisme de
  \(\Gamma_{\rm APP}\).
  \item Une \emph{route observable} est une lecture de ce morphisme qui
  conserve conjointement les feuillets de somme et de produit.
  \item Une \emph{extension TPK} est une élévation de la route observable en
  une histoire \(h_m\in\mathcal H_m\), avec ses données de résidu, de quotient,
  de signature, d'orientation, de frontière, de retenue et de registre.
  \item Une \emph{extension survivante} est une histoire finie qui admet
  des prolongations compatibles à tout horizon.
  \item Une \emph{section globale} est une famille compatible
  \(x=(h_m)_{m\geq0}\) d'extensions survivantes.
\end{enumerate}
En particulier, la route est l'ombre de l'extension ; ce n'est pas un objet extérieur
que l'extension choisirait ensuite.
\end{definition}

Pour \(h\in\mathcal H_m\), son cylindre de prolongations est
\[
 Z_m(h)=
 \{x=(h_n)_n\in\varprojlim_n\mathcal H_n:x_m=h\}.
\]
En particulier, toute famille de \(Z_m(h)\) satisfait en outre
\(p_{n,\ell}(h_n)=h_\ell\) pour \(n\geq\ell\) ; on n'admet pas de collections de
troncatures incompatibles.
Définissons également
\[
 \mathcal S_m^{(N)}=p_{N,m}(\mathcal H_N),
 \qquad
 \mathcal S_m^\infty=\bigcap_{N\geq m}\mathcal S_m^{(N)}.
\]

\begin{theorem}[Critère conditionnel d'extension survivante]
\label{thm:criterio-extension-superviviente}
Fixons \(h\in\mathcal H_m\) et, pour \(N\geq m\), écrivons
\[
 \mathcal P_N(h)=\{g\in\mathcal H_N:p_{N,m}(g)=h\}.
\]
Supposons que les troncatures sont cohérentes et que chaque
\(\mathcal P_N(h)\) est fini. Sont équivalentes :
\begin{enumerate}
  \item \(\mathcal P_N(h)\ne\varnothing\) pour tout \(N\geq m\), c'est-à-dire
  \(h\in\mathcal S_m^\infty\) ;
  \item l'arbre des prolongations finies de \(h\) est à ramification finie
  et possède un niveau non vide à toute profondeur ;
  \item \(Z_m(h)\ne\varnothing\).
\end{enumerate}
Par conséquent, dès lors que les troncatures survivantes finies sont non
vides, l'ensemble
\[
 \Omega_{\rm H}^{\rm surv}
 =\varprojlim_m\mathcal S_m^\infty
\]
est exactement l'espace des sections globales survivantes. Le théorème
n'affirme pas que ces troncatures sont non vides pour toute entrée APP : cette
propriété doit être démontrée au moyen du protocole de survie.
\end{theorem}

\begin{proof}
L'implication de la troisième condition vers les deux premières s'obtient par
troncature. Supposons maintenant la première condition et ordonnons les éléments de
\(\bigsqcup_{N\geq m}\mathcal P_N(h)\) par la relation de préfixe. Chaque niveau
est fini et non vide, et la cohérence des troncatures fait de cette union
un arbre à ramification finie avec des nœuds à toute profondeur ; cela donne la
deuxième condition. Réciproquement, à partir de la deuxième condition, le lemme de König
fournit une branche infinie. Ses troncatures forment une famille compatible
appartenant à \(Z_m(h)\), ce qui est la troisième condition. La même construction
identifie la limite inverse indiquée là où ses troncatures finies sont
non vides.
\end{proof}

Le théorème résout l'éventuelle discordance entre « prolongeable jusqu'à chaque
horizon séparément » et « possède une prolongation globale » : la finitude des
branches est l'hypothèse qui permet de choisir une unique chaîne compatible. Il
n'affirme pas que tous les chemins APP appartiennent à
\(\operatorname{sh}(\Omega_{\rm H}^{\rm surv})\). Cette couverture exige de démontrer
que les fermetures enrichies ne vident pas les fibres correspondantes.

\section{Retour de phase et relation de monodromie}

Le retour du quotient phase--longueur est formalisé séparément dans
\cref{def:odometro-9-12,prop:proyeccion-fase-longitud} :
\[
(b,r)\xrightarrow{\;9\;}(b+1,r),
\qquad
(b,r)\xrightarrow{\;108\;}(b,r).
\]
Ces égalités appartiennent à l'odomètre \(9\times12\). Elles ne constituent pas la
démonstration de \(F_{\rm obs}^{108}=I\), qui utilise l'état observable et apparaît plus
loin.

Soit \(\phi\) la coordonnée de phase à valeurs dans \(\mathbb Z/9\mathbb Z\).
Un retour visible satisfait
\[
 \phi(h_{k+9})=\phi(h_k).
\]
Le registre complet peut cependant satisfaire
\[
 B(h_{k+9})\ne B(h_k).
\]
Cette inégalité n'est pas un défaut de périodicité : c'est la mémoire du tour.
Une action ordinaire de \(C_9\) sur l'état complet imposerait que la
neuvième puissance soit l'identité et effacerait précisément cette information.

Pour typer le retour, soit \(H\) une hexade et
\(P_2(H)\) son ensemble de quinze paires. Soient \(\mathcal B_k\) les états de
registre qui survivent à la phase \(k\), soit
\(\mathcal T_k\subseteq\mathcal B_k\times\mathcal B_{k+1}\) la relation de
transport publiée et supposons donné un encodeur de fibre
\[
 c_k:\mathcal B_k\longrightarrow P_2(H).
\]
Ce n'est qu'après avoir déclaré ces encodeurs que l'on définit la relation observable
\[
 R_k=(c_k\times c_{k+1})(\mathcal T_k)
 \subseteq P_2(H)\times P_2(H).
\]
La monodromie relationnelle d'un tour est
\[
 \mathcal M_k
 =R_{k+8}\circ R_{k+7}\circ\cdots\circ R_k.
\]

\begin{proposition}[Classification du retour nonadique étant donné l'encodeur]
\label{prop:clasificacion-retorno-nonadico}
On distingue trois niveaux logiquement distincts.
\begin{enumerate}
  \item Sans démonstration de fonctionnalité, le retour est la relation
  \(\mathcal M_k\).
  \item Si chaque \(R_k\) est le graphe d'une application
  \(\mu_k:P_2(H)\to P_2(H)\), le transport est le produit biaisé
  \[
   T(k,p)=(k+1,\mu_k(p)),
  \]
  et
  \[
   T^9(k,p)=(k,M_k(p)),
   \qquad
   M_k=\mu_{k+8}\circ\cdots\circ\mu_k.
  \]
  \item Si les \(\mu_k\) sont bijectives, \(M_k\) est une holonomie
  automorphe. Le transport descend en une action ordinaire de
  \(C_9\) sur l'état complet si et seulement si \(M_k=I\) dans toutes les
  fibres.
\end{enumerate}
\end{proposition}

\begin{proof}
La première affirmation est la définition de l'image et de la composition relationnelle une
fois fixés \((c_k)_k\). Sous
fonctionnalité, neuf applications consécutives maintiennent la première
coordonnée fixe et composent leurs actions sur la seconde. La composition de
bijections est une bijection. Enfin, \(T^9=I\) équivaut exactement à
\(M_k=I\) pour chaque phase et chaque point de la fibre.
\end{proof}

\begin{corollary}[Non-retour ponctuel sous transport accumulé]
\label{pf84:E02-29}
Soient \(q:X\to V\), \(T:X\to X\) et \(x\in X\). Si
\[
 q(T^9x)=q(x),\qquad
 T^9x=\mathcal M_x\,x,\qquad
 \mathcal M_x\,x\ne x,
\]
alors la projection revient, mais pas l'état :
\[
 q(T^9x)=q(x),\qquad T^9x\ne x.
\]
L'hypothèse ponctuelle \(\mathcal M_x\,x\ne x\) est indispensable ; l'inégalité
globale \(\mathcal M_x\ne I\) n'exclut pas que \(x\) soit un point fixe.
\end{corollary}

\begin{proof}
La première égalité est une hypothèse. La seconde et la séparation ponctuelle donnent
\(T^9x=\mathcal M_x\,x\ne x\). L'énoncé serait réfuté par un cas qui
satisferait les trois prémisses et aurait \(T^9x=x\) ; il n'autorise pas à remplacer la
troisième prémisse par \(\mathcal M_x\ne I\).
\end{proof}

\begin{theorem}[Suffisance fonctionnelle du tuple enrichi]
\label{pf84:E02-30}
Soit
\[
 E:\mathscr X_{\rm TPK}^{\rm enr}\longrightarrow\mathscr S_E
\]
le tuple qui conserve le bloc visible, les blocs nouveaux et la retenue, l'état
d'élévation, les cylindres, la signature de frontière, l'orientation, l'observable, la phase et le
registre dodécaphasique. Soient
\[
 T:\mathscr X_{\rm TPK}^{\rm enr}\to
 \mathscr X_{\rm TPK}^{\rm enr},
 \qquad
 O:\mathscr X_{\rm TPK}^{\rm enr}\to Y.
\]
Lorsqu'il existe des applications
\[
 \widehat T:\mathscr S_E\to\mathscr X_{\rm TPK}^{\rm enr},
 \qquad
 \widehat O:\mathscr S_E\to Y
\]
tels que
\[
 T=\widehat T\circ E,\qquad O=\widehat O\circ E,
\]
l'égalité \(E(x)=E(y)\) implique \(T(x)=T(y)\), \(O(x)=O(y)\) et, par
déterminisme, les mêmes continuations et sorties ultérieures.
\end{theorem}

\begin{proof}
Par factorisation,
\[
 T(x)=\widehat T(E(x))=\widehat T(E(y))=T(y),
 \qquad
 O(x)=\widehat O(E(x))=\widehat O(E(y))=O(y).
\]
La première égalité permet d'itérer l'argument. C'est une condition
suffisante, non une équivalence non démontrée. Deux états ayant le même tuple et
des transitions ou des sorties différentes, l'accès à une coordonnée absente ou la
consultation d'une valeur cible sont des réfutateurs. La microfibre visible de
cardinal cinq de \cref{pf84:E02-22} contrôle uniquement la non-fidélité du
mot \(w_6\) ; elle ne démontre pas cette factorisation positive.
\end{proof}

La fenêtre finie publiée contient neuf retours de phase avec changement de
l’état holonomique intégral et une sélection locale de trois continuations à une lors de
l’augmentation de deux horizons. Cela prouve le retour avec mémoire et réfute
l’identification du registre à un cycle sans mémoire. L’extraction globale des
neuf codeurs \(c_k\) et, à travers eux, des relations
\(R_k\), à partir de toutes les histoires survivantes est une opération distincte.
La proposition précédente classe le retour une fois donnée l’application
registre--paires ; elle ne dérive pas cette application de l’état holonomique intégral de TPK.

L'épisode concret de la neuvième phase, avec ses trois continuations, le
bloc \(B_{11}\), la frontière \((502,662,638)\) et les inégalités entières
de cylindres, est démontré une seule fois dans
\cref{thm:seleccion-novena-fase}. Ici, seule sa
conséquence typée est utilisée :
\[
\#\operatorname{Surv}_{1}(B_9)=3,\qquad
\#\operatorname{Surv}_{2}(B_9)=1.
\]

Les deux autres mots ne peuvent précéder la même queue ternaire qu'en changeant
l'extension profonde de Hensel. Ils partagent une partie du résidu visible,
mais non la généalogie de l'état.

\section{Le catalogue de 468 émissions comme produit de canaux}

Le catalogue fini part de l'ensemble
\(\Omega_{\rm seed}\) de \(104\,976\) présentations microscopiques. Soit
\(\mathcal U_6\) l'ensemble des émissions sextuples distinctes et soit
\[
 F:\Omega_{\rm seed}\twoheadrightarrow\mathcal U_6.
\]
Le recensement exact est
\[
 |\mathcal U_6|=468,
 \qquad
 \#\{U:|F^{-1}(U)|=144,432,1944\}=432,18,18.
\]
Ces décomptes appartiennent à \(F\). Les comparer à des histoires TPK de
profondeur six exigerait une application explicite du domaine des histoires
vers \(\Omega_{\rm seed}\) ; aucun résultat de cette section ne présuppose cette
application ni n’identifie les présentations du catalogue à des états holonomiques intégraux.
Le catalogue fini fournit, pour chaque émission, une signature additive
\(s\) et une signature multiplicative \(e\). L’application conjointe est une
bijection
\[
 \mathcal U_6\xrightarrow{\sim}
 \mathcal S_+\times\mathcal S_\times,
 \qquad
 |\mathcal S_+|=18,
 \quad
 |\mathcal S_\times|=26.
\]
Ainsi,
\[
 \boxed{468=18\cdot26}
\]
est une factorisation de l'objet et non seulement de son cardinal. La première
coordonnée conserve la manière dont l'émission additionne ; la seconde, la manière dont elle accumule le
produit. Cette séparation est la forme finie explicite de
l'incompatibilité somme--produit : les deux lectures coexistent dans l'état, mais
chaque projection oublie la signature complémentaire.

La première coordonnée admet une paramétrisation supplémentaire qui sera utile pour
la comparer à la frontière nonadique. Si un germe additif a pour coordonnées
\((i,j)\in(\mathbb Z/9\mathbb Z)^2\) et pour orientation
\(\varepsilon=+1\) pour les directions est--sud et \(-1\) pour
nord--ouest, sa signature complète est déterminée bijectivement par
\[
 \bigl(i+j\bmod9,\varepsilon\bigr)
 \in(\mathbb Z/9\mathbb Z)\times\{\pm1\}.
\]
Par conséquent,
\[
 \boxed{\mathcal S_+\simeq
 (\mathbb Z/9\mathbb Z)\times\{\pm1\}.}
\]
La coordonnée produit ne possède pas une uniformité analogue : ses \(26\) signatures
ont des fibres de germes de tailles \(8\), \(24\) et \(108\), avec l'histogramme
\(24\times8+1\times24+1\times108=324\). Cette asymétrie est structurelle. Le
facteur \(18\) contient déjà la phase et l'orientation ; toute descente ultérieure
du facteur produit vers une classe de paires doit conserver une mémoire supplémentaire ou
accepter des fibres non uniformes.

Définissons sur \(\ell^2(\mathcal S_+\times\mathcal S_\times)\) les
espérances conditionnelles
\[
\begin{aligned}
 (E_+f)(s,e)&=\frac1{18}\sum_{s'\in\mathcal S_+}f(s',e),\\
 (E_\times f)(s,e)&=\frac1{26}\sum_{e'\in\mathcal S_\times}f(s,e').
\end{aligned}
\]
Le premier opérateur renouvelle la signature additive et conserve la multiplicative ;
le second réalise l'opération duale. Le troisième canal retient les deux.

\section{La trajectoire unique et le transfert de chemins}

À la différence du retour de l'odomètre défini dans
\cref{def:odometro-9-12,prop:proyeccion-fase-longitud}, la proposition
suivante démontre le retour de l'opérateur observable \(F_{\rm obs}\).

La loi observable publiée met à jour un seul curseur à chaque pas. Aux
phases \(+1\), le curseur additif avance ; aux phases \(-1\), le
multiplicatif ; aux phases \(0\), les deux restent immobiles. Après les
pas \(27,54,81,108\), les deux orientations sont inversées. Notons cette
évolution déterministe \(F_{\rm obs}\).

\begin{proposition}[Retour de la chronologie observable]
\label{prop:retorno-fobs-identidad}
Pour tout choix des deux positions et orientations initiales, les
positions et les orientations sont restaurées après \(54\) pas. En
restaurant également la phase, on obtient
\[
 F_{\rm obs}^{108}=I.
\]
Par conséquent, le retour de Poincaré de \(108\) pas induit l'identité sur
\(\mathcal U_6\) et ne sélectionne pas une mesure stationnaire unique.
\end{proposition}

\begin{proof}
Dans chaque bloc de \(27\) pas, chaque curseur effectue exactement neuf
avancements dans une orientation fixe. Comme l'espace des positions est
\(X_9=(\mathbb Z/9\mathbb Z)^2\), son déplacement est \(9d=0\). À la fin
du bloc, \(d\) est inversé. Deux blocs restaurent la position et l'orientation ;
quatre restaurent aussi l'indice dans \(\Theta_{108}\). Le retour induit
est donc l'identité. Toute probabilité est stationnaire pour
l'identité, de sorte qu'aucune n'est sélectionnée par cette chronologie.
\end{proof}

Le problème apparaît même avant le retour si l'on tente d'oublier la
mémoire de route. La signature produit
\[
 e_0=(5,5,5,5,5,5)
\]
possède vingt-quatre représentants. Après un avancement dans leur propre orientation,
ces représentants se séparent, huit par huit, entre
\[
 (5,5,5,7,1,7),\qquad
 (5,5,5,7,7,1),\qquad
 (5,5,5,1,7,7).
\]
Ainsi, l'avancement produit ne descend pas en une application des vingt-six
signatures : le représentant qui avait été oublié redevient pertinent. Ce
témoin explique pourquoi une route individuelle ne peut pas devenir, par simple
projection, le renouvellement d'une fibre complète.

La lecture APP est différente. Au lieu de choisir un représentant, elle conserve toutes
les continuations compatibles. Pour la formaliser dans le catalogue fini,
considérons les sous-algèbres de fonctions qui ne dépendent que de l'une des deux
signatures. Par rapport au comptage uniforme, \(E_+\) et \(E_\times\) sont les
espérances conditionnelles qui oublient exactement la coordonnée active. Ce sont les
seuls projecteurs de Markov qui satisfont simultanément :
\begin{enumerate}
  \item la préservation de la coordonnée inactive ;
  \item l'oubli complet de la coordonnée active ;
  \item la préservation de l'état de comptage uniforme.
\end{enumerate}
En effet, une ligne qui a oublié la coordonnée de départ doit être
constante sur les \(18\), respectivement \(26\), destinations compatibles ; la
normalisation impose les valeurs \(1/18\) et \(1/26\). C'est la canonicité
relative du renouvellement : on ne postule pas de préférence entre des branches qu'APP
déclare également possibles.

Soit \(C_9=\mathbb Z/9\mathbb Z\), avec le mot de phase
\[
 \tau=(+1,-1,0,+1,-1,0,+1,-1,0).
\]
Écrivons
\[
 P_{+1}=E_+,
 \qquad P_{-1}=E_\times,
 \qquad P_0=I.
\]

\begin{definition}[Transfert TPK biaisé par la phase]
\label{def:transferencia-tpk-fase}
Sur \(\mathcal U_6\times C_9\), nous définissons
\[
 \mathcal K_{\rm ph}\bigl((u,r),(v,r+1)\bigr)
 =P_{\tau(r)}(u,v),
\]
et nous annulons les autres entrées. Cette définition est relative au lecteur de
tous les chemins compatibles : une phase active renouvelle uniformément la fibre
que cette phase lit et la phase neutre retient l'émission.
\end{definition}

\begin{theorem}[Retour nonadique uniforme]
\label{thm:retorno-nonadico-uniforme}
Le noyau \(\mathcal K_{\rm ph}\) est doublement stochastique et irréductible,
a une période exactement égale à neuf et possède un unique état stationnaire,
\[
 \pi_{\rm ph}(u,r)=\frac1{9\cdot468}.
\]
Dans chaque fibre de phase, son retour de neuf pas est
\[
 \mathcal K_{\rm ph}^{\,9}\big|_r=E_+E_\times,
 \qquad
 (E_+E_\times)(u,v)=\frac1{468}.
\]
Si la phase initiale est distribuée uniformément, l'ombre d'un pas sur
les émissions est
\[
 \overline K_{\rm ph}=\frac13(I+E_++E_\times).
\]
\end{theorem}

\begin{proof}
Chaque \(P_{\tau(r)}\) est doublement stochastique et le déplacement
\(r\mapsto r+1\) est une permutation ; leur produit biaisé conserve les deux sommes.
Les trois opérateurs sont idempotents et commutent. Comme chacun apparaît trois
fois dans le mot nonadique,
\[
 \prod_{r\in C_9}P_{\tau(r)}
 =E_+^3E_\times^3I^3=E_+E_\times.
\]
La première espérance oublie la signature additive et la seconde la multiplicative,
de sorte que leur composition attribue à chacune des \(18\cdot26=468\)
destinations la probabilité \(1/468\). Ainsi, en neuf pas, deux
états quelconques d'une même phase communiquent ; les pas restants font communiquer les
phases, et le noyau est irréductible. Tout retour doit avoir une longueur multiple
de neuf, tandis que la diagonale de \(\mathcal K_{\rm ph}^9\) est positive ;
la période est neuf. La double stochasticité et l'irréductibilité donnent
l'unique état uniforme. Enfin, moyenner les neuf phases compte trois
copies de chaque canal et produit la dernière formule.
\end{proof}

Oublier la phase n'est pas un quotient de Markov fort : à partir d'une même émission,
deux phases distinctes appliquent \(I\), \(E_+\) ou \(E_\times\), qui ont des lignes
distinctes. L'opérateur \(\overline K_{\rm ph}\) est une ombre moyennée sur la
phase stationnaire, non le successeur chronologique d'une émission sans phase. Cette
distinction conserve à la fois la dynamique et le sens de la mesure.

\begin{definition}[Noyau moyenné induit par le transfert]
\label{def:nucleo-tritico-refresco}
Nous appelons noyau moyenné du catalogue l'ombre de phase
\[
 \boxed{
 K_{\rm cat}=\frac13(I+E_++E_\times).
 }
\]
Les poids sont induits par le calendrier TPK, qui contient trois occurrences
de chaque signe. Les opérateurs de renouvellement sont induits relativement à la
normalisation uniforme de toutes les continuations compatibles de la fibre
active. Aucune des deux affirmations n'identifie \(K_{\rm cat}\) à
\(F_{\rm obs}\).
\end{definition}

\begin{proposition}[Famille convexe de noyaux stochastiques à mesure uniforme commune]
\label{prop:familia-refrescos-abc}
Pour \(a,b,c\geq0\), \(a+b+c=1\), la famille
\[
 K_{a,b,c}=aI+bE_++cE_\times
\]
est formée de noyaux symétriques et doublement stochastiques. Si \(b,c>0\),
ils sont irréductibles et, ayant une diagonale positive, apériodiques. Leur spectre est
\[
 \operatorname{Spec}(K_{a,b,c})=
 \{1^{[1]},(a+c)^{[17]},(a+b)^{[25]},a^{[425]}\}.
\]
En particulier, il existe un continuum de dynamiques ayant la même factorisation et la
même mesure uniforme. Le choix
\(K_{\rm cat}=K_{1/3,1/3,1/3}\) n'est pas caractérisé par le seul recensement. Il est
sélectionné conjointement par le mot de phase TPK et le transfert
normalisé de \cref{def:transferencia-tpk-fase}.
\end{proposition}

\begin{proof}
Les trois termes sont des noyaux symétriques et stochastiques ; par convexité, il en va
de même pour \(K_{a,b,c}\), et la symétrie transforme la stochasticité par lignes
en stochasticité par colonnes. Si \(b,c>0\), deux pas permettent de renouveler
successivement les deux coordonnées et de faire communiquer toute paire d'états. De plus,
\[
 K_{a,b,c}((s,e),(s,e))=a+\frac b{18}+\frac c{26}>0,
\]
de sorte que le noyau est apériodique. Dans la décomposition orthogonale
\[
 \mathbf1\oplus\ell_0^2(\mathcal S_+)
 \oplus\ell_0^2(\mathcal S_\times)
 \oplus\bigl(\ell_0^2(\mathcal S_+)\otimes
              \ell_0^2(\mathcal S_\times)\bigr),
\]
les paires \((E_+,E_\times)\) agissent, respectivement, comme
\((1,1),(0,1),(1,0),(0,0)\). On obtient ainsi les quatre valeurs propres et les
multiplicités \(1,17,25,425\). Pour \(b,c>0\), l'irréductibilité rend unique
la mesure stationnaire uniforme, ce qui montre la non-unicité même sous
cette exigence.
\end{proof}

\begin{theorem}[Spectre exact du renouvellement équipondéré]
\label{thm:dinamica-cociente-468}
Le noyau \(K_{\rm cat}\) est symétrique, doublement stochastique,
irréductible et apériodique. Son spectre, avec multiplicités, est
\[
 \boxed{
 \operatorname{Spec}(K_{\rm cat})
 =\left\{1^{[1]},
 \left(\frac23\right)^{[42]},
 \left(\frac13\right)^{[425]}\right\}.
 }
\]
Par conséquent, son écart spectral est (1/3) et son unique état
stationnaire est
\[
 \boxed{
 \tau_{468}(f)=\frac1{468}\sum_{U\in\mathcal U_6}f(U).
 }
\]
\end{theorem}

\begin{proof}
Les opérateurs \(E_+\) et \(E_\times\) sont des projecteurs orthogonaux de
moyennage et commutent. Les lignes et les colonnes de chacun ont pour somme un ; il en va de même
pour leur moyenne. De \((s,e)\), on atteint \((s',e')\) en deux
renouvellements, et la contribution de \(I\) fournit une boucle. Cela
démontre l'irréductibilité et l'apériodicité.

L'espace se décompose orthogonalement en quatre secteurs :
\[
 \mathbf1,
 \qquad
 \ell_0^2(\mathcal S_+),
 \qquad
 \ell_0^2(\mathcal S_\times),
 \qquad
 \ell_0^2(\mathcal S_+)\otimes
 \ell_0^2(\mathcal S_\times).
\]
Leurs dimensions sont \(1,17,25,17\cdot25\). Dans le secteur constant, les
trois termes agissent comme l'identité ; dans chaque secteur d'une seule coordonnée,
ils agissent comme (I+I+0) ; dans le secteur d'interaction, comme (I+0+0).
Diviser par trois donne les valeurs propres et les multiplicités annoncées. Un noyau
fini irréductible et apériodique possède un unique état stationnaire ; comme il est
doublement stochastique, celui-ci est uniforme.
\end{proof}

Le théorème est exact relativement à la factorisation
\(\mathcal U_6\simeq\mathcal S_+\times\mathcal S_\times\), au calendrier
TPK et au comptage uniforme des continuations compatibles de chaque fibre. La
famille \(K_{a,b,c}\) démontre que la dynamique ne se déduit pas du seul
recensement. \cref{prop:retorno-fobs-identidad} distingue le transfert de
tous les chemins de la chronologie de l'un d'entre eux.

\section{Élévation compensée construite à partir du quotient}

Écrivons \(m(U)=|F^{-1}(U)|\). À partir du noyau déjà choisi dans le quotient,
nous définissons le noyau de présentations par
\[
 \boxed{
 K_{\rm seed}(x,y)
 =\frac{K_{\rm cat}(F(x),F(y))}{m(F(y))}
 }
\]
Ce noyau est stochastique. En sommant sur une fibre de destination, on obtient
\[
 \sum_{F(y)=V}K_{\rm seed}(x,y)
 =K_{\rm cat}(F(x),V),
\]
indépendamment du représentant \(x\). Cette identité montre que le
relèvement construit se projette exactement sur \(K_{\rm cat}\). Elle ne démontre pas
qu’une dynamique microscopique préexistante de l’état holonomique intégral de TPK descende
au catalogue.

\begin{theorem}[Élévation réversible construite à partir du quotient]
\label{thm:elevacion-reversible-cociente}
La probabilité
\[
 \widetilde\tau(x)
 =\frac1{468\,m(F(x))}
\]
est stationnaire et réversible pour \(K_{\rm seed}\), et
\(F_\#\widetilde\tau=\tau_{468}\). Si
\[
 J_F\delta_U
 =m(U)^{-1/2}\sum_{F(x)=U}\delta_x,
 \qquad
 D=\operatorname{diag}(m(U)),
\]
alors
\[
 K_{\rm seed}J_F
 =J_F\bigl(D^{1/2}K_{\rm cat}D^{-1/2}\bigr).
\]
En particulier, l'entrelacement hilbertien sans conjugaison par \(D\)
a lieu si et seulement si \(K_{\rm cat}\) commute avec les multiplicités des
fibres.
\end{theorem}

\begin{proof}
La normalisation découle du fait que chacune des 468 fibres reçoit une masse
totale (1/468). En utilisant la symétrie de \(K_{\rm cat}\), pour
\(F(x)=U\) et \(F(y)=V\), on a
\[
 \widetilde\tau(x)K_{\rm seed}(x,y)
 =\frac{K_{\rm cat}(U,V)}{468m(U)m(V)}
 =\widetilde\tau(y)K_{\rm seed}(y,x).
\]
Cela démontre l'équilibre détaillé, la stationnarité et l'identité de la
mesure image.
Appliquer les deux membres de la dernière identité à un vecteur de base
\(\delta_V\) donne, sur la fibre \(U\), le coefficient
\(K_{\rm cat}(U,V)/\sqrt{m(V)}\) ; c'est exactement le coefficient de
l'opérateur conjugué présenté.
\end{proof}

\begin{theorem}[Élévation de phase fidèle à l'immobilité]
\label{thm:elevacion-fase-semillas}
Soit \(X=\Omega_{\rm seed}\), \(m(u)=|F^{-1}(u)|\), et définissons
\[
\begin{aligned}
 L_0(x,y)&=\mathbf1_{\{x=y\}},\\
 L_+(x,y)&=\frac{E_+(F(x),F(y))}{m(F(y))},\\
 L_\times(x,y)&=\frac{E_\times(F(x),F(y))}{m(F(y))}.
\end{aligned}
\]
Sur \(X\times C_9\), soit
\[
 \widehat{\mathcal K}_{\rm ph}\bigl((x,r),(y,r+1)\bigr)
 =L_{\tau(r)}(x,y).
\]
Alors \((F,\mathrm{id}_{C_9})\) satisfait la propriété
d'\emph{agrégabilité forte} (\emph{strong lumpability}) de
\(\widehat{\mathcal K}_{\rm ph}\) vers \(\mathcal K_{\rm ph}\). Le noyau
microscopique est irréductible, a pour période neuf et son unique état
stationnaire est
\[
 \widetilde\pi_{\rm ph}(x,r)
 =\frac1{9\cdot468\,m(F(x))}.
\]
Son retour nonadique satisfait
\[
 \widehat{\mathcal K}_{\rm ph}^{\,9}
 \bigl((x,r),(y,r)\bigr)
 =\frac1{468\,m(F(y))}.
\]
\end{theorem}

\begin{proof}
La somme de \(L_+\), respectivement \(L_\times\), sur une fibre de destination
est \(E_+\), respectivement \(E_\times\), et \(L_0\) se projette sur l'identité.
Cela démontre l'agrégabilité à chaque phase. Les facteurs \(m(F(y))\) s'annulent
lors de la composition : on sélectionne d'abord uniformément la signature active, puis un
représentant de l'émission obtenue. Ainsi, \(L_+\) et \(L_\times\) sont des
projecteurs commutants et
\[
 L_+L_\times(x,y)=\frac1{468\,m(F(y))}.
\]
La positivité de cette entrée fait communiquer tous les représentants en neuf
pas. La phase impose que les longueurs de retour soient des multiples de neuf,
et l'entrée diagonale du retour est positive. La formule stationnaire se
vérifie en sommant d'abord les \(m(u)\) représentants de chaque émission ; chacune
des \(9\cdot468\) classes phase--émission reçoit la même masse.
\end{proof}

Cette élévation améliore la lecture du noyau moyenné
\(K_{\rm seed}\) : dans la phase neutre, elle conserve le représentant microscopique,
au lieu de remélanger silencieusement sa fibre. En moyennant et en oubliant la phase,
les deux constructions ont la même ombre \(K_{\rm cat}\), mais pas la même
dynamique cachée. La différence est nécessaire pour respecter le fait que le zéro de TRIT
est une mémoire d'immobilité.

\section{Congruences minimales stables par avancement, demi-tour et 4e de tour : \(28\) et \(162\) classes de mémoire}

Le témoin \(e_0=(5,5,5,5,5,5)\) montre que la signature produit est trop
grossière pour transporter une route. Il est possible de calculer exactement la mémoire
minimale manquante. Soit
\[
 \mathcal X_\times=X_9\times\mathscr D_{\rm card}
\]
l'espace des \(324\) germes produit. Notons \(P\) l'avancement d'une
case dans l'orientation stockée, \(H\) le demi-tour de cette
orientation et \(Q\) le quart de tour
\[
 N\longmapsto E\longmapsto S\longmapsto O\longmapsto N.
\]

\begin{theorem}[Congruences d'avancement et de rotation]
\label{thm:congruencias-avance-giro}
La congruence minimale qui raffine les vingt-six signatures produit et est stable
par \(P\) et \(H\) possède \(28\) classes : vingt-sept sont de taille huit et une
est de taille \(108\). L'unique scission nouvelle est
\[
 (5,5,5,5,5,5)\longrightarrow
 \mathcal L_0\sqcup\mathcal L_1\sqcup\mathcal L_2,
 \qquad |\mathcal L_j|=8.
\]
Ses orbites sous \(P,H\) ont pour tailles
\[
 9,\quad9,\quad9,\quad1.
\]
En combinant ce raffinement avec les dix-huit signatures additives, on obtient
\[
 18\cdot28=504=468+36.
\]

La congruence minimale stable par \(P\) et \(Q\) possède \(162\) classes de deux
germes. Chaque classe est une orbite de l'involution
\[
 J(i,j,d)=\bigl(7-i,\,7-j,\,\overline d\bigr)\pmod9.
\]
Les opérateurs induits par \(P,Q\) agissent transitivement sur les
\(162\) classes.
\end{theorem}

\begin{proof}
On part de la partition par les vingt-six signatures et l'on raffine
itérativement une classe dès que deux de ses éléments ont des images dans des
classes distinctes par l'un des générateurs. Le processus se termine car
\(\mathcal X_\times\) est fini. Le recensement exhaustif produit, pour
\((P,H)\), l'histogramme \(27\times8+1\times108=324\), les quatre orbites
indiquées et trois classes au-dessus de la fibre exceptionnelle de taille vingt-quatre.
Pour \((P,Q)\), il produit \(162\times2=324\). La vérification directe démontre
que \(J\) préserve chaque signature, commute avec les deux générateurs et apparie
exactement les deux éléments de chaque classe ; le graphe de classes engendré par
\(P,Q\) est connexe.
\end{proof}

L'excès \(36\) et l'objet \(R_{36}\) sont de types distincts. La frontière
de profondeur trente-six est
\[
 R_{36}=
 \begin{pmatrix}
 222220\\021222\\102011
 \end{pmatrix},
\]
c'est-à-dire trois blocs ordonnés de six trits, un par canal. Ce n'est pas un
ensemble de trente-six états. De son côté, le terme
\(504-468=36\) compte une mémoire supplémentaire de feuillet : le demi-tour qui rend
fonctionnel le transport produit ajoute exactement deux feuillets nouveaux à une
seule signature et, en les répétant sur les dix-huit signatures additives, produit
\(18(3-1)=36\). La coïncidence de l'entier (36) exprime deux graduations
du transport, profondeur de frontière et multiplicité de feuillet, mais
n'autorise pas à identifier leurs objets.

Le quart de tour \(Q\) n'était pas non plus caché dans \(F_{\rm obs}\). La
chronologie publiée contient des avancements et des demi-tours programmés ; incorporer
\(Q\) est une extension explicite de l'alphabet des routes. Une fois incorporé,
une marche symétrique paresseuse dans \(P^{\pm1},Q^{\pm1}\) est irréductible et
apériodique sur les \(162\) classes, et sa mesure uniforme est unique. C'est la
forme mathématique exacte de « compter les rotations » : la rotation est une coordonnée de
transport et non une correction décimale ultérieure.

\section{Cocycles de mode et d'orientation}
\label{sec:cociclos-modo-orientacion}

La même règle permet de définir les compteurs opératoires sans introduire la
mécanique dimensionnelle. Ajoutons à l'état le dernier mode actif
\(m\in\{\Sigma,\Pi\}\). Une signature \(+1\) fixe \(m=\Sigma\), une signature
\(-1\) fixe \(m=\Pi\), et une signature \(0\) conserve le mode précédent. Pour une
flèche \(a\), définissons le cocycle signé
\[
 \omega_{\rm sw}(a)=
 \begin{cases}
  +1,&\Pi\longrightarrow\Sigma,\\
  -1,&\Sigma\longrightarrow\Pi,\\
  0,&\text{le mode est conservé}.
 \end{cases}
\]
Sa variation positive est \(s(a)=|\omega_{\rm sw}(a)|\). Définissons en outre
\(g(a)=1\) lorsque la paire ordonnée des orientations cardinales change et
\(g(a)=0\) lorsqu'elle est conservée. En maintenant aux extrémités le mode et
l'orientation terminale, les sommes
\[
 \Omega_{\rm sw}(\gamma)=\sum_{a\subset\gamma}\omega_{\rm sw}(a),
 \qquad
 S(\gamma)=\sum_{a\subset\gamma}s(a),
 \qquad
 G(\gamma)=\sum_{a\subset\gamma}g(a)
\]
sont additives par concaténation.

Dans un tour nonadique fermé, la suite des modes est
\[
 (\Sigma,\Pi,\Pi,\Sigma,\Pi,\Pi,\Sigma,\Pi,\Pi).
\]
Par conséquent,
\[
 \Omega_{\rm sw}(C_9)=0,\qquad S(C_9)=6:
\]
il y a trois changements \(\Pi\to\Sigma\), trois changements
\(\Sigma\to\Pi\) et trois maintiens. Dans le cycle de \(108\) pas,
\[
 \#(+1)=\#(-1)=\#(0)=36,\qquad
 \Omega_{\rm sw}=0,\qquad S=72,\qquad G=4.
\]
Le dernier nombre compte les quatre inversions simultanées programmées dans
\(27,54,81,108\). Ces cocycles proviennent directement du calendrier TPK\@.

\subsection{Quotient des modes et correspondance aréale--volumétrique}
\label{sec:cociente-modos-fav}

Le signe du calendrier, l'étiquette de mode et le feuillet géométrique sont trois
objets typés. En particulier, on n'identifie pas le signe de phase à
l'observable angulaire sur les chiffres. Pour la correspondance géométrique, nous fixons le
lecteur HMT typé
\[
 \mathfrak g_{\rm av}(\Sigma)=\mathscr H_{\rm ar},
 \qquad
 \mathfrak g_{\rm av}(\Pi)=\mathscr H_{\rm vol}.
\]
La première attribution exprime la lecture additive ou de bord ; la seconde, la
lecture multiplicative ou d'intérieur. L'application \(\mathfrak g_{\rm av}\) fait
partie du typage des feuillets. Le résultat suivant démontre qu'une fois
ce typage fixé, la dynamique TPK impose l'incidence entre eux.

\begin{theorem}[Descente nécessaire de l'incidence des feuillets]
\label{thm:descenso-necesario-fav}
Soit \(m_r\) le mode obtenu après application de la phase \(r\) du bloc
primitif \((+1,-1,0)\), avec la règle
\[
 +1:\ m\mapsto\Sigma,\qquad
 -1:\ m\mapsto\Pi,\qquad
 0:\ m\mapsto m.
\]
L'orbite cyclique des modes est \((\Sigma,\Pi,\Pi)\). Si
\[
 N_3(i,j)
 :=
 \#\{r\in\mathbb Z/3\mathbb Z:(m_r,m_{r+1})=(i,j)\},
\]
alors, dans l'ordre \((\Sigma,\Pi)\),
\[
 \boxed{
 N_3=
 \begin{pmatrix}0&1\\1&1\end{pmatrix}.}
\]
Par répétition du calendrier,
\[
 N_9=3N_3,\qquad N_{27}=9N_3,\qquad N_{108}=36N_3.
\]
Après application de \(\mathfrak g_{\rm av}\), la matrice primitive
d'incidence est, dans l'ordre
\((\mathscr H_{\rm ar},\mathscr H_{\rm vol})\),
\[
 \boxed{
 F_{\rm av}=
 \begin{pmatrix}0&1\\1&1\end{pmatrix}.}
\]
Ainsi, les deux transitions croisées, l'unique autorétention volumétrique,
l'absence d'autorétention aréale et les multiplicités primitives unitaires
sont des conséquences du quotient des modes de TPK ; ce ne sont pas quatre postulats
ajoutés.
\end{theorem}

\begin{proof}
Les trois couples consécutifs, avec fermeture cyclique, sont
\[
 \Sigma\longrightarrow\Pi,\qquad
 \Pi\longrightarrow\Pi,\qquad
 \Pi\longrightarrow\Sigma.
\]
Chacun apparaît une fois et le couple
\(\Sigma\to\Sigma\) n'apparaît pas. Cela détermine les quatre entrées de
\(N_3\). Les calendriers de longueurs \(9,27,108\) contiennent,
respectivement, \(3,9,36\) copies du bloc primitif, de sorte que leurs
matrices d'incidence sont les multiples indiqués. Le plus grand commun diviseur
des entrées non nulles de chaque matrice longue élimine uniquement cette
répétition globale et restitue \(N_3\). Enfin,
\(\mathfrak g_{\rm av}\) transporte \(\Sigma,\Pi\) vers les feuillets aréal et
volumétrique sans modifier l'incidence.
\end{proof}

La réalisation hyperbolique réelle qui utilise cette matrice et prouve son
invariance lorentzienne est formulée dans
\cref{sec:tres-generadores-concretos-trit} ; on y utilise \(F_{\rm av}\)
comme antécédent, sans répéter sa dérivation.

Ce théorème doit être distingué d'une affirmation de Markov plus forte.
L'oubli de la phase n'est pas un quotient de Markov fort de
\(\mathcal K_{\rm ph}\) : le mode visible ne détermine pas lequel des trois
opérateurs \(I,E_+,E_\times\) agit à l'itération élémentaire suivante. Ce qui descend
sans hypothèse supplémentaire est le multigraphe des arêtes d'un bloc TPK\@.

Il existe cependant une complétion à mémoire un déterminée de manière
unique par cette descente. Soit \(X_{\rm av}\) le plus petit sous-décalage de type fini
sur
\(\{\mathscr H_{\rm ar},\mathscr H_{\rm vol}\}\) qui contient tous les
mots obtenus en oubliant la phase et qui ne conserve que des couples consécutifs.
Tout sous-décalage à un pas ayant cette propriété doit admettre exactement les trois
couples observés ; le plus petit d'entre eux exclut l'unique couple non observé. Par
conséquent, sa matrice d'adjacence est \(F_{\rm av}\).

\begin{corollary}[Mesure de Parry de l'enveloppe à mémoire un]
\label{cor:parry-envolvente-fav}
Le sous-décalage \(X_{\rm av}\) est primitif, a pour entropie topologique
\(\log\varphi\) et possède une unique mesure d'entropie maximale. Ses poids
stationnaires de sommet satisfont
\[
 \boxed{
 \frac{\mu_{\rm ar}}{\mu_{\rm vol}}=\varphi^{-2}.}
\]
\end{corollary}

\begin{proof}
La racine de Perron de \(F_{\rm av}\) est \(\varphi\), et un vecteur de Perron
positif est \((1,\varphi)\). Comme la matrice est symétrique, les poids de sommet
de Parry sont proportionnels au produit des vecteurs gauche et droit,
c'est-à-dire à \((1,\varphi^2)\).
\end{proof}

La mesure de Parry appartient à l'enveloppe des chemins à mémoire un. Elle n'est pas
l'image de \(\tau_{468}\), ni la fréquence temporelle de l'orbite périodique
des modes. Cette dernière est
\[
 \nu_{\rm cyc}(\mathscr H_{\rm ar})=\frac13,\qquad
 \nu_{\rm cyc}(\mathscr H_{\rm vol})=\frac23,
\]
tandis que \(\tau_{468}\) est uniforme sur les émissions et que la mesure de Parry
pondère des histoires de renouvellement de longueur arbitraire. Les trois mesures
répondent à des questions distinctes. Le rapport irrationnel
\(\varphi^{-2}\) est ainsi déduit canoniquement de l'enveloppe des chemins,
sans l'attribuer à un quotient de cardinaux finis.

\subsubsection{Du comptage chronologique au retour marqué par \texorpdfstring{\(\pi\)}{π}}
\label{sec:retorno-marcado-pi-phi2}

Les deux formules
\[
 \left(\frac13,\frac23\right)
 \qquad\text{et}\qquad
 \frac{\mu_{\rm ar}}{\mu_{\rm vol}}=\varphi^{-2}
\]
ne sont pas des approximations l'une de l'autre. La première compte les trois instants du
calendrier primitif \((\Sigma,\Pi,\Pi)\) ; la seconde pondère toutes les
histoires admissibles de l'enveloppe à mémoire un. C'est pourquoi la fréquence
chronologique est rationnelle et le rapport de Parry est irrationnel. Dans ce qui suit,
nous écrivons
\[
 \frac{\pi}{\varphi^2}
 =\pi\varphi^{-2};
\]
il ne s'agit pas de diviser par \(\varphi^{-2}\).

L'apparition de \(\varphi^{-2}\) se voit directement dans la dynamique
linéaire. Si \(F_n\) désigne le \(n\)-ième nombre de Fibonacci,
\[
 F_{\rm av}^{\,n}
 =
 \begin{pmatrix}
  F_{n-1}&F_n\\
  F_n&F_{n+1}
 \end{pmatrix}
 \qquad(n\geq1).
\]
Les valeurs propres de \(F_{\rm av}\) sont
\(\varphi\) et \(-\varphi^{-1}\). La branche contractée inverse l'orientation en
une itération. Lors du passage à la mémoire quadratique, c'est-à-dire à
\(\operatorname{Sym}^2(\mathbb R^2)\), cette inversion est enregistrée avant la
disparition du signe. Dans la base \((x^2,2xy,y^2)\), en représentant par
lignes les images des vecteurs de base, l'opérateur induit est
\[
 \operatorname{Sym}^2(F_{\rm av})=
 \begin{pmatrix}
  0&0&1\\
  0&1&2\\
  1&1&1
 \end{pmatrix},
 \qquad
 \chi(t)=(t+1)(t^2-3t+1).
\]
Ses trois valeurs propres sont
\[
 \varphi^2,\qquad -1,\qquad \varphi^{-2}.
\]
Ainsi, \(\varphi^{-2}\) est le mode stable positif de la mémoire du second
ordre. Cette étape explique pourquoi le nombre d'or apparaît ici comme loi de
retour et non comme moyenne statistique ajoutée.

La clôture \(\pi\) appartient à un autre niveau du même état HMT : c'est la sortie
coinductive du lecteur régional déjà construit, non une nouvelle valeur propre de
\(F_{\rm av}\). Soient
\[
 P_{\rm ar}=
 \begin{pmatrix}1&0\\0&0\end{pmatrix},
 \qquad
 P_{\rm vol}=
 \begin{pmatrix}0&0\\0&1\end{pmatrix}.
\]
Parmi les opérateurs positifs qui préservent les deux feuillets, il en existe un et un
seul dont la normalisation volumétrique vaut un et dont la marque aréale est la clôture
engendrée \(\pi\) :
\[
 D_\pi={\pi}P_{\rm ar}+P_{\rm vol}.
\]
L'unicité est relative à ces trois clauses explicites : positivité,
diagonalité par rapport aux feuillets et les deux normalisations. Avec
\(e_{\rm ar}=(1,0)^T\) et \(e_{\rm vol}=(0,1)^T\), le retour marqué de
longueur \(n\) est
\[
 \Theta_n
 :=
 \frac{\langle e_{\rm ar},
 D_\pi F_{\rm av}^{\,n}e_{\rm ar}\rangle}
 {\langle e_{\rm vol},
 F_{\rm av}^{\,n}e_{\rm vol}\rangle}
 =
 \pi\frac{F_{n-1}}{F_{n+1}}.
\]
Par conséquent,
\[
 \boxed{\displaystyle
 \lim_{n\to\infty}\Theta_n
 =\frac{\pi}{\varphi^2}.}
\]
La marque \(\pi\) est appliquée une seule fois, au retour de frontière ;
elle n'est pas réinjectée à chaque transition. La combinatoire finie produit le mode
stable \(\varphi^{-2}\), le lecteur coinductif produit la clôture
\(\pi\), et \(D_\pi\) compose les deux données sans inverser leur ordre
causal.

La longueur native \(108\) offre un témoin exact de cette convergence :
\[
 (F_{\rm av}^{108})_{\rm ar,ar}
 =F_{107}=10284720757613717413913,
\]
\[
 (F_{\rm av}^{108})_{\rm vol,vol}
 =F_{109}=26925748508234281076009.
\]
En conséquence,
\[
 \left|
 \frac{F_{107}}{F_{109}}-\varphi^{-2}
 \right|
 =
 6.1684979916614407936\ldots\times10^{-46},
\]
et, après l'unique marque de clôture,
\[
 \left|
 \pi\frac{F_{107}}{F_{109}}
 -\frac{\pi}{\varphi^2}
 \right|
 =
 1.9378907974286976068\ldots\times10^{-45}.
\]
Ces erreurs mesurent l'approximation finie du retour ; ce ne sont pas des ajustements de
\(\pi\) ni de \(\varphi\).

\paragraph{Lecteur cylindrique du rapport.}
Soient \(I_n^\pi=[\underline\pi_n,\overline\pi_n]\) et
\(I_n^\varphi=[\underline\varphi_n,\overline\varphi_n]\) les cylindres
positifs et emboîtés produits par les deux lecteurs HMT\@. L'opération sur les
intervalles
\[
 \mathcal Q(I,J)
 :=
 \left[
 \frac{\inf I}{(\sup J)^2},
 \frac{\sup I}{(\inf J)^2}
 \right]
\]
est le plus petit intervalle qui contient \(x/y^2\) pour tout
\((x,y)\in I\times J\). Elle préserve l'emboîtement : si les cylindres d'entrée sont
raffinés, \(\mathcal Q(I_n^\pi,I_n^\varphi)\) est également raffiné. Comme leurs
diamètres tendent vers zéro et que la limite de \(I_n^\varphi\) est positive,
\[
 \bigcap_n\mathcal Q(I_n^\pi,I_n^\varphi)
 =
 \left\{\frac{\pi}{\varphi^2}\right\}.
\]
Avec les douze blocs en base mille déjà publiés, l'intervalle quotient impose
trente-cinq décimales :
\[
 \frac{\pi}{\varphi^2}
 =
 1.19998161486432666111577775331680692\ldots.
\]
Cette construction conserve la provenance de chaque chiffre : le numérateur
provient du cylindre de clôture et le dénominateur de l'autoéchelle.

\paragraph{Séparation algébrique entre lecture aréale, lecture volumétrique et leurs codomaines respectifs}
Le rapport complet ne peut provenir de la seule algèbre rationnelle de
\(F_{\rm av}\). En effet, ses fonctions rationnelles spectrales appartiennent à
\(\mathbb Q(\sqrt5)\) ; on y trouve \(\varphi^{-2}\), mais pas le nombre
transcendant \(\pi\). Par conséquent, toute dérivation qui prétendrait
obtenir \(\pi/\varphi^2\) uniquement à partir de la matrice \(2\times2\) serait
mal typée. La donnée transcendante doit provenir du lecteur coinductif de
clôture, ou d'un cocycle de frontière équivalent démontré
indépendamment. Ce contrôle négatif fixe la répartition des tâches entre
la correspondance finie et le lecteur de nombre.

\subsubsection{Involution d'échange du quotient}
\label{sec:espejo-pi-phi-electron-hmt}

La correspondance admet deux orientations exactes d'une même loi. Si
\[
 \mathscr R(x,y)=\frac{x}{y^2},
 \qquad
 \mathsf J(x,y)=(y,x),
\]
alors \(\mathsf J^2=I\) et
\[
 \mathscr R(\pi,\varphi)=\frac{\pi}{\varphi^2},
 \qquad
 (\mathscr R\circ\mathsf J)(\pi,\varphi)
 =\frac{\varphi}{\pi^2}.
\]
On vérifie en outre les identités
\[
 \mathscr R(\pi,\varphi)\mathscr R(\varphi,\pi)
 =\frac1{\pi\varphi},
 \qquad
 \frac{\mathscr R(\pi,\varphi)}
 {\mathscr R(\varphi,\pi)}
 =\left(\frac{\pi}{\varphi}\right)^3.
\]
En écrivant
\(\mathcal B_{\rm av}^{(\pi)}=\mathscr R(\pi,\varphi)\), on obtient
également la reparamétrisation exacte
\[
 \frac{\varphi}{\pi^2}
 =
 \frac1{\varphi^3
 \bigl(\mathcal B_{\rm av}^{(\pi)}\bigr)^2},
 \qquad
 e^{\varphi/\pi^2}
 =
 \exp\!\left[
 \frac1{\varphi^3
 \bigl(\mathcal B_{\rm av}^{(\pi)}\bigr)^2}
 \right].
\]
Ainsi, le changement d'orientation entre clôture et croissance transporte le
rapport de frontière vers sa lecture échangée :
\[
 \frac{\pi}{\varphi^2}
 \xleftrightarrow{\ \mathsf J\ }
 \frac{\varphi}{\pi^2}.
\]
La première orientation lit la clôture aréale marquée sur la mémoire
volumétrique quadratique ; la seconde lit la croissance intérieure sur la
clôture quadratique. Cette correspondance est un théorème du lecteur HMT. La
mécanique dimensionnelle applique ensuite la seconde orientation dans la
construction électronique, sans faire de l'électron un antécédent de
\(\pi\), de \(\varphi\) ou de la correspondance.

\subsubsection{Principe d'Ambivalence des lecteurs aréal et volumétrique}
\label{sec:ambivalencia-grav-inercial-pi-phi2}

Le Principe d'Ambivalence affirme que la lecture aréale et la lecture
volumétrique sont deux fonctionnelles coordonnées, mais non interchangeables, d'un
même état avec mémoire. Leur coïncidence sur la diagonale est
\[
 \mathcal L_{\rm ar}^{\rm diag}(x)
 =\mathcal L_{\rm vol}^{\rm diag}(x).
\]
L'égalité diagonale n'efface pas l'ambivalence en dehors de celle-ci.

Sur la frontière aréale--volumétrique, les deux lecteurs se factorisent à travers
une même amplitude positive de mémoire \(\mathcal M(x)\) :
\[
 \mathcal L_{\rm vol}(x)=\mu_{\rm vol}\mathcal M(x),
 \qquad
 \mathcal L_{\rm ar}(x)=
 \pi\mu_{\rm ar}\mathcal M(x).
\]
Alors le théorème de Parry et la marque de retour donnent
\[
 \boxed{\displaystyle
 \frac{\mathcal L_{\rm ar}(x)}
 {\mathcal L_{\rm vol}(x)}
 =\frac{\pi}{\varphi^2}.}
\]
Avec la normalisation volumétrique, la différence relative des lecteurs est
\[
 \frac{\mathcal L_{\rm ar}-\mathcal L_{\rm vol}}
 {\mathcal L_{\rm vol}}
 =
 \frac{\pi}{\varphi^2}-1
 =
 0.1999816148643266611\ldots.
\]
La lecture physique de cette correspondance est effectuée plus tard, dans le chapitre de
Cartan--Holst, après la construction de l'holonomie, de la cotétrade, de la connexion, de la torsion
et de la courbure. Ici, le résultat est entièrement HMT : un rapport exact entre
deux lecteurs de la même mémoire de feuillet.

Il convient de les distinguer de l'extracteur enrichi utilisé ensuite en
mécanique dimensionnelle. Là, une histoire complète conserve le comptage d'action
\(n_A\), les douze événements orientés \(e_0,\ldots,e_{11}\) et la cochaîne
torsionnelle au niveau du pas ; on en obtient
\[
 s_C=\sum_{i=0}^{5}(e_i+e_{i+6}),
 \qquad
 W_\pm=6n_A\pm s_C.
\]
Les canaux déjà engendrés par HMT satisfont alors l'identité exacte
\[
 \exp\!\left(n_AA_{\rm rad}+\frac{s_C}{6}C^\ast_{\rm rad}\right)
 =\left(q_+^{-W_+}q_-^{-W_-}\right)^{1/12}.
\]
Par conséquent, le passage chemin--caractère angulaire n'est pas absent : il s'effectue sur
le registre dodécaphasique enrichi. Ce que ce calcul de phase prouve à lui
seul, ce sont \(\Omega_{\rm sw},S,G\) ; il ne les remplace pas et ne permet pas de reconstruire le
registre complet après l'avoir projeté.

La mesure uniforme sur les \(104\,976\) présentations descend vers des atomes
de masses \(1/729\), \(1/243\) et \(1/54\), mais pas vers \(1/468\). Dans la
microfibre quintuple de \(\pi\), les deux probabilités sont
\[
 \tau_{468}(\mathscr R_\pi)=\frac5{468},
 \qquad
 F_\#\mu_{\rm seed}(\mathscr R_\pi)=\frac7{729}.
\]
Le relèvement moyenné comme le relèvement de phase réalisent la première : ils ne
transportent pas la mesure uniforme sur les germes, mais attribuent la même masse
totale à chaque classe observable. Le relèvement de phase conserve en outre
l'immobilité microscopique dans les pas neutres.

Pour relier le relèvement de phase aux histoires enrichies de
profondeur six, il faut un noyau \(\mathcal K_6\) sur
\(\mathcal H_6^{\rm cat}\times C_9\). Celui-ci descend vers
\(\widehat{\mathcal K}_{\rm ph}\) à travers
\(\rho_6\times\mathrm{id}_{C_9}\) si et seulement s'il satisfait, à chaque phase, la
condition d'agrégabilité forte
\[
 \sum_{\rho_6(g)=y}\mathcal K_6((h,r),(g,r+1))
 =L_{\tau(r)}(\rho_6(h),y),
\]
indépendamment de \(h\) dans chaque fibre de \(\rho_6\). Dans ce cas,
la composition avec \(F\) se projette sur \(\mathcal K_{\rm ph}\). La construction
d'une famille \((\mathcal K_m)_{m\geq6}\) qui satisfait cette identité et
commute avec toutes les troncatures est l'étape restante vers un
transfert de la limite inverse.

\begingroup
\setlength{\parskip}{1pt}
\setlength{\abovedisplayskip}{3pt}\setlength{\belowdisplayskip}{3pt}
\setlength{\abovedisplayshortskip}{2pt}\setlength{\belowdisplayshortskip}{2pt}
\section{Cocycle d'échelle et conservation cylindrique}
\label{sec:cociclo-parry-escala}

Soit \(r=(1,\varphi)\) un vecteur de Perron--Frobenius de
\(F_{\rm av}\). La matrice de transition de Parry est
\[
P_{ij}=\frac{(F_{\rm av})_{ij}r_j}{\varphi r_i}
=
\begin{pmatrix}
0&1\\
\varphi^{-2}&\varphi^{-1}
\end{pmatrix}.
\]
Pour une route admissible \(\gamma:i\to j\) de longueur \(n\),
\[
\mathbb P(\gamma\mid i)=\frac{r_j}{\varphi^nr_i},
\qquad
\boxed{\chi_P(\gamma)=\varphi^n\frac{r_i}{r_j}.}
\]
La télescopie démontre
\[
\chi_P(\gamma_2\circ\gamma_1)
=\chi_P(\gamma_2)\chi_P(\gamma_1).
\]
Sur une route fermée, \(\chi_P=\varphi^n\) ; son carré et son inverse
engendrent les combinaisons
\[
C_n=\frac{\chi_P^2+\chi_P^{-2}}2,\qquad
S_n=\frac{\chi_P^2-\chi_P^{-2}}2,
\]
qui satisfont
\[
x_{n+1}=3x_n-x_{n-1}.
\]

Comme \(F_{\rm av}\) est symétrique, soit
\[
\ell=\frac{r}{1+\varphi^2},\qquad \ell^{\mathsf T}r=1.
\]
La mesure d’un cylindre est
\[
\mu[x_0\cdots x_n]=\ell_{x_0}r_{x_n}\varphi^{-n}.
\]
Par conséquent,
\[
\boxed{
V_n(x)
=\frac{\mu[x_0]}{\mu[x_0\cdots x_n]}
=\varphi^n\frac{r_{x_0}}{r_{x_n}}
=\chi_P(x_0\cdots x_n).}
\]
Pour une échelle de dimension typée \(D>0\),
\[
a_n^{(D)}(x)=V_n(x)^{1/D}
\]
satisfait la loi de conservation cylindrique
\[
\boxed{
\bigl(a_n^{(D)}(x)\bigr)^D
\mu[x_0\cdots x_n]=\mu[x_0].}
\]

En rang spatial trois et sur des retours fermés, nous écrivons
\(a_n:=a_n^{(3)}\) et obtenons
\[
\boxed{
\frac{a_9}{a_0}=\varphi^3,\qquad
\frac{a_{27}}{a_0}=\varphi^9,\qquad
\frac{a_{108}}{a_0}=\varphi^{36}.}
\]
Si \(A_k:=a_{9k}\) pour \(k\geq1\), alors, pour \(k\geq2\),
\[
A_{k+1}-2A_k+A_{k-1}
=A_{k-1}(\varphi^3-1)^2>0.
\]
L'inégalité démontre la convexité en profondeur de retour ; une accélération
physique requiert en outre une carte de durée.

Le calcul premier de
\(\operatorname{Sym}^2(F_{\rm av})\) dans
\cref{sec:retorno-marcado-pi-phi2} démontre que son spectre est
\(\{\varphi^2,-1,\varphi^{-2}\}\). Nous conservons ici uniquement le corollaire
pour le cocycle d'échelle : sur une route fermée avec \(V_n=\varphi^n\), le
mode stable obéit à
\[
\boxed{M_n=\varphi^{-2n}=V_n^{-2}.}
\]
L'identification de ces trois espaces propres aux trois feuillets du TRIT
exige un opérateur d'entrelacement ; la coïncidence de cardinalité ne le remplace pas. Le
corollaire serait réfuté si la route fermée ne satisfaisait pas
\(V_n=\varphi^n\) ou si le mode stable n'avait pas pour valeur propre
\(\varphi^{-2}\).

\section{Théorème de prolongation dynamique : hypothèses globales et domaine certifié}

Les objets suivants du problème fini ont été construits et séparés
explicitement :
\[
 F,
 \qquad
 F_{\rm obs},
 \qquad
 \mathcal K_{\rm ph},
 \qquad
 \widehat{\mathcal K}_{\rm ph},
 \qquad
 K_{\rm cat},
 \qquad
 \tau_{468},
 \qquad
 F_{\rm av},
 \qquad
 \operatorname{Sym}^2(F_{\rm av}),
 \qquad
 D_\pi,
 \qquad
 \mathcal Q.
\]
Pour \(F_{\rm obs}\), le retour déterministe a été démontré et un
témoin explicite de non-descente produit a été exhibé. Pour \(\mathcal K_{\rm ph}\), on a
démontré la double stochasticité, l'irréductibilité, la période neuf, le
retour nonadique uniforme et l'unicité de
\(\pi_{\rm ph}\). Pour \(\widehat{\mathcal K}_{\rm ph}\), on a démontré
l'agrégabilité forte, la conservation microscopique de l'immobilité et la mesure
stationnaire compensée. \(K_{\rm cat}\) n'est plus une moyenne sans
provenance : c'est l'ombre d'un pas en moyennant la phase uniforme.
Le quotient de modes du TPK induit \(F_{\rm av}\) ; son élévation quadratique
sépare les modes \(\varphi^2,-1,\varphi^{-2}\). L'opérateur \(D_\pi\)
marque une seule fois le retour de frontière avec la clôture engendrée, et le
lecteur d'intervalles \(\mathcal Q\) conserve l'emboîtement des cylindres de
\(\pi\) et de \(\varphi\). Ces quatre objets construisent
\(\pi/\varphi^2\) sans confondre la fréquence temporelle,
la mesure uniforme et la mesure de Parry.

La force de cette induction est relative et concrète. Le mot TPK impose les
poids \(1/3\). L'interprétation APP de toutes les continuations compatibles,
avec le comptage uniforme, impose \(E_+\) et \(E_\times\). La trajectoire
déterministe \(F_{\rm obs}\) n'impose pas ces opérateurs et, de fait, ne les
produit pas. Ce n'est pas un défaut du résultat, mais la différence mathématique
entre l'évolution d'un chemin et le transfert d'un espace de chemins.

\ifdefined\HMTInternalTraceability
\begin{criterion}[Élévation enrichie du renversement]
\label{pf84:SYN31-RHO-LIFT}
On cherche un sous-domaine \(D_\rho\subset X_{\rm TPK}^{\rm enr}\) et une application
\[
 \Theta_\rho:D_\rho\longrightarrow X_{\rm TPK}^{\rm enr}
\]
qui relève le renversement entier \(27\mapsto72\) de
\cref{pf84:SYN30-RHO-DESC,pf84:SYN30-CROSS-WITNESS}. Elle doit commuter avec les
lecteurs résiduel, entier et généalogique, quantifier le quotient et la
retenue, et préserver l'orientation, la route, le retour et, lorsque le transport est
inversible, la monodromie, la région et la survie.
Ce n'est que si \(D_\rho\) est totale et invariante sous \(\Theta_\rho\) que l'on peut exiger
\(\Theta_\rho^2=I\) et la nommer involution.

Cette construction ni une démonstration d'involutivité n'existent encore. Le
programme est réfuté si l'application n'existe pas, si l'un des carrés de
lecture échoue, si l'on perd de la mémoire sans la quantifier ou si
\(\Theta_\rho^2\ne I\) dans un domaine où l'on revendique une involution.
\end{criterion}

\begin{criterion}[Élévation enrichie du carré]
\label{pf84:SYN31-SQ-LIFT}
On cherche un sous-domaine \(D_s\subset X_{\rm TPK}^{\rm enr}\) et une application
\[
 \Theta_s:D_s\longrightarrow X_{\rm TPK}^{\rm enr}
\]
qui relève \(s(n)=n^2\), en particulier \(27\mapsto729\), et commute avec les
lecteurs résiduel, entier et généalogique de
\cref{pf84:SYN30-SQ-DESC,pf84:SYN30-CROSS-WITNESS}. Elle doit quantifier le
quotient et la retenue et préserver l'orientation, la route, le retour et, lorsque le
transport est inversible, la monodromie, la région et la survie. Cette application n'est
pas involutive et on ne lui impose pas
\(\Theta_s^2=I\).

La démonstration reste absente. L'inexistence de \(\Theta_s\), l'échec
d'un carré de lecture ou la perte non contrôlée de l'une quelconque des
coordonnées enrichies réfuteraient le programme. L'ablation
\(81\mapsto72\) lors du changement de \(q^\times(9,9)\) est un contrôle indépendant :
ce n'est une branche ni de \(\Theta_\rho\) ni de \(\Theta_s\), et elle ne démontre aucune
des deux applications.
\end{criterion}
\fi

La conclusion finie est complète dans son domaine : le retour uniforme, la
mesure stationnaire, l'agrégabilité à profondeur six et la séparation
entre l'évolution d'une route et le transfert de l'espace des routes sont
démontrés par les objets construits. En particulier,
\(\widehat{\mathcal K}_{\rm ph}\) ne s'identifie pas à la chronologie de
\(F_{\rm obs}\), comme le montre \cref{prop:retorno-fobs-identidad}.
\par\endgroup
 \end{HMTScientificOwnerSubordinated}
}
 
\chapter{Construction corrélative et inséparable du continu}
\chaptermark{Structure discrète du continu}
\label{chap:83-04}

Les cinq constructions sont engendrées corrélativement sur un unique état
APP--TRIT--TPK. Elles constituent des opérations consubstantielles du continu et non une
énumération exhaustive de ses propriétés, cinq branches indépendantes ni une
sélection rétrospective déterminée par des applications ultérieures.

\begin{HMTScientificOwnerSubordinated}
% Corrección exclusivamente editorial del segundo epígrafe de 04b1. El
% fuente teoremática canónico permanece inmutable y se conserva íntegro; aquí se
% traduce la nomenclatura interna «levantamiento TRIT» a una formulación
% pública que nombra la operación y el resultado matemático.
\newcount\HMTContinuoCommonSection
\HMTContinuoCommonSection=0
\let\HMTContinuoCommonOriginalSection\section
\renewcommand{\section}[1]{%
  \advance\HMTContinuoCommonSection by 1
  \ifnum\HMTContinuoCommonSection=2
    \HMTContinuoCommonOriginalSection{Extension tritique de la donnée d'incidence : orientation équilibrée et cocycle de retenue généalogique}%
  \else
    \HMTContinuoCommonOriginalSection{#1}%
  \fi
}
\chapter{État commun, arbre des prolongements et mesure généalogique}
\label{chap:continuo-estado-arbol-medida-comun}

\begin{thesisbox}[title={Un seul état holonomique intégral avant toute discrimination interne}]
L’objet de ce chapitre est l’état fini complet produit par
\(\APP\), orienté par TRIT et transporté par \(\TPK\).  À partir de lui sont
construits le groupoïde des symétries réversibles, l’arbre de toutes les
prolongations admissibles, la multisection des histoires complètes et sa mesure
canonique.  Aucun domaine spécialisé n’est introduit : le domaine est
l’ensemble total des états qu’engendre effectivement le même mécanisme
\(\APP\)--TRIT--\(\TPK\).  On ne sélectionne pas non plus un feuillet, une continuation ou
une lecture ultérieure.
\end{thesisbox}

\section{Données finies produites par APP}
\label{sec:continuo-comun-app-finito}

Pour chaque profondeur \(k\geq0\), nous notons \(\mathsf A_k\) l'ensemble fini des préfixes cellulaires produits par APP, et
\begin{equation}
\label{eq:continuo-comun-app-truncamiento}
 \rho^{\APP}_{k+1,k}:\mathsf A_{k+1}\longrightarrow\mathsf A_k
\end{equation}
l'effacement du dernier raffinement. Cet effacement conserve le préfixe déjà construit. La cellule neutre d'APP fournit un prolongement
\begin{equation}
\label{eq:continuo-comun-app-degeneracion}
 \eta^{\APP}_k:\mathsf A_k\longrightarrow\mathsf A_{k+1},
 \qquad
 \rho^{\APP}_{k+1,k}\eta^{\APP}_k=I_{\mathsf A_k}.
\end{equation}
Écrivons
\begin{equation}
\label{eq:continuo-comun-emision-neutra-app}
 \varepsilon^0_k(a):a\longrightarrow\eta^{\APP}_k(a)
\end{equation}
pour l'émission qui réalise ce prolongement. Ainsi restent distincts sa cible, \(\eta^{\APP}_k(a)\), et l'opération qui la produit, \(\varepsilon^0_k(a)\). On n'interprète pas \(\eta^{\APP}_k\) comme une absence d'histoire : elle enregistre un raffinement de longueur un à émission neutre et conserve matériellement l'état précédent.

Chaque \(a\in\mathsf A_k\) porte simultanément les deux feuillets d'APP,
\begin{equation}
\label{eq:continuo-comun-dos-hojas}
 \operatorname{Sh}_k(a)
 =\bigl(\Sigma_k(a),\Pi_k(a)\bigr)
 \in\mathsf H_k^{\Sigma}\times\mathsf H_k^{\Pi},
\end{equation}
et non un élément choisi dans la paire. Pour une émission élémentaire \(\epsilon:a\to a'\), APP produit en outre deux divisions euclidiennes typées,
\begin{equation}
\label{eq:continuo-comun-residuo-cociente-app}
 \delta^{\diamond}(\epsilon)
 =\operatorname{res}^{\diamond}(\epsilon)
  +9\operatorname{quo}^{\diamond}(\epsilon),
 \qquad
 \operatorname{res}^{\diamond}(\epsilon)\in\{0,\ldots,8\},
 \quad
 \operatorname{quo}^{\diamond}(\epsilon)\in\mathbb Z,
\end{equation}
où \(\diamond\in\{\Sigma,\Pi\}\). L'étiquette \(\diamond\) indique uniquement le feuillet en cours de lecture ; les deux divisions demeurent dans le même registre.

La source et la cible cellulaires de \(\epsilon\) déterminent sa frontière orientée
\begin{equation}
\label{eq:continuo-comun-frontera-elemental}
 \partial\epsilon:=t(\epsilon)-s(\epsilon).
\end{equation}
L'émission neutre satisfait
\begin{equation}
\label{eq:continuo-comun-neutral-app-datos}
 \operatorname{res}^{\diamond}(\varepsilon^0_k(a))=0,
 \qquad
 \operatorname{quo}^{\diamond}(\varepsilon^0_k(a))=0,
 \qquad
 \partial\varepsilon^0_k(a)=0.
\end{equation}

\begin{definition}[Graphe orienté des émissions APP]
\label{def:continuo-comun-grafo-app}
Soit \(\mathsf E_k(a)\) l'ensemble fini des émissions APP de longueur un issues de \(a\in\mathsf A_k\). Nous conservons les émissions avec leur multiplicité : deux opérations ayant la même cible, mais des résidus, quotients ou provenances différents, sont des éléments distincts de \(\mathsf E_k(a)\). Par \eqref{eq:continuo-comun-app-degeneracion},
\begin{equation}
\label{eq:continuo-comun-out-no-vacio-app}
 \varepsilon^0_k(a)\in\mathsf E_k(a),
 \qquad
 1\leq \#\mathsf E_k(a)<\infty.
\end{equation}
\end{definition}

\section{Relèvement TRIT : orientation et retenue généalogique produites}
\label{sec:continuo-comun-trit-lift}

L'incidence signée d'une émission APP est un entier \(\Delta(\epsilon)\). Le représentant équilibré modulo trois donne de manière unique
\begin{equation}
\label{eq:continuo-comun-trit-euclid}
 \Delta(\epsilon)
 =\operatorname{or}(\epsilon)+3\kappa(\epsilon),
 \qquad
 \operatorname{or}(\epsilon)\in\{-1,0,+1\},
 \quad
 \kappa(\epsilon)\in\mathbb Z.
\end{equation}
Ainsi, l'orientation et la retenue ne sont pas des conditions imposées à une histoire : ce sont des fonctions totales de la donnée APP.

Pour rendre la composition visible, soit \(\mathsf C_{\epsilon}:\mathsf M_{s(\epsilon)}\to
\mathsf M_{t(\epsilon)}\) le transport de mémoire induit par l'émission. La retenue native d'une composition satisfait
\begin{equation}
\label{eq:continuo-comun-cociclo-carry}
 \kappa(\epsilon_2\epsilon_1)
 =\kappa(\epsilon_2)
  +\mathsf C_{\epsilon_2}\kappa(\epsilon_1),
\end{equation}
dès que \(t(\epsilon_1)=s(\epsilon_2)\). Pour l'identité et l'émission neutre,
\begin{equation}
\label{eq:continuo-comun-carry-unidad}
 \kappa(I_a)=0,
 \qquad
 \kappa(\varepsilon^0_k(a))=0,
 \qquad
 \mathsf C_{I_a}=I_{\mathsf M_a}.
\end{equation}

\begin{lemma}[Associativité du registre de retenue]
\label{lem:continuo-comun-asociatividad-carry}
Pour trois émissions composables, l'évaluation de \(\kappa(\epsilon_3\epsilon_2\epsilon_1)\) ne dépend pas du parenthésage.
\end{lemma}

\begin{proof}
En appliquant deux fois \eqref{eq:continuo-comun-cociclo-carry} au parenthésage gauche, on obtient
\[
 \kappa(\epsilon_3)
 +\mathsf C_{\epsilon_3}\kappa(\epsilon_2)
 +\mathsf C_{\epsilon_3}\mathsf C_{\epsilon_2}
  \kappa(\epsilon_1).
\]
Le parenthésage droit produit la même expression, car le transport de mémoire est fonctoriel : \(\mathsf C_{\epsilon_3\epsilon_2}
=\mathsf C_{\epsilon_3}\mathsf C_{\epsilon_2}\).
\end{proof}

La réversion orientée d'une émission réversible se note \(\overline\epsilon\). TRIT agit par
\begin{equation}
\label{eq:continuo-comun-reversion-trit}
 \operatorname{or}(\overline\epsilon)
 =-\operatorname{or}(\epsilon),
 \qquad
 \partial\overline\epsilon=-\partial\epsilon,
 \qquad
 \mathsf C_{\overline\epsilon}=\mathsf C_\epsilon^{-1}.
\end{equation}
La retenue inverse est déterminée, et non choisie, par
\begin{equation}
\label{eq:continuo-comun-carry-inverso}
 \kappa(\overline\epsilon)
 =-\mathsf C_\epsilon^{-1}\kappa(\epsilon),
\end{equation}
puisque \(\kappa(\overline\epsilon\epsilon)=0\).

\section{Transport TPK des \(2\) feuillets et de la mémoire}
\label{sec:continuo-comun-tpk-transporte}

La chronologie du TPK fait partie du type de chaque émission et ne se reconstruit pas a posteriori en comptant des lettres. Nous retrouvons ici l'action de l'odomètre de \cref{def:odometro-9-12,prop:proyeccion-fase-longitud}. Soit
\begin{equation}
\label{eq:continuo-comun-ciclo-fase-tpk}
 C_9:=\mathbb Z/9\mathbb Z,\qquad
 \widetilde{\mathcal O}_{9}:=\mathbb Z_{\geq0}\times C_9,\qquad
 \mathcal O_{9,12}:=\mathbb Z/12\mathbb Z\times C_9.
\end{equation}
L'état déroulé de phase et de mémoire est \(\mathsf Q_k=(w_k,r_k)\in\widetilde{\mathcal O}_{9}\) ; sa réduction \((w,r)\mapsto([w]_{12},r)\) restitue \(\mathcal O_{9,12}\). Le germe canonique de l'horloge est
\begin{equation}
\label{eq:continuo-comun-semilla-fase-tpk}
 \mathsf Q_0(a_0):=(0,[0]_9)\qquad(a_0\in\mathsf A_0);
\end{equation}
les autres choix d'origine sont ses translations cohérentes par \(C_9\) et ne modifient pas la loi d'actualisation. Pour \(\bar r\in\{0,\ldots,8\}\), nous posons
\begin{equation}
\label{eq:continuo-comun-carry-fase-tpk}
 \chi_9(r):=\left\lfloor\frac{\bar r+1}{9}\right\rfloor.
\end{equation}
Toute émission \(\epsilon\) faisant passer de la profondeur \(k\) à la profondeur \(k+1\) agit sur l'odomètre par
\begin{equation}
\label{eq:continuo-comun-accion-fase-tpk}
 \sigma_9(w,r):=(w+\chi_9(r),r+[1]_9),\qquad
 \mathsf Q_{k+1}(\epsilon_*\widehat x_k)
 =\sigma_9\mathsf Q_k(\widehat x_k).
\end{equation}
Sa composante de phase est, par définition typée,
\begin{equation}
\label{eq:continuo-comun-fase-emision-tpk}
 \operatorname{ph}(\epsilon):r_k\longmapsto r_k+[1]_9,
 \qquad
 \operatorname{ph}(\epsilon_n\cdots\epsilon_1):
 r_k\longmapsto r_k+[n]_9.
\end{equation}
Nous écrivons \(g_k:=r_k\in C_9\) pour la phase exprimée et \(q^{(9)}_k:=w_k\) pour la mémoire verticale, ainsi que \(\beta_k:=[w_k]_{12}\) pour sa réduction dodécaphasique. L'itération de \eqref{eq:continuo-comun-accion-fase-tpk} est
\begin{equation}
\label{eq:continuo-comun-iteracion-fase-tpk}
 \sigma_9^n(w,r)=
 \left(w+\left\lfloor\frac{\bar r+n}{9}\right\rfloor,
       r+[n]_9\right).
\end{equation}
En effet, les termes \(\chi_9(r),\chi_9(r+[1]_9),\ldots,\chi_9(r+[n-1]_9)\) comptent exactement les passages du représentant \(8\) au représentant \(0\), au nombre de \(\lfloor(\bar r+n)/9\rfloor\). En particulier,
\begin{equation}
\label{eq:continuo-comun-retorno-fase-avance-memoria}
 \sigma_9^9(w,r)=(w+1,r),\qquad
 g_{k+9}=g_k,\qquad q^{(9)}_{k+9}=q^{(9)}_k+1.
\end{equation}
Même l'émission neutre d'APP réalise \(g_k\mapsto g_{k+1}\) : elle est neutre en résidu, quotient, retenue et frontière, mais enregistre qu'une extension a eu lieu. Le mot de phase
\begin{equation}
\label{eq:continuo-comun-palabra-fase-tpk}
 \boldsymbol\tau_k^{(N)}:=(g_k,g_{k+1},\ldots,g_N),
 \qquad g_{j+1}=g_j+[1]_9,
\end{equation}
et la mémoire \(q^{(9)}\) typent par des objets distincts le retour de phase et le changement d'état.

\subsection{Renouvellement simultané de la fibre complète}
\label{sec:continuo-comun-kph}

L'odomètre précédent conserve la chronologie de chaque extension. Le même tour nonadique possède en outre la réalisation markovienne interne, déjà construite sur le catalogue APP--TPK, qui renouvelle simultanément les deux coordonnées de la fibre complète. Il est essentiel de conserver les deux réalisations et de ne pas identifier une arête générique de l'arbre à une émission du catalogue sans application de typage.

\begin{proposition}[Réalisation markovienne interne de la connexion nonadique]
\label{prop:continuo-comun-kph-heredada}
\label{pII-full:thm:seleccion-novena-fase}
Soient le domaine des germes APP, son émission observable et ses multiplicités
\[
 T_{\min}:\Omega_{\rm seed}\twoheadrightarrow\mathcal U_6,
 \qquad
 m(u):=\#T_{\min}^{-1}(u),
\]
et la factorisation produite par le TPK
\[
 \mathcal U_6\simeq\mathcal S_+\times\mathcal S_\times,
 \qquad |\mathcal S_+|=18,\qquad |\mathcal S_\times|=26.
\]
Sur les fonctions définies sur \(\mathcal U_6\), les deux espérances conditionnelles agissent par
\[
 (E_+f)(s,e)=\frac1{18}\sum_{s'\in\mathcal S_+}f(s',e),
 \qquad
 (E_\times f)(s,e)=\frac1{26}
 \sum_{e'\in\mathcal S_\times}f(s,e').
\]
Pour
\[
 \tau=(+1,-1,0,+1,-1,0,+1,-1,0),
 \qquad
 P_{+1}=E_+,\quad P_{-1}=E_\times,\quad P_0=I,
\]
le transfert de catalogue et son relèvement à tous les germes sont
\begin{align}
 \mathcal K_{\rm ph}\bigl((u,r),(v,r+1)\bigr)
 &=P_{\tau(r)}(u,v),
 \label{eq:continuo-comun-kph-catalogo}\\
 L_0(\omega,\omega')&=\mathbf1_{\{\omega=\omega'\}},\nonumber\\
 L_+(\omega,\omega')
 &=\frac{E_+(T_{\min}\omega,T_{\min}\omega')}
         {m(T_{\min}\omega')},\nonumber\\
 L_\times(\omega,\omega')
 &=\frac{E_\times(T_{\min}\omega,T_{\min}\omega')}
         {m(T_{\min}\omega')},\nonumber\\
 \widehat{\mathcal K}_{\rm ph}
 \bigl((\omega,r),(\omega',r+1)\bigr)
 &=L_{\tau(r)}(\omega,\omega').
 \label{eq:continuo-comun-kph-elevada}
\end{align}
Pour \(u=T_{\min}\omega\), on a phase par phase
\begin{equation}
\label{eq:continuo-comun-kph-lumpability}
 \sum_{\omega':\,T_{\min}\omega'=v}
 \widehat{\mathcal K}_{\rm ph}
 \bigl((\omega,r),(\omega',r+1)\bigr)
 =
 \mathcal K_{\rm ph}\bigl((u,r),(v,r+1)\bigr).
\end{equation}
En particulier,
\begin{equation}
\label{eq:continuo-comun-kph-novena-potencia}
 \boxed{\;
 \mathcal K_{\rm ph}^{\,9}\big|_r=E_+E_\times,
 \qquad
 (E_+E_\times)(u,v)=\frac1{468}.
 \;}
\end{equation}
\end{proposition}

\begin{proof}
Dans une phase \(+\) ou \(\times\), la somme de \eqref{eq:continuo-comun-kph-elevada} sur \(T_{\min}^{-1}(v)\) annule exactement le dénominateur \(m(v)\) ; dans une phase zéro, \(L_0\) se projette sur l'identité. Cela démontre \eqref{eq:continuo-comun-kph-lumpability}. Les opérateurs \(E_+\), \(E_\times\) et \(I\) sont idempotents et commutent, et chacun apparaît trois fois dans le mot nonadique. Par conséquent
\[
 \prod_{r\in C_9}P_{\tau(r)}
 =E_+^3E_\times^3I^3=E_+E_\times.
\]
La composition renouvelle d'abord une coordonnée de cardinal \(18\), puis l'autre de cardinal \(26\), de sorte que chaque cible reçoit \(1/(18\cdot26)=1/468\).
\end{proof}

\begin{remark}[Deux réalisations internes, sans identification fictive]
\label{rem:continuo-comun-kph-no-identificacion-aristas}
Le transfert \(\mathcal K_{\rm ph}\) est la réalisation globale du catalogue complet, et \(\sigma_9\) est l'action chronologique sur chaque état de l'arbre. L'une ne se substitue pas à l'autre. Une égalité arête par arête nécessiterait une application supplémentaire
\[
 \operatorname{Cat}_{k,x}:\operatorname{Out}(x)\longrightarrow\mathcal U_6
\]
vérifiant l'identité correspondante de somme sur les fibres. Cette application n'est pas utilisée dans ce volume. La monodromie complète \eqref{eq:continuo-comun-kph-novena-potencia} provient du catalogue APP--TPK, tandis que la généalogie de chaque histoire demeure dans l'arbre et son registre.
\end{remark}

TPK associe à chaque émission admissible un transport typé de la paire de feuillets,
\begin{equation}
\label{eq:continuo-comun-transporte-hojas}
 U_\epsilon:
 \mathsf H_{s(\epsilon)}^{\Sigma}\times
 \mathsf H_{s(\epsilon)}^{\Pi}
 \longrightarrow
 \mathsf H_{t(\epsilon)}^{\Sigma}\times
 \mathsf H_{t(\epsilon)}^{\Pi},
\end{equation}
avec
\begin{equation}
\label{eq:continuo-comun-transporte-hojas-funtorial}
 U_{I_a}=I,
 \qquad
 U_{\epsilon_2\epsilon_1}=U_{\epsilon_2}U_{\epsilon_1}.
\end{equation}
Le codomaine de \(U_\epsilon\) demeure la paire complète. Aucune opération de ce chapitre ne remplace \((\Sigma,\Pi)\) par l'une de ses coordonnées.

Soit \(m\in\mathsf M_{s(\epsilon)}\) une mémoire accumulée. L'action TPK sur celle-ci est la transformation affine
\begin{equation}
\label{eq:continuo-comun-accion-memoria}
 \mathsf U_\epsilon(m)
 :=\mathsf C_\epsilon m+\kappa(\epsilon).
\end{equation}

\begin{lemma}[Action TPK sur la mémoire]
\label{lem:continuo-comun-accion-memoria}
Les transformations \(\mathsf U_\epsilon\) satisfont
\begin{equation}
\label{eq:continuo-comun-accion-memoria-composicion}
 \mathsf U_{I_a}=I,
 \qquad
 \mathsf U_{\epsilon_2\epsilon_1}
 =\mathsf U_{\epsilon_2}\mathsf U_{\epsilon_1}.
\end{equation}
\end{lemma}

\begin{proof}
L'identité découle de \eqref{eq:continuo-comun-carry-unidad}. Pour la composition,
\begin{align*}
 \mathsf U_{\epsilon_2}\mathsf U_{\epsilon_1}(m)
 &=\mathsf C_{\epsilon_2}
   \bigl(\mathsf C_{\epsilon_1}m+\kappa(\epsilon_1)\bigr)
   +\kappa(\epsilon_2)\\
 &=\mathsf C_{\epsilon_2\epsilon_1}m
   +\kappa(\epsilon_2)
   +\mathsf C_{\epsilon_2}\kappa(\epsilon_1)\\
 &=\mathsf C_{\epsilon_2\epsilon_1}m
   +\kappa(\epsilon_2\epsilon_1),
\end{align*}
où la dernière égalité est \eqref{eq:continuo-comun-cociclo-carry}.
\end{proof}

Un parcours fini est un mot composable
\begin{equation}
\label{eq:continuo-comun-ruta}
 \gamma=\epsilon_1\epsilon_2\cdots\epsilon_k.
\end{equation}
Sa frontière, son transport, sa retenue et sa mémoire sont
\begin{align}
 \partial\gamma
 &:=\sum_{j=1}^{k}\partial\epsilon_j,
 \label{eq:continuo-comun-frontera-ruta}\\
 U_\gamma
 &:=U_{\epsilon_k}\cdots U_{\epsilon_1},
 \label{eq:continuo-comun-transporte-ruta}\\
 \kappa(\gamma)
 &:=\kappa(\epsilon_k)
   +\sum_{j=1}^{k-1}
    \mathsf C_{\epsilon_k}\cdots\mathsf C_{\epsilon_{j+1}}
    \kappa(\epsilon_j),
 \label{eq:continuo-comun-carry-ruta}\\
 m(\gamma;m_0)
 &:=\mathsf U_{\epsilon_k}\cdots
       \mathsf U_{\epsilon_1}(m_0).
 \label{eq:continuo-comun-memoria-ruta}
\end{align}
Dans \eqref{eq:continuo-comun-frontera-ruta}, les extrémités intérieures s'annulent et il reste
\begin{equation}
\label{eq:continuo-comun-frontera-telescopica}
 \partial\gamma=t(\epsilon_k)-s(\epsilon_1).
\end{equation}

\section{L'état commun et son domaine total}
\label{sec:continuo-comun-estado-total}

\begin{definition}[Registre élémentaire]
\label{def:continuo-comun-ledger-elemental}
Pour une émission située entre deux préfixes enrichis consécutifs, nous notons \(\mathsf Q^-(\epsilon)=(w^-(\epsilon),r^-(\epsilon))\) l'odomètre du préfixe source et
\begin{equation}
\label{eq:continuo-comun-ledger-fase-extremos}
 \mathsf Q^+(\epsilon):=\sigma_9\mathsf Q^-(\epsilon)
 =(w^-(\epsilon)+\chi_9(r^-(\epsilon)),r^-(\epsilon)+[1]_9)
\end{equation}
celui du préfixe cible. Ainsi, l'odomètre n'est pas appliqué à une cellule APP nue. Le registre de l'émission est
\begin{equation}
\label{eq:continuo-comun-ledger-elemental}
 \lambda(\epsilon):=\bigl(
 s(\epsilon),t(\epsilon),
 \operatorname{res}^{\Sigma}(\epsilon),
 \operatorname{quo}^{\Sigma}(\epsilon),
 \operatorname{res}^{\Pi}(\epsilon),
 \operatorname{quo}^{\Pi}(\epsilon),
 \operatorname{or}(\epsilon),
 \kappa(\epsilon),
 \partial\epsilon,
 \mathsf Q^-(\epsilon),\mathsf Q^+(\epsilon)\bigr).
\end{equation}
Le registre de \(\gamma\) est le mot ordonné
\begin{equation}
\label{eq:continuo-comun-ledger-ruta}
 \operatorname{Led}(\gamma)
 :=\bigl(\lambda(\epsilon_1),\ldots,
          \lambda(\epsilon_k)\bigr).
\end{equation}
Conserver ce mot, et non seulement sa somme, distingue deux histoires possédant la même sortie terminale.
\end{definition}

\subsection{Lecteur complet de l'état TPK}
\label{sec:continuo-comun-lector-estado-tpk}

La profondeur \(K\) des blocs TPK et la profondeur \(k\) de l'arbre des extensions sont des types distincts. Dans cette sous-section, \(K\) désigne exclusivement la profondeur de l'état TPK déjà engendré par le mécanisme opératoriel ; elle ne s'identifie ni à \(k\) ni à une arête générique.

\begin{definition}[Lecteur compact de mémoire nonadique]
\label{def:continuo-comun-lector-compacto-nonadico}
L’état holonomique intégral TPK de profondeur \(K\) est
\begin{equation}
\label{eq:continuo-comun-estado-tpk-completo}
 v_K=
 \bigl(
 B_K,\boldsymbol\delta_K,\boldsymbol x_K,
 \boldsymbol{\mathcal C}_K,\Sigma_K,
 \omega_K,\ell_K,g(K),\Lambda_K
 \bigr),
 \qquad
 B_K=(u_K^\pi,u_K^e,u_K^\varphi).
\end{equation}
Son lecteur compact n'ajoute pas de coordonnées : il les exprime à partir de ce même état,
\begin{equation}
\label{eq:continuo-comun-lector-compacto-nonadico}
 \mathfrak r_9(v_K)
 :=\bigl(B_K,C_K,p_K,c_K,\Sigma_K,W_K,g(K)\bigr),
\end{equation}
où
\begin{equation}
\label{eq:continuo-comun-lector-compacto-componentes}
\begin{aligned}
 C_K&:=\boldsymbol{\mathcal C}_K
      =\bigl(C_K^\chi\bigr)_{\chi\in\{\pi,e,\varphi\}},\\
 C_K^\chi&:=(N_K^\chi,L_K^\chi,r_K^\chi,
             D_K^\chi,A_K^\chi,B_K^\chi),\\
 p_K&:=(B_0,\ldots,B_K),\\
 c_K&:=(\boldsymbol\delta_K,\boldsymbol x_K,
       \kappa(\gamma^{\rm TPK}_{0:K})),
\end{aligned}
\end{equation}
et la projection orientée de Witt est conservée sous la forme
\begin{equation}
\label{eq:continuo-comun-sombra-witt}
 W_K:=
 \bigl((u_K^\chi,u_K^\chi A_W)\bigr)_
       {\chi\in\{\pi,e,\varphi\}},
 \qquad
 A_W=
 \begin{pmatrix}
 0&1&1&1&1&1\\
 1&0&1&2&2&1\\
 1&1&0&1&2&2\\
 1&2&1&0&1&2\\
 1&2&2&1&0&1\\
 1&1&2&2&1&0
 \end{pmatrix},
 \qquad A_W^2=-I_6.
\end{equation}
Le parcours \(\gamma^{\rm TPK}_{0:K}\) est le mot des extensions codé par \(\Lambda_K\) ; son préfixe de longueur \(j\) est celui qui produit \(v_j\). La signature est la sortie de la primitive interne explicite
\begin{equation}
\label{eq:continuo-comun-primitiva-firma-tpk}
 \Sigma_B:
 \mathsf B_{\rm TPK}\times\mathsf C_{\rm TPK}
 \times\mathbb Z[\mathsf A]\times\Omega_{12}\times\mathsf L_{\rm TPK}
 \longrightarrow\mathsf S_{\rm TPK},
 \quad
 (B,c,\partial\gamma,\omega,\ell)\longmapsto
 \Sigma_B(B,c,\partial\gamma,\omega,\ell),
\end{equation}
avec \(\Sigma_B(g\cdot-)=g\cdot\Sigma_B(-)\) pour chaque symétrie TPK. Sa structure initiale explicite est \((B_K,c_K,\partial\gamma^{\rm TPK}_{0:K},\omega_K,\ell_K)\). Elle n'est pas remplacée par une signature extérieure ni inférée rétrospectivement à partir d'un lecteur numérique.
\end{definition}

L'action du pas \(K\mapsto K+1\) s'écrit sans omettre l'état profond. Pour chaque \(\chi\in\{\pi,e,\varphi\}\),
\begin{equation}
\label{eq:continuo-comun-update-hensel-tpk}
 y_K^\chi=x_K^\chi A_H,\qquad
 u_{K+1}^\chi=y_K^\chi\bmod3,\qquad
 x_{K+1}^\chi=\frac{y_K^\chi-u_{K+1}^\chi}{3},
\end{equation}
avec
\[
 A_H=
 \begin{pmatrix}
 2&2&2&1&2&1\\
 2&1&2&2&1&1\\
 1&0&0&1&0&2\\
 0&0&1&1&1&0\\
 2&2&2&0&2&1\\
 2&1&2&1&0&0
 \end{pmatrix},
 \qquad \det A_H=1.
\]
Si
\[
 \operatorname{ind}_3(u_1,\ldots,u_6)
 :=\sum_{j=1}^{6}\overline u_j\,3^{6-j}
 \in\{0,\ldots,728\},
\]
l'actualisation du préfixe et de ses \(729\) sous-cylindres est
\begin{align}
 N_{K+1}^\chi
 &=729N_K^\chi+\operatorname{ind}_3(u_{K+1}^\chi),
 &L_{K+1}^\chi&=L_K^\chi+6,
 \label{eq:continuo-comun-update-cilindro-indice}\\
 I_{K,u}^\chi
 &:=
 \left[
  \frac{729N_K^\chi+u}{3^{L_K^\chi+6}},
  \frac{729N_K^\chi+u+1}{3^{L_K^\chi+6}}
 \right),
 &0&\leq u<729.
 \label{eq:continuo-comun-update-cilindro-intervalo}
\end{align}
Le cylindre décimal et la retenue associée s'actualisent selon les deux cas entiers du même opérateur \(\mathcal C_{3,1000}\). Si aucune nouvelle triade décimale n'est exprimée,
\begin{equation}
\label{eq:continuo-comun-update-cilindro-sin-triada}
 r_{K+1}^\chi=r_K^\chi,\qquad D_{K+1}^\chi=D_K^\chi,\qquad
 A_{K+1}^\chi=729A_K^\chi+
 \operatorname{ind}_3(u_{K+1}^\chi)1000^{r_K^\chi},\qquad
 B_{K+1}^\chi=A_{K+1}^\chi+1000^{r_K^\chi}.
\end{equation}
La primitive \(\mathcal C_{3,1000}\) produit la donnée optionnelle
\[
 d_{K+1,\chi}^{\rm dec}\in
 \{\bot\}\sqcup\{0,\ldots,999\};
\]
dans le premier cas, elle vaut \(d_{K+1,\chi}^{\rm dec}=\bot\). Si une triade apparaît, nous écrivons \(d_{K+1,\chi}^{\rm dec}=\delta_{K+1}^\chi\in\{0,\ldots,999\}\) et
\begin{equation}
\label{eq:continuo-comun-update-cilindro-con-triada}
 \begin{aligned}
 r_{K+1}^\chi&=r_K^\chi+1,\qquad
 D_{K+1}^\chi=1000D_K^\chi+\delta_{K+1}^\chi,\\
 A_{K+1}^\chi&=729000A_K^\chi
 +\operatorname{ind}_3(u_{K+1}^\chi)1000^{r_K^\chi+1}
 -729\delta_{K+1}^\chi3^{L_K^\chi},\\
 B_{K+1}^\chi&=A_{K+1}^\chi+1000^{r_K^\chi+1}.
 \end{aligned}
\end{equation}
Le registre \(\boldsymbol\delta_{K+1}\) ajoute ces triades et les deux cocycles conjoints
\begin{equation}
\label{eq:continuo-comun-update-triadas-carry-conjunto}
 \begin{gathered}
 q_{K+1}=(u_{K+1}^\pi+u_{K+1}^e+u_{K+1}^\varphi)\bmod3,
 \quad
 c_{K+1}^{\rm joint}=\frac{u_{K+1}^\pi+u_{K+1}^e+u_{K+1}^\varphi-q_{K+1}}3,\\
 a_{K+1}=(u_{K+1}^\pi+u_{K+1}^e-u_{K+1}^\varphi)\bmod3,
 \quad
 h_{K+1}^{\rm joint}=\frac{u_{K+1}^\pi+u_{K+1}^e-u_{K+1}^\varphi-a_{K+1}}3,\\
 \boldsymbol\delta_{K+1}
 =\boldsymbol\delta_K\star
 \bigl((d_{K+1,\chi}^{\rm dec})_\chi,
       c_{K+1}^{\rm joint},h_{K+1}^{\rm joint}\bigr).
 \end{gathered}
\end{equation}
Ainsi, \(C_{K+1}\), \(\boldsymbol\delta_{K+1}\) et \(\boldsymbol x_{K+1}\) possèdent des actions concrètes, non seulement des noms. Les autres coordonnées s'actualisent simultanément par
\begin{align}
 p_{K+1}&=p_K\star B_{K+1},
 \label{eq:continuo-comun-update-prefijo-tpk}\\
 \kappa(\gamma^{\rm TPK}_{0:K+1})
 &=\kappa(\varepsilon^{\rm TPK}_{K+1})
   +\mathsf C_{\varepsilon^{\rm TPK}_{K+1}}
    \kappa(\gamma^{\rm TPK}_{0:K}),
 \label{eq:continuo-comun-update-carry-lector-tpk}\\
 c_{K+1}^{\rm comp}
 &:=(\boldsymbol\delta_{K+1},\boldsymbol x_{K+1},
       \kappa(\gamma^{\rm TPK}_{0:K+1})),
 \label{eq:continuo-comun-update-carry-compacto-tpk}\\
 \Sigma_{K+1}
 &=\Sigma_B\bigl(
 B_{K+1},c_{K+1}^{\rm comp},
 \partial\gamma^{\rm TPK}_{0:K}
 +\partial\varepsilon^{\rm TPK}_{K+1},
 \omega_{K+1},\ell_{K+1}\bigr),
 \label{eq:continuo-comun-update-firma-tpk}\\
 W_{K+1}
 &=\bigl((u_{K+1}^\chi,u_{K+1}^\chi A_W)\bigr)_\chi,
 &g(K+1)&=g(K)+[1]_9.
 \label{eq:continuo-comun-update-witt-fase-tpk}
\end{align}
Puisque \(A_H\) est inversible modulo trois et que \(A_W^2=-I_6\), les actualisations résiduelle et de projection ne confondent pas des blocs distincts. La troncature conserve les suites complètes
\[
 (B_0,\ldots,B_K),\ (C_0,\ldots,C_K),\
 (p_0,\ldots,p_K),\ (c_0,\ldots,c_K),\
 (\Sigma_0,\ldots,\Sigma_K),\ (W_0,\ldots,W_K)
\]
et supprime simultanément leur dernière entrée. Le lecteur \(\mathfrak r_9\) commute donc avec l'effacement d'un bloc.

\begin{proposition}[Portée universelle et portée certifiée de l'absence de réinitialisation]
\label{prop:continuo-comun-no-reinicio-alcance}
Dans tout tour horizontal de neuf pas sont conservés
\[
 g(K+9)=g(K),\qquad
 q_{\rm mem}(K+9)=q_{\rm mem}(K)+1,
\]
et le préfixe se prolonge par neuf émissions, si bien que l'état ne se réinitialise pas. Dans la section conjointe certifiée \(\chi\in\{\pi,e,\varphi\}\), \(0\leq K<387\), on a en outre
\begin{equation}
\label{eq:continuo-comun-no-reinicio-certificado}
 (p_K,C_K)\neq(p_{K+9},C_{K+9}),\qquad
 (B_K,\mathbf N_K,W_K)\neq
 (B_{K+9},\mathbf N_{K+9},W_{K+9}),
 \qquad
 \mathbf N_K:=(N_K^\pi,N_K^e,N_K^\varphi).
\end{equation}
\end{proposition}

\begin{proof}
Les deux identités universelles sont l'itération exacte de \eqref{eq:continuo-comun-accion-fase-tpk} ; le préfixe conserve les neuf entrées ajoutées par \eqref{eq:continuo-comun-update-prefijo-tpk}. Pour la section conjointe, la récurrence \eqref{eq:continuo-comun-update-cilindro-indice} compare les \(387\) paires nonadiques de chacun des trois canaux et produit \eqref{eq:continuo-comun-no-reinicio-certificado} ; c'est l'affirmation certifiée dans \cref{thm:generacion-monodromica-conjunta-rev7}. Enfin, \(u\mapsto(u,uA_W)\) est injective, de sorte que le changement de bloc conserve également celui de sa projection complète. Cette dernière inégalité n'est pas étendue à un prolongement neutre arbitraire hors de la section certifiée.
\end{proof}

\begin{definition}[État APP--TRIT--TPK de profondeur \(k\)]
\label{def:continuo-comun-estado}
Pour un chemin \(\gamma=\epsilon_1\cdots\epsilon_k\) qui naît dans une cellule
APP \(a_0\), l’état holonomique intégral commun est
\begin{equation}
\label{eq:continuo-comun-estado}
 \widehat x_k(\gamma):=\Bigl(
  a_0,\ldots,a_k;
  \operatorname{Sh}_0,\ldots,\operatorname{Sh}_k;
  \mathbf{res}_k,\mathbf{quo}_k;
  \mathbf{or}_k,\mathbf\kappa_k;
  \mathsf Q_k;g_k;q^{(9)}_k;\beta_k;m_k;\partial\gamma;\gamma;
  \operatorname{Led}(\gamma)
 \Bigr),
\end{equation}
où
\begin{align*}
 \mathbf{res}_k
 &:=\bigl(
  (\operatorname{res}^{\Sigma}(\epsilon_j),
   \operatorname{res}^{\Pi}(\epsilon_j))
  \bigr)_{1\leq j\leq k},\\
 \mathbf{quo}_k
 &:=\bigl(
  (\operatorname{quo}^{\Sigma}(\epsilon_j),
   \operatorname{quo}^{\Pi}(\epsilon_j))
  \bigr)_{1\leq j\leq k},\\
 \mathbf{or}_k
 &:=\bigl(\operatorname{or}(\epsilon_j)\bigr)_{1\leq j\leq k},
 \qquad
 \mathbf\kappa_k
 :=\bigl(\kappa(\epsilon_j)\bigr)_{1\leq j\leq k},\\
 m_k&:=m(\gamma;m_0).
\end{align*}
Les feuillets intermédiaires s'obtiennent par \(\operatorname{Sh}_{j}=U_{\epsilon_j}\operatorname{Sh}_{j-1}\).
\end{definition}

La présence simultanée du parcours et du registre empêche d'identifier des mots distincts qui se terminent dans la même cellule. L'état n'est pas une liste de cardinaux : chaque coordonnée possède les équations de type et de composition des sections précédentes.

\begin{definition}[Domaine total de profondeur finie]
\label{def:continuo-comun-dominio-total}
Soit
\begin{equation}
\label{eq:continuo-comun-dominio-total}
 \widehat{\mathcal X}^{\rm tot}_k
 :=\left\{
  \widehat x_k(\epsilon_1\cdots\epsilon_k):
  \begin{array}{l}
   a_0\in\mathsf A_0,\\
   \epsilon_j\in\mathsf E_{j-1}(a_{j-1}),\\
   s(\epsilon_j)=a_{j-1},\ t(\epsilon_j)=a_j
  \end{array}
 \right\}.
\end{equation}
Aucune condition d'une expression ultérieure n'est ajoutée à \eqref{eq:continuo-comun-dominio-total}. Appartenir au domaine signifie exactement avoir été produit par les actions totales \eqref{eq:continuo-comun-residuo-cociente-app}, \eqref{eq:continuo-comun-trit-euclid} et \eqref{eq:continuo-comun-accion-fase-tpk}--\eqref{eq:continuo-comun-accion-memoria}.
\end{definition}

La troncature commune efface la dernière émission ainsi que toutes les entrées du registre qu'elle a produites, et elles seules :
\begin{equation}
\label{eq:continuo-comun-truncamiento-total}
 \rho_{k+1,k}:\widehat{\mathcal X}^{\rm tot}_{k+1}
 \longrightarrow\widehat{\mathcal X}^{\rm tot}_k,
 \qquad
 \rho_{k+1,k}\widehat x_{k+1}
 =\widehat x_k.
\end{equation}
Soit \(\varepsilon^0_k(x)\) l'émission neutre \eqref{eq:continuo-comun-emision-neutra-app} appliquée à l'extrémité APP de \(x\). Le prolongement neutre définit
\begin{equation}
\label{eq:continuo-comun-seccion-total}
 \eta_k:\widehat{\mathcal X}^{\rm tot}_k
 \longrightarrow\widehat{\mathcal X}^{\rm tot}_{k+1},
 \qquad
 \eta_k(x):=(\varepsilon^0_k(x))_*x,
 \qquad
 \rho_{k+1,k}\eta_k=I.
\end{equation}

\begin{proposition}[Finitude, non-vacuité et surjectivité]
\label{prop:continuo-comun-total-finito}
Pour tout \(k\), l'ensemble \(\widehat{\mathcal X}^{\rm tot}_k\) est fini et non vide, et \(\rho_{k+1,k}\) est surjectif.
\end{proposition}

\begin{proof}
APP produit un ensemble fini non vide à la profondeur zéro. À chaque pas, chaque état possède un ensemble fini d'émissions par \eqref{eq:continuo-comun-out-no-vacio-app} ; par récurrence, il n'existe qu'un nombre fini de mots de longueur \(k\). La suite des prolongements neutres produit au moins un mot de chaque longueur. Enfin, \eqref{eq:continuo-comun-seccion-total} est un inverse à droite de la troncature, qui est donc surjective.
\end{proof}

\section{Groupoïde fini des symétries réversibles}
\label{sec:continuo-comun-groupoide}

Les opérations réversibles d'APP, de TRIT et du TPK agissent sur l'état entier. Soit \(\mathsf U\) le groupe des systèmes cohérents \(g=(g_k)_{k\geq0}\) engendrés par ces opérations, où
\begin{equation}
\label{eq:continuo-comun-simetria-coherente}
 \rho_{k+1,k}g_{k+1}=g_k\rho_{k+1,k},
 \qquad
 g_{k+1}\eta_k=\eta_kg_k.
\end{equation}
Nous notons \(\mathsf U_k\) son image dans \(\operatorname{Sym}(\widehat{\mathcal X}^{\rm tot}_k)\). Cette image est finie puisqu'elle agit sur un ensemble fini. L'action de \(g_k\) sur
\begin{equation*}
 \widehat x=(a_{\leq k};\operatorname{Sh}_{\leq k};
 \mathbf{res},\mathbf{quo};\mathbf{or},\mathbf\kappa;
 \mathsf Q_k;g_k;q_k^{(9)};\beta_k;m;\partial\gamma;\gamma;
 \operatorname{Led}(\gamma))
\end{equation*}
est
\begin{equation}
\label{eq:continuo-comun-accion-reversible}
\begin{aligned}
 g_k\cdot\widehat x:=\bigl(&
 g_k a_{\leq k};
 U_{g_k}\operatorname{Sh}_{\leq k};
 \mathbf{res}(g_k\gamma),\mathbf{quo}(g_k\gamma);
 \mathbf{or}(g_k\gamma),\mathbf\kappa(g_k\gamma);\\
 &\mathsf Q(g_k\gamma);g(g_k\gamma);q^{(9)}(g_k\gamma);
 \beta(g_k\gamma);
 \mathsf U_{g_k}(m);
 (g_k)_*\partial\gamma;
 g_k\gamma;
 \operatorname{Led}(g_k\gamma)\bigr).
\end{aligned}
\end{equation}
Tous les termes du membre de droite sont recalculés selon les règles APP, TRIT et TPK ; ils ne sont pas copiés comme des étiquettes indépendantes.

\begin{definition}[Groupoïde d'action commun]
\label{def:continuo-comun-groupoide}
Le groupoïde fini de profondeur \(k\) est
\begin{equation}
\label{eq:continuo-comun-groupoide}
 \mathfrak G_k
 :=\mathsf U_k\ltimes\widehat{\mathcal X}^{\rm tot}_k.
\end{equation}
Ses objets sont les états communs. Une flèche \((g_k,\widehat x)\) a pour source \(\widehat x\), pour cible \(g_k\cdot\widehat x\), pour identité \((I,\widehat x)\) et pour inverse \((g_k^{-1},g_k\cdot\widehat x)\).
\end{definition}

\begin{lemma}[L'action conserve toutes les équations d'état]
\label{lem:continuo-comun-groupoide-conserva}
L'action \eqref{eq:continuo-comun-accion-reversible} conserve la paire de feuillets, les divisions APP, l'orientation TRIT, la retenue, l'odomètre de phase, la mémoire verticale, sa réduction dodécaphasique, la mémoire affine, la frontière, le parcours et le registre, dans leurs types respectifs. En outre,
\begin{equation}
\label{eq:continuo-comun-groupoide-accion}
 I\cdot\widehat x=\widehat x,
 \qquad
 (g_{2,k}g_{1,k})\cdot\widehat x
 =g_{2,k}\cdot(g_{1,k}\cdot\widehat x).
\end{equation}
\end{lemma}

\begin{proof}
Les divisions APP sont uniques, de sorte que l'action sur une émission détermine à nouveau ses résidus et quotients. La décomposition équilibrée \eqref{eq:continuo-comun-trit-euclid} est également unique. La composition de la retenue est associative par \cref{lem:continuo-comun-asociatividad-carry} ; l'action sur la mémoire l'est par \cref{lem:continuo-comun-accion-memoria} ; et la frontière est fonctorielle par \eqref{eq:continuo-comun-frontera-telescopica}. Enfin, l'action sur le mot ordonné induit l'action sur son registre entrée par entrée. La cohérence avec l'odomètre recalcule \(Q\), \(g\) et \(q^{(9)}\) à partir du parcours transformé et conserve \eqref{eq:continuo-comun-accion-fase-tpk}. Ces identités prouvent \eqref{eq:continuo-comun-groupoide-accion}.
\end{proof}

\section{Arbre fini de tous les prolongements}
\label{sec:continuo-comun-arbol}

Pour \(x\in\widehat{\mathcal X}^{\rm tot}_k\), soit \(\operatorname{Out}(x)\) l'ensemble de toutes les émissions TPK admissibles de longueur un issues de l'extrémité de \(x\). Une émission appartient à \(\operatorname{Out}(x)\) uniquement lorsque son action actualise conjointement toutes les coordonnées de \eqref{eq:continuo-comun-estado}. En particulier,
\begin{equation}
\label{eq:continuo-comun-out-finito}
 1\leq d(x):=\#\operatorname{Out}(x)<\infty,
 \qquad
 \varepsilon^0_k(x)\in\operatorname{Out}(x).
\end{equation}

\begin{definition}[Arbre tronqué des prolongements]
\label{def:continuo-comun-arbol-finito}
Étant donnés \(N\geq k\) et \(x\in\widehat{\mathcal X}^{\rm tot}_k\), nous définissons \(\mathscr T_{k,N}(x)\) comme l'arbre enraciné dont les sommets de hauteur \(j\), \(0\leq j\leq N-k\), sont les mots
\begin{equation}
\label{eq:continuo-comun-vertice-arbol}
 h=(\epsilon_{k+1},\ldots,\epsilon_{k+j}),
 \qquad
 \epsilon_{n+1}\in\operatorname{Out}(x_n),
 \qquad
 x_{n+1}=(\epsilon_{n+1})_*x_n,
\end{equation}
avec pour racine le mot vide. Une arête ajoute une émission à la fin. Deux mots atteignant le même état terminal restent des sommets distincts si leurs registres diffèrent.
\end{definition}

La dernière condition est essentielle : l'arbre enregistre des généalogies, non le quotient qui ne retient que la cible. Son niveau terminal est
\begin{equation}
\label{eq:continuo-comun-terminal-finito}
 \operatorname{Term}_{k,N}(x)
 :=\{h\in\mathscr T_{k,N}(x):|h|=N-k\}.
\end{equation}
Chaque \(h\) conserve l'état terminal complet
\begin{equation}
\label{eq:continuo-comun-registro-terminal}
 \operatorname{ter}_{N}(h)
 :=\bigl(x_N,m_N,\partial(\gamma h),
          \gamma h,\operatorname{Led}(\gamma h)\bigr).
\end{equation}

\begin{proposition}[Finitude et non-vacuité à tout horizon]
\label{prop:continuo-comun-arbol-finito-no-vacio}
Pour tous \(x\), \(k\) et \(N\geq k\), l'arbre \(\mathscr T_{k,N}(x)\) est fini et \(\operatorname{Term}_{k,N}(x)\neq\varnothing\).
\end{proposition}

\begin{proof}
La racine possède un nombre fini et strictement positif de successeurs par \eqref{eq:continuo-comun-out-finito}. En appliquant le même argument à chaque niveau, on obtient par récurrence la finitude à tout horizon fini. Le mot formé d'émissions neutres \((\varepsilon^0_k,\varepsilon^0_{k+1},\ldots,
\varepsilon^0_{N-1})\) est un sommet terminal, de sorte que le niveau \eqref{eq:continuo-comun-terminal-finito} n'est pas vide.
\end{proof}

L'effacement de la dernière émission définit
\begin{equation}
\label{eq:continuo-comun-proyeccion-terminal}
 \pi_{N+1,N}:\operatorname{Term}_{k,N+1}(x)
 \longrightarrow\operatorname{Term}_{k,N}(x).
\end{equation}

\begin{lemma}[Surjectivité terminale]
\label{lem:continuo-comun-terminal-sobreyectivo}
L'application \(\pi_{N+1,N}\) est surjective pour tout \(N\geq k\).
\end{lemma}

\begin{proof}
Soit \(h\in\operatorname{Term}_{k,N}(x)\), d'état terminal \(v_h\). L'ajout de l'émission \(\varepsilon_N^0(v_h)\) produit \(h'\in\operatorname{Term}_{k,N+1}(x)\) et \(\pi_{N+1,N}(h')=h\).
\end{proof}

\section{Connexion nonadique : phase, holonomie et mémoire}
\label{sec:continuo-comun-gamma9-total}

Un mot de neuf émissions n'est pas appelé holonomie au seul motif qu'il est de longueur neuf. On conserve d'abord la connexion horizontale déterminée par le foncteur de phase. Pour \(r\in C_9\), nous définissons
\begin{equation}
\label{eq:continuo-comun-relacion-horizontal-fase}
 \begin{split}
 \mathscr R_{r,k}:=\bigl\{
 (x,\epsilon,\epsilon_*x):\;&
 x\in\widehat{\mathcal X}^{\rm tot}_k,\quad
 \epsilon\in\operatorname{Out}(x),\\
 &\mathsf Q_k(x)=(w,r),\quad
 \mathsf Q_{k+1}(\epsilon_*x)=\sigma_9(w,r)
 \bigr\}.
 \end{split}
\end{equation}
La dernière égalité ne restreint pas le domaine : c'est l'action totale \eqref{eq:continuo-comun-accion-fase-tpk} de toute émission TPK élémentaire. La connexion transporte simultanément feuillet, mémoire et registre :
\begin{equation}
\label{eq:continuo-comun-conexion-tpk-generador}
 \nabla_{\rm TPK}(\epsilon):=
 \bigl(\operatorname{ph}(\epsilon),U_\epsilon,
       \mathsf U_\epsilon,\lambda(\epsilon)\bigr).
\end{equation}
Sa composition se calcule au moyen de \eqref{eq:continuo-comun-transporte-hojas-funtorial}, de \eqref{eq:continuo-comun-accion-memoria-composicion} et de la concaténation du registre.

Le sommet du tour nonadique est désormais l'ensemble des relèvements horizontaux compatibles avec la phase
\begin{equation}
\label{eq:continuo-comun-gamma9-apice}
 \begin{split}
 \mathsf Z^{\Gamma}_{9,k}:=\coprod_{r\in C_9}\bigl\{
 (x,\epsilon_1,\ldots,\epsilon_9):\;&x_0=x,\quad
 x_i=(\epsilon_i)_*x_{i-1},\\
 &(x_{i-1},\epsilon_i,x_i)
 \in\mathscr R_{r+i-1,k+i-1},\quad 1\leq i\leq9
 \bigr\},
 \end{split}
\end{equation}
 où les indices de \(r+i-1\) se lisent dans \(C_9\). Le terme est déterminé de manière unique par \(r=g_k(x)\), de sorte que le coproduit ne duplique pas les témoins. Les applications source et cible sont
\begin{equation}
\label{eq:continuo-comun-gamma9-span}
 \widehat{\mathcal X}^{\rm tot}_k
 \xleftarrow{\ s_{9,k}\ }
 \mathsf Z^{\Gamma}_{9,k}
 \xrightarrow{\ t_{9,k}\ }
 \widehat{\mathcal X}^{\rm tot}_{k+9},
 \qquad
 s_{9,k}(x,\boldsymbol\epsilon)=x,\quad
 t_{9,k}(x,\boldsymbol\epsilon)
 =(\epsilon_9\cdots\epsilon_1)_*x.
\end{equation}
L'holonomie relationnelle est
\begin{equation}
\label{eq:continuo-comun-gamma9-relacion}
 \Gamma_{9,k}:=
 \{(s_{9,k}z,t_{9,k}z):z\in\mathsf Z^{\Gamma}_{9,k}\}
 =\operatorname{Hol}_{C_9}(\nabla_{\rm TPK}).
\end{equation}

\begin{lemma}[Le niveau neuf de l'arbre réalise l'holonomie]
\label{lem:continuo-comun-gamma9-arbol}
L'application qui oublie les états intermédiaires déjà déterminés par les émissions induit une bijection
\begin{equation}
\label{eq:continuo-comun-gamma9-apice-arbol}
 \mathsf Z^{\Gamma}_{9,k}
 \xrightarrow{\ \cong\ }
 \coprod_{x\in\widehat{\mathcal X}^{\rm tot}_k}
 \operatorname{Term}_{k,k+9}(x).
\end{equation}
\end{lemma}

\begin{proof}
Chaque arête de \(\operatorname{Out}(x_i)\) satisfait \eqref{eq:continuo-comun-accion-fase-tpk} ; ainsi, un mot terminal de longueur neuf détermine les neuf éléments de \eqref{eq:continuo-comun-relacion-horizontal-fase}. Réciproquement, un relèvement horizontal de \eqref{eq:continuo-comun-gamma9-apice} est un mot composable de l'arbre. Les deux opérations conservent les émissions et sont inverses. Ainsi, \eqref{eq:continuo-comun-gamma9-apice-arbol} est un théorème dérivé de la connexion, non sa définition.
\end{proof}

\begin{proposition}[Retour de phase, progression de la mémoire et absence de réinitialisation]
\label{prop:continuo-comun-gamma9-retorno-memoria}
Pour tout \(z=(x,\epsilon_1,\ldots,\epsilon_9)\) du sommet,
\begin{align}
 g(t_{9,k}z)&=g(s_{9,k}z),
 \label{eq:continuo-comun-gamma9-retorno-fase}\\
 q^{(9)}(t_{9,k}z)&=q^{(9)}(s_{9,k}z)+1,
 \label{eq:continuo-comun-gamma9-avance-memoria}\\
 \operatorname{Led}(t_{9,k}z)
 &=\operatorname{Led}(s_{9,k}z)\star
   \bigl(\lambda(\epsilon_1),\ldots,\lambda(\epsilon_9)\bigr).
 \label{eq:continuo-comun-gamma9-ledger}
\end{align}
En particulier, l’identification de phase
\(g(t_{9,k}z)=g(s_{9,k}z)\) n’induit pas une identification des états : après
transport des deux registres vers le même type de comparaison, leurs mémoires
verticales diffèrent d’une unité. C’est la phase publiée qui revient, et non l’état
holonomique intégral.
\end{proposition}

\begin{proof}
Les deux premières identités sont \eqref{eq:continuo-comun-retorno-fase-avance-memoria}. La troisième est la définition de la concaténation du registre. Puisque la mémoire verticale diffère d'une unité, les états complets sont distincts, même si une coordonnée visible supplémentaire pouvait coïncider.
\end{proof}

\begin{proposition}[Compression isométrique de la même holonomie]
\label{prop:continuo-comun-gamma9-compresion-isometrica}
Soit \(U_{\Gamma_9}:\mathcal H_{\rm full}\to\mathcal H_{\rm full}\) une réalisation isométrique de l'holonomie précédente, et soit \(J:\mathcal H_{\rm vis}\to\mathcal H_{\rm full}\) l'inclusion isométrique du registre visible, avec \(E:=J^*\) et \(P:=JE\). Nous définissons
\begin{equation}
\label{eq:continuo-comun-gamma9-compresion-datos}
 T:=EU_{\Gamma_9}J,\qquad
 \eta:=(I-P)U_{\Gamma_9}J.
\end{equation}
Alors
\begin{equation}
\label{eq:continuo-comun-gamma9-compresion}
 \boxed{\;
 U_{\Gamma_9}J=JT+\eta,\qquad
 I-T^*T=\eta^*\eta.
 \;}
\end{equation}
\end{proposition}

\begin{proof}
En insérant \(I=P+(I-P)\) devant \(U_{\Gamma_9}J\), on obtient
\[
 U_{\Gamma_9}J
 =JEU_{\Gamma_9}J+(I-P)U_{\Gamma_9}J
 =JT+\eta.
\]
Les deux termes appartiennent à des sous-espaces orthogonaux. Puisque \(U_{\Gamma_9}\) et \(J\) sont des isométries,
\[
 I=(U_{\Gamma_9}J)^*(U_{\Gamma_9}J)
   =T^*T+\eta^*\eta,
\]
ce qui constitue la seconde identité.
\end{proof}

\begin{proposition}[Totalité et équivariance de l'holonomie]
\label{prop:continuo-comun-gamma9-total-natural}
L'application \(s_{9,k}\) est surjective. Les neuf prolongements neutres forment la section
\begin{equation}
\label{eq:continuo-comun-gamma9-seccion-neutra}
 \begin{gathered}
 x_0:=x,\qquad
 \varepsilon^0_{k+i}:=\varepsilon^0_{k+i}(x_i),\qquad
 x_{i+1}:=(\varepsilon^0_{k+i})_*x_i\quad(0\leq i<9),\\
 \eta^{(9)}_k(x):=
 (x,\varepsilon^0_k,\ldots,\varepsilon^0_{k+8})
 \in\mathsf Z^{\Gamma}_{9,k},
 \qquad s_{9,k}\eta^{(9)}_k=I.
 \end{gathered}
\end{equation}
Toute symétrie cohérente \(a\) dont l'action sur l'odomètre satisfait \(\mathsf Q\,a=\bar a\,\mathsf Q\) et \(\bar a\,\sigma_9=\sigma_9\,\bar a\) induit
\[
 a^Z(x,\epsilon_1,\ldots,\epsilon_9)
 :=(a x,a\epsilon_1,\ldots,a\epsilon_9)
\]
et vérifie
\begin{equation}
\label{eq:continuo-comun-gamma9-naturalidad}
 s_{9,k}a^Z=a_ks_{9,k},
 \qquad
 t_{9,k}a^Z=a_{k+9}t_{9,k}.
\end{equation}
\end{proposition}

\begin{proof}
Le mot neutre existe en toute source et satisfait l'action de phase \eqref{eq:continuo-comun-accion-fase-tpk} ; cela prouve la section et la surjectivité. La condition \(\bar a\sigma_9=\sigma_9\bar a\) envoie les relations horizontales sur des relations horizontales. L'action entrée par entrée conserve composition, retenue, mémoire, frontière et registre ; lire le premier ou le dernier état avant ou après l'action donne \eqref{eq:continuo-comun-gamma9-naturalidad}.
\end{proof}

\section{Survie totale et multisection complète}
\label{sec:continuo-comun-supervivencia}

\begin{definition}[Survie produite]
\label{def:continuo-comun-supervivencia}
Un état \(x\in\widehat{\mathcal X}^{\rm tot}_k\) survit lorsque
\begin{equation}
\label{eq:continuo-comun-supervivencia}
 \operatorname{Term}_{k,N}(x)\neq\varnothing
 \qquad\text{pour tout }N\geq k.
\end{equation}
Nous notons \(\operatorname{Surv}_k\) l'ensemble de ces états.
\end{definition}

\begin{theorem}[Le domaine total est survivant]
\label{thm:continuo-comun-total-superviviente}
À toute profondeur,
\begin{equation}
\label{eq:continuo-comun-surv-igual-total}
 \operatorname{Surv}_k=\widehat{\mathcal X}^{\rm tot}_k.
\end{equation}
En particulier, la survie ne définit pas un sous-ensemble supposé avant la construction de l'état.
\end{theorem}

\begin{proof}
L'inclusion \(\operatorname{Surv}_k\subseteq
\widehat{\mathcal X}^{\rm tot}_k\) appartient à la définition. L'inclusion opposée est exactement la non-vacuité à tout horizon démontrée dans \cref{prop:continuo-comun-arbol-finito-no-vacio}.
\end{proof}

\begin{definition}[Multisection des prolongements complets]
\label{def:continuo-comun-multiseccion}
La multisection complète au-dessus de \(x\) est la limite inverse
\begin{equation}
\label{eq:continuo-comun-multiseccion}
 \operatorname{Msect}(x)
 :=\varprojlim_{N\geq k}
 \bigl(\operatorname{Term}_{k,N}(x),\pi_{N+1,N}\bigr).
\end{equation}
Ses éléments sont toutes les histoires infinies compatibles qui prolongent \(x\). Aucune ne reçoit de priorité au sein de \eqref{eq:continuo-comun-multiseccion}.
\end{definition}

\begin{theorem}[Non-vacuité concrète de la multisection]
\label{thm:continuo-comun-multiseccion-no-vacia}
Pour tout \(x\in\widehat{\mathcal X}^{\rm tot}_k\),
\begin{equation}
\label{eq:continuo-comun-multiseccion-no-vacia}
 \operatorname{Msect}(x)\neq\varnothing.
\end{equation}
En outre, chaque coordonnée finie d'une histoire de \(\operatorname{Msect}(x)\) conserve les deux feuillets, l'orientation, la retenue, la mémoire, la frontière, le parcours et le registre de l'état qu'elle prolonge.
\end{theorem}

\begin{proof}
Les ensembles terminaux sont finis et non vides par \cref{prop:continuo-comun-arbol-finito-no-vacio}, et leurs applications de liaison sont surjectives par \cref{lem:continuo-comun-terminal-sobreyectivo}. On peut construire récursivement un élément compatible : on part de la racine et, à chaque pas, on conserve un prolongement du préfixe précédent. Le prolongement neutre garantit que le processus ne s'arrête jamais. La compatibilité des coordonnées découle du fait que chaque arête de l'arbre agit par l'unique action d'état définie dans \eqref{eq:continuo-comun-estado}.
\end{proof}

\section{Mesure canonique engendrée par l'arbre}
\label{sec:continuo-comun-medida}

La mesure n'est pas introduite comme un poids extérieur sur les sorties. Elle s'obtient en comptant les émissions admissibles à chaque sommet avant de choisir une continuation.

\begin{definition}[Poids local et mesure à horizon donné]
\label{def:continuo-comun-medida-finita}
Pour \(v\in\widehat{\mathcal X}^{\rm tot}_j\) et \(\epsilon\in\operatorname{Out}(v)\), nous posons
\begin{equation}
\label{eq:continuo-comun-peso-local}
 p_v(\epsilon):=\frac{1}{d(v)}\in\mathbb Q_{>0}.
\end{equation}
Si \(h=(\epsilon_{k+1},\ldots,\epsilon_N)\in
\operatorname{Term}_{k,N}(x)\) et si \(x_j\) est le préfixe précédant \(\epsilon_{j+1}\), nous définissons
\begin{equation}
\label{eq:continuo-comun-medida-horizonte}
 \mu_{x,N}(h)
 :=\prod_{j=k}^{N-1}p_{x_j}(\epsilon_{j+1})
 =\prod_{j=k}^{N-1}\frac{1}{d(x_j)}.
\end{equation}
\end{definition}

\begin{lemma}[Normalisation]
\label{lem:continuo-comun-medida-normalizada}
Pour tout horizon \(N\geq k\),
\begin{equation}
\label{eq:continuo-comun-medida-suma-uno}
 \sum_{h\in\operatorname{Term}_{k,N}(x)}\mu_{x,N}(h)=1.
\end{equation}
\end{lemma}

\begin{proof}
À la hauteur un, \(\sum_{\epsilon\in\operatorname{Out}(x)}d(x)^{-1}=1\). En supposant la normalisation à la hauteur \(N-k\), chaque préfixe terminal \(h\), d'état final \(v_h\), contribue au niveau suivant par
\[
 \sum_{\epsilon\in\operatorname{Out}(v_h)}
 \mu_{x,N}(h)p_{v_h}(\epsilon)
 =\mu_{x,N}(h).
\]
La sommation sur tous les préfixes prouve l'identité par récurrence.
\end{proof}

\begin{proposition}[Cohérence projective]
\label{prop:continuo-comun-medida-consistente}
Les mesures finies satisfont
\begin{equation}
\label{eq:continuo-comun-medida-pushforward}
 (\pi_{N+1,N})_*\mu_{x,N+1}=\mu_{x,N}.
\end{equation}
\end{proposition}

\begin{proof}
Pour \(h\in\operatorname{Term}_{k,N}(x)\), la masse de sa fibre est
\begin{align*}
 \sum_{\pi_{N+1,N}(h')=h}\mu_{x,N+1}(h')
 &=\mu_{x,N}(h)
   \sum_{\epsilon\in\operatorname{Out}(v_h)}p_{v_h}(\epsilon)\\
 &=\mu_{x,N}(h).
\end{align*}
C'est précisément l'égalité des mesures ponctuelles équivalente à \eqref{eq:continuo-comun-medida-pushforward}.
\end{proof}

Les cylindres de la multisection sont
\begin{equation}
\label{eq:continuo-comun-cilindro-msect}
 [h]_{x,N}
 :=\{\omega\in\operatorname{Msect}(x):
       \omega|_N=h\}.
\end{equation}

\begin{theorem}[Mesure canonique des prolongements]
\label{thm:continuo-comun-medida-canonica}
Il existe une unique mesure de cylindres \(\mu_x\) sur \(\operatorname{Msect}(x)\) telle que
\begin{equation}
\label{eq:continuo-comun-medida-cilindro}
 \mu_x([h]_{x,N})=\mu_{x,N}(h).
\end{equation}
Elle est normalisée, son support est la multisection entière, et elle s'obtient uniquement à partir des ensembles finis de prolongements APP--TRIT--TPK.
\end{theorem}

\begin{proof}
La cohérence de \cref{prop:continuo-comun-medida-consistente} rend \eqref{eq:continuo-comun-medida-cilindro} indépendante de l'horizon dans lequel un cylindre est représenté. La normalisation provient de \cref{lem:continuo-comun-medida-normalizada}. Deux cylindres sont soit disjoints, soit emboîtés ; les masses finies définissent donc de manière unique une mesure additive sur l'algèbre des cylindres. La multisection est une limite d'ensembles finis discrets et les cylindres constituent sa base compacte ; la mesure se prolonge de manière unique à la \(\sigma\)-algèbre qu'ils engendrent. Enfin, \(p_v(\epsilon)>0\) pour toute émission admissible, si bien que tout cylindre non vide possède une masse positive et que le support est la multisection complète.
\end{proof}

\section{Naturalité sous les symétries et les troncatures}
\label{sec:continuo-comun-naturalidad}

Un système cohérent \(g=(g_j)_{j\geq0}\in\mathsf U\) transporte une émission sortante de profondeur \(j\) au moyen de
\begin{equation}
\label{eq:continuo-comun-accion-out}
 (g_{j+1},g_j)_*:\operatorname{Out}(x)\longrightarrow
      \operatorname{Out}(g_j\cdot x),
 \qquad
 \epsilon\longmapsto g_{j+1}\epsilon g_j^{-1}.
\end{equation}
Les relations TPK font de \eqref{eq:continuo-comun-accion-out} une bijection préservant les multiplicités. L'action sur chaque entrée donne un isomorphisme d'arbres
\begin{equation}
\label{eq:continuo-comun-isomorfismo-arbol}
 g_{*,N}:\mathscr T_{k,N}(x)
 \xrightarrow{\ \cong\ }
 \mathscr T_{k,N}(g_k\cdot x).
\end{equation}

\begin{theorem}[Naturalité de l'arbre, de la multisection et de la mesure]
\label{thm:continuo-comun-naturalidad}
Pour toute symétrie réversible cohérente \(g=(g_j)_{j\geq0}\),
\begin{align}
 \pi_{N+1,N}g_{*,N+1}
 &=g_{*,N}\pi_{N+1,N},
 \label{eq:continuo-comun-naturalidad-terminal}\\
 g_*\operatorname{Msect}(x)
 &=\operatorname{Msect}(g_k\cdot x),
 \label{eq:continuo-comun-naturalidad-msect}\\
 (g_{*,N})_*\mu_{x,N}
 &=\mu_{g_k\cdot x,N},
 \qquad
 (g_*)_*\mu_x=\mu_{g_k\cdot x}.
 \label{eq:continuo-comun-naturalidad-medida}
\end{align}
\end{theorem}

\begin{proof}
La première égalité exprime que l'action entrée par entrée commute avec l'effacement de la dernière entrée. Le passage à la limite inverse donne \eqref{eq:continuo-comun-naturalidad-msect}. Puisque \eqref{eq:continuo-comun-accion-out} est bijective, \(d(g_j\cdot v)=d(v)\) ; par \eqref{eq:continuo-comun-medida-horizonte}, un mot et son image ont le même poids. Cela prouve la première égalité de mesures de \eqref{eq:continuo-comun-naturalidad-medida}. La seconde se déduit sur les cylindres au moyen de \eqref{eq:continuo-comun-medida-cilindro}, et les cylindres déterminent la mesure.
\end{proof}

Les troncatures de l'état et celles de l'arbre forment également le carré
\begin{equation}
\label{eq:continuo-comun-cuadrado-truncamiento}
\begin{CD}
 \operatorname{Term}_{k,N+1}(x)
 @>{\pi_{N+1,N}}>>
 \operatorname{Term}_{k,N}(x)\\
 @V{\operatorname{ter}_{N+1}}VV
 @VV{\operatorname{ter}_{N}}V\\
 \widehat{\mathcal X}^{\rm tot}_{N+1}
 @>{\rho_{N+1,N}}>>
 \widehat{\mathcal X}^{\rm tot}_{N}.
\end{CD}
\end{equation}
La commutativité est littérale : les deux parcours effacent la dernière émission et sa dernière entrée de registre.

\section{Limite projective \(\widehat\Omega\) des histoires complètes et état holonomique intégral commun APP–TRIT–TPK}
\label{sec:continuo-comun-limite}

Les surjections \eqref{eq:continuo-comun-truncamiento-total} produisent l'objet des histoires complètes
\begin{equation}
\label{eq:continuo-comun-omega}
 \widehat\Omega
 :=\varprojlim_k
 \bigl(\widehat{\mathcal X}^{\rm tot}_k,\rho_{k+1,k}\bigr).
\end{equation}
Un élément \(\widehat x=(\widehat x_k)_{k\geq0}\) contient, à chaque profondeur, la même paire de feuillets transportée, les orientations et retenues produites, la mémoire accumulée, les frontières, le parcours et le registre complet. Sa fibre de continuations depuis tout préfixe est \(\operatorname{Msect}(\widehat x_k)\), munie de \(\mu_{\widehat x_k}\).

\begin{theorem}[État holonomique intégral commun APP--TRIT--TPK]
\label{thm:continuo-comun-fuente teoremática}
Le système
\begin{equation}
\label{eq:continuo-comun-fuente teoremática}
 \mathfrak C^{\rm gen}:=\Biggl(
 \begin{aligned}
 &\{\widehat{\mathcal X}^{\rm tot}_k,\rho_{k+1,k},\eta_k\}_{k\geq0},
   \widehat\Omega,\{\mathfrak G_k\}_{k\geq0},
   \{\mathscr T_{k,N}\}_{N\geq k},\\
 &\operatorname{Msect},\operatorname{Surv},\{\mu_x\},
   \{\mathfrak r_9(v_K)\}_{K\geq0},
   \mathcal K_{\rm ph},\widehat{\mathcal K}_{\rm ph},\\
 &\operatorname{Sh},\operatorname{or},\kappa,
   \mathsf Q,g,q^{(9)},\beta,m,\partial,\gamma,\operatorname{Led},\\
 &\{\mathsf Z^\Gamma_{9,k},s_{9,k},t_{9,k},\Gamma_{9,k}\}_{k\geq0}
 \end{aligned}
 \Biggr)
\end{equation}
est défini pour tout état produit par APP--TRIT--TPK\@. Il est non vide, naturel sous les symétries réversibles et compatible avec le raffinement. En particulier :
\begin{enumerate}
 \item les deux feuillets sont transportés ensemble ;
 \item l'orientation et la retenue sont calculées par la division TRIT ;
 \item la mémoire obéit à l'action affine TPK ;
 \item la phase obéit à l'odomètre TPK, revient après neuf émissions et la mémoire verticale avance d'une unité ;
 \item la réalisation markovienne de ce tour satisfait \(\mathcal K_{\rm ph}^9=E_+E_\times\) et renouvelle uniformément la fibre complète de \(468=18\cdot26\) émissions ;
 \item le lecteur compact conserve bloc, cylindre, préfixe, retenue, signature, projection de Witt et phase, avec leurs actualisations et troncatures ;
 \item frontière, parcours et registre se composent sans effacer les extrémités ni l'ordre des émissions ;
 \item tout état survit et possède une multisection non vide de prolongements complets ;
 \item la mesure canonique est projective et possède un support sur tous ces prolongements.
\end{enumerate}
\end{theorem}

\begin{proof}
La totalité et la finitude découlent de \cref{prop:continuo-comun-total-finito}. Les identités des feuillets, de la retenue et de la mémoire sont démontrées dans \cref{lem:continuo-comun-asociatividad-carry,lem:continuo-comun-accion-memoria} ; la phase, la mémoire verticale et l'holonomie nonadique, y compris sa réalisation sur la fibre et son lecteur d'état, sont démontrées dans \cref{prop:continuo-comun-gamma9-retorno-memoria,prop:continuo-comun-gamma9-total-natural,prop:continuo-comun-kph-heredada,prop:continuo-comun-gamma9-compresion-isometrica,prop:continuo-comun-no-reinicio-alcance} ; la frontière conserve ses extrémités par \eqref{eq:continuo-comun-frontera-telescopica}. La survie totale et la non-vacuité des multisections sont \cref{thm:continuo-comun-total-superviviente,thm:continuo-comun-multiseccion-no-vacia}. L'existence, la cohérence et le support de la mesure sont \cref{prop:continuo-comun-medida-consistente,thm:continuo-comun-medida-canonica}. Enfin, la compatibilité avec les symétries et les troncatures est démontrée dans \cref{thm:continuo-comun-naturalidad} et \eqref{eq:continuo-comun-cuadrado-truncamiento}.
\end{proof}

\section{Nécessité structurelle et provenance}
\label{sec:continuo-comun-falsadores}

\begin{proposition}[Falsificateurs de l’état holonomique intégral commun]
\label{prop:continuo-comun-falsadores}
Chacune des opérations suivantes remplace \(\mathfrak C^{\rm gen}\) par un autre objet et bloque le transfert des théorèmes précédents :
\begin{enumerate}
 \item choisir \(\Sigma\) ou \(\Pi\) et écarter l'autre feuillet ;
 \item conserver le résidu et effacer le quotient ;
 \item conserver le signe TRIT et effacer sa retenue ;
 \item remplacer le registre ordonné par sa valeur terminale ;
 \item identifier deux mots parce qu'ils atteignent la même cellule ;
 \item choisir une continuation et la substituer à la multisection ;
 \item déclarer la survie sans exhiber des prolongements à chaque horizon ;
 \item normaliser uniformément le seul dernier niveau, en ignorant les degrés des préfixes et en perdant la cohérence projective ;
 \item utiliser une lecture ultérieure pour décider quels états entrent dans \(\widehat{\mathcal X}^{\rm tot}_k\).
\end{enumerate}
\end{proposition}

\begin{proof}
Les six premiers points suppriment des coordonnées ou des relations présentes dans la définition de l'état, de l'arbre ou de la multisection. Le septième supprime la preuve de \cref{thm:continuo-comun-total-superviviente}. Le huitième peut violer \eqref{eq:continuo-comun-medida-pushforward} lorsque les degrés terminaux ne sont pas constants. Le dernier remplace le domaine total \eqref{eq:continuo-comun-dominio-total} par une sélection rétrospective.
\end{proof}

Les cellules APP, la paire de feuillets, le relèvement TRIT, les extensions
TPK, le cocycle de retenue, la mémoire, la frontière, le chemin et la conservation
du registre constituent les structures de départ. Sur elles sont réunis dans
un domaine unique le groupoïde fini, l’arbre des mots, la multisection et
la mesure locale normalisée. Les théorèmes de survie totale, de non-vacuité,
de consistance et de naturalité certifient cette construction conjointe.
 \let\section\HMTContinuoCommonOriginalSection
\chapter{Opérations intrinsèques corrélatives du même continu}
\label{chap:pdfv2-operaciones-intrinsecas-correlativas}

\begin{thesisbox}
Les symboles \(\mathrm{sp},\mathrm{sol},\mathrm{gau},\mathrm{coh}\) et \(\mathrm{det}\) qui apparaissent dans ce chapitre marquent cinq opérations corrélatives sur un unique état APP--TRIT--TPK\@. Chaque opération reçoit la même extension survivante, la même orientation, la même retenue, la même mémoire, la même frontière et le même parcours. Les sections suivantes constituent des étapes mathématiques successives d'une seule construction, non des exposés autonomes.
\end{thesisbox}

\section{Arête génératrice et état généalogique commun}

Nous ne répétons ici ni l'état, ni l'arbre, ni la mesure construits dans le chapitre précédent. Nous prenons \(\widehat x=\widehat x_k(\gamma)\in
\widehat{\mathcal X}^{\rm tot}_k\) de \eqref{eq:continuo-comun-estado} et toutes ses émissions admissibles \(\alpha\in\operatorname{Out}(\widehat x)\). Si \(\widehat y=\alpha_*\widehat x\), nous écrivons
\begin{equation}
\label{eq:pdfv2-cor-hijos-supervivientes}
 \operatorname{Ch}_k(\widehat x):=
 \{(\alpha,\alpha_*\widehat x):
      \alpha\in\operatorname{Out}(\widehat x)\},
 \qquad
 d_k(\widehat x):=\#\operatorname{Ch}_k(\widehat x).
\end{equation}
La finitude et la non-vacuité de cet ensemble sont déjà \eqref{eq:continuo-comun-out-finito} ; la survie à tous les horizons, la multisection complète et sa mesure sont \cref{thm:continuo-comun-total-superviviente,thm:continuo-comun-multiseccion-no-vacia,thm:continuo-comun-medida-canonica}. Aucune seconde non-vacuité n'est supposée ni démontrée à nouveau.

Chaque paire \((\alpha,\widehat y)\) de \eqref{eq:pdfv2-cor-hijos-supervivientes} conserve le registre unique
\begin{equation}
\label{eq:pdfv2-cor-registro-extension}
 \begin{aligned}
 \mathfrak e_\alpha:=\bigl(&
 \alpha,\operatorname{Ext}_{\gamma_\alpha};
 \operatorname{Sh}_{s(\alpha)},\operatorname{Sh}_{t(\alpha)};\\
 &
 \operatorname{res}^{\Sigma,\Pi}(\alpha),
 \operatorname{quo}^{\Sigma,\Pi}(\alpha);\\
 &\operatorname{or}(\alpha),\kappa(\alpha),\mathsf C_\alpha;
 \mathsf Q(s(\alpha)),\mathsf Q(t(\alpha)),g(s(\alpha)),g(t(\alpha)),
 q^{(9)}(s(\alpha)),q^{(9)}(t(\alpha));\\
 &
 m(\gamma\alpha;m_0),\partial\alpha,\gamma\alpha,
 \operatorname{Led}(\gamma\alpha);\\
 &\{\operatorname{Term}_{k+1,N}(\widehat y)\}_{N\geq k+1},
 \operatorname{Msect}(\widehat y),\mu_{\widehat y}\bigr).
 \end{aligned}
\end{equation}
Ici, \(\operatorname{Ext}_{\gamma_\alpha}\) désigne la relation complète déjà satisfaite par l'émission admissible ; elle n'introduit pas un autre filtre. Les opérations ultérieures sont des applications coordonnées de \(\mathfrak e_\alpha\) et ne reconstruisent pas le registre à partir de leurs sorties.

\section{Actions intrinsèques corrélatives sur le générateur commun}

Les marques internes sont fixées sans changer de source : \(\mathrm{sp}\) désigne le transport et la graduation des deux feuillets ; \(\mathrm{sol}\), l'accroissement signé, le bilan et la mémoire ; \(\mathrm{gau}\), l'incidence et la courbure de retenue ; \(\mathrm{coh}\), les ports, la différence rectangulaire et le complexe cellulaire ; et \(\mathrm{det}\), le complexe muni de bases, le remplissage corrélatif et la droite déterminant. Les formules suivantes montrent leurs couplages sur le même \(\mathfrak e_\alpha\) ; leurs actions partagent littéralement source, cible, parcours et registre.

La double évaluation APP fournit, sur le même support, les étiquettes \(\Sigma\) et \(\Pi\). Pour une extension \(\alpha\), leurs multiplicités se calculent à partir du registre complet, sans compteur extérieur :
\begin{equation}
\label{eq:pdfv2-cor-multiplicidades-hojas}
 N_\diamond(\alpha):=
 \sum_{\epsilon\ {\rm en}\ \gamma\alpha}
 \mathbf1_{\{\operatorname{res}^{\diamond}(\epsilon)\neq0\}},
 \qquad \diamond\in\{\Sigma,\Pi\}.
\end{equation}
Soit en outre le coefficient orienté total
\[
 \Delta_{\gamma_\alpha}:=
 \sum_{\epsilon\ {\rm en}\ \gamma\alpha}\operatorname{or}(\epsilon)
 \in\mathbb Z.
\]
L'action signée est
\begin{equation}
\label{eq:pdfv2-cor-b-alpha}
 b_\alpha:=\Delta_{\gamma_\alpha}
 \bigl(N_\Sigma(\alpha)-N_\Pi(\alpha)\bigr),
 \qquad
 b_\alpha^\pm:=\frac{|b_\alpha|\pm b_\alpha}{2}.
\end{equation}
Par construction,
\begin{equation}
\label{eq:pdfv2-cor-b-identidades}
 b_\alpha=b_\alpha^+-b_\alpha^-,\qquad
 |b_\alpha|=b_\alpha^++b_\alpha^-,\qquad
 b_\alpha^+b_\alpha^-=0.
\end{equation}

Le même ensemble d'extensions produit simultanément les ports
\begin{equation}
\label{eq:pdfv2-cor-puertos-comunes}
 I_k(\widehat x_k):=
 \operatorname{Ch}_k(\widehat x_k)\times\{\Sigma\},
 \qquad
 J_k(\widehat x_k):=
 \operatorname{Ch}_k(\widehat x_k)\times\{\Pi\}.
\end{equation}
On ne suppose pas qu'un feuillet visible contienne davantage de prolongements que l'autre : chaque extension survivante est évaluée dans les deux feuillets du même support. Par \eqref{eq:continuo-comun-out-finito}, les deux ensembles de ports sont non vides.

Pour \(A_N=\mathbb Z/3^N\mathbb Z\), l'action cellulaire sur les générateurs est
\begin{equation}
\label{eq:pdfv2-cor-kappa-celda}
 \begin{aligned}
 \kappa_{k,N}:A_N^{I_k}\oplus A_N^{J_k}
 &\longrightarrow A_N^{I_k\times J_k},\\
 \kappa_{k,N}(u,v)_{(\alpha,\Sigma),(\beta,\Pi)}
 &:=u_{(\alpha,\Sigma)}-v_{(\beta,\Pi)}.
 \end{aligned}
\end{equation}
Après choix des premiers ports suivant l'ordre TPK, la différence rectangulaire
\begin{equation}
\label{eq:pdfv2-cor-delta-celda}
 \delta(w)_{ij}:=w_{ij}-w_{ij_0}-w_{i_0j}+w_{i_0j_0}
\end{equation}
réalise la suite scindée
\begin{equation}
\label{eq:pdfv2-cor-sucesion-celda}
 0\longrightarrow A_N\xrightarrow{\Delta}
 A_N^{I_k}\oplus A_N^{J_k}
 \xrightarrow{\kappa_{k,N}}A_N^{I_k\times J_k}
 \xrightarrow{\delta}A_N^{(|I_k|-1)(|J_k|-1)}
 \longrightarrow0.
\end{equation}
C'est une identité de la même fibre APP : elle n'ajoute aucune cellule étrangère au générateur \(\mathfrak e_\alpha\).

L'action cellulaire locale utilise les deux coordonnées entières accumulées par ce même registre. Pour \(\widehat x=\widehat x_k(\gamma)\) et \(\widehat y=\alpha_*\widehat x=\widehat x_{k+1}(\gamma\alpha)\), nous posons
\begin{equation}
\label{eq:pdfv2-cor-c-v}
 \begin{aligned}
 c_{\widehat x}&:=\kappa(\gamma), &
 c_\alpha:=c_{\widehat y}&:=\kappa(\gamma\alpha),\\
 v_{\widehat x}&:=\sum_{\epsilon\ {\rm en}\ \gamma}
                  \operatorname{or}(\epsilon), &
 v_{\widehat y}&:=\sum_{\epsilon\ {\rm en}\ \gamma\alpha}
                  \operatorname{or}(\epsilon),\\
 v_\alpha&:=v_{\widehat y}-v_{\widehat x}
           =\operatorname{or}(\alpha).&&
 \end{aligned}
\end{equation}
Sur les générateurs \(f,r,e,b,g\), on définit
\begin{equation}
\label{eq:pdfv2-cor-complejo-local}
 \partial f=3c_\alpha e+b,\qquad
 \partial r=0,\qquad
 \partial e=g,\qquad
 \partial b=-3c_\alpha g,\qquad
 \partial g=0.
\end{equation}
L'annulation
\begin{equation}
\label{eq:pdfv2-cor-dos-cero}
 \partial^2f=3c_\alpha g-3c_\alpha g=0
\end{equation}
est littérale et utilise la même retenue avec les deux signes imposés par l'orientation.

Le germe d'incidence s'obtient à partir du module libre des deux feuillets,
\begin{equation}
\label{eq:pdfv2-cor-v-tpk}
 V_{\rm TPK}:=\mathbf F_3e_\Sigma\oplus\mathbf F_3e_\Pi,
 \qquad
 \mathscr I:=\mathbf P^1(V_{\rm TPK}),
 \qquad |\mathscr I|=4,
\end{equation}
et de ses paires non ordonnées
\begin{equation}
\label{eq:pdfv2-cor-incidencias}
 \mathscr E:=\binom{\mathscr I}{2},
 \qquad |\mathscr E|=6,
 \qquad
 \mathscr B:=\mathscr I\times\{+,-\},
 \qquad |\mathscr B|=8.
\end{equation}
L'application d'incidence est déterminée sur la base par
\begin{equation}
\label{eq:pdfv2-cor-b-incidencia-generadores}
 B[e_L]=\sum_{L'\neq L}e_{\{L,L'\}},
 \qquad L\in\mathscr I,
\end{equation}
et, par linéarité,
\begin{equation}
\label{eq:pdfv2-cor-b-incidencia-coordenadas}
 (Ba)_{\{L,L'\}}=a_L+a_{L'}.
\end{equation}
Si \(s(q)=\sum_{E\in\mathscr E}q_E\), chaque \(L\) apparaît trois fois et
\begin{equation}
\label{eq:pdfv2-cor-sb-cero}
 s(Ba)=0.
\end{equation}
Par conséquent, la fonction
\begin{equation}
\label{eq:pdfv2-cor-wc}
 W_c(q):=\mathbf1_{\{s(q)=0\}}-\frac13
\end{equation}
satisfait \(W_c(q+Ba)=W_c(q)\). Les identités \eqref{eq:pdfv2-cor-v-tpk}--\eqref{eq:pdfv2-cor-wc} sont internes et exactes. Le relèvement concret de \(q\) à partir des cellules et retenues communes est traité dans la couche d'assemblage ; cette identité d'incidence ne se confond pas avec la construction de ce relèvement.

\section{Transport unitaire, polarisation et bilan sur la fibre commune}

La non-vacuité de \eqref{eq:pdfv2-cor-hijos-supervivientes} permet de fixer, sans sélectionner un prolongement,
\begin{equation}
\label{eq:pdfv2-cor-peso-conteo}
 p_k(\alpha,\widehat y\mid\widehat x)
 :=\frac1{d_k(\widehat x)}.
\end{equation}
Soit l'espace des deux feuillets du même état
\begin{equation}
\label{eq:pdfv2-cor-hilbert-hojas}
 \mathscr H_k:=
 \bigoplus_{\widehat x\in\widehat{\mathcal X}^{\rm tot}_k}
 \bigl(\mathsf H_{\widehat x}^{\Sigma}
       \oplus\mathsf H_{\widehat x}^{\Pi}\bigr),
\end{equation}
et soit \(\iota_{\widehat x}\) l'inclusion de la fibre indiquée. Le transport TPK de \eqref{eq:continuo-comun-transporte-hojas} donne l'action exacte
\begin{equation}
\label{eq:pdfv2-cor-u-comun}
 U_k\iota_{\widehat x}v:=
 \sum_{(\alpha,\widehat y)\in\operatorname{Ch}_k(\widehat x)}
 p_k(\alpha,\widehat y\mid\widehat x)^{1/2}
 \iota_{\widehat y}U_\alpha v.
\end{equation}
Les successeurs de deux prédécesseurs distincts sont disjoints puisque la troncature est une fonction. En outre, deux témoins \(\alpha\) distincts ne se confondent pas dans le même facteur direct : le parcours complet \(\gamma\), y compris son dernier pas, est une coordonnée de \(\widehat y\). En conséquence, l'identité de Gram est
\begin{equation}
\label{eq:pdfv2-cor-u-gram}
 \left.U_k^*U_k\right|_{\mathsf H_{\widehat x}^{\Sigma}
       \oplus\mathsf H_{\widehat x}^{\Pi}}
 =
 \sum_{(\alpha,\widehat y)\in\operatorname{Ch}_k(\widehat x)}
 p_k(\alpha,\widehat y\mid\widehat x)\,U_\alpha^*U_\alpha.
\end{equation}
Donc \(U_k^*U_k=I\) exactement lorsque les transports de feuillet
\(U_\alpha\) publiés par TPK sont des isométries. La normalisation des poids,
à elle seule, ne peut remplacer cette vérification.
Pour \(\ell>k\), l’histoire complète est conservée au moyen de
\begin{equation}
\label{eq:pdfv2-cor-u-iterado}
 U_{\ell,k}:=U_{\ell-1}\cdots U_k,
 \qquad U_{k,k}:=I.
\end{equation}

La graduation agit sur les deux coordonnées déjà présentes dans chaque fibre de \eqref{eq:pdfv2-cor-hilbert-hojas} :
\begin{equation}
\label{eq:pdfv2-cor-sigma}
 \sigma_k\iota_{\widehat x}(v_\Sigma,v_\Pi)
 :=\iota_{\widehat x}(v_\Sigma,-v_\Pi),
\qquad
 P_k:=\frac{I+\sigma_k}{2},\quad
 Q_k:=\frac{I-\sigma_k}{2}.
\end{equation}
Le seuil TRIT appartient à l'orientation et à la retenue de \(\widehat x\) ; il n'efface aucun des deux feuillets APP et ne nécessite pas de troisième cas dans \eqref{eq:pdfv2-cor-sigma}.

La partie qui conserve le feuillet et celle qui le change sont deux composantes du même \(U_k\) :
\begin{align}
 A_k&:=P_{k+1}U_kP_k+Q_{k+1}U_kQ_k,
 \label{eq:pdfv2-cor-a-comun}\\
 \eta_k&:=Q_{k+1}U_kP_k+P_{k+1}U_kQ_k.
 \label{eq:pdfv2-cor-eta-comun}
\end{align}
Par orthogonalité de \(P_{k+1}\) et de \(Q_{k+1}\),
\begin{equation}
\label{eq:pdfv2-cor-balance-hojas}
 U_k=A_k+\eta_k,\qquad
 A_k^*A_k+\eta_k^*\eta_k=U_k^*U_k.
\end{equation}
Le dernier membre est \(I\) sous, et uniquement sous, la condition isométrique
identifiée après \eqref{eq:pdfv2-cor-u-gram}.
Cette identité et l’incrément \eqref{eq:pdfv2-cor-b-alpha} ne sont pas deux
transports : \(A_k,\eta_k\) agissent sur la base des histoires et
\(b_\alpha\) lit l’arête de cette même base.

\subsection{Polarisation signée, télescopie non autonome et retour sans réinitialisation}
On n'introduit pas ici de branche spectrale indépendante. Les notations suivantes sont restreintes au même espace \(\mathscr H_k\) et au même transport \(U_k\) déjà construits
\[
 \mathcal H_k^{\mathrm{sp}}:=\mathscr H_k,
 \qquad
 \Lambda(\gamma_{\widehat x})
 :=\bigl(\operatorname{or},\operatorname{Carr},\operatorname{Mem},
          \partial,\gamma\bigr)(\widehat x),
\]
et \(\Sigma_{\rm reg}\Lambda(\gamma_{\widehat x})\) désigne sa famille de
troncatures compatibles. De cette manière, le lecteur signé, le défaut et le
terminal suivants sont des actions du constructeur corrélatif commun, et non une
sélection rétrospective d’histoires.
\begingroup
  \let\section\subsection
  % Extracción quirúrgica del fuente teoremática V2 05a: líneas 293--394.
% El resto permanece en este archivo como procedencia, pero no se publica.
\iffalse
\chapter{Macrostructure spectrale--polaire interne du continu}
\label{chap:continuo-macro-sp-interna}

\section{Portée strictement interne}

La construction de ce chapitre se termine avant toute reconnaissance
externe. Ses seules données sont le support APP à deux feuillets, l'orientation
TRIT, les extensions et la retenue TPK, le système projectif d'histoires, la
mesure cylindrique et le ledger signé. Le résultat est une macrostructure
graduée avec transport, expression de changement de feuillet, lecteur signé et
télescopage non autonome. Aucune fonction extérieure n'est utilisée pour définir
l'une quelconque de ses applications.

\section{Domaine survivant de la branche}

Pour \(x_k\in\mathcal A_k\), soit
\begin{equation}
\label{eq:continuo-sp-msect-n}
 \operatorname{Msect}^{\mathrm{sp},(N)}_k(x_k)
 =\left\{(x_j)_{j=k}^{N}:
 \tau_{j+1,j}x_{j+1}=x_j,
 \ (x_j,x_{j+9})\in\Gamma_{9,j},
 \ \text{et le faisceau est conservé }\{\Sigma,\Pi\}\right\}.
\end{equation}
On définit
\begin{equation}
\label{eq:continuo-sp-dominio-superviviente}
 \mathcal C_k^{\mathrm{sp}}
 =\left\{x_k:
 \operatorname{Msect}^{\mathrm{sp},(N)}_k(x_k)\neq\varnothing
 \text{ pour tout }N\geq k\right\},
 \qquad
 \mathcal C_{\rm cont}^{\mathrm{sp}}
 =\varprojlim_k\mathcal C_k^{\mathrm{sp}}.
\end{equation}
Le seuil TRIT n'est pas exclu de ce domaine. Le faisceau conserve les deux feuillets
arithmétiques même si une carte visible n'en publie aucun lors d'un
événement concret.

\section{Paquet de treize coordonnées}

Pour \(x_k\in\mathcal C_k^{\mathrm{sp}}\), le paquet dépendant est
\begin{equation}
\label{eq:continuo-sp-d13}
\begin{aligned}
 D_{\mathrm{sp},k}(x_k)=\bigl(&
 \mathcal F_{\mathrm{sp},k}(x_k),
 \operatorname{Ext}_{\Gamma,\mathrm{sp},k}(x_k),
 \operatorname{Rel}_{\mathrm{sp},k}(x_k),
 \mathcal N_{\mathrm{sp},k}(x_k),
 \operatorname{or}_{\mathrm{sp},k}(x_k),\\
 &\operatorname{Carr}_{\mathrm{sp},k}(x_k),
 \operatorname{Mem}_{\mathrm{sp},k}(x_k),
 \partial_{\mathrm{sp},k}(x_k),
 \gamma_{\mathrm{sp},k}(x_k),
 \operatorname{Msect}_{\mathrm{sp},k}(x_k),\\
 &\operatorname{Surv}_{\mathrm{sp},k}(x_k),
 \operatorname{Term}_{\mathrm{sp},k}(x_k),
 \mathcal L^{\rm sgn}_{\mathrm{sp},k}(x_k)\bigr).
\end{aligned}
\end{equation}
Ses composantes sont déterminées corrélativement comme suit.

La fibre \(\mathcal F_{\mathrm{sp},k}\) contient l'état TPK et le faisceau
ordonné \(\{\Sigma,\Pi\}\). L'extension est la restriction complète de
\(\operatorname{Ext}_\Gamma\) aux histoires qui satisfont
\eqref{eq:continuo-sp-dominio-superviviente}. Les relations contiennent les
lois d'identité et de composition, le relèvement résidu--quotient, le cocycle de
retenue, la compatibilité de frontière et la naturalité de la troncature. La
coordonnée numérique contient profondeur, phase, résidus, quotients et signature
interne. L'orientation stocke séparément feuillet et orientation TRIT. La
mémoire contient le quotient, le ledger et l'incrément de tour. La
frontière et le chemin restent non contractés. La multisection est le système
complet des prolongements de \eqref{eq:continuo-sp-msect-n}. La
survie est l'appartenance stable de
\eqref{eq:continuo-sp-dominio-superviviente}. Le terminal et le lecteur sont
construits dans les sections suivantes.

Les troncatures de toutes les coordonnées forment une application
\(u^{\mathrm{sp}}_{\ell,k}\) qui satisfait
\begin{equation}
\label{eq:continuo-sp-d-naturalidad}
 u^{\mathrm{sp}}_{\ell,k}D_{\mathrm{sp},\ell}
 =D_{\mathrm{sp},k}\tau_{\ell,k}.
\end{equation}

\section{Graphe de la restriction et section dépendante}

On définit
\begin{equation}
\label{eq:continuo-sp-g-graph}
 \mathcal G_{\mathrm{sp},k}
 =\operatorname{Graph}(D_{\mathrm{sp},k})
 =\{(x_k,D_{\mathrm{sp},k}x_k):
     x_k\in\mathcal C_k^{\mathrm{sp}}\}.
\end{equation}
La restriction totale est
\begin{equation}
\label{eq:continuo-sp-res-graph}
 \operatorname{Res}_{\mathrm{sp}}=s_{\mathrm{sp}}:
 \mathcal C_{\rm cont}^{\mathrm{sp}}
 \longrightarrow\mathcal G_{\mathrm{sp}},
 \qquad
 s_{\mathrm{sp}}(x)=(x,D_{\mathrm{sp}}x).
\end{equation}
La projection \(\pi_1(x,d)=x\) vérifie
\begin{equation}
\label{eq:continuo-sp-res-inversa}
 \pi_1s_{\mathrm{sp}}=I,
 \qquad
 s_{\mathrm{sp}}\pi_1=I_{\mathcal G_{\mathrm{sp}}}.
\end{equation}

Soit \(r_{\mathrm{sp},k}(x_k)\) la multisection des feuillets et des chemins qui appartient
déjà à \(D_{\mathrm{sp},k}(x_k)\). On définit
\begin{equation}
\label{eq:continuo-sp-x-dependiente}
 X_{\chi,\mathrm{sp},k}
 =\{(r_{\mathrm{sp},k}x_k,x_k,D_{\mathrm{sp},k}x_k):
     x_k\in\mathcal C_k^{\mathrm{sp}}\},
\end{equation}
\begin{equation}
\label{eq:continuo-sp-pr-x}
 \operatorname{pr}_{X,\mathrm{sp}}(x,D_{\mathrm{sp}}x)
 =(r_{\mathrm{sp}}x,x,D_{\mathrm{sp}}x).
\end{equation}
Oublier la première coordonnée de
\eqref{eq:continuo-sp-x-dependiente} est l'inverse de
\(\operatorname{pr}_{X,\mathrm{sp}}\).

\section{Faisceau de feuillets et mesure projective}

Soit
\begin{equation}
\label{eq:continuo-sp-haz-hojas}
 \widetilde{\mathcal C}_k^{\mathrm{sp}}
 =\{(x_k,h):x_k\in\mathcal C_k^{\mathrm{sp}},
              h\in\{\Sigma,\Pi\}\}.
\end{equation}
La mesure projective du continu induit des probabilités
\(\widetilde\mu_k\) sur ce faisceau. On se restreint au support positif ; pour
\(z\) dans ce support,
\begin{equation}
\label{eq:continuo-sp-consistencia-medida}
 \widetilde\mu_k(z)
 =\sum_{\widetilde\tau_{k+1,k}(z')=z}
  \widetilde\mu_{k+1}(z').
\end{equation}
La linéarisation de niveau est
\begin{equation}
\label{eq:continuo-sp-hilbert-nivel}
 \mathcal H_k^{\mathrm{sp}}
 =\ell^2(\operatorname{supp}\widetilde\mu_k),
\end{equation}
de base orthonormale \(e_z\).

\section{Transport de toutes les extensions et retenue}

Pour \(z'\) qui prolonge \(z\), nous définissons le poids conditionnel positif
\begin{equation}
\label{eq:continuo-sp-peso-condicional}
 w_k(z'\mid z)
 =\left(\frac{\widetilde\mu_{k+1}(z')}
               {\widetilde\mu_k(z)}\right)^{1/2}.
\end{equation}
Le transport est
\begin{equation}
\label{eq:continuo-sp-u-accion}
 U_ke_z
 =\sum_{\substack{z':\widetilde\tau_{k+1,k}(z')=z\\
            (z,z')\in\operatorname{Ext}_{\Gamma,\mathrm{sp},k}}}
 w_k(z'\mid z)e_{z'}.
\end{equation}
La somme contient toutes les extensions admissibles. Elle ne dépend pas d'un sélecteur.

Par \eqref{eq:continuo-sp-consistencia-medida},
\[
 \sum_{z':\widetilde\tau(z')=z}|w_k(z'\mid z)|^2=1.
\]
Les ensembles de descendants de deux parents distincts sont disjoints. Par
conséquent,
\begin{equation}
\label{eq:continuo-sp-u-isometria}
 U_k^*U_k=I_{\mathcal H_k^{\mathrm{sp}}}.
\end{equation}

Chaque base \(e_{z'}\) inclut la retenue de l'extension correspondante. Pour
des chemins composables, on conserve
\begin{align}
 \kappa(\gamma_2\gamma_1)
 &=\kappa(\gamma_2)+
   \mathsf C(\gamma_2)\kappa(\gamma_1),
 \label{eq:continuo-sp-carry-nativo}\\
 \operatorname{Carr}(\gamma_2\gamma_1)
 &=\operatorname{Carr}(\gamma_2)H(\gamma_1)
  +Y(\gamma_2)\operatorname{Carr}(\gamma_1).
 \label{eq:continuo-sp-carry-operador}
\end{align}
Les poids conditionnels se multiplient le long de l'unique chaîne de
troncatures d'un descendant. Les équations
\eqref{eq:continuo-sp-carry-nativo}--
\eqref{eq:continuo-sp-carry-operador} garantissent que les étiquettes de retenue
coïncident. Par conséquent,
\begin{equation}
\label{eq:continuo-sp-u-composicion}
 U_{\ell,k}=U_{\ell-1}\cdots U_k.
\end{equation}

\section{Graduation, projecteurs et décomposition du transport}

La graduation de feuillet est l'involution
\begin{equation}
\label{eq:continuo-sp-sigma}
 \sigma_ke_{(x,\Sigma)}=+e_{(x,\Sigma)},
 \qquad
 \sigma_ke_{(x,\Pi)}=-e_{(x,\Pi)}.
\end{equation}
Ainsi,
\[
 \sigma_k^*=\sigma_k,
 \qquad
 \sigma_k^2=I.
\]
Les projecteurs internes sont
\begin{equation}
\label{eq:continuo-sp-p-q}
 P_k=\frac{I+\sigma_k}{2},
 \qquad
 Q_k=\frac{I-\sigma_k}{2}.
\end{equation}
Elles satisfont
\begin{equation}
\label{eq:continuo-sp-p-q-relaciones}
 P_k^2=P_k,
 \quad Q_k^2=Q_k,
 \quad P_kQ_k=0,
 \quad P_k+Q_k=I.
\end{equation}

Nous séparons la partie du transport qui conserve le feuillet de celle qui le change :
\begin{align}
 A_k&=P_{k+1}U_kP_k+Q_{k+1}U_kQ_k,
 \label{eq:continuo-sp-a}\\
 \eta_k&=Q_{k+1}U_kP_k+P_{k+1}U_kQ_k.
 \label{eq:continuo-sp-eta}
\end{align}
Alors
\begin{equation}
\label{eq:continuo-sp-u-a-eta}
 U_k=A_k+\eta_k,
\end{equation}
et la graduation vérifie
\begin{equation}
\label{eq:continuo-sp-paridad-a-eta}
 \sigma_{k+1}A_k=A_k\sigma_k,
 \qquad
 \sigma_{k+1}\eta_k=-\eta_k\sigma_k.
\end{equation}

\begin{proposition}[Bilan local des feuillets]
\label{prop:continuo-sp-balance-local}
Pour chaque niveau,
\begin{equation}
\label{eq:continuo-sp-balance-local}
 A_k^*A_k+\eta_k^*\eta_k=I.
\end{equation}
\end{proposition}

\begin{proof}
L'orthogonalité de \(P_{k+1}\) et \(Q_{k+1}\) donne
\begin{align*}
 A_k^*A_k
 &=P_kU_k^*P_{k+1}U_kP_k
   +Q_kU_k^*Q_{k+1}U_kQ_k,\\
 \eta_k^*\eta_k
 &=P_kU_k^*Q_{k+1}U_kP_k
   +Q_kU_k^*P_{k+1}U_kQ_k.
\end{align*}
En sommant et en utilisant \(P_{k+1}+Q_{k+1}=I\),
\[
 A_k^*A_k+\eta_k^*\eta_k
 =P_kU_k^*U_kP_k+Q_kU_k^*U_kQ_k
 =P_k+Q_k=I.
\]
\end{proof}

L'opérateur \(\eta_k\) enregistre l'abandon du canal qui conserve le feuillet. Il
n'est pas la retenue entière : les coordonnées
\eqref{eq:continuo-sp-carry-nativo}--
\eqref{eq:continuo-sp-carry-operador} restent dans le paquet
\eqref{eq:continuo-sp-d13}.

\fi
\section{Lecteur signé avec registre généalogique}

Pour \(v\in\mathcal H_k^{\mathrm{sp}}\), la polarisation signée est
\begin{equation}
\label{eq:continuo-sp-polarizacion-firmada}
 \operatorname{pol}_k(v)
 =\langle v,\sigma_kv\rangle
 =\|P_kv\|^2-\|Q_kv\|^2.
\end{equation}
Le lecteur spécialisé conserve en outre la source :
\begin{equation}
\label{eq:continuo-sp-lector-firmado}
 \mathcal L^{\rm sgn}_{\mathrm{sp},k}(x_k;v)
 =\bigl(
   \Lambda(\gamma_{x_k}),
   \Sigma_{\rm reg}(\Lambda(\gamma_{x_k})),
   \|P_kv\|^2,
   \|Q_kv\|^2,
   \operatorname{pol}_k(v)\bigr).
\end{equation}
L'orientation TRIT est conservée dans \(\Lambda(\gamma_{x_k})\) ; elle ne
s'identifie pas au signe de feuillet. À l'horizon \(N\), le registre garde la
famille des lecteurs pour tous les niveaux \(k\leq N\). Tronquer signifie
supprimer uniquement les entrées ultérieures.

\section{Produit non autonome et terminal}

On définit
\begin{equation}
\label{eq:continuo-sp-f-producto}
 F_{0,0}=I_{\mathcal H_0^{\mathrm{sp}}},
 \qquad
 F_{0,k}=A_{k-1}\cdots A_1A_0
 \quad(k\geq1).
\end{equation}
Par conséquent,
\begin{equation}
\label{eq:continuo-sp-f-recursion}
 F_{0,k+1}=A_kF_{0,k}.
\end{equation}
Le terme publié lors de l'abandon du canal à l'étape \(k\) et le terminal à
l'horizon \(N\) sont
\begin{equation}
\label{eq:continuo-sp-defecto-terminal}
 \mathcal D_k=F_{0,k}^*\eta_k^*\eta_kF_{0,k},
 \qquad
 \mathcal T_N=F_{0,N}^*F_{0,N}.
\end{equation}

\begin{theorem}[Télescopage spectral--polaire non autonome]
\label{thm:continuo-sp-telescopia-no-autonoma}
Pour tout horizon fini \(N\),
\begin{equation}
\label{eq:continuo-sp-telescopia-no-autonoma}
 \boxed{
 I=\sum_{k=0}^{N-1}
 F_{0,k}^*\eta_k^*\eta_kF_{0,k}
 +F_{0,N}^*F_{0,N}.}
\end{equation}
De plus,
\[
 0\preceq\mathcal T_{N+1}
 \preceq\mathcal T_N\preceq I,
\]
de sorte qu'il existe la limite forte positive
\[
 \mathcal T_\infty
 =\operatorname{s-lim}_{N\to\infty}\mathcal T_N
\]
et
\begin{equation}
\label{eq:continuo-sp-telescopia-infinita}
 I=\sum_{k=0}^{\infty}
 F_{0,k}^*\eta_k^*\eta_kF_{0,k}
 +\mathcal T_\infty
\end{equation}
en topologie forte. Le théorème n'affirme pas que
\(\mathcal T_\infty\) soit nul.
\end{theorem}

\begin{proof}
En multipliant \eqref{eq:pdfv2-cor-balance-hojas} par
\(F_{0,k}^*\) et \(F_{0,k}\), et en utilisant
\eqref{eq:continuo-sp-f-recursion}, on obtient
\[
 F_{0,k}^*F_{0,k}
 =F_{0,k+1}^*F_{0,k+1}
  +F_{0,k}^*\eta_k^*\eta_kF_{0,k}.
\]
La somme de \(k=0\) à \(N-1\) annule tous les termes
intermédiaires et laisse précisément
\eqref{eq:continuo-sp-telescopia-no-autonoma}. Comme
\(A_k^*A_k\preceq I\),
\[
 \mathcal T_{N+1}
 =F_{0,N}^*A_N^*A_NF_{0,N}
 \preceq F_{0,N}^*F_{0,N}=\mathcal T_N.
\]
La monotonie et la borne positive fournissent la limite forte. Le passage à la
limite dans l'identité finie donne
\eqref{eq:continuo-sp-telescopia-infinita}.
\end{proof}

\iffalse
\section{Assembleur comme graphe}

Pour
\[
 z_N=(r_{\mathrm{sp},N}x_N,x_N,D_{\mathrm{sp},N}x_N)
 \in X_{\chi,\mathrm{sp},N},
\]
on définit l'assembleur matériel
\begin{equation}
\label{eq:continuo-sp-a-accion}
\begin{aligned}
 a_{\mathrm{sp},N}(z_N)=\bigl(&
 (\mathcal H_k^{\mathrm{sp}},\sigma_k,P_k,Q_k)_{k\leq N},
 (U_k,A_k,\eta_k)_{k<N},\\
 &(F_{0,k},\mathcal D_k,\mathcal T_k,
   \mathcal L^{\rm sgn}_{\mathrm{sp},k})_{k\leq N}\bigr).
\end{aligned}
\end{equation}
La macrostructure est
\begin{equation}
\label{eq:continuo-sp-m-graph}
 \mathfrak M_{\mathrm{sp},N}
 =\operatorname{Graph}(a_{\mathrm{sp},N}),
 \qquad
 \mathcal A_{\mathrm{sp},N}(z_N)
 =(z_N,a_{\mathrm{sp},N}z_N).
\end{equation}
La projection sur \(z_N\) est l'inverse de
\(\mathcal A_{\mathrm{sp},N}\). La troncature supprime les entrées dont
l'indice dépasse le nouvel horizon et vérifie
\begin{equation}
\label{eq:continuo-sp-a-naturalidad}
 m^{\mathrm{sp}}_{\ell,k}a_{\mathrm{sp},\ell}
 =a_{\mathrm{sp},k}\xi^{\mathrm{sp}}_{\ell,k}.
\end{equation}

\section{Publicateur comme graphe}

Le publicateur ne contracte pas la macrostructure en un scalaire. Il renvoie le
certificat complet
\begin{equation}
\label{eq:continuo-sp-p-accion}
 p_{\mathrm{sp},N}(m_N)
 =\left(
 \mathcal L^{\rm sgn}_{\mathrm{sp},k},
 \mathcal D_k,
 \mathcal T_k,
 I=\sum_{j=0}^{k-1}\mathcal D_j+\mathcal T_k
 \right)_{0\leq k\leq N}.
\end{equation}
On définit
\begin{equation}
\label{eq:continuo-sp-c-graph}
 \mathcal C_{\mathrm{sp},N}
 =\operatorname{Graph}(p_{\mathrm{sp},N}),
 \qquad
 \mathcal P_{\mathrm{sp},N}(m_N)
 =(m_N,p_{\mathrm{sp},N}m_N).
\end{equation}
La projection sur \(m_N\) est son inverse. Tronquer le certificat supprime seulement
les identités ultérieures, de sorte que
\begin{equation}
\label{eq:continuo-sp-p-naturalidad}
 c^{\mathrm{sp}}_{\ell,k}p_{\mathrm{sp},\ell}
 =p_{\mathrm{sp},k}m^{\mathrm{sp}}_{\ell,k}.
\end{equation}

\section{Théorème de la chaîne interne complète}

\begin{theorem}[Macrostructure spectrale--polaire engendrée par le continu]
\label{thm:continuo-sp-cadena-interna}
Les actions
\eqref{eq:continuo-sp-res-graph},
\eqref{eq:continuo-sp-pr-x},
\eqref{eq:continuo-sp-m-graph} et
\eqref{eq:continuo-sp-c-graph} forment la chaîne naturelle
\begin{equation}
\label{eq:continuo-sp-cadena-interna}
 \mathcal C_{\rm cont}^{\rm disc}
 \supset\mathcal C_{\rm cont}^{\mathrm{sp}}
 \xrightarrow[\cong]{\operatorname{Res}_{\mathrm{sp}}}
 \mathcal G_{\mathrm{sp}}
 \xrightarrow[\cong]{\operatorname{pr}_{X,\mathrm{sp}}}
 X_{\chi,\mathrm{sp}}
 \xrightarrow[\cong]{\mathcal A_{\mathrm{sp}}}
 \mathfrak M_{\mathrm{sp}}
 \xrightarrow[\cong]{\mathcal P_{\mathrm{sp}}}
 \mathcal C_{\mathrm{sp}}.
\end{equation}
Son contenu publié est le faisceau à deux feuillets, la graduation
\(\sigma_k\), les projecteurs \(P_k,Q_k\), le transport total \(U_k\),
la séparation \(U_k=A_k+\eta_k\), la retenue et la mémoire, le lecteur signé
et le télescopage
\eqref{eq:continuo-sp-telescopia-no-autonoma} avec terminal explicite. Chaque
objet conserve une projection canonique vers tous ses antécédents.
\end{theorem}

\begin{proof}
Les identités
\eqref{eq:continuo-sp-res-inversa} prouvent la première équivalence. La
définition dépendante \eqref{eq:continuo-sp-x-dependiente} prouve la
deuxième. Les graphes
\eqref{eq:continuo-sp-m-graph} et
\eqref{eq:continuo-sp-c-graph} prouvent les deux restantes. La naturalité des
quatre étapes est
\eqref{eq:continuo-sp-d-naturalidad},
\eqref{eq:continuo-sp-a-naturalidad} et
\eqref{eq:continuo-sp-p-naturalidad}, avec la naturalité de
\(r_{\mathrm{sp},k}\). Le contenu opératoriel est prouvé par
\eqref{eq:continuo-sp-u-isometria}, le
\cref{prop:continuo-sp-balance-local} et le
\cref{thm:continuo-sp-telescopia-no-autonoma}. Toutes les extensions apparaissent
dans \eqref{eq:continuo-sp-u-accion} ; aucun choix de
chemin n'intervient donc.
\end{proof}

\section{Retour de phase et évolution de l'état}

La branche conserve expressément
\begin{equation}
\label{eq:continuo-sp-no-reset}
 g(k+9)=g(k),
 \qquad
 x_{k+9}\neq x_k\quad\text{en général}.
\end{equation}
Le second membre enregistre l'accroissement de mémoire et les changements possibles
de feuillet, de retenue, de cylindre et de frontière. Ni \(\sigma_k\), ni les projecteurs, ni
le lecteur signé ne contractent cette différence.

\section{Provenance}

Les deux feuillets APP, l’état tritique holonomique, les fibres et extensions TPK,
le carry, le retour sans réinitialisation, la survie et le ledger sont
\texttt{ARCHITECTURE\_AUCTORIALE\_PRÉEXISTANTE}. L’obligation de conserver le
terminal fini est un \texttt{RÉSULTAT\_RETROUVÉ}. Le faisceau explicite de feuillets,
l’involution \(\sigma_k\), les projecteurs, la linéarisation qui somme toutes
les extensions, la décomposition \(U_k=A_k+\eta_k\), le produit non
autonome \(F_{0,k}\) et les graphes dépendants sont
\texttt{FORMALISATION\_NOUVELLE}. Le
\cref{thm:continuo-sp-telescopia-no-autonoma,thm:continuo-sp-cadena-interna}
constitue le \texttt{CERTIFICAT\_NOUVEAU} de cette formalisation.

\section{Falsadores internos}

\begin{proposition}[Falsificateurs de la macrostructure]
\label{prop:continuo-sp-falsadores}
La conclusion du
\cref{thm:continuo-sp-cadena-interna} n'est pas transférable si l'une quelconque
des substitutions suivantes se produit :
\begin{enumerate}
 \item choisir un prolongement dans
       \eqref{eq:continuo-sp-u-accion} et effacer les autres ;
 \item remplacer \(\mathcal G_{\mathrm{sp}}\) par l'image extérieure de
       \(D_{\mathrm{sp}}\) ;
 \item effacer \((x,D_{\mathrm{sp}}x)\) lors de la formation de
       \(X_{\chi,\mathrm{sp}}\) ;
 \item identifier le feuillet à l'orientation TRIT ou effacer le seuil ;
 \item rompre la cohérence
       \eqref{eq:continuo-sp-consistencia-medida} ;
 \item effacer la retenue ou enfreindre
       \eqref{eq:continuo-sp-carry-nativo}--
       \eqref{eq:continuo-sp-carry-operador} ;
 \item identifier \(\eta_k\) à toute la retenue ;
 \item remplacer la famille \((A_k)_k\) par un seul opérateur sans démontrer
       la stationnarité ;
 \item supprimer \(F_{0,N}^*F_{0,N}\) à un horizon fini ;
 \item affirmer \(\mathcal T_\infty=0\) sans preuve supplémentaire ;
 \item reconstruire une histoire à partir de la polarisation signée ;
 \item utiliser \(g(k+9)=g(k)\) pour affirmer \(x_{k+9}=x_k\) ;
 \item introduire un domaine, une fonction ou un critère extérieur pour
       définir l'une quelconque des flèches internes ;
 \item introduire une sortie ultérieure comme antécédent de la branche.
\end{enumerate}
\end{proposition}
\fi
   % Extracción quirúrgica del fuente teoremática V2 05a: líneas 512--523.
% El resto permanece en este archivo como procedencia, pero no se publica.
\iffalse
\chapter{Macrostructure spectrale--polaire interne du continu}
\label{chap:continuo-macro-sp-interna}

\section{Portée strictement interne}

La construction de ce chapitre se termine avant toute reconnaissance
externe. Ses seules données sont le support APP à deux feuillets, l'orientation
TRIT, les extensions et la retenue TPK, le système projectif d'histoires, la
mesure cylindrique et le ledger signé. Le résultat est une macrostructure
graduée avec transport, expression de changement de feuillet, lecteur signé et
télescopage non autonome. Aucune fonction extérieure n'est utilisée pour définir
l'une quelconque de ses applications.

\section{Domaine survivant de la branche}

Pour \(x_k\in\mathcal A_k\), soit
\begin{equation}
\label{eq:continuo-sp-msect-n}
 \operatorname{Msect}^{\mathrm{sp},(N)}_k(x_k)
 =\left\{(x_j)_{j=k}^{N}:
 \tau_{j+1,j}x_{j+1}=x_j,
 \ (x_j,x_{j+9})\in\Gamma_{9,j},
 \ \text{et le faisceau est conservé }\{\Sigma,\Pi\}\right\}.
\end{equation}
On définit
\begin{equation}
\label{eq:continuo-sp-dominio-superviviente}
 \mathcal C_k^{\mathrm{sp}}
 =\left\{x_k:
 \operatorname{Msect}^{\mathrm{sp},(N)}_k(x_k)\neq\varnothing
 \text{ pour tout }N\geq k\right\},
 \qquad
 \mathcal C_{\rm cont}^{\mathrm{sp}}
 =\varprojlim_k\mathcal C_k^{\mathrm{sp}}.
\end{equation}
Le seuil TRIT n'est pas exclu de ce domaine. Le faisceau conserve les deux feuillets
arithmétiques même si une carte visible n'en publie aucun lors d'un
événement concret.

\section{Paquet de treize coordonnées}

Pour \(x_k\in\mathcal C_k^{\mathrm{sp}}\), le paquet dépendant est
\begin{equation}
\label{eq:continuo-sp-d13}
\begin{aligned}
 D_{\mathrm{sp},k}(x_k)=\bigl(&
 \mathcal F_{\mathrm{sp},k}(x_k),
 \operatorname{Ext}_{\Gamma,\mathrm{sp},k}(x_k),
 \operatorname{Rel}_{\mathrm{sp},k}(x_k),
 \mathcal N_{\mathrm{sp},k}(x_k),
 \operatorname{or}_{\mathrm{sp},k}(x_k),\\
 &\operatorname{Carr}_{\mathrm{sp},k}(x_k),
 \operatorname{Mem}_{\mathrm{sp},k}(x_k),
 \partial_{\mathrm{sp},k}(x_k),
 \gamma_{\mathrm{sp},k}(x_k),
 \operatorname{Msect}_{\mathrm{sp},k}(x_k),\\
 &\operatorname{Surv}_{\mathrm{sp},k}(x_k),
 \operatorname{Term}_{\mathrm{sp},k}(x_k),
 \mathcal L^{\rm sgn}_{\mathrm{sp},k}(x_k)\bigr).
\end{aligned}
\end{equation}
Ses composantes sont déterminées corrélativement comme suit.

La fibre \(\mathcal F_{\mathrm{sp},k}\) contient l'état TPK et le faisceau
ordonné \(\{\Sigma,\Pi\}\). L'extension est la restriction complète de
\(\operatorname{Ext}_\Gamma\) aux histoires qui satisfont
\eqref{eq:continuo-sp-dominio-superviviente}. Les relations contiennent les
lois d'identité et de composition, le relèvement résidu--quotient, le cocycle de
retenue, la compatibilité de frontière et la naturalité de la troncature. La
coordonnée numérique contient profondeur, phase, résidus, quotients et signature
interne. L'orientation stocke séparément feuillet et orientation TRIT. La
mémoire contient le quotient, le ledger et l'incrément de tour. La
frontière et le chemin restent non contractés. La multisection est le système
complet des prolongements de \eqref{eq:continuo-sp-msect-n}. La
survie est l'appartenance stable de
\eqref{eq:continuo-sp-dominio-superviviente}. Le terminal et le lecteur sont
construits dans les sections suivantes.

Les troncatures de toutes les coordonnées forment une application
\(u^{\mathrm{sp}}_{\ell,k}\) qui satisfait
\begin{equation}
\label{eq:continuo-sp-d-naturalidad}
 u^{\mathrm{sp}}_{\ell,k}D_{\mathrm{sp},\ell}
 =D_{\mathrm{sp},k}\tau_{\ell,k}.
\end{equation}

\section{Graphe de la restriction et section dépendante}

On définit
\begin{equation}
\label{eq:continuo-sp-g-graph}
 \mathcal G_{\mathrm{sp},k}
 =\operatorname{Graph}(D_{\mathrm{sp},k})
 =\{(x_k,D_{\mathrm{sp},k}x_k):
     x_k\in\mathcal C_k^{\mathrm{sp}}\}.
\end{equation}
La restriction totale est
\begin{equation}
\label{eq:continuo-sp-res-graph}
 \operatorname{Res}_{\mathrm{sp}}=s_{\mathrm{sp}}:
 \mathcal C_{\rm cont}^{\mathrm{sp}}
 \longrightarrow\mathcal G_{\mathrm{sp}},
 \qquad
 s_{\mathrm{sp}}(x)=(x,D_{\mathrm{sp}}x).
\end{equation}
La projection \(\pi_1(x,d)=x\) vérifie
\begin{equation}
\label{eq:continuo-sp-res-inversa}
 \pi_1s_{\mathrm{sp}}=I,
 \qquad
 s_{\mathrm{sp}}\pi_1=I_{\mathcal G_{\mathrm{sp}}}.
\end{equation}

Soit \(r_{\mathrm{sp},k}(x_k)\) la multisection des feuillets et des chemins qui appartient
déjà à \(D_{\mathrm{sp},k}(x_k)\). On définit
\begin{equation}
\label{eq:continuo-sp-x-dependiente}
 X_{\chi,\mathrm{sp},k}
 =\{(r_{\mathrm{sp},k}x_k,x_k,D_{\mathrm{sp},k}x_k):
     x_k\in\mathcal C_k^{\mathrm{sp}}\},
\end{equation}
\begin{equation}
\label{eq:continuo-sp-pr-x}
 \operatorname{pr}_{X,\mathrm{sp}}(x,D_{\mathrm{sp}}x)
 =(r_{\mathrm{sp}}x,x,D_{\mathrm{sp}}x).
\end{equation}
Oublier la première coordonnée de
\eqref{eq:continuo-sp-x-dependiente} est l'inverse de
\(\operatorname{pr}_{X,\mathrm{sp}}\).

\section{Faisceau de feuillets et mesure projective}

Soit
\begin{equation}
\label{eq:continuo-sp-haz-hojas}
 \widetilde{\mathcal C}_k^{\mathrm{sp}}
 =\{(x_k,h):x_k\in\mathcal C_k^{\mathrm{sp}},
              h\in\{\Sigma,\Pi\}\}.
\end{equation}
La mesure projective du continu induit des probabilités
\(\widetilde\mu_k\) sur ce faisceau. On se restreint au support positif ; pour
\(z\) dans ce support,
\begin{equation}
\label{eq:continuo-sp-consistencia-medida}
 \widetilde\mu_k(z)
 =\sum_{\widetilde\tau_{k+1,k}(z')=z}
  \widetilde\mu_{k+1}(z').
\end{equation}
La linéarisation de niveau est
\begin{equation}
\label{eq:continuo-sp-hilbert-nivel}
 \mathcal H_k^{\mathrm{sp}}
 =\ell^2(\operatorname{supp}\widetilde\mu_k),
\end{equation}
de base orthonormale \(e_z\).

\section{Transport de toutes les extensions et retenue}

Pour \(z'\) qui prolonge \(z\), nous définissons le poids conditionnel positif
\begin{equation}
\label{eq:continuo-sp-peso-condicional}
 w_k(z'\mid z)
 =\left(\frac{\widetilde\mu_{k+1}(z')}
               {\widetilde\mu_k(z)}\right)^{1/2}.
\end{equation}
Le transport est
\begin{equation}
\label{eq:continuo-sp-u-accion}
 U_ke_z
 =\sum_{\substack{z':\widetilde\tau_{k+1,k}(z')=z\\
            (z,z')\in\operatorname{Ext}_{\Gamma,\mathrm{sp},k}}}
 w_k(z'\mid z)e_{z'}.
\end{equation}
La somme contient toutes les extensions admissibles. Elle ne dépend pas d'un sélecteur.

Par \eqref{eq:continuo-sp-consistencia-medida},
\[
 \sum_{z':\widetilde\tau(z')=z}|w_k(z'\mid z)|^2=1.
\]
Les ensembles de descendants de deux parents distincts sont disjoints. Par
conséquent,
\begin{equation}
\label{eq:continuo-sp-u-isometria}
 U_k^*U_k=I_{\mathcal H_k^{\mathrm{sp}}}.
\end{equation}

Chaque base \(e_{z'}\) inclut la retenue de l'extension correspondante. Pour
des chemins composables, on conserve
\begin{align}
 \kappa(\gamma_2\gamma_1)
 &=\kappa(\gamma_2)+
   \mathsf C(\gamma_2)\kappa(\gamma_1),
 \label{eq:continuo-sp-carry-nativo}\\
 \operatorname{Carr}(\gamma_2\gamma_1)
 &=\operatorname{Carr}(\gamma_2)H(\gamma_1)
  +Y(\gamma_2)\operatorname{Carr}(\gamma_1).
 \label{eq:continuo-sp-carry-operador}
\end{align}
Les poids conditionnels se multiplient le long de l'unique chaîne de
troncatures d'un descendant. Les équations
\eqref{eq:continuo-sp-carry-nativo}--
\eqref{eq:continuo-sp-carry-operador} garantissent que les étiquettes de retenue
coïncident. Par conséquent,
\begin{equation}
\label{eq:continuo-sp-u-composicion}
 U_{\ell,k}=U_{\ell-1}\cdots U_k.
\end{equation}

\section{Graduation, projecteurs et décomposition du transport}

La graduation de feuillet est l'involution
\begin{equation}
\label{eq:continuo-sp-sigma}
 \sigma_ke_{(x,\Sigma)}=+e_{(x,\Sigma)},
 \qquad
 \sigma_ke_{(x,\Pi)}=-e_{(x,\Pi)}.
\end{equation}
Ainsi,
\[
 \sigma_k^*=\sigma_k,
 \qquad
 \sigma_k^2=I.
\]
Les projecteurs internes sont
\begin{equation}
\label{eq:continuo-sp-p-q}
 P_k=\frac{I+\sigma_k}{2},
 \qquad
 Q_k=\frac{I-\sigma_k}{2}.
\end{equation}
Elles satisfont
\begin{equation}
\label{eq:continuo-sp-p-q-relaciones}
 P_k^2=P_k,
 \quad Q_k^2=Q_k,
 \quad P_kQ_k=0,
 \quad P_k+Q_k=I.
\end{equation}

Nous séparons la partie du transport qui conserve le feuillet de celle qui le change :
\begin{align}
 A_k&=P_{k+1}U_kP_k+Q_{k+1}U_kQ_k,
 \label{eq:continuo-sp-a}\\
 \eta_k&=Q_{k+1}U_kP_k+P_{k+1}U_kQ_k.
 \label{eq:continuo-sp-eta}
\end{align}
Alors
\begin{equation}
\label{eq:continuo-sp-u-a-eta}
 U_k=A_k+\eta_k,
\end{equation}
et la graduation vérifie
\begin{equation}
\label{eq:continuo-sp-paridad-a-eta}
 \sigma_{k+1}A_k=A_k\sigma_k,
 \qquad
 \sigma_{k+1}\eta_k=-\eta_k\sigma_k.
\end{equation}

\begin{proposition}[Bilan local des feuillets]
\label{prop:continuo-sp-balance-local}
Pour chaque niveau,
\begin{equation}
\label{eq:continuo-sp-balance-local}
 A_k^*A_k+\eta_k^*\eta_k=I.
\end{equation}
\end{proposition}

\begin{proof}
L'orthogonalité de \(P_{k+1}\) et \(Q_{k+1}\) donne
\begin{align*}
 A_k^*A_k
 &=P_kU_k^*P_{k+1}U_kP_k
   +Q_kU_k^*Q_{k+1}U_kQ_k,\\
 \eta_k^*\eta_k
 &=P_kU_k^*Q_{k+1}U_kP_k
   +Q_kU_k^*P_{k+1}U_kQ_k.
\end{align*}
En sommant et en utilisant \(P_{k+1}+Q_{k+1}=I\),
\[
 A_k^*A_k+\eta_k^*\eta_k
 =P_kU_k^*U_kP_k+Q_kU_k^*U_kQ_k
 =P_k+Q_k=I.
\]
\end{proof}

L'opérateur \(\eta_k\) enregistre l'abandon du canal qui conserve le feuillet. Il
n'est pas la retenue entière : les coordonnées
\eqref{eq:continuo-sp-carry-nativo}--
\eqref{eq:continuo-sp-carry-operador} restent dans le paquet
\eqref{eq:continuo-sp-d13}.

\section{Lecteur signé avec ledger}

Pour \(v\in\mathcal H_k^{\mathrm{sp}}\), la polarisation signée est
\begin{equation}
\label{eq:continuo-sp-polarizacion-firmada}
 \operatorname{pol}_k(v)
 =\langle v,\sigma_kv\rangle
 =\|P_kv\|^2-\|Q_kv\|^2.
\end{equation}
Le lecteur spécialisé conserve en outre la source :
\begin{equation}
\label{eq:continuo-sp-lector-firmado}
 \mathcal L^{\rm sgn}_{\mathrm{sp},k}(x_k;v)
 =\bigl(
   \Lambda(\gamma_{x_k}),
   \Sigma_{\rm reg}(\Lambda(\gamma_{x_k})),
   \|P_kv\|^2,
   \|Q_kv\|^2,
   \operatorname{pol}_k(v)\bigr).
\end{equation}
L'orientation TRIT est conservée dans \(\Lambda(\gamma_{x_k})\) ; elle ne
s'identifie pas au signe de feuillet. À l'horizon \(N\), le registre garde la
famille des lecteurs pour tous les niveaux \(k\leq N\). Tronquer signifie
supprimer uniquement les entrées ultérieures.

\section{Produit non autonome et terminal}

On définit
\begin{equation}
\label{eq:continuo-sp-f-producto}
 F_{0,0}=I_{\mathcal H_0^{\mathrm{sp}}},
 \qquad
 F_{0,k}=A_{k-1}\cdots A_1A_0
 \quad(k\geq1).
\end{equation}
Par conséquent,
\begin{equation}
\label{eq:continuo-sp-f-recursion}
 F_{0,k+1}=A_kF_{0,k}.
\end{equation}
Le terme publié lors de l'abandon du canal à l'étape \(k\) et le terminal à
l'horizon \(N\) sont
\begin{equation}
\label{eq:continuo-sp-defecto-terminal}
 \mathcal D_k=F_{0,k}^*\eta_k^*\eta_kF_{0,k},
 \qquad
 \mathcal T_N=F_{0,N}^*F_{0,N}.
\end{equation}

\begin{theorem}[Télescopage spectral--polaire non autonome]
\label{thm:continuo-sp-telescopia-no-autonoma}
Pour tout horizon fini \(N\),
\begin{equation}
\label{eq:continuo-sp-telescopia-no-autonoma}
 \boxed{
 I=\sum_{k=0}^{N-1}
 F_{0,k}^*\eta_k^*\eta_kF_{0,k}
 +F_{0,N}^*F_{0,N}.}
\end{equation}
De plus,
\[
 0\preceq\mathcal T_{N+1}
 \preceq\mathcal T_N\preceq I,
\]
de sorte qu'il existe la limite forte positive
\[
 \mathcal T_\infty
 =\operatorname{s-lim}_{N\to\infty}\mathcal T_N
\]
et
\begin{equation}
\label{eq:continuo-sp-telescopia-infinita}
 I=\sum_{k=0}^{\infty}
 F_{0,k}^*\eta_k^*\eta_kF_{0,k}
 +\mathcal T_\infty
\end{equation}
en topologie forte. Le théorème n'affirme pas que
\(\mathcal T_\infty\) soit nul.
\end{theorem}

\begin{proof}
En multipliant \eqref{eq:continuo-sp-balance-local} par
\(F_{0,k}^*\) et \(F_{0,k}\), et en utilisant
\eqref{eq:continuo-sp-f-recursion}, on obtient
\[
 F_{0,k}^*F_{0,k}
 =F_{0,k+1}^*F_{0,k+1}
  +F_{0,k}^*\eta_k^*\eta_kF_{0,k}.
\]
La somme de \(k=0\) à \(N-1\) annule tous les termes
intermédiaires et laisse précisément
\eqref{eq:continuo-sp-telescopia-no-autonoma}. Comme
\(A_k^*A_k\preceq I\),
\[
 \mathcal T_{N+1}
 =F_{0,N}^*A_N^*A_NF_{0,N}
 \preceq F_{0,N}^*F_{0,N}=\mathcal T_N.
\]
La monotonie et la borne positive fournissent la limite forte. Le passage à la
limite dans l'identité finie donne
\eqref{eq:continuo-sp-telescopia-infinita}.
\end{proof}

\section{Assembleur comme graphe}

Pour
\[
 z_N=(r_{\mathrm{sp},N}x_N,x_N,D_{\mathrm{sp},N}x_N)
 \in X_{\chi,\mathrm{sp},N},
\]
on définit l'assembleur matériel
\begin{equation}
\label{eq:continuo-sp-a-accion}
\begin{aligned}
 a_{\mathrm{sp},N}(z_N)=\bigl(&
 (\mathcal H_k^{\mathrm{sp}},\sigma_k,P_k,Q_k)_{k\leq N},
 (U_k,A_k,\eta_k)_{k<N},\\
 &(F_{0,k},\mathcal D_k,\mathcal T_k,
   \mathcal L^{\rm sgn}_{\mathrm{sp},k})_{k\leq N}\bigr).
\end{aligned}
\end{equation}
La macrostructure est
\begin{equation}
\label{eq:continuo-sp-m-graph}
 \mathfrak M_{\mathrm{sp},N}
 =\operatorname{Graph}(a_{\mathrm{sp},N}),
 \qquad
 \mathcal A_{\mathrm{sp},N}(z_N)
 =(z_N,a_{\mathrm{sp},N}z_N).
\end{equation}
La projection sur \(z_N\) est l'inverse de
\(\mathcal A_{\mathrm{sp},N}\). La troncature supprime les entrées dont
l'indice dépasse le nouvel horizon et vérifie
\begin{equation}
\label{eq:continuo-sp-a-naturalidad}
 m^{\mathrm{sp}}_{\ell,k}a_{\mathrm{sp},\ell}
 =a_{\mathrm{sp},k}\xi^{\mathrm{sp}}_{\ell,k}.
\end{equation}

\section{Publicateur comme graphe}

Le publicateur ne contracte pas la macrostructure en un scalaire. Il renvoie le
certificat complet
\begin{equation}
\label{eq:continuo-sp-p-accion}
 p_{\mathrm{sp},N}(m_N)
 =\left(
 \mathcal L^{\rm sgn}_{\mathrm{sp},k},
 \mathcal D_k,
 \mathcal T_k,
 I=\sum_{j=0}^{k-1}\mathcal D_j+\mathcal T_k
 \right)_{0\leq k\leq N}.
\end{equation}
On définit
\begin{equation}
\label{eq:continuo-sp-c-graph}
 \mathcal C_{\mathrm{sp},N}
 =\operatorname{Graph}(p_{\mathrm{sp},N}),
 \qquad
 \mathcal P_{\mathrm{sp},N}(m_N)
 =(m_N,p_{\mathrm{sp},N}m_N).
\end{equation}
La projection sur \(m_N\) est son inverse. Tronquer le certificat supprime seulement
les identités ultérieures, de sorte que
\begin{equation}
\label{eq:continuo-sp-p-naturalidad}
 c^{\mathrm{sp}}_{\ell,k}p_{\mathrm{sp},\ell}
 =p_{\mathrm{sp},k}m^{\mathrm{sp}}_{\ell,k}.
\end{equation}

\section{Théorème de la chaîne interne complète}

\begin{theorem}[Macrostructure spectrale--polaire engendrée par le continu]
\label{thm:continuo-sp-cadena-interna}
Les actions
\eqref{eq:continuo-sp-res-graph},
\eqref{eq:continuo-sp-pr-x},
\eqref{eq:continuo-sp-m-graph} et
\eqref{eq:continuo-sp-c-graph} forment la chaîne naturelle
\begin{equation}
\label{eq:continuo-sp-cadena-interna}
 \mathcal C_{\rm cont}^{\rm disc}
 \supset\mathcal C_{\rm cont}^{\mathrm{sp}}
 \xrightarrow[\cong]{\operatorname{Res}_{\mathrm{sp}}}
 \mathcal G_{\mathrm{sp}}
 \xrightarrow[\cong]{\operatorname{pr}_{X,\mathrm{sp}}}
 X_{\chi,\mathrm{sp}}
 \xrightarrow[\cong]{\mathcal A_{\mathrm{sp}}}
 \mathfrak M_{\mathrm{sp}}
 \xrightarrow[\cong]{\mathcal P_{\mathrm{sp}}}
 \mathcal C_{\mathrm{sp}}.
\end{equation}
Son contenu publié est le faisceau à deux feuillets, la graduation
\(\sigma_k\), les projecteurs \(P_k,Q_k\), le transport total \(U_k\),
la séparation \(U_k=A_k+\eta_k\), la retenue et la mémoire, le lecteur signé
et le télescopage
\eqref{eq:continuo-sp-telescopia-no-autonoma} avec terminal explicite. Chaque
objet conserve une projection canonique vers tous ses antécédents.
\end{theorem}

\begin{proof}
Les identités
\eqref{eq:continuo-sp-res-inversa} prouvent la première équivalence. La
définition dépendante \eqref{eq:continuo-sp-x-dependiente} prouve la
deuxième. Les graphes
\eqref{eq:continuo-sp-m-graph} et
\eqref{eq:continuo-sp-c-graph} prouvent les deux restantes. La naturalité des
quatre étapes est
\eqref{eq:continuo-sp-d-naturalidad},
\eqref{eq:continuo-sp-a-naturalidad} et
\eqref{eq:continuo-sp-p-naturalidad}, avec la naturalité de
\(r_{\mathrm{sp},k}\). Le contenu opératoriel est prouvé par
\eqref{eq:continuo-sp-u-isometria}, le
\cref{prop:continuo-sp-balance-local} et le
\cref{thm:continuo-sp-telescopia-no-autonoma}. Toutes les extensions apparaissent
dans \eqref{eq:continuo-sp-u-accion} ; aucun choix de
chemin n'intervient donc.
\end{proof}

\fi
\section{Retour de phase et évolution non réinitialisée de l'état commun}

La branche conserve expressément
\begin{equation}
\label{eq:continuo-sp-no-reset}
 g(k+9)=g(k),
 \qquad
 x_{k+9}\neq x_k\quad\text{en général}.
\end{equation}
Le second membre enregistre l'accroissement de mémoire et les changements possibles
de feuillet, de retenue, de cylindre et de frontière. Ni \(\sigma_k\), ni les projecteurs, ni
le lecteur signé ne contractent cette différence.

\iffalse
\section{Provenance}

Les deux feuillets APP, l’état tritique holonomique, les fibres et extensions TPK,
le carry, le retour sans réinitialisation, la survie et le ledger sont
\texttt{ARCHITECTURE\_AUCTORIALE\_PRÉEXISTANTE}. L’obligation de conserver le
terminal fini est un \texttt{RÉSULTAT\_RETROUVÉ}. Le faisceau explicite de feuillets,
l’involution \(\sigma_k\), les projecteurs, la linéarisation qui somme toutes
les extensions, la décomposition \(U_k=A_k+\eta_k\), le produit non
autonome \(F_{0,k}\) et les graphes dépendants sont
\texttt{FORMALISATION\_NOUVELLE}. Le
\cref{thm:continuo-sp-telescopia-no-autonoma,thm:continuo-sp-cadena-interna}
constitue le \texttt{CERTIFICAT\_NOUVEAU} de cette formalisation.

\section{Falsadores internos}

\begin{proposition}[Falsificateurs de la macrostructure]
\label{prop:continuo-sp-falsadores}
La conclusion du
\cref{thm:continuo-sp-cadena-interna} n'est pas transférable si l'une quelconque
des substitutions suivantes se produit :
\begin{enumerate}
 \item choisir un prolongement dans
       \eqref{eq:continuo-sp-u-accion} et effacer les autres ;
 \item remplacer \(\mathcal G_{\mathrm{sp}}\) par l'image extérieure de
       \(D_{\mathrm{sp}}\) ;
 \item effacer \((x,D_{\mathrm{sp}}x)\) lors de la formation de
       \(X_{\chi,\mathrm{sp}}\) ;
 \item identifier le feuillet à l'orientation TRIT ou effacer le seuil ;
 \item rompre la cohérence
       \eqref{eq:continuo-sp-consistencia-medida} ;
 \item effacer la retenue ou enfreindre
       \eqref{eq:continuo-sp-carry-nativo}--
       \eqref{eq:continuo-sp-carry-operador} ;
 \item identifier \(\eta_k\) à toute la retenue ;
 \item remplacer la famille \((A_k)_k\) par un seul opérateur sans démontrer
       la stationnarité ;
 \item supprimer \(F_{0,N}^*F_{0,N}\) à un horizon fini ;
 \item affirmer \(\mathcal T_\infty=0\) sans preuve supplémentaire ;
 \item reconstruire une histoire à partir de la polarisation signée ;
 \item utiliser \(g(k+9)=g(k)\) pour affirmer \(x_{k+9}=x_k\) ;
 \item introduire un domaine, une fonction ou un critère extérieur pour
       définir l'une quelconque des flèches internes ;
 \item introduire une sortie ultérieure comme antécédent de la branche.
\end{enumerate}
\end{proposition}
\fi
 \endgroup

L'espérance conditionnelle associée au même poids est
\begin{equation}
\label{eq:pdfv2-cor-esperanza-comun}
 (E_kF)(\widehat x):=
 \sum_{(\alpha,\widehat y)\in\operatorname{Ch}_k(\widehat x)}
 p_k(\alpha,\widehat y\mid\widehat x)F(\alpha,\widehat y).
\end{equation}
Pour l'accroissement fin \(b_f\), on pose
\begin{equation}
\label{eq:pdfv2-cor-kappa-jensen}
 b_c:=E_kb_f,\qquad
 \kappa_k^{\rm can}:=\frac{E_k|b_f|-|b_c|}{2}\geq0.
\end{equation}
Alors
\begin{equation}
\label{eq:pdfv2-cor-jensen-hojas}
 E_kb_f^+=b_c^++\kappa_k^{\rm can},\qquad
 E_kb_f^-=b_c^-+\kappa_k^{\rm can}.
\end{equation}
L'exposant \({\rm can}\) distingue cette retenue d'annulation de la retenue entière \(\kappa_\alpha\) ; toutes deux proviennent du même transport, mais sont de types distincts.

Les mémoires signées ne nécessitent pas de condition initiale extérieure : pour \(\gamma=(\epsilon_1,\ldots,\epsilon_k)\), elles sont définies par la somme du registre lui-même,
\begin{equation}
\label{eq:pdfv2-cor-memorias}
 \begin{aligned}
 M_0^\pm&:=0,\\
 M_k^\pm(\gamma)&:=\sum_{j=1}^k b_{\epsilon_j}^\pm,\\
 M_{k+1}^\pm(\gamma\alpha)&:=M_k^\pm(\gamma)+b_\alpha^\pm,\\
 D_k^0&:=M_k^+-M_k^-, &
 H_k^0&:=M_k^++M_k^-.
 \end{aligned}
\end{equation}
Ainsi, le bilan signé, la mémoire totale et la polarisation \(\langle v,\sigma_kv\rangle\) sont trois lectures coordonnées du même préfixe, non trois états reconstruits séparément.

\subsection{Raffinement nonadique, pertes orthogonales et mémoire terminale}
L’indice \(j\) du développement suivant est le niveau du même arbre de
prolongations et son registre à treize coordonnées est la restriction du
registre commun de \eqref{eq:pdfv2-cor-registro-extension}. On fixe donc
\(D_{{\rm sol},j}:=D_j^{(13)}\) et aucun domaine supplémentaire n’est postulé.
\begingroup
  \let\section\subsection
  % Extracción quirúrgica del fuente teoremática V2 05b: líneas 100--359.
% El resto permanece en este archivo como procedencia, pero no se publica.
\iffalse
% Macroestructura solenoidal estrictamente interna del continuo discreto.
% Este módulo no introduce ninguna ecuación clásica ni un problema exterior.

\chapter{Macrostructure interne de bilan, mémoire et télescopage}
\label{chap:pdfv2-sol-interna}

\begin{statusbox}
L'unique entrée de ce chapitre est la structure discrète du continu.
La branche \(\mathrm{sol}\) se construit avant toute évaluation scalaire
extérieure et avant toute formulation différentielle classique. Sa chaîne
complète est
\[
 \mathcal C_{\rm cont}^{\rm disc}
 \supset\mathcal C_{\rm cont}^{\rm sol}
 \xrightarrow{\operatorname{Res}_{\rm sol}}
 \mathcal G_{\rm sol}
 \xrightarrow{\operatorname{pr}_{X,{\rm sol}}}
 X_{\chi,{\rm sol}}
 \xrightarrow{\mathcal A_{\rm sol}}
 \mathfrak M_{\rm sol}
 \xrightarrow{\mathcal P_{\rm sol}}
 \mathcal C_{\rm sol}^{\rm HMT}.
\]
Chaque flèche conserve comme coordonnées l'histoire et l'appareil qui la
produit. Aucune image de cette chaîne n'est un quotient qui oublie sa source.
\end{statusbox}

\section{Domaine propre avant toute évaluation ultérieure}
\label{sec:pdfv2-sol-dominio}

Soit
\[
 x=(x_0\xrightarrow{e_0}x_1\xrightarrow{e_1}\cdots)
 \in\mathcal C_{\rm cont}^{\rm disc}
\]
une histoire survivante. Chaque arête \(e_j\) conserve ses deux feuillets,
son orientation, son résidu, son quotient et sa retenue entière. Soient
\(N_j^+(e_j)\) et \(N_j^-(e_j)\) les multiplicités des deux feuillets et soit
\(\Delta_{\gamma_j}\in\mathbb Z\) le coefficient orienté du chemin, avec la
retenue de la composition déjà incorporée. L'accroissement signé interne est
\begin{equation}\label{eq:pdfv2-sol-b}
 b_j=b(e_j):=
 \Delta_{\gamma_j}\bigl(N_j^+(e_j)-N_j^-(e_j)\bigr)\in\mathbb Z.
\end{equation}
Il ne s'obtient pas en évaluant une trajectoire extérieure : c'est une coordonnée du
ledger de l'arête elle-même.

Ses parties positive et négative se construisent par
\begin{equation}\label{eq:pdfv2-sol-bpm}
 b_j^+:=\frac{|b_j|+b_j}{2},
 \qquad
 b_j^-:=\frac{|b_j|-b_j}{2}.
\end{equation}
Par conséquent,
\begin{equation}\label{eq:pdfv2-sol-bpm-identidades}
 b_j=b_j^+-b_j^-,
 \qquad
 |b_j|=b_j^++b_j^-,
 \qquad
 b_j^+b_j^-=0.
\end{equation}

Les deux mémoires primitives et leurs deux lecteurs sont
\begin{align}
 M_0^+=M_0^-&=0,
 &M_j^+&=\sum_{0\leq i<j}b_i^+,
 &M_j^-&=\sum_{0\leq i<j}b_i^-;
 \label{eq:pdfv2-sol-memorias}\\
 D_j^0&:=M_j^+-M_j^-,
 &H_j^0&:=M_j^++M_j^-.
 \label{eq:pdfv2-sol-dh}
\end{align}
En conséquence,
\begin{equation}\label{eq:pdfv2-sol-dh-diferencias}
 D_{j+1}^0-D_j^0=b_j,
 \qquad
 H_{j+1}^0-H_j^0=|b_j|.
\end{equation}
La mémoire signée \(D^0\) et la mémoire totale \(H^0\) restent des
objets distincts ; conserver seulement leur différence finale détruirait la
généalogie.

La sous-structure \(\mathcal C_{\rm cont}^{\rm sol}\) se compose des histoires
pour lesquelles, à tout préfixe et à tout raffinement nonadique :
\begin{enumerate}
 \item les accroissements \eqref{eq:pdfv2-sol-b} et les deux mémoires
       \eqref{eq:pdfv2-sol-memorias} existent ;
 \item les cylindres fins se regroupent en fibres de neuf extensions avec la
       même loi de probabilité projective ;
 \item les projections de profondeur et de résolution définies ci-dessous sont
       compatibles avec les extensions d'histoire ;
 \item chaque ligne de pertes et chaque colonne terminale est finie à horizon
       fini ;
 \item tout préfixe admet au moins une extension compatible à tout
       horizon ultérieur.
\end{enumerate}
Ce sont des conditions d'appartenance à la branche, non des prémisses tacites
attribuées à une histoire arbitraire.

\fi
\section{Raffinement nonadique commun : inclusion, espérance et mémoire de bilan}
\label{sec:pdfv2-sol-e9j9}

Soit \(Y_j(x)\) l'ensemble fini des cylindres compatibles avec le préfixe
\(x_{\leq j}\), muni de la mesure projective \(\mu_j\). La troncature
\[
 \tau_{j+1,j}:Y_{j+1}(x)\longrightarrow Y_j(x)
\]
a neuf extensions admissibles au-dessus de chaque cylindre. On considère les
espaces finis
\[
 \mathscr H_j(x):=\ell^2\bigl(Y_j(x),\mu_j\bigr).
\]
L'inclusion et l'espérance nonadiques sont
\begin{align}
 J_{9,j}:\mathscr H_j&\longrightarrow\mathscr H_{j+1},
 &(J_{9,j}f)(y)&=f(\tau_{j+1,j}y),
 \label{eq:pdfv2-sol-j9}\\
 E_{9,j}:\mathscr H_{j+1}&\longrightarrow\mathscr H_j,
 &(E_{9,j}g)(z)&=\frac1{9}
   \sum_{\tau_{j+1,j}y=z}g(y).
 \label{eq:pdfv2-sol-e9}
\end{align}
La compatibilité projective de \(\mu_{j+1}\) et \(\mu_j\) donne
\begin{equation}\label{eq:pdfv2-sol-ej-projector}
 E_{9,j}J_{9,j}=I_{\mathscr H_j},
 \qquad
 J_{9,j}^*=E_{9,j},
 \qquad
 P_{9,j}:=J_{9,j}E_{9,j}=P_{9,j}^*=P_{9,j}^2.
\end{equation}
Le projecteur microscopique complémentaire est
\begin{equation}\label{eq:pdfv2-sol-qmicro-preliminar}
 Q_{9,j}:=I-P_{9,j}.
\end{equation}

Si \(b_f\) est l'accroissement sur la fibre fine et
\begin{equation}\label{eq:pdfv2-sol-bcoarse}
 b_c:=E_{9,j}b_f,
\end{equation}
la convexité élémentaire de \(|\cdot|\) construit, et ne postule pas, la retenue de
compensation
\begin{equation}\label{eq:pdfv2-sol-kappa}
 \kappa_j
 :=\frac{E_{9,j}|b_f|-|b_c|}{2}\geq0.
\end{equation}
En effet, en utilisant \(t^\pm=(|t|\pm t)/2\), on obtient
\begin{align}
 E_{9,j}b_f^+
 &=\frac{E_{9,j}|b_f|+E_{9,j}b_f}{2}
   =b_c^++\kappa_j,
 \label{eq:pdfv2-sol-jensen-plus}\\
 E_{9,j}b_f^-
 &=\frac{E_{9,j}|b_f|-E_{9,j}b_f}{2}
   =b_c^-+\kappa_j.
 \label{eq:pdfv2-sol-jensen-minus}
\end{align}
Les deux feuillets perdent la même quantité lors de la compression ; c'est pourquoi la différence
signée est conservée tandis que la mémoire totale enregistre la compensation :
\begin{equation}\label{eq:pdfv2-sol-jensen-dh}
 E_{9,j}(b_f^+-b_f^-)=b_c,
 \qquad
 E_{9,j}(b_f^++b_f^-)=|b_c|+2\kappa_j.
\end{equation}

\section{Factorisation orthogonale des degrés de profondeur et des couches de résolution dans \(\mathscr H_j^{\mathrm{depth}}\widehat\otimes\mathscr H_j^{\mathrm{res}}\)}
\label{sec:pdfv2-sol-micro-shell}

Pour empêcher qu'une même perte soit comptée à la fois comme profondeur et
comme résolution, l'espace de chaque niveau est typé comme
\begin{equation}\label{eq:pdfv2-sol-factor-hilbert}
 \mathscr H_j
 =\mathscr H_j^{\rm depth}\widehat\otimes
  \mathscr H_j^{\rm res}.
\end{equation}
Pour \(j\geq1\), le premier facteur porte
\[
 P_j^{\rm depth}:=J_{9,j-1}E_{9,j-1}
 \quad\text{sur }\mathscr H_j^{\rm depth},
\]
et l'on fixe \(P_0^{\rm depth}=I\). Le second facteur porte la partition
finie des shells héritée de la frontière des cylindres. Si
\(\sigma_j:Y_j^{\rm res}\to S_j\) est cette partition, ses opérateurs sont
\begin{align}
 J_j^{\rm res}f&=f\circ\sigma_j,
 &E_j^{\rm res}g(s)&=
 \frac1{\#\sigma_j^{-1}(s)}
 \sum_{\sigma_j(y)=s}g(y),
 \label{eq:pdfv2-sol-shell-je}\\
 P_j^{\rm res}&:=J_j^{\rm res}E_j^{\rm res}.
 \label{eq:pdfv2-sol-shell-p}
\end{align}

On définit les trois projecteurs mutuellement orthogonaux
\begin{align}
 P_j^{\rm term}
 &:=P_j^{\rm depth}\otimes P_j^{\rm res},
 \label{eq:pdfv2-sol-pterm}\\
 Q_j^{\rm micro}
 &:=(I-P_j^{\rm depth})\otimes I,
 \label{eq:pdfv2-sol-qmicro}\\
 Q_j^{\rm shell}
 &:=P_j^{\rm depth}\otimes(I-P_j^{\rm res}).
 \label{eq:pdfv2-sol-qshell}
\end{align}
Alors
\begin{equation}\label{eq:pdfv2-sol-ortogonalidad}
 Q_j^{\rm micro}Q_j^{\rm shell}=0,
 \qquad
 I=P_j^{\rm term}+Q_j^{\rm micro}+Q_j^{\rm shell},
\end{equation}
et, pour tout \(\xi\in\mathscr H_j\),
\begin{equation}\label{eq:pdfv2-sol-pitagoras}
 \|\xi\|^2
 =\|P_j^{\rm term}\xi\|^2
  +\|Q_j^{\rm micro}\xi\|^2
  +\|Q_j^{\rm shell}\xi\|^2.
\end{equation}
Le terme
\((I-P_j^{\rm depth})\otimes(I-P_j^{\rm res})\) n'apparaît pas une seconde fois : il est déjà
contenu exactement une fois dans \(Q_j^{\rm micro}\).

\section{Extensions de cylindre et ligne complète de pertes}
\label{sec:pdfv2-sol-gamma-loss}

Une extension de cylindres du niveau \(j\) au niveau \(j+1\) induit l'application
\begin{equation}\label{eq:pdfv2-sol-gamma}
 \Gamma_j:\mathscr H_j\longrightarrow\mathscr H_{j+1},
 \qquad
 \Gamma_j f=f\circ\tau_{j+1,j},
\end{equation}
avec toutes les étiquettes de feuillet, d'orientation, de retenue, de mémoire et de chemin
transportées dans la base des cylindres. Pour \(n>j\), on écrit
\begin{equation}\label{eq:pdfv2-sol-gamma-composition}
 \Gamma_{j:n}:=\Gamma_{n-1}\cdots\Gamma_j,
 \qquad
 \Gamma_{j:j}:=I.
\end{equation}

Sur \(\mathscr H_j\), les accroissements non négatifs définissent des opérateurs de
multiplication
\begin{equation}\label{eq:pdfv2-sol-loss-multipliers}
 B_j^+:=M_{\sqrt{b_j^+}},
 \qquad
 B_j^-:=M_{\sqrt{b_j^-}},
 \qquad
 K_j:=M_{\sqrt{2\kappa_j}}.
\end{equation}
Chaque racine est prise point par point sur l'ensemble fini des cylindres ; si une
coordonnée est nulle, son opérateur est nul. L'espace des pertes est la somme
orthogonale externe
\[
 \mathscr E_j
 :=\mathscr H_j^{(+)}\oplus\mathscr H_j^{(-)}
   \oplus\mathscr H_j^{(\kappa)}
   \oplus\mathscr H_j^{({\rm micro})}
   \oplus\mathscr H_j^{({\rm shell})},
\]
et la ligne complète est
\begin{equation}\label{eq:pdfv2-sol-loss-row}
 L_j:\mathscr H_j\longrightarrow\mathscr E_j,
 \qquad
 L_j\xi
 =\bigl(
 B_j^+\xi,B_j^-\xi,K_j\xi,
 Q_j^{\rm micro}\xi,Q_j^{\rm shell}\xi
 \bigr).
\end{equation}
Par conséquent,
\begin{equation}\label{eq:pdfv2-sol-loss-norm}
 \|L_j\xi\|^2
 =\|B_j^+\xi\|^2+\|B_j^-\xi\|^2+\|K_j\xi\|^2
  +\|Q_j^{\rm micro}\xi\|^2
  +\|Q_j^{\rm shell}\xi\|^2.
\end{equation}
Les parties positive, négative, retenue de compensation, microstructure et shell
apparaissent chacune une fois. Elles ne fusionnent pas en une constante anonyme et ne se
répètent pas dans deux termes de la somme.

\section{Terminal d'histoire complète}
\label{sec:pdfv2-sol-terminal}

Soit \(\operatorname{Hist}_{J}^{\rm sol}\) l'ensemble des histoires
survivantes jusqu'à l'horizon \(J\), chacune avec son paquet de treize
coordonnées. L'espace terminal est défini comme l'espace libre de ces
registres complets :
\begin{equation}\label{eq:pdfv2-sol-terminal-space}
 \mathscr T_J^{\rm sol}
 :=\ell^2\!\left(
 \left\{(h,D_{{\rm sol},J}(h)):
 h\in\operatorname{Hist}_{J}^{\rm sol}\right\}
 \right).
\end{equation}
L'évaluation terminale
\begin{equation}\label{eq:pdfv2-sol-terminal-evaluation}
 E_J^{\rm term}:\mathscr H_J\longrightarrow\mathscr T_J^{\rm sol}
\end{equation}
envoie chaque vecteur de base de cylindre au registre complet de l'histoire qui
l'engendre et se prolonge linéairement. En particulier, le terminal retient
fibre, relations, nombres, orientation, retenue, mémoire, frontière, chemin,
multisection, survie et lecteur. Il n'est pas remplacé par la dernière valeur
visible du chemin.

\section{Colonne d'observation, Gram et télescopage}
\label{sec:pdfv2-sol-gram}

Pour \(0\leq j\leq J\), la colonne future complète est
\begin{equation}\label{eq:pdfv2-sol-g-column}
 \begin{aligned}
 G_{j;J}:\mathscr H_j&\longrightarrow
 \mathscr T_J^{\rm sol}\oplus
 \bigoplus_{n=j}^{J-1}\mathscr E_n,\\
 G_{j;J}\xi
 &:={}
 E_J^{\rm term}\Gamma_{j:J}\xi
 \ \oplus\ 
 \bigoplus_{n=j}^{J-1}L_n\Gamma_{j:n}\xi.
 \end{aligned}
\end{equation}
Pour \(j=J\), la somme des pertes est vide et
\(G_{J;J}=E_J^{\rm term}\). Le Gram terminal est
\begin{equation}\label{eq:pdfv2-sol-u-gram}
 U_{j;J}:=G_{j;J}^*G_{j;J}\succeq0.
\end{equation}

En réordonnant la somme orthogonale du codomaine, on obtient l'identité littérale
\begin{equation}\label{eq:pdfv2-sol-g-recursion}
 G_{j;J}\xi
 \simeq
 \bigl(L_j\xi,\,G_{j+1;J}\Gamma_j\xi\bigr).
\end{equation}
En conséquence,
\begin{equation}\label{eq:pdfv2-sol-stein}
 \boxed{U_{j;J}
 =L_j^*L_j+\Gamma_j^*U_{j+1;J}\Gamma_j}
\end{equation}
et
\begin{equation}\label{eq:pdfv2-sol-stein-difference}
 U_{j;J}-\Gamma_j^*U_{j+1;J}\Gamma_j
 =L_j^*L_j\succeq0.
\end{equation}
Il n'y a pas de mutation de signe : le terme de perte est du côté positif
de l'identité parce qu'il provient d'une somme de carrés.

En itérant \eqref{eq:pdfv2-sol-stein}, on obtient le télescopage exact
\begin{equation}\label{eq:pdfv2-sol-telescopia-operatorial}
 U_{j;J}
 =\Gamma_{j:J}^*(E_J^{\rm term})^*E_J^{\rm term}\Gamma_{j:J}
  +\sum_{n=j}^{J-1}
   \Gamma_{j:n}^*L_n^*L_n\Gamma_{j:n},
\end{equation}
ou, sous forme quadratique,
\begin{equation}\label{eq:pdfv2-sol-telescopia-normas}
 \langle U_{j;J}\xi,\xi\rangle
 =\|E_J^{\rm term}\Gamma_{j:J}\xi\|^2
  +\sum_{n=j}^{J-1}\|L_n\Gamma_{j:n}\xi\|^2.
\end{equation}
Le premier terme conserve le futur terminal complet ; la somme conserve,
niveau par niveau, toutes les pertes et toutes les retenues qui y ont conduit.

\iffalse
\section{Les treize coordonnées de la restriction}
\label{sec:pdfv2-sol-trece}

Pour un préfixe \(x_{\leq j}\), on définit
\begin{equation}\label{eq:pdfv2-sol-d-package}
 \begin{aligned}
 D_{{\rm sol},j}(x):=\bigl(&
 \mathcal F_{{\rm sol},j},
 \operatorname{Ext}_{\Gamma,{\rm sol},j},
 \operatorname{Rel}_{{\rm sol},j},
 \mathcal N_{{\rm sol},j},
 \operatorname{or}_{{\rm sol},j},
 \operatorname{Carr}_{{\rm sol},j},
 \operatorname{Mem}_{{\rm sol},j},\\
 &\partial_{{\rm sol},j},
 \gamma_{{\rm sol},j},
 \operatorname{Msect}_{{\rm sol},j},
 \operatorname{Surv}_{{\rm sol},j},
 \operatorname{Term}_{{\rm sol},j},
 \mathcal L_{{\rm sol},j}\bigr).
 \end{aligned}
\end{equation}
Ses composantes sont les suivantes.
\begin{enumerate}
 \item \(\mathcal F_{{\rm sol},j}\) contient les cylindres \(Y_j\), leur mesure
       projective, \(\mathscr H_j^{\rm depth}\),
       \(\mathscr H_j^{\rm res}\), leurs produits tensoriels et les cinq
       espaces de pertes.
 \item \(\operatorname{Ext}_{\Gamma,{\rm sol},j}\) contient
       \(J_{9,j},E_{9,j},\Gamma_j,\Gamma_{j:n}\), les inclusions de shell et
       toutes les applications de troncature et de raffinement.
 \item \(\operatorname{Rel}_{{\rm sol},j}\) contient
       \eqref{eq:pdfv2-sol-bpm-identidades},
       \eqref{eq:pdfv2-sol-ej-projector},
       \eqref{eq:pdfv2-sol-jensen-plus}--
       \eqref{eq:pdfv2-sol-jensen-minus},
       \eqref{eq:pdfv2-sol-ortogonalidad} et
       \eqref{eq:pdfv2-sol-stein}.
 \item \(\mathcal N_{{\rm sol},j}\) conserve
       \(j,J,9,b_j,b_j^+,b_j^-,\kappa_j\), les multiplicités de feuillet, les tailles
       des fibres et les dimensions des shells.
 \item \(\operatorname{or}_{{\rm sol},j}\) conserve
       \(\Delta_{\gamma_j}\), l'ordre du chemin et l'action d'inversion
       \(b_j\mapsto-b_j\), qui échange \(b_j^+\) et \(b_j^-\) sans changer
       \(|b_j|\) ni \(\kappa_j\).
 \item \(\operatorname{Carr}_{{\rm sol},j}\) conserve la retenue TPK entière,
       la retenue de composition et la retenue de compensation \(\kappa_j\) comme
       coordonnées différentes.
 \item \(\operatorname{Mem}_{{\rm sol},j}\) est
       \((M_j^+,M_j^-,D_j^0,H_j^0)\) avec le ledger ordonné de tous les
       accroissements antérieurs.
 \item \(\partial_{{\rm sol},j}\) contient la frontière orientée des cylindres,
       les accroissements \(b_j\), la ligne \(L_j\) et les différences de Stein.
 \item \(\gamma_{{\rm sol},j}\) est le chemin complet
       \(e_0e_1\cdots e_{j-1}\), et non seulement ses extrémités.
 \item \(\operatorname{Msect}_{{\rm sol},j}\) contient toutes les extensions
       nonadiques et de shell compatibles ; il ne choisit pas un prolongement.
 \item \(\operatorname{Surv}_{{\rm sol},j}\) enregistre les horizons
       \(J\geq j\) pour lesquels il existe \(G_{j;J}\), \(U_{j;J}\) et
       des terminaux compatibles.
 \item \(\operatorname{Term}_{{\rm sol},j}\) contient
       \(E_J^{\rm term}\Gamma_{j:J}\), le registre complet d'histoire et
       les treize coordonnées à chaque horizon.
 \item \(\mathcal L_{{\rm sol},j}\) publie en interne
       \((b^\pm,M^\pm,D^0,H^0,\kappa,Q^{\rm micro},Q^{\rm shell},
       L,G,U)\) et ses identités.
\end{enumerate}

Pour une extension \(u:x_j\to x_{j'}\), le transport coordonné
\(u_{{\rm sol},j',j}\) satisfait
\begin{equation}\label{eq:pdfv2-sol-naturality-d}
 u_{{\rm sol},j',j}D_{{\rm sol},j'}(x_{j'})
 =D_{{\rm sol},j}(\tau_{j',j}x_{j'}).
\end{equation}
L'égalité inclut les cinq composantes de \(L\), le terminal complet et
la composition de \(\Gamma\) ; ce n'est pas une égalité des seuls cardinaux.

\section{Graphe de restriction et objet dépendant}
\label{sec:pdfv2-sol-graphs}

On définit
\begin{equation}\label{eq:pdfv2-sol-g-graph}
 \mathcal G_{\rm sol}:=\operatorname{Graph}(D_{\rm sol}),
 \qquad
 \operatorname{Res}_{\rm sol}(x):=(x,D_{\rm sol}x).
\end{equation}
La première projection est inverse sur le graphe :
\begin{equation}\label{eq:pdfv2-sol-res-inverse}
 \pi_1\operatorname{Res}_{\rm sol}=I,
 \qquad
 \operatorname{Res}_{\rm sol}\pi_1=I_{\mathcal G_{\rm sol}}.
\end{equation}

Soit \(\chi_{\rm sol}(x)\) la multisection complète formée de tous les
cylindres, raffinements, shells, lignes de pertes et terminaux compatibles
avec \(x\). Alors
\begin{equation}\label{eq:pdfv2-sol-xchi}
 X_{\chi,{\rm sol}}
 :=\{(\chi_{\rm sol}(x),x,D_{\rm sol}x):
 x\in\mathcal C_{\rm cont}^{\rm sol}\},
\end{equation}
et
\begin{equation}\label{eq:pdfv2-sol-prx}
 \operatorname{pr}_{X,{\rm sol}}(x,D_{\rm sol}x)
 :=(\chi_{\rm sol}(x),x,D_{\rm sol}x).
\end{equation}
De nouveau, la projection qui conserve \((x,D_{\rm sol}x)\) est inverse sur
cette image dépendante.

\section{Assembleur et publicateur internes}
\label{sec:pdfv2-sol-assembler-publisher}

L'espace ambiant \(\mathscr M_{\rm sol}^{\rm amb}\) se compose de systèmes
compatibles
\[
 \bigl(
 \{b,b^\pm,M^\pm,D^0,H^0,\kappa\},
 \{E_9,J_9,P^{\rm term},Q^{\rm micro},Q^{\rm shell}\},
 \{\Gamma,L,E^{\rm term},G,U\}
 \bigr)
\]
avec toutes les identités précédentes. L'assembleur brut est
\begin{equation}\label{eq:pdfv2-sol-assembler}
 \begin{aligned}
 a_{\rm sol}:X_{\chi,{\rm sol}}&\longrightarrow
 \mathscr M_{\rm sol}^{\rm amb},\\
 a_{\rm sol}(\chi(x),x,Dx)&:=
 \bigl(
 \{b,b^\pm,M^\pm,D^0,H^0,\kappa\}_x,
 \{E_9,J_9,P,Q\}_x,
 \{\Gamma,L,E^{\rm term},G,U\}_x
 \bigr).
 \end{aligned}
\end{equation}
La macrostructure conserve son antécédent par
\begin{equation}\label{eq:pdfv2-sol-macro-graph}
 \mathfrak M_{\rm sol}:=\operatorname{Graph}(a_{\rm sol}),
 \qquad
 \mathcal A_{\rm sol}(z):=(z,a_{\rm sol}z).
\end{equation}

Soit \(\mathscr Y_{\rm sol}^{\rm int}\) le type des certificats formés par
les identités
\begin{equation}\label{eq:pdfv2-sol-certificate-list}
 \begin{gathered}
 b=b^+-b^-,\qquad |b|=b^++b^-,\\
 E_9J_9=I,\qquad J_9E_9=P_9,\\
 E_9b_f^\pm=b_c^\pm+\kappa,\qquad\kappa\geq0,\\
 I=P^{\rm term}+Q^{\rm micro}+Q^{\rm shell},
 \qquad Q^{\rm micro}Q^{\rm shell}=0,\\
 U_{j;J}=L_j^*L_j+\Gamma_j^*U_{j+1;J}\Gamma_j,\\
 \langle U_{j;J}\xi,\xi\rangle
 =\|E_J^{\rm term}\Gamma_{j:J}\xi\|^2
  +\sum_{n=j}^{J-1}\|L_n\Gamma_{j:n}\xi\|^2,
 \end{gathered}
\end{equation}
avec la naturalité de toutes les flèches. Le publicateur brut
\begin{equation}\label{eq:pdfv2-sol-publisher}
 p_{\rm sol}^{\rm raw}:\mathfrak M_{\rm sol}
 \longrightarrow\mathscr Y_{\rm sol}^{\rm int}
\end{equation}
publie ce certificat sans supprimer le macro-objet. Enfin,
\begin{equation}\label{eq:pdfv2-sol-output-graph}
 \mathcal C_{\rm sol}^{\rm HMT}
 :=\operatorname{Graph}(p_{\rm sol}^{\rm raw}),
 \qquad
 \mathcal P_{\rm sol}(m):=(m,p_{\rm sol}^{\rm raw}m).
\end{equation}

\section{Théorème interne de génération solénoïdale}
\label{sec:pdfv2-sol-theorem}

\begin{theorem}[Génération interne de la macrostructure de bilan]
\label{thm:pdfv2-sol-internal-generation}
Pour tout \(x\in\mathcal C_{\rm cont}^{\rm sol}\), la chaîne
\[
 \mathcal C_{\rm cont}^{\rm sol}
 \xrightarrow{\operatorname{Res}_{\rm sol}}
 \mathcal G_{\rm sol}
 \xrightarrow{\operatorname{pr}_{X,{\rm sol}}}
 X_{\chi,{\rm sol}}
 \xrightarrow{\mathcal A_{\rm sol}}
 \mathfrak M_{\rm sol}
 \xrightarrow{\mathcal P_{\rm sol}}
 \mathcal C_{\rm sol}^{\rm HMT}
\]
construit naturellement une macrostructure qui contient les deux mémoires,
la retenue de compensation, la séparation orthogonale micro/shell, la ligne complète
de pertes, la colonne future, le terminal d'histoire complète, le Gram et le
télescopage. Les treize coordonnées du continu peuvent être récupérées par
projections successives depuis n'importe quel point de la chaîne.
\end{theorem}

\begin{proof}
Les identités \eqref{eq:pdfv2-sol-bpm-identidades} et
\eqref{eq:pdfv2-sol-dh-diferencias} s'obtiennent directement à partir de la
décomposition positive et négative de l'accroissement entier ; les
mémoires ne sont donc pas des données ajoutées. Les formules
\eqref{eq:pdfv2-sol-j9}--\eqref{eq:pdfv2-sol-e9}, avec la mesure
projective, prouvent \(E_9J_9=I\), \(J_9^*=E_9\) et que \(P_9\) est un
projecteur orthogonal.

L'inégalité triangulaire donne
\(E_9|b_f|\geq|E_9b_f|=|b_c|\), de sorte que
\(\kappa\geq0\). Substituer
\(b^\pm=(|b|\pm b)/2\) prouve
\eqref{eq:pdfv2-sol-jensen-plus} et
\eqref{eq:pdfv2-sol-jensen-minus} sans hypothèse supplémentaire.

Les définitions tensorielles
\eqref{eq:pdfv2-sol-pterm}--\eqref{eq:pdfv2-sol-qshell} prouvent que
\(P^{\rm term}\), \(Q^{\rm micro}\) et \(Q^{\rm shell}\) sont orthogonaux et
ont pour somme l'identité. Par conséquent, la ligne \(L_j\) compte chaque secteur une
seule fois.

La définition \eqref{eq:pdfv2-sol-g-column} sépare orthogonalement le terminal
et les pertes de chaque niveau. Séparer le premier terme de pertes produit
\eqref{eq:pdfv2-sol-g-recursion}. Prendre l'adjoint et composer prouve l'identité
de Stein \eqref{eq:pdfv2-sol-stein} ; l'itérer prouve
\eqref{eq:pdfv2-sol-telescopia-operatorial} et
\eqref{eq:pdfv2-sol-telescopia-normas}. Tout terme est une norme au carré,
de sorte que le signe opposé ne peut apparaître.

La compatibilité de la troncature des cylindres, de l'espérance, des projections et des
extensions \(\Gamma\) donne la naturalité
\eqref{eq:pdfv2-sol-naturality-d}. Enfin,
\(\mathcal G_{\rm sol}\), \(\mathfrak M_{\rm sol}\) et
\(\mathcal C_{\rm sol}^{\rm HMT}\) sont des graphes, et
\(X_{\chi,{\rm sol}}\) est un objet dépendant qui conserve explicitement
son antécédent. La première projection de chaque étape récupère l'étape
précédente. Ainsi aucune composition n'efface l'histoire, la mémoire, la multisection,
le terminal ni le lecteur.
\end{proof}

\section{Provenance probatoire et falsificateurs}
\label{sec:pdfv2-sol-provenance-falsifiers}

\paragraph{Résultat récupéré.}
Appartiennent à l'architecture déjà construite l'accroissement signé avec retenue,
la séparation \(b=b^+-b^-\), les mémoires \(M^\pm,D^0,H^0\), le couple
\(E_9,J_9\), la retenue de Jensen \(\kappa\), l'extension des cylindres, la
ligne de pertes, le terminal, le Gram, Stein et le télescopage.

\paragraph{Formalisation nouvelle.}
Sont des formalisations explicites de cette architecture le typage préalable à toute
évaluation extérieure, la factorisation
\(\mathscr H^{\rm depth}\widehat\otimes\mathscr H^{\rm res}\), la séparation
orthogonale qui empêche le double comptage, le terminal libre sur des histoires à
treize coordonnées et les quatre mises en paquet non ablatives par graphes.

\begin{criterion}[Falsificateurs internes de la branche \(\mathrm{sol}\)]
\label{crit:pdfv2-sol-falsifiers}
La construction échoue si l'un quelconque des faits suivants se produit :
\begin{enumerate}
 \item \(b\neq b^+-b^-\) o \(|b|\neq b^++b^-\);
 \item \(E_9J_9\neq I\), la mesure n'est pas projective ou \(P_9\) n'est pas un
       projecteur ;
 \item \(\kappa<0\) ou
       \(E_9b_f^\pm=b_c^\pm+\kappa\) échouent ;
 \item \(Q^{\rm micro}Q^{\rm shell}\neq0\), ou un même secteur apparaît dans
       les deux termes de la ligne \(L\) ;
 \item l'une de \(b^+,b^-,\kappa,Q^{\rm micro},Q^{\rm shell}\) est omise ou
       comptée deux fois ;
 \item le terminal conserve seulement la valeur visible et non l'histoire avec ses
       treize coordonnées ;
 \item l'identité de Stein échoue ou le télescopage contient un signe
       contraire à une somme de carrés ;
 \item une extension ne commute pas avec \(E_9,J_9,L,G,U\) ;
 \item la multisection des continuations est remplacée par un prolongement
       choisi rétrospectivement ;
 \item l'une quelconque des treize coordonnées est omise.
\end{enumerate}
Dans le dernier cas, le diagnostic contraignant est : \emph{objet substitué ;
conclusion non transférable}.
\end{criterion}

\begin{warningbox}
Ce chapitre conclut uniquement la construction interne de
\(\mathcal C_{\rm sol}^{\rm HMT}\). Il ne contient pas d'équation différentielle
classique, de données initiales classiques, d'application à partir de données extérieures ni de
conclusion sur un problème extérieur. Ces entrelaceurs appartiennent à un autre
document et ne gouvernent pas l'existence de cette macrostructure.
\end{warningbox}
\fi
 \endgroup

\section{Orientation, retenue et composition simultanée des trajectoires}

L'inversion du parcours agit une seule fois sur \(\mathfrak e_\alpha\). Soit \(\mathsf J(v_\Sigma,v_\Pi):=(v_\Pi,v_\Sigma)\) l'échange des feuillets, et soit \(P_{\rm tr}\) la transposition des matrices de ports. Ses cinq expressions locales sont
\begin{equation}
\label{eq:pdfv2-cor-orientacion-cruzada}
 \begin{gathered}
 \mathsf J\sigma_k=-\sigma_k\mathsf J,
 \qquad b_{\bar\alpha}=-b_\alpha,
 \qquad (b_{\bar\alpha}^+,b_{\bar\alpha}^-)
       =(b_\alpha^-,b_\alpha^+),\\
 \vartheta_L(K,\epsilon):=
 \begin{cases}
  (L,-\epsilon),&K=L,\\
  (K,\epsilon),&K\neq L,
 \end{cases}
 \qquad
 \kappa_{J,I}(v,u)=-P_{\rm tr}\kappa_{I,J}(u,v),\\
 \kappa(\bar\alpha)
 =-\mathsf C_\alpha^{-1}\kappa(\alpha),
 \qquad v_{\bar\alpha}=-v_\alpha.
 \end{gathered}
\end{equation}
Ici, \(\mathsf J\) échange les deux feuillets, \(P_{\rm tr}\) transpose la matrice d'incidence et \(\vartheta_L\) échange exactement les deux faces de la direction \(L\) en fixant les six autres. Dans la base libre de \(\mathscr B=\mathscr I\times\{+,-\}\), les huit vecteurs sont distincts, \(\vartheta_L^2=I\) et \(\vartheta_L\neq\vartheta_{L'}\) pour \(L\neq L'\). La même orientation explique tous les signes de \eqref{eq:pdfv2-cor-orientacion-cruzada}.

Pour des extensions composables \(\alpha:x\to y\) et \(\beta:y\to z\), la retenue native satisfait
\begin{equation}
\label{eq:pdfv2-cor-carry-nativo}
 \kappa(\beta\alpha)=
 \kappa(\beta)+\mathsf C(\beta)\kappa(\alpha).
\end{equation}
Le relèvement matriciel ne se réduit pas à un nom. Dans la base ordonnée \((f,r,e,b,g)\), soit \(\widetilde L_N(\alpha)\in\operatorname{Mat}_5(\mathbb Z)\) le représentant entier centré de la matrice de \(\Psi_\alpha\) de \eqref{eq:pdfv2-cor-psi-generadores} modulo \(3^N\), et définissons
\begin{equation}
\label{eq:pdfv2-cor-carry-matricial-def}
 c_N(\beta,\alpha):=
 \frac{\widetilde L_N(\beta)\widetilde L_N(\alpha)
       -\widetilde L_N(\beta\alpha)}{3^N}
 \in\operatorname{Mat}_5(\mathbb Z).
\end{equation}
La divisibilité provient de l'égalité des deux compositions modulo \(3^N\). En écrivant \(L_N=[\widetilde L_N]_{3^N}\), les relèvements entiers satisfont
\begin{equation}
\label{eq:pdfv2-cor-carry-matricial}
 \widetilde L_N(\beta)\widetilde L_N(\alpha)
 =\widetilde L_N(\beta\alpha)+3^Nc_N(\beta,\alpha)
\end{equation}
et l'associativité produit le cocycle
\begin{equation}
\label{eq:pdfv2-cor-cociclo-matricial}
 c_N(\delta,\beta\alpha)+\widetilde L_N(\delta)c_N(\beta,\alpha)
 =c_N(\delta\beta,\alpha)
  +c_N(\delta,\beta)\widetilde L_N(\alpha).
\end{equation}

Le même accroissement de retenue agit sur le complexe local au moyen de
\begin{equation}
\label{eq:pdfv2-cor-psi-generadores}
 \begin{gathered}
 \Psi_\alpha(r_x)=r_y,\quad
 \Psi_\alpha(e_x)=e_y,\quad
 \Psi_\alpha(g_x)=g_y,\quad
 \Psi_\alpha(f_x)=f_y,\\
 \Psi_\alpha(b_x)=b_y+3(c_y-c_x)e_y.
 \end{gathered}
\end{equation}
Le changement de frontière orientée est
\begin{equation}
\label{eq:pdfv2-cor-k-alpha}
 K_\alpha:=(v_y-v_x)f_y,
 \qquad
 K_{\beta\alpha}=K_\beta+\Psi_\beta K_\alpha.
\end{equation}
Les équations \eqref{eq:pdfv2-cor-carry-nativo}, \eqref{eq:pdfv2-cor-cociclo-matricial} et \eqref{eq:pdfv2-cor-k-alpha} sont trois degrés du même registre de composition. Aucune n'autorise à effacer les deux autres.

\subsection{Nerf des prolongements et diagonale d'Alexander--Whitney}
Pour une histoire du registre commun, on écrit
\(A_n:=\mathbf Z/3^n\mathbf Z\) et \(\mathsf T_m(x)\) pour la catégorie
finie de \emph{tous} ses préfixes TPK de profondeur \(m\), avec les
morphismes d’extension déjà construits. La diagonale et la counité qui suivent
s’appliquent à ce même nerf ; elles ne définissent pas une cinquième branche autonome.
\begingroup
  \let\section\subsection
  % Extracción quirúrgica del fuente teoremática V2 05e: líneas 68--110.
% El resto permanece en este archivo como procedencia, pero no se publica.
\iffalse
\chapter[Macrostructure déterminantielle interne]{Macrostructure déterminantielle interne engendrée par le continu discret}
\label{chap:detint-macroestructura}

\begin{thesisbox}
La branche \({\rm det}\) part exclusivement d'histoires survivantes de
\(\mathcal C_{\rm cont}^{\rm disc}\). Le nerf des extensions, la diagonale
Alexander--Whitney, le complexe local, l'unité, les deux publications, le
filler, la réduction terminale et la tour \(3\)-adique se construisent avant
toute évaluation ultérieure. Les conditions de bilan et d'acyclicité sont
incorporées au domaine typé et ne sont jamais tacitement présupposées.
\end{thesisbox}

\section{Domaine déterminantiel survivant}

Soit
\begin{equation}\label{eq:detint-an}
 A_n:=\mathbf Z/3^n\mathbf Z,
 \qquad
 \rho_{n+1,n}:A_{n+1}\longrightarrow A_n.
\end{equation}
Pour \(x\in\mathcal C_{\rm cont}^{\rm disc}\), soit
\(\mathsf T_m(x)\) la catégorie finie des préfixes TPK de profondeur \(m\)
compatibles avec \(x\). Chaque objet conserve feuillet, orientation, résidu,
quotient, retenue entière, mémoire, frontière, chemin, multisection, survie
et terminal ; ses morphismes sont les extensions admissibles qui conservent ces
coordonnées.

De chaque objet \(y\in\mathsf T_m(x)\), on extrait les entiers internes
\begin{equation}\label{eq:detint-c-v-lifts}
 \widetilde c_y,\widetilde v_y\in\mathbf Z,
 \qquad
 c_{y,n}:=\widetilde c_y\bmod3^n,
 \qquad
 v_{y,n}:=\widetilde v_y\bmod3^n.
\end{equation}
La première coordonnée est la retenue non orientée et la seconde, le coefficient
orienté de frontière. Le renversement satisfait
\begin{equation}\label{eq:detint-or-cv}
 c_{\bar y,n}=c_{y,n},
 \qquad
 v_{\bar y,n}=-v_{y,n}.
\end{equation}

On écrit
\begin{equation}\label{eq:detint-z3-k3}
 \mathbf Z_3:=\varprojlim_n A_n,
 \qquad
 K_3:=\operatorname{Frac}(\mathbf Z_3).
\end{equation}
La restriction déterminantielle est définie par
\begin{equation}\label{eq:detint-dominio}
\begin{split}
 \mathcal C_{\rm cont}^{\rm det}:=
 \bigl\{x\in\mathcal C_{\rm cont}^{\rm disc}:{}&
 x\text{ survit à toute profondeur};\\
 &\widehat P_x^{\rm red}\text{ est fini et basé};\\
 &\chi(\widehat P_x^{\rm red})=0;\\
 &H_*(\widehat P_x^{\rm red}\otimes_{\mathbf Z_3}K_3)=0;\\
 &\text{la réduction canonique produit un bloc carré}\\
 &\text{stable par raffinement}\bigr\}.
\end{split}
\end{equation}
Les objets \(\widehat P_x^{\rm red}\) et le bloc carré seront construits
dans \cref{sec:detint-terminal}. Par conséquent,
\eqref{eq:detint-dominio} est une condition d'appartenance vérifiable et non
une conclusion insérée dans la définition des applications.

\fi
\section{Nerf des chemins, diagonale d'Alexander--Whitney et counité}

On définit le complexe normalisé
\begin{equation}\label{eq:detint-gmn}
 G_{m,n}(x):=
 N_*^{\rm norm}\!\left(N\mathsf T_m(x);A_n\right).
\end{equation}
Pour un simplexe non dégénéré
\(\sigma=[y_0\to y_1\to\cdots\to y_q]\), la diagonale est
\begin{equation}\label{eq:detint-aw}
 \Delta_{\rm AW}(\sigma)
 =\sum_{j=0}^q
 [y_0\to\cdots\to y_j]\otimes
 [y_j\to\cdots\to y_q].
\end{equation}
La counité est donnée par
\begin{equation}\label{eq:detint-counidad-g}
 \epsilon_G([y])=1,
 \qquad
 \epsilon_G(\sigma)=0\quad(q>0).
\end{equation}
En particulier,
\begin{equation}\label{eq:detint-vertice-grouplike}
 \Delta_{\rm AW}[y]=[y]\otimes[y],
 \qquad
 \epsilon_G([y])=1.
\end{equation}

Si \(m'\ge m\) et \(n'\ge n\), la troncature et la réduction des coefficients
induisent
\begin{equation}\label{eq:detint-rmn}
 R_{m',m}^{n',n}:G_{m',n'}(x)\longrightarrow G_{m,n}(x),
\end{equation}
et l'on vérifie
\begin{equation}\label{eq:detint-aw-natural}
 R\partial=\partial R,
 \qquad
 (R\otimes R)\Delta_{\rm AW}=\Delta_{\rm AW}R,
 \qquad
 \epsilon_GR=\rho_{n',n}\epsilon_G.
\end{equation}
Le chemin ne se réduit pas à ses extrémités : le simplexe entier reste dans
\(G_{m,n}(x)\).

\iffalse
\section{Complexe local interne}

Pour \(c\in A_n\), on définit
\begin{equation}\label{eq:detint-p0-grados}
 P^0_{n,1}(c)=A_nf,
 \qquad
 P^0_{n,0}(c)=A_nr\oplus A_ne\oplus A_nb,
 \qquad
 P^0_{n,-1}(c)=A_ng,
\end{equation}
avec différentielle
\begin{equation}\label{eq:detint-p0-diferencial}
 \partial f=3ce+b,
 \qquad
 \partial r=0,
 \qquad
 \partial e=g,
 \qquad
 \partial b=-3cg.
\end{equation}
L'identité de chaîne est littérale :
\begin{equation}\label{eq:detint-d2}
 \partial^2f=3c\,\partial e+\partial b=3cg-3cg=0,
 \qquad
 \partial^2e=\partial^2b=\partial^2r=0.
\end{equation}

L'augmentation
\begin{equation}\label{eq:detint-epsilon-p}
 \epsilon_P:P_n^0(c)\longrightarrow A_n[0]
\end{equation}
est fixée par
\[
 \epsilon_P(r)=1,
 \qquad
 \epsilon_P(e)=\epsilon_P(b)=\epsilon_P(f)=\epsilon_P(g)=0.
\]
Le complexe possède la coalgèbre différentielle
\begin{equation}\label{eq:detint-coalgebra-p}
 \Delta_P(r)=r\otimes r,
 \qquad
 \Delta_P(q)=q\otimes r+r\otimes q
 \quad(q\in\{e,b,f,g\}).
\end{equation}
Ainsi, \(r\) est group-like et les quatre autres générateurs sont primitifs
relativement à \(r\). Les équations
\eqref{eq:detint-p0-diferencial} prouvent que \(\Delta_P\) commute avec la
différentielle tensorielle.

\section{Système local d'extensions et ledger de mémoire}

Pour une extension \(\alpha:y\to y'\), nous définissons
\begin{equation}\label{eq:detint-psi}
 \Psi_\alpha:P_n^0(c_y)\longrightarrow P_n^0(c_{y'})
\end{equation}
par
\begin{equation}\label{eq:detint-psi-generadores}
\begin{gathered}
 \Psi_\alpha(r_y)=r_{y'},\quad
 \Psi_\alpha(e_y)=e_{y'},\quad
 \Psi_\alpha(g_y)=g_{y'},\quad
 \Psi_\alpha(f_y)=f_{y'},\\
 \Psi_\alpha(b_y)=b_{y'}+3(c_{y'}-c_y)e_{y'}.
\end{gathered}
\end{equation}
C'est une application de complexes parce que
\begin{equation}\label{eq:detint-psi-b-chain}
 \partial\Psi_\alpha(b_y)
 =-3c_{y'}g+3(c_{y'}-c_y)g
 =-3c_yg
 =\Psi_\alpha(\partial b_y),
\end{equation}
et
\begin{equation}\label{eq:detint-psi-f-chain}
 \Psi_\alpha(\partial f_y)
 =3c_ye+b+3(c_{y'}-c_y)e
 =3c_{y'}e+b
 =\partial\Psi_\alpha(f_y).
\end{equation}
De plus,
\begin{equation}\label{eq:detint-psi-functorial}
 \Psi_\beta\Psi_\alpha=\Psi_{\beta\alpha}.
\end{equation}

Le changement de la coordonnée orientée n'est pas effacé. Il est enregistré par
\begin{equation}\label{eq:detint-k-alpha}
 K_\alpha:=(v_{y'}-v_y)f_{y'}.
\end{equation}
Pour des extensions composables,
\begin{equation}\label{eq:detint-k-coherencia}
 K_{\beta\alpha}=K_\beta+\Psi_\beta K_\alpha.
\end{equation}
Cette égalité est le ledger exact de l'accroissement de frontière et rend visible
la mémoire qu'une composition purement nominale perdrait.

Le diagramme local complet est la construction de Grothendieck
\begin{equation}\label{eq:detint-grothendieck-local}
 \int_{y\in\mathsf T_m(x)}P_n^0(c_y).
\end{equation}
Dans celle-ci, une flèche au-dessus de \(\alpha:y\to y'\) porte comme composante linéaire
\(\Psi_\alpha\) et comme cohérence \(K_\alpha\). Lorsqu'elle est appliquée à un
élément supporté en un sommet, on utilise la notation
\begin{equation}\label{eq:detint-psi-soporte-vertice}
 \Psi_\alpha(p,[y]):=(\Psi_\alpha p,[y']).
\end{equation}
L'arête \([y\to y']\) reste dans la coordonnée de chemin du nerf. Cette
convention sera celle utilisée dans les formules de naturalité de
\(z_\pm\) et \(h_*\) ; on n'affirme pas une application globale qui efface les autres
simplexes de \(G_{m,n}(x)\).

\section{Pullback, unité, racine, publications et filler}

Le pullback de complexes est
\begin{equation}\label{eq:detint-pgen}
 P^{\rm gen}_{y;m,n}
 :=P_n^0(c_y)\times_{A_n[0]}G_{m,n}(x)
 =\{(p,\gamma):\epsilon_P(p)=\epsilon_G(\gamma)\}.
\end{equation}
Sa différentielle est
\begin{equation}\label{eq:detint-pgen-d}
 \partial(p,\gamma)=(\partial p,\partial\gamma).
\end{equation}
L'unité commune associée au sommet \(y\) est
\begin{equation}\label{eq:detint-unidad}
 \mathbf1_y:=(r,[y]).
\end{equation}
Ses deux projections sont group-like par
\eqref{eq:detint-vertice-grouplike} et
\eqref{eq:detint-coalgebra-p} ; c'est le sens précis d'unité
group-like compatible du pullback.

La projection racine est l'application de complexes
\begin{equation}\label{eq:detint-proyeccion-root}
 \varpi_{\rm root}:P^{\rm gen}_{y;m,n}\longrightarrow
 \operatorname{Root}_{y;m,n},
 \qquad
 \varpi_{\rm root}(p,\gamma):=(\epsilon_P(p)r,\gamma).
\end{equation}
On définit les deux publications internes
\begin{equation}\label{eq:detint-zpm}
 z_-(y):=(r,[y]),
 \qquad
 z_+(y):=(r+3c_yv_ye+v_yb,[y]).
\end{equation}
Toutes deux sont des cycles :
\begin{equation}\label{eq:detint-zpm-ciclos}
 \partial z_-(y)=0,
 \qquad
 \partial z_+(y)=3c_yv_yg-3c_yv_yg=0.
\end{equation}
Leurs racines coïncident :
\begin{equation}\label{eq:detint-raiz-comun}
 \varpi_{\rm root}z_-(y)
 =(r,[y])
 =\varpi_{\rm root}z_+(y).
\end{equation}

Le filler se construit directement :
\begin{equation}\label{eq:detint-hstar}
 h_*(y):=(-v_yf,0)\in P^{\rm gen}_{y;m,n,1}.
\end{equation}
Alors
\begin{equation}\label{eq:detint-filler-identidad}
 \partial h_*(y)
 =(-3c_yv_ye-v_yb,0)
 =z_-(y)-z_+(y).
\end{equation}
En particulier,
\begin{equation}\label{eq:detint-clases-iguales}
 [z_-(y)]=[z_+(y)]
 \quad\text{dans}\quad H_0(P^{\rm gen}_{y;m,n}),
\end{equation}
et l'égalité conserve sa chaîne témoin.

Sous une extension \(\alpha:y\to y'\), on a
\begin{equation}\label{eq:detint-z-natural}
 z_-(y')=\Psi_\alpha z_-(y),
 \qquad
 z_+(y')-\Psi_\alpha z_+(y)=\partial K_\alpha,
\end{equation}
et
\begin{equation}\label{eq:detint-h-natural}
 h_*(y')-\Psi_\alpha h_*(y)=-K_\alpha.
\end{equation}
Par conséquent, la naturalité n'identifie pas des valeurs distinctes de \(v\) : elle enregistre
leur différence au moyen de l'homotopie cohérente \(K_\alpha\).

\section{Orientation interne}

Sur \(P_n^0(c)\), on définit
\begin{equation}\label{eq:detint-omega-p}
 \Omega(r)=r,
 \qquad
 \Omega(q)=-q\quad(q\in\{e,b,f,g\}).
\end{equation}
C'est un automorphisme de complexes et de la coalgèbre relative. Dans le nerf,
l'inversion orientée est
\begin{equation}\label{eq:detint-omega-g}
 \mathfrak o_G[y_0\to\cdots\to y_q]
 :=(-1)^{q(q+1)/2}
 [\bar y_q\to\cdots\to\bar y_0].
\end{equation}
Avec \eqref{eq:detint-or-cv}, ces actions donnent
\begin{equation}\label{eq:detint-or-z-h}
 \Omega z_-(y)=z_-(\bar y),
 \qquad
 \Omega z_+(c_y,v_y)=z_+(c_{\bar y},v_{\bar y}),
 \qquad
 \Omega h_*(y)=h_*(\bar y).
\end{equation}
L'orientation agit sur le chemin, le complexe, les publications, le filler et la droite
déterminantielle ; ce n'est pas une étiquette sans action.

\section{Réduction terminale et condition de survie}
\label{sec:detint-terminal}

Seule la racine commune est supprimée :
\begin{equation}\label{eq:detint-p-reducido}
 \overline P_{y;m,n}
 :=P^{\rm gen}_{y;m,n}/A_n\mathbf1_y,
 \qquad
 \widehat P_{y;m}^{\rm red}:=\varprojlim_n\overline P_{y;m,n}.
\end{equation}
La base est ordonnée d'abord selon l'ordre TPK des chemins, puis selon
\begin{equation}\label{eq:detint-orden-base}
 f\prec e\prec b\prec g.
\end{equation}
On exécute l'algorithme interne suivant.
\begin{enumerate}
 \item Chercher les coefficients qui sont des unités de \(\mathbf Z_3\).
 \item Choisir le premier selon \eqref{eq:detint-orden-base} et l'ordre des
 chemins.
 \item Annuler le couple correspondant et enregistrer les opérations
 élémentaires, leur signe et leur contribution à l'orientation.
 \item Répéter jusqu'à ce qu'il ne reste aucun pivot unitaire annulable.
\end{enumerate}
La condition de survie de \eqref{eq:detint-dominio} exige que,
cofinalement en \(m\), le résultat soit un bloc à deux termes
\begin{equation}\label{eq:detint-sx}
 \mathbf Z_3^{r_x}\xrightarrow{S_x}\mathbf Z_3^{r_x},
 \qquad
 \det S_x\ne0\text{ dans }K_3,
\end{equation}
et qu'un raffinement n'ajoute que des blocs contractiles unitaires et des changements
de base enregistrés. Si ces conditions échouent, l'histoire n'appartient pas à
\(\mathcal C_{\rm cont}^{\rm det}\) ; on ne force pas un terminal inexistant.

Le chemin identité fournit un témoin élémentaire. Après le retrait de la racine,
le complexe local est
\begin{equation}\label{eq:detint-testigo-contractil}
 A_nf\xrightarrow{\binom{3c}{1}}
 A_ne\oplus A_nb\xrightarrow{(1,-3c)}A_ng.
\end{equation}
Il est contractile par
\begin{equation}\label{eq:detint-contraccion-local}
 s(g)=e,
 \qquad
 s(ae+\beta b)=\beta f,
 \qquad
 s(f)=0.
\end{equation}
Ainsi, les conditions terminales admettent au moins le témoin identité
interne et ne décrivent pas un sous-domaine formellement vide.

\section{Smith, ordre \texorpdfstring{\(3\)}{3}-adique et unité initiale}

Dans des bases orientées, une diagonalisation par valuation a la forme
\begin{equation}\label{eq:detint-smith}
 U_xS_xV_x
 =\operatorname{diag}(3^{a_1}u_1,\ldots,3^{a_r}u_r),
 \qquad
 a_i\in\mathbf N,quad u_i\in\mathbf Z_3^\times.
\end{equation}
On définit
\begin{equation}\label{eq:detint-orden-d}
 d_x:=v_3(\det S_x)=\sum_{i=1}^r a_i
\end{equation}
et l'unité initiale
\begin{equation}\label{eq:detint-leading-unit}
 \lambda_x:=3^{-d_x}\det S_x\in\mathbf Z_3^\times.
\end{equation}
La coordonnée d'orientation trivialise la droite déterminantielle. Si une
transition produit
\begin{equation}\label{eq:detint-estabilidad-bloques}
 S_x\oplus I_q=U S_{x'}V,
\end{equation}
l'orientation est transportée par le facteur inverse de \(\det U\det V\).
L'unité orientée \(\lambda_x^{\rm or}\) est alors naturelle et ne change pas
lors de l'ajout de blocs contractiles.

En précision finie,
\begin{equation}\label{eq:detint-smith-finito}
 S_{x,n}=S_x\bmod3^n,
 \qquad
 a_i^{(n)}=\min(a_i,n).
\end{equation}
Le terminal est donc la famille compatible de toutes les précisions et
non un entier isolé.

\section{Tour de Bockstein et retenue d'ordre supérieur}

Soit
\[
 C_x^1\xrightarrow{S_x}C_x^0
\]
le bloc terminal. Si une classe \([\bar y]\) admet un relèvement
\(y\in C_x^1\) avec \(S_xy\in3^qC_x^0\), nous définissons
\begin{equation}\label{eq:detint-beta-q}
 \beta_q[\bar y]
 :=\left[3^{-q}S_xy\bmod3\right].
\end{equation}
Le changement de relèvement modifie le représentant par un bord de la page
correspondante, de sorte que \(\beta_q\) est bien défini. La retenue
supérieure est
\begin{equation}\label{eq:detint-carry-superior}
 \operatorname{Carr}_q(y):=3^{-q}S_xy.
\end{equation}
Elle ne remplace pas la retenue TPK de \eqref{eq:detint-c-v-lifts} ; elle enregistre son
prolongement dans la tour des coefficients.

Pour un bloc \(3^{a_i}u_i\), la classe persiste jusqu'à la page \(a_i\),
et à cette page, la différentielle non nulle est la multiplication par
\(\bar u_i\in\mathbf F_3^\times\). Par conséquent,
\begin{equation}\label{eq:detint-bock-longitudes}
 \{\text{longueurs de survie}\}=\{a_1,\ldots,a_r\},
 \qquad
 d_x=\sum_i a_i.
\end{equation}
Le produit orienté des trivialisations de page satisfait
\begin{equation}\label{eq:detint-bock-leading}
 \Theta_{\rm Bock}(x)=\overline{\lambda_x^{\rm or}}
 \in\mathbf F_3^\times.
\end{equation}
La preuve se fait dans chaque bloc diagonal et se transporte par les changements
unimodulaires enregistrés. Smith, Bockstein et l'unité initiale sont ainsi trois
lectures du même terminal.

\section{Les treize coordonnées de la restriction déterminantielle}

Pour chaque \((m,n)\), on définit
\begin{equation}\label{eq:detint-d-trece}
 D_{{\rm det},m,n}(x):=
 (\mathcal F,\operatorname{Ext}_\Gamma,\operatorname{Rel},
 \mathcal N,\operatorname{or},\operatorname{Carr},\operatorname{Mem},
 \partial,\gamma,\operatorname{Msect},\operatorname{Surv},
 \operatorname{Term},\mathcal L)_{{\rm det},m,n}(x),
\end{equation}
avec le contenu suivant.
\begin{enumerate}
 \item La fibre est
 \[
  \mathcal F=(A_n,\mathsf T_m(x),
  \{P_n^0(c_y)\}_y,\{P^{\rm gen}_{y;m,n}\}_y).
 \]
 \item \(\operatorname{Ext}_\Gamma\) contient tous les morphismes
 \(\alpha\), les applications \(\Psi_\alpha\), les homotopies \(K_\alpha\) et
 les applications \(R_{m',m}^{n',n}\).
 \item \(\operatorname{Rel}\) contient \(\partial^2=0\), l'équation de
 pullback, AW, la counité, \(\Psi_\beta\Psi_\alpha=\Psi_{\beta\alpha}\) et
 \(K_{\beta\alpha}=K_\beta+\Psi_\beta K_\alpha\).
 \item
 \(\mathcal N=(m,n,\widetilde c_y,\widetilde v_y,c_{y,n},v_{y,n},
 r_x,a_1,\ldots,a_{r_x},d_x)\).
 \item \(\operatorname{or}\) contient \(y\mapsto\bar y\),
 \(\mathfrak o_G\), \(\Omega\) et l'orientation de la droite
 déterminantielle.
 \item \(\operatorname{Carr}\) conserve séparément la retenue entière,
 sa réduction \(c_{y,n}\) et toutes les retenues supérieures
 \(\operatorname{Carr}_q\).
 \item \(\operatorname{Mem}\) contiene
 \[
  y_0\to\cdots\to y_m,
  \quad\{(\widetilde c_{y_j},\widetilde v_{y_j},\mathbf1_{y_j})\}_{j=0}^m,
  \quad\{K_{\alpha_j}\}.
 \]
 \item
 \(\partial=(\partial_G,\partial_P,\partial_{P^{\rm gen}},
 \partial_{\overline P},h_*)\).
 \item \(\gamma=[y_0\to\cdots\to y_q]\) conserve le simplexe orienté
 complet.
 \item \(\operatorname{Msect}\) contient toutes les extensions
 compatibles, tous les relèvements \(3\)-adiques et toutes les présentations
 unimodulaires du même terminal orienté ; il n'en choisit aucune a posteriori.
 \item \(\operatorname{Surv}\) enregistre la profondeur arbitraire,
 la compatibilité des coefficients, le bilan, l'acyclicité sur \(K_3\),
 la stabilité par blocs unitaires et la permanence de la racine.
 \item
 \[
  \operatorname{Term}=(\widehat P_x^{\rm red},S_x,
  \{a_i,u_i\},d_x,\lambda_x^{\rm or},
  \{\beta_q\},\Theta_{\rm Bock}).
 \]
 \item Le lecteur interne est
 \[
  \mathcal L_{\rm det}^{\rm int}
  =(\mathbf1_y,
  \varpi_{\rm root}z_-=\varpi_{\rm root}z_+,
  [z_-]=[z_+],d_x,\lambda_x^{\rm or},\Theta_{\rm Bock}).
 \]
\end{enumerate}

Les transitions satisfont, coordonnée par coordonnée,
\begin{equation}\label{eq:detint-naturalidad-trece}
 u^{\rm det}_{(m',n'),(m,n)}D_{{\rm det},m',n'}(x')
 =D_{{\rm det},m,n}(\tau_{m',m}x').
\end{equation}

\section{Assembleur, publicateur et chaîne non ablative}

On définit
\begin{equation}\label{eq:detint-graph-d}
 \mathcal G_{\rm det}:=\operatorname{Graph}(D_{\rm det}),
 \qquad
 \operatorname{Res}_{\rm det}(x):=(x,D_{\rm det}x).
\end{equation}
Soit \(\chi_{\rm det}(x)\) la multisection complète de germes
\[
 (y,\widetilde c_y,\widetilde v_y,\operatorname{or}_y,\gamma_y)
\]
compatibles à toute profondeur. L'objet dépendant est
\begin{equation}\label{eq:detint-x-dependiente}
 X_{\chi,{\rm det}}
 :=\{(\chi_{\rm det}(x),x,D_{\rm det}x):
 x\in\mathcal C_{\rm cont}^{\rm det}\},
\end{equation}
et
\begin{equation}\label{eq:detint-prx}
 \operatorname{pr}_{X,{\rm det}}(x,D_{\rm det}x)
 :=(\chi_{\rm det}(x),x,D_{\rm det}x).
\end{equation}

L'assembleur brut
\begin{equation}\label{eq:detint-ensamblador-crudo}
 a_{\rm det}:X_{\chi,{\rm det}}\longrightarrow
 \mathscr M_{\rm det}^{\rm amb}
\end{equation}
produit
\begin{equation}\label{eq:detint-ensamblador-salida}
\begin{split}
 a_{\rm det}(z)=\bigl(&
 \{G_{m,n},\Delta_{\rm AW},\epsilon_G\}_{m,n};\\
 &\{P_n^0(c),\Psi_\alpha,K_\alpha,P^{\rm gen},
 \mathbf1,z_-,z_+,h_*\}_{m,n};\\
 &\{\widehat P^{\rm red},S,d,\lambda^{\rm or},
 \beta_q,\Theta_{\rm Bock}\}\bigr).
\end{split}
\end{equation}
Le macro-objet conserve son entrée :
\begin{equation}\label{eq:detint-graph-a}
 \mathfrak M_{\rm det}:=\operatorname{Graph}(a_{\rm det}),
 \qquad
 \mathcal A_{\rm det}(z):=(z,a_{\rm det}z).
\end{equation}

Le publicateur brut
\begin{equation}\label{eq:detint-publicador-crudo}
 p_{\rm det}:\mathfrak M_{\rm det}\longrightarrow\mathscr Y_{\rm det}
\end{equation}
publie conjointement le certificat
\begin{equation}\label{eq:detint-certificado-publicado}
\begin{gathered}
 \partial^2=0,quad \text{AW et counité},quad
 \mathbf1_y\text{ group-like compatible},\\
 \partial z_\pm=0,quad
 \varpi_{\rm root}z_-=\varpi_{\rm root}z_+,quad
 \partial h_*=z_--z_+,\\
 \Psi_\beta\Psi_\alpha=\Psi_{\beta\alpha},quad
 K_{\beta\alpha}=K_\beta+\Psi_\beta K_\alpha,\\
 d=v_3(\det S)=\sum_i a_i,quad
 \Theta_{\rm Bock}=\overline{\lambda^{\rm or}},quad
 \text{naturalité complète par raffinement}.
\end{gathered}
\end{equation}
Enfin,
\begin{equation}\label{eq:detint-graph-p}
 \mathcal C_{\rm det}^{\rm HMT}:=\operatorname{Graph}(p_{\rm det}),
 \qquad
 \mathcal P_{\rm det}(M):=(M,p_{\rm det}M).
\end{equation}
La chaîne intégrale est
\begin{equation}\label{eq:detint-cadena-completa}
 \mathcal C_{\rm cont}^{\rm disc}\supset
 \mathcal C_{\rm cont}^{\rm det}
 \xrightarrow{\operatorname{Res}_{\rm det}}\mathcal G_{\rm det}
 \xrightarrow{\operatorname{pr}_{X,{\rm det}}}X_{\chi,{\rm det}}
 \xrightarrow{\mathcal A_{\rm det}}\mathfrak M_{\rm det}
 \xrightarrow{\mathcal P_{\rm det}}\mathcal C_{\rm det}^{\rm HMT}.
\end{equation}
Chaque flèche possède un inverse sur son image donné par une projection de
coordonnées. On ne remplace pas l'histoire par une valeur, la multisection par un sélecteur,
le complexe par un déterminant ni la preuve par un nom.

\section{Théorème de réalisation déterminantielle interne}

\begin{theorem}[Réalisation déterminantielle interne complète]
\label{thm:detint-realizacion}
Pour tout \(x\in\mathcal C_{\rm cont}^{\rm det}\), la chaîne
\eqref{eq:detint-cadena-completa} construit un macro-objet naturel en
profondeur et en précision. Il contient deux cycles à racine commune, un filler
explicite qui les identifie, une unité group-like compatible et un terminal
\(3\)-adique dont l'ordre, la forme de Smith et la tour de Bockstein produisent le même
lecteur orienté. Les treize coordonnées de \eqref{eq:detint-d-trece}
restent récupérables en sortie.
\end{theorem}

\begin{proof}
Les équations \eqref{eq:detint-p0-diferencial}--
\eqref{eq:detint-d2} prouvent que \(P_n^0(c)\) est un complexe. Les deux
augmentations sont des applications de complexes, de sorte que le pullback est typé.
Alexander--Whitney est naturel et rend chaque sommet group-like ; avec
\(r\), cela donne l'unité commune. Le calcul direct
\eqref{eq:detint-zpm-ciclos}--\eqref{eq:detint-filler-identidad} prouve que
les publications sont des cycles, ont la même racine et sont reliées par
\(h_*\). Les formules de \(\Psi_\alpha\) prouvent la fonctorialité, et
\(K_\alpha\) prouve la naturalité homotopique en conservant l'accroissement de
mémoire. La définition de \(\mathcal C_{\rm cont}^{\rm det}\) garantit que
la réduction terminale produit un bloc carré sans le postuler hors du
domaine. Smith sur \(\mathbf Z_3\) donne \(d=\sum_i a_i\) ; le calcul de chaque
bloc \(3^{a_i}u_i\) donne
\(\Theta_{\rm Bock}=\overline{\lambda^{\rm or}}\). Troncature, réduction et
ajout de blocs unitaires conservent ces identités. Enfin, les
trois constructions par graphe retiennent l'entrée comme première coordonnée,
de sorte que la composition complète est non ablative.
\end{proof}

\section{Provenance et falsificateurs}

\paragraph{Provenance.}
Le nerf TPK, Alexander--Whitney, la counité, l'unité commune, la forme
normale \(\partial f=3ce+b\), \(\partial e=g\),
\(\partial b=-3cg\), la racine commune, le filler et l'architecture de tour
ont le statut \texttt{ARCHITECTURE\_AUTORALE\_PRÉEXISTANTE}. Le calcul
littéral des deux cycles, de leur racine et de l'identité
\(\partial h_*=z_--z_+\) a le statut \texttt{RÉSULTAT\_RÉCUPÉRÉ}. Le
système local \(\Psi_\alpha/K_\alpha\), le sous-domaine survivant
équilibré, le terminal exclusivement \(3\)-adique et la mise en paquet par
graphes ont le statut \texttt{FORMALISATION\_NOUVELLE} ; ils formalisent sans changer
la provenance conceptuelle de l'appareil.

\begin{ablation}[Falsificateurs internes de la branche déterminantielle]
\label{abl:detint-falsadores}
La construction est falsifiée au niveau correspondant si l'un des
faits suivants se produit.
\begin{enumerate}
 \item \(\partial^2=0\) échoue ; le complexe local n'existe alors pas.
 \item Alexander--Whitney ne commute pas avec le raffinement, la counité échoue ou
 le sommet cesse d'être group-like.
 \item \(z_\pm\) ne sont pas des cycles, leurs racines diffèrent ou
 \(\partial h_*=z_--z_+\) échoue.
 \item \(\Psi_\alpha\) n'est pas une application de complexes ou
 \(K_{\beta\alpha}=K_\beta+\Psi_\beta K_\alpha\) échoue ; dans ce cas, on a
 perdu de la retenue ou de la mémoire.
 \item Le complexe réduit n'est pas équilibré ou n'est pas acyclique sur
 \(K_3\) ; l'histoire reste hors du domaine terminal typé.
 \item Les longueurs de Bockstein ne coïncident pas avec les \(a_i\), ou
 \(\Theta_{\rm Bock}\ne\overline{\lambda^{\rm or}}\).
 \item Le raffinement modifie \(d\) ou \(\lambda^{\rm or}\) en dehors de
 blocs unitaires et de changements d'orientation enregistrés.
 \item L'une des treize coordonnées est omise : \emph{objet substitué ;
 conclusion non transférable}.
\end{enumerate}
Ces falsificateurs vérifient l'objet interne construit ; ils ne le remplacent pas
par un contrôle ultérieur.
\end{ablation}
\fi
 \endgroup

\section{Incidence cellulaire et remplissage corrélatif d'une même trajectoire}

La suite \eqref{eq:pdfv2-cor-sucesion-celda} associe à chaque préfixe le complexe à deux termes
\begin{equation}
\label{eq:pdfv2-cor-cono-celda}
 C_{k,N}(\widehat x):=
 \left[A_N^{I_k}\oplus A_N^{J_k}
 \xrightarrow{\kappa_{k,N}}A_N^{I_k\times J_k}\right].
\end{equation}
Un raffinement agit sur les ports par troncature de l'extension et sur les coefficients par \(A_{N'}\twoheadrightarrow A_N\). Lorsque \(\alpha_u:I_{k'}\to I_k\) et \(\beta_u:J_{k'}\to J_k\) sont ces applications, leur action sur les générateurs est
\begin{equation}
\label{eq:pdfv2-cor-refinamiento-celda}
 (p,q)\longmapsto(p\circ\alpha_u,q\circ\beta_u),
 \qquad
 w\longmapsto w\circ(\alpha_u\times\beta_u),
\end{equation}
et vérifie
\begin{equation}
\label{eq:pdfv2-cor-refinamiento-cadena}
 \kappa_{k',N}\Gamma(u)_1=\Gamma(u)_0\kappa_{k,N}.
\end{equation}

L'incidence n'est pas fabriquée en étiquetant quatre copies d'une émission neutre. Elle est engendrée par la complétion cellulaire libre du module des deux feuillets que porte déjà le même état. Pour
\[
 V_{\widehat x}:=\mathbf F_3e_{\Sigma,\widehat x}
 \oplus\mathbf F_3e_{\Pi,\widehat x},
\]
les quatre droites internes sont
\begin{equation}
\label{eq:pdfv2-cor-cuatro-rectas-tpk}
 \mathscr I_{\widehat x}:=
 \{L_0,L_\infty,L_+,L_-\}
 =\mathbf P^1(V_{\widehat x}),\qquad
 \begin{aligned}
 L_0&=\langle e_\Sigma\rangle,&
 L_\infty&=\langle e_\Pi\rangle,\\
 L_+&=\langle e_\Sigma+e_\Pi\rangle,&
 L_-&=\langle e_\Sigma-e_\Pi\rangle.
 \end{aligned}
\end{equation}
Si \(v_L\) est respectivement \((1,0)^t,(0,1)^t,(1,1)^t,(1,-1)^t\), l'action de la direction \(L\) sur les générateurs de feuillet est le projecteur
\begin{equation}
\label{eq:pdfv2-cor-proyectores-cuatro-direcciones}
 P_L:=v_L(v_L^tv_L)^{-1}v_L^t,
\end{equation}
c'est-à-dire, sur \(\mathbf F_3\),
\[
 P_0=\begin{pmatrix}1&0\\0&0\end{pmatrix},\quad
 P_\infty=\begin{pmatrix}0&0\\0&1\end{pmatrix},\quad
 P_+=\begin{pmatrix}2&2\\2&2\end{pmatrix},\quad
 P_-=\begin{pmatrix}2&1\\1&2\end{pmatrix}.
\]
Le calcul direct donne
\begin{equation}
\label{eq:pdfv2-cor-proyectores-no-colapso}
 P_L^2=P_L,\qquad \operatorname{rk}P_L=1,\qquad
 \operatorname{Im}P_L=L,\qquad
 P_L=P_{L'}\Longrightarrow L=L'.
\end{equation}
Par conséquent, les quatre directions sont des sorties distinctes du module \(V_{\widehat x}\) ; elles ne sont pas quatre noms attribués à la même action.

\begin{definition}[Complétion cellulaire libre des feuillets TPK]
\label{def:pdfv2-cor-completacion-celular-tpk}
Soit \(\operatorname{Cell}_{\rm TPK}(\widehat x)\) le complexe libre
\[
 C_2(\widehat x)\xrightarrow{\partial_2}
 C_1(\widehat x)\xrightarrow{\partial_1}C_0(\widehat x)
\]
avec
\begin{align*}
 C_0(\widehat x)
 &:=
 \mathbf F_3[v_\star]\oplus
 \bigoplus_{L\in\mathscr I}\mathbf F_3[v_L]\oplus
 \bigoplus_{(L,\epsilon)\in\mathscr B}
       \mathbf F_3[v_{L,\epsilon}]\oplus
 \bigoplus_{E\in\mathscr E}\mathbf F_3[v_E],\\
 C_1(\widehat x)
 &:=
 \bigoplus_{L\in\mathscr I}\mathbf F_3[u_L]\oplus
 \bigoplus_{(L,\epsilon)\in\mathscr B}
       \mathbf F_3[u_{L,\epsilon}]\oplus
 \bigoplus_{E=\{L,L'\}\in\mathscr E}
 \bigl(\mathbf F_3[u_{L'\mid L}]
       \oplus\mathbf F_3[u_{L\mid L'}]\bigr),\\
 C_2(\widehat x)
 &:=
 \bigoplus_{E\in\mathscr E}\mathbf F_3[\omega_E].
\end{align*}
Son action sur les générateurs est
\begin{align}
 \partial u_L&=v_L-v_\star,
 &
 \partial u_{L'\mid L}&=v_E-v_L,
 \label{eq:pdfv2-cor-cell-frontera-aristas}\\
 \partial u_{L,\epsilon}&=v_{L,\epsilon}-v_L,
 &
 \partial\omega_{\{L,L'\}}
 &=u_L+u_{L'\mid L}-u_{L'}-u_{L\mid L'}.
 \label{eq:pdfv2-cor-cell-frontera-celda}
\end{align}
\end{definition}

Les huit générateurs \(v_{L,\epsilon}\) réalisent concrètement l'écran \(\mathscr B\). L'involution précédente s'étend au complexe par
\begin{equation}
\label{eq:pdfv2-cor-pantalla-ocho-generadores}
 \vartheta_Lv_{K,\epsilon}:=v_{\vartheta_L(K,\epsilon)},\qquad
 \vartheta_Lu_{K,\epsilon}:=u_{\vartheta_L(K,\epsilon)},
\end{equation}
et fixe les autres générateurs. La base étant libre,
\begin{equation}
\label{eq:pdfv2-cor-pantalla-ocho-no-colapso}
 |\{v_{L,\epsilon}:(L,\epsilon)\in\mathscr B\}|=8,\qquad
 \vartheta_L^2=I,\qquad
 \vartheta_L\vartheta_{L'}=\vartheta_{L'}\vartheta_L.
\end{equation}
En outre, \(\partial\vartheta_L=\vartheta_L\partial\), car les deux parcours envoient \(u_{L,\epsilon}\) sur \(v_{L,-\epsilon}-v_L\) et fixent les frontières restantes.

Les deux mots
\[
 p_{L,L'}:=u_{L'\mid L}u_L,\qquad
 p_{L',L}:=u_{L\mid L'}u_{L'}
\]
sont des chemins différents du complexe libre. Leur différence est la frontière de \(\omega_{\{L,L'\}}\), et
\begin{equation}
\label{eq:pdfv2-cor-cell-dos-cero}
 \partial^2\omega_E
 =(v_L-v_\star)+(v_E-v_L)
 -(v_{L'}-v_\star)-(v_E-v_{L'})=0.
\end{equation}
La représentation du TPK sur les feuillets est fixée par
\[
 \varrho(u_L)=P_L,\qquad
 \varrho(u_{L'\mid L})=P_{L'},\qquad
 \varrho(p_{L,L'})=P_{L'}P_L.
\]
Même si deux matrices exprimées coïncidaient dans une spécialisation, les deux chemins ne se confondraient pas : ils restent des générateurs libres distincts de \(C_1(\widehat x)\), reliés par une 2--cellule explicite.

La coordonnée complète d'incidence est le groupe des 2--cochaînes
\begin{equation}
\label{eq:pdfv2-cor-q-incidencia-total}
 Q_{\rm inc}(\widehat x):=
 C^2(\operatorname{Cell}_{\rm TPK}(\widehat x);\mathbf F_3)
 =\bigoplus_{E\in\mathscr E}\mathbf F_3\omega_E^*
 \simeq\mathbf F_3^6.
\end{equation}
L'opération interne renvoie tout \(Q_{\rm inc}(\widehat x)\), non un élément sélectionné. Chaque \(q\) possède une expansion unique \(q=\sum_Eq_E\omega_E^*\). Pour \(a\in\mathbf F_3^{\mathscr I}\), son changement de potentiel est
\begin{equation}
\label{eq:pdfv2-cor-collar-variacion}
 q\longmapsto q+Ba,\qquad (Ba)_{\{L,L'\}}=a_L+a_{L'}.
\end{equation}
Puisque chaque droite appartient à trois paires,
\[
 s(Ba)=3\sum_{L\in\mathscr I}a_L=0,
\]
et donc
\begin{equation}
\label{eq:pdfv2-cor-wc-invariante-celular}
 W_c(q+Ba)=W_c(q)\qquad(q\in Q_{\rm inc}(\widehat x)).
\end{equation}
L'unité de la complétion ne fait qu'ancrer le complexe dans l'état :
\begin{equation}
\label{eq:pdfv2-cor-cell-unidad}
 \eta^{\rm cell}_{\widehat x}:V_{\widehat x}\longrightarrow
 V_{\widehat x}\otimes C_0(\operatorname{Cell}_{\rm TPK}(\widehat x)),\qquad
 e_{\diamond,\widehat x}\longmapsto
 e_{\diamond,\widehat x}\otimes v_\star,
\end{equation}
et marque l'origine \(q=0\) ; elle n'affirme pas que l'orbite de cette origine soit tout \(Q_{\rm inc}\).

La construction est fonctorielle relativement à la troncature commune. Si \(\tau_{\ell,k}\widehat x_\ell=\widehat x_k\), nous définissons
\begin{align*}
 \tau_V&:V_{\widehat x_\ell}\longrightarrow V_{\widehat x_k},&
 \tau_Ve_{\diamond,\widehat x_\ell}&=e_{\diamond,\widehat x_k},\\
 \tau_a&:\mathbf F_3^{\mathscr I_{\widehat x_\ell}}
          \longrightarrow\mathbf F_3^{\mathscr I_{\widehat x_k}},&
 (\tau_aa)_{\tau_{\mathscr I}L}&=a_L,\\
 \tau_Q&:Q_{\rm inc}(\widehat x_\ell)
          \longrightarrow Q_{\rm inc}(\widehat x_k),&
 (\tau_Qq)_{\tau_{\mathscr E}E}&=q_E,
\end{align*}
où \(\tau_{\mathscr I}[v]=[\tau_Vv]\) et \(\tau_{\mathscr E}\{L,L'\}=\{\tau_{\mathscr I}L,
\tau_{\mathscr I}L'\}\). L'application de chaînes \(\operatorname{Cell}(\tau):C_\bullet(\widehat x_\ell)
\to C_\bullet(\widehat x_k)\) agit sur les générateurs par
\begin{equation}
\label{eq:pdfv2-cor-cell-truncamiento-generadores}
 v_\bullet(\widehat x_\ell)\mapsto v_\bullet(\widehat x_k),\quad
 u_\bullet(\widehat x_\ell)\mapsto u_\bullet(\widehat x_k),\quad
 \omega_E(\widehat x_\ell)\mapsto\omega_E(\widehat x_k).
\end{equation}
De \eqref{eq:pdfv2-cor-cell-frontera-aristas}--\eqref{eq:pdfv2-cor-cell-frontera-celda}, on obtient
\begin{equation}
\label{eq:pdfv2-cor-collar-naturalidad}
\begin{aligned}
 \operatorname{Cell}(\tau)\partial
 &=\partial\operatorname{Cell}(\tau),\\
 \tau_VP_L&=P_{\tau_{\mathscr I}L}\tau_V,\qquad
 \tau_QB=B\tau_a,\qquad W_c\circ\tau_Q=W_c,\\
 \operatorname{Cell}(\tau)\vartheta_L
 &=\vartheta_{\tau_{\mathscr I}L}\operatorname{Cell}(\tau),\\
 (\tau_V\otimes\operatorname{Cell}_0(\tau))
 \eta^{\rm cell}_{\widehat x_\ell}
 &=\eta^{\rm cell}_{\widehat x_k}\tau_V.
\end{aligned}
\end{equation}
L'identité et la composition se vérifient sur chaque générateur. Ainsi, \(\operatorname{Cell}_{\rm TPK}\) est la complétion cellulaire fonctorielle de la même fibre APP--TRIT--TPK, non un support ajouté à partir d'une expression ultérieure.

\subsection{Mesure d'incidence, courbure de retenue et naturalité adjointe}
Dans le développement suivant, on identifie, exclusivement comme abréviations de la complétion commune,
\[
 I_{\rm gau}:=\mathscr I_{\widehat x},\qquad
 \mathscr B_{\rm gau}:=\mathscr B_{\widehat x},\qquad
 Q_{\rm gau}:=Q_{\rm inc}(\widehat x).
\]
Avec la moyenne uniforme de comptage sur \(Q_{\rm gau}\), la fonction \(W_c\) de \eqref{eq:pdfv2-cor-wc} satisfait
\begin{equation}
\label{eq:pdfv2-cor-wc-momentos}
 \mathbb E_QW_c=0,\qquad
 \mathbb E_QW_c^2=\frac29,\qquad
 \mathbb E_QW_c^3=\frac{2}{27}.
\end{equation}
Ces identités sont des propriétés de la même cellule et de la même retenue, et non
des données d’une théorie extérieure.
\begingroup
  \let\section\subsection
  % Extracción quirúrgica del fuente teoremática V2 05c: líneas 99--335.
% El resto permanece en este archivo como procedencia, pero no se publica.
\iffalse
\chapter[Macrostructure de jauge interne]{Macrostructure de jauge interne engendrée par le continu discret}
\label{chap:gauint-macroestructura}

\begin{thesisbox}
La macrostructure de ce chapitre se construit uniquement à partir d'une
restriction survivante de \(\mathcal C_{\rm cont}^{\rm disc}\). Les nombres
\(4\), \(6\) et \(8\), l'écran, les quatre involutions, la mesure
projective, l'opérateur de transfert, le mode d'incidence et le terminal
cellulaire se déduisent au sein de cette restriction. Aucune évaluation ultérieure
ne fait partie du domaine ni des flèches qui suivent.
\end{thesisbox}

\section{Le germe projectif et la dérivation interne de \texorpdfstring{\(4/6/8\)}{4/6/8}}

Soit \(V_{\rm TPK}=\mathbf F_3^2\), avec la base ordonnée héritée des deux
feuillets de l'état TPK. Nous définissons l'ensemble des directions
\begin{equation}\label{eq:gauint-cuatro-direcciones}
 I_{\rm gau}:=\mathbf P^1(\mathbf F_3)
 =\{L_\infty,L_0,L_1,L_{-1}\},
\end{equation}
où
\[
 L_\infty=[1:0],\qquad L_0=[0:1],\qquad
 L_1=[1:1],\qquad L_{-1}=[1:-1].
\]
Par conséquent,
\begin{equation}\label{eq:gauint-cardinal-cuatro}
 |I_{\rm gau}|=\frac{3^2-1}{3-1}=4.
\end{equation}
On n'introduit pas un ensemble de quatre indices : il apparaît comme le quotient des
huit vecteurs non nuls par l'action de \(\mathbf F_3^\times\).

Les incidences non orientées sont
\begin{equation}\label{eq:gauint-seis-incidencias}
 E_{\rm gau}:=\binom{I_{\rm gau}}2,
 \qquad |E_{\rm gau}|=\binom42=6,
 \qquad Q_{\rm gau}:=\mathbf F_3^{E_{\rm gau}}.
\end{equation}
L'écran orienté double chaque direction, et non chaque incidence :
\begin{equation}\label{eq:gauint-ocho-caras}
 \mathscr B_{\rm gau}:=I_{\rm gau}\times\{+,-\},
 \qquad |\mathscr B_{\rm gau}|=2|I_{\rm gau}|=8.
\end{equation}
Ainsi, \(4\), \(6\) et \(8\) proviennent respectivement de directions, de couples de
directions et de faces orientées d'un même germe.

Pour chaque \(L\in I_{\rm gau}\), soit \(\vartheta_L\) l'involution de
l'écran qui échange \((L,+)\) et \((L,-)\) et laisse fixes les six autres
faces. Alors
\begin{equation}\label{eq:gauint-cuatro-involuciones}
 \vartheta_L^2=I,
 \qquad
 \vartheta_L\vartheta_{L'}=\vartheta_{L'}\vartheta_L
 \quad(L\ne L').
\end{equation}
Les quatre involutions ont été obtenues avant de définir une mesure ou un opérateur.

\section{Incidence, mode interne et témoin non constant}

L'application d'incidence est
\begin{equation}\label{eq:gauint-incidencia-b}
 B:\mathbf F_3^{I_{\rm gau}}\longrightarrow Q_{\rm gau},
 \qquad
 (Ba)_{\{L,L'\}}:=a_L+a_{L'}.
\end{equation}
Si
\[
 s(q):=\sum_{e\in E_{\rm gau}}q_e\in\mathbf F_3,
\]
chaque direction apparaît exactement trois fois dans la somme des six
incidences, et donc
\begin{equation}\label{eq:gauint-incidencia-invariante}
 s(Ba)=3\sum_{L\in I_{\rm gau}}a_L=0,
 \qquad
 s(q+Ba)=s(q).
\end{equation}

Sur \(Q_{\rm gau}\), on utilise exclusivement la moyenne de comptage
normalisée. Le mode d'incidence est la fonction rationnelle
\begin{equation}\label{eq:gauint-wc-definicion}
 W_c(q):=\mathbf 1_{\{s(q)=0\}}-\frac13.
\end{equation}
Comme \(s:Q_{\rm gau}\to\mathbf F_3\) est surjective, chaque fibre a
\(3^5\) éléments. On en obtient, par comptage exact,
\begin{equation}\label{eq:gauint-wc-momentos}
 \mathbb E_{Q}W_c=0,
 \qquad
 \mathbb E_{Q}W_c^2=\frac29,
 \qquad
 \mathbb E_{Q}W_c^3=\frac{2}{27}.
\end{equation}
En outre, \eqref{eq:gauint-incidencia-invariante} donne
\begin{equation}\label{eq:gauint-wc-invariancia}
 W_c(q+Ba)=W_c(q).
\end{equation}
Donc \(W_c\) est non nul, centré et invariant par l'incidence
engendrée en interne.

\fi
\section{Incidence, courbure de retenue et domaine commun de naturalité}

Pour \(x\in\mathcal C_{\rm cont}^{\rm disc}\), notons
\(Y_m(x)\) l'ensemble fini des préfixes de jauge de profondeur \(m\) qui
conservent simultanément
\begin{equation*}
\begin{gathered}
 \text{fibre, relations d’extension, nombres, orientation, retenue, mémoire,}\\
 \text{frontière, route, multisection, survie, terminal et lecteur}.
\end{gathered}
\end{equation*}
Il existe une application de collier
\begin{equation}\label{eq:gauint-collar}
 q_m:Y_m(x)\longrightarrow Q_{\rm gau}
\end{equation}
qui enregistre les six incidences sans effacer l'histoire qui les produit.
Pour \(\ell\ge m\), la troncature s'écrit
\[
 \tau_{\ell,m}:Y_\ell(x)\longrightarrow Y_m(x).
\]

Pour \(y,z\in Y_m(x)\), soit
\begin{equation}\label{eq:gauint-todas-extensiones}
 \operatorname{Ext}^{\rm adm}_{m}(y,z)
\end{equation}
l'ensemble de \emph{toutes} les extensions TPK admissibles de \(y\) à \(z\).
Une extension appartient à cet ensemble seulement si elle transporte la fibre entière,
les relations \(\operatorname{Ext}_\Gamma\), la retenue native et opératorielle,
l'orientation, la mémoire accumulée, la frontière, le chemin et la condition de
survie. Nous définissons
\begin{equation}\label{eq:gauint-conteo-kernel}
 a_m(y,z):=\#\operatorname{Ext}^{\rm adm}_m(y,z),
 \qquad
 d_m(y):=\sum_{z\in Y_m(x)}a_m(y,z),
\end{equation}
et, lorsque \(d_m(y)>0\),
\begin{equation}\label{eq:gauint-kernel-normalizado}
 k_m(y,z):=\frac{a_m(y,z)}{d_m(y)}.
\end{equation}
Ainsi, \(k_m\) est un comptage normalisé d'extensions complètes, non une
fonction introduite après la projection de l'état.

La restriction \(\mathcal C_{\rm cont}^{\rm gau}\) est formée des
histoires survivantes \(x\) pour lesquelles les conditions
suivantes sont satisfaites à toute profondeur.
\begin{enumerate}
 \item Les ensembles \(Y_m(x)\) sont non vides, \(d_m(y)>0\), et les
 troncatures \(\tau_{\ell,m}\) sont surjectives.
 \item Il existe une multisection non vide de poids rationnels strictement
 positifs \(\mu_m:Y_m(x)\to\mathbf Q_{>0}\) telle que
 \begin{equation}\label{eq:gauint-medida-proyectiva}
  \sum_{y\in Y_m}\mu_m(y)=1,
  \qquad
  \mu_m(y)=\sum_{\tau_{\ell,m}\widetilde y=y}\mu_\ell(\widetilde y).
 \end{equation}
 \item Chaque \(\mu_m\) est stationnaire :
 \begin{equation}\label{eq:gauint-estacionariedad}
  \sum_{y\in Y_m}\mu_m(y)k_m(y,z)=\mu_m(z).
 \end{equation}
 \item La marginale de \((q_m)_*\mu_m\) est le comptage uniforme de
 \(Q_{\rm gau}\). Cette condition rend littéralement applicables les calculs
 de \eqref{eq:pdfv2-cor-wc-momentos} au relèvement \(W_c\circ q_m\).
 \item La lumpabilité directe est satisfaite
 \begin{equation}\label{eq:gauint-lump-directa}
  \sum_{\tau_{\ell,m}\widetilde z=z}
   k_\ell(\widetilde y,\widetilde z)
  =k_m(\tau_{\ell,m}\widetilde y,z)
 \end{equation}
 pour tout \(\widetilde y\in Y_\ell\) et \(z\in Y_m\).
 \item Pour le noyau adjoint
 \begin{equation}\label{eq:gauint-kernel-adjunto}
  k_m^*(z,y):=\frac{\mu_m(y)k_m(y,z)}{\mu_m(z)},
 \end{equation}
 la lumpabilité adjointe est satisfaite
 \begin{equation}\label{eq:gauint-lump-adjunta}
  \sum_{\tau_{\ell,m}\widetilde y=y}
   k_\ell^*(\widetilde z,\widetilde y)
  =k_m^*(\tau_{\ell,m}\widetilde z,y).
 \end{equation}
 \item Les catégories de préfixes possèdent la racine TPK comme objet initial,
 de sorte que leur nerf possède la contraction terminale explicite de
 \eqref{eq:gauint-sdr} ci-dessous.
\end{enumerate}
Ces conditions définissent le domaine typé ; en dehors de celui-ci, on ne publie pas une
mesure ou un opérateur naturel par simple déclaration.

\section{Mesure projective, noyau et naturalité directe et adjointe}

Une branche de la multisection \(\{\mu_m\}_m\) étant fixée, soit
\begin{equation}\label{eq:gauint-hm}
 H_m:=\operatorname{Fun}(Y_m,\mathbf Q),
 \qquad
 \langle f,g\rangle_m:=\sum_{y\in Y_m}\mu_m(y)f(y)g(y).
\end{equation}
L'opérateur engendré par toutes les extensions est
\begin{equation}\label{eq:gauint-km}
 (K_mf)(y):=\sum_{z\in Y_m}k_m(y,z)f(z).
\end{equation}
Les équations de Markov et de stationnarité donnent
\begin{equation}\label{eq:gauint-markov-estacionario}
 K_m\mathbf1=\mathbf1,
 \qquad
 K_m^*\mathbf1=\mathbf1,
 \qquad
 \|K_mf\|_m\le\|f\|_m.
\end{equation}

Le relèvement et l'espérance de troncature sont
\begin{equation}\label{eq:gauint-lift-esperanza}
 J_{\ell,m}f:=f\circ\tau_{\ell,m},
 \qquad
 E_{m,\ell}:=J_{\ell,m}^*.
\end{equation}
La projectivité de \(\mu\) implique
\begin{equation}\label{eq:gauint-ej-identidad}
 E_{m,\ell}J_{\ell,m}=I_{H_m}.
\end{equation}
Les deux conditions de lumpabilité établissent, respectivement,
\begin{equation}\label{eq:gauint-naturalidad-directa-adjunta}
 K_\ell J_{\ell,m}=J_{\ell,m}K_m,
 \qquad
 K_\ell^*J_{\ell,m}=J_{\ell,m}K_m^*.
\end{equation}
La seconde égalité ne se déduit pas de la première sans la condition adjointe :
les deux sont enregistrées dans le domaine.

On définit l'opérateur interne positif
\begin{equation}\label{eq:gauint-transferencia-t}
 T_m:=K_m^*K_m.
\end{equation}
Alors
\begin{equation}\label{eq:gauint-t-positivo-natural}
 \langle f,T_mf\rangle_m=\|K_mf\|_m^2\ge0,
 \qquad
 T_\ell J_{\ell,m}=J_{\ell,m}T_m.
\end{equation}

\section{Loi commune de \(8\) faces et \(4\) carrés de Fubini}

Sur \(Y_m^{\mathscr B_{\rm gau}}\), on définit la masse rationnelle
\begin{equation}\label{eq:gauint-ley-ocho-caras}
 \Omega_m^{(8)}(\mathbf y)
 :=\sum_{z\in Y_m}\mu_m(z)
 \prod_{(L,\varepsilon)\in\mathscr B_{\rm gau}}
 k_m(z,y_{L,\varepsilon}).
\end{equation}
C'est une probabilité parce que chaque facteur est markovien. Toutes les faces proviennent
du même état latent \(z\) et du même noyau ; ce ne sont pas huit mesures
indépendantes ajoutées.

Pour des fonctions \(f_{L,+},f_{L,-}\in H_m\), Fubini fini donne
\begin{equation}\label{eq:gauint-fubini-ocho}
\begin{aligned}
 &\sum_{\mathbf y}\Omega_m^{(8)}(\mathbf y)
  \prod_{L\in I_{\rm gau}}
  f_{L,+}(y_{L,+})f_{L,-}(y_{L,-})\\
 &\qquad=
 \sum_{z\in Y_m}\mu_m(z)
 \prod_{L\in I_{\rm gau}}
 (K_mf_{L,+})(z)(K_mf_{L,-})(z).
\end{aligned}
\end{equation}
En prenant \(f_{L,+}=f_{L,-}=f_L\), les quatre
carrés apparaissent simultanément
\begin{equation}\label{eq:gauint-cuatro-cuadrados}
 \sum_{z\in Y_m}\mu_m(z)
 \prod_{L\in I_{\rm gau}}(K_mf_L(z))^2.
\end{equation}
La marginale de tout couple opposé satisfait
\begin{equation}\label{eq:gauint-marginal-t}
 \sum_{y_+,y_-}\Omega_{m,L}^{(2)}(y_+,y_-)
 f(y_+)g(y_-)
 =\langle K_mf,K_mg\rangle_m
 =\langle f,T_mg\rangle_m.
\end{equation}
La commutativité du produit dans \eqref{eq:gauint-ley-ocho-caras} prouve
\begin{equation}\label{eq:gauint-involuciones-ley}
 (\vartheta_L)_*\Omega_m^{(8)}=\Omega_m^{(8)}
 \qquad(L\in I_{\rm gau}).
\end{equation}

\section{Opérateur interne exact et témoin propre}

La projection sur le secteur constant et son complément sont
\begin{equation}\label{eq:gauint-pq}
 P_{\Omega,m}f:=\langle\mathbf1,f\rangle_m\mathbf1,
 \qquad
 Q_m:=I-P_{\Omega,m}.
\end{equation}
L'opérateur cellulaire interne est défini par
\begin{equation}\label{eq:gauint-d-hmt}
 D_m^{\rm HMT}:=3Q_m.
\end{equation}
Le coefficient \(3\) est le cardinal du corps des résidus utilisé dans
\eqref{eq:pdfv2-cor-proyectores-cuatro-direcciones} ; il ne représente pas une
échelle ajoutée.
Pour le relèvement
\[
 W_{c,m}:=W_c\circ q_m
\]
on a, par \eqref{eq:pdfv2-cor-wc-momentos},
\begin{equation}\label{eq:gauint-w-eigen}
 P_{\Omega,m}W_{c,m}=0,
 \qquad
 D_m^{\rm HMT}W_{c,m}=3W_{c,m},
 \qquad
 \|W_{c,m}\|_m^2=\frac29.
\end{equation}
Ainsi, le secteur complémentaire possède un témoin non nul construit par
l'incidence des six arêtes.

\section{Terminal cellulaire et rétracte fort}

Soit \(C_*(N\mathsf Y_m;A_n)\) le complexe normalisé du nerf de la
catégorie des préfixes, et soit \(y_*\) sa racine initiale. L'inclusion et
l'augmentation sont
\[
 \iota_m:A_n[0]\longrightarrow C_*(N\mathsf Y_m;A_n),
 \quad \iota_m(1)=[y_*],
 \qquad
 \epsilon_m:C_*(N\mathsf Y_m;A_n)\longrightarrow A_n[0].
\]
L'objet initial fournit la dégénérescence supplémentaire
\begin{equation}\label{eq:gauint-homotopia-celular}
 H_m[y_0\to\cdots\to y_q]
 :=[y_*\to y_0\to\cdots\to y_q],
\end{equation}
avec la convention normalisée lorsqu'une dégénérescence apparaît. On vérifie
\begin{equation}\label{eq:gauint-sdr}
 \epsilon_m\iota_m=I,
 \qquad
 \partial H_m+H_m\partial=I-\iota_m\epsilon_m,
 \qquad
 H_m^2=H_m\iota_m=\epsilon_mH_m=0.
\end{equation}
Ce terminal cellulaire est distinct de la projection constante
\(P_{\Omega,m}\) : le premier contracte le nerf et la seconde sépare les
secteurs de l'espace des fonctions. Les deux sont conservés.

\iffalse
\section{Les treize coordonnées de la restriction de jauge}

Pour chaque \((m,n)\), on définit
\begin{equation}\label{eq:gauint-d-trece}
 D_{{\rm gau},m,n}(x):=
 (\mathcal F,\operatorname{Ext}_\Gamma,\operatorname{Rel},
 \mathcal N,\operatorname{or},\operatorname{Carr},\operatorname{Mem},
 \partial,\gamma,\operatorname{Msect},\operatorname{Surv},
 \operatorname{Term},\mathcal L)_{{\rm gau},m,n}(x),
\end{equation}
avec le contenu suivant.
\begin{enumerate}
 \item \(\mathcal F=(V_{\rm TPK},I_{\rm gau},E_{\rm gau},
 Q_{\rm gau},\mathscr B_{\rm gau},Y_m,H_m)\).
 \item \(\operatorname{Ext}_\Gamma\) contient tous les ensembles
 \(\operatorname{Ext}^{\rm adm}_m(y,z)\), les troncatures, les relèvements,
 les espérances et leurs compositions.
 \item \(\operatorname{Rel}\) contient \(B\), \(sB=0\), les quatre
 involutions, Markov, la stationnarité et les deux lumpabilités.
 \item \(\mathcal N=(3,4,6,8,m,n,|Y_m|,a_m,d_m)\) ; ces nombres ne
 remplacent aucune des autres coordonnées.
 \item \(\operatorname{or}\) conserve les faces \(+/-\), les quatre
 inversions \(\vartheta_L\) et le renversement de chaque chemin.
 \item \(\operatorname{Carr}\) contient la retenue complète de chaque
 extension comptée dans \(a_m(y,z)\), y compris la retenue déjà propagée.
 \item \(\operatorname{Mem}\) conserve le préfixe complet, l'état
 commun \(z\) de \(\Omega_m^{(8)}\) et la mémoire de ses huit extensions.
 \item \(\partial=(q_m,\partial_{N\mathsf Y_m},Q_m,H_m)\) conserve
 la frontière de collier, la frontière cellulaire et le complément.
 \item \(\gamma\) est le simplexe orienté du nerf qui enregistre le chemin
 et non seulement ses extrémités.
 \item \(\operatorname{Msect}\) est la famille de toutes les mesures
 stationnaires projectives compatibles et de toutes les extensions qui
 interviennent dans les comptages ; aucune branche rétrospective n'est sélectionnée.
 \item \(\operatorname{Surv}\) enregistre la non-vacuité, la profondeur
 arbitraire, la projectivité, les deux lumpabilités et la persistance de la
 marginale uniforme de collier.
 \item \(\operatorname{Term}=(\iota_m,\epsilon_m,H_m;
 P_{\Omega,m},Q_m)\) conserve séparément les deux terminaux.
 \item \(\mathcal L=(\Omega_m^{(8)},K_m,T_m,D_m^{\rm HMT},W_{c,m})\)
 est le lecteur interne postérieur à la construction de l'histoire.
\end{enumerate}

Pour \(\ell\ge m\) et \(n'\ge n\), toutes les coordonnées satisfont
\begin{equation}\label{eq:gauint-naturalidad-trece}
 u^{\rm gau}_{(\ell,n'),(m,n)}D_{{\rm gau},\ell,n'}(x')
 =D_{{\rm gau},m,n}(\tau_{\ell,m}x').
\end{equation}

\section{Assembleur, publicateur et chaîne non ablative}

Soit
\begin{equation}\label{eq:gauint-graph-d}
 \mathcal G_{\rm gau}:=\operatorname{Graph}(D_{\rm gau}),
 \qquad
 \operatorname{Res}_{\rm gau}(x):=(x,D_{\rm gau}x).
\end{equation}
La multisection dépendante \(\chi_{\rm gau}(x)\) contient toutes les branches
compatibles de collier, de mesure et d'extension. Nous définissons
\begin{equation}\label{eq:gauint-x-dependiente}
 X_{\chi,{\rm gau}}
 :=\{(\chi_{\rm gau}(x),x,D_{\rm gau}x):
 x\in\mathcal C_{\rm cont}^{\rm gau}\},
\end{equation}
\begin{equation}\label{eq:gauint-prx}
 \operatorname{pr}_{X,{\rm gau}}(x,D_{\rm gau}x)
 :=(\chi_{\rm gau}(x),x,D_{\rm gau}x).
\end{equation}

L'assembleur brut
\begin{equation}\label{eq:gauint-ensamblador-crudo}
 a_{\rm gau}:X_{\chi,{\rm gau}}\longrightarrow
 \mathscr M_{\rm gau}^{\rm amb}
\end{equation}
produit la famille complète
\[
 \bigl(I,E,\mathscr B,B,W_c;
 Y_m,\mu_m,k_m,J_{\ell,m},E_{m,\ell};
 K_m,T_m,\Omega_m^{(8)},P_{\Omega,m},Q_m,D_m^{\rm HMT};
 \iota_m,\epsilon_m,H_m\bigr)_{m,n}.
\]
Le macro-objet conserve son entrée par
\begin{equation}\label{eq:gauint-graph-a}
 \mathfrak M_{\rm gau}:=\operatorname{Graph}(a_{\rm gau}),
 \qquad
 \mathcal A_{\rm gau}(z):=(z,a_{\rm gau}z).
\end{equation}

Le publicateur brut
\begin{equation}\label{eq:gauint-publicador-crudo}
 p_{\rm gau}:\mathfrak M_{\rm gau}\longrightarrow
 \mathscr Y_{\rm gau}
\end{equation}
publie conjointement les identités
\[
 \begin{gathered}
 |I|=4,quad |E|=6,quad |\mathscr B|=8,quad sB=0,
 \quad W_c(q+Ba)=W_c(q),\\
 K_\ell J=JK_m,quad K_\ell^*J=JK_m^*,
 \quad T_m=K_m^*K_m,\\
 (\vartheta_L)_*\Omega_m^{(8)}=\Omega_m^{(8)},
 \quad \text{les quatre factorisations quadratiques de Fubini},\\
 D_m^{\rm HMT}=3Q_m,quad D_m^{\rm HMT}W_{c,m}=3W_{c,m},
 \quad \partial H_m+H_m\partial=I-\iota_m\epsilon_m.
 \end{gathered}
\]
Enfin,
\begin{equation}\label{eq:gauint-graph-p}
 \mathcal C_{\rm gau}^{\rm HMT}:=\operatorname{Graph}(p_{\rm gau}),
 \qquad
 \mathcal P_{\rm gau}(M):=(M,p_{\rm gau}M).
\end{equation}
La chaîne complète est
\begin{equation}\label{eq:gauint-cadena-completa}
 \mathcal C_{\rm cont}^{\rm disc}\supset
 \mathcal C_{\rm cont}^{\rm gau}
 \xrightarrow{\operatorname{Res}_{\rm gau}}\mathcal G_{\rm gau}
 \xrightarrow{\operatorname{pr}_{X,{\rm gau}}}X_{\chi,{\rm gau}}
 \xrightarrow{\mathcal A_{\rm gau}}\mathfrak M_{\rm gau}
 \xrightarrow{\mathcal P_{\rm gau}}\mathcal C_{\rm gau}^{\rm HMT}.
\end{equation}
Chaque flèche possède un inverse sur son image donné par une projection de
coordonnées. Aucune sortie ne remplace son histoire génératrice.

\section{Théorème de réalisation de jauge interne}

\begin{theorem}[Réalisation interne complète et naturelle]
\label{thm:gauint-realizacion}
Pour tout \(x\in\mathcal C_{\rm cont}^{\rm gau}\), la chaîne
\eqref{eq:gauint-cadena-completa} construit, à partir de la structure
discrète du continu, une macrostructure naturelle avec quatre directions,
six incidences, huit faces, quatre involutions, une mesure projective, un
noyau de comptage complet, son transfert positif, une loi commune de huit
faces, un témoin d'incidence non constant et un terminal cellulaire. Les
treize coordonnées de \eqref{eq:gauint-d-trece} restent récupérables en
sortie.
\end{theorem}

\begin{proof}
Les équations \eqref{eq:gauint-cuatro-direcciones}--
\eqref{eq:gauint-ocho-caras} déduisent \(4/6/8\). La somme des incidences
\eqref{eq:gauint-incidencia-invariante} prouve l'invariance de \(W_c\), et
le comptage uniforme prouve \eqref{eq:gauint-wc-momentos}. Les extensions
complètes produisent \(k_m\) par \eqref{eq:gauint-kernel-normalizado} ; les
conditions explicites du domaine prouvent Markov, la stationnarité et les deux
égalités de naturalité. Donc \(T_m=K_m^*K_m\) est positif et naturel.
Fubini appliqué une seule fois à \eqref{eq:gauint-ley-ocho-caras} donne les quatre
formes quadratiques sans changer de mesure. La marginale uniforme centre
\(W_{c,m}\), d'où \eqref{eq:gauint-w-eigen}. L'objet initial du nerf
donne la dégénérescence supplémentaire et l'identité SDR \eqref{eq:gauint-sdr}. Enfin,
les trois constructions par graphe retiennent leurs entrées comme premières
coordonnées, donc la composition n'oublie aucune des treize entrées.
\end{proof}

\section{Provenance et falsificateurs}

\paragraph{Provenance.}
Le germe TPK, l'écran orienté, les extensions survivantes, la
mesure commune et la contraction terminale ont le statut
\texttt{ARCHITECTURE\_AUTORALE\_PRÉEXISTANTE}. Le transfert par comptage,
le mode d'incidence et le secteur exact \(3Q\) sont présentés comme
\texttt{RÉSULTAT\_RÉCUPÉRÉ}. La dérivation explicite
\(\mathbf P^1(\mathbf F_3)\to4/6/8\), les deux conditions de lumpabilité et
la mise en paquet non ablative par graphes ont le statut
\texttt{FORMALISATION\_NOUVELLE} ; elles formalisent l'architecture antérieure sans
revendiquer une origine conceptuelle nouvelle.

\begin{ablation}[Falsificateurs internes de la branche de jauge]
\label{abl:gauint-falsadores}
La construction est falsifiée au niveau indiqué si l'un quelconque des
faits suivants se produit.
\begin{enumerate}
 \item On remplace \(I_{\rm gau}\), \(E_{\rm gau}\) ou
 \(\mathscr B_{\rm gau}\) par une liste sans les applications qui les déduisent.
 \item \(sB=0\) échoue ou \(W_c\) cesse d'être invariant et non constant.
 \item Le comptage \(a_m(y,z)\) omet la retenue, l'orientation, la mémoire, la frontière ou
 le chemin : dans ce cas, on a compté un autre objet.
 \item La projectivité, la stationnarité ou l'une des deux
 lumpabilités échoue ; on ne peut alors publier la naturalité correspondante.
 \item Les huit faces ne proviennent pas d'une unique \(\mu_m\) et d'un unique
 \(k_m\), ou l'une des \(\vartheta_L\) ne préserve pas \(\Omega_m^{(8)}\).
 \item \(P_{\Omega,m}W_{c,m}\ne0\) o
 \(D_m^{\rm HMT}W_{c,m}\ne3W_{c,m}\).
 \item L'identité SDR \(\partial H+H\partial=I-\iota\epsilon\) échoue.
 \item L'une des treize coordonnées est omise : \emph{objet substitué ;
 conclusion non transférable}.
\end{enumerate}
Ces falsificateurs contrôlent l'appartenance et la naturalité ; ils ne dégradent pas une
histoire qui satisfait toutes les identités.
\end{ablation}
\fi
 \endgroup

La même règle définit son action sur une émission élémentaire \(\alpha:x\to y\). Puisque \(U_\alpha\) transporte la paire \emph{ordonnée} de feuillets, son application sur les étiquettes est \(\overline U_\alpha e_{\Sigma,x}=e_{\Sigma,y}\) et \(\overline U_\alpha e_{\Pi,x}=e_{\Pi,y}\). Sa projectivisation \(\alpha_{\mathscr I}[v]:=[\overline U_\alpha v]\) agit sur \(\mathbf P^1(V_x)\) et, par conséquent,
\begin{equation}
\label{eq:pdfv2-cor-cell-accion-emision}
 \begin{gathered}
 \operatorname{Cell}_{\rm TPK}(\alpha)v_\star(x)=v_\star(y),\qquad
 \operatorname{Cell}_{\rm TPK}(\alpha)v_L(x)
   =v_{\alpha_{\mathscr I}L}(y),\\
 \operatorname{Cell}_{\rm TPK}(\alpha)v_{\{L,L'\}}(x)
   =v_{\{\alpha_{\mathscr I}L,\alpha_{\mathscr I}L'\}}(y),\qquad
 \operatorname{Cell}_{\rm TPK}(\alpha)v_{L,\epsilon}(x)
   =v_{\alpha_{\mathscr I}L,\epsilon}(y),\\
 \operatorname{Cell}_{\rm TPK}(\alpha)u_L(x)
   =u_{\alpha_{\mathscr I}L}(y),\qquad
 \operatorname{Cell}_{\rm TPK}(\alpha)u_{L,\epsilon}(x)
   =u_{\alpha_{\mathscr I}L,\epsilon}(y),\\
 \operatorname{Cell}_{\rm TPK}(\alpha)u_{L'\mid L}(x)
   =u_{\alpha_{\mathscr I}L'\mid\alpha_{\mathscr I}L}(y),\qquad
 \operatorname{Cell}_{\rm TPK}(\alpha)\omega_{\{L,L'\}}(x)
   =\omega_{\{\alpha_{\mathscr I}L,
                 \alpha_{\mathscr I}L'\}}(y).
 \end{gathered}
\end{equation}
Dans la fibre des feuillets et sur son écran, on a en outre
\begin{equation}
\label{eq:pdfv2-cor-cell-entrelaza-proyectores-pantalla}
 \overline U_\alpha P_L=P_{\alpha_{\mathscr I}L}\overline U_\alpha,
 \qquad
 \operatorname{Cell}_{\rm TPK}(\alpha)\vartheta_L
 =\vartheta_{\alpha_{\mathscr I}L}
  \operatorname{Cell}_{\rm TPK}(\alpha).
\end{equation}
Les formules de frontière donnent
\begin{equation}
\label{eq:pdfv2-cor-cell-accion-emision-natural}
 \partial\operatorname{Cell}_{\rm TPK}(\alpha)
 =\operatorname{Cell}_{\rm TPK}(\alpha)\partial,\qquad
 \operatorname{Cell}_{\rm TPK}(\beta\alpha)
 =\operatorname{Cell}_{\rm TPK}(\beta)
  \operatorname{Cell}_{\rm TPK}(\alpha).
\end{equation}
Ainsi, l'application de complexes utilisée dans le registre terminal est l'action libre induite par la même émission, non un symbole ajouté.

\subsection{Module des parcours, totalisation et cône du retour nonadique}
Soit \(J_d(x)\) la catégorie complète des préfixes de profondeur \(d\) de
l’arbre commun, \(\operatorname{Term}_d(x)\) son ensemble terminal et
\(\Gamma_{d,N}(\nu):=C_{d,N}(\nu)\) le complexe cellulaire de
\eqref{eq:pdfv2-cor-cono-celda}. Le module, la totalisation et le cône qui
suivent conservent simultanément toutes les continuations ; ils ne sélectionnent pas
une histoire ni ne séparent une construction cohomologique du continu commun.
\begingroup
  \let\section\subsection
  % Extracción quirúrgica del fuente teoremática V2 05d: líneas 286--512.
% El resto permanece en este archivo como procedencia, pero no se publica.
\iffalse
% Macroestructura de coherencia estrictamente interna del continuo discreto.
% No usa geometrías, clases ni conclusiones clásicas.

\chapter{Macrostructure interne de cellules, cohérence et complétude}
\label{chap:pdfv2-coh-interna}

\begin{statusbox}
L'unique entrée de ce chapitre est la structure discrète du continu. La
branche \(\mathrm{coh}\) se construit sur les anneaux
\(A_N=\mathbb Z/3^N\mathbb Z\), les cellules APP, leurs histoires et leurs
terminaux. Sa chaîne non ablative est
\[
 \mathcal C_{\rm cont}^{\rm disc}
 \supset\mathcal C_{\rm cont}^{\rm coh}
 \xrightarrow{\operatorname{Res}_{\rm coh}}
 \mathcal G_{\rm coh}
 \xrightarrow{\operatorname{pr}_{X,{\rm coh}}}
 X_{\chi,{\rm coh}}
 \xrightarrow{\mathcal A_{\rm coh}}
 \mathfrak M_{\rm coh}
 \xrightarrow{\mathcal P_{\rm coh}}
 \mathcal C_{\rm coh}^{\rm HMT}.
\]
Aucun objet classique n'apparaît en entrée ni en sortie. Chaque image
conserve littéralement son histoire génératrice.
\end{statusbox}

\section{Domaine typé et séparation des indices}
\label{sec:pdfv2-coh-domain}

Pour ne pas identifier profondeur et précision, on utilise l'ensemble dirigé
\begin{equation}\label{eq:pdfv2-coh-lambda}
 \Lambda:=\mathbb N_{\rm prof}\times\mathbb N_{\rm prec},
 \qquad
 \lambda=(d,N),
 \qquad
 A_N:=\mathbb Z/3^N\mathbb Z.
\end{equation}
Si \(\lambda\leq\lambda'=(d',N')\), la transition
\begin{equation}\label{eq:pdfv2-coh-rho-lambda}
 \rho_{\lambda',\lambda}:
 \mathcal C_{{\rm cont},\lambda'}^{\rm disc}
 \longrightarrow
 \mathcal C_{{\rm cont},\lambda}^{\rm disc}
\end{equation}
tronque l'histoire après la profondeur \(d\) et réduit les coefficients par
\(A_{N'}\twoheadrightarrow A_N\), en conservant feuillet, orientation, résidu,
quotient, retenue, mémoire, frontière, chemin, survie et terminal.

La sous-structure \(\mathcal C_{\rm cont}^{\rm coh}\) se compose des histoires
\(x\in\mathcal C_{\rm cont}^{\rm disc}\) qui satisfont, à tout niveau :
\begin{enumerate}
 \item le préfixe de \(x\) porte les cellules orientées et basées définies dans
       la section suivante ;
 \item chaque raffinement induit une application du complexe de la cellule ;
 \item troncature, réduction des coefficients, orientation et extension
       nonadique commutent ;
 \item les multisections d'extension survivante sont non vides à toute
       profondeur ;
 \item le retour nonadique conserve la différence terminale entre l'histoire
       et sa continuation.
\end{enumerate}
Ces conditions définissent le domaine de la branche. Elles ne sont pas attribuées à une
histoire arbitraire sans être vérifiées.

\section{Cellule APP sur \texorpdfstring{\(A_N\)}{AN}}
\label{sec:pdfv2-coh-cell}

Une cellule orientée est
\[
 \nu=(h_\nu,I_\nu,J_\nu,\epsilon_\nu),
\]
où \(h_\nu\) est une histoire finie, \(I_\nu,J_\nu\) sont des ensembles
ordonnés non vides de ports,
\[
 a_\nu:=|I_\nu|\geq1,
 \qquad
 b_\nu:=|J_\nu|\geq1,
\]
et \(\epsilon_\nu\in\{+1,-1\}\) est son orientation.

On définit
\begin{equation}\label{eq:pdfv2-coh-kappa}
 \begin{aligned}
 \kappa_{\nu,N}:
 A_N^{I_\nu}\oplus A_N^{J_\nu}
 &\longrightarrow A_N^{I_\nu\times J_\nu},\\
 \kappa_{\nu,N}(u,v)_{ij}&:=u_i-v_j.
 \end{aligned}
\end{equation}
Les ports de base \(i_0\in I_\nu\), \(j_0\in J_\nu\) étant choisis, soit
\begin{equation}\label{eq:pdfv2-coh-delta}
 \begin{aligned}
 \delta_{\nu,N}:A_N^{I_\nu\times J_\nu}
 &\longrightarrow
 A_N^{(I_\nu\setminus i_0)\times(J_\nu\setminus j_0)},\\
 \delta_{\nu,N}(w)_{ij}
 &:={}
 w_{ij}-w_{i j_0}-w_{i_0j}+w_{i_0j_0}.
 \end{aligned}
\end{equation}

\begin{proposition}[Suite exacte de la cellule]
\label{prop:pdfv2-coh-cell-exact}
Pour toute cellule \(\nu\) et toute précision \(N\), la suite
\begin{equation}\label{eq:pdfv2-coh-exact-sequence}
 0\longrightarrow A_N
 \xrightarrow{\ \Delta\ }
 A_N^{I_\nu}\oplus A_N^{J_\nu}
 \xrightarrow{\ \kappa_{\nu,N}\ }
 A_N^{I_\nu\times J_\nu}
 \xrightarrow{\ \delta_{\nu,N}\ }
 A_N^{(a_\nu-1)(b_\nu-1)}
 \longrightarrow0,
\end{equation}
où
\[
 \Delta(t)=\bigl((t)_{i\in I_\nu},(t)_{j\in J_\nu}\bigr),
\]
est exacte et scindée sur \(A_N\).
\end{proposition}

\begin{proof}
Si \(\kappa_{\nu,N}(u,v)=0\), alors \(u_i=v_j\) pour tout \(i,j\),
de sorte que \((u,v)\) est diagonale. Ainsi,
\(\ker\kappa_{\nu,N}=\operatorname{im}\Delta\). La substitution directe donne
\(\delta_{\nu,N}\kappa_{\nu,N}=0\).

Si \(\delta_{\nu,N}(w)=0\), posons
\[
 u_i=w_{i j_0},
 \qquad
 v_j=w_{i_0j_0}-w_{i_0j}.
\]
Alors
\[
 u_i-v_j=w_{i j_0}+w_{i_0j}-w_{i_0j_0}=w_{ij},
\]
de sorte que \(\ker\delta_{\nu,N}=\operatorname{im}\kappa_{\nu,N}\). Pour
prouver la surjectivité de \(\delta_{\nu,N}\), on annule la ligne et la
colonne de base et l'on place arbitrairement les coordonnées souhaitées aux
positions restantes. Ces mêmes formules fournissent des scindements.
\end{proof}

La cellule complexe, aux degrés \(1\to0\), est le cône homologique de
\(\kappa_{\nu,N}\) avec la convention qui place sa source au degré \(1\) :
\begin{equation}\label{eq:pdfv2-coh-cell-complex}
 C_{\nu,N}:=\operatorname{Cone}(\kappa_{\nu,N})
 :=\left[
 A_N^{I_\nu}\oplus A_N^{J_\nu}
 \xrightarrow{\ \kappa_{\nu,N}\ }
 A_N^{I_\nu\times J_\nu}
 \right].
\end{equation}
La proposition précédente donne
\begin{equation}\label{eq:pdfv2-coh-cell-homology}
 H_1(C_{\nu,N})\simeq A_N,
 \qquad
 H_0(C_{\nu,N})\simeq A_N^{(a_\nu-1)(b_\nu-1)}.
\end{equation}
Le défaut de cellule est ainsi construit, non ajouté comme un cardinal sans
application.

\section{Orientation et retenue de la cellule}
\label{sec:pdfv2-coh-orientation-carry}

Soient
\[
 \tau_1(u,v)=(v,u),
 \qquad
 P(w)_{ji}=w_{ij}.
\]
L'échange orienté est l'application de complexes
\begin{equation}\label{eq:pdfv2-coh-theta}
 \Theta_{\nu,N}:=(\tau_1,-P):
 C_{a_\nu,b_\nu,N}\longrightarrow C_{b_\nu,a_\nu,N},
\end{equation}
parce que
\begin{equation}\label{eq:pdfv2-coh-orientation}
 \kappa_{b_\nu,a_\nu,N}\tau_1
 =-P\kappa_{a_\nu,b_\nu,N}.
\end{equation}
En outre,
\(\Theta_{\nu^{\rm op},N}\Theta_{\nu,N}=I\). Le chemin opposé ne
s'identifie pas à l'original : il change le signe de degré zéro et conserve l'acte
d'inversion.

Pour \(r\in A_N\), soit
\(\langle r\rangle_N\in\{0,\ldots,3^N-1\}\) son représentant canonique.
L'emprunt/retenue local est
\begin{equation}\label{eq:pdfv2-coh-beta}
 \beta_N(r,s):=
 \frac{\langle r-s\rangle_N-\langle r\rangle_N+\langle s\rangle_N}
 {3^N}\in\{0,1\}.
\end{equation}
Ainsi,
\begin{equation}\label{eq:pdfv2-coh-carry-cell}
 \langle u_i-v_j\rangle_N
 =\langle u_i\rangle_N-\langle v_j\rangle_N
  +3^N\beta_N(u_i,v_j).
\end{equation}
La matrice
\[
 \operatorname{Carr}_{\nu,N}(u,v)
 :=\bigl(\beta_N(u_i,v_j)\bigr)_{ij}
\]
est conservé avec \(\kappa_{\nu,N}(u,v)\) : résidu et retenue sont deux
coordonnées distinctes.

Lorsque les tailles des ports proviennent des deux feuillets APP, l'admissibilité
de la cellule inclut l'identité entière
\begin{equation}\label{eq:pdfv2-coh-app-defect}
 (a_\nu-1)(b_\nu-1)
 =\rho_{\Pi,\nu}-\rho_{\Sigma,\nu}+1
  +9(c_{\Pi,\nu}-c_{\Sigma,\nu}).
\end{equation}
L'égalité est exigée uniquement pour une cellule ayant cette généalogie ; elle n'est pas affirmée
pour des entiers dépourvus de feuillets APP.

Si \(L_N(u)\) est le relèvement entier canonique de la matrice d'une flèche
\(u\), on définit la retenue de composition par
\begin{equation}\label{eq:pdfv2-coh-carry-composition}
 L_N(v)L_N(u)=L_N(vu)+3^N c_N(v,u).
\end{equation}
L'associativité de trois flèches donne le cocycle exact
\begin{equation}\label{eq:pdfv2-coh-carry-cocycle}
 c_N(w,vu)+L_N(w)c_N(v,u)
 =c_N(wv,u)+c_N(w,v)L_N(u).
\end{equation}

\section{Catégorie d'histoires et foncteur cellulaire}
\label{sec:pdfv2-coh-category}

Pour un préfixe \(x_{\leq d}\), soit \(J_d(x)\) la catégorie finie dont les
objets sont toutes ses cellules \(\nu\). Ses flèches sont engendrées par :
\begin{enumerate}
 \item des raffinements \(u:\nu\to\nu'\) avec des applications surjectives de
       ports
       \[
        \alpha_u:I_{\nu'}\twoheadrightarrow I_\nu,
        \qquad
        \beta_u:J_{\nu'}\twoheadrightarrow J_\nu;
       \]
 \item les échanges d'orientation \(\Theta_\nu\) ;
 \item les extensions nonadiques \(\Gamma_9\nu\) ;
 \item l'identité et la composition, soumises aux relations de frontière,
       de chemin, d'orientation et de retenue de l'histoire.
\end{enumerate}
La finitude se rapporte au préfixe ; le système de tous les préfixes n'est pas
déclaré fini.

Le foncteur cellulaire est
\begin{equation}\label{eq:pdfv2-coh-gamma-functor}
 \Gamma_{d,N}:J_d(x)\longrightarrow
 \operatorname{Ch}^{[0,1]}(A_N),
 \qquad
 \Gamma_{d,N}(\nu)=C_{\nu,N}.
\end{equation}
Pour une flèche de raffinement \(u:\nu\to\nu'\), il agit par
\begin{align}
 \Gamma_{d,N}(u)_1(p,q)
 &:=(p\circ\alpha_u,q\circ\beta_u),
 \label{eq:pdfv2-coh-gamma-one}\\
 \Gamma_{d,N}(u)_0(w)
 &:=w\circ(\alpha_u\times\beta_u).
 \label{eq:pdfv2-coh-gamma-zero}
\end{align}
On vérifie coordonnée par coordonnée
\begin{equation}\label{eq:pdfv2-coh-gamma-chain-map}
 \kappa_{\nu',N}\Gamma_{d,N}(u)_1
 =\Gamma_{d,N}(u)_0\kappa_{\nu,N}.
\end{equation}
Sur une inversion d'orientation, il agit par
\(\Theta_{\nu,N}\), et sur les compositions, il agit par composition, en gardant
séparément \(c_N(v,u)\).

La réduction \(r_{N+1,N}:A_{N+1}\twoheadrightarrow A_N\) satisfait
\begin{equation}\label{eq:pdfv2-coh-kappa-reduction}
 r_{N+1,N}^{I\times J}\kappa_{\nu,N+1}
 =\kappa_{\nu,N}
  \bigl(r_{N+1,N}^{I}\oplus r_{N+1,N}^{J}\bigr).
\end{equation}
Par conséquent, le diagramme des cellules est naturel simultanément en histoire et en
précision.

\fi
\section{Module des chemins, totalisation cellulaire et terminal nonadique}
\label{sec:pdfv2-coh-route-module}

Soit \(\operatorname{Term}_d(x)\) l'ensemble des terminaux de mémoire du
préfixe. Pour \(\nu\in J_d(x)\), définissons
\begin{equation}\label{eq:pdfv2-coh-w-module}
 W_{d,N}(\nu):=
 A_N\bigl[
 \operatorname{Path}_{J_d(x)}(\nu,\operatorname{Term}_d(x))
 \bigr].
\end{equation}
C'est le module libre sur tous les chemins terminaux ; il n'en sélectionne aucun. Si
\(u:\nu\to\nu'\), l'action contravariante est
\begin{equation}\label{eq:pdfv2-coh-w-action}
 W_{d,N}(u):W_{d,N}(\nu')\longrightarrow W_{d,N}(\nu),
 \qquad
 [\lambda]\longmapsto[\lambda\circ u].
\end{equation}
Chaque générateur conserve orientation, retenue, frontière, chemin et mémoire.

\section{Complexe bar bilatéral et totalisation de la bichaîne}
\label{sec:pdfv2-coh-bar}

Pour \(q\geq0\) et \(r\in\{0,1\}\), soit
\begin{equation}\label{eq:pdfv2-coh-bar-bigraded}
 B_{q,r}^{d,N}(x):=
 \bigoplus_{\nu_0\xrightarrow{u_1}\cdots\xrightarrow{u_q}\nu_q}
 W_{d,N}(\nu_q)\otimes_{A_N}\Gamma_{d,N}(\nu_0)_r.
\end{equation}
La différentielle bar est définie sans abréviation par
\begin{align}
 d_{\rm bar}(w;u_1,\ldots,u_q;c)
 :={}&(w;u_2,\ldots,u_q;\Gamma(u_1)c)
 \nonumber\\
 &+\sum_{i=1}^{q-1}(-1)^i
 (w;u_1,\ldots,u_{i+1}u_i,\ldots,u_q;c)
 \nonumber\\
 &+(-1)^q
 (W(u_q)w;u_1,\ldots,u_{q-1};c).
 \label{eq:pdfv2-coh-bar-differential}
\end{align}
Pour \(q=1\), c'est exactement la relation
\begin{equation}\label{eq:pdfv2-coh-coend-relation}
 d_{\rm bar}(w;u;c)
 =w\otimes\Gamma(u)c-W(u)w\otimes c.
\end{equation}
On identifie ainsi une incidence transportée à la même incidence écrite
dans la carte raffinée ; la relation ne reste pas cachée sous le nom du
totalisateur.

Le complexe total est
\begin{equation}\label{eq:pdfv2-coh-total-complex}
 K_{d,N}(x):=\operatorname{Tot}B_{\bullet,\bullet}^{d,N}(x),
 \qquad
 D_{\rm tot}:=d_{\rm bar}+(-1)^q\partial_\Gamma
 \quad\text{sur }B_{q,r}^{d,N}.
\end{equation}
La fonctorialité de \(W\) et \(\Gamma\), avec
\(\partial_\Gamma^2=0\), donne
\begin{equation}\label{eq:pdfv2-coh-total-square-zero}
 D_{\rm tot}^2=0.
\end{equation}
En notation compacte, mais désormais établie par
\eqref{eq:pdfv2-coh-bar-bigraded}--
\eqref{eq:pdfv2-coh-bar-differential},
\begin{equation}\label{eq:pdfv2-coh-derived-coend}
 K_{d,N}(x)
 \simeq
 W_{d,N}\otimes^{\mathbf L}_{A_N[J_d(x)]}\Gamma_{d,N}.
\end{equation}

\section{Colimite de profondeur et terminal nonadique}
\label{sec:pdfv2-coh-terminal}

Les préfixes forment la colimite
\begin{equation}\label{eq:pdfv2-coh-k-infinity}
 K_{\infty,N}(x):=\varinjlim_d K_{d,N}(x).
\end{equation}
L'extension de neuf phases fait passer de la profondeur \(d\) à la profondeur \(d+9\).
Ce n'est qu'après la formation de \eqref{eq:pdfv2-coh-k-infinity} qu'elle induit un
endomorphisme bien typé
\begin{equation}\label{eq:pdfv2-coh-u-gamma9}
 U_{\Gamma_9,N}:K_{\infty,N}(x)
 \longrightarrow K_{\infty,N}(x).
\end{equation}
Le terminal est
\begin{equation}\label{eq:pdfv2-coh-terminal-cone}
 \mathcal T_{{\rm coh},N}(x)
 :=\operatorname{Cone}\bigl(I-U_{\Gamma_9,N}\bigr).
\end{equation}
Avec la convention homologique
\[
 \operatorname{Cone}(f)_q
 =K_{\infty,N,q-1}\oplus K_{\infty,N,q},
\]
sa différentielle est
\begin{equation}\label{eq:pdfv2-coh-cone-differential}
 \partial_{\rm Cone}(a,b)
 =\bigl(-\partial_Ka,\partial_Kb+f(a)\bigr).
\end{equation}

Bien que la phase visible de \(h\) et \(\Gamma_9h\) coïncide, les générateurs
\([h]\) et \([\Gamma_9h]\) sont différents parce que leurs ledgers et leurs profondeurs
sont différents. Le terminal conserve
\begin{equation}\label{eq:pdfv2-coh-terminal-difference}
 (I-U_{\Gamma_9,N})[h]=[h]-[\Gamma_9h].
\end{equation}
Les réductions de précision commutent avec le terminal :
\begin{equation}\label{eq:pdfv2-coh-u-reduction}
 r_{N+1,N}U_{\Gamma_9,N+1}
 =U_{\Gamma_9,N}r_{N+1,N}.
\end{equation}

\section{Lecteur cohomologique terminal : cycle d'incidence et classe dans \(H_0(\operatorname{Cone}(I-U_{\Gamma_9,N}))\)}
\label{sec:pdfv2-coh-reader}

Chaque cellule de \(x_{\leq d}\) porte un vecteur d'incidence
\[
 \omega_{\nu,N}(x)
 \in A_N^{I_\nu\times J_\nu}=C_{\nu,N,0}
\]
et le chemin terminal complet \(\lambda_{\nu,x}\). Le cycle brut est
\begin{equation}\label{eq:pdfv2-coh-z-cycle}
 z_{d,N}(x):=
 \sum_{\nu\in\operatorname{Cell}_d(x)}
 \epsilon_\nu[\lambda_{\nu,x}]
 \otimes\omega_{\nu,N}(x)
 \in B_{0,0}^{d,N}(x).
\end{equation}
Comme le complexe total est homologique non négatif,
\begin{equation}\label{eq:pdfv2-coh-z-closed}
 D_{\rm tot}z_{d,N}(x)=0.
\end{equation}
On conserve aussi bien le cycle que sa classe
\[
 [z_{d,N}(x)]\in H_0(K_{d,N}(x)).
\]
Après son inclusion dans \(K_{\infty,N}\), sa classe terminale est
\[
 [(0,z_{d,N}(x))]
 \in H_0(\mathcal T_{{\rm coh},N}(x)).
\]
Le lecteur complet est
\begin{equation}\label{eq:pdfv2-coh-reader-full}
 \mathcal L_{{\rm coh},d,N}(x)
 :=\bigl(
 z_{d,N}(x),[z_{d,N}(x)],
 U_{\Gamma_9,N}z_{d,N}(x),
 [(0,z_{d,N}(x))]
 \bigr).
\end{equation}
On ne publie pas uniquement une classe quotient : le représentant, le
chemin qui l'a produit, sa continuation et le terminal subsistent.

\section{Fibres survivantes et corrections sans sélecteur}
\label{sec:pdfv2-coh-surviving-fibers}

Soit \(\mathscr S_\lambda^{\rm coh}\) le \(A_N\)-module libre engendré par les
histoires finies admissibles au niveau \(\lambda=(d,N)\). Le lecteur se
prolonge linéairement à
\begin{equation}\label{eq:pdfv2-coh-e-lambda}
 e_\lambda:\mathscr S_\lambda^{\rm coh}
 \longrightarrow\mathscr O_\lambda^{\rm coh},
 \qquad
 \mathscr O_\lambda^{\rm coh}
 :=H_0(K_{d,N})\oplus H_0(\mathcal T_{{\rm coh},N}).
\end{equation}
Le codomaine survivant est
\begin{equation}\label{eq:pdfv2-coh-surviving-output}
 \mathscr O_\infty^{\rm surv}
 :=\operatorname{Im}\!\left(
 e_\infty:\mathcal C_{\rm cont}^{\rm coh}
 \longrightarrow\varprojlim_\lambda\mathscr O_\lambda^{\rm coh}
 \right).
\end{equation}
Pour \(a=(a_\lambda)_\lambda\in\mathscr O_\infty^{\rm surv}\), on définit
\begin{equation}\label{eq:pdfv2-coh-fiber}
 \mathscr F_\lambda(a):=
 \left\{s_\lambda\in\mathscr S_\lambda^{\rm coh}:
 \begin{array}{l}
 e_\lambda(s_\lambda)=a_\lambda,\\
 \exists y\in\mathcal C_{\rm cont}^{\rm coh}\text{ avec }
 y|_\lambda=s_\lambda\text{ et }e_\infty(y)=a
 \end{array}
 \right\}.
\end{equation}
Ces fibres sont non vides par la définition du codomaine survivant. Si
\(\lambda'\geq\lambda\), la restriction est surjective :
\begin{equation}\label{eq:pdfv2-coh-fiber-surjection}
 \rho_{\lambda',\lambda}:\mathscr F_{\lambda'}(a)
 \twoheadrightarrow\mathscr F_\lambda(a).
\end{equation}
En effet, une histoire complète qui témoigne de
\(s_\lambda\in\mathscr F_\lambda(a)\) fournit par troncature un
antécédent à tout niveau ultérieur. Cet antécédent n'est pas choisi comme
partie de l'objet.

La multisection complète est
\begin{equation}\label{eq:pdfv2-coh-extension-multisection}
 \Sigma_{\lambda',\lambda}^{\rm coh}(s_\lambda;a)
 :=\{s_{\lambda'}\in\mathscr F_{\lambda'}(a):
 \rho_{\lambda',\lambda}s_{\lambda'}=s_\lambda\}.
\end{equation}
Pour un prolongement provisoire \(\widetilde s_{\lambda'}\), soit
\begin{equation}\label{eq:pdfv2-coh-discrepancy}
 \delta_{\lambda'}:=
 a_{\lambda'}-e_{\lambda'}(\widetilde s_{\lambda'}).
\end{equation}
La relation de toutes les corrections est
\begin{equation}\label{eq:pdfv2-coh-correction-relation}
 \operatorname{Corr}_{\lambda',\lambda}
 (\widetilde s_{\lambda'},\delta_{\lambda'})
 :=\left\{\eta\in\mathscr S_{\lambda'}^{\rm coh}:
 \begin{array}{l}
 \rho_{\lambda',\lambda}\eta=0,\\
 e_{\lambda'}(\eta)=\delta_{\lambda'},\\
 \eta\text{ est supportée sur de nouvelles cellules},\\
 \eta\text{ respecte l’orientation, la retenue, la mémoire et le terminal}
 \end{array}
 \right\}.
\end{equation}
Si deux relèvements sont comparables, leur différence
\(\eta=s_{\lambda'}-\widetilde s_{\lambda'}\) appartient à cette relation
exactement lorsqu'elle satisfait les quatre conditions précédentes. On n'affirme pas
que \eqref{eq:pdfv2-coh-correction-relation} soit non vide pour un
écart arbitraire : la non-vacuité déterminante est celle de
\eqref{eq:pdfv2-coh-extension-multisection} sur le domaine survivant.

\iffalse
\section{Les treize coordonnées de la restriction}
\label{sec:pdfv2-coh-thirteen}

Pour \(x_\lambda\in\mathcal C_{{\rm cont},\lambda}^{\rm coh}\), soit
\begin{equation}\label{eq:pdfv2-coh-d-package}
 \begin{aligned}
 D_{{\rm coh},\lambda}(x_\lambda):=\bigl(&
 \mathcal F_{{\rm coh},\lambda},
 \operatorname{Ext}_{\Gamma,{\rm coh},\lambda},
 \operatorname{Rel}_{{\rm coh},\lambda},
 \mathcal N_{{\rm coh},\lambda},
 \operatorname{or}_{{\rm coh},\lambda},
 \operatorname{Carr}_{{\rm coh},\lambda},
 \operatorname{Mem}_{{\rm coh},\lambda},\\
 &\partial_{{\rm coh},\lambda},
 \gamma_{{\rm coh},\lambda},
 \operatorname{Msect}_{{\rm coh},\lambda},
 \operatorname{Surv}_{{\rm coh},\lambda},
 \operatorname{Term}_{{\rm coh},\lambda},
 \mathcal L_{{\rm coh},\lambda}\bigr).
 \end{aligned}
\end{equation}
Ses coordonnées sont :
\begin{enumerate}
 \item \(\mathcal F_{{\rm coh},\lambda}\) : tous les modules
       \(A_N^{I_\nu}\), \(A_N^{J_\nu}\),
       \(A_N^{I_\nu\times J_\nu}\), les complexes \(C_{\nu,N}\), les
       modules \(W_{d,N}(\nu)\) et \(\mathscr S_\lambda^{\rm coh}\).
 \item \(\operatorname{Ext}_{\Gamma,{\rm coh},\lambda}\) : toutes les applications
       \(\Gamma(u)\), les réductions de coefficients, les troncatures, les extensions
       \(\Gamma_9\) et les relations \(\Sigma_{\lambda',\lambda}\).
 \item \(\operatorname{Rel}_{{\rm coh},\lambda}\) :
       \(\kappa\), \(\delta\), la suite exacte, les faces bar, la
       relation de cofin, l'orientation et les lois de composition.
 \item \(\mathcal N_{{\rm coh},\lambda}\):
       \[
       (d,N,3^N,a_\nu,b_\nu,(a_\nu-1)(b_\nu-1),
       \rho_{\Sigma,\nu},\rho_{\Pi,\nu},c_{\Sigma,\nu},c_{\Pi,\nu}).
       \]
 \item \(\operatorname{or}_{{\rm coh},\lambda}\) : les signes
       \(\epsilon_\nu\), les ordres des ports et les applications
       \(\Theta_{\nu,N}\).
 \item \(\operatorname{Carr}_{{\rm coh},\lambda}\) : les matrices
       \(\beta_N\), les cocycles \(c_N(v,u)\), les résidus et les quotients,
       conservés séparément.
 \item \(\operatorname{Mem}_{{\rm coh},\lambda}\) : le ledger hérité et le
       mot ordonné des raffinements, retenues, corrections et avancées
       \(\Gamma_9\), non seulement leur valeur agrégée.
 \item \(\partial_{{\rm coh},\lambda}\) :
       \(\kappa_{\nu,N}\), \(D_{\rm tot}\),
       \(\partial_{\rm Cone}\) et les frontières des cellules.
 \item \(\gamma_{{\rm coh},\lambda}\) : toutes les chaînes composables
       \(\nu_0\to\cdots\to\nu_q\) et leurs continuations terminales.
 \item \(\operatorname{Msect}_{{\rm coh},\lambda}\) : toutes les cellules,
       les chemins terminaux, les extensions \(\Sigma\) et les corrections
       \(\operatorname{Corr}\), sans sélecteur.
 \item \(\operatorname{Surv}_{{\rm coh},\lambda}\) : le système des fibres
       \(\mathscr F_\lambda(a)\) et des restrictions surjectives.
 \item \(\operatorname{Term}_{{\rm coh},\lambda}\) :
       \(\operatorname{Cone}(I-U_{\Gamma_9,N})\), avec \(z\), \(Uz\)
       et leur différence.
 \item \(\mathcal L_{{\rm coh},\lambda}\) :
       \((z,[z],Uz,[(0,z)])\) et toutes ses compatibilités de réduction.
\end{enumerate}

Pour \(\lambda'\geq\lambda\), le transport coordonné satisfait
\begin{equation}\label{eq:pdfv2-coh-d-naturality}
 u_{\lambda',\lambda}^{\rm coh}
 D_{{\rm coh},\lambda'}(x_{\lambda'})
 =D_{{\rm coh},\lambda}
  (\rho_{\lambda',\lambda}x_{\lambda'}).
\end{equation}

\section{Graphe de restriction et objet dépendant}
\label{sec:pdfv2-coh-graphs}

La restriction totale est le graphe
\begin{equation}\label{eq:pdfv2-coh-g-graph}
 \mathcal G_{\rm coh}:=\operatorname{Graph}(D_{\rm coh}),
 \qquad
 \operatorname{Res}_{\rm coh}(x):=(x,D_{\rm coh}x).
\end{equation}
Ainsi,
\begin{equation}\label{eq:pdfv2-coh-res-inverse}
 \pi_1\operatorname{Res}_{\rm coh}=I,
 \qquad
 \operatorname{Res}_{\rm coh}\pi_1=I_{\mathcal G_{\rm coh}}.
\end{equation}

L'extracteur complet est
\begin{equation}\label{eq:pdfv2-coh-chi}
 \chi_{\rm coh}(x):=
 \bigl(
 \{C_{\nu,N}\},
 \{z_\lambda,[z_\lambda]\},
 \{\mathscr F_\lambda(e_\infty x)\},
 \{\Sigma_{\lambda',\lambda}\},
 \{\operatorname{Corr}_{\lambda',\lambda}\}
 \bigr).
\end{equation}
On définit
\begin{equation}\label{eq:pdfv2-coh-xchi}
 X_{\chi,{\rm coh}}
 :=\{(\chi_{\rm coh}(x),x,D_{\rm coh}x):
 x\in\mathcal C_{\rm cont}^{\rm coh}\}
\end{equation}
et
\begin{equation}\label{eq:pdfv2-coh-prx}
 \operatorname{pr}_{X,{\rm coh}}(x,D_{\rm coh}x)
 :=(\chi_{\rm coh}(x),x,D_{\rm coh}x).
\end{equation}
La projection sur \((x,D_{\rm coh}x)\) est inverse sur cet objet
dépendant.

\section{Assembleur et publicateur internes}
\label{sec:pdfv2-coh-assembler-publisher}

L'espace ambiant \(\mathscr M_{\rm coh}^{\rm amb}\) se compose de systèmes
compatibles de suites exactes de cellules, de catégories, de foncteurs, de modules de
chemins, de totalisateurs, de terminaux, de lecteurs, de fibres et de multisections.
L'assembleur brut est
\begin{equation}\label{eq:pdfv2-coh-assembler}
 \begin{aligned}
 a_{\rm coh}:X_{\chi,{\rm coh}}&\longrightarrow
 \mathscr M_{\rm coh}^{\rm amb},\\
 a_{\rm coh}(\chi(x),x,Dx)
 &:={}
 \bigl(
 \{C_{\nu,N},\kappa_{\nu,N},\delta_{\nu,N}\},
 \{J_d,\Gamma,W,K\},\\
 &\hspace{31mm}
 \{U_{\Gamma_9,N},\mathcal T_N\},
 \{e_\lambda,\mathscr F_\lambda,\Sigma,\operatorname{Corr}\}
 \bigr)_x.
 \end{aligned}
\end{equation}
La macrostructure est un autre graphe :
\begin{equation}\label{eq:pdfv2-coh-macro-graph}
 \mathfrak M_{\rm coh}:=\operatorname{Graph}(a_{\rm coh}),
 \qquad
 \mathcal A_{\rm coh}(z):=(z,a_{\rm coh}z).
\end{equation}

Soit \(\mathscr Y_{\rm coh}^{\rm int}\) le type de certificats qui contient
\begin{equation}\label{eq:pdfv2-coh-certificate-list}
 \begin{gathered}
 \text{la sucesión exacta \eqref{eq:pdfv2-coh-exact-sequence}},\\
 \kappa_{b,a}\tau=-P\kappa_{a,b},
 \qquad D_{\rm tot}^2=0,
 \qquad\partial_{\rm Cone}^2=0,\\
 rU_{\Gamma_9}=U_{\Gamma_9}r,
 \qquad
 e_\lambda\rho_{\lambda',\lambda}
 =r_{\lambda',\lambda}e_{\lambda'},\\
 \rho_{\lambda',\lambda}:\mathscr F_{\lambda'}(a)
 \twoheadrightarrow\mathscr F_\lambda(a),
 \qquad
 \Sigma_{\lambda',\lambda}(s_\lambda;a)\neq\varnothing,
 \end{gathered}
\end{equation}
avec les cycles bruts, leurs classes et leurs terminaux. Le publicateur
\begin{equation}\label{eq:pdfv2-coh-publisher}
 p_{\rm coh}^{\rm raw}:\mathfrak M_{\rm coh}
 \longrightarrow\mathscr Y_{\rm coh}^{\rm int}
\end{equation}
produit ce certificat sans supprimer la macrostructure. Enfin,
\begin{equation}\label{eq:pdfv2-coh-output-graph}
 \mathcal C_{\rm coh}^{\rm HMT}
 :=\operatorname{Graph}(p_{\rm coh}^{\rm raw}),
 \qquad
 \mathcal P_{\rm coh}(m):=(m,p_{\rm coh}^{\rm raw}m).
\end{equation}

\section{Théorème interne de génération cohérente}
\label{sec:pdfv2-coh-theorem}

\begin{theorem}[Génération interne de la macrostructure de cohérence]
\label{thm:pdfv2-coh-internal-generation}
Sur le domaine explicite \(\mathcal C_{\rm cont}^{\rm coh}\), la chaîne
\[
 \mathcal C_{\rm cont}^{\rm coh}
 \xrightarrow{\operatorname{Res}_{\rm coh}}
 \mathcal G_{\rm coh}
 \xrightarrow{\operatorname{pr}_{X,{\rm coh}}}
 X_{\chi,{\rm coh}}
 \xrightarrow{\mathcal A_{\rm coh}}
 \mathfrak M_{\rm coh}
 \xrightarrow{\mathcal P_{\rm coh}}
 \mathcal C_{\rm coh}^{\rm HMT}
\]
construit naturellement les cellules, leurs cônes, le diagramme d'histoires, le
bar bilatéral, la totalisation, le terminal nonadique et les fibres
survivantes. Elle conserve les treize coordonnées et toutes les corrections comme
multisections. Chaque étape admet, sur son image, une projection qui récupère
l'étape précédente.
\end{theorem}

\begin{proof}
La \cref{prop:pdfv2-coh-cell-exact} construit chaque cellule sur \(A_N\), avec
noyau diagonal et défaut \((a_\nu-1)(b_\nu-1)\). L'identité
\eqref{eq:pdfv2-coh-orientation} fait de l'échange orienté une application de
complexes involutive. Les formules
\eqref{eq:pdfv2-coh-gamma-one}--\eqref{eq:pdfv2-coh-gamma-zero} prouvent que
chaque raffinement est une application de complexes, tandis que
\eqref{eq:pdfv2-coh-kappa-reduction} prouve la compatibilité en précision.

La précomposition des chemins rend \(W\) contravariant. Les identités
fonctorielles de \(W\) et \(\Gamma\) impliquent
\(d_{\rm bar}^2=0\) ; leur commutation avec la différentielle cellulaire donne
\(D_{\rm tot}^2=0\). Ainsi, \(K_{d,N}\) est construit par les formules bar,
non par une invocation nominale du cofin.

L'extension \(\Gamma_9\) augmente la profondeur. Elle n'induit donc
l'endomorphisme \(U_{\Gamma_9,N}\) qu'après le passage à la colimite
\(K_{\infty,N}\). Le cône
\(\operatorname{Cone}(I-U_{\Gamma_9,N})\) conserve alors la différence
\([h]-[\Gamma_9h]\) et n'identifie pas la répétition de phase à la répétition
d'état.

Le cycle \(z_{d,N}(x)\) se forme avec toutes les cellules, tous les signes, chemins et
incidences de l'histoire. Troncature et réduction commutent avec le lecteur. Si
\(s_\lambda\in\mathscr F_\lambda(a)\), l'histoire complète qui apparaît dans
la définition même de la fibre fournit un antécédent à tout niveau
\(\lambda'\geq\lambda\) ; cela prouve la surjectivité
\eqref{eq:pdfv2-coh-fiber-surjection}. L'objet conserve l'ensemble entier
\(\Sigma_{\lambda',\lambda}\), non un choix. Les différences entre
relèvements sont enregistrées par
\eqref{eq:pdfv2-coh-correction-relation}, avec support, orientation, retenue,
mémoire et terminal explicites.

Enfin, \(\mathcal G_{\rm coh}\), \(\mathfrak M_{\rm coh}\) et
\(\mathcal C_{\rm coh}^{\rm HMT}\) sont des graphes, tandis que
\(X_{\chi,{\rm coh}}\) conserve comme coordonnées \(x\) et \(D_{\rm coh}x\).
La première projection de chaque objet récupère le précédent. Par conséquent, la
composition complète ne remplace pas l'histoire par une classe, la multisection par un
sélecteur ni le terminal par la phase visible.
\end{proof}

\section{Provenance probatoire et falsificateurs}
\label{sec:pdfv2-coh-provenance-falsifiers}

\paragraph{Résultat récupéré.}
Appartiennent à l'architecture déjà construite la cellule \(\kappa_{a,b}\), son
noyau diagonal et son défaut, l'orientation
\(\kappa_{b,a}\tau=-P\kappa_{a,b}\), la retenue APP, le cône de la cellule, la
catégorie d'histoires, la totalisation par cofin, l'action \(\Gamma_9\), le
terminal \(\operatorname{Cone}(I-U_{\Gamma_9})\) et l'architecture
d'extension survivante.

\paragraph{Formalisation nouvelle.}
Sont des formalisations explicites la suite exacte
\eqref{eq:pdfv2-coh-exact-sequence} sur \(A_N\), la séparation entre
profondeur et précision, les faces complètes du bar, la formation du
terminal seulement après la colimite, la mise en paquet des treize coordonnées,
les corrections comme relation sans sélecteur et les quatre objets dépendants
ou graphes qui empêchent la perte de généalogie.

\begin{criterion}[Falsificateurs internes de la branche \(\mathrm{coh}\)]
\label{crit:pdfv2-coh-falsifiers}
La construction échoue si l'un quelconque des faits suivants se produit :
\begin{enumerate}
 \item la suite
       \eqref{eq:pdfv2-coh-exact-sequence} n'est pas exacte ;
 \item \(\kappa_{b,a}\tau=-P\kappa_{a,b}\) échoue ;
 \item un raffinement ne satisfait pas
       \(\kappa_{\nu'}\Gamma(u)_1=\Gamma(u)_0\kappa_\nu\) ;
 \item la retenue de composition ne satisfait pas
       \eqref{eq:pdfv2-coh-carry-cocycle} ;
 \item \(D_{\rm tot}^2\neq0\);
 \item la réduction, le lecteur ou le terminal ne commutent pas avec \(U_{\Gamma_9}\) ;
 \item on forme \(\operatorname{Cone}(I-U_{\Gamma_9})\) en identifiant d'abord
       les profondeurs \(d\) et \(d+9\) ;
 \item une fibre déclarée survivante est vide ou une restriction n'est pas
       surjective ;
 \item on remplace \(\Sigma_{\lambda',\lambda}\) par un unique
       relèvement ;
 \item on publie seulement \([z]\) et l'on supprime \(z\), le chemin, la retenue, la mémoire ou
       le terminal ;
 \item l'une quelconque des treize coordonnées est omise.
\end{enumerate}
Dans le dernier cas, le diagnostic contraignant est : \emph{objet substitué ;
conclusion non transférable}.
\end{criterion}

\begin{warningbox}
Ce chapitre conclut uniquement la construction interne de
\(\mathcal C_{\rm coh}^{\rm HMT}\). Il n'introduit pas d'espace classique, de
classe classique ni de publication classique. Tout entrelaceur ultérieur
appartient à un autre document et ne gouverne pas l'existence de cette
macrostructure.
\end{warningbox}
\fi
 \endgroup

Le complexe local \eqref{eq:pdfv2-cor-complejo-local} et le sommet \(y_\alpha\), extrémité de cette même extension dans le nerf, forment le produit fibré augmenté. On y définit
\begin{equation}
\label{eq:pdfv2-cor-zpm-h}
 \mathbf1_\alpha=(r,[y_\alpha]),\qquad
 z_-(\alpha)=(r,[y_\alpha]),
\end{equation}
\begin{equation}
\label{eq:pdfv2-cor-zplus-h}
 z_+(\alpha)=
 (r+3c_\alpha v_\alpha e+v_\alpha b,[y_\alpha]),
 \qquad
 h_*(\alpha)=(-v_\alpha f,0).
\end{equation}
L'action sur les générateurs donne, sans affirmation d'existence abstraite supplémentaire,
\begin{equation}
\label{eq:pdfv2-cor-filler}
 \partial z_-=\partial z_+=0,\qquad
 \partial h_*=z_--z_+.
\end{equation}
Le remplissage corrélatif, la cellule et le collier conservent le même \((\alpha,\operatorname{or}(\alpha),\kappa(\alpha),\partial\alpha,
\gamma\alpha,m(\gamma\alpha;m_0))\) ; ils appliquent seulement des lecteurs internes distincts au registre commun.

Pour le même complexe fini muni de bases, on définit, sans exiger de bilan ni d'acyclicité, la droite déterminant totale
\begin{equation}
\label{eq:pdfv2-cor-linea-determinante-total}
 \operatorname{Det}(C_{k,N}):=
 \bigotimes_q
 \left(\bigwedge^{\rm top}C_{k,N,q}\right)^{(-1)^q}.
\end{equation}
Le raffinement induit le morphisme déterminantal de l'application de complexes et transporte simultanément les modules d'homologie, les formes normales de Smith et les pages de Bockstein comme champs du registre. Ni l'existence de \eqref{eq:pdfv2-cor-linea-determinante-total} ni son transport n'exigent l'annulation de l'homologie. Si le cône du raffinement est contractile et orienté, le morphisme se trivialise canoniquement ; dans le cas général, le cône est conservé sans être remplacé par un isomorphisme fictif.

\subsection{Forme normale de Smith de l'opérateur terminal et décomposition primaire de son conoyau}
Pour chaque préfixe \(y\), la retenue non orientée et le coefficient orienté de frontière du registre commun admettent les relèvements
\begin{equation}
\label{eq:detint-c-v-lifts}
 \widetilde c_y,\widetilde v_y\in\mathbf Z,\qquad
 c_{y,n}:=\widetilde c_y\bmod3^n,\qquad
 v_{y,n}:=\widetilde v_y\bmod3^n.
\end{equation}
La sortie déterminantale n'est ni une entrée ni une branche séparée : c'est le secteur produit par la même réduction terminale lorsque les conditions typées sont satisfaites
\begin{equation}
\label{eq:detint-dominio}
\begin{split}
 \mathcal C_{\rm cont}^{\rm det}:=\bigl\{x\in\mathcal C_{\rm cont}^{\rm disc}:{}&
 x\text{ survit à toute profondeur};\\
 &\widehat P_x^{\rm red}\text{ est fini et basé};\\
 &\chi(\widehat P_x^{\rm red})=0;\\
 &H_*(\widehat P_x^{\rm red}\otimes_{\mathbf Z_3}K_3)=0;\\
 &\text{la réduction produit un bloc carré stable par raffinement}
 \bigr\}.
\end{split}
\end{equation}
\begingroup
  \let\section\subsection
  % Extracción quirúrgica del fuente teoremática V2 05e: líneas 325--447.
% El resto permanece en este archivo como procedencia, pero no se publica.
\iffalse
\chapter[Macrostructure déterminantielle interne]{Macrostructure déterminantielle interne engendrée par le continu discret}
\label{chap:detint-macroestructura}

\begin{thesisbox}
La branche \({\rm det}\) part exclusivement d'histoires survivantes de
\(\mathcal C_{\rm cont}^{\rm disc}\). Le nerf des extensions, la diagonale
Alexander--Whitney, le complexe local, l'unité, les deux publications, le
filler, la réduction terminale et la tour \(3\)-adique se construisent avant
toute évaluation ultérieure. Les conditions de bilan et d'acyclicité sont
incorporées au domaine typé et ne sont jamais tacitement présupposées.
\end{thesisbox}

\section{Domaine déterminantiel survivant}

Soit
\begin{equation}\label{eq:detint-an}
 A_n:=\mathbf Z/3^n\mathbf Z,
 \qquad
 \rho_{n+1,n}:A_{n+1}\longrightarrow A_n.
\end{equation}
Pour \(x\in\mathcal C_{\rm cont}^{\rm disc}\), soit
\(\mathsf T_m(x)\) la catégorie finie des préfixes TPK de profondeur \(m\)
compatibles avec \(x\). Chaque objet conserve feuillet, orientation, résidu,
quotient, retenue entière, mémoire, frontière, chemin, multisection, survie
et terminal ; ses morphismes sont les extensions admissibles qui conservent ces
coordonnées.

De chaque objet \(y\in\mathsf T_m(x)\), on extrait les entiers internes
\begin{equation}\label{eq:detint-c-v-lifts}
 \widetilde c_y,\widetilde v_y\in\mathbf Z,
 \qquad
 c_{y,n}:=\widetilde c_y\bmod3^n,
 \qquad
 v_{y,n}:=\widetilde v_y\bmod3^n.
\end{equation}
La première coordonnée est la retenue non orientée et la seconde, le coefficient
orienté de frontière. Le renversement satisfait
\begin{equation}\label{eq:detint-or-cv}
 c_{\bar y,n}=c_{y,n},
 \qquad
 v_{\bar y,n}=-v_{y,n}.
\end{equation}

On écrit
\begin{equation}\label{eq:detint-z3-k3}
 \mathbf Z_3:=\varprojlim_n A_n,
 \qquad
 K_3:=\operatorname{Frac}(\mathbf Z_3).
\end{equation}
La restriction déterminantielle est définie par
\begin{equation}\label{eq:detint-dominio}
\begin{split}
 \mathcal C_{\rm cont}^{\rm det}:=
 \bigl\{x\in\mathcal C_{\rm cont}^{\rm disc}:{}&
 x\text{ survit à toute profondeur};\\
 &\widehat P_x^{\rm red}\text{ est fini et basé};\\
 &\chi(\widehat P_x^{\rm red})=0;\\
 &H_*(\widehat P_x^{\rm red}\otimes_{\mathbf Z_3}K_3)=0;\\
 &\text{la réduction canonique produit un bloc carré}\\
 &\text{stable par raffinement}\bigr\}.
\end{split}
\end{equation}
Les objets \(\widehat P_x^{\rm red}\) et le bloc carré seront construits
dans \cref{sec:detint-terminal}. Par conséquent,
\eqref{eq:detint-dominio} est une condition d'appartenance vérifiable et non
une conclusion insérée dans la définition des applications.

\section{Nerf, Alexander--Whitney et counité}

On définit le complexe normalisé
\begin{equation}\label{eq:detint-gmn}
 G_{m,n}(x):=
 N_*^{\rm norm}\!\left(N\mathsf T_m(x);A_n\right).
\end{equation}
Pour un simplexe non dégénéré
\(\sigma=[y_0\to y_1\to\cdots\to y_q]\), la diagonale est
\begin{equation}\label{eq:detint-aw}
 \Delta_{\rm AW}(\sigma)
 =\sum_{j=0}^q
 [y_0\to\cdots\to y_j]\otimes
 [y_j\to\cdots\to y_q].
\end{equation}
La counité est donnée par
\begin{equation}\label{eq:detint-counidad-g}
 \epsilon_G([y])=1,
 \qquad
 \epsilon_G(\sigma)=0\quad(q>0).
\end{equation}
En particulier,
\begin{equation}\label{eq:detint-vertice-grouplike}
 \Delta_{\rm AW}[y]=[y]\otimes[y],
 \qquad
 \epsilon_G([y])=1.
\end{equation}

Si \(m'\ge m\) et \(n'\ge n\), la troncature et la réduction des coefficients
induisent
\begin{equation}\label{eq:detint-rmn}
 R_{m',m}^{n',n}:G_{m',n'}(x)\longrightarrow G_{m,n}(x),
\end{equation}
et l'on vérifie
\begin{equation}\label{eq:detint-aw-natural}
 R\partial=\partial R,
 \qquad
 (R\otimes R)\Delta_{\rm AW}=\Delta_{\rm AW}R,
 \qquad
 \epsilon_GR=\rho_{n',n}\epsilon_G.
\end{equation}
Le chemin ne se réduit pas à ses extrémités : le simplexe entier reste dans
\(G_{m,n}(x)\).

\section{Complexe local interne}

Pour \(c\in A_n\), on définit
\begin{equation}\label{eq:detint-p0-grados}
 P^0_{n,1}(c)=A_nf,
 \qquad
 P^0_{n,0}(c)=A_nr\oplus A_ne\oplus A_nb,
 \qquad
 P^0_{n,-1}(c)=A_ng,
\end{equation}
avec différentielle
\begin{equation}\label{eq:detint-p0-diferencial}
 \partial f=3ce+b,
 \qquad
 \partial r=0,
 \qquad
 \partial e=g,
 \qquad
 \partial b=-3cg.
\end{equation}
L'identité de chaîne est littérale :
\begin{equation}\label{eq:detint-d2}
 \partial^2f=3c\,\partial e+\partial b=3cg-3cg=0,
 \qquad
 \partial^2e=\partial^2b=\partial^2r=0.
\end{equation}

L'augmentation
\begin{equation}\label{eq:detint-epsilon-p}
 \epsilon_P:P_n^0(c)\longrightarrow A_n[0]
\end{equation}
est fixée par
\[
 \epsilon_P(r)=1,
 \qquad
 \epsilon_P(e)=\epsilon_P(b)=\epsilon_P(f)=\epsilon_P(g)=0.
\]
Le complexe possède la coalgèbre différentielle
\begin{equation}\label{eq:detint-coalgebra-p}
 \Delta_P(r)=r\otimes r,
 \qquad
 \Delta_P(q)=q\otimes r+r\otimes q
 \quad(q\in\{e,b,f,g\}).
\end{equation}
Ainsi, \(r\) est group-like et les quatre autres générateurs sont primitifs
relativement à \(r\). Les équations
\eqref{eq:detint-p0-diferencial} prouvent que \(\Delta_P\) commute avec la
différentielle tensorielle.

\section{Système local d'extensions et ledger de mémoire}

Pour une extension \(\alpha:y\to y'\), nous définissons
\begin{equation}\label{eq:detint-psi}
 \Psi_\alpha:P_n^0(c_y)\longrightarrow P_n^0(c_{y'})
\end{equation}
par
\begin{equation}\label{eq:detint-psi-generadores}
\begin{gathered}
 \Psi_\alpha(r_y)=r_{y'},\quad
 \Psi_\alpha(e_y)=e_{y'},\quad
 \Psi_\alpha(g_y)=g_{y'},\quad
 \Psi_\alpha(f_y)=f_{y'},\\
 \Psi_\alpha(b_y)=b_{y'}+3(c_{y'}-c_y)e_{y'}.
\end{gathered}
\end{equation}
C'est une application de complexes parce que
\begin{equation}\label{eq:detint-psi-b-chain}
 \partial\Psi_\alpha(b_y)
 =-3c_{y'}g+3(c_{y'}-c_y)g
 =-3c_yg
 =\Psi_\alpha(\partial b_y),
\end{equation}
et
\begin{equation}\label{eq:detint-psi-f-chain}
 \Psi_\alpha(\partial f_y)
 =3c_ye+b+3(c_{y'}-c_y)e
 =3c_{y'}e+b
 =\partial\Psi_\alpha(f_y).
\end{equation}
De plus,
\begin{equation}\label{eq:detint-psi-functorial}
 \Psi_\beta\Psi_\alpha=\Psi_{\beta\alpha}.
\end{equation}

Le changement de la coordonnée orientée n'est pas effacé. Il est enregistré par
\begin{equation}\label{eq:detint-k-alpha}
 K_\alpha:=(v_{y'}-v_y)f_{y'}.
\end{equation}
Pour des extensions composables,
\begin{equation}\label{eq:detint-k-coherencia}
 K_{\beta\alpha}=K_\beta+\Psi_\beta K_\alpha.
\end{equation}
Cette égalité est le ledger exact de l'accroissement de frontière et rend visible
la mémoire qu'une composition purement nominale perdrait.

Le diagramme local complet est la construction de Grothendieck
\begin{equation}\label{eq:detint-grothendieck-local}
 \int_{y\in\mathsf T_m(x)}P_n^0(c_y).
\end{equation}
Dans celle-ci, une flèche au-dessus de \(\alpha:y\to y'\) porte comme composante linéaire
\(\Psi_\alpha\) et comme cohérence \(K_\alpha\). Lorsqu'elle est appliquée à un
élément supporté en un sommet, on utilise la notation
\begin{equation}\label{eq:detint-psi-soporte-vertice}
 \Psi_\alpha(p,[y]):=(\Psi_\alpha p,[y']).
\end{equation}
L'arête \([y\to y']\) reste dans la coordonnée de chemin du nerf. Cette
convention sera celle utilisée dans les formules de naturalité de
\(z_\pm\) et \(h_*\) ; on n'affirme pas une application globale qui efface les autres
simplexes de \(G_{m,n}(x)\).

\section{Pullback, unité, racine, publications et filler}

Le pullback de complexes est
\begin{equation}\label{eq:detint-pgen}
 P^{\rm gen}_{y;m,n}
 :=P_n^0(c_y)\times_{A_n[0]}G_{m,n}(x)
 =\{(p,\gamma):\epsilon_P(p)=\epsilon_G(\gamma)\}.
\end{equation}
Sa différentielle est
\begin{equation}\label{eq:detint-pgen-d}
 \partial(p,\gamma)=(\partial p,\partial\gamma).
\end{equation}
L'unité commune associée au sommet \(y\) est
\begin{equation}\label{eq:detint-unidad}
 \mathbf1_y:=(r,[y]).
\end{equation}
Ses deux projections sont group-like par
\eqref{eq:detint-vertice-grouplike} et
\eqref{eq:detint-coalgebra-p} ; c'est le sens précis d'unité
group-like compatible du pullback.

La projection racine est l'application de complexes
\begin{equation}\label{eq:detint-proyeccion-root}
 \varpi_{\rm root}:P^{\rm gen}_{y;m,n}\longrightarrow
 \operatorname{Root}_{y;m,n},
 \qquad
 \varpi_{\rm root}(p,\gamma):=(\epsilon_P(p)r,\gamma).
\end{equation}
On définit les deux publications internes
\begin{equation}\label{eq:detint-zpm}
 z_-(y):=(r,[y]),
 \qquad
 z_+(y):=(r+3c_yv_ye+v_yb,[y]).
\end{equation}
Toutes deux sont des cycles :
\begin{equation}\label{eq:detint-zpm-ciclos}
 \partial z_-(y)=0,
 \qquad
 \partial z_+(y)=3c_yv_yg-3c_yv_yg=0.
\end{equation}
Leurs racines coïncident :
\begin{equation}\label{eq:detint-raiz-comun}
 \varpi_{\rm root}z_-(y)
 =(r,[y])
 =\varpi_{\rm root}z_+(y).
\end{equation}

Le filler se construit directement :
\begin{equation}\label{eq:detint-hstar}
 h_*(y):=(-v_yf,0)\in P^{\rm gen}_{y;m,n,1}.
\end{equation}
Alors
\begin{equation}\label{eq:detint-filler-identidad}
 \partial h_*(y)
 =(-3c_yv_ye-v_yb,0)
 =z_-(y)-z_+(y).
\end{equation}
En particulier,
\begin{equation}\label{eq:detint-clases-iguales}
 [z_-(y)]=[z_+(y)]
 \quad\text{dans}\quad H_0(P^{\rm gen}_{y;m,n}),
\end{equation}
et l'égalité conserve sa chaîne témoin.

Sous une extension \(\alpha:y\to y'\), on a
\begin{equation}\label{eq:detint-z-natural}
 z_-(y')=\Psi_\alpha z_-(y),
 \qquad
 z_+(y')-\Psi_\alpha z_+(y)=\partial K_\alpha,
\end{equation}
et
\begin{equation}\label{eq:detint-h-natural}
 h_*(y')-\Psi_\alpha h_*(y)=-K_\alpha.
\end{equation}
Par conséquent, la naturalité n'identifie pas des valeurs distinctes de \(v\) : elle enregistre
leur différence au moyen de l'homotopie cohérente \(K_\alpha\).

\section{Orientation interne}

Sur \(P_n^0(c)\), on définit
\begin{equation}\label{eq:detint-omega-p}
 \Omega(r)=r,
 \qquad
 \Omega(q)=-q\quad(q\in\{e,b,f,g\}).
\end{equation}
C'est un automorphisme de complexes et de la coalgèbre relative. Dans le nerf,
l'inversion orientée est
\begin{equation}\label{eq:detint-omega-g}
 \mathfrak o_G[y_0\to\cdots\to y_q]
 :=(-1)^{q(q+1)/2}
 [\bar y_q\to\cdots\to\bar y_0].
\end{equation}
Avec \eqref{eq:detint-or-cv}, ces actions donnent
\begin{equation}\label{eq:detint-or-z-h}
 \Omega z_-(y)=z_-(\bar y),
 \qquad
 \Omega z_+(c_y,v_y)=z_+(c_{\bar y},v_{\bar y}),
 \qquad
 \Omega h_*(y)=h_*(\bar y).
\end{equation}
L'orientation agit sur le chemin, le complexe, les publications, le filler et la droite
déterminantielle ; ce n'est pas une étiquette sans action.

\fi
\section{Relèvement du conoyau au moyen de la tour de Bockstein et conservation de la retenue}
\label{sec:detint-terminal}

Seule la racine commune est supprimée :
\begin{equation}\label{eq:detint-p-reducido}
 \overline P_{y;m,n}
 :=P^{\rm gen}_{y;m,n}/A_n\mathbf1_y,
 \qquad
 \widehat P_{y;m}^{\rm red}:=\varprojlim_n\overline P_{y;m,n}.
\end{equation}
La base est ordonnée d'abord selon l'ordre TPK des chemins, puis selon
\begin{equation}\label{eq:detint-orden-base}
 f\prec e\prec b\prec g.
\end{equation}
On exécute l'algorithme interne suivant.
\begin{enumerate}
 \item Chercher les coefficients qui sont des unités de \(\mathbf Z_3\).
 \item Choisir le premier selon \eqref{eq:detint-orden-base} et l'ordre des
 chemins.
 \item Annuler le couple correspondant et enregistrer les opérations
 élémentaires, leur signe et leur contribution à l'orientation.
 \item Répéter jusqu'à ce qu'il ne reste aucun pivot unitaire annulable.
\end{enumerate}
La condition de survie de \eqref{eq:detint-dominio} exige que,
cofinalement en \(m\), le résultat soit un bloc à deux termes
\begin{equation}\label{eq:detint-sx}
 \mathbf Z_3^{r_x}\xrightarrow{S_x}\mathbf Z_3^{r_x},
 \qquad
 \det S_x\ne0\text{ dans }K_3,
\end{equation}
et qu'un raffinement n'ajoute que des blocs contractiles unitaires et des changements
de base enregistrés. Si ces conditions échouent, l'histoire n'appartient pas à
\(\mathcal C_{\rm cont}^{\rm det}\) ; on ne force pas un terminal inexistant.

Le chemin identité fournit un témoin élémentaire. Après le retrait de la racine,
le complexe local est
\begin{equation}\label{eq:detint-testigo-contractil}
 A_nf\xrightarrow{\binom{3c}{1}}
 A_ne\oplus A_nb\xrightarrow{(1,-3c)}A_ng.
\end{equation}
Il est contractile par
\begin{equation}\label{eq:detint-contraccion-local}
 s(g)=e,
 \qquad
 s(ae+\beta b)=\beta f,
 \qquad
 s(f)=0.
\end{equation}
Ainsi, les conditions terminales admettent au moins le témoin identité
interne et ne décrivent pas un sous-domaine formellement vide.

\section{Smith, ordre \texorpdfstring{\(3\)}{3}-adique et unité initiale}

Dans des bases orientées, une diagonalisation par valuation a la forme
\begin{equation}\label{eq:detint-smith}
 U_xS_xV_x
 =\operatorname{diag}(3^{a_1}u_1,\ldots,3^{a_r}u_r),
 \qquad
 a_i\in\mathbf N,quad u_i\in\mathbf Z_3^\times.
\end{equation}
On définit
\begin{equation}\label{eq:detint-orden-d}
 d_x:=v_3(\det S_x)=\sum_{i=1}^r a_i
\end{equation}
et l'unité initiale
\begin{equation}\label{eq:detint-leading-unit}
 \lambda_x:=3^{-d_x}\det S_x\in\mathbf Z_3^\times.
\end{equation}
La coordonnée d'orientation trivialise la droite déterminantielle. Si une
transition produit
\begin{equation}\label{eq:detint-estabilidad-bloques}
 S_x\oplus I_q=U S_{x'}V,
\end{equation}
l'orientation est transportée par le facteur inverse de \(\det U\det V\).
L'unité orientée \(\lambda_x^{\rm or}\) est alors naturelle et ne change pas
lors de l'ajout de blocs contractiles.

En précision finie,
\begin{equation}\label{eq:detint-smith-finito}
 S_{x,n}=S_x\bmod3^n,
 \qquad
 a_i^{(n)}=\min(a_i,n).
\end{equation}
Le terminal est donc la famille compatible de toutes les précisions et
non un entier isolé.

\section{Tour de Bockstein et retenue d'ordre supérieur}

Soit
\[
 C_x^1\xrightarrow{S_x}C_x^0
\]
le bloc terminal. Si une classe \([\bar y]\) admet un relèvement
\(y\in C_x^1\) avec \(S_xy\in3^qC_x^0\), nous définissons
\begin{equation}\label{eq:detint-beta-q}
 \beta_q[\bar y]
 :=\left[3^{-q}S_xy\bmod3\right].
\end{equation}
Le changement de relèvement modifie le représentant par un bord de la page
correspondante, de sorte que \(\beta_q\) est bien défini. La retenue
supérieure est
\begin{equation}\label{eq:detint-carry-superior}
 \operatorname{Carr}_q(y):=3^{-q}S_xy.
\end{equation}
Elle ne remplace pas la retenue TPK de \eqref{eq:detint-c-v-lifts} ; elle enregistre son
prolongement dans la tour des coefficients.

Pour un bloc \(3^{a_i}u_i\), la classe persiste jusqu'à la page \(a_i\),
et à cette page, la différentielle non nulle est la multiplication par
\(\bar u_i\in\mathbf F_3^\times\). Par conséquent,
\begin{equation}\label{eq:detint-bock-longitudes}
 \{\text{longueurs de survie}\}=\{a_1,\ldots,a_r\},
 \qquad
 d_x=\sum_i a_i.
\end{equation}
Le produit orienté des trivialisations de page satisfait
\begin{equation}\label{eq:detint-bock-leading}
 \Theta_{\rm Bock}(x)=\overline{\lambda_x^{\rm or}}
 \in\mathbf F_3^\times.
\end{equation}
La preuve se fait dans chaque bloc diagonal et se transporte par les changements
unimodulaires enregistrés. Smith, Bockstein et l'unité initiale sont ainsi trois
lectures du même terminal.

\iffalse
\section{Les treize coordonnées de la restriction déterminantielle}

Pour chaque \((m,n)\), on définit
\begin{equation}\label{eq:detint-d-trece}
 D_{{\rm det},m,n}(x):=
 (\mathcal F,\operatorname{Ext}_\Gamma,\operatorname{Rel},
 \mathcal N,\operatorname{or},\operatorname{Carr},\operatorname{Mem},
 \partial,\gamma,\operatorname{Msect},\operatorname{Surv},
 \operatorname{Term},\mathcal L)_{{\rm det},m,n}(x),
\end{equation}
avec le contenu suivant.
\begin{enumerate}
 \item La fibre est
 \[
  \mathcal F=(A_n,\mathsf T_m(x),
  \{P_n^0(c_y)\}_y,\{P^{\rm gen}_{y;m,n}\}_y).
 \]
 \item \(\operatorname{Ext}_\Gamma\) contient tous les morphismes
 \(\alpha\), les applications \(\Psi_\alpha\), les homotopies \(K_\alpha\) et
 les applications \(R_{m',m}^{n',n}\).
 \item \(\operatorname{Rel}\) contient \(\partial^2=0\), l'équation de
 pullback, AW, la counité, \(\Psi_\beta\Psi_\alpha=\Psi_{\beta\alpha}\) et
 \(K_{\beta\alpha}=K_\beta+\Psi_\beta K_\alpha\).
 \item
 \(\mathcal N=(m,n,\widetilde c_y,\widetilde v_y,c_{y,n},v_{y,n},
 r_x,a_1,\ldots,a_{r_x},d_x)\).
 \item \(\operatorname{or}\) contient \(y\mapsto\bar y\),
 \(\mathfrak o_G\), \(\Omega\) et l'orientation de la droite
 déterminantielle.
 \item \(\operatorname{Carr}\) conserve séparément la retenue entière,
 sa réduction \(c_{y,n}\) et toutes les retenues supérieures
 \(\operatorname{Carr}_q\).
 \item \(\operatorname{Mem}\) contiene
 \[
  y_0\to\cdots\to y_m,
  \quad\{(\widetilde c_{y_j},\widetilde v_{y_j},\mathbf1_{y_j})\}_{j=0}^m,
  \quad\{K_{\alpha_j}\}.
 \]
 \item
 \(\partial=(\partial_G,\partial_P,\partial_{P^{\rm gen}},
 \partial_{\overline P},h_*)\).
 \item \(\gamma=[y_0\to\cdots\to y_q]\) conserve le simplexe orienté
 complet.
 \item \(\operatorname{Msect}\) contient toutes les extensions
 compatibles, tous les relèvements \(3\)-adiques et toutes les présentations
 unimodulaires du même terminal orienté ; il n'en choisit aucune a posteriori.
 \item \(\operatorname{Surv}\) enregistre la profondeur arbitraire,
 la compatibilité des coefficients, le bilan, l'acyclicité sur \(K_3\),
 la stabilité par blocs unitaires et la permanence de la racine.
 \item
 \[
  \operatorname{Term}=(\widehat P_x^{\rm red},S_x,
  \{a_i,u_i\},d_x,\lambda_x^{\rm or},
  \{\beta_q\},\Theta_{\rm Bock}).
 \]
 \item Le lecteur interne est
 \[
  \mathcal L_{\rm det}^{\rm int}
  =(\mathbf1_y,
  \varpi_{\rm root}z_-=\varpi_{\rm root}z_+,
  [z_-]=[z_+],d_x,\lambda_x^{\rm or},\Theta_{\rm Bock}).
 \]
\end{enumerate}

Les transitions satisfont, coordonnée par coordonnée,
\begin{equation}\label{eq:detint-naturalidad-trece}
 u^{\rm det}_{(m',n'),(m,n)}D_{{\rm det},m',n'}(x')
 =D_{{\rm det},m,n}(\tau_{m',m}x').
\end{equation}

\section{Assembleur, publicateur et chaîne non ablative}

On définit
\begin{equation}\label{eq:detint-graph-d}
 \mathcal G_{\rm det}:=\operatorname{Graph}(D_{\rm det}),
 \qquad
 \operatorname{Res}_{\rm det}(x):=(x,D_{\rm det}x).
\end{equation}
Soit \(\chi_{\rm det}(x)\) la multisection complète de germes
\[
 (y,\widetilde c_y,\widetilde v_y,\operatorname{or}_y,\gamma_y)
\]
compatibles à toute profondeur. L'objet dépendant est
\begin{equation}\label{eq:detint-x-dependiente}
 X_{\chi,{\rm det}}
 :=\{(\chi_{\rm det}(x),x,D_{\rm det}x):
 x\in\mathcal C_{\rm cont}^{\rm det}\},
\end{equation}
et
\begin{equation}\label{eq:detint-prx}
 \operatorname{pr}_{X,{\rm det}}(x,D_{\rm det}x)
 :=(\chi_{\rm det}(x),x,D_{\rm det}x).
\end{equation}

L'assembleur brut
\begin{equation}\label{eq:detint-ensamblador-crudo}
 a_{\rm det}:X_{\chi,{\rm det}}\longrightarrow
 \mathscr M_{\rm det}^{\rm amb}
\end{equation}
produit
\begin{equation}\label{eq:detint-ensamblador-salida}
\begin{split}
 a_{\rm det}(z)=\bigl(&
 \{G_{m,n},\Delta_{\rm AW},\epsilon_G\}_{m,n};\\
 &\{P_n^0(c),\Psi_\alpha,K_\alpha,P^{\rm gen},
 \mathbf1,z_-,z_+,h_*\}_{m,n};\\
 &\{\widehat P^{\rm red},S,d,\lambda^{\rm or},
 \beta_q,\Theta_{\rm Bock}\}\bigr).
\end{split}
\end{equation}
Le macro-objet conserve son entrée :
\begin{equation}\label{eq:detint-graph-a}
 \mathfrak M_{\rm det}:=\operatorname{Graph}(a_{\rm det}),
 \qquad
 \mathcal A_{\rm det}(z):=(z,a_{\rm det}z).
\end{equation}

Le publicateur brut
\begin{equation}\label{eq:detint-publicador-crudo}
 p_{\rm det}:\mathfrak M_{\rm det}\longrightarrow\mathscr Y_{\rm det}
\end{equation}
publie conjointement le certificat
\begin{equation}\label{eq:detint-certificado-publicado}
\begin{gathered}
 \partial^2=0,quad \text{AW et counité},quad
 \mathbf1_y\text{ group-like compatible},\\
 \partial z_\pm=0,quad
 \varpi_{\rm root}z_-=\varpi_{\rm root}z_+,quad
 \partial h_*=z_--z_+,\\
 \Psi_\beta\Psi_\alpha=\Psi_{\beta\alpha},quad
 K_{\beta\alpha}=K_\beta+\Psi_\beta K_\alpha,\\
 d=v_3(\det S)=\sum_i a_i,quad
 \Theta_{\rm Bock}=\overline{\lambda^{\rm or}},quad
 \text{naturalité complète par raffinement}.
\end{gathered}
\end{equation}
Enfin,
\begin{equation}\label{eq:detint-graph-p}
 \mathcal C_{\rm det}^{\rm HMT}:=\operatorname{Graph}(p_{\rm det}),
 \qquad
 \mathcal P_{\rm det}(M):=(M,p_{\rm det}M).
\end{equation}
La chaîne intégrale est
\begin{equation}\label{eq:detint-cadena-completa}
 \mathcal C_{\rm cont}^{\rm disc}\supset
 \mathcal C_{\rm cont}^{\rm det}
 \xrightarrow{\operatorname{Res}_{\rm det}}\mathcal G_{\rm det}
 \xrightarrow{\operatorname{pr}_{X,{\rm det}}}X_{\chi,{\rm det}}
 \xrightarrow{\mathcal A_{\rm det}}\mathfrak M_{\rm det}
 \xrightarrow{\mathcal P_{\rm det}}\mathcal C_{\rm det}^{\rm HMT}.
\end{equation}
Chaque flèche possède un inverse sur son image donné par une projection de
coordonnées. On ne remplace pas l'histoire par une valeur, la multisection par un sélecteur,
le complexe par un déterminant ni la preuve par un nom.

\section{Théorème de réalisation déterminantielle interne}

\begin{theorem}[Réalisation déterminantielle interne complète]
\label{thm:detint-realizacion}
Pour tout \(x\in\mathcal C_{\rm cont}^{\rm det}\), la chaîne
\eqref{eq:detint-cadena-completa} construit un macro-objet naturel en
profondeur et en précision. Il contient deux cycles à racine commune, un filler
explicite qui les identifie, une unité group-like compatible et un terminal
\(3\)-adique dont l'ordre, la forme de Smith et la tour de Bockstein produisent le même
lecteur orienté. Les treize coordonnées de \eqref{eq:detint-d-trece}
restent récupérables en sortie.
\end{theorem}

\begin{proof}
Les équations \eqref{eq:detint-p0-diferencial}--
\eqref{eq:detint-d2} prouvent que \(P_n^0(c)\) est un complexe. Les deux
augmentations sont des applications de complexes, de sorte que le pullback est typé.
Alexander--Whitney est naturel et rend chaque sommet group-like ; avec
\(r\), cela donne l'unité commune. Le calcul direct
\eqref{eq:detint-zpm-ciclos}--\eqref{eq:detint-filler-identidad} prouve que
les publications sont des cycles, ont la même racine et sont reliées par
\(h_*\). Les formules de \(\Psi_\alpha\) prouvent la fonctorialité, et
\(K_\alpha\) prouve la naturalité homotopique en conservant l'accroissement de
mémoire. La définition de \(\mathcal C_{\rm cont}^{\rm det}\) garantit que
la réduction terminale produit un bloc carré sans le postuler hors du
domaine. Smith sur \(\mathbf Z_3\) donne \(d=\sum_i a_i\) ; le calcul de chaque
bloc \(3^{a_i}u_i\) donne
\(\Theta_{\rm Bock}=\overline{\lambda^{\rm or}}\). Troncature, réduction et
ajout de blocs unitaires conservent ces identités. Enfin, les
trois constructions par graphe retiennent l'entrée comme première coordonnée,
de sorte que la composition complète est non ablative.
\end{proof}

\section{Provenance et falsificateurs}

\paragraph{Provenance.}
Le nerf TPK, Alexander--Whitney, la counité, l'unité commune, la forme
normale \(\partial f=3ce+b\), \(\partial e=g\),
\(\partial b=-3cg\), la racine commune, le filler et l'architecture de tour
ont le statut \texttt{ARCHITECTURE\_AUTORALE\_PRÉEXISTANTE}. Le calcul
littéral des deux cycles, de leur racine et de l'identité
\(\partial h_*=z_--z_+\) a le statut \texttt{RÉSULTAT\_RÉCUPÉRÉ}. Le
système local \(\Psi_\alpha/K_\alpha\), le sous-domaine survivant
équilibré, le terminal exclusivement \(3\)-adique et la mise en paquet par
graphes ont le statut \texttt{FORMALISATION\_NOUVELLE} ; ils formalisent sans changer
la provenance conceptuelle de l'appareil.

\begin{ablation}[Falsificateurs internes de la branche déterminantielle]
\label{abl:detint-falsadores}
La construction est falsifiée au niveau correspondant si l'un des
faits suivants se produit.
\begin{enumerate}
 \item \(\partial^2=0\) échoue ; le complexe local n'existe alors pas.
 \item Alexander--Whitney ne commute pas avec le raffinement, la counité échoue ou
 le sommet cesse d'être group-like.
 \item \(z_\pm\) ne sont pas des cycles, leurs racines diffèrent ou
 \(\partial h_*=z_--z_+\) échoue.
 \item \(\Psi_\alpha\) n'est pas une application de complexes ou
 \(K_{\beta\alpha}=K_\beta+\Psi_\beta K_\alpha\) échoue ; dans ce cas, on a
 perdu de la retenue ou de la mémoire.
 \item Le complexe réduit n'est pas équilibré ou n'est pas acyclique sur
 \(K_3\) ; l'histoire reste hors du domaine terminal typé.
 \item Les longueurs de Bockstein ne coïncident pas avec les \(a_i\), ou
 \(\Theta_{\rm Bock}\ne\overline{\lambda^{\rm or}}\).
 \item Le raffinement modifie \(d\) ou \(\lambda^{\rm or}\) en dehors de
 blocs unitaires et de changements d'orientation enregistrés.
 \item L'une des treize coordonnées est omise : \emph{objet substitué ;
 conclusion non transférable}.
\end{enumerate}
Ces falsificateurs vérifient l'objet interne construit ; ils ne le remplacent pas
par un contrôle ultérieur.
\end{ablation}
\fi
 \endgroup

\section{Naturalité de l'assemblage et compatibilité sous raffinement}

Avant de former le registre d'horizon, nous explicitons l'action totale des cinq opérations développées ici sur \emph{la même} arête \(\mathfrak e_\alpha\). Si \(P_x,Q_x\) désignent les restrictions de \eqref{eq:pdfv2-cor-sigma} à la fibre de \(x=s(\alpha)\), et de même dans \(y=t(\alpha)\), nous posons d'abord
\begin{equation}
\label{eq:pdfv2-cor-a-eta-emision}
 A_\alpha:=P_yU_\alpha P_x+Q_yU_\alpha Q_x,
 \qquad
 \eta_\alpha:=Q_yU_\alpha P_x+P_yU_\alpha Q_x.
\end{equation}
Les cinq lectures constituent des systèmes dépendants indexés par \(\mathfrak e_\alpha\) elle-même. L'indice fixe source, cible, parcours et registre, mais n'est pas recopié dans la valeur. Leurs types sont les registres dépendants \(\operatorname{LedgerEntry}:=\operatorname{cod}(\lambda)\) et
\begin{align*}
 \mathsf Y_{\rm sp}[\alpha]&:=
 \operatorname{Rec}[U:\mathsf H_x\to\mathsf H_y;
 \sigma_x\in\operatorname{End}\mathsf H_x;
 \sigma_y\in\operatorname{End}\mathsf H_y;A,\eta:\mathsf H_x\to\mathsf H_y],\\
 \mathsf Y_{\rm sol}[\alpha]&:=
 \operatorname{Rec}[b,b^+,b^-\in\operatorname{LedgerEntry};
 \operatorname{upd}:\mathsf{Mem}_x\to\mathsf{Mem}_y],\\
 \mathsf Y_{\rm gau}[\alpha]&:=
 \mathcal P_{\rm fin}\operatorname{Rec}[
 \operatorname{Cell}(\alpha);q\in Q_{\rm inc}(y);B;W_c(q)],\\
 \mathsf Y_{\rm coh}[\alpha]&:=
 \operatorname{Rec}[I_y,J_y,\kappa_{y,N},\delta_y,C_{y,N};
 F_\alpha:C_{y,N}\to C_{x,N}],\\
 \mathsf Y_{\rm det}[\alpha]&:=
 \operatorname{Rec}[\Psi_\alpha,K_\alpha,\mathbf1_\alpha,z_\pm,h_*;
 \operatorname{Det},H_*,\operatorname{Smith},\operatorname{Boc}].
\end{align*}
Par conséquent, le jugement commun est
\[
 \mathfrak e_\alpha:\operatorname{Ext}_k\ \vdash\
 \mathcal O_\bullet(\mathfrak e_\alpha)
 \in\mathsf Y_\bullet[\alpha],
 \qquad
 \bullet\in\{\mathrm{sp,sol,gau,coh,det}\}.
\]
\begin{align}
 \mathcal O_{\rm sp}(\mathfrak e_\alpha)
 &:=
 \bigl(U_\alpha,\sigma_x,\sigma_y,
 A_\alpha,\eta_\alpha\bigr),
 \label{eq:pdfv2-cor-accion-sp}\\
 \mathcal O_{\rm sol}(\mathfrak e_\alpha)
 &:=
 \bigl(b_\alpha,b_\alpha^+,b_\alpha^-;\,
 (M^+,M^-,D,H)\mapsto
 (M^++b_\alpha^+,M^-+b_\alpha^-,
  D+b_\alpha,H+|b_\alpha|)\bigr),
 \label{eq:pdfv2-cor-accion-sol}\\
 \mathcal O_{\rm gau}(\mathfrak e_\alpha)
 &:=
 \left\{
 \bigl(\operatorname{Cell}_{\rm TPK}(\alpha),
 q,B,W_c(q)\bigr):
 q\in Q_{\rm inc}(y)
 \right\},
 \label{eq:pdfv2-cor-accion-gau}\\
 \mathcal O_{\rm coh}(\mathfrak e_\alpha)
 &:=
 \bigl(I_y,J_y,\kappa_{y,N},\delta_y,C_{y,N},
 F_\alpha:C_{y,N}\longrightarrow C_{x,N}\bigr),
 \label{eq:pdfv2-cor-accion-coh}\\
 \mathcal O_{\rm det}(\mathfrak e_\alpha)
 &:=
 \bigl(\Psi_\alpha,K_\alpha,\mathbf1_\alpha,z_\pm(\alpha),
 h_*(\alpha),\operatorname{Det}(C_{y,N}),H_*(C_{y,N}),
 \operatorname{Smith}(C_{y,N}),\operatorname{Boc}(C_{y,N})\bigr).
 \label{eq:pdfv2-cor-accion-det}
\end{align}

La lecture archimédienne interne est la première composante de cette action simultanée ; elle n'introduit encore aucune forme classique :
\begin{equation}
\label{pII-full:eq:continuo-lector-arquimediano}
 \mathcal L_{\rm arq}^{\rm int}(\mathfrak e_\alpha)
 :=\mathcal O_{\rm sp}(\mathfrak e_\alpha).
\end{equation}
Pour toute valeur de ces systèmes, y compris chaque élément de la multisection \(\mathcal O_{\rm gau}\), son support est l'indice dépendant
\begin{equation}
\label{eq:pdfv2-cor-soporte-dependiente}
 \begin{aligned}
 \mathsf B_\alpha&:=
 \{(s(\alpha),t(\alpha),\gamma\alpha,
          \operatorname{Led}(\gamma\alpha))\},\\
 \operatorname{supp}_\alpha:\mathsf Y_\bullet[\alpha]&\longrightarrow
 \mathsf B_\alpha,\qquad
 o\longmapsto(s(\alpha),t(\alpha),\gamma\alpha,
          \operatorname{Led}(\gamma\alpha)).
 \end{aligned}
\end{equation}
Ainsi sont typés, sans projection reconstruisant l'arête,
\begin{equation}
\label{eq:pdfv2-cor-soporte-cinco-lecturas}
 \begin{gathered}
 s_\alpha(\mathcal O_\bullet):=s(\alpha),\qquad
 t_\alpha(\mathcal O_\bullet):=t(\alpha),\\
 \gamma_\alpha(\mathcal O_\bullet):=\gamma\alpha,
 \qquad
 \operatorname{Led}_\alpha(\mathcal O_\bullet)
 :=\operatorname{Led}(\gamma\alpha).
 \end{gathered}
\end{equation}
Ici, \(\operatorname{Cell}_{\rm TPK}(\alpha)\) est l'application de complexes \eqref{eq:pdfv2-cor-cell-accion-emision} induite par le transport de la même fibre de feuillets, et \(F_\alpha\) est \eqref{eq:pdfv2-cor-refinamiento-celda} pour cette extension. En particulier, \(\mathcal O_{\rm gau}\) conserve la multisection complète des \(q\) ; il ne choisit ni l'origine ni une orbite.

Les cinq formules précédentes constituent une seule action dépendante
\begin{equation}
\label{eq:pdfv2-cor-accion-unica-dependiente}
 \mathcal O(\mathfrak e_\alpha)
 :=\left\langle
 \mathcal O_{\rm sp},\mathcal O_{\rm sol},\mathcal O_{\rm gau},
 \mathcal O_{\rm coh},\mathcal O_{\rm det};
 \operatorname{Cpl}_\alpha
 \right\rangle,
\end{equation}
où \(\operatorname{Cpl}_\alpha\) est le certificat formé par les identités
\begin{equation}
\label{eq:pdfv2-cor-acoplamientos-unicos}
 \begin{gathered}
 \operatorname{supp}_\alpha(\mathcal O_{\rm sp})
 =\operatorname{supp}_\alpha(\mathcal O_{\rm sol})
 =\operatorname{supp}_\alpha(\mathcal O_{\rm gau})
 =\operatorname{supp}_\alpha(\mathcal O_{\rm coh})
 =\operatorname{supp}_\alpha(\mathcal O_{\rm det}),\\
 \mathsf Q(t(\alpha))=\sigma_9\mathsf Q(s(\alpha)),\quad
 g(t(\alpha))=g(s(\alpha))+[1]_9,\quad
 q^{(9)}(t(\alpha))=q^{(9)}(s(\alpha))+\chi_9(g(s(\alpha))),\\
 A_\alpha^*A_\alpha+\eta_\alpha^*\eta_\alpha=U_\alpha^*U_\alpha,\qquad
 (M^+,M^-,D,H)'=(M^++b_\alpha^+,M^-+b_\alpha^-,
 D+b_\alpha,H+|b_\alpha|),\\
 Q_{\rm inc}=C^2(\operatorname{Cell}_{\rm TPK};\mathbf F_3),\qquad
 sB=0,\qquad W_c(q+Ba)=W_c(q),\\
 C_{y,N}\text{ est le même complexe dans }
 \mathcal O_{\rm coh}\text{ et }\mathcal O_{\rm det},\qquad
 \partial h_*=z_--z_+,\qquad
 K_{\beta\alpha}=K_\beta+\Psi_\beta K_\alpha.
 \end{gathered}
\end{equation}
\Needspace{18\baselineskip}
Le symbole \(\bullet\) parcourt les cinq lectures de \eqref{eq:pdfv2-cor-accion-unica-dependiente} ; toutes conservent les mêmes source, cible, parcours et registre.

Nous réunissons, sans séparer l'état, les opérations construites jusqu'ici :
\begin{equation}
\label{eq:pdfv2-cor-constructor-unico}
 \begin{aligned}
 \mathfrak D_{k,N}(\widehat x)
 &:=\bigl(\{\mathfrak e_\alpha\}_\alpha;\\
 &\quad U_k,\sigma_k,P_k,Q_k,A_k,\eta_k;\\
 &\quad b_\alpha,b_\alpha^\pm,M_k^\pm,D_k^0,H_k^0,
 \kappa_k^{\rm can};\\
 &\quad V_{\rm TPK},\mathscr I,\mathscr E,\mathscr B,B;\\
 &\quad \operatorname{Cell}_{\rm TPK}(\widehat x),Q_{\rm inc}(\widehat x),W_c;\\
 &\quad I_k,J_k,\kappa_{k,N},\delta,C_{k,N};\\
 &\quad c_\alpha,v_\alpha,\Psi_\alpha,K_\alpha,
 \mathbf1_\alpha,z_\pm(\alpha),h_*(\alpha);\\
 &\quad
 \operatorname{Det}(C_{k,N}),H_*(C_{k,N}),
 \operatorname{Smith}(C_{k,N}),\operatorname{Boc}(C_{k,N});\\
 &\quad \{\mathcal O(\mathfrak e_\alpha)\}_\alpha
 \bigr).
 \end{aligned}
\end{equation}
Le point-virgule ne sépare que des niveaux de lecture d'un même tuple dépendant ; il ne définit pas de structures indépendantes. L'état d'entrée n'est pas enregistré comme première coordonnée ni récupéré par une projection. Chaque opération se démontre par son action sur les générateurs \(\mathfrak e_\alpha\), qui conservent le même support, le même parcours et le même registre.

Pour \(\ell\geq k\) et \(N'\geq N\), le transport est défini composante par composante par troncature des histoires, réduction des coefficients et précomposition des ports. L'affirmation de naturalité à démontrer est unique :
\begin{equation}
\label{eq:pdfv2-cor-naturalidad-unica}
 u_{(\ell,N'),(k,N)}\mathfrak D_{\ell,N'}(\widehat x_\ell)
 =\mathfrak D_{k,N}(\tau_{\ell,k}\widehat x_\ell).
\end{equation}
La naturalité se formule sur l'histoire engendrée, sans prétendre qu'une arête fine isolée doive rester une arête après troncature de ses deux extrémités. Pour \(\gamma=(\varepsilon_1,\ldots,\varepsilon_\ell)\), nous définissons
\[
 \mathsf E_{\leq\ell}(\gamma)
 :=(\mathfrak e_{\varepsilon_i})_{1\leq i\leq\ell},
 \qquad
 \mathcal O_{\leq\ell}(\gamma)
 :=(\mathcal O(\mathfrak e_{\varepsilon_i}))_{1\leq i\leq\ell}.
\]
Les deux transports sont les applications explicites de préfixe
\begin{align}
 u^{\mathfrak e}_{\ell,k}:
 \mathsf E_{\leq\ell}(\gamma)&\longrightarrow
 \mathsf E_{\leq k}(\gamma),&
 (\mathfrak e_{\varepsilon_i})_{i\leq\ell}
 &\longmapsto(\mathfrak e_{\varepsilon_i})_{i\leq k},
 \label{eq:pdfv2-cor-transporte-arista}\\
 u^{\mathcal O}_{\ell,k}:
 \mathcal O_{\leq\ell}(\gamma)&\longrightarrow
 \mathcal O_{\leq k}(\gamma),&
 (\mathcal O(\mathfrak e_{\varepsilon_i}))_{i\leq\ell}
 &\longmapsto
 (\mathcal O(\mathfrak e_{\varepsilon_i}))_{i\leq k}.
 \label{eq:pdfv2-cor-transporte-accion}
\end{align}
Dans chaque entrée conservée, la composante \(\mathrm{sp}\) entrelace \(U,\sigma,P,Q,A,\eta\) ; la composante \(\mathrm{sol}\) applique l'espérance conditionnelle à \(b,b^\pm,\kappa\) ; la composante \(\mathrm{gau}\) applique \(\operatorname{Cell}(\tau)\oplus\tau_Q\) ; la composante \(\mathrm{coh}\) précompose les ports avec \(F\) ; et la composante \(\mathrm{det}\) applique \(\Psi,K,H_*,\operatorname{Smith},\operatorname{Boc}\) à la même application de complexes. Le préfixe agit simultanément sur ces cinq composantes. Ses composantes déjà établies sont
\begin{equation}
\label{eq:pdfv2-cor-naturalidad-componentes}
 \begin{gathered}
 U_{m,k}=U_{m,\ell}U_{\ell,k}\quad(m\geq\ell\geq k),\qquad
 E_kb_f^\pm=b_c^\pm+\kappa_k^{\rm can},\\
 r_{N',N}\kappa_{k,N'}=
 \kappa_{k,N}(r_{N',N}\oplus r_{N',N}),\\
 \Psi_\beta\Psi_\alpha=\Psi_{\beta\alpha},
 \qquad K_{\beta\alpha}=K_\beta+\Psi_\beta K_\alpha,\\
 u^{\mathcal O}_{\ell,k}\mathcal O_{\leq\ell}(\gamma)
 =\mathcal O_{\leq k}(\gamma)
 =\operatorname{map}(\mathcal O)\bigl(u^{\mathfrak e}_{\ell,k}
                    \mathsf E_{\leq\ell}(\gamma)\bigr).
 \end{gathered}
\end{equation}
La première égalité utilise la multiplication des poids conditionnels sur le même système de cylindres communs ; les autres sont les formules explicites précédentes. Aucune projection visible ne peut servir à reconstruire ou à sélectionner \(\widehat x\).

Pour compléter la dernière composante, notons \(F_{(\ell,N'),(k,N)}:C_{\ell,N'}\to C_{k,N}\) l'application de complexes de \eqref{eq:pdfv2-cor-refinamiento-celda}. Son cône est conservé comme partie de la transition. La fonctorialité de l'homologie et de la suite de Bockstein donne
\begin{equation}
\label{eq:pdfv2-cor-naturalidad-homologia-bockstein}
 H_*(F_{m,k})=H_*(F_{\ell,k})H_*(F_{m,\ell}),
 \qquad
 \operatorname{Boc}(F_{m,k})
 =\operatorname{Boc}(F_{\ell,k})
  \operatorname{Boc}(F_{m,\ell}).
\end{equation}
La droite déterminant est transportée au moyen de l'identité du triangle exact
\begin{equation}
\label{eq:pdfv2-cor-naturalidad-determinante}
 \operatorname{Det}(C_{k,N})
 \simeq
 \operatorname{Det}(C_{\ell,N'})\otimes
 \operatorname{Det}\!\left(\operatorname{Cone}
 F_{(\ell,N'),(k,N)}\right).
\end{equation}
Cette formule enregistre également les raffinements qui ajoutent de l'homologie. La lecture de Smith est conservée sans feindre l'invariance : à chaque transition est associé le certificat total
\begin{equation}
\label{eq:pdfv2-cor-transicion-smith}
 \mathcal S(F):=
 \bigl(\operatorname{Smith}(C_{\ell,N'}),
       \operatorname{Smith}(C_{k,N}),
       \operatorname{Smith}(\operatorname{Cone}F)\bigr),
\end{equation}
qui détermine littéralement quels facteurs se conservent, apparaissent ou disparaissent.

L'action d'une symétrie cohérente \(g\) sur l'arbre commun induit des isométries de permutation \(\widetilde g_k\) dans \(\mathscr H_k\). La bijection des successeurs et la conservation des poids de \eqref{eq:continuo-comun-naturalidad-medida} prouvent sur les vecteurs de base
\begin{equation}
\label{eq:pdfv2-cor-naturalidad-transporte-hojas}
 \widetilde g_{k+1}U_k=U_k\widetilde g_k,
 \qquad
 \widetilde g_k\sigma_k=\sigma_k\widetilde g_k,
 \qquad
 \widetilde g_kP_k=P_k\widetilde g_k,
 \qquad
 \widetilde g_kQ_k=Q_k\widetilde g_k.
\end{equation}

\begin{theorem}[Totalité et naturalité du constructeur corrélatif]
\label{thm:pdfv2-cor-totalidad-naturalidad}
Pour tout état total de \eqref{eq:continuo-comun-estado}, tout horizon fini et tout niveau de coefficients, l'expression \eqref{eq:pdfv2-cor-constructor-unico} est définie. Ses cinq caractéristiques sont des opérations simultanées sur le même registre d'extension et satisfont l'unique identité \eqref{eq:pdfv2-cor-naturalidad-unica}. Le collier, le complexe et la droite déterminant restent non vides même lorsque leurs lecteurs ultérieurs ne sont ni acycliques ni inversibles.
\end{theorem}

\begin{proof}
La non-vacuité de l'arbre et de ses ports est \cref{prop:continuo-comun-arbol-finito-no-vacio} ; la mesure et l'espérance sont totales par \cref{thm:continuo-comun-medida-canonica}. Les formules \eqref{eq:pdfv2-cor-u-comun}--\eqref{eq:pdfv2-cor-memorias} définissent transport, polarité, bilan et mémoire pour tous les générateurs. Les quatre directions de \eqref{eq:pdfv2-cor-cuatro-rectas-tpk} sont distinctes par \eqref{eq:pdfv2-cor-proyectores-no-colapso} ; la complétion cellulaire \cref{def:pdfv2-cor-completacion-celular-tpk} contient leurs deux ordres de chemin et six cellules, tandis que \eqref{eq:pdfv2-cor-collar-naturalidad} prouve sa naturalité. La suite cellulaire est exacte par \eqref{eq:pdfv2-cor-sucesion-celda}, le complexe local satisfait \(\partial^2=0\) et le remplissage corrélatif est explicite par \eqref{eq:pdfv2-cor-filler}. Toute construction finie munie de bases possède la droite \eqref{eq:pdfv2-cor-linea-determinante-total} ; homologie, Bockstein et certificat de Smith sont transportés par \eqref{eq:pdfv2-cor-naturalidad-homologia-bockstein}--\eqref{eq:pdfv2-cor-transicion-smith}. Enfin, les identités se vérifient sur chaque générateur \(\mathfrak e_\alpha\) au moyen de \eqref{eq:pdfv2-cor-naturalidad-componentes}, \eqref{eq:pdfv2-cor-collar-naturalidad} et \eqref{eq:pdfv2-cor-naturalidad-transporte-hojas} ; par linéarité et composition fonctorielle des parcours, on obtient \eqref{eq:pdfv2-cor-naturalidad-unica} pour le registre entier.
\end{proof}

\section{Distinction causale entre identités, déductions et réalisations du constructeur corrélatif}

\paragraph{Non-vacuité et identités exactes du constructeur corrélatif}
La non-vacuité importée de \eqref{eq:continuo-comun-out-finito} et les identités \eqref{eq:pdfv2-cor-b-identidades}, \eqref{eq:pdfv2-cor-sucesion-celda}, \eqref{eq:pdfv2-cor-sb-cero}, \eqref{eq:pdfv2-cor-u-gram}, \eqref{eq:pdfv2-cor-balance-hojas}, \eqref{eq:pdfv2-cor-jensen-hojas}, \eqref{eq:pdfv2-cor-dos-cero}, \eqref{eq:pdfv2-cor-cociclo-matricial} et \eqref{eq:pdfv2-cor-filler} se déduisent des définitions et des relations APP--TRIT--TPK exposées. Leur statut est \texttt{EXACT\_INTERNE}.

\paragraph{Hypothèses initiales pour l'isométrie de \(U_k\), la suite cellulaire et la composition \(\Psi/K\)}
La construction de \(U_k\) utilise comme structure de départ la finitude des
fibres, la survie à tout horizon, la troncature fonctionnelle et les
transports des deux feuillets déjà produits par TPK\@. L’égalité
\(U_k^*U_k=I\) exige en outre de vérifier \(U_\alpha^*U_\alpha=I\) pour chaque
émission ; sans cette vérification, l’identité de Gram
\eqref{eq:pdfv2-cor-u-gram} demeure, et elle est bien inconditionnelle. La suite cellulaire utilise
les deux ports APP ordonnés. La composition \(\Psi/K\) utilise le carry entier,
l’orientation et la frontière. Sous ces prémisses énumérées, les conclusions
indiquées sont affirmatives.

\paragraph{Données supplémentaires pour les réalisations isométriques, les rétractes équilibrés et la trivialisation de la droite déterminant}
Le constructeur interne est complet avec le Gram positif de \eqref{eq:pdfv2-cor-u-gram}, l'espace total \(Q_{\rm inc}\), la suite cellulaire et la droite déterminant munie de son cône de transition. Une expression spécialisée peut exiger des données supplémentaires : l'isométrie de chaque \(U_\alpha\), une identification numérique concrète du défaut APP, un rétracte équilibré et acyclique ou la trivialisation d'une droite déterminant. Ces données sont des entrelaceurs du codomaine exprimé ; elles ne font partie ni du domaine ni de l'existence de \(\mathfrak D_{k,N}\). Leur absence empêche exclusivement la promotion qui les mentionne.

\begin{criterion}[Falsificateurs de la corrélation intrinsèque]
\label{crit:pdfv2-cor-falsadores}
La construction conjointe échoue au premier niveau où survient l'un des faits suivants :
\begin{enumerate}
 \item une opération utilise une extension, une orientation, une retenue, une mémoire, une frontière ou un parcours différent de celui enregistré dans \(\mathfrak e_\alpha\) ;
 \item un successeur de \(\operatorname{Ch}_k(\widehat x_k)\) est choisi et les autres sont supprimés ;
 \item le poids de \eqref{eq:pdfv2-cor-peso-conteo} ne somme pas à un, on n’a pas
 \eqref{eq:pdfv2-cor-u-gram}, ou l’on déclare \(U_k^*U_k=I\) sans prouver
 \(U_\alpha^*U_\alpha=I\) pour toute émission ;
 \item \(\kappa_k^{\rm can}\) est identifié à la retenue entière \(\kappa_\alpha\) ;
 \item \(sB=0\) échoue, \(Q_{\rm inc}(\widehat x)\) est contracté à son origine, l'action \(q\mapsto q+Ba\) est modifiée, ou les deux ordres de chemin sont identifiés avant l'application du lecteur ;
 \item un raffinement rompt \eqref{eq:pdfv2-cor-refinamiento-cadena} ou le cocycle \eqref{eq:pdfv2-cor-cociclo-matricial} ;
 \item \(\partial^2=0\) échoue, \(z_\pm\) ne sont pas des cycles ou \(\partial h_*\neq z_--z_+\) ;
 \item un scalaire déterminantal non nul est affirmé sans le rétracte, le bilan et l'acyclicité exigés par cette expression ;
 \item l'une quelconque des cinq opérations est présentée comme un objet séparé du constructeur \(\mathfrak D_{k,N}\).
\end{enumerate}
Dans tous ces cas, le diagnostic est : \emph{objet commun remplacé ; conclusion non transférable}.
\end{criterion}
 \end{HMTScientificOwnerSubordinated}
 \label{chap:83-05}

\begin{HMTScientificOwnerSubordinated}
\chapter[Terminal, naturalité et limite]{Terminal multivalué, naturalité unique et passage à la limite}
\label{chap:pdfv2-terminal-naturalidad-limite}

\begin{thesisbox}
Le continu discret produit un seul système d'histoires, d'extensions, de mémoires et de terminaux. Les discriminants internes \(\mathrm{sp},\mathrm{sol},\mathrm{gau},\mathrm{coh},\mathrm{det}\) désignent cinq caractéristiques simultanées de ce même système : ils ne fixent pas d'entrées distinctes et n'autorisent pas à choisir des histoires différentes. Dans ce chapitre, le terminal conserve toutes les continuations, la connexion nonadique reste une correspondance et la naturalité s'exprime par une seule identité.
\end{thesisbox}

\section{Histoires survivantes avec témoins d'extension}

Soit \(\widehat x_k\in\widehat{\mathcal X}^{\rm tot}_k\) l'état total de \eqref{eq:continuo-comun-dominio-total}. Pour \(N\geq k\), nous définissons sans autre prédicat d'appartenance
\begin{equation}
\label{eq:pdfv2-terminal-historias-finita}
 \mathscr H_{k,N}(\widehat x_k)
 :=\operatorname{Term}_{k,N}(\widehat x_k).
\end{equation}
Chaque \(h=(\varepsilon_{k+1},\ldots,\varepsilon_N)\) du membre de droite contient, par \eqref{eq:continuo-comun-vertice-arbol}, ses états intermédiaires, les témoins de \(\operatorname{Ext}\), les retenues, les frontières, les parcours et le registre. Chaque arête satisfait l'action de phase \eqref{eq:continuo-comun-accion-fase-tpk} ; tout sous-chemin de neuf pas détermine donc, avec ses neuf émissions et non seulement sa longueur, un élément de \(\mathsf Z^\Gamma_{9,j}\) et un témoin de \(\Gamma_{9,j}\) par \eqref{eq:continuo-comun-gamma9-apice}--\eqref{eq:continuo-comun-gamma9-relacion}. Ainsi, Ext et l'holonomie nonadique naissent du même mot compatible avec la phase ; on n'ajoute à \eqref{eq:pdfv2-terminal-historias-finita} ni filtre de survie ni seconde condition \(\Gamma_9\).

Pour \(k\leq\ell\leq N\), soit
\begin{equation}
\label{eq:pdfv2-terminal-prefijos-intermedios}
 \mathscr P_{\ell\mid k}(\widehat x_k)
 :=\{h_{[k,\ell]}:h\in\mathscr H_{k,N}(\widehat x_k)
       \text{ pour un certain }N\geq\ell\}.
\end{equation}
Pour \(p\in\mathscr P_{\ell\mid k}(\widehat x_k)\), nous utilisons l'abréviation
\begin{equation}
\label{eq:pdfv2-terminal-continuaciones-prefijo}
 \mathscr H_{\ell,N}(p)
 :=\{q:p\star q\in\mathscr H_{k,N}(\widehat x_k)\}.
\end{equation}
Le prolongement neutre de \eqref{eq:continuo-comun-seccion-total} étend chaque préfixe à tout horizon. En conséquence,
\begin{equation}
\label{eq:pdfv2-terminal-no-vacuidad-historias}
 \mathscr H_{k,N}(\widehat x_k)\neq\varnothing,
 \qquad
 \mathscr P_{\ell\mid k}(\widehat x_k)\neq\varnothing,
\end{equation}
par \cref{prop:continuo-comun-arbol-finito-no-vacio}.

La troncature
\begin{equation}
\label{eq:pdfv2-terminal-truncamiento-historias}
 \pi_{N+1,N}^{\mathscr H}:
 \mathscr H_{k,N+1}(\widehat x_k)
 \longrightarrow\mathscr H_{k,N}(\widehat x_k)
\end{equation}
est exactement \eqref{eq:continuo-comun-proyeccion-terminal} et est surjectif par \cref{lem:continuo-comun-terminal-sobreyectivo}. La survie est, par \eqref{eq:continuo-comun-surv-igual-total}, une sortie démontrée du domaine total, non une sélection préalable.

\section{Registre terminal engendré et ses \(13\) coordonnées}

Fixons une histoire
\[
 h=(\varepsilon_{j+1}:x_j\longrightarrow x_{j+1})_{j=k}^{N-1},
 \qquad
 \gamma_{k,k}:=I_{x_k},\qquad
 \gamma_{k,j+1}:=\gamma_{k,j}\star\varepsilon_{j+1},
\]
La notation opératorielle \(U_{\gamma_{k,j+1}}
=U_{\varepsilon_{j+1}}U_{\gamma_{k,j}}\) conserve l'ordre d'application, tandis que \(\star\) conserve l'ordre chronologique du mot. Nous fixons en outre un niveau de coefficients \(r\geq1\). Pour abréger, nous écrivons
\[
 C_{x_j,r}:=C_{j,r}(x_j),\qquad
 F_{\varepsilon_{j+1}}:C_{x_{j+1},r}\longrightarrow C_{x_j,r}
\]
pour le complexe et l'application de complexes de \eqref{eq:pdfv2-cor-cono-celda} et de \eqref{eq:pdfv2-cor-refinamiento-celda}. Le type du registre de niveau \(j\), noté \(\mathfrak R_{j,r}\), est un registre dépendant possédant les treize champs suivants. Pour chaque champ, nous indiquons le type, le germe et l'action d'une arête \(\varepsilon:=\varepsilon_{j+1}\).

L'état \(x_k\) contient déjà le mot complet \(\gamma_{0,k}=\varepsilon_1\star\cdots\star\varepsilon_k\). Ainsi, un ``germe en \(k\)'' ne réinitialise jamais la généalogie : il est le pliage des mêmes actualisations depuis la racine jusqu'à \(k\). Nous notons
\begin{equation}
\label{eq:pdfv2-terminal-prefijo-heredado}
 \begin{gathered}
 M_k^\pm:=\sum_{i=1}^{k}b_{\varepsilon_i}^\pm,qquad
 D_k:=\sum_{i=1}^{k}b_{\varepsilon_i},qquad
 H_k:=\sum_{i=1}^{k}|b_{\varepsilon_i}|,\\
 \mathbf c_{\leq k}:=
  \bigl(c_r(\varepsilon_i,\gamma_{0,i-1})\bigr)_{1\leq i\leq k},
 \qquad
 \mathbf k_{\leq k}:=
  \bigl(\kappa_{i-1}^{\rm can}\bigr)_{1\leq i\leq k},\\
 \mathbf F_{\leq k}:=
  \bigl(F_{\varepsilon_i}:C_{x_i,r}\to C_{x_{i-1},r}\bigr)_{1\leq i\leq k}.
 \end{gathered}
\end{equation}
Toutes les listes vides ci-dessous ne sont à entendre littéralement qu'en \(k=0\) ; pour \(k>0\), elles sont remplacées par les préfixes hérités de \eqref{eq:pdfv2-terminal-prefijo-heredado} et du registre de \(\gamma_{0,k}\). Ainsi, le registre ne conserve pas de copie opaque de \(x_k\), mais ne perd pas non plus son passé APP--TRIT--TPK.

\paragraph{Fibre \(\mathsf F\) : domaine local et conservation sous prolongement}
La fibre élémentaire est
\begin{equation}
\label{eq:pdfv2-terminal-tipo-fibra}
 \operatorname{Fib}_r(x):=
 \bigl(\mathsf H_x^\Sigma\oplus\mathsf H_x^\Pi,\,
       \mathsf M_x,\,
       \operatorname{Cell}_{\rm TPK}(x),\,C_{x,r}\bigr).
\end{equation}
Le champ \(\mathsf F_{k,j}\) est un mot dans la catégorie de ces fibres. Son germe et son actualisation sont
\begin{align}
 \mathsf F_{k,k}
 &:=\bigl(\operatorname{Fib}_r(x_0),I\bigr)\star
 \mathop{\star}_{i=1}^{k}
 \bigl(\operatorname{Fib}_r(x_i),U_{\varepsilon_i},
 \mathsf C_{\varepsilon_i},\operatorname{Cell}_{\rm TPK}(\varepsilon_i),
 F_{\varepsilon_i}\bigr),
 \label{eq:pdfv2-terminal-semilla-f}\\
 \operatorname{Upd}^{\mathsf F}_{\varepsilon}\mathsf F_{k,j}
 &:=
 \mathsf F_{k,j}\star
 \bigl(\operatorname{Fib}_r(x_{j+1}),U_\varepsilon,
 \mathsf C_\varepsilon,\operatorname{Cell}_{\rm TPK}(\varepsilon),
 F_\varepsilon\bigr).
 \label{eq:pdfv2-terminal-update-f}
\end{align}
Elle est non vide puisque les deux feuillets, la mémoire, l'unité cellulaire \eqref{eq:pdfv2-cor-cell-unidad} et le complexe des ports sont déjà construits. La composition de ses cinq applications prouve la naturalité.

\paragraph{Extensions \(\mathsf E\) : morphismes de prolongement compatible de l'état}
Le codomaine est la catégorie des mots de relations typées. Nous posons
\begin{align}
 \mathsf E_{k,k}
 &:=(I_{x_0},\operatorname{Ext}_{I_{x_0}})\star
 \mathop{\star}_{i=1}^{k}
 \bigl(\varepsilon_i,\operatorname{Ext}_{\varepsilon_i},
 \operatorname{Ext}_{\varepsilon_i}\circ
 \operatorname{Ext}_{\gamma_{0,i-1}}
 \subseteq\operatorname{Ext}_{\gamma_{0,i}}\bigr),
 \label{eq:pdfv2-terminal-semilla-e}\\
 \operatorname{Upd}^{\mathsf E}_{\varepsilon}\mathsf E_{k,j}
 &:=
 \mathsf E_{k,j}\star
 \bigl(\varepsilon,\operatorname{Ext}_{\varepsilon},
 \operatorname{Ext}_{\varepsilon}\circ
 \operatorname{Ext}_{\gamma_{0,j}}
 \subseteq\operatorname{Ext}_{\gamma_{0,j+1}}\bigr).
 \label{eq:pdfv2-terminal-update-e}
\end{align}
L'identité et l'émission neutre prouvent la non-vacuité ; l'associativité de la composition relationnelle prouve la naturalité.

\paragraph{Relations \(\mathsf R\) : incidence, orientation et compatibilité de composition}
Soit \(\operatorname{RelBlock}_{j,r}\) le type des blocs d'équations étiquetés par l'arête et le préfixe sur lesquels ils sont vérifiés, et soit \(\operatorname{Pres}_{j,r}:=\operatorname{List}
(\operatorname{RelBlock}_{j,r})\). Le bloc unité et le germe sont
\begin{equation}
\label{eq:pdfv2-terminal-semilla-r}
 \mathsf B_I:=(I_{x_0};U_I=I,\ \mathsf U_I=I,\ \kappa(I)=0,\ \partial I=0),
 \qquad
 \mathsf R_{k,k}:=(\mathsf B_I)\star
 \mathop{\star}_{i=1}^{k}
 \mathsf B_R(\varepsilon_i;\gamma_{0,i-1}).
\end{equation}
Pour une arête \(\varepsilon:x_j\to x_{j+1}\), le nouveau bloc est le tuple étiqueté
\begin{equation}
\label{eq:pdfv2-terminal-bloque-r}
 \begin{split}
 \mathsf B_R(\varepsilon;\gamma_{0,j}):=\bigl(
 &\varepsilon,\gamma_{0,j};\ 
 \delta^\diamond(\varepsilon)
 =\operatorname{res}^\diamond(\varepsilon)
  +9\operatorname{quo}^\diamond(\varepsilon),\\
 &\Delta(\varepsilon)=\operatorname{or}(\varepsilon)
  +3\kappa(\varepsilon),\quad
 U_{\gamma_{0,j+1}}=U_\varepsilon U_{\gamma_{0,j}},\\
 &\mathsf U_{\gamma_{0,j+1}}
 =\mathsf U_\varepsilon\mathsf U_{\gamma_{0,j}},\quad
 \kappa(\gamma_{0,j+1})
 =\kappa(\varepsilon)+\mathsf C_\varepsilon\kappa(\gamma_{0,j}),\\
 &\partial\gamma_{0,j+1}
 =\partial\varepsilon+\partial\gamma_{0,j},\quad
 \operatorname{Cpl}_\varepsilon
 \text{ de }\eqref{eq:pdfv2-cor-acoplamientos-unicos}
 \bigr).
 \end{split}
\end{equation}
L'actualisation concatène le bloc sans perdre ni date ni multiplicité :
\begin{equation}
\label{eq:pdfv2-terminal-update-r}
 \operatorname{Upd}^{\mathsf R}_{\varepsilon}\mathsf R_{k,j}
 :=\mathsf R_{k,j}\star\mathsf B_R(\varepsilon;\gamma_{0,j}).
\end{equation}
La troncature \(\operatorname{tr}^{\mathsf R}_{j+1,j}\) supprime le dernier bloc étiqueté. En outre, \(a_*\mathsf B_R(\varepsilon;\gamma)
=\mathsf B_R(a\varepsilon;a\gamma)\), car chaque équation est naturelle sur la même arête et le même préfixe.

\paragraph{Nombres \(\mathsf N\) : expression de la valeur avec conservation de la généalogie}
Pour une arête \(\varepsilon:x_j\to x_{j+1}\), nous définissons la donnée numérique
\begin{equation}
\label{eq:pdfv2-terminal-numero-arista}
 \nu(\varepsilon):=
 \bigl(j+1,\mathsf Q_{j+1},g_{j+1},q^{(9)}_{j+1},\beta_{j+1},
 \operatorname{res}^{\Sigma},\operatorname{quo}^{\Sigma},
 \operatorname{res}^{\Pi},\operatorname{quo}^{\Pi},
 \operatorname{or},\kappa,\Delta\bigr)(\varepsilon).
\end{equation}
Son type est
\[
 \mathbb N\times\widetilde{\mathcal O}_9\times C_9\times\mathbb N
 \times\mathbb Z/12\mathbb Z
 \times\{0,\ldots,8\}^2\times\mathbb Z^5.
\]
Le germe contient la liste numérique complète du préfixe \(0\to k\) ; l'actualisation ajoute exactement la donnée de l'arête suivante :
\begin{equation}
\label{eq:pdfv2-terminal-update-n}
 \mathsf N_{k,k}:=
 ((0,\mathsf Q_0,g_0,q^{(9)}_0,\beta_0;0,\ldots,0))
 \star\mathop{\star}_{i=1}^{k}\nu(\varepsilon_i),\qquad
 \operatorname{Upd}^{\mathsf N}_{\varepsilon}\mathsf N_{k,j}
 :=\mathsf N_{k,j}\star\nu(\varepsilon).
\end{equation}
La totalité provient des divisions APP, de la division TRIT et de l'odomètre TPK.

\paragraph{Orientation \(\mathsf O\) : transport du signe et sélection du feuillet}
Le type d'un état d'orientation est
\[
 \mathbb Z\times
 (\mathsf H^\Sigma_x\oplus\mathsf H^\Pi_x)
 \times\mathbb Z[\mathsf A],
\]
avec pour germe et actualisation
\begin{align}
 \mathsf O_{k,k}&:=
 \left(\sum_{i=1}^{k}\operatorname{or}(\varepsilon_i),
       \operatorname{Sh}(x_k),\partial\gamma_{0,k}\right),
 \label{eq:pdfv2-terminal-semilla-o}\\
 \operatorname{Upd}^{\mathsf O}_{\varepsilon}
 (\mathfrak o,\operatorname{Sh},\mathfrak b)
 &:=
 \bigl(\mathfrak o+\operatorname{or}(\varepsilon),
 U_\varepsilon\operatorname{Sh},
 \mathfrak b+\partial\varepsilon\bigr).
 \label{eq:pdfv2-terminal-update-o}
\end{align}
Les feuillets APP prouvent la non-vacuité. La fonctorialité de \(U\), l'additivité de \(\operatorname{or}\) et la télescopie de la frontière prouvent la naturalité.

\paragraph{Retenue \(\mathsf K\) : actualisation du quotient et mémoire résiduelle}
Nous utilisons les relèvements entiers typés dans \eqref{eq:pdfv2-cor-carry-matricial-def}. Le champ est de type
\[
 \mathsf M_{x_j}\times
 \operatorname{Hom}(\mathsf M_{x_0},\mathsf M_{x_j})
 \times\operatorname{Mat}_5(\mathbb Z)
 \times\operatorname{List}(\operatorname{Mat}_5(\mathbb Z))
 \times\operatorname{List}(\mathbb Q_{\geq0}).
\]
Pour \(\gamma_{k,j}=\varepsilon_j\cdots\varepsilon_{k+1}\), nous fixons explicitement
\begin{equation}
\label{eq:pdfv2-terminal-lift-entero-ruta}
 \widetilde L_{\gamma_{k,k}}:=I_5,\qquad
 \widetilde L_{\gamma_{k,j}}
 :=\widetilde L_r(\varepsilon_j)\cdots
   \widetilde L_r(\varepsilon_{k+1}).
\end{equation}
Pour le préfixe hérité, nous fixons également
\[
 \widetilde L_{\gamma_{0,k}}
 :=\widetilde L_r(\varepsilon_k)\cdots
   \widetilde L_r(\varepsilon_1).
\]
Le germe et l'actualisation sont
\begin{align}
 \mathsf K_{k,k}&:=
 \bigl(\kappa(\gamma_{0,k}),\mathsf C_{\gamma_{0,k}},
       \widetilde L_{\gamma_{0,k}},
       \mathbf c_{\leq k},\mathbf k_{\leq k}\bigr),
 \label{eq:pdfv2-terminal-semilla-k}\\
 \operatorname{Upd}^{\mathsf K}_{\varepsilon}
 (\kappa_\gamma,\mathsf C_\gamma,\widetilde L_\gamma,
   \mathbf c,\mathbf k)
 &:=
 \bigl(\kappa(\varepsilon)+\mathsf C_\varepsilon\kappa_\gamma,\,
 \mathsf C_\varepsilon\mathsf C_\gamma,\,
 \widetilde L_r(\varepsilon)\widetilde L_\gamma,\,
 \mathbf c\star c_r(\varepsilon,\gamma),\,
 \mathbf k\star\kappa_j^{\rm can}\bigr).
 \label{eq:pdfv2-terminal-update-k}
\end{align}
L'identité fournit le germe ; les deux cocycles \eqref{eq:continuo-comun-cociclo-carry} et \eqref{eq:pdfv2-cor-cociclo-matricial} prouvent l'indépendance du parenthésage et la naturalité.

\paragraph{Mémoire \(\mathsf M\) : registre du parcours et compatibilité sous troncature}
Son type conserve simultanément mémoire affine, registre et mémoires signées :
\begin{equation}
\label{eq:pdfv2-terminal-tipo-ledger-entry}
 \operatorname{LedgerEntry}:=
 \{\lambda(\varepsilon):\varepsilon\text{ émission élémentaire typée}\}.
\end{equation}
\[
 \mathsf M_x\times\operatorname{List}(\operatorname{LedgerEntry})
 \times\mathbb N^2\times\mathbb Z\times\mathbb N.
\]
Nous posons
\begin{align}
 \mathsf M_{k,k}&:=
 \bigl(m_k,\operatorname{Led}(\gamma_{0,k}),
       M_k^+,M_k^-,D_k,H_k\bigr),
 \label{eq:pdfv2-terminal-semilla-m}\\
 \operatorname{Upd}^{\mathsf M}_{\varepsilon}
 (m,\mathbf{Led},M^+,M^-,D,H)
 &:=
 \bigl(\mathsf C_\varepsilon m+\kappa(\varepsilon),\,
 \mathbf{Led}\star\lambda(\varepsilon),\,
 M^++b_\varepsilon^+,\,M^-+b_\varepsilon^-,\,
 D+b_\varepsilon,\,H+|b_\varepsilon|\bigr).
 \label{eq:pdfv2-terminal-update-m}
\end{align}
L'action affine et \eqref{eq:pdfv2-cor-memorias} prouvent la totalité et la naturalité ; le passé n'est pas recalculé à partir de la dernière valeur.

\paragraph{Frontière \(\mathsf B\) : données de raccordement et survie de l'extension}
Le type est une chaîne de frontière accompagnée d'un mot de complexes munis de bases et d'applications de chaînes. Son germe et son actualisation sont
\[
 \mathsf B_{k,j}\in
 \operatorname{Rec}[\mathfrak b\in\mathbb Z[\mathsf A_j];
 C_{x_j,r}\in\operatorname{Ch}_{\rm bas};
 \mathbf F\in\operatorname{Path}_{\operatorname{Ch}_{\rm bas}}
 (C_{x_j,r},C_{x_0,r})].
\]
\begin{align}
 \mathsf B_{k,k}&:=
 \bigl(\partial\gamma_{0,k},C_{x_k,r},\mathbf F_{\leq k}\bigr),
 \label{eq:pdfv2-terminal-semilla-b}\\
 \operatorname{Upd}^{\mathsf B}_{\varepsilon}
 (\mathfrak b,C_{x_j,r},\mathbf F)
 &:=
 \bigl(\mathfrak b+\partial\varepsilon,\,
 C_{x_{j+1},r},\,\mathbf F\star F_\varepsilon\bigr).
 \label{eq:pdfv2-terminal-update-b}
\end{align}
La non-vacuité provient des ports communs ; \(\partial^2=0\), \(F_{\beta\varepsilon}=F_\varepsilon F_\beta\) et \eqref{eq:continuo-comun-frontera-telescopica} prouvent la compatibilité.

\paragraph{Parcours \(\mathsf G\) : histoire observable de l'extension survivante}
Le champ est la liste des préfixes dans la catégorie des chemins :
\begin{equation}
\label{eq:pdfv2-terminal-update-g}
 \mathsf G_{k,k}:=(\gamma_{0,i})_{0\leq i\leq k},\qquad
 \operatorname{Upd}^{\mathsf G}_{\varepsilon}\mathsf G_{k,j}
 :=\mathsf G_{k,j}\star\gamma_{0,j+1}.
\end{equation}
L'identité, la concaténation associative et \(a(\varepsilon\gamma)=a(\varepsilon)a(\gamma)\) prouvent les trois propriétés requises.

\paragraph{Multisection \(\mathsf X\) : fibre finie stable et obstruction à la sélection ponctuelle}
Pour \(M\geq j\), nous définissons par action
\begin{equation}
\label{eq:pdfv2-terminal-multiseccion-interna}
 \mathsf X_{k,j}(M):=
 \{\gamma_{k,j}\star q:
 q\in\operatorname{Term}_{j,M}(x_j)\}.
\end{equation}
Ainsi
\[
 \mathsf X_{k,k}(M)=\operatorname{Term}_{k,M}(x_k),
\]
et
\begin{equation}
\label{eq:pdfv2-terminal-update-x}
 \operatorname{Upd}^{\mathsf X}_{\varepsilon}\mathsf X_{k,j}(M)
 :=
 \{\gamma_{k,j}\star\varepsilon\star q:
 q\in\operatorname{Term}_{j+1,M}(x_{j+1})\}.
\end{equation}
Le prolongement neutre prouve la non-vacuité. L'effacement du dernier pas et l'action d'une symétrie commutent littéralement avec la concaténation.

\paragraph{Survie \(\mathsf S\) : persistance sous raffinement et passage à la limite}
Ici, on n'enregistre pas une proposition, mais des témoins. Soit
\begin{equation}
\label{eq:pdfv2-terminal-testigo-neutro}
 s_{j,j}:=I_{x_j},\qquad
 s_{j,M+1}:=s_{j,M}\star
 \varepsilon_M^0(t(s_{j,M})).
\end{equation}
Le germe est \(\mathsf S_{k,k}:=(s_{k,M})_{M\geq k}\), et
\begin{equation}
\label{eq:pdfv2-terminal-update-s}
 \operatorname{Upd}^{\mathsf S}_{\varepsilon}\mathsf S_{k,j}
 :=(\gamma_{k,j}\star\varepsilon\star s_{j+1,M})_{M\geq j+1}.
\end{equation}
Par construction, \(\pi_{M+1,M}s_{j,M+1}=s_{j,M}\) ; la naturalité de l'émission neutre prouve l'équivariance.

\paragraph{Terminal \(\mathsf T\) : universalité de l'objet limite et factorisation des lecteurs}
Le terminal est le système inverse, non une queue choisie :
\begin{equation}
\label{eq:pdfv2-terminal-sistema-inverso-campo}
 \mathsf T(x_j):=
 \bigl((\operatorname{Term}_{j,M}(x_j))_{M\geq j},
       (\pi_{M+1,M})_{M\geq j}\bigr).
\end{equation}
Pour chaque \(M\geq j\), l'application de préfixe typée est
\begin{equation}
\label{eq:pdfv2-terminal-update-t}
 \operatorname{pref}_{\gamma_{k,j},M}:
 \operatorname{Term}_{j,M}(x_j)\longrightarrow
 \operatorname{Term}_{k,M}(x_k),\qquad
 q\longmapsto\gamma_{k,j}\star q.
\end{equation}
Nous définissons l'état du champ et son actualisation par
\begin{align}
 \mathsf T_{k,j}
 &:=
 \left(\mathsf T(x_j),
   (\operatorname{pref}_{\gamma_{k,j},M})_{M\geq j}\right),
 \label{eq:pdfv2-terminal-estado-t}\\
 \operatorname{Upd}^{\mathsf T}_{\varepsilon}\mathsf T_{k,j}
 &:=
 \left(\mathsf T(x_{j+1}),
   (\operatorname{pref}_{\gamma_{k,j+1},M})_{M\geq j+1}\right),
 \qquad \gamma_{k,j+1}:=\gamma_{k,j}\star\varepsilon.
 \label{eq:pdfv2-terminal-update-t-explicita}
\end{align}
En particulier, \(\mathsf T_{k,k}\) utilise \(\gamma_{k,k}=I_{x_k}\). Les témoins \eqref{eq:pdfv2-terminal-testigo-neutro} et la surjectivité de \(\pi\) prouvent que le système est non vide ; niveau par niveau,
\begin{equation}
\label{eq:pdfv2-terminal-prefijo-natural}
 \pi^{\mathscr H}_{M+1,M}
 \operatorname{pref}_{\gamma_{k,j},M+1}
 =\operatorname{pref}_{\gamma_{k,j},M}
  \pi^{\mathscr H}_{M+1,M},
\end{equation}
car les deux parcours suppriment seulement la dernière émission de \(q\).

\paragraph{13. Lecteur terminal \(\mathsf L\) : évaluation de l'extension compatible et conservation de la généalogie}
Le lecteur ne contient aucun objet classique. Son type est la catégorie des chemins de l'unique constructeur dépendant \(\mathfrak D_{j,r}\) de \eqref{eq:pdfv2-cor-constructor-unico}. Nous posons
\begin{equation}
\label{eq:pdfv2-terminal-transicion-lector-tipado}
 u^{\mathfrak D}_{j+1,j}
 :=u_{(j+1,r),(j,r)}:
 \mathfrak D_{j+1,r}(x_{j+1})
 \longrightarrow\mathfrak D_{j,r}(x_j),
\end{equation}
la composante du transport commun \eqref{eq:pdfv2-cor-naturalidad-unica} qui tronque la dernière émission, réduit les coefficients et précompose les ports. Par composition,
\begin{equation}
\label{eq:pdfv2-terminal-transicion-lector-composicion}
 u^{\mathfrak D}_{m,k}
 =u^{\mathfrak D}_{\ell,k}u^{\mathfrak D}_{m,\ell}
 \qquad(m\geq\ell\geq k).
\end{equation}
\begin{align}
 \mathsf L_{k,k}
 &:=(\mathfrak D_{0,r}(x_0))\star
 \mathop{\star}_{i=1}^{k}
 \bigl(\mathfrak D_{i,r}(x_i),u^{\mathfrak D}_{i,i-1},
       \mathcal O(\mathfrak e_{\varepsilon_i})\bigr),
 \label{eq:pdfv2-terminal-semilla-l}\\
 \operatorname{Upd}^{\mathsf L}_{\varepsilon}\mathsf L_{k,j}
 &:=
 \mathsf L_{k,j}\star
 \bigl(\mathfrak D_{j+1,r}(x_{j+1}),
 u^{\mathfrak D}_{j+1,j},\mathcal O(\mathfrak e_\varepsilon)\bigr).
 \label{eq:pdfv2-terminal-update-l}
\end{align}
La totalité de \(\mathfrak D\) et \eqref{eq:pdfv2-terminal-transicion-lector-composicion} prouvent la non-vacuité et la naturalité.

Les treize germes précédents forment un élément unique
\begin{equation}
\label{eq:pdfv2-terminal-semilla-trece}
 \mathfrak r^0_{k,r}:=
 \langle
 \mathsf F_{k,k};\mathsf E_{k,k};\mathsf R_{k,k};\mathsf N_{k,k};
 \mathsf O_{k,k};\mathsf K_{k,k};\mathsf M_{k,k};\mathsf B_{k,k};
 \mathsf G_{k,k};\mathsf X_{k,k};\mathsf S_{k,k};\mathsf T_{k,k};
 \mathsf L_{k,k}
 \rangle\in\mathfrak R_{k,r}.
\end{equation}
L'action d'une arête est l'application totale
\begin{equation}
\label{eq:pdfv2-terminal-actualizacion-unica}
 \operatorname{Upd}_\varepsilon:
 \mathfrak R_{j,r}(\gamma_{0,j})
 \longrightarrow\mathfrak R_{j+1,r}(\gamma_{0,j+1}),\qquad
 \operatorname{Upd}_\varepsilon:=
 \langle\operatorname{Upd}^{\mathsf F}_\varepsilon,\ldots,
 \operatorname{Upd}^{\mathsf L}_\varepsilon\rangle,
\end{equation}
avec les treize actions \eqref{eq:pdfv2-terminal-update-f}--\eqref{eq:pdfv2-terminal-update-l}. Enfin, le registre terminal est \emph{engendré} par le pliage
\begin{equation}
\label{eq:pdfv2-terminal-registro-trece}
 \boxed{
 \mathfrak r_{k,N}(h):=
 \operatorname{Upd}_{\varepsilon_N}\circ\cdots\circ
 \operatorname{Upd}_{\varepsilon_{k+1}}(\mathfrak r^0_{k,r}).}
\end{equation}
Ce n'est pas le dépaquetage d'une copie intacte de l'état. Sur l'image de ce pliage, on définit, composante par composante, la troncature
\begin{equation}
\label{eq:pdfv2-terminal-truncamiento-trece-def}
 \operatorname{tr}_{13;j+1,j}
 \mathfrak r_{k,j+1}(h\star\varepsilon)
 :=\mathfrak r_{k,j}(h).
\end{equation}
Elle est bien définie : les champs \(\mathsf E\) et \(\mathsf G\) conservent la liste ordonnée des émissions et des préfixes, si bien que l'égalité de deux registres implique l'égalité de leurs histoires étiquetées. Dans \(\mathsf F,\mathsf E,\mathsf R,\mathsf N,\mathsf G,\mathsf T,\mathsf L\), la dernière entrée ou le dernier bloc est supprimé ; dans \(\mathsf O,\mathsf K,\mathsf M,
\mathsf B\), la formule cumulative est appliquée au préfixe retrouvé ; et dans \(\mathsf X,\mathsf S\), le cône des queues est restreint. On n'inverse pas une union idempotente et l'on ne sélectionne pas une continuation.

\begin{proposition}[Totalité et naturalité des treize actualisations]
\label{prop:pdfv2-terminal-trece-total-natural}
Tous les germes de \eqref{eq:pdfv2-terminal-semilla-trece} existent, chaque \(\operatorname{Upd}_\varepsilon\) est total et, pour toute symétrie cohérente \(a\) et toute troncature d'horizon,
\begin{equation}
\label{eq:pdfv2-terminal-trece-naturalidad}
 a_*\operatorname{Upd}_\varepsilon
 =\operatorname{Upd}_{a\varepsilon}a_*,
 \qquad
 \operatorname{tr}_{13;j+1,j}
 \operatorname{Upd}_\varepsilon\mathfrak r_{k,j}(h)
 =\mathfrak r_{k,j}(h).
\end{equation}
\end{proposition}

\begin{proof}
L'existence de chaque germe a été exposée dans les treize paragraphes. Dans \(\mathsf F,\mathsf E,\mathsf G,\mathsf X,\mathsf S,\mathsf T,\mathsf L\), la formule est une concaténation avec une application déjà typée. Dans \(\mathsf R,\mathsf N,\mathsf O,\mathsf K,\mathsf M,\mathsf B\), la formule est, respectivement, la concaténation d'un bloc d'équations étiqueté, l'adjonction d'une donnée numérique, le transport de feuillets, la composition de cocycles, l'action affine et l'application de chaînes. Les identités de naturalité citées après chaque formule prouvent la première égalité de \eqref{eq:pdfv2-terminal-trece-naturalidad} composante par composante. La seconde est exactement \eqref{eq:pdfv2-terminal-truncamiento-trece-def} ; sa bonne définition utilise l'histoire ordonnée conservée par \(\mathsf E\) et \(\mathsf G\).
\end{proof}

\begin{proposition}[Certificat conjoint des cinq développements internes]
\label{prop:pdfv2-terminal-cinco-desarrollos-integrales}
Les actualisations des treize coordonnées transportent conjointement : la télescopie spectrale--polaire et le retour sans réinitialisation de \cref{thm:continuo-sp-telescopia-no-autonoma,eq:continuo-sp-no-reset} ; le bilan orthogonal, la colonne d'observation et l'identité de Stein de \cref{eq:pdfv2-sol-ortogonalidad,eq:pdfv2-sol-stein} ; la naturalité directe et adjointe ainsi que le rétracte cellulaire de \cref{eq:gauint-naturalidad-directa-adjunta,eq:gauint-sdr} ; la totalisation, le cône du retour nonadique et les fibres complètes de \cref{eq:pdfv2-coh-total-square-zero,eq:pdfv2-coh-terminal-cone,eq:pdfv2-coh-extension-multisection} ; et la diagonale d'Alexander--Whitney, la réduction de Smith et la tour de Bockstein de \cref{eq:detint-aw-natural,eq:detint-smith,eq:detint-bock-leading}. Toutes ces actions ont pour domaine le même registre engendré par \eqref{eq:pdfv2-terminal-registro-trece} ; aucune n'introduit de sous-domaine initial, d'histoire privilégiée ou de sélection motivée par une application ultérieure.
\end{proposition}

\begin{proof}
Chaque identité citée agit sur une coordonnée du registre commun, et ses applications de troncature sont des restrictions de \eqref{eq:pdfv2-terminal-truncamiento-trece-def}. La naturalité composante par composante démontrée dans \eqref{eq:pdfv2-terminal-trece-naturalidad} permet de les assembler en une seule actualisation. L'égalité des domaines et la conservation de l'histoire ordonnée empêchent de remplacer la génération corrélative par cinq projections rétrospectives.
\end{proof}

\section{Terminal multivalué sans sélecteur}

Le terminal fini engendré par \(\widehat x_k\) est la multisection
\begin{equation}
\label{eq:pdfv2-terminal-multivaluado}
 \mathfrak T_{k,N}(\widehat x_k)
 :=\left\{\mathfrak r_{k,N}(h):
            h\in\mathscr H_{k,N}(\widehat x_k)\right\}.
\end{equation}
L'expression du membre de droite renvoie l'ensemble entier. Elle ne contient aucune règle distinguant une histoire privilégiée.

L'espace ambiant terminal du niveau s'obtient en réunissant ces multisections sans identifier les registres :
\begin{equation}
\label{eq:pdfv2-terminal-ambiente-nivel}
 \mathfrak T_{k,N}
 :=\bigcup_{\widehat x_k\in\widehat{\mathcal X}^{\rm tot}_k}
   \mathfrak T_{k,N}(\widehat x_k).
\end{equation}

La version linéaire, lorsqu'il faut sommer des contributions, se définit dans le module libre par
\begin{equation}
\label{eq:pdfv2-terminal-evaluacion-todas}
 \begin{aligned}
 \mathsf E_{k,N}^{\rm all}:
 \mathbb Z[\widehat{\mathcal X}^{\rm tot}_k]&\longrightarrow
 \mathbb Z[\mathfrak T_{k,N}],\\
 [\widehat x_k]&\longmapsto
 \sum_{h\in\mathscr H_{k,N}(\widehat x_k)}
 [\mathfrak r_{k,N}(h)].
 \end{aligned}
\end{equation}
Chaque témoin apparaît avec sa multiplicité. S'il existe une mesure projective positive \(\mu\), la version isométrique utilise toutes les probabilités conditionnelles :
\begin{equation}
\label{eq:pdfv2-terminal-evaluacion-ponderada}
 \widehat{\mathsf E}_{k,N}^{\mu}e_{\widehat x_k}
 :=\sum_{h\in\mathscr H_{k,N}(\widehat x_k)}
 \mu_{\widehat x_k,N}(h)^{1/2}
 e_{\mathfrak r_{k,N}(h)}.
\end{equation}
La projectivité donne
\begin{equation}
\label{eq:pdfv2-terminal-isometria}
 \left\|\widehat{\mathsf E}_{k,N}^{\mu}e_{\widehat x_k}\right\|^2
 =\sum_{h\in\mathscr H_{k,N}(\widehat x_k)}
   \mu_{\widehat x_k,N}(h)=1.
\end{equation}
La somme n'est pas remplacée par ``l'histoire qui engendre le cylindre''.

\section{La connexion nonadique comme correspondance}

L'holonomie nonadique a été engendrée par l'odomètre et la connexion TPK, sans hypothèses d'appartenance, dans \eqref{eq:continuo-comun-gamma9-apice}--\eqref{eq:continuo-comun-gamma9-relacion}. Dans ce chapitre, nous utilisons
\begin{equation}
\label{eq:pdfv2-terminal-apice-gamma9}
 \mathsf Z_{9,k}:=\mathsf Z^{\Gamma}_{9,k}.
\end{equation}
Les applications
\begin{equation}
\label{eq:pdfv2-terminal-span-gamma9}
 \widehat{\mathcal X}^{\rm tot}_k
 \xleftarrow{\ s_{9,k}\ }
 \mathsf Z^{\Gamma}_{9,k}
 \xrightarrow{\ t_{9,k}\ }
 \widehat{\mathcal X}^{\rm tot}_{k+9}
\end{equation}
sont ceux de \eqref{eq:continuo-comun-gamma9-span} ; son application source est surjective par \cref{prop:continuo-comun-gamma9-total-natural}. Chaque point du sommet conserve les neuf émissions compatibles avec la phase. Il n'est pas remplacé par la paire de ses extrémités.

Pour \(N\geq k+9\), nous définissons le sommet d'horizon avant linéarisation :
\begin{equation}
\label{eq:pdfv2-terminal-apice-horizonte-gamma9}
 \mathsf Z^\Gamma_{9;k,N}:=
 \coprod_{z=(\widehat x_k,\boldsymbol\omega)
            \in\mathsf Z^\Gamma_{9,k}}
 \operatorname{Term}_{k+9,N}(t_{9,k}z).
\end{equation}
Un élément s'écrit \((z,q)\), où \(\boldsymbol\omega\) est le mot nonadique de neuf émissions et \(q\) sa queue jusqu'à \(N\). Les applications sur les histoires utilisent le codomaine étiqueté
\begin{equation}
\label{eq:pdfv2-terminal-gamma9-codominio-colas}
 \mathscr D^\Gamma_{9;k,N}:=
 \coprod_{z\in\mathsf Z^\Gamma_{9,k}}
 \operatorname{Term}_{k+9,N}(t_{9,k}z).
\end{equation}
Ainsi, les applications sont
\begin{align}
 S_{k,N}(z,q)&:=\boldsymbol\omega\star q,
 \label{eq:pdfv2-terminal-gamma9-s-horizonte}\\
 T_{k,N}(z,q)&:=(z,q)\in\mathscr D^\Gamma_{9;k,N}.
 \label{eq:pdfv2-terminal-gamma9-t-horizonte}
\end{align}
La première aboutit dans \(\coprod_x\mathscr H_{k,N}(x)\) ; la seconde conserve expressément le facteur étiqueté par \(z\).

L'application de liaison du sommet est l'action explicite
\begin{equation}
\label{eq:pdfv2-terminal-gamma9-enlace-apice}
 \pi^Z_{N+1,N}(z,q\star\varepsilon):=(z,q).
\end{equation}
Elle est surjective : étant donné \((z,q)\), on ajoute à la queue l'émission neutre de l'extrémité de \(q\). Le codomaine des queues possède la liaison
\begin{equation}
\label{eq:pdfv2-terminal-gamma9-enlace-codominio}
 \pi^{\mathscr D}_{N+1,N}(z,q\star\varepsilon):=(z,q).
\end{equation}
Nous notons \(\pi^{\mathscr H,j}_{N+1,N}\) l'application \eqref{eq:pdfv2-terminal-truncamiento-historias} de niveau initial \(j\). Alors
\begin{align}
 \pi^{\mathscr H,k}_{N+1,N}S_{k,N+1}
 &=S_{k,N}\pi^Z_{N+1,N},
 \label{eq:pdfv2-terminal-gamma9-cuadrado-fuente}\\
 \pi^{\mathscr D}_{N+1,N}T_{k,N+1}
 &=T_{k,N}\pi^Z_{N+1,N}.
 \label{eq:pdfv2-terminal-gamma9-cuadrado-destino}
\end{align}
Dans le premier carré, les deux parcours envoient \(\boldsymbol\omega\star q\star\varepsilon\) sur \(\boldsymbol\omega\star q\) ; dans le second, ils envoient \((z,q\star\varepsilon)\) sur \((z,q)\). Par décomposition unique d'un mot en ses neuf premières émissions et sa queue, \(S_{k,N}\) est bijectif. Le second carré est cartésien par la définition en coproduit fibré de \eqref{eq:pdfv2-terminal-apice-horizonte-gamma9}. Ce sont les carrés de Beck--Chevalley qui conservent témoins et multiplicités.

Le sommet correspondant des registres, désormais engendré par les treize actualisations, est
\begin{equation}
\label{eq:pdfv2-terminal-apice-terminal-gamma9}
 \begin{split}
 \mathsf Z^{\mathfrak T}_{9;k,N}:=\bigl\{
 (&z,q;
   \mathfrak r_{k,N}(\boldsymbol\omega\star q),
   \mathfrak r_{k+9,N}(q)):\\
 &z=(\widehat x_k,\boldsymbol\omega)\in
       \mathsf Z^\Gamma_{9,k},\quad
 q\in\operatorname{Term}_{k+9,N}(t_{9,k}z)
 \bigr\}.
 \end{split}
\end{equation}
Ses applications source et cible lisent respectivement le premier ou le second registre, mais ne suppriment jamais de leur domaine le témoin \((z,q)\) :
\begin{equation}
\label{eq:pdfv2-terminal-span-terminal}
 \mathfrak T_{k,N}
 \xleftarrow{\ s^{\mathfrak T}_{9;k,N}\ }
 \mathsf Z^{\mathfrak T}_{9;k,N}
 \xrightarrow{\ t^{\mathfrak T}_{9;k,N}\ }
 \mathfrak T_{k+9,N}.
\end{equation}
Si \(\zeta=(z,q;r_s,r_t)\), son action est
\begin{equation}
\label{eq:pdfv2-terminal-span-terminal-accion}
 s^{\mathfrak T}_{9;k,N}(\zeta):=r_s,
 \qquad
 t^{\mathfrak T}_{9;k,N}(\zeta):=r_t.
\end{equation}
Sa liaison est définie par
\begin{equation}
\label{eq:pdfv2-terminal-gamma9-enlace-registros}
 \begin{split}
 \pi^{Z,\mathfrak T}_{N+1,N}
 \bigl(&z,q\star\varepsilon;
       \mathfrak r_{k,N+1}
          (\boldsymbol\omega\star q\star\varepsilon),
       \mathfrak r_{k+9,N+1}(q\star\varepsilon)\bigr)\\
 &:=(z,q;\mathfrak r_{k,N}(\boldsymbol\omega\star q),
       \mathfrak r_{k+9,N}(q)).
 \end{split}
\end{equation}
La paire \((z,q)\) est une coordonnée du sommet, non une représentation choisie d'un triplet ; la liaison est donc bien définie même si deux registres terminaux coïncident. La compatibilité de \(\operatorname{Upd}\) avec l'effacement de la dernière arête prouve que \eqref{eq:pdfv2-terminal-gamma9-enlace-registros} commute avec les deux applications de \eqref{eq:pdfv2-terminal-span-terminal}.

Si l'on souhaite une action sur des modules libres, cette correspondance est linéarisée en sommant tous les témoins :
\begin{equation}
\label{eq:pdfv2-terminal-linealizacion-span}
 [\Gamma_9]_{k,N}[r]
 :=\sum_{\zeta\in(s^{\mathfrak T}_{9;k,N})^{-1}(r)}
 [t^{\mathfrak T}_{9;k,N}(\zeta)].
\end{equation}
La formule conserve les multiplicités. L'égalité des phases après neuf pas ne contracte ni les états, ni les témoins, ni leurs registres.

\section{Unicité de la décomposition par préfixes et naturalité du système de troncatures}

Pour \(k\leq\ell\leq N\), toute histoire de \(\mathscr H_{k,N}(\widehat x_k)\) possède un unique préfixe \(p\in\mathscr P_{\ell\mid k}(\widehat x_k)\). Il existe donc une décomposition disjointe par préfixes, exprimée sans en choisir aucun :
\begin{equation}
\label{eq:pdfv2-terminal-particion-historias}
 \mathscr H_{k,N}(\widehat x_k)
 =\bigsqcup_{p\in\mathscr P_{\ell\mid k}(\widehat x_k)}
 \left\{p\star q:q\in\mathscr H_{\ell,N}(p)\right\}.
\end{equation}
Soit \(\operatorname{Conc}_p\) l'opération qui place en tête le registre du préfixe \(p\) et applique successivement \eqref{eq:pdfv2-terminal-actualizacion-unica}. Nous définissons la transition commune sur un système complet de terminaux par
\begin{equation}
\label{eq:pdfv2-terminal-transicion-comun}
 \mathfrak u_{\ell,k}
 \left(
   \bigl(\mathfrak T_{\ell,N}(p)\bigr)_{p\in\mathscr P_{\ell\mid k}(\widehat x_k)}
 \right)
 :=\bigcup_{p\in\mathscr P_{\ell\mid k}(\widehat x_k)}
   \operatorname{Conc}_p\mathfrak T_{\ell,N}(p).
\end{equation}

En écrivant \(\mathbf G_{k,N}(\widehat x_k):=\mathfrak T_{k,N}(\widehat x_k)\), toute la naturalité se concentre dans l'identité unique
\begin{equation}
\label{eq:pdfv2-terminal-naturalidad-unica}
 \boxed{
 \mathfrak u_{\ell,k}
 \left(
   \bigl(\mathbf G_{\ell,N}(p)\bigr)_{p\in\mathscr P_{\ell\mid k}(\widehat x_k)}
 \right)
 =\mathbf G_{k,N}(\widehat x_k).}
\end{equation}
Il n'existe pas de loi distincte pour chaque discriminant interne. Les cinq caractéristiques obéissent à \eqref{eq:pdfv2-terminal-naturalidad-unica} puisque chacune s'actualise au moyen du même préfixe, de la même extension et du même registre.

La correspondance nonadique est compatible avec cette identité au moyen des applications déjà construites, non d'une condition implicite. À tout horizon, \eqref{eq:pdfv2-terminal-gamma9-cuadrado-fuente} et \eqref{eq:pdfv2-terminal-gamma9-cuadrado-destino} conservent le témoin nonadique avec sa multiplicité, et \eqref{eq:pdfv2-terminal-gamma9-enlace-registros} fait de même dans les treize coordonnées. Si \(a\) est une symétrie cohérente de l'état commun, son action
\[
 a^Z(z,q):=(a^Zz,aq)
\]
satisfait
\begin{equation}
\label{eq:pdfv2-terminal-gamma9-equivarianza-completa}
 S a^Z=aS,\qquad T a^Z=aT,\qquad
 \pi^Za^Z=a^Z\pi^Z,\qquad
 s^{\mathfrak T}a^Z=as^{\mathfrak T},\qquad
 t^{\mathfrak T}a^Z=at^{\mathfrak T}.
\end{equation}
Ces égalités se vérifient entrée par entrée au moyen de \eqref{eq:continuo-comun-gamma9-naturalidad} et de la naturalité des treize actualisations. Elles constituent la condition concrète permettant d'utiliser \eqref{eq:pdfv2-terminal-linealizacion-span} ; une simple égalité de valeurs visibles ne la remplace pas.

\section{Passage à la limite sans intervertir image et limite}

La troncature d'une histoire et de chacun des treize champs définit
\begin{equation}
\label{eq:pdfv2-terminal-transicion-horizonte}
 \pi^{\mathfrak T}_{N+1,N}:
 \mathfrak T_{k,N+1}(\widehat x_k)\longrightarrow
 \mathfrak T_{k,N}(\widehat x_k),
 \qquad
 \pi^{\mathfrak T}_{N+1,N}\mathfrak r_{k,N+1}(h)
 =\mathfrak r_{k,N}(h_{\leq N}).
\end{equation}
Elle est surjective par \eqref{eq:pdfv2-terminal-truncamiento-historias}. Le terminal complet est alors
\begin{equation}
\label{eq:pdfv2-terminal-limite-inverso}
 \mathfrak T_{k,\infty}(\widehat x_k)
 :=\left\{(r_N)_{N\geq k}:
 r_N\in\mathfrak T_{k,N}(\widehat x_k),\quad
 \pi^{\mathfrak T}_{N+1,N}r_{N+1}=r_N\right\}.
\end{equation}
Les ensembles finis de \eqref{eq:pdfv2-terminal-multivaluado} sont non vides et leurs transitions sont surjectives ; ainsi
\begin{equation}
\label{eq:pdfv2-terminal-limite-no-vacio}
 \mathfrak T_{k,\infty}(\widehat x_k)\neq\varnothing.
\end{equation}

\begin{corollary}[Persistance simultanée des cinq développements]
\label{cor:pdfv2-terminal-limite-cinco-desarrollos}
Dans chaque élément de \(\mathfrak T_{k,\infty}(\widehat x_k)\) persistent simultanément la limite forte \(\mathcal T_\infty\) de la télescopie polaire, le système de Gram et de pertes, la mesure d'incidence avec ses deux naturalités, le cône \(\operatorname{Cone}(I-U_{\Gamma_9,N})\) et le système compatible des formes de Smith et des pages de Bockstein. Aucune de ces données ne s'obtient en reconstruisant la seule valeur terminale visible.
\end{corollary}

\begin{proof}
La surjectivité des troncatures conserve les témoins complets. Les équations de télescopie, de naturalité et de réduction de \cref{prop:pdfv2-terminal-cinco-desarrollos-integrales} commutent avec ces troncatures ; par universalité de la limite inverse, elles définissent un unique système compatible sur chaque histoire survivante.
\end{proof}

Le sommet nonadique à la limite est défini, non inféré d'une image :
\begin{equation}
\label{eq:pdfv2-terminal-gamma9-apice-limite-raw}
 \mathsf Z^\Gamma_{9;k,\infty}
 :=\varprojlim_{N\geq k+9}
 \bigl(\mathsf Z^\Gamma_{9;k,N},\pi^Z_{N+1,N}\bigr).
\end{equation}
Puisque les liaisons ne modifient pas le témoin fini \(z\) et tronquent seulement sa queue, il existe l'isomorphisme explicite
\begin{equation}
\label{eq:pdfv2-terminal-gamma9-apice-limite-msect}
 \mathsf Z^\Gamma_{9;k,\infty}
 \xrightarrow{\ \cong\ }
 \coprod_{z\in\mathsf Z^\Gamma_{9,k}}
 \operatorname{Msect}(t_{9,k}z),\qquad
 ((z,q_N))_N\longmapsto(z,(q_N)_N).
\end{equation}
L'inverse prend un témoin \(z\) et une queue compatible \((q_N)_N\), puis forme \(((z,q_N))_N\). Les deux compositions sont des identités coordonnée par coordonnée. Les applications limites sont
\begin{equation}
\label{eq:pdfv2-terminal-gamma9-mapas-limite-raw}
 S_{k,\infty}((z_N)_N):=(S_{k,N}z_N)_N,\qquad
 T_{k,\infty}((z_N)_N):=(T_{k,N}z_N)_N;
\end{equation}
sont bien définies précisément par \eqref{eq:pdfv2-terminal-gamma9-cuadrado-fuente} et \eqref{eq:pdfv2-terminal-gamma9-cuadrado-destino}.

On définit de même le sommet des registres
\begin{equation}
\label{eq:pdfv2-terminal-gamma9-apice-registros-limite}
 \mathsf Z^{\mathfrak T}_{9;k,\infty}
 :=\varprojlim_{N\geq k+9}
 \bigl(\mathsf Z^{\mathfrak T}_{9;k,N},
       \pi^{Z,\mathfrak T}_{N+1,N}\bigr).
\end{equation}
Les correspondances \eqref{eq:pdfv2-terminal-span-terminal} forment alors, composante par composante, le span
\begin{equation}
\label{eq:pdfv2-terminal-span-limite}
 \mathfrak T_{k,\infty}
 \xleftarrow{\ s^{\mathfrak T}_{9;k,\infty}\ }
 \mathsf Z^{\mathfrak T}_{9;k,\infty}
 \xrightarrow{\ t^{\mathfrak T}_{9;k,\infty}\ }
 \mathfrak T_{k+9,\infty}.
\end{equation}
L'image d'une limite n'a pas été identifiée à la limite d'images : \eqref{eq:pdfv2-terminal-limite-inverso}, \eqref{eq:pdfv2-terminal-gamma9-apice-limite-raw} et \eqref{eq:pdfv2-terminal-gamma9-apice-registros-limite} construisent directement les trois systèmes compatibles et leurs applications.

\section{Reconstruction interne de l’état holonomique intégral à partir de ses \(13\) coordonnées généalogiques}

La reconstruction ne dépaquette pas des noms : elle répète le pliage qui a engendré le registre. Pour \(h=(\varepsilon_{k+1},\ldots,\varepsilon_N)\), nous définissons
\begin{equation}
\label{eq:pdfv2-terminal-recuperacion-trece}
 \boxed{
 \operatorname{Rec}_{13;k,N}(h):=
 \operatorname{Upd}_{\varepsilon_N}\circ\cdots\circ
 \operatorname{Upd}_{\varepsilon_{k+1}}
 (\mathfrak r^0_{k,r}).}
\end{equation}
Par \eqref{eq:pdfv2-terminal-registro-trece}, \(\operatorname{Rec}_{13;k,N}(h)=\mathfrak r_{k,N}(h)\), mais l'égalité est désormais une identité entre deux compositions d'actions, non la projection d'une entrée préservée. Si \(\operatorname{pr}_{\mathsf A}\) désigne le projecteur du registre sur le champ \(\mathsf A\in\{\mathsf F,\mathsf E,\mathsf R,\mathsf N,\mathsf O,
\mathsf K,\mathsf M,\mathsf B,\mathsf G,\mathsf X,\mathsf S,\mathsf T,
\mathsf L\}\), alors
\begin{equation}
\label{eq:pdfv2-terminal-recuperacion-campo-fold}
 \operatorname{pr}_{\mathsf A}\operatorname{Rec}_{13;k,N}(h)
 =\operatorname{Upd}^{\mathsf A}_{\varepsilon_N}\circ\cdots\circ
  \operatorname{Upd}^{\mathsf A}_{\varepsilon_{k+1}}
  (\mathsf A_{k,k}).
\end{equation}
La preuve se fait par récurrence sur \(N-k\) : pour \(N=k\), les deux membres sont le germe ; l'étape de récurrence applique la composante \(\mathsf A\) de \(\operatorname{Upd}_{\varepsilon_N}\).

De plus,
\begin{equation}
\label{eq:pdfv2-terminal-recuperacion-compatible}
 \operatorname{Rec}_{13;k,N}(h_{\leq N})
 =\operatorname{tr}_{13,N+1,N}
   \bigl(\operatorname{Rec}_{13;k,N+1}(h)\bigr),
\end{equation}
où \(\operatorname{tr}_{13,N+1,N}\) supprime seulement la dernière actualisation et restreint son cône de queues. En conséquence, un élément de \(\mathfrak T_{k,\infty}(\widehat x_k)\) détermine un système compatible des treize coordonnées à toute profondeur.

\section{Terminalité multivaluée, naturalité et limite compatible du continu}

\begin{theorem}[Terminal multivalué, naturalité et limite du continu]
\label{thm:pdfv2-terminal-conjunto}
Pour tout \(\widehat x_k\in\widehat{\mathcal X}^{\rm tot}_k\), les règles \eqref{eq:pdfv2-terminal-historias-finita}--\eqref{eq:pdfv2-terminal-recuperacion-compatible} construisent, à partir de l'unique continu discret :
\begin{enumerate}
 \item une multisection terminale finie non vide à chaque horizon ;
 \item une évaluation qui conserve toutes les histoires et multiplicités ;
 \item la connexion nonadique comme correspondance munie de témoins ;
 \item l'unique identité de naturalité \eqref{eq:pdfv2-terminal-naturalidad-unica} ;
 \item un terminal inverse non vide et une correspondance nonadique à la limite ;
 \item la reconstruction compatible de la fibre, des extensions, des relations, des nombres, de l'orientation, de la retenue, de la mémoire, de la frontière, du parcours, de la multisection, de la survie, du terminal et du lecteur.
\end{enumerate}
Les cinq caractéristiques intrinsèques restent corrélées au sein de chaque registre et de chaque témoin d'extension.
\end{theorem}

\begin{proof}
La non-vacuité de \(\mathscr H_{k,N}(\widehat x_k)\) est non vide par \cref{prop:continuo-comun-arbol-finito-no-vacio}, et sa troncature est surjective par \cref{lem:continuo-comun-terminal-sobreyectivo}. La règle \eqref{eq:pdfv2-terminal-actualizacion-unica} applique les identités APP, TRIT et TPK sur chaque témoin \(\varepsilon_j\) ; par récurrence sur \(N-k\), elle produit tous les champs de \eqref{eq:pdfv2-terminal-registro-trece} et conserve leurs relations.

Prendre l'ensemble complet dans \eqref{eq:pdfv2-terminal-multivaluado}, ou la somme complète dans \eqref{eq:pdfv2-terminal-evaluacion-todas}, démontre que le terminal ne contient pas de sélecteur. La version pondérée est isométrique par l'égalité projective \eqref{eq:pdfv2-terminal-isometria}.

Le témoin \(z=(\widehat x_k,\boldsymbol\omega)\) porte l'action de phase de chacune de ses neuf émissions ; par \cref{prop:continuo-comun-gamma9-retorno-memoria}, il définit les deux applications de \eqref{eq:pdfv2-terminal-span-gamma9} sans transformer la relation en fonction ni assimiler phase et état. Les formules \eqref{eq:pdfv2-terminal-gamma9-s-horizonte}--\eqref{eq:pdfv2-terminal-gamma9-enlace-registros} prouvent les deux carrés d'horizon et la même construction sur les registres. Sommer l'espace source de chaque témoin donne \eqref{eq:pdfv2-terminal-linealizacion-span} sans effacer les multiplicités.

La partition \eqref{eq:pdfv2-terminal-particion-historias} est disjointe puisque toute histoire possède un préfixe déterminé de longueur \(\ell\). Concaténer chaque préfixe avec toutes ses continuations reconstruit une fois et une seule chaque histoire de \(\mathscr H_{k,N}(\widehat x_k)\). Cela prouve directement \eqref{eq:pdfv2-terminal-naturalidad-unica} pour le registre entier ; cinq démonstrations disjointes ne sont pas nécessaires.

La surjectivité de \eqref{eq:pdfv2-terminal-transicion-horizonte} et la finitude de chaque niveau donnent la non-vacuité de \eqref{eq:pdfv2-terminal-limite-inverso}. Les liaisons explicites \eqref{eq:pdfv2-terminal-gamma9-enlace-apice} et \eqref{eq:pdfv2-terminal-gamma9-enlace-registros}, avec leurs carrés, construisent les limites \eqref{eq:pdfv2-terminal-gamma9-apice-limite-raw} et \eqref{eq:pdfv2-terminal-gamma9-apice-registros-limite} ; de celles-ci découle \eqref{eq:pdfv2-terminal-span-limite}. Enfin, \eqref{eq:pdfv2-terminal-recuperacion-trece} exécute à nouveau les actualisations successives, et \eqref{eq:pdfv2-terminal-recuperacion-compatible} démontre que ces lectures passent conjointement à la limite.
\end{proof}

\section{Obstructions à la terminalité de l’état holonomique intégral : perte de fibre, d’orientation, de mémoire ou de lecteur}

\begin{proposition}[Falsificateurs du terminal conjoint]
\label{prop:pdfv2-terminal-falsadores}
Chacune des opérations suivantes remplace l'objet construit et bloque la conclusion du \cref{thm:pdfv2-terminal-conjunto} :
\begin{enumerate}
 \item envoyer un cylindre sur une seule histoire alors qu'il en admet plusieurs ;
 \item choisir une queue du champ terminal \(\mathsf T\) et supprimer les autres ;
 \item remplacer \(\Gamma_9\) par une fonction sans démontrer l'unicité de chaque cible et de son témoin ;
 \item linéariser une correspondance en ignorant les multiplicités ;
 \item identifier le retour de phase au retour d'état ;
 \item conserver l'entrée intacte comme première coordonnée et utiliser sa récupérabilité comme une génération supposée ;
 \item omettre l'un quelconque des treize champs de \eqref{eq:pdfv2-terminal-registro-trece} ;
 \item permettre à deux caractéristiques intrinsèques d'utiliser des histoires, retenues, frontières ou terminaux différents ;
 \item remplacer \eqref{eq:pdfv2-terminal-naturalidad-unica} par cinq affirmations indépendantes ;
 \item affirmer la compatibilité opératorielle de \(\Gamma_9\) sans conserver l'espace des témoins et leurs multiplicités ;
 \item intervertir sans preuve la formation des images et le passage à la limite ;
 \item présenter une liste de treize noms sans les actions \eqref{eq:pdfv2-terminal-actualizacion-unica} qui les produisent.
 \item omettre la télescopie polaire, la ligne des pertes, la naturalité adjointe, le module de toutes les continuations ou la réduction Smith--Bockstein établis dans \cref{prop:pdfv2-terminal-cinco-desarrollos-integrales} ;
 \item transformer l'un quelconque de ces cinq développements en branche autonome ou utiliser une application classique ultérieure pour sélectionner son domaine.
\end{enumerate}
\end{proposition}

\begin{proof}
Les cinq premiers actes contractent la multisection ou la correspondance. Les quatre suivants rompent le support et la naturalité communs. Les trois derniers suppriment la preuve générative ou modifient le passage à la limite. Dans chaque cas, l'objet résultant cesse de satisfaire au moins l'une de \eqref{eq:pdfv2-terminal-multivaluado}, \eqref{eq:pdfv2-terminal-span-terminal}, \eqref{eq:pdfv2-terminal-naturalidad-unica} ou \eqref{eq:pdfv2-terminal-recuperacion-compatible} ; la conclusion n'est donc pas transférable.
\end{proof}

\begin{statusbox}
L'organisation du terminal complet, la naturalité unique et le passage à la limite sont établis conjointement. La multisection survivante, les retenues, le registre, la frontière, le parcours et la connexion nonadique sont leurs antécédents structurels.\quad
La linéarisation sur les modules libres est démontrée par \eqref{eq:pdfv2-terminal-linealizacion-span} et ses carrés de Beck--Chevalley. Une promotion ultérieure à un opérateur borné sur une complétion analytique exige en outre de déclarer cette norme et de démontrer la borne ; la correspondance \eqref{eq:pdfv2-terminal-span-limite} ne dépend pas de cette interface.
\end{statusbox}
 \end{HMTScientificOwnerSubordinated}

\HMTDeferredContinuumConsequences
 
\chapter{Mesure projective et conservation évolutive de l'unité}
\label{chap:83-06}

\section{Mesure projective et conservation évolutive de l'unité}
\label{p12-83-c6-s1:chap:medida-proyectiva}

La topologie des cylindres permet de parler de proximité ; une théorie des chemins
a en outre besoin de poids. La conservation évolutive de l'unité acquiert ici
une forme probabiliste exacte : chaque cylindre répartit son poids entre ses
extensions et la somme totale demeure égale à un. L'unité ne reste pas fixe dans
une cellule ; elle est transportée à travers le raffinement.

\subsection{Poids cohérents sur les niveaux finis}

Soit \(\mu_n\) une probabilité sur l'ensemble fini
\(\mathcal S_n\). La famille est \emph{projectivement cohérente} lorsque
\begin{equation}
 \mu_n(a)
 =\sum_{\substack{b\in\mathcal S_{n+1}\\
                   \rho_{n+1,n}(b)=a}}
   \mu_{n+1}(b)
 \qquad(a\in\mathcal S_n).
 \label{p12-83-c6-s1:eq:consistencia-pesos}
\end{equation}
Cette égalité exprime la conservation locale de l'unité : le poids d'une
histoire tronquée est la somme des poids de tous ses enfants.

\begin{theorem}[Mesure généalogique sur les cylindres]
\label{p12-83-c6-s1:thm:medida-genealogica}
Toute famille de probabilités finies qui satisfait
\eqref{p12-83-c6-s1:eq:consistencia-pesos} détermine une unique probabilité borélienne
\(\mu_\infty\) sur \(\Omega_\infty\) telle que
\begin{equation}
 \mu_\infty([a]_n)=\mu_n(a).
 \label{p12-83-c6-s1:eq:medida-cilindro}
\end{equation}
\end{theorem}

\begin{proof}
\[
 \mathscr A
 :=\{\text{réunions finies de cylindres}\}
\]
est une algèbre : l'intersection de deux cylindres est vide ou est un cylindre du
niveau le plus profond. Tout élément de \(\mathscr A\) peut être raffiné en une
réunion disjointe de cylindres d'un même niveau. Si
\(A=\bigsqcup_{a\in I}[a]_n\), nous posons
\[
 \mu_0(A)=\sum_{a\in I}\mu_n(a).
\]
Itérer \eqref{p12-83-c6-s1:eq:consistencia-pesos} démontre que la valeur ne dépend pas du
niveau commun choisi.

Si un cylindre \([a]_n\) est vide, la ramification finie et le lemme de König
\cite{Koenig1936}
impliquent qu'il existe une profondeur sans descendants de \(a\). La
cohérence itérée donne alors \(\mu_n(a)=0\). C'est pourquoi \(\mu_0\) est bien
définie même pour les représentations contenant des cylindres vides et est
finiment additive.

Les cylindres sont ouverts et fermés. Si une suite décroissante
d'éléments de \(\mathscr A\) avait une intersection vide, la compacité
impliquerait que l'un d'eux est déjà vide ; c'est pourquoi \(\mu_0\) est continue en
l'ensemble vide et constitue une prémesure dénombrablement additive. Le théorème d'extension de
Carathéodory ---de manière équivalente, l'extension de Kolmogorov pour ce système
projectif fini--- la prolonge à la \(\sigma\)-algèbre engendrée par les
cylindres \cite{Caratheodory1914,Kolmogorov1933}. Celle-ci est la tribu borélienne de
\(\Omega_\infty\), et la finitude de la
mesure garantit l'unicité.
\end{proof}

La mesure peut provenir de transitions locales. Si
\(P_n(b\mid a)\ge0\) satisfait
\[
 \sum_{\rho_{n+1,n}(b)=a}P_n(b\mid a)=1,
\]
alors
\(\mu_{n+1}(b)=\mu_n(a)P_n(b\mid a)\) est cohérente. L'uniformité est un
choix possible, non une obligation. Les poids peuvent dépendre du feuillet,
de la retenue, de l'orientation, d'un coût d'action non négatif ou d'un registre pourvu
qu'ils préservent la somme. Par exemple,
\[
 P_n(b\mid a)
 =\frac{e^{-\beta a_n(b\mid a)}}
        {\sum_{b'\mapsto a}e^{-\beta a_n(b'\mid a)}},
 \qquad a_n\geq0,
\]
définit des poids positifs normalisés. La phase orthogonale d'action appartient au
noyau du \cref{chap:integral-genealogica} ; elle n'est pas insérée dans ces
probabilités positives.

\subsection{Désintégration de la mesure généalogique sur les fibres de même valeur archimédienne}

L'évaluation archimédienne transporte la mesure vers \(K\) :
\begin{equation}
 \nu:=\Val_*\mu_\infty,
 \qquad
 \nu(B)=\mu_\infty(\Val^{-1}(B)).
 \label{p12-83-c6-s1:eq:pushforward-arquimediano}
\end{equation}
La probabilité visible d'une région additionne le poids de toutes les généalogies qui
la publient. Lorsque deux sections possèdent la même valeur, leurs contributions
s'agrègent dans \(\nu\), tout en demeurant séparées dans \(\mu_\infty\).

\begin{proposition}[Désintégration par valeurs]
Si \(K\) est un espace borélien standard, il existe une famille de mesures
conditionnelles \(\{\mu_y\}_{y\in K}\), définie pour \(\nu\)-presque tout \(y\),
telle que \(\mu_y\) est portée par \(\mathfrak G_y=\Val^{-1}(y)\) et
\[
 \int_{\Omega_\infty}F\dd\mu_\infty
 =\int_K\left(\int_{\mathfrak G_y}F\dd\mu_y\right)\dd\nu(y)
\]
pour toute fonction intégrable \(F\).
\end{proposition}

Cette formule sépare valeur et généalogie au niveau de la mesure. La physique qui
ne dépend que de la valeur intègre d'abord toute la fibre ; une observable de mémoire
peut distinguer des distributions \(\mu_y\) différentes sur la même
coordonnée réelle.

\subsection{Conditionnement comme élagage des futurs}

Soit \(A\subseteq\Omega_\infty\) un événement observable avec
\(\mu_\infty(A)>0\). La mesure conditionnelle est
\begin{equation}
 \mu_\infty(B\mid A)
 =\frac{\mu_\infty(B\cap A)}{\mu_\infty(A)}.
 \label{p12-83-c6-s1:eq:condicionamiento-genealogico}
\end{equation}
Lorsque \(A\) est la lecture d'un registre, son intersection avec chaque cylindre
élimine les extensions incompatibles. La normalisation ne crée pas de nouvelles branches :
elle redistribue l'unité entre celles qui ont survécu.

\begin{proposition}[Mise à jour compatible]
Si \(A\) est une réunion de cylindres au niveau \(m\), alors pour tout \(n\ge m\)
la distribution conditionnelle sur \(\mathcal S_n\) s'obtient en supprimant les
états dont les troncatures n'appartiennent pas à \(A\) et en renormalisant les poids
restants. Les distributions conditionnelles continuent de satisfaire
\eqref{p12-83-c6-s1:eq:consistencia-pesos}.
\end{proposition}

\begin{proof}
L'identité est la règle de Bayes sur les ensembles finis. La somme des
poids des enfants conditionnés coïncide avec le poids conditionné du parent,
car intersection et réunion disjointe commutent.
\end{proof}

\subsection{Conservation combinatoire et complétion dimensionnelle}

La même loi apparaît dans la complétion binomiale :
\[
 1=\sum_{r=0}^{d}\binom dr
 \left(\frac9{10}\right)^{d-r}
 \left(\frac1{10}\right)^r.
\]
En dimension trois,
\[
 1000=729+243+27+1,
\]
de sorte que l'intérieur, les nouvelles faces, les arêtes extrêmes et le sommet final forment
une partition de l'unité décimale. Le passage
\(729+270=999\to729+271=1000\) distingue l'état antérieur à la clôture de
l'unité complétée. Dans le système de chemins, l'identité analogue est
\[
 \mu([a]_n)=\sum_{b\mapsto a}\mu([b]_{n+1}).
\]
L'unité évolue en se distribuant dans les strates et se retrouve en additionnant
la généalogie complète.

\begin{figure}[htbp]
\centering
\begin{tikzpicture}[>=Latex,node distance=12mm and 17mm,
  box/.style={draw=hmtblue,rounded corners,fill=hmtlightblue,align=center,
              minimum width=28mm,minimum height=8mm},
  leaf/.style={draw=hmtgreen,rounded corners,fill=hmtlightgreen,align=center,
               minimum width=22mm,minimum height=7mm}]
\node[box] (a) {$[a]_n$\\poids $\mu_n(a)$};
\node[leaf,below left=of a] (b1) {$[b_1]_{n+1}$\\$\mu_{n+1}(b_1)$};
\node[leaf,below=of a] (b2) {$[b_2]_{n+1}$\\$\mu_{n+1}(b_2)$};
\node[leaf,below right=of a] (b3) {$[b_3]_{n+1}$\\$\mu_{n+1}(b_3)$};
\draw[->] (a)--(b1); \draw[->] (a)--(b2); \draw[->] (a)--(b3);
\node[below=20mm of b2,align=center,text width=.78\textwidth] (eq)
 {$\mu_n(a)=\mu_{n+1}(b_1)+\mu_{n+1}(b_2)+\mu_{n+1}(b_3)$ :\\
  l'unité se conserve lorsque la généalogie est raffinée.};
\end{tikzpicture}
\caption{Conservation projective du poids sur les cylindres.}
\label{p12-83-c6-s1:fig:conservacion-peso-cilindros}
\end{figure}

\begin{statusbox}
La mesure projective se construit sur le système projectif HMT.
Le théorème d'extension relève des mathématiques classiques ; l'identification de ses
cylindres aux états TPK est exacte une fois les poids de transition fixés.
La sélection d'une mesure physique concrète appartient au modèle dynamique et ne
se déduit pas uniquement du recensement des branches.
\end{statusbox}

La mesure projective est caractérisée par la positivité, la conservation du
poids entre parents et enfants et le conditionnement au sein de l'événement observé. Une
réalisation expérimentale complète ajoute la préparation et le protocole
d'échantillonnage ; avec ces données, la même mesure produit des probabilités de lecture
comparables à des fréquences.
 
\chapter{Morphismes archimédien et exceptionnel à partir d’un domaine commun}
\label{chap:83-17}

\section{Morphismes archimédien et exceptionnel à partir d'un domaine enrichi commun}
\label{p12-83-c17-s1:chap:doble}
\index{double projection}
\index{Moonshine}

La double projection compare deux évaluations du domaine enrichi commun
\(X_\infty^{\rm cal}\) déjà construit dans le
\cref{p12-83-c17-s2:thm:lector-excepcional-global-rev6}. L'une produit des valeurs
archimédiennes ; l'autre, avec un étalonnage de son codomaine, conserve
l'incidence combinatoire qui conduit aux codes, aux plans et aux groupes
d'automorphismes. Les deux évaluations partent de l'unique état initial
\(X_\infty^{\rm cal}\) et préservent la provenance commune dans des codomaines
différenciés.

\subsection{Domaines des \(2\) projections}

Sur \(X_\infty^{\rm cal}\), les cylindres se contractent en un point et la
frontière admet la prolongation d'incidence. Dans cet espace de résidence, nous fixons
\(\mathbf{Arch}:=\mathbb R\). L'évaluation archimédienne et l'évaluation
exceptionnelle sont respectivement
\[
\mathfrak R:X_\infty^{\rm cal}\longrightarrow\mathbf{Arch},
\qquad
\mathfrak E:X_\infty^{\rm cal}\times\mathscr C_{\rm exc}
\longrightarrow\mathbf{Exc}.
\]
Ici, \(\mathscr C_{\rm exc}\) contient l'ordre projectif, la jauge de
Paley, la dodécade, l'alignement d'incidence et la chambre de réseau
spécifiés dans chaque construction. La première application produit des valeurs ; la
seconde, des codes, des plans, des réseaux et des automorphismes. Le théorème global
précédent est la source de la compatibilité et permet de comparer
des invariants obtenus par des opérations distinctes sans intersecter de nouveau
des domaines définis séparément. La comparaison finie de cet ouvrage
s'effectue dans le drapeau
\(B_K\subset H^-_\alpha\).

\begin{figure}[H]
\centering
\resizebox{\textwidth}{!}{\begin{tikzpicture}[font=\small,
  big/.style={draw,rounded corners=3mm,align=center,minimum width=3.7cm,minimum height=1.25cm},
  arr/.style={-{Stealth[length=2.6mm]},very thick}]
  \node[big,draw=hmtgold,fill=hmtlightgold,font=\bfseries] (x) at (0,0) {état holonomique intégral\\\(\XTPK\)};
  \node[big,draw=hmtblue,fill=hmtlightblue] (arq) at (5,2.6) {évaluation archimédienne\\cylindres \(\to\RR\)};
  \node[big,draw=hmtgreen,fill=hmtlightgreen] (exc) at (5,0) {projection exceptionnelle\\Witt, 3-voisin et Leech};
  \draw[arr,hmtblue] (x)--node[above left]{\(\mathfrak R\)} (arq);
  \draw[arr,hmtgreen] (x)--node[above]{\(\mathfrak E\)} (exc);
  \node[big,draw=hmtgray,fill=hmtpaper] (value) at (10,2.6) {valeur réelle\\classe d'évaluation};
  \node[big,draw=hmtgray,fill=hmtpaper] (sym) at (10,0) {automorphisme\\de l'incidence};
  \draw[arr,hmtblue] (arq)--(value);
  \draw[arr,hmtgreen] (exc)--(sym);
  \draw[dashed,hmtred,thick,<->] (value.south west) to[bend right=18] node[fill=white,inner sep=1pt]{invariantes comunes} (sym.north west);
\end{tikzpicture}
 }
\caption{Schéma typé des deux projections. L'évaluation archimédienne
identifie des états d'une même fibre ; l'évaluation exceptionnelle, relative
à son étalonnage, conserve la structure d'incidence et les automorphismes.}
\label{p12-83-c17-s1:fig:doble-proyeccion}
\end{figure}

\subsection{Topologie de l'espace des histoires}

La formulation au moyen de cylindres admet une réalisation standard comme
limite projective. Soient \(\mathsf S_n\) les ensembles finis d'états
admissibles à la profondeur \(n\), munis de la topologie discrète, et soient
\[
 \tau_{n+1,n}:\mathsf S_{n+1}\longrightarrow\mathsf S_n
\]
les applications de troncature. Nous supposerons dans cette section qu’elles sont
surjectives et définirons
\[
 \mathcal X=\varprojlim_n\mathsf S_n.
\]
Si \(x\ne y\), soit \(m(x,y)\) la première profondeur à laquelle leurs
composantes diffèrent. La fonction
\[
 d(x,y)=3^{-m(x,y)},\qquad d(x,x)=0,
\]
est une ultramétrique qui induit la topologie de la limite projective.

\begin{theorem}[Espace compact des histoires]
\label{p12-83-c17-s1:thm:compacto-historias}
\label{thm:compacto-historias}
L'espace \(\mathcal X\) est compact et totalement discontinu. Lorsque
toute histoire admet des bifurcations à des profondeurs arbitrairement grandes,
\(\mathcal X\) n'a pas de points isolés et est homéomorphe à l'espace de
Cantor.
\end{theorem}

\begin{proof}
Le produit \(\prod_n\mathsf S_n\) est compact. Les équations
\(\tau_{n+1,n}(x_{n+1})=x_n\) définissent un sous-ensemble fermé, de sorte que
\(\mathcal X\) est compact. Les ensembles d'histoires ayant un préfixe fixe
sont à la fois ouverts et fermés et forment une base ; cela prouve la
discontinuité totale. Sous l'hypothèse de bifurcation arbitrairement profonde,
aucun cylindre ne contient d'histoire isolée. La caractérisation des
espaces compacts, métriques, parfaits et totalement discontinus conclut
la dernière affirmation.
\end{proof}

Associons à chaque \(s\in\mathsf S_n\) un intervalle fermé \(I_n(s)\), de manière compatible avec la troncature, et supposons que
\[
 \operatorname{diam}I_n(s)\leq C\lambda^n,
 \qquad 0<\lambda<1.
\]
L'intersection des intervalles correspondant à une histoire contient
alors un unique point.

\begin{theorem}[Évaluation archimédienne]
\label{p12-83-c17-s1:thm:evaluacion-holder}
\label{thm:evaluacion-holder}
L'application
\[
 \operatorname{val}:\mathcal X\longrightarrow\mathbb R,
 \qquad
 \{\operatorname{val}(x)\}=\bigcap_n I_n(x_n),
\]
est continue et höldérienne par rapport à \(d\). Si
\(x\sim_{\mathrm{Arch}}y\) équivaut à
\(\operatorname{val}(x)=\operatorname{val}(y)\), l'application induite
est un homéomorphisme
\[
 \mathcal X/{\sim_{\mathrm{Arch}}}
 \;\cong\;\operatorname{val}(\mathcal X).
\]
\end{theorem}

\begin{proof}
Deux histoires ayant un préfixe commun de longueur \(n\) ont des valeurs dans un
même intervalle de diamètre inférieur ou égal à \(C\lambda^n\). Puisque
\(d(x,y)=3^{-n}\), cette estimation est une borne de Hölder. L'application
induite sur le quotient est une bijection continue d'un compact sur
un espace séparé et, par conséquent, un homéomorphisme.
\end{proof}

Le théorème formalise la projection archimédienne : la valeur retient la classe d’évaluation et élimine la distinction entre trajectoires situées dans une même fibre. Son image est \(\operatorname{val}(\mathcal X)\) ; son identification à la totalité de \(\mathbb R\) exige de démontrer également la surjectivité.

\subsection{Séparation par les \(2\) projections}

\begin{theorem}[Séparation conjointe]
\label{p12-83-c17-s1:thm:separacion-conjunta}
\label{thm:separacion-conjunta}
Soient \(P_1:\mathcal X\to Y_1\) et \(P_2:\mathcal X\to Y_2\) des applications
continues à codomaines séparés. Si
\[
 x\ne y\quad\Longrightarrow\quad
 P_1(x)\ne P_1(y)\ \text{ou}\ P_2(x)\ne P_2(y),
\]
l'application diagonale
\[
 (P_1,P_2):\mathcal X\longrightarrow Y_1\times Y_2
\]
est un plongement topologique.
\end{theorem}

\begin{proof}
L'hypothèse exprime exactement l'injectivité de l'application diagonale.
Une application continue et injective d'un compact dans un espace séparé
est un homéomorphisme sur son image.
\end{proof}

Une jauge déclarée \(c\in\mathscr C_{\rm exc}\) étant fixée, soit
\[
 \mathfrak E_c:X_\infty^{\rm cal}\longrightarrow\mathbf{Exc},
 \qquad
 \mathfrak E_c(x):=\mathfrak E(x,c).
\]
Appliqué à HMT, le théorème prend
\[
 P_1=\mathfrak R,\qquad P_2=\mathfrak E_c
\]
sur le domaine commun, ou sur le quotient par les changements de coordonnées déclarés. L’énoncé global de séparation exige de vérifier qu’aucun couple de trajectoires non équivalentes ne demeure dans une même fibre des deux applications.

À un niveau fini \(S\), cette condition est décidable. Étant donnés
\(p:S\to A\), \(q:S\to B\) et une relation d'équivalence des coordonnées,
on énumère les couples \(x\not\sim y\) qui satisfont simultanément
\[
 p(x)=p(y),\qquad q(x)=q(y).
\]
L'absence de tels couples équivaut à l'injectivité conjointe à ce
niveau. Le calcul doit effectuer les deux applications séparément : lorsque
\(q=f\circ p\), la seconde possède les mêmes fibres que la première et ne
récupère pas l'information éliminée par \(p\).

Au niveau développé ici, l'égalité
\[
 B_K=H_e\cap P_\pi
\]
constitue un témoin fini de compatibilité entre les deux branches. Elle n’établit pas à elle seule la séparation conjointe à toutes les profondeurs ; cet énoncé correspond exactement à l’hypothèse du \cref{p12-83-c17-s1:thm:separacion-conjunta}.

\subsection{Quotient archimédien par les fibres d'évaluation}

Une section HMT contient bloc, orientation, retenue, étiquette d'observable
et coordonnée d'extension.
L'évaluation \(\val\) contracte ses cylindres en un point. Si deux sections
\(x,y\) satisfont
\[
\val(x)=\val(y),
\]
l'évaluation archimédienne les identifie même si leur transport est distinct. En
ce sens précis,
\[
\boxed{\text{une valeur réelle est une classe d’équivalence de l’évaluation}.}
\]
Dans tout sous-domaine où somme, produit et évaluation passent au quotient de manière compatible, l’image hérite de ces opérations. Cette section ne construit pas pour autant tout \(\RR\) et ne prouve pas la surjectivité sur \(\RR\).

Une identification globale à \(\RR\) exigerait l’existence de la limite, la détermination de son image, le passage au quotient de la somme et du produit, ainsi que la caractérisation des fibres. Les fenêtres finies et le théorème conditionnel de limite projective fournissent la structure ; le prolongement à toute profondeur exige les hypothèses de survie unitaire et de compatibilité.

\subsection{Invariant commun avec le système de Witt}

Le vecteur \(K\) définit le support supérieur
\[
B_K=\{m:K_m\ge729\}=\{4,6,9,10\}.
\]
Dans l'atlas de Witt apparaissent les supports
\[
P_\pi=\{2,4,5,6,7,8,9,10,11\},
\]
\[
H_e=\{1,3,4,6,9,10\},
\qquad
H_\varphi=\{1,2,3,4,7,9\}.
\]
Le support des retenues négatives de \(\alpha\) est
\[
H_\alpha^-=\{4,5,6,7,9,10\}.
\]
On vérifie les identités de supports
\[
\boxed{
B_K=H_e\cap P_\pi=H_e\cap H_\alpha^-,}
\]
et
\[
\boxed{H_\alpha^-=B_K\sqcup\{5,7\}.}
\]
La polarité de Witt envoie les quatre points de \(B_K\) sur les paires
\[
\{1,3\},\quad\{2,11\},\quad\{8,12\},\quad\{5,7\}.
\]
Trois reçoivent l’étiquette positive ; le dernier complète le support négatif. Ces identités sont exactes dans la jauge déclarée, mais constituent des contrôles conditionnels : le vecteur, le seuil, les mots et la jauge de Paley ne sont pas des entrées statistiquement indépendantes. L’égalité est obtenue par évaluation directe et ne constitue pas une estimation probabiliste.

\subsection{Plan résiduel sur une dodécade et relèvement des hexades}
\label{p12-83-c17-s1:sec:w12-w24-elevacion}

La relation entre les deux plans de Witt s'obtient au sein du code
binaire étendu de Golay et non par une identification de cardinalités.
Soit
\[
\mathcal G_{24}\subset\mathbb F_2^{24}
\]
le code binaire étendu de Golay, et notons \(\mathcal O_{24}\) la famille
des supports de ses mots de poids huit. C'est un résultat classique que
\((\{1,\ldots,24\},\mathcal O_{24})\) est le système de Steiner
\(W_{24}=S(5,8,24)\) \cite{MacWilliamsSloane1977,ConwaySloane1999}.
Fixons un mot de poids douze et écrivons \(D\) pour son support. Nous définissons
\begin{equation}
 \mathcal H(D)
 :=\{O\cap D:O\in\mathcal O_{24},\ |O\cap D|=6\}.
 \label{p12-83-c17-s1:eq:hexadas-residuales-dodecada}
\end{equation}
La famille \(\mathcal H(D)\) est intrinsèque au couple
\((\mathcal G_{24},D)\) : elle n'exige pas de choisir des coordonnées dans \(D\).

\begin{theorem}[Plan résiduel d'une dodécade]
\label{p12-83-c17-s1:thm:dodecada-w12}
Le système d'incidence
\[
 (D,\mathcal H(D))
\]
est un système de Steiner \(S(5,6,12)\). En particulier,
\(|\mathcal H(D)|=132\), et tout sous-ensemble de cinq points de \(D\) est
contenu dans une unique hexade de \(\mathcal H(D)\).
\end{theorem}

\begin{proof}
Soit \(A\subset D\), avec \(|A|=5\). Puisque les octades forment
\(S(5,8,24)\), il existe une unique octade \(O_A\) qui contient \(A\). Si
\(j=|O_A\cap D|\), la somme binaire des mots caractéristiques de
\(O_A\) et \(D\) appartient à \(\mathcal G_{24}\) et a pour poids
\[
 \operatorname{wt}(\mathbf1_{O_A}+\mathbf1_D)=8+12-2j=20-2j.
\]
Les poids du code binaire étendu de Golay sont
\(0,8,12,16,24\). Puisque \(j\leq8\), il s'ensuit que
\(j\in\{2,4,6\}\). L'inclusion \(A\subset O_A\cap D\) donne \(j\geq5\),
donc \(j=6\). Par conséquent, \(O_A\cap D\in\mathcal H(D)\) est une hexade qui
contient \(A\).

Si deux membres de \(\mathcal H(D)\) contenaient \(A\), leurs octades
d'origine contiendraient le même ensemble de cinq points ; la propriété de
Steiner de \(W_{24}\) les ferait coïncider. Cela prouve l'unicité. Le
nombre de blocs s'obtient en comptant les incidences avec des sous-ensembles de cinq
points :
\[
 |\mathcal H(D)|
 =\frac{\binom{12}{5}}{\binom65}=132.\qedhere
\]
\end{proof}

\begin{theorem}[Relèvement unique d'une hexade]
\label{p12-83-c17-s1:thm:elevacion-unica-hexada}
Pour chaque \(H\in\mathcal H(D)\), il existe une unique octade
\(O_H\in\mathcal O_{24}\) telle que
\[
 O_H\cap D=H.
\]
En conséquence, l'application
\[
 \iota_D:\mathcal H(D)\longrightarrow\mathcal O_{24},
 \qquad H\longmapsto O_H,
\]
est bien définie et injective.
\end{theorem}

\begin{proof}
L'existence fait partie de la définition de \(\mathcal H(D)\). Si deux
octades \(O_1,O_2\) avaient pour intersection \(H\) avec \(D\), toutes deux
contiendraient n'importe quel sous-ensemble de cinq points de \(H\). L'unicité de
l'octade incidente à ces cinq points implique \(O_1=O_2\).
L'injectivité découle d'une nouvelle intersection avec \(D\).
\end{proof}

Le relèvement possède une loi locale qui sera utilisée pour comparer les cinq
régions de \(\pi\). Dans un système \(S(t,k,v)\), le nombre de blocs qui
contiennent un sous-ensemble fixé de \(s\leq t\) points est
\[
 \lambda_s=\frac{\binom{v-s}{t-s}}{\binom{k-s}{t-s}}.
\]
Par conséquent,
\begin{equation}
 \lambda_4(W_{12})
 =\frac{\binom81}{\binom21}=4,
 \qquad
 \lambda_4(W_{24})
 =\frac{\binom{20}1}{\binom41}=5.
 \label{p12-83-c17-s1:eq:lambda4-w12-w24}
\end{equation}

\begin{corollary}[Décomposition \(4+1\) sur une face]
\label{p12-83-c17-s1:cor:cuatro-mas-uno-w24}
Soit \(B\subset D\), avec \(|B|=4\). Les cinq octades de \(W_{24}\) qui
contiennent \(B\) se décomposent de manière unique en
\[
 \{O_H:H\in\mathcal H(D),\ B\subset H\}
 \;\sqcup\;\{O_B^\perp\}.
\]
Le premier ensemble contient quatre octades et satisfait
\(|O_H\cap D|=6\) ; l'octade restante satisfait
\[
 O_B^\perp\cap D=B.
\]
\end{corollary}

\begin{proof}
La première égalité de \eqref{p12-83-c17-s1:eq:lambda4-w12-w24} fournit quatre hexades
résiduelles contenant \(B\), et le
\cref{p12-83-c17-s1:thm:elevacion-unica-hexada} fournit quatre octades distinctes. La
seconde égalité de \eqref{p12-83-c17-s1:eq:lambda4-w12-w24} montre qu'il existe exactement
une cinquième octade \(O_B^\perp\).

Puisque \(B\subset O_B^\perp\cap D\), le spectre des intersections employé dans
la preuve du \cref{p12-83-c17-s1:thm:dodecada-w12} ne laisse que les possibilités quatre et
six. Une intersection de taille six serait une cinquième hexade résiduelle qui
contient \(B\), ce qui contredit \(\lambda_4(W_{12})=4\). L'intersection
est donc de taille quatre et coïncide avec \(B\).
\end{proof}

\subsubsection{Alignement avec les \(12\) positions HMT}

Le plan résiduel est défini sur les points de \(D\), tandis que le
plan obtenu à partir du code ternaire \(C_W\) est défini sur les douze
positions dodécaphasiques. Pour ne pas confondre les deux ensembles, nous appellerons
\emph{alignement binaire HMT} une bijection
\begin{equation}
 \ell:\{1,\ldots,12\}_{\rm HMT}\xrightarrow{\sim}D
 \quad\text{tel que}\quad
 \{\ell(H):H\in W_{12}^{\rm HMT}\}=\mathcal H(D).
 \label{p12-83-c17-s1:eq:alineamiento-hmt-w24}
\end{equation}
Ici, \(W_{12}^{\rm HMT}\) désigne le plan des supports de poids six
obtenu à partir du code ternaire \(C_W\). L'unicité du système de Witt
\(S(5,6,12)\) à isomorphisme près garantit l'existence de \(\ell\). Deux
alignements distincts diffèrent par des automorphismes des plans de départ et
d'arrivée ; c'est pourquoi \(\ell\) est une donnée de coordonnées et non un invariant
supplémentaire.
Le couple \((D,\ell)\) est la donnée d'étalonnage nécessaire pour transporter
le relèvement \(\iota_D\) vers le système de coordonnées HMT\@. Une fois celui-ci fixé, chaque
hexade HMT \(H\) détermine sans ambiguïté l'octade
\[
 \widehat H:=\iota_D(\ell(H)).
\]
La dodécade et l'alignement sont des données du codomaine binaire ; ils ne s'identifient pas
à une sortie décimale du noyau de transport.

Pour la face distinguée
\[
 B_K=\{4,6,9,10\},
\]
les quatre hexades incidentes sont
\[
 B_K\sqcup\{1,3\},\quad
 B_K\sqcup\{2,11\},\quad
 B_K\sqcup\{5,7\},\quad
 B_K\sqcup\{8,12\}.
\]
La troisième coïncide avec
\(H_\alpha^-=B_K\sqcup\{5,7\}\). L'alignement \(\ell\) relève ces
quatre hexades en les quatre octades résiduelles sur \(\ell(B_K)\) ; le
\cref{p12-83-c17-s1:cor:cuatro-mas-uno-w24} détermine en outre une unique octade transversale.

\subsubsection{Correspondance d'incidence entre la microfibre de \texorpdfstring{\(\pi\)}{pi} et les octades résiduelles : multiplicité \texorpdfstring{\(432+4\cdot144=1008\)}{432+4x144=1008}}

La microfibre exacte du mot
\(w_\pi=010211\) contient cinq régions décimales distinctes :
\[
\begin{array}{c|c|c}
\text{multiplicité}&U_6& e(U_6)\\
\hline
432&501|614|498|169|272|272&(5,5,5,5,5,5)\\
144&810|923|870|169|272|272&(3,3,9,5,5,5)\\
144&810|983|810|169|272|272&(3,9,3,5,5,5)\\
144&870|923|810|169|272|272&(9,3,3,5,5,5)\\
144&870|923|810|418|674|521&(9,3,3,7,1,7).
\end{array}
\]
Ici, \(e(U_6)\) est la signature multiplicative récupérée à partir de chaque bloc
décimal \(\overline{d_1d_2d_3}\) au moyen de
\(e=d_2-d_1\pmod {10}\). La première région est l'unique région de multiplicité
432 et de signature constante. Les quatre autres sont de multiplicité 144 ; parmi
elles, trois partagent le bloc de retour \(169|272|272\), tandis que la
quatrième possède le retour \(418|674|521\). Le catalogue distingue les cinq
lignes au moyen d'un prédicat reproductible : les trois lignes de type
\(\mathrm{ES}\) forment le triplet positif ; parmi les deux lignes de type
\(\mathrm{NO}\), celle de multiplicité maximale et de signature constante est appelée
transversale, et l'autre est appelée négative. Relativement à ce prédicat de
rôle, la partition est
\[
 \boxed{1_\perp+3_{++}+1_{--}.}
\]
Les noms de signe sont une convention ; ce qui est intrinsèque au catalogue est la
partition obtenue à partir de ses colonnes, non une polarité binaire préalable.

L'alignement régional complète la donnée \((D,\ell)\) par une bijection
qui envoie la région transversale sur \(O_{\ell(B_K)}^\perp\), la région
négative sur \(\widehat{H_\alpha^-}\) et les trois régions positives sur les
relèvements de
\[
 B_K\sqcup\{1,3\},\qquad
 B_K\sqcup\{2,11\},\qquad
 B_K\sqcup\{8,12\}.
\]
L'affectation ordonnée entre les trois régions positives et ces trois hexades
est un choix de coordonnées. Les deux lignes distinguées et le triplet sont
reproductibles au moyen du prédicat de catalogue précédent ; leur dénomination comme
négative, transversale et positives relève de l'alignement. Ainsi, la loi \(4+1\) n'identifie pas par la seule
cardinalité les cinq régions à cinq octades : elle fournit l'espace
d'incidence, tandis que \((D,\ell)\) et l'alignement régional spécifient
le transport entre les deux objets.

\subsection{Prolongation en réseau et écrans dimensionnels}

Le relèvement des hexades ouvre deux branches qu'il ne faut pas confondre. La première
conserve l'incidence binaire et prolonge \(W_{12}\) vers \(W_{24}\). La
seconde conserve le transport ternaire :
\[
C_{12}^{(3)}
\longrightarrow N(A_2^{12})
\longrightarrow\Lambda_{24}
\longrightarrow V^\natural
\longrightarrow\mathbb M
\longrightarrow\text{Moonshine}.
\]
La preuve primaire du recollement, du \(3\)-voisin sans racines et du corridor
jusqu'à Moonshine réside dans le
\cref{chap:bandera-excepcional-acumulativa,thm:corredor-hmt-moonshine}.
Il n'existe ni flèche canonique \(W_{24}\to N(A_2^{12})\), ni action
directe du Monstre sur les mots TPK : les deux branches partent du même
code étalonné et préservent des structures distinctes.

Les réductions \(243/11\), \(12\to11\to10\) et \(26\to24\) ne sont pas non plus des
identités entre numéraux. Le \cref{chap:pantallas-dualidad} construit leurs
espaces, quotients et applications : le quotient parfait ternaire, l'écran
marqué par \(K\) et la réduction isotrope du réseau lorentzien. Cette
section conserve ainsi une seule fonction : relier l'incidence finie déjà
démontrée à ses deux prolongations, sans répéter les preuves de référence.

\subsection{Involution résidu–quotient et conservation de l'information de l'état}

Pour \(n=9Q+r\), nous définissons l'énergie de réseau
\[
E_R(r,Q)=\frac{r^2}{R^2}+Q^2R^2.
\]
Alors
\[
\boxed{E_R(r,Q)=E_{1/R}(Q,r).}
\]
L'échange \((r,Q,R)\leftrightarrow(Q,r,1/R)\) est une involution exacte
du modèle de réseau : le résidu et le quotient échangent leurs fonctions.

\subsection{Compatibilité finie et prolongation conjointe}

Le témoin fini \(B_K\), le motif de retenues et le cadre de Witt fixent la
compatibilité locale. Sa prolongation n'est plus laissée à l'état de schéma
abstrait : le \cref{p12-83-c17-s2:thm:lector-excepcional-global-rev6} construit bloc par
bloc le lecteur exceptionnel avec une unique jauge initiale transportée, et le
\cref{p12-83-c17-s2:thm:doble-realizacion-global-rev6} démontre que cette lecture et
l'évaluation archimédienne se globalisent sur la même section TPK. Tronquer
l'histoire avant de lire équivaut dans les deux cas à tronquer la sortie après
lecture.

L'évaluation archimédienne identifie les histoires par leur valeur limite.
L'évaluation exceptionnelle conserve le flux visible et son incidence
Paley--Witt ; elle ne récupère pas pour autant les mémoires cachées que l'émetteur visible
ne publie pas. La double projection est donc une construction conjointe avec
deux codomaines distincts, non une identification entre le continu et la branche
exceptionnelle.

\begin{remark}[Compatibilité des deux projections]
Le support \(B_K\), le motif de retenues et le cadre de Witt fournissent les
invariants communs. La chaîne de réseaux prolonge cette compatibilité jusqu'à un
voisin sans racines et, par classification classique, jusqu'au réseau de Leech. La
branche binaire vers \(W_{24}\) demeure séparée : son morphisme d'incidence est
\(\iota_D\), défini après fixation de la dodécade et de l'alignement
\((D,\ell)\), et n'a pour codomaine ni un réseau de Niemeier ni le réseau
de Leech.
\end{remark}


\section{Projection exceptionnelle de l’état holonomique intégral : incidence compatible à toute profondeur}
\label{p12-83-c17-s2:chap:lector-excepcional-global-rev6}
\index{lecteur exceptionnel}\index{double projection}
\index{Paley--Witt}\index{système projectif}

Le témoin fini de la double projection acquiert une portée globale lorsque la
lecture exceptionnelle est appliquée à chaque bloc visible et que la jauge est
transportée, au lieu d'être choisie de nouveau, à mesure que la profondeur augmente. Cette
section construit cette famille. Le domaine demeure l'histoire TPK
enrichie ; la lecture exceptionnelle est fidèle sur son flux visible et ne
prétend pas se substituer aux champs de mémoire que ce flux ne publie pas.

\subsection{L'opération locale sur les 729 mots}

Dans la carte de Paley--Witt, on fixe
\[
A_W=
\begin{pmatrix}
0&1&1&1&1&1\\
1&0&1&2&2&1\\
1&1&0&1&2&2\\
1&2&1&0&1&2\\
1&2&2&1&0&1\\
1&1&2&2&1&0
\end{pmatrix}
\in M_6(\mathbb F_3).
\]
La multiplication matricielle donne
\[
A_W^{\mathsf T}=A_W,\qquad A_W^2=-I_6.
\]
Pour tout opérateur conjugué \(A=fA_Wf^{-1}\), nous définissons
\[
\Gamma_A:\mathbb F_3^6\longrightarrow\mathbb F_3^{12},
\qquad
\Gamma_A(w)=(w,wA).
\]
Son image est une copie étalonnée \(C_W(A)\) du code ternaire de Witt.
La projection sur les six premières coordonnées est un inverse à gauche :
\[
\operatorname{pr}_1\Gamma_A=\operatorname{id}_{\mathbb F_3^6}.
\]
Par conséquent, \(\Gamma_A\) est totale et injective sur les \(729\) mots et
n'a pas besoin de distinguer à l'avance \(\pi\), \(e\), \(\varphi\) ni aucune
autre sortie.

\subsection{Prolongations compatibles et unicité de la jauge initiale}

Soit \(X_n\) l'ensemble des préfixes TPK admissibles de profondeur \(n\), avec
troncatures
\[
p_{n,m}:X_n\longrightarrow X_m,\qquad m\leq n.
\]
Fixons un cadre exceptionnel initial \(f_0\). Puisque les changements \(g_j\) font
partie du transport enregistré par l'histoire, un préfixe étalonné est le
couple \((x,f_0)\) ; nous écrivons
\[
 X_n^{\rm cal}:=X_n\times\{f_0\},\qquad
 p_{n,m}^{\rm cal}(x,f_0):=(p_{n,m}x,f_0),
\]
et
\[
 X_\infty^{\rm cal}
 :=\varprojlim_n(X_n^{\rm cal},p_{n,m}^{\rm cal}).
\]
Dans la suite, nous abrégeons \(p_{n,m}^{\rm cal}\) en \(p_{n,m}\).
L'émission visible compatible est
\[
\varepsilon_n:X_n^{\rm cal}\longrightarrow(\mathbb F_3^6)^n,
\qquad
\varepsilon_n(x)=(w_0,\ldots,w_{n-1}),
\]
et satisfait
\[
\operatorname{pr}_{n,m}\varepsilon_n
=\varepsilon_m p_{n,m}.
\]
Cette application distingue ses trois types :
\[
|\mathcal U_6|=468,\qquad
|\operatorname{im}w_6|=243,\qquad
|\mathbb F_3^6|=729.
\]
L'émission \(w_j\) ne s'identifie pas à l'état complet qui conserve
feuillet, quotient, retenue, orientation, phase et survie.

La jauge exceptionnelle se compose d'une carte de Witt, d'une dodécade, d'un
alignement, d'une origine de réseau et d'une chambre orientée. Elle est fixée une seule
fois dans le bloc initial. Pour
\(x\in X_n^{\rm cal}\), le transport de l'histoire fournit
\[
 g_j(x)\in\operatorname{GL}_6(\mathbb F_3),
 \qquad 0\le j<n-1,
\]
et \(g_j(x)\) ne dépend que du préfixe
\(p_{n,j+1}(x)\). Le cadre est transporté par
\[
 f_{j+1}(x)=g_j(x)f_j(x),\qquad
 A_j=f_jA_Wf_j^{-1}.
\]
Ainsi, la famille \((A_j)\) est déterminée par l'histoire et par la jauge
initiale ; ce n'est pas une suite de choix indépendants.

Soit \(\mathcal F_{\rm cal}\subseteq\operatorname{GL}_6(\mathbb F_3)\)
l'ensemble des cadres accessibles à partir de \(f_0\). L'émission qui conserve la
jauge transportée est
\[
 \widetilde\varepsilon_n:X_n^{\rm cal}
 \longrightarrow(\mathbb F_3^6\times\mathcal F_{\rm cal})^n,
 \qquad
 \widetilde\varepsilon_n(x)
 =\bigl((w_0,f_0),\ldots,(w_{n-1},f_{n-1})\bigr).
\]
La dépendance par préfixes implique
\(\operatorname{pr}_{n,m}\widetilde\varepsilon_n
=\widetilde\varepsilon_m p_{n,m}\). Sa projection sur les mots est
\(\varepsilon_n\) ; l'oubli des cadres ne suffit pas à reconstruire la
lecture exceptionnelle.

\subsection{Naturalité et compatibilité à toute profondeur}

L'opération locale est covariante avec la convention des vecteurs-lignes. Pour
tout changement de coordonnées \(f\in\operatorname{GL}_6(\mathbb F_3)\),
\[
A^f=fAf^{-1}
\quad\Longrightarrow\quad
\Gamma_{A^f}(wf^{-1})
=\Gamma_A(w)(f^{-1}\oplus f^{-1}).
\]
Le changement de cadre transporte simultanément le mot, l'opérateur et
l'atlas. Il n'est pas exigé qu'il demeure dans le stabilisateur d'une unique copie
étiquetée.

L'ensemble des cadres accessibles \(\mathcal F_{\rm cal}\) étant déjà défini,
nous construisons l'atlas étalonné fixe
\[
 \mathscr C_W^{\rm cal}
 :=
 \left\{
 \bigl(f,\Gamma_{fA_Wf^{-1}}(w)\bigr):
 f\in\mathcal F_{\rm cal},\ w\in\mathbb F_3^6
 \right\}.
\]
L'étiquette \(f\) conserve la copie étalonnée dans laquelle vit chaque mot.
Nous définissons
\[
Q_n(x)=
\bigl(
\bigl(f_0,\Gamma_{A_0}(w_0)\bigr),\ldots,
\bigl(f_{n-1},\Gamma_{A_{n-1}}(w_{n-1})\bigr)
\bigr)\in\bigl(\mathscr C_W^{\rm cal}\bigr)^n.
\]
Soit
\[
 r_{n,m}:
 \bigl(\mathscr C_W^{\rm cal}\bigr)^n
 \longrightarrow
 \bigl(\mathscr C_W^{\rm cal}\bigr)^m
\]
la troncature des \(n-m\) derniers blocs.

\begin{theorem}[Compatibilité projective de l'incidence exceptionnelle]
\label{p12-83-c17-s2:thm:lector-excepcional-global-rev6}
\label{thm:lector-excepcional-global-rev6}
Pour tout \(m\leq n\),
\[
\boxed{r_{n,m}\circ Q_n=Q_m\circ p_{n,m}.}
\]
En conséquence, il existe une unique application
\[
\boxed{
Q_\infty:X_\infty^{\rm cal}
\longrightarrow
\bigl(\mathscr C_W^{\rm cal}\bigr)^{\mathbb N}}
\]
dont les projections finies sont les \(Q_n\). L'application est fidèle sur le
flux visible.
\end{theorem}

\begin{proof}
Écrivons
\[
\varepsilon_n(x)
=(w_0,\ldots,w_{m-1},w_m,\ldots,w_{n-1}).
\]
La compatibilité de \(\varepsilon\) montre que le flux visible de
\(p_{n,m}(x)\) est \((w_0,\ldots,w_{m-1})\). Pour \(j<m\), le cadre
\(f_j\) ne dépend que de la jauge initiale et des changements antérieurs à
\(j\), tous contenus dans ce préfixe. Les \(m\) premiers opérateurs
\(A_j\) sont donc les mêmes avant et après troncature. Appliquer
\(\Gamma_{A_j}\) bloc par bloc donne l'identité de l'énoncé. La propriété
universelle de la limite projective produit \(Q_\infty\). Enfin,
\(\operatorname{pr}_1\Gamma_{A_j}=I\) récupère chaque \(w_j\), ce qui prouve
la fidélité sur le flux visible.
\end{proof}

La dernière affirmation conserve son type. Deux histoires ayant le même flux
\((w_j)_j\) et des mémoires cachées distinctes peuvent avoir la même image
exceptionnelle. Le lecteur global préserve exactement ce que publie chaque bloc
visible et l'incidence ajoutée à ce bloc.

La compatibilité des émissions enrichies induit une unique application
\[
 \widetilde\varepsilon_\infty:X_\infty^{\rm cal}
 \longrightarrow
 (\mathbb F_3^6\times\mathcal F_{\rm cal})^{\mathbb N}.
\]
Soient \(\operatorname{pr}_w\) la projection sur le flux de mots et
\[
 \widetilde\Gamma_\infty
 \bigl((w_j,f_j)_{j\ge0}\bigr)
 :=\bigl(
   (f_j,\Gamma_{f_jA_Wf_j^{-1}}(w_j))
 \bigr)_{j\ge0}.
\]
Alors
\(Q_\infty=\widetilde\Gamma_\infty
\circ\widetilde\varepsilon_\infty\), avec domaines et codomaines déclarés.

\subsection{Bifurcation d'un préfixe survivant en valeur archimédienne et incidence exceptionnelle}

\(m\in\mathbb Z\) étant fixé, notons
\[
 \operatorname{flat}:(\mathbb F_3^6)^{\mathbb N}
 \longrightarrow\mathbb F_3^{\mathbb N}
\]
la concaténation des blocs dans l'ordre des coordonnées déclaré, et
\[
 \operatorname{ev}_3((d_k)_{k\ge1})
 :=\sum_{k\ge1}\frac{d_k}{3^k}\in[0,1],
 \qquad
 \tau_m(t):=m+t,
\]
l'évaluation ternaire et la translation par la partie entière. Le préfixe visible
\((d_1,\ldots,d_{6n})\) détermine le cylindre
\[
I_n(x)=
\left[
m+\sum_{k=1}^{6n}\frac{d_k}{3^k},
m+\sum_{k=1}^{6n}\frac{d_k}{3^k}+3^{-6n}
\right].
\]
Une prolongation restreint le cylindre et
\(\operatorname{diam}I_n=3^{-6n}\to0\). Par conséquent, toute histoire infinie
compatible détermine un unique point réel
\(\mathcal P_{\rm arq}(x)\).

\begin{theorem}[Double réalisation globale]
\label{p12-83-c17-s2:thm:doble-realizacion-global-rev6}
\label{thm:doble-realizacion-global-rev6}
Toute section compatible \(x=(x_n)_n\in X_\infty^{\rm cal}\) détermine,
à partir des mêmes préfixes, deux sorties compatibles :
\[
\mathcal P_{\rm arq}(x)\in\mathbb R,
\qquad
\mathcal P_{\rm exc}(x)=Q_\infty(x)
\in\bigl(\mathscr C_W^{\rm cal}\bigr)^{\mathbb N}.
\]
Tronquer avant de lire équivaut à tronquer chaque lecture après lecture. La
première sortie identifie les histoires par leur valeur limite ; la seconde conserve
chaque bloc visible et lui adjoint sa rotation de Paley--Witt.
\end{theorem}

\begin{proof}
L'emboîtement des cylindres et la convergence de leurs diamètres prouvent la
première affirmation. Le
\cref{p12-83-c17-s2:thm:lector-excepcional-global-rev6} prouve la seconde et la
compatibilité avec les troncatures. Toutes deux se factorisent par la même émission :
\[
\mathcal P_{\rm arq}
=\tau_m\circ\operatorname{ev}_3\circ\operatorname{flat}
 \circ\operatorname{pr}_w
 \circ\widetilde\varepsilon_\infty,
\qquad
\mathcal P_{\rm exc}
=\widetilde\Gamma_\infty\circ\widetilde\varepsilon_\infty.
\]
\end{proof}

La double projection est ainsi construite sur des histoires admissibles et une
jauge initiale transportée. Le continu est l'évaluation archimédienne de
la section ; la branche exceptionnelle est la réalisation compatible de son flux
visible. À partir de \(C_W\) s'ouvrent, avec les alignements déjà déclarés, les
branches vers \(W_{12}\), \(W_{24}\) et vers
\(N(A_2^{12})\), Leech, \(V^\natural\) et le Monstre. Ces branches ne transforment pas
les chiffres en un ensemble sur lequel le Monstre agirait directement :
elles transportent l'incidence que l'évaluation par valeur cesse de voir.

\subsection{Couplage postgénératif avec la clôture opératorielle d'Euler}
\label{p12-83-c17-s2:sec:acoplamiento-euler-doble-proyeccion-rev3}

Soit \(x_\pi\in X_\infty^{\rm cal}\) la section survivante du canal de
\(\pi\). La double réalisation du
\cref{p12-83-c17-s2:thm:doble-realizacion-global-rev6} produit
\[
 \mathcal P_{\rm arq}(x_\pi)\in\mathbb R,
 \qquad
 \mathcal P_{\rm exc}(x_\pi)=Q_\infty(x_\pi).
\]
La première branche publie l'angle engendré ; la seconde conserve le flux
visible et réalise dans chaque bloc la relation quadratique de Paley--Witt.

\begin{theorem}[Couplage postgénératif Euler--double projection]
\label{p12-83-c17-s2:thm:acoplamiento-euler-doble-proyeccion-rev3}
Supposons que la construction coinductive du canal de \(\pi\) détermine
\[
 \mathcal P_{\rm arq}(x_\pi)=\pi,
\]
et que la branche exceptionnelle utilise la jauge transportée avec
\(A_W^2=-I_6\). Soit \(J_{\mathrm e}\) la réalisation réelle de la source
entière du \cref{thm:dos-realizaciones-fuente-cuadratica-rev3}. Alors
\begin{equation}
 \boxed{
 \operatorname{Exp}\!\bigl(
   \mathcal P_{\rm arq}(x_\pi)J_{\mathrm e}
 \bigr)+I=0.}
 \label{p12-83-c17-s2:eq:acoplamiento-euler-doble-proyeccion-rev3}
\end{equation}
\end{theorem}

\begin{proof}
D'après le \cref{thm:dos-realizaciones-fuente-cuadratica-rev3}, \(A_W\) et
\(J_{\mathrm e}\) réalisent la relation \(u^2+1=0\) dans des catégories distinctes.
Dans la branche réelle,
\[
 \operatorname{Exp}(\theta J_{\mathrm e})
 =\cos\theta\,I+\sin\theta\,J_{\mathrm e}.
\]
La substitution de la sortie précédemment engendrée
\(\theta=\mathcal P_{\rm arq}(x_\pi)=\pi\) produit \(-I\), et
l'addition de \(I\) donne l'identité annoncée.
\end{proof}

\paragraph{Ordre causal et absence de circularité.}
La dépendance démontrée est
\[
 \text{émetteur}
 \longrightarrow x_\pi
 \longrightarrow\mathcal P_{\rm arq}(x_\pi)
 \longrightarrow\pi
 \longrightarrow
 \operatorname{Exp}({\pi}J_{\mathrm e})+I.
\]
La branche exceptionnelle fixe le type quadratique ; elle ne sélectionne ni l'angle ni ses
chiffres. L'identité d'Euler reconnaît ensuite la compatibilité des deux
réalisations et n'intervient pas comme entrée du générateur de
\(\pi\).


\section{Globalisation naturelle des projections archimédienne et exceptionnelle sur un système projectif commun}
\label{p12-83-c17-s3:chap:globalizacion-doble-proyeccion}

La thèse de la double projection acquiert un contenu mathématique lorsque les deux
lectures sont définies sur un même système d'états et commutent avec la
restriction de profondeur. L'évaluation archimédienne contracte une famille de
cylindres en une valeur ; l'évaluation exceptionnelle conserve des relations
d'incidence que la première peut oublier. Il ne s'agit pas de deux expressions
numériques juxtaposées, mais de deux transformations naturelles de domaine
commun.

\subsection{Systèmes compatibles à profondeur finie}

Soient \((X_n,p_{n+1,n})\), \((A_n,a_{n+1,n})\) et \((E_n,e_{n+1,n})\) trois systèmes projectifs. Le premier contient des trajectoires TPK ; le deuxième, des cylindres ou des approximants archimédiens ; le troisième, des données d’incidence exceptionnelle. Supposons données des applications
\[
 R_n:X_n\longrightarrow A_n,
 \qquad
 Q_n:X_n\times C_n\longrightarrow E_n,
\]
où \(C_n\) réunit les données de coordonnées du lecteur exceptionnel. On exige que les jauges forment également un système compatible et que, pour tout \(n\), les diagrammes commutent
\[
 a_{n+1,n}\circ R_{n+1}=R_n\circ p_{n+1,n},
\]
\[
 e_{n+1,n}\circ Q_{n+1}
 =Q_n\circ(p_{n+1,n}\times c_{n+1,n}).
\]

\begin{theorem}[Globalisation des deux évaluations]
\label{p12-83-c17-s3:thm:globalizacion-evaluaciones}
Sous les compatibilités précédentes, il existe des applications canoniques
\[
 R_\infty:\varprojlim X_n\longrightarrow\varprojlim A_n,
 \qquad
 Q_\infty:\varprojlim X_n\times\varprojlim C_n
            \longrightarrow\varprojlim E_n.
\]
Leur application diagonale
\[
 (R_\infty,Q_\infty):X_\infty\times C_\infty
 \longrightarrow A_\infty\times E_\infty
\]
est unique parmi les applications dont les projections coordonnées sont
\(R_n\) et \(Q_n\).
\end{theorem}

\begin{proof}
Soit \(x=(x_n)_n\in\varprojlim X_n\). La première équation de commutativité
implique
\[
 a_{n+1,n}(R_{n+1}(x_{n+1}))=R_n(x_n),
\]
de sorte que \((R_n(x_n))_n\) appartient à \(\varprojlim A_n\). Nous définissons
\(R_\infty(x)\) au moyen de cette famille. Le même argument, appliqué à
\((x_n,c_n)_n\), définit \(Q_\infty\). L'unicité découle du fait qu'un élément
d'une limite projective est déterminé par toutes ses coordonnées.
\end{proof}

Le théorème identifie la preuve concrète qu’exige une réalisation HMT : il ne suffit pas de disposer d’un résultat archimédien à une profondeur et d’un code à une autre ; tous deux doivent provenir d’états compatibles et respecter les applications de troncature.

\subsection{Contraction archimédienne de cylindres compatibles et réalisation de la droite réelle}

À chaque \(x_n\in X_n\) est associé un intervalle compact non vide
\(I_n(x_n)\subset\mathbb R\). Si
\[
 p_{n+1,n}(x_{n+1})=x_n
 \quad\Longrightarrow\quad
 I_{n+1}(x_{n+1})\subseteq I_n(x_n)
\]
et s'il existe \(\varepsilon_n\to0\) avec
\(\operatorname{diam}I_n(x_n)\leq\varepsilon_n\), l'intersection
\[
 \bigcap_{n\geq0}I_n(x_n)
\]
contient un unique réel. On obtient ainsi
\[
 \operatorname{ev}_{\rm arq}:X_\infty^{\rm arq}\longrightarrow\mathbb R.
\]
La phrase « les cylindres tendent vers zéro » signifie ici que leur diamètre
converge vers zéro, non que le nombre évalué tende vers zéro.

La relation
\[
 x\sim_{\rm arq}y
 \quad\Longleftrightarrow\quad
 \operatorname{ev}_{\rm arq}(x)=\operatorname{ev}_{\rm arq}(y)
\]
décrit exactement l'information éliminée par le lecteur. La valeur réelle
est une classe d'évaluation ; l'histoire conserve en outre orientation,
retenue, feuillet et signature de frontière.

\subsection{Projection exceptionnelle par conservation de l'incidence et de la frontière}

La seconde évaluation ne prend pas ses valeurs dans \(\mathbb R\). Son codomaine est une
catégorie de configurations d'incidence. Dans la réalisation finie en vigueur,
la donnée de coordonnées contient :
\[
 \begin{aligned}
 C_{\rm exc}=(\,&\text{ordre projectif},\text{calibre de Paley},
 \text{dodécade},\\
 &\text{alignement},\text{origine réticulaire},\text{chambre}).
 \end{aligned}
\]
Après fixation de cette donnée, la chaîne contient deux descendances qui doivent
demeurer séparées :
\[
 C_W\longrightarrow W_{12}\longrightarrow W_{24},
\]
\[
 C_W\longrightarrow N(A_2^{12})
 \longrightarrow\Lambda_{24}
 \longrightarrow V^\natural
 \longrightarrow\mathbb M.
\]
La première transporte des plans ; la seconde conduit, au moyen de résultats
classiques ultérieurs, au réseau de Leech, à l'algèbre d'opérateurs de
sommets et au groupe Monstre. L'holonomie centrale et le corridor
Schläfli--\(E_6\) forment une extraction parallèle de frontière ; ils ne s'obtiennent pas
en remplaçant l'une des deux descendances précédentes.

\subsection{Compatibilité finie du domaine commun}

Le sceau dodécaphasique
\[
 K=(234,543,140,729,659,824,621,58,914,794,146,601)
\]
détermine le support
\[
 B_K=\{m:K_m\geq729\}=\{4,6,9,10\}.
\]
Dans la carte de Witt publiée,
\[
 P_\pi=\{2,4,5,6,7,8,9,10,11\},
 \qquad
 H_e=\{1,3,4,6,9,10\},
\]
et, par évaluation indépendante des deux supports,
\[
 \boxed{B_K=H_e\cap P_\pi.}
\]
En outre, pour le support de retenue négative
\[
 H_\alpha^-=\{4,5,6,7,9,10\}
\]
on a
\[
 B_K=H_e\cap H_\alpha^-,
 \qquad H_\alpha^-=B_K\sqcup\{5,7\}.
\]
Ces identités constituent un témoin de compatibilité finie : le même
sous-ensemble est retrouvé à partir du sceau arithmétique et de l'incidence de
Witt. Ce n'est pas une égalité de cardinalités, puisqu'elle identifie des positions
concrètes.

\subsection{Séparation typée des publications archimédienne et exceptionnelle de l'état commun}

\begin{theorem}[Critère de fidélité conjointe]
Soient \(X\) compact et \(Y,Z\) des espaces séparés. Si deux applications
continues \(R:X\to Y\) et \(Q:X\to Z\) satisfont
\[
 x\neq y\Longrightarrow R(x)\neq R(y)\ \text{ou}\ Q(x)\neq Q(y),
\]
alors \((R,Q):X\to Y\times Z\) est un plongement topologique.
\end{theorem}

\begin{proof}
La condition est l'injectivité de l'application diagonale. Toute application
continue et injective d'un compact dans un espace séparé est un
homéomorphisme sur son image.
\end{proof}

À chaque profondeur finie, la condition est décidable par énumération des couples \(x\neq y\) qui demeurent simultanément dans une fibre de \(R_n\) et de \(Q_n\). La globalisation exige la compatibilité de ces contrôles lors du passage de \(n+1\) à \(n\). Le témoin \(B_K\) démontre que les deux lectures partagent déjà un invariant non trivial ; la fidélité conjointe exige le contrôle de toutes les fibres et ne peut être remplacée par ce seul témoin.

\subsection{Formulation précise de la thèse sur la droite réelle}

L'objet enrichi associé à la lecture archimédienne est la flèche
\[
 \mathbb R_{\mathfrak H}:=
 \bigl(X_\infty^{\rm arq}
 \xrightarrow{\operatorname{ev}_{\rm arq}}
 \operatorname{Im}(\operatorname{ev}_{\rm arq})\subseteq\mathbb R\bigr).
\]
Le domaine conserve les généalogies ; l'image conserve les valeurs. La couverture des
cylindres à toute échelle et la prolongeabilité compatible sont démontrées dans
la construction du lecteur archimédien
\eqref{eq:continuo-lector-arquimediano} : pour tout compact
\(K\subset\mathbb R\),
\(\operatorname{ev}_{\mathbb R}(\Omega_{\infty,K})=K\), et la réunion dirigée
de ces intervalles couvre \(\mathbb R\). Par conséquent,
\[
 \operatorname{Im}(\operatorname{ev}_{\rm arq})=\mathbb R,
 \qquad
 X_\infty^{\rm arq}/\!\sim_{\mathbb R}\ \cong\mathbb R.
\]
Le corps réel n'agit pas comme point de départ de la construction, mais comme
quotient archimédien d'un espace de sections avec mémoire. La descente de la
somme et du produit est vérifiée opération par opération et ne conditionne pas cette
surjectivité déjà démontrée.

La double lecture s'exprime alors par
\[
 X_\infty^{\rm com}\times C_\infty
 \xrightarrow{\ (\operatorname{ev}_{\rm arq},Q_\infty)\ }
 \mathbb R\times E_\infty.
\]
Cette formule est la traduction mathématique de l'affirmation conceptuelle : une
même histoire peut se contracter jusqu'à une valeur et, dans la lecture transversale
de sa frontière, conserver l'incidence qui conduit aux symétries
exceptionnelles.


% Fuente normativa activa de la obstrucción de factorización.
\section{Dépendance du lecteur d’incidence à l’égard de l’état holonomique intégral}
\label{p12-83-c17-s4:sec:obstruccion-lector-excepcional-minimo}

La compatibilité finie des deux évaluations n'implique pas que n'importe quel
lecteur d'incidence raffine la partition produite par le lecteur visible. Le
catalogue régional permet de trancher cette question exactement pour la
réalisation minimale de Paley--Witt. Soit \(\mathcal U_{468}\) l'ensemble
des 468 états régionaux publiés et soit
\[
 \Psi:\mathcal U_{468}\longrightarrow\mathcal W_{243}\subset\mathbb F_3^6,
 \qquad
 \Psi(U_1,\ldots,U_6)=(-U_1,\ldots,-U_6)\pmod 3.
\]
Pour \(w\in\mathbb F_3^6\), l’application
\[
 c(w)=(w,wA_W)\in C_W\subset\mathbb F_3^{12}
\]
est le codeur ternaire fixé dans la section de Paley--Witt. Nous définissons
le lecteur d'incidence minimal par
\[
 Q_{\min}=\operatorname{Supp}\circ c\circ\Psi:
 \mathcal U_{468}\longrightarrow\mathcal P([12]).
\]

Pour une application \(f:X\to Y\), notons
\[
 \operatorname{Eq}(f)=\{(x,x')\in X^2:f(x)=f(x')\}
\]
sa relation d'équivalence par fibres.

\begin{theorem}[Obstruction du lecteur exceptionnel minimal]
\label{p12-83-c17-s4:thm:obstruccion-lector-excepcional-minimo}
Dans le catalogue complet de 468 états, on vérifie
\[
 \operatorname{Eq}(\Psi,Q_{\min})=\operatorname{Eq}(\Psi).
\]
Dans la notation abrégée habituelle des partitions par fibres,
\[
 \ker(\Psi,Q_{\min})=\ker\Psi.
\]
Par conséquent, \((\Psi,Q_{\min})\) n'est pas injective et le lecteur
d'incidence minimal ne distingue aucun couple d'états déjà
identifié par \(\Psi\).
\end{theorem}

\begin{proof}
La factorisation \(Q_{\min}=S\circ\Psi\), avec
\(S=\operatorname{Supp}\circ c\), implique
\[
 \Psi(u)=\Psi(v)\Longrightarrow Q_{\min}(u)=Q_{\min}(v).
\]
L'implication réciproque entre les relations d'équivalence est immédiate,
puisque \(\Psi\) est la première composante de l'application diagonale.
L'égalité des relations s'obtient donc sans hypothèses statistiques.
L'énumération exhaustive confirme en outre que les deux partitions possèdent le
même histogramme :
\[
 132\,[1]+66\,[2]+18\,[3]+3\,[4]+6\,[5]+18\,[6],
\]
où \(a\,[m]\) signifie qu'il existe \(a\) fibres de cardinal \(m\).
\end{proof}

La microfibre de \(w_\pi=010211\) fournit le témoin distingué. Ses
cinq régions sont
\[
\begin{split}
 &(501,614,498,169,272,272),\quad
 (810,923,870,169,272,272),\\
 &(810,983,810,169,272,272),\quad
 (870,923,810,169,272,272),\\
 &(870,923,810,418,674,521).
\end{split}
\]
Les cinq ont la même image visible et le même support exceptionnel
\[
 P_\pi=\{2,4,5,6,7,8,9,10,11\}.
\]
Au niveau des 243 mots, on trouve 213 supports distincts : 185 ont
un antécédent, 27 en ont deux et un en a quatre. La prise de support introduit
donc un second quotient en plus de celui imposé par \(\Psi\).

Le théorème délimite le lecteur minimal, non l'évaluation exceptionnelle
enrichie. Un lecteur capable de séparer les régions d'un même mot
doit agir avant le quotient visible et utiliser au moins l'une des données
que celui-ci élimine : orientation, feuillet, signature, retour ou état de frontière.
L'entrelaceur dodécaphasique développé dans la partie suivante appartient à cette
classe enrichie et ne contredit pas l'obstruction précédente.

L'énumération complète des fibres produit l'histogramme précédent et
démontre que la prise de support introduit un second quotient.


\section{Publication archimédienne et d'incidence}

Pour chaque niveau, soit \(\mathcal E_n\) un objet d'incidence et soient
\begin{equation}
 e_n:S_n\longrightarrow\mathcal E_n,
 \qquad
 \sigma_{n,m}:\mathcal E_n\longrightarrow\mathcal E_m
 \label{p12-83-c17-s5:eq:datos-lector-incidencial-u024}
\end{equation}
des applications qui satisfont
\begin{equation}
 \sigma_{n,m}\circ e_n=e_m\circ r_{n,m}.
 \label{p12-83-c17-s5:eq:naturalidad-lector-incidencial-u024}
\end{equation}
La propriété universelle construit
\begin{equation}
 \operatorname{ev}_{\rm inc}:\Omega_\infty
 \longrightarrow\mathcal E_\infty,
 \qquad
 \mathcal E_\infty:=\varprojlim_n\mathcal E_n.
 \label{p12-83-c17-s5:eq:evaluacion-incidencial-limite-u024}
\end{equation}
La double publication est
\begin{equation}
 \mathcal P_\infty
 :=(\operatorname{ev}_{\mathbb R},\operatorname{ev}_{\rm inc}):
 \Omega_\infty\longrightarrow\mathbb R\times\mathcal E_\infty.
 \label{p12-83-c17-s5:eq:doble-publicacion-continuo-u024}
\end{equation}

\begin{figure}[H]
\centering
\begin{tikzcd}[column sep=10em,row sep=5em,
  cells={nodes={font=\large}}]
&\Omega_\infty
 \arrow[dl,"\operatorname{ev}_{\mathbb R}"']
 \arrow[dr,"\operatorname{ev}_{\rm inc}"]&\\
\mathbb R
 \arrow[dr,"\operatorname{coord}"']&&
\mathcal E_\infty
 \arrow[dl,"\operatorname{pub}_{\rm inc}"]\\
&\mathcal O&
\end{tikzcd}
\caption{Deux publications d'une même histoire. La commutativité dans
\(\mathcal O\) n'identifie pas les codomaines \(\mathbb R\) et
\(\mathcal E_\infty\).}
\label{p12-83-c17-s5:fig:doble-publicacion-continuo-u024}
\end{figure}

Si l’on démontre
\begin{equation}
 \operatorname{pub}_{\rm inc}\circ\operatorname{ev}_{\rm inc}
 =\operatorname{coord}\circ\operatorname{ev}_{\mathbb R},
 \label{p12-83-c17-s5:eq:compatibilidad-doble-publicacion-u024}
\end{equation}
on obtient une comparaison typée des deux expressions. L’identité n’affirme pas qu’une branche soit construite à partir de l’autre. Toutes deux partent de \(\Omega_\infty\).


\section{Jauge de la publication exceptionnelle}

La seconde projection n'est pas un lecteur unaire dépourvu de données
d'incidence. Soit
\[
 \mathscr C_{\rm exc}
 =\mathscr C_{\rm ord}\times\mathscr C_{\rm Paley}
  \times\mathscr C_{\rm dod}\times\mathscr C_{\rm inc}
  \times\mathscr C_{\rm ret}
\]
l'espace des jauges exceptionnelles, dont les coordonnées fixent
respectivement l'ordre projectif, la jauge de Paley, la dodécade,
l'alignement d'incidence et la chambre de réseau. On fixe
\[
 \chi_{\rm exc}
 =(\chi_{\rm ord},\chi_{\rm Paley},\chi_{\rm dod},
   \chi_{\rm inc},\chi_{\rm ret})\in\mathscr C_{\rm exc}.
\]
Le lecteur typé a pour domaine
\[
 \mathfrak E:
 X_\infty^{\rm cal}\times\mathscr C_{\rm exc}
 \longrightarrow\mathbf{Exc},
 \qquad
 \mathfrak E_{\chi_{\rm exc}}(u):=\mathfrak E(u,\chi_{\rm exc}).
\]
Désormais, toute abréviation \(\mathfrak E(u)\) signifie
\(\mathfrak E_{\chi_{\rm exc}}(u)\) pour cette jauge explicitement fixée.
La dépendance à l'égard de \(\chi_{\rm exc}\) n'affecte pas l'existence du
continu ni sa publication archimédienne : elle type la conservation exceptionnelle
de l'incidence une fois l'état commun construit.

\begin{falsifier}[Perte de la jauge exceptionnelle]
La comparaison demeure incomplète si l'on change l'ordre projectif, la dodécade,
l'alignement d'incidence ou la chambre de réseau sans transporter les autres
coordonnées de la jauge, ou si l'on utilise \(\mathfrak E\) comme lecteur unaire avant
de fixer \(\chi_{\rm exc}\). Une telle défaillance invalide cette comparaison exceptionnelle ;
elle ne modifie pas le quotient archimédien et ne régénère pas le continu.
\end{falsifier}


La seconde lecture ne naît pas du nombre réel. Soit
\[
 \mathfrak E:
 \mathscr O_{\mathbb R,\leftrightarrow}^{\rm HMT}
 \longrightarrow\mathbf{Exc}
\]
le lecteur continu d'incidence, à codomaine séparé, qui prolonge
la chaîne de Paley--Witt, Golay,
Leech, \(V^\natural\), Monstre et Moonshine. Les deux applications partent de
la même histoire :
\begin{equation}
 \mathfrak D(u)=(\mathfrak R(u),\mathfrak E(u)).
 \label{p12-83-c17-s7:hmtch:eq-doble-proyeccion-two-way}
\end{equation}
La première contracte les fibres compatibles jusqu'à leur moyenne archimédienne ; la
seconde conserve incidence, orientation et mémoire de bord.

L'origine du calendrier dodécaphasique étant fixée, soit
\[
 R_+:=\{k\in\mathbb Z:1+(k\bmod 9)\in\{1,4,7\}\}
\]
l'ensemble bi-infini des retours complets \((\tau=+1)\). Il est syndétique :
ses lacunes sont uniformément bornées. Pour
\(u=(m,(x_k)_{k\in\mathbb Z})\), nous définissons
\[
 \operatorname{Ret}_+(u)
 :=\bigl(m,(x_k)_{k\in R_+}\bigr)
 \in\mathbb Z\times\prod_{k\in R_+}X,
\]
mais chaque \(x_k\) est ici l'état entier
\eqref{p12-83-c19-s1:hmtch:eq-celula-trit-completa} avec feuillet, retenue, route et incidence,
non le signe isolé.

\begin{proposition}[Échantillonnage sans perte du registre généalogique de retour]
\label{p12-83-c17-s7:hmtch:prop-retorno-reconstruye}
L'application \(\operatorname{Ret}_+\), munie de la topologie induite dans le
produit, est un homéomorphisme sur son image.
En conséquence, il existe une unique application continue
\(\mathfrak E_+\) sur cette image telle que
\begin{equation}
 \mathfrak E=\mathfrak E_+\circ\operatorname{Ret}_+.
 \label{p12-83-c17-s7:hmtch:eq-factorizacion-excepcional-retorno}
\end{equation}
\end{proposition}

\begin{proof}
Entre deux retours consécutifs, il y a un nombre uniformément borné de pas
du calendrier. L'état complet lors d'un retour contient la transition et
le registre généalogique qui déterminent les pas intermédiaires ; dans l'autre direction, on
utilise la flèche inverse de la complétion en groupoïde. Toute
l'orbite est ainsi reconstruite. L'application est continue et, lorsqu'on la restreint à une carte
compacte, une injection continue dans un espace séparé est un
homéomorphisme sur son image. La coordonnée entière est conservée
explicitement. La factorisation est définie en transportant \(\mathfrak E\) par
l'inverse et est continue carte par carte.
\end{proof}

C'est pourquoi le \(+1\) participe au corridor exceptionnel comme mémoire complète
de retour : l'échantillonnage d'états entiers ne perd pas la sortie exceptionnelle.
Cette proposition démontre la récupération et la factorisation, non une génération
du corridor à partir du scalaire \(+1\) isolé. La sortie archimédienne et la
sortie exceptionnelle sont deux résultantes coordonnées du même antécédent.
 
\chapter{Reconstruction archimédienne par cylindres semi-fermés avec mémoire}
\label{chap:83-18}

La limite projective \(\Omega_\infty\) conserve la trajectoire complète. Pour obtenir une coordonnée réelle, il faut ajouter un système compatible d’intervalles et démontrer trois faits distincts : chaque trajectoire détermine une valeur unique ; toute valeur du compact déclaré est atteinte ; et les mesures finies passent au quotient de manière cohérente. Aucun de ces faits ne découle de la simple croissance du nombre d’états.

La publication archimédienne conserve également une approximation de
l'identité par des espérances conditionnelles cylindriques. La concentration de masse
unitaire sur des préfixes compatibles fournit le passage exact des
cylindres nonadiques à la mesure de Dirac, sans identifier densité,
mesure du support et masse totale.

\begin{HMTScientificOwnerSubordinated}
\chapter{Cylindres nonadiques, identité et noyau de Dirac}

\section{Système inverse de cylindres nonadiques compatibles avec mémoire généalogique}

Soit
\[
X=(\Znine)^{\mathbb N}
\]
avec la topologie produit et la mesure uniforme \(\mu\). Pour
\(x=(x_1,x_2,\ldots)\), le cylindre de profondeur \(n\) est
\[
C_n(x)=\{y\in X:y_j=x_j\text{ pour }1\le j\le n\}.
\]
Les cylindres sont simultanément ouverts et fermés, forment une base et
satisfont
\[
\mu(C_n(x))=9^{-n},\qquad
C_{n+1}(x)\subset C_n(x),\qquad
\bigcap_n C_n(x)=\{x\}.
\]
Ajouter un chiffre ne crée pas un nouveau nombre : cela restreint le cylindre des
prolongements admissibles. Cette assertion donne un contenu topologique à la thèse
généalogique du nombre.

\section{Concentration cylindrique de l'identité et construction du noyau de Dirac}

\begin{theorem}[Concentration unitaire nonadique]
Définissons
\[
f_{n,x}=9^n\mathbf1_{C_n(x)}.
\]
Alors
\[
\mu(C_n(x))\to0,\qquad f_{n,x}(x)\to\infty,\qquad
\int_Xf_{n,x}\,\dd\mu=1,
\]
et les mesures \(f_{n,x}\mu\) convergent faiblement-* vers \(\delta_x\).
\end{theorem}

\begin{proof}
La première identité est immédiate. La hauteur de \(f_{n,x}\) sur son support
est \(9^n\), de sorte qu'elle diverge. L'intégrale vaut
\(9^n9^{-n}=1\). Pour \(g\in C(X)\), l'oscillation de \(g\) dans \(C_n(x)\)
tend vers zéro par continuité uniforme ; par conséquent
\[
\left|\int g f_{n,x}\,\dd\mu-g(x)\right|
\le \sup_{y\in C_n(x)}|g(y)-g(x)|\longrightarrow0.
\]
C'est la convergence faible-* vers \(\delta_x\).
\end{proof}

Le théorème n'affirme pas \(1/0=\infty\). Il exhibe trois évaluations d'un même
objet : mesure du support, densité et masse totale. Par blocs de trois chiffres,
\[
\mu(C_{3m})=729^{-m},\qquad
f_{3m}=729^m\mathbf1_{C_{3m}},\qquad
729^{-m}729^m=1.
\]
La conservation évolutive de l'unité prend ici une forme analytique exacte.

\section{Convergence du noyau de Kronecker vers la distribution de Dirac}

Au niveau fini \(X_n=(\Znine)^n\), soit
\[
K_n(x,y)=9^n\mathbf1_{\{x=y\}}.
\]
Avec la mesure uniforme, l'opérateur
\[
(E_nf)(x)=\sum_{y\in X_n}K_n(x,y)f(y)9^{-n}
\]
est l'identité. Dans la limite cylindrique, la condition \(x=y\) est remplacée par
« même préfixe de longueur \(n\) » et \(E_n\) est l'espérance conditionnelle sur
l'algèbre cylindrique. On a
\[
E_n^2=E_n,\qquad E_n^*=E_n,\qquad E_nf\to f
\quad\text{dans }L^2(X,\mu).
\]

Le pont exact est donc
\[
\delta_{ij}
\longrightarrow K_n(x,y)
\longrightarrow\delta_x.
\]
Le delta de Kronecker est un noyau matriciel ; le delta de Dirac est une mesure ou
une distribution. Aucun n'est l'opérateur de Dirac
\(D=\gamma^\mu\partial_\mu\). Le passage aux spineurs et à Clifford exige
un second opérateur d'entrelacement et sera traité ultérieurement.

\section{Loi d'échelle, dimension effective et densité de la mesure cylindrique}

Supposons qu'une contraction de rapport \(\lambda\) agisse de manière indépendante
dans \(d\) directions et que la mesure produit varie comme \(\lambda^d\). Pour
conserver une masse égale à un, la densité doit varier comme
\[
\rho\longmapsto\lambda^{-d}\rho.
\]
Les exposants \((-1,-2,-3)\) apparaissent alors comme des poids de densité en ligne,
en aire et en volume. Cette proposition est générale et exacte sous l'hypothèse de
produit. Elle ne permet pas d'identifier automatiquement \(\lambda\) à
\(\alpha\) ; cette identification nécessite une carte métrique.

\section{Monodromie nonadique et calculabilité de l'évaluation archimédienne}

L'existence d'une branche infinie dans un arbre à branchement fini peut
découler du lemme de König, mais König ne produit pas de sélecteur calculable. Il existe
des arbres binaires récursifs infinis sans branche récursive. La connexion nonadique
conjointe ne repose pas sur cette existence abstraite : elle possède un sélecteur uniforme
certifié qui engendre à profondeur finie arbitraire et n'ouvre les témoins
décimaux qu'après la génération.

\begin{falsadorBox}[title={Trois confusions interdites}]
On n'écrira pas \(1/0\) ; on ne représentera pas le delta de Dirac comme une fonction
ordinaire avec « un un infini » ; on n'utilisera pas König comme algorithme. Le résultat
fort est la compatibilité exacte entre cylindres, densité unitaire et limite
faible, ainsi que le sélecteur constructif indépendant de la comparaison décimale.
\end{falsadorBox}
 \end{HMTScientificOwnerSubordinated}

\section{Système d'intervalles associé aux états finis}

Soit \((S_n,r_{n,m})\) le système projectif de la
\cref{chap:sistema-inverso-continuidad-genealogica}. À chaque
\(a\in S_n\), nous attribuons un intervalle compact non vide
\begin{equation}
 I_n(a)=[\ell_n(a),u_n(a)]\subset\mathbb R.
 \label{p12-83-c18-s1:eq:intervalo-asociado-estado}
\end{equation}
L'affectation est \emph{compatible avec le raffinement} lorsque
\begin{equation}
 r_{n+1,n}(b)=a
 \quad\Longrightarrow\quad
 I_{n+1}(b)\subseteq I_n(a).
 \label{p12-83-c18-s1:eq:anidamiento-intervalos-estados}
\end{equation}
Elle est \emph{uniformément contractante} lorsqu'il existe une suite
\(\delta_n\downarrow0\) telle que
\begin{equation}
 \operatorname{diam}I_n(a)\leq\delta_n
 \qquad(a\in S_n).
 \label{p12-83-c18-s1:eq:contraccion-uniforme-intervalos}
\end{equation}

Pour une histoire \(x=(x_n)_n\), les inclusions
\eqref{p12-83-c18-s1:eq:anidamiento-intervalos-estados} produisent une chaîne
\begin{equation}
 I_0(x_0)\supseteq I_1(x_1)\supseteq I_2(x_2)\supseteq\cdots.
 \label{p12-83-c18-s1:eq:cadena-intervalos-historia}
\end{equation}

\begin{theorem}[Évaluation d'une histoire]
\label{p12-83-c18-s1:thm:evaluacion-unica-historia-u024}
Si l'affectation est compatible et uniformément contractante, alors pour
chaque \(x\in\Omega_\infty\) il existe un unique nombre réel
\begin{equation}
 \operatorname{ev}_{\mathbb R}(x)
 \in\bigcap_{n\geq0}I_n(x_n).
 \label{p12-83-c18-s1:eq:evaluacion-real-historia-u024}
\end{equation}
L'application
\begin{equation}
 \operatorname{ev}_{\mathbb R}:\Omega_\infty\longrightarrow\mathbb R
 \label{p12-83-c18-s1:eq:aplicacion-evaluacion-real-u024}
\end{equation}
est continue.
\end{theorem}

\begin{proof}
Les intervalles de \eqref{p12-83-c18-s1:eq:cadena-intervalos-historia} sont compacts,
non vides et emboîtés. Le théorème des intervalles emboîtés garantit que
leur intersection est non vide. Si \(s,t\) lui appartiennent, alors
\[
 |s-t|\leq\operatorname{diam}I_n(x_n)\leq\delta_n
\]
pour tout \(n\). En faisant \(n\to\infty\), on obtient \(s=t\).

Pour la continuité, soit \(\varepsilon>0\) et choisissons \(n\) avec
\(\delta_n<\varepsilon\). Si \(y\in[x_n]_n\), alors
\(I_n(y_n)=I_n(x_n)\). Les deux évaluations appartiennent à cet intervalle et,
par conséquent,
\[
 |\operatorname{ev}_{\mathbb R}(x)
  -\operatorname{ev}_{\mathbb R}(y)|<\varepsilon.
\]
Le cylindre \([x_n]_n\) est un voisinage ouvert de \(x\).
\end{proof}

La construction n'identifie pas \(x\) à
\(\operatorname{ev}_{\mathbb R}(x)\). Une même coordonnée peut avoir
plusieurs histoires, et une histoire contient davantage de données que n'importe quel point de
la droite.

\section{Couverture cohérente du compact visible}

Soit \(K\subset\mathbb R\) un compact non vide. Outre l’emboîtement et la contraction, nous exigeons
\begin{align}
 K&=\bigcup_{a\in S_0}I_0(a),
 \label{p12-83-c18-s1:eq:cobertura-inicial-compacto-u024}\\
 I_n(a)&=
 \bigcup_{\substack{b\in S_{n+1}\\r_{n+1,n}(b)=a}}
 I_{n+1}(b)
 \qquad(a\in S_n).
 \label{p12-83-c18-s1:eq:cobertura-extensiones-compatible-u024}
\end{align}
La seconde égalité est une condition géométrique supplémentaire. La surjectivité de \(r_{n+1,n}\) garantit seulement l’existence d’états au-dessus de \(a\) ; elle ne garantit pas que leurs intervalles recouvrent tout \(I_n(a)\).

\begin{theorem}[Surjectivité de l'évaluation]
\label{p12-83-c18-s1:thm:sobreyectividad-evaluacion-compacto}
Si \eqref{p12-83-c18-s1:eq:cobertura-inicial-compacto-u024} et \eqref{p12-83-c18-s1:eq:cobertura-extensiones-compatible-u024} sont satisfaites, alors
\begin{equation}
 \operatorname{ev}_{\mathbb R}(\Omega_\infty)=K.
 \label{p12-83-c18-s1:eq:imagen-evaluacion-compacto}
\end{equation}
\end{theorem}

\begin{proof}
L'inclusion de gauche à droite découle de
\(I_n(x_n)\subseteq I_0(x_0)\subseteq K\).

Fixons \(t\in K\). Définissons
\begin{equation}
 T_n(t):=\{a\in S_n:t\in I_n(a)\}.
 \label{p12-83-c18-s1:eq:estados-que-cubren-punto}
\end{equation}
La couverture initiale rend \(T_0(t)\) non vide. Si \(a\in T_n(t)\),
l'identité \eqref{p12-83-c18-s1:eq:cobertura-extensiones-compatible-u024} fournit au
moins un \(b\in T_{n+1}(t)\) avec \(r_{n+1,n}(b)=a\). Le graphe de ces
extensions est à ramification finie et possède des niveaux non vides à toute
profondeur. Il existe donc une chaîne compatible \(x=(x_n)_n\) avec
\(t\in I_n(x_n)\) pour tout \(n\). L'unicité de l'intersection impose
\(\operatorname{ev}_{\mathbb R}(x)=t\).
\end{proof}

\section{Quotient généalogique archimédien et homéomorphisme avec \texorpdfstring{\(\mathbb R\)}{R}}

Nous définissons la relation
\begin{equation}
 x\sim_{\mathbb R}y
 \quad\Longleftrightarrow\quad
 \operatorname{ev}_{\mathbb R}(x)
 =\operatorname{ev}_{\mathbb R}(y).
 \label{p12-83-c18-s1:eq:relacion-equivalencia-arquimediana}
\end{equation}
Soit
\begin{equation}
 q_{\mathbb R}:\Omega_\infty
 \longrightarrow\Omega_\infty/\!\sim_{\mathbb R}
 \label{p12-83-c18-s1:eq:proyeccion-cociente-arquimediano}
\end{equation}
la projection canonique.

\begin{theorem}[Reconstruction du compact comme quotient]
\label{p12-83-c18-s1:thm:homeomorfismo-cociente-compacto-u024}
Sous les hypothèses précédentes, il existe un unique homéomorphisme
\begin{equation}
 \overline{\operatorname{ev}}_{\mathbb R}:
 \Omega_\infty/\!\sim_{\mathbb R}
 \xrightarrow{\ \cong\ }K
 \label{p12-83-c18-s1:eq:homeomorfismo-cociente-compacto-u024}
\end{equation}
tel que
\begin{equation}
 \operatorname{ev}_{\mathbb R}
 =\overline{\operatorname{ev}}_{\mathbb R}\circ q_{\mathbb R}.
 \label{p12-83-c18-s1:eq:factorizacion-evaluacion-cociente}
\end{equation}
La fibre au-dessus de \(t\in K\) est
\begin{equation}
 \mathfrak G_t
 :=\operatorname{ev}_{\mathbb R}^{-1}(t),
 \label{p12-83-c18-s1:eq:fibra-genealogias-mismo-valor}
\end{equation}
l'ensemble complet des généalogies qui publient \(t\).
\end{theorem}

\begin{proof}
L'application \(\operatorname{ev}_{\mathbb R}\) est continue, surjective et
constante exactement sur les classes de \(\sim_{\mathbb R}\). Par la
propriété universelle du quotient, elle induit une application continue et bijective
\(\overline{\operatorname{ev}}_{\mathbb R}\). Le domaine est compact, car il est
l'image continue du compact \(\Omega_\infty\), et \(K\) est séparé.
Une bijection continue d'un compact vers un espace séparé est un homéomorphisme.
\end{proof}

Le théorème prouve deux choses et les maintient séparées :
\begin{equation}
 \boxed{
 \begin{array}{c}
 \text{l’espace des histoires est profini et conserve la mémoire,}\\
 \text{son quotient par égalité d’évaluation est le compact visible.}
 \end{array}}
 \label{p12-83-c18-s1:eq:dos-niveles-continuo-u024}
\end{equation}
Il n'est pas affirmé que \(\Omega_\infty\) soit homéomorphe à \(K\).

\section{Atlas généalogique de la droite réelle}

Pour \(m\geq1\), soit \(K_m=[-m,m]\). Supposons donné pour chaque \(m\) un système \(\Omega_\infty^{(m)}\) qui reconstruit \(K_m\), accompagné d’inclusions fermées
\begin{equation}
 j_m:\Omega_\infty^{(m)}\hookrightarrow
 \Omega_\infty^{(m+1)}
 \label{p12-83-c18-s1:eq:inclusion-atlas-genealogico}
\end{equation}
qui satisfont
\begin{equation}
 \operatorname{ev}_{m+1}\circ j_m
 =\iota_m\circ\operatorname{ev}_m,
 \label{p12-83-c18-s1:eq:compatibilidad-atlas-real}
\end{equation}
où \(\iota_m:K_m\hookrightarrow K_{m+1}\) est l'inclusion ordinaire.

\begin{proposition}[Exhaustion de la droite]
\label{p12-83-c18-s1:prop:agotamiento-recta-real-u024}
Le quotient archimédien de la colimite stricte
\begin{equation}
 \Omega_\infty^{\mathbb R}
 :=\varinjlim_m\Omega_\infty^{(m)}
 \label{p12-83-c18-s1:eq:colimite-historias-recta}
\end{equation}
est
\begin{equation}
 \varinjlim_m[-m,m]
 =\bigcup_{m\geq1}[-m,m]
 =\mathbb R
 \label{p12-83-c18-s1:eq:colimite-compactos-recta}
\end{equation}
munie de sa topologie usuelle.
\end{proposition}

\begin{proof}
Chaque pièce possède pour quotient \(K_m\) d'après le
\cref{p12-83-c18-s1:thm:homeomorfismo-cociente-compacto-u024} ;
\eqref{p12-83-c18-s1:eq:compatibilidad-atlas-real} entrelace les quotients. La topologie
finale de l'exhaustion par des compacts coïncide avec la topologie usuelle de
\(\mathbb R\) : un sous-ensemble \(U\subseteq\mathbb R\) est ouvert si et seulement
si \(U\cap[-m,m]\) est ouvert dans \([-m,m]\) pour tout \(m\).
\end{proof}

Le passage d'un intervalle borné à la droite n'efface pas les fibres
\(\mathfrak G_t\). Les inclusions \(j_m\) conservent l'histoire et ne font
qu'élargir la carte visible disponible.

\section{Cartes ternaire équilibrée et décimale par blocs}

La représentation ternaire équilibrée d'un entier a la forme
\begin{equation}
 n=\sum_{j=0}^{J}\epsilon_j3^j,
 \qquad
 \epsilon_j\in\{-1,0,+1\},
 \label{p12-83-c18-s1:eq:carta-ternaria-balanceada-u024}
\end{equation}
avec une règle de retenue fixée. Cette carte conserve le signe local et le
quotient de la réduction ; elle ne doit pas être confondue avec le mot non négatif dans
\(\mathbb F_3\).

La carte décimale par blocs utilise
\begin{equation}
 D_N=1000D_{N-1}+u_N,
 \qquad
 u_N\in\{0,\ldots,999\},
 \qquad D_0=0.
 \label{p12-83-c18-s1:eq:concatenacion-base-mil-u024}
\end{equation}
Le bloc \(u_N\), sa classe modulo neuf et la retenue sont des champs distincts.
Puisque \(1000\equiv1\pmod9\),
\begin{equation}
 D_N\equiv\sum_{j=1}^{N}u_j\pmod9.
 \label{p12-83-c18-s1:eq:base-mil-compatible-mod9-u024}
\end{equation}
Cette congruence ne récupère pas l'ordre des blocs ; la séquence
\((u_1,\ldots,u_N)\) demeure dans le registre.

Pour \(M\geq1\), les cellules décimales semi-ouvertes sont
\begin{equation}
 \mathcal D_{M,q}
 :=[q10^{-M},(q+1)10^{-M}),
 \qquad q\in\mathbb Z.
 \label{p12-83-c18-s1:eq:celda-decimal-semiabierta-u024}
\end{equation}
Un intervalle semi-ouvert \([\ell,h)\) est contenu dans une unique cellule si
et seulement si
\begin{equation}
 \left\lfloor10^M\ell\right\rfloor
 =\left\lceil10^Mh\right\rceil-1.
 \label{p12-83-c18-s1:eq:criterio-exacto-celda-decimal-u024}
\end{equation}
La formule avec plafond à l'extrémité droite couvre également le cas où
\(h\) coïncide exactement avec une frontière décimale.

\begin{definition}[Troncature et arrondi]
Pour \(x\geq0\), nous définissons
\begin{align}
 \operatorname{Tr}_M(x)
 &:=10^{-M}\left\lfloor10^Mx\right\rfloor,
 \label{p12-83-c18-s1:eq:truncamiento-decimal-u024}\\
 \operatorname{Rd}_M(x)
 &:=10^{-M}\left\lfloor10^Mx+\frac12\right\rfloor.
 \label{p12-83-c18-s1:eq:redondeo-decimal-u024}
\end{align}
\end{definition}
Ces opérations coïncident sauf lorsque la fraction écartée franchit le
seuil d'une demi-unité à la dernière position. La distinction est sémantique et
arithmétique ; on ne peut substituer l'une à l'autre pour forcer la compatibilité des
préfixes.

\section{Descente des opérations par les fibres d'évaluation}

Soit \(D_\star\subseteq\Omega_\infty\times\Omega_\infty\) le domaine d'une
opération partielle continue
\begin{equation}
 \star:D_\star\longrightarrow\Omega_\infty.
 \label{p12-83-c18-s1:eq:operacion-historias-u024}
\end{equation}
Nous supposons que \(D_\star\) est saturé par les fibres de
\(\operatorname{ev}_{\mathbb R}\times\operatorname{ev}_{\mathbb R}\).

\begin{theorem}[Critère exact de descente]
\label{p12-83-c18-s1:thm:criterio-descenso-operaciones-u024}
Il existe une unique opération visible
\begin{equation}
 \overline\star:
 (\operatorname{ev}_{\mathbb R}\times\operatorname{ev}_{\mathbb R})
 (D_\star)\longrightarrow K
 \label{p12-83-c18-s1:eq:operacion-descendida-visible}
\end{equation}
qui satisfait
\begin{equation}
 \operatorname{ev}_{\mathbb R}(x\star y)
 =\operatorname{ev}_{\mathbb R}(x)
  \ \overline\star\ 
  \operatorname{ev}_{\mathbb R}(y)
 \label{p12-83-c18-s1:eq:entrelazamiento-operacion-descendida}
\end{equation}
si et seulement si
\begin{equation}
 \begin{split}
 &\operatorname{ev}_{\mathbb R}(x)=
  \operatorname{ev}_{\mathbb R}(x'),\\
 &\operatorname{ev}_{\mathbb R}(y)=
  \operatorname{ev}_{\mathbb R}(y')
 \end{split}
 \quad\Longrightarrow\quad
 \operatorname{ev}_{\mathbb R}(x\star y)=
 \operatorname{ev}_{\mathbb R}(x'\star y')
 \label{p12-83-c18-s1:eq:constancia-operacion-sobre-fibras}
\end{equation}
pour tous les couples du domaine.
\end{theorem}

\begin{proof}
Si \(\overline\star\) existe, le membre de droite de
\eqref{p12-83-c18-s1:eq:entrelazamiento-operacion-descendida} ne dépend que des deux
évaluations, et l'on obtient
\eqref{p12-83-c18-s1:eq:constancia-operacion-sobre-fibras}. Réciproquement, cette
constance permet de définir
\[
 s\ \overline\star\ t
 :=\operatorname{ev}_{\mathbb R}(x\star y)
\]
pour tous représentants \(x,y\) d'évaluations \(s,t\). La
condition prouve que la définition ne dépend pas des représentants.
L'unicité est immédiate.
\end{proof}

Les feuillets additif et multiplicatif d'APP fournissent des opérations sur un
support commun, mais leurs retenues et leurs mémoires ne s'identifient pas. Le
théorème précédent permet à la somme et au produit de descendre vers la carte visible
sans exiger que leurs réalisations généalogiques coïncident.

\section{Correspondance entre cardinalité discrète, mesure cylindrique et publication archimédienne}

À chaque profondeur apparaissent trois opérations qui doivent demeurer
séparées :
\begin{align}
 \operatorname{Count}_n(A)&:=|A|,
 \label{p12-83-c18-s1:eq:operacion-conteo-u024}\\
 \operatorname{Measure}_n(A)&:=\mu_n(A),
 \label{p12-83-c18-s1:eq:operacion-medida-u024}\\
 \operatorname{Publish}_n(a)&:=L_n(a).
 \label{p12-83-c18-s1:eq:operacion-publicacion-u024}
\end{align}
Le dénombrement dépend du nombre d'états ; la mesure dépend de poids
compatibles ; la publication dépend du lecteur. Deux fibres de même
cardinal peuvent avoir des poids distincts, et deux états de même publication
peuvent appartenir à des cylindres distincts. L'égalité de l'une de ces trois
sorties n'implique pas l'égalité des deux autres.

\section{Obstructions de l'évaluation archimédienne : emboîtement, continuité, quotient, mesures projectives et descente des opérations}

\begin{criterion}[Réfutation de l'évaluation et de la mesure]
La construction est réfutée si l'un des faits
suivants est vérifié :
\begin{enumerate}
 \item une extension compatible reçoit un intervalle qui n'est pas contenu
       dans l'intervalle de sa restriction ;
 \item les diamètres ne tendent pas vers zéro, mais l'unicité de
       l'intersection est affirmée ;
 \item l'évaluation n'est pas continue dans la topologie des cylindres ;
 \item on affirme que tout point de \(K\) apparaît sans démontrer la couverture
       cohérente \eqref{p12-83-c18-s1:eq:cobertura-extensiones-compatible-u024} ;
 \item l'application induite par le quotient n'est pas bijective ;
 \item l'espace profini est identifié à son quotient visible ;
 \item les inclusions de l'atlas n'entrelacent pas leurs évaluations ;
 \item on utilise \(\lfloor10^Mh\rfloor\) au lieu de
       \(\lceil10^Mh\rceil-1\) pour un intervalle semi-ouvert et l'on perd
       le cas de frontière ;
 \item on confond troncature et arrondi ;
 \item une opération descendue dépend du représentant choisi dans une
       fibre ;
 \item une famille \((\mu_n)_n\) viole
       \eqref{eq:consistencia-proyectiva-medidas-u024} ;
 \item la mesure d'un cylindre dépend du niveau auquel il est raffiné ;
 \item une mesure conditionnelle cesse d'avoir une somme égale à un ou perd sa cohérence ;
 \item la mesure visible \(\nu\) est identifiée aux mesures internes
       \(\mu_t\) ;
 \item on déduit une flèche entre \(\mathbb R\) et
       \(\mathcal E_\infty\) au seul motif que les deux publications partagent
       un antécédent ;
 \item on utilise le dénombrement d'une fibre comme s'il déterminait sa mesure ou sa
       valeur publiée.
\end{enumerate}
\end{criterion}


\section{De l'espace des histoires à la droite réelle}

Soit
\[
 \Psi:\mathbb Z_3^6\xrightarrow{\ \sim\ }
       (\mathbb F_3^6)^{\mathbb N}
\]
la conjugaison de Hensel entre la mémoire initiale et le ruban complet de
blocs de six trits. Concaténer les six coordonnées de chaque bloc donne
un ruban ternaire
\(
 (a_n)_{n\geq 1}\in\{0,1,2\}^{\mathbb N}
\), sur lequel agit l'évaluation
\[
 \operatorname{ev}_3((a_n))
 :=\sum_{n\geq1}\frac{a_n}{3^n}\in[0,1].
\]
Les deux développements d'une extrémité ternaire représentent la même valeur et
demeurent différenciés avant l'application de l'évaluation. Cela ne constitue pas
une ambiguïté de l'objet généalogique : c'est précisément une fibre du lecteur.

On définit
\[
 X_{\mathrm{raw}}:=\mathbb Z\times\mathbb Z_3^6,
 \qquad
 P(m,x):=m+\operatorname{ev}_3(\operatorname{flat}(\Psi x)).
\]
La coordonnée entière fournit l'atlas dénombrable de la droite et la
coordonnée 3-adique conserve l'histoire complète de la partie fractionnaire.

\begin{theorem}[Reconstruction exacte du quotient archimédien]
\label{p12-83-c18-s2:hmtch:thm-cociente-real}
L'application \(P:X_{\mathrm{raw}}\to\mathbb R\) est surjective. Si
\[
 x\sim_P y
 \quad\Longleftrightarrow\quad
 P(x)=P(y),
\]
alors
\[
 \overline P:X_{\mathrm{raw}}/\!\sim_P\;\longrightarrow\mathbb R,
 \qquad [x]\longmapsto P(x),
\]
est une bijection. En conséquence,
\[
 \left|X_{\mathrm{raw}}/\!\sim_P\right|
 =|\mathbb R|
 =2^{\aleph_0}.
\]
\end{theorem}

\begin{proof}
Toute suite ternaire se regroupe de manière unique en blocs consécutifs de
six chiffres. La conjugaison \(\Psi\) attribue à ce ruban une unique mémoire
3-adique initiale. Étant donné \(r\in\mathbb R\), on prend
\(m=\lfloor r\rfloor\) et un développement ternaire de
\(r-m\in[0,1)\) ; la reconstruction précédente produit
\(x\in\mathbb Z_3^6\) avec \(P(m,x)=r\). Ainsi, \(P\) est surjective.

L'application \(\overline P\) est bien définie par la définition de
\(\sim_P\), elle est surjective puisque \(P\) l'est, et injective puisque deux
classes de même image appartiennent à une même fibre. Enfin,
\[
 |\mathbb Z_3^6|
 =\left|(\mathbb F_3^6)^{\mathbb N}\right|
 =729^{\aleph_0}
 =2^{\aleph_0},
\]
et multiplier par la coordonnée dénombrable \(\mathbb Z\) ne modifie pas ce
cardinal. Le calcul coïncide avec celui du quotient, déjà identifié à
\(\mathbb R\).
\end{proof}

\begin{remark}[Version survivante]
\label{p12-83-c18-s2:hmtch:rem-version-superviviente}
Le théorème précédent est inconditionnel pour la complétion brute. Pour un sous-espace survivant \(X_{\mathrm{surv}}\), le même résultat est obtenu lorsque les recouvrements de chaque niveau sont cohérents, que les diamètres se contractent uniformément et que toute cellule visible possède un prolongement. Sous ces hypothèses, le lemme de la branche infinie produit une section au-dessus de chaque point et l’évaluation induit, compact par compact, un homéomorphisme du quotient avec l’intervalle exprimé. L’exhaustion par \([-m,m]\) reconstruit ensuite la droite entière.
\end{remark}

\section{Moyenne archimédienne et double projection du même état}

Soit \(u=(m,\mathbf u)\in
\mathscr O_{\mathbb R,\leftrightarrow}^{\rm HMT}\). Un préfixe
\(\mathbf u_n\) détermine dans la carte fractionnaire un cylindre orienté
semi-ouvert
\[
 I_n(\mathbf u_n)=[a_n,b_n),
\]
avec la convention complémentaire aux extrémités à double développement. Pour
appliquer la compacité, on utilise son adhérence
\(K_n=[a_n,b_n]\). Les compatibilités donnent
\[
 K_{n+1}\subseteq K_n,
 \qquad
 \delta_n:=b_n-a_n\longrightarrow0.
\]
La moyenne archimédienne globale du cylindre est
\[
 \mu_n(u)=m+\frac{a_n+b_n}{2}.
\]
Si \(q\geq n\), alors
\[
 |\mu_q(u)-\mu_n(u)|\leq\delta_n,
\]
de sorte que
\begin{equation}
 \mathfrak R(u):=\lim_{n\to\infty}\mu_n(u)
 \label{p12-83-c18-s2:hmtch:eq-media-arquimediana}
\end{equation}
existe. Ce qui tend vers zéro est le diamètre du cylindre ; la valeur réelle est
la résultante de l'emboîtement, non ce qui s'annule.

\begin{theorem}[Quotient two-way et droite réelle]
\label{p12-83-c18-s2:hmtch:thm-cociente-two-way-real}
Soit
\[
 u\sim_{\mathfrak R}v
 \quad\Longleftrightarrow\quad
 \mathfrak R(u)=\mathfrak R(v).
\]
Sous la couverture cohérente des cartes HMT, l'évaluation induit un
homéomorphisme
\begin{equation}
 \overline{\mathfrak R}:
 \mathscr O_{\mathbb R,\leftrightarrow}^{\rm HMT}/\!
 \sim_{\mathfrak R}
 \xrightarrow{\ \sim\ }\mathbb R.
 \label{p12-83-c18-s2:hmtch:eq-homeomorfismo-two-way-real}
\end{equation}
En particulier, il s'agit d'un morphisme HMT continu : non du résultat du choix
d'une histoire isolée dans chaque fibre.
\end{theorem}

\begin{proof}
La borne des cylindres prouve la continuité de \(\mathfrak R\). La couverture
produit une histoire au-dessus de chaque réel et donc la surjectivité. En passant
au quotient par les fibres, l'application induite est bijective. Dans chaque
carte compacte, une bijection continue vers un intervalle séparé est un
homéomorphisme. La topologie globale du quotient est la topologie finale du
recollement de ces cartes ; la fibre d'évaluation identifie l'extrémité
\(m+1\) d'une carte à l'extrémité \(0\) translatée de la carte suivante.
Le recouvrement résultant est localement fini, de sorte que les
homéomorphismes de carte se recollent en un homéomorphisme global vers
\(\mathbb R\).
\end{proof}
  \part{Caractères mathématiques du prolongement coinductif}

\chapter{Germes, complétion cubique et synchronisation des caractères fondamentaux}
\section{Les deux projections mécaniques et la conservation exacte de l’unité de bord}

L’affirmation selon laquelle l’unité se conserve au cours du prolongement admet désormais une formulation arithmétique indépendante d’une analogie physique. Soit

\[
 \delta:=\log_{10}\!\left(\frac{10}{9}\right),
 \qquad 0<\delta<1,
\]

l’écart entre la capacité triadique \(3^6=729\) et la complétion décimale \(10^3=1000\). L’irrationalité de \(\delta\) est déjà démontrée dans la construction par valuations premières ; sa transcendance a été prouvée dans ce cahier par Gelfond--Schneider. Considérons les deux codages unilatéraux de la même droite irrationnelle :

\[
 z_n^-:=\lfloor(n+1)\delta\rfloor-\lfloor n\delta\rfloor,
 \qquad
 z_n^+:=\lceil(n+1)\delta\rceil-\lceil n\delta\rceil,
 \qquad n\geq0.
\]

Tous deux prennent leurs valeurs dans \(\{0,1\}\). Il ne s’agit ni de deux densités ni de deux dynamiques : ce sont les deux conventions de fermeture d’une même coupure irrationnelle. Puisque \(n\delta\notin\mathbb Z\) pour tout \(n\geq1\), on a

\[
 z_0^-=0,\qquad z_0^+=1,\qquad z_n^-=z_n^+\quad(n\geq1).
\]

Par télescopage, leurs dénombrements dans le préfixe de longueur \(N>0\) sont

\[
 Z_N^-:=\sum_{n=0}^{N-1}z_n^-=\lfloor N\delta\rfloor,
 \qquad
 Z_N^+:=\sum_{n=0}^{N-1}z_n^+=\lceil N\delta\rceil.
\]

On obtient donc le théorème de bord

\[
 \boxed{Z_N^+-Z_N^-=1\qquad(N>0).}
\]

L’unité n’est pas créée à la fin du bloc : elle était localisée dans le choix du demi-bord appartenant au cylindre. Si l’on définit les deux polarisations

\[
 r_N^-:=N\delta-\lfloor N\delta\rfloor,
 \qquad
 r_N^+:=\lceil N\delta\rceil-N\delta,
\]

alors

\[
 0<r_N^\pm<1,
 \qquad
 \boxed{r_N^-+r_N^+=1.}
\]

C’est la première forme exacte de l’\textbf{évolution conservative de l’unité} : les deux projections holographiques répartissent une unité de frontière et leur somme la retrouve à toute profondeur finie. Aucune conservation d’énergie n’est encore affirmée. La promotion thermodynamique exige en outre que la réalisation physique de l’opérateur soit réversible. Dans l’état TPK complet, l’actualisation bijective implique déjà \(H(X\mid U(X))=0\) ; cet énoncé logique se distingue tant de l’entropie topologique nulle du mot sturmien que d’une réalisation thermique concrète.

La loi est également bidirectionnelle. La lecture inférieure fixe le demi-bord de sortie et la supérieure celui d’entrée. Inverser l’orientation échange les deux conventions et transporte l’unité, sans modifier la pente ni l’ordre à longue portée. C’est le contenu minimal que doit conserver tout lecteur \emph{two-way} : l’inversion change la localisation de l’unité de bord, sans l’éliminer.

\textbf{Critère de réfutation.} Un \(N>0\) tel que \(Z_N^+-Z_N^-\neq1\), ou une formulation qui transforme cette identité arithmétique, sans transducteur, en égalité énergétique.

\section{Cocycle de concaténation, mémoire de retenue et restriction exacte du TPK}

Pour une phase initiale \(k\in\mathbb Z\), le nombre de lacunes d’un bloc de longueur \(N\) est

\[
 V_N(k):=\lfloor(k+N)\delta\rfloor-\lfloor k\delta\rfloor.
\]

Le mot est équilibré, de sorte que

\[
 V_N(k)\in
 \left\{\lfloor N\delta\rfloor,\lceil N\delta\rceil\right\}
\]

pour toute phase. L’identité de composition est plus importante pour la dynamique enrichie

\[
 \boxed{V_{M+N}(k)=V_M(k)+V_N(k+M).}
\]

Le second terme ne peut pas être réévalué à l’origine : il reçoit la phase transportée par le premier bloc. Oublier cette phase fait perdre exactement la retenue

\[
 \kappa_\delta(M,N)
 :=\lfloor(M+N)\delta\rfloor
   -\lfloor M\delta\rfloor-\lfloor N\delta\rfloor
 \in\{0,1\}.
\]

Ainsi,

\[
 \lfloor(M+N)\delta\rfloor
 =\lfloor M\delta\rfloor+\lfloor N\delta\rfloor
  +\kappa_\delta(M,N).
\]

L’associativité de la somme impose l’équation de cocycle

\[
 \kappa_\delta(L,M)+\kappa_\delta(L+M,N)
 =\kappa_\delta(M,N)+\kappa_\delta(L,M+N).
\]

La retenue n’est donc pas une correction numérique ajoutée à la fin du calcul : c’est la donnée qui rend associative la concaténation des blocs lorsque leur trajectoire est conservée. C’est précisément cette structure qu’une projection vers de simples cardinaux dissimule.

La restriction de la dynamique TPK au sous-système de lacunes peut s’écrire sans remplacer la fibre complète. Sur

\[
 \mathcal Y_{108}:=C_{108}\times\mathbb T\times\mathbb Z
\]

on définit

\[
 \mathcal F(r,\theta,m)
 =\bigl(r+1,\{\theta+\delta\},
         m+\lfloor\theta+\delta\rfloor\bigr),
 \qquad 0\leq\theta<1.
\]

La première composante est l’horloge finie ; la deuxième, la phase irrationnelle ; la troisième, la mémoire intégrée. L’application est bijective, d’inverse

\[
 \mathcal F^{-1}(r,\theta,m)
 =\bigl(r-1,\{\theta-\delta\},
 m-\lfloor\{\theta-\delta\}+\delta\rfloor\bigr).
\]

Dans le langage opératoriel directeur, cette formule est un \textbf{facteur} de \(U_t=\operatorname{Upd}_t\circ\operatorname{Tra}_t\circ
\operatorname{Sel}_t\) : le sélecteur distingue événement et non-événement ; le transport fait avancer la phase nonadique et la phase irrationnelle ; l’actualisation incorpore l’impulsion au registre de mémoire. Elle ne remplace pas la fibre TPK, qui conserve en outre feuillet, orientation, résidu, quotient, signature, cylindre, frontière, parcours et survie.

Après \(108\) itérations, on a

\[
 \mathcal F^{108}(r,\theta,m)
 =\bigl(r,\{\theta+108\delta\},m+V_{108}(\theta)\bigr),
\]

avec \(V_{108}(\theta)\in\{4,5\}\). La première horloge revient ; ni la phase irrationnelle ni la mémoire ne reviennent. Cette égalité donne un type exact à l’expression « la phase revient, non l’état » au sein du sous-système apériodique.

\textbf{Critère de réfutation.} L’échec de l’équation de cocycle, la non-inversibilité de \(\mathcal F\) ou l’omission, par la projection du facteur, d’un champ utilisé ensuite pour distinguer des états.

\section{Le spectre renormalisé des cardinaux et ses deux autoréférences exactes}

Définissons le lecteur équilibré de longueur

\[
 \mathcal W_\delta(N)
 :=\{\lfloor N\delta\rfloor,\lceil N\delta\rceil\}.
\]

Cet opérateur associe à chaque échelle entière les deux seuls dénombrements permis par la coupure irrationnelle. Il ne consulte ni \(\pi\), ni \(\varphi\), ni \(e\), ni \(\alpha\), ni une masse, ni une constante expérimentale. Son application à des cardinaux déjà produits par APP--TRIT--TPK fait apparaître les égalités exactes suivantes :

\[
 \begin{aligned}
 \mathcal W_\delta(54)&=\{2,3\},&
 \mathcal W_\delta(81)&=\{3,4\},\\
 \mathcal W_\delta(90)&=\{4,5\},&
 \mathcal W_\delta(120)&=\{5,6\},\\
 \mathcal W_\delta(137)&=\{6,7\},&
 \mathcal W_\delta(202)&=\{9,10\},\\
 \mathcal W_\delta(210)&=\{9,10\},&
 \mathcal W_\delta(243)&=\{11,12\},\\
 \mathcal W_\delta(271)&=\{12,13\},&
 \mathcal W_\delta(324)&=\{14,15\},\\
 \mathcal W_\delta(369)&=\{16,17\},&
 \mathcal W_\delta(468)&=\{21,22\},\\
 \mathcal W_\delta(729)&=\{33,34\},&
 \mathcal W_\delta(1000)&=\{45,46\}.
 \end{aligned}
\]

Deux lignes sont autoréférentielles en un sens précis. L’entier \(137\), obtenu auparavant comme trace de la matrice de synchronisation, reproduit, lorsqu’il est lu par \(\mathcal W_\delta\), les retours secondaires :

\[
 \boxed{\mathcal W_\delta(137)=\{6,7\}.}
\]

Le dénombrement démontré \(468\) reproduit de même les lacunes primaires :

\[
 \boxed{\mathcal W_\delta(468)=\{21,22\}.}
\]

Il ne s’agit pas d’arrondir \(\alpha^{-1}\) à \(137\). La trace \(137\) a été construite auparavant à partir de la synchronisation entière, et le lecteur de lacunes retrouve ensuite \(6/7\). On n’introduit pas non plus \(21/22\) pour sélectionner \(468\) : le dénombrement appartient au prolongement nonadique et le même lecteur le restitue comme spectre de fenêtres. Ces deux identités constituent une fermeture de renormalisation au sein de la généalogie, non une définition rétrospective.

Il existe une troisième identité structurelle particulièrement nette. Puisque

\[
 210=90+120,
\]

la signature ordonnée de l’intervalle est

\[
 \bigl(V_{90}(0),V_{120}(90)\bigr)=(4,5),
\]

et donc

\[
 Z^-_{210}=4+5=9,
 \qquad
 Z^+_{210}=1+4+5=10.
\]

Les canaux \(90/120\), engendrés dans le TPK avant leurs évaluateurs parangulaires, condensent ainsi la transition \(9\to10\) au moyen de la même unité de bord. L’égalité n’affirme pas que \(90\) et \(120\) soient « les nombres » neuf et dix ; elle affirme que le lecteur de capacité \(729/1000\), appliqué à leur concaténation causale, produit exactement ces deux coordonnées.

\textbf{Critère de réfutation.} Un seul écart dans l’un quelconque des planchers ou plafonds exacts, ou l’utilisation d’une valeur reconnue extérieurement pour choisir les longueurs.

\section{La complétion cubique contient conjointement les échelles 34 et 11}

La complétion

\[
 10^3=(9+1)^3=729+243+27+1
\]

possède désormais une image exacte sous le cocycle de lacunes. En conservant l’ordre des quatre termes, le mot mécanique inférieur produit la signature

\[
 \boxed{(33,11,1,0),}
\]

parce que

\[
 \begin{aligned}
 V_{729}(0)&=33,\\
 V_{243}(729)&=11,\\
 V_{27}(972)&=1,\\
 V_1(999)&=0.
 \end{aligned}
\]

La somme est \(45\), tandis que la lecture supérieure place l’unité de bord dans le premier bloc et donne

\[
 \boxed{(34,11,1,0),\qquad 34+11+1+0=46.}
\]

On obtient ainsi une forme affinée de la loi \(45/46\) : elle ne compte pas seulement les événements dans une fenêtre de mille positions, mais conserve aussi la provenance de chaque événement au sein de la complétion \(729+243+27+1\). L’ordre des termes fait partie des données. Une permutation peut redistribuer la retenue entre les blocs, même si la somme totale reste fixe ; la signature enrichie contient donc davantage d’information que le cardinal \(45\) ou \(46\).

L’échelle \(729\) se décompose à son tour en trois blocs visibles de longueur \(243\) :

\[
 \bigl(V_{243}(0),V_{243}(243),V_{243}(486)\bigr)
 =(11,11,11).
\]

En conséquence,

\[
 Z^-_{729}=11+11+11=33,
 \qquad
 \boxed{Z^+_{729}=1+11+11+11=34.}
\]

Ce résultat situe \(11\) et \(34\) dans une seule hiérarchie apériodique avant toute interprétation dimensionnelle. En outre,

\[
 \begin{aligned}
 324&=243+81,& (V_{243}(0),V_{81}(243))&=(11,3),\\
 369&=243+126,&(V_{243}(0),V_{126}(243))&=(11,5),
 \end{aligned}
\]

d'où

\[
 Z^-_{324}=14,\quad Z^+_{324}=15,
 \qquad
 Z^-_{369}=16,\quad Z^+_{369}=17.
\]

Deux connexions avec des objets déjà démontrés sont exactes au niveau cardinal :

\[
 \boxed{Z^+_{324}=15
 =\det\begin{pmatrix}21&22\\6&7\end{pmatrix},}
\]

et

\[
 \boxed{Z^-_{369}=16=|\det H_4|=|\operatorname{coker}H_4|.}
\]

Combinées au théorème déterminantal actif,

\[
 \operatorname{ord}_{10}
 \bigl(\varphi\alpha^{16}\bigr)=-34,
\]

elles produisent la factorisation scalaire

\[
 \boxed{
 -\operatorname{ord}_{10}
 \bigl(\varphi\alpha^{Z^-_{369}}\bigr)
 =Z^+_{729}.}
\]

La formulation préliminaire de ce résultat le décrivait comme un carré réduit aux cardinaux. Les sections 95--96 achèvent son typage : le domaine pertinent n’est pas une énumération arbitraire de seize lacunes, mais la signature ordonnée de la complétion cubique et son lecteur de profondeur. La branche d’action se factorise exactement sous la forme

\[
 \boxed{
 \mathcal V^-_{369}
 \xrightarrow{\;Z^-\;}16
 =|\operatorname{coker}(H_4)|
 \xrightarrow{\;\mathcal N_\alpha\;}
 \alpha^{16}
 \xrightarrow{\;\varphi\cdot\;}
 \Sigma_H
 \xrightarrow{\;-\operatorname{ord}_{10}\;}34
 =Z^+_{729}.}
\]

Ici, \(\mathcal N_\alpha(n)=\alpha^{n}\) est le lecteur de norme de l’exposant déterminantal. La forme normale de Smith de \(H_4\) produit le conoyau d’ordre \(16\) ; la section \(\alpha\) le transporte vers la norme \(\alpha^{16}\) ; l’autoéchelle \(\varphi\) forme \(\Sigma_H\) ; et l’ordre décimal détermine la profondeur \(34\). L’égalité \(Z^-_{369}=|\operatorname{coker}(H_4)|\) n’est donc pas utilisée comme une bijection élément par élément : elle intervient comme égalité d’invariants dans la composition qui démontre la profondeur d’action. C’est précisément le type d’énoncé exigé par le théorème.

De même,

\[
 Z^-_{243}=11
\]

est la seconde coordonnée du même lecteur cubique. Sa factorisation typée est

\[
 \boxed{
 \mathcal V^-_{243}
 \xrightarrow{\;Z^-\;}11
 =\operatorname{pr}_{243}\!\left(
   \mathcal D_{\rm cub}(\sigma^-,\sigma^+)\right)
 \xrightarrow{\;\operatorname{Pub}_G\;}
 \mathcal D_{10}(G)=11.}
\]

L’assemblage avec la branche d’action donne, sans consultation de valeurs cibles,

\[
 \boxed{
 \bigl(\mathcal D_{10}(\hbar),\mathcal D_{10}(G)\bigr)=(34,11),
 \qquad
 \mathcal D_{10}(G\otimes\hbar)=34+11=45,
 \qquad 45+1=46.}
\]

L’identité \(3^6\cdot243=3^{11}\) conserve son domaine propre comme lecteur de dimension ; la coordonnée \(11\) de la complétion et la profondeur gravitationnelle sont reliées par \(\mathcal D_{\rm cub}\) et \(\operatorname{Pub}_G\), non par une identification tacite de leurs éléments. Le produit \(G_a\hbar_a\), démontré dans la section 95, ferme le carré dynamique : le retour two-way inverse réciproquement les deux branches et conserve leur produit.

\textbf{Critère de réfutation.} L’absence de télescopage des signatures ordonnées, un déterminant ou un conoyau différent de la valeur déclarée, l’impossibilité pour \(\mathcal D_{\rm cub}\) de produire \((34,11,45,46)\), ou une réalisation dimensionnelle dépendant de \(\hbar\), de \(G\), de CODATA ou de la position décimale qu’elle doit précisément engendrer.

\section{La scission électronique de l’espace ambiant et l’apparition interne de 9, 10 et 24}

Lors de la première réouverture conjointe, le canal électronique contient \(527\) extensions d’un espace ambiant de \(729\). Son déficit est

\[
 202=729-527=6+28\cdot7.
\]

La décomposition ordonnée de l’espace ambiant

\[
 729=202+527
\]

est exactement compatible avec le lecteur apériodique :

\[
 \bigl(V_{202}(0),V_{527}(202)\bigr)=(9,24).
\]

Par conséquent,

\[
 \boxed{Z^-_{729}=9+24=33,\qquad
        Z^+_{729}=10+24=34.}
\]

Le même déficit qui singularise la fenêtre électronique reproduit, sous \(\mathcal W_\delta\), la frontière nonadique--décimale :

\[
 \boxed{\mathcal W_\delta(202)=\{9,10\}.}
\]

L’unité de bord agit ici sur le secteur déficitaire ; le secteur survivant conserve vingt-quatre événements dans la phase canonique. Cette décomposition affine l’égalité précédente \(34=1+3\cdot11\) : la même coordonnée se lit désormais comme \(34=10+24\). Les deux factorisations coexistent parce qu’elles enregistrent des coupures distinctes du même cylindre, et non parce que l’une remplace l’autre.

L’entier \(24\) coïncide avec le rang du réseau de Leech, mais l’égalité cardinale n’identifie pas les \(527\) extensions électroniques à des vecteurs du réseau. Son contenu exact est une connexion de rang dans la seconde projection. L’entrelacement exceptionnel est construit, avec son incidence et sa forme quadratique propres, par la chaîne Paley--Witt--Golay--Leech ; aucune action du groupe Monstre n’est attribuée aux chiffres ni aux lacunes.

\textbf{Critère de réfutation.} Un dénombrement électronique différent de \(527\), un bloc complémentaire ne contenant pas vingt-quatre événements, ou l’absence de tout transport préservant l’incidence exceptionnelle exigée.

\section{Les deux horloges finies et la suspension hélicoïdale apériodique}

L’horloge \(108\) se divise en deux demi-cycles de \(54\) pas. Le mot inférieur possède la signature

\[
 \bigl(V_{54}(0),V_{54}(54)\bigr)=(2,2),
\]

de sorte que

\[
 Z^-_{108}=2+2=4,
 \qquad Z^+_{108}=1+2+2=5.
\]

Le retour \(54+54\) ferme la phase finie et conserve un choix de demi-bord. La triple répétition de l’horloge ne répète toutefois pas la même signature :

\[
 \bigl(V_{108}(0),V_{108}(108),V_{108}(216)\bigr)=(4,5,5),
\]

et donc

\[
 Z^-_{324}=4+5+5=14,
 \qquad Z^+_{324}=1+4+5+5=15.
\]

C’est une manifestation directe de l’ordre apériodique : trois blocs de même longueur possèdent des dénombrements équilibrés, mais non identiques. La séquence conserve un ordre à longue portée et un écart inférieur à un, sans acquérir de période de translation.

Le cardinal \(324\) admet deux lectures internes déjà présentes dans la construction : \(324=81\cdot4\) pour les parcours orientés et \(324=108\cdot3\) pour un relèvement tritique de l’horloge. La seconde factorisation définit le modèle cyclique minimal

\[
 0\longrightarrow C_3\longrightarrow C_{324}
 \longrightarrow C_{108}\longrightarrow0.
\]

L’extension n’est pas scindée : \(C_{324}\) contient un élément d’ordre \(324\), tandis que \(C_3\times C_{108}\) a pour exposant \(108\). Dans l’horloge relevée, une translation de \(108\) pas revient à la même coordonnée de \(C_{108}\), mais avance au feuillet suivant du noyau \(C_3\) ; la coordonnée relevée ne revient qu’après \(324\) pas. C’est le plus petit modèle algébrique d’une hélice à trois feuillets au-dessus de l’horloge \(108\).

En décorant cette horloge par la rotation irrationnelle de pente \(\delta\), on obtient une suspension hélicoïdale sturmienne : la base finie peut revenir, le feuillet relevé peut revenir avec une période plus longue, et la phase irrationnelle ne revient jamais exactement. La mémoire accumulée distingue ces trois faits. L’hélice ascendante utilise \(\mathcal F\) ; l’hélice descendante utilise \(\mathcal F^{-1}\). Ce sont des lectures conjuguées de la même dynamique réversible, non deux séries indépendantes. Le miroir \(\pi/e\) reste une autre involution : il agit sur les fenêtres de sortie et fixe \(\varphi\) ; il n’inverse pas à lui seul l’orientation temporelle.

\textbf{Critère de réfutation.} Une section de groupes de \(C_{324}\to C_{108}\), une répétition \((4,4,4)\) ou une conjugaison qui efface la mémoire.

\section{La première supercellule qui ferme l’horloge 108}

La fraction continue de \(\delta\) commence par

\[
 [0;21,1,5,1,6,2,3,1,1,8,1,2,\ldots].
\]

Le premier convergent dont le dénominateur est divisible par \(108\) est

\[
 \boxed{\frac{7813}{170748}.}
\]

Aucun dénominateur de convergent antérieur ne possède cette propriété. En outre, le certificat rationnel prouve que

\[
 \left\lfloor n\delta\right\rfloor
 =\left\lfloor\frac{7813n}{170748}\right\rfloor
 \qquad(0\leq n<170748),
\]

tandis qu’à l’extrémité finale

\[
 \lfloor170748\delta\rfloor=7812,
 \qquad
 \frac{7813\cdot170748}{170748}=7813.
\]

Le mot apériodique et son approximant périodique de Christoffel coïncident sur tout préfixe propre de la supercellule et ne diffèrent qu’à la fermeture. L’approximant périodique ne corrige pas un nuage d’erreurs : il transporte exactement l’unité de bord de la projection inférieure vers la supérieure.

Le dénominateur possède la factorisation

\[
 \boxed{
 170748=108\cdot1581
       =324\cdot527
       =2\cdot54\cdot3\cdot527.}
\]

Ici, \(527\) est exactement le dénombrement du canal électronique lors de la première réouverture et \(324\) le cardinal des parcours orientés. Puisque \(\gcd(324,527)=1\), il existe la décomposition canonique d’horloges

\[
 C_{170748}\cong C_{324}\times C_{527}.
\]

Le théorème chinois des restes justifie cette décomposition de groupes. Le second facteur possède le cardinal de la fibre électronique et le premier celui de l’atlas orienté ; cette coïncidence de cardinaux n’identifie pas les objets et ne transfère pas automatiquement leurs actions.

Le couple d’horloges fournit une seconde donnée inattendue. Dans la branche inférieure, après \(170748\) blocs,

\[
 7812\equiv36\pmod{108},
 \qquad
 170748-7812=162936\equiv72\pmod{108}.
\]

Par conséquent, bien que l’horloge des blocs se ferme,

\[
 \boxed{(r_{\rm vac},r_{\rm cap})=(36,72),\qquad36+72=108.}
\]

La branche supérieure transfère l’unité de bord :

\[
 (r_{\rm vac}^+,r_{\rm cap}^+)=(37,71),
 \qquad37+71=108.
\]

L’apparition de \(36\) dans la branche inférieure coïncide avec l’indice du stade \(R_{36}\) du prolongement coinductif. Cette congruence n’identifie pas à elle seule le résidu de lacunes à la coordonnée de \(R_{36}\) ; elle démontre que la supercellule ferme l’horloge des blocs, conserve la répartition de l’unité et laisse une mémoire résiduelle non triviale dans l’horloge de capacité.

Le numérateur possède une autre factorisation interne :

\[
 7813=13\cdot601,
 \qquad
 7812=36\cdot217,
 \qquad217=2\cdot108+1.
\]

Il en découle le repère unimodulaire

\[
 \boxed{
 M_\star=
 \begin{pmatrix}
 601&217\\
 36&13
 \end{pmatrix}\in\operatorname{SL}_2(\mathbb Z),
 \qquad
 601\cdot13-217\cdot36=1.}
\]

Les quatre entiers sont localisés dans l’architecture : \(601\) est la douzième coordonnée du sceau dodécaphasique \(K\) ; \(13\) est le nombre de multisections de l’atlas ; \(36\) est le stade \(R_{36}\) ; et \(216\) est le retour cinématique étendu. L’égalité est exacte et n’utilise aucune valeur physique. Elle définit un repère unimodulaire de leurs cardinaux ; elle ne postule pas, pour cette seule raison, une équivalence catégorique entre les quatre objets.

\textbf{Critère de réfutation.} Un convergent antérieur de dénominateur multiple de \(108\), un écart dans un préfixe propre, une factorisation fausse ou un transport qui ne conjugue pas les actions.

\section{La seconde dérivation de \(\alpha\) comme renormalisation d’un seul mot}

Le polynôme \(P_9\) de la seconde voie de \(\alpha\) contient les couples \(21/22\), \(6/7\) et \(45/46\). Jusqu’ici, leurs coefficients avaient été reconstruits par la matrice de synchronisation, les retours de Beatty et la frontière \(7{:}9\). Le nouveau lecteur démontre que les trois couples sont aussi des images du \textbf{même} mot mécanique à trois échelles internes :

\[
 \boxed{
 \mathcal W_\delta(468)=\{21,22\},\qquad
 \mathcal W_\delta(137)=\{6,7\},\qquad
 \mathcal W_\delta(1000)=\{45,46\}.}
\]

La première longueur est le dénombrement des émissions décimales, la deuxième la trace entière de synchronisation et la troisième la complétion cubique. La généalogie de la seconde voie peut donc s’écrire sans introduire ses coefficients comme une liste indépendante :

\[
 \begin{aligned}
 &\mathrm{APP}\longrightarrow\mathrm{TRIT}\longrightarrow\mathrm{TPK}
 \longrightarrow X^{\rm enr},\\
 &729\longrightarrow1000\longrightarrow\delta
 \longrightarrow\mathcal W_\delta,\\
 &(468,137,1000)
 \longmapsto
 \bigl(\{21,22\},\{6,7\},\{45,46\}\bigr),\\
 &\text{signature de synchronisation et de bord }7{:}9
 \longrightarrow P_9
 \longrightarrow\alpha.
 \end{aligned}
\]

Cette chaîne ne remplace pas la première voie dodécaphasique. Toutes deux partent de l’état coinductif qui détermine également \(\pi\), \(\varphi\) et \(e\) ; elles conservent ensuite leurs lecteurs propres. La racine de \(P_9\) est calculée à la fin, une fois la signature construite. La valeur de \(\alpha\) ne sélectionne jamais \(468\), \(137\), \(1000\) ni leurs fenêtres.

Le numérateur \(7813=13\cdot601\) situe en outre une fermeture possible entre la supercellule apériodique et le sceau dodécaphasique : \(601=K_{12}\) appartient à la première voie, tandis que \(13\) appartient à l’atlas ultérieur. L’identité unimodulaire de la partie précédente permet de poser correctement la question : il faut déterminer si le module de mémoire sur lequel agissent \(D_3K\), \(R_{36}\), le retour \(216\) et les treize multisections admet le changement de base \(M_\star\). La coïncidence des facteurs ne doit pas être transformée en preuve, mais elle ne doit pas non plus être écartée comme une simple curiosité : elle est localisée dans la première supercellule sélectionnée par l’horloge \(108\), non dans un entier recherché après coup.

\textbf{Critère de réfutation.} Un couple ne provenant pas de \(\mathcal W_\delta\), l’intervention de \(P_9\) dans la sélection de ses propres entrées, ou un module proposé pour \(M_\star\) ne conservant pas le registre de mémoire.

\section{Approfondissement des germes et séparation rigoureuse des échelles physiques}

L’espace d’un germe orienté possède le cardinal

\[
 |\mathcal S|=9^2\cdot4=324,
\]

et le balayage des couples comporte

\[
 |\mathcal S\times\mathcal S|=324^2=104976.
\]

L’application du lecteur apériodique au cardinal total donne

\[
 \mathcal W_\delta(104976)=\{4803,4804\}.
\]

Cette égalité est un dénombrement de la droite mécanique, pas encore un dénombrement de germes. Pour en faire une propriété du balayage APP, il faut fixer un ordre canonique

\[
 \operatorname{ord}_{\mathcal S^2}:
 \mathcal S\times\mathcal S\overset\sim\longrightarrow
 \{0,\ldots,104975\}
\]

compatible avec les deux feuillets, l’action diédrique, le sélecteur TRIT et l’actualisation TPK. Alors seulement peut-on demander quels germes occupent les lacunes et comment le dénombrement change sous raffinement intérieur ou expansion extérieure. La seule cardinalité ne doit pas sélectionner cet ordre.

La nouvelle analyse détermine cependant déjà où chercher les échelles \(10^{-11}\) et \(10^{-34}\). Elles ne s’obtiennent pas en élevant arbitrairement un germe à une puissance décimale. Le cristal produit d’abord les entiers

\[
 Z^-_{243}=11,
 \qquad
 Z^+_{729}=34,
\]

au sein de la même signature de complétion. Deux transducteurs distincts peuvent ensuite réaliser ces ordres sur des droites dimensionnelles :

\[
 \begin{aligned}
 \mathscr T_G&:\bigl(X^{\rm enr},\mathcal V^-_{243}\bigr)
 \longrightarrow 10^{-11}\mathcal L_G,\\
 \mathscr T_\hbar&:\bigl(X^{\rm enr},\mathcal V^+_{729},G_H,\alpha,\varphi\bigr)
 \longrightarrow10^{-34}\mathcal L_S.
 \end{aligned}
\]

Le second transducteur possède déjà une voie démontrée par \(|G_H|=16\) et \(\varphi\alpha^{16}\). Le premier conserve, dans la source faisant autorité, une carte décimale et des critères de sélection, et doit être concilié avec le lecteur de lacunes sans utiliser CODATA. L’expression « fractaliser vers l’intérieur jusqu’à \(10^{-11}\) et \(10^{-34}\) » acquiert ainsi un sens mathématique précis : raffiner le même état jusqu’à deux cardinaux de fenêtres, puis seulement appliquer la transduction dimensionnelle correspondante. Les puissances de dix appartiennent à la carte de la droite dimensionnelle, non au germe nu.

La conservation de l’unité se maintient pendant ce passage si le transducteur préserve simultanément les deux projections et le cocycle. L’exigence est un carré de naturalité

\[
 \mathscr T\circ\operatorname{Upd}
 =\operatorname{Upd}_{\mathcal L}\circ\mathscr T
\]

dans lequel le membre gauche actualise d’abord la mémoire discrète et le membre droit transporte d’abord la section dimensionnelle. Cette équation, et non une égalité décimale isolée, est la preuve requise qu’une constante physique émerge « à l’échelle appropriée ».

La structure discrète du continu doit recevoir simultanément toutes ces fibres. Le raffinement des germes, la résolubilité des extensions, la construction de la mesure, la cohérence et le lecteur déterminantal ne constituent pas cinq objets séparables : ce sont cinq aspects inséparables du même passage à la limite. Le cristal apériodique fournit le calendrier équilibré ; APP conserve somme, produit, résidu, quotient et retenue ; TRIT type l’orientation et les trois courbures ; TPK conserve la trajectoire pendant la sélection, le transport et l’actualisation. Seul leur assemblage produit la droite réelle comme sortie. La droite \(\mathbb R\) n’est pas, dans cette généalogie, le support présupposé pour construire le protocole : elle est l’une de ses réalisations terminales.

\textbf{Critère de réfutation.} Un ordre non naturel vis-à-vis d’APP, de TRIT et de TPK ; une puissance décimale choisie à partir de la valeur physique ; ou un transducteur ne préservant pas l’unité de bord et le registre de mémoire.

\section{La couronne cellulaire du continu contient l’horloge 108 et la signature TRIT}

La révision de la construction cellulaire de référence apporte un élément reliant la complétion, l’horloge et la double projection avec davantage de rigidité qu’une coïncidence de chiffres. Pour le couple cubique

\[
 Q_{9,3}\subset Q_{10,3},
\]

le complexe relatif possède le vecteur de cardinaux

\[
 \mathbf c=(c_0,c_1,c_2,c_3)=(271,756,702,217).
\]

Sa caractéristique d’Euler est nulle :

\[
 \chi(\mathbf c)=271-756+702-217=0.
\]

Cette nullité n’est toutefois pas une annulation opaque. Chaque coordonnée admet une décomposition en cardinaux appartenant déjà à l’architecture :

\[
 \begin{aligned}
 271&=243+27+1,\\
 756&=729+27,\\
 702&=729-27,\\
 217&=216+1=2\cdot108+1.
 \end{aligned}
\]

L’appariement des degrés extrêmes et intérieurs fait apparaître deux défauts égaux :

\[
 \boxed{c_0-c_3=271-217=54,}
 \qquad
 \boxed{c_1-c_2=756-702=54.}
\]

Par conséquent,

\[
 \boxed{(c_0-c_3)+(c_1-c_2)=54+54=108.}
\]

L’horloge \(108\) n’a pas été imposée à la couronne par un calendrier extérieur : son double demi-cycle est contenu dans la différence entre les degrés de bord et dans celle entre les degrés intérieurs. La caractéristique nulle se réécrit

\[
 \chi(\mathbf c)=(c_0-c_3)-(c_1-c_2)=54-54=0.
\]

Cette identité distingue deux sens. Le premier \(54\) mesure la séparation entre les extrémités du complexe ; le second mesure la séparation interne. Leur égalité rend la fermeture compatible : ce qui sort par un appariement revient par l’autre sans effacer le degré cellulaire. La somme \(54+54\) mesure la longueur du parcours bidirectionnel ; la différence \(54-54\) certifie sa fermeture d’Euler. Ce sont deux opérations distinctes sur le même couple, qui ne doivent pas être réduites au seul nombre \(108\).

On applique maintenant à chaque degré le lecteur apériodique

\[
 \mathcal W_\delta(N)=
 \{\lfloor N\delta\rfloor,\lceil N\delta\rceil\}.
\]

Les deux projections de lacunes de la couronne sont

\[
 \mathbf v^-=(12,34,32,9),
 \qquad
 \mathbf v^+=(13,35,33,10).
\]

Leurs compléments de capacité sont

\[
 \mathbf b^+=\mathbf c-\mathbf v^-
 =(259,722,670,208),
\]

\[
 \mathbf b^-=\mathbf c-\mathbf v^+
 =(258,721,669,207).
\]

Dans les deux fermetures unilatérales, on vérifie composante par composante

\[
 \mathbf c=\mathbf v^-+\mathbf b^+
           =\mathbf v^++\mathbf b^-.
\]

L’information décisive apparaît lorsqu’on prend la caractéristique alternée :

\[
 \boxed{\chi(\mathbf v^-)=\chi(\mathbf v^+)=+1,}
\]

\[
 \boxed{\chi(\mathbf b^-)=\chi(\mathbf b^+)=-1.}
\]

La couronne de charge totale nulle se polarise donc exactement comme

\[
 \boxed{0=(+1)+(-1).}
\]

L’indépendance à l’égard de la convention de demi-bord est structurelle :

\[
 \mathbf v^+-\mathbf v^-=(1,1,1,1),
 \qquad
 \chi(1,1,1,1)=0.
\]

Déplacer l’unité de bord entre les deux fermetures ne modifie pas la charge d’Euler de la branche. On obtient ainsi une réalisation interne de l’alphabet TRIT :

\[
 \boxed{
 \text{projection de vacance}\mapsto+1,
 \quad
 \text{couronne complète}\mapsto0,
 \quad
 \text{projection de capacité}\mapsto-1.}
\]

Cela ne remplace pas la définition générale du TRIT et n’identifie pas à soi seul les trois charges à des champs physiques. Un fait plus élémentaire et antérieur est prouvé : la double projection d’un objet cellulaire unique engendre, sans ajout de symboles, les trois classes orientées (+1,0,-1). Les algèbres

\[
 \mathbb A_\tau=\mathbb R[J_\tau]/(J_\tau^2+\tau I)
\]

peuvent ensuite agir sur ces classes. Le type elliptique, parabolique ou hyperbolique n’est plus sélectionné à partir d’une géométrie continue présupposée : la signature qui l’indexe a émergé de la polarisation discrète du complexe.

\textbf{Critère de réfutation.} Une modification de l’un des quatre cardinaux, un défaut ne valant plus \(54\), des projections ne possédant pas les charges (+1/-1), ou l’utilisation d’une géométrie classique pour sélectionner rétrospectivement le signe.

\section{Repère polarisé, racine de bord et évolution conservative de l’unité}

Le résultat précédent est un cas d’une loi exacte valable pour toute longueur. Soit

\[
 a_N=\lfloor N\delta\rfloor,
 \qquad N>0.
\]

Les deux conventions de fermeture déterminent les vecteurs

\[
 u_N^-=(a_N,N-a_N),
 \qquad
 u_N^+=(a_N+1,N-a_N-1).
\]

Chaque ligne conserve exactement le bloc :

\[
 (1,1)u_N^-=(1,1)u_N^+=N.
\]

La différence est indépendante de (N) :

\[
 \boxed{u_N^+-u_N^-=(1,-1).}
\]

L’unité n’apparaît ni ne disparaît ; elle est transportée de la coordonnée de capacité à celle de lacune. Le vecteur ((1,-1)) engendre le noyau de l’augmentation

\[
 \mathbb Z^2\xrightarrow{\ (x,y)\mapsto x+y\ }\mathbb Z,
\]

c’est-à-dire la racine \((A_1)\). La double projection possède donc une direction de frontière canonique. Il n’est pas nécessaire de choisir une métrique physique pour la définir.

Le repère des deux polarisations possède le déterminant

\[
 \boxed{
 \det\begin{pmatrix}
 a_N&N-a_N\\
 a_N+1&N-a_N-1
 \end{pmatrix}=-N.}
\]

La démonstration est immédiate :

\[
 a_N(N-a_N-1)-(N-a_N)(a_N+1)=-N.
\]

Ainsi, le transfert d’une seule unité de bord balaie un parallélogramme entier d’aire orientée (N). Après normalisation par (N), les deux lignes appartiennent au simplexe (x+y=1), sont séparées par ((1,-1)/N) et convergent vers

\[
 (\delta,1-\delta)=(\delta,\rho).
\]

C’est une forme précise de l’évolution conservative de l’unité dimensionnelle : à chaque échelle, il existe une unité normalisée ; le raffinement ne la multiplie pas, mais réduit la taille de son quantum de frontière tout en conservant la direction entière qui permet de le retrouver.

La trajectoire complète s’écrit dans le réseau

\[
 \gamma(N)=
 \bigl(\lfloor N\delta\rfloor,
       N-\lfloor N\delta\rfloor\bigr).
\]

Chaque pas incrémente exactement une coordonnée :

\[
 \gamma(N+1)-\gamma(N)\in\{(1,0),(0,1)\}.
\]

De plus,

\[
 \left\|\gamma(N)-N(\delta,\rho)\right\|_\infty<1.
\]

Dans sa section réticulaire, l’hélice apériodique est une droite discrète d’écart uniformément borné. Elle possède un ordre à longue portée parce qu’elle ne s’écarte jamais de plus d’une cellule de la direction limite, mais elle n’a pas de période de translation puisque \(\delta\) est irrationnel. Le mot sturmien n’orne pas cette trajectoire : il est la suite de décisions indiquant laquelle des deux coordonnées reçoit l’unité suivante.

La même structure d’augmentation réapparaît dans les trois canaux. Le vecteur de réouverture

\[
 v_\Omega=(3,-7,4)
\]

satisfait (3-7+4=0) et appartient donc au réseau

\[
 A_2=\ker\bigl(\mathbb Z^3\xrightarrow{x+y+z}\mathbb Z\bigr).
\]

La racine \((A_1)\) régit le transfert entre les deux projections ; le défaut \((A_2)\) régit la redistribution entre les trois fenêtres corrélées. Ce n’est pas le même vecteur et ils ne relèvent pas du même rang, mais constituent deux degrés successifs d’une seule loi de conservation par augmentation. Cela explique pourquoi double lecture et orientation ternaire peuvent coexister sans transformer la théorie en codage binaire : la conservation fixe la diagonale, tandis que la donnée HMT réside dans son complément orienté.

\textbf{Critère de réfutation.} Une échelle dont le repère n’a pas pour déterminant (-N), un pas modifiant simultanément les deux coordonnées, ou une identification de \((A_1)\) et \((A_2)\) ne déclarant pas l’application de changement de rang.

\section{Charge topologico-temporelle de bord et sélection non destructive de courbure}

La dérivée tritique du mot de lacunes,

\[
 \tau_n=z_{n-1}-z_n\in\{+1,0,-1\},
\]

possède une loi de conservation qui complète son interprétation comme unité d’information topologico-temporelle. Pour tout intervalle fini \((a\leq n\leq b)\),

\[
 \boxed{
 Q_{[a,b]}:=\sum_{n=a}^{b}\tau_n
 =z_{a-1}-z_b\in\{-1,0,+1\}.}
\]

La charge orientée de l’intérieur se télescope entièrement et reste supportée aux deux extrémités. Le TRIT visible n’accumule pas une charge extensive arbitraire : il enregistre la manière dont le bloc s’ouvre, persiste ou se ferme relativement à sa frontière. Le registre de mémoire du TPK conserve les termes intérieurs que la somme alternée dissimule ; le télescopage n’équivaut donc pas à une perte de trajectoire.

Puisque \((0<\delta<1/2)\), deux lacunes ne sont pas consécutives. Les symboles non nuls de \(((\tau_n))\) alternent donc strictement :

\[
 -1,+1,-1,+1,\ldots
\]

lorsque les zéros intermédiaires sont supprimés. Une ouverture doit se fermer avant qu’une autre puisse se produire. Les zéros ne signifient pas une absence de dynamique : ce sont les plateaux pendant lesquels le régime se conserve tandis que l’horloge des blocs continue d’avancer. Lorsque la dynamique est induite sur les événements, ces plateaux possèdent les longueurs \(6\) ou \(7\), comme cela a déjà été démontré.

Chaque symbole sélectionne l’une des trois exponentielles

\[
 \exp(tJ_\tau)=C_\tau(t)I+S_\tau(t)J_\tau,
 \qquad J_\tau^2=-\tau I.
\]

Le mot ordonne ainsi rotations elliptiques, transports nilpotents et ouvertures hyperboliques. La sélection ne détruit pas l’apériodicité car l’application différentielle est inversible dès lors qu’un bit de frontière est conservé :

\[
 z_n=z_0-\sum_{k=1}^{n}\tau_k.
\]

C’est la raison exacte pour laquelle le TRIT peut lire trois courbures sans remplacer le mot qui les engendre. L’unité tritique est le changement visible ; le TPK conserve la donnée de frontière, le feuillet APP, la retenue et l’ordre des opérateurs nécessaires à l’inversion de la lecture.

Le renversement temporel

\[
 C_{\rm rev}(\tau_1,\ldots,\tau_n)
 =(-\tau_n,\ldots,-\tau_1)
\]

échange ouverture et fermeture, préserve le seuil et satisfait \((C_{\rm rev}^2=I)\). Dans la fibre matricielle, il inverse aussi l’ordre du produit

\[
 \mathcal A_N=exp(t_NJ_{\tau_N})\cdots
               \exp(t_1J_{\tau_1}).
\]

La double direction temporelle n’est donc pas la lecture d’une liste à l’envers. Elle est la conjugaison d’un mot, de ses signes et de l’ordre non commutatif de ses transports, en conservant la frontière permettant de reconstruire l’état.

\textbf{Critère de réfutation.} Deux ouvertures consécutives dans le mot induit, une charge d’intervalle non portée par le bord, ou une reconstruction se passant de \((z_0)\).

\section{Tour \(S\)-adique non stationnaire des deux horloges}

La fraction continue de la pente interne commence par

\[
 \delta=[0;21,1,5,1,6,2,3,1,1,8,1,2,\ldots].
\]

Pour chaque coefficient positif (a), définissons l’opérateur de continuant

\[
 A(a)=\begin{pmatrix}a&1\\1&0\end{pmatrix},
 \qquad\det A(a)=-1.
\]

Si \((p_n/q_n)\) est le (n)-ième convergent, alors

\[
 A(a_1)\cdots A(a_n)
 =\begin{pmatrix}q_n&q_{n-1}\\p_n&p_{n-1}\end{pmatrix}.
\]

Les premiers stades sont

\[
 \begin{aligned}
 A(21)&=\begin{pmatrix}21&1\\1&0\end{pmatrix},\\
 A(21)A(1)&=\begin{pmatrix}22&21\\1&1\end{pmatrix},\\
 A(21)A(1)A(5)&=\begin{pmatrix}131&22\\6&1\end{pmatrix},\\
 A(21)A(1)A(5)A(1)&=\begin{pmatrix}153&131\\7&6\end{pmatrix}.
 \end{aligned}
\]

Cela démontre que \(21/22\) et \(6/7\) ne sont pas deux régularités juxtaposées. Les entiers (21,22) sont les premiers dénominateurs du squelette diophantien ; (6,7) sont les numérateurs qui apparaissent deux inductions plus tard :

\[
 \begin{aligned}
 131&=5\cdot22+21,&6&=5\cdot1+1,\\
 153&=1\cdot131+22,&7&=1\cdot6+1.
 \end{aligned}
\]

Les lacunes primaires \(21/22\) et les retours secondaires \(6/7\) appartiennent donc à la même renormalisation. Les convergents suivants sont

\[
 \frac{48}{1049},\quad
 \frac{103}{2251},\quad
 \frac{357}{7802},\quad
 \frac{460}{10053},\quad
 \frac{817}{17855},\quad
 \frac{6996}{152893},\quad
 \frac{7813}{170748}.
\]

Chaque couple consécutif conserve une cellule primitive du réseau :

\[
 \boxed{p_nq_{n-1}-p_{n-1}q_n=(-1)^{n-1}.}
\]

Le premier dénominateur de cette tour divisible par \(108\) est \(170748\), et son apparition conserve la récurrence complète :

\[
 \boxed{
 170748=152893+17855,
 \qquad
 7813=6996+817.}
\]

À ce stade,

\[
 \begin{pmatrix}
 170748&152893\\
 7813&6996
 \end{pmatrix}
\in GL_2(\mathbb Z),
 \qquad\det=-1.
\]

La tour conserve donc l’unité réticulaire même lorsque la supercellule devient assez grande pour fermer l’horloge finie. La fermeture de \(C_{108}\) ne rend pas périodique la direction irrationnelle : l’approximant rationnel coïncide avec le mot sur tout préfixe propre et ne transfère une unité qu’au dernier demi-bord.

La renormalisation est en outre véritablement ordonnée. Pour \(a\ne b\),

\[
 \boxed{
 A(a)A(b)-A(b)A(a)
 =(a-b)\begin{pmatrix}0&1\\-1&0\end{pmatrix}.}
\]

Par exemple, \(a=21\) et \(b=5\) produisent le défaut orienté

\[
 16\begin{pmatrix}0&1\\-1&0\end{pmatrix}.
\]

La rotation \((x\mapsto x+\delta)\) sur le cercle reste abélienne. La non-commutativité apparaît lorsqu’on ordonne les opérateurs de renormalisation et, dans la fibre complète, les feuillets APP et les générateurs de courbure. Cette distinction élimine deux erreurs possibles : qualifier la rotation de base de non abélienne, ou réduire toute la dynamique à une rotation abélienne en oubliant le cocycle.

Puisque \(\delta\) est transcendant, ce n’est pas un irrationnel quadratique. D’après le théorème de Lagrange, sa fraction continue n’est pas ultimement périodique. La directive de continuants ne se stabilise pas en une unique substitution périodique : le cristal est \(S\)-adique et non stationnaire. Il conserve néanmoins une complexité linéaire, une entropie nulle et des déterminants unimodulaires. C’est une définition du « cristal apériodique » beaucoup plus forte que la simple absence de répétition : il existe une renormalisation à toutes les échelles, mais aucune cellule finale de renormalisation ne se répète indéfiniment.

\textbf{Critère de réfutation.} Un écart dans les coefficients, un déterminant différent de \(\pm1\), un dénominateur de convergent antérieur divisible par \(108\), ou une directive ultimement périodique.

\section{Généalogie cardinale de la couronne, de la frontière \(369\) et du sceau dodécaphasique}

La tour permet de déterminer les cardinaux structurels sans les sélectionner par ressemblance décimale. Les couples suivants proviennent du même lecteur et possèdent des dérivations indépendantes :

\[
 \begin{aligned}
 \mathcal W_\delta(217)&=\{9,10\},\\
 \mathcal W_\delta(271)&=\{12,13\},\\
 \mathcal W_\delta(601)&=\{27,28\},\\
 \mathcal W_\delta(973)&=\{44,45\}.
 \end{aligned}
\]

L’entier \(217\) est le cardinal des \(3\)-cellules de la couronne, et aussi \(2\cdot108+1\). Sa lecture détermine la frontière entre la base nonadique et la complétion décimale. L’entier \(271\) est la couronne cubique \(1000-729\) ; sa lecture détermine \(12/13\), précisément le couple formé par la dodécaphase et le nombre de multisections de l’atlas. L’entier \(601\) est la douzième coordonnée du sceau \(K\) ; sa lecture détermine \(27/28\), la frontière cubique et le cardinal de la microfibre. Enfin,

\[
 973=271+702=756+217
\]

est la somme commune des degrés pairs et impairs de la couronne, et

\[
 \mathcal W_\delta(973)=\{44,45\}.
\]

Ici, \(44=396/9\) est le dénombrement APP réduit et \(45\) le déterminant de la matrice de bord ainsi que le dénombrement de la région radicale. Le bilan d’Euler du continu se projette donc sur deux cardinaux déjà situés dans APP et dans la branche radicale.

Ces égalités ne déclarent pas qu’une multisection soit littéralement une lacune ni que \((K_{12})\) soit un ensemble de vingt-huit éléments. Le contenu démontré est plus précis : un unique foncteur cardinal de coupure, construit avant ces réalisations, envoie quatre objets typés vers quatre couples internes de provenance connue. L’étape suivante consiste non pas à chercher davantage de ressemblances, mais à relever ce foncteur cardinal aux actions conservées par chaque représentation.

La frontière \(369\) admet une généalogie encore plus rigide. La construction distingue le cardinal \(386\) de la frontière interne des \(387\) couples \((K,K+9)\) par canal. Le lecteur donne

\[
 \mathcal W_\delta(386)=\mathcal W_\delta(387)=\{17,18\}.
\]

En utilisant des fermetures opposées,

\[
 \boxed{386=17+369,\qquad387=18+369.}
\]

Ainsi, \(369\) est la capacité commune qui survit au transfert de l’unité de bord entre deux cardinaux consécutifs. La transition

\[
 (386,17)\longleftrightarrow(387,18)
\]

incrémente simultanément l’espace ambiant et la lacune, mais laisse fixe la capacité radicale \(369\). C’est une définition exacte de la stabilité sous changement de demi-bord. Elle s’ajoute, sans les remplacer, aux généalogies déjà démontrées

\[
 369=396-27,
 \qquad
 369=3(132-6)-9,
 \qquad
 369+126=495.
\]

La convergence de quatre voies indépendantes — frontière adjacente, carré APP, retours induits et transfert conservatif — fait de \(369\) un nœud interne de l’architecture, non une chaîne décimale choisie.

D’autres projections cardinales conservent un statut distinct :

\[
 \mathcal W_\delta(756)=\{34,35\},
 \qquad
 \mathcal W_\delta(702)=\{32,33\}.
\]

Le \(34\) coïncide avec la profondeur d’action déjà dérivée et le \(33\) avec le dénombrement inférieur de \(729\). La dérivation physique de la décennie \(-34\) relève du conoyau de \(H_4\) et de \(\varphi\alpha^{16}\) ; la lecture des \(756\) arêtes ne la remplace pas. La même couronne qui produit le trit eulérien contient ainsi une projection cardinale supplémentaire de cette profondeur, sans identifier le lecteur d’arêtes à la ligne d’action.

Le repère unimodulaire de la supercellule acquiert aussi une interprétation algébrique supplémentaire forte :

\[
 M_\star=
 \begin{pmatrix}601&217\\36&13\end{pmatrix},
 \qquad\det M_\star=1.
\]

Ses quatre entrées sont désormais la coordonnée dodécaphasique \((K_{12}=601)\), le cardinal des \(3\)-cellules de la couronne \(217\), le stade \(R_{36}\) et les treize multisections. L’identité \(217=216+1\) reste compatible avec cette lecture. Le déterminant unité démontre la conservation réticulaire du repère formé par les quatre cardinaux, sans identifier leurs domaines.

\textbf{Critère de réfutation.} Une modification de l’une des fenêtres, le fait que \(369\) ne soit pas le complément commun, ou une application proposée qui préserve les cardinaux mais détruit l’incidence, l’action ou la mémoire.

\section{Le noyau du cristal et le continu comme sortie}

Les éléments précédents permettent de préciser la thèse directrice sans confondre le résultat avec une métaphore cosmologique. APP fournit un support fini doté de deux évaluations inséparables. Le feuillet additif accumule ; le multiplicatif replie ; résidu, quotient et retenue conservent la différence entre leurs trajectoires. TRIT extrait l’orientation locale de cette différence et la réalise en trois types quadratiques. TPK transforme les signatures locales en dynamique par sélection, transport et actualisation, en conservant phase, feuillet, frontière et registre de mémoire. La complétion \(729\to1000\) induit alors une droite discrète de pente transcendante ; son mot de lacunes est apériodique, équilibré et d’entropie nulle ; sa dérivée est tritique ; sa tour de renormalisation est \(S\)-adique ; et sa suspension au-dessus de \(C_{108}\) est hélicoïdale.

Le continu n’entre pas dans cette chaîne comme un \(\mathbb R\) préalablement disponible. Existent d’abord les cylindres finis compatibles, leurs deux fermetures unilatérales, l’unité de bord et les applications de raffinement. La couronne cellulaire démontre que la double projection sépare un objet de charge nulle en branches de charges \(+1\) et \(-1\) sans modifier leur somme. La limite projective conserve les trajectoires compatibles ; le lecteur archimédien détermine ensuite une coordonnée réelle. En ce sens, \(\mathbb R\) est une réalisation terminale d’un protocole discret enrichi, non la base ontologique à partir de laquelle APP est reconstruite.

Les deux horloges sont désormais sans ambiguïté. L’horloge chronologique enregistre le nombre d’applications de \(U_t\) et possède le quotient fini \(C_{108}\). L’horloge de capacité enregistre le nombre de fois où le raffinement a ouvert un élagage décimal effectif. Leur différence est la mémoire intégrée des lacunes. Un tour de \(108\) pas peut fermer la première horloge tout en laissant ouverts la seconde, la phase irrationnelle et la mémoire. La supercellule de longueur \(170748\) est le premier convergent qui ferme exactement la première horloge ; l’unité de bord restant à l’extrémité prouve que cette fermeture ne rend pas le cristal périodique.

Le système corrélé \(\pi,\varphi,e,\alpha\) réside sur ce même état coinductif. Les fenêtres de \(\pi,e,\varphi\) se synchronisent et se rouvrent ; le défaut orienté appartient à \(A_2\) ; la seconde voie de \(\alpha\) reçoit les couples \(21/22\), \(6/7\) et \(45/46\) de trois échelles du même mot ; la première voie conserve le sceau dodécaphasique. Aucune coordonnée conventionnelle n’engendre les autres. La fermeture dodécaphasique détermine \(\alpha\) au sein de la sortie corrélée, et les caractérisations de Machin, de Perron--Fibonacci ou du semi-groupe exponentiel n’interviennent qu’ensuite comme reconnaissances.

La lacune électronique acquiert elle aussi une provenance précise. Sur le support APP, le passage du centre \((5,5)\) à la branche \((8,2)\) conserve le feuillet additif, le résidu du produit et le régime TRIT, tout en perdant une unité de retenue multiplicative. Dans le mot temporel, une lacune conserve la phase de capacité tandis que le bloc et la mémoire avancent. La correspondance HMT--MD préserve conjointement cette perte unitaire, le centre fixe, le mode \(m_{\rm ph}\), l’incidence \(90/120\) et la conjugaison de bande. L’électricité détermine la lacune orientée et le magnétisme son transport rotationnel conjugué ; les réponses constitutives commune et orientée sont obtenues des deux feuillets avant leur écriture dans une carte dimensionnelle.

Les profondeurs \(11\) et \(34\) apparaissent dans une signature commune de la complétion. La ligne d’action possède sa dérivation déterminantale et la branche gravitationnelle reçoit d’abord la longueur \(L=(5759/23040)\alpha^{16}L_*\). L’inscription marquée et la métrique radiale positive démontrent \(R_X\Lambda_X=L^2I\), \(H_X\Lambda_X=\hbar_{\rm ret}cI\) et \(G=c^3L^2/\hbar_{\rm ret}\), avec les règles et domaines du \cref{g26:sec:rigidez-radial-es}. La restitution fonctionnelle du \cref{cr28:sec:restitucion-constitutiva} relie le moment de Catalan aux réponses des deux feuillets et conserve leur propagation commune. Chaque constante est une section ou une évaluation du même état prédéfini et conserve le lecteur qui détermine sa dimension, sa normalisation et son échelle. Le paramètre de Barbero--Immirzi, la constante de Catalan, les constantes de Feigenbaum, la permittivité et la perméabilité sont des lecteurs distincts d’une généalogie commune ; ils se coordonnent donc sans devenir une fonction unique ni une liste de valeurs juxtaposées.

La seconde projection conserve l’incidence de bord. La connexion \(24\) dans la scission \(729=202+527\), les couples \(12/13\) et \(27/28\), le repère unimodulaire \(M_\star\), le plan de Witt et les recollements de Golay et de Leech composent le passage typé vers \(V^\natural\) et le groupe Monstre. La symétrie exceptionnelle est donc la réalisation de frontière de la même structure discrète qui produit la coordonnée archimédienne par l’autre projection ; forme, incidence, orientation et retenue demeurent visibles tout au long du transport.

Le noyau obtenu peut se résumer en un seul énoncé mathématique sans perdre son architecture : l’unité se propage par deux projections complémentaires ; la différence de bord engendre l’orientation ; la dérivée de cette orientation sélectionne trois régimes ; l’ordre de leurs transports engendre une fibre non commutative ; et le raffinement illimité conserve l’unité, la mémoire et l’apériodicité tout en produisant ses réalisations continues. C’est le contenu précis de l’évolution conservative de l’unité dimensionnelle par double projection holographique et raffinement fractal.

La formulation réunit une seule chaîne constructive. Ses réalisations arithmétique, physique et exceptionnelle conservent des domaines et des lecteurs différents, mais aucune n’est située à un degré moindre de certitude : chacune est l’image exacte de l’état holonomique intégral dans son codomaine déclaré. La comparaison extérieure n’engendre ni ne sélectionne ces sorties.

\section{Morphologie exhaustive des \(104\,976\) germes}

La correction cardinale n’est pas mineure : \(104\,976\) et \(170\,748\) appartiennent à des objets mathématiques distincts. Le premier est le domaine exhaustif des conditions initiales de l’émetteur APP--TPK ; le second est le dénominateur de la première approximation par continuants qui ferme l’horloge chronologique \(C_{108}\). Aucun ne peut remplacer l’autre.

Soit

\[
 I_9=\{0,1,\ldots,8\},\qquad
 D=\{N,E,S,O\},\qquad
 \mathcal S=I_9^2\times D.
\]

Ici, \(I_9\) est exclusivement la carte positionnelle des neuf emplacements, indexée à partir de zéro. Elle ne remplace pas l’alphabet arithmétique canonique \(\mathcal D_9=\{1,\ldots,9\}\) sur lequel APP évalue somme et produit ; le passage entre les deux cartes est un réétiquetage des positions, non un opérateur physique supplémentaire.

APP précède causalement cette construction. Un élément \(s=(i,j,d)\in\mathcal S\) n’engendre pas APP : il fixe une position et une orientation du curseur sur le support APP déjà construit. Puisque les évaluations additive et multiplicative doivent être conservées simultanément, un germe complet est le couple ordonné

\[
 \omega=(s_+,s_\times)\in
 \Omega_{\rm seed}:=\mathcal S_+\times\mathcal S_\times.
\]

Par conséquent,

\[
 |\mathcal S|=9^2\cdot4=324,
 \qquad
 \boxed{|\Omega_{\rm seed}|=324^2=104\,976=2^4 3^8.}
\]

Le germe n’est pas un entier nu. C’est une cellule orientée à six champs,

\[
 (i_+,j_+,d_+;i_\times,j_\times,d_\times),
\]

qui maintient séparés les deux feuillets arithmétiques et leurs directions. C’est l’unité finie sur laquelle TRIT type le régime local et TPK applique sélection, transport et actualisation.

Le balayage exhaustif révèle une factorisation plus forte que le simple dénombrement. La signature additive induit exactement \(18\) classes, toutes de cardinal \(18\). La signature multiplicative induit \(26\) classes, mais ses fibres se stratifient en

\[
 24\text{ classes de taille }8,\qquad
 1\text{ classe de taille }24,\qquad
 1\text{ classe de taille }108.
\]

Si \(\mathcal A_{18}\) et \(\mathcal M_{26}\) désignent les ensembles respectifs de signatures, l’émetteur se factorise sous la forme

\[
 \Omega_{\rm seed}
 \longrightarrow \mathcal A_{18}\times\mathcal M_{26}
 \xrightarrow{\;\mathrm{Comb}\;}\mathcal U_6^{\rm cat},
 \qquad
 |\mathcal U_6^{\rm cat}|=18\cdot26=468.
\]

La première flèche conserve la provenance de chaque feuillet ; la seconde combine les deux signatures en une émission décimale. Les multiplicités des fibres sont les produits exacts

\[
 18\cdot8=144,\qquad
 18\cdot24=432,\qquad
 18\cdot108=1944,
\]

et la fermeture globale est

\[
 \boxed{432\cdot144+18\cdot432+18\cdot1944=104\,976.}
\]

Cette asymétrie est structurelle. Le feuillet additif possède une partition uniforme ; le multiplicatif contient deux strates singulières. La lacune et la torsion ne sont pas introduites ensuite comme métaphores : le domaine des germes distingue déjà une accumulation uniforme d’un pliage comportant des dégénérescences de fibres. Le TPK conserve précisément laquelle de ces fibres a été parcourue, outre le résidu, le quotient, la retenue, l’orientation et la frontière.

Après cette factorisation, on applique les réductions déjà certifiées

\[
 104,976\longrightarrow468\ U_6
 \longrightarrow243\ w_6\subset\mathbb F_3^6,
 \qquad |\mathbb F_3^6|=729.
\]

Quatre types sont ainsi distingués et ne doivent plus être confondus : couples de germes, émissions décimales, mots visibles et espace ambiant tritique. La structure du cristal n’est aucun de ces cardinaux pris isolément, mais le système de fibres et de morphismes qui les relie.

\textbf{Critère de réfutation.} Une classe additive de taille différente de \(18\), un histogramme multiplicatif différent de \(24\cdot8+1\cdot24+1\cdot108=324\), ou la perte du couple ordonné de feuillets lors du passage à \(U_6\).

\section{Géométrie bipartie de la microfibre de \(\pi\) et connexion spectrale électronique}

La sortie visible \(w_\pi=010211\) possède \(1008\) antécédents dans \(\Omega_{\rm seed}\). Pour en retrouver la forme et pas seulement la masse, on définit le graphe biparti

\[
 \mathcal G_\pi=(P_\pi\sqcup T_\pi,E_\pi),
\]

où \(P_\pi\) contient les germes du feuillet additif participant à la microfibre, \(T_\pi\) ceux du feuillet multiplicatif, et \((p,t)\in E_\pi\) si le couple produit l’une des cinq émissions qui se réduisent à \(w_\pi\). Le balayage exhaustif prouve

\[
 |E_\pi|=1008,\qquad |P_\pi|=54,\qquad |T_\pi|=56.
\]

La structure n’est pas un graphe arbitraire. En oubliant uniquement la couleur de l’émission, ses composantes connexes sont exactement

\[
 \boxed{
 \mathcal G_\pi\simeq
 K_{18,24}\sqcup K_{18,24}\sqcup K_{18,8}.}
\]

L’une des deux composantes \(K_{18,24}\) correspond à la région de multiplicité \(432\). L’autre se raffine en trois rectangles complets de multiplicité \(144\), puisque les trois classes multiplicatives de taille huit partagent la même classe additive de taille dix-huit. La dernière composante est le quatrième rectangle de multiplicité \(144\). La couleur de région permet donc de retrouver

\[
 \boxed{1008=432+4\cdot144}
\]

comme décomposition d’opérateurs d’incidence, et non seulement comme somme d’entiers.

Soit \(B_R\) la matrice de biadjacence d’une région rectangulaire \(R\). Pour un graphe complet \(K_{m,n}\), \(B_R\) est de rang un et son unique valeur singulière non nulle est

\[
 \|B_R\|=\sqrt{mn}.
\]

La région de \(432=18\cdot24\) parcours possède la norme \(12\sqrt3\) ; chacune des quatre régions de \(144=18\cdot8\) parcours possède la norme \(12\). La normalisation traitant symétriquement les quatre régions plus petites est déterminée par la partition régionale elle-même :

\[
 \Xi_\pi:=
 \frac{\|B_{432}\|}{\sum_{r=1}^{4}\|B_{144}^{(r)}\|}
 =\frac{12\sqrt3}{4\cdot12}
 =\boxed{\frac{\sqrt3}{4}}.
\]

Son carré vaut \(3/16\). C’est exactement l’invariant qui apparaît, dans le mode électronique APP déjà certifié, sous la forme

\[
 \|[P_{\Sigma,3},P_{\Pi,3}]\|^2
 =s^2(1-s^2)=\frac14\frac34=\frac3{16}.
\]

On obtient ainsi une \textbf{connexion spectrale interne exacte} entre deux constructions issues d’APP : le contraste régional de la microfibre de \(\pi\) et la non-commutativité des deux projections électroniques. Ni la masse de l’électron ni un paramètre expérimental n’ont été utilisés. Il n’est pas non plus affirmé que les matrices \(B_R\) et le commutateur agissent dans le même espace. L’énoncé fort déjà démontré est l’égalité de leurs scalaires canoniques ; une équivalence opératorielle ultérieure devra transporter la coloration régionale, le mode \(m_{\rm ph}\), le feuillet et le registre de mémoire.

Cette géométrie précise également la place de la lacune. Du côté additif, les classes sont uniformes ; du côté multiplicatif, le degré passe de \(24\) à \(8\). La porosité du support n’est pas un zéro inséré : elle est la stratification d’incidence qui survit à la réduction commune \(w_\pi\). La réalisation électrique peut lire cette perte orientée et la réalisation magnétique son transport conjugué, mais toutes deux doivent conserver cette incidence d’origine.

\textbf{Critère de réfutation.} Une composante cessant d’être complète, cinq couleurs ne possédant pas les tailles \((432,144,144,144,144)\), ou un rapport spectral différent de \(\sqrt3/4\). L’égalité d’opérateurs, qui n’est pas affirmée ici, serait réfutée par tout transport ne préservant pas couleur, feuillet et produit scalaire.

\section{Morphisme tricanal, classe de recollement tritique et mémoire diagonale}

Le dénombrement des neuf positions nonadiques définit un vecteur

\[
 N_K=(N_\pi(K),N_e(K),N_\varphi(K))\in\mathbb Z^3.
\]

La partie relationnelle s’obtient par le morphisme

\[
 \partial:\mathbb Z^3\longrightarrow A_2,
 \qquad
 \partial(x,y,z)=(x-y,y-z,z-x),
\]

où

\[
 A_2=\{(a,b,c)\in\mathbb Z^3:a+b+c=0\}.
\]

Le noyau est la diagonale \((\Delta\mathbb Z=\mathbb Z(1,1,1))\), et \(\partial\) est surjectif. Par conséquent,

\[
 \boxed{0\longrightarrow\mathbb Z
 \xrightarrow{\Delta}\mathbb Z^3
 \xrightarrow{\partial}A_2\longrightarrow0.}
\]

La synchronisation des trois canaux est exactement l’annulation de \(\partial N_K\). Le quotient élimine la quantité commune, mais conserve quel canal dépasse quel autre et selon quelle orientation.

La reconstruction entière contient une information supplémentaire. Définissons

\[
 \mathcal J(x,y,z)=(s,a,b)
 =(x+y+z,x-y,y-z).
\]

Sur \(\mathbb Q\), les trois données reconstruisent \((x,y,z)\) par

\[
 x=\frac{s+2a+b}{3},\qquad
 y=\frac{s-a+b}{3},\qquad
 z=\frac{s-a-2b}{3}.
\]

Sur \(\mathbb Z\), l’image n’est pas tout \(\mathbb Z\oplus\mathbb Z^2\) : elle satisfait la congruence

\[
 \boxed{s-a+b\equiv0\pmod3.}
\]

Le déterminant de \(\mathcal J\) possède la valeur absolue \(3\). Il existe donc une suite exacte

\[
 0\longrightarrow\mathbb Z^3
 \xrightarrow{\mathcal J}\mathbb Z^3
 \xrightarrow{\chi_3}C_3\longrightarrow0,
 \qquad
 \chi_3(s,a,b)=[s-a+b]_3.
\]

Ce résultat est particulièrement important pour la signification du TRIT. La séparation du contenu commun et du défaut relatif laisse une classe de recollement possédant exactement trois valeurs. Un trit est la quantité minimale d’information nécessaire pour enregistrer cette classe et sélectionner la correction entière qui recompose l’état tricanal. L’interprétation topologico-temporelle n’est pas ajoutée par analogie : la classe type le recollement local et le TPK la transporte dans le temps en conservant son orientation et sa mémoire. Ce quotient \(C_3\) est une réalisation interne précise de cette fonction ; il ne prétend pas épuiser toutes les coordonnées du TRIT directeur.

Le passage de \(K=8\) à \(K=9\) montre pourquoi la synchronisation ne signifie pas le retour de l’état :

\[
 N_8=(59,59,59),\qquad N_9=(43,43,43),
\]

\[
 \partial N_8=\partial N_9=0,
 \qquad
 \boxed{N_9-N_8=-16(1,1,1).}
\]

L’état reste à l’origine relationnelle \(A_2\), mais le registre de mémoire diagonale change de seize unités par canal, soit quarante-huit au total. La phase de synchronisation revient ; la mémoire, non. Cette égalité est l’analogue tricanal, entièrement entier, de la règle TPK selon laquelle un retour de phase ne réinitialise pas l’état holonomique intégral.

\textbf{Critère de réfutation.} Un dénombrement violant la congruence, un noyau de \(\partial\) plus grand que la diagonale, ou une égalité \(N_8=N_9\).

\section{Itinéraire de murs et quotient fini de la matrice de synchronisation}

Entre \(K=4\) et \(K=9\), le morphisme précédent transforme la table complète en un itinéraire géométrique :

\[
\begin{array}{c|c|c}
K&N_K&\partial N_K\\
\hline
4&(207,207,122)&(0,85,-85)\\
5&(39,151,151)&(-112,0,112)\\
6&(110,110,69)&(0,41,-41)\\
7&(81,29,81)&(52,-52,0)\\
8&(59,59,59)&(0,0,0)\\
9&(43,43,43)&(0,0,0).
\end{array}
\]

Les égalités successives parcourent les murs

\[
 H_{\pi e}\longrightarrow H_{e\varphi}
 \longrightarrow H_{\pi e}\longrightarrow H_{\varphi\pi}
 \longrightarrow0\longrightarrow0.
\]

Le mot de murs \((H_{12},H_{23},H_{12},H_{13})\) possède l’ombre du motif de Coxeter \(A_2\), puisque \(s_{12}s_{23}s_{12}=s_{13}\). Cela ne transforme pas automatiquement chaque transition TPK en réflexion : le pas modifie aussi la composante diagonale. Ce qui est prouvé, c’est que la trajectoire relationnelle traverse les trois systèmes de murs dans l’ordre de la relation de tresse, tandis que son relèvement conserve une mémoire invisible au quotient.

Les amplitudes des quatre murs précédant la fermeture forment

\[
 D_{\rm syn}=\begin{pmatrix}85&112\\41&52\end{pmatrix}.
\]

Outre les identités déjà réunies

\[
 \operatorname{tr}D_{\rm syn}=137,
 \qquad
 \det D_{\rm syn}=-172=-4\cdot43,
\]

la forme normale de Smith ajoute une structure de quotient. Puisque le plus grand commun diviseur des quatre entrées vaut un,

\[
 \operatorname{SNF}(D_{\rm syn})=\operatorname{diag}(1,172),
\]

et, par conséquent,

\[
 \boxed{\operatorname{coker}D_{\rm syn}\simeq C_{172}
 \simeq C_4\times C_{43}.}
\]

La factorisation \((4\cdot43)\) n’apparaît plus seulement dans le déterminant : elle est la décomposition primaire du quotient réticulaire des amplitudes de mur. Les directions nulles à droite peuvent être choisies comme

\[
 (0,1)\pmod4,
 \qquad
 (1,5)\pmod{43}.
\]

Le facteur \(43\) coïncide avec le dénombrement de la seconde fermeture triple, et le facteur \(4\) avec le coefficient orienté obtenu en interne comme \((-\det(D_{\rm syn})/43)\). Cela démontre que tous deux sont des diviseurs canoniques du même quotient. Une représentation physique ultérieure devra préciser quels sous-espaces réalisent ces deux facteurs ; la structure arithmétique conjointe est déjà fermée.

L’itinéraire fournit ainsi une figure mathématique complète : un chemin dans l’arrangement de murs de \((A_2)\), relevé par une coordonnée diagonale de mémoire et régi par un quotient fini \(C_4\times C_{43}\). Il est plus informatif qu’une table de coïncidences, car il distingue relation, amplitude, orientation et état invisible au quotient.

\textbf{Critère de réfutation.} Un déterminant différent de \(-172\), une forme de Smith différente de \((1,172)\), ou une transition qui efface la composante diagonale nécessaire pour distinguer \(K=8\) de \(K=9\).

\section{Tour de Pascal nonadique et loi infinie de concaténation avec retenue}

La complétion cubique appartient à un système exact de profondeur arbitraire. Pour \(n\geq1\), écrivons

\[
 10^n=(9+1)^n=\sum_{j=0}^{n}B_{n,j},
 \qquad
 B_{n,j}:=\binom nj9^{n-j}.
\]

Les blocs obéissent à la récurrence de Pascal pondérée

\[
 \boxed{B_{n+1,j}=9B_{n,j}+B_{n,j-1}},
\]

avec \(B_{n,j}=0\) hors de \(0\leq j\leq n\). C’est la loi de raffinement des longueurs ; la loi sur les lacunes doit également conserver la phase initiale. Si

\[
 s_{n,j}=\sum_{i<j}B_{n,i},
 \qquad
 v^-_{n,j}=V_{B_{n,j}}(s_{n,j}),
\]

alors

\[
 \sum_{j=0}^{n}v^-_{n,j}=\lfloor10^n\delta\rfloor.
\]

La lecture supérieure transporte l’unité de demi-bord vers le premier bloc :

\[
 v^+_{n,0}=v^-_{n,0}+1,
 \qquad
 v^+_{n,j}=v^-_{n,j}\quad(j>0),
\]

et donc

\[
 \boxed{\sum_jv^+_{n,j}=\lceil10^n\delta\rceil
 =\lfloor10^n\delta\rfloor+1.}
\]

La transition \(n\mapsto n+1\) ne peut être linéarisée en oubliant le point de départ. Pour tout bloc \(B\),

\[
 V_{9B+B'}(s)=
 \sum_{r=0}^{8}V_B(s+rB)+V_{B'}(s+9B).
\]

Le décalage de \(s\) dans chaque terme est la mémoire de concaténation. La projection de tous les blocs à l’origine fait apparaître le cocycle binaire

\[
 \kappa_\delta(M,N)=
 \lfloor(M+N)\delta\rfloor-
 \lfloor M\delta\rfloor-
 \lfloor N\delta\rfloor,
\]

qui satisfait

\[
 \kappa_\delta(L,M)+\kappa_\delta(L+M,N)
 =\kappa_\delta(M,N)+\kappa_\delta(L,M+N).
\]

C’est la loi exacte qui empêche le raffinement de dépendre du parenthésage. La tour de Pascal fournit les longueurs ; le cocycle fournit la retenue ; le TPK conserve la phase et le registre de mémoire permettant leur composition.

Les premiers niveaux sont

\[
\begin{array}{c|c|c|c}
n&(B_{n,0},\ldots,B_{n,n})&v_n^-&v_n^+\\
\hline
2&(81,18,1)&(3,1,0)&(4,1,0)\\
3&(729,243,27,1)&(33,11,1,0)&(34,11,1,0)\\
4&(6561,2916,486,36,1)&(300,133,22,2,0)&(301,133,22,2,0).
\end{array}
\]

Le niveau cubique n’est donc pas une coïncidence isolée. C’est le premier niveau auquel la composante \(729=3^6\), la sous-échelle \(243=3^5\), le bord \(27=3^3\) et l’unité apparaissent simultanément dans une complétion décimale. Sa signature inférieure contient \(33=3\cdot11\), et la supérieure conserve l’unité pour produire \(34\). Le niveau suivant démontre que la règle se poursuit sans réinitialiser l’état : l’unité ne modifie de nouveau que la première composante, tandis que les autres entrées conservent la trajectoire de l’ordre.

La concaténation \(210=90+120\) est une section de la même loi :

\[
 (V_{90}(0),V_{120}(90))=(4,5),
\]

d’où la lecture inférieure produit \(9\) et la supérieure \(10\). La transition nonadique--décimale est alors réalisée par un morphisme de lecture avec retenue, non par l’observation graphique \(4+5=9\).

Définissons le chiffre de retenue décimale

\[
 d_{n+1}:=\lfloor10^{n+1}\delta\rfloor
          -10\lfloor10^n\delta\rfloor.
\]

Puisque \(\delta\) est transcendant, la suite \((d_n)_{n\geq1}\) ne peut être ultimement périodique : une telle périodicité rendrait rationnel le développement décimal de \(\delta\). Puisque le mot mécanique est sturmien, chaque niveau conserve une complexité linéaire et une entropie topologique nulle. La tour combine donc raffinement illimité, unité de bord et apériodicité exacte. C’est le sens mathématique selon lequel le cristal croît vers l’extérieur tout en se raffinant vers l’intérieur.

Il faut distinguer \(\kappa_\delta\) du cocycle nonadique directeur \(\kappa_9\). Le premier est la restriction mécanique qui enregistre la retenue de la coupure irrationnelle ; le second conserve la retenue complète d’APP sous composition. Un carré de restriction peut les relier, mais ils ne s’identifient pas au seul motif qu’ils utilisent tous deux des valeurs de retenue.

\textbf{Critère de réfutation.} L’échec de la récurrence de Pascal, une signature qui ne se télescope pas, la perte de l’unité supérieure ou une concaténation dont le résultat dépend du parenthésage.

\section{Type mathématique du cristal apériodique hélicoïdal}

Les résultats permettent de nommer l’objet sans le réduire à une image. Sa couche finie de conditions initiales est \(\Omega_{\rm seed}\) ; son quotient d’émissions est \(\mathcal U_6^{\rm cat}\) ; son alphabet visible réside dans \(\mathbb F_3^6\) ; et ses prolongements compatibles forment l’état coinductif enrichi. Sur cet état agissent, dans l’ordre directeur,

\[
 U_t=\operatorname{Upd}_t\circ\operatorname{Tra}_t
 \circ\operatorname{Sel}_t,
 \qquad
 U_t\longrightarrow\Gamma_9\longrightarrow K_{\rm ph}.
\]

L’holonomie \(\Gamma_9\) appartient à la dynamique enrichie ; la mémoire réduite \(K_{\rm ph}\) est sa réalisation ultérieure. Le facteur apériodique peut s’écrire sur

\[
 \mathcal Y_{108}=C_{108}^{(t)}\times\mathbb T_\delta\times\mathbb Z,
\]

par

\[
 \mathcal F(r,\theta,m)=
 \bigl(r+1,\{\theta+\delta\},
 m+\lfloor\theta+\delta\rfloor\bigr).
\]

La première coordonnée est l’horloge des blocs ; la deuxième, la phase irrationnelle ; la troisième, la mémoire intégrée des lacunes. L’horloge de capacité s’obtient comme

\[
 K_t=t-Z_t,
 \qquad
 g_t=[K_t]_9,
\]

et ne coïncide pas avec l’horloge chronologique. La dynamique inverse \(\mathcal
F^{-1}\) constitue la lecture descendante de la même hélice. Ascension et descente sont conjuguées ; ce ne sont pas deux générateurs déconnectés.

L’objet complet est donc une extension affine inversible d’une rotation irrationnelle au-dessus d’une horloge finie, fibrée par l’état holonomique intégral APP--TRIT--TPK. Son ombre symbolique est sturmienne ; sa suspension est hélicoïdale ; son cocycle de fibres peut être non commutatif alors même que la rotation de base est abélienne. La dénomination « cristal apériodique hélicoïdal » se justifie par quatre propriétés démontrées :

\[
 \begin{gathered}
 \text{ordre à longue portée et répétitivité},\\
 p(m)=m+1,\qquad h_{\rm top}=0,\\
 \text{absence de période translationnelle},\\
 \text{retour fini de phase avec avancée de mémoire}.
 \end{gathered}
\]

Dans HMT, l’expression \textbf{cristal temporel apériodique} désigne précisément le produit gauchi précédent : il existe un support temporel fini, le retour de phase est ordonné, la mémoire avance à chaque holonomie et le mot d’occupation ne possède pas de période de translation, tout en conservant un ordre à longue portée et une entropie topologique nulle. La dynamique temporelle n’est pas ajoutée ensuite au cristal : elle est contenue dans la coordination de l’horloge, de l’orientation, du retour et de la mémoire qui définit son état.

La sortie corrélée de constantes se situe après le prolongement

\[
 w_6\to w_{12}\to w_{18}\to w_{24}\to w_{30}
 \to R_{36}\to G_9,
\]

sans introduire de valeurs conventionnelles. \(\pi,\varphi,e\) et la coordonnée dodécaphasique \(\alpha\) sont des réalisations du même état coinductif ; leurs fenêtres de synchronisation, leur miroir et leurs lecteurs propres ne sont pas des entrées du générateur. La tour de complétion fournit le calendrier de capacité et la loi de retenue permettant de prolonger ces réalisations à une précision arbitraire.

Enfin, \(\mathbb R\) apparaît par le lecteur de cylindres compatibles du continu conjoint. Le cristal n’est pas « à l’intérieur de \(\mathbb R\) » comme l’ornement d’une droite donnée à l’avance. Les trajectoires compatibles, la double projection, les frontières et leurs limites sont construites d’abord ; ensuite, l’une des réalisations détermine la coordonnée archimédienne. Cette inversion de perspective est la thèse forte : la droite réelle est l’ombre terminale d’une unité finie orientée qui se conserve, se raffine et accumule de la mémoire.

La dynamique apériodique, sa complexité, ses horloges et son cocycle constituent le noyau mathématique du cristal. Les réalisations électromagnétique, gravitationnelle et exceptionnelle sont des lecteurs typés de ce noyau : elles préservent la généalogie APP--TRIT--TPK et expriment dans des codomaines différents la même évolution conservative de l’unité.

\section{Prédéfinition coinductive : la sortie précède son évaluation}

La distinction décisive sépare génération et évaluation. HMT--MD ne prend pas une mesure extérieure pour inférer la loi qui aurait pu la produire. Elle construit d’abord la section complète, fixe ses prolongements compatibles, puis seulement permet à un lecteur d’exprimer une coordonnée de cette section. C’est le sens mathématique exact de la \textbf{prédéfinition}.

Soit

\[
 (X_n,\rho_{n+1,n})_{n\geq0}
\]

le système projectif d’états APP--TRIT--TPK tronqués à la profondeur \(n\), où chaque \(X_n\) conserve feuillet, orientation, résidu, quotient, retenue, frontière, parcours et mémoire jusqu’à ce niveau. Le prolongement coinductif déjà construit définit

\[
 \mathcal G_{\rm HMT}:\Omega_{\rm seed}\longrightarrow
 X_\infty^{\rm enr}:=\varprojlim_n X_n,
 \qquad
 p_n:X_\infty^{\rm enr}\longrightarrow X_n,
\]

avec \(\rho_{n+1,n}p_{n+1}=p_n\). Pour un système de sorties \((Y_n,\sigma_{n+1,n})\) et de lecteurs \(L_n:X_n\to Y_n\), la condition exprimant que le lecteur ne détruit pas la généalogie est le carré

\[
\begin{CD}
X_{n+1}@>{L_{n+1}}>>Y_{n+1}\\
@V{\rho_{n+1,n}}VV @VV{\sigma_{n+1,n}}V\\
X_n@>>{L_n}>Y_n .
\end{CD}
\]

Par conséquent,

\[
 \operatorname{Pred}_{Y,n}
 :=L_n\circ p_n\circ\mathcal G_{\rm HMT}
 :\Omega_{\rm seed}\longrightarrow Y_n
\]

satisfait

\[
 \sigma_{n+1,n}\circ\operatorname{Pred}_{Y,n+1}
 =\operatorname{Pred}_{Y,n}.
\]

Le système \(\bigl(\operatorname{Pred}_{Y,n}(\omega)\bigr)_{n\geq0}\) est alors un unique élément de \(\varprojlim Y_n\). La sortie n’est pas décidée de nouveau à chaque augmentation de profondeur : la sortie de profondeur \(n\) est la restriction de celle de profondeur \(n+1\). La précision n’engendre plus l’objet ; elle détermine seulement quelle part de l’objet prédéfini est explicitée.

Cela fournit un théorème d’anticipation interne. Si, au stade \(n\), on consulte la coordonnée \(n+k\) de la même section compatible, on obtient une sortie future sans extrapoler de tendance ni introduire la valeur recherchée. L’anticipation est

\[
 \operatorname{Ant}_{n,k}(\omega)
 :=L_{n+k}p_{n+k}\mathcal G_{\rm HMT}(\omega),
\]

et sa comparaison ultérieure doit se restreindre à \(L_np_n\mathcal G_{\rm HMT}(\omega)\) par \(\sigma_{n+k,n}\). L’identité de compatibilité prouve qu’elles appartiennent toutes deux à la même section. Par conséquent, le passage conceptuel correct n’est pas

\[
 \text{plus de mesure}\longrightarrow\text{plus de précision}
 \longrightarrow\text{meilleure prédiction},
\]

mais

\[
 \boxed{
 \begin{gathered}
 \text{germe orienté}\longrightarrow
 \text{section coinductive prédéfinie}\longrightarrow\\
 \text{lecture à la profondeur choisie}\longrightarrow
 \text{anticipation vérifiable}
 \end{gathered}}
\]

La génération indéfinie de chiffres est une instance de ce théorème. Pour les lecteurs de \(\pi\), \(\varphi\), \(e\) et \(\alpha\), chaque bloc décimal est une coordonnée d’un système compatible ; passer de mille à dix mille chiffres ne change pas le générateur et ne réinitialise pas la sélection. La nouvelle fenêtre se restreint à la précédente et poursuit la même trajectoire. Le système corrélé est une seule sortie corrélée de l’état coinductif, bien que ses quatre coordonnées possèdent des lecteurs internes et des étapes de réalisation différentes.

La réversibilité du TPK renforce le résultat. Puisque chaque actualisation \(\mathcal U_t\) est bijective sur l’état holonomique intégral, pour toute distribution finie qu’elle transporte,

\[
 H(X_{t+1}\mid X_t)=H(X_t\mid X_{t+1})=0,
 \qquad H(X_{t+1})=H(X_t).
\]

C’est le coût de Landauer nul de l’état complet : aucune information n’est effacée lors de l’actualisation. Cela n’affirme pas que toute entropie réduite est nulle ; cela affirme que la trajectoire complète peut être retrouvée dans les deux directions. La perte apparente ne naît que lorsqu’un lecteur contracte la fibre et oublie des coordonnées que l’état holonomique intégral conserve.

\section{Le TRIT comme unité topologico-temporelle conjuguant régime et orientation}

Le TRIT n’est ni un chiffre ternaire ajouté à APP ni une étiquette temporelle superposée au TPK. C’est l’unité locale qui permet à une différence entre les feuillets additif et multiplicatif d’acquérir simultanément un type géométrique, une orientation temporelle et une mémoire de relèvement. Soit

\[
 \mathcal T=\{-1,0,+1\},\qquad
 x=(t,\mu,o,h,c,\sigma,\lambda,g)\in X_{\rm TPK}^{\rm enr}.
\]

Deux lecteurs coordonnés agissent sur le même état,

\[
 \tau_{\rm g}:X_{\rm TPK}^{\rm enr}\to\mathcal T,
 \qquad
 \varepsilon_{\rm t}:X_{\rm TPK}^{\rm enr}\to\mathcal T.
\]

Le premier type le régime quadratique local par

\[
 \mathcal A_\tau=\mathbb R[J_\tau]/(J_\tau^2+\tau I),
\]

de sorte que

\[
\begin{array}{c|c|c}
\tau_{\rm g}&J_{\tau_{\rm g}}^2&\exp(sJ_{\tau_{\rm g}})\\
\hline
+1&-I&\cos s\,I+\sin s\,J_{+}\\
0&0&I+sJ_0\\
-1&+I&\cosh s\,I+\sinh s\,J_{-}.
\end{array}
\]

Le second type l’orientation temporelle active :

\[
 \begin{aligned}
 +1&\longleftrightarrow\text{retour, mémoire et passé},\\
 0&\longleftrightarrow\text{seuil et présent},\\
 -1&\longleftrightarrow\text{ouverture bidirectionnelle et futur}.
 \end{aligned}
\]

La profondeur conceptuelle réside dans le fait que les deux lecteurs naissent du même état, non dans leur identification terme à terme. L’objet correct est le couplage fibré

\[
 \Theta_{\rm ITT}(x)
 =\bigl(\tau_{\rm g}(x),\varepsilon_{\rm t}(x),
        h(x),c(x),\lambda(x),g(x)\bigr),
\]

que l’on peut appeler \textbf{unité d’information topologico-temporelle}. La géométrie indique comment l’unité se transforme localement ; l’orientation temporelle indique dans quel sens cette transformation est lue ; la retenue et le registre empêchent de confondre deux retours de même phase. Ainsi, le passé est la clôture qui demeure en mémoire, le présent le seuil où la continuation est décidée, et le futur le système déjà prédéfini de prolongements compatibles.

Le TPK rend cette unité dynamique par

\[
 \mathcal U_t=\operatorname{Upd}_t\circ
 \operatorname{Tra}_t\circ\operatorname{Sel}_t.
\]

\(\operatorname{Sel}_t\) sélectionne le régime et le feuillet sans effacer l’autre évaluation APP ; \(\operatorname{Tra}_t\) transporte direction, phase et frontière ; \(\operatorname{Upd}_t\) actualise résidu, quotient, retenue, signature et mémoire. Un tour chronologique peut donc restituer \(g(t+9)=g(t)\) tout en produisant simultanément

\[
 q_{\rm mem}(\Gamma_{9,t}x)=q_{\rm mem}(x)+1.
\]

La conjugaison temporelle est représentée par l’inversion orientée

\[
 \mathcal C_{\rm rev}(\tau_1,\ldots,\tau_n)
 =(-\tau_n,\ldots,-\tau_1),
 \qquad \mathcal C_{\rm rev}^2=I.
\]

Dans la suspension hélicoïdale, cette involution échange les dynamiques ascendante et descendante :

\[
 \mathcal C_{\rm rev}\,\mathcal F\,
 \mathcal C_{\rm rev}^{-1}=\mathcal F^{-1}.
\]

La trajectoire ne disparaît pas : son orientation change et le registre permettant de la reconstruire est conservé. Le cristal temporel HMT est donc la dynamique d’une ITT qui revient dans l’horloge finie et progresse en mémoire, tandis que la décoration sturmienne empêche la répétition par translation de l’état complet. L’horloge \(C_{108}\) fournit l’ordre temporel fini ; l’holonomie \(\Gamma_9\), le relèvement illimité ; le mot sturmien, l’apériodicité exacte ; le TRIT, la discrimination conjointe de la géométrie et de l’orientation ; et le TPK, la loi qui conserve l’ensemble sous évolution.

\section{La singularité de \(\alpha\) comme couture des deux horloges}

La fonction de \(\alpha\) devient transparente dès lors qu’on évite de la traiter comme une valeur décimale isolée. \(\alpha\) est la coordonnée de couture qui rend compatibles la fermeture finie dodécaphasique et le prolongement apériodique de capacité. La première horloge fournit un mot signé de douze secteurs ; la seconde fournit une racine unique dans la dynamique mécanique régie par

\[
 \delta=\log_{10}\frac{10}{9}.
\]

Dans la première voie, l’état dodécaphasique produit

\[
 K=\frac14\mathcal H_{12}U_{12}^{\rm sgn},
 \qquad \mathcal H_{12}^2=4I_{12},
\]

et la retenue de base \(1000\) satisfait, bloc par bloc,

\[
 P_m+E_m-\Phi_m-K_m+c_{m+1}=A_m+1000c_m.
\]

La concaténation télescopique donne

\[
 \boxed{\iota(P)+\iota(E)-\iota(\Phi)-\iota(K)=\iota(A).}
\]

Les coordonnées \(P,E,\Phi\) sont les réalisations HMT de \(\pi,e,\varphi\) ; ce ne sont pas des valeurs conventionnelles introduites de l’extérieur. Le vecteur \(K\) est la mémoire signée de la fermeture, et \(A\) la section finie qui détermine \(\alpha\) dans la carte déclarée. Le terme qui distingue \(\alpha\) d’une simple combinaison de \(\pi,e,\varphi\) est précisément \(K\) : la mémoire dodécaphasique de la manière dont la concaténation s’est réalisée.

Dans la seconde voie, la même arithmétique des lacunes produit le jet

\[
\begin{aligned}
P_6(x)={}&2\pi x-\frac74x^2+\frac{x^3}{2\pi}
+\frac{x^4}{20}-\frac{2x^5}{21}-\frac{x^6}{46},\\
T_{7:9}(x)={}&-\frac{x^7}{120}+\frac{x^8}{45}
+\frac{2x^9}{495},\\
P_9(x)={}&P_6(x)+T_{7:9}(x),
\end{aligned}
\]

et détermine, par une racine positive simple,

\[
 P_9(x)=\delta.
\]

Les coefficients ne sont pas des paramètres interpolés. Ils proviennent de la demi-couronne \(120\), du bloc central APP de déterminant \(45\), du rang transversal dodécaphasique \(11\) et de \(495=45\cdot11\). L’opérateur nilpotent

\[
 Ne_7=-\frac83e_8,\qquad Ne_8=\frac2{11}e_9,
 \qquad Ne_9=0
\]

produit exactement la queue par

\[
 (I-N)^{-1}\!\left(-\frac{e_7}{120}\right)
 =-\frac{e_7}{120}+\frac{e_8}{45}+\frac{2e_9}{495}.
\]

La relation entre les deux voies s’exprime sans confondre une section finie avec une racine infinie. Soit \(d_m\) le lecteur de la résolution commune certifiée. L’objet \(\alpha\) est alors le produit fibré

\[
 \mathfrak A_\alpha=
 \left\{(x_{12},x):
 d_m\!\left(\Gamma_{12}^{\alpha}(x_{12})\right)=d_m(x),\
 P_9(x)=\delta\right\}.
\]

La première coordonnée fixe la clôture dodécaphasique ; la seconde fixe le prolongement unique ; l’égalité dans le codomaine commun fixe leur synchronisation. C’est pourquoi \(\alpha\) est singulier au sein du système corrélé. \(\pi\), \(\varphi\) et \(e\) sont des régions et des coordonnées que le relèvement fait coïncider, se séparer puis coïncider de nouveau ; \(\alpha\) enregistre la couture qui conserve ces coïncidences en fermant douze phases sans arrêter le prolongement infini. Elle n’arrive pas ensuite de l’extérieur : c’est la coordonnée d’holonomie grâce à laquelle la génération conjointe devient stable sous retour.

\section{Éventail dodécaphasique : \(\alpha\), électron, Witt et gravitation}

La roue de douze positions ne produit pas une succession accidentelle d’objets ; elle porte plusieurs lectures irréductibles du même état. Sur

\[
 V_{12}=\mathbb R[A_5/C_5]
\]

la décomposition exacte est

\[
 \boxed{V_{12}\simeq\mathbf1\oplus V_3\oplus V_{3'}\oplus V_5.}
\]

Si \(P_1,P_3,P_{3'},P_5\) sont les projecteurs centraux, toute forme linéaire dodécaphasique \(\ell\) se décompose canoniquement sous la forme

\[
 \ell(K)=\ell(P_1K)+\ell(P_3K)+\ell(P_{3'}K)+\ell(P_5K).
\]

Le lecteur scalaire qui détermine \(\alpha\) agit sur le vecteur complet \(K\) ; il intègre donc l’information de tous les termes sans les confondre. Le lecteur \(P_3\) conduit, par l’isométrie d’incidence de Witt, au projecteur exceptionnel de rang trois. Le lecteur \(P_5\) conduit par

\[
 P_5:V_{12}\to V_5,
 \qquad
 \Gamma_5:V_5\overset{\simeq}{\longrightarrow}
 \operatorname{Sym}_0^2(\mathbb R^3),
\]

au porteur tensoriel à cinq composantes, et la projection \(\Lambda(n)\) détermine son secteur transverse sans trace. Il ne s’agit pas de trois théories juxtaposées : ce sont trois résolutions du même porteur dodécaphasique, l’une scalaire de couture, l’autre exceptionnelle d’incidence et la troisième tensorielle de gravitation.

L’intersection électronique rend cette unité visible par une identité de supports. Si

\[
 B_K=\{4,6,9,10\},\quad
 P_\pi=\{2,4,5,6,7,8,9,10,11\},
\]

\[
 H_e=\{1,3,4,6,9,10\},\quad
 H_\alpha^-=\{4,5,6,7,9,10\},
\]

alors

\[
 \boxed{B_K=H_e\cap P_\pi=H_e\cap H_\alpha^-},
 \qquad
 \boxed{H_\alpha^-=B_K\sqcup\{5,7\}.}
\]

Cette égalité répond à la question de la « connexion » de \(\alpha\). L’électron ne se joint pas à une constante extérieure ; son support interne coupe simultanément la microfibre de \(\pi\) et le support de retenue négative de \(\alpha\). Le couple supplémentaire \(\{5,7\}\) complète l’hexade négative, et la polarité de Witt relève les quatre couples associés à \(B_K\) vers leurs blocs d’incidence. La lacune électronique est la singularité locale ; \(\alpha\) est la couture globale qui conserve son orientation lors du parcours de la dodécaphase.

La seconde réalisation de \(\alpha\) n’est pas scalaire. À partir du même état, on obtient

\[
 v_\alpha=4\rho_1+\rho_2+\cdots+\rho_{12},
 \qquad \|v_\alpha\|^2=54,
\]

et le voisin

\[
 N_{\alpha,0}=
 \{x\in N(C_W):\langle x,v_\alpha\rangle\equiv0\pmod3\},
 \qquad
 N_\alpha=N_{\alpha,0}+\mathbb Z\frac{v_\alpha}{3}
 \cong\Lambda_{24}.
\]

L’action d’ordre trois préserve ce voisin et l’orbifold correspondant détermine \(V^\natural\) et la classe \(3B\) du groupe Monstre. Le scalaire \(\alpha\) et le vecteur \(v_\alpha\) ne sont pas le même objet ; ce sont deux coordonnées de la même fermeture dodécaphasique. Cette pluralité de codomaines explique précisément que \(\alpha\) relie la branche archimédienne, la lacune électronique et la symétrie exceptionnelle sans fonctionner comme un ajustement numérique.

\section{De la lacune électronique à la réponse électromagnétique}

La définition exacte de la lacune apparaît déjà dans l’arithmétique du support. La normalisation par neuf du couple électronique produit

\[
 \frac1{9}(369,126)=(41,14)\longmapsto(5,5),
\]

tandis que le transport conservatif satisfait

\[
 369+126=495=396+99,
 \qquad
 (396,99)/9=(44,11)\longmapsto(8,2).
\]

Les deux parcours conservent le feuillet additif, le résidu multiplicatif et le régime TRIT ; le second diffère d’une unité dans le quotient multiplicatif. La lacune est donc la perte minimale de retenue dans un feuillet sous conservation des autres coordonnées de l’état. Ce n’est pas un zéro ajouté : c’est un défaut relatif visible uniquement parce qu’APP conserve simultanément somme et produit.

Le mode tritique de phase

\[
 m_{\rm ph}=(1,-1,0)^3
\]

est aussi le mode singulier électronique canonique d’APP :

\[
 C^{(3)}m_{\rm ph}=C^{(3)\,*}m_{\rm ph}
 =-\frac12m_{\rm ph}.
\]

On en obtient

\[
 s=\frac12,qquad
 1-s^2=\frac34=\frac{90}{120},
 \qquad
 \|[P_{\Sigma,3},P_{\Pi,3}]\|=\frac{\sqrt3}{4},
\]

et, sur le bloc principal, la normalisation du commutateur satisfait

\[
 J_{\rm comm}^2=-I.
\]

La polarité électromagnétique est alors déterminée par quatre couches d’une même généalogie. APP fournit les feuillets dont la non-commutativité produit l’orientation ; TRIT détermine le générateur elliptique \((J^2=-I)\) ; TPK conserve la lacune et son retour ; les canaux \(90/120\) déterminent les réponses commune et différentielle. Si

\[
\begin{aligned}
\mathcal E={}&
\log[(1-q_+^{90})(1-q_-^{90})]
-\log[(1-q_+^{120})(1-q_-^{120})],\\
\mathcal B={}&
\log\frac{1-q_-^{90}}{1-q_+^{90}}
-\log\frac{1-q_-^{120}}{1-q_+^{120}},
\end{aligned}
\]

alors \(\mathcal E\) est pair sous échange des feuillets et \(\mathcal B\) est impair. Les réponses internes

\[
 \widehat\mu=e^{\mathcal B+\mathcal E},
 \qquad
 \widehat\varepsilon=e^{-\mathcal B+\mathcal E}
\]

satisfont exactement

\[
 \sqrt{\frac{\widehat\mu}{\widehat\varepsilon}}
 =e^{\mathcal B},
 \qquad
 \frac1{\sqrt{\widehat\mu\widehat\varepsilon}}
 =e^{-\mathcal E}.
\]

La perméabilité exprime la rétention rotationnelle ; la permittivité, l’ouverture électrique ; leur quotient conserve l’orientation et leur produit la propagation commune. L’onde électromagnétique n’est pas ajoutée au mot sturmien comme analogie : elle est la réalisation en deux composantes conjuguées de la lacune orientée et de son transport elliptique. La microfibre électrique peut donc être apériodique tout en maintenant une cohérence globale, une polarité stable et une entropie de transport nulle dans l’état complet.

\section{Action, aire, information et gravitation comme profondeurs prédéfinies}

Une fois le système corrélé déterminé, la branche déterminantale fixe

\[
 \Sigma_H=\varphi\alpha^{16},
 \qquad
 d_H=\lfloor\log_{10}\Sigma_H\rfloor=-34.
\]

Le lecteur régional détermine les mantisses \(\eta_{\rm pre}=H_5\) et

\[
 \eta_{\rm ret}=H_5-\frac{90}{\pi}\mathscr C_\pi,
\]

de sorte que

\[
 \hbar_a=\eta_a10^{-34}\mathcal U_S,
 \qquad h_a=2\pi\hbar_a.
\]

L’échelle \(10^{-34}\) ne provient pas d’une mesure de Planck : c’est la profondeur prédéfinie à laquelle le lecteur d’action trouve sa section. La phase \(\theta\) agit ensuite par

\[
 S(\theta)=\hbar_{\rm ret}\theta_{\rm rad}
 =h_{\rm ret}\frac{\theta_{\rm deg}}{360}.
\]

Le facteur \(2\pi\) ne corrige pas une unité choisie : il relie l’action par radian à l’action par tour. La carte angulaire et la droite d’action déterminent deux coordonnées du même mouvement orienté.

La transition hexade--octade évalue ensuite les canaux et produit le paramètre de Barbero--Immirzi. Sur le bord, si

\[
 N=N_{90}+N_{120},\qquad D=90N_{90}+120N_{120},
\]

alors

\[
 N_{90}=4N-\frac{D}{30},
 \qquad
 N_{120}=\frac{D}{30}-3N.
\]

Le lecteur conjoint \((N,D)\) conserve intégralement la provenance sectorielle :

\[
 H((N_{90},N_{120})\mid(N,D))=0.
\]

C’est la formulation exacte du coût de Landauer nul dans le passage aire--information. L’entropie d’un état réduit peut être positive ; ce qui vaut zéro est la perte conditionnelle lorsque le couple reconstructif complet est conservé. Le spectre d’aire

\[
 \frac{\mathcal A_{\rm phys}}{\ell_P^2}
 =\gamma_{\rm HMT}a_0(N_{90}+N_{120})
\]

et le hamiltonien modulaire

\[
 K_A=cI+\Delta_{\rm mod}(90N_{90}+120N_{120})
\]

sont ainsi deux lecteurs complémentaires : le premier compte l’aire totale ; le second conserve sa répartition entre les deux secteurs.

La gravitation occupe une autre résolution de la même fermeture dodécaphasique. Avec

\[
 \omega_{108}=\frac{2\pi}{108t_0},
 \qquad R_{108}=\frac{54}{\pi}\ell_0,
\]

et

\[
 G_a(\vartheta)=\vartheta\frac{c_{\rm int}^5t_0^2}{\hbar_a},
\]

la condition de coïncidence des trois longueurs

\[
 R_{108}=\bar\lambda_{C,a}=r_{g,a}
\]

sélectionne de manière unique

\[
 \vartheta_{108}=\left(\frac{54}{\pi}\right)^2
\]

et produit

\[
 \ell_P=\sqrt{\frac{G_a\hbar_a}{c_{\rm int}^3}}=R_{108}.
\]

La profondeur \(10^{-11}\) de la branche gravitationnelle et la profondeur \(10^{-34}\) de la branche d’action ne sont pas des échelles observées introduites dans le cristal. Ce sont des positions du raffinement où des lecteurs distincts de la même section prédéfinie déterminent gravitation et action. L’identité \(G_a\hbar_a=\vartheta_{108}c_{\rm int}^5t_0^2\) montre que les deux orientations d’action modifient réciproquement la gravitation et conservent la longueur de Planck. L’horloge, l’action et la gravitation sont donc des coordonnées synchronisées d’une évolution dimensionnelle unique.

\section{Loi directrice de l’article : évolution conservative de l’unité dimensionnelle}

Les résultats précédents permettent d’énoncer la thèse centrale sans en faire une métaphore ni un catalogue de constantes.

\textbf{Théorème d’évolution conservative et de prédéfinition holographique.} Soit \(\omega\in\Omega_{\rm seed}\) un germe APP orienté et soit \(\mathcal G_{\rm HMT}(\omega)\in X_\infty^{\rm enr}\) son prolongement par TRIT et TPK. Alors :

\begin{enumerate}
\item l’actualisation conserve de manière inversible feuillet, résidu, quotient, retenue, orientation, frontière, parcours et mémoire ;
\item l’horloge finie revient et l’holonomie avance, de sorte que l’orbite est hélicoïdale ;
\item la décoration mécanique irrationnelle possède une entropie topologique nulle, une complexité linéaire et aucune période de translation ;
\item tout lecteur compatible définit un système prédéfini dans une limite projective, de sorte que l’approfondissement révèle sans recalibrer ;
\item la double projection détermine, à partir de la même section, la droite réelle et la branche exceptionnelle Witt--Golay--Leech--\(V^\natural\)--Monstre ;
\item le système corrélé \((\pi,\varphi,e,\alpha)\) est une sortie corrélée et \(\alpha\) sa couture dodécaphasique avec l’horloge apériodique ;
\item la lacune électronique, l’action de Planck, le paramètre de Barbero--Immirzi, les grandeurs constitutives du vide et la gravitation sont des réalisations de lecteurs typés sur cette même généalogie.
\end{enumerate}

La démonstration est la composition des carrés déjà écrits. La compatibilité des troncatures prouve la prédéfinition ; l’inversion de \(\mathcal U_t\) prouve la conservation ; le cocycle sturmien prouve l’apériodicité d’entropie nulle ; l’identité des supports prouve la connexion électron--\(\pi\)--\(\alpha\) ; la décomposition de \(V_{12}\) prouve la séparation coordonnée des lecteurs exceptionnel et tensoriel ; les lecteurs d’action, d’aire et d’information prouvent la persistance de l’unité lors du changement de dimension.

La forme compacte du théorème est

\[
\boxed{
\begin{CD}
\Omega_{\rm seed}@>{\mathrm{APP\to TRIT\to TPK}}>>
X_\infty^{\rm enr}\\
@. @V{(\mathcal P_{\rm arq},\mathcal P_{\rm exc},
\mathcal P_{\rm dim})}VV\\
@. \mathbb R\times\mathscr E_{\rm exc}\times\mathscr D_{\rm phys},
\end{CD}}
\]

où les trois réalisations conservent la même provenance malgré des codomaines distincts. L’unité ne demeure pas immobile : elle se conserve en se raffinant, en changeant de carte, en acquérant de la mémoire et en augmentant sa dimension. Le continu n’est pas le support préalable de ce processus, mais sa projection terminale. C’est pourquoi le cristal apériodique hélicoïdal constitue le noyau de HMT--MD : il prédéfinit, dans une seule dynamique réversible, la généalogie des nombres et des constantes qui apparaissent ensuite comme mathématiques et physique.
 
\chapter{Génération corrélée de \(\pi\), \(e\) et \(\varphi\)}
\label{chap:83-09}

% Unidad U012. Biografía estructural completa de pi.

% La biografía de pi que seguía en esta residencia es exactamente la ya
% publicada en el capítulo 8, salvo el namespace de sus etiquetas. Se conserva
% en la fuente por su procedencia, pero no se imprime dos veces.
\iffalse

\section{Construction du caractère de clôture et des trajectoires généalogiques structurelles de \(\pi\)}
\label{p12-83-c9-s1:chap:biografia-pi-autonoma}

\subsection{Séparation causale entre germe, caractère et évaluation}

La construction se divise en trois strates :
\[
 \text{sélection finie}
 \longrightarrow
 \text{construction du caractère de clôture}
 \longrightarrow
 \text{évaluation archimédienne}.
\]
Le développement décimal n’intervient pas dans la première strate. La chaîne des types
est
\begin{equation}
 104\,976
 \longrightarrow468\,\mathcal U_6
 \xrightarrow{\Psi}243\,\mathcal W_6
 \lhook\joinrel\longrightarrow\mathbb F_3^6,
 \qquad|\mathbb F_3^6|=729.
 \label{p12-83-c9-s1:eq:cadena-tipica-pi-autonoma}
\end{equation}
Les cardinalités comptent, respectivement, les conditions initiales
enrichies, les émissions distinctes, les mots ternaires visibles et les éléments
de la fibre ambiante complète.

Pour \(U_6=(u_1,\ldots,u_6)\), la projection est
\[
 \Psi(U_6)=((-u_1)\bmod3,\ldots,(-u_6)\bmod3).
\]
Le recensement des \(43\) orbites de \(D_3\) démontre que l’orbite de
clôture est la seule dont les six orientations possèdent à la fois :
\[
 \text{cinq relèvements},\qquad
 \text{masse }1008,\qquad
 1008=432+4\cdot144.
\]
Le calibre
\[
 \operatorname{diag}_c=6,
 \qquad
 \operatorname{card}=ES
\]
sélectionne un unique représentant :
\begin{equation}
 \boxed{w_6^\pi=010211.}
 \label{p12-83-c9-s1:eq:semilla-pi-autonoma}
\end{equation}
La sélection utilise le catalogue, les multiplicités, l’action orbitale et l’orientation ;
elle ne consulte pas un préfixe de \(\pi\).

\subsection{Relèvements \(L_0,L_1\) et 1ʳᵉ frontière coinductive}

Sur les vecteurs-lignes de \(\mathbb F_3^6\), soient
\[
L_0=
\begin{pmatrix}
2&2&2&1&2&1\\
2&1&2&2&1&1\\
1&0&0&1&0&2\\
0&0&1&1&1&0\\
2&2&2&0&2&1\\
2&1&2&1&0&0
\end{pmatrix},
\quad
L_1=
\begin{pmatrix}
2&1&2&0&2&2\\
1&2&1&1&2&0\\
1&2&1&1&1&2\\
0&1&0&0&2&1\\
2&0&2&1&0&1\\
0&2&2&2&1&1
\end{pmatrix}.
\]
Si \(b_1=w_6^\pi\), la récurrence des blocs est
\[
 b_{k+1}=b_kL_{\varepsilon_k},
 \qquad
 (\varepsilon_1,\varepsilon_2,\varepsilon_3,\varepsilon_4)=(0,0,1,1),
\]
et le mot accumulé est formé par concaténation :
\[
 w_{6(k+1)}=w_{6k}\mathbin{\|}b_{k+1}.
\]
On ne multiplie pas un mot de longueur \(6k\) par une matrice \(6\times6\).

Pour le canal de clôture :
\[
\begin{array}{c|c}
\text{profondeur}&\text{nouveau bloc}\\ \hline
6&010211\\
12&012222\\
18&010211\\
24&002111\\
30&110221
\end{array}
\]
et
\begin{equation}
 \boxed{
 w_{30}^{\pi}
 =010211|012222|010211|002111|110221.}
 \label{p12-83-c9-s1:eq:w30-pi-autonomo}
\end{equation}
La frontière suivante est
\[
 u_5^\pi=222220.
\]
Avec les autres canaux,
\[
 R_{36}=
 \begin{pmatrix}
 222220\\021222\\102011
 \end{pmatrix}.
\]
Le symbole \(R_{36}\) enregistre le sixième bloc de trois histoires et ouvre une
frontière ; ce n’est pas un terminal.

La branche certifiée de vingt blocs commence par
\begin{align*}
\pi:\;&010211|012222|010211|002111|110221|222220|111201|\\
&212121|200121|100100|101222|022212|012012|111210|\\
&121011|200220|120210|000101|022010|020111.
\end{align*}
C’est pourquoi ni \(w_{30}\) ni \(R_{36}\) ne terminent la généalogie.

\subsection{Décomposition orientée de la microfibre de \texorpdfstring{\(\pi\)}{pi} : région centrale, \(4\) régions latérales et multiplicité \texorpdfstring{\(432+4\cdot144=1008\)}{432+4x144=1008}}

La microfibre
\[
 \mathscr R_\pi^{\rm micro}
 :=\{U\in\mathcal U_6:\Psi(U)=010211\}
\]
contient exactement cinq régions :
\[
\begin{array}{c|c|c|c}
U_6&E_{\rm sig}&\operatorname{mult}&\text{rôle}\\ \hline
501|614|498|169|272|272&555555&432&\perp\\
810|923|870|169|272|272&339555&144&++\\
810|983|810|169|272|272&393555&144&++\\
870|923|810|169|272|272&933555&144&++\\
870|923|810|418|674|521&933717&144&--
\end{array}
\]
En conséquence,
\begin{equation}
 \boxed{
 5=1_\perp+3_{++}+1_{--},
 \qquad
 1008=432+4\cdot144.}
 \label{p12-83-c9-s1:eq:descomposiciones-cinco-regiones-pi}
\end{equation}
La première identité classe des fonctions d’orientation. La seconde compte
des multiplicités de provenance. Elles ne comptent pas cinq copies indiscernables.

La région transversale est
\[
 501|614|498|169|272|272,
 \qquad E_{\rm sig}=555555,
 \qquad\operatorname{mult}=432.
\]
Le retour muté est
\[
 870|923|810|418|674|521,
 \qquad E_{\rm sig}=933717,
 \qquad\operatorname{mult}=144.
\]
Les intervertir conserverait le recensement \(3/1/1\), mais fausserait la structure
régionale complète.

\subsection{Descente d'un bloc nonadique complet au multigraphe de transfert de phase}

Pour \(U_6=(U_1,\ldots,U_6)\), soit
\[
 A(U_6)=(U_1,U_2,U_3),
 \qquad
 B(U_6)=(U_4,U_5,U_6).
\]
Chaque émission définit une arête entre \(A(U_6)\) et \(B(U_6)\). Le multigraphe
brut possède \(96\) sommets, répartis en composantes de tailles
\[
 28,28,28,4,4,4.
\]
On passe au quotient de la phase interne par
\[
 (a,b,c)\sim(b,c,a)\sim(c,a,b).
\]
Le représentant est la rotation lexicographiquement minimale. Le graphe
quotient possède \(32\) sommets et deux composantes de tailles \(28\) et \(4\).

Les ancrages sont
\[
\begin{aligned}
U_{\rm APP}&=903|850|850|252|109|252,\\
U_e&=850|963|890|149|252|212,\\
U_\varphi&=830|943|830|109|252|252.
\end{aligned}
\]
Nous numérotons les régions
\[
\begin{aligned}
U_\pi^{(1)}&=810|923|870|169|272|272,\\
U_\pi^{(2)}&=810|983|810|169|272|272,\\
U_\pi^{(3)}&=870|923|810|169|272|272,\\
U_\pi^{(4)}&=870|923|810|418|674|521,\\
U_\pi^{(5)}&=501|614|498|169|272|272.
\end{aligned}
\]

\subsection{Les \(5\) cycles orientés de clôture}

\subsubsection{Clôture de \texorpdfstring{\(U_\pi^{(1)}\)}{U-pi 1}}
\begin{align*}
&252|109|252|903|850|850\to
903|850|850|252|109|252\to\\
&252|109|252|923|870|810\to
810|923|870|169|272|272\to\\
&272|169|272|963|890|850\to
850|963|890|149|252|212\to\\
&212|149|252|943|830|830\to
830|943|830|109|252|252\to
252|109|252|903|850|850.
\end{align*}

\subsubsection{Clôture de \texorpdfstring{\(U_\pi^{(2)}\)}{U-pi 2}}
\begin{align*}
&903|850|850|252|109|252\to
272|169|272|903|850|850\to\\
&272|169|272|983|810|810\to
810|983|810|169|272|272\to\\
&272|169|272|963|890|850\to
850|963|890|149|252|212\to\\
&212|149|252|943|830|830\to
830|943|830|109|252|252\to
903|850|850|252|109|252.
\end{align*}

\subsubsection{Clôture de \texorpdfstring{\(U_\pi^{(3)}\)}{U-pi 3}}
\begin{align*}
&903|850|850|252|109|252\to
292|189|232|903|850|850\to\\
&292|189|232|923|810|870\to
870|923|810|169|272|272\to\\
&272|169|272|963|890|850\to
850|963|890|149|252|212\to\\
&212|149|252|943|830|830\to
830|943|830|109|252|252\to
903|850|850|252|109|252.
\end{align*}

\subsubsection{Clôture de \texorpdfstring{\(U_\pi^{(4)}\)}{U-pi 4}}
\begin{align*}
&903|850|850|252|109|252\to
272|169|272|903|850|850\to\\
&272|169|272|963|890|850\to
850|963|890|149|252|212\to\\
&212|149|252|923|810|870\to
870|923|810|418|674|521\to\\
&943|830|830|521|418|674\to
830|943|830|109|252|252\to
903|850|850|252|109|252.
\end{align*}

\subsubsection{Clôture de \texorpdfstring{\(U_\pi^{(5)}\)}{U-pi 5}}
\begin{align*}
&252|109|252|903|850|850\to
903|850|850|252|109|252\to\\
&614|498|501|252|109|252\to
501|614|498|169|272|272\to\\
&272|169|272|963|890|850\to
850|963|890|149|252|212\to\\
&212|149|252|943|830|830\to
830|943|830|109|252|252\to
252|109|252|903|850|850.
\end{align*}

\begin{theorem}[Minimalité des cinq clôtures]
\label{p12-83-c9-s1:thm:minimalidad-cierres-pi-autonoma}
Pour chaque \(i\in\{1,\ldots,5\}\), la longueur minimale d’une marche fermée
qui parcourt simultanément
\[
 U_\pi^{(i)},\quad U_e,\quad U_\varphi,\quad U_{\rm APP}
\]
dans le multigraphe quotient est huit.
\end{theorem}

\begin{proof}
Pour \(i\) fixé, soit
\[
 E_*=\{U_\pi^{(i)},U_e,U_\varphi,U_{\rm APP}\}.
\]
On considère l’espace fini
\[
 \mathscr S_i=\{(v,M):v\in V(G),\ M\subseteq E_*\}.
\]
La première composante enregistre le sommet et la seconde les arêtes cibles
déjà parcourues. Si une arête \(e\) mène de \(v\) à \(v'\), la transition est
\[
 (v,M)\longmapsto
 \bigl(v',M\cup(E_*\cap\{e\})\bigr).
\]
La recherche en largeur énumère les longueurs par ordre croissant. Le premier
retour à \((v,E_*)\) apparaît à la longueur huit pour les cinq régions.
Les marches imprimées réalisent la borne ; l’optimalité de la recherche exclut
une longueur inférieure.
\end{proof}

La fonction de clôture est encodée par
\[
 \mathfrak C_\pi=
 \bigl(\mathscr R_\pi^{\rm micro},
 w_\pi,\rho_\pi,m_\pi,\ell_{\rm cl}\bigr),
\]
\[
 \rho_\pi=(\perp,++ ,++ ,++ ,--),
 \qquad
 m_\pi=(432,144,144,144,144),
 \qquad
 \ell_{\rm cl}=8.
\]

\subsection{Opérateur d’incidence et demi-tour orienté}

L’ombre d’incidence est
\[
 A_W=
 \begin{pmatrix}
 0&1&1&1&1&1\\
 1&0&1&2&2&1\\
 1&1&0&1&2&2\\
 1&2&1&0&1&2\\
 1&2&2&1&0&1\\
 1&1&2&2&1&0
 \end{pmatrix}.
\]
La multiplication dans \(\mathbb F_3\) donne
\begin{equation}
 \boxed{A_W^2=2I_6=-I_6.}
 \label{p12-83-c9-s1:eq:aw-cuadrado-menos-identidad-pi}
\end{equation}
Les produits entre lignes et colonnes distinctes s’annulent modulo trois,
tandis que chaque produit diagonal vaut deux. Une application réalise un quart
de tour ternaire ; deux applications réalisent l’inversion orientée.

\subsection{La pente 1/5 imposée par les \(5\) régions}

Soit
\[
 \lambda=(\lambda_1,\ldots,\lambda_5)^{\mathsf T}
 \in\mathbb Q_{\geq0}^5,
 \qquad
 \mathbf1^{\mathsf T}\lambda=1.
\]
Oublier l’étiquette régionale applique le symétriseur
\[
 P_5=\frac15\mathbf1\mathbf1^{\mathsf T}.
\]
Pour tout vecteur normalisé,
\[
 P_5\lambda=\frac15\mathbf1.
\]
L’invariance \(P_5\lambda=\lambda\) a donc une unique solution :
\begin{equation}
 \boxed{\lambda=\frac15\mathbf1,\qquad t_1=\frac15.}
 \label{p12-83-c9-s1:eq:pendiente-primitiva-pi}
\end{equation}
L’entier cinq provient de la microfibre régionale ; il n’est pas reconnu dans une
formule de \(\pi\).

\subsection{Composition rationnelle et compensateur \texorpdfstring{\(239\)}{239}}

La loi de composition des pentes est
\[
 t\star s:=\frac{t+s}{1-ts}.
\]
Le premier doublement donne
\[
 t_2=t_1\star t_1
 =\frac{2/5}{1-1/25}
 =\frac5{12}.
\]
Le second donne
\[
 t_4=t_2\star t_2
 =\frac{10/12}{1-25/144}
 =\frac{120}{119}.
\]
Le retour à la diagonale unitaire exige \(q>0\) tel que
\[
 \frac{120/119-1/q}{1+(120/119)/q}=1.
\]
La multiplication par \(119q\) produit
\[
 120q-119=119q+120,
\]
et impose
\begin{equation}
 \boxed{q=239.}
 \label{p12-83-c9-s1:eq:compensador-239-pi}
\end{equation}

\subsection{Identité gaussienne et 1ʳᵉ demi-révolution}

On définit
\begin{equation}
 \Theta_\pi
 :=16\arctan\frac15-4\arctan\frac1{239}.
 \label{p12-83-c9-s1:eq:theta-pi-autonoma}
\end{equation}
Les carrés successifs sont
\[
\begin{array}{c|r@{}l}
\text{puissance}&\multicolumn{2}{c}{\text{entier gaussien}}\\ \hline
(5+i)^2&24&+10i\\
(5+i)^4&476&+480i\\
(5+i)^8&-3824&+456960i\\
(5+i)^{16}&-208797818624&-3494830080i\\
(239-i)^2&57120&-478i\\
(239-i)^4&3262465916&-54606720i
\end{array}
\]
et la multiplication finale donne
\begin{equation}
 \boxed{
 (5+i)^{16}(239-i)^4
 =-681386607803576157184.}
 \label{p12-83-c9-s1:eq:identidad-gaussiana-pi-autonoma}
\end{equation}
La partie imaginaire est nulle et la partie réelle est négative. Par conséquent,
\[
 \Theta_\pi\equiv\pi\pmod{2\pi}.
\]
Para \(0<x<1\),
\[
 x-\frac{x^3}{3}<\arctan x<x.
\]
D’où
\[
 \Theta_\pi>
 16\left(\frac15-\frac1{3\cdot5^3}\right)-\frac4{239}>3
\]
et
\[
 \Theta_\pi<\frac{16}{5}<4.
\]
L’unique orientation négative dans cet intervalle est le premier demi-tour
positif :
\begin{equation}
 \boxed{\Theta_\pi=\pi.}
 \label{p12-83-c9-s1:eq:reconocimiento-pi-exacto}
\end{equation}

L’ordre de la reconnaissance archimédienne ultérieure est
\[
 \text{cinq régions}
 \to\frac15
 \to\frac5{12}
 \to\frac{120}{119}
 \to239
 \to\Theta_\pi
 \to\pi.
\]

\subsection{Encadrements rationnels de précision arbitraire}

Pour \(0<x<1\), soit
\[
 S_n(x)=\sum_{r=0}^{n}(-1)^r\frac{x^{2r+1}}{2r+1}.
\]
Le critère de Leibniz donne
\[
 S_{2m+1}(x)<\arctan x<S_{2m}(x).
\]
On définit
\[
 a_m^-=S_{2m+1}(1/5),\quad a_m^+=S_{2m}(1/5),
\]
\[
 b_m^-=S_{2m+1}(1/239),\quad b_m^+=S_{2m}(1/239),
\]
et
\begin{equation}
 \ell_{\pi,m}=16a_m^- -4b_m^+,
 \qquad
 h_{\pi,m}=16a_m^+ -4b_m^-.
 \label{p12-83-c9-s1:eq:horquilla-pi-autonoma}
\end{equation}
Alors
\[
 \ell_{\pi,m}<\pi<h_{\pi,m}
\]
et
\begin{equation}
 h_{\pi,m}-\ell_{\pi,m}
 =
 \frac{16}{(4m+3)5^{4m+3}}
 +\frac{4}{(4m+3)239^{4m+3}}
 \longrightarrow0.
 \label{p12-83-c9-s1:eq:anchura-horquilla-pi}
\end{equation}
Les bornes inférieures croissent et les bornes supérieures décroissent.

Comme \(3<\pi<4\), le lecteur ternaire agit sur
\[
 r_\pi:=\pi-3\in(0,1),
\]
avec encadrement
\[
 [\ell_{\pi,m}-3,h_{\pi,m}-3].
\]

\subsection{Certification archimédienne ultérieure de la fibre déjà engendrée}

Supposons que la relation interne \(\operatorname{Ext}\) et la survie
coinductive aient déjà produit, sans utiliser Machin ni un développement tabulé,
l’extension compatible \(N_{K+1}\) du préfixe \(N_K\), de longueur
\(L_K=6K\). Les \(729\) sous-cylindres de sa fibre ambiante sont
\[
 I_{K,u}=
 \left[
 \frac{729N_K+u}{3^{L_K+6}},
 \frac{729N_K+u+1}{3^{L_K+6}}
 \right),
 \qquad0\leq u<729.
\]
L’ordre \(m=m(K)\) est ensuite augmenté jusqu’à ce que l’encadrement de Leibniz
localise le sous-cylindre déjà produit. Soient
\[
 \ell_{\pi,K}:=\ell_{\pi,m(K)}-3,
 \qquad
 h_{\pi,K}:=h_{\pi,m(K)}-3.
\]
Les extrémités entières compatibles sont
\[
 j_{\min}=
 \max\!\left(
 729N_K,\left\lfloor
 \ell_{\pi,K}3^{L_K+6}\right\rfloor
 \right),
\]
\[
 j_{\max}=
 \min\!\left(
 729N_K+728,\left\lceil
 h_{\pi,K}3^{L_K+6}\right\rceil-1
 \right).
\]
La condition de localisation unitaire est
\begin{equation}
 \boxed{j_{\min}=j_{\max},\qquad N_{K+1}=j_{\min}.}
 \label{p12-83-c9-s1:eq:seleccion-unitaria-subcilindro-pi}
\end{equation}
La première égalité certifie la localisation unitaire ; la seconde vérifie
que cette localisation coïncide avec l’extension engendrée intérieurement. Aucune
des deux ne définit la transition TPK ni ne sélectionne parmi les \(729\) états.
L’existence et l’unicité procèdent de \(\operatorname{Ext}\) et de la
survie ; Machin et les encadrements de Leibniz sont des reconnaissances
ultérieures.

\subsection{Prolongement coinductif de la section de \(\pi\) par le système inverse des histoires survivantes}

Soit \(\mathcal H_m^\pi\) l’ensemble des histoires enrichies à la profondeur
\(m\), et soit
\[
 p_{n,m}:\mathcal H_n^\pi\longrightarrow\mathcal H_m^\pi
\]
la troncature. Nous définissons
\[
 \operatorname{Surv}_m^{(N),\pi}
 :=p_{N,m}(\mathcal H_N^\pi),
\]
\[
 \operatorname{Surv}_m^\pi
 :=\bigcap_{N\geq m}\operatorname{Surv}_m^{(N),\pi}.
\]
Lorsque la relation \(\operatorname{Ext}\) est non vide, univoque et
compatible à tous les niveaux, les ensembles sont des singletons et satisfont
\[
 p_{m+1,m}(\operatorname{Surv}_{m+1}^{\pi})
 =\operatorname{Surv}_{m}^{\pi}.
\]
Dans ce domaine de prolongeabilité, il existe donc une unique section
\[
 x_\pi\in
 X_\pi:=\varprojlim_m\operatorname{Surv}_m^\pi.
\]
Les cylindres sont emboîtés et leurs diamètres sont \(3^{-6K}\), qui tend vers
zéro. Leur intersection contient un unique point, identifié par le caractère
angulaire à \(r_\pi=\pi-3\).

Pour toute profondeur décimale \(M\), on peut choisir un niveau \(K(M)\) dont le
cylindre est contenu dans une seule cellule décimale de diamètre \(10^{-M}\).
Les préfixes profonds se tronquent exactement aux précédents.

\subsection{Caractérisation du canal coinductif de \(\pi\) : sélection de \(w_6^\pi\), \(5\) régions, inversion orientée et unicité de la section limite}

\begin{theorem}[Caractérisation du canal de clôture]
\label{p12-83-c9-s1:thm:caracterizacion-final-pi-autonoma}
À partir du catalogue, de l’action \(D_3\), de la jauge orientée, des
opérateurs \(L_0,L_1,A_W\), de la loi des pentes et de la filtration rationnelle :
\begin{enumerate}
 \item on sélectionne \(w_6^\pi=010211\) ;
 \item on construit cinq régions avec les identités de
       \eqref{p12-83-c9-s1:eq:descomposiciones-cinco-regiones-pi} ;
 \item chaque région appartient à une clôture minimale de longueur huit ;
 \item \(A_W^2=-I_6\) fixe l’inversion orientée ;
 \item la symétrisation régionale force \(1/5\) ;
 \item la composition force \(120/119\) ;
 \item le retour unitaire force \(239\) ;
 \item l’identité gaussienne reconnaît le premier demi-tour ;
 \item les encadrements et la fibre de \(729\) éléments produisent une section
       compatible à toute profondeur finie.
\end{enumerate}
Son évaluation est
\[
 \boxed{\pi=16\arctan\frac15-4\arctan\frac1{239}.}
\]
\end{theorem}

La chaîne complète est
\[
 \mathrm{APP}\to\mathrm{TRIT}\to\mathrm{TPK}
 \to104\,976\to468\to243\to w_6^\pi
 \to w_{30}^\pi\to R_{36}\to\Gamma_9
 \to X_\pi\to\pi.
\]

\begin{criterion}[Réfutation du canal de clôture]
La dérivation est réfutée si :
\begin{enumerate}
 \item l’orbite de clôture cesse d’être unique ;
 \item l’une des cinq régions disparaît ou son identité complète change ;
 \item l’une des cinq clôtures a une longueur inférieure à huit ;
 \item \(A_W^2\ne2I_6\);
 \item le symétriseur admet une autre distribution normalisée ;
 \item les compositions ne produisent pas \(5/12\), \(120/119\) et \(239\) ;
 \item l’identité gaussienne échoue ;
 \item un encadrement cesse de contenir le caractère angulaire ;
 \item un niveau ne parcourt pas les \(729\) éléments de la fibre ambiante ou ne
       localise pas exactement une extension ;
 \item un développement cible intervient avant la reconnaissance.
\end{enumerate}
\end{criterion}
\fi

La construction complète du caractère de clôture, des cinq régions et de la
biographie structurelle de \(\pi\) a été établie dans le
\cref{chap:83-08}. À partir de ce même caractère déjà construit, le présent
chapitre prolonge maintenant le raffinement vers \(e\) et \(\varphi\), sans
réinitialiser la généalogie ni répéter sa première étape.


% Unidad U013. Biografía estructural completa del carácter exponencial.

\section{Propagation, retour avec mémoire et construction du caractère exponentiel}
\label{p12-83-c9-s2:chap:biografia-completa-e}

\subsection{Sélecteur orbital fini : existence et unicité de l’émission compatible}

Soit \(\mathcal W=\mathbb F_3^6\). Pour une émission
\(U=(u_1,\ldots,u_6)\), le mot visible est
\[
 \Psi(U)=((-u_1)\bmod3,\ldots,(-u_6)\bmod3).
\]
Sur \(\mathcal W\), l’action diédrale est donnée par
\[
\begin{aligned}
 r(u_1,u_2,u_3,u_4,u_5,u_6)
 &=(u_3,u_1,u_2,u_5,u_6,u_4),\\
 s(u_1,u_2,u_3,u_4,u_5,u_6)
 &=(u_4,u_5,u_6,u_1,u_2,u_3).
\end{aligned}
\]
Pour \(w\in\mathcal W\), nous définissons
\[
 n(w)=\#\Psi^{-1}(w),
 \qquad
 M(w)=\sum_{U\in\Psi^{-1}(w)}\operatorname{mult}(U),
\]
et
\[
 \nu(w)=
 \bigl(
 \#\{j:w_j=0\},
 \#\{j:w_j=1\},
 \#\{j:w_j=2\}
 \bigr).
\]

\begin{theorem}[Sélection structurelle du canal exponentiel]
\label{p12-83-c9-s2:thm:seleccion-estructural-e}
Parmi les \(43\) orbites du catalogue visible, il existe une unique orbite dont les
fibres satisfont simultanément :
\[
 n(w)=1,\qquad M(w)=144,\qquad
 \mathfrak D(w)=\{0,3,6\},\qquad
 \nu(w)=(2,3,1).
\]
Le calibre
\[
 \operatorname{diag}_c=6,
 \qquad \operatorname{card}=ES
\]
y sélectionne
\[
 \boxed{w_6^e=201101.}
\]
\end{theorem}

\begin{proof}
Les quatre conditions sont évaluées sur les \(243\) mots du catalogue,
non sur un échantillon. Une unique orbite satisfait le prédicat conjoint. Dans
cette orbite, une seule émission possède à la fois la diagonale \(6\) et l’orientation
\(ES\). L’énumération utilise uniquement des champs du catalogue APP--TRIT et ne
consulte pas de développement de \(e\).
\end{proof}

\subsection{Relèvement visible et chaîne de blocs}

La récurrence \(L_0,L_0,L_1,L_1\) produit
\[
 w_{30}^{e}
 =201101|121221|102011|012222|102011,
\]
et la frontière conjointe apporte
\[
 u_5^e=021222.
\]
Le registre de trente-six trits devient
\[
 w_{36}^{e}
 =201101|121221|102011|012222|102011|021222.
\]

\subsection{Signature décimale et carré positionnel}

La lecture initiale de douze chiffres est
\[
 \mathbf d_e^{(12)}
 =718281828459
 =7\,\underbrace{1828\,1828}_{\text{carré local}}\,459.
\]
L’égalité centrale est la concaténation
\[
 1828\mathbin{\|}1828,
\]
non un produit arithmétique ni le début supposé d’une période.

\begin{definition}[Carré positionnel]
Un mot \(d_1\cdots d_{12}\) contient un carré de longueur \(\ell\)
à la position \(p\) lorsque
\[
 d_p\cdots d_{p+\ell-1}
 =d_{p+\ell}\cdots d_{p+2\ell-1}.
\]
La position et les contextes latéraux appartiennent à la signature.
\end{definition}

\begin{table}[htbp]
\centering
\caption{Recensement exhaustif des carrés dans les lectures de douze chiffres.}
\label{p12-83-c9-s2:tab:censo-cuadrados-e}
\begin{tabular}{c*{6}{r}}
\toprule
Longueur \(\ell\)&1&2&3&4&5&6\\
\midrule
Occurrences&240&21&3&2&0&0\\
Mots qui les contiennent&153&19&3&2&0&0\\
\bottomrule
\end{tabular}
\end{table}

L’autre mot contenant un carré de longueur quatre est
\[
 575157518472
 =\underbrace{5751\,5751}_{\text{position initiale }1}\,8472.
\]

\begin{proposition}[Unicité positionnelle]
\label{p12-83-c9-s2:prop:cuadrado-posicional-e}
La lecture du canal \(e\) est la seule qui possède un carré de longueur
quatre après un préfixe de longueur un et avec les contextes \(7\) et \(459\).
\end{proposition}

\begin{proof}
Le recensement ne laisse que deux mots contenant des carrés de longueur quatre. L’un
commence son carré à la première position ; l’autre après le préfixe \(7\). La
position et les contextes séparent les deux possibilités.
\end{proof}

\subsection{Retour visible et rupture de mémoire}

Dans le relèvement
\[
 201101|121221|102011|012222|102011
\]
les troisième et cinquième blocs coïncident :
\[
 B_3=B_5=102011.
\]
Leurs préétats modulo \(27\) sont cependant
\[
 p_3=(25,11,7,17,8,16),
 \qquad
 p_5=(7,8,4,23,26,7),
\]
et les quotients ultérieurs :
\[
 q_3=(6,13,9,2,13,1),
 \qquad
 q_5=(15,0,3,19,23,25).
\]
On vérifie \(p_3\ne p_5\) et \(q_3\ne q_5\).

\begin{figure}[htbp]
\centering
\begin{tikzpicture}[
 node distance=11mm and 7mm,
 every node/.style={font=\small},
 block/.style={draw,rounded corners,minimum width=17mm,minimum height=8mm},
 state/.style={draw,dashed,rounded corners,align=center,inner sep=3pt},
 >=stealth
]
\node[block] (b1) {\(201101\)};
\node[block,right=of b1] (b2) {\(121221\)};
\node[block,right=of b2] (b3) {\(102011\)};
\node[block,right=of b3] (b4) {\(012222\)};
\node[block,right=of b4] (b5) {\(102011\)};
\draw[->] (b1)--(b2);
\draw[->] (b2)--(b3);
\draw[->] (b3)--(b4);
\draw[->] (b4)--(b5);
\node[state,below=of b3] (p3)
 {\(p_3=(25,11,7,17,8,16)\)\\
  \(q_3=(6,13,9,2,13,1)\)};
\node[state,below=of b5] (p5)
 {\(p_5=(7,8,4,23,26,7)\)\\
  \(q_5=(15,0,3,19,23,25)\)};
\draw[->] (p3)--(b3);
\draw[->] (p5)--(b5);
\draw[<->,thick] (b3) to[bend left=28]
 node[above]{même mot visible} (b5);
\end{tikzpicture}
\caption{Retour visible de \(102011\) avec renouvellement du préétat et du
quotient.}
\label{p12-83-c9-s2:fig:ruptura-memoria-e}
\end{figure}

\begin{theorem}[Retour visible sans périodicité de l’état]
\label{p12-83-c9-s2:thm:retorno-visible-no-periodico-e}
\(B_3=B_5\) est un retour de la projection visible, non une orbite périodique de l’état holonomique intégral.
\end{theorem}

\begin{proof}
Le retour de l’état complet exigerait l’égalité de toutes ses coordonnées.
Les inégalités du préétat et du quotient l’excluent. De même,
\(1828|1828\) ne commence pas une période décimale : après le deuxième bloc apparaît
\(4\), tandis qu’une continuation périodique exigerait \(1\).
\end{proof}

\subsection{Certification factorielle ultérieure de la section exponentielle coinductive}

Soit \(\mathcal A\) une algèbre commutative unitaire sur \(\mathbb Q\), et
soit \(P_t\in\mathcal A[[t]]\) le propagateur formel. La concaténation
temporelle impose
\[
 P_{s+t}=P_sP_t,
 \qquad
 P_0=1.
\]
Écrivons
\[
 P_t=\sum_{n=0}^{\infty}a_nt^n.
\]
L’unité fixe \(a_0=1\), et le transport élémentaire normalisé fixe
\(a_1=1\).

\begin{theorem}[Récurrence factorielle]
\label{p12-83-c9-s2:thm:recurrencia-factorial-e}
\[
 (n+1)a_{n+1}=a_n
 \qquad(n\geq0),
\]
et, par conséquent,
\[
 a_n=\frac1{n!},
 \qquad
 P_t=\sum_{n=0}^{\infty}\frac{t^n}{n!}.
\]
En \(t=1\),
\[
 \boxed{P_1=e.}
\]
\end{theorem}

\begin{proof}
Le coefficient de \(s^nt\) dans \(P_{s+t}\) est
\[
 \binom{n+1}{n}a_{n+1}=(n+1)a_{n+1}.
\]
Dans \(P_sP_t\), il est \(a_na_1=a_n\). L’égalité des séries produit la
récurrence. Une récurrence à partir de \(a_0=1\) donne \(a_n=1/n!\).
\end{proof}

\begin{corollary}[Équation différentielle]
\[
 \frac{d}{dt}P_t=P_t,
 \qquad P_0=1.
\]
\end{corollary}

\begin{proof}
\[
 P_t'=\sum_{n\geq0}(n+1)a_{n+1}t^n
 =\sum_{n\geq0}a_nt^n=P_t.
\]
\end{proof}

\subsection{Emboîtement cofinal des encadrements rationnels et convergence vers la valeur publiée}

Soit
\[
 S_N=\sum_{n=0}^{N}\frac1{n!}.
\]

\begin{theorem}[Borne factorielle explicite]
\label{p12-83-c9-s2:thm:cota-factorial-e}
Pour \(N\geq1\),
\begin{equation}
 \boxed{
 S_N<e<
 S_N+\frac{N+2}{(N+1)(N+1)!}.}
 \label{p12-83-c9-s2:eq:horquilla-factorial-e}
\end{equation}
\end{theorem}

\begin{proof}
Le reste est
\[
 R_N=\sum_{k=0}^{\infty}\frac1{(N+1+k)!}>0.
\]
Para \(k\geq0\),
\[
 (N+1+k)!
 \geq(N+1)!(N+2)^k.
\]
Par conséquent,
\[
 R_N
 \leq\frac1{(N+1)!}
 \sum_{k=0}^{\infty}\frac1{(N+2)^k}
 =\frac{N+2}{(N+1)(N+1)!}.
\]
L’inégalité est stricte à partir de \(k=2\).
\end{proof}

On définit
\[
 J_N^{(e)}=
 \left[
 S_N,\,
 S_N+\frac{N+2}{(N+1)(N+1)!}
 \right]
\]
et, pour la partie fractionnaire,
\[
 \widetilde J_N^{(e)}=J_N^{(e)}-2.
\]
La largeur converge vers zéro. Pour chaque \(M\), il existe \(N\) tel que
\[
 \operatorname{diam}J_N^{(e)}<10^{-M}.
\]
Si les extrémités appartiennent à une même cellule décimale, les \(M\) premiers
chiffres sont certifiés par arithmétique rationnelle.

\subsection{Prolongement coinductif de la section exponentielle par les sous-cylindres ternaires du Topological Phase Kernel}

Au niveau \(K\), TPK et la relation \(\operatorname{Ext}\) ont déjà engendré
un unique sous-cylindre compatible. On prend le plus petit \(N=N_e(K)\) pour lequel
l’encadrement factoriel est contenu dans ce sous-cylindre. \(N_e(K)\) est un
ordre de certification analytique, non un indice générateur ni une règle de
sélection de la fibre.

Les cylindres TPK sont compatibles par troncature indépendamment du
semi-groupe. Le théorème du semi-groupe identifie ensuite leur évaluation
archimédienne à \(e-2\) ; la partie entière retrouve \(e\). Si l’orbite,
\(\operatorname{Ext}\) ou la troncature échouent, la génération échoue ; si la
récurrence ou la borne factorielle échouent, c’est ce certificat analytique qui échoue, non la
généalogie APP--TRIT--TPK.

\subsection{Certificat adversarial de la génération de \(e\) : unicité orbitale, récurrence factorielle, encadrements rationnels et prolongation cylindrique}

\begin{criterion}[Réfutation du canal exponentiel]
La construction est réfutée si :
\begin{enumerate}
 \item une autre orbite satisfait le prédicat fini complet ;
 \item la signature \(7[1828\,1828]459\) n’est pas positionnellement unique ;
 \item les préétats ou les quotients des deux blocs \(102011\) coïncident ;
 \item le semi-groupe ne force pas \((n+1)a_{n+1}=a_n\) ;
 \item la borne \eqref{p12-83-c9-s2:eq:horquilla-factorial-e} échoue ;
 \item un niveau ne localise pas de manière unitaire une extension compatible ;
 \item un développement tabulé intervient avant la reconnaissance.
\end{enumerate}
\end{criterion}


% Unidad U014. Biografía estructural completa del carácter de autoescala.

\section{Incidence modale, demi-droite de Perron et construction du caractère d'auto-échelle}
\label{p12-83-c9-s3:chap:biografia-completa-phi}

\subsection{Sélection finie du germe d’auto-échelle}

Soit \(\mathcal U_6\subset\{0,\ldots,999\}^6\) le catalogue des
\(468\) émissions distinctes construit dans le
\cref{chap:semillas-emisor-54}, et soit
\[
 \Psi:\mathcal U_6\longrightarrow\mathbb F_3^6,
 \qquad
 \Psi(U)=((-U_j)\bmod3)_{j=1}^{6}.
\]
Pour \(w\in\Psi(\mathcal U_6)\), rappelons les invariants
\[
 n(w)=\#\Psi^{-1}(w),
 \qquad
 M(w)=\sum_{U\in\Psi^{-1}(w)}\operatorname{mult}(U),
\]
et le vecteur d’occupation
\[
 \nu(w)=
 \bigl(\#\{j:w_j=0\},\#\{j:w_j=1\},\#\{j:w_j=2\}\bigr).
\]

\begin{theorem}[Sélection orbitale de l’auto-échelle]
\label{p12-83-c9-s3:thm:seleccion-orbital-phi}
L’action diédrale sur \(\Psi(\mathcal U_6)\) contient une unique orbite
qui satisfait le prédicat d’auto-échelle
\[
 n(w)=1,
 \qquad
 M(w)=144,
 \qquad
 \nu(w)=(2,2,2).
\]
La jauge interne
\[
 \operatorname{diag}_c=6,
 \qquad
 \operatorname{card}=ES
\]
sélectionne dans cette orbite le mot
\[
 \boxed{w_6^{\varphi}=121200.}
\]
Son unique émission est
\[
 \boxed{U_{\varphi}=(830,943,830,109,252,252),}
\]
de multiplicité \(144\), de signature multiplicative \(555933\), d’orientation
additive \(ES\), de support multiplicatif de cardinal six et de type résiduel
\(A\).
\end{theorem}

\begin{proof}
La projection directe donne
\[
 U_{\varphi}\bmod3=(2,1,2,1,0,0),
 \qquad
 -U_{\varphi}\bmod3=(1,2,1,2,0,0).
\]
L’énumération est effectuée sur les \(243\) mots visibles et leurs
\(43\) orbites, non sur une sélection d’exemples. Le prédicat conjoint
\(n=1\), \(M=144\), \(\nu=(2,2,2)\) laisse une orbite. En son sein, la
diagonale et l’orientation déclarées laissent une émission. À cette étape,
ni l’équation \(x^2-x-1=0\), ni une table décimale, ni le nom
conventionnel de la constante n’ont été introduits.
\end{proof}

\subsection{Relèvement par blocs et 1ʳᵉ frontière coinductive}

Soit \(b_1=w_6^{\varphi}\), et soit
\(\varepsilon=(0,0,1,1)\). Avec les matrices \(L_0,L_1\) construites dans le
\cref{hmtctx:p12:chap:orbitas-elevacion-r36}, nous définissons
\[
 b_{k+1}=b_kL_{\varepsilon_k}\pmod3
 \qquad(1\leq k\leq4),
\]
et concaténons
\[
 w_{6(k+1)}^{\varphi}=w_{6k}^{\varphi}\mathbin{\|}b_{k+1}.
\]
Le calcul matriciel produit, bloc par bloc,
\[
\begin{aligned}
 b_1&=121200,\\
 b_2&=112202,\\
 b_3&=121020,\\
 b_4&=010210,\\
 b_5&=010200,
\end{aligned}
\]
et, par conséquent,
\begin{equation}
 \boxed{
 w_{30}^{\varphi}
 =121200|112202|121020|010210|010200.}
 \label{p12-83-c9-s3:eq:w30-phi}
\end{equation}
La neutralité duale de la frontière ajoute
\[
 u_5^{\varphi}=102011,
\]
de sorte que
\begin{equation}
 \boxed{
 w_{36}^{\varphi}
 =121200|112202|121020|010210|010200|102011.}
 \label{p12-83-c9-s3:eq:w36-phi}
\end{equation}

\begin{remark}[Deux usages distincts d’un mot visible]
Le mot \(102011\) est également visible dans la trajectoire de \(e\). L’égalité d’une projection n’identifie pas les états holonomiques intégraux : le canal, le préfixe, le quotient, la retenue, la phase et la mémoire appartiennent au type de l’objet. Cette distinction est obligatoire avant de construire tout caractère réel.
\end{remark}

\subsection{Du calendrier tritique au multigraphe d’incidence}

\subsubsection{Incidence modale primitive \(F_{\rm av}\) et dynamique de mémoire d'ordre \(1\)}

Le calendrier tritique primitif est
\[
 (+1,-1,0).
\]
Nous fixons deux modes visibles, \(\Sigma\) et \(\Pi\), et la règle
\begin{equation}
 +1:m\longmapsto\Sigma,
 \qquad
 -1:m\longmapsto\Pi,
 \qquad
 0:m\longmapsto m.
 \label{p12-83-c9-s3:eq:regla-modal-primitiva-phi}
\end{equation}
Avec la condition de clôture cyclique, l’orbite modale produite par
\eqref{p12-83-c9-s3:eq:regla-modal-primitiva-phi} est
\[
 \boxed{(\Sigma,\Pi,\Pi).}
\]
Les trois paires consécutives sont
\[
 \Sigma\longrightarrow\Pi,
 \qquad
 \Pi\longrightarrow\Pi,
 \qquad
 \Pi\longrightarrow\Sigma.
\]

\begin{definition}[Matrice d’incidence chronologique]
Pour \(i,j\in\{\Sigma,\Pi\}\), soit
\[
 N_3(i,j)=
 \#\{r\in\mathbb Z/3\mathbb Z:(m_r,m_{r+1})=(i,j)\}.
\]
L’indice \(3\) compte les instants du calendrier primitif ; il ne désigne pas une
puissance matricielle.
\end{definition}

\begin{theorem}[Descente nécessaire de l’incidence modale]
\label{p12-83-c9-s3:thm:descenso-incidencia-modal-phi}
Dans l’ordre \((\Sigma,\Pi)\),
\begin{equation}
 \boxed{
 N_3=
 \begin{pmatrix}
 0&1\\
 1&1
 \end{pmatrix}.}
 \label{p12-83-c9-s3:eq:N3-phi}
\end{equation}
Par répétition exacte du calendrier,
\begin{equation}
 \boxed{
 N_9=3N_3,
 \qquad
 N_{27}=9N_3,
 \qquad
 N_{108}=36N_3.}
 \label{p12-83-c9-s3:eq:incidencias-9-27-108-phi}
\end{equation}
Ces trois égalités sont des égalités de dénombrements d’arêtes ; elles n’affirment
ni \(N_9=N_3^3\), ni \(N_{27}=N_3^9\), ni \(N_{108}=N_3^{36}\).
\end{theorem}

\begin{proof}
\(\Sigma\to\Sigma\) n’apparaît pas ; chacune des trois autres paires apparaît
une fois. Cela détermine les quatre entrées de \(N_3\). Les calendriers de
longueurs \(9\), \(27\) et \(108\) contiennent respectivement \(3\), \(9\)
et \(36\) copies du bloc primitif. Chaque répétition multiplie tous les
dénombrements par le même entier, ce qui prouve
\eqref{p12-83-c9-s3:eq:incidencias-9-27-108-phi}.
\end{proof}

\subsubsection{Publication surfacique–volumétrique et rapport stable \(\varphi^{-2}\) sous la mesure de Parry}

L’étiquette modale et le feuillet géométrique sont des objets distincts. La carte
de publication est
\[
 \mathfrak g_{\rm av}(\Sigma)=\mathscr H_{\rm ar},
 \qquad
 \mathfrak g_{\rm av}(\Pi)=\mathscr H_{\rm vol}.
\]
En transportant \eqref{p12-83-c9-s3:eq:N3-phi} par \(\mathfrak g_{\rm av}\), nous obtenons
l’opérateur d’incidence
\begin{equation}
 \boxed{
 F_{\rm av}=
 \begin{pmatrix}
 0&1\\
 1&1
 \end{pmatrix}.}
 \label{p12-83-c9-s3:eq:Fav-derivada-phi}
\end{equation}
Ainsi, les deux transitions croisées, l’unique autorétention volumétrique et
l’absence d’autorétention aréale découlent du calendrier ; ce ne sont pas quatre
entrées indépendantes.

\begin{remark}[Portée de la descente]
L’oubli de la phase élémentaire ne produit pas, à lui seul, un quotient de Markov
fort du Topological Phase Kernel, parce que le mode visible ne détermine pas
lequel des trois opérateurs internes agira au pas suivant. Ce qui
descend en revanche sans hypothèse supplémentaire est le multigraphe des arêtes d’un
bloc complet. Son enveloppe minimale de mémoire un a pour matrice
d’adjacence \(F_{\rm av}\).
\end{remark}

\subsection{Demi-droite positive invariante et unicité de l’auto-échelle}

\begin{theorem}[Demi-droite de Perron dérivée]
\label{p12-83-c9-s3:thm:rayo-perron-phi}
Le cône positif de \(F_{\rm av}\) contient une unique demi-droite propre. Si nous la
normalisons comme \(v=(1,x)^{\mathsf T}\), alors
\[
 x^2-x-1=0,
 \qquad
 x>0,
\]
et, par conséquent,
\begin{equation}
 \boxed{x=\frac{1+\sqrt5}{2}=\varphi.}
 \label{p12-83-c9-s3:eq:phi-perron}
\end{equation}
\end{theorem}

\begin{proof}
L’équation \(F_{\rm av}v=\lambda v\) donne, coordonnée par coordonnée,
\[
 x=\lambda,
 \qquad
 1+x=\lambda x.
\]
En éliminant \(\lambda\), on obtient \(x^2-x-1=0\). Ses racines sont
\((1\pm\sqrt5)/2\), et seule la positive appartient à l’intérieur du cône. La
matrice \(F_{\rm av}\) est primitive, puisque \(F_{\rm av}^2\) a toutes ses
entrées positives ; le théorème de Perron--Frobenius garantit que cette demi-droite
positive est simple et unique. Le nom \(\varphi\) est attribué après cette
construction.
\end{proof}

\begin{corollary}[Équation scalaire d’auto-échelle]
\[
 \varphi=1+\frac1\varphi,
 \qquad
 \varphi^2=\varphi+1,
 \qquad
 \varphi^{-1}=\varphi-1.
\]
Ces identités sont des conséquences de \eqref{p12-83-c9-s3:eq:phi-perron} ; elles ne sélectionnent
ni le germe ni la matrice.
\end{corollary}

\subsection{Récurrence de Fibonacci de la demi-droite de Perron et emboîtement diophantien des encadrements rationnels}

Soit \(F_0=0\), \(F_1=1\) et
\[
 F_{n+1}=F_n+F_{n-1}\qquad(n\geq1).
\]

\begin{theorem}[Puissances exactes de l’incidence]
\label{p12-83-c9-s3:thm:potencias-fibonacci-phi}
Pour tout \(n\geq1\),
\begin{equation}
 \boxed{
 F_{\rm av}^{\,n}=
 \begin{pmatrix}
 F_{n-1}&F_n\\
 F_n&F_{n+1}
 \end{pmatrix}.}
 \label{p12-83-c9-s3:eq:potencias-Fav-fibonacci}
\end{equation}
\end{theorem}

\begin{proof}
Pour \(n=1\), l’identité est la définition de \(F_{\rm av}\). Si elle vaut
pour \(n\), la multiplication à droite par \(F_{\rm av}\) donne
\[
 \begin{pmatrix}
 F_n&F_{n-1}+F_n\\
 F_{n+1}&F_n+F_{n+1}
 \end{pmatrix}
 =
 \begin{pmatrix}
 F_n&F_{n+1}\\
 F_{n+1}&F_{n+2}
 \end{pmatrix},
\]
ce qui est le cas \(n+1\).
\end{proof}

\begin{theorem}[Identité de Cassini]
\label{p12-83-c9-s3:thm:cassini-phi}
Pour tout \(n\geq1\),
\begin{equation}
 \boxed{F_{n+1}F_{n-1}-F_n^2=(-1)^n.}
 \label{p12-83-c9-s3:eq:cassini-phi}
\end{equation}
\end{theorem}

\begin{proof}
Le déterminant de \(F_{\rm av}\) est \(-1\). En prenant les déterminants dans
\eqref{p12-83-c9-s3:eq:potencias-Fav-fibonacci},
\[
 F_{n-1}F_{n+1}-F_n^2
 =\det(F_{\rm av}^{\,n})=(-1)^n.
\]
\end{proof}

Nous définissons \(r_n=F_{n+1}/F_n\) pour \(n\geq1\). De Cassini découle
\[
 r_n-r_{n-1}
 =\frac{F_{n+1}F_{n-1}-F_n^2}{F_nF_{n-1}}
 =\frac{(-1)^n}{F_nF_{n-1}}.
\]
Par conséquent, les approximations alternent autour de la demi-droite de Perron.

\begin{proposition}[Encadrements rationnels de Perron]
\label{p12-83-c9-s3:prop:horquillas-perron-phi}
Pour \(m\geq1\),
\begin{equation}
 \frac{F_{2m+2}}{F_{2m+1}}
 <\varphi<
 \frac{F_{2m+1}}{F_{2m}},
 \label{p12-83-c9-s3:eq:horquilla-fibonacci-phi}
\end{equation}
et la largeur est
\begin{equation}
 \boxed{
 \frac{F_{2m+1}}{F_{2m}}
 -\frac{F_{2m+2}}{F_{2m+1}}
 =\frac{1}{F_{2m}F_{2m+1}}.}
 \label{p12-83-c9-s3:eq:anchura-fibonacci-phi}
\end{equation}
\end{proposition}

\begin{proof}
L’alternance précédente place les quotients pairs en dessous et les impairs
au-dessus de \(\varphi\). La formule de la largeur est
\eqref{p12-83-c9-s3:eq:cassini-phi} avec \(n=2m+1\), après réduction des deux fractions au
même dénominateur. Comme \(F_n\to\infty\), la largeur tend vers zéro.
\end{proof}

\subsection{Incidence de mémoire d’ordre \(1\) et mesure de Parry de l’enveloppe symbolique}

Soit \(X_{\rm av}\) le plus petit sous-décalage de type fini sur
\(\{\mathscr H_{\rm ar},\mathscr H_{\rm vol}\}\) qui admet exactement
les trois paires observées dans le calendrier modal. Sa matrice d’adjacence est
\(F_{\rm av}\).

\begin{theorem}[Transition et mesure de Parry]
\label{p12-83-c9-s3:thm:parry-phi}
Soit \(r=(1,\varphi)^{\mathsf T}\). La matrice stochastique de Parry est
\begin{equation}
 \boxed{
 P_{ij}=
 \frac{(F_{\rm av})_{ij}r_j}{\varphi r_i}
 =
 \begin{pmatrix}
 0&1\\
 \varphi^{-2}&\varphi^{-1}
 \end{pmatrix}.}
 \label{p12-83-c9-s3:eq:parry-matriz-phi}
\end{equation}
La mesure stationnaire des sommets est proportionnelle à
\((1,\varphi^2)\), donc
\begin{equation}
 \boxed{
 \frac{\mu_{\rm ar}}{\mu_{\rm vol}}=\varphi^{-2}.}
 \label{p12-83-c9-s3:eq:parry-razon-phi}
\end{equation}
\end{theorem}

\begin{proof}
La somme de la ligne \(i\) de \(P\) est
\[
 \frac{(F_{\rm av}r)_i}{\varphi r_i}=1.
\]
Comme \(F_{\rm av}\) est symétrique, les vecteurs propres gauche et droit
coïncident. Les poids stationnaires sont proportionnels à leurs produits
coordonnée par coordonnée, c’est-à-dire à \((1,\varphi^2)\). Le rapport entre les
deux poids est \(\varphi^{-2}\).
\end{proof}

\begin{remark}[Trois mesures à ne pas confondre]
La fréquence chronologique de l’orbite primitive est
\[
 \nu_{\rm cyc}(\mathscr H_{\rm ar})=\frac13,
 \qquad
 \nu_{\rm cyc}(\mathscr H_{\rm vol})=\frac23.
\]
La mesure uniforme \(\tau_{468}\) vit sur le catalogue des émissions. La
mesure de Parry vit sur les histoires admissibles de longueur arbitraire. La
première compte les instants, la deuxième compte les émissions et la troisième pondère
les cylindres dynamiques. Aucune n’est une approximation des autres.
\end{remark}

\subsection{Caractère multiplicatif et conservation cylindrique}

Pour une route admissible \(\gamma:i\to j\) de longueur \(n\), nous définissons
\begin{equation}
 \chi_P(\gamma)=\varphi^n\frac{r_i}{r_j}.
 \label{p12-83-c9-s3:eq:caracter-parry-phi}
\end{equation}

\begin{proposition}[Multiplicativité par concaténation]
\label{p12-83-c9-s3:prop:multiplicatividad-parry-phi}
Si \(\gamma_1:i\to j\) et \(\gamma_2:j\to k\), alors
\[
 \chi_P(\gamma_2\circ\gamma_1)
 =\chi_P(\gamma_2)\chi_P(\gamma_1).
\]
Sur une route fermée de longueur \(n\),
\[
 \chi_P(\gamma)=\varphi^n.
\]
\end{proposition}

\begin{proof}
La longueur s’additionne et la coordonnée intermédiaire se simplifie par télescopage :
\[
 \varphi^{n_1+n_2}\frac{r_i}{r_k}
 =\left(\varphi^{n_1}\frac{r_i}{r_j}\right)
  \left(\varphi^{n_2}\frac{r_j}{r_k}\right).
\]
Si \(i=k\), le quotient terminal est égal à un.
\end{proof}

Nous normalisons le vecteur gauche comme
\[
 \ell=\frac{r}{1+\varphi^2},
 \qquad
 \ell^{\mathsf T}r=1.
\]
La mesure d’un cylindre est
\[
 \mu[x_0\cdots x_n]
 =\ell_{x_0}r_{x_n}\varphi^{-n}.
\]
Par conséquent,
\begin{equation}
 V_n(x)
 :=\frac{\mu[x_0]}{\mu[x_0\cdots x_n]}
 =\varphi^n\frac{r_{x_0}}{r_{x_n}}
 =\chi_P(x_0\cdots x_n).
 \label{p12-83-c9-s3:eq:volumen-cilindrico-phi}
\end{equation}
Pour une dimension typée \(D>0\), l’échelle
\[
 a_n^{(D)}(x)=V_n(x)^{1/D}
\]
satisfait la loi exacte
\begin{equation}
 \boxed{
 \bigl(a_n^{(D)}(x)\bigr)^D
 \mu[x_0\cdots x_n]=\mu[x_0].}
 \label{p12-83-c9-s3:eq:conservacion-cilindrica-phi}
\end{equation}

En dimension trois et sur des retours fermés,
\[
 \frac{a_9}{a_0}=\varphi^3,
 \qquad
 \frac{a_{27}}{a_0}=\varphi^9,
 \qquad
 \frac{a_{108}}{a_0}=\varphi^{36}.
\]
Si \(A_k=a_{9k}\), alors
\begin{equation}
 A_{k+1}-2A_k+A_{k-1}
 =A_{k-1}(\varphi^3-1)^2>0.
 \label{p12-83-c9-s3:eq:convexidad-retornos-phi}
\end{equation}
C’est une convexité en profondeur de retour ; elle n’est pas interprétée comme une
grandeur cinématique sans une carte supplémentaire de durée.

\subsection{Mémoire quadratique et branche contractée}

Les valeurs propres de \(F_{\rm av}\) sont
\[
 \varphi,
 \qquad
 -\varphi^{-1}.
\]
La branche contractée change d’orientation à chaque itération. Pour conserver
ce signe avant d’élever au carré, nous passons à
\(\operatorname{Sym}^2(\mathbb R^2)\). Dans la base
\((x^2,2xy,y^2)\),
\begin{equation}
 \operatorname{Sym}^2(F_{\rm av})=
 \begin{pmatrix}
 0&0&1\\
 0&1&2\\
 1&1&1
 \end{pmatrix}.
 \label{p12-83-c9-s3:eq:sym2-Fav-phi}
\end{equation}

\begin{theorem}[Spectre de la mémoire de second ordre]
\label{p12-83-c9-s3:thm:espectro-sym2-phi}
\[
 \det\bigl(tI-\operatorname{Sym}^2(F_{\rm av})\bigr)
 =(t+1)(t^2-3t+1),
\]
et
\begin{equation}
 \boxed{
 \operatorname{Spec}\bigl(\operatorname{Sym}^2(F_{\rm av})\bigr)
 =\{\varphi^2,-1,\varphi^{-2}\}.}
 \label{p12-83-c9-s3:eq:espectro-sym2-phi}
\end{equation}
L’unique mode stable positif est \(\varphi^{-2}\).
\end{theorem}

\begin{proof}
Les valeurs propres d’une puissance symétrique de degré deux sont les produits
\(\lambda_1^2\), \(\lambda_1\lambda_2\) et \(\lambda_2^2\). Avec
\(\lambda_1=\varphi\) et \(\lambda_2=-\varphi^{-1}\), on obtient
\(\varphi^2\), \(-1\) et \(\varphi^{-2}\). La première dépasse un, la
deuxième a pour module un et la troisième appartient à \((0,1)\).
\end{proof}

Sur une route fermée avec \(V_n=\varphi^n\), le mode stable satisfait
\begin{equation}
 \boxed{M_n=\varphi^{-2n}=V_n^{-2}.}
 \label{p12-83-c9-s3:eq:modo-estable-phi}
\end{equation}

\subsection{Réalisation hyperbolique normalisée de l’auto-échelle : \(\exp(2\log\varphi\,J_h)=F_{\mathrm{av}}^2\)}

Soit
\[
 A=F_{\rm av}^{\,2}=
 \begin{pmatrix}1&1\\1&2\end{pmatrix},
 \qquad
 J_h=\frac{2A-3I}{\sqrt5}.
\]
Alors
\[
 J_h^2=I,
 \qquad
 A=\frac32I+\frac{\sqrt5}{2}J_h.
\]
Comme
\[
 \cosh(2\log\varphi)=\frac32,
 \qquad
 \sinh(2\log\varphi)=\frac{\sqrt5}{2},
\]
on obtient
\begin{equation}
 \boxed{
 \exp(2\log\varphi\,J_h)=F_{\rm av}^{\,2}.}
 \label{p12-83-c9-s3:eq:realizacion-hiperbolica-phi}
\end{equation}
La constante d’auto-échelle est donc simultanément la valeur propre de
Perron de l’incidence et le paramètre exponentiel de sa réalisation
hyperbolique normalisée.

\subsection{Caractérisation, encadrements et prolongement}

La construction finie détermine la demi-droite positive de \(F_{\rm av}\). La
caractérisation analytique ultérieure est
\[
 x^2-x-1=0,
 \qquad x>0.
\]
Les encadrements \eqref{p12-83-c9-s3:eq:horquilla-fibonacci-phi} ont des extrémités rationnelles
et une largeur explicite \eqref{p12-83-c9-s3:eq:anchura-fibonacci-phi}. Pour chaque niveau
nonadique \(K\), nous choisissons le plus petit \(m=m_{\varphi}(K)\) dont l’encadrement est
contenu dans un unique sous-cylindre parmi les \(729\) éléments de la fibre
ambiante. La minimalité fait de la localisation une fonction, non un
choix tacite.

Les cylindres sélectionnés sont compatibles avec les troncatures et leurs
diamètres tendent vers zéro. Leur intersection contient un unique réel ; le
\cref{p12-83-c9-s3:thm:rayo-perron-phi} l’identifie à \(\varphi\). Les quotients de
Fibonacci sont donc des certificats rationnels de la localisation, non des données
décimales employées pour décider du germe.

\subsection{Certificat interne d'unicité, de compatibilité et de convergence de l'auto-échelle}

\begin{criterion}[Réfutation de la biographie d'auto-échelle]
La construction est réfutée si l'un des faits
suivants se produit :
\begin{enumerate}
 \item le prédicat fini sélectionne plus d’une orbite ou plus d’une émission ;
 \item la récurrence \(L_0,L_0,L_1,L_1\) ne produit pas
       \eqref{p12-83-c9-s3:eq:w30-phi} ;
 \item le calendrier ne produit pas exactement les trois paires utilisées dans
       \eqref{p12-83-c9-s3:eq:N3-phi} ;
 \item un multiple de dénombrement \(N_{9},N_{27},N_{108}\) est confondu avec une
       puissance de \(N_3\) ;
 \item \(F_{\rm av}\) possède une autre demi-droite propre positive ;
 \item \eqref{p12-83-c9-s3:eq:potencias-Fav-fibonacci} ou l’identité de Cassini échoue ;
 \item les encadrements rationnels ne sont pas emboîtés ou ne contractent pas leur largeur ;
 \item la fréquence chronologique, la mesure uniforme du
       catalogue ou la mesure de Parry sont identifiées ;
 \item une table décimale intervient avant la construction de la demi-droite de Perron ;
 \item un cylindre sélectionné ne se tronque pas au cylindre précédent.
\end{enumerate}
\end{criterion}


% Unidad U015. Acoplamiento de horquillas racionales, refinamiento TPK y
% certificación a profundidad arbitraria.

\section{Certification diagonale, limite projective et génération à une précision arbitraire}
\label{p12-83-c9-s4:chap:precision-arbitraria-constantes}

\subsection{Séparation typée des \(4\) indices du prolongement nonadique}

Cette unité a le statut de certification ultérieure. Les préfixes et leurs
extensions compatibles procèdent de la relation interne
\(\operatorname{Ext}\), de la survie coinductive et des lois de
transition APP--TRIT--TPK. Les encadrements rationnels qui suivent ne choisissent pas
parmi les \(729\) éléments d’une fibre : ils localisent dans la carte archimédienne
l’extension déjà engendrée et certifient sa précision à une profondeur arbitraire.

Pour chaque caractère \(\chi\in\{\pi,e,\varphi\}\) interviennent quatre
paramètres de natures différentes :
\begin{enumerate}
 \item \(n\), ordre de l’encadrement analytique rationnel ;
 \item \(K\), niveau de raffinement du Topological Phase Kernel ;
 \item \(L_K=6K\), longueur ternaire du préfixe publié ;
 \item \(M\), nombre de positions décimales demandé.
\end{enumerate}
On ne postule ni \(n=K\), ni \(n=L_K\), ni \(K=M\), ni une proportionnalité fixe
entre eux. Le couplage est construit par des inclusions d’intervalles
décidées par arithmétique entière.

Les parties entières sont
\[
 m_\pi=3,
 \qquad
 m_e=2,
 \qquad
 m_\varphi=1,
\]
et les caractères fractionnaires sont
\[
 \rho_\pi=\pi-3,
 \qquad
 \rho_e=e-2,
 \qquad
 \rho_\varphi=\varphi-1.
\]
Les trois familles d’encadrements semi-ouverts s’écrivent
\[
 J_{\chi,n}=[\ell_{\chi,n},h_{\chi,n}).
\]
Leurs extrémités sont rationnelles, contiennent \(\rho_\chi\) et satisfont
\[
 \ell_{\chi,n}\nearrow\rho_\chi,
 \qquad
 h_{\chi,n}\searrow\rho_\chi,
 \qquad
 h_{\chi,n}-\ell_{\chi,n}\longrightarrow0.
\]
Pour \(\pi\), les extrémités proviennent des sommes alternées de l’identité
de Machin démontrée dans le \cref{chap:biografia-pi-autonoma} ; pour \(e\),
des sommes factorielles et de leur reste rationnel du
\cref{p12-83-c9-s2:chap:biografia-completa-e} ; pour \(\varphi\), des quotients de
Fibonacci et de Cassini du \cref{p12-83-c9-s3:chap:biografia-completa-phi}.

\subsection{Cylindres ternaires et fibre ambiante immédiate}

Un préfixe ternaire de longueur \(L_K=6K\) est représenté par un entier
\(N_K\) avec
\[
 0\leq N_K<3^{L_K}.
\]
Son cylindre semi-ouvert est
\begin{equation}
 C(N_K,L_K)=
 \left[
 \frac{N_K}{3^{L_K}},
 \frac{N_K+1}{3^{L_K}}
 \right).
 \label{p12-83-c9-s4:eq:cilindro-nivel-K}
\end{equation}
L’émission suivante ajoute six trits. La fibre ambiante immédiate contient
\(3^6=729\) éléments, indexés par \(u\in\{0,\ldots,728\}\) :
\begin{equation}
 C_{K,u}=
 \left[
 \frac{729N_K+u}{3^{L_K+6}},
 \frac{729N_K+u+1}{3^{L_K+6}}
 \right).
 \label{p12-83-c9-s4:eq:729-subcilindros}
\end{equation}
Les \(C_{K,u}\) sont disjoints et
\[
 C(N_K,L_K)=\bigsqcup_{u=0}^{728}C_{K,u}.
\]
Le nombre \(729\) compte les coordonnées de cette fibre ambiante ; il ne compte
ni les germes, ni les émissions distinctes, ni les orbites, ni les états survivants.

\subsubsection{Comparaison diophantienne des extrémités rationnelles de l'encadrement}

Soit \(J=[\ell,h)\subset C(N_K,L_K)\), avec
\(\ell,h\in\mathbb Q\). Nous écrivons \(S_K=3^{L_K+6}\) et définissons
\begin{align}
 j_{\min}(J,K)
 &=\max\!\left(729N_K,\left\lfloor S_K\ell\right\rfloor\right),
 \label{p12-83-c9-s4:eq:jmin-diagonal}\\
 j_{\max}(J,K)
 &=\min\!\left(729N_K+728,\left\lceil S_Kh\right\rceil-1\right).
 \label{p12-83-c9-s4:eq:jmax-diagonal}
\end{align}

\begin{lemma}[Critère de localisation unitaire]
\label{p12-83-c9-s4:lem:criterio-localizacion-unitaria}
L’encadrement \(J\) est contenu dans un unique élément de la fibre
\eqref{p12-83-c9-s4:eq:729-subcilindros} si et seulement si
\begin{equation}
 \boxed{j_{\min}(J,K)=j_{\max}(J,K).}
 \label{p12-83-c9-s4:eq:jmin-igual-jmax}
\end{equation}
Dans ce cas, si \(j\) est la valeur commune,
\[
 u=j-729N_K,
 \qquad
 N_{K+1}=j.
\]
\end{lemma}

\begin{proof}
Les entiers \(j\) dont les intervalles élémentaires de longueur \(S_K^{-1}\)
rencontrent \(J\) forment l’intervalle entier
\[
 \left[
 \lfloor S_K\ell\rfloor,
 \lceil S_Kh\rceil-1
 \right]\cap
 [729N_K,729N_K+728].
\]
Celui-ci contient exactement un élément si et seulement si ses extrémités coïncident.
L’égalité identifie le numérateur du sous-cylindre et, en tronquant ses six
derniers trits, on retrouve \(N_K\).
\end{proof}

\subsection{Certification diagonale minimale}

La génération ne fixe pas à l’avance un même ordre analytique pour tous les
niveaux. Une fois l’extension compatible produite, l’ordre de
l’encadrement n’est augmenté que jusqu’à sa certification unitaire.

\begin{definition}[Récurrence diagonale de certification]
\label{p12-83-c9-s4:def:recurrencia-diagonal-constantes}
Une fois fixée par \(\operatorname{Ext}\) une famille de cylindres compatibles
\(C(N_K,L_K)\) dont la publication contient \(\rho_\chi\), on choisit un ordre
initial \(n_\chi(K_0-1)\geq1\) et nous définissons récursivement
\begin{equation}
 n_\chi(K)=
 \min\left\{
 n\geq n_\chi(K-1):
 j_{\min}(J_{\chi,n},K)=j_{\max}(J_{\chi,n},K)
 \right\},
 \label{p12-83-c9-s4:eq:nchi-K-minimo}
\end{equation}
et nous vérifions
\begin{equation}
 N_{\chi,K+1}
 =j_{\min}(J_{\chi,n_\chi(K)},K).
 \label{p12-83-c9-s4:eq:Nchi-K-mas-uno}
\end{equation}
L’ensemble
\[
 \Delta_\chi=
 \{(n_\chi(K),K):K\geq K_0\}
\]
est la diagonale analytique--nonadique du caractère \(\chi\).
\end{definition}

\begin{theorem}[Existence et unicité de la diagonale de certification]
\label{p12-83-c9-s4:thm:existencia-unicidad-diagonal}
Pour \(\chi\in\{\pi,e,\varphi\}\), la récurrence
\eqref{p12-83-c9-s4:eq:nchi-K-minimo} est définie à tout niveau fini. L’indice
\(n_\chi(K)\) et la coordonnée \(N_{\chi,K+1}\) sont uniques.
\end{theorem}

\begin{proof}
Supposons engendrée intérieurement l’extension cylindrique de niveau \(K\) qui contient
\(\rho_\chi\). Comme \(\rho_\chi\) est irrationnel, il n’appartient à aucune
frontière ternaire \(3^{-L}\mathbb Z\). Il est donc à l’intérieur d’un
unique sous-cylindre \(C_{K,u_*}\), à une distance positive de ses deux extrémités.
Les encadrements \(J_{\chi,n}\) contiennent \(\rho_\chi\) et leur diamètre
tend vers zéro ; il existe donc un \(n\) pour lequel l’encadrement est
contenu dans \(C_{K,u_*}\). L’ensemble sur lequel est pris le minimum dans
\eqref{p12-83-c9-s4:eq:nchi-K-minimo} est non vide et, par le bon ordre de
\(\mathbb N\), possède un plus petit élément. Le
\cref{p12-83-c9-s4:lem:criterio-localizacion-unitaria} prouve que l’encadrement la
localise de manière unitaire et que le numérateur localisé coïncide avec celui déjà
engendré. La récurrence sur \(K\) complète l’argument.
\end{proof}

\begin{corollary}[Absence de paramètres décimaux]
La diagonale est calculée avec des sommes, des produits, des puissances, des divisions euclidiennes,
des parties entières inférieures et supérieures d’entiers. Elle ne reçoit aucun chiffre décimal de \(\pi\), de \(e\) ou de
\(\varphi\) comme donnée de certification.
\end{corollary}

\subsection{Compatibilité avec la connexion nonadique conjointe}

Soit \(\widehat x_{\chi,K}\) l’état holonomique intégral du canal \(\chi\) au niveau \(K\). Ses coordonnées comprennent le préfixe, le cylindre, la phase, le quotient, la retenue, la mémoire verticale non bornée, le terme terminal et le registre de provenance. La phase résiduelle est
\[
 g_0(K)=[K]_9\in\mathbb Z/9\mathbb Z,
\]
tandis que l’étiquette publique est
\[
 g(K)=1+((K-1)\bmod9)\in\{1,\ldots,9\}.
\]
La mémoire verticale \(q_{\rm mem}(K)\in\mathbb Z_{\geq0}\) n’est pas réduite
modulo \(12\) ; seule l’est sa projection dodécaphasique
\[
 \beta(K)=q_{\rm mem}(K)\bmod12
\]
le fait.

Pour chaque \(K\), soit
\[
 \Gamma_{9,K}:\widehat X_K\longrightarrow\widehat X_{K+9}
\]
l’holonomie d’un tour de phase. Alors
\[
 g_0(K+9)=g_0(K),
 \qquad
 q_{\rm mem}(K+9)=q_{\rm mem}(K)+1,
 \qquad
 \beta(K+9)=\beta(K)+1\pmod{12}.
\]

\begin{theorem}[Retour de phase avec conservation généalogique]
\label{p12-83-c9-s4:thm:retorno-fase-conservacion-genealogica}
L’égalité de phase ne réinitialise pas l’état :
\[
 g_0(K+9)=g_0(K),
 \qquad
 \widehat x_{\chi,K+9}\ne\widehat x_{\chi,K}.
\]
Les profondeurs \(27\), \(54\) et \(108\) sont les compositions typées
\begin{align*}
 \Gamma_{27,K}
 &=\Gamma_{9,K+18}\circ\Gamma_{9,K+9}\circ\Gamma_{9,K},\\
 \Gamma_{54,K}
 &=\Gamma_{9,K+45}\circ\cdots\circ\Gamma_{9,K},\\
 \Gamma_{108,K}
 &=\Gamma_{9,K+99}\circ\cdots\circ\Gamma_{9,K}.
\end{align*}
En aucun cas il ne s’agit de réinitialisations du préfixe, de la retenue, du terminal ou de la
mémoire.
\end{theorem}

\begin{proof}
La première coordonnée se calcule dans \(\mathbb Z/9\mathbb Z\), tandis que \(q_{\rm mem}\) appartient à un monoïde non borné. Chaque application de \(\Gamma_{9,K}\) ajoute une unité de mémoire, six blocs par niveau intermédiaire et la composition correspondante de cocycles. La phase peut donc coïncider sans que l’état holonomique intégral coïncide. Les formules de profondeur indiquent les indices de chaque application et empêchent de composer comme si toutes étaient des endomorphismes du même espace.
\end{proof}

\subsubsection{Compression, expression et terme terminal}

La réalisation opératorielle d’un tour satisfait
\begin{equation}
 U_\Gamma J=JT+\eta,
 \qquad
 J^*\eta=0,
 \qquad
 I-T^*T=\eta^*\eta.
 \label{p12-83-c9-s4:eq:compresion-expresion-precision}
\end{equation}
La contraction visible stricte est réservée à
\[
 M_9=\frac59V_9;
\]
n’est pas identifié à l’opérateur de compression \(T\). Le télescopage de
profondeur \(L\) conserve le terme terminal :
\begin{equation}
 S_0=
 \sum_{r=0}^{L-1}M_9^{*r}D_9^*D_9M_9^r
 +M_9^{*L}S_0M_9^L.
 \label{p12-83-c9-s4:eq:telescopia-terminal-precision}
\end{equation}
Omettre le dernier terme avant de prouver sa disparition effacerait
de l’information de frontière et transformerait le retour de phase en une fausse
réinitialisation.

\subsection{Système inverse des histoires compatibles}

Soit \(\mathcal H_K^\chi\) l’ensemble des histoires enrichies admissibles
du canal \(\chi\) à la profondeur \(K\), et soit
\[
 p_{N,K}:\mathcal H_N^\chi\longrightarrow\mathcal H_K^\chi
 \qquad(N\geq K)
\]
la troncature. Nous définissons
\[
 \operatorname{Surv}_K^{(N)}(\chi)
 =p_{N,K}(\mathcal H_N^\chi)
\]
et
\begin{equation}
 \operatorname{Surv}_K(\chi)
 =\bigcap_{N\geq K}\operatorname{Surv}_K^{(N)}(\chi).
 \label{p12-83-c9-s4:eq:supervivencia-todas-profundidades}
\end{equation}
L’objet coinductif est
\begin{equation}
 X_\chi=
 \varprojlim_K
 \bigl(\operatorname{Surv}_K(\chi),p_{K+1,K}\bigr).
 \label{p12-83-c9-s4:eq:limite-inverso-caracteres}
\end{equation}

\begin{theorem}[Section compatible sous prolongeabilité enrichie]
\label{p12-83-c9-s4:thm:seccion-compatible-caracteres}
Supposons en outre que, sur chaque sous-cylindre numérique engendré, la
relation \(\operatorname{Ext}\) soit non vide, univoque et compatible avec les
troncatures. Alors \(\operatorname{Ext}\) et la survie produisent une famille
\[
 x_\chi=(\widehat x_{\chi,K})_{K\geq K_0}
\]
qui satisfait
\[
 p_{K+1,K}(\widehat x_{\chi,K+1})
 =\widehat x_{\chi,K}.
\]
Par conséquent, \(x_\chi\in X_\chi\).
\end{theorem}

\begin{proof}
Le sous-cylindre de niveau \(K+1\) a pour numérateur
\(N_{\chi,K+1}=729N_{\chi,K}+u\). La troncature élimine ses six derniers
trits et retrouve \(N_{\chi,K}\). Les autres coordonnées sont obtenues par les
lois fonctorielles de retenue, de mémoire, de terminal et de phase. L’hypothèse
d’univocité de \(\operatorname{Ext}\) fixe le sous-cylindre et l’extension
enrichie ; l’unicité locale du
\cref{p12-83-c9-s4:lem:criterio-localizacion-unitaria} certifie ensuite sa
localisation archimédienne, et la compatibilité de \(\operatorname{Ext}\)
fournit l’égalité de troncature.
\end{proof}

\subsection{Cellules décimales et précision prescrite}

Pour \(M\geq1\) et \(q\in\{0,\ldots,10^M-1\}\), soit
\[
 D_{M,q}=
 \left[
 \frac{q}{10^M},
 \frac{q+1}{10^M}
 \right).
\]
Un encadrement \(J=[\ell,h)\) est contenu dans une unique cellule décimale si
et seulement si
\begin{equation}
 \boxed{
 \left\lfloor10^M\ell\right\rfloor
 =\left\lceil10^Mh\right\rceil-1=q.}
 \label{p12-83-c9-s4:eq:criterio-celda-decimal}
\end{equation}

\begin{definition}[Niveau minimal de certification décimale]
Pour \(\chi\in\{\pi,e,\varphi\}\), nous définissons \(K_\chi(M)\) comme le plus petit
niveau pour lequel il existe un ordre analytique \(n\) et un entier \(q\) tels
que
\begin{enumerate}
 \item \(J_{\chi,n}\) satisfait
       \eqref{p12-83-c9-s4:eq:criterio-celda-decimal} ;
 \item \(J_{\chi,n}\) est contenu dans le cylindre engendré par
       l’état \(\widehat x_{\chi,K}\) ;
 \item la transition vers le niveau suivant satisfait
       \eqref{p12-83-c9-s4:eq:jmin-igual-jmax}.
\end{enumerate}
La paire \((K,n)\) est minimisée lexicographiquement.
\end{definition}

\begin{theorem}[Certification à précision décimale arbitraire des sorties coinductives]
\label{p12-83-c9-s4:thm:precision-decimal-arbitraria}
Pour chaque \(\chi\in\{\pi,e,\varphi\}\) et chaque \(M\geq1\), il existe
\(K_\chi(M)<\infty\). La section engendrée détermine un unique préfixe décimal
de \(M\) positions, et l’encadrement reconnaît et certifie ce préfixe. Si
\(M'<M\), le préfixe d’ordre \(M\) se tronque à celui
d’ordre \(M'\), et les cellules correspondantes sont emboîtées.
\end{theorem}

\begin{proof}
Les trois caractères sont irrationnels : \(\varphi\) l’est par son polynôme
quadratique à discriminant non carré ; l’irrationalité de \(e\) et de
\(\pi\) est tirée de leurs théorèmes classiques de caractérisation, postérieurs à
la construction finie. Aucun n’appartient donc à une frontière décimale ou
ternaire rationnelle. Pour un \(M\) fixé, la distance du caractère à l’ensemble
fini des frontières décimales pertinentes est positive. Les encadrements
rationnels contractent leur diamètre jusqu’à être contenus dans une unique cellule décimale.
Le \cref{p12-83-c9-s4:thm:existencia-unicidad-diagonal} couple cet encadrement à un
cylindre TPK et à une extension unique. Le minimum lexicographique existe par
bon ordre. La compatibilité de la section prouve l’emboîtement lorsque
\(M\) augmente.
\end{proof}

\begin{corollary}[Unicité réelle]
Les cellules emboîtées ont des diamètres \(10^{-M}\to0\). Leur intersection est
un singleton, identifié par les caractérisations précédentes à
\(\pi\), \(e\) ou \(\varphi\), respectivement.
\end{corollary}

\subsection{Équivalence de \(2\) réalisations finies de la loi uniforme}

Les témoins de \(1,000\) et de \(10,000\) positions ne définissent pas les
caractères et ne limitent pas le prolongement. Ce sont des exécutions de la même récurrence
uniforme à deux profondeurs différentes.

\begin{table}[htbp]
\centering
\caption{Dimensions exactes des exécutions certifiées, par canal.}
\label{p12-83-c9-s4:tab:dimensiones-ejecuciones-1000-10000}
\begin{tabular}{lrr}
\toprule
Objet compté & \(1,000\) positions & \(10,000\) positions\\
\midrule
Partitions exhaustives de la fibre ambiante & 396 & 3501\\
Trits transportados & 2376 & 21006\\
Paires de même phase \((K,K+9)\) & 387 & 3492\\
Tours nonadiques complets & 43 & 388\\
Éléments examinés par partition & 729 & 729\\
Extensions compatibles conservées par partition & 1 & 1\\
\bottomrule
\end{tabular}
\end{table}

Para \(10,000\) posiciones,
\[
 B=3501=389\cdot9,
 \qquad
 6B=21006,
 \qquad
 3^{6B}>10^{10000}.
\]
Les \(396\) premiers blocs de l’exécution longue coïncident avec
l’exécution de \(1,000\) positions. Dans chacune des \(3492\) paires de
même phase, la coordonnée de phase est préservée et le préfixe, la
profondeur, le cylindre et la mémoire sont prolongés. Le certificat enregistre aussi les
répétitions éventuelles de projections locales sans les confondre avec un retour
de l’état complet.

\subsection{Séparation entre génération et vérification}

\begin{definition}[Générateur structurel]
Le générateur reçoit le catalogue fini, \(L_0,L_1\), l’état initial, les
lois de connexion et la relation interne \(\operatorname{Ext}\). À chaque niveau :
\begin{enumerate}
 \item il énumère les \(729\) éléments de la fibre ambiante ;
 \item il applique \(\operatorname{Ext}\) et conserve l’unique extension compatible ;
 \item il actualise la phase, le quotient, la retenue, la mémoire et le terminal ;
 \item écrit le bloc ternaire, le registre de sélection et le registre généalogique.
\end{enumerate}
\end{definition}

\begin{definition}[Vérification ultérieure indépendante]
La vérification ultérieure indépendante reçoit les sorties closes du
générateur et les trois familles d'encadrements rationnels ; elle calcule
\(j_{\min}\) et \(j_{\max}\), certifie leur égalité avec l'extension déjà
produite, reconstruit les
partitions, vérifie les carrés de troncature, les couples
\((K,K+9)\), les sommes de partition
\[
 \operatorname{killed}_{\mathrm{izq}}+1+\operatorname{killed}_{\mathrm{der}}=729,
\]
l'intégrité des sorties et, à la fin, la concordance avec les développements de référence.
\end{definition}

\begin{theorem}[Isolement des données cibles]
\label{p12-83-c9-s4:thm:aislamiento-datos-objetivo-precision}
Les développements décimaux de référence et les données métrologiques
n’interviennent pas dans la sélection d’une extension compatible. La comparaison est
effectuée après l’écriture des trois préfixes.
\end{theorem}

\begin{proof}
L'entrée mathématique de chaque décision générative consiste en le catalogue
fini, \(L_0,L_1\), l'état hérité et la relation interne
\(\operatorname{Ext}\). Les extrémités rationnelles, puissances de \(3\), planchers et
plafonds relèvent de la vérification ultérieure. Le générateur termine avant
que celle-ci n'ouvre les encadrements ou n'exécute les contrôles décimaux. Le sens
du flux de données est
\[
 \text{structure}\longrightarrow
 \text{préfixes et registres}\longrightarrow
 \text{vérification ultérieure indépendante};
\]
il n’existe pas de flèche en sens contraire.
\end{proof}

\subsection{Obstructions différenciées de la génération et de sa certification à précision arbitraire}

\begin{criterion}[Réfutation de la génération à une précision arbitraire]
Les défaillances d'orbite, d'extension, de troncature ou de retour réfutent la génération ;
les défaillances d'encadrement, d'ordre ou de contrôle décimal réfutent le certificat
analytique correspondant ; et une divergence entre les deux réfute
l'identification conjointe de la sortie publiée. Avec cette séparation des types,
l'affirmation complète est réfutée si :
\begin{enumerate}
 \item on identifie l’un quelconque des indices \(n,K,L_K,M\) sans un
       théorème de liaison ;
 \item un encadrement cesse de contenir son caractère ou ne contracte pas son diamètre ;
 \item une étape satisfait \(j_{\min}\ne j_{\max}\) après l’ordre
       déclaré minimal ;
 \item une extension ne se tronque pas au cylindre du niveau précédent ;
 \item le retour de phase réinitialise le préfixe, la retenue, le terminal ou la mémoire ;
 \item le terme terminal est omis dans
       \eqref{p12-83-c9-s4:eq:telescopia-terminal-precision} sans que son annulation soit démontrée ;
 \item les \(396\) premiers blocs de l’exécution longue diffèrent de
       l’exécution courte ;
 \item un contrôle décimal ou métrologique est lu avant la clôture de la sortie ;
 \item un chiffre imprimé diffère de la section mathématique certifiée ;
 \item la preuve dépend du fait que la précision demandée soit \(1,000\) ou
       \(10,000\), au lieu d’un \(M\) arbitraire.
\end{enumerate}
\end{criterion}
 
\chapter{Clôture dodécaphasique et deux dérivations internes de \(\alpha\)}
\label{chap:83-10}

% Unidad U016. Registro firmado, lectura terminal y reconstruccion de K.
% Redaccion autonoma desde la autoridad material 1063 y su sucesor V2.

\section{Génération initiale de l'état dodécaphasique signé et reconstruction du sceau \(K\)}
\label{p12-83-c10-s1:chap:registro-dodecafasico-hadamard-k}

La construction du sceau entier ne commence ni par \(\alpha\), ni par son inverse, ni par une table métrologique. Son domaine est l’état holonomique terminal du Topological Phase Kernel obtenu après la chaîne déjà construite
\begin{equation}
 \boxed{
 \mathrm{APP}\longrightarrow\mathrm{TRIT}\longrightarrow\mathrm{TPK}
 \longrightarrow\Omega_{\rm seed}
 \xrightarrow{T_{\min}}\mathcal U_6^{\rm cat}
 \xrightarrow{\rho_3}W_{243}
 \longrightarrow w_{30}\longrightarrow R_{36}
 \longrightarrow G_9\longrightarrow x_{\rm term}.}
 \label{p12-83-c10-s1:eq:cadena-upstream-registro-dodecafasico}
\end{equation}
Les flèches précédentes ont été définies avec leurs domaines respectifs :
paires de germes, signatures, émissions, mots ternaires, frontières,
histoires compatibles et holonomie nonadique. Cette unité ne choisit à nouveau
aucune de ces données. Elle construit les cartes entières du même état
terminal et démontre leur équivalence.

\subsection{Traces orientées sur des fenêtres de longueur \(108\)}

Soit
\[
 \Theta_{108}=\mathbb Z/108\mathbb Z
\]
le calendrier observable. Ses douze fenêtres nonadiques sont
\[
 E_m=\{9m+1,\ldots,9m+9\},
 \qquad 0\leq m<12.
\]
La partition est disjointe et exhaustive :
\[
 \{1,\ldots,108\}=\bigsqcup_{m=0}^{11}E_m.
\]
L’ensemble des codes de direction qui publient un événement est fixé comme
\[
 D_{\rm evt}=\{2,4,6,8\}.
\]
L’orientation des douze fenêtres est
\begin{equation}
 \Sigma_{12}=(+,+,+,-,-,-,+,+,+,-,-,-).
 \label{p12-83-c10-s1:eq:signatura-doce-ventanas}
\end{equation}
Si \(d_t\) est la direction codée au pas \(t\), l’événement signé
de la fenêtre \(m\) est
\begin{equation}
 e_m=\Sigma_{12}(m)\,
 \#\{t\in E_m:d_t\in D_{\rm evt}\}.
 \label{p12-83-c10-s1:eq:evento-firmado-ventana}
\end{equation}
La formule distingue le nombre non négatif d’événements et le signe de
l’orientation. Elle ne remplace pas une direction cardinale par un signe : les deux
coordonnées restent présentes dans l’état TPK.

\begin{definition}[Registre enrichi dodécaphasique]
\label{p12-83-c10-s1:def:registro-enriquecido-dodecafasico}
Le registre transporté par la trace est
\[
 \mathcal R_{12}(x)=
 \bigl((E_m,\tau_m,\sigma_m,\kappa_m,\Lambda_m)\bigr)_{m=0}^{11},
\]
où \(\tau_m\) est l’orientation TRIT, \(\sigma_m\) la signature de
frontière, \(\kappa_m\) la retenue et \(\Lambda_m\) le segment de registre généalogique
qui conserve la route. Il existe une chaîne de projections
\[
 \mathcal X_{\rm TPK}^{\rm enr}
 \longrightarrow\mathcal R_{12}
 \longrightarrow\Theta_{108},
\]
mais aucune des deux flèches n’est une identification.
\end{definition}

En particulier, \(\Theta_{108}\) conserve uniquement la position du calendrier.
Le registre \(\mathcal R_{12}\) conserve en outre l’orientation, la signature,
la retenue et la provenance. L’état complet conserve encore le cylindre,
la frontière, le quotient, le terminal et la mémoire verticale.

\subsection{Carte entière de l’opposition dodécaphasique}

Pour \(0\leq i<6\), les positions opposées déterminent
\begin{equation}
 p_i=e_i+e_{i+6},
 \qquad
 a_i=e_i-e_{i+6}.
 \label{p12-83-c10-s1:eq:pares-opuestos-registro}
\end{equation}
La charge centrale et les cinq marges relatives sont
\begin{equation}
 s_C=\sum_{i=0}^{5}p_i,
 \qquad
 M_i=p_i-p_5\quad(0\leq i<5).
 \label{p12-83-c10-s1:eq:carga-margenes-registro}
\end{equation}

\begin{definition}[Carte intégrale \(1+5+6\)]
\label{p12-83-c10-s1:def:carta-integral-156}
On définit
\begin{equation}
 \mathcal L_{12}(e_0,\ldots,e_{11})
 =\bigl(s_C,M_0,\ldots,M_4,a_0,\ldots,a_5\bigr).
 \label{p12-83-c10-s1:eq:carta-integral-156}
\end{equation}
La partition des coordonnées est
\[
 1+5+6=12.
\]
\end{definition}

\begin{theorem}[Inversion exacte de la carte entière]
\label{p12-83-c10-s1:thm:inversion-carta-integral-156}
Un tuple
\[
 \ell=(s_C,M_0,\ldots,M_4,a_0,\ldots,a_5)\in\mathbb Z^{12}
\]
appartient à l’image de \(\mathcal L_{12}\) si et seulement si
\begin{align}
 s_C-\sum_{i=0}^{4}M_i&\equiv0\pmod6,
 \label{p12-83-c10-s1:eq:congruencia-seis-carta}\\
 p_i&\equiv a_i\pmod2
 \quad(0\leq i<6),
 \label{p12-83-c10-s1:eq:congruencia-paridad-carta}
\end{align}
où
\begin{align}
 p_5&=\frac{s_C-\sum_{i=0}^{4}M_i}{6},
 \label{p12-83-c10-s1:eq:inversa-p5-carta}\\
 p_i&=p_5+M_i\quad(0\leq i<5).
 \label{p12-83-c10-s1:eq:inversa-pi-carta}
\end{align}
Dans ce cas, l’unique antécédent entier est
\begin{equation}
 e_i=\frac{p_i+a_i}{2},
 \qquad
 e_{i+6}=\frac{p_i-a_i}{2}.
 \label{p12-83-c10-s1:eq:inversa-eventos-carta}
\end{equation}
\end{theorem}

\begin{proof}
En sommant \(M_i=p_i-p_5\) pour \(0\leq i<5\), nous obtenons
\[
 \sum_{i=0}^{4}M_i
 =\sum_{i=0}^{4}p_i-5p_5
 =s_C-6p_5,
\]
ce qui impose \eqref{p12-83-c10-s1:eq:inversa-p5-carta} et la congruence modulo six.
Une fois les six \(p_i\) reconstruits, le système
\[
 p_i=e_i+e_{i+6},
 \qquad a_i=e_i-e_{i+6}
\]
admet l’unique solution \eqref{p12-83-c10-s1:eq:inversa-eventos-carta}. Cette solution est
entière exactement lorsque \(p_i\) et \(a_i\) ont la même parité.
Les deux familles de congruences sont donc nécessaires et suffisantes.
\end{proof}

Ce théorème est antérieur à la clôture numérique de \(\alpha\). Il explique comment la
trace orientée produit une carte entière réversible, et pourquoi il ne suffit pas
de conserver le calendrier réduit.

\subsection{Sélection de l’état terminal commun}

Les trois sections compatibles atteignent la profondeur terminale avec des catalogues
de cardinalités
\[
 |\mathcal T_\pi|=196,
 \qquad
 |\mathcal T_e|=197,
 \qquad
 |\mathcal T_\varphi|=196.
\]
Le produit contient
\begin{equation}
 196\cdot197\cdot196=7\,567\,952
 \label{p12-83-c10-s1:eq:cardinal-producto-terminal}
\end{equation}
triplets ordonnés.

La marge de colonnes du lecteur terminal est
\[
 q_\partial=(1,2,2,1,2,0)\in\mathbb F_3^6,
\]
et l’orientation des lignes du canal conjoint est
\[
 \sigma=(1,1,-1)=(1,1,2)\in\mathbb F_3^3.
\]
La première donnée agit sur les colonnes ; la seconde sur les lignes. L’une ne se déduit pas
de l’autre.

Les signatures locales de Gram, de norme, de poids et de distance réduisent
\eqref{p12-83-c10-s1:eq:cardinal-producto-terminal} à \(331\) triplets. La marge
\(q_\partial\) laisse deux indices :
\[
 (120,197,157),
 \qquad
 (129,197,148).
\]
Leurs matrices sont
\[
 B_0=
 \begin{pmatrix}
 0&2&0&2&2&0\\
 0&0&1&2&2&0\\
 1&0&1&0&1&0
 \end{pmatrix},
 \qquad
 B_1=
 \begin{pmatrix}
 0&2&1&0&2&0\\
 0&0&1&2&2&0\\
 1&0&0&2&1&0
 \end{pmatrix}.
\]
Les sommes des lignes sont
\[
 (0,2,0),
 \qquad
 (2,2,1).
\]
Comme
\[
 \langle\sigma\rangle
 =\{(0,0,0),(1,1,2),(2,2,1)\},
\]
seul \(B_1\) satisfait la condition axiale.

\begin{theorem}[Unicité de la publication terminale visible]
\label{p12-83-c10-s1:thm:unicidad-publicacion-terminal-visible}
La sélection conjointe détermine
\begin{equation}
 B_{\rm term}=
 \begin{pmatrix}
 021020\\001220\\100210
 \end{pmatrix},
 \label{p12-83-c10-s1:eq:matriz-terminal-visible}
\end{equation}
correspondant au triplet \((129,197,148)\). Aucune évaluation décimale de
\(\alpha\) n’intervient dans le recensement ni dans les deux conditions de sélection.
\end{theorem}

\begin{proof}
Le recensement et la marge laissent exactement \(B_0,B_1\). La charge de \(B_0\)
n’appartient pas à \(\langle\sigma\rangle\), tandis que celle de \(B_1\) est
\(2\sigma\). La seconde candidate est donc la seule qui satisfasse
simultanément les deux exigences typées.
\end{proof}

\subsection{Publications archimédienne et exceptionnelle de l'état commun : codomaines non équivalents et généalogie partagée}

Soit
\[
 x_{\rm term}\in
 \mathcal X_{\rm TPK}^{\rm term,enr}
 \subseteq\mathcal X_{\rm TPK}^{[108]}.
\]
La lecture visible est
\[
 \pi_{\rm term}(x_{\rm term})=B_{\rm term}.
\]
La carte signée conserve une autre coordonnée du même état.

\begin{definition}[Sous-espace signé terminal]
\label{p12-83-c10-s1:def:subespacio-firmado-terminal}
Soit \(\mathcal H_{12}\) la transformation entière qui sera construite dans la section suivante. Définissons
\begin{equation}
 \mathcal X_{\rm TPK}^{[108],\rm sgn}
 :=\{(x_{\rm core},K,U):U=\mathcal H_{12}K\}
 \label{p12-83-c10-s1:eq:subespacio-firmado-terminal}
\end{equation}
et le module d’expression
\begin{equation}
 \mathcal U_{12}^{\rm sgn}
 :=\operatorname{im}
 \bigl(\mathcal H_{12}:\mathbb Z^{12}\to\mathbb Z^{12}\bigr).
 \label{p12-83-c10-s1:eq:modulo-publicacion-firmada}
\end{equation}
La lecture signée est
\begin{equation}
 R_{\rm sgn}:
 \mathcal X_{\rm TPK}^{[108],\rm sgn}
 \longrightarrow\mathcal U_{12}^{\rm sgn},
 \qquad
 R_{\rm sgn}(x_{\rm core},K,U)=U.
 \label{p12-83-c10-s1:eq:lector-firmado-terminal}
\end{equation}
\end{definition}

L’écriture conjointe correcte est
\begin{equation}
 x_{\rm term}
 \xrightarrow{(\pi_{\rm term},R_{\rm sgn})}
 (B_{\rm term},U_{12}^{\rm sgn}).
 \label{p12-83-c10-s1:eq:doble-publicacion-terminal}
\end{equation}
Aucune flèche \(B_{\rm term}\to U_{12}^{\rm sgn}\) n’est interposée : la matrice visible a précisément oublié les coordonnées d’orientation, d’opposition, de quotient, de retenue et de registre généalogique utilisées par la carte signée. La source des deux expressions est l’état holonomique intégral produit par \eqref{p12-83-c10-s1:eq:cadena-upstream-registro-dodecafasico}.

L’évaluation du registre signé donne
\begin{equation}
 \boxed{
 U=(2378,1406,2479,-452,998,-551,
 -668,-204,-371,-322,-28,-997).}
 \label{p12-83-c10-s1:eq:vector-u-firmado-autonomo}
\end{equation}

\subsection{Naturalité par rapport à la profondeur}

Pour \(N'\geq N\geq108\), les troncatures du système signé sont
\begin{equation}
 \tau_{N',N}^{\rm sgn}
 (x_{\rm core}^{[N']},K,U)
 =\bigl(\tau_{N',N}x_{\rm core}^{[N']},K,U\bigr).
 \label{p12-83-c10-s1:eq:truncamiento-preserva-ku-autonomo}
\end{equation}

\begin{theorem}[Carré de naturalité de la lecture signée]
\label{p12-83-c10-s1:thm:naturalidad-lector-firmado-autonomo}
On a
\begin{equation}
 \boxed{
 R_{{\rm sgn},N}\circ\tau_{N',N}^{\rm sgn}
 =\operatorname{id}_{\mathcal U_{12}^{\rm sgn}}
  \circ R_{{\rm sgn},N'}.}
 \label{p12-83-c10-s1:eq:cuadrado-naturalidad-lector-firmado}
\end{equation}
\end{theorem}

\begin{proof}
Les deux membres, appliqués à \((x_{\rm core}^{[N']},K,U)\), renvoient littéralement \(U\). La troncature modifie la trajectoire ultérieure mais conserve la coordonnée terminale déjà fixée. La commutativité est une égalité d’applications, non une analogie entre diagrammes.
\end{proof}

\subsection{Orbites de pas \(3\) et transformée de Hadamard}

L’ordre cyclique du sceau se répartit dans les trois orbites
\begin{align}
 K^{(1)}&=(K_1,K_4,K_7,K_{10}),
 \label{p12-83-c10-s1:eq:orbita-k-1}\\
 K^{(2)}&=(K_2,K_5,K_8,K_{11}),
 \label{p12-83-c10-s1:eq:orbita-k-2}\\
 K^{(3)}&=(K_3,K_6,K_9,K_{12}).
 \label{p12-83-c10-s1:eq:orbita-k-3}
\end{align}
Soit
\begin{equation}
 H_4=
 \begin{pmatrix}
 1&1&1&1\\
 1&1&-1&-1\\
 1&-1&1&-1\\
 1&-1&-1&1
 \end{pmatrix}.
 \label{p12-83-c10-s1:eq:matriz-hadamard-cuatro-autonoma}
\end{equation}
Si \(P_{\rm orb}\) réordonne les douze coordonnées conformément à
\eqref{p12-83-c10-s1:eq:orbita-k-1}--\eqref{p12-83-c10-s1:eq:orbita-k-3}, la transformation globale est
\begin{equation}
 \mathcal H_{12}
 =P_{\rm orb}^{-1}
 (H_4\oplus H_4\oplus H_4)P_{\rm orb}.
 \label{p12-83-c10-s1:eq:hadamard-global-doce}
\end{equation}

Les trois blocs de \eqref{p12-83-c10-s1:eq:vector-u-firmado-autonomo}, lus par
orbites, sont
\begin{align*}
 U^{(1)}&=(2378,-452,-668,-322),\\
 U^{(2)}&=(1406,998,-204,-28),\\
 U^{(3)}&=(2479,-551,-371,-997).
\end{align*}

\begin{lemma}[Quart de tour algébrique de Hadamard]
\label{p12-83-c10-s1:lem:cuarto-inversa-hadamard}
La matrice \(H_4\) satisfait
\begin{equation}
 H_4^2=4I_4,
 \qquad
 H_4^{-1}=\frac14H_4,
 \qquad
 \det H_4=-16.
 \label{p12-83-c10-s1:eq:identidades-hadamard-cuatro}
\end{equation}
\end{lemma}

\begin{proof}
Les lignes de \(H_4\) sont orthogonales et chacune possède une norme au carré égale à quatre.
Comme la matrice est symétrique, \(H_4^{\mathsf T}H_4=H_4^2=4I_4\). La
formule de l’inverse et le déterminant s’en déduisent immédiatement.
\end{proof}

\begin{theorem}[Reconstruction entière unique du sceau]
\label{p12-83-c10-s1:thm:reconstruccion-entera-k-hadamard}
L’inversion
\begin{equation}
 K^{(r)}=\frac14H_4U^{(r)}
 \qquad(r=1,2,3)
 \label{p12-83-c10-s1:eq:inversion-bloques-hadamard}
\end{equation}
produit
\begin{align*}
 K^{(1)}&=(234,729,621,794),\\
 K^{(2)}&=(543,659,058,146),\\
 K^{(3)}&=(140,824,914,601).
\end{align*}
En défaisant l’ordre orbital, on obtient
\begin{equation}
 \boxed{
 K=(234,543,140,729,659,824,621,058,914,794,146,601).}
 \label{p12-83-c10-s1:eq:sello-k-doce-autonomo}
\end{equation}
C'est l'unique préimage rationnelle de \(U\), et elle appartient à \(\mathbb Z^{12}\).
\end{theorem}

\begin{proof}
Par \eqref{p12-83-c10-s1:eq:identidades-hadamard-cuatro}, l’antécédent rationnel est unique.
Les multiplications explicites sont
\begin{align*}
 H_4U^{(1)}&=(936,2916,2484,3176),\\
 H_4U^{(2)}&=(2172,2636,232,584),\\
 H_4U^{(3)}&=(560,3296,3656,2404).
\end{align*}
Toutes les coordonnées sont divisibles par quatre. La division et le réentrelacement
produisent \eqref{p12-83-c10-s1:eq:sello-k-doce-autonomo}.
\end{proof}

Le facteur \(1/4\) n’est pas choisi pour ajuster une sortie : c’est l’inverse
imposé par \(H_4^2=4I_4\). Ce n’est pas non plus la racine carrée de \(-1\). La
structure complexe interne \(A_W^2=-I\) apparaîtra après la construction de la
matrice de Paley ; ce sont deux faits de types différents.

\subsection{Structure entière et forme normale de Smith}

\begin{theorem}[Image entière de \(H_4\)]
\label{p12-83-c10-s1:thm:smith-hadamard-autonomo}
La forme normale de Smith est
\begin{equation}
 \operatorname{SNF}(H_4)=\operatorname{diag}(1,2,2,4).
 \label{p12-83-c10-s1:eq:snf-hadamard-autonoma}
\end{equation}
En conséquence,
\begin{equation}
 \operatorname{coker}\mathcal H_{12}
 \cong(\mathbb Z/2\mathbb Z)^6
 \oplus(\mathbb Z/4\mathbb Z)^3,
 \label{p12-83-c10-s1:eq:cociente-hadamard-doce}
\end{equation}
et l’image de \(\mathcal H_{12}\) a pour indice
\[
 16^3=4096
\]
dans \(\mathbb Z^{12}\). Pour \(u\in\mathbb Z^4\),
\begin{equation}
 u\in H_4\mathbb Z^4
 \quad\Longleftrightarrow\quad
 H_4u\equiv0\pmod4.
 \label{p12-83-c10-s1:eq:criterio-imagen-hadamard}
\end{equation}
\end{theorem}

\begin{proof}
Le plus grand commun diviseur des entrées est un. Les mineurs d’ordre deux sont
pairs et l’un d’eux a pour valeur absolue deux ; ceux d’ordre trois sont des multiples de
quatre et l’un d’eux a pour valeur absolue quatre ; le déterminant a pour module
seize. Les diviseurs déterminantaux sont \(1,2,4,16\), d’où l’on obtient
les facteurs invariants \(1,2,2,4\). La somme directe triple répète les
facteurs et multiplie les indices. Enfin,
\(H_4^{-1}=H_4/4\), de sorte que l’antécédent est entier exactement sous la
congruence \eqref{p12-83-c10-s1:eq:criterio-imagen-hadamard}.
\end{proof}

\subsection{2e système de coordonnées : différences de pas \(3\) et \(4\)}

Avec des indices cycliques modulo douze, nous définissons
\begin{equation}
 (D_3K)_m=K_m-K_{m+3},
 \qquad
 (D_4K)_m=K_m-K_{m+4},
 \qquad
 Q(K)=\sum_{m=1}^{12}K_m.
 \label{p12-83-c10-s1:eq:extractor-diferencias-k}
\end{equation}
Pour \eqref{p12-83-c10-s1:eq:sello-k-doce-autonomo},
\begin{align*}
 D_3K={}&(-495,-116,-684,108,601,-90,\\
         &\qquad-173,-88,313,560,-397,461),\\
 D_4K={}&(-425,-281,-481,671,-255,30,\\
         &\qquad475,-543,680,251,6,-128),\\
 Q(K)={}&6263.
\end{align*}

\begin{theorem}[Complétude de l’extracteur transversal]
\label{p12-83-c10-s1:thm:completitud-extractor-transversal-k}
L'opérateur
\[
 \mathcal E_{\rm dod}=(D_3,D_4,Q):
 \mathbb Z^{12}\longrightarrow\mathbb Z^{25}
\]
a rang douze et pour forme normale de Smith
\begin{equation}
 \operatorname{SNF}(\mathcal E_{\rm dod})
 =\operatorname{diag}(\underbrace{1,\ldots,1}_{11},12).
 \label{p12-83-c10-s1:eq:snf-extractor-transversal-k}
\end{equation}
En particulier, \((D_3K,D_4K,Q(K))\) détermine un unique \(K\).
\end{theorem}

\begin{proof}
Si \(D_3v=D_4v=0\), alors \(v\) est constant le long des pas
trois et quatre. Comme \(\gcd(3,4)=1\), ces pas engendrent tout
\(\mathbb Z/12\mathbb Z\), de sorte que
\[
 \ker D_3\cap\ker D_4=\langle\mathbf1\rangle.
\]
La charge totale ne s’annule pas sur cette droite, car \(Q(\mathbf1)=12\). Le rang
est douze. Les lignes de \((D_3,D_4)\) sont les incidences orientées d’un graphe
de Cayley connexe ; un arbre couvrant apporte onze facteurs invariants
unitaires. Tout mineur maximal doit contenir la ligne \(Q\), et les opérations
unimodulaires réduisent sa dernière contribution à douze. Il existe un mineur maximal de
module douze ; le dernier facteur est donc exactement douze.
\end{proof}

La reconstruction de \(U\) à partir de ces coordonnées ne nécessite pas une nouvelle multiplication par \(H_4\). Pour \(r=1,2,3\), soit
\begin{equation}
 B_r=\sum_{j=0}^{3}(D_4K)_{r+3j}.
 \label{p12-83-c10-s1:eq:br-diferencias-k}
\end{equation}
Alors
\[
 (B_1,B_2,B_3)=(972,-1073,101).
\]
Les sommes orbitales sont
\begin{align}
 A_1&=\frac{Q+2B_1+B_2}{3}=2378,
 \label{p12-83-c10-s1:eq:a1-reconstruccion-u}\\
 A_2&=A_1-B_1=1406,
 \label{p12-83-c10-s1:eq:a2-reconstruccion-u}\\
 A_3&=A_2-B_2=2479.
 \label{p12-83-c10-s1:eq:a3-reconstruccion-u}
\end{align}
Si \(d_{r,j}=(D_3K)_{r+3j}\), le bloc signé de l'orbite \(r\) est
\begin{equation}
 \bigl(A_r,
 d_{r,1}-d_{r,3},
 d_{r,0}+d_{r,2},
 d_{r,0}-d_{r,2}\bigr),
 \label{p12-83-c10-s1:eq:bloque-firmado-desde-diferencias}
\end{equation}
qui reproduit les trois blocs \(U^{(r)}\).

\subsection{Tétrade de seuil du sceau}

La complétion décimale locale satisfait
\[
 1000=729+271.
\]
Le support des coordonnées du sceau qui atteignent le seuil \(729\) est
\begin{equation}
 \boxed{
 B_K:=\{m:K_m\geq729\}=\{4,6,9,10\}.}
 \label{p12-83-c10-s1:eq:tetrada-umbral-k}
\end{equation}
Cette tétrade s’obtient après la construction de \(K\). Dans cette unité, elle n’est pas
encore appelée étoile de Witt : un tel objet ne sera défini qu’après
la construction du code ternaire, de ses hexades et du système de Steiner.

\subsection{Indépendance de la construction généalogique à l'égard de données métrologiques ou physiques externes}

\begin{theorem}[Isolement de \(U\) et \(K\)]
\label{p12-83-c10-s1:thm:aislamiento-u-k-alpha-codata}
La composition
\begin{equation}
 \Omega_{\rm seed}\longrightarrow x_{\rm term}
 \xrightarrow{R_{\rm sgn}}U
 \xrightarrow{\mathcal H_{12}^{-1}}K
 \label{p12-83-c10-s1:eq:composicion-upstream-u-k}
\end{equation}
ne reçoit en entrée ni \(\alpha\), ni \(\alpha^{-1}\), ni CODATA, ni une suite
décimale cible. Le sceau \(K\) est déterminé avant la définition de la
récurrence qui produira \(\alpha\).
\end{theorem}

\begin{proof}
La première partie de la composition utilise uniquement des germes, des émissions, des lecteurs ternaires, des opérateurs \(L_0,L_1\), des relations de survie, du transport de phase et les cartes de l’état TPK. La lecture \(R_{\rm sgn}\) est une projection du registre terminal, et la dernière flèche est la transformation linéaire \(\mathcal H_{12}^{-1}=\mathcal H_{12}/4\) sur son image entière. Aucune formule ne contient de variable physique ni de table métrologique. La récurrence en base mille n’est introduite que dans l’unité suivante, après fixation de \(K\).
\end{proof}

\subsection{Obstructions du registre dodécaphasique et du sceau \(K\) : naturalité, intégrité de Hadamard et isolement métrologique}

\begin{criterion}[Réfutation du registre et du sceau]
La construction est réfutée si l'un des faits
suivants est vérifié :
\begin{enumerate}
 \item les fenêtres \(E_m\) ne partitionnent pas les cent huit pas ;
 \item la carte \(\mathcal L_{12}\) ne respecte pas son inverse ou ses congruences ;
 \item le recensement terminal ne produit pas \(331\) triplets, deux matrices et une unique
       charge axiale ;
 \item on tente d’obtenir \(U\) à partir de \(B_{\rm term}\) après avoir oublié les
       coordonnées signées ;
 \item le carré de naturalité
       \eqref{p12-83-c10-s1:eq:cuadrado-naturalidad-lector-firmado} échoue ;
 \item \(H_4^2\neq4I_4\), un antécédent n’est pas entier ou le vecteur
       reconstruit diffère de \eqref{p12-83-c10-s1:eq:sello-k-doce-autonomo} ;
 \item la forme normale de Smith diffère de
       \eqref{p12-83-c10-s1:eq:snf-hadamard-autonoma} ;
 \item l'extracteur \((D_3,D_4,Q)\) perd du rang ou ne reproduit pas \(U\) ;
 \item \(B_K\) est introduit avant la construction de \(K\), ou on lui attribue une
       structure de Witt avant la définition de \(W_{12}\) ;
 \item un chiffre de \(\alpha\) ou une référence métrologique intervient dans
       l’une quelconque des flèches de \eqref{p12-83-c10-s1:eq:composicion-upstream-u-k}.
\end{enumerate}
\end{criterion}


% Unidad U017. Cierre afín, prolongación y semántica de frontera de alpha.
% Redacción autónoma desde la autoridad material 1063 y su sucesor V2.

\section{Construction initiale de la frontière \texorpdfstring{\(662/663\)}{662/663} au moyen du cocycle dodécaphasique de retenue}
\label{p12-83-c10-s2:chap:acarreo-alpha-frontera-662-663}

\subsection{Domaine arithmétique de la clôture}

La construction reçoit quatre coordonnées dodécaphasiques obtenues dans les
unités précédentes :
\[
\begin{aligned}
 P={}&(141,592,653,589,793,238,462,643,383,279,502,884),\\
 E={}&(718,281,828,459,045,235,360,287,471,352,662,497),\\
 \Phi={}&(618,033,988,749,894,848,204,586,834,365,638,117),\\
 K={}&(234,543,140,729,659,824,621,058,914,794,146,601).
\end{aligned}
\]
Les trois premiers vecteurs sont les douze premières triades des lectures archimédiennes fractionnaires de \(\pi\), \(e\) et \(\varphi\). Le quatrième est la coordonnée dodécaphasique de l’état holonomique terminal du Topological Phase Kernel, construite avant l’ouverture de l’opération de retenue.

Nous fixons la base
\[
 B:=10^3=1000
\]
et la concaténation
\begin{equation}
 \iota(v_1,\ldots,v_{12})
 :=\sum_{m=1}^{12}v_mB^{12-m}.
 \label{p12-83-c10-s2:eq:concatenacion-base-mil-alpha}
\end{equation}
\Needspace{8\baselineskip}
En particulier,
\[
\begin{aligned}
 \iota(P)&=141592653589793238462643383279502884,\\
 \iota(E)&=718281828459045235360287471352662497,\\
 \iota(\Phi)&=618033988749894848204586834365638117,\\
 \iota(K)&=234543140729659824621058914794146601.
\end{aligned}
\]

L’opération étudiée dans cette unité est l’application
\begin{equation}
 \mathfrak C_{12}:
 (P,E,\Phi,K,c_{13})
 \longmapsto(A,C),
 \label{p12-83-c10-s2:eq:mapa-cierre-dodecafasico-alpha}
\end{equation}
où \(A=(A_1,\ldots,A_{12})\) est le mot de sortie et
\[
 C=(c_1,\ldots,c_{13})
\]
est le système complet des treize quotients euclidiens : douze quotients
internes et une condition à la frontière droite.

\subsection{Division euclidienne orientée de droite à gauche}

Pour \(m=12,11,\ldots,1\), on définit
\begin{equation}
\begin{aligned}
 s_m&:=P_m+E_m-\Phi_m-K_m+c_{m+1},\\
 A_m&:=s_m\bmod B,
       \qquad 0\leq A_m<B,\\
 c_m&:=\frac{s_m-A_m}{B}
      =\left\lfloor\frac{s_m}{B}\right\rfloor .
\end{aligned}
\label{p12-83-c10-s2:eq:recurrencia-corta-alpha}
\end{equation}
La deuxième égalité pour \(c_m\) est également valable lorsque \(s_m<0\),
car on utilise la division euclidienne avec reste dans
\(\{0,\ldots,B-1\}\). C'est pourquoi,
\[
 -431=569-1000,
 \qquad
 -717=283-1000,
 \qquad
 -1200=800-2\cdot1000,
\]
et les emprunts négatifs n’introduisent aucune ambiguïté.

\begin{theorem}[Existence et unicité de la retenue orientée]
\label{p12-83-c10-s2:thm:unicidad-acarreo-orientado-alpha}
Une fois \(P,E,\Phi,K\) et la retenue terminale \(c_{13}\) fixés, la récurrence
\eqref{p12-83-c10-s2:eq:recurrencia-corta-alpha} détermine de manière unique
\[
 (A_1,\ldots,A_{12};c_1,\ldots,c_{12}).
\]
\end{theorem}

\begin{proof}
La ligne \(m=12\) est déterminée par la division euclidienne de
\[
 P_{12}+E_{12}-\Phi_{12}-K_{12}+c_{13}
\]
par \(B\). Une fois \(c_{12}\) connu, la ligne \(m=11\) est déterminée par
la même propriété. La récurrence descendante atteint \(m=1\). À chaque
étape, l’unicité du quotient et du reste euclidiens empêche deux sorties
distinctes.
\end{proof}

\subsection{Fenêtre dodécaphasique de l'état signé et clôture réversible du cycle}

La fenêtre finie est définie par la condition terminale
\[
 c_{13}^{\rm fin}=0.
\]
L’évaluation complète est la suivante. Le tableau est imprimé dans l’ordre
croissant, bien qu’il doive être lu causalement de la dernière ligne vers la
première.

\begingroup
\scriptsize
\setlength{\tabcolsep}{3.2pt}
\begin{longtable}{@{}r|rrrr|r|r|r|r@{}}
\caption{Les douze divisions euclidiennes et les treize retenues de la fenêtre
dodécaphasique finie.}
\label{p12-83-c10-s2:tab:acarreo-alpha-finito-autonomo}\\
\toprule
\(m\)&\(P_m\)&\(E_m\)&\(\Phi_m\)&\(K_m\)&
\(c_{m+1}\)&\(s_m\)&\(A_m\)&\(c_m\)\\
\midrule
\endfirsthead
\toprule
\(m\)&\(P_m\)&\(E_m\)&\(\Phi_m\)&\(K_m\)&
\(c_{m+1}\)&\(s_m\)&\(A_m\)&\(c_m\)\\
\midrule
\endhead
1  &141&718&618&234& \(0\)& \(7\)&007& \(0\)\\
2  &592&281&033&543& \(0\)& \(297\)&297& \(0\)\\
3  &653&828&988&140& \(-1\)& \(352\)&352& \(0\)\\
4  &589&459&749&729& \(-1\)& \(-431\)&569& \(-1\)\\
5  &793&045&894&659& \(-2\)& \(-717\)&283& \(-1\)\\
6  &238&235&848&824& \(-1\)& \(-1200\)&800& \(-2\)\\
7  &462&360&204&621& \(0\)& \(-3\)&997& \(-1\)\\
8  &643&287&586&058& \(-1\)& \(285\)&285& \(0\)\\
9  &383&471&834&914& \(-1\)& \(-895\)&105& \(-1\)\\
10 &279&352&365&794& \(0\)& \(-528\)&472& \(-1\)\\
11 &502&662&638&146& \(0\)& \(380\)&380& \(0\)\\
12 &884&497&117&601& \(0\)& \(663\)&663& \(0\)\\
\bottomrule
\end{longtable}
\endgroup

Le mot complet des treize retenues est
\[
 \boxed{
 C^{\rm fin}
 =(c_1,\ldots,c_{13})
 =(0,0,0,-1,-1,-2,-1,0,-1,-1,0,0,0).}
\]
La sortie dodécaphasique est
\[
 \boxed{
 A^{\rm fin}
 =(007,297,352,569,283,800,997,285,105,472,380,663).}
\]
Nous définissons le rationnel fini
\begin{equation}
 \boxed{
 \alpha_{12}^{\rm fin}
 :=10^{-36}\iota(A^{\rm fin})
 =\frac{7297352569283800997285105472380663}{10^{36}}.}
 \label{p12-83-c10-s2:eq:alpha-finita-autonoma}
\end{equation}
Par conséquent,
\[
 \alpha_{12}^{\rm fin}
 =0.007297352569283800997285105472380663.
\]

\subsection{Reconnaissance métrologique ultérieure : CODATA 2018}

La comparaison métrologique n’est effectuée qu’après clôture de la récurrence,
de la frontière droite et du rationnel précédent. L’identité entière donne
\begin{equation}
 \boxed{
 \alpha_{12}^{\rm fin}
 =10^{-36}
 \left\lfloor\frac{10^{36}}{137.035999084}\right\rfloor.}
 \label{p12-83-c10-s2:eq:alpha-finita-codata-2018}
\end{equation}
En conséquence,
\begin{equation}
 (\alpha_{12}^{\rm fin})^{-1}
 =137.03599908400000000000000000000001066588\ldots,
 \label{p12-83-c10-s2:eq:alpha-inversa-codata-2018}
\end{equation}
et reproduit exactement la suite centrale publiée par CODATA 2018
\cite{Tiesinga2021CODATA} :
\[
 \boxed{137.035999084.}
\]
Le mot \emph{exactement} renvoie à la suite centrale publiée, et non
à une prétendue mesure expérimentale à trente-six chiffres. Les chiffres
ultérieurs de \eqref{p12-83-c10-s2:eq:alpha-inversa-codata-2018} sont des conséquences du
rationnel HMT déjà fixé. CODATA ne sélectionne pas l’état, ne détermine pas (K),
ne décide pas du sens de la retenue et n’intervient pas dans le tableau dodécaphasique.

\subsection{Identité affine complète}

Chaque ligne du tableau satisfait exactement
\begin{equation}
 P_m+E_m-\Phi_m-K_m+c_{m+1}
 =A_m+Bc_m.
 \label{p12-83-c10-s2:eq:identidad-local-acarreo-alpha}
\end{equation}
En introduisant
\[
 C^-:=(c_1,\ldots,c_{12}),
 \qquad
 C^+:=(c_2,\ldots,c_{13}),
\]
les douze identités locales se réunissent en
\begin{equation}
 \boxed{
 P+E-\Phi-K+C^+-BC^-=A.}
 \label{p12-83-c10-s2:eq:identidad-afin-vectorial-alpha}
\end{equation}
Le terme
\[
 \partial_BC:=C^+-BC^-
\]
est le cocycle de retenue de la chaîne orientée. Par conséquent,
\(P+E-\Phi-K=A\) n’est pas une égalité coordonnée par coordonnée dans
\(\mathbb Z^{12}\) : elle ne devient correcte qu’en incorporant
\(\partial_BC\).

\begin{theorem}[Télescopage des treize retenues]
\label{p12-83-c10-s2:thm:telescopado-acarreos-alpha}
Pour toute condition terminale \(c_{13}\), la récurrence satisfait
\begin{equation}
 \boxed{
 \iota(P)+\iota(E)-\iota(\Phi)-\iota(K)-\iota(A)
 =B^{12}c_1-c_{13}.}
 \label{p12-83-c10-s2:eq:telescopado-general-alpha}
\end{equation}
Dans la fenêtre finie, \(c_1=c_{13}=0\), et donc
\begin{equation}
 \boxed{
 \iota(P)+\iota(E)-\iota(\Phi)-\iota(K)
 =\iota(A^{\rm fin}).}
 \label{p12-83-c10-s2:eq:identidad-entera-alpha-finita}
\end{equation}
\end{theorem}

\begin{proof}
Nous multiplions \eqref{p12-83-c10-s2:eq:identidad-local-acarreo-alpha} par
\(B^{12-m}\) et sommons sur \(m=1,\ldots,12\). La contribution des
retenues est
\[
 \sum_{m=1}^{12}
 (Bc_m-c_{m+1})B^{12-m}.
\]
Les composantes intérieures s’annulent :
\begin{align*}
 \sum_{m=1}^{12}(Bc_m-c_{m+1})B^{12-m}
 &=\sum_{m=1}^{12}c_mB^{13-m}
   -\sum_{m=1}^{12}c_{m+1}B^{12-m}\\
 &=B^{12}c_1-c_{13}.
\end{align*}
Cela prouve \eqref{p12-83-c10-s2:eq:telescopado-general-alpha}. La spécialisation
\(c_1=c_{13}=0\) donne \eqref{p12-83-c10-s2:eq:identidad-entera-alpha-finita}.
\end{proof}

\subsection{Inversion exacte du sceau \texorpdfstring{\(K\)}{K}}

La récurrence locale est réversible. En isolant \(K_m\) dans
\eqref{p12-83-c10-s2:eq:identidad-local-acarreo-alpha},
\begin{equation}
 \boxed{
 K_m=P_m+E_m-\Phi_m+c_{m+1}-A_m-Bc_m.}
 \label{p12-83-c10-s2:eq:inversion-local-k-desde-alpha}
\end{equation}
Par concaténation,
\begin{equation}
 \boxed{
 \iota(K)
 =\iota(P)+\iota(E)-\iota(\Phi)-\iota(A)
  -B^{12}c_1+c_{13}.}
 \label{p12-83-c10-s2:eq:inversion-concatenada-k-desde-alpha}
\end{equation}
À la frontière finie,
\[
 \iota(K)
 =\iota(P)+\iota(E)-\iota(\Phi)-\iota(A^{\rm fin}).
\]

La réversibilité algébrique n’inverse pas l’ordre causal. L’ordre effectif
est
\[
 x_{\rm term}
 \longrightarrow K
 \longrightarrow(P,E,\Phi,K,c_{13})
 \longrightarrow(A,C).
\]
La formule \eqref{p12-83-c10-s2:eq:inversion-local-k-desde-alpha} est un contrôle de
commutativité : elle démontre que la sortie et les retenues récupèrent le sceau
qui était déjà fixé. Elle ne constitue pas une règle permettant de sélectionner
rétrospectivement \(K\) à partir d’une valeur souhaitée de \(\alpha\).

\subsection{Prolongement coinductif du sceau dodécaphasique et récurrence inverse de la retenue dans \(\mathcal C=\{-2,-1,0,1,2\}\)}

Soient
\[
 (P_k)_{k\geq1},
 \qquad
 (E_k)_{k\geq1},
 \qquad
 (\Phi_k)_{k\geq1}
\]
les suites infinies de triades produites par les trois branches
coinductives, et soit
\[
 \kappa(k):=1+((k-1)\bmod12).
\]
Le prolongement du sceau est périodique :
\[
 K_k^{\rm per}:=K_{\kappa(k)}.
\]
Pour une profondeur finie \(N\) et une condition terminale
\(t=c_{N+1}\), la récurrence est
\begin{equation}
\begin{aligned}
 s_k^{(N,t)}
 &=P_k+E_k-\Phi_k-K_{\kappa(k)}+c_{k+1}^{(N,t)},\\
 A_k^{(N,t)}
 &=s_k^{(N,t)}\bmod1000,\\
 c_k^{(N,t)}
 &=\left\lfloor\frac{s_k^{(N,t)}}{1000}\right\rfloor,
 \qquad k=N,\ldots,1.
\end{aligned}
\label{p12-83-c10-s2:eq:recurrencia-larga-alpha}
\end{equation}

L’ensemble des retenues terminales
\[
 \mathcal C:=\{-2,-1,0,1,2\}
\]
est invariant pour les entrées certifiées : pour chaque triade et chaque
\(c_{k+1}\in\mathcal C\), le quotient \(c_k\) appartient à nouveau à
\(\mathcal C\). On ne choisit pas une condition terminale favorable. On propage
les cinq conditions et l’on ne publie que le préfixe commun à toutes.

L’exécution à profondeur certifiée emploie
\[
\begin{array}{c|r}
\text{quantité}&\text{valeur}\\ \hline
\text{blocs ternaires par canal}&3501\\
\text{trits par canal}&3501\cdot6=21006\\
\text{chiffres décimaux de garde}&10020\\
\text{triades de garde}&3340\\
\text{conditions terminales propagées}&5\\
\text{triades communes}&3339\\
\text{chiffres communs}&10017\\
\text{chiffres publiés}&10000
\end{array}
\]
Les cinq branches terminales coalescent avant le préfixe publié. En
particulier, aucune condition terminale extraite d’une valeur cible de
\(\alpha\) n’intervient dans la sortie.

\begin{theorem}[Compatibilité des sorties stabilisées]
\label{p12-83-c10-s2:thm:compatibilidad-prolongacion-alpha}
Soient \(A^{[N]}\) et \(A^{[N']}\), avec \(N'>N\), deux exécutions dont les
cinq conditions terminales ont coalescé avant la position \(N\).
Alors la troncature de \(A^{[N']}\) à ses \(N\) premiers blocs
coïncide avec \(A^{[N]}\). Les préfixes stabilisés définissent une unique
section dans la limite inverse des mots finis.
\end{theorem}

\begin{proof}
La récurrence est déterministe de droite à gauche une fois connue la
retenue entrant dans une position. La coalescence affirme que toutes les
conditions distantes produisent la même retenue avant la frontière du
préfixe. À partir de ce point, les mêmes triades \(P_k,E_k,\Phi_k\) et le même
sceau périodique produisent, par unicité euclidienne, les mêmes restes et
quotients à toutes les positions précédentes. C’est exactement la
commutativité avec la troncature.
\end{proof}

\subsection{Frontière provenant de la queue distante}

Le prolongement détermine à la frontière visible
\[
 c_{13}^{\infty}=-1.
\]
Les douze retenues intérieures restent identiques à celles de la fenêtre
finie :
\[
 \boxed{
 C_{\leq13}^{\infty}
 =(0,0,0,-1,-1,-2,-1,0,-1,-1,0,0,-1).}
\]
Seule change l’information qui entre depuis la treizième triade. Les
deux divisions pertinentes sont
\[
\begin{aligned}
 s_{12}^{\infty}
 &=884+497-117-601-1=662,\\
 A_{12}^{\infty}&=662,
 \qquad c_{12}^{\infty}=0,
\end{aligned}
\]
et
\[
\begin{aligned}
 s_{13}^{\infty}
 &=197+757-720-234-1=-1,\\
 A_{13}^{\infty}&=999,
 \qquad c_{13}^{\infty}=-1.
\end{aligned}
\]
Par conséquent, le début du mot prolongé est
\[
 \boxed{
 A^{\infty}
 =(007,297,352,569,283,800,997,285,105,472,380,662,999,\ldots).}
\]
Le développement correspondant commence par
\begin{equation}
 \boxed{
 \alpha_{\infty}
 =0.007297352569283800997285105472380662999564172539642\ldots.}
 \label{p12-83-c10-s2:eq:alpha-infinita-autonoma}
\end{equation}

La différence entre les deux frontières est locale et exacte :
\[
 c_{13}^{\rm fin}-c_{13}^{\infty}=1,
 \qquad
 A_{12}^{\rm fin}-A_{12}^{\infty}=1,
\]
tandis que
\[
 A_m^{\rm fin}=A_m^\infty,
 \qquad 1\leq m\leq11.
\]

En appliquant le théorème de télescopage aux douze premiers blocs du
prolongement,
\[
 \iota(P)+\iota(E)-\iota(\Phi)-\iota(K)
 -\iota(A^\infty_{\leq12})
 =-c_{13}^{\infty}=1.
\]
Donc
\begin{equation}
 \boxed{
 \iota(A^{\rm fin})
 =\iota(A^\infty_{\leq12})+1.}
 \label{p12-83-c10-s2:eq:relacion-unidad-fronteras-alpha}
\end{equation}
La terminaison \(663\) ne contredit pas la terminaison \(662\) : la première
est le représentant fini obtenu en fermant la fenêtre avec
\(c_{13}=0\) ; la seconde conserve la retenue \(-1\) provenant de la queue
distante.

\begin{figure}[htbp]
\centering
\begin{tikzpicture}[
  >=Stealth,
  node distance=12mm and 16mm,
  box/.style={draw,rounded corners=1.5mm,align=center,
    inner xsep=3.2mm,inner ysep=2.5mm},
  arr/.style={->,thick}
]
\node[box] (shared) {données partagées\\
 \(P,E,\Phi,K\)};
\node[box,below left=of shared] (finite)
 {frontière finie\\\(c_{13}=0\)};
\node[box,below right=of shared] (remote)
 {queue distante\\\(c_{13}=-1\)};
\node[box,below=of finite] (a663)
 {\(A_{12}=663\)\\clausura finita};
\node[box,below=of remote] (a662)
 {\(A_{12}=662\), \(A_{13}=999\)\\préfixe prolongé};
\draw[arr] (shared)--(finite);
\draw[arr] (shared)--(remote);
\draw[arr] (finite)--(a663);
\draw[arr] (remote)--(a662);
\draw[<->,densely dashed] (a663)--node[below,align=center]
 {différence entière exacte \(1\)} (a662);
\end{tikzpicture}
\caption{Bifurcation de frontière d’une même récurrence. On ne compare pas
deux générateurs distincts : seule change la donnée qui entre depuis la
droite.}
\label{p12-83-c10-s2:fig:bifurcacion-frontera-alpha}
\end{figure}

\subsection{Troncature et arrondi à \(36\) chiffres}

Pour \(x\geq0\), nous définissons
\[
 \operatorname{Tr}_{36}(x)
 :=10^{-36}\left\lfloor10^{36}x\right\rfloor
\]
et
\[
 \operatorname{Rd}_{36}(x)
 :=10^{-36}\left\lfloor10^{36}x+\frac12\right\rfloor.
\]
Soit
\[
 a:=7297352569283800997285105472380662.
\]
Le développement \eqref{p12-83-c10-s2:eq:alpha-infinita-autonoma} peut s’écrire
\[
 \alpha_\infty=\frac{a+\rho}{10^{36}},
 \qquad
 \rho=0.999564172539642709\ldots,
\]
d’où \(0.999<\rho<1\).

\begin{theorem}[Résolution exacte de la frontière \(662/663\)]
\label{p12-83-c10-s2:thm:frontera-alpha-662-663}
Le prolongement coinductif satisfait
\[
 \boxed{
 \operatorname{Tr}_{36}(\alpha_\infty)
 =0.007297352569283800997285105472380662,}
\]
tandis que son arrondi à trente-six chiffres satisfait
\[
 \boxed{
 \operatorname{Rd}_{36}(\alpha_\infty)
 =0.007297352569283800997285105472380663
 =\alpha_{12}^{\rm fin}.}
\]
\end{theorem}

\begin{proof}
Comme \(0\leq\rho<1\),
\[
 \left\lfloor10^{36}\alpha_\infty\right\rfloor=a.
\]
Cela donne la troncature se terminant par \(662\). Comme \(\rho>1/2\),
\[
 \left\lfloor10^{36}\alpha_\infty+\frac12\right\rfloor=a+1,
\]
qui se termine par \(663\). Par \eqref{p12-83-c10-s2:eq:alpha-finita-autonoma},
\((a+1)10^{-36}\) est exactement \(\alpha_{12}^{\rm fin}\).
\end{proof}

De plus,
\[
 0<\alpha_{12}^{\rm fin}-\alpha_\infty
 =\frac{1-\rho}{10^{36}}<10^{-39}.
\]
La différence n’est pas une erreur arithmétique du tableau : elle quantifie la
différence entre deux opérateurs de frontière distincts.

\subsection{Isolement causal à l’égard de \texorpdfstring{\(\alpha\)}{alpha} et des données métrologiques}

Pour la profondeur \(N\), le domaine complet de l’opération est
\[
 \mathscr D_N
 =\bigl((P_k)_{k=1}^{N},(E_k)_{k=1}^{N},
        (\Phi_k)_{k=1}^{N},K,c_{N+1}\bigr).
\]
L’application de retenue
\[
 \mathfrak C_N:\mathscr D_N
 \longrightarrow
 \{0,\ldots,999\}^{N}\times\mathcal C^{N}
\]
est définie exclusivement par \eqref{p12-83-c10-s2:eq:recurrencia-larga-alpha}. Son
domaine ne contient ni valeur cible de \(\alpha\), ni son inverse, ni tableau
métrologique, ni suite décimale cible, ni retenue terminale choisie par
comparaison.

\begin{theorem}[Non-interférence de la cible métrologique]
\label{p12-83-c10-s2:thm:no-interferencia-metrologica-alpha}
Soit \(\mathscr Y\) un espace quelconque de données externes contenant des valeurs
métrologiques ou des tableaux de comparaison. L’extension formelle
\[
 \widetilde{\mathfrak C}_N:
 \mathscr D_N\times\mathscr Y
 \longrightarrow
 \{0,\ldots,999\}^{N}\times\mathcal C^{N},
 \qquad
 \widetilde{\mathfrak C}_N(d,y):=\mathfrak C_N(d),
\]
se factorise par la projection
\[
 \operatorname{pr}_{\mathscr D_N}:
 \mathscr D_N\times\mathscr Y\longrightarrow\mathscr D_N.
\]
Par conséquent, pour tous \(y,y'\in\mathscr Y\),
\[
 \widetilde{\mathfrak C}_N(d,y)
 =\widetilde{\mathfrak C}_N(d,y').
\]
Aucune intervention sur une valeur de \(\alpha\) ou sur un tableau
métrologique n’altère \(K\), les quotients, les restes ni le préfixe produit.
\end{theorem}

\begin{proof}
L’affirmation résulte de
\[
 \widetilde{\mathfrak C}_N
 =\mathfrak C_N\circ\operatorname{pr}_{\mathscr D_N}.
\]
Toutes les opérations de \(\mathfrak C_N\) sont des additions, des soustractions et des divisions
euclidiennes des éléments de \(\mathscr D_N\). Aucune coordonnée de
\(\mathscr Y\) n’y apparaît.
\end{proof}

L’ordre causal complet est
\[
 \text{branches enrichies de }\pi,e,\varphi
 \longrightarrow(P,E,\Phi),
\]
\[
 x_{\rm term}\longrightarrow K,
\]
\[
 (P,E,\Phi,K,\text{frontière})
 \longrightarrow(A,C)
 \longrightarrow\alpha_\infty,
\]
\[
 \alpha_\infty
 \longrightarrow\text{comparaison externe}.
\]
La dernière flèche ne peut pas être inversée pour sélectionner l’une quelconque des
précédentes.

\subsection{Obstructions de la prolongation affine de \(\alpha\) : télescopage, inversion de \(K\), frontière \(662/663\) et isolement des données cibles}

\begin{criterion}[Réfutation de la clôture affine et de sa prolongation]
Le bloc est réfuté dans son domaine si l'un des faits
suivants se produit :
\begin{enumerate}
 \item une ligne de la \cref{p12-83-c10-s2:tab:acarreo-alpha-finito-autonomo}
       ne satisfait pas \(s_m=A_m+1000c_m\) ;
 \item le vecteur publié des treize retenues ne reproduit pas les douze lignes ;
 \item l’identité télescopique échoue
       \[
       \iota(P)+\iota(E)-\iota(\Phi)-\iota(K)-\iota(A)
       =1000^{12}c_1-c_{13};
       \]
 \item l’inversion \eqref{p12-83-c10-s2:eq:inversion-local-k-desde-alpha} ne récupère pas
       les douze coordonnées de \(K\) ;
 \item la propagation distante produit \(c_{13}\neq-1\) pour le préfixe
       certifié ;
 \item les cinq conditions terminales de
       \(\mathcal C=\{-2,-1,0,1,2\}\) ne coalescent pas avant les chiffres
       publiés ;
 \item la triade suivant \(662\) n’est pas \(999\), auquel cas l’identification
       de l’arrondi à \(663\) cesse d’être démontrée ;
 \item le générateur ouvre un développement cible de \(\alpha\), son inverse,
       un tableau métrologique ou une condition terminale choisie pour imposer
       la coïncidence.
\end{enumerate}
\end{criterion}

Deux identités de frontière distinctes sont ainsi démontrées :
\[
 \boxed{
 663=\text{clôture finie avec }c_{13}=0
     =\text{arrondi à 36 chiffres de }\alpha_\infty,}
\]
\[
 \boxed{
 662=\text{troncature à 36 chiffres de }\alpha_\infty
 \quad\text{avec}\quad c_{13}=-1,\ A_{13}=999.}
\]
Le résultat métrologique central est finalement publié dans la même
résidence que la dérivation :
\begin{equation}
 \boxed{
 \begin{gathered}
  (\alpha_{12}^{\rm fin})^{-1}
  =137.03599908400000000000000000000001066588\ldots,\\
  \text{chaîne centrale CODATA 2018 :}\qquad137.035999084.
 \end{gathered}}
 \label{p12-83-c10-s2:eq:resultado-final-alpha-codata-2018}
\end{equation}
L’implication exprime une comparaison ultérieure avec la suite centrale
publiée ; elle n’introduit pas cette suite dans le constructeur.
 % Insercion editor-ready para el capitulo 31.

\section{Prolongement nonadique gradué du noyau \texorpdfstring{\(E_{21/22}\)}{E21/22} : opérateur \texorpdfstring{\(T_{7:9}\)}{T7:9}, généalogie \texorpdfstring{\(120/45/11/495\)}{120/45/11/495} et résidu d’ordre \texorpdfstring{\(9\)}{9}}
\label{sec:l9s102:alpha-dos-vias}

Cette section développe deux constructions internes relatives à la coordonnée \(\alpha\) de la
\hyperref[sec:tesis-inaugural-helicoidal-hmt-md]{tesis inaugural}. La voie A opère par la clôture dodécaphasique réversible de l’état signé et exprime \(\alpha_{12}\). La voie B, une fois le noyau analytique \(E_{21/22}\) et ses coefficients construits en interne, exprime dans sa carte analytique une racine positive simple et unique, notée \(\alpha_E:=x_{21/22}\). Les deux sorties partagent une généalogie ; leur concordance démontrée appartient au codomaine décimal typé précisé ci-dessous. \(\alpha_E\) n’est pas identifié par définition à \(\alpha_{12}\). Leurs domaines, opérations et certificats restent distincts :
\begin{equation}
 X_{12}^{\mathrm{sig}}
 \xrightarrow{\ \Gamma_{12}^{\alpha}\ }
 \alpha_{12},
 \qquad
 E_{21/22}:I_\alpha\longrightarrow\mathbb R,
 \quad E_{21/22}(x_{21/22})=0,
 \quad E_{21/22}'(x_{21/22})\ne0.
 \label{eq:l9s102:alpha-dos-mapas}
\end{equation}
La clôture entière exprime \(\alpha_{12}\) par retenue, tandis que l’équation analytique détermine \(\alpha_E=x_{21/22}\) par existence et unicité de la racine dans l’intervalle certifié. La concordance actuellement démontrée des deux sorties réside exactement dans le codomaine du morphisme
\begin{equation}
 \operatorname{Tr}_9(y):=10^{-9}\lfloor10^9y\rfloor,
 \qquad
 \operatorname{Tr}_9(x_{21/22}^{-1})
 =137.035999084
 =\operatorname{Tr}_9(\alpha_{12}^{-1}).
 \label{eq:l9s102:concordancia-tipica-alpha}
\end{equation}
L’identité analytique exacte correspond à \(E_{21/22}(\alpha_E)=0\), tandis que \eqref{eq:l9s102:concordancia-tipica-alpha} est l’égalité exacte des deux expressions après application de \(\operatorname{Tr}_9\). La promotion plus forte \(\alpha_E=\alpha_{12}\) dans \(\mathbb R_{\rm HMT}\) exige de démontrer \(E_{21/22}(\alpha_{12})=0\), ou un morphisme injectif relevant l’égalité tronquée ; le résidu non nul de la \cref{tab:l9s102:alpha-residuos} ne fournit pas cette identité. Aucune voie n’introduit la valeur cible comme donnée d’entrée ; l’identification métrologique relève uniquement de la reconnaissance conventionnelle ultérieure.

\subsection{Noyau analytique \texorpdfstring{\(P_6\)}{P6} et racine unique de l’équation \texorpdfstring{\(E_{21/22}\)}{E21/22}}
\label{sec:l9s102:alpha-p6}

La caractérisation des lacunes emploie
\begin{equation}
 \begin{aligned}
 P_6(x)={}&2\pi x-\frac74x^2
 +\frac{x^3}{2\pi}
 +\frac{x^4}{20}-\frac{2x^5}{21}-\frac{x^6}{46},\\
 P_6(x)={}&\log_{10}\frac{10}{9}.
 \end{aligned}
 \label{eq:l9s102:p6}
\end{equation}
La racine positive certifiée de cette équation satisfait
\begin{equation}
 x_6^{-1}=137.03599908400002702009\ldots .
 \label{eq:l9s102:raíz-p6}
\end{equation}
Le polynôme \(P_6\) fournit le jet initial de la caractérisation analytique ; l’holonomie nonadique prolonge maintenant ses degrés \(7,8,9\).

\subsection{Généalogie APP--TRIT--TPK des invariants \texorpdfstring{\(120,45,11\)}{120, 45, 11} et \texorpdfstring{\(495\)}{495}}
\label{sec:l9s102:alpha-genealogia-7-9}

La demi-couronne de phase du TPK construit
\begin{equation}
 |\mathcal F_8^{\mathrm{ph}}|=120,
 \qquad |L_{\mathrm{ph}}|=45.
 \label{eq:l9s102:alpha-120-45}
\end{equation}
APP apporte le bloc central
\begin{equation}
 B_c=\begin{pmatrix}7&2\\2&7\end{pmatrix},
 \qquad
 \det B_c=45,
 \qquad
 \operatorname{Spec}(B_c)=\{9,5\}.
 \label{eq:l9s102:alpha-bc}
\end{equation}
L'extracteur transversal dodécaphasique fournit
\begin{equation}
 r_{\mathrm{dod}}=\operatorname{rank}(D_3,D_4)=11,
 \qquad
 (D_3K)_1=-495,
 \qquad
 495=45\cdot11=\det(B_c)\,r_{\mathrm{dod}}.
 \label{eq:l9s102:alpha-495}
\end{equation}
Les coefficients du jet prolongé proviennent donc de trois
invariants préalables d'APP--TRIT--TPK :
\begin{equation}
 -\frac1{120}\longmapsto\frac1{45}
 \longmapsto\frac2{45\cdot11}.
 \label{eq:l9s102:alpha-coeficientes}
\end{equation}

\subsection{Résolvante nilpotente \texorpdfstring{\((I-N)^{-1}\)}{(I-N)^-1} et construction exacte de la queue \texorpdfstring{\(T_{7:9}\)}{T7:9}}
\label{sec:l9s102:alpha-resolvente}

Sur la base \(e_7,e_8,e_9\), nous définissons
\begin{equation}
 Ne_7=-\frac83e_8,
 \qquad
 Ne_8=\frac2{11}e_9,
 \qquad
 Ne_9=0,
 \qquad
 v_7=-\frac1{120}e_7.
 \label{eq:l9s102:alpha-nilpotente}
\end{equation}
Comme \(N^3=0\), sa résolvante finie agit exactement par
\begin{equation}
 (I-N)^{-1}v_7
 =v_7+Nv_7+N^2v_7
 =-\frac{e_7}{120}+\frac{e_8}{45}+\frac{2e_9}{495}.
 \label{eq:l9s102:alpha-resolvente-finito}
\end{equation}
L'évaluation \(e_n\mapsto x^n\) produit
\begin{equation}
 \boxed{
 T_{7:9}(x)=-\frac{x^7}{120}+\frac{x^8}{45}
 +\frac{2x^9}{495}.}
 \label{eq:l9s102:alpha-t79}
\end{equation}

\subsection{Jet nonadique \texorpdfstring{\(P_9\)}{P9}, racine prolongée et échelle de résidus certifiés}
\label{sec:l9s102:alpha-jet-p9}

Le jet nonadique d'ordre \(9\) est défini par
\begin{equation}
 P_9(x)=P_6(x)+T_{7:9}(x),
 \qquad
 P_9(x)=\log_{10}\frac{10}{9}.
 \label{eq:l9s102:alpha-p9}
\end{equation}
Sa racine positive satisfait
\begin{equation}
 x_9^{-1}
 =137.0359990840000000000000111926291900\ldots .
 \label{eq:l9s102:alpha-raíz-p9}
\end{equation}

\begin{table}[htbp]
\centering
\caption{Décroissance certifiée du résidu de la caractérisation analytique}
\label{tab:l9s102:alpha-residuos}
\begin{tabular}{@{}lc@{}}
\toprule
Noyau évalué en \(\alpha_{12}\) & Résidu \\
\midrule
\(P_6\) & \(9.0038886\times10^{-18}\) \\
\(P_6-x^7/120\) & \(-1.7893115\times10^{-19}\) \\
\(P_6-x^7/120+x^8/45\) & \(-2.3708601\times10^{-22}\) \\
\(P_9\) & \(3.7297132\times10^{-27}\) \\
\bottomrule
\end{tabular}
\end{table}

\subsection{Opérateur dodécaphasique alternatif \texorpdfstring{\(A_{\mathrm{dod}}\)}{A_dod} et démonstration de la non-unicité de la sélection}
\label{sec:l9s102:alpha-control}

L'opérateur de contrôle
\begin{equation}
 A_{\mathrm{dod}}
 =I-\frac14(D_3^*D_3+D_4^*D_4)
 \label{eq:l9s102:alpha-control-alternativo}
\end{equation}
produit une suite différente de coefficients. Ce contrôle délimite la force sélective de la résolvante \cref{eq:l9s102:alpha-resolvente-finito,eq:l9s102:alpha-t79} : un opérateur dodécaphasique générique ne détermine pas à lui seul le prolongement \(T_{7:9}\). Il ne diminue pas la portée de la voie des lacunes, dont la condition probatoire est l’existence et l’unicité de la racine du noyau déclaré. Le morphisme de contrôle \(\Gamma_{E21}^{(9)}\) est réservé exclusivement à l’unicité canonique du prolongement nonadique \(7{:}9\) ; il ne certifie ni ne conditionne le statut de la voie B.

\subsection{Portée de l’opérateur nilpotent \texorpdfstring{\(N\)}{N} sur le jet nonadique \texorpdfstring{\(P_9\)}{P9} et contrôle d’unicité du prolongement \texorpdfstring{\(7{:}9\)}{7:9}}
\label{sec:l9s102:alpha-alcance}

L’arithmétique de \(120,45,11,495\), l’opérateur nilpotent, la résolvante et l’identité \cref{eq:l9s102:alpha-t79} possèdent une exactitude interne. La racine positive simple et unique du noyau \(E_{21/22}\) constitue une seconde construction interne exacte, \(\alpha_E\), pour le domaine et les coefficients déclarés. Son identification à la sortie dodécaphasique est démontrée exactement après \(\operatorname{Tr}_9\) ; son élévation à une égalité dans \(\mathbb R_{\rm HMT}\) exige l’identité ou le relèvement injectif énoncés au début de la section. Indépendamment, l’unicité canonique du prolongement \(7{:}9\) est typée par le morphisme
\begin{equation}
 \Gamma_{E21}^{(9)}:
 X_{\mathrm{APP\text{-}TRIT\text{-}TPK}}
 \longrightarrow \mathbb Q(\pi)[x]/(x^{10}),
 \label{eq:l9s102:alpha-gamma-e21}
\end{equation}
dont l’action détermine conjointement le germe orienté, les deux transferts, la graduation \(7,8,9\) et le quotient de jets \(F^7/F^{10}\). Le réfutateur distingue donc deux énoncés : il invaliderait la voie B en cas d’échec de la monotonie, du changement de signe ou de l’isolement de la racine ; il invaliderait la promotion canonique de \(T_{7:9}\) si un autre lecteur du même domaine produisait la même graduation sans se factoriser par \(\Gamma_{E21}^{(9)}\).
 
% Selección íntegra de los cuerpos matemáticos de la Parte III rectora. Las
% cartas dimensionales se reservan para la parte siguiente, de modo que una
% constante física no pueda retroactuar como generador de estos caracteres.
% Punto de montaje redefinible por el compilador único.
% La ruta por defecto es válida cuando el módulo se consume desde
% ``manuscrito/``; el editor único puede redefinirla antes del \input.
\providecommand{\HMTParteIIIRoot}{../colaboracion/parte_iii}
\providecommand{\HMTParteIIISource}{\HMTParteIIIRoot/source/public_final}
\providecommand{\HMTParteIIICapitulos}{\HMTParteIIIRoot/tex/capitulos_finales}
\providecommand{\APP}{\mathrm{APP}}
\providecommand{\TRIT}{\mathrm{TRIT}}
\providecommand{\TPK}{\mathrm{TPK}}
\providecommand{\etaRet}{\eta_{\mathrm{ret}}^{(\mathrm{deg})}}
\providecommand{\etaRetrad}{\eta_{\mathrm{ret}}^{(\mathrm{rad})}}
\providecommand{\alphaAngle}{\alpha_{\angle}}
 
% \newcommand{\HMTMonographChapterInput}[2]{%
%   \chapter{#1}%
%   \begingroup
%   \renewcommand{\chapter}[2][]{}%
%   \input{#2}%
%   \endgroup}

\chapter{Classification algébrique et généalogique des constantes fondamentales}
\begingroup
\renewcommand{\chapter}[2][]{}%
\chapter{Algèbre et classification généalogique des constantes fondamentales}
\label{chap:p3-final:27:algebra_clasificacion}
\label{chap:p3-83-23-algebra_clasificacion}


\label{int:chap:producto-nativo}
\index{producto nativo}

L'extension euclidienne résidu--quotient d'APP contient deux applications
locales : l'application additive normalise les collisions et la multiplicative
produit des produits partiels. Dans cette section, on construit la multiplication
globale des mots en utilisant exclusivement ces deux applications. La correspondance avec la multiplication entière sera une
conséquence de l'entrelaceur, et non la définition de l'opération.

\section{Coordonnées euclidiennes de la cellule locale}

\begin{definition}[Applications additive et multiplicative sur un même support]
Soit $\Dnine=\{1,\ldots,9\}$. Pour $a,b\in\Dnine$, on définit par division
euclidienne avec reste dans $\Dnine$ :
\begin{align}
 a+b&=9q^+(a,b)+R^+(a,b),\label{int:eq:local-plus}\\
 ab&=9q^\times(a,b)+R^\times(a,b).\label{int:eq:local-times}
\end{align}
En particulier, le reste nul est représenté par $9$ et le quotient est réduit
d'une unité. Le quadruplet
\[
 \CellRQ(a,b)=
 (R^+(a,b),q^+(a,b);R^\times(a,b),q^\times(a,b))
\]
est l'extension arithmétique locale par résidu et quotient.
\end{definition}

L'agrégation selon les trois classes modulo trois du tableau visible
\(R^\times\) définit une matrice distincte des applications locales :
\[
(M_\Pi)_{rs}
=\frac19
\sum_{\substack{a\equiv r\;(3)\\b\equiv s\;(3)}}R^\times(a,b),
\qquad r,s\in\{1,2,0\},
\]
et, dans cet ordre des classes,
\[
M_\Pi=
\begin{pmatrix}4&5&6\\5&4&6\\6&6&9\end{pmatrix}.
\]
Ainsi, \(q^+\) et \(q^\times\) sont des coordonnées de quotient d'une cellule,
tandis que \(M_\Pi\) est un opérateur agrégé sur les classes résiduelles.

Les $81$ cellules satisfont exactement
\[
\begin{array}{c|rrrr}
 &\sum R&\sum q&\sum(9q+R)\\\hline
 +&405&45&810\\
 \times&459&174&2025
\end{array}
\]
et donc
\[
 2025-810=(459-405)+9(174-45)=54+9\cdot129=1215.
\]
Le quotient $q^\times$ est nécessaire à la reconstruction. Par exemple,
\[
 2\cdot5=9\cdot1+1.
\]
Si l'on ne conserve que $R^\times$, le produit $10$ se confond avec la valeur
visible $1$.

\subsection{Cocycles d'associativité et de distributivité}

La cohérence du transducteur peut déjà être décidée dans la cellule. Pour tout
$a,b,c\in\Dnine$, les trois identités sont vérifiées
\begin{align}
q^+(a,b)+q^+(R^+(a,b),c)
&=q^+(b,c)+q^+(a,R^+(b,c)),\label{int:eq:cocycle-plus}\\
q^\times(R^\times(a,b),c)+cq^\times(a,b)
&=q^\times(a,R^\times(b,c))+aq^\times(b,c),
\label{int:eq:cocycle-times}\\
q^\times(a,R^+(b,c))+aq^+(b,c)
&=q^+(R^\times(a,b),R^\times(a,c))\notag\\
&\quad+q^\times(a,b)+q^\times(a,c).
\label{int:eq:cocycle-distributive}
\end{align}

\begin{theorem}[Cohérence locale de l'extension euclidienne]
Les identités \eqref{int:eq:cocycle-plus}--\eqref{int:eq:cocycle-distributive},
avec les égalités correspondantes de résidus, expriment
respectivement l'associativité additive, l'associativité multiplicative et la
distributivité dans l'extension résidu--quotient. Les trois familles présentent
zéro écart sur les $9^3=729$ triplets.
\end{theorem}

\begin{proof}
En développant $(a+b)+c$ au moyen de \eqref{int:eq:local-plus}, le coefficient de
neuf est le membre de gauche de \eqref{int:eq:cocycle-plus} ; en développant
$a+(b+c)$, on obtient celui de droite. L'unicité de la division avec reste dans
$\Dnine$ égalise aussi bien les résidus que les quotients. Le même argument appliqué à
$(ab)c=a(bc)$ donne \eqref{int:eq:cocycle-times}. Enfin, les deux développements de
$a(b+c)=ab+ac$ produisent \eqref{int:eq:cocycle-distributive}. L'énumération
directe des \(729\) cas de chaque identité confirme le calcul
algébrique.
\end{proof}

\section{Bijectivité de la numération}

Le symbole $9$ n'est pas identifié à zéro. Dans $\CellRQ$, $9$ représente un
franchissement complet de la frontière nonadique. Le remplacement par les chiffres
$0,\ldots,8$ éliminerait le quotient euclidien qui intervient dans
l'entrelacement.

\begin{definition}[Mots canoniques]
Soit
\[
 \Wnine^+=\bigsqcup_{\ell\geq1}\Dnine^\ell,
 \qquad \Wnine=\{\varnothing\}\sqcup\Wnine^+.
\]
Les mots s'écrivent à partir du chiffre le moins significatif. Nous définissons
\[
 \operatorname{ev}_9(u_0,\ldots,u_{\ell-1})
 =\sum_{j=0}^{\ell-1}u_j9^j,
 \qquad \operatorname{ev}_9(\varnothing)=0.
\]
\end{definition}

\begin{theorem}[Unicité de la forme normale]
L'évaluation $\operatorname{ev}_9$ est une bijection de $\Wnine^+$ sur
$\N_{\geq1}$, et de $\Wnine$ sur $\N_0$.
\end{theorem}

\begin{proof}
Les mots de longueur $\ell$ couvrent l'intervalle
\[
 \frac{9^\ell-1}{8}
 \leq \operatorname{ev}_9(u)
 \leq 9\frac{9^\ell-1}{8}.
\]
Le minimum de l'intervalle suivant est supérieur d'une unité au maximum du
précédent. À l'intérieur de chaque intervalle, la division de $n-1$ par neuf détermine
de manière unique le chiffre
\[
 d=(n-1)\bmod9+1
\]
et le quotient qui est codé récursivement. Le processus se termine parce que le
quotient décroît strictement.
\end{proof}

Exemples :
\[
 9\longleftrightarrow(9),\qquad
 10\longleftrightarrow(1,1),\qquad
 81\longleftrightarrow(9,8),\qquad
 1000\longleftrightarrow(1,3,3,1).
\]

\section{Fibres des projections résiduelles dans les relèvements entiers}

Pour ce bloc, nous distinguons deux alphabets. La forme normale
\(\Wnine\) utilise les chiffres \(\{1,\ldots,9\}\), écrits à partir du moins
significatif. En revanche, \(\mathscr B_9\) désigne les écritures positionnelles
conventionnelles en base neuf, avec les chiffres \(\{0,\ldots,8\}\) et le plus
significatif à gauche. Les deux types ne sont pas identifiés.

Soit
\[
 \mathscr W_{10}^{(2)}=\{0,\ldots,9\}^2,\qquad
 \rho_{10}(a,b)=(b,a),\qquad
 \operatorname{ev}_{10}(a,b)=10a+b,
\]
et soit \(q_9:\mathbb N\to\mathbb Z/9\mathbb Z\) la réduction résiduelle.

\begin{lemma}[Descente résiduelle du renversement décimal]
\label{int:pf84:SYN30-RHO-DESC}
Le renversement sur les mots décimaux de longueur deux est involutif et son
ombre modulo neuf est l'identité :
\[
 \rho_{10}^2=I,\qquad
 q_9\!\left(\operatorname{ev}_{10}(\rho_{10}(a,b))\right)
 =
 q_9\!\left(\operatorname{ev}_{10}(a,b)\right).
\]
La longueur fixe conserve les zéros initiaux.
\end{lemma}

\begin{proof}
La double permutation retrouve \((a,b)\). Comme \(10\equiv1\pmod9\),
\[
 10b+a\equiv a+b\equiv10a+b\pmod9.
\]
Un mot dont le double renversement ne serait pas lui-même ou dont la somme des chiffres
changerait modulo neuf réfuterait le lemme.
\end{proof}

\begin{lemma}[Descente résiduelle du carré]
\label{int:pf84:SYN30-SQ-DESC}
L'application \(s(n)=n^2\) induit l'endoapplication résiduelle
\[
 [n]_9\longmapsto[n]_9^2,
 \qquad q_9(s(n))=q_9(n)^2.
\]
\end{lemma}

\begin{proof}
La compatibilité d'une congruence avec le produit donne
\(n\equiv m\pmod9\Rightarrow n^2\equiv m^2\pmod9\). Un entier pour lequel
\(q_9(n^2)\ne q_9(n)^2\) réfuterait l'énoncé.
\end{proof}

\begin{proposition}[Fibres résiduelles du renversement entier]
\label{int:pf84:SYN30-RHO-NONFAITH}
L'ombre résiduelle ne reconstruit pas le relèvement entier
\(\widehat\rho_{10}=\operatorname{ev}_{10}\circ\rho_{10}\). En effet,
\[
 18\equiv27\pmod9,\qquad
 \widehat\rho_{10}(1,8)=81\ne72=\widehat\rho_{10}(2,7).
\]
\end{proposition}

\begin{proof}
Les deux entiers appartiennent à la classe résiduelle zéro, mais leurs renversements sont
distincts. Il n'existe donc pas de
\(\overline\rho:\mathbb Z/9\mathbb Z\to\mathbb N\) qui factorise
simultanément ces deux valeurs.
\end{proof}

\begin{proposition}[Fibres résiduelles du carré entier]
\label{int:pf84:SYN30-SQ-NONFAITH}
La même ombre ne reconstruit pas non plus le carré entier :
\[
 18\equiv27\pmod9,\qquad
 18^2=324\ne729=27^2.
\]
\end{proposition}

\begin{proof}
Une fonction \(\overline s:\mathbb Z/9\mathbb Z\to\mathbb N\) qui
factoriserait \(s\) devrait attribuer à la même classe zéro les deux valeurs
\(324\) et \(729\), ce qui est impossible.
\end{proof}

\begin{corollary}[Témoin croisé \(27\)--\(72\)--\(729\)]
\label{int:pf84:SYN30-CROSS-WITNESS}
\[
 \widehat\rho_{10}(2,7)=72,\qquad
 27^2=729=9^3=3^6,
\]
et, dans l'alphabet conventionnel \(\mathscr B_9\),
\[
 27=30_9,\qquad72=80_9,\qquad729=1000_9.
\]
Les trois entiers ont une ombre résiduelle zéro et une racine numérique positive neuf.
\end{corollary}

\begin{proof}
Les deux opérations se calculent directement ; les divisions successives par
neuf produisent les trois écritures. Le résultat échouerait si l'on confondait
\(\mathscr B_9\) avec \(\Wnine\), si l'une des écritures était incorrecte
ou si l'on affirmait que la descente résiduelle fidèle reconstruit l'opération et sa
généalogie complète.
\end{proof}

\section{Additionneur natif et propagation du quotient}

\begin{algorithm}[Somme de mots]
On parcourt deux mots par positions, avec une retenue initiale \(c_0=0\). Si, à
une position, il y a deux chiffres, on applique \eqref{int:eq:local-plus} ; s'il n'y en a qu'un,
ce chiffre est pris comme résidu provisoire et le quotient provisoire est zéro.
Lorsque \(c_j\ne0\), on applique une seconde cellule additive au résidu provisoire
et à \(c_j\). On publie le nouveau résidu et l'on transporte la somme des
quotients à la position suivante. S'il reste à la fin une retenue non nulle, elle est
publiée comme dernier chiffre. De cette manière, les cellules locales ne reçoivent jamais le
symbole absent \(0\), qui n'appartient pas à \(\Dnine\).
\end{algorithm}

\begin{lemma}[Invariant de l'additionneur]
Si $u\boxplus_\#v$ est le mot produit par l'algorithme précédent,
\[
 \operatorname{ev}_9(u\boxplus_\#v)
 =\operatorname{ev}_9(u)+\operatorname{ev}_9(v).
\]
La retenue intermédiaire est dans $\{0,1,2\}$.
\end{lemma}

\begin{proof}
Après traitement des positions $0,\ldots,k-1$, soit $c_k$ le quotient en attente.
Les identités locales impliquent
\[
 \sum_{j<k}(u_j+v_j)9^j
 =\sum_{j<k}w_j9^j+c_k9^k,
\]
en incluant comme zéro les chiffres absents. Chaque nouvelle cellule conserve cette
égalité. Le maximum local est $9+9+2=20$, de sorte que $c_{k+1}\leq2$. À la
fin, le dernier quotient est publié, et l'identité globale demeure.
\end{proof}

\section{Multiplication par un chiffre}

\begin{algorithm}[Transducteur scalaire]
$d\in\Dnine$ étant fixé, chaque chiffre $u_j$ produit
\[
 u_jd=9q_j+r_j,
 \qquad
 (r_j,q_j)=\bigl(R^\times(u_j,d),q^\times(u_j,d)\bigr).
\]
Soit $c_j\in\{0,\ldots,9\}$ le quotient entrant. Comme $0\notin\Dnine$,
la transition est définie par branches :
\[
 (w_j,c_{j+1})=
 \begin{cases}
 (r_j,q_j),&c_j=0,\\[1mm]
 \bigl(R^+(r_j,c_j),q_j+q^+(r_j,c_j)\bigr),&c_j\in\Dnine.
 \end{cases}
\]
Ainsi, aucune application locale n'est évaluée hors de son domaine. La première
branche exprime l'absence de quotient en attente ; la seconde normalise la
collision au moyen de l'application additive.
\end{algorithm}

\begin{lemma}[Invariant scalaire]
Le mot $d\odot_\#u$ produit par le transducteur satisfait
\[
 \operatorname{ev}_9(d\odot_\#u)
 =d\operatorname{ev}_9(u).
\]
Le quotient multiplicatif appartient à $\{0,\ldots,9\}$.
\end{lemma}

\begin{proof}
L'argument est celui de l'additionneur, en remplaçant chaque terme local par
$u_jd$. Si $c_j=0$, l'identité
$u_jd=9q_j+r_j$ conserve directement l'égalité partielle. Si
$c_j\in\Dnine$, la seconde division euclidienne donne
\[
 r_j+c_j=9q^+(r_j,c_j)+R^+(r_j,c_j),
\]
et conserve la même égalité après transport de
$q_j+q^+(r_j,c_j)$. Le maximum $9\cdot9+9=90$ et la forme du reste dans
$\Dnine$ maintiennent le quotient dans la borne déclarée. Le dernier quotient est
publié comme chiffre.
\end{proof}

\section{Produit global par Horner}

\begin{definition}[Produit natif]
Soit $v=(v_0,\ldots,v_{m-1})$. En partant de $a_m=\varnothing$, on définit pour
$j=m-1,\ldots,0$
\[
 a_j=(9\odot_\#a_{j+1})\boxplus_\#(v_j\odot_\#u).
\]
Le résultat est
\[
 u\boxtimes_\#v:=a_0.
\]
\end{definition}

Cette définition ne contient pas
\[
 \operatorname{dec}_9(\operatorname{ev}_9(u)\operatorname{ev}_9(v)).
\]
L'évaluation n'apparaîtra que dans la démonstration de l'entrelacement.

\begin{theorem}[Multiplicativité native]
Pour tous les mots $u,v\in\Wnine$,
\begin{equation}
 \label{int:eq:ev-multiplicative}
 \boxed{
 \operatorname{ev}_9(u\boxtimes_\#v)
 =\operatorname{ev}_9(u)\operatorname{ev}_9(v).}
\end{equation}
Le mot vide est absorbant et $(1)$ est l'unité.
\end{theorem}

\begin{proof}
Par les deux lemmes précédents,
\[
 \operatorname{ev}_9(a_j)
 =9\operatorname{ev}_9(a_{j+1})
  +v_j\operatorname{ev}_9(u).
\]
L'itération de Horner donne
\[
 \operatorname{ev}_9(a_0)
 =\operatorname{ev}_9(u)
  \sum_{j=0}^{m-1}v_j9^j.
\]
Le dernier facteur est $\operatorname{ev}_9(v)$. L'unicité de la forme normale
identifie le mot résultant.
\end{proof}

\begin{corollary}[Monoïde multiplicatif de mots]
La structure
\[
\WordMonoid=(\Wnine,\boxtimes_\#,(1))
\]
est un monoïde commutatif ; le mot vide est son élément absorbant. Sa
construction n'utilise pas la multiplication globale des entiers.
\end{corollary}

L'exemple de frontière le plus court est
\[
 (9)\boxtimes_\#(9)=(9,8),
 \qquad 9+8\cdot9=81.
\]
Les \(17\) cellules locales de la multiplication produisent
\[
 (9,8)\boxtimes_\#(9,8)=(9,8,8,8),
 \qquad \operatorname{ev}_9(9,8,8,8)=6561.
\]

\section{Linéarisation matricielle de l'état généalogique et construction de l'entrelaceur complètement positif}

\begin{definition}[Secteurs de Hilbert]
Soient
\[
 \Krad=\ell^2(\Wnine^+),
 \qquad
 \Hplus=\ell^2(\N_{\geq1}).
\]
Nous définissons sur les bases canoniques
\[
 W_\times e_u=e_{\operatorname{ev}_9(u)}.
\]
Dans \(\Hplus\), l'opérateur nombre est noté
\[
 \mathsf N e_n=n e_n,
 \qquad \Spec(\mathsf N)=\N_{\geq1}.
\]
Le vide peut être ajouté séparément comme vecteur de vide $e_0$ ; il ne se
confond pas avec $\Hplus$, dont le spectre commence à un.
\end{definition}

\begin{theorem}[Unitarité et carré multiplicatif]
$W_\times$ est unitaire. Dans le cœur algébrique
$\mathbb C[\Wnine^+\times\Wnine^+]$, soit
\[
 \mu_\#(e_u\otimes e_v)=e_{u\boxtimes_\#v},
 \qquad
 \mu_{\N}(e_a\otimes e_b)=e_{ab}.
\]
Alors on a le carré
\[
\begin{CD}
\Krad\otimes_{\mathrm{alg}}\Krad @>{W_\times\otimes W_\times}>>
\Hplus\otimes_{\mathrm{alg}}\Hplus\\
@V{\mu_\#}VV @VV{\mu_{\N}}V\\
\Krad @>>{W_\times}> \Hplus
\end{CD}
\]
et, par conséquent, le carré commute.
\end{theorem}

\begin{proof}
La bijection des formes normales envoie une base orthonormale sur une autre et prouve
l'unitarité. Sur un tenseur de base, les deux chemins produisent
$e_{\operatorname{ev}_9(u)\operatorname{ev}_9(v)}$ par
\eqref{int:eq:ev-multiplicative}. La linéarité donne l'égalité dans le cœur.
\end{proof}

Les applications $\mu_\#$ et $\mu_{\N}$, définies initialement sur les cœurs
algébriques respectifs, sont densément définies et non bornées. Pour
\(c=(c_{u,v})\in\ell^2(\Wnine^+\times\Wnine^+)\), chaque somme sur une fibre de produit
est finie. L'adjoint de \(\mu_\#\) contient tous les vecteurs à support
fini dans son domaine, de sorte que \(\mu_\#\) est fermable. Le domaine de sa
fermeture est
\[
 \Dom(\overline{\mu_\#})=
 \left\{c\in\ell^2(\Wnine^+\times\Wnine^+):
 \sum_{w\in\Wnine^+}
 \left|\sum_{u\boxtimes_\#v=w}c_{u,v}\right|^2<\infty
 \right\}.
\]
Chaque fibre est finie parce que $\operatorname{ev}_9$ l'identifie aux couples
ordonnés de diviseurs d'un entier. La même description vaut pour
\(\overline{\mu_{\N}}\) après application de \(W_\times\otimes W_\times\), et
$W_\times$ entrelace les graphes des deux fermetures.

\section{Coproduit et sélecteur natifs}

L'ouverture multiplicative est maintenant définie avant de regarder le codomaine :
\[
 \Delta_\#e_w
 =\sum_{u\boxtimes_\#v=w}e_u\otimes e_v.
\]
La somme s'obtient en énumérant les produits du transducteur. On ne demande pas si
un entier ordinaire en divise un autre.
Dans le secteur entier, nous définissons, sur le cœur algébrique correspondant,
\[
 \Delta_\times e_n
 =\sum_{\substack{a,b\in\N_{\geq1}\\ab=n}}e_a\otimes e_b.
\]
Les deux coproduits énumèrent des couples ordonnés.

\begin{theorem}[Coassociativité native]
Dans le cœur algébrique,
\[
 (\Delta_\#\otimes I)\Delta_\#
 =(I\otimes\Delta_\#)\Delta_\#.
\]
De plus,
\[
 (W_\times\otimes W_\times)\Delta_\#
 =\Delta_\times W_\times.
\]
\end{theorem}

\begin{proof}
Les deux membres de la première égalité énumèrent une fois chaque triplet natif
$a\boxtimes_\#b\boxtimes_\#c=w$. La seconde égalité découle de la bijection
entre fibres induite par \eqref{int:eq:ev-multiplicative}.
\end{proof}

\begin{theorem}[Coproduit arithmétique complet]
\label{int:pf84:C06}
Dans le cœur algébrique
\(c_{00}(\mathbb N_{\geq1})\subset\ell^2(\mathbb N_{\geq1})\), soit
\[
 \Delta_\times e_n=\sum_{ab=n}e_a\otimes e_b.
\]
Sa fermeture a pour domaine
\[
 \operatorname{Dom}(\overline{\Delta_\times})
 =
 \left\{x=(x_n):
 \sum_{n\geq1}d(n)|x_n|^2<\infty\right\},
\]
est coassociative dans le cœur et, pour \(Qe_n=ne_n\), satisfait
\[
 (\log Q\otimes I+I\otimes\log Q)\Delta_\times
 =\Delta_\times\log Q.
\]
\end{theorem}

\begin{proof}
Les supports de \(\Delta_\times e_n\) sont orthogonaux pour des entiers distincts
et leur norme au carré est le nombre \(d(n)\) de diviseurs, ce qui donne le domaine
pondéré de la fermeture. Les deux coassociations réindexent la même somme
sur \(abc=n\). Enfin, \(\log a+\log b=\log n\) dans chaque terme
prouve l'entrelacement. Une triple factorisation comptée de manière différente
ou l'échec de cette identité réfuterait le théorème.
\end{proof}

\begin{theorem}[Coproduit réduit et sélecteur d'irréductibles]
\label{int:pf84:C07}
Dans le secteur entier, le coproduit réduit
\[
 \Delta_\times^{\rm red}e_n
 =\sum_{\substack{ab=n\\a,b>1}}e_a\otimes e_b
\]
est fermable. L'opérateur
\[
 N_\times^{\rm red}
 =\bigl(\overline{\Delta_\times^{\rm red}}\bigr)^*
  \overline{\Delta_\times^{\rm red}}
\]
est positif autoadjoint et diagonal :
\[
 N_\times^{\rm red}e_n=d_{\rm red}(n)e_n.
\]
Son noyau est engendré par \(e_1\) et par les \(e_p\) avec \(p\) irréductible ;
par conséquent,
\[
 P_{\Irr}
 =\mathbf1_{\{0\}}(N_\times^{\rm red})
  -\lvert e_1\rangle\langle e_1\rvert
\]
sélectionne les irréductibles et exclut l'unité.

Dans le secteur natif, le coproduit correspondant est
\[
 \widetilde\Delta_\#e_w
 =\sum_{\substack{u\boxtimes_\#v=w\\u,v\ne(1)}}e_u\otimes e_v.
\]
L'entrelaceur \(W_\times\) transporte
\(\Delta_\times^{\rm red}\), \(N_\times^{\rm red}\) et \(P_{\Irr}\) vers
\(\widetilde\Delta_\#\),
\[
 A_\#:=
 \bigl(\overline{\widetilde\Delta_\#}\bigr)^*
 \overline{\widetilde\Delta_\#},
 \qquad
 P_{\Irr,\#}
 =(I-P_{(1)})\mathbf1_{\{0\}}(A_\#),
 \qquad
 P_{(1)}=\lvert e_{(1)}\rangle\langle e_{(1)}\rvert.
\]
\end{theorem}

\begin{proof}
Les supports des images de deux vecteurs de base distincts sont orthogonaux.
La norme au carré de
\(\Delta_\times^{\rm red}e_n\) est \(d_{\rm red}(n)\), qui s'annule
exactement pour \(n=1\) ou \(n\) irréductible. Cela prouve la diagonalisation
et le noyau déclaré ; soustraire le projecteur de \(e_1\) élimine l'unité.
La multiplicativité de \(W_\times\) identifie les fibres entières aux
fibres natives et transporte les trois opérateurs. Un composé dans le noyau,
un irréductible en dehors de celui-ci ou l'inclusion de l'unité dans \(P_{\Irr}\)
réfuteraient le théorème.
\end{proof}

\begin{proposition}[Coïncidence nonadique des nombres premiers et composés]
\label{int:pf84:B05}
Soit \(\gcd(n,90)=1\),
\(\tau_n=\operatorname{ord}_n(10)\) et
\[
 M_n=\frac{10^{\tau_n}-1}{n}.
\]
Alors \(M_n\equiv0\pmod9\), aussi bien pour les nombres premiers que pour les composés. En
particulier, la clôture nonadique et la période décimale prises isolément ne sont pas un
sélecteur d'irréductibles.
\end{proposition}

\begin{proof}
Comme \(9\mid10^{\tau_n}-1\) et \(n\) est inversible modulo neuf, la
divisibilité par neuf survit au quotient. Les nombres premiers \(7,13\) et les
composés \(77,91,143\) partagent la période six ; les nombres de Carmichael
apportent des contrôles supplémentaires. L'énoncé ne serait réfuté que si la
congruence ou la période séparaient tous les nombres premiers de tous les composés
dans le domaine déclaré.
\end{proof}

Cette affirmation est strictement arithmétique : elle identifie des fibres de
factorisation et les opérateurs qu'elles induisent sur le secteur des formes
normales.

\ifdefined\HMTInternalTraceability
\section{Extension monoïdale du produit natif aux histoires compatibles}
\label{int:pf84:SYN-05D-HIST}

Le résultat distingue l'espace \(\Krad\) des formes normales canoniques et
l'espace \(\mathcal K_{\#,\mathrm{hist}}\) des histoires complètes de
cellules et de quotients propagés.
De nombreuses histoires peuvent produire la même forme normale. Par conséquent,
$W_\times$ est unitaire sur $\Krad$, mais la projection d'une histoire sur
sa valeur empêche qu'il soit injectif sur toutes les classes de transport.

Soit \(\mathscr X_n^{\rm hist}\) l'ensemble des histoires enrichies de
profondeur \(n\) et soit
\(\mathcal H_n=\ell^2(\mathscr X_n^{\rm hist})\). Notons \(P_n\) la
projection orthogonale sur le sous-espace engendré par les histoires
survivantes et supposons définie une multiplication \(\mu_n\) sur un
domaine algébrique de
\(\mathcal H_n\otimes_{\mathrm{alg}}\mathcal H_n\). Une extension monoïdale
du produit de mots au système inverse consiste en une famille
d'applications \(\iota_n:\Krad\to\mathcal H_n\) qui commute avec les
restrictions de profondeur et satisfait, pour chaque \(n\),
\[
 P_n\,\iota_n\mu_\#
 =P_n\,\mu_n
 (\iota_n\otimes\iota_n)
\]
dans le cœur algébrique. Cette égalité est la condition de compatibilité qui
définit l'extension. Les théorèmes précédents n'établissent ni l'existence ni
l'unicité d'une famille \((\iota_n)_n\) ; ils démontrent le produit dans le
monoïde \(\WordMonoid\) et déterminent l'équation qu'une extension à l'espace
des histoires doit satisfaire. Le transfert conserve donc le
statut de \emph{programme ouvert} : il n'hérite pas de l'exactitude du produit dans
\(\WordMonoid\). L'inexistence d'une famille compatible, l'échec de la
commutation avec les troncatures, la perte de survie, de quotient,
de retenue, d'orientation ou de mémoire, ou le non-respect de l'identité précédente
sont ses falsificateurs.
\fi

\section{Énumération exhaustive et nécessité des coordonnées structurelles}

L'énumération finie comprend :
\[
\begin{array}{@{}lr@{}}
\toprule
\text{Test}&\text{Résultat}\\\midrule
\text{Cellules par application locale}&81\\
\text{Parcours réversibles de numérotation}&59050\\
\text{Paires de l’entrelaceur}&532900\\
\text{Erreurs de l’entrelaceur}&0\\
\text{Triplets associatifs}&531441\\
\text{Erreurs associatives}&0\\
\text{Fenêtre de coproduit}&1\text{ à }500\\
\text{Irréductibles obtenus}&95\\
\bottomrule
\end{array}
\]

\begin{definition}[Quotient multiplicatif]
Conserver seulement $R^\times$ transforme $2\cdot5=10$ en la valeur visible $1$.
Perturber $q^\times(9,9)$ de $8$ à $7$ transforme $81$ en $72$. Le quotient
local est nécessaire.
\end{definition}

\begin{definition}[Application additive]
L'application multiplicative produit des produits partiels, mais ne résout pas leurs collisions.
Éliminer $q^+$ rompt la normalisation globale ; par exemple, elle cesse de reproduire
les quotients intermédiaires de $10\cdot18=180$.
\end{definition}

\begin{definition}[Coordonnées et symbole de bord]
Le zéro ordinaire n'est pas un chiffre de la construction. Le mot vide est
$0$ et le symbole $9$ conserve sa valeur et son quotient. Mélanger les coordonnées
positionnelles indexées à partir de zéro avec les étiquettes $1,\ldots,9$ modifie les
cellules.
\end{definition}

\begin{proposition}[Portée exacte des ablations et du recensement fini]
\label{int:pf84:D10}
Les ablations précédentes prouvent, dans les conventions déclarées :
\begin{enumerate}
  \item sans \(q^\times\), \(2\cdot5=10\) se réduit au résidu visible \(1\) ;
  \item remplacer \(q^\times(9,9)=8\) par \(7\) transforme \(81\) en \(72\) ;
  \item sans \(q^+\), la normalisation des collisions additives échoue.
\end{enumerate}
L'énumération exhaustive comprend les \(2\cdot81=162\) cellules des deux
applications, les \(729\) triplets locaux, \(532\,900\) paires de mots et
les \(95\) irréductibles inférieurs ou égaux à \(500\). Dans ces domaines finis,
les identités indiquées sont démontrées, mais la survie
coinductive de toute histoire TPK n'est pas prouvée.
\end{proposition}

\begin{proof}
Les trois contre-exemples se calculent directement avec
\eqref{int:eq:local-plus} et \eqref{int:eq:local-times}. L'énumération des quatre
domaines finis produit les cardinaux écrits dans l'énoncé et n'ajoute
aucune identité coinductive à celles qui ont déjà été démontrées dans le monoïde
de mots.
\end{proof}

\section{Structure algébrique du monoïde de mots}

La construction démontre, dans le monoïde \(\WordMonoid\) :
\begin{enumerate}
 \item une forme normale unique pour chaque entier non négatif ;
 \item une somme et un produit globaux construits uniquement avec les deux applications locales ;
 \item la conservation symbolique de la valeur sous la propagation du quotient ;
 \item un unitaire $W_\times$ et le carré multiplicatif exact ;
 \item un coproduit diviseur et un sélecteur irréductible définis nativement ;
 \item la nécessité de conserver la classe de transport lorsque l'on quitte le
 secteur des formes normales.
\end{enumerate}

Le passage à l’espace holonomique se réduit désormais à un problème d’existence : construire un système \((\iota_n)_n\) compatible avec les restrictions de profondeur et démontrer qu’il satisfait l’équation de la section précédente. Cette question se distingue de l’entrelaceur exact sur \(\WordMonoid\), déjà déterminé par la forme normale et les applications locales.
 

\label{p3-83-c23-s1:ch:foundation}

\section{Problème de détermination des constantes à partir de l’état généalogique}

Une constante physique est représentée comme une section d’une droite de
grandeur munie d’une règle d’évaluation, d’une orientation et d’une normalisation.
Sa coordonnée numérique dépend de la carte choisie, tandis que sa
généalogie reste invariante. Par conséquent, le problème de
détermination comprend :

\begin{enumerate}
\item l’identification de l’état HMT antécédent ;
\item la spécification de l’opérateur ou du caractère de sélection ;
\item la détermination de l’invariant conservé ;
\item la droite de grandeur correspondante ;
\item la carte dans laquelle s’exprime la valeur terminale.
\end{enumerate}

La chaîne causale du travail est

\begin{equation}
\begin{aligned}
\APP&\longrightarrow\TRIT\longrightarrow\TPK
\longrightarrow\text{réalisation interne}
\longrightarrow\text{invariant sans dimension}\\
&\longrightarrow\text{section de grandeur}
\longrightarrow\R.
\end{aligned}
\label{p3-83-c23-s1:eq:master-chain}
\end{equation}

Le terme \(\R\) de la chaîne représente l’évaluation terminale de la
section ; il ne constitue pas le domaine originaire à partir duquel la route est
reconstruite rétrospectivement.

\section{Chaîne antécédente APP--TRIT--TPK et domaine des opérateurs d’extraction}

APP fournit la géométrie des positions, des restes, des quotients et des
frontières. TRIT incorpore l’orientation \((-1,0,+1)\) et distingue
l’ouverture, le seuil et le retour. TPK complète l’état par la retenue, la route,
le feuillet, la mémoire et la survie. En particulier, la transition \(9\to1\)
conserve la phase et modifie l’état de mémoire ; elle ne constitue donc pas
une réinitialisation dépourvue de généalogie.

Les réalisations développées ci-dessous conservent cet antécédent sans le remplacer. Le vide s’obtient par une fonction spectrale de l’opérateur bicouche ; la gravité sélectionne une section au moyen de la structure chronologique \(108\) ; la criticité utilise une famille dodécaphasique ; les constantes de Gauss proviennent de l’évaluation de cylindres ; et le Hamiltonien modulaire est défini sur une frontière qui conserve l’incidence \(6\to8\). La production de valeurs scalaires n’établit pas une identité fonctionnelle entre ces réalisations.

\begin{proposition}[Insuffisance de la coïncidence scalaire]
Soient \(F_1:X_1\to L_1\) et \(F_2:X_2\to L_2\) deux fonctionnelles typées.
L’égalité des coordonnées \(c_1(F_1(x_1))=c_2(F_2(x_2))\) n’induit à elle
seule ni une identification \(X_1\simeq X_2\) ni un entrelaceur
\(F_2I=JF_1\).
\end{proposition}

En particulier, une proximité numérique telle que
\(I_{\rm or}\approx\gammaH a_0\) ne peut être promue en identité sans
construire l’entrelaceur correspondant.

\section{Droites de grandeur, cartes orientées et systèmes de coordonnées dimensionnelles}

Soit \(L_Q\) une droite réelle orientée représentant la grandeur \(Q\).
Une unité \(u_Q\in L_Q\setminus\{0\}\) détermine une coordonnée
provisoire \(L_Q\simeq\R\), sans remplacer la structure intrinsèque de
la droite. La relation constitutive

\[
L_\mu=L_S L_C^{-1}L_Q^{-2},\qquad
L_\varepsilon=L_S^{-1}L_C^{-1}L_Q^2
\]

est antérieure au choix SI. De même,

\[
L_Z=L_SL_Q^{-2},\qquad L_G=L_C^5L_T^2L_S^{-1}
\]

typent l’impédance et la constante gravitationnelle avant d’exprimer leurs grandeurs en
ohms ou en unités SI. Dans cet ouvrage, nous écrirons

\[
Q_{\rm int}=\widehat Q\,u_Q
\]

pour distinguer le caractère adimensionnel \(\widehat Q\) de la section
de référence \(u_Q\).

\section{Torseur de repères et covariance des évaluations dimensionnelles}

Les cartes admissibles forment un torseur et, par conséquent, aucune ne possède un caractère
canonique par rapport aux autres. Le changement de repère
\(u_Q\mapsto\lambda_Q u_Q\) induit la transformation contragrédiente
de la coordonnée. Les identités physiques doivent être homogènes par rapport
à cette action. Par exemple,

\[
\mu\varepsilon=c^{-2},\qquad \frac\mu\varepsilon=Z^2
\]

sont valides dans tout repère compatible. L’unification mathématique et
dimensionnelle n’élimine pas les unités, mais les reconstruit par la
covariance des droites et permet de placer l’invariant avant son expression
en coordonnées.

\section{Évaluation métrologique terminale et interdiction de sélection rétrospective}

Chaque entrée de l’atlas final prend la forme

\begin{equation}
\boxed{\text{antécédent}\to\text{opérateur}\to\text{invariant}
\to\text{signification}\to\text{valeur terminale}.}
\label{p3-83-c23-s1:eq:atlas-rule}
\end{equation}

Cet ordre exprime une dépendance causale : la valeur numérique n’est
déterminée qu’après fixation de l’antécédent, de l’opérateur, de l’invariant et de
l’interprétation. L’omission de l’un de ces éléments réduit le
résultat à une reconstruction numérique et non à une dérivation
généalogique.

\section{Domaine exact de la génération généalogique et de sa carte métrologique terminale}

On adopte comme résultats de départ la génération de
\(\pi,e,\varphi\), la clôture dodécaphasique de \(\alpha\),
la section d’action et la sortie \(\gammaH\) de Barbero--Immirzi. Le
présent développement ne recalibre pas ces constructions ; il détermine les constantes
et les invariants qui résultent de leur application aux réalisations
constitutive, thermique, gravitationnelle, modulaire et spectrale.

La masse de l’électron n’est pas non plus incluse comme problème directeur, puisqu’elle
relève de la Loi générale des masses. Toute apparition de \(m_e\) dans une
identité atomique est traitée comme une donnée provenant de ce domaine et non comme
une nouvelle sélection du présent développement.

\subsection*{Portée logique et domaine de validité}
L’architecture APP--TRIT--TPK et les droites dimensionnelles constituent la
structure de départ. La règle uniforme à cinq champs de l’atlas explicite
sa discipline causale dans chaque entrée.



\section{Taxonomie fonctionnelle de \(114\) expressions numériques et de \(106\) lectures opératorielles}
\label{p3-83-c23-s2:sec:atlas-taxonomico-114-constantes-rev11}
\index{constantes!atlas taxonomique de 114 entrées}

Soit l’ensemble fini indexé
\[
 E_{114}=\{\varepsilon_1,\ldots,\varepsilon_{114}\}
\]
d’expressions numériques et soit \(\mathcal O\) l’espace des lectures du
lecteur triadique. L’application
\[
 T:E_{114}\longrightarrow\mathcal O,
 \qquad \varepsilon_i\longmapsto T_{\rm tri}(\varepsilon_i),
 \qquad |\operatorname{im}T|=106,
\]
classe les expressions selon leur rôle opératoire. Chaque
\(\varepsilon_i\), y compris les cibles comparées, est fixé comme
entrée de \(T\) ; cette application ne génère donc pas les valeurs qu’elle
évalue. La partition fonctionnelle est
\[
\begin{array}{l|r}
\text{famille fonctionnelle}&\text{entrées}\\ \hline
\text{noyau structurel}&5\\
\text{couplage et clôture}&1\\
\text{ponts entre modes}&7\\
\text{angulaire et clôture}&22\\
\text{exponentielle--logarithmique}&12\\
\text{spectral et régularisation}&14\\
\text{hyperbolique et modulaire}&2\\
\text{algébrique et métrique}&50\\
\text{divers typés}&1
\end{array}
\qquad
\sum n_i=114.
\]
La centralité de réseau, les chemins minimaux et les communautés calculées sur
le graphe induit par \(T\) décrivent la connectivité au sein de la classification ;
ils n’établissent ni causalité mathématique ni généalogie HMT.

Sept fibres de \(T\) réunissent quinze expressions ayant une lecture commune. Une
fibre partagée n’identifie pas non plus leurs valeurs. En particulier,
\(\varphi^{-1}\), \(\varphi\) et \(\varphi^2\) sont des réels distincts que
l’application \(T\) regroupe selon une même orbite d’auto-échelle ; la projection
commune conserve la fonction dynamique et oublie la coordonnée archimédienne qui les
sépare.

L’énumération exhaustive a parcouru les \(114\) entrées de \(E_{114}\),
quatre cent soixante-huit candidats directs \(U_{12}\) et cent quatre
mille neuf cent soixante-seize germes pondérés sans trouver, dans
ce domaine fini, de rival reproduisant la correction dodécaphasique
orientée de \(\alpha\). Le domaine de la recherche est ainsi
borné par ces trois ensembles. Le recensement démontre l’absence d’un
rival dans la classe énumérée ; il ne démontre pas l’unicité parmi toutes les
fonctionnelles possibles et n’intervient pas comme prémisse du théorème de
\(\alpha\).

Pour chaque \(\varepsilon_i\), la taxonomie attribue le quadruplet
\((\text{objet},\text{lecteur},\text{domaine},\text{fonction structurelle})\)
utilisé dans l’énumération.
 \section{Confluence spectrale, informationnelle et thermodynamique de \(\log3\)}

L’identité

\[
\log3=C_1+C_2-2C_0
=S_{\TRIT}/k_B
=E_P/(k_B\Theta_{\rm clk}),
\]

relie trois réalisations distinctes. Le premier membre est un défaut
résiduel, le deuxième une entropie de décision et le troisième un rapport
entre énergie et température. Leur égalité démontre la conservation du
scalaire sous le changement de réalisation, sans identifier les espaces
d'états.

La première égalité se déduit des trois canaux
résiduels
\[
Z_r(s)=\sum_{\substack{n\ge1\\n\equiv r\ ({\rm mod}\ 3)}}n^{-s},
\qquad r=0,1,2.
\]
Alors
\[
Z_0(s)=3^{-s}\zeta(s),\quad
Z_1(s)+Z_2(s)=(1-3^{-s})\zeta(s),\quad
Z_1(s)-Z_2(s)=L(s,\chi_{-3}).
\]
Si \(C_r\) désigne la partie finie de \(Z_r\) en \(s=1\), le développement
de \(3^{-s}\) donne
\begin{equation}
C_0+C_1+C_2=\gamma,qquad
-2C_0+C_1+C_2=\log3,qquad
C_1-C_2=\frac{\pi}{3\sqrt3}.
\label{eq:residual-three-covectors}
\end{equation}
Les trois covecteurs sont linéairement indépendants. Ils retrouvent, à partir d'une seule
décomposition résiduelle, le terme fini total, la dilatation de la
classe zéro et l'orientation entre les deux classes non nulles. \(\log3\) est
donc un invariant de changement d'échelle avant d'être interprété comme
entropie ou température.

\section{Classification épistémologique des constantes dérivées}

Les contributions du chapitre sont classées selon trois modalités :

\begin{enumerate}
\item attribution d'un opérateur interne à une constante connue, comme dans
le cas de Catalan ;
\item évaluation composée exacte de deux caractères internes, comme
\(L(1,\chi_{12})\) ;
\item obtention d'une constante à partir d'un corps ou d'une famille
construits en interne, comme \(\gamma_{\mathbb Q(\sqrt3)}\) ou \(A_k\).
\end{enumerate}

La provenance mathématique détermine le statut de récupération ou de
nouveauté ; la proximité décimale n'intervient pas dans cette classification.

\begin{remark}[Contrôle négatif]
Le caractère \(\chi_{12}\) doit avoir exactement le motif
\((+,-,-,+)\) dans \((1,5,7,11)\) ; tout autre motif invalide
\cref{eq:L1chi12,eq:L2chi12}. L'identité
\cref{eq:odd-zeta-bendersky} et la factorisation
\cref{eq:dedekind-sqrt3} sont des contrôles indépendants.
\end{remark}

\subsection*{Portée logique et domaine de validité}
\(J^2=-I\), Catalan, \(\chi_{12}\), ses valeurs \(L\), Euler--Kronecker et la tour de Bendersky--Glaisher : \emph{Théorème interne exact}. La réunion en un réseau \(3/4/12\) est une \emph{Formalisation nouvelle} qui conserve les provenances mathématiques.

 Tous ces résultats sont unis parce qu’ils transforment la valeur terminale en l’expression quantitative d’une structure préalablement construite. L’exactitude numérique accompagne une conception différente de ce qu’est une constante fondamentale. La connexion nonadique n’engendre pas d’abord trois constantes pour en ajouter ensuite une quatrième : le prolongement \(w_6\to w_{12}\to w_{18}\to w_{24}\to w_{30}\to R_{36}\to G_9\) stabilise un état coinductif commun à partir duquel \(\pi\), \(e\), \(\varphi\) et \(\alpha\) sont publiés de manière corrélée. Leurs lecteurs distinguent clôture, propagation, auto-échelle et couplage ; la coordonnée \(\alpha\) conserve, au sein de cette génération conjointe, sa réalisation signée et dodécaphasique et ses deux voies internes de dérivation.\par

Dans la description habituelle, une constante est souvent présentée comme un nombre placé devant une équation : une quantité dont la théorie a besoin, mais dont la valeur doit être reçue de l’extérieur. Dans l’holographie modulaire triadique, c’est l’inverse qui se produit. La constante apparaît à la fin, lorsqu’un même état a été parcouru par plusieurs lecteurs et que l’on vérifie quelle relation leur a tous survécu.\par

C’est pourquoi les constantes du document sont des grandeurs dotées d’une généalogie : des restes invariants d’une transformation. Chacune dit ce qui a été conservé lorsque quelque chose a changé :\par

\begin{itemize}[leftmargin=*,itemsep=0.35em]
\item un feuillet a été inversé ;
\item un cycle est revenu à sa phase initiale ;
\item une mémoire a traversé la couture ;
\item une description intérieure a été publiée au bord ;
\item une grandeur adimensionnelle a reçu une réalisation physique ;
\item un système a changé d’échelle en conservant intégralement sa généalogie.
\end{itemize}

La mécanique dimensionnelle ajoute ensuite les droites de grandeur et les unités. Mais la structure qui sélectionne la constante est antérieure. La valeur réelle apparaît à la fin parce que \(\mathbb R\) fonctionne ici comme évaluation terminale : elle est l’ombre quantitative d’une construction qui était déjà complète avant de devenir un nombre décimal.\par

Dans cette perspective, les neuf développements considérés forment une seule narration sur l’orientation, la mémoire, la géométrie, l’information et l’échelle.\par
  \endgroup
% Rectificación quirúrgica: generación conjunta, profundidad arbitraria y
% estatuto tipado de los transportes ternarios.  Este bloque pertenece al
% capítulo 27 y precede a la separación expositiva de los cuatro lectores.

\section{Théorème conjoint de génération coinductive et de profondeur arbitraire}
\label{sec:c27-generacion-coinductiva-profundidad-arbitraria}

La structure de départ de ce théorème n'est pas l'ensemble nu de neuf
symboles. C'est le système typé
\[
 \mathfrak H_9=
 \bigl(\mathrm{APP},\mathrm{TRIT},\mathrm{TPK};
       \widehat{\mathcal X}^{\mathrm{enr}},\mathcal U,
       \Gamma_9,\operatorname{Ext},\operatorname{tr}\bigr),
\]
construit sur la double réalisation positionnelle et arithmétique de
\(\mathcal D_9=\{1,\ldots,9\}\). APP accumule et replie au moyen de ses feuillets
\(\Sigma\) et \(\Pi\), en conservant résidu et quotient ; TRIT oriente chaque
transition dans l'un de ses trois régimes ; et TPK sélectionne, transporte,
mémorise et revient sur la fibre enrichie. En particulier,
\[
 \mathcal U_t
 =\operatorname{Rec}_t\circ\operatorname{Mem}_t
  \circ T_{d(t)}\circ M_{\mathrm{ph}}(g(t)).
\]
Cette composition est la donnée dynamique effective : la carte à neuf symboles est son support fini, non un substitut à la dynamique.

Pour \(r\in\{6,12,18,24,30,36\}\), soient
\(\widehat{\mathcal X}^{\mathrm{enr}}_r\) les états tronqués avec feuillet,
orientation, résidu, quotient, retenue, voie, bord et mémoire, et soient
\[
 \operatorname{Ext}_{r+6,r}:
 \widehat{\mathcal X}^{\mathrm{enr}}_r
 \longrightarrow
 \widehat{\mathcal X}^{\mathrm{enr}}_{r+6},
 \qquad
 \operatorname{tr}_{r+6,r}:
 \widehat{\mathcal X}^{\mathrm{enr}}_{r+6}
 \longrightarrow
 \widehat{\mathcal X}^{\mathrm{enr}}_r
\]
le prolongement admissible et la troncature. La compatibilité \(\operatorname{tr}_{r+6,r}\operatorname{Ext}_{r+6,r}=\mathrm{id}\) se prolonge à tout niveau. La notation
\[
 w_6\longrightarrow w_{12}\longrightarrow w_{18}\longrightarrow w_{24}
 \longrightarrow w_{30}\longrightarrow R_{36}\longrightarrow G_9
\]
désigne ces applications et leurs changements de type, non une liste d'étiquettes.
L'holonomie \(\Gamma_9\) satisfait
\[
 g(k+9)=g(k),
 \qquad q^{(9)}(k+9)=q^{(9)}(k)+1,
\]
de sorte que la phase revient tandis qu'avancent le quotient, la voie et la
mémoire de l'état.

\begin{theorem}[Construction inaugurale à partir des germes]
\label{thm:c27-construccion-inaugural-semillas}
Soit
\[
 I_9=\{0,\ldots,8\},\qquad D=\{N,E,S,O\},\qquad
 \mathcal S=I_9^2\times D,
\]
et soit \(\Omega_{\mathrm{seed}}=\mathcal S_+\times\mathcal S_\times\) le
domaine des deux curseurs qui réalisent les feuillets additif et multiplicatif
d'APP. Alors
\[
 |\Omega_{\mathrm{seed}}|=(9^2\cdot4)^2=104\,976,
\]
et la récurrence de \(54\) pas et la réduction ternaire définissent
\[
 \Omega_{\mathrm{seed}}
 \xrightarrow{\ T_{54}^{\rm seed}\ }
 \mathcal U_6^{\mathrm{cat}}
 \xrightarrow{\ \rho_3\ }
 \mathcal W_6
 \hookrightarrow\mathbb F_3^6,
\]
avec
\[
 |\mathcal U_6^{\mathrm{cat}}|=468,
 \qquad |\mathcal W_6|=243,
 \qquad |\mathbb F_3^6|=729.
\]
Les fibres satisfont exactement
\[
 432\cdot144+18\cdot432+18\cdot1944=104\,976.
\]
La relation d’extension du TPK prolonge les mots survivants en une section compatible \(x_\infty\) de profondeur infinie. Sur cette section sont construites conjointement l’ombre archimédienne et la réalisation exceptionnelle du \cref{thm:c27-tesis-rectora-doble-proyeccion}.
\end{theorem}

\begin{proof}
Le produit cartésien précédent énumère exhaustivement les positions et orientations des deux curseurs. L’application \(T_{54}^{\rm seed}\) s’obtient en composant à chaque pas la sélection du feuillet, le transport cardinal, l’actualisation du quotient et la mémoire de retenue. L’énumération de son image contient \(468\) sextuplets distincts ; la somme des cardinaux de ses fibres est l’identité exposée. L’application composante par composante \(\rho_3\) contient \(243\) mots distincts. Enfin, la non-vacuité, l’univocité sur le sous-domaine survivant et la compatibilité de \(\operatorname{Ext}\) avec les troncatures produisent la section de la limite projective. Chaque énoncé est une égalité finie ou une conséquence de la compatibilité projective ; aucun n’utilise en entrée un nombre réel, une constante reconnue ou une table décimale.
\end{proof}

\begin{theorem}[Thèse directrice : un continu et deux projections]
\label{thm:c27-tesis-rectora-doble-proyeccion}
Soit \(\Omega_\Gamma^{\mathrm{cal}}\subset
X_\infty^{\mathrm{cal}}\) l'espace des sections survivantes des
tours enrichis du TPK. La retenue des trois tours
\(\gamma_-,\gamma_0,\gamma_+\) définit, sans évaluation archimédienne préalable,
l'homéomorphisme
\[
 \beta:\Omega_\Gamma^{\mathrm{cal}}\xrightarrow{\ \sim\ }
 (\mathbb F_3^6)^{\mathbb N}=:\mathcal W_\infty.
\]
Les formes équilibrées normalisées
\[
 a_k+c_k=\rho_k+3c_{k+1},\qquad
 \rho_k\in\{-1,0,+1\},
\]
construisent l'anneau ordonné \(\mathbb Z_{\rm HMT}\), son corps des
fractions \(\mathbb Q_{\rm HMT}\) et la complétion
\[
 \mathbb R_{\rm HMT}:=
 \operatorname{Cau}(\mathbb Q_{\rm HMT})/{\asymp}.
\]
La structure discrète conjointe est le produit fibré
\[
\mathcal C_{\mathrm{cont}}^{\mathrm{disc}}
=\bigl(\mathbb Z_{\rm HMT}\times
\Omega_\Gamma^{\mathrm{cal}}\bigr)
\times_{\mathcal W_\infty}X_\infty^{\mathrm{cal}}.
\]
Ses deux applications naturelles sont
\[
 P_{\mathrm{ar}}:\mathbb Z_{\rm HMT}\times
 \Omega_\Gamma^{\mathrm{cal}}\longrightarrow\mathbb R_{\rm HMT},
 \qquad
 Q_\infty:X_\infty^{\mathrm{cal}}\longrightarrow
 (\mathscr C_W^{\mathrm{cal}})^{\mathbb N}.
\]
Si \(\beta(\xi)=(b_q)_{q\geq0}\),
\(\nu(b)=\sum_{i=1}^{6}b_i3^{6-i}\) et
\[
 q_n(m,\xi)=m+\sum_{q=0}^{n-1}
 \frac{\nu(b_q)}{729^{q+1}},
\]
alors \((q_n)\) est de Cauchy et \(P_{\mathrm{ar}}(m,\xi)\) est sa classe
dans \(\mathbb R_{\rm HMT}\). Les cylindres rationnels
\[
 I_n(m;b_0,\ldots,b_{n-1})=
 \left[q_n,q_n+729^{-n}\right]
\]
sont compatibles, ont un diamètre tendant vers zéro et la partition récurrente
en \(729\) sous-cylindres construit, pour toute classe de Cauchy, un ruban
\((b_q)\) qui la représente. En conséquence,
\[
 P_{\mathrm{ar}}\ \text{est surjective},\qquad
 (\mathbb Z_{\rm HMT}\times\Omega_\Gamma^{\mathrm{cal}})
 /{\sim_{\rm ar}}\ \cong\ \mathbb R_{\rm HMT}.
\]
Ce n'est qu'après cette construction interne que l'unicité du corps ordonné
complet archimédien reconnaît \(\mathbb R_{\rm HMT}\cong\mathbb R\).

La seconde conserve l’incidence calibrée de chaque niveau ; dans les coordonnées \(x=((w_j,f_j))_{j\geq1}\), elle agit par
\[
 Q_\infty(x)
 =\bigl(f_j,\Gamma_{f_jA_Wf_j^{-1}}(w_j)\bigr)_{j\geq1}.
\]
Ainsi, \(\mathbb R_{\rm HMT}\), reconnu ultérieurement comme \(\mathbb R\), est l’ombre archimédienne obtenue en contractant la généalogie complète des cylindres. La projection exceptionnelle conserve cette généalogie comme incidence de frontière et admet la réalisation successive Paley--Witt--Golay--Leech--\(V^\natural\)--Monstre. Les deux branches proviennent du même état et aucune n’est reconstruite rétrospectivement à partir de l’autre.
\end{theorem}

\begin{proof}
Les trois tours enrichis existent après chaque préfixe et leur retenue
les distingue par \(-1,0,+1\). La concaténation et la troncature sont
inverses, ce qui prouve la bijectivité et la continuité de \(\beta\). Pour une
classe de Cauchy de \(\mathbb Q_{\rm HMT}\), on choisit d'abord son cylindre
unitaire ordonné, puis, récursivement, l'unique sous-cylindre semi-ouvert
de la partition interne en \(729\) morceaux qui contient la classe ; sur un
bord commun, les deux écritures sont identifiées par \(\sim_{\rm ar}\).
La suite des extrémités gauches est précisément \((q_n)\), converge vers la
classe donnée et prouve la surjectivité sans présupposer un développement de
\([0,1]\) ni un point de \(\mathbb R\). L'identification au corps réel
est le théorème d'unicité appliqué ensuite à \(\mathbb R_{\rm HMT}\).

Les sous-espaces satisfaisant la fermeture d’un caractère donné sont des filtres de survie au sein de la capacité globale. Leur fonction est d’individualiser les sections de \(\pi,e,\varphi,\alpha\) ; la surjectivité du continu procède de la complétion de \(\mathbb Q_{\rm HMT}\). Cette séparation des domaines empêche de transformer une condition de sélection de constantes en prémisse du recouvrement de \(\mathbb R\).

Dans la seconde branche, la compatibilité
\(r_{n,m}Q_n=Q_m p_{n,m}\) permet de passer à la limite et définit \(Q_\infty\).
La conjugaison par \(f_j\) transporte le calibre sans effacer l'incidence.
Les foncteurs ultérieurs de code, de réseau et d'algèbre de vertex réalisent
la chaîne exceptionnelle. L'égalité des rubans dans le produit fibré
construit la source conjointe et conserve des informations complémentaires : valeur
archimédienne dans la première ; incidence, bord et généalogie dans la seconde.
\end{proof}

\begin{corollary}[Développement infini comme objet mathématique]
\label{cor:c27-expansion-infinita-objeto}
Pour chaque caractère numérique \(\chi\) publié par l'état coinductif, la
lecture décimale est une séquence complète
\[
 \operatorname{Dec}_\chi(x_\infty)
 =(d_{\chi,1},d_{\chi,2},\ldots)
 \in\{0,\ldots,9\}^{\mathbb N}.
\]
Son préfixe de longueur \(N\) est la projection
\(\rho_N\operatorname{Dec}_\chi(x_\infty)\). Ainsi, mille, dix mille ou un
million de chiffres sont des coupes du même objet infini ; la génération ne
consiste pas à relever successivement une borne expérimentale.
\end{corollary}

\begin{theorem}[Système corrélé à quatre coordonnées de l’état holonomique coinductif intégral]
\label{thm:c27-cuadrirrelacion-generada}
L’unique composition APP--TRIT--TPK du système \(\mathfrak H_9\) détermine un état compatible
\[
 x_\infty\in
 \widehat\Omega:=
 \varprojlim_r\widehat{\mathcal X}^{\mathrm{enr}}_r
\]
et quatre caractères corrélés
\[
 \bigl(\Pi_{\mathrm{cl}},\Pi_{\mathrm{prop}},
       \Pi_{\mathrm{aut}},\Pi_{\mathrm{dod}}\bigr)(x_\infty)
 =\bigl(\pi,e,
        \varphi,\alpha\bigr).
\]
Les trois premiers caractères proviennent respectivement de l’unique orbite de clôture, de l’unique orbite de propagation et de l’unique orbite d’autoéchelle du catalogue enrichi. Le quatrième provient de la coordonnée primitive signée du même état terminal et de sa fermeture dodécaphasique. Les quatre coordonnées sont donc des sorties d’un seul prolongement coinductif, bien que leurs applications d’expression soient distinctes.
\end{theorem}

\begin{proof}
Le recensement exhaustif
\(104,976\to468\to243\subset\mathbb F_3^6\) décompose l'image réalisée
en \(43\) orbites diédriques. Les invariants de fibre laissent une seule orbite
dans chacune des classes de clôture, de propagation et d'auto-échelle ; le calibre
interne y oriente, sans consultation décimale, les représentants
\[
 w_6^{\mathrm{cl}}=010211,
 \qquad w_6^{\mathrm{prop}}=201101,
 \qquad w_6^{\mathrm{aut}}=121200.
\]
Le prolongement conserve les prédicats définissant les trois classes. Dans le canal de clôture, la composition pentarégionale à retour unité produit les rapports \(1/5\), \(1/239\) et la relation \(120/119\) ; l’identité de Machin reconnaît exactement le demi-tour. Dans le canal de propagation, la loi de semi-groupe normalisée impose \((n+1)a_{n+1}=a_n\), donc \(a_n=1/n!\) et \(P_1=e\). Dans le canal d’autoéchelle, l’incidence
\[
 F_{\mathrm{av}}=\begin{pmatrix}0&1\\1&1\end{pmatrix}
\]
possède un unique rayon positif et son équation propre impose
\(x^2-x-1=0\), \(x>0\), c'est-à-dire \(x=\varphi\).

Dans la phase terminale de l’unique structure APP--TRIT--TPK, la lecture \(x_{\mathrm{term}}\mapsto(B_{\mathrm{term}},U_{12}^{\mathrm{sgn}})\) exprime simultanément la carte des trois canaux et la coordonnée primitive signée. Cette dernière s’identifie réversiblement au sceau \(K\) et à ses différences dodécaphasiques. La récurrence euclidienne avec retenue
\[
 s_m=P_m+E_m-\Phi_m-K_m+c_{m+1},\qquad
 A_m=s_m\bmod1000,\qquad
 c_m=(s_m-A_m)/1000
\]
produit \(\alpha\) sans la recevoir en entrée. La caractérisation analytique indépendante par le noyau \(E_{21/22}\) détermine une racine simple dans \((0,10^{-2})\), dont l’inverse coïncide avec celui de la première voie après application de \(\operatorname{Tr}_9\). La séparation des lecteurs ne modifie donc pas l’origine commune des quatre coordonnées.
\end{proof}

\begin{proposition}[Registre dodécaphasique de l'état complet antérieur à \(\alpha\)]
\label{prop:c27-registro-u-k-producido}
L’état holonomique terminal possède la double expression
\[
x_{\mathrm{term}}
\xrightarrow{(\pi_{\mathrm{term}},R_{\mathrm{sgn}})}
(B_{\mathrm{term}},U),
\]
où
\[
U=(2378,1406,2479,-452,998,-551,
-668,-204,-371,-322,-28,-997).
\]
En ordonnant les coordonnées selon les trois orbites de pas trois, la transformée
\[
\mathcal H_{12}
=P_{\mathrm{orb}}^{-1}(H_4\oplus H_4\oplus H_4)P_{\mathrm{orb}},
\qquad H_4^2=4I_4,
\]
possède l'inverse \(\mathcal H_{12}/4\) sur son image entière et détermine
de manière univoque
\[
K=(234,543,140,729,659,824,621,058,914,794,146,601).
\]
Par conséquent, au sein de l’état holonomique intégral,
\[
x_{\mathrm{term}}\xrightarrow{R_{\mathrm{sgn}}}U
\xrightarrow{\mathcal H_{12}^{-1}}K
\]
est une composition exacte antérieure à la retenue qui exprime \(\alpha\). La matrice visible \(B_{\mathrm{term}}\) et le registre signé \(U\) sont deux coordonnées du même état ; le second conserve les champs d’orientation, d’opposition, de quotient et de registre que la première contracte. Il n’existe ici ni seconde APP ni profil réduit concurrent : APP produit les feuillets et le registre arithmétique, TRIT les oriente et TPK accumule dans son registre la fenêtre signée que \(R_{\mathrm{sgn}}\) projette. Exiger qu’APP isolée exprime une coordonnée appartenant au codomaine ultérieur du TPK constitue une erreur de type. La composition complète ne reçoit pas \(U\), \(K\), leurs différences, \(\alpha\) ni des développements cibles en entrée.
\end{proposition}

\begin{theorem}[Clôture bidirectionnelle ultérieure du caractère \(\alpha\)]
\label{thm:c27-alpha-segunda-via-bidireccional}
Soit \(\alpha=\alpha\) la coordonnée déjà exprimée par la retenue dodécaphasique. L’évaluation du jet de lacunes de degrés \(7{:}9\) sur cette sortie détermine la fonctionnelle
\[
\begin{aligned}
B_9(\alpha):={}&\log_{10}\!\left(\frac{10}{9}\right)
+\frac74\alpha^2-\frac{\alpha^4}{20}
+\frac{2\alpha^5}{21}+\frac{\alpha^6}{46}\\
&+\frac{\alpha^7}{120}-\frac{\alpha^8}{45}
-\frac{2\alpha^9}{495}.
\end{aligned}
\]
En relevant la rotation à l'inconnue positive \(T\), la résolvante équivaut à
\[
\boxed{\alpha T^2-B_9(\alpha)T+\alpha^3=0.}
\]
Avec
\(\xi=\operatorname{arcosh}(B_9(\alpha)/(2\alpha^2))\), les solutions sont
\[
T_\pm=\alpha e^{\pm\xi},
\qquad T_+T_-=\alpha^2,
\qquad
J_\alpha(T)=\frac{\alpha^2}{T},
\qquad J_\alpha(T_\pm)=T_\mp.
\]
La chambre \(6<T<7\) sélectionne \(T_+\) ; la clôture finie \(\pi_9=T_+/2\) est compatible, à la résolution déclarée, avec le canal de clôture que \(\Gamma_9\) prolonge indépendamment jusqu’à \(\pi\). L’identité \(\alpha=\sqrt{T_+T_-}\) fixe la moyenne géométrique des branches conjuguées. Le jet \(7{:}9\), les deux racines et l’involution constituent une clôture bidirectionnelle exacte \emph{sur} le caractère déjà engendré. Ce théorème ne reçoit pas de constante conventionnelle, mais ne constitue pas non plus une seconde dérivation indépendante de \(\alpha\) : son domaine contient \(\alpha\) et son résultat est un certificat fonctionnel ultérieur.
\end{theorem}

\begin{theorem}[Expression exceptionnelle corrélée au secteur \(\alpha\)]
\label{thm:c27-alpha-rama-excepcional}
Dans la carte d'incidence du même registre dodécaphasique, les racines radiales
de \(A_2^{12}\) déterminent
\[
v_\alpha=4\rho_1+\rho_2+\cdots+\rho_{12},
\qquad v_\alpha^2=54,
\]
et le voisin d'indice trois
\[
N_{\alpha,0}
=\{x\in N(C_W):\langle x,v_\alpha\rangle\equiv0\pmod3\},
\qquad
N_\alpha=N_{\alpha,0}+\mathbb Z\frac{v_\alpha}{3}.
\]
Le réseau \(N_\alpha\) est pair, unimodulaire, de rang vingt-quatre et sans
racines, de sorte que \(N_\alpha\cong\Lambda_{24}\). La chaîne
\[
C_W\longrightarrow N(A_2^{12})
\xrightarrow{v_\alpha}\Lambda_{24}
\longrightarrow V^\natural
\xrightarrow{\operatorname{Aut}}\mathbb M
\]
expose la charnière exceptionnelle. Le scalaire \(\alpha\) et le vecteur \(v_\alpha\) appartiennent à des codomaines distincts et partagent l’état dodécaphasique d’origine. Le résultat exact de cette branche est la construction \(N_\alpha\cong\Lambda_{24}\to V^\natural\) et l’action ultérieure de \(\operatorname{Aut}(V^\natural)\cong\mathbb M\) sur \(V^\natural\). Le scalaire n’est pas identifié au vecteur et la corrélation ne permet pas de déduire une action du Monstre sur les chiffres ou les parcours ; une telle promotion exigerait un morphisme équivariant de domaine et codomaine explicites.
\end{theorem}

\begin{proposition}[Statut des transports \(L_0,L_1\)]
\label{prop:c27-estatuto-l0-l1}
Les endomorphismes \(L_0,L_1\in\operatorname{End}(\mathbb F_3^6)\) et le calendrier \((L_0,L_0,L_1,L_1)\) appartiennent à la loi de transport de l’état holonomique intégral. Les six transitions exprimées de rang plein fixent chaque matrice de manière unique par \(L=X^{-1}Y\) dans \(\mathbb F_3\), et le calendrier est l’unique parmi les \(16\) calendriers binaires qui conserve simultanément les trois trajectoires structurelles. Leur fonction est de prolonger les représentants déjà sélectionnés ; l’évaluation archimédienne et la nomenclature conventionnelle s’appliquent ensuite.
\end{proposition}

Le domaine de cette proposition est l’unique structure \(\mathfrak H_9\). Ses projections d’oubli sont des lecteurs ultérieurs et ne constituent pas des profils générateurs alternatifs. Ce typage empêche de confondre une ombre modale avec l’état qui la produit.

\begin{theorem}[Profondeur arbitraire et détermination des développements complets]
\label{thm:c27-profundidad-arbitraria}
Pour chaque \(\chi\in\{\pi,e,
\varphi\}\) et chaque \(N\ge1\), il existe un unique préfixe décimal \(p_{\chi,N}\) de longueur \(N\), produit par une troncature finie de l’état holonomique intégral, tel que
\[
 \rho_{N+1,N}(p_{\chi,N+1})=p_{\chi,N}.
\]
La famille compatible est la famille des projections d'un unique élément
\[
 p_{\chi,\infty}
 \in\varprojlim_N\{0,\ldots,9\}^{N}
 \cong\{0,\ldots,9\}^{\mathbb N},
\]
et l'intersection de ses cylindres archimédiens est constituée d'un unique réel.
Par conséquent, la connexion engendre le développement complet ; chaque précision
finie est une projection de cette sortie infinie.
\end{theorem}

\begin{proof}
Pour une précision \(N\), les encadrements rationnels de clôture et de propagation ainsi que les quotients consécutifs de Perron sont raffinés jusqu’à être contenus dans un unique cylindre décimal de diamètre \(10^{-N}\). L’application \(\operatorname{Ext}\), le cocycle de retenue et la troncature TPK font que le cylindre d’ordre \(N+1\) se projette sur celui d’ordre \(N\). Les cylindres sont emboîtés et leurs diamètres tendent vers zéro ; leur limite projective détermine un point unique. Puisque \(N\) est arbitraire et que la relation d’extension n’a pas de profondeur terminale, la conclusion vaut pour toute longueur finie et, conjointement, pour le développement infini.
\end{proof}

\begingroup\fontsize{10}{11.4}\selectfont
\setlength{\parskip}{1.5pt}\setlength{\abovedisplayskip}{4pt}\setlength{\belowdisplayskip}{4pt}
L’exécution de \(396\) niveaux, \(2376\) trits, \(43\) tours, \(387\) couples \((K,K+9)\) par canal et \(1000\) chiffres vérifie une instance matérielle étendue du théorème. L’exécution de \(10,000\) chiffres vérifie une autre instance. Le quantificateur \(\forall N\) et la compatibilité \(\rho_{N+1,N}p_{N+1}=p_N\), avec l’existence de l’élément \(p_\infty\) de la limite projective, démontrent la sortie décimale infinie ; les calculs finis sont des certificats reproductibles de ses projections.

La coordonnée \(\alpha\) appartient au même état limite. Sa voie entière se prolonge par \(\mathsf T_{\alpha,\infty}\), tandis que la racine simple de \(E_{21/22}\) fournit un second lecteur interne. La stabilité vérifiée à \(10,000\) chiffres certifie la propagation lointaine de la retenue ; l’objet mathématique exprimé par les deux lecteurs est fixé avant toute comparaison métrologique.

\begin{proposition}[Prolongement complet de la voie entière de \(\alpha\)]
\label{prop:c27-alpha-prolongacion-completa}
Soit \(B=1000\), soient \((P_k),(E_k),(\Phi_k)\) les trois suites de blocs
publiées par l'état coinductif et soit \((K_k)\) la suite
dodécaphasique prolongée. Les quatre séries
\[
 p=\sum_{k\geq1}P_kB^{-k},\qquad
 e_0=\sum_{k\geq1}E_kB^{-k},\qquad
 \phi_0=\sum_{k\geq1}\Phi_kB^{-k},\qquad
 \kappa=\sum_{k\geq1}K_kB^{-k}
\]
convergent absolument et déterminent
\[
 \alpha_A=p+e_0-\phi_0-\kappa.
\]
La récurrence de retenue est l'écriture en base \(B\) de cette égalité.
Pour chaque \(M\), une troncature suffisamment profonde détermine les
premiers \(M\) chiffres du développement canonique de \(\alpha_A\) ; l'erreur de
queue est bornée par une constante multipliée par \(B^{-N}\). Par
conséquent, \(\mathsf T_{\alpha,\infty}\) publie une suite décimale
complète, et ses exécutions à \(3000\) ou \(10\,000\) chiffres sont des projections
finies de cette sortie.
\end{proposition}

\begin{proof}
Chaque bloc appartient à un ensemble entier borné, de sorte que les quatre
séries sont dominées par une série géométrique. La somme partielle jusqu'à \(N\)
diffère de la somme complète de \(O(B^{-N})\). La division euclidienne de
droite à gauche transporte exactement cette queue au moyen de la retenue ;
par unicité du développement canonique, tout préfixe qui reste séparé d'un
bord de cylindre se stabilise lorsque \(N\) augmente, et sur un bord
on applique la convention canonique non terminale. Cela définit simultanément
tous les chiffres, non une suite d'objectifs numériques indépendants.
\end{proof}

\begingroup\fontsize{10.2}{11.7}\selectfont
\setlength{\parskip}{2pt}\setlength{\abovedisplayskip}{5pt}\setlength{\belowdisplayskip}{5pt}
\begin{corollary}[Séparation entre génération et reconnaissance]
\label{cor:c27-generacion-reconocimiento}
Les identités de Machin, la série exponentielle, le rayon de Perron et la carte décimale sont des réalisations exactes de caractères préalablement sélectionnés par \(\mathfrak H_9\). Les noms conventionnels \(\pi,e,\varphi,\alpha\) et les tables métrologiques reconnaissent les sorties ; ils ne déterminent ni orbites, ni transports, ni retenues, ni préfixes.
\end{corollary}
 

\par\endgroup
\par\endgroup
\chapter{Construction du générateur elliptique \texorpdfstring{\(i=\sqrt{-1}\)}{i=sqrt(-1)} et géométrie de phase}
\begingroup
\renewcommand{\chapter}[2][]{}%
\chapter{Construction interne de \(\sqrt{-1}\) : rotation d’ordre \(4\), opérateur complexe et géométrie de phase}
\label{chap:p3-final:32:raiz_menos_uno}
\label{chap:p3-83-26-raiz_menos_uno}

\section*{La racine interne de \texorpdfstring{\(-1\)}{-1} et la famille exponentielle tritique}

La racine de \(-1\) n’est pas introduite comme un symbole formel sans antécédent. Les
six unités modulo neuf sont munies de coordonnées sur
\(\mathbb P^1(\mathbb F_5)\), et le caractère quadratique de
\(\mathbb F_5\) détermine la matrice de conférence de Paley. Sa réduction
ternaire est
\[
 A_W=
 \begin{pmatrix}
 0&1&1&1&1&1\\
 1&0&1&2&2&1\\
 1&1&0&1&2&2\\
 1&2&1&0&1&2\\
 1&2&2&1&0&1\\
 1&1&2&2&1&0
 \end{pmatrix}\in M_6(\mathbb F_3).
\]
L’identité de conférence se réduit à
\begin{equation}\label{p3-83-c26-s1:eq:constantes-raiz-interna-menos-uno}
 \boxed{A_W^2=2I_6=-I_6.}
\end{equation}
Par conséquent, \(A_W\) est un opérateur de quart de tour dans l’espace ternaire.
Après extension des scalaires à \(\mathbb F_9\), il se décompose en les espaces
propres conjugués des racines de \(X^2+1\). La source entière commune
\[
 \mathcal Q_{\mathbb Z}=\mathbb Z[u]/(u^2+1)
\]
possède une réalisation finie \(u\mapsto A_W\) et une réalisation réelle
\(u\mapsto J_{\rm e}\). La source coordonne la relation quadratique ; elle n’identifie pas
des corps de caractéristiques distinctes.

Les trois régimes de TRIT se réalisent par des opérateurs
\[
 J_{+1}^2=-I,\qquad J_0^2=0,\qquad J_{-1}^2=I.
\]
En séparant les puissances paires et impaires de l’exponentielle, on obtient une
seule famille :
\begin{equation}\label{p3-83-c26-s1:eq:constantes-euler-extendido}
 \boxed{
 \exp(tJ_\tau)=C_\tau(t)I+S_\tau(t)J_\tau,
 \qquad \tau\in\{+1,0,-1\},}
\end{equation}
où
\[
 (C_{+1},S_{+1})=(\cos,\sin),\qquad
 (C_0,S_0)=(1,t),\qquad
 (C_{-1},S_{-1})=(\cosh,\sinh).
\]
Ses trois spécialisations directrices sont
\begin{equation}\label{p3-83-c26-s1:eq:constantes-euler-tres-ramas}
 \exp(\pi J_{+1})+I=0,\qquad
 \exp(J_0)=I+J_0,\qquad
 \exp\!\bigl(2(\log\varphi)J_{-1}\bigr)=F_{\rm av}^{\,2}.
\end{equation}
La branche elliptique contient la clôture circulaire complexe ; la parabolique, le
transport de seuil ; l’hyperbolique, l’auto-échelle dorée. L’équation
d’Euler apparaît ainsi comme une spécialisation d’une famille qui coordonne rotation,
frontière et expansion hyperbolique.


\section{Réalisation elliptique de \(J^2=-I\) par la fibration de Hopf et l’holonomie de Möbius}

\subsection{Structure complexe et fibration de Hopf}

Soit \(S^3\subset\mathbb C^2\) et
\[
h(z_0,z_1)=
\bigl(2\Re(z_0\overline z_1),
2\Im(z_0\overline z_1),|z_0|^2-|z_1|^2\bigr)
\in S^2.
\]
Dans le tore de Hopf
\[
T_\eta=\{(z_0,z_1):|z_0|=\cos\eta,
|z_1|=\sin\eta\},
\]
les courbes
\[
\gamma_+(t)=(\cos\eta\,e^{it},\sin\eta\,e^{it}),
\]
\[
\gamma_-(t)=(\cos\eta\,e^{it},\sin\eta\,e^{-it})
\]
présentent des comportements distincts : \(\gamma_+\) est une fibre de Hopf et
\(\gamma_-\) se projette avec un degré deux. Sous projection stéréographique, les
classes homologiques \((1,1)\) et \((1,-1)\) produisent les deux familles diagonales
de Villarceau.

La conjugaison complexe totale
\[
K(z_0,z_1)=(\overline z_0,\overline z_1)
\]
préserve chaque famille non orientée et inverse son orientation. Elle n’échange pas
les deux familles. La conjugaison partielle peut les échanger, mais ce n’est ni une
application complexe linéaire ni l’antiunitaire standard global. Cette correction
est essentielle pour ne pas fabriquer une symétrie particule--antiparticule à partir du
dessin toroïdal.

\subsection{Revêtement double et retour spinoriel}

L’application \(SU(2)\to SO(3)\) est un revêtement double. Un tour de
\(2\pi\) dans \(SO(3)\) peut se relever en \(-I\) dans \(SU(2)\) ; le tour de
\(4\pi\) récupère \(I\). La structure de signe ne transforme pas un cercle en
ruban de Möbius. Pour obtenir un ruban, il faut un fibré réel en droites
dont le caractère d’holonomie soit \(-1\).

Soit \(L_\chi\to S^1\) le fibré associé à
\[
\chi:\pi_1(S^1)\longrightarrow\{\pm1\},
\qquad \chi(1)=-1.
\]
Son espace total est un ruban de Möbius. La route \(C_M=-\operatorname{rev}\)
apporte le candidat de signe ; le cercle de Villarceau apporte le lacet ;
l’entrelaceur physique doit prouver que le premier est l’holonomie du second.

\subsection{Réalisation toroïdale de l’opérateur \texorpdfstring{\(B_{\mathrm{mem}}^{\varepsilon}\)}{Bmem} et détermination de l’angle de Villarceau}

Avec \(B_c=7I+2S\), définissons l’opérateur orienté de mémoire
\[
B_{\rm mem}^{\varepsilon}=5I+2\varepsilon S,
\qquad\varepsilon=\pm1.
\]
Pour \(\varepsilon=+1\),
\[
\Spec(B_{\rm mem}^{+})=\{7,3\},
\qquad\det B_{\rm mem}^{+}=21.
\]
Interpréter les valeurs \((7,3)\) comme les bords équatoriaux extérieur et intérieur
dans une carte positive commune produit
\[
R+r=7\ell,\qquad R-r=3\ell,
\]
et donc
\[
R=5\ell,\qquad r=2\ell,\qquad
\sin\beta=\frac rR=\frac25.
\]
L’angle de Villarceau est
\[
\beta=\arcsin\frac25
=23.5781784782\ldots^\circ.
\]

Le théorème algébrique \(\{7,3\}\leftrightarrow(5,2)\) est exact. La prémisse
géométrique qui interprète le spectre comme les deux bords d’un tore est une
interface déclarée. Elle ne doit pas être dissimulée et l’angle ne doit pas être employé comme cible.

La réalisation du bloc complet \(B_c\) conserve une autre lecture :
\(R:r=7:2\) et \(\beta=\arcsin(2/7)\). Ce ne sont pas des valeurs rivales ; elles relèvent de
deux foncteurs : bloc total et mode de mémoire. La priorité physique du second
doit être décidée par l’entrelaceur avec la dynamique, et non par une proximité décimale avec
une autre constante.

\subsection{Action de l’opérateur \texorpdfstring{\(J^2=-I\)}{J2=-I} sur les coordonnées angulaires conjuguées}

L’angle natif de HMT est une phase \(u\in\mathbb R/\mathbb Z\). Les lecteurs
publient
\[
\theta_{\rm rad}=2\pi u,
\qquad
\theta_{\rm deg}=360u.
\]
Radians et degrés sont des cartes distinctes du même élément circulaire. Le radian
n’est pas plus « ontologique » parce qu’il est adimensionnel ; sa naturalité provient de la
longueur d’arc. Les degrés fournissent une autre coordonnée absolue du
tour. Toute relation avec \(\alpha^{-1}\), \(\varphi\) ou CKM doit
d’abord être formulée comme un quotient ou une holonomie invariante, puis seulement
être évaluée dans une carte angulaire.

L’identité exacte déjà disponible
\[
\tanh^2(2\log\varphi)=\frac59
\]
relie le canal Perron--Fibonacci au quotient spectral de \(B_c\). Son
complément est \(4/9\), le saut normalisé. Cette relation est structurelle.
Elle n’implique pas que tout angle construit avec \(\varphi\) soit un angle physique.

\subsection{Brouwer, Möbius et degré}

Le théorème du point fixe de Brouwer, les transformations fractionnaires de
Möbius et le ruban de Möbius sont trois objets distincts. Ils ne sont reliés
qu’après fixation du domaine, de l’action et du fibré. Les orbites elliptiques d’avance et de
retard peuvent être modélisées comme deux sections d’un fibré à holonomie de
signe ; la transformation de Möbius décrit l’action conforme dans la carte ;
le théorème de Brouwer garantit un point fixe pour les applications continues du disque.
Aucune de ces phrases ne remplace le diagramme de réalisation.

\begin{remark}[Résultat et frontière]
On démontre le degré double de la courbe de Hopf, l’orientation de la
conjugaison totale et la construction algébrique \(7,3\mapsto5,2\). La lecture
toroïdale (5{:}2), l’holonomie de Möbius et l’interprétation
particule--antiparticule forment une chaîne d’interfaces qui doit partager un
unique entrelaceur. Le volume n’impose pas cet entrelaceur au moyen d’une
coïncidence angulaire.
\end{remark}
  \endgroup

\chapter{Prolongement de Hensel--Witt de \texorpdfstring{\(-1/12\)}{-1/12} et mémoire modulaire}
\begingroup
\renewcommand{\chapter}[2][]{}%
\chapter{Dérivation Hensel--Witt de \(-1/12\) : prolongement compatible, régularisation et mémoire modulaire}
\label{chap:p3-final:33:menos_un_doceavo}
\label{chap:p3-83-27-menos_un_doceavo}

\section{Domaine Hensel--Witt, applications de Teichmüller et cible rationnelle \(-1/12\)}

Le NTHW est le paquet de données et d’applications
\[
\NTHW=
\bigl(
\mathscr R_\pi^{\rm micro},
\Gnine,
\Cdoce,
K,
\mathcal W_{12\to24},
\mathcal I_{6\to8}
\bigr),
\]
avec les compatibilités suivantes :
\begin{enumerate}
  \item les régions de Hensel partagent un mot quotient et conservent
  des mémoires distinctes ;
  \item \(\Gnine\) transporte l’orientation et produit une monodromie ;
  \item \(\Cdoce\) organise l’observable signé et le sceau \(K\) ;
  \item la lecture de Witt transforme les supports du même bord en hexades et en
  étoiles ;
  \item l’extension \(W_{12}\to W_{24}\) produit l’interface
  hexade--octade ;
  \item les lecteurs archimédien, exceptionnel et dimensionnel reçoivent
  des invariants du même état.
\end{enumerate}

\begin{figure}[!htbp]
\centering
\resizebox{\textwidth}{!}{\begin{tikzpicture}[font=\small,node distance=1.2cm and 1cm,
  nodebox/.style={draw,rounded corners=2mm,align=center,minimum width=3cm,minimum height=1.05cm},
  arr/.style={-{Stealth[length=2.4mm]},thick}]
  \node[nodebox,draw=hmtblue,fill=hmtlightblue] (hensel) at (-5,2.3) {cylindres henséliens\\cocycle de retenue et \(\MonNine\)};
  \node[nodebox,draw=hmtblue,fill=hmtlightblue] (regions) at (-5,-.1) {cinq régions de \(\pi\)\\régions de \(e,\varphi\)};
  \node[nodebox,draw=hmtgold,fill=hmtlightgold,minimum width=3.8cm,minimum height=2.05cm,font=\bfseries] (core) at (0,1.1) {\(\mathscr T_{\rm HW}\)\\Correspondances\\Hensel--Witt};
  \node[nodebox,draw=hmtviolet,fill=hmtlightviolet] (c12) at (0,-1.8) {\(\Cdoce\), \(\Usgn\), \(K\)\\différences \(D_3,D_4,Q\)};
  \node[nodebox,draw=hmtgreen,fill=hmtlightgreen] (witt) at (5,2.3) {Witt \(W_{12}\)\\supports et \(M_{12}\)};
  \node[nodebox,draw=hmtgreen,fill=hmtlightgreen] (oct) at (5,-.1) {inclusion d’incidence \(6\to8\)\\canaux \(90/120\), normalisation \(-1/12\)};
  \node[nodebox,draw=hmtred,fill=hmtlightred] (readers) at (5,-2.5) {lecteurs coordonnés\\valeur et symétrie};
  \draw[arr,hmtblue] (hensel)--(core);
  \draw[arr,hmtblue] (regions)--(core);
  \draw[arr,hmtgold] (core)--(c12);
  \draw[arr,hmtgreen] (core)--(witt);
  \draw[arr,hmtgreen] (core)--(oct);
  \draw[arr,hmtred] (c12)--(readers);
  \draw[arr,hmtred] (witt.east)--++(1.9,0)|-(readers.east);
  \draw[arr,hmtred] (oct)--(readers);
  \node[align=center,hmtgray,font=\footnotesize] at (0,-3.45) {Compatibilité des cylindres, de la dodécaphase et des incidences sur un même état holonomique intégral.};
\end{tikzpicture}
 }
\caption{Le NTHW comme interface typée. La tâche mathématique centrale consiste à faire
commuter le diagramme, et non à accumuler des coïncidences autour d’un nom.}
\label{p3-83-c27-s1:fig:nthw}
\end{figure}

Le NTHW donne un nom formel au lieu auparavant décrit comme « pierre de
Rosette ». La substitution est conceptuelle : une pierre traduit par ressemblance ; une
transduction précise l’entrée, la sortie, la donnée conservée et la perte
d’information.

\section{Application de Teichmüller du code ternaire à l’anneau de Witt}

La frontière réversible d’APP sélectionne, avec des jauges déclarées, une matrice
\(A_W\) sur \(\FF_3\) qui satisfait
\[
A_WA_W^T=2I\pmod3.
\]
Le code graphe
\[
C_W=\{(q,qA_W):q\in\FF_3^6\}
\]
est le code de Golay ternaire étendu. Son énumérateur de poids est
\[
1+264y^6+440y^9+24y^{12}.
\]
Les \(264\) mots de poids six, identifiés au signe près, donnent \(132\)
supports. Chaque sous-ensemble de cinq points appartient à une unique hexade :
\[
W_{12}=S(5,6,12).
\]

Le champ de \(495\) involutions étoile engendre un groupe d’ordre
\(95\,040\), reconnu comme \(M_{12}\). Une étoile isolée n’engendre pas le
groupe. Le repère HMT ordonné possède un stabilisateur trivial et fixe donc une
jauge dans l’action.

\section{Incidence hexade--octade et canaux modulaires \(90/120\)}

Soit \(H\) une hexade et \(O_H\) une octade qui la contient. En marquant un point
et une paire de points de l’hexade, on obtient
\[
F_6(H)=H\times\binom H2,
\qquad
|F_6|=6\binom62=90.
\]
Si le point marqué peut parcourir l’octade :
\[
F_8(H)=O_H\times\binom H2,
\qquad
|F_8|=8\binom62=120.
\]
L’inclusion \(H\hookrightarrow O_H\) induit \(F_6\hookrightarrow F_8\).
Son défaut normalisé est
\[
\boxed{
\frac{90-120}{24\binom62}
=\frac{6-8}{24}
=-\frac1{12}.}
\]
Ce théorème combinatoire précède la lecture angulaire de Barbero. Les
deux objets partagent l’interface \(6\to8\), mais ne sont pas identiques.

\section{4 routes compatibles vers \texorpdfstring{\(-1/12\)}{-1/12}}

La construction contient quatre réalisations qui doivent être exposées ensemble et non
comptées comme quatre preuves indépendantes d’une même provenance.

\subsection{Réalisation angulaire par l’incrément élémentaire dodécaphasique}

Une dent de \(\Cdoce\) mesure \(30^\circ\). Alors
\[
\frac{30}{120}-\frac{30}{90}
=\frac14-\frac13
=-\frac1{12}.
\]
C’est l’ombre angulaire du motif \(90/120\).

\subsection{Pseudo-inverse orientée de l’opérateur de Gram sur le sous-espace survivant}

Dans le saut certifié \(60\to81\), deux fibres ont pour tailles \(0\) et
\(12\). L’opérateur de Gram est
\[
K_{60,81}=\operatorname{diag}(0,12).
\]
Sur le sous-espace vivant, sa pseudo-inverse vaut \(1/12\) ; avec l’orientation
négative de bord,
\[
E_{\rm surv}=-K_{60,81}^{+}
\]
produit \(-1/12\).

\subsection{Extraction d’échelle par 2e différence triadique}

La réalisation par un extracteur d’échelle triadique est la construction par
seconde différence développée intégralement dans
\cref{sec:c27-segunda-diferencia-triadica}. Cette première apparition fixe sa
fonction dans l’inventaire des réalisations ; la résidence indiquée conserve
les définitions de \(T_L\), \(a(L)\), \(A_1(L)\), \(C(L)\), l’indépendance
du résidu à l’égard de \(L\) et le statut exact de la normalisation.

\subsection{Réalisation modulaire par le quotient \(\eta(\tau)/\eta(3\tau)\) et la fonction \(t_3\)}

La fonction êta satisfait
\[
\frac{\eta(\tau)}{\eta(3\tau)}
=q^{-1/12}\prod_{n\ge1}(1+q^n+q^{2n})^{-1}.
\]
L’exposant provient de \(1/24-3/24=-1/12\). En prenant douze copies :
\[
t_3(q)=\left(\frac{\eta(\tau)}{\eta(3\tau)}\right)^{12}
=q^{-1}-12+54q-76q^2-243q^3+\cdots.
\]
Le coefficient \(54\) coïncide avec le défaut APP. Le corridor
\[
j=\frac{(t_3+27)(t_3+243)^3}{t_3^3}
\]
se relie ensuite à la branche modulaire classique.

\begin{remark}[Ce qui a été démontré dans le corridor]
Le défaut APP \(54\), l’exposant êta \(-1/12\), le défaut
hexade--octade et les identités modulaires une fois \(t_3\) fixé sont exacts.
L’affirmation forte est qu’ils soient tous des images naturelles d’un unique opérateur
TPK. Cette naturalité doit être démontrée au moyen du diagramme du NTHW ; elle ne se
déduit pas de la seule répétition des mêmes entiers.
\end{remark}


\section{5 réalisations algébriques de \texorpdfstring{\(-1/12\)}{-1/12} et séparation de leurs domaines}

Outre le défaut normalisé de l’inclusion hexade--octade, apparaissent
cinq réalisations opératorielles du même rationnel. Leurs domaines sont
précisés séparément ; l’égalité des valeurs n’identifie pas les
opérateurs sans un entrelaceur supplémentaire.

\subsection{Différence de périodicités dodécaphasiques}

Une dent de \(\Cdoce\) mesure \(30^\circ\). Alors
\[
\frac{30}{120}-\frac{30}{90}
=\frac14-\frac13
=-\frac1{12}.
\]
C’est la différence orientée entre les périodes de \(90^\circ\) et
\(120^\circ\) dans cette réalisation.

\subsection{Pseudo-inverse de l’opérateur de Gram sur la composante vivante}

Dans le saut exact \(60\to81\), deux fibres ont pour tailles \(0\) et
\(12\). L’opérateur de Gram est
\[
K_{60,81}=\operatorname{diag}(0,12).
\]
Sur la composante non nulle de son image, la pseudo-inverse vaut \(1/12\) ; avec l’orientation
négative de bord,
\[
E_{\rm surv}=-K_{60,81}^{+}
\]
produit \(-1/12\).

\subsection{2e différence entre échelles triadiques}
\label{sec:c27-segunda-diferencia-triadica}

Définissons
\[
T_L=\sum_{n=1}^{L-1}n(L-n)=\frac{L(L^2-1)}6,
\qquad
a(L)=-\frac{T_L}{L^2}=-\frac L6+\frac1{6L}.
\]
La première soustraction d’échelle est
\[
A_1(L)=a(L)-\frac{a(3L)}3=\frac4{27L},
\]
et la seconde
\[
C(L)=LA_1(L)-\frac{3L}{3}A_1(3L)=\frac8{81}.
\]
Le résidu \(8/81\) est indépendant de \(L\). Pour obtenir
\(-1/12\), on applique la normalisation déclarée
\[
R(L)=-\frac{27}{32}C(L)=-\frac1{12}.
\]
Le calcul dérive \(8/81\) ; il ne dérive pas à lui seul le facteur \(-27/32\).

\subsection{Chaîne modulaire}

La fonction êta satisfait
\[
\frac{\eta(\tau)}{\eta(3\tau)}
=q^{-1/12}\prod_{n\ge1}(1+q^n+q^{2n})^{-1}.
\]
L’exposant provient de \(1/24-3/24=-1/12\). En prenant douze copies :
\[
t_3(q)=\left(\frac{\eta(\tau)}{\eta(3\tau)}\right)^{12}
=q^{-1}-12+54q-76q^2-243q^3+\cdots.
\]
Le coefficient \(54\) coïncide avec le défaut APP\@. L’identité rationnelle
\[
j=\frac{(t_3+27)(t_3+243)^3}{t_3^3}
\]
établit la connexion avec la branche modulaire classique.

\Needspace{32\baselineskip}
\subsection{Défaut des projecteurs de différence cyclique}

Soit \(S\) le décalage cyclique sur \(\QQ^{12}\) et définissons
\[
 D_r=I-S^r.
\]
Le noyau de \(D_r\) est formé des vecteurs constants sur chaque orbite
de \(S^r\) ; par conséquent,
\[
 \dim\ker D_r=\gcd(12,r).
\]
Soient \(\Pi_{\ker D_3}\) et \(\Pi_{\ker D_4}\) les projecteurs rationnels sur
\(\ker D_3\) et \(\ker D_4\), et soit
\(\tau(T)=\operatorname{Tr}(T)/12\) la trace normalisée. Alors
\[
 \boxed{
 \tau(\Pi_{\ker D_3}-\Pi_{\ker D_4})
 =\frac{\operatorname{rank}\Pi_{\ker D_3}
       -\operatorname{rank}\Pi_{\ker D_4}}{12}
 =\frac{3-4}{12}
 =-\frac1{12}.}
\]
Cette réalisation exprime le rationnel comme défaut transversal des pas
trois et quatre dans le cycle à douze phases.

\begin{remark}[Domaines des réalisations]
Le défaut APP \(54\), l’exposant êta \(-1/12\), le défaut
combinatoire du relèvement hexade--octade et les
identités modulaires une fois \(t_3\) fixé sont exacts. La valeur
de la pseudo-inverse, l’extracteur normalisé et le défaut des projecteurs sont également exacts.
Chaque égalité conserve son domaine et son application correspondante.
\end{remark}
  \endgroup

\chapter{Caractères quadratiques et constantes spectrales : Catalan, Apéry, Euler--Mascheroni et Stieltjes}
\begingroup
\renewcommand{\chapter}[2][]{}%
\chapter{Caractères quadratiques et constantes spectrales : Catalan, Apéry et tour d’Euler--Stieltjes}
\label{chap:p3-final:41:catalan_apery_euler}
\label{ch:spectral}

\section{Représentation opératorielle des caractères de Dirichlet}

Les constantes de ce chapitre sont regroupées parce qu'elles partagent une construction
spectrale indépendante de leurs développements décimaux et exempte de valeurs
cibles. Soit

\[
Ne_n=ne_n,\qquad n\in\N,
\]

dans \(\ell^2(\N)\). Un caractère résiduel \(\chi\) définit l'opérateur diagonal \(Q_\chi e_n=\chi(n)e_n\) ; pourvu que la puissance soit de classe trace,

\begin{equation}
\Tr(N^{-s}Q_\chi)=L(s,\chi).
\label{eq:L-trace}
\end{equation}

L'équation conserve le module, l'orientation et la provenance de chaque
caractère ; par conséquent, elle n'identifie pas des objets résiduels distincts.

\section{Construction HMT des caractères résiduels}

L'opérateur diagonal précédent ne constitue pas la donnée algébrique initiale. La
chaîne de sélection est

\[
\begin{aligned}
\APP&\longrightarrow\TRIT\longrightarrow\TPK
\longrightarrow\{\text{résidu modulo }3,\ C_P\}\\
&\longrightarrow\{\chi_{-3},\ J=C_P\bmod3\}
\longrightarrow\{\chi_{-3},\ J^2=-I,\chi_{-4}\}\\
&\longrightarrow Q_3,Q_4.
\end{aligned}
\]

La fonctionnelle résiduelle triadique détermine la différence entre les deux classes
orientées non nulles. Dans le secteur d'incidence, on emploie la matrice de conférence
symétrique de Paley d'ordre six,
\[
C_P^{\mathsf T}=C_P,\qquad C_PC_P^{\mathsf T}=5I_6.
\]
Sa réduction \(J=C_P\bmod3\) satisfait
\[
J^2=2I_6=-I_6
\quad\text{dans }\operatorname{End}_{\mathbb F_3}(\mathbb F_3^6).
\]
Le quart de torsion détermine alors la phase cyclique
\(1,\ii,-1,-\ii\). L'application du théorème chinois des restes conserve
les deux données résiduelles et produit \(Q_{12}=Q_3Q_4\). La trace
\(\Tr(N^{-s}Q)\) est évaluée à l'étape terminale ; ni la série de Dirichlet
ni son expression décimale n'interviennent dans la sélection du caractère.

La construction est explicitée par les tableaux
\begin{equation}
\chi_{-3}(n)=
\begin{cases}
0,&3\mid n,\\
1,&n\equiv1\pmod3,\\
-1,&n\equiv2\pmod3,
\end{cases}
\qquad
\chi_{-4}(n)=
\begin{cases}
0,&2\mid n,\\
1,&n\equiv1\pmod4,\\
-1,&n\equiv3\pmod4.
\end{cases}
\label{eq:characters34-tables}
\end{equation}
En \(\ell^2(\N)\),
\[
Q_3e_n=\chi_{-3}(n)e_n,
\qquad Q_4e_n=\chi_{-4}(n)e_n.
\]
Comme ils sont diagonaux, ils commutent ; le théorème chinois des restes donne
\begin{equation}
(Q_3Q_4)e_n=\chi_{-3}(n)\chi_{-4}(n)e_n
=\chi_{12}(n)e_n.
\label{eq:q12-product-proof}
\end{equation}
Le produit conserve simultanément le résidu triadique et le quart de
torsion. Par conséquent, le caractère modulo douze est déterminé
avant d'effectuer toute évaluation de ses valeurs \(L\).

\begin{remark}[Contrôle négatif]
La chaîne est invalidée si \(Q_3\) ne distingue pas les deux classes orientées,
si \(J^2\ne-I\), si \(Q_3Q_4\) n'a pas le tableau CRT déclaré ou si
l'opérateur est sélectionné à partir de la valeur de la constante.
\end{remark}

\section{Structure complexe interne et constante de Catalan}

La structure complexe interne satisfait

\begin{equation}
J^2=-I.
\label{eq:J-square}
\end{equation}

L'égalité \(J^2=-I\) définit une représentation
\(\rho:C_4\to\operatorname{Aut}_{\mathbb F_3}(\mathbb F_3^6)\),
\(\rho(j)=J\). Pour la réaliser spectralement, on prend le caractère unitaire
\(\upsilon:C_4\to U(1)\), \(\upsilon(j)=\ii\), et l'application de phase
\(n\mapsto j^n\). Sur \(\ell^2(\N)\), la représentation diagonale induite
est
\[
U_4e_n=\upsilon(j^n)e_n=\ii^ne_n.
\]
Sa partie orientée est
\[
Q_4=\frac{U_4-U_4^*}{2\ii},\qquad
Q_4e_n=\sin\!\left(\frac{\pi n}{2}\right)e_n
=\chi_{-4}(n)e_n.
\]
C'est le relèvement explicite entre le générateur fini d'ordre quatre et l'opérateur
résiduel ; on n'identifie pas \(\mathbb F_9\) à \(\mathbb C\).

Sur les entiers impairs, le même caractère peut s'écrire

\[
\chi_J(n)=\ii^{n-1}=\chi_{-4}(n).
\]

De manière équivalente,

\begin{equation}
Q_4=\frac{U_4-U_4^*}{2\ii},
\qquad Q_4e_n=\chi_{-4}(n)e_n.
\label{eq:Q4}
\end{equation}

La constante de Catalan apparaît comme

\begin{equation}
G=\Tr(N^{-2}Q_4)=L(2,\chi_{-4})
=\sum_{n\ge0}\frac{(-1)^n}{(2n+1)^2}.
\label{eq:catalan}
\end{equation}

Le relèvement unitaire du quart de torsion engendre exactement le
caractère dont le moment quadratique est la constante de Catalan.

\begin{proof}[Preuve opératorielle]
L'opérateur \(N^{-2}Q_4\) est de classe trace parce que
\(\sum_{n\ge1}|\chi_{-4}(n)|n^{-2}<\infty\). Sa trace est donc la
somme des éléments diagonaux. Le tableau \cref{eq:characters34-tables}
annule les pairs et alterne les impairs :
\[
\Tr(N^{-2}Q_4)
=1-\frac1{3^2}+\frac1{5^2}-\frac1{7^2}+\cdots.
\]
L'identité avec \(L(2,\chi_{-4})\) est la définition de la série de
Dirichlet ; l'égalité avec \(\operatorname{Im}\operatorname{Li}_2(\ii)\)
s'obtient en prenant la partie imaginaire de
\(\sum_{n\ge1}\ii^n/n^2\).
\end{proof}

La valeur terminale est
\[
G=0.9159655941772190150546035149\ldots.
\]
La structure déterminante est l'identité \(J^2=-I\), qui induit
exactement le signe alterné des classes impaires ; le développement décimal
est son évaluation terminale.

\section{Composition résiduelle triadique--quaternaire}

La fonctionnelle résiduelle modulo trois fournit \(\chi_{-3}\). Son produit
CRT avec \(\chi_{-4}\) définit

\begin{equation}
\chi_{12}=\chi_{-3}\chi_{-4}.
\label{eq:chi12}
\end{equation}

Sur les unités modulo douze,

\[
\chi_{12}(1)=1,\quad
\chi_{12}(5)=-1,\quad
\chi_{12}(7)=-1,\quad
\chi_{12}(11)=1,
\]

et s'annule en dehors de celles-ci. Deux évaluations exactes sont

\begin{align}
L(1,\chi_{12})&=\frac{\log(2+\sqrt3)}{\sqrt3},
\label{eq:L1chi12}\\
L(2,\chi_{12})&=\frac{\pi^2}{6\sqrt3}.
\label{eq:L2chi12}
\end{align}

La première contient l'unité hyperbolique

\begin{equation}
\sqrt3L(1,\chi_{12})=\arcosh2=\log(2+\sqrt3).
\label{eq:pell2}
\end{equation}

Avec \(\log(3+2\sqrt2)=2\log(1+\sqrt2)\) et
\(\log(30+\sqrt{899})\), cette identité détermine une famille d'échelles
de Pell dont les membres conservent leurs provenances algébriques respectives.

Les évaluations terminales certifiées sont

\[
L(1,\chi_{12})=0.7603459963009463475310942549\ldots,
\]

\[
L(2,\chi_{12})=0.9497031262940093952634984917\ldots.
\]

\begin{proof}[Évaluations dodécaphasiques]
En regroupant par classes résiduelles,
\begin{align*}
L(1,\chi_{12})
&=\sum_{n\ge0}\left(
\frac1{12n+1}-\frac1{12n+5}
-\frac1{12n+7}+\frac1{12n+11}\right)\\
&=\int_0^1
\frac{1-x^4-x^6+x^{10}}{1-x^{12}}\,dx
=\int_0^1\frac{1-x^4}{1+x^6}\,dx.
\end{align*}
Une décomposition réelle dans les facteurs quadratiques de \(1+x^6\)
produit
\(\log(2+\sqrt3)/\sqrt3\). Pour le deuxième moment, si
\(\psi_1\) est la fonction trigamma,
\begin{align*}
L(2,\chi_{12})
&=\frac1{144}\bigl[
\psi_1(1/12)-\psi_1(5/12)
-\psi_1(7/12)+\psi_1(11/12)\bigr]\\
&=\frac{\pi^2}{144}\left(
\csc^2\frac\pi{12}-\csc^2\frac{5\pi}{12}\right)
=\frac{\pi^2}{6\sqrt3},
\end{align*}
où l'on a utilisé
\(\psi_1(z)+\psi_1(1-z)=\pi^2\csc^2(\pi z)\). Les deux preuves emploient
le tableau CRT, et non une valeur décimale cible.
\end{proof}

\section{Constantes d'Euler et de Stieltjes et transformée thêta--Mellin}

Soit
\[
\vartheta(t)=\sum_{n\in\Z}e^{-\pi n^2t},\qquad t>0.
\]
Pour \(\Re s>1\), la fonctionnelle thêta--Mellin retrouve l'identité convergente

\[
\pi^{-s/2}\Gamma\!\left(\frac{s}{2}\right)\zeta(s)
=\frac12\int_0^\infty
\bigl(\vartheta(t)-1\bigr)t^{s/2-1}\,dt.
\]

Le prolongement méromorphe s'obtient en scindant l'intégrale en \(t=1\), en utilisant
\(\vartheta(t)=t^{-1/2}\vartheta(1/t)\) et en séparant les termes polaires. Son développement en \(s=1\),

\begin{equation}
\zeta(s)=\frac1{s-1}+\sum_{j\ge0}
\frac{(-1)^j\gamma_j}{j!}(s-1)^j,
\label{eq:stieltjes}
\end{equation}

situe \(\gamma=\gamma_0\) et les constantes de Stieltjes dans une même
famille de moments renormalisés. Cette organisation exprime une
provenance analytique commune, sans attribuer une interprétation physique
identique à tous ses niveaux.

 La constante de Catalan apparaît habituellement comme une série alternée :\par

\[
G
=
1-\frac1{3^2}+\frac1{5^2}-\frac1{7^2}+\cdots.
\]

Cette formule donne sa valeur et la généalogie HMT explique pourquoi cette alternance particulière existe.\par

Dans HMT, l’origine se trouve dans le quart de torsion.\par

Il existe un opérateur interne \(J\) tel que\par

\[
J^2=-I.
\]

La racine de \(-1\) émerge de deux applications successives d’un quart de tour. La structure complexe est la mémoire algébrique d’une rotation interne.\par

Cette rotation engendre un caractère d’ordre quatre. Lorsqu’il est publié sur les entiers positifs, il produit exactement le motif résiduel qui alterne entre \(+1\), \(0\), \(-1\), \(0\).\par

La constante de Catalan apparaît alors comme une trace spectrale :\par

\[
G
=
L(2,\chi_{-4})
=
\operatorname{Tr}(N^{-2}Q_4).
\]

Cela change sa signification. Catalan est simultanément la somme d’une série particulière et la réponse spectrale des entiers à un quart de torsion.\par

Chaque terme de la série enregistre la manière dont la rotation interne agit sur une classe résiduelle. L’alternance est l’ombre arithmétique de la rotation.\par

La structure triadique possède son propre caractère modulo \(3\). En le combinant au caractère d’ordre quatre, le théorème chinois des restes produit un caractère modulo \(12\) :\par

\[
\chi_{12}
=
\chi_{-3}\chi_{-4}.
\]

Le module \(12\) apparaît comme l’espace commun minimal dans lequel l’orientation triadique et le quart de torsion peuvent coexister.\par

La dodécaphase est, en ce sens, une réconciliation de \(3\) et de \(4\).\par

Du caractère combiné émergent\par

\[
L(1,\chi_{12})
=
\frac{\log(2+\sqrt3)}{\sqrt3},
\]

y\par

\[
L(2,\chi_{12})
=
\frac{\pi^2}{6\sqrt3}.
\]

La cantidad\par

\[
2+\sqrt3
\]

est l’unité fondamentale du corps réel \(\mathbb Q(\sqrt3)\). Son logarithme représente un déplacement hyperbolique. De nouveau, une structure finie de résidus ouvre une dynamique multiplicative infinie.\par

Ici se rencontrent rotation, torsion, caractère arithmétique, unité de Pell et déplacement hyperbolique.\par

La constante d’Euler–Kronecker du corps \(\mathbb Q(\sqrt3)\) prolonge cette structure. Elle mesure le terme régulier qui demeure après la séparation du pôle principal de la fonction zêta du corps. Dans la généalogie HMT, ce terme constitue l’invariant régulier de mémoire arithmétique fine qui subsiste après le retrait de la divergence universelle.\par

La famille Bendersky–Glaisher prolonge le mécanisme en une tour. Les dérivées spectrales de zêta engendrent cette tour. Apéry et d’autres constantes associées à des valeurs impaires s’organisent comme membres d’une même famille fonctionnelle.\par

L’apport nouveau réside dans la généalogie qui réunit les constantes classiques :\par

\begin{itemize}[leftmargin=*,itemsep=0.35em]
\item le quart de torsion produit la structure complexe ;
\item le résidu triadique produit la structure modulo \(3\) ;
\item leur composition produit le caractère modulo \(12\) ;
\item le caractère engendre des valeurs \(L\) ;
\item les valeurs \(L\) ouvrent le corps quadratique et ses constantes spectrales.
\end{itemize}

La torsion géométrique est devenue arithmétique en conservant intégralement son orientation.\par
  \endgroup

\chapter{Euler--Kronecker, Meissel--Mertens et distributions premières dans la double lecture des feuillets}
\begingroup
\renewcommand{\chapter}[2][]{}%
\chapter{Constantes d’Euler--Kronecker et distributions de nombres premiers : caractères dodécaphasiques, Meissel--Mertens et séries de nombres premiers jumeaux}
\label{chap:p3-final:42:euler_kronecker_primos}
\section{Constante d'Euler--Kronecker du caractère dodécaphasique}

La factorisation de Dedekind

\begin{equation}
\zeta_{\mathbb Q(\sqrt3)}(s)=\zeta(s)L(s,\chi_{12})
\label{eq:dedekind-sqrt3}
\end{equation}

donne, en comparant les termes constants en \(s=1\),

\begin{equation}
\gamma_{\mathbb Q(\sqrt3)}
=\gamma+\frac{L'(1,\chi_{12})}{L(1,\chi_{12})}
=1.0539656082484825800\ldots.
\label{eq:euler-kronecker}
\end{equation}

Cette constante représente le terme fini normalisé au pôle de la
fonction zêta du corps réel dodécaphasique. Elle ne constitue pas une variante de
\(\gamma\) d'Euler, mais la correction arithmétique induite par
\(\chi_{12}\).

En effet, en écrivant \(u=s-1\),
\[
\zeta(s)=u^{-1}+\gamma+O(u),
\qquad
L(s,\chi_{12})=L_1+L_1'u+O(u^2),
\]
on obtient
\begin{equation}
\zeta_{\mathbb Q(\sqrt3)}(s)
=\frac{L_1}{u}+\bigl(\gamma L_1+L_1'\bigr)+O(u).
\label{eq:ek-laurent-derivation}
\end{equation}
La constante d'Euler--Kronecker est le quotient du terme fini par
le résidu ; diviser \cref{eq:ek-laurent-derivation} par \(L_1\) prouve
\cref{eq:euler-kronecker}. Cette dérivation explique simultanément le
résidu
\[
\operatorname*{Res}_{s=1}\zeta_{\mathbb Q(\sqrt3)}(s)
=\frac{\log(2+\sqrt3)}{\sqrt3}
\]
et la correction finie. La valeur numérique est certifiée en évaluant
\(L(1,\chi_{12})\) et \(L'(1,\chi_{12})\) avec le même tableau résiduel et une
borne de queue ; elle n'est pas importée d'un tableau de corps quadratiques.

\section{Constantes arithmétiques associées aux nombres premiers}

Une fois les cycles arithmétiques irréductibles sélectionnés, on obtient deux
familles distinctes :

\begin{align}
B_1&=\lim_{x\to\infty}\left(
\sum_{p\le x}\frac1p-\log\log x
\right),
\label{eq:meissel-mertens}\\
C_2&=\prod_{p>2}\frac{p(p-2)}{(p-1)^2}.
\label{eq:twin-prime-constant}
\end{align}

\(B_1\) régularise une somme locale ; \(C_2\) est une densité multiplicative de paires jumelles. Les deux exigent un contrôle de queue. Ils ne se déduisent pas l'un de l'autre parce qu'ils partagent les nombres premiers.

  \endgroup

\chapter{Produits réguliers et système de Bendersky--Glaisher--Kinkelin}
\begingroup
\renewcommand{\chapter}[2][]{}%
\chapter{Famille Bendersky--Glaisher--Kinkelin : produits réguliers, dérivées de la fonction zêta et hiérarchies de normalisation}
\label{chap:p3-final:43:bendersky_glaisher}
\section{Génération de la famille Bendersky--Glaisher--Kinkelin par régularisation zêta}

Soit \(\mathcal Z_3(s):=\zeta(s)\) la réalisation scalaire de la fonctionnelle
zêta triadique dans la normalisation de ce volume,
\(H_k=\sum_{j=1}^k j^{-1}\) (avec \(H_0=0\)) et \(B_{k+1}\) le nombre de
Bernoulli de convention \(B_1=-1/2\). Nous définissons

\begin{equation}
A_k=\exp\!\left(
\frac{H_kB_{k+1}}{k+1}-\mathcal Z_3'(-k)
\right),\qquad k\ge0.
\label{eq:bendersky}
\end{equation}

Les premiers niveaux incluent

\[
A_0=\sqrt{2\pi},
\]

et la constante de Glaisher--Kinkelin en \(A_1\). Pour \(n\ge1\),

\begin{equation}
\zeta(2n+1)=(-1)^{n+1}
\frac{2(2\pi)^{2n}}{(2n)!}\log A_{2n}.
\label{eq:odd-zeta-bendersky}
\end{equation}

La constante d'Apéry, \(\zeta(3)\), est le premier niveau pair non trivial.
Sa présence dans le gaz photonique du \cref{ch:metrology} constitue une
réalisation physique ultérieure de cette valeur spectrale.

La première partie de la tour peut se voir de manière synoptique :
\begin{center}
\small
\begin{tabular}{@{}cll@{}}
\toprule
\(k\) & identidad estructural & valeur terminale\\
\midrule
0 & \(A_0=\exp[-\zeta'(0)]=\sqrt{2\pi}\)
  & \(2.5066282746310005024\ldots\)\\
1 & \(\log A_1=1/12-\zeta'(-1)\)
  & \(1.2824271291006226369\ldots\)\\
2 & \(\log A_2=-\zeta'(-2)=\zeta(3)/(4\pi^2)\)
  & \(1.0309167521973921\ldots\)\\
3 & \(\log A_3=-11/720-\zeta'(-3)\)
  & \(0.9795555269428446\ldots\)\\
4 & \(\log A_4=-\zeta'(-4)\)
  & \(0.9920479745250403\ldots\)\\
\bottomrule
\end{tabular}
\end{center}

\begin{proof}[Relation avec les valeurs impaires]
L'équation fonctionnelle de \(\zeta\), développée en \(s=-2n\), donne
\begin{equation}
\zeta'(-2n)=(-1)^n
\frac{(2n)!}{2(2\pi)^{2n}}\zeta(2n+1).
\label{eq:zeta-derivative-even-negative}
\end{equation}
Comme \(B_{2n+1}=0\) pour \(n\ge1\), la définition
\cref{eq:bendersky} réduit \(\log A_{2n}=-\zeta'(-2n)\). Isoler
\(\zeta(2n+1)\) prouve \cref{eq:odd-zeta-bendersky}. L'indice \(2n\)
conserve le niveau de dérivation : ce n'est pas une étiquette éditoriale.
\end{proof}

  \endgroup

\chapter{Dynamique de Gauss et constantes de Khinchin, Lévy et Lochs}
\begingroup
\renewcommand{\chapter}[2][]{}%
\chapter{Dynamique de Gauss et constantes de Khinchin, Lévy et Lochs : cofinalité des cylindres, opérateur de transfert et spectre de Gauss--Kuzmin--Wirsing}
\label{chap:p3-final:44:gauss_khinchin_levy_lochs}
\label{ch:gauss}

\section{Cofinalité entre cylindres HMT et cylindres de fractions continues}

Une section HMT irrationnelle détermine une famille emboîtée de cylindres compacts. Son évaluation archimédienne est un point \(x\in(0,1)\setminus\mathbb Q\). La fraction continue de \(x\) détermine des cylindres classiques

\[
I(a_1,\ldots,a_n).
\]

\begin{proposition}[Cofinalité]
Pour tout cylindre fini \(I(a_1,\ldots,a_n)\) qui contient l'évaluation \(x\), il existe une profondeur HMT \(m(n)\) telle que le cylindre HMT de niveau \(m\ge m(n)\) soit contenu dans \(I(a_1,\ldots,a_n)\).
\end{proposition}

\begin{proof}
\(n\) étant fixé, le cylindre de fraction continue \(I(a_1,\ldots,a_n)\) est un voisinage ouvert de l'irrationnel \(x\). Il existe alors \(\varepsilon_n>0\) avec \((x-\varepsilon_n,x+\varepsilon_n)\subset I(a_1,\ldots,a_n)\). Les cylindres HMT \(C_m\) contiennent \(x\), sont emboîtés et satisfont \(\operatorname{diam}C_m\to0\). Choisissons \(m(n)\) avec \(\operatorname{diam}C_{m(n)}<\varepsilon_n\). Comme \(x\in C_{m(n)}\), tout point de ce cylindre est à une distance de \(x\) inférieure à \(\varepsilon_n\), donc \(C_{m(n)}\subset I(a_1,\ldots,a_n)\). L'emboîtement donne l'inclusion pour tout niveau ultérieur.
\end{proof}

La démonstration utilise la compacité, l'emboîtement et l'unicité de
l'évaluation. Le résultat établit le transport des approximations, mais
n'implique pas la conjugaison du décalage de mémoire TPK avec l'application de
Gauss sur toutes les fibres profondes.

La généalogie complète de cette réalisation est
\[
\APP\to\TRIT\to\TPK
\to (C_m,\rho_{m+1,m})_{m\ge0}
\to x=\bigcap_m C_m
\to \bigl(I(a_1,\ldots,a_n)\bigr)_{n\ge1}.
\]
Les cylindres HMT conservent la route, la phase et la retenue ; l'application terminale projette
ces coordonnées sur les préfixes de fraction continue. La cofinalité
démontre que cette projection converge vers le même point. Un entrelaceur
dynamique exigerait en outre un relèvement \(\widetilde T_G\) avec des
carrés commutatifs sur les fibres de valuation, propriété qui ne se
déduit pas de la cofinalité.

\section{Mesure de Gauss et observables de Khinchin et de Lévy}

L'application de Gauss \cite{Khinchin,IosifescuKraaikamp}

\[
T_G(x)=\frac1x-\left\lfloor\frac1x\right\rfloor
\]

préserve

\begin{equation}
d\mu_G(x)=\frac{dx}{\log2(1+x)}.
\label{eq:gauss-measure}
\end{equation}

La constante de Khinchin est définie par

\begin{equation}
\log K=\int_0^1\log a_1(x)\,d\mu_G(x).
\label{eq:khinchin}
\end{equation}

La distribution du premier chiffre se calcule directement :
\begin{equation}
\mu_G(a_1=k)
=\frac1{\log2}\int_{1/(k+1)}^{1/k}\frac{dx}{1+x}
=\log_2\!\left(\frac{(k+1)^2}{k(k+2)}\right).
\label{eq:gauss-digit-law}
\end{equation}
Par le théorème ergodique,
\begin{equation}
K=\prod_{k\ge1}
k^{\,\log_2((k+1)^2/[k(k+2)])}
=2.685452001065306445309714835\ldots
\label{eq:khinchin-product}
\end{equation}
pour \(\mu_G\)-presque toute section. Le produit ne s'applique pas à chaque point :
la section \(\varphi^{-1}\), par exemple, a tous ses chiffres égaux
à un et constitue un critère de réfutation immédiat d'une lecture point par point. Une
certification de \cref{eq:khinchin-product} tronque la somme des logarithmes et
borne la queue au moyen de l'estimation
\[
\mu_G(a_1=k)=O(k^{-2}),
\qquad
\sum_{k>N}\mu_G(a_1=k)\log k
=O\!\left(\frac{\log N+1}{N}\right).
\]

La constante de Lévy est

\begin{equation}
L=-\int_0^1\log x\,d\mu_G(x)
=\frac{\pi^2}{12\log2}.
\label{eq:levy}
\end{equation}

En effet,

\[
\frac1{1+x}=\sum_{n\ge0}(-1)^nx^n,\qquad
\int_0^1(-\log x)x^n\,dx=\frac1{(n+1)^2},
\]

de sorte que l'intégrale est \(\eta(2)/\log2=\pi^2/(12\log2)\).

La constante de Khinchin détermine la moyenne logarithmique des quotients
partiels, tandis que la constante de Lévy détermine le taux exponentiel
de croissance des dénominateurs. Ce sont deux observables distinctes d'un
même système mesuré.

\section{Entropie métrique, exposant de Lyapunov et constante de Lochs}

Comme

\[
\log|T_G'(x)|=-2\log x,
\]

l'exposant de Lyapunov et l'entropie métrique sont

\begin{equation}
\chi_G=h_G=2L=\frac{\pi^2}{6\log2}.
\label{eq:gauss-entropy}
\end{equation}

Le rapport de Lochs entre les fractions continues et un développement en base \(b\) est

\begin{equation}
\Lambda_b=\frac{\log b}{h_G}
=\frac{6\log2\log b}{\pi^2}.
\label{eq:lochs}
\end{equation}

En particulier,

\begin{align*}
\Lambda_{10}&=0.9702701143920339257402560192\ldots,\\
\Lambda_{729}&=\frac{36\log2\log3}{\pi^2}
=2.777618966374305777705782976\ldots.
\end{align*}

La base \(729=3^6\) est pertinente pour l'espace ambiant \(\mathbb F_3^6\), mais \cref{eq:lochs} reste un théorème presque sûr de \(\mu_G\) ; et non une égalité point par point pour toute route HMT.

La valeur \(\Lambda_{729}\) représente le taux asymptotique typique de
quotients partiels obtenus par bloc complet de six trits. Elle n'identifie pas
l'espace ambiant de \(729\) mots aux autres objets de cardinal \(729\) du
corpus. L'égalité est un taux d'information entre partitions, et non une
bijection entre états.

\section{Opérateur de transfert de Gauss--Kuzmin--Wirsing}

L'opérateur de Gauss--Kuzmin--Wirsing agit par

\begin{equation}
(\mathcal L f)(x)=\sum_{n\ge1}\frac1{(x+n)^2}
f\!\left(\frac1{x+n}\right).
\label{eq:gkw-operator}
\end{equation}

Sa valeur propre dominante correspond à la densité invariante. La constante GKW est le module de la valeur propre sous-dominante sur le sous-espace de moyenne nulle. Un schéma HMT compatible utilise des partitions de \(9^m\) cylindres, des matrices de transfert avec intervalles et une borne de queue.

Pour promouvoir le calcul, il faut certifier : l'isolement de la valeur propre, l'invariance du sous-espace, l'erreur de troncature, la convergence des projecteurs spectraux et l'absence d'une autre valeur propre de même module. La valeur décimale connue ne peut pas être introduite comme centre de l'intervalle initial.

Une déformation conjointe contient Khinchin et Lévy dans un seul objet :
\begin{equation}
(\mathcal L_{s,t}f)(x)
=\sum_{n\ge1}n^t(n+x)^{-2s}
 f\!\left(\frac1{n+x}\right),
\qquad
P(s,t)=\log r(\mathcal L_{s,t}).
\label{eq:gauss-pressure-operator}
\end{equation}
Au point \((s,t)=(1,0)\), la perturbation spectrale donne
\begin{equation}
\partial_tP(1,0)=\log K,
\qquad
-\partial_sP(1,0)=\chi_G=2L.
\label{eq:pressure-derivatives}
\end{equation}
La hessienne de \(P\) détermine les variances et covariances asymptotiques de
\(\log a_1\) et \(-2\log x\). Par conséquent, les constantes de Khinchin et de
Lévy sont deux dérivées transversales de la même fonction de pression.

La reconnaissance numérique de la valeur propre sous-dominante est
\[
\lambda_2\simeq-0.3036630028987326586,\qquad
\kappa_{\rm GKW}=|\lambda_2|,\qquad
-\log\kappa_{\rm GKW}\simeq1.19183673556055.
\]
Dans ce volume, ces décimales constituent des valeurs de référence pour la
comparaison, et non un certificat spectral.
La preuve HMT proposée utilise Galerkin ou Ulam sur les cylindres cofinaux de
profondeur \(9^m\), une inégalité de Lasota--Yorke, l'arithmétique
d'intervalles pour isoler \(\lambda_2\) et une borne rigoureuse de la queue. Un
intervalle contenant deux valeurs propres de même module réfute la
promotion, sans pour autant affecter Khinchin, Lévy ou Lochs.

 \section{Certificat cofinal du spectre de Gauss--Kuzmin--Wirsing sur la hiérarchie nonadique}
\label{sec:p3-final:cierre-gkw}

Pour Gauss--Kuzmin--Wirsing, soit \(P_N\) la famille de troncatures
certifiée sur l’espace fonctionnel déclaré, et définissons
\[
 \Pi_m=P_{9^m}.
\]
La hiérarchie \((\Pi_m)\) est une sous-suite cofinale ; les bornes uniformes, la convergence des projecteurs et l’absence de pollution spectrale sont héritées du système \((P_N)\). Une discrétisation cylindrique différente devra fournir son entrelaceur avec \((P_{9^m})\).
\endgroup

\chapter{Criticité unimodale, constantes de Feigenbaum et renormalisation hyperbolique}
\begingroup
\renewcommand{\chapter}[2][]{}%
\chapter{Criticité unimodale de la famille dodécaphasique : profil analytique exact, approximants de Feigenbaum et réduction de l’universalité au certificat hyperbolique}
\label{chap:p3-final:45:feigenbaum}
\label{mab:rm:ch:chaos}

\section{Construction intrinsèque de la famille dodécaphasique}

Les constantes de Feigenbaum caractérisent l’universalité du
doublement de période. Leur dérivation dans HMT exige de construire au préalable
une famille analytique unimodale, de démontrer sa criticité quadratique et
d’étudier ensuite l’action de l’opérateur de renormalisation.

L’application initiale ne procède pas d’un choix harmonique externe. Le retour
de \(\TPK\) détermine la roue dodécaphasique orientée
\(U_{12}^{\rm sign}\). Dans le secteur uniforme \(\eta=0\), la trace
normalisée de ses douze secteurs fixe \(w_m=1/12\), et la fonctionnelle de phase
\begin{equation}
\mathcal P_{\rm crit}:U_{12}^{\rm sign}\longrightarrow\R/2\pi\Z,
\qquad
m\longmapsto \frac{(30m-10)\pi}{180}
=\pi\frac{3m-1}{18}
\label{mab:rm:eq:critical-phase-reader}
\end{equation}
conserve l’ordre et l’orientation. La composition causale est
\[
\APP\to\TRIT\to\TPK\to U_{12}^{\rm sign}
\xrightarrow{\mathcal P_{\rm crit}}
\{(\phi_m,w_m)\}_{m=1}^{12}
\to u_0\to f_\lambda.
\]
Toute modification de \(\mathcal P_{\rm crit}\), des poids ou de l’ordre des
phases altère la famille et constitue un critère de réfutation. Ces
choix doivent être fixés avant la confrontation avec \(\delta_F\).

La famille dodécaphasique résultante part de

\[
a_m=\frac{3m-1}{18},\qquad
\phi_m=\pi a_m,\qquad m=1,\ldots,12,
\]

et du deuxième moment

\[
D_0=\frac1{12}\sum_{m=1}^{12}a_m^2.
\]

L’identité

\begin{equation}
\sum_{m=1}^{12}(3m-1)^2=5394=6\cdot899
\label{mab:rm:eq:chaos-sum2}
\end{equation}

produit

\[
D_0=\frac{899}{648},\qquad
s_0=\frac{36}{\pi\sqrt{899}}.
\]

Par conséquent,

\begin{equation}
s_0\phi_m=\frac{2(3m-1)}{\sqrt{899}}.
\label{mab:rm:eq:pi-cancellation}
\end{equation}

Le facteur \(\pi\) s’annule avant d’effectuer les itérations. La famille ne
dépend ni de \(\delta_F\) ni de \(\alpha_F\) comme données d’entrée.

L’annulation est une propriété structurelle de la normalisation. Avant de
normaliser,
le deuxième moment angulaire est
\[
\frac1{12}\sum_{m=1}^{12}\phi_m^2
=\pi^2D_0.
\]
La condition de criticité quadratique fixe
\(s_0^2\pi^2D_0=2\), et chaque phase normalisée dépend donc seulement de
l’arithmétique entière de \(3m-1\). La fonctionnelle conserve la provenance angulaire de
la roue, mais l’itération reçoit une famille réelle intrinsèque. C’est
précisément ce qui permet de distinguer une dérivation d’un étalonnage
par \(\pi\).

\section{Représentation analytique exacte et expression trigonométrique fermée}

On définit

\begin{equation}
u_0(x)=1-\frac1{12}\sum_{m=1}^{12}
\cos\!\left(\frac{2(3m-1)}{\sqrt{899}}x\right).
\label{mab:rm:eq:chaos-profile}
\end{equation}

Le développement à l’origine est

\begin{equation}
u_0(x)=x^2-\frac{715769}{2424603}x^4+O(x^6).
\label{mab:rm:eq:chaos-taylor}
\end{equation}

La somme finie admet l’expression fermée suivante. Si \(y=x/\sqrt{899}\),

\begin{equation}
u_0(x)=1-\frac{\cos(37y)\sin(36y)}{12\sin(3y)}
=1-\frac{\sin(73y)-\sin y}{24\sin(3y)}.
\label{mab:rm:eq:chaos-closed}
\end{equation}

La singularité apparente en \(y=0\) est éliminable et reproduit \cref{mab:rm:eq:chaos-taylor}.

\begin{proof}[Dérivation de la forme fermée]
Les douze arguments forment une progression arithmétique :
\[
2(3m-1)y=(6m-2)y,\qquad m=1,\ldots,12.
\]
La somme des cosinus d’une progression satisfait
\[
\sum_{m=1}^{N}\cos(a+md)
=\frac{\sin(Nd/2)}{\sin(d/2)}
 \cos\!\left(a+\frac{N+1}{2}d\right).
\]
En prenant \(N=12\), \(d=6y\) et \(a=-2y\), on obtient
\[
\sum_{m=1}^{12}\cos((6m-2)y)
=\frac{\sin36y}{\sin3y}\cos37y.
\]
L’identité de transformation produit--en--somme
\(2\cos37y\sin36y=\sin73y-\sin y\) donne la deuxième expression de
\cref{mab:rm:eq:chaos-closed}.
\end{proof}

Le coefficient quartique n’est pas non plus ajusté. Les sommes entières
\begin{equation}
\sum_{m=1}^{12}(3m-1)^2=5394,\qquad
\sum_{m=1}^{12}(3m-1)^4=4294614
\label{mab:rm:eq:chaos-moments24}
\end{equation}
insérées dans le développement de \(\cos z\) produisent exactement
\[
\frac1{24\cdot12}\sum_m
\left(\frac{2(3m-1)}{\sqrt{899}}\right)^4
=\frac{715769}{2424603}.
\]
Ainsi, tant la normalisation quadratique que la première correction non
linéaire sont des moments finis de la même roue.

\section{Appartenance à la classe analytique unimodale quadratique}

Considérons

\[
f_\lambda(x)=1-\lambda u_0(x),\qquad \lambda>0.
\]

Pour \(0<x\le1\), tous les arguments

\[
\frac{2(3m-1)}{\sqrt{899}}x
\]

appartiennent à \((0,\pi)\). Par conséquent,

\[
u_0'(x)>0,\qquad u_0'''(x)<0.
\]

Comme \(u_0\) est paire, \(x=0\) est son unique point critique et \(u_0''(0)=2\). La dérivée schwarzienne de \(f_\lambda\) est

\begin{equation}
S(f_\lambda)
=\frac{u_0'''}{u_0'}
-\frac32\left(\frac{u_0''}{u_0'}\right)^2<0
\qquad(0<|x|\le1).
\label{mab:rm:eq:negative-schwarzian}
\end{equation}

\begin{theorem}[Appartenance exacte à la classe unimodale]
Pour
\[
0<\lambda\le \frac{2}{u_0(1)},
\]
la famille \(f_\lambda=1-\lambda u_0\) est une application réelle--analytique de \([-1,1]\) dans lui-même, paire, avec un unique point critique quadratique à l’origine et une dérivée schwarzienne négative en dehors de ce point. Elle appartient donc, sans utiliser les valeurs de Feigenbaum comme valeur cible, à la classe \(S\)-unimodale quadratique sur laquelle agit le renormalisateur de doublement de période \cite{rm:Feigenbaum1978,rm:Lanford1982}. Pour \(\lambda\) en dehors de cet intervalle, le germe analytique et sa forme locale sont conservés, mais l’application de l’intervalle dans lui-même n’est pas automatiquement affirmée.
\end{theorem}

Le théorème établit la contribution intrinsèque de HMT préalable au théorème
d’universalité : le germe analytique et son type de criticité.

\begin{proof}[Détails du signe et de l’application dans soi-même]
En dérivant \cref{mab:rm:eq:chaos-profile},
\begin{align*}
u_0'(x)&=\frac1{12}\sum_m k_m\sin(k_mx),\\
u_0'''(x)&=-\frac1{12}\sum_m k_m^3\sin(k_mx),
\qquad k_m=\frac{2(3m-1)}{\sqrt{899}}.
\end{align*}
Pour \(0<x\le1\), \(0<k_mx<\pi\), donc tous les termes de la première
ligne sont positifs et ceux de la seconde négatifs. La parité étend la
monotonie aux deux côtés de l’origine. Comme \(u_0(0)=0\) et
\(0\le u_0(x)\le u_0(1)\), la borne
\(0<\lambda\le2/u_0(1)\) implique
\(-1\le1-\lambda u_0(x)\le1\). Enfin,
\(u_0''(0)=1/12\sum k_m^2=2\), de sorte que le point critique est
quadratique.
\end{proof}

\section{Suites finies de paramètres superstables}

La suite superstable certifiée jusqu’au niveau huit produit

\[
\delta_8=4.66921218915\ldots,
\qquad
\alpha_8=2.50296847988\ldots.
\]

Son statut est celui de \emph{Certificat fini exempt de valeurs cibles} : chaque
approximant est obtenu à partir de la famille préalablement définie, sans introduire
comme donnée la limite universelle. La coïncidence décimale fournit un
indicateur de convergence, mais ne remplace pas le théorème d’hyperbolicité.

Une extension numérique indépendante de la même suite, avec un pas
adapté à la séparation entre les racines et un contrôle de la période exacte,
atteint provisoirement :

\begin{center}
\small
\begin{tabular}{@{}ccll@{}}
\toprule
Niveau & \(\lambda_n\) & \(\delta_n\) & \(\alpha_n\)\\
\midrule
11 & 1.87403828195439908045 & 4.66920170694424982745 & 2.50290847517165181764\\
12 & 1.87403836411023460451 & 4.66920163036308158848 & 2.50290800379001546915\\
13 & 1.87403838170549870372 & 4.66920161361839219913 & 2.50290790268748193108\\
14 & 1.87403838547386545431 & 4.66920161007472591023 & 2.50290788100969754274\\
\bottomrule
\end{tabular}
\end{center}

Les niveaux 11--14 conservent leur statut d’évaluations numériques au-delà du certificat persistant jusqu’au niveau huit. L’extrapolation d’Aitken donne $\lambda_\infty^{\HMT}\simeq1.8740383865008918080$ et $\alpha_F\simeq2.50290787509$. Le théorème de \ref{sec:p3-final:cierre-feigenbaum} établit la loi d’échelle paramétrique des premières racines ; l’évaluation spatiale conserve sa portée propre.

Pour éviter une procédure numérique non explicitée, la méthode est formulée comme
suit. Définissez
\begin{equation}
F_n(\lambda)=f_\lambda^{\,2^n}(0).
\label{mab:rm:eq:superstable-equation}
\end{equation}
La racine superstable \(\lambda_n\) est la première racine de \(F_n\) à
droite de \(\lambda_{n-1}\) qui satisfait simultanément
\[
f_{\lambda_n}^{\,2^j}(0)\ne0
\quad(0\le j<n).
\]
Cette inégalité exclut les périodes diviseurs. Avec des intervalles disjoints
\(I_n\ni\lambda_n\), on calcule
\begin{equation}
\delta_n=
\frac{\lambda_{n-1}-\lambda_{n-2}}
     {\lambda_n-\lambda_{n-1}},
\qquad
\alpha_n=-\frac{d_{n-1}}{d_n},
\label{mab:rm:eq:feigenbaum-approximants}
\end{equation}
où \(d_n=f_{\lambda_n}^{\,2^{n-1}}(0)\) est la distance orientée du
dernier point nouveau au point critique. Un certificat persistant doit conserver
les intervalles de racine, l’exclusion des périodes inférieures et la propagation
de l’erreur de \cref{mab:rm:eq:feigenbaum-approximants} ; un simple nombre décimal imprimé ne
remplit pas ces trois conditions.

La règle de discrétisation constitue également un contrôle négatif. Une
borne inférieure fixe supérieure à la séparation
\(|\lambda_n-\lambda_{n-1}|\) peut manquer la première racine et produire
néanmoins une suite numériquement plausible. C’est pourquoi l’incrément
doit être proportionné à la séparation précédente et chaque candidat doit être vérifié
par sa période exacte.

\section{Conditions de certification de l’opérateur de renormalisation}

Sur un espace de fonctions paires analytiques, nous définissons

\begin{equation}
\mathcal R(f)(x)=-a(f)^{-1}f\bigl(f(-a(f)x)\bigr),
\qquad a(f)=-f(1).
\label{mab:rm:eq:renormalization}
\end{equation}

La clôture complète de \(\delta_F\) et de \(\alpha_F\) exige de certifier conjointement :

\begin{enumerate}
\item l’existence d’un point fixe \(g=\mathcal R(g)\) dans une boule analytique déclarée ;
\item l’unicité locale de ce point fixe ;
\item une unique valeur propre réelle instable de \(D\mathcal R_g\), égale à \(\delta_F\) ;
\item un rayon spectral strictement inférieur à un sur le complément stable ;
\item la transversalité de la courbe \(\lambda\mapsto f_\lambda\) par rapport à la variété stable de \(g\).
\end{enumerate}

Une implémentation HMT naturelle emploie des troncatures de Chebyshev d’ordre
\(9^m\), de l’arithmétique d’intervalles et des polynômes de rayons. La preuve doit
fournir une boule, un inverse approché, une borne du défaut et une borne
de Lipschitz satisfaisant \(Y+Z(r)<r\). Ainsi, la profondeur
nonadique paramètre un certificat de point fixe et non une simple suite
d’approximations décimales.

Dans une base paire de Chebyshev,
\[
g(x)=\sum_{j=0}^{N}g_{2j}T_{2j}(x),\qquad N=9^m,
\]
le résidu fini est \(F_N(g)=\Pi_N(\mathcal Rg-g)\). Si \(A_N\) approche
\((DF_N(\bar g))^{-1}\), les quantités à certifier sont
\begin{align*}
Y&=\|A_NF_N(\bar g)\|,\\
Z_0&=\|I-A_NDF_N(\bar g)\|,\\
Z_1(r)&=\sup_{\|h\|\le r}
\|A_N(DF_N(\bar g+h)-DF_N(\bar g))\|.
\end{align*}
Une inégalité
\(Y+(Z_0+Z_1(r))r<r\) enferme un unique point fixe. Le bloc spectral
ultérieur doit isoler une direction instable et contracter le complément. L’usage
de \(9^m\) intervient ainsi comme loi de raffinement du certificat, non comme
décoration numérologique.

\begin{corollary}[Confluence exacte du mode trente]
Le mode \(30\) admet trois réalisations simultanées et typologiquement
distinctes :
\begin{align*}
30&\longmapsto(3\cdot30,4\cdot30)=(90,120)
&&\text{dans la complétion constitutive},\\
30&\longmapsto30^2-1=899=29\cdot31
&&\text{dans la normalisation critique},\\
30&\longmapsto
M_{30}=\begin{pmatrix}30&899\\1&30\end{pmatrix}
&&\text{dans la transformation hyperbolique de Pell}.
\end{align*}
Les trois identités sont exactes et sont obtenues sans prescrire de valeurs cibles. L’invariant conservé est
l’entier central \(30\) avec son orientation : extension \(3\to4\),
défaut quadratique \(-1\) et unité de norme un. Ce qui n’est pas affirmé est
l’égalité entre l’opérateur de vide et le renormalisateur de doublement.
\end{corollary}

\begin{remark}[Contrôle négatif]
Une deuxième valeur propre en dehors du disque unité, une boule de rayons sans solution, une perte d’unicité ou la tangence de la famille HMT à la variété stable réfutent la promotion universelle. Aucune ne réfute le germe exact ni le théorème de la dérivée schwarzienne.
\end{remark}

\begin{proposition}[Résultat principal]
La signification structurelle de la constante de doublement réside dans
l’annulation dodécaphasique de \(\pi\), dans l’identité
\(899=30^2-1\), dans la forme trigonométrique fermée et dans l’appartenance
démontrée à la classe analytique sur laquelle agit le renormalisateur. Son
développement décimal constitue exclusivement l’évaluation terminale.
\end{proposition}

\subsection*{Portée logique et domaine de validité}
Le profil, la normalisation, l’unité de Pell, l’unimodalité et la schwarzienne conservent les démonstrations précédentes. Les huit niveaux certifiés et les extrapolations des niveaux 11--14 conservent leurs portées numériques respectives. Le développement de \ref{sec:p3-final:cierre-feigenbaum} démontre la loi d’échelle paramétrique de la suite initiale des premières racines dans le secteur uniforme, par identification des centres et contrôle du reste. Cette conclusion ne transfère pas au rapport spatial une démonstration portant sur le rapport paramétrique.

\subsection{Déformation harmonique pondérée et sélection de la branche superstable}
\label{mab:sec:extension-dodecafasica-ponderada-armonica}

Le profil uniforme possède une extension paramétrique qui conserve les douze
phases et sépare la déformation harmonique du paramètre dynamique. Pour
\(1\le m\le12\), soient

\begin{equation}
 \phi_m=\frac{(30m-10)\pi}{180},
 \qquad
 p(\phi)=12\cos(3\phi)-\cos(4\phi).
 \label{mab:eq:fases-y-deformacion-armonica-armonizada}
\end{equation}

Pour tout \(\eta\) réel tel que les poids restent positifs, nous définissons

\begin{equation}
 w_m(\eta)=\frac{1+\eta p(\phi_m)}{12}.
 \label{mab:eq:pesos-dodecafasicos-armonizados}
\end{equation}

Les sommes des harmoniques de rang trois et quatre sur les douze phases s’annulent ;
par conséquent,

\[
 \sum_{m=1}^{12}w_m(\eta)=1.
\]

Soit

\begin{equation}
 s_\eta^2=\frac{2}{\sum_{m=1}^{12}w_m(\eta)\phi_m^2},
 \qquad
 u_\eta(x)=1-\sum_{m=1}^{12}w_m(\eta)
 \cos(s_\eta\phi_m x).
 \label{mab:eq:familia-ponderada-armonizada}
\end{equation}

Le développement du cosinus et la définition de \(s_\eta\) donnent

\[
 u_\eta(x)=x^2+O(x^4)\qquad(x\to0).
\]

Par conséquent, \(u_\eta\) est réelle--analytique, paire et conserve un point
critique quadratique à l’origine. Nous introduisons le paramètre dynamique par

\begin{equation}
 f_{\lambda,\eta}(x)=1-\lambda u_\eta(x).
 \label{mab:eq:familia-feigenbaum-ponderada-armonizada}
\end{equation}

Un paramètre superstable de période exacte \(2^n\) satisfait

\[
 f_{\lambda,\eta}^{\,2^n}(0)=0,
 \qquad
 f_{\lambda,\eta}^{\,2^k}(0)\ne0\quad(0\le k<n).
\]

L’équation peut posséder plusieurs branches. On fixe une branche
\(\mathscr B\) en parcourant, immédiatement après le paramètre du
niveau précédent, un maillage adaptatif ; on prend le premier intervalle avec changement
de signe et l’on applique une dichotomie. La notation
\(\lambda_n^{\mathscr B}(\eta)\) incorpore ce sélecteur et n’attribue pas
d’unicité globale à l’équation superstable. Le premier niveau est déterminé
exactement :

\[
 \lambda_1^{\mathscr B}(\eta)=\frac1{u_\eta(1)},
\]

si \(u_\eta(1)>0\), parce que
\(f_{\lambda,\eta}^{\,2}(0)=f_{\lambda,\eta}(1)=0\).
Étant donnés trois niveaux consécutifs et

\[
 x_n^{\mathscr B}(\eta)
 :=f_{\lambda_n^{\mathscr B}(\eta),\eta}^{\,2^{n-1}}(0),
\]

on forme

\begin{equation}
 \delta_n^{\mathscr B}(\eta)
 =\frac{\lambda_{n-1}^{\mathscr B}(\eta)-
         \lambda_{n-2}^{\mathscr B}(\eta)}
        {\lambda_n^{\mathscr B}(\eta)-
         \lambda_{n-1}^{\mathscr B}(\eta)},
 \qquad
 \alpha_n^{\mathscr B}(\eta)
 =\frac{|x_{n-1}^{\mathscr B}(\eta)|}
        {|x_n^{\mathscr B}(\eta)|}.
 \label{mab:eq:aproximantes-ponderados-armonizados}
\end{equation}

Pour \(\eta=0\), on retrouve exactement le profil uniforme. L’extension
pondérée introduit un paramètre transversal : elle permet d’étudier
la transversalité de la famille par rapport à la variété stable du
renormalisateur sans transformer une coïncidence finie en une preuve
d’universalité.

\begin{remark}[Contrôle négatif]
L’extension est réfutée si les poids perdent leur positivité dans le domaine
déclaré, si la normalisation cesse de produire un coefficient quadratique
unitaire, si le sélecteur n’exclut pas les périodes diviseurs ou si une branche est choisie
en utilisant comme entrée les valeurs universelles terminales.
\end{remark}
 \section{Sélection des premières racines et loi d’échelle paramétrique de Feigenbaum}
\label{sec:p3-final:cierre-feigenbaum}

Le lecteur dodécaphasique uniforme détermine également une suite de premières racines de retour critique. On démontre que le rapport de leurs écarts converge vers la constante de Feigenbaum, en conservant la sélection initiale, l’orientation et la mémoire de la construction.

% Focal insertion; hierarchy inherited from the host.
\begingroup
\let\HMTFGHostSection\section
\edef\HMTFGSavedSecnumdepth{\number\value{secnumdepth}}
\setcounter{secnumdepth}{4}
\let\HMTFGHostSubsection\subsection
\let\HMTFGHostSubsubsection\subsubsection
\let\section\HMTFGHostSubsection
\let\subsection\HMTFGHostSubsubsection
\let\subsubsection\paragraph
% Fragmento maestro: incluir desde el anfitrion, con \HMTFGRoot fijada a la raiz del paquete.
% Requiere amsmath, amssymb/mathrsfs o unicode-math, amsthm, booktabs, longtable, hyperref, needspace.
% El anfitrion conserva sus entornos theorem, proposition, lemma y proof.
\providecommand{\HMTFGRoot}{.}
\begingroup
\providecommand{\ZZ}{}\renewcommand{\ZZ}{\mathbb Z}
\providecommand{\RR}{}\renewcommand{\RR}{\mathbb R}
\providecommand{\FF}{}\renewcommand{\FF}{\mathbb F}
\providecommand{\APP}{}\renewcommand{\APP}{\mathrm{APP}}
\providecommand{\TPK}{}\renewcommand{\TPK}{\mathrm{TPK}}
\providecommand{\TRIT}{}\renewcommand{\TRIT}{\{-1,0,+1\}}
% unicode-math usa lBrack/rBrack; newtxmath ya proporciona llbracket/rrbracket.
\providecommand{\llbracket}{\lBrack}
\providecommand{\rrbracket}{\rBrack}
% Fragmento de inserción; se incluye desde la raíz de este paquete editorial.
% Las fuentes completas y los cortes materiales constan en ANTECEDENTES_Y_DEPENDENCIAS.md.
% Procedencia: XI REV11 ES, sections/nucleo.tex, tpk_desarrollo_integrado.tex,
% ch_tres_vueltas_ext.tex, ch_elevacion_catalogal.tex, ch_cierre_cardinal.tex,
% ch_descenso_por_fibras.tex; perfil integral public_final/feigenbaum/c37_feigenbaum.tex;
% Grassmannianos REV07 ES, sections/jacobi_spectral.tex;
% TRANSVERSALIDAD_Y_ESCALADO_LOCAL_FEIGENBAUM.md, sección 14.
% Ampliación expositiva de demostraciones existentes; no introduce resultados nuevos.
\providecommand{\ZZ}{\mathbb Z}
\providecommand{\RR}{\mathbb R}
\providecommand{\FF}{\mathbb F}
\providecommand{\APP}{\mathrm{APP}}
\providecommand{\TPK}{\mathrm{TPK}}
\providecommand{\TRIT}{\{-1,0,+1\}}

\section{Antécédents operatoriels de la famille de criticité}
\label{hmtfg:antecedentes-operatorios}

La famille de criticité est construite à partir d’une lecture du calendrier TPK. Les deux feuilles arithmétiques, l’orientation tritique et le transport avec mémoire déterminent d’abord les états et leurs prolongements. La lecture ordonnée du continu permet ensuite d’évaluer des fonctions analytiques, de comparer des paramètres et de sélectionner des racines. Ce développement conserve l’état généalogique et ses deux lectures : la carte archimédienne fournit l’espace d’évaluation, tandis que la lecture d’incidence conserve les relations de frontière. La structure discrète conjointe du continu maintient ses cinq constructions consubstantielles. La présente exposition utilise son lecteur ordonné et conserve l’antécédent commun dont il provient.

\subsection{Cellules arithmétiques, orientation et transport}
\label{hmtfg:app-trit-transporte}

% XI REV11, sections/nucleo.tex: definiciones de APP y TRIT;
% tpk_desarrollo_integrado.tex:22--165, transporte levantado y calendario.
Soient \(D_9=\{1,\ldots,9\}\) et \(X_9=(\ZZ/9\ZZ)^2\). Pour
\(n\in\ZZ\), la division avec représentant positif est
\[
 \rho_9(n)=1+((n-1)\bmod9),\qquad
 q_9(n)=\frac{n-\rho_9(n)}9,\qquad
 n=\rho_9(n)+9q_9(n).
\]
La cellule APP associée à \((p,q)\in D_9^2\) est
\[
 a_{p,q}=(p,q;R^\Sigma,q^\Sigma,R^\Pi,q^\Pi),
 \quad
 p+q=R^\Sigma+9q^\Sigma,\qquad
 pq=R^\Pi+9q^\Pi.
\]
Les deux couples \((R^\Sigma,q^\Sigma)\) et \((R^\Pi,q^\Pi)\) conservent respectivement l’évaluation additive et l’évaluation multiplicative. La formule de division reconstruit les deux évaluations avant toute réduction ultérieure. La sélection d’une feuille active est enregistrée dans une coordonnée de l’état et conserve la cellule complète.

Le représentant tritique équilibré \(r_3(n)\in\{-1,0,+1\}\) satisfait
\(n\equiv r_3(n)\pmod3\). Sa retenue est
\[
 \chi(a,b)=\frac{a+b-r_3(a+b)}3,\qquad
 a,b\in\{-1,0,+1\}.
\]
L’égalité
\[
 \chi(a,b)+\chi(r_3(a+b),c)
 =\chi(b,c)+\chi(a,r_3(b+c))
\]
s’obtient en écrivant les deux membres sous la forme
\((a+b+c-r_3(a+b+c))/3\). Cette identité conserve la retenue lors du regroupement d’une somme tritique. Le sélecteur opératoire est
\[
 M_{\rm ph}(r)=
 \begin{cases}
 +1,&r\in\{1,4,7\},\\
 -1,&r\in\{2,5,8\},\\
 0,&r\in\{3,6,9\}.
 \end{cases}
\]
La valeur \(+1\) active la feuille additive, la valeur \(-1\) la feuille multiplicative, et la valeur \(0\) conserve la position des deux curseurs pendant l’événement neutre. Les trois événements appartiennent à la chronologie complète et sont enregistrés dans la route.

Soient \(\mathcal D=\{N,E,S,O\}\), avec les vecteurs
\(\nu_N=(-1,0)\), \(\nu_E=(0,1)\), \(\nu_S=(1,0)\) et
\(\nu_O=(0,-1)\), et \(\Theta_{108}=\ZZ/108\ZZ\). La configuration observable est
\[
 q=(x_+,d_+,x_\times,d_\times,\theta)
 \in\mathcal Q_{\rm obs}=(X_9\times\mathcal D)^2\times\Theta_{108}.
\]
Pour le représentant \(\bar\theta\in\{0,\ldots,107\}\), on définit
\[
 r(\theta)=1+(\bar\theta\bmod9),\qquad
 \beta(\theta)=\left[\left\lfloor\frac{\bar\theta}{9}\right\rfloor\right]_{12},
\]
\[
 a_+(\theta)=\mathbf1_{\{1,4,7\}}(r(\theta)),\qquad
 a_\times(\theta)=\mathbf1_{\{2,5,8\}}(r(\theta)).
\]
La transition des positions est
\[
 x_+'=x_++a_+(\theta)\nu_{d_+},\qquad
 x_\times'=x_\times+a_\times(\theta)\nu_{d_\times}.
\]
L’involution \(\jmath\) échange \(N,S\) et \(E,O\). Avec
\[
 h(\theta)=\mathbf1_{\{0,27,54,81\}}\bigl((\bar\theta+1)\bmod108\bigr),
\]
l’opérateur observable complet est fixé par
\[
 F_{\rm obs}(q)=
 (x_+',\jmath^{h(\theta)}d_+,x_\times',
   \jmath^{h(\theta)}d_\times,\theta+1).
\]

\begin{proposition}[Conservation du transport relevé et retour observable]
\label{hmtfg:transporte-calendario}
Le transport conserve le déplacement entier au moyen du registre des franchissements de frontière. L’opérateur \(F_{\rm obs}\) est inversible et satisfait
\(F_{\rm obs}^{108}=I\).
\end{proposition}
\begin{proof}
Soit \(s:X_9\to D_9^2\) la section des représentants positifs. Pour
\(x'=x+a\nu_d\), \(a\in\{0,1\}\), on définit
\[
 b_d(x,a)=\frac{s(x)+a\nu_d-s(x')}{9}\in\ZZ^2.
\]
Si \(z'=z+b_d(x,a)\), alors
\[
 s(x')+9z'=s(x)+9z+a\nu_d.
\]
La somme de cette égalité le long d’une route donne
\[
 s(x_n)+9z_n=s(x_0)+9z_0+\sum_{t=0}^{n-1}a_t\nu_{d_t}.
\]
La partition de la somme conserve l’identité par concaténation. Pour inverser \(F_{\rm obs}\), on soustrait d’abord une unité à la phase ; cette phase détermine l’inversion des directions et le curseur déplacé. On récupère les directions précédentes et on soustrait le vecteur correspondant. Chaque bloc de \(27\) pas contient neuf activations de chaque curseur. Le bloc suivant inverse les directions et annule leurs déplacements entiers. En \(54\) pas, les positions et les directions sont rétablies ; en \(108\), la coordonnée du calendrier est également rétablie.
\end{proof}

L’état complet conserve, outre cette configuration, les deux évaluations arithmétiques, les retenues, l’orientation, la trajectoire, le préfixe, le cylindre compatible, sa frontière, le degré de mémoire et le registre de survie. Son actualisation se factorise sous la forme
\[
 U_t=\mathsf{Mem}_t\circ\mathsf{Tra}_t\circ\mathsf{Sel}_t.
\]
L’holonomie \(\Gamma_9\) compose neuf transitions complètes. Le retour de phase et l’avancement de mémoire sont des coordonnées différentes de cette composition. La construction suivante fixe les transitions et leur actualisation champ par champ.

% Recuperación íntegra del propietario de 855 líneas, con etiquetas prefijadas.
\input{inserciones_20260930/feigenbaum/00a_transiciones_nonadicas.tex}

\subsection{Raffinement des cylindres et lecture ordonnée}
\label{hmtfg:refinamiento-ordenado}

% XI REV11, ch_elevacion_catalogal.tex: bloques de 54 aristas, beta_blk,
% partición efectiva de 729 subcilindros; ch_cierre_cardinal.tex:375--497.
Pour \(\mathbf b=(b_1,\ldots,b_6)\in\FF_3^6\), soit
\(\upsilon:\FF_3\to\{-1,0,+1\}\) le représentant équilibré. Les trois prolongements construits arête par arête définissent
\[
 \Gamma^{[54]}_{\mathbf b,k}
 =\gamma_{\upsilon(b_6),k+45}\circ\cdots\circ
   \gamma_{\upsilon(b_1),k}.
\]
En partant de \(Y_0=\{x_\star\}\), on pose
\[
 Y_{N+1}
 =\{\Gamma^{[54]}_{\mathbf b,1+54N}(x):
          x\in Y_N,\ \mathbf b\in\FF_3^6\}.
\]
La troncature \(\rho^{\rm blk}_{N+1,N}\) supprime les cinquante-quatre dernières arêtes, en récupérant l’état antérieur au moyen du registre de transition.

\begin{proposition}[Codage effectif des blocs de prolongement]
\label{hmtfg:bloques-729}
La lecture des six couples tritiques de chaque bloc définit des bijections
\[
 \beta_N:Y_N\xrightarrow{\ \cong\ }(\FF_3^6)^N
\]
qui commutent avec la troncature. Chaque état de \(Y_N\) possède exactement \(729\) prolongements de bloc dans \(Y_{N+1}\).
\end{proposition}
\begin{proof}
À chaque tour, le registre conserve le couple tritique et sa retenue. Les trois couples produisent les trois sorties distinctes \(\chi(\varepsilon,\varepsilon)=\varepsilon\), et le registre permet de récupérer la branche utilisée, y compris celle de retenue nulle. Les six concaténations produisent donc tous les mots de \(\FF_3^6\). Deux mots différents se distinguent au premier couple enregistré où ils diffèrent. La récupération du préfixe prouve la compatibilité avec la troncature. Par récurrence sur \(N\), chaque mot de longueur \(N\) possède un unique état antécédent ; le mot suivant peut être n’importe laquelle des \(3^6=729\) possibilités.
\end{proof}

Soit \(\Omega_{\Gamma}^{\rm cal}=\varprojlim_NY_N\). Les bijections précédentes induisent
\[
 \Omega_{\Gamma}^{\rm cal}\cong(\FF_3^6)^{\mathbb N}.
\]
La correspondance et son inverse sont continues : chaque coordonnée finie dépend d’un préfixe fini. Les cylindres sont ouverts et fermés. Chaque cylindre de profondeur \(N\) se partitionne en ses \(729\) sous-cylindres de profondeur \(N+1\) ; chacun contient un prolongement infini, obtenu, par exemple, en poursuivant par la branche nulle dans tous les blocs ultérieurs.

% Antecedentes efectivos del continuo conjunto y del universo interno.
% Se imprimen en ambos destinos; la etiqueta de universo no reemplaza su regla.
\input{inserciones_20260930/feigenbaum/00b_continuo_universo.tex}

\subsection{Lecture ordonnée des cylindres de bloc}
\label{hmtfg:publicacion-cilindros-bloque}

Pour le représentant \([b]_3\in\{0,1,2\}\), on fixe
\[
 \nu(\mathbf b)=\sum_{i=1}^6[b_i]_3\,3^{6-i}\in\{0,\ldots,728\}.
\]
Dans la complétion ordonnée HMT, un préfixe
\(s=(m;\mathbf b_0,\ldots,\mathbf b_{N-1})\), \(m\in\ZZ\), détermine
\[
 a_N(s)=m+\sum_{j=0}^{N-1}\nu(\mathbf b_j)\,729^{-j-1},\qquad
 I_N(s)=[a_N(s),a_N(s)+729^{-N})_{\rm HMT}.
\]
La retenue conserve le raccord entre les deux représentants d’une frontière. Si \(C_N(s)=\overline{I_N(s)}\), les relations précédentes donnent
\[
 I_N(s)=\bigsqcup_{\mathbf b\in\FF_3^6}I_{N+1}(s;\mathbf b),
 \qquad \operatorname{diam}C_N(s)=729^{-N}\longrightarrow0.
\]

\begin{proposition}[Couverture de la lecture archimédienne]
\label{hmtfg:cobertura-arquimediana}
L'application
\[
 P_{\rm ar}:\ZZ\times\Omega_\Gamma^{\rm cal}
       \longrightarrow\mathbb R_{\rm HMT},\qquad
 \{P_{\rm ar}(m,\xi)\}
    =\bigcap_N C_N(m;\beta_N(\rho_{\infty,N}\xi))
\]
est bien définie, est surjective et détermine
\[
 (\ZZ\times\Omega_\Gamma^{\rm cal})/\!\sim_{P_{\rm ar}}
       \ \cong\ \mathbb R_{\rm HMT}.
\]
\end{proposition}
\begin{proof}
Les extrémités gauches forment une suite de Cauchy dans la complétion HMT, puisque les intervalles sont emboîtés et que leurs diamètres tendent vers zéro. Leur limite appartient à tous les intervalles fermés et est unique. Pour tout élément de la complétion, on choisit d’abord l’intervalle entier qui le contient, puis l’unique sous-intervalle semi-ouvert de chaque partition qui le contient. La proposition précédente réalise chacun de ces choix par un prolongement TPK effectif. Les préfixes compatibles définissent \(\xi\), dont la lecture est l’élément choisi. Aux frontières, la retenue identifie les écritures équivalentes après leur construction. Le quotient par égalité de lecture est, par définition et surjectivité, en bijection avec la complétion.
\end{proof}

La construction est également effectuée dans l’univers généalogique interne \(\mathfrak G=\mathfrak G_{\rm HMT}(h,m_\infty)\). Ici, \(h\) code les règles operatorielles et \(m_\infty\) l’histoire compatible ; l’univers s’obtient par clôture transitive initiale, définissabilité à chaque successeur et réunion aux stades limites :
\[
 G_0=\operatorname{TC}\{\omega,u_{h,m},h,m_\infty\},\qquad
 G_{\beta+1}=\operatorname{Def}_{\Sigma_{\rm HMT}}(G_\beta),\qquad
 G_\lambda=\bigcup_{\beta<\lambda}G_\beta .
\]
Les relations de \(\Sigma_{\rm HMT}\) sont celles des opérations déjà construites et de leurs codes. Le domaine interne est donc maintenu dans toutes les affirmations de couverture. Notons \(\mathbb R_{\rm HMT}^{\mathfrak G}\) sa complétion et \(\mathbb R^{\mathfrak G}\) la réalisation ordonnée au sein du même univers.

\begin{proposition}[Carte ordonnée interne]
\label{hmtfg:carta-interna}
Il existe un unique isomorphisme de corps ordonnés archimédiens complets
\[
 \iota:\mathbb R_{\rm HMT}^{\mathfrak G}
             \xrightarrow{\ \cong\ }\mathbb R^{\mathfrak G}
\]
qui fixe l’unité. Il conserve les extrémités rationnelles, les racines carrées positives et les limites de suites convergentes.
\end{proposition}
\begin{proof}
L’unité fixe les entiers et les rationnels. À chaque élément de la complétion HMT est associée la coupure des rationnels qui lui sont inférieurs ; la complétude de la réalisation ordonnée fournit l’unique élément ayant cette coupure. La densité rationnelle prouve l’injectivité et la conservation stricte de l’ordre. Le même procédé en sens inverse prouve la surjectivité. Les approximations rationnelles par excès et par défaut conservent la somme et le produit, et donc la structure de corps. L’élément positif dont le carré est \(a>0\) est transporté vers l’élément positif dont le carré est \(\iota(a)\). Si une suite converge, pour chaque tolérance rationnelle positive ses derniers termes sont dans le voisinage rationnel de la limite ; la conservation de l’ordre transporte cette condition. Toutes ces opérations s’effectuent dans \(\mathfrak G\). L’unicité découle de la conservation des coupures rationnelles.
\end{proof}

\subsection{Descente par les fibres et conservation du sélecteur}
\label{hmtfg:descenso-selector}

% XI REV11, ch_descenso_por_fibras.tex:10--39; prueba completa pertinente.
\begin{proposition}[Critère de descente d’une opération]
\label{hmtfg:criterio-descenso}
Soient \(q:\Omega\twoheadrightarrow J\) et
\(\star:D\subseteq\Omega^2\to\Omega\), avec \(D\) saturé par les fibres de \(q\times q\). Il existe une unique opération \(\bar\star\) sur \((q\times q)(D)\) telle que
\[
 q(x\star y)=q(x)\,\bar\star\,q(y)
\]
si et seulement si
\[
 q(x)=q(x'),\quad q(y)=q(y')
 \ \Longrightarrow\ q(x\star y)=q(x'\star y')
\]
pour les couples du domaine.
\end{proposition}
\begin{proof}
Si \(\bar\star\) existe, les deux résultats sont l’évaluation du même couple d’arguments. Réciproquement, on choisit des antécédents des arguments visibles et on définit le résultat comme la lecture de leur opération. La saturation fixe le domaine ; la condition sur les fibres rend le résultat indépendant des antécédents choisis. La surjectivité prouve l’unicité.
\end{proof}

Dans une application au moyen de \(P_{\rm ar}\), ce critère identifie la condition précise permettant de transporter une opération depuis les états complets jusqu’à leur lecture. Dans la carte \(\iota\), la structure de corps et les limites satisfont déjà cette condition. Le profil analytique est ensuite construit à partir du calendrier, et cette compatibilité est utilisée pour transporter son itération et ses racines.

\subsection{Lecture dodécaphasique et normalisation quadratique}
\label{hmtfg:perfil-dodecafasico}

Le profil, ses phases et sa normalisation sont ceux déjà construits en \eqref{mab:rm:eq:critical-phase-reader}--\eqref{mab:rm:eq:chaos-profile}. Sa forme fermée et ses moments sont démontrés en \eqref{mab:rm:eq:chaos-closed} et \eqref{mab:rm:eq:chaos-moments24}. La somme finie de cosinus détermine le prolongement aux singularités apparentes en chaque zéro du dénominateur de la forme fermée. Le moment angulaire conserve des représentants réels orientés : les modifier modulo $2\pi$ change le moment. Le lecteur et son secteur uniforme constituent la structure de départ explicite de cette réalisation.

\subsection{Transport des itérées, des racines et des minima}
\label{hmtfg:transporte-raices}

% Reunión existente: TRANSVERSALIDAD_Y_ESCALADO_LOCAL_FEIGENBAUM.md, sección 14.1.
Dans \(\mathbb R_{\rm HMT}^{\mathfrak G}\), la série convergente
\[
 \cos_{\rm HMT}(t)=\sum_{j=0}^{\infty}
                   \frac{(-1)^j t^{2j}}{(2j)!}
\]
définit la réalisation analytique correspondante. Soient
\[
 u_{\rm HMT}(x)=\frac1{12}\sum_{m=1}^{12}
  \left[1-\cos_{\rm HMT}
   \left(\frac{2(3m-1)x}{\sqrt{899}_{\rm HMT}}\right)\right],
 \qquad F_a(x)=1-a\,u_{\rm HMT}(x),
\]
\[
 S_n^{\rm HMT}(a)=F_a^{\,2^n}(0),\qquad
 S_n(\lambda)=f_\lambda^{\,2^n}(0).
\]

\begin{proposition}[Conjugaison ordonnée de la famille et du sélecteur de racines]
\label{hmtfg:conjugacion-selector}
Pour les arguments internes des expressions précédentes,
\[
 \iota\circ F_a=f_{\iota(a)}\circ\iota,\qquad
 \iota(S_n^{\rm HMT}(a))=S_n(\iota(a)).
\]
La carte transporte exactement les conditions de retour, de période exacte et d’ordre des paramètres. Si un ensemble de racines admissibles possède un minimum, l’image de ce minimum est le minimum de l’ensemble image.
\end{proposition}
\begin{proof}
La carte fixe les rationnels, transporte la racine carrée positive et conserve les limites des sommes partielles du cosinus. Ainsi, \(\iota(u_{\rm HMT}(x))=u_0(\iota(x))\). La compatibilité avec la somme et le produit prouve la première égalité ; la récurrence sur le nombre d’itérations prouve la seconde. Une racine se caractérise par l’égalité à zéro, conservée dans les deux sens. La période exacte \(2^n\) ajoute les inégalités \(F_a^{2^j}(0)\ne0\), \(0\le j<n\), également conservées par l’injectivité de \(\iota\). Si l’on exige que le paramètre dépasse un centre précédent, cette inégalité est transportée par conservation de l’ordre. Enfin, \(a\le b\) pour tout \(b\) de l’ensemble équivaut à \(\iota(a)\le\iota(b)\) pour tout élément de son image. La surjectivité sur les ensembles ainsi définis complète l’affirmation relative au minimum.
\end{proof}

Le domaine de ce transport est le continu interne indiqué. L’existence de chaque minimum et son identification à une branche dynamique précise sont des résultats de la construction de criticité qui suit ; la carte conserve ces propriétés lorsqu’elles ont été établies.

\subsection{Représentation spectrale finie du lecteur uniforme}
\label{hmtfg:jacobi-finito}

% Grassmannianos REV07, sections/jacobi_spectral.tex:35--136,
% mecanismo de Hankel, Gram--Schmidt, recurrencia y representación espectral.
% Especialización atómica ya reunida en la sección 14.2 del desarrollo Feigenbaum.
Les fréquences au carré du profil définissent la mesure positive
\[
 s_m=k_m^2=\frac{4(3m-1)^2}{899},\qquad
 \nu_0=\frac1{12}\sum_{m=1}^{12}\delta_{s_m},\qquad
 M_r=\int s^r\,d\nu_0(s).
\]
Ses douze nœuds sont distincts et positifs. Pour \(1\le r\le12\), la matrice de Hankel \(H_r=(M_{i+j})_{0\le i,j<r}\) est définie positive. En effet, pour \(p\) de degré inférieur à \(r\),
\[
 \mathbf p^{\mathsf T}H_r\mathbf p
   =\frac1{12}\sum_{m=1}^{12}p(s_m)^2>0
\]
sauf si \(p=0\), car un polynôme de degré au plus onze qui s’annule en douze nœuds distincts est nul. Au degré douze apparaît le polynôme \(\prod_m(s-s_m)\), de norme nulle pour cette mesure. La représentation spectrale correspondante est de dimension douze.

Soient \(p_0,\ldots,p_{11}\) les polynômes orthogonaux moniques,
\(h_j=\int p_j^2\,d\nu_0>0\) et
\(\psi_j=p_j/\sqrt{h_j}\). Comme \(\nu_0\) est une probabilité,
\(\psi_0=1\). La multiplication par \(s\) dans \(L^2(\nu_0)\) est autoadjointe et possède une matrice tridiagonale \(J_{12}\) dans cette base.

\begin{proposition}[Récurrence et entrelacement spectral exacts]
\label{hmtfg:jacobi-entrelazamiento}
Avec
\[
 a_j=\frac{\int s p_j(s)^2\,d\nu_0(s)}{h_j},\qquad
 \beta_j=\frac{h_j}{h_{j-1}}>0\quad(1\le j\le11),
\]
la matrice \(J_{12}\) a pour diagonale \(a_j\) et pour sous-diagonale
\(\sqrt{\beta_j}\). Si
\[
 Q_{mj}=\frac{\psi_j(s_m)}{\sqrt{12}}
 \quad(1\le m\le12,\ 0\le j\le11),\qquad
 D=\operatorname{diag}(s_1,\ldots,s_{12}),
\]
alors
\[
 Q^{\mathsf T}Q=I,\qquad DQ=QJ_{12},\qquad
 J_{12}=Q^{\mathsf T}DQ,
\]
et le profil satisfait
\begin{equation}
 u_0(x)=e_0^{\mathsf T}
          \bigl[I-\cos(x\sqrt{J_{12}})\bigr]e_0.
 \label{hmtfg:perfil-espectral}
\end{equation}
\end{proposition}
\begin{proof}
La positivité des systèmes de Hankel fournit de manière unique les polynômes orthogonaux moniques. En développant \(s p_j\) dans la base orthogonale, les coefficients de \(p_i\), \(i\le j-2\), s’annulent car
\(\langle sp_j,p_i\rangle=\langle p_j,sp_i\rangle=0\).
Le coefficient de \(p_{j+1}\) est un par monicité, celui de \(p_j\) est \(a_j\) et celui de \(p_{j-1}\) est \(h_j/h_{j-1}\). Au dernier degré, la même égalité s’entend dans \(L^2(\nu_0)\), où le polynôme monique de degré douze qui s’annule aux nœuds représente le vecteur nul. La normalisation donne la sous-diagonale positive indiquée.

L’orthonormalité équivaut à \(Q^{\mathsf T}Q=I\). L’égalité de récurrence à chaque nœud donne \(DQ=QJ_{12}\). Comme \(Q\) est carrée, on a également \(QQ^{\mathsf T}=I\), et l’on obtient la diagonalisation de \(J_{12}\). Ses valeurs propres positives permettent de définir sa racine carrée positive. Pour toute fonction définie en ces douze nœuds,
\[
 e_0^{\mathsf T}\Phi(J_{12})e_0
   =\frac1{12}\sum_{m=1}^{12}\Phi(s_m),
\]
car \(Qe_0=(1,\ldots,1)^{\mathsf T}/\sqrt{12}\).
Le choix \(\Phi(s)=1-\cos(x\sqrt{s})\) prouve
\eqref{hmtfg:perfil-espectral}.
\end{proof}

La représentation conserve exactement la fonctionnelle uniforme, ses moments et son profil. Ses premiers coefficients sont
\[
 a_0=M_1=2,\qquad
 \beta_1=M_2-M_1^2=\frac{2493348}{808201}.
\]
La généalogie complète demeure dans l’état TPK dont provient le lecteur ; \(J_{12}\) représente sa lecture spectrale particulière. La construction de Gram--Schmidt et la récurrence de Jacobi s’appliquent à cette mesure atomique, avec leur terminaison en dimension douze.

La chaîne operatorielle utilisée dans les sections suivantes est ainsi fixée : les cellules APP et l’orientation TRIT alimentent la transition TPK ; ses prolongements produisent les cylindres et leur lecture ordonnée ; les douze fenêtres orientées et la fonctionnelle uniforme produisent \(u_0\) ; la famille \(f_\lambda\) détermine les retours et les racines ; la carte interne conserve l’ordre de leurs sélecteurs. La réalisation spectrale décrit le même profil au moyen d’une matrice autoadjointe finie et d’un vecteur cyclique.

% Fragmento integrable. Preambulo anfitrion: amsmath, amssymb, longtable, hyperref.
\section{Profil dodécaphasique et retour quadratique}
\label{hmtfg:fragmento:perfil-dodecafasico-y-retorno-cuadratico}

\subsection{Généalogie et objet de la démonstration}
\label{hmtfg:desarrollo:1}

L’objet initial est la lecture critique de la roue dodécaphasique orientée déjà construite dans le corpus. APP évalue la somme et le produit sur les deux feuilles du réseau de chiffres, en conservant quotient et reste ; TRIT détermine le régime et l’orientation ; le transport TPK sélectionne, transporte et actualise l’état avec ses registres de retenue, de frontière et de mémoire. La connexion nonadique prolonge cet état, avec retour de phase et avancement de mémoire. Ses lecteurs agissent sur la même structure discrète du continu.

Le lecteur focal de la source est
\[
\mathcal P_{\mathrm{crit}}:U_{12}^{\mathrm{sign}}\longrightarrow
\mathbb R/2\pi\mathbb Z,\qquad
m\longmapsto\phi_m=\frac{(3m-1)\pi}{18},\quad 1\le m\le12.
\]
Pour évaluer le moment angulaire, on conserve le représentant orienté de cette carte : le remplacer par un autre représentant modulo \(2\pi\) modifierait le moment. La trace uniforme fixe \(w_m=1/12\). La chaîne focale est
\[\begin{aligned}
\text{APP–TRIT–TPK}&\longrightarrow
\text{structure discrète du continu}\\ &\longrightarrow U_{12}^{\mathrm{sign}}
\xrightarrow{\mathcal P_{\mathrm{crit}}}(\phi_m,w_m)
\\ &\longrightarrow u_0\longrightarrow f_\lambda
\longrightarrow\text{retours et renormalisation}.
\end{aligned}\]

L’écriture abrégée conserve comme antécédent l’état complet et ses prolongements ; la famille scalaire est une lecture de cet antécédent. Le lecteur déjà fixé, avec son ordre, son orientation et son secteur uniforme, est pris comme structure de départ explicite. Les conclusions portent sur ce lecteur précis ; l’unicité de ce lecteur parmi toutes les lectures possibles du TPK reste hors de ces preuves.

La coordonnée \(\pi_{\mathrm{HMT}}\) utilisée par la carte appartient aux sorties internes du corpus. Son facteur s’annule exactement lors de la normalisation du deuxième moment. Les valeurs de Feigenbaum interviennent exclusivement dans la reconnaissance ultérieure, jamais dans la définition de la famille ni dans la certification des racines.


\subsection{Profil exact et criticité quadratique}
\label{hmtfg:desarrollo:2}

On conserve les notations $s_m=3m-1$, $\kappa_0=s_0$, $k_m=2s_m/\sqrt{899}$ et $f_{\lambda,0}=f_\lambda$ du profil déjà démontré. Son prolongement entier possède la série rationnelle
\[
u_0(z)=\sum_{j\ge1}(-1)^{j+1}
\frac{4^j\sum_m s_m^{2j}}{12\,899^j(2j)!}z^{2j}
=z^2-\frac{715769}{2424603}z^4
+\frac{1353832978}{32695771455}z^6+\cdots.
\]
Ainsi, le registre fini de douze secteurs produit un profil entier, dont la criticité quadratique et l’échelle sont fixées avant l’itération.


\subsection{Robustesse structurelle avec bornes uniformes}
\label{hmtfg:desarrollo:3}

La déformation est celle définie en \eqref{mab:eq:pesos-dodecafasicos-armonizados}--\eqref{mab:eq:familia-feigenbaum-ponderada-armonizada}. On écrit $M_{2,\eta}=\sum_mw_m\phi_m^2$ et $\kappa_\eta=s_\eta=\sqrt{2/M_{2,\eta}}$, en conservant les mêmes poids, phases et profils.
Les notations sont $p_m=p(\phi_m)=12\cos(3\phi_m)-\cos(4\phi_m)$ et $k_{m,\eta}=\kappa_\eta\phi_m$ ; dans les formules pondérées suivantes, $k_m$ abrège $k_{m,\eta}$, tandis que $k_{m,0}$ correspond au secteur uniforme.

\textbf{Théorème de robustesse structurelle.} Pour \(|\eta|\le1/40\), les poids sont positifs et leur somme vaut un. Le profil est pair, entier, satisfait \(u_\eta''(0)=2\) et est strictement croissant sur \((0,1]\). Pour \(0<\lambda\le2/u_\eta(1)\), l’application envoie \([-1,1]\) dans lui-même, possède un unique point critique dans cet intervalle et satisfait
\[
S f_{\lambda,\eta}(x)\le-\frac{640}{47647}<0,\qquad0<|x|\le1.
\]

\textbf{Démonstration.} Les sommes des cosinus \(3\phi_m\) et \(4\phi_m\) s’annulent par cycles de longueurs quatre et trois. Comme \(|p_m|\le13\),
\[
\frac9{160}\le w_m\le\frac{53}{480},\qquad
\frac{27}{40}M_{2,0}\le M_{2,\eta}\le\frac{53}{40}M_{2,0}.
\]
Il en découle que
\[
k_{\max,\eta}^2\le\frac{196000}{24273}<9<\pi^2,\qquad
k_{\min,\eta}^2\ge\frac{640}{47647}.
\]
Pour \(0<x\le1\), tous les sinus \(\sin(k_{m,\eta}x)\) sont positifs, d’où
\[
u_\eta'(x)=\sum_mw_mk_m\sin(k_mx)>0,\qquad
u_\eta'''(x)=-\sum_mw_mk_m^3\sin(k_mx)<0.
\]
Le quotient \(u_\eta'''/u_\eta'\) est l’opposé d’une moyenne positive des \(k_m^2\). L’invariance de la schwarzienne sous les transformations affines de la valeur donne
\[
S f_{\lambda,\eta}
=\frac{u_\eta'''}{u_\eta'}-\frac32
\left(\frac{u_\eta''}{u_\eta'}\right)^2
\le-k_{\min,\eta}^2.
\]
La parité couvre \(x<0\), et l’image de l’intervalle est
\([1-\lambda u_\eta(1),1]\). \(\square\)

La largeur \(1/40\) est un choix permettant d’obtenir des bornes uniformes, distinct d’une constante fondamentale ou d’un seuil optimal. Le secteur original reste \(\eta=0\).


\subsection{Coordonnée holomorphe quadratique globale sur une bande}
\label{hmtfg:desarrollo:4}

\textbf{Théorème de factorisation holomorphe.} Pour \(|\eta|\le1/40\), dans
\[
\mathcal S=\{z\in\mathbb C:|\operatorname{Re}z|<\pi/3\}
\]
il existe une fonction holomorphe, impaire et univalente \(b_\eta\), avec \(b_\eta'(0)=1\), telle que
\[
\boxed{u_\eta(z)=b_\eta(z)^2,\qquad
f_{\lambda,\eta}(z)=1-\lambda b_\eta(z)^2.}
\]

\textbf{Démonstration.} Pour \(z=x+iy\) dans la demi-bande droite,
\[
\operatorname{Re}u_\eta'(z)
=\sum_mw_mk_m\sin(k_mx)\cosh(k_my)>0.
\]
L’intégrale de \(u_\eta'\) sur le segment entre deux points, divisée par la différence des extrémités, a une partie réelle positive. La convexité de la demi-bande y prouve l’injectivité. La parité donne le résultat correspondant à gauche.

De plus,
\[
\operatorname{Im}u_\eta(x+iy)
=\sum_mw_m\sin(k_mx)\sinh(k_my)
\]
a le signe de \(y\) pour \(x>0\). Sur l’axe réel positif, \(u_\eta>0\) ; sur l’axe imaginaire,
\[
u_\eta(iy)=\sum_mw_m(1-\cosh(k_my))
\]
est strictement négatif hors de zéro et décroît strictement avec \(|y|\). Ces séparations, l’injectivité de chaque demi-bande et la parité prouvent que les fibres de \(u_\eta\) dans \(\mathcal S\) sont exactement \(\{z,-z\}\), sauf la fibre de l’origine.

La fonction \(u_\eta(z)/z^2\), prolongée par la valeur un en zéro, est holomorphe et ne possède aucun zéro dans la bande simplement connexe. Elle possède un logarithme holomorphe \(L_\eta\) avec \(L_\eta(0)=0\), qui est pair par unicité de la normalisation. Définissons
\[
b_\eta(z)=z\exp\!\left(\frac12L_\eta(z)\right).
\]
Alors \(b_\eta^2=u_\eta\), \(b_\eta\) est impaire et \(b_\eta'(0)=1\). L’égalité \(b_\eta(z)=b_\eta(w)\) impose \(w=z\) ou \(w=-z\). Le second cas ne donne l’égalité qu’à l’origine. Cela prouve son univalence. \(\square\)

Les douze phases déterminent ainsi une déformation conforme explicite du pli quadratique. La conjugaison dynamique complète conserve cette déformation :
\[
b_\eta\circ f_{\lambda,\eta}\circ b_\eta^{-1}(w)
=b_\eta(1-\lambda w^2),
\]
sur leurs domaines de composition. La factorisation démontrée et une conjugaison à toute la famille quadratique sont des affirmations distinctes.


\subsection{Retour, changement d’échelle et mémoire}
\label{hmtfg:desarrollo:5}

Soient \(a=-f(1)>0\), \(H_a(x)=-ax\), \(J=H_a([-1,1])\). Lorsque \(J\subset[-1,1]\) est admissible pour le retour, \(f^2(J)\subset J\) et que les compositions sont définies,
\[
\mathcal Rf=H_a^{-1}f^2H_a,\qquad
\mathcal Rf(x)=-a^{-1}f(f(-ax)).
\]
Si, de plus, \(f(J)\subset(0,1]\), un unique maximum quadratique central est conservé :
\[
(\mathcal Rf)(0)=1,\qquad
(\mathcal Rf)''(0)=-a f'(1)f''(0)<0.
\]
La composition des schwarziennes donne, aux points réguliers,
\[
S(\mathcal Rf)(x)=a^2
\left[(Sf)(f(-ax))f'(-ax)^2+(Sf)(-ax)\right]<0.
\]
Les inclusions de domaines font partie des hypothèses et sont vérifiées à chaque échelle.

Soit \(v_\eta=(u_\eta|_{[0,1]})^{-1}\). Les branches inverses sont
\[
\beta_\sigma(y)=\sigma v_\eta((1-y)/\lambda),\quad
\sigma\in\{-1,+1\},
\]
avec le registre critique \(\beta_0(1)=0\). Le pas
\[
(x,w)\longmapsto(f(x),w\mathbin{\|}\sigma(x))
\]
s’inverse sur son image en lisant la dernière feuille et en appliquant \(\beta_\sigma\). Le mot conserve les décisions de branche ; la valeur antérieure se reconstruit par l’équation de l’application.

Le retour regroupe deux transitions. À partir de l’extrémité \(x'\) et de ses feuilles \((\sigma_0,\sigma_1)\), on récupère
\[
x=H_a^{-1}\beta_{\sigma_0}
\!\left(\beta_{\sigma_1}(H_ax')\right).
\]
Les points intermédiaires sont également reconstruits. Les registres APP–TRIT–TPK de feuille, de retenue, de route, de frontière et de mémoire sont concaténés en conservant leur identité. Le signe de la branche analytique est un registre supplémentaire : il ne remplace pas ces données par une étiquette binaire.

Si \(C_N\) regroupe les histoires de longueur \(2N\) par couples, les restrictions satisfont
\[
\pi^{\mathrm{ret}}_{N,M}C_N
=C_M\pi^f_{2N,2M},\qquad M\le N.
\]
Les deux membres suppriment les mêmes couples terminaux et conservent les mêmes branches inverses. Le transport compatible des histoires finies est ainsi démontré.


\subsection{Échelles accumulées et dérivée de l’opérateur}
\label{hmtfg:desarrollo:6}

Soient \(c_n=f^{2^n}(0)\), \(c_0=1\) et \(H_n(x)=c_nx\). Tant que les compositions sont admissibles et que \(c_n\ne0\),
\[
\boxed{\mathcal R^nf=H_n^{-1}f^{2^n}H_n},\qquad
(\mathcal R^nf)(1)=\frac{c_{n+1}}{c_n}=-a_n.
\]
La récurrence découle de \(H_nH_{a_n}=H_{n+1}\). La convention \(a_n>0\) exige l’alternance du signe de \(c_n\). À une racine superstable de période \(2^N\), \(c_N=0\) : cette normalisation devient singulière et il faut utiliser les niveaux antérieurs ou une limite contrôlée.

Pour \(z=-ax\), \(w=f(z)\) et une perturbation \(h\), la dérivée complète est
\[
\boxed{
D\mathcal R_f[h](x)=
-\frac{h(w)+f'(w)h(z)}a+
\frac{h(1)}a
\left[\mathcal Rf(x)-xf'(w)f'(z)\right].
}
\]
On dérive en utilisant \(\dot a=-h(1)\), \(\dot z=xh(1)\),
\(\dot w=h(z)+xf'(z)h(1)\) et la variation du dénominateur. Si \(h(0)=0\), on obtient \(D\mathcal R_f[h](0)=0\). La direction paramétrique initiale HMT est \(h=-u_\eta\).

Cette formule inclut la variation d’échelle ; figer \(a\) donnerait un autre opérateur linéarisé.


\subsection{Racines et itinéraires certifiés jusqu’à la période 256}
\label{hmtfg:desarrollo:7}

Soit \(F_n(\lambda)=f_{\lambda,0}^{2^n}(0)\). On calcule
\[
x_0=p_0=0,\qquad x_{j+1}=1-\lambda u_0(x_j),\qquad
p_{j+1}=-u_0(x_j)-\lambda u_0'(x_j)p_j.
\]
Alors \(F_n=x_{2^n}\) et \(F_n'=p_{2^n}\).

Le programme joint utilise des intervalles rationnels avec arrondi entier vers l’extérieur, une racine carrée encadrée par une racine entière et des sinus/cosinus avec reste de Taylor explicite. La borne exacte \(|u_0''|\le2\) permet une évaluation centrée avec reste quadratique.

Les approximations historiques proposent seulement des boîtes de rayon \(10^{-42}\). On certifie des signes opposés à leurs extrémités, \(F_n'\) séparé de zéro dans toute la boîte et l’exclusion des retours de périodes qui divisent proprement \(2^n\). Le théorème des valeurs intermédiaires et la monotonie prouvent l’existence et l’unicité locales ; l’exclusion des diviseurs prouve la période minimale.

\begingroup\small
\begin{center}\begin{tabular}{p{.12\linewidth}p{.17\linewidth}p{.62\linewidth}}
\hline
Niveau & Période minimale & Centre de la boîte, arrondi \\
\hline
1 & 2 & 1.345913990471832792210949 \\
2 & 4 & 1.763379787846910127290794 \\
3 & 8 & 1.850413796950136464530567 \\
4 & 16 & 1.868979756606798125150574 \\
5 & 32 & 1.872955034435659740004205 \\
6 & 64 & 1.873806367467030119412262 \\
7 & 128 & 1.873988695221365362307169 \\
8 & 256 & 1.874027744154450672897008 \\
\hline\end{tabular}\end{center}
\endgroup

La première racine possède l’expression exacte \(\lambda_1=1/u_0(1)\). Les boîtes complètes et leurs vérifications figurent dans le certificat JSON.

Les 502 états de préclôture de ces huit niveaux, tous séparés de zéro, ainsi que leurs mots de signes complets, ont été encadrés. En particulier,
\[
\operatorname{sign}f_{\lambda_n,0}^{2^j}(0)=(-1)^j,\qquad0\le j<n.
\]
Les orientations finies des changements d’échelle antérieurs à la clôture sont établies. Les inclusions d’intervalles restrictifs à toutes les échelles constituent une vérification supplémentaire.

Pour \(\delta_{\mathrm F,n}=(\lambda_{n-1}-\lambda_{n-2})/(\lambda_n-\lambda_{n-1})\), un intervalle décimal extérieur abrégé est
\[
\boxed{\begin{gathered}
4.66921218915020415455609621588966392804<\delta_{\mathrm F,8},\\
\delta_{\mathrm F,8}<4.66921218915020415455609621588966392863.
\end{gathered}}
\]
La largeur de l’intervalle complet est inférieure à \(5.808\cdot10^{-37}\). Elle mesure l’erreur numérique du rapport fini \(\delta_{\mathrm F,8}\), distincte de l’erreur d’approximation de la limite universelle.


\subsection{Persistance analytique des racines déformées}
\label{hmtfg:desarrollo:8}

Pour chaque niveau certifié, \(F_n'(\lambda_n)\ne0\). Le théorème des fonctions implicites produit une branche analytique locale \(\lambda_n(\eta)\) avec
\[
f_{\lambda_n(\eta),\eta}^{2^n}(0)=0.
\]
Par continuité, les retours propres restent séparés de zéro pour \(\eta\) suffisamment petit ; la période minimale et l’itinéraire sont préservés. L’existence de ces voisinages locaux est démontrée. Leur largeur commune n’est pas encore identifiée à toute la bande structurelle \(1/40\).

Avec \(M=M_{2,\eta}\), la réponse du profil est
\[
v_\eta(x)=\partial_\eta u_\eta(x)
=\sum_m\frac{p_m}{12}(1-\cos(k_{m,\eta}x))
-\frac{M'}{2M}\,x u_\eta'(x).
\]
Le second terme conserve la normalisation du deuxième moment. Si \(q_j=\partial_\eta x_j\) à \(\lambda\) fixé,
\[
q_0=0,\qquad q_{j+1}=-\lambda v_\eta(x_j)-\lambda u_\eta'(x_j)q_j,
\qquad
\boxed{\lambda_n'(\eta)=-q_{2^n}/p_{2^n}}.
\]
La chaîne poids–échelle–réponse–paramètre superstable est ainsi explicitée. Cette non-dégénérescence finie est distincte de la transversalité asymptotique à la feuille stable du point fixe universel.


\subsection{Queues analytiques et propagation de l’erreur}
\label{hmtfg:desarrollo:9}

Pour \(\eta=0\), soient
\[
\mu_{2j}=\frac1{12}\sum_mk_m^{2j},\quad
p_N(z)=\sum_{j=1}^N(-1)^{j+1}\frac{\mu_{2j}}{(2j)!}z^{2j},
\quad\Omega=\frac{70}{\sqrt{899}}.
\]
Si \(q_N=\Omega^2r^2/[(2N+3)(2N+4)]<1\),
\[
\|u_0-p_N\|_r\le
\frac{\mu_{2N+2}r^{2N+2}}{(2N+2)!(1-q_N)}.
\]
L’inégalité \(\mu_{2j+2}\le\Omega^2\mu_{2j}\) majore la queue par une série géométrique. Pour \(0\le s\le2N+2\),
\[
\|(u_0-p_N)^{(s)}\|_r\le
\frac{\mu_{2N+2}r^{2N+2-s}}{(2N+2-s)!}
\left(1-\frac{\Omega^2r^2}{(2N+3-s)(2N+4-s)}\right)^{-1},
\]
lorsque le dénominateur est positif.

Pour transporter les erreurs par un retour, on contrôle également ses domaines : si \(\|f-g\|\le\varepsilon\), \(a_f,a_g\ge a_*>0\), que les compositions restent dans un domaine convexe commun et que l’on y a \(\|f'\|\le L,\ \|g\|\le M\), alors
\[
\|\mathcal Rf-\mathcal Rg\|_r
\le\varepsilon\left[\frac{1+L+L^2r}{a_*}+\frac{M}{a_*^2}\right].
\]
La preuve additionne la perturbation extérieure, la perturbation intérieure, la modification de l’argument due à \(|a_f-a_g|\le\varepsilon\) et la variation du dénominateur. Ces bornes évitent de réutiliser la somme initiale de cosinus comme si tous les profils renormalisés conservaient cette même forme.

% Fragmento integrable. Preambulo anfitrion: amsmath, amssymb, longtable, hyperref.
\section{Profondeur nonadique et persistance de l’itinéraire}
\label{hmtfg:fragmento:profundidad-nonadica-y-persistencia-del-itinerario}

\subsection{Croissance extérieure et résolution intérieure}
\label{hmtfg:prolongacion:2}

La loi de raffinement de l’article XI, de base entière \(B\ge2\), est
\[
L_c(x,y)=(Bx-c,By+c),\qquad 0\le c<B.
\]
Pour un mot \(c_1\ldots c_m\), son registre est
\[
C_m=\sum_{j=1}^m c_jB^{m-j},\qquad
(x_m,y_m)=(B^mx_0-C_m,B^my_0+C_m).
\]
La somme extérieure vaut \(B^mM_0\), avec \(M_0=x_0+y_0\). La lecture normalisée du registre appartient au cylindre
\[
I_m(C_m)=\left[\frac{C_m}{B^m},\frac{C_m+1}{B^m}\right],
\qquad \operatorname{diam}I_m=B^{-m}.
\]
Le couple entier et l’échelle restituent les quotients, les restes et les zéros initiaux du mot. Les extrémités à double développement conservent leur marque de frontière. Ainsi se coordonnent la croissance du support et le raffinement de la lecture sans identifier des histoires distinctes par leur seule image limite.

La composition est exacte :
\[
L_D^{[b]}\circ L_C^{[a]}=L_{B^bC+D}^{[a+b]}.
\]
Par associativité, regrouper une histoire de dix-huit événements en neuf blocs de deux ou en deux blocs de neuf produit le même registre. Au niveau de retour \(2^N\), l’horizon commun est \(9\,2^N\). La preuve est une égalité de composition sur la même histoire ordonnée ; elle conserve ses retenues et ses sous-blocs récupérables.

La restriction des états et le prolongement de leurs lecteurs s’expriment par les limites suivantes :
\[
X_\infty=\varprojlim X_m,\qquad
r^*a=a\circ r,\qquad
r^*e_x=\sum_{r(y)=x}e_y,
\qquad \mathfrak H_{\rm cyl}=\varinjlim\mathfrak H_m.
\]
Les états se restreignent ; les lecteurs se prolongent. L’image réciproque préserve les produits, l’unité et les idempotents. Une famille exactement compatible induit un lecteur sur la limite inductive. Les approximants qui ne sont compatibles qu’à une erreur près exigent en outre une estimation de convergence dans leur complétion. La section suivante fournit cette estimation pour le retour critique.

\subsubsection{Capacité et vacances du registre}

Le lecteur de capacité de XI compare des blocs de \(729\) et de \(1000\) possibilités. Avec \(\rho=\log_{1000}729=\log_{10}9\), il conserve
\[
\mathcal C_k=\lfloor(k+1)\rho\rfloor,\qquad
\mathcal C_k-\mathcal C_{k-1}=1-z_k,
\]
\[
\mathcal C_b-\mathcal C_a+\sum_{k=a+1}^b z_k=b-a,
\qquad T_{\rm cap}(a)=\lceil a/\rho\rceil-1\quad(a\ge1).
\]
Chaque vacance enregistre une transition pendant laquelle la capacité se maintient ; la longueur chronologique est conservée dans le bilan. La phase de capacité et la phase chronologique sont des lecteurs distincts.

Pour stocker un préfixe de \(m\) chiffres en base neuf, une capacité décimale suffisante est
\[
a(m)=\left\lceil\frac{m\rho}{3}\right\rceil.
\]
Le quotient et le reste permettent de convertir l’entier \(C_m\) entre bases, en conservant séparément la profondeur \(m\). On a
\[
9^m\le1000^{a(m)}\le729^{T_{\rm cap}(a(m))+1}.
\]
Ces inégalités se décident également au moyen de puissances entières, sans arrondir les logarithmes. Le budget de lecture de la section suivante se traduit ainsi en capacité du registre en conservant les unités. Ce calcul dimensionne le stockage ; les règles de survie déterminent en outre quelles histoires natives sont admissibles.


\subsection{Budget explicite de profondeur pour chaque retour}
\label{hmtfg:prolongacion:3}

\textbf{Théorème.} Pour \(|\eta|\le1/40\) fixé, normalisons
\[
t=\lambda u_\eta(1)/2\in[0,1],\qquad z=(x+1)/2,
\]
\[
F_{t,\eta}(z)=1-t\frac{u_\eta(2z-1)}{u_\eta(1)},\qquad z\in[0,1].
\]
Pour tout entier \(q\ge1\) et deux paramètres \(t,s\),
\[
\left|F_{t,\eta}^{q}(1/2)-F_{s,\eta}^{q}(1/2)\right|
\le S_q|t-s|,
\qquad S_q=\sum_{j=0}^{q-1}(6\sqrt2)^j<9^q.
\]
Par conséquent, un cylindre de paramètres en base neuf, de profondeur \(m=2^N+r\), enferme la lecture du retour \(q=2^N\) dans un intervalle de diamètre inférieur à \(9^{-r}\).

\textbf{Démonstration.} La fonction \((1-\cos v)/v^2\) décroît pour \(0<v<\pi\) : le signe de sa dérivée est celui de \(v\sin v-2(1-\cos v)\), négatif car \(\tan(v/2)>v/2\). Par le deuxième moment,
\[
u_\eta(1)\ge\frac{1-\cos3}{9}\sum_jw_jk_j^2
=\frac{2(1-\cos3)}9>\frac13.
\]
La dernière inégalité se vérifie, par exemple, en majorant \(\cos3\) par sa somme de Taylor alternée jusqu’au degré huit, qui est inférieure à \(-1/2\). De plus, Cauchy–Schwarz donne \(|u_\eta'(x)|\le\sqrt2\) sur la droite réelle. On en déduit
\[
\operatorname{Lip}_zF_{t,\eta}\le6\sqrt2<9,\qquad
\operatorname{Lip}_tF_{t,\eta}\le1.
\]
Si \(d_j\) est la différence entre les deux orbites, alors \(d_0=0\) et \(d_{j+1}\le6\sqrt2\,d_j+|t-s|\). La récurrence prouve la somme géométrique. Pour \(|t-s|\le9^{-m}\), le diamètre est inférieur à \(9^{q-m}\). La substitution de \(q=2^N\) et de \(m=q+r\) achève la preuve. \(\square\)

La borne contrôle l’incertitude du paramètre pour \(\eta\) fixé. Une implémentation ajoute et propage séparément l’erreur d’évaluation des cosinus et l’erreur d’arrondi ; le certificat rationnel antérieur fournit ces opérations. La borne est suffisante et conservatrice, avec un coût qui croît avec l’horizon.

\textbf{Lecteur limite.} Réalisons \(t\) à partir de \((x_0,y_0)=(1,0)\), \(B=9\), de sorte que
\[
(x_m,y_m)=(9^m-C_m,C_m),\qquad t_m=C_m/9^m.
\]
Les lecteurs \(E_{q,m}=F_{t_m,\eta}^{q}(1/2)\), ramenés à un niveau commun, satisfont
\[
\|E_{q,m+a}-r^*E_{q,m}\|_\infty\le S_q9^{-m}.
\]
Ils convergent uniformément vers un lecteur continu \(E_q\) dans la complétion de l’algèbre cylindrique. La version ensembliste
\[
J_{q,m}(C)=\{F_{t,\eta}^{q}(1/2):t\in I_m(C)\}
\]
est exactement emboîtée et ses diamètres tendent vers zéro. Ce résultat donne le passage à la limite de la \textbf{lecture à horizon fixé}, avec contrôle de profondeur à tout horizon demandé.


\subsection{Mot de doublement et lecteur nonadique de la chronologie}
\label{hmtfg:prolongacion:4}

Les huit itinéraires certifiés dans le travail antérieur obéissent à
\[
W_1=+,\qquad W_{n+1}=W_n\,(-1)^n\,W_n.
\]
Son unique prolongement préfixiel est
\[
\tau_j=(-1)^{v_2(j)},\qquad j\ge1;
\qquad \tau_{2j}=-\tau_j,\quad\tau_{2j+1}=+.
\]
La preuve utilise \(v_2(2^n+r)=v_2(r)\) pour \(0<r<2^n\). C’est une récurrence à toute profondeur, et la concordance avec les huit retours calculés constitue sa vérification finie dans la famille.

Dans l’ordre lexicographique à parité d’une application à maximum critique, le facteur d’orientation précédant la position \(k\) est \(\prod_{j<k}(-\tau_j)\). Le mot \(\tau\) est strictement supérieur à chacun de ses décalages positifs. En effet, le premier désaccord avec un décalage impair se produit en \(k=1\) ou \(2\) ; un décalage pair réduit la comparaison de moitié et double \(k\). Le premier désaccord se situe donc en \(k=2^r\). À cet endroit
\[
\prod_{j<2^r}(-\tau_j)=(-1)^r=\tau_{2^r},
\]
et la différence orientée de signe opposé est positive. Cette démonstration prouve également l’apériodicité.

Le registre chronologique conserve
\[
K=\rho_9(K)+9q_9(K).
\]
Son lecteur dyadique est
\[
\ell_n(B_K)=[\rho_9(K)+9q_9(K)]_{2^n}=[K]_{2^n}.
\]
Sur la carte des histoires où \(\Gamma_9 B_K=B_{K+9}\),
\[
\ell_n(\Gamma_9B_K)=\ell_n(B_K)+9\pmod{2^n},\qquad
\pi_{n+1,n}\ell_{n+1}=\ell_n.
\]
Comme neuf est impair, l’avancement de neuf visite les \(2^n\) positions. Sa conjugaison avec l’avancement unitaire est \(j\mapsto9j\). Pour \(0<j<2^n\),
\[
v_2(9j\bmod2^n)=v_2(j),
\]
de sorte que le rééchantillonnage nonadique conserve également le mot orienté dans chaque cycle certifié.

Ce lecteur utilise le reste \textbf{et} le quotient. Les phases de \(K=1\) et de \(K=10\) coïncident modulo neuf, tandis que leurs signes dyadiques sont opposés. La mémoire distingue les deux histoires. La propriété arithmétique du rééchantillonnage vaut pour tout multiplicateur impair ; le calendrier HMT fixe ici le choix de neuf.


\subsection{Mémoire bilatérale compatible à profondeur infinie}
\label{hmtfg:prolongacion:5}

Soient \(\mathcal H_n=\ell^2(\mathbb Z/2^n\mathbb Z)\),
\[
C_ne_j=e_{j+1},\qquad
J_ne_j=(e_j+e_{j+2^n})/\sqrt2.
\]
Une vérification sur la base donne \(J_n^*J_n=I\) et \(C_{n+1}J_n=J_nC_n\), y compris les termes de bord.

L’article XI fournit les projecteurs
\[
P_0=\frac19\begin{pmatrix}8&-\sqrt8\\-\sqrt8&1\end{pmatrix},\qquad
P_1=\frac19\begin{pmatrix}1&\sqrt8\\\sqrt8&8\end{pmatrix}.
\]
Ils sont orthogonaux et complémentaires. Par conséquent
\[
\mathbb U_n=P_0\otimes I+P_1\otimes C_n
=\begin{pmatrix}T_n&D_n\\D_n&R_n\end{pmatrix}
\]
est unitaire, avec
\[
T_n=(8I+C_n)/9,\quad D_n=\sqrt8(C_n-I)/9,\quad R_n=(I+8C_n)/9.
\]
La naturalité déjà démontrée implique
\[
\mathbb U_{n+1}(I_2\otimes J_n)=(I_2\otimes J_n)\mathbb U_n,
\quad \mathbb U_n^a=P_0\otimes I+P_1\otimes C_n^a
\quad(a\in\mathbb Z).
\]
La limite hilbertienne s’identifie à \(L^2(\mathbb Z_2,\mu_{\rm Haar})\), en envoyant \(e_j^{(n)}\) sur \(2^{n/2}\mathbf1_{j+2^n\mathbb Z_2}\). Les identités se prolongent par densité et bornitude. Pour tout vecteur \(v\) et tout temps entier \(a\),
\[
\boxed{\|T_{\infty,a}v\|^2+\|D_{\infty,a}v\|^2=\|v\|^2,}
\]
\[
\boxed{v=T_{\infty,a}^*(T_{\infty,a}v)+D_{\infty,a}^*(D_{\infty,a}v).}
\]
Ici \(T_{\infty,a}=(8I+C_\infty^a)/9\). La composition bilatérale conserve l’archive entre les pas ; composer seulement \(T_n\) décrit une autre opération.

La permutation \(S_ne_j=e_{9j\bmod2^n}\) satisfait
\[
S_{n+1}J_n=J_nS_n,\qquad S_nC_nS_n^{-1}=C_n^9.
\]
On obtient ainsi à la limite
\[
S_\infty C_\infty S_\infty^{-1}=C_\infty^9,
\quad (I_2\otimes S_\infty)\mathbb U_\infty
(I_2\otimes S_\infty^{-1})=\mathbb U_\infty^9.
\]
La représentation temporelle est fidèle : \(C_\infty^a=C_\infty^b\) impose \(2^n\mid a-b\) pour tout \(n\), d’où \(a=b\). Des préparations particulières peuvent conserver des symétries temporelles. L’identité operatorielle conserve cette distinction.

Le point projectif \(([K]_{2^n})_n\) et un vecteur de \(L^2\) sont de types différents : une histoire ponctuelle induit une mesure de Dirac, tandis que les amplitudes normalisables appartiennent à l’espace de Hilbert. Le résultat construit un lecteur équivariant de la chronologie ; les autres coordonnées natives du TPK accompagnent le registre complet.


\subsection{Existence et séparation des obligations asymptotiques}
\label{hmtfg:prolongacion:6}

\subsubsection{Réalisation du mot infini sur toute la bande}
\label{hmtfg:prolongacion:6.1}

\textbf{Théorème d’existence.} Pour chaque \(|\eta|\le1/40\) il existe au moins un
\[
\lambda_\infty(\eta)\in
\left[\frac1{2u_\eta(1)},\frac2{u_\eta(1)}\right]
\]
tel que
\[
\boxed{\operatorname{sign}f_{\lambda_\infty(\eta),\eta}^{j}(0)
=(-1)^{v_2(j)}\quad\text{pour tout }j\ge1.}
\]
La conclusion réalise l’itinéraire infini de doublement au sein de la famille fixée. Le quantificateur d’existence reste séparé d’une sélection unique et de la loi des distances entre paramètres superstables.

\textbf{Démonstration.} Nous fixons \(\eta\). Nous écrivons \(K_\lambda\) pour le mot de signes de l’orbite critique ; lorsque celle-ci revient pour la première fois à zéro, le mot se termine par ce zéro. Nous utilisons l’ordre unimodal à maximum, dont le facteur avant la comparaison du symbole \(j\) est le produit des signes \(-K_i\), \(i<j\).

À l’extrémité gauche, l’image de tout l’intervalle est dans \([1/2,1]\) ; donc \(K_-=+++\cdots\). À l’extrémité droite, l’orbite est \(0\mapsto1\mapsto-1\mapsto-1\), d’où \(K_+=+---\cdots\). La comparaison des deuxième et troisième symboles, respectivement, donne
\[
K_-<\tau<K_+.
\]

Si \(K_{\lambda_0}\) est infini et distinct de \(\tau\), la première différence se produit à une position finie. La continuité des itérées correspondantes rend localement constant le côté de la comparaison. Il reste à examiner un paramètre superstable de première période \(p\), de mot \(W0\), où \(|W|=p-1\).

Pour \(\lambda\) proche de \(\lambda_0\), soient \(g_\lambda=f_{\lambda,\eta}^p\) et \(a=g_\lambda(0)\). On a \(g_\lambda'(0)=0\), et une borne locale uniforme de la dérivée seconde donne
\[
g_\lambda(x)=a+O(x^2).
\]
Pour \(|a|\) suffisamment petit et non nul, \(g_\lambda\) envoie \([-2|a|,2|a|]\) dans \([a-|a|/2,a+|a|/2]\). Les retours conservent le signe de \(a\), et les \(p-1\) positions intermédiaires conservent \(W\). Les mots voisins sont donc \((W+)^\infty\) ou \((W-)^\infty\) ; les paramètres avec \(a=0\) conservent \(W0\).

Si \(\tau\) diffère de \(W\) avant la position \(p\), cette différence fixe le même côté pour les trois mots. Si elle partage \(W\), posons
\[
s=\tau_p,\qquad e=\prod_{j<p}(-\tau_j),\qquad A=Ws.
\]
Le signe de la comparaison de \(\tau\) avec \(W0\) et avec \((W,-s)^\infty\) est \(es\). Pour la comparer à \(A^\infty\), soit \(q\) la première position où \(\tau_{p+q}\ne\tau_q\). Elle existe par apériodicité. Jusqu’à la position \(p+q-1\), les deux mots coïncident, et le signe d’orientation se factorise en \(O_pO_{q-1}\), avec \(O_p=-es\). La maximalité stricte déjà démontrée donne
\[
O_{q-1}(\tau_q-\tau_{p+q})>0.
\]
Par conséquent, la comparaison de \(\tau\) avec \(A^\infty\) a pour signe \(-O_p=es\). Les trois possibilités locales sont du même côté de \(\tau\).

Si aucun paramètre ne réalisait \(\tau\), les ensembles \(\{\lambda:K_\lambda<\tau\}\) et \(\{\lambda:K_\lambda>\tau\}\) seraient ouverts, disjoints, non vides et couvriraient l’intervalle des paramètres. La connexité de l’intervalle exclut cette partition. Le paramètre annoncé existe donc. \(\square\)

Cet argument de continuité des itinéraires s’applique après fixation du générateur HMT et de sa famille analytique ; ses hypothèses sont vérifiées ici directement. La construction conserve le mot infini sans introduire comme donnée un paramètre de la famille quadratique ni la valeur universelle.

\subsubsection{Marges empêchant de perdre l’itinéraire lors du passage à la limite}
\label{hmtfg:prolongacion:6.2}

Écrivons \(c_p=f_{\lambda,\eta}^p(0)\). Une borne rationnelle uniforme suffisante est
\[
|f'|,|f''|\le16,\qquad
B_p:=16^p\frac{16^p-1}{15}\ge\|(f^p)''\|_{[-1,1]}.
\]
En effet, \(1-\cos v\ge v^2/2-v^4/24\ge v^2/8\) pour \(|v|\le3\), de sorte que \(u_\eta(1)\ge1/4\), \(\lambda\le8\) et les deux dérivées sont majorés par \(16\), en utilisant les moments. La règle de dérivation en chaîne prouve par récurrence la borne \(B_p\).

Le mot \(\tau\) satisfait \(\tau_{2p}=-\tau_p\) pour tout \(p\). Pour chacune de ses réalisations, \(c_{2p}\) et \(c_p\) sont de signes opposés. Comme \((f^p)'(0)=0\), Taylor donne
\[
|c_{2p}-c_p|\le\tfrac12B_p|c_p|^2.
\]
Le membre de gauche est supérieur à \(|c_p|\), et donc
\[
\boxed{|c_p|>2/B_p>0.}
\]
Pour chaque horizon fixé, il existe ainsi une marge uniforme qui empêche une suite de réalisations de perdre son signe en convergeant. Dans les coordonnées compactes \((\eta,b)\), \(b=\lambda u_\eta(1)\in[1/2,2]\), l’ensemble des réalisations de \(\tau\) est fermé et compact, et sa projection sur toute la bande de \(\eta\) est surjective.

Le même argument vaut pour un arbre de préfixes : si tous les préfixes de \(\tau\) sont réalisés, une sous-suite dans l’ensemble compact conserve chaque signe, car les préfixes suffisamment longs comprennent simultanément les positions \(p\) et \(2p\). Le théorème d’existence fournit maintenant la non-vacuité nécessaire à cet argument.

\subsubsection{Intervalles invariants de retour à tous les niveaux}
\label{hmtfg:prolongacion:6.3}

\textbf{Théorème des retours emboîtés.} Pour chacune des réalisations du théorème d’existence de la section \ref{hmtfg:prolongacion:6.1}, soient \(c_j=f^j(0)\) et
\[
J_n=\operatorname{conv}\{c_{2^n},c_{2^{n+1}}\},\qquad n\ge1.
\]
Alors \(0\in\operatorname{int}J_n\), \(J_{n+1}\subset J_n\),
\[
f^{2^n}(J_n)=J_n,
\]
et les \(2^n\) images \(J_n,f(J_n),\ldots,f^{2^n-1}(J_n)\) sont disjointes. Les retours normalisés sont unimodaux à maximum, de valeur critique un et de même mot \(\tau\). Ces intervalles sont des noyaux invariants de retour ; une normalisation exigeant des extrémités périodiques peut élargir le noyau à l’intervalle restrictif maximal correspondant et doit préciser cette étape.

\textbf{Démonstration du premier retour.} Soit \(F\) une application unimodale à maximum, avec \(F(0)=1\) et de mot \(\tau\), et \(d_j=F^j(0)\). Ses premiers signes sont \(+,-,+,+,+,-,+\). L’intervalle \(I=[d_2,1]\) est invariant : \(F(d_2)=d_3>0\), \(F(1)=d_2<0\), et le maximum vaut un. Posons \(J=[d_2,d_4]\). Comme \(F\) décroît sur les arguments positifs et
\[
F(d_3)=d_4>0> d_6=F(d_5),
\]
on a \(d_3<d_5\), et donc \(F(J)=[d_3,1]\).

De plus, \(d_3>d_4\). Pour le vérifier, \(F^2\) est strictement croissante entre les deux points positifs \(d_3,d_4\), car leurs images \(d_4,d_5\) sont positives. Ses valeurs sont \(F^2(d_3)=d_5>0>d_6=F^2(d_4)\), ce qui impose cet ordre. Ainsi \(J\) et \(F(J)\) sont séparés par un intervalle ouvert, et
\[
F^2(J)=F([d_3,1])=[d_2,d_4]=J.
\]
L’application \(F^2|_J\) possède un unique minimum critique : l’image \(F(J)\) reste positive. La conjugaison \(h(x)=d_2x\) inverse l’orientation et normalise le retour en un maximum de valeur un. Son itinéraire est
\[
\operatorname{sign}\frac{d_{2j}}{d_2}
=-\tau_{2j}=\tau_j.
\]
Le même argument peut être répété à toute profondeur.

\textbf{Disjonction par récurrence.} Supposons disjointes les \(p=2^n\) images du parent \(J_n\), et soit \(g=f^p\). Le premier retour appliqué à l’application normalisée produit \(B=J_{n+1}\), \(A=g(B)\), avec \(A\cap B=\varnothing\), \(g(B)=A\), \(g(A)=B\). Pour \(0\le j<p\), une intersection de \(f^j(A)\) et de \(f^j(B)\) produirait, par application de \(f^{p-j}\), une intersection de \(B\) et de \(A\). Ces deux images sont donc disjointes. Les images correspondant à des \(j\) distincts appartiennent à des images disjointes du parent. Les \(2p\) images disjointes et le retour \(f^{2p}(B)=B\) sont ainsi démontrés. La formule des extrémités résulte de la composition des normalisations \(x\mapsto c_{2^n}x\). \(\square\)

Les ensembles compacts
\[
\Lambda_n=\bigcup_{0\le j<2^n}f^j(J_n)
\]
sont non vides et emboîtés. Leur intersection \(\Lambda_\infty\) possède un lecteur continu sur l’odomètre dyadique : à chaque point est attribuée l’étiquette de l’unique composante qu’il occupe à chaque niveau. Toute suite compatible d’étiquettes possède une préimage par compacité, et l’avancement par \(f\) ajoute un aux étiquettes. Cette application est une semi-conjugaison surjective. Sa promotion en conjugaison ponctuelle exige de démontrer que les diamètres des composantes tendent vers zéro ; cette promotion métrique reste séparée.

% Fragmento integrable. Preambulo anfitrion: amsmath, amssymb, longtable, hyperref.
\section{Compatibilité des lecteurs et restitution du paramètre}
\label{hmtfg:fragmento:compatibilidad-de-lectores-y-restitucion-del-parametro}

\subsection{Compatibilité des lecteurs, des minima et restitution de l’histoire}
\label{hmtfg:escalado:13}

Les opérations de prolongement et les lecteurs ordonnés conservent l’information requise pour la sélection des racines. Les identités suivantes fixent leurs domaines et leur composition.

\subsubsection{Règle forward restituée}
\label{hmtfg:escalado:13.1}

L’actualisation déterministe part de l’état de frontière R36 avec le caractère \(\chi\) et le reste interne. L’actualisation est
\begin{equation}
\mathcal U_K^\chi=
\operatorname{Upd}_K^\chi\circ
\operatorname{Tra}_K^\chi\circ
\operatorname{Sel}_K^\chi,
\qquad
N_{K+1}=729N_K+d_K^\chi.
\label{hmtfg:escalado:eq:26}
\end{equation}
Le chiffre et l’actualisation dépendent de l’état précédent. Deux sections forward ayant le même état complet R36 coïncident par induction. La sous-relation \(\operatorname{Ext}^{\chi,\rightarrow}\) est le graphe de cette actualisation ; la relation totale \(\operatorname{Ext}\) admet une multiplicité de continuations selon les données et les caractères.

Les prolongements compatibles transportent les lecteurs régionaux et la signature conjointe sur les mêmes préfixes.

Ces règles transportent l’orientation, le reste, le quotient, la retenue, la frontière et la mémoire. Leur application à la cascade doit également conserver la sélection ordonnée du paramètre du lecteur critique.

\subsubsection{Indépendance des constantes antécédentes à l’égard du paramètre libre}
\label{hmtfg:escalado:13.2}

Soit \(B\) l’état HMT qui fournit la roue et ses constantes antécédentes. Notons \(\mathcal C(B)\) l’ensemble de ces lectures et \(u_B\) le profil critique. Dans le domaine du lecteur critique, la construction ultérieure est
\begin{equation}
(B,\lambda)\longmapsto f_{B,\lambda}(x)=1-\lambda u_B(x),
\qquad 0<\lambda\le 2/u_B(1).
\label{hmtfg:escalado:eq:27}
\end{equation}
Pour un même \(B\), la lecture des constantes antérieures se factorise par la projection sur la première composante :
\begin{equation}
\widetilde{\mathcal C}(B,\lambda)
=\mathcal C\circ\operatorname{pr}_B(B,\lambda)
=\mathcal C(B).
\label{hmtfg:escalado:eq:28}
\end{equation}
Par conséquent, deux valeurs distinctes de \(\lambda\) conservent les mêmes lectures antécédentes. L’égalité est une propriété de cette construction causale : le profil est fixé avant de faire varier \(\lambda\), et le facteur \(\pi_{\mathrm{HMT}}\) se simplifie dans sa normalisation quadratique.

\textbf{Lemme.} Supposons qu’une condition \(A\) soit évaluée uniquement sur ces lectures antécédentes. Pour tout \(B\) fixé,
\begin{equation}
\{\lambda\in\Lambda:A(\widetilde{\mathcal C}(B,\lambda))\}
=
\begin{cases}
\Lambda,&A(\mathcal C(B))\text{ est satisfaite},\\
\varnothing,&A(\mathcal C(B))\text{ échoue}.
\end{cases}
\label{hmtfg:escalado:eq:29}
\end{equation}
\textbf{Démonstration.} La substitution de \eqref{hmtfg:escalado:eq:28} rend le prédicat indépendant de \(\lambda\). Par conséquent, tous les paramètres du domaine le satisfont, ou aucun. Cela inclut une comparaison des constantes antécédentes avec des intervalles métrologiques.

L’unicité de la section \(B\) et la restriction ultérieure \(f_{B,\lambda}^{2^n}(0)=0\) produisent une fibre de solutions dans \(\lambda\). La règle de première racine agit sur cette fibre par le retour et son itinéraire. La sélection et son prolongement sont démontrés dans les résultats suivants.

Ce résultat maintient l’ordre HMT : la métrologie confronte des sorties préalablement construites et la sélection de la racine se démontre dans le lecteur qui la produit.

\subsubsection{Lemme de transport ordonné du sélecteur}
\label{hmtfg:escalado:13.3}

Soit \(P_n\) un ensemble ordonné de candidats généalogiques de minimum \(p_n^*\). Soit \(Z_n\) l’ensemble initial des racines candidates du corpus, muni de l’ordre réel, et soit \(L_n:P_n\to Z_n\) un lecteur monotone. Son image est dite \textbf{coinitiale} lorsque
\begin{equation}
\forall z\in Z_n\quad\exists p\in P_n:
L_n(p)\le z.
\label{hmtfg:escalado:eq:30}
\end{equation}
Alors
\begin{equation}
\boxed{L_n(p_n^*)=\min Z_n
\quad\Longleftrightarrow\quad
L_n(P_n)\text{ est coinitial dans }Z_n.}
\label{hmtfg:escalado:eq:31}
\end{equation}
\textbf{Démonstration.} Si l’image du minimum est le minimum global, on prend \(p=p_n^*\) dans \eqref{hmtfg:escalado:eq:30}. Réciproquement, étant donné \(z\in Z_n\), la coinitialité fournit \(p\) avec \(L_n(p)\le z\). La monotonie donne \(L_n(p_n^*)\le L_n(p)\le z\). Puisque \(L_n(p_n^*)\in Z_n\), c’est le minimum global. La surjectivité de \(L_n\) constitue une condition suffisante, plus forte que \eqref{hmtfg:escalado:eq:30}.

L’isomorphisme ordonné de la carte du continu satisfait la surjectivité du lecteur. L’application concrète au retour et à son ensemble complet de zéros est démontrée dans la section de transport ordonné ; la classification et l’isolement déterminent ensuite sa continuation dynamique.

\subsubsection{Sélection minimale par survie à profondeur arbitraire}
\label{hmtfg:escalado:13.6}

La composition avec la survie produit également un sélecteur global précis. On conserve la famille uniforme \(\eta=0\), construite dans la section~\ref{hmtfg:escalado:1}, et l’on utilise les théorèmes de réalisation de l’itinéraire infini et de conservation de ses signes. Soient
\[
I=\left[\frac1{2u_0(1)},\frac2{u_0(1)}\right],\qquad
c_j(\lambda)=f_\lambda^j(0),\qquad
\tau_j=(-1)^{v_2(j)}.
\]
L’ensemble
\[
\Lambda_\tau=\{\lambda\in I:\operatorname{sign}c_j(\lambda)=\tau_j
\text{ pour tout }j\ge1\}
\]
est compact et non vide. Le contrôle des dérivées de cette source fournit
\begin{equation}
B_j=\frac{16^j(16^j-1)}{15},\qquad
|c_j(\lambda)|>\frac2{B_j}\quad(\lambda\in\Lambda_\tau).
\label{hmtfg:escalado:eq:32}
\end{equation}
Définissons
\begin{equation}
\varepsilon_j=\frac1{B_j},\qquad
A_N=\{\lambda\in I:\tau_jc_j(\lambda)\ge\varepsilon_j,\
1\le j\le N\}.
\label{hmtfg:escalado:eq:33}
\end{equation}
Chaque \(A_N\) est compact par continuité des itérées et non vide parce qu’il contient \(\Lambda_\tau\). De plus,
\begin{equation}
A_{N+1}\subseteq A_N,\qquad
\bigcap_{N\ge1}A_N=\Lambda_\tau.
\label{hmtfg:escalado:eq:34}
\end{equation}
L’inclusion de droite à gauche découle de \eqref{hmtfg:escalado:eq:32}. L’inclusion inverse découle des marges strictement positives de \eqref{hmtfg:escalado:eq:33}, qui fixent chaque signe.

\textbf{Théorème du minimum survivant.} Les minima \(a_N=\min A_N\) satisfont
\begin{equation}
\boxed{a_N\uparrow a_*=\min\Lambda_\tau.}
\label{hmtfg:escalado:eq:35}
\end{equation}
\textbf{Démonstration.} L’inclusion \eqref{hmtfg:escalado:eq:34} rend croissante la suite des minima, bornée dans \(I\) ; elle admet donc une limite \(a\). Pour chaque \(r\), tous les termes \(a_N\), \(N\ge r\), appartiennent au fermé \(A_r\) ; ainsi \(a\in A_r\). Par \eqref{hmtfg:escalado:eq:34}, \(a\in\Lambda_\tau\). Pour tout \(\lambda\in\Lambda_\tau\), on a \(a_N\le\lambda\) à chaque niveau et, en passant à la limite, \(a\le\lambda\). Cela prouve \eqref{hmtfg:escalado:eq:35}.

Dans le lecteur cylindrique HMT, l’objet \(a_*\) possède des préfixes compatibles à toute profondeur, avec conservation de l’orientation, du reste et de la mémoire. Son choix utilise l’ordre de la lecture paramétrique et la survie complète. La compacité démontre l’existence et la convergence de ces minima.

Les paramètres \(a_N\) résolvent des contraintes de signes avec marge. Leurs frontières peuvent satisfaire \(\tau_jc_j=\varepsilon_j\), tandis que les racines superstables satisfont \(c_{2^n}=0\). Ainsi, \eqref{hmtfg:escalado:eq:35} établit un sélecteur global de réalisation infinie et préserve, comme objet distinct, le sélecteur de premières racines du corpus.

\subsubsection{Restitution du paramètre à partir de l’histoire critique complète}
\label{hmtfg:escalado:13.7}

Le deuxième terme de l’histoire critique fournit un lecteur inverse exact :
\begin{equation}
c_1(\lambda)=1,\qquad
c_2(\lambda)=1-\lambda u_0(1),\qquad
\boxed{\lambda=\frac{1-c_2(\lambda)}{u_0(1)}.}
\label{hmtfg:escalado:eq:36}
\end{equation}
La positivité \(u_0(1)>0\), démontrée pour le profil, rend la division valide. Par conséquent, deux histoires critiques complètes identiques ont le même paramètre. Le mot de signes \(\tau\) constitue une lecture réduite de cette histoire ; il conserve sa combinatoire et omet la grandeur de \(c_2\).

Les formules \eqref{hmtfg:escalado:eq:35}–\eqref{hmtfg:escalado:eq:36} unissent la sélection par survie et la restitution du paramètre à partir de l’histoire critique. L’égalité de sa limite avec le croisement stable est démontrée par première sortie, et l’identification des centres par le transport analytique de leurs classes hybrides.

% Fragmento integrable. Preambulo anfitrion: amsmath, amssymb, longtable, hyperref.
\section{Transport ordonné et représentation de Jacobi}
\label{hmtfg:fragmento:transporte-ordenado-y-representacion-de-jacobi}

\subsection{Transport complet depuis le continu et réalisation spectrale du profil}
\label{hmtfg:escalado:14}

Le transport ordonné du continu et la représentation spectrale du lecteur se composent avant la sélection paramétrique.

\subsubsection{Couverture de toutes les racines et transport de leur minimum}
\label{hmtfg:escalado:14.1}

La construction ordonnée du continu développée dans les travaux antérieurs produit la lecture archimédienne surjective à partir de cylindres compatibles et de leur quotient ordonné. Dans son univers déclaré \(\mathfrak G\), elle démontre l’isomorphisme de corps ordonnés complets
\begin{equation}
\iota:\mathbb R_{\rm HMT}^{\mathfrak G}
   \xrightarrow{\cong}\mathbb R^{\mathfrak G}.
\label{hmtfg:escalado:eq:37}
\end{equation}
Le nom \(\eta\) du document de référence est remplacé ici par \(\iota\), afin de le distinguer du paramètre de pondération, fixé à zéro. La seconde projection conserve l’incidence, la frontière et la mémoire sur le même antécédent ; le quotient d’évaluation et l’histoire complète conservent leurs domaines respectifs.

Construisons dans le corps HMT, avant de sélectionner une racine,
\[
\cos_{\rm H}(t)=\sum_{k=0}^{\infty}\frac{(-1)^kt^{2k}}{(2k)!},
\qquad
u_{\rm H}(x)=\frac1{12}\sum_{m=1}^{12}
\left[1-\cos_{\rm H}\!\left(\frac{2(3m-1)x}{\sqrt[\rm H]{899}}\right)\right],
\]
\begin{equation}
F_a(x)=1-a\,u_{\rm H}(x),\qquad S_n^{\rm H}(a)=F_a^{\,2^n}(0).
\label{hmtfg:escalado:eq:38}
\end{equation}
Les entiers et le deuxième moment sont ceux déjà dérivés du lecteur dodécaphasique. La racine carrée est la racine positive du corps ordonné. La série convergente réalise analytiquement ce lecteur ; elle conserve le profil et le coefficient normalisateur de la section~\ref{hmtfg:escalado:1}.

\textbf{Théorème de transport du retour.} Pour tout paramètre du domaine interne de \eqref{hmtfg:escalado:eq:38},
\begin{equation}
\iota(F_a(x))=f_{\iota(a)}(\iota(x)),\qquad
\iota(S_n^{\rm H}(a))=S_n(\iota(a)).
\label{hmtfg:escalado:eq:39}
\end{equation}
\textbf{Démonstration.} L’isomorphisme ordonné fixe les rationnels et préserve la racine positive. Il préserve les limites convergentes : chaque voisinage de rayon rationnel positif est transporté vers un autre de même rayon. Appliqué aux sommes partielles de \eqref{hmtfg:escalado:eq:38}, il transporte le cosinus et \(u_{\rm H}\). La première égalité résulte de la conservation de la somme et du produit ; la seconde, d’une récurrence sur le nombre d’itérations.

Un centre antérieur \(a_{n-1}\) étant fixé, définissons \(Z_n^{\rm H}\) par un retour nul, une période critique exacte \(2^n\) et \(a>a_{n-1}\), dans le domaine paramétrique déclaré. Soit \(Z_n\) l’ensemble défini par les mêmes conditions pour \(f_\lambda\), à droite de \(\iota(a_{n-1})\). Alors
\begin{equation}
\boxed{\iota[Z_n^{\rm H}]=Z_n,\qquad
\iota(\min Z_n^{\rm H})=\min Z_n.}
\label{hmtfg:escalado:eq:40}
\end{equation}
La seconde égalité est affirmée lorsque le minimum existe ; la première préserve également l’absence de candidats.

\textbf{Démonstration.} \eqref{hmtfg:escalado:eq:39} préserve le retour nul et les retours antérieurs qui déterminent la période exacte. L’ordre conserve la position par rapport au centre précédent. La surjectivité de \eqref{hmtfg:escalado:eq:37}, appliquée à tout paramètre racine, prouve la couverture inverse. Si \(a_*\) est le minimum, tout \(z\in Z_n\) est de la forme \(\iota(a)\), avec \(a\in Z_n^{\rm H}\), d’où \(\iota(a_*)\le z\). L’affirmation réciproque utilise \(\iota^{-1}\).

La surjectivité couvre tous les paramètres de la carte, y compris ceux dont les racines n’ont pas été isolées par un calcul fini. L’isomorphisme conserve l’univers d’interprétation déclaré et l’ordre de sélection.

\subsubsection{Représentation de Jacobi du profil dodécaphasique complet}
\label{hmtfg:escalado:14.2}

Définissons la mesure atomique du lecteur,
\begin{equation}
s_m=\frac{4(3m-1)^2}{899},\qquad
\nu_0=\frac1{12}\sum_{m=1}^{12}\delta_{s_m},\qquad
M_r=\int s^r\,d\nu_0(s).
\label{hmtfg:escalado:eq:41}
\end{equation}
Les douze nœuds sont distincts et positifs. Pour un polynôme \(p\ne0\) de degré inférieur à \(r\le12\),
\begin{equation}
\sum_{i,j=0}^{r-1}p_iM_{i+j}p_j
=\frac1{12}\sum_{m=1}^{12}p(s_m)^2>0.
\label{hmtfg:escalado:eq:42}
\end{equation}
Un polynôme de degré inférieur à douze ne s’annule pas aux douze nœuds. Au degré douze apparaît \(\prod_m(s-s_m)\), de norme nulle.

L’orthogonalisation par rapport à cette mesure produit des polynômes orthonormaux \(\psi_0,\ldots,\psi_{11}\), avec \(\psi_0=1\). Soient
\[
Q_{mj}=\psi_j(s_m)/\sqrt{12},\quad
D=\operatorname{diag}(s_m),\quad J_{12}=Q^{\mathsf T}DQ.
\]
Les relations suivantes sont satisfaites
\begin{equation}
Q^{\mathsf T}Q=I,\quad DQ=QJ_{12},\qquad
\boxed{u_0(x)=e_0^{\mathsf T}
\bigl[I-\cos(x\sqrt{J_{12}})\bigr]e_0.}
\label{hmtfg:escalado:eq:43}
\end{equation}
\textbf{Démonstration.} L’orthogonalité découle de \eqref{hmtfg:escalado:eq:42}. La multiplication par \(s\) produit une récurrence à trois termes : les éléments séparés de plus d’une position s’annulent par orthogonalité et degré. Fixer des coefficients dominants positifs rend les sous-diagonales positives. Par calcul fonctionnel,
\(\cos(x\sqrt{J_{12}})=Q^{\mathsf T}\cos(x\sqrt D)Q\).
Le vecteur \(Qe_0\) est constant, de valeur \(1/\sqrt{12}\) ; sa substitution restitue la somme de la section~\ref{hmtfg:escalado:1}.

L’article sur les grassmanniennes développe ce mécanisme pour les mesures marquées et les raffinements. Il est appliqué ici à \eqref{hmtfg:escalado:eq:41} ; le rang douze appartient à ce lecteur. Les marques généalogiques sont conservées avec le lecteur, conformément à la fidélité marquée du document de référence. \(J_{12}\) représente le profil analytique ; l’état complet conserve en outre les routes, l’orientation et la mémoire.

Le vérificateur de cette extension contrôle, avec des fractions exactes, les douze degrés de norme positive, l’annulation au degré douze, l’entrelaceur \(VT=DV\) dans la base monique et l’autoadjonction pondérée \(HT=T^{\mathsf T}H\). Il vérifie les moments de degré zéro à trente ; l’identité à tout degré découle de l’entrelaceur. Ses premiers coefficients sont
\begin{equation}
a_0=2,\qquad \beta_1=\frac{2493348}{808201}.
\label{hmtfg:escalado:eq:44}
\end{equation}

\subsubsection{Sensibilité orientée du retour}
\label{hmtfg:escalado:14.3}

Pour un retour exact de période \(p=2^n\), soient
\[
c_j=f_\lambda^j(0),\quad q_j=\partial_\lambda c_j,\quad
D_j=\prod_{k=1}^j f_\lambda'(c_k),\quad D_0=1.
\]
Avant le retour critique, les facteurs de \(D_{p-1}\) sont non nuls. La différentiation et le télescopage donnent
\begin{equation}
q_{j+1}=-u_0(c_j)+f_\lambda'(c_j)q_j,\qquad
\boxed{\frac{q_p}{D_{p-1}}
=-\sum_{j=1}^{p-1}\frac{u_0(c_j)}{D_j}.}
\label{hmtfg:escalado:eq:45}
\end{equation}
Dans le mot canonique, \(\operatorname{sign}c_j=(-1)^{v_2(j)}\). Le signe de \(f_\lambda'(c_j)\) est opposé ; par conséquent,
\begin{equation}
\operatorname{sign}D_j=(-1)^{j+\sum_{k=1}^jv_2(k)}
=(-1)^{s_2(j)},
\label{hmtfg:escalado:eq:46}
\end{equation}
en utilisant \(\sum_{k=1}^jv_2(k)=j-s_2(j)\), où \(s_2(j)\) compte les uns binaires. La sensibilité normalisée est
\begin{equation}
\mathcal T_n=-\sum_{j=1}^{2^n-1}(-1)^{s_2(j)}t_j,\qquad
t_j=\frac{u_0(c_j)}{|D_j|}>0.
\label{hmtfg:escalado:eq:47}
\end{equation}
La positivité nodale de \eqref{hmtfg:escalado:eq:42} détermine \(u_0\) ; \eqref{hmtfg:escalado:eq:47} conserve en outre les produits orbitaux et leurs orientations.

La représentation supplémentaire par moments exigerait une mesure positive \(\rho\) sur \([0,1]\) avec \(t_j=\int z^{j-1}\,d\rho\). Sous cette condition,
\begin{equation}
\mathcal T_n=\int_0^1
\frac{1-\prod_{r=0}^{n-1}(1-z^{2^r})}{z}\,d\rho(z)>0
\label{hmtfg:escalado:eq:48}
\end{equation}
pour \(\rho\ne0\), en prolongeant l’intégrande par sa valeur \(1\) en zéro. L’identité s’obtient en développant le produit binaire.

Le contrôle rationnel à la racine de période quatre précédemment certifiée obtient
\[
t_1\approx0.40425224267600139730,\quad
t_2\approx0.04925249167432646403,\quad
t_3\approx0.15740870098294833737
\]
et l’intervalle extérieur rigoureux
\begin{equation}
\begin{gathered}
0.296096033367379523967764444238832345360107<\mathcal T_2,\\
\mathcal T_2<0.296096033367379523967764444238832345360112.
\end{gathered}
\label{hmtfg:escalado:eq:49}
\end{equation}
Il certifie également \(t_3>t_2\). Une mesure positive sur \([0,1]\) imposerait \(t_3\le t_2\) ; ce relèvement précis des poids orbitaux en moments positifs est exclu. Le signe positif de \eqref{hmtfg:escalado:eq:49} reste certifié. Ce contrôle distingue l’échec d’une preuve proposée de l’échec de la propriété que l’on cherchait à prouver.

% Fragmento integrable. Preambulo anfitrion: amsmath, amssymb, longtable, hyperref.
\section{Classification et prolongement du sélecteur des premières racines}
\label{hmtfg:fragmento:clasificacion-y-continuacion-del-selector-de-primeras-raices}

\subsection{Objet et provenance}
\label{hmtfg:clasificacion:1}

Cette annexe conserve la famille issue de la lecture dodécaphasique HMT :

\[
u_0(x)=\frac1{12}\sum_{m=1}^{12}\left(1-\cos\frac{2(3m-1)x}{\sqrt{899}}\right),
\qquad f_\lambda(x)=1-\lambda u_0(x).
\]

La famille provient des antécédents operatoriels et analytiques exposés dans ce développement. L’existence et la compacité de ses réalisations de l’itinéraire infini ont été démontrées dans les sous-sections consacrées à la réalisation du mot et à la conservation de ses signes.

L’argument suivant est propre à ce développement et utilise exclusivement ces propriétés et les opérations de retour de la famille fixée. Il n’utilise comme entrée ni un paramètre d’une autre famille ni la valeur d’une constante universelle. Les feuilles des itinéraires utilisés ici sont les deux branches monotones du lecteur critique ; les autres coordonnées et la mémoire HMT conservent leur résidence dans l’état dont provient le lecteur.

Écrivons
\[
c_j(\lambda)=f_\lambda^j(0),\qquad
\tau_j=(-1)^{v_2(j)},\qquad
W_n=\tau_1\cdots\tau_{2^n-1}.
\]
L’itinéraire d’une orbite critique qui revient pour la première fois à l’origine se termine par le symbole zéro. L’ordre des itinéraires est l’ordre unimodal à maximum : avant la comparaison du symbole d’indice \(j\), le facteur d’orientation est
\[
O_{j-1}(K)=\prod_{i=1}^{j-1}(-K_i).
\]
L’ordre ordinaire des symboles est \(-1<0<1\).


\subsection{Classification avant le mot infini}
\label{hmtfg:clasificacion:2}

\textbf{Théorème de classification.} Soit \(F:I\to I\), avec \(I=[\ell,1]\), \(\ell<0\), une application continue, strictement croissante à gauche de zéro et strictement décroissante à sa droite, dont le maximum satisfait \(F(0)=1\). Supposons que zéro soit de période exacte \(2^n\), avec \(n\ge1\), et que son itinéraire critique \(K\) satisfasse \(K<\tau\). Alors
\[
\boxed{K=W_n0.}
\]
La conclusion porte sur l’itinéraire et permet que des paramètres distincts d’une famille possèdent ce même itinéraire.

\textbf{Démonstration.} Notons \(d_j=F^j(0)\). Pour la période deux, l’itinéraire est \(+0=W_10\). Considérons une période exacte \(p=2^n\ge4\).

On a \(d_2<0\). En effet, \(d_2=0\) donnerait la période deux. Si \(d_2>0\), la monotonie de la branche positive donne
\[
F([d_2,1])=[d_2,d_3]\subset[d_2,1],
\]
et l’orbite critique demeure dans un intervalle strictement positif, incompatible avec le retour à zéro.

Les signes initiaux sont alors \(+,-\). Le facteur d’orientation pour comparer le troisième signe est \(-1\) ; par conséquent, \(K<\tau\) impose \(d_3>0\). Un zéro à cette position produirait la période trois.

On ne peut pas non plus avoir \(d_4<0\). Sous cette hypothèse,
\[
d_2<d_4<0,
\qquad F([d_2,d_4])=[d_3,d_5]\subset(0,1],
\]
car \(d_4=F(d_3)>F(1)=d_2\) et \(d_5=F(d_4)>F(d_2)=d_3\). La branche positive est décroissante, de sorte que
\[
F^2([d_2,d_4])=[d_6,d_4]\subset[d_2,d_4].
\]
Ici \(d_6\ge d_2\), car \(d_5\le1\), et \(d_6<d_4\), car \(d_5>d_3\). Cet intervalle strictement négatif piège les retours pairs de l’orbite critique et exclut le retour à zéro. Donc \(d_4\ge0\). Si \(p=4\), on obtient précisément \(+-+0=W_20\).

Supposons maintenant \(p\ge8\). Alors \(d_4>0\). Les facteurs d’orientation pour les cinquième et sixième signes sont, respectivement, \(-1\) et \(+1\). Comme \(\tau_5=+1\) et \(\tau_6=-1\), l’inégalité \(K<\tau\) impose
\begin{equation}
\operatorname{sign}(d_1,\ldots,d_6)=(+,-,+,+,+,-).
\label{hmtfg:clasificacion:eq:1}
\end{equation}
Les signes nuls à ces positions sont exclus par la période exacte.

Nous construisons l’intervalle de retour
\[
J=[d_2,d_4].
\]
L’inégalité \(F(d_3)=d_4>0>d_6=F(d_5)\), jointe à \(d_3,d_5>0\), implique \(d_3<d_5\) ; donc
\[
F(J)=[d_3,1].
\]
De plus, \(d_4<d_3\). Pour le vérifier, l’image par \(F\) de l’intervalle d’extrémités \(d_3,d_4\) est strictement positive, puisque les images de ses extrémités sont \(d_4,d_5>0\). Sur cet intervalle, \(F^2\) est strictement croissante. Ses valeurs aux extrémités satisfont
\[
F^2(d_3)=d_5>0>d_6=F^2(d_4),
\]
ce qui impose \(d_3>d_4\). Par conséquent,
\begin{equation}
J\cap F(J)=\varnothing,
\qquad F^2(J)=F([d_3,1])=[d_2,d_4]=J.
\label{hmtfg:clasificacion:eq:2}
\end{equation}

L’application \(F^2|_J\) possède un unique minimum en zéro. La conjugaison par \(h(x)=d_2x\), qui inverse l’orientation, définit
\begin{equation}
G=h^{-1}\circ F^2\circ h:
\left[\frac{d_4}{d_2},1\right]\longrightarrow
\left[\frac{d_4}{d_2},1\right].
\label{hmtfg:clasificacion:eq:3}
\end{equation}
Cette application satisfait les mêmes hypothèses de maximum que \(F\), avec \(G(0)=1\), et son point critique a pour période exacte \(p/2\). Son itinéraire est
\begin{equation}
K'_j=-K_{2j}.
\label{hmtfg:clasificacion:eq:4}
\end{equation}
Tous les symboles d’indice impair de \(K\) sont positifs, car les positions paires appartiennent à \(J\) et les positions impaires à \(F(J)\subset(0,1]\).

La transformation \eqref{hmtfg:clasificacion:eq:4} conserve l’ordre de comparaison avec \(\tau\). En effet, \(\tau_{2j-1}=+1\) et \(\tau_{2j}=-\tau_j\). Si \(q\) est la première position où \(K'\) et \(\tau\) diffèrent, la première position de différence entre \(K\) et \(\tau\) est \(2q\), et
\begin{equation}
O_{2q-1}(K)=-O_{q-1}(K'),
\qquad K_{2q}-\tau_{2q}=-(K'_q-\tau_q).
\label{hmtfg:clasificacion:eq:5}
\end{equation}
Le produit qui détermine l’ordre est identique. Les identités restent valides lorsque le premier symbole distinct est le zéro final. Ainsi \(K'<\tau\).

La récurrence appliquée à \(G\) donne \(K'=W_{n-1}0\). Les positions impaires positives et \eqref{hmtfg:clasificacion:eq:4} reconstruisent exactement \(K=W_n0\). Le théorème est démontré. \(\square\)

Le théorème démontre également que chaque application classifiée possède les retours restrictifs de doublement antérieurs à son retour superstable, avec les domaines de \eqref{hmtfg:clasificacion:eq:2}–\eqref{hmtfg:clasificacion:eq:3}. Il ne présuppose pas de monotonie du paramètre d’une famille.

\textbf{Normalisation symétrique de la famille fixée.} Lorsque \(F\) est paire et définie sur \([-1,1]\), les inégalités précédentes donnent en outre
\[
F(d_4)=d_5>d_3=F(d_2)=F(-d_2).
\]
Comme les deux arguments \(d_4,-d_2\) sont positifs et que \(F\) décroît sur cette branche,
\begin{equation}
0<d_4<-d_2<1.
\label{hmtfg:clasificacion:eq:5a}
\end{equation}
Donc \(F^2([d_2,-d_2])=[d_2,d_4]\), et la même conjugaison \(h(x)=d_2x\) définit sur tout \([-1,1]\) une application paire à maximum dont l’image est \([d_4/d_2,1]\subset[-1,1]\). Elle coïncide exactement avec l’opérateur \(RF=-F(F(-ax))/a\), \(a=-F(1)\), utilisé dans le corpus. Cette extension symétrique conserve le retour critique, sa période et son itinéraire.


\subsection{Premier paramètre d’itinéraire infini}
\label{hmtfg:clasificacion:3}

Nous travaillons sur l’intervalle compact déjà utilisé dans la démonstration d’existence,
\[
L=\left[\frac1{2u_0(1)},\frac2{u_0(1)}\right].
\]
Soit
\[
\Lambda_\tau=\{\lambda\in L:
\operatorname{sign}c_j(\lambda)=\tau_j\text{ pour tout }j\ge1\}.
\]
D’après les démonstrations d’existence et de séparation des signes du prolongement, cet ensemble est non vide et compact. Il existe donc
\begin{equation}
a=\min\Lambda_\tau.
\label{hmtfg:clasificacion:eq:6}
\end{equation}

\textbf{Lemme de comparaison.} Pour tout \(\lambda\in[1/(2u_0(1)),a)\), on a \(K_\lambda<\tau\).

\textbf{Démonstration.} À l’extrémité gauche, l’itinéraire est \(+^\infty<\tau\). L’argument de séparation de la réalisation du mot infini démontre que les ensembles de comparaison stricte avec \(\tau\) sont ouverts : dans un itinéraire infini distinct de \(\tau\), la première différence finie persiste par continuité ; à un paramètre superstable, les deux mots périodiques voisins et le mot à zéro final se situent du même côté de \(\tau\), par sa maximalité stricte. L’intervalle antérieur à \(a\) ne contient, par définition, aucun paramètre d’itinéraire \(\tau\). Sa connexité maintient la comparaison initiale sur tout cet intervalle. \(\square\)


\subsection{Existence de chaque première nouvelle racine}
\label{hmtfg:clasificacion:4}

Nous conservons le sélecteur original. La première racine est
\[
\lambda_1=\frac1{u_0(1)},
\]
et, pour \(n\ge2\), \(\lambda_n\) est la première racine de \(c_{2^n}\) située à droite de \(\lambda_{n-1}\) qui satisfait
\begin{equation}
c_{2^j}(\lambda_n)\ne0\quad(1\le j<n).
\label{hmtfg:clasificacion:eq:7}
\end{equation}
Comme toute première période divisant \(2^n\) est une puissance de deux, \eqref{hmtfg:clasificacion:eq:7} équivaut à exiger une période critique exacte \(2^n\).

\textbf{Lemme d’existence.} Toutes les racines de ce sélecteur existent et satisfont
\begin{equation}
\lambda_1<\lambda_2<\cdots<\lambda_n<a.
\label{hmtfg:clasificacion:eq:8}
\end{equation}

\textbf{Démonstration.} Le signe \(c_2(a)<0\) implique \(\lambda_1<a\). Supposons \(\lambda_{n-1}<a\) construite, et posons \(p=2^{n-1}\). La fonction analytique \(c_p(\lambda)\) n’est pas identiquement nulle, puisque \(c_p(a)\ne0\). Son ensemble de zéros dans \([\lambda_{n-1},a]\) est fini, non vide et ne contient pas \(a\). Soit \(z<a\) son dernier zéro.

Sur \((z,a]\), le signe de \(c_p\) est constant et égal à \(s=\tau_p\). Pour \(g_\lambda=f_\lambda^p\), on a
\[
g_\lambda(0)=c_p(\lambda),\qquad g'_\lambda(0)=0.
\]
Une borne locale uniforme de la dérivée seconde donne
\begin{equation}
c_{2p}(\lambda)
=g_\lambda(c_p(\lambda))
=c_p(\lambda)+O(c_p(\lambda)^2).
\label{hmtfg:clasificacion:eq:9}
\end{equation}
Lorsque \(\lambda\) tend vers \(z\) par la droite, le terme du premier ordre domine ; donc \(c_{2p}(\lambda)\) a pour signe \(s\). En \(a\), son signe est
\[
\tau_{2p}=-\tau_p=-s.
\]
Le théorème des valeurs intermédiaires fournit un zéro \(\beta\in(z,a)\) de \(c_{2p}\). Comme \(c_p\ne0\) sur cet intervalle, la période critique en \(\beta\) est exactement \(2p\).

L’analyticité et \(c_{2p}(a)\ne0\) impliquent que les zéros de \(c_{2p}\) dans \([\lambda_{n-1},a]\) sont en nombre fini. Parmi ses zéros de période exacte \(2p\) situés à droite de \(\lambda_{n-1}\), il en existe donc un premier. C’est le sélecteur original \(\lambda_n\), et il satisfait \(\lambda_n\le\beta<a\). \(\square\)

L’utilisation du dernier zéro auxiliaire \(z\) appartient uniquement à la preuve d’existence. Le paramètre sélectionné reste la \textbf{première nouvelle racine} \(\lambda_n\).


\subsection{Convergence du sélecteur global vers le premier itinéraire infini}
\label{hmtfg:clasificacion:5}

\textbf{Théorème de prolongement global.} Pour la famille HMT fixée et le sélecteur original, tous les itinéraires sélectionnés sont canoniques et
\begin{equation}
\boxed{K_{\lambda_n}=W_n0,\qquad
\lambda_n\nearrow\min\Lambda_\tau.}
\label{hmtfg:clasificacion:eq:10}
\end{equation}

\textbf{Démonstration.} Les lemmes de comparaison et d’existence placent chaque paramètre sélectionné au-dessous de \(a\), avec un itinéraire inférieur à \(\tau\). Le théorème de classification identifie cet itinéraire à \(W_n0\). La suite est croissante et majorée par \(a\) ; soit \(b\le a\) sa limite.

Fixons \(p\ge1\). Pour \(n\) suffisamment grand, \(2p<2^n\). Alors les signes de \(c_p(\lambda_n)\) et de \(c_{2p}(\lambda_n)\) coïncident avec \(\tau_p\) et \(\tau_{2p}=-\tau_p\), respectivement. La borne uniforme des dérivées secondes,
\[
\|(f_\lambda^p)''\|\le B_p,
\qquad B_p=16^p\frac{16^p-1}{15},
\]
et l’identité \((f_\lambda^p)'(0)=0\) impliquent
\[
|c_{2p}(\lambda_n)-c_p(\lambda_n)|
\le\tfrac12B_p|c_p(\lambda_n)|^2.
\]
Comme les deux termes du membre de gauche sont de signes opposés,
\begin{equation}
|c_p(\lambda_n)|>2/B_p.
\label{hmtfg:clasificacion:eq:11}
\end{equation}
Par continuité, \(c_p(b)\) conserve le signe \(\tau_p\) et satisfait \(|c_p(b)|\ge2/B_p>0\). Cela vaut pour chaque \(p\), de sorte que \(b\in\Lambda_\tau\). La minimalité de \(a\) donne \(b\ge a\), et donc \(b=a\). \(\square\)

La preuve complète conserve le sélecteur global et produit ses itinéraires et sa limite. L’information des huit premiers niveaux est compatible avec le résultat, mais la démonstration de \eqref{hmtfg:clasificacion:eq:10} n’extrapole pas ces huit niveaux : elle utilise la récurrence, des intervalles invariants et la séparation uniforme de chaque signe.


\subsection{Contrôles logiques et falsificateurs}
\label{hmtfg:clasificacion:7}

\par\noindent\textbf{1.}\enspace Le théorème de classification exige un maximum unimodal strict et un intervalle invariant. Si l’une des deux branches perd sa monotonie, les intervalles \eqref{hmtfg:clasificacion:eq:2} peuvent cesser d’être restrictifs et la classification doit être vérifiée à nouveau.
\par\noindent\textbf{2.}\enspace La condition \(K<\tau\) est essentielle. Au-delà du mot limite, il peut exister des centres de période \(2^n\) avec d’autres combinatoires ; le théorème ne les identifie pas à \(W_n0\).
\par\noindent\textbf{3.}\enspace L’analyticité en le paramètre assure l’isolement et la finitude des zéros sur le compact. Pour une famille seulement continue, l’existence d’une première racine strictement postérieure exige un argument supplémentaire.
\par\noindent\textbf{4.}\enspace Une limite de racines critiques peut, en général, perdre des signes en atteignant un retour de période inférieure. L’inégalité \eqref{hmtfg:clasificacion:eq:11} exclut précisément ce phénomène pour cette suite et chaque horizon fixé.
\par\noindent\textbf{5.}\enspace La convergence des premières racines vers le premier paramètre d’itinéraire \(\tau\) n’équivaut pas à l’unicité de tous les paramètres de cet itinéraire et ne détermine pas à elle seule un taux métrique. Le croisement transverse et l’identification de la branche apportent ces conclusions ultérieures.

% Fragmento integrable. Preambulo anfitrion: amsmath, amssymb, longtable, hyperref.
\section{Entrée transverse et changement d’échelle local}
\label{hmtfg:fragmento:entrada-transversal-y-escalado-local}

\subsection{Provenance et résultat}
\label{hmtfg:escalado:1}

APP évalue la somme et le produit sur les deux feuilles du réseau, en conservant quotient et reste. TRIT détermine le régime et l’orientation. Le transport TPK sélectionne, transporte et actualise l’état avec ses registres de retenue, de frontière et de mémoire. Ses lecteurs agissent sur la même structure discrète du continu. Le prolongement \texttt{w6→w12→w18→w24→w30→R36→G9} conserve cette provenance et distingue le retour de phase de l’incrément de mémoire.

Le lecteur critique de la roue dodécaphasique, construit dans les antécédents operatoriels et transporté par le raffinement compatible, produit
\[
\phi_j=\frac{(3j-1)\pi_{\mathrm{HMT}}}{18},\qquad
w_j=\frac1{12},\qquad
k_j=\sqrt{\frac2{\sum_{\ell=1}^{12}w_\ell\phi_\ell^2}}\,\phi_j
=\frac{2(3j-1)}{\sqrt{899}}.
\]
La coordonnée \(\pi_{\mathrm{HMT}}\) conserve son origine APP–TRIT–TPK ; son facteur commun s’annule dans ce deuxième moment. Le lecteur résultant est
\begin{equation}
u_0(x)=\frac1{12}\sum_{j=1}^{12}(1-\cos k_jx),\qquad
f_\lambda(x)=1-\lambda u_0(x).
\label{hmtfg:escalado:eq:1}
\end{equation}
Le secteur focal est le secteur uniforme, \(\eta=0\). L’extension antérieure conserve séparément la famille pondérée \(|\eta|\le1/40\). Les nouvelles bornes numériques correspondent à \eqref{hmtfg:escalado:eq:1}.

La famille est paire, entière, avec \(f_\lambda(0)=1\), un point critique quadratique et un deuxième moment \(\sum_jw_jk_j^2=2\). Le paramètre \(\lambda\) parcourt une famille postérieure au générateur HMT. La constante cible de Feigenbaum reste extérieure aux entrées de \eqref{hmtfg:escalado:eq:1}.

\textbf{Résultat principal.} Il existe un unique
\begin{equation}
\lambda_*\in[1.874038,1.874039]
\label{hmtfg:escalado:eq:2}
\end{equation}
pour lequel le quatrième retour normalisé de \eqref{hmtfg:escalado:eq:1} appartient à la feuille stable locale du point fixe réel de doublement. Le croisement de la famille avec cette feuille est transverse, avec une séparation quantifiée. Pour tout \(n\ge0\),
\begin{equation}
\boxed{\|R^{4+n}f_{\lambda_*}-g\|_{\mathcal A}
<0.001666\,(0.83996)^n.}
\label{hmtfg:escalado:eq:3}
\end{equation}
La cascade superstable locale associée à ce croisement possède des paramètres \(\mu_j\) et des constantes \(C\ne0\), \(0<\theta<1\), tels que
\begin{equation}
\mu_j-\lambda_*=C\delta_{\mathrm F}^{-j}\bigl[1+O(\theta^j)\bigr],
\qquad
\frac{\mu_{j-1}-\mu_{j-2}}{\mu_j-\mu_{j-1}}\longrightarrow\delta_{\mathrm F}.
\label{hmtfg:escalado:eq:4}
\end{equation}
Ici \(\delta_{\mathrm F}>1\) est la valeur propre instable de l’opérateur de doublement. La sélection de \(\mu_j\) s’effectue par les retours locaux de la section~\ref{hmtfg:escalado:8}. L’identification avec le sélecteur global est démontrée dans les sections~\ref{hmtfg:escalado:15}--\ref{hmtfg:escalado:17} : tous les termes globaux conservent la combinatoire et les centres coïncident à tout niveau suffisamment élevé lorsque leurs périodes sont égalisées.


\subsection{Carte analytique et retour effectif}
\label{hmtfg:escalado:2}

On utilise la carte de fonctions paires
\[
s=x^2,\qquad z=\frac{s-1}{5/2},\qquad
f(x)=1-x^2V(z),\qquad V(z)=\sum_{j\ge0}v_jz^j.
\]
Les coordonnées de l’espace de Banach sont
\begin{equation}
u=10v_0,\qquad \nu=(v_1,v_2,\ldots),\qquad
\|(\Delta u,\Delta\nu)\|_{\mathcal A}
=|\Delta u|+\sum_{j\ge1}|\Delta v_j|.
\label{hmtfg:escalado:eq:5}
\end{equation}
En particulier, le coefficient constant de \(V\) a un poids de dix. Cette normalisation coïncide avec la carte \(u/10\) de Hertling–Spandl, section 2 de la source.

Soit \(a=-f(1)=v_0-1\). Pour les retours réels admissibles,
\begin{equation}
(Rf)(x)=-\frac{f(f(-ax))}{a}.
\label{hmtfg:escalado:eq:6}
\end{equation}
Le programme évalue une factorisation exacte de \eqref{hmtfg:escalado:eq:6}. En définissant
\[
S=\frac{a^2s-1}{5/2},\qquad b=V(S),\qquad
h=1-a^2sb,\qquad t=\frac{h^2-1}{5/2},
\]
\[
\operatorname{shift}V(z)=\sum_{j\ge0}v_{j+1}z^j,
\]
on obtient
\begin{equation}
\boxed{V_R=a\,b\,(h+1)\left[V(t)+\frac25\operatorname{shift}V(t)\right].}
\label{hmtfg:escalado:eq:7}
\end{equation}
Pour la vérifier, on utilise \(v_0=1+a\),
\(h^2-1=-a^2sb(h+1)\) et
\(V(t)-v_0=t\operatorname{shift}V(t)\). Ainsi,
\(a+1-h^2V(t)=a^2sb(h+1)[V(t)+(2/5)\operatorname{shift}V(t)]\),
ce qui est précisément \(asV_R\).

Les moments initiaux sont rationnels :
\begin{equation}
\frac1{12}\sum_jk_j^{2m}
=\frac1{12}\left(\frac4{899}\right)^m
\sum_{j=1}^{12}(3j-1)^{2m}.
\label{hmtfg:escalado:eq:8}
\end{equation}
Ainsi, tant \eqref{hmtfg:escalado:eq:7} que la construction initiale utilisent des intervalles rationnels ; les fonctions trigonométriques initiales sont représentées par leurs séries avec une borne extérieure du reste.


\subsection{Bornes externes utilisées comme reconnaissance ultérieure}
\label{hmtfg:escalado:3}

Une fois \eqref{hmtfg:escalado:eq:1} construite, le théorème d’hyperbolicité de Lanford est utilisé comme certificat analytique de l’opérateur \eqref{hmtfg:escalado:eq:6}. Soient \(g_0\) le polynôme central publié et \(g\) le point fixe. Les estimations utilisées sont
\[
\|g-g_0\|_{\mathcal A}<\varepsilon_0=4\,10^{-5},
\]
et, dans la boule \(\|f-g_0\|_{\mathcal A}<0.01\), les blocs de \(DR\), notés \(R=(U,V)\), satisfont
\begin{equation}
U_u\ge a_0=4.521,\quad
\|U_\nu\|\le b_0=0.560,\quad
\|V_u\|\le c_0=0.756,\quad
\|V_\nu\|\le d_0=0.719.
\label{hmtfg:escalado:eq:9}
\end{equation}
Le point fixe possède une unique direction instable, de valeur propre réelle positive \(\delta_{\mathrm F}>1\) ; le reste du spectre se situe à l’intérieur du disque unité. Les chiffres de \eqref{hmtfg:escalado:eq:9} proviennent des estimations publiées, et non d’une nouvelle reproduction de la preuve de Lanford dans ce dossier. Les coefficients de \(g_0\) sont explicitement incorporés au polynôme central de référence du programme, conformément au tableau 2 de Hertling–Spandl.

Sources : \href{https://repo-archives.ihes.fr/A_GARDER/20180409/Test/I_Prepublications/LANFORD/1968-2013/P_81_17/P_81_17_web.pdf}{Lanford, estimations d’hyperbolicité}; \href{https://arxiv.org/pdf/1410.3277}{Hertling–Spandl, carte analytique et tableau 2}.


\subsection{Certificat rationnel d’entrée et de direction}
\label{hmtfg:escalado:4}

Le certificat d’entrée dans la renormalisation divise \eqref{hmtfg:escalado:eq:2} en seize intervalles rationnels égaux. Il propage quatre retours de \eqref{hmtfg:escalado:eq:7} et la dérivée exacte par rapport à \(\lambda\), en utilisant la différentiation de la composition et l’arrondi extérieur à cent décimales.

Chaque série se compose d’un polynôme de degré 32 à coefficients intervalles et d’une erreur analytique arbitraire \(E\), \(\|E\|_1\le\epsilon\). Cette erreur peut également modifier les coefficients de bas degré. Pour \(\|q\|_1<1\),
\begin{equation}
\|E\circ q\|_1\le\epsilon,\qquad
\|E'\circ q\|_1\le\frac{\epsilon}{(1-\|q\|_1)^2},\qquad
\|\operatorname{shift}E\|_1\le\epsilon.
\label{hmtfg:escalado:eq:10}
\end{equation}
Les produits comprennent les termes polynôme–erreur et erreur–erreur. La queue initiale est majorée par son premier terme en valeur absolue et une raison géométrique, en utilisant \(\|1+(5/2)z\|_1=7/2\). Chaque argument composé reste strictement à l’intérieur de la boule unité de cette algèbre.

Pour \(\Gamma(\lambda)=R^4f_\lambda=(u_\lambda,\nu_\lambda)\), le certificat donne les bornes conservatrices suivantes, uniformes en \eqref{hmtfg:escalado:eq:2} :

\begingroup\small
\begin{center}\begin{tabular}{p{0.465\linewidth}p{0.465\linewidth}}
\hline
Quantité & Borne extérieure certifiée \\
\hline
\(\|\Gamma(\lambda)-g_0\|_{\mathcal A}\) & \(<0.004993128\) \\
\(\|\nu_\lambda-\nu_{g_0}\|_1\) & \(<0.001396077\) \\
\(u'_\lambda\) & \(>3821.289115\) \\
\(\|\nu'_\lambda\|_1\) & \(<551.757000\) \\
\(\|\nu'_\lambda\|_1/u'_\lambda\) & \(<0.144386<1/4\) \\
Erreur analytique de la fonction & \(<8.482\,10^{-19}\) \\
Erreur analytique de la tangente & \(<4.569\,10^{-15}\) \\
\hline\end{tabular}\end{center}
\endgroup

Les décimales complètes, les seize sous-intervalles et chaque retour intermédiaire sont conservés dans le certificat rationnel d’entrée.

Le cône \(\|\dot\nu\|_1\le|\dot u|/4\) est invariant par \eqref{hmtfg:escalado:eq:9} :
\[
|\dot u_{\rm nuevo}|\ge(4.521-0.560/4)|\dot u|
=4.381|\dot u|,
\]
\begin{equation}
\|\dot\nu_{\rm nueva}\|_1\le(0.756+0.719/4)|\dot u|
=0.93575|\dot u|<\tfrac14(4.381)|\dot u|.
\label{hmtfg:escalado:eq:11}
\end{equation}
La direction paramétrique certifiée entre donc dans un cône uniformément expansif de l’opérateur.


\subsection{Construction quantitative de la feuille stable}
\label{hmtfg:escalado:5}

Translatons \(g\) à l’origine et écrivons les coordonnées sous la forme \((x,y)=(u-u_g,\nu-\nu_g)\). Fixons
\begin{equation}
L=0.16,\qquad r=0.004,\qquad
A=a_0-Lc_0=4.40004,\qquad q=d_0+c_0L=0.83996.
\label{hmtfg:escalado:eq:12}
\end{equation}
Le cylindre \(|x|\le Lr\), \(\|y\|_1\le r\) se trouve à l’intérieur de la boule de \eqref{hmtfg:escalado:eq:9}, puisque
\((1+L)r+\varepsilon_0=0.00468<0.01\).

Considérons l’espace complet des fonctions \(\sigma:B_r\to\mathbb R\) avec \(\sigma(0)=0\) et de constante de Lipschitz au plus \(L\), muni de la distance uniforme. Pour chaque \(y\), on définit \((\mathcal G\sigma)(y)\) comme la racine dans \([-L\|y\|_1,L\|y\|_1]\) de
\begin{equation}
F_\sigma(x,y)=U(x,y)-\sigma(V(x,y)).
\label{hmtfg:escalado:eq:13}
\end{equation}
Sur cet intervalle, \(\|V(x,y)\|_1\le q\|y\|_1<r\). Lorsque \(x\) augmente, \eqref{hmtfg:escalado:eq:9} montre que \(F_\sigma\) croît avec une pente au moins égale à \(A\). À ses extrémités,
\[
F_\sigma(L\|y\|_1,y)
\ge[AL-b_0-Ld_0]\|y\|_1,
\]
et l’inégalité opposée vaut à l’extrémité négative. La marge est
\begin{equation}
AL-b_0-Ld_0=0.0289664>0.
\label{hmtfg:escalado:eq:14}
\end{equation}
On obtient une racine unique. En comparant \eqref{hmtfg:escalado:eq:13} pour deux valeurs de \(y\),
\[
\operatorname{Lip}(\mathcal G\sigma)
\le\frac{b_0+Ld_0}{A}
=\frac{0.67504}{4.40004}<0.16.
\]
En comparant deux graphes,
\begin{equation}
\|\mathcal G\sigma-\mathcal G\widetilde\sigma\|_\infty
\le A^{-1}\|\sigma-\widetilde\sigma\|_\infty,
\qquad A^{-1}<0.228.
\label{hmtfg:escalado:eq:15}
\end{equation}
Le théorème de contraction produit un unique graphe fixe \(x=\sigma_*(y)\). Sur celui-ci,
\begin{equation}
\|y_n\|_1\le q^n\|y_0\|_1,
\qquad
\|(x_n,y_n)\|_{\mathcal A}\le(1+L)q^n\|y_0\|_1.
\label{hmtfg:escalado:eq:16}
\end{equation}
De plus, la distance verticale \(d(x,y)=x-\sigma_*(y)\) satisfait
\[
|d(R(x,y))|\ge A|d(x,y)|
\]
tant que l’orbite reste dans le cylindre. Toute orbite qui y demeure indéfiniment appartient au graphe. Celui-ci identifie exactement la feuille stable locale et fournit l’estimation \eqref{hmtfg:escalado:eq:16}. Sa régularité locale coïncide avec celle de la variété stable analytique du point fixe hyperbolique.


\subsection{Existence, unicité et transversalité du paramètre}
\label{hmtfg:escalado:6}

Pour \(\lambda\) dans \eqref{hmtfg:escalado:eq:2}, définissons
\[
H(\lambda)=u_\lambda-u_g-
\sigma_*(\nu_\lambda-\nu_g).
\]
Le certificat assure \(\|\nu_\lambda-\nu_g\|_1<0.001436077<r\). Pour \(\lambda_2>\lambda_1\),
\begin{equation}
H(\lambda_2)-H(\lambda_1)
\ge\int_{\lambda_1}^{\lambda_2}
\bigl[u'_t-L\|\nu'_t\|_1\bigr]dt
>3733.007995\,(\lambda_2-\lambda_1).
\label{hmtfg:escalado:eq:17}
\end{equation}
Pour déterminer les signes aux extrémités, on utilise \(\|g-g_0\|_{\mathcal A}<\varepsilon_0\) et \(|\sigma_*(y)|\le L\|y\|_1\). L’évaluation ponctuelle rationnelle produit
\begin{equation}
H(1.874038)<-0.0011856110945,
\qquad H(1.874039)>0.0021866053399.
\label{hmtfg:escalado:eq:18}
\end{equation}
Par continuité, il existe une racine ; \eqref{hmtfg:escalado:eq:17} prouve son unicité et la transversalité du croisement. À cette racine, \eqref{hmtfg:escalado:eq:16} et
\((1+L)(0.001396077+0.00004)<0.001666\)
démontrent \eqref{hmtfg:escalado:eq:3}.

La borne est uniforme en le nombre de retours : elle convertit l’avancement en profondeur en une erreur operatorielle géométrique explicite.


\subsection{Admissibilité réelle de tous les retours}
\label{hmtfg:escalado:7}

Les quatre retours initiaux s’accompagnent, sur chaque sous-intervalle rationnel, des inégalités
\begin{equation}
0<a=-f(1)<1,\qquad b=f(a)>a,\qquad f(b)<a.
\label{hmtfg:escalado:eq:19}
\end{equation}
Le programme conserve \(a,b,f(b)\) avant chaque application de \eqref{hmtfg:escalado:eq:7}. L’unimodalité, la parité et le point critique quadratique sont conservés par ces retours normalisés.

Pour les itérations suivantes, une vérification uniforme dans la boule de \eqref{hmtfg:escalado:eq:9} suffit. Si \(E=V-V_{g_0}\), la norme \eqref{hmtfg:escalado:eq:5} donne, pour \(-0.4\le z\le0\),
\[
|E(z)|\le0.004,\qquad |E'(z)|\le0.01.
\]
La seconde inégalité utilise \(\sup_{j\ge1}j(0.4)^{j-1}=1\). L’évaluation rationnelle des coefficients centraux fournit
\begin{equation}
\frac65<V(z)<\frac85,\qquad |V'(z)|<\frac12,
\qquad 0.39<a<0.41.
\label{hmtfg:escalado:eq:20}
\end{equation}
En conséquence,
\[
f'(x)=-2x\left[V(z)+\frac25x^2V'(z)\right],
\]
et l’expression entre crochets dépasse un pour \(0\le x\le1\). L’application est unimodale et préserve \([-1,1]\). De plus,
\[
b=f(a)>1-(0.41)^2\frac85=\frac{4569}{6250}>a,
\]
\begin{equation}
f(b)<1-\left(\frac{4569}{6250}\right)^2\frac65
=\frac{35028967}{97656250}<0.39<a.
\label{hmtfg:escalado:eq:21}
\end{equation}
Les orbites de la feuille stable satisfont \eqref{hmtfg:escalado:eq:19} à toute profondeur. Ainsi, les itérations analytiques de l’opérateur correspondent à des retours réels de doublement, dont les domaines et les orientations sont conservés.


\subsection{Changement d’échelle de la cascade superstable locale}
\label{hmtfg:escalado:8}

La conclusion métrique exige deux croisements : celui de la famille avec la feuille stable, démontré dans la section~\ref{hmtfg:escalado:6}, et celui de la variété instable avec la cible superstable. Le second est démontré par Eckmann–Wittwer pour
\(\Sigma_2=\{\psi:\psi(1)=0\}\), correspondant au cycle superstable de période deux. La normalisation \eqref{hmtfg:escalado:eq:6} coïncide avec celle de l’opérateur réel pair utilisé dans ce travail.

Le théorème général de Collet–Eckmann–Lanford, section 6 de la source citée, proposition 6.1 et théorèmes 6.2–6.3, s’applique à une application d’un espace de Banach \(C^2\), un point fixe possédant une unique valeur propre instable \(\delta_{\mathrm F}>1\), une courbe transverse à sa feuille stable et une cible transverse à sa variété instable. Les préimages locales de la cible rencontrent la courbe en des points qui satisfont
\begin{equation}
\delta_{\mathrm F}^j(\mu_j-\lambda_*)\longrightarrow C\ne0.
\label{hmtfg:escalado:eq:22}
\end{equation}
L’indice \(j\) compte les retours locaux. Comme la courbe initiale est ici \(R^4f_\lambda\), la période physique est \(2^{j+5}\) lorsque la cible est prise dans \(\Sigma_2\) ; tout nombre fixe de pas utilisé pour ramener la cible dans le voisinage local est incorporé à un décalage fixe de l’indice et à \(C\).

Les conditions réelles de la section~\ref{hmtfg:escalado:7} et la branche réelle de doublement de la cible fixent la période exacte et l’itinéraire. La suite est sélectionnée par ce prolongement local orienté. En prenant les différences dans \eqref{hmtfg:escalado:eq:22}, on obtient la limite des rapports de \eqref{hmtfg:escalado:eq:4}.

Sources des deux résultats utilisés : \href{https://link.springer.com/article/10.1007/BF01013368}{Eckmann–Wittwer, \emph{A complete proof of the Feigenbaum conjectures}}; \href{https://people.math.harvard.edu/~knill/history/lanford/papers/ColletEckmannLanford.pdf}{Collet–Eckmann–Lanford, section 6}.

\subsubsection{Reste asymptotique en puissance}

L’analyticité de notre famille permet de préciser \eqref{hmtfg:escalado:eq:22}. Dans la construction de coordonnées normales de Collet–Eckmann–Lanford, seules la variété stable, la variété instable et l’hypersurface cible sont aplaties. Ces trois sous-variétés admettent des changements locaux \(C^2\) à dérivées bornées. La coupure lisse employée s’effectue dans la coordonnée scalaire instable ; le feuilletage complet s’obtient ensuite par le changement construit.

Dans ces coordonnées, la correction s’écrit comme une série d’incréments \(w_j\) dont la norme \(C^1\) satisfait
\[
\|w_j\|_{C^1}\le C_1\rho^j,\qquad 0<\rho<1.
\]
La règle de dérivation en chaîne pour l’application munie d’une coupure, avec des dérivées d’ordres un et deux uniformément bornées, fournit
\[
\|w_j\|_{C^2}\le C_2B^j
\]
pour un certain \(B\ge1\). La norme pondérée de la construction contrôle la norme \(C^1\) ordinaire sur le domaine considéré. Par interpolation,
\[
[Dw_j]_\beta\le C_3
\left(\rho^{1-\beta}B^\beta\right)^j.
\]
On choisit
\begin{equation}
0<\beta<\min\left\{1,
\frac{-\log\rho}{\log B-\log\rho}\right\}.
\label{hmtfg:escalado:eq:23}
\end{equation}
La série converge dans \(C^{1,\beta}\). Soit \(z\) la coordonnée normale obtenue, normalisée de sorte que les préimages de la cible satisfassent \(z=\delta_{\mathrm F}^{-j}\), en absorbant le décalage fixe de l’indice. Comme notre courbe est analytique et transverse,
\[
z(\Gamma(\lambda))=c(\lambda-\lambda_*)
+O(|\lambda-\lambda_*|^{1+\beta}),\qquad c\ne0.
\]
L’inversion locale donne
\begin{equation}
\mu_j-\lambda_*=C\delta_{\mathrm F}^{-j}
+O\left(\delta_{\mathrm F}^{-(1+\beta)j}\right).
\label{hmtfg:escalado:eq:24}
\end{equation}
Cela démontre le reste de \eqref{hmtfg:escalado:eq:4}, avec \(\theta=\delta_{\mathrm F}^{-\beta}\). L’énoncé affirme l’existence des constantes du reste paramétrique. Le taux operatoriel explicite \(0.83996\) de \eqref{hmtfg:escalado:eq:3} contrôle, quant à lui, la convergence des applications renormalisées.

\subsubsection{Conversion du reste en erreur sur les rapports}

Si \(\mu_j-\lambda_*=C\delta_{\mathrm F}^{-j}(1+e_j)\), avec \(|e_j|\le M\theta^j\), on écrit
\[
\mu_{j-1}-\mu_j=C(\delta_{\mathrm F}-1)\delta_{\mathrm F}^{-j}(1+r_j),
\quad
r_j=\frac{\delta_{\mathrm F} e_{j-1}-e_j}{\delta_{\mathrm F}-1}.
\]
Avec \(K=M(\delta_{\mathrm F}+\theta)/(\delta_{\mathrm F}-1)\), on a \(|r_j|\le K\theta^{j-1}\). Par conséquent,
\begin{equation}
\boxed{
\left|\frac{\mu_{j-2}-\mu_{j-1}}{\mu_{j-1}-\mu_j}-\delta_{\mathrm F}\right|
\le\frac{\delta_{\mathrm F} K(1+\theta)\theta^{j-2}}
{1-K\theta^{j-1}}}
\label{hmtfg:escalado:eq:25}
\end{equation}
lorsque \(K\theta^{j-1}<1\). Cette formule sépare la précision d’évaluation de chaque racine de l’erreur de troncature de la cascade.


\subsection{Prolongement certifié des huit racines initiales}
\label{hmtfg:escalado:9}

Le certificat de prolongement fini conserve la famille \eqref{hmtfg:escalado:eq:1} et le certificat antérieur de huit racines. Pour \(2\le n\le8\), il étudie
\[
S_n(\lambda)=f_\lambda^{2^n}(0)
\]
depuis la boîte certifiée de \(\lambda_{n-1}\) jusqu’à \(1.874039\). Le recouvrement rationnel complet contient exactement la racine héritée \(\lambda_{n-1}\) et la nouvelle racine \(\lambda_n\). En particulier, \(\lambda_n\) est l’unique nouvelle racine dans \((\lambda_{n-1},1.874039]\).

Les boîtes sont exclues par le signe du retour ou par une dérivée locale de signe certifié. Les boîtes d’ancrage conservent l’existence et l’unicité par les signes aux extrémités et la monotonie locale. Leurs extrémités couvrent exactement chaque intervalle ; les adjacences sont vérifiées. L’évaluation utilise les moments \eqref{hmtfg:escalado:eq:8}, l’évaluation de Horner d’ordre 32, des queues géométriques extérieures et \(|u_0''|\le2\). La queue uniforme du potentiel est inférieure à \(1.294\,10^{-58}\) sur \(|x|\le3/2\).

Le certificat final contient 650 boîtes traitées et 332 feuilles de recouvrement. Le premier niveau possède l’expression exacte \(\lambda_1=1/u_0(1)\). Les approximations suivantes sont présentées uniquement à titre indicatif ; leurs boîtes complètes sont conservées dans les certificats :

\begingroup\small
\begin{center}\begin{tabular}{p{0.465\linewidth}p{0.465\linewidth}}
\hline
Niveau & Paramètre superstable \\
\hline
1 & 1.345913990471832792210949095… \\
2 & 1.763379787846910127290794030… \\
3 & 1.850413796950136464530566778… \\
4 & 1.868979756606798125150574309… \\
5 & 1.872955034435659740004205128… \\
6 & 1.873806367467030119412262311… \\
7 & 1.873988695221365362307168780… \\
8 & 1.874027744154450672897007573… \\
\hline\end{tabular}\end{center}
\endgroup

% Fragmento integrable. Preambulo anfitrion: amsmath, amssymb, longtable, hyperref.
\section{Isolement de la limite et changement d’échelle du sélecteur global}
\label{hmtfg:fragmento:aislamiento-del-limite-y-escalado-del-selector-global}

\subsection{Continuation indéfinie de la première racine et exclusion de l’intervalle intermédiaire}
\label{hmtfg:escalado:15}

La classification complète démontrée précédemment conserve la famille et le sélecteur de première racine nouvelle. Le transport ordonné conserve les zéros et leur minimum ; le mot de retours conserve la lecture chronologique nonadique.

\subsubsection{Classification et existence à toute profondeur}
\label{hmtfg:escalado:15.1}

Écrivons \(W_n=\tau_1\cdots\tau_{2^n-1}\), avec \(\tau_j=(-1)^{v_2(j)}\), et \(a_*=\min\Lambda_\tau\). L’existence et la compacité de \(\Lambda_\tau\) sont démontrées dans le développement de prolongation, \ref{hmtfg:prolongacion:6.1}–\ref{hmtfg:prolongacion:6.2} ; son minimum existe par compacité et ordre.

\textbf{Classification.} Toute application strictement unimodale à maximum, sur un intervalle invariant, dont le point critique a pour période exacte \(2^n\) et dont l’itinéraire \(K\) satisfait \(K<\tau\), a pour itinéraire \(W_n0\).

L’induction s’effectue au moyen du retour effectif. Pour une période d’au moins huit, la comparaison avec \(\tau\), jointe aux intervalles qui piégeraient des orbites de signes incompatibles, impose le préfixe \((+,-,+,+,+,-)\). On en déduit
\[
J_1=[d_2,d_4],\quad F(J_1)=[d_3,1],\quad
d_4<d_3,\quad F^2(J_1)=J_1.
\]
Le retour normalisé par \(h(x)=d_2x\) a pour maximum un, pour période la moitié de la période initiale et pour itinéraire \(K'_j=-K_{2j}<\tau\). L’ordre est conservé parce que
\[
O_{2q-1}(K)=-O_{q-1}(K'),\qquad
K_{2q}-\tau_{2q}=-(K'_q-\tau_q).
\]
Les cas initiaux sont \(+0\) et \(+-+0\) ; l’induction reconstruit \(W_n0\). Pour les applications paires, on a aussi \(0<d_4<-d_2<1\), de sorte que le retour s’étend à \([-1,1]\) et coïncide avec \eqref{hmtfg:escalado:eq:6}.

Avant \(a_*\), tous les itinéraires sont inférieurs à \(\tau\) : les ensembles de comparaison stricte sont ouverts par la séparation démontrée dans \ref{hmtfg:prolongacion:6.1} ; l’intervalle précédant \(a_*\) est connexe et commence par \(+^\infty<\tau\).

\textbf{Existence de chaque première racine.} Supposons \(\lambda_{n-1}<a_*\) construite, et soit \(p=2^{n-1}\). Les zéros de \(c_p\) dans \([\lambda_{n-1},a_*]\) forment un ensemble fini non vide : la fonction est analytique et \(c_p(a_*)\ne0\). Soit \(z<a_*\) son dernier zéro. Dans \((z,a_*]\), \(c_p\) est de signe \(\tau_p\). Pour \(g_\lambda=f_\lambda^p\),
\[
g_\lambda'(0)=0,\qquad
c_{2p}(\lambda)=c_p(\lambda)+O(c_p(\lambda)^2).
\]
Au voisinage de \(z\), à droite, \(c_{2p}\) est de signe \(\tau_p\) ; en \(a_*\), son signe est opposé. Il existe un zéro de \(c_{2p}\) dans \((z,a_*)\), et sa période est exactement \(2p\), car on y a \(c_p\ne0\). La finitude des zéros fournit la première racine nouvelle dans \((\lambda_{n-1},a_*)\). Le dernier zéro auxiliaire prouve l’existence ; le choix effectif reste la première racine. Le cas de base est \(\lambda_1=1/u_0(1)<a_*\).

\subsubsection{Limite de la suite initiale}
\label{hmtfg:escalado:15.2}

\textbf{Théorème.} Le sélecteur initial existe à tous les niveaux et satisfait
\begin{equation}
\boxed{K_{\lambda_n}=W_n0,\qquad
\lambda_n\nearrow a_*=\min\Lambda_\tau.}
\label{hmtfg:escalado:eq:50}
\end{equation}
\textbf{Démonstration.} Le résultat d’existence précédent donne une suite croissante majorée par \(a_*\). La classification fixe tous ses itinéraires. Soit \(b\le a_*\) sa limite. Pour chaque \(p\), les termes suffisamment avancés conservent les signes opposés \(\tau_p,\tau_{2p}\). La formule de Taylor et la borne uniforme \(B_p=16^p(16^p-1)/15\) donnent
\[
|c_{2p}(\lambda_n)-c_p(\lambda_n)|
\le\tfrac12 B_p|c_p(\lambda_n)|^2,\qquad
|c_p(\lambda_n)|>2/B_p.
\]
Le passage à la limite conserve \(\operatorname{sign}c_p(b)=\tau_p\), avec une valeur absolue d’au moins \(2/B_p>0\). Pour tout \(p\), cela implique \(b\in\Lambda_\tau\) ; la minimalité donne \(b=a_*\).

Le résultat établit l’existence indéfinie, la conservation de la cascade symbolique et la sélection de son premier paramètre limite pour la règle globale initiale. La complétion HMT transporte cette suite et sa limite en conservant l’ordre et la lecture.

\subsubsection{Certificat d’exclusion de l’intervalle intermédiaire}
\label{hmtfg:escalado:15.3}

Le programme de pont certifie
\begin{equation}
\Lambda_\tau\cap[\lambda_{8,\rm inf},1.874038]=\varnothing,
\label{hmtfg:escalado:eq:51}
\end{equation}
où
\[
\lambda_{8,\rm inf}
=1.874027744154450672897007573595092408435133414523352839687903472138975
\]
est l’extrémité inférieure rationnelle de l’intervalle certifié de \(\lambda_8\).

Le recouvrement contient 374 intervalles exactement adjacents et aucune région en suspens : 319 avec \(c_{512}>0\), 10 avec \(c_{1024}<0\), deux avec \(c_{2048}>0\), tous opposés au signe de \(\tau\), et 43 proches de centres de périodes 256, 512 ou 1024, exclus par leur deuxième retour.

Dans ce dernier cas, \(g=f_\lambda^p\), \(d=g(0)\) et une borne uniforme \(|g''|\le M\) entre zéro et \(d\) satisfont \(M|d|<2\). Alors
\[
|g(d)-d|\le\tfrac12M|d|^2<|d|\quad(d\ne0).
\]
Les deux retours ont le même signe ; pour \(d=0\), ils sont tous deux nuls. Les deux situations excluent \(\tau_{2p}=-\tau_p\) avec des signes stricts.

L’évaluation utilise des moments rationnels, des queues géométriques du potentiel et de sa dérivée seconde, et \(|u_0'''|<6\). Le centrage paramétrique intersecte l’enveloppe naturelle avec
\[
c_j(\lambda_{\rm medio})+
\partial_\lambda c_j(I)\,(I-\lambda_{\rm medio});
\]
la dérivée est propagée sur la boîte entière. La propriété d’autoapplication de \([-1,1]\) est démontrée avant d’y restreindre les enveloppes.

Le certificat complet conserve le recouvrement et vérifie son pavage exact, ainsi que chaque inégalité de signe ou de Taylor.

\subsubsection{Localisation conjointe}
\label{hmtfg:escalado:15.4}

Le croisement local \(\lambda_*\) réalise \(\tau\) par ses retours réels à toute profondeur, de sorte que \(a_*\le\lambda_*\). Par \eqref{hmtfg:escalado:eq:50}, \(\lambda_8<a_*\) ; avec \eqref{hmtfg:escalado:eq:51},
\begin{equation}
\boxed{1.874038<a_*\le\lambda_*<1.874039.}
\label{hmtfg:escalado:eq:52}
\end{equation}
La localisation conserve l’itinéraire canonique et situe son premier paramètre limite dans le même intervalle que le croisement stable.

\subsection{Isolement par première sortie et coïncidence des limites}
\label{hmtfg:escalado:16}

Cette section complète l’isolement identifié dans la section~\ref{hmtfg:escalado:15.4}. Elle conserve la famille HMT \eqref{hmtfg:escalado:eq:1}, la sélection minimale \eqref{hmtfg:escalado:eq:50}, la carte \eqref{hmtfg:escalado:eq:5} et la feuille stable des sections~\ref{hmtfg:escalado:5}--\ref{hmtfg:escalado:7}. La preuve utilise l’ordre \(a_*\le\lambda_*\) : il suffit d’exclure une séparation vers le côté négatif du cône expansif.

\subsubsection{Bornes supérieure et inférieure du transport transversal}
\label{hmtfg:escalado:16.1}

Outre \eqref{hmtfg:escalado:eq:9}, le lemme 2.1 de \href{https://arxiv.org/pdf/1410.3277}{Hertling–Spandl, section 2} fournit \(\|D\Phi\|<0.9\) dans la même boule de rayon \(0.01\), où
\[
\Phi_u(u,\nu)=u-\frac{U(u,\nu)-u}{3.669}.
\]
Par conséquent,
\begin{equation}
\left|1-\frac{U_u-1}{3.669}\right|<0.9,\qquad
U_u<1+1.9(3.669)=7.9711.
\label{hmtfg:escalado:eq:53}
\end{equation}
Le nombre \(3.669\) appartient au préconditionneur publié de ce certificat analytique ; la famille générée \eqref{hmtfg:escalado:eq:1} conserve ses coefficients antérieurs.

Prenons \(\kappa=0.21\). Pour deux points de la boule joints par une sécante avec
\(\Delta u<0\), \(\|\Delta\nu\|_1\le\kappa|\Delta u|\), l’intégration des blocs de \(DR\) sur le segment donne
\begin{equation}
4.4034|\Delta u|
\le|\Delta u'|\le8.0887|\Delta u|,\qquad \Delta u'<0,
\label{hmtfg:escalado:eq:54}
\end{equation}
\begin{equation}
\|\Delta\nu'\|_1
\le(0.756+0.719\kappa)|\Delta u|
=0.90699|\Delta u|
<\kappa(4.4034)|\Delta u|.
\label{hmtfg:escalado:eq:55}
\end{equation}
Ici \(4.4034=4.521-0.560\kappa\) et \(8.0887=7.9711+0.560\kappa\). Ainsi, le cône des sécantes conserve son orientation et se dilate tant que le segment reste dans la boule.

\subsubsection{Réduction d’une séparation infinie à une région finie}
\label{hmtfg:escalado:16.2}

Supposons \(a_*<\lambda_*\), et posons
\[
f_n=R^{4+n}f_{a_*},\qquad
s_n=R^{4+n}f_{\lambda_*},\qquad
(\Delta u_n,\Delta\nu_n)=f_n-s_n.
\]
La borne positive de \(u'\) et \(\|\nu'\|_1/u'<0.144386<\kappa\) sur \(\Gamma(J)\) placent la sécante initiale dans le cône négatif. La feuille stable est écrite par rapport à \(g\), avec
\begin{equation}
s_n=g+e_n,\quad
\|e_n\|_{\mathcal A}<0.001666(0.83996)^n,\quad
|(e_n)_u|\le0.16\|(e_n)_\nu\|_1.
\label{hmtfg:escalado:eq:56}
\end{equation}
Les contrôles d’entrée impliquent
\begin{equation}
|\Delta u_0|
<0.004993128+0.00004
 +0.16(0.001396077+0.00004)
=0.00526290032< h,\qquad h=0.006.
\label{hmtfg:escalado:eq:57}
\end{equation}
Tant que \(|\Delta u_n|<h\), \eqref{hmtfg:escalado:eq:55}–\eqref{hmtfg:escalado:eq:56} donnent
\begin{equation}
\|f_n-g_0\|_{\mathcal A}
<(1+\kappa)h+0.001666+0.00004
=0.008966<0.01.
\label{hmtfg:escalado:eq:58}
\end{equation}
Le point \(s_n\) et le segment entre les deux restent également dans cette boule convexe. L’expansion inférieure \eqref{hmtfg:escalado:eq:54} impose une première sortie finie \(N\), et la borne supérieure appliquée à l’étape précédente fournit
\begin{equation}
0.006\le-\Delta u_N<0.0485322,\qquad
\|\Delta\nu_N\|_1\le0.21|\Delta u_N|.
\label{hmtfg:escalado:eq:59}
\end{equation}
On utilise l’inégalité \(8.0887h=0.0485322\) ; le saut complet, et non seulement la surface de première sortie, est inclus.

Par conséquent, \(f_N\) appartient à la région suivante :
\begin{equation}
\begin{split}
\mathcal C_-=\{\,g+(d_u,d_\nu)+e:\;&
-0.0485322\le d_u\le-0.006,\\
&\|d_\nu\|_1\le0.21|d_u|,\quad
|e_u|+\|e_\nu\|_1\le0.001666,\\
&|e_u|\le0.16\|e_\nu\|_1\,\}.
\end{split}
\label{hmtfg:escalado:eq:60}
\end{equation}
Puisque \(a_*\) réalise \(\tau\), tous ses retours normalisés sont des autoapplications unimodales de \([-1,1]\) ayant le même itinéraire. Cette propriété découle de l’induction des retours, indépendamment de la distance de \(f_N\) à \(g_0\). C’est pourquoi le certificat suivant traite \(\mathcal C_-\) intersectée avec cette classe d’autoapplications.

\subsubsection{Exclusion certifiée de la région de première sortie}
\label{hmtfg:escalado:16.3}

Le certificat rationnel d’exclusion par première sortie, reproduit dans le supplément computationnel, couvre exactement l’intervalle de \eqref{hmtfg:escalado:eq:59} par 64 bandes rationnelles adjacentes. Toutes sont exclues ; aucune subdivision supplémentaire ni région en suspens n’intervient dans le résultat.

Parmi les 64 bandes, 54 présentent un signe contraire et dix un retour contractant de période 16 ou 32. La plus grande borne de Lipschitz de ces retours est inférieure à \(0.475146491\) ; l’horizon maximal d’exclusion est l’itérée 64.

L’évaluation conserve comme variables partagées les quarante premiers coefficients de la perturbation, ainsi qu’une borne explicite de leur queue infinie. Pour chaque bande, le centre est \(g_0\) déplacé dans la coordonnée \(u\). Si \(a,b\ge0\) bornent les sensibilités duales dans \(u,\nu\), le support de l’incertitude du point fixe et du résidu stable est
\begin{equation}
0.00004\max(a,b)
+0.001666\max\left(b,\frac{0.16a+b}{1.16}\right).
\label{hmtfg:escalado:eq:61}
\end{equation}
La deuxième expression s’obtient en maximisant \(a|e_u|+b\|e_\nu\|_1\) sous les deux contraintes de \eqref{hmtfg:escalado:eq:56}. On utilise ainsi l’orientation de la feuille stable démontrée précédemment, en plus de sa norme.

Soit \(x_j\) l’orbite du centre et \(L_j\) sa sensibilité linéaire à tous les coefficients partagés. Le modèle de Taylor prend la forme
\[
c_j=x_j+L_j(\Delta v)+r_j,\qquad |r_j|\le E_j.
\]
Si \(\rho_j\) borne \(|L_j(\Delta v)|+E_j\), la récurrence extérieure employée est
\begin{equation}
E_{j+1}\le |f_0'(x_j)|E_j
+\tfrac12\sup|f_0''|\rho_j^2
+\sup|\Delta f'|\rho_j.
\label{hmtfg:escalado:eq:62}
\end{equation}
Les dérivées sont bornées sur le segment des deux états. L’orbite centrale est vérifiée à l’intérieur de \([-1,1]\) ; les autres orbites appartiennent à cet intervalle par la propriété d’autoapplication incluse dans le domaine certifié. Les calculs utilisent un arrondi rationnel vers l’extérieur, y compris pour les queues géométriques.

Chaque bande est exclue par l’une de deux inégalités :

\par\noindent\textbf{1.}\enspace Une valeur critique a un signe strictement opposé à \(\tau_j\).
\par\noindent\textbf{2.}\enspace Pour \(p=2^k\), le retour \(H=f^p\) a pour constante de Lipschitz \(L_p<1\) sur le segment entre \(0\) et \(H(0)\). Alors
\[
|H(H(0))-H(0)|\le L_p|H(0)|.
\]
Si \(H(0)\ne0\), les deux valeurs ont le même signe ; si \(H(0)=0\), elles s’annulent toutes deux. Dans les deux cas, l’itinéraire strict \(\tau_{2p}=-\tau_p\) est exclu.

La constante \(L_p\) s’obtient en propageant le segment initial \([-r,r]\), \(r\ge|H(0)|\), avec les boîtes de l’orbite critique et en multipliant les bornes extérieures des dérivées. On borne les applications de la bande entière, en maintenant \eqref{hmtfg:escalado:eq:61}.

\subsubsection{Théorème de coïncidence du paramètre d’accumulation}
\label{hmtfg:escalado:16.4}

\textbf{Théorème.} Pour la famille uniforme HMT \eqref{hmtfg:escalado:eq:1}, la limite de la sélection initiale des premières racines coïncide avec le croisement stable certifié :
\begin{equation}
\boxed{\lambda_n\nearrow a_*=\lambda_*,
\qquad 1.874038<\lambda_*<1.874039.}
\label{hmtfg:escalado:eq:63}
\end{equation}

\textbf{Démonstration.} \eqref{hmtfg:escalado:eq:50}–\eqref{hmtfg:escalado:eq:52} donnent \(\lambda_n\uparrow a_*\le\lambda_*\), les deux paramètres étant dans \(J\). Si l’inégalité était stricte, \eqref{hmtfg:escalado:eq:54}–\eqref{hmtfg:escalado:eq:59} produiraient un retour \(f_N\in\mathcal C_-\) réalisant \(\tau\). Le recouvrement complet de la section~\ref{hmtfg:escalado:16.3} exclut précisément cette possibilité. Donc \(a_*=\lambda_*\).

La preuve contrôle toutes les profondeurs au moyen d’une première sortie finie et d’un recouvrement uniforme de sa région. L’identité \eqref{hmtfg:escalado:eq:63} relie la limite du sélecteur global au croisement transversal.


\subsection{Identification des centres initiaux et changement d’échelle global}
\label{hmtfg:escalado:17}

La dernière étape conserve la sélection de \eqref{hmtfg:escalado:eq:50}. La continuation locale et les paramètres initiaux sont comparés par leur période exacte et leur classe hybride ; la coïncidence s’obtient par une unicité dans un voisinage paramétrique fixe.

\subsubsection{Réalisation analytique de degré deux et transversalité}
\label{hmtfg:escalado:17.1}

Le point fixe \(g\) de la section~\ref{hmtfg:escalado:3} est pair, analytique, à point critique quadratique et de combinatoire stationnaire de doublement. Le théorème 2.1 de \href{https://annals.math.princeton.edu/wp-content/uploads/annals-v164-n3-p01.pdf}{de Faria–de Melo–Pinto}, appliqué à cette unique combinatoire, donne un ensemble limite à un seul élément, avec une extension holomorphe propre de degré deux entre disques topologiques emboîtés. Son attraction sur l’intervalle identifie cet élément à \(g\), puisque \(Rg=g\).

Le lemme 3.2 de la même source étend un nombre fixe de retours à un voisinage analytique complet de \(g\), avec des images de degré deux et un module positif. La norme \eqref{hmtfg:escalado:eq:5} contrôle la norme uniforme sur un voisinage complexe suffisamment petit de \([-1,1]\), car on y a \(|(x^2-1)/2.5|<1\). La convergence \eqref{hmtfg:escalado:eq:3} entre dans ce voisinage ; la dépendance analytique des retours fournit, pour un certain entier fixe \(d\), une famille
\begin{equation}
\mathcal F_\lambda=R^d f_\lambda
\label{hmtfg:escalado:eq:64}
\end{equation}
d’applications de degré deux, analytique dans un voisinage complexe de \(\lambda_*\). Les disques et le nombre de retours sont fixés avant de faire varier \(\lambda\) dans ce voisinage.

La tangente paramétrique appartient au cône expansif, tandis que la feuille stable satisfait la borne de pente de la section~\ref{hmtfg:escalado:5}. Les retours conservent cette séparation. La réalisation de degré deux identifie localement la feuille stable à la classe hybride de l’application infiniment renormalisable ; \eqref{hmtfg:escalado:eq:64} est donc transversale à cette classe.

Le changement d’espace analytique est justifié en détail dans \ref{hmtfg:clasificacion:8.2}, consacrée au transport de la transversalité. Une direction tangente à la classe hybride produirait des tangentes renormalisées décroissantes, par la convergence uniforme sur un disque hybride et l’estimation de Cauchy. La direction certifiée satisfait, en revanche, \(|\dot f(1)|=|\dot u|/10\) et croît sous les retours grâce au cône expansif. Cette évaluation commune des deux représentants prouve la transversalité sans supposer l’équivalence de leurs normes.

\subsubsection{Unicité dans un voisinage fixe}
\label{hmtfg:escalado:17.2}

Pour \(n>d\), la classification de la section~\ref{hmtfg:escalado:15} donne
\begin{equation}
K(\mathcal F_{\lambda_n})=W_{n-d}0.
\label{hmtfg:escalado:eq:65}
\end{equation}
Le redressement d’une application de degré deux conserve l’orbite critique ; \eqref{hmtfg:escalado:eq:65} identifie sa classe hybride à celle du centre quadratique canonique de période \(2^{n-d}\). Les bornes a priori de la suite réelle de doublement sont celles employées dans le théorème d’universalité de Lyubich.

Le \href{https://www.maths.tcd.ie/EMIS/journals/Annals/149_2/lyubich.pdf\#page=69}{théorème 7.4 de Lyubich, p. 387–388} donne une seule intersection avec chacune de ces classes, aux niveaux suffisamment élevés, sur la transversale analytique et dans un voisinage fixe du paramètre d’accumulation.

Par \eqref{hmtfg:escalado:eq:63}, \(\lambda_n\) entre finalement dans ce voisinage. La continuation locale de la section~\ref{hmtfg:escalado:8}, indexée par la même période, appartient à la même classe hybride. L’unicité implique
\begin{equation}
\boxed{\lambda_n=\widehat\mu_n\quad(n\ge n_0),}
\label{hmtfg:escalado:eq:66}
\end{equation}
où \(\widehat\mu_n\) désigne le centre local de période physique \(2^n\). Il s’agit d’une réindexation de \(\mu_j\) par un décalage fixe déterminé par les retours et la cible de la section~\ref{hmtfg:escalado:8}. La sélection initiale de \(\lambda_n\) reste intacte.

L’entier \(n_0\) est existentiel dans cette preuve. La continuation à tous les niveaux et le mot \(W_n0\) sont déjà démontrés dans \eqref{hmtfg:escalado:eq:50} ; \eqref{hmtfg:escalado:eq:66} apporte l’identité métrique nécessaire à la limite.

\subsubsection{Théorème de changement d’échelle pour les premières racines HMT}
\label{hmtfg:escalado:17.3}

\textbf{Théorème.} Pour la famille uniforme \eqref{hmtfg:escalado:eq:1}, soient \(\lambda_n\) les premières racines nouvelles successives de période critique exacte \(2^n\). Il existe \(C>0\), \(0<\theta<1\) et \(M<\infty\) tels que, pour tout \(n\) suffisamment grand,
\begin{equation}
\lambda_*-\lambda_n=C\delta_{\mathrm F}^{-n}(1+e_n),
\qquad |e_n|\le M\theta^n,
\label{hmtfg:escalado:eq:67}
\end{equation}
et
\begin{equation}
\boxed{\lim_{n\to\infty}
\frac{\lambda_{n-1}-\lambda_{n-2}}
{\lambda_n-\lambda_{n-1}}=\delta_{\mathrm F}.}
\label{hmtfg:escalado:eq:68}
\end{equation}

\textbf{Démonstration.} \eqref{hmtfg:escalado:eq:66} transporte le changement d’échelle local \eqref{hmtfg:escalado:eq:24} aux racines initiales ; le décalage fini des indices est absorbé dans \(C,M\). Son signe est positif parce que \(\lambda_n<\lambda_*\).

Le reste en puissance admet aussi une justification directe par l’holonomie des classes hybrides. Le lemme 7.3 de Lyubich donne une régularité \(C^{1+\beta}\), avec une dérivée non nulle, sur la transversale au point d’accumulation. La dynamique instable unidimensionnelle est analytique et a pour multiplicateur \(\delta_{\mathrm F}>1\) ; sa coordonnée linéarisante place les centres successifs à une distance proportionnelle à \(\delta_{\mathrm F}^{-n}\). La composition avec l’inverse de l’holonomie donne un terme linéaire non nul et un reste \(O(\delta_{\mathrm F}^{-(1+\beta)n})\). On peut ainsi prendre \(\theta=\delta_{\mathrm F}^{-\beta}\).

En écrivant
\[
r_n=\frac{\delta_{\mathrm F} e_{n-1}-e_n}{\delta_{\mathrm F}-1},
\qquad K=\frac{M(\delta_{\mathrm F}+\theta)}{\delta_{\mathrm F}-1},
\]
on obtient
\[
\lambda_n-\lambda_{n-1}
=C(\delta_{\mathrm F}-1)\delta_{\mathrm F}^{-n}(1+r_n),\qquad
|r_n|\le K\theta^{n-1}.
\]
Par conséquent,
\begin{equation}
\left|
\frac{\lambda_{n-1}-\lambda_{n-2}}
{\lambda_n-\lambda_{n-1}}-\delta_{\mathrm F}
\right|
\le
\frac{\delta_{\mathrm F} K(1+\theta)\theta^{n-2}}
{1-K\theta^{n-1}},
\quad K\theta^{n-1}<1.
\label{hmtfg:escalado:eq:69}
\end{equation}
Le membre de droite tend vers zéro, ce qui démontre \eqref{hmtfg:escalado:eq:68}.

\subsubsection{Enchaînement et portée de la clôture}
\label{hmtfg:escalado:17.4}

La composition effective est : lecteur dodécaphasique APP–TRIT–TPK ; profil \eqref{hmtfg:escalado:eq:1} ; transport ordonné des retours \eqref{hmtfg:escalado:eq:38}–\eqref{hmtfg:escalado:eq:40} ; classification et existence des premières racines \eqref{hmtfg:escalado:eq:50} ; localisation \eqref{hmtfg:escalado:eq:52} ; isolement par première sortie \eqref{hmtfg:escalado:eq:63} ; identité des centres à partir d’un certain rang \eqref{hmtfg:escalado:eq:66} ; changement d’échelle \eqref{hmtfg:escalado:eq:67}–\eqref{hmtfg:escalado:eq:69}. La représentation de Jacobi \eqref{hmtfg:escalado:eq:43} conserve le profil complet de ce lecteur et ses marques restent dans l’antécédent.

L’expansion spatiale du lecteur HMT et le raffinement nonadique apportent la généalogie et l’évaluation à profondeur arbitraire. Le changement d’échelle entre paramètres se démontre par la dynamique du retour, la transversalité et l’isolement que nous venons de composer. Les preuves analytiques citées interviennent sur la famille déjà construite ; leurs valeurs de comparaison ne sélectionnent ni ses phases ni ses coefficients.

La clôture concerne le sélecteur global initial du secteur uniforme \(\eta=0\), avec les dépendances théorématiques énumérées. La loi asymptotique conserve ses quantificateurs existentiels. La largeur d’une boîte de racines finies mesure la précision d’évaluation ; l’estimation \eqref{hmtfg:escalado:eq:69} exprime la convergence vers la limite.


Les programmes et certificats du supplément de reproduction permettent de reproduire les parties finies de la preuve. La continuation à toute profondeur repose sur les inductions, l’argument uniforme de première sortie et les théorèmes d’identification, avec leurs domaines explicites.

% Fragmento integrable. Preambulo anfitrion: amsmath, amssymb, longtable, hyperref.
\section{Identification analytique des centres initiaux}
\label{hmtfg:fragmento:identificacion-analitica-de-los-centros-originales}

\subsection{Réalisation de degré deux et identification métrique des centres}
\label{hmtfg:clasificacion:8}

L’égalité entre le premier paramètre de l’itinéraire infini et le croisement stable, démontrée par première sortie, fournit
\begin{equation}
\boxed{a_*:=\min\Lambda_\tau=\lambda_*.}
\label{hmtfg:clasificacion:eq:H}
\end{equation}

La famille, le profil dodécaphasique, les retours et le sélecteur restent exactement tels qu’ils sont dans les sections~\ref{hmtfg:clasificacion:1}--\ref{hmtfg:clasificacion:5}. Les résultats externes employés ci-dessous reconnaissent la dynamique de cet objet déjà construit. Leurs valeurs universelles n’interviennent pas dans la sélection des coefficients ni des paramètres de la famille HMT.

\textbf{Théorème de transfert métrique.} L’identification \eqref{hmtfg:clasificacion:eq:H}, jointe aux bornes d’entrée, d’admissibilité et de transversalité démontrées précédemment, implique que la suite initiale des premières racines nouvelles satisfait, pour certaines constantes \(C>0\) et \(0<\theta<1\),
\begin{equation}
\boxed{\lambda_*-\lambda_n=C\delta_{\mathrm F}^{-n}
\bigl(1+O(\theta^n)\bigr),}
\label{hmtfg:clasificacion:eq:14}
\end{equation}
où \(\delta_{\mathrm F}>1\) est la valeur propre instable de l’opérateur de doublement au point fixe reconnu. En particulier,
\begin{equation}
\frac{\lambda_{n-1}-\lambda_{n-2}}
     {\lambda_n-\lambda_{n-1}}
=\delta_{\mathrm F}+O(\theta^n).
\label{hmtfg:clasificacion:eq:15}
\end{equation}
La suite coïncide à partir d’un certain rang avec les centres de la continuation locale, une fois que les deux indices désignent la même période physique. La preuve conserve le sélecteur de première racine de la section~\ref{hmtfg:clasificacion:4}.

\subsubsection{Obtention d’un domaine commun de type quadratique}
\label{hmtfg:clasificacion:8.1}

Notons \(g\) le point fixe de la renormalisation et
\(\Gamma(\lambda)=R^4f_\lambda\). La carte analytique utilisée dans le certificat est
\[
f(x)=1-x^2V\!\left(\frac{x^2-1}{2.5}\right),\qquad
V(z)=\frac{u}{10}+\sum_{k\ge1}\nu_kz^k,
\]
avec la norme des différences \(|\Delta u|+\|\Delta\nu\|_1\). Les inégalités d’admissibilité prouvent que \(g\) est paire, analytique, unimodale, quadratique et de type doublement ; de plus, \(g(x)=\varphi(x^2)\) avec \(\varphi'<0\) dans \([0,1]\).

Le théorème 2.1 de de Faria–de Melo–Pinto, appliqué à l’ensemble de combinatoires constitué uniquement du doublement, fournit un ensemble limite à un seul point avec une extension de type quadratique. Puisque \(Rg=g\), la convergence réelle de ce théorème identifie ce point à \(g\). Le lemme 3.2 fournit, dans un voisinage analytique de \(g\), une itérée finie de renormalisation avec un domaine de type quadratique et un module uniformément positif. Les références précises sont les sections 2.1–2.3, p. 736–738, et la section 3.1, lemme 3.2, p. 745, de \href{https://annals.math.princeton.edu/wp-content/uploads/annals-v164-n3-p01.pdf}{de Faria–de Melo–Pinto, \emph{Global hyperbolicity of renormalization for \(C^r\) unimodal mappings}}.

Pour vérifier l’appartenance au voisinage de ce lemme, soit \(\Omega_\epsilon\) un voisinage complexe suffisamment petit de \([-1,1]\), avec
\[
\rho_\epsilon:=\sup_{x\in\Omega_\epsilon}
\left|\frac{x^2-1}{2.5}\right|<1,
\qquad M_\epsilon:=\sup_{x\in\Omega_\epsilon}|x|^2<\infty.
\]
La restriction de la carte certifiée satisfait
\begin{equation}
\|\Delta f\|_{\Omega_\epsilon}
\le M_\epsilon\left(\frac{|\Delta u|}{10}
          +\rho_\epsilon\|\Delta\nu\|_1\right)
\le C_\epsilon\|\Delta f\|_{\mathcal A}.
\label{hmtfg:clasificacion:eq:16}
\end{equation}
Le certificat donne
\[
\|R^m\Gamma(\lambda_*)-g\|_{\mathcal A}
<0.001666(0.83996)^m.
\]
On choisit donc un entier fixe \(m\) qui place cette image dans le voisinage du lemme 3.2. La continuité de la composition finie fournit un intervalle \(I_*\) autour de \(\lambda_*\) où les mêmes opérations sont définies. Si \(N\) est l’itérée fournie par ce lemme, nous posons
\begin{equation}
d=4+m+N,\qquad
\mathcal F_\lambda=R^df_\lambda,\quad \lambda\in I_*.
\label{hmtfg:clasificacion:eq:17}
\end{equation}
La famille \(\mathcal F\) possède des représentants de type quadratique de degré deux, sur un domaine commun et avec un module uniformément positif. Sa dépendance analytique réelle admet une complexification locale du paramètre. Le nombre \(d\) est fixe et compte les doublements effectués avant d’appliquer le résultat paramétrique.

Cette construction fixe un domaine commun et un nombre fini de retours \(d\), suffisants pour appliquer le théorème paramétrique.

\subsubsection{Transport de la transversalité certifiée}
\label{hmtfg:clasificacion:8.2}

La direction de \(\Gamma\) satisfait la condition de cône
\[
\|\dot\nu\|_1\le\tfrac14|\dot u|,\qquad \dot u\ne0.
\]
Les bornes différentielles certifiées impliquent l’invariance de ce cône et
\begin{equation}
|\dot u_{k+1}|\ge4.381|\dot u_k|
\label{hmtfg:clasificacion:eq:18}
\end{equation}
le long de l’orbite de \(\Gamma(\lambda_*)\). Tous ses points restent dans le domaine certifié. En particulier, la direction transportée par les étapes finies de \eqref{hmtfg:clasificacion:eq:17} conserve une composante expansive non nulle.

La classe hybride du point fixe coïncide avec sa variété stable dans l’espace des germes de type quadratique : \href{https://www.maths.tcd.ie/EMIS/journals/Annals/149_2/lyubich.pdf}{Lyubich, \emph{Feigenbaum–Coullet–Tresser universality and Milnor's hairiness conjecture}, théorème 6.1, p. 371–372}. La transversalité reste valable lors du passage à cet espace, comme on peut le vérifier sans identifier des normes de Banach différentes. En effet, si \(\partial_\lambda\mathcal F_{\lambda_*}\) était tangente à la classe hybride, elle serait la tangente d’un disque analytique contenu dans cette classe. La convergence exponentielle uniforme de la renormalisation sur un sous-disque compact, suivie de l’estimation de Cauchy en son paramètre, ferait converger vers zéro l’évaluation en \(x=1\) de ses tangentes renormalisées. Cependant, dans la carte précédente,
\[
\dot f(1)=-\dot u/10,
\]
et \eqref{hmtfg:clasificacion:eq:18} fait croître sa valeur absolue. Les deux dérivées représentent le même germe et ont la même évaluation réelle. La contradiction démontre
\begin{equation}
\partial_\lambda\mathcal F_{\lambda_*}
\notin T_{\mathcal F_{\lambda_*}}\mathcal H_{c_\infty}.
\label{hmtfg:clasificacion:eq:19}
\end{equation}
Ainsi, \(S=\{\mathcal F_\lambda\}\) est une transversale analytique complexe à la classe hybride de la limite de doublement.

\subsubsection{Identification à partir d’un certain rang des premières racines par leur période}
\label{hmtfg:clasificacion:8.3}

Notons \(c_r\) le centre quadratique canonique de période \(2^r\), obtenu par \(r\) doublements successifs du centre de période un. Ce symbole désigne le paramètre de la famille quadratique de reconnaissance ; il se distingue des fonctions orbitales \(c_j(\lambda)=f_\lambda^j(0)\) des sections précédentes. On écrit \(c_r^{\rm quad}\) pour conserver cette distinction.

Le théorème de continuation globale et \eqref{hmtfg:clasificacion:eq:H} donnent \(\lambda_n\to\lambda_*\), d’où \(\lambda_n\in I_*\) pour \(n\) suffisamment grand. La classification et les intervalles restrictifs de la section~\ref{hmtfg:clasificacion:2} montrent que, lorsque \(n>d\),
\begin{equation}
\mathcal F_{\lambda_n}=R^df_{\lambda_n}
\quad\text{a une période critique exacte }2^{n-d}
\quad\text{et itinéraire }W_{n-d}0.
\label{hmtfg:clasificacion:eq:20}
\end{equation}
Son orbite critique reste dans l’intervalle réel invariant, contenu dans le domaine de type quadratique ; son ensemble de Julia rempli est donc connexe. Le redressement conserve l’orbite postcritique et les branches réelles. En parcourant ses retours restrictifs, on obtient \(n-d\) doublements canoniques et, à la fin, un point critique fixe. Ce dernier se redresse au centre \(0\) ; les inverses successifs du redressement des copies de doublement donnent l’unique centre \(c_{n-d}^{\rm quad}\). Il s’agit de l’application concrète du paramétrage par combinatoire et de l’homéomorphisme de redressement des copies, décrits dans Lyubich, sections 5.1–5.3, p. 362–364. On conclut
\begin{equation}
\mathcal F_{\lambda_n}\in\mathcal H_{c_{n-d}^{\rm quad}}.
\label{hmtfg:clasificacion:eq:21}
\end{equation}

Le théorème 7.4 de Lyubich, p. 387–388, donne une unique intersection locale de \(S\) avec chaque classe \(\mathcal H_{c_r^{\rm quad}}\) pour \(r\) suffisamment grand. Ses hypothèses de bornes complexes a priori sont satisfaites pour la cascade quadratique réelle de doublement, conformément au théorème 5.6, p. 367, du même travail. Notons \(\zeta_r\) le paramètre de cette intersection. Les équations \eqref{hmtfg:clasificacion:eq:20}–\eqref{hmtfg:clasificacion:eq:21}, la convergence vers le croisement et l’unicité locale impliquent
\begin{equation}
\boxed{\lambda_n=\zeta_{n-d}\quad\text{pour tout }n\text{ suffisamment grand}.}
\label{hmtfg:clasificacion:eq:22}
\end{equation}
L’unicité utilisée ici concerne \textbf{chaque classe hybride canonique} dans un voisinage fixe du croisement. Elle ne se déduit ni de la seule existence d’une cascade locale ni de l’unicité du croisement avec la feuille stable. L’équation \eqref{hmtfg:clasificacion:eq:22} identifie le sélecteur initial, déjà défini globalement, aux intersections locales ; elle n’introduit pas de sélecteur de substitution.

En particulier, soit \(\mu_j\) la cascade locale de la section \ref{hmtfg:escalado:8}, et soit \(s\) l’entier fixe déterminé par sa période physique exacte :
\[
\operatorname{per}_{\rm crit}(f_{\mu_j})=2^{j+s}.
\]
Dans la présentation avec la courbe \(R^4f_\lambda\) et la cible de période deux, \(s=5\) ; les étapes fixes supplémentaires utilisées pour localiser la cible modifient \(s\). Par \eqref{hmtfg:clasificacion:eq:21} et la même unicité,
\begin{equation}
\mu_j=\zeta_{j+s-d},\qquad
\boxed{\lambda_n=\mu_{n-s}\quad(n\text{ suffisamment grand}).}
\label{hmtfg:clasificacion:eq:23}
\end{equation}
Le décalage \(d\) de la préparation de type quadratique et le décalage \(s\) des périodes de la cible sont des objets distincts. L’identification s’effectue par la période physique, et non en égalant ces deux entiers.

\subsubsection{Reste en puissance par holonomie transversale}
\label{hmtfg:clasificacion:8.4}

Le lemme 7.3 de Lyubich, p. 384, avec démonstration p. 384–387, apporte une régularité \(C^{1+\beta}\)-conforme, pour un certain \(\beta>0\), de l’holonomie entre transversales à la classe hybride du point de Feigenbaum. Sa définition p. 384 concerne les lieux de connexité sur les transversales ; ils contiennent notamment les centres considérés ici. Nous pouvons réduire l’exposant de sorte que \(0<\beta\le1\).

Soit \(h\) l’holonomie de \(S\) vers la variété instable et soit \(z\) une coordonnée analytique qui linéarise la restriction unidimensionnelle de la renormalisation :
\begin{equation}
z(g)=0,\qquad z(Rq)=\delta_{\mathrm F} z(q).
\label{hmtfg:clasificacion:eq:24}
\end{equation}
L’existence de cette coordonnée s’obtient en appliquant la linéarisation d’un germe holomorphe répulsif, de multiplicateur \(\delta_{\mathrm F}>1\), à la variété instable analytique. La direction expansive de \eqref{hmtfg:clasificacion:eq:18}, l’invariance du cône et le retour stationnaire identifient ce multiplicateur à la valeur propre instable positive employée dans le certificat.

Si \(q_r\) est le centre de la classe \(\mathcal H_{c_r^{\rm quad}}\) situé sur la variété instable, la compatibilité du redressement et de la renormalisation donne \(Rq_r=q_{r-1}\). Cette identité est celle utilisée dans la démonstration du théorème 7.4, p. 388. Par \eqref{hmtfg:clasificacion:eq:24}, il existe \(b\ne0\), qui absorbe tout indice initial fixe, tel que
\begin{equation}
z(q_r)=b\delta_{\mathrm F}^{-r}
\label{hmtfg:clasificacion:eq:25}
\end{equation}
pour tout \(r\) suffisamment grand.

La régularité de \(h\) et la transversalité \eqref{hmtfg:clasificacion:eq:19} donnent, sur le lieu de connexité proche de \(\lambda_*\),
\begin{equation}
H(\lambda):=z\bigl(h(\mathcal F_\lambda)\bigr)
=A(\lambda-\lambda_*)
+O\!\left(|\lambda-\lambda_*|^{1+\beta}\right),
\qquad A\ne0.
\label{hmtfg:clasificacion:eq:26}
\end{equation}
Puisque \(H(\zeta_r)=b\delta_{\mathrm F}^{-r}\), l’estimation inférieure
\(|H(\lambda)|\ge |A|\,|\lambda-\lambda_*|/2\) dans un voisinage suffisamment petit implique d’abord
\(|\zeta_r-\lambda_*|=O(\delta_{\mathrm F}^{-r})\). En substituant cette borne dans le reste de \eqref{hmtfg:clasificacion:eq:26}, on obtient
\begin{equation}
\zeta_r-\lambda_*=\frac bA\delta_{\mathrm F}^{-r}
+O\!\left(\delta_{\mathrm F}^{-(1+\beta)r}\right).
\label{hmtfg:clasificacion:eq:27}
\end{equation}
Nous appliquons \eqref{hmtfg:clasificacion:eq:22} avec \(r=n-d\). Le décalage fini modifie le coefficient principal et la constante du reste. Puisque \(\lambda_n<\lambda_*\), on obtient \eqref{hmtfg:clasificacion:eq:14}, avec
\[
C=-\frac bA\delta_{\mathrm F}^d>0,
\qquad \theta=\delta_{\mathrm F}^{-\beta}\in(0,1).
\]
Enfin,
\[
\lambda_n-\lambda_{n-1}
=C(\delta_{\mathrm F}-1)\delta_{\mathrm F}^{-n}\bigl(1+O(\theta^n)\bigr),
\]
ce qui démontre \eqref{hmtfg:clasificacion:eq:15}. \(\square\)


\subsubsection{Portée quantitative du théorème}
\label{hmtfg:clasificacion:8.5}

L’identification des limites démontrée par première sortie active le transfert métrique pour les premières racines globales du secteur uniforme. Le taux opératoriel \(0.83996\) contrôle les applications renormalisées sur la feuille stable. Le taux paramétrique \(\theta=\delta_{\mathrm F}^{-\beta}\) découle de la régularité de l’holonomie transversale. Les deux estimations conservent leur variable et leur domaine.

Le théorème établit l’existence de \(\beta,\theta,C\), du voisinage \(I_*\), du nombre fini de retours \(d\) et d’un indice \(n_0\) à partir duquel les centres coïncident. Leurs valeurs numériques individuelles ne sont pas attribuées dans cet énoncé asymptotique. Cette formulation complète la limite et son reste en puissance ; l’évaluation numérique de ces témoins constitue une spécialisation quantitative distincte.

\section{Retour, mémoire et universalité de la cascade dodécaphasique}
\label{sec:hmt-feig-interpretacion}

\subsection{La structure génératrice et la lecture critique}

La construction de Feigenbaum en Holographie Modulaire Triadique relie une
organisation arithmétique orientée à une loi asymptotique de changement d’échelle.
Son point de départ est la composition effective d’APP, TRIT et TPK. APP évalue
la somme et le produit sur les deux feuilles du réseau de chiffres et conserve les
quotients et les restes de ces opérations. TRIT détermine le régime local et
l’orientation de sa lecture. TPK sélectionne et transporte les états,
actualise leurs registres et prolonge leurs routes avec retenue, frontière et mémoire.
La structure discrète conjointe du continu fournit la résidence commune
de ces états et des lecteurs qui agissent sur eux.

Le prolongement
\[
w_6\longrightarrow w_{12}\longrightarrow w_{18}
\longrightarrow w_{24}\longrightarrow w_{30}
\longrightarrow R_{36}\longrightarrow G_9
\]
organise cette persistance. La graine initiale est relevée en conservant ses
préfixes ; les relèvements successifs changent le régime sans perdre l’orientation
ni la mémoire ; la frontière coinductive permet de poursuivre la construction, et les
neuf cartes réunissent ses lectures locales. L’holonomie nonadique exprime le
retour de phase accompagné d’un accroissement de mémoire. Ainsi, un tour
visible conserve sa position dans l’histoire complète. L’opération de retour
acquiert dès l’origine deux composantes inséparables :
réitération d’une règle et conservation de sa provenance.

Le lecteur critique fixé sur la roue dodécaphasique orientée convertit cette
structure en douze phases avec une trace uniforme. Leurs représentants sont
\[
\phi_m=\frac{(3m-1)\pi_{\mathrm{HMT}}}{18},
\qquad w_m=\frac1{12},\qquad 1\le m\le12.
\]
Le choix du représentant orienté importe parce que le deuxième moment
utilise sa grandeur. Le quotient angulaire enregistre la phase, tandis que le
représentant conserve l’information nécessaire à sa normalisation. La
coordonnée \(\pi_{\mathrm{HMT}}\) provient des lecteurs coinductifs du même
noyau APP--TRIT--TPK. Dans son image corrélée, les coordonnées
\(\pi,\varphi,e\) et la clôture dodécaphasique de \(\alpha_{\mathrm{HMT}}\)
conservent leur ordre constructif et leurs domaines propres. Dans le lecteur critique
considéré ici, on utilise la coordonnée angulaire déjà générée.

La normalisation est déterminée par une somme entière :
\[
\sum_{m=1}^{12}(3m-1)^2=5394=6\cdot899,
\qquad k_m=\frac{2(3m-1)}{\sqrt{899}}.
\]
Le facteur commun \(\pi_{\mathrm{HMT}}\) se simplifie lors de la fixation du deuxième moment.
Cette simplification conserve la provenance des phases et produit une
représentation particulièrement précise de leur lecture normalisée. On obtient
\[
u_0(x)=\frac1{12}\sum_{m=1}^{12}\bigl(1-\cos(k_mx)\bigr),
\qquad f_\lambda(x)=1-\lambda u_0(x),
\qquad \sum_mw_mk_m^2=2.
\]
Le profil est pair et entier ; il satisfait \(u_0(0)=u_0'(0)=0\) et
\(u_0''(0)=2\). Pour \(0<\lambda\le2/u_0(1)\), la famille envoie
\([-1,1]\) dans lui-même et possède un maximum critique quadratique avec
\(f_\lambda''(0)=-2\lambda\). Le paramètre \(\lambda\) parcourt cette famille
après fixation du lecteur, de l’orientation, des poids et de l’échelle. Le nombre
universel qui sera reconnu à la fin est ainsi séparé de la sélection
des coefficients.

\subsection{Douze phases, moments et représentation spectrale}

Le profil admet une deuxième expression qui réunit ses phases dans un opérateur
fini. Les nœuds
\[
s_m=k_m^2=\frac{4(3m-1)^2}{899}
\]
déterminent la mesure positive \(\mu_0=12^{-1}\sum_m\delta_{s_m}\). Ses
moments \(M_r=\int s^r\,d\mu_0(s)\) définissent des formes de Hankel strictement
positives jusqu’au rang douze. Cette positivité provient de l’évaluation
d’un polynôme sur douze nœuds distincts : un polynôme de degré inférieur à
douze conserve au moins une évaluation non nulle. Au degré douze, le produit
des douze facteurs nodaux annule la norme et détermine exactement le rang
du lecteur.

L’orthogonalisation de ces moments produit une matrice de Jacobi
\(J_{12}\), réelle symétrique, à coefficients sous-diagonaux positifs. Son calcul fonctionnel
restitue le profil complet :
\[
u_0(x)=e_0^{\mathsf T}
\bigl[I-\cos(x\sqrt{J_{12}})\bigr]e_0.
\]
La somme de cosinus et l’opérateur de Jacobi sont deux réalisations du même
lecteur. Les moments maintiennent toutes les fréquences et leurs poids ; les
marques d’orientation, de route et de mémoire restent dans l’état qui les
engendre. Le rang douze identifie la dimension spectrale de cette lecture
particulière. La profondeur de ses prolongements relève d’une autre dimension
de la construction : la longueur des histoires compatibles et la résolution
avec laquelle elles sont évaluées.

La géométrie analytique est elle aussi explicitée. Pour la déformation
pondérée \(u_\eta\), dans la bande \(|\eta|\le1/40\), les poids restent
positifs et le deuxième moment demeure normalisé. Dans
\(|\operatorname{Re}z|<\pi/3\), il existe une fonction holomorphe impaire et univalente
\(b_\eta\), normalisée par \(b_\eta'(0)=1\), telle que
\(u_\eta=b_\eta^2\). La roue produit donc une déformation conforme
du pli quadratique. Sa conjugaison dynamique conserve cette déformation
à travers \(b_\eta(1-\lambda w^2)\). Le théorème de changement d’échelle paramétrique
développé ici concerne le membre uniforme \(\eta=0\) ; les
propriétés structurelles de la bande pondérée et la persistance locale de
ses racines conservent leurs quantificateurs spécifiques.

\subsection{Profondeur de lecture et restitution des histoires}

La profondeur nonadique affine la lecture d’un paramètre, tandis que le
doublement dynamique change l’horizon de retour. Le raffinement avec
retenue exprime la première opération par
\[
L_c(x,y)=(Bx-c,By+c),\qquad
C_m=\sum_{j=1}^m c_jB^{m-j},
\]
\[
(x_m,y_m)=(B^mx_0-C_m,B^my_0+C_m).
\]
Le support extérieur croît comme \(B^m\), et le cylindre normalisé du registre
a pour diamètre \(B^{-m}\). La paire entière, la profondeur et la marque de
frontière conservent la reconstruction des quotients, des restes et des préfixes.
Croissance extérieure et résolution intérieure décrivent ainsi des opérations
compatibles sur une même histoire.

La composition exacte
\(L_D^{[b]}L_C^{[a]}=L_{B^bC+D}^{[a+b]}\) permet de regrouper une histoire en
blocs sans altérer ses données. Un retour de doublement réunit deux
transitions ; le calendrier nonadique en réunit neuf. Les deux regroupements sont
comparés sur un horizon commun de \(9\,2^N\) événements, en conservant
l’ordre des sous-blocs et leurs retenues.

La contribution de la mémoire peut s’observer directement dans les branches
inverses du profil. Si \(v=(u_0|_{[0,1]})^{-1}\), les deux branches sont
\[
\beta_\sigma(y)=\sigma v\!\left(\frac{1-y}{\lambda}\right),
\qquad \sigma\in\{-1,+1\}.
\]
La valeur finale, jointe à la dernière feuille, reconstruit la valeur précédente. Une
suite de feuilles permet de restituer toutes les étapes intermédiaires. Le mot
des branches accompagne les registres APP--TRIT--TPK de route, de frontière, de retenue et de
mémoire ; chaque registre conserve sa fonction. Cette restitution concrète
explique quelle information exige l’inversion de la lecture critique.

La résolution de chaque retour dispose également d’une borne explicite. Avec
\(t=\lambda u_\eta(1)/2\), \(z=(x+1)/2\) et
\[
F_{t,\eta}(z)=1-t\frac{u_\eta(2z-1)}{u_\eta(1)},
\]
on obtient
\[
\left|F_{t,\eta}^{q}(1/2)-F_{s,\eta}^{q}(1/2)\right|
\le S_q|t-s|,\qquad S_q=\sum_{j=0}^{q-1}(6\sqrt2)^j<9^q.
\]
Un cylindre de base neuf et de profondeur \(m=2^N+r\) permet donc de lire
le retour d’horizon \(2^N\) avec un diamètre inférieur à \(9^{-r}\). Cette
estimation donne un contenu quantitatif à la profondeur arbitraire : pour chaque
horizon et chaque précision demandés, elle détermine une profondeur suffisante.
L’évaluation des fonctions analytiques et l’arrondi possèdent leurs propres
restes, qui se propagent avec l’incertitude du paramètre.

Les vacances interviennent dans le bilan de capacité du registre. La
comparaison entre blocs de 729 et 1000 possibilités distingue la capacité
accumulée de la longueur chronologique, en conservant les étapes où la première
reste constante. Cette comptabilité détermine l’espace requis pour
conserver une lecture et sa profondeur. L’admissibilité des histoires
découle, quant à elle, des règles de survie. Les deux opérations
collaborent à la construction et maintiennent des objets mathématiques distincts.

\subsection{Chronologie nonadique et combinatoire de doublement}

Le mot critique de doublement possède une expression à toute profondeur :
\[
\tau_j=(-1)^{v_2(j)},\qquad
\tau_{2j}=-\tau_j,\qquad \tau_{2j+1}=+1.
\]
Ses préfixes satisfont \(W_{n+1}=W_n(-1)^nW_n\). La règle détermine comment
les deux copies du retour sont orientées et où apparaît leur séparation. La
lecture nonadique de la chronologie conserve l’entier complet
\(K=\rho_9(K)+9q_9(K)\). Lorsqu’on le réduit modulo \(2^n\), l’holonomie qui
avance de neuf événements induit la translation de neuf sur ce cycle. Puisque
neuf est impair, cette translation visite toutes ses positions. La
permutation \(j\mapsto9j\pmod{2^n}\), qui la conjugue à l’avancée unitaire,
conserve la valuation dyadique de chaque position non nulle. Le reste et le quotient
chronologiques soutiennent conjointement cette compatibilité.

La mémoire bilatérale possède également une réalisation opératorielle compatible
avec le raffinement. Sur les cycles de longueur \(2^n\), l’opérateur
d’avancée \(C_n\) est incorporé dans
\[
\mathbb U_n=P_0\otimes I+P_1\otimes C_n,
\]
avec les projecteurs complémentaires du lecteur nonadique. À la limite, ses
composantes visible et complémentaire satisfont
\[
\|T_{\infty,a}v\|^2+\|D_{\infty,a}v\|^2=\|v\|^2,
\qquad
v=T_{\infty,a}^*T_{\infty,a}v+D_{\infty,a}^*D_{\infty,a}v.
\]
Ces identités décrivent la conservation et la restitution dans l’espace
et la représentation déclarés. La contraction des applications renormalisées,
qui apparaîtra ensuite, mesure une autre opération : le rapprochement de leurs profils
vers une forme stable. L’histoire conservée par le relèvement et la
convergence d’une lecture fonctionnelle s’articulent à des niveaux différents.

\subsection{La carte ordonnée et la sélection globale des racines}

La lecture archimédienne du continu HMT transporte les opérations
par un isomorphisme de corps ordonnés complets, dans l’univers
déclaré \(\mathfrak G\) :
\[
\iota:\mathbb R_{\mathrm{HMT}}^{\mathfrak G}
\longrightarrow\mathbb R^{\mathfrak G}.
\]
La construction interne du cosinus par sa série, de la racine positive de 899,
du profil et de ses itérées précède la sélection des zéros. La conservation
de la somme, du produit, de l’ordre et des limites convergentes donne
\[
\iota\bigl(S_n^{\mathrm H}(a)\bigr)=S_n\bigl(\iota(a)\bigr),
\qquad S_n(\lambda)=f_\lambda^{2^n}(0).
\]
La surjectivité couvre toutes les racines du domaine de cette carte ; l’ordre
conserve les retours antérieurs, la période exacte et la condition de
se trouver à droite de la racine précédente. On transporte ainsi
également le minimum de chaque ensemble de candidats. La généalogie atteint
la sélection effective du paramètre, la deuxième projection et ses
marques étant conservées sur le même antécédent.

Le sélecteur est fixé par \(\lambda_1=1/u_0(1)\) et par la première racine
nouvelle de période critique exacte \(2^n\) située à droite de
\(\lambda_{n-1}\). La classification des itinéraires démontre que chaque racine
sélectionnée a pour mot \(W_n0\). Les retours restrictifs réduisent sa
période de moitié en conservant l’ordre unimodal, et l’induction restitue
le mot complet. L’analyticité des retours non identiquement nuls
isole leurs zéros dans chaque compact ; la
comparaison des signes avant le premier paramètre d’itinéraire infini
garantit l’existence d’une racine nouvelle à tous les niveaux.

Soit \(\Lambda_\tau\) l’ensemble compact non vide des paramètres qui réalisent
le mot infini, et soit \(a_*\) son minimum. On démontre
\[
\lambda_n\nearrow a_*.
\]
Le passage à la limite conserve chaque signe au moyen d’une borne uniforme à horizon
fixe. Si \(c_p=f_\lambda^p(0)\), les signes opposés de \(c_p\) et
\(c_{2p}\), avec \((f_\lambda^p)'(0)=0\), impliquent
\[
|c_p|>\frac2{B_p},\qquad
B_p=16^p\frac{16^p-1}{15}.
\]
Ainsi, les préfixes de périodes croissantes préservent le mot infini en
s’accumulant. La preuve utilise la suite initiale des premières racines dès
sa définition globale.

\subsection{La feuille stable et l’identification de la limite sélectionnée}

Le retour normalisé
\[
(Rf)(x)=-\frac{f(f(-ax))}{a},\qquad a=-f(1)>0,
\]
réunit deux transitions, change l’échelle et conserve l’orientation nécessaire
pour restituer un maximum de valeur un. Les domaines restrictifs établissent
son interprétation comme retour réel. La dérivée de cet opérateur inclut
la variation de l’échelle \(a\) ; cette contribution détermine la sensibilité
paramétrique de la transformation complète.

Dans la carte \(f(x)=1-x^2V((x^2-1)/(5/2))\), de coordonnées \(u=10v_0\)
et \(\nu=(v_1,v_2,\ldots)\), le quatrième retour de la famille entre dans un
voisinage quantifié du point fixe \(g\). La courbe traverse transversalement
sa feuille stable en un unique paramètre local
\[
1.874038<\lambda_*<1.874039,
\]
et satisfait
\[
\|R^{4+n}f_{\lambda_*}-g\|_{\mathcal A}
<0.001666(0.83996)^n.
\]
La feuille stable est un graphe de pente au plus \(0.16\). Les
sécantes paramétriques appartiennent à un cône transversal expansif. Les deux
propriétés distinguent la diminution des déformations stables de la
croissance de la séparation entre paramètres.

L’identification de \(a_*\) et de \(\lambda_*\) réunit cette géométrie locale
et le minimum global. On exclut d’abord l’intervalle entre la huitième racine
et l’extrémité inférieure de la boîte locale. Ensuite, l’ordre
\(a_*\le\lambda_*\) permet d’étudier une séparation hypothétique au moyen
du cône négatif. Tant que sa coordonnée transversale a une valeur absolue inférieure
à \(0.006\), la sécante reste dans le domaine analytique et croît d’un
facteur compris entre \(4.4034\) et \(8.0887\). Une séparation persistante
produirait donc une première sortie d’amplitude comprise entre
\(0.006\) et \(0.0485322\).

Le résidu de l’orbite stable conserve une information plus précise que sa
norme : ses composantes satisfont \(|e_u|\le0.16\|e_\nu\|_1\). Cette
relation permet de borner conjointement leurs effets sur l’orbite critique.
Le recouvrement de toute la région de première sortie montre dans chaque cas un
signe incompatible avec \(\tau\), ou un retour contractant qui impose à deux
valeurs critiques dyadiques successives d’avoir le même signe. Le mot
infini exige des signes opposés. La contradiction identifie les paramètres :
\[
\lambda_n\nearrow a_*=\lambda_*.
\]
La démonstration combine une induction valable à toute profondeur avec une
exclusion finie uniforme de la région qu’atteindrait toute
séparation. La portée infinie repose sur cette composition.

\subsection{Changement d’échelle universel et signification de la précision}

La réalisation holomorphe de degré deux des retours et la transversalité
établie permettent de comparer leurs classes hybrides. Pour chaque période
suffisamment élevée, la classe canonique possède une unique intersection locale
avec la famille. Les premières racines globales convergent déjà vers le croisement et
conservent la combinatoire de doublement ; elles coïncident donc à partir d’un certain rang
avec les centres locaux lorsque leurs périodes physiques sont égalées. La preuve
transporte la sélection initiale jusqu’à la loi métrique universelle.

Il existe \(C>0\), \(M<\infty\) et \(0<\theta<1\) tels que
\[
\lambda_*-\lambda_n=C\delta_{\mathrm F}^{-n}(1+e_n),
\qquad |e_n|\le M\theta^n,
\]
et, par conséquent,
\[
\lim_{n\to\infty}
\frac{\lambda_{n-1}-\lambda_{n-2}}
{\lambda_n-\lambda_{n-1}}=\delta_{\mathrm F}.
\]
Le facteur \(\delta_{\mathrm F}\) est la valeur propre expansive de l’opérateur de doublement
reconnu à la fin de la construction. La famille particulière conserve ses
douze phases, ses moments et son histoire de provenance ; la limite de ses
rapports paramétriques appartient à la même classe universelle de criticité
quadratique. Cette relation explique conjointement la spécificité du
générateur et l’universalité de la sortie.

La précision possède ici trois résidences. Les intervalles des racines
contrôlent leur évaluation finie. L’estimation
\(0.001666(0.83996)^n\) contrôle la convergence opératorielle sur la feuille
stable. Le terme \(M\theta^n\) contrôle le transfert de cette géométrie à
l’échelle paramétrique au moyen de l’holonomie transversale. Chaque estimation
correspond à sa variable et à son opération. Le huitième quotient est encadré
autour de \(4.669212189150204154556096215889663928\), avec une largeur inférieure
à \(5.808\cdot10^{-37}\) ; cette largeur mesure l’évaluation de \(\delta_{\mathrm F,8}\).
La distance de \(\delta_{\mathrm F,8}\) à la limite exige de quantifier les constantes
du reste paramétrique, dont l’existence est déjà démontrée.

La notation maintient séparés \(\delta_{\mathrm F}\), facteur de changement d’échelle paramétrique
de Feigenbaum ; \(\alpha_{\mathrm F}\), constante spatiale associée à son
problème de renormalisation ; et \(\alpha_{\mathrm{HMT}}\), constante de
structure fine du corpus. Le présent résultat identifie \(\delta_{\mathrm F}\)
au moyen des premières racines du secteur uniforme \(\eta=0\). La bande
pondérée conserve ses théorèmes structurels et d’existence, et son extension
métrique quantifiée constitue un énoncé spécifique distinct.

L’apport de l’enchaînement réside dans la continuité entre génération,
lecture, mémoire, sélection et limite. L’arithmétique orientée détermine le
profil ; la représentation spectrale conserve ses douze composantes ; les
cylindres permettent de le résoudre avec une précision arbitraire à chaque horizon ;
le transport des histoires restitue ses parcours ; la carte ordonnée
conserve toutes les racines et leurs minima ; la classification prolonge le
sélecteur ; et l’expansion transversale, jointe à la contraction stable,
détermine sa loi asymptotique. L’universalité apparaît comme une propriété
démontrée de cette construction complète et de sa réalisation analytique.

\endgroup

% Fragmento compartido por VII y CRISTAL; incluir despues del cuerpo Feigenbaum.
% Procedencia: resultados recuperados, con demostraciones elementales reunidas.
% Perfil, Jacobi y escalado: 00_antecedentes_operatorios.tex, 03, 06 y 07
% del presente paquete; no se cambia su selector ni su dominio uniforme.
% Constitutiva: CRISTAL_ES/main_autosuficiente.tex:58738--58840,58864--58894.
% Inversion: CRISTAL_ES/sections/cr28_restitucion_constitutiva.tex:18--59.
% Catalan: el mismo propietario, 277--373; identidad integral, unicidad y serie.
% Corte material: CUMPLIMIENTO_EDITORIAL_20260928/fuentes/CRISTAL_ES.
% No contiene una nueva calibracion ni una nueva campana de calculo.
\section{Articulation des lecteurs critique, spectral et constitutif}
\label{hmtfg:articulacion:seccion}

La composition APP--TRIT--TPK précède les lecteurs de cette section.
APP conserve dans chaque cellule les évaluations additive et multiplicative,
avec leurs restes et leurs quotients ; la réduction tritique détermine le régime,
l’orientation et la retenue de la transition. La composition effective
de la sélection, du transport et de l’actualisation de la mémoire produit l’état
enrichi et ses prolongements, développés dans
\ref{hmtfg:app-trit-transporte} et
\ref{hmtfg:subsubsec:tres-vueltas-ext-ch}.
L’holonomie nonadique compose ces transitions : le retour de phase
fait avancer la mémoire et conserve l’histoire. Les lecteurs agissent sur la
même structure discrète conjointe du continu, avec ses cinq
constructions consubstantielles et les deux lectures du même
antécédent : évaluation ordonnée et incidence de frontière.

La roue orientée et les canaux angulaires sont des sorties de cette
composition. Leurs coordonnées HMT déjà générées sont utilisées dans des opérations
ultérieures, avec leurs marques et leurs domaines ; ce ne sont pas des nombres métrologiques
introduits pour sélectionner une réponse. La lecture ordonnée
permet d’évaluer les formules suivantes, tandis que l’état de provenance
conserve la feuille, la route, le reste, le quotient, la retenue, la frontière et la mémoire.
Les restitutions démontrées retrouvent l’objet de leur domaine
déclaré ; la représentation évaluée reste accompagnée de ces
registres.

\subsection{La séparation critique comme lecture asymptotique intrinsèque}
\label{hmtfg:articulacion:separacion}

Dans le secteur uniforme de la roue dodécaphasique, les représentants
orientés et leurs poids sont
\[
 \phi_m=\frac{(3m-1)\pi_{\mathrm{HMT}}}{18},\qquad
 w_m=\frac1{12},\qquad 1\le m\le12.
\]
Le deuxième moment est normalisé par
\(\sum_{m=1}^{12}(3m-1)^2=5394=6\cdot899\).
Ainsi, le facteur angulaire commun se simplifie et il reste
\begin{equation}
 k_m=\frac{2(3m-1)}{\sqrt{899}},\qquad
 u_0(z)=\frac1{12}\sum_{m=1}^{12}\bigl(1-\cos(k_mz)\bigr),
 \qquad f_\lambda(z)=1-\lambda u_0(z).
 \label{hmtfg:articulacion:perfil}
\end{equation}
La variable \(\lambda\) parcourt la famille après fixation du
lecteur, de ses poids et de la normalisation. Dans
\(0<\lambda\le2/u_0(1)\), le domaine dynamique est \([-1,1]\).
La première racine est \(\lambda_1=1/u_0(1)\) ; pour \(n\ge2\),
\(\lambda_n\) est la plus petite racine supérieure à \(\lambda_{n-1}\)
dont le point critique a pour période exacte \(2^n\).
La classification, l’existence de chaque minimum et l’identification
métrique de ces centres ont été démontrées dans
\ref{hmtfg:escalado:15}--\ref{hmtfg:escalado:17}.

\begin{proposition}[Lecture structurelle de la constante de changement d’échelle]
\label{hmtfg:articulacion:delta}
Pour le lecteur uniforme \eqref{hmtfg:articulacion:perfil} et son sélecteur
de premières racines, la constante de séparation critique est
\begin{equation}
 \delta_{\mathrm{HMT,crit}}
 :=\lim_{n\to\infty}
 \frac{\lambda_n-\lambda_{n-1}}{\lambda_{n+1}-\lambda_n}
 =\delta_{\mathrm F}.
 \label{hmtfg:articulacion:limite}
\end{equation}
C’est l’invariant asymptotique de séparation des paramètres des retours
critiques de ce lecteur, dont l’ordre de sélection est conservé.
\end{proposition}
\begin{proof}
Le théorème de changement d’échelle \eqref{hmtfg:escalado:eq:67} fournit
\(\lambda_*-\lambda_n=C\delta_{\mathrm F}^{-n}(1+e_n)\),
avec \(C>0\) et \(e_n\to0\). En soustrayant deux termes consécutifs,
\[
 \lambda_n-\lambda_{n-1}
 =C\delta_{\mathrm F}^{-n}
   \bigl[\delta_{\mathrm F}-1+
          \delta_{\mathrm F}e_{n-1}-e_n\bigr].
\]
Le quotient de cette égalité et de celle d’indice \(n+1\) converge vers
\(\delta_{\mathrm F}\). L’identité à partir d’un certain rang des centres locaux
et des premières racines, démontrée dans
\eqref{hmtfg:escalado:eq:66}, permet d’appliquer le changement d’échelle précisément
à la suite sélectionnée, et non à une autre cascade choisie après coup.
\end{proof}

Cette formulation conserve plus d’information que le numéral terminal :
elle identifie la roue, sa lecture, la famille et le sélecteur dont la limite
est évaluée. Son interprétation est intrinsèque à cette composition HMT ;
la reconnaissance de la constante universelle utilise ensuite le théorème
analytique appliqué au lecteur déjà construit.

\subsection{La mesure finie et la restitution spectrale du profil}
\label{hmtfg:articulacion:jacobi}

Les fréquences au carré produisent la mesure de probabilité atomique
\begin{equation}
 s_m=k_m^2=\frac{4(3m-1)^2}{899},\qquad
 \nu_0=\frac1{12}\sum_{m=1}^{12}\delta_{s_m},\qquad
 M_j=\int s^j\,d\nu_0(s).
 \label{hmtfg:articulacion:medida-finita}
\end{equation}
Ici, \(\delta_{s_m}\) est la masse de Dirac en un nœud, distincte
de la constante \(\delta_{\mathrm F}\). La mesure conserve les
douze fréquences et leurs poids. Pour un polynôme non nul de degré inférieur
à \(r\le12\),
\[
 \int p(s)^2\,d\nu_0(s)
 =\frac1{12}\sum_{m=1}^{12}p(s_m)^2>0,
\]
car un polynôme de degré inférieur à douze ne peut s’annuler aux
douze nœuds distincts. Ainsi, les matrices de Hankel
\((M_{i+j})_{0\le i,j<r}\) sont définies positives.
Le polynôme \(\prod_m(s-s_m)\) est de norme nulle ; l’espace
\(L^2(\nu_0)\) est de dimension douze.

\begin{proposition}[Représentation exacte du lecteur critique]
\label{hmtfg:articulacion:representacion}
Soient \(\psi_0,\ldots,\psi_{11}\) les polynômes orthonormaux
à coefficient principal positif pour \(\nu_0\), avec
\(\psi_0=1\). La multiplication par \(s\) a pour matrice de
Jacobi \(J_{12}\), réelle symétrique à coefficients sous-diagonaux positifs.
Avec \(e_0=(1,0,\ldots,0)^{\mathsf T}\),
\begin{equation}
 u_0(z)=\int\bigl(1-\cos(z\sqrt{s})\bigr)\,d\nu_0(s)
 =e_0^{\mathsf T}\bigl[I-\cos(z\sqrt{J_{12}})\bigr]e_0.
 \label{hmtfg:articulacion:perfil-jacobi}
\end{equation}
\end{proposition}
\begin{proof}
La positivité précédente permet d’orthogonaliser jusqu’au degré onze.
Si \(i<j-1\), l’autoadjonction de la multiplication donne
\(\langle s\psi_j,\psi_i\rangle
=\langle\psi_j,s\psi_i\rangle=0\).
Par symétrie, il ne reste que la diagonale et ses deux voisines ; le rapport
positif des coefficients principaux fixe la positivité de la
sous-diagonale. Définissons
\[
 Q_{mj}=\frac{\psi_j(s_m)}{\sqrt{12}},\qquad
 D=\operatorname{diag}(s_1,\ldots,s_{12}).
\]
L’orthonormalité donne \(Q^{\mathsf T}Q=I\), et l’évaluation de la
récurrence aux nœuds donne \(DQ=QJ_{12}\). Puisque \(Q\) est
carrée, \(J_{12}=Q^{\mathsf T}DQ\). Ses valeurs propres sont
positives, de sorte que la racine positive et le cosinus se calculent dans
cette diagonalisation. Enfin,
\(Qe_0=(1,\ldots,1)^{\mathsf T}/\sqrt{12}\), ce qui prouve
\eqref{hmtfg:articulacion:perfil-jacobi}.
\end{proof}

La réalisation par \((J_{12},e_0)\) restitue le même profil ;
ainsi, en conservant la famille et le sélecteur, elle restitue également
ses retours et ses premières racines. La positivité spectrale
justifie cette représentation. La transversalité paramétrique
utilise en outre la dérivée du retour et ses produits orbitaux,
développés dans \ref{hmtfg:escalado:14.3} et
\ref{hmtfg:escalado:6}. Ce sont des propriétés d’opérations différentes :
la démonstration du changement d’échelle conserve les deux et ne remplace pas le
contrôle de la tangente par la positivité des moments.

\subsection{Les canaux constitutifs et leur inversion étiquetée}
\label{hmtfg:articulacion:constitutiva}

Le lecteur constitutif agit sur un autre objet typé du même
support : le transport angulaire à deux feuilles. Ses coordonnées déjà
générées sont \(x=A_{\rm rad}\) et \(y=C^*_{\rm rad}\), dans la
chambre contractante \(x>|y|\). L’involution des feuilles
\(R=R^*=R^{-1}\) détermine les projecteurs marqués
\(P_\pm=(I\pm R)/2\). La représentation du transport est
\begin{equation}
 q_+=\exp[-(x+y)],\qquad q_-=\exp[-(x-y)],\qquad
 T=q_+P_++q_-P_-,\qquad 0<q_\pm<1.
 \label{hmtfg:articulacion:canales}
\end{equation}
L’évaluation exponentielle réalise ces canaux après leur
construction angulaire. Les incidences \(90\) et \(120\) du
transport fixent la réponse
\[
 \mathcal V=(I-T^{90})(I-T^{120})^{-1}.
\]
L’inverse existe parce que \(1-q_\pm^{120}>0\).
En écrivant \(s=q^{30}\), avec \(30=\gcd(90,120)\), on obtient
\begin{equation}
 f(s)=\frac{1-s^3}{1-s^4}
     =\frac{1+s+s^2}{1+s+s^2+s^3},\qquad
 r_\pm=f(q_\pm^{30}),\qquad
 \mathcal V=r_+P_++r_-P_-.
 \label{hmtfg:articulacion:respuesta}
\end{equation}
Dans cette sous-section, \(f\) désigne la réponse constitutive ;
\(f_\lambda\) continue de désigner l’application dynamique de
\eqref{hmtfg:articulacion:perfil}. Leurs arguments et leurs domaines
sont distincts.

\begin{proposition}[Restitution des canaux et de l’orientation]
\label{hmtfg:articulacion:inversion}
La fonction constitutive \(f:(0,1)\to(3/4,1)\) est une bijection
décroissante. Son inverse \(\psi\) associe à \(r\) l’unique racine
\(s\in(0,1)\) de
\begin{equation}
 r s^3+(r-1)(1+s+s^2)=0.
 \label{hmtfg:articulacion:cubica}
\end{equation}
Les lectures constitutives normalisées, avec leurs feuilles étiquetées,
\[
 \widehat\varepsilon=r_+^2,\qquad
 \widehat\mu=r_-^2,\qquad
 \widehat Z=\frac{r_-}{r_+},\qquad
 \widehat c=\frac1{r_+r_-},
\]
avec \(\widehat\varepsilon,\widehat\mu\in(9/16,1)\), permettent de restituer
\begin{equation}
 \begin{aligned}
 r_+&=\sqrt{\widehat\varepsilon},&
 r_-&=\sqrt{\widehat\mu},\\
 s_\pm&=\psi(r_\pm),& q_\pm&=s_\pm^{1/30}.
 \end{aligned}
 \label{hmtfg:articulacion:recuperacion-canales}
\end{equation}
\begin{equation}
 x=-\frac12\log(q_+q_-),\qquad
 y=\frac12\log\frac{q_-}{q_+}.
 \label{hmtfg:articulacion:recuperacion-angular}
\end{equation}
\end{proposition}
\begin{proof}
La dérivation du quotient rationnel donne
\begin{equation}
 f'(s)=-\frac{s^2(3+2s+s^2)}{(1+s+s^2+s^3)^2}<0
 \qquad(0<s<1).
 \label{hmtfg:articulacion:monotonia}
\end{equation}
Les limites \(f(0^+)=1\) et \(f(1^-)=3/4\), jointes à la
continuité, prouvent la bijectivité. Multiplier \(f(s)=r\) par
le dénominateur positif produit \eqref{hmtfg:articulacion:cubica}.
La positivité de \(r_\pm\) détermine les racines carrées
et celle de \(s_\pm\) détermine les racines trentièmes.
Par \eqref{hmtfg:articulacion:canales},
\(\log q_+=-x-y\) et \(\log q_-=-x+y\) ;
l’addition et la soustraction restituent \(x,y\).
Réciproquement, les racines restituées appartiennent à \((0,1)\),
de sorte que \(x\pm y=-\log q_\pm>0\) ; la paire obtenue
reste dans la chambre \(x>|y|\).
\end{proof}

En particulier,
\(\widehat\mu\widehat\varepsilon=\widehat c^{-2}\) et
\(\widehat\mu/\widehat\varepsilon=\widehat Z^2\).
L’échange des étiquettes permute \(q_+\) et \(q_-\),
conserve \(x\) et change le signe de \(y\). La restitution
exige cette orientation : le produit isolé \(r_+r_-\)
conserve la propagation commune, mais ne détermine pas les deux réponses
étiquetées. Les racines positives et l’inversion cubique restituent
les canaux à l’intérieur du lecteur fixé ; le registre complet de leur
génération reste dans l’état enrichi.

\subsection{Le moment de Catalan et la restitution fonctionnelle}
\label{hmtfg:articulacion:catalan}

Les mêmes incidences déterminent une relation entre fonctions,
avant l’évaluation de l’un ou l’autre canal. Pour \(0<s<1\), soit
\(\mathcal D=s\,d/ds\), et définissons le moment orienté
\begin{equation}
 \mathcal C_2(s)
 =\int_0^s\log\frac{s}{t}\,
     \frac{1+t}{1+t+t^2}f(t)\,dt
 =\int_0^s\frac{\log(s/t)}{1+t^2}\,dt.
 \label{hmtfg:articulacion:integral}
\end{equation}
La deuxième égalité découle de la factorisation
\(1+t+t^2+t^3=(1+t)(1+t^2)\). Le poids continu et positif
\(dt/(1+t^2)\) est ainsi produit par la composition rationnelle
du lecteur constitutif, et non par une identification à
\(\nu_0\).

\begin{proposition}[Inversion différentielle, unicité et lecture à l’extrémité]
\label{hmtfg:articulacion:restitucion-catalan}
Le moment \(\mathcal C_2\) est l’unique fonction
\(F\in C^2((0,1))\) qui satisfait
\begin{equation}
 \mathcal D^2F(s)=\frac{s(1+s)}{1+s+s^2}f(s),\qquad
 \lim_{s\downarrow0}F(s)=0.
 \label{hmtfg:articulacion:problema-diferencial}
\end{equation}
La réponse constitutive est restituée par
\begin{equation}
 \mathcal D^2\mathcal C_2(s)=\frac{s}{1+s^2},\qquad
 f(s)=\frac{1+s+s^2}{s(1+s)}\,
          \mathcal D^2\mathcal C_2(s).
 \label{hmtfg:articulacion:inversa-diferencial}
\end{equation}
Son extension continue à l’intervalle fermé possède la série
\begin{equation}
 \mathcal C_2(s)=\sum_{j=0}^{\infty}
       \frac{(-1)^j s^{2j+1}}{(2j+1)^2},\qquad 0\le s\le1,
 \qquad
 \mathcal C_2(1)=\sum_{j=0}^{\infty}\frac{(-1)^j}{(2j+1)^2}
 =G_{\rm Cat}.
 \label{hmtfg:articulacion:serie-catalan}
\end{equation}
\end{proposition}
\begin{proof}
L’intégrande de \eqref{hmtfg:articulacion:integral} est positif ou nul
et
\[
 0\le\mathcal C_2(s)\le\int_0^s\log(s/t)\,dt=s.
\]
Cela prouve l’intégrabilité en zéro et la limite requise.
L’identité \(\log(s/t)=\int_t^s du/u\), suivie de
l’interversion d’intégrales positives, fournit
\[
 \mathcal C_2(s)
 =\int_0^s\frac1u\left(\int_0^u\frac{dt}{1+t^2}\right)du.
\]
Sur tout intervalle compact du domaine ouvert, on applique le
théorème fondamental du calcul intégral. Par conséquent,
\[
 \mathcal D\mathcal C_2(s)=\int_0^s\frac{dt}{1+t^2},
 \qquad
 \mathcal D^2\mathcal C_2(s)=\frac{s}{1+s^2}.
\]
La factorisation rationnelle utilisée dans
\eqref{hmtfg:articulacion:integral} donne
\eqref{hmtfg:articulacion:problema-diferencial} ; isoler \(f\)
donne \eqref{hmtfg:articulacion:inversa-diferencial}, dont les
dénominateurs sont positifs dans \((0,1)\).

Si deux solutions ont la limite nulle requise, leur différence
\(H\) satisfait \(\mathcal D^2H=0\). Le changement \(z=\log s\)
transforme cette équation en \(d^2H(e^z)/dz^2=0\), d’où
\(H(s)=a\log s+b\). L’existence de la limite nulle impose
\(a=b=0\) et prouve l’unicité.

Pour \(s<1\), le développement géométrique de \((1+t^2)^{-1}\)
peut être intégré terme à terme avec le poids logarithmique, car
\[
 \int_0^s t^{2j}\log(s/t)\,dt
 =\frac{s^{2j+1}}{(2j+1)^2},\qquad
 \sum_{j\ge0}\frac{s^{2j+1}}{(2j+1)^2}<\infty.
\]
Cela prouve la série à l’intérieur. La borne
\(1/(2j+1)^2\) donne la convergence absolue et uniforme de cette
série sur \([0,1]\). Dans l’intégrale, l’intégrande prolongé
par zéro pour \(t>s\) est dominé par \(\log(1/t)\)
sur \((0,1)\). La convergence dominée permet d’atteindre
\(s=1\) et coïncide avec le prolongement de la série. Sa valeur
à l’extrémité est précisément la série de Catalan indiquée.
\end{proof}

L’inversion différentielle utilise la famille \(\mathcal C_2(s)\),
et non seulement sa valeur à l’extrémité \(G_{\rm Cat}\). En particulier,
\(\mathcal D^2\mathcal C_2(s)\to1/2\) lorsque \(s\uparrow1\).
La série dérivée deux fois avec \(\mathcal D\) y possède
une limite radiale d’Abel ; la série du moment lui-même est
absolument convergente à la frontière. Cette distinction maintient
les conditions précises de chaque passage à la limite.

\subsection{Composition des restitutions et portée conjointe}
\label{hmtfg:articulacion:conclusion}

Les interrelations précédentes ont des opérations et des domaines
explicites. Les moments de la mesure atomique \(\nu_0\)
produisent \(J_{12}\), et son calcul fonctionnel restitue le profil
\(u_0\). Le profil et le sélecteur de premières racines déterminent
la suite dont le changement d’échelle fournit la lecture \(\delta_{\mathrm F}\).
D’autre part, la fonction constitutive \(f\) produit
\(\mathcal C_2\) au moyen du moment intégral, et
\eqref{hmtfg:articulacion:inversa-diferencial} restitue \(f\)
à partir de cette famille. Ses évaluations étiquetées en
\(s_\pm=q_\pm^{30}\) déterminent les réponses constitutives ;
l’inversion cubique restitue les deux canaux et leurs coordonnées
angulaires. La lecture du moment à l’extrémité donne \(G_{\rm Cat}\).

La mesure \(\nu_0\) est positive, atomique et de rang douze.
Le moment de Catalan intègre un poids positif continu et possède
un développement alterné. Les signes de ce développement sont des
coefficients de la série géométrique ; ce ne sont pas des poids négatifs de
\(\nu_0\), et ils ne convertissent pas l’intégrale positive en la même
représentation spectrale. Chaque lecteur conserve son support,
son orientation, sa normalisation et son opération d’évaluation.

L’unité réside dans la structure APP--TRIT--TPK qui engendre et
transporte leurs antécédents ; les relations fonctionnelles sont
établies par les compositions démontrées. L’égalité
d’origine n’identifie pas des opérateurs de domaines différents.
En particulier, cette articulation ne postule pas de formule scalaire
qui détermine \(\delta_{\mathrm F}\) à partir de la constante de Catalan et
d’autres constantes physiques. Elle situe le changement d’échelle critique aux côtés de
restitutions effectives entre lecteurs du même support,
en conservant l’information requise par chaque direction.

\setcounter{secnumdepth}{\HMTFGSavedSecnumdepth}
\endgroup

\endgroup

\chapter{Identités de Rogers--Ramanujan et relèvement nonadique de Pell}
\begingroup
\renewcommand{\chapter}[2][]{}%
\chapter{Arithmétique du mode \(30\) : identités de Rogers--Ramanujan, normalisateur \(899\) et relèvement nonadique de Pell}
\label{chap:p3-final:46:modo_30_pell}
\section{Normalisateur \(899\), unité de Pell et structure arithmétique du mode \(30\)}

Le normalisateur possède la factorisation arithmétique

\begin{equation}
899=30^2-1=29\cdot31.
\label{mab:rm:eq:899}
\end{equation}

Le même mode \(30\) qui engendre les secteurs \(90\) et \(120\) détermine
simultanément la paire de nombres premiers jumeaux \(29,31\). De plus,

\begin{equation}
\varepsilon_{30}=30+\sqrt{899},\qquad
\varepsilon_{30}^{-1}=30-\sqrt{899},\qquad
\log\varepsilon_{30}=\arcosh30.
\label{mab:rm:eq:pell30}
\end{equation}

Cette unité de Pell détermine l’échelle hyperbolique associée au
normalisateur et relie la phase dodécaphasique, le mode \(30\), les nombres premiers
jumeaux et la géométrie hyperbolique. Cette relation n’identifie pas
l’opérateur constitutif du vide au renormalisateur de doublement ; une telle
identification exigerait un entrelaceur explicite.

L’unité agit de manière entière par
\begin{equation}
M_{30}=\begin{pmatrix}30&899\\1&30\end{pmatrix},
\qquad \det M_{30}=1.
\label{mab:rm:eq:pell-matrix30}
\end{equation}
Sur \(\Z[\sqrt{899}]\), cette matrice représente la multiplication par
\(30+\sqrt{899}\). Ses valeurs propres sont \(\varepsilon_{30}^{\pm1}\) ; par
conséquent, \(\log\varepsilon_{30}\) est la longueur de translation de la transformation
hyperbolique associée. Le même entier central \(30\) réapparaît dans le vide
parce que \(90=3\cdot30\) et \(120=4\cdot30\), mais les deux opérateurs continuent
d’avoir des domaines distincts. La confluence démontrée est celle du mode et de ses
complétions ; l’identité des dynamiques n’est pas postulée.

 \section{Relèvement nonadique de l'unité de Pell}

Pour \(k\ge1\), soit l'anneau quadratique non ramifié
\begin{equation}
\mathcal R_k=(\Z/3^k\Z)[r]/(r^2-2),
\label{eq:pell-ring}
\end{equation}
et soit

\[
u=3+2r.
\]

La conjugaison \(r\mapsto-r\) donne
\[
N_{\mathcal R_k/(\Z/3^k\Z)}(u)
=(3+2r)(3-2r)=9-8=1,
\]
de sorte que \(u\) est une unité de norme un. En outre,
\[
u^2=17+12r,\qquad
u^4-1=576+408r=3(192+136r).
\]
Le second facteur est une unité modulo trois parce qu'il se réduit à \(r\ne0\). Par
conséquent, la valuation \(3\)-adique satisfait
\(v_3(u^4-1)=1\). Le lemme de relèvement pour une unité congrue à un
modulo trois produit, pour tout \(j\ge0\),
\[
v_3\!\left(u^{4\cdot3^j}-1\right)=j+1.
\]
Comme \(u\bmod3=2r\) et \((2r)^2=-1\), son ordre modulo trois est quatre.
En conséquence,

\begin{equation}
\operatorname{ord}(u)=4\cdot3^{k-1},
\label{eq:pell-order}
\end{equation}

et la clôture profinie des groupes cycliques compatibles est
\begin{equation}
\varprojlim_k\langle u\bmod3^k\rangle
\simeq C_4\times\mathbb Z_3.
\label{eq:pell-profinite}
\end{equation}
Sa réalisation archimédienne est

\[
3+2\sqrt2=(1+\sqrt2)^2.
\]

Cette construction démontre la conservation de la profondeur nonadique
par une unité hyperbolique. Les unités \(3+2\sqrt2\),
\(2+\sqrt3\) et \(30+\sqrt{899}\) forment une famille de Pell de
provenances distinctes ; toute identification entre elles exige des
applications explicites entre les extensions correspondantes.

 Le mode \(30\) est l’un des points où l’architecture commence à montrer une unité déterminée par la généalogie opératorielle et arithmétique.\par

Il apparaît d’abord dans le vide parce que\par

\[
90=3\cdot30,
\qquad
120=4\cdot30.
\]

C’est la cadence élémentaire qui permet de comparer les secteurs \(90\) et \(120\). Le vide lit cette cadence comme une complétion \(3\to4\).\par

Mais le même \(30\) réapparaît dans la construction du germe de criticité. Là, l’identité centrale est\par

\[
899=30^2-1=29\cdot31.
\]

La normalisation des douze phases produit \(\sqrt{899}\), et le profil analytique peut s’écrire sous forme fermée. Ce qui est réellement frappant, c’est que \(\pi\) disparaît avant le début de l’itération.\par

Cela doit être lu attentivement. La géométrie angulaire de \(\pi\) est absorbée dans l’organisation des phases. Une fois la configuration dodécaphasique normalisée, le germe dynamique est gouverné par l’arithmétique interne de \(899\).\par

C’est pourquoi la construction commence par engendrer intérieurement un objet dynamique propre :\par

\begin{itemize}[leftmargin=*,itemsep=0.35em]
\item il est pair ;
\item il est analytique ;
\item il a un point critique quadratique ;
\item il appartient à la classe unimodale correspondante ;
\item il possède une dérivée schwarzienne négative ;
\item il peut être soumis à l’opérateur de renormalisation.
\end{itemize}

Les rapports superstables commencent alors à s’approcher des constantes universelles. Le germe a été produit en premier et l’universalité apparaît comme comportement terminal de sa dynamique.\par

C’est une différence épistémologique profonde. Une coïncidence numérique dirait : « notre calcul ressemble à Feigenbaum ». La construction HMT dit : « voici l’objet exact qui entre dans la classe d’universalité, et voici les quotients qui émergent lorsqu’on l’itère ».\par

Le \(30\) réapparaît encore une troisième fois dans la mémoire modulaire. Si\par

\[
N=N_{90}+N_{120}
\]

mesure l’occupation totale et\par

\[
D=90N_{90}+120N_{120}
\]

mesure l’occupation pondérée par la provenance, la matrice qui transforme \((N_{90},N_{120})\) en \((N,D)\) a pour déterminant \(30\).\par

Cela permet d’inverser le processus :\par

\[
N_{90}=4N-\frac D{30},
\qquad
N_{120}=\frac D{30}-3N.
\]

De sorte que le même \(30\) remplit trois fonctions :\par

\begin{itemize}[leftmargin=*,itemsep=0.35em]
\item dans le vide, il rend commensurables les propagations \(90\) et \(120\) ;
\item dans la criticité, il détermine la normalisation \(899=30^2-1\) ;
\item dans la mémoire modulaire, il est le déterminant qui permet de retrouver l’origine de chaque occupation.
\end{itemize}

L’opérateur du vide, le germe de chaos et le hamiltonien modulaire restent typologiquement distincts. La confluence est plus subtile : tous trois conservent une partition commune de la généalogie.\par

Nous pourrions dire que le mode \(30\) est un engrenage. Dans une machine, il mesure la propagation ; dans une autre, il normalise la criticité ; dans une autre, il rend réversible la perte apparente de provenance. L’engrenage est le même, mais chaque machine remplit une fonction différente.\par

L’apparition de l’unité de Pell\par

\[
30+\sqrt{899}
\]

renforce cette interprétation. Le \(30\) organise simultanément une somme finie et une croissance hyperbolique. La même arithmétique qui clôt le germe ouvre une dynamique d’échelle.\par

C’est peut-être la raison profonde pour laquelle le résultat est si beau : le mode qui rend possible la complétion finie du vide est aussi le mode qui prépare une universalité infinie de la criticité.\par

Fini et infini s’articulent par la répétition cohérente d’une structure finie qui conserve la mémoire ; de celle-ci naît l’infini dynamique.\par

  \endgroup
\begin{HMTSpiralChapterBody}
% Corte mecánico íntegro: parte_ii.tex:323--419.
% Residencia material del ensamblaje sucesor de 102 capítulos.
% Residencia causal: capítulo 46.

\chapter{Vis spectrale de Pell et suspension logarithmique}

\section{Loi de vis}

Le corpus contient l'identité exacte
\begin{equation}
 V^{-1}\mathscr S V=(2+\sqrt3)\ii\,\mathscr S.
 \label{espiral:eq:tornillo-pell}
\end{equation}
Soit \(\varepsilon=2+\sqrt3\). L'orbite transverse
\[
 z_n=z_0(\varepsilon\ii)^n
\]
satisfait
\[
 r_n=r_0\varepsilon^n,
 \qquad
 \theta_n=\theta_0+\frac{\pi n}{2}.
\]
En éliminant \(n\),
\begin{equation}
 r(\theta)=r_0
 \exp\!\left[
 \frac{2\log(2+\sqrt3)}{\pi}
 (\theta-\theta_0)
 \right].
 \label{espiral:eq:espiral-log-pell}
\end{equation}
La spirale logarithmique est une sortie de l'opérateur : son équation n'a pas été choisie pour ajuster ensuite \(\varepsilon\).

\begin{proof}[Dérivation de \cref{espiral:eq:espiral-log-pell}]
De \(\theta_n-\theta_0=n\pi/2\) découle \(n=2(\theta_n-\theta_0)/\pi\). En substituant dans \(r_n=r_0\varepsilon^n\),
\[
 r_n=r_0\exp\left(
 \frac{2\log\varepsilon}{\pi}(\theta_n-\theta_0)
 \right).
\]
L'interpolation continue est l'équation indiquée.
\end{proof}

\section{Involution conjuguée de contraction}

L'extension bilatérale \(n\in\Z\) produit
\[
 n\to+\infty:\ r_n\to\infty,
 \qquad
 n\to-\infty:\ r_n\to0.
\]
L'inversion \(n\mapsto-n\) échange expansion et contraction. Comme \(\varepsilon^{-1}=2-\sqrt3\), les deux branches appartiennent à la même unité de Pell. Ce ne sont pas deux spirales indépendantes, mais deux orientations de la même vis spectrale.

\begin{figure}[htbp]
\centering
\begin{tikzpicture}[scale=.82]
 \draw[->,HMTGray] (-5.1,0)--(5.1,0) node[right] {$x$};
 \draw[->,HMTGray] (0,-3.1)--(0,3.1) node[above] {$y$};
 \draw[HMTBlue,very thick,domain=-7:5.2,samples=240,smooth,variable=\t]
   plot ({.27*exp(.22*\t)*cos(deg(\t))},{.27*exp(.22*\t)*sin(deg(\t))});
 \draw[HMTRed,thick,domain=-7:5.2,samples=240,smooth,variable=\t]
   plot ({.27*exp(.22*\t)*cos(deg(-\t))},{.27*exp(.22*\t)*sin(deg(-\t))});
 \foreach \k in {-4,-3,...,4}{
   \pgfmathsetmacro{\ang}{1.57079632679*\k}
   \fill[HMTNavy]
     ({.27*exp(.22*\ang)*cos(deg(\ang))},
      {.27*exp(.22*\ang)*sin(deg(\ang))}) circle (.8pt);
 }
 \node[HMTBlue,font=\small\sffamily] at (3.5,2.4) {expansion};
 \node[HMTRed,font=\small\sffamily] at (3.5,-2.4) {orientation conjuguée};
\end{tikzpicture}
\caption{Projection transverse de la vis de Pell. Les marques représentent des quarts de tour.}
\label{espiral:fig:tornillo-pell}
\end{figure}

\section{Covariance d'échelle et conservation énergétique}

La dilatation \(\varepsilon\) ne doit pas être interprétée comme une croissance gratuite de l'amplitude énergétique d'une onde dans le vide. Sa réalisation cohérente agit sur l'échelle spatiale ou spectrale :
\[
 (r,\varphi,\lambda)
 \longmapsto
 (\varepsilon r,\varphi+\pi/2,\varepsilon\lambda).
\]
Cette distinction sera essentielle dans la réalisation électromagnétique de la \cref{espiral:chap:realizacion-em}.

\begin{resultbox}[title={Compatibilité exacte entre la vis spectrale de Pell et la suspension phase--mémoire}]
Le relèvement phase--mémoire minimal est une hélice circulaire. La vis spectrale combine un quart de tour et une dilatation de Pell ; sa projection est une spirale logarithmique. Tous deux proviennent du registre HMT, mais ne sont pas la même périodicité et ne s'obtiennent pas l'un à partir de l'autre par oubli des types.
\end{resultbox}
 \end{HMTSpiralChapterBody}

\chapter{Théorie arithmétique des lacunes nonadiques et torsion de frontière}
\begingroup
\renewcommand{\chapter}[2][]{}%
\chapter{Théorie arithmétique des lacunes nonadiques : défaut \(10^3-3^6\), réseau radical \(3\to6\to9\), loi puissance--lacune et torsion de bord}
\label{chap:p3-final:47:vacancias_nonadicas}
\section{Classification des cycles arithmétiques irréductibles du réseau radical nonadique}

Dans la réalisation arithmétique, un nombre premier correspond à un cycle
irréductible qui ne se factorise pas en cycles positifs plus petits. La périodicité
nonadique peut sélectionner des candidats et organiser des classes résiduelles ; la
primalité reste définie par la propriété multiplicative exacte

\[
p=ab\Longrightarrow a=1\ \text{ou}\ b=1.
\]

Une fois le coproduit des cycles arithmétiques premiers sélectionné, on
obtient \cref{eq:meissel-mertens,eq:twin-prime-constant}. La structure
hélicoïdale et la racine numérique paramètrent la phase et la retenue, mais ne remplacent pas
le critère de factorisation.

\section{Densité, terme régularisé et holonomie des lacunes}

La densité positive de base est

\begin{equation}
\varepsilon_{\rm vac}=\log_{10}\frac{10}{9}.
\label{eq:vacancy-density}
\end{equation}

La séquence se déduit de la complétion
\(1000=729+271\) détermine
\begin{equation}
\varepsilon_{\rm vac}
=1-\log_{1000}(729)
=\log_{10}\frac{10}{9},
\label{eq:vacancy-density-two-forms}
\end{equation}
et le mot mécanique exact
\begin{equation}
z_n=\lfloor(n+1)\varepsilon_{\rm vac}\rfloor
-\lfloor n\varepsilon_{\rm vac}\rfloor,
\qquad n\ge1.
\label{eq:vacancy-mechanical-word}
\end{equation}
Comme \(\varepsilon_{\rm vac}\) est irrationnel, le mot est sturmien :
il a une complexité \(p(N)=N+1\), un équilibre un et une entropie topologique nulle.
Cette dynamique diffère du décalage complet sur l'alphabet
triadique : la lacune détermine une frontière ordonnée d'entropie nulle.

La somme télescopique produit
\begin{equation}
\sum_{n=1}^{N}z_n
=\lfloor(N+1)\varepsilon_{\rm vac}\rfloor,
\qquad
\sum_{n=1}^{N}(z_n-\varepsilon_{\rm vac})
=\varepsilon_{\rm vac}
-\{(N+1)\varepsilon_{\rm vac}\}.
\label{eq:vacancy-discrepancy}
\end{equation}
Les sommes partielles de la discrépance ont un module inférieur à un. C'est
l'estimation qui soutient les constantes harmoniques ultérieures.

La fraction continue commence
\[
\varepsilon_{\rm vac}
=[0;21,1,5,1,6,2,3,\ldots].
\]
C'est pourquoi les uns de \(z_n\) sont séparés par \(21\) ou \(22\) pas.
Si
\[
\sigma=22-\varepsilon_{\rm vac}^{-1}
=1-T_G(\varepsilon_{\rm vac}),
\]
alors \(a_1(\sigma)=6\) et les retours de la séparation courte sont de
longueur \(6\) ou \(7\). Les couples \(21/22\) et \(6/7\) sont deux échelles d'une
même rotation, et non des coïncidences d'entiers.

Pour le mot exact \cref{eq:vacancy-mechanical-word}, la constante régularisée est

\begin{equation}
\gamma_{\rm vac}^{+}
=\lim_{N\to\infty}\left(
\sum_{n\le N}\frac{z_n}{n}
-\varepsilon_{\rm vac}\log N
\right).
\label{eq:vacancy-gamma}
\end{equation}

L'holonomie associée est

\begin{equation}
\rho_\varepsilon(k)=\exp(2\pi\ii k\varepsilon_{\rm vac}).
\label{eq:vacancy-holonomy}
\end{equation}

Les trois quantités ont des types distincts : densité, terme fini et phase. La tour de Stieltjes des lacunes prolonge \cref{eq:vacancy-gamma} :

\begin{equation}
\gamma_{\rm vac,j}^{+}
=\varepsilon_{\rm vac}\gamma_j
+\sum_{n\ge1}\frac{(z_n-\varepsilon_{\rm vac})(\log n)^j}{n}.
\label{eq:vacancy-stieltjes}
\end{equation}

Pour \(j=0\), on retrouve \(\gamma_{\rm vac}^{+}\). La convergence exige la borne de discrépance du développement intégral des lacunes.

La convergence est prouvée ici. Pour \(j\ge0\), la suite
\(b_n=(\log n)^j/n\) tend vers zéro et est décroissante à partir de
\(n>e^j\). La borne uniforme \cref{eq:vacancy-discrepancy} et le critère de
Dirichlet prouvent la convergence de
\[
\sum_{n\ge1}(z_n-\varepsilon_{\rm vac})b_n.
\]
Pour \(j=0\), séparer
\(\varepsilon_{\rm vac}\sum_{n\le N}n^{-1}\) donne exactement
\cref{eq:vacancy-gamma}. Une sommation d'Abel borne de manière conservative la
queue après \(N\) par \(2/(N+1)\) : un terme provient du bord et
un autre de la variation totale de \(1/n\). L'exécution exhaustive déclarée utilise
\(N=10^8\), énumère les positions de Beatty des uns, vérifie les
ensembles d'écarts \(\{21,22\}\) et de retours \(\{6,7\}\), et obtient comme
partie finie
\[
-0.11207107505876366224207105242526397929\ldots.
\]
En combinant la borne analytique avec les gardes d'arrondi et de frontière
enregistrées, le certificat fini encadre
\begin{equation}
-0.112071094143613634
<\gamma_{\rm vac}^{+}
<-0.112071055516338732.
\label{eq:vacancy-gamma-interval}
\end{equation}
Cet intervalle est une valeur terminale avec erreur contrôlée. Il ne définit pas le
mot ; le mot et sa discrépance le précèdent.

Une seconde évaluation par arithmétique multiprécision et contrôle dirigé
de l'arrondi fournit, à dix-huit décimales, l'intervalle certifié
\[
-0.112071086058763562
<\gamma_{\rm vac}^{+}
<-0.112071064058763762,
\]
strictement contenu dans \cref{eq:vacancy-gamma-interval}. L'emboîtement des
deux intervalles vérifie la stabilité de la partie finie par rapport à la
troncature, aux gardes de frontière et au régime de précision. Le
résultat dépend uniquement des récurrences, des bornes de queue et de
l'arithmétique d'intervalles exposées dans ce chapitre.

\section{Contenu opératoriel de la représentation hélicoïdale}

La représentation hélicoïdale code une phase périodique accompagnée d'un
incrément de mémoire. Elle possède un contenu opératoriel lorsqu'il existe un cocycle

\[
c(\gamma\eta)=c(\gamma)+c(\eta)
\]

qui transporte la retenue et dont la monodromie s'entrelace avec l'opérateur
arithmétique. En l'absence de cette application, l'hélice ne constitue qu'une
paramétrisation de phase ; avec l'entrelaceur, elle distingue le retour de phase
du retour d'état.

\begin{remark}[Contrôle négatif]
La cofinalité des cylindres n'implique pas la conjugaison dynamique. Un
contre-exemple à la commutation entre le décalage TPK et \(T_G\)
empêche l'entrelaceur, mais n'affecte pas
\cref{eq:levy,eq:gauss-entropy}. Dans le secteur des lacunes, la divergence
de \cref{eq:vacancy-stieltjes} sous la borne déclarée réfute la famille de
constantes régularisées.
\end{remark}

\subsection*{Portée logique et domaine de validité}
Cofinalité, identités de Gauss, Lévy, entropie et Lochs : \emph{Théorème interne exact dans le système transporté}. GKW et entrelaceur dynamique : \emph{Programme ouvert avec critère de réfutation}. Densité, constante et holonomie de lacune : \emph{Théorème interne exact avec contrôle de discrépance}.
 \section{Théorème directeur du réseau radical et de la loi puissance--lacune}

L’apparition de \(369\) et de \(693\) ne se réduit pas à une ressemblance graphique. Deux mécanismes exacts convergent :

\begin{enumerate}
\item \textbf{mécanisme algébrique nonadique :} \(3,6,9\) sont les représentants visuels du radical nilpotent \(\mathfrak m=(3)\subset\mathbb Z/9\mathbb Z\) ; les mots \(369\), \(693\) et \(936\) sont les trois phases de son orbite multiplicative ;
\item \textbf{mécanisme archimédien des lacunes :} la différence entre la longueur décimale de \(9^n\) et la frontière \(10^n\) est gouvernée exactement par la même rotation irrationnelle \(\delta_{\rm vac}=\log_{10}(10/9)\) qui organise déjà le calendrier des lacunes.
\end{enumerate}

Pour \(n=9^9\), les deux mécanismes se rencontrent dans l’identité

\[
9^9-\left\lceil 9^9\log_{10}\frac{10}{9}\right\rceil
=369\,693\,099,
\]

de sorte que

\[
\#\operatorname{cifras}\bigl(9^{9^9}\bigr)=369\,693\,100.
\]

Le bloc \(369\mid693\) est donc une sortie exacte du lecteur logarithmique de la puissance nonadique ; ce n’est pas le préfixe décimal de \(9^{9^9}\). L’unité finale provient de l’opérateur universel qui transforme l’ordre décimal en longueur décimale.

Cette identité ouvre une lecture structurelle de plus grande portée : somme, produit, puissance, soustraction, division, racine et logarithme peuvent être organisés comme une hiérarchie de transports et de quotients. Chaque montée d’opération compactifie des histoires ; le TPK conserve les coordonnées qui permettent de les reconstruire.


\section{Structure radicale de l’APP sur \(\mathbb Z/9\mathbb Z\)}

\subsection{Anneau local nonadique \(\mathbb Z/9\mathbb Z\) et radical nilpotent}

Soit

\[
R=\mathbb Z/9\mathbb Z.
\]

Son groupe des unités et son idéal maximal sont

\[
R^\times=\{1,2,4,5,7,8\},
\qquad
\mathfrak m=(3)=\{0,3,6\}=\{9,3,6\}_{\rm vis}.
\]

De plus,

\[
\mathfrak m^2=0.
\]

L’égalité précédente est une égalité d’idéaux : tout produit de deux éléments de \(\mathfrak m\) s’annule modulo neuf. Elle ne doit pas être confondue avec le produit cartésien \(\mathfrak m\times\mathfrak m\), qui contient neuf positions de l’APP.

\subsection{Suite exacte du multiplicateur \(3\)}

La multiplication par trois,

\[
M_3:R\longrightarrow R,
\qquad x\longmapsto3x,
\]

a

\[
\ker M_3=\mathfrak m,
\qquad
\operatorname{im}M_3=\mathfrak m.
\]

Elle induit donc un isomorphisme d’espaces sur \(\mathbb F_3\),

\[
\overline M_3:R/\mathfrak m\xrightarrow{\sim}\mathfrak m,
\qquad
\overline x\longmapsto3x.
\]

En représentants visuels,

\[
\overline0\longmapsto9,
\qquad
\overline1\longmapsto3,
\qquad
\overline2\longmapsto6.
\]

Ainsi, la bande \(3,6,9\) n’est pas un ornement décimal. C’est une copie interne du quotient ternaire logée dans le radical nonadique.

\subsection{Orbite cyclique \(369\to693\to936\to369\)}

La ligne de la table multiplicative correspondant à \(3\) est

\[
3,6,9,3,6,9,3,6,9.
\]

Le changement de phase \(j\mapsto j+1\) produit

\[
369\longmapsto693\longmapsto936\longmapsto369.
\]

Ces trois mots sont les trois lectures cycliques de l’application

\[
\mathbb F_3\longrightarrow\mathfrak m,
\qquad
\bar j\longmapsto3j.
\]

La ligne \(6\), en satisfaisant \(6\equiv-3\pmod9\), inverse l’orientation du même cycle. La marque \(9\) remplit deux fonctions typées :

\[
9\bmod9=0,
\qquad
\operatorname{dr}_9(9)=9.
\]

Le module publie la classe nulle ; la racine numérique positive conserve la marque de clôture. C’est précisément cette différence de lecteurs qui permet à \(9\) d’être à la fois zéro résiduel et clôture visible.

\subsection{Croix radicale de l’APP et réduction coordonnée à \(\mathbb F_3^2\)}

Sur

\[
X_9=R^2,
\]

l’union

\[
\mathscr C_{369}
=(\mathfrak m\times R)\cup(R\times\mathfrak m)
\]

est la croix formée par les lignes et les colonnes \(3,6,9\). Elle contient

\[
3\cdot9+3\cdot9-3\cdot3=45
\]

positions. Son intersection centrale \(\mathfrak m\times\mathfrak m\) contient neuf positions. Dans l’observable multiplicative, la somme de la croix est

\[
297=216+81=3(72+27).
\]

La réduction coordonnée produit

\[
X_9/(\mathfrak m\times\mathfrak m)\cong\mathbb F_3^2
\]

comme ensemble quotient typé par les classes modulo trois. Le relèvement inverse exige de conserver la coordonnée à l’intérieur de chaque fibre de trois éléments.

\subsection{Noyau \(7227\), décomposition spectrale et lien avec la croix radicale}

Le bloc central de l’APP est

\[
B_c=
\begin{pmatrix}
7&2\\
2&7
\end{pmatrix}
=7I+2R
=9P_++5P_-.
\]

Sa lecture par lignes produit \(72\mid27\) ; sa lecture par diagonales produit \(77\mid22\). On vérifie

\[
72+27=77+22=99,
\]

\[
72-27=45=\det B_c,
\]

\[
77-22=55=54+1.
\]

La relation avec la croix radicale n’est pas nominale :

\[
3\cdot72=216,
\qquad
3\cdot27=81,
\qquad
3\cdot99=297.
\]

Le bloc unitaire central et le réseau non unitaire \(3,6,9\) sont ainsi reliés par une identité locale–globale exacte. La hiérarchie de la différence somme–produit conserve ses types :

\[
4\quad\text{(saut spectral local)},\qquad
2\quad\text{(couplage hors diagonale)},
\]

\[
6\quad\text{(différence comprimée modulo trois)},\qquad
54\quad\text{(déploiement bidimensionnel)}.
\]


\section{Loi puissance--lacune pour les échelles \texorpdfstring{\(9^n\)}{9 à la puissance n} et \texorpdfstring{\(10^n\)}{10 à la puissance n}}

\subsection{Fonctions de longueur décimale et de lacune cumulée}

Pour \(n\ge1\), soit

\[
D_9(n)=\#\operatorname{cifras}(9^n)
=\left\lfloor n\log_{10}9\right\rfloor+1.
\]

La frontière décimale \(10^n\) possède \(n+1\) chiffres. Nous définissons la lacune cumulée de la puissance nonadique par

\[
V_9(n)=(n+1)-D_9(n).
\]

Nous introduisons les deux nombres complémentaires

\[
\rho_{\rm vac}=\log_{10}9,
\qquad
\delta_{\rm vac}=1-\rho_{\rm vac}
=\log_{10}\frac{10}{9}.
\]

\subsection{Identité exacte entre longueur décimale et rotation des lacunes}

\textbf{Théorème.} Pour tout entier \(n\ge1\),

\[
\boxed{V_9(n)=\left\lceil n\delta_{\rm vac}\right\rceil},
\]

et, par conséquent,

\[
\boxed{D_9(n)=n-\left\lceil n\delta_{\rm vac}\right\rceil+1}.
\]

\textbf{Démonstration.} Comme \(\delta_{\rm vac}\) est irrationnel, \(n\delta_{\rm vac}\notin\mathbb Z\). Alors

\[
\begin{aligned}
D_9(n)
&=\lfloor n(1-\delta_{\rm vac})\rfloor+1\\
&=\lfloor n-n\delta_{\rm vac}\rfloor+1\\
&=n-\lceil n\delta_{\rm vac}\rceil+1.
\end{aligned}
\]

En soustrayant de \(n+1\), on obtient la première identité. \(\square\)

\subsection{Encodage mécanique des lacunes par les puissances nonadiques}

Les accroissements de la lacune sont

\[
Z_n=V_9(n+1)-V_9(n)
=\lceil(n+1)\delta_{\rm vac}\rceil
-\lceil n\delta_{\rm vac}\rceil
\in\{0,1\}.
\]

C’est exactement le mot mécanique de pente \(\delta_{\rm vac}\) utilisé par le calendrier HMT des lacunes. Ce calendrier admet donc une lecture supplémentaire entièrement exacte :

\begin{quote}
\(Z_n=1\) si et seulement si, en passant de \(9^n\) à \(9^{n+1}\), on perd un nouveau rang décimal par rapport à la frontière \(10^n\to10^{n+1}\).
\end{quote}

La puissance n’ajoute pas une coïncidence à la théorie des lacunes : elle fournit sa réalisation naturelle comme compteur de perte de rang décimal.

\subsection{Loi \(21/22\) comme cas du théorème puissance--lacune}

Puisque

\[
\delta_{\rm vac}^{-1}
=21.85434532678283256\ldots,
\]

les indices pour lesquels \(Z_n=1\) sont séparés par \(21\) ou \(22\). Les vingt et un premiers exposants satisfont

\[
D_9(n)=n,\qquad1\le n\le21,
\]

tandis que

\[
D_9(22)=21.
\]

En particulier,

\[
9^9=387\,420\,489
\]

possède exactement neuf chiffres. Cette conservation initiale du rang n’est pas une règle indéfinie : le calendrier des lacunes détermine exactement quand elle cesse d’être vérifiée.

\subsection{Évaluation exacte en \(n=9^9\) et longueur de \(9^{9^9}\)}

Soit

\[
n_*=9^9=387\,420\,489.
\]

Alors

\[
n_*\delta_{\rm vac}
=17\,727\,389.3684296412564569\ldots,
\]

et, par conséquent,

\[
\boxed{V_9(n_*)=17\,727\,390}.
\]

La coordonnée conservée après déduction des lacunes est

\[
n_*-V_9(n_*)
=387\,420\,489-17\,727\,390
=369\,693\,099.
\]

Par conséquent,

\[
\boxed{\operatorname{ord}_{10}\bigl(9^{9^9}\bigr)=369\,693\,099},
\]

\[
\boxed{D_9(9^9)=369\,693\,100}.
\]

On obtient en outre les congruences

\[
9^9\equiv0\pmod9,\qquad
V_9(9^9)\equiv0\pmod9,
\]

\[
369\,693\,099\equiv0\pmod9,\qquad
369\,693\,100\equiv1\pmod9.
\]

Les sommes de chiffres correspondantes sont

\[
s_{10}(9^9)=45,\qquad
s_{10}(V_9(9^9))=36,
\]

\[
s_{10}(369\,693\,099)=54,\qquad
s_{10}(369\,693\,100)=37.
\]

Le retour \(9\to1\) apparaît ici avec des types précis :

\begin{enumerate}
\item \(9^9\) et la lacune cumulée appartiennent à la classe nulle modulo neuf ;
\item leur différence, l’ordre décimal, reste dans la classe nulle et possède une somme de chiffres \(54\) ;
\item l’opérateur universel \(D=\operatorname{ord}+1\) ajoute l’unité de clôture et publie la classe positive \(1\).
\end{enumerate}

\subsection{Domaine de validité de la loi puissance--lacune}

Il est démontré :

\begin{itemize}
\item le théorème puissance–lacune pour tout \(n\) ;
\item l’identité avec le mot mécanique des lacunes ;
\item l’évaluation exacte en \(n=9^9\) ;
\item l’ordre et la longueur décimale de \(9^{9^9}\) ;
\item les congruences et sommes de chiffres précédentes.
\end{itemize}

On n’affirme pas que \(9\) soit l’unique base ayant une propriété analogue dans tout domaine, que chaque apparition externe de \(369\) provienne de cette tour, ni que la tour remplace la construction des neuf phases ou le certificat de génération de \(\pi,e,\varphi\).


\section{Hiérarchie opératorielle, réversibilité et mémoire généalogique}

\subsection{Translation, dilatation et exponentiation comme niveaux opératoriels}

Les opérations élémentaires peuvent être vues comme des actions :

\[
T_a(x)=x+a,\qquad
M_a(x)=ax,\qquad
P_k(x)=x^k.
\]

Leurs inverses partielles sont, respectivement, la soustraction, la division et l’extraction de racines. Le passage d’une opération à la suivante modifie la structure des fibres.

\subsection{Réversibilité stratifiée dans \(\mathbb Z/9\mathbb Z\)}

\textbf{Proposition.} Dans l’anneau nonadique :

\begin{enumerate}
\item toute translation \(T_a\) est bijective ;
\item \(M_a\) est bijective si et seulement si \(a\in R^\times\) ;
\item si \(a\in\mathfrak m\), la multiplication possède des fibres non triviales et son image est contenue dans \(\mathfrak m\) ;
\item pour \(x\in\mathfrak m\), \(x^k=0\) pour tout \(k\ge2\) ;
\item sur \(R^\times\cong C_6\), l’application \(P_k\) est un automorphisme si et seulement si \(\gcd(k,6)=1\).
\end{enumerate}

\textbf{Démonstration.} La première affirmation utilise la structure de groupe additif. La deuxième est la caractérisation des unités. La troisième découle de \(aR\subseteq\mathfrak m\). La quatrième procède de \(\mathfrak m^2=0\). La cinquième est la condition pour que la multiplication par \(k\) soit un automorphisme du groupe cyclique d’ordre six. \(\square\)

Le résultat produit une hiérarchie exacte :

\[
\text{somme : réversibilité globale},
\]

\[
\text{produit : réversibilité restreinte aux unités},
\]

\[
\text{puissance : dynamique cyclique sur les unités et effondrement nilpotent dans }\mathfrak m.
\]

Les racines ne sont des fonctions que lorsque la puissance correspondante est injective sur le domaine choisi ; sur le domaine total, ce sont des correspondances à plusieurs branches ou des relations vides.

\subsection{Lecture du Topological Phase Kernel sur la hiérarchie opératorielle}

Cette stratification suggère une définition fonctionnelle :

\begin{quote}
Le TPK n’ajoute pas une opération extérieure à la somme, au produit et à la puissance. Il conserve le feuillet, le quotient, la retenue, l’orientation, la route et la frontière que chaque évaluation ordinaire élimine lorsqu’elle identifie plusieurs histoires à une même sortie.
\end{quote}

La somme et le produit sont les deux premières évaluations sur l’APP. La puissance est une composition répétée du produit ; la racine sélectionne ou reconstruit des branches de cette composition. Le registre distingue la valeur publiée de l’histoire opératorielle qui la produit.

\subsection{Descente logarithmique entre niveaux consécutifs de la hiérarchie}

Dans un domaine positif,

\[
\log(xy)=\log x+\log y,\qquad
\log(x^k)=k\log x,\qquad
\log(a^x)=x\log a.
\]

Le logarithme transforme le produit en somme, la puissance en produit et l’exponentielle en dilatation. Le défaut apparaît lorsque le lecteur ultérieur discrétise :

\[
\mathcal L_{10}(a^n)
=\lfloor n\log_{10}a\rfloor+1.
\]

La partie entière introduit le calendrier des franchissements de frontière. Pour \(a=9\), ce calendrier est précisément le mot des lacunes.

\subsection{Conjugaison affine des translations et des dilatations}

Les translations et les dilatations satisfont

\[
M_aT_bM_a^{-1}=T_{ab}
\]

lorsque \(a\) est inversible. Cette identité exprime que le produit réorganise les unités de somme. Lorsque \(a\) cesse d’être inversible, la conjugaison se rompt parce que des branches disparaissent. Le radical \(3,6,9\) localise exactement cet échec.


\section{Algèbre des chemins : composition multiplicative, action additive et superposition}

Dans une algèbre des chemins, la somme agrège les alternatives et le produit concatène les trajectoires. Pour des poids locaux \(w(e)\),

\[
K(x,y)
=\sum_{\gamma:x\to y}
\prod_{e\in\gamma}w(e).
\]

Si l’on fixe une branche logarithmique cohérente,

\[
\prod_{e\in\gamma}w(e)
=\exp\!\left(\sum_{e\in\gamma}\log w(e)\right),
\]

et, avec l’action discrète appropriée,

\[
K(x,y)=\sum_{\gamma:x\to y}e^{-S(\gamma)/\hbar}.
\]

L’équation réunit les opérations :

\begin{itemize}
\item le produit compose les pas successifs ;
\item le logarithme accumule l’action locale ;
\item l’exponentielle reconstruit le poids de la route ;
\item la somme agrège toutes les alternatives.
\end{itemize}

L’APP est le support minimal sur lequel somme et produit agissent sur les mêmes histoires. Le TPK ajoute les données qui empêchent la somme sur les chemins de devenir une somme sur des valeurs sans généalogie.

Dans ce cadre, un chiffre est la sortie d’un lecteur ; une valeur est la projection archimédienne d’une section ; un nombre HMT conserve la famille compatible de chemins qui soutient cette valeur ; une particule conserve cette généalogie sous une réalisation spectrale persistante.


\section{Racine interne de \(-1\), structure exponentielle et géométries elliptique, parabolique et hyperbolique}

\subsection{Racine interne de \(-1\) dans la branche finie}

La matrice d’incidence de Paley fournit un opérateur \(A_W\) avec

\[
A_W^2=-I.
\]

Cela fait de \(\mathbb F_3^6\) un espace tridimensionnel sur

\[
\mathbb F_9=\mathbb F_3[\iota]/(\iota^2+1).
\]

La racine de \(-1\) n’est pas adjointe de l’extérieur comme une notation complexe. Elle émerge comme endomorphisme d’un support ternaire déjà construit.

\subsection{Réalisation continue du type quadratique induit}

Dans la branche réelle elliptique, le générateur \(J_{\rm e}\) satisfait

\[
J_{\rm e}^2=-I,
\]

et

\[
e^{tJ_{\rm e}}=\cos t\,I+\sin t\,J_{\rm e}.
\]

En particulier,

\[
e^{\frac\pi2J_{\rm e}}=J_{\rm e},\qquad
e^{\pi J_{\rm e}}=-I,\qquad
e^{2\pi J_{\rm e}}=I.
\]

Ainsi, \(J_{\rm e}\) est le quart de tour interne et l’équation d’Euler est sa clôture au demi-tour.

\subsection{Support opératoriel commun de \(3\) lectures géométriques}

Soit

\[
\mathbb D=\{z\in\mathbb C:|z|<1\}.
\]

Ce même ensemble contient le diamètre réel \((-1,1)\), l’intérieur complexe, la frontière trigonométrique \(S^1\) et la métrique hyperbolique de Poincaré

\[
ds_{\rm P}^2=\frac{4\ell^2|dz|^2}{(1-|z|^2)^2}.
\]

L'application

\[
R_{\pi/2}(z)=iz
\]

préserve \(|z|\) et \(|dz|\). Par conséquent, c’est simultanément :

\begin{itemize}
\item une multiplication complexe par \(\sqrt{-1}\) ;
\item une rotation trigonométrique de \(90^\circ\) ;
\item une isométrie euclidienne du disque ;
\item une isométrie hyperbolique de Poincaré ;
\item un automorphisme de sa frontière circulaire.
\end{itemize}

La coïncidence est littérale : la même application agit sur le même porteur ; c’est le lecteur métrique qui change.

\subsection{Géométries elliptique, parabolique et hyperbolique de l’équation d’Euler étendue}

Les générateurs

\[
J_{+1}^2=-I,\qquad
J_0^2=0,\qquad
J_{-1}^2=I
\]

produisent

\[
e^{tJ_{+1}}=\cos t\,I+\sin t\,J_{+1},
\]

\[
e^{tJ_0}=I+tJ_0,
\]

\[
e^{tJ_{-1}}=\cosh t\,I+\sinh t\,J_{-1}.
\]

La même série exponentielle se sépare en trois régimes selon le carré du générateur : elliptique, parabolique et hyperbolique. La spécialisation

\[
e^{2\log\varphi\,J_{-1}}=F_{\rm av}^{\,2}
\]

intègre l’auto-échelle de Fibonacci. Euler et Fibonacci apparaissent comme des spécialisations d’une famille exponentielle quadratique commune.

\subsection{Porteur commun, opérateur commun et famille exponentielle quadratique}

La fusion géométrique se formule à trois niveaux :

\begin{enumerate}
\item \textbf{porteur commun :} le disque unité contient le diamètre réel, l’intérieur complexe et le bord circulaire ;
\item \textbf{opérateur commun :} \(z\mapsto iz\) est un quart de tour, une multiplication complexe et une isométrie hyperbolique ;
\item \textbf{famille commune :} l’équation exponentielle tritique réunit les branches elliptique, parabolique et hyperbolique.
\end{enumerate}

Le résultat n’identifie pas les métriques euclidienne et hyperbolique. Il affirme que toutes deux sont des réalisations distinctes sur un même support et partagent l’automorphisme issu de la racine interne de \(-I\).


\section{Restriction de la rotation intérieure, holonomie de bord et torsion minimale}

La construction contient deux objets positifs et différenciés :

\begin{enumerate}
\item le générateur interne \(J_{\rm e}\), avec \(J_{\rm e}^2=-I\) ;
\item l’invariant global
\end{enumerate}
   \[
   \Delta_4=
   \frac{\pi-e\log\pi}{270}
   \left(1+\frac{\pi}{729}\right),
   \]
   appelé torsion minimale et utilisé avant la sélection d’espèce.

La thèse de l’auteur ajoute que la dynamique rotationnelle interne, lorsqu’elle est publiée au bord, est absorbée comme torsion minimale. La relation ne doit pas s’écrire comme une simple égalité \(J_{\rm e}=\Delta_4\) : les objets ont des types différents. Sa forme naturelle exige

\[
\text{rotation intérieure}
\xrightarrow{\operatorname{Res}_{\partial}}
\text{holonomie de bord}
\xrightarrow{\operatorname{Def}}
\text{défaut scalaire}.
\]

Soit \(q_\partial\) le lecteur de frontière, \(\nabla_\partial\) la connexion induite et \(e_\partial\) le repère discret. La torsion de bord a pour type

\[
T_\partial=d_{\nabla_\partial}e_\partial.
\]

L’affirmation de l’auteur acquiert une forme probatoire au moyen d’une fonctionnelle

\[
\Lambda_\partial:
\operatorname{Tors}(\partial X)
\longrightarrow\mathbb R
\]

tel que

\[
\Lambda_\partial(T_\partial)=\Delta_4
\]

et d’un théorème qui fasse procéder \(T_\partial\) de la restriction de \(J_{\rm e}\) par les opérateurs de bord du TPK.

La racine interne de \(-I\), le quart de tour, l’équation d’Euler étendue et la formule de \(\Delta_4\) sont démontrés dans leurs domaines. Le passage de la rotation intérieure à la torsion minimale de bord exige le morphisme \(\operatorname{Res}_{\partial}\) et la fonctionnelle \(\Lambda_{\partial}\) déclarés ; il n’est pas identifié par une égalité scalaire.


\section{Triplicité générationnelle, conjugaison et mélange sur l’état holonomique intégral}

La phrase d’Isidor Isaac Rabi peut servir de référence historique, mais elle ne doit pas régir le titre ni constituer une section autonome. La question scientifique est :

\begin{quote}
Quel mécanisme interne distingue trois ordres stables de réalisation fermionique et conserve cette distinction sous couleur, conjugaison et mélange ?
\end{quote}

La réponse HMT n’est pas « parce que TRIT a trois valeurs ». La structure sépare

\[
\operatorname{flav}(X)
=\bigl(\sigma_{ud}(X),g(X),\chi_{\rm fina}(X)\bigr),
\]

où \(\sigma_{ud}\) distingue les branches, \(g\) distingue la génération physique et \(\chi_{\rm fina}\) conserve l’orientation et la mémoire au sein d’une branche. Dans les quarks, l’orbite de couleur est ajoutée après la saveur et avant le caractère.

La triplicité est obtenue comme conséquence de la structure enrichie ; ses effets sur la couleur, la conjugaison et le mélange sont développés ci-dessous :

\begin{enumerate}
\item couleur, saveur, génération et masse sont des lecteurs distincts du même état holonomique intégral ;
\item la conjugaison particule–antiparticule agit sur l’orientation sans effacer la classe générationnelle ;
\item CKM et PMNS transportent des bases d’états déjà typés ; ils ne créent pas la triplicité ;
\item la classification neutrinique doit découler de la représentation, de l’occupation et du mélange, non d’un changement de nom des profondeurs \(G0,\ldots,G4\).
\end{enumerate}

Le fait que les neutrinos soient des leptons de spin demi-entier et d’obéissance fermionique est une classification physique connue. L’apport structurel HMT consiste à obtenir intérieurement pourquoi le secteur neutre occupe cette représentation, pourquoi il admet trois réalisations physiques et comment il se couple à l’opérateur de mélange sans lire ses angles externes.


\section{Mémoire opératorielle du TPK : conservation du résidu, du quotient, de la retenue, du feuillet, de l’orientation, de la route et de la frontière sous projection}

\subsection{Perte de réversibilité et relèvement généalogique}

Dans l’APP, différentes histoires peuvent produire le même résidu. Le produit replie davantage d’histoires que la somme ; la puissance en replie davantage que le produit ; la racine tente de reconstruire des branches déjà identifiées. Le TPK conserve la donnée qui permet d’inverser cette perte :

\[
\text{résidu},\quad
\text{quotient},\quad
\text{retenue},\quad
\text{feuillet},\quad
\text{orientation},\quad
\text{route},\quad
\text{frontière},\quad
\text{mémoire}.
\]

L’opération algébrique peut être caractérisée non seulement par sa valeur, mais par le type de généalogie qu’elle conserve ou détruit.

\subsection{Projection des histoires compatibles sur le continu archimédien}

La droite réelle apparaît lorsqu’une famille de cylindres compatibles se contracte sous l’évaluation archimédienne. La valeur réelle est l’ombre de la section ; non sa généalogie complète. La projection exceptionnelle conserve une autre partie comme incidence et symétrie.

La somme sur les chemins exprime le même principe :

\[
\text{alternatives}\xrightarrow{\sum}\text{amplitude globale},
\qquad
\text{étapes successives}\xrightarrow{\prod}\text{poids de route}.
\]

Le logarithme transforme la composition multiplicative en action additive ; l’exponentielle la ramène à une amplitude. La structure discrète du continu peut être formulée comme conservation généalogique de toutes les routes avant leur projection sur \(\mathbb R\).

\subsection{Distinction entre section numérique et réalisation persistante de particule}

Un nombre et une particule partagent des invariants ontologiques minimaux : mémoire, orientation, persistance, clôture, signature et position dans une famille de prolongements.

Le nombre publie une coordonnée archimédienne. La particule ajoute une réalisation spectrale, une occupation, une propagation et une stabilité. Le passage HMT \(\to\) MD change le lecteur appliqué à une même généalogie enrichie ; il ne change pas arbitrairement d’objet.


\section{Clôture conjointe du réseau radical, de la loi puissance--lacune, de la géométrie d’Euler et de la hiérarchie opératorielle}

\subsection{Théorème du réseau radical nonadique}

Le réseau \(3,6,9\) de l’APP est la représentation visuelle du radical nilpotent \((3)\subset\mathbb Z/9\mathbb Z\). Ses phases \(369\), \(693\) et \(936\) réalisent l’isomorphisme \(R/(3)\cong(3)\) induit par la multiplication par trois. Les bandes localisent exactement le défaut d’inversibilité du produit.

\subsection{Théorème puissance--lacune}

Pour tout \(n\ge1\), le nombre de rangs décimaux perdus par \(9^n\) par rapport à \(10^n\) est \(\lceil n\log_{10}(10/9)\rceil\). Ses accroissements forment le mot mécanique des lacunes. En \(n=9^9\), la coordonnée retenue est \(369\,693\,099\) et la longueur résultante est \(369\,693\,100\).

\subsection{Théorème géométrique de la famille exponentielle d’Euler}

La racine interne de \(-I\) fournit un quart de tour. Dans le disque unité, la même application \(z\mapsto iz\) est une multiplication complexe, une rotation trigonométrique et une isométrie de Poincaré. L’équation exponentielle tritique prolonge ce régime aux régimes parabolique et hyperbolique, en incorporant dans le dernier l’auto-échelle de Fibonacci.

\subsection{Principe de conservation opératorielle entre niveaux}

La hiérarchie somme–produit–puissance est une hiérarchie de fibres et de perte de réversibilité. Le TPK est le relèvement qui conserve les coordonnées nécessaires pour distinguer les histoires identifiées par les lecteurs ordinaires.


\section{Hypothèses de typage et contrôles négatifs}

\begin{enumerate}
\item Ne pas identifier \(369\,693\,100\) à un mot TPK sans déclarer le lecteur entre les domaines.
\item Ne pas confondre l’idéal \(\mathfrak m^2=0\) avec \(\mathfrak m\times\mathfrak m\).
\item Ne pas confondre modulo neuf et racine numérique : \(9\mapsto0\) dans le premier et \(9\mapsto9\) dans la seconde.
\item Ne pas utiliser la triplicité de TRIT comme démonstration suffisante de trois générations physiques.
\item Ne pas transformer \(J^2=-I\) en \(\Delta_4\) sans l’application de bord et la fonctionnelle de torsion.
\item Ne pas présenter \(369\) comme un préfixe de \(9^{9^9}\) : il appartient à son ordre et à sa longueur décimale.
\item Ne pas réduire les dix mille chiffres à une preuve de grandeur : ils certifient des troncatures étendues d’un processus paramétrable dont la thèse est le prolongement indéfini.
\end{enumerate}

\section{Conséquences conjointes pour le réseau radical, la mémoire opératorielle et la torsion de bord}

La découverte principale n’est pas qu’une tour contienne visuellement les chiffres \(3,6,9\). La puissance nonadique et le calendrier des lacunes sont deux présentations du même défaut logarithmique, tandis que l’APP localise \(3,6,9\) comme le radical où la multiplication perd son inversibilité. La tour relie les deux extrémités :

\[
\text{radical nonadique}
\longrightarrow
\text{puissance}
\longrightarrow
\text{logarithme}
\longrightarrow
\text{vacance}
\longrightarrow
\text{ordre décimal}
\longrightarrow
\text{clôture }9\to1.
\]

L’équation d’Euler étendue ajoute la seconde direction : la racine interne de \(-I\) transforme la structure discrète en dynamique de rotation et réunit, sur un support commun, la droite réelle, le cercle trigonométrique et le disque hyperbolique. Le passage ultérieur à la torsion minimale doit conserver l’application de bord déjà développée dans le corpus, non la remplacer par une identification nominale.
\endgroup

\chapter{Radion, phase angulaire et équivalence entre les réalisations cyclique et linéaire}
\begingroup
\renewcommand{\chapter}[2][]{}%
% Integración científica del capítulo del radión. La numeración y los rótulos
% se adaptan a la monografía; los cuerpos matemáticos permanecen íntegros.
% Adaptador de compatibilidad del fuente teoremática íntegro del radión.
%
% El tratado rector ya carga tipografía, geometría, hipervínculos, teoremas,
% tablas, TikZ y tcolorbox.  Este archivo aporta únicamente los nombres y
% entornos exclusivos del fuente teoremática, sin importar su preámbulo autónomo ni
% alterar la jerarquía tipográfica general.

\makeatletter
\@ifundefined{color@RadNavy}{\definecolor{RadNavy}{HTML}{102A43}}{}
\@ifundefined{color@RadBlue}{\definecolor{RadBlue}{HTML}{1F5E7A}}{}
\@ifundefined{color@RadTeal}{\definecolor{RadTeal}{HTML}{168C8C}}{}
\@ifundefined{color@RadCopper}{\definecolor{RadCopper}{HTML}{B66A3C}}{}
\@ifundefined{color@RadGold}{\definecolor{RadGold}{HTML}{D4A73A}}{}
\@ifundefined{color@RadInk}{\definecolor{RadInk}{HTML}{1A202C}}{}
\@ifundefined{color@RadMist}{\definecolor{RadMist}{HTML}{EDF4F7}}{}
\@ifundefined{color@RadCream}{\definecolor{RadCream}{HTML}{FAF7F0}}{}
\@ifundefined{color@RadRose}{\definecolor{RadRose}{HTML}{8D3F55}}{}
\@ifundefined{color@RadGreen}{\definecolor{RadGreen}{HTML}{2F7D5D}}{}
\@ifundefined{color@RadGray}{\definecolor{RadGray}{HTML}{66788A}}{}

\providecommand{\TRIT}{\ensuremath{\mathrm{TRIT}}}
\providecommand{\HMT}{\ensuremath{\mathrm{HMT}}}
\providecommand{\Rstate}{\mathcal R_{12}}
\providecommand{\epsR}{\varepsilon_R}
\providecommand{\epsone}{\varepsilon_R^{(1)}}
\providecommand{\pih}{\pi}
\providecommand{\alphah}{\alpha}
\providecommand{\hret}{\hbar_{\mathrm{ret}}}
\providecommand{\hfull}{h_{\mathrm{ret}}}
\providecommand{\diff}{\mathop{}\!\mathrm d}
\providecommand{\e}{\mathrm e}
\providecommand{\R}{\mathbb R}
\providecommand{\C}{\mathbb C}
\providecommand{\F}{\mathbb F}
\providecommand{\Z}{\mathbb Z}
\providecommand{\dd}{\,\mathrm d}
\providecommand{\status}[1]{\textsf{\bfseries #1}}
\providecommand{\sourcepath}[1]{\nolinkurl{#1}}

\@ifundefined{typebox}{%
  \newtcolorbox{typebox}[1][]{enhanced,breakable,colback=white,
    colframe=RadBlue,boxrule=.7pt,arc=1.5mm,left=4mm,right=4mm,
    top=3mm,bottom=3mm,title={Reçu de types},fonttitle=\bfseries,
    coltitle=RadNavy,#1}%
}{}
\@ifundefined{narrativebox}{%
  \newtcolorbox{narrativebox}[1][]{enhanced,breakable,colback=white,
    colframe=RadGold,borderline west={2.4pt}{0pt}{RadGold},
    boxrule=.25pt,arc=0mm,left=5mm,right=4mm,top=3mm,bottom=3mm,#1}%
}{}

% Las once portadillas internas del fuente teoremática autónomo se convierten en
% divisiones científicas ordinarias dentro del capítulo 52.  Se evita así
% introducir partes autónomas, gigantismo tipográfico o páginas vacías.
\providecommand{\partopener}[3]{\section{#2}\noindent\emph{#3}\par\medskip}

\providecommand{\BigRadionEquation}{%
  \begin{equation*}
  \boxed{\displaystyle
  \frac{180^{\circ}}{\pih}
  =50^{\circ}+A^{\circ}-\frac{20}{81}\,\Delta_4^{\circ}+\epsR^{\circ}}
  \end{equation*}}
\makeatother
 
\chapter{Le radion angulaire : raccordement avec retenue des sections directe et torsionnelle et fermeture exacte de l’unité de phase, d’action et de mémoire}
\label{chap:sucesor-102-radion}

\section{Composition APP--TRIT--TPK de l'état de phase : génération de l'unité orientée avec mémoire}
% Fragmento público generado desde fuentes teoremáticas íntegros.
% Residencia: 52.1.
% Fuente: /Users/ruben/Documents/HMT2/EL_CIERRE_HOLOGRAFICO_DEL_INFINITO_SUCESOR_102_CAPITULOS_2026-08-28_EN_TRABAJO/incoming/radion_monografia_20260827/manuscrito/sections/01_resumen_y_guia.tex:1-128
\paragraph{Synthèse mathématique de l'objet et de sa chaîne causale}

Cette monographie reconstruit le radian à partir de la chaîne générative
\[
\APP\longrightarrow\TRIT\longrightarrow\TPK
\longrightarrow X_{\mathrm{enr}}
\longrightarrow\text{structure discrète du continu}
\longrightarrow\MD,
\]
sans prendre comme entrée ni la valeur métrologique de \(180/\pi\), ni le petit
résidu angulaire que l'on souhaite expliquer. Le résultat est une unité enrichie,
le \emph{radion}, défini comme l'unité orientée qui parcourt un radian de
phase et balaie exactement une action \(\hret\) dans l'ellipse positive associée
au même état.

La clôture directrice est
\BigRadionEquation
avec
\[
S_E=2\pih\left(\frac5{36}+\frac{25}{9}\alphah\right),
\qquad
\Delta_4=
\frac{\pih-e\log\pih}{270}
\left(1+\frac{\pih}{729}\right),
\]
\[
\rho_{\mathrm{tw}}=\frac{2\cdot90}{729}=\frac{20}{81},
\qquad
\epsone=\frac{20}{81}\Delta_4-(S_E-1).
\]
La division APP prouve \(1<S_E<2\), de sorte que le quotient est exactement
un et le reste est \(D_E=S_E-1\). L'égalité
\[
1=S_E-\frac{20}{81}\Delta_4+\epsone
\]
est alors une identité interne. La carte angulaire ultérieure produit
\[
\epsR^\circ
=5.13830167585752820716851646226459474899085744\ldots
\times10^{-8}\,{}^\circ.
\]

Le texte démontre en outre que l'ellipse d'action
\[
Q=a_S\cos\theta,\qquad P=b_S\sin\theta,\qquad
a_Sb_S=2\hret
\]
satisfait
\[
\operatorname{Area}(0,\theta)=\hret\theta.
\]
Par conséquent, un radion balaie \(\hret\) et le tour \(2\pi\) enferme
\(h_{\mathrm{ret}}=2\pi\hret\). La même ellipse possède un bord symplectique
fini et un intérieur d'aire infinie lorsqu'elle est munie de la métrique
de Hilbert--Beltrami--Klein. Telle est la forme exacte de la thèse physique qui guide
l'ouvrage : la surface visible clôt l'action ; la profondeur intérieure
conserve un développement non épuisable par la projection plane.

L'exposition maintient typées quatre réalisations du même état :

\begin{enumerate}[label=\textbf{\arabic*.}]
\item le cercle \(S^1\), qui publie la phase ;
\item l'ellipse positive, qui publie l'action et l'anisotropie ;
\item l'intérieur de Klein, qui mesure la profondeur au moyen du birapport ;
\item l'hélice nonadique, qui conserve retour, mémoire et contraction.
\end{enumerate}

Il ne s’agit pas de quatre géométries superposées sans critère. Ce sont quatre lecteurs de domaines, codomaines et données conservées différents, appliqués à un unique état holonomique intégral. Le secteur \(90/120\), la dodécaphase, la retenue du radion, la queue de Lambert, le noyau de Barbero et les tours globales sont également présentés comme des objets liés mais non interchangeables.

\begin{thesisbox}
Le radion est l'unité orientée qui balaie \(\hret\) dans l'ellipse d'action.
Le cercle publie sa phase, Klein mesure sa profondeur, l'hélice conserve sa
mémoire et la retenue coinductive raccorde exactement ses sections directe et
torsionnelle.
\end{thesisbox}

\paragraph{Ordre d'exposition et dépendances de lecture}

L'ordre n'est pas ornemental. Les Parties I et II produisent l'objet et sa clôture
à partir d'APP--TRIT--TPK. Ce n'est qu'ensuite que sont introduites les réalisations dans le langage
de Hilbert, LC, Maxwell, Minkowski, Hodge, Lorentz, de la modularité et des orbites. Chaque
chapitre distingue trois niveaux :

\begin{description}[style=nextline,leftmargin=4.2cm,labelwidth=3.9cm]
\item[Résultat interne.] Égalité ou construction démontrée dans le domaine
HMT avec ses primitives déclarées.
\item[Formalisation réunie.] Nouvelle application articulant des résultats déjà existants
sans utiliser la cible comme sélecteur.
\item[Interface physique.] Traduction dimensionnelle ultérieure reconnaissant une
structure conventionnelle sans en faire l'origine de la sélection.
\end{description}

Les encadrés de \emph{statut et falsificateur} indiquent quelle observation concrète
invaliderait chaque promotion forte. Cette discipline n'affaiblit pas la thèse ; elle la
rend auditable. L'objectif de la narration étendue est qu'aucun mot
comme \emph{cercle}, \emph{ellipse}, \emph{torsion}, \emph{charge} ou
\emph{mémoire} ne masque un changement de type.

\paragraph{Résultats directeurs et portée probatoire}

\begin{enumerate}[label=\textbf{C\arabic*.}]
\item Clôture indépendante de la valeur finale du radion par deux publications du même état et
soustraction avec retenue APP.
\item Dérivation du coefficient \(20/81\) comme deux feuillets du secteur actif
\(90\), normalisés par la fibre \(3^6=729\).
\item Égalisateur typé entre la couture \(\Re s=1/2\), la face APP de cent
marques et la phase dodécaphasique \(\theta_2=50^\circ\).
\item Théorème d'aire du radion : \(\operatorname{Area}(\theta)=\hret\theta\).
\item Étude focale exacte de l'ellipse et de ses invariants euclidiens,
hyperboliques, symplectiques et modulaires.
\item Théorème bord--intérieur : action totale finie \(h\) et intérieur de Klein
d'aire hyperbolique infinie.
\item Réalisation exacte par covariance minimale et oscillateur LC, suivie de
la publication électromagnétique \(90/120\).
\item Atlas de non-identifications pour tous les objets \(90/120\), y compris
la séparation entre \(\epsR\), queue spectrale et noyau de Barbero \(12:-1\).
\item Caractérisation structurelle de \(\pi\) dans la carte du radion, distincte
du générateur coinductif primaire.
\item Synthèse ontologique : dimension comme rang minimal de mémoires qu'une
projection de phase ne peut plus conserver.
\end{enumerate}
% Fuente: /Users/ruben/Documents/HMT2/EL_CIERRE_HOLOGRAFICO_DEL_INFINITO_SUCESOR_102_CAPITULOS_2026-08-28_EN_TRABAJO/incoming/radion_monografia_20260827/manuscrito/sections/01b_mapa_resultados.tex:1-133
\paragraph{Correspondance entre résultats, types et résidences causales}

La correspondance suivante ordonne les affirmations directrices et leurs démonstrations,
en maintenant visible la hiérarchie qui empêche que
la richesse physique efface la clôture mathématique qui la soutient.

% HMT_RESULT_BEGIN:RDN-01-GENEALOGIA:theorem
\begin{exactbox}[title={Généalogie APP--TRIT--TPK du radion}]
Le radion naît de la composition effective
\(\APP\to\TRIT\to\TPK\to X_{\mathrm{enr}}\to\MD\) ; il ne s'obtient pas
en redéfinissant rétrospectivement le radian conventionnel.
\end{exactbox}
% HMT_RESULT_END:RDN-01-GENEALOGIA:theorem

% HMT_RESULT_BEGIN:RDN-02-CIERRE:theorem
\begin{exactbox}[title={Clôture exacte et indépendance prospective}]
La retenue APP produit intérieurement
\[
1=S_E-\frac{20}{81}\Delta_4+\varepsilon_R^{(1)},
\qquad
\varepsilon_R^{(1)}=\frac{20}{81}\Delta_4-(S_E-1),
\]
et la carte angulaire transporte cette même identité vers
\(180^\circ/\pi\) sans introduire le résidu cible.
\end{exactbox}
% HMT_RESULT_END:RDN-02-CIERRE:theorem

% HMT_RESULT_BEGIN:RDN-03-BISAGRA-50:theorem
\begin{exactbox}[title={Coordonnée critique de \(50^\circ\)}]
La coordonnée critique \(\Re s=1/2\), publiée sur la face APP de cent
marques et sur la roue dodécaphasique, sélectionne la phase
\(\theta_2=50^\circ\). C'est une égalité de lecteurs typés du même état.
\end{exactbox}
% HMT_RESULT_END:RDN-03-BISAGRA-50:theorem

% HMT_RESULT_BEGIN:RDN-04-COEFICIENTE-2081:theorem
\begin{exactbox}[title={Densité bidirectionnelle nonadique \(20/81\)}]
Le coefficient torsionnel n'est pas ajusté sur le réseau \(1/81\) : il est construit comme
\(\rho_{\mathrm{tw}}=2N_6/3^6=2\cdot90/729=20/81\), en conservant feuillet,
orientation et normalisation de la fibre nonadique.
\end{exactbox}
% HMT_RESULT_END:RDN-04-COEFICIENTE-2081:theorem

% HMT_RESULT_BEGIN:RDN-05-AREA-ACCION:theorem
\begin{exactbox}[title={Loi surfacique du radion}]
Pour l'ellipse positive
\(Q=a_S\cos\theta\), \(P=b_S\sin\theta\),
\(a_Sb_S=2\hbar_{\mathrm{ret}}\), l'aire orientée balayée est
\(\mathcal A(\theta)=\hbar_{\mathrm{ret}}\theta\). Un radion balaie exactement
\(\hbar_{\mathrm{ret}}\), et le tour \(2\pi\) enferme
\(h_{\mathrm{ret}}=2\pi\hbar_{\mathrm{ret}}\).
\end{exactbox}
% HMT_RESULT_END:RDN-05-AREA-ACCION:theorem

% HMT_RESULT_BEGIN:RDN-06-INVARIANTES-FOCALES:theorem
\begin{exactbox}[title={Invariants focaux de l'opérateur positif}]
Le couple publié \(A\pm C^\ast\) fixe les demi-axes, l'excentricité, la
séparation focale et le paramètre de compression symplectique. En particulier,
\((d_1/d_0)^2=C^\ast/2\), identité adimensionnelle qui sépare la forme de l'échelle.
\end{exactbox}
% HMT_RESULT_END:RDN-06-INVARIANTES-FOCALES:theorem

% HMT_RESULT_BEGIN:RDN-07-BORDE-INTERIOR:theorem
\begin{exactbox}[title={Finitude symplectique du bord et divergence de l'aire hyperbolique intérieure}]
Le bord symplectique de l'ellipse enferme l'action finie
\(h_{\mathrm{ret}}\), tandis que l'aire de Hilbert--Beltrami--Klein de l'intérieur
diverge lorsqu'elle atteint la frontière. Phase visible et profondeur ne sont pas la même
mesure.
\end{exactbox}
% HMT_RESULT_END:RDN-07-BORDE-INTERIOR:theorem

% HMT_RESULT_BEGIN:RDN-08-HELICE-MEMORIA:theorem
\begin{exactbox}[title={Relèvement hélicoïdal nonadique avec retour de phase et avance de mémoire}]
La monodromie nonadique ramène la phase sans réinitialiser l'état. L'hélice
relève le cercle en conservant retenue et mémoire ; la clôture \(2\pi/4\pi\)
distingue retour projectif, inversion d'orientation et retour complet.
\end{exactbox}
% HMT_RESULT_END:RDN-08-HELICE-MEMORIA:theorem

% HMT_RESULT_BEGIN:RDN-09-OSCILADOR-EM:development
\begin{narrativebox}[title={Réalisation oscillatoire et constitutive électromagnétique}]
La même anisotropie se réalise par une covariance minimale et un oscillateur
LC : une coordonnée stocke l'aspect électrique-élastique et sa conjuguée
l'aspect magnétique-rotationnel. La fréquence et l'impédance sont typées
séparément.
\end{narrativebox}
% HMT_RESULT_END:RDN-09-OSCILADOR-EM:development

% HMT_RESULT_BEGIN:RDN-10-TIPOS-90120:theorem
\begin{exactbox}[title={Distinction typée des opérateurs \(90/120\)}]
Extracteur dodécaphasique, retenue du radion, queue de Lambert, fonctionnelle hexade--octade, noyau de Barbero et semi-groupe global des tours sont des objets distincts. Le poids \(12:-1\) n’est imposé qu’une fois fixés le pôle \(1/24\) et le contre-terme entier minimal.
\end{exactbox}
% HMT_RESULT_END:RDN-10-TIPOS-90120:theorem

% HMT_RESULT_BEGIN:RDN-11-GEOMETRIAS-CAMPO:development
\begin{narrativebox}[title={Réalisations complexe, lorentzienne et duale de Hodge}]
Les régimes elliptique, parabolique et hyperbolique du TRIT sont publiés comme
rotation, seuil et transformation hyperbolique. L'étoile de Hodge, les lois constitutives et la force
de Lorentz agissent ensuite sur l'état déjà orienté ; elles ne sélectionnent pas ses
constantes.
\end{narrativebox}
% HMT_RESULT_END:RDN-11-GEOMETRIAS-CAMPO:development

% HMT_RESULT_BEGIN:RDN-12-MODULAR-ORBITAL:development
\begin{narrativebox}[title={Caractères modulaire, orbital et arithmétique de frontière}]
L'involution modulaire, les orbites elliptiques, l'avance-retard et l'arithmétique
3--6--9 reconnaissent différentes mémoires du même état. La fonction modulaire et
les vacances sont utilisées comme lecteurs ultérieurs, jamais comme substituts du
générateur APP--TRIT--TPK.
\end{narrativebox}
% HMT_RESULT_END:RDN-12-MODULAR-ORBITAL:development

% HMT_RESULT_BEGIN:RDN-13-SINTESIS-ONTOLOGICA:conclusion
\begin{thesisbox}[title={Synthèse opératoire de phase, d'action, de profondeur et de mémoire}]
Le cercle est la projection de phase ; l'ellipse est la section d'action ; Klein
mesure la profondeur ; l'hélice conserve l'histoire. Le radion est l'unité
orientée qui maintient ces publications couplées sans effacer le résidu qui
les distingue.
\end{thesisbox}
% HMT_RESULT_END:RDN-13-SINTESIS-ONTOLOGICA:conclusion

% HMT_RESULT_BEGIN:RDN-14-CONTROL-ALCANCE:control
\begin{statusbox}[title={Délimitation des domaines et continuité causale}]
Les lois complètes des masses, le raffinement électronique, CKM/CP et le
couloir exceptionnel conservent leurs propres démonstrations directrices. Cette monographie les
cite seulement là où ils partagent un antécédent ou une carte ; elle ne transfère pas au radion leurs
sélecteurs spécifiques et ne confond pas leurs résidences probatoires.
\end{statusbox}
% HMT_RESULT_END:RDN-14-CONTROL-ALCANCE:control
 
\subsection{Support bivarié d’APP, orientation ternaire de TRIT et actualisation de l’état holonomique intégral par TPK}
% Fragmento público generado desde fuentes teoremáticas íntegros.
% Residencia: 52.1.1.
% Fuente: /Users/ruben/Documents/HMT2/EL_CIERRE_HOLOGRAFICO_DEL_INFINITO_SUCESOR_102_CAPITULOS_2026-08-28_EN_TRABAJO/incoming/radion_monografia_20260827/manuscrito/sections/02_genealogia.tex:1-180
\noindent La généalogie opératoire précède la mesure angulaire et détermine les variables conservées par chaque lecteur.\par\medskip

\subsubsection{Support APP--TRIT--TPK et construction de l’état holonomique intégral}

\paragraph{Antériorité de l’état holonomique intégral par rapport à l’expression circulaire}

La définition scolaire du radian commence par un cercle déjà donné, un centre déjà
donné et une longueur d'arc déjà mesurable. Si le rayon et l'arc ont la
même longueur, l'angle sous-tendu mesure un radian. La définition est correcte
au sein de la géométrie euclidienne, mais ne répond pas à une question antérieure :
quelle structure permet à un tour de posséder phase, orientation, action et
mémoire, et pourquoi l'unité naturelle est-elle liée à \(2\pi\) ?

HMT inverse l'ordre causal. Elle construit d'abord l'état fini, puis sa
prolongation coinductive, ensuite le tour et seulement à la fin sa publication
angulaire. Le radion ne corrige pas le radian conventionnel ; il le relève vers son domaine
généalogique.

\begin{typebox}
\[
\APP\to\TRIT\to\TPK\to X_{\mathrm{enr}}
\to\mathcal C^{\mathrm{disc}}_{\mathrm{cont}}
\to\bigl(\pih,e,\varphi,\alphah,A,C^*,\Delta_4,\hret\bigr)
\to\mathfrak r_\sigma.
\]
Les mathématiques conventionnelles n'apparaissent qu'ensuite comme langage de réalisation :
\(S^1\), ellipse symplectique, métrique de Klein, espace de Hilbert, circuit
LC, écran de Minkowski et équations de champ.
\end{typebox}

\paragraph{Support bivarié d'APP : feuillets, résidus, quotients et retenue}

La plateforme d'arithmétique positionnelle (APP) conserve simultanément deux
réalisations des symboles \(1,\dots,9\) : un feuillet additif et un feuillet
multiplicatif. Un nombre n'est pas réduit à sa valeur publiée. Avec la valeur
subsistent feuillet, orientation, quotient, reste, retenue, frontière, chemin et
survie.

Le recensement de base comporte trois strates qui ne doivent jamais être identifiées par simple
coïncidence de cardinalité :
\[
104,976\longrightarrow468\;U_6\longrightarrow243\;w_6
\subset\F_3^6,\qquad |\F_3^6|=729.
\]
Les \(468\) émissions décimales, les \(243\) mots visibles et l'espace ambiant
de \(729\) mots sont des objets distincts. L'espace ambiant sera décisif pour la
normalisation \(20/81\).

En base cubique \(B=1000\), la soustraction orientée de deux mots de triades
s'effectue de droite à gauche :
\begin{equation}
u_i=t_i-v_i+c_{i+1},\qquad
r_i=u_i\bmod1000,\qquad
c_i=\frac{u_i-r_i}{1000}.
\label{eq:subacarreo}
\end{equation}
Les opérandes et la retenue distante étant fixés, quotient et reste sont uniques. Tel est
l'opérateur qui, plus loin, publiera les chiffres de \(\epsone\). Il n'y a
aucune décimale choisie après avoir vu le radian.

La complétion
\[
10^3=729+243+27+1
\]
n'est pas une curiosité décimale. Elle exprime la conservation de l'unité lors du passage
de l'espace ambiant ternaire à une chambre de mille positions. La différence
\(1000-729=271\) est la couronne complétante ; l'identité correcte est
\(999=729+270\), non \(729+271\).

\paragraph{Orientation ternaire et classification des 3 régimes opératoires du TRIT}

Le TRIT attribue à chaque état local l'un de trois régimes orientés. Dans la
convention temporelle active,
\[
\begin{aligned}
+1&\longleftrightarrow\text{retour, mémoire, passé},\\
0&\longleftrightarrow\text{seuil, présent},\\
-1&\longleftrightarrow\text{ouverture bidirectionnelle, futur}.
\end{aligned}
\]
La réalisation quadratique distingue :
\begin{equation}
J^2=-I,\qquad N^2=0,\qquad R^2=I.
\label{eq:tres-regimenes}
\end{equation}
Ainsi apparaissent, respectivement, la rotation elliptique, le seuil parabolique et
la séparation hyperbolique. On ne mélange pas trois métriques. On conserve trois
actions typées sur un antécédent commun.

La racine interne de \(-1\) n'est pas importée en supposant d'avance le plan
complexe. Dans le bloc APP, deux projections \(P,Q\) de rapport
\(r=3/4=90/120\) produisent
\[
J_{\mathrm{comm}}=\frac{[P,Q]}{\sqrt{r(1-r)}},\qquad J_{\mathrm{comm}}^2=-I,
\]
tandis que
\[
R_{\mathrm{comm}}=2P-I,\qquad R_{\mathrm{comm}}^2=I,\qquad
R_{\mathrm{comm}}J_{\mathrm{comm}}=-J_{\mathrm{comm}}R_{\mathrm{comm}}.
\]
Le couple génère la forme minimale de \(\mathrm{Cl}_{1,1}\simeq M_2(\R)\) :
rotation et ouverture émergent ensemble, mais restent des opérations différentes.

\paragraph{Composition dynamique du TPK : sélection, transport, mémoire et retour}

Le Topological Phase Kernel n'est pas le nom d'une boîte d'objets. C'est la
composition qui sélectionne la phase, transporte la cardinalité, relève les feuillets,
enregistre la mémoire et exécute le retour. La dépendance fondamentale est
\[
U_t\longrightarrow\Gamma_9\longrightarrow K_{\mathrm{ph}}.
\]
L’holonomie nonadique \(\Gamma_9\) agit sur l’état holonomique intégral ; l’expression \(K_{\mathrm{ph}}\) est ultérieure et conserve une mémoire réduite. L’identité \(K_{\mathrm{ph}}^9|_r=E_+E_\times\) ne transforme pas \(K_{\mathrm{ph}}\) en générateur.

Le certificat nonadique conserve la chaîne complète
\[
w_6\to w_{12}\to w_{18}\to w_{24}\to w_{30}
\to R_{36}\to G_9,
\]
avec \(396\) niveaux, \(2376\) trits, \(43\) tours, \(387\) paires
\((K,K+9)\) par canal et cinq régions de \(\pi\). La phase satisfait
\[
g(k+9)=g(k),
\]
mais l'état complet ne revient pas :
\[
B_{k+9}\ne B_k.
\]
Tel est le germe mathématique de l'hélice : la projection angulaire se ferme ; la
mémoire avance.

\paragraph{Extension temporelle non scindée \(C_{12}\to C_{108}\to C_9\)}

La mémoire de douze pas s'organise au moyen de
\begin{equation}
0\longrightarrow C_{12}\longrightarrow C_{108}
\longrightarrow C_9\longrightarrow0.
\label{eq:extension-108}
\end{equation}
Le cocycle
\[
c(r,r')=\left[\left\lfloor\frac{r+r'}9\right\rfloor\right]_{12}
\]
mesure la retenue entre retour de phase et avance de mémoire. Neuf pas ferment
la phase observable \(C_9\) ; ils ne réinitialisent pas le registre dodécaphasique.

Il convient de séparer quatre objets :
\begin{center}
\begin{tabularx}{.95\textwidth}{>{\bfseries}l X}
\toprule
\(C_{12}\) & groupe cyclique de douze positions ;\\
\(\Rstate\) & registre enrichi \((E_m,\tau_m,\sigma_m,\kappa_m,\Lambda_m)_{m=1}^{12}\) ;\\
\(\Theta_{108}\) & publication brute dans le calendrier de 108 ;\\
\(\Gamma_{108}\) & histoire ou holonomie relevée.\\
\bottomrule
\end{tabularx}
\end{center}
C'est pourquoi l'expression vague \enquote{post-\(C_{12}\)} sera remplacée par
\enquote{postérieur au registre dodécaphasique \(\Rstate\), dans sa lecture
sexadique} lorsque telle sera l'application effective.

\paragraph{Émergence corrélative des 5 constructions du continu à partir de l'état unique}

La structure discrète du continu ne s'obtient pas en assemblant cinq théories
indépendantes. Ses cinq aspects internes ---spectral, solénoïdal, de jauge,
cohomologique et déterminantal--- émergent corrélativement du même état
APP--TRIT--TPK. Cet ouvrage ne démontre pas de nouveau le traité intégral, mais
préserve le fait causal dont le radion a besoin : la section coinductive qui
génère \(\pi,e,\varphi\) et \(\alpha\) ne vit pas sur une droite réelle nue ; elle vit
dans un objet doté de fibres, de relations, d'orientation, de retenue, de mémoire, de terminal et de
lecteur.

\begin{statusbox}
\textbf{Falsificateur de généalogie.} Si une formule ultérieure reçoit
\(180/\pi\), \(\epsR\), une masse ou une table métrologique pour choisir une phase,
un coefficient ou une branche, elle cesse d'être une génération HMT-first. La confrontation
externe peut reconnaître une sortie ; elle ne peut pas la sélectionner.
\end{statusbox}
 
\subsection{Indépendance prospective de la clôture par rapport à la valeur finale}
% Fragmento público generado desde fuentes teoremáticas íntegros.
% Residencia: 52.1.2.
% Fuente: /Users/ruben/Documents/HMT2/EL_CIERRE_HOLOGRAFICO_DEL_INFINITO_SUCESOR_102_CAPITULOS_2026-08-28_EN_TRABAJO/incoming/radion_monografia_20260827/manuscrito/sections/02_genealogia.tex:181-272

\subsubsection{Généalogie des constantes et construction du couple angulaire}

\paragraph{Section coinductive et unicité de la valeur générée de \(\pi\)}

Soit \(x_\infty\) la section globale du système projectif. Les lecteurs
coinductifs produisent
\[
\operatorname{ev}_\pi(x_\infty)=\pih,\qquad
\operatorname{ev}_e(x_\infty)=e,\qquad
\operatorname{ev}_\varphi(x_\infty)=\varphi.
\]
Ces évaluations déterminent les constantes universelles elles-mêmes, non des variantes dépendant du procédé qui les construit. En particulier,
\[
\pih=3.14159265358979323846264338327950288419716939937510\ldots.
\]
Le tour positif est
\begin{equation}
T_+(x_\infty)=2\pih.
\label{eq:vuelta-plus}
\end{equation}
Sa publication angulaire d'une unité est
\begin{equation}
R_{\HMT}^{\circ}=\frac{360^\circ}{T_+}=\frac{180^\circ}{\pih}.
\label{eq:radian-publicado}
\end{equation}
L'équation ne définit pas encore le radion enrichi ; elle définit la carte visible
du radian une fois le tour généré.

\paragraph{Bidirectionnalité des chemins et prolongations conjuguées}

La bidirectionnalité ne s'énonce pas comme un slogan. L'involution concrète
\begin{equation}
J_\alpha(T)=\frac{\alphah^2}{T},\qquad
T_-=J_\alpha(T_+),\qquad T_+T_-=\alphah^2
\label{eq:bidireccional-turn}
\end{equation}
conserve un produit et inverse le feuillet. Expansion et contraction sont deux
prolongations conjuguées du même état. La racine finie d'un jet
quadratique peut approcher \(T_+\), mais ne remplace pas la section
coinductive : la différence certifiée est
\[
T_{P_9}-T_+=-5.1110566290335\ldots\times10^{-25}.
\]
Le contrôle fini est un témoin ; le fondement est la limite.

\paragraph{Dépendance causale entre \(\alpha\), action et publication angulaire}

La clôture dodécaphasique génère \(\alphah\) à partir de l'état signé et de ses
cartes réversibles. Ensuite, le lecteur déterminantal local fixe la décennie
d'action au moyen de
\[
\Sigma_H=\varphi\alphah^{16},
\qquad \operatorname{ord}_{10}\Sigma_H=-34.
\]
La section de retour produit \(\hret\). Ni \(\alpha\) ni \(\hbar\) ne se
\enquote{mesurent en degrés}. Ils sont publiés ultérieurement au moyen du couple angulaire
\begin{equation}
P(\alpha,\eta_{\mathrm{ret}})=
\left([10^3\alpha]_{360},[2(\eta_{\mathrm{ret}}+\alpha)]_{360}\right)
=(A,C^*).
\label{eq:parangular}
\end{equation}
Dans la branche canonique,
\[
A=7.297352569283800997285105472380663\ldots{}^\circ,
\]
\[
C^*=2.1237383389686457242619513803933607937\ldots{}^\circ.
\]
Alors seulement, la carte analytique
\[
x=\frac{\pi A}{180},\qquad y=\frac{\pi C^*}{180}
\]
ouvre les canaux \(q_\pm=e^{-x\mp y}\).

L’ordre causal complet est
\[
\alpha\to(\eta\parallel-34)\to\hbar
\to P\to\kappa_{360}\to q_\pm
\to\text{réalisations MD}.
\]
La relation \(C^*/2-\alpha\) peut servir de vérification inverse. Elle n'est pas utilisée
pour générer \(\hret\) rétroactivement.

\begin{narrativebox}
La carte angulaire ne transforme pas une constante adimensionnelle ou une action en un
angle par décret. Elle publie deux coordonnées de l'état sur une roue finie.
Le degré est le langage de cette publication ; il n'est pas la substance ontologique de
\(\alpha\) ni de \(\hbar\).
\end{narrativebox}
% Fuente: /Users/ruben/Documents/HMT2/EL_CIERRE_HOLOGRAFICO_DEL_INFINITO_SUCESOR_102_CAPITULOS_2026-08-28_EN_TRABAJO/incoming/radion_monografia_20260827/manuscrito/sections/03_cierre_exacto.tex:254-285
\paragraph{Indépendance prospective par rapport à la valeur finale}

Le vérificateur de l'opérateur de profondeur interdit expressément trois entrées :
\[
\frac{180}{\pi},\qquad \epsR,\qquad
\text{le décimal cible}.
\]
Il construit d'abord \(S_E\) et \(S_T\), puis exécute la soustraction APP, et ne
compare la clôture qu'à la fin. Les contrôles donnent
\[
\texttt{target\_epsilon\_used=False},\qquad
\texttt{target\_180\_over\_pi\_used=False}.
\]
De même :
\begin{itemize}
\item un seul feuillet donne un défaut d'ordre \(10^{-5}\) ;
\item remplacer \(\theta_2\) par \(\theta_1\) ou \(\theta_3\) déplace le
résultat d'environ \(\pm\pi/6\) ;
\item éliminer le canal \(e\) laisse un défaut d'ordre \(10^{-3}\) ;
\item la variante historique \(\Delta_4(A,C^*)\) produit un autre nombre ;
\item la queue spectrale postérieure à \(\Rstate\) et l'opérateur harmonique à
quatre termes ne coïncident pas non plus avec \(\epsone\).
\end{itemize}

\begin{statusbox}
\textbf{Statut.} La clôture est une \status{formalisation nouvelle} démontrée
avec une structure de départ explicite et un certificat informatique indépendant de la valeur finale.
L'égalité est exacte à la limite coinductive. Le profil court de 36 chiffres
de \(\alpha\) génère une variante qui ne diverge qu'après de nombreux chiffres ;
il est conservé comme témoin historique, non comme valeur canonique.
\end{statusbox}

 
\section{Sections directe et torsionnelle : retenue APP et clôture exacte du radion}
\subsection{Construction indépendante des sections \(S_E\) et \(S_T\) avant la définition du cocycle \(\varepsilon_R=S_T-S_E\)}
% Fragmento público generado desde fuentes teoremáticas íntegros.
% Residencia: 52.2.1.
% Fuente: /Users/ruben/Documents/HMT2/EL_CIERRE_HOLOGRAFICO_DEL_INFINITO_SUCESOR_102_CAPITULOS_2026-08-28_EN_TRABAJO/incoming/radion_monografia_20260827/manuscrito/sections/03_cierre_exacto.tex:1-50
\noindent La clôture exacte s'obtient au moyen de deux sections construites avant le cocycle qui les relie.\par\medskip

\subsubsection{Construction indépendante des sections directe et torsionnelle}

\paragraph{Section directe du système temporel dodécaphasique}

La roue dodécaphasique a pour centres
\begin{equation}
\theta_m=30m-10^\circ,\qquad m=1,\ldots,12.
\label{eq:rueda-dodecafasica}
\end{equation}
La deuxième phase est \(\theta_2=50^\circ\). Avec la publication
\(A=10^3\alphah\), elle forme la section directe
\begin{align}
S_E
&=\frac{T_+}{360}\,(50+10^3\alphah)\\
&=2\pih\left(\frac5{36}+\frac{25}{9}\alphah\right).
\label{eq:seccion-electrica}
\end{align}
L'indice \(E\) signifie ici \emph{directe ou commune} ; sa réalisation
électrique viendra après le lecteur constitutif. On ne postule pas
littéralement \(A=E\).

Les cylindres coinductifs prouvent
\begin{equation}
1<S_E<2.
\label{eq:SE-intervalo}
\end{equation}
La division APP impose alors
\begin{equation}
q_E=\lfloor S_E\rfloor=1,
\qquad
D_E=\operatorname{rem}_1(S_E)=S_E-1.
\label{eq:DE}
\end{equation}
L'entier un n'est pas le nombre de tours de l'hélice. C'est le quotient unique
de la section directe au sein de la cellule d'unité.

\paragraph{Section torsionnelle déterminée par l'opérateur \(\Delta_4\)}

La seconde section utilise une sortie coinductive antérieure :
\begin{equation}
\Delta_4=
\frac{\pih-e\log\pih}{270}
\left(1+\frac{\pih}{729}\right).
\label{eq:delta4}
\end{equation}
Le mot \emph{torsion} désigne ici le quantum global de la carte HMT. Il ne
l'identifie pas automatiquement à la torsion de Cartan
\(T^I=D_\omega e^I\). Le pont vers Cartan exige une autre flèche typée.
 \subsection{Densité orientée du secteur actif et détermination unique de \(20/81\)}
% Fragmento público generado desde fuentes teoremáticas íntegros.
% Residencia: 52.2.2.
% Fuente: /Users/ruben/Documents/HMT2/EL_CIERRE_HOLOGRAFICO_DEL_INFINITO_SUCESOR_102_CAPITULOS_2026-08-28_EN_TRABAJO/incoming/radion_monografia_20260827/manuscrito/sections/03_cierre_exacto.tex:51-103

\paragraph{Densité orientée \(20/81\) du secteur actif \(90/120\)}

Le sélecteur tritique de phase sépare
\[
H_+=\{1,4,7\},\qquad H_-=\{2,5,8\},\qquad H_0=\{3,6,9\}.
\]
Il y a six phases actives et quinze paires non ordonnées :
\[
H_{\mathrm{act}}=H_+\cup H_-,\qquad
\left|\binom{H_{\mathrm{act}}}2\right|=\binom62=15.
\]
Le mode actif et son extension à huit positions sont
\begin{equation}
N_6=6\binom62=90,\qquad
N_8=8\binom62=120.
\label{eq:90-120-incidencia}
\end{equation}
Deux feuillets conjugués du secteur actif, normalisés par la fibre ambiante,
donnent
\begin{equation}
\rho_{\mathrm{tw}}
=\frac{2N_6}{|\F_3^6|}
=\frac{2\cdot90}{729}
=\frac{180}{729}
=\boxed{\frac{20}{81}}.
\label{eq:rho-twoway}
\end{equation}

\begin{proposition}[Unicité relative de la densité]
Le secteur actif \(N_6=90\), les deux feuillets conjugués et la
normalisation par \(|\F_3^6|=729\) étant fixés, la densité orientée est nécessairement
\(20/81\).
\end{proposition}
\begin{proof}
La cardinalité transportée est \(2N_6=180\). La normalisation par l'espace ambiant
donne \(180/729\) ; diviser le numérateur et le dénominateur par \(9\) produit \(20/81\).
Aucune comparaison avec \(180/\pi\) ni avec \(\epsR\) n'intervient.
\end{proof}

La lecture historique comme projection sur un réseau \(1/81\) est un contrôle de
rigidité utile, mais n'est plus le fondement : le coefficient s'obtient avant
la clôture. L'ablation d'un feuillet produit \(10/81\), et la sortie change à
l'ordre \(10^{-5}\) ; le facteur deux n'est donc pas décoratif.

La section torsionnelle est
\begin{equation}
\Phi_T=\rho_{\mathrm{tw}}\Delta_4,
\qquad
S_T=1+\Phi_T.
\label{eq:seccion-torsional}
\end{equation}

 \subsection{La retenue \(\varepsilon_R\) comme cocycle de transition entre sections}
% Fragmento público generado desde fuentes teoremáticas íntegros.
% Residencia: 52.2.3.
% Fuente: /Users/ruben/Documents/HMT2/EL_CIERRE_HOLOGRAFICO_DEL_INFINITO_SUCESOR_102_CAPITULOS_2026-08-28_EN_TRABAJO/incoming/radion_monografia_20260827/manuscrito/sections/03_cierre_exacto.tex:104-253
\subsubsection{Cocycle de transition entre les sections directe et torsionnelle}

\paragraph{Définition opératoire de la retenue \(\varepsilon_R\)}

\begin{definition}[Retenue réelle du radion]
La sortie unitaire du radion est le cocycle de transition
\begin{equation}
\epsone:=S_T-S_E
=\Phi_T-D_E
=\frac{20}{81}\Delta_4-operatorname{rem}_1(S_E).
\label{eq:epsilon-def}
\end{equation}
\end{definition}

L'équation contient une soustraction, mais ce n'est pas un ajustement à la cible. Ses deux
opérandes \(S_T\) et \(S_E\) sont générés séparément. L'indépendance
démontrée est une indépendance par rapport à \(180/\pi\) et à un
\(\varepsilon\) cible ; on n'affirme pas d'indépendance algébrique entre une
différence et ses termes.

\begin{theorem}[Clôture exacte unitaire]
Avec \(S_E\), \(\Phi_T\) et \(\epsone\) définis par
\eqref{eq:seccion-electrica}, \eqref{eq:delta4}, \eqref{eq:rho-twoway} et
\eqref{eq:epsilon-def}, on a exactement
\begin{equation}
\boxed{1=S_E-\Phi_T+\epsone.}
\label{eq:cierre-unitario}
\end{equation}
\end{theorem}
\begin{proof}
Par \eqref{eq:SE-intervalo}, \(D_E=S_E-1\). Par définition,
\(\epsone=\Phi_T-D_E\). Alors
\[
S_E-\Phi_T+\epsone
=S_E-\Phi_T+\Phi_T-(S_E-1)=1.
\]
La force de la preuve ne réside pas dans cette annulation élémentaire, mais dans le fait que la
généalogie préalable fixe les deux opérandes sans utiliser la clôture comme entrée.
\end{proof}

\paragraph{Publication de la clôture dans la carte angulaire}

La carte du tour multiplie une fraction unitaire par
\(360/T_+=180/\pih\). En particulier,
\[
\epsR^\circ=\frac{360^\circ}{T_+}\epsone.
\]
Multiplier \eqref{eq:cierre-unitario} par \(360^\circ/T_+\) et utiliser
\[
\frac{360}{T_+}S_E
=360\left(\frac5{36}+\frac{25}{9}\alphah\right)
=50+1000\alphah=50+A
\]
produit la clôture angulaire.

\begin{theorem}[Équation exacte du radion]
\BigRadionEquation
\end{theorem}

L'équation dit quelque chose de plus précis que \enquote{le radian est
\(57.2957\ldots{}^\circ\)}. Elle décompose cette publication en quatre opérations
avec démonstration de référence :
\begin{center}
\begin{tabularx}{.96\textwidth}{>{\bfseries}l X}
\toprule
\(50^\circ\) & deuxième phase de la roue ; charnière critique dans une carte déclarée ;\\
\(A^\circ\) & coordonnée commune de la publication parangulaire ;\\
\(-\frac{20}{81}\Delta_4^\circ\) & transport orienté de la torsion globale ;\\
\(+\epsR^\circ\) & retenue réelle qui raccorde les deux relèvements.\\
\bottomrule
\end{tabularx}
\end{center}

\paragraph{Convergence selon la profondeur du raffinement}

À la profondeur \(N\), le générateur fournit des cylindres rationnels
\(p_N\to\pih\) et \(e_N\to e\). On définit
\[
S_{E,N}=2p_N\left(\frac5{36}+\frac{25}{9}\alphah\right),
\]
\[
\Delta_{4,N}=\frac{p_N-e_N\log p_N}{270}
\left(1+\frac{p_N}{729}\right),\qquad
S_{T,N}=1+\frac{20}{81}\Delta_{4,N},
\]
et
\[
\mathcal E_N=S_{T,N}-S_{E,N}.
\]
La continuité de la fonctionnelle en \(3<p<4\), \(2<e<3\), jointe aux
dérivées monotones, transporte chaque couple de cylindres vers un intervalle
certifié. À la profondeur de \(45\) blocs, sa largeur est inférieure à
\(4.81\times10^{-130}\).

Chaque retour de neuf niveaux contient six trits par niveau. C'est pourquoi le taux
de raffinement est
\begin{equation}
729^{-9}=3^{-54}.
\label{eq:contraccion-54}
\end{equation}
Ce chiffre contrôle la largeur des intervalles ; ce n'est pas la valeur de
\(\epsR\).

La soustraction de triades stabilise
\begin{center}
\ttfamily
000|000|000|896|802|822|044|562|981|999|050|695|020|651|286|413
\end{center}
et le préfixe de retenues
\begin{center}
\ttfamily
0,0,0,0,0,-1,0,-1,-1,-1,0,-1,0.
\end{center}
Tel est le mécanisme concret de mémoire décimale : non une queue choisie, mais
la récurrence \eqref{eq:subacarreo} appliquée à des opérandes préalablement
construits.

\paragraph{Évaluation des valeurs canoniques}

La sortie coinductive est
\begin{align}
\epsone
&=8.9680282204456298199905069502065128641372736205009\ldots
\times10^{-10},\\
\epsR^\circ
&=5.1383016758575282071685164622645947489908574400054\ldots
\times10^{-8}\,{}^\circ.
\label{eq:epsilon-valores}
\end{align}
La publication complète donne
\[
\frac{180^\circ}{\pih}
=57.2957795130823208767981548141051703324054724665643\ldots{}^\circ.
\]

\begin{exactbox}
\begin{align*}
50^\circ+A^\circ
&=57.29735256928380099728510547238066299956\ldots{}^\circ,\\
\frac{20}{81}\Delta_4^\circ
&=0.00157310758449687906223272996065728980465\ldots{}^\circ,\\
50^\circ+A^\circ-\frac{20}{81}\Delta_4^\circ
&=57.29577946169930411822287274242000570976\ldots{}^\circ,\\
\epsR^\circ
&=0.00000005138301675857528207168516462264594749\ldots{}^\circ,\\
\text{somme finale}
&=57.29577951308232087679815481410517033241\ldots{}^\circ.
\end{align*}
\end{exactbox}

 
\section{Coordonnée critique et système temporel dodécaphasique orienté : détermination affine de la phase de \(50^\circ\)}
\subsection{Application affine \(j_{100}\) de la carte APP et unicité de la phase \(m=2\) dans \(j_{100}(1/2)=\theta_m=50^\circ\)}
% Fragmento público generado desde fuentes teoremáticas íntegros.
% Residencia: 52.3.1.
% Fuente: /Users/ruben/Documents/HMT2/EL_CIERRE_HOLOGRAFICO_DEL_INFINITO_SUCESOR_102_CAPITULOS_2026-08-28_EN_TRABAJO/incoming/radion_monografia_20260827/manuscrito/sections/04_cincuenta_critica.tex:1-105
\noindent La coordonnée critique relie la carte de Cayley, le registre APP et la phase dodécaphasique par des applications typées.\par\medskip

\subsubsection{Deux réalisations exactes de la coordonnée angulaire \(50^\circ\)}

\paragraph{Deuxième phase du système temporel dodécaphasique}

Par \eqref{eq:rueda-dodecafasica}, la suite des centres commence
\[
20^\circ,\ 50^\circ,\ 80^\circ,\ 110^\circ,\ldots.
\]
Résoudre
\[
30m-10=50
\]
donne uniquement \(m=2\). Par conséquent, \(50^\circ\) n'est ni un arrondi de \(A\),
ni un demi-tour ni le centre d'un cercle : c'est une phase typée du calendrier
dodécaphasique. La fraction de tour est
\[
\frac{50}{360}=\frac5{36},
\]
exactement le premier terme de \(S_E/T_+\).

\paragraph{Égalisation de la coordonnée critique de Cayley}

L'involution fonctionnelle
\[
\rho\longmapsto1-\bar\rho
\]
fixe
\[
\operatorname{Fix}(\rho\mapsto1-\bar\rho)
=\{\rho:\Re\rho=1/2\}.
\]
Le nombre \(1/2\) est la coordonnée réelle autoduale de l'intervalle fonctionnel
\(0\leq\Re s\leq1\). Ce n'est pas \enquote{la moitié d'une droite} au sens
longitudinal, et il ne contient pas la hauteur spectrale \(\gamma\). Pour
\[
\rho_\gamma=\frac12+i\gamma,
\]
la transformation de Cayley peut s'écrire
\begin{equation}
\tau_\gamma
=\frac{\rho_\gamma-1}{\rho_\gamma}
=\frac{\gamma+i/2}{\gamma-i/2}
\in S^1,
\qquad
\gamma=\frac12\cot\frac{\theta_\gamma}{2}.
\label{eq:cayley-critica}
\end{equation}
Ainsi, \(1/2\) fixe la couture, \(\gamma\) détermine le caractère angulaire et
\(\Gamma_9\) conserve la profondeur de ses itérations.

\paragraph{Injection affine \(j_{100}\) du registre APP dans la carte angulaire}

La complétion cubique dispose une chambre de dix marques par coordonnée. Une
face contient \(10^2=100\) positions. Nous déclarons la carte affine
\begin{equation}
j_{100}:[0,1]\longrightarrow[0,100^\circ],
\qquad j_{100}(\sigma)=100\sigma^\circ.
\label{eq:j100}
\end{equation}
Alors
\[
j_{100}\left(\frac12\right)=50^\circ.
\]
En exigeant simultanément \(j_{100}(1/2)=\theta_m\), il vient
\[
100\cdot\frac12=30m-10
\quad\Longleftrightarrow\quad m=2.
\]

\begin{theorem}[Égalisateur critique--dodécaphasique]
Dans la carte APP de cent marques, la coordonnée fixe de l'involution critique
se publie exactement dans la deuxième phase de la roue :
\begin{equation}
\boxed{j_{100}(1/2)=\theta_2=50^\circ.}
\label{eq:igualador-50}
\end{equation}
La phase \(m=2\) est unique.
\end{theorem}

Le théorème est exact une fois \eqref{eq:j100} déclarée. Son statut est
\status{formalisation nouvelle}: le demi critique, la complétion \(10^3\), la
face de cent marques et \(\theta_2\) étaient déjà construits ; le choix de cette
face comme carte critique affine réunit pour la première fois leurs domaines.

\paragraph{Compatibilité de la valeur critique avec la carte de Cayley}

La transformation de Cayley n'envoie pas \(1/2\) sur \(50^\circ\). Si
\(\gamma=0\), \eqref{eq:cayley-critica} donne
\[
\tau_0=-1=e^{i\pi},
\]
c'est-à-dire \(180^\circ\). Les \(50^\circ\) appartiennent à la carte APP de cent
marques et au calendrier dodécaphasique ; l'angle de Cayley dépend de \(\gamma\).
Cette distinction empêche de réduire tous les zéros de la fonction zêta à une seule
phase.

\begin{falsifier}
L'identification \(1/2\leftrightarrow50^\circ\) échoue si l'on élimine la
carte \eqref{eq:j100}. L'égalité scalaire \(100(1/2)=50\) ne suffit pas
à identifier à elle seule deux objets. Le théorème dépend de la face APP
distinguée et de la roue dodécaphasique.
\end{falsifier}

 \subsection{Compatibilité de la phase critique avec le retour nonadique et l'avance de mémoire}
% Fragmento público generado desde fuentes teoremáticas íntegros.
% Residencia: 52.3.2.
% Fuente: /Users/ruben/Documents/HMT2/EL_CIERRE_HOLOGRAFICO_DEL_INFINITO_SUCESOR_102_CAPITULOS_2026-08-28_EN_TRABAJO/incoming/radion_monografia_20260827/manuscrito/sections/04_cincuenta_critica.tex:106-171
\subsubsection{Séparation typée des 2 objets de valeur \(1/2\)}

APP contient en outre un mode singulier avec
\[
s=\frac12,\qquad s^2=\frac14,\qquad
1-s^2=\frac34=\frac{90}{120},\qquad
s\sqrt{1-s^2}=\frac{\sqrt3}{4}.
\]
Ce \(s\) naît des angles principaux des feuillets somme et produit. La
couture critique \(\Re\rho=1/2\) naît d'une involution fonctionnelle. Que les deux
coordonnées valent \(1/2\) est une coïncidence scalaire dotée d'un contenu possible,
mais n'autorise pas à les identifier sans un carré typé.

La manière correcte de conserver la découverte est un diagramme d'égalisateurs :
\[
\begin{tikzpicture}[baseline=(current bounding box.center),node distance=14mm and 23mm,
 every node/.style={font=\small}]
\node[draw=RadBlue,rounded corners,fill=RadMist] (crit) {costura crítica \(\Re s=1/2\)};
\node[draw=RadTeal,rounded corners,fill=RadCream,right=of crit] (face) {carta \(j_{100}\)};
\node[draw=RadCopper,rounded corners,fill=RadCream,right=of face] (theta) {fase \(\theta_2=50^\circ\)};
\node[draw=RadRose,rounded corners,fill=white,below=of face] (mode) {modo APP \(s=1/2\), \(1-s^2=90/120\)};
\draw[-{Latex},thick,RadBlue] (crit)--node[above]{\(j_{100}\)}(face);
\draw[-{Latex},thick,RadTeal] (face)--node[above]{igualador}(theta);
\draw[-{Latex},dashed,RadRose] (mode)--node[right,align=left]{misma cifra;\\dominio distinto}(face);
\end{tikzpicture}
\]
Le diagramme ne simule pas la flèche absente. Il montre ce qui est démontré, ce qui a
été formalisé et quelle égalité nécessite encore un transport supplémentaire si
l'on vise une identification opératoire globale.

\subsubsection{Compatibilité entre cercle spectral et mémoire hélicoïdale}

Le lecteur de Cayley--Pontryagin transporte une fibre spectrale au moyen de
\[
\widetilde U_{\Gamma_9^n}J\psi_\gamma
=\tau_\gamma^nJ\psi_\gamma.
\]
La publication circulaire est \(\tau_\gamma^n\in S^1\). Le relèvement
\[
(\tau_\gamma^n,n)\in S^1\times\Z
\]
conserve le nombre de retours. Sur le feuillet contractant apparaît
\[
\widetilde M_9^nJ\psi_\gamma
=\left(\frac59\right)^n\tau_\gamma^nJ\psi_\gamma.
\]
Nous avons donc trois images typées :
\[
\text{cercle de phase},\qquad
\text{hélice de mémoire},\qquad
\text{spirale contractante}.
\]

Un zéro n'est pas un tour, et un neuf n'est pas un zéro particulier de la fonction
zêta. Le zéro fixe un caractère \(\tau_\gamma\) qui persiste à travers les
tours ; le retour nonadique incrémente la mémoire. Le neuf fonctionne comme zéro
résiduel dans \(\Z/9\Z\), opération arithmétique distincte de l'énumération
des zéros spectraux.

\begin{narrativebox}
La droite critique devient un cercle lorsque nous oublions la hauteur linéaire et
conservons le caractère unitaire. Elle devient une hélice lorsque nous réinsérons
la mémoire du nombre de fois où le caractère a été transporté. La profondeur
n'est pas une troisième coordonnée décorative : c'est l'information que la projection
circulaire ne peut pas retenir.
\end{narrativebox}
 
\section{Couple angulaire et action réduite : construction de l'ellipse positive et de l'aire par radion}
\subsection{Loi surfacique du radion : \(\mathcal A(\theta)=\hbar_{\mathrm{ret}}\theta\) et \(\mathcal A(2\pi)=h_{\mathrm{ret}}\)}
% Fragmento público generado desde fuentes teoremáticas íntegros.
% Residencia: 52.4.1.
% Fuente: /Users/ruben/Documents/HMT2/EL_CIERRE_HOLOGRAFICO_DEL_INFINITO_SUCESOR_102_CAPITULOS_2026-08-28_EN_TRABAJO/incoming/radion_monografia_20260827/manuscrito/sections/05_geometria_accion.tex:1-217
\noindent La géométrie spectrale détermine la forme et la loi surfacique détermine l'action transportée par la phase.\par\medskip

\subsubsection{Équivalence entre les cartes spectrale et géométrique}

\paragraph{Opérateur spectral positif et valeurs propres conjuguées}

Le couple angulaire produit l'opérateur bicouche
\begin{equation}
B_1=AI+C^*R,\qquad R^2=I,\qquad
P_\pm=\frac{I\pm R}{2}.
\label{eq:B1}
\end{equation}
Ses valeurs propres sont \(A\pm C^*\). Comme \(0<C^*<A\), les deux sont positives.
Dans la carte \(\kappa=\pi/180\), l'ellipse spectrale a pour demi-axes
\begin{equation}
a_1^2=\kappa(A+C^*),\qquad
b_1^2=\kappa(A-C^*).
\label{eq:semiejes-espectrales}
\end{equation}
La décomposition est réversible :
\begin{equation}
A=\frac{a_1^2+b_1^2}{2\kappa},\qquad
C^*=\frac{a_1^2-b_1^2}{2\kappa}.
\label{eq:reconstruccion-AC}
\end{equation}
L'ellipse n'est pas dessinée pour illustrer un nombre. Elle reconstruit exactement le
couple qui l'a générée.

Définissons
\begin{equation}
\rho_1=\sqrt{A^2-C^{*2}},\qquad
\eta=\operatorname{artanh}\frac{C^*}{A}
=\frac12\log\frac{A+C^*}{A-C^*}.
\label{eq:eta}
\end{equation}
Alors
\[
A\pm C^*=\rho_1e^{\pm\eta},
\]
et
\begin{equation}
a_1=\sqrt{\kappa\rho_1}\,e^{\eta/2},
\qquad
b_1=\sqrt{\kappa\rho_1}\,e^{-\eta/2}.
\label{eq:semiejes-eta}
\end{equation}
Le produit conserve l'échelle ; le quotient conserve la forme :
\[
a_1b_1=\kappa\rho_1,
\qquad
\frac{b_1}{a_1}=e^{-\eta}
=\sqrt{\frac{A-C^*}{A+C^*}}.
\]

\paragraph{Invariants focaux de l'opérateur positif}

La demi-distance focale est
\begin{equation}
c_1^2=a_1^2-b_1^2=2\kappa C^*,
\qquad
F_\pm=(\pm c_1,0).
\label{eq:focos}
\end{equation}
La distance totale vaut
\begin{equation}
d_1=2c_1=2\sqrt{2\kappa C^*}
=0.544545509325626623406299779866023\ldots.
\label{eq:d1}
\end{equation}
Le nombre isolé n'est pas la thèse. L'identité \(c_1^2=2\kappa C^*\) dit que
\(C^*\) est littéralement la scission focale quadratique dans la carte
distinguée. Les foyers sont les extrémités conjuguées de l'ouverture spectrale ;
ils ne sont pas, à eux seuls, les pôles électrique et magnétique. Cette lecture nécessite
le transducteur constitutif.

L'excentricité est
\begin{equation}
e_f=\frac{c_1}{a_1},\qquad
e_f^2=1-e^{-2\eta}=\frac{2C^*}{A+C^*}.
\label{eq:excentricidad}
\end{equation}
Numériquement,
\[
\begin{aligned}
\eta&=0.299689683921630799684926562863927\ldots,\\
e_f&=0.671451895550191986916277055864332\ldots.
\end{aligned}
\]

\paragraph{Bloc APP de référence et transport normalisé}

Le bloc central
\[
B_0=7I+2R=9P_++5P_-
\]
possède
\[
\rho_0=\sqrt{45},\qquad
\eta_0=\frac12\log\frac95,\qquad
e_0=\frac23.
\]
Sa distance focale est
\[
d_0=4\sqrt\kappa
=0.528443639680801468446003835349358\ldots.
\]
La comparaison exacte est
\begin{equation}
\boxed{\left(\frac{d_1}{d_0}\right)^2=\frac{C^*}{2}},
\qquad
d_1=d_0\sqrt{\frac{C^*}{2}}.
\label{eq:invariante-focal-relativo}
\end{equation}
Telle est la généalogie défendable de \(0.544545\ldots\) : une échelle focale APP
de référence multipliée par un invariant relatif. Sa proximité visuelle avec
\(0.54\) n'est pas une identité avec le défaut APP \(54\).

Le transport depuis \(B_0\) se sépare en échelle et anisotropie :
\begin{equation}
T_{\mathrm{rel}}=\frac{\rho_1}{\sqrt{45}}
\exp\bigl((\eta-\eta_0)R\bigr).
\label{eq:transporte-rel}
\end{equation}
Sur les feuillets,
\[
t_+=\frac{A+C^*}{9}=1.0467878786947163\ldots,
\qquad
t_-=\frac{A-C^*}{5}=1.0347228460630311\ldots.
\]
La déformation n'est pas une seule décimale. Sa partie isotrope est
\(\sqrt{t_+t_-}\) ; sa partie hyperbolique est
\(\tfrac12\log(t_+/t_-)\).

\subsubsection{Ellipse positive associée au couple angulaire}

\paragraph{Construction au moyen de demi-axes conjugués}

L'échelle spectrale précédente est dépourvue d'unités d'action. La mécanique
dimensionnelle publie la même forme sur la droite interne d'action :
\begin{equation}
a_S=\sqrt{2\hret e^\eta},
\qquad
b_S=\sqrt{2\hret e^{-\eta}}.
\label{eq:semiejes-accion}
\end{equation}
Immédiatement,
\begin{equation}
a_Sb_S=2\hret,\qquad
\frac{a_S}{b_S}=e^\eta.
\label{eq:producto-forma}
\end{equation}
Le produit fixe l'action ; le quotient fixe l'anisotropie. L'ellipse est
paramétrée par
\begin{equation}
Q(\theta)=a_S\cos\theta,\qquad
P(\theta)=b_S\sin\theta.
\label{eq:elipse-param}
\end{equation}

Les ellipses spectrale et d'action sont homothétiques :
\begin{equation}
\lambda_{\mathrm{act}}^2=\frac{2\hret}{\kappa\rho_1},
\qquad
(a_S,b_S,c_S)=\lambda_{\mathrm{act}}(a_1,b_1,c_1).
\label{eq:homotecia}
\end{equation}
Elles partagent la forme \(e^{-\eta}\), non la dimension physique.

\paragraph{Loi surfacique de l'action par radion}

La forme symplectique canonique est \(\omega=\diff Q\wedge\diff P\). L'aire
orientée balayée de \(0\) à \(\theta\) est
\begin{align}
\mathcal A(\theta)
&=\frac12\int_0^\theta
\bigl(Q(t)P'(t)-P(t)Q'(t)\bigr)\,\diff t\\
&=\frac12a_Sb_S\int_0^\theta
(\cos^2t+\sin^2t)\,\diff t.
\end{align}
En utilisant \eqref{eq:producto-forma}, nous obtenons le résultat central.

\begin{theorem}[Action par radion]
Pour l'ellipse \eqref{eq:elipse-param},
\begin{equation}
\boxed{\mathcal A(\theta)=\hret\theta.}
\label{eq:area-theorem}
\end{equation}
En particulier,
\begin{equation}
\boxed{\mathcal A(1)=\hret},
\qquad
\boxed{\mathcal A(2\pi)=2\pi\hret=\hfull}.
\label{eq:area-radion-vuelta}
\end{equation}
\end{theorem}

Cette égalité permet de définir le radion sans faire d'abord appel à la longueur
d'arc.

\begin{definition}[Radion orienté]
\begin{equation}
\mathfrak r_\sigma=(\theta=1,\sigma),
\qquad \sigma\in\{+1,-1\},
\qquad
\mathcal A(\mathfrak r_\sigma)=\sigma\hret.
\label{eq:radion-oriented}
\end{equation}
\end{definition}

Le radian est la projection qui oublie \(\sigma\), le feuillet, la mémoire et
l'aire. Le radion conserve cette information. D'où le nom \emph{rad--ion} :
l'orientation précède la charge physique, et la charge apparaît plus tard comme
\[
Q(x)=n_Q(x)e_{\mathrm{int}},\qquad
e_{\mathrm{int}}=\sqrt{\frac{2\alphah\hfull}{Z_{\mathrm{int}}}}.
\]

 \subsection{Demi-axes conjugués, invariants focaux et orientation de l'anisotropie}
% Fragmento público generado desde fuentes teoremáticas íntegros.
% Residencia: 52.4.2.
% Fuente: /Users/ruben/Documents/HMT2/EL_CIERRE_HOLOGRAFICO_DEL_INFINITO_SUCESOR_102_CAPITULOS_2026-08-28_EN_TRABAJO/incoming/radion_monografia_20260827/manuscrito/sections/05_geometria_accion.tex:218-290
\paragraph{Correspondance entre phase, action, aire et orientation}

\begin{figure}[H]
\centering
\begin{tikzpicture}[scale=1.12]
  \def\aa{4.0}
  \def\bb{2.95}
  \def\cc{2.70}
  \draw[->,RadGray] (-4.7,0)--(4.7,0) node[right]{\(Q\)};
  \draw[->,RadGray] (0,-3.5)--(0,3.5) node[above]{\(P\)};
  \fill[RadTeal!8] (0,0) ellipse[x radius=\aa,y radius=\bb];
  \draw[RadTeal,line width=1.35pt] (0,0) ellipse[x radius=\aa,y radius=\bb];
  \fill[RadCopper] (-\cc,0) circle (2pt) node[below=2mm]{\(F_-\)};
  \fill[RadCopper] (\cc,0) circle (2pt) node[below=2mm]{\(F_+\)};
  \draw[RadGold,line width=1.1pt,-{Latex}] (\aa,0)
     arc[start angle=0,end angle=57.3,x radius=\aa,y radius=\bb];
  \node[RadGold] at (3.7,2.3) {un radion};
  \draw[dashed,RadBlue] (0,0)--(\aa,0) node[midway,below]{\(a_S\)};
  \draw[dashed,RadBlue] (0,0)--(0,\bb) node[midway,left]{\(b_S\)};
  \node[align=center,RadNavy] at (0,-3.25)
    {aire totale \(\pi a_Sb_S=2\pi\hret=\hfull\)};
\end{tikzpicture}
\caption{L'ellipse d'action. Un arc paramétrique d'un radian balaie
\(\hret\) ; un tour enferme \(\hfull\). Les foyers codent l'ouverture
\(C^*\) dans la carte distinguée.}
\label{fig:elipse-accion}
\end{figure}

\paragraph{Relation entre l'action réduite, l'incertitude et l'aire}

La relation \(\hbar=h/(2\pi)\) est souvent présentée comme une commodité
de notation : lorsqu'on décrit des oscillations, des rotations et des transformées de Fourier,
les facteurs \(2\pi\) disparaissent si l'on utilise \(\hbar\). L'ellipse d'action
révèle une lecture plus forte. La division par \(2\pi\) est exactement le
passage de l'action d'un tour à l'action par unité de phase :
\[
h=\mathcal A(2\pi),\qquad \hbar=\mathcal A(1).
\]
La constante réduite n'est pas seulement \(h\) avec une notation plus commode ; c'est la
densité surfacique d'action par rapport à la phase.

Heisenberg, Dirac et Jordan ont établi le langage des variables conjuguées,
des commutateurs et des transformations canoniques. On ne leur attribue pas ici une
\enquote{ellipse ontologique} littérale sans source. La contribution HMT--MD consiste à
identifier la figure positive, à construire ses demi-axes à partir de \((A,C^*)\) et à
démontrer \eqref{eq:area-theorem}. La généalogie historique ouvre le langage ;
le théorème actuel fixe l'objet.

\begin{narrativebox}
L'action ne s'ajoute pas ensuite au mouvement. Dans l'ellipse du radion,
l'action est l'aire que le mouvement génère. Un tour ne transporte pas une
constante préfabriquée : il l'enferme.
\end{narrativebox}

\Needspace{10\baselineskip}
\paragraph{Contrôles algébriques des invariants focaux}

Les ressemblances décimales sont soumises à un contrôle :
\begin{center}
\begin{tabularx}{.96\textwidth}{l r X}
\toprule
Comparaison & Résidu & Verdict\\
\midrule
\(d_1\) face à \(0.54\) & \(4.5455\times10^{-3}\) & pas d'égalité ; le \(54\) ne se déduit pas ;\\
\(d_1^2\) face à \(8/27\) & \(2.3352\times10^{-4}\) & proximité sans application ;\\
\(\eta\) frente a \(3/10\) & \(-3.1032\times10^{-4}\) & no igualdad;\\
\(e_f\) face à \(2/3\) & \(4.7852\times10^{-3}\) & \(2/3\) appartient exactement à \(B_0\) ;\\
\(e^{-\eta}\) face à \(20/27\) & \(3.0740\times10^{-4}\) & no igualdad.\\
\bottomrule
\end{tabularx}
\end{center}
Les invariants solides sont \eqref{eq:eta}, \eqref{eq:excentricidad} et
\eqref{eq:invariante-focal-relativo}. Le critère HMT ne consiste pas à récompenser
une ressemblance décimale, mais à exiger objet, opération et conservation.
 
\section{Retour nonadique et intérieur projectif : disque de Klein et profondeur hélicoïdale de mémoire}
\subsection{Normalisation au disque et intérieur hyperbolique de Hilbert--Beltrami--Klein}
% Fragmento público generado desde fuentes teoremáticas íntegros.
% Residencia: 52.5.1.
% Fuente: /Users/ruben/Documents/HMT2/EL_CIERRE_HOLOGRAFICO_DEL_INFINITO_SUCESOR_102_CAPITULOS_2026-08-28_EN_TRABAJO/incoming/radion_monografia_20260827/manuscrito/sections/06_interior_helice.tex:1-108
\noindent La géométrie projective distingue l'action finie de la profondeur hyperbolique et de la mémoire du retour.\par\medskip

\subsubsection{Intérieur hyperbolique de Hilbert--Beltrami--Klein}

\paragraph{Normalisation projective dans le disque de Klein}

L'application affine
\[
u=\frac{Q}{a_S},\qquad v=\frac{P}{b_S}
\]
envoie l'ellipse d'action sur le disque unité
\[
D=\{(u,v):u^2+v^2<1\}.
\]
Sur chaque corde, la distance de Hilbert est définie par le birapport
de ses extrémités de frontière. La métrique infinitésimale résultante est la forme
de Beltrami--Klein
\begin{equation}
\diff s_K^2=
\frac{(1-r^2)(\diff u^2+\diff v^2)+(u\diff u+v\diff v)^2}
{(1-r^2)^2},
\qquad r^2=u^2+v^2.
\label{eq:klein-metric}
\end{equation}
Avec cette normalisation, la courbure gaussienne est constante :
\begin{equation}
K=-1.
\label{eq:klein-curvature}
\end{equation}
Les géodésiques sont des cordes euclidiennes, mais leur longueur hyperbolique diverge
à l'approche de la frontière.

\paragraph{Finitude de la frontière et infinitude de la profondeur hyperbolique}

Le déterminant métrique de \eqref{eq:klein-metric} produit
\begin{equation}
\diff A_K=\frac{\diff u\,\diff v}{(1-r^2)^{3/2}}.
\label{eq:klein-area-density}
\end{equation}
L'aire du sous-disque \(r<R<1\) est
\begin{align}
A_K(R)
&=\int_0^{2\pi}\int_0^R
\frac{r\,\diff r\,\diff\theta}{(1-r^2)^{3/2}}\\
&=2\pi\left(\frac1{\sqrt{1-R^2}}-1\right).
\label{eq:klein-area-R}
\end{align}
\Needspace{4\baselineskip}
Par conséquent,
\[
\lim_{R\to1^-}A_K(R)=+\infty.
\]
Simultanément, la même frontière enferme l'aire symplectique
\[
\pi a_Sb_S=2\pi\hret=\hfull<\infty.
\]

\begin{theorem}[Finitude de l'action et infinitude de la profondeur]
L'ellipse du radion possède une action de tour finie \(\hfull\), tandis que
son intérieur, muni de la métrique de Hilbert--Beltrami--Klein, a une aire
hyperbolique infinie :
\begin{equation}
\boxed{
\operatorname{Area}_{\mathrm{simp}}(\partial\mathbb E_S)=\hfull<\infty,
\qquad
\operatorname{Area}_{K}(\mathbb E_S)=+\infty.}
\label{eq:finite-infinite}
\end{equation}
\end{theorem}

Il n'y a pas de contradiction : \(\operatorname{Area}_{\mathrm{simp}}\) et
\(\operatorname{Area}_K\) sont des mesures différentes. La première compte l'action
orientée en phase ; la seconde mesure la profondeur projective intérieure. La thèse
\enquote{bord fini, intérieur infini} est désormais une égalité et une limite,
non une métaphore.

\paragraph{Coordonnées focales dans l'intérieur de Klein}

Dans le disque normalisé, les foyers se situent en \((\pm e_f,0)\). La distance
du centre à un foyer est
\begin{equation}
u_f=\operatorname{artanh}e_f
=\operatorname{arcosh}(e^\eta)
=0.813382331175342333760889731088228\ldots.
\label{eq:klein-focal}
\end{equation}
La distance foyer--foyer est
\[
d_K(F_-,F_+)=2u_f
=1.62676466235068466752177946217646\ldots.
\]
On a
\begin{equation}
e_f=\tanh u_f,\qquad
e^{-\eta}=\operatorname{sech}u_f,\qquad
\eta=\log\cosh u_f.
\label{eq:focal-hyperbolic-chain}
\end{equation}
La région jusqu'au rayon focal a pour aire
\[
A_K(e_f)=2\pi(e^\eta-1)
=2.19559620884061869541206115980360\ldots.
\]

Toutes les ellipses ouvertes sont projectivement isométriques en tant que domaines de
Hilbert. C'est pourquoi la courbure \(-1\) seule ne mémorise pas l'excentricité. La
forme HMT réside dans la structure conjointe : carte affine distinguée, forme
symplectique, étiquetage des feuillets et paramétrisation de frontière.
 \subsection{Différence entre retour de phase et retour de l’état holonomique intégral}
% Fragmento público generado desde fuentes teoremáticas íntegros.
% Residencia: 52.5.2.
% Fuente: /Users/ruben/Documents/HMT2/EL_CIERRE_HOLOGRAFICO_DEL_INFINITO_SUCESOR_102_CAPITULOS_2026-08-28_EN_TRABAJO/incoming/radion_monografia_20260827/manuscrito/sections/06_interior_helice.tex:109-201
\subsubsection{Relèvement hélicoïdal du retour nonadique}

\paragraph{Décomposition nonadique et quotient de mémoire}

L'opération élémentaire
\begin{equation}
n=9q_9(n)+\rho_9(n),
\qquad
H_9(n)=(\rho_9(n),q_9(n))
\label{eq:H9}
\end{equation}
satisfait
\begin{equation}
H_9(n+9)=H_9(n)+(0,1).
\label{eq:H9-helix}
\end{equation}
La première coordonnée se ferme ; la seconde s'élève. Cette identité est la
version arithmétique minimale d'une vis.

Pour une fibre spectrale,
\[
n\longmapsto(\tau_\gamma^n,n)
\]
est une hélice sur \(S^1\). Ajouter la contraction \((5/9)^n\) produit une
spirale intérieure. Le radion participe aux deux lectures : une unité de phase
à la frontière et une unité orientée d'action dans le relèvement.

\paragraph{Rang de mémoire comme lecteur dimensionnel}

Supposons qu'une projection visible conserve la valeur de phase, mais perde
\(d\) cocycles de mémoire linéairement indépendants. Tout relèvement fidèle
nécessite au moins \(d\) coordonnées supplémentaires pour les reconstruire. Cela
motive une lecture précise de la mécanique dimensionnelle :
\begin{equation}
\begin{aligned}
\text{dimension émergente}
&=\text{rang minimal de mémoire indépendante}\\
&\phantom{=}\text{non publiable sur le feuillet visible}.
\end{aligned}
\label{eq:dimension-memory}
\end{equation}
On n'affirme pas que toute dimension physique soit un compteur. On affirme que, dans le
domaine TPK, chaque mémoire indépendante qui doit survivre impose une
coordonnée du relèvement.

\begin{narrativebox}
La dimension apparaît lorsque la projection ne peut plus conserver toute
l'histoire. Le cercle suffit à dire où se trouve la phase ; il ne
suffit pas à dire combien de fois elle est revenue, quel feuillet elle a transporté, quel intervalle
s'est contracté ou quelle retenue a survécu.
\end{narrativebox}

\paragraph{Assemblage de cylindres, de spirales et de trajectoires hélicoïdales}

Les cylindres coinductifs ne sont pas des tubes physiques ajoutés à une spirale. Ce sont
des intervalles emboîtés qui conservent les préfixes et resserrent la valeur. La spirale
décrit la contraction transversale ; l'hélice décrit la mémoire longitudinale ;
le cylindre décrit l'ensemble des valeurs compatibles à une profondeur
finie. À la limite, l'intersection des cylindres sélectionne une section
unique.

L'image spatiale réunit, sans confondre :
\[
\begin{array}{ccl}
\text{angle} &\leftrightarrow& \text{phase visible},\\
\text{hauteur} &\leftrightarrow& \text{retour/mémoire},\\
\text{rayon de cylindre} &\leftrightarrow& \text{incertitude de raffinement},\\
\text{feuillet} &\leftrightarrow& \text{orientation bidirectionnelle}.
\end{array}
\]
Le taux \(3^{-54}\) contracte le cylindre à chaque retour ; il ne s'ajoute pas à l'angle et
ne remplace pas la retenue réelle \(\epsR\).

\subsubsection{Clôtures de phase d'ordres \(2\pi\) et \(4\pi\)}

Dans la projection circulaire, \(2\pi\) restaure la phase. Dans un relèvement orienté
à deux feuillets, un tour peut se fermer sur le feuillet conjugué et nécessiter un
second tour pour restaurer l'orientation :
\[
U_{2\pi}=-I,\qquad U_{4\pi}=I.
\]
La structure coïncide avec la périodicité spinorielle lorsque l'on applique le
foncteur à une représentation \(SU(2)\) ; on n'identifie pas le relèvement TPK à
\(SU(2)\) sans cette application.

La lecture surfacique est compatible :
\[
\mathcal A(2\pi)=\hfull,\qquad
\mathcal A(4\pi)=2\hfull.
\]
Sur un ruban de Möbius, un circuit visible inverse le feuillet et deux circuits rétablissent l’orientation. Le \(2\pi/4\pi\) n’est pas un doublement accidentel ; c’est la différence entre retour de phase et retour de l’état holonomique intégral.
 
\section{Ellipse orientée et variables conjuguées : incertitude minimale et oscillation électromagnétique}
\subsection{Variables symplectiques conjuguées et détermination de l'incertitude minimale}
% Fragmento público generado desde fuentes teoremáticas íntegros.
% Residencia: 52.6.1.
% Fuente: /Users/ruben/Documents/HMT2/EL_CIERRE_HOLOGRAFICO_DEL_INFINITO_SUCESOR_102_CAPITULOS_2026-08-28_EN_TRABAJO/incoming/radion_monografia_20260827/manuscrito/sections/07_incertidumbre_oscilador.tex:1-111
\noindent La forme symplectique coordonne action, incertitude et dynamique oscillatoire sans modifier la généalogie de l'état.\par\medskip

\subsubsection{Réalisation de l'état dans un espace de Hilbert}

\paragraph{Matrice de covariance minimale}

Une fois construite l'ellipse d'action, elle peut se réaliser dans un couple
d'opérateurs conjugués
\[
[\widehat Q,\widehat P]=i\hret.
\]
On définit
\begin{equation}
V_\eta=\frac{\hret}{2}
\begin{pmatrix}
e^\eta&0\\[2pt]0&e^{-\eta}
\end{pmatrix}.
\label{eq:covariance}
\end{equation}
Son déterminant est
\begin{equation}
\det V_\eta=\frac{\hret^2}{4}.
\label{eq:cov-det}
\end{equation}
Par conséquent,
\begin{equation}
(\Delta Q)^2=\frac{\hret}{2}e^\eta,
\qquad
(\Delta P)^2=\frac{\hret}{2}e^{-\eta},
\qquad
\Delta Q\,\Delta P=\frac{\hret}{2}.
\label{eq:uncertainty}
\end{equation}
L'inégalité de Robertson--Schrödinger est saturée.

L'annihilateur comprimé
\begin{equation}
\widehat a_\eta=
\frac{e^{-\eta/2}\widehat Q+i e^{\eta/2}\widehat P}
{\sqrt{2\hret}}
\label{eq:a-eta}
\end{equation}
satisfait \([\widehat a_\eta,\widehat a_\eta^\dagger]=1\). Son vide a pour
fonction d'onde
\begin{equation}
\psi_\eta(Q)=
(\pi\hret e^\eta)^{-1/4}
\exp\left(-\frac{Q^2}{2\hret e^\eta}\right).
\label{eq:gaussian}
\end{equation}

\paragraph{Ellipse comme niveau de la fonctionnelle de Mahalanobis}

L’ensemble
\begin{equation}
\mathbb E_{\mathrm{act}}
=\{x\in\R^2:x^{\mathsf T}V_\eta^{-1}x\le4\}
\label{eq:mahalanobis}
\end{equation}
a pour demi-axes
\begin{equation}
a_S=2\Delta Q,\qquad b_S=2\Delta P.
\label{eq:two-sigma}
\end{equation}
Son aire est
\[
4\pi\sqrt{\det V_\eta}=2\pi\hret=\hfull.
\]
L'ellipse HMT construite précédemment coïncide exactement avec le contour à deux
écarts de l'état gaussien minimal.

\begin{proposition}[Chaîne unique du quotient de forme]
Après avoir déclaré les interfaces spectrale, symplectique et de Hilbert,
\begin{equation}
\boxed{
e^{-\eta}
=\sqrt{\frac{A-C^*}{A+C^*}}
=\frac{b_1}{a_1}
=\frac{b_S}{a_S}
=\frac{\Delta P}{\Delta Q}.}
\label{eq:shape-chain-hilbert}
\end{equation}
\end{proposition}

L'ellipse d'action correspond au niveau classique \(H=\hret\omega\) et au
contour \(2\sigma\) du vide comprimé. Elle ne doit pas être confondue avec l'énergie
moyenne du vide, \(\hret\omega/2\). Cette différence d'un facteur deux est un
contrôle physique obligatoire.

\paragraph{Conditions d'invariance de l'ellipse comme objet opératoire}

Dans la formulation conventionnelle, une ellipse d'incertitude peut être considérée comme
un outil graphique pour représenter une matrice de covariance. Ici, l'ordre
causal est inversé : l'ellipse positive existe déjà comme publication de
\((A,C^*,\hret)\) ; l'espace de Hilbert la reconnaît au moyen de
\eqref{eq:covariance}. La réalité ontologique revendiquée ne signifie pas qu'une
particule classique parcourt physiquement un ovale dans le plan \((Q,P)\). Elle signifie
que l'état conserve une forme positive, une action totale et un quotient
d'orientation avant qu'une théorie probabiliste ne les interprète.

La fonction d'onde ne crée pas la forme ; elle la représente. L'incertitude ne crée pas
\(\hbar\) ; elle publie l'impossibilité de comprimer simultanément les deux
directions sans modifier le produit symplectique.

\begin{narrativebox}
La lecture probabiliste est une réalisation possible de l'ellipse, non sa cause.
L'objet premier est la conservation conjointe du produit et de l'orientation. La
probabilité apparaît lorsque cette géométrie est publiée comme covariance d'un
état sur Hilbert.
\end{narrativebox}

 \subsection{Échange \(\mu\leftrightarrow\varepsilon\), inversion \(Z\leftrightarrow Z^{-1}\) et conservation de la vitesse causale}
% Fragmento público generado desde fuentes teoremáticas íntegros.
% Residencia: 52.6.2.
% Fuente: /Users/ruben/Documents/HMT2/EL_CIERRE_HOLOGRAFICO_DEL_INFINITO_SUCESOR_102_CAPITULOS_2026-08-28_EN_TRABAJO/incoming/radion_monografia_20260827/manuscrito/sections/07_incertidumbre_oscilador.tex:156-197
\paragraph{Échange énergétique entre les secteurs électrique et magnétique}

Les deux termes de \eqref{eq:HLC} deviennent
\begin{equation}
U_E(\theta)=\hret\omega\cos^2\theta,
\qquad
U_B(\theta)=\hret\omega\sin^2\theta.
\label{eq:energy-exchange}
\end{equation}
Par conséquent,
\[
U_E+U_B=\hret\omega.
\]
La polarité électrique--magnétique ne consiste pas à étiqueter arbitrairement
les foyers. Elle consiste en deux réserves conjuguées qui s'échangent pendant
l'orbite tout en maintenant constantes la fréquence et l'action totale.

Le quotient commun se prolonge :
\begin{equation}
\boxed{
e^{-\eta}
=\frac{b_S}{a_S}
=\frac{\Delta P}{\Delta Q}
=\frac{Z_\eta}{Z_*}
=\sqrt{\frac{A-C^*}{A+C^*}}.}
\label{eq:shape-chain-LC}
\end{equation}
Les produits conservés sont
\[
a_Sb_S=2\hret,\qquad
\det V_\eta=\frac{\hret^2}{4},
\qquad
L_\eta C_\eta=\omega^{-2}.
\]

\begin{exactbox}
L'oscillateur LC est la réalisation physique minimale de la double lecture : le
produit conserve l'action et la fréquence ; le quotient conserve l'orientation et
l'impédance. Un cycle complet échange électricité et magnétisme sans
perdre l'aire \(\oint P\,\diff Q=\hfull\).
\end{exactbox}

 \subsection{Couple inductif--capacitif \((L_\eta,C_\eta)\) : invariance de \(L_\eta C_\eta=\omega^{-2}\) et orientation de l'impédance}
% Fragmento público generado desde fuentes teoremáticas íntegros.
% Residencia: 52.6.3.
% Fuente: /Users/ruben/Documents/HMT2/EL_CIERRE_HOLOGRAFICO_DEL_INFINITO_SUCESOR_102_CAPITULOS_2026-08-28_EN_TRABAJO/incoming/radion_monografia_20260827/manuscrito/sections/07_incertidumbre_oscilador.tex:112-155
\subsubsection{Oscillateur inductif--capacitif associé au couple angulaire}

\paragraph{Construction indépendante des 2 réservoirs conjugués}

Étant donnés \(\omega>0\) et une impédance de référence \(Z_*>0\), nous définissons
\begin{equation}
L_\eta=\frac{Z_*}{\omega}e^{-\eta},
\qquad
C_\eta=\frac{1}{Z_*\omega}e^\eta.
\label{eq:LC-def}
\end{equation}
Alors
\begin{equation}
L_\eta C_\eta=\omega^{-2},
\qquad
\sqrt{\frac{L_\eta}{C_\eta}}=Z_*e^{-\eta}.
\label{eq:LC-invariants}
\end{equation}
La fréquence reste fixe ; l'impédance conserve l'orientation de la
forme.

L'hamiltonien est
\begin{equation}
H_{LC}=\frac{q^2}{2C_\eta}+\frac{\Phi^2}{2L_\eta}.
\label{eq:HLC}
\end{equation}
Avec les variables normalisées
\begin{equation}
Q=\sqrt{Z_*}\,q,\qquad
P=\frac{\Phi}{\sqrt{Z_*}},
\label{eq:LC-map}
\end{equation}
on obtient
\begin{equation}
H_{LC}=\frac\omega2
\left(e^{-\eta}Q^2+e^\eta P^2\right).
\label{eq:HLC-QP}
\end{equation}
Sur \eqref{eq:elipse-param},
\begin{equation}
H_{LC}=\hret\omega.
\label{eq:LC-level}
\end{equation}

% Fuente: /Users/ruben/Documents/HMT2/EL_CIERRE_HOLOGRAFICO_DEL_INFINITO_SUCESOR_102_CAPITULOS_2026-08-28_EN_TRABAJO/incoming/radion_monografia_20260827/manuscrito/sections/07_incertidumbre_oscilador.tex:198-239
\paragraph{Oscillateur LC comme unité dynamique élémentaire}

L'oscillateur harmonique traverse presque toute la physique parce qu'il linéarise le
voisinage d'un équilibre, décompose les champs en modes, organise les quanta et
permet de lire la propagation comme une superposition de phases. Dans le radion, cette
universalité reçoit une généalogie concrète : le mode ne commence pas par une
équation différentielle externe, mais par une forme positive dont l'action de tour est
fixe. La dynamique conventionnelle apparaît lorsqu'on choisit une horloge \(\omega\) et un
transducteur \(Z_*\).

La fréquence ne modifie pas l'ellipse d'action : elle change l'énergie du niveau
au moyen de \(E=\hret\omega\). Cette séparation est cruciale. La géométrie fixe
l'action ; l'horloge convertit l'action en énergie.

\subsubsection{Réversion temporelle et propagations conjuguées}

Soit \(H=H^*\) et soit \(J_E\) une structure complexe avec \(J_E^2=-I\).
Supposons un opérateur involutif \(C\) tel que
\[
CJ_EC^{-1}=-J_E,\qquad CHC^{-1}=H.
\]
L'évolution
\begin{equation}
U(t)=\exp\left(-\frac{J_EHt}{\hret}\right)
\label{eq:evolution-J}
\end{equation}
satisfait exactement
\begin{equation}
CU(t)C^{-1}=U(-t).
\label{eq:time-reversal}
\end{equation}
Telle est la forme opératoire minimale de la bidirectionnalité : la conjugaison
échange les prolongations temporellement opposées.

La théorie de l'absorbeur de Wheeler--Feynman et la lecture des antiparticules
comme trajectoires temporellement conjuguées offrent une réalisation historique
ultérieure. L'identité \eqref{eq:time-reversal} est démontrée ; la transformer
en une théorie complète des propagateurs avancés et retardés exige de déclarer
l'action, les conditions de frontière, la causalité et les interactions de l'absorbeur.
La monographie conserve l'intuition physique sans transformer une correspondance
en une égalité de théories.

 
\section{Les trois régimes de TRIT : polarités électrique--magnétique et élastique--rotationnelle}
\subsection{Coordination des réalisations elliptique, parabolique et hyperbolique}
% Fragmento público generado desde fuentes teoremáticas íntegros.
% Residencia: 52.7.1.
% Fuente: /Users/ruben/Documents/HMT2/EL_CIERRE_HOLOGRAFICO_DEL_INFINITO_SUCESOR_102_CAPITULOS_2026-08-28_EN_TRABAJO/incoming/radion_monografia_20260827/manuscrito/sections/08_complejo_electromagnetismo.tex:1-58
\noindent Les 3 régimes du TRIT admettent des réalisations complexe, électromagnétique et mécanique avec des types différenciés.\par\medskip

\subsubsection{Opérateur quadratique de racine interne de \(-1\)}

\paragraph{Représentation exponentielle du régime elliptique \(J^2=-I\)}

Avec \(J^2=-I\),
\begin{equation}
T_{\mathrm{ell}}(x,y)=e^{-xI-yJ}
=e^{-x}(\cos y\,I-\sin y\,J).
\label{eq:elliptic-exp}
\end{equation}
Sous l'identification générée par \(J\),
\[
T_{\mathrm{ell}}(x,y)\longleftrightarrow e^{-x-iy}.
\]
Le scalaire \(x\) contrôle la dilatation ; \(y\) contrôle la rotation orientée. La forme
\(a+bi\) se reconnaît comme élasticité plus rotation, non comme deux quantités
collées arbitrairement.

\paragraph{Représentation exponentielle du régime parabolique \(J^2=0\)}

Avec \(N^2=0\), l'exponentielle se tronque :
\begin{equation}
T_{\mathrm{par}}(x,y)=e^{-x}(I-yN).
\label{eq:parabolic-exp}
\end{equation}
Le régime décrit le seuil où il n'y a ni rotation périodique ni ouverture
exponentielle. Dans la lecture temporelle active, c'est le présent euclidien, mais
il conserve la mémoire des branches qui le délimitent.

\paragraph{Représentation exponentielle du régime hyperbolique \(J^2=+I\)}

Avec \(R^2=I\) et
\(P_\pm=(I\pm R)/2\),
\begin{equation}
T_{\mathrm{hyp}}(x,y)=e^{-xI-yR}
=q_+P_++q_-P_-,
\qquad q_\pm=e^{-x\mp y}.
\label{eq:hyperbolic-exp}
\end{equation}
Les invariants élémentaires sont
\begin{equation}
q_+q_-=e^{-2x},
\qquad
\frac{q_-}{q_+}=e^{2y}.
\label{eq:q-product-ratio}
\end{equation}
Le produit conserve le mode commun ; le quotient conserve la séparation
orientée.

Appliqué au couple angulaire,
\[
x=\frac{\pi A}{180},\qquad y=\frac{\pi C^*}{180}.
\]
C'est pourquoi \(A\) est la coordonnée commune de la publication et \(C^*\) son
contraste orienté. Ce n'est qu'après l'application d'un lecteur physique qu'ils sont reconnus
comme pôles électriques et magnétiques.

 \subsection{Séparation typée des polarités constitutives et mécaniques}
% Fragmento público generado desde fuentes teoremáticas íntegros.
% Residencia: 52.7.2.
% Fuente: /Users/ruben/Documents/HMT2/EL_CIERRE_HOLOGRAFICO_DEL_INFINITO_SUCESOR_102_CAPITULOS_2026-08-28_EN_TRABAJO/incoming/radion_monografia_20260827/manuscrito/sections/08_complejo_electromagnetismo.tex:59-158
\subsubsection{Opérateur constitutif sur les canaux \(90/120\)}

\paragraph{Décomposition en composantes commune et orientée}

On définit
\begin{align}
\mathcal E
&=\log\bigl[(1-q_+^{90})(1-q_-^{90})\bigr]
-\log\bigl[(1-q_+^{120})(1-q_-^{120})\bigr],
\label{eq:Ecommon}\\
\mathcal B
&=\log\frac{1-q_-^{90}}{1-q_+^{90}}
-\log\frac{1-q_-^{120}}{1-q_+^{120}}.
\label{eq:Borient}
\end{align}
Sous \(C^*\mapsto-C^*\), \(q_+\leftrightarrow q_-\) s'échangent. Par
conséquent, \(\mathcal E\) est pair et \(\mathcal B\) est impair.

Les lois constitutives normalisées sont
\begin{equation}
\widehat\mu=e^{\mathcal E+\mathcal B},
\qquad
\widehat\varepsilon=e^{\mathcal E-\mathcal B},
\qquad
\widehat Z=e^{\mathcal B},
\qquad
\widehat c=e^{-\mathcal E}.
\label{eq:constitutives}
\end{equation}
On vérifie exactement
\begin{equation}
\widehat Z=\sqrt{\frac{\widehat\mu}{\widehat\varepsilon}},
\qquad
\widehat c=(\widehat\mu\widehat\varepsilon)^{-1/2}.
\label{eq:constitutive-identities}
\end{equation}
La conjugaison produit
\begin{equation}
C^*\mapsto-C^*:\qquad
\widehat\mu\leftrightarrow\widehat\varepsilon,
\quad
\widehat Z\leftrightarrow\widehat Z^{-1},
\quad
\widehat c\mapsto\widehat c.
\label{eq:EM-conjugation}
\end{equation}

\begin{theorem}[Polarité constitutive]
Le lecteur \(90/120\) transforme le mode commun et le contraste orienté du
couple \((A,C^*)\) en deux lois constitutives conjuguées. L'inversion de feuillet
échange perméabilité et permittivité, inverse l'impédance et conserve la
vitesse normalisée.
\end{theorem}

Telle est la formulation précise de l'intuition \enquote{partie électrique et
partie magnétique}. On n'écrit ni \(A=E\) ni \(C^*=B\). On écrit l'application
\[
(A,C^*)\to(x,y)\to(q_+,q_-)
\to(\mathcal E,\mathcal B)
\to(\widehat\varepsilon,\widehat\mu,\widehat Z,\widehat c).
\]

\paragraph{Réalisation élastique--rotationnelle de la polarité constitutive}

La même information peut être publiée dans le régime complexe au moyen de
\begin{equation}
Z_{\mathrm{em}}=e^{\mathcal E/2}
\left(\cos\frac{\mathcal B}{2}\,I
-\sin\frac{\mathcal B}{2}\,J\right).
\label{eq:Zem}
\end{equation}
Le facteur \(e^{\mathcal E/2}\) est une dilatation élastique ; le facteur généré
par \(J\) est une rotation. L'écriture complexe n'ajoute pas une dimension
imaginaire mystérieuse à un monde réel préalablement donné : elle enregistre deux modes
opératoires que le TRIT a déjà séparés.

\paragraph{Transduction de l'orientation en charge effective}

Le feuillet \(\sigma\) du radion fixe une orientation avant d'introduire une
unité de charge. La mécanique dimensionnelle dispose ensuite le transducteur
\[
e_{\mathrm{int}}=\sqrt{\frac{2\alphah\hfull}{Z_{\mathrm{int}}}},
\qquad
Q=n_Qe_{\mathrm{int}}.
\]
La chaîne causale est
\[
\text{feuillet}\to\text{orientation}\to\text{constitutive}
\to\text{transducteur dimensionnel}\to\text{charge physique}.
\]
Appeler l'unité \emph{radion} n'identifie pas prématurément l'orientation
au coulomb ; cela rend visible le lieu où le signe de charge peut résider.

\begin{narrativebox}
L'électricité et le magnétisme ne sont pas collés au cercle depuis l'extérieur. Ils émergent
lorsque le même état est lu à travers le produit et le quotient : le produit
retient le mode commun, le quotient retient l'orientation. L'oscillateur LC
réalise l'échange ; le lecteur \(90/120\) publie les lois constitutives ; la
carte dimensionnelle attribue des unités.
\end{narrativebox}
 
\section{Incidence des secteurs \(90/120\) : registre, spectre, noyau de Barbero--Immirzi et tours}
\subsection{Distinction typée des réalisations d’incidence du secteur \(90/120\)}
% Fragmento público generado desde fuentes teoremáticas íntegros.
% Residencia: 52.8.1.
% Fuente: /Users/ruben/Documents/HMT2/EL_CIERRE_HOLOGRAFICO_DEL_INFINITO_SUCESOR_102_CAPITULOS_2026-08-28_EN_TRABAJO/incoming/radion_monografia_20260827/manuscrito/sections/09_sector_90120.tex:176-244
\subsubsection{Hiérarchie harmonique globale et semi-groupe \(\langle90,120\rangle\)}

Pour éviter une collision de signes, nous fixons
\[
R_{\mathrm{tow}}:=-R_{\mathrm{ch}}.
\]
Alors
\begin{equation}
T=e^{-xI+yR_{\mathrm{tow}}}
=q_-P_+^{\mathrm{tow}}+q_+P_-^{\mathrm{tow}}.
\label{eq:tower-operator}
\end{equation}
Les tours paires et impaires sont
\begin{align}
c_n&=\operatorname{Tr}(T^n)
=q_-^n+q_+^n
=2e^{-nx}\cosh(ny),
\label{eq:cn}\\
s_n&=\operatorname{Tr}(R_{\mathrm{tow}}T^n)
=q_-^n-q_+^n
=2e^{-nx}\sinh(ny).
\label{eq:sn}
\end{align}
Sous \(C^*\mapsto-C^*\), \(c_n\) est pair et \(s_n\) est impair. On a
\begin{equation}
c_n^2-s_n^2=4e^{-2nx},
\label{eq:tower-invariant}
\end{equation}
et les deux suites satisfont la récurrence d'ordre deux déterminée par les
valeurs propres \(q_\pm\).

Le semi-groupe
\begin{equation}
\langle90,120\rangle=30\langle3,4\rangle
\label{eq:semigroup}
\end{equation}
est la hiérarchie harmonique d'un unique opérateur bicouche. Ce n'est pas une liste de huit
corrections latérales. Sa complétude signifie que tout rapport positif
admet un développement convergent ; elle ne signifie pas qu'un développement obtenu à
partir de la cible soit un sélecteur physique prospectif.

\subsubsection{Classification typée et séparation des réalisations \(90/120\)}

\begin{longtable}{>{\bfseries}p{.19\textwidth} p{.24\textwidth} p{.45\textwidth}}
\toprule
Objet & Opération & Ce qu'il conserve / pourquoi ce n'est pas un autre objet\\
\midrule
\endfirsthead
\toprule
Objet & Opération & Ce qu'il conserve / pourquoi ce n'est pas un autre objet\\
\midrule
\endhead
Incidence \(N_6,N_8\) & dénombrement fini sur les phases actives & cardinalité et orientation ; produit \(90,120\), non une série ;\\
\(-1/12\) & différence normalisée \eqref{eq:minus-one-twelve} & scalaire d'incidence ; non le poids de Barbero ;\\
\(E_{\mathrm{dod}}\) & extracteur harmonique de la roue & phase dodécaphasique ; non la retenue réelle ;\\
\(h=C^*/6\) & reconstruction de six paires & ouverture par paire ; non une normalisation métrologique ;\\
\(F_{\mathrm{spec}}\) & queue logarithmique \(j\ge2\) & publication spectrale complète ; diffère de \(\epsR^\circ\) ;\\
\(\epsone\) & soustraction APP entre \(S_T\) et \(S_E\) & raccord réel de deux relèvements ; non une série de Lambert ;\\
\(K_{6\to8}\) & \(12\Delta S_{90}-\Delta S_{120}\) & pôle \(1/24\) et saut marqué ; sortie de Barbero ;\\
\(12:-1\) & solution diophantienne de \eqref{eq:pole-condition} & poids du noyau ; non \(-1/12\) ;\\
\(\langle90,120\rangle\) & semi-groupe des puissances de \(T\) & hiérarchie harmonique globale ; non un résidu par espèce.\\
\bottomrule
\end{longtable}

\begin{thesisbox}
Tous ces objets descendent du même état APP--TRIT--TPK. Cette généalogie
commune explique pourquoi ils partagent des nombres et des symétries. Leur typage explique pourquoi
ils ne peuvent pas se remplacer mutuellement.
\end{thesisbox}
 \subsection{Conservation de l'incidence et de l'orientation entre le noyau tritique et ses réalisations ultérieures}
% Fragmento público generado desde fuentes teoremáticas íntegros.
% Residencia: 52.8.2.
% Fuente: /Users/ruben/Documents/HMT2/EL_CIERRE_HOLOGRAFICO_DEL_INFINITO_SUCESOR_102_CAPITULOS_2026-08-28_EN_TRABAJO/incoming/radion_monografia_20260827/manuscrito/sections/09_sector_90120.tex:1-30
\noindent La hiérarchie typée sépare incidence, registre, évaluation spectrale, fonctionnelle de Barbero et tours harmoniques.\par\medskip

\subsubsection{Incidence tritique, registre dodécaphasique et lecture sexadique}

\paragraph{Dérivation combinatoire primaire de \(90\) et \(120\)}

Les dénombrements
\[
90=6\binom62,
\qquad
120=8\binom62
\]
naissent du sélecteur de phase et de ses incidences actives. L'orientation
normalisée satisfait
\begin{equation}
\frac{30}{120}-\frac{30}{90}=-\frac1{12}.
\label{eq:minus-one-twelve}
\end{equation}
Ce scalaire est un résultat d'incidence. Ce n'est pas le poids \(12:-1\) du
noyau de Barbero et ce n'est pas une retenue dodécaphasique.

L'extracteur dodécaphasique peut être publié au moyen d'une combinaison
d'harmoniques telle que
\[
E_{\mathrm{dod}}(\phi)=12\cos(3\phi)-\cos(4\phi).
\]
Son domaine est la phase de la roue ; sa sortie est un sélecteur harmonique. Bien que
les nombres \(12\) et \(-1\) apparaissent, il n'opère pas sur des séries de Lambert et ne
produit pas \(\epsR\).

 \subsection{Unicité diophantienne relative du poids \(12:-1\) sous la condition de pôle \(1/24\) et le contre-terme entier minimal}
% Fragmento público generado desde fuentes teoremáticas íntegros.
% Residencia: 52.8.3.
% Fuente: /Users/ruben/Documents/HMT2/EL_CIERRE_HOLOGRAFICO_DEL_INFINITO_SUCESOR_102_CAPITULOS_2026-08-28_EN_TRABAJO/incoming/radion_monografia_20260827/manuscrito/sections/09_sector_90120.tex:115-175
\subsubsection{Fonctionnelle hexade--octade de Barbero--Immirzi}

\paragraph{Séries orientées et saut d'incidence}

On définit
\begin{equation}
S_{90}(q)=\frac{q^{90}}{1-q^{270}},
\qquad
S_{120}(q)=\frac{q^{120}}{1-q^{360}},
\label{eq:barbero-series}
\end{equation}
et
\[
\Delta S_m=S_m(q_-)-S_m(q_+).
\]
La fonctionnelle d'inclusion marquée est
\begin{equation}
K_{6\to8}=12\Delta S_{90}-\Delta S_{120}.
\label{eq:barbero-núcleo}
\end{equation}

\paragraph{Condition diophantienne pour le poids \(12:-1\)}

Considérons \(aS_{90}-bS_{120}\). La condition de pôle déclarée est
\begin{equation}
\frac{a}{270}-\frac{b}{360}=\frac1{24}.
\label{eq:pole-condition}
\end{equation}
Multiplier par \(1080\) donne
\[
4a-3b=45.
\]
Si le contre-terme est l’entier minimal, \(|b|=1\). Pour \(b=1\), \(a=12\) ; pour \(b=-1\), \(a=21/2\) n’est pas entier. Par conséquent,
\begin{equation}
\boxed{a:b=12:1,\qquad aS_{90}-bS_{120}=12S_{90}-S_{120}.}
\label{eq:12minus1}
\end{equation}

\begin{proposition}[Unicité diophantienne relative]
Le poids signé \(12:-1\) est unique parmi les coefficients entiers lorsque sont fixées la condition de pôle \(1/24\), la positivité du canal principal et le contre-terme minimal \(|b|=1\).
\end{proposition}

L’incidence hexade--octade fournit les canaux \(90/120\). Elle n’impose pas à elle seule la série de Lambert ni la condition \eqref{eq:pole-condition}. La force de \(12:-1\) est relative à ces prémisses analytiques explicites.

La chaîne de Barbero est ultérieure :
\[
(A,C^*)\to(x,y)\to q_\pm
\to K_{6\to8}\to\gamma_{\mathrm{BI}},
\]
avec l'évaluation
\[
\gamma_{\mathrm{BI}}=0.274007672694459\ldots.
\]
Barbero--Immirzi ne génère pas rétroactivement \(C^*\), \(\Delta_4\) ni
\(\epsR\).

 \subsection{Extraction de la composante finie ultérieure au système dodécaphasique orienté \(R_{12}\) : \(C^*\mapsto C^*/6\), reste \(j\ge2\) et reconstruction de six sections}
% Fragmento público generado desde fuentes teoremáticas íntegros.
% Residencia: 52.8.4.
% Fuente: /Users/ruben/Documents/HMT2/EL_CIERRE_HOLOGRAFICO_DEL_INFINITO_SUCESOR_102_CAPITULOS_2026-08-28_EN_TRABAJO/incoming/radion_monografia_20260827/manuscrito/sections/09_sector_90120.tex:31-54
\paragraph{Extraction de la composante \(C^*/6\) du registre dodécaphasique}

Les douze secteurs sont regroupés en six paires opposées :
\begin{equation}
p_i=e_i+e_{i+6},\qquad
a_i=e_i-e_{i+6},\qquad i=1,\ldots,6.
\label{eq:sexadic-pairs}
\end{equation}
La reconstruction a pour profil \(1+5+6=12\). L'ouverture par paire est
\begin{equation}
h=\frac{C^*}{6}.
\label{eq:h-C6}
\end{equation}
Ce n'est ni une normalisation extérieure ni le sixième d'une cible : c'est l'indice
interne de la reconstruction sexadique du registre \(\Rstate\).

On définit
\begin{equation}
q_{h,-}=\exp\left[-\frac\pi{180}(A-h)\right],
\qquad
q_{h,+}=\exp\left[-\frac\pi{180}(A+h)\right].
\label{eq:q-h}
\end{equation}

 \subsection{Décomposition quantitative des réalisations de retenue et spectrale : \(F_{\mathrm{spec}}=\varepsilon_R^\circ+\chi_R\)}
% Fragmento público generado desde fuentes teoremáticas íntegros.
% Residencia: 52.8.5.
% Fuente: /Users/ruben/Documents/HMT2/EL_CIERRE_HOLOGRAFICO_DEL_INFINITO_SUCESOR_102_CAPITULOS_2026-08-28_EN_TRABAJO/incoming/radion_monografia_20260827/manuscrito/sections/09_sector_90120.tex:55-114
\subsubsection{Publication spectrale de la queue d'ordres \(j\ge2\)}

\paragraph{Définition de la fonctionnelle spectrale}

Pour \(m\in\{90,120\}\), soit
\begin{equation}
L_m^{(2)}(q)=\sum_{j\ge2}\frac{q^{mj}}{j}
=-\log(1-q^m)-q^m.
\label{eq:L2}
\end{equation}
La lecture d'une paire est
\begin{equation}
T_h^{(2)}=\frac{180}{\pi}
\Bigl[
L_{90}^{(2)}(q_{h,-})-L_{90}^{(2)}(q_{h,+})
-L_{120}^{(2)}(q_{h,-})+L_{120}^{(2)}(q_{h,+})
\Bigr].
\label{eq:T-h}
\end{equation}
La reconstruction des six paires donne
\begin{equation}
F_{\mathrm{spec}}=6T_h^{(2)}
=5.1499235153094574410\ldots\times10^{-8}\,{}^\circ.
\label{eq:Fspec}
\end{equation}

La retenue du radion est
\[
F_{\mathrm{acarreo}}=\epsR^\circ
=5.1383016758575282072\ldots\times10^{-8}\,{}^\circ.
\]
Ils ne coïncident pas :
\begin{equation}
\chi_R:=F_{\mathrm{spec}}-F_{\mathrm{acarreo}}
=1.16218394519292338\ldots\times10^{-10}\,{}^\circ.
\label{eq:chiR}
\end{equation}

L’égalité
\[
F_{\mathrm{acarreo}}=F_{\mathrm{spec}}-\chi_R
\]
est un changement de carte valable une fois les deux sorties construites. Il ne doit pas
être présenté comme génération indépendante de \(\epsR^\circ\) si \(\chi_R\) est
défini par la soustraction elle-même. La queue est clôturée et significative ; ce qui est faux,
c'est de l'identifier directement à la retenue.

\begin{falsifier}
Après conversion de la queue exprimée en degrés dans l'unité angulaire interne,
\[
\frac{\pi}{180^\circ}F_{\mathrm{spec}}-\epsone
=2.0283936357433839\ldots\times10^{-12}\ne0.
\]
Par conséquent, ni le log complet ni la queue \(j\ge2\), même après application de la
carte d'unités, ne sont \(\epsone\). Une
seconde factorisation spectrale autonome nécessiterait de générer \(\chi_R\) sans
utiliser la différence comme définition ; la clôture APP--TPK du radion ne dépend pas
de cette promotion.
\end{falsifier}

 
\section{Holonomie orientée et géométrie physique : réalisations de Minkowski, Hodge, Lorentz et Cartan--Holst}
% Fragmento público generado desde fuentes teoremáticas íntegros.
% Residencia: 52.9.
% Fuente: /Users/ruben/Documents/HMT2/EL_CIERRE_HOLOGRAFICO_DEL_INFINITO_SUCESOR_102_CAPITULOS_2026-08-28_EN_TRABAJO/incoming/radion_monografia_20260827/manuscrito/sections/10_minkowski_hodge_lorentz.tex:1-55
\noindent L'holonomie orientée admet des réalisations causales et géométriques ultérieures dans Minkowski, Hodge, Lorentz et Cartan--Holst.\par\medskip

\paragraph{Construction de la métrique lorentzienne à partir du bloc APP}

\paragraph{Matrice APP de paramètres \(7\) et \(2\)}

La matrice centrale
\begin{equation}
B_c=\begin{pmatrix}7&2\\2&7\end{pmatrix}
=9P_++5P_-
\label{eq:Bc}
\end{equation}
est normalisée par son déterminant :
\begin{equation}
M_c=\frac{B_c}{\sqrt{45}}
=\exp\left(\frac12\log\frac95\,R\right)
\in SO^+(1,1).
\label{eq:Mc}
\end{equation}
La rapidité interne est
\[
\eta_c=\frac12\log\frac95.
\]
La publication projective associée satisfait
\begin{equation}
P_B^n=P_++\left(\frac59\right)^nP_-.
\label{eq:PB-contract}
\end{equation}
Le même rapport \(5/9\) qui apparaît comme contraction spectrale est reconnu ici
comme séparation d'un mode stable et d'un autre contractant. L'identité est
opératoire ; elle ne transforme pas toutes les occurrences de \(5/9\) en la même application.

\paragraph{Incidence \(6\to8\) et quotient de signature causale}

Soit \(M\) l'incidence typée entre six faces et huit sommets. La relation
spectrale récupérée est
\begin{equation}
MM^{\mathsf T}=12P_u+4P_-,
\label{eq:MMt}
\end{equation}
où \(P_u\) projette sur le mode uniforme et \(P_-\) sur le sous-espace
orienté. Le quotient survivant de dimension quatre se décompose comme
\begin{equation}
Q_4\simeq\R u\oplus E_-.
\label{eq:radion-Q4}
\end{equation}
En déclarant négatif le mode uniforme et positifs les trois modes orientés,
on obtient
\begin{equation}
B_{\mathrm{caus}}=\operatorname{diag}(-1,1,1,1).
\label{eq:minkowski}
\end{equation}
L'écran de Minkowski ne sélectionne pas l'état HMT. C'est la réalisation
causale ultérieure du quotient que l'incidence a déjà construit.

 \subsection{Dualité de Hodge sur les \(2\)-formes lorentziennes}
% Fragmento público generado desde fuentes teoremáticas íntegros.
% Residencia: 52.9.1.
% Fuente: /Users/ruben/Documents/HMT2/EL_CIERRE_HOLOGRAFICO_DEL_INFINITO_SUCESOR_102_CAPITULOS_2026-08-28_EN_TRABAJO/incoming/radion_monografia_20260827/manuscrito/sections/10_minkowski_hodge_lorentz.tex:56-100
\subsubsection{Dualité de Hodge et opérateur de quart de tour}

La signature \((-+++ )\) et une orientation étant fixées, l'étoile de Hodge sur
les deux-formes satisfait
\begin{equation}
\star:\Lambda^2Q_4^*\to\Lambda^2Q_4^*,
\qquad \star^2=-I.
\label{eq:hodge-star}
\end{equation}
De cette manière, le quart de tour qui, dans le plan de phase, s'exprimait au moyen de
\(J^2=-I\) réapparaît dans l'espace des champs comme structure complexe sur
\(\Lambda^2\). Les domaines sont différents, mais la généalogie d'orientation
est conservée par le foncteur.

Soit \(A_{\mathrm{em}}\) un potentiel et
\begin{equation}
F=\diff A_{\mathrm{em}}.
\label{eq:F-dA}
\end{equation}
La décomposition de \(F\) par rapport à un observateur temporel produit des champs
électrique et magnétique ; \(\star F\) échange leurs rôles avec le signe fixé
par l'orientation et la signature. La polarité constitutive de
\eqref{eq:EM-conjugation} acquiert ainsi une réalisation différentielle.

\paragraph{Chaîne de transduction électromagnétique et causale}

Le transport qui doit rester visible est
\begin{equation}
\Gamma_9
\longrightarrow\text{hélice de mémoire}
\longrightarrow(Q_4,B_{\mathrm{caus}},\star)
\longrightarrow A_{\mathrm{em}}
\longrightarrow F=\diff A_{\mathrm{em}}
\longrightarrow
m\nabla_u u=q(\iota_uF)^\sharp.
\label{eq:lorentz-chain}
\end{equation}
Chaque flèche ajoute de la structure : l'holonomie conserve le retour ; le quotient fixe
la causalité ; Hodge fixe la dualité ; le potentiel fixe le champ ; le couplage
avec la charge produit la dynamique.

L'équation finale est la force de Lorentz sous forme géométrique. L'hélice TPK
est son antécédent de mémoire, non une orbite physique déjà mue par un champ. La
dynamique n'apparaît qu'à l'introduction de \(F\), \(q\) et de la connexion \(\nabla\).

 \subsection{Séparation entre mémoire torsionnelle interne et réalisation géométrique de Cartan--Holst}
% Fragmento público generado desde fuentes teoremáticas íntegros.
% Residencia: 52.9.2.
% Fuente: /Users/ruben/Documents/HMT2/EL_CIERRE_HOLOGRAFICO_DEL_INFINITO_SUCESOR_102_CAPITULOS_2026-08-28_EN_TRABAJO/incoming/radion_monografia_20260827/manuscrito/sections/10_minkowski_hodge_lorentz.tex:101-186
\subsubsection{Réalisation lorentzienne des trajectoires hélicoïdales}

Dans un champ magnétique uniforme, une particule de charge \(q\), de masse \(m\) et de
facteur relativiste \(\gamma\) possède
\begin{equation}
\omega_c=\frac{|q|B}{\gamma m},
\qquad
r_L=\frac{p_\perp}{|q|B},
\qquad
\Delta z=\frac{2\pi p_\parallel}{|q|B}.
\label{eq:lorentz-helix}
\end{equation}
La phase cyclotronique se ferme tandis que le mouvement parallèle avance : l'orbite
est hélicoïdale. Le parallélisme avec \eqref{eq:H9-helix} est structurel et peut
être formalisé par un transducteur qui préserve phase, orientation et pas. On
n'identifie pas le quotient \(q_9\) à la coordonnée spatiale \(z\) ; on reconnaît
la même architecture de retour plus avance.

La force de Lorentz ne se dérive pas de l'image d'une spirale par ressemblance.
Elle se dérive en complétant \eqref{eq:lorentz-chain}. L'image aide à comprendre
pourquoi phase circulaire et mémoire longitudinale sont compatibles ; la forme de
champ détermine la dynamique.

\subsubsection{Confluence entre action, aire, information et gravité}

\paragraph{Réalisation Cartan--Holst de l'holonomie}

La promotion gravitationnelle s'exprime au moyen d'un morphisme d'holonomies
\begin{equation}
\operatorname{Hol}_{\mathcal A}
\bigl(\mathfrak R_{\mathrm{grav}}(\gamma)\bigr)
=\rho_C\bigl(\operatorname{Hol}_{\TPK}(\gamma)\bigr).
\label{eq:holonomy-gravity}
\end{equation}
Le membre de droite contient la mémoire de chemin ; \(\rho_C\) la réalise dans une
connexion gravitationnelle. La torsion de Cartan
\[
T^I=D_\omega e^I
\]
apparaît dans le codomaine. La torsion globale \(\Delta_4\) est une coordonnée
en amont qui peut alimenter l'application, mais n'est pas renommée \(T^I\) sans
\eqref{eq:holonomy-gravity}.

\paragraph{Transduction de l'aire d'action en aire informationnelle}

Le radion fixe un quantum d'action par phase. Le lecteur hexade--octade de
Barbero utilise ultérieurement les canaux \(q_\pm\), l'incidence \(90/120\) et
la fonctionnelle \eqref{eq:barbero-núcleo}. La direction causale est
\[
(A,C^*,\hret)\to q_\pm\to K_{6\to8}
\to\gamma_{\mathrm{BI}}\to\text{aire/information}.
\]
L'alphabet tritique apporte l'unité d'entropie \(\log3\) et une coupure de
\(81\) canaux apporte \(81\log3\). Cette entropie ne se confond ni avec
\(\log2\), ni avec \(\log\varphi\), ni avec une entropie microcanonique, ni avec une
entropie de von Neumann sans spécifier l'état.

La synthèse physique est affirmative mais typée : le radion apporte l'unité
d'action orientée qu'une holonomie peut transporter ; Barbero publie la
normalisation d'aire ; l'écran causal et Hodge permettent de réaliser des champs ;
Cartan--Holst reçoit l'holonomie en géométrie gravitationnelle.

\begin{statusbox}
L'équation scalaire du radion n'est pas à elle seule une équation d'Einstein. Sa
contribution à la gravité consiste à être une section exacte de l'architecture
d'action, d'orientation et d'holonomie que la réalisation gravitationnelle consomme.
\end{statusbox}

\subsubsection{Transformation de Wick comme changement de lecteur analytique}

Dans la formulation conventionnelle, la rotation de Wick relie une variable
temporelle lorentzienne à une coordonnée euclidienne par un prolongement
complexe. Dans HMT, la racine \(J^2=-I\), l'opérateur involutif \(R^2=I\) et le seuil
\(N^2=0\) existent déjà auparavant. Le changement de lecteur
\[
J\longleftrightarrow R
\]
transporte un couple commun/orienté entre publication elliptique et
publication hyperbolique. Le prolongement complexe reconnaît ensuite cette
opération en coordonnées analytiques.

Cela explique pourquoi cercle, ellipse d'action et intérieur hyperbolique peuvent
appartenir au même état sans être la même métrique. Le cercle utilise \(J\) pour
la phase ; la forme de l'ellipse utilise une compression symplectique générée par \(R\) ; Klein mesure
l'intérieur par le birapport ; l'écran de Minkowski réalise la causalité
dans un autre codomaine.
 
\section{Ellipse, réseau et phase de chemin : caractère modulaire, orbites et propagation}
\subsection{Réseau rectangulaire de périodes, involution \(\tau\mapsto-1/\tau\) et invariance de \(j\)}
% Fragmento público generado desde fuentes teoremáticas íntegros.
% Residencia: 52.10.1.
% Fuente: /Users/ruben/Documents/HMT2/EL_CIERRE_HOLOGRAFICO_DEL_INFINITO_SUCESOR_102_CAPITULOS_2026-08-28_EN_TRABAJO/incoming/radion_monografia_20260827/manuscrito/sections/11_modularidad_orbitas.tex:1-84
\noindent Les caractères modulaires, orbitaux et de chemin conservent des invariants distincts d'une même forme orientée.\par\medskip

\subsubsection{Tore complexe associé à l'ellipse d'action}

\paragraph{Réseau rectangulaire de périodes}

L'ellipse à elle seule n'est pas un tore. Nous déclarons le réseau
\begin{equation}
\Lambda_\eta=a_S\Z+i,b_S\Z
\label{eq:lattice}
\end{equation}
et le quotient \(\C/\Lambda_\eta\). Son module est
\begin{equation}
\tau_\eta=i\frac{b_S}{a_S}=ie^{-\eta}
=i\,0.741048144159375160795483663369300\ldots.
\label{eq:tau-eta}
\end{equation}
L'inversion bidirectionnelle \(\eta\mapsto-\eta\) échange les axes :
\begin{equation}
\tau_{-\eta}=ie^\eta=-\frac1{\tau_\eta}.
\label{eq:modular-S}
\end{equation}
C'est exactement la transformation modulaire \(S:\tau\mapsto-1/\tau\).

\paragraph{Invariance de la fonction modulaire \(j\)}

Par invariance modulaire,
\begin{equation}
\boxed{j(\tau_{-\eta})=j(\tau_\eta).}
\label{eq:j-invariance}
\end{equation}
Dans la normalisation \(j(i)=1728\),
\[
j(\tau_\eta)
=5597.43843932408729830097089254768\ldots.
\]
La fonction moonshine normalisée \(J=j-744\) donne
\[
J(\tau_\eta)
=4853.43843932408729830097089254768\ldots.
\]
Pour le bloc APP basal,
\[
\tau_0=i\frac{\sqrt5}{3},
\qquad
j(\tau_0)=5369.48832024674306021916882626352\ldots.
\]

\begin{proposition}[Ce que conserve la classe modulaire]
L'involution de feuillet inverse \(\eta\), mais \(j\) conserve la classe du
tore. Les foyers et l'étiquetage des axes retiennent l'orientation que \(j\)
oublie.
\end{proposition}

Ce résultat relie de manière exacte la double voie à la modularité une
fois le quotient déclaré. Il n'autorise pas à faire agir directement le
Monstre sur la charge, les chiffres ou les foyers. Le couloir exceptionnel
\[
\mathrm{APP}\to\mathrm{Paley}\to\mathrm{Witt}\to\mathrm{Golay}
\to A_2^{12}\to\mathrm{Leech}\to V^\natural
\to\mathrm{Monster}\to\mathrm{Moonshine}
\]
est une autre projection du même état holonomique intégral. L’identification d’une action concrète sur le radion exige un entrelaceur préservant ses coordonnées pertinentes.

\paragraph{Chaîne d'invariants du caractère de forme}

En réunissant les réalisations déclarées,
\begin{equation}
\boxed{
r_\eta=e^{-\eta}
=\sqrt{\frac{A-C^*}{A+C^*}}
=\frac{b_1}{a_1}
=\frac{b_S}{a_S}
=\frac{\Delta P}{\Delta Q}
=\frac{Z_\eta}{Z_*}
=\Im\tau_\eta.}
\label{eq:shape-master}
\end{equation}
Un seul rapport parcourt géométrie spectrale, action, covariance, impédance et
module. Chaque égalité repose sur une application explicite ; aucune ne dépend d'une
ressemblance décimale.

 \subsection{Compatibilité entre le caractère modulaire et la propagation sur des chemins orientés}
% Fragmento público generado desde fuentes teoremáticas íntegros.
% Residencia: 52.10.2.
% Fuente: /Users/ruben/Documents/HMT2/EL_CIERRE_HOLOGRAFICO_DEL_INFINITO_SUCESOR_102_CAPITULOS_2026-08-28_EN_TRABAJO/incoming/radion_monografia_20260827/manuscrito/sections/11_modularidad_orbitas.tex:85-187
\subsubsection{Conservation surfacique dans les orbites de Kepler et de Sommerfeld}

\paragraph{Équivalence surfacique des réalisations orbitale et oscillatoire}

L'ellipse de Kepler vit dans l'espace de configuration et répond à un potentiel
central. L'ellipse du radion vit dans l'espace des phases et répond à une forme
d'action. Ce sont des objets différents. Elles partagent une loi surfacique parce que les deux
réalisations possèdent une deux-forme conservée.

Pour une orbite képlérienne,
\begin{equation}
r(\phi)=\frac{p}{1+e\cos(\phi-\phi_0)},
\qquad
\frac{\diff A}{\diff t}=\frac{\ell}{2m}.
\label{eq:kepler}
\end{equation}
La troisième loi prend la forme
\[
T^2=\frac{4\pi^2\mu}{K}a^3.
\]
Pour l'oscillateur en phase,
\begin{equation}
J=\oint p\,\diff q=\frac{2\pi E}{\omega}.
\label{eq:osc-action}
\end{equation}
Si la cellule élémentaire est \(J=\hfull\), alors
\begin{equation}
E=\hret\omega,\qquad \frac{J}{2\pi}=\hret.
\label{eq:E-hbarw}
\end{equation}
Telle est la convergence nette entre action de Planck, fréquence angulaire et
radion.

Sommerfeld quantifie les intégrales d’action sur les cycles. La lecture HMT ajoute la mémoire du cycle, du feuillet et du retour qui survivent. Sur un ruban de Möbius, une boucle visible peut porter \(\hfull/2\) et deux boucles rétablir \(\hfull\). Ce facteur topologique ne s’ajoute pas à l’indice de Maslov d’une quantification EBK ; ils relèvent de corrections fondées sur des sources théorématiques distinctes.

\paragraph{Holonomie orbitale et propagations avancée et retardée}

Une orbite qui accumule de l'holonomie satisfait une condition du type
\begin{equation}
\epsilon_\gamma\exp\left(\frac{iJ_\gamma}{\hret}\right)=1.
\label{eq:holonomy-action}
\end{equation}
Le signe \(\epsilon_\gamma\) conserve le feuillet ; \(J_\gamma\) conserve
l'action. L'avance ou le retard apsidal s'interprète comme découplage entre
la clôture angulaire visible et le retour complet. Le radion fournit
l'unité orientée avec laquelle cette différence est mesurée.

\subsubsection{Composition des amplitudes de Schrödinger et somme des chemins}

\paragraph{Quatre réalisations d'une même phase}

Une fois \(\hret\) fixé, différentes équations conventionnelles reconnaissent
des aspects du même transport :
\begin{enumerate}
\item Schrödinger publie la phase au moyen de
\(i\hret\partial_t\psi=H\psi\).
\item Pauli ajoute le double feuillet de spin et son couplage magnétique.
\item Dirac linéarise la relation énergie--quantité de mouvement et organise particule et
antiparticule dans une représentation causale.
\item La somme des chemins attribue à chaque histoire la phase
\(\exp(iS[\gamma]/\hret)\).
\end{enumerate}
Ces théories n’engendrent pas rétrospectivement \(\hret\). Elles reçoivent l’unité d’action et réalisent différents contenus de l’état holonomique intégral.

\paragraph{Fonctionnelle de phase sur l'espace des chemins}

Dans la formulation de Feynman,
\begin{equation}
\mathcal K(b,a)=\int_{\gamma:a\to b}
\exp\left(\frac{i}{\hret}S[\gamma]\right)\mathcal D\gamma.
\label{eq:path-integral}
\end{equation}
Un incrément d'action \(\hret\) change la phase d'un radian. Un incrément
\(\hfull\) la change d'un tour. Le théorème
\(\mathcal A(\theta)=\hret\theta\) transforme cette relation en géométrie
d'aire : la phase d'un chemin peut se lire comme action balayée en unités de
radion.

L'inversion \(C_M=-\operatorname{rev}\) sur les mots dodécaphasiques et
l'identité \eqref{eq:time-reversal} organisent des chemins conjugués. Dans une
réalisation d'absorbeur, on sépare
\[
G_{\mathrm{sym}}=\tfrac12(G_{\mathrm{ret}}+G_{\mathrm{adv}}),
\qquad
G_{\mathrm{odd}}=G_{\mathrm{ret}}-G_{\mathrm{adv}}.
\]
Le premier secteur est pair sous l'échange retardé--avancé ; le second est
impair. Cela reproduit la séparation entre énergie/forme paire et action
orientée impaire du radion lorsque le transducteur préserve la frontière et la forme
symplectique.

\begin{statusbox}
La théorie classique de l'absorbeur et l'interprétation de Feynman--Stückelberg des
antiparticules sont des réalisations historiques différentes. La première utilise une
interaction temporellement symétrique ; la seconde réorganise la propagation dans la
théorie relativiste. HMT apporte un antécédent commun de réversion, mais ne les
fusionne pas sans une application de propagateurs.
\end{statusbox}
 
\section{Radical nonadique et frontière décimale : conservation de la retenue du radion}
\subsection{La retenue du radion dans la complétion décimale}
% Fragmento público generado desde fuentes teoremáticas íntegros.
% Residencia: 52.11.1.
% Fuente: /Users/ruben/Documents/HMT2/EL_CIERRE_HOLOGRAFICO_DEL_INFINITO_SUCESOR_102_CAPITULOS_2026-08-28_EN_TRABAJO/incoming/radion_monografia_20260827/manuscrito/sections/12_aritmetica_borde.tex:1-159
\noindent La frontière nonadique conserve le défaut somme--produit, la vacance et la retenue de complétion.\par\medskip

\subsubsection{Radical \(3\! -\! 6\! -\! 9\) du bord}

\paragraph{Génération et nilpotence du radical nonadique}

Dans l'anneau \(\Z/9\Z\), l'idéal
\begin{equation}
\mathfrak m=(3)=\{0,3,6\}
\label{eq:radical-369}
\end{equation}
satisfait
\begin{equation}
\mathfrak m^2=(9)=0.
\label{eq:m-square-zero}
\end{equation}
C'est le secteur nilpotent du bord nonadique. L'orbite visible de ses trois
marques se publie comme
\[
369\longrightarrow693\longrightarrow936\longrightarrow369.
\]
Le retour n'implique pas que tous les états soient égaux : la permutation
visible se ferme tandis que les cocycles de chemin peuvent avancer.

La double convention
\[
9\equiv0\pmod9,
\qquad
\operatorname{dr}_9(9)=9
\]
exprime deux lecteurs : zéro résiduel algébrique et clôture visible de la racine
numérique. C'est pourquoi un neuf peut absorber une torsion de bord sans devenir
un zéro spécifique de la fonction zêta.

\paragraph{Défaut somme--produit \(459-405=54\)}

Les feuillets APP produisent des totaux complémentaires
\[
S_+=405,\qquad S_\times=459,\qquad S_\times-S_+=54.
\]
Le \(54\) est un défaut exact des tables, et réapparaît comme exposant de
contraction par l'identité \(9\cdot6=54\). Cependant, chaque occurrence
nécessite son application. La distance focale \(d_1=0.544545\ldots\) n'est pas égale à
\(0.54\), et l'analyse des invariants n'a pas produit de flèche qui
l'identifie au défaut \(54\).

La puissance
\[
9^9=387,420,489
\]
génère une échelle où
\begin{equation}
\operatorname{ord}_{10}\bigl(9^{9^9}\bigr)=369,693,099,
\qquad
\#\,\text{chiffres}=369,693,100.
\label{eq:369693}
\end{equation}
Le bloc \(369|693\) apparaît dans le couple ordre--longueur ; on n'affirme pas qu'il soit
le préfixe décimal de la puissance.

\subsubsection{Théorie arithmétique des vacances décimales}

\paragraph{Pente asymptotique de vacance}

La séparation entre la décennie et la base nonadique est
\begin{equation}
\delta_{\mathrm{vac}}=\log_{10}\frac{10}{9}.
\label{eq:delta-vac}
\end{equation}
Le compteur cumulé
\begin{equation}
V_9(n)=\left\lceil n\delta_{\mathrm{vac}}\right\rceil
\label{eq:vacancy-count}
\end{equation}
enregistre le nombre de positions de correction nécessaires à la publication décimale
d'une évolution nonadique. Ses sauts forment un mot binaire
\[
z_n=V_9(n)-V_9(n-1)\in\{0,1\}.
\]
Les lacunes se répartissent avec des longueurs \(21/22\), les retours avec des
longueurs \(6/7\), et le déterminant combiné est \(15\). Ces nombres
proviennent du lecteur de vacance, non du noyau de Barbero, même si le \(15\)
coïncide avec \(\binom62\).

\paragraph{Polarisation et courant arithmétique du réseau nonadique}

On définit
\begin{equation}
P_N=\sum_{n=1}^N(z_n-\delta_{\mathrm{vac}}),
\qquad
J_N=P_N-P_{N-1}=z_N-\delta_{\mathrm{vac}}.
\label{eq:arith-current}
\end{equation}
La séquence équilibrée satisfait \(|P_N|<1\). Le \enquote{courant}
\(J_N\) est ici un courant arithmétique : il mesure la création et l'absorption de
vacances par rapport à la pente moyenne. Seul un transducteur dimensionnel
ultérieur peut le publier comme densité de courant électromagnétique.

L'intuition physique est puissante parce que la même forme algébrique sépare une
moyenne paire et une fluctuation orientée. La rigueur exige de conserver le type :
\[
J_N\in\R\quad\text{sans dimension},
\qquad
j^\mu\in\text{droite physique de courant}.
\]

\subsubsection{Complétion \(729\to1000\) et conservation de la retenue}

\paragraph{Décomposition \(1000=729+243+27+1\)}

La complétion
\[
1000=729+243+27+1
\]
conserve l'unité au moyen de quatre échelles ternaires. La couronne totale est
\[
1000-729=271,
\]
tandis que le dernier entier avant la clôture satisfait
\[
999=729+270.
\]
La différence entre \(270\) et \(271\) distingue frontière occupée et unité
complétante. Le même soin évite d'identifier une vacance à un zéro ou une
retenue à un résidu réel.

\paragraph{Résidence de la retenue du radion dans la couronne décimale}

Les opérandes du radion commencent, en blocs de trois chiffres, par
\[
\Phi_T:\quad000|027|455|906|837|565|\cdots,
\]
\[
D_E:\quad000|027|455|010|034|743|\cdots.
\]
La soustraction APP donne
\[
\epsone:\quad000|000|000|896|802|822|\cdots.
\]
La petitesse du résultat ne s'obtient pas en tronquant les termes inférieurs.
Elle s'obtient parce que deux publications distinctes partagent un long préfixe et que
l'opérateur \(\operatorname{SubAcarreo}_{1000}\) conserve les emprunts nécessaires
pour les soustraire exactement.

Le mot \emph{mémoire} a ici deux niveaux :
\begin{enumerate}
\item \(q_{\mathrm{mem}}\) enregistre la profondeur des retours nonadiques ;
\item \(c_i\) enregistre les emprunts locaux de la soustraction décimale.
\end{enumerate}
Le premier organise l'hélice ; le second publie la décimale. \(\epsR\) utilise
les deux parce que ses opérandes proviennent de cylindres en profondeur et que sa soustraction
s'effectue avec retenue, mais il n'est égal à aucun des compteurs.

\begin{thesisbox}
Le bord n'est pas un lieu où la théorie perd de l'information. C'est le lieu où
l'information change de représentation : phase qui revient, mémoire qui
s'élève, cylindre qui se contracte, vacance qui s'enregistre et retenue qui
transporte la différence.
\end{thesisbox}
 \subsection{Compatibilité des sections directe et torsionnelle : caractérisation structurelle de \(\pi\)}
% Fragmento público generado desde fuentes teoremáticas íntegros.
% Residencia: 52.11.2.
% Fuente: /Users/ruben/Documents/HMT2/EL_CIERRE_HOLOGRAFICO_DEL_INFINITO_SUCESOR_102_CAPITULOS_2026-08-28_EN_TRABAJO/incoming/radion_monografia_20260827/manuscrito/sections/03_cierre_exacto.tex:286-309
\subsubsection{Caractérisation structurelle de \(\pi\) par compatibilité des sections}

Réordonner \eqref{eq:cierre-unitario} donne
\begin{equation}
\boxed{
2\pih=
\frac{1+\Phi_T-\epsone}
{\frac5{36}+\frac{25}{9}\alphah}.}
\label{eq:pi-radion}
\end{equation}
Cette égalité peut être appelée une nouvelle définition de \(\pi\) au sein de la
carte du radion si l'on entend \emph{définition structurelle} comme une
caractérisation exacte par compatibilité des sections. Ce n'est pas un second
générateur indépendant : \(\Delta_4\) et l'opérateur de retenue reçoivent déjà la
section coinductive de \(\pi\). La nouveauté réside dans le fait que le tour est
caractérisé par le raccord entre publication directe, torsion et mémoire.

\begin{narrativebox}
L'équation \eqref{eq:pi-radion} ne dit pas que nous ignorons d'abord \(\pi\) et
le devinons en fermant un angle. Elle dit que le \(\pi\) généré par la connexion
nonadique est exactement celui qui rend compatibles deux lectures internes du
même état. C'est une identité de reconnaissance intrinsèque, non un
étalonnage métrologique.
\end{narrativebox}
 
\section{Invariance évolutive de l’unité holographique à travers les dimensions}
% Conversión TeX fiel del cuerpo científico del delta narrativo.
% Fuente congelada: /Users/ruben/Documents/HMT2/CONTINUIDAD_EDITORIAL/RECONSTRUCCION_ARMONICA_MACROTRATADO_20260822/REVISION_CONJUNTA_PARTES_I_II/auditar_pdf_rigor_academico/docs/DELTA_NARRATIVO_UNIDAD_RADION_AUTOCONSERVACION_AUTOINERCIA_GRAVEDAD_PENTADICA_20260828.md:12-366.
\subsection{Séparation typée entre le régime de courbure \(\tau\in\{+1,0,-1\}\), la réponse radiale \(p\) et l'holonomie globale de spin}

La biographie du radion peut être représentée par une courbe relevée

\[
\gamma(\theta)
=
\bigl(r(\theta)\cos\theta,\,
r(\theta)\sin\theta,\,
m(\theta)\bigr),
\]

où les deux premières coordonnées appartiennent au plan euclidien \(\mathbb R^2\) et \(m(\theta)\) enregistre la mémoire ou la profondeur du relèvement. Cette équation permet de lire conjointement cercle, hélice et spirale.

Lorsque l'angle évolue et que le rayon et la mémoire restent constants, la publication plane est un cercle. Lorsque l'angle et la mémoire évoluent tandis que le rayon conserve sa valeur, la trajectoire complète est une hélice de rayon constant. Lorsque l'angle et le rayon évoluent sur une mémoire fixe, une spirale apparaît dans \(\mathbb R^2\). Lorsque les trois coordonnées évoluent, la trajectoire est une hélice radiale : une biographie dans laquelle phase, amplitude et mémoire se transforment conjointement.

Le caractère radial \(p\) classe la loi de réponse du rayon. Une famille utile est

\[
r_p(\theta)
=
r_0\bigl[1+p\lambda(\theta-\theta_0)\bigr]^{1/p},
\qquad p\neq 0,
\]

avec pour limite logarithmique

\[
r_0\exp\bigl(\lambda(\theta-\theta_0)\bigr)
\qquad\text{lorsque }p\to0.
\]

Dans cette paramétrisation, \(p=2\) publie le régime de Fermat, \(p=1\) le régime archimédien, \(p=0\) le logarithmique, \(p=-1\) l'hyperbolique et \(p=-2\) celui du lituus, avec les choix de domaine et d'orientation correspondants. Les valeurs fractionnaires décrivent des régimes intermédiaires. Ainsi, \(p=\tfrac12\), \(p=\tfrac32\) ou \(p=-\tfrac12\) expriment des réponses radiales qui interpolent entre les familles entières et élargissent la classification géométrique disponible.

Le caractère peut aussi changer lors du franchissement d'un seuil. L'écriture différentielle

\[
\frac{dx}{d\theta}
=
\beta(m)\,x^{1-p(m)},
\qquad x=\frac r{r_0},
\]

décrit une biographie par régimes : \(p(m)\) reste stable au sein d'une phase et change lorsque la mémoire atteint une frontière de publication. La trajectoire résultante relie des secteurs circulaires, hélicoïdaux et spiraux par une loi commune. La continuité de l'état et le registre du seuil font de ces secteurs les phases d'une même évolution.

Trois invariants occupent ici des niveaux distincts et coordonnés. Le signe de courbure

\[
\tau\in\{+1,0,-1\}
\]

classe les régimes elliptique, parabolique et hyperbolique. Le caractère \(p\) classe la réponse radiale. Le spin enregistre l'holonomie d'orientation accumulée par le relèvement. Courbure, réponse radiale et spin forment un système de lecture : la courbure détermine le type local d'espace, \(p\) détermine comment la trajectoire occupe cet espace et le spin conserve l'orientation acquise en le parcourant. Les courbes fermées et ouvertes apparaissent alors comme des biographies géométriques distinctes de la même unité orientée.

\subsection{Particule comme section persistante d'une orbite enrichie : phase, chemin, fibre, multisection et mémoire}

La mécanique dimensionnelle peut formuler la particule comme une section persistante de cette extension géométrique et dynamique. Une notation conceptuelle compacte est

\[
\mathfrak P
=
\operatorname{Sec}_{\mathrm{pers}}
\bigl[
\tau,p,\operatorname{Hol};
\text{feuillet},\text{fibre},\text{route},\text{mémoire}
\bigr].
\]

La particule conserve un chemin reconnaissable à travers les changements de phase, de courbure et de dimension. Son spin publie l'holonomie d'orientation ; sa charge publie la polarité de phase ; sa couleur publie une orientation interne de la fibre de jauge ; sa saveur distingue des familles de chemins ou de multisections ; sa masse exprime le caractère de transduction dimensionnelle par lequel l'action surfacique acquiert une lecture volumique. Chaque propriété est un lecteur de l'état persistant et toutes renvoient à une généalogie commune.

L'électron occupe le seuil basal de cette construction parce qu'il offre le premier lecteur stable dans lequel phase, charge, spin et masse sont simultanément publiés. La masse électronique directrice

\[
m_e^\beta=0.510998950690000808608\ldots\ \mathrm{MeV}
\]

fixe une référence de transduction ; les états de base antérieurs conservent leur fonction d’étapes généalogiques. À partir de l’électron, les quotients \(K_D/K_\Omega\), les tours \(90/120\) et le domaine \(81/56/13/324\) organisent l’expression des systèmes massifs. La loi de masse lit donc la manière dont chaque section persistante traverse l’espace holonomique et transforme la densité surfacique d’action en présence volumique.

Cette interprétation relie les spirales aux particules de manière précise. L'indice \(p\) apporte la réponse radiale ; \(\tau\) apporte le régime de courbure ; l'holonomie apporte le spin ; le chemin et la fibre apportent la saveur et la couleur ; la transduction aire–volume apporte le caractère massif. L'atlas des particules devient ainsi une classification des modes stables de traversée dimensionnelle.

\subsection{Transduction de la phase orientée en polarisation électrique, réponse magnétique et émission radiative de frontière}

L'électricité se présente comme transport dirigé de phase et d'action à travers une structure de vacances. Une vacance est une transition admissible du support : elle ouvre un canal, fixe un seuil et permet à l'actualisation d'avancer. La composante électrique publie le transport tangentiel le long du chemin. La composante magnétique publie la courbure transversale induite par ce transport. Leur perpendicularité exprime la relation géométrique entre avance et rotation.

Lorsque deux composantes transversales sinusoïdales maintiennent le déphasage approprié, le vecteur du champ tourne. La propagation longitudinale relève cette rotation en une hélice. La variation d'amplitude transforme sa projection frontale en une spirale. La lecture géométrique est alors entièrement coordonnée : l'angle porte la phase ; la vitesse angulaire porte la fréquence ; le rayon porte l'amplitude et la densité d'action ; l'orientation porte la polarisation et le sens ; la coordonnée verticale porte la propagation et la mémoire ; \(p\) porte le régime radial ; le seuil porte l'émission, l'absorption ou le changement d'état.

L'électromagnétisme apparaît comme une dynamique de phase–action–mémoire avec des publications complémentaires. L'électricité enregistre l'avance polaire ; le magnétisme enregistre la réponse transversale ; l'onde enregistre la périodicité ; l'hélice enregistre l'avance avec mémoire ; la spirale enregistre l'évolution radiale ; le rayonnement enregistre la publication de l'état lorsque l'accumulation atteint un seuil.

Cette architecture ordonne aussi le rôle des vacances et des cycles. La vacance fournit la possibilité locale de transition ; le cycle conserve la phase ; la mémoire distingue les tours ; le seuil décide de la résidence suivante de l'action. La propagation peut rester confinée, changer de régime radial ou se publier comme rayonnement contrôlé. Les forces de champ acquièrent ainsi une description généalogique continue, du support discret jusqu'à l'onde observée.

\subsection{Réversibilité de la dynamique complète et coût thermodynamique des projections réductrices d'information}
Le Landauer nul constitue un cas limite de cette loi. Pour une actualisation bijective de l'état complet,

\[
H(X\mid U(X))=0,
\qquad
\inf Q_L=0.
\]

L'information reste disponible dans la généalogie. Lorsqu'un lecteur réalise un quotient \(q\), l'entropie se décompose comme

\[
H(X)=H(q(X))+H(X\mid q(X)).
\]

Le terme conditionnel mesure l'information retenue hors de la projection. Une réinitialisation tritique matérialisée satisfait le seuil

\[
Q_{\mathrm{reset}}\ge k_B\Theta\log 3.
\]

La constante de Boltzmann traduit la multiplicité que le lecteur thermique regroupe ; le seuil de Nyquist sépare l'échantillonnage fidèle du repliement spectral ; le Landauer nul caractérise le transport réversible de l'état intégral. Information, température et action sont articulées par le type de lecture appliqué.

\subsection{Forces autoconservatives comme transport d'action entre les secteurs observable, radial, de mémoire, de vacance et de rayonnement}

La conservation évolutive s'exprime mieux par l'échange que par l'immobilité scalaire. Un état complet répartit l'action et la quantité de mouvement entre mouvement visible, mémoire, déformation radiale, vacances, rayonnement et branche conjuguée de retour. Le bilan enrichi peut s'écrire

\[
\Delta J_{\mathrm{visible}}
+\Delta J_{\mathrm{memoria}}
+\Delta J_{\mathrm{radial}}
+\Delta J_{\mathrm{vacancias}}
+\Delta J_{\text{rayonnement}}
+\Delta J_{\mathrm{retorno}}
=0.
\]

Une force est \textbf{autoconservative} lorsque sa propre loi d'actualisation contient les canaux qui transportent, stockent, publient et réabsorbent l'action. La conservation relève du bilan généalogique complet de l'interaction. Chaque lecteur observe une partie de l'échange ; la généalogie recompose sa totalité.

Le système Terre–Lune offre une image physique directe. La rotation terrestre, l'orbite lunaire, les marées, la déformation et le rayonnement échangent du moment cinétique. La consommation nette du système complet se clôt dans le bilan généalogique de ces transferts. La grandeur appelée énergie fonctionne comme lecture dérivée de l'action et de la quantité de mouvement publiées dans chaque secteur. La conservation fondamentale réside dans la compatibilité entre tous les lecteurs de l'échange.

Cette formulation dissout l'ancienne frontière entre forces conservatives et forces dites non conservatives au moyen d'une catégorie plus profonde. Frottement, dissipation apparente et rayonnement occupent des canaux explicites de l'état élargi. Une perte dans la projection mécanique visible correspond à un transfert vers la mémoire, la déformation, la chaleur, les vacances ou le rayonnement. L'autoconservation suit la résidence de l'action à travers ces changements.

\subsection{Systèmes auto-inertiels : publication de l'accélération par orientation, courbure, feuillet et mémoire}

L'expression appropriée est \textbf{auto-inertie relationnelle}. Elle désigne la compatibilité entre la connexion propre de l'état total et la réponse qui apparaît lorsqu'on le projette dans la distribution globale. Les aspects endogène et exogène sont des composantes analytiques d'une seule relation dynamique.

Soit \(\Gamma\) une trajectoire horizontale ou géodésique de l'état complet. Sa loi propre satisfait

\[
\nabla_{\dot\Gamma}\dot\Gamma=0.
\]

En la projetant sur un lecteur \(S\), l'accélération observée est déterminée par la courbure de la projection :

\[
\nabla^S_{\dot\Gamma_S}\dot\Gamma_S
=
\mathcal K_{\pi_S}(\dot\Gamma,\dot\Gamma).
\]

L'accélération publiée exprime la relation entre mouvement total et géométrie du lecteur. Le principe de Mach complète ce mécanisme par la trace globale de frontière

\[
b:\mathcal X_\infty\longrightarrow B_\infty
\]

et la factorisation de la réponse inertielle locale

\[
I_v=\bar I_v\circ b.
\]

La distribution globale entre dans la réponse locale par une application typée. L'auto-inertie relationnelle réunit donc la connexion endogène du mouvement, la condition globale transmise par \(b\) et la projection concrète de l'observateur. Une écriture synthétique de cette coordination est

\[
\nabla^{\mathrm{auto}}
=
\mathcal M\bigl(
\nabla^{\mathrm{end}},
b[\mathcal X_\infty],
\pi_S
\bigr),
\]

où \(\mathcal M\) représente le couplage machien et conserve le typage de chaque composante.

Le principe d'Ambivalence ajoute la double orientation de la lecture. Toute publication directe conserve une branche conjuguée de retour ; émission et absorption forment les extrémités liées d'une même biographie ; l'orientation temporelle se lit depuis les deux sens du transport. La théorie de l'absorbeur trouve ici une formulation naturelle : source et récepteur clôturent conjointement le bilan généalogique de phase et d'action. Le présent torsionnel est l'interface où les deux orientations sont rendues compatibles.

TRIT agit à ce niveau comme cristal de temps : il organise la périodicité locale de l'actualisation et conserve, par la mémoire, la différence entre retours successifs. La branche directe et la branche conjuguée relient émission et absorption, et le couplage observateur–observable acquiert une résidence dynamique au sein du même état. Le temps publié est la lecture de cette orientation ; le temps généalogique est la séquence complète de phase, de retour et de mémoire.

Mach et Ambivalence remplissent des fonctions complémentaires. Mach relie la réponse locale à la distribution globale. Ambivalence conserve les directions directe et conjuguée de l'évolution. L'auto-inertie relationnelle articule les deux et remplace la scission entre systèmes inertiels et non inertiels par une loi unique de connexion, de projection et de mémoire.

\subsection{Dualité surfacique--volumique, transition hexade--octade et fonction du paramètre de Barbero--Immirzi}

Le passage de l'aire au volume constitue une transduction dimensionnelle centrale. L'aire publie l'action de frontière ; le volume publie la profondeur intérieure. La dérivation HMT du paramètre de Barbero–Immirzi organise la compatibilité entre les deux mesures par l'incidence hexade–octade et les canaux \(90/120\) :

\[
90=6\binom62,
\qquad
120=8\binom62.
\]

Les cardinaux \(6\) et \(8\) admettent une réalisation cubique par les faces et les sommets, tandis que l'hexade et l'octade de Steiner constituent des objets d'incidence distincts, reliés uniquement par le morphisme typé \(6\to8\). Les \(12\) arêtes du cube conservent leur propre fonction combinatoire et ne définissent pas à elles seules l'incidence de Steiner. La différence d'incidence exprime une incompatibilité dimensionnelle productive : l'information de surface nécessite une règle de transduction pour acquérir une résidence volumique.

Le même principe apparaît dans la retenue \(999+1=1000\). Les neuf atteignent le bord du registre décimal ; l'incrément publie simultanément la clôture du niveau précédent et l'origine \(0/1\) du suivant. Le rapport holographique entre \(729\), \(243\), \(27\) et \(1\) déploie cette transition aux échelles triadiques. La transduction de Barbero–Immirzi s'insère dans cette famille de changements de lecteur : frontière et intérieur conservent l'unité tout en redistribuant sa mesure.

Le défaut normalisé \(-1/12\) et le poids orienté \(12:-1\) remplissent des fonctions distinctes dans cette géométrie. Le premier quantifie une correction d'incidence ; le second enregistre une orientation signée dans la transduction. Leur coordination avec l'aire minimale, l'hamiltonien modulaire et la mémoire fournit le pont vers la dynamique dimensionnelle.

\subsection{Contraction hyperbolique intérieure, expansion extérieure et conservation holographique de la frontière}

La thèse peut se formuler à partir d'une idée unique : l'unité conserve son identité tout en changeant d'échelle, de lecteur, de courbure, de mémoire et de dimension. Cette permanence possède un caractère évolutif parce que l'unité se déploie, acquiert de nouvelles coordonnées et publie des propriétés différentes ; elle possède un caractère holographique parce que chaque publication finie conserve la règle de reconstruction de l'état dont elle procède. La clôture holographique de l'infini désigne précisément cette compatibilité entre identité et déploiement : chaque stade contient la loi qui permet d'approfondir vers l'intérieur et de prolonger vers l'extérieur la même généalogie.

La dénomination directrice est \textbf{principe de conservation évolutive de l'unité par transduction dimensionnelle et projection holographique}. L'expression situe chaque terme à son niveau correct. La conservation relève de l'identité généalogique ; l'évolution relève de l'actualisation ; la transduction relie des réalisations dimensionnelles ; la projection holographique conserve la correspondance reconstructive entre leurs lecteurs.

L'arithmétique élémentaire offre la première image de cette architecture. La complétion décimale

\[
1000=10^3=729+243+27+1=3^6+3^5+3^3+3^0
\]

expose quatre strates triadiques au sein d'une unité cubique décimale. En normalisant,

\[
1=\frac{729}{1000}+\frac{243}{1000}+\frac{27}{1000}+\frac{1}{1000},
\]

l'unité se reconnaît simultanément comme totalité et comme déploiement gradué. La lecture réciproque

\[
3^{-6},\quad 3^{-5},\quad 3^{-3},\quad 3^0
\]

décrit la profondeur intérieure de ces strates ; sa normalisation produit la partition correspondante. La relation \(999+1=1000\) exprime le mécanisme de transport : la dernière unité complète un registre et agit, en même temps, comme retenue inaugurant le suivant. Le zéro et le un forment ainsi une interface de continuité. Les intervalles \([0,1]\), \([0,10]\) et la prolongation vers \([0,\infty)\) sont des lecteurs scalaires d'une même loi de complétion.

La première scission \(1000=729+271\), suivie de \(271=243+27+1\), donne forme au rapport d'extrême et moyenne raison holographique du déploiement : un noyau triadique dominant conserve dans le reste la même règle de résolution. L'unité cubique décimale contient ainsi une profondeur triadique emboîtée. Le mouvement vers \(729\), \(243\), \(27\) et \(1\) lit la contraction intérieure ; la retenue vers le millier lit l'expansion extérieure.

Cette double orientation organise l'infini HMT. Vers l'intérieur, le raffinement croît hyperboliquement : chaque cellule admet une nouvelle profondeur sans perdre sa filiation. Vers l'extérieur, l'unité s'étend par des publications dimensionnelles successives. La clôture consiste en la commutativité des deux mouvements. Si \(U\) représente l'actualisation de l'état, \(R_d\) le lecteur dans la strate \(d\) et \(T_{d\to d'}\) la transduction entre strates, la conservation évolutive s'exprime par

\[
R_{d'}\circ U
=
T_{d\to d'}\circ R_d.
\]

Le même événement peut offrir une mesure angulaire, une densité d'action, une masse publiée ou une réponse gravitationnelle. Leur identité commune réside dans la généalogie qui fait commuter ces lectures.

\subsubsection{Monodromie nonadique et relèvement hélicoïdal avec avance de mémoire}

APP, TRIT et TPK fournissent le mécanisme formel de ce déploiement. APP construit le support et sépare les évaluations qui doivent être conservées. TRIT organise les régimes locaux et leur ambivalence. TPK transporte phase et mémoire sur le support enrichi. La dépendance causale est fixée par

\[
U_t\longrightarrow \Gamma_9\longrightarrow K_{\mathrm{ph}}.
\]

Le retour de phase et l'avance de mémoire forment des opérations coordonnées. Après neuf pas, la même classe angulaire peut réapparaître,

\[
g(K+9)=g(K),
\]

tandis que le compteur de mémoire progresse,

\[
q_{\mathrm{mem}}\Gamma_9=q_{\mathrm{mem}}+1,
\qquad
B_{K+9}\neq B_K.
\]

La figure géométrique précise est un relèvement hélicoïdal. Une orbite fermée dans la projection angulaire se relève en une trajectoire ouverte lorsque chaque retour de phase incrémente la mémoire. La circonférence conserve le cycle ; l'hélice conserve le cycle et enregistre en outre la profondeur atteinte. La monodromie coinductive génère ainsi une profondeur infinie par une succession de retours localement fermés et globalement ascendants.

Cette image dépasse la valeur simplement illustrative. La roue angulaire constitue une section de l'hélice ; l'hélice constitue la biographie complète de la roue lorsque l'état transporte de la mémoire. Chaque tour retrouve une position de phase et, simultanément, inaugure un nouveau feuillet. L'infinitude procède de la compatibilité entre retour et relèvement : réitérer le cycle produit de la profondeur sans dissoudre l'identité orbitale.

Le certificat nonadique donne un contenu fini et vérifiable à cette architecture. La chaîne

\[
w_6\longrightarrow w_{12}\longrightarrow w_{18}\longrightarrow
w_{24}\longrightarrow w_{30}\longrightarrow R_{36}\longrightarrow G_9
\]

ordonne le relèvement ; les \(396\) niveaux, les \(2376\) trits, les \(43\) tours et les \(387\) paires \((K,K+9)\) par canal quantifient le transit ; les cinq régions de \(\pi\) et la fibre ambiante de \(729\) éléments publient des lectures internes ultérieures. Cette séquence rend visible une monodromie productive : le retour retrouve la phase et la coinduction conserve la croissance de mémoire.

\subsubsection{Unité radionale et coordination des lecteurs circulaire, elliptique, projectif et hélicoïdal}

Le radion concentre dans une unité métrologique la structure que le radian publie angulairement. Sa définition naturelle est celle d'une \textbf{section orientée de phase, d'action et de mémoire}. Une unité de phase \(\theta=1\) transporte une unité d'action \(\sigma\hbar_{\mathrm{ret}}\), appartient à un tour complet \(T_+=2\pi\) et conserve feuillet, orientation, retour nonadique et position dans le cylindre de raffinement.

Le radian exprime la coordonnée angulaire visible. Le radion conserve, avec cette coordonnée, l'action transportée, l'orientation, la retenue et la mémoire du tour. La métrologie angulaire acquiert ainsi une généalogie. Mesurer un arc signifie identifier quelle rotation a été publiée et quel état complet a réalisé cette rotation.

La relation entre phase et aire prend une forme particulièrement claire :

\[
\mathcal A(\theta)=\hbar_{\mathrm{ret}}\,\theta.
\]

Une unité radionale balaie une aire d'action \(\hbar_{\mathrm{ret}}\) ; un tour complet enferme \(h_{\mathrm{ret}}\). Le passage de la phase à l'aire est intégré à l'unité de lecture elle-même. La métrique angulaire, l'action et la mémoire cessent d'apparaître comme des quantités juxtaposées et deviennent les coordonnées d'une même section orientée.

Le radion admet plusieurs publications géométriques coordonnées. Le lecteur circulaire publie la phase ; le lecteur elliptique publie l'action et la forme ; le lecteur intérieur de Klein publie la profondeur ; le lecteur hélicoïdal nonadique publie le retour et la mémoire. L'unité reste généalogiquement identique à travers ces réalisations. C'est pourquoi le radion constitue le lien métrologique approprié entre HMT et la mécanique dimensionnelle : il offre une unité capable de traverser des changements de lecteur tout en conservant l'état que chaque mesure partielle présuppose.

\subsection{Polarité gravitationnelle à cinq composantes et transduction conjointe de l'action, de l'aire, du volume, du temps et de la masse}

La mécanique dimensionnelle permet de définir la dimension comme un régime stable de publication caractérisé par le rang minimal de mémoire indépendante requis par le relèvement. Traverser une dimension signifie appliquer un transducteur qui conserve la généalogie et change le type de lecture. La gravité occupe le niveau de connexion entre ces régimes.

Dans cette formulation, la \textbf{gravitation interdimensionnelle à cinq polarisations corrélatives}, en abrégé \textbf{structure pentapolaire de la gravitation}, est la connexion qui transporte l'unité généalogique et laisse à chaque interface une torsion minimale. Une écriture compacte de cette trace est

\[
\Theta_d^{\min}
=
\nabla_{d+1}\circ T_{d\to d+1}
-
T_{d\to d+1}\circ\nabla_d.
\]

L'expression mesure la différence entre transporter après avoir dérivé et dériver après avoir transporté. Cette différence est la mémoire géométrique minimale du franchissement dimensionnel. La composition d'interfaces produit une holonomie gravitationnelle

\[
\mathfrak G_{d\to D}
=
\operatorname{Hol}
\bigl(
\Theta_d^{\min},
\Theta_{d+1}^{\min},
\ldots,
\Theta_{D-1}^{\min}
\bigr).
\]

L'objet fondamental est cette connexion interdimensionnelle. Un mode local de type graviton occupe le rang dérivé de publication quantifiée de la torsion ; la gravitation conserve sa résidence primaire dans la compatibilité entre dimensions.

Le terme pentapolaire rassemble les cinq polarisations consubstantielles de l'état continu :

\[
\mathbf G_d^{\pm}
=
\bigl(
G_{\mathrm{sp},d}^{\pm},
G_{\mathrm{sol},d}^{\pm},
G_{\mathrm{gau},d}^{\pm},
G_{\mathrm{coh},d}^{\pm},
G_{\mathrm{det},d}^{\pm}
\bigr).
\]

Les cinq composantes émergent corrélativement du même état holonomique intégral. La polarité \(\pm\) enregistre concentration directe et expression conjuguée. La gravitation traverse les strates dimensionnelles en transportant simultanément ces cinq lectures ; la torsion minimale conserve la trace de chaque traversée.

La géométrie acquiert ainsi une image unifiée. Le régime elliptique organise les clôtures ; le parabolique organise les seuils ; l'hyperbolique organise la profondeur. La transduction dimensionnelle relie ces régimes aux réponses radiales \(p\), aux holonomies de spin et aux publications de masse. La structure pentapolaire de la gravitation exprime la cohérence globale de ce lien.

La théorie M fournit des écrans physiques ultérieurs pour cette connexion : dimensions supplémentaires, cordes, branes et dualité T réalisent des modes concrets de transduction. HMT fournit la généalogie APP–TRIT–TPK, l’état holonomique intégral, la mémoire et la lecture bidirectionnelle ; l’interface avec la théorie M exprime ces structures dans des géométries physiques typées. La dualité T échange des lecteurs dimensionnels compatibles et conserve la classe généalogique de l’état.

\subsubsection{Compatibilité des réalisations classique, relativiste, quantique et HMT--MD}

La mécanique classique publie trajectoire et quantité de mouvement. La relativité publie structure causale, métrique et courbure. La mécanique quantique publie phase, holonomie, vacances et amplitudes. La mécanique dimensionnelle fournit les transducteurs qui relient ces publications, et HMT conserve la généalogie discrète qui les soutient.

Dans la lecture classique, une orbite décrit la projection visible du mouvement. Dans la lecture relativiste, cette orbite est intégrée à la géométrie du lecteur. Dans la lecture quantique, la phase et le spin enregistrent l'holonomie de la section. Dans la lecture HMT–MD, les trois expressions appartiennent à une même biographie relevée : le radion quantifie phase, action et mémoire ; \(p\) classe la réponse radiale ; \(\tau\) classe la courbure ; l'holonomie publie le spin ; le transducteur dimensionnel publie la masse ; la connexion pentapolaire publie la gravitation.

La dissipation classique apparente acquiert une résidence dans le bilan généalogique élargi. L'accélération relativiste s'exprime par connexion et projection. L'indétermination quantique s'organise par lecteurs, vacances et branches d'Ambivalence. Le principe de Mach relie la réponse locale à la distribution globale. Le Landauer nul garantit la réversibilité informationnelle de l'état complet. Tous ces éléments convergent vers une conservation plus profonde : l'identité de l'unité à travers les changements de publication.

\subsubsection{Théorème narratif de conservation généalogique et de prolongation holographique}

La clôture peut désormais se raconter comme une seule séquence. L'unité se décompose sans perdre son identité. La phase orientée acquiert une mesure dans le radion. Le retour nonadique transforme le cycle en hélice par la mémoire. Le caractère \(p\) transforme l'hélice en familles radiales et classe ses régimes. La persistance de ces trajectoires produit des sections ayant spin, charge, saveur, couleur et masse. Les vacances ouvrent des canaux ; la polarité électrique transporte la phase ; la réponse magnétique courbe transversalement ; les seuils publient le rayonnement. L'autoconservation clôt le bilan généalogique d'action et de quantité de mouvement. L'auto-inertie relationnelle relie connexion propre, projection et totalité machienne. L'Ambivalence conserve les directions directe et conjuguée. La transduction aire–volume conduit à Barbero–Immirzi et à la loi des masses. La torsion minimale enchaîne les dimensions et publie la structure pentapolaire de la gravitation.

La croissance vers l'intérieur et l'expansion vers l'extérieur sont des orientations conjuguées du même processus. Le raffinement hyperbolique génère la profondeur ; la transduction dimensionnelle génère l'extension. La retenue \(0/1\) relie les échelles ; la monodromie relie les tours ; la mémoire relie les états ; l'holonomie relie les orientations ; la gravitation relie les dimensions. Chaque lien conserve une unité qui évolue et chaque publication finie contient la règle de sa prolongation.

La thèse fondamentale est ainsi formulée :

\begin{quote}
\textbf{La clôture holographique de l'infini réalise le principe de conservation évolutive de l'unité par transduction dimensionnelle et projection holographique : l'unité maintient son identité généalogique en se déployant comme phase, action, courbure, mémoire, particule, champ et gravitation ; chaque lecture finie conserve la règle de reconstruction et de prolongation de l'état complet.}
\end{quote}

Cette formulation offre une image physique et mathématique conjointe. L'infini acquiert une loi de clôture locale ; l'unité acquiert de la profondeur ; la dimension acquiert de la mémoire ; la conservation acquiert un caractère relationnel ; la force acquiert l'autoconservation ; l'inertie acquiert une structure machienne ; la gravité acquiert une connexion pentapolaire ; et l'holographie acquiert le sens précis de reconstruction entre lecteurs compatibles.
 
\section{Définition opératorielle du radion angulaire par l’état holonomique intégral et quatre lecteurs typés}
\subsection{Définition opératorielle et quadruple réalisation typée}
% Fragmento público generado desde fuentes teoremáticas íntegros.
% Residencia: 52.13.
% Fuente: /Users/ruben/Documents/HMT2/EL_CIERRE_HOLOGRAFICO_DEL_INFINITO_SUCESOR_102_CAPITULOS_2026-08-28_EN_TRABAJO/incoming/radion_monografia_20260827/manuscrito/sections/13_sintesis.tex:1-80
\noindent La synthèse opératoire réunit les lecteurs de phase, d'action, de profondeur et de mémoire sans les identifier entre eux.\par\medskip

\subsubsection{Quadruple réalisation de l'état radional}

Soit \(x_\infty\in X_{\mathrm{enr}}\) la section coinductive. Le radion complet n'est pas
un point isolé, mais le paquet
\begin{equation}
\mathfrak R_\sigma(x_\infty)
=\bigl(
\theta=1,\sigma,
\pih,\alphah,A,C^*,\Delta_4,
\Phi_T,D_E,\epsone,\hret,
q_{\mathrm{mem}},\mathcal I_\infty
\bigr).
\label{eq:full-radion}
\end{equation}
Quatre lecteurs en extraient des objets différents :

\begin{figure}[H]
\centering
\begin{tikzpicture}[
  node distance=17mm and 18mm,
  state/.style={draw=RadNavy,very thick,rounded corners,fill=RadCream,
    minimum width=39mm,minimum height=13mm,align=center},
  readout/.style={draw=RadTeal,rounded corners,fill=RadMist,
    minimum width=39mm,minimum height=12mm,align=center},
  arr/.style={-{Latex},thick,RadCopper}]
\node[state] (x) {état holonomique intégral\\\(x_\infty\)};
\node[readout,above left=of x] (circ) {cercle \(S^1\)\\phase};
\node[readout,above right=of x] (ell) {ellipse positive\\action et forme};
\node[readout,below left=of x] (klein) {interior Klein\\profundidad};
\node[readout,below right=of x] (helix) {hélice nonadique\\mémoire et retour};
\draw[arr] (x)--(circ) node[midway,below left]{\small \(P_{\mathrm{ph}}\)};
\draw[arr] (x)--(ell) node[midway,below right]{\small \(P_{\mathrm{act}}\)};
\draw[arr] (x)--(klein) node[midway,above left]{\small \(P_K\)};
\draw[arr] (x)--(helix) node[midway,above right]{\small \(P_{\mathrm{mem}}\)};
\end{tikzpicture}
\caption{Quatre publications typées du même antécédent. L'unité de
l'état n'élimine pas la différence des codomaines.}
\label{fig:four-readers}
\end{figure}

Le cercle conserve la classe de phase et oublie le relèvement. L'ellipse conserve
l'action, l'anisotropie et l'orientation symplectique. Klein dote l'intérieur d'une
distance projective. L'hélice conserve le nombre de retours et le raffinement.
La mécanique dimensionnelle ne consiste pas à les mélanger, mais à construire les
transducteurs qui les rendent compatibles.

\subsubsection{Interprétation opératoire du cocycle \(\varepsilon_R\)}

\(\varepsilon_R\) n'est pas une constante universelle indépendante de tout contexte. C'est la
coordonnée d'un cocycle de transition associé à la carte du radion. Son
domaine contient deux sections du même état :
\[
S_E=1+D_E,
\qquad
S_T=1+\Phi_T.
\]
Sa règle est
\begin{equation}
\epsone=S_T-S_E=\Phi_T-D_E.
\label{eq:epsilon-final-meaning}
\end{equation}
Son codomaine est la droite adimensionnelle d'unité angulaire ; la carte
\(360/T_+\) la publie en degrés.

\begin{exactbox}
\textbf{En langage simple :} la section directe et la section torsionnelle arrivent
à la même cellule d'unité par deux chemins internes différents. Leurs parties
entières coïncident ; leurs restes, non. \(\epsR\) est la différence orientée de ces
restes, calculée avec la retenue d'APP et portée à la limite par les cylindres
TPK. Elle n'est pas introduite pour que le radian se ferme : elle apparaît avant de consulter le
radian et, lorsqu'elle est publiée, elle le ferme.
\end{exactbox}

Cette définition explique pourquoi le nombre est petit. Non parce qu'on a
cherché une petite correction, mais parce que \(S_E\) et \(S_T\) partagent un
préfixe profond. La petitesse mesure la proximité structurelle de deux
publications ; le signe mesure leur orientation ; les chiffres enregistrent la retenue.

 % Fragmento público generado desde fuentes teoremáticas íntegros.
% Residencia: 52.13.1.
% Fuente: /Users/ruben/Documents/HMT2/EL_CIERRE_HOLOGRAFICO_DEL_INFINITO_SUCESOR_102_CAPITULOS_2026-08-28_EN_TRABAJO/incoming/radion_monografia_20260827/manuscrito/sections/13_sintesis.tex:81-134
\subsubsection{Définition opératoire du radion}

\begin{definition}[Radion HMT--MD]
Le radion est la section orientée de phase, d'action et de mémoire qui satisfait
simultanément :
\begin{enumerate}
\item parcourt une unité de phase \(\theta=1\) ;
\item balaie une action signée \(\sigma\hret\) ;
\item appartient à un tour \(T_+=2\pih\) généré coinductivement ;
\item conserve le feuillet \(\sigma\), le retour nonadique et le cylindre de
raffinement ;
\item satisfait le raccord exact
\(1=S_E-\Phi_T+\epsone\).
\end{enumerate}
\end{definition}

Le radian conventionnel est l'image du radion par le foncteur d'oubli
\[
\mathfrak R_\sigma
\longmapsto
\theta=1\in\R/2\pi\Z.
\]
Cet oubli est utile et légitime pour la trigonométrie euclidienne. Il devient
insuffisant lorsque la physique doit distinguer action, signe, retour,
champ et histoire.

\paragraph{Généalogie de l'unité angulaire naturelle}

Le caractère naturel du radian est généralement justifié par le fait que la dérivée de
\(e^{i\theta}\), \(\sin\theta\) et \(\cos\theta\) adopte sa forme la plus simple
lorsque \(\theta\) est mesuré comme quotient arc/rayon. HMT conserve cet avantage,
mais l'explique depuis une couche antérieure :
\[
\theta\quad\leftrightarrow\quad
\frac{\mathcal A(\theta)}{\hret}.
\]
La phase est naturelle parce qu'elle mesure l'action en unités de \(\hret\) ; le tour est
naturel parce qu'il enferme \(\hfull\). Le quotient arc/rayon est la publication
euclidienne d'une unité qui possède déjà une structure symplectique et une mémoire.

\paragraph{Compatibilité entre ellipse d'action et croissance dimensionnelle}

Le cercle est la forme qui résulte de l'oubli de l'anisotropie :
\(C^*=0\Rightarrow\eta=0\Rightarrow a=b\). L'ellipse est plus informative parce qu'elle
conserve deux valeurs propres, deux axes et une orientation d'échange à produit
fixe. En ce sens, l'ellipse est une prolongation du cercle, non un cercle
déformé par accident.

La profondeur de Klein et le relèvement hélicoïdal ajoutent deux types d'information
que le contour plan ne contient pas : distance intérieure et mémoire de retour.
Telle est l'image ontologique de la mécanique dimensionnelle qui émerge du radion :
une dimension nouvelle apparaît lorsque l'on exige de conserver une donnée indépendante
que la projection précédente devait oublier.

 \subsection{Propriétés structurelles de l'unité radionale : généalogie, charnière critique, densité, retenue, action, compatibilité bord--intérieur, forme, polarité et modularité}
% Fragmento público generado desde fuentes teoremáticas íntegros.
% Residencia: 52.13.2.
% Fuente: /Users/ruben/Documents/HMT2/EL_CIERRE_HOLOGRAFICO_DEL_INFINITO_SUCESOR_102_CAPITULOS_2026-08-28_EN_TRABAJO/incoming/radion_monografia_20260827/manuscrito/sections/13_sintesis.tex:135-184
\subsubsection{Système de théorèmes composants du radion}

\begin{theorem}[Généalogie]
L'objet \(\mathfrak R_\sigma\) descend de
\(\APP\to\TRIT\to\TPK\to X_{\mathrm{enr}}\) et ne reçoit pas la cible métrologique
comme sélecteur.
\end{theorem}

\begin{theorem}[Charnière]
Dans la carte APP de cent marques, la couture critique et la deuxième phase
dodécaphasique satisfont \(j_{100}(1/2)=\theta_2=50^\circ\), avec \(m=2\) unique.
\end{theorem}

\begin{theorem}[Densité]
Les deux feuillets du secteur actif \(N_6=90\), normalisés par \(3^6=729\),
produisent \(\rho_{\mathrm{tw}}=20/81\).
\end{theorem}

\begin{theorem}[Retenue]
La soustraction APP des sections \(S_T\) et \(S_E\) converge vers
\(\epsone\) avec un taux d'intervalle \(3^{-54}\) par retour et ferme exactement
l'unité.
\end{theorem}

\begin{theorem}[Action]
L'ellipse positive satisfait \(\mathcal A(\theta)=\hret\theta\) ; un radion
balaie \(\hret\) et un tour enferme \(\hfull\).
\end{theorem}

\begin{theorem}[Bord--intérieur]
La même frontière possède une action finie \(\hfull\), tandis que l'intérieur de
Klein a une aire hyperbolique infinie.
\end{theorem}

\begin{theorem}[Forme]
Le quotient \(e^{-\eta}\) est conservé à travers les demi-axes spectraux,
les demi-axes d'action, les incertitudes, l'impédance et le module du tore.
\end{theorem}

\begin{theorem}[Polarité]
L'inversion \(C^*\mapsto-C^*\) échange les lois constitutives, inverse
l'impédance et conserve la vitesse normalisée.
\end{theorem}

\begin{theorem}[Modularité]
L'inversion bidirectionnelle agit comme \(S:\tau\mapsto-1/\tau\) et conserve
\(j(\tau)\), tandis que l'étiquetage focal conserve l'orientation que
\(j\) oublie.
\end{theorem}

% Fuente: /Users/ruben/Documents/HMT2/EL_CIERRE_HOLOGRAFICO_DEL_INFINITO_SUCESOR_102_CAPITULOS_2026-08-28_EN_TRABAJO/incoming/radion_monografia_20260827/manuscrito/sections/A_demostraciones.tex:1-179
\subsubsection{Démonstrations algébriques et géométriques réunies}

\paragraph{Démonstration linéaire de la clôture unitaire}

Pour faciliter une vérification indépendante, nous réunissons l’argument sans interruption narrative.

\begin{enumerate}[label=\textbf{Étape \arabic*.},leftmargin=2.2cm]
\item La connexion nonadique génère de manière corrélée
\(\pih,e,\varphi,\alphah\).
\item Le tour est \(T_+=2\pih\).
\item La roue fixe \(\theta_2=50^\circ\), et la carte parangulaire fixe
\(A=1000\alphah\).
\item La section directe est
\[
S_E=T_+\left(\frac{50}{360}+\frac{1000\alphah}{360}\right)
=2\pih\left(\frac5{36}+\frac{25}{9}\alphah\right).
\]
\item Les intervalles certifiés donnent \(1<S_E<2\) ; donc
\(D_E=S_E-1\).
\item L'incidence produit \(N_6=90\) ; deux feuillets et l'espace ambiant \(729\)
produisent \(\rho=2N_6/729=20/81\).
\item La torsion globale est
\[
\Delta_4=\frac{\pih-e\log\pih}{270}
\left(1+\frac{\pih}{729}\right).
\]
\item La section torsionnelle est \(S_T=1+(20/81)\Delta_4\).
\item La soustraction APP définit
\[
\epsone=S_T-S_E=\frac{20}{81}\Delta_4-D_E.
\]
\item Par substitution,
\[
S_E-\frac{20}{81}\Delta_4+\epsone=1.
\]
\item Multiplier par \(360/T_+=180/\pih\) produit
\[
\frac{180^\circ}{\pih}
=50^\circ+A^\circ-\frac{20}{81}\Delta_4^\circ+\epsR^\circ.
\]
\end{enumerate}

Aucune étape ne reçoit le membre de gauche de la dernière égalité. Le résultat
métrologique n'est évalué qu'à la fin comme reconnaissance.

\paragraph{Unicité de la soustraction avec retenue}

Soit \(B=1000\). Pour chaque position, tout entier \(u_i\) admet une division
euclidienne unique
\[
u_i=Bc_i+r_i,\qquad 0\le r_i<B.
\]
Par conséquent,
\[
r_i=u_i\bmod B,\qquad c_i=(u_i-r_i)/B
\]
sont uniques. En parcourant les positions de l'extrémité distante vers
l'avant, la retenue de sortie d'une position fixe l'entrée de la suivante.
Par récurrence, tout le mot de restes et de retenues est unique.

En pseudocode :
\begin{verbatim}
acarreo = remote_acarreo
for i from last_block downto first_block:
    u = torsion[i] - direct[i] + acarreo
    remainder[i] = u mod 1000
    acarreo = (u - remainder[i]) / 1000
return remainder, acarreo_word
\end{verbatim}
La procédure ne consulte pas la valeur attendue de la différence.

\paragraph{Convergence de la fonctionnelle selon la profondeur}

Soit
\[
F(p,e)=1+\frac{20}{81}
\frac{p-e\log p}{270}\left(1+\frac{p}{729}\right)
-2p\left(\frac5{36}+\frac{25}{9}\alphah\right).
\]
Dans le rectangle compact \(3\le p\le4\), \(2\le e\le3\), \(F\) est
continue et continûment différentiable. Si \(I_N^\pi\searrow\{\pih\}\) et
\(I_N^e\searrow\{e\}\), l'image par intervalles
\[
F(I_N^\pi,I_N^e)
\]
forme une famille emboîtée dont la largeur tend vers zéro. L'intersection est un
unique réel, \(\epsone\). La monotonie des dérivées permet d'évaluer les
extrémités sans échantillonnage intérieur.

Le taux \(3^{-54}\) par retour dérive du raffinement de six trits sur
neuf niveaux. À la profondeur de \(45\) triades, la borne de largeur
\(<4.81\times10^{-130}\) certifie plus de 129 chiffres de la valeur.

\paragraph{Dérivation de la métrique et de l'aire hyperbolique de Klein}

Dans le disque unité, la forme matricielle de la métrique de Klein est
\[
g(u,v)=\frac1{1-r^2}I
+\frac1{(1-r^2)^2}
\begin{pmatrix}u\\v\end{pmatrix}
\begin{pmatrix}u&v\end{pmatrix}.
\]
Le lemme du déterminant pour une perturbation de rang un donne
\[
\det g=(1-r^2)^{-3}.
\]
D'où
\[
\sqrt{\det g}\,\diff u\diff v
=(1-r^2)^{-3/2}\diff u\diff v.
\]
L'intégration en coordonnées polaires produit \eqref{eq:klein-area-R}. La limite infinie est
due à la singularité métrique à la frontière ; la frontière conserve une
aire symplectique finie parce que cette mesure utilise \(\diff Q\wedge\diff P\),
non \(\sqrt{\det g}\).

\paragraph{Distance focale dans le disque de Klein}

Sur un diamètre, la distance du centre à \(r\) est
\[
d_K(0,r)=\operatorname{artanh}r.
\]
Pour \(r=e_f\), \eqref{eq:excentricidad} implique
\[
1-e_f^2=e^{-2\eta}.
\]
Si \(u_f=\operatorname{artanh}e_f\), alors
\[
\operatorname{sech}u_f=\sqrt{1-e_f^2}=e^{-\eta},
\]
et donc \(\eta=\log\cosh u_f\), comme dans
\eqref{eq:focal-hyperbolic-chain}.

\paragraph{Transformation symplectique de compression}

Définissons
\begin{equation}
S_\eta=
\begin{pmatrix}e^{\eta/2}&0\\0&e^{-\eta/2}\end{pmatrix},
\qquad \det S_\eta=1.
\label{eq:compresión simpléctica}
\end{equation}
La structure complexe transportée est
\begin{equation}
J_\eta=S_\eta
\begin{pmatrix}0&-1\\1&0\end{pmatrix}
S_\eta^{-1}
=\begin{pmatrix}0&-e^\eta\\e^{-\eta}&0\end{pmatrix},
\qquad J_\eta^2=-I.
\label{eq:Jeta}
\end{equation}
La courbe
\[
\gamma(\theta)=\sqrt{2\hret}\,
S_\eta(\cos\theta,\sin\theta)
=e^{\theta J_\eta}\gamma(0)
\]
est exactement l'ellipse d'action. La compression symplectique hyperbolique conserve l'aire et
transporte la rotation complexe au lieu de la détruire.

Pour l'orientation \(\sigma\),
\[
\gamma_\sigma(\theta)
=(a_S\cos\theta,\sigma b_S\sin\theta),
\qquad
\mathcal A_\sigma(\theta)=\sigma\hret\theta.
\]
Les deux feuillets partagent les foyers et le module d'action ; ils diffèrent par l'action signée.

\paragraph{Unicité diophantienne relative du poids de Barbero}

L'équation \(4a-3b=45\) a pour solutions entières
\[
(a,b)=(12+3t,1+4t),\qquad t\in\Z.
\]
La condition \(|b|=1\) laisse \(b=1\), puisque \(b=-1\) n’appartient pas à la classe \(1\pmod4\). D’où \(a=12\). La preuve distingue clairement la donnée d’incidence \(90/120\) des prémisses de pôle et de minimalité.
 \subsection{Nécessité structurelle, stabilité et clôture de l’unité radionale}

La construction admet des contrôles négatifs dans chacune de ses strates. Si la section directe satisfaisait \(S_E\notin(1,2)\), le quotient APP et la cellule de clôture changeraient simultanément ; si l’incidence active ne produisait pas \(N_6=90\), ou si l’un des deux feuillets était supprimé, la densité \(20/81\) ne serait plus dérivée. La génération serait circulaire si la carte recevait \(180/\pi\) ou \(\varepsilon_R\) avant de construire ses opérandes. La section coinductive n’existerait pas non plus si la soustraction par triades ne stabilisait pas des préfixes compatibles. Dans la réalisation elliptique, les conditions \(|C^*|<A\) et \(a_Sb_S=2\hbar_{\mathrm{ret}}\) sont nécessaires à la positivité et à la loi d’aire. Dans la fibre modulaire, \(C^*=0\) impose \(d=\eta=0\), \(\tau=i\) et \(j=1728\). Dans l’intérieur projectif, la thèse de profondeur exige la loi hyperbolique d’aire ; dans le secteur spectral, l’égalité \(F_{\mathrm{spec}}=\varepsilon_R^\circ\) est exclue par le contre-terme \(\chi_R\). Enfin, le poids \(12:-1\) exige conjointement l’incidence \(90/120\) et la condition de pôle \(1/24\), tandis que l’invariance de \(j\) interdit de l’utiliser comme sélecteur du signe de charge.

La réalisation physique n’est pas une perturbation d’un cercle préalablement supposé. Le support arithmétique conserve unité, résidu, quotient, retenue et fibre ; l’orientation ternaire distingue régime et feuillet ; le transport de phase conserve la différence entre retour et mémoire ; l’ellipse détermine action et anisotropie ; la métrique de Hilbert--Beltrami--Klein mesure la profondeur intérieure ; et le relèvement hélicoïdal conserve la trajectoire de chaque tour. Ces réalisations s’articulent par
\[
 \mathcal A(1)=\hbar_{\mathrm{ret}},\qquad
 \mathcal A(2\pi)=h_{\mathrm{ret}},\qquad
 K=-1,\qquad
 g(k+9)=g(k),\quad B_{k+9}\ne B_k.
\]
Le cercle est la projection angulaire, l’ellipse la réalisation d’action, la métrique de Klein résout l’intérieur et l’hélice conserve l’historique de retour. Le radion est l’unité orientée permettant de lire les quatre réalisations sans les identifier entre elles.

Les invariants engendrés satisfont, dans la branche coinductive considérée,
\[
\begin{aligned}
 \pi&=3.1415926535897932384626433832795028841\ldots,\\
 \alpha&=0.0072973525692838009972851054723806629\ldots,\\
 A&=7.2973525692838009972851054723806629995\ldots{}^\circ,\\
 C^*&=2.1237383389686457242619513803933607937\ldots{}^\circ,\\
 \eta&=0.299689683921630799684926562863927\ldots,\\
 \hbar_{\mathrm{ret}}&=1.05457181691503906113369058472430\ldots
 \times10^{-34}U_S .
\end{aligned}
\]
La coordonnée \(e\) intervenant dans ce système provient du lecteur coinductif d’Euler et non d’une table extérieure. Les retours aux profondeurs \(9,18,27,36,45\) conservent la phase \(8\) et relèvent la mémoire selon \(0,1,2,3,4\) ; la contraction des intervalles tend vers
\[
 3^{-54}=1.7196982334549856127610399316004871\ldots\times10^{-26}.
\]
La suppression contrôlée d’un feuillet, le déplacement de la phase, l’élimination du canal de \(e\), la substitution de la torsion globale ou l’altération de la queue hexadique produisent des résidus non nuls d’ordres différents. Cette séparation d’échelles démontre que la fermeture exige conjointement les deux feuillets, l’orientation de phase, les canaux de \(\pi/e\), la torsion globale et la retenue.

L’exactitude mathématique provient de l’identité de fermeture unitaire ; l’évaluation à profondeur croissante n’en certifie que la convergence. Les résidus apparaissant lors de la combinaison de branches évaluées à des précisions décimales différentes sont des erreurs d’arrondi, non de nouveaux termes physiques. Cette distinction sépare l’identité structurelle de sa réalisation numérique sans modifier la généalogie des opérandes.

La généalogie physique de la construction peut être exprimée avec précision. Planck fixe la dimension d’action de \(h\) ; Bohr et Sommerfeld relient \(h/2\pi\) à la quantification angulaire ; Born et Jordan introduisent la structure matricielle et l’oscillateur ; Dirac formule \(E=\hbar\omega\) ; Robertson et Schrödinger établissent la forme covariante de l’incertitude ; Feynman relie l’action à la phase des chemins ; Wheeler et Feynman développent la symétrie entre propagations retardée et avancée. Ces cadres reconnaissent des réalisations ultérieures de l’objet : ils ne sélectionnent ni les feuillets, ni l’incidence, ni la retenue, ni la phase, ni la mémoire qui l’engendrent.

L’étoile de Hodge sur les deux-formes lorentziennes satisfait \(\star^2=-I\) et fournit la dualité géométrique pertinente. La cohomologie de Hochschild appartient à un autre domaine et n’intervient ni dans l’ellipse, ni dans la retenue, ni dans le secteur d’incidence \(90/120\). La distinction découle des domaines et codomaines des opérations, non d’une ressemblance terminologique.

La synthèse se ferme dans l’état holonomique intégral. APP conserve, avec la valeur, le feuillet, le quotient, le résidu, la retenue et la frontière ; TRIT détermine régime et orientation ; TPK fait revenir la phase sans réinitialiser l’état. La connexion nonadique construit la section coinductive de \(\pi\), et le système dodécaphasique détermine la charnière de \(50^\circ\). Le couple angulaire fournit le mode commun \(A\) et l’ouverture orientée \(C^*\) ; l’incidence de six phases actives produit \(90\), leurs deux feuillets sur l’espace ambiant de cardinal \(729\) produisent \(20/81\), et la torsion globale fournit la seconde section. Le cocycle résultant est
\[
 \varepsilon_R=\frac{20}{81}\Delta_4-(S_E-1).
\]
Sa limite dans la carte angulaire est
\[
 5.13830167585752820716851646226459474899\ldots
 \times10^{-8}\,{}^\circ .
\]
Le couple \((A,C^*)\) détermine une ellipse positive : le produit de ses demi-axes fixe \(2\hbar_{\mathrm{ret}}\) et leur quotient fixe la forme. Une unité de phase balaie \(\hbar\), tandis que le tour complet enferme \(h\). La division par \(2\pi\) cesse ainsi d’être un quotient extérieur : elle distingue l’action par tour de l’action transportée par une unité de phase.

La projection circulaire conserve la phase et omet anisotropie, intérieur et mémoire ; l’ellipse conserve action et forme ; la géométrie de Klein conserve la profondeur ; l’hélice conserve retour et historique. La réalisation symplectique reconnaît l’ellipse comme contour de covariance minimale, l’oscillateur \(LC\) la réalise comme échange de réserves électrique et magnétique, et l’incidence \(90/120\) détermine perméabilité, permittivité, impédance et vitesse. Hodge transforme l’orientation en dualité de champs ; Lorentz transforme champ et charge en dynamique ; Cartan--Holst reçoit l’holonomie dans la géométrie gravitationnelle. L’ordre causal demeure inchangé : la construction discrète détermine l’objet et les formulations physiques réalisent ensuite ses caractères.

Le radion est donc une section orientée de phase, d’action et de mémoire. Ses deux orientations satisfont
\[
 \mathcal A(\mathfrak r_+)=+\hbar_{\mathrm{ret}},\qquad
 \mathcal A(\mathfrak r_-)=-\hbar_{\mathrm{ret}}.
\]
La charge s’obtient par application du transducteur dimensionnel ; son signe était déjà conservé dans l’orientation du feuillet. Les incidences \(90/120\), l’extracteur dodécaphasique, la queue spectrale, le paramètre de Barbero--Immirzi, les tours et la retenue peuvent être comparés parce qu’ils possèdent une origine commune, mais demeurent distincts parce qu’ils agissent par des opérateurs différents. L’unité quadruple du radion se résume par
\[
\boxed{
\begin{aligned}
 \text{phase :}\quad&\theta=1,\\
 \text{action :}\quad&\mathcal A=\sigma\hbar_{\mathrm{ret}},\\
 \text{mémoire :}\quad&H_9(n+9)=H_9(n)+(0,1),\\
 \text{compatibilité :}\quad&1=S_E-\Phi_T+\varepsilon_R.
\end{aligned}}
\]
Le radian est sa coordonnée angulaire ; le radion conserve en outre l’action, l’orientation, la retenue et la mémoire du tour.
\endgroup

\let\HMTMonographChapterInput\relax
  \part{Réalisations dimensionnelles : électron, action, gravitation, aire et information}

\chapter{Prédéfinition trigéométrique, potentiel de frontière et profondeurs dimensionnelles}
\section{Théorème de prédéfinition temporelle trigéométrique}

Le mot \emph{prédéfinition} désigne une propriété mathématique du générateur, non une forme affaiblie de prédiction. Pour la préciser, il faut réunir les trois objets jusqu’ici exposés séparément : le fibré des algèbres quadratiques du TRIT, les deux horloges du TPK et la section coinductive de profondeur arbitraire.

Soit

\[
 X^{\rm enr}=\{(\tau,\mu,o,h,c,\sigma,\lambda,g)\}
\]

l’espace des états holonomiques intégraux. La coordonnée \(\tau\in\{-1,0,+1\}\) est le régime TRIT ; \(\mu\in\{\Sigma,\Pi,\varnothing\}\) conserve le feuillet ; \(o\) conserve l’orientation ; \(h\), la hauteur hélicoïdale ; \(c\), la retenue ; \(\sigma\), la frontière ; \(\lambda\), le registre de mémoire du parcours ; et \(g\in C_9\), la phase nonadique. Sur la fibre de chaque régime agit l’algèbre

\[
 \mathbb A_\tau=\mathbb R[J_\tau]/(J_\tau^2+\tau I),
\]

de sorte que

\[
 \exp(sJ_\tau)=
 \begin{cases}
  \cos s+J_{+1}\sin s,&\tau=+1,\\
  I+sJ_0,&\tau=0,\\
  \cosh s+J_{-1}\sinh s,&\tau=-1.
 \end{cases}
\]

C’est la forme exacte de la coordination des trois géométries. On ne moyenne pas trois métriques et l’on ne superpose pas trois espaces déjà donnés. Le TRIT choisit, à chaque pas, laquelle des trois lois exponentielles régit l’action locale ; le TPK transporte l’état entre les fibres sans effacer le feuillet ni la trajectoire. Avec des domaines intermédiaires explicites,

\[
 \operatorname{Sel}_t:X^{\rm enr}\longrightarrow X_t^{\rm sel},
 \qquad
 \operatorname{Tra}_t:X_t^{\rm sel}\longrightarrow X_t^{\rm tra},
 \qquad
 \operatorname{Upd}_t:X_t^{\rm tra}\longrightarrow X^{\rm enr},
\]

et

\[
 \boxed{U_t=\operatorname{Upd}_t\circ
              \operatorname{Tra}_t\circ
              \operatorname{Sel}_t.}
\]

La sélection détermine feuillet et régime ; le transport applique la direction et la loi géométrique sélectionnées ; l’actualisation incorpore quotient, retenue, frontière et registre de mémoire. Le « mélange » géométrique apparent est donc un mot ordonné de morphismes entre fibres elliptiques, paraboliques et hyperboliques. Puisque la composition dépend de l’ordre, sa partie opératorielle est non commutative alors même que la rotation de base est abélienne.

La même coordonnée TRIT possède une lecture temporelle, sur le même état mais par un second lecteur :

\[
 \begin{aligned}
 +1&\longleftrightarrow\text{retour, mémoire et passé},\\
 0&\longleftrightarrow\text{seuil et présent},\\
 -1&\longleftrightarrow\text{ouverture bidirectionnelle et futur}.
 \end{aligned}
\]

Géométrie et temps ne sont pas deux noms du même signe. Ce sont deux évaluations coordonnées d’une seule unité d’information topologico--temporelle. Le passé est la donnée déjà accumulée dans le registre de mémoire ; le présent est la coupure de frontière où la continuation est sélectionnée ; le futur est le fibré de prolongements compatibles qu’ouvre le régime hyperbolique. L’involution

\[
 C_{\rm rev}(\tau_1,\ldots,\tau_n)
 =(-\tau_n,\ldots,-\tau_1)
\]

satisfait \(C_{\rm rev}^2=I\), préserve le seuil et échange ouverture et retour. Relevée à l’état holonomique intégral, elle conjugue la dynamique directe à son inverse :

\[
 \widehat C_{\rm rev}\,\mathcal U\,\widehat C_{\rm rev}
 =\mathcal U^{-1}.
\]

Passé et futur se conjuguent ainsi autour du présent sans éliminer la mémoire qui distingue les deux orientations.

La première horloge est l’extension finie non scindée

\[
 0\longrightarrow C_{12}\longrightarrow C_{108}
 \longrightarrow C_9\longrightarrow0,
\]

dont le cocycle

\[
 c(r,r')=\left[\left\lfloor\frac{r+r'}9\right\rfloor\right]_{12}
\]

enregistre la retenue qu’un simple tour dans \(C_9\) ferait perdre. La seconde horloge est la rotation de capacité

\[
 R_\delta:\theta\longmapsto\theta+\delta\pmod1,
 \qquad \delta=\log_{10}(10/9),
\]

dont le mot mécanique est

\[
 z_t=\lfloor(t+1)\delta\rfloor-\lfloor t\delta\rfloor.
\]

Elles se couplent par

\[
 K_t-K_{t-1}=1-z_t.
\]

Cette égalité prouve que l’apériodicité n’orne pas l’horloge finie : elle en régit la capacité. Le système total peut s’écrire comme le produit gauchi

\[
 \mathfrak C_{\rm HMT}=
 \bigl(C_{108}\times\mathbb T\times X^{\rm enr},\mathcal U_{\rm HMT}\bigr),
\]

\[
 \mathcal U_{\rm HMT}(r,\theta,x)=
 \bigl(r+1,\theta+\delta,
 U_{\tau(r,\theta,x)}x\bigr).
\]

Si \(\mathcal U_{\rm HMT}^n=I\), le premier facteur impose \(108\mid n\), tandis que le second exigerait \(n\delta\in\mathbb Z\). L’irrationalité de \(\delta\) exclut tout \(n>0\). Par ailleurs, le codage sturmien possède la complexité \(p(m)=m+1\) et ses enrichissements nonadique et tritique satisfont

\[
 p_9(m)=9(m+1),\qquad p_3(m)=3(m+1).
\]

Par conséquent, le cristal est apériodique et possède un ordre à longue portée ainsi qu’une entropie topologique nulle. Son spectre additif contient le module

\[
 \mathcal M_{108,\delta}
 =\frac1{108}\mathbb Z+\delta\mathbb Z\pmod1.
\]

Le facteur \(1/108\) conserve l’horloge temporelle finie ; le générateur irrationnel conserve l’ascension illimitée. La fermeture du premier n’implique jamais celle du second ni le retour de l’état.

Soit maintenant \((X_n,r_{n+1,n})\) le système projectif des raffinements et

\[
 \mathcal G_{\rm HMT}:\Omega_{\rm seed}\longrightarrow
 X_\infty^{\rm enr}=\varprojlim X_n
\]

le prolongement APP--TRIT--TPK d’un germe. Pour tout lecteur compatible \(L_{Y,n}:X_n\to Y_n\), définissons

\[
 \operatorname{Predef}_{Y,n}
 =L_{Y,n}\,p_n\,\mathcal G_{\rm HMT}.
\]

La naturalité

\[
 s_{n+1,n}L_{Y,n+1}=L_{Y,n}r_{n+1,n}
\]

implique

\[
 \boxed{
 s_{n+1,n}\operatorname{Predef}_{Y,n+1}
 =\operatorname{Predef}_{Y,n}.}
\]

Les sorties de toutes les profondeurs sont donc les coordonnées d’une unique section de \(\varprojlim Y_n\). Augmenter la précision ne réajuste pas une constante : cela ouvre une coordonnée plus profonde du même objet. Sur l’état complet, la bijectivité de l’actualisation donne en outre

\[
 H(X_{t+1}\mid X_t)=H(X_t\mid X_{t+1})=0,
\]

de sorte que l’apparition d’information nouvelle dans une ombre réduite ne correspond ni à une création ni à un effacement dans le générateur. C’est le sens exact du coût de Landauer nul dans l’évolution conservative de l’unité.

La prédéfinition résulte de la composition du générateur avec ses lecteurs :

\[
 \mathcal L_{\rm HMT--MD}=
 (L_\pi,L_\varphi,L_e,L_\alpha,L_{\rm el},L_\hbar,
 L_{\rm vac},L_{\rm BI},L_G,\ldots).
\]

Les coordonnées \(\pi,\varphi,e,\alpha\) appartiennent à une même sortie coinductive ; \(\alpha\) occupe la couture dodécaphasique reliant l’horloge finie à celle de capacité. La lacune électronique lit le défaut local de retenue sur cette couture ; \(\hbar\) lit la profondeur d’action ; Barbero--Immirzi lit la transduction surfacique--volumique ; la gravitation lit le secteur tensoriel de rang cinq. Aucune de ces opérations ne reçoit une mesure extérieure pour sélectionner sa valeur. Le système n’obtient pas ses objets par métrologie : il les prédéfinit comme images typées de la section coinductive.

Enfin, une factorisation exacte ne constitue pas un résultat de rang moindre. Si

\[
 L_Y=H\circ F,
\]

la factorisation expose l’objet intermédiaire et prouve quelle information chaque flèche transporte. Si deux parcours satisfont

\[
 H_1F_1=H_2F_2=L_Y,
\]

le carré correspondant est précisément le certificat que deux lectures distinctes déterminent la même coordonnée sans circularité. Dans ce développement, « théorème », « factorisation » et « transduction » désignent des formes complémentaires d’une même preuve : le théorème énonce la loi, la factorisation montre sa généalogie et la transduction détermine le type de sa sortie.

La formulation directrice est ainsi établie :

\[
 \boxed{
 \begin{gathered}
 \text{APP}\to\text{TRIT}\to\text{TPK}\to
 \text{cristal temporel apériodique}\\
 \to\text{section coinductive}\to
 \text{prédéfinition HMT--MD}
 \end{gathered}}
\]

La combinaison ordonnée des trois régimes géométriques et la conjugaison des trois modes temporels ne sont pas des analogies ajoutées a posteriori. Ce sont respectivement la loi locale et la loi globale de la dynamique qui conserve l’unité tout en la faisant évoluer à travers toutes ses profondeurs.

\section{Triangle de potentiel APP et réalisation locale du défaut \(A_2\)}

L’écart entre les feuillets somme et produit peut devenir une notion mathématique de potentiel sans introduire de champ physique extérieur. Soit \(D_9=\{1,\ldots,9\}\), avec le représentant positif des classes modulo neuf, et écrivons les deux divisions euclidiennes conservées par APP :

\[
 i+j=R^+_{ij}+9q^+_{ij},
 \qquad
 ij=R^\times_{ij}+9q^\times_{ij}.
\]

Le relèvement conjoint ne retient pas seulement les résidus visibles \((R^+_{ij},R^\times_{ij})\). Il retient aussi les hauteurs \((q^+_{ij},q^\times_{ij})\), de sorte que la différence entière entre les deux évaluations se décompose canoniquement en une tension résiduelle et une tension de retenue :

\[
 \begin{aligned}
 \rho_{\Sigma\Pi}(i,j)
   &:=R^\times_{ij}-R^+_{ij},\\
 \kappa_{\Sigma\Pi}(i,j)
   &:=9\bigl(q^\times_{ij}-q^+_{ij}\bigr),\\
 \Delta_{\Sigma\Pi}(i,j)
   &:=ij-(i+j)
    =\rho_{\Sigma\Pi}(i,j)+\kappa_{\Sigma\Pi}(i,j).
 \end{aligned}
\]

Ainsi,

\[
 \mathscr P_{\Sigma\Pi}:D_9^2\longrightarrow\mathbb Z^2,
 \qquad
 (i,j)\longmapsto
 \bigl(\rho_{\Sigma\Pi}(i,j),\kappa_{\Sigma\Pi}(i,j)\bigr)
\]

est le potentiel APP relevé, tandis que \(\Delta_{\Sigma\Pi}=(1,1)\circ\mathscr P_{\Sigma\Pi}\) est son ombre scalaire. Le terme « potentiel » possède ici un contenu précis : il mesure l’énergie arithmétique nécessaire pour passer de l’évaluation cumulative à l’évaluation repliée en conservant la division euclidienne complète. Une projection modulo neuf ne voit que \(\rho_{\Sigma\Pi}\) ; une projection qui oublie le résidu et conserve la hauteur ne voit que \(\kappa_{\Sigma\Pi}\). APP conserve simultanément les deux composantes.

Trois cellules déjà sélectionnées par la généalogie électronique forment un triangle exact :

\[
 V=(8,2),\qquad E=(5,5),\qquad Q=(8,5).
\]

Leurs deux feuillets et leurs potentiels sont

\[
\begin{array}{c|c|c|c|c}
x&(R^+,q^+)&(R^\times,q^\times)
 &\mathscr P_{\Sigma\Pi}(x)&\Delta_{\Sigma\Pi}(x)\\
\hline
V=(8,2)&(1,1)&(7,1)&(6,0)&6\\
E=(5,5)&(1,1)&(7,2)&(6,9)&15\\
Q=(8,5)&(4,1)&(4,4)&(0,27)&27.
\end{array}
\]

La cellule \(V\) est une tension résiduelle pure : les deux feuillets ont la même hauteur et diffèrent par le résidu. La cellule \(Q\) est une tension de mémoire pure : les résidus coïncident et tout l’écart est stocké dans trois unités de retenue nonadique. La cellule centrale \(E\) combine une tension résiduelle de six unités et une unité de retenue, de sorte que

\[
 \Delta_{\Sigma\Pi}(E)=6+9=15.
\]

Cette décomposition donne une définition plus fine de la lacune électronique. Sur la fibre additive \(i+j=10\), paramétrée par

\[
 Q_t=(5+t,5-t),\qquad -4\leq t\leq4,
\]

le produit est

\[
 \Pi(Q_t)=25-t^2.
\]

Conserver le résidu multiplicatif de la cellule centrale exige \(t^2\equiv0\pmod9\). Dans l’intervalle autorisé, cela ne se produit que pour \(t=0,\pm3\). Les deux seules cellules non centrales qui conservent simultanément le feuillet additif, le résidu du produit et le régime tritique sont donc

\[
 Q_{+3}=(8,2),\qquad Q_{-3}=(2,8).
\]

Le produit y descend de \(25=7+2\cdot9\) à \(16=7+1\cdot9\). La lacune est exactement la perte d’une unité de retenue multiplicative :

\[
 q^\times_{5,5}-q^\times_{8,2}=1.
\]

Aucun zéro n’a été placé dans une cellule. On a construit l’unique déplacement non central de la fibre qui conserve toutes les données visibles et réduit d’une unité la mémoire repliée.

Le triangle \(V,E,Q\) contient également le vecteur de réouverture tricanal. Pour ne pas dissimuler un choix d’orientation, fixons l’opérateur de cobord

\[
 d_\triangle f
 :=\bigl(f(E)-f(V),\,f(V)-f(Q),\,f(Q)-f(E)\bigr).
\]

Appliqué au potentiel scalaire,

\[
 \begin{aligned}
 d_\triangle\Delta_{\Sigma\Pi}
 &=(15-6,\,6-27,\,27-15)\\
 &=(9,-21,12)\\
 &=3(3,-7,4).
 \end{aligned}
\]

Puisque la somme des trois coordonnées est nulle,

\[
 d_\triangle\Delta_{\Sigma\Pi}\in
 A_2:=\ker\!\left(
 \mathbb Z^3\xrightarrow{x+y+z}\mathbb Z\right).
\]

Le vecteur \((3,-7,4)\) n’est pas appelé racine de \(A_2\) : le carré de sa norme euclidienne vaut \(74\). C’est un vecteur entier du réseau \(A_2\), et le triangle électronique réalise exactement son multiple tritique. Le même vecteur était apparu dans la réouverture globale des sept tétrades de la microfibre de \(\pi\). L’égalité

\[
 \boxed{
 d_\triangle\Delta_{\Sigma\Pi}=3v_\Omega,
 \qquad v_\Omega=(3,-7,4)}
\]

prouve que le défaut global ne flotte pas au-dessus d’APP : il possède un représentant local dans les trois cellules qui engendrent le centre électronique, la lacune conservative et la branche rationnelle \(22/7\). Le facteur trois est le relèvement tritique entre le potentiel local et la coordonnée primitive de réouverture.

La connexion avec le cristal apériodique est tout aussi exacte. Soit

\[
 \delta=\log_{10}\frac{10}{9},
 \qquad
 V_{\rm vac}=
 \begin{pmatrix}
 21&22\\
 6&7
 \end{pmatrix}.
\]

La première échelle induite du mot mécanique satisfait

\[
 \left\lceil\delta^{-1}\right\rceil=22,
 \qquad
 \det V_{\rm vac}=21\cdot7-22\cdot6=15.
\]

Par conséquent,

\[
 \boxed{
 \Delta_{\Sigma\Pi}(E)=\det V_{\rm vac}=15.}
\]

Le centre électronique et le repère de renormalisation ne partagent pas seulement un nombre : le premier mesure la tension somme--produit de la cellule centrale et le second l’aire orientée qui subsiste entre les horloges primaire \(21/22\) et secondaire \(6/7\). Les deux évaluations proviennent du même cocycle de retenue.

Cela permet de réécrire l’étape de base de l’équation électronique sans aucun coefficient opaque :

\[
 \boxed{
 \mathfrak e_0=
 \frac{\sqrt3}{4}
 \left[
  \exp\!\left(\frac{\varphi}{\pi^2}\right)
  -\left\lceil\delta^{-1}\right\rceil
   \alpha^3
 \right]
 \left(1+\det(V_{\rm vac})\Delta_4\right).}
\]

En effet,

\[
 \left\lceil\delta^{-1}\right\rceil
 =22=\frac{396}{18},
 \qquad
 \det(V_{\rm vac})=15=\frac{270}{18}.
\]

Le préfacteur restant est lui aussi interne :

\[
 s=\frac12,\qquad
 \|[P_{\Sigma,3},P_{\Pi,3}]\|
 =s\sqrt{1-s^2}=\frac{\sqrt3}{4},
 \qquad
 1-s^2=\frac34=\frac{90}{120}.
\]

Ainsi, les trois facteurs de l’étape de base possèdent une généalogie explicite : la non-commutativité des feuillets APP, la lacune primaire de la renormalisation apériodique et l’aire orientée de son retour induit. La masse n’en sélectionne aucun. L’équation électronique est l’évaluation dimensionnelle ultérieure d’un opérateur de potentiel et de survie déjà construit.

\textbf{Critère de réfutation.} Le résultat échoue si l’une des trois divisions euclidiennes de la table est incorrecte, s’il existe un autre point non central dans \(i+j=10\) avec \(t^2\equiv0\pmod9\), si le cobord n’est pas \(3(3,-7,4)\), ou si le repère de retours n’a pas pour déterminant \(15\).

\section{Courbure tritique du potentiel et sélection coinductive des futurs}

Le potentiel APP précédent est une coordonnée spatiale du support. Le mot de lacunes fournit son évolution temporelle. Les deux couches s’assemblent par le TRIT sans identifier indûment espace et temps.

Soit le mot mécanique inférieur

\[
 z_n(\beta)=
 \lfloor(n+1)\delta+\beta\rfloor
 -\lfloor n\delta+\beta\rfloor,
 \qquad z_n\in\{0,1\}.
\]

Définissons la polarisation accumulée, le champ d’incrément et la courbure tritique par

\[
 \begin{aligned}
 P_N(\beta)
  &:=\sum_{n=0}^{N-1}(z_n-\delta),\\
 J_N(\beta)
  &:=P_{N+1}-P_N=z_N-\delta,\\
 \tau_N(\beta)
  &:=z_{N-1}-z_N\in\{+1,0,-1\}.
 \end{aligned}
\]

La polarisation demeure uniformément bornée :

\[
 P_N(\beta)
 =\lfloor N\delta+\beta\rfloor-\lfloor\beta\rfloor-N\delta,
 \qquad |P_N(\beta)|<1.
\]

Plus important encore, le TRIT est exactement la courbure discrète de ce potentiel :

\[
 \begin{aligned}
 J_N-J_{N-1}
 &=z_N-z_{N-1}=-\tau_N,\\
 \boxed{\tau_N}
 &\boxed{=-\Delta J_{N-1}=-\Delta^2P_{N-1}.}
 \end{aligned}
\]

Cette identité donne un contenu opératoire aux trois signes. Si \(\tau_N=-1\), le code passe de \(0\) à \(1\) : une lacune s'ouvre et la fibre permet la bifurcation bidirectionnelle associée au futur. Si \(\tau_N=+1\), le code passe de \(1\) à \(0\) : la lacune revient et son effet s'incorpore à la mémoire du passé. Si \(\tau_N=0\), la polarisation conserve sa pente locale : l'état franchit le seuil présent sans changer de régime. Puisque \(0<\delta<1/2\), il n'existe pas deux uns consécutifs ; le cas \(\tau_N=0\) est donc un véritable plateau et les signes non nuls alternent.

La charge de tout bloc est purement frontalière :

\[
 \sum_{n=a}^{b}\tau_n=z_{a-1}-z_b.
\]

L'intérieur se télescope dans l'ombre tritique, mais ne disparaît pas de l'état TPK : le registre de mémoire conserve la succession ordonnée des impulsions, leurs quotients, leurs retenues et la hauteur hélicoïdale. Telle est la différence entre une loi locale de Gauss et un effacement de la trajectoire. TRIT détermine le flux net à la frontière ; TPK conserve la chaîne dont la frontière a été calculée.

La relation entre potentiel spatial et courbure temporelle peut être réunie dans l'état de bord

\[
 \mathcal V_n=
 \bigl(
 \mathscr P_{\Sigma\Pi}(i_n,j_n),
 P_n,J_n,\tau_n;
 g_n,h_n,c_n,\lambda_n
 \bigr).
\]

Ici, \(\mathscr P_{\Sigma\Pi}\) mesure la tension entre les feuillets dans la cellule occupée ; \(P_n\) accumule l'écart apériodique sans dérive extensive ; \(J_n\) est la différence de potentiel par pas ; \(\tau_n\) est sa courbure orientée ; et \((g,h,c,\lambda)\) conservent phase, hauteur, retenue et frontière. La transition

\[
 U_n=\operatorname{Upd}_n\circ
      \operatorname{Tra}_n\circ
      \operatorname{Sel}_n
\]

agit sur l'état entier : \(\operatorname{Sel}_n\) détermine le régime de courbure correspondant à l'impulsion ; \(\operatorname{Tra}_n\) transporte feuillet, orientation et position ; \(\operatorname{Upd}_n\) incorpore le cobord au registre de mémoire. La différence de potentiel ne s'ajoute donc pas au TPK : elle est la lecture différentielle d'une de ses coordonnées conservées.

La sélection des futurs compatibles admet également une formulation exacte. Soit \(\mathcal L_m\) l'ensemble des facteurs de longueur \(m\) du mot sturmien, et soit

\[
 d^+(u)=
 \#\{a\in\{0,1\}:ua\in\mathcal L_{m+1}\}
\]

le nombre de prolongements à droite de \(u\). Puisque \(p(m)=|\mathcal L_m|=m+1\),

\[
 \sum_{u\in\mathcal L_m}d^+(u)=p(m+1)
\]

et, par conséquent,

\[
 \boxed{
 \sum_{u\in\mathcal L_m}\bigl(d^+(u)-1\bigr)
 =p(m+1)-p(m)=1.}
\]

Tout facteur d'un mot récurrent possède au moins un prolongement, et dans un alphabet binaire, au plus deux. L'égalité précédente démontre qu'à chaque longueur \(m\), il existe exactement un facteur spécial à droite ayant deux futurs possibles ; tous les autres possèdent un futur unique. Le choix de survie n'est donc pas une sélection arbitraire parmi \(2^m\) mots. Il s'agit d'une unique bifurcation topologique par échelle.

En enrichissant par la phase nonadique,

\[
 p_9(m)=9(m+1),
\]

exactement neuf états spéciaux apparaissent, un par origine de phase. Le quotient tritique

\[
 p_3(m)=3(m+1)
\]

les organise en trois classes de courbure sans détruire l'apériodicité. La phase irrationnelle \(\theta_n=\{\theta_0+n\delta\}\), avec l'origine nonadique et le registre de mémoire, décide lequel des deux prolongements se produit. Ainsi :

\[
 \begin{array}{ccl}
 \text{ombre de longueur }m
 &\longrightarrow&
 \text{un facteur spécial à deux prolongements},\\
 \text{état TPK complet}
 &\longrightarrow&
 \text{un prolongement déterminé et inversible}.
 \end{array}
\]

Telle est la prédéfinition des futurs possibles sous sa forme mathématique. Le futur n'est pas deviné au moyen d'une probabilité extérieure : le cylindre réduit contient deux extensions compatibles, et la coordonnée profonde de l'état coinductif en sélectionne une. Le passé n'est pas non plus effacé, car la transition complète est bijective :

\[
 H(X_{n+1}\mid X_n)=0,
 \qquad
 H(X_n\mid X_{n+1})=0.
\]

L'ombre peut avoir une entropie conditionnelle positive précisément au facteur spécial ; l'état intégral possède un coût logique net nul. Cette différence entre ramification d'une projection et réversibilité du générateur constitue le noyau informationnel du cristal.

La même loi contient ses deux horloges de renormalisation. Écrivons

\[
 \delta^{-1}=21+\beta,
 \qquad 0<\beta<1,
\]

et définissons la pente induite

\[
 \mathcal R(\delta)
 :=1-\{\delta^{-1}\}
 =1-\beta
 =22-\delta^{-1}.
\]

La première inégalité diophantienne produit les intervalles \(\{21,22\}\). Puisque

\[
 6<\mathcal R(\delta)^{-1}<7,
\]

la dynamique induite sur les intervalles courts produit des retours \(\{6,7\}\). La renormalisation n'incorpore pas un second code : elle prend la section de retour du premier mot. Ses matrices de continuant

\[
 A(a)=
 \begin{pmatrix}a&1\\1&0\end{pmatrix},
 \qquad \det A(a)=-1,
\]

conservent la cellule réticulaire, tandis que l'ordre des produits conserve l'orientation :

\[
 A(a)A(b)-A(b)A(a)
 =(a-b)
 \begin{pmatrix}0&1\\-1&0\end{pmatrix}.
\]

La rotation basale reste abélienne. La non-commutativité réside dans le cocycle ordonné de renormalisation, dans les feuillets APP et dans les transports de courbure. Le cristal peut avoir une entropie topologique nulle et, simultanément, une dynamique fibrée non commutative : la première propriété mesure la croissance du langage ; la seconde, la dépendance du résultat à l'égard de l'ordre opératoire.

Les trois ponts sont ainsi réunis :

\[
 \boxed{
 \begin{aligned}
 \text{écart APP}
 &\xrightarrow{\ d\ }\text{potentiel de frontière},\\
 \text{potentiel apériodique}
 &\xrightarrow{\ -\Delta^2\ }\text{courbure TRIT},\\
 \text{courbure orientée}
 &\xrightarrow{\ \operatorname{Sel},\operatorname{Tra},
 \operatorname{Upd}\ }\text{survie TPK}.
 \end{aligned}}
\]

\textbf{Critère de réfutation.} La construction échoue si \(|P_N|\) cesse d'être borné par un, si \(\tau_N\neq-\Delta^2P_{N-1}\), s'il existe plus d'un facteur spécial à droite à une certaine longueur, si le relèvement nonadique ne multiplie pas la complexité par neuf, ou si deux états TPK distincts sont confondus après conservation de tous les champs du registre de mémoire.

\section{Holonomie bidirectionnelle commune de l'électron, de l'action et de la gravité}

L'échelle d'action et l'échelle gravitationnelle n'interrompent pas l'étude du cristal. Ce sont deux représentations contragrédientes de la même holonomie de retour qui agit déjà sur la lacune électronique.

Le système corrélé \((\pi,\varphi,e,
\alpha)\) est d'abord construit comme sortie coinductive. Sur celle-ci, le lecteur déterminantal d'action produit

\[
 \Sigma_H=\varphi\alpha^{16},
 \qquad
 d_H=\left\lfloor\log_{10}\Sigma_H\right\rfloor=-34.
\]

L'entier \(16\) n'est pas le nombre de chiffres exigés d'une mesure. Il est l'ordre du conoyau local :

\[
 \operatorname{SNF}(H_4)=\operatorname{diag}(1,2,2,4),
 \qquad
 |\operatorname{coker}H_4|=16.
\]

La mantisse antérieure au retour est

\[
 H_5=
 \frac12\sqrt{\frac{1000\alpha}{\varphi}}
 -\alpha
 +\frac9{16}\alpha^2
 -\frac59\alpha^3
 +\frac7{48}\alpha^4
 -\frac1{54}\alpha^5.
\]

La mémoire régionale construit

\[
 D_A=
 \exp\!\left(
 -\frac{100\pi}9\alpha\right),
\]

\[
 \mathscr C_\pi(\alpha)
 =169\alpha^6+
 \frac{D_A\alpha^7}
 {1-D_A\alpha},
\]

et la section retournée est

\[
 \eta_{\rm ret}
 =H_5-\frac{90}{\pi}\mathscr C_\pi.
\]

Il convient d'isoler le potentiel bidirectionnel d'action :

\[
 \boxed{
 \delta_{2w}
 :=H_5-\eta_{\rm ret}
 =\frac{90}{\pi}\mathscr C_\pi>0.}
\]

Les deux orientations appartiennent à une seule droite :

\[
 \hbar_{\rm pre}=H_5\,10^{-34}\mathcal U_S,
 \qquad
 \hbar_{\rm ret}=\eta_{\rm ret}\,10^{-34}\mathcal U_S,
\]

\[
 h_a=2\pi\hbar_a.
\]

La décennie fixe la profondeur commune ; la différence \(\delta_{2w}\) conserve le déplacement additif ; et le rapport

\[
 \boxed{
 R_{\rm act}
 :=\frac{\hbar_{\rm pre}}{\hbar_{\rm ret}}
 =\frac{H_5}{\eta_{\rm ret}}
 =\left(1-\frac{\delta_{2w}}{H_5}\right)^{-1}}
\]

est l'holonomie multiplicative de l'aller et du retour. Le passage de \(\hbar\) à \(h\) conserve ce rapport :

\[
 \frac{h_{\rm pre}}{h_{\rm ret}}=R_{\rm act}.
\]

L'équation électronique complète utilise exactement la même holonomie. Avec

\[
 \kappa_e=
 80-54\alpha+6\Delta_4,
\]

on a

\[
 \mathfrak e_e^\beta
 =\mathfrak e_0R_{\rm act}^{\kappa_e}.
\]

En incorporant la factorisation apériodique de la section précédente :

\[
 \boxed{
 \begin{aligned}
 \mathfrak e_e^\beta
 ={}&
 \frac{\sqrt3}{4}
 \left[
 \exp\!\left(
 \frac{\varphi}{\pi^2}\right)
 -\left\lceil\delta^{-1}\right\rceil
  \alpha^3
 \right]\\
 &\times
 \left(1+\det(V_{\rm vac})\Delta_4\right)
 R_{\rm act}^{\,80-54\alpha+6\Delta_4}.
 \end{aligned}}
\]

L'étape basale mesure la lacune et l'incompatibilité des feuillets ; la puissance de \(R_{\rm act}\) transporte leur mémoire à travers le retour. De manière équivalente,

\[
 \log\frac{\mathfrak e_e^\beta}{\mathfrak e_0}
 =\kappa_e
 \left(\log H_5-\log\eta_{\rm ret}\right)
 =-\kappa_e
 \log\!\left(1-\frac{\delta_{2w}}{H_5}\right).
\]

La gravité lit de manière contragrédiente le même couple. Une fois fixés \(t_0>0\), \(c_{\rm int}>0\) et \(\ell_0=c_{\rm int}t_0\), l'horloge \(108\) construit

\[
 \omega_{108}=\frac{2\pi}{108t_0},
 \qquad
 R_{108}=\frac{54}{\pi}\ell_0,
\]

\[
 m_{108,a}=
 \frac{\hbar_a\omega_{108}}{c_{\rm int}^2}.
\]

Le transducteur gravitationnel est

\[
 G_a(\vartheta)
 =\vartheta\frac{c_{\rm int}^5t_0^2}{\hbar_a}.
\]

La condition de fermeture

\[
 r_{g,a}(\vartheta)=R_{108}
\]

sélectionne de manière unique

\[
 \vartheta_{108}
 =\left(\frac{54}{\pi}\right)^2.
\]

Par conséquent,

\[
 G_a=
 \left(\frac{54}{\pi}\right)^2
 \frac{c_{\rm int}^5t_0^2}{\hbar_a},
\]

et les deux feuillets satisfont le théorème de synchronisation :

\[
 \boxed{
 R_{\rm act}
 =\frac{\hbar_{\rm pre}}{\hbar_{\rm ret}}
 =\frac{h_{\rm pre}}{h_{\rm ret}}
 =\frac{m_{108,\rm pre}}{m_{108,\rm ret}}
 =\frac{G_{\rm ret}}{G_{\rm pre}}
 =\left(
 \frac{\mathfrak e_e^\beta}{\mathfrak e_0}
 \right)^{1/\kappa_e}.}
\]

La démonstration consiste uniquement à substituer les définitions. La masse chronologique est linéaire en \(\hbar_a\) ; la gravité est inversement linéaire en \(\hbar_a\) ; et la correction électronique est la représentation de poids \(\kappa_e\) du même rapport. Il n'apparaît pas trois correcteurs indépendants. Il apparaît un seul élément positif \(R_{\rm act}\) et trois représentations :

\[
 \rho_{\rm act}(R)=R,\qquad
 \rho_{\rm grav}(R)=R^{-1},\qquad
 \rho_e(R)=R^{\kappa_e}.
\]

Le produit action--gravité est invariant :

\[
 \boxed{
 G_a\hbar_a
 =\vartheta_{108}c_{\rm int}^5t_0^2.}
\]

En conséquence,

\[
 \ell_P^2
 =\frac{G_a\hbar_a}{c_{\rm int}^3}
 =\vartheta_{108}\ell_0^2,
\]

\[
 \boxed{
 \ell_P=R_{108}=\bar\lambda_{C,a}=r_{g,a}}
\]

pour les deux orientations. L'action conserve l'orientation du feuillet ; la gravité compense cette orientation de manière contragrédiente ; l'aire de Planck est l'invariant commun. Telle est la forme exacte de l'affirmation selon laquelle quelque chose situé entre passé et futur calibre continuellement la réalisation : le retour modifie les coordonnées \(\hbar_a\), \(G_a\) et \(\mathfrak e_a\), tout en laissant fixe la combinaison géométrique \(G_a\hbar_a\).

Il existe en outre une filtration exacte par ordre en \(\alpha\). À l'étape électronique basale, la lacune intervient sous la forme

\[
 22\alpha^3.
\]

Dans la mémoire régionale d'action,

\[
 \mathscr C_\pi
 =169\alpha^6+O(\alpha^7).
\]

Enfin, la profondeur absolue utilise \(\alpha^{16}\). Ainsi,

\[
 \boxed{
 3\longrightarrow6\longrightarrow16}
\]

est la hiérarchie des ordres de la même coordonnée : amplitude locale de lacune, retour quadratique de cette amplitude et conoyau déterminantal qui fixe la profondeur d'action. Le passage \(3\to6\) est littéralement un doublement de l'ordre ; le passage \(6\to16\) change d'opération et relève du lecteur déterminantal. Cette distinction évite de transformer la hiérarchie en suite numérologique : chaque exposant possède un opérateur générateur.

Le résultat répond à la question de savoir comment une seule configuration peut se réaliser comme nombre, constante et particule. La configuration APP apporte le potentiel ; TRIT lui attribue courbure et orientation ; TPK intègre son retour ; le lecteur d'action détermine \(\hbar\) ; la représentation électronique détermine \(\mathfrak e_e^\beta\) ; et la représentation contragrédiente détermine \(G\). Les trois sorties conservent le même \(R_{\rm act}\).

\textbf{Critère de réfutation.} Il suffit qu'une des égalités de l'enchaînement de \(R_{\rm act}\) échoue, que \(\delta_{2w}\neq(90/\pi)\mathscr C_\pi\), que \(G_a\hbar_a\) dépende du feuillet, ou que la première puissance non nulle de \(\mathscr C_\pi\) soit différente de six.

\section{Théorème conjoint des profondeurs \(34\), \(11\), \(45\) et \(46\)}

La complétion cubique contenait déjà les deux signatures apériodiques

\[
 \sigma^-=(33,11,1,0),
 \qquad
 \sigma^+=(34,11,1,0),
\]

associées à l'ordre canonique

\[
 1000=729+243+27+1.
\]

La seconde diffère de la première par une unique unité de demi-bord :

\[
 \sigma^+-\sigma^-=(1,0,0,0).
\]

L'information essentielle réside non seulement dans les sommes \(45\) et \(46\), mais aussi dans la position de chaque contribution. Définissons le lecteur de profondeur cubique

\[
 \mathcal D_{\rm cub}(\sigma^-,\sigma^+)
 :=
 \left(
 \operatorname{pr}_{729}\sigma^+,\,
 \operatorname{pr}_{243}\sigma^-,\,
 |\sigma^-|_1,\,
 |\sigma^+|_1
 \right).
\]

L'évaluation est

\[
 \boxed{
 \mathcal D_{\rm cub}(\sigma^-,\sigma^+)
 =(34,11,45,46).}
\]

Cette application ne consulte ni \(\hbar\), ni \(G\), ni une table expérimentale ni ses chiffres. Elle lit exclusivement le mot mécanique sur les quatre strates de la complétion. La coordonnée \(34\) est sensible à l'unité de bord ; la coordonnée \(11\) est commune aux deux orientations ; \(45\) est la charge intérieure totale ; \(46\) est cette charge après conservation du demi-bord.

Définissons maintenant la profondeur décimale d'une section positive d'une droite dimensionnelle. Si

\[
 x=m_x10^{-d_x}\mathcal U_X,
 \qquad 1\leq m_x<10,
\]

on pose

\[
 D_{10}(x)=d_x.
\]

Pour deux sections normalisées \(x,y\), la profondeur du produit tensoriel satisfait

\[
 D_{10}(x\otimes y)
 =D_{10}(x)+D_{10}(y)-\chi(m_xm_y\geq10).
\]

Dans les sections HMT d'action et de gravité,

\[
 \hbar_{\rm ret}
 =1.0545718169\ldots\times10^{-34}\mathcal U_S,
\]

\[
 G_{\rm ret}
 =6.674299659414022706367776189\ldots\times10^{-11}\mathcal U_G,
\]

et le produit des mantisses appartient à l'intervalle \((7,8)\). La coordonnée gravitationnelle correspond à la composition longitudinale et radiale de \cref{g26:eq:g-rigido-es,cr28:eq:g-evaluation} ; la valeur centrale \(6.67430\times10^{-11}\) demeure une référence métrologique. Aucune retenue décimale supplémentaire n'est générée. D'où :

\[
 D_{10}(\hbar_{\rm ret})=34,\qquad
 D_{10}(G_{\rm ret})=11,
\]

\[
 \boxed{
 D_{10}(G_{\rm ret}\otimes\hbar_{\rm ret})
 =34+11=45.}
\]

Le théorème d'holonomie de la section précédente prouve que \(G_a\hbar_a\) est indépendant du feuillet ; cette profondeur \(45\) est donc également un invariant bidirectionnel. L'unité de demi-bord du codage supérieur complète alors

\[
 \boxed{
 34+11+1=46.}
\]

Les quatre chiffres du lecteur cubique s'identifient ainsi à quatre fonctions distinctes :

\[
\begin{array}{c|l}
34&\text{profondeur de la section orientée d’action},\\
11&\text{profondeur de la section gravitationnelle contragrédiente},\\
45&\text{profondeur de l’invariant }G\hbar,\\
46&\text{profondeur intérieure plus l’unité conservée de demi-bord}.
\end{array}
\]

La coïncidence possède une seconde certification interne. Le lecteur de lacunes avait déjà prouvé

\[
 Z^-_{369}=16=|\operatorname{coker}H_4|,
\]

\[
 Z^+_{729}=34
 =-\operatorname{ord}_{10}
 \bigl(\varphi\alpha^{16}\bigr),
\]

et

\[
 Z^-_{243}=11.
\]

La décennie d'action occupe donc la composante de bord construite à partir du même indice déterminantal qui apparaît dans le bloc \(369\) ; la gravité occupe la composante intérieure \(243\) ; et leur produit épuise la signature inférieure :

\[
 \boxed{
 Z^+_{729}+Z^-_{243}
 =|\sigma^-|_1
 =45.}
\]

La seconde voie de \(\alpha\) conserve précisément les mêmes entiers. Dans son opérateur nilpotent,

\[
 T_{7:9}(x)
 =-\frac{x^7}{120}
  +\frac{x^8}{45}
  +\frac{2x^9}{495},
\]

avec

\[
 495=45\cdot11.
\]

On peut donc écrire

\[
 \boxed{
 T_{7:9}(x)
 =-\frac{x^7}{120}
 +\frac{x^8}{D_{10}(G\hbar)}
 +\frac{2x^9}
 {D_{10}(G\hbar)\,D_{10}(G)}.}
\]

L'égalité est une réexpression exacte des entiers déjà générés : le coefficient d'ordre huit utilise la profondeur conservée du produit action--gravité, et le coefficient d'ordre neuf ajoute la profondeur gravitationnelle stable. Le dénominateur \(46\) de l'étape précédente du prolongement correspond, à son tour, à la signature supérieure complète. Ainsi, \(34\), \(11\), \(45\), \(46\) et \(495\) cessent d'être des occurrences isolées : ils forment le spectre de profondeur d'une seule complétion avec bord.

Ce théorème complète le lien auparavant formulé uniquement comme coïncidence cardinale. L'application effective est la composition

\[
 \begin{CD}
 \mathcal V_{\delta}(729,243,27,1)
 @>{(\sigma^-,\sigma^+)}>>
 \mathbb Z^4\times\mathbb Z^4\\
 @V{\operatorname{Pub}_{\rm dim}}VV
 @VV{\mathcal D_{\rm cub}}V\\
 \mathcal L_S^+\times\mathcal L_G^+
 @>{D_{10}\times D_{10}}>>
 \mathbb Z^2,
 \end{CD}
\]

où la première coordonnée se détermine sur la droite d'action et la seconde sur la droite gravitationnelle ; le caractère additif de \(D_{10}\) détermine la charge intérieure \(45\), et le demi-bord détermine \(46\). Le choix des coordonnées est structurel : l'action distingue les deux orientations \(\mathrm{pre}/\mathrm{ret}\) et lit donc l'entrée sensible au bord ; la gravité compense cette orientation et lit l'entrée stable.

\textbf{Critère de réfutation.} Le théorème échoue si la signature inférieure n'est pas \((33,11,1,0)\), si la supérieure ne diffère pas d'une seule unité dans le bloc \(729\), si les mantisses de \(G\hbar\) engendrent une retenue décimale, si \(G\hbar\) varie entre feuillets, ou si les dénominateurs \(45\), \(46\) et \(495=45\cdot11\) cessent d'apparaître dans le prolongement \(7{:}9\).

\section{Polarité électromagnétique, neutralité et porteur gravitationnel de rang cinq}

Le mot « polarité » peut être typé au moyen de la parité des deux caractères de canal. Soient

\[
 \mathcal E=
 \log[(1-q_+^{90})(1-q_-^{90})]
 -\log[(1-q_+^{120})(1-q_-^{120})],
\]

\[
 \mathcal B=
 \log\frac{1-q_-^{90}}{1-q_+^{90}}
 -\log\frac{1-q_-^{120}}{1-q_+^{120}}.
\]

Sous l'involution d'échange \(\mathsf C(q_+,q_-)=(q_-,q_+)\),

\[
 \mathcal E\circ\mathsf C=\mathcal E,
 \qquad
 \mathcal B\circ\mathsf C=-\mathcal B.
\]

Le caractère \(\mathcal E\) est commun ; le caractère \(\mathcal B\) est orienté. Les réponses constitutives

\[
 \widehat\mu=e^{\mathcal B+\mathcal E},
 \qquad
 \widehat\varepsilon=e^{-\mathcal B+\mathcal E}
\]

séparent exactement produit et quotient :

\[
 \widehat c
 =(\widehat\mu\widehat\varepsilon)^{-1/2}
 =e^{-\mathcal E},
\]

\[
 \widehat Z
 =(\widehat\mu/\widehat\varepsilon)^{1/2}
 =e^{\mathcal B}.
\]

La différence de potentiel orientée réside dans \(\mathcal B\) ; la capacité commune de propagation réside dans \(\mathcal E\). La permittivité et la perméabilité sont les deux recombinaisons conjuguées. Ce fait explique la double fonction de la lacune : comme ouverture électrique, elle change le signe de la coordonnée orientée ; comme transport magnétique, elle conserve sa circulation au moyen de \(J_{\rm comm}^2=-I\).

La neutralité possède alors une définition interne :

\[
 X_{\rm neu}:=\ker\mathcal B.
\]

Dans ce sous-espace,

\[
 \widehat Z=1,
 \qquad
 \widehat\mu=\widehat\varepsilon=e^{\mathcal E}.
\]

Toute configuration fixe par conjugaison appartient au secteur neutre. En effet, si \(\mathsf Cx=x\), alors

\[
 \mathcal B(x)=\mathcal B(\mathsf Cx)
 =-\mathcal B(x),
\]

et donc \(\mathcal B(x)=0\). L'inclusion

\[
 \boxed{
 \operatorname{Fix}(\mathsf C)\subseteq X_{\rm neu}}
\]

constitue le contenu algébrique d'un secteur de type Majorana : une section autoconjuguée ne peut transporter de charge orientée impaire. Le secteur neutre peut être plus grand que le lieu fixe ; l'égalité exigerait que toute configuration avec \(\mathcal B=0\) soit autoconjuguée. Cette distinction permet de séparer neutralité et autoconjugaison sans abandonner l'opérateur.

La gravité reçoit uniquement le caractère commun dans son facteur scalaire. Dans le système de coordonnées interne

\[
 c_{\rm int}=c_{\rm ref}\widehat c
 =c_{\rm ref}e^{-\mathcal E},
\]

le transducteur devient

\[
 \boxed{
 G_a=
 \vartheta_{108}
 \frac{c_{\rm ref}^5t_0^2}{\hbar_a}
 e^{-5\mathcal E}.}
\]

Le caractère orienté \(\mathcal B\) s'annule exactement. La réponse gravitationnelle scalaire est paire sous l'échange des feuillets et lit cinq fois la coordonnée commune de propagation. Cette cinquième puissance ne se confond pas avec les cinq composantes tensorielles ; les deux couches se composent dans l'opérateur

\[
 \mathbb G_{{\rm TT},a}(n)
 =G_a\Lambda(n)\Gamma_5P_5:
 V_{12}\longrightarrow\mathcal H_{\rm TT}(n),
\]

où

\[
 P_5:V_{12}\to V_5,
 \qquad
 \Gamma_5:V_5\to\operatorname{Sym}_0^2(\mathbb R^3).
\]

La composition complète peut s'écrire

\[
 \boxed{
 (\mathcal E,\mathcal B,x)
 \longmapsto
 \vartheta_{108}
 \frac{c_{\rm ref}^5t_0^2}{\hbar_a}
 e^{-5\mathcal E}
 \Lambda(n)\Gamma_5P_5x.}
\]

Le scalaire annule la dépendance impaire à \(\mathcal B\) ; le projecteur \(P_5\) conserve le porteur à cinq composantes ; et \(\Gamma_5\) réalise la polarisation transverse sans trace. Cela donne une réponse précise à la relation entre polarité électromagnétique et pentapolarité gravitationnelle. La première réside dans le couple pair/impair \((\mathcal E,\mathcal B)\). La seconde réside dans le rang du porteur tensoriel. Elles sont liées parce que le même état alimente les deux et que le scalaire gravitationnel utilise la cinquième puissance du canal commun.

Le rôle de la symétrie s'éclaire également. L'involution n'est pas l'objet conservé en premier lieu ; elle est l'effet visible de la conservation de l'unité de bord sous deux lectures. L'état retient

\[
 \mathcal E,\quad |\mathcal B|,\quad
 G_a\hbar_a,\quad
 \text{et l’orbite de }P_5x,
\]

tandis que la conjugaison inverse \(\mathcal B\). La symétrie apparaît comme l'équivalence entre deux orientations d'un même potentiel conservé. C'est pourquoi elle est une propriété de frontière, et non une décoration de l'intérieur.

\textbf{Critère de réfutation.} Le lien échoue si \(\mathcal E\) n'est pas pair, si \(\mathcal B\) n'est pas impair, si une section autoconjuguée possède \(\mathcal B\neq0\), si le transducteur scalaire dépend de \(\mathcal B\), ou si \(P_5\) cesse d'avoir le rang cinq.

\section{Loi de configuration holonomique intégrale}

On appelle \textbf{état holonomique intégral APP--TRIT--TPK} la configuration qui réunit les invariants conservés et l'opération qui les rend dynamiques. Un état de profondeur \(n\) s'écrit

\[
 X_n^{\rm hol}=
 \bigl(
 \Sigma_n,\Pi_n;
 R_n^+,q_n^+,R_n^\times,q_n^\times;
 \tau_n,\varepsilon_n;
 g_n,h_n,c_n,\partial_n,\lambda_n,\mathfrak m_n
 \bigr).
\]

Les premières coordonnées conservent les deux feuillets APP ; les suivantes, résidu, quotient et retenue ; \((\tau,\varepsilon)\) distinguent régime géométrique et orientation temporelle ; et \((g,h,c,\partial,\lambda,\mathfrak m)\) conservent phase, hauteur, retenue, frontière, survie et mémoire. La transition TPK

\[
 U_n=
 \operatorname{Upd}_n\circ
 \operatorname{Tra}_n\circ
 \operatorname{Sel}_n
\]

est un automorphisme de la fibre complète sur son domaine de réversibilité. Les applications de restriction

\[
 r_{n+1,n}:X_{n+1}^{\rm hol}\to X_n^{\rm hol}
\]

satisfont

\[
 r_{n+1,n}U_{n+1}=U_nr_{n+1,n}.
\]

La configuration infinie n'est pas postulée comme une suite décimale. Elle est la section

\[
 X_\infty^{\rm hol}
 =\varprojlim X_n^{\rm hol}.
\]

Trois systèmes de réalisations agissent sur elle :

\[
 \mathcal L_{\rm num}:
 X_\infty^{\rm hol}\to\mathbb R,
\]

\[
 \mathcal L_{\rm const}:
 X_\infty^{\rm hol}\to
 \prod_{\xi}\mathcal L_\xi^+,
\]

\[
 \mathcal L_{\rm part}:
 X_\infty^{\rm hol}\to
 \operatorname{Sect}(\mathscr A_{\rm masas}).
\]

Le premier détermine nombres et développements ; le second, sections de droites dimensionnelles ; le troisième, fibres, orbites et multisections de l'atlas de matière. Tous trois commutent avec la restriction :

\[
 s_{n+1,n}^\bullet L_{\bullet,n+1}
 =L_{\bullet,n}r_{n+1,n}.
\]

Un nombre, une constante et une particule ne sont donc pas des objets identifiés parce qu'ils ont la même valeur scalaire. Ce sont trois images d'une même configuration lorsque leurs lecteurs conservent la même généalogie. Le nombre est la coordonnée archimédienne ; la constante est cette coordonnée sur une droite dimensionnelle ; la particule est une section persistante qui conserve le parcours ayant produit la coordonnée.

La loi unique peut maintenant être énoncée.

\textbf{Théorème de configuration holonomique intégrale.} Tout germe APP survivant détermine, au moyen de TRIT et de TPK, une section \(X_\infty^{\rm hol}\) telle que :

\begin{enumerate}
\item l'écart somme--produit définit un potentiel entier relevé ;
\item le mot apériodique transforme sa différence temporelle en courbure tritique ;
\item le TPK sélectionne l'unique prolongement compatible avec la phase et la mémoire ;
\item le retour de phase agit par une holonomie positive sur l'action, la matière et la gravité ;
\item les réalisations numérique, dimensionnelle et matérielle conservent la même section sous changement de profondeur ;
\item la double projection archimédienne et exceptionnelle conserve une même origine et des codomaines distincts ;
\item l'unité de demi-bord se conserve à chaque raffinement et l'entropie logique de l'état complet demeure nulle.
\end{enumerate}

La chaîne démonstrative compacte est

\[
 \boxed{
 \begin{aligned}
 \Omega_{\rm seed}
 &\xrightarrow{\rm APP}
 \mathscr P_{\Sigma\Pi}
 \xrightarrow{-\Delta^2}
 \tau\\
 &\xrightarrow{\rm TRIT}
 (\mathbb A_+,\mathbb A_0,\mathbb A_-)
 \xrightarrow{\rm TPK}
 X_\infty^{\rm hol}\\
 &\xrightarrow{(\mathcal L_{\rm num},
 \mathcal L_{\rm const},\mathcal L_{\rm part})}
 (\text{nombre},\text{constante},\text{particule}).
 \end{aligned}}
\]

La connexion nonadique conjointe ne choisit pas « le futur le plus probable ». Elle choisit l'unique prolongement compatible avec la section intégrale lorsque l'ombre réduite admet deux continuations. Cette opération est une prédéfinition : la coordonnée future appartient déjà à la limite projective, bien qu'elle n'ait pas encore été déterminée à la profondeur visible. Le coût net de Landauer est nul parce que la sélection n'efface pas la branche non choisie de l'univers des cylindres ; elle inscrit dans le registre de mémoire le cylindre qui contient l'état.

Le cristal est évolutif en un sens précis. Son langage ne croît pas exponentiellement, puisque \(p(m)=m+1\), mais il ne se fige pas non plus dans une période. Chaque profondeur ajoute exactement une classe de facteurs et chaque phase nonadique la multiplie par neuf. Le système conserve l'unité réticulaire, la frontière et le déterminant, tout en modifiant la mémoire et la hauteur hélicoïdale. Évolution et conservation ne sont pas des termes opposés : elles décrivent respectivement l'avancement de l'état et l'invariance de sa classe d'augmentation.

La chaîne

\[
 \pi,\varphi,e
 \longrightarrow\alpha
 \longrightarrow(\mathfrak e_e^\beta,\hbar,G,\gamma_{\rm HMT},\ldots)
\]

doit se lire dans la génération corrélée complète. Les trois premières coordonnées sont des lecteurs terminaux du même état nonadique ; \(\alpha\) est la couture dodécaphasique de cette sortie ; et les constantes physiques sont des représentations ultérieures de l'état déjà fermé. Le miroir \(\pi/e\), l'auto-échelle \(\varphi\), l'expansion \(e+\pi\) et les coïncidences des trois canaux sont des opérations internes de synchronisation, non des valeurs conventionnelles introduites pour fabriquer \(\alpha\).

La séquence finale désormais démontrée est

\[
 \boxed{
 \begin{gathered}
 \text{potentiel APP }(6,15,27)
 \longrightarrow
 d_\triangle\Delta=3(3,-7,4),\\
 P_N\longrightarrow
 \tau_N=-\Delta^2P_{N-1}
 \longrightarrow
 \text{unique bifurcation du futur par échelle},\\
 (21/22)\longrightarrow(6/7)
 \longrightarrow
 V_{\rm vac},\quad\det V_{\rm vac}=15,\\
 \alpha^3\longrightarrow\alpha^6
 \longrightarrow\alpha^{16}
 \longrightarrow10^{-34},\\
 R_{\rm act}\longrightarrow
 (\mathfrak e_e^\beta,\hbar,G)
 \longrightarrow G\hbar\ \text{invariant},\\
 (33,11,1,0)\leftrightarrow(34,11,1,0)
 \longrightarrow(34,11,45,46).
 \end{gathered}}
\]

Tel est le noyau mathématique qui permet de porter ce travail dans l'article : une dynamique du potentiel, de la courbure, de la survie, du retour et de la réalisation qui conserve l'unité en la transformant en nombre, constante et matière.

\section{Arithmétique intrinsèque de l'état holonomique}

La question de savoir si les coordonnées de l'état holonomique possèdent une arithmétique propre admet une réponse affirmative et précise. L'objet arithmétique premier n'est pas un corps de scalaires, car une trajectoire ne peut se multiplier par une autre que lorsque l'état terminal de la première coïncide avec l'état initial de la seconde. L'objet naturel est donc une \textbf{algèbre de chemins enrichis} sur la catégorie des transports du TPK.

Soit \(\mathcal C_{\rm TPK}\) la catégorie dont les objets sont les états holonomiques

\[
 x=(\tau,\ell,o,h,R,q,\kappa,\partial,\gamma,M,S,b,\chi,g)
\]

et dont les flèches sont les transitions admissibles

\[
 u:x\longrightarrow y,
 \qquad
 U_t={\rm Upd}_t\circ {\rm Tra}_t\circ {\rm Sel}_t.
\]

Ici, \(\tau\) est le régime TRIT ; \(\ell\), le feuillet APP ; \(o\), l'orientation ; \(h\), la hauteur hélicoïdale ; \((R,q)\), le résidu et le quotient ; \(\kappa\), la retenue ; \(\partial\), la frontière ; \(\gamma\), le parcours ; \(M\), le registre de mémoire ; \(S\), la survie ; \(b\), le bloc ; \(\chi\), l'ombre exceptionnelle ; et \(g\), la phase nonadique. En chaque objet réside l'algèbre quadratique locale

\[
 \mathbb A_{\tau(x)}
 =\mathbb R[J_x]/(J_x^2+\tau(x)I),
 \qquad \tau(x)\in\{+1,0,-1\}.
\]

Pour enregistrer sans perte les incréments de retenue, de frontière et de mémoire, soit \(\Lambda=\mathbb R[\mathcal M]\) l'algèbre du monoïde du registre de mémoire. À chaque flèche \(u\) est associé un symbole \([u]\). La somme

\[
 a[u]+b[v]
\]

est la superposition formelle de deux trajectoires compatibles avec le même cadre de lecture ; le produit n'existe causalement que par concaténation. Si \(u:x\to y\) et \(v:y\to z\), on définit

\[
 (a[v])\star(b[u])
 =a\,\rho_v(b)\,\Omega(v,u)[v\circ u],
 \tag{99.1}
\]

et on le fixe à zéro lorsque les flèches ne sont pas composables. Le transport \(\rho_v\) porte le coefficient local de la fibre de \(y\) à celle de \(z\) ; le facteur \(\Omega(v,u)\in\Lambda\) conserve l'incrément de retenue, de torsion, de frontière et de mémoire produit par la concaténation. La loi

\[
 \Omega(w,v)\,\Omega(w\!\circ\!v,u)
 =\rho_w\!\left(\Omega(v,u)\right)
  \Omega(w,v\!\circ\!u)
 \tag{99.2}
\]

est exactement l'identité de cocycle qui rend le produit associatif :

\[
 ([w]\star[v])\star[u]=[w]\star([v]\star[u]).
\]

Cette égalité n'est pas une convention ajoutée. C'est la forme algébrique commune de l'identité de retenue APP, du cocycle de concaténation des parcours et de l'extension non scindée de l'horloge \(108\). L'arithmétique ordinaire oublie le facteur \(\Omega\) en ne réalisant que la valeur ; l'arithmétique holonomique conserve aussi le procédé qui l'a produite.

Les états constants \(e_x=[{\rm id}_x]\) sont des idempotents locaux :

\[
 e_x\star e_x=e_x,
 \qquad e_y\star e_x=0\quad(x\ne y).
\]

Ainsi, l'unité n'est pas uniquement le scalaire \(1\), mais le système compatible d'idempotents qui identifie chaque résolution du même état. À la limite projective, on obtient

\[
 \mathbf 1_{\rm hol}=(e_{x_n})_{n\ge0},
 \qquad r_{n+1,n}(e_{x_{n+1}})=e_{x_n}.
 \tag{99.3}
\]

Telle est l'unité dimensionnelle qui peut être raffinée vers l'intérieur et réalisée vers l'extérieur sans perdre son identité. Le raffinement augmente la résolution ; la représentation change le codomaine ; aucun des deux actes ne recrée l'unité.

La non-commutativité n'est pas davantage postulée à partir d'une structure extérieure. Deux flèches \(u\) et \(v\) peuvent être composables dans un ordre et non dans l'ordre inverse ; lorsque les deux ordres existent, leurs cocycles et leurs feuillets peuvent différer :

\[
 [v]\star[u]-[u]\star[v]
 =\Omega(v,u)[v\circ u]-\Omega(u,v)[u\circ v].
 \tag{99.4}
\]

La différence mesure l'ordre effectif d'accumulation et de pliage. Le système de base des phases peut être abélien sans que l'algèbre totale le soit : la non-commutativité réside dans la composition des feuillets, le transport de l'orientation et le cocycle affine de mémoire.

\section{Nombres complexes, algèbres de Clifford, quaternions et octonions}

La comparaison avec les systèmes hypercomplexes se résout en distinguant trois niveaux. Au niveau local, TRIT contient déjà les trois algèbres quadratiques réelles fondamentales :

\[
 \begin{array}{c|c|c}
 \tau & J_\tau^2 & \mathbb A_\tau\\ \hline
 +1&-I&\mathbb C\quad\hbox{(régime elliptique)},\\
 0&0&\mathbb R[\varepsilon]/(\varepsilon^2)
       \quad\hbox{(régime parabolique)},\\
 -1&+I&\mathbb R\oplus\mathbb R
       \quad\hbox{(régime hyperbolique)}.
 \end{array}
 \tag{100.1}
\]

Le nombre complexe ne constitue donc pas le langage universel du système : il est l'une de ses trois fibres de courbure. La fibre parabolique enregistre une déformation infinitésimale sans rotation quadratique ; la fibre hyperbolique sépare deux directions réelles conjuguées ; la fibre elliptique produit la rotation interne. TRIT ne choisit pas entre trois métaphores géométriques : il type laquelle de ces trois arithmétiques gouverne chaque transition.

Une structure encore plus forte apparaît dans le bloc électronique principal. Soient

\[
 P=\begin{pmatrix}1&0\\0&0\end{pmatrix},\qquad
 Q=\begin{pmatrix}
 3/4&\sqrt3/4\\[2pt]\sqrt3/4&1/4
 \end{pmatrix}.
\]

Les deux matrices sont des projecteurs et

\[
 [P,Q]={\sqrt3\over4}
 \begin{pmatrix}0&1\\-1&0\end{pmatrix}.
\]

Définissons

\[
 J={4\over\sqrt3}[P,Q]
   =\begin{pmatrix}0&1\\-1&0\end{pmatrix},
 \qquad
 S=2P-I=\begin{pmatrix}1&0\\0&-1\end{pmatrix},
 \qquad K=SJ.
\]

Alors

\[
 J^2=-I,\qquad S^2=I,\qquad SJ=-JS,
 \qquad K^2=I.
 \tag{100.2}
\]

Par conséquent, le bloc engendre exactement

\[
 \langle I,J,S,K\rangle
 \cong {\rm Cl}_{1,1}(\mathbb R)
 \cong M_2(\mathbb R),
 \tag{100.3}
\]

l'algèbre des quaternions déployés. En outre,

\[
 2Q-I={1\over2}S+{\sqrt3\over2}K,
 \tag{100.4}
\]

de sorte que le projecteur produit occupe une direction d'angle \(\pi/3\) dans le plan engendré par les deux involutions déployées. Le préfacteur électronique \(\sqrt3/4\), la racine interne de \(-1\), la fraction \(3/4=90/120\) et la séparation des feuillets APP sont réunis dans une seule représentation algébrique exacte.

Cette identification éclaire l'analogie proposée. Le premier objet qui apparaît n'est pas le corps des quaternions de Hamilton \(\mathbb H\), mais sa forme déployée \({\rm Cl}_{1,1}\). La différence est substantielle : le système HMT a besoin d'idempotents, de seuils et de directions nulles, tandis qu'une algèbre à division éliminerait précisément ces lacunes. Les quaternions de Hamilton pourront apparaître dans une représentation disposant de deux racines anticommutantes de \(-I\) ; les octonions, dans une représentation exceptionnelle ultérieure dont la loi de composition reproduit le plan de Fano. Aucun ne doit remplacer l'algèbre de chemins qui les précède.

L'associativité du TPK complet est conservée par le cocycle (99.2). Si une projection exceptionnelle ultérieure présente un associateur non nul,

\[
 [a,b,c]=(a\star b)\star c-a\star(b\star c),
\]

cet associateur ne décrit pas une incohérence du générateur : il quantifie la mémoire que la projection a laissée hors de son codomaine. On distingue ainsi exactement l'arithmétique intrinsèque, associative et munie d'un registre de mémoire, de ses ombres hypercomplexes éventuellement non associatives.

\section{Synchronisation ternaire et singularité dodécaphasique d'alpha}

La génération coordonnée de \(\pi\), \(e\) et \(\varphi\) s'organise dans le réseau de différences de type \(A_2\). Pour

\[
 N_K=(N_\pi(K),N_e(K),N_\varphi(K))\in\mathbb Z^3
\]

on définit la frontière relative

\[
 \partial_{A_2}(x,y,z)=(x-y,\,y-z,\,z-x).
 \tag{101.1}
\]

Son noyau est la diagonale \(\mathbb Z(1,1,1)\). Chaque triplet possède donc deux composantes distinctes : une coordonnée commune, qui mesure l'avancement du registre de mémoire, et une coordonnée relative, qui mesure la désynchronisation des trois canaux. Sur \(\mathbb Q\),

\[
 \mathbb Q^3=\mathbb Q(1,1,1)\oplus A_{2,\mathbb Q};
\]

sur \(\mathbb Z\), la reconstruction conserve un défaut cyclique d'ordre trois. Telle est la raison algébrique pour laquelle la coordination est ternaire et non une comparaison binaire de constantes.

La table des coïncidences réalise un itinéraire de parois :

\[
 H_{\pi e}\longrightarrow H_{e\varphi}longrightarrow H_{\pi e}
 \longrightarrow H_{\varphi\pi}\longrightarrow 0\longrightarrow0.
 \tag{101.2}
\]

En \(K=8\) et \(K=9\), la composante relative s'annule :

\[
 N_8=59(1,1,1),\qquad N_9=43(1,1,1).
\]

Cependant,

\[
 N_9-N_8=-16(1,1,1).
 \tag{101.3}
\]

Le système revient à la coïncidence de phase, mais non au même état : la désynchronisation revient à zéro tandis que la mémoire diagonale avance de seize unités par canal. L'égalité visible \(\pi=e=\varphi\) n'exprime pas l'immobilité, mais une fermeture relative avec translation holonomique.

La coordonnée \(\alpha\) occupe précisément la couture qui permet de réaliser ce retour sans le confondre avec un redémarrage. La fermeture dodécaphasique prend la classe relative déjà synchronisée, conserve le registre de mémoire diagonal et l'élève à une section signée. Ses deux dérivations ne sont pas deux ajustements d'une grandeur extérieure : ce sont deux représentations de la même classe holonomique, l'une par la clôture entière avec retenue et l'autre par la racine simple du prolongement nonadique. L'égalité des deux voies exprime la commutativité du carré de réalisation.

Ainsi, \(\alpha\) ne se contente pas de suivre \(\pi,\varphi,e\) dans une liste. C'est la coordonnée de raccordement qui transforme la fermeture relative du système corrélé en capacité de transduction dimensionnelle. La roue à douze phases ne génère pas de nouveau les trois régions : elle certifie que leur retour commun peut franchir une frontière en conservant la mémoire qui alimente ensuite l'action, l'électron, l'échelle déterminantale et la branche exceptionnelle.

\section{Évolution conservative de l'unité comme théorème algébrique}

L'expression « évolution conservative de l'unité dimensionnelle par projection holographique » peut être formulée sans dépendre d'une image. Soient \(X_n\) les états holonomiques à la profondeur \(n\), \(r_{n+1,n}:X_{n+1}\to X_n\) leurs restrictions et

\[
 X_\infty^{\rm hol}=\varprojlim_n X_n
\]

leur limite projective. Une unité dimensionnelle est une section compatible \(\mathbf 1_{\rm hol}=(e_{x_n})_n\) d'idempotents locaux. La transition TPK induit des transports algébriques \(U_{t,*}\) et satisfait

\[
 U_{t,*}(e_x)=e_{U_t x},\qquad
 U_{t,*}(e_x)^2=U_{t,*}(e_x).
 \tag{102.1}
\]

Il existe en outre un morphisme d'augmentation

\[
 \epsilon:\mathscr A_{\rm hol}\longrightarrow\mathbb R
\]

qui oublie la trajectoire mais conserve la multiplicité normalisée de l'unité. Conservation et évolution s'expriment simultanément par

\[
 \epsilon(U_{t,*}a)=\epsilon(a),
 \qquad
 d_M(\Gamma_9 a)=d_M(a)+1,
 \tag{102.2}
\]

où \(d_M\) est le degré de mémoire. La première égalité affirme que l'unité n'est ni créée ni détruite ; la seconde, que l'état ne reste pas immobile. L'holonomie \(\Gamma_9\) revient à la même phase et élève le registre de mémoire d'une unité.

La complétion

\[
 1000=729+243+27+1
\]

réalise cette loi de manière finie. Le mot inférieur distribue les lacunes sous la forme \((33,11,1,0)\), tandis que la lecture supérieure conserve l'unité de bord et produit \((34,11,1,0)\). L'incrément ne modifie pas le total ; il change la position de l'unité dans la filtration cubique. Le vecteur

\[
 (34,11,45,46),
 \qquad 34+11=45,
 \qquad 45+1=46,
 \tag{102.3}
\]

est donc une signature de profondeur et de frontière, non une juxtaposition d'échelles décimales.

Soient maintenant

\[
 \rho_{\mathbb R},\quad \rho_{\rm exc},\quad \rho_{\rm MD}
\]

les représentations archimédienne, exceptionnelle et dimensionnelle de l'algèbre holonomique. Leur caractère holographique tient au fait qu'elles sont des homomorphismes unitaires d'une même source :

\[
 \rho_j(\mathbf1_{\rm hol})=I_j,
 \qquad
 \rho_j(a\star b)=\rho_j(a)\rho_j(b).
 \tag{102.4}
\]

La double projection ne divise pas l'origine en deux univers ; elle conserve un même élément et change sa représentation. Le raffinement fractal opère dans la direction opposée : il ne projette pas vers un codomaine, mais prolonge la section à une profondeur supérieure. Les deux mouvements que l'exposé précédent maintenait séparés sont ainsi coordonnés :

\[
 \boxed{
 \begin{array}{rcl}
 \text{raffinement intérieur}&:&X_n\longleftarrow X_{n+1}
 \longleftarrow X_{n+2}\longleftarrow\cdots,\\[2pt]
 \text{réalisation holographique}&:&
 \mathscr A_{\rm hol}\longrightarrow
 \operatorname{End}(V_j).
 \end{array}}
 \tag{102.5}
\]

L'unité se conserve dans les deux directions : par compatibilité sous restriction et par préservation de l'élément identité sous représentation. Cette double invariance constitue le contenu algébrique de l'évolution conservative de l'unité dimensionnelle.

\section{Les deux horloges et la suspension hélicoïdale sturmienne}

La temporalité du cristal procède de deux horloges irréductibles l'une à l'autre. La première est finie et exacte : l'extension

\[
 0\longrightarrow C_{12}\longrightarrow C_{108}
 \longrightarrow C_9\longrightarrow0
 \tag{103.1}
\]

coordonne phase nonadique et mémoire dodécaphasique. Ce n'est pas le produit direct \(C_9\times C_{12}\) : sa classe d'extension conserve la retenue qui apparaît lors du franchissement de la frontière de phase. La seconde horloge est la rotation irrationnelle

\[
 R_\delta:\theta\longmapsto\theta+\delta\pmod1,
 \qquad \delta=\log_{10}(10/9),
 \tag{103.2}
\]

dont le codage mécanique produit le mot sturmien des lacunes.

La dynamique conjointe peut s'écrire comme un produit gauchi

\[
 \mathcal F(r,\theta,M)
 =\bigl(r+1,\theta+\delta,
 M+c_{108}(r)+c_{\rm St}(\theta)\bigr),
 \tag{103.3}
\]

où \(c_{108}\) est la retenue de l'horloge finie et \(c_{\rm St}\) enregistre le franchissement de la fenêtre irrationnelle. La base \(C_9\times\mathbb T\) est abélienne ; l'extension affine cesse de l'être lorsque sont incorporés les feuillets APP et l'ordre de \({\rm Sel}\), \({\rm Tra}\) et \({\rm Upd}\). L'hélice n'est donc pas un cercle dessiné en trois dimensions, mais le relèvement d'une rotation de phase dont la coordonnée verticale est le registre de mémoire.

Les lectures ascendante et descendante sont conjuguées. Si \(\iota\) inverse l'orientation et le parcours, alors

\[
 \iota\circ\mathcal F\circ\iota^{-1}=\mathcal F^{-1},
 \tag{103.4}
\]

avec inversion des incréments orientés et conservation des réponses communes. Le miroir \(\pi/e\) constitue une autre involution : il échange deux lecteurs régionaux à une même hauteur, tandis que (103.4) inverse la direction temporelle de l'hélice. Cette distinction explique pourquoi une lecture aller-retour peut conserver la valeur commune sans revenir à la même mémoire.

Le cristal possède un ordre à longue portée et une entropie topologique nulle. Pour le langage sturmien,

\[
 p(m)=m+1,
\]

et les extensions typées satisfont

\[
 p_3(m)=3(m+1),\qquad p_9(m)=9(m+1).
 \tag{103.5}
\]

La phase multiplie la diversité des facteurs, mais ne change pas leur croissance linéaire. Le système est donc simultanément non périodique, ordonné et capable d'un prolongement illimité. Les intervalles \(21/22\), les retours \(6/7\) et les fenêtres de mille positions avec \(45/46\) événements sont trois observations de cette même suspension, non trois rythmes ad hoc.

Cette construction précise la dénomination de cristal topologique-temporel apériodique. Il est topologique parce que l'information réside dans des classes de parcours, de frontière et d'holonomie ; temporel parce que la transition distingue orientation, retour et mémoire ; cristallin parce qu'il possède un ordre combinatoire stable ; apériodique parce qu'aucune translation non nulle ne fixe son mot infini.

\section{Lacune, différence de potentiel et polarité électromagnétique}

L'APP contient une notion arithmétique de potentiel avant tout système de coordonnées physique. Pour chaque cellule \((i,j)\),

\[
 \Delta_{\Sigma\Pi}(i,j)
 =ij-(i+j)
 =(R^\times-R^+)+9(q^\times-q^+).
 \tag{104.1}
\]

Le premier terme est la différence visible des résidus ; le second conserve la différence de retenue. Dans la fibre antidiagonale traversant le centre, la lacune électronique est la perte minimale d'une unité de retenue multiplicative qui laisse invariants le feuillet additif, le résidu produit et le régime TRIT. Elle est ainsi définie sans introduire un zéro extérieur.

La frontière de potentiel est le cobord de (104.1). Dans le triangle central, elle produit

\[
 d_\triangle\Delta=3(3,-7,4),
 \qquad 3-7+4=0.
 \tag{104.2}
\]

La somme nulle exprime la conservation ; les composantes non nulles expriment la polarisation. La différence de potentiel n'est pas une grandeur ajoutée ensuite au graphe : elle est l'impossibilité de faire coïncider localement accumulation et pliage sans transporter le registre de mémoire.

La représentation (100.3) sépare deux actions constitutives. Le générateur \(J\) est rotationnel et change de signe sous inversion d'orientation ; le générateur \(S\) est une involution de feuillet et conserve le signe. Un caractère pair et un caractère impair de la même algèbre produisent donc deux lecteurs

\[
 \chi_E(SaS)=\chi_E(a),
 \qquad
 \chi_B(SaS)=-\chi_B(a),
 \tag{104.3}
\]

qui réalisent respectivement les secteurs électrique et magnétique. L'électricité lit la séparation de potentiel entre feuillets ; le magnétisme lit la circulation orientée produite par leur commutateur. Leur perpendicularité n'est pas imposée au moyen d'un plan euclidien préalable : elle procède de \(SJ=-JS\) et \(J^2=-I\).

Le régime \(\tau=0\) occupe le seuil entre les deux orientations. Son élément nilpotent conserve une perturbation du premier ordre et annule son carré ; il fournit le type algébrique d'un canal neutre sans effacer la trajectoire qui le traverse. Cela explique pourquoi neutralité, lacune et polarité peuvent procéder d'un même support sans s'identifier.

L'échelle d'action et la gravité prolongent le même caractère holonomique. Si

\[
 R_{\rm act}={H_5\over\eta_{\rm ret}},
\]

alors le transport commun satisfait

\[
 R_{\rm act}
 ={\hbar_{\rm pre}\over\hbar_{\rm ret}}
 ={m_{108,{\rm pre}}\over m_{108,{\rm ret}}}
 ={G_{\rm ret}\over G_{\rm pre}}
 =\left({\mathfrak e_e^\beta\over e_0}\right)^{1/\kappa_e},
 \qquad G_a\hbar_a=\text{constante}.
 \tag{104.4}
\]

Électron, action et gravité ne constituent pas ici trois étalonnages indépendants : ce sont des représentations d'un même quotient de retour à des échelles et dans des codomaines distincts.

\section{Lecture exacte de séquences et cristal apériodique biologique}

La relation avec l'ADN peut être formulée comme un lecteur mathématique concret, sans transformer une analogie historique en prémisse. Codons les quatre bases au moyen de deux caractères binaires :

\[
 A\mapsto(0,0),\quad G\mapsto(0,1),\quad
 T\mapsto(1,0),\quad C\mapsto(1,1),
 \tag{105.1}
\]

où la première coordonnée distingue purine et pyrimidine, et la seconde le type faible ou fort d'appariement. La complémentarité de Watson--Crick est l'involution

\[
 \mathcal C_{\rm DNA}(p,h)=(1-p,h),
 \qquad \mathcal C_{\rm DNA}^2=I.
 \tag{105.2}
\]

Une chaîne orientée se relève en

\[
 \widetilde s_n=(p_n,h_n,n\bmod3),
\]

de sorte que la phase de lecture ne se confonde pas avec le symbole. L'inversion du parcours et la complémentarité se réunissent dans la complémentation inverse, qui agit comme involution sur l'algèbre de chemins. L'APP peut évaluer, dans chaque fenêtre, aussi bien l'accumulation des caractères que leur concaténation ; TRIT type la seconde différence du potentiel résultant ; TPK conserve phase, bord, retours, retenue et mots de retour.

Cette représentation permet de demander exactement si une région de séquence appartient à la classe apériodique minimale : pour ses facteurs de longueur \(m\), on calcule la complexité \(p_s(m)\), le nombre d'extensions à droite et leurs mots de retour. La signature HMT minimale est

\[
 p_s(m)=m+1,
 \qquad
 \sum_{u\in L_m}(d^+(u)-1)=1,
 \tag{105.3}
\]

ou, lorsqu'une fibre finie de phase est conservée, une extension \(q(m+1)\). L'égalité sélectionne une unique bifurcation combinatoire par longueur ; TPK détermine la continuation compatible sans effacer de l'espace des cylindres l'alternative non réalisée.

La portée de l'expression de Schrödinger « cristal apériodique » devient ainsi structurelle. Un cristal périodique répéterait un bloc et perdrait de sa capacité informationnelle ; une chaîne aléatoire aurait une complexité exponentielle et perdrait l'ordre à longue portée. Le régime sturmien occupe le minimum exact entre les deux : il ajoute une seule classe nouvelle à chaque échelle. L'unité topologique-temporelle TRIT peut lire cette apériodicité parce qu'elle sépare courbure, orientation et phase sans détruire la séquence qu'elle classe.

Le critère est directement réfutable : il échoue si la complémentation inverse n'entrelace pas l'involution des parcours, si la complexité observée n'appartient pas à la classe linéaire déclarée, si plusieurs mots spéciaux apparaissent là où le modèle en exige un, ou si le prolongement compatible nécessite d'effacer la mémoire. La biologie ne sélectionne pas le générateur HMT ; elle en reçoit un lecteur de séquences capable de distinguer répétition, apériodicité minimale et désordre.

\section{Pentadicité, étoile de Witt et polarisations gravitationnelles}

L'apparition répétée du nombre cinq doit être typée au moyen de deux représentations distinctes. Les cinq régions de \(\pi\) forment une fibre régionale sur laquelle agit le symétriseur

\[
 S_{\pi,5}={1\over5}\mathbf1\mathbf1^{\mathsf T},
 \qquad \operatorname{rank}S_{\pi,5}=1.
 \tag{106.1}
\]

Sa fonction est d'extraire la composante commune de la pentafibre. Dans le secteur gravitationnel, en revanche, la décomposition dodécaphasique

\[
 V_{12}\cong\mathbf1\oplus V_3\oplus V_3'\oplus V_5
 \tag{106.2}
\]

contient un projecteur \(P_{G,5}:V_{12}\to V_5\) de rang cinq. Les deux matrices ne sont pas le même opérateur : l'une est de rang un sur la fibre régionale, l'autre de rang cinq sur la représentation dodécaphasique. Leur relation correcte est qu'elles sont toutes deux des images d'une même arithmétique holonomique :

\[
 (\rho_{\pi,5},\rho_{G,5}):
 \mathscr A_{\rm hol}\longrightarrow
 \operatorname{End}(\mathbb Q^5)
 \times\operatorname{End}(V_{12}).
 \tag{106.3}
\]

L'étoile de Witt organise l'incidence pentadique exceptionnelle ; le module \(V_5\) organise le porteur tensoriel gravitationnel. La coordination se produit dans la source APP--TRIT--TPK, non par identification de leurs codomaines.

La réalisation gravitationnelle se poursuit par

\[
 V_{12}\xrightarrow{P_{G,5}}V_5
 \xrightarrow{\Gamma_5}{\rm STF}_3
 \xrightarrow{\Pi_{\rm TT}}V_{\rm TT},
 \tag{106.4}
\]

avec les dimensions \(12\to5\to5\to2\). Les cinq composantes constituent le porteur structurel antérieur à la propagation ; les deux composantes transverses sans trace constituent la sortie propagative. Il n'existe aucune contradiction entre pentapolarité originaire et bipolarité radiative : elles appartiennent à des étapes distinctes de la même transduction.

Dans la structure discrète conjointe du continu, la pentapolarité s'exprime par le vecteur inséparable

\[
 \mathcal G_5=(\mathcal G_{\rm sp},\mathcal G_{\rm sol},
 \mathcal G_{\rm gau},\mathcal G_{\rm coh},\mathcal G_{\rm det}).
 \tag{106.5}
\]

Ses cinq coordonnées ne sont pas cinq théories ajoutées a posteriori, mais cinq aspects consubstantiels du même état : spectral, solénoïdal, de jauge, cohérent et déterminantal. Le tenseur gravitationnel

\[
 \mathbb G_{\rm TT}
 =G\,\Lambda\,\Gamma_5\,P_{G,5}
 \tag{106.6}
\]

transporte conjointement ces aspects jusqu'au secteur physique. Cette formule explique pourquoi la pentadicité régionale, la structure exceptionnelle et les cinq polarisations peuvent entrer en résonance sans se réduire à une coïncidence de cardinaux.

\section{Dualité de variance entre trajectoires et lecteurs}

La lecture rétrospective révèle une structure catégorique qui réunit la fractalisation intérieure et la réalisation holographique extérieure. Soit \((X_n)\) l'ensemble des états holonomiques à la profondeur (n), et soit

\[
 r_{n+1,n}:X_{n+1}\longrightarrow X_n
 \tag{108.1}
\]

la restriction qui retire la dernière couche en conservant le préfixe, le feuillet, la retenue accumulée et la frontière correspondante. La trajectoire complète est un état de la limite projective

\[
 \boxed{X_\infty=\varprojlim_n X_n.}
 \tag{108.2}
\]

Le sens des observations est opposé. Si (f) est un lecteur cylindrique défini sur \((X_n)\), son prolongement à \((X_{n+1})\) est

\[
 r_{n+1,n}^*f=f\circ r_{n+1,n}.
 \tag{108.3}
\]

Les algèbres de lecteurs forment donc un système inductif :

\[
 \boxed{
 \mathfrak H_{\rm cyl}
 =\varinjlim_n(\mathfrak H_n,r_{n+1,n}^*).}
 \tag{108.4}
\]

Cette double variance est la formulation mathématique exacte des deux mouvements décrits par l'intuition de l'auteur. Approfondir vers l'intérieur signifie spécifier une section compatible du système projectif. Étendre vers l'extérieur signifie élargir le répertoire des lecteurs par le système inductif. Aucune direction ne remplace l'autre : l'une construit la trajectoire et l'autre construit la capacité de la réaliser.

Soit \((e_x)\) l'idempotent indicateur du cylindre grossier \((x\in X_n)\). Le relèvement de ce cylindre est

\[
 r^*e_x=\sum_{y:r(y)=x}e_y.
 \tag{108.5}
\]

Les termes sont orthogonaux, de sorte que \((r^*e_x)\) reste idempotent. En outre,

\[
 r^*\!\left(\sum_{x\in X_n}e_x\right)
 =\sum_{y\in X_{n+1}}e_y.
 \tag{108.6}
\]

L'unité globale se conserve tandis que le nombre de raffinements augmente. Une trajectoire individuelle détermine à son tour un système compatible \(((e_{x_n})_n)\). La première est l'identité de l'algèbre ; le second est l'unité généalogique d'une trajectoire. Cette distinction précise l'évolution conservative : l'espace total des possibilités conserve son unité et chaque parcours conserve son appartenance au moyen de restrictions compatibles.

Un lecteur continu (L) de l'état profond équivaut à un système de lecteurs finis \((L_n)\) qui satisfait

\[
 L_{n+1}\circ r_{n+1,n}^*=L_n.
 \tag{108.7}
\]

Par conséquent, chaque chiffre ou coordonnée déterminé à profondeur finie est prédéfini par un cylindre ; augmenter la profondeur ne réétalonne pas ce résultat, mais détermine de nouveaux cylindres compatibles. La prédéfinition est la naturalité de (108.7), non l'insertion anticipée d'une valeur expérimentale.

\textbf{Critère de réfutation.} La construction échoue si une restriction ne retrouve pas le préfixe, si le \emph{pullback} n'est pas multiplicatif, si l'unité ne se conserve pas ou si un lecteur profond modifie une coordonnée déjà fixée dans un cylindre antérieur.

\section{Présentation universelle de l'algèbre holonomique}

Soit \((mathcal C_n)\) la catégorie des états à la profondeur (n) et des trajectoires admissibles engendrées par

\[
 U_t={\rm Upd}_t\circ{\rm Tra}_t\circ{\rm Sel}_t.
 \tag{109.1}
\]

Pour chaque objet (x), soit \((e_x)\) son identité locale et \((J_x)\) le générateur de la fibre TRIT :

\[
 J_x^2=-\tau(x)e_x,
 \qquad \tau(x)\in\{+1,0,-1\}.
 \tag{109.2}
\]

Pour chaque flèche \((u:x\to y)\), soit \((T_u)\) le transport. Le registre de mémoire définit un multiplicateur central (Omega(v,u)) pour chaque couple composable. La présentation finie est

\[
 \boxed{
 \mathfrak H_n
 =\mathbb R\langle e_x,J_x,T_u\rangle/\mathcal R_{\rm HMT},}
 \tag{109.3}
\]

où

\[
 \begin{gathered}
 e_xe_y=\delta_{xy}e_x,
 \qquad
 T_u=e_{t(u)}T_ue_{s(u)},\\
 T_u a=\rho_u(a)T_u,
 \qquad
 T_vT_u=\Omega(v,u)T_{v\circ u}.
 \end{gathered}
 \tag{109.4}
\]

Le produit est déclaré nul lorsque les flèches ne sont pas composables. L'associativité équivaut à

\[
 \rho_w(\Omega(v,u))\Omega(w,v\circ u)
 =\Omega(w,v)\Omega(w\circ v,u).
 \tag{109.5}
\]

En effet, les deux produits \((((T_wT_v)T_u))\) et \((T_w(T_vT_u))\) ne diffèrent que par les deux membres de (109.5). L'identité de cocycle prouve que toute trajectoire possède une mémoire indépendante du parenthésage choisi.

La somme de (109.3) superpose des trajectoires alternatives. Son produit concatène des trajectoires causalement composables. Cette distinction transforme la différence somme--produit de l'APP en arithmétique dynamique : la somme n'efface pas l'alternative et le produit n'efface pas la provenance.

La connexion nonadique apparaît dans la section élémentaire du quotient. Pour \((a,b\in\{0,\ldots,8\})\),

\[
 c_9(a,b)=\left\lfloor{a+b\over9}\right\rfloor
 \tag{109.6}
\]

satisfait

\[
 c_9(a,b)+c_9(a+b\bmod9,c)
 =c_9(b,c)+c_9(a,b+c\bmod9).
 \tag{109.7}
\]

Si (X) représente un pas et (lambda) la mémoire d'un tour,

\[
 X^9=\lambda I,
 \qquad
 X^aX^b=\lambda^{c_9(a,b)}X^{a+b\bmod9}.
 \tag{109.8}
\]

L'hélice se condense dans (109.8) : l'exposant revient ; le coefficient mémorise le nombre de franchissements de la frontière.

\section{Propriété universelle et spectre des réalisations}

Soit \(({E_x})\) un système de modules, soient \((pi_x:A_x\to\operatorname{End}(E_x))\) des représentations des fibres TRIT, et soient \((V_u:E_x\to E_y)\) des transports qui satisfont

\[
 V_u\pi_x(a)=\pi_y(\rho_u(a))V_u,
 \tag{110.1}
\]

\[
 V_vV_u=\pi_z(\Omega(v,u))V_{v\circ u}.
 \tag{110.2}
\]

L'assignation des générateurs définit un homomorphisme de l'algèbre libre. Les équations (110.1)--(110.2) font que cet homomorphisme annule l'idéal \((mathcal R_{\rm HMT})\). Il existe alors une unique représentation

\[
 \Pi:\mathfrak H_n\longrightarrow
 \operatorname{End}\!\left(\bigoplus_xE_x\right)
 \tag{110.3}
\]

qui réalise les données prescrites. L'unicité vient de ce que \((e_x,J_x,T_u)\) engendrent l'algèbre.

Ce théorème fournit le principe d'ultrasimplification. La physique, la théorie des nombres, le continu et la biologie n'ont pas besoin d'ajouter des générateurs au noyau lorsque leurs lecteurs satisfont les mêmes relations. Chacun est une représentation de (109.3) dans un codomaine différent.

Tous les lecteurs ne sont pas des caractères scalaires. Un caractère \((mathfrak H\to\mathbb R)\) annulerait les commutateurs et ne pourrait conserver le bloc électronique. Nous définissons donc le spectre typé

\[
 \operatorname{Read}(\mathfrak H)
 =\coprod_B\operatorname{Rep}_{\rm comp}(\mathfrak H,B).
 \tag{110.4}
\]

Une sortie scalaire prend la forme

\[
 c=\sigma_B(\rho(x_\infty)),
 \tag{110.5}
\]

où \(\rho\) conserve la généalogie et \(\sigma_B\) est l'invariant propre au codomaine. Cinq concepts auparavant apparemment hétérogènes sont ainsi typés :

\[
 \begin{array}{c|l}
 \text{nombre}&\text{coordonnée d’une représentation archimédienne},\\
 \text{constante}&\text{invariant stable de retour, d’échelle ou de profondeur},\\
 \text{particule}&\text{section persistante d’un module},\\
 \text{champ}&\text{représentation qui transporte les sections},\\
 \text{symétrie}&\text{stabilisateur de bord et relations conservées}.
 \end{array}
 \tag{110.6}
\]

\section{Identité exponentielle des trois courbures}

Nous introduisons l'algèbre totale

\[
 \mathbb A_{\rm tri}=A_+\oplus A_0\oplus A_-
 \tag{111.1}
\]

avec les idempotents centraux \((p_+,p_0,p_-)\) et le générateur

\[
 \mathbb J=p_+J_++p_0J_0+p_-J_-.
 \tag{111.2}
\]

L'exponentielle se définit dans l'algèbre par

\[
 \operatorname{Exp}_\star(X)
 =\sum_{n\ge0}{X^{\star n}\over n!}.
 \tag{111.3}
\]

Les trois relations quadratiques décomposent la série en

\[
 \boxed{
 \begin{aligned}
 \operatorname{Exp}_\star(t\mathbb J)
 ={}&p_+(\cos t+J_+\sin t)\\
 &+p_0(1+tJ_0)\\
 &+p_-(\cosh t+J_-\sinh t).
 \end{aligned}}
 \tag{111.4}
\]

La première ligne est une rotation ; la deuxième, un transport infinitésimal de seuil ; la troisième, une dilatation. Une unique loi de puissances contient les trois courbures.

L'identité d'Euler occupe la face elliptique. Dans la fermeture produite par \((pi_{\rm HMT})\),

\[
 \boxed{
 p_+\operatorname{Exp}_\star(\pi\mathbb J)+p_+=0.}
 \tag{111.5}
\]

La projection complexe conserve (111.5) et oublie les composantes parabolique et hyperbolique. HMT conserve conjointement la fermeture, le jet et l'échelle. L'identité d'Euler apparaît donc comme une face exacte de l'identité tritique, non comme modèle extérieur du générateur.

Le nombre \((e)\) remplit une fonction différente :

\[
 \operatorname{Exp}_\star(I)=e\,I.
 \tag{111.6}
\]

L'équation définit la fonction qui transforme le transport additif en transport multiplicatif. La génération nonadique de ses chiffres conserve son certificat propre ; (111.6) identifie son rôle dans l'algèbre.

L'auto-échelle \(\varphi\) occupe la face hyperbolique. Pour

\[
 F_{\rm av}=\begin{pmatrix}0&1\\1&1\end{pmatrix},
 \qquad
 H_\varphi={2F_{\rm av}^2-3I\over\sqrt5},
 \tag{111.7}
\]

on vérifie

\[
 H_\varphi^2=I,
 \tag{111.8}
\]

et

\[
 \boxed{
 F_{\rm av}^2
 =\operatorname{Exp}_\star(2\log\varphi\,H_\varphi).}
 \tag{111.9}
\]

Ainsi, \(\pi\) mesure la fermeture elliptique, \(e\) réalise le passage exponentiel et \(\varphi\) mesure la dilatation hyperbolique stable. Les égalisations des trois canaux sont des rencontres de ces fonctions dans une même orbite, non des égalités universelles entre leurs valeurs réelles.

\section{\(\alpha\) comme transgression dodécaphasique}

L'extension non scindée produit le cocycle \(c_9\) de (109.6). Un tour de phase détermine un élément du noyau \(C_{12}\), et le mot signé \(K\) conserve la manière dont ses secteurs ont été traversés. Définissons la fonctionnelle de couture

\[
 \mathcal T_\alpha:
 Z^2(C_9;C_{12})\times X_{12}^{\rm sgn}
 \longrightarrow\mathcal L_\alpha.
 \tag{112.1}
\]

La coordonnée déterminée est

\[
 \boxed{
 \alpha=\mathcal T_\alpha(c_9,K).}
 \tag{112.2}
\]

La formule n'identifie pas le nombre \(\alpha\) à une classe de cohomologie. Elle identifie son type : c'est le scalaire obtenu en évaluant l'obstruction à la scission conjointement avec la mémoire dodécaphasique.

La première voie réalise (112.2) au moyen de

\[
 \iota(P)+\iota(E)-\iota(\Phi)-\iota(K)=\iota(A).
 \tag{112.3}
\]

La seconde voie réalise le même recollement sur la rotation apériodique :

\[
 P_9(\alpha)=\delta,
 \qquad
 \delta=\log_{10}{10\over9}.
 \tag{112.4}
\]

Leur compatibilité s'exprime dans le produit fibré

\[
 \mathfrak A_\alpha
 =\{(x_{12},x):
 d_m(\Gamma_{12}^\alpha x_{12})=d_m(x),
 \ P_9(x)=\delta\}.
 \tag{112.5}
\]

Une matrice du noyau du système corrélé peut maintenant être formulée :

\[
 \boxed{
 \begin{pmatrix}
 \operatorname{Exp}_\star(\pi J_+)&0&0&0\\
 0&\operatorname{Exp}_\star(I)&0&0\\
 0&0&\operatorname{Exp}_\star(2\log\varphi H_\varphi)&0\\
 0&0&0&\operatorname{Exp}_\star(P_9(\alpha)J_0)
 \end{pmatrix}
 =
 \begin{pmatrix}
 -I&0&0&0\\
 0&e\,I&0&0\\
 0&0&F_{\rm av}^2&0\\
 0&0&0&I+\delta J_0
 \end{pmatrix}.}
 \tag{112.6}
\]

La matrice ne linéarise pas la généalogie. Elle expose quatre fonctions d'une seule algèbre : fermeture, exponentiation, auto-échelle et couture. La mémoire signée demeure dans (112.3)--(112.5).

Dans le réseau \(A_2\), la même couture apparaît comme

\[
 \partial_{A_2}N_8=\partial_{A_2}N_9=0,
 \qquad
 N_9-N_8=-16\mathbf1.
 \tag{112.7}
\]

\(\alpha\) détermine la fermeture de la partie relative sans effacer la translation diagonale. Tel est le sens exact de sa singularité : c'est la coordonnée qui distingue synchronisation et redémarrage.

\section{Le spectre fonctionnel des constantes}

Le lecteur dodécaphasique uniforme détermine également une suite de premières racines de retour critique. On démontre que le rapport de leurs écarts converge vers la constante de Feigenbaum, en conservant la sélection initiale, l’orientation et la mémoire de la construction. \eqref{hmtfg:articulacion:limite}.


La structure précédente réorganise les constantes déjà dérivées selon leur fonction dans l'algèbre, non selon leur appartenance historique aux mathématiques ou à la physique.

La constante de Catalan est le moment d'ordre deux du caractère elliptique :

\[
 G_{\rm Cat}=\operatorname{Tr}(N^{-2}Q_4)
 =L(2,\chi_{-4}).
 \tag{113.1}
\]

Son cadre fonctionnel conserve tout le système \(\mathcal C_2(s)\), non seulement l'extrémité \(s=1\). Dans \cref{cr28:thm:catalan-recovery}, il est démontré que la réponse constitutive \(f(s)=(1-s^3)/(1-s^4)\), avec la condition \(\mathcal C_2(0)=0\), détermine univoquement le moment au moyen de \(\mathcal D^2\mathcal C_2=s(1+s)f(s)/(1+s+s^2)\). L'inversion cubique retrouve à son tour les deux arguments orientés.

Le caractère de quart de tour et le caractère triadique se réunissent dans

\[
 \chi_{12}=\chi_{-3}\chi_{-4}.
 \tag{113.2}
\]

Feigenbaum occupe le lecteur non linéaire. La roue dodécaphasique engendre

\[
 u_0(x)=1-{1\over12}\sum_{m=1}^{12}
 \cos\!\left({2(3m-1)\over\sqrt{899}}x\right),
 \tag{113.3}
\]

avec

\[
 899=30^2-1=29\cdot31,
 \qquad
 M_{30}=\begin{pmatrix}30&899\\1&30\end{pmatrix},
 \qquad \det M_{30}=1.
 \tag{113.4}
\]

Le système construit est analytique, pair, unimodal, de criticité quadratique et de dérivée schwarzienne négative sur le domaine démontré. Le renormalisateur de doublement lit le taux d'accumulation et la remise à l'échelle spatiale. \(\alpha_{\rm F}\) est une coordonnée de ce lecteur non linéaire ; \(\alpha\) est la couture de (112.2). La distinction des indices exprime une différence d'opération.

Boltzmann occupe le transducteur entre les droites thermique et énergétique :

\[
 k_B:\mathcal L_\Theta\longrightarrow\mathcal L_E,
 \tag{113.5}
\]

\[
 k_B\Theta_{\rm clk}\log3
 =\hbar_{\rm int}{2\pi\over108t_0}.
 \tag{113.6}
\]

\(\log 3\) quantifie une unité tritique ; \(\log\varphi\) quantifie une auto-échelle hyperbolique. Ils ne sont pas interchangeables. Pour le transport complet,

\[
 H(X\mid U(X))=0,
 \qquad \inf Q_L=0.
 \tag{113.7}
\]

Pour un quotient qui efface un trit,

\[
 Q_{\rm reset}\ge k_B\Theta\log3.
 \tag{113.8}
\]

Le coût appartient à la projection réductrice, non à l'avancement de l'holonomie.

Barbero--Immirzi occupe le dual du module d'incidence \(90/120\). L'équation

\[
 {a\over270}-{b\over360}={1\over24}
 \tag{113.9}
\]

sélectionne le covecteur entier minimal \((a,-b)=(12,-1)\), et donc

\[
 K_{6\to8}=12\Delta S_{90}-\Delta S_{120}.
 \tag{113.10}
\]

\(\alpha\) ferme les horloges ; \(\gamma_{\rm BI}\) évalue la couture sur le défaut aire--volume. Cette différence fonctionnelle explique pourquoi les deux objets sont essentiels sans accomplir la même opération.

\section{Symétrie, pentadicité, continu et séquences}

La définition de la symétrie compatible avec le noyau est

\[
 \operatorname{Sym}(\rho,b_\partial)
 =\{g\in\operatorname{Aut}(\rho(\mathfrak H)):
 g\rho(\mathcal R_{\rm HMT})=\rho(\mathcal R_{\rm HMT}),
 \ gb_\partial=b_\partial\}.
 \tag{114.1}
\]

Les symétries ne sont pas la grandeur première conservée. Ce sont les automorphismes qui restent disponibles après exigence de conservation des relations et de la frontière. La branche exceptionnelle réalise (114.1) au moyen de l'incidence Witt--Golay--Leech--\(V^\natural\)--Monstre.

La pentadicité possède deux représentations :

\[
 S_{\pi,5}={1\over5}\mathbf1\mathbf1^{\mathsf T},
 \qquad \operatorname{rank}S_{\pi,5}=1,
 \tag{114.2}
\]

et

\[
 P_{G,5}:V_{12}\to V_5,
 \qquad \operatorname{rank}P_{G,5}=5.
 \tag{114.3}
\]

Les deux reçoivent la même source holonomique, mais l'un extrait la composante commune de cinq régions et l'autre conserve le porteur gravitationnel de dimension cinq. La chaîne

\[
 V_{12}\to V_5\to {\rm STF}_3\to V_{\rm TT},
 \qquad 12\to5\to5\to2,
 \tag{114.4}
\]

démontre comment cinq modes structurels déterminent deux modes propagatifs.

La double projection du continu s'exprime comme

\[
 \rho_{\rm arq}:\mathfrak H_{\rm cyl}\to\mathcal A_{\mathbb R},
 \qquad
 \rho_{\rm exc}:\mathfrak H_{\rm cyl}\to\mathcal A_{\rm exc}.
 \tag{114.5}
\]

La première contracte des cylindres compatibles ; la seconde conserve l'incidence de frontière. Les cinq aspects spectral, solénoïdal, de jauge, cohérent et déterminantal appartiennent au même domaine de (114.5). \(\mathbb R\) est la sortie d'une représentation commutative ; l'arithmétique holonomique conserve en outre les commutateurs et le registre de mémoire que cette sortie ne voit pas.

La lecture de séquences constitue une autre représentation. Avec

\[
 A\mapsto(0,0),\quad G\mapsto(0,1),\quad
 T\mapsto(1,0),\quad C\mapsto(1,1),
 \tag{114.6}
\]

et la phase \((n\bmod3)\), la complémentation inverse définit une involution sur la catégorie des mots. L'APP évalue accumulation et concaténation ; TRIT type la courbure ; TPK conserve retours et frontière. La condition sturmienne

\[
 p_s(m)=m+1
 \tag{114.7}
\]

fournit le seuil exact entre périodicité et complexité exponentielle. Ainsi, l'idée de cristal apériodique appliquée aux séquences biologiques devient un lecteur falsifiable de facteurs, d'extensions et de mots de retour.

\section{Théorème fondamental de l'arithmétique holonomique HMT}

Les sections précédentes permettent de condenser la généalogie sans la réduire. Définissons

 \[
 \boxed{
 \mathfrak H_{\rm HMT}^{\rm alg}
 =\left[\varinjlim_n
 \left(
 \mathbb R\langle e_x,J_x,T_u\rangle/
 \mathcal R_{\rm HMT}
 \right)\right]
 \rtimes_{\Gamma_9}\mathbb N.}
 \tag{115.1}
 \]

Dans le secteur réversible, le dernier facteur est \(\mathbb Z\). L'espace des états est \(X_\infty=\varprojlim X_n\). Pour une représentation compatible et bornée \(\rho\), la complétion \(\mathfrak H_{\rm HMT}^{\rho}
=\overline{\rho(\mathfrak H_{\rm HMT}^{\rm alg})}^{\|\cdot\|}\) appartient au lecteur ; le générateur universel ne présuppose pas de norme.

\begin{quote}
\textbf{Théorème fondamental.} La double évaluation APP, l'orientation trigéométrique TRIT et la dynamique TPK engendrent une arithmétique associative non commutative des trajectoires, tordue par le cocycle de retenue et graduée par phase et mémoire. Le raffinement conserve l'unité, l'holonomie ramène la phase tout en avançant la mémoire, et le spectre typé des représentations détermine le système corrélé \((\pi,\varphi,e,\alpha)\), la structure discrète conjointe du continu, l'électron, les échelles d'action et de gravité, les paramètres géométriques, les constantes dérivées et les symétries de frontière.
\end{quote}

La démonstration s'obtient en concaténant cinq faits déjà établis :

\begin{enumerate}
\item les équations APP conservent les deux feuillets, le résidu et le quotient ;
\item \(J_\tau^2=-\tau I\) conserve les trois courbures dans un seul TRIT ;
\item l'identité de cocycle prouve l'associativité du produit des trajectoires ;
\item l'extension non scindée prouve que retour de phase et retour d'état sont distincts ;
\item la propriété universelle prouve que tout lecteur compatible se factorise par (115.1).
\end{enumerate}

La première conséquence est que \(\mathbb R\) ne constitue pas le support fondamental. C'est une représentation terminale d'une structure qui conserve plus d'information que sa coordonnée réelle. La deuxième est qu'une constante n'est pas un nombre enregistré : c'est un invariant d'un lecteur. La troisième est qu'une particule n'est pas un nombre : c'est une section persistante d'un module. La quatrième est qu'une symétrie n'est pas postulée : elle apparaît comme stabilisateur d'une frontière conservée.

La forme la plus brève de l'évolution conservative est

\[
 \boxed{
 \underbrace{\varprojlim X_n}_{\text{histoire qui s’approfondit}}
 \quad\longleftrightarrow\quad
 \underbrace{\varinjlim\mathfrak H_n}_{\text{lecteurs qui s’étendent}}.}
 \tag{115.2}
\]

Le cristal apériodique hélicoïdal coordonne les deux limites. L'identité trigéométrique (111.4) apporte son calcul local ; \(\alpha\) apporte la couture de mémoire ; le bloc \({\rm Cl}_{1,1}\) apporte son coin électronique ; l'invariant \((G\hbar)\) apporte sa conservation dimensionnelle ; Barbero--Immirzi apporte son covecteur de frontière ; et la double projection apporte ses sorties archimédienne et exceptionnelle. Telle est la structure minimale à partir de laquelle peut se déployer l'article autonome sans dépendre, dans son exposé, des (2\,099) pages et sans perdre aucun de ses opérateurs essentiels.

\textbf{Critère de réfutation.} Le théorème est réfuté si l'identité de cocycle échoue, si le raffinement cesse d'être unitaire, si l'holonomie réinitialise la mémoire, si une fibre TRIT disparaît, si une réalisation ne satisfait pas les relations qu'elle déclare représenter ou si une sortie utilise sa propre valeur conventionnelle pour sélectionner la trajectoire qui devait la produire.

\section{Irréductibilité triadique et différence de potentiel dans les lecteurs de séquences}

La composition APP--TRIT--TPK peut être démontrée irréductible au moyen de trois quotients d'oubli. Soit \(q_{\rm APP}\) le quotient qui identifie les feuillets additif et multiplicatif. Alors

\[
q_{\rm APP}(\Delta_{\Sigma\Pi})=0,
\qquad
q_{\rm APP}([P_\Sigma,P_\Pi])=0.
\tag{116.1}
\]

Soit \(q_{\rm TRIT}\) le quotient qui identifie les trois idempotents de régime. Le triplet de carrés est alors perdu

\[
(J_+^2,J_0^2,J_-^2)=(-I,0,I).
\tag{116.2}
\]

Soit \(q_{\rm TPK}\) le quotient qui trivialise le cocycle de retenue et la graduation de mémoire. Dans ce quotient,

\[
q_{\rm TPK}\!\left(d_M(\Gamma_9x)-d_M(x)\right)=0,
\tag{116.3}
\]

bien que, dans l'état holonomique complet, la différence vaille \(1\). Les trois invariants (116.1)--(116.3) sont de types distincts et aucun ne peut être reconstruit à partir des deux autres après application de son quotient. L'APP apporte l'écart et l'ordre des feuillets ; TRIT, la signature de courbure ; TPK, la composition, l'holonomie et la mémoire. Telle est la forme algébrique exacte de l'irréductibilité ternaire du noyau.

La même séparation clarifie l'application aux séquences. Dans l'alphabet \(\{A,G,T,C\}\), définissons

\[
\chi_p(A,G,T,C)=(0,0,1,1),
\qquad
\chi_h(A,G,T,C)=(0,1,0,1),
\tag{116.4}
\]

et l'involution de complément \(\mathcal C_{\rm DNA}(A,G,T,C)=(T,C,A,G)\). On vérifie

\[
\chi_p(\mathcal C_{\rm DNA}b)=1-\chi_p(b),
\qquad
\chi_h(\mathcal C_{\rm DNA}b)=\chi_h(b).
\tag{116.5}
\]

Par conséquent,

\[
\mathcal V_n=\chi_p(b_{n+1})-\chi_p(b_n)
\tag{116.6}
\]

est une différence de potentiel symbolique orientée, tandis que \(\chi_h\) conserve la classe transversale de liaison. Le relèvement

\[
b_n\longmapsto
\bigl(\chi_p(b_n),\chi_h(b_n),n\bmod3,\lambda_n,c_n,m_n\bigr)
\tag{116.7}
\]

ajoute régime TRIT, frontière, retenue et mémoire TPK. Ainsi se formule la raison structurelle pour laquelle polarité, potentiel et mot apériodique peuvent être lus dans la même arithmétique sans faire de la biologie le générateur du formalisme.

\textbf{Critère de réfutation.} L'irréductibilité échoue si l'un des trois invariants survit au quotient qui doit l'effacer ou peut être reconstruit à partir des deux autres sans réintroduire la donnée éliminée. Le lecteur de séquences échoue si le complément n'inverse pas \(\chi_p\), ne conserve pas \(\chi_h\), ou si le relèvement n'entrelace pas la phase ternaire avec la dynamique TPK.

\section{La lacune électronique comme paire nulle de Witt}

L'arithmétique du coin électronique permet de compléter localement les trois régimes sans introduire une autre algèbre. Avec

\[
S=2P-I,\qquad
J=\frac{[P,Q]}{\sqrt3/4},
\tag{117.1}
\]

on a

\[
S^2=I,\qquad J^2=-I,\qquad SJ=-JS.
\tag{117.2}
\]

Les combinaisons

\[
n_+=\frac{S+J}{2},\qquad
n_-=\frac{S-J}{2}
\tag{117.3}
\]

satisfont

\[
n_+^2=n_-^2=0,\qquad
n_+n_-+n_-n_+=I.
\tag{117.4}
\]

Par conséquent,

\[
\{J,n_\pm,S\}
\quad\longmapsto\quad
\{\text{elliptique},\text{parabolique},\text{hyperbolique}\}
\tag{117.5}
\]

est une réalisation locale complète du typage TRIT dans \({\rm Cl}_{1,1}\cong M_2(\mathbb R)\). La lacune se caractérise par deux descriptions compatibles : dans l'APP, c'est la perte unitaire du quotient-produit qui conserve le feuillet additif, le résidu produit et le régime ; dans le coin de Clifford, c'est une direction nulle parabolique. Le couple \((n_+,n_-)\) est la paire de Witt locale des deux orientations. Cette décomposition offre le domaine algébrique précis à partir duquel l'incidence Paley--Witt ultérieure peut réaliser la même généalogie dans son propre module combinatoire.

L'affirmation est plus forte que l'analogie avec les quaternions. Les quaternions de Hamilton possèdent trois générateurs de carré \(-I\) et forment une algèbre à division. Le bloc (117.2) est déployé : il contient des diviseurs de zéro, des idempotents et des directions nulles, précisément les objets requis par feuillet, lacune et seuil. Les octonions augmentent la dimension en sacrifiant l'associativité ; l'arithmétique holonomique maintient l'associativité au moyen du cocycle et élargit la structure par les parcours, la mémoire et le raffinement.

\textbf{Critère de réfutation.} Il suffit qu'une des quatre identités (117.2)--(117.4) échoue, que la lacune ne conserve pas ses données APP déclarées, ou que le transport TPK ne conserve pas l'idempotent de régime.

\section{Relations de Weyl de l'horloge finie et de l'horloge irrationnelle}

Soit \(T\) le pas de l'orbite, soit \(\delta=\log_{10}(10/9)\), soient \(\zeta_9\) et \(\omega_\delta\) les caractères internes de phase nonadique et de capacité, et définissons sur \(\lvert t,m\rangle\)

\[
V_9\lvert t,m\rangle=\zeta_9^t\lvert t,m\rangle,
\qquad
W_\delta\lvert t,m\rangle=\omega_\delta^t\lvert t,m\rangle.
\tag{118.1}
\]

L'avancement temporel satisfait

\[
V_9T=\zeta_9TV_9,\qquad
W_\delta T=\omega_\delta TW_\delta.
\tag{118.2}
\]

Le tour \(\Gamma_9=T^9\) produit

\[
V_9\Gamma_9=\Gamma_9V_9,\qquad
W_\delta\Gamma_9=\omega_\delta^9\Gamma_9W_\delta.
\tag{118.3}
\]

Les deux horloges se distinguent ainsi algébriquement. La première possède une période de neuf ; la seconde conserve une phase irrationnelle. Le mot mécanique des lacunes est le codage par seuil de l'orbite de la seconde horloge. Si \(D_M\) est l'opérateur de degré de mémoire, le relèvement TPK vérifie

\[
[D_M,\Gamma_9]=\Gamma_9.
\tag{118.4}
\]

La représentation elliptique ultérieure reconnaît ces caractères au moyen de

\[
\rho_+(\zeta_9)
=\operatorname{Exp}_+\!\left(\frac{2\pi}9J_+\right),
\qquad
\rho_+(\omega_\delta)
=\operatorname{Exp}_+\!\left(2\pi\delta J_+\right).
\tag{118.5}
\]

Les caractères, et non leurs expressions complexes reconnues, constituent les données opératoires de l'horloge.

L'équation (118.4) affirme exactement que \(\Gamma_9\) est un opérateur d'élévation : la phase revient et l'état monte d'une unité dans l'hélice. L'action de translation est abélienne ; le produit croisé de cette action avec les observables \(V_9,W_\delta,D_M\) est non commutatif par (118.2)--(118.4). La structure temporelle n'exige donc pas de choisir entre « horloge finie » et « cristal apériodique » : ce sont deux caractères conjugués d'un unique système gauchi.

L'entropie topologique nulle procède de la rotation irrationnelle et de la complexité \(p(m)=m+1\) ; la croissance de mémoire procède de (118.4). Une entropie nulle et une trajectoire infinie coexistent parce que la seconde mesure la profondeur reconstructible, tandis que la première mesure la croissance exponentielle des facteurs.

\textbf{Critère de réfutation.} Le modèle échoue si \(V_9\) ne commute pas avec \(\Gamma_9\), si \(W_\delta\) acquiert une période finie, si le commutateur de mémoire diffère de \(\Gamma_9\), ou si le mot de seuil cesse d'avoir une complexité sturmienne.

\section{Équation minimale du générateur trigéométrique}

Soit

\[
\mathbb J=p_+J_++p_0J_0+p_-J_-,
\qquad p_++p_0+p_-=I.
\tag{119.1}
\]

Les trois relations de courbure équivalent à

\[
\boxed{\mathbb J^6=\mathbb J^2.}
\tag{119.2}
\]

En effet, \(K=\mathbb J^2\) satisfait \(K^3=K\), et les trois idempotents sont retrouvés sans données extérieures :

\[
p_+=\frac{K^2-K}{2},\qquad
p_0=I-K^2,\qquad
p_-=\frac{K^2+K}{2}.
\tag{119.3}
\]

L'identité exponentielle intrinsèque est

\[
\operatorname{Exp}_\star(t\mathbb J)=
p_+(\cos t+\mathbb J\sin t)+
p_0(I+t\mathbb J)+
p_-(\cosh t+\mathbb J\sinh t).
\tag{119.4}
\]

Les équations (118.4) et (119.4) se réunissent dans le propagateur holonomique

\[
\boxed{
\mathcal E_{\rm HMT}(t)
=\Gamma_9\operatorname{Exp}_\star(t\mathbb J),
\qquad
[D_M,\Gamma_9]=\Gamma_9.}
\tag{119.5}
\]

Telle est la compression opératoire recherchée. La première coordonnée de (119.5) gouverne la réponse locale dans les trois courbures ; la seconde gouverne l'élévation de mémoire. L'APP intervient par les idempotents de feuillet, leurs projecteurs et le cocycle qui définit \(\Gamma_9\) ; TRIT intervient par \(\mathbb J\) ; TPK intervient par la composition, le relèvement et \(D_M\). Le système corrélé, l'électron, l'action, la gravité, Barbero--Immirzi, les constantes d'échelle, les symétries et les lecteurs de séquences s'obtiennent comme représentations typées de l'algèbre universelle contenant (119.5).

La fonction de \(\alpha\) devient transparente dans cette présentation. \(\operatorname{Exp}_\star\) contient la dynamique locale et \(\Gamma_9\) le tour avec mémoire ; la fermeture dodécaphasique évalue leur compatibilité et détermine \(\alpha\) comme couture. Ainsi, \(\alpha\) relie la synchronisation de \((\pi,e,\varphi)\) aux lecteurs dimensionnels parce qu'il mesure une classe de fermeture du propagateur holonomique, non parce qu'il s'ajoute à trois constantes préexistantes.

\textbf{Critère de réfutation.} L'équation minimale échoue si \(\mathbb J^6\ne\mathbb J^2\), si les polynômes (119.3) cessent d'être des idempotents orthogonaux, si (119.4) ne se restreint pas aux trois exponentielles certifiées, ou si \(\alpha\) ne coïncide pas dans ses deux voies typées à l'intérieur du carré de compatibilité existant.
 
% Realizaciones dimensionales completas tomadas de la autoridad material viva.
% El orden respeta la precedencia electrón estructural -> acción -> constitución
% del vacío -> gravitación -> área e información.
% \newcommand{\HMTMonographChapterInput}[2]{%
%   \chapter{#1}%
%   \begingroup
%   \renewcommand{\chapter}[2][]{}%
%   \input{#2}%
%   \endgroup}

\chapter{L'électron comme section persistante de la lacune orientée}
\begingroup
\renewcommand{\chapter}[2][]{}%
\chapter{L’électron comme 1re réalisation matérielle : étape basale, fermeture bêta et structure atomistique de l’électricité}
\label{chap:parte-iv-48}

L’état généalogique acquiert ici sa première réalisation matérielle complète. La localisation centrale, le ruban électron--positron et la lecture bidirectionnelle sont construits avant toute comparaison métrologique ; l’étape basale est distinguée de la fermeture bêta directrice.

\section{Clôture \(\beta\) de la section électronique et expression dimensionnelle de son énergie au repos}
\label{chap:biografia-electron-rev6}
\index{électron!construction holographique}
\index{positron!bande conjuguée}
\index{masse!loi électronique}

L'électron apparaît dans HMT comme le point où s'assemblent six
constructions antérieures et mutuellement typées : l'incompatibilité entre
les deux feuillets arithmétiques d'APP, l'orientation induite par l'involution
de la charnière surfacique--volumétrique, deux quotients entiers du recensement APP--TPK,
la clôture
dodécaphasique de \(\alpha\), le défaut global
\(\Delta_4\) et le centre unique de l'atlas nonadique. La bande de route ajoute
ensuite la conjugaison électron--positron. La carte en MeV et la comparaison
avec CODATA sont postérieures à toutes ces constructions.

Cette antériorité fait partie du contenu mathématique : les facteurs, leurs
types et leur ordre causal sont déterminés avant l'ouverture de la carte tabulée, qui
n'intervient que dans la comparaison métrologique ultérieure.

\subsection{Scalaire électronique \(\mathfrak e\) et graphe causal de ses facteurs APP–TPK, \(\alpha\), \(\Delta_4\) et atlas nonadique}

La sortie interne est le scalaire adimensionnel
\begin{equation}
 \label{eq:electron-holografico-rev6}
 \boxed{
 \mathfrak e
 =
 \frac{\sqrt3}{4}
 \left[
 \exp\!\left(\frac{\varphi}{\pi^2}\right)
 -\frac{396}{18}\,\alpha^3
 \right]
 \left[
 1+\frac{270}{18}\,\Delta_4
 \right].}
\end{equation}
Les identités
\[
 \frac{396}{18}=22=27-5,
 \qquad
 \frac{270}{18}=15=\binom62=3\cdot5
 \]
permettent de retrouver l'écriture abrégée
\[
 \mathfrak e
 =
 \frac{\sqrt3}{4}
 \left(e^{\varphi/\pi^2}-22\alpha^3\right)
 (1+15\Delta_4).
\]
Ici, la lettre \(e\) à l'intérieur de l'exponentielle est le nombre d'Euler ;
l'indice de \(\mathfrak e\) désigne la section électronique.

Le graphe causal de l'équation est
\[
\begin{array}{ccl}
\text{projecteurs APP en dimension trois}
&\longrightarrow& \sqrt3/4,\\[2mm]
\text{calendrier TPK et feuillets aréal--volumétrique}
&\longrightarrow&
F_{\mathrm{av}},\ \pi/\varphi^2,\
\varphi/\pi^2,\\[2mm]
\text{clôture dodécaphasique}
&\longrightarrow&\alpha,\\[2mm]
\text{recensement APP--TPK et famille électronique}
&\longrightarrow&396,\ 18,\ 270,\ 22,\ 15,\ 5,\\[2mm]
\text{lectures globales de }\pi,e\text{ et registre}
&\longrightarrow&\Delta_4,\\[2mm]
\text{atlas nonadique et anti-section}
&\longrightarrow&(5,5),\
\text{bande électron--positron}.
\end{array}
\]
Aucune flèche de ce diagramme ne reçoit comme argument la masse électronique
de référence.

\subsection{1ᵉʳ facteur : incompatibilité exacte des feuillets APP}

Pour \(d\geq2\), soit
\[
 \mathscr A_d=\{1,\ldots,9\}^d
\]
muni de la mesure uniforme, et soient
\[
 \Sigma_d(x)=\operatorname{dr}(x_1+\cdots+x_d),
 \qquad
 \Pi_d(x)=\operatorname{dr}(x_1\cdots x_d)
\]
les observables additive et multiplicative. Leurs espérances conditionnelles
orthogonales sur \(\ell^2(\mathscr A_d)\) sont
\(P_{\Sigma,d}\) et \(P_{\Pi,d}\). La
\cref{prop:app-no-conmutatividad} démontre par énumération exacte que
\[
 \left\|[P_{\Sigma,3},P_{\Pi,3}]\right\|
 =\frac{\sqrt3}{4}.
\]

\begin{proposition}[Contenu APP du préfacteur électronique]
\label{prop:prefactor-app-electron-rev6}
Le préfacteur \(\sqrt3/4\) de
\cref{eq:electron-holografico-rev6} est la norme exacte de la
non-commutativité entre les feuillets somme et produit en dimension trois.
Ce n'est pas un facteur géométrique introduit à partir de la formule électronique.
\end{proposition}

\begin{proof}
Le tableau conjoint de \(\Sigma_3\) et \(\Pi_3\) détermine les angles
principaux entre
\(\operatorname{Ran}P_{\Sigma,3}\) et
\(\operatorname{Ran}P_{\Pi,3}\). Si \(s\) est l'une de ses valeurs
singulières, la contribution correspondante à la norme du commutateur est
\(s\sqrt{1-s^2}\). L'énumération des \(9^3\) mots donne pour maximum
\[
 s^2(1-s^2)=\frac3{16},
 \qquad s^2=\frac14.
\]
La norme maximale est donc \(\sqrt{3/16}=\sqrt3/4\), avant
d'introduire \(\alpha\), \(\Delta_4\), une unité d'énergie
ou une masse de référence.
\end{proof}

La fonction physique du facteur est ainsi fixée : il mesure l'incompatibilité
maximale des deux lectures arithmétiques sur le même support ternaire.
L'électron n'hérite pas de deux nombres indépendants, l'un additif et l'autre
multiplicatif ; il hérite de l'obstruction exacte à rendre simultanément
compatibles les deux feuillets. Le préfacteur conserve dans la loi de masse la
mémoire locale du désaccord APP dont naît la route.

\subsection{2ᵉ facteur : l'involution d'échange entre \(\pi\) et \(\varphi\)}

Le calendrier primitif des modes
\((\Sigma,\Pi,\Pi)\) induit, dans les feuillets surfacique et volumétrique, la matrice
\[
 F_{\mathrm{av}}
 =
 \begin{pmatrix}
 0&1\\
 1&1
 \end{pmatrix}.
\]
Le \cref{thm:descenso-necesario-fav} l'obtient en comptant les transitions
\(\Sigma\to\Pi\), \(\Pi\to\Pi\) et \(\Pi\to\Sigma\), chacune une fois.
La matrice n'est pas postulée à partir du nombre d'or : ses valeurs propres
\(\varphi\) et \(-\varphi^{-1}\) apparaissent après la construction de
l'incidence.

Dans la mémoire de deuxième ordre,
\(\operatorname{Sym}^2(F_{\mathrm{av}})\) a pour spectre
\[
 \varphi^2,\qquad -1,\qquad\varphi^{-2}.
\]
Ainsi, \(\varphi^{-2}\) est le mode stable positif de la charnière. La marque
de clôture \(\pi\), appliquée une seule fois au retour surfacique, produit
\[
 \frac{\pi}{\varphi^2}.
\]
Ce rapport conserve l'ordre clôture--auto-échelle. La loi électronique utilise son
orientation opposée.

\begin{theorem}[Involution exacte de la charnière électronique]
\label{thm:espejo-electron-rev6}
Soient
\[
 \mathscr R(x,y)=\frac{x}{y^2},
 \qquad
 \mathsf J(x,y)=(y,x).
\]
Alors \(\mathsf J^2=I\) et
\[
 \mathscr R(\pi,\varphi)=\frac{\pi}{\varphi^2},
 \qquad
 (\mathscr R\circ\mathsf J)(\pi,\varphi)
 =\frac{\varphi}{\pi^2}.
\]
Par conséquent,
\[
 \frac{\varphi}{\pi^2}
 =
 \frac{1}{
 \varphi^3\bigl(\pi/\varphi^2\bigr)^2},
 \qquad
 e^{\varphi/\pi^2}
 =
 \exp\!\left[
 \frac{1}{
 \varphi^3\bigl(\pi/\varphi^2\bigr)^2}
 \right].
\]
\end{theorem}

\begin{proof}
Le premier couple d'identités s'obtient en appliquant
\(\mathscr R\) avant et après l'involution \(\mathsf J\). Pour le
reparamétrage,
\[
 \varphi^3
 \left(\frac{\pi}{\varphi^2}\right)^2
 =\frac{\pi^2}{\varphi},
\]
et son inverse est \(\varphi/\pi^2\). La dernière égalité découle de l'application de
l'exponentielle.
\end{proof}

La présence simultanée de \(\pi\) et \(\varphi\) possède donc une
signification structurelle précise. \(\pi/\varphi^2\) lit la clôture surfacique
sur la mémoire volumétrique quadratique ; \(\varphi/\pi^2\) lit
l'auto-croissance intérieure sur la clôture quadratique. Ce ne sont pas deux quotients
semblables choisis par commodité. Ce sont les deux orientations d'un même
lecteur, échangées par une involution exacte. Le terme
\(e^{\varphi/\pi^2}\) publie dans la section électronique l'orientation
intérieure de cette charnière.

\subsection{3ᵉ et 4ᵉ facteurs : les entiers 22, 15 et le centre 5}

Le recensement de la carte APP employée par la famille électronique fournit
trois sommes entières :
\[
 \Sigma(\mathrm{APP})=396,
 \qquad
 \Sigma(\mathrm{centro}\ 2\times2)=18,
 \qquad
 \Sigma(\mathrm{canal})=270.
\]
La première admet la décomposition en anneaux
\[
 396=18+72+108+198.
\]
Le noyau \(2\times2\) est donc une sous-structure explicite de la
carte complète, et l'entier 270 appartient au canal global du même recensement
APP--TPK.
La somme des anneaux pairs est
\[
 18+108=126,
 \qquad
 \frac{396}{126}=\frac{22}{7}.
\]
Ce dernier quotient est le sceau rationnel de clôture de la carte, non une
identification de \(22/7\) au nombre \(\pi\) engendré par le lecteur
coinductif. Il montre bien que le numérateur \(22\) remplit déjà une fonction
interne avant d'apparaître dans la loi électronique.

\begin{theorem}[Provenance interne des coefficients électroniques]
\label{thm:coeficientes-electron-rev6}
Dans la famille électronique fixée par la carte APP et son canal global,
\[
 \boxed{
 22=\frac{\Sigma(\mathrm{APP})}
 {\Sigma(\mathrm{centro})}
 =\frac{396}{18}=27-5,}
\]
\[
 \boxed{
 15=\frac{\Sigma(\mathrm{canal})}
 {\Sigma(\mathrm{centro})}
 =\frac{270}{18}=3\cdot5=\binom62.}
\]
En particulier, les coefficients sont déterminés sans consulter
\(Z_0\), la masse de l'électron ni une carte d'unités.
\end{theorem}

\begin{proof}
Les deux premières égalités sont des divisions exactes du recensement fini :
\[
 396=22\cdot18,
 \qquad
 270=15\cdot18.
\]
En définissant
\[
 p:=27-22,
\]
on obtient \(p=5\), et alors
\[
 15=3p=3\cdot5=\binom62.
\]
Toutes les quantités appartiennent au domaine entier de la carte et de la
famille avant la transduction métrologique.
\end{proof}

\subsubsection{Sélecteur de persistance dans le retour central}

Le même couple possède une construction dynamique qui fixe en outre son signe sans
consulter une masse. Les quatre orientations du germe central
\((5,5)\) reviennent pour la première fois à cette cellule après vingt-sept pas.
L'espace typé du retour est
\[
 \mathcal K_e=\mathbb Z/9\mathbb Z\times\mathbb F_3,
 \qquad |\mathcal K_e|=27.
\]
Par ailleurs, la microfibre de \(\pi\) contient les cinq régions
généalogiquement distinctes
\[
 \mathscr R_\pi=3_{++}+1_{--}+1_{\perp}.
\]
La jauge régionale \(\operatorname{diag}_c\) et, au sein d'une phase répétée,
l'ordre interne de \(U_6\) définissent l'injection
\[
 \iota:\mathscr R_\pi\hookrightarrow\mathcal K_e
\]
au moyen des cinq places distinctes
\[
 (4,0),\quad(5,0),\quad(6,0),\quad(6,1),\quad(6,2).
\]
Les deux premières correspondent respectivement à la région transversale et
au retour muté ; les trois dernières séparent les régions positives par leur
feuillet tritique. Le mot visible commun n'intervient pas comme inverse de l'application.

\begin{theorem}[Sélecteur interne des coefficients électroniques]
\label{thm:selector-persistencia-electron-rev7}
Une fois fixés le retour central et les cinq régions typées, les coefficients de
l'équation électronique sont
\[
 \boxed{
 n_e=-\bigl|\mathcal K_e\setminus\iota(\mathscr R_\pi)\bigr|=-22,
 \qquad
 m_e^{\rm tor}=\bigl|\mathscr R_\pi\times\mathbb F_3\bigr|=+15.}
\]
Le premier signe enregistre le déficit de frontière ; le second, la mémoire
torsionnelle retenue.
\end{theorem}

\begin{proof}
Les cinq places précédentes sont distinctes, de sorte que \(\iota\) est
injective et
\[
 |\mathcal K_e\setminus\iota(\mathscr R_\pi)|=27-5=22.
\]
Chacune des cinq régions conserve trois feuillets tritiques, donc
\[
 |\mathscr R_\pi\times\mathbb F_3|=5\cdot3=15.
\]
La règle d'orientation attribue un signe négatif aux états retirés de la
coquille et un signe positif à la mémoire retenue. La construction n'ouvre que le
centre, le retour de vingt-sept états et les identités régionales ; elle
n'ouvre ni une masse observée, ni une impédance, ni une signature étalonnée.
\end{proof}

Le nombre \(22\) mesure combien de noyaux centraux de poids \(18\) recompose la
somme APP \(396\). Le nombre \(15\) mesure combien de ces mêmes noyaux
recompose le canal \(270\), et compte simultanément les paires d'une
hexade. Le nombre \(5\) n'est pas un paramètre supplémentaire :
\[
 5=27-22,\qquad 15=3\cdot5.
\]
Les trois écritures montrent la même clôture depuis la profondeur 27,
depuis la normalisation centrale et depuis la combinatoire hexadique.

\subsubsection{Conservation des coefficients 22 et 15 dans la réalisation matérielle de l'électron}

L'exploration historique a d'abord localisé \(p=5\) au moyen d'une comparaison
finie avec \(Z_0\), et a écrit à partir de là
\((27-p,3p)=(22,15)\). Cette procédure conserve sa valeur comme histoire de la
découverte et comme contrôle externe. Ce n'est cependant pas la démonstration
qui gouverne la loi : le dénombrement ultérieur
\[
 \frac{396}{18}=22,
 \qquad
 \frac{270}{18}=15
\]
construit les deux entiers au sein d'APP--TPK et retrouve \(p=5\) comme
conséquence. L'histoire explique comment le couple a été trouvé ; le recensement
explique pourquoi il appartient à l'équation.

\subsection{5ᵉ facteur : le cube de la clôture \(\alpha\)}

La clôture de \(\alpha\) précède la mécanique dimensionnelle.
Comme il est établi dans le \cref{chap:alpha}, les douze blocs en base mille
de \(\pi\), \(e\), \(\varphi\) et du sceau \(K\) se combinent au moyen de la
retenue réversible
\[
 P_m+E_m-\Phi_m-K_m+c_{m+1}
 =A_m+1000c_m,
 \qquad c_{13}=c_1=0.
\]
La sortie est
\[
 A_\alpha=
 (007,297,352,569,283,800,997,285,105,472,380,663)
\]
et, par concaténation, la fenêtre dodécaphasique exacte
\[
 \alpha_{12}=\pi_{12}(\alpha)
 =
 \frac{
 7297352569283800997285105472380663
 }{10^{36}}.
\]
Le télescopage des retenues démontre l'identité entière complète ; la
prolongation compatible construit
\(\alpha=\varprojlim_N A^{(N)}\), et \(\alpha_{12}\) est sa
projection rationnelle à trente-six chiffres. La
comparaison avec la constante de structure fine s'effectue ensuite et
n'intervient pas dans sa construction.

La loi électronique prend la troisième puissance
\[
 \alpha^3.
\]
Mathématiquement, il s'agit de l'évaluation multiplicative de la troisième
puissance symétrique du même scalaire déjà clos ; on n'introduit ni trois
copies métrologiques ni trois paramètres différents. Sa contribution est
\[
 22\alpha^3
 =
 0.00000854906599200551727014979807598228\ldots.
\]

\begin{proposition}[Statut du canal cubique électronique]
\label{prop:alpha-cubica-electron-rev6}
Le terme \(22\alpha^3\) de la loi électronique dépend
uniquement de la clôture préalable de \(\alpha\), de la
multiplication interne et de l'entier \(22\) du
\cref{thm:coeficientes-electron-rev6}. Il ne dépend pas d'une lecture de la
masse électronique.
\end{proposition}

\begin{proof}
\(\alpha_{12}\) est le rationnel fixé par l'identité de concaténation et de
retenue, et la famille compatible de fenêtres détermine la section
\(\alpha\). La projection \(\alpha_{12}^3\) certifie les
chiffres imprimés de \(\alpha^3\).
Le coefficient \(22\) est déterminé par \(396/18\). La composition des
deux opérations ne contient que des données construites antérieurement à la
carte d'énergie.
\end{proof}

Cette proposition fixe exactement ce qui est utilisé dans l'électron. La puissance
trois est le degré cubique de ce canal dans la loi publiée. Il n'est pas nécessaire
de postuler ici une règle universelle selon laquelle toute occurrence
\(\alpha^d\) représente automatiquement une dimension \(d\) ; une
affirmation universelle de ce type exigerait son propre théorème de
multiplicativité entre degrés. La généalogie électronique est complète
sans transformer cette généralisation en prémisse.

\subsection{6ᵉ facteur : torsion minimale et défaut global \(\Delta_4\)}
\index{torsion minimale!\(\Delta_4\)}

Les dénominations \emph{torsion minimale} et \emph{défaut global}
désignent ici un unique invariant, non deux corrections distinctes. L'objet
employé par la loi des masses est
\begin{equation}
 \label{eq:delta4-electron-rev6}
 \boxed{
 \Delta_4
 =
 \frac{\pi-e\log\pi}{270}
 \left(1+\frac{\pi}{729}\right).}
\end{equation}
Ses entrées sont les lectures HMT déjà fixées de \(\pi\) et \(e\), le canal
global \(270\) et la cellule intérieure \(729\). Par conséquent,
\[
 \Delta_4
 =
 0.00011119642269214005405113478612067909\ldots.
\]
Le premier rapport mesure le défaut signé
\(\pi-e\log\pi\) en unités du canal 270 ; le second facteur couple ce
rapport à la cellule intérieure au moyen de \(\pi/729\). Cette lecture n'ajoute pas de
paramètres à la formule : elle sépare ses deux opérations constitutives.
Le \cref{p3-20260826:chap:medida-accion-positiva} montre qu'il s'agit d'une entrée mathématique
prédimensionnelle et la distingue de l'opérateur qui la réalise éventuellement
comme changement de feuillet.

\begin{proposition}[Globalité du défaut électronique]
\label{prop:delta4-global-electron-rev6}
\(\Delta_4\) est un invariant global antérieur à l'espèce. Dans la loi
électronique, il intervient au moyen de
\[
 1+15\Delta_4
 =
 1.00166794634038210081076702179181019\ldots.
\]
Ce n'est pas un correcteur choisi après avoir pris connaissance du résidu de la masse de
l'électron.
\end{proposition}

\begin{proof}
La \cref{eq:delta4-electron-rev6} ne contient ni étiquette d'espèce, ni
masse, ni unité. L'entier \(15\) est déjà construit par \(270/18\).
La composition \(1+15\Delta_4\) est donc déterminée par les données
globales et le recensement interne.
\end{proof}

Il faut distinguer \(\Delta_4\) des autres blocs du caractère de masse et des
autres notions de torsion. En particulier, ce n'est ni \(s_{270}\), ni
l'invariant quadratique \(135\Delta_4^2\), ni le résidu angulaire
\(\varepsilon_{\rm res}\), ni la torsion géométrique de Cartan--Holst.
Substituer l'un à l'autre répéterait ou changerait le niveau de la correction. Dans
l'électron, \(\Delta_4\) conserve la mémoire globale du raffinement ; le facteur
\(15\) déclare combien de fois la famille électronique la lit. Le qualificatif
\emph{minimale} conserve le nom historique de cet invariant et n'est pas utilisé
comme substitut à un théorème variationnel de minimalité.

\subsection{Réalisation matérielle du point fixe (5,5) dans la clôture électronique}

Soit
\[
 \mathcal C_{81}=\{1,\ldots,9\}^2
\]
l'atlas nonadique et
\[
 \rho(r,c)=(10-r,10-c)
\]
son inversion centrale. Le \cref{thm:censo-atlas-81} prouve que l'atlas
contient une unique cellule de génération zéro et que celle-ci porte l'étiquette
positionnelle \(e_{\mathrm{core}}\).

\begin{proposition}[Unicité du centre électronique]
\label{prop:centro-electron-rev6}
L'involution \(\rho\) possède un unique point fixe :
\[
 \boxed{(r,c)=(5,5).}
\]
Son unicité est combinatoire et ne dépend pas de la valeur de
\(\mathfrak e\).
\end{proposition}

\begin{proof}
L'équation \((r,c)=\rho(r,c)\) équivaut à
\[
 2r=10,\qquad2c=10.
\]
Dans \(\{1,\ldots,9\}^2\), elle a pour unique solution \(r=c=5\). L'évaluation
de la masse n'apparaît pas dans l'argument.
\end{proof}

\begin{corollary}[Universalité de la section centrale]
\label{cor:universalidad-electron-central-rev6}
Tout automorphisme admissible de l'atlas qui entrelace l'inversion centrale
préserve la cellule \((5,5)\).
\end{corollary}

\begin{proof}
Si \(f\rho=\rho f\) et \(\rho(5,5)=(5,5)\), alors
\(\rho f(5,5)=f\rho(5,5)=f(5,5)\). Par conséquent, \(f(5,5)\) est un autre point fixe
de \(\rho\) ; l'unicité impose \(f(5,5)=(5,5)\).
\end{proof}

C'est la raison structurelle pour laquelle la classe électronique est la même dans
toutes les réalisations qui préservent l'atlas et son inversion : un nombre décimal local ne se répète pas
par hasard ; toute carte admissible doit au contraire
transporter le même point fixe et la même chaîne d'invariants. La lecture
physique de cette unicité identifie chaque électron à une réalisation de cette
classe centrale ; la proposition démontre l'universalité interne de l'atlas.

Dans la famille électronique fixée par la construction, ce point est l'ancrage
géométrique de la section. Il convient de ne pas confondre quatre objets :
\[
 (5,5),\qquad
 e_{\mathrm{core}},\qquad
 v_e=(0,0,0,0,0),\qquad
 \text{la signature brute de la cellule centrale}.
\]
Le premier est une position ; le deuxième, son étiquette interne ; le troisième,
l'origine choisie dans le réseau fin d'action ; le quatrième conserve les
données brutes de l'atlas et n'est pas nécessairement le vecteur nul. Leur
coordination explique pourquoi l'électron occupe le centre sans effacer la
généalogie des représentations qui le décrivent.

\paragraph{Routes électroniques orientées \(e_{++},e_{--}\), enveloppe \(e_{\mathrm{env}}\), origine d'action et signatures de la cellule centrale}
La cellule centrale porte deux histoires orientées de douze événements. En écrivant
\(v^2\) pour la concaténation de deux copies d'un sextuplet, la route active
de l'orientation \({++}\), son histoire opposée et l'enveloppe historique
sont respectivement
\[
\begin{aligned}
 e_{\rm act}=e_{++}
   &=(2,2,5,-4,-3,-2)^2,\\
 e_{--}
   &=(4,3,2,-2,-2,-5)^2,\\
 e_{\rm env}
   &=(4,3,5,-4,-3,-5)^2.
\end{aligned}
\tag{E-centro-1}\label{eq:electron-rutas-finas-rev8}
\]
Les deux histoires publient le même mot
\[
 \tau_{12}=\texttt{+++---+++---},
\]
mais ne sont pas la même route. L'enveloppe prend le maximum absolu
coordonnée par coordonnée en conservant le signe ; elle enregistre les amplitudes, mais
aplatit l'ordre généalogique et ne se substitue pas à \(e_{\rm act}\).

Les autres lecteurs du centre vivent dans des codomaines distincts :
\[
\begin{array}{rcll}
 v_e&=&(0,0,0,0,0),
   &\text{origine affine du réseau d’action},\\
 R(5,5)&=&(24,0,12,0,-12),
   &\text{signature APP brute de la cellule},\\
 \operatorname{Sig}^{\Gamma_{108}}_{\rm masa}
   &=&(-12,-4,12,-6,12),
   &\text{projection paire de masse},\\
 \operatorname{Sig}^{\Gamma_{108}}_{\rm carga}
   &=&(0,0,0,0),
   &\text{projection impaire de charge}.
\end{array}
\tag{E-centro-2}\label{eq:electron-firmas-tipadas-rev8}
\]
Le fait que le mot de charge ait une signature impaire nulle ne transforme pas la route en
zéro et n'identifie pas la signature brute à l'origine affine. Position, route
active, route opposée, enveloppe, origine et les deux signatures sont des objets
coordonnés, non six notations pour un même vecteur.

\subsection{Conjugaison particule–antiparticule et persistance de la bande électronique}

L'identité électronique ne s'arrête ni au point central ni à la valeur de la
masse. La conjugaison matérielle se réalise sur des mots dodécaphasiques. Soit
\[
 C_M(\mathbf t)=-\operatorname{rev}(\mathbf t)
\]
l'anti-section orientée et soit
\(\chi\in\{-1,+1\}\) l'orientation de charge. La bande est la classe
\[
 [(\mathbf t,\chi)]_M
 =
 \{(\mathbf t,\chi),(C_M\mathbf t,-\chi)\}
\]
de la \cref{def:banda-dodecafasica-ruta}. La bande conserve ses composantes
paires et présente deux sections de composantes impaires opposées.

\begin{theorem}[Électron et positron comme sections d'une bande]
\label{thm:electron-positron-biografia-rev6}
Pour le caractère de bande du
\cref{thm:principio-conjugacion-masa-banda},
\[
 m(C_M\mathbf t,-\chi)=m(\mathbf t,\chi).
\]
Dans le canal électronique,
\[
 (n,\chi)=(-22,+1)
 \quad\longleftrightarrow\quad
 (+22,-1)
\]
et les deux couples produisent
\[
 \chi n=-22.
\]
L'électron et le positron ont donc la même masse et des orientations de
charge opposées.
\end{theorem}

\begin{proof}
Si \(E_+\) et \(E_-\) sont respectivement les composantes paire et impaire du
caractère par rapport à \(C_M\), alors
\[
 E_+(C_M\mathbf t)=E_+(\mathbf t),
 \qquad
 E_-(C_M\mathbf t)=-E_-(\mathbf t).
\]
D'où,
\[
 E_{\mathrm{band}}(C_M\mathbf t,-\chi)
 =
 E_{\mathrm{band}}(\mathbf t,\chi),
\]
et l'égalité de masse s'obtient en appliquant l'exponentielle. Dans le terme
électronique,
\[
 (+1)(-22)=(-1)(+22)=-22.
\]
\end{proof}

La clôture finie de la \cref{prop:cierre-antiseccion-ct108} renforce cette
interprétation : sur les quatre-vingt-une routes CT--108, ni le signe pur ni
la réversion pure ne préservent la taxonomie, tandis que
\(-\operatorname{rev}\) la préserve bien, avec
\[
 81=54+9+9+9.
\]
L'antiparticule n'est donc ni une masse négative ni une copie
indépendante. C'est l'autre section orientée de la même bande : elle conserve les
invariants pairs qui déterminent la masse et inverse les invariants impairs qui
déterminent charge et orientation.

\subsection{Théorème de fermeture électronique}

\begin{theorem}[Composition holographique de la section électronique]
\label{thm:cierre-electron-holografico-rev6}
Une fois fixés les objets construits dans les sections précédentes, la section
électronique adimensionnelle est
\[
 \mathfrak e
 =
 \left\|[P_{\Sigma,3},P_{\Pi,3}]\right\|
 \left[
 \exp\bigl((\mathscr R\circ\mathsf J)(\pi,\varphi)\bigr)
 -\frac{\Sigma(\mathrm{APP})}
 {\Sigma(\mathrm{centro})}\,
 \alpha^3
 \right]
 \left[
 1+
 \frac{\Sigma(\mathrm{canal})}
 {\Sigma(\mathrm{centro})}\,
 \Delta_4
 \right].
\]
De manière équivalente,
\[
 \boxed{
 \mathfrak e
 =
 \frac{\sqrt3}{4}
 \left(e^{\varphi/\pi^2}
 -22\alpha^3\right)
 (1+15\Delta_4).}
\]
Son évaluation est
\[
 \boxed{
 \mathfrak e
 =
 0.51099892248806607250987087719134943763\ldots.}
\]
\end{theorem}

\begin{proof}
La \cref{prop:prefactor-app-electron-rev6} apporte
\[
 \left\|[P_{\Sigma,3},P_{\Pi,3}]\right\|=\sqrt3/4.
\]
Le \cref{thm:espejo-electron-rev6} apporte
\[
 (\mathscr R\circ\mathsf J)(\pi,\varphi)=\varphi/\pi^2.
\]
Le \cref{thm:coeficientes-electron-rev6} apporte
\[
 \frac{\Sigma(\mathrm{APP})}{\Sigma(\mathrm{centro})}=22,
 \qquad
 \frac{\Sigma(\mathrm{canal})}{\Sigma(\mathrm{centro})}=15.
\]
La clôture dodécaphasique apporte la fenêtre rationnelle \(\alpha_{12}\) de la
section coinductive \(\alpha\), et
\cref{eq:delta4-electron-rev6} apporte le défaut global. La substitution
produit la formule abrégée. L'évaluation successive de
\[
\begin{aligned}
 e^{\varphi/\pi^2}
 &=1.17814494259630678081160095680888910\ldots,\\
 22\alpha^3
 &=0.00000854906599200551727014979807598228\ldots,\\
 1+15\Delta_4
 &=1.00166794634038210081076702179181019\ldots
\end{aligned}
\]
donne la valeur déclarée.
\end{proof}

Le théorème est holographique en un sens vérifiable : chaque facteur local
retient une couche distincte du système complet. Le commutateur retient la
dualité des feuillets ; l'exponentielle retient l'orientation intérieure de la
charnière \(\pi\)--\(\varphi\) ; \(22\) et \(15\) retiennent le recensement global et
son noyau central ; \(\alpha^3\) retient la clôture fine ; et
\(\Delta_4\) retient le raffinement global. Le centre détermine où
s'ancre la section et la bande détermine comment elle survit sous conjugaison.
La valeur ne remplace pas cette structure : elle est leur évaluation commune.

\subsubsection{Dédoublement de l'action antérieure et postérieure au retour}
\label{sec:electron-desdoblamiento-two-way-rev9}

La réalisation régionale de la coordonnée angulaire conjuguée \(C^*\) distingue deux coordonnées sur la
même droite interne d'action. La coordonnée antérieure au retour est \(H_5\) et
la coordonnée postérieure est
\[
\etaRet
 =H_5-\frac{90}{\pi}\mathscr C_\pi(\alpha).
\]
La notation \(\etaRet\) est la coordonnée sexagésimale canonique définie dans la
branche d'action par \(\etaRet=\Cstar/2-\alphaAngle\). L'égalité précédente est
évaluée dans cette carte ; elle n'introduit pas une seconde coordonnée de retour.
Les cartes additive et multiplicative du dédoublement orienté
sont donc
\begin{equation}
 \boxed{
 \delta_{2w}:=H_5-\etaRet
 =\frac{90}{\pi}\mathscr C_\pi(\alpha),
 \qquad
 R_{\rm act}:=\frac{\hbar^{\rm pre}_{\rm HMT}}
 {\hbar^{\rm ret}_{\rm HMT}}
 =\frac{H_5}{\etaRet}.}
 \label{eq:electron-desdoblamiento-two-way-rev9}
\end{equation}
Les deux coordonnées proviennent d'un unique transport pré-retour/post-retour.
Ce ne sont pas deux constantes de Planck : le facteur commun
\(10^{-34}\mathcal U_S\) fixe la droite dimensionnelle et se simplifie exactement dans
\(R_{\rm act}\). L'ellipse d'action réalise ensuite
\(\operatorname{Area}(\theta)=\hbar_{\rm HMT}^{\rm ret}\theta\), de sorte que
le dédoublement précède la lecture d'holonomie.

Le couplage
\begin{equation}
 \kappa_e^{\beta}
 =80-54\alpha+6\Delta_4,
 \qquad
 \mathfrak e_e^{\beta}
 =\mathfrak eR_{\rm act}^{\kappa_e^{\beta}}
 =0.5109989506900008086\ldots
 \label{eq:electron-beta-two-way-rev9}
\end{equation}
est la clôture quantitative bidirectionnelle de l'électron. Les deux lecteurs
indépendants du registre électronique construits dans le chapitre suivant
sélectionnent, sans consulter la masse de comparaison,
\[
 K_D(\Gamma_e)=K_\Omega(\Gamma_e)=(80,54,6),
\]
et dérivent exactement \(\kappa_e^\beta\). L'entier \(80\) provient de la
première coordonnée de ces lecteurs ; le mot \(w_e=201101\) conserve son
espace de résidence propre dans le canal de la constante d'Euler.

\begin{statusbox}[title={Clôture interne du raffinement bidirectionnel}]
L'identité \eqref{eq:electron-desdoblamiento-two-way-rev9}, la réduction
\((80,54,6)\) et l'équation
\eqref{eq:electron-beta-two-way-rev9} constituent des résultats internes exacts
et sans cible sur le registre électronique. La propagation aux fibres de
rang supérieur se formule au moyen des opérateurs de multisection.
\end{statusbox}

\subsection{Carte énergétique et confrontation de la clôture bidirectionnelle}

Soit \(\mathcal U_E\) une droite réelle de dimension énergétique et soit
\(u_E\) une base choisie. La transduction métrologique est
\[
 \varsigma_{u_E}:\mathbb R\longrightarrow\mathcal U_E,
 \qquad
 \varsigma_{u_E}(x)=x\,u_E.
\]
La carte conserve deux étapes d'une même généalogie. L'étape de base est
\[
 E_e^{(0)}
 =
 \varsigma_{u_E}(\mathfrak e).
\]
Avec \(u_E=1\,\mathrm{MeV}\),
\[
 E_e^{(0)}
 =
 0.5109989224880660725\ldots\ \mathrm{MeV}.
\]
Le rapport d'action et l'exposant dérivé complètent l'étape de référence,
\[
 E_e^\beta=E_e^{(0)}R_{\rm act}^{\kappa_e^\beta}
 =0.5109989506900008086082571648\ldots\ \mathrm{MeV},
 \qquad
 m_e^\beta=E_e^\beta/c^2.
\]
L'unité n'apparaît pas dans la
\cref{eq:electron-holografico-rev6} ; elle est la base d'une droite dimensionnelle
appliquée à un résultat déjà construit.

La référence CODATA 2022 pour l'énergie équivalente de l'électron est
\cite{CODATA2022}
\[
 E_e^{\mathrm{CODATA}}
 =
 0.51099895069(16)\ \mathrm{MeV}.
\]
Dans cette même carte, le résultat bidirectionnel satisfait
\[
 E_e^\beta-E_e^{\mathrm{CODATA}}
 =
 8.086082571648\times10^{-16}\ \mathrm{MeV},
\]
avec pour différence relative
\[
 1.582406883758\times10^{-15}.
\]
La comparaison est effectuée après fixation de l'évaluation interne. Les
identités du sélecteur de base demeurent
\[
 22=\frac{396}{18},
 \qquad
 15=\frac{270}{18},
 \qquad
 p=27-22=5.
\]
et reçoivent en outre la démonstration combinatoire
\(-22=-|K_e\setminus\iota(R_\pi)|\),
\(+15=|R_\pi\times\mathbb F_3|\). La confrontation métrologique vérifie le
résultat ; le parcours, les lecteurs et le rapport d'action l'engendrent.

\subsection{Conséquences structurelles de la clôture électronique}

La lecture structurelle précédente impose plusieurs distinctions :

\begin{enumerate}
\item \(\sqrt3/4\) est la norme du commutateur APP en dimension trois. Sa
coïncidence avec d'autres formules géométriques ne remplace pas cette provenance.

\item \(\pi/\varphi^2\) et \(\varphi/\pi^2\) ne sont pas interchangeables. Le
premier rapport est l'orientation de clôture ; le second, obtenu au moyen de
\(\mathsf J\), est l'orientation électronique d'autocroissance.

\item \(22\) et \(15\) procèdent de \(396/18\) et de \(270/18\).
Les présenter comme des valeurs logiquement sélectionnées par \(Z_0\) inverserait
la dépendance démontrée.

\item \(\alpha^3\) utilise la clôture préalable de
\(\alpha\). La puissance cubique appartient à la loi
électronique ; elle n'autorise pas, à elle seule, une règle universelle
\(\alpha^d\leftrightarrow d\) sans théorème supplémentaire.

\item \(\Delta_4\) est global. Il ne doit pas être remplacé par
\(s_{270}\), par \(135\Delta_4^2\) ou par un résidu particulier de
l'espèce.

\item La position \((5,5)\), l'étiquette \(e_{\mathrm{core}}\), l'origine
fine \(v_e=0\) et la signature brute centrale sont des objets coordonnés, et non
des synonymes.

\item Électron et positron ne se distinguent pas par deux lois de masse : ils sont les
deux sections orientées d'une même bande, de masse paire et de charge impaire.

\item \(0.5109989224\ldots\) est d'abord un scalaire adimensionnel. MeV,
\(\mathrm{MeV}/c^2\), CODATA et \(Z_0\) relèvent de la comparaison
ultérieure.
\end{enumerate}

Ces distinctions expliquent pourquoi l'équation électronique est une
miniature du système HMT--MD. Dans une seule section réapparaissent support,
feuillets, orientation, dénombrement, clôture fine, raffinement, centre, persistance
et conjugaison. Tel est le contenu du mot \emph{holographique} : non
qu'une valeur numérique représente poétiquement le tout, mais que chaque facteur
conserve une dépendance vérifiable à l'égard du tout qui la produit.
 
\section{Dérivation bidirectionnelle de l'exposant électronique à partir du parcours enrichi}
\label{sec:cierre-electron-two-way-rev11}
\index{électron!lecteur bidirectionnel}
\index{calendrier d'ordre 108!exposant électronique}

Le parcours enrichi du point central donne deux représentations
complémentaires de sa mémoire : le mot tritique complet
\(\gamma_{108}\in\mathbb F_3^{108}\) et l'enregistrement dodécaphasique signé
\(\tau_{12}\in\{-1,0,+1\}^{12}\). Les deux lecteurs définis ci-dessous
sont calculés exclusivement à partir de ces représentations, du calendrier et des
constantes HMT déjà construites.

\subsection{Lecteur des déplacements sur le calendrier cyclique d'ordre 108}

Soit \(S\) le déplacement cyclique dans \(\mathbb Z/108\mathbb Z\). Pour
\(k\in\mathbb Z/108\mathbb Z\), nous définissons
\begin{equation}
 D_k(\Gamma):=
 \#\{t\in\mathbb Z/108\mathbb Z:\gamma_t\ne\gamma_{t+k}\}
 =\lVert(I-S^k)\gamma\rVert_0.
 \label{eq:lector-dk-ct108-rev11}
\end{equation}
Le triplet entier du lecteur est
\begin{equation}
 K_D(\Gamma)=
 \left(D_{27}+D_{36},\ D_1,\ \frac{D_4-D_3}{2}\right)
 =:(A_D,B_D,C_D),
 \label{eq:KD-rev11}
\end{equation}
et son évaluation scalaire est
\begin{equation}
 \kappa_D(\Gamma)=A_D-{\alpha}B_D+\Delta_4 C_D.
 \label{eq:kappaD-rev11}
\end{equation}
Les déplacements (27) et (36) sont les relèvements chronologiques de
\(90^\circ\) et de \(120^\circ\) parce que neuf pas discrets représentent
\(30^\circ\). Le terme \(D_1\) mesure la frontière locale immédiate et la
différence \(D_4-D_3\) compare les deux généalogies avant leur quotient de
hauteur. Dans le domaine des 324 parcours, on a
\(\gamma_{t+54}=\gamma_t\) ; tous les \(D_k\) sont pairs et \(C_D\) est entier.

\subsection{Lecteur dodécaphasique des différences de corrélation}

Pour le mot \(\tau=(\tau_r)_{r\in\mathbb Z/12\mathbb Z}\), soient
\[
 C_k^{\rm corr}(\tau)=\sum_{r=0}^{11}\tau_r\tau_{r+k},
 \qquad
 A_\ell(\tau)=\sum_{r=0}^{11}
 \left(\sum_{j=0}^{\ell-1}\tau_{r+j}\right)^2.
\]
L'obstruction généalogique primitive et la mémoire antipodale sont
\begin{align}
 \Omega_{90\mid120}(\tau)
 &=4A_3-3A_4
 =-2C_1^{\rm corr}-4C_2^{\rm corr}-6C_3^{\rm corr},
 \label{eq:omega-two-way-rev11}\\
 P_6(\tau)&=\sum_{r=0}^{5}\tau_r\tau_{r+6}.
 \label{eq:p6-two-way-rev11}
\end{align}
Le symbole \(P_6\) évite la collision avec le cycle court du corridor des
constantes. Le second triplet et son évaluation sont
\begin{equation}
 K_\Omega(\Gamma)=\bigl(\Omega_{90\mid120}(\tau_\Gamma),54,
 P_6(\tau_\Gamma)\bigr),
 \qquad
 \kappa_\Omega=\Omega_{90\mid120}-54\alpha+P_6\Delta_4.
 \label{eq:kappa-omega-rev11}
\end{equation}

\HMTClaim{MD-R-ELECTRON-OMEGA-PRIMITIVE}
\begin{lemma}[Unicité linéaire entière primitive]
\label{lem:unicidad-omega-rev11}
Soit \(L=uA_3+vA_4\), avec \(u,v\in\mathbb Z\), un contraste qui s'annule
lorsque les généalogies descendent à une même hauteur,
\(A_3=3a\), \(A_4=4a\). Alors
\((u,v)=n(4,-3)\). Le signe \(90-120\) étant fixé,
\(\Omega_{90\mid120}=4A_3-3A_4\) est le générateur primitif.
\end{lemma}

\begin{proof}
L'annulation pour tout \(a\) exige \(3u+4v=0\). Puisque
\(\gcd(3,4)=1\), les solutions entières sont
\((u,v)=n(4,-3)\). La primitivité correspond à \(|n|=1\).
\end{proof}

Le lemme sélectionne \(\Omega\) dans la classe des contrastes linéaires
entiers des énergies \(A_3,A_4\). Son domaine d'unicité est cette classe
fonctionnelle ; les invariants non linéaires relèvent d'un problème de sélection
distinct.

\subsection{Évaluation du parcours central}

Le parcours électronique \(\Gamma_e=(5,5,++)\) fournit, avant la
confrontation métrologique,
\[
 (D_1,D_3,D_4,D_{27},D_{36})=(54,56,68,48,32),
\]
et son enregistrement dodécaphasique signé est
\[
 \tau_e=(+,+,+,-,-,-,+,+,+,-,-,-).
\]
Par conséquent,
\[
 K_D(\Gamma_e)=(48+32,54,(68-56)/2)=(80,54,6),
\]
\[
 K_\Omega(\Gamma_e)=(80,54,6).
\]

\HMTClaim{MD-R-ELECTRON-TWOWAY-REDUCTION}
\begin{theorem}[Réduction électronique commune des lecteurs bidirectionnels]
\label{thm:reduccion-electron-two-way-rev11}
Les lecteurs \(K_D\) et \(K_\Omega\) coïncident sur le parcours central et
déterminent
\begin{equation}
 \kappa_e=80-54\alpha+6\Delta_4.
 \label{eq:kappa-electron-demostrado-rev11}
\end{equation}
\end{theorem}

\begin{proof}
Les deux évaluations précédentes donnent le même triplet entier. La substitution
de ses composantes dans
\(A-{\alpha}B+\Delta_4C\) prouve
\eqref{eq:kappa-electron-demostrado-rev11}. Toute quantité employée dans cette
réduction appartient au parcours, à ses deux enregistrements ou aux constantes HMT
préalablement construites.
\end{proof}

\paragraph{Réalisation ultérieure dans la carte énergétique.}
Avec \(R_{\rm act}=H_5/\etaRet\), l'action positive sur la valeur de base
\(\mathfrak e_0\) produit
\begin{equation}
 \mathfrak e_e^\beta
 =\mathfrak e_0R_{\rm act}^{\kappa_e}.
 \label{eq:masa-electron-adimensional-two-way-rev11}
\end{equation}
Après application de la carte actuelle vers les MeV, on obtient
\begin{equation}
 \varsigma_{u_E}(\mathfrak e_e^\beta)\big|_{u_E=1\,\mathrm{MeV}}
 =0.5109989506900008086082571648\ldots\ \mathrm{MeV}.
 \label{eq:masa-electron-two-way-rev11}
\end{equation}
La première égalité appartient au théorème interne conditionné par la valeur de base et
la droite d'action ; la seconde est sa réalisation dimensionnelle. La valeur
externe est réservée à la confrontation métrologique.

L'entier (80) procède de \(D_{27}+D_{36}\) et, dans le lecteur dodécaphasique,
de \(\Omega_{90\mid120}\). La coïncidence avec d'autres cardinaux internes est
un contrôle de compatibilité et ne remplace pas cette provenance opératorielle. Le
parcours d'orientation \((--)\) est une translation temporelle du parcours \((++)\)
et conserve le même profil, conformément à la parité de masse sous conjugaison.

\subsection{Séparation des lecteurs et des échelles de retour}

La coïncidence numérique des composantes internes reste subordonnée à leur
type. Au point central interviennent notamment les objets
\[
 R_{\rm atlas}(5,5)=(24,0,12,0,-12),
 \qquad
 L_{\kappa_0}(\Gamma_e^{++})=(24,0,0,0,-12),
\]
\[
 \begin{gathered}
 k_{\Delta_4}^{\rm dod},k_{\Delta_4}^{\rm pro}:
 \mathscr R_{324}\longrightarrow\mathbb Z,\\
 k_{\Delta_4}^{\rm dod}(\Gamma_e^{++})=12,
 \qquad
 k_{\Delta_4}^{\rm pro}(\Gamma_e^{++})=4,\\
 K_D(\Gamma_e^{++})=K_\Omega(\Gamma_e^{++})=(80,54,6).
 \end{gathered}
\]
Ici, \(\mathscr R_{324}\) désigne le domaine des 324 parcours orientés.
Le premier est une signature statique de famille ; le deuxième, un extracteur de
généalogie ; les troisième et quatrième termes sont des lecteurs torsionnels définis,
respectivement, par la lecture dodécaphasique et par le prolongement prospectif ;
le dernier objet est la paire de lecteurs bidirectionnels. Aucune égalité
partielle n'identifie ces morphismes ni ne permet de transférer un coefficient entre
lecteurs.

On distingue également quatre échelles temporelles : \(27\) est le premier retour
de classe, d'orbite et de cellule ; \(54=D_1(\Gamma_e)\) est le retour cinématique
observable dans le cycle ; \(108\) clôt le calendrier ; et \(216\) clôt l'état cinématique
orienté dans l'extension. Ce typage maintient distincts retour
positionnel, observation et holonomie orientée.

\begingroup\fontsize{10.2}{11.7}\selectfont
\setlength{\parskip}{2pt}\setlength{\abovedisplayskip}{5pt}\setlength{\belowdisplayskip}{5pt}
\paragraph{Chaîne causale du résultat électronique}
Le parcours et ses deux enregistrements déterminent \(K_D\), \(K_\Omega\), leur réduction
commune et l'évaluation de \(\kappa_e\). La réalisation en MeV applique ensuite
la carte énergétique ; la masse externe n'intervient que dans la confrontation. La paire
\(\mathbf K=(K_D,K_\Omega)\) est le résultat opératoriel complet de HMT sur
chaque parcours. Dans une fibre non centrale, sa représentation par un unique scalaire
physique exige un morphisme supplémentaire entre l'opérateur de fibre et la carte
dimensionnelle ; cette exigence ne modifie pas la loi interne et n'autorise pas à choisir
une cellule par proximité métrologique.
 
\paragraph{Passage de la réalisation de Cartan–Holst à la confluence horizon–aire–information par conservation de l’action} La fermeture électronique fixe une réalisation matérielle ; l’étape suivante identifie la grammaire qui convertit action, horloge et masse en grandeurs physiques sans utiliser cette réalisation comme donnée d’entrée.

\par\endgroup
 \endgroup

\chapter{Dérivation déterminantale de l'échelle absolue d'action et valuation \texorpdfstring{\(10^{-34}\)}{10^-34}}
\begingroup
\renewcommand{\chapter}[2][]{}%
\chapter{Dérivation déterminantale de l’échelle d’action : forme normale de Smith, conoyau d’ordre \(16\) et valuation \(10^{-34}\)}
\label{chap:p3-final:34:escala_accion}
\label{ch:vacuum}

\section{Dérivation de l'échelle d'action à partir de l'état dodécaphasique}

La décennie d'action est déterminée intrinsèquement, sans l'adopter comme
unité héritée. La fonctionnelle déterminantale locale s'applique à la sortie HMT
préalablement construite \(\alpha\) et définit
\begin{equation}
\Sigma_H=\varphi\alpha^{16}.
\label{eq:action-determinantal-reader}
\end{equation}
L'exposant \(16\) est l'ordre du conoyau local de la clôture signée ; sa valuation décimale fixe
\begin{equation}
\operatorname{ord}_{10}(\Sigma_H)=-34.
\label{eq:action-decade}
\end{equation}

L'origine de l'exposant peut être vérifiée dans ce volume. Sur
l'orbite locale à quatre feuillets agit le bloc de Hadamard
\begin{equation}
H_4=
\begin{pmatrix}
1&1&1&1\\
1&1&-1&-1\\
1&-1&1&-1\\
1&-1&-1&1
\end{pmatrix}.
\label{eq:action-hadamard}
\end{equation}

\begin{lemma}[Indice local de la droite d'action]
La forme normale de Smith de \(H_4\) est
\begin{equation}
\operatorname{SNF}(H_4)=\operatorname{diag}(1,2,2,4).
\label{eq:action-smith}
\end{equation}
En conséquence,
\[
\operatorname{coker}(H_4)\simeq
\Z/2\Z\oplus\Z/2\Z\oplus\Z/4\Z,
\qquad
|\operatorname{coker}(H_4)|=16.
\]
\end{lemma}

\begin{proof}
Le plus grand commun diviseur des entrées est un ; celui des mineurs
\(2\times2\) non nuls est deux ; celui des mineurs \(3\times3\) est quatre ; et
\(|\det H_4|=16\). Les quotients successifs de ces diviseurs déterminantaux
sont \(1,2,2,4\), qui prouvent \cref{eq:action-smith}. Le produit des
invariants donne l'ordre du conoyau.
\end{proof}

La norme de la section \(\alpha\) sur ces seize feuillets
produit \(\alpha^{16}\) ; l'auto-échelle \(\varphi\) agit
une seule fois sur la base, et non une fois par feuillet. D'où
\cref{eq:action-determinantal-reader}. Les comparaisons certifiées
\begin{equation}
\alpha^{16}<10^{-34}
<\varphi\alpha^{16}<10^{-33}
\label{eq:action-decade-bounds}
\end{equation}
démontrent, sans introduire l'unité SI, que l'ordre décimal est
exactement \(-34\). L'évaluation terminale commence par
\begin{equation}
\Sigma_H
=1.046243838104756580184229208303089\ldots\times10^{-34}.
\label{eq:action-sigma-value}
\end{equation}
Le remplacement du conoyau par trois copies avec exposant \(16^3\), ou
l'application réitérée de \(\varphi\) dans chaque feuillet, modifierait la
généalogie. Les deux opérations constituent des critères de réfutation et non
des conventions équivalentes.
La chaîne causale qui est conservée est donc
\[
\APP\to\TRIT\to\TPK\to U_{12}^{\rm sign}
\to\alpha\to\Sigma_H
\to10^{-34}\U_S\to
\{\hbar_{\rm pre},\hbar_{\rm ret}\}.
\]
Ce volume ne redérive pas \(\alpha\), mais développe la géométrie
métrologique postérieure à la sélection du type et de la décennie par
\(\Sigma_H\). L'entier \(-34\) n'est ni une mantisse ni un ajustement au Système
international, mais la valuation terminale de la fonctionnelle locale.

\section{Sections orientées de la droite d'action}

Les réalisations antérieure et postérieure au retour constituent deux sections
d'une même droite d'action :

\begin{equation}
\hbar_{\rm pre}=H_5\,10^{-34}\U_S,\qquad
\hbar_{\rm ret}=\eta_{\rm ret}\,10^{-34}\U_S.
\label{eq:action-twoway}
\end{equation}

La décennie commune détermine la droite, tandis que l'information
orientée s'exprime par

\begin{equation}
R_{\rm act}=\frac{\hbar_{\rm pre}}{\hbar_{\rm ret}}
=\frac{H_5}{\eta_{\rm ret}},
\qquad
\xi_{\rm act}=\frac12\log R_{\rm act}.
\label{eq:action-rapidity}
\end{equation}

La première quantité fournit une coordonnée multiplicative et la seconde
une coordonnée additive de type rapidité. Elles ne représentent pas deux constantes de
Planck indépendantes, mais les deux orientations d'une même section
par rapport au retour.

Si \(h_a=2\pi\hbar_a\), \(a\in\{\mathrm{pre},\mathrm{ret}\}\), une grandeur proportionnelle à \(h_a^{1/2}\) varie comme \(R_{\rm act}^{1/2}\), une grandeur proportionnelle à \(h_a^{-1/2}\) varie comme \(R_{\rm act}^{-1/2}\), et un quotient où l'action s'annule est un invariant bidirectionnel.

  \endgroup
% Distribución causal exacta del delta Ley 9; contenido conservado sin síntesis.
\Needspace{28\baselineskip}
\section{Bifurcation causale postérieure à \texorpdfstring{\(\alpha\)}{alpha HMT}}
\label{sec:l9s102:bifurcacion-accion-angular}

La généalogie de l'action et du couple angulaire forme un diagramme bifurqué.
La branche régionale produit la mantisse réduite \(\eta_{\mathrm{ret}}\) ; la branche
déterminantale produit l'ordre décimal \(\varepsilon_H=-34\). Les deux
convergent dans la droite d'action. Le couple angulaire est construit dans une branche
sœur à partir de \(\alpha\) et de
\(\eta_{\mathrm{ret}}\) :
\begin{equation}
\begin{tikzcd}[column sep=small,row sep=normal]
\alpha,\pi,\varphi
  \arrow[r] \arrow[d]
& H_5,\mathscr C_\pi
  \arrow[d]
&& \alpha,\varphi,H_4
  \arrow[d] & {}\\
x_A,D_A,\rho_\pi
  \arrow[r]
& \eta_{\mathrm{ret}}
  \arrow[r] \arrow[d]
& \hbar_{\mathrm{ret}}\in\mathcal L_S
& \Sigma_H,\varepsilon_H=-34
  \arrow[l] & {}\\
& (A_{\deg},C^*_{\deg})
  \arrow[r]
& (A_{\mathrm{rad}},C^*_{\mathrm{rad}})
  \arrow[r]
& (q_+,q_-)
  \arrow[r]
& V_{90,120}.
\end{tikzcd}
\label{diag:l9s102:dag-accion-angular}
\end{equation}
L'ordre d'exposition situe la dérivation de Planck avant la publication
angulaire complète ; le diagramme conserve simultanément la dépendance
mathématique exacte de chaque sortie.

\section{Conoyau local et dérivation de l'ordre absolu \texorpdfstring{\(10^{-34}\)}{10^-34}}
\label{sec:l9s102:orden-accion}

La forme normale de Smith du bloc dodécaphasique local est
\begin{equation}
 \operatorname{SNF}(H_4)=\operatorname{diag}(1,2,2,4),
 \qquad
 \operatorname{coker}(H_4)
 \cong\mathbb Z/2\mathbb Z\oplus\mathbb Z/2\mathbb Z
 \oplus\mathbb Z/4\mathbb Z.
 \label{eq:l9s102:snf-h4}
\end{equation}
Le conoyau local est d'ordre \(16\). Pour la section constante
\(s_\alpha(g)=\alpha\), la norme locale et une action de
\(\varphi\) produisent
\begin{equation}
 \boxed{
 \Sigma_H=\varphi\mathcal N_{G_H}(s_\alpha)
 =\varphi\alpha^{16}.}
 \label{eq:l9s102:sigma-h}
\end{equation}
Les inégalités exactes
\begin{equation}
 \alpha^{16}<10^{-34}
 <\varphi\alpha^{16}<10^{-33}
 \label{eq:l9s102:orden-34}
\end{equation}
déterminent
\begin{equation}
 \boxed{
 \varepsilon_H=\left\lfloor\log_{10}\Sigma_H\right\rfloor=-34.}
 \label{eq:l9s102:epsilon-h}
\end{equation}
 
\chapter{Caractère régional de degré cinq et mesure positive d'action}
\begingroup
\renewcommand{\chapter}[2][]{}%
\chapter{Mesure de quotient et caractère positif de l’action : hilbertisation canonique, défaut régional--dodécaphasique et loi fonctorielle sur les chemins}
\label{chap:p3-final:35:medida_accion_positiva}
\label{p3-20260826:chap:medida-accion-positiva}
\index{mesure!catalogue d’émissions}
\index{action!cocycle positif}
\index{trace normalisée}

La section précédente identifie isométriquement un secteur positif de rang trois dans
l’interface orientée \(t=7\), la dodécaphase et le module d’incidence de Witt. La construction
suivante dérive l’autre coefficient de la loi angulaire, \(5/468\), et
démontre pourquoi les deux termes se combinent par une somme positive. La
question exige de séparer
trois opérations :
\begin{enumerate}
  \item compter les présentations microscopiques ;
  \item compter les émissions distinctes après formation du quotient TPK ;
  \item attribuer une action additive à une route.
\end{enumerate}
Les deux premières opérations produisent des mesures différentes. La troisième se
caractérise par une propriété universelle et non par la proximité ultérieure
d’un nombre décimal.

Nous noterons \(\Gtr\) le graphe orienté à \(1944\) états du lecteur de
chemins enrichis. Sa catégorie libre de chemins sera notée
\(\mathsf P(\Gtr)\). La définition locale de ses arêtes et de son action est
développée dans la section correspondante.

\section{Espace de germes, quotient d’émissions et fibres régionales}

Soit
\[
 \Omega_{\rm seed}
 =\bigl((\mathbb Z/9\mathbb Z)^2\times D_4\bigr)^2,
 \qquad
 |\Omega_{\rm seed}|=104\,976,
\]
l’ensemble exécutable des paires de germes de l’émetteur fini, et soit
\[
 F:\Omega_{\rm seed}\twoheadrightarrow\mathcal U_6,
 \qquad
 |\mathcal U_6|=468,
\]
l’application qui conserve l’émission sextuple distincte. Le recensement exact des
fibres est
\[
\begin{array}{c|c|c}
|F^{-1}(U)|&\#\{U\text{ de cette taille}\}&\text{germes apportés}\\
\hline
144&432&62\,208\\
432&18&7\,776\\
1\,944&18&34\,992.
\end{array}
\]
Les trois contributions ont pour somme \(104\,976\). Par conséquent, \(F\) n’est pas un
revêtement régulier : des émissions différentes possèdent des nombres différents de
présentations microscopiques.

La microfibre de \(\pi\) contient cinq émissions distinctes. Quatre ont
pour multiplicité \(144\) et une a pour multiplicité \(432\). Cette donnée suffit
à démontrer que la mesure uniforme sur les germes et la mesure uniforme
sur les émissions ne sont pas identiques.

\section{\(2\) mesures exactes sur le quotient des émissions et des mots}

Si chaque germe reçoit le poids \(1/104\,976\), la mesure image par \(F\) est
\[
 \nu_{\rm seed}(U)
 =\frac{|F^{-1}(U)|}{104\,976}.
\]
Ses atomes possibles sont
\[
 \frac1{729},\qquad\frac1{243},\qquad\frac1{54}.
\]
\Needspace{5\baselineskip}
En particulier,
\[
 \nu_{\rm seed}(\mathscr R_\pi^{\rm micro})
 =\frac{4\cdot144+432}{104\,976}
 =\boxed{\frac7{729}}.
\]

Si l’état élémentaire est, en revanche, une émission distincte
\(U\in\mathcal U_6\), l’algèbre diagonale des observables est
\[
 \mathcal D_{468}=\ell^\infty(\mathcal U_6)
\]
et son état de comptage normalisé est
\[
 \tau_{468}(f)
 =\frac1{468}\sum_{U\in\mathcal U_6}f(U).
\]
Soit \(\Pi_\pi^{(468)}\) la fonction caractéristique des cinq régions de la
microfibre. Alors
\[
 \boxed{\tau_{468}(\Pi_\pi^{(468)})=\frac5{468}.}
\]
La différence exacte entre les deux lectures est
\[
 \frac5{468}-\frac7{729}
 =\frac{41}{37\,908}.
\]
Ce ne sont pas deux approximations. La première mesure compte les présentations ; la
seconde compte les classes d’émission.

\begin{theorem}[Unicité de la mesure uniforme du catalogue]
\label{p3-20260826:thm:unicidad-medida-468}
L’unique probabilité sur \(\mathcal U_6\) invariante sous tout
réétiquetage de ses \(468\) atomes est la mesure uniforme.
De manière équivalente, l’unique trace positive et unitaire sur
\(M_{468}(\mathbb C)\) est
\[
 \tau_{468}(A)=\frac1{468}\operatorname{Tr}A.
\]
\end{theorem}

\begin{proof}
Les transpositions envoient tout atome sur tout autre. L’invariance
impose qu’ils aient tous la même masse, et la normalisation la fixe à
\(1/468\). Dans la version matricielle, l’invariance unitaire égalise la valeur
de tous les projecteurs de rang un. La décomposition spectrale et la
linéarité reconstruisent la trace normalisée.
\end{proof}

L’hypothèse décisive est manifeste : le lecteur travaille avec le
\emph{catalogue quotient}. Le théorème n’affirme pas que la mesure uniforme des
germes devienne uniforme lorsqu’on oublie les fibres. La symétrie utilisée
est la naturalité de l’objet fini sans étiquettes, et non une propriété
ergodique attribuée sans preuve à la dynamique microscopique.

\section{Hilbertisation canonique du quotient}

La relation entre germes et émissions peut être formulée sans ambiguïté.
Soient
\[
 H_{\rm seed}=\ell^2(\Omega_{\rm seed}),
 \qquad
 H_U=\ell^2(\mathcal U_6).
\]
L’émetteur détermine
\[
 \boxed{
 J_F\delta_U
 =|F^{-1}(U)|^{-1/2}
 \sum_{\omega\in F^{-1}(U)}\delta_\omega.
 }
\]
Les fibres distinctes sont disjointes et chaque colonne a pour norme un ; par
conséquent,
\[
 J_F^*J_F=I_{H_U}.
\]
De plus, pour tout observable \(f\) du catalogue,
\[
 M_{f\circ F}J_F=J_FM_f.
\]
L’image \(J_FH_U\) est exactement le sous-espace des fonctions constantes
sur chaque fibre.

Si \(F^\sharp:H_U\to H_{\rm seed}\) désigne l’application induite par
précomposition, sans normalisation,
alors
\[
 J_F
 =F^\sharp\bigl((F^\sharp)^*F^\sharp\bigr)^{-1/2}.
\]
C’est le facteur isométrique de la décomposition polaire de cette application. Cette
expression démontre sa naturalité sous les isomorphismes et les réétiquetages de
l’émetteur. Elle n’implique pas une fonctorialité fausse sous toute composition de
surjections non régulières : la normalisation dépend des fibres de
l’application concrète.

Dans le langage des mesures, le relèvement de la trace du quotient qui est
uniforme à l’intérieur de chaque fibre vaut
\[
 \widetilde\tau(\omega)
 =\frac1{468\,|F^{-1}(F\omega)|}.
\]
Chaque fibre reçoit une masse totale \(1/468\). La compensation inverse par
la multiplicité, et non une invocation informelle des orbites, constitue la preuve exacte.

\begin{corollary}[Canonicité de \(5/468\)]
\label{p3-20260826:cor:canonicidad-cinco-468}
Une fois le lecteur régional défini sur les émissions distinctes, la masse
de la microfibre de \(\pi\) est imposée :
\[
 \tau_{468}(\Pi_\pi^{(468)})
 =\frac{\operatorname{rank}\Pi_\pi^{(468)}}{468}
 =\frac5{468}.
\]
\end{corollary}

\section{Défaut positif entre les mesures régionale et dodécaphasique}

Écrivons
\[
 H_{\rm ang}=H_U\otimes V_{11}.
\]
Le scalaire interne
\[
 \Delta_4
 =\frac{\pi-e\log\pi}{270}
 \left(1+\frac\pi{729}\right)
 =0.0001111964226921402\ldots
\]
est construit avec les lectures HMT déjà fixées de \(\pi\) et \(e\), et avec les
entiers \(270\) et \(729\) du registre. Il constitue ici une entrée mathématique
prédimensionnelle : ni unités ni observables physiques n’interviennent.

La construction interne du scalaire ne détermine pas à elle seule quel opérateur
il doit multiplier. Pour garder cette prémisse visible, nous déclarons :
\begin{description}
  \item[\HDef --- défaut de changement de feuillet.] La coordonnée dodécaphasique du défaut
  de changement de feuillet est \(\Delta_4P_3\).
\end{description}
\HAdd fixera la forme additive de l’action, mais ne sélectionne pas \HDef : pour tout
\(c\ge0\), le même argument admettrait la paire
\((\Pi_\pi^{(468)},cP_3)\). Démontrer \HDef directement à partir d’une transition
totale du noyau topologique de phase enrichi est une condition supplémentaire.

Nous définissons la somme de Kronecker
\[
 \boxed{
 D_{\rm switch}^{\rm dod}
 =\Pi_\pi^{(468)}\otimes I_{11}
 +\Delta_4 I_{468}\otimes P_3.
 }
\]
Les deux termes sont positifs, commutent et agissent sur des facteurs différents.
Par conséquent,
\[
 D_{\rm switch}^{\rm dod}\ge0.
\]
L’entrelaceur \(U_W\) du
\cref{thm:entrelazador-dodecafase-witt} produit la carte de Witt :
\[
\begin{aligned}
 D_{\rm switch}^{\rm Witt}
 &=(I_{468}\otimes U_W)
 D_{\rm switch}^{\rm dod}
 (I_{468}\otimes U_W)^*\\
 &=\Pi_\pi^{(468)}\otimes I_{E_{30}}
 +\Delta_4I_{468}\otimes Q_3.
\end{aligned}
\]
Ainsi, le même opérateur positif se lit dans la dodécaphase et dans le module
d’incidence.

La chaîne isométrique orientée identifie également les secteurs actifs
\[
 I_{L_{\rm or}}\otimes P_G,\qquad P_3,\qquad Q_3.
\]
On n’affirme pas une équivalence des compléments complets des trois
modules : \(W_{GK}\) est une équivalence partielle des rangs
sélectionnés. C’est exactement l’information nécessaire pour identifier
la seconde contribution de la trace.

\begin{theorem}[Portée exacte de l’entrelacement et de l’induction de l’action]
\label{p3-20260826:thm:alcance-induccion-accion}
Soient \(W_{GK}\) et \(U_W\) les isométries partielles du
\cref{thm:identificacion-isometrica-tres-proyectores}. Alors :
\begin{enumerate}
  \item pour tout \(c\geq0\),
  \[
   W_{GK}^{*}(cP_3)W_{GK}
   =c(I_{L_{\rm or}}\otimes P_G),
   \qquad
   U_W(cP_3)U_W^{*}=cQ_3;
  \]
  par conséquent, l’entrelacement transporte le coefficient, mais ne le
  sélectionne pas ;
  \item si \(P_{\rm sgn}\) et \(P_{\rm std}\) sont les projecteurs centraux
  de la décomposition
  \[
   P_3V_{11}\big|_{H_*}
   \simeq \operatorname{sgn}_{H_*}\oplus\operatorname{std},
   \qquad
   \operatorname{rank}P_{\rm sgn}=1,
   \quad
   \operatorname{rank}P_{\rm std}=2,
  \]
  tout opérateur positif autoadjoint supporté sur \(P_3V_{11}\) et
  \(H_*\)-équivariant a la forme
  \[
   D_{a,b}=aP_{\rm sgn}+bP_{\rm std},
   \qquad a,b\geq0.
  \]
  Même l’imposition de la trace de \(\Delta_4P_3\) laisse la famille
  \[
   a+2b=3\Delta_4.
  \]
  L’isotropie \(a=b\) ne se déduit pas de la symétrie résiduelle \(S_3\) ;
  \item la probabilité uniforme du catalogue,
  \(\mu_U(U)=1/468\), est la trace canonique du quotient fini et son
  relèvement compensé est celui que détermine \(J_F\). Ce n’est pas la
  mesure image de la probabilité uniforme sur les germes, dont les atomes sont
  \(1/729,1/243,1/54\). En particulier,
  \[
   \mu_U(\mathscr R_\pi^{\rm micro})=\frac5{468},
   \qquad
   F_\#\mu_{\rm seed}(\mathscr R_\pi^{\rm micro})=\frac7{729};
  \]
  \item la loi complète
  \[
   \delta_*=\frac5{468}+\frac3{11}\Delta_4,
   \qquad
   \lambda_*=1-\delta_*,
   \qquad
   a_*(e)=g(e)+\lambda_*s(e)
  \]
  est un théorème relatif aux cinq clauses \HTr--\HAdd--\HDef--\HRes--\HLoc. Le moteur
  exécutable du noyau topologique de phase réduit n’induit pas, par revêtement orienté, l’opérateur ramifié
  \Gtr : la dynamique associée revient à l’identité après \(54\) itérations,
  tandis que chaque état de
  \Gtr possède quatre successeurs distincts et aucune boucle.
\end{enumerate}
Par conséquent, les résultats précédents démontrent la canonicité des identifications
isométriques, de la trace sur le catalogue quotient et du caractère additif
dans le cadre de leurs prémisses. L’affirmation plus forte selon laquelle une transition totale
du noyau topologique de phase enrichi induit nécessairement \(\Delta_4P_3\), la mesure uniforme et la
loi locale exige encore le théorème dynamique supplémentaire spécifié dans
l’analyse de l’induction.
\end{theorem}

\begin{proof}
Les identités du premier point s’obtiennent en multipliant
\(W_{GK}^{*}W_{GK}=I_{L_{\rm or}}\otimes P_G\),
\(W_{GK}W_{GK}^{*}=P_3\) et \(U_WP_3U_W^{*}=Q_3\) par le scalaire \(c\).
Pour le second point, la restriction à \(H_*\simeq S_3\) est une somme sans
multiplicités de deux représentations irréductibles non isomorphes. Le lemme de
Schur diagonalise tout opérateur autoadjoint équivariant dans ces deux blocs ;
la positivité équivaut à \(a,b\geq0\). Sa trace normalisée est
\((a+2b)/11\), ce qui prouve la liberté résiduelle même après fixation de la
trace. L’action complète de \(A_5\) sur le triplet irréductible réduirait
la famille à \(cP_3\), mais ne fixerait pas non plus la valeur de \(c\).

Le troisième point découle du recensement des fibres : \(F\) a pour multiplicités
\(144,432,1944\), de sorte que la mesure image pondère une émission par son
nombre de présentations. La trace \(\tau_{468}\), en revanche, attribue le
même poids à chaque classe ; \(J_F\) réalise exactement son relèvement
normalisé. Le quatrième point est l’enchaînement logique déjà démontré :
\HTr fixe la trace du quotient, \HAdd la somme des traces, \HDef déclare
\(\Delta_4P_3\), \HRes définit la capacité résiduelle et \HLoc la porte sur les arêtes.
L’énumération finie prouve l’obstruction entre le noyau réduit et
\Gtr : un revêtement orienté préserverait l’étoile sortante et une
semi-conjugaison du retour identité exigerait des boucles dans l’image,
en contradiction avec les deux recensements de \Gtr.
\end{proof}

Le spectre complet de \(D_{\rm switch}^{\rm dod}\) est
\[
\begin{array}{c|c}
\text{valeur propre}&\text{multiplicité}\\ \hline
0&(468-5)(11-3)=3704\\
\Delta_4&(468-5)3=1389\\
1&5(11-3)=40\\
1+\Delta_4&5\cdot3=15.
\end{array}
\]
Par conséquent,
\[
\begin{aligned}
 \tau_{468\cdot11}(D_{\rm switch})
 &=\tau_{468}(\Pi_\pi^{(468)})+\Delta_4\tau_{11}(P_3)\\
 &=\boxed{\frac5{468}+\frac3{11}\Delta_4}.
\end{aligned}
\]
La même identité vaut dans la carte de Witt. Dans la carte du secteur orienté, le terme
\(3/11\) est la trace de \(P_G\) sur l’espace homogène à onze branches.

\section{Propriété universelle de la fonctionnelle positive d’action}

Le mot \emph{action} est employé au sens additif : la concaténation de deux
routes reçoit la somme de leurs incréments. Pour séparer ce contenu d’une
formule choisie, nous définissons d’abord la catégorie dans laquelle la loi sera unique.

\begin{definition}[Caractère positif d’action HMT]
\label{p3-20260826:def:caracter-accion-hmt}
Un caractère local d’action sur le produit
régional--dodécaphasique est une application
\[
 \delta:
 M_{468}(\mathbb C)_+\oplus M_{11}(\mathbb C)_+
 \longrightarrow\mathbb R_{\ge0}
\]
qui satisfait :
\begin{enumerate}
  \item additivité exacte dans le cône produit ;
  \item homogeneidad positiva;
  \item invariance par conjugaison unitaire indépendante dans les deux
  facteurs ;
  \item normalisation bifactorielle,
  \[
  \delta(I_{468},0)=\delta(0,I_{11})=1.
  \]
\end{enumerate}
\end{definition}

L’additivité traduit la concaténation ; l’invariance élimine les
coordonnées ; la normalisation bifactorielle fixe une échelle dans chaque terme.
Il ne s’agit pas d’une fonctionnelle unitaire sur l’algèbre somme directe : son unité
est \((I_{468},I_{11})\) et l’additivité donne
\(\delta(I_{468},I_{11})=2\). Aucune clause ne contient la valeur d’un
angle.

\begin{theorem}[Unicité du caractère positif et capacité résiduelle relative]
\label{p3-20260826:thm:ley-accion-hmt}
Dans la classe de la définition précédente, il existe un unique caractère positif
d’action HMT\@. C’est
\[
 \boxed{\delta(A,B)=\tau_{468}(A)+\tau_{11}(B).}
\]
Appliqué au défaut déclaré par \HDef,
\((\Pi_\pi^{(468)},\Delta_4P_3)\), il produit
\[
 \delta_*=\frac5{468}+\frac3{11}\Delta_4.
\]
Si l’on déclare en outre \HRes, selon laquelle la capacité autoduale de changement de
feuillet est normalisée à un et le défaut la réduit par
\(\lambda=1-\delta\), la capacité résiduelle dans le cadre de \HTr--\HAdd--\HDef--\HRes est
\[
 \boxed{
 \lambda_*
 =1-\delta_*
 =1-\frac5{468}-\frac3{11}\Delta_4.
 }
\]
\end{theorem}

\begin{proof}
L’additivité sépare les facteurs :
\[
 \delta(A,B)=\delta(A,0)+\delta(0,B).
\]
Chaque restriction est positive, homogène et invariante par conjugaison. Pour
un projecteur de rang un, l’invariance unitaire donne une valeur commune.
L’additivité sur une décomposition spectrale impose que chaque restriction
soit un multiple de la trace. La normalisation bifactorielle fixe ces multiples à
\(1/468\) et \(1/11\). Il en résulte la somme des traces. La dernière formule
applique la normalisation déclarée de la capacité résiduelle.
\end{proof}

Numériquement,
\[
 \boxed{\lambda_*=0.9892859130191415\ldots},
 \qquad
 0<\lambda_*<1.
\]
La positivité découle déjà de
\(0<\Delta_4<10^{-3}\) et \(5/468<1\). Elle ne dépend pas d’une comparaison
métrologique.

\section{Foncteur d’action positive sur la catégorie des chemins compatibles}

Rappelons le graphe orienté \(\Gtr\) et sa catégorie libre
\(\mathsf P(\Gtr)\). Pour une arête
\(e:x\to y\), nous définissons
\[
 g(e)=\mathbf1_{\{d_y\ne d_x\}},
 \qquad
 s(e)=\mathbf1_{\{m_y\ne m_x\}},
\]
où \(g\) enregistre une rotation et \(s\) un changement de feuillet. La flèche du
scalaire \(\lambda_*\) vers les arêtes ne se déduit pas du
théorème des traces. Elle est déclarée par la clause \HLoc : une rotation a un coût
unitaire, un changement de feuillet a pour coût \(\lambda_*\), et les deux indicateurs
s’additionnent lorsqu’ils apparaissent sur la même arête. Sous \HLoc, l’action locale est
\[
 \boxed{a_*(e)=g(e)+\lambda_*s(e).}
\]
Pour une route \(\gamma=e_1\cdots e_n\),
\[
 \mathcal A_*(\gamma)
 =\sum_{j=1}^na_*(e_j)
 =\#\operatorname{giros}(\gamma)
 +\lambda_*\#\operatorname{cambios\ de\ hoja}(\gamma).
\]
Pour la route vide \(\mathbf1_x\), on fixe
\(\mathcal A_*(\mathbf1_x)=0\). Alors
\[
 \boxed{
 \mathcal A_*(\gamma_1\gamma_2)
 =\mathcal A_*(\gamma_1)+\mathcal A_*(\gamma_2).
 }
\]
Par conséquent,
\[
 \mathcal A_*:\mathsf P(\Gtr)\longrightarrow(\mathbb R_{\ge0},+)
\]
est un foncteur monoïdal, ou cocycle additif de chemins.

Les indicateurs \(g\) et \(s\) sont insensibles au sens de parcours.
L’orientation n’est pas perdue : elle réside dans la droite \(L_{\rm or}\) et dans
l’involution étoile, tandis que le coût positif réside dans les
projecteurs. Confondre ces deux données introduirait des signes négatifs dans
l’action.

Pour \(\beta>0\), l’opérateur de transfert associé est
\[
 (\mathcal L_{\beta,*}f)(y)
 =\sum_{e:x\to y}e^{-\beta a_*(e)}f(x).
\]
Tous ses poids sont strictement positifs. Comme le potentiel dépend seulement
des deux indicateurs locaux, il commute avec l’action grossière
\((\mathbb Z/3\mathbb Z)^3\) et respecte la décomposition en \(27\)
secteurs de Bloch utilisée dans la section suivante.

Une dynamique déterministe finie non pondérée agit par des matrices de
permutation et possède des valeurs singulières égales à un. Le cocycle
\(\mathcal A_*\) transforme l’inventaire des chemins en un opérateur de
transfert qui préserve le cône positif, c’est-à-dire en une matrice à
poids strictement positifs sur les arêtes autorisées. Cette propriété ne
démontre pas à elle seule la contraction, la simplicité de la valeur propre dominante ni
la séparation spectrale ; ces conclusions exigent une normalisation ou un
argument spectral supplémentaire.

\ifdefined\HMTInternalTraceability
\section{Induction par le quotient et extension dynamique}

L’expression \emph{induite par le noyau topologique de phase enrichi} possède deux sens qui doivent rester
séparés.

\begin{enumerate}
  \item \textbf{Induction par l'objet quotient.} L'émetteur exécutable détermine \(F\), son facteur polaire \(J_F\), l'espace \(H_U\) et la trace normalisée. L'état holonomique intégral apporte \(\Pi_\pi^{(468)}\), \(P_3\), \(K\), \(B_K\), le repère de Witt et \(\Delta_4\). En ce sens catégorique, \(D_{\rm switch}\) se construit uniquement avec des données internes APP--TRIT--TPK\@.
  \item \textbf{Induction par une transition totale.} Cette lecture exigerait
  une transition finie exécutable sur toutes les coordonnées de
  \(\mathcal X_{\mathrm{TPK}}^{\mathrm{enr}}\) et une preuve que sa
  mesure image dynamique
  préserve \(J_FH_U\) et produit le potentiel \(a_*\). Cette transition totale
  n’appartient pas au domaine exécutable construit dans cette section.
\end{enumerate}

Le \cref{p3-20260826:thm:ley-accion-hmt} et la construction des chemins closent la première
lecture au moyen de cinq clauses explicites du lecteur :
\begin{description}
  \item[\HTr --- trace du quotient.] Lorsque la multiplicité de présentation
  reste dans la fibre oubliée, une émission distincte est un atome et
  l’état du catalogue est la trace normalisée de son algèbre finie.
  \item[\HAdd --- cocycle d’action.] Le défaut d’une route est le caractère
  positif, normalisé bifactoriellement, additif et invariant du produit
  régional--dodécaphasique.
  \item[\HDef --- défaut de changement de feuillet.] Le second argument du caractère est
  \(\Delta_4P_3\).
  \item[\HRes --- capacité résiduelle.] La capacité autoduale est normalisée à
  un et le défaut la réduit selon \(\lambda=1-\delta\).
  \item[\HLoc --- action locale.] Sur les arêtes de \Gtr, on fixe
  \(a(e)=g(e)+\lambda s(e)\), et l’action d’une route est la somme sur ses
  arêtes.
\end{description}
\HTr impose \(5/468\). \HAdd impose la somme des traces et exclut les termes
mixtes. \HDef fixe le scalaire qui occupe la coordonnée de commutation. \HRes transforme
le défaut en capacité résiduelle et \HLoc transporte cette capacité vers les arêtes.
Ainsi, \(\delta_*\) est une conséquence de \HTr--\HAdd--\HDef, \(\lambda_*\) l’est de
\HTr--\HAdd--\HDef--\HRes et \(\mathcal A_*\) l’est de
\HTr--\HAdd--\HDef--\HRes--\HLoc. \HTr et \HAdd possèdent des caractérisations universelles ; \HDef, \HRes et
\HLoc sont des identifications constitutives internes ; elles ne se déduisent pas de la
transition finie précédente.

\section{Familles compatibles avec les symétries partielles}

\subsection{Symétrie du cycle \(1/2\) et famille compatible de mesures}

L’échange des deux moitiés d’une émission décompose les \(468\)
émissions en \(234\) paires. La mesure uniforme du catalogue comme les
mesures images des germes sont invariantes sous cet échange. Il en va de même
pour toute attribution constante sur chaque paire. L’invariance du retour,
à elle seule, ne prouve pas l’uniformité.

\subsection{Symétrie de \texorpdfstring{\Gtr}{Gamma tr} et famille de coûts}

Les \(7776\) arêtes de \Gtr se répartissent en \(288\) orbites de taille
\(27\) sous \((\mathbb Z/3\mathbb Z)^3\). Les \(2160\) arêtes de
changement de feuillet forment \(80\) orbites :
\[
 80=5\cdot16.
\]
La symétrie permet d’attribuer des coûts différents à ces orbites. Le coût
scalaire commun provient de \HAdd, et non de la cardinalité.

\subsection{Liberté scalaire de la coordonnée \(2^{\mathrm a}\)}

Pour tout \(c\ge0\), l’opérateur
\[
D_c=\Pi_\pi^{(468)}\otimes I_{11}
+cI_{468}\otimes P_3
\]
est positif et le caractère unique de \HAdd donne
\[
\tau(D_c)=\frac5{468}+\frac3{11}c.
\]
La formule interne de \(\Delta_4\) construit un candidat distingué, mais
la propriété universelle de la trace ne démontre pas à elle seule
\(c=\Delta_4\). Cette égalité est précisément le contenu de \HDef.

\subsection{Positivité, linéarité et termes d’ordre \(2\)}

Les fonctions
\[
 1-u-v+cuv,\qquad(1-u)(1-v),\qquad e^{-(u+v)}
\]
peuvent partager la normalisation et le premier ordre. L’absence de \(uv\) se
déduit de l’additivité exacte. Sans \HAdd, la loi linéaire ne serait pas unique.

\section{Hypothèses d’existence, de positivité, de naturalité et d’extension dynamique}

Le résultat dépend des conditions vérifiables suivantes :
\begin{enumerate}
  \item le catalogue ne contient pas \(468\) émissions ou la microfibre n’a pas
  rang cinq ;
  \item le recensement des fibres n’est pas
  \(432\times144,\ 18\times432,\ 18\times1944\) ;
  \item \(J_F\) n’est pas isométrique ou n’entrelace pas les observables du quotient ;
  \item l’un des \(P_G,P_3,Q_3\) cesse d’être un projecteur de rang trois ;
  \item il existe un autre caractère positif, normalisé bifactoriellement, additif et
  invariant qui n’est pas la somme des traces ;
  \item une transition totale du noyau topologique de phase enrichi, une fois construite, induit un
  coefficient de commutation différent de \(\Delta_4\) ;
  \item \(D_{\rm switch}\) acquiert une valeur propre négative ;
  \item l’action accumulée cesse d’être additive sous concaténation.
\end{enumerate}

Les certificats reconstruisent les fibres, les deux mesures, le facteur polaire, le spectre du défaut, les orbites d'arêtes et les ablations. La valeur de \(\lambda_*\) apparaît à la fin comme conséquence de ces objets. \fi \endgroup
% Distribución causal exacta del delta Ley 9; contenido conservado sin síntesis.
\section{Caractère régional de degré \texorpdfstring{\(5\)}{5} et mantisse réduite d'action}
\label{sec:l9s102:eta-ret}

Le bloc central
\begin{equation}
 \mathcal B_c=\begin{pmatrix}7&2\\2&7\end{pmatrix}
 \label{eq:l9s102:bc-accion}
\end{equation}
possède les valeurs propres \(9\) et \(5\). La longueur \(4\) des orbites de pas
\(3\), le rang \(12\) du cycle dodécaphasique et
\(104976/1944=54\) fixent le caractère
\begin{equation}
 \boxed{
 H_5(\alpha,\varphi)
 =\frac12\sqrt{\frac{1000\alpha}{\varphi}}-\alpha
 +\frac9{16}\alpha^2-\frac59\alpha^3
 +\frac7{48}\alpha^4-\frac1{54}\alpha^5.}
 \label{eq:l9s102:h5}
\end{equation}
La complétion nonadique construit la coordonnée analytique et le déterminant
régional
\begin{equation}
 x_A=\frac{50\pi}{9}\alpha,
 \qquad
 D_A=e^{-2x_A}
 =\exp\!\left(-\frac{100\pi}{9}\alpha\right).
 \label{eq:l9s102:xa-da}
\end{equation}
La microfibre apporte \(\rho_\pi=169\), et sa série absolument
convergente est
\begin{equation}
 \mathscr C_\pi(\alpha)
 =169\alpha^6+\frac{D_A\alpha^7}{1-D_A\alpha}.
 \label{eq:l9s102:cpi}
\end{equation}
La mantisse réduite postérieure au retour est alors construite par
\begin{equation}
 \boxed{
 \eta_{\mathrm{ret}}
 =H_5-\frac{90}{\pi}\mathscr C_\pi(\alpha)
 \in\mathbb R_{>0}.}
 \label{eq:l9s102:eta-ret}
\end{equation}
Son membre de droite contient exclusivement des objets engendrés auparavant
par APP--TRIT--TPK et par la clôture de \(\alpha\).

 
\chapter{Dérivation autoduale de la constante de Planck et de sa droite d'action}
\begingroup
\renewcommand{\chapter}[2][]{}%
\chapter{Dérivation autoduale de la constante de Planck : périodicité temporelle de \(108\) phases, action et longueur de Compton}
\label{chap:p3-final:36:planck_autodual}
\label{chap:planck-autodual}

\section{Réalisation de Compton associée à la périodicité temporelle de \(108\) phases}

Sur le cycle de \cref{eq:realizacion-circular-ct108}, fixons une section positive d'action \(\hbar\). Le quantum angulaire associé est
\begin{equation}
E_{108}=\hbar\omega_{108},
\qquad
m_{108}=\frac{E_{108}}{c^2}
=\frac{\hbar\omega_{108}}{c^2}.
\label{eq:cuanto-ct108}
\end{equation}
\(G\) n’a pas encore été utilisé. La géométrie du cycle et la quantification
angulaire suffisent à fixer ses deux lecteurs de Compton.

\begin{theorem}[Identité entre la périodicité temporelle de \(108\) phases et la longueur de Compton]
\label{thm:ct108-compton}
Pour le quantum \eqref{eq:cuanto-ct108},
\begin{equation}
\lambdabarC(m_{108})=\frac{\hbar}{m_{108}c}=R_{108},
\qquad
\lambdaC(m_{108})=\frac{h}{m_{108}c}=C_{108}.
\label{eq:identidad-ct108-compton}
\end{equation}
Par conséquent, le rayon du cycle est sa longueur de Compton réduite et la
longueur d’un tour est sa longueur de Compton complète.
\end{theorem}

\begin{proof}
En substituant \(m_{108}=\hbar\omega_{108}/c^2\),
\[
\frac{\hbar}{m_{108}c}
=\frac{\hbar c^2}{\hbar\omega_{108}c}
=\frac{c}{\omega_{108}}=R_{108}.
\]
Comme \(h=2\pi\hbar\), la deuxième égalité est \(2\pi R_{108}=C_{108}\).
\end{proof}

La coïncidence n’est pas un ajustement numérique. C’est une identité homogène : elle vaut
pour tout \(\hbar,c,t_0>0\) et dépend seulement du fait que la même période détermine
la fréquence, la longueur du cycle et le quantum d’action.

\section{Transducteur gravitationnel et équation de sélection autoduale}

La droite gravitationnelle HMT--MD contient le transducteur
\begin{equation}
G(\theta)=\theta\frac{c^5t_0^2}{\hbar},
\qquad \theta>0.
\label{eq:transductor-g-planck-anexo}
\end{equation}
L’homogénéité est exacte, mais ne sélectionne pas à elle seule \(\theta\). Le
cycle CT108 apporte une condition interne supplémentaire : le lecteur gravitationnel
doit renvoyer la même échelle radiale que le lecteur géométrique et le lecteur de
Compton.

Nous définissons le rayon gravitationnel réduit
\begin{equation}
\rg(\theta):=\frac{G(\theta)m_{108}}{c^2}.
\label{eq:radio-gravitatorio-reducido}
\end{equation}
En substituant \eqref{eq:cuanto-ct108} et
\eqref{eq:transductor-g-planck-anexo},
\begin{equation}
\rg(\theta)
=\theta\frac{\pi}{54}\ell_0.
\label{eq:rg-theta-lineal}
\end{equation}

\begin{theorem}[Sélecteur autodual de la périodicité temporelle de \(108\) phases]
\label{thm:selector-autodual-ct108}
L’équation
\[
\rg(\theta)=R_{108}
\]
possède une unique solution positive :
\begin{equation}
\boxed{\thetaStar=\left(\frac{54}{\pi}\right)^2.}
\label{eq:theta-star}
\end{equation}
\end{theorem}

\begin{proof}
Par \eqref{eq:realizacion-circular-ct108} et
\eqref{eq:rg-theta-lineal}, l’équation est
\[
\theta\frac{\pi}{54}\ell_0
=\frac{54}{\pi}\ell_0.
\]
Comme \(\ell_0>0\), on simplifie et l’on obtient
\(\theta=(54/\pi)^2\). La fonction
\(\theta\mapsto\rg(\theta)\) est strictement croissante sur
\(\RR_{>0}\), de sorte qu’il n’existe aucune autre solution positive.
\end{proof}

\begin{corollary}[Unité de Planck interne]
\label{cor:unidad-planck-interna}
Avec \(G_*=G(\thetaStar)\), les sections de Planck satisfont
\begin{equation}
\ellP:=\sqrt{\frac{\hbar G_*}{c^3}}=R_{108},
\qquad
\tP:=\sqrt{\frac{\hbar G_*}{c^5}}=\frac1{\omega_{108}},
\label{eq:planck-long-time-ct108}
\end{equation}
\begin{equation}
\mP:=\sqrt{\frac{\hbar c}{G_*}}=m_{108},
\qquad
\EP:=\mP c^2=E_{108}.
\label{eq:planck-mass-energy-ct108}
\end{equation}
En particulier,
\begin{equation}
\boxed{
R_{108}=\lambdabarC=\rg=\ellP,
\qquad
C_{108}=\lambdaC=2\pi\rg=2\pi\ellP.}
\label{eq:tres-lectores-planck}
\end{equation}
\end{corollary}

\begin{proof}
On substitue \(G_*=(54/\pi)^2c^5t_0^2/\hbar\). Alors
\[
\ellP=\frac{54}{\pi}ct_0=R_{108},
\qquad
\tP=\frac{54}{\pi}t_0=\frac1{\omega_{108}},
\]
et
\[
\mP=\frac{\pi}{54}\frac{\hbar}{c^2t_0}
=\frac{\hbar\omega_{108}}{c^2}=m_{108}.
\]
Le reste découle du \cref{thm:ct108-compton}.
\end{proof}

\begin{remark}[Trois lecteurs, une section]
La géométrie circulaire publie \(R_{108}\) et \(C_{108}\) ; la quantification de
Compton publie \(\lambdabarC\) et \(\lambdaC\) ; la réalisation gravitationnelle
publie \(\rg\) et \(2\pi\rg\). En \(\thetaStar\), ces trois paires sont les
coordonnées d’une unique section. C’est la raison précise pour laquelle
l’unité de Planck cesse d’apparaître comme une combinaison dimensionnelle accidentelle
dans cette formalisation.
\end{remark}

\section{Équivalence des normalisations sur l’identité autoduale}

Le transducteur \eqref{eq:transductor-g-planck-anexo} admet deux manières de
fixer l’échelle élémentaire, qui ne doivent pas être confondues.

\begin{enumerate}[label=\textup{(\roman*)}]
\item Si l’on fixe \(\theta=1\), alors
\(\sqrt{\hbar G/c^3}=\ell_0\). La transition élémentaire \(ct_0\) est identifiée
à la longueur de Planck de cette carte.
\item Si l’on exige que l’\emph{orbite complète} CT108 ait le rayon de Planck,
la condition sélectionne \(\thetaStar=(54/\pi)^2\) et
\(R_{108}=\ellP\).
\end{enumerate}

Ce ne sont pas des valeurs rivales de \(G\). Ce sont deux jauges de la même équation
homogène : l’une attribue l’unité de Planck au pas élémentaire ; l’autre, au rayon
du cycle complet. L’annexe utilise la seconde parce que sa thèse concerne
l’identité de la clôture CT108.

\section{Conventions gravitationnelles : rayon réduit et rayon de Schwarzschild}

La notation gravitationnelle est fixée sans équivoque :
\[
\rg=\frac{Gm}{c^2},
\qquad
\rs=\frac{2Gm}{c^2}=2\rg.
\]
Le théorème précédent emploie \(\rg\), non \(\rs\). Dans la convention habituelle de
la masse de Planck, l’égalité \(\hbar/(mc)=Gm/c^2\) se produit en \(m=\mP\),
tandis que
\begin{equation}
\frac{\hbar}{mc}=\frac{2Gm}{c^2}
\quad\Longleftrightarrow\quad
m=\frac{\mP}{\sqrt2}.
\label{eq:cruce-schwarzschild-factor-dos}
\end{equation}
C’est pourquoi l’expression « échelle gravitationnelle » désigne dans la thèse directrice le rayon
réduit. Le rayon d’horizon conventionnel demeure un lecteur distinct.

\section{Dédoublement bidirectionnel de la droite d’action}

La droite interne d’action possède deux sections orientées :
\begin{equation}
\hbarpre=H_5\,10^{-34}\mathcal U_S,
\qquad
\hbarret=\eta_{\mathrm{ret}}\,10^{-34}\mathcal U_S,
\qquad
\Ract=\frac{H_5}{\eta_{\mathrm{ret}}}.
\label{eq:hbar-pre-ret-anexo}
\end{equation}
La décade \(10^{-34}\) fixe l’échelle absolue ; le quotient \(\Ract\)
transporte la différence orientée. Cette séparation permet de demander si
l’autodualité de Planck subsiste dans les deux coordonnées.

\begin{theorem}[Covariance réciproque de l’action et de la gravité]
\label{thm:accion-gravedad-reciproca}
Pour \(a\in\{\mathrm{pre},\mathrm{ret}\}\), soit
\[
G_a(\theta)=\theta\frac{c^5t_0^2}{\hbar_a},
\qquad
m_{108,a}=\frac{\hbar_a\omega_{108}}{c^2}.
\]
Alors
\begin{equation}
\frac{G_a(\theta)m_{108,a}}{c^2}
=\theta\frac{\pi}{54}\ell_0
\label{eq:invariante-pre-ret-planck}
\end{equation}
est indépendant de \(a\). De plus,
\begin{equation}
\frac{m_{108,\mathrm{pre}}}{m_{108,\mathrm{ret}}}=\Ract,
\qquad
\frac{G_{\mathrm{pre}}}{G_{\mathrm{ret}}}=\Ract^{-1},
\label{eq:reciprocidad-pre-ret}
\end{equation}
et la longueur de Planck
\(\sqrt{\hbar_aG_a/c^3}=\sqrt\theta\,\ell_0\) est la même dans les deux
orientations.
\end{theorem}

\begin{proof}
Les facteurs \(\hbar_a\) s’annulent dans le produit \(G_am_{108,a}\).
Les deux rapports de \eqref{eq:reciprocidad-pre-ret} sont obtenus directement
à partir des définitions. L’invariance de la longueur de Planck découle de
\(\hbar_aG_a=\theta c^5t_0^2\).
\end{proof}

Une paramétrisation symétrique rend la géométrie visible :
\begin{equation}
\hbar_{\mathrm{geom}}:=\sqrt{\hbarpre\hbarret},
\qquad
\xi:=\frac12\log\Ract,
\qquad
\hbarpre=\hbar_{\mathrm{geom}}e^\xi,
\quad
\hbarret=\hbar_{\mathrm{geom}}e^{-\xi}.
\label{eq:rapidez-accion-pre-ret}
\end{equation}
La mémoire bidirectionnelle réside dans \(\xi\) ; la géométrie de Planck
réside dans le produit invariant. Cela permet une différence réelle entre les
deux coordonnées d’action sans exiger une différence de vitesse de
propagation.

\begin{remark}[Portée physique du théorème]
La réciprocité de \cref{thm:accion-gravedad-reciproca} est exacte au sein du
transducteur. Pour identifier \(G_{\mathrm{pre}}\) et \(G_{\mathrm{ret}}\) à deux
mesures gravitationnelles orientées, il faut spécifier le protocole,
l’observable et la symétrie qui les relie. Le théorème prouve la
compatibilité structurelle ; il n’anticipe pas à lui seul une violation de Lorentz
ni une variation observable de \(G\).
\end{remark}
 \label{mdg:sec:identidad-autodual-planck-ct108-rev3}
\index{Planck!identité autoduale}
\index{calendrier d’ordre 108!réalisation circulaire}
\index{Compton!lectures réduite et complète}

La droite d’action construite dans la section précédente admet une
réalisation circulaire compatible avec la périodicité temporelle d’ordre 108.
Cette réalisation coordonne trois objets qui possèdent déjà leurs types propres : la
section d’action, la géométrie d’un tour et la transduction
gravitationnelle. La coordination qui en résulte fournit une identité de
Planck interne et prépare la droite positive sur laquelle agira ensuite
l’opérateur général de masse.

\subsection{Cartes orientées d’action et réalisation circulaire déclarée}

Soient
\[
 \mathcal L_S^+,\quad \mathcal L_T^+,\quad \mathcal L_C^+,
 \quad \mathcal L_L^+,\quad \mathcal L_M^+,\quad \mathcal L_G^+
\]
les demi-droites positives d’action, de temps, de vitesse, de longueur et de masse.
Les coordonnées antérieure et postérieure au retour s’écrivent
\begin{equation}
 \hbar_{\rm pre}=H_5\,10^{-34}\mathcal U_S,
 \qquad
 \hbar_{\rm ret}=\eta_{\rm ret}\,10^{-34}\mathcal U_S,
 \qquad
 R_{\rm act}=\frac{\hbar_{\rm pre}}{\hbar_{\rm ret}}
 =\frac{H_5}{\eta_{\rm ret}}.
 \label{mdg:eq:planck-ct108-cartas-pre-ret-rev3}
\end{equation}
L’indice \(a\in\{\mathrm{pre},\mathrm{ret}\}\) désignera l’une quelconque
des deux cartes. Dans chacune, la distinction exacte est conservée
\begin{equation}
 h_a=2\pi\hbar_a .
 \label{mdg:eq:planck-ct108-h-hbar-rev3}
\end{equation}

Fixons une durée élémentaire \(t_0\in\mathcal L_T^+\), une vitesse
interne \(c_{\rm int}\in\mathcal L_C^+\) et
\(\ell_0=c_{\rm int}t_0\). La réalisation circulaire déclarée de la
périodicité temporelle d’ordre 108 est déterminée par
\begin{equation}
 T_{108}=108t_0,
 \qquad
 \omega_{108}=\frac{2\pi}{T_{108}},
 \qquad
 R_{108}=\frac{c_{\rm int}}{\omega_{108}}
 =\frac{54}{\pi}\ell_0,
 \label{mdg:eq:planck-ct108-realizacion-circular-rev3}
\end{equation}
et par la longueur du tour complet
\begin{equation}
 C_{108}=2\pi R_{108}=c_{\rm int}T_{108}=108\ell_0.
 \label{mdg:eq:planck-ct108-perimetro-rev3}
\end{equation}
Ainsi, \(R_{108}\) est le rayon de la réalisation et \(C_{108}\) son
périmètre. La fréquence angulaire, la fréquence cyclique et les deux
constantes d’action sont reliées par
\(\omega_{108}=2\pi/T_{108}\) et \(h_a=2\pi\hbar_a\).

\subsection{Correspondance entre la périodicité d’ordre 108 et les lectures de Compton}

Dans la carte \(a\), le quantum circulaire et sa section de masse sont
\begin{equation}
 E_{108,a}=\hbar_a\omega_{108},
 \qquad
 m_{108,a}=\frac{E_{108,a}}{c_{\rm int}^{2}}
 =\frac{\hbar_a\omega_{108}}{c_{\rm int}^{2}}.
 \label{mdg:eq:planck-ct108-cuanto-circular-rev3}
\end{equation}

\begin{theorem}[Identité entre la périodicité circulaire et les lectures de Compton]
\label{mdg:thm:planck-ct108-compton-rev3}
Pour chaque carte \(a\in\{\mathrm{pre},\mathrm{ret}\}\), les lectures de
Compton réduite et complète satisfont
\begin{equation}
 \boxed{
 \bar\lambda_{\rm C}(m_{108,a})
 =\frac{\hbar_a}{m_{108,a}c_{\rm int}}=R_{108},
 \qquad
 \lambda_{\rm C}(m_{108,a})
 =\frac{h_a}{m_{108,a}c_{\rm int}}=C_{108}.}
 \label{mdg:eq:planck-ct108-identidad-compton-rev3}
\end{equation}
\end{theorem}

\begin{proof}
La substitution de
\eqref{mdg:eq:planck-ct108-cuanto-circular-rev3} produit
\[
 \frac{\hbar_a}{m_{108,a}c_{\rm int}}
 =\frac{c_{\rm int}}{\omega_{108}}=R_{108}.
\]
La deuxième égalité résulte de l’application simultanée de
\eqref{mdg:eq:planck-ct108-h-hbar-rev3} et de
\eqref{mdg:eq:planck-ct108-perimetro-rev3}.
\end{proof}

Le théorème est homogène en
\((\hbar_a,c_{\rm int},t_0)\) : la même périodicité détermine le rayon,
le périmètre et les deux cartes de Compton. La réalisation circulaire est la
primitive géométrique déclarée ; les égalités précédentes sont ses
conséquences exactes.

\subsection{Sélecteur circulaire et unité interne de Planck}

La droite gravitationnelle est coordonnée à la carte \(a\) par le transducteur
homogène
\begin{equation}
 G_a(\theta)=\theta\,
 \frac{c_{\rm int}^{5}t_0^2}{\hbar_a}
 \in\mathcal L_G^+,
 \qquad \theta\in\mathbb R_{>0}.
 \label{mdg:eq:planck-ct108-transductor-rev3}
\end{equation}
Son rayon gravitationnel réduit, évalué sur le quantum circulaire, est
\begin{equation}
 r_{{\rm g},a}(\theta)
 :=\frac{G_a(\theta)m_{108,a}}{c_{\rm int}^{2}}
 =\theta\frac{\pi}{54}\ell_0.
 \label{mdg:eq:planck-ct108-radio-gravitatorio-rev3}
\end{equation}

\begin{proposition}[Sélecteur circulaire de Planck]
\label{mdg:prop:selector-circular-planck-ct108-rev3}
Au sein de la réalisation circulaire
\eqref{mdg:eq:planck-ct108-realizacion-circular-rev3}, la condition
\(r_{{\rm g},a}(\theta)=R_{108}\) sélectionne une unique coordonnée
positive,
\begin{equation}
 \boxed{\theta_{108}^{\circ}=\left(\frac{54}{\pi}\right)^2.}
 \label{mdg:eq:theta-circular-planck-ct108-rev3}
\end{equation}
La coordonnée est indépendante de la carte orientée \(a\).
\end{proposition}

\begin{proof}
Les équations
\eqref{mdg:eq:planck-ct108-realizacion-circular-rev3} et
\eqref{mdg:eq:planck-ct108-radio-gravitatorio-rev3} réduisent la condition à
\[
 \theta\frac{\pi}{54}\ell_0=\frac{54}{\pi}\ell_0.
\]
La positivité de \(\ell_0\) et la stricte monotonie de l’application
\(\theta\mapsto r_{{\rm g},a}(\theta)\) fournissent la solution et son
unicité.
\end{proof}

Soit \(G_{108,a}^{\circ}:=G_a(\theta_{108}^{\circ})\). Les unités internes
associées à cette carte satisfont
\begin{align}
 \ell_{{\rm P},a}^{\circ}
 &:=\sqrt{\frac{\hbar_a G_{108,a}^{\circ}}{c_{\rm int}^{3}}}
 =R_{108},
 &
 t_{{\rm P},a}^{\circ}
 &:=\sqrt{\frac{\hbar_a G_{108,a}^{\circ}}{c_{\rm int}^{5}}}
 =\omega_{108}^{-1},
 \label{mdg:eq:planck-ct108-longitud-tiempo-rev3}\\
 m_{{\rm P},a}^{\circ}
 &:=\sqrt{\frac{\hbar_a c_{\rm int}}{G_{108,a}^{\circ}}}
 =m_{108,a},
 &
 E_{{\rm P},a}^{\circ}
 &:=m_{{\rm P},a}^{\circ}c_{\rm int}^{2}=E_{108,a}.
 \label{mdg:eq:planck-ct108-masa-energia-rev3}
\end{align}
L’égalité des rayons est donc publiée comme
\begin{equation}
 R_{108}=\bar\lambda_{\rm C}=r_{\rm g}=\ell_{\rm P}^{\circ},
 \qquad
 C_{108}=\lambda_{\rm C}=2\pi r_{\rm g}=2\pi\ell_{\rm P}^{\circ}.
 \label{mdg:eq:planck-ct108-tres-lectores-rev3}
\end{equation}

La convention de ce théorème est le rayon gravitationnel réduit
\(r_{\rm g}=Gm/c_{\rm int}^{2}\). Le rayon de Schwarzschild satisfait
\(r_{\rm s}=2r_{\rm g}\) ; si cette seconde convention est adoptée, son croisement avec
\(\bar\lambda_{\rm C}\) se produit en
\(m=m_{{\rm P},a}^{\circ}/\sqrt2\). La déclaration du rayon fait donc
partie de la carte et évite de transférer par inadvertance un facteur deux
entre les deux identités.

La coordonnée \(\theta_{108}^{\circ}\) appartient à la spécialisation
circulaire déclarée. Le sélecteur gravitationnel général, formulé sur
l’holonomie enrichie et sa réponse spectrale, conserve sa résidence dans
la \cref{mdg:sec:selector-gravitatorio-rev11}. Ainsi, l’identité
circulaire fixe une normalisation interne de Planck et le critère général
fixe les conditions que doit satisfaire la réalisation physique de la
section gravitationnelle.

\subsection{Covariance pré-retour/post-retour}

Le produit des sections réciproques d’action et de gravité conserve la
géométrie circulaire :
\begin{equation}
 \frac{m_{108,\rm pre}}{m_{108,\rm ret}}=R_{\rm act},
 \qquad
 \frac{G_{108,\rm pre}^{\circ}}{G_{108,\rm ret}^{\circ}}
 =R_{\rm act}^{-1},
 \qquad
 \hbar_a G_{108,a}^{\circ}
 =\theta_{108}^{\circ}c_{\rm int}^{5}t_0^2.
 \label{mdg:eq:planck-ct108-covariancia-pre-ret-rev3}
\end{equation}
En particulier, \(\ell_{{\rm P},\rm pre}^{\circ}\) et
\(\ell_{{\rm P},\rm ret}^{\circ}\) coïncident avec \(R_{108}\). La mémoire
bidirectionnelle réside dans le quotient des masses et dans son réciproque
gravitationnel ; l’échelle géométrique réside dans le produit invariant.

La réciprocité admet une paramétrisation symétrique qui sépare l’échelle
commune de l’orientation du retour. Définissons
\begin{equation}
 \begin{aligned}
 \hbar_{\rm geom}
 &:=\sqrt{\hbar_{\rm pre}\hbar_{\rm ret}},
 &
 G_{108,\rm geom}^{\circ}
 &:=\sqrt{G_{108,\rm pre}^{\circ}G_{108,\rm ret}^{\circ}},
 &
 \xi_{\rm act}&:=\frac12\log R_{\rm act}.
 \end{aligned}
 \label{mdg:eq:planck-ct108-parametrizacion-simetrica-rev3}
\end{equation}
Alors
\[
 \hbar_{\rm pre}=\hbar_{\rm geom}e^{\xi_{\rm act}},
 \qquad
 \hbar_{\rm ret}=\hbar_{\rm geom}e^{-\xi_{\rm act}},
\]
tandis que les sections gravitationnelles conjuguées satisfont
\[
 G_{108,\rm pre}^{\circ}
 =G_{108,\rm geom}^{\circ}e^{-\xi_{\rm act}},
 \qquad
 G_{108,\rm ret}^{\circ}
 =G_{108,\rm geom}^{\circ}e^{\xi_{\rm act}}.
\]
La coordonnée \(\xi_{\rm act}\) conserve l’orientation
pré-retour/post-retour ; le produit \(\hbar_aG_{108,a}^{\circ}\)
conserve l’échelle géométrique. Cette paramétrisation des cartes d’action
n’introduit pas à elle seule des vitesses de propagation distinctes.

\subsection{Droite autoduale déployée des masses}

Une carte \(a\) étant fixée, soit l’algèbre déployée
\[
 \mathbb D_{\rm sp}:=\mathbb R[j]/(j^2-1),
 \qquad
 P_\pm:=\frac{I\pm j}{2},
\]
et considérons sa décomposition spectrale
\[
 \mathbb D_{\rm sp}
 =\mathbb R P_+\oplus\mathbb R P_-,
 \qquad
 P_\pm^2=P_\pm,\quad
 P_+P_-=0,\quad
 P_++P_-=I.
\]
La conjugaison déployée échange les projecteurs,
\[
 \overline{uP_++vP_-}=vP_++uP_-,
\]
et sa norme est
\[
 N_{\rm sp}(uP_++vP_-):=uv.
\]

Pour \(m\in\mathcal L_M^+\), définissons la section
\begin{equation}
 Z_a(m)
 :=\frac{m_{{\rm P},a}^{\circ}}{m}P_+
   +\frac{m}{m_{{\rm P},a}^{\circ}}P_- .
 \label{mdg:eq:planck-ct108-seccion-partida-masa-rev3}
\end{equation}
L’unité \(m_{{\rm P},a}^{\circ}\) détermine également l’involution
\begin{equation}
 \iota_{{\rm P},a}(m)
 =\frac{(m_{{\rm P},a}^{\circ})^2}{m}.
 \label{mdg:eq:planck-ct108-involucion-masas-rev3}
\end{equation}

\begin{theorem}[Droite autoduale déployée des masses]
\label{mdg:thm:planck-ct108-recta-partida-masa-rev3}
Pour tout \(m\in\mathcal L_M^+\),
\begin{equation}
 \boxed{
 N_{\rm sp}(Z_a(m))=1,
 \qquad
 \overline{Z_a(m)}
 =Z_a\!\left(\iota_{{\rm P},a}(m)\right).}
 \label{mdg:eq:planck-ct108-norma-conjugacion-rev3}
\end{equation}
La conjugaison possède un unique point fixe positif,
\[
 m=m_{{\rm P},a}^{\circ},
 \qquad
 Z_a(m_{{\rm P},a}^{\circ})=I.
\]
\end{theorem}

\begin{proof}
La norme est le produit des deux coordonnées réciproques :
\[
 \frac{m_{{\rm P},a}^{\circ}}m
 \frac m{m_{{\rm P},a}^{\circ}}=1.
\]
La conjugaison échange les deux coordonnées, ce qui équivaut à remplacer
\(m\) par \((m_{{\rm P},a}^{\circ})^2/m\). Le point fixe positif satisfait
\(m^2=(m_{{\rm P},a}^{\circ})^2\).
\end{proof}

Avec la rapidité
\begin{equation}
 x_a(m)=\log\frac{m}{m_{{\rm P},a}^{\circ}},
 \label{mdg:eq:planck-ct108-rapidez-masa-rev3}
\end{equation}
la section et sa conjuguée s’écrivent
\[
 Z_a(m)=e^{-x_a}P_++e^{x_a}P_-,
 \qquad
 \overline{Z_a(m)}
 =e^{x_a}P_++e^{-x_a}P_-.
\]
Leurs lecteurs pair et impair sont
\begin{equation}
 Q_a(m):=\frac{e^{x_a}+e^{-x_a}}2=\cosh x_a,
 \qquad
 A_a(m):=\frac{e^{x_a}-e^{-x_a}}2=\sinh x_a.
 \label{mdg:eq:planck-ct108-lectores-par-impar-rev3}
\end{equation}
Le premier conserve la distance au point autodual et le second conserve son
orientation.

Les lecteurs physiques réciproques associés à la même coordonnée sont
\[
 r_{{\rm g},a}(m)=\frac{G_{108,a}^{\circ}m}{c_{\rm int}^{2}},
 \qquad
 \bar\lambda_{{\rm C},a}(m)=\frac{\hbar_a}{mc_{\rm int}},
\]
et satisfont
\begin{equation}
 \frac{r_{{\rm g},a}(m)}{\bar\lambda_{{\rm C},a}(m)}
 =e^{2x_a(m)},
 \qquad
 \frac{\bar\lambda_{{\rm C},a}(m)}{r_{{\rm g},a}(m)}
 =e^{-2x_a(m)}.
 \label{mdg:eq:planck-ct108-lectores-exponenciales-rev3}
\end{equation}
La conjugaison échange les deux lecteurs ; l’identité de Planck est leur
équilibre unitaire.

\begin{remark}[Identité de Planck et clôture électronique]
L’identité \(m=m_{{\rm P},a}^{\circ}\) fixe l’origine autoduale de la droite
positive des masses. La section électronique fixe une position différente sur
cette même droite par une généalogie propre. La
\cref{mdg:sec:ley-general-masa-accion-autonoma} construit l’opérateur universel
qui transporte l’échelle absolue, le caractère, les tours et la lecture bidirectionnelle ;
la \cref{mdg:sec:cierre-electron-two-way-rev11} réalise ensuite sa réduction
centrale de rang un. Ainsi, l’unité de Planck fournit l’identité de
la coordonnée et l’électron fournit la première transition matérielle
complète sur celle-ci.
\end{remark}

\begin{proposition}[Réalisation généalogique du centre électronique]
\label{mdg:prop:planck-ct108-realizacion-centro-electron-rev3}
Soit
\[
 F_e^+=\mathbb R_{>0}(P_++P_-)
\]
la demi-droite positive de la fibre centrale de rang un, soit
\(\mathcal G_e\) l’espace des généalogies
électroniques enrichies et soit \(\mathcal L_M^+\) la droite positive des
masses. Un morphisme de réalisation
\[
 \mathfrak R_e:
 F_e^+\times\mathcal G_e\longrightarrow\mathcal L_M^+
\]
peut envoyer l’identité \(I=P_++P_-\) sur une masse différente de
\(m_{{\rm P},a}^{\circ}\) sans modifier le
\cref{mdg:thm:planck-ct108-recta-partida-masa-rev3} : l’identité de Planck fixe
l’origine autoduale de la carte, tandis que la généalogie électronique fixe un
déplacement sur cette carte.
\end{proposition}

\begin{proof}
Si \(\chi_e:\mathcal G_e\to\mathbb R_{>0}\) est le caractère positif construit
par le registre électronique, la règle
\[
 \mathfrak R_e(\lambda I,g)
 =\lambda\,m_{{\rm P},a}^{\circ}\chi_e(g)
\]
est positive et multiplicative dans la coordonnée généalogique. Pour la généalogie
neutre, \(\chi_e=1\), on retrouve l’origine de Planck ; pour une généalogie avec
\(\chi_e(g)\ne1\), la même identité interne occupe une position différente sur
la droite. La clôture électronique concrète conserve son développement antécédent dans la
\cref{mdg:sec:cierre-electron-two-way-rev11}.
\end{proof}
 \section{Autodualité de l'unité interne de Planck sous inversion d'orientation}

Le sélecteur CT108 construit

\begin{equation}
\theta_{108}=\left(\frac{54}{\pi}\right)^2,
\qquad
\ell_P=\frac{54}{\pi}\ell_0,
\qquad
t_P=\frac{54}{\pi}t_0,
\label{eq:planck-length-time}
\end{equation}

et

\begin{equation}
E_P=\frac{\hbar_{\rm int}}{t_P}
=\frac{\pi\hbar_{\rm int}}{54t_0}.
\label{eq:planck-energy}
\end{equation}

En comparant \cref{eq:boltzmann-clock,eq:planck-energy}, on obtient

\begin{equation}
\boxed{E_P=k_B\Theta_{\rm clk}\log3.}
\label{eq:planck-trit}
\end{equation}

La relation ne dépend pas d'une coïncidence entre unités : les deux membres
proviennent de la même structure chronologique \(108\).

Lorsque la constante gravitationnelle du \cref{ch:gravity} est introduite, la famille de Planck satisfait

\begin{align}
\ell_P^2&=\frac{G\hbar}{c^3}=\theta_{108}\ell_0^2,
\label{eq:planck-area}\\
F_P&=\frac{c^4}{G},
\qquad
P_P=\frac{c^5}{G},
\label{eq:planck-force-power}\\
\rho_P&=\frac{\hbar}{\theta_{108}^2c^5t_0^4}.
\label{eq:planck-density}
\end{align}

Bien que \(G\) et \(\hbar\) se transforment réciproquement entre les sections antérieure et postérieure au retour, leur produit dans \(\ell_P^2\) conserve la géométrie. Cette invariance détermine l'aire comme interface entre les réalisations géométrique, informationnelle et d'intrication. \endgroup
% Distribución causal exacta del delta Ley 9; contenido conservado sin síntesis.
\section{Droite d'action et constante de Planck}
\label{sec:l9s102:recta-planck}

Soit \(\mathcal L_S\) la droite réelle d'action de base positive
\(\mathcal U_S\). La transduction
\begin{equation}
 \tau_S:\mathbb R_{>0}\longrightarrow\mathcal L_S,
 \qquad u\longmapsto u\,10^{\varepsilon_H}\mathcal U_S
 \label{eq:l9s102:transductor-accion}
\end{equation}
publie les sections antérieure et postérieure au retour :
\begin{equation}
 \hbar_{\mathrm{pre}}=H_5\,10^{-34}\mathcal U_S,
 \qquad
 \boxed{
 \hbar_{\mathrm{ret}}=\eta_{\mathrm{ret}}\,10^{-34}\mathcal U_S.}
 \label{eq:l9s102:hbar-pre-ret}
\end{equation}
Les actions de cycle complet sont
\begin{equation}
 h_{\mathrm{pre}}=2\pi\hbar_{\mathrm{pre}},
 \qquad
 h_{\mathrm{ret}}=2\pi\hbar_{\mathrm{ret}}.
 \label{eq:l9s102:h-pre-ret}
\end{equation}
Pour une phase \(\theta\), les cartes angulaire et d'action commutent :
\begin{equation}
 S(\theta)
 =\hbar_{\mathrm{ret}}\theta_{\mathrm{rad}}
 =\hbar_{\mathrm{ret}}\frac{\pi}{180}\theta_{\deg}
 =h_{\mathrm{ret}}\frac{\theta_{\deg}}{360}.
 \label{eq:l9s102:accion-fase}
\end{equation}

 
\chapter{Opérateur de phase, rayon spectral et factorisation de la double projection}
\begingroup
\renewcommand{\chapter}[2][]{}%
\chapter{Opérateur angulaire prédimensionnel : décomposition en \(27\) secteurs, rayon spectral et composition effective de la double projection}
\label{chap:p3-final:37:operador_angular}
\label{p3-20260826:chap:operador-angular-dictamen}
\index{opérateur angulaire}
\index{phase angulaire!sortie spectrale}
\index{cierre!pre-dimensional}

Les deux sections précédentes ont réuni trois éléments qui apparaissaient auparavant
séparés : une identification isométrique entre projecteurs, une mesure de
quotient et un cocycle positif d’action. La présente section étudie la conséquence
spectrale de cette loi et situe le résultat dans l’architecture complète
de HMT\@.

La séparation des statuts est précise. L’identification triple est un théorème
exact relatif à la nouvelle règle variationnelle interne et à un ancrage orienté.
La mesure et le caractère positif sont des théorèmes exacts de \HTr--\HAdd ; le défaut
est en outre relatif à \HDef, la capacité à \HRes et le coût local à \HLoc. La valeur
angulaire est une sortie certifiée par intervalles de l’opérateur \Gtr. Sa
proximité avec la coordonnée angulaire conjuguée canonique est extraordinaire, mais l’intervalle de
la différence exclut l’égalité.

\section{Décomposition de l’opérateur angulaire en \(27\) secteurs de \(72\) états}

Le lecteur de tous les chemins \Gtr possède \(1944\) états microscopiques.
L’action de translation grossière de \((\mathbb Z/3\mathbb Z)^3\) permet la
décomposition
\[
 1944=27\cdot72.
\]
Chaque bloc de \(72\) états est paramétré par
\[
 (r,s,m,d)
 \in(\mathbb Z/3\mathbb Z)^2\times\{0,1\}\times D_4,
\]
où \(m\) est le feuillet et \(d\) la direction précédente. De chaque état
partent quatre continuations cardinales. Le TRIT de la cellule d’arrivée
décide du feuillet suivant.

Avec l’action du \cref{p3-20260826:chap:medida-accion-positiva}, le poids d’une arête
est
\[
 \exp[-g(e)-\lambda_*s(e)].
\]
La transformée en caractères réduit l’opérateur à des matrices complexes
\(L_{k_a,k_b}\) de taille \(72\times72\). Le programme énumère les
arêtes de chaque bloc ; il ne reçoit pas une matrice spectrale précalculée.
L’itération des puissances est conservée à titre exploratoire, mais le résultat qui
suit ne repose plus sur elle : les rayons spectraux sont isolés au moyen d’une
certification rationnelle complète.

Soient \(\rho_{01}\) et \(\rho_{12}\) les rayons spectraux des secteurs
orientés \((0,1)\) et \((1,2)\). Nous définissons la \emph{phase de bande}
\[
 \boxed{y_{\rm band}=\log\rho_{12}-\log\rho_{01}.}
\]
Pour
\[
 \lambda_*
 =1-\frac5{468}-\frac3{11}\Delta_4,
\]
la certification donne
\[
\begin{aligned}
1.460655120862011111
&<\rho_{01}<1.460655120862404024,\\
1.515812008270944328
&<\rho_{12}<1.515812008271195416,
\end{aligned}
\]
et, par propagation par intervalles,
\[
 0.037066226434427529
 <y_{\rm band}<
 0.037066226434862172,
\]
\[
 \boxed{
 2.123738337168943^\circ
 <C_{\rm band}:=\frac{180}{\pi}y_{\rm band}
 <2.123738337193847^\circ.}
\]
La quantité primaire est \(y_{\rm band}\), un logarithme adimensionnel de quotient
spectral. Les degrés apparaissent ensuite lorsque l’on applique la carte
\(180/\pi\). L’opérateur produit d’abord une phase interne et seulement ensuite
une représentation angulaire archimédienne.

\subsection{Calcul exact du rayon spectral de l’opérateur angulaire}

Les blocs sont non normaux. Par conséquent, un petit résidu
\(\|Lv-\mu v\|\) ne suffit pas à démontrer que \(\mu\) est la valeur propre de
plus grand module. La preuve utilisée ici comporte quatre étapes.

On observe d’abord douze lignes structurellement nulles dans chaque bloc. En
développant le polynôme caractéristique suivant ces lignes, on obtient un facteur
\(z^{12}\) et un mineur principal d’ordre soixante. Dans ce dernier apparaissent huit
autres lignes structurellement nulles. Par conséquent,
\[
 \det(zI_{72}-L)=z^{20}\det(zI_{52}-A),
\]
et toute valeur propre non nulle est contenue, avec sa multiplicité, dans un bloc
réduit \(A\) d’ordre cinquante-deux.

Deuxièmement, le certificat fournit des matrices gaussiennes rationnelles
\(S,R\), de dénominateur \(10^{16}\). Le vérificateur calcule exactement
\(E=I-RS\) et obtient
\[
 \|E\|_\infty\le8.5664\times10^{-14}
 \quad\hbox{et}\quad
 \|E\|_\infty\le3.7628\times10^{-14}
\]
pour les secteurs \((0,1)\) et \((1,2)\), respectivement. La série de
Neumann démontre alors que \(S\) est inversible.

Troisièmement, si \(A_0\) est le point médian rationnel du bloc d’intervalles, l’identité
\[
 S^{-1}AS=(I-E)^{-1}RAS
\]
borne la distance à \(C_0=RA_0S\) par
\[
 1.25125\times10^{-13}
 \quad\hbox{et}\quad
 5.7037\times10^{-14}.
\]
Les disques de Gershgorin de \(C_0\), agrandis de ces bornes, contiennent
les disques de la similitude exacte. Dans chaque bloc, un disque est séparé
des cinquante et un autres ; le second théorème de Gershgorin lui attribue
exactement une valeur propre. Son module inférieur dépasse le module supérieur de
tous les autres disques. C’est la preuve que la valeur propre isolée est le
rayon spectral.

Quatrièmement, \(\pi,e,\log\pi,e^{-1},e^{-\lambda_*}\) et \(\sqrt3\) ne sont pas
acceptés comme nombres décimaux de bibliothèque. Le vérificateur les encadre au moyen de
la formule de Machin, des séries factorielle et logarithmique avec reste, de la
série alternée de l’exponentielle et d’un encadrement rationnel quadratique pour
\(\sqrt3\). La vérification emploie uniquement des entiers et des fractions ;
l’argument est clos par la réduction rationnelle, la série de Neumann et
la séparation de Gershgorin exposées dans cette partie. Les développements
décimaux reproductibles à une profondeur de dix mille sont publiés, comme témoins
finis indépendants, dans
\hyperref[app:desarrollos-10000-rev9]{l’annexe des développements}.

\section{Opérateur effectif de bande \texorpdfstring{\(T_{\rm band}\)}{Tband} et isolement de la phase angulaire}

Soit
\[
 x=\frac{50\pi}{9}\alpha
 =0.127362829012869966\ldots
\]
et soit
\[
 R=\begin{pmatrix}1&0\\0&-1\end{pmatrix}.
\]
L’opérateur bicouche effectif associé exclusivement à la phase de bande est
\[
 \boxed{T_{\rm band}=\exp(-xI+y_{\rm band}R).}
\]
Dans la base des secteurs orientés,
\[
 T_{\rm band}
 =\begin{pmatrix}
 e^{-x+y_{\rm band}}&0\\
 0&e^{-x-y_{\rm band}}
 \end{pmatrix}.
\]
On vérifie
\[
 \det T_{\rm band}=e^{-2x},
 \qquad
 \frac{(T_{\rm band})_{11}}{(T_{\rm band})_{22}}
 =e^{2y_{\rm band}}.
\]
Par conséquent,
\[
 -\frac9{100\pi}\log\det T_{\rm band}=\alpha,
 \qquad
 \frac{90}{\pi}
 \log\frac{(T_{\rm band})_{11}}{(T_{\rm band})_{22}}
 =C_{\rm band}.
\]

Les deux récupérations ont des statuts distincts. La seconde restitue la
sortie produite par l’écart. La première relit le
\(\alpha\) entré dans \(x\) ; elle démontre la réversibilité de la
carte, et non une seconde dérivation indépendante de \(\alpha\).

La forme exponentielle sépare deux fonctions. Le facteur commun \(e^{-x}\)
décrit la contraction partagée par les deux feuillets. Les facteurs
\(e^{\pm y_{\rm band}}\) distinguent les orientations en conservant le déterminant. Comme
\(I\) et \(R\) commutent, la paire
\[
 \left(\det T_{\rm band},
 \frac{(T_{\rm band})_{11}}{(T_{\rm band})_{22}}\right)
\]
reconstruit univoquement \((x,y_{\rm band})\).

\section{Comparaison ultérieure et échelle de précision}

La section suivante construit une autre phase, la \emph{phase régionale
régularisée} \(y_{\rm reg}\) ---également notée \(y_*\) dans cette section---. Pour
formuler ici une comparaison avec ce résultat ultérieur, nous anticipons la
notation
\[
 C_{\rm reg}:=\frac{180}{\pi}y_{\rm reg}
 =2.123738338968646^\circ.
\]
Cette ligne ne constitue ni une seconde définition ni une entrée de l’opérateur
spectral : elle cite la conclusion du
\cref{p3-20260826:thm:contraangulo-regional}, dont la construction est développée dans la
partie suivante.
Les deux paires d’objets ne sont pas identifiées :
\[
 (y_{\rm band},C_{\rm band})
 \ne
 (y_{\rm reg},C_{\rm reg}).
\]
La différence certifiée par intervalles de la loi linéaire est
\[
 -1.799703\times10^{-9}\ {}^\circ
 <C_{\rm band}-C_{\rm reg}
 <-1.774799\times10^{-9}\ {}^\circ,
\]
avec une différence relative encadrée par
\[
 -8.475\times10^{-10}
 <\frac{C_{\rm band}-C_{\rm reg}}{C_{\rm reg}}
 <-8.356\times10^{-10}.
\]
La loi qui produit \(C_{\rm band}\) ne reçoit ni \(C_{\rm reg}\), ni CKM, ni Barbero,
ni masses ni données SI\@. Cette indépendance à l’égard de la cible fait de la proximité un
candidat de reconnaissance fort et falsifiable. La certification spectrale
élimine l’ancienne lacune numérique ; la différence par intervalles, qui exclut
zéro, fixe désormais la limite exacte : \(C_{\rm band}\) est une sortie spectrale
prédimensionnelle distincte de la réalisation régionale
\(C_{\rm reg}\).

Le terme supplémentaire
\[
 -\frac{\Delta_4^2}{27}
\]
en \(\lambda\) produce
\[
 C_{\rm num}\approx2.12373833894104980{}^\circ,
\]
à \(2.76\times10^{-11}\) degrés de \(C_{\rm reg}\). L’entier \(27\) est interne,
car il compte les caractères de \((\mathbb Z/3\mathbb Z)^3\), mais ce
fait n’impose pas le coefficient du second moment. L’additivité exacte de
\HAdd exclut les termes quadratiques. La variante est conservée comme ablation et non
comme loi principale.

\section{Dépendance fonctionnelle du rayon spectral à l’égard des composantes de l’opérateur}

Toutes les lignes suivantes utilisent le même opérateur ; seule une entrée
de la loi change. Ce sont des diagnostics numériques d’ablation ; l’intervalle certifié
précédent correspond exclusivement à la loi complète
\HTr--\HAdd--\HDef--\HRes--\HLoc (et, en son sein, le défaut positif est fixé dans la
strate \HTr--\HAdd--\HDef).
\[
\begin{array}{l|c|c}
\text{règle}&\lambda&C_{\rm num}\text{ en degrés}\\ \hline
\text{autodual}&1&2.082755794784418\\
\text{microfibre seule}&1-5/468&2.123621809570142\\
\text{trace }5/468+(3/11)\Delta_4
&0.989285913019141&2.123738337181414\\
4/468\text{ au lieu de }5/468
&0.991422665155894&2.115535228138553\\
3/12\text{ au lieu de }3/11
&0.989288440210566&2.123728626433667\\
5/243\text{ au lieu de }5/468
&0.979393542015659&2.161907060464779\\
5/729\text{ au lieu de }5/468
&0.993110963140488&2.109064222037112\\
7/729\text{, mesure des germes}
&0.990367478915522&2.119584299852700.
\end{array}
\]
Les ablations démontrent que la sortie distingue la microfibre complète,
la dimension réduite onze et la mesure des émissions. Elles n’attribuent pas à elles
seules une probabilité au hasard : une telle probabilité exigerait de définir au préalable un
espace de modèles alternatifs et une mesure sur celui-ci.

\section{Composition typée des morphismes archimédien, exceptionnel et dimensionnel}

La construction complète peut s’écrire comme une succession d’applications
aux domaines visibles :
\[
\begin{aligned}
\Omega_{\rm seed}
&\xrightarrow{F}\mathcal U_6
\xrightarrow{J_F}H_U,\\
\mathscr R_\pi^{++}
&\xrightarrow{\beta_\pi}P_GH_G,\\
L_{\rm or}\otimes P_GH_G
&\xrightarrow{W_{GK}}P_3V_{11}
\xrightarrow{U_W}Q_3E_{30},\\
(\Pi_\pi^{(468)},\Delta_4P_3)
&\longmapsto D_{\rm switch}
\xrightarrow{\tau}\delta_*
\xrightarrow{\rm \HRes}\lambda_*,\\
\lambda_*
&\xrightarrow{\rm \HLoc}a_*
\xrightarrow{\mathcal L_{1,*}}y_{\rm band}
\xrightarrow{180/\pi}C_{\rm band},\\
(\alpha,y_{\rm band})
&\xrightarrow{\exp(-xI+y_{\rm band}R)}T_{\rm band}.
\end{aligned}
\]
Chaque flèche déclare ce qu’elle conserve :
\begin{itemize}
  \item \(F\) conserve la classe d’émission et laisse dans la fibre sa
  multiplicité de présentation ;
  \item \(J_F\) compense cette multiplicité dans la norme ;
  \item \(\beta_\pi\) conserve la position exceptionnelle ;
  \item \(W_{GK}\) conserve la représentation orientée ;
  \item \(U_W\) conserve l’incidence et la métrique ;
  \item la trace oublie les coordonnées et conserve les rangs ;
  \item l’opérateur de transfert somme les chemins avec le poids de l’action ;
  \item la carte \(180/\pi\) représente une phase interne en degrés.
\end{itemize}

\section{Diagramme des dépendances de l’opérateur angulaire et de ses publications}

\begin{longtable}{>{\raggedright\arraybackslash}p{.31\textwidth}>{\raggedright\arraybackslash}p{.22\textwidth}>{\raggedright\arraybackslash}p{.37\textwidth}}
\toprule
Objet&Nature mathématique&Contenido demostrado\\
\midrule
\endfirsthead
\toprule
Objet&Nature mathématique&Contenido demostrado\\
\midrule
\endhead
\(U_W=I/6:V_{11}\to E_{30}\)
&théorème exact
&isométrie, surjectivité, incidence équivariante et unicité positive\\
\(P_3\leftrightarrow Q_3\)
&théorème exact
&projecteurs conjugués, rang trois et trace \(3/11\)\\
étoile \(R_{B_K}\)
&théorème exact typé
&spectre \(1^{[7]},(-1)^{[4]}\), indice \(3/11\) et obstruction de conjugaison\\
\(P_G\) dans \(t=7\)
&théorème fini exact
&filtration \(11\to3\to1\), projecteur positif et action \(S_3\)\\
\(\mathscr R_\pi^{++}\leftrightarrow P_G\)
&théorème exact
&bijection \(S_3\)-équivariante par position exceptionnelle\\
identification isométrique orientée
&nouveau théorème relatif
&sélecteur unique par \(B_K,u_K\), Reynolds, polaire et réduction \(240\to1\)\\
mesure \(1/468\)
&théorème exact de \HTr
&trace normalisée unique et application polaire induite par les fibres\\
\(D_{\rm switch}\) y \(\delta_*\)
&théorème exact relatif à \HTr--\HAdd--\HDef
&positivité, spectre et trace \(5/468+(3/11)\Delta_4\)\\
\(\lambda_*\)
&théorème interne relatif à \HTr--\HAdd--\HDef--\HRes
&capacité résiduelle obtenue à partir du défaut par la règle \HRes\\
\(\mathcal A_*\)
&théorème interne relatif à \HTr--\HAdd--\HDef--\HRes--\HLoc
&cocycle additif de chemins de coût local \(g+\lambda_*s\)\\
\(T_{\rm band}\) y \(C_{\rm band}\)
&théorème computationnel par intervalles
&réduction \(72\to60\to52\), similitude rationnelle, isolement de
Gershgorin et sortie dont la distance certifiée est comprise entre
\(1.774799\times10^{-9}\) et \(1.799703\times10^{-9}\) degrés de
\(C_{\rm reg}\)\\
\(C_{\rm band}=C_{\rm reg}\)
&affirmation réfutée pour la loi linéaire
&l’intervalle certifié de la différence exclut zéro\\
transition totale
\(\mathcal X_{\mathrm{TPK}}^{\mathrm{enr}}\to \Gtr\)
&réalisateur global requis
&\HTr--\HAdd--\HDef--\HRes--\HLoc définit le lecteur relatif ; la transition
totale exige en outre une mesure image dynamique sur \(\Gtr\)\\
\bottomrule
\end{longtable}

\section{Action prédimensionnelle au centre de la double projection de l’état généalogique}

La double projection part du même état holonomique intégral et produit deux lectures. Dans la branche archimédienne, les cylindres compatibles se raffinent et leurs diamètres décroissent ; l'évaluation fixe une valeur limite et oublie une partie de la généalogie. Dans la branche exceptionnelle, les incidences se prolongent vers Witt, Golay, Niemeier--Leech, \(V^\natural\), le Monstre et le moonshine.

Le nouveau résultat ajoute un objet commun plus fort qu’une coïncidence de
cardinaux. Le module à onze modes, son secteur positif de rang trois et la
trace de l’action passent de la dodécaphase à Witt par \(U_W\). Le secteur orienté
\(t=7\) s’incorpore par \(W_{GK}\). La microfibre de
\(\pi\) apporte le projecteur régional \(\Pi_\pi^{(468)}\), et la mesure uniforme du
catalogue fixe son poids.

L’action positive se situe ainsi au centre des deux lectures :
\[
 \text{région}
 \quad\oplus\quad
 \text{polarisation orientée}
 \quad\longmapsto\quad
 \text{défaut positif}.
\]
L’évaluation archimédienne et l’évaluation exceptionnelle restent des
applications différentes ; l’entrelaceur identifie un secteur commun, il ne
confond pas leurs codomaines.

\section{Compatibilité de l’opérateur angulaire avec les cylindres du continu discret}

Une histoire HMT détermine des cylindres emboîtés
\[
 C_1\supset C_2\supset\cdots
\]
avec des préfixes compatibles. Lorsque le lecteur archimédien satisfait
\[
 \operatorname{diam}(C_n)\longrightarrow0,
\]
la complétude de \(\mathbb R\) fixe au plus une valeur dans l’intersection.
La nouvelle action ajoute une structure avant cette limite : elle distingue le coût
de la route, la classe d’émission et la multiplicité des présentations que la
valeur finale ne montre plus.

C’est le sens mathématique précis de l’affirmation selon laquelle le continu se lit comme
projection de généalogies discrètes. Le système des cylindres et ses
évaluations se prolongent à toute échelle ; le théorème global de couverture
démontre
\(\operatorname{ev}_{\mathbb R}(\Omega_{\infty,K})=K\) pour chaque compact
\(K\subset\mathbb R\) et, en dirigeant les intervalles, que l’image est tout
\(\mathbb R\). Le théorème fini de cette partie contrôle l’action angulaire
à chaque niveau et s’intègre dans cette couverture globale ; il ne la rouvre ni ne la
conditionne. La descente algébrique de chaque opération conserve sa vérification
propre.

\section{Compatibilité avec le corridor exceptionnel, le Moonshine et les écrans de la théorie M}

L’entrelaceur \(U_W\) aboutit dans le module d’incidence de \(W_{12}\).
À partir du code ternaire, les deux continuations exceptionnelles restent
séparées :
\[
 C_W\longrightarrow W_{12}
 \xrightarrow[\text{relatif à }(D,\ell)]{}W_{24},
\]
et
\[
 C_W\xrightarrow{\ \phi\ } N(A_2^{12})
\xrightarrow{\ v_\alpha\ }\Lambda_{24}
\longrightarrow V_\Lambda
\longrightarrow V^\natural
\longrightarrow\mathbb M
\longrightarrow\text{moonshine}.
\]
L’action positive apporte à cette branche un secteur transporté et une droite
d’orientation. Les flèches de Leech au moonshine sont les constructions
classiques appliquées à l’objet fini obtenu. L’action du Monstre se
réalise sur \(V^\natural\) ; les mots HMT ne sont pas identifiés ici à un
module direct pour cette action.

Les écrans \(12\to11\to10\), \(26\to24\) et l’involution T-duale forment
une interface mathématique exacte avec les dimensions et les dualités utilisées en
théorie M. L’identification physique complète relève de la mécanique
dimensionnelle et est typée dans la seconde partie par l’application supplémentaire vers
\(g_{MN}\), \(C_{MNP}\), \(\psi_M\), l’action et la supersymétrie. Elle n’est pas réservée
à un autre ouvrage et ne se déduit pas de l’action du Monstre.

\section{Factorisation spectrale et isolement des secteurs dominants}

La ligne prédimensionnelle est organisée en quatre strates.
\begin{enumerate}
  \item \textbf{Noyau fini.} APP, TRIT, l’émetteur TPK, le catalogue
  \(104\,976\to468\to243\), les cinq régions de \(\pi\), la
  dodécaphase, \(K\), \(\alpha\), Witt et la branche exceptionnelle possèdent
  des constructions explicites et des certificats.
  \item \textbf{Pont opératoriel.} L’incidence construit \(U_W\) ;
  le secteur orienté \(t=7\) construit \(P_G\) ; \(B_K\) et \(u_K\) sélectionnent
  l’identification isométrique orientée ; les trois projecteurs de rang trois
  sont identifiés relativement à cette règle.
  \item \textbf{Extension d’action.} \HTr fixe la mesure \(1/468\), \HAdd
  fixe le caractère additif et \HDef déclare le couplage
  \(\Delta_4P_3\). \HRes transforme le défaut en capacité et \HLoc porte la
  capacité sur les arêtes. Les cinq clauses produisent la loi locale.
  \item \textbf{Conséquence spectrale.} \Gtr produit
  \(T_{\rm band}\) et une phase angulaire prédimensionnelle certifiée, avec
  une différence relative comprise entre \(8.356\times10^{-10}\) et
  \(8.475\times10^{-10}\) par rapport à \(C_{\rm reg}\).
\end{enumerate}

Cela permet une clôture forte et typée de cette route : l’émetteur fini de la
section coinductive infinie, l’incidence, l’identification des
projecteurs, la mesure et l’action sont clos
relativement aux primitives déclarées, à la règle variationnelle
\(\mathcal E_B\) et aux clauses d’action \HTr--\HAdd--\HDef--\HRes--\HLoc. \HAdd
sélectionne la forme linéaire ; \HDef sélectionne sa seconde coordonnée ; \HRes et \HLoc
précisent comment elle se lit comme coût. Le théorème global de
couverture de tout \(\mathbb R\), l’induction à partir d’une transition totale de
\(\mathcal X_{\mathrm{TPK}}^{\mathrm{enr}}\) et l’équivalence physique avec
la théorie M sont formulés comme des extensions des résultats précédents.

\section{Compatibilité entre les réalisations matricielle, angulaire et réticulaire}

Les représentations construites satisfont simultanément les conditions
de compatibilité suivantes :
\begin{enumerate}
  \item cinq régions de \(\pi\), dont trois positives, une négative et une
  transversale ;
  \item \(P_G,P_3,Q_3\) comme projecteurs positifs et \(R_{B_K}\) comme
  involution d’orientation ;
  \item le secteur homogène orienté \(t=7\), et non le quotient hétérogène de \(t=8\) ;
  \item la droite \(L_{\rm or}\) dans l’identification isométrique triple ;
  \item le caractère relatif de la règle variationnelle qui sélectionne
  \(H_*,c_*,r_*\) ;
  \item \(5/468\) sur les émissions et \(7/729\) sur les germes, sans
  les intervertir ;
  \item la linéarité du caractère comme conséquence de \HAdd et le coefficient
  \(\Delta_4\) comme contenu séparé de \HDef ;
  \item \((y_{\rm band},C_{\rm band})\) comme sortie spectrale certifiée,
  séparée de \((y_{\rm reg},C_{\rm reg})\), la réalisation régionale
  canonique de la section suivante ;
  \item \(\alpha\) comme entrée de \(x\), de sorte que la lecture
  du déterminant est réversible et non indépendante ;
  \item les branches \(W_{12}\to W_{24}\) et
  \(A_2^{12}\to\Lambda_{24}\) comme constructions séparées ;
  \item le lecteur relatif \HTr--\HAdd--\HDef--\HRes--\HLoc distingué d’une
  transition totale, dont la mesure image dynamique constitue une
  construction supplémentaire.
\end{enumerate}

Ces conditions ne sont pas des notes éditoriales. Elles font partie de la structure logique du résultat. \endgroup

\chapter{Équivalence des systèmes de coordonnées global et analytique du couple angulaire}
La mesure positive postérieure au retour s'obtient avant le contre-angle. La branche déterminantale détermine la valuation d'action et les deux constructions convergent dans la section de la droite dimensionnelle. Alors seulement se forment l'angle, son complément orienté et les caractères \(q_{+}\) et \(q_{-}\), qui évaluent les incidences déjà engendrées par le noyau.
% Adaptador único del capítulo 38 del sucesor 102.
% Sustituye exclusivamente la jerarquía de encabezados. Los cuerpos de los
% dos fuentes teoremáticas científicos y del delta Ley 9 se conservan línea a línea,
% reordenados conforme al DAG causal aprobado. Véase MAPA_SEGMENTOS_CAPITULO_38.tsv.

% HMT_SOURCE_SEGMENT_BEGIN id=B00 source=colaboracion/parte_iii/source/public_final/base_83/c30_body.tex body_lines=1-2
\label{chap:p3-83-30-angulo_contraangulo}

% HMT_SOURCE_SEGMENT_END id=B00
% HMT_SOURCE_SEGMENT_BEGIN id=A00 source=colaboracion/parte_iii/source/public_final/ascensos_20260826/16f_realizacion_analitica_contraangulo.tex body_lines=1-4
\label{p3-20260826:chap:realizacion-analitica-contraangulo}
\index{contre-angle}
\index{microfibre de pi@microfibre de \(\pi\)!caractère analytique}
\index{caractère déterminantal}
% HMT_SOURCE_SEGMENT_END id=A00
% HMT_SOURCE_SEGMENT_BEGIN id=D00 source=manuscrito/sucesor_102/deltas_ley9/c38_par_angular_cartas_y_canales.tex body_lines=1
% Distribución causal exacta del delta Ley 9; contenido conservado sin síntesis.
% HMT_SOURCE_SEGMENT_END id=D00

\section{Phase angulaire intrinsèque \texorpdfstring{\(\alpha\)}{alpha HMT}, relèvement nonadique \texorpdfstring{\(10^3\)}{10^3} et précurseur \texorpdfstring{\(x_A\)}{x_A}}
\label{sec:s102:c38:1}

\subsection{Construction du précurseur analytique \texorpdfstring{\(x_A\)}{x_A}, de la contraction \texorpdfstring{\(D_A\)}{D_A} et de la correction régionale \texorpdfstring{\(\mathscr C_\pi(\alpha)\)}{C_pi(alpha)}}
\label{sec:s102:c38:1:1}

% HMT_SOURCE_SEGMENT_BEGIN id=B01 source=colaboracion/parte_iii/source/public_final/base_83/c30_body.tex body_lines=4-32
\label{p3-83-c30-s1:chap:realizacion-analitica-contraangulo}
\index{coordonnée angulaire conjuguée \(C^*\)}
\index{microfibre de pi@microfibre de \(\pi\)!caractère analytique}
\index{caractère déterminantal}

Le catalogue fini et une évaluation analytique répondent à des questions de types
distincts. Le catalogue détermine quelles régions existent, lesquelles partagent un
retour et quelle information chacune conserve. Une évaluation analytique doit
en outre décider comment la longueur d'un mot se transforme en degré, comment
une branche se prolonge après son retour et avec quel caractère scalaire cette
prolongation se lit. Ce chapitre construit cette seconde opération de manière
explicite.

La construction part de quatre données déjà obtenues dans les chapitres
précédents : la microfibre à cinq régions de \(\pi\), la longueur six de
chaque mot régional, la clôture dodécaphasique de \(\alpha\) et le transport
bicouche associé à la coordonnée commune
\(A=10^3\alpha\). Aucune valeur de la coordonnée angulaire conjuguée \(C_*^{\rm HMT}\), de la composante réduite
d'action, de la fonctionnelle de Barbero--Immirzi ou d'une constante métrologique n'est
employée dans l'évaluation. Le résultat est un caractère régional convergent,
une récurrence exacte et une phase orientée dont la carte sexagésimale reproduit
la coordonnée angulaire conjuguée publiée.

La séparation entre la partie finie et la partie analytique est essentielle. L'entier
\(169\) sera un théorème du catalogue. La série qui le prolonge sera le
théorème d'un lecteur analytique dont les règles sont déclarées avant de l'évaluer.
Ainsi, les prémisses supplémentaires ne restent pas dissimulées dans une coïncidence
décimale.

% HMT_SOURCE_SEGMENT_END id=B01
\section{Construction préangulaire de \texorpdfstring{\(\eta_{\mathrm{ret}}\)}{eta_ret} au moyen de \texorpdfstring{\(H_5\)}{H5}, \texorpdfstring{\(D_A\)}{D_A} et \texorpdfstring{\(\mathscr C_\pi(\alpha)\)}{C_pi(alpha)}}
\label{sec:s102:c38:2}

\subsection{Dérivation indépendante de \texorpdfstring{\(\eta_{\mathrm{ret}}=H_5-(90/\pi)\mathscr C_\pi(\alpha)\)}{eta_ret=H5-(90/pi)C_pi(alpha)} au moyen du lecteur analytique régional}
\label{sec:s102:c38:2:1}

% HMT_SOURCE_SEGMENT_BEGIN id=A01 source=colaboracion/parte_iii/source/public_final/ascensos_20260826/16f_realizacion_analitica_contraangulo.tex body_lines=5-38,40-47

Le catalogue fini et une évaluation analytique répondent à des questions de types
distincts. Le catalogue détermine quelles régions existent, lesquelles partagent un
retour et quelle information chacune conserve. Une évaluation analytique doit
décider, en outre, comment la longueur d'un mot se transforme en degré, comment
une branche se prolonge après son retour et avec quel caractère scalaire on lit
cette prolongation. Cette section construit cette seconde opération de manière
explicite.

La construction part de la microfibre de cinq régions de \(\pi\), de la
longueur six de chaque mot régional, de la clôture dodécaphasique de \(\alpha\)
et de l'insertion angulaire déclarée
\[
 \iota_\angle(x):=10^3x\,u_\circ,\qquad
 \Adeg:=10^3\alpha .
\]
Ici, \(u_\circ\) est la base sexagésimale dont le tour complet a pour
coordonnée \(360\).
Le facteur \(10^3\) appartient à cette normalisation de carte et ne se déduit pas de la
clôture arithmétique de \(\alpha\). Le transport bicouche est construit ensuite
à partir de \(\Adeg\). La fonction qui évalue le lecteur ne reçoit comme
argument ni le contre-angle, ni la coordonnée postérieure au retour, ni la fonctionnelle
hexade--octade, ni une constante métrologique. Le résultat est un caractère
régional convergent, une récurrence exacte et une phase orientée dans la carte
sexagésimale déclarée.

La séparation entre la partie finie et la partie analytique est essentielle. L'entier
\(169\) est un théorème du catalogue. La série qui le prolonge appartient
à un lecteur analytique dont les règles sont déclarées avant son évaluation.

La phase est notée \(y_*\) et, lorsque l'on souligne son origine régionale,
\(y_{\rm reg}\) ; les deux symboles désignent le même lecteur construit
ci-après.


L'évaluation reçoit le catalogue de \(468\) émissions, la valeur exacte de
\(\alpha\), \(\varphi=(1+\sqrt5)/2\), la formule fixée de \(H_5\),
la carte angulaire et les cinq règles du lecteur régional. Avec ces données,
elle calcule \(169\), \(D_A\), la récurrence et \(C^*\). Le théorème est exact dans
cette classe déclarée de lecteurs ; il n'utilise pas \(C^*\) comme entrée et n'affirme pas que
les règles de \(H_5\) constituent l'unique classe naturelle possible.

% HMT_SOURCE_SEGMENT_END id=A01
\subsection{Persistance du retour dans les \texorpdfstring{\(5\)}{5} régions de la microfibre}
\label{sec:s102:c38:2:2}

% HMT_SOURCE_SEGMENT_BEGIN id=B02 source=colaboracion/parte_iii/source/public_final/base_83/c30_body.tex body_lines=34-155

Soit
\[
\mathcal F_\pi=\Psi^{-1}(010211)\subset\mathcal U_6
\]
la microfibre relative au catalogue. Rappelons ses cinq éléments, maintenant séparés en bloc initial et bloc terminal :
\[
\begin{aligned}
 &(501,614,498\mid169,272,272),\\
 &(810,923,870\mid169,272,272),\\
 &(810,983,810\mid169,272,272),\\
 &(870,923,810\mid169,272,272),\\
 &(870,923,810\mid418,674,521).
\end{aligned}
\]
Nous définissons la projection de retour
\[
p_{\rm ret}:\mathbb Z^6\longrightarrow\mathbb Z^3,
\qquad
p_{\rm ret}(u_1,\ldots,u_6)=(u_4,u_5,u_6).
\]

\begin{proposition}[Bloc terminal persistant]
\label{p3-83-c30-s1:prop:bloque-terminal-persistente}
Le multiensemble \(p_{\rm ret}(\mathcal F_\pi)\) possède un unique élément de
multiplicité maximale :
\[
b_\pi=(169,272,272),
\qquad
\operatorname{mult}_{\mathcal F_\pi}(b_\pi)=4.
\]
L'unique bloc terminal restant est
\[
b_\pi^-=(418,674,521)
\]
et a une multiplicité de un.
\end{proposition}

\begin{proof}
L'affirmation s'obtient par énumération exacte des cinq lignes du catalogue qui satisfont \(\Psi(U)=010211\). Quatre lignes possèdent le bloc terminal \((169,272,272)\) ; la cinquième, qui porte le retour muté, possède \((418,674,521)\). La multiplicité maximale est quatre et n'est atteinte qu'une seule fois.
\end{proof}

La différence orientée entre le retour muté et le retour persistant est
\[
 \omega_\pi=b_\pi^- - b_\pi=(249,402,249),
 \qquad
 \langle(1,1,1),\omega_\pi\rangle=900.
\]
Cette identité sépare deux objets qu'il ne faut pas confondre. L'entier
\(169\) est une coordonnée d'une donnée terminale, c'est-à-dire une donnée de degré
zéro. Le vecteur \(\omega_\pi\) peut être lu comme une cochaîne de degré un
après orientation de l'arête qui relie les deux blocs. L'appeler cocycle exigerait,
en outre, de fixer un complexe de cochaînes et de vérifier que son cobord s'annule.
La prolongation analytique construite plus loin ne présuppose pas cette
identification.

La sélection de \(169\) n'exige pas de privilégier la première position du bloc.
Le groupe symétrique \(\mathfrak S_3\) agit en permutant ses trois coordonnées.
Le stabilisateur de \(b_\pi\) échange les deux occurrences de \(272\) et
fixe l'unique coordonnée de multiplicité un. Notons
\(\operatorname{sing}(b)\) cette coordonnée lorsque le bloc est de type
\((r,s,s)\), avec \(r\ne s\). De manière équivalente,
\[
\operatorname{sing}(b)
=\sum_{j=1}^3b_j-2\operatorname{moda}(b).
\]
En particulier,
\[
\boxed{\operatorname{sing}(b_\pi)=169.}
\]

\begin{definition}[Sélecteur résiduel de la microfibre]
\label{p3-83-c30-s1:def:selector-residual-pi}
Le sélecteur \(\rho_\pi\) applique successivement deux opérations :
\begin{enumerate}
  \item sélectionne le bloc terminal de multiplicité maximale au sein de \(p_{\rm ret}(\mathcal F_\pi)\) ;
  \item retient la coordonnée fixe sous l'action du stabilisateur de ce bloc.
\end{enumerate}
Par la \cref{p3-83-c30-s1:prop:bloque-terminal-persistente},
\[
\boxed{\rho_\pi=169.}
\]
\end{definition}

\begin{theorem}[Unicité équivariante du retour]
\label{p3-83-c30-s1:thm:unicidad-retorno-169}
Parmi les sélecteurs qui
\begin{enumerate}
  \item sont invariants sous les permutations des cinq régions ;
  \item sont équivariants sous les permutations des trois coordonnées
  terminales ;
  \item sélectionnent le bloc de multiplicité maximale ;
  \item retiennent une coordonnée fixe sous l'action de son stabilisateur,
\end{enumerate}
la valeur de retour est déterminée de manière unique et vaut \(169\).
\end{theorem}

\begin{proof}
L'invariance sous \(\mathfrak S_5\) empêche d'utiliser l'ordre des lignes. La
troisième condition sélectionne \(b_\pi\), qui est l'unique mode d'après la
\cref{p3-83-c30-s1:prop:bloque-terminal-persistente}. Son stabilisateur dans
\(\mathfrak S_3\) possède deux orbites sur les coordonnées : la paire de
valeurs \(272\) et la valeur singleton \(169\). La quatrième condition sélectionne la
seconde orbite. L'équivariance conserve sa valeur lorsque sa position est déplacée.
\end{proof}

Cet argument améliore la lecture positionnelle du retour. Le \(169\) n'est pas
« la première coordonnée » d'une ligne choisie : c'est un invariant du
multiensemble régional et de l'action de ses symétries de réordonnancement.
Il n'est pas non plus sélectionné pour sa rareté globale. Dans les \(468\) émissions du
catalogue, \(169\) apparaît \(84\) fois, uniformément réparti avec
\(14\) occurrences dans chacune des six positions. L'unicité du
\cref{p3-83-c30-s1:thm:unicidad-retorno-169} est donc relationnelle et régionale : elle réside
dans la combinaison de la microfibre, de la persistance modale et du stabilisateur, non dans la
fréquence isolée de l'entier.

% HMT_SOURCE_SEGMENT_END id=B02
\subsection{Construction du lecteur analytique régional et de son domaine de convergence}
\label{sec:s102:c38:2:3}

% HMT_SOURCE_SEGMENT_BEGIN id=B05 source=colaboracion/parte_iii/source/public_final/base_83/c30_body.tex body_lines=261-301

L'extension analytique est fixée par les règles suivantes.

\begin{definition}[Lecteur régional déterminantal]
\label{p3-83-c30-s1:def:lector-regional-determinantal}
Le lecteur régional déterminantal satisfait :
\begin{enumerate}
  \item \emph{compatibilité de degré} : la longueur cylindrique est le degré
  analytique;
  \item \emph{décomposition de renouvellement} : l'impulsion finie de retour et
  la continuation stricte appartiennent à des termes additifs distincts ;
  \item \emph{normalisation de continuation} : après l'enregistrement, une seule
  fois, du résidu \(\rho_\pi\), la branche survivante commence à l'unité de
  la droite déterminant ;
  \item \emph{covariance par concaténation} : chaque prolongation élémentaire
  agit au moyen de
  \(\bigwedge^2\mathcal T=D_A\) ;
  \item \emph{orientation} : dans la convention cardinale fixée pour APP, le retour régional appartient au secteur compensateur et est soustrait de la phase d'action de cinquième degré.
\end{enumerate}
\end{definition}

La troisième règle a un contenu mathématique. La valeur \(169\) est enregistrée
une fois comme amplitude du retour ; elle n'est pas multipliée à nouveau à chaque
prolongation. La continuation stricte compte l'unique droite qui se poursuit et
la normalise par son unité. C'est la différence entre une récurrence de
renouvellement et l'itération homogène de l'impulsion initiale.

La nécessité de déclarer cette normalisation peut être prouvée. Pour tout
\(\lambda\in\mathbb R\), la famille
\[
F_\lambda(z)
=169z^6+\lambda\frac{D_Az^7}{1-D_Az}
\]
conserve le même retour fini et le même rapport déterminantal entre
coefficients consécutifs. Le catalogue et le rapport seuls ne distinguent pas
\(\lambda\). L'unité de la droite déterminant fixe
\(\lambda=1\) avant d'évaluer \(z=\alpha\). De cette manière, la normalisation
ne s'obtient pas en ajustant la valeur finale ; elle fait partie de la définition du
caractère.

% HMT_SOURCE_SEGMENT_END id=B05
\subsection{Résolution analytique de la récurrence de retour}
\label{sec:s102:c38:2:4}

% HMT_SOURCE_SEGMENT_BEGIN id=B06 source=colaboracion/parte_iii/source/public_final/base_83/c30_body.tex body_lines=303-364

Nous définissons les coefficients \(\kappa_n\) par
\[
\kappa_6=169,
\qquad
\kappa_7=D_A,
\qquad
\kappa_{n+1}=D_A\kappa_n\quad(n\ge7).
\]
La première égalité est l'impulsion régionale ; les deux suivantes décrivent la
continuation normalisée.

\begin{theorem}[Réalisation analytique graduée]
\label{p3-83-c30-s1:thm:realizacion-analitica-graduada}
Le germe analytique compatible avec les règles de la
\cref{p3-83-c30-s1:def:lector-regional-determinantal} est unique et vaut
\[
\boxed{
\mathscr C_\pi(z)
=169z^6+\frac{D_Az^7}{1-D_Az}.}
\]
Ses coefficients satisfont
\[
\kappa_n=D_A^{n-6}\qquad(n\ge7).
\]
\end{theorem}

\begin{proof}
Écrivons
\[
F(z)=169z^6+\sum_{k\ge1}c_kz^{6+k}.
\]
La normalisation de continuation et la covariance déterminantale donnent
\[
c_1=D_A,
\qquad
c_{k+1}=D_Ac_k.
\]
Par récurrence, \(c_k=D_A^k\). La somme géométrique produit
\[
\sum_{k\ge1}D_A^kz^{6+k}
=\frac{D_Az^7}{1-D_Az}.
\]
Aucun coefficient libre ne subsiste.
\end{proof}

Comme \(0<\alpha<1\) et
\(0<D_A<1\),
\[
0<D_A\alpha<1.
\]
En conséquence, \(\mathscr C_\pi(\alpha)\) converge absolument et son
rayon de convergence est \(D_A^{-1}>1\). Numériquement,
\[
D_A=0.7751291192471564301888840964802\ldots,
\]
\[
D_A\alpha=0.00565639046986492673885079095469\ldots.
\]
La rapidité de cette contraction explique pourquoi la somme complète et sa
troncature de degré sept possèdent les mêmes quinze premiers chiffres dans la
carte finale.

% HMT_SOURCE_SEGMENT_END id=B06
\subsection{Persistance du retour dans les \texorpdfstring{\(5\)}{5} régions et compatibilité de la composante \texorpdfstring{\(\eta_{\mathrm{ret}}\)}{eta_ret}}
\label{sec:s102:c38:2:5}

% HMT_SOURCE_SEGMENT_BEGIN id=A02 source=colaboracion/parte_iii/source/public_final/ascensos_20260826/16f_realizacion_analitica_contraangulo.tex body_lines=49-170

Soit
\[
\mathcal F_\pi=\Psi^{-1}(010211)\subset\mathcal U_6
\]
la microfibre relative au catalogue. Rappelons ses cinq éléments, maintenant séparés en bloc initial et bloc terminal :
\[
\begin{aligned}
 &(501,614,498\mid169,272,272),\\
 &(810,923,870\mid169,272,272),\\
 &(810,983,810\mid169,272,272),\\
 &(870,923,810\mid169,272,272),\\
 &(870,923,810\mid418,674,521).
\end{aligned}
\]
Nous définissons la projection de retour
\[
p_{\rm ret}:\mathbb Z^6\longrightarrow\mathbb Z^3,
\qquad
p_{\rm ret}(u_1,\ldots,u_6)=(u_4,u_5,u_6).
\]

\begin{proposition}[Bloc terminal persistant]
\label{p3-20260826:prop:bloque-terminal-persistente}
Le multiensemble \(p_{\rm ret}(\mathcal F_\pi)\) possède un unique élément de
multiplicité maximale :
\[
b_\pi=(169,272,272),
\qquad
\operatorname{mult}_{\mathcal F_\pi}(b_\pi)=4.
\]
L'unique bloc terminal restant est
\[
b_\pi^-=(418,674,521)
\]
et a une multiplicité de un.
\end{proposition}

\begin{proof}
L'affirmation s'obtient par énumération exacte des cinq lignes du catalogue qui satisfont \(\Psi(U)=010211\). Quatre lignes possèdent le bloc terminal \((169,272,272)\) ; la cinquième, qui porte le retour muté, possède \((418,674,521)\). La multiplicité maximale est quatre et n'est atteinte qu'une seule fois.
\end{proof}

La différence orientée entre le retour muté et le retour persistant est
\[
 \omega_\pi=b_\pi^- - b_\pi=(249,402,249),
 \qquad
 \langle(1,1,1),\omega_\pi\rangle=900.
\]
Cette identité sépare deux objets qu'il ne faut pas confondre. L'entier
\(169\) est une coordonnée d'une donnée terminale, c'est-à-dire une donnée de degré
zéro. Le vecteur \(\omega_\pi\) peut être lu comme une cochaîne de degré un
après orientation de l'arête qui relie les deux blocs. L'appeler cocycle exigerait,
en outre, de fixer un complexe de cochaînes et de vérifier que son cobord s'annule.
La prolongation analytique construite plus loin ne présuppose pas cette
identification.

La sélection de \(169\) n'exige pas de privilégier la première position du bloc.
Le groupe symétrique \(\mathfrak S_3\) agit en permutant ses trois coordonnées.
Le stabilisateur de \(b_\pi\) échange les deux occurrences de \(272\) et
fixe l'unique coordonnée de multiplicité un. Notons
\(\operatorname{sing}(b)\) cette coordonnée lorsque le bloc est de type
\((r,s,s)\), avec \(r\ne s\). De manière équivalente,
\[
\operatorname{sing}(b)
=\sum_{j=1}^3b_j-2\operatorname{moda}(b).
\]
En particulier,
\[
\boxed{\operatorname{sing}(b_\pi)=169.}
\]

\begin{definition}[Sélecteur résiduel de la microfibre]
\label{p3-20260826:def:selector-residual-pi}
Le sélecteur \(\rho_\pi\) applique successivement deux opérations :
\begin{enumerate}
  \item sélectionne le bloc terminal de multiplicité maximale au sein de \(p_{\rm ret}(\mathcal F_\pi)\) ;
  \item retient la coordonnée fixe sous l'action du stabilisateur de ce bloc.
\end{enumerate}
Par la \cref{p3-20260826:prop:bloque-terminal-persistente},
\[
\boxed{\rho_\pi=169.}
\]
\end{definition}

\begin{theorem}[Unicité équivariante du retour]
\label{p3-20260826:thm:unicidad-retorno-169}
Parmi les sélecteurs qui
\begin{enumerate}
  \item sont invariants sous les permutations des cinq régions ;
  \item sont équivariants sous les permutations des trois coordonnées
  terminales ;
  \item sélectionnent le bloc de multiplicité maximale ;
  \item retiennent une coordonnée fixe sous l'action de son stabilisateur,
\end{enumerate}
la valeur de retour est déterminée de manière unique et vaut \(169\).
\end{theorem}

\begin{proof}
L'invariance sous \(\mathfrak S_5\) empêche d'utiliser l'ordre des lignes. La
troisième condition sélectionne \(b_\pi\), qui est l'unique mode d'après la
\cref{p3-20260826:prop:bloque-terminal-persistente}. Son stabilisateur dans
\(\mathfrak S_3\) possède deux orbites sur les coordonnées : la paire de
valeurs \(272\) et la valeur singleton \(169\). La quatrième condition sélectionne la
seconde orbite. L'équivariance conserve sa valeur lorsque sa position est déplacée.
\end{proof}

Cet argument améliore la lecture positionnelle du retour. Le \(169\) n'est pas
« la première coordonnée » d'une ligne choisie : c'est un invariant du
multiensemble régional et de l'action de ses symétries de réordonnancement.
Il n'est pas non plus sélectionné pour sa rareté globale. Dans les \(468\) émissions du
catalogue, \(169\) apparaît \(84\) fois, uniformément réparti avec
\(14\) occurrences dans chacune des six positions. L'unicité du
\cref{p3-20260826:thm:unicidad-retorno-169} est donc relationnelle et régionale : elle réside
dans la combinaison de la microfibre, de la persistance modale et du stabilisateur, non dans la
fréquence isolée de l'entier.

% HMT_SOURCE_SEGMENT_END id=A02
\subsection{Application du lecteur régional au couple angulaire conjugué}
\label{sec:s102:c38:2:6}

% HMT_SOURCE_SEGMENT_BEGIN id=A05 source=colaboracion/parte_iii/source/public_final/ascensos_20260826/16f_realizacion_analitica_contraangulo.tex body_lines=277-317

L'extension analytique est fixée par les règles suivantes.

\begin{definition}[Lecteur régional déterminantal]
\label{p3-20260826:def:lector-regional-determinantal}
Le lecteur régional déterminantal satisfait :
\begin{enumerate}
  \item \emph{compatibilité de degré} : la longueur cylindrique est le degré
  analytique;
  \item \emph{décomposition de renouvellement} : l'impulsion finie de retour et
  la continuation stricte appartiennent à des termes additifs distincts ;
  \item \emph{normalisation de continuation} : après l'enregistrement, une seule
  fois, du résidu \(\rho_\pi\), la branche survivante commence à l'unité de
  la droite déterminant ;
  \item \emph{covariance par concaténation} : chaque prolongation élémentaire
  agit au moyen de
  \(\bigwedge^2\mathcal T=D_A\) ;
  \item \emph{orientation} : dans la convention cardinale fixée pour APP, le retour régional appartient au secteur compensateur et est soustrait de la phase d'action de cinquième degré.
\end{enumerate}
\end{definition}

La troisième règle a un contenu mathématique. La valeur \(169\) est enregistrée
une fois comme amplitude du retour ; elle n'est pas multipliée à nouveau à chaque
prolongation. La continuation stricte compte l'unique droite qui se poursuit et
la normalise par son unité. C'est la différence entre une récurrence de
renouvellement et l'itération homogène de l'impulsion initiale.

La nécessité de déclarer cette normalisation peut être prouvée. Pour tout
\(\lambda\in\mathbb R\), la famille
\[
F_\lambda(z)
=169z^6+\lambda\frac{D_Az^7}{1-D_Az}
\]
conserve le même retour fini et le même rapport déterminantal entre
coefficients consécutifs. Le catalogue et le rapport seuls ne distinguent pas
\(\lambda\). L'unité de la droite déterminant fixe
\(\lambda=1\) avant d'évaluer \(z=\alpha\). De cette manière, la normalisation
ne s'obtient pas en ajustant la valeur finale ; elle fait partie de la définition du
caractère.

% HMT_SOURCE_SEGMENT_END id=A05
\subsection{Somme fermée et publication de la composante postérieure au retour}
\label{sec:s102:c38:2:7}

% HMT_SOURCE_SEGMENT_BEGIN id=A06 source=colaboracion/parte_iii/source/public_final/ascensos_20260826/16f_realizacion_analitica_contraangulo.tex body_lines=319-380

Nous définissons les coefficients \(\kappa_n\) par
\[
\kappa_6=169,
\qquad
\kappa_7=D_A,
\qquad
\kappa_{n+1}=D_A\kappa_n\quad(n\ge7).
\]
La première égalité est l'impulsion régionale ; les deux suivantes décrivent la
continuation normalisée.

\begin{theorem}[Réalisation analytique graduée]
\label{p3-20260826:thm:realizacion-analitica-graduada}
Le germe analytique compatible avec les règles de la
\cref{p3-20260826:def:lector-regional-determinantal} est unique et vaut
\[
\boxed{
\mathscr C_\pi(z)
=169z^6+\frac{D_Az^7}{1-D_Az}.}
\]
Ses coefficients satisfont
\[
\kappa_n=D_A^{n-6}\qquad(n\ge7).
\]
\end{theorem}

\begin{proof}
Écrivons
\[
F(z)=169z^6+\sum_{k\ge1}c_kz^{6+k}.
\]
La normalisation de continuation et la covariance déterminantale donnent
\[
c_1=D_A,
\qquad
c_{k+1}=D_Ac_k.
\]
Par récurrence, \(c_k=D_A^k\). La somme géométrique produit
\[
\sum_{k\ge1}D_A^kz^{6+k}
=\frac{D_Az^7}{1-D_Az}.
\]
Aucun coefficient libre ne subsiste.
\end{proof}

Comme \(0<\alpha<1\) et
\(0<D_A<1\),
\[
0<D_A\alpha<1.
\]
En conséquence, \(\mathscr C_\pi(\alpha)\) converge absolument et son
rayon de convergence est \(D_A^{-1}>1\). Numériquement,
\[
D_A=0.7751291192471564301888840964802\ldots,
\]
\[
D_A\alpha=0.00565639046986492673885079095469\ldots.
\]
La rapidité de cette contraction explique pourquoi la somme complète et sa
troncature de degré sept possèdent les mêmes quinze premiers chiffres dans la
carte finale.

% HMT_SOURCE_SEGMENT_END id=A06
\section{Confluence de \texorpdfstring{\(\eta_{\mathrm{ret}}\)}{eta_ret} avec la valuation \texorpdfstring{\(\varepsilon_H=-34\)}{epsilon_H=-34} dans la grandeur \texorpdfstring{\(\hbar_{\mathrm{ret}}\)}{hbar_ret}}
\label{sec:s102:c38:3}

\subsection{Dérivation du caractère déterminantal du transport parangulaire}
\label{sec:s102:c38:3:1}

% HMT_SOURCE_SEGMENT_BEGIN id=B04 source=colaboracion/parte_iii/source/public_final/base_83/c30_body.tex body_lines=196-259

La fermeture dodécaphasique fournit \(\alpha\). La complétion des trois
couples nonadiques fournit l'échelle \(10^3\), de sorte que
\[
A=10^3\alpha.
\]
La carte angulaire archimédienne étant choisie, la coordonnée commune en radians est
\[
x=\frac{\pi A}{180}=\frac{50\pi}{9}\alpha.
\]
Soit \(R\) une involution réelle de trace nulle en dimension deux,
\[
R^2=I_2,
\qquad
\operatorname{tr}R=0,
\]
et considérons le transport formel bicouche
\[
\mathcal T(x,y)=\exp(-xI_2-yR).
\]
La variable orientée \(y\) n'a pas encore été évaluée. Cependant, l'action
de \(\mathcal T\) sur la droite déterminante en est indépendante :
\[
\begin{aligned}
\det\mathcal T(x,y)
&=\exp\!\left(\operatorname{tr}(-xI_2-yR)\right)\\
&=e^{-2x}.
\end{aligned}
\]
On définit
\[
\boxed{
D_A=e^{-2x}=e^{-100\pi\alpha/9}.}
\]
L'indice rappelle qu'il s'agit du déterminant de la contraction commune associée à \(A\), et non d'une fonction de \(C_*^{\rm HMT}\).

\begin{proposition}[Caractère de la droite déterminante]
\label{p3-83-c30-s1:prop:caracter-determinante}
L'application
\[
\delta_A:(\mathbb N_0,+)\longrightarrow(\mathbb R_{>0},\cdot),
\qquad
\delta_A(k)=D_A^k,
\]
est l'unique caractère multiplicatif dont la valeur sur une prolongation
élémentaire est \(D_A\).
\end{proposition}

\begin{proof}
La multiplicativité donne
\[
\delta_A(k)
=\delta_A(\underbrace{1+\cdots+1}_{k\text{ fois}})
=\delta_A(1)^k=D_A^k.
\]
La même identité prouve l'existence et l'unicité.
\end{proof}

Le déterminant reste invariant par conjugaison de \(\mathcal T\) et par échange des feuillets \(R\mapsto-R\). La queue construite à partir de \(D_A\) appartient donc au mode commun. L'orientation apparaîtra dans le signe avec lequel le caractère régional corrige la phase.

% HMT_SOURCE_SEGMENT_END id=B04
\subsection{Caractère d'action de degré \texorpdfstring{\(5\)}{5} dans la construction régionale}
\label{sec:s102:c38:3:2}

% HMT_SOURCE_SEGMENT_BEGIN id=B07 source=colaboracion/parte_iii/source/public_final/base_83/c30_body.tex body_lines=366-405

Le bloc central
\[
\mathcal B_c=\begin{pmatrix}7&2\\2&7\end{pmatrix}
\]
possède les valeurs propres \(9\) et \(5\). Le cycle dodécaphasique apporte le rang
\(12\) ; ses orbites de pas trois possèdent une longueur de quatre ; et le recensement complet
apporte le quotient exact
\[
\frac{104976}{1944}=54.
\]
À partir de ces invariants, on fixe le caractère d'action de cinquième degré
\[
\boxed{
H_5(\alpha,\varphi)
=\frac12\sqrt{\frac{1000\alpha}{\varphi}}-\alpha
+\frac9{16}\alpha^2-\frac59\alpha^3
+\frac7{48}\alpha^4-\frac1{54}\alpha^5.}
\]
Les quatre coefficients rationnels peuvent s'écrire
\[
\frac9{16}=\frac{\lambda_+}{4^2},
\qquad
\frac59=\frac{\lambda_-}{\lambda_+},
\qquad
\frac7{48}=\frac{\tfrac12\operatorname{tr}\mathcal B_c}{4\cdot12},
\qquad
\frac1{54}=\frac{1944}{104976},
\]
avec \((\lambda_+,\lambda_-)=(9,5)\). Ces identités certifient la
provenance des nombres utilisés. L'affectation successive de ces
invariants aux degrés deux, trois, quatre et cinq fait partie du caractère
analytique ; elle ne se déduit pas de la seule existence du catalogue fini.

Cette précision élimine une ambiguïté habituelle. Une combinaison
d'invariants de la construction est une définition interne et reproductible ; sa
naturalité exige les règles du lecteur qui spécifient l'ordre, les signes et la
normalisation. Ici, ces règles sont visibles et sont soumises aux mêmes
contrôles négatifs que le reste de la construction.

% HMT_SOURCE_SEGMENT_END id=B07
\subsection{Décomposition exacte de la droite d'action en sections pré-retour et post-retour}
\label{sec:s102:c38:3:3}

% HMT_SOURCE_SEGMENT_BEGIN id=B12 source=colaboracion/parte_iii/source/public_final/base_83/c30_body.tex body_lines=571-617

Le caractère de cinquième degré et la composante d'action postérieure au retour
sont deux valeurs distinctes :
\[
H_5
=1.05457181764615446962582182943306\ldots,
\]
\[
\bar h_{\rm HMT}^{\rm ret}
=1.05457181691503906113369058472430\ldots.
\]
La première est l'action locale avant l'incorporation du caractère régional ; la
seconde est l'action orientée qui reste après le retour. La correction
qui engendre la coordonnée angulaire conjuguée empêche de les identifier.

La comparaison métrologique doit respecter cette distinction. La première valeur donne
la mantisse
\[
2\pi H_5
=6.62607014999998792600111034989947\ldots,
\]
tandis que l'action postérieure donne
\[
2\pi\bar h_{\rm HMT}^{\rm ret}
=6.62607014540625433351074984759288\ldots.
\]
Dans le Système international, la valeur
\[
h=6.62607015\times10^{-34}\;\mathrm{J\,s}
\]
est exacte par définition \cite{BIPMSI}. C'est pourquoi,
\[
2\pi H_5-6.62607015
=-1.2073998889650101\ldots\times10^{-14},
\]
et ce n'est pas une égalité exacte. La proximité de la valeur antérieure au retour est
un contrôle numérique remarquable ; elle n'autorise pas à remplacer la constante définitoire
du SI par la composante postérieure ni à confondre les deux actions HMT\@.

Ce résultat détermine avec précision la situation mathématique. Le lecteur
régional engendre \(C_*^{\rm HMT}\) et, à partir de celui-ci, la fonctionnelle
\(\Gamma_{6\to8}\). Le même calcul sépare l'approximation de Planck
produite par \(H_5\) de l'action réduite qui intervient dans
la coordonnée angulaire conjuguée. La séparation n'affaiblit pas la construction : elle élimine une
identification incompatible avec les chiffres exacts et transforme chaque
comparaison en une affirmation réfutable de type bien défini.

% HMT_SOURCE_SEGMENT_END id=B12
\subsection{Factorisation du caractère déterminantal dans les \texorpdfstring{\(2\)}{2} cartes de phase}
\label{sec:s102:c38:3:4}

% HMT_SOURCE_SEGMENT_BEGIN id=A04 source=colaboracion/parte_iii/source/public_final/ascensos_20260826/16f_realizacion_analitica_contraangulo.tex body_lines=211-275

La clôture dodécaphasique fournit \(\alpha\). L'insertion angulaire déclarée
emploie l'échelle \(10^3\), compatible avec le regroupement de trois paires
nonadiques, de sorte que
\[
\Adeg=10^3\alpha\,{}^\circ.
\]
La carte angulaire archimédienne étant choisie, la coordonnée commune en radians est
\[
x=: \Arad=\frac{\pi\Adeg}{180}=\frac{50\pi}{9}\alpha.
\]
Soit \(R\) une involution réelle de trace nulle en dimension deux,
\[
R^2=I_2,
\qquad
\operatorname{tr}R=0,
\]
et considérons le transport formel bicouche
\[
\mathcal T(x,y)=\exp(-xI_2-yR).
\]
La variable orientée \(y\) n'a pas encore été évaluée. Cependant, l'action
de \(\mathcal T\) sur la droite déterminante en est indépendante :
\[
\begin{aligned}
\det\mathcal T(x,y)
&=\exp\!\left(\operatorname{tr}(-xI_2-yR)\right)\\
&=e^{-2x}.
\end{aligned}
\]
On définit
\[
\boxed{
D_A=e^{-2x}=e^{-100\pi\alpha/9}.}
\]
L'indice rappelle qu'il s'agit du déterminant de la contraction commune
associée à \(A\), non d'une fonction du contre-angle.

\begin{proposition}[Caractère de la droite déterminante]
\label{p3-20260826:prop:caracter-determinante}
L'application
\[
\delta_A:(\mathbb N_0,+)\longrightarrow(\mathbb R_{>0},\cdot),
\qquad
\delta_A(k)=D_A^k,
\]
est l'unique caractère multiplicatif dont la valeur sur une prolongation
élémentaire est \(D_A\).
\end{proposition}

\begin{proof}
La multiplicativité donne
\[
\delta_A(k)
=\delta_A(\underbrace{1+\cdots+1}_{k\text{ fois}})
=\delta_A(1)^k=D_A^k.
\]
La même identité prouve l'existence et l'unicité.
\end{proof}

Le déterminant reste invariant par conjugaison de \(\mathcal T\) et par échange des feuillets \(R\mapsto-R\). La queue construite à partir de \(D_A\) appartient donc au mode commun. L'orientation apparaîtra dans le signe avec lequel le caractère régional corrige la phase.

% HMT_SOURCE_SEGMENT_END id=A04
\subsection{Publication du caractère d'action de degré \texorpdfstring{\(5\)}{5} dans la carte parangulaire}
\label{sec:s102:c38:3:5}

% HMT_SOURCE_SEGMENT_BEGIN id=A07 source=colaboracion/parte_iii/source/public_final/ascensos_20260826/16f_realizacion_analitica_contraangulo.tex body_lines=382-418

Le bloc central
\[
\mathcal B_c=\begin{pmatrix}7&2\\2&7\end{pmatrix}
\]
possède les valeurs propres \(9\) et \(5\). Le cycle dodécaphasique apporte le rang
\(12\) ; ses orbites de pas trois possèdent une longueur de quatre ; et le recensement complet
apporte le quotient exact
\[
\frac{104976}{1944}=54.
\]
Le lecteur en vigueur déclare le caractère d'action de cinquième degré
\[
\boxed{
H_5(\alpha,\varphi)
=\frac12\sqrt{\frac{1000\alpha}{\varphi}}-\alpha
+\frac9{16}\alpha^2-\frac59\alpha^3
+\frac7{48}\alpha^4-\frac1{54}\alpha^5.}
\]
Les quatre coefficients rationnels peuvent s'écrire
\[
\frac9{16}=\frac{\lambda_+}{4^2},
\qquad
\frac59=\frac{\lambda_-}{\lambda_+},
\qquad
\frac7{48}=\frac{\tfrac12\operatorname{tr}\mathcal B_c}{4\cdot12},
\qquad
\frac1{54}=\frac{1944}{104976},
\]
avec \((\lambda_+,\lambda_-)=(9,5)\). Ces identités fixent la
représentabilité arithmétique interne des coefficients du caractère. Leur
ordre, leurs signes et leur normalisation appartiennent à la règle graduée du lecteur
régional et sont appliqués avant toute carte d'unités. Les ablations
certifient que chaque terme intervient effectivement dans la phase résultante.
En particulier, ni la mantisse SI de l'action ni la valeur sexagésimale du
contre-angle ne font partie du domaine de la fonction génératrice.

% HMT_SOURCE_SEGMENT_END id=A07
\section{Application parangulaire \texorpdfstring{\(P(\alpha,\eta_{\mathrm{ret}})\)}{P(alpha,eta_ret)} et clôture bicouche \texorpdfstring{\(C^*=2(\eta_{\mathrm{ret}}+\alpha)\)}{C*=2(eta_ret+alpha)}}
\label{sec:s102:c38:4}

\subsection{Phase orientée et coordonnée angulaire conjuguée}
\label{sec:s102:c38:4:1}

% HMT_SOURCE_SEGMENT_BEGIN id=B08 source=colaboracion/parte_iii/source/public_final/base_83/c30_body.tex body_lines=407-472

La phase de cinquième degré antérieure au retour est
\[
y_5=\frac\pi{90}\bigl(H_5+\alpha\bigr).
\]
Le caractère régional complet produit la correction
\(\mathscr C_\pi(\alpha)\). Avec l'orientation fixée dans la
\cref{p3-83-c30-s1:def:lector-regional-determinantal}, nous définissons
\[
\boxed{
y_*=\frac\pi{90}\bigl(H_5+\alpha\bigr)
-169\alpha^6
-\frac{D_A\alpha^7}{1-D_A\alpha}.}
\]
C'est la formule primaire. L'unité sexagésimale n'intervient qu'à l'application de la carte
\[
C_*^{\rm HMT}=\frac{180}{\pi}y_*.
\]

\begin{theorem}[Coordonnée angulaire conjuguée du lecteur régional déterminantal]
\label{p3-83-c30-s1:thm:contraangulo-regional}
Une fois fixés la clôture dodécaphasique de \(\alpha\),
\(\varphi=(1+\sqrt5)/2\) et les règles du lecteur régional déterminantal, la
phase \(y_*\) et sa coordonnée angulaire conjuguée sont déterminées sans introduire la valeur
recherchée. Leur évaluation est
\[
\boxed{
y_*=0.03706622646583826398493779409230638\ldots,}
\]
\[
\boxed{
C_*^{\rm HMT}
=2.12373833896864572426195138039336\ldots{}^\circ.}
\]
Donc, à quinze décimales,
\[
\boxed{C_*^{\rm HMT}=2.123738338968646^\circ.}
\]
\end{theorem}

\begin{proof}
Le \cref{p3-83-c30-s1:thm:unicidad-retorno-169} détermine \(169\). La clôture de
\(\alpha\) détermine \(A\), \(x\) et
\(D_A=e^{-100\pi\alpha/9}\). Le \cref{p3-83-c30-s1:thm:realizacion-analitica-graduada}
détermine la série complète. La formule de \(H_5\) contient uniquement
\(\alpha\), \(\varphi\) et des invariants structurels déclarés. Substituer
ces données dans l'expression de \(y_*\) donne la valeur imprimée. La valeur décimale
de \(C_*^{\rm HMT}\) n'apparaît dans aucun des membres de droite.
\end{proof}

La composante réduite d'action postérieure au retour est désormais une
sortie :
\[
\boxed{
\bar h_{\rm HMT}^{\rm ret}
=\frac{C_*^{\rm HMT}}2-\alpha
=H_5-\frac{90}{\pi}\mathscr C_\pi(\alpha).}
\]
Numériquement,
\[
\bar h_{\rm HMT}^{\rm ret}
=1.05457181691503906113369058472430\ldots.
\]
L'identité inverse n'a pas été utilisée pour introduire \(C_*^{\rm HMT}\) ; elle
s'obtient après la génération de \(y_*\).

% HMT_SOURCE_SEGMENT_END id=B08
\subsection{Coordonnée conjuguée \texorpdfstring{\(C^*\)}{C*} et réversion de son orientation sous l'involution bicouche}
\label{sec:s102:c38:4:2}

% HMT_SOURCE_SEGMENT_BEGIN id=A08 source=colaboracion/parte_iii/source/public_final/ascensos_20260826/16f_realizacion_analitica_contraangulo.tex body_lines=420-503

La phase de cinquième degré antérieure au retour est
\[
y_5=\frac\pi{90}\bigl(H_5+\alpha\bigr).
\]
Le caractère régional complet produit la correction
\(\mathscr C_\pi(\alpha)\). Avec l'orientation fixée dans la
\cref{p3-20260826:def:lector-regional-determinantal}, nous définissons
\[
\boxed{
\Crad:=y_*=\frac\pi{90}\bigl(H_5+\alpha\bigr)
-169\alpha^6
-\frac{D_A\alpha^7}{1-D_A\alpha}.}
\]
Désormais, nous écrivons aussi
\[
 \boxed{y_{\rm reg}:=y_*}
\]
lorsqu'il est nécessaire de distinguer cette réalisation de la phase spectrale de
\Gtr. Les deux notations désignent ici le même objet régional canonique.
Telle est la formule primaire. L'unité sexagésimale n'intervient qu'à l'application de la carte
\[
\boxed{\Cdeg:=\frac{180}{\pi}\Crad.}
\]
Les symboles \(\Crad\) et \(\Cdeg\) ne désignent pas deux contre-angles : ce sont les
coordonnées radiale et sexagésimale du même élément angulaire.

\begin{theorem}[Contre-angle du lecteur régional déterminantal]
\label{p3-20260826:thm:contraangulo-regional}
La clôture dodécaphasique de \(\alpha\),
\(\varphi=(1+\sqrt5)/2\), le germe \(H_5\) et les règles du lecteur régional
déterminantal étant fixés, la phase \(\Crad\) et son contre-angle sont déterminés sans
introduire la valeur recherchée dans la fonction d'évaluation. Son évaluation est
\[
\boxed{
\Crad=0.03706622646583826398493779409230638\ldots,}
\]
\[
\boxed{
\Cdeg
=2.12373833896864572426195138039336\ldots{}^\circ.}
\]
Donc, à quinze décimales,
\[
\boxed{\Cdeg=2.123738338968646^\circ.}
\]
\end{theorem}

\begin{proof}
Le \cref{p3-20260826:thm:unicidad-retorno-169} détermine \(169\). La clôture de
\(\alpha\) détermine \(\Adeg\), \(\Arad\) et
\(D_A=e^{-100\pi\alpha/9}\). Le \cref{p3-20260826:thm:realizacion-analitica-graduada}
détermine la série complète. La formule fixée de \(H_5\) contient uniquement
\(\alpha\), \(\varphi\) et des invariants structurels déclarés. Substituer
ces données dans l'expression de \(\Crad\) donne la valeur imprimée. La valeur décimale
du contre-angle n'apparaît dans aucun des membres de droite de cette
évaluation. La comparaison avec une mantisse métrologique d'action est effectuée
uniquement après l'obtention de la phase et ne modifie pas la fonction génératrice.
\end{proof}

Soit \(u_\circ\) la base sexagésimale dont le tour complet a pour coordonnée
\(360\), et introduisons la coordonnée angulaire
\(\alpha_{\angle,{\rm deg}}:=\alpha\). La sortie postérieure au
retour est alors typée par
\[
\boxed{
\etaRet
=\frac{\Cdeg}2-\alpha_{\angle,{\rm deg}}
=H_5-\frac{90}{\pi}\mathscr C_\pi(\alpha).}
\]
Sa coordonnée radiale est
\[
 \etaRetrad=\frac{\pi}{180}\etaRet
 =\frac{\Crad}{2}-\frac{\pi}{180}\alpha.
\]
Numériquement,
\[
\etaRet
=1.05457181691503906113369058472430\ldots.
\]
L'identité inverse n'a pas été utilisée pour introduire \(\Cdeg\) ; elle s'obtient
après la génération de \(\Crad\). Le symbole \(\hbar\) reste réservé à
l'élément de la droite d'action construit dans la section suivante.

% HMT_SOURCE_SEGMENT_END id=A08
\subsection{Théorème direct \texorpdfstring{\(C^*=2(\eta_{\mathrm{ret}}+\alpha)\)}{C*=2(eta_ret+alpha)} et caractère inverse de l'identité \texorpdfstring{\(\eta_{\mathrm{ret}}=C^*/2-\alpha\)}{eta_ret=C*/2-alpha}}
\label{sec:s102:c38:4:3}

% HMT_SOURCE_SEGMENT_BEGIN id=A13 source=colaboracion/parte_iii/source/public_final/ascensos_20260826/16f_realizacion_analitica_contraangulo.tex body_lines=750-797

Le caractère de cinquième degré et la coordonnée postérieure au retour
sont deux valeurs distinctes :
\[
H_5
=1.05457181764615446962582182943306\ldots,
\]
\[
\etaRet
=1.05457181691503906113369058472430\ldots.
\]
Le premier est le caractère local avant l'intégration du caractère régional ; le
second est la coordonnée orientée qui subsiste après le retour. La correction
qui génère le contre-angle empêche de les identifier.

La comparaison métrologique doit respecter cette distinction. La première valeur donne
la mantisse
\[
2\pi H_5
=6.62607014999998792600111034989947\ldots,
\]
tandis que la coordonnée ultérieure donne
\[
2\pi\etaRet
=6.62607014540625433351074984759288\ldots.
\]
À titre de contrôle externe uniquement, nous retenons la mantisse \(6.62607015\) de la
valeur définitoire du Système international \cite{BIPMSI}. On n'intègre
pas encore son exposant ni son unité à la droite HMT : la décennie est dérivée dans le
\cref{chap:orden-determinantal-accion} et l'unité appartient à une carte
métrologique ultérieure. La comparaison des mantisses donne
\[
2\pi H_5-6.62607015
=-1.2073998889650101\ldots\times10^{-14},
\]
et ce n'est pas une égalité exacte. La proximité de la valeur antérieure au retour est
un contrôle numérique remarquable ; elle n'autorise pas à remplacer la constante définitoire
du SI par la coordonnée ultérieure ni à confondre une coordonnée angulaire avec
l'élément d'action qui sera construit lors de l'intégration de la décennie.

Ce résultat détermine avec précision la situation mathématique. Le lecteur
régional génère \((\Crad,\Cdeg)\) et, à partir de cet unique élément angulaire, la fonctionnelle
\(\Gamma_{6\to8}\). Le même calcul sépare la proximité externe de la
mantisse produite par \(H_5\) de la coordonnée de retour qui intervient dans le
contre-angle. La séparation n'affaiblit pas la construction : elle élimine une
identification incompatible avec les chiffres exacts et transforme chaque
comparaison en une affirmation falsifiable de type bien défini.

% HMT_SOURCE_SEGMENT_END id=A13
% HMT_SOURCE_SEGMENT_BEGIN id=D01 source=manuscrito/sucesor_102/deltas_ley9/c38_par_angular_cartas_y_canales.tex body_lines=3-20
\label{sec:l9s102:par-angular}

La complétion cubique
\begin{equation}
 10^3=(9+1)^3=9^3+3\cdot9^2+3\cdot9+1
 =729+243+27+1
 \label{eq:l9s102:base-1000}
\end{equation}
détermine la carte globale de la coordonnée directe. La clôture bicouche apporte
le facteur \(2\) de la coordonnée conjuguée :
\begin{equation}
 \boxed{
 A_{\deg}=10^3\alpha,
 \qquad
 C^*_{\deg}=2(\eta_{\mathrm{ret}}+\alpha)
 \quad\text{dans }\mathbb R/(360\mathbb Z).}
 \label{eq:l9s102:a-cstar-deg}
\end{equation}
% HMT_SOURCE_SEGMENT_END id=D01
% HMT_SOURCE_SEGMENT_BEGIN id=D04 source=manuscrito/sucesor_102/deltas_ley9/c38_par_angular_cartas_y_canales.tex body_lines=60-66
L'identité inverse
\begin{equation}
 \eta_{\mathrm{ret}}=\frac12C^*_{\deg}-\alpha
 \label{eq:l9s102:eta-corolario-inverso}
\end{equation}
est établie comme corollaire algébrique de la construction directe.

% HMT_SOURCE_SEGMENT_END id=D04
\section{Commutativité exacte des cartes en tours, degrés et radians}
\label{sec:s102:c38:5}

\subsection{Dérivation de la longueur cylindrique et du degré analytique du précurseur parangulaire}
\label{sec:s102:c38:5:1}

% HMT_SOURCE_SEGMENT_BEGIN id=B03 source=colaboracion/parte_iii/source/public_final/base_83/c30_body.tex body_lines=157-194

Soit \(\mathcal P\) le module libre engendré par les mots régionaux et
soit \(F^n\mathcal P\) sa filtration par longueur. La graduation associée
enregistre le premier niveau auquel apparaît chaque mot. Pour un paramètre
\(z\), le caractère ordinaire de longueur est
\[
\chi_z([w])=z^{|w|}.
\]
La multiplicativité
\[
\chi_z([uv])=\chi_z([u])\chi_z([v])
\]
exprime que la concaténation additionne les longueurs. Lors de l'évaluation en
\(z=\alpha\), un mot de longueur \(n\) reçoit le degré \(\alpha^n\).

\begin{definition}[Compatibilité de degré]
\label{p3-83-c30-s1:def:compatibilidad-grado}
Le lecteur analytique régional est compatible avec la filtration cylindrique
lorsqu'il envoie la composante homogène de longueur \(n\) sur le degré analytique
\(n\) :
\[
\operatorname{gr}_n\mathcal P\longrightarrow
\alpha^n\mathbb R.
\]
\end{definition}

Les cinq régions appartiennent à \(\mathcal U_6\) ; c'est pourquoi la donnée de
retour apparaît en degré six :
\[
\rho_\pi[U_6]\longmapsto169\alpha^6.
\]
Ce six est la longueur du mot régional. Ce n'est pas la longueur des
marches fermées du \cref{chap:cinco-ciclos-pi}. Ces marches parcourent à
la fois les quatre ancrages
\(U_\pi,U_e,U_\varphi,U_{\rm APP}\) et possèdent une longueur minimale de huit. Un
mot, une arête du graphe de blocs et une marche fermée sont des objets de
types distincts ; la graduation utilisée ici appartient au premier.

% HMT_SOURCE_SEGMENT_END id=B03
\subsection{Publication du degré analytique dans la carte globale de phase}
\label{sec:s102:c38:5:2}

% HMT_SOURCE_SEGMENT_BEGIN id=A03 source=colaboracion/parte_iii/source/public_final/ascensos_20260826/16f_realizacion_analitica_contraangulo.tex body_lines=172-209

Soit \(\mathcal P\) le module libre engendré par les mots régionaux et
soit \(F^n\mathcal P\) sa filtration par longueur. La graduation associée
enregistre le premier niveau auquel apparaît chaque mot. Pour un paramètre
\(z\), le caractère ordinaire de longueur est
\[
\chi_z([w])=z^{|w|}.
\]
La multiplicativité
\[
\chi_z([uv])=\chi_z([u])\chi_z([v])
\]
exprime que la concaténation additionne les longueurs. Lors de l'évaluation en
\(z=\alpha\), un mot de longueur \(n\) reçoit le degré \(\alpha^n\).

\begin{definition}[Compatibilité de degré]
\label{p3-20260826:def:compatibilidad-grado}
Le lecteur analytique régional est compatible avec la filtration cylindrique
lorsqu'il envoie la composante homogène de longueur \(n\) sur le degré analytique
\(n\) :
\[
\operatorname{gr}_n\mathcal P\longrightarrow
\alpha^n\mathbb R.
\]
\end{definition}

Les cinq régions appartiennent à \(\mathcal U_6\) ; c'est pourquoi la donnée de
retour apparaît en degré six :
\[
\rho_\pi[U_6]\longmapsto169\alpha^6.
\]
Ce six est la longueur du mot régional. Ce n'est pas la longueur des
marches fermées du \cref{chap:cinco-ciclos-pi}. Ces marches parcourent à
la fois les quatre ancrages
\(U_\pi,U_e,U_\varphi,U_{\rm APP}\) et possèdent une longueur minimale de huit. Un
mot, une arête du graphe de blocs et une marche fermée sont des objets de
types distincts ; la graduation utilisée ici appartient au premier.

% HMT_SOURCE_SEGMENT_END id=A03
% HMT_SOURCE_SEGMENT_BEGIN id=D02 source=manuscrito/sucesor_102/deltas_ley9/c38_par_angular_cartas_y_canales.tex body_lines=21-34
Le calendrier global présente les trois factorisations
\begin{equation}
 360=12\cdot30=4\cdot90=3\cdot120.
 \label{eq:l9s102:360-factorizaciones}
\end{equation}
La carte analytique s'obtient au moyen de
\begin{equation}
 A_{\mathrm{rad}}=\frac\pi{180}A_{\deg}
 =\frac{50\pi}{9}\alpha,
 \qquad
 C^*_{\mathrm{rad}}=\frac\pi{180}C^*_{\deg}
 =\frac\pi{90}(\eta_{\mathrm{ret}}+\alpha).
 \label{eq:l9s102:a-cstar-rad}
\end{equation}
% HMT_SOURCE_SEGMENT_END id=D02
% HMT_SOURCE_SEGMENT_BEGIN id=D03 source=manuscrito/sucesor_102/deltas_ley9/c38_par_angular_cartas_y_canales.tex body_lines=35-59

\begin{theorem}[Commutativité des factorisations régionale et d'action]
\label{thm:l9s102:conmutatividad-cstar}
Les deux chemins vers la phase conjuguée produisent la même sortie :
\begin{equation}
 \boxed{
 \frac\pi{180}C^*_{\deg}
 =\frac\pi{90}(\eta_{\mathrm{ret}}+\alpha)
 =\frac\pi{90}(H_5+\alpha)
 -\mathscr C_\pi(\alpha)=y_*.}
 \label{eq:l9s102:cuadrado-cstar}
\end{equation}
\end{theorem}

\begin{proof}
En substituant \cref{eq:l9s102:eta-ret} dans le second membre, on obtient
\[
 \frac\pi{90}
 \left(H_5-\frac{90}{\pi}\mathscr C_\pi+\alpha\right)
 =\frac\pi{90}(H_5+\alpha)-\mathscr C_\pi.
\]
La dernière expression est la définition régionale de \(y_*\), tandis que la
première égalité procède de \cref{eq:l9s102:a-cstar-deg}.
\end{proof}

% HMT_SOURCE_SEGMENT_END id=D03
\section{Décomposition spectrale \texorpdfstring{\(q_\pm\)}{q+/-} et calcul fonctionnel sur les canaux \texorpdfstring{\(90/120\)}{90/120} préalablement générés}
\label{sec:s102:c38:6}

\subsection{Reconstruction spectrale du couple angulaire et transition hexade--octade}
\label{sec:s102:c38:6:1}

% HMT_SOURCE_SEGMENT_BEGIN id=B11 source=colaboracion/parte_iii/source/public_final/base_83/c30_body.tex body_lines=523-569

La phase commune et la phase orientée déterminent les deux canaux
\[
q_-=e^{-x+y_*},
\qquad
q_+=e^{-x-y_*}.
\]
On retrouve exactement les invariants
\[
q_-q_+=D_A,
\qquad
\alpha=-\frac9{100\pi}\log D_A,
\qquad
C_*^{\rm HMT}=\frac{90}{\pi}\log\frac{q_-}{q_+}.
\]
La première égalité lit la contraction commune ; la troisième, la séparation
orientée des feuillets. L'échange \(R\mapsto-R\) permute
\(q_+\) et \(q_-\), conserve \(D_A\) et inverse le signe de \(y_*\).

Pour
\[
S_{90}(q)=\frac{q^{90}}{1-q^{270}},
\qquad
S_{120}(q)=\frac{q^{120}}{1-q^{360}},
\]
la fonctionnelle de transition hexade--octade s'écrit intégralement dans la
carte des phases :
\[
\Gamma_{6\to8}
=\frac{180}{\pi}xy_*
+12\bigl(S_{90}(q_-)-S_{90}(q_+)\bigr)
-\bigl(S_{120}(q_-)-S_{120}(q_+)\bigr).
\]
L'évaluation produite par la coordonnée \(C_*^{\rm HMT}\) du
\cref{p3-83-c30-s1:thm:contraangulo-regional} est
\[
\boxed{
\Gamma_{6\to8}
=0.274007672694459274747255260774\ldots.}
\]
Ainsi, la valeur interne reconnue en mécanique dimensionnelle comme fonctionnelle de
Barbero--Immirzi cesse de recevoir \(C_*^{\rm HMT}\) comme littéral indépendant :
toutes deux descendent du même couple de phases engendré dans ce chapitre.
L'interprétation géométrique et physique de \(\Gamma_{6\to8}\) appartient à l'œuvre
de mécanique dimensionnelle ; son évaluation mathématique appartient déjà à l'interface
précédemment construite.

% HMT_SOURCE_SEGMENT_END id=B11
\subsection{Transporteur relatif et disjonction spectrale des blocs}
\label{sec:s102:c38:6:2}

% HMT_SOURCE_SEGMENT_BEGIN id=B13 source=colaboracion/parte_iii/source/public_final/base_83/c30_body.tex body_lines=619-675
\label{p3-83-c30-s1:sec:transportador-relativo-split}

Dans la base ordonnée des canaux, soit
\[
 R=\begin{pmatrix}0&1\\1&0\end{pmatrix},
 \qquad P_\pm=\frac{I\pm R}{2},
 \qquad
 B_0=7I+2R,
 \qquad B_1=\Adeg I+\Cdeg R.
\]
La décomposition scindée
\[
 xI+yR=\rho(x,y)e^{\eta(x,y)R},
 \quad \rho(x,y)=\sqrt{x^2-y^2},
 \quad \eta(x,y)=\operatorname{artanh}(y/x)
\]
construit l'unique transporteur
\[
 T_{\rm rel}
 =\frac{\sqrt{(\Adeg)^2-(\Cdeg)^2}}{\sqrt{45}}
   \exp\!\left[
   \left(\operatorname{artanh}\frac{\Cdeg}{\Adeg}
   -\operatorname{artanh}\frac27\right)R\right],
 \qquad T_{\rm rel}B_0=B_1.
\]
Dans la base spectrale, il agit par
\(9\mapsto\Adeg+\Cdeg\) et
\(5\mapsto\Adeg-\Cdeg\).

\begin{theorem}[Disjonction spectrale des réalisations locale et angulaire]
\label{p3-83-c30-s1:thm:no-go-entrelazador-bloques-split}
Les spectres
\[
 \operatorname{Spec}(B_0)=\{9,5\},\qquad
 \operatorname{Spec}(B_1)=\{\Adeg+\Cdeg,\Adeg-\Cdeg\}
\]
sont disjoints. En conséquence,
\[
 \operatorname{Hom}_{B_0,B_1}
 :=\{\Phi:\Phi B_0=B_1\Phi\}=\{0\}.
\]
\end{theorem}

\begin{proof}
Les bornes rationnelles construites pour les coordonnées angulaires donnent
\(7.29<\Adeg<7.30\) et \(2.12<\Cdeg<2.13\) ; donc,
\(9.41<\Adeg+\Cdeg<9.43\) et
\(5.16<\Adeg-\Cdeg<5.18\). Dans la base qui diagonalise \(R\), chaque entrée
d'un entrelaceur est multipliée par la différence entre une valeur propre
de \(B_0\) et une de \(B_1\) ; la disjonction impose qu'elles soient toutes nulles.
\end{proof}

Le transporteur relatif compare deux blocs au sein de l'algèbre scindée ;
les entrelaceurs dodécaphasique--Witt et d'incidence agissent entre des
représentations équivalentes dans leurs espaces de résidence propres. Cette séparation
type positivement les deux opérations et conserve leurs domaines.

% HMT_SOURCE_SEGMENT_END id=B13
\subsection{Évaluation spectrale \texorpdfstring{\(V_{90,120}\)}{V_90,120} des canaux d'incidence préalablement générés par le TPK}
\label{sec:s102:c38:6:3}

% HMT_SOURCE_SEGMENT_BEGIN id=A11 source=colaboracion/parte_iii/source/public_final/ascensos_20260826/16f_realizacion_analitica_contraangulo.tex body_lines=554-572

La phase commune et la phase orientée déterminent les deux canaux
\[
q_-=e^{-\Arad+\Crad},
\qquad
q_+=e^{-\Arad-\Crad}.
\]
On retrouve exactement les invariants
\[
q_-q_+=D_A,
\qquad
\alpha=-\frac9{100\pi}\log D_A,
\qquad
\Cdeg=\frac{90}{\pi}\log\frac{q_-}{q_+}.
\]
La première égalité lit la contraction commune ; la troisième, la séparation
orientée des feuillets. L'échange \(R\mapsto-R\) permute
\(q_+\) et \(q_-\), conserve \(D_A\) et inverse le signe de \(\Crad\).

% HMT_SOURCE_SEGMENT_END id=A11
% HMT_SOURCE_SEGMENT_BEGIN id=D05 source=manuscrito/sucesor_102/deltas_ley9/c38_par_angular_cartas_y_canales.tex body_lines=68-97
\label{sec:l9s102:qpm-posteriores}

Soient \(H^2=I\) et \(P_\pm=(I\pm H)/2\). Le bloc angulaire
\begin{equation}
 B=A_{\mathrm{rad}}I+C^*_{\mathrm{rad}}H
 =(A_{\mathrm{rad}}+C^*_{\mathrm{rad}})P_+
 +(A_{\mathrm{rad}}-C^*_{\mathrm{rad}})P_-
 \label{eq:l9s102:b-angular}
\end{equation}
se publie multiplicativement par
\begin{equation}
 e^{-B}=q_+P_++q_-P_-,
 \qquad
 q_\pm=\exp[-(A_{\mathrm{rad}}\pm C^*_{\mathrm{rad}})].
 \label{eq:l9s102:qpm-angular}
\end{equation}
L'involution \(C^*\mapsto-C^*\) permute \(q_+\) et \(q_-\), et conserve
\begin{equation}
 q_+q_-=e^{-2A_{\mathrm{rad}}}=D_A.
 \label{eq:l9s102:det-angular}
\end{equation}
Les cardinaux \(90/120\), construits préalablement par l'incidence du
TPK, reçoivent maintenant leur évaluation spectrale :
\begin{equation}
 V_{90,120}
 =(I-T^{90})(I-T^{120})^{-1}
 =\sum_{\sigma\in\{+,-\}}
 \frac{1-q_\sigma^{90}}{1-q_\sigma^{120}}P_\sigma.
 \label{eq:l9s102:v90120-posterior}
\end{equation}
% HMT_SOURCE_SEGMENT_END id=D05
\subsubsection{Transporteur relatif scindé}
\label{sec:s102:c38:6:3:1}

% HMT_SOURCE_SEGMENT_BEGIN id=A12 source=colaboracion/parte_iii/source/public_final/ascensos_20260826/16f_realizacion_analitica_contraangulo.tex body_lines=574-748
\label{p3-20260826:sec:transportador-relativo-split}

La comparaison avec le bloc central d'APP peut s'effectuer sans identifier
des représentations distinctes. Fixons la base ordonnée des canaux et
l'involution
\[
 R=\begin{pmatrix}0&1\\1&0\end{pmatrix},
 \qquad R^2=I,
 \qquad
 P_\pm=\frac{I\pm R}{2}.
\]
Dans l'algèbre commutative
\[
 \mathscr B_R:=\{xI+yR:x,y\in\mathbb R\}
\]
considérons, dans la carte sexagésimale déjà déclarée,
\[
 B_0:=7I+2R,
 \qquad
 B_1:=\Adeg I+\Cdeg R.
\]
Le changement orthogonal de base
\[
 S=\frac1{\sqrt2}
 \begin{pmatrix}1&1\\1&-1\end{pmatrix}
\]
diagonalise simultanément l'involution et les deux blocs. Pour \(x>|y|\),
définissons
\[
 \rho(x,y):=\sqrt{x^2-y^2},
 \qquad
 \eta(x,y):=\operatorname{artanh}\!\left(\frac yx\right).
\]
Alors
\[
 xI+yR=\rho(x,y)e^{\eta(x,y)R}.
\]
Si
\[
 \rho_0=\sqrt{45},\quad
 \eta_0=\operatorname{artanh}(2/7),\quad
 \rho_1=\rho(\Adeg,\Cdeg),\quad
 \eta_1=\eta(\Adeg,\Cdeg),
\]
l'unique multiplication à gauche qui envoie \(B_0\) sur \(B_1\) est
\[
 \boxed{
 T_{\rm rel}
 :=\frac{\rho_1}{\sqrt{45}}
 e^{(\eta_1-\eta_0)R},
 \qquad
 T_{\rm rel}B_0=B_1.}
\]

\begin{proposition}[Transporteur relatif des deux blocs]
\label{p3-20260826:prop:transportador-relativo-split}
Une fois fixés \(R\), l'ordre des canaux, la branche positive et la carte
sexagésimale, \(T_{\rm rel}\) est l'unique élément de
\(\mathscr B_R\) qui satisfait \(TB_0=B_1\). Dans la base spectrale, il agit par
\[
 9\longmapsto\Adeg+\Cdeg,
 \qquad
 5\longmapsto\Adeg-\Cdeg.
\]
\end{proposition}

\begin{proof}
La décomposition polaire scindée donne la formule imprimée et prouve l'existence.
Comme \(\det B_0=45\ne0\), toute solution matricielle de \(TB_0=B_1\) satisfait
\(T=B_1B_0^{-1}\) ; elle est donc unique. La diagonalisation au moyen de \(S\)
produit les deux quotients spectraux
\((\Adeg+\Cdeg)/9\) et \((\Adeg-\Cdeg)/5\), équivalents à l'expression
exponentielle de \(T_{\rm rel}\).
\end{proof}

Ce transporteur n'est pas un entrelaceur représentationnel. En effet, un
entrelaceur \(\Phi\) entre les endomorphismes \(B_0\) et \(B_1\) devrait
satisfaire
\[
 \Phi B_0=B_1\Phi.
\]

\begin{theorem}[Obstruction spectrale à l'entrelacement]
\label{p3-20260826:thm:no-go-entrelazador-bloques-split}
Pour les valeurs de \(\Adeg\) et \(\Cdeg\) construites dans cette section,
\[
 \boxed{\Phi B_0=B_1\Phi\quad\Longrightarrow\quad\Phi=0.}
\]
\end{theorem}

\begin{proof}
Les spectres sont
\[
 \operatorname{Spec}(B_0)=\{9,5\},
 \qquad
 \operatorname{Spec}(B_1)
 =\{\Adeg+\Cdeg,\Adeg-\Cdeg\}.
\]
Les bornes rationnelles préalablement démontrées fournissent les intervalles
disjoints
\[
 7.29<\Adeg<7.30,
 \qquad
 2.12<\Cdeg<2.13,
\]
et, par conséquent,
\[
 9.41<\Adeg+\Cdeg<9.43,
 \qquad
 5.16<\Adeg-\Cdeg<5.18.
\]
Dans la base qui diagonalise \(R\), chaque entrée \(\phi_{ij}\) de \(\Phi\)
est multipliée par la différence entre une valeur propre de \(B_0\) et une de
\(B_1\). Toutes ces différences sont non nulles ; donc toutes les entrées de
\(\Phi\) s'annulent.
\end{proof}

Numériquement,
\[
 \rho_1=6.9814819335172378048\ldots,
 \qquad
 \frac{\rho_1}{\sqrt{45}}=1.0407378791354140779\ldots,
\]
\[
 \eta_1-\eta_0
 =0.005796351470571295590\ldots.
\]
La formule conserve l'algèbre scindée, les projecteurs \(P_\pm\) et la forme
lorentzienne à un facteur conforme \(\rho_1^2/45\) près. Elle dépend cependant
de \(B_1\) déjà construit : elle ne constitue pas une dérivation antérieure de
\(\Adeg\) ou \(\Cdeg\) et ne doit pas être confondue avec les entrelaceurs
équivariants entre réalisations d'un même module.

Pour
\[
S_{90}(q)=\frac{q^{90}}{1-q^{270}},
\qquad
S_{120}(q)=\frac{q^{120}}{1-q^{360}},
\]
la fonctionnelle de transition hexade--octade s'écrit intégralement dans la
carte des phases :
\[
\Gamma_{6\to8}
=\frac{180}{\pi}\Arad\Crad
+12\bigl(S_{90}(q_-)-S_{90}(q_+)\bigr)
-\bigl(S_{120}(q_-)-S_{120}(q_+)\bigr).
\]
L'évaluation produite par le contre-angle du
\cref{p3-20260826:thm:contraangulo-regional} est
\[
\boxed{
\Gamma_{6\to8}
=0.274007672694459274747255260774\ldots.}
\]
Les formes équivalentes de la fonctionnelle mathématique interne sont
\[
\boxed{
\Gamma_{6\to8}
=\Arad\Cdeg+K_{6\to8}
=\frac{180}{\pi}\Arad\Crad+K_{6\to8}
=\frac{\pi\Adeg\Cdeg}{180}+K_{6\to8}.}
\]
Ces identités évaluent \(\Gamma_{6\to8}\) après la génération de
\((\Crad,\Cdeg)\) ; elles ne sont pas utilisées en sens inverse pour sélectionner le
contre-angle. La lecture physique dans la connexion d'Ashtekar--Barbero est
représentée ensuite au moyen de
\[
 \Gamma_{6\to8}\dashrightarrow\gamma_{\rm BI}.
\]
La flèche n'introduit pas de troisième fonctionnelle et ne rétroagit pas sur \(C^*\). Les
séries \(S_{90}\), \(S_{120}\), les coefficients \(12\) et \(-1\) et leur
affectation au pas \(6\to8\) appartiennent à la construction mathématique
déclarée. L'identification physique exige en outre la normalisation propre
de la connexion.

% HMT_SOURCE_SEGMENT_END id=A12
\section{Naturalité sous raffinement et séparation d'avec la reconnaissance métrologique}
\label{sec:s102:c38:7}

\subsection{Estimation exacte du reste de la prolongation analytique}
\label{sec:s102:c38:7:1}

% HMT_SOURCE_SEGMENT_BEGIN id=B09 source=colaboracion/parte_iii/source/public_final/base_83/c30_body.tex body_lines=474-499

L'expression initialement trouvée,
\[
y_7=\frac\pi{90}(H_5+\alpha)-169\alpha^6-D_A\alpha^7,
\]
est la troncature de degré sept du caractère complet. Elle produit
\[
C_7=2.12373833896864600265403827103919\ldots{}^\circ.
\]
Le reste omis n'est pas un terme indéterminé ; c'est la queue géométrique
exacte
\[
\mathscr C_\pi(\alpha)
-169\alpha^6-D_A\alpha^7
=\frac{D_A^2\alpha^8}{1-D_A\alpha}.
\]
Par conséquent,
\[
\left|C_7-C_*^{\rm HMT}\right|
=\frac{180}{\pi}\frac{D_A^2\alpha^8}{1-D_A\alpha}
=2.7839208689064583\ldots\times10^{-16}\;{}^\circ.
\]
La récurrence transforme ainsi une approximation de septième degré en une
réalisation analytique complète et fournit une borne close pour chaque
troncature ultérieure.

% HMT_SOURCE_SEGMENT_END id=B09
\subsection{Dépendance fonctionnelle du sélecteur régional et du prolongement analytique}
\label{sec:s102:c38:7:2}

% HMT_SOURCE_SEGMENT_BEGIN id=B10 source=colaboracion/parte_iii/source/public_final/base_83/c30_body.tex body_lines=501-521

La construction change lorsque l'une quelconque de ses composantes actives est retirée.
Les valeurs suivantes s'obtiennent sans réétalonner les autres termes :
\[
\begin{array}{@{}lc@{}}
\toprule
\text{variante}&\text{coordonnée angulaire conjuguée en degrés}\\
\midrule
\rho_\pi=168&2.1237383389772976863904487\ldots\\
\rho_\pi=169&2.1237383389686457242619514\ldots\\
\rho_\pi=170&2.1237383389599937621334541\ldots\\
\text{sans queue déterminantale}&2.1237383389686949415301675\ldots\\
D_A\text{ remplacé par }1&2.1237383389686313409965826\ldots\\
\bottomrule
\end{array}
\]
Ces ablations ne constituent pas à elles seules une estimation probabiliste.
Elles démontrent que le sélecteur équivariant et le caractère déterminantal sont
des composantes effectives de la sortie et qu'ils ne peuvent être supprimés tout en maintenant la
même réalisation.

% HMT_SOURCE_SEGMENT_END id=B10
\subsection{Interprétation opératoire de la coordonnée angulaire conjuguée}
\label{sec:s102:c38:7:3}

% HMT_SOURCE_SEGMENT_BEGIN id=B14 source=colaboracion/parte_iii/source/public_final/base_83/c30_body.tex body_lines=677-715

La procédure complète peut être résumée par une chaîne d'applications :
\[
\begin{aligned}
\mathcal F_\pi
&\longmapsto b_\pi
\longmapsto\rho_\pi=169,\\
\alpha
&\longmapsto A
\longmapsto x
\longmapsto D_A,\\
(\rho_\pi,D_A)
&\longmapsto\mathscr C_\pi(\alpha),\\
(H_5,\mathscr C_\pi(\alpha))
&\longmapsto y_*
\longmapsto C_*^{\rm HMT},\\
(x,y_*)
&\longmapsto(q_+,q_-)
\longmapsto\Gamma_{6\to8}.
\end{aligned}
\]
La première ligne appartient au catalogue fini ; la deuxième, à la clôture commune de
\(\alpha\) ; la troisième, à la prolongation déterminantale ; la quatrième produit la
coordonnée orientée ; la cinquième la transporte vers l'interface dimensionnelle.

Dans cette organisation, \(\alpha\) mesure la contraction commune du transport et
\(C_*^{\rm HMT}\) mesure la séparation orientée de ses deux feuillets. Le retour
régional n'ajoute pas de constante externe : il sélectionne une amplitude finie et la
prolonge au moyen du déterminant déjà fixé par \(\alpha\). La carte en
degrés est ultérieure. L'objet primaire est le couple de phases \((x,y_*)\).

Une vérification exhaustive reconstruit les cinq régions directement
à partir du catalogue, parcourt les actions de
\(\mathfrak S_5\) et \(\mathfrak S_3\), vérifie la récurrence et somme la
série. L'entrée de \(\alpha\) est la sortie exacte de sa clôture
dodécaphasique, non un chiffre métrologique transcrit. La vérification maintient
séparés les théorèmes du catalogue, les règles du lecteur et la comparaison
externe ; c'est pourquoi aucune comparaison expérimentale n'intervient dans la fonction
génératrice.
% HMT_SOURCE_SEGMENT_END id=B14
\subsection{Certification de convergence de la troncature parangulaire}
\label{sec:s102:c38:7:4}

% HMT_SOURCE_SEGMENT_BEGIN id=A09 source=colaboracion/parte_iii/source/public_final/ascensos_20260826/16f_realizacion_analitica_contraangulo.tex body_lines=505-530

L'expression initialement trouvée,
\[
y_7=\frac\pi{90}(H_5+\alpha)-169\alpha^6-D_A\alpha^7,
\]
est la troncature de degré sept du caractère complet. Elle produit
\[
C_7=2.12373833896864600265403827103919\ldots{}^\circ.
\]
Le reste omis n'est pas un terme indéterminé ; c'est la queue géométrique
exacte
\[
\mathscr C_\pi(\alpha)
-169\alpha^6-D_A\alpha^7
=\frac{D_A^2\alpha^8}{1-D_A\alpha}.
\]
Par conséquent,
\[
\left|C_7-\Cdeg\right|
=\frac{180}{\pi}\frac{D_A^2\alpha^8}{1-D_A\alpha}
=2.7839208689064583\ldots\times10^{-16}\;{}^\circ.
\]
La récurrence transforme ainsi une approximation de septième degré en une
réalisation analytique complète et fournit une borne close pour chaque
troncature ultérieure.

% HMT_SOURCE_SEGMENT_END id=A09
\subsection{Naturalité du sélecteur régional sous raffinement et changement de carte}
\label{sec:s102:c38:7:5}

% HMT_SOURCE_SEGMENT_BEGIN id=A10 source=colaboracion/parte_iii/source/public_final/ascensos_20260826/16f_realizacion_analitica_contraangulo.tex body_lines=532-552

La construction change lorsque l'une quelconque de ses composantes actives est retirée.
Les valeurs suivantes s'obtiennent sans réétalonner les autres termes :
\[
\begin{array}{@{}lc@{}}
\toprule
\text{variante}&\text{contre-angle en degrés}\\
\midrule
\rho_\pi=168&2.1237383389772976863904487\ldots\\
\rho_\pi=169&2.1237383389686457242619514\ldots\\
\rho_\pi=170&2.1237383389599937621334541\ldots\\
\text{sans queue déterminantale}&2.1237383389686949415301675\ldots\\
D_A\text{ remplacé par }1&2.1237383389686313409965826\ldots\\
\bottomrule
\end{array}
\]
Ces ablations ne constituent pas à elles seules une estimation probabiliste.
Elles démontrent que le sélecteur équivariant et le caractère déterminantal sont
des composantes effectives de la sortie et qu'ils ne peuvent être supprimés tout en maintenant la
même réalisation.

% HMT_SOURCE_SEGMENT_END id=A10
\subsection{Indépendance métrologique de la construction parangulaire}
\label{sec:s102:c38:7:6}

% HMT_SOURCE_SEGMENT_BEGIN id=A14 source=colaboracion/parte_iii/source/public_final/ascensos_20260826/16f_realizacion_analitica_contraangulo.tex body_lines=799-831

La procédure complète peut être résumée par une chaîne d'applications :
\[
\begin{aligned}
\mathcal F_\pi
&\longmapsto b_\pi
\longmapsto\rho_\pi=169,\\
\alpha
&\longmapsto \Adeg
\longmapsto \Arad
\longmapsto D_A,\\
(\rho_\pi,D_A)
&\longmapsto\mathscr C_\pi(\alpha),\\
(H_5,\mathscr C_\pi(\alpha))
&\longmapsto \Crad
\longmapsto \Cdeg,\\
(\Arad,\Crad)
&\longmapsto(q_+,q_-)
\longmapsto\Gamma_{6\to8}.
\end{aligned}
\]
La première ligne appartient au catalogue fini ; la deuxième, à la clôture commune de
\(\alpha\) ; la troisième, à la prolongation déterminantale ; la quatrième produit la
coordonnée orientée ; la cinquième la transporte vers l'interface dimensionnelle.

Dans cette organisation, \(\alpha\) mesure la contraction commune du transport et
le contre-angle mesure la séparation orientée de ses deux feuillets. Dans
l'évaluation en vigueur, le retour régional sélectionne une amplitude finie et la
prolonge au moyen du déterminant déjà fixé par \(\alpha\), sans insérer de
constante externe supplémentaire. La carte en degrés est ultérieure. L'objet
premier est le couple de phases \((\Arad,\Crad)\). La récurrence, la somme
régionale et les ablations dépendent du catalogue et de
\(\alpha\), mais ne reçoivent pas de constante métrologique supplémentaire.
% HMT_SOURCE_SEGMENT_END id=A14

 
\chapter{Opérateur constitutif positif du vide et dérivation conjointe de \texorpdfstring{\(\varepsilon,\mu,Z,c\)}{epsilon, mu, Z, c}}
\begingroup
\renewcommand{\chapter}[2][]{}%
\chapter{Opérateur constitutif positif du vide : complétion \(3\to4\), reconstruction spectrale et dérivation conjointe de \(\varepsilon\), de \(\mu\), de \(Z\) et de \(c\)}
\label{chap:p3-final:39:vacio_constitutivo}
\section{Décomposition spectrale bicouche de l'opérateur angulaire}

Soit \(R=R^*=R^{-1}\) l'involution active et

\[
P_\pm=\frac{I\pm R}{2},
\qquad
T=q_+P_+ +q_-P_-,
\qquad 0<q_+<q_-<1.
\]

Avec
\[
x=A_{\rm rad}=\frac{\pi A}{180},\qquad
y=C^*_{\rm rad}=\frac{\pi C^*}{180},
\]
les deux canaux sont construits avant d'évaluer le vide :
\begin{equation}
q_+=e^{-x-y},\qquad q_-=e^{-x+y}.
\label{eq:vacuum-angular-channels}
\end{equation}
L'échange de feuillets inverse \(C^*\), c'est-à-dire \(y\mapsto-y\), et permute \(q_+\leftrightarrow q_-\). Le couple est ainsi la lecture spectrale d'un même antécédent angulaire, et non deux paramètres ajustés séparément.

\begin{definition}[Réponse constitutive]
\begin{equation}
\mathcal V_{90,120}=(I-T^{90})(I-T^{120})^{-1}.
\label{eq:vacuum-operator}
\end{equation}
\end{definition}

Comme \(\|T\|<1\), le dénominateur est positif et inversible. Le calcul fonctionnel donne

\begin{equation}
\mathcal V_{90,120}=r_+P_+ +r_-P_-,
\qquad
r_\pm=\frac{1-q_\pm^{90}}{1-q_\pm^{120}}.
\label{eq:vacuum-eigenvalues}
\end{equation}

L'affirmation d'inversibilité mérite d'être explicitée. Pour chaque feuillet,
\(0<q_\pm^{120}<1\) ; par conséquent
\[
I-T^{120}=(1-q_+^{120})P_++(1-q_-^{120})P_-
\]
est strictement positif et
\[
(I-T^{120})^{-1}
=\frac{P_+}{1-q_+^{120}}+
 \frac{P_-}{1-q_-^{120}}.
\]
La réponse \(\mathcal V_{90,120}\) ne s'obtient pas en résolvant quatre
équations scalaires séparées : elle est le résultat du calcul fonctionnel d'une
seule contraction positive. Cette observation empêche que perméabilité,
permittivité, impédance et propagation perdent l'antécédent qu'elles partagent.

\begin{theorem}[Détermination spectrale conjointe des quatre relations constitutives]
Les quatre caractères constitutifs
\begin{equation}
\widehat\varepsilon=r_+^2,\qquad
\widehat\mu=r_-^2,\qquad
\widehat Z=\frac{r_-}{r_+},\qquad
\widehat c=\frac1{r_+r_-}
\label{eq:vacuum-four}
\end{equation}
sont des fonctions spectrales de \(\mathcal V_{90,120}\). Ils satisfont
\begin{equation}
\widehat\mu\widehat\varepsilon=\widehat c^{-2},
\qquad
\frac{\widehat\mu}{\widehat\varepsilon}=\widehat Z^2.
\label{eq:vacuum-relations}
\end{equation}
\end{theorem}

\begin{proof}
Les identités s'obtiennent en substituant \cref{eq:vacuum-four}. Aucune carte dimensionnelle ni aucune valeur extérieure n'intervient.
\end{proof}

\begin{corollary}[Reconstruction des deux feuillets]
Les quatre caractères ne contiennent pas quatre degrés de liberté. N'importe lequel
des couples \((\widehat\mu,\widehat\varepsilon)\),
\((\widehat Z,\widehat c)\) ou \((\mathcal E,\mathcal B)\) reconstruit les
deux valeurs propres positives :
\begin{align}
r_-&=\sqrt{\widehat\mu}
=\sqrt{\widehat Z/\widehat c},
&r_+&=\sqrt{\widehat\varepsilon}
=\frac1{\sqrt{\widehat Z\widehat c}},
\label{eq:vacuum-reconstruct-r}\\
\log r_-&=\frac{\mathcal E+\mathcal B}{2},
&\log r_+&=\frac{\mathcal E-\mathcal B}{2}.
\label{eq:vacuum-reconstruct-log}
\end{align}
En particulier, l'orientation se perd si l'on conserve uniquement le produit,
mais se retrouve en ajoutant le quotient.
\end{corollary}

\begin{proof}
La première ligne est la racine positive des identités
\cref{eq:vacuum-four,eq:vacuum-relations}. La seconde inverse le système
linéaire
\(\mathcal E=\log r_-+\log r_+\),
\(\mathcal B=\log r_--\log r_+\).
\end{proof}

Ce corollaire exprime une propriété centrale de HMT : produit et quotient sont
des fonctionnelles complémentaires d'un même état. Le produit détermine la
propagation commune, tandis que le quotient conserve l'orientation
relative des feuillets intérieur et extérieur.

Dans le certificat V2,

\begin{align*}
r_+&=0.9999996285482493631786952281\ldots,\\
r_-&=0.9997241371895183641315165853\ldots,\\
\widehat\varepsilon&=0.9999992570966367027604416155\ldots,\\
\widehat\mu&=0.9994483504793269350899815559\ldots,\\
\widehat Z&=0.9997245085389372154555543466\ldots,\\
\widehat c&=1.0002763104861575261063862689\ldots.
\end{align*}

Ces numéraux sont la reconnaissance terminale de \cref{eq:vacuum-operator} ; ils ne définissent pas ses exposants \(90,120\).

\section{Factorisation harmonique des secteurs \(90\) et \(120\)}

Le semi-groupe harmonique est

\[
\langle90,120\rangle=30\langle3,4\rangle.
\]

Si \(s=q^{30}\), alors

\begin{equation}
\frac{1-q^{90}}{1-q^{120}}
=\frac{1-s^3}{1-s^4}
=\frac{1+s+s^2}{1+s+s^2+s^3}.
\label{eq:three-four-completion}
\end{equation}

L'identité sépare également la réponse de son défaut. Si
\[
F_{3\to4}(s)=\frac{1+s+s^2}{1+s+s^2+s^3},
\qquad 0<s<1,
\]
alors
\begin{equation}
d_{3\to4}(s):=1-F_{3\to4}(s)
=\frac{s^3}{1+s+s^2+s^3}>0.
\label{eq:vacuum-defect34}
\end{equation}
Le terme supplémentaire est conservé exactement comme défaut
positif. En outre,
\begin{equation}
F_{3\to4}'(s)
=-\frac{s^2(3+2s+s^2)}{(1+s+s^2+s^3)^2}<0.
\label{eq:vacuum-monotonic34}
\end{equation}
Par conséquent, l'ordre des canaux angulaires détermine univoquement
l'ordre des réponses constitutives. L'attribution de
\(\widehat\mu\) et \(\widehat\varepsilon\) aux feuillets respectifs s'obtient
analytiquement, sans recourir à leurs développements décimaux.

Chaque feuillet du vide constitue la complétion exacte de trois à quatre
termes d'un même mode élémentaire \(30\). En conséquence, perméabilité
et permittivité s'interprètent comme deux fonctionnelles quadratiques d'une même
complétion évaluée sur des feuillets conjugués, et non comme des paramètres
constitutifs indépendants.

Comme la fonction \((1-q^{90})/(1-q^{120})\) décroît dans \(0<q<1\) et \(q_+<q_-\), on obtient

\[
0<r_-<r_+<1,\qquad
0<\widehat\mu<\widehat\varepsilon<1,\qquad
0<\widehat Z<1<\widehat c.
\]

Les limites de la fonctionnelle déterminent sa géométrie. Lorsque \(q\to0^+\),
\(F(q)\to1\) : les quatre lectures s'approchent de l'unité et le défaut de
mémoire disparaît. Lorsque \(q\to1^-\),
\[
F(q)\longrightarrow\frac{90}{120}=\frac34.
\]
La réponse prend donc ses valeurs dans l'intervalle positif \((3/4,1)\). La
borne \(3/4\) n'est pas une constante ajustée : elle est le quotient des deux
longueurs de canal et la réalisation continue de la complétion
trois--à--quatre.

\section{Involution des feuillets et conservation de l'orientation}

L'échange de feuillets \(C^*\mapsto-C^*\) agit par

\begin{equation}
\widehat\mu\longleftrightarrow\widehat\varepsilon,\qquad
\widehat Z\longmapsto\widehat Z^{-1},
\qquad
\widehat c\longmapsto\widehat c.
\label{eq:vacuum-sheet}
\end{equation}

La propagation commune est paire ; l'impédance conserve l'orientation.
Cette distinction réapparaîtra dans Barbero et CKM : toutes les fonctionnelles
associées au changement de feuillet n'ont pas la même parité.

\section{Caractères logarithmiques et compatibilité traciale}

Définissons

\[
\mathcal E=\log(r_+r_-),\qquad
\mathcal B=\log(r_-/r_+).
\]

Avec la trace normalisée \(\tau=\Tr/2\),

\begin{equation}
2\tau(\log\mathcal V)=\mathcal E,
\qquad
-2\tau(R\log\mathcal V)=\mathcal B.
\label{eq:vacuum-traces}
\end{equation}

Le couple \((\mathcal E,\mathcal B)\) constitue un système de coordonnées linéaire du logarithme
opératoriel :
\begin{equation}
\log\mathcal V
=\frac{\mathcal E}{2}I-\frac{\mathcal B}{2}R.
\label{eq:vacuum-log-decomposition}
\end{equation}
La composante scalaire est invariante sous l'échange de feuillets ; la
composante orientée change de signe. Cette décomposition démontre
directement la parité de la propagation et le caractère orienté de
l'impédance.

L'amplification nonadique

\[
\mathcal V_m=\mathcal V\otimes I_{9^m}
\]

est compatible avec les espérances \(E_m=\mathrm{id}\otimes\tau_9\) :

\[
E_m(\mathcal V_{m+1})=\mathcal V_m.
\]

Par conséquent, \(\mathcal E\) et \(\mathcal B\) ne sont pas des particularités d'une matrice \(2\times2\) ; ils subsistent dans la tour UHF et dans sa clôture traciale \(II_1\).

Plus précisément, pour tout polynôme \(p\) puis, par continuité,
pour toute fonction continue sur le spectre,
\begin{equation}
E_m\bigl(p(\mathcal V_{m+1})\bigr)=p(\mathcal V_m).
\label{eq:vacuum-functional-refinement}
\end{equation}
La compatibilité ne préserve pas seulement deux nombres : elle préserve toute
l'algèbre fonctionnelle engendrée par la réponse. En particulier, elle conserve ses
projecteurs spectraux, ses puissances, son logarithme et les formes quadratiques
qui mesurent les défauts \cref{eq:vacuum-defect34}.

\section{Réalisation dimensionnelle des relations constitutives}

La mécanique dimensionnelle détermine

\begin{align}
Z_{\rm vac}&=\widehat Z\,u_Su_Q^{-2},\\
c_{\rm vac}&=\widehat c\,u_C,\\
\mu_{\rm vac}&=\widehat\mu\,u_Su_C^{-1}u_Q^{-2},\\
\varepsilon_{\rm vac}&=\widehat\varepsilon\,u_S^{-1}u_C^{-1}u_Q^2.
\label{eq:vacuum-sections}
\end{align}

Les identités constitutives se maintiennent dans les droites :

\begin{equation}
\mu_{\rm vac}\varepsilon_{\rm vac}=c_{\rm vac}^{-2},
\qquad
\mu_{\rm vac}/\varepsilon_{\rm vac}=Z_{\rm vac}^2.
\end{equation}

La carte SI est postérieure. Il n'est pas nécessaire de déclarer \(c\) primitif et de résoudre ensuite \(\mu\varepsilon=c^{-2}\) : les trois objets descendent conjointement de \(\mathcal V\) et des droites déclarées.

La séparation des niveaux s'exprime par le diagramme
\[
\begin{array}{c}
(A,C^*)\longrightarrow T\longrightarrow\mathcal V
\longrightarrow
(\widehat\mu,\widehat\varepsilon,\widehat Z,\widehat c)
\\[2mm]
\Big\downarrow\text{ cartes de droite}\hspace{36mm}
\Big\downarrow\text{ évaluation finale}
\\[2mm]
(u_S,u_Q,u_C)\longrightarrow
(\mu_{\rm vac},\varepsilon_{\rm vac},Z_{\rm vac},c_{\rm vac})
\longrightarrow\text{ coordonnées SI}.
\end{array}
\]
La première ligne est adimensionnelle et opératorielle ; la deuxième détermine les types de
grandeur ; la troisième sélectionne une carte métrologique. Modifier la carte finale
ne modifie ni les valeurs propres ni les identités constitutives. En ce
sens précis, l'unification mathématique--physique consiste à distinguer
l'invariant généalogique de son système de coordonnées, et non à supprimer la
structure dimensionnelle.

\begin{proposition}[Minimalité spectrale]
Soit \(W\) un opérateur positif à deux feuillets de spectre simple
\(\{r_+,r_-\}\). Si quatre caractères positifs satisfont
\(\mu\varepsilon=c^{-2}\) et \(\mu/\varepsilon=Z^2\), si la propagation
commune est déterminée par \(c^{-1}=\det W=r_-r_+\), et si son orientation
fixe \(Z=r_-/r_+\), alors nécessairement
\[
(\mu,\varepsilon,Z,c)
=\left(r_-^2,r_+^2,\frac{r_-}{r_+},\frac1{r_-r_+}\right).
\]
Par conséquent, \cref{eq:vacuum-four} n'est pas une paramétrisation parmi d'autres :
c'est la réalisation positive minimale compatible avec le produit, le quotient et
l'orientation.
\end{proposition}

\begin{proof}
De \(Z=r_-/r_+\) et \(c^{-1}=\det W=r_-r_+\), on obtient de manière unique les
racines positives. Les deux identités constitutives imposent ensuite
\(\mu=r_-^2\) et \(\varepsilon=r_+^2\).
\end{proof}

 La première inversion conceptuelle consiste à concevoir le vide comme une réponse relationnelle positive.\par

Dans HMT, le vide est la réponse dynamique minimale du système lorsque ses deux feuillets doivent coexister, se propager et clore leur différence en conservant l’orientation qui les distingue.\par

Les deux feuillets constituent deux lectures conjuguées du même antécédent angulaire. L’un conserve une orientation et l’autre conserve l’orientation opposée. C’est pourquoi apparaissent deux canaux, \(q_+\) et \(q_-\), comme les deux faces spectrales corrélées d’un même état.\par

L'opérateur\par

\[
\mathcal V_{90,120}
=
(I-T^{90})(I-T^{120})^{-1}
\]

compare deux profondeurs de propagation : \(90\) et \(120\). Ces profondeurs sont déterminées par le même mode élémentaire :\par

\[
90=3\cdot30,
\qquad
120=4\cdot30.
\]

C’est pourquoi la réponse de chaque feuillet peut se réécrire comme\par

\[
\frac{1-q^{90}}{1-q^{120}}
=
\frac{1+s+s^2}{1+s+s^2+s^3},
\qquad
s=q^{30}.
\]

Et c’est ici que l’interprétation devient réellement puissante. Le numérateur contient trois états du mode \(30\) ; le dénominateur contient ces mêmes trois états et un de plus. Le vide est donc une complétion exacte \(3\to4\).\par

Ce quatrième terme intègre et conserve les trois précédents. Il les complète. Ce qui est ajouté est enregistré comme un résidu positif de complétion :\par

\[
1-\frac{1+s+s^2}{1+s+s^2+s^3}
=
\frac{s^3}{1+s+s^2+s^3}>0.
\]

La mémoire du vide est précisément l’accroissement exact qui permet à trois termes de devenir quatre. Elle constitue l’empreinte positive et nécessaire de la complétion.\par

Des deux valeurs propres de cet unique opérateur découlent les quatre relations constitutives :\par

\[
\widehat\varepsilon=r_+^2,
\qquad
\widehat\mu=r_-^2,
\qquad
\widehat Z=\frac{r_-}{r_+},
\qquad
\widehat c=\frac1{r_+r_-}.
\]

Cela change complètement l’interprétation de la perméabilité et de la permittivité.\par

La permittivité est la réponse quadratique d’un feuillet. La perméabilité est la réponse quadratique du feuillet conjugué. L’impédance conserve la relation orientée entre les deux. La vitesse de propagation dépend de leur produit commun.\par

Ainsi, produit et quotient accomplissent deux tâches différentes :\par

\begin{itemize}[leftmargin=*,itemsep=0.35em]
\item le produit oublie quel était chaque feuillet et conserve la propagation commune ;
\item le quotient conserve quelle était leur orientation relative.
\end{itemize}

Le produit répond à la question : « que partagent les deux feuillets ? ». Le quotient répond à la question : « comment se distinguent-ils ? ».\par

Por eso\par

\[
\widehat\mu\,\widehat\varepsilon=\widehat c^{-2}
\]

y\par

\[
\frac{\widehat\mu}{\widehat\varepsilon}
=
\widehat Z^2
\]

sont les deux manières complémentaires de lire la même paire spectrale.\par

L’aspect holographique apparaît ici de façon très nette : quatre grandeurs visibles ne contiennent en réalité que deux degrés de liberté internes. Chacune de certaines paires adéquates permet de reconstruire les deux feuillets. La publication constitutive conserve la généalogie lorsque produit et quotient restent conjointement disponibles.\par

Telle est la réparation profonde de l’ablation de \(\mu\) et de \(\varepsilon\). La réparation réunit leurs valeurs et leur origine opératorielle commune.\par

La carte du Système international peut ensuite transformer ces sections intrinsèques en \(\mu_0\), \(\varepsilon_0\), \(Z_0\) et \(c\). Sa fonction consiste à exprimer, au sein d’une convention métrologique, une structure dont la provenance a déjà été sélectionnée par l’opérateur.\par

Le vide cesse ainsi d’être un réservoir de quatre constantes. Il devient une unique opération spectrale avec deux mémoires conjuguées.\par

Et cette vision a une conséquence philosophique importante : le vide est une relation de compatibilité entre deux manières orientées de réaliser un même état.\par

% BEGIN CR28_CONSTITUTIVE_INVERSE
\section{Récupération fonctionnelle, composition et moment de Catalan}
\label{cr28:sec:restitucion-constitutiva}

La réponse des incidences \(90/120\) conserve plus d'information qu'une valeur scalaire isolée. Dans la même notation que \eqref{eq:three-four-completion}, écrivons
\[
 f(s)=\frac{1+s+s^2}{1+s+s^2+s^3},\qquad
 \mathcal D=s\,\frac{d}{ds},\qquad
 \mathcal V=f(T^{30}).
\]
Les marques \(P_\pm\) maintiennent l'orientation des canaux. Le transport angulaire est l'antécédent de la réponse ; sa reconstruction fonctionnelle retrouve cette représentation sans remplacer l'archive de mémoire dont elle provient.

\begin{theorem}[Inversion de la réponse constitutive]
\label{cr28:thm:inversion-constitutiva} La fonction \(f\) est une bijection décroissante de \((0,1)\) sur \((3/4,1)\). Son inverse \(\psi\) associe à \(r\) l'unique racine \(s\in(0,1)\) de
\begin{equation}
 r s^3+(r-1)(1+s+s^2)=0.
 \label{cr28:eq:cubic}
\end{equation}
Les deux réponses étiquetées \(r_\pm\) déterminent
\begin{equation}
 s_\pm=\psi(r_\pm),\quad q_\pm=s_\pm^{1/30},\quad
 x=-\frac12\log(q_+q_-),\quad
 y=\frac12\log\frac{q_-}{q_+}.
 \label{cr28:eq:angular-recovery}
\end{equation}
Pour l'opérateur complet \(S=\mathcal V^2\), positif et de spectre dans \((9/16,1)\), le calcul fonctionnel fournit
\begin{equation}
 T=\bigl[\psi(S^{1/2})\bigr]^{1/30},\qquad
 \det T=\det\bigl[\psi(S^{1/2})\bigr]^{1/30}.
 \label{cr28:eq:operator-recovery}
\end{equation}
\end{theorem}
\begin{proof}
La dérivée négative de \eqref{eq:vacuum-monotonic34} et les limites \(f(0^+)=1,\ f(1^-)=3/4\) prouvent la bijection. Multiplier \(f(s)=r\) par son dénominateur positif donne \eqref{cr28:eq:cubic}. La racine positive d'ordre \(30\) retrouve \(q_\pm\) ; la somme et la différence de \(\log q_\pm=-x\mp y\) donnent les deux coordonnées angulaires. Par positivité, \(S^{1/2}=\mathcal V\). Appliquer \(\psi\) à chaque projecteur spectral retrouve \(T^{30}\) et la racine positive donne \(T\). Le calcul de son déterminant conclut.
\end{proof}

Cette récupération de l'opérateur complet ne s'étend pas à un déterminant isolé. En effet,
\[
 T_1=\operatorname{diag}(1/5,1/3),\qquad
 T_2=\operatorname{diag}(1/6,2/5)
\]
ont pour déterminant \(1/15\), mais possèdent des spectres distincts dans la même chambre contractive. L'injectivité démontrée implique \(f(T_1^{30})^2\ne f(T_2^{30})^2\). Les deux canaux et leurs marques préservent la récupération ; leur réduction à un produit perd de l'information.
% END CR28_CONSTITUTIVE_INVERSE

% BEGIN CR28_TRANSPORT_COMPOSITION
% Recuperación expositiva de INTERRELACIONES_ESTRUCTURALES_VACIO.md, §§6--8.
% Insertar después de la inversión cúbica y antes de las coordenadas logarítmicas.
\subsection{Composition transportée des réponses constitutives}
\label{cr28:comp-sec}

Le transport angulaire et les incidences $90$ et $120$ proviennent de la
même structure APP--TRIT--TPK : les feuillets conservent les évaluations et
leurs retenues ; le régime tritique conserve l’orientation ; le calendrier
TPK transporte la phase et la mémoire jusqu’aux canaux angulaires. La réponse
constitutive applique à ces canaux la fonction $f$ de
\eqref{eq:vacuum-monotonic34}. Son inversion permet d’exprimer, dans les coordonnées
de réponse, la composition des transports du même repère. On conserve
ainsi l’ordre constructif : transport, composition et évaluation spectrale.
L’opération obtenue agit sur cette représentation de l’état enrichi ;
son inverse retrouve les canaux, tandis que l’histoire complète de mémoire
demeure dans la structure de provenance.

Soient $s=q^{30}$ et $\psi=f^{-1}$. L’exposant $30$ est le plus grand commun
diviseur des incidences $90$ et $120$. L’inversion cubique détermine
$\psi:(3/4,1)\longrightarrow(0,1)$. En adjoignant la réponse de frontière
$f(1)=3/4$, nous étendons la bijection à
\[
 \psi:D\longrightarrow(0,1],\qquad D=[3/4,1).
\]
Sur cette frontière, on utilise la fonction rationnelle $f$ ; le quotient
$(1-q^{90})/(1-q^{120})$ y a une singularité éliminable.

\begin{proposition}[Monoïde des réponses et coordonnée additive]
\label{cr28:comp-monoid}
L’opération
\begin{equation}
 r\odot t=f\bigl(\psi(r)\psi(t)\bigr),\qquad r,t\in D,
 \label{cr28:comp-law}
\end{equation}
fait de $D$ un monoïde commutatif et simplifiable d’élément neutre $3/4$.
L’application
\begin{equation}
 \ell(r)=-\frac1{30}\log\psi(r)
 \label{cr28:comp-character}
\end{equation}
est un isomorphisme de ce monoïde avec $(\mathbb R_{\geq0},+)$.
L’intervalle $(3/4,1)$ est stable par $\odot$ ; le neutre appartient à
l’extension de frontière. Le seul élément inversible dans $D$
est le neutre.
\end{proposition}
\begin{proof}
Le produit de deux éléments de $(0,1]$ appartient au même intervalle et
\[
 \psi(r\odot t)=\psi(r)\psi(t).
\]
Pour trois réponses, les deux ordres d’association ont pour image
$\psi(r)\psi(t)\psi(u)$ ; l’injectivité de $\psi$ prouve
l’associativité. La commutativité provient du produit scalaire. Si
$r\odot t=r\odot u$, la positivité de $\psi(r)$ permet de simplifier et
d’obtenir $\psi(t)=\psi(u)$, donc $t=u$. Le paramètre $1$ correspond à
$3/4$ et fournit l’identité. Le logarithme transforme le produit en
somme, et son image dans $(0,1]$ prouve que $\ell$ est une bijection sur
$\mathbb R_{\geq0}$. En particulier, pour la réponse d’un canal $q$,
$\ell(r)=-\log q$. Si les deux paramètres sont strictement inférieurs à
$1$, leur produit l’est aussi. Enfin, deux nombres de $(0,1]$
ont pour produit $1$ uniquement s’ils sont tous deux $1$.
\end{proof}

\begin{proposition}[Simplification admissible et prolongement inverse]
\label{cr28:comp-inverse}
Pour $r,u\in D$, l’équation $r\odot t=u$ a une solution $t\in D$
exactement lorsque $u\geq r$. Cette solution est unique et donnée par
\begin{equation}
 t=f\!\left(\frac{\psi(u)}{\psi(r)}\right).
 \label{cr28:comp-cancellation}
\end{equation}
La même fonction $f$ est une bijection de $(0,\infty)$ sur $(0,1)$.
Avec cette extension de $\psi$, la loi~\eqref{cr28:comp-law} transporte le
groupe multiplicatif des paramètres positifs sur l’intervalle $(0,1)$.
L’inversion du groupe satisfait
\begin{equation}
 r^{\ominus}=\psi(r)r,\qquad
 r\odot r^{\ominus}=\frac34.
 \label{cr28:comp-group-inverse}
\end{equation}
\end{proposition}
\begin{proof}
L’équation équivaut à
$\psi(t)=\psi(u)/\psi(r)$. Ce quotient appartient à $(0,1]$
exactement lorsque $\psi(u)\leq\psi(r)$, ce qui, par monotonie décroissante,
équivaut à $u\geq r$. La bijection prouve l’existence et l’unicité.
La dérivée de \eqref{eq:vacuum-monotonic34} est négative pour tout $s>0$ ;
les limites de $f$ en $0$ et à l’infini sont $1$ et $0$, respectivement.
Cela établit la bijection étendue. En multipliant le numérateur et le
dénominateur par $s^3$, on obtient
\[
 f(s^{-1})=\frac{s^3+s^2+s}{s^3+s^2+s+1}=s f(s).
\]
L’inverse multiplicatif de $s=\psi(r)$ a donc pour réponse
$\psi(r)r$. Sa composition avec $r$ correspond au paramètre $1$.
\end{proof}

Le prolongement à $s>1$, ou de manière équivalente à $q>1$, est l’extension
algébrique du lecteur. Le secteur physique contractant du chapitre conserve
son domaine déclaré. La coordonnée $\ell$ se prolonge en un isomorphisme
avec $(\mathbb R,+)$ et change de signe sous l’inversion du groupe.

\subsubsection{Réalisation opératorielle dans un repère de feuillets commun}

\begin{proposition}[Composition avec projecteurs marqués]
\label{cr28:comp-operator}
Soient $\mathcal R=\mathcal R^*$, $\mathcal R^2=I$ et
$P_\pm=(I\pm\mathcal R)/2$. Pour $j=a,b$, soient
\[
 T_j=q_{j,+}P_++q_{j,-}P_-,\qquad
 \mathcal V_j=f(T_j^{30}),\qquad 0<q_{j,\pm}<1.
\]
La composition $T_{a\circ b}=T_aT_b$ satisfait
\begin{equation}
 \mathcal V_{a\circ b}
 =f\bigl(\psi(\mathcal V_a)\psi(\mathcal V_b)\bigr).
 \label{cr28:comp-operator-law}
\end{equation}
Ses deux réponses sont $r_{a,\pm}\odot r_{b,\pm}$. Si
$T_j=\exp[-(x_jI+y_j\mathcal R)]$ avec $x_j>|y_j|$, alors
\begin{equation}
 T_aT_b=
 \exp[-((x_a+x_b)I+(y_a+y_b)\mathcal R)],
 \label{cr28:comp-angular-addition}
\end{equation}
y $x_a+x_b>|y_a+y_b|$.
\end{proposition}
\begin{proof}
L’orthogonalité $P_+P_-=0$ donne
\[
 T_aT_b=q_{a,+}q_{b,+}P_++q_{a,-}q_{b,-}P_-.
\]
En particulier, $T_a$ et $T_b$ commutent. Le calcul fonctionnel dans ces
projecteurs fournit
$\psi(\mathcal V_j)=T_j^{30}$ et
$(T_aT_b)^{30}=T_a^{30}T_b^{30}$ ; appliquer $f$ prouve
\eqref{cr28:comp-operator-law}. Pour toute fonction $g$ définie sur les
deux valeurs propres, on a, plus explicitement,
\[
 P_\pm g(T_j)=g(q_{j,\pm})P_\pm.
\]
Comme $q_{j,\pm}=\exp[-(x_j\pm y_j)]$, leurs produits prouvent
\eqref{cr28:comp-angular-addition}. L’inégalité triangulaire conclut
$x_a+x_b>|y_a|+|y_b|\geq|y_a+y_b|$.
\end{proof}

Les coordonnées retrouvées peuvent s’écrire directement comme
\[
 x=\frac{\ell(r_+)+\ell(r_-)}2,\qquad
 y=\frac{\ell(r_+)-\ell(r_-)}2.
\]
La composition additionne ces paires. L’échange simultané de feuillets
permute les deux canaux de chaque transport, agit comme
$(x,y)\mapsto(x,-y)$ et est un automorphisme de la loi composante par
composante. L’inversion des deux transports dans l’extension algébrique
agit comme $(x,y)\mapsto(-x,-y)$. Les deux involutions commutent et
conservent des fonctions distinctes : l’une échange les marques d’orientation ;
l’autre inverse le transport.

Pour représenter l’échange par conjugaison, on utilise un
opérateur $L$ satisfaisant $L\mathcal R L^{-1}=-\mathcal R$. Dans la
représentation
\[
 \mathcal R=\begin{pmatrix}0&1\\1&0\end{pmatrix},\qquad
 L=\begin{pmatrix}1&0\\0&-1\end{pmatrix},
\]
la multiplication directe donne $LP_\pm L^{-1}=P_\mp$ et, par conséquent,
$LT_jL^{-1}=q_{j,-}P_++q_{j,+}P_-$. La conjugaison par
$\mathcal R$ conserve ses propres projecteurs. Une trace qui réunit les
canaux ne remplace pas non plus les projecteurs dans cette identité fonctionnelle :
l’information de leurs étiquettes intervient dans la composition.

L’hypothèse de repère commun est substantielle. Par exemple, les opérateurs
positifs définis
\[
 A_0=\begin{pmatrix}1/2&0\\0&1/3\end{pmatrix},\qquad
 B_0=\begin{pmatrix}1/2&1/6\\1/6&1/2\end{pmatrix}
\]
ont des valeurs propres positives, mais
\[
 A_0B_0=\begin{pmatrix}1/4&1/12\\1/18&1/6\end{pmatrix}
\]
n’est pas autoadjoint. La proposition spécifie la composition des
transports dans les projecteurs communs ; son domaine ne comprend pas un mélange
arbitraire de milieux matériels. La réalisation d’une concaténation de
chemins HMT conserve en outre ses données de mémoire : la fermeture algébrique du
cône de paramètres décrit la représentation et ne sélectionne pas à elle seule
de nouveaux chemins primaires.

\subsubsection{Produit de canaux et composition de transports}

La vitesse normalisée d’un état reste déterminée par
\[
 \widehat c=(r_+r_-)^{-1}.
\]
Ici, le produit ordinaire réunit les deux réponses de ce même état.
La loi $\odot$ évalue une autre opération : la composition de deux transports
dans chaque feuillet. Le calcul exact suivant distingue leurs arguments :
\begin{equation}
 f(1/2)=\frac{14}{15},\qquad
 \frac{14}{15}\odot\frac{14}{15}=f(1/4)=\frac{84}{85}
 \ne\frac{196}{225}=\left(\frac{14}{15}\right)^2.
 \label{cr28:comp-rational-example}
\end{equation}
Les fractions de cet exemple sont des paramètres algébriques de vérification.
L’expression de $\widehat c$ et les identités constitutives précédentes
restent inchangées. Le contenu ajouté est une loi explicite pour
transporter la composition et retrouver sa coordonnée angulaire additive.

% END CR28_TRANSPORT_COMPOSITION

% BEGIN CR28_CATALAN_RECOVERY
\subsection{Restitution fonctionnelle du moment orienté}
\label{cr28:sec:catalan-recovery}

La relation différentielle~\eqref{cr28:eq:catalan-problem} admet une inverse qui conserve le système complet des moments. Son domaine est la fonction constitutive produite par les secteurs du cycle, avec l'orientation de~\eqref{eq:vacuum-angular-channels}. Les incidences $90=3\cdot30$ et $120=4\cdot30$ fixent cette fonction avant son évaluation dans les canaux $s_\pm=q_\pm^{30}$. La condition à l'origine correspond à la disparition du moment lorsque son argument contractif s'annule. Cette condition détermine également les constantes d'intégration.

\begin{theorem}[Restitution intégrale et unicité]
\label{cr28:thm:catalan-recovery}
Soit $f(s)=(1-s^3)/(1-s^4)$, pour $0<s<1$, la réponse
de~\eqref{eq:three-four-completion}, et soit $\mathcal D=s\,d/ds$.
Il existe une unique fonction $F\in C^2((0,1))$ qui satisfait
\begin{equation}
 \mathcal D^2F(s)=\frac{s(1+s)}{1+s+s^2}f(s),
 \qquad \lim_{s\downarrow0}F(s)=0.
 \label{cr28:eq:catalan-problem}
\end{equation}
Elle est donnée par
\begin{align}
 F(s)&=\int_0^s\log\frac{s}{t}\,
          \frac{1+t}{1+t+t^2}f(t)\,dt
       =\int_0^s\frac{\log(s/t)}{1+t^2}\,dt
       =\mathcal C_2(s),
 \label{cr28:eq:catalan-integral}\\
 f(s)&=\frac{1+s+s^2}{s(1+s)}\mathcal D^2F(s).
 \label{cr28:eq:catalan-inverse}
\end{align}
L’extension continue à $s=1$ retrouve $F(1)=G_{\rm Cat}$.
\end{theorem}
\begin{proof}
La factorisation $1+t+t^2+t^3=(1+t)(1+t^2)$ donne
\[
 \frac{1+t}{1+t+t^2}f(t)=\frac1{1+t^2}.
\]
L’intégrande de~\eqref{cr28:eq:catalan-integral} est
intégrable en zéro, parce que
\[
 0\leq F(s)\leq\int_0^s\log(s/t)\,dt=s.
\]
Par conséquent, $F(s)\to0$. La différentiation sur l’intervalle ouvert,
d’abord sur un intervalle tronqué, puis par convergence dominée,
produit
\[
 \mathcal DF(s)=\int_0^s\frac{dt}{1+t^2},
 \qquad \mathcal D^2F(s)=\frac{s}{1+s^2}
     =\frac{s(1+s)}{1+s+s^2}f(s).
\]
En particulier, $F$ est de classe $C^2$ et satisfait l’équation.
Si $F_1,F_2$ sont deux solutions, la différence $H=F_1-F_2$
satisfait $\mathcal D^2H=0$. Au moyen de $z=\log s$, on obtient
$d^2H(e^z)/dz^2=0$, donc $H(s)=a\log s+b$.
L’existence de la limite nulle lorsque $s\downarrow0$ impose d’abord
$a=0$, puis $b=0$. La condition sur $F$ détermine donc
la solution et également la limite $\mathcal DF(s)\to0$.

Pour $s<1$, le développement géométrique de $(1+t^2)^{-1}$ et
l’intégrabilité absolue de ses termes pondérés permettent d’intégrer :
\[
 \int_0^s t^{2k}\log(s/t)\,dt
     =\frac{s^{2k+1}}{(2k+1)^2},\qquad
 F(s)=\sum_{k\geq0}\frac{(-1)^ks^{2k+1}}{(2k+1)^2}.
\]
La série représente le système intégral~\eqref{cr28:eq:catalan-integral}. La convergence absolue et uniforme de cette dernière série sur $[0,1]$ autorise le passage à l'extrémité $s=1$. Dans l'intégrale, la majoration intégrable $\log(1/t)$ l'autorise également, en prolongeant l'intégrande par zéro pour $t>s$. Enfin, isoler $f$ dans l'équation différentielle donne \eqref{cr28:eq:catalan-inverse}, avec des dénominateurs positifs sur $(0,1)$.
\end{proof}

La limite du moment et celle de sa dérivée seconde sont prises
dans leurs classes de convergence respectives. En particulier,
\[
 \lim_{s\uparrow1}\mathcal D^2\mathcal C_2(s)=\frac12
\]
est la limite radiale d’Abel de la série dérivée. L’évaluation
de $\mathcal C_2(1)$ utilise directement sa série absolument
convergente. Cette distinction conserve le même opérateur de profondeur
à l’intérieur et dans sa lecture de frontière.

\subsection{Compatibilité de la restitution avec les canaux et la vitesse}
\label{cr28:sec:channels-velocity}

Le théorème restitue une fonction à partir d'une autre fonction préalablement construite. L'inversion cubique~\eqref{cr28:eq:cubic} agit en revanche sur les évaluations $r_\pm$ : elle retrouve les arguments $s_\pm$ dans ce système fixé. La composition des deux opérations conserve l'orientation du quart de tour, les étiquettes $P_\pm$ et les incidences du transport. Ainsi, $G_{\rm Cat}=\mathcal C_2(1)$ est la lecture à l'extrémité d'un objet fonctionnel qui participe aussi aux réponses intérieures $r_\pm=f(s_\pm)$.

Pour $\mathcal V_m=\mathcal V\otimes I_{9^m}$ et
$\tau_m=\operatorname{Tr}/(2\cdot9^m)$, la lecture de vitesse
maintient exactement la normalisation de
\eqref{eq:vacuum-traces} :
\begin{equation}
 \exp[-2\tau_m(\log\mathcal V_m)]
   =\frac1{r_+r_-}
   =(\det\mathcal V)^{-1}=\widehat c.
 \label{cr28:eq:velocity-normalized}
\end{equation}
En effet, chaque valeur propre $r_\pm$ se répète $9^m$ fois, de sorte
que $2\tau_m(\log\mathcal V_m)=\log r_++\log r_-$.
L’identité utilise le déterminant de la réalisation à deux feuillets
et la trace normalisée de son amplification. Les sections
dimensionnelles déjà fixées donnent alors
\[
 c_{\rm int}=c_{\rm vac}=\widehat c\,u_C,
 \qquad \ell_0=c_{\rm vac}t_0,
\]
avec le même état, le même système de coordonnées et la même durée élémentaire que~\eqref{cr28:eq:same-velocity}. La restitution fonctionnelle conserve cette section de vitesse lorsqu'elle la compose avec l'action et la longueur ; la multiplicité spectrale est absorbée par la normalisation déclarée.

% END CR28_CATALAN_RECOVERY

% BEGIN CR28_CONSTITUTIVE_RADIAL_LINK
\subsection{Une même section constitutive dans l'action, la longueur et la gravité}
\label{cr28:sec:constitutive-radial} Les relations \eqref{eq:vacuum-sections} et la restitution précédente maintiennent une unique section de vitesse :
\begin{equation}
 c_{\rm int}=c_{\rm vac}=\widehat c\,u_C,\qquad
 \ell_0=c_{\rm int}t_0,\qquad
 c_{\rm int}=(\mu_{\rm vac}\varepsilon_{\rm vac})^{-1/2}.
 \label{cr28:eq:same-velocity}
\end{equation}
La dernière égalité est une identité de sections dans des bases dimensionnelles compatibles. La démonstration radiale du \cref{g26:sec:rigidez-radial-es} construit d'abord \(L=(5759/23040)\alpha^{16}L_*\), puis détermine
\begin{equation}
 G_{\rm ret}=\frac{c_{\rm int}^3L^2}{\hbar_{\rm ret}}
 =\frac{L^2}
 {\hbar_{\rm ret}(\mu_{\rm vac}\varepsilon_{\rm vac})^{3/2}}.
 \label{cr28:eq:constitutive-gravity}
\end{equation}
La seconde expression résulte de la substitution de l'identité constitutive dans la première ; elle n'introduit aucun nouvel étalonnage des deux canaux. En coordonnées, si \(c_{\rm int}=c_*\widehat c\), alors \(G_{\rm ret}=c_*^3\widehat c^{\,3}L^2/\hbar_{\rm ret}\). Le facteur de propagation n'intervient qu'une seule fois. La récupération de Catalan, le transport composé et la réponse radiale conservent ainsi la même origine angulaire, ses marques et sa normalisation.
% END CR28_CONSTITUTIVE_RADIAL_LINK



  \endgroup

\chapter{Charge, résistance de von Klitzing, conductance, constante de Josephson et flux magnétique}
\begingroup
\renewcommand{\chapter}[2][]{}%
\chapter{Dérivation de la charge et des quanta électriques : résistance de von Klitzing, conductance, constante de Josephson et flux magnétique}
\label{chap:p3-final:40:cuantos_electricos}
\label{ch:metrology}

\section{Dérivation de la section de charge à partir du vide}

\(Z_{\rm vac}\), \(h_a=2\pi\hbar_a\) et \(\alpha\) étant fixés, la charge du feuillet \(a\in\{\mathrm{pre},\mathrm{ret}\}\) est

\begin{equation}
e_a=\sqrt{\frac{2\alphah_a}{Z_{\rm vac}}}.
\label{eq:charge}
\end{equation}

La section \(e_a\) n'est pas introduite au moyen d'une coordonnée du Système
international destinée à reconstruire \(\alpha\). L'ordre causal procède de
\(\alpha\), de l'action et du vide vers la section de charge.

\section{Constantes de von Klitzing et de Josephson, conductance et flux magnétique}

De \cref{eq:charge} se déduisent

\begin{align}
R_K&=\frac{h_a}{e_a^2}=\frac{Z_{\rm vac}}{2\alpha},
\label{eq:rk}\\
G_0&=\frac{2e_a^2}{h_a}=\frac{4\alpha}{Z_{\rm vac}},
\label{eq:g0}\\
\Phi_{0,a}&=\frac{h_a}{2e_a}
=\sqrt{\frac{h_aZ_{\rm vac}}{8\alpha}},
\label{eq:phi0}\\
K_{J,a}&=\frac{2e_a}{h_a}=\Phi_{0,a}^{-1}.
\label{eq:kj}
\end{align}

En conséquence,

\begin{equation}
R_KG_0=2,\qquad G_0Z_{\rm vac}=4\alpha,
\qquad K_{J,a}\Phi_{0,a}=1.
\label{eq:quantum-electrical-identities}
\end{equation}

La covariance bidirectionnelle prend la forme

\[
R_K^{\rm pre}=R_K^{\rm ret},\qquad
G_0^{\rm pre}=G_0^{\rm ret},
\]

\[
\frac{\Phi_{0,\rm pre}}{\Phi_{0,\rm ret}}=R_{\rm act}^{1/2},
\qquad
\frac{K_{J,\rm pre}}{K_{J,\rm ret}}=R_{\rm act}^{-1/2}.
\]

Les relations de résistance et de conductance sont indépendantes de
l'orientation de l'action. Le couple flux--fréquence conserve cette
orientation et transforme réciproquement ses deux composantes.

  \endgroup

\chapter{Constantes thermiques et rayonnement : Boltzmann, Wien et Stefan--Boltzmann}
\begingroup
\renewcommand{\chapter}[2][]{}%
\chapter{Constantes thermiques et rayonnement : Boltzmann, Wien, Stefan--Boltzmann et thermodynamique spectrale du gaz photonique}
\label{chap:p3-final:48:constantes_termicas}
\section{Relation thermique intrinsèque de l'échelle \(108\)}

La structure chronologique \(108\) fixe

\begin{equation}
k_B\Theta_{\rm clk}\log3
=\frac{h_{\rm int}}{108t_0}
=\frac{2\pi\hbar_{\rm int}}{108t_0}.
\label{eq:boltzmann-clock}
\end{equation}

Cette égalité ne détermine pas séparément une coordonnée numérique de \(k_B\) et une autre de \(\Theta_{\rm clk}\) sans normalisation supplémentaire. Elle détermine bien leur produit comme section d'énergie. Le facteur \(\log3\) a trois antécédents compatibles :

\begin{equation}
\log3=C_1+C_2-2C_0
=\frac{S_{\TRIT}}{k_B}
=\frac{E_P}{k_B\Theta_{\rm clk}}.
\label{eq:log3-triple}
\end{equation}

Les trois égalités correspondent aux réalisations résiduelle, informationnelle et
thermodynamique d'un même scalaire. La coïncidence de valeur n'identifie pas
leurs domaines respectifs.

\section{Formulations spectrales de la loi de déplacement de Wien}

Pour la densité spectrale par longueur d'onde, le maximum satisfait

\[
5(1-\ee^{-x_5})=x_5,
\qquad
x_5=5+W_0(-5\ee^{-5}).
\]

Avec \cref{eq:boltzmann-clock},

\begin{equation}
\lambda_{\max}(\Theta_{\rm clk})
=\frac{108\log3}{x_5}\,\ell_0.
\label{eq:wien-lambda}
\end{equation}

Pour la densité par fréquence,

\[
3(1-\ee^{-x_3})=x_3,
\qquad
x_3=3+W_0(-3\ee^{-3}),
\]

et

\begin{equation}
\nu_{\max}(\Theta_{\rm clk})
=\frac{x_3}{108t_0\log3}.
\label{eq:wien-frequency}
\end{equation}

Les maxima ne sont pas réciproques, puisque les densités spectrales se
transforment avec le jacobien correspondant. Le quotient \(x_5/x_3\)
appartient à la structure de la paramétrisation et ne représente pas une
incohérence.

\section{Loi de Stefan--Boltzmann et thermodynamique du gaz de photons}

Évaluer la loi de rayonnement en \(\Theta_{\rm clk}\) produit le flux intrinsèque

\begin{equation}
j_\star(\Theta_{\rm clk})
=\frac{2\pi^5}{15\,108^4(\log3)^4}\,
\frac{h_{\rm int}}{\ell_0^2t_0^2},
\label{eq:stefan-hmt}
\end{equation}

dans la normalisation déclarée. La constante de Stefan--Boltzmann n'est pas insérée comme primitive : elle résulte de l'intégration du spectre bosonique une fois l'action, la propagation et la température fixées.
Les lois spectrales employées sont celles de Planck et de Wien \cite{Planck,Wien} ; HMT fixe ici la section sur laquelle elles sont évaluées, et non leurs constantes par comparaison rétrospective.

En unités de l'échelle de Planck,

\begin{align}
n_\gamma\ell_P^3
&=\frac{2\zeta(3)}{\pi^2\log^3 3},
\label{eq:photon-number}\\
\frac{u_\gamma\ell_P^3}{E_P}
&=\frac{\pi^2}{15\log^4 3},
\label{eq:photon-energy}\\
\frac{s_\gamma\ell_P^3}{k_B}
&=\frac{4\pi^2}{45\log^3 3}.
\label{eq:photon-entropy}
\end{align}

\(\zeta(3)\) représente le moment bosonique d'ordre trois, tandis que
\(\log3\) détermine l'énergie de décision TRIT. L'équation conserve la
distinction entre les deux invariants.

La conductance thermique quantique s'exprime comme

\begin{equation}
\frac{\kappa_Qt_0}{k_B}=\frac{\pi^2}{324\log3}.
\label{eq:thermal-conductance}
\end{equation}

\section{Confluence métrologique des invariants thermiques et radiatifs}

\begin{longtable}{@{}p{0.18\textwidth}p{0.27\textwidth}p{0.45\textwidth}@{}}
\toprule
Objet & Équation interne & Significado estructural\\
\midrule
\endhead
\(R_K\) & \(Z_{\rm vac}/(2\alpha)\) & résistance quantique invariante bidirectionnelle\\
\(G_0\) & \(4\alpha/Z_{\rm vac}\) & conductancia complementaria\\
\(\Phi_0\) & \(\sqrt{hZ/(8\alpha)}\) & flux qui conserve l'orientation de l'action\\
\(K_J\) & \(\Phi_0^{-1}\) & relation réciproque entre fréquence et flux\\
\(E_P\) & \(k_B\Theta\log3\) & énergie du cycle et contenu informationnel TRIT\\
Wien & \(108\log3/x_5\) & extremum spectral à l'échelle chronologique\\
Apéry photonique & \(2\zeta(3)/(\pi^2\log^3 3)\) & densité de nombre bosonique\\
\bottomrule
\end{longtable}

\begin{remark}[Contrôle négatif]
Le bloc est réfuté si un changement de variable spectrale conserve
indûment le même extremum de Wien, si \(\zeta(3)\) remplace
\(\log3\), ou si des valeurs séparées sont attribuées à \(k_B\) et \(\Theta\) sans
spécifier de carte. Les identités
\cref{eq:quantum-electrical-identities,eq:planck-trit} fournissent des contrôles
algébriques directs.
\end{remark}

\subsection*{Portée logique et domaine de validité}
Quanta électriques et confluence Planck--TRIT : \emph{Théorème interne exact}. Wien, Stefan--Boltzmann et gaz photonique : \emph{Composition mathématique exacte} avec le continu bosonique déclaré. La carte SI reste une \emph{Comparaison métrologique externe}.

  \endgroup

\chapter{Constantes électrofaibles, relation \texorpdfstring{\(W/Z\)}{W/Z} et mélange angulaire}
\begingroup
\renewcommand{\chapter}[2][]{}%
\chapter{Constantes électrofaibles et mélange : déterminant angulaire, relation W/Z, angle de Weinberg et réflexion rationnelle CKM}
\label{chap:p3-final:49:sector_electrodebil}
\section{Déterminant angulaire et relation électrofaible de Weinberg}

Le déterminant de \(T\) est

\begin{equation}
\det T=q_+q_-=D_A.
\label{eq:detT}
\end{equation}

Le secteur constitutif du vide évalue \(F(T)^2\), où
\(F(z)=(1-z^{90})/(1-z^{120})\), tandis que le secteur angulaire
électrofaible évalue \(\det T\). Il s'agit de deux fonctions d'un
antécédent commun, et non d'une factorisation horizontale entre constantes.
Cette distinction est développée dans le \cref{ch:ew}.

\begin{remark}[Contrôle négatif]
Le bloc est réfuté si l'une quelconque de ces identités exactes est violée :
positivité de \(I-T^{120}\), égalité
\cref{eq:three-four-completion}, relations \cref{eq:vacuum-relations} ou
compatibilité \(E_m(\mathcal V_{m+1})=\mathcal V_m\). Une coïncidence SI
ne peut pas compenser une défaillance opératorielle.
\end{remark}

\subsection*{Portée logique et domaine de validité}
Opérateur, spectre, complétion \(3\to4\), parité, amplification et sections constitutives : \emph{Théorème interne exact}. La formulation qui les réunit en un seul opérateur est une \emph{Formalisation nouvelle} ; les données angulaires et les droites sont un \emph{Résultat antécédent démontré} du développement HMT précédent et sont déclarées ici comme structure de départ explicite.

 \label{ch:ew}

\section{Fonctions constitutive et déterminantale de l'opérateur angulaire}

Rappelons l'opérateur à deux feuillets du \cref{ch:vacuum},
\begin{equation}
\label{eq:ew-angular-operator}
T=q_+P_+ +q_-P_-,
\qquad
q_-=\ee^{-(A_{\rm rad}-C^*_{\rm rad})},
\qquad
q_+=\ee^{-(A_{\rm rad}+C^*_{\rm rad})}.
\end{equation}
Le vide électromagnétique applique la fonction
\begin{equation}
\label{eq:ew-vacuum-reader}
F(T)^2,
\qquad
F(z)=\frac{1-z^{90}}{1-z^{120}},
\end{equation}
tandis que le secteur angulaire électrofaible applique le déterminant
\begin{equation}
\label{eq:ew-determinant-reader}
\boxed{
\det T=q_+q_-=\ee^{-2A_{\rm rad}}=:D_A.}
\end{equation}
Ce sont deux réalisations fonctionnelles de l'antécédent commun \((A,C^*,T)\). La génération part de \(T\) ; retrouver cette représentation à partir de l'opérateur constitutif complet est en outre possible dans le domaine contractif. Par \cref{cr28:thm:inversion-constitutiva},
\[
 T=\bigl[\psi((F(T)^2)^{1/2})\bigr]^{1/30},
\]
et, par conséquent, \(\det T\) est également retrouvé. Le déterminant isolé ne reconstruit pas les deux canaux : le contre-exemple positif du même théorème conserve leur produit et modifie leurs réponses. L'unité structurelle réside dans \(T\), ses projecteurs et la mémoire de provenance ; la récupération fonctionnelle conserve cet ordre et n'identifie pas entre elles les différentes lectures de perméabilité, de permittivité et de \(\theta_W\).

Dans l'involution des feuillets \(C^*\mapsto-C^*\), les valeurs propres
\(q_+,q_-\) s'échangent. C'est pourquoi \(D_A\) et la vitesse constitutive
sont paires, tandis que \((\widehat\mu,\widehat\varepsilon)\) s'échange et
\(\widehat Z\) s'inverse. La relation électrofaible qui suit utilise
l'invariant pair.

\section{Relation exacte entre les secteurs \(W\) et \(Z\)}

La coordonnée angulaire de durée--perduration est
\begin{equation}
\label{eq:ew-angular-character}
\chi_{\rm dur}(n_A,s_C)
=\exp\!\left(n_AA_{\rm rad}
+\frac{s_C}{6}C^*_{\rm rad}\right).
\end{equation}
Cette coordonnée n'est pas introduite comme une étiquette isolée. Sa provenance est
\begin{equation}
\label{eq:ew-signature-provenance}
\begin{aligned}
\APP&\longrightarrow\TRIT\longrightarrow\TPK
\longrightarrow\text{extension survivante}
\longrightarrow\text{route}\\
&\longrightarrow\text{registre}
\longrightarrow(n_A,s_C,\ldots)
\longrightarrow\chi_{\rm dur}.
\end{aligned}
\end{equation}
Les signatures de \(W\), \(Z\) et \(H\) sont des grandeurs typées obtenues au moyen de
cette chaîne ; le présent chapitre les adopte sans les sélectionner à partir de
leurs masses terminales. Dans la réalisation angulaire établie, les secteurs
chargé et neutre satisfont
\begin{equation}
\label{eq:ew-WZ-signatures}
(n_A,s_C)_W=(94,-1),
\qquad
(n_A,s_C)_Z=(95,-1).
\end{equation}
Cette affirmation concerne la composante angulaire commune. Les
coordonnées supplémentaires de raffinement de la loi générale des masses ne sont pas
éliminées à moins que le transducteur physique ne démontre leur annulation.

\begin{theorem}[Défaut de Weinberg dans le schéma des masses physiques de la carte angulaire]
\label{thm:ew-weinberg}
Si \(W\) et \(Z\) sont évalués dans une même section d'échelle et partagent
le raffinement qui accompagne \cref{eq:ew-WZ-signatures}, alors
\begin{equation}
\label{eq:ew-WZ-ratio}
\boxed{
\frac{m_W^{\angle}}{m_Z^{\angle}}
=\ee^{-A_{\rm rad}}=\sqrt{D_A},
\qquad
\left(\frac{m_W^{\angle}}{m_Z^{\angle}}\right)^2=D_A.}
\end{equation}
Dans le schéma des masses physiques,
\begin{equation}
\label{eq:ew-sin2thetaW}
\boxed{
\sin^2\theta_W^{\rm OS,HMT}=1-D_A
=0.2248708807528435698111159035\ldots .}
\end{equation}
De plus,
\begin{equation}
\label{eq:ew-WZ-terminal-ratio}
\frac{m_W^{\angle}}{m_Z^{\angle}}
=0.8804141748331613670128885459\ldots .
\end{equation}
\end{theorem}

\begin{proof}
Par \cref{eq:ew-angular-character,eq:ew-WZ-signatures},
\[
\frac{\chi_{\rm dur}(94,-1)}{\chi_{\rm dur}(95,-1)}
=\exp(-A_{\rm rad});
\]
la coordonnée \(C^*\) s'annule parce que les deux secteurs partagent
\(s_C=-1\). En élevant au carré et en utilisant
\cref{eq:ew-determinant-reader}, on obtient \(D_A\). La définition
des masses physiques \(s_W^2=1-m_W^2/m_Z^2\) produit
\cref{eq:ew-sin2thetaW}.
\end{proof}

Le contenu principal est l'identité entre deux constructions
internes : le déterminant de l'opérateur angulaire et la différence d'un pas
dans la réalisation \(W/Z\). La promotion en sélecteur physique de jauge exige
un entrelaceur qui transporte les feuillets de \(T\) vers la base neutre
\((W^3,B)\) et qui démontre quels raffinements sont conservés ou annulés.

\section{Couplages de jauge en normalisation rationalisée}

Nous définissons le couplage adimensionnel de jauge
\begin{equation}
\label{eq:ew-egauge}
e_g:=\sqrt{4\pi\alpha}
=0.3028221208717526427662442689\ldots .
\end{equation}
Il ne faut pas confondre \(e_g\) avec une section dimensionnelle de charge
\(e_{\rm int}\) : la première est un couplage en unités naturelles ; la
seconde vit dans la droite de charge de la mécanique dimensionnelle.

Soit
\begin{equation}
\label{eq:ew-sw-cw}
s_W^2=1-D_A,
\qquad
c_W^2=D_A.
\end{equation}
La convention électrofaible rationalisée au niveau de l'arbre donne
\begin{align}
g&=\frac{e_g}{s_W}
=2\sqrt{\frac{\pi\alpha}{1-D_A}}
=0.6385883427334574283647003809\ldots,
\label{eq:ew-g}\\
g'&=\frac{e_g}{c_W}
=2\sqrt{\frac{\pi\alpha}{D_A}}
=0.3439541633108502429362319257\ldots,
\label{eq:ew-gprime}\\
\frac{g'}{g}
&=\sqrt{D_A^{-1}-1}
=0.5386164141966093594996610933\ldots .
\label{eq:ew-gprime-over-g}
\end{align}

\begin{proposition}[Relation custodiale à l'arbre]
\label{prop:ew-rho}
Dans la même interface,
\begin{equation}
\label{eq:ew-rho}
\rho_{\rm EW}^{(0)}
:=\frac{(m_W^{\angle})^2}
{(m_Z^{\angle})^2c_W^2}=1.
\end{equation}
\end{proposition}

\begin{proof}
Par \cref{eq:ew-WZ-ratio,eq:ew-sw-cw}, le numérateur relatif comme
\(c_W^2\) sont \(D_A\).
\end{proof}

Les équations \cref{eq:ew-g,eq:ew-gprime} sont des compositions exactes une
fois déclarés le schéma des masses physiques, la rationalisation et l'échelle. Elles
n'affirment pas qu'un changement de schéma de renormalisation laisse leurs
valeurs numériques inchangées.

\section{Constante de Fermi dans l'approximation à l'arbre}

En unités naturelles, l'intégration du propagateur chargé à basse
énergie produit
\begin{equation}
\label{eq:ew-GF-def}
\frac{G_F^{(0)}}{\sqrt2}=\frac{g^2}{8m_W^2}.
\end{equation}
En substituant \cref{eq:ew-g}, on obtient la composition
\begin{equation}
\label{eq:ew-GF}
\boxed{
G_F^{(0)}
=\frac{\pi\alpha}
{\sqrt2(1-D_A)m_W^2}.}
\end{equation}
En adoptant la masse obtenue précédemment
\begin{equation}
\label{eq:ew-W-output}
m_W^{\HMT}=80.381592550221\ldots\ \mathrm{GeV},
\end{equation}
l'évaluation terminale est
\begin{equation}
\label{eq:ew-GF-value}
\boxed{
G_F^{(0)}
=1.1157162817659203186621784\times10^{-5}\ \mathrm{GeV}^{-2}.}
\end{equation}

La différence par rapport à une constante de Fermi renormalisée n'est pas utilisée
pour choisir un coefficient. Dans la convention
\begin{equation}
\label{eq:ew-Delta-r}
G_F
=\frac{\pi\alpha}
{\sqrt2(1-D_A)m_W^2(1-\Delta r)},
\end{equation}
l'obligation restante est de dériver \(\Delta r\) à partir du secteur interne des
boucles, des états et de l'échelle. Consulter la valeur observée de \(G_F\) pour définir
\(\Delta r\) serait une reconstruction calibrée, et non une prédiction.

\section{Dépendance orientée du couplage de Higgs}

Dans la carte angulaire, les routes pertinentes sont
\begin{equation}
\label{eq:ew-HW-signatures}
(n_A,s_C)_H=(97,9),
\qquad
(n_A,s_C)_W=(94,-1).
\end{equation}
Par conséquent,
\begin{equation}
\label{eq:ew-HW-ratio}
\boxed{
\frac{m_H^{\angle}}{m_W^{\angle}}
=\frac{\chi_{\rm dur}(97,9)}{\chi_{\rm dur}(94,-1)}
=\exp\!\left(3A_{\rm rad}+\frac53C^*_{\rm rad}\right)
=1.55872087212325548564\ldots .}
\end{equation}
Ici, \(C^*\) ne s'annule pas : le secteur scalaire conserve l'orientation
que le déterminant \(D_A=\ee^{-2A_{\rm rad}}\) ne distingue pas.

Prenons la convention du potentiel
\begin{equation}
\label{eq:ew-Higgs-potential}
V(H)=-\mu_H^2H^\dagger H
+\lambda_H(H^\dagger H)^2,
\end{equation}
pour laquelle
\begin{equation}
\label{eq:ew-Higgs-tree-relations}
m_H^2=2\lambda_Hv^2,
\qquad
m_W^2=\frac{g^2v^2}{4}.
\end{equation}

\begin{corollary}[Autocouplage orienté de Higgs]
\label{cor:ew-Higgs-lambda}
Dans l'interface à l'arbre déclarée,
\begin{equation}
\label{eq:ew-Higgs-lambda}
\boxed{
\lambda_H
=\frac{\pi\alpha}{2(1-D_A)}
\exp\!\left(6A_{\rm rad}+\frac{10}{3}C^*_{\rm rad}\right)
=0.1238479115482466804662\ldots .}
\end{equation}
De plus,
\begin{equation}
\label{eq:ew-GF-Higgs}
G_F^{(0)}m_H^2=\sqrt2\,\lambda_H.
\end{equation}
\end{corollary}

\begin{proof}
Diviser les deux équations de
\cref{eq:ew-Higgs-tree-relations} donne
\(m_H^2/m_W^2=8\lambda_H/g^2\). Avec
\(g^2=4\pi\alpha/(1-D_A)\) et le carré de
\cref{eq:ew-HW-ratio}, on obtient \cref{eq:ew-Higgs-lambda}. D'autre part,
\(G_F^{(0)}=1/(\sqrt2v^2)\) et
\(m_H^2=2\lambda_Hv^2\) impliquent
\cref{eq:ew-GF-Higgs}.
\end{proof}

La valeur de \(\lambda_H\) est antérieure à l'incorporation des corrections
radiatives et dépend de
la convention et de l'échelle déclarées. Son intérêt structurel est qu'elle
sépare deux mémoires : le secteur commun de jauge emploie l'invariant pair \(D_A\),
tandis que la route de Higgs retient explicitement \(C^*\).

\section{Clôture déterminantale du secteur électrofaible : mélange, couplages et spectre constitutif}

La chaîne complète de ce chapitre est
\begin{equation}
\label{eq:ew-causal-diagram}
\boxed{
\begin{aligned}
(A,C^*)&\longrightarrow T
\longrightarrow
\begin{cases}
F(T)^2&\longrightarrow
(\widehat\mu,\widehat\varepsilon,\widehat Z,\widehat c),\\
\det T=D_A&\longrightarrow
(s_W^2,c_W^2),
\end{cases}\\
(D_A,\alpha)&\longrightarrow(e_g,g,g'),\\
(D_A,\alpha,m_W)&\longrightarrow G_F^{(0)},\\
(A,C^*,D_A,\alpha)&\longrightarrow
\left(\frac{m_H}{m_W},\lambda_H\right).
\end{aligned}}
\end{equation}
Le chiffre apparaît toujours à la fin. Ni \(\sin^2\theta_W\), ni
\(G_F\), ni \(\lambda_H\) ne sont introduits comme valeurs cibles pour sélectionner
\(A,C^*\) ou les signatures de route.

\begin{proposition}[Résultat principal]
Le même opérateur angulaire détermine le vide par une fonction rationnelle
de ses feuillets et le défaut électrofaible par son déterminant. La
différence d'un pas entre \(W\) et \(Z\) transforme ce déterminant en
\(\sin^2\theta_W^{\rm OS}=1-D_A\) ; la route de Higgs, en revanche, conserve
le contre-angle et conserve l'orientation éliminée par le déterminant.
\end{proposition}

\begin{remark}[Contrôle négatif]
La relation de Weinberg est réfutée dans son domaine si les signatures
angulaires \(W/Z\) ne diffèrent pas uniquement par \(n_A:94\to95\), si une coordonnée
de raffinement ne s'annule pas dans l'interface déclarée ou si
\((m_W^{\angle}/m_Z^{\angle})^2\neq D_A\). La promotion de jauge est
bloquée s'il n'existe pas d'entrelaceur avec la base \((W^3,B)\). Les chiffres
de \(g,g',G_F,\lambda_H\) changent si l'on modifie le schéma ou l'échelle sans
l'enregistrer. Utiliser les valeurs observées de \(G_F\) ou \(m_H\) pour sélectionner
\(\Delta r\), les signatures ou \(C^*\) introduit une rétroaction par la
valeur cible et invalide la dérivation. Confondre \(e_g\) avec la charge
dimensionnelle constitue une erreur de type.
\end{remark}

\subsection*{Portée logique et domaine de validité}
Déterminant angulaire, quotient \(W/Z\) et défaut complémentaire de Weinberg :
\emph{Théorème interne exact dans la carte angulaire conditionnelle}. Couplages \(g,g'\),
\(\rho_{\rm EW}^{(0)}\), Fermi et Higgs : \emph{Composition exacte dans
l'interface électrofaible à l'arbre}. L'entrelaceur de base de jauge et la
correction interne \(\Delta r\) : \emph{Frontière ouverte avec critère de réfutation}.
Fermi et Higgs ne sont pas des constantes indépendantes sélectionnées sans valeurs
cibles : ce sont des compositions exactes qui emploient les signatures, la grandeur
\(m_W\) et la
convention à l'arbre déclarée. Les valeurs conventionnelles ultérieures
sont une \emph{Comparaison métrologique externe}.

 \section{Involution des feuillets dans \(4\) représentations}

L'involution \(C^*\mapsto-C^*\) organise quatre réponses distinctes :
\begin{equation}
\label{eq:modular-sheet-parities}
\begin{aligned}
S_{\rm bi}(A,-C^*)&=S_{\rm bi}(A,C^*),\\
\Gamma_{6\to8}(A,-C^*)&=-\Gamma_{6\to8}(A,C^*),\\
(\widehat\mu,\widehat\varepsilon)&\longleftrightarrow
(\widehat\varepsilon,\widehat\mu),\\
\widehat c&\longmapsto\widehat c,
\qquad D_A\longmapsto D_A.
\end{aligned}
\end{equation}
L'information bicouche est paire ; l'orientation aréale est impaire ; les deux
constitutives forment un doublet, et les déterminants restent fixes. Ce tableau
des parités précède toute interprétation comme symétrie physique.

\section{Transformation rationnelle involutive dans l'espace des paramètres CKM}

Soit
\begin{equation}
\label{eq:modular-Mckm}
M_{\rm CKM}=
\begin{pmatrix}
2&-5&28\\
0&7&-25/2\\
1&-20&-2/3
\end{pmatrix},
\qquad
\det M_{\rm CKM}=-\frac{3857}{6}.
\end{equation}
L'inverse existe sur \(\mathbb Q\) et est
\begin{equation}
\label{eq:modular-Mckm-inverse}
M_{\rm CKM}^{-1}=
\begin{pmatrix}
1528/3857&3380/3857&801/3857\\
75/3857&176/3857&-150/3857\\
6/551&-30/551&-12/551
\end{pmatrix}.
\end{equation}
Avec
\begin{equation}
\label{eq:modular-ckm-coordinates}
h=\frac{C^*}{6},
\qquad
\Delta=\frac{180}{\pi}\Delta_4,
\qquad
\boldsymbol\theta=M_{\rm CKM}(A,h,\Delta)^T,
\end{equation}
la coordonnée impaire entre dans la direction
\begin{equation}
\label{eq:modular-ch}
c_h=\begin{pmatrix}-5\\7\\-20\end{pmatrix}.
\end{equation}

Si \(c_A=(2,0,1)^T\) et
\(c_\Delta=(28,-25/2,-2/3)^T\), les trois colonnes
\((c_A,c_h,c_\Delta)\) sont précisément la base transportée par
\(M_{\rm CKM}\). La deuxième ligne de
\cref{eq:modular-Mckm-inverse} est le covecteur \(\ell_h\) qui extrait la
coordonnée impaire :
\begin{equation}
\label{eq:modular-dual-coordinate-check}
\ell_hc_A=0,
\qquad
\ell_hc_h=1,
\qquad
\ell_hc_\Delta=0.
\end{equation}

\begin{theorem}[Réflexion CKM induite par le feuillet]
\label{thm:modular-ckm-reflection}
Soit
\begin{equation}
\label{eq:modular-Dh-lh}
D_h=\operatorname{diag}(1,-1,1),
\qquad
\ell_h=\frac1{3857}\begin{pmatrix}75&176&-150\end{pmatrix}.
\end{equation}
L'involution induite dans l'espace de \(\boldsymbol\theta\) est
\begin{equation}
\label{eq:modular-Rckm}
\boxed{
\mathcal R_{\rm CKM}
=M_{\rm CKM}D_hM_{\rm CKM}^{-1}
=I-2c_h\ell_h.}
\end{equation}
Satisfait
\begin{equation}
\label{eq:modular-Rckm-properties}
\mathcal R_{\rm CKM}^2=I,
\qquad
\det\mathcal R_{\rm CKM}=-1,
\qquad
\mathcal R_{\rm CKM}M_{\rm CKM}=M_{\rm CKM}D_h.
\end{equation}
\end{theorem}

\begin{proof}
La multiplication de
\cref{eq:modular-Mckm,eq:modular-Mckm-inverse} donne l'identité ; en
particulier, \cref{eq:modular-dual-coordinate-check} prouve que
\(c_h\ell_h\) est le projecteur de rang un sur \(\mathbb Qc_h\) le
long de \(\operatorname{span}\{c_A,c_\Delta\}\). Le produit exact
\(\ell_hc_h=1\) implique
\[
(I-2c_h\ell_h)^2
=I-4c_h\ell_h+4c_h(\ell_hc_h)\ell_h=I.
\]
L'opérateur fixe point par point l'hyperplan
\(\ker\ell_h=\operatorname{span}\{c_A,c_\Delta\}\) et change le signe de
\(c_h\). Ses valeurs propres sont donc \(1,1,-1\) ; d'où
\(\det\mathcal R_{\rm CKM}=-1\) et
\(\Tr\mathcal R_{\rm CKM}=1\). Enfin,
\[
\mathcal R_{\rm CKM}M_{\rm CKM}
=M_{\rm CKM}D_hM_{\rm CKM}^{-1}M_{\rm CKM}
=M_{\rm CKM}D_h,
\]
ce qui prouve l'identité d'entrelacement.
\end{proof}

La matrice explicite est
\begin{equation}
\label{eq:modular-Rckm-matrix}
\mathcal R_{\rm CKM}=
\begin{pmatrix}
4607/3857&1760/3857&-1500/3857\\
-150/551&199/551&300/551\\
3000/3857&7040/3857&-2143/3857
\end{pmatrix}.
\end{equation}
Bien que \(\mathcal R_{\rm CKM}\) ne soit pas une réflexion orthogonale pour le
produit euclidien des trois coordonnées angulaires, elle l'est pour la
métrique rationnelle transportée
\begin{equation}
\label{eq:modular-ckm-metric}
G_{\rm CKM}=(M_{\rm CKM}^{-1})^TM_{\rm CKM}^{-1},
\qquad
\mathcal R_{\rm CKM}^TG_{\rm CKM}\mathcal R_{\rm CKM}=G_{\rm CKM}.
\end{equation}
En effet, dans la base \((c_A,c_h,c_\Delta)\), la métrique est l'identité
et l'involution est \(D_h\). C'est pourquoi le terme ``réflexion'' a un
contenu métrique précis et ne décrit pas seulement le fait que
\(\mathcal R_{\rm CKM}^2=I\).
La phase \(\delta_{\rm CP}=9A\) reste fixe sous cette réflexion. C'est
pourquoi \(\mathcal R_{\rm CKM}\) n'est pas la conjugaison CP physique, mais le
transport de l'inversion de feuillet HMT à l'espace de mélange.

L'inverse rationnel de \(M_{\rm CKM}\) retrouve
\begin{align}
A&=\frac{1528\theta_{12}+3380\theta_{23}+801\theta_{13}}{3857},
\label{eq:modular-A-from-ckm}\\
C^*&=\frac{6(75\theta_{12}+176\theta_{23}-150\theta_{13})}{3857}.
\label{eq:modular-C-from-ckm}
\end{align}
La troisième coordonnée se retrouve également sans perte :
\begin{equation}
\label{eq:modular-Delta-from-ckm}
\Delta
=\frac{6\theta_{12}-30\theta_{23}-12\theta_{13}}{551}.
\end{equation}
Ces équations sont des changements de coordonnées. Elles ne font pas de CKM la cause
rétrospective de \(A,C^*\) ou de \(\gammaH\).

Comme reconnaissance terminale, le constructeur préalable détermine
\begin{equation}
\label{eq:modular-ckm-values}
(\theta_{12},\theta_{23},\theta_{13},\delta_{\rm CP})
=(13.003313589509^\circ,
2.398056157332^\circ,
0.213977382244^\circ,
65.676173123554^\circ).
\end{equation}
La réflexion rationnelle est démontrée avant l'utilisation de ces décimales.

\begin{proposition}[Résultat principal]
La transition \(6\to8\) est d'abord formalisée par l'inclusion
\(\iota_{6,8}:\mathcal F_6\hookrightarrow\mathcal F_8\), dont on
obtient
\(90,120\) et le défaut \(-1/12\). Le mot de six trits fournit
ensuite la bijection \(729\leftrightarrow729\), et
\(\mathcal J_{6\to8}\) transporte les deux modes collectifs au bord. Sur
cet antécédent, l'aire et le hamiltonien modulaire reconstruisent exactement
les deux familles ; la purification thermochamp double détermine l'entropie
d'intrication, et la première loi
finie conserve le défaut d'entropie relative. Dans CKM, l'inversion de
feuillet devient une réflexion rationnelle de rang impair un.
\end{proposition}

\begin{remark}[Contrôle négatif]
La généalogie initiale échoue si \(\iota_{6,8}\) n'est pas injective, si les
décomptes ne sont pas \(90/120\), si \(\delta_{6,8}^{\rm inc}\neq-1/12\), si
\(\beta_{729}\) n'est pas bijective ou si elle est présentée à tort comme un
isomorphisme additif. Le transport collectif échoue si
\(\mathcal J_{6\to8}D_{\rm inc}\neq
D_{\rm area}\mathcal J_{6\to8}\). La reconstruction échoue si le
déterminant des observables de bord n'est pas trente ou si
\cref{eq:modular-inverse-channels} ne retrouve pas les occupations.
La TFD échoue si sa trace partielle n'est pas \(\rho_A\). La première loi ne
s'étend pas aux changements finis sans l'entropie relative de
\cref{eq:modular-finite-first-law}. Le type
\(III_{\lambda_\partial}\) est invalidé si un canal ne persiste pas, s'il n'y a pas de
produit compatible ou si le groupe des rapports change. La réflexion CKM
est réfutée par n'importe laquelle des conditions
\(\mathcal R^2\neq I\), \(\det\mathcal R\neq-1\) ou
\(\mathcal RM\neq MD_h\). Une formule
\(\boldsymbol\theta=f(\gammaH)\) sans entrelaceur viole l'indépendance
typée des deux réalisations fonctionnelles de l'antécédent commun.
\end{remark}

\subsection*{Portée logique et domaine de validité}
Inclusion \(H\subset O\), décomptes \(90/120\), deux normalisations du
défaut \(-1/12\), bijection de six trits \(729\leftrightarrow729\),
entrelaceur collectif, échelle \(a_0\), opérateurs finis, spectre aréal,
TFD, défaut orienté, première loi locale et finie et réflexion CKM :
\emph{Théorème interne exact avec ses structures de départ explicites}.
Classification
\(III_{\lambda_\partial}\) : \emph{Théorème avec hypothèses explicites de
persistance}. Interprétation de la copie thermochamp double comme extérieur physique :
\emph{Réalisation physique conditionnelle}. La grandeur \(\gammaH\) est adoptée
comme \emph{Résultat antécédent démontré} ; elle n'est pas redérivée dans ce volume.

 Le lien électrofaible part du même opérateur angulaire qui a produit les constitutives du vide, mais utilise une autre fonctionnelle.\par

Pour obtenir \(\mu\), \(\varepsilon\), \(Z\) et \(c\), il fallait conserver le spectre complet. Il fallait distinguer les deux feuillets.\par

Le déterminant fait l’inverse :\par

\[
D_A=\det T=q_+q_-.
\]

En multipliant les valeurs propres, le déterminant intègre les deux orientations dans une grandeur angulaire commune et invariante. Il publie la partie paire et réserve l’identité du feuillet dans la mémoire orientée.\par

C’est pourquoi \(D_A\) est pair sous l’inversion\par

\[
C^*\longmapsto-C^*.
\]

Le secteur \(W/Z\) possède précisément cette structure. Les routes de \(W\) et de \(Z\) sont adjacentes dans la coordonnée angulaire principale et partagent la même coordonnée orientée. Lorsque l’on forme le quotient, la contribution de \(C^*\) s’annule.\par

Il reste\par

\[
\left(\frac{m_W}{m_Z}\right)^2
=
D_A,
\]

et donc\par

\[
\sin^2\theta_W
=
1-D_A.
\]

L’identité révèle que l’angle de Weinberg est le défaut complémentaire du déterminant angulaire.\par

Le déterminant mesure la partie commune qui subsiste lorsque les deux feuillets sont rabattus sur leur produit. \(1-D_A\) mesure la partie complémentaire nécessaire au mélange.\par

Nous pourrions le dire ainsi : le vide conserve toute la réponse ; le secteur électrofaible utilise la compression paire de cette réponse.\par

Le même antécédent angulaire produit deux fonctions différentes :\par

\begin{itemize}[leftmargin=*,itemsep=0.35em]
\item le calcul fonctionnel complet engendre les relations constitutives ;
\item le déterminant engendre l’invariant de mélange.
\end{itemize}

Cela explique pourquoi le vide et l’angle de Weinberg sont deux réalisations fonctionnelles distinctes d’un même antécédent angulaire.\par

À partir de là, les couplages \(g\) et \(g'\) sont, au niveau déclaré, les deux projections du même couplage électromagnétique rationalisé sur les directions sinus et cosinus du mélange :\par

\[
g^2
=
\frac{4\pi\alpha}{1-D_A},
\qquad
g'^2
=
\frac{4\pi\alpha}{D_A}.
\]

L’identité custodiale à l’arbre\par

\[
\rho_{\rm EW}^{(0)}=1
\]

exprime que la même partition angulaire organise de manière cohérente masse et couplage.\par

La constante de Fermi apparaît ensuite comme la réponse effective à basse énergie du canal chargé :\par

\[
G_F^{(0)}
=
\frac{\pi\alpha}
{\sqrt2(1-D_A)m_W^2}.
\]

La constante de Fermi apparaît comme une composition de la structure fine, du défaut angulaire et de l’échelle du porteur \(W\).\par

Le secteur de Higgs introduit une différence décisive. Sa route possède une coordonnée orientée différente de celle de la route de \(W\), de sorte que \(C^*\) persiste. L’autocouplage scalaire publie l’orientation que le déterminant réserve.\par

Le secteur de jauge est pair : il utilise \(D_A\).\par
Le secteur scalaire est orienté : il retient \(C^*\).\par

Cette division est physiquement suggestive. Les champs de jauge organisent la partie commune de la transformation ; le secteur scalaire conserve la mémoire de l’orientation brisée.\par

Les corrections radiatives appartiennent à un niveau ultérieur et doivent être construites prospectivement à partir des boucles, états et échelles internes, avec \(G_F\) comme sortie et évaluation terminale. Le squelette causal est déjà présent : vide, mélange, couplages et réponse de Fermi procèdent de différentes lectures du même état angulaire enrichi.\par

 \section{Clôture angulaire électrofaible et conditions de compatibilité du transducteur de jauge}
\label{sec:p3-final:cierre-electrodebil}

Soient
\[
 c_W=\sqrt{D_A},
 \qquad
 s_W=\sqrt{1-D_A}.
\]
La rotation
\[
 \begin{pmatrix}Z\\A\end{pmatrix}
 =
 \begin{pmatrix}c_W&-s_W\\s_W&c_W\end{pmatrix}
 \begin{pmatrix}W^3\\B\end{pmatrix}
\]
est orthogonale et préserve l’orientation. Les identités
\[
 \left(\frac{m_W}{m_Z}\right)^2=D_A,
 \qquad
 \sin^2\theta_W=1-D_A
\]
sont closes dans la carte angulaire. La promotion de jauge exige que le
transducteur préserve la forme cinétique, les charges et les identités de Ward ;
la correction \(\Delta r\) appartient à la dynamique radiative et ne s’obtient pas
en l’ajustant à \(G_F\).

  \endgroup

\chapter{Dérivation intrinsèque de la gravitation et réciprocité contragrédiente \texorpdfstring{\(G\leftrightarrow\hbar\)}{G<->hbar}}
\begingroup
\renewcommand{\chapter}[2][]{}%
\chapter{Dérivation intrinsèque de \(G\) et réciprocité \(G\leftrightarrow\hbar\) : clôture de phase, longueurs de Compton et de Planck et invariance aréale}
\label{chap:p3-final:50:gravitacion_accion}
\label{ch:gravity}
% BEGIN CR28_RADIAL_FULL
% Sección incorporable; requiere amsmath y amssymb.
% Fuente principal: RIGIDEZ_SIMULTANEA_DEL_LECTOR_RADIAL_20260926.md.
% Complementos demostrativos: COMPOSICION_METRICA_MEMORIA_Y_ACOPLAMIENTO.md
% y ACCION_ACOPLADA_FASE_Y_MEMORIA_DE_FRONTERA.md.
% La integración, la compilación y la revisión visual corresponden al editor.
\section{Rigidité simultanée du lecteur longitudinal et du couplage radial}
\label{g26:sec:rigidez-radial-es}

La construction de l'holographie modulaire triadique transmet à la réalisation
radiale un état enrichi avec ses coordonnées, son orientation, sa retenue,
sa mémoire et ses incidences. La suite
\[
 \begin{gathered}
 \mathrm{APP}\longrightarrow\mathrm{TRIT}\longrightarrow\mathrm{TPK}
 \longrightarrow\text{état enrichi}\\
 \longrightarrow\text{structure discrète conjointe du continu}
 \end{gathered}
\]
conserve la projection archimédienne et la projection d'incidence, ainsi que les cinq constructions consubstantielles du continu. Le prolongement
\[
 w_6\longrightarrow w_{12}\longrightarrow w_{18}\longrightarrow
 w_{24}\longrightarrow w_{30}\longrightarrow R_{36}\longrightarrow G_9
\]
s'articule avec l'ordre opératoriel
\(U_t\longrightarrow\Gamma_9\longrightarrow K_{\rm ph}\) : la phase revient
et le registre de mémoire s'incrémente. Les coordonnées régionales, le
registre sélectionné \(K\), la coordonnée \(\alpha\) et l'action de retour
sont reçus avec cette provenance commune et leurs lecteurs préalablement
construits.

Le résultat de cette section est une classification de réalisations qui conservent simultanément les données générées et les règles longitudinales, métriques et temporelles précisées ci-après. Pour chaque caractère positif, la réponse radiale est déterminée ; le coefficient qui relie cette réponse à la masse est commun à tous les caractères admissibles. L'interprétation gravitationnelle est formulée après la construction des deux lectures sur le même espace.

\subsection{Registre des incidences et contraintes constitutives}

On considère des espaces de Hilbert complexes \(V_{\rm in}\) et
\(V_{\rm out}\) de dimension \(N\), avec leurs produits scalaires fixés.
Le cardinal marqué spécifié dans R1 est positif ; en particulier,
les deux fibres sont non nulles. Les adjoints sont pris par rapport à ces produits. Le symbole
\(\mathsf U\) désignera le transport normalisé de l'archive et
\(U_t\) conservera sa signification d'évolution de l'état enrichi.

\paragraph{R1. Registre marqué.}
L'ensemble des incidences est
\begin{equation}
 \Omega=\left\{(j,k,q,u,v):
 \begin{array}{l}
 j,k\in\mathbb Z/12\mathbb Z,\quad k-j\in\{0,3,6,9\},\\
 1\leq q\leq8,\quad u<v,\quad
 u,v\in\{1,2,4,5,7,8\}
 \end{array}\right\}.
 \label{g26:eq:omega-rigidez-es}
\end{equation}
Les douze phases, les quatre positions relatives, les huit positions de l'
indice \(q\) et les quinze paires du dernier ensemble donnent
\begin{equation}
 N=|\Omega|=12\cdot4\cdot8\binom62
   =12\cdot4\cdot120=5760.
 \label{g26:eq:cardinal-rigidez-es}
\end{equation}
Chaque incidence possède sa marque individuelle. Dans la base correspondante \((e_i)_{i\in\Omega}\) de \(V_{\rm out}\), soient \(P_i=e_ie_i^*\). Le mode commun \(\boldsymbol u\) est normalisé par \(|u_i|^2=1/N\) ; il s'écrit \(P=\boldsymbol u\boldsymbol u^*\). L'algèbre marquée est \(\mathcal A_\Omega=\operatorname{span}\{P_i:i\in\Omega\}\).

\paragraph{R2. Normalisation et archive complète.}
L'inverseur généré \(B:V_{\rm in}\to V_{\rm out}\) satisfait
\begin{equation}
 B^*B=\tfrac14 I_{\rm in},\qquad
 BB^*=\tfrac14 I_{\rm out}.
 \label{g26:eq:normalizacion-b-es}
\end{equation}
Les deux composantes sont conservées
\begin{equation}
 C_0=(I-P)B,\qquad M_0=PB,\qquad
 C_0+M_0=B,\qquad \mathsf U=2B.
 \label{g26:eq:archivo-completo-es}
\end{equation}
La première enregistre la variation par rapport au mode commun ; la seconde
conserve ce mode. Leurs images sont orthogonales et leur somme restitue
le transport complet. D'après \eqref{g26:eq:normalizacion-b-es},
\(\mathsf U\) est unitaire.

\paragraph{R3. Inscription individuelle.}
L'application \(\mathcal E:\operatorname{End}(V_{\rm out})\to\mathcal A_\Omega\) est une espérance conditionnelle linéaire qui fixe chaque \(P_i\) et satisfait
\begin{equation}
 \mathcal E(A_1TA_2)=A_1\mathcal E(T)A_2,
 \qquad A_1,A_2\in\mathcal A_\Omega.
 \label{g26:eq:bimodularidad-es}
\end{equation}
L'inscription du résultat conserve, dans un registre complémentaire,
les données complètes de la construction. L'espérance détermine
l'observable inscrit ; la conservation de l'archive détermine sa
reconstruction ultérieure.

\paragraph{R4. Composition longitudinale.}
On fixe la coordonnée générée \(\alpha>0\) et la section de référence longitudinale \(L_*>0\). Le transport longitudinal est défini par
\begin{equation}
 D=L_*\alpha^{16}\mathcal E(C_0C_0^*)\mathsf U:
 V_{\rm in}\longrightarrow V_{\rm out}.
 \label{g26:eq:regla-longitudinal-es}
\end{equation}
Cette règle applique linéairement le coefficient inscrit à la section
longitudinale. Son type fixe la position de la puissance \(\alpha^{16}\)
et du coefficient de retour dans la composition. La lecture d'
une racine d'intensité constitue une opération différente et possède sa
propre règle de transformation.

\paragraph{R5. Métrique radiale positive.}
Soit \(X=X^*>0\) le caractère complet sur \(V_{\rm in}\). On définit
la lecture radiale positive \(R_X\) sur \(V_{\rm out}\) par
\begin{equation}
 \|R_Xv\|^2=\|XD^*v\|^2
 \qquad(v\in V_{\rm out}).
 \label{g26:eq:metrica-radial-es}
\end{equation}
Cette égalité fixe la métrique sur chaque vecteur de la fibre. En notation polaire, \(R_X\) est le module d'arrivée \(\lvert(DX)^*\rvert=\sqrt{DX^2D^*}\).

\paragraph{R6. Action et horloge.}
La section d'action générée \(\hbar_{\rm ret}>0\) et le relèvement
réel de phase satisfont
\begin{equation}
 S(\theta)=\hbar_{\rm ret}\theta.
 \label{g26:eq:accion-fase-es}
\end{equation}
La vitesse constitutive \(c>0\) est conservée. Une fois obtenue la longueur intensive \(L\) de \eqref{g26:eq:regla-longitudinal-es}, l'horloge conjuguée satisfait \(\omega L=c\). Pour \(\theta(t)=\omega t\), le taux d'action détermine l'échelle énergétique
\begin{equation}
 E_{L}=\frac{dS}{dt}=\hbar_{\rm ret}\omega
   =\frac{\hbar_{\rm ret}c}{L}.
 \label{g26:eq:escala-energetica-es}
\end{equation}
Le relèvement et ses raffinements restent associés à la mémoire
de la route.

\paragraph{R7. Lectures du même caractère.}
Sur \(Y=\mathsf U X\mathsf U^*>0\), on définit
\begin{equation}
 H_X=E_{L}Y,\qquad M_X=\frac{H_X}{c^2},\qquad
 \Lambda_X=LY^{-1}.
 \label{g26:eq:lecturas-conjugadas-es}
\end{equation}
L'énergie, la masse et la longueur conjuguée conservent le même caractère
et le même transport. L'inverse agit dans le secteur positif ; le secteur
nul conserve son traitement propre.

\paragraph{R8. Rayon réduit et sections dimensionnelles.}
L'identification physique de cette réalisation est
\begin{equation}
 R_X=\frac{G}{c^2}M_X.
 \label{g26:eq:radio-reducido-es}
\end{equation}
Ainsi, \(R_X\) représente le rayon gravitationnel réduit. Le rayon de Schwarzschild correspondant est \(2R_X\). On maintient conjointement \(L_*\), la section d'action et la section constitutive de vitesse ; un changement d'unités transforme leurs coordonnées selon leurs dimensions.

\paragraph{Provenance des restrictions.}
Le transport de Hadamard, la conservation de l'archive, l'inscription,
l'action de phase et la réciprocité possèdent des antécédents dans le corpus.
Le registre élargi \eqref{g26:eq:omega-rigidez-es} et les compositions
\eqref{g26:eq:regla-longitudinal-es}--\eqref{g26:eq:metrica-radial-es} sont
incorporés explicitement dans la construction réunie. R4 et R5 sont des règles
constitutives de cette réalisation ; la démonstration suivante établit
leur conséquence conjointe avec les autres antécédents. R8 précise
l'identification physique qui fixe le sens du coefficient.

\subsection{Théorème de détermination simultanée}

\paragraph{Théorème.}
Pour les mêmes données générées \(B\), \((P_i)\), \(P\),
\(\alpha\), \(L_*\), \(\hbar_{\rm ret}\), \(c\) et le même caractère
\(X\), les contraintes R1--R8 déterminent univoquement l'inscription,
le transport longitudinal et la réponse radiale. Pour tous les
caractères positifs admissibles, le coefficient \(G\) est commun. On a
\begin{align}
 \mathcal E&=\Delta_\Omega,
 &\Delta_\Omega(T)&=\sum_{i\in\Omega}P_iTP_i,\label{g26:eq:delta-unica-es}\\
 r_N&=\frac{N-1}{4N}=\frac{5759}{23040},
 &L&=L_*\alpha^{16}r_N,\label{g26:eq:longitud-unica-es}\\
 D&=L\mathsf U,
 &R_X&=L\mathsf U X\mathsf U^*,\label{g26:eq:radial-unica-es}\\
 R_X\Lambda_X&=L^2I,
 &H_X\Lambda_X&=\hbar_{\rm ret}cI.\label{g26:eq:productos-radiales-es}
\end{align}
Le couplage est déterminé par
\begin{equation}
 \boxed{G=\frac{c^3L^2}{\hbar_{\rm ret}}
 =\frac{c^3L_*^2}{\hbar_{\rm ret}}
   \left(\frac{5759}{23040}\right)^2\alpha^{32}.}
 \label{g26:eq:g-rigido-es}
\end{equation}

\paragraph{Démonstration : unicité de l'inscription.}
Soient \(E_{ij}=e_ie_j^*\) les unités matricielles. Par bimodularité,
\[
 \mathcal E(E_{ij})=P_i\mathcal E(E_{ij})P_j.
\]
Pour \(i\ne j\), l'appartenance de \(\mathcal E(E_{ij})\) à l'algèbre diagonale implique \(P_i\mathcal E(E_{ij})P_j=0\). Pour \(i=j\), la fixation de \(P_i\) donne \(\mathcal E(E_{ii})=P_i\). La linéarité détermine alors \eqref{g26:eq:delta-unica-es} sur toutes les matrices. L'inscription marquée est ainsi fixée par ses propriétés algébriques.

\paragraph{Démonstration : coefficient de retour et archive.}
La normalisation de \(B\) et \(P^2=P=P^*\) fournissent
\[
 C_0C_0^*=\tfrac14(I-P),\qquad M_0M_0^*=\tfrac14P.
\]
Comme \(P_iPP_i=|u_i|^2P_i=P_i/N\), on obtient
\begin{equation}
 \Delta_\Omega(C_0C_0^*)=\frac{N-1}{4N}I,
 \qquad
 \Delta_\Omega(M_0M_0^*)=\frac{1}{4N}I.
 \label{g26:eq:retorno-y-memoria-es}
\end{equation}
Les deux lectures ont pour somme \(I/4\). La première est scalaire, de sorte que
toute distribution normalisée \((p_i)\) sur les incidences,
avec \(p_i\geq0\) et \(\sum_i p_i=1\), produit
\(\sum_i p_i r_N=r_N\). Le même coefficient est donc préservé
sous une pondération diagonale arbitraire.

\paragraph{Démonstration : longueur et métrique.}
La substitution de \eqref{g26:eq:retorno-y-memoria-es} dans R4 donne
\(D=L\mathsf U\), avec \(L\) comme dans \eqref{g26:eq:longitud-unica-es}.
L'unitarité de \(\mathsf U\) implique \(DD^*=L^2I\).
D'après R5, les formes quadratiques de \(R_X^2\) et \(DX^2D^*\) coïncident.
L'identité de polarisation détermine
\[
 R_X^2=DX^2D^*=L^2\mathsf U X^2\mathsf U^*.
\]
L'opérateur \(L\mathsf U X\mathsf U^*\) est positif et son carré est l'opérateur du second membre. L'unicité de la racine carrée positive démontre \eqref{g26:eq:radial-unica-es}. En particulier, \(R_I=LI\).

\paragraph{Démonstration : action, réciprocité et couplage.}
R6 et R7 donnent
\[
 H_X=\frac{\hbar_{\rm ret}c}{L}Y,
 \qquad M_X=\frac{\hbar_{\rm ret}}{cL}Y,
 \qquad \Lambda_X=LY^{-1}.
\]
La multiplication produit les deux produits de
\eqref{g26:eq:productos-radiales-es}. Ensuite, R8 équivaut à
\[
 LY=\frac{G\hbar_{\rm ret}}{c^3L}Y.
\]
L'inversibilité de \(Y\) détermine \eqref{g26:eq:g-rigido-es}.
La même substitution vérifie l'égalité pour chaque caractère positif.
Si \(G_1\) et \(G_2\) satisfont l'identité avec les mêmes données,
\((G_1-G_2)Y=0\), et donc \(G_1=G_2\).
Cette conclusion distingue la réponse dépendant de l'état
\(R_X\) du coefficient commun \(G\). \hfill\(\square\)

\subsection{Réciprocité et rigidité de la phase}

L'antécédent puissance--logarithmique utilise, pour \(x>0\),
\[
 \ell_p(x)=\frac{x^p-1}{p}\quad(p\ne0),
 \qquad \ell_0(x)=\log x.
\]
La substitution directe démontre
\(\ell_{-p}(x^{-1})=-\ell_p(x)\). Dans le domaine positif de leurs
inverses, cette identité équivaut à
\[
 \exp_p(s)\exp_{-p}(-s)=1.
\]
La composition spectrale sur un caractère positif transmet ce produit à la paire \(LY\), \(LY^{-1}\). Dans la réalisation actuelle, R5 peut s'exprimer de manière équivalente, une fois R1--R4 et R7 fixées, par \(R_X\Lambda_X=L^2I\) : multiplier par \(\Lambda_X^{-1}\) détermine \(R_X=L\mathsf U X\mathsf U^*\), et cette expression satisfait R5.

Les deux produits de \eqref{g26:eq:productos-radiales-es} donnent en outre une
forme directe de la rigidité :
\begin{equation}
 R_X=\frac{L^2}{\hbar_{\rm ret}c}H_X.
 \label{g26:eq:proporcion-operadores-es}
\end{equation}
Pour un changement d'échelle positif \(s\), le couple
\((sR_X,\Lambda_X/s)\) conserve le produit radial ; le produit
\(H_X(\Lambda_X/s)=\hbar_{\rm ret}cI/s\), à énergie et action
fixées, exige \(s=1\). La conservation conjointe détermine l'échelle.
De même, tout opérateur radial positif ayant la même métrique R5
coïncide avec \(R_X\) par l'unicité de la racine positive.

L'action relevée possède aussi une détermination exacte. Pour
\(\theta\ne0\), l'égalité \(S/\hbar_{\rm ret}=\theta\) fixe \(S\).
Même lorsque les rotations de toutes les subdivisions
\(\theta/9^n\) sont conservées, un changement d'échelle réel fixe \(a\) qui satisfait
\[
 \exp(-ia\theta/9^n)=\exp(-i\theta/9^n)
 \qquad(n\geq0)
\]
vérifie \((a-1)\theta/9^n\in2\pi\mathbb Z\). Cette suite tend vers zéro ; pour \(n\) suffisamment grand, sa seule valeur possible est zéro. De \(\theta\ne0\), on déduit \(a=1\).

\subsection{Échelles de Planck et spécialisation du caractère unitaire}

Nous écrivons \(L_\alpha=L\) ; dans la section de retour, \(E_{P,\rm ret}=E_L\).

La longueur construite $L_\alpha>0$, la section d'action $\hbar_a>0$ et la vitesse $c>0$ déterminent les échelles

\[
t_P=\frac{L_\alpha}{c},\qquad
m_{P,a}=\frac{\hbar_a}{cL_\alpha},\qquad
E_{P,a}=\frac{\hbar_a c}{L_\alpha}.
\]

En effet, substituer $G_a=c^3L_\alpha^2/\hbar_a$ dans les expressions
$\sqrt{\hbar_aG_a/c^5}$, $\sqrt{\hbar_ac/G_a}$ et
$\sqrt{\hbar_ac^5/G_a}$ donne respectivement ces trois identités.
La positivité fixe dans chaque cas la racine. La longueur précède la
reconnaissance de ces échelles ; leurs expressions conventionnelles sont
des identités de la réalisation déjà construite.

Pour un mode de caractère $y>0$, les lectures massique et longitudinale sont

\[
m(y)=m_{P,a}y,\qquad r_g(y)=L_\alpha y,\qquad
\bar\lambda(y)=\frac{L_\alpha}{y}.
\]

Par conséquent, $r_g(y)\bar\lambda(y)=L_\alpha^2$ pour tous les modes
positifs. L'égalité des deux longueurs caractérise le mode $y=1$,
puisque $y=y^{-1}$ et $y>0$ impliquent $y=1$. Dans ce mode, la masse est
$m_{P,a}$ et les deux longueurs coïncident avec $L_\alpha$. On distingue ainsi
la réciprocité générale de sa spécialisation au caractère unitaire.

Le pas chronologique et le temps de Planck satisfont $t_0=(\pi_{\rm HMT}/54)t_P$. La période complète de phase de 108 pas est $108t_0=2\pi_{\rm HMT}t_P$ et sa fréquence angulaire est $\omega_{108}=1/t_P=c/L_\alpha$. Un horizon de rayon $R_h=\zeta_hL_\alpha$, avec $\zeta_h>0$, conserve sa propre variable géométrique. Le rayon de Schwarzschild associé à une masse est $2r_g$.

Si l'on compare les deux sections d'action à $L_\alpha$ et $c$ fixés,
le quotient $\rho=\hbar_b/\hbar_a$ donne
$G_b=G_a/\rho$, $m_{P,b}=\rho m_{P,a}$ et
$E_{P,b}=\rho E_{P,a}$ ; $L_\alpha$ et $t_P$ restent identiques.
Cette comparaison entre sections maintient la distinction avec une
évolution temporelle des constantes.

\subsection{Covariance de l'archive et transformation des unités}

Soient \(T:V_{\rm in}\to V'_{\rm in}\) et
\(S:V_{\rm out}\to V'_{\rm out}\) unitaires. Le transport conjoint
des données est
\[
 B'=SBT^*,\quad P'=SPS^*,\quad P_i'=SP_iS^*,\quad X'=TXT^*.
\]
L'espérance transportée \(\mathcal E'(A')=S\mathcal E(S^*A'S)S^*\) fixe \(P_i'\) et est bimodulaire dans l'algèbre marquée transformée. Par substitution,
\begin{align*}
 \mathsf U'&=S\mathsf UT^*,&D'&=SDT^*,\\
 R'_{X'}&=SR_XS^*,&H'_{X'}&=SH_XS^*,&
 \Lambda'_{X'}&=S\Lambda_XS^*.
\end{align*}
La conjugaison conserve les égalités de R1--R8 et le même coefficient
\(G\). Les marques, orientations et registres sont transportés avec les
bases. Les permutations d'incidences constituent un cas particulier.

L'incorporation d'un espace auxiliaire de mémoire \(\mathcal F\)
transporte \(P\) vers \(P\otimes I_{\mathcal F}\),
\(\mathsf U\) vers \(\mathsf U\otimes I_{\mathcal F}\) et
\(\Delta_\Omega\) vers \(\Delta_\Omega\otimes\operatorname{id}\).
L'inscription de \(C_0C_0^*\otimes I_{\mathcal F}\) est
\(r_N I\otimes I_{\mathcal F}\) ; le coefficient conserve le cardinal
marqué original.

Pour préciser la covariance dimensionnelle, soient \(\ell\), \(\tau\) et \(m\) les facteurs positifs des nouvelles unités de longueur, de temps et de masse par rapport aux précédentes. Les coordonnées numériques satisfont
\[
 L'=L/\ell,\qquad c'=c\tau/\ell,\qquad
 \hbar'_{\rm ret}=\hbar_{\rm ret}\tau/(m\ell^2).
\]
Ainsi,
\[
 \frac{(c')^3(L')^2}{\hbar'_{\rm ret}}
 =G\frac{m\tau^2}{\ell^3},
\]
qui est la transformation d'une grandeur de dimension
\(\mathrm{longitud}^3\mathrm{masa}^{-1}\mathrm{tiempo}^{-2}\).
La section \(L_*\) se transforme comme \(L\), tandis que \(\alpha\)
et \(r_N\) restent sans dimension.

\subsection{Élimination exacte de la mémoire couplée}

Soient \(\mathcal K\) et \(\mathcal H\) des espaces de Hilbert finis,
\(W=W^*>0\) un poids global et
\(\mathcal C:\mathcal K\to\mathcal H\) surjective. L'action
des résidus \(r\), avec un défaut de frontière \(z\), est
\[
 \mathcal S(r)=\mathfrak a\,r^*Wr,
 \qquad \mathcal Cr=z,\qquad \mathfrak a>0.
\]
Le poids global conserve ses blocs croisés entre incidences. Définissons \(\mathcal R=\mathcal CW^{-1}\mathcal C^*\). Pour \(v\ne0\), l'injectivité de \(\mathcal C^*\) donne
\[
 v^*\mathcal Rv=\|W^{-1/2}\mathcal C^*v\|^2>0.
\]
Par conséquent, \(\mathcal R\) est inversible. Le résidu d'action minimale est
\[
 r_*(z)=W^{-1}\mathcal C^*\mathcal R^{-1}z.
\]
On vérifie \(\mathcal Cr_*=z\). Chaque résidu admissible s'écrit \(r=r_*+k\), avec \(\mathcal Ck=0\), et \(k^*Wr_*=k^*\mathcal C^*\mathcal R^{-1}z=0\). En conséquence,
\begin{equation}
 \mathcal S(r)=\mathfrak a\,z^*\mathcal R^{-1}z
             +\mathfrak a\,k^*Wk.
 \label{g26:eq:accion-residuo-es}
\end{equation}
La positivité de \(W\) prouve l'unicité du minimum. Le couple
\((z,k)\) conserve et reconstruit le résidu complet.

Pour une route avec les transports \(U_j\), les résidus
\(r_j=f_j-U_jf_{j-1}\) se reconstruisent par
\(f_j=U_jf_{j-1}+r_j\). L'application \(\mathcal C\) réunit les
transports ultérieurs de chaque résidu vers l'extrémité. Si ses
blocs sont \(\mathcal C_j\), alors
\[
 \mathcal R=\sum_{j,k}\mathcal C_j(W^{-1})_{jk}\mathcal C_k^*.
\]
Cette expression conserve les corrélations entre tronçons de la route. Pour un transport longitudinal inversible \(D\), la frontière \(y=Dz\) a pour réponse
\[
 \mathcal R_L=D\mathcal RD^*,\qquad
 \mathcal S_{{\rm ef},L}(y)
 =\mathfrak a\,y^*(D\mathcal RD^*)^{-1}y
 =\mathfrak a\,z^*\mathcal R^{-1}z.
\]
On obtient la même action en éliminant d'abord la mémoire ou en transportant
d'abord la frontière. Le préfacteur conserve la normalisation de l'action
source ; lorsque cette action porte \(\mathfrak a=\hbar_{\rm ret}\), la
réduction transmet la même section.

\subsection{Réduction conjointe du rayon, de l'énergie et de la longueur conjuguée}

Écrivons le même caractère positif dans une décomposition en frontière
et mémoire :
\[
 Y=\begin{pmatrix}A&F\\F^*&C\end{pmatrix},\qquad C>0,
 \qquad Y_{\rm ef}=A-FC^{-1}F^*.
\]
L’identité
\[
 \begin{pmatrix}b\\y\end{pmatrix}^{\!*}Y
 \begin{pmatrix}b\\y\end{pmatrix}
 =b^*Y_{\rm ef}b+\eta^*C\eta,
 \qquad \eta=y+C^{-1}F^*b,
\]
prouve \(Y_{\rm ef}>0\), caractérise l'extension d'énergie minimale et conserve la reconstruction \(y=\eta-C^{-1}F^*b\). Pour \(s>0\), la substitution dans le complément de Schur donne \(\operatorname{Schur}(sY)=sY_{\rm ef}\). En conséquence,
\[
 R_{\rm ef}=LY_{\rm ef},\qquad
 H_{\rm ef}=E_{L}Y_{\rm ef}.
\]
Pour calculer la longueur conjuguée comprimée, on résout
\(Y(x,y)=(f,0)\). La seconde équation donne
\(y=-C^{-1}F^*x\), et la première donne
\(x=Y_{\rm ef}^{-1}f\). Le bloc de frontière de \(Y^{-1}\) est
donc \(Y_{\rm ef}^{-1}\). Ainsi,
\begin{equation}
 \Lambda_{\rm b}=LY_{\rm ef}^{-1},\qquad
 R_{\rm ef}\Lambda_{\rm b}=L^2I,\qquad
 R_{\rm ef}=\frac{L}{E_{L}}H_{\rm ef}.
 \label{g26:eq:schur-reciproco-es}
\end{equation}
La réduction agit sur les mêmes blocs dans les deux lectures et
conserve le résidu \(\eta\). Le caractère effectif peut dépendre de
la mémoire sous forme matricielle et le coefficient \(L/E_{L}\) reste fixe.

\subsection{Mémoire dynamique et norme temporelle induite}

L'élimination dynamique conserve une dépendance spectrale supplémentaire. Dans une représentation stationnaire, soit
\[
 H=\begin{pmatrix}H_{\rm bb}&V\\V^*&H_{\rm ii}\end{pmatrix},
 \qquad H_{\rm ii}>0.
\]
Pour \(z\) dans le domaine résolvant pertinent, l'élimination de la
composante intérieure dans \((H-zI)(b,y)=(f,0)\) donne
\begin{equation}
 \mathcal K_H(z)=H_{\rm bb}-zI
       -V(H_{\rm ii}-zI)^{-1}V^*.
 \label{g26:eq:resolvente-memoria-es}
\end{equation}
Lorsque les deux inverses existent, le bloc de frontière de
\((H-zI)^{-1}\) est \(\mathcal K_H(z)^{-1}\). Pour
\(|z|<\lambda_{\min}(H_{\rm ii})\), la série convergente de la
résolvante fournit
\begin{align}
 \mathcal K_H(z)&=H_{\rm est}-zZ-z^2VH_{\rm ii}^{-3}V^*-\cdots,
 \label{g26:eq:expansion-resolvente-es}\\
 H_{\rm est}&=H_{\rm bb}-VH_{\rm ii}^{-1}V^*,\qquad
 Z=I+VH_{\rm ii}^{-2}V^*\geq I.
 \label{g26:eq:norma-z-es}
\end{align}
La positivité découle de \(Z-I=(H_{\rm ii}^{-1}V^*)^*(H_{\rm ii}^{-1}V^*)\). L'extension statique \((b,-H_{\rm ii}^{-1}V^*b)\) a pour carré de norme \(b^*Zb\) ; l'énergie et la norme temporelle proviennent de la même extension. Le paramètre physique est \(z=\hbar_{\rm ret}\omega\).

Au premier ordre temporel, \(T=Z^{1/2}\) et \(\psi=Tb\) transforment
le couple \((H_{\rm est},Z)\) en
\((T^{-*}H_{\rm est}T^{-1},I)\). La section
\(\hbar_{\rm ret}\) demeure fixée. Si \(T\) dépend du temps,
la transformation conserve également le terme \(\dot T b\) de
\(\dot\psi=\dot T b+T\dot b\).

La proportionnalité radiale est conservée dans la résolvante complète.
En écrivant \(\chi_{\rm rad}=L/E_{L}\) et \(R=\chi_{\rm rad} H\), la substitution de
chaque bloc démontre
\begin{equation}
 \mathcal K_R(\chi_{\rm rad} z)=\chi_{\rm rad}\mathcal K_H(z).
 \label{g26:eq:resolvente-proporcional-es}
\end{equation}
Cette identité transporte simultanément l'opérateur et le paramètre spectral, et conserve chaque ordre de la mémoire dynamique. Pour les formes effectives de \eqref{g26:eq:schur-reciproco-es}, la normalisation canonique exige le transport dual de la longueur inverse :
\begin{align*}
 H_{\rm can}&=T^{-*}H_{\rm ef}T^{-1},&
 R_{\rm can}&=T^{-*}R_{\rm ef}T^{-1},&
 \Lambda_{\rm can}&=T\Lambda_{\rm b}T^*.
\end{align*}
Multiplier ces expressions donne
\[
 R_{\rm can}=\frac{L}{E_{L}}H_{\rm can},\qquad
 R_{\rm can}\Lambda_{\rm can}=L^2I.
\]
La démonstration admet \(T\), \(R_{\rm ef}\) et
\(\Lambda_{\rm b}\) non commutatifs. La norme temporelle induite et
la longueur duale conservent conjointement le même \(G\).

\subsection{Raffinement, domaines spectraux et portée}

Soient \(J\) et \(K\) des inscriptions isométriques entre deux niveaux qui
satisfont \(X'J=JX\) et \(\mathsf U'J=K\mathsf U\),
avec la même section intensive \(L\). Alors
\[
 R'K=L\mathsf U'X'\mathsf U'^*K
     =KL\mathsf U X\mathsf U^*=KR,
\]
et le même calcul démontre \(H'K=KH\). L'entrelacement des
inverses se formule sur les domaines correspondants. Lorsque le système
compatible possède une limite autoadjointe positive et injective \(X\),
le calcul spectral définit \(Y=\mathsf U X\mathsf U^*\),
\(R=LY\), \(H=E_{L}Y\) et \(\Lambda=LY^{-1}\). La proportionnalité
\(R=(L/E_{L})H\) vaut sur \(\operatorname{Dom}(Y)\), tandis que
\(R\Lambda=L^2I\) et \(H\Lambda=\hbar_{\rm ret}cI\) valent sur
\(\operatorname{Dom}(Y^{-1})\), avec un prolongement borné de leurs
produits. Les coupures spectrales
\(\mathbf1_{[\varepsilon,M]}(Y)\) vérifient les identités avant
le passage à la limite dans ces domaines. Une longueur d'arête raffinée
\(L/9\) incorpore son propre facteur d'échelle ; la section intensive
se conserve selon l'entrelacement déclaré.

La classification distingue les transformations qui préservent la réalisation de celles qui modifient une règle constitutive. Les pondérations diagonales, les équivalences unitaires, l'archive auxiliaire et la réduction conjointe conservent \(G\). Le changement de résolution marquée modifie R1 ; le remplacement de \(r_N\) par \(\sqrt{r_N}\), de \(\alpha^{16}\) par son carré ou de la composante inscrite par la norme totale modifie R4. La remise à l'échelle isolée du rayon change R5, et celle de l'action change R6. R8 maintient la distinction entre rayon réduit, rayon de Schwarzschild et changement d'unités.

On obtient ainsi l'unicité mathématique de la réalisation R1--R8 : les données
générées et le caractère fixent la réponse, et la composition commune fixe
\(G\). Les règles longitudinale et métrique figurent explicitement dans
cette conclusion comme compositions incorporées ; l'identification
gravitationnelle conserve son statut physique. La preuve couvre la classe
déclarée avec ses définitions constitutives et ses transformations
compatibles. Les vérifications rationnelles du développement vérifient
des cas finis des identités exposées ; la démonstration générale
réside dans les équations et les arguments précédents.
% END CR28_RADIAL_FULL


% BEGIN CR28_RADIAL_CONSTITUTIVE_EVALUATION
\section{Continuité constitutive et évaluation de la section gravitationnelle}
\label{cr28:sec:radial-evaluation}

La construction longitudinale détermine \(L\) avant d'utiliser l'identification gravitationnelle. La composition des règles R1--R8 conserve la section de vitesse de \eqref{cr28:eq:same-velocity} et l'action de retour préalablement engendrée. Avec
\[
 L=\frac{5759}{23040}\alpha^{16}L_*,
 \qquad \hbar_{\rm ret}=\eta_{\rm ret}S_*,
 \qquad c_{\rm int}=c_*\widehat c,
\]
le résultat de \eqref{g26:eq:g-rigido-es} s'exprime comme
\begin{equation}
 G_{\rm ret}
 =\frac{c_*^3L_*^2}{S_*}
   \left(\frac{5759}{23040}\right)^2
   \frac{\alpha^{32}\widehat c^{\,3}}{\eta_{\rm ret}}
 =\frac{L^2}
 {\hbar_{\rm ret}(\mu_{\rm vac}\varepsilon_{\rm vac})^{3/2}}.
 \label{cr28:eq:g-constitutive-evaluation}
\end{equation}
Les canaux sont les mêmes que ceux de \eqref{eq:vacuum-angular-channels} :
\[
 q_\pm=\exp\!\left[-\frac{\pi}{180}(A_{\rm deg}\pm C^*_{\rm deg})\right],
 \qquad r_\pm=\frac{1-q_\pm^{90}}{1-q_\pm^{120}},
 \qquad \widehat c=(r_+r_-)^{-1}.
\]
Les bases constitutives satisfont \(\varepsilon=\varepsilon_*r_+^2\), \(\mu=\mu_*r_-^2\) et \(c_*=(\mu_*\varepsilon_*)^{-1/2}\). Leur produit prouve \(c_{\rm int}=(\mu\varepsilon)^{-1/2}\). Si la coordonnée dimensionnelle de \(c_{\rm int}\) est déclarée directement, sa base se retrouve comme \(c_*=c_{\rm int}/\widehat c\). Le coefficient constitutif n'est pas multiplié de nouveau : les deux expressions représentent la même section.

Dans la représentation \(L_*=1\,\mathrm m\), \(S_*=10^{-34}\,\mathrm{J\,s}\) et \(c_{\rm int}=299792458\,\mathrm{m\,s^{-1}}\), l'évaluation des lecteurs et de la composition radiale donne
\begin{equation}
 \boxed{
 G_{\rm ret}
 =6.674299659414022706367776189\ldots\times10^{-11}
 \,\mathrm{m^3\,kg^{-1}\,s^{-2}}.}
 \label{cr28:eq:g-evaluation}
\end{equation}
La valeur de comparaison \(G_{\rm CODATA\,2022}=6.67430(15)\times10^{-11}\,
\mathrm{m^3\,kg^{-1}\,s^{-2}}\) conserve son incertitude expérimentale. Dans les mêmes unités, l'écart de l'évaluation à la valeur centrale est \(-3.40585977\ldots\times10^{-18}\), soit environ \(-0.00227057\) incertitudes-types. Les chiffres de l'évaluation correspondent aux lecteurs et aux bases déclarés ; l'incertitude appartient au résultat mesuré. Aucun chiffre de la référence ne sélectionne \(N\), \(r_N\), la règle longitudinale ou la norme radiale.

\subsection{Réalisation chronologique du même rayon construit}
\label{cr28:sec:radial-clock} Avec la durée élémentaire
\begin{equation}
 t_0=\frac{\pi L}{54c_{\rm int}},
 \qquad \ell_0=c_{\rm int}t_0=\frac{\pi L}{54},
 \label{cr28:eq:clock-from-length}
\end{equation}
le calendrier de \(108\) positions satisfait
\[
 T_{108}=108t_0=\frac{2\pi L}{c_{\rm int}},\qquad
 \omega_{108}=\frac{c_{\rm int}}{L},\qquad
 R_{108}=\frac{c_{\rm int}}{\omega_{108}}=L.
\]
Par substitution,
\begin{equation}
 \left(\frac{54}{\pi}\right)^2
 \frac{c_{\rm int}^5t_0^2}{\hbar_a}
 =\frac{c_{\rm int}^3L^2}{\hbar_a}.
 \label{cr28:eq:circular-radial-agreement}
\end{equation}
Cette identité place le sélecteur circulaire conservé ci-dessous dans la même normalisation longitudinale. Son théorème d'unicité conserve la prémisse géométrique qu'il énonce ; la construction précédente spécifie en outre la longueur utilisée dans cette réalisation.

Le caractère positif \(Y=I\) donne simultanément \(R_X=\Lambda_X=L\) et \(M_X=\hbar_a/(c_{\rm int}L)\). Pour un caractère général \(Y>0\), les longueurs sont réciproques : \(R_X=LY\), \(\Lambda_X=LY^{-1}\). Leur produit demeure \(L^2I\) sans que leurs valeurs individuelles doivent coïncider. L'identification des quatre longueurs du cas circulaire correspond donc à un secteur précis de la preuve radiale et ne supprime pas la dépendance au caractère.

La proportionnalité \(R_X=(G/c_{\rm int}^{2})M_X\) est commune à tous les caractères du domaine. Sa conservation sous changement de base, raffinement, minimisation avec mémoire, complément de Schur et normalisation dynamique a été démontrée dans \cref{g26:sec:rigidez-radial-es}. Le porteur à cinq composantes et sa réduction à deux hélicités, développés ensuite, décrivent la représentation tensorielle ; ils ne remplacent pas la démonstration du coefficient radial commun.
% END CR28_RADIAL_CONSTITUTIVE_EVALUATION


\section{Chaîne constitutive du secteur gravitationnel : échelle chronologique \(108\), réduction tensorielle, holonomie et évaluation métrologique de \(G\)}

La gravité s'organise ici en quatre strates qu'il ne faut pas fusionner :

\begin{equation}
\label{eq:gravity-causal-chain}
\begin{aligned}
\text{inscription et norme radiale R1--R8}
&\longrightarrow L
\longrightarrow G=c_{\rm int}^3L^2/\hbar_a,\\
\text{incidence icosaédrique}
&\longrightarrow V_{12}\xrightarrow{P_5}V_5
\xrightarrow{\Gamma_5}\operatorname{STF}_3
\xrightarrow{\Lambda(n)}\mathcal H_{\rm TT}(n),\\
\text{calendrier CT108 avec }t_0=\pi L/(54c_{\rm int})
&\longrightarrow R_{108}=L,\quad
\theta_{108}^{\rm circ}=(54/\pi)^2,\\
\text{carte métrologique}
&\longrightarrow 6.6742996594140227\ldots\times10^{-11}\,
\mathrm{m^3\,kg^{-1}\,s^{-2}}.
\end{aligned}
\end{equation}

La première ligne est la détermination démontrée dans \cref{g26:sec:rigidez-radial-es} : la longueur précède le couplage. La deuxième construit le porteur et sa réduction transverse sans trace. La troisième conserve le calendrier comme réalisation de cette même longueur, selon \eqref{cr28:eq:circular-radial-agreement}. La quatrième évalue la section dans les bases de \cref{cr28:sec:radial-evaluation}. Les branches alternatives d'holonomie et le lecteur décimal sont conservés avec leurs domaines propres ; ils ne remplacent pas la preuve radiale.

\section{Droite de grandeur gravitationnelle et échelle chronologique \(108\)}

Soient \(c_{\rm int}>0\), \(t_0>0\),
\(\ell_0=c_{\rm int}t_0\) et une section positive
\(\hbar_a\) de la droite d'action, avec
\(a\in\{\mathrm{pre},\mathrm{ret}\}\). Le transducteur dimensionnel est

\begin{equation}
\label{eq:gravity-transducer}
G_a(\theta)
=\theta\frac{c_{\rm int}^{5}t_0^2}{\hbar_a}
=\theta\frac{c_{\rm int}^{3}\ell_0^2}{\hbar_a},
\qquad \theta>0.
\end{equation}

L'égalité des deux expressions utilise seulement \(\ell_0=c_{\rm int}t_0\) et \([G]=L^3M^{-1}T^{-2}\). La dimension est fixée par la droite ; le caractère adimensionnel \(\theta\) distingue les réalisations de ce système. Dans la composition longitudinale déjà démontrée, \(t_0=\pi L/(54c_{\rm int})\) et \(\theta=(54/\pi)^2\) reproduisent exactement \(G_a=c_{\rm int}^3L^2/\hbar_a\).

Le calendrier CT108 définit

\begin{equation}
\label{eq:gravity-clock108}
T_{108}=108t_0,
\qquad
\omega_{108}=\frac{2\pi}{108t_0},
\qquad
R_{108}=\frac{c_{\rm int}}{\omega_{108}}
=\frac{54}{\pi}\ell_0.
\end{equation}

Pour chaque feuillet d'action, soient

\begin{equation}
\label{eq:gravity-clock-energy-mass}
E_{108,a}=\hbar_a\omega_{108},
\qquad
m_{108,a}=\frac{E_{108,a}}{c_{\rm int}^2}.
\end{equation}

Alors, la longueur de Compton réduite coïncide avec le rayon du cycle :

\begin{equation}
\label{eq:gravity-compton-clock}
\bar\lambda_{C,a}
=\frac{\hbar_a}{m_{108,a}c_{\rm int}}
=R_{108}.
\end{equation}

\begin{theorem}[Unicité du sélecteur sous la condition de périodicité HMT]
\label{thm:gravity-circular-selector}
Avec la convention réduite
\(r_g=Gm/c_{\rm int}^2\) et la prémisse de périodicité HMT qui identifie le rayon
gravitationnel du quantum CT108 à son rayon de phase, la condition de clôture
\begin{equation}
\label{eq:gravity-closing-condition}
r_{g,a}(\theta)=R_{108}
\end{equation}
sélectionne une unique solution positive,
\begin{equation}
\label{eq:gravity-theta108}
\boxed{
\theta_{108}^{\rm circ}=\left(\frac{54}{\pi}\right)^2
=295.4525715010569415303525\ldots .}
\end{equation}
Dans cette réalisation,
\begin{equation}
\label{eq:gravity-four-lengths}
R_{108}=\bar\lambda_{C,a}=r_{g,a}=\ell_P^{\rm circ}.
\end{equation}
\end{theorem}

\begin{proof}
Par \cref{eq:gravity-transducer,eq:gravity-clock-energy-mass},
\[
r_{g,a}(\theta)
=\frac{G_a(\theta)m_{108,a}}{c_{\rm int}^2}
=\theta c_{\rm int}t_0^2\omega_{108}
=\theta\frac{\pi}{54}\ell_0.
\]
En le comparant à
\(R_{108}=(54/\pi)\ell_0\), on obtient nécessairement
\(\theta=(54/\pi)^2\). L'équation est linéaire en \(\theta\) et tous les
facteurs sont positifs, donc la solution est unique.
\end{proof}

L'égalité \(r_g=R_{108}\) ne se déduit pas de la dimensionnalité de
\cref{eq:gravity-transducer}. C'est la condition géométrique explicite de cette
réalisation HMT : action, Compton et retour périodique doivent déterminer une même
longueur. Le théorème démontre que, cette condition structurelle étant acceptée, le
caractère \(\theta\) est fixé sans utiliser la valeur SI de \(G\).

\begin{remark}[Convention de rayon]
Le théorème utilise le rayon gravitationnel réduit \(Gm/c^2\). Si on le remplace
par le rayon de Schwarzschild \(2Gm/c^2\), le sélecteur change d'un facteur
deux. Les deux expressions correspondent à des conditions géométriques distinctes,
c'est pourquoi la convention doit être déclarée avant le calcul.
\end{remark}

\section{Réciprocité \(G\leftrightarrow\hbar\) et conservation de l'aire de Planck}

\begin{corollary}[Famille de Planck associée à la périodicité temporelle de \(108\) phases]
\label{cor:gravity-planck-family}
Pour \(G_a=G_a(\theta_{108}^{\rm circ})\),
\begin{align}
\ell_P&=\frac{54}{\pi}\ell_0,
&t_P&=\frac{54}{\pi}t_0,
\label{eq:gravity-planck-length-time}\\
E_{P,a}&=\frac{\hbar_a}{t_P}
=\frac{\pi\hbar_a}{54t_0}
=\frac{h_a}{108t_0},
&\ell_P^2&=\frac{G_a\hbar_a}{c_{\rm int}^3}
=\theta_{108}^{\rm circ}\ell_0^2.
\label{eq:gravity-planck-energy-area}
\end{align}
De plus,
\begin{equation}
\label{eq:gravity-planck-force-power-density}
F_{P,a}=\frac{c_{\rm int}^4}{G_a},
\qquad
P_{P,a}=\frac{c_{\rm int}^5}{G_a},
\qquad
\rho_{P,a}=\frac{\hbar_a}
{(\theta_{108}^{\rm circ})^2c_{\rm int}^5t_0^4}.
\end{equation}
\end{corollary}

\begin{proof}
Les trois premières identités découlent de
\(\ell_P=R_{108}\), \(t_P=\ell_P/c_{\rm int}\) et
\(h_a=2\pi\hbar_a\). Pour l'aire,
\[
\frac{G_a\hbar_a}{c_{\rm int}^3}
=\theta_{108}^{\rm circ}c_{\rm int}^2t_0^2
=\theta_{108}^{\rm circ}\ell_0^2.
\]
Les expressions restantes sont les compositions dimensionnelles de force,
de puissance et de masse par volume avec la même section.
\end{proof}

Soit
\begin{equation}
\label{eq:gravity-Ract}
R_{\rm act}:=\frac{\hbar_{\rm pre}}{\hbar_{\rm ret}}.
\end{equation}
Comme l'action apparaît au numérateur de la masse chronologique et au
dénominateur de \(G\),
\begin{equation}
\label{eq:gravity-twoway-covariance}
\frac{m_{108,\rm pre}}{m_{108,\rm ret}}=R_{\rm act},
\qquad
\frac{G_{\rm pre}}{G_{\rm ret}}=R_{\rm act}^{-1},
\qquad
G_a\hbar_a
=\theta_{108}^{\rm circ}c_{\rm int}^5t_0^2.
\end{equation}
Par conséquent, \(\ell_P^2=G_a\hbar_a/c^3\) ne change pas de feuillet. La gravité et l'action s'orientent réciproquement, tandis que l'aire conserve la géométrie commune. Cette annulation constitue l'interface de convergence vers l'opérateur d'aire de \cref{ch:modular}. La seconde relation fondamentale apparaît lorsque l'on compare action et gravité.\par

Dans HMT, l’action mesure combien il se produit et conserve simultanément l’orientation de la réalisation. Il existe une section antérieure au retour et une autre postérieure. La phase peut revenir accompagnée d’une histoire de mémoire différente. Cette différence est stockée dans le rapport bidirectionnel des deux sections d’action.\par

La gravité répond de manière réciproque.\par

Le transducteur gravitationnel a la forme\par

\[
G_a(\theta)
=
\theta\frac{c_{\rm int}^{5}t_0^2}{\hbar_a}.
\]

L’action apparaît au dénominateur. Par conséquent, si, lors du changement de feuillet, \(\hbar\) est multiplié par un rapport \(R_{\rm act}\), la constante gravitationnelle est multipliée par le rapport inverse :\par

\[
\frac{\hbar_{\rm pre}}{\hbar_{\rm ret}}
=
R_{\rm act},
\qquad
\frac{G_{\rm pre}}{G_{\rm ret}}
=
R_{\rm act}^{-1}.
\]

L’action et la gravité forment ici une paire contragrédiente : lorsque l’une avance dans une direction, l’autre compense exactement ce déplacement.\par

Ce qui reste invariant est\par

\[
G_a\hbar_a.
\]

Et c’est précisément la combinaison qui détermine l’aire de Planck :\par

\[
\ell_P^2=\frac{G_a\hbar_a}{c_{\rm int}^3}.
\]

L’action mémorise le feuillet. La gravité répond à cette mémoire. L’aire reste invariante.\par

C’est une idée extraordinairement féconde. L’aire de Planck émerge comme l’objet qui subsiste lorsque l’action et la gravité se transforment réciproquement.\par

On pourrait le dire ainsi : \(\hbar\) conserve l’orientation dynamique ; \(G\) conserve la compensation géométrique ; \(\ell_P^2\) conserve ce sur quoi les deux réalisations s’accordent.\par

L’aire est donc une monnaie d’échange entre physique quantique et géométrie, parce que la transformation bidirectionnelle de \(\hbar\) et de \(G\) sélectionne exactement cette combinaison comme invariant.\par

Cette réciprocité permet de comprendre autrement la relation entre information et gravité. L’information exige une distinction entre les états ; l’action mesure le coût de la transition entre eux. Mais si cette transition doit acquérir une géométrie, la réponse gravitationnelle change de manière inverse pour maintenir stable le quantum d’aire.\par

La géométrie incorpore l’information de phase par une compensation exacte.\par

Selon une lecture machienne, c’est déjà un résultat très suggestif. Une grandeur locale de géométrie est déterminée relationnellement par l’orientation et le retour de l’état complet. L’aire invariante restreint cette dépendance relationnelle et en fixe la cohérence.\par

La version HMT du principe relationnel consiste en quelque chose de plus précis que la maxime « la matière détermine la géométrie » : l’orientation de l’action et la réponse gravitationnelle se codéterminent de telle manière que la géométrie élémentaire reste invariante.\par

L’action contient l’histoire. La gravité ajuste la réponse. L’aire conserve l’accord.\par

 La dérivation de \(G\) porte cette réciprocité à une condition géométrique beaucoup plus concrète.\par

La mécanique dimensionnelle détermine d’abord quel type d’objet \(G\) peut être. Elle fixe sa droite de grandeur et son homogénéité dimensionnelle. La clôture chronologique sélectionne ensuite sa valeur : la dimension détermine la classe de grandeur et la clôture détermine son coefficient adimensionnel.\par

La sélection procède de la périodicité temporelle de (108) phases.\par

Le cycle \(108\) détermine une fréquence, un rayon de phase et une énergie. À partir de cette énergie apparaît une masse chronologique. Et en construisant la longueur de Compton réduite de cette masse, on obtient exactement le même rayon du cycle :\par

\[
\bar\lambda_C=R_{108}.
\]

La phase et la matière expriment déjà la même longueur avant l’introduction de \(G\).\par

La gravité intervient alors. On construit le rayon gravitationnel réduit\par

\[
r_g=\frac{Gm}{c^2},
\]

et l’on exige que la clôture gravitationnelle identifie le rayon de phase, la longueur de Compton et le rayon gravitationnel comme une seule échelle structurelle :\par

\[
r_g=R_{108}=\bar\lambda_C.
\]

Cette condition sélectionne de manière univoque\par

\[
\theta_{108}
=
\left(\frac{54}{\pi}\right)^2.
\]

Ainsi,\par

\[
R_{108}
=
\bar\lambda_C
=
r_g
=
\ell_P.
\]

La fonction structurelle de \(\theta_{108}\) détermine la signification de son nombre décimal terminal : elle montre ce que fait \(G\).\par

\(G\) apparaît comme le coefficient nécessaire pour faire coïncider les quatre longueurs publiées par ces langages :\par

\begin{itemize}[leftmargin=*,itemsep=0.35em]
\item le langage de la phase produit \(R_{108}\) ;
\item le langage quantique produit \(\bar\lambda_C\) ;
\item le langage gravitationnel produit \(r_g\) ;
\item le langage de Planck produit \(\ell_P\).
\end{itemize}

La constante gravitationnelle est le médiateur qui empêche ces quatre lecteurs de se séparer.\par

Dit de manière plus physique : \(G\) exprime combien la géométrie doit répondre pour qu’un quantum dont la masse provient de l’horloge \(108\) se referme sur sa propre phase.\par

C’est la relation qui rend possible que matière, phase et géométrie soient trois publications compatibles du même événement.\par

Ici, la lecture machienne devient plus forte. L’échelle gravitationnelle locale est sélectionnée par une condition de clôture globale du cycle. Le rayon local d’un quantum acquiert un sens en s’intégrant dans une périodicité complète. L’inertie, la phase et la géométrie se déterminent relationnellement.\par

La portée démontrée fournit un mécanisme mathématique concret pour une lecture machienne : une propriété locale — le couplage gravitationnel — est sélectionnée par la cohérence de l’état global. Les autres formulations historiques du principe de Mach conservent leur domaine propre.\par

Ensuite apparaît la dimension cinq.\par

Le scalaire \(G\) fixe l’intensité gravitationnelle et le porteur spécifie sa structure représentationnelle. HMT construit un espace tensoriel symétrique de trace nulle et de dimension cinq. Ce porteur admet cinq poids hélicoïdaux :\par

\[
+2,\quad +1,\quad 0,\quad -1,\quad -2.
\]

Ce sont les cinq degrés de liberté naturels d’un tenseur symétrique sans trace en trois dimensions avant l’imposition des contraintes dynamiques d’une théorie particulière.\par

Cela clarifie un point souvent confondu. « Cinq composantes » désigne le porteur complet, tandis que les deux hélicités désignent la réduction transverse de masse nulle en relativité générale.\par

Cela signifie que le porteur complet possède cinq coordonnées irréductibles. Ensuite, la dynamique décide quelle partie de ce porteur est physique :\par

\begin{itemize}[leftmargin=*,itemsep=0.35em]
\item dans une réalisation massive de spin deux, les cinq poids peuvent subsister ;
\item dans une réalisation sans masse avec invariance de jauge, la projection transverse sans trace élimine les secteurs \(\pm1\) et \(0\) ;
\item les hélicités \(\pm2\) subsistent.
\end{itemize}

Ainsi, HMT fournit un espace antérieur à la décision dynamique. Elle construit le porteur complet à cinq composantes et contient, par la projection correspondante, sa réduction physique à deux hélicités.\par

La beauté réside dans la coexistence ordonnée de deux affirmations :\par

\begin{itemize}[leftmargin=*,itemsep=0.35em]
\item la géométrie de spin deux possède un porteur à cinq composantes ;
\item le rayonnement gravitationnel sans masse occupe un sous-espace bidimensionnel.
\end{itemize}

Les deux affirmations forment une réduction cohérente et précise.\par

Et cette réduction conserve un autre enseignement central : le scalaire \(G\) fixe l’intensité commune et l’information d’orientation réside dans le porteur tensoriel. L’intensité réside dans la constante, les hélicités résident dans la représentation et la généalogie complète réunit grandeur et porteur.\par

  \endgroup
% Secciones 1--3 del delta gravitatorio íntegro.
\section{Domaine généalogique de la réalisation gravitationnelle}
\label{sec:l9s102:gravedad-dominio-genealogico}

La réalisation gravitationnelle reçoit comme objet de départ une section du
continu discret déjà construite par la composition
\begin{equation}
 \mathrm{APP}\longrightarrow\mathrm{TRIT}\longrightarrow\mathrm{TPK}
 \longrightarrow X_{\mathrm{TPK}}^{\mathrm{enr}}
 \longrightarrow\mathcal C_{\mathrm{cont}}^{\mathrm{disc}}.
 \label{eq:l9s102:gravedad-cadena-genealogica}
\end{equation}
APP conserve les évaluations additive et multiplicative au moyen du résidu,
du quotient et de la retenue ; TRIT détermine le régime et l'orientation ; le TPK
incorpore le sélecteur tritique de phase sans éliminer aucun des deux feuillets
APP, effectue le transport cardinal, conserve la mémoire et réalise le retour
de phase sans réinitialiser l'état. Le transducteur
dimensionnel agit ensuite sur la section survivante du continu :
\begin{equation}
 \mathfrak G_{108}:\mathcal C_{\mathrm{cont}}^{\mathrm{disc}}
 \times\mathcal L_S^+
 \longrightarrow\mathcal G_{\mathrm{HMT\text{-}MD}},
 \qquad
 (x,\hbar_a)\longmapsto
 \bigl(G_a,R_{108},m_{108,a},\ell_{P,a}^{,2}\bigr).
 \label{eq:l9s102:gravedad-mapa-realizacion}
\end{equation}
Les cardinaux \(90/120\), la périodicité \(108\), l'action et le couple
angulaire entrent avec leur généalogie conservée. Aucune identification
gravitationnelle ultérieure ne sélectionne rétroactivement ces antécédents.

\section{Transducteur gravitationnel et sélecteur chronologique de périodicité}
\label{sec:l9s102:gravedad-transductor}

Soient \(c_{\mathrm{int}}>0\), \(t_0>0\),
\(\ell_0=c_{\mathrm{int}}t_0\) et
\(\hbar_a\in\mathcal L_S^+\), avec
\(a\in\{\mathrm{pre},\mathrm{ret}\}\). Le transducteur dimensionnel est
\begin{equation}
 G_a(\theta)
 =\theta\frac{c_{\mathrm{int}}^5t_0^2}{\hbar_a}
 =\theta\frac{c_{\mathrm{int}}^3\ell_0^2}{\hbar_a},
 \qquad \theta>0.
 \label{eq:l9s102:g-transductor}
\end{equation}
La chronologie cyclique de \(108\) états détermine
\begin{equation}
 T_{108}=108t_0,
 \qquad
 \omega_{108}=\frac{2\pi}{108t_0},
 \qquad
 R_{108}=\frac{c_{\mathrm{int}}}{\omega_{108}}
 =\frac{54}{\pi}\ell_0.
 \label{eq:l9s102:g-reloj108}
\end{equation}
Le quantum chronologique et sa masse associée sont
\begin{equation}
 E_{108,a}=\hbar_a\omega_{108},
 \qquad
 m_{108,a}=\frac{E_{108,a}}{c_{\mathrm{int}}^2}.
 \label{eq:l9s102:g-cuanto-masa}
\end{equation}
La longueur de Compton réduite satisfait
\begin{equation}
 \bar\lambda_{C,a}
 =\frac{\hbar_a}{m_{108,a}c_{\mathrm{int}}}=R_{108}.
 \label{eq:l9s102:g-compton}
\end{equation}

\begin{theorem}[Unicité du sélecteur gravitationnel de périodicité]
\label{thm:l9s102:g-selector}
Avec la convention réduite \(r_g=Gm/c_{\mathrm{int}}^2\), la condition
\(r_{g,a}(\theta)=R_{108}\) possède l'unique solution positive
\begin{equation}
 \boxed{
 \theta_{108}^{\mathrm{circ}}=\left(\frac{54}{\pi}\right)^2.}
 \label{eq:l9s102:g-theta}
\end{equation}
La réalisation correspondante identifie structurellement
\begin{equation}
 R_{108}=\bar\lambda_{C,a}=r_{g,a}=\ell_P^{\mathrm{circ}}.
 \label{eq:l9s102:g-cuatro-longitudes}
\end{equation}
\end{theorem}

\begin{proof}
Par \cref{eq:l9s102:g-transductor,eq:l9s102:g-cuanto-masa},
\[
 r_{g,a}(\theta)
 =\frac{G_a(\theta)m_{108,a}}{c_{\mathrm{int}}^2}
 =\theta\frac{\pi}{54}\ell_0.
\]
L'égalité avec \(R_{108}=(54/\pi)\ell_0\) impose
\(\theta=(54/\pi)^2\). La linéarité en \(\theta\) et la positivité des
facteurs déterminent l'unicité.
\end{proof}

Dans la carte où \(c_{\mathrm{int}}=(\mu\varepsilon)^{-1/2}\), le
transducteur adopte la forme
\begin{equation}
 \boxed{
 G_a=\left(\frac{54}{\pi}\right)^2
 \frac{t_0^2}{\hbar_a(\mu\varepsilon)^{5/2}}.}
 \label{eq:l9s102:g-mu-epsilon}
\end{equation}
Le produit constitutif \(\mu\varepsilon\) gouverne le cône causal et
l'échelle gravitationnelle ; le quotient \(Z^2=\mu/\varepsilon\) conserve
l'orientation d'impédance.

\section{Covariance bidirectionnelle de la gravité et de l'action}
\label{sec:l9s102:g-hbar}

On définit
\begin{equation}
 R_{\mathrm{act}}=\frac{\hbar_{\mathrm{pre}}}
 {\hbar_{\mathrm{ret}}}.
 \label{eq:l9s102:r-act}
\end{equation}
La masse chronologique et le couplage gravitationnel se transforment
de manière contragrédiente :
\begin{equation}
 \frac{m_{108,\mathrm{pre}}}{m_{108,\mathrm{ret}}}=R_{\mathrm{act}},
 \qquad
 \frac{G_{\mathrm{pre}}}{G_{\mathrm{ret}}}=R_{\mathrm{act}}^{-1}.
 \label{eq:l9s102:g-hbar-reciprocidad}
\end{equation}
Leur produit demeure fixe :
\begin{equation}
 \boxed{
 G_a\hbar_a
 =\theta_{108}^{\mathrm{circ}}c_{\mathrm{int}}^5t_0^2,
 \qquad
 \ell_P^2=\frac{G_a\hbar_a}{c_{\mathrm{int}}^3}
 =\theta_{108}^{\mathrm{circ}}\ell_0^2.}
 \label{eq:l9s102:g-hbar-invariante}
\end{equation}
L'action conserve l'orientation du feuillet ; l'aire de Planck conserve la
géométrie commune aux deux orientations.

\begin{theorem}[Indépendance de carte des sections \(h\) et \(G\)]
\label{thm:l9s102:covariancia-unidades-g}
Soient \(u_L,u_T,u_S\) des bases positives des droites dimensionnelles de
longueur, de temps et d'action, et remplaçons-les par
\[
 u_L'=\lambda_Lu_L,\qquad
 u_T'=\lambda_Tu_T,\qquad
 u_S'=\lambda_Su_S,
 \qquad \lambda_L,\lambda_T,\lambda_S>0.
\]
Les coordonnées des mêmes sections vérifient
\begin{equation}
 c_{\rm int}'=\frac{\lambda_T}{\lambda_L}c_{\rm int},\qquad
 t_0'=\frac{t_0}{\lambda_T},\qquad
 \hbar_a'=\frac{\hbar_a}{\lambda_S},\qquad
 G_a'=\frac{\lambda_S\lambda_T^3}{\lambda_L^5}G_a.
 \label{eq:l9s102:cambio-carta-g}
\end{equation}
Pour la base induite
\(u_G=u_L^5u_T^{-3}u_S^{-1}\), on obtient
\begin{equation}
 \hbar_a'u_S'=\hbar_au_S,
 \qquad G_a'u_G'=G_au_G.
 \label{eq:l9s102:secciones-h-g-invariantes}
\end{equation}
Par conséquent,
\(h_a=2\pi\hbar_a\),
\(G_a=\theta_{108}^{\rm circ}c_{\rm int}^5t_0^2/\hbar_a\) et
\(G_a\hbar_a=\theta_{108}^{\rm circ}c_{\rm int}^5t_0^2\)
sont des identités de sections, non des égalités dépendant d'un choix
d'unités. La carte métrologique publie leurs coordonnées numériques après
la construction HMT--MD.
\end{theorem}

\begin{proof}
Les trois premières transformations de
\eqref{eq:l9s102:cambio-carta-g} sont la loi contragrédiente des coordonnées
dans \(L\otimes T^{-1}\), \(T\) et \(S\). Leur substitution dans
\eqref{eq:l9s102:g-transductor} produit la quatrième. La base de la droite
gravitationnelle se transforme par le facteur inverse, ce qui prouve
\eqref{eq:l9s102:secciones-h-g-invariantes}. Les scalaires
\(\pi\) et \(\theta_{108}^{\rm circ}\) sont adimensionnels et
demeurent invariants.
\end{proof}
 
\chapter{Paramètre de Barbero--Immirzi, incidence hexadique--octadique et spectre de l'opérateur d'aire}
\begingroup
\renewcommand{\chapter}[2][]{}%
\chapter{Dérivation HMT du paramètre de Barbero--Immirzi et de l’aire minimale : incidence hexade--octade, canaux \(90/120\) et spectre de l’opérateur d’aire}
\label{chap:p3-final:51:barbero_area}
\section{Morphisme d’incidence hexade--octade et conservation des canaux orientés \(90/120\)}

La dodécade et l’alignement du
\cref{sec:w12-w24-elevacion} étant fixés, soit \(H\) une hexade dans les coordonnées HMT,
soit \(\ell(H)\subset D\) son image et soit
\(O_H=\iota_D(\ell(H))\) son unique relèvement. En marquant un point
et une paire de points de l’hexade résiduelle, on obtient
\[
F_6(H)=\ell(H)\times\binom{\ell(H)}2,
\qquad
|F_6|=6\binom62=90.
\]
Si le point marqué peut parcourir l’octade :
\[
F_8(H)=O_H\times\binom{\ell(H)}2,
\qquad
|F_8|=8\binom62=120.
\]
L’inclusion \(\ell(H)\hookrightarrow O_H\) induit le morphisme d’incidence
\[
 \mathcal I_{6\to8}:F_6(H)\hookrightarrow F_8(H).
\]
En normalisant la différence par le facteur
choisi \(24\binom62\), on obtient
\[
\boxed{
\frac{90-120}{24\binom62}
=\frac{6-8}{24}
=-\frac1{12}.}
\]
L'égalité est un défaut normalisé de l'interface combinatoire \(6\to8\). Le facteur \(24\binom62\) fait partie de la définition de cette normalisation ; l'inclusion seule détermine la différence \(-30\). L'opérateur \(\mathcal I_{6\to8}\) a pour domaine les configurations marquées d'incidence ; ce n'est ni l'extracteur transverse de l'état dodécaphasique ni le système des moments angulaires. \label{mdg:chap:barbero}

La transition d’une hexade à une octade combine trois structures de types
différents : une incidence finie, deux séries de Lambert associées aux
canaux \(q_\pm\) et une identification au paramètre réel de l’action de
Holst. Les trois s’articulent par l’incidence, la fonctionnelle angulaire et
le dictionnaire aréal vers la géométrie de la connexion
d’Ashtekar--Barbero développé dans cette même résidence.

\section{Cardinalités d’incidence 90 et 120}

Soit \(H\) un ensemble de six points et soit \(O\) un ensemble de huit points
contenant une copie distinguée de \(H\). Considérons
\[
 \mathcal F_6=H\times\binom H2,
 \qquad
 \mathcal F_8=O\times\binom H2.
\]
L’inclusion distinguée \(H\hookrightarrow O\) induit
\[
 \iota_{6,8}:\mathcal F_6\hookrightarrow\mathcal F_8,
 \qquad
 (h,\{h_1,h_2\})\longmapsto(h,\{h_1,h_2\}).
\]
Ce morphisme d’incidences est l’antécédent combinatoire des
cardinaux \(90/120\) ; ce n’est ni l’extracteur dodécaphasique ni une série angulaire.

\begin{theorem}[Cardinalités de l’inclusion]
\[
 |\mathcal F_6|=6\binom62=90,
 \qquad
 |\mathcal F_8|=8\binom62=120.
\]
La différence normalisée par vingt-quatre coordonnées est
\[
 \frac{90-120}{24\binom62}=-\frac1{12}.
\]
\end{theorem}
\begin{proof}
Chaque élément de la première coordonnée peut être combiné avec n’importe laquelle des
quinze paires de \(H\). Le principe multiplicatif donne les deux cardinalités ;
la dernière identité résulte de la substitution de \(\binom62=15\).
\end{proof}

\begin{figure}[htbp]
\centering
\begin{tikzpicture}[
  font=\small,
  setnode/.style={circle,draw,minimum size=6mm,inner sep=0pt,
    font=\scriptsize\bfseries},
  box/.style={draw,rounded corners=1.5mm,align=center,
    text width=5.0cm,inner xsep=6pt,inner ysep=5pt},
  arr/.style={-{Stealth[length=2.2mm]},line width=.7pt}]

  % Inclusión geométrica. Las flechas parten del contorno y no atraviesan nodos.
  \node[setnode,draw=hmtblue,fill=hmtlightblue] (h1) at (-2.3,1.25) {1};
  \node[setnode,draw=hmtblue,fill=hmtlightblue] (h2) at (-1.1,1.9) {2};
  \node[setnode,draw=hmtblue,fill=hmtlightblue] (h3) at (.1,1.25) {3};
  \node[setnode,draw=hmtblue,fill=hmtlightblue] (h4) at (.1,-.15) {4};
  \node[setnode,draw=hmtblue,fill=hmtlightblue] (h5) at (-1.1,-.8) {5};
  \node[setnode,draw=hmtblue,fill=hmtlightblue] (h6) at (-2.3,-.15) {6};
  \node[setnode,draw=hmtgreen,fill=hmtlightgreen] (o7) at (-3.2,.55) {7};
  \node[setnode,draw=hmtgreen,fill=hmtlightgreen] (o8) at (1.0,.55) {8};
  \node[draw=hmtblue,dashed,rounded corners=4mm,
    fit=(h1)(h2)(h3)(h4)(h5)(h6),inner sep=4mm,
    label={[hmtblue,font=\bfseries]above:{hexada \(H\), \(|H|=6\)}}] (H) {};
  \node[draw=hmtgreen,line width=.8pt,rounded corners=6mm,
    fit=(H)(o7)(o8),inner sep=5mm,
    label={[hmtgreen,font=\bfseries]below:{octada \(O\), \(|O|=8\), \(H\subset O\)}}] (O) {};

  \node[box,draw=hmtblue,fill=hmtlightblue] (f6) at (5.3,1.35)
    {Incidences définies sur \(H\)\\[1mm] \(\mathcal F_6=H\times\binom H2\),\quad \(|\mathcal F_6|=6\binom62=90\)};
  \node[box,draw=hmtgreen,fill=hmtlightgreen] (f8) at (5.3,-1.0)
    {Extensión de incidencias a \(O\)\\[1mm]
     \(\mathcal F_8=O\times\binom H2\),\quad
     \(|\mathcal F_8|=8\binom62=120\)};
  \draw[arr,hmtblue] (O.east) -- ++(.45,0) |- (f6.west);
  \draw[arr,hmtgreen] (O.east) -- ++(.45,0) |- (f8.west);
  \draw[arr,hmtgold] (f6.south) -- node[right,align=left,font=\footnotesize]
    {inclusión \(\iota_{6,8}\)} (f8.north);

  \node[box,draw=hmtgold,fill=hmtlightgold,text width=10.9cm]
    (def) at (1.45,-3.80)
    {Defecto combinatorio orientado\\[-.2mm]
     con la normalización declarada:\\[1mm]
     \(\displaystyle
       \frac{|\mathcal F_6|-|\mathcal F_8|}{24\binom62}
       =\frac{90-120}{24\cdot15}=-\frac1{12}.\)};
  \draw[arr,hmtgold] (f8.south) -- (f8.south |- def.north);

  \node[align=center,text=hmtgray,font=\footnotesize,text width=11.2cm]
    at (1.45,-5.50)
    {Esta identidad pertenece al mapa de incidencias. No identifica el
     extractor dodecafásico, el semigrupo de torres ni el funcional angular.};
\end{tikzpicture}
 \caption{Inclusión de un subconjunto distinguido de seis elementos en uno de
ocho elementos y conteos de incidencia. Los cardinales \(90\) y \(120\) no fusionan el extractor
dodecafásico, el funcional angular ni el semigrupo de torres.}
\label{mdg:fig:incidencia-seis-ocho-rev9}
\end{figure}

Los índices iniciales de la jerarquía angular quedan así asociados a dos
espacios de incidencia: 90 configuraciones sobre la hexada y 120 sobre la
octada ambiente. El cociente \(-1/12\) es una identidad combinatoria de esta
normalización.

\section{Series de Lambert asociadas a los 2 sectores de incidencia}

Para \(0<q<1\), definimos simultáneamente
\[
 S_{90}(q)=\frac{q^{90}}{1-q^{270}},
 \qquad
 S_{120}(q)=\frac{q^{120}}{1-q^{360}},
 \qquad
 \Delta S_m=S_m(q_-)-S_m(q_+).
\]
Les développements géométriques
\[
 \Delta S_{90}=\sum_{j\geq0}s_{90+270j},
 \qquad
 \Delta S_{120}=\sum_{j\geq0}s_{120+360j}
\]
situent les deux termes dans le semi-groupe global du
\cref{mdg:chap:torres}.

La fonctionnelle angulaire associée à la transition \(6\to8\) est
\[
 K_{6\to8}=12\Delta S_{90}-\Delta S_{120},
\]
et ne se déduit pas des seules cardinalités d'incidence. Ses poids ne sont déterminés qu'après fixation des deux séries, du résidu de pôle \(1/24\), de la positivité entière et du critère de minimalité démontré ci-dessous. La coordonnée angulaire résultante est
\[
\Gamma_{6\to8}
=\frac{\pi A\Cstar}{180}+K_{6\to8}.
\]
Au point HMT,
\[
\begin{aligned}
 \Delta S_{90}&=2.95169439297457621139847987861\times10^{-4},\\
 \Delta S_{120}&=1.96835112502951720093738706217\times10^{-5},\\
 \frac{\pi A\Cstar}{180}&=0.270485322934140078465586458790\ldots,
\end{aligned}
\]
et, par conséquent,
\[
 \boxed{\Gamma_{6\to8}
 =0.27400767269445927474725526077398041940\ldots.}
\]
Les tableaux calculés avec la coordonnée angulaire conjuguée \(C^*\) arrondie à quinze décimales donnent
\(0.2740076726944593\) en arithmétique à double précision ; la valeur précédente
est l’évaluation multiprécision de la phase régionale complète.

\section{Équation diophantienne des poids}

L’identité
\[
 \lim_{q\to1^-}(1-q)\frac{q^m}{1-q^n}=\frac1n
\]
associe à la combinaison \(a\Delta S_{90}-b\Delta S_{120}\) le résidu normalisé
\[
 \frac a{270}-\frac b{360}.
\]

\begin{theorem}[Poids entiers de résidu \texorpdfstring{\(1/24\)}{1/24}]
Les solutions entières positives de
\[
 \frac a{270}-\frac b{360}=\frac1{24}
\]
sont
\[
 (a,b)=(12+3t,1+4t),\qquad t\in\ZZ_{\geq0}.
\]
La paire \((12,1)\) est l’unique solution minimale coordonnée par coordonnée et
l’unique de norme \(\ell^1\) minimale.
\end{theorem}
\begin{proof}
La multiplication par \(1080\) produit \(4a-3b=45\). Modulo trois, on obtient
\(a\equiv0\pmod3\) ; en écrivant \(a=3u\), il en résulte \(4u-b=15\). La
positivité implique \(u=4+t\), \(b=1+4t\), avec \(t\geq0\). Par conséquent,
\(a+b=13+7t\), ce qui démontre la minimalité.
\end{proof}

La primitivité \(\gcd(a,b)=1\) ne sélectionne pas à elle seule le couple minimal, car il existe d'autres solutions primitives. Le poids \(12:-1\) résulte de la combinaison de la condition de résidu \(1/24\) avec la minimalité. Le choix de ce résidu est une condition de normalisation de la construction ; il ne se déduit pas de l'évaluation décimale de \(\Gamma_{6\to8}\).

\newpage
\section{Monotonie de la fonctionnelle par rapport à la coordonnée angulaire conjuguée}

Para \(m\in\{90,120\}\),
\[
 S_m'(q)=\frac{m q^{m-1}(1+2q^{3m})}{(1-q^{3m})^2}>0.
\]
De plus,
\[
 \frac{\dd q_-}{\dd C}=\frac\pi{180}q_-,
 \qquad
 \frac{\dd q_+}{\dd C}=-\frac\pi{180}q_+.
\]

\begin{proposition}[Monotonie stricte]
Pour \(1^\circ\leq C\leq3^\circ\), la fonction
\(C\mapsto\Gamma_{6\to8}(A,C)\) est strictement croissante. Par conséquent,
toute valeur de son image possède au plus un antécédent dans cet intervalle.
\end{proposition}
\begin{proof}
La dérivée de chaque différence \(\Delta S_m\) est positive. Dans l’intervalle
indiqué, on obtient
\(\Delta S'_{120}<5.35\times10^{-4}\), tandis que
\(\pi A/180>0.1273\). Par conséquent,
\(\Gamma'(C)>0.1267\). En \(C=\Cstar\),
\[
 \Gamma'(\Cstar)=0.1328995105618907\ldots.
\]
\end{proof}

L’injectivité locale exclut une seconde coordonnée angulaire proche donnant la
même évaluation de la fonctionnelle. Ce contrôle ne construit ni ne sélectionne
\(C^*\) : la coordonnée angulaire conjuguée a déjà été obtenue dans la branche régionale précédente.
La proposition quantifie uniquement la sensibilité de l’évaluation descendante
une fois \(A\) et \(C^*\) fixés ; elle n’autorise pas à inverser
\(\Gamma_{6\to8}\) ni une valeur physique de \(\gamma_{\mathrm{BI}}\) pour
les redéfinir.

 \label{ch:modular}

\section{Inclusion hexade--octade et réalisations fonctionnelles associées}

Fixons un ensemble \(H\) de six points et un ensemble \(O\) de huit
points avec une inclusion distinguée \(j:H\hookrightarrow O\). Dans la
résidence de Witt, \(H\) est une hexade d'une dodécade et \(O\) son octade
relevée ; dans ce chapitre, on utilise seulement l'inclusion déjà fixée, et non la
valeur terminale de Barbero--Immirzi. Nous écrivons
\begin{equation}
\label{eq:modular-pairs-H}
\binom H2:=\{P\subset H:|P|=2\},
\qquad \left|\binom H2\right|=\binom62=15,
\end{equation}
et définissons deux espaces d'incidences de types distincts :
\begin{equation}
\label{eq:modular-F6-F8}
\mathcal F_6:=H\times\binom H2,
\qquad
\mathcal F_8:=O\times\binom H2.
\end{equation}
L'inclusion d'incidences est
\begin{equation}
\label{eq:modular-iota68}
\iota_{6,8}:\mathcal F_6\hookrightarrow\mathcal F_8,
\qquad
(h,P)\longmapsto(j(h),P).
\end{equation}
Elle est injective parce que \(j\) l'est et conserve la paire marquée \(P\).

\begin{theorem}[Décomptes d'incidence et défaut orienté]
\label{thm:modular-incidence-90-120}
Les deux espaces de \cref{eq:modular-F6-F8} satisfont
\begin{equation}
\label{eq:modular-counts-90-120}
|\mathcal F_6|=6\binom62=90,
\qquad
|\mathcal F_8|=8\binom62=120,
\qquad
|\mathcal F_8\setminus\iota_{6,8}(\mathcal F_6)|=30.
\end{equation}
La normalisation par les vingt-quatre coordonnées du couple
dodécade--vingt-quatre et par les quinze paires de \(H\) produit
\begin{equation}
\label{eq:modular-incidence-defect}
\boxed{
\delta_{6,8}^{\rm inc}
=\frac{|\mathcal F_6|-|\mathcal F_8|}
{24\binom62}=-\frac1{12}.}
\end{equation}
La même fraction s'obtient par la fonctionnelle réciproque associée
à la mémoire incorporée,
\begin{equation}
\label{eq:modular-reciprocal-defect}
\boxed{
\delta_{6,8}^{\rm rec}
=\frac{30}{120}-\frac{30}{90}=-\frac1{12}.}
\end{equation}
Ce sont deux réalisations exactes du même défaut orienté, mais elles ont
des domaines distincts et ne sont pas identifiées à \(\zeta(-1)\) sans un
entrelaceur supplémentaire.
\end{theorem}

\begin{proof}
Le principe multiplicatif appliqué à
\cref{eq:modular-F6-F8} donne \(6\cdot15=90\) et \(8\cdot15=120\).
L'image de \(\iota_{6,8}\) a quatre-vingt-dix éléments, de sorte que le
complémentaire en a trente. Substituer dans
\cref{eq:modular-incidence-defect} donne
\(-30/(24\cdot15)=-1/12\), tandis que
\cref{eq:modular-reciprocal-defect} est
\(1/4-1/3=-1/12\).
\end{proof}

La transition hexade--octade admet deux réalisations fonctionnelles qui ne
se déduisent pas l'une de l'autre. Le diagramme algébrique, formulé sur les espaces
qui justifient les décomptes, est
\begin{equation}
\label{eq:modular-68-two-readers}
\begin{aligned}
(H\xhookrightarrow{\,j\,}O)
&\longmapsto
\left(|\mathcal F_6|,|\mathcal F_8|\right)=(90,120)
\longmapsto\Gamma_{6\to8}=\gammaH,\\
(h_6,o_8,t)=(6,8,3)
&\longmapsto-\frac{o_8-h_6}{t}=-\frac23,
\end{aligned}
\end{equation}
où \(-2/3\) est le coefficient de la troisième coordonnée de la ligne
\(13\) du constructeur CKM. La première réalisation transforme
l'incidence en aire et en mémoire modulaire ; la seconde transforme la même
transition en paramètres de mélange angulaire.
On ne postule pas \(\boldsymbol\theta_{\rm CKM}=f(\gammaH)\) : les deux sorties
conservent un antécédent commun.

Le morphisme qui détermine les secteurs \(90/120\) est \(\iota_{6,8}\). La grandeur \(\gammaH\) est adoptée comme théorème déjà démontré dans la construction hexade--octade et n’est pas reconstruite à partir de sa représentation décimale. Les conséquences opératorielles ne commencent qu’après fixation de l’échelle élémentaire \(a_0\), définie ci-dessous sans utiliser l’entropie ni le produit terminal \(\gammaH a_0\).

\section{Correspondance entre l'intérieur nonadique et les degrés de liberté de bord}

Soit \(\mathbb T_3=\{0,1,2\}\). Tout \(x\in\mathbb Z/9\mathbb Z\) possède
une écriture en chiffres unique
\(x=x_0+3x_1\), avec \(x_i\in\mathbb T_3\), et tout
\(y\in\mathbb Z/27\mathbb Z\) une écriture unique
\(y=y_0+3y_1+9y_2\). C'est pourquoi
\begin{equation}
\label{eq:modular-bulk-boundary-729}
\mathcal B=(\mathbb Z/9\mathbb Z)^3,
\qquad
\mathcal B_\partial=(\mathbb Z/27\mathbb Z)^2,
\qquad
|\mathcal B|=|\mathcal B_\partial|=3^6=729.
\end{equation}
L'égalité des cardinaux s'accompagne de l'application explicite. Si
\(x_j=r_{2j-1}+3r_{2j}\) pour \(j=1,2,3\), nous définissons
\begin{equation}
\label{eq:modular-beta729}
\beta_{729}(x_1,x_2,x_3)
=\bigl(r_1+3r_2+9r_3,\;r_4+3r_5+9r_6\bigr).
\end{equation}

\begin{proposition}[Bijection généalogique \(729\leftrightarrow729\)]
\label{prop:modular-beta729}
L'application \(\beta_{729}:\mathcal B\to\mathcal B_\partial\) est bijective et
préserve le mot ordonné de six trits. Ce n'est pas, en général, un
isomorphisme de groupes additifs : regrouper les chiffres change l'endroit où
la retenue est enregistrée.
\end{proposition}

\begin{proof}
Extraire les six chiffres de trois coordonnées nonadiques et les regrouper en
deux triplets est une permutation des coordonnées de \(\mathbb T_3^6\). La
décomposition positionnelle est unique des deux côtés, de sorte que
l'opération possède un inverse. La mise en garde additive se vérifie, par exemple,
en additionnant des mots qui engendrent une retenue entre \(r_3\) et \(r_4\) : cette
frontière sépare les deux coordonnées d'ordre vingt-sept, mais traverse
deux coordonnées distinctes dans le regroupement nonadique.
\end{proof}

Sur le tore de bord \(\mathcal B_\partial\), nous prenons le bloc de sommets
\begin{equation}
\label{eq:modular-block-A}
A=\{0,1,\ldots,8\}\times\{0,1,\ldots,8\}.
\end{equation}
Ses liens orientés qui traversent le bord se séparent par direction :
\begin{align}
\partial A_{90}
&=\{((-1,j),(0,j)),((8,j),(9,j)):0\leq j\leq8\},
\label{eq:modular-boundary90}\\
\partial A_{120}
&=\{((i,-1),(i,0)),((i,8),(i,9)):0\leq i\leq8\},
\label{eq:modular-boundary120}
\end{align}
où toutes les coordonnées se lisent modulo vingt-sept. Ainsi
\begin{equation}
\label{eq:modular-boundary-counts}
\partial A=\partial A_{90}\sqcup\partial A_{120},
\qquad
|\partial A_{90}|=|\partial A_{120}|=18.
\end{equation}
Les indices \(90,120\) étiquettent la provenance du canal ; ils ne disent pas que
le bord possède quatre-vingt-dix ou cent vingt liens.

\begin{theorem}[Loi aréale de la coupe triadique]
\label{thm:modular-triadic-cut}
Dans le graphe cubique \(N\times N\times N\) à capacité unitaire, toute
coupe entre deux faces opposées possède une capacité d'au moins \(N^2\), et une
couche transversale atteint cette borne. Pour \(N=9\),
\begin{equation}
\label{eq:modular-mincut-entropy}
\mathcal A_{\min}=81,
\qquad
S_{\rm tri}=81\log3.
\end{equation}
\end{theorem}

\begin{proof}
Il existe \(N^2\) droites d'arêtes parallèles à la direction normale, deux à deux
disjointes. Toute coupe doit intercepter chacune d'elles, de sorte que sa capacité
est d'au moins \(N^2\). Une seule couche transversale contient exactement
\(N^2\) arêtes et réalise le minimum. Si chaque lien coupé porte un trit
maximalement mélangé, son entropie est \(\log3\) ; l'additivité sur les
quatre-vingt-un liens donne \cref{eq:modular-mincut-entropy}.
\end{proof}

L'entropie \(S_{\rm tri}\) mesure la capacité de la coupe cubique.
L'entropie modulaire des trente-six liens de
\cref{eq:modular-boundary-counts} mesure un autre état et ne doit pas lui être
égalisée.

\section{Entrelaceur collectif d'incidence entre intérieur et bord}

Dans la somme
\(\ell^2(\mathcal F_6)\oplus\ell^2(\mathcal F_8)\), nous définissons
\begin{equation}
\label{eq:modular-u90-u120}
u_{90}=\frac1{\sqrt{90}}\sum_{f\in\mathcal F_6}(\delta_f,0),
\qquad
u_{120}=\frac1{\sqrt{120}}\sum_{f\in\mathcal F_8}(0,\delta_f).
\end{equation}
Dans
\(\ell^2(\partial A_{90})\oplus\ell^2(\partial A_{120})\), soient
\begin{equation}
\label{eq:modular-v90-v120}
v_{90}=\frac1{\sqrt{18}}\sum_{e\in\partial A_{90}}(\delta_e,0),
\qquad
v_{120}=\frac1{\sqrt{18}}\sum_{e\in\partial A_{120}}(0,\delta_e).
\end{equation}
Les couples \((u_{90},u_{120})\) et \((v_{90},v_{120})\) sont orthonormaux.
Il existe donc une unique isométrie entre les sous-espaces collectifs
qui satisfait
\begin{equation}
\label{eq:modular-J68}
\mathcal J_{6\to8}u_{90}=v_{90},
\qquad
\mathcal J_{6\to8}u_{120}=v_{120}.
\end{equation}
Définissons
\begin{align}
D_{\rm inc}
&=90|u_{90}\rangle\langle u_{90}|
+120|u_{120}\rangle\langle u_{120}|,
\label{eq:modular-Dinc}\\
D_{\rm area}
&=90|v_{90}\rangle\langle v_{90}|
+120|v_{120}\rangle\langle v_{120}|.
\label{eq:modular-Darea}
\end{align}

\begin{theorem}[Transport exact des deux étiquettes]
\label{thm:modular-J68}
L'application \(\mathcal J_{6\to8}\) est unitaire entre les deux plans
collectifs et entrelace les opérateurs d'incidence et d'aire :
\begin{equation}
\label{eq:modular-J-intertwining}
\boxed{
\mathcal J_{6\to8}D_{\rm inc}
=D_{\rm area}\mathcal J_{6\to8}.}
\end{equation}
Par conséquent, toute combinaison polynomiale des étiquettes \(90,120\), en
particulier la combinaison primitive \(12:-1\), est transportée avec les
mêmes coefficients.
\end{theorem}

\begin{proof}
L'orthonormalité montre que \(\mathcal J_{6\to8}\) envoie une base
orthonormale sur une autre. En \(u_{90}\), les deux membres de
\cref{eq:modular-J-intertwining} valent \(90v_{90}\) ; en \(u_{120}\),
ils valent tous deux \(120v_{120}\). L'égalité sur une base prouve l'identité.
Le calcul fonctionnel polynomial conserve l'entrelacement.
\end{proof}

Cette application ne prétend pas injecter quatre-vingt-dix incidences individuelles dans
dix-huit liens : elle transporte exactement les deux modes uniformes et leurs
étiquettes spectrales. L'étendre à des fluctuations non uniformes exigerait
une nouvelle application d'incidence ; l'opérateur modulaire défini ici n'utilise que
le sous-espace collectif qui vient d'être construit.

\section{Détermination de l'échelle élémentaire du spectre de l'opérateur d'aire}

Soit \(C^*\) la coordonnée angulaire conjuguée, exprimée en degrés, et fixons
\begin{equation}
\label{eq:modular-theta-clock-a0}
\theta_{\rm clk}=\frac{C^*}{6}\frac{\pi}{180},
\qquad
\boxed{
a_0=\frac{1+2\cos(10^\circ)}{3}\,\theta_{\rm clk}.}
\end{equation}
Le facteur qui précède \(\theta_{\rm clk}\) n'est pas une correction
décimale. C'est le caractère réel normalisé du triplet de phases
\begin{equation}
\label{eq:modular-c10-character}
\frac13\Tr\operatorname{diag}
\left(1,\ee^{\ii\pi/18},\ee^{-\ii\pi/18}\right)
=\frac{1+2\cos(10^\circ)}3.
\end{equation}
Ainsi, \(a_0\) est la moyenne spectrale de la structure angulaire sur le trit
orienté. Elle est positive, est fixée avant l'opérateur d'aire et n'est pas
sélectionnée à partir de l'entropie. L'identification ultérieure
\(\gammaH=\Gamma_{6\to8}\) réalise la sortie HMT dans la connexion
d'Ashtekar--Barbero \cite{Barbero,Immirzi} ; elle ne rouvre pas sa dérivation.

\begin{remark}[Portée du symbole \(a_0\)]
Dans ce chapitre, \(a_0\equiv a_0^{\rm area}\) désigne exclusivement la
normalisation de \cref{eq:modular-theta-clock-a0}. Ce n'est ni le rayon de Bohr
ni le facteur d'échelle cosmologique initial, bien que ces objets s'écrivent
traditionnellement avec le même symbole. Aucune identité de ce chapitre
ne permet de substituer l'un à l'autre.
\end{remark}
 \section{Loi aréale de la coupe triadique : capacité \(81\log 3\), opérateur d’aire, spectre et entrelaceur collectif des canaux \(90/120\)}
\label{hap:sec:area-entrelazamiento-barbero-rev11}
\index{Barbero--Immirzi!opérateur d’aire}
\index{entropie modulaire!Hamiltonien modulaire}

L’identification HMT de la fonctionnelle \(\Gamma_{6\to8}\) à la coordonnée
de Barbero--Immirzi s’effectue sur le même bord triadique qui porte
l’incidence (90/120). L’intérieur et son feuillet de bord ont la même cardinalité,
\[
 \mathcal B=(\mathbb Z/9\mathbb Z)^3,
 \quad |\mathcal B|=729,
 \qquad
 \partial\mathcal B=(\mathbb Z/27\mathbb Z)^2,
 \quad |\partial\mathcal B|=729,
\]
par le regroupement de six trits en trois coordonnées nonadiques ou deux
coordonnées d’ordre vingt-sept.

\begin{theorem}[Loi aréale de la coupe triadique]
\label{hap:thm:ley-areal-corte-triadico-rev11}
Dans le graphe cubique \(N\times N\times N\) de capacité unitaire, la coupe
minimale entre deux faces opposées a pour capacité \(N^2\). Pour \(N=9\),
\[
 \mathcal A_{\min}=81,
 \qquad S_{\rm tri}=\mathcal A_{\min}\log3=81\log3.
\]
\end{theorem}
\begin{proof}
Les \(N^2\) droites parallèles à la direction normale sont des chemins disjoints en
arêtes, de sorte que toute coupe les intercepte. Une couche transversale contient
exactement \(N^2\) arêtes et atteint la borne. Chaque liaison triadique apporte
\(\log3\) à l’entropie d’un état maximalement mélangé.
\end{proof}

Dans le tore de bord \(27\times27\), soit \(A\) le bloc canonique
\(9\times9\). Sa frontière se décompose exactement en
\begin{equation}
 \partial A=\partial A_{90}\sqcup\partial A_{120},
 \qquad |\partial A_{90}|=|\partial A_{120}|=18.
 \label{hap:eq:corte-90-120-rev11}
\end{equation}
Les liaisons horizontales représentent le canal (90), et les verticales le
canal (120). Cette décomposition transforme la paire d’incidences en une
décomposition opératorielle du bord.

L’entropie aréale \(S_{\rm tri}=81\log3\) appartient à l’état maximalement
mélangé des liaisons de la coupe. L’entropie bicouche
\(S_{\rm bi}(y)\) de la \cref{hap:def:entropia-bicapa-rev6} appartient à la
distribution normalisée entre les deux feuillets ; les deux objets conservent
des domaines distincts.

Soit \(N=\operatorname{diag}(0,1,2)\) l’opérateur nombre d’une liaison tritique.
La réalisation aréale déclare l’échelle positive
\[
 \theta_{\rm clk}=\frac{C^*}{6}\frac{\pi}{180},
 \qquad
 a_0=\frac{1+2\cos(10^\circ)}{3}\,\theta_{\rm clk},
 \qquad \gamma_{\rm HMT}=\Gamma_{6\to8}.
\]
La formule de \(a_0\) constitue la structure de départ de cette
réalisation ; une fois cette échelle fixée, l’opérateur et son spectre se déduisent
exactement.

\begin{definition}[Opérateur d’aire HMT]
Sur le feuillet de bord
\(\mathcal H_{\partial}=\bigotimes_{e\in\partial A}\mathbb C^3\), on définit
\begin{equation}
 \widehat{\mathcal A}_{\rm HMT}
 =\gamma_{\rm HMT}a_0\sum_{e\in\partial A}N_e.
 \label{hap:eq:operador-area-hmt-rev11}
\end{equation}
Son spectre dans la coupe canonique est
\[
 \operatorname{Spec}\widehat{\mathcal A}_{\rm HMT}
 =\{n\gamma_{\rm HMT}a_0:0\leq n\leq72\},
\]
avec des multiplicités données par les coefficients de
\((1+z+z^2)^{36}\). Le prolongement par coupes compatibles produit la
famille protégée \(\mathcal A_n^{\rm prot}=n\gamma_{\rm HMT}a_0\).
\end{definition}

L’entrelaceur qui fixe la normalisation se construit dans le sous-espace des
canaux collectifs. Soient \(u_{90},u_{120}\) les vecteurs unitaires uniformes
des incidences hexadique et octadique, et soient \(v_{90},v_{120}\) les
vecteurs unitaires uniformes des deux ensembles de liaisons de
\eqref{hap:eq:corte-90-120-rev11}. L’application
\begin{equation}
 \begin{aligned}
 \mathcal J_{6\to8}:{}
 &\operatorname{span}\{u_{90},u_{120}\}
 \longrightarrow
 \operatorname{span}\{v_{90},v_{120}\},\\
 &\mathcal J_{6\to8}u_{90}=v_{90},
 \qquad
 \mathcal J_{6\to8}u_{120}=v_{120}
 \end{aligned}
 \label{hap:eq:entrelazador-area-barbero-rev11}
\end{equation}
est une isométrie entre ces sous-espaces collectifs. Si
\[
 D_{\rm inc}=90|u_{90}\rangle\langle u_{90}|
 +120|u_{120}\rangle\langle u_{120}|,
 \qquad
 D_{\rm area}=90|v_{90}\rangle\langle v_{90}|
 +120|v_{120}\rangle\langle v_{120}|,
\]
alors
\[
 \boxed{\mathcal J_{6\to8}D_{\rm inc}
 =D_{\rm area}\mathcal J_{6\to8}.}
\]
En particulier, la combinaison primitive \(12:-1\) agit
avec les mêmes coefficients avant et après \(\mathcal J_{6\to8}\). Cette
égalité fixe le dictionnaire de normalisation qui insère
\(\gamma_{\rm HMT}=\Gamma_{6\to8}\) dans la connexion
d’Ashtekar--Barbero.

 Le passage de l’hexade à l’octade est une transformation de la structure d’incidence.\par

Lors du passage de \(6\) à \(8\), les paires disponibles, les relations internes et la manière dont l’orientation peut se conserver changent. Cette transition engendre un défaut orienté. La variation cardinale réalise une rotation du mode d’organisation.\par

La même inversion de feuillet s’exprime par\par

\[
C^*\longmapsto -C^*.
\]

Mais des lecteurs différents répondent de manières différentes à cette inversion.\par

La fonctionnelle hexade–octade qui conduit au paramètre de Barbero–Immirzi est impaire : lorsque le feuillet est inversé, elle change de signe. L’orientation géométrique s’inverse.\par

L’entropie bicouche, en revanche, est paire : elle attribue la même entropie à \(C^*\) et à \(-C^*\) lorsque l’on observe la distribution de probabilités. L’information totale peut rester identique même si l’orientation a été inversée.\par

L’espace CKM réalise une troisième réponse. L’inversion induit sur les trois angles une réflexion rationnelle exacte :\par

\[
R_{\rm CKM}
=
M_{\rm CKM}
\operatorname{diag}(1,-1,1)
M_{\rm CKM}^{-1},
\]

con\par

\[
R_{\rm CKM}^2=I,
\qquad
\det R_{\rm CKM}=-1.
\]

Cela signifie que le changement de feuillet est transporté dans l’espace de mélange comme une véritable réflexion. Il existe une direction impaire de rang un et un plan qui reste fixe.\par

Le contre-angle \(C^*\) est précisément la coordonnée qui change. La phase construite à partir de \(A\) reste invariante. La transformation sépare ainsi deux mémoires :\par

\begin{itemize}[leftmargin=*,itemsep=0.35em]
\item une mémoire d’orientation ;
\item une mémoire de phase.
\end{itemize}

Cela offre une lecture très profonde de CKM. Les angles de mélange sont les coordonnées d’un espace sur lequel agit la même involution qui organisait déjà la transition hexade–octade.\par

La rotation \(6\to8\) fournit l’opération d’orientation que CKM réalise dans son propre codomaine ; la causalité physique appartient au transducteur correspondant.\par

Nous avons alors trois publications sœurs du même passage :\par

\begin{enumerate}[leftmargin=*,itemsep=0.65em]
\item \textbf{Publication géométrique.}
Le paramètre de Barbero–Immirzi oriente le quantum d’aire.\par

\item \textbf{Publication informationnelle.}
L’état réduit et son hamiltonien modulaire mesurent quelle information de cette orientation demeure au bord.\par

\item \textbf{Publication de mélange.}
La réflexion CKM transforme les coordonnées de mélange et isole leur composante impaire.\par
\end{enumerate}

\(\gamma_{\rm BI}\), l’entropie et les angles CKM demeurent des réalisations distinctes d’une charnière généalogique commune. L’unité réside dans l’antécédent antérieur à la ramification. Ils partagent la charnière, mais chacun en conserve une partie différente.\par

C’est l’une des plus belles expressions de la méthodologie HMT : l’unité fondamentale est retrouvée en remontant les généalogies jusqu’à trouver l’opération commune qui engendre des résultats terminaux différents.\par

La gravité conserve la rotation comme orientation de l’aire.\par
L’information conserve la rotation comme mémoire de l’état.\par
Le mélange conserve la rotation comme réflexion des coordonnées.\par

La même transformation « respire » de trois manières.\par

 \section{Séparation causale entre le spectre d’aire, la mémoire modulaire et la réalisation Cartan--Holst}
La dérivation de $\gamma_{\mathrm{BI}}$, de l'aire minimale et du spectre de l'opérateur d'aire s'achève dans ce chapitre. La reconstruction conjointe de l'occupation et de la provenance au moyen du Hamiltonien modulaire de bord, ainsi que la réalisation Palatini--Cartan--Holst, conservent leurs domaines, opérateurs et preuves propres dans la Partie IV. \endgroup
% Secciones 4--7 del delta gravitatorio íntegro.
\section{Incidence hexade--octade, paramètre de Barbero--Immirzi et spectre d'aire}
\label{sec:l9s102:barbero-area}

La dérivation surfacique reçoit comme antécédents l'inclusion d'incidence
\(6\to8\), les cardinaux TPK \(90/120\), le défaut normalisé
\(-1/12\) et les canaux spectraux \(q_\pm\). La fonctionnelle
hexade--octade détermine le paramètre de Barbero--Immirzi et la
normalisation de l'opérateur d'aire. Cet emplacement contient la dérivation
HMT et le spectre ; la réalisation Palatini--Cartan--Holst occupe un emplacement
ultérieur dans la mécanique dimensionnelle.

Soient \(H\subset O\), avec \(|H|=6\) et \(|O|=8\), et conservons le couple
marqué de l'hexade dans l'inclusion
\begin{equation}
 \iota_{6,8}:\mathcal F_6=H\times\binom H2
 \lhook\joinrel\longrightarrow
 \mathcal F_8=O\times\binom H2.
 \label{eq:l9s102:barbero-inclusion-incidencia}
\end{equation}
Le comptage interne donne
\begin{equation}
 |\mathcal F_6|=6\binom62=90,
 \qquad |\mathcal F_8|=8\binom62=120,
 \qquad
 \frac{90-120}{24\binom62}=-\frac1{12}.
 \label{eq:l9s102:barbero-conteos}
\end{equation}
Une fois \(\pi\), l'angle \(A\) et le
contre-angle unique \(C^*\) publiés, les deux canaux sont
\begin{equation}
 q_\pm=\exp\!\left[-\frac{\pi}{180}(A\pm C^*)\right],
 \qquad 0<q_+<q_-<1,
 \label{eq:l9s102:barbero-qpm}
\end{equation}
et leurs séries de retour orienté sont
\begin{equation}
 S_{90}(q)=\frac{q^{90}}{1-q^{270}},\qquad
 S_{120}(q)=\frac{q^{120}}{1-q^{360}},\qquad
 \Delta S_m=S_m(q_-)-S_m(q_+).
 \label{eq:l9s102:barbero-series}
\end{equation}

\begin{theorem}[Détermination diophantienne de la fonctionnelle de Barbero--Immirzi]
\label{thm:l9s102:barbero-funcional} Parmi les combinaisons positives \(a\Delta S_{90}-b\Delta S_{120}\), la condition de résidu \(1/24\) et la minimalité coordonnée par coordonnée déterminent l'unique couple primitif \((a,b)=(12,1)\). La sortie adimensionnelle est donc
\begin{equation}
 \boxed{
 \gamma_{\rm HMT}=\Gamma_{6\to8}
 =\frac{{\pi}AC^*}{180}
   +12\Delta S_{90}-\Delta S_{120}.}
 \label{eq:l9s102:barbero-gamma}
\end{equation}
Cette formule reçoit le couple angulaire déjà généré et n'utilise pas en entrée
la valeur conventionnelle du paramètre de Barbero--Immirzi.
\end{theorem}

\begin{proof}
Le résidu de \(S_{90}\) en \(q=1\) est \(1/270\), et celui de \(S_{120}\) est \(1/360\). La condition requise équivaut à
\[
 \frac a{270}-\frac b{360}=\frac1{24}
 \quad\Longleftrightarrow\quad 4a-3b=45.
\]
Ses solutions entières positives sont \((a,b)=(12+3t,1+4t)\), \(t\ge0\) ; l'unique solution minimale est \((12,1)\). Substituer ce couple dans la fonctionnelle d'incidence prouve \eqref{eq:l9s102:barbero-gamma}.
\end{proof}

La séquence causale est fixée sans échange des rôles :
\begin{equation}
 (\mathcal F_6\hookrightarrow\mathcal F_8)
 \longrightarrow (90,120,-1/12)
 \longrightarrow (A,C^*,q_\pm)
 \longrightarrow \Gamma_{6\to8}=\gamma_{\rm HMT}
 \longrightarrow \widehat{\mathcal A}_{\rm phys}.
 \label{eq:l9s102:barbero-cadena-completa}
\end{equation}
La réalisation ultérieure dans l'action de Holst reconnaît la même section adimensionnelle ; elle n'intervient pas dans sa sélection.

Pour les deux secteurs de bord, soient
\begin{equation}
 N_m=\sum_{e\in\partial A_m}N_e,
 \qquad m\in\{90,120\},
 \qquad \operatorname{Spec}(N_e)=\{0,1,2\}.
 \label{eq:l9s102:n90-n120}
\end{equation}
L'opérateur physique d'aire s'écrit
\begin{equation}
 \boxed{
 \frac{\widehat{\mathcal A}_{\mathrm{phys}}}{\ell_P^2}
 =\gamma_{\mathrm{HMT}}a_0(N_{90}+N_{120}).}
 \label{eq:l9s102:area-operador}
\end{equation}
Sur \(36\) liens qutrit,
\begin{equation}
 \operatorname{Spec}(N_{90}+N_{120})=\{0,1,\ldots,72\},
 \label{eq:l9s102:area-numero-spectrum}
\end{equation}
et la multiplicité de la valeur propre \(n\) est le coefficient de \(z^n\) dans
\((1+z+z^2)^{36}\). Il en résulte
\begin{equation}
 \operatorname{Spec}\!\left(
 \frac{\widehat{\mathcal A}_{\mathrm{phys}}}{\ell_P^2}\right)
 =\{n\gamma_{\mathrm{HMT}}a_0:0\le n\le72\}.
 \label{eq:l9s102:area-spectrum}
\end{equation}

\section{Géométrie totale et mémoire de provenance}
\label{sec:l9s102:area-hamiltoniano}

Le hamiltonien modulaire du bord est
\begin{equation}
 \boxed{
 K_A=-\log\rho_A
 =cI+\Delta_{\mathrm{mod}}(90N_{90}+120N_{120}).}
 \label{eq:l9s102:hamiltoniano-modular}
\end{equation}
Soient
\begin{equation}
 N=N_{90}+N_{120},
 \qquad
 D=90N_{90}+120N_{120}.
 \label{eq:l9s102:n-d}
\end{equation}
La matrice de passage
\(\left(\begin{smallmatrix}1&1\\90&120\end{smallmatrix}\right)\)
a pour déterminant \(30\), et son inverse reconstruit les secteurs :
\begin{equation}
 \boxed{
 N_{90}=4N-\frac{D}{30},
 \qquad
 N_{120}=\frac{D}{30}-3N.}
 \label{eq:l9s102:reconstruccion-sectores}
\end{equation}
L'aire publie le nombre total d'excitations ; le hamiltonien modulaire
publie leur secteur de provenance. Le couple conjoint reconstruit l'occupation
complète du bord.

\section{État bicouche et information orientée}
\label{sec:l9s102:estado-bicapa}

Soit \(R=R^*=R^{-1}\) l'involution des feuillets et
\(y=C^*_{\mathrm{rad}}\). L'état normalisé est
\begin{equation}
 \boxed{
 \rho_y=\frac{e^{-yR}}{2\cosh y}.}
 \label{eq:l9s102:rho-y}
\end{equation}
Son hamiltonien modulaire et son entropie sont
\begin{equation}
 K_y=-\log\rho_y=\log(2\cosh y)I+yR,
 \label{eq:l9s102:ky}
\end{equation}
\begin{equation}
 S(\rho_y)=\log(2\cosh y)-y\tanh y.
 \label{eq:l9s102:entropia-y}
\end{equation}
L'inversion \(y\mapsto-y\) laisse l'entropie fixe et produit la distance
informationnelle orientée
\begin{equation}
 \boxed{
 D(\rho_y\Vert\rho_{-y})=2y\tanh y.}
 \label{eq:l9s102:entropia-relativa-hojas}
\end{equation}
La composante paire mesure la distribution spectrale ; la composante impaire mesure
le coût informationnel de l'inversion de l'orientation.

\section{Quantum informationnel d'énergie et cellule gravitationnelle de demi-action}
\label{sec:l9s102:horizonte-media-accion}

La périodicité de \(108\) états détermine
\begin{equation}
 t_P=\frac{54}{\pi}t_0,
 \qquad
 E_P=\frac{\hbar}{t_P}=\frac{h}{108t_0}.
 \label{eq:l9s102:ep-tp}
\end{equation}
La réalisation thermique de l'horloge tritique satisfait
\begin{equation}
 \boxed{
 E_P=k_B\Theta_{\mathrm{clk}}\log3.}
 \label{eq:l9s102:planck-trit}
\end{equation}
Pour l'horizon de rayon \(R\),
\begin{equation}
 E_\Lambda(R)=\frac{c^4R}{2G},
 \qquad
 E_\Lambda(R)=T_H(R)S_H(R).
 \label{eq:l9s102:energia-horizonte}
\end{equation}
Dans la cellule de Planck,
\begin{equation}
 \boxed{
 E_\Lambda
 =T_HS_H
 =\frac{E_P}{2}
 =\frac{k_B\Theta_{\mathrm{clk}}\log3}{2}.}
 \label{eq:l9s102:confluencia-horizonte}
\end{equation}
Par conséquent,
\begin{equation}
 \boxed{
 E_\Lambda t_P=\frac{\hbar}{2},
 \qquad
 E_\Lambda T_{108}=\frac h2.}
 \label{eq:l9s102:media-accion}
\end{equation}
L'énergie du vide, l'entropie surfacique, la température chronologique et
l'action partagent une même cellule de confluence avec des domaines typés.

Lorsque \(S_H/k_B=\pi\), la multiplicité effective satisfait
\begin{equation}
 \Omega_{\mathrm{eff}}=e^{S_H/k_B}=e^\pi.
 \label{eq:l9s102:omega-e-pi}
\end{equation}
L'opérateur hyperbolique d'Euler possède
\begin{equation}
 \operatorname{Spec}(e^{\pi H})=\{e^\pi,e^{-\pi}\}.
 \label{eq:l9s102:e-pi-spectrum}
\end{equation}
L'égalité des scalaires établit la confluence spectrale ;
l'opérateur d'entrelacement entre multiplicité microscopique d'horizon et dynamique de
Perron relève de son emplacement physique spécifique.
 
\chapter{Réversibilité exacte, coût d'effacement et limite nulle de Landauer}
\begingroup
\renewcommand{\chapter}[2][]{}%
\chapter{Mise à jour réversible, coût d’effacement et domaine exact de la limite nulle de Landauer}
\label{chap:parte-iv-57}

Le transport complet conserve l’information de fibre et peut avoir un coût logique nul ; la projection ou la réinitialisation qui écarte la généalogie est soumise au coût de Landauer. Les deux domaines sont explicitement séparés.

\section{Réversibilité logique et principe de Landauer dans les fibres de TPK}
\label{sec:informacion-fibra-landauer-rev7}
\index{Landauer}\index{information!fibre}

Soit \((\Omega,2^\Omega,\mu)\) un espace probabilisé fini, soit \(X:\Omega\to\mathcal X\) une variable aléatoire à valeurs dans l’espace \(\mathcal X\) des états complets et soit \(q:\mathcal X\to Y\) un lecteur. La règle de chaîne donne
\[
 \mathsf H(X)
 =
 \mathsf H(q(X))+\mathsf H(X\mid q(X)).
\]
Le deuxième terme mesure l’information de la fibre que le lecteur n’exprime pas.

\begin{theorem}[Coût de Landauer nul pour le transport complet et coût de la projection]
\label{thm:landauer-relativo-lector-rev7}
Si \(U:\mathcal X\to\mathcal X\) est une mise à jour bijective de l’état HMT complet, alors
\[
 \boxed{\mathsf H(X\mid U(X))=0.}
\]
Pour une sortie projetée \(q(U(X))\), l’information écartée est
\[
 \mathsf H(X\mid q(U(X))).
\]
Une réinitialisation physique qui identifie des états distincts de la mémoire doit comptabiliser cette information dans l’environnement.
\end{theorem}

\begin{proof}
La bijectivité implique \(X=U^{-1}(U(X))\) ; \(X\) est donc déterminé par \(U(X)\) et l’entropie conditionnelle est nulle. Pour le lecteur \(q\), la règle de chaîne sépare l’entropie exprimée de l’entropie de la fibre. Si l’opération de réinitialisation n’est pas injective, la sortie ne permet plus de reconstruire l’état de mémoire. Le principe thermodynamique de Landauer borne le coût de l’implémentation en fonction de cette entropie effacée ; la mise à jour réversible ne fournit pas à elle seule de terme d’effacement.
\end{proof}

En nats et en bits, respectivement,
\[
 Q\ge k_B\Theta\,\mathsf H_{\rm nat},
 \qquad
 Q\ge k_B\Theta\log 2\,\mathsf H_{\rm bit}.
\]
Pour le transport bijectif de l’état complet, l’entropie d’effacement est nulle ; le terme minimal de Landauer associé exclusivement à l’effacement satisfait donc
\[
\boxed{Q_{\rm L}^{\min}[U]=k_B\Theta\,
 \mathsf H(X\mid U(X))=0.}
\]
Pour la réinitialisation d’un TRIT uniforme, en revanche,
\[
 \boxed{Q_{\rm reset}\ge k_B\Theta\log3.}
\]
Ces deux équations situent le coût dans la perte d’information : le coût apparaît lorsque la projection ou la réinitialisation oublie la fibre. La préparation, la projection et la réinitialisation peuvent produire de l’entropie ; le transport réversible conserve l’information généalogique. La valeur nulle est la limite logique de Landauer pour cette mise à jour complète, tandis que les coûts dynamiques de contrôle, de vitesse finie et de couplage à l’environnement relèvent de la réalisation physique du dispositif.

\begin{statusbox}[title={Portée du coût nul de Landauer}] L'égalité \(\mathsf H(X\mid U(X))=0\) est exacte pour une mise à jour bijective de l'état holonomique intégral. La borne \(Q_{\rm reset}\ge k_B\Theta\log3\) est la spécialisation de Landauer à l'effacement d'une décision triadique uniforme. La réalisation des deux formules en joules utilise le transducteur de Boltzmann et une température physique déclarée.
\end{statusbox}
 
\paragraph{Transduction de l’information de bord en capacité de la coupe triadique minimale.} L'information conservée par le bord réapparaît géométriquement dans la coupe triadique minimale et sa capacité. \endgroup

\chapter{Aire minimale de la coupe triadique et capacité \texorpdfstring{\(81\log 3\)}{81 log 3}}
\begingroup
\renewcommand{\chapter}[2][]{}%
\chapter{Aire minimale de la coupe triadique, capacité \texorpdfstring{\(81\log 3\)}{81 log 3} et transport collectif d’incidence}
\label{chap:parte-iv-58}

La quantification aréale part de l’intérieur nonadique, de son bord et du transport collectif d’incidence. La coupe minimale contient \(81\) unités et sa capacité qutrit est \(81\log 3\) ; le résultat précède sa réalisation gravitationnelle.

\section{Quantification aréale, dynamique modulaire et transformation rationnelle CKM}
\label{rm:ch:modular}

\subsection{Inclusion hexade--octade et réalisations fonctionnelles associées}

Fixons un ensemble \(H\) de six points et un ensemble \(O\) de huit
points avec une inclusion distinguée \(j:H\hookrightarrow O\). Dans la
résidence de Witt, \(H\) est une hexade d'une dodécade et \(O\) son octade
relevée ; dans ce chapitre, on utilise seulement l'inclusion déjà fixée, et non la
valeur terminale de Barbero--Immirzi. Nous écrivons
\begin{equation}
\label{rm:eq:modular-pairs-H}
\binom H2:=\{P\subset H:|P|=2\},
\qquad \left|\binom H2\right|=\binom62=15,
\end{equation}
et définissons deux espaces d'incidences de types distincts :
\begin{equation}
\label{rm:eq:modular-F6-F8}
\mathcal F_6:=H\times\binom H2,
\qquad
\mathcal F_8:=O\times\binom H2.
\end{equation}
L'inclusion d'incidences est
\begin{equation}
\label{rm:eq:modular-iota68}
\iota_{6,8}:\mathcal F_6\hookrightarrow\mathcal F_8,
\qquad
(h,P)\longmapsto(j(h),P).
\end{equation}
Elle est injective parce que \(j\) l'est et conserve la paire marquée \(P\).

\begin{theorem}[Décomptes d'incidence et défaut orienté]
\label{rm:thm:modular-incidence-90-120}
Les deux espaces de \cref{rm:eq:modular-F6-F8} satisfont
\begin{equation}
\label{rm:eq:modular-counts-90-120}
|\mathcal F_6|=6\binom62=90,
\qquad
|\mathcal F_8|=8\binom62=120,
\qquad
|\mathcal F_8\setminus\iota_{6,8}(\mathcal F_6)|=30.
\end{equation}
La normalisation par les vingt-quatre coordonnées du couple
dodécade--vingt-quatre et par les quinze paires de \(H\) produit
\begin{equation}
\label{rm:eq:modular-incidence-defect}
\boxed{
\delta_{6,8}^{\rm inc}
=\frac{|\mathcal F_6|-|\mathcal F_8|}
{24\binom62}=-\frac1{12}.}
\end{equation}
La même fraction s'obtient par la fonctionnelle réciproque associée
à la mémoire incorporée,
\begin{equation}
\label{rm:eq:modular-reciprocal-defect}
\boxed{
\delta_{6,8}^{\rm rec}
=\frac{30}{120}-\frac{30}{90}=-\frac1{12}.}
\end{equation}
Ce sont deux réalisations exactes du même défaut orienté, mais elles ont
des domaines distincts et ne sont pas identifiées à \(\zeta(-1)\) sans un
entrelaceur supplémentaire.
\end{theorem}

\begin{proof}
Le principe multiplicatif appliqué à
\cref{rm:eq:modular-F6-F8} donne \(6\cdot15=90\) et \(8\cdot15=120\).
L'image de \(\iota_{6,8}\) a quatre-vingt-dix éléments, de sorte que le
complémentaire en a trente. Substituer dans
\cref{rm:eq:modular-incidence-defect} donne
\(-30/(24\cdot15)=-1/12\), tandis que
\cref{rm:eq:modular-reciprocal-defect} est
\(1/4-1/3=-1/12\).
\end{proof}

La transition hexade--octade admet deux réalisations fonctionnelles qui ne
se déduisent pas l'une de l'autre. Le diagramme algébrique, formulé sur les espaces
qui justifient les décomptes, est
\begin{equation}
\label{rm:eq:modular-68-two-readers}
\begin{aligned}
(H\xhookrightarrow{\,j\,}O)
&\longmapsto
\left(|\mathcal F_6|,|\mathcal F_8|\right)=(90,120)
\longmapsto\Gamma_{6\to8}=\gammaH,\\
(h_6,o_8,t)=(6,8,3)
&\longmapsto-\frac{o_8-h_6}{t}=-\frac23,
\end{aligned}
\end{equation}
où \(-2/3\) est le coefficient de la troisième coordonnée de la ligne
\(13\) du constructeur CKM. La première réalisation transforme
l'incidence en aire et en mémoire modulaire ; la seconde transforme la même
transition en paramètres de mélange angulaire.
On ne postule pas \(\boldsymbol\theta_{\rm CKM}=f(\gammaH)\) : les deux sorties
conservent un antécédent commun.

Le morphisme qui détermine les secteurs \(90/120\) est \(\iota_{6,8}\). La grandeur \(\gammaH\) est adoptée comme théorème déjà démontré dans la construction hexade--octade et n’est pas reconstruite à partir de sa représentation décimale. Les conséquences opératorielles ne commencent qu’après fixation de l’échelle élémentaire \(a_0\), définie ci-dessous sans utiliser l’entropie ni le produit terminal \(\gammaH a_0\).

\subsection{Correspondance entre l'intérieur nonadique et les degrés de liberté de bord}

Soit \(\mathbb T_3=\{0,1,2\}\). Tout \(x\in\mathbb Z/9\mathbb Z\) possède
une écriture en chiffres unique
\(x=x_0+3x_1\), avec \(x_i\in\mathbb T_3\), et tout
\(y\in\mathbb Z/27\mathbb Z\) une écriture unique
\(y=y_0+3y_1+9y_2\). C'est pourquoi
\begin{equation}
\label{rm:eq:modular-bulk-boundary-729}
\mathcal B=(\mathbb Z/9\mathbb Z)^3,
\qquad
\mathcal B_\partial=(\mathbb Z/27\mathbb Z)^2,
\qquad
|\mathcal B|=|\mathcal B_\partial|=3^6=729.
\end{equation}
L'égalité des cardinaux s'accompagne de l'application explicite. Si
\(x_j=r_{2j-1}+3r_{2j}\) pour \(j=1,2,3\), nous définissons
\begin{equation}
\label{rm:eq:modular-beta729}
\beta_{729}(x_1,x_2,x_3)
=\bigl(r_1+3r_2+9r_3,\;r_4+3r_5+9r_6\bigr).
\end{equation}

\begin{proposition}[Bijection généalogique \(729\leftrightarrow729\)]
\label{rm:prop:modular-beta729}
L'application \(\beta_{729}:\mathcal B\to\mathcal B_\partial\) est bijective et
préserve le mot ordonné de six trits. Ce n'est pas, en général, un
isomorphisme de groupes additifs : regrouper les chiffres change l'endroit où
la retenue est enregistrée.
\end{proposition}

\begin{proof}
Extraire les six chiffres de trois coordonnées nonadiques et les regrouper en
deux triplets est une permutation des coordonnées de \(\mathbb T_3^6\). La
décomposition positionnelle est unique des deux côtés, de sorte que
l'opération possède un inverse. La mise en garde additive se vérifie, par exemple,
en additionnant des mots qui engendrent une retenue entre \(r_3\) et \(r_4\) : cette
frontière sépare les deux coordonnées d'ordre vingt-sept, mais traverse
deux coordonnées distinctes dans le regroupement nonadique.
\end{proof}

Sur le tore de bord \(\mathcal B_\partial\), nous prenons le bloc de sommets
\begin{equation}
\label{rm:eq:modular-block-A}
A=\{0,1,\ldots,8\}\times\{0,1,\ldots,8\}.
\end{equation}
Ses liens orientés qui traversent le bord se séparent par direction :
\begin{align}
\partial A_{90}
&=\{((-1,j),(0,j)),((8,j),(9,j)):0\leq j\leq8\},
\label{rm:eq:modular-boundary90}\\
\partial A_{120}
&=\{((i,-1),(i,0)),((i,8),(i,9)):0\leq i\leq8\},
\label{rm:eq:modular-boundary120}
\end{align}
où toutes les coordonnées se lisent modulo vingt-sept. Ainsi
\begin{equation}
\label{rm:eq:modular-boundary-counts}
\partial A=\partial A_{90}\sqcup\partial A_{120},
\qquad
|\partial A_{90}|=|\partial A_{120}|=18.
\end{equation}
Les indices \(90,120\) étiquettent la provenance du canal ; ils ne disent pas que
le bord possède quatre-vingt-dix ou cent vingt liens.

\begin{theorem}[Loi aréale de la coupe triadique]
\label{rm:thm:modular-triadic-cut}
Dans le graphe cubique \(N\times N\times N\) à capacité unitaire, toute
coupe entre deux faces opposées possède une capacité d'au moins \(N^2\), et une
couche transversale atteint cette borne. Pour \(N=9\),
\begin{equation}
\label{rm:eq:modular-mincut-entropy}
\mathcal A_{\min}=81,
\qquad
S_{\rm tri}=81\log3.
\end{equation}
\end{theorem}

\begin{proof}
Il existe \(N^2\) droites d'arêtes parallèles à la direction normale, deux à deux
disjointes. Toute coupe doit intercepter chacune d'elles, de sorte que sa capacité
est d'au moins \(N^2\). Une seule couche transversale contient exactement
\(N^2\) arêtes et réalise le minimum. Si chaque lien coupé porte un trit
maximalement mélangé, son entropie est \(\log3\) ; l'additivité sur les
quatre-vingt-un liens donne \cref{rm:eq:modular-mincut-entropy}.
\end{proof}

L'entropie \(S_{\rm tri}\) mesure la capacité de la coupe cubique.
L'entropie modulaire des trente-six liens de
\cref{rm:eq:modular-boundary-counts} mesure un autre état et ne doit pas lui être
égalisée.

\subsection{Entrelaceur collectif d'incidence entre intérieur et bord}

Dans la somme
\(\ell^2(\mathcal F_6)\oplus\ell^2(\mathcal F_8)\), nous définissons
\begin{equation}
\label{rm:eq:modular-u90-u120}
u_{90}=\frac1{\sqrt{90}}\sum_{f\in\mathcal F_6}(\delta_f,0),
\qquad
u_{120}=\frac1{\sqrt{120}}\sum_{f\in\mathcal F_8}(0,\delta_f).
\end{equation}
Dans
\(\ell^2(\partial A_{90})\oplus\ell^2(\partial A_{120})\), soient
\begin{equation}
\label{rm:eq:modular-v90-v120}
v_{90}=\frac1{\sqrt{18}}\sum_{e\in\partial A_{90}}(\delta_e,0),
\qquad
v_{120}=\frac1{\sqrt{18}}\sum_{e\in\partial A_{120}}(0,\delta_e).
\end{equation}
Les couples \((u_{90},u_{120})\) et \((v_{90},v_{120})\) sont orthonormaux.
Il existe donc une unique isométrie entre les sous-espaces collectifs
qui satisfait
\begin{equation}
\label{rm:eq:modular-J68}
\mathcal J_{6\to8}u_{90}=v_{90},
\qquad
\mathcal J_{6\to8}u_{120}=v_{120}.
\end{equation}
Définissons
\begin{align}
D_{\rm inc}
&=90|u_{90}\rangle\langle u_{90}|
+120|u_{120}\rangle\langle u_{120}|,
\label{rm:eq:modular-Dinc}\\
D_{\rm area}
&=90|v_{90}\rangle\langle v_{90}|
+120|v_{120}\rangle\langle v_{120}|.
\label{rm:eq:modular-Darea}
\end{align}

\begin{theorem}[Transport exact des deux étiquettes]
\label{rm:thm:modular-J68}
L'application \(\mathcal J_{6\to8}\) est unitaire entre les deux plans
collectifs et entrelace les opérateurs d'incidence et d'aire :
\begin{equation}
\label{rm:eq:modular-J-intertwining}
\boxed{
\mathcal J_{6\to8}D_{\rm inc}
=D_{\rm area}\mathcal J_{6\to8}.}
\end{equation}
Par conséquent, toute combinaison polynomiale des étiquettes \(90,120\), en
particulier la combinaison primitive \(12:-1\), est transportée avec les
mêmes coefficients.
\end{theorem}

\begin{proof}
L'orthonormalité montre que \(\mathcal J_{6\to8}\) envoie une base
orthonormale sur une autre. En \(u_{90}\), les deux membres de
\cref{rm:eq:modular-J-intertwining} valent \(90v_{90}\) ; en \(u_{120}\),
ils valent tous deux \(120v_{120}\). L'égalité sur une base prouve l'identité.
Le calcul fonctionnel polynomial conserve l'entrelacement.
\end{proof}

Cette application ne prétend pas injecter quatre-vingt-dix incidences individuelles dans
dix-huit liens : elle transporte exactement les deux modes uniformes et leurs
étiquettes spectrales. L'étendre à des fluctuations non uniformes exigerait
une nouvelle application d'incidence ; l'opérateur modulaire défini ici n'utilise que
le sous-espace collectif qui vient d'être construit.

\subsection{Quantum élémentaire et échelle minimale de l’opérateur d’aire}

Soit \(C^*\) la coordonnée angulaire conjuguée, exprimée en degrés, et fixons
\begin{equation}
\label{rm:eq:modular-theta-clock-a0}
\theta_{\rm clk}=\frac{C^*}{6}\frac{\pi}{180},
\qquad
\boxed{
a_0=\frac{1+2\cos(10^\circ)}{3}\,\theta_{\rm clk}.}
\end{equation}
Le facteur qui précède \(\theta_{\rm clk}\) n'est pas une correction
décimale. C'est le caractère réel normalisé du triplet de phases
\begin{equation}
\label{rm:eq:modular-c10-character}
\frac13\Tr\operatorname{diag}
\left(1,\ee^{\ii\pi/18},\ee^{-\ii\pi/18}\right)
=\frac{1+2\cos(10^\circ)}3.
\end{equation}
Ainsi, \(a_0\) est la moyenne spectrale de la structure angulaire sur le trit
orienté. Elle est positive, est fixée avant l'opérateur d'aire et n'est pas
sélectionnée à partir de l'entropie. L'identification ultérieure
\(\gammaH=\Gamma_{6\to8}\) réalise la sortie HMT dans la connexion
d'Ashtekar--Barbero \cite{Barbero1995,Immirzi1997} ; elle ne rouvre pas sa dérivation.

\begin{remark}[Portée du symbole \(a_0\)]
Dans ce chapitre, \(a_0\equiv a_0^{\rm area}\) désigne exclusivement la
normalisation de \cref{rm:eq:modular-theta-clock-a0}. Ce n'est ni le rayon de Bohr
ni le facteur d'échelle cosmologique initial, bien que ces objets s'écrivent
traditionnellement avec le même symbole. Aucune identité de ce chapitre
ne permet de substituer l'un à l'autre.
\end{remark}
 
\paragraph{Construction de l’opérateur d’aire à partir de la coupe triadique et séparation des canaux.} La coupe et sa capacité déterminent un opérateur d'aire ; son spectre doit conserver séparément l'occupation totale et la provenance du canal. \endgroup

\chapter{Spectre de l'opérateur d'aire et reconstruction de la provenance du canal}
\begingroup
\renewcommand{\chapter}[2][]{}%
\chapter{Spectre de l’opérateur d’aire et reconstruction conjointe de l’occupation et de la provenance}
\label{chap:parte-iv-59}

Les secteurs \(90\) et \(120\) définissent des opérateurs compatibles d’occupation et de mémoire. L’inversion du système linéaire reconstruit les deux canaux et empêche l’aire totale d’effacer la provenance.

\section{Opérateurs aréal et modulaire dans les secteurs \texorpdfstring{\(90\)}{90} et \texorpdfstring{\(120\)}{120}}

Dans chacune des deux familles de \cref{rm:eq:modular-boundary-counts}, il existe dix-huit canaux qutrit. Pour \(m\in\{90,120\}\), définissons
\begin{equation}
\label{rm:eq:modular-number-operators}
N_m=\sum_{e\in\partial A_m}N_e,
\qquad \Spec(N_e)=\{0,1,2\}.
\end{equation}
La sortie de Barbero--Immirzi est utilisée ici comme résultat antérieur, sans être redérivée. Avec la normalisation aréale construite dans \cref{rm:eq:modular-theta-clock-a0}, il vient
\begin{equation}
\label{rm:eq:modular-area-operator}
\boxed{
\frac{\widehat{\mathcal A}_{\rm phys}}{\ell_P^2}
=\gammaH a_0(N_{90}+N_{120}),}
\qquad
\gammaH a_0=0.0016755940749224991.
\end{equation}

\begin{theorem}[Spectre fini du quantum d’aire]
\label{rm:thm:modular-area-spectrum}
Sur les trente-six liens de la coupe,
\begin{equation}
\label{rm:eq:modular-total-number-spectrum}
\Spec(N_{90}+N_{120})=\{0,1,\ldots,72\}.
\end{equation}
La multiplicité de la valeur propre \(n\) est le coefficient de \(z^n\) dans \((1+z+z^2)^{36}\). Par conséquent,
\begin{equation}
\label{rm:eq:modular-area-spectrum}
\boxed{
\Spec\!\left(\frac{\widehat{\mathcal A}_{\rm phys}}{\ell_P^2}\right)
=\{n\gammaH a_0:0\leq n\leq72\}.}
\end{equation}
Les coupes compatibles prolongent la famille \(\mathcal A_n^{\rm prot}=n\gammaH a_0\ell_P^2\).
\end{theorem}

\begin{proof}
Chaque opérateur \(N_e\) a pour fonction génératrice spectrale \(1+z+z^2\). Les trente-six facteurs commutent et agissent sur des facteurs tensoriels distincts ; la fonction génératrice du total est donc \((1+z+z^2)^{36}\), dont les degrés parcourent tous les entiers de zéro à soixante-douze. La multiplication par le scalaire positif \(\gammaH a_0\ell_P^2\) donne \cref{rm:eq:modular-area-spectrum}.
\end{proof}

Il faut distinguer deux procédés de quantification. La loi de coupe \cref{rm:thm:modular-triadic-cut} compte quatre-vingt-une capacités unitaires et détermine \(81\log3\) ; l’opérateur \cref{rm:eq:modular-area-operator} quantifie les occupations de trente-six liens de réponse, jusqu’à soixante-douze trits de nombre. Ils partagent le même feuillet de bord, mais pas la même observable.

La torsion modulaire produit
\begin{equation}
\label{rm:eq:modular-delta-value}
\Dmod
=0.0001111970500599773249414566131933856\ldots,
\end{equation}
et le Hamiltonien de bord
\begin{equation}
\label{rm:eq:modular-hamiltonian}
\boxed{
\HHM^{(A)}=-\log\rho_A
=cI+\Dmod(90N_{90}+120N_{120}).}
\end{equation}
Le qualificatif \emph{holographique} ne désigne pas une analogie : l’incidence \(6\to8\) est transportée vers les ensembles de bord \(\partial A_{90}\) et \(\partial A_{120}\), et l’état réduit réside dans l’algèbre engendrée par leurs opérateurs nombre.

\begin{theorem}[Reconstruction conjointe de la géométrie et de la mémoire]
\label{rm:thm:modular-area-memory}
Soient
\begin{equation}
\label{rm:eq:modular-ND}
N=N_{90}+N_{120},
\qquad
D=90N_{90}+120N_{120}.
\end{equation}
Alors
\begin{equation}
\label{rm:eq:modular-inverse-channels}
\boxed{
N_{90}=4N-\frac{D}{30},
\qquad
N_{120}=\frac{D}{30}-3N.}
\end{equation}
En particulier, \([\widehat{\mathcal A}_{\rm phys},\HHM^{(A)}]=0\), et le couple \((\widehat{\mathcal A}_{\rm phys},\HHM^{(A)}-cI)\) reconstruit les deux secteurs d’occupation.
\end{theorem}

\begin{proof}
Les opérateurs \(N_{90}\) et \(N_{120}\) agissent sur des facteurs distincts et commutent. L’application linéaire qui envoie \((N_{90},N_{120})\) sur \((N,D)\) a pour matrice
\[
\begin{pmatrix}1&1\\90&120\end{pmatrix},
\qquad \det=30.
\]
Son inverse est exactement \cref{rm:eq:modular-inverse-channels}. Puisque \cref{rm:eq:modular-area-operator,rm:eq:modular-hamiltonian} sont des combinaisons linéaires des mêmes opérateurs nombre, le commutateur s’annule.
\end{proof}

L’interprétation opératorielle est précise : l’observable d’aire conserve le nombre total d’excitations, tandis que le Hamiltonien modulaire conserve leur secteur d’origine. Aucune des deux observables ne retient à elle seule les deux coordonnées ; leur couple conjoint les détermine.

\section{Confluence de l’aire, de l’information et de l’énergie}

L’échelle angulaire \(\theta_{\rm clk}\) de \cref{rm:eq:modular-theta-clock-a0} et la température chronologique \(\Theta_{\rm clk}\) sont des objets distincts ; l’emploi de la majuscule fait partie du typage. La réalisation thermodynamique de \cref{rm:ch:metrology} démontre
\begin{equation}
\label{rm:eq:modular-planck-trit-bridge}
E_P=k_B\Theta_{\rm clk}\log3,
\end{equation}
tandis que la réalisation aréale construit l’incrément élémentaire
\begin{equation}
\label{rm:eq:modular-area-quantum}
\delta\mathcal A
=\gammaH a_0\ell_P^2.
\end{equation}
Elles ne sont pas égalisées : \cref{rm:eq:modular-planck-trit-bridge} fixe le quantum informationnel d’énergie et \cref{rm:eq:modular-area-quantum} le quantum géométrique d’aire. Leur quotient définit la tension de conversion typée
\begin{equation}
\label{rm:eq:modular-information-area-tension}
\boxed{
\mathcal T_{I/A}
=\frac{E_P}{\delta\mathcal A}
=\frac{k_B\Theta_{\rm clk}\log3}
{\gammaH a_0\ell_P^2}.}
\end{equation}
L’expression n’introduit pas de constante supplémentaire : elle détermine les normalisations qui doivent être conservées lors de la transformation de l’information en géométrie.

\enlargethispage{7\baselineskip} Dans la coupe minimale de quatre-vingt-un liens, la capacité informationnelle physique et l'énergie d'un trit par lien sont
\begin{equation}
\label{rm:eq:modular-cut-info-energy}
S_{\rm cut}^{\rm phys}=81k_B\log3,
\qquad
E_{\rm cut}^{(1)}=81E_P
=\Theta_{\rm clk}S_{\rm cut}^{\rm phys}.
\end{equation}
L’égalité finale est une équation d’état de cette réalisation, non une loi universelle d’équipartition. En particulier, l’état de Gibbs des trente-six liens de bord possède une entropie différente et est traité ci-dessous.

Cette séparation détermine le contenu structurel de la confluence : Barbero--Immirzi oriente le spectre d’aire ; \(\log3\) mesure la décision irréductible du TRIT ; \(\Theta_{\rm clk}\) convertit cette information en énergie ; et \(\ell_P^2\) relie le quantum d’aire à la famille gravitationnelle. La gravité est décrite comme confluence de réalisations fonctionnelles d’un antécédent commun, et non comme une égalité imposée entre leurs valeurs.
 
\paragraph{Transport de la mémoire de canal vers le Hamiltonien modulaire et la purification thermochamp double.} La mémoire de canal est représentée ci-dessous au moyen du Hamiltonien modulaire et d'une purification thermochamp double. \endgroup

\chapter{Hamiltonien modulaire de frontière, purification thermochamp double et information orientée}
\begingroup
\renewcommand{\chapter}[2][]{}%
\chapter{Hamiltonien modulaire de bord, purification thermochamp double et déficit informationnel orienté}
\label{chap:parte-iv-60}

L’état qutrit fidèle admet une purification thermochamp double. Le Hamiltonien modulaire enregistre l’orientation des canaux et le déficit informationnel mesure la séparation finie entre états sans s’identifier à l’aire.

\section{État qutrit, purification thermochamp double et déficit informationnel}

Pour un lien de poids \(w>0\), soient
\begin{equation}
\label{rm:eq:modular-link-state}
Z(w)=1+\ee^{-w}+\ee^{-2w},
\qquad
\rho(w)=\frac1{Z(w)}\sum_{n=0}^2\ee^{-wn}|n\rangle\langle n|.
\end{equation}
Les poids des deux familles sont
\begin{align}
w_{90}&=90\Dmod
=0.0100077345053979592447310951874\ldots,
\label{rm:eq:modular-w90}\\
w_{120}&=120\Dmod
=0.0133436460071972789929747935832\ldots,
\qquad \frac{w_{120}}{w_{90}}=\frac43.
\label{rm:eq:modular-w120}
\end{align}
L’état fini de bord est
\begin{equation}
\label{rm:eq:modular-product-state}
\rho_A=\rho(w_{90})^{\otimes18}
\otimes\rho(w_{120})^{\otimes18}.
\end{equation}
Son logarithme reproduit \cref{rm:eq:modular-hamiltonian}, avec \(c=18\log Z(w_{90})+18\log Z(w_{120})\).

\begin{theorem}[Purification thermochamp double]
\label{rm:thm:modular-tfd}
Le vecteur
\begin{equation}
\label{rm:eq:modular-tfd-link}
|\operatorname{TFD}(w)\rangle
=\frac1{\sqrt{Z(w)}}\sum_{n=0}^2
\ee^{-wn/2}|n\rangle_A\otimes|n\rangle_{\bar A}
\end{equation}
purifie \(\rho(w)\). Le produit de dix-huit copies pour chacun des poids \(w_{90},w_{120}\) purifie \(\rho_A\) ; dans le système dupliqué, \(S(\rho_A)\) est exactement l’entropie d’intrication.
\end{theorem}

\begin{proof}
En prenant la trace sur le deuxième facteur, l’orthogonalité de \(|n\rangle_{\bar A}\) élimine les termes croisés et laisse
\[
\operatorname{Tr}_{\bar A}
|\operatorname{TFD}(w)\rangle\langle\operatorname{TFD}(w)|
=\rho(w).
\]
La trace partielle commute avec les produits tensoriels, ce qui prouve l’assertion pour les trente-six liens.
\end{proof}

Pour un lien,
\begin{equation}
\label{rm:eq:modular-link-entropy}
S(w)=\log Z(w)
+w\frac{\ee^{-w}+2\ee^{-2w}}{Z(w)}.
\end{equation}
En conséquence,
\begin{equation}
\label{rm:eq:modular-entropy-value}
S(\rho_A)=18S(w_{90})+18S(w_{120})
=39.54837320881807994957319730\ldots\ \text{nats}.
\end{equation}
L’état maximalement mélangé de trente-six qutrits possède l’entropie \(36\log3\). Définissons le défaut orienté
\begin{equation}
\label{rm:eq:modular-oriented-information}
\boxed{
I_{\rm or}=36\log3-S(\rho_A)
=0.001669183233868940655631225913\ldots\ \text{nats}.}
\end{equation}
Sa positivité découle de la maximalité entropique de l’état uniforme. La proximité de \(I_{\rm or}\) et de \(\gammaH a_0\) ne constitue pas une identité et n’est pas utilisée pour sélectionner l’une ou l’autre des deux valeurs.
 
\section{Information paire et transport géométrique impair dans la bicouche}

La distribution normalisée entre feuillets sépare l’entropie paire du transport hexade--octade impair. Cette parité bicouche conserve l’information de bord enregistrée par le Hamiltonien modulaire et évite d’identifier entropie, aire et orientation.

\begingroup
\let\chapter\section
\let\section\subsection
\let\subsection\subsubsection
\let\subsubsection\paragraph
\let\clearpage\relax
\let\cleardoublepage\relax
\chapter{La bicouche : information paire et transport géométrique impair}
\label{chap:paridad-bicapa-barbero-rev6}
\index{entropie!bicouche}
\index{Barbero--Immirzi!parité}
\index{feuillet!échange}

Barbero--Immirzi et l'entropie des deux feuillets ne sont pas deux quantités
juxtaposées. Ils naissent des mêmes canaux \(q_\pm\), mais conservent des données
distinctes. La normalisation probabiliste oublie l'échelle commune et mesure la
répartition entre orientations ; la fonctionnelle hexade--octade conserve le
signe de cette orientation. La première lecture est paire et la seconde impaire.

\section{Distribution normalisée entre feuillets}

Écrivons
\[
 x=A_{\rm rad},
 \qquad
 y=C^*_{\rm rad},
 \qquad
 q_-=e^{-x+y},
 \qquad
 q_+=e^{-x-y}.
\]
La somme \(q_-+q_+=2e^{-x}\cosh y\) permet de définir
\[
 p_-=\frac{q_-}{q_-+q_+}
     =\frac{e^y}{2\cosh y},
 \qquad
 p_+=\frac{q_+}{q_-+q_+}
     =\frac{e^{-y}}{2\cosh y}.
\]
Le facteur commun \(e^{-x}\) disparaît. La distribution conserve l'asymétrie
entre feuillets, mais non l'échelle angulaire commune.

\begin{definition}[Entropie bicouche]
\label{def:entropia-bicapa-rev6}
L'entropie de la distribution orientée est
\[
 S_{\rm bi}(y)
 :=
 -p_-\log p_--p_+\log p_+
 =
 \log(2\cosh y)-y\tanh y.
\]
\end{definition}

\begin{proposition}[Parité et maximum de l'entropie bicouche]
\label{prop:paridad-entropia-bicapa-rev6}
La fonction \(S_{\rm bi}\) est paire et satisfait
\[
 S_{\rm bi}(-y)=S_{\rm bi}(y),
 \qquad
 S_{\rm bi}'(y)=-y\,\operatorname{sech}^2y,
 \qquad
 0<S_{\rm bi}(y)\leq\log2.
\]
Le maximum \(\log2\) est atteint uniquement en \(y=0\), lorsque les deux feuillets
ont le même poids.
\end{proposition}

\begin{proof}
La parité découle du fait que \(\cosh\) est paire et que \(y\tanh y\) l'est également.
La dérivation directe produit
\[
 \frac{\dd}{\dd y}\log(2\cosh y)=\tanh y,
 \qquad
 \frac{\dd}{\dd y}(y\tanh y)
 =\tanh y+y\operatorname{sech}^2y.
\]
Par conséquent \(S_{\rm bi}'(y)=-y\operatorname{sech}^2y\) : la fonction croît sur
\((-\infty,0)\), décroît sur \((0,\infty)\) et prend à l'origine la valeur
\(\log2\). La positivité est celle de l'entropie d'une distribution à deux
poids strictement positifs.
\end{proof}

\section{Parité du transport hexade--octade}

Soit
\[
 s_n(q_-,q_+)=q_-^n-q_+^n.
\]
L'échange des feuillets est l'involution
\[
 \mathsf J_{\rm sh}:(x,y)\longmapsto(x,-y),
\]
qui permute \(q_-\leftrightarrow q_+\). D'où,
\[
 s_n\circ\mathsf J_{\rm sh}=-s_n.
\]
La fonctionnelle construite dans le \cref{mdg:chap:barbero} est
\[
 \Gamma_{6\to8}(A,C^*)
 =
 \frac{\pi A C^*}{180}
{}+12\bigl(S_{90}(q_-)-S_{90}(q_+)\bigr)
-\bigl(S_{120}(q_-)-S_{120}(q_+)\bigr).
\]

\begin{theorem}[Théorème de parité bicouche]
\label{thm:paridad-bicapa-barbero-rev6}
Sous l'échange \(C^*\mapsto-C^*\),
\[
 \boxed{S_{\rm bi}(-C^*_{\rm rad})
        =S_{\rm bi}(C^*_{\rm rad}),}
\]
\[
 \boxed{\Gamma_{6\to8}(A,-C^*)
        =-\Gamma_{6\to8}(A,C^*).}
\]
La même bicouche produit donc une lecture informationnelle paire et un
transport géométrique impair.
\end{theorem}

\begin{proof}
La première égalité est la
\cref{prop:paridad-entropia-bicapa-rev6}. Dans la seconde, le terme
\(\pi AC^*/180\) change de signe et chaque différence de Lambert change de
signe parce que ses deux arguments sont échangés. La combinaison complète est
impaire.
\end{proof}

\section{Interprétation de l'opérateur constitutif comme couplage entre vide, champ et réponse}

L'entropie bicouche et \(\Gamma_{6\to8}\) ne sont pas deux noms pour le même
nombre. La première quantifie le poids que partagent les deux orientations
après annulation de l'échelle commune. Le second conserve l'orientation et
l'incidence du passage \(6\to8\). Au point HMT,
\[
 S_{\rm bi}
 =0.69246069960358548\ldots\ {\rm nats},
 \qquad
 \Gamma_{6\to8}
 =0.27400767269445927\ldots.
\]
Le lien entre information et géométrie consiste en ce que les deux lectures
se factorisent par les mêmes canaux \(q_\pm\), et non à égaler leurs valeurs.

La bande particule--antiparticule de \cref{chap:banda-particula-virtual} reproduit la même distinction de parité : la masse et la quantité de corrélation descendent au quotient de bande, tandis que charge, orientation et transport changent de signe. Dans l'action de Cartan--Holst, \(\Gamma_{6\to8}\) occupe la coordonnée orientée ; dans la lecture informationnelle, \(S_{\rm bi}\) compte la distribution entre feuillets. Cela explique pourquoi une seule structure peut conserver l'information et inverser l'orientation sans confondre ces deux opérations. \endgroup

\paragraph{Variation tangente de l’état modulaire et dérivation de la loi aréale.} La variation tangente de l'état modulaire produit la première loi aréale ; les variations finies conservent le terme d'entropie relative. \endgroup

\chapter{Première loi modulaire aréale et extension finie par l'entropie relative}
\begingroup
\renewcommand{\chapter}[2][]{}%
\chapter{1re loi modulaire aréale et extension finie par l’entropie relative}
\label{chap:parte-iv-61}

La première loi modulaire relie les variations d’entropie et d’aire dans le système fini déclaré. Son extension non linéaire conserve le défaut positif d’entropie relative et ne confond donc pas égalité tangente et identité globale.

\section{1re loi modulaire pour les variations de l’aire}

Définissons les deux aires normalisées
\begin{equation}
\label{rm:eq:modular-normalized-area-channels}
\mathcal A_m=\gammaH a_0N_m,
\qquad m\in\{90,120\}.
\end{equation}
Alors
\begin{equation}
\label{rm:eq:modular-beta-definition}
\HHM^{(A)}-cI
=\beta_{90}\mathcal A_{90}+\beta_{120}\mathcal A_{120},
\qquad
\beta_m=\frac{m\Dmod}{\gammaH a_0}.
\end{equation}
Les valeurs terminales sont
\begin{align}
\beta_{90}&=5.9726485401070929914651798889\ldots,
\label{rm:eq:modular-beta90}\\
\beta_{120}&=7.9635313868094573219535731852\ldots,
\qquad
\frac{\beta_{120}}{\beta_{90}}=\frac43.
\label{rm:eq:modular-beta120}
\end{align}

\begin{theorem}[Première loi modulaire en variables aréales]
\label{rm:thm:modular-first-law}
Pour une famille différentiable d’états normalisés \(\rho(\varepsilon)\) passant par \(\rho_A\), tout en maintenant fixes \(\Dmod,\gammaH,a_0\),
\begin{equation}
\label{rm:eq:modular-first-law}
\boxed{
\delta S
=\beta_{90}\,\delta\langle\mathcal A_{90}\rangle
+\beta_{120}\,\delta\langle\mathcal A_{120}\rangle.}
\end{equation}
\end{theorem}

\begin{proof}
La première loi d’intrication dans l’état de référence est \(\delta S=\delta\langle-\log\rho_A\rangle\). Le terme scalaire \(cI\) ne contribue pas puisque \(\delta\Tr\rho=0\). Substituer \cref{rm:eq:modular-beta-definition} prouve l’identité.
\end{proof}

La formule est locale dans l’espace des états. Pour des changements finis, l’entropie relative apparaît ; l’omettre transformerait une identité tangente en une assertion fausse de linéarité globale.

\Needspace{22\baselineskip}
\begin{theorem}[Identité finie avec terme de défaut]
\label{rm:thm:modular-finite-first-law}
Pour tout état \(\sigma\) dont le support est contenu dans celui de \(\rho_A\),
\begin{equation}
\label{rm:eq:modular-finite-first-law}
\boxed{
S(\sigma)-S(\rho_A)
=\beta_{90}\Delta_\sigma\langle\mathcal A_{90}\rangle
+\beta_{120}\Delta_\sigma\langle\mathcal A_{120}\rangle
-D(\sigma\Vert\rho_A),}
\end{equation}
où \(D(\sigma\Vert\rho_A)=\Tr\sigma(\log\sigma-\log\rho_A)\geq0\).
En particulier, la variation linéaire de \cref{rm:thm:modular-first-law} est le terme tangent de cette identité.
\end{theorem}

\begin{proof}
Par définition et par \(\HHM^{(A)}=-\log\rho_A\),
\[
D(\sigma\Vert\rho_A)
=-S(\sigma)+\Tr(\sigma\HHM^{(A)}).
\]
Pour \(\sigma=\rho_A\), le membre de gauche est nul et \(S(\rho_A)=\Tr(\rho_A\HHM^{(A)})\). La soustraction des deux identités donne
\[
S(\sigma)-S(\rho_A)
=\Delta_\sigma\langle\HHM^{(A)}\rangle
-D(\sigma\Vert\rho_A).
\]
Le scalaire \(cI\) s’annule entre états normalisés ; substituer \cref{rm:eq:modular-beta-definition} produit \cref{rm:eq:modular-finite-first-law}. La positivité de l’entropie relative est l’inégalité de Klein.
\end{proof}

La correction \(D(\sigma\Vert\rho_A)\) mesure la distance informationnelle entre l’état réel et l’état modulaire de référence. Ainsi, la première loi ne perd pas son contenu au-delà du régime infinitésimal : elle devient une identité exacte avec un défaut positif contrôlé.
 
\paragraph{Persistance inductive des poids modulaires et classification factorielle de la limite.} La persistance inductive des poids modulaires conduit à la fermeture opératorielle infinie et à sa classification factorielle. \endgroup

\chapter{Persistance des incidences \texorpdfstring{\(90/120\)}{90/120} et limite factorielle hyperfinie}
\begingroup
\renewcommand{\chapter}[2][]{}%
\chapter{Persistance des canaux \texorpdfstring{\(90/120\)}{90/120} et classification de la limite comme facteur hyperfini \texorpdfstring{\(III_{\lambda_\partial}\)}{III(lambda partiel)}}
\label{chap:parte-iv-62}

Les produits locaux fidèles conservent les rapports des canaux \(90\) et \(120\) sous raffinement. La représentation GNS et la limite ITPFI déterminent le facteur hyperfini de type III avec un paramètre de bord explicite.

\section{Limite locale et classification modulaire du système infini}

Dans une UHF infinie, il n’existe généralement pas de matrice densité globale à trace à l’intérieur de l’algèbre. L’objet correct est le réseau
\begin{equation}
\label{rm:eq:modular-local-net}
\HHM^{(m)}=-\log\rho_m
\end{equation}
avec des restrictions d’états compatibles, ou le générateur modulaire dans la représentation GNS. L’espérance traciale \(\mathrm{id}\otimes\tau_9\) ne préserve pas automatiquement un état de Gibbs non tracial et ne se substitue pas à la compatibilité des états.

\begin{theorem}[Type modulaire sous persistance fidèle des deux canaux]
\label{rm:thm:modular-typeIII}
Supposons que :
\begin{enumerate}[label=(\roman*)]
\item chaque coupe est un produit d’états qutrit fidèles de poids \(w_{90}\) et \(w_{120}\) ;
\item les inclusions préservent l’état et l’étiquette de canal ;
\item une infinité de facteurs de chaque type apparaît, les deux types ayant une densité asymptotique positive ;
\item aucun quotient spectral supplémentaire n’est introduit et la représentation produit est factorielle.
\end{enumerate}
Alors le groupe additif des logarithmes des quotients de valeurs propres est
\begin{equation}
\label{rm:eq:modular-ratio-group}
90\Dmod\Z+120\Dmod\Z=30\Dmod\Z.
\end{equation}
La fermeture GNS ITPFI est le facteur hyperfini
\begin{equation}
\label{rm:eq:modular-typeIII-lambda}
\boxed{
III_{\lambda_\partial},
\qquad
\lambda_\partial=\ee^{-30\Dmod}
=0.9966696464689573786108299769\ldots .}
\end{equation}
\end{theorem}

\begin{proof}
Dans chaque lien, les quotients de valeurs propres sont des puissances de \(\ee^{-w_m}\). La persistance infinie des deux types engendre le sous-groupe \(w_{90}\Z+w_{120}\Z\). Comme \(\gcd(90,120)=30\), on obtient \cref{rm:eq:modular-ratio-group}. Le critère de rapport asymptotique d’Araki--Woods pour le produit fidèle donne le facteur ITPFI hyperfini de paramètre \(\ee^{-30\Dmod}\) \cite{rm:ArakiWoods1968,rm:Takesaki}.
\end{proof}

Le théorème déclare ses hypothèses effectives. Si l’un des canaux disparaît, si le produit n’est pas compatible ou si le groupe des rapports change, le type doit être recalculé.
 
\paragraph{Expression gravitationnelle de la fermeture modulaire au moyen du support à cinq composantes.} La fermeture modulaire prépare l'expression gravitationnelle : on construit d'abord le porteur à cinq composantes à partir de l'incidence icosaédrique. \endgroup

\chapter{Descente icosaédrique et porteur gravitationnel à cinq composantes}
\begingroup
\renewcommand{\chapter}[2][]{}%
\chapter{Descente icosaédrique \texorpdfstring{\(A_5/C_5\)}{A₅/C₅} et sélection du support gravitationnel à cinq composantes}
\label{chap:parte-iv-63}

La droite gravitationnelle conserve la longueur construite et sa réalisation chronologique de \cref{g26:sec:rigidez-radial-es,cr28:sec:radial-clock} ; son porteur descend ensuite au quotient icosaédrique. Le projecteur irréductible de rang \(5\) sélectionne le porteur tensoriel sans supposer encore une réalisation massive ou radiative.

\section{Constante gravitationnelle intrinsèque et représentation tensorielle}
\label{rm:ch:gravity}

\subsection{Chaîne constitutive du secteur gravitationnel : échelle chronologique \texorpdfstring{\(108\)}{108}, réduction tensorielle, holonomie et évaluation métrologique de \texorpdfstring{\(G\)}{G}}

La gravité s'organise ici en quatre strates qu'il ne faut pas fusionner :

\begin{equation}
\label{rm:eq:gravity-causal-chain}
\begin{aligned}
\text{inscription et norme radiale R1--R8}
&\longrightarrow L
\longrightarrow G=c_{\rm int}^3L^2/\hbar_a,\\
\text{incidence icosaédrique}
&\longrightarrow V_{12}\xrightarrow{P_5}V_5
\xrightarrow{\Gamma_5}\operatorname{STF}_3
\xrightarrow{\Lambda(n)}\mathcal H_{\rm TT}(n),\\
\text{calendrier CT108 avec }t_0=\pi L/(54c_{\rm int})
&\longrightarrow R_{108}=L,\quad
\theta_{108}^{\rm circ}=(54/\pi)^2,\\
\text{carte métrologique}
&\longrightarrow 6.6742996594140227\ldots\times10^{-11}\,
\mathrm{m^3\,kg^{-1}\,s^{-2}}.
\end{aligned}
\end{equation}

La première ligne conserve la détermination de la longueur et du couplage au moyen des règles R1--R8 de \cref{g26:sec:rigidez-radial-es}. La deuxième construit le porteur et sa réduction transverse sans trace. La troisième est la réalisation chronologique de la longueur déjà construite, et la quatrième évalue la même section dimensionnelle. La réduction \(12\to5\to5\to2\) ne modifie pas le coefficient radial commun : intensité et représentation conservent leur provenance et leurs fonctions respectives.

\subsection{Droite de grandeur gravitationnelle et échelle chronologique \texorpdfstring{\(108\)}{108}}

Soient \(c_{\rm int}>0\), \(t_0>0\),
\(\ell_0=c_{\rm int}t_0\) et une section positive
\(\hbar_a\) de la droite d'action, avec
\(a\in\{\mathrm{pre},\mathrm{ret}\}\). Le transducteur dimensionnel est

\begin{equation}
\label{rm:eq:gravity-transducer}
G_a(\theta)
=\theta\frac{c_{\rm int}^{5}t_0^2}{\hbar_a}
=\theta\frac{c_{\rm int}^{3}\ell_0^2}{\hbar_a},
\qquad \theta>0.
\end{equation}

L'égalité des deux expressions utilise seulement \(\ell_0=c_{\rm int}t_0\) et \([G]=L^3M^{-1}T^{-2}\). La dimension est fixée par la droite ; le caractère adimensionnel \(\theta\) distingue les réalisations de ce système. Dans la composition longitudinale déjà démontrée, \(t_0=\pi L/(54c_{\rm int})\) et \(\theta=(54/\pi)^2\) reproduisent exactement \(G_a=c_{\rm int}^3L^2/\hbar_a\).

Le calendrier CT108 définit

\begin{equation}
\label{rm:eq:gravity-clock108}
T_{108}=108t_0,
\qquad
\omega_{108}=\frac{2\pi}{108t_0},
\qquad
R_{108}=\frac{c_{\rm int}}{\omega_{108}}
=\frac{54}{\pi}\ell_0.
\end{equation}

Pour chaque feuillet d'action, soient

\begin{equation}
\label{rm:eq:gravity-clock-energy-mass}
E_{108,a}=\hbar_a\omega_{108},
\qquad
m_{108,a}=\frac{E_{108,a}}{c_{\rm int}^2}.
\end{equation}

Alors, la longueur de Compton réduite coïncide avec le rayon du cycle :

\begin{equation}
\label{rm:eq:gravity-compton-clock}
\bar\lambda_{C,a}
=\frac{\hbar_a}{m_{108,a}c_{\rm int}}
=R_{108}.
\end{equation}

\begin{theorem}[Unicité du sélecteur sous la condition de périodicité HMT]
\label{rm:thm:gravity-circular-selector}
Avec la convention réduite
\(r_g=Gm/c_{\rm int}^2\) et la prémisse de périodicité HMT qui identifie le rayon
gravitationnel du quantum CT108 à son rayon de phase, la condition de clôture
\begin{equation}
\label{rm:eq:gravity-closing-condition}
r_{g,a}(\theta)=R_{108}
\end{equation}
sélectionne une unique solution positive,
\begin{equation}
\label{rm:eq:gravity-theta108}
\boxed{
\theta_{108}^{\rm circ}=\left(\frac{54}{\pi}\right)^2
=295.4525715010569415303525\ldots .}
\end{equation}
Dans cette réalisation,
\begin{equation}
\label{rm:eq:gravity-four-lengths}
R_{108}=\bar\lambda_{C,a}=r_{g,a}=\ell_P^{\rm circ}.
\end{equation}
\end{theorem}

\begin{proof}
Par \cref{rm:eq:gravity-transducer,rm:eq:gravity-clock-energy-mass},
\[
r_{g,a}(\theta)
=\frac{G_a(\theta)m_{108,a}}{c_{\rm int}^2}
=\theta c_{\rm int}t_0^2\omega_{108}
=\theta\frac{\pi}{54}\ell_0.
\]
En le comparant à
\(R_{108}=(54/\pi)\ell_0\), on obtient nécessairement
\(\theta=(54/\pi)^2\). L'équation est linéaire en \(\theta\) et tous les
facteurs sont positifs, donc la solution est unique.
\end{proof}

L'égalité \(r_g=R_{108}\) ne se déduit pas de la dimensionnalité de
\cref{rm:eq:gravity-transducer}. C'est la condition géométrique explicite de cette
réalisation HMT : action, Compton et retour périodique doivent déterminer une même
longueur. Le théorème démontre que, cette condition structurelle étant acceptée, le
caractère \(\theta\) est fixé sans utiliser la valeur SI de \(G\).

\begin{remark}[Convention de rayon]
Le théorème utilise le rayon gravitationnel réduit \(Gm/c^2\). Si on le remplace
par le rayon de Schwarzschild \(2Gm/c^2\), le sélecteur change d'un facteur
deux. Les deux expressions correspondent à des conditions géométriques distinctes,
c'est pourquoi la convention doit être déclarée avant le calcul.
\end{remark}

\subsection{Système de Planck et covariance bidirectionnelle}

\begin{corollary}[Famille de Planck associée à la structure chronologique CT108]
\label{rm:cor:gravity-planck-family}
Pour \(G_a=G_a(\theta_{108}^{\rm circ})\),
\begin{align}
\ell_P&=\frac{54}{\pi}\ell_0,
&t_P&=\frac{54}{\pi}t_0,
\label{rm:eq:gravity-planck-length-time}\\
E_{P,a}&=\frac{\hbar_a}{t_P}
=\frac{\pi\hbar_a}{54t_0}
=\frac{h_a}{108t_0},
&\ell_P^2&=\frac{G_a\hbar_a}{c_{\rm int}^3}
=\theta_{108}^{\rm circ}\ell_0^2.
\label{rm:eq:gravity-planck-energy-area}
\end{align}
De plus,
\begin{equation}
\label{rm:eq:gravity-planck-force-power-density}
F_{P,a}=\frac{c_{\rm int}^4}{G_a},
\qquad
P_{P,a}=\frac{c_{\rm int}^5}{G_a},
\qquad
\rho_{P,a}=\frac{\hbar_a}
{(\theta_{108}^{\rm circ})^2c_{\rm int}^5t_0^4}.
\end{equation}
\end{corollary}

\begin{proof}
Les trois premières identités découlent de
\(\ell_P=R_{108}\), \(t_P=\ell_P/c_{\rm int}\) et
\(h_a=2\pi\hbar_a\). Pour l'aire,
\[
\frac{G_a\hbar_a}{c_{\rm int}^3}
=\theta_{108}^{\rm circ}c_{\rm int}^2t_0^2
=\theta_{108}^{\rm circ}\ell_0^2.
\]
Les expressions restantes sont les compositions dimensionnelles de force,
de puissance et de masse par volume avec la même section.
\end{proof}

Soit
\begin{equation}
\label{rm:eq:gravity-Ract}
R_{\rm act}:=\frac{\hbar_{\rm pre}}{\hbar_{\rm ret}}.
\end{equation}
Comme l'action apparaît au numérateur de la masse chronologique et au
dénominateur de \(G\),
\begin{equation}
\label{rm:eq:gravity-twoway-covariance}
\frac{m_{108,\rm pre}}{m_{108,\rm ret}}=R_{\rm act},
\qquad
\frac{G_{\rm pre}}{G_{\rm ret}}=R_{\rm act}^{-1},
\qquad
G_a\hbar_a
=\theta_{108}^{\rm circ}c_{\rm int}^5t_0^2.
\end{equation}
Par conséquent, \(\ell_P^2=G_a\hbar_a/c^3\) ne change pas de feuillet. La gravité et
l'action s'orientent réciproquement, tandis que l'aire conserve la
géométrie commune. Cette annulation est l'interface de confluence vers l'opérateur
d'aire du \cref{rm:ch:modular}.

\subsection{Composante icosaédrique de dimension \(5\)}

L’antécédent n’est ni \(\R\) ni un choix ultérieur de cinq polarisations. L’objet chargé de transporter l’incidence est défini explicitement avant l’introduction de la réalisation exceptionnelle.

Soit \(X_{\rm TPK}^{\rm exc}\) le sous-espace de l’état \(\TPK\) enrichi
qui conserve, outre la sortie arithmétique, l’orientation, la route,
la retenue et l’incidence exceptionnelle. Pour expliciter l’espace géométrique d’arrivée,
soit \(\varphi_g=(1+\sqrt5)/2\), \(\nu=\sqrt{1+\varphi_g^2}\), et définissons
\begin{equation}
\label{rm:eq:gravity-icosahedron-vertices}
\Omega_{12}^{\rm ico}
=\frac1\nu\bigl(
\{0\}\times\{\pm1\}\times\{\pm\varphi_g\}
\;\cup\;
\{\pm1\}\times\{\pm\varphi_g\}\times\{0\}
\;\cup\;
\{\pm\varphi_g\}\times\{0\}\times\{\pm1\}
\bigr)\subset S^2.
\end{equation}
Ses douze éléments sont les classes orientées de sommet d’un icosaèdre.
L’incidence ne reste pas verbale : elle est fixée par
\begin{equation}
\label{rm:eq:gravity-icosahedron-incidence}
\mathcal I_{12}^{\rm ico}
=\{(v,w)\in\Omega_{12}^{\rm ico}\times\Omega_{12}^{\rm ico}:
v\neq w,\ \langle v,w\rangle=1/\sqrt5\}.
\end{equation}
Chaque sommet a cinq voisins dans cette relation. Le pont de lecture
est typé comme une application
\begin{equation}
\label{rm:eq:gravity-exceptional-reader-map}
\pi_{\rm ico}:X_{\rm TPK}^{\rm exc}\longrightarrow\Omega_{12}^{\rm ico}.
\end{equation}
La cardinalité douze de l’image de \(\pi_{\rm ico}\) ne suffit pas. Pour définir une réalisation d’incidence, l’application doit satisfaire simultanément :
\begin{enumerate}[label=(\roman*)]
\item être surjective et constante exactement sur les fibres qui écartent
la carte arithmétique mais conservent la route exceptionnelle ;
\item entrelacer l’action des automorphismes orientés de l’antécédent
avec une action transitive sur \(\Omega_{12}^{\rm ico}\) ;
\item transporter l’incidence TPK déclarée vers
\(\mathcal I_{12}^{\rm ico}\), et non
ses seuls cardinaux ;
\item conserver l’involution de feuillet et le retour nonadique dans le
stabilisateur de la fibre correspondante.
\end{enumerate}

Ces conditions séparent trois assertions : l’existence de l’application \(\pi_{\rm ico}\), l’identification du groupe qui agit et la représentation linéaire construite ensuite. La deuxième et la troisième ne peuvent servir à définir rétrospectivement la première.

\begin{theorem}[Descente de l’incidence au quotient icosaédrique]
\label{rm:thm:gravity-exceptional-descent}
Supposons que \(\pi_{\rm ico}\) satisfait (i)--(iv) et que l’action
orientée induite sur \(\Omega_{12}^{\rm ico}\) a pour groupe
\(G_{\rm ico}\simeq A_5\). Alors, pour toute classe
\(v_0\in\Omega_{12}^{\rm ico}\),
\begin{equation}
\label{rm:eq:gravity-icosahedral-quotient}
H_{v_0}=\operatorname{Stab}_{G_{\rm ico}}(v_0)\simeq C_5,
\qquad
\Omega_{12}^{\rm ico}\simeq G_{\rm ico}/H_{v_0}\simeq A_5/C_5.
\end{equation}
L’application
\begin{equation}
\label{rm:eq:gravity-coset-map}
\kappa_{v_0}:G_{\rm ico}/H_{v_0}\longrightarrow\Omega_{12}^{\rm ico},
\qquad gH_{v_0}\longmapsto g\cdot v_0,
\end{equation}
est une bijection équivariante et préserve l’incidence si l’on déclare sur le quotient
\begin{equation}
\label{rm:eq:gravity-coset-incidence}
gH_{v_0}\sim g'H_{v_0}
\quad\Longleftrightarrow\quad
(g\cdot v_0,g'\cdot v_0)\in\mathcal I_{12}^{\rm ico}.
\end{equation}
Par conséquent, la composée \(\kappa_{v_0}^{-1}\pi_{\rm ico}\) est l’application
typée de l’antécédent TPK vers le quotient \(A_5/C_5\).
\end{theorem}

\begin{proof}
La liste de \cref{rm:eq:gravity-icosahedron-vertices} a douze éléments.
Par transitivité, \(|G_{\rm ico}\cdot v_0|=12\), et le théorème
orbite--stabilisateur donne \(|H_{v_0}|=60/12=5\). Tout groupe d’ordre
premier cinq est cyclique ; donc \(H_{v_0}\simeq C_5\). La formule de
\(\kappa_{v_0}\) ne dépend pas du représentant : si
\(gH_{v_0}=g'H_{v_0}\), alors \(g^{-1}g'\in H_{v_0}\) et
\(g'v_0=gv_0\). La surjectivité découle de la transitivité, et
l’égalité des cardinaux donne l’injectivité. L’équivariance est immédiate :
\(\kappa_{v_0}(agH_{v_0})=a\kappa_{v_0}(gH_{v_0})\). Enfin,
\cref{rm:eq:gravity-coset-incidence} est précisément le transport de
structure par cette bijection ; l’hypothèse (iii) fait commuter ce
transport avec \(\pi_{\rm ico}\).
\end{proof}

Le théorème ne remplace pas la construction de \(\pi_{\rm ico}\) : il démontre ce qui est obtenu une fois ses quatre conditions vérifiées. Dans le profil HMT complet, le corridor d’incidence qui mène à Paley--Witt--Golay--Leech--\(V^\natural\)--Monstre--Moonshine est le domaine de définition de cette application. Dans ce volume autonome, l’entrée exacte est l’application \cref{rm:eq:gravity-exceptional-reader-map} assortie de ses conditions ; toute la chaîne depuis \(A_5/C_5\) jusqu’à la réduction TT est prouvée ci-dessous sans consulter à nouveau une valeur gravitationnelle.

La chaîne de provenance s’exprime au moyen de tous ses morphismes :
\begin{equation}
\label{rm:eq:gravity-exceptional-provenance}
\APP\longrightarrow\TRIT\longrightarrow\TPK
\longrightarrow X_{\rm TPK}^{\rm exc}
\xrightarrow{\ \pi_{\rm ico}\ }\Omega_{12}^{\rm ico}
\xleftarrow[\simeq]{\ \kappa_{v_0}\ }A_5/C_5.
\end{equation}
Le corridor exceptionnel complet est démontré dans son développement propre. L’hypothèse que ce chapitre ne peut remplacer par un dénombrement est l’existence de \(\pi_{\rm ico}\) satisfaisant (i)--(iv). Cette entrée étant fixée, la descente au quotient est \cref{rm:thm:gravity-exceptional-descent}, et l’on construit à partir d’elle \(P_5\), \(\Gamma_5\) et \(\Lambda(n)\). La réalisation du support résultant comme champ gravitationnel demeure une interface physique typée, et non une conséquence de cardinalité.

Soit \(G=A_5\), soit \(H\simeq C_5\) le stabilisateur d’un sommet
orienté de l’icosaèdre et considérons
\begin{equation}
\label{rm:eq:gravity-v12}
V_{12}=\R[G/H],
\qquad |G/H|=\frac{60}{5}=12.
\end{equation}
La représentation de permutation \(\rho_{12}\) se décompose comme
\begin{equation}
\label{rm:eq:gravity-v12-decomposition}
V_{12}\simeq\mathbf1\oplus V_3\oplus V_{3'}\oplus V_5.
\end{equation}
L’idempotent central du facteur de caractère
\(\chi_5=(5,1,-1,0,0)\) est
\begin{equation}
\label{rm:eq:gravity-p5}
\boxed{
P_5=\frac5{60}\sum_{g\in A_5}\chi_5(g^{-1})\rho_{12}(g)
=\frac1{12}\left(
5I+\sum_{g\in2A}\rho_{12}(g)
-\sum_{g\in3A}\rho_{12}(g)
\right).}
\end{equation}

\begin{theorem}[Projecteur et couplage gravitationnel de rang cinq]
\label{rm:thm:gravity-p5-coupling}
L'opérateur \(P_5\) satisfait
\begin{equation}
\label{rm:eq:gravity-p5-properties}
P_5^2=P_5=P_5^*,
\qquad \Tr P_5=\rank P_5=5.
\end{equation}
Pour
\begin{equation}
\label{rm:eq:gravity-G5}
\mathbb G_{5,a}:=G_a(\theta_{108}^{\rm circ})P_5
\end{equation}
on a
\begin{equation}
\label{rm:eq:gravity-G5-spectrum}
\Spec(\mathbb G_{5,a})
=\{G_a(\theta_{108}^{\rm circ})\text{ avec multiplicité }5,
\;0\text{ avec multiplicité }7\}.
\end{equation}
\end{theorem}

\begin{proof}
L’orthogonalité des caractères fait de \cref{rm:eq:gravity-p5}
l’idempotent central qui agit comme l’identité sur l’unique copie de
\(V_5\) et comme zéro sur les trois autres facteurs de
\cref{rm:eq:gravity-v12-decomposition}. La réalité du caractère et
l’orthogonalité de \(\rho_{12}\) donnent l’autoadjonction. Sa trace et son rang
sont cinq. Multiplier un projecteur de rang cinq par le scalaire positif
\(G_a\) produit immédiatement le spectre indiqué.
\end{proof}

\begingroup\fontsize{10.2}{11.7}\selectfont \setlength{\parskip}{2pt}\setlength{\abovedisplayskip}{5pt}\setlength{\belowdisplayskip}{5pt} L'expression \(\mathbb G_{5,a}\) n'ajoute pas cinq valeurs différentes de la constante gravitationnelle. C'est une même section de \(G\) agissant sur le facteur tensoriel sélectionné et annulant les sept autres modes du module de permutation.
 
\paragraph{Réduction équivariante du support au secteur tensoriel sans trace.} Le support est entrelacé avec les tenseurs symétriques sans trace et réduit de manière équivariante avant d’imposer la transversalité.

\par\endgroup
 \endgroup

\chapter{Entrelaceur tensoriel et réduction transverse \texorpdfstring{\(12\to5\to5\to2\)}{12->5->5->2}}
\begingroup
\renewcommand{\chapter}[2][]{}%
\chapter{Entrelaceur avec les tenseurs symétriques sans trace et réduction transverse \texorpdfstring{\(12\to5\to5\to2\)}{12→5→5→2}}
\label{chap:parte-iv-64}

Le repère icosaédrique construit une isométrie équivariante vers l’espace des tenseurs symétriques sans trace. La projection transverse réalise \(12\to5\to5\to2\), conserve le support à cinq composantes et localise le sous-espace radiatif de dimension \(2\).

\section{Entrelaceur avec l’espace des tenseurs symétriques sans trace}

Soit \(\mathcal V\subset S^2\) l’ensemble des douze sommets unitaires de l’icosaèdre. Pour chaque \(v\in\mathcal V\), définissons
\begin{equation}
\label{rm:eq:gravity-Qv}
Q_v=vv^T-\frac13I_3\in
\operatorname{STF}_3:=\operatorname{Sym}^2_0(\R^3).
\end{equation}
L’application de repère
\begin{equation}
\label{rm:eq:gravity-Gamma0}
\Gamma_0\!\left(\sum_{v\in\mathcal V}x_ve_v\right)
=\sum_{v\in\mathcal V}x_vQ_v
\end{equation}
est \(A_5\)-équivariante et satisfait
\begin{equation}
\label{rm:eq:gravity-tight-frame}
\Gamma_0\Gamma_0^*=\frac85I_{\operatorname{STF}_3},
\qquad
\Gamma_0=\Gamma_0P_5.
\end{equation}

\begin{theorem}[Entrelaceur tensoriel normalisé]
\label{rm:thm:gravity-gamma5}
L’application
\begin{equation}
\label{rm:eq:gravity-gamma5}
\boxed{
\Gamma_5:=\sqrt{\frac58}\,\Gamma_0|_{V_5}:
V_5\xrightarrow{\ \simeq\ }\operatorname{STF}_3}
\end{equation}
est une isométrie \(A_5\)-équivariante. Son adjoint est
\begin{equation}
\label{rm:eq:gravity-gamma5-adjoint}
\Gamma_5^*Q=\sqrt{\frac58}
\sum_{v\in\mathcal V}(v^TQv)e_v.
\end{equation}
\end{theorem}

\begin{proof}
L’équivariance découle de \(Q_{gv}=gQ_vg^T\). L’opérateur
\(\Gamma_0\Gamma_0^*\) commute avec la représentation irréductible de
\(A_5\) sur \(\operatorname{STF}_3\) ; il est donc scalaire. Sa trace
est
\[
\sum_{v\in\mathcal V}\|Q_v\|_F^2
=12\cdot\frac23=8;
\]
en divisant par la dimension cinq, on obtient \(8/5\). D’où
\(\Gamma_5\Gamma_5^*=I\). Comme le domaine et le codomaine sont de dimension
cinq, on a aussi \(\Gamma_5^*\Gamma_5=I\). La formule de l’adjoint
découle de \(\langle Q,Q_v\rangle_F=v^TQv\).
\end{proof}

La normalisation au moyen du repère icosaédrique élimine une coïncidence purement cardinale : elle construit l’application qui transporte le facteur fini vers le support continu de spin deux.

\section{Projection transversale sans trace}

Pour \(n\in S^2\), soit \(\Pi(n)=I-nn^T\). Nous définissons
\begin{equation}
\label{rm:eq:gravity-lambda-tt}
\boxed{
\Lambda(n)Q
=\Pi Q\Pi-\frac12\Tr(\Pi Q\Pi)\Pi.}
\end{equation}

\begin{theorem}[Réduction de cinq composantes à deux]
\label{rm:thm:gravity-tt}
La restriction de \(\Lambda(n)\) à \(\operatorname{STF}_3\) est le
projecteur orthogonal sur
\begin{equation}
\label{rm:eq:gravity-htt}
\mathcal H_{\rm TT}(n)
=\{Q\in\operatorname{STF}_3:Qn=0\}.
\end{equation}
En particulier,
\begin{equation}
\label{rm:eq:gravity-lambda-properties}
\Lambda(n)^2=\Lambda(n)=\Lambda(n)^*,
\qquad \rank\Lambda(n)=2.
\end{equation}
Pour \(n=e_3\), l’image est engendrée par
\begin{equation}
\label{rm:eq:gravity-plus-cross}
E_+=\frac1{\sqrt2}\operatorname{diag}(1,-1,0),
\qquad
E_\times=\frac1{\sqrt2}
\begin{pmatrix}0&1&0\\1&0&0\\0&0&0\end{pmatrix}.
\end{equation}
\end{theorem}

\begin{proof}
Comme \(\Pi n=0\), la sortie est transversale. Sa trace vaut
\[
\Tr(\Pi Q\Pi)-\frac12\Tr(\Pi)\Tr(\Pi Q\Pi)=0,
\]
car \(\Tr\Pi=2\). Si \(Qn=0\) et \(\Tr Q=0\), alors
\(\Pi Q\Pi=Q\) et \(\Tr(\Pi Q\Pi)=0\) ; donc
\(\Lambda(n)Q=Q\). La formule est autoadjointe pour le produit de
Frobenius. Dans le repère \(n=e_3\), sa matrice dans une base STF adaptée est
\(\operatorname{diag}(1,1,0,0,0)\), ce qui prouve le rang.
\end{proof}

La chaîne complète devient donc
\begin{equation}
\label{rm:eq:gravity-12-5-5-2}
\boxed{
V_{12}\xrightarrow{P_5}V_5
\xrightarrow{\Gamma_5}\operatorname{STF}_3
\xrightarrow{\Lambda(n)}\mathcal H_{\rm TT}(n),
\qquad 12\longrightarrow5\longrightarrow5\longrightarrow2.}
\end{equation}
Le couplage réduit
\begin{equation}
\label{rm:eq:gravity-effective-map}
\mathbb G_{{\rm TT},a}(n)
=G_a(\theta_{108}^{\rm circ})\Lambda(n)\Gamma_5P_5
\end{equation}
est de rang deux.
 
\paragraph{Décomposition en hélicités du spectre et séparation des modes massif et sans masse.} La réduction permet de distinguer le spectre complet d'hélicité d'une réalisation massive et ses deux hélicités sans masse. \endgroup

\chapter{Décomposition en hélicités du porteur de spin deux}
\begingroup
\renewcommand{\chapter}[2][]{}%
\chapter{Décomposition en hélicités du support de spin \texorpdfstring{\(2\)}{2} et distinction entre réalisations massive et sans masse}
\label{chap:parte-iv-65}

Le support réel de dimension \(5\) se décompose en \(5\) poids d’hélicité. La réalisation massive conserve les \(5\) composantes ; la réduction transverse sans masse ne conserve que les hélicités radiatives \(\pm2\), sans transformer une possibilité collective en particule fondamentale postulée.

\section{Décomposition en \(5\) modes de polarisation}

La dimension cinq acquiert un contenu tensoriel lorsque l’on choisit une direction
unitaire de propagation \(n\) et un repère orthonormé orienté
\((e_1,e_2,n)\). Nous définissons les cinq tenseurs réels
\begin{align}
E_{+}&=\frac1{\sqrt2}(e_1e_1^T-e_2e_2^T),
&E_{\times}&=\frac1{\sqrt2}(e_1e_2^T+e_2e_1^T),
\label{rm:eq:gravity-helicity2-real}\\
E_{1}&=\frac1{\sqrt2}(e_1n^T+ne_1^T),
&E_{2}&=\frac1{\sqrt2}(e_2n^T+ne_2^T),
\label{rm:eq:gravity-helicity1-real}\\
E_{0}&=\frac1{\sqrt6}(2nn^T-e_1e_1^T-e_2e_2^T).
\label{rm:eq:gravity-helicity0-real}
\end{align}
Ils sont tous symétriques et sans trace. Les deux premiers sont transversaux ; les
deux suivants mêlent le plan transversal à la direction de
propagation, et le dernier est le mode scalaire longitudinal sans trace.

\begin{theorem}[Décomposition hélicoïdale complète du porteur]
\label{rm:thm:gravity-five-helicities}
Les tenseurs de \cref{rm:eq:gravity-helicity2-real,rm:eq:gravity-helicity1-real,rm:eq:gravity-helicity0-real} forment une base orthonormale de \(\operatorname{STF}_3\). Dans la complexification, les combinaisons
\begin{equation}
\label{rm:eq:gravity-complex-helicities}
E_{\pm2}=\frac1{\sqrt2}(E_+\mp\ii E_\times),
\qquad
E_{\pm1}=\frac1{\sqrt2}(E_1\mp\ii E_2),
\qquad E_0=E_0
\end{equation}
sont des vecteurs de poids \(\pm2,\pm1,0\), respectivement, sous les rotations
autour de \(n\). De plus,
\begin{equation}
\label{rm:eq:gravity-tt-on-five}
\Lambda(n)E_{\pm2}=E_{\pm2},
\qquad
\Lambda(n)E_{\pm1}=0,
\qquad
\Lambda(n)E_0=0.
\end{equation}
\end{theorem}

\begin{proof}
Dans le repère \((e_1,e_2,n)\), chaque tenseur a une norme de Frobenius égale à un.
Leurs motifs d’entrées non diagonales sont disjoints, et les trois diagonales
de \(E_+,E_0\) ont un produit intérieur nul ; les cinq tenseurs
sont donc orthonormés. Comme \(\dim\operatorname{STF}_3=6-1=5\), ils forment une base.

Une rotation d’angle \(\psi\) autour de \(n\) envoie
\((e_1,e_2)\) sur
\((\cos\psi\,e_1+\sin\psi\,e_2,
-\sin\psi\,e_1+\cos\psi\,e_2)\). En substituant dans les formules, on obtient
une rotation d’angle \(2\psi\) sur
\(\operatorname{span}\{E_+,E_\times\}\), une d’angle \(\psi\) sur
\(\operatorname{span}\{E_1,E_2\}\), et \(E_0\) fixe. D’où les poids de
\cref{rm:eq:gravity-complex-helicities}. Enfin,
\(\Pi E_{+}\Pi=E_+\) et \(\Pi E_\times\Pi=E_\times\), tandis que
\(\Pi E_1\Pi=\Pi E_2\Pi=0\). Pour \(E_0\),
\(\Pi E_0\Pi=-(e_1e_1^T+e_2e_2^T)/\sqrt6\), qui est purement de trace dans
le plan, et la soustraction de \cref{rm:eq:gravity-lambda-tt} l’annule.
\end{proof}

Le théorème précise le sens de ``cinq polarisations'' sans leur attribuer
à toutes le même statut dynamique. Avant d’imposer les équations de champ,
le porteur de spin deux a cinq coordonnées irréductibles. Dans une
réalisation massive de Fierz--Pauli, les contraintes laissent précisément
les cinq poids de \cref{rm:eq:gravity-complex-helicities} ; dans une réalisation
sans masse avec invariance de jauge, les secteurs \(\pm1\) et \(0\) sont éliminés et
il reste \(\pm2\). HMT fournit ici le porteur, la base et le projecteur ;
le choix de la dynamique massive ou sans masse relève de l’interface
physique. Cinq ne signifie donc pas ``cinq ondes observées'', mais cinq
degrés tensoriels disponibles avant la réduction.

L’orientation reste également visible : remplacer
\((e_1,e_2,n)\) par \((e_1,-e_2,-n)\) conserve \(E_+,E_0\) et échange
les phases des paires hélicoïdales. Cette information ne se trouve pas dans le
scalaire \(G_a\) ; elle réside dans le porteur sur lequel agit
\(\mathbb G_{5,a}\).

\section{Chaîne représentationnelle \(12\to5\to5\to2\) : support \(V_5\), tenseur sans trace, polarisations de Fierz–Pauli et hélicités transverses}

La dimension cinq apparaît dans des objets distincts, et seul
\cref{rm:eq:gravity-12-5-5-2} autorise le transport entre eux :

\begin{enumerate}
\item \(V_5\) est le facteur irréductible de la représentation sur les
douze sommets. C’est un résultat fini de théorie des représentations.
\item \(\operatorname{STF}_3\) est le porteur réel de dimension cinq des
poids \(m=-2,-1,0,1,2\) de spin deux.
\item Une théorie massive de Fierz--Pauli peut réaliser ces cinq poids
comme cinq états physiques, après fixation de l’action et des contraintes.
\item La relativité générale sans masse conserve deux hélicités radiatives,
\(\pm2\), représentées par \(E_+\) et \(E_\times\).
\item Les classes \(E(2)\) d’ondes métriques, comme \(III_5\), classent
les amplitudes possibles sous d’autres hypothèses ; elles ne sont pas automatiquement
\(V_5\), Fierz--Pauli ou la relativité générale.
\end{enumerate}

En particulier, la formulation correcte est : \emph{HMT construit un support tensoriel à cinq composantes avant les contraintes ; la réduction TT de la réalisation sans masse est de rang deux}. Il n’est pas affirmé que la relativité générale possède cinq polarisations propagatives.
 
\paragraph{Conditions d’existence du mode par sélection holonomique naturelle.} L'existence physique du mode requiert encore un sélecteur holonomique naturel et falsifiable. \endgroup

\chapter{Sélecteur gravitationnel naturel et transducteur dimensionnel}
\section[Critères structurels et carte métrologique de \(G\)]{Critères de
sélection de la coordonnée gravitationnelle : holonomie enrichie, réponse
spectrale et carte métrologique}
\label{sec:selector-gravitatorio-rev11}
\index{constante gravitationnelle!transducteur et sélecteur}
\index{gravité!coordonnée métrologique}
\index{holonomie!sélecteur gravitationnel}

L'action de Palatini--Cartan--Holst et son Hamiltonien contraint fixent d'abord la dynamique. Sur celle-ci, la constante gravitationnelle s'organise en trois niveaux typés : la droite dimensionnelle qui abrite \(G\), le transducteur qui convertit une coordonnée adimensionnelle en section de cette droite et le sélecteur structurel de cette coordonnée. Le système de coordonnées décimal et l'incertitude expérimentale sont présentés à la fin comme représentation et comparaison.

\subsection{Droite dimensionnelle et transducteur interne}

Dans le réseau dimensionnel de la mécanique dimensionnelle,
\[
 \mathbf G=-\mathbf S+5\mathbf C+2\mathbf T,
 \qquad [G]=L^3M^{-1}T^{-2}.
\]
Soient
\[
 c_{\rm int}\in\mathcal L_C^+,
 \qquad t_0\in\mathcal L_T^+,
 \qquad \ell_0=c_{\rm int}t_0,
 \qquad \hbar_{\rm int}\in\mathcal L_S^+.
\]

\HMTClaim{MD-R-G-TRANSDUCER}
\begin{definition}[Transducteur gravitationnel]
\label{def:transductor-gravitatorio-rev11}
Pour \(\theta_G>0\), on définit
\begin{equation}
 \boxed{
 G_{\rm int}(\theta_G)
 =\theta_G\frac{c_{\rm int}^{5}t_0^2}{\hbar_{\rm int}}
 =\theta_G\frac{c_{\rm int}^{3}\ell_0^2}{\hbar_{\rm int}}
 \in\mathcal L_G^+.}
 \label{eq:transductor-gravitatorio-rev11}
\end{equation}
\end{definition}

L'égalité des deux expressions et le type de
\eqref{eq:transductor-gravitatorio-rev11} sont exacts. Si
\[
 L_G^2:=\frac{\hbar_{\rm int}G_{\rm int}}{c_{\rm int}^3},
\]
alors
\begin{equation}
 \theta_G=\left(\frac{L_G}{\ell_0}\right)^2.
 \label{eq:thetaG-razon-longitudes-rev11}
\end{equation}
La coordonnée gravitationnelle est donc un rapport quadratique de longueurs
ou une réponse d'holonomie normalisée. La décade
\(10^{-34}\mathcal U_S\) appartient déjà à \(\hbar_{\rm int}\) : elle fixe la section
absolue d'action et ne constitue pas un facteur relatif ajouté à \(G\).

\subsection{Torsion observable et caractères positifs}

Le retour observable de \(27\) pas
\[
 \mathcal K_{27}:=F_{\rm obs}^{27}
\]
satisfait
\[
 \mathcal K_{27}^{4}=F_{\rm obs}^{108}=I.
\]
Dans la coordonnée \(\theta\in\mathbb Z/108\mathbb Z\), l'ordre est
exactement quatre, car
\(\mathcal K_{27}(\theta)=\theta+27\),
\(\mathcal K_{27}^{2}(\theta)=\theta+54\neq\theta\) et
\(\mathcal K_{27}^{4}(\theta)=\theta\).

\HMTClaim{MD-R-G-POSITIVE-CHARACTER-NOGO}
\begin{theorem}[Trivialité du caractère positif observable]
\label{thm:caracter-positivo-G-rev11}
Tout homomorphisme multiplicatif
\[
 \chi:\langle\mathcal K_{27}\rangle
 \longrightarrow\mathbb R_{>0}^{\times}
\]
cumple \(\chi(\mathcal K_{27})=1\).
\end{theorem}

\begin{proof}
La multiplicativité implique
\(\chi(\mathcal K_{27})^4=\chi(I)=1\).
Le seul réel positif dont la puissance quatrième est un est \(1\).
\end{proof}

Le même cycle transporte bien les caractères unitaires
\[
 \chi_j(\mathcal K_{27})=\exp(2\pi ij/4),
 \qquad j=0,1,2,3.
\]
Le quotient observable admet une phase, tandis que sa torsion élimine une
échelle positive multiplicative non triviale. Une continuation multiplicative doit donc
agir sur un relèvement \(\widetilde{\mathcal K}_{27}\) dont la classe dans
l'abélianisation du groupe d'isotropie soit non de torsion, et doit sélectionner
naturellement tant le caractère que sa normalisation.

\subsection{Mémoire chronologique et réponse spectrale}
\label{sec:respuesta-espectral-G-rev11}

Le revêtement chronologique dirigé réalise le retour de \(27\) pas comme la
translation
\[
 (G_{\rm or},S)\longmapsto(G_{\rm or},S)+(1,18),
\]
d'ordre de semi-groupe infini. Cette voie conserve l'accumulation, mais sa
réponse spectrale requiert une interférence entre histoires.

Soit \(\Gamma=(V,E)\) un multigraphe orienté fini muni d'une mesure positive,
de poids \(a_\ast(e)\geq0\) et d'un cocycle \(c(e)\in\mathbb Z^2\). L'opérateur
tordu
\[
 (\mathcal L_{\beta,\vartheta}f)(y)
 =\sum_{e:x\to y}
 e^{-\beta a_\ast(e)}e^{i\langle\vartheta,c(e)\rangle}f(x)
\]
induit l'opérateur positif
\(\mathcal H_{\beta,\vartheta}
=\mathcal L_{\beta,\vartheta}^{*}\mathcal L_{\beta,\vartheta}\).

\HMTClaim{MD-R-G-PHASE-INTERFERENCE}
\begin{proposition}[Critère d'interférence pour la réponse de phase]
\label{prop:interferencia-selector-G-rev11}
Dans une orbite déterministe avec une unique arête entrante par sommet,
\(\mathcal H_{\beta,\vartheta}\) est indépendant de \(\vartheta\). Toute
réponse spectrale de phase non nulle requiert des histoires cohérentes
multiples ou une différentielle magnétique dont l'holonomie subsiste dans les cycles
enrichis.
\end{proposition}

\begin{proof}
Dans l'orbite déterministe, l'unique phase de chaque terme est multipliée par sa
conjuguée dans
\(\mathcal L_{\beta,\vartheta}^{*}\mathcal L_{\beta,\vartheta}\)
et s'annule. L'absence de termes croisés élimine toute dépendance à
\(\vartheta\).
\end{proof}

Si une bande simple \(\lambda_0(\vartheta)\) possède un minimum à l'origine,
sa hessienne \(Q\) est semi-définie positive. Avec
\[
 v_{27}=(1,18),
 \qquad
 \mathcal R_{27}=v_{27}^{\mathsf T}Qv_{27},
\]
la condition \(\mathcal R_{27}\neq0\) constitue le contrôle minimal de
sensibilité à la mémoire accumulée. Une normalisation naturelle
\(\mathfrak N\) permettrait de définir
\[
 \Lambda_G=\ell_0\mathfrak N(Q,v_{27}),
 \qquad
 G_{\rm spec}
 =\frac{c_{\rm int}^{3}}{\hbar_{\rm int}}
  \bigl(\delta_\theta\Lambda_G\bigr)^2,
 \qquad
 \delta_\theta=\frac{A_{\rm deg}-C_{\rm deg}^{\ast}}{360}.
\]
L'homogénéité de ce gabarit est exacte. Sa réalisation physique est
caractérisée par quatre conditions positives : multigraphe enrichi,
réponse non nulle, stabilité par raffinement et indépendance vis-à-vis du changement d'échelle,
avec une normalisation \(\mathfrak N\) fixée avant la confrontation
métrologique.

\subsection{Fonctionnelle géométrique sur l'holonomie enrichie}

La seconde continuation définit directement une fonctionnelle de longueur
\[
 \mathscr L:
 \operatorname{Loop}(\mathcal X_{\rm TPK}^{\rm enr})
 \longrightarrow\mathcal L_L^+\cup\{0\},
\]
invariante par changement cyclique du point de base, conjugaison, réétiquetage
admissible et subdivision. Pour un relèvement enrichi, on écrit
\[
 \Lambda_{27}:=\mathscr L(\widetilde{\mathcal K}_{27}),
 \qquad
 \rho_{27}:=\frac{\Lambda_{27}}{c_{\rm int}t_0},
 \qquad
 \theta_G^{\rm hol}:=(\delta_\theta\rho_{27})^2.
\]
La transduction produit le gabarit homogène
\begin{equation}
 \boxed{
 G_{\rm hol}
 =\frac{c_{\rm int}^{3}}{\hbar_{\rm int}}
  \bigl(\delta_\theta\Lambda_{27}\bigr)^2.}
 \label{eq:plantilla-holonomia-G-rev11}
\end{equation}
La loi de concaténation de cette fonctionnelle sur le cycle de torsion fait
partie de sa définition et de sa preuve. La sélection physique est complétée
en construisant \(\widetilde{\mathcal K}_{27}\) et \(\mathscr L\), en fixant une
normalisation unique et en vérifiant la stabilité de la sortie. La valeur
\(G_{\rm CODATA\,2022}\) intervient exclusivement ensuite, dans la confrontation
métrologique.

Ces conditions décrivent ce modèle d'holonomie. La réalisation longitudinale et radiale déjà spécifiée dans \cref{g26:sec:rigidez-radial-es} utilise le registre marqué, son espérance conditionnelle et la norme positive ; sa normalisation et la valeur commune de \(G\) y sont démontrées. On n'identifie pas \(\delta_\theta\Lambda_{27}\) à \(L\) par ressemblance formelle : le modèle alternatif conserve l'obligation de construire sa fonctionnelle avant d'effectuer cette identification.

\begin{proposition}[Séparation des deux sélecteurs gravitationnels]
\label{prop:separacion-selectores-G}
Le sélecteur circulaire
\(\theta_{108}^{\rm circ}=(54/\pi)^2\), construit par la fermeture
action--Compton--gravité, et le sélecteur holonomique
\(\theta_G^{\rm hol}=(\delta_\theta\rho_{27})^2\) sont deux réalisations
typées distinctes du transducteur \(G_{\rm int}\). Ils ne s'identifient pas l'un à
l'autre et la carte décimale \(\mathscr R_G\) ne constitue pas une évaluation de
\(G_{\rm hol}\).
\end{proposition}
\begin{proof}
Le premier sélecteur dépend uniquement de la fermeture circulaire d'ordre 108. Le
second contient, en outre, la longueur spectrale
\(\Lambda_{27}=\mathscr L(\widetilde{\mathcal K}_{27})\). Une égalité
entre les deux exigerait de calculer \(\Lambda_{27}\) et de prouver la coïncidence des
normalisations ; aucune opération du lecteur décimal n'effectue ce calcul.
Par conséquent, leurs domaines et données d'entrée les distinguent avant toute
confrontation numérique.
\end{proof}

\subsection{Critère de sélection structurelle}

\HMTClaim{MD-R-G-SELECTOR-CRITERION}
\begin{criterion}[Sélecteur structurel de la constante gravitationnelle]
\label{crit:selector-G-rev11}
Un sélecteur HMT de \(G\) est identifié lorsqu'il satisfait conjointement :
\begin{enumerate}
 \item domaine enrichi, composition et orientation déclarés ;
 \item descente à une classe abélienne non de torsion, ou fonctionnelle géométrique
 invariante par point de base, conjugaison, réétiquetage et subdivision ;
 \item réponse positive non triviale et stable sous raffinement ;
 \item normalisation unique et indépendante de l'échelle globale de l'opérateur ;
 \item sélection interne de la coordonnée et de sa position décimale ;
 \item gel du code et des entrées avant d'ouvrir la valeur de
 confrontation ;
 \item réalisation finale comme section de \(\mathcal L_G\) dans une carte
 dimensionnelle commune.
\end{enumerate}
\end{criterion}

Les rapports adimensionnels
\[
 \frac{\mathcal G_{\rm tor}}{\mathcal I_{\rm tor}}=1,
 \qquad
 \frac{\mathcal G_{\rm av}}{\mathcal I_{\rm av}}
 =\frac{\pi}{\varphi^2}
\]
décrivent des lecteurs gravitationnels--inertiels distincts et demeurent
compatibles avec une seule section dimensionnelle \(G\). La construction de cette
section utilise en outre
\((\hbar_{\rm int},c_{\rm int},\ell_0)\) et le sélecteur \(\theta_G\).

\subsection{Évaluation de la carte décimale gravitationnelle et confrontation métrologique ultérieure}

La représentation décimale conservée dans le corpus est
\begin{equation}
 \mathscr R_G([n,k,d]_3;t,e)
 :=10^{-n}\left(k+\frac{t}{10^3}+\frac{e}{10^d}\right).
 \label{eq:lector-decimal-G-rev11}
\end{equation}
Pour la carte
\([n,k,d]_3=[11,6,5]_3\), \(t=674\) et \(e=30\), l'évaluation rationnelle est
\begin{equation}
 \boxed{
 \mathscr R_G
 =10^{-11}\left(6+\frac{674}{1000}+\frac{30}{10^5}\right)
 =6.67430\times10^{-11}.}
 \label{eq:evaluacion-decimal-G-rev11}
\end{equation}
Le triplet, les deux blocs et la position décimale constituent les données déclarées de ce système de coordonnées. Son évaluation est exacte ; sa sélection structurelle est régie par \cref{crit:selector-G-rev11}. Cette décomposition est conservée comme témoin décimal historique. La branche radiale de \eqref{g26:eq:g-rigido-es} obtient sa coordonnée au moyen de la longueur et de l'action et donne \eqref{cr28:eq:g-evaluation} ; elle n'utilise pas les blocs \(674\) et \(30\) pour définir le couplage.

Après choix de la base SI de \(\mathcal L_G\), la confrontation en vigueur est
\begin{equation}
 \boxed{
 G_{\rm CODATA\,2022}=6.67430(15)\times10^{-11}
 \ \mathrm{m^3\,kg^{-1}\,s^{-2}}.}
 \label{eq:G-CODATA-2022-rev11}
\end{equation}

\begin{center}
\begin{tabularx}{0.94\textwidth}{@{}L{0.27\textwidth}Y@{}}
\toprule
Champ & Déclaration métrologique\\
\midrule
Valeur centrale &
\(6.67430\times10^{-11}\ \mathrm{m^3\,kg^{-1}\,s^{-2}}\)\\
Incertitude type &
\(u(G)=0.00015\times10^{-11}
=1.5\times10^{-15}\ \mathrm{m^3\,kg^{-1}\,s^{-2}}\)\\
Incertidumbre relativa & \(u_r(G)\simeq2.2\times10^{-5}\)\\
Source et temporalité &
Valeurs recommandées CODATA 2022, publiées en 2024
\cite{CODATA2022}\\
Fonction probatoire &
Confrontation externe ultérieure ; l'unité et l'incertitude n'interviennent pas dans
l'évaluation rationnelle de \eqref{eq:evaluacion-decimal-G-rev11}.\\
\bottomrule
\end{tabularx}
\end{center}

\paragraph{Reproductibilité du lecteur gravitationnel.}
Une énumération indépendante vérifie conjointement le type dimensionnel, le
lecteur décimal et les identités de monodromie sur le domaine fini
déclaré.

\begin{proposition}[Chaîne causale de la constante gravitationnelle]
L'action covariante, le hamiltonien contraint et l'équation
Einstein--Schrödinger gouvernent la dynamique gravitationnelle. Le type dimensionnel
et le transducteur \eqref{eq:transductor-gravitatorio-rev11} sont exacts ;
l'évaluation \eqref{eq:evaluacion-decimal-G-rev11} est exacte pour sa carte
déclarée. La valeur, l'unité et l'incertitude de
\eqref{eq:G-CODATA-2022-rev11} constituent une confrontation externe CODATA 2022. La
sélection autonome de \(\theta_G\) est typée par le relèvement
enrichi, la réponse interférente et la fonctionnelle d'holonomie soumis
au \cref{crit:selector-G-rev11}.
\end{proposition}
Dans la réalisation radiale des règles R1--R8, la sélection est déterminée par \eqref{g26:eq:longitud-unica-es} et \eqref{g26:eq:g-rigido-es} ; la compatibilité avec le calendrier est démontrée dans \eqref{cr28:eq:circular-radial-agreement}. Le modèle d'holonomie conserve séparément les conditions de son propre domaine.
 % Integración por delta de la adenda de realizaciones físicas.
% Residencia de procedencia: tipo dimensional de G, selector de holonomía y carta decimal.

\chapter{Transduction dimensionnelle et sélection holonomique de la constante
gravitationnelle}
\label{chap:transductor-gravitatorio-G}
\index{constante gravitationnelle!transducteur dimensionnel}
\index{naturalité!sélecteur gravitationnel}
\index{lecteur décimal!contrôle de G}

\section{Droite gravitationnelle et forme dimensionnelle}
\label{sec:recta-gravitatoria-interna}

Dans le réseau du \cref{chap:funtor-dimensional-hmt}, la dimension de la
constante gravitationnelle est
\[
 \mathbf G=-\mathbf S+5\mathbf C+2\mathbf T,
 \qquad
 [G]=L^3M^{-1}T^{-2}.
\]
Soient
\[
 c_{\rm int}\in\mathcal L_C^+,\qquad
 t_0\in\mathcal L_T^+,\qquad
 \ell_0:=c_{\rm int}t_0\in\mathcal L_L^+,
 \qquad
 \hbar_{\rm int}\in\mathcal L_S^+.
\]
La section \(\hbar_{\rm int}\) n'est obtenue qu'après le lecteur
déterminantal de décade \(-34\) et le choix d'une base de la droite
d'action. Ce n'est pas le caractère angulaire adimensionnel antérieur au retour.

Pour tout caractère adimensionnel \(\theta_G>0\), nous définissons
\[
 \boxed{
 G_{\rm int}(\theta_G)
 =\theta_G
 \frac{c_{\rm int}^{5}t_0^2}{\hbar_{\rm int}}
 =\theta_G
 \frac{c_{\rm int}^{3}\ell_0^2}{\hbar_{\rm int}}
 \in\mathcal L_G^+.}
\]
L'égalité des deux expressions et leur type sont
\textsc{exacts internes} dans le réseau dimensionnel. Ils ne sélectionnent pas
\(\theta_G\). Si
\[
 L_G^2:=\frac{\hbar_{\rm int}G_{\rm int}}{c_{\rm int}^3},
\]
alors
\[
 \theta_G=\left(\frac{L_G}{\ell_0}\right)^2.
\]
Cette dernière identité localise la donnée sélectrice : le rapport des longueurs doit être produit par une structure d'holonomie ou de réponse spectrale antérieure au système de coordonnées métrologique. La réalisation de \cref{g26:sec:rigidez-radial-es} satisfait cette exigence par l'inscription et la norme radiale, avec \(L_G=L\). La suite examine d'autres modèles de retour et d'holonomie, en conservant les obstructions propres à leurs quotients.

\section{Obstruction du caractère positif sur le retour observable}
\label{sec:selector-holonomia-G}

Le \cref{sec:jerarquia-retornos-27-54-108} définit
\[
 \mathcal K_{27}:=F_{\rm obs}^{27}.
\]
La proposition primaire \cref{prop:retorno-fobs-identidad} démontre
\[
 F_{\rm obs}^{108}=I;
\]
par conséquent, dans le quotient observable,
\[
 \mathcal K_{27}^{4}=I.
\]

\begin{theorem}[Obstruction de torsion pour le sélecteur positif observable]
\label{thm:no-go-chi-G-observable}
Tout homomorphisme multiplicatif
\[
 \chi:\langle\mathcal K_{27}\rangle
 \longrightarrow\mathbb R_{>0}^{\times}
\]
est trivial. En particulier,
\[
 \chi(\mathcal K_{27})=1.
\]
\end{theorem}

\begin{proof}
La multiplicativité et \(\mathcal K_{27}^{4}=I\) impliquent
\[
 \chi(\mathcal K_{27})^4=\chi(I)=1.
\]
Le seul élément réel positif dont la puissance quatrième est un est \(1\).
\end{proof}

L'ordre observable est exactement quatre, et non simplement un diviseur de quatre :
dans la coordonnée \(\theta\in\mathbb Z/108\mathbb Z\),
\[
 \mathcal K_{27}(\theta)=\theta+27,
 \qquad
 \mathcal K_{27}^{2}(\theta)=\theta+54\neq\theta,
 \qquad
 \mathcal K_{27}^{4}(\theta)=\theta.
\]
Le même groupe cyclique admet des caractères unitaires
\[
 \chi_k(\mathcal K_{27})
 =\exp\!\left(\frac{2\pi i k}{4}\right),
 \qquad k=0,1,2,3,
\]
et, en particulier, les valeurs de phase \(\pm i\). La conclusion typée est
que le cycle observable peut transporter une phase, mais non une échelle positive
non triviale. Ces caractères de phase ne s'identifient pas à l'opérateur
\(J_W\) de Paley : ils satisfont une relation quadratique analogue sur des domaines
distincts.

L'obstruction ne dépend pas d'une comparaison métrologique. Elle affecte la forme
même du sélecteur : un caractère positif sur le cycle observable ne peut
produire un \(\theta_G\) non trivial. En outre, un simple ensemble de classes de
conjugaison ne constitue pas, en général, un groupe et ne peut être déclaré domaine
d'un caractère sans que sa composition soit définie au préalable.

\subsection{Unique continuation multiplicative admissible}

Une voie multiplicative ne peut être rouverte que sur un objet enrichi qui
n'hérite pas de la torsion observable. Pour un objet \(x\) du groupoïde de mémoire,
le domaine correctement typé serait le groupe d'isotropie
\[
 H_x:=\operatorname{Aut}_{\mathcal G_{\rm mem}}(x)
\]
et, par multiplicativité vers un groupe abélien, son abélianisation
\[
 H_x^{\rm ab}\longrightarrow\mathbb R_{>0}^{\times}.
\]
Tout caractère positif annule la partie de torsion de \(H_x^{\rm ab}\). Ainsi,
la proposition doit construire un relèvement
\[
 \widetilde{\mathcal K}_{27}\in H_x,
 \qquad
 p(\widetilde{\mathcal K}_{27})=\mathcal K_{27},
\]
prouver que sa classe abélienne est non de torsion et sélectionner naturellement
un caractère \(\widetilde\chi_G\). La voie positive sur le cycle observable
est ainsi fermée par obstruction : sans cette donnée non de torsion, il n'existe pas de
sélecteur multiplicatif non trivial. Si tout relèvement admissible
satisfaisait
\(\widetilde{\mathcal K}_{27}^{\,4}=I\), la voie multiplicative serait
également écartée dans le domaine enrichi.

\section{Voie chronologique de longueur, d'action et de spectre}
\label{sec:via-cronologica-espectral-G}

Le \cref{thm:cobertura-acumulativa-dirigida} fournit une troisième voie,
distincte tant du caractère observable obstrué que de l'éventuel
relèvement dans un groupe d'isotropie. Sur
\[
 \widetilde{\mathcal O}_{108}^{+}
 =\mathcal O_{108}\times\mathbb N^2,
\]
le retour de vingt-sept pas ajoute \((1,18)\) et sa quatrième puissance ajoute \((4,72)\). Cette action appartient à une catégorie chronologique dirigée. Ce n'est ni un relèvement du groupoïde ni un champ démontré natif de l'état holonomique intégral complet.

La fonctionnelle positive déjà construite,
\[
 \mathcal A_\ast(\gamma)
 =G_{\rm or}(\gamma)+\lambda_\ast S(\gamma),
\]
est additive par concaténation d'histoires dirigées. Ce n'est pas un caractère du
groupe d'isotropie, car ses indicateurs ne changent pas de signe sous une
inversion formelle. La complétion
\(\mathbb N^2\to\mathbb Z^2\) permet d'introduire des caractères unitaires pour
l'analyse harmonique, mais n'altère pas ce fait.

\begin{equation}
\label{eq:fraccion-angular-orientada-G}
 \delta_\theta
 :=\frac{A_{\rm deg}-C_{\rm deg}^{\ast}}{360}.
\end{equation}
Cette fraction angulaire orientée est définie une fois construits l'angle et le
contre-angle uniques. Elle sera utilisée dans les deux gabarits dimensionnels
ultérieurs et n'intervient pas rétroactivement dans l'obtention de
\(A_{\rm deg}\) ni de \(C_{\rm deg}^{\ast}\).

\subsection{Opérateur tordu et contrôle négatif déterministe}

Pour le contrôle exact, soit \(\Gamma=(V,E)\) un multigraphe orienté
\emph{fini} enrichi, muni d'une mesure
\(\mu:V\to\mathbb R_{>0}\), de poids \(a_\ast(e)\geq0\), d'un paramètre
\(\beta\geq0\) et d'un cocycle
\(c(e)\in\mathbb Z^2\). Pour une normalisation fixée, considérons
\[
 (\mathcal L_{\beta,\vartheta}f)(y)
 =
 \sum_{e:x\to y}
 e^{-\beta a_\ast(e)}
 e^{i\langle\vartheta,c(e)\rangle}f(x),
 \qquad
 \vartheta\in\mathbb T^2,
\]
et prenons pour domaine tout l'espace de Hilbert fini
\(\ell^2(V,\mu)\). Alors \(\mathcal L_{\beta,\vartheta}\) est borné et
\[
 \mathcal H_{\beta,\vartheta}
 :=\mathcal L_{\beta,\vartheta}^{*}
   \mathcal L_{\beta,\vartheta}\geq0
\]
est autoadjoint. Une extension à un graphe infini exigerait, en outre,
des opérateurs bornés ou bien des opérateurs fermés densément définis tels
que \(\mathcal L_{\beta,\vartheta}^{*}
\mathcal L_{\beta,\vartheta}\) soit autoadjoint sur un domaine commun ;
cette extension n'est pas présupposée ici.

\begin{proposition}[Annulation de phase dans l'orbite déterministe]
\label{prop:control-negativo-espectral-determinista}
Si \(\Gamma\) est uniquement l'orbite déterministe observable et que chaque sommet
possède une seule arête entrante, alors
\(\mathcal H_{\beta,\vartheta}\) est indépendant de \(\vartheta\). En
particulier, la hessienne de chacune de ses bandes par rapport à
\(\vartheta\) est nulle.
\end{proposition}

\begin{proof}
Dans chaque composante, \(\mathcal L_{\beta,\vartheta}\) multiplie l'unique
terme entrant par une phase unitaire. Dans
\(\mathcal L_{\beta,\vartheta}^{*}\mathcal L_{\beta,\vartheta}\), cette phase
est multipliée par sa conjuguée et s'annule. Il ne reste pas de termes croisés
entre histoires distinctes.
\end{proof}

Ce contrôle empêche d'attribuer une rigidité spectrale au simple cycle déterministe.
Une réponse de phase non triviale exige au moins l'une des constructions
suivantes :
\begin{enumerate}
 \item un multigraphe enrichi dans lequel deux histoires cohérentes ou plus
 contribuent à une même amplitude ; ou
 \item une différentielle magnétique
 \[
  (d_\vartheta f)(e)
  =
  e^{i\langle\vartheta,c(e)\rangle/2}f(t(e))
  -
  e^{-i\langle\vartheta,c(e)\rangle/2}f(s(e)),
 \]
 avec un laplacien
 \(\Delta_\vartheta=d_\vartheta^*d_\vartheta\), dont la phase dépende de
 l'holonomie de cycles enrichis.
\end{enumerate}
Dans les deux cas, l'espace de Hilbert, le domaine, la mesure,
les poids et la normalisation doivent être déclarés, et il faut démontrer que la famille n'est pas
unitairement équivalente à la famille sans phase.

Supposons que
\(\vartheta\mapsto\mathcal H_{\beta,\vartheta}\) soit une famille
autoadjointe analytique sur un domaine commun et que
\(\lambda_0(\vartheta)\) soit une bande simple et isolée dans un voisinage de
\(\vartheta=0\). Si \(0\) est un point critique et un minimum local,
\[
 Q_{ij}
 =
 \left.
 \frac{\partial^2\lambda_0(\vartheta)}
      {\partial\vartheta_i\partial\vartheta_j}
 \right|_{\vartheta=0}
\]
est semi-définie positive. Sans l'hypothèse de minimum, \(Q\) représente seulement
la réponse quadratique. Pour
\[
 v_{27}=(1,18),
 \qquad
 \mathcal R_{27}=v_{27}^{\mathsf T}Qv_{27},
\]
la condition \(\mathcal R_{27}\neq0\) est une preuve minimale de sensibilité
spectrale à la mémoire accumulée.

\subsection{Normalisation typée de la réponse spectrale et construction de \texorpdfstring{\(\theta_G^{\mathrm{spec}}\)}{theta G spec}}

Si une normalisation naturelle et unique \(\mathfrak N\) transforme la réponse
précédente en un rapport positif indépendant de l'échelle globale de
l'opérateur, on peut définir
\[
 \Lambda_G
 :=\ell_0\,\mathfrak N(Q,v_{27}),
 \qquad
 \theta_G^{\rm spec}
 :=
 \left(
 \delta_\theta\frac{\Lambda_G}{\ell_0}
 \right)^2,
\]
et, par le transducteur dimensionnel,
\[
 \boxed{
 G_{\rm spec}
 =
 \frac{c_{\rm int}^{3}}{\hbar_{\rm int}}
 \bigl(\delta_\theta\Lambda_G\bigr)^2.}
\]
L'homogénéité du gabarit est exacte. La classe de ses évaluations est
caractérisée par les conditions de sélection canonique de l'opérateur,
de réponse non nulle, de naturalité et de stabilité sous raffinement, d'indépendance
vis-à-vis du changement d'échelle, d'unicité de \(\mathfrak N\) et de scellement préalable à toute consultation
de la valeur métrologique de \(G\). Sans ces conditions, l'invariance par
changement d'échelle démontre la non-unicité du sélecteur ; avec elles, l'évaluation est
déterminée dans le domaine normalisé déclaré.

\section{Fonctionnelle non homomorphe sur les voies enrichies}
\label{sec:funcional-longitud-K27-G}

La seconde continuation possible abandonne le caractère multiplicatif et utilise
une fonctionnelle géométrique, spectrale ou d'action. Soit
\[
 \mathscr L:
 \operatorname{Loop}(\mathcal X_{\rm TPK}^{\rm enr})
 \longrightarrow\mathcal L_L^+\cup\{0\}
\]
invariante par changement cyclique du point de base, réétiquetage admissible,
conjugaison et subdivision. Sa loi de concaténation doit être déclarée
explicitement ---par exemple, comme norme sous-additive ou fonctionnelle spectrale--- ;
on ne présume pas qu'elle soit un homomorphisme du groupe de torsion. Pour un
relèvement enrichi candidat, nous écrivons
\[
 \Lambda_{27}:=\mathscr L(\widetilde{\mathcal K}_{27}),
 \qquad
 \rho_{27}:=\frac{\Lambda_{27}}{c_{\rm int}t_0}.
\]
La fraction angulaire orientée de
\eqref{eq:fraccion-angular-orientada-G} produit le candidat adimensionnel
\[
 \theta_G^{\rm hol}:=(\delta_\theta\rho_{27})^2.
\]
Sa spécialisation dans le transducteur donne
\[
 \boxed{
 G_{\rm hol}
 =\frac{c_{\rm int}^{3}}{\hbar_{\rm int}}
  \bigl(\delta_\theta\Lambda_{27}\bigr)^2.}
\]
La forme temporelle est la même identité :
\[
 T_{27}:=\frac{\Lambda_{27}}{c_{\rm int}},
 \qquad
 t_\ast:=\delta_\theta T_{27}.
\]
En particulier, il ne convient pas de remplacer \(t_\ast\) par
\(T_{27}/\delta_\theta\).

L'homogénéité de ces formules est exacte. Leur statut n'en devient pas pour autant celui d'une dérivation de \(G\) par ce modèle : sa sélection reste ouverte jusqu'à la construction de \(\widetilde{\mathcal K}_{27}\) et de \(\mathscr L\) sans consultation de la valeur extérieure, à la démonstration de leur naturalité et à la fixation d'une normalisation unique. L'usage en aval de \(A\) et de \(C^\ast\) dans \(\delta_\theta\) n'autorise pas à utiliser \(G\) pour les redéfinir. La preuve radiale de \cref{g26:sec:rigidez-radial-es} constitue une autre réalisation spécifiée : ses données, sa norme, sa longueur et son couplage sont construits explicitement, et sa réduction conserve la mémoire. La limitation du modèle précédent ne se transfère pas à cette démonstration.

\section{Transduction de la carte décimale gravitationnelle sur la droite de grandeur et séparation causale de la reconnaissance métrologique}
\label{sec:lector-decimal-independiente-G}

Soit
\[
 \mathscr R_G([n,k,d]_3;t,e)
 :=10^{-n}\left(k+\frac{t}{10^3}+\frac{e}{10^d}\right),
\]
où l'indice de \([n,k,d]_3\) enregistre le triplet du lecteur et ne désigne pas
un numéral ternaire. Pour les données déclarées
\[
 [n,k,d]_3=[11,6,5]_3,\qquad t=674,\qquad e=30,
\]
l'évaluation rationnelle est
\[
 \boxed{
 \mathscr R_G
 =10^{-11}\left(6+\frac{674}{1000}+\frac{30}{10^5}\right)
 =6.67430\times10^{-11}.}
\]
L'égalité arithmétique est exacte une fois les paramètres déclarés. Le triplet,
les blocs \(674\), \(30\) et leur position décimale ont le statut de données
de la carte, de sorte que cette branche est conservée comme contrôle paramétré et
reconnaissance ultérieure. Le nombre est adimensionnel jusqu'au choix d'une base
\(u_G\in\mathcal L_G^+\) ; alors seulement la section est définie
\[
 G_{\rm dec}:=\mathscr R_G\,u_G.
\]
La coïncidence de sa coordonnée avec une publication métrologique constitue une \textsc{reconnaissance extérieure}. Une promotion ultérieure exigerait que les entiers \(674\), \(30\), le triplet \([11,6,5]\) et la position décimale soient imposés de manière unique par des symétries établies et fixés avant de consulter la valeur extérieure.

\section{Hiérarchie probatoire des constructions gravitationnelles}
\label{sec:cuadrado-contraste-G}

La réalisation longitudinale et radiale détermine le couplage dans le domaine explicite R1--R8. Les modèles alternatifs et les contrôles précédents conservent leurs portées propres ; ils ne constituent pas à eux seuls d'autres généalogies démontrées de cette réalisation. L'organisation est :
\begin{center}
\begin{tabularx}{0.96\textwidth}{@{}L{0.25\textwidth}Y L{0.26\textwidth}@{}}
\toprule
Objet & Résultat disponible & Verdict\\
\midrule Inscription et norme radiale R1--R8 &
\(L=(5759/23040)\alpha^{16}L_*\),
\(R_X\Lambda_X=L^2I\), \(H_X\Lambda_X=\hbar_{\rm ret}cI\),
\(G=c^3L^2/\hbar_{\rm ret}\). &
Détermination unique dans le domaine déclaré ; mémoire, réduction et raffinement conservés.\\
Réalisation circulaire CT108 &
\(t_0=\pi L/(54c)\), \(R_{108}=L\) y
\(\theta_{108}^{\rm circ}=(54/\pi)^2\). &
Compatibilité exacte avec la longueur construite.\\
Transductor dimensional &
Type exactement \(G_{\rm int}(\theta_G)\), sans fixer \(\theta_G\). &
Théorème dimensionnel ; non sélecteur.\\
Caractère de l'observable \(\mathcal K_{27}\) &
\(\mathcal K_{27}^4=I\) impose
\(\chi(\mathcal K_{27})=1\). &
Formulation obstruée.\\
Levantamiento enriquecido &
Possible seulement avec une classe abélienne non de torsion, un caractère naturel et une
normalisation unique. &
Obstruction de canonicité démontrée sans ces données.\\
Revêtement chronologique et réponse spectrale &
La translation cumulative est d'ordre de semi-groupe infini ; dans l'orbite
déterministe, la phase s'annule. &
Réponse nulle démontrée dans l'orbite déterministe ; classe non triviale
caractérisée par l'interférence et la normalisation naturelle.\\
Fonctionnelle de voie &
La formule de \(G_{\rm hol}\) est homogène une fois
\(\Lambda_{27}\) fixé. &
Famille paramétrée exacte ; il n'existe pas de sélecteur univoque sans spécifier
\(\mathscr L\) et sa normalisation.\\
Lecteur décimal &
Reproduit exactement la coordonnée déclarée, mais reçoit ses blocs et sa
position. &
Control parametrizado.\\
\bottomrule
\end{tabularx}
\end{center}

\begin{criterion}[Sélecteur de la constante gravitationnelle] \label{crit:selector-G} Pour les modèles alternatifs de relèvement, de réponse spectrale et de lecteur décimal examinés dans ce chapitre, une promotion est rejetée si :
\begin{enumerate}
 \item le relèvement enrichi conserve uniquement de la torsion ou ne descend pas
 à une classe abélienne bien définie ;
 \item l'opérateur chronologique ne conserve pas l'interférence, produit
 \(\mathcal R_{27}=0\), est dépourvu de domaine ou de mesure, ou sa réponse change sous
 un raffinement admissible ;
 \item \(\mathscr L\) dépend du point de base, du réétiquetage, de
 l'orientation choisie, d'une subdivision ou d'une normalisation libre ;
 \item la normalisation \(\mathfrak N\) dépend de l'échelle globale de
 l'opérateur, n'est pas unique ou est fixée en consultant \(G\) ;
 \item dans le lecteur décimal, \(674\), \(30\) ou la position décimale admettent des alternatives de même légitimité structurelle ;
 \item le candidat consulte la valeur métrologique avant d'être scellé ;
 \item la valeur n'est pas stable sous raffinement ou ne définit pas une section de
 \(\mathcal L_G\) dans une carte dimensionnelle commune.
\end{enumerate}
\end{criterion}

\section{Rapports gravitationnel--inertiel et constante dimensionnelle}
\label{sec:razones-no-rivales-G}

La normalisation locale de la carte torsionnelle,
\[
 \frac{\mathcal G_{\rm tor}}{\mathcal I_{\rm tor}}=1,
\]
et l'interface surfacique--volumétrique du
\cref{sec:ambivalencia-grav-inercial-pi-phi2},
\[
 \frac{\mathcal G_{\rm av}}{\mathcal I_{\rm av}}
 =\frac{\pi}{\varphi^{2}},
\]
sont des rapports adimensionnels de lecteurs. Ce ne sont pas deux valeurs rivales de la
constante dimensionnelle \(G\). La première exprime une restriction diagonale
locale ; la seconde est une interface conditionnée par la factorisation à travers une
mémoire commune. La construction d'une section de
\(\mathcal L_G\) requiert, en outre, le repère
\((\hbar_{\rm int},c_{\rm int},\ell_0)\) et un sélecteur
\(\theta_G\).
 
\chapter{Dynamique de Perron, mémoire torsionnelle et réalisation Einstein--Cartan}
\begingroup
\renewcommand{\chapter}[2][]{}%
\chapter{Dynamique de Perron, dilution \texorpdfstring{\(a^{-6}\)}{a⁻⁶} de la mémoire et réalisation d’Einstein–Cartan}
\label{chap:parte-iv-67}

L’opérateur de Perron fixe le mode stable de mémoire ; exprimée comme densité de spin, sa contribution se dilue selon la sixième puissance inverse du facteur d’échelle. La réalisation d’Einstein--Cartan conserve cette provenance et ses conditions de validité.

\section{Dynamique de Perron et mode stable de mémoire}

Sur un parcours fermé, le cocycle cylindrique de Perron satisfait

\begin{equation}
V_n=\varphi^n.
\label{rm:eq:cosmo-perron-V}
\end{equation}

En dimension spatiale \(D=3\), l’échelle déterminée par le cocycle est

\begin{equation}
\frac{a_n}{a_0}=V_n^{1/3}=\varphi^{n/3},
\label{rm:eq:cosmo-perron-a}
\end{equation}

tandis que le mode stable de l’action quadratique est

\begin{equation}
M_n=\varphi^{-2n}=V_n^{-2}.
\label{rm:eq:cosmo-perron-M}
\end{equation}

Éliminer \(n\) entre \cref{rm:eq:cosmo-perron-a,rm:eq:cosmo-perron-M} donne le résultat central :

\begin{theorem}[Loi cylindrique de sixième ordre pour la mémoire]
Pour tout niveau fermé de la réalisation de Perron dans \(D=3\),
\begin{equation}
\boxed{
M_n=\left(\frac{a_n}{a_0}\right)^{-6}.}
\label{rm:eq:cosmo-minus-six}
\end{equation}
L’exposant \(-6\) n’est pas ajusté : c’est le quotient de l’exposant stable \(-2\) par la racine dimensionnelle \(1/3\).
\end{theorem}

Aux retours distingués, on obtient

\[
\frac{a_9}{a_0}=\varphi^3,
\qquad
\frac{a_{27}}{a_0}=\varphi^9,
\qquad
\frac{a_{108}}{a_0}=\varphi^{36},
\]

et les mêmes niveaux vérifient exactement \(M_n=(a_n/a_0)^{-6}\). Si l’on déclare le paramètre temporel additif \(t_n=nt_0\), le taux logarithmique est constant :

\begin{equation}
H_\varphi
:=\frac{\Delta\log a_n}{\Delta t}
=\frac{\log\varphi}{3t_0},
\qquad
H_\varphi t_0
=0.1604039416865344824992529711\ldots.
\label{rm:eq:cosmo-Hphi}
\end{equation}

Le symbole \(H_\varphi\) désigne le taux de cette échelle discrète. Sa lecture comme paramètre de Hubble exige le transducteur cosmographique ; l’identité interne ne dépend pas de cette promotion.

\section{Réalisation dans le système d’Einstein--Cartan}

Dans un fluide de spin d’Einstein--Cartan, la densité quadratique de spin obéit à

\begin{equation}
s^2\propto a^{-6}.
\label{rm:eq:cosmo-EC-scaling}
\end{equation}

\cref{rm:eq:cosmo-minus-six} fournit l’exposant et le mode stable sans consulter un rebond. Pour effectuer la lecture physique, il faut déclarer une section positive \(s_*^2\) et une application

\begin{equation}
\mathcal T_{\rm EC}:M_n\longmapsto
s_n^2:=s_*^2M_n
=s_*^2\left(\frac{a_n}{a_0}\right)^{-6}.
\label{rm:eq:cosmo-EC-transducer}
\end{equation}

L’application conserve la loi d’échelle, mais ne détermine pas à elle seule l’amplitude \(s_*^2\), le courant de spin, le signe du terme de contorsion ni l’équation d’état. Avec une densité de masse, une équation effective prend la forme

\begin{equation}
H^2=\frac{8\pi G}{3}
\bigl(\rho_{\rm m}-\rho_{\rm spin,m}\bigr)+\cdots,
\label{rm:eq:cosmo-EC-mass}
\end{equation}

tandis qu’avec une densité d’énergie, il faut écrire

\begin{equation}
H^2=\frac{8\pi G}{3c^2}
\bigl(\rho_{\rm E}-\rho_{\rm spin,E}\bigr)+\cdots.
\label{rm:eq:cosmo-EC-energy}
\end{equation}

Les deux conventions ne se mélangent pas. Un candidat au rebond exige, outre l’égalité des densités, régularité, matière et données aux limites.

\begin{falsifier}
La transduction d’Einstein--Cartan échoue si la fonctionnelle d’échelle produit un exposant différent de \(-6\), si \(V_n\ne\varphi^n\) sur le parcours déclaré, ou si l’application physique ne conserve pas la positivité et les dimensions. Choisir \(s_*^2\) rétrospectivement pour imposer un seuil de rebond ne démontre pas le transducteur.
\end{falsifier}

\section{Comparaison structurelle entre le cycle de 108 états et la dynamique de Perron}

\begin{center}
\begin{tabularx}{\textwidth}{>{\bfseries}p{29mm}XX}
\toprule
& CT108--de Sitter & Perron--Einstein--Cartan\\
\midrule
Antécédent & retour circulaire d’ordre \(108\) & cocycle cylindrique et mode stable\\
Sélecteur & \(r_g=R_{108}\) & \(V_n=\varphi^n\), \(D=3\)\\
Magnitudes determinadas & rayon, courbure, température et entropie & expansion, taux et dilution de la mémoire\\
Invariant & \(G_a\hbar_a\) et rayon circulaire & \(M_na_n^6\)\\
Statut & interface de Sitter conditionnée & loi interne exacte ; lecture EC conditionnée\\
Critère de réfutation & rayon ou équation de fond incompatibles & perte de l’exposant \(-6\) ou du transducteur\\
\bottomrule
\end{tabularx}
\end{center}

\begin{statusbox}
Les identités CT108, le sélecteur \(\theta_{108}\) et la loi \(M_n=(a_n/a_0)^{-6}\) sont \status{Théorème interne exact}. La lecture du même rayon comme rayon de Sitter et l’identification de \(M_n\) à une densité de spin sont
\status{Réalisations physiques conditionnées}, chacune dotée de son propre critère de réfutation.
\end{statusbox}
 
\paragraph{Coordination de la torsion et de l’aire dans l’action de Palatini--Cartan--Holst.} La torsion et l'aire sont ensuite coordonnées dans l'action du premier ordre de Palatini--Cartan--Holst. \endgroup

\chapter{Réalisation Palatini--Cartan--Holst du paramètre de Barbero--Immirzi}
\begingroup
\renewcommand{\chapter}[2][]{}%
\chapter{Réalisation de Palatini–Cartan–Holst du paramètre de Barbero–Immirzi et transduction physique du spectre d’aire}
\label{chap:parte-iv-68}

La dérivation structurelle du paramètre et du spectre d’aire précède ce chapitre. On construit ici la tétrade, la connexion, la courbure et l’action du premier ordre, et ce n’est qu’ensuite que le paramètre de Holst est spécialisé comme réalisation physique du lecteur aréal.

\section{Courbure, torsion et action d’Einstein--Cartan--Holst}
\label{chap:gravedad-cartan-holst}

La structure de transport HMT distingue la variable de mémoire, l’holonomie et le défaut de fermeture. Sa réalisation gravitationnelle se formule au moyen d’une tétrade et d’une connexion. Cette section établit la géométrie du premier ordre et spécifie le type de l’application entre celle-ci et l’espace des états TPK\@.

\subsection{Identité discrète de Gauss–Stokes et transport de l’incidence vers la frontière}

Soit \(K\) un complexe cellulaire, soit \(W\) un groupe abélien et soit \(\Omega\in C^1(K;W)\) une cochaîne. Nous notons \(\delta:C^1(K;W)\to C^2(K;W)\) le cobord cellulaire et \(\langle\,\cdot\,,\,\cdot\,\rangle\) l’accouplement entre cochaînes à valeurs dans \(W\) et chaînes cellulaires entières.

\begin{theorem}[Identité cellulaire universelle de Stokes]
Pour toute \(2\)-chaîne finie \(z\in C_2(K;\mathbb Z)\),
\[
 \boxed{\langle\delta\Omega,z\rangle
 =\langle\Omega,\partial z\rangle.}
\]
\end{theorem}
\begin{proof}
C’est l’identité qui définit le cobord comme opérateur adjoint du bord cellulaire. La finitude de \(z\) garantit que les deux accouplements sont des sommes finies dans \(W\).
\end{proof}

\begin{corollary}[Gauss--Stokes sur une pseudovariété cellulaire]
Soit \(K\) une pseudovariété cellulaire régulière, bidimensionnelle, finie et compacte, orientée de manière cohérente, éventuellement à bord. Si \(z_K\) est la somme de ses \(2\)-cellules avec l’orientation cohérente, alors
\[
 \sum_{f\in K^{(2)}}(\delta\Omega)(f)
 =\sum_e[\partial z_K:e]\,\Omega(e).
\]
En particulier, la régularité permet d’identifier les incidences non annulées de \(\partial z_K\) aux arêtes du bord géométrique.
\end{corollary}
\begin{proof}
On applique le théorème à \(z_K\). L’orientation cohérente fait apparaître chaque arête intérieure deux fois dans \(\partial z_K\), avec des coefficients opposés ; les incidences non annulées forment la chaîne de bord.
\end{proof}

La version non abélienne ne s’obtient pas en remplaçant simplement les sommes par des produits. Dans cet ouvrage, elle n’est utilisée que lorsque les transports vers une base commune, un ordre des holonomies et la conjugaison correspondante ont été fixés. En dehors de ce sous-domaine, aucune identité de Stokes non abélienne n’est affirmée.

\subsection{Capacité aréale de la cellule qutrit : \(81\log 3\) et transport collectif d’incidence}

\begin{theorem}[Coupe minimale d’un cube discret]
Dans le graphe cubique \(N\times N\times N\) de capacité unitaire, toute coupe séparant deux faces opposées a une capacité au moins égale à \(N^2\), et il en existe une de capacité \(N^2\).
\end{theorem}
\begin{proof}
Il existe \(N^2\) chemins disjoints en arêtes qui traversent le cube dans la direction normale aux deux faces. Toute coupe doit les intercepter, donc sa capacité est au moins \(N^2\). Une couche transversale contient exactement \(N^2\) arêtes et atteint la borne.
\end{proof}

Pour \(N=9\), la coupe vaut 81 et
\[
 9^3=729=27^2.
\]
La dernière égalité permet de regrouper six trits sous forme d’un volume \(9^3\) ou d’une grille de cardinalité \(27^2\). C’est une bijection d’ensembles ; ce n’est pas un isomorphisme entre \((\ZZ/9\ZZ)^3\) et \((\ZZ/27\ZZ)^2\), dont les exposants diffèrent.

\subsection{Tétrade, connexion et \(2\) défauts}

Soient \(e^I\) une tétrade et \(\omega^{IJ}=-\omega^{JI}\) une connexion de Lorentz. Définissons
\[
 T_{\rm tor}^I=D_\omega e^I
 =\dd e^I+\omega^I{}_J\wedge e^J,
\]
\[
 F^{IJ}=\dd\omega^{IJ}+\omega^I{}_K\wedge\omega^{KJ},
 \qquad
 g=\eta_{IJ}e^I\otimes e^J.
\]
La décomposition
\[
 \omega^{IJ}=\mathring\omega^{IJ}(e)+K^{IJ}
\]
sépare la connexion de Levi--Civita et la contorsion. Alors
\[
 T_{\rm tor}^I=K^I{}_J\wedge e^J.
\]
Courbure et torsion sont des défauts différents : \(F\) mesure la fermeture du transport des vecteurs ; \(T_{\rm tor}\), la fermeture de la tétrade. La mémoire \(T\) du principe d’Ambivalence n’est pas non plus une torsion géométrique : seul un foncteur de réalisation peut les relier.

\subsection{Action du 1er ordre}

Soit \(M\) une variété lisse orientée de dimension quatre. Si \(\partial M\ne\varnothing\), nous supposons des variations à support compact dans l’intérieur ou, de manière équivalente pour le problème variationnel considéré, un terme de bord et des conditions aux limites annulant la contribution frontalière. Lorsque \(\Psi\) contient des champs spinoriels, nous supposons en outre une structure de spin et des champs compatibles avec elle. Soient
\[
 \Sigma^{IJ}=e^I\wedge e^J,
 \qquad
 (\star X)^{IJ}=\frac12\epsilon^{IJ}{}_{KL}X^{KL},
 \qquad
 \kappa=\frac{8\pi G}{c^4},
 \qquad
 \gamma\in\mathbb R\setminus\{0\}.
\]
L’action de Palatini--Cartan--Holst est
\[
\boxed{
 S[e,\omega,\Psi]
 =\frac1{2\kappa c}\int_M
 \Sigma_{IJ}\wedge\left(\star F^{IJ}+\frac1\gamma F^{IJ}\right)
 -\frac\Lambda{\kappa c}\int_M\operatorname{vol}_e
 +S_{\rm m}[e,\omega,\Psi].}
\]
La variation de cette action détermine les corrections torsionnelles \cite{Sciama1964,Kibble1961,Hehl1976,Holst1996}.

Définissons le courant de spin par
\[
 \delta_\omega S_{\rm m}
 =-\frac12\int_M\delta\omega^{IJ}\wedge\sigma_{IJ}.
\]

\begin{theorem}[Équation de Cartan--Holst]
Sous les hypothèses géométriques et de bord précédentes, la variation par rapport à \(\omega\) donne
\[
 D_\omega\!\left[(\star+\gamma^{-1}I)\Sigma^{IJ}\right]
 =\kappa c\,\sigma^{IJ},
\]
avec la normalisation globale fixée par la définition précédente de \(\sigma^{IJ}\).
\end{theorem}
\begin{proof}
On utilise \(\delta F^{IJ}=D_\omega\delta\omega^{IJ}\), on intègre par parties et l’on annule le coefficient de la variation arbitraire \(\delta\omega^{IJ}\). Avec la convention de signe fixée dans la définition de \(\sigma^{IJ}\), le terme de matière passe au membre de droite avec un signe positif. Le facteur \(1/(2\kappa c)\) produit le coefficient \(\kappa c\). Le terme cosmologique ne dépend pas de \(\omega\), et les conditions aux limites éliminent le terme engendré par l’intégration par parties.
\end{proof}

En signature lorentzienne, \(\star^2=-I\) sur les bivecteurs et, pour \(\gamma\in\RR\setminus\{0\}\),
\[
 (\star+\gamma^{-1}I)^{-1}
 =\frac{\gamma^2}{1+\gamma^2}(-\star+\gamma^{-1}I).
\]

\begin{corollary}[Découplage torsionnel]
Pour une cotétrade non dégénérée et \(\gamma\in\mathbb R\setminus\{0\}\),
\[
 \sigma^{IJ}=0\Longrightarrow T_{\rm tor}^I=0
 \Longrightarrow\omega=\mathring\omega(e).
\]
\end{corollary}

Pour expliciter la normalisation, définissons les sources normalisées par l’action au moyen de
\[
 \delta_g S_{\rm m}
 =-\frac12\int_M\operatorname{vol}_e\,
 \mathcal T^{S}_{\mu\nu}\,\delta g^{\mu\nu}.
\]
Lorsque les coordonnées de la tétrade ont la dimension d’une longueur, le tenseur énergie--impulsion physique conventionnel est \(T^{\rm phys}_{\mu\nu}:=c\,\mathcal T^{S}_{\mu\nu}\). Notons de même \(\mathcal U^{S,\rm spin}\) la contribution normalisée par l’action et \(U^{\rm phys,\rm spin}:=c\,\mathcal U^{S,\rm spin}\) sa convention physique. En éliminant algébriquement la connexion, l’équation métrique prend alors les deux formes équivalentes
\[
 G_{\mu\nu}+\Lambda g_{\mu\nu}
 =\kappa c\,
 \bigl(\mathcal T_{\mu\nu}^{S,\rm LC}
       +\mathcal U_{\mu\nu}^{S,\rm spin}\bigr)
 =\kappa\,
 \bigl(T_{\mu\nu}^{\rm phys,LC}
       +U_{\mu\nu}^{\rm phys,spin}\bigr),
\]
où la contribution du spin est quadratique dans le courant correspondant et dépend du couplage de matière. Le coefficient de cette contribution ne peut être fixé sans spécifier l’action fermionique et les conventions de dualité.

\subsection{Identité de Bianchi avec torsion et bilan de flux dans la connexion de Cartan}

Pour toute connexion,
\[
 D_\omega T_{\rm tor}^I=F^I{}_J\wedge e^J,
 \qquad
 D_\omega F^{IJ}=0.
\]
Après élimination de la torsion,
\[
\nabla^\mu(T_{\mu\nu}^{\rm phys,LC}
          +U_{\mu\nu}^{\rm phys,spin})=0.
\]
Ces identités sont des tests de cohérence obligatoires : tout terme gravitationnel HMT doit provenir d’une action compatible ou satisfaire le bilan de Noether correspondant.

\subsection{Application de réalisation et spécialisation du paramètre de Holst}

L’insertion définie dans \cref{mdg:chap:barbero} prend
\[
 \gamma=\Gamma_{6\to8}.
\]
Une fois cette substitution effectuée, toutes les conséquences variationnelles précédentes sont des théorèmes de l’action spécialisée. La sélection physique de cette action par le transport HMT s’exprime au moyen d’une application de réalisation
\[
 \mathfrak R_{\rm grav}:\mathcal X_{\rm TPK}
 \dashrightarrow\{(e,\omega,\Psi)\}
\]
dotée de trois propriétés : entrelacer le transport TPK avec celui de \(\omega\), identifier les classes d’holonomie et fournir les échelles de \(G\) et de \(\Lambda\). Ces propriétés fixent le domaine précis d’une réalisation gravitationnelle du formalisme.

Pour une cotétrade non dégénérée, \(\gamma\) réel non nul et une matière sans courant de spin, le secteur torsionnel se réduit à la formulation du premier ordre de la relativité générale avec la même constante cosmologique, sous les conditions aux limites déjà déclarées. Dans un fluide de spin, la loi d’échelle \(s^2\propto a^{-6}\) peut produire un rebond d’Einstein--Cartan ; obtenir sa densité et ses conditions initiales à partir de l’état TPK est une question distincte de la cohérence variationnelle de l’action.
 
\section{Retour nonadique, groupoïde de mémoire et condition d’entrelacement de Cartan}

La hiérarchie \(27\to54\to108\), le recouvrement cumulatif dirigé et le groupoïde de mémoire fournissent la condition exacte sous laquelle le retour du TPK admet une connexion de Cartan compatible. Ce pont précède la spécialisation de Palatini--Cartan--Holst et conserve intégralement son domaine et ses critères de réfutation.

\begingroup
\let\chapter\section
\let\section\subsection
\let\subsection\subsubsection
\let\subsubsection\paragraph
\let\clearpage\relax
\let\cleardoublepage\relax
\chapter{Retour nonadique, mémoire et holonomie de Cartan}
\label{chap:puente-retorno-holonomia-cartan}
\index{retour!27--54--108}\index{holonomie!TPK--Cartan}
\index{mémoire!distinction par rapport à la torsion}

La réalisation gravitationnelle ne commence pas par identifier une variable discrète
à une grandeur géométrique. Elle commence par distinguer trois niveaux : le retour
du transport observable, la mémoire qui différencie des histoires ayant la même
projection et l'holonomie d'une connexion. Le pont formulé ici
reçoit les retours et les cocycles déjà démontrés dans HMT ; il ne les reconstruit pas
à partir de l'action de Holst, d'une masse ou d'une constante métrologique.

\section{Hiérarchie exacte des retours}
\label{sec:jerarquia-retornos-27-54-108}

Soit \(F_{\rm obs}\) la chronologie observable définie dans le noyau HMT et soit
\[
 \mathcal K_{27}:=F_{\rm obs}^{27}.
\]
Le symbole \(\mathcal K_{27}\) désigne exclusivement la composition de vingt-sept mises à jour de la boucle positionnelle et orientée ; il ne désigne ni une fenêtre générique, ni une trace de \(108\) pas, ni l'état holonomique intégral complet. Notons \(p\) la projection sur les positions des deux curseurs et \(\mathsf J_{\rm or}\) l'inversion simultanée de leurs orientations. La démonstration primaire se trouve dans \cref{prop:retorno-fobs-identidad} ; ses conséquences typées sont
\[
 p(\mathcal K_{27}x)=p(x),\qquad
 o(\mathcal K_{27}x)=\mathsf J_{\rm or}\,o(x),
 \qquad \mathsf J_{\rm or}^{\,2}=I,
\]
\[
 F_{\rm obs}^{54}:
 \text{rétablit les positions et les orientations},
 \qquad
 F_{\rm obs}^{108}=I:
 \text{rétablit aussi la phase observable}.
\]
Par conséquent, retour de position, retour d'orientation et retour de phase sont
des affirmations différentes. Aucune d'elles n'implique à elle seule le retour de
tous les champs d'une extension enrichie.

Le quotient de modes et les cocycles d'orientation sont établis dans
\cref{sec:cociclos-modo-orientacion,thm:descenso-necesario-fav}. Dans l'ordre
surfacique--volumétrique,
\[
 N_{27}=9F_{\rm av},
 \qquad
 F_{\rm av}=
 \begin{pmatrix}0&1\\1&1\end{pmatrix}.
\]
Si
\[
 \Omega_{\rm sw}(\gamma)=\sum_{a\subset\gamma}\omega_{\rm sw}(a),
 \qquad
 S(\gamma)=\sum_{a\subset\gamma}|\omega_{\rm sw}(a)|,
 \qquad
 G_{\rm or}(\gamma)=\sum_{a\subset\gamma}g(a),
\]
alors un bloc de vingt-sept pas contient trois tours nonadiques et une
inversion d'orientation :
\[
 \Omega_{\rm sw}(\mathcal K_{27})=0,\qquad
 S(\mathcal K_{27})=18,\qquad
 G_{\rm or}(\mathcal K_{27})=1.
\]
Dans le retour observable de \(108\) pas,
\[
 \Omega_{\rm sw}(C_{108})=0,\qquad
 S(C_{108})=72,\qquad
 G_{\rm or}(C_{108})=4.
\]
Ainsi, pour la fonctionnelle additive déclarée
\[
 \mathcal A_\ast(\gamma)
 :=G_{\rm or}(\gamma)+\lambda_\ast S(\gamma),
\]
on obtient, sans introduire d'échelle dimensionnelle,
\[
 \mathcal A_\ast(\mathcal K_{27})=1+18\lambda_\ast,
 \qquad
 \mathcal A_\ast(C_{108})=4+72\lambda_\ast.
\]
Ces valeurs sont des comptages du registre. Leur utilisation thermique ou
gravitationnelle ultérieure n'altère pas leur provenance.

\subsection{Système dirigé de revêtements, raffinement et limite de la connexion de Cartan}
\label{sec:cobertura-acumulativa-dirigida}

Soit \(\mathcal O_{108}\) l'orbite observable et, pour chaque arête chronologique
\(e\), soit
\[
 c(e)=\bigl(g(e),s(e)\bigr)\in\mathbb N^2
\]
le couple formé par les indicateurs d'inversion d'orientation et de changement
effectif de feuillet. Pour une histoire dirigée
\(\gamma=e_1\cdots e_n\), on définit
\[
 c(\gamma)=\sum_{j=1}^{n}c(e_j).
\]
La loi de concaténation prend la forme
\[
 c_{m+n}(x)
 =c_m(x)+c_n\!\left(F_{\rm obs}^{m}x\right).
\]
Les comptages précédents expriment, uniformément sur l'orbite,
\[
 c_{27}(x)=(1,18),
 \qquad
 c_{108}(x)=(4,72).
\]

\begin{theorem}[Revêtement cumulatif de la chronologie observable]
\label{thm:cobertura-acumulativa-dirigida}
Sur
\[
 \widetilde{\mathcal O}_{108}^{+}
 :=\mathcal O_{108}\times\mathbb N^2
\]
définissons
\[
 \widetilde F^{+}(x,m)
 :=\bigl(F_{\rm obs}x,m+c(x\to F_{\rm obs}x)\bigr).
\]
Alors le transport de vingt-sept pas satisfait
\[
 \widetilde{\mathcal K}_{27}^{+}(x,m)
 =\bigl(\mathcal K_{27}x,m+(1,18)\bigr),
\]
et
\[
 \left(\widetilde{\mathcal K}_{27}^{+}\right)^4(x,m)
 =\bigl(x,m+(4,72)\bigr)\neq(x,m).
\]
Par conséquent, sa projection observable est d'ordre quatre, tandis que
l'action cumulative du monoïde chronologique possède un ordre de semi-groupe infini.
\end{theorem}

\begin{proof}
La loi cocyclique et la constance de \(c_{27}\) donnent
\[
 c_{108}(x)
 =\sum_{j=0}^{3}c_{27}(\mathcal K_{27}^{j}x)
 =4(1,18)=(4,72).
\]
La première coordonnée revient par
\(\mathcal K_{27}^{4}=I\) ; la seconde ne revient pas. Après \(q>0\) itérations,
l'incrément est \(q(4,72)\), non nul.
\end{proof}

\begin{remark}[Type de l'extension cumulative]
La construction précédente étend les histoires dirigées au moyen de compteurs
positifs. Ce n'est pas un relèvement automorphe du groupoïde
\(\mathcal G_{\rm mem}\), et elle ne démontre pas que la coordonnée illimitée appartienne
déjà à l'état \(\mathcal X_{\mathrm{TPK}}^{\mathrm{enr}}\). Sa complétion de Grothendieck
\(\mathbb Z^2\) peut être utilisée pour l'analyse harmonique, mais elle ne transforme pas
le coût positif
\(\mathcal A_\ast=G_{\rm or}+\lambda_\ast S\) en un caractère du groupe
d'isotropie : lorsqu'on inverse formellement une voie, ses indicateurs ne prennent pas
le signe opposé.
\end{remark}

\paragraph{Domaine de validité de l'extension cumulative.}
La hiérarchie \(27\to54\to108\), les incidences de \(F_{\rm av}\) et les cocycles précédents sont des identités exactes dans la chronologie et les lecteurs déclarés. L'affirmation se restreint à l'état observable conservé par ces lecteurs ; l'étendre à l'état holonomique intégral complet exigerait de prouver le retour de chaque champ oublié.

\section{Groupoïde de mémoire}
\label{sec:grupoide-memoria-tpk}

Soit \(\mathsf P_{\rm TPK}^{\rm enr}\) le groupoïde des chemins enrichis :
ses objets sont des états TPK survivants et ses flèches sont des voies admissibles
avec origine et destination déclarées, composition par concaténation et inversion
par renversement de voie. Soit \(\sim_{\rm mem}\) une équivalence qui
n'identifie deux flèches que si elle conserve les champs de mémoire requis par
le lecteur considéré et qui est congruente avec les identités, l'inversion et la
concaténation. Nous définissons
\[
 \mathcal G_{\rm mem}
 :=\mathsf P_{\rm TPK}^{\rm enr}/\!\sim_{\rm mem}.
\]
Cette définition ne transforme pas la mémoire en torsion. La mémoire appartient à
l'état discret et distingue des histoires ; la torsion appartient à une connexion
géométrique et mesure \(D_\omega e\). Seule une application entre les deux domaines peut
les relier.

Une réalisation géométrique exige, comme données supplémentaires, une longueur
\(\ell_0>0\), une représentation du groupoïde
\[
 \rho_{\rm C}:
 \mathcal G_{\rm mem}
 \longrightarrow
 \operatorname{Spin}(1,3)\ltimes\RR^{1,3},
\]
et une application de champs
\[
 \mathfrak R_{\rm grav}:
 \mathsf P_{\rm TPK}^{\rm enr}
 \dashrightarrow
 \mathsf{Cartan}_{1,3}.
\]
Pour définir la seconde application, il faut spécifier son action sur les objets et les
voies, sa compatibilité avec la concaténation et l'inversion et la régularité de la
connexion résultante. Aucune de ces propriétés ne se déduit du simple comptage
des retours.

\section{Connexion conjointe et condition d'entrelacement}
\label{sec:conexion-cartan-conjunta}

Sous les hypothèses géométriques du \cref{chap:gravedad} ---variété
orientée de dimension quatre, cotétrade \(e^I\), connexion de Lorentz
\(\omega^{IJ}\) et conditions au bord qui y sont déclarées---, la connexion de
Cartan conjointe s'écrit
\[
 \mathcal A_{\ell_0}
 =
 \begin{pmatrix}
  \omega&\ell_0^{-1}e\\
  0&0
 \end{pmatrix}.
\]
Sa courbure est l'identité algébrique
\[
 \mathcal F_{\ell_0}
 =\dd\mathcal A_{\ell_0}
  +\mathcal A_{\ell_0}\wedge\mathcal A_{\ell_0}
 =
 \begin{pmatrix}
  F_\omega&\ell_0^{-1}T_{\rm tor}\\
  0&0
 \end{pmatrix},
\]
où
\[
 F_\omega^{IJ}
 =\dd\omega^{IJ}+\omega^I{}_K\wedge\omega^{KJ},
 \qquad
 T_{\rm tor}^{I}=D_\omega e^I.
\]
La courbure et la torsion apparaissent comme des composantes distinctes d'un même défaut
de transport. L'égalité matricielle découle algébriquement des données
explicites \(e,\omega,\ell_0\), mais ne construit pas ces entrées à partir de TPK.

Le contenu fort du pont est la commutativité des holonomies :
\[
 \boxed{
 \Hol_{\mathcal A_{\ell_0}}
 \bigl(\mathfrak R_{\rm grav}(\gamma)\bigr)
 =
 \rho_{\rm C}\!\left(
 \Hol_{\rm TPK}(\gamma)
 \right)}
 \qquad
 (\gamma\in\mathsf P_{\rm TPK}^{\rm enr}).
\]
L'égalité doit être compatible avec les identités, l'inversion, la concaténation,
la subdivision, le changement de point de base et la conjugaison. Alors seulement la mémoire
discrète est représentée dans une connexion ; elle n'est pas rebaptisée
torsion ou courbure.

\begin{criterion}[Conditions d'existence du pont TPK--Cartan]
\label{crit:puente-tpk-cartan}
Il n'existe pas de réalisation possédant la commutativité et la compatibilité déclarées
si l'un quelconque des faits suivants se produit :
\begin{enumerate}
 \item deux représentants équivalents d'une même voie produisent
 des holonomies géométriques non conjuguées ;
 \item l'image ne préserve pas les identités, les inverses ou la concaténation ;
 \item une subdivision admissible altère l'holonomie ;
 \item l'égalité des holonomies échoue pour \(\mathcal K_{27}\) ou pour ses
 concaténations de profondeurs \(54\) et \(108\) ;
 \item l'application identifie mémoire et torsion sans exhiber un morphisme qui
 préserve leurs types ;
 \item un choix descendant ---l'action de Holst, \(G\), une masse ou
 une carte SI--- est utilisé pour redéfinir \(A\), \(C^\ast\) ou le retour
 TPK d'origine.
\end{enumerate}
\end{criterion}

\paragraph{Antériorité de la fonctionnelle \(\Gamma_{6\to8}\) par rapport à la réalisation Palatini--Cartan--Holst.}
Le \cref{mdg:chap:barbero} fixe, une fois construits \(A,C^\ast,q_\pm\), la fonctionnelle interne \(\Gamma_{6\to8}\). Le \cref{chap:gravedad} étudie ensuite l'action de Palatini--Cartan--Holst et ses équations variationnelles. Le présent pont n'introduit aucune voie inverse entre les deux étapes. \endgroup

% Owner íntegro de la realización discreta del retorno, la holonomía y la
% respuesta colectiva. Se sitúa después del puente que declara sus hipótesis.
\begin{HMTScientificOwnerSubordinated}
\chapter{Gravité : holonomie, torsion et réponse collective}
\label{chap:gravedad}
\index{gravité}\index{Cartan--Holst}\index{holonomie}
\index{torsion}\index{graviton!non primitif}

La gravité n'est pas introduite ici au moyen d'une formule décimale pour \(G\).
Sa construction réunit deux branches déjà typées. Dans la première, une voie conserve
orientation et mémoire ; une connexion compare des repères en des points distincts ; son
holonomie mesure le retour et sa courbure, le défaut de fermeture. Dans la seconde,
la carte dimensionnelle construit \(G_{\rm int}(\theta_G)\) et
\(\kappa_{\rm int}\) à partir de l'échelle interne d'action, de la vitesse et
de la durée élémentaire. Ce n'est qu'après avoir disposé des deux branches que l'on écrit
l'action Palatini--Cartan--Holst. Aucune valeur métrologique externe ne participe à
cette antériorité.

\section{Retours : position, orientation, phase et mémoire}

Soit \(F_{\rm obs}\) la chronologie observable du transport et
\[
 \mathcal K_{27}=F_{\rm obs}^{27}.
\]
La hiérarchie exacte distingue trois retours :
\[
 \begin{array}{c|c}
 27\ \text{pas}&
 \text{retour de position et inversion d’orientation},\\
 54\ \text{pas}&
 \text{retour de position et d’orientation},\\
 108\ \text{pas}&
 \text{retour complet de la phase observable}.
 \end{array}
\]
Le retour observable n'efface pas le registre accumulé. Si \(g(e)\) compte
une inversion d'orientation et \(s(e)\) un changement effectif de feuillet, alors
\[
 c(\gamma)=\sum_{e\subset\gamma}(g(e),s(e))
\]
satisfait la loi cocyclique de concaténation. Dans les blocs canoniques,
\[
 c_{27}=(1,18),\qquad c_{108}=(4,72).
\]
Par conséquent, la projection observable de \(\mathcal K_{27}\) est d'ordre
quatre, tandis que la coordonnée cumulative continue de croître. Revenir à la
même phase ne signifie pas revenir à la même histoire.

\begin{theorem}[Revêtement cumulatif du retour]
\label{thm:cobertura-retorno-grav-rev6}
Sur \(\mathcal O_{108}\times\mathbb N^2\), le relèvement
\[
 \widetilde F(x,m)=
 \bigl(F_{\rm obs}x,m+c(x\to F_{\rm obs}x)\bigr)
\]
satisfait
\[
 \widetilde{\mathcal K}_{27}(x,m)
 =\bigl(\mathcal K_{27}x,m+(1,18)\bigr)
\]
et
\[
 \widetilde{\mathcal K}_{27}^{\,4}(x,m)
 =\bigl(x,m+(4,72)\bigr)\neq(x,m).
\]
\end{theorem}

\begin{proof}
La constance de \(c_{27}\) sur l'orbite et la loi de concaténation donnent
\[
 c_{108}=\sum_{j=0}^{3}c_{27}(\mathcal K_{27}^jx)
 =4(1,18).
\]
La première coordonnée revient ; la seconde conserve la mémoire accumulée.
\end{proof}

Telle est la donnée qu'une réalisation géométrique doit préserver. On n'identifie pas
la mémoire discrète à la torsion : on construit une représentation
qui permette de les comparer.

\section{Du groupoïde des chemins au groupe de Cartan}

Soit \(\mathcal G_{\rm mem}\) le groupoïde des voies enrichies, avec
composition par concaténation et inverse par renversement. Fixons un élément
\(R_4\in\operatorname{Spin}^+(1,3)\) d'ordre quatre et un vecteur temporel
unitaire \(u\in\mathbb R^{1,3}\) fixé par \(R_4\). Si
\[
c(\gamma)=\bigl(c_g(\gamma),c_s(\gamma)\bigr)\in\mathbb Z^2
\]
est le cocycle additif d'inversion et de changement de feuillet, nous définissons
\[
\rho_{\rm C}^{\rm disc}(\gamma)
=
\left(
R_4^{\,c_g(\gamma)},
\ell_0c_s(\gamma)u
\right)
\in\operatorname{Spin}^+(1,3)\ltimes\mathbb R^{1,3}.
\]

\begin{theorem}[Réalisation discrète du cocycle de retour]
\label{thm:realizacion-discreta-cartan-rev6}
\(\rho_{\rm C}^{\rm disc}\) est un foncteur du groupoïde des voies dans le groupe de
Cartan. En particulier, il conserve l'identité, la concaténation et l'inversion, et
transporte le retour \(27\to54\to108\) sans effacer la coordonnée cumulative.
\end{theorem}

\begin{proof}
La loi du produit semi-direct est
\[
(R,a)(S,b)=(RS,a+Rb).
\]
Comme \(R_4u=u\) et que \(c\) est additif,
\[
\begin{aligned}
\rho_{\rm C}^{\rm disc}(\gamma_2)
\rho_{\rm C}^{\rm disc}(\gamma_1)
&=
\left(
R_4^{c_g(\gamma_2)+c_g(\gamma_1)},
\ell_0\bigl(c_s(\gamma_2)+c_s(\gamma_1)\bigr)u
\right)\\
&=\rho_{\rm C}^{\rm disc}(\gamma_2\gamma_1).
\end{aligned}
\]
La voie identité a un cocycle nul et le renversement change son signe, de sorte
que l'identité et l'inverse sont préservés. Les formules de retour découlent de
\(c_{27}=(1,18)\) et de \(c_{108}=4c_{27}\).
\end{proof}

Ce théorème construit le pont discret. Une réalisation lisse ultérieure
consiste en une représentation
\[
 \rho_{\rm C}:\mathcal G_{\rm mem}
 \longrightarrow\operatorname{Spin}(1,3)\ltimes\mathbb R^{1,3}
\]
et en des champs \(e^I,\omega^{IJ}\) dont l'holonomie entrelace cette représentation.
Le diagramme qui fixe le type est
\[
 \operatorname{Hol}_{\mathcal A_{\ell_0}}
   \bigl(\mathfrak R_{\rm C}(\gamma)\bigr)
 =
 \rho_{\rm C}\bigl(\operatorname{Hol}_{\rm TPK}(\gamma)\bigr).
\]
Les identités, l'inversion, la concaténation et la subdivision doivent être conservées. Cette
équation ne rebaptise pas la retenue courbure : elle exige que le transport
discret déjà représenté et le transport géométrique réalisent la même composition.

Une échelle positive \(\ell_0\) étant fixée, la connexion conjointe est
\[
 \mathcal A_{\ell_0}
 =
 \begin{pmatrix}
 \omega&\ell_0^{-1}e\\
 0&0
 \end{pmatrix}.
\]

\begin{proposition}[Courbure conjointe de Cartan]
\label{prop:curvatura-conjunta-cartan-rev6}
La courbure de \(\mathcal A_{\ell_0}\) est
\[
 \mathcal F_{\ell_0}
 =\dd\mathcal A_{\ell_0}
  +\mathcal A_{\ell_0}\wedge\mathcal A_{\ell_0}
 =
 \begin{pmatrix}
 F_\omega&\ell_0^{-1}T_{\rm tor}\\
 0&0
 \end{pmatrix},
\]
où
\[
 F_\omega^{IJ}
 =\dd\omega^{IJ}+\omega^I{}_K\wedge\omega^{KJ},
 \qquad
 T_{\rm tor}^{I}=D_\omega e^I.
\]
\end{proposition}

\begin{proof}
On multiplie les matrices de formes en utilisant le produit extérieur. Le
bloc diagonal produit \(F_\omega\) ; le bloc de translation produit
\(\dd e+\omega\wedge e=D_\omega e\).
\end{proof}

La courbure et la torsion sont deux composantes d'un même défaut de transport,
mais remplissent des fonctions distinctes. La courbure mesure le retour des repères ;
la torsion mesure la fermeture de la cotétrade.

\section{Gauss--Stokes et capacité surfacique}

Dans un complexe cellulaire fini, pour toute \(1\)-cochaîne \(\Omega\) et toute
\(2\)-chaîne \(z\),
\[
 \boxed{\quad
 \langle\delta\Omega,z\rangle
 =\langle\Omega,\partial z\rangle.
 \quad}
\]
L'identité est la définition du cobord comme adjoint du bord.
Lorsque \(z\) somme les faces orientées de manière cohérente d'une
pseudovariété, les arêtes intérieures s'annulent et seul subsiste le
bord géométrique. La même égalité a déjà été lue comme
vorticité--circulation et comme courbure--holonomie.

L'échelle surfacique possède également un modèle exact. Dans le graphe cubique
\(N\times N\times N\), toute coupe entre deux faces opposées contient au moins
\(N^2\) arêtes et une couche transversale atteint la borne. Pour \(N=9\),
\[
 9^3=729=27^2.
\]
Cette égalité cardinale relie volume et coupe minimale ; elle n'identifie pas les
groupes \((\mathbb Z/9)^3\) et \((\mathbb Z/27)^2\), qui ont des exposants
distincts.

\section{L'action Palatini--Cartan--Holst}

Soit \(M\) une variété lisse orientée de dimension quatre, munie d'une cotétrade non
dégénérée \(e^I\), d'une connexion de Lorentz
\(\omega^{IJ}=-\omega^{JI}\), et de conditions au bord annulant le terme
variationnel de bord. Écrivons
\[
 \Sigma^{IJ}=e^I\wedge e^J,\qquad
 (\star X)^{IJ}=\frac12\epsilon^{IJ}{}_{KL}X^{KL},
 \qquad
 G_{\rm int}=G_{\rm int}(\theta_G),\qquad
 \kappa_{\rm int}=\frac{8\pi G_{\rm int}}{c_{\rm int}^{4}}.
\]
Ces deux dernières quantités ne sont pas définies au moyen de l'action : elles proviennent de
la construction constitutive de la \cref{eq:G-interno}. Le
\cref{thm:homogeneidad-pch} démontre auparavant que la combinaison suivante
appartient à la droite d'action.
L'action du premier ordre est
\[
\boxed{
 S[e,\omega,\Psi]
 =\frac1{2\kappa_{\rm int}c_{\rm int}}\int_M
 \Sigma_{IJ}\wedge
 \left(\star F^{IJ}+\frac1\gamma F^{IJ}\right)
 -\frac{\Lambda_{\rm int}}{\kappa_{\rm int}c_{\rm int}}
  \int_M\operatorname{vol}_e
 +S_{\rm m}[e,\omega,\Psi].}
\]
Le paramètre interne construit au chapitre précédent spécialise
\[
 \gamma=\Gamma_{6\to8}.
\]
La substitution est effectuée après avoir obtenu
\(\Gamma_{6\to8}\) à partir de \(A,C^\ast,q_\pm\) et de l'inclusion
hexade--octade ; l'action n'est pas utilisée pour redéfinir ces entrées.

\begin{theorem}[Équation de Cartan--Holst]
\label{thm:cartan-holst-rev6}
Si le courant de spin est normalisé au moyen de
\[
 \delta_\omega S_{\rm m}
 =-\frac12\int_M\delta\omega^{IJ}\wedge\sigma_{IJ},
\]
la variation par rapport à \(\omega\) produit
\[
 \boxed{\quad
 D_\omega
 \left[(\star+\gamma^{-1}I)\Sigma^{IJ}\right]
 =\kappa_{\rm int}c_{\rm int}\,\sigma^{IJ}.
 \quad}
\]
\end{theorem}

\begin{proof}
On utilise \(\delta F^{IJ}=D_\omega\delta\omega^{IJ}\), on intègre par parties et
on annule le coefficient de la variation arbitraire. Le terme
cosmologique ne dépend pas de \(\omega\).
\end{proof}

En signature lorentzienne, \(\star^2=-I\) sur les bivecteurs. Pour
\(\gamma\in\mathbb R\setminus\{0\}\),
\[
 (\star+\gamma^{-1}I)^{-1}
 =\frac{\gamma^2}{1+\gamma^2}
  (-\star+\gamma^{-1}I).
\]
Par conséquent, si le courant de spin s'annule,
\[
 T_{\rm tor}^I=0,\qquad
 \omega=\mathring\omega(e).
\]
En présence de matière spinorielle, la torsion est déterminée algébriquement par le
courant. Elle ne constitue pas, dans l'action minimale, un degré propagatif
indépendant.

\section{Équation métrique, spin et bilan géométrique}

Pour fixer sans ambiguïté la normalisation, nous définissons les sources mesurées dans
la droite d'action par
\[
 \delta_g S_{\rm m}
 =-\frac12\int_M\operatorname{vol}_e\,
 \mathcal T^{S}_{\mu\nu}\,\delta g^{\mu\nu}.
\]
Lorsque les coordonnées de la tétrade ont la dimension d'une longueur, nous écrivons
\[
 T^{\rm phys}_{\mu\nu}:=c_{\rm int}\mathcal T^{S}_{\mu\nu},
 \qquad
 U^{\rm phys,spin}_{\mu\nu}
 :=c_{\rm int}\mathcal U^{S,\rm spin}_{\mu\nu}.
\]
Après élimination algébrique de la connexion, l'équation métrique adopte les
deux formes équivalentes
\[
 \boxed{
 G_{\mu\nu}+\Lambda_{\rm int}g_{\mu\nu}
 =\kappa_{\rm int}c_{\rm int}
 \bigl(\mathcal T_{\mu\nu}^{S,\rm LC}
       +\mathcal U_{\mu\nu}^{S,\rm spin}\bigr)
 =\kappa_{\rm int}
 \bigl(T_{\mu\nu}^{\rm phys,LC}
       +U_{\mu\nu}^{\rm phys,spin}\bigr).}
\]
La contribution de spin est quadratique en le courant correspondant et
dépend du couplage de matière ; son coefficient n'est pas fixé sans déclarer
l'action fermionique et les conventions de dualité.

Les identités géométriques
\[
 D_\omega T_{\rm tor}^{I}=F^{I}{}_{J}\wedge e^{J},
 \qquad
 D_\omega F^{IJ}=0
\]
impliquent, après élimination de la torsion,
\[
 \nabla^\mu\bigl(T_{\mu\nu}^{\rm phys,LC}
             +U_{\mu\nu}^{\rm phys,spin}\bigr)=0.
\]
Ce bilan de Bianchi--Noether est une condition de cohérence du
transducteur gravitationnel : aucun terme supplémentaire ne peut être incorporé sans
provenir d'une action compatible ou satisfaire le même bilan.

\section{Correspondance surfacique–volumétrique \(\pi/\varphi^2\)}

Le calendrier \((\Sigma,\Pi,\Pi)\) produit la matrice d'incidence
\[
 F_{\rm av}=\begin{pmatrix}0&1\\1&1\end{pmatrix}.
\]
Sa mesure de Parry donne
\(\mu_{\rm ar}/\mu_{\rm vol}=\varphi^{-2}\), et le lecteur régional marque une
seule fois la fermeture \(\pi\). Le retour marqué satisfait
\[
 \Theta_n
 =\pi\frac{F_{n-1}}{F_{n+1}}
 \longrightarrow
 \frac{\pi}{\varphi^2}.
\]
Ce rapport n'est pas une autre version de \(G\). C'est la correspondance adimensionnelle entre
lecture surfacique et lecture volumétrique. Sous une même amplitude positive
\(\mathcal M\),
\[
 \mathcal I_{\rm av}=\mu_{\rm vol}\mathcal M,\qquad
 \mathcal G_{\rm av}=\pi\mu_{\rm ar}\mathcal M
\]
donnent exactement
\[
 \frac{\mathcal G_{\rm av}}{\mathcal I_{\rm av}}
 =\frac{\pi}{\varphi^2}.
\]
Sur le feuillet plan, la restriction diagonale est
\(\mathcal G_{\rm tor}=\mathcal I_{\rm tor}\). L'équivalence locale et
l'ambivalence de bord sont donc deux régimes du lecteur, non deux
valeurs incompatibles d'une masse.

\section{Typage de \(G\) comme transducteur gravitationnel et ordre causal de ses lecteurs}

Le type et l'homogénéité de \(G\) appartiennent à la construction
constitutive préalable. La \cref{eq:G-interno} fixe une seule fois la famille
covariante \(G_{\rm int}(\theta_G)\), et le
\cref{thm:homogeneidad-pch} démontre que son insertion dans l'action possède
la dimension d'une action. Ce chapitre ne redéfinit pas cette famille : il utilise sa
section avec l'holonomie déjà construite. La causalité n'est pas linéaire,
mais convergente :
\[
\begin{aligned}
 \mathcal B_{\rm geom}&:
 \text{retour et cocycle}\to\text{représentation de Cartan}
 \to(e,\omega),\\
 \mathcal B_{\rm dim}&:
 (\hbar_{\rm int},c_{\rm int},t_0,\theta_G)
 \to G_{\rm int}\to\kappa_{\rm int},
\end{aligned}
\]
\[
 (\mathcal B_{\rm geom},\mathcal B_{\rm dim})
 \to S_{\rm PCH}
 \to\text{torsion algébrique, équation métrique et bilan}.
\]
La comparaison métrologique intervient à la fin et ne sélectionne aucune des
deux branches.

\Needspace{22\baselineskip}
\section{Origine collective du mode gravitationnel dans l'holonomie des voies}

L'action précédente a pour objets fondamentaux la cotétrade et la connexion.
La torsion de Cartan est algébrique ; la courbure possède le secteur propagatif qui
contient les ondes gravitationnelles. Le formalisme n'a donc pas besoin de
postuler un graviton élémentaire pour définir l'interaction. Une éventuelle
quantification des perturbations peut produire un mode collectif
d'hélicités \(\pm2\) ; ce mode serait une excitation du transport
géométrique, non une donnée ajoutée au niveau d'APP, de TRIT ou de TPK.

Cette distinction exclut également une confusion fréquente. L'existence
d'ondes classiques démontre une propagation de courbure ; elle ne démontre pas à elle seule
un quantum fondamental. Réciproquement, décrire le champ au moyen de la
géométrie n'élimine pas ses degrés radiatifs. HMT--MD situe la question au
bon endroit : on construit d'abord la mémoire et son holonomie ; on étudie ensuite
le spectre de ses excitations.
 \end{HMTScientificOwnerSubordinated}

\paragraph{Comparaison typée du quantum d’aire avec l’énergie, l’entropie et l’information d’horizon.} L'action du premier ordre permet de comparer le quantum d'aire à l'énergie, à l'entropie et à l'information d'horizon sans confondre ces observables. \endgroup

\chapter{Convergence de l'horizon, de l'aire, de l'information, de l'entropie et de l'action}
\begingroup
\renewcommand{\chapter}[2][]{}%
\chapter{Confluence horizon–aire–information–entropie–action}
\label{chap:parte-iv-69}

La région d’horizon coordonne énergie, température, entropie et action au moyen d’identités typées. La cellule de demi-action et la fermeture chronologique sont présentées avec leurs contrôles de Schwarzschild et de Sitter, sans transformer une interface en définition rétroactive.

\section{Énergie d’horizon et cellule de demi-action}
\label{rm:chap:horizonte-media-accion}

\subsection{Conventions géométriques de la région de Sitter}

Considérons un espace de Sitter quadridimensionnel de rayon aréal \(R>0\), de vitesse causale \(c\) et de constante gravitationnelle \(G\). Adoptons la convention
\begin{equation}
 H_R=\frac cR,
 \qquad
 \Lambda_R=\frac3{R^2},
 \qquad
 \ell_P^2=\frac{G\hbar}{c^3}.
 \label{rm:eq:horizon-ds-conventions}
\end{equation}
La densité d’énergie associée au terme cosmologique est
\begin{equation}
 \varepsilon_\Lambda
 :=\frac{\Lambda_Rc^4}{8\pi G}.
 \label{rm:eq:horizon-vacuum-density}
\end{equation}
La section spatiale plate de la région contient, jusqu’au rayon \(R\), le volume
\begin{equation}
 V_R=\frac{4\pi}{3}R^3.
 \label{rm:eq:horizon-flat-volume}
\end{equation}
Ces conventions fixent positivement l’énergie quasi locale utilisée dans ce chapitre. Le choix d’un autre signe pour la charge de Killing de Sitter relève d’une construction différente et doit être traduit avant toute comparaison des formules.

\subsection{Énergie du vide et énergie de Misner--Sharp}

Dans une géométrie à symétrie sphérique, soit
\begin{equation}
 \dd s^2=g_{ab}(x)\dd x^a\dd x^b
 +\mathcal R(x)^2\dd\Omega_2^2.
 \label{rm:eq:horizon-spherical-metric}
\end{equation}
L’énergie de Misner--Sharp \cite{rm:MisnerSharp1964}, dans la convention positive adoptée, est définie par
\begin{equation}
 E_{\rm MS}(\mathcal R)
 :=\frac{c^4\mathcal R}{2G}
 \left(1-g^{ab}\nabla_a\mathcal R\nabla_b\mathcal R\right).
 \label{rm:eq:horizon-Misner-Sharp}
\end{equation}
Une sphère marginale satisfait
\begin{equation}
 g^{ab}\nabla_a\mathcal R\nabla_b\mathcal R=0.
 \label{rm:eq:horizon-marginal-sphere}
\end{equation}

\begin{theorem}[Confluence volumétrique et quasi locale]
\label{rm:thm:horizon-vacuum-MS}
À l’horizon de Sitter \(\mathcal R=R\), l’énergie intégrée du vide et l’énergie de Misner--Sharp coïncident exactement :
\begin{equation}
 \boxed{
 E_\Lambda
 :=\varepsilon_\Lambda V_R
 =E_{\rm MS}(R)
 =\frac{c^4R}{2G}.}
 \label{rm:eq:horizon-E-Lambda}
\end{equation}
\end{theorem}

\begin{proof}
Les équations \cref{rm:eq:horizon-ds-conventions,rm:eq:horizon-vacuum-density,rm:eq:horizon-flat-volume} donnent
\[
 \varepsilon_\Lambda V_R
 =\frac{3c^4}{8\pi GR^2}\frac{4\pi R^3}{3}
 =\frac{c^4R}{2G}.
\]
La condition marginale \cref{rm:eq:horizon-marginal-sphere} appliquée à \cref{rm:eq:horizon-Misner-Sharp} produit le même résultat.
\end{proof}

L’égalité réunit deux expressions géométriques. La première intègre une densité volumétrique ; la seconde évalue le défaut de piégeage de la sphère aréale. Toutes deux conservent le rayon \(R\) et la compliance gravitationnelle \(G/c^4\).

\subsection{Température, entropie et produit d’horizon}

La gravité de surface de Sitter a pour module \(c^2/R\). La température de Gibbons--Hawking \cite{rm:GibbonsHawking1977} et l’entropie de Bekenstein--Hawking \cite{rm:Bekenstein1973,rm:Hawking1975} sont
\begin{align}
 k_BT_H&=\frac{\hbar c}{2\pi R},
 \label{rm:eq:horizon-temperature}\\
 \frac{S_H}{k_B}
 &=\frac{A_H}{4\ell_P^2}
 =\frac{4\pi R^2}{4\ell_P^2}
 =\frac{\pi R^2}{\ell_P^2}.
 \label{rm:eq:horizon-entropy}
\end{align}

\begin{theorem}[Identité thermodynamique de Sitter]
\label{rm:thm:horizon-TS}
Pour tout rayon \(R\) de la famille déclarée,
\begin{equation}
 \boxed{
 T_HS_H=\frac{c^4R}{2G}=E_\Lambda.}
 \label{rm:eq:horizon-TS-E}
\end{equation}
\end{theorem}

\begin{proof}
La multiplication de \cref{rm:eq:horizon-temperature,rm:eq:horizon-entropy} donne
\[
 T_HS_H
 =\frac{\hbar c}{2\pi Rk_B}
  \frac{k_B\pi R^2}{\ell_P^2}
 =\frac{\hbar cR}{2\ell_P^2}.
\]
L’identité \(\ell_P^2=G\hbar/c^3\) transforme le dernier membre en \(c^4R/(2G)\), égal à \cref{rm:eq:horizon-E-Lambda}.
\end{proof}

Le facteur \(1/2\) possède une décomposition géométrique transparente. Le facteur aréal \(4\pi\), la normalisation holographique par quatre et la périodicité thermique \((2\pi)^{-1}\) satisfont
\begin{equation}
 \frac{4\pi}{4}\frac1{2\pi}=\frac12.
 \label{rm:eq:horizon-half-factor}
\end{equation}
La moitié est donc la contraction conjointe entre quantification aréale et périodicité euclidienne.

\subsection{Quantum complet de l’horizon et force de Planck}

Définissons la tension ou force de Planck et le quantum complet associé au rayon par
\begin{equation}
 \mathscr F_P:=\frac{c^4}{G},
 \qquad
 \boxed{
 \mathcal Q_H(R):=2T_HS_H
 =\frac{c^4R}{G}=\mathscr F_P\,R.}
 \label{rm:eq:horizon-full-quantum}
\end{equation}
L’énergie quasi locale occupe une demi-cellule du quantum complet :
\begin{equation}
 \boxed{E_\Lambda=T_HS_H=\frac{\mathcal Q_H(R)}2.}
 \label{rm:eq:horizon-half-quantum-general}
\end{equation}

\begin{proposition}[Compliance radiale]
\label{rm:prop:horizon-compliance}
La dérivée du quantum complet par rapport au rayon est la force de Planck, et la dérivée de l’énergie quasi locale en est la moitié :
\begin{equation}
 \frac{\dd\mathcal Q_H}{\dd R}=\mathscr F_P,
 \qquad
 \frac{\dd E_\Lambda}{\dd R}=\frac{\mathscr F_P}{2}.
 \label{rm:eq:horizon-radial-derivatives}
\end{equation}
La compliance gravitationnelle \(\mathscr C_P=G/c^4\) satisfait
\begin{equation}
 R=\mathscr C_P\mathcal Q_H.
 \label{rm:eq:horizon-compliance-law}
\end{equation}
\end{proposition}

\begin{proof}
Les trois identités sont des dérivations ou des inversions directes de \cref{rm:eq:horizon-full-quantum}.
\end{proof}

L’équation \cref{rm:eq:horizon-compliance-law} exprime la réponse radiale : le quantum complet applique une tension et la géométrie exprime un rayon au moyen de la compliance universelle. Le produit \(T_HS_H\) occupe la moitié quasi locale de cette réponse.

\subsection{Spécialisation à la fermeture chronologique de 108 états}

La fermeture chronologique de \cref{rm:ch:cosmo} détermine
\begin{equation}
 R_{108}=\ell_P,
 \qquad
 t_P=\frac{\ell_P}{c},
 \qquad
 T_{108}=2\pi t_P,
 \label{rm:eq:horizon-CT108-scales}
\end{equation}
et le quantum de Planck
\begin{equation}
 E_P=\frac{\hbar}{t_P}
 =k_B\Theta_{\rm clk}\log3.
 \label{rm:eq:horizon-Planck-trit}
\end{equation}

\begin{theorem}[Cellule CT108 de demi-action]
\label{rm:thm:horizon-half-action}
Dans la spécialisation \(R=R_{108}=\ell_P\),
\begin{align}
 \mathcal Q_H(\ell_P)
 &=\frac{c^4\ell_P}{G}=E_P,
 \label{rm:eq:horizon-QH-EP}\\
 \boxed{
 E_\Lambda=T_HS_H
 =\frac{E_P}{2}
 =\frac{k_B\Theta_{\rm clk}\log3}{2},}
 \label{rm:eq:horizon-CT108-half-energy}\\
 \boxed{E_\Lambda t_P=\frac{\hbar}{2},}
 \qquad
 \boxed{E_\Lambda T_{108}=\frac h2.}
 \label{rm:eq:horizon-half-actions}
\end{align}
\end{theorem}

\begin{proof}
La première égalité s’obtient en substituant \(\ell_P=\sqrt{G\hbar/c^3}\) et \(E_P=\sqrt{\hbar c^5/G}\). \cref{rm:eq:horizon-half-quantum-general} produit \(E_\Lambda=E_P/2\), et \cref{rm:eq:horizon-Planck-trit} fournit la carte informationnelle. Enfin,
\[
 E_\Lambda t_P=\frac{E_Pt_P}{2}=\frac\hbar2
\]
et, puisque \(T_{108}=2\pi t_P\) et \(h=2\pi\hbar\),
\[
 E_\Lambda T_{108}
 =\frac{E_P}{2}(2\pi t_P)=\pi\hbar=\frac h2.
\]
\end{proof}

La première égalité de \cref{rm:eq:horizon-half-actions} exprime une demi-action réduite sur le temps de Planck ; la seconde exprime une demi-action complète sur la période CT108. Toutes deux représentent une même cellule dans les cartes angulaire et circulaire.

La température et l’entropie sont également déterminées :
\begin{equation}
 k_BT_H=\frac{E_P}{2\pi},
 \qquad
 \boxed{\frac{S_H}{k_B}=\pi.}
 \label{rm:eq:horizon-CT108-TS}
\end{equation}
La température convertit la fréquence d’horizon en énergie ; l’entropie convertit l’aire de la sphère de Planck en information. Leur produit retrouve le même demi-quantum.

\subsection{Poids effectif de l’horizon et exponentielle hyperbolique}

Définissons le poids thermodynamique effectif
\begin{equation}
 \Omega_{\rm eff}(R):=\exp\!\left(\frac{S_H(R)}{k_B}\right).
 \label{rm:eq:horizon-effective-weight}
\end{equation}
Au rayon CT108,
\begin{equation}
 \boxed{\Omega_{\rm eff}(\ell_P)=\ee^\pi.}
 \label{rm:eq:horizon-effective-e-pi}
\end{equation}
Le poids \(\Omega_{\rm eff}\) est une grandeur positive définie par l’entropie. Sa valeur réelle exprime une multiplicité effective ; l’interpréter comme cardinal entier de microétats exige une discrétisation supplémentaire.

L’opérateur d’incidence aréale--volumétrique est
\begin{equation}
 F_{\rm av}=\begin{pmatrix}0&1\\1&1\end{pmatrix},
 \qquad
 J_\varphi:=\frac{2F_{\rm av}^2-3I}{\sqrt5}.
 \label{rm:eq:horizon-J-phi}
\end{equation}
Un calcul direct donne
\begin{equation}
 J_\varphi^2=I,
 \qquad
 \Spec(J_\varphi)=\{+1,-1\}.
 \label{rm:eq:horizon-J-involution}
\end{equation}
Le flot hyperbolique au temps de fermeture \(\pi\) est
\begin{equation}
 \mathcal E_\pi
 :=\ee^{\pi J_\varphi}
 =\cosh\pi\,I+\sinh\pi\,J_\varphi,
 \label{rm:eq:horizon-hyperbolic-flow}
\end{equation}
et possède le spectre
\begin{equation}
 \boxed{
 \Spec(\mathcal E_\pi)=\{\ee^\pi,\ee^{-\pi}\},
 \qquad
 \rho(\mathcal E_\pi)=\ee^\pi.}
 \label{rm:eq:horizon-hyperbolic-spectrum}
\end{equation}

\begin{corollary}[Confluence scalaire horizon--Perron]
\label{rm:cor:horizon-Perron-scalar}
À la fermeture CT108,
\begin{equation}
 \boxed{
 \frac{S_H}{k_B}
 =\log\Omega_{\rm eff}
 =\log\rho(\mathcal E_\pi)
 =\pi.}
 \label{rm:eq:horizon-Perron-log}
\end{equation}
\end{corollary}

\begin{proof}
Les équations \cref{rm:eq:horizon-CT108-TS,rm:eq:horizon-effective-e-pi,rm:eq:horizon-hyperbolic-spectrum} donnent les trois membres.
\end{proof}

L'identité scalaire est exacte et conserve deux généalogies distinctes. L'horizon produit \(\ee^\pi\) par exponentiation de l'entropie aréale ; la dynamique de transfert, construite auparavant, reconnaît le même scalaire comme valeur propre expansive d'une involution hyperbolique. L'égalité terminale fixe une condition de bord pour l'entrelaceur et préserve la distinction des domaines.

\subsection{Cohérence du transducteur gravitationnel avec l’horizon de Schwarzschild}

Un horizon de Schwarzschild de même rayon aréal satisfait
\begin{equation}
 R_H=\frac{2GE_H}{c^4},
 \qquad
 E_H=\frac{c^4R_H}{2G}.
 \label{rm:eq:horizon-Schwarzschild-energy}
\end{equation}
L’entropie aréale coïncide avec \cref{rm:eq:horizon-entropy}, tandis que la température est
\begin{equation}
 k_BT_{\rm Sch}=\frac{\hbar c}{4\pi R_H}
 =\frac12k_BT_H\big|_{R=R_H}.
 \label{rm:eq:horizon-Schwarzschild-temperature}
\end{equation}
En conséquence,
\begin{equation}
 E_H=2T_{\rm Sch}S_H,
 \qquad
 T_H^{\rm dS}=2T_{\rm Sch}
 \label{rm:eq:horizon-Schwarzschild-Smarr}
\end{equation}
pour un même rayon. Dans \(R_H=\ell_P\),
\begin{equation}
 E_H=\frac{E_P}{2},
 \qquad
 T_{\rm Sch}S_H=\frac{E_P}{4}.
 \label{rm:eq:horizon-Schwarzschild-Planck}
\end{equation}
La comparaison sépare les deux géométries : elles partagent l’énergie quasi locale, l’aire et l’entropie, et leurs gravités de surface produisent des températures dans un rapport de deux.

Avec \(\beta=(k_BT_{\rm Sch})^{-1}\), l’action euclidienne satisfait
\begin{equation}
 \frac{I_E}{\hbar}
 =\frac{E_H}{k_BT_{\rm Sch}}-\frac{S_H}{k_B}
 =2\pi-\pi=\pi
 \label{rm:eq:horizon-Schwarzschild-action}
\end{equation}
au rayon de Planck. Le même \(\pi\) apparaît comme entropie réduite, action euclidienne réduite, temps du flot hyperbolique et logarithme du poids effectif.

\subsection{Entrelaceur horizon–Perron : compatibilité spectrale entre énergie de frontière et mémoire de retour}

La coïncidence \(\Omega_{\rm eff}=\rho(\mathcal E_\pi)\) fixe une donnée spectrale et laisse ouverte la relation entre les objets complets. Notons \(\mathcal A_H^{(m)}\) l’algèbre des observables d’horizon au niveau nonadique \(m\), \(\omega_H^{(m)}\) son état, et \(\mathcal P_\varphi^{(m)}\) le système de transfert de Perron.

\begin{definition}[Entrelaceur thermodynamique--dynamique]
\label{rm:def:horizon-Perron-intertwiner}
Un entrelaceur admissible est une famille d’applications linéaires unitales complètement positives (UCP)
\begin{equation}
 \mathfrak I_m:\mathcal A_H^{(m)}\longrightarrow
 \mathcal A_\varphi^{(m)}
 \label{rm:eq:horizon-intertwiner}
\end{equation}
qui satisfait :
\begin{enumerate}[label=\roman*)]
 \item préservation de l’état, en plus de l’unité et de la positivité complète déjà intégrées à la condition UCP ;
 \item compatibilité avec le raffinement nonadique ;
 \item transport de l’involution intérieur--extérieur vers \(J_\varphi\mapsto-J_\varphi\) ;
 \item commutation entre le semi-groupe modulaire d’horizon et le semi-groupe de transfert sur le sous-espace déclaré ;
 \item égalité du rayon spectral au temps \(\pi\), donnée par \cref{rm:eq:horizon-Perron-log}.
\end{enumerate}
\end{definition}

La cinquième condition est nécessaire et scalaire ; les quatre premières fournissent le contenu algébrique qu’une coïncidence numérique ne contient pas. Un entrelaceur construit de cette manière identifierait la multiplicité thermodynamique effective à une dynamique d’expansion et de contraction tout en conservant l’orientation, l’état et le raffinement.

\begin{proofstatusbox}
Les identités \(E_\Lambda=E_{\rm MS}=T_HS_H=c^4R/(2G)\), la définition \(\mathcal Q_H=\mathscr F_P\,R\), la spécialisation CT108, les deux cellules de demi-action et l’égalité \(S_H/k_B=\log\rho(\ee^{\pi J_\varphi})=\pi\) sont des résultats exacts sous les conventions et la réalisation de Sitter déclarées. L’identification de \(\Omega_{\rm eff}\) à un cardinal microscopique et l’entrelaceur \(\mathfrak I_m\) relèvent d’une interface physique délimitée.
\end{proofstatusbox}

\enlargethispage{4\baselineskip}
\begin{falsifierbox}
La fermeture géométrique échoue si \(\Lambda_R\ne3/R^2\), si la sphère de rayon \(R\) ne satisfait pas la condition marginale ou si l’énergie intégrée diffère de \(E_{\rm MS}(R)\). La fermeture thermodynamique échoue si la température ou la loi d’aire produisent \(T_HS_H\ne c^4R/(2G)\). La spécialisation CT108 échoue si \(R_{108}\ne\ell_P\), \(T_{108}\ne2\pi t_P\) ou \(E_P\ne k_B\Theta_{\rm clk}\log3\). L’interface de Perron est réfutée par une obstruction de positivité, d’état, d’orientation ou de raffinement, même si la valeur propre \(\ee^\pi\) coïncide scalairement.
\end{falsifierbox}
 
\paragraph{Factorisation machienne de la mémoire de transport à partir de la fermeture globale et de la réponse locale.} La relation entre fermeture globale et réponse locale conduit au principe de Mach formulé sur la mémoire de transport. \endgroup

\let\HMTMonographChapterInput\relax
  \part{Incidence de frontière et stabilisateurs exceptionnels}

\chapter{Entrelaceur dodécaphasique--Witt}
\section{Construction de l'entrelaceur d'incidence point--hexade et transport de l'action résiduelle de \texorpdfstring{\(S_3\)}{S3}}
\label{chap:entrelazador-accion}
\index{opérateur d'entrelacement!dodécaphase--Witt}
\index{Witt!module d'incidence}
\index{survie!projecteur de rang trois}

L'apparition répétée du quotient \(3/11\) dans l'écran dodécaphasique,
dans la polarité d'une étoile de Witt et dans la filtration nonadique pose
une question plus forte que l'égalité arithmétique de trois fractions. La
question correcte est opératorielle : déterminer si ces trois apparitions sont
des traces du même secteur positif vu dans des réalisations différentes et, si elles le
sont, construire les applications qui les entrelacent sans introduire de donnée
métrologique ultérieure.

La réponse comporte quatre niveaux. Premièrement, l'incidence point--hexade construit une isométrie canonique entre le module dodécaphasique de dimension onze et un espace propre de dimension onze du plan de Witt. Deuxièmement, l'étoile de Witt agit sur ce même module, mais comme involution signée ; son \(3/11\) est un indice, non la trace d'un projecteur positif. Troisièmement, le secteur orienté de survie apparaît dans \(t=7\), où un unique espace à onze branches contient un secteur de rang trois. Quatrièmement, la droite d'orientation de ce triplet, avec les témoins internes \(B_K\) et \(u_K\), sélectionne variationnellement un unique normalisateur \(S_3\), fixe son orientation et produit, par moyenne de Reynolds et décomposition polaire, l'entrelacement isométrique secteur--dodécaphase. Le passage dodécaphase--Witt est une seconde isométrie d'incidence ; on ne présuppose pas que les deux flèches réalisent une unique représentation de \(S_3\). Les moyennes de Reynolds, les décompositions en représentations et le facteur polaire utilisés dans cette section sont les constructions standard \cite{Serre1977Representations,HornJohnson2013}.

\begin{figure}[htbp]
\centering
\begin{tikzpicture}[
  font=\footnotesize,
  box/.style={draw,rounded corners=2mm,align=center,minimum width=2.35cm,minimum height=.98cm},
  arr/.style={-{Stealth[length=2.2mm]},thick},
  link/.style={<->,thick,dashed}
]
  \node[box,draw=hmtblue,fill=hmtlightblue] (seed) at (0,2.2)
    {germes\\\(104\,976\)};
  \node[box,draw=hmtblue,fill=hmtlightblue] (u) at (3,2.2)
    {émissions\\\(\mathcal U_6,\ |\mathcal U_6|=468\)};
  \node[box,draw=hmtgold,fill=hmtlightgold] (ppi) at (6,2.2)
    {microfibra\\\(\Pi_\pi^{(468)},\ \operatorname{rank}=5\)};
  \node[box,draw=hmtred,fill=hmtlightred] (d) at (9,2.2)
    {défaut positif\\\(D_{\rm switch}\)};
  \node[box,draw=hmtviolet,fill=hmtlightviolet] (lam) at (12,2.2)
    {capacidad residual\\\(\lambda_*=1-\delta_*\)};

  \draw[arr,hmtblue] (seed)--node[above=16pt,font=\scriptsize,fill=white,inner sep=1pt]{\(F\)} (u);
  \draw[arr,hmtgold] (u)--node[above=16pt,font=\scriptsize,fill=white,inner sep=1pt]{\(\tau_{468}\)} (ppi);
  \draw[arr,hmtred] (ppi)--node[above=16pt,font=\scriptsize,fill=white,inner sep=1pt]{\(D_{\rm switch}\)} (d);
  \draw[arr,hmtviolet] (d)--node[above=16pt,font=\scriptsize,fill=white,inner sep=1pt]{\(\tau+\mathrm{CR}\)} (lam);

  \node[box,draw=hmtgold,fill=hmtlightgold] (g) at (1.5,-.2)
    {secteur orienté \(t=7\)\\\(L_{\rm or}\otimes(H_G,P_G)\)};
  \node[box,draw=hmtgreen,fill=hmtlightgreen] (v) at (5,-.2)
    {dodécaphase\\\(V_{11},P_3\)};
  \node[box,draw=hmtgreen,fill=hmtlightgreen] (e) at (8.5,-.2)
    {incidence de Witt\\\(E_{30},Q_3\)};
  \node[box,draw=hmtgray,fill=hmtpaper] (t) at (12,-.2)
    {lecteur spectral\\\(\mathcal L_{\beta,*}\to T_{\rm band}\)};

  \draw[arr,hmtgold] (g)--node[above=16pt,font=\scriptsize,fill=white,inner sep=1pt]{\(W_{GK}\)} (v);
  \draw[arr,hmtgreen] (v)--node[above=16pt,font=\scriptsize,fill=white,inner sep=1pt]{\(U_W=I/6\)} (e);
  \draw[arr,hmtgray] (lam)--(t);
  \draw[arr,hmtred] (v) to[bend left=20]
    node[below=3pt,font=\scriptsize,fill=white,inner sep=1pt]{\((3/11)\Delta_4\)} (d);
\end{tikzpicture}
 \caption{Architecture de la fermeture prédimensionnelle. La ligne supérieure construit la mesure et l'action à partir du quotient des émissions ; la ligne inférieure transporte le secteur orienté \(t=7\) vers la dodécaphase, puis la dodécaphase vers Witt. Les flèches continues sont des isométries de projecteurs démontrées. La première est fixée par la sélection variationnelle du \cref{thm:identificacion-isometrica-tres-proyectores} ; pour la seconde, aucune équivariance relativement au \(H_*\) sélectionné n'est affirmée.}
\label{fig:entrelazador-accion}
\end{figure}

\section{Typage des objets associés à \texorpdfstring{\(3/11\)}{3/11}}

Soit
\[
 V_{12}=\mathbb R^{12},
 \qquad
 V_{11}=\mathbf 1^\perp
 =\left\{v\in\mathbb R^{12}:\sum_{p=1}^{12}v_p=0\right\}.
\]
L'action icosaédrique déclarée sur les douze positions se décompose comme
\[
 V_{12}=V_1\oplus V_3\oplus V_{3'}\oplus V_5.
\]
Notons \(P_3\) le projecteur orthogonal sur le premier terme
tridimensionnel. Alors
\[
 P_3^2=P_3=P_3^*,
 \qquad
 \operatorname{rank}P_3=3,
 \qquad
 \tau_{11}(P_3)
 =\frac1{11}\operatorname{Tr}_{V_{11}}P_3
 =\frac3{11}.
\]

Il y a trois objets proches et un quatrième objet de type distinct :
\begin{enumerate}
  \item l'incidence point--hexade transporte \(P_3\) vers un projecteur
  \(Q_3\) de rang trois sur un module de Witt ;
  \item la survie dans \(t=7\) distingue trois branches parmi onze et
  produit un projecteur \(P_G\) de trace normalisée \(3/11\) ;
  \item un opérateur d'entrelacement orienté identifie
  \(P_3V_{11}\) à la réalisation orientée de \(P_GH_G\) ;
  \item une étoile de Witt a sept modes pairs et quatre impairs,
  de sorte que sa trace normalisée vaut également \(3/11\), mais elle contient un spectre
  négatif et n'est pas un projecteur.
\end{enumerate}
La distinction entre projecteur et involution sera essentielle pour que la loi
d'action de la section suivante soit positive.

\section{La matrice d'incidence point--hexade}

Soit \(\mathcal H\) l'ensemble des \(132\) hexades du plan
\[
 W_{12}=S(5,6,12).
\]
On définit
\[
 I:\mathbb R^{12}\longrightarrow\mathbb R^{\mathcal H},
 \qquad
 (Iv)(H)=\sum_{p\in H}v_p,
\]
ou, de manière équivalente,
\[
 I_{H,p}=\mathbf 1_{\{p\in H\}}.
\]
Cet opérateur n'enregistre que l'incidence. Il ne contient ni \(\alpha\), ni unités,
ni masses, ni CKM, ni Barbero, ni aucun objectif physique.

\begin{lemma}[Identité métrique de l'incidence]
\label{lem:incidencia-metrica-witt}
On a
\[
 \boxed{I^*I=36I_{12}+30J_{12}.}
\]
Par conséquent, pour tout \(v\in V_{11}\),
\[
 \|Iv\|^2=36\|v\|^2.
\]
\end{lemma}

\begin{proof}
Dans un système \(S(5,6,12)\), chaque point appartient à
\[
 r=\frac{\binom{11}{4}}{\binom54}=66
\]
hexades, et chaque paire de points appartient à
\[
 \lambda_2=\frac{\binom{10}{3}}{\binom43}=30
\]
hexades. Les entrées diagonales de \(I^*I\) valent \(66\) et les entrées non
diagonales valent \(30\). Donc
\[
 I^*I=(66-30)I_{12}+30J_{12}.
\]
Comme \(J_{12}\) s'annule sur \(\mathbf1^\perp\), on obtient la seconde
identité.
\end{proof}

Pour identifier la branche spectrale sélectionnée par \(I\), nous introduisons
\[
 G_{H,H'}=\binom{|H\cap H'|}{4}.
\]

\begin{lemma}[Identité spectrale de l'incidence]
\label{lem:incidencia-espectral-witt}
Avec \(J_{132,12}\) la matrice rectangulaire de uns,
\[
 \boxed{GI=30I+15J_{132,12}.}
\]
En particulier, si \(v\in V_{11}\), alors
\[
 G(Iv)=30Iv.
\]
\end{lemma}

\begin{proof}
L'entrée \((H,p)\) de \(GI\) est
\[
 \sum_{H'\ni p}\binom{|H\cap H'|}{4}.
\]
Le recensement exhaustif de \(W_{12}\) donne \(45\) lorsque \(p\in H\) et \(15\)
lorsque \(p\notin H\). Par conséquent, l'entrée vaut
\(30I_{H,p}+15\). La reconstruction des \(132\) hexades à partir du code
de Golay ternaire étendu vérifie les \(132\cdot12\) entrées comme des
entiers.
\end{proof}

Définissons
\[
 E_{30}^{\rm inc}:=I(V_{11})\subseteq\ker(G-30I).
\]
Sa dimension est onze par le \cref{lem:incidencia-metrica-witt}. Le calcul
modulo \(1009\) obtient
\[
 \operatorname{rank}_{\mathbb F_{1009}}(G-30I)=121.
\]
Comme la réduction modulo \(1009\) donne une borne inférieure du rang
rationnel, on en déduit
\(\dim_{\mathbb Q}\ker(G-30I)\le11\). L'inclusion précédente fournit
onze vecteurs linéairement indépendants dans ce noyau. Par conséquent,
\[
 E_{30}^{\rm inc}=\ker(G-30I).
\]
Nous écrirons désormais \(E_{30}\).

\begin{theorem}[Opérateur d'entrelacement canonique dodécaphase--Witt]
\label{thm:entrelazador-dodecafase-witt}
L'opérateur
\[
 \boxed{
 U_W=\frac16 I\big|_{V_{11}}:V_{11}\longrightarrow E_{30}
 }
\]
est une isométrie surjective. Pour tout
\(g\in\operatorname{Aut}(W_{12})\cong M_{12}\),
\[
 U_W\rho_{12}(g)=\rho_{\mathcal H}(g)U_W.
\]
Parmi les noyaux point--hexade équivariants, la condition d'isométrie et
le signe positif sur les incidences déterminent \(U_W\) de manière unique.
\end{theorem}

\begin{proof}
L'isométrie découle de \(I^*I=36I\) sur \(V_{11}\). Son image est de
dimension onze et est contenue dans \(E_{30}\), également de dimension onze ;
elle est donc surjective. L'équivariance traduit
\[
 p\in H\quad\Longleftrightarrow\quad g(p)\in g(H).
\]

Pour l'unicité, \(M_{12}\) est transitif sur les drapeaux
\((p,H)\), avec \(p\in H\), et sur les antidrapeaux. Tout noyau
équivariant a deux valeurs, \(a\) sur les drapeaux et \(b\) sur les
antidrapeaux :
\[
 T=aI+b(J-I)=(a-b)I+bJ.
\]
Le terme \(J\) s'annule sur \(V_{11}\), de sorte que \(T\) y est un
multiple de \(I\). L'isométrie impose \(|a-b|=1/6\) et la normalisation
positive sélectionne \(a-b=1/6\).
\end{proof}

\section{Transport positif du projecteur tritique}

L'opérateur d'entrelacement permet de définir
\[
 \boxed{Q_3=U_WP_3U_W^*.}
\]
On obtient
\[
 Q_3^2=Q_3=Q_3^*,
 \qquad
 \operatorname{rank}Q_3=3,
\]
et le diagramme
\[
\begin{CD}
V_{11} @>{U_W}>> E_{30}\\
@V{P_3}VV @VV{Q_3}V\\
V_{11} @>{U_W}>> E_{30}
\end{CD}
\]
commute. En particulier,
\[
 \boxed{\tau_{V_{11}}(P_3)=\tau_{E_{30}}(Q_3)=\frac3{11}.}
\]

La canonicité de \(U_W\) appartient au plan étiqueté et à son action
d'automorphismes. Il faut la distinguer de l'écran \(A_5/C_5\) utilisé pour
définir \(P_3\) : dans la réalisation étiquetée construite, seule l'identité
de cette action concrète de \(A_5\) préserve l'ensemble fixe des hexades. Ainsi,
\(U_W\) transporte canoniquement le \(P_3\) déjà donné vers
\(Q_3\), mais ne transforme pas à lui seul l'action \(A_5/C_5\) en une
sous-action de \(M_{12}\). Le pont avec le secteur orienté sera construit séparément
à partir des données marquées \(K\) et \(B_K\).

\section{Indice signé de l'étoile de Witt}

Soit \(B\) une tétrade. Les quatre hexades qui la contiennent s'écrivent
\[
 H_i=B\sqcup P_i,\qquad |P_i|=2,
\]
où les quatre pétales \(P_i\) partitionnent le complémentaire de \(B\).
L'involution d'étoile \(R_B\) fixe \(B\) et échange les deux points
de chaque pétale. Son type de cycles sur douze points est \(1^4\,2^4\). Dans
\(V_{11}\), et au moyen de \(U_W\) également dans \(E_{30}\), son spectre est
\[
 \operatorname{Spec}(R_B)=\{1^{[7]},(-1)^{[4]}\}.
\]
Par conséquent,
\[
 \boxed{\tau_{11}(R_B)=\frac{7-4}{11}=\frac3{11}.}
\]

\begin{theorem}[Obstruction projecteur--étoile]
\label{thm:no-go-proyector-estrella}
Il n'existe pas d'isomorphisme \(F:V_{11}\to E_{30}\) tel que
\[
 FP_3=R_BF.
\]
\end{theorem}

\begin{proof}
S'il existait, \(P_3\) et \(R_B\) seraient semblables. Cependant,
\[
 \operatorname{Spec}(P_3)=\{1^{[3]},0^{[8]}\},
 \qquad
 \operatorname{Spec}(R_B)=\{1^{[7]},(-1)^{[4]}\}.
\]
Leurs polynômes minimaux sont \(x(x-1)\) et
\((x-1)(x+1)\), respectivement. La similitude est impossible.
\end{proof}

L'égalité \(3/11\) possède ainsi deux lectures sur le même module. Pour
\(Q_3\), c'est la masse d'un secteur positif de rang trois. Pour \(R_B\), c'est
un indice de parité. Dans une action positive, \(Q_3\) doit intervenir ;
l'étoile conserve l'orientation et distingue les feuillets, mais ne remplace pas le
projecteur.

\section{Réduction homogène \texorpdfstring{\(11\to3\to1\)}{11-3-1} en \texorpdfstring{\(t=7\)}{t=7}}

La lecture précédente de \(3/11\) dans \(t=8\) comparait trois racines locales
à onze continuations d'une seule racine. Le numérateur et le dénominateur
n'appartenaient pas au même espace. Le tableau détaillé situe l'objet correct
en \(t=7\) : il y a onze branches locales ; les onze survivent à un horizon, trois
survivent à deux et une survit à trois,
\[
 \boxed{11\longrightarrow3\longrightarrow1.}
\]

\begin{center}
\small
\begin{tabularx}{\textwidth}{cYcc}
\toprule
indice&triplet de mots&\(h=2\)&\(h=3\)\\
\midrule
0&\texttt{200112|212221|110110}&3&0\\
1&\texttt{200121|212212|110110}&3&1\\
2&\texttt{200211|212122|110110}&3&0\\
3&\texttt{202122|210211|110110}&0&0\\
4&\texttt{202212|210121|110110}&0&0\\
5&\texttt{202221|210112|110110}&0&0\\
6&\texttt{210111|202222|110110}&0&0\\
7&\texttt{210222|202111|110110}&0&0\\
8&\texttt{212112|200221|110110}&0&0\\
9&\texttt{212121|200212|110110}&0&0\\
10&\texttt{212211|200122|110110}&0&0\\
\bottomrule
\end{tabularx}
\end{center}

Soit \(\mathcal B_7\) l'ensemble de ces onze branches et
\[
 H_G=\ell^2(\mathcal B_7).
\]
On définit
\[
 P_G\delta_b=
 \begin{cases}
 \delta_b,&b\text{ admet un prolongement jusqu’à l’horizon deux},\\
 0,&\text{sinon}.
 \end{cases}
\]
Alors
\[
 P_G^2=P_G=P_G^*,
 \qquad
 \operatorname{rank}P_G=3,
 \qquad
 \boxed{\tau_{H_G}(P_G)=\frac3{11}.}
\]
C'est la troisième réalisation positive.

\section{L'action résiduelle de \texorpdfstring{\(S_3\)}{S3} et les régions positives de \texorpdfstring{\(\pi\)}{pi}}

Une permutation des trois positions finales, appliquée simultanément aux
deux premiers mots de chaque branche en laissant le troisième fixe, préserve
\(\mathcal B_7\). Cette règle définit une action de \(S_3\). Ses orbites
ont pour tailles
\[
 3+3+1+1+3=11.
\]
Le support de \(P_G\), formé des indices \(0,1,2\), est une orbite
complète, de sorte que \(P_G\) commute avec \(S_3\).

Le caractère de la représentation sur les onze branches, lu sur l'identité,
la transposition et le cycle d'ordre trois, est
\[
 \chi_{\mathcal B_7}=(11,5,2)
 =5\chi_{\rm triv}+3\chi_{\rm std}.
\]
Sur l'image de \(P_G\),
\[
 \chi_G=(3,1,0)
 =\chi_{\rm triv}+\chi_{\rm std}.
\]

Les trois régions positives de la microfibre de \(\pi\) ont pour signatures
\[
 339555,\qquad393555,\qquad933555.
\]
Les queues du premier mot des trois branches survivantes sont
\[
 112,\qquad121,\qquad211.
\]
La substitution en coordonnées \(1\mapsto3\), \(2\mapsto9\) donne
\[
\begin{array}{c|c|c}
\text{branche}&\text{queue}&\text{région positive}\\ \hline
0&112&339555\\
1&121&393555\\
2&211&933555.
\end{array}
\]

\begin{theorem}[Pont horizontal--vertical de la microfibre positive]
\label{thm:puente-pi-puerta-s3}
La règle précédente est une bijection \(S_3\)-équivariante entre les trois
régions positives de \(\pi\) et les trois branches sélectionnées par \(P_G\).
L'application conserve la position de l'unique symbole exceptionnel.
\end{theorem}

\begin{proof}
Les deux ensembles sont les trois mots obtenus en plaçant un unique symbole
exceptionnel dans l'une des trois positions. La substitution s'applique coordonnée par
coordonnée et commute donc avec toute permutation des positions. La
position exceptionnelle retrouve de manière univoque l'élément d'origine.
\end{proof}

\section{Dépendance à l'orientation dans l'opérateur d'entrelacement dodécaphasique--Witt}

Pour tout normalisateur \(H=N_{A_5}(C_3)\simeq S_3\) d'un sous-groupe
d'ordre trois dans \(A_5\), nous définissons le caractère abstrait du quotient
\[
 \operatorname{sgn}_H:H\longrightarrow H/C_3\simeq C_2
 \longrightarrow\{+1,-1\},
 \qquad
 \operatorname{sgn}_H(h)=
 \begin{cases}
 +1,&h\in C_3,\\
 -1,&h\in H\setminus C_3.
 \end{cases}
\]
Ce caractère n'est pas le signe ambiant d'une permutation de cinq lettres :
tout élément de \(A_5\) a pour signe ambiant \(+1\). C'est le caractère
intrinsèque de la réflexion dans le quotient \(H/C_3\). Avec cette convention, le
caractère du secteur icosaédrique restreint à \(H\) est
\[
 \chi_{P_3V_{11}|H}=(3,-1,0)
 =\chi_{\operatorname{sgn}_H}+\chi_{\rm std}.
\]
En revanche,
\[
 \chi_{\operatorname{im}P_G}=(3,1,0)
 =\chi_{\rm triv}+\chi_{\rm std}.
\]

\begin{corollary}[Obstruction à l'opérateur d'entrelacement non orienté]
\label{cor:obstruccion-entrelazador-incidencia-p3}
Il n'existe pas d'isomorphisme \(S_3\)-équivariant
\[
 \operatorname{im}P_G\longrightarrow P_3V_{11}
\]
pour les actions non tordues.
\end{corollary}

\begin{proof}
Des représentations isomorphes possèdent le même caractère. Les valeurs sur une
transposition sont \(1\) et \(-1\).
\end{proof}

Soit \((e_0,e_1,e_2)\) la base positionnelle de
\(\mathbb R^3_{\rm perm}\). Nous définissons
\[
 L_{\rm or}
 =\bigwedge\nolimits^3\mathbb R^3_{\rm perm},
 \qquad
 \ell_{\rm or}=e_0\wedge e_1\wedge e_2
\]
la droite d'orientation du triplet, munie de l'ancrage positif
\(\ell_{\rm or}\). L'action de \(S_3\) sur cette droite est précisément son
caractère de signe abstrait. En tensorisant par elle,
\[
 \chi_{L_{\rm or}\otimes\operatorname{im}P_G}
 =(3,-1,0),
\]
de sorte que l'obstruction des caractères disparaît. L'égalité des
caractères démontre l'existence d'une classe d'isomorphismes, mais n'en
sélectionne pas encore un. La sélection s'obtient à partir de deux témoins déjà présents dans la
dodécaphase.

\section{Sélection variationnelle du normalisateur \texorpdfstring{\(S_3\)}{S3}}

Soit
\[
 K=(234,543,140,729,659,824,621,58,914,794,146,601)
\]
la carte décimale du sceau. Le témoin commun minimal des deux lectures est
\[
 B_K=H_e\cap P_\pi
 =\{p:K_p\ge729\}
 =\{4,6,9,10\},
\]
dans la numérotation mathématique \(1,\ldots,12\). Dans la convention équivalente
avec des indices commençant à zéro, le même ensemble est \(\{3,5,8,9\}\). Nous définissons deux
vecteurs du secteur \(P_3\) :
\[
 b_K=P_3\left(\mathbf1_{B_K}-\frac{|B_K|}{12}\mathbf1\right),
 \qquad
 u_K=P_3\left(K-\overline K\,\mathbf1\right).
\]
Le premier satisfait exactement
\[
 \|b_K\|^2=1.
\]

Le groupe \(A_5\) contient dix sous-groupes \(C_3\). Pour chacun, soit
\[
 H(C_3)=N_{A_5}(C_3)\simeq S_3
\]
et soit le projecteur de signe
\[
 e_{{\rm sgn},H}
 =\frac16\sum_{h\in H}\operatorname{sgn}_H(h)\rho(h).
\]
Ici \(\operatorname{sgn}_H\) est le caractère abstrait de \(H/C_3\)
défini ci-dessus, non le signe ambiant —trivial— de \(A_5\).
L'énergie d'orientation associée au témoin commun est
\[
 \mathcal E_B(H)=\|e_{{\rm sgn},H}b_K\|^2.
\]
Le recensement exact dans \(\mathbb Q(\sqrt5)\) produit un maximum unique :
\[
 \mathcal E_B(H_*)
 =\frac23+\frac{2}{15}\sqrt5.
\]
Le deuxième niveau est
\[
 \frac13+\frac{2}{15}\sqrt5,
\]
de sorte que l'écart exact est
\[
 \boxed{\mathcal E_B(H_*)-\mathcal E_B(H_{\rm segundo})=\frac13.}
\]
Le même \(H_*\) est le maximum unique lorsque \(b_K\) est remplacé par \(u_K\).
Le sélecteur ne dépend donc pas d'un seul codage du témoin.

Dans cet étiquetage,
\[
 H_*=N_{A_5}\bigl(\langle c_*\rangle\bigr),
\qquad
 c_*=(2\,3\,4).
\]
La règle est naturelle sous un réétiquetage simultané par \(A_5\) : en
transportant \(B_K\), \(K\) et \(P_3\), le maximum est conjugué par le même
élément.

\section{Sélection de l'orientation et présentation de \texorpdfstring{\(S_3\)}{S3}}

Dans \(H_*\), considérons
\[
 s(g)=\langle b_K,\rho(g)u_K\rangle.
\]
Le maximum parmi les deux éléments d'ordre trois sélectionne de manière univoque
\[
 c_*=(0,1,3,4,2),
\]
dans la notation des images de la permutation de cinq points, et le maximum parmi
les trois involutions sélectionne
\[
 r_*=(1,0,4,3,2).
\]
On vérifie exactement
\[
 r_*c_*r_*=c_*^{-1}.
\]
La présentation du secteur orienté au moyen du cycle de positions
\(c=(0\,1\,2)\) et de la réflexion \(r=(0\,2)\), qui fixe la position centrale,
détermine alors un unique isomorphisme
\[
 \theta:S_3^{\rm or}\xrightarrow{\sim}H_*,
 \qquad
 \theta(c)=c_*,
 \qquad
 \theta(r)=r_*.
\]
Ici \(S_3^{\rm or}\) désigne le groupe des réétiquetages du secteur orienté
\(t=7\), engendré par \(c\) et \(r\) ; il n'introduit pas de classe opératoire
supplémentaire.
L'unique branche qui atteint l'horizon trois est précisément la branche centrale
\(1\) ; cette donnée fixe laquelle des trois positions reste immobile sous la
réflexion.

\section{Moyenne de Reynolds et facteur polaire}

Soit
\[
 W_{\rm or}=L_{\rm or}\otimes\mathbb R^3_{\rm perm}
\]
le triplet du secteur orienté. L'ancrage étant déjà fixé, nous prenons
\[
 \boxed{w=\ell_{\rm or}\otimes e_1,}
\]
où \(e_1\) représente la position centrale, la seule qui atteint l'horizon
trois. Nous définissons
\[
 A
 =\sum_{g\in S_3}
 \rho_{W_{\rm or}}(g)\,
 |w\rangle\langle b_K|\,
 \rho_{V_{11}}(\theta(g))^{-1}.
\]
Par construction,
\[
 A\rho_{V_{11}}(\theta(g))
 =\rho_{W_{\rm or}}(g)A,
\qquad
 A=AP_3.
\]
Le calcul exact dans \(\mathbb Q(\sqrt5)\) obtient
\[
 \operatorname{rank}A=3.
\]
De plus, \(G_A=AA^*\) a toutes ses entrées diagonales égales à
\[
 3+\frac25\sqrt5
\]
et toutes ses entrées non diagonales égales à
\[
 \frac52+\frac35\sqrt5.
\]
Ses deux valeurs propres isotypiques sont
\[
 \lambda_{\rm sgn}=8+\frac85\sqrt5,
 \qquad
 \lambda_{\rm std}=\frac12-\frac15\sqrt5.
\]
Toutes deux sont strictement positives. En particulier, \(G_A^{-1/2}\) existe et
la décomposition polaire
\[
 \boxed{U_{GK}=G_A^{-1/2}A}
\]
est une isométrie équivariante de \(P_3V_{11}\) sur \(W_{\rm or}\) :
\[
 U_{GK}^*U_{GK}=P_3,
 \qquad
 U_{GK}U_{GK}^*=I_{W_{\rm or}}.
\]

Soit
\[
 J_G:W_{\rm or}\longrightarrow L_{\rm or}\otimes H_G
\]
l'inclusion isométrique dans les trois branches sélectionnées. Alors
\[
 J_GJ_G^*=I_{L_{\rm or}}\otimes P_G.
\]
L'isométrie partielle secteur--dodécaphase est notée
\[
 \boxed{
 W_{GK}=U_{GK}^*J_G^*
 :L_{\rm or}\otimes H_G\longrightarrow V_{11}
 }
\]
satisfait
\[
 W_{GK}^*W_{GK}
 =I_{L_{\rm or}}\otimes P_G,
\qquad
 W_{GK}W_{GK}^*=P_3,
\]
et, par conséquent,
\[
 P_3W_{GK}
 =W_{GK}
 =W_{GK}(I_{L_{\rm or}}\otimes P_G).
\]

\begin{theorem}[Identification isométrique relative de trois projecteurs]
\label{thm:identificacion-isometrica-tres-proyectores}
Relativement à l'action \(A_5/C_5\), au projecteur \(P_3\), au sceau
\(K\), au témoin doublement construit \(B_K\), au secteur fini orienté
\(t=7\), au caractère abstrait \(\operatorname{sgn}_{H_*}\), à l'ancrage
\(\ell_{\rm or}\) et à la règle variationnelle \(\mathcal E_B\), il existe une
unique isométrie partielle normalisée \(W_{GK}\) qui identifie \(P_G\) à
\(P_3\). Indépendamment, l'isométrie d'incidence \(U_W\) identifie
\(P_3\) à \(Q_3\). Par conséquent, on obtient la chaîne suivante
d'identifications isométriques de projecteurs :
\[
\begin{CD}
L_{\rm or}\otimes H_G
 @>{W_{GK}}>>
 V_{11}
 @>{U_W}>>
 E_{30}\\
@V{I_{L_{\rm or}}\otimes P_G}VV
 @VV{P_3}V
 @VV{Q_3}V\\
L_{\rm or}\otimes H_G
 @>{W_{GK}}>>
 V_{11}
 @>{U_W}>>
 E_{30}.
\end{CD}
\]
Les traces typées satisfont
\[
 \tau_{L_{\rm or}\otimes H_G}(I_{L_{\rm or}}\otimes P_G)
 =\tau_{V_{11}}(P_3)
 =\tau_{E_{30}}(Q_3)
 =\frac3{11}.
\]
L'énoncé identifie isométriquement trois projecteurs ; il n'affirme pas que leurs trois espaces soient une même représentation d'un groupe commun.
\end{theorem}

\begin{proof}
L'énergie \(\mathcal E_B\) sélectionne un unique normalisateur \(H_*\) ; le critère de score \(s(g)\) sélectionne un unique couple \((c_*,r_*)\) et, avec lui, un unique isomorphisme \(\theta\). La moyenne de Reynolds est alors équivariante pour les actions de \(S_3^{\rm or}\) et de \(H_*\) reliées par \(\theta\), et possède le rang trois. La partie positive de sa décomposition polaire fixe la normalisation sans introduire de base. Les identités partielles précédentes démontrent que \(W_{GK}\) identifie \(P_G\) et \(P_3\). D'autre part, le \cref{thm:entrelazador-dodecafase-witt} identifie \(P_3\) et \(Q_3\) comme projecteurs au moyen de \(U_W\). Ces deux identités opératorielles rendent le diagramme commutatif sans étendre l'équivariance de la première flèche à la seconde.
\end{proof}

\begin{proposition}[Obstruction à l'équivariance triple dans le calibre fixé]
\label{prop:no-go-equivarianza-hstar-witt}
La construction précédente ne démontre pas que \(U_W\) soit
\(H_*\)-équivariante. Dans l'étiquetage déclaré, l'action fixe de
\(A_5/C_5\) sur les douze points ne préserve l'ensemble des hexades de
\(W_{12}\) que pour l'élément identité. Ainsi,
l'entrelacement \(H_*\)-équivariant de \(W_{GK}\) ne se prolonge pas, avec ce
calibre, en une représentation commune sur \(E_{30}\).
\end{proposition}

\begin{proof}
Le premier énoncé est une séparation des domaines d'équivariance :
\(W_{GK}\) utilise \(H_*\subset A_5\), tandis que le théorème d'incidence
prouve l'équivariance de \(U_W\) pour
\(\operatorname{Aut}(W_{12})\). Le recensement exact reproductible applique les soixante
éléments de l'action \(A_5/C_5\) au système d'hexades étiqueté ; seule
l'identité conserve ce système. Comme \(H_*\) contient six éléments, son
action ne peut pas entrer par restriction de cette action fixe dans
\(\operatorname{Aut}(W_{12})\).
\end{proof}

Pour promouvoir la chaîne de projecteurs en représentation commune, il faut construire explicitement un monomorphisme
\[
 \boxed{
 \iota:H_*\hookrightarrow\operatorname{Aut}(W_{12})
 }
\]
tel que, pour tout \(h\in H_*\),
\[
 \boxed{
 U_W\,\rho_{V_{11}}(h)
 =\rho_{E_{30}}\!\bigl(\iota(h)\bigr)\,U_W.
 }
\]
Cette équation est également une condition suffisante pour l'équivariance de la seconde flèche. Aucun \(\iota\) possédant cette propriété n'a été construit dans la jauge actuelle ; l'équivariance démontrée s'arrête donc à \(P_3\) et ne se prolonge pas par simple égalité des rangs.

Avant l'utilisation de \(B_K\), de \(u_K\), de l'orientation et du facteur polaire, il existait \(10\cdot6\cdot4=240\) applications étiquetées, ou \(40\) après identification des réétiquetages internes du secteur. Les règles précédentes réduisent cet ensemble à un seul élément. Tel est le contenu vérifiable du mot \emph{canonique} pour \(W_{GK}\), dans les règles déclarées. Cela n'étend pas la canonicité à une action commune de \(H_*\) sur le module de Witt.

La règle variationnelle \(\mathcal E_B\) n'utilise que des données HMT antérieures
et fixe la classe dans laquelle l'unicité est obtenue. Elle n'identifie pas l'involution
d'étoile à \(P_3\), car leurs spectres restent distincts.

\section{Comparaison des lectures de \texorpdfstring{\(t=7\) et \(t=8\)}{t=7 et t=8}}

En \(t=8\), il y a trois racines locales et onze extensions à l'horizon neuf de
l'unique racine qui se poursuit. Les tailles des trois fibres sont
\[
 (11,0,0).
\]
La matrice naturelle branche--chaîne a une ligne de onze uns et deux lignes
nulles, de sorte que son rang est un. En \(t=9\), de plus, le quotient visible
passe de \(3/11\) à \(3/6\). Cette fraction n'était pas la trace d'un
projecteur stable. La correction n'élimine pas le secteur : elle déplace la lecture
vers l'objet homogène \(P_G\) de \(t=7\), où le rang et la dimension
appartiennent au même espace.

\section{Conséquence pour une action positive}

La loi d'action peut maintenant utiliser une unique classe isométrique de projecteurs
positifs dans trois espaces de support :
\[
 P_G
 \xleftrightarrow{\ W_{GK}\ }
 P_3
 \xleftrightarrow{\ U_W\ }
 Q_3.
\]
L'étoile \(R_{B_K}\) reste réservée à la parité et à l'orientation. Le projecteur \(P_G\) enregistre la survie ; \(P_3\) enregistre la polarisation dodécaphasique ; \(Q_3\) enregistre sa réalisation d'incidence dans Witt. Cette séparation des fonctions constitue la prémisse mathématique de la mesure et de la loi positive de la section suivante. La chaîne exprime une équivalence isométrique de projecteurs, non une représentation commune : cette dernière exigerait le monomorphisme \(\iota\) et l'identité d'équivariance formulés plus haut.
 
\chapter{Du saut spectral au réseau de Leech}
\chaptermark{Publication exceptionnelle du continu}
\label{chap:83-20}

\section{Du saut spectral \texorpdfstring{\(4\to2\to6\to54\)}{4 à 2 à 6 à 54} : couplage local, compression modulaire et intégration bidimensionnelle}
\label{sec:sucesor83-salto-espectral-4-2-6-54}

Pour deux représentants consécutifs \((k,k+1)\), considérons
\[
B_k=
\begin{pmatrix}
k^2&k(k+1)\\
k(k+1)&(k+1)^2
\end{pmatrix}\pmod 9.
\]

\begin{theorem}[Unicité du bloc diagonal adjacent]
\label{thm:sucesor83-app-bloque-central}
Il existe un unique bloc \(B_k\) dont les entrées diagonales coïncident. Il est
déterminé par \(k=4\) et vaut
\[
B_c=\begin{pmatrix}7&2\\2&7\end{pmatrix}.
\]
\end{theorem}

\begin{proof}
L'égalité des diagonales équivaut à
\[
k^2\equiv(k+1)^2\pmod 9,
\]
c'est-à-dire \(2k+1\equiv0\pmod 9\). Comme \(2\) est une unité de
\(\mathbb Z/9\mathbb Z\), la congruence possède l'unique solution
\(k\equiv4\pmod 9\).
\end{proof}

Par conséquent,
\[
B_c=7I+2R,
\qquad
R=\begin{pmatrix}0&1\\1&0\end{pmatrix},
\qquad R^2=I.
\]
Avec \(P_\pm=(I\pm R)/2\),
\[
B_c=9P_++5P_-.
\]

\begin{theorem}[Hiérarchie de la séparation entre somme et produit]
\label{thm:sucesor83-app-jerarquia-4-2-6-54}
Le bloc central a pour saut spectral local \(9-5=4\). Son coefficient
hors diagonale est \(2\). La compression par les \(3\) classes modulo
\(3\) produit la différence agrégée \(51-45=6\), et le déploiement sur les
\(9\) positions produit le défaut global
\[
9\cdot6=459-405=54.
\]
Les entiers \(4,2,6,54\) appartiennent, respectivement, au saut spectral,
au couplage local, à la compression modulaire et à l'intégration
bidimensionnelle.
\end{theorem}

\begin{proof}
La décomposition \(B_c=7I+2R=9P_++5P_-\) donne le saut et le coefficient
local. Les matrices agrégées \(M_\Pi,M_\Sigma\) satisfont
\(\sum M_\Pi=51\) et \(\sum M_\Sigma=45\). Chaque entrée agrégée représente
\(9\) positions, si bien que le déploiement multiplie la différence par
\(9\) et récupère les totaux \(459\) et \(405\).
\end{proof}

Leurs concaténations ordonnées produisent les blocs \(72\) et \(27\) :
\[
72+27=99,
\qquad
72-27=45=\det B_c.
\]
La concaténation des entrées diagonales et antidiagonales produit,
respectivement, \(77\) et \(22\). Par conséquent,
\[
72+27=77+22=99,
\qquad
77-22=55=54+1.
\]
Ces égalités expriment deux partitions exactes du même bloc ; elles
n'identifient pas à elles seules une constante externe.
 
% Unidad U018. Paley, extension cuadratica finita y codigo ternario.
% Redaccion autonoma desde la autoridad material 1063.

\section{Forme de conférence de Paley, extension quadratique finie et code ternaire autodual}
\label{p12-83-c20-s1:chap:paley-extension-cuadratica-codigo-ternario}

La tétrade \(B_K\) construite dans l'unité précédente ne possède pas encore
d'étoile de Witt : pour la définir, il faut d'abord construire le système
d'incidence qui contient ses hexades. Cette unité réalise cette étape à partir de la
jauge projective des six unités modulo neuf, obtient une racine
quadratique interne de \(-I_6\) et construit le code ternaire en graphe de
longueur douze. Aucune identification ne se déduit d'une simple égalité de
cardinaux.

\subsection{Jauge projective des \(6\) unités modulo \(9\)}

Soit
\[
 \mathcal U_9=(\mathbb Z/9\mathbb Z)^\times
 =\{1,2,4,8,7,5\},
\]
ordonné par les puissances successives de \(2\) modulo \(9\). Nous fixons la
bijection de coordonnées
\begin{equation}
 \theta:\mathcal U_9\longrightarrow\mathbb P^1(\mathbb F_5),
 \qquad
 \begin{array}{c|cccccc}
 u&1&2&4&8&7&5\\ \hline
 \theta(u)&\infty&0&1&2&3&4
 \end{array}.
 \label{p12-83-c20-s1:eq:calibre-unidades-paley}
\end{equation}
L'application \(\theta\) est une jauge de coordonnées. Il n'est pas affirmé
qu'elle soit un homomorphisme entre la multiplication de \(\mathcal U_9\) et une
opération sur \(\mathbb P^1(\mathbb F_5)\).

Soit
\[
 \chi:\mathbb F_5\longrightarrow\{-1,0,1\}
\]
le caractère quadratique :
\[
 \chi(0)=0,
 \qquad
 \chi(1)=\chi(4)=1,
 \qquad
 \chi(2)=\chi(3)=-1.
\]
Dans l'ordre \((\infty,0,1,2,3,4)\), nous définissons
\begin{align}
 (C_P)_{\infty,x}&=(C_P)_{x,\infty}=1,
 &&x\in\mathbb F_5,\label{p12-83-c20-s1:eq:paley-fila-infinito}\\
 (C_P)_{x,y}&=\chi(x-y),
 &&x,y\in\mathbb F_5, x\neq y,\label{p12-83-c20-s1:eq:paley-entradas-finitas}\\
 (C_P)_{x,x}&=0.
 \label{p12-83-c20-s1:eq:paley-diagonal}
\end{align}
Ainsi,
\begin{equation}
 C_P=
 \begin{pmatrix}
 0& 1& 1& 1& 1& 1\\
 1& 0& 1&-1&-1& 1\\
 1& 1& 0& 1&-1&-1\\
 1&-1& 1& 0& 1&-1\\
 1&-1&-1& 1& 0& 1\\
 1& 1&-1&-1& 1& 0
 \end{pmatrix}.
 \label{p12-83-c20-s1:eq:matriz-conferencia-paley-autonoma}
\end{equation}

\begin{lemma}[Somme de corrélation quadratique]
\label{p12-83-c20-s1:lem:correlacion-cuadratica-f5}
Pour \(a\neq b\) dans \(\mathbb F_5\),
\begin{equation}
 \sum_{t\in\mathbb F_5}
 \chi(t-a)\chi(t-b)=-1.
 \label{p12-83-c20-s1:eq:correlacion-cuadratica-f5}
\end{equation}
\end{lemma}

\begin{proof}
La substitution \(u=(t-a)/(b-a)\) est une bijection de \(\mathbb F_5\).
Puisque \(\chi((b-a)^2)=1\), la somme se réduit à
\[
 \sum_{u\in\mathbb F_5}\chi(u(u-1)).
\]
L'évaluation de \(u=0,1,2,3,4\) donne respectivement
\(0,0,1,-1,-1\), dont la somme est \(-1\).
\end{proof}

\begin{proposition}[Identité de conférence]
\label{p12-83-c20-s1:prop:paley-identidad-conferencia-autonoma}
La matrice \(C_P\) est symétrique et satisfait
\begin{equation}
 \boxed{C_PC_P^{\mathsf T}=5I_6.}
 \label{p12-83-c20-s1:eq:paley-identidad-conferencia-autonoma}
\end{equation}
\end{proposition}

\begin{proof}
Chaque ligne contient cinq entrées de module un et une entrée nulle ; le carré de sa
norme est cinq. Pour deux lignes finies distinctes, la contribution
de la coordonnée \(\infty\) est \(+1\), tandis que
\cref{p12-83-c20-s1:lem:correlacion-cuadratica-f5} donne \(-1\) pour les cinq coordonnées
finies. Leur produit scalaire est nul. La somme d'un caractère non trivial
sur \(\mathbb F_5\) est nulle, de sorte que la ligne correspondant à
\(\infty\) est également orthogonale à chaque ligne finie.
\end{proof}

\subsection{Réduction ternaire et racine quadratique interne de \texorpdfstring{\(-I_6\)}{-I₆}}

Nous réduisons \(C_P\) modulo trois, avec \(-1\equiv2\pmod3\) :
\begin{equation}
 A_W=
 \begin{pmatrix}
 0&1&1&1&1&1\\
 1&0&1&2&2&1\\
 1&1&0&1&2&2\\
 1&2&1&0&1&2\\
 1&2&2&1&0&1\\
 1&1&2&2&1&0
 \end{pmatrix}
 \in M_6(\mathbb F_3).
 \label{p12-83-c20-s1:eq:matriz-paley-witt-ternaria}
\end{equation}

\begin{theorem}[Racine quadratique interne de \(-I_6\)]
\label{p12-83-c20-s1:thm:aw-cuadrado-menos-identidad-autonomo}
On a
\begin{equation}
 \boxed{A_W^2=-I_6.}
 \label{p12-83-c20-s1:eq:aw-cuadrado-menos-identidad-autonomo}
\end{equation}
En particulier,
\[
 A_W^{-1}=-A_W.
\]
\end{theorem}

\begin{proof}
La \cref{p12-83-c20-s1:prop:paley-identidad-conferencia-autonoma} donne
\[
 A_WA_W^{\mathsf T}
 \equiv5I_6
 \equiv2I_6
 =-I_6
 \pmod3.
\]
Puisque \(A_W=A_W^{\mathsf T}\), on obtient \(A_W^2=-I_6\).
\end{proof}

L'identité précédente est l'espace de résidence algébrique de la racine de \(-1\)
dans ce corridor. Elle ne provient pas du facteur \(1/4\) de l'inversion de Hadamard :
celui-ci est un scalaire rationnel imposé par \(H_4^2=4I_4\) ; celle-là est une
structure quadratique en caractéristique trois.

\subsection{Extension quadratique à \texorpdfstring{\(\mathbb F_9\)}{F₉}}

On définit
\begin{equation}
 \mathbb F_9=\mathbb F_3[\iota]/(\iota^2+1),
 \qquad \iota^2=-1.
 \label{p12-83-c20-s1:eq:definicion-f9-autonoma}
\end{equation}
Le polynôme \(X^2+1\) n'a pas de racines dans \(\mathbb F_3\), de sorte que le
quotient est un corps à neuf éléments. Pour
\(V=\mathbb F_3^6\), soit
\[
 V_9=V\otimes_{\mathbb F_3}\mathbb F_9.
\]
Les opérateurs
\begin{equation}
 P_\iota=\frac12(I-\iota A_W),
 \qquad
 P_{-\iota}=\frac12(I+\iota A_W)
 \label{p12-83-c20-s1:eq:proyectores-f9-paley}
\end{equation}
sont bien définis puisque \(2\) est inversible dans \(\mathbb F_3\).

\begin{proposition}[Projecteurs conjugués]
\label{p12-83-c20-s1:prop:proyectores-conjugados-paley}
On a
\[
 P_\iota^2=P_\iota,
 \qquad
 P_{-\iota}^2=P_{-\iota},
 \qquad
 P_\iota P_{-\iota}=0,
 \qquad
 P_\iota+P_{-\iota}=I.
\]
\end{proposition}

\begin{proof}
En utilisant \(\iota^2=-1\) et \(A_W^2=-I\),
\[
 (\iota A_W)^2=I.
\]
Les quatre identités sont alors les identités élémentaires des
projecteurs \((I\mp\iota A_W)/2\) d'une involution.
\end{proof}

\begin{theorem}[Décomposition spectrale conjuguée]
\label{p12-83-c20-s1:thm:descomposicion-espectral-f9-paley}
Soit
\[
 E_{\pm\iota}=\ker(A_W\mp\iota I).
\]
Alors
\begin{equation}
 V_9=E_\iota\oplus E_{-\iota},
 \qquad
 \dim_{\mathbb F_9}E_\iota
 =\dim_{\mathbb F_9}E_{-\iota}=3.
 \label{p12-83-c20-s1:eq:descomposicion-espectral-f9-paley}
\end{equation}
\end{theorem}

\begin{proof}
Le polynôme minimal de \(A_W\) divise \(X^2+1\), qui se scinde sur
\(\mathbb F_9\) avec des racines simples \(\pm\iota\). Les projecteurs de
\cref{p12-83-c20-s1:prop:proyectores-conjugados-paley} produisent la somme directe. Puisque
\(A_W\) a ses coefficients dans \(\mathbb F_3\), l'automorphisme de
Frobenius échange les deux espaces propres ; leurs dimensions sont égales et
leur somme vaut six.
\end{proof}

Cette construction dote \(\mathbb F_3^6\) d'une structure de module de rang
trois sur \(\mathbb F_9\), au moyen de
\begin{equation}
 (a+b\iota)\cdot v:=av+b(vA_W).
 \label{p12-83-c20-s1:eq:accion-f9-sobre-f3-seis}
\end{equation}
L'égalité \(3^6=9^3\) est ainsi accompagnée d'une action explicite ; elle n'est pas
utilisée comme argument cardinal.

\subsection{Code ternaire en graphe de longueur \(12\)}

On définit
\begin{equation}
 c:\mathbb F_3^6\longrightarrow\mathbb F_3^{12},
 \qquad
 c(q)=(q,qA_W),
 \label{p12-83-c20-s1:eq:aplicacion-codigo-grafico-witt}
\end{equation}
et
\begin{equation}
 C_W:=\operatorname{im}c.
 \label{p12-83-c20-s1:eq:definicion-codigo-witt-ternario}
\end{equation}
Sa matrice génératrice est
\[
 G_W=(I_6\mid A_W).
\]

\begin{theorem}[Injectivité et autodualité]
\label{p12-83-c20-s1:thm:cw-inyectivo-autodual}
L'application \(c\) est injective, \(\dim C_W=6\) et
\begin{equation}
 \boxed{C_W=C_W^\perp.}
 \label{p12-83-c20-s1:eq:cw-autodual}
\end{equation}
\end{theorem}

\begin{proof}
La première moitié de \(c(q)\) est \(q\), de sorte que
\(\operatorname{pr}_1\circ c=\operatorname{id}\) ; cela prouve
l'injectivité et la dimension. En outre,
\[
 G_WG_W^{\mathsf T}
 =I_6+A_WA_W^{\mathsf T}
 =I_6+2I_6=0
\]
dans \(\mathbb F_3\). Par conséquent, \(C_W\subseteq C_W^\perp\). Les deux espaces
sont de dimension six, donc ils sont égaux.
\end{proof}

Pour \(r=0,\ldots,12\), soit
\[
 N_r:=\#\{q\in\mathbb F_3^6:
            \operatorname{wt}(q,qA_W)=r\}.
\]
L'énumération des \(3^6=729\) entrées produit
\begin{equation}
 N_0=1,
 \qquad
 N_6=264,
 \qquad
 N_9=440,
 \qquad
 N_{12}=24,
 \label{p12-83-c20-s1:eq:censo-pesos-cw}
\end{equation}
et \(N_r=0\) pour les autres poids.

\begin{theorem}[Énumérateur et distance minimale]
\label{p12-83-c20-s1:thm:enumerador-cw-autonomo}
Le code \(C_W\) possède les paramètres
\[
 \boxed{C_W=[12,6,6]_3}
\]
et l'énumérateur des poids
\begin{equation}
 \boxed{W_{C_W}(y)=1+264y^6+440y^9+24y^{12}.}
 \label{p12-83-c20-s1:eq:enumerador-pesos-cw}
\end{equation}
\end{theorem}

\begin{proof}
La dimension et l'autodualité ont déjà été démontrées. Le recensement
\eqref{p12-83-c20-s1:eq:censo-pesos-cw} parcourt les \(729\) mots sans échantillonnage. Son
plus petit poids non nul est six, qui est la distance minimale d'un code linéaire.
La somme \(1+264+440+24=729\) contrôle qu'aucun mot n'a été omis.
\end{proof}

\subsection{Nature de la structure obtenue}

Trois objets distincts et trois flèches typées ont été construits :
\[
 \mathcal U_9
 \xrightarrow{\theta}\mathbb P^1(\mathbb F_5)
 \xrightarrow{\chi}C_P
 \xrightarrow{\bmod3}A_W,
\]
\[
 A_W^2=-I_6
 \quad\Longrightarrow\quad
 \mathbb F_3^6\cong\mathbb F_9^3
 \quad\Longrightarrow\quad
 C_W=\{(q,qA_W):q\in\mathbb F_3^6\}.
\]
La première flèche est une jauge, la deuxième utilise le caractère quadratique et la
troisième est une réduction des coefficients. La structure sur
\(\mathbb F_9\) est finie ; elle ne s'identifie pas à une structure complexe
réelle ou analytique.

\subsection{Obstructions du corridor Paley–\(C_W\) : bijection, \(A_W^2=-I_6\), autodualité et recensement de \(729\) mots}

\begin{criterion}[Réfutation du corridor Paley--\(C_W\)]
La construction est réfutée si l'un des faits
suivants se produit :
\begin{enumerate}
 \item \(\theta\) n'est pas bijective ou est utilisée comme homomorphisme sans preuve ;
 \item la somme \eqref{p12-83-c20-s1:eq:correlacion-cuadratica-f5} diffère de \(-1\) ;
 \item \(C_PC_P^{\mathsf T}\neq5I_6\);
 \item \(A_W^2\neq-I_6\);
 \item les projecteurs \eqref{p12-83-c20-s1:eq:proyectores-f9-paley} ne sont pas
       complémentaires ;
 \item l'application \(q\mapsto(q,qA_W)\) perd l'injectivité ou
       l'autodualité ;
 \item la somme du recensement ne vaut pas \(729\), diffère de
       \eqref{p12-83-c20-s1:eq:censo-pesos-cw}, ou présente un mot de poids compris entre un
       et cinq ;
 \item on infère \(\mathbb F_3^6\cong\mathbb F_9^3\) uniquement de
       \(3^6=9^3\), en omettant l'action \eqref{p12-83-c20-s1:eq:accion-f9-sobre-f3-seis}.
\end{enumerate}
\end{criterion}


% Unidad U019. Sistema de Witt, campo de estrellas y grupo M12.
% Redaccion autonoma desde la autoridad material 1063.

\section{Système de Steiner ternaire, involutions d'étoile de Witt et génération de \texorpdfstring{\(M_{12}\)}{M12}}
\label{p12-83-c20-s2:chap:sistema-witt-estrellas-m12}

L'unité précédente a construit le code autodual
\(C_W=[12,6,6]_3\). On projectivise maintenant ses mots de poids six,
on prouve la propriété de Steiner, et c'est alors seulement que l'on définit l'étoile
associée à la tétrade \(B_K\). Cet ordre empêche de transformer une figure
d'incidence en un objet ajouté par ressemblance visuelle.

\subsection{Supports projectifs de poids \(6\)}

Soit
\begin{equation}
 \mathcal H_W
 :=\{\operatorname{supp}(c):c\in C_W,
                              \operatorname{wt}(c)=6\}.
 \label{p12-83-c20-s2:eq:familia-soportes-peso-seis}
\end{equation}
Chaque support apparaît pour la paire de mots \(\{c,-c\}\). En effet, si deux
mots de poids six ayant le même support n'étaient pas proportionnels, on pourrait
multiplier l'un d'eux par un scalaire pour annuler une coordonnée en les soustrayant ; on obtiendrait
un mot non nul de \(C_W\) de poids inférieur à six, ce qui contredit la distance
minimale. Puisqu'il n'existe pas non plus de mot non nul égal à son opposé en
caractéristique trois, chaque orbite projective possède exactement deux éléments.
Par conséquent,
\begin{equation}
 |\mathcal H_W|=\frac{264}{2}=132.
 \label{p12-83-c20-s2:eq:numero-hexadas-cw}
\end{equation}

Pour \(T\in\binom{[12]}5\), nous définissons
\begin{equation}
 \nu(T):=\#\{H\in\mathcal H_W:T\subseteq H\}.
 \label{p12-83-c20-s2:eq:incidencia-cinco-hexada}
\end{equation}
L'énumération exacte des
\[
 \binom{12}{5}=792
\]
sous-ensembles de cinq points produit
\begin{equation}
 \nu(T)=1
 \qquad
 \text{pour tout }T\in\binom{[12]}5.
 \label{p12-83-c20-s2:eq:unicidad-cinco-en-hexada}
\end{equation}

\begin{theorem}[Système de Witt ternaire]
\label{p12-83-c20-s2:thm:w12-autonomo}
Le système d'incidence
\[
 W_{12}:=([12],\mathcal H_W)
\]
est le système de Steiner
\begin{equation}
 \boxed{W_{12}=S(5,6,12).}
 \label{p12-83-c20-s2:eq:w12-steiner-autonomo}
\end{equation}
\end{theorem}

\begin{proof}
L'égalité \eqref{p12-83-c20-s2:eq:unicidad-cinco-en-hexada} est précisément la
propriété de Steiner. Le nombre de blocs se retrouve également par le double
comptage des couples \((T,H)\) avec \(T\subset H\) :
\[
 |\mathcal H_W|\binom65=\binom{12}{5},
 \qquad
 |\mathcal H_W|=\frac{792}{6}=132.
\]
\end{proof}

\subsection{Résidu d'une tétrade et décomposition en \(4\) composantes incidentes}

Soit \(B\in\binom{[12]}4\). Dans un système \(S(5,6,12)\), le nombre
d'hexades qui contiennent \(B\) est
\begin{equation}
 \lambda_4
 =\frac{\binom{12-4}{5-4}}{\binom{6-4}{5-4}}
 =\frac82=4.
 \label{p12-83-c20-s2:eq:lambda-cuatro-w12}
\end{equation}
Nous écrivons ces quatre hexades sous la forme
\[
 H_i=B\sqcup P_i,
 \qquad |P_i|=2,
 \qquad i=1,\ldots,4.
\]

\begin{lemma}[Partition résiduelle d'une tétrade]
\label{p12-83-c20-s2:lem:particion-residual-tetrada-witt}
Les quatre paires \(P_1,\ldots,P_4\) sont disjointes et
\begin{equation}
 [12]\setminus B
 =P_1\sqcup P_2\sqcup P_3\sqcup P_4.
 \label{p12-83-c20-s2:eq:particion-residual-tetrada-witt}
\end{equation}
\end{lemma}

\begin{proof}
Si deux paires résiduelles partageaient un point \(p\), les hexades
correspondantes contiendraient le même sous-ensemble \(B\cup\{p\}\) de cinq points,
en contradiction avec l'unicité de Steiner. Les quatre paires sont donc
disjointes ; elles contiennent huit points et épuisent le complémentaire de \(B\).
\end{proof}

\begin{definition}[Étoile et involution d'étoile de Witt]
\label{p12-83-c20-s2:def:estrella-witt-involucion}
L'étoile centrée en \(B\) est
\begin{equation}
 \mathscr S_B
 :=\{H\in\mathcal H_W:B\subseteq H\}
 =\{B\sqcup P_1,\ldots,B\sqcup P_4\}.
 \label{p12-83-c20-s2:eq:estrella-witt-general}
\end{equation}
Si \(P_i=\{p_i^-,p_i^+\}\), son involution est
\begin{equation}
 \sigma_B:=\prod_{i=1}^{4}(p_i^-\ p_i^+).
 \label{p12-83-c20-s2:eq:involucion-estrella-witt-general}
\end{equation}
La permutation \(\sigma_B\) fixe point par point \(B\) et échange les
extrémités de chaque pétale résiduel.
\end{definition}

Le terme mathématique employé désormais est \emph{étoile de Witt} ou
\emph{involution d'étoile de Witt}. Il ne s'agit pas d'une étoile
euclidienne.

\subsection{Étoile sélectionnée par la tétrade du sceau}

L'unité dodécaphasique a produit, sans utiliser le plan de Steiner,
\[
 B_K=\{4,6,9,10\}.
\]
Maintenant que \(W_{12}\) est construit, son étoile est bien définie. Les
quatre hexades incidentes sont
\begin{align}
 H_{13}&=B_K\sqcup\{1,3\}
       =\{1,3,4,6,9,10\},
 \label{p12-83-c20-s2:eq:hexada-petalo-13}\\
 H_{2,11}&=B_K\sqcup\{2,11\}
       =\{2,4,6,9,10,11\},
 \label{p12-83-c20-s2:eq:hexada-petalo-211}\\
 H_{5,7}&=B_K\sqcup\{5,7\}
       =\{4,5,6,7,9,10\}
       =H_\alpha^-,
 \label{p12-83-c20-s2:eq:hexada-petalo-57}\\
 H_{8,12}&=B_K\sqcup\{8,12\}
       =\{4,6,8,9,10,12\}.
 \label{p12-83-c20-s2:eq:hexada-petalo-812}
\end{align}
En conséquence,
\begin{equation}
 \boxed{
 \sigma_{B_K}=(1\ 3)(2\ 11)(5\ 7)(8\ 12).}
 \label{p12-83-c20-s2:eq:involucion-estrella-witt-bk}
\end{equation}

\begin{figure}[htbp]
\centering
\begin{tikzpicture}[
  >=Stealth,
  font=\small,
  core/.style={draw,double,rounded corners=2mm,align=center,
    minimum width=4.0cm,minimum height=1.3cm,fill=black!4},
  petal/.style={draw,rounded corners=2mm,align=center,
    minimum width=3.55cm,minimum height=1.08cm,fill=black!2},
  negative/.style={petal,line width=1.1pt,double,fill=black!10},
  arr/.style={->,thick}
]
\node[core] (B) at (0,0)
 {\(\displaystyle B_K=\{4,6,9,10\}\)\\
  points fixés par \(\sigma_{B_K}\)};
\node[petal] (P13) at (0,3.0)
 {\(P_1=\{1,3\}\), \((1\ 3)\)\\
  \(H_{13}=B_K\sqcup P_1\)};
\node[petal] (P211) at (5.25,0)
 {\(P_2=\{2,11\}\), \((2\ 11)\)\\
  \(H_{2,11}=B_K\sqcup P_2\)};
\node[negative] (P57) at (0,-3.0)
 {\(P_3=\{5,7\}\), \((5\ 7)\)\\
  \(H_{5,7}=H_\alpha^-\)};
\node[petal] (P812) at (-5.25,0)
 {\(P_4=\{8,12\}\), \((8\ 12)\)\\
  \(H_{8,12}=B_K\sqcup P_4\)};
\draw[arr] (B)--(P13);
\draw[arr] (B)--(P211);
\draw[arr] (B)--(P57);
\draw[arr] (B)--(P812);
\node[align=center,font=\footnotesize] at (0,-4.35)
 {\(\mathscr S_{B_K}
   =\{H_{13},H_{2,11},H_{5,7},H_{8,12}\}\).\\
  Le diagramme représente l'incidence, non la distance géométrique.};
\end{tikzpicture}
\caption{Étoile de Witt centrée en \(B_K\). Le centre est la tétrade
fixe ; les quatre pétales sont les paires résiduelles qui déterminent
l'involution.}
\label{p12-83-c20-s2:fig:estrella-witt-bk-autonoma}
\end{figure}

\subsection{Champ global des 495 étoiles}

Il existe une étoile pour chaque tétrade, de sorte que
\begin{equation}
 \#\{\mathscr S_B:B\in\tbinom{[12]}4\}
 =\binom{12}{4}=495.
 \label{p12-83-c20-s2:eq:numero-estrellas-witt}
\end{equation}
Le recensement exact vérifie
\begin{equation}
 \sigma_B(\mathcal H_W)=\mathcal H_W
 \qquad
 \text{pour toute }B\in\binom{[12]}4.
 \label{p12-83-c20-s2:eq:estrellas-conservan-hexadas}
\end{equation}
On définit
\begin{equation}
 \Gamma_\star
 :=\langle\sigma_B:B\in\tbinom{[12]}4\rangle
 \leq\operatorname{Aut}(W_{12}).
 \label{p12-83-c20-s2:eq:grupo-generado-estrellas-witt}
\end{equation}

Un ensemble générateur de six tétrades est
\[
 1234,
 \quad1235,
 \quad1236,
 \quad1245,
 \quad1345,
 \quad2345.
\]
Les ordres des sous-groupes obtenus en ajoutant successivement leurs
involutions sont
\begin{equation}
 1, 2, 6, 18, 360, 7920, 95040.
 \label{p12-83-c20-s2:eq:cadena-ordenes-generadores-m12}
\end{equation}

\begin{theorem}[Génération du groupe de Mathieu \(M_{12}\)]
\label{p12-83-c20-s2:thm:495-estrellas-m12-autonomo}
On a
\begin{equation}
 \boxed{
 \langle\sigma_B:B\in\tbinom{[12]}4\rangle
 =\operatorname{Aut}(W_{12})
 \cong M_{12}.}
 \label{p12-83-c20-s2:eq:estrellas-generan-m12}
\end{equation}
\end{theorem}

\begin{proof}
Chaque \(\sigma_B\) conserve les \(132\) hexades, donc
\(\Gamma_\star\leq\operatorname{Aut}(W_{12})\). Le recensement de Schreier des
six générateurs indiqués produit l'ordre \(95040\). Le groupe
d'automorphismes de \(S(5,6,12)\) est \(M_{12}\) et possède ce même ordre ;
l'inclusion est donc une égalité.
\end{proof}

Une étoile isolée engendre uniquement
\[
 \langle\sigma_B\rangle\cong C_2.
\]
La conclusion concernant \(M_{12}\) appartient au champ global des
\(495\) involutions, non à \(\sigma_{B_K}\) considérée isolément.

\subsection{Action sur le module d'augmentation}

Soit
\[
 V_{11}=\mathbf1^\perp\subset\mathbb R^{12}.
\]
L'involution \(\sigma_B\) fixe les quatre coordonnées de \(B\) et agit
par quatre transpositions sur son complémentaire. Dans \(\mathbb R^{12}\),
elle possède huit directions de valeur propre \(+1\) et quatre de valeur propre \(-1\).
La direction constante appartient au premier espace propre. Par conséquent,
\begin{equation}
 \operatorname{Spec}(\sigma_B|_{V_{11}})
 =\{1^{[7]},(-1)^{[4]}\},
 \qquad
 \frac1{11}\operatorname{Tr}(\sigma_B|_{V_{11}})=\frac3{11}.
 \label{p12-83-c20-s2:eq:espectro-involucion-estrella-aumentacion}
\end{equation}
C'est une involution signée ; ce n'est pas un projecteur positif de rang trois.

\subsection{Obstructions de \(W_{12}\), du champ de \(495\) involutions d'étoile et de la génération de \(M_{12}\)}

\begin{criterion}[Réfutation de \(W_{12}\), de ses étoiles et de \(M_{12}\)]
La construction est réfutée si l'un des faits
suivants se produit :
\begin{enumerate}
 \item la projectivisation des \(264\) mots de poids six ne produit pas
       \(132\) supports ;
 \item il existe un sous-ensemble de cinq points contenu dans zéro hexade ou dans plus
       d'une hexade ;
 \item une tétrade n'appartient pas exactement à quatre hexades ;
 \item les quatre paires résiduelles ne partitionnent pas le complémentaire de la
       tétrade ;
 \item les pétales de \(B_K\) ne sont pas
       \(\{1,3\},\{2,11\},\{5,7\},\{8,12\}\) ;
 \item l'une des \(495\) involutions ne conserve pas
       \(\mathcal H_W\) ;
 \item le groupe engendré n'a pas pour ordre \(95040\) ;
 \item on attribue \(M_{12}\) à une unique étoile au lieu du champ
       global des involutions ;
 \item l'étoile est introduite avant de construire \(W_{12}\).
\end{enumerate}
\end{criterion}


% Unidad U020. Prolongacion binaria W12--W24 y M24.
% Redaccion autonoma desde la autoridad material 1063.

\section{Extension binaire relative à une dodécade, système \texorpdfstring{\(W_{24}\)}{W24} et groupe \texorpdfstring{\(M_{24}\)}{M24}}
\label{p12-83-c20-s3:chap:dodecada-w24-m24}

Le système \(W_{12}\) construit sur douze phases admet une prolongation
binaire lorsqu'on déclare un code de Golay étendu, une dodécade et un
isomorphisme d'incidence. Ces données sont des entrées relatives de cette branche :
elles ne se déduisent pas uniquement de \(C_W\) et n'interviennent pas dans la branche de réseaux
ternaire qui conduira au réseau de Leech.

\subsection{Code binaire de Golay et octades}

Soit
\[
 \mathcal G_{24}\subset\mathbb F_2^{24}
\]
le code binaire étendu de Golay. Ses poids sont
\begin{equation}
 0, 8, 12, 16, 24.
 \label{p12-83-c20-s3:eq:pesos-golay-binario-extendido}
\end{equation}
Les supports des mots de poids huit forment le système de Steiner
\begin{equation}
 W_{24}=S(5,8,24).
 \label{p12-83-c20-s3:eq:w24-steiner-octadas}
\end{equation}
Ses blocs sont appelés \emph{octades}. Fixons une dodécade \(D\), c'est-à-
dire le support d'un mot de poids douze, et soit \(\mathcal O_{24}\)
la famille des octades.

On définit
\begin{equation}
 \mathcal H(D)
 :=\{O\cap D:O\in\mathcal O_{24},\ |O\cap D|=6\}.
 \label{p12-83-c20-s3:eq:hexadas-residuales-dodecada}
\end{equation}

\begin{theorem}[Plan résiduel d'une dodécade]
\label{p12-83-c20-s3:thm:diseno-residual-dodecada-autonomo}
Le système résiduel est
\begin{equation}
 \boxed{(D,\mathcal H(D))\cong S(5,6,12).}
 \label{p12-83-c20-s3:eq:diseno-residual-dodecada-w12}
\end{equation}
En outre, chaque \(H\in\mathcal H(D)\) possède une unique octade \(O_H\) telle que
\begin{equation}
 O_H\cap D=H.
 \label{p12-83-c20-s3:eq:elevacion-unica-hexada-octada}
\end{equation}
\end{theorem}

\begin{proof}
Soit \(A\subset D\), avec \(|A|=5\). Dans \(W_{24}=S(5,8,24)\), il existe une
unique octade \(O_A\) qui contient \(A\). Si \(j=|O_A\cap D|\), la somme
binaire des mots caractéristiques de \(O_A\) et \(D\) appartient à
\(\mathcal G_{24}\) et a pour poids
\[
 8+12-2j=20-2j.
\]
Puisque \(A\subseteq O_A\cap D\), on a \(5\leq j\leq8\). Les quatre
poids possibles \(10,8,6,4\), combinés à
\eqref{p12-83-c20-s3:eq:pesos-golay-binario-extendido}, imposent \(j=6\). Ainsi, tout
sous-ensemble de cinq points de \(D\) appartient à une unique hexade résiduelle.

Si deux octades avaient la même intersection hexadique avec \(D\), toutes deux
contiendraient les mêmes cinq points de cette hexade, ce qui contredirait
l'unicité dans \(S(5,8,24)\). Cela prouve également l'unicité du
relèvement.
\end{proof}

\subsection{Jauge relative entre le plan HMT et la dodécade}

Pour relier la construction précédente au plan obtenu à partir de
\(C_W\), il faut déclarer explicitement
\begin{equation}
 (D,\ell),
 \qquad
 \ell:W_{12}^{\rm HMT}
 \xrightarrow{\sim}(D,\mathcal H(D)),
 \label{p12-83-c20-s3:eq:dato-relativo-dodecada-alineamiento}
\end{equation}
où \(\ell\) est un isomorphisme d'incidence. Si \(\ell_*\) désigne son
action sur les blocs, nous définissons
\begin{equation}
 \widehat\iota_{D,\ell}
 :=\iota_D\circ\ell_*:
 W_{12}^{\rm HMT}\longrightarrow W_{24},
 \qquad
 \widehat\iota_{D,\ell}(H)=O_{\ell(H)}.
 \label{p12-83-c20-s3:eq:elevacion-tipificada-w12-w24}
\end{equation}
La composition \(\widehat\iota_{D,\ell}\) évite d'appliquer
\(\iota_D\), dont le domaine est \(\mathcal H(D)\), directement à une hexade
qui vit encore dans \([12]_{\rm HMT}\).

\begin{proposition}[Injectivité du relèvement relatif]
\label{p12-83-c20-s3:prop:inyectividad-elevacion-w12-w24}
L'application \(\widehat\iota_{D,\ell}\) est injective et conserve
l'incidence entre points et hexades après application de \(\ell\).
\end{proposition}

\begin{proof}
L'isomorphisme \(\ell\) est bijectif sur les blocs. D'après le
\cref{p12-83-c20-s3:thm:diseno-residual-dodecada-autonomo}, chaque hexade résiduelle possède une
unique octade de relèvement. Deux hexades HMT ayant la même image auraient la même
hexade résiduelle et, par bijectivité de \(\ell\), seraient égales.
\end{proof}

Il n'existe pas de choix absolu de \((D,\ell)\) antérieur à cette jauge.
L'objet invariant est la classe de la construction sous changement simultané de
dodécade et d'alignement.

\subsection{Loi d'incidence \texorpdfstring{\(4+1\)}{4+1} sur une tétrade}

Pour une tétrade \(B\subset D\), les paramètres dérivés des deux
systèmes de Steiner sont
\begin{equation}
 \lambda_4(W_{12})
 =\frac{\binom81}{\binom21}=4,
 \qquad
 \lambda_4(W_{24})
 =\frac{\binom{20}1}{\binom41}=5.
 \label{p12-83-c20-s3:eq:ley-cuatro-mas-uno-witt}
\end{equation}
Par conséquent, les cinq octades qui contiennent \(B\) se décomposent de manière
unique en
\begin{equation}
 \{O_H:H\in\mathcal H(D),\ B\subset H\}
 \sqcup\{O_B^\perp\}.
 \label{p12-83-c20-s3:eq:descomposicion-cinco-octadas}
\end{equation}
Les quatre premières satisfont \(|O_H\cap D|=6\) ; la cinquième satisfait
\begin{equation}
 O_B^\perp\cap D=B.
 \label{p12-83-c20-s3:eq:octada-transversal-tetrada}
\end{equation}

Pour \(B=B_K\), les quatre relèvements résiduels correspondent aux
pétales
\[
 \{1,3\},
 \qquad
 \{2,11\},
 \qquad
 \{5,7\},
 \qquad
 \{8,12\},
\]
et la cinquième octade est la transversale \(O_{B_K}^\perp\). Cette loi explique
la transition de quatre hexades à cinq octades sans identifier une face
d'un cube à une hexade ni un sommet à une octade.

\subsection{Correspondance des strates et liberté du triplet positif}

La structure à cinq régions du canal \(\pi\) a pour types
\[
 1_\perp+1_{--}+3_{++}.
\]
L'incidence sur \(B_K\) a pour types
\[
 1_{\rm transversal}+1_{H^-}+3_{\rm restantes}
\]
après désignation de l'hexade négative. Il existe donc une
correspondance naturelle entre \emph{strates de rôle} :
\begin{equation}
 1_\perp\longleftrightarrow O_{B_K}^\perp,
 \qquad
 1_{--}\longleftrightarrow
 \widehat\iota_{D,\ell}(H_\alpha^-),
 \qquad
 3_{++}\longleftrightarrow
 \mathcal O_{B_K}^{+,3}.
 \label{p12-83-c20-s3:eq:correspondencia-estratos-cinco-cinco}
\end{equation}
Ici, \(\mathcal O_{B_K}^{+,3}\) est l'ensemble non ordonné des trois
autres octades résiduelles. Une bijection individuelle entre les trois régions
positives et ces trois octades exige de déclarer une jauge supplémentaire de
\(S_3\). La coïncidence du motif \(1+1+3\) ne détermine pas à elle seule cet
ordonnancement.

\subsection{Groupe d'automorphismes du système d'octades}

\begin{theorem}[Groupe de Mathieu \(M_{24}\)]
\label{p12-83-c20-s3:thm:automorfismos-w24-m24}
Le groupe d'automorphismes du système de Steiner des octades est
\begin{equation}
 \boxed{
 \operatorname{Aut}(W_{24})\cong M_{24},
 \qquad
 |M_{24}|=244823040.}
 \label{p12-83-c20-s3:eq:automorfismos-w24-m24}
\end{equation}
\end{theorem}

Le théorème identifie le groupe classique d'automorphismes de la donnée binaire
d'entrée. Il n'affirme pas que \(M_{24}\) soit engendré à partir d'une seule
dodécade, ni que le code ternaire \(C_W\) détermine à lui seul
\(\mathcal G_{24}\).

\subsection{Séparation catégorique par rapport à la branche de réseaux}

La branche construite dans cette unité est
\begin{equation}
 (W_{12}^{\rm HMT};\mathcal G_{24},D,\ell)
 \longrightarrow W_{24}
 \longrightarrow\operatorname{Aut}(W_{24})\cong M_{24}.
 \label{p12-83-c20-s3:eq:rama-binaria-w24-m24}
\end{equation}
La branche de réseaux qui sera construite ensuite est
\begin{equation}
 C_W\longrightarrow N(A_2^{12})
 \xrightarrow{\text{voisin d’indice trois}}\Lambda_{24}.
 \label{p12-83-c20-s3:eq:rama-ternaria-niemeier-leech-anunciada}
\end{equation}
Leurs constructeurs, domaines et codomaines sont distincts. En particulier,
\begin{equation}
 \boxed{
 \text{aucune flèche n’est introduite }
 W_{24}\longrightarrow\Lambda_{24}.}
 \label{p12-83-c20-s3:eq:no-flecha-w24-leech}
\end{equation}

\begin{figure}[htbp]
\centering
\begin{tikzpicture}[
  >=Stealth,
  node distance=13mm and 16mm,
  box/.style={draw,rounded corners=1.5mm,align=center,
    inner xsep=3mm,inner ysep=2.2mm},
  arr/.style={->,thick}
]
\node[box] (cw) {\(C_W\)};
\node[box,right=of cw] (w12) {\(W_{12}\)};
\node[box,right=of w12] (w24) {\(W_{24}\)};
\node[box,right=of w24] (m24) {\(M_{24}\)};
\node[box,below=of cw] (cw2) {\(C_W\)};
\node[box,right=of cw2] (n) {\(N(A_2^{12})\)};
\node[box,right=of n] (leech) {\(\Lambda_{24}\)};
\draw[arr] (cw)--(w12);
\draw[arr] (w12)--node[above,font=\scriptsize]
 {\((\mathcal G_{24},D,\ell)\)} (w24);
\draw[arr] (w24)--(m24);
\draw[arr] (cw2)--node[below,font=\scriptsize]
 {pegado ternario} (n);
\draw[arr] (n)--node[below,font=\scriptsize]
 {\(3\)-voisin} (leech);
\node[below=5mm of w24,font=\footnotesize,align=center]
 {L'alignement vertical n'est pas une flèche.};
\end{tikzpicture}
\caption{Bifurcation typée des prolongations binaire et ternaire.}
\label{p12-83-c20-s3:fig:bifurcacion-w24-leech-autonoma}
\end{figure}

\subsection{Obstructions de la prolongation binaire \(W_{12}\to W_{24}\) : code de Golay, octades, alignement et séparation de la branche de Leech}

\begin{criterion}[Réfutation de la prolongation binaire]
La construction est réfutée si l'un des faits
suivants se produit :
\begin{enumerate}
 \item le spectre des poids de \(\mathcal G_{24}\) contient un poids non
       répertorié dans \eqref{p12-83-c20-s3:eq:pesos-golay-binario-extendido} ;
 \item les supports de poids huit ne forment pas \(S(5,8,24)\) ;
 \item une intersection résiduelle d'une octade avec \(D\) qui contient cinq
       points de \(D\) a un cardinal différent de six ;
 \item une hexade résiduelle admet zéro octade de relèvement ou plus d'une ;
 \item on applique \(\iota_D\) directement à une hexade HMT sans
       l'alignement \(\ell\) ;
 \item la loi d'incidence sur une tétrade n'est pas \(4+1\) ;
 \item on ordonne le triplet positif sans déclarer la jauge \(S_3\) ;
 \item \(\operatorname{Aut}(W_{24})\) n'a pas pour ordre \(244823040\) ;
 \item on affirme que \(W_{24}\) naît de \(W_{12}\) sans la donnée binaire
       \(\mathcal G_{24}\), la dodécade et l'alignement ;
 \item on sérialise les deux branches au moyen d'une flèche
       \(W_{24}\to\Lambda_{24}\).
\end{enumerate}
\end{criterion}


% Unidad U021. Pegado discriminante ternario y Niemeier A2^12.
% Redaccion autonoma desde la autoridad material 1063.

\section{Recollement discriminant ternaire et construction du réseau de Niemeier de type \texorpdfstring{\(A_2^{12}\)}{A₂¹²}}
\label{p12-83-c20-s4:chap:pegado-niemeier-a2-doce}

Cette unité commence la branche de réseaux directement à partir du code ternaire
autodual \(C_W\). Le système binaire \(W_{24}\), la dodécade et le groupe
\(M_{24}\) ne sont pas des entrées de cette construction. La succession typée est
\[
 C_W
 \longrightarrow
 (A_2^*/A_2)^{12}
 \longrightarrow
 N(C_W)
 \cong N(A_2^{12}).
\]

\subsection{Réseau de racines \texorpdfstring{\(A_2\)}{A₂}}

Soit
\[
 A_2=\mathbb Z\alpha_1\oplus\mathbb Z\alpha_2
\]
le réseau de racines dont la matrice de Gram, dans la base ordonnée
\((\alpha_1,\alpha_2)\), est
\begin{equation}
 G_{A_2}=
 \begin{pmatrix}
 2&-1\\
 -1&2
 \end{pmatrix}.
 \label{p12-83-c20-s4:eq:gram-a2-autonomo}
\end{equation}
Pour \(x=x_1\alpha_1+x_2\alpha_2\),
\begin{equation}
 \langle x,x\rangle
 =2(x_1^2-x_1x_2+x_2^2).
 \label{p12-83-c20-s4:eq:norma-a2-autonoma}
\end{equation}
En particulier, \(A_2\) est pair, défini positif et
\[
 \det A_2=3.
\]

On définit
\[
 \rho:=\alpha_1+\alpha_2.
\]
Alors
\begin{equation}
 \rho^2=2,
 \qquad
 \langle\rho,\alpha_1\rangle
 =\langle\rho,\alpha_2\rangle=1.
 \label{p12-83-c20-s4:eq:identidades-rho-a2}
\end{equation}

\subsection{Module et forme discriminants}

Le réseau dual est
\[
 A_2^*
 :=\{x\in A_2\otimes_{\mathbb Z}\mathbb Q:
       \langle x,A_2\rangle\subseteq\mathbb Z\}.
\]
Nous choisissons les représentants
\begin{equation}
 \mu_0=0,
 \qquad
 \mu_1=\frac{2\alpha_1+\alpha_2}{3},
 \qquad
 \mu_2=\frac{\alpha_1+2\alpha_2}{3}.
 \label{p12-83-c20-s4:eq:representantes-discriminantes-a2}
\end{equation}
Leurs accouplements avec la base simple sont
\[
\begin{array}{c|cc}
 &\alpha_1&\alpha_2\\ \hline
 \mu_1&1&0\\
 \mu_2&0&1
\end{array},
\]
de sorte que \(\mu_1,\mu_2\in A_2^*\). En outre,
\[
 3\mu_1=2\alpha_1+\alpha_2\in A_2,
 \qquad
 3\mu_2=\alpha_1+2\alpha_2\in A_2,
\]
tandis que \(\mu_1,\mu_2\notin A_2\). On obtient
\begin{equation}
 A_2^*/A_2
 =\{\overline\mu_0,\overline\mu_1,\overline\mu_2\}
 \cong\mathbb Z/3\mathbb Z
 \cong\mathbb F_3.
 \label{p12-83-c20-s4:eq:discriminante-a2-f3}
\end{equation}

La somme des représentants se réduit explicitement au moyen de
\[
 \mu_1+\mu_1-\mu_2=\alpha_1,
 \qquad
 \mu_1+\mu_2=\rho,
 \qquad
 \mu_2+\mu_2-\mu_1=\alpha_2.
\]
Par conséquent,
\begin{equation}
 \mu_r+\mu_s-\mu_{r+s\,({\rm mod}\,3)}\in A_2
 \qquad(r,s\in\mathbb F_3).
 \label{p12-83-c20-s4:eq:adicion-representantes-discriminantes-a2}
\end{equation}

\begin{proposition}[Forme discriminante de \(A_2\)]
\label{p12-83-c20-s4:prop:forma-discriminante-a2}
Sous \eqref{p12-83-c20-s4:eq:discriminante-a2-f3}, les formes bilinéaire et quadratique
discriminantes sont
\begin{align}
 b(\overline\mu_r,\overline\mu_s)
 &\equiv\frac{2rs}{3}\pmod{\mathbb Z},
 \label{p12-83-c20-s4:eq:bilineal-discriminante-a2}\\
 q(\overline\mu_r)
 &\equiv\frac{2r^2}{3}\pmod{2\mathbb Z}.
 \label{p12-83-c20-s4:eq:cuadratica-discriminante-a2}
\end{align}
\end{proposition}

\begin{proof}
La formule \eqref{p12-83-c20-s4:eq:norma-a2-autonoma} donne
\[
 \mu_1^2=\mu_2^2=\frac23,
 \qquad
 \langle\mu_1,\mu_2\rangle=\frac13.
\]
Pour \(r,s\in\{0,1,2\}\), ces valeurs coïncident, modulo
\(\mathbb Z\), avec \(2rs/3\). La formule quadratique résulte du choix
\(r=s\) en rappelant que le réseau est pair.
\end{proof}

\subsection{\(12\) facteurs et application de recollement}

Soit
\[
 L_0:=A_2^{12}.
\]
Alors
\[
 L_0^*=(A_2^*)^{12},
 \qquad
 L_0^*/L_0\cong\mathbb F_3^{12}.
\]
Pour \(c=(c_1,\ldots,c_{12})\in\mathbb F_3^{12}\), nous définissons
\begin{equation}
 \phi(c):=(\mu_{c_1},\ldots,\mu_{c_{12}})\in L_0^*.
 \label{p12-83-c20-s4:eq:aplicacion-pegado-cw}
\end{equation}
L'extension déterminée par \(C_W\) est
\begin{equation}
 \begin{aligned}
 N(C_W)
 &:=\bigcup_{c\in C_W}(L_0+\phi(c))\\
 &=\{x\in L_0^*:x\bmod L_0\in C_W\}.
 \end{aligned}
 \label{p12-83-c20-s4:eq:definicion-niemeier-cw}
\end{equation}

\begin{lemma}[Clôture additive du recollement]
\label{p12-83-c20-s4:lem:clausura-pegado-cw}
L'ensemble \(N(C_W)\) est un sous-groupe discret de \(L_0^*\) de rang
vingt-quatre.
\end{lemma}

\begin{proof}
Soient \(x=\ell+\phi(c)\) et \(y=m+\phi(d)\), avec \(\ell,m\in L_0\) et
\(c,d\in C_W\). D'après \eqref{p12-83-c20-s4:eq:adicion-representantes-discriminantes-a2},
\[
 \phi(c)+\phi(d)-\phi(c+d)\in L_0.
\]
Puisque \(C_W\) est linéaire, \(c+d\in C_W\), et alors
\(x+y\in L_0+\phi(c+d)\). La même identité appliquée à l'opposé donne
les inverses. Les inclusions
\[
 L_0\subseteq N(C_W)\subseteq L_0^*
\]
prouvent la discrétion et le rang vingt-quatre.
\end{proof}

\subsection{Intégralité induite par l'autodualité}

La forme bilinéaire discriminante de \(L_0\) est la somme orthogonale de douze
copies de \eqref{p12-83-c20-s4:eq:bilineal-discriminante-a2}. Pour
\(c,d\in\mathbb F_3^{12}\),
\begin{equation}
 b(c,d)
 \equiv\frac23\sum_{i=1}^{12}c_id_i
 \pmod{\mathbb Z}.
 \label{p12-83-c20-s4:eq:bilineal-discriminante-a2-doce}
\end{equation}

\begin{lemma}[Intégralité du recollement]
\label{p12-83-c20-s4:lem:integralidad-pegado-cw}
Le réseau \(N(C_W)\) est entier.
\end{lemma}

\begin{proof}
L'autodualité \(C_W=C_W^\perp\) implique
\[
 \sum_{i=1}^{12}c_id_i\equiv0\pmod3
 \qquad(c,d\in C_W).
\]
D'après \eqref{p12-83-c20-s4:eq:bilineal-discriminante-a2-doce}, la forme discriminante
s'annule sur \(C_W\times C_W\). Par conséquent, l'accouplement de deux classes
recollées quelconques est entier.
\end{proof}

\subsection{Parité induite par les poids du code}

Sur une coordonnée non nulle du code, la forme quadratique discriminante
vaut \(2/3\) modulo \(2\mathbb Z\). Par conséquent,
\begin{equation}
 q(c)\equiv\frac23\operatorname{wt}(c)\pmod{2\mathbb Z}.
 \label{p12-83-c20-s4:eq:cuadratica-discriminante-cw}
\end{equation}

\begin{lemma}[Parité du recollement]
\label{p12-83-c20-s4:lem:paridad-pegado-cw}
Le réseau \(N(C_W)\) est pair.
\end{lemma}

\begin{proof}
Les seuls poids de \(C_W\) sont \(0,6,9,12\), et
\[
 \frac23\{0,6,9,12\}=\{0,4,6,8\}\subseteq2\mathbb Z.
\]
Toute classe de recollement est isotrope pour la forme quadratique discriminante,
de sorte que toutes les normes de \(N(C_W)\) sont paires.
\end{proof}

Soient \(g_1,\ldots,g_6\) les lignes de \(G_W=(I_6\mid A_W)\) et
\(\gamma_i=\phi(g_i)\). La multiplication exacte dans \(A_2^*\) donne
\begin{equation}
 (\langle\gamma_i,\gamma_j\rangle)_{i,j=1}^{6}
 =
 \begin{pmatrix}
 4&2&2&2&2&2\\
 2&4&2&2&2&2\\
 2&2&4&2&2&2\\
 2&2&2&4&2&2\\
 2&2&2&2&4&2\\
 2&2&2&2&2&4
 \end{pmatrix}.
 \label{p12-83-c20-s4:eq:gram-generadores-pegado-cw}
\end{equation}
Cette matrice fournit un contrôle direct supplémentaire de l'intégralité et de la
parité pour les six générateurs de recollement.

\subsection{Indice, déterminant et unimodularité}

Les classes latérales de \(L_0\) qui apparaissent dans
\eqref{p12-83-c20-s4:eq:definicion-niemeier-cw} sont indexées par les
\(|C_W|=729=3^6\) éléments du code et sont distinctes. En conséquence,
\begin{equation}
 [N(C_W):L_0]=3^6.
 \label{p12-83-c20-s4:eq:indice-niemeier-sobre-a2-doce}
\end{equation}
Comme
\[
 \det L_0=(\det A_2)^{12}=3^{12},
\]
la formule du déterminant pour une extension finie donne
\begin{equation}
 \det N(C_W)
 =\frac{\det L_0}{[N(C_W):L_0]^2}
 =\frac{3^{12}}{(3^6)^2}=1.
 \label{p12-83-c20-s4:eq:determinante-niemeier-cw}
\end{equation}

\subsection{Système de racines induit par l'incidence de Witt et classification de ses orbites}

Dans chacune des deux classes discriminantes non nulles de
\(A_2^*/A_2\), la norme minimale est
\[
 \mu_1^2=\mu_2^2=\frac23.
\]
Si \(c\in C_W\setminus\{0\}\), alors
\(\operatorname{wt}(c)\geq6\). Par orthogonalité des douze facteurs,
tout vecteur de \(L_0+\phi(c)\) a une norme au moins égale à
\begin{equation}
 6\cdot\frac23=4.
 \label{p12-83-c20-s4:eq:cota-clases-no-triviales-niemeier}
\end{equation}
Par conséquent, aucune classe de recollement non triviale ne contient de vecteurs de norme
deux. Les racines de \(N(C_W)\) sont exactement les racines déjà présentes dans
\(L_0=A_2^{12}\).

\begin{theorem}[Réalisation du Niemeier de type \(A_2^{12}\)]
\label{p12-83-c20-s4:thm:niemeier-a2-doce-autonomo}
Le réseau \(N(C_W)\) est défini positif, pair et unimodulaire, et son système
de racines est \(A_2^{12}\). En conséquence,
\begin{equation}
 \boxed{N(C_W)\cong N(A_2^{12}).}
 \label{p12-83-c20-s4:eq:identificacion-niemeier-a2-doce}
\end{equation}
\end{theorem}

\begin{proof}
La positivité provient de \(L_0^*\). La
\cref{p12-83-c20-s4:lem:integralidad-pegado-cw,p12-83-c20-s4:lem:paridad-pegado-cw} démontre
l'intégralité et la parité ; \eqref{p12-83-c20-s4:eq:determinante-niemeier-cw} démontre
l'unimodularité. La borne \eqref{p12-83-c20-s4:eq:cota-clases-no-triviales-niemeier}
identifie le système de racines. La classification de Niemeier associe de
manière unique à ce système de racines le réseau indiqué.
\end{proof}

\subsection{Obstructions du recollement ternaire \(C_W\to N(A_2^{12})\) : parité, unimodularité et système de racines}

\begin{criterion}[Réfutation du recollement ternaire]
La construction est réfutée si l'un des faits
suivants se produit :
\begin{enumerate}
 \item l'un des représentants \(\mu_0,\mu_1,\mu_2\) n'appartient pas à
       \(A_2^*\), deux de leurs classes coïncident, ou
       \eqref{p12-83-c20-s4:eq:adicion-representantes-discriminantes-a2} échoue ;
 \item la forme discriminante n'est pas celle donnée dans
       \cref{p12-83-c20-s4:prop:forma-discriminante-a2} ;
 \item \(N(C_W)\) n'est pas stable par addition ou passage à l'inverse ;
 \item il existe \(c,d\in C_W\) avec
       \(\frac23\sum_i c_id_i\notin\mathbb Z\) ;
 \item un mot de \(C_W\) produit
       \(\frac23\operatorname{wt}(c)\notin2\mathbb Z\) ;
 \item l'indice n'est pas \(3^6\) ou le déterminant n'est pas un ;
 \item une classe de recollement non triviale contient un vecteur de norme deux ;
 \item on utilise \(W_{24}\), une dodécade ou une donnée binaire dans l'une
       quelconque des opérations de recollement précédentes.
\end{enumerate}
\end{criterion}


% Unidad U022. Origen incidencial, vecino ternario y Leech.
% Redaccion autonoma desde la autoridad material 1063.

\section{Origine d'incidence, problème radial typé et construction du voisin unimodulaire sans racines}
\label{p12-83-c20-s5:chap:origen-incidencial-vecino-leech}

Le réseau \(N(A_2^{12})\) est déjà construit. Cette unité détermine un
point de base à partir d'un cadre ordonné d'incidences, résout un
problème variationnel dans un domaine déclaré, vérifie les trois conditions
du voisin d'indice trois et démontre que le résultat est le réseau de
Leech. Le vecteur du voisin est noté \(v_{\mathcal F_*}\) : il dépend du
cadre d'incidence et n'intervient pas dans la retenue numérique de \(\alpha\).

\subsection{Cadre ordonné d'incidences}

Considérons les quatre hexades
\begin{align}
 H^-&=\{4,5,6,7,9,10\},
 \label{p12-83-c20-s5:eq:hexada-marco-negativa}\\
 H_e&=\{1,3,4,6,9,10\},
 \label{p12-83-c20-s5:eq:hexada-marco-e}\\
 H_\varphi&=\{1,2,3,4,7,9\},
 \label{p12-83-c20-s5:eq:hexada-marco-phi}\\
 H_{\rm APP}&=\{2,3,5,10,11,12\}.
 \label{p12-83-c20-s5:eq:hexada-marco-app}
\end{align}
La donnée est le cadre ordonné
\begin{equation}
 \mathcal F:=(H^-,H_e,H_\varphi,H_{\rm APP}).
 \label{p12-83-c20-s5:eq:marco-incidencial-ordenado}
\end{equation}
On définit
\begin{equation}
 o(\mathcal F)
 :=(H_e\cap H_\varphi)\setminus(H^-\cup H_{\rm APP}).
 \label{p12-83-c20-s5:eq:operador-origen-incidencial}
\end{equation}

\begin{theorem}[Reconstruction équivariante de l'origine]
\label{p12-83-c20-s5:thm:origen-incidencial-equivariante}
Pour le cadre canonique \(\mathcal F_*\),
\begin{equation}
 \boxed{o(\mathcal F_*)=\{1\}.}
 \label{p12-83-c20-s5:eq:origen-incidencial-singleton}
\end{equation}
En outre, pour toute permutation simultanée \(g\in M_{12}\),
\begin{equation}
 o(g\mathcal F)=g\,o(\mathcal F).
 \label{p12-83-c20-s5:eq:equivarianza-origen-incidencial}
\end{equation}
\end{theorem}

\begin{proof}
On a
\[
 H_e\cap H_\varphi=\{1,3,4,9\}
\]
et
\[
 (H^-\cup H_{\rm APP})\cap\{1,3,4,9\}=\{3,4,9\}.
\]
La différence est \(\{1\}\). Pour l'équivariance, on applique les
identités
\[
 g(A\cap B)=gA\cap gB,
 \quad
 g(A\cup B)=gA\cup gB,
 \quad
 g(A\setminus B)=gA\setminus gB,
\]
valables pour toute bijection.
\end{proof}

La tétrade \(B_K=\{4,6,9,10\}\) ne détermine pas à elle seule l'origine. Le
singleton s'obtient à partir du cadre ordonné complet ; ce n'est pas une coordonnée
choisie après observation du résultat.

\subsection{Chambre radiale positive du système de racines et sélection du cône dominant}

Soit \(\rho_i\) la copie de \(\rho=\alpha_1+\alpha_2\) dans le
\(i\)-ième facteur de \(A_2^{12}\). Nous considérons le domaine
\begin{equation}
 \mathscr R_+
 :=\left\{
 v(a)=\sum_{i=1}^{12}a_i\rho_i:
 a_i=1+3b_i, b_i\in\mathbb Z_{\geq0}
 \right\}.
 \label{p12-83-c20-s5:eq:camara-radial-positiva}
\end{equation}
Puisque les facteurs sont orthogonaux et \(\rho_i^2=2\), si
\(S=\sum_i b_i\), alors
\begin{align}
 v(a)^2
 &=2\sum_{i=1}^{12}(1+3b_i)^2\notag\\
 &=24+12S+18\sum_{i=1}^{12}b_i^2.
 \label{p12-83-c20-s5:eq:norma-radial-expandida}
\end{align}
En réduisant modulo dix-huit,
\[
 v(a)^2\equiv6+12S\pmod{18},
\]
et donc
\begin{equation}
 v(a)^2\equiv0\pmod{18}
 \quad\Longleftrightarrow\quad
 S\equiv1\pmod3.
 \label{p12-83-c20-s5:eq:congruencia-radial-suma-b}
\end{equation}

\begin{proposition}[Minima dans la chambre radiale déclarée]
\label{p12-83-c20-s5:prop:minimos-camara-radial}
Dans le sous-ensemble de \(\mathscr R_+\) défini par
\(v(a)^2\equiv0\pmod{18}\), la norme minimale est \(54\). Elle est atteinte
exactement en douze vecteurs : pour chaque \(j\in\{1,\ldots,12\}\), le
coefficient \(a_j\) vaut quatre et les onze autres valent un.
\end{proposition}

\begin{proof}
La condition \eqref{p12-83-c20-s5:eq:congruencia-radial-suma-b} et la non-négativité des
\(b_i\) impliquent que la plus petite valeur possible de \(S\) est un. Cela
impose un unique indice \(j\) avec \(b_j=1\) et \(b_i=0\) pour
\(i\neq j\). En substituant dans \eqref{p12-83-c20-s5:eq:norma-radial-expandida},
\[
 v(a)^2=24+12+18=54.
\]
Il y a douze choix de \(j\). La valeur admissible suivante de \(S\) est
quatre et produit une norme strictement supérieure.
\end{proof}

L'origine \(o(\mathcal F_*)=\{1\}\) sélectionne
\begin{equation}
 \boxed{
 v_{\mathcal F_*}
 =4\rho_1+\rho_2+\cdots+\rho_{12},
 \qquad
 v_{\mathcal F_*}^2=54.}
 \label{p12-83-c20-s5:eq:vector-marco-incidencial}
\end{equation}

\begin{remark}[Portée exacte de la minimalité]
La norme \(54\) est minimale dans la chambre radiale positive
\eqref{p12-83-c20-s5:eq:camara-radial-positiva}, non dans toute la classe résiduelle. Le vecteur
\[
 u=(\alpha_1-2\alpha_2,\rho,\ldots,\rho)
\]
satisfait
\[
 \alpha_1-2\alpha_2-\rho=-3\alpha_2,
\]
de sorte que
\[
 u\equiv(\rho,\ldots,\rho)
 \equiv v_{\mathcal F_*}\pmod{3A_2^{12}}.
\]
Cependant,
\[
 (\alpha_1-2\alpha_2)^2
 =2(1+2+4)=14,
 \qquad
 u^2=14+11\cdot2=36.
\]
Ce témoin empêche d'étendre l'affirmation de minimalité hors du domaine
de \cref{p12-83-c20-s5:prop:minimos-camara-radial}.
\end{remark}

\subsection{Caractère et noyau de congruence}

Écrivons
\[
 N:=N(A_2^{12}),
 \qquad
 v:=v_{\mathcal F_*}.
\]
On définit
\begin{equation}
 \chi_v:N\longrightarrow\mathbb Z/3\mathbb Z,
 \qquad
 \chi_v(x)=\langle x,v\rangle\bmod3,
 \label{p12-83-c20-s5:eq:caracter-vecino-leech}
\end{equation}
son noyau
\begin{equation}
 N_0:=\ker\chi_v
 =\{x\in N:\langle x,v\rangle\equiv0\pmod3\},
 \label{p12-83-c20-s5:eq:nucleo-congruencia-vecino}
\end{equation}
et l'extension
\begin{equation}
 N_v:=N_0+\mathbb Z\frac v3.
 \label{p12-83-c20-s5:eq:definicion-vecino-indice-tres}
\end{equation}

\subsection{Condition I1 : surjectivité et exclusion des racines}

Par construction, \(v\in A_2^{12}\subseteq N\). Les racines d'un facteur
\(A_2\) sont
\[
 \pm\alpha_1,
 \quad
 \pm\alpha_2,
 \quad
 \pm\rho.
\]
Leurs accouplements avec \(\rho\) sont \(\pm1\) ou \(\pm2\). Tous les
coefficients de \(v\) sont congrus à un modulo trois :
\[
 4\equiv1\pmod3,
 \qquad
 1\equiv1\pmod3.
\]
Par conséquent, aucune racine de \(A_2^{12}\) n'a d'accouplement nul avec
\(v\) modulo trois. En particulier,
\[
 \chi_v(\alpha_{1,1})=1,
\]
de sorte que \(\chi_v\) est surjectif et
\begin{equation}
 [N:N_0]=3.
 \label{p12-83-c20-s5:eq:indice-nucleo-vecino}
\end{equation}
En outre, aucune racine de \(N\) n'appartient à \(N_0\).

\subsection{Condition I2 : divisibilité de la norme}

L'identité \eqref{p12-83-c20-s5:eq:vector-marco-incidencial} donne
\begin{equation}
 v^2=54\equiv0\pmod{18},
 \qquad
 \left(\frac v3\right)^2=6.
 \label{p12-83-c20-s5:eq:condicion-i2-vecino}
\end{equation}
En particulier, \(\langle v,v\rangle\equiv0\pmod3\), donc
\(v\in N_0\).

\subsection{Condition I3 : classe duale d'ordre \(3\)}

Si \(x\in N_0\), il existe \(k\in\mathbb Z\) avec
\(\langle x,v\rangle=3k\). Par conséquent,
\[
 \left\langle x,\frac v3\right\rangle=k\in\mathbb Z,
\]
et
\begin{equation}
 \frac v3\in N_0^*.
 \label{p12-83-c20-s5:eq:v-tercio-en-dual-nucleo}
\end{equation}
D'autre part,
\[
 \left\langle\frac v3,\alpha_{1,1}\right\rangle=\frac43.
\]
Puisque \(\alpha_{1,1}\in N\) et que \(N\) est entier, cela démontre
\begin{equation}
 \frac v3\notin N.
 \label{p12-83-c20-s5:eq:v-tercio-no-en-niemeier}
\end{equation}
Enfin, \(3(v/3)=v\in N_0\). La classe de \(v/3\) dans
\(N_0^*/N_0\) est non triviale et a pour ordre exactement trois. En
conséquence,
\begin{equation}
 [N_v:N_0]=3.
 \label{p12-83-c20-s5:eq:indice-extension-vecino}
\end{equation}

\subsection{Parité, intégralité et déterminant du voisin}

\begin{proposition}[Structure unimodulaire du voisin]
\label{p12-83-c20-s5:prop:estructura-unimodular-vecino}
Le réseau \(N_v\) est entier, pair et unimodulaire.
\end{proposition}

\begin{proof}
Puisque \(N\) est unimodulaire et \([N:N_0]=3\),
\[
 \det N_0=[N:N_0]^2\det N=9.
\]
Comme \([N_v:N_0]=3\),
\[
 \det N_v=\frac{\det N_0}{[N_v:N_0]^2}=1.
\]

Soit \(y=x+m(v/3)\), avec \(x\in N_0\), et écrivons
\(\langle x,v\rangle=3k\). Alors
\begin{align*}
 y^2
 &=x^2+2m\left\langle x,\frac v3\right\rangle
       +m^2\left(\frac v3\right)^2\\
 &=x^2+2mk+6m^2.
\end{align*}
Les trois termes forment un entier pair. La forme polarisée prouve
l'intégralité, et le déterminant un prouve l'unimodularité.
\end{proof}

\subsection{\(2\) classes globales et \(6\) classes locales}

Globalement, \eqref{p12-83-c20-s5:eq:definicion-vecino-indice-tres} n'ajoute que deux classes
latérales non triviales :
\[
 N_0+\frac v3,
 \qquad
 N_0-\frac v3.
\]
Localement, dans chaque facteur \(A_2\), une coordonnée de ces classes appartient
à l'une des six classes
\begin{equation}
 \mathscr C_{j,\pm}
 :=A_2+\mu_j\pm\frac\rho3,
 \qquad j\in\{0,1,2\}.
 \label{p12-83-c20-s5:eq:seis-clases-locales-vecino}
\end{equation}
Dans la première coordonnée, \(4\rho/3\equiv\rho/3\pmod{A_2}\), de sorte que
la même liste vaut pour les douze facteurs.

\begin{table}[htbp]
\centering
\caption{Représentants minimaux des six classes latérales locales.}
\label{p12-83-c20-s5:tab:seis-clases-locales-vecino}
\begin{tabular}{cclc}
\toprule
classe&représentant minimal&\((u,v)\)&
\(2(u^2-uv+v^2)/9\)\\
\midrule
\(\mathscr C_{0,-}\)&\(-\rho/3\)&\((-1,-1)\)&\(2/9\)\\
\(\mathscr C_{0,+}\)&\( \rho/3\)&\((1,1)\)&\(2/9\)\\
\(\mathscr C_{1,-}\)&\( \mu_1-\rho/3\)&\((1,0)\)&\(2/9\)\\
\(\mathscr C_{1,+}\)&\( \mu_1+\rho/3-\rho\)&\((0,-1)\)&\(2/9\)\\
\(\mathscr C_{2,-}\)&\( \mu_2-\rho/3\)&\((0,1)\)&\(2/9\)\\
\(\mathscr C_{2,+}\)&\( \mu_2+\rho/3-\rho\)&\((-1,0)\)&\(2/9\)\\
\bottomrule
\end{tabular}
\end{table}

Ici, \((u,v)\) signifie le représentant
\((u\alpha_1+v\alpha_2)/3\).

\begin{lemma}[Minimum exact des six classes]
\label{p12-83-c20-s5:lem:minimo-seis-clases-locales}
Chaque classe de \eqref{p12-83-c20-s5:eq:seis-clases-locales-vecino} a pour norme minimale
exactement \(2/9\).
\end{lemma}

\begin{proof}
La norme de \((u\alpha_1+v\alpha_2)/3\) est
\[
 \frac29(u^2-uv+v^2).
\]
La forme entière \(u^2-uv+v^2\) est définie positive. Dans une classe résiduelle
non nulle, elle ne prend pas la valeur zéro, donc elle vaut au moins un. Les six
représentants de la \cref{p12-83-c20-s5:tab:seis-clases-locales-vecino} atteignent un.
\end{proof}

\subsection{Critère d'absence de racines dans le réseau de Leech et conséquence pour la distance minimale}

\begin{theorem}[Absence de vecteurs de norme deux]
\label{p12-83-c20-s5:thm:vecino-sin-raices}
Le réseau \(N_v\) ne contient pas de vecteurs de norme deux.
\end{theorem}

\begin{proof}
Toutes les racines de \(N\) sont celles de \(A_2^{12}\). La condition I1 prouve
qu'aucune n'appartient à \(N_0\). La classe \(N_0\) ne contient donc pas
de vecteurs de norme deux.

Dans une nouvelle classe \(N_0\pm v/3\), chacune des douze coordonnées
appartient à l'une des six classes locales. D'après
\cref{p12-83-c20-s5:lem:minimo-seis-clases-locales}, chaque coordonnée contribue pour au moins
\(2/9\). L'orthogonalité donne
\[
 \left\|x\pm\frac v3\right\|^2
 \geq12\cdot\frac29=\frac83>2.
\]
Les nouvelles classes ne contiennent pas non plus de racines.
\end{proof}

\subsection{Identification au réseau de Leech}

\begin{theorem}[Construction du réseau de Leech]
\label{p12-83-c20-s5:thm:identificacion-vecino-leech}
Le réseau
\begin{equation}
 N_v
 =\ker(\langle\,\cdot\,,v_{\mathcal F_*}\rangle\bmod3)
  +\mathbb Z\frac{v_{\mathcal F_*}}3
 \label{p12-83-c20-s5:eq:formula-final-vecino-leech}
\end{equation}
est isométrique au réseau de Leech :
\begin{equation}
 \boxed{N_v\cong\Lambda_{24}.}
 \label{p12-83-c20-s5:eq:identificacion-leech-autonoma}
\end{equation}
\end{theorem}

\begin{proof}
Le réseau \(N_v\) est de rang vingt-quatre et défini positif. La
\cref{p12-83-c20-s5:prop:estructura-unimodular-vecino} démontre qu'il est pair et unimodulaire ;
le \cref{p12-83-c20-s5:thm:vecino-sin-raices} démontre qu'il est dépourvu de racines. Par la
classification de Niemeier, il existe une unique classe d'isométrie de réseaux
définis positifs, pairs et unimodulaires de rang vingt-quatre sans racines :
le réseau de Leech.
\end{proof}

\subsection{Les 243 syndromes du code ternaire poinçonné}

Soit \(G_{11}\) le code ternaire de Golay de paramètres
\([11,6,5]_3\). Son espace de syndromes a pour cardinal
\begin{equation}
 3^{11-6}=3^5=243.
 \label{p12-83-c20-s5:eq:numero-sindromes-golay-ternario}
\end{equation}
Le nombre d'erreurs de poids au plus deux est
\begin{equation}
 1+2\binom{11}{1}+2^2\binom{11}{2}
 =1+22+220=243.
 \label{p12-83-c20-s5:eq:bola-radio-dos-golay-ternario}
\end{equation}

\begin{proposition}[Correspondance parfaite erreur--syndrome]
\label{p12-83-c20-s5:prop:correspondencia-error-sindrome-golay}
Chaque syndrome de \(G_{11}\) possède un unique représentant d'erreur de poids
au plus deux.
\end{proposition}

\begin{proof}
La distance minimale cinq implique que deux erreurs distinctes de poids au
plus deux ne peuvent différer d'un mot du code : leur différence aurait un
poids au plus quatre. L'application erreur--syndrome est donc
injective sur la boule de rayon deux. Les deux ensembles ont pour cardinal
\(243\) d'après \eqref{p12-83-c20-s5:eq:numero-sindromes-golay-ternario} et
\eqref{p12-83-c20-s5:eq:bola-radio-dos-golay-ternario} ; l'application est bijective.
\end{proof}

Ces \(243\) syndromes ne sont pas les \(243\) mots visibles obtenus à partir du
catalogue TPK. Ils partagent un cardinal, mais possèdent des domaines, des équivalences et
des fonctions distincts.

\subsection{Obstructions du voisin ternaire \(N_v\cong\Lambda_{24}\) : classe d'ordre \(3\), absence de racines et séparation des cardinaux \(243\)}

\begin{criterion}[Réfutation du voisin et de l'identification de Leech]
La construction est réfutée si l'un des faits
suivants se produit :
\begin{enumerate}
 \item l'opérateur \(o(\mathcal F)\) ne produit pas \(\{1\}\) ou n'est pas
       équivariant ;
 \item il existe dans la chambre radiale un vecteur admissible de norme inférieure à
       \(54\), ou le nombre de minimiseurs n'est pas douze ;
 \item on affirme une minimalité globale malgré le témoin de norme \(36\) ;
 \item \(\chi_v\) n'est pas surjectif ou une racine appartient à \(N_0\) ;
 \item \(v^2\not\equiv0\pmod{18}\) o \((v/3)^2\neq6\);
 \item \(v/3\notin N_0^*\), \(v/3\in N\), ou sa classe n'est pas d'ordre trois ;
 \item le déterminant de \(N_v\) n'est pas un ou le voisin cesse d'être pair ;
 \item l'une des six classes locales a un minimum inférieur à \(2/9\) ;
 \item l'une des deux nouvelles classes globales contient une racine ;
 \item on confond les deux classes globales avec les six classes locales ;
 \item on utilise \(W_{24}\) comme antécédent du recollement ou du voisin ;
 \item les \(243\) syndromes sont identifiés aux \(243\) mots
       visibles au motif qu'ils partagent le même cardinal.
\end{enumerate}
\end{criterion}


\section{Holonomie finie, configuration de Schläfli et symétrie \texorpdfstring{\(E_6\)}{E6}}
\label{p12-83-c20-s6:chap:holonomia-schlafli-e6-nuevo}
\index{holonomie finie}
\index{Schläfli}
\index{E6@\(E_6\)}

La branche de Witt n'épuise pas les structures exceptionnelles accessibles à partir d'APP. La branche parallèle de ce chapitre part des deux polarisations mixtes somme--produit et produit--somme, avec une orientation de plaquette et une section centrale déclarées. La courbure calculée détermine la forme alternée ; celle-ci définit l'extension centrale. Une extension n'est pas inférée de la seule cardinalité 27.

\subsection{Courbures mixtes d’APP}

Sur \(T_9=(\mathbb Z/9\mathbb Z)^2\), avec
\(S(x,y)=x+y\) et \(P(x,y)=xy\), soient
\[
 \Delta_xf=f(x+1,y)-f(x,y),\qquad
 \Delta_yf=f(x,y+1)-f(x,y).
\]
Les connexions mixtes sont
\[
 A^{SP}=(\Delta_xS,\Delta_yP)=(1,x),\qquad
 A^{PS}=(\Delta_xP,\Delta_yS)=(y,1).
\]
Leur courbure discrète est
\[
 F_A=\Delta_xA_y-\Delta_yA_x.
\]

\begin{proposition}[Courbures orientées opposées]
\label{p12-83-c20-s6:prop:curvaturas-app-e6-nuevo}
\[
 F_{A^{SP}}=1,\qquad F_{A^{PS}}=-1\pmod9.
\]
La circulation sur un rectangle orienté de côtés \(r,s\) est,
respectivement, \(rs\) et \(-rs\).
\end{proposition}

\begin{proof}
\(\Delta_xx-\Delta_y1=1\) et
\(\Delta_x1-\Delta_yy=-1\). Un rectangle contient \(rs\) plaquettes ; lors de la
sommation, les arêtes intérieures s’annulent et il reste la circulation du bord.
\end{proof}

Pour des déplacements \(u=(a,b)\), \(v=(c,d)\), la classe alternée est
\begin{equation}
 \sigma(u,v)=ad-bc\pmod9.
 \label{p12-83-c20-s6:eq:forma-alternante-e6-nueva}
\end{equation}
L’orientation fixe le signe de \(\sigma\). La section centrale employée
ci-dessous appartient à la jauge de cette réalisation ; la modifier par un cobord
produit une extension isomorphe, non un nouvel invariant.

\subsection{Extension centrale et réduction ternaire}

Nous définissons le groupe de Heisenberg nonadique
\begin{equation}
 \mathcal H_9=\{(a,b,z):a,b,z\in\mathbb Z/9\mathbb Z\},
 \quad
 (a,b,z)(c,d,w)=(a+c,b+d,z+w+bc).
 \label{p12-83-c20-s6:eq:heisenberg-9-nuevo}
\end{equation}
Le commutateur de deux déplacements horizontaux est
\[
 [(a,b,0),(c,d,0)]=(0,0,bc-ad),
\]
qui enregistre \(-\sigma(u,v)\). En réduisant chaque coordonnée modulo trois, on
obtient
\begin{equation}
 H_3=\{(a,b,z):a,b,z\in\mathbb F_3\},\qquad |H_3|=27,
 \label{p12-83-c20-s6:eq:heisenberg-3-nuevo}
\end{equation}
avec la même loi. Le noyau de \(\mathcal H_9\to H_3\) est \((3\mathbb Z/9\mathbb Z)^3\), de cardinal 27.

\subsection{Représentation de Weyl--Heisenberg et phase centrale}

L’extension nonadique admet une réalisation unitaire explicite. Soit
\[
 \zeta_9:=\exp(2\pi i/9),
 \qquad
 \mathscr H_9^{\rm Sch}:=\mathbb C^9,
\]
de base \(\{|n\rangle:n\in\mathbb Z/9\mathbb Z\}\). Nous définissons
\begin{equation}
 X|n\rangle:=|n+1\rangle,
 \qquad
 Z|n\rangle:=\zeta_9^n|n\rangle.
 \label{p12-83-c20-s6:eq:weyl-desplazamiento-fase-u029}
\end{equation}
On a
\[
 ZX=\zeta_9XZ
\]
et, par multiplication,
\begin{equation}
 (X^aZ^b)(X^cZ^d)
 =\zeta_9^{bc}X^{a+c}Z^{b+d}.
 \label{p12-83-c20-s6:eq:cociclo-weyl-nonadico-u029}
\end{equation}

\begin{theorem}[Représentation fidèle de Weyl--Heisenberg]
\label{p12-83-c20-s6:thm:representacion-weyl-heisenberg-u029}
L'application
\begin{equation}
 \rho_{\rm WH}:\mathcal H_9\longrightarrow
 U(\mathscr H_9^{\rm Sch}),
 \qquad
 \rho_{\rm WH}(a,b,z):=\zeta_9^zX^aZ^b,
 \label{p12-83-c20-s6:eq:representacion-weyl-heisenberg-u029}
\end{equation}
est une représentation unitaire fidèle. Pour
\(u=(a,b)\) et \(v=(c,d)\), son commutateur satisfait
\begin{equation}
 [\rho_{\rm WH}(u,0),\rho_{\rm WH}(v,0)]
 =\zeta_9^{\,bc-ad}I
 =\zeta_9^{-\sigma(u,v)}I.
 \label{p12-83-c20-s6:eq:conmutador-weyl-holonomia-u029}
\end{equation}
\end{theorem}

\begin{proof}
L’identité \eqref{p12-83-c20-s6:eq:cociclo-weyl-nonadico-u029} reproduit le terme
\(bc\) de \eqref{p12-83-c20-s6:eq:heisenberg-9-nuevo} ; d’où la propriété
d’homomorphisme. Si \(\zeta_9^zX^aZ^b=I\), l’action sur la base impose
successivement \(a=0\), \(b=0\) et \(z=0\). Cela prouve la fidélité. Comparer
les produits dans les deux ordres donne
\(\zeta_9^{bc-ad}\), qui est la phase alternée de
\eqref{p12-83-c20-s6:eq:forma-alternante-e6-nueva}.
\end{proof}

La représentation distingue deux observables conjugués. Soient
\[
 Q_{\rm H}:=\operatorname{diag}(0,1,\ldots,8),
 \qquad
 F_{jk}:=\frac13\zeta_9^{jk},
 \qquad
 P_{\rm H}:=F^\dagger Q_{\rm H}F.
\]
Les bases propres de \(Q_{\rm H}\) et \(P_{\rm H}\) sont mutuellement non biaisées, puisque \(|F_{jk}|^2=1/9\). Pour tout vecteur unitaire \(\psi\in\mathscr H_9^{\rm Sch}\),
\begin{align}
 \operatorname{Var}_\psi(Q_{\rm H})
 \operatorname{Var}_\psi(P_{\rm H})
 &\ge
 \frac14\left|
 \langle\psi,[Q_{\rm H},P_{\rm H}]\psi\rangle
 \right|^2,
 \label{p12-83-c20-s6:eq:incertidumbre-weyl-u029}\\
 H_{Q_{\rm H}}(\psi)+H_{P_{\rm H}}(\psi)
 &\ge\log9.
 \label{p12-83-c20-s6:eq:incertidumbre-entropica-weyl-u029}
\end{align}
Ces opérateurs sont adimensionnels. Une carte physique d'unités est une
réalisation ultérieure et n'intervient pas dans l'extension centrale.

\subsection{Action de \(10\) déplacements sur le graphe de Cayley}

Pour une forme linéaire \(\ell(a,b)=pa+qb\), nous posons
\[
 D_\ell(a,b)=\left(a,b,\frac{ab}{2}+\ell(a,b)\right),
 \qquad Z=(0,0,1),
\]
où \(2^{-1}=2\) dans \(\mathbb F_3\). L'ensemble stable par inversion
\begin{equation}
 S_\ell=
 \{D_\ell(v):v\in\mathbb F_3^2\setminus\{0\}\}
 \cup\{Z,Z^{-1}\}
 \label{p12-83-c20-s6:eq:conjunto-conexion-e6}
\end{equation}
possède dix éléments et ne contient pas l'identité.

\begin{theorem}[Indépendance de la jauge linéaire]
\label{p12-83-c20-s6:thm:independencia-calibre-e6}
Les neuf fonctionnelles \(\ell\) produisent des graphes de Cayley isomorphes. Les
neuf ensembles \(S_\ell\) forment une seule orbite sous
\(\operatorname{Aut}(H_3)\).
\end{theorem}

\begin{proof}
L'application
\[
 \phi_\ell(a,b,z)=(a,b,z+\ell(a,b))
\]
est un automorphisme central et satisfait
\(\phi_\ell(S_0)=S_\ell\). Pour écrire la famille complète dans les
coordonnées polarisées de \cref{p12-83-c20-s6:eq:heisenberg-9-nuevo}, soit
\(c(v,w)=b c\) pour \(v=(a,b)\), \(w=(c,d)\). Étant donné
\(M\in\operatorname{GL}(2,3)\), le défaut

\[
 B_M(v,w)=c(Mv,Mw)-\det(M)c(v,w)
\]
est symétrique, car les parties alternées des deux termes coïncident. Nous définissons la correction quadratique
\begin{equation}
 q_M(v)=2B_M(v,v),
 \qquad
 q_M(v+w)-q_M(v)-q_M(w)=B_M(v,w),
 \label{p12-83-c20-s6:eq:correccion-cuadratica-automorfismos-h3}
\end{equation}
où \(2=2^{-1}\) dans \(\mathbb F_3\). La famille complète d'automorphismes s'écrit alors
\begin{equation}
 (v,t)\longmapsto
 \bigl(Mv,\det(M)t+q_M(v)+\ell(v)\bigr),
 \qquad M\in\operatorname{GL}(2,3),\
 \ell\in(\mathbb F_3^2)^*.
 \label{p12-83-c20-s6:eq:automorfismos-h3-completos}
\end{equation}
L’identité polaire de \cref{p12-83-c20-s6:eq:correccion-cuadratica-automorfismos-h3}
prouve directement que le produit central est conservé. La famille possède
\(48\cdot9=432\) éléments. L’énumération exhaustive des
\(\binom{13}{5}=1287\) sous-ensembles stables par inversion de cardinal dix
trouve exactement ces neuf sous-ensembles avec les paramètres dérivés
ci-après.
\end{proof}

Écrivons \(\underline S=\sum_{s\in S}s\) dans
\(\mathbb Z[H_3]\). Le dénombrement des produits ordonnés donne
\begin{equation}
 \underline S^{\,2}
 =10e+\underline S+5(\underline{H_3}-e-\underline S)
 =5\underline{H_3}+5e-4\underline S.
 \label{p12-83-c20-s6:eq:identidad-grupo-schlafli-nueva}
\end{equation}
En effet, l'identité reçoit dix factorisations ; un élément de \(S\) en reçoit une ; chacun des seize autres en reçoit cinq.

\begin{theorem}[Graphe des vingt-sept droites]
\label{p12-83-c20-s6:thm:grafo-27-rectas-nuevo}
La matrice d’adjacence \(A_\ell\) de
\(\Gamma_\ell=\operatorname{Cay}(H_3,S_\ell)\) satisfait
\begin{equation}
 A_\ell^2=5I-4A_\ell+5J.
 \label{p12-83-c20-s6:eq:identidad-schlafli-nueva}
\end{equation}
Par conséquent,
\[
 \Gamma_\ell=\operatorname{srg}(27,10,1,5),
 \qquad
 \operatorname{Spec}(A_\ell)=\{10^1,1^{20},(-5)^6\}.
\]
Il est isomorphe au graphe d'intersection des vingt-sept droites d'une surface cubique lisse.
\end{theorem}

\begin{proof}
La représentation régulière de
\cref{p12-83-c20-s6:eq:identidad-grupo-schlafli-nueva} produit
\cref{p12-83-c20-s6:eq:identidad-schlafli-nueva}. Ses entrées indiquent que tout sommet
possède dix voisins, qu’une paire adjacente possède un voisin commun et qu’une paire non
adjacente en possède cinq. Sur la droite constante, \(A_\ell\) agit par 10 ;
dans son complément, \(J=0\) et
\((A_\ell-I)(A_\ell+5I)=0\). La trace nulle fixe les multiplicités 20 et 6.

Dans le modèle d’une cubique obtenue en éclatant six points de
\(\mathbb P^2\), les 27 classes sont
\[
 E_i,\qquad L_{ij}=H-E_i-E_j,\qquad
 Q_i=2H-\sum_{j\ne i}E_j.
\]
Avec \(H^2=1\), \(E_i^2=-1\), \(H\cdot E_i=0\) et
\(E_i\cdot E_j=0\) pour \(i\ne j\), sa matrice d'intersection possède les
mêmes paramètres. L'identification à la configuration classique utilise le
théorème d'unicité de cette réalisation des 27 droites
\cite{Manivel2005} ; elle n'est pas attribuée à un certificat inexistant.
\end{proof}

L’assertion d’isomorphie peut être rendue entièrement effective par
une bijection calculée. L’étiquetage n’est pas absolu : une autre section linéaire
ou un automorphisme de la surface produit un autre tableau de la même classe.
\begin{table}[!htbp]
\centering\fontsize{9.4}{10.8}\selectfont
\setlength{\tabcolsep}{4pt}\renewcommand{\arraystretch}{1.04}
\caption{Bijection explicite entre \(H_3\) et les vingt-sept classes de
droites.}
\label{p12-83-c20-s6:tab:h3-27-rectas-u029}
\begin{tabular}{@{}c@{\,}c@{\,}c@{\qquad}l@{}}

\toprule
\(a\)&\(b\)&\(z\)&classe de droite\\
\midrule

0&0&0&\(E_1\)\\
0&0&1&\(L_{12}\)\\
0&0&2&\(Q_2\)\\
0&1&0&\(L_{13}\)\\
0&1&1&\(Q_1\)\\
0&1&2&\(E_3\)\\
0&2&0&\(Q_3\)\\
0&2&1&\(E_2\)\\
0&2&2&\(L_{23}\)\\
1&0&0&\(L_{14}\)\\
1&0&1&\(L_{35}\)\\
1&0&2&\(L_{26}\)\\
1&1&0&\(L_{36}\)\\
1&1&1&\(L_{24}\)\\
1&1&2&\(L_{15}\)\\
1&2&0&\(L_{25}\)\\
1&2&1&\(L_{16}\)\\
1&2&2&\(L_{34}\)\\
2&0&0&\(Q_4\)\\
2&0&1&\(L_{46}\)\\
2&0&2&\(E_6\)\\
2&1&0&\(E_5\)\\
2&1&1&\(Q_6\)\\
2&1&2&\(L_{56}\)\\
2&2&0&\(L_{45}\)\\
2&2&1&\(E_4\)\\
2&2&2&\(Q_5\)\\
\bottomrule
\end{tabular}
\end{table}

Comme contrôle local, les dix voisins de \(e=(0,0,0)\) sont les éléments de
\(S_0\). Le tableau les transforme en les dix classes qui rencontrent \(E_1\) :
les cinq \(L_{1j}\) et les cinq \(Q_j\), avec \(2\le j\le6\). La
vérification complète compare les 729 entrées des deux matrices
d’adjacence.

Le complément \(B_\ell=J-I-A_\ell\) vérifie
\begin{equation}
 B_\ell^2=8I+2B_\ell+8J,
 \qquad
 \operatorname{srg}(27,16,10,8),
 \qquad
 \operatorname{Spec}(B_\ell)=\{16^1,4^6,(-2)^{20}\}.
 \label{p12-83-c20-s6:eq:complemento-schlafli}
\end{equation}
C'est le graphe de Schläfli dans la convention d'adjacence complémentaire ; \(A_\ell\) est le graphe d'intersection de paramètres \((27,10,1,5)\).

Chaque arête appartient à un unique triangle, puisque \(\lambda=1\). Le nombre
de triangles est
\[
 \frac{27\cdot10}{6}=45.
\]
Le groupe complet des automorphismes est d’ordre
\[
 27\cdot1920=51840
\]
et, par son action sur la configuration de droites, il est identifié à \(W(E_6)\), et pas seulement par égalité des ordres \cite{Manivel2005}.

\subsection{Le réseau de racines qui réalise \texorpdfstring{\(E_6\)}{E6}}

Pour expliciter l’objet sur lequel agit ce groupe, soit
\begin{equation}
 \operatorname{Pic}(X)
 =\mathbb ZH\oplus\bigoplus_{i=1}^6\mathbb ZE_i,
 \qquad
 H^2=1,\qquad E_i^2=-1,\qquad H\cdot E_i=E_i\cdot E_j=0\ (i\ne j),
 \label{p12-83-c20-s6:eq:picard-cubica-e6}
\end{equation}
le réseau de Picard de l’éclatement de six points généraux de
\(\mathbb P^2\). Sa classe canonique est
\[
 K_X=-3H+E_1+\cdots+E_6.
\]
Le complément orthogonal
\begin{equation}
 L_{E_6}=K_X^\perp\subset\operatorname{Pic}(X),
 \qquad (x,y)_{E_6}:=-x\cdot y,
 \label{p12-83-c20-s6:eq:reticulo-raices-e6}
\end{equation}
est défini positif de rang six. Une base de racines simples est
\begin{equation}
 \begin{aligned}
 \alpha_1&=E_1-E_2,& \alpha_2&=E_2-E_3,&
 \alpha_3&=E_3-E_4,\\
 \alpha_4&=E_4-E_5,& \alpha_5&=E_5-E_6,&
 \alpha_6&=H-E_1-E_2-E_3.
 \end{aligned}
 \label{p12-83-c20-s6:eq:raices-simples-e6}
\end{equation}
Chaque \(\alpha_i\) est orthogonal à \(K_X\), est de norme deux et leurs produits
sont \(-1\) exactement pour les paires adjacentes du diagramme
\(E_6\) : la chaîne \(1-2-3-4-5\), avec le nœud \(6\) relié au nœud \(3\).
Par conséquent, leur matrice de Gram est la matrice de Cartan de \(E_6\).

\begin{proposition}[Les soixante-douze racines et leurs réflexions]
\label{p12-83-c20-s6:prop:setenta-dos-raices-e6}
Les racines de \(L_{E_6}\) sont
\begin{equation}
 \begin{aligned}
 &E_i-E_j &&(i\ne j),\\
 &\pm(H-E_i-E_j-E_k) &&(1\le i<j<k\le6),\\
 &\pm(2H-E_1-\cdots-E_6),
 \end{aligned}
 \label{p12-83-c20-s6:eq:enumeracion-raices-e6}
\end{equation}
au total \(30+40+2=72\). Pour chaque racine \(\alpha\),
\begin{equation}
 s_\alpha(x)=x+(x\cdot\alpha)\alpha
 \label{p12-83-c20-s6:eq:reflexion-raiz-e6}
\end{equation}
préserve l’intersection et \(K_X\). Les six réflexions simples engendrent
\(W(E_6)\) et permutent fidèlement les 27 classes
\(E_i,L_{ij},Q_i\) de \cref{p12-83-c20-s6:thm:grafo-27-rectas-nuevo}.
\end{proposition}

\begin{proof}
Substituer chaque classe de \cref{p12-83-c20-s6:eq:enumeracion-raices-e6} dans \cref{p12-83-c20-s6:eq:picard-cubica-e6} donne \(K_X\cdot\alpha=0\) et \(\alpha^2=-2\). Le décompte des classes est celui indiqué. La formule de réflexion est la réflexion orthogonale pour la forme \(-(\cdot)\) ; elle préserve \(K_X\) car \(K_X\cdot\alpha=0\). Le système simple de \cref{p12-83-c20-s6:eq:raices-simples-e6} engendre toutes les racines et son groupe de Coxeter est \(W(E_6)\). Son action sur les 27 classes de courbes exceptionnelles conserve l'intersection, de sorte qu'elle coïncide avec l'action sur le graphe déjà construit \cite{Manivel2005}.
\end{proof}

\subsection{Stratification des 729 mots}

Nous écrivons un mot de six trits sous la forme
\[
 w=(\ell_1,h_1,\ell_2,h_2,\ell_3,h_3),
\]
et nous définissons deux éléments de \(H_3\),
\[
 L(w)=(\ell_1,\ell_2,\ell_3),\qquad
 H(w)=(h_1,h_2,h_3),
\]
en utilisant la loi centrale dans ses produits. Selon
\(L(w)^{-1}H(w)\), les 729 mots se divisent en
\begin{align*}
 \mathcal D_\ell&=\{w:L(w)^{-1}H(w)=e\},\\
 \mathcal I_\ell&=\{w:L(w)^{-1}H(w)\in S_\ell\},\\
 \mathcal S_\ell&=\{w:L(w)^{-1}H(w)\notin\{e\}\cup S_\ell\}.
\end{align*}

\begin{theorem}[Partition Hensel--Schläfli]
\label{p12-83-c20-s6:thm:particion-hensel-schlafli-nueva}
\begin{equation}
 729=27+270+432,
 \label{p12-83-c20-s6:eq:particion-729-e6-nueva}
\end{equation}
correspondant à la diagonale, à l’incidence et au gauchissement.
\end{theorem}

\begin{proof}
Pour chacune des 27 valeurs de \(L(w)\), il y a une unique valeur de
\(H(w)\) dont la différence est l’identité, dix dont la différence est dans \(S_\ell\) et
seize autres. On multiplie \(27\) par \(1,10,16\).
\end{proof}

Le cardinal \(432\) de la strate gauchie coïncide numériquement avec la
multiplicité de la région transversale de \(\pi\) dans
\cref{eq:descomposiciones-cinco-regiones-pi}. Cela n’identifie pas les deux objets : ici,
on compte des mots classés par \(L(w)^{-1}H(w)\), tandis que là,
on compte des préimages d’un registre \(U_6\). Sans opérateur d’entrelacement explicite et
équivariant entre ces domaines, \(432=432\) n’est qu’un contrôle cardinal.

\subsection{Nécessité de la phase centrale}

L’énumération de tous les sous-ensembles stables par inversion de cardinal dix
dans les cinq groupes d’ordre 27 donne
\[
\begin{array}{c|c|c}
\text{groupe}&\text{abélien}&
\#\operatorname{srg}(27,10,1,5)\\ \hline
C_{27}&\text{oui}&0\\
C_9\times C_3&\text{oui}&0\\
C_3^3&\text{oui}&0\\
C_9\rtimes C_3&\text{non}&9\\
H_3&\text{non}&9.
\end{array}
\]
Ainsi, conserver 27 étiquettes mais abélianiser la composition détruit la configuration. Ni la cardinalité ni le degré ne suffisent.

Si \(\mathcal T\) est l’ensemble des 45 triangles, le cubique
\[
 \mathcal C_\Gamma(x_1,\ldots,x_{27})
 =\sum_{\{i,j,k\}\in\mathcal T}x_ix_jx_k
\]
récupère l'adjacence, car chaque arête appartient à un triangle unique. Son stabilisateur par permutations est exactement
\[
 \operatorname{Aut}(\Gamma_\ell)\cong W(E_6).
\]

\begin{figure}[htbp]
\centering
\[
 \text{courbures APP }\pm1
 \longrightarrow\sigma
 \longrightarrow\mathcal H_9
 \longrightarrow H_3
 \longrightarrow\Gamma_{27}
 \longrightarrow\mathcal T_{45}
 \longrightarrow W(E_6).
\]
\caption{Chaîne causale de la branche d’holonomie mixte. Elle est parallèle à
\(C_W\to W_{12}\to M_{12}\) et n’en est pas un maillon.}
\label{p12-83-c20-s6:fig:cadena-holonomia-e6-u029}
\end{figure}

\begin{criterion}[Critères de réfutation de la branche \(E_6\)]
\label{p12-83-c20-s6:crit:refutacion-rama-e6-u029}
La branche est réfutée si les courbures ne sont pas opposées, si le cocycle
central est effacé, si la correction quadratique
\eqref{p12-83-c20-s6:eq:correccion-cuadratica-automorfismos-h3} ne conserve pas le produit,
si \eqref{p12-83-c20-s6:eq:identidad-grupo-schlafli-nueva} ne produit pas les coefficients
\((10,1,5)\), si un groupe abélien produit la même configuration, si la
partition n’a pas pour somme 729, si le tableau
\ref{p12-83-c20-s6:tab:h3-27-rectas-u029} n’entrelace pas les matrices d’adjacence ou si
\(W(E_6)\) est identifié uniquement par son ordre et non par son action sur les
vingt-sept classes et les 72 racines.
\end{criterion}


\section{Couplage orienté APP–Fock et rigidité transversale de \(12\) fibres}
\label{p12-83-c20-s7:chap:acoplamiento-app-fock}
\index{APP--Fock!couplage orienté}
\index{Fock!degré deux}
\index{orientation de feuillet}

Le déterminant cyclotomique ordinaire ne distingue pas \(C\) de \(C^{-1}\).
L'information somme--produit est conservée en introduisant explicitement la
parité de feuillet et une droite d'orientation. Le couplage résultant
produit l'indice \(54\) à partir de l'agrégat APP. Une fois fixées les quatre
conditions du couplage, \(54\), \(J\) et \(196884\) ne sont pas des entrées de
l'évaluation. Cette indépendance évaluative n'affirme ni l'unicité universelle
de l'ansatz ni l'indépendance historique de sa conception.

\subsection{Cadre triangulaire dans chaque orbite}

Soit
\[
M=\mathbb Z[\mathbb Z/9\mathbb Z]
\]
et soit \(M_{\mathrm{bal}}^{C_3}\) le sous-module des éléments invariants
\(\nu\) dont les six coefficients unitaires coïncident :
\[
\nu_u=c(\nu)
\qquad
\bigl(u\in(\mathbb Z/9\mathbb Z)^\times\bigr).
\]
Les coefficients de \(e_0,e_3,e_6\) demeurent libres.

Pour une orbite libre
\(\mathcal O=\{u,4u,7u\}\), posons
\[
b_v=e_v-e_{4v}\in A(\mathcal O).
\]
Si \(b_v^\flat\) est le covecteur défini par la forme de \(A_2\), les trois
racines du cadre triangulaire satisfont
\[
\boxed{
\sum_{v\in\mathcal O}b_v\otimes b_v^\flat
=3I_{A(\mathcal O)}.}
\]
En conséquence,
\[
\sum_{v\in\mathcal O}\nu_v\,b_v\otimes b_v^\flat
=3c(\nu)I.
\]
Le facteur trois provient de la résolution de l'identité par le cadre de
racines.

\subsection{Couplage orienté \(\mathcal H_{\Sigma,\Pi}^{(m)}\) entre le bilan \(C_3\), la parité de feuillet et le degré \(2\) de l'espace de Fock}

Pour \(m\in2\mathbb N_{>0}\), soit
\[
\begin{aligned}
\mathcal V_m&=\mathcal V_{\Sigma,m}\oplus\mathcal V_{\Pi,m},\\
\mathcal V_{\Sigma,m}&=(A_2\otimes\mathbb C)^{\oplus m/2},\qquad
\mathcal V_{\Pi,m}=(A_2\otimes\mathbb C)^{\oplus m/2},\\
g_m&=C^{\oplus m/2}\oplus(C^{-1})^{\oplus m/2},\\
\mathsf E_m&=I_{\mathcal V_{\Sigma,m}}
\oplus(-I_{\mathcal V_{\Pi,m}}).
\end{aligned}
\]
La spécialisation \(m=12\), avec les étiquettes de paire, de feuillet et d'orbite,
est \(\mathcal V_{12}=W_{\mathrm{APP}}\otimes\mathbb C\). Cette notation
évite de confondre l'espace de représentation avec le système de Witt
\(W_{12}=S(5,6,12)\). Le degré deux de l'espace de
Fock et l'action induite sont
\[
\mathcal F_2(\mathcal V_m)
=\mathcal V_m[2]\oplus\operatorname{Sym}^2(\mathcal V_m[1]),
\]
\[
\Gamma_2(g_m)
=g_m|_{\mathcal V_m[2]}
\oplus\operatorname{Sym}^2(g_m|_{\mathcal V_m[1]}).
\]

Soit \(o_{\Sigma\Pi}\) la droite d'orientation de signe
\(\Pi-\Sigma\), et soit \(M_{\mathrm{bal}}^{C_3,-}\) la copie orientée dans
laquelle l'échange de feuillets agit par \(\nu\mapsto-\nu\). Nous définissons
\[
\boxed{
\begin{aligned}
\mathcal H_{\Sigma,\Pi}^{(m)}&:
M_{\mathrm{bal}}^{C_3,-}\longrightarrow
\operatorname{End}(\mathcal F_2(\mathcal V_m))\otimes o_{\Sigma\Pi},\\
\mathcal H_{\Sigma,\Pi}^{(m)}(\nu)
&=3c(\nu)
\left(0_{\mathcal V_m[2]}\oplus\operatorname{Sym}^2\mathsf E_m\right)
\otimes o_{\Sigma\Pi}.
\end{aligned}}
\]
La transformation est linéaire en \(c\), est portée par
\(\operatorname{Sym}^2(\mathcal V_m[1])\), utilise la parité de feuillet et adopte la
normalisation du cadre triangulaire.

Soit \(\sigma\) l'échange de feuillets. Alors
\[
\nu\longmapsto-\nu,\qquad
o_{\Sigma\Pi}\longmapsto-o_{\Sigma\Pi},\qquad
\mathsf E_m\longmapsto-\mathsf E_m,
\]
et
\[
\operatorname{Sym}^2(-\mathsf E_m)
=\operatorname{Sym}^2(\mathsf E_m).
\]
Par conséquent,
\[
\mathcal H_{\Sigma,\Pi}^{(m)}(\sigma\nu)
=\sigma\mathcal H_{\Sigma,\Pi}^{(m)}(\nu).
\]
Lors de l'inversion de l'orientation, le covecteur dual change également de signe ; ainsi
\((-o_{\Sigma\Pi}^{\vee})(-o_{\Sigma\Pi})=1\), et l'évaluation orientée
demeure normalisée.

Une fois
\(o_{\Sigma\Pi}^{\vee}(o_{\Sigma\Pi})=1\) fixé, la trace orientée est
\[
\operatorname{Tr}_{o,\mathcal F_2}
(T\otimes o_{\Sigma\Pi})
=\operatorname{Tr}_{\mathcal F_2}(T).
\]

\begin{theorem}[Indice orienté APP--Fock]
\label{p12-83-c20-s7:thm:indice-orientado-app-fock}
Supposons que \(m\) fibres \(A_2\), avec \(m\) pair, se répartissent également
entre les orientations \(C\) et \(C^{-1}\). Alors
\[
\boxed{
\operatorname{Tr}_{o,\mathcal F_2}
\left(
\mathcal H_{\Sigma,\Pi}^{(m)}(\nu)\Gamma_2(g_m)
\right)
=-\frac{3m\,c(\nu)}2.}
\]
Pour l'agrégat produit par APP,
\[
\delta=\nu_\Pi-\nu_\Sigma
=12e_0+3e_3+3e_6
-3\sum_{u\in(\mathbb Z/9\mathbb Z)^\times}e_u,
\]
on a \(c(\delta)=-3\) et
\[
m=3\ \text{paires}\cdot2\ \text{feuillets}\cdot2\ \text{orbites}=12.
\]
Par conséquent,
\[
\boxed{
\operatorname{Tr}_{o,\mathcal F_2}
\left(
\mathcal H_{\Sigma,\Pi}^{(12)}(\delta)\Gamma_2(g_{12})
\right)=54.}
\]
\end{theorem}

\begin{proof}
Soit \(A=\mathsf E_mg_m\). Les contributions des deux feuillets s'annulent :
\(\operatorname{tr}A=0\). Puisque \(\mathsf E_m\) commute avec \(g_m\) et
\(\mathsf E_m^2=I\),
\[
\operatorname{tr}(A^2)=\operatorname{tr}(g_m^2)=-m.
\]
L’identité
\[
\operatorname{tr}(\operatorname{Sym}^2A)
=\frac{(\operatorname{tr}A)^2+\operatorname{tr}(A^2)}2
\]
donne \(-m/2\) ; le cadre apporte \(3c(\nu)\).
\end{proof}

\begin{corollary}[Minimalité du degré bosonique détecteur]
\label{p12-83-c20-s7:cor:minimalidad-grado-dos-app-fock-u030}
Dans la classe de couplages déclarée, le degré un ne détecte pas l'agrégat
orienté et le degré deux est le premier degré bosonique positif qui le
détecte :
\[
 \operatorname{tr}(\mathsf E_mg_m)=0,
 \qquad
 \operatorname{tr}\!\left(
 \operatorname{Sym}^2(\mathsf E_mg_m)
 \right)=-\frac m2\neq0.
\]
\end{corollary}

\begin{proof}
La première égalité est l'annulation des deux feuillets. La seconde a été
obtenue dans la preuve du \cref{p12-83-c20-s7:thm:indice-orientado-app-fock}. Puisque le degré
zéro ne contient que l'unité et que le degré un donne une trace nulle, le degré deux est
le premier à fournir une réponse non nulle.
\end{proof}

\subsection{Rigidité transversale du nombre de fibres}

Pour un entier \(m\ge1\), nous définissons la série formelle
\begin{equation}
 T_m(q):=
 q^{-1}\prod_{n\ge1}(1+q^n+q^{2n})^{-m}
 \in q^{-1}\mathbb Z[[q]]
 \label{p12-83-c20-s7:eq:familia-fock-m-fibras-u030}
\end{equation}
et, dans la somme orthogonale \(A_2^m\), le vecteur radial
\begin{equation}
 v_m:=4\rho_1+\rho_2+\cdots+\rho_m,
 \qquad
 \|\rho_j\|^2=2.
 \label{p12-83-c20-s7:eq:familia-radial-m-fibras-u030}
\end{equation}
Nous construisons deux fonctions entières sans utiliser préalablement \(m=12\) :
\[
 F(m):=[q]T_m(q),
 \qquad
 R(m):=\|v_m\|^2.
\]

\begin{proposition}[Coefficient de Fock et norme radiale]
\label{p12-83-c20-s7:prop:dos-funciones-rigidez-doce-u030}
Pour tout \(m\ge1\),
\begin{equation}
 F(m)=\frac{m(m-3)}2,
 \qquad
 R(m)=2(m+15).
 \label{p12-83-c20-s7:eq:funciones-rigidez-doce-u030}
\end{equation}
\end{proposition}

\begin{proof}
Jusqu'à l'ordre deux,
\[
 (1+q+q^2)^{-m}
 =1-mq+\frac{m(m-1)}2q^2+O(q^3),
\]
\[
 (1+q^2+q^4)^{-m}=1-mq^2+O(q^4).
\]
Les facteurs avec \(n\ge3\) ne contribuent pas au coefficient de \(q^2\).
Puisque le facteur extérieur de \(T_m\) est \(q^{-1}\),
\[
 F(m)=\frac{m(m-1)}2-m=\frac{m(m-3)}2.
\]
Par orthogonalité des \(m\) copies de \(A_2\),
\[
 R(m)=4^2\|\rho_1\|^2+
 \sum_{j=2}^{m}\|\rho_j\|^2
 =32+2(m-1)=2(m+15).
\]
\end{proof}

\begin{theorem}[Rigidité cyclotomique--radiale]
\label{p12-83-c20-s7:thm:rigidez-transversal-doce-u030}
L'égalité \(F(m)=R(m)\) équivaut à
\begin{equation}
 (m-12)(m+5)=0.
 \label{p12-83-c20-s7:eq:ecuacion-rigidez-doce-u030}
\end{equation}
Dans le domaine \(m\in\mathbb Z_{>0}\), l'unique solution est
\[
 \boxed{m=12,\qquad F(12)=R(12)=54.}
\]
\end{theorem}

\begin{proof}
Égaler les deux expressions de
\eqref{p12-83-c20-s7:eq:funciones-rigidez-doce-u030} produit
\[
 m(m-3)=4(m+15),
\]
c’est-à-dire,
\[
 m^2-7m-60=(m-12)(m+5)=0.
\]
La seconde racine est négative et se situe hors du domaine.
\end{proof}

Cette rigidité appartient à la famille cyclotomique--radiale déclarée. Elle n'utilise pas la
valeur \(54\) pour choisir \(m\) : elle construit d'abord \(F(m)\) et \(R(m)\), puis
résout leur égalité.

\subsection{Factorisation du couplage APP–Fock : hypothèses, invariants et conditions nécessaires de validité et obstructions}

Le couplage se factorise par
\[
M_{\mathrm{bal}}^{C_3,-}/\ker c
\cong\mathbb Z\otimes o_{\Sigma\Pi},
\]
et l'application induite est injective puisque
\(\operatorname{Sym}^2\mathsf E_m\neq0\).
La formule est sélectionnée par les quatre conditions déclarées :
linéarité en \(c\), support dans
\(\operatorname{Sym}^2(\mathcal V_m[1])\), parité
de feuillet et normalisation par le cadre triangulaire. On n'affirme pas l'unicité
parmi tous les couplages naturels ni l'indépendance historique de
l'ansatz. Une fois ces conditions fixées, \(54\), \(J\) et \(196884\) ne
sont pas des entrées de l'évaluation.

Trois contrôles montrent que le résultat n'est pas fixé par une
normalisation rétrospective :
\[
(m,c)=(12,-2)\longmapsto36,\qquad
(m,c)=(10,-3)\longmapsto45,
\]
et
\[
\delta_k=\delta+ke_0
\]
laisse invariant \(\mathcal H_{\Sigma,\Pi}^{(12)}\) tandis que le poids numérique
passe de \(54\) à \(54+9k\). En particulier, pour \(k=1\),
\[
\operatorname{Tr}_{o,\mathcal F_2}
\left(
\mathcal H_{\Sigma,\Pi}^{(12)}(\delta+e_0)\Gamma_2(g_{12})
\right)=54,
\qquad
w(\delta+e_0)=63.
\]

\begin{proposition}[Nécessité de la donnée d'orientation]
\label{p12-83-c20-s7:prop:orientacion-app-fock}
Si l'on oublie \(o_{\Sigma\Pi}\), tout scalaire naturel linéaire en
\(\nu_\Pi-\nu_\Sigma\) et ne dépendant que de la classe de la représentation
s'annule.
\end{proposition}

\begin{proof}
\(C\) et \(C^{-1}\) sont conjugués, de sorte que la classe de représentation est
paire sous l'échange de feuillets. L'agrégat
\(\nu_\Pi-\nu_\Sigma\) est impair. Sans la droite d'orientation, une
transformation simultanément paire et impaire vaut zéro.
\end{proof}

Le théorème établit ainsi le transport exact de l'agrégat vers le degré deux de
Fock. Le défaut local état par état conserve l'obstruction de descente
du \cref{p12-83-c1-s5:cor:obstruccion-isotipica-app-autonoma} ; les deux énoncés appartiennent à
des domaines distincts et sont compatibles.


% Unidad U031. Alineamiento APP--Witt, automorfia de orden tres,
% construccion VOA y Moonshine.
 \section{Du saut spectral 4 aux symétries du Monstre : déploiement APP
\texorpdfstring{$2\to6\to54$}{2--6--54}, réseau de Leech,
\texorpdfstring{$V^\natural$}{V natural} et corridor Moonshine}
\label{sec:cierre-equivalencias-corredor-moonshine-20260825}

Le corridor n'est pas une succession de coïncidences numériques. Il part du bloc
central d'APP, conserve l'orientation somme--produit, se ramifie dans les
couplages APP--Fock et APP--Witt, et se coordonne de nouveau par la
graduation. L'ordre causal contraignant est
\[\begin{aligned}
 B_c=9P_++5P_-
 &\longrightarrow 4\longrightarrow2\longrightarrow6\longrightarrow54\\
 &\longrightarrow
 \left\{
 \begin{array}{l}
  \text{degré deux de Fock et }[q]t_3=54,\\
  \text{alignement APP--Witt et }N_\alpha\cong\Lambda_{24}
 \end{array}
 \right.\\
 &\longrightarrow V^\natural
 \longrightarrow
 \bigl(\operatorname{Aut},\operatorname{ch},g\mapsto T_g\bigr)
 \longrightarrow\text{Moonshine}.
\end{aligned}\]
Les deux branches intermédiaires possèdent des codomaines distincts. La clé commune \(54\)
est un invariant gradué transporté par le corridor ; ce n'est pas une application
linéaire d'un scalaire vers une algèbre d'opérateurs de vertex.

\subsection{Hiérarchie typée et \(2\) réalisations de l'indice \(54\)}

\begin{theorem}[Hiérarchie spectrale et déploiement APP]
\label{thm:cierre-vtm-jerarquia-app}
Soit \(B_c=7I+2R=9P_++5P_-\) le bloc central d'APP.\@ Si
\(M_\Pi,M_\Sigma\) sont les compressions modulaires définies, alors
\[
 9-5=4,
 \qquad (B_c)_{i\ne j}=2,
 \qquad \sum M_\Pi-\sum M_\Sigma=51-45=6,
\]
et le déploiement sur les neuf positions donne
\[
 9\cdot6=459-405=54.
\]
Les quatre entiers appartiennent respectivement au saut spectral, au
couplage local, à la compression modulaire et à l'intégration bidimensionnelle.
\end{theorem}

\begin{proof}
La décomposition spectrale de \(B_c\) donne \(9\) et \(5\), tandis que
\(7I+2R\) fixe le coefficient non diagonal. Les deux sommes agrégées sont \(51\)
et \(45\). Chaque entrée comprimée représente neuf positions, de sorte que le
déploiement multiplie la différence par neuf. Aucun des quatre entiers
ne s'obtient en identifiant rétrospectivement leurs types.
\end{proof}

La droite d'orientation \(o_{\Sigma\Pi}\), la parité de feuillet et la
normalisation triangulaire étant fixées, le couplage orienté
\[
 \mathcal H_{\Sigma,\Pi}^{(m)}:
 M_{\mathrm{bal}}^{C_3,-}
 \longrightarrow
 \operatorname{End}\!\bigl(\mathcal F_2(\mathcal V_m)\bigr)
 \otimes o_{\Sigma\Pi}
\]
satisfait
\[
 \operatorname{Tr}_{o,\mathcal F_2}
 \left(\mathcal H_{\Sigma,\Pi}^{(m)}(\nu)\Gamma_2(g_m)\right)
 =-\frac{3m\,c(\nu)}2.
\]
Pour \(m=12\) et \(c(\delta)=-3\), la trace vaut \(54\). D'autre part,
\[
 t_3(q)=q^{-1}\prod_{n\ge1}(1+q^n+q^{2n})^{-12}
       =q^{-1}-12+54q-76q^2-243q^3+\cdots,
\]
d'où \([q]t_3=54\).

\begin{theorem}[Cône de compatibilité de l'indice commun]
\label{thm:cierre-vtm-indice-comun-54}
Les structures de départ de chaque branche étant conservées, on a
\[
 \boxed{
 D_{\mathrm{APP}}
 =\operatorname{Tr}_{o,\mathcal F_2}
   \left(\mathcal H_{\Sigma,\Pi}^{(12)}(\delta)\Gamma_2(g_{12})\right)
 =[q]t_3
 =v_\alpha^2
 =\frac{196884}{5\cdot729+1}
 =54.}
\]
Cette égalité constitue un cône d'évaluations vers le même entier. Elle n'
identifie pas \(W_{\mathrm{APP}}\), l'espace de Fock, le réseau voisin et
\(V^\natural\).
\end{theorem}

\begin{proof}
Le recensement APP donne \(459-405=54\) ; la trace orientée donne \(54\) ; le développement
du produit cyclotomique donne \([q]t_3=54\) ; et le vecteur radial du voisin
satisfait \(v_\alpha^2=4^2\cdot2+11\cdot2=54\). Enfin,
\(5\cdot729+1=3646\) et \(54\cdot3646=196884\). Les égalités sont comparées
après la construction de chaque objet et ne constituent donc pas une sélection numérique du
corridor.
\end{proof}

\paragraph{Statut exact de la hiérarchie \(4\to2\to6\to54\) et de l'identification réticulaire.}
La hiérarchie \(4\to2\to6\to54\), la trace APP--Fock et l'extraction
\([q]t_3\) sont des identités internes exactes. L'identification du voisin à
Leech utilise la chambre, le vecteur de voisinage et la classification réticulaire
déclarés.
%
 
\chapter{Algèbre d'opérateurs de vertex et Moonshine monstrueux}
\section{Alignement cyclotomique, orbifold d'ordre \(3\) et Moonshine monstrueux}
\label{p12-83-c20-s8:chap:alineamiento-orbifold-moonshine}

Le réseau de Leech construit dans la
\cref{p12-83-c20-s5:thm:identificacion-vecino-leech} permet de passer de l'incidence
finie à une structure graduée infinie. Cette transition ne s'effectuera pas
au moyen d'une chaîne de noms. Seront construits, dans un ordre typé :
\begin{enumerate}
 \item l'entrelacement des douze fibres \(A_2\) ;
 \item une isométrie de réseau sans points fixes d'ordre trois ;
 \item sa trace graduée et son relèvement de type zéro ;
 \item l'orbifold cyclique et son automorphisme dual de classe \(3B\) ;
 \item l'algèbre d'opérateurs de sommets lunaire ;
 \item par des publications distinctes, son groupe d'automorphismes, son caractère
       gradué et les séries de McKay--Thompson.
\end{enumerate}
Moonshine sera le théorème qui coordonne l'action et les séries. On n'utilisera pas
le raccourci non typé
\(V^\natural\to\mathbb M\to\text{Moonshine}\).

\subsection{Alignement des \(12\) fibres APP et Witt}

Les douze fibres de \(W_{\mathrm{APP}}\cong A_2^{12}\) portent trois
étiquettes indépendantes :
\[
 i\in\mathbb Z/3\mathbb Z
 \quad\text{(paire de coordonnées)},\qquad
 \mu\in\{0,1\}
 \quad\text{(feuillet)},\qquad
 \varepsilon\in\{0,1\}
 \quad\text{(orbite cyclotomique)}.
\]
L'origine dodécaphasique construite dans le TPK étant fixée, nous définissons
\begin{equation}
 \lambda(i,\mu,\varepsilon)
 :=4i+2\mu+\varepsilon\pmod{12}.
 \label{p12-83-c20-s8:eq:alineamiento-doce-fibras-u031}
\end{equation}

\begin{lemma}[Bijectivité de l'alignement]
\label{p12-83-c20-s8:lem:biyectividad-alineamiento-u031}
L'application \(\lambda\) de
\eqref{p12-83-c20-s8:eq:alineamiento-doce-fibras-u031} est une bijection sur
\(\mathbb Z/12\mathbb Z\).
\end{lemma}

\begin{proof}
Pour chaque \(i\), les quatre valeurs \(4i+2\mu+\varepsilon\) forment le bloc
\(\{4i,4i+1,4i+2,4i+3\}\). Les trois blocs pour \(i=0,1,2\) sont disjoints
et couvrent les douze classes.
\end{proof}

Soit \(C\) le Coxeter de \(A_2\) et soit \(\mathsf R\) la réflexion qui
satisfait
\begin{equation}
 \mathsf R C\mathsf R=C^{-1}.
 \label{p12-83-c20-s8:eq:reflexion-conjuga-coxeter-u031}
\end{equation}
Les trivialisations \(\iota_b\) et les cadres orientés \(P_b\) étant fixés,
nous posons
\begin{align}
 g_b&:=\iota_b^{-1}P_bCP_b^{-1}\iota_b,
 \label{p12-83-c20-s8:eq:coxeter-calibrado-fibra-u031}\\
 T_{i\mu\varepsilon}
 &:=\iota_{\lambda(i,\mu,\varepsilon)}^{-1}
 P_{\lambda(i,\mu,\varepsilon)}\mathsf R^\mu.
 \label{p12-83-c20-s8:eq:entrelazador-fibra-u031}
\end{align}

\begin{theorem}[Entrelacement APP--Witt]
\label{p12-83-c20-s8:thm:alineamiento-app-witt-u031}
Pour chaque triplet \((i,\mu,\varepsilon)\),
\begin{equation}
 g_{\lambda(i,\mu,\varepsilon)}T_{i\mu\varepsilon}
 =T_{i\mu\varepsilon}C^{(-1)^\mu}.
 \label{p12-83-c20-s8:eq:identidad-entrelazamiento-app-witt-u031}
\end{equation}
La somme directe définit un isomorphisme \(C_3\)-équivariant
\begin{equation}
 T_\lambda:
 W_{\mathrm{APP}}\otimes\mathbb C
 \xrightarrow{\;\cong\;}
 W_c:=\bigoplus_{b=0}^{11}
 (A_{2,b}\otimes\mathbb C).
 \label{p12-83-c20-s8:eq:isomorfismo-app-witt-u031}
\end{equation}
\end{theorem}

\begin{proof}
Pour \(\mu=0\), la définition de \(g_b\) donne directement la conjugaison
par \(C\). Pour \(\mu=1\), on insère
\(\mathsf RC\mathsf R=C^{-1}\). Le
\cref{p12-83-c20-s8:lem:biyectividad-alineamiento-u031} garantit que la somme directe
utilise exactement une fois chaque position de Witt et conserve paire, feuillet et orbite.
\end{proof}

Le théorème est relatif au transport simultané de l'ordre des paires, des feuillets et des
orbites, de l'origine de Witt, des cadres et de l'orientation. Changer une
seule de ces coordonnées ne définit pas le même foncteur.

\subsection{Trace graduée du générateur ternaire}

Soit
\[
 \rho_W:=\bigoplus_{b=0}^{11}g_b
\]
et considérons l'enveloppe bosonique
\begin{equation}
 \mathcal H_{\rm bos}(W_c)
 :=
 \operatorname{Sym}\!\left(
 \bigoplus_{n\ge1}W_c[n]
 \right).
 \label{p12-83-c20-s8:eq:fock-bosonico-doce-fibras-u031}
\end{equation}
Si \(N\) est l'opérateur de degré, nous définissons
\begin{equation}
 Z_{\rm Fock}(q)
 :=
 q^{-1}\operatorname{Tr}_{\mathcal H_{\rm bos}(W_c)}
 \left(\Gamma_{\rm bos}(\rho_W)q^N\right).
 \label{p12-83-c20-s8:eq:traza-graduada-fock-u031}
\end{equation}

\begin{proposition}[Produit cyclotomique de niveau trois]
\label{p12-83-c20-s8:prop:producto-tres-doce-fibras-u031}
Comme série formelle de Laurent,
\begin{equation}
 \boxed{
 Z_{\rm Fock}(q)
 =q^{-1}\prod_{n\ge1}(1+q^n+q^{2n})^{-12}
 =:t_3(q).}
 \label{p12-83-c20-s8:eq:t3-producto-u031}
\end{equation}
Son développement commence par
\begin{equation}
 t_3(q)
 =q^{-1}-12+54q-76q^2-243q^3+1188q^4+O(q^5).
 \label{p12-83-c20-s8:eq:t3-expansion-u031}
\end{equation}
\end{proposition}

\begin{proof}
Chaque \(g_b\) est conjugué à \(C\), de sorte que
\[
 \det(I-g_bq^n)=1+q^n+q^{2n}.
\]
L’identité
\[
 \sum_{k\ge0}
 \operatorname{Tr}(\operatorname{Sym}^k\rho_W)t^k
 =\det(I-\rho_Wt)^{-1}
\]
appliquée à chaque degré produit le produit. La multiplication formelle jusqu'au
degré cinq donne le développement. En particulier, le coefficient de \(q\)
s'obtient comme
\[
 [q^2]\,
 (1+q+q^2)^{-12}(1+q^2+q^4)^{-12}
 =66-12=54.
\]
\end{proof}

La valeur \(54\) apparaît ici comme coefficient d'une trace construite à partir de
douze fibres. L'unité U030 l'a également obtenue comme indice orienté de
l'agrégat APP et a prouvé que douze est l'unique solution positive de la rigidité
cyclotomique--radiale. L'égalité des deux valeurs est une compatibilité de
deux réalisations ; aucune n'utilise le coefficient de \(J\) comme entrée.

\subsection{Mot de support total et isométrie sans points fixes}

Soit
\[
 q_*=(1,1,1,1,1,1)\in\mathbb F_3^6.
\]
La multiplication par la matrice de Paley--Witt donne
\begin{equation}
 q_*A_W=(2,1,1,1,1,1).
 \label{p12-83-c20-s8:eq:producto-qstar-aw-u031}
\end{equation}
Par conséquent,
\begin{equation}
 c_*:=(q_*,q_*A_W)
 =(1,1,1,1,1,1,2,1,1,1,1,1)
 \in C_W
 \label{p12-83-c20-s8:eq:palabra-soporte-total-u031}
\end{equation}
et toutes ses coordonnées sont non nulles.

Pour chaque position \(b\), nous fixons
\begin{equation}
 P_b(c_*)=
 \begin{cases}
  \mathsf R,&(c_*)_b=1,\\
  I,&(c_*)_b=2,
 \end{cases}
 \qquad
 g_c:=\bigoplus_{b=0}^{11}P_b(c_*)CP_b(c_*)^{-1}.
 \label{p12-83-c20-s8:eq:isometria-orden-tres-u031}
\end{equation}

\begin{proposition}[Ordre, absence de points fixes et recollement]
\label{p12-83-c20-s8:prop:isometria-fijo-libre-pegado-u031}
L'opérateur \(g_c\) est d'ordre trois, n'a pas de vecteurs fixes non nuls et
préserve le réseau de Niemeier \(N(A_2^{12})\) construit dans U021.
\end{proposition}

\begin{proof}
Chaque bloc est conjugué au Coxeter \(C\), dont le polynôme minimal est
\(X^2+X+1\). Il est donc d'ordre trois et ne possède pas la valeur propre un. La somme
directe conserve les deux propriétés. Le Coxeter agit trivialement sur
\(A_2^*/A_2\) ; c'est pourquoi chaque bloc conserve la classe discriminante et la
réunion de classes latérales qui définit le recollement.
\end{proof}

\subsection{Préservation du voisin de Leech}

Écrivons
\[
 v:=v_{\mathcal F_*}
 =4\rho_1+\rho_2+\cdots+\rho_{12},
 \qquad
 v^2=54,
\]
pour le vecteur sélectionné par l'origine d'incidence de
\eqref{p12-83-c20-s5:eq:vector-marco-incidencial}. Soient
\[
 N:=N(A_2^{12}),
 \qquad
 N_0:=\ker(\langle\,\cdot\,,v\rangle\bmod3),
 \qquad
 N_v:=N_0+\mathbb Z\frac v3\cong\Lambda_{24}.
\]

\begin{theorem}[Préservation orientée du voisin]
\label{p12-83-c20-s8:thm:preservacion-vecino-orden-tres-u031}
Les déplacements
\begin{equation}
 y_*:=\frac{g_cv-v}{3},
 \qquad
 z_*:=\frac{g_c^{-1}v-v}{3}
 \label{p12-83-c20-s8:eq:desplazamientos-vecino-u031}
\end{equation}
appartiennent à \(N_0\). En conséquence,
\[
 g_cN_0=N_0,
 \qquad
 g_cN_v=N_v.
\]
\end{theorem}

\begin{proof}
Les mots discriminants de \(y_*\) et \(z_*\) sont \(c_*\) et
\(-c_*\), tous deux dans \(C_W\) ; d'où \(y_*,z_*\in N\). Le calcul dans chaque
facteur \(A_2\) donne
\begin{equation}
 \langle y_*,v\rangle
 =\langle z_*,v\rangle
 =-(4^2+11)=-27\equiv0\pmod3.
 \label{p12-83-c20-s8:eq:emparejamientos-desplazamientos-u031}
\end{equation}
Ils sont donc tous deux dans \(N_0\). Pour \(x\in N_0\),
\[
 \langle g_cx,v\rangle
 =\langle x,g_c^{-1}v\rangle
 =\langle x,v\rangle+3\langle x,z_*\rangle
 \equiv0\pmod3.
\]
Le même calcul avec \(g_c^{-1}\) donne l'égalité du noyau. Enfin,
\[
 g_c(v/3)=v/3+y_*,
\]
de sorte que la classe qui engendre le voisin est préservée.
\end{proof}

\subsection{\(24\) orientations et naturalité des réseaux}

L’énumérateur de poids de \eqref{p12-83-c20-s1:eq:enumerador-pesos-cw} contient
\(24y^{12}\). Par conséquent, l’ensemble
\[
 C_W^\times:=
 \{u\in C_W:u_b\in\{1,2\}\text{ pour tout }b\}
\]
possède 24 éléments. Pour chaque \(u\in C_W^\times\), nous définissons
\[
 P_b(u)=
 \begin{cases}
  \mathsf R,&u_b=1,\\
  I,&u_b=2,
 \end{cases}
 \qquad
 g_u:=\bigoplus_{b=0}^{11}P_b(u)CP_b(u)^{-1}.
\]

\begin{proposition}[Robustesse des orientations à support total]
\label{p12-83-c20-s8:prop:robustez-orientaciones-u031}
Pour toute \(u\in C_W^\times\),
\[
 \frac{g_uv-v}{3},
 \quad
 \frac{g_u^{-1}v-v}{3}
 \in N_0.
\]
En particulier, chaque \(g_u\) préserve le voisin construit avec
l’orientation transportée.
\end{proposition}

\begin{proof}
En coordonnées de racines simples,
\[
 C\rho-\rho=(-2,-1),
 \qquad
 C^{-1}\rho-\rho=(-1,-2).
\]
Les classes discriminantes des deux déplacements sont \(u\) et \(-u\).
Pour le coefficient radial \(a_b\in\{4,1\}\),
\[
 \left\langle
 \frac{(C^{\pm1}-I)a_b\rho}{3},a_b\rho
 \right\rangle=-a_b^2.
\]
La somme sur les douze facteurs donne de nouveau \(-27\). La preuve du
\cref{p12-83-c20-s8:thm:preservacion-vecino-orden-tres-u031} s’applique composante par
composante.
\end{proof}

Pour typer le recalibrage, soit
\(h\in\operatorname{Aut}(C_W)\) monomial : une permutation \(\sigma\) des
douze coordonnées suivie de multiplicateurs
\(\epsilon_b\in\{1,2\}\). Nous définissons son relèvement réticulaire
\begin{equation}
 \widetilde h
 :=P_\sigma\circ
 \bigoplus_{b=0}^{11}\mathsf R^{\nu_b},
 \qquad
 \nu_b=
 \begin{cases}
 0,&\epsilon_b=1,\\
 1,&\epsilon_b=2.
 \end{cases}
 \label{p12-83-c20-s8:eq:elevacion-monomial-u031}
\end{equation}
Son action sur le discriminant est exactement \(h\). Si le changement transporte
simultanément le code, le drapeau, l’origine incidencielle, le vecteur \(v\), le voisin et les
repères, alors
\begin{equation}
 g_{hu}=\widetilde h\,g_u\,\widetilde h^{-1}.
 \label{p12-83-c20-s8:eq:naturalidad-isometrias-u031}
\end{equation}
La matrice ternaire \(h\) n’agit sans élévation ni sur \(A_2^{12}\) ni sur
Leech.

\subsection{Trace de réseau et calcul du type \(0\)}

Soit \(\widehat g_c\) le relèvement standard de \(g_c\) à la VOA de réseau
\[
 V_{N_v}=M(1)\otimes\mathbb C_\varepsilon[N_v].
\]
Pour \(j=1,2\), la décomposition oscillateur--réseau produit
\begin{equation}
 \begin{split}
 Z_{0,j}(\tau)
 &:=
 \operatorname{Tr}_{V_{N_v}}
 \left(\widehat g_c^{\,j}q^{L_0-1}\right)\\
 &=q^{-1}
 \left(
 \sum_{\lambda\in N_v^{g_c^j}}
 \varepsilon_j(\lambda)
 q^{\langle\lambda,\lambda\rangle/2}
 \right)
 \prod_{n\ge1}\det(I-g_c^jq^n)^{-1}.
 \end{split}
 \label{p12-83-c20-s8:eq:factorizacion-traza-reticular-u031}
\end{equation}

\begin{proposition}[Trace des secteurs non triviaux]
\label{p12-83-c20-s8:prop:traza-reticular-t3-u031}
On a
\[
 Z_{0,1}(\tau)=Z_{0,2}(\tau)=t_3(\tau).
\]
\end{proposition}

\begin{proof}
L'absence de points fixes implique
\[
 N_v^{g_c}=N_v^{g_c^2}=\{0\},
\]
de sorte que la somme de réseau de
\eqref{p12-83-c20-s8:eq:factorizacion-traza-reticular-u031} vaut un. Dans chaque facteur,
\(g_c^j\) a pour valeurs propres \(\zeta_3,\zeta_3^2\) ; par conséquent,
\[
 \det(I-g_c^jq^n)=1+q^n+q^{2n}.
\]
Les douze facteurs reproduisent \eqref{p12-83-c20-s8:eq:t3-producto-u031}.
\end{proof}

\begin{theorem}[Poids tordu et type du relèvement]
\label{p12-83-c20-s8:thm:peso-torcido-tipo-cero-u031}
Le relèvement standard est d'ordre trois, de poids tordu
\begin{equation}
 \rho_{g_c}=\frac43
 \label{p12-83-c20-s8:eq:peso-torcido-u031}
\end{equation}
et de type zéro.
\end{theorem}

\begin{proof}
Dans
\(\mathfrak h=N_v\otimes_{\mathbb Z}\mathbb C\), chacun des douze
facteurs apporte un vecteur propre de valeur propre \(\zeta_3\) et un autre de valeur propre
\(\zeta_3^2\). Par conséquent,
\[
 \dim\mathfrak h_{\zeta_3}
 =\dim\mathfrak h_{\zeta_3^2}=12.
\]
La formule du poids du vide tordu d'ordre trois donne
\[
 \rho_{g_c}
 =\frac1{4\cdot3^2}
 \bigl(1\cdot2\cdot12+2\cdot1\cdot12\bigr)
 =\frac43.
\]
Pour une isométrie de réseau d'ordre impair, le relèvement standard conserve
l'ordre. Le type est la classe de
\[
 3^2\rho_{g_c}=12\equiv0\pmod3.
\]
Ainsi, les deux hypothèses initialement indépendantes sont calculées ici.
\end{proof}

\subsection{Orbifold cyclique et classe duale \texorpdfstring{\(3B\)}{3B}}

\begin{theorem}[Orbifold triadique du voisin de Leech]
\label{p12-83-c20-s8:thm:orbifold-triadico-vnatural-u031}
Les théorèmes classiques d'orbifold cyclique, appliqués au relèvement standard
du \cref{p12-83-c20-s8:thm:peso-torcido-tipo-cero-u031}, donnent
\begin{equation}
 \operatorname{Orb}_{\mathbb Z/3}
 (V_{N_v},\widehat g_c)
 \cong V^\natural.
 \label{p12-83-c20-s8:eq:orbifold-triadico-vnatural-u031}
\end{equation}
L'automorphisme dual appartient à la classe \(3B\) et sa série est
\begin{equation}
 T_{3B}(\tau)=t_3(\tau)+12.
 \label{p12-83-c20-s8:eq:serie-3b-u031}
\end{equation}
\end{theorem}

\begin{proof}
L'isométrie est sans points fixes et d'ordre trois ; le poids tordu et le type zéro ont été
calculés, non supposés. Les théorèmes d'orbifold cyclique identifient
l'extension holomorphe au module lunaire et classent l'automorphisme dual
comme \(3B\)
\cite{ChenLamShimakura2016,AbeLamYamada2017,Carnahan2017}.
L'identification utilise ces théorèmes après la construction de l'objet de réseau ;
elle ne se déduit pas uniquement d'une égalité de séries.
\end{proof}

Si \(V^{(i)}\) est le module \(\widehat g_c^{\,i}\)-tordu et
\[
 Z_{i,j}(\tau):=
 \operatorname{Tr}_{V^{(i)}}\!
 \left(
 \phi_i(\widehat g_c)^j q^{L_0-1}
 \right),
\]
le caractère de l'orbifold est la moyenne des neuf secteurs :
\begin{equation}
 \operatorname{ch}_{V^\natural}(\tau)
 =\frac13\sum_{i,j\in\mathbb Z/3}Z_{i,j}(\tau)
 =J(\tau).
 \label{p12-83-c20-s8:eq:promedio-nueve-sectores-u031}
\end{equation}
L'identité rationnelle de niveau trois est
\begin{equation}
 J
 =t_3+12+\frac{10\,3^9}{t_3}
 +\frac{4\,3^{14}}{t_3^2}
 +\frac{3^{18}}{t_3^3}.
 \label{p12-83-c20-s8:eq:j-desde-t3-u031}
\end{equation}
Comme \(t_3^{-1}=q+O(q^2)\),
\begin{equation}
 [q]J=54+10\cdot3^9=196884
 =54(5\cdot729+1).
 \label{p12-83-c20-s8:eq:coeficiente-196884-u031}
\end{equation}
C'est une compatibilité graduée au sein du corridor construit, non un
morphisme scalaire \(54\to V^\natural\).

\subsection{Construction de réseau d'ordre \(2\)}

La route d'ordre trois précédente et la construction classique de
Frenkel--Lepowsky--Meurman produisent deux présentations de la même VOA lunaire.
Pour déclarer la seconde, nous fixons l'extension centrale
\begin{equation}
 1\longrightarrow\{\pm1\}
 \longrightarrow\widehat\Lambda_{24}
 \longrightarrow\Lambda_{24}
 \longrightarrow1
 \label{p12-83-c20-s8:eq:extension-central-leech-u031}
\end{equation}
de commutateur
\((-1)^{\langle\lambda,\mu\rangle}\), et un cocycle
\(\varepsilon(\lambda,\mu)\) qui la représente. L'algèbre de groupe
tordue est \(\mathbb C_\varepsilon[\Lambda_{24}]\), et
\begin{equation}
 V_\Lambda
 :=M(1)\otimes\mathbb C_\varepsilon[\Lambda_{24}]
 \label{p12-83-c20-s8:eq:voa-reticular-leech-u031}
\end{equation}
est la VOA de réseau de charge centrale 24.

La négation \(\lambda\mapsto-\lambda\) possède un relèvement
\(\theta\) compatible avec \eqref{p12-83-c20-s8:eq:extension-central-leech-u031}. Sa trace
est
\begin{equation}
 t_2(\tau)
 :=
 \operatorname{Tr}_{V_\Lambda}
 (\theta q^{L_0-1})
 =q^{-1}\prod_{n\ge1}(1+q^n)^{-24}.
 \label{p12-83-c20-s8:eq:t2-leech-u031}
\end{equation}
En effet, la négation apparie \(e^\lambda\) avec \(e^{-\lambda}\), et
l'absence de torsion ne laisse que \(\lambda=0\) dans la trace de réseau ; chacun
des 24 oscillateurs apporte \((1+q^n)^{-1}\).

Soit \(V_\Lambda^T\) le module irréductible tordu par \(\theta\), et
\((V_\Lambda^T)^+\) son secteur de signe positif. La construction FLM est
\begin{equation}
 V^\natural
 =(V_\Lambda)^+\oplus(V_\Lambda^T)^+.
 \label{p12-83-c20-s8:eq:vnatural-flm-u031}
\end{equation}
La route \(\mathbb Z/2\) de
\eqref{p12-83-c20-s8:eq:vnatural-flm-u031} et la route \(\mathbb Z/3\) de
\eqref{p12-83-c20-s8:eq:orbifold-triadico-vnatural-u031} sont des constructions parallèles
identifiées par des théorèmes classiques ; aucune ne se substitue à l'autre.

\subsection{Groupe d'automorphismes, caractère et séries graduées}

Une fois \(V^\natural\) construit, deux publications distinctes apparaissent :
\begin{align}
 \operatorname{Aut}(V^\natural)&\cong\mathbb M,
 \label{p12-83-c20-s8:eq:aut-vnatural-monstruo-u031}\\
 \operatorname{ch}_{V^\natural}(\tau)
 &=J(\tau)=j(\tau)-744
 =q^{-1}+196884q+21493760q^2+\cdots.
 \label{p12-83-c20-s8:eq:caracter-vnatural-j-u031}
\end{align}
La première identifie le groupe de symétries ; la seconde compte les dimensions
graduées. La droite de vide est
\[
 (V^\natural)_0=\mathbb C\mathbf1,
\]
et l'absence de racines de \(\Lambda_{24}\) donne
\[
 (V^\natural)_1=0.
\]
En degré deux,
\[
 196884=1+196883=196560+324=54(5\cdot729+1).
\]
Seule la première égalité est une décomposition en dimensions de
représentations du Monstre ; les autres sont des contrôles arithmétiques.

Pour \(g\in\mathbb M\), on définit la série de McKay--Thompson
\begin{equation}
 T_g(\tau):=
 \operatorname{Tr}_{V^\natural}
 \left(gq^{L_0-1}\right).
 \label{p12-83-c20-s8:eq:serie-mckay-thompson-u031}
\end{equation}

\begin{theorem}[Moonshine monstrueux]
\label{p12-83-c20-s8:thm:moonshine-monstruoso-u031}
L'action graduée de \(\mathbb M\) sur \(V^\natural\) produit la famille
\(\{T_g\}_{g\in\mathbb M}\). Les théorèmes de Moonshine identifient ces
séries à des fonctions modulaires de genre zéro et prouvent leurs identités de
réplicabilité
\cite{ConwayNorton1979,Borcherds1992,FLM1988}.
\end{theorem}

Ce théorème a pour entrée le couple
\[
 \bigl(V^\natural,\operatorname{Aut}(V^\natural)\bigr)
\]
et sa graduation. Ce n'est pas la conséquence d'une flèche dont le domaine serait le
groupe abstrait \(\mathbb M\).

Pour \(J=T_{1A}\), soient \(P_n\) les polynômes de Faber normalisés par
\[
 P_n(J)=q^{-n}+O(q).
\]
L'identité de Hecke--Faber est
\begin{equation}
 P_n(J)=n\,T_nJ.
 \label{p12-83-c20-s8:eq:hecke-faber-u031}
\end{equation}
Elle contrôle tous les degrés ; le coefficient 196884 n'est que la première
manifestation non triviale.

\subsection{Graphe de dépendances du corridor Paley–Witt–Golay–Leech–Moonshine}

\begin{figure}[htbp]
\centering
\begin{tikzcd}[column sep=huge,row sep=large]
 &C_W
   \arrow[dl,"W_{12}"']
   \arrow[d,"N(A_2^{12})"',font=\scriptsize]
   \arrow[dr,phantom,"\substack{\text{branches}\\\text{distinctes}}",
          description,pos=.72,font=\scriptsize]&\\
 M_{12}
 &\Lambda_{24}
   \arrow[d,"V_\Lambda"]
 &(\mathcal G_{24},D,\ell)
   \arrow[d,"W_{24}"]\\
 &V^\natural
   \arrow[dl,"\operatorname{Aut}"']
   \arrow[d,"\operatorname{ch}"]
   \arrow[dr,"g\mapsto T_g"]
 &M_{24}\\
 \mathbb M
 &J
 &\{T_g\}_{g\in\mathbb M}
\end{tikzcd}
\caption{Bifurcation exceptionnelle et trois publications de
\(V^\natural\). L'entrée binaire \((\mathcal G_{24},D,\ell)\) est
indépendante ; il n'existe pas de flèche \(W_{24}\to\Lambda_{24}\).}
\label{p12-83-c20-s8:fig:grafo-excepcional-moonshine-u031}
\end{figure}

\begin{criterion}[Critères de réfutation du corridor lunaire]
\label{p12-83-c20-s8:crit:refutacion-corredor-lunar-u031}
La construction est réfutée si l'un des faits suivants se produit :
\begin{enumerate}
 \item \(\lambda\) n'est pas bijective ou
       \eqref{p12-83-c20-s8:eq:identidad-entrelazamiento-app-witt-u031} échoue ;
 \item \(c_*\notin C_W\), \(g_c\) possède des points fixes ou ne préserve pas le
       voisin ;
 \item les déplacements de
       \eqref{p12-83-c20-s8:eq:desplazamientos-vecino-u031} n'appartiennent pas à \(N_0\) ;
 \item le relèvement de réseau \(\widetilde h\) ne réalise pas l'automorphisme
       monomial \(h\) dans le discriminant ;
 \item le poids tordu n'est pas \(4/3\) ou le relèvement n'est pas de type zéro ;
 \item l'orbifold est identifié à \(V^\natural\) sans vérification des
       hypothèses du théorème classique ;
 \item on utilise \(W_{24}\) comme constructeur de Leech ;
 \item une égalité de dimensions est présentée comme construction de
       \(\mathbb M\) ;
 \item on sérialise indûment les publications
       \(\operatorname{Aut}\), \(\operatorname{ch}\) et \(g\mapsto T_g\).
\end{enumerate}
\end{criterion}
 \section{Contrôle thêta du réseau de Leech : identité \texorpdfstring{\(\Theta_{\Lambda}=E_4^3-720\Delta\)}{Theta Lambda = E4^3 - 720 Delta} et séparation entre coïncidence arithmétique et morphisme structurel}
\label{sec:control-theta-leech}

À partir de \(\Lambda_{24}\), la chaîne
\[
\Lambda_{24}
\longrightarrow V^\natural
\longrightarrow\mathbb M
\longrightarrow\text{Moonshine}
\]
relève des mathématiques classiques : algèbre d'opérateurs de vertex, groupe Monstre et
fonctions de McKay--Thompson. La contribution vérifiée ici se situe dans l'entrée
finie, la sélection du bord et le \(3\)-voisin sans racines, non dans une nouvelle
appropriation de ces théorèmes.

Trois identités numériques accompagnent cette chaîne :
\[
\Theta_{\Lambda}=E_4^3-720\Delta,
\]
\[
196884=196560+324,
\]
\[
196884=54(3645+1)=54(5\cdot729+1).
\]
Ce sont des identités arithmétiques exactes, mais sans morphisme ni décomposition
de caractères, elles ne constituent pas à elles seules des signes mathématiques de
compatibilité. Aucune ne remplace la construction de \(V^\natural\) ni
ne fournit une action directe du Monstre sur les mots décimaux HMT.
 \section{Uniformisation d'ordre \texorpdfstring{\(2\)}{2} dans le Moonshine : fonctions \texorpdfstring{\(T_{2A}\)}{T2A} et \texorpdfstring{\(j\)}{j}, structure modulaire et contrôle du corridor micro--macro}
\label{sec:voa-moonshine-orden-dos-rev11}
\index{VOA!réseau de Leech}
\index{Moonshine!classe 2A}

La branche d'incidence construit d'abord le réseau de Leech
\(\Lambda_{24}\). À partir de cette donnée, la théorie classique des algèbres
d'opérateurs de vertex fournit un prolongement canonique. Si
\(\mathfrak h=\Lambda_{24}\otimes_{\mathbb Z}\mathbb C\) et \(M(1)\) est
le module de Fock de son algèbre de Heisenberg, le VOA de réseau est
\begin{equation}
 V_{\Lambda}=M(1)\otimes\mathbb C[\Lambda_{24}],
 \label{eq:voa-reticular-leech-rev11}
\end{equation}
de charge centrale \(24\). Pour \(\alpha\in\Lambda_{24}\), ses opérateurs
de vertex s'écrivent
\[
 Y(e^\alpha,z)=E^-(-\alpha,z)E^+(-\alpha,z)e^\alpha z^{\alpha_0}.
\]
L'absence de racines dans \(\Lambda_{24}\) implique l'annulation du secteur
de poids un dans l'orbifold de Moonshine.

L'involution
\[
 \theta:\Lambda_{24}\longrightarrow\Lambda_{24},
 \qquad \theta(\lambda)=-\lambda,
\]
est intrinsèque et se prolonge à \(V_{\Lambda}\). Sa trace graduée s'obtient
par comptage direct des modes oscillateurs.

\begin{theorem}[Trace de l'involution de négation]
\label{thm:traza-theta-leech-rev11}
\begin{equation}
 t_2(\tau)
 :=\operatorname{Tr}_{V_\Lambda}
 \bigl(\theta q^{L_0-1}\bigr)
 =q^{-1}\prod_{n\geq1}(1+q^n)^{-24}.
 \label{eq:t2-producto-leech-rev11}
\end{equation}
\end{theorem}
\begin{proof}
L'involution apparie \(e^\lambda\) et \(e^{-\lambda}\) ; comme le réseau
est sans torsion, seul \(\lambda=0\) contribue à la trace du secteur
réticulaire. Chacun des vingt-quatre oscillateurs d'énergie \(n\) apporte
\(\sum_{k\geq0}(-1)^kq^{nk}=(1+q^n)^{-1}\). Le décalage du vide
produit le facteur \(q^{-1}\).
\end{proof}

La complétion de Fricke associée à cette trace est
\begin{equation}
 T_{2A}(\tau)=t_2(\tau)+24+\frac{2^{12}}{t_2(\tau)}
 =q^{-1}+4372q+96256q^2+1240002q^3+\cdots.
 \label{eq:T2A-hmt-rev11}
\end{equation}
Le terme \(24\) annule le terme constant de \(t_2\) ; le facteur
\(2^{12}\) est la dégénérescence du secteur tordu de vingt-quatre bosons.
La même série de bord détermine l'invariant modulaire au moyen de
\begin{equation}
 \boxed{j(\tau)=\frac{(t_2(\tau)+256)^3}{t_2(\tau)^2}}.
 \label{eq:j-desde-t2-rev11}
\end{equation}
Les équations \eqref{eq:t2-producto-leech-rev11}--
\eqref{eq:j-desde-t2-rev11} construisent la branche canonique d'ordre deux sans
paramètres continus.

L'orbifold de Frenkel--Lepowsky--Meurman est défini par
\[
 V^\natural=(V_\Lambda)^+\oplus(V_\Lambda^T)^+.
\]
C'est un VOA holomorphe de charge centrale \(24\), qui satisfait
\((V^\natural)_1=0\), possède pour caractère
\[
 \operatorname{ch}_{V^\natural}(\tau)=j(\tau)-744
\]
et son groupe d'automorphismes est le groupe Monstre
\cite{FLM1988,ConwayNorton1979,Borcherds1992}. Dans cette chaîne, HMT fixe le
réseau de Leech et l'involution \(\theta\) ; les théorèmes classiques de VOA et
de FLM identifient le prolongement résultant à \(V^\natural\) et à son groupe
d'automorphismes. L'énoncé n'exige pas d'action du Monstre sur les
chiffres : l'évaluation archimédienne et la branche exceptionnelle descendent d'un
antécédent HMT commun et conservent des invariants différents.

La fonction voisine
\(T_{3A}=t_3+12+729/t_3\) appartient à la théorie classique des Hauptmoduln
d'ordre trois. La réalisation HMT active dans le corridor tritique est la classe
\(3B\) ; une identification à \(3A\) exige de déclarer le relèvement spécifique
qui change de classe.
 \subsection{Entrelacement APP--Witt, voisin de Leech et orbifold d'ordre \(3\)}

Les douze fibres portent les étiquettes
\[
 (i,\mu,\varepsilon)\in
 \mathbb Z/3\mathbb Z\times\{0,1\}\times\{0,1\},
 \qquad
 \lambda(i,\mu,\varepsilon)=4i+2\mu+\varepsilon\pmod{12}.
\]
L'application \(\lambda\) est bijective et, les trivialisations et les cadres étant fixés,
les entrelaceurs satisfont
\[
 g_{\lambda(i,\mu,\varepsilon)}T_{i\mu\varepsilon}
 =T_{i\mu\varepsilon}C^{(-1)^\mu}.
\]
Leur somme directe est un isomorphisme \(C_3\)-équivariant
\[
 T_\lambda:W_{\mathrm{APP}}\otimes\mathbb C
 \xrightarrow{\ \cong\ }
 \bigoplus_{b=0}^{11}(A_{2,b}\otimes\mathbb C).
\]

Soit \(N=N(A_2^{12})\), soit \(v=v_\alpha\) le vecteur radial de norme \(54\),
et définissons
\[
 N_0=\ker\bigl(\langle\,\cdot\,,v\rangle\bmod3\bigr),
 \qquad N_v=N_0+\mathbb Z\frac v3.
\]
Le calcul réticulaire vérifie que \(N_v\cong\Lambda_{24}\), que l'
isométrie \(g_c\) est d'ordre trois et ne possède aucun point fixe non nul, et que
\[
 \frac{g_cv-v}{3},\quad\frac{g_c^{-1}v-v}{3}\in N_0.
\]
Par conséquent, \(g_cN_v=N_v\). Son relèvement standard au VOA réticulaire possède un poids
tordu \(4/3\) et un type zéro. Les théorèmes classiques d'orbifold cyclique donnent
\[
 V_{(3)}:=\operatorname{Orb}_{\mathbb Z/3}
       (V_{N_v},\widehat g_c)\cong V^\natural,
 \qquad T_{3B}=t_3+12.
\]
Le caractère se retrouve par la moyenne des neuf secteurs et par l'
uniformisation de niveau trois :
\[
 \operatorname{ch}_{V_{(3)}}
 =\frac13\sum_{i,j\in\mathbb Z/3}Z_{i,j}=J,
\]
\[
 J=t_3+12+\frac{10\,3^9}{t_3}
       +\frac{4\,3^{14}}{t_3^2}+\frac{3^{18}}{t_3^3}.
\]
En particulier,
\[
 [q]J=54+10\cdot3^9=196884=54(5\cdot729+1).
\]

\subsection{Constructions d'ordres \(2\) et \(3\) de l'algèbre d'opérateurs de vertex \texorpdfstring{\(V^\natural\)}{V natural}}

Fixons l'extension centrale et le cocycle de Leech
\[
 1\longrightarrow\{\pm1\}\longrightarrow\widehat\Lambda_{24}
 \longrightarrow\Lambda_{24}\longrightarrow1,
 \qquad
 V_\Lambda=M(1)\otimes\mathbb C_\varepsilon[\Lambda_{24}].
\]
La négation \(\lambda\mapsto-\lambda\) se relève en \(\theta\), et
\[
 t_2(\tau)=\operatorname{Tr}_{V_\Lambda}
  \left(\theta q^{L_0-1}\right)
 =q^{-1}\prod_{n\ge1}(1+q^n)^{-24}.
\]
La présentation FLM est
\[
 V_{(2)}:=(V_\Lambda)^+\oplus(V_\Lambda^T)^+\cong V^\natural.
\]

\begin{theorem}[Équivalence typée des présentations d'ordres \(3\) et \(2\)]
\label{thm:cierre-vtm-equivalencia-voa-23}
Soient
\[
 \phi_3:V_{(3)}\xrightarrow{\ \cong\ }V^\natural,
 \qquad
 \phi_2:V_{(2)}\xrightarrow{\ \cong\ }V^\natural
\]
les isomorphismes gradués fournis par les théorèmes d'orbifold
appliqués aux données précédentes. Alors
\[
 \boxed{\Phi_{23}:=\phi_2^{-1}\circ\phi_3:
 V_{(3)}\xrightarrow{\ \cong\ }V_{(2)}}
\]
est un isomorphisme gradué d'algèbres d'opérateurs de vertex. De plus,
\[
 a\longmapsto\Phi_{23}a\Phi_{23}^{-1}
\]
identifie leurs groupes d'automorphismes et conserve les traces graduées.
\end{theorem}

\begin{proof}
La composition d'isomorphismes de VOA est un isomorphisme de VOA et préserve le
vide, le vecteur conforme, les opérateurs de vertex et la graduation par
\(L_0\). La conjugaison transporte les automorphismes. Pour chaque automorphisme
\(a\), l'invariance de la trace par conjugaison donne l'égalité des
séries graduées correspondantes.
\end{proof}

L'isomorphisme \(\Phi_{23}\) dépend des identifications \(\phi_2\) et
\(\phi_3\) ; aucune canonicité absolue n'est affirmée. La construction d'ordre deux
conserve en outre le contenu exclusif
\[
 T_{2A}=t_2+24+\frac{2^{12}}{t_2},
 \qquad
 \boxed{j=\frac{(t_2+256)^3}{t_2^2}},
\]
qui complète l'uniformisation modulaire d'ordre deux.

Une fois \(V^\natural\) construit, les trois publications s'énoncent sans
les sérialiser indûment :
\[
 \operatorname{Aut}(V^\natural)\cong\mathbb M,
 \qquad
 \operatorname{ch}_{V^\natural}=J,
 \qquad
 g\longmapsto T_g(\tau)
 =\operatorname{Tr}_{V^\natural}\!\left(gq^{L_0-1}\right).
\]
Le théorème de Moonshine coordonne l'action graduée avec les fonctions de
genre zéro et leurs identités de réplicabilité. L'action du Monstre vit
dans \(V^\natural\) ; le fait qu'elle ne soit pas une permutation ponctuelle de chiffres est un contrôle
final des types, non la thèse du corridor.

 \section{Chaîne probatoire continue de la publication exceptionnelle : du déploiement APP à Paley--Witt,
\texorpdfstring{\(V^\natural\)}{V natural} et Moonshine}
\label{sec:prueba-corredor-excepcional}

La publication exceptionnelle part du lecteur d'incidence
\(\mathfrak E_{\rm exc}\) de \eqref{eq:frontal-doble-proyeccion}. Elle ne prend pas en
entrée un code, un réseau ou un groupe exceptionnel déjà donnés : elle prend le
bord, l'orientation, le feuillet et la mémoire d'une histoire enrichie et
les exprime dans le calibre projectif fixé par l'holonomie nonadique. Ce n'est
qu'ensuite que les théorèmes classiques des codes, des réseaux et des algèbres
d'opérateurs de vertex sont appliqués aux objets produits.

\subsection{\texorpdfstring{Déploiement APP et cône gradué de l'indice \(54\)}{Déploiement APP et cône gradué de l'indice 54}}

Soit \(R\) l'involution qui échange les feuillets APP et soient
\(P_+=(I+R)/2\), \(P_-=(I-R)/2\) ses projecteurs. Le bloc central
\begin{equation}
 B_c=7I+2R=9P_++5P_-
 \label{eq:exc-bloque-central-app}
\end{equation}
produit quatre objets de types distincts :
\begin{equation}
 9-5=4,\qquad (B_c)_{i\ne j}=2,\qquad
 \sum M_\Pi-\sum M_\Sigma=51-45=6,\qquad
 9\cdot6=459-405=54.
 \label{eq:exc-despliegue-42654}
\end{equation}
Le \(4\) est le saut spectral entre feuillets, le \(2\) le couplage local,
le \(6\) la différence des compressions modulaires et le \(54\)
l'intégration de cette différence sur les neuf positions. Par conséquent,
\(4\to2\to6\to54\) désigne une hiérarchie de morphismes et d'évaluations, non
une chaîne d'égalités entre objets.

L'orientation \(o_{\Sigma\Pi}\), la parité de feuillet et la
normalisation triangulaire étant fixées, la branche de Fock possède le couplage
\begin{equation}
 \mathcal H_{\Sigma,\Pi}^{(m)}:
 M_{\rm bal}^{C_3,-}\longrightarrow
 \operatorname{End}\!\bigl(\mathcal F_2(\mathcal V_m)\bigr)
 \otimes o_{\Sigma\Pi},
 \label{eq:exc-acoplamiento-fock}
\end{equation}
et sa trace orientée satisfait
\begin{equation}
 \operatorname{Tr}_{o,\mathcal F_2}
 \bigl(\mathcal H_{\Sigma,\Pi}^{(m)}(\nu)\Gamma_2(g_m)\bigr)
 =-\frac{3m\,c(\nu)}2.
 \label{eq:exc-traza-fock}
\end{equation}
Pour \(m=12\) et \(c(\delta)=-3\), on a \(54\). Indépendamment, le
produit cyclotomique
\begin{equation}
 t_3(q)=q^{-1}\prod_{n\ge1}(1+q^n+q^{2n})^{-12}
 =q^{-1}-12+54q-76q^2-243q^3+\cdots
 \label{eq:exc-t3}
\end{equation}
donne \([q]t_3=54\).

\begin{theorem}[Cône typé de l'indice commun]
\label{thm:exc-cono-54}
Une fois construits séparément le recensement APP, la branche de Fock, le produit
cyclotomique, le voisin réticulaire de
\cref{eq:exc-vecino-indice-tres} et le module lunaire de
\cref{eq:exc-flm}, leurs lecteurs vérifient
\begin{equation}
 D_{\rm APP}
 =\operatorname{Tr}_{o,\mathcal F_2}
  \bigl(\mathcal H_{\Sigma,\Pi}^{(12)}(\delta)\Gamma_2(g_{12})\bigr)
 =[q]t_3=v^2=\frac{196884}{5\cdot729+1}=54.
 \label{eq:exc-cono-54}
\end{equation}
L'égalité réunit des évaluations dans un cône vers \(\mathbb Z\) ; elle
n'identifie pas leurs domaines et n'utilise pas \(196884\) pour sélectionner le corridor.
\end{theorem}

\begin{proof}
Les trois premières égalités sont
\eqref{eq:exc-despliegue-42654}, \eqref{eq:exc-traza-fock} et
\eqref{eq:exc-t3}. Pour le vecteur radial
\(v=4\rho_1+\rho_2+\cdots+\rho_{12}\) de
\eqref{eq:exc-vector-vecino}, l'orthogonalité des facteurs \(A_2\) donne
\(v^2=4^2\cdot2+11\cdot2=54\). Enfin,
\(5\cdot729+1=3646\) et \(54\cdot3646=196884\). Chaque identité est prouvée dans
son propre codomaine avant leur comparaison.
\end{proof}

\subsection{\texorpdfstring{Forme de Paley et racine quadratique interne de \(-I_6\)}{Forme de Paley et racine quadratique interne de -I6}}

Soit
\[
 \mathcal U_9=(\mathbb Z/9\mathbb Z)^\times=\{1,2,4,8,7,5\},
\]
dans l'ordre des puissances de \(2\), et soit
\begin{equation}
 \theta:\mathcal U_9\xrightarrow{\sim}\mathbb P^1(\mathbb F_5),
 \qquad
 \begin{array}{c|cccccc}
 u&1&2&4&8&7&5\\ \hline
 \theta(u)&\infty&0&1&2&3&4
 \end{array}.
 \label{eq:exc-calibre-proyectivo}
\end{equation}
Cette bijection est un calibre de coordonnées du lecteur ; elle n'est pas utilisée comme
homomorphisme. Pour le caractère quadratique
\(\chi:\mathbb F_5\to\{-1,0,1\}\), définissons
\begin{equation}
 C_P=
 \begin{pmatrix}
 0& 1& 1& 1& 1& 1\\
 1& 0& 1&-1&-1& 1\\
 1& 1& 0& 1&-1&-1\\
 1&-1& 1& 0& 1&-1\\
 1&-1&-1& 1& 0& 1\\
 1& 1&-1&-1& 1& 0
 \end{pmatrix},
 \label{eq:exc-matriz-paley}
\end{equation}
où les entrées finies hors diagonale sont \(\chi(x-y)\) et les
entrées incidentes à \(\infty\) valent un.

\begin{lemma}[Corrélation quadratique]
\label{lem:exc-correlacion-cuadratica}
Si \(a\ne b\) dans \(\mathbb F_5\), alors
\[
 \sum_{t\in\mathbb F_5}\chi(t-a)\chi(t-b)=-1.
\]
\end{lemma}

\begin{proof}
La substitution \(u=(t-a)/(b-a)\) réduit la somme à
\(\sum_u\chi(u(u-1))\). Pour \(u=0,1,2,3,4\), les termes sont
\(0,0,1,-1,-1\).
\end{proof}

\begin{proposition}[Identité de conférence et réduction ternaire]
\label{prop:exc-paley-witt}
La matrice \(C_P\) est symétrique et satisfait
\[
 C_PC_P^{\mathsf T}=5I_6.
\]
Sa réduction \(A_W=C_P\bmod3\) vérifie
\begin{equation}
 \boxed{A_W^2=-I_6\quad\text{dans }M_6(\mathbb F_3).}
 \label{eq:exc-aw-cuadrado}
\end{equation}
\end{proposition}

\begin{proof}
Chaque ligne de \(C_P\) a une norme au carré égale à cinq. Deux lignes finies
distinctes ont pour produit scalaire \(1-1=0\) par le
\cref{lem:exc-correlacion-cuadratica} ; la somme d'un caractère non trivial
prouve aussi l'orthogonalité avec la ligne de \(\infty\). En réduisant modulo
trois, on obtient
\(A_WA_W^{\mathsf T}=5I_6=2I_6=-I_6\). Comme \(A_W\) est symétrique,
\eqref{eq:exc-aw-cuadrado} en découle.
\end{proof}

L'identité \eqref{eq:exc-aw-cuadrado} construit la structure complexe
finie du corridor. En effet,
\[
 \mathbb F_9=\mathbb F_3[\iota]/(\iota^2+1),
 \qquad
 P_{\pm\iota}=\frac12(I\mp\iota A_W)
\]
sont des projecteurs complémentaires de rang trois et munissent
\(\mathbb F_3^6\) d'une structure de \(\mathbb F_9\)-module de rang trois.
L'action est explicite :
\begin{equation}
 (a+b\iota)\cdot v=av+b(vA_W).
 \label{eq:exc-accion-f9}
\end{equation}
Par conséquent, \(3^6=9^3\) n'est pas utilisé comme substitut cardinal d'une action.

\subsection{Code ternaire autodual et système de Witt}

Définissons
\begin{equation}
 c:\mathbb F_3^6\longrightarrow\mathbb F_3^{12},
 \qquad c(q)=(q,qA_W),
 \qquad C_W:=\operatorname{im}c.
 \label{eq:exc-codigo-grafico}
\end{equation}

\begin{theorem}[Code ternaire produit par l'incidence]
\label{thm:exc-codigo-ternario}
L'application \(c\) est injective et
\begin{equation}
 C_W=C_W^\perp,
 \qquad
 W_{C_W}(y)=1+264y^6+440y^9+24y^{12}.
 \label{eq:exc-enumerador-cw}
\end{equation}
En particulier, \(C_W\) a pour paramètres \([12,6,6]_3\).
\end{theorem}

\begin{proof}
La première projection de \(c(q)\) est \(q\), donc \(c\) est injective.
Avec \(G_W=(I_6\mid A_W)\),
\[
 G_WG_W^{\mathsf T}=I_6+A_WA_W^{\mathsf T}=0
\]
dans \(\mathbb F_3\). Ainsi \(C_W\subseteq C_W^\perp\), et l'égalité des
dimensions donne l'autodualité. L'énumérateur est la somme finie exacte
\[
 \sum_{q\in\mathbb F_3^6}
 y^{\operatorname{wt}(q)+\operatorname{wt}(qA_W)}.
\]
Évaluer les \(3^6=729\) lignes de la matrice explicite
\eqref{eq:exc-matriz-paley} donne, respectivement, les effectifs
\(1,264,440,24\) aux poids \(0,6,9,12\), et zéro aux autres. La somme
des quatre effectifs est \(729\), de sorte qu'aucun mot ne reste hors
du certificat fini.
\end{proof}

Les \(264\) vecteurs de poids six apparaissent par paires \(\{x,-x\}\) ; leurs
\(132\) supports forment le système de Steiner
\begin{equation}
 W_{12}=S(5,6,12).
 \label{eq:exc-w12}
\end{equation}
C'est la reconnaissance, après construction de \(C_W\), du code ternaire
étendu de Golay et de son plan de Witt
\citep{MacWilliamsSloane1977,ConwaySloane1999}. Le plan n'est pas introduit
pour sélectionner la matrice \(A_W\).

Le prolongement binaire demeure séparé. Si
\(\mathcal G_{24}\subset\mathbb F_2^{24}\) est le code binaire étendu
de Golay, \(D\) une dodécade et
\(\ell:\{1,\ldots,12\}_{\rm HMT}\to D\) un isomorphisme de plans, les
supports de poids huit forment \(W_{24}=S(5,8,24)\). Pour une tétrade fixe,
\[
 \lambda_4(W_{12})=4,
 \qquad
 \lambda_4(W_{24})=5,
\]
de sorte que quatre octades relèvent les quatre hexades incidentes et qu'une cinquième
est transversale. La donnée \((\mathcal G_{24},D,\ell)\) appartient au codomaine
binaire ; il n'existe pas de flèche canonique \(W_{24}\to\Lambda_{24}\).

\subsection{Recollement ternaire et voisin de Leech}

Soit \(Q=A_2^{12}\), avec \(\det Q=3^{12}\). Le discriminant de chaque
facteur \(A_2\) est \(\mathbb F_3\), et le recollement défini par
\eqref{eq:exc-codigo-grafico} est
\begin{equation}
 N:=\{x\in Q^*:x\bmod Q\in C_W\}.
 \label{eq:exc-niemeier-pegado}
\end{equation}

\begin{theorem}[Réseau de Niemeier produit par le code]
\label{thm:exc-niemeier}
Le réseau \(N\) est défini positif, pair et unimodulaire de rang
vingt-quatre, et son système de racines est \(A_2^{12}\). Par conséquent,
\(N\cong N(A_2^{12})\).
\end{theorem}

\begin{proof}
L'autodualité de \(C_W\) rend le recollement entier. Tous ses poids sont
des multiples de trois, ce qui rend l'extension paire. En outre,
\([N:Q]=|C_W|=3^6\), et donc
\[
 \det N=\frac{\det Q}{[N:Q]^2}
 =\frac{3^{12}}{3^{12}}=1.
\]
Une classe discriminante non nulle de \(A_2\) a pour norme minimale \(2/3\).
Comme tout mot non nul de \(C_W\) a un poids au moins égal à six, tout vecteur
de recollement extérieur à \(Q\) a une norme au moins égale à quatre. Les vecteurs de
norme deux sont exactement les racines des douze facteurs \(A_2\). La
classification des réseaux pairs unimodulaires de rang vingt-quatre
identifie le type indiqué \citep{ConwaySloane1999}.
\end{proof}

Soit \(\rho_i\) la plus haute racine du facteur \(A_2^{(i)}\), normalisée par
\(\rho_i^2=2\), et définissons
\begin{equation}
 v=4\rho_1+\rho_2+\cdots+\rho_{12},
 \qquad v^2=54.
 \label{eq:exc-vector-vecino}
\end{equation}
Le repère ordonné du lecteur exceptionnel fixe laquelle des douze fibres occupe
la première position ; aucune valeur archimédienne n'intervient dans cette sélection.
Posons
\begin{equation}
 \chi_v(x)=\langle x,v\rangle\bmod3,
 \quad N_0=\ker\chi_v,
 \quad N_v=N_0+\mathbb Z\frac v3.
 \label{eq:exc-vecino-indice-tres}
\end{equation}

\begin{lemma}[Structure du voisin]
\label{lem:exc-estructura-vecino}
Le caractère \(\chi_v\) est surjectif, \([N:N_0]=3\),
\(v/3\in N_0^*\), et \(N_v\) est pair et unimodulaire.
\end{lemma}

\begin{proof}
Pour toute racine d'un facteur \(A_2\), le produit avec la plus haute racine est
\(\pm1\) ou \(\pm2\) ; les coefficients de \(v\) sont congrus à un
modulo trois. Par conséquent, aucune racine n'appartient à \(N_0\) et \(\chi_v\) est
non trivial, donc surjectif. La définition du noyau donne l'indice trois et
\(\det N_0=9\). Pour \(x\in N_0\),
\(\langle x,v/3\rangle\in\mathbb Z\), tandis que
\((v/3)^2=6\). Ainsi, l'extension par \(v/3\) est entière et paire. Comme
\([N_v:N_0]=3\),
\(\det N_v=9/3^2=1\).
\end{proof}

\begin{lemma}[Minimum des nouvelles classes]
\label{lem:exc-minimo-clases}
Chaque coordonnée d'une classe locale apparaissant dans
\(N_0\pm v/3\) a pour norme minimale \(2/9\). En conséquence,
\[
 \|x\pm v/3\|^2\ge12\cdot\frac29=\frac83
 \qquad(x\in N_0).
\]
\end{lemma}

\begin{proof}
En écrivant une coordonnée de \(A_2^*\) dans la base des racines simples, les
six classes orientées pertinentes sont les translations de
\(\pm\rho/3\) par les trois classes discriminantes. Compléter les carrés dans la
forme de Gram
\(\bigl(\begin{smallmatrix}2&-1\\-1&2\end{smallmatrix}\bigr)\)
donne le minimum \(2/9\) dans les six classes. La somme orthogonale des douze
coordonnées donne l'inégalité.
\end{proof}

\begin{theorem}[Construction du réseau de Leech]
\label{thm:exc-leech}
Le voisin \(N_v\) de \eqref{eq:exc-vecino-indice-tres} est isométrique au
réseau de Leech :
\begin{equation}
 \boxed{N_v\cong\Lambda_{24}.}
 \label{eq:exc-leech}
\end{equation}
\end{theorem}

\begin{proof}
Les racines de \(N\) n'appartiennent pas à \(N_0\), et les deux nouvelles classes ne
contiennent pas de vecteurs de norme deux par le
\cref{lem:exc-minimo-clases}. Le \cref{lem:exc-estructura-vecino} prouve que
\(N_v\) est défini positif, pair, unimodulaire et de rang vingt-quatre.
L'unicité du réseau possédant ces propriétés et dépourvu de racines l'identifie à
\(\Lambda_{24}\) \citep{ConwaySloane1999}.
\end{proof}

\subsection{Isométrie d'ordre trois, orbifold et module lunaire}

La rotation de Coxeter de chaque facteur \(A_2\) induit une isométrie
\(g_c\) d'ordre trois. La stabilité du code \(C_W\) prouve que
\(g_cN=N\). Pour le repère orienté qui fixe \(v\), les translations
\begin{equation}
 y_*:=\frac{g_cv-v}{3},
 \qquad
 z_*:=\frac{g_c^{-1}v-v}{3}
 \label{eq:exc-desplazamientos-vecino}
\end{equation}
appartiennent à \(N_0\) : leurs mots discriminants sont un couple
\(c_*,-c_*\in C_W\), et
\(\langle y_*,v\rangle=\langle z_*,v\rangle=-27\). Par conséquent,
\begin{equation}
 g_cN_v=N_v.
 \label{eq:exc-preservacion-vecino}
\end{equation}
L'action n'a pas de points fixes dans \(N_v\otimes\mathbb R\), car chaque
facteur \(A_2\) possède les deux valeurs propres primitives d'ordre trois.

Soit \(\widehat g_c\) le relèvement standard à la VOA de réseau
\(V_{\Lambda_{24}}\). Ses deux espaces propres non triviaux sont de dimension
douze. L'énergie du vide du module tordu est donc
\begin{equation}
 \rho_{g_c}
 =\frac1{4\cdot3^2}
 \sum_{k=1}^{2}k(3-k)\dim\mathfrak h_{(k)}
 =\frac43.
 \label{eq:exc-peso-torcido}
\end{equation}
La congruence \(3^2\rho_{g_c}=12\equiv0\pmod3\) prouve que le relèvement
est de type zéro.

\begin{theorem}[Orbifold triadique]
\label{thm:exc-orbifold}
L'orbifold cyclique du relèvement précédent satisfait
\begin{equation}
 \operatorname{Orb}_{\mathbb Z/3}
 (V_{\Lambda_{24}},\widehat g_c)\cong V^\natural.
 \label{eq:exc-orbifold}
\end{equation}
L'automorphisme dual appartient à la classe \(3B\).
\end{theorem}

\begin{proof}
Les trois hypothèses qui ne peuvent être remplacées par une
égalité de séries ont été vérifiées : isométrie sans point fixe d'ordre trois, préservation du
voisin et type zéro. Les théorèmes d'orbifold cyclique identifient alors
l'extension holomorphe au module lunaire et classent l'automorphisme dual
\citep{ChenLamShimakura2016,AbeLamYamada2017,Carnahan2017}.
\end{proof}

La construction d'ordre deux est une réalisation parallèle. Pour le relèvement
\(\theta\) de la négation sur \(\Lambda_{24}\),
\begin{equation}
 V^\natural=(V_{\Lambda_{24}})^+
 \oplus(V_{\Lambda_{24}}^T)^+.
 \label{eq:exc-flm}
\end{equation}
Les voies d'ordre deux et trois sont identifiées par les théorèmes de VOA ; elles ne sont pas
mises en série causalement l'une dans l'autre \citep{FLM1988}.

\begin{theorem}[Publications du module lunaire]
\label{thm:exc-moonshine}
Une fois \(V^\natural\) construit, on obtient
\begin{align}
 \operatorname{Aut}(V^\natural)&\cong\mathbb M,
 \label{eq:exc-aut-monster}\\
 \operatorname{ch}_{V^\natural}(\tau)
 &=J(\tau)=j(\tau)-744
 =q^{-1}+196884q+21493760q^2+\cdots,
 \label{eq:exc-caracter-j}\\
 T_g(\tau)&=\operatorname{Tr}_{V^\natural}
 \bigl(gq^{L_0-1}\bigr),\qquad g\in\mathbb M.
 \label{eq:exc-mckay-thompson}
\end{align}
Les séries \(T_g\) sont les fonctions modulaires de genre zéro du théorème
de Moonshine et satisfont les identités de réplicabilité.
\end{theorem}

\begin{proof}
La première identification et la construction graduée proviennent de la
réalisation FLM. La coordination entre l'action graduée, les séries de
McKay--Thompson et leurs propriétés modulaires est le théorème de Moonshine de
Conway--Norton et de Borcherds
\citep{FLM1988,ConwayNorton1979,Borcherds1992}. En particulier,
\((V^\natural)_0=\mathbb C\mathbf1\) et
\((V^\natural)_1=0\) ; le coefficient \(196884=1+196883\) est la première
décomposition non triviale en dimensions de représentations du Monstre.
\end{proof}

\begin{falsifier}[Corridor exceptionnel]
\label{fals:exc-corredor}
La publication est réfutée si l'une des étapes effectives échoue :
\eqref{eq:exc-despliegue-42654}, la trace
\eqref{eq:exc-traza-fock}, le coefficient \([q]t_3=54\),
\(C_PC_P^{\mathsf T}=5I_6\), \(A_W^2=-I_6\), l'autodualité ou
l'énumérateur de \(C_W\), la parité ou l'unimodularité du recollement, l'absence
de racines du voisin, sa préservation par \(g_c\), le poids tordu
\(4/3\) ou le type zéro. Est également réfuté un récit qui utilise
\(W_{24}\) pour construire \(\Lambda_{24}\), qui confond les deux ensembles
de \(243\) par leur cardinal, ou qui fait agir le Monstre directement
sur les chiffres et les voies. L'action de \(\mathbb M\) réside dans \(V^\natural\).
\end{falsifier}
 
\chapter{Image spectrale--polaire et réduction exacte de Riemann--Weil}
\begingroup
\renewcommand{\chapter}[2][]{}%
\chapter{Image spectrale--polaire et réduction exacte de Riemann--Weil}

\section{Point de départ : la structure spectrale--polaire complète}

L'entrée de la résolution n'est ni la forme de Weil isolée ni une liste de zéros. C'est la structure spectrale--polaire complète reproduite comme donnée dans ce volume. La chaîne causale propre à ce chapitre est
\begin{equation}\label{eq:pdf2-rh-cadena-procedencia}
 \mathfrak G_{\rm sp}
 \xrightarrow{\ \operatorname{pr}_{X,{\rm sp}}\ }X_{\chi_{\rm sp}}
 \xrightarrow{\ \mathcal A_{\rm sp}\ }\mathfrak M_{\rm sp}
 \xrightarrow{\ \mathcal P_{\rm sp}\ }\mathcal C_{\rm GW}.
\end{equation}
Le codomaine classique n'apparaît que dans la dernière flèche ; il ne sélectionne pas
rétrospectivement l'état, les cartes, la projection ni le terminal.

La structure de départ contient l'objet à treize coordonnées
\begin{equation}\label{eq:pdf2-rh-res-sp-completa}
\begin{aligned}
 \mathfrak G_{\rm sp}^{(13)}=\bigl(&
 \mathcal F_{\rm sp},
 \operatorname{Ext}_{\Gamma,{\rm sp}},
 \operatorname{Rel}_{\rm sp},
 \mathcal N_{\rm sp},
 \operatorname{or}_{\rm sp},
 \operatorname{Carr}_{\rm sp},
 \operatorname{Mem}_{\rm sp},\\
 &\partial_{\rm sp},
 \gamma_{\rm sp},
 \Sigma_{\rm sp}^{\rm mult},
 \operatorname{Surv}_{\rm sp},
 \tau_{\rm sp}^{\rm term},
 \mathcal L_{\rm sp}\bigr).
\end{aligned}
\end{equation}
Chaque coordonnée intervient matériellement :
\begin{enumerate}
 \item \(\mathcal F_{\rm sp}\) conserve les deux feuillets, leur orientation, la
 paire résidu--quotient et la signature de la publication double.
 \item \(\operatorname{Ext}_{\Gamma,{\rm sp}}\) transporte cette fibre par la
 classe dirigée de coupures et de niveaux nonadiques.
    \item \(\operatorname{Rel}_{\rm sp}\) contient le cocycle de retenue,
    l'orthogonalité \(P=P^*=P^2\), les deux publications et le complément
    structurel \(B_{\rm sp}^{\rm str}:=\Xi^*(I-P)\Xi\). L'identification
    de ce complément à la forme de Guinand--Weil est une flèche distincte
    et fait l'objet d'un audit séparé dans ce chapitre.
 \item \(\mathcal N_{\rm sp}\) enregistre le niveau \(m\), la fenêtre \(R\),
 le quotient \(q=5/9\), les moments des deux feuillets et les indices des
 canaux premier, gamma et polaire.
 \item \(\operatorname{or}_{\rm sp}=\widehat R=(R,K)\) maintient distinctes
 les deux cartes orientées.
 \item \(\operatorname{Carr}_{\rm sp}=\operatorname{Carr}_R\) est la
 retenue opératorielle et ne s'annule pas lors de la publication d'une translation.
 \item \(\operatorname{Mem}_{\rm sp}=(S_0,M,\mathcal O_9)\) conserve
 l'identité finie de Stein et son terminal.
 \item \(\partial_{\rm sp}=(I-P)\Xi\) est la frontière qui n'apparaît pas dans la
 compression \(P\Xi\) et dont l'énergie définit le complément structurel. Son
 identification à la forme de Weil n'est pas incorporée à cette définition.
 \item \(\gamma_{\rm sp}\) est le chemin cofinal de fenêtres, avec un unique
 régulateur pour tous les canaux.
 \item \(\Sigma_{\rm sp}^{\rm mult}\) est la multisection des deux
 publications et de leurs composantes premier--gamma--polaire ; ce n'est pas la signature d'un
 unique état.
 \item \(\operatorname{Surv}_{\rm sp}\) conserve la famille sous
 \(m\mapsto m+1\) et \(\mathcal D_R\hookrightarrow\mathcal D_{R'}\).
 \item \(\tau_{\rm sp}^{\rm term}\) conserve le terminal
 \(M^{*N}S_0M^N\), qui, dans la carte miroir, se publie au moyen de \(E_R\).
    \item \(\mathcal L_{\rm sp}=(\Xi,P,B_{\rm sp}^{\rm str})\) est le lecteur final ; il s'applique
 après antécédent, cartes, mémoire, chemin et terminal.
\end{enumerate}

L'assembleur et le publicateur sont
\begin{equation}\label{eq:pdf2-rh-ensamblador-sp}
\begin{aligned}
 \mathcal A_{\rm sp}(\operatorname{pr}_{X,{\rm sp}}\widehat x)
  ={}&(\widehat X_\infty,\Gamma_9,\widehat R,
      \operatorname{Carr}_R,S_0,M,\mathcal O_9,
      \Xi,P,\mathcal D_R^{\rm cof}),\\
 \mathfrak M_{\rm sp}={}&\operatorname{Im}\mathcal A_{\rm sp},\\
 \mathcal P_{\rm sp}(\mathfrak M_{\rm sp})
  ={}&(B_{\rm sp}^{\rm str},\Xi^*(I-P)\Xi,
       \mathscr I_{\rm GW},\Delta_{\rm bind},
       \text{famille cofinale premier--gamma--polaire}),\\
 \Delta_{\rm bind}:={}&B_{\rm sp}^{\rm str}-\mathscr I_{\rm GW}.
\end{aligned}
\end{equation}
La mémoire finie de l'objet de départ satisfait
\begin{equation}\label{eq:pdf2-rh-stein-finito}
 S_0=\sum_{r=0}^{N-1}M^{*r}\mathcal O_9^*\mathcal O_9M^r
     +M^{*N}S_0M^N.
\end{equation}
Le dernier terme reste visible jusqu'à ce que son extinction soit démontrée. Par
conséquent, remplacer \(\mathfrak G_{\rm sp}\) par l'ombre publiée
\(B_{\rm sp}^{\rm str}\) change l'objet de départ ; ce n'est pas une abréviation
du même théorème. Le résiduel \(\Delta_{\rm bind}\) est conservé jusqu'à ce que la
colligation source--first démontre matériellement son annulation dans
\eqref{eq:pdf2-rh-binding-source-first}.

\section{Domaine cofinal et forme cible}

Pour $R>0$, on fixe
\[
 \mathcal D_R=C_c^\infty((-R/2,R/2)),\qquad
 \widehat\psi(t)=\int_{\mathbb R}\psi(x)e^{-itx}\,dx,
\]
et, pour $\psi,\phi\in\mathcal D_R$,
\[
 h_{\psi,\phi}(z)=\widehat\psi(z)
 \overline{\widehat\phi(\overline z)},\qquad
 \check h_{\psi,\phi}(a)=\langle\psi,T_a\phi\rangle.
\]
L'union $\bigcup_R\mathcal D_R=C_c^\infty(\mathbb R)$ est exacte. Pour tout
ensemble fini de points complexes distincts, les évaluations
$\psi\mapsto(\widehat\psi(z_1),\ldots,\widehat\psi(z_n))$ sont conjointement
surjectives lorsque $R$ est suffisamment grand. Par conséquent, cette classe cofinale
sépare les évaluations requises par le critère de Weil.

\begin{lemma}[Preuve de la séparation conjointe]
\label{lem:pdf2-rh-separacion}
L'application
\[
 \mathcal D_R\longrightarrow\mathbb C^n,
 \qquad
 \psi\longmapsto
 (\widehat\psi(z_1),\ldots,\widehat\psi(z_n))
\]
est surjective pour tout ensemble de points complexes distincts
\(z_1,\ldots,z_n\), dès que \(R\) contient un intervalle non vide.
\end{lemma}

\begin{proof}
Si les fonctionnelles d'évaluation étaient linéairement dépendantes, il existerait
\(c_1,\ldots,c_n\), non tous nuls, tels que
\[
 \int_{\mathbb R}\psi(x)
 \left(\sum_{j=1}^nc_je^{-iz_jx}\right)\,dx=0
 \qquad(\psi\in\mathcal D_R).
\]
La fonction analytique \(\sum_jc_je^{-iz_jx}\) s'annulerait dans
\((-R/2,R/2)\), donc dans tout le plan. L'indépendance des exponentielles
d'exposants distincts impose \(c_j=0\) pour tout \(j\), contradiction.
Les fonctionnelles sont indépendantes ; comme le codomaine est de dimension finie,
l'application est de rang \(\mathbb C^n\).
\end{proof}

\section{Les deux colonnes et les trois canaux calculés avant le signe}

Soit $\vartheta_L$ un régulateur commun. Avec
\[
 \ell_{p,k}=k\log p,\qquad w_{p,k}=(\log p)p^{-k/2},\qquad
 \mu_\Gamma(a)=\frac{e^{-a/2}}{1-e^{-2a}},
\]
on utilise l'inclusion uniforme $j_m$ et le canal de raccordement
$\widehat\delta_{a,m}$, qui satisfont
\[
 j_m^*j_m=I,\qquad
 \widehat\delta_{a,m}^*\widehat\delta_{b,m}
 =(I-T_a)^*(I-T_b).
\]
Les colonnes régularisées sont
\begin{align*}
 J_{+,m,R,L}\psi={}&
 \Bigl(
  (\sqrt{\vartheta_L(\ell_{p,k})w_{p,k}}\,
       \widehat\delta_{\ell_{p,k},m}\psi)_{\ell_{p,k}\leq R},\\
 &\hspace{9mm}a\mapsto
  \sqrt{\vartheta_L(a)\mu_\Gamma(a)}\,
       \widehat\delta_{a,m}\psi,\ \zeta_+b_+(\psi)
 \Bigr),\\
 J_{-,m,R,L}\psi={}&
 \Bigl(
  (\sqrt{2\vartheta_L(\ell_{p,k})w_{p,k}}\,j_m\psi)_{\ell_{p,k}\leq R},
  \sqrt{-c_\Gamma}\,j_m\psi,\ \zeta_-b_-(\psi)
 \Bigr),
\end{align*}
où $c_\Gamma=\psi_\Gamma(1/4)-\log\pi<0$ et
$b_\pm=(\widehat\psi(i/2)\pm\widehat\psi(-i/2))/\sqrt2$.

Un développement terme à terme, sans utiliser de zéros ni le signe recherché, donne
\begin{align*}
 &\langle J_{+,m,R,L}\psi,J_{+,m,R,L}\phi\rangle
 -\langle J_{-,m,R,L}\psi,J_{-,m,R,L}\phi\rangle
 =\mathscr I_{{\rm GW},L}(\psi,\phi),\\
 \mathscr I_{{\rm GW},L}(\psi,\phi)
 ={}&h_{\psi,\phi}(i/2)+h_{\psi,\phi}(-i/2)
 +c_\Gamma\check h_{\psi,\phi}(0)\\
 &+\int_0^\infty\vartheta_L(a)\mu_\Gamma(a)
 [2\check h(0)-\check h(a)-\check h(-a)]\,da\\
 &-\sum_{p,k}\vartheta_L(\ell_{p,k})w_{p,k}
 [\check h(\ell_{p,k})+\check h(-\ell_{p,k})].
\end{align*}
La troncature des deux colonnes en \(\ell_{p,k}\leq R\) est obligatoire, non
notationnelle. Si \(\psi,\phi\in\mathcal D_R\), alors
\(\check h_{\psi,\phi}(\pm\ell)=0\) pour \(\ell>R\) ; la contribution de chaque
paire de coordonnées omise à la \emph{différence} des Grams est donc nulle.
La colonne négative non tronquée aurait, en revanche, une norme infinie car
\(\sum_p(\log p)p^{-1/2}=\infty\). Dans l'ensemble fini
\(\ell_{p,k}\leq R\), les coordonnées premières convergent individuellement lorsque
\(L\to\infty\). Dans le canal gamma, l'intégrande est \(O(a)\) en zéro et
\(\mu_\Gamma(a)=O(e^{-a/2})\) à l'infini, de sorte que la convergence se fait aussi
en norme. Nous définissons donc
\begin{equation}\label{eq:pdf2-rh-columnas-reducidas}
 J_{\pm,m,R}:=\lim_{L\to\infty}J_{\pm,m,R,L}
 \quad\text{dans leurs espaces de Hilbert de coordonnées tronquées},
\end{equation}
et supprimons \(m\) lorsqu'il reste fixe. Par convergence dominée,
\[
 \mathscr I_{{\rm GW},L}\longrightarrow\mathscr I_{\rm GW}
 \qquad(L\to\infty),
\]
et les identités sont naturelles sous $m\mapsto m+1$ et $R\le R'$.

Les colonnes réduites \(J_{+,R}\) et \(J_{-,R}\) sont maintenant des opérateurs
matériels de \(\mathcal D_R\) vers Hilbert. Le développement explicite prouve, sans
consulter de zéros ni de signe,
\begin{equation}\label{eq:pdf2-rh-gram-difference-gw}
 \langle J_{+,R}f,J_{+,R}g\rangle
 -\langle J_{-,R}f,J_{-,R}g\rangle
 =\mathscr I_{\rm pr}(f,g)+\mathscr I_\Gamma(f,g)
  +\mathscr I_{\rm pol}(f,g)
 =\mathscr I_{\rm GW}(f,g).
\end{equation}
Les horloges \(\ell_{p,k}\) produisent le terme premier ; l'intégrale avec
\(\mu_\Gamma\) et \(c_\Gamma\check h(0)\), le terme gamma ; et
\[
 \langle b_+(f),b_+(g)\rangle-
 \langle b_-(f),b_-(g)\rangle
 =h_{f,g}(i/2)+h_{f,g}(-i/2)
\]
produit le terme polaire. Cette partie du calcul est matérielle et indépendante
du pont structurel.

\section{Colligation source--first de l'image spectrale--polaire}

Nous fixons d'abord un niveau nonadique \(m\) ; cet indice est supprimé dans
\(J_{\pm,R}\), \(E_R\) et \(\Sigma_R\) jusqu'à la construction du codeur. La
décomposition des codomaines explicites de
\eqref{eq:pdf2-rh-columnas-reducidas} est
\[
 I_+=\{\mathrm{pr},\Gamma,\mathrm{pol}\},\qquad
 I_-(R)=\{\Gamma0,\mathrm{pol}\}
 \sqcup\{(p,k):\ell_{p,k}\leq R\},
\]
\[
 \mathscr H_{+,R}^{\rm amb}=\bigoplus_{i\in I_+}H_{+,i,R},
 \qquad
 \mathscr H_{-,R}^{\rm amb}=\bigoplus_{j\in I_-(R)}H_{-,j,R}.
\]
Nous n'identifions pas un sous-espace atteignable à une somme de ses projections
coordonnées. Nous définissons, les inclusions ambiantes étant sous-entendues,
\[
 H_{+,R}:=\overline{\operatorname{Ran}J_{+,R}},\qquad
 H_{-,R}:=\overline{\operatorname{Ran}J_{-,R}},
\]
comme sous-espaces fermés des deux espaces de Hilbert ambiants. La structure \(\mathfrak G_{\rm sp}\) contient, parmi ses données mathématiques et avant l'expression classique, un espace d'états antispéculaire \(\mathcal H_R^{\rm sh}\) et la ligne isométrique source--first
\begin{equation}\label{eq:pdf2-rh-coligacion-source-first}
 \Sigma_R=
 \begin{pmatrix}\mathcal K_R^{\rm src}\\ \mathcal F_R^{\rm src}\end{pmatrix}:
 H_{+,R}\longrightarrow
 H_{-,R}\oplus\mathcal H_R^{\rm sh},
 \qquad
 \Sigma_R^*\Sigma_R
 = (\mathcal K_R^{\rm src})^*\mathcal K_R^{\rm src}
  +(\mathcal F_R^{\rm src})^*\mathcal F_R^{\rm src}=I_{H_{+,R}}.
\end{equation}
Son action atteignable est fixée sur la colonne positive complète par
les relations des deux publications qui font partie de
\(\operatorname{Rel}_{\rm sp}\). Pour éviter toute substitution de lecteur,
nous notons
\((\mathcal X_R^{\rm str},\Xi_R^{\rm str},P_R^{\rm str})\) la restriction à
\(\mathcal D_R\) du lecteur structurel de
\eqref{eq:pdf2-rh-res-sp-completa}. Ses deux moments sont, littéralement,
\begin{equation}\label{eq:pdf2-rh-momentos-lector-estructural}
\begin{aligned}
 \langle\Xi_R^{\rm str}f,\Xi_R^{\rm str}g\rangle
  &=\langle J_{+,R}f,J_{+,R}g\rangle,\\
 \langle P_R^{\rm str}\Xi_R^{\rm str}f,
          P_R^{\rm str}\Xi_R^{\rm str}g\rangle
  &=\langle J_{-,R}f,J_{-,R}g\rangle,
 \qquad (P_R^{\rm str})^*= (P_R^{\rm str})^2=P_R^{\rm str}.
\end{aligned}
\end{equation}
La coordonnée de frontière est accompagnée de l'isométrie structurelle
\[
 \jmath_R^{\rm sh}:
 \overline{\operatorname{Ran}((I-P_R^{\rm str})\Xi_R^{\rm str})}
 \longrightarrow\mathcal H_R^{\rm sh}.
\]
Par définition de la ligne source--first, non par une identification ultérieure,
\begin{equation}\label{eq:pdf2-rh-source-state-output}
 E_R:=\mathcal F_R^{\rm src}J_{+,R}
     =\jmath_R^{\rm sh}(I-P_R^{\rm str})\Xi_R^{\rm str}:
 \mathcal D_R\longrightarrow\mathcal H_R^{\rm sh},
 \qquad
 \Omega_{0,R}:=\mathcal K_R^{\rm src}J_{+,R}-J_{-,R}.
\end{equation}
Il n'y a pas une colligation par canal : la même ligne reçoit simultanément les
horloges premier--puissance, l'intégrale gamma, le mode scalaire et les deux lignes
polaires. En particulier, elle conserve tous les produits croisés qui apparaissent en
polarisant \eqref{eq:pdf2-rh-gram-difference-gw}.
Avec les projections ambiantes \(\pi_j^-:\mathscr H_{-,R}^{\rm amb}
\to H_{-,j,R}\), les lignes de sortie correctement typées sont
\begin{equation}\label{eq:pdf2-rh-bloques-coligacion}
 \mathcal K_{j,R}^{\rm src}
 :=\pi_j^-\mathcal K_R^{\rm src}:
 H_{+,R}\longrightarrow H_{-,j,R},
 \qquad
 \mathcal F_R^{\rm src}:H_{+,R}\longrightarrow\mathcal H_R^{\rm sh}.
\end{equation}
Ainsi, on n'applique pas \(\mathcal K_R^{\rm src}\) à une coordonnée positive qui
pourrait ne pas appartenir à elle seule à \(H_{+,R}\). Tous les croisements entre
canaux restent dans
\((\mathcal K_R^{\rm src})^*\mathcal K_R^{\rm src}\) et
\((\mathcal F_R^{\rm src})^*\mathcal F_R^{\rm src}\) ; ils ne sont pas orthogonalisés
avant l'identité d'énergie.

La coordonnée de survie de \(\mathfrak G_{\rm sp}\) fournit des espaces résiduels \(\mathcal R_{n,R}^{\rm sh}\), des continuations \(\Omega_{n,R}:\mathcal D_R\to\mathcal R_{n,R}^{\rm sh}\) et des isométries antispéculaires \(V_{n,R}:\mathcal R_{n+1,R}^{\rm sh}\to\mathcal R_{n,R}^{\rm sh}\). Son premier espace et son premier opérateur sont, par définition typée,
\[
 \mathcal R_{0,R}^{\rm sh}:=H_{-,R},\qquad
 \Omega_{0,R}:=\mathcal K_R^{\rm src}J_{+,R}-J_{-,R}.
\]
Sa récurrence d'un tour complet est
\begin{equation}\label{eq:pdf2-rh-recurrencia-residual}
 q=\frac59,\qquad
 \Omega_{n,R}=qV_{n,R}\Omega_{n+1,R},
 \qquad V_{n,R}^*V_{n,R}=I,
 \qquad
 \sup_{n\geq0}\|\Omega_{n,R}f\|<\infty
 \quad(f\in\mathcal D_R).
\end{equation}
Cette récurrence appartient à la colligation enrichie : elle transporte la
différence de sortie, non la forme de Weil ni son signe.

\begin{proposition}[Annulation coinductive de la différence de sortie]
\label{prop:pdf2-rh-omega-cero}
Pour tout \(R>0\) et tout \(f\in\mathcal D_R\),
\begin{equation}\label{eq:pdf2-rh-output-binding-source}
 \Omega_{n,R}f=0\quad(n\geq0),
 \qquad
 \boxed{\mathcal K_R^{\rm src}J_{+,R}=J_{-,R}}.
\end{equation}
\end{proposition}

\begin{proof}
Itérer \eqref{eq:pdf2-rh-recurrencia-residual} \(N\) fois et utiliser le fait que chaque
\(V_{n,R}\) est isométrique donne
\[
 \|\Omega_{n,R}f\|
 =q^N\|\Omega_{n+N,R}f\|
 \leq q^N\sup_{j\geq n}\|\Omega_{j,R}f\|.
\]
Le membre droit converge vers zéro car \(0<q=5/9<1\). Par conséquent,
\(\Omega_{n,R}f=0\) ; le cas \(n=0\) et
\eqref{eq:pdf2-rh-source-state-output} produisent l'identité encadrée.
\end{proof}

L'action ligne par ligne de la sortie n'est plus implicite. Les projections sont
les restrictions à \(H_{-,R}\) des projections du Hilbert ambiant, et
\begin{equation}\label{eq:pdf2-rh-salidas-generador-a-generador}
\begin{aligned}
 \mathcal K_{\Gamma0,R}^{\rm src}J_{+,R}f
  &=\sqrt{-c_\Gamma}\,j_mf,\\
 \mathcal K_{(p,k),R}^{\rm src}J_{+,R}f
  &=\sqrt{2w_{p,k}}\,j_mf
  &&(\ell_{p,k}\leq R),\\
 \mathcal K_{{\rm pol},R}^{\rm src}J_{+,R}f
  &=\zeta_-b_-(f).
\end{aligned}
\end{equation}
Ces lignes épuisent \(J_{-,R}\), tandis que
\begin{equation}\label{eq:pdf2-rh-frontera-generador-a-generador}
 \mathcal F_R^{\rm src}J_{+,R}f=E_Rf
\end{equation}
conserve, avec tous ses croisements, la coordonnée de frontière.
Évaluer \eqref{eq:pdf2-rh-coligacion-source-first} en \(J_{+,R}f\) et
\(J_{+,R}g\), et utiliser \eqref{eq:pdf2-rh-output-binding-source}, produit
l'identité d'énergie antérieure au signe
\begin{equation}\label{eq:pdf2-rh-energia-source-first}
 \boxed{
 \langle J_{+,R}f,J_{+,R}g\rangle
 =\langle J_{-,R}f,J_{-,R}g\rangle
  +\langle E_Rf,E_Rg\rangle.}
\end{equation}
En la comparant au calcul indépendant
\eqref{eq:pdf2-rh-gram-difference-gw}, on obtient, sans introduire une racine de
la forme cible,
\begin{equation}\label{eq:pdf2-rh-binding-source-first}
 \boxed{E_R^*E_R=\mathscr I_{\rm GW}\big|_{\mathcal D_R}.}
\end{equation}
En particulier,
\(\Delta_{{\rm bind},R}
:=E_R^*E_R-\mathscr I_{\rm GW}|_{\mathcal D_R}=0\) ; cette annulation est la
sortie de \eqref{eq:pdf2-rh-coligacion-source-first}--
\eqref{eq:pdf2-rh-energia-source-first}, non une prémisse insérée dans le
publicateur.

\section{Complément structurel, feuillet miroir et télescopage terminal}

L'état \(E_R=\mathcal F_R^{\rm src}J_{+,R}\) construit par la ligne
source--first est la coordonnée de frontière de
\(\mathfrak G_{\rm sp}\). Nous définissons sa forme structurelle par
\begin{equation}\label{eq:pdf2-rh-bw-estructural}
 \mathcal B_{{\rm sp},R}^{\rm str}
 :=E_R^*E_R
 =\mathscr I_{\rm GW}\big|_{\mathcal D_R},
\end{equation}
où la deuxième égalité est
\eqref{eq:pdf2-rh-binding-source-first}, déjà déduite de la colligation et du
développement générateur par générateur.

Soient \(P_+^{\rm sh}\) et \(P_-^{\rm sh}\) les projecteurs orthogonaux
complémentaires des deux feuillets. On fixe
\begin{equation}\label{eq:pdf2-rh-q-a-d}
 q=\frac59,\qquad
 A=P_+^{\rm sh}+qP_-^{\rm sh},\qquad
 D=\sqrt{1-q^2}\,P_-^{\rm sh}.
\end{equation}
L'orthogonalité donne l'identité exacte
\begin{equation}\label{eq:pdf2-rh-conservacion-local}
 A^*A+D^*D
 =P_+^{\rm sh}+q^2P_-^{\rm sh}
 +(1-q^2)P_-^{\rm sh}=I.
\end{equation}
Les neuf phases coordonnent un seul tour avec mémoire :
\begin{equation}\label{eq:pdf2-rh-monodromia}
 M_9=\Phi A_8\cdots A_0=qV_9,
 \qquad V_9^*V_9=I.
\end{equation}
La contraction par tour est \(q\), non \(q^9\) ; les phases internes conservent
l'ordre et la retenue.

La colligation type
\(E_R:\mathcal D_R\to\operatorname{Ran}P_-^{\rm sh}\). Par conséquent,
\begin{equation}\label{eq:pdf2-rh-binding-espejo}
 P_+^{\rm sh}E_R=0,
 \qquad AE_R=qE_R,
 \qquad
 E_R^*E_R=\mathcal B_{{\rm sp},R}^{\rm str}
 =\mathscr I_{\rm GW}\big|_{\mathcal D_R}.
\end{equation}

\begin{proposition}[Télescopage fini avec terminal conservé]
\label{prop:pdf2-rh-telescopia}
Pour chaque \(N\geq1\),
\begin{equation}\label{eq:pdf2-rh-telescopia-finita}
 \mathcal B_{{\rm sp},R}^{\rm str}
 =\sum_{n=0}^{N-1}(DA^nE_R)^*(DA^nE_R)
 +(A^NE_R)^*(A^NE_R).
\end{equation}
De manière équivalente, évaluée en \(f,g\in\mathcal D_R\), l'identité finie de
la colligation est
\begin{equation}\label{eq:pdf2-rh-balance-finito-columnas}
\boxed{
\begin{aligned}
 &\langle J_{+,R}f,J_{+,R}g\rangle
  -\langle J_{-,R}f,J_{-,R}g\rangle\\
 &\quad=\sum_{n=0}^{N-1}
  \langle DA^nE_Rf,DA^nE_Rg\rangle
  +\langle A^NE_Rf,A^NE_Rg\rangle.
\end{aligned}}
\end{equation}
Le dernier terme satisfait
\begin{equation}\label{eq:pdf2-rh-terminal-extincion}
 (A^NE_R)^*(A^NE_R)
 =q^{2N}\mathcal B_{{\rm sp},R}^{\rm str}
 \xrightarrow[N\to\infty]{}0.
\end{equation}
\end{proposition}

\begin{proof}
En appliquant \eqref{eq:pdf2-rh-conservacion-local} à \(A^nE_R\),
\[
 (A^nE_R)^*(A^nE_R)
 =(DA^nE_R)^*(DA^nE_R)
 +(A^{n+1}E_R)^*(A^{n+1}E_R).
\]
La somme de \(n=0\) à \(N-1\) annule uniquement les termes intermédiaires et
conserve \((A^NE_R)^*(A^NE_R)\). Par
\eqref{eq:pdf2-rh-binding-espejo}, \(A^NE_R=q^NE_R\) ; le terminal
est donc \(q^{2N}E_R^*E_R=q^{2N}\mathcal B_{{\rm sp},R}^{\rm str}\). Comme
\(0<q=5/9<1\), il converge géométriquement vers zéro.
\end{proof}

Ce n'est qu'après avoir démontré \eqref{eq:pdf2-rh-terminal-extincion} que l'on définit
\begin{equation}\label{eq:pdf2-rh-l-r}
 \mathscr L_Rf=(DA^nE_Rf)_{n\geq0}.
\end{equation}
La limite du télescopage fini donne
\begin{equation}\label{eq:pdf2-rh-lstar-l}
 \boxed{
 \mathscr L_R^*\mathscr L_R
 =\mathcal B_{{\rm sp},R}^{\rm str}
 =\mathscr I_{\rm GW}\big|_{\mathcal D_R}\succeq0.}
\end{equation}

\section{Codeur commun, projection et double publication explicites}

Fixons \(m\geq1\) et deux vecteurs orthonormaux
\(u_m,\eta_m\in\mathbb C^{9^m}\), où \(u_m\) est le mode uniforme et
\(\eta_m\) le mode de retenue. L'espace et la projection du codeur sont
\begin{equation}\label{eq:pdf2-rh-hilbert-comun}
 \mathcal X_{m,R}^{\rm sp}
 :=\mathbb C^{9^m}\widehat\otimes
   (H_{-,R}\oplus\mathcal H_R^{\rm sh}),
 \qquad
 P_{m,R}:=|u_m\rangle\langle u_m|\otimes
 I_{H_{-,R}\oplus\mathcal H_R^{\rm sh}},
 \qquad Q_{m,R}:=I-P_{m,R}.
\end{equation}
L'action complète du codeur des cinq canaux est
\begin{equation}\label{eq:pdf2-rh-xi-comun-explicita}
 \boxed{
 \begin{aligned}
  \mathfrak C_{m,R}f
  &:=u_m\otimes(J_{-,R}f,0)
    +\eta_m\otimes(0,E_Rf),\\
  \Xi_{m,R}f&:=\mathfrak C_{m,R}f.
 \end{aligned}}
\end{equation}
On ne prend pas une racine de \(\mathscr I_{\rm GW}\) : les deux entrées sont la
sortie et l'état de la colligation source--first déjà définis dans
\eqref{eq:pdf2-rh-source-state-output}.

\begin{proposition}[Les deux moments s'obtiennent à partir de l'action du codeur]
\label{prop:pdf2-rh-vmas-isometria}
Pour tout \(f,g\in\mathcal D_R\),
\begin{equation}\label{eq:pdf2-rh-dos-momentos-construidos}
\begin{aligned}
 \langle\mathfrak C_{m,R}f,\mathfrak C_{m,R}g\rangle
  &=\langle J_{+,R}f,J_{+,R}g\rangle,\\
 \langle P_{m,R}\mathfrak C_{m,R}f,
         P_{m,R}\mathfrak C_{m,R}g\rangle
  &=\langle J_{-,R}f,J_{-,R}g\rangle,\\
 \langle Q_{m,R}\Xi_{m,R}f,Q_{m,R}\Xi_{m,R}g\rangle
  &=\mathscr I_{\rm GW}(f,g).
\end{aligned}
\end{equation}
\end{proposition}

\begin{proof}
L'orthogonalité \(u_m\perp\eta_m\) et
\eqref{eq:pdf2-rh-energia-source-first} donnent
\[
 \langle\mathfrak C_{m,R}f,\mathfrak C_{m,R}g\rangle
 =\langle J_{-,R}f,J_{-,R}g\rangle
  +\langle E_Rf,E_Rg\rangle
 =\langle J_{+,R}f,J_{+,R}g\rangle.
\]
La projection \(P_{m,R}\) conserve le terme uniforme et élimine celui de
retenue ; \(Q_{m,R}\) fait l'inverse. Les deux autres identités découlent de
\eqref{eq:pdf2-rh-binding-source-first}.
\end{proof}

Le codeur qui vient d'être écrit est un représentant explicite du lecteur
structurel initial, non un lecteur de substitution. Sur les sous-espaces atteignables, on
définit
\begin{equation}\label{eq:pdf2-rh-isomorfismo-lector-estructural}
\begin{aligned}
 \mathcal U_{m,R}^{\rm str}
  (P_R^{\rm str}\Xi_R^{\rm str}f)
   &:=u_m\otimes(J_{-,R}f,0),\\
 \mathcal U_{m,R}^{\rm str}
  ((I-P_R^{\rm str})\Xi_R^{\rm str}f)
   &:=\eta_m\otimes(0,E_Rf).
\end{aligned}
\end{equation}
Les deux sources sont orthogonales. La première égalité de Gram de chaque source
est \eqref{eq:pdf2-rh-momentos-lector-estructural} ; pour la deuxième, on utilise
\(E_R=\jmath_R^{\rm sh}(I-P_R^{\rm str})\Xi_R^{\rm str}\). Ainsi, les deux
règles sont bien définies, sont isométriques et s'étendent à la somme orthogonale
de leurs adhérences. Dans cette adhérence atteignable, elles satisfont
\begin{equation}\label{eq:pdf2-rh-entrelazamiento-lector-estructural}
 \mathcal U_{m,R}^{\rm str}\Xi_R^{\rm str}=\Xi_{m,R},
 \qquad
 \mathcal U_{m,R}^{\rm str}P_R^{\rm str}
 =P_{m,R}\mathcal U_{m,R}^{\rm str},
 \qquad
 (\Xi_R^{\rm str})^*(I-P_R^{\rm str})\Xi_R^{\rm str}=E_R^*E_R.
\end{equation}
Ainsi, \(B_{{\rm sp},R}^{\rm str}\), le complément de la structure de
départ, est exactement le complément du codeur construit ; on n'a pas
changé d'objet en passant de \eqref{eq:pdf2-rh-res-sp-completa} à
\eqref{eq:pdf2-rh-xi-comun-explicita}.

La règle
\begin{equation}\label{eq:pdf2-rh-vmas-rango}
 V_{+,m,R}^{0}(J_{+,R}f):=\Xi_{m,R}f
\end{equation}
est bien définie et isométrique par la première égalité de
\eqref{eq:pdf2-rh-dos-momentos-construidos} ; elle s'étend de manière unique en
une isométrie
\(V_{+,m,R}:H_{+,R}\to\mathcal X_{m,R}^{\rm sp}\). La publication négative
est l'isométrie
\begin{equation}\label{eq:pdf2-rh-vmenos}
 V_{-,m,R}:H_{-,R}\longrightarrow\mathcal X_{m,R}^{\rm sp},
 \qquad V_{-,m,R}y=u_m\otimes(y,0).
\end{equation}
Par construction,
\begin{equation}\label{eq:pdf2-rh-doble-publicacion-material}
 \boxed{
 V_{+,m,R}J_{+,R}=\Xi_{m,R},
 \qquad
 V_{-,m,R}J_{-,R}=P_{m,R}\Xi_{m,R}.}
\end{equation}
Le complément conserve exactement l'état antispéculaire et sa colonne
d'observation :
\begin{equation}\label{eq:pdf2-rh-complemento-comun}
 \boxed{
 \Xi_{m,R}^*Q_{m,R}\Xi_{m,R}
 =E_R^*E_R
 =\mathscr L_R^*\mathscr L_R
 =\mathscr I_{\rm GW}\big|_{\mathcal D_R}.}
\end{equation}

\begin{proposition}[Connecteur induit et coïncidence source--first]
\label{prop:pdf2-rh-conector}
L'opérateur
\begin{equation}\label{eq:pdf2-rh-k-r}
 \mathcal K_R
 :=V_{-,m,R}^*P_{m,R}V_{+,m,R}:
 H_{+,R}\longrightarrow H_{-,R}
\end{equation}
est contractif, coïncide avec \(\mathcal K_R^{\rm src}\) sur tout
\(H_{+,R}\) et satisfait
\begin{equation}\label{eq:pdf2-rh-k-binding}
 \|\mathcal K_R\|\leq1,
 \qquad
 \mathcal K_RJ_{+,R}=J_{-,R}.
\end{equation}
\end{proposition}

\begin{proof}
Les deux cartes sont des isométries et \(P_{m,R}\) est une projection orthogonale,
donc \(\|\mathcal K_R\|\leq1\). Sur l'image dense de \(J_{+,R}\),
\[
 \mathcal K_RJ_{+,R}f
 =V_{-,m,R}^*P_{m,R}\Xi_{m,R}f
 =V_{-,m,R}^*V_{-,m,R}J_{-,R}f
 =J_{-,R}f.
\]
La même identité vaut pour \(\mathcal K_R^{\rm src}\) par
\eqref{eq:pdf2-rh-output-binding-source} ; la continuité et la densité
prouvent que les deux opérateurs coïncident.
\end{proof}

La différence de Gram est ainsi publiée dans toutes ses réalisations
coïncidentes :
\begin{equation}\label{eq:pdf2-rh-tres-publicaciones}
\boxed{
\begin{aligned}
 J_{+,R}^*J_{+,R}-J_{-,R}^*J_{-,R}
 &=\Xi_{m,R}^*Q_{m,R}\Xi_{m,R}\\
 &=E_R^*E_R\\
 &=\mathscr L_R^*\mathscr L_R\\
 &=\mathcal B_{{\rm sp},R}^{\rm str}=\mathscr I_{\rm GW}.
\end{aligned}}
\end{equation}

\section{Naturalité cofinale et critère réversible de Weil}

Pour \(R\leq R'\), soit
\(i_{R,R'}:\mathcal D_R\hookrightarrow\mathcal D_{R'}\). Les colonnes
individuelles ne forment pas un système isométrique : lorsqu'une nouvelle longueur
\(R<\ell_{p,k}\leq R'\) apparaît, les deux acquièrent la même énergie positive. Pour
\(f,g\in\mathcal D_R\), la séparation des supports donne exactement
\begin{equation}\label{eq:pdf2-rh-naturalidad}
 \left\langle(I-T_{\ell_{p,k}})f,
                    (I-T_{\ell_{p,k}})g\right\rangle
 =2\langle f,g\rangle
 =\left\langle\sqrt2j_mf,\sqrt2j_mg\right\rangle .
\end{equation}
Par conséquent, l'incrément premier de
\(J_{+,R'}^*J_{+,R'}\) coïncide avec celui de
\(J_{-,R'}^*J_{-,R'}\) et s'annule dans la différence, mais on ne déclare pas une
fausse isométrie entre les colonnes. La compatibilité cofinale correcte est celle
des formes :
\begin{equation}\label{eq:pdf2-rh-naturalidad-p}
 \boxed{
 \mathscr I_{{\rm GW},R}
 =i_{R,R'}^*\mathscr I_{{\rm GW},R'}i_{R,R'},
 \qquad
 \mathcal B_{{\rm sp},R}^{\rm str}
 =i_{R,R'}^*\mathcal B_{{\rm sp},R'}^{\rm str}i_{R,R'}.}
\end{equation}
Comme \(E_R^*E_R=\mathscr L_R^*\mathscr L_R=\mathscr I_{{\rm GW},R}\), les
réalisations minimales de Kolmogorov admettent bien des isométries uniques sur leurs
images atteignables, définies par
\begin{equation}\label{eq:pdf2-rh-naturalidad-bw}
 \iota_{R,R'}^E(E_Rf):=E_{R'}f,
 \qquad
 \iota_{R,R'}^L(\mathscr L_Rf):=\mathscr L_{R'}f
 \quad(f\in\mathcal D_R).
\end{equation}
Ces deux règles sont bien définies précisément par
\eqref{eq:pdf2-rh-naturalidad-p} ; elles ne s'étendent pas par décret à
\(J_{\pm,R}\), \(\Xi_{m,R}\) ou \(\mathcal K_R\). Le raffinement nonadique
\(m\mapsto m+1\) reste, à fenêtre fixe, dans les isométries structurelles
qui transportent \(j_m\), \(\widehat\delta_{a,m}\), \(u_m\) et \(\eta_m\).

Pour ne pas comprimer la dernière étape, nous fixons aussi ses objets et
normalisations. Soit
\[
 \xi(s):=\frac12s(s-1)\pi^{-s/2}\Gamma(s/2)\zeta(s)
\]
et soit \(\mathscr Z\) le multiensemble de ses zéros, chaque zéro étant répété selon
sa multiplicité. L'équation fonctionnelle et la réalité des coefficients
préservent \(\mathscr Z\) par
\(\rho\mapsto1-\rho\), \(\rho\mapsto\overline\rho\) et
\(\rho\mapsto1-\overline\rho\). Avec notre convention de Fourier, on définit
\begin{equation}\label{eq:pdf2-rh-transformada-weil}
 \Phi_f(s):=\widehat f\!\left(i(s-\tfrac12)\right),
 \qquad
 \widetilde f(x):=\overline{f(-x)}.
\end{equation}
La formule explicite complète calculée dans
\eqref{eq:pdf2-rh-gram-difference-gw}, y compris les termes gamma et polaire
qui retirent les pôles et les zéros triviaux, a la forme spectrale
\begin{equation}\label{eq:pdf2-rh-formula-espectral-weil}
 \mathscr I_{\rm GW}(f,g)
 =\sum_{\rho\in\mathscr Z}
   \Phi_f(\rho)\,
   \overline{\Phi_g(1-\overline\rho)}.
\end{equation}
La somme compte les multiplicités. Elle converge absolument : pour
\(f,g\in C_c^\infty(\mathbb R)\), Paley--Wiener donne une décroissance plus rapide
que toute puissance dans les bandes verticales, tandis que le nombre de zéros avec
\(|\operatorname{Im}\rho|\leq T\) est \(O(T\log T)\).

\begin{proposition}[Critère réversible de Weil avec quantificateurs complets]
\label{prop:pdf2-rh-criterio-weil-desplegado}
Sous la normalisation
\eqref{eq:pdf2-rh-transformada-weil} et en comptant tous les zéros non triviaux
avec leur multiplicité,
\begin{equation}\label{eq:pdf2-rh-criterio-weil}
 \left[
  \mathscr I_{\rm GW}(f,f)\geq0
  \text{ pour tout }f\in C_c^\infty(\mathbb R;\mathbb C)
 \right]
 \Longleftrightarrow
 \left[
  \operatorname{Re}\rho=\frac12
  \text{ pour tout }\rho\in\mathscr Z
 \right].
\end{equation}
\end{proposition}

L'implication réciproque ne s'obtient pas en interpolant un nombre croissant de
zéros : cette opération ne donnerait pas de contrôle uniforme sur la queue. On utilise le
critère classique complet de Weil dans la formulation de Bombieri, dont l'espace
de test et la normalisation sont transposés ici sans les abréger.\footnote{%
E.~Bombieri, \emph{Remarks on Weil's quadratic functional in the theory of
prime numbers, I}, Rendiconti Lincei, Matematica e Applicazioni, ser.~9,
vol.~11 (2000), pp.~183--233. Le critère qui y est fixé affirme que la fonctionnelle
quadratique de la formule explicite est semi-définie positive sur tout
\(C_c^\infty(0,\infty)\) si et seulement si l'hypothèse de Riemann est vraie.}

\begin{proof}
Nous écrivons matériellement le changement de coordonnées qui identifie les deux
critères. Pour \(f\in C_c^\infty(\mathbb R)\), soit
\begin{equation}\label{eq:pdf2-rh-cambio-mellin}
 (\mathcal M f)(x):=x^{-1/2}f(\log x),
 \qquad x>0.
\end{equation}
L’application
\[
 \mathcal M:C_c^\infty(\mathbb R)
 \longrightarrow C_c^\infty(0,\infty)
\]
est bijectif, avec pour inverse
\begin{equation}\label{eq:pdf2-rh-cambio-mellin-inverso}
 (\mathcal M^{-1}g)(u)=e^{u/2}g(e^u).
\end{equation}
Si
\(\widehat g_{\rm Mel}(s):=\int_0^\infty g(x)x^{s-1}\,dx\), le changement
\(x=e^u\) donne, sans facteur résiduel,
\begin{equation}\label{eq:pdf2-rh-mellin-fourier}
 \begin{aligned}
 \widehat{\mathcal M f}_{\rm Mel}(s)
 &=\int_{\mathbb R}f(u)e^{(s-1/2)u}\,du\\
 &=\widehat f\!\left(i(s-\tfrac12)\right)
 =\Phi_f(s).
 \end{aligned}
\end{equation}
De plus,
\begin{equation}\label{eq:pdf2-rh-mellin-involucion}
 \widehat{\overline{\mathcal M f}}_{\rm Mel}(1-\rho)
 =\overline{
   \widehat{\mathcal M f}_{\rm Mel}(1-\overline\rho)}
 =\overline{\Phi_f(1-\overline\rho)}.
\end{equation}
Par conséquent, la fonctionnelle du critère de Weil--Bombieri est exactement
\begin{equation}\label{eq:pdf2-rh-bombieri-identificado}
 \sum_{\rho\in\mathscr Z}
 \widehat{\mathcal M f}_{\rm Mel}(\rho)\,
 \widehat{\overline{\mathcal M f}}_{\rm Mel}(1-\rho)
 =
 \sum_{\rho\in\mathscr Z}
 \Phi_f(\rho)\overline{\Phi_f(1-\overline\rho)}
 =\mathscr I_{\rm GW}(f,f).
\end{equation}
La classe des tests n'est pas modifiée, aucune multiplicité n'est perdue et
l'ensemble des zéros n'est pas tronqué, car \(\mathcal M\) est une bijection. Le
critère classique cité appliqué à
\eqref{eq:pdf2-rh-bombieri-identificado} donne précisément l'équivalence
\eqref{eq:pdf2-rh-criterio-weil}. Dans le sens direct, si
\(\operatorname{Re}\rho=1/2\), la même identité se réduit explicitement à
\[
 \mathscr I_{\rm GW}(f,f)
 =\sum_{\rho\in\mathscr Z}|\Phi_f(\rho)|^2\geq0.
\]
\end{proof}

L'égalité \(\bigcup_R\mathcal D_R=C_c^\infty(\mathbb R)\) est exacte ; chaque
fonction employée dans l'argument appartient à une fenêtre et la naturalité
cofinale y transporte la positivité déjà démontrée.

\begin{theorem}[Résolution de Riemann--Weil à partir de l'image spectrale--polaire]
\label{thm:pdf2-rh-resolucion}
La structure complète
\(\mathfrak G_{\rm sp}\) étant fixée, la forme
complète de Guinand--Weil est positive sur la classe cofinale et séparatrice. En
conséquence, tout zéro non trivial \(\rho\) de la fonction zêta de Riemann
satisfait
\[
 \boxed{\operatorname{Re}\rho=\frac12.}
\]
\end{theorem}

\begin{proof}
Le développement premier--gamma--polaire identifie la différence des deux colonnes
à \(\mathscr I_{\rm GW}\) en utilisant un unique régulateur et le retire par
convergence dominée. La récurrence antispéculaire
\eqref{eq:pdf2-rh-recurrencia-residual} annule la différence de sortie avant de
former une inégalité ; l'isométrie \(\Sigma_R\) donne donc
\eqref{eq:pdf2-rh-energia-source-first}. Comparer cette identité au
développement calculé produit
\(E_R^*E_R=\mathscr I_{\rm GW}\). L'identité locale
\(A^*A+D^*D=I\) est itérée sans effacer le terminal et donne
\eqref{eq:pdf2-rh-telescopia-finita} ; le terminal est exactement
\(q^{2N}E_R^*E_R\) et ne s'éteint qu'à la limite. D'où
\(\mathscr L_R^*\mathscr L_R=\mathscr I_{\rm GW}\succeq0\).
L'action
\eqref{eq:pdf2-rh-xi-comun-explicita} construit alors le codeur, ses
deux moments, la projection et les deux publications, et
\eqref{eq:pdf2-rh-tres-publicaciones} identifie toutes les réalisations du
même carré. La naturalité transporte cette identité à toute l'union
cofinale. Enfin,
\eqref{eq:pdf2-rh-criterio-weil} conclut l'affirmation.
\end{proof}

\section{Contrôle de Redheffer local ultérieur et non directeur}

Le déploiement suivant vérifie que la passivité locale peut être obtenue à partir de
la retenue sans utiliser la forme de Weil. C'est un contrôle de non-circularité, non une
prémisse du \cref{thm:pdf2-rh-resolucion}. Ce transfert tronqué n'est pas
la colligation complète \(\Sigma_R\) de
\eqref{eq:pdf2-rh-coligacion-source-first}.

Pour \(N=9^m\), soit \(U_{a,N}\) l'odomètre unitaire
\[
 U_{a,N}(e_j\otimes f)=e_{j+1}\otimes f\ (j<N-1),
 \qquad U_{a,N}(e_{N-1}\otimes f)=e_0\otimes T_af.
\]
Avec \(P_N=|u_N\rangle\langle u_N|\otimes I\), \(Q_N=I-P_N\),
\(A_{a,N}=P_NU_{a,N}P_N\) et \(C_{a,N}=Q_NU_{a,N}P_N\),
\begin{equation}\label{eq:pdf2-rh-source-energia}
 A_{a,N}^*A_{a,N}+C_{a,N}^*C_{a,N}=P_N,
 \qquad
 C_{a,N}^*C_{a,N}
 =\frac{N-1}{N^2}(2I-T_a-T_a^*).
\end{equation}
Pour \(0<s<1\), le transfert
\begin{equation}\label{eq:pdf2-rh-source-redheffer}
 \mathscr R_{a,N}(s)=(sI-A_{a,N})(I-sA_{a,N})^{-1}
\end{equation}
satisfait
\begin{equation}\label{eq:pdf2-rh-source-defecto}
\begin{aligned}
 I-\mathscr R_{a,N}(s)^*\mathscr R_{a,N}(s)
 ={}&(1-s^2)(I-sA_{a,N}^*)^{-1}\\
 &\quad\cdot C_{a,N}^*C_{a,N}(I-sA_{a,N})^{-1}\succeq0.
\end{aligned}
\end{equation}
Avec
\[
 s_{p,N}=\frac{Np^{-1/2}}{1+(N-1)p^{-1/2}},
 \qquad s_\Gamma(a)=e^{-a/2},
\]
la somme sur les horloges premières et l'intégrale directe gamma conservent une norme au
plus égale à un. Explicitement,
\begin{equation}\label{eq:pdf2-rh-source-relojes}
 \mathscr R_{N,R}
 =\bigoplus_{\log p\leq R}
   \mathscr R_{\log p,N}(s_{p,N})
 \ \oplus\!
 \int_0^\infty{}^\oplus
   \mathscr R_{a,N}(s_\Gamma(a))\,da,
 \qquad \|\mathscr R_{N,R}\|\leq1.
\end{equation}
La carte finie des rôles
\[
 \mathcal H_+=\operatorname{span}\{f_3,f_6,f_9\},\quad
 \mathcal H_-=\operatorname{span}\{f_4,f_8\},\quad
 C_\angle=|f_4\rangle\langle f_3|+|f_8\rangle\langle f_9|
\]
satisfait \(I-C_\angle^*C_\angle=P_6\). Avec les cartes énergétiques
two-way, on fixe
\begin{equation}\label{eq:pdf2-rh-source-cartas}
\begin{aligned}
 \mathcal K_q&=K_2\otimes(\mathbb C^9)^{\otimes q},\\
 \mathsf B_{\pm,q}:\mathcal M_{\pm,q,R}
 &\longrightarrow
 \mathcal K_q\widehat\otimes\mathcal H_\pm
 \widehat\otimes\mathcal H_{\mathfrak C_R},\\
 \mathsf B_{\pm,q}^*\mathsf B_{\pm,q}&=\mathsf G_{\pm,q},
 \qquad
 \mathsf B_{\pm,q}^{\sharp}
 =\mathsf G_{\pm,q}^{-1}\mathsf B_{\pm,q}^*.
\end{aligned}
\end{equation}
Le transfert
\begin{equation}\label{eq:pdf2-rh-source-transferencia}
 \mathscr C_{q,R}^{\rm src}
 =\mathsf B_{-,q}^{\sharp}
  (I\otimes C_\angle\otimes\mathscr R_{N,R})\mathsf B_{+,q}
\end{equation}
a pour défaut
\begin{align}
 \Delta_{q,R}^{\rm src}
 ={}&I\otimes\bigl[P_6\otimes I
 +(P_3+P_9)\otimes
 (I-\mathscr R_{N,R}^*\mathscr R_{N,R})\bigr]\succeq0,
 \label{eq:pdf2-rh-source-delta}\\
 I-(\mathscr C_{q,R}^{\rm src})^\sharp\mathscr C_{q,R}^{\rm src}
 ={}&\mathsf B_{+,q}^{\sharp}\Delta_{q,R}^{\rm src}
 \mathsf B_{+,q}\succeq0.
 \label{eq:pdf2-rh-source-contractividad}
\end{align}
Ce contrôle conserve chemin, horloge, rôle, retenue et canaux avant de former
une différence. Il n'est pas identifié à \(\mathcal K_R\) et ne remplace pas la
publication commune explicite de
\eqref{eq:pdf2-rh-doble-publicacion-material} ; sa portée locale ne
rouvre donc pas la conclusion cofinale.

\section{Réfutateurs d'une ombre amputée}

\begin{falsifier}[Substitutions qui changent l'objet de départ]
L'argument précédent ne se transfère pas à une ombre qui élimine l'une des
treize coordonnées de \eqref{eq:pdf2-rh-res-sp-completa}, construise \(J_+\) et
\(J_-\) à partir d'antécédents distincts, remplace une projection orthogonale par
un idempotent non autoadjoint, omette un canal, utilise des régulateurs incompatibles,
choisisse \(P_R\) après le signe, efface le terminal avant la limite ou réduise
la classe à une famille non cofinale.
\end{falsifier}

\begin{proof}
Chaque opération rompt, respectivement, le domaine de
\(\mathfrak G_{\rm sp}\), la bonne définition de
\eqref{eq:pdf2-rh-vmas-rango}, le carré orthogonal, l'identité fonctionnelle,
la direction causale, le télescopage fini ou le critère réversible. Le
réfutateur affecte l'ombre substituée ; il ne rétrograde pas l'image structurelle
complète employée dans le \cref{thm:pdf2-rh-resolucion}.
\end{proof}

\paragraph{Provenance et force probante.}
La restriction spectrale--polaire, la colligation source--first, son complément et le feuillet miroir sont une \texttt{ARCHITECTURE\_AUCTORIALE\_PRÉEXISTANTE}. L'annulation coinductive de \(\Omega\), l'identité finie avec terminal et le codeur d'action \eqref{eq:pdf2-rh-xi-comun-explicita} constituent une \texttt{FORMALISATION\_NOUVELLE} de cette structure complète. L'identification Guinand--Weil et la fermeture cofinale constituent un \texttt{RÉSULTAT\_RETROUVÉ}. Les équations \eqref{eq:pdf2-rh-coligacion-source-first}--\eqref{eq:pdf2-rh-tres-publicaciones} exposent l'objet, ses domaines, son action, le terminal et la limite sans introduire de racine abstraite de la forme ; leur portée est \emph{démontrée avec une structure de départ explicite}. \endgroup
 \part{Synthèse théorématique}

\chapter{L'arithmétique antérieure à la réalisation archimédienne}
\section{VI. La structure discrète du continu}

\subsection{Le continu est une réalisation de cylindres compatibles}

La limite projective \(X_\infty\) conserve toutes les trajectoires compatibles. Ses cylindres semi-fermés déterminent une topologie ; le raffinement détermine la forme ; le cocycle détermine l'holonomie ; l'orientation détermine les bords ; et l'algèbre inductive détermine les lecteurs. Les cinq aspects internes spectral, solénoïdal, de jauge, cohérent et déterminantal émergent conjointement de cette construction. Les séparer détruirait l'objet qui leur confère une naturalité commune.

La double projection est un couple de représentations de la même algèbre :

\[
\rho_{\rm arq}:\mathfrak H_{\rm cyl}\to\mathcal A_{\mathbb R},
\qquad
\rho_{\rm exc}:\mathfrak H_{\rm cyl}\to\mathcal A_{\rm exc}.
\]

La première contracte des cylindres compatibles vers la lecture archimédienne. La seconde conserve l'incidence de frontière vers Witt--Golay--Leech--\(V^\natural\)--Monster. \(\mathbb R\) et la branche exceptionnelle ne sont pas deux univers d'origine : ce sont deux quotients représentatifs d'une source plus riche.

L'abélianisation oublie nécessairement le commutateur électronique et une partie du registre de mémoire. La droite réelle peut donc être complète comme continu scalaire tout en demeurant insuffisante comme description de la généalogie. L'arithmétique holonomique précède \(\mathbb R\) comme une trajectoire précède l'une de ses coordonnées.

\subsection{Prédéfinition des futurs compatibles}

Soit

\[
\mathcal S_n(x)=\{y\in X_{n+1}:r_{n+1,n}(y)=x
\text{ et }y\text{ est admissible}\}.
\]

L'arbre coinductif contient les prolongements possibles. La loi sturmienne garantit un unique mot spécial à droite par longueur ; TRIT type sa courbure et TPK sélectionne la continuation en conservant la branche complémentaire à la frontière. La prédéfinition ne signifie pas qu'une valeur extérieure a été enregistrée. Elle signifie que chaque sortie finie est fixée par un cylindre et que le système de cylindres est compatible à toute profondeur.

Ainsi se précise l'expression « choix des futurs de survie ». Le futur n'est pas deviné à partir d'une mesure : il est sélectionné dans l'ensemble déjà typé des prolongements qui conservent les relations de l'état. Passé, présent et futur sont des fonctions de l'opérateur de retour, du seuil et de l'ouverture, non de simples étiquettes verbales.

\section{VII. Théorème fondamental et conséquences}

\subsection{Énoncé minimal}

Nous définissons l'objet fondamental

\[
\boxed{
\mathfrak H_{\rm HMT}^{\rm alg}
=\left[\varinjlim_n
\left(
\mathbb R\langle e_x,J_x,T_u\rangle/
\mathcal R_{\rm HMT}
\right)\right]
\rtimes_{\Gamma_9}\mathbb N.}
\]

Dans le secteur réversible, \(\mathbb N\) est remplacé par \(\mathbb Z\). Son espace d'états infinis est

\[
X_\infty=\varprojlim_nX_n.
\]

L'expression est algébrique. Une représentation bornée \(\rho\) détermine sa propre complétion \(\mathfrak H_{\rm HMT}^{\rho}
=\overline{\rho(\mathfrak H_{\rm HMT}^{\rm alg})}^{\|\cdot\|}\) ; la topologie du codomaine appartient donc au lecteur et ne précède pas le générateur.

\textbf{Théorème fondamental de l'évolution conservative de l'unité.} L'arithmétique APP, orientée par les fibres \(J_\tau^2=-\tau I\) et transportée par le cocycle TPK, engendre une algèbre associative non commutative dont le raffinement préserve l'unité et dont les retours élèvent la mémoire. Le spectre typé des représentations de cette algèbre détermine corrélativement le système corrélé \((\pi,\varphi,e,\alpha)\), la structure conjointe du continu, la lacune électronique, les échelles d'action et de gravité, les paramètres géométriques, les constantes du catalogue et les symétries de frontière. Chaque réalisation conserve son domaine, son lecteur et son degré.

La démonstration se divise en six identités consécutives :

\begin{enumerate}
\item les équations de l'APP prouvent la conservation du résidu et de la retenue ;
\item les relations tritiques prouvent la coexistence des trois courbures ;
\item l'identité de cocycle prouve l'associativité de la concaténation ;
\item la non-scission de \(C_{108}\) prouve que le retour de phase élève la mémoire ;
\item le couple limite projective/limite inductive prouve la conservation de l'unité sous raffinement et lecture ;
\item la propriété universelle prouve que chaque sortie typée est une représentation de la même source, non une entrée conventionnelle dissimulée.
\end{enumerate}

\subsection{Conséquence conceptuelle : la réalité scalaire est une ombre fidèle mais incomplète}

La représentation archimédienne peut être fidèle à la valeur déterminée sans être fidèle à la généalogie. Deux trajectoires peuvent avoir la même coordonnée réelle et des retenues, hauteurs ou mémoires différentes. La structure discrète du continu résout cette multiplicité en conservant les fibres avant de les contracter. La phrase « un point devient deux » se traduit ainsi : une coordonnée scalaire peut admettre deux relèvements holonomiques distincts, et la bidirectionnalité réside dans la fibre, non dans une duplication arbitraire du point.

\subsection{Conséquence physique : l'unité ne conserve pas d'abord les symétries ; elle les produit}

Ce qui se conserve en premier lieu est la possibilité de reconstruction : feuillet, résidu, quotient, retenue, orientation, frontière et mémoire. Une symétrie apparaît lorsqu'une représentation possède des automorphismes qui préservent ces relations. Ainsi, les symétries exceptionnelles, la polarité électromagnétique et la pentapolarité gravitationnelle sont des formes différentes de stabilisation d'une frontière déjà engendrée. La conservation de la symétrie est une conséquence de la conservation généalogique.

\subsection{Conséquence informationnelle : croissance infinie à coût interne nul}

La complexité linéaire \(p(m)=m+1\) permet de produire une information non périodique illimitée sans croissance exponentielle de l'entropie. Le raffinement ajoute de la mémoire, mais la mise à jour intégrale est réversible. Le coût apparaît lorsqu'un lecteur décide d'effacer la fibre qui permettrait de reconstruire le pas. Cette combinaison — apériodicité, entropie topologique nulle, mémoire croissante et unité conservée — constitue la définition mathématique opératoire du cristal topologique-temporel HMT.

\section{VIII. Contributions autonomes issues de la lecture rétrospective}

\subsection{Première contribution : dualité limite projective--limite inductive}

La distinction entre la limite projective des états et la limite inductive des observables n'avait pas été formulée au niveau du noyau dans les développements antérieurs. Elle explique pourquoi approfondissement et projection sont des opérations conjuguées sans être des inverses triviaux ; elle corrige l'identification entre unité globale et section individuelle et fournit le cadre naturel de la prédéfinition.

\subsection{Deuxième contribution : algèbre universelle des trajectoires}

La présentation par générateurs et relations réunit les fibres TRIT, l'ordre APP, les transports TPK et le cocycle du registre de mémoire. Sa propriété universelle démontre que les lecteurs ultérieurs n'ont pas besoin de nouveaux générateurs s'ils préservent les relations.

\subsection{Troisième contribution : identité exponentielle trigéométrique}

L'exponentielle commune de \(J_+,J_0,J_-\) montre que la rotation complexe, le jet nilpotent et la dilatation hyperbolique sont trois composantes d'une seule loi. L'identité d'Euler est sa projection elliptique ; \(\varphi\) est son échelle hyperbolique ; \(e\) est la loi exponentielle ; et \(\alpha\) est la couture qui recolle le jet à l'holonomie dodécaphasique.

\subsection{Quatrième contribution : constantes comme spectre typé de représentations}

La notion de caractère scalaire s'étend à un spectre de représentations et de fonctionnelles. Cela permet de réunir constantes, particules et champs sans confondre leurs codomaines. La même définition situe Catalan comme moment, Feigenbaum comme remise à l'échelle non linéaire, Boltzmann comme transducteur, Barbero--Immirzi comme covecteur d'incidence, Planck comme section d'action et la gravité comme caractère contragrédient.

La récupération fonctionnelle explicite une relation supplémentaire : la réponse constitutive d'incidences \(90/120\), avec sa régularité à l'origine, détermine le moment de Catalan complet ; son inversion retrouve les canaux angulaires étiquetés. La composition des transports s'exprime alors par la loi \(r\odot t=f(\psi(r)\psi(t))\), en conservant produit, orientation et trace normalisée. Sur cette même section de vitesse, l'inscription et la norme radiale déterminent d'abord la longueur \(L\), puis les produits \(R_X\Lambda_X=L^2I\) et \(H_X\Lambda_X=\hbar_{\rm ret}cI\), et enfin le coefficient commun \(G=c^3L^2/\hbar_{\rm ret}\). Les preuves de \cref{cr28:sec:restitucion-constitutiva,g26:sec:rigidez-radial-es} conservent la mémoire et les hypothèses constitutives de chaque réalisation.

La constante de Feigenbaum reçoit une définition structurelle précise : elle est l’invariant asymptotique de séparation des paramètres sélectionnés par les retours critiques du lecteur dodécaphasique uniforme. La construction conserve sa provenance APP--TRIT--TPK, le continu conjoint et l’ordre des premières racines. L’identification des centres sélectionnés aux centres analytiques permet d’appliquer le théorème d’échelle à cette même suite ; les démonstrations de \ref{hmtfg:escalado:15}--\ref{hmtfg:escalado:17} et \eqref{hmtfg:articulacion:limite} établissent le résultat paramétrique. La représentation de Jacobi restitue le profil à partir de sa mesure à douze nœuds, tandis que l’inversion constitutive restitue les canaux orientés et le moment dont l’extrémité est la constante de Catalan. Ces opérations manifestent les relations entre lecteurs d’un même antécédent, en conservant le domaine et l’information propres à chacun.

\subsection{Cinquième contribution : symétrie comme stabilisateur de reconstruction}

L'intuition selon laquelle les symétries résultent de la conservation devient une définition du stabilisateur de bord et de relations. Cette formulation relie la double projection, le corridor exceptionnel et les représentations physiques sans postuler la symétrie comme source.

\subsection{Sixième contribution : irréductibilité triadique du noyau}

L'image du nœud borroméen proposée par l'auteur admet une formulation exacte au moyen de trois quotients d'oubli. Si les deux feuillets APP sont identifiés, le potentiel \(\Delta_{\Sigma\Pi}\) et le commutateur \([P_\Sigma,P_\Pi]\) disparaissent simultanément. Si les trois régimes TRIT sont identifiés, le spectre des signes disparaît

\[
J_+^2=-I,\qquad J_0^2=0,\qquad J_-^2=I.
\]

Si le cocycle TPK est trivialisé, la relation \(d_M(\Gamma_9x)-d_M(x)=1\) devient nulle et le retour de phase devient indiscernable du retour d'état. Chaque quotient détruit un invariant que les deux autres facteurs ne peuvent reconstruire. APP, TRIT et TPK sont donc trois générateurs fonctionnellement irréductibles de l'objet fondamental : l'APP apporte l'écart et l'ordre des feuillets ; TRIT apporte la signature locale de courbure ; TPK apporte composition, holonomie et mémoire.

Cette irréductibilité explique pourquoi le noyau ne peut se réduire ni à une table arithmétique, ni à une algèbre quadratique, ni à une dynamique des phases. La structure n'apparaît que dans leur composition.

\begingroup\fontsize{10.2}{11.7}\selectfont
\setlength{\parskip}{2pt}\setlength{\abovedisplayskip}{5pt}\setlength{\belowdisplayskip}{5pt}
\fontsize{9.5}{10.7}\selectfont
\setlength{\parskip}{1.5pt}\setlength{\abovedisplayskip}{2pt}\setlength{\belowdisplayskip}{2pt}
\subsection{Septième contribution : l'électron contient la paire nulle de Witt}

La comparaison rétrospective du bloc électronique avec les trois régimes TRIT révèle une clôture locale demeurée implicite. À partir de

\[
S^2=I,\qquad J^2=-I,\qquad SJ=-JS
\]

on forme

\[
n_\pm=\frac{S\pm J}{2}.
\]

Alors

\[
n_\pm^2=0,\qquad n_+n_-+n_-n_+=I.
\]

Le coin \({\rm Cl}_{1,1}\cong M_2(\mathbb R)\) contient donc les trois classes locales : \(J\) est elliptique, \(S\) est hyperbolique et \(n_\pm\) sont paraboliques. La lacune électronique est simultanément un défaut minimal de retenue dans l'APP et une direction nulle de l'algèbre déployée. TRIT n'introduit pas le seuil comme correction extérieure à l'électron : il reconnaît son type et le conserve sous transport TPK. La dénomination de Witt désigne ici la décomposition nulle locale ; la structure d'incidence Paley--Witt s'obtient ensuite comme représentation combinatoire de frontière de la même généalogie.

\subsection{Huitième contribution : relations de Weyl des deux horloges}

La phrase « l'hélice contient deux horloges » admet une présentation opératorielle exacte. Soit \(T\) le pas, \(V_9\) le caractère de phase nonadique, \(W_\delta\) le caractère de capacité irrationnelle, et \(\delta=\log_{10}(10/9)\). Si \(\zeta_9\) et \(\omega_\delta\) désignent leurs caractères internes, on a

\[
V_9T=\zeta_9TV_9,\qquad
W_\delta T=\omega_\delta TW_\delta.
\]

Pour \(\Gamma_9=T^9\),

\[
V_9\Gamma_9=\Gamma_9V_9,\qquad
W_\delta\Gamma_9=\omega_\delta^9\Gamma_9W_\delta.
\]

L'horloge de phase se ferme et l'horloge de capacité continue. Si \(D_M\) mesure la hauteur hélicoïdale,

\[
[D_M,\Gamma_9]=\Gamma_9.
\]

Cette dernière identité est la forme algébrique la plus brève de la différence entre retour de phase et retour d'état. Elle localise aussi précisément la non-commutativité : la translation basale appartient à un groupe abélien, mais son produit croisé avec les observables de phase et de capacité satisfait des relations de Weyl non commutatives. Le cristal apériodique HMT est ainsi une suspension ordonnée d'une horloge finie par un caractère irrationnel, relevée par un opérateur de mémoire.

Dans la représentation elliptique ultérieure, \(\zeta_9\) et \(\omega_\delta\) sont reconnus comme \(\operatorname{Exp}_+(2\pi{}J_+/9)\) et \(\operatorname{Exp}_+(2\pi{}\delta J_+)\). Le caractère HMT précède cette écriture complexe.

\subsection{Neuvième contribution : une équation polynomiale pour les trois géométries}

Les trois générateurs locaux peuvent être réunis dans

\[
\mathbb J=p_+J_++p_0J_0+p_-J_-.
\]

L’identité

\[
\boxed{\mathbb J^6=\mathbb J^2}
\]

contient exactement le spectre des carrés \(\{-I,0,I\}\). Si \(K=\mathbb J^2\), les régimes se retrouvent à partir de l'opérateur lui-même :

\[
p_+=\frac{K^2-K}{2},\qquad
p_0=I-K^2,\qquad
p_-=\frac{K^2+K}{2}.
\]

Cette formule dissipe l'apparence selon laquelle TRIT ne serait qu'une étiquette ternaire. Le régime est une décomposition spectrale intrinsèque. L'identité exponentielle correspondante

\[
\operatorname{Exp}_\star(t\mathbb J)=
p_+(\cos t+\mathbb J\sin t)+
p_0(I+t\mathbb J)+
p_-(\cosh t+\mathbb J\sinh t)
\]

contient comme faces la clôture elliptique d'Euler, le seuil nilpotent et l'auto-échelle hyperbolique. Sa composition avec \(\Gamma_9\), jointe à \([D_M,\Gamma_9]=\Gamma_9\), constitue l'identité opératorielle minimale à partir de laquelle peut être exposée l'évolution conservative de l'unité dimensionnelle.
  

\par\endgroup
\appendix
\label{app:desarrollos-10000-rev9}
% Copia mecánica íntegra de apendices.tex:1--761.
% Residencia material del ensamblaje sucesor de 102 capítulos.
\part{Appendices démonstratifs et tables de classification : certificats algébriques, arithmétiques et coinductifs}

\chapter{Calcul affine ordonné des feuillets APP}

\section{Produit semi-direct affine et composition des voies}

Soit \(U_{\APP}\) le module radial antérieur à son évaluation archimédienne. L'enveloppe correcte est le monoïde
\[
 \AffEnd(U_{\APP})
 =\operatorname{End}(U_{\APP})\ltimes U_{\APP},
\]
car les feuillets de seuil ou de pliage peuvent contenir des endomorphismes non inversibles. L'action de \((\Lambda,C)\) sur \(u\in U_{\APP}\) est
\[
 (\Lambda,C)\cdot u=\Lambda u+C.
\]
La composition se calcule sans commuter les facteurs :
\begin{align}
 (\Lambda_2,C_2)(\Lambda_1,C_1)\cdot u
 &=\Lambda_2(\Lambda_1u+C_1)+C_2 \notag\\
 &=(\Lambda_2\Lambda_1)u+(C_2+\Lambda_2C_1).
 \label{espiral:eq:ap-comp-afin}
\end{align}
Par conséquent,
\[
 (\Lambda_2,C_2)(\Lambda_1,C_1)
 =(\Lambda_2\Lambda_1,C_2+\Lambda_2C_1).
\]
La seconde coordonnée n'est pas une somme ordinaire de déplacements : c'est un cocycle tordu par l'action linéaire accumulée.

\begin{proposition}[Formule des préfixes]
Pour une voie de \(N\) facteurs \(h_j=(\lambda_j,C_j)\), ordonnés de \(j=0\) à \(j=N-1\), l'holonomie préfixe est
\[
 H_N=h_{N-1}\cdots h_1h_0=(\Lambda_N,\mathcal C_N),
\]
où
\[
 \Lambda_N=\prod_{j=0}^{N-1}\lambda_j
\]
et
\[
 \mathcal C_N
 =C_{N-1}+\lambda_{N-1}C_{N-2}
 +\lambda_{N-1}\lambda_{N-2}C_{N-3}
 +\cdots+
 \lambda_{N-1}\cdots\lambda_1C_0.
\]
\end{proposition}
\begin{proof}
Pour \(N=1\), la formule est immédiate. Si elle est valable pour \(N\), alors
\[
 h_NH_N
 =(\lambda_N\Lambda_N,C_N+\lambda_N\mathcal C_N),
\]
ce qui produit exactement l'expression de longueur \(N+1\). La récurrence préserve l'ordre de la voie et n'identifie pas les parcours ayant les mêmes extrémités.
\end{proof}

\section{Différence somme--produit comme courbure d'ordre}

Soient \(A_a=(I,a)\) et \(M_\lambda=(\lambda,0)\). Dans le secteur inversible,
\[
 A_aM_\lambda A_a^{-1}M_\lambda^{-1}
 =A_{a(1-\lambda)}.
\]
La démonstration directe sur \(u\) est
\[
 u\xmapsto{M_\lambda^{-1}}\lambda^{-1}u
 \xmapsto{A_a^{-1}}\lambda^{-1}u-a
 \xmapsto{M_\lambda}u-\lambda a
 \xmapsto{A_a}u+a(1-\lambda).
\]
Ainsi, la différence entre accumuler et plier est publiée sous forme de translation. Lorsque \(\lambda=1\), le commutateur s'annule ; lorsque \(\lambda\neq1\), la voie conserve une holonomie qui ne peut être reconstruite à partir du seul point final.

\begin{example}[Deux voies avec la même donnée visible]
Fixons \(u_0=0\), \(a=1\), \(\lambda=2\). Les voies
\[
 A_1M_2:u_0\longmapsto1,
 \qquad
 M_2A_{1/2}:u_0\longmapsto1
\]
publient la même valeur finale. Cependant, leurs registres contiennent des facteurs, déplacements et préfixes distincts. L'évaluation ponctuelle identifie les voies ; l'holonomie affine enrichie les sépare.
\end{example}

\section{Tableau exhaustif de la carte locale à 2 orientations}

Le tableau suivant ne postule pas que tout degré radial global provient d'une seule paire. Il enregistre la carte élémentaire obtenue en composant deux orientations TRIT. Le degré d'une voie générale est le cocycle entier de \cref{espiral:eq:cociclos-radiales} ; cette carte fournit ses neuf générateurs locaux.

\begin{longtable}{rrrrrrp{.27\textwidth}}
\caption{Les \(9\) paires orientées de la carte TRIT locale.}
\label{espiral:tab:ap-nueve-pares}\\
\toprule
\(\tau_+\)&\(\tau_-\)&\(p\)&\(a\)&\(r_3\)&\(c_3\)&Lectura local\\
\midrule
\endfirsthead
\toprule
\(\tau_+\)&\(\tau_-\)&\(p\)&\(a\)&\(r_3\)&\(c_3\)&Lectura local\\
\midrule
\endhead
\( 1\)&\( 1\)&\( 2\)&\( 0\)&\(-1\)&\( 1\)&ouverture radiale double avec retenue positive\\
\( 1\)&\( 0\)&\( 1\)&\( 1\)&\( 1\)&\( 0\)&ouverture et seuil\\
\( 1\)&\(-1\)&\( 0\)&\( 2\)&\( 0\)&\( 0\)&équilibre à orientation positive\\
\( 0\)&\( 1\)&\( 1\)&\(-1\)&\( 1\)&\( 0\)&seuil et ouverture\\
\( 0\)&\( 0\)&\( 0\)&\( 0\)&\( 0\)&\( 0\)&umbral doble\\
\( 0\)&\(-1\)&\(-1\)&\( 1\)&\(-1\)&\( 0\)&seuil et contraction\\
\(-1\)&\( 1\)&\( 0\)&\(-2\)&\( 0\)&\( 0\)&équilibre à orientation négative\\
\(-1\)&\( 0\)&\(-1\)&\(-1\)&\(-1\)&\( 0\)&contraction et seuil\\
\(-1\)&\(-1\)&\(-2\)&\( 0\)&\( 1\)&\(-1\)&contraction double avec retenue négative\\
\bottomrule
\end{longtable}

Ici \(p=r_3+3c_3\), avec le résidu équilibré \(r_3\in\{-1,0,1\}\). La retenue distingue \(p=2\) de son résidu \(-1\) et \(p=-2\) de son résidu \(1\). La paire \((p,a)\) distingue également les trois généalogies locales de degré \(0\).

\chapter{Certificat nonadique et géométrie de la mémoire}

\section{Distinction typée entre catalogue décimal, mots ternaires et fibre ambiante}

\begin{longtable}{p{.18\textwidth}p{.18\textwidth}p{.52\textwidth}}
\caption{Strates du certificat de la connexion nonadique conjointe.}
\label{espiral:tab:certificado-nonadico}\\
\toprule
Objet&Cardinalité&Fonction mathématique\\
\midrule
\endfirsthead
\toprule
Objet&Cardinalité&Fonction mathématique\\
\midrule
\endhead
Catalogue des triplets d'origine&\(104\,976\)&Domaine fini antérieur à la sélection conjointe ; ce n'est pas un tableau décimal.\\
Émissions décimales \(U_6\)&\(468\)&Émissions compatibles de longueur \(6\) dans la publication décimale.\\
Mots visibles \(w_6\)&\(243\)&Mots ternaires publiés après le typage ; ils ne constituent pas l'espace ambiant complet.\\
Espace ambiant \(\F_3^6\)&\(729\)&Espace de tous les mots ternaires de longueur \(6\) ; il contient les \(243\) visibles.\\
Niveaux certifiés&\(396\)&Prolongement fini jusqu'aux bords déclarés.\\
Trits recorridos&\(2\,376\)&Longitud total \(396\cdot6\).\\
Vueltas completas&\(43\)&Cycles nonadiques complets dans le segment certifié.\\
Paires \((K,K+9)\)&\(387\)&Paires de même phase par canal ; le registre peut différer.\\
\bottomrule
\end{longtable}

La chaîne de prolongement se conserve dans l'ordre
\[
 w_6\longrightarrow w_{12}\longrightarrow w_{18}
 \longrightarrow w_{24}\longrightarrow w_{30}
 \longrightarrow R_{36}\longrightarrow G_9.
\]
Aucun de ces objets ne peut être remplacé par le quotient modal ultérieur. L'holonomie \(\Gamma_9\) agit sur la dynamique enrichie ; la publication de mémoire réduite ne conserve que l'information spécifiée par son codomaine.

\section{Retour de phase, avancée de mémoire et état complet}

Avec
\[
 \rho_9(N)=1+((N-1)\bmod9),
 \qquad
 q_9(N)=\left\lfloor\frac{N-1}{9}\right\rfloor,
\]
on a
\[
 \rho_9(N+9)=\rho_9(N),
 \qquad q_9(N+9)=q_9(N)+1.
\]
Par conséquent,
\[
 H_9(N+9)=H_9(N)+(0,0,1).
\]
L'égalité de la première coordonnée de phase n'implique pas l'égalité de l'état :
\[
 g(K+9)=g(K),
 \qquad B_{K+9}\neq B_K
\]
en général. La formule hélicoïdale est précisément la représentation minimale de cette distinction.

\section{Extension dodécaphasique non scindée}

La suite
\[
 0\longrightarrow C_{12}\longrightarrow C_{108}
 \longrightarrow C_9\longrightarrow0
\]
couple le calendrier dodécaphasique au retour nonadique. S'il existait une section homomorphe globale \(s:C_9\to C_{108}\), le générateur de \(C_9\) devrait se relever en un élément d'ordre divisant \(9\) qui se projette sur le générateur. Le relèvement temporel actif conserve cependant le déplacement dodécaphasique accumulé et n'identifie pas l'état après un tour de phase. La couture est enregistrée sous forme de retenue ; géométriquement, elle se manifeste par une translation axiale de l'hélice.

\section{5 régions de \texorpdfstring{\(\pi\)}{pi} et non-fidélité de l'évaluation}

La microfibre de \(\pi\) contient cinq régions orientées de multiplicité totale
\[
 432+4\cdot144=1008.
\]
Les régions publient la même valeur réelle, mais conservent des généalogies distinctes. Ce résultat fournit le précédent interne de la classification spirale : une courbe visible ne détermine pas à elle seule la région, le feuillet, la retenue ni la mémoire qui l'ont produite.

\chapter{Algèbre puissance--logarithmique et 5 publications canoniques}

\section[Coordonnée puissance--logarithmique et inverse]{La coordonnée \(L_p\) et son inverse}

Pour \(x>0\), on définit
\[
 L_p(x)=
 \begin{cases}
 (x^p-1)/p,&p\neq0,\\
 \log x,&p=0.
 \end{cases}
\]
Son inverse est
\[
 \exp_p(u)=
 \begin{cases}
 (1+pu)^{1/p},&p\neq0,\ 1+pu>0,\\
 \e^u,&p=0.
 \end{cases}
\]
La dérivée
\[
 \frac{\dd}{\dd x}L_p(x)=x^{p-1}>0
\]
prouve que \(L_p\) est strictement croissante sur son domaine et donc inversible.

\begin{proposition}[Loi somme--produit]
Pour \(x,y>0\),
\[
 L_p(xy)=L_p(x)+L_p(y)+pL_p(x)L_p(y).
\]
\end{proposition}
\begin{proof}
Si \(p\neq0\),
\begin{align*}
 L_p(x)+L_p(y)+pL_p(x)L_p(y)
 &=\frac{x^p-1}{p}+\frac{y^p-1}{p}
   +\frac{(x^p-1)(y^p-1)}p\\
 &=\frac{x^py^p-1}{p}=L_p(xy).
\end{align*}
Pour \(p=0\), l'égalité se réduit à \(\log(xy)=\log x+\log y\).
\end{proof}

\begin{corollary}[Groupe déformé sur le domaine positif]
L'opération
\[
 u\boxplus_pv=u+v+puv
\]
satisfait
\[
 1+p(u\boxplus_pv)=(1+pu)(1+pv).
\]
Son neutre est \(0\) et l'inverse de \(u\) est
\[
 \ominus_pu=-\frac{u}{1+pu},
\]
si \(1+pu\neq0\).
\end{corollary}

\section{Dérivation des 5 courbes canoniques}

Soit \(m\) la profondeur publiée et supposons
\[
 L_p(r/r_0)=\beta m,
 \qquad
 \theta=\theta_0+\omega m,
 \qquad\omega\neq0.
\]
En éliminant \(m=(\theta-\theta_0)/\omega\),
\[
 r(\theta)=r_0
 \left[1+p\frac{\beta}{\omega}
 (\theta-\theta_0)\right]^{1/p},
 \qquad p\neq0,
\]
et
\[
 r(\theta)=r_0
 \exp\!\left[\frac{\beta}{\omega}
 (\theta-\theta_0)\right]
\]
pour \(p=0\). Les changements d'origine et les changements d'échelle positifs produisent les formes conventionnelles :

\begin{longtable}{rllp{.31\textwidth}}
\caption{Publications canoniques du secteur translationnel.}
\label{espiral:tab:cinco-publicaciones-detalle}\\
\toprule
\(p\)&Équation normale&Nom conventionnel&Domaine et caractéristique\\
\midrule
\endfirsthead
\toprule
\(p\)&Équation normale&Nom conventionnel&Domaine et caractéristique\\
\midrule
\endhead
\(2\)&\(r^2=a^2\theta\)&Fermat&\(\theta\ge0\) ; accroissement linéaire de l'aire radiale quadratique.\\
\(1\)&\(r=a\theta\)&Archimède&\(\theta\ge0\) ; accroissement radial constant par unité angulaire.\\
\(0\)&\(r=r_0\e^{b\theta}\)&logarithmique&\(\theta\in\R\) ; dilatation constante par unité angulaire.\\
\(-1\)&\(r=a/\theta\)&hyperbolique&\(\theta>0\) ; produit \(r\theta\) constant.\\
\(-2\)&\(r^2=a^2/\theta\)&lituus&\(\theta>0\) ; produit \(r^2\theta\) constant.\\
\bottomrule
\end{longtable}

\section{Inversion des feuillets}

Para \(p\neq0\),
\[
 L_{-p}(x^{-1})
 =\frac{x^p-1}{-p}=-L_p(x).
\]
L'identité vaut aussi pour \(p=0\). C'est pourquoi l'inversion radiale échange \(p\) avec \(-p\) et change le signe de la coordonnée linéarisée. Cette relation regroupe Fermat avec le lituus et Archimède avec l'hyperbolique comme orientations conjuguées du lecteur, sans les identifier comme courbes égales.

\section{Secteur affine général}

Si la coordonnée radiale satisfait
\[
 u_{n+1}=\Lambda u_n+C,
\]
l'itération interne exacte est
\[
 u_n=\Lambda^n u_0+\sum_{k=0}^{n-1}\Lambda^kC.
\]
Si \(I-\Lambda\) est inversible, \(u_*=(I-\Lambda)^{-1}C\) et \(u_n-u_*=\Lambda^n(u_0-u_*)\). Ce n'est que dans une réalisation scalaire positive, \(\chi\Lambda=\lambda\chi\) avec \(\lambda>0\), que l'on définit \(u_*=C/(1-\lambda)\). Avec \(\theta_n=\theta_0+n\omega\), \(\omega\neq0\), l'interpolation publiée dans cette réalisation est
\[
 r(\theta)=r_0\exp_p\!\left[
 u_*+\lambda^{(\theta-\theta_0)/\omega}(u_0-u_*)
 \right].
\]
La taxonomie complète n'est donc pas une liste de cinq noms. Elle contient au moins le degré \(p\), l'orientation \(a\), la classe affine \([\Lambda,C]\), le cocycle de retenue, le bord, l'incrément angulaire et le registre de mémoire.

\chapter{Naturalité, équivariance et non-fidélité}

\section{Le lecteur de préfixes}

Pour une histoire tronquée
\[
 \widehat\gamma_N=\widehat e_N\cdots\widehat e_1,
\]
le lecteur enrichi enregistre
\[
 \mathcal P_{\rad,N}^{\enr}(\widehat\gamma_N)
 =
 \left(
 (H_0,\Led_0);
 \bigl(p(\widehat e_j),a(\widehat e_j),P_j,A_j,
 h_{e_j},H_j,\Led_j\bigr)_{j=1}^{N}
 \right).
\]
Le terme \(H_j=h_{e_j}\cdots h_{e_1}\) conserve les préfixes. Le registre \(\Led_j\) conserve, selon la voie, le feuillet, le résidu, le quotient, la retenue, le bord, l'orientation, le cylindre, la signature, la survie et la mémoire. Les coordonnées réelles s'obtiennent ensuite :
\[
 \mathcal P_{\arq}^{\mathrm{esp}}
 =\ev_{\R^2}\circ\mathcal P_{\rad}^{\enr}.
\]

\section{Preuve complète de naturalité sous troncature}

Soit \(\pi_{N,M}\) l'effacement du suffixe de longueur \(N-M\), avec \(M<N\). Soit \(\operatorname{pr}_{N,M}\) la projection qui conserve l'enregistrement initial et les \(M\) premiers enregistrements d'arête. Alors
\[
 \operatorname{pr}_{N,M}\,
 \mathcal P_{\rad,N}^{\enr}(\widehat\gamma_N)
 =
 \left(
 (H_0,\Led_0);
 \bigl(p(\widehat e_j),a(\widehat e_j),P_j,A_j,
 h_{e_j},H_j,\Led_j\bigr)_{j=1}^{M}
 \right).
\]
Par définition de la troncature,
\[
 \pi_{N,M}(\widehat\gamma_N)=\widehat e_M\cdots\widehat e_1,
\]
et en appliquant de nouveau le lecteur, on obtient la même séquence. Par conséquent,
\[
 \operatorname{pr}_{N,M}\circ
 \mathcal P_{\rad,N}^{\enr}
 =
 \mathcal P_{\rad,M}^{\enr}\circ\pi_{N,M}.
\]
La démonstration n'utilise pas la valeur de la courbe. Elle utilise l'identité des préfixes, qui est l'information que le lecteur réduit à sept coordonnées aurait perdue.

\section{Raffinements non triviaux}

Pour un raffinement \(r:\gamma\to\gamma'\), la naturalité requiert une application \(\widehat r\) sur les registres et l'égalité locale
\[
 h_e
 =
 h_{e'_k}\cdots h_{e'_1}
\]
pour chaque arête \(e\) remplacée par la chaîne raffinée \((e'_1,\ldots,e'_k)\). Si de plus la retenue, le bord et la mémoire de la chaîne raffinée composent les enregistrements de l'étape originale, alors
\[
 \widehat r\circ\mathcal P_{\rad}^{\enr}
 =
 \mathcal P_{\rad}^{\enr}\circ r.
\]
Cette condition est un carré vérifiable ; elle ne doit pas être supposée pour un raffinement arbitraire sans présenter les facteurs.

\section{Équivariance, non invariance}

L'action de \(\Gamma_9\) ramène la phase et prolonge l'état. Dans le codomaine du lecteur agit \(\widehat\Gamma_9^{\rad}\), qui déplace les enregistrements et actualise le dernier préfixe. L'identité correcte est
\[
 \mathcal P_{\rad}^{\enr}(\Gamma_9\gamma)
 =
 \widehat\Gamma_9^{\rad}
 \mathcal P_{\rad}^{\enr}(\gamma).
\]
Exiger l'invariance effacerait précisément la mémoire que la connexion conserve. Dans le secteur réversible, l'inversion des voies satisfait
\[
 H_{C_{\mathrm{rev}}\gamma}=H_\gamma^{-1}
\]
lorsque tous les facteurs appartiennent au sous-groupe affine inversible. Cette affirmation typée ne s'étend pas automatiquement à une involution globale sur n'importe quel quotient ultérieur.

\section{3 démonstrations de non-fidélité}

\begin{proposition}[Cercle intrinsèque et cercle projeté]
Il existe deux histoires ayant la même courbe visible
\[
 c(\theta)=(\cos\theta,\sin\theta)
\]
et des registres distincts.
\end{proposition}
\begin{proof}
La première histoire possède une mémoire axiale constante \(m=0\). La seconde est l'hélice
\[
 h(\theta)=(\cos\theta,\sin\theta,\theta/2\pi),
\]
dont le registre enregistre \(m(\theta)=\theta/2\pi\). Toutes deux satisfont \(\pi_{xy}\circ h=c\). La projection est donc non injective.
\end{proof}

\begin{proposition}[Même extrémité, ordre affine différent]
Si \(a(1-\lambda)\neq0\), les chemins \(A_aM_\lambda\) et \(M_\lambda A_a\) ne sont pas équivalents dans le lecteur enrichi, bien qu'ils puissent être calibrés pour publier une même extrémité.
\end{proposition}
\begin{proof}
Leur commutateur est \(A_{a(1-\lambda)}\neq I\). La donnée du commutateur appartient au registre de préfixes et disparaît lorsqu'on ne conserve que l'extrémité.
\end{proof}

\begin{proposition}[Même courbe, région généalogique différente]
L'égalité de la valeur réelle publiée n'identifie pas les cinq régions de la microfibre de \(\pi\) ; elle ne suffit donc pas non plus à identifier la classe spirale lorsque le lecteur conserve la région, le bord ou la retenue.
\end{proposition}

\chapter{Taxonomie stratifiée des correspondances spirales}

\section{Classification par niveau mathématique}

\begin{longtable}{p{.20\textwidth}p{.30\textwidth}p{.39\textwidth}}
\caption{Niveaux qui ne doivent pas être confondus en une seule notion de courbure.}
\label{espiral:tab:niveles-curvatura}\\
\toprule
Niveau&Invariant&Signification\\
\midrule
\endfirsthead
\toprule
Niveau&Invariant&Signification\\
\midrule
\endhead
APP&\(F(S,P)=1,\ F(P,S)=-1\)&Courbure discrète de l'ordre d'accumulation et de pliage.\\
TRIT&\(\tau\in\{-1,0,1\}\), \(J_\tau^2=-\tau I\)&Régime local elliptique, parabolique ou hyperbolique du transport.\\
TPK&\(\Gamma_9\), retenue, bord, voie et registre&Holonomie de phase et mémoire sur l'état holonomique intégral.\\
Radial&\(p,a,[\Lambda,C]\)&Caractères qui déterminent la loi radiale avant de reconnaître une courbe.\\
Archimédien&\(r(\theta)\)&Courbe publiée dans \(\R^2\) ; elle peut oublier la généalogie.\\
Frenet&\(\kappa_{\mathrm{F}},\mathcal T_{\mathrm{F}}\)&Courbure et torsion différentielles de la courbe ou de la suspension visible.\\
Dimensionnel&champs, échelles et représentations&Réalisation physique ultérieure, avec ses propres prémisses.\\
\bottomrule
\end{longtable}

\section{Classification par rotation, rayon et mémoire}

\begin{longtable}{ccccl}
\caption{Taxonomie cinématique complète des cas dégénérés minimaux.}
\label{espiral:tab:taxonomia-cinematica-ap}\\
\toprule
\(\Delta\theta\)&\(\Delta u_{\rad}\)&\(\Delta m\)&Salida visible&
Information généalogique\\
\midrule
\endfirsthead
\toprule
\(\Delta\theta\)&\(\Delta u_{\rad}\)&\(\Delta m\)&Salida visible&
Information généalogique\\
\midrule
\endhead
\(0\)&\(0\)&\(0\)&point&état stationnaire publié\\
\(0\)&\(\neq0\)&\(0\)&recta radial&changement d'échelle\\
\(0\)&\(0\)&\(\neq0\)&fibra axial&mémoire sans déplacement transverse\\
\(0\)&\(\neq0\)&\(\neq0\)&feuillet radial suspendu&échelle et mémoire\\
\(\neq0\)&\(0\)&\(0\)&cercle&fermeture angulaire intrinsèque\\
\(\neq0\)&\(0\)&\(\neq0\)&hélice cylindrique&phase fermée, mémoire relevée\\
\(\neq0\)&\(\neq0\)&\(0\)&espiral plana&caractère radial visible\\
\(\neq0\)&\(\neq0\)&\(\neq0\)&suspension hélicoïdale radiale&rotation, rayon et profondeur conservés\\
\bottomrule
\end{longtable}

\section{Classification par holonomie affine}

Le tableau suivant appartient à une réalisation scalaire réelle de la partie linéaire, \(\chi\Lambda=\lambda\chi\) ; pour en alléger la lecture, on écrit \(\Lambda\) dans la colonne des conditions. Il ne classe pas des endomorphismes abstraits au moyen d'inégalités numériques.

\begin{longtable}{p{.18\textwidth}p{.22\textwidth}p{.48\textwidth}}
\caption{Types d'holonomie sur la coordonnée radiale.}
\label{espiral:tab:tipos-holonomia}\\
\toprule
Condition&Orbite de \(u\)&Interprétation\\
\midrule
\endfirsthead
\toprule
Condition&Orbite de \(u\)&Interprétation\\
\midrule
\endhead
\(\Lambda=1,\ C=0\)&\(u_n=u_0\)&rayon constant ; la mémoire peut être nulle ou cachée par projection.\\
\(\Lambda=1,\ C\neq0\)&\(u_n=u_0+nC\)&secteur translationnel qui publie la famille \(L_p\).\\
\(0<\Lambda<1\)&\(u_n\to u_*\)&contraction vers une section fixe.\\
\(\Lambda>1\)&\(|u_n-u_*|\) croît&expansion à partir d'une section fixe.\\
\(\Lambda<0\)&alternance d'orientation&requiert de contrôler le domaine de \(\exp_p\) à chaque étape.\\
\(\Lambda=0\)&\(u_{n+1}=C\)&pliage non inversible ; il appartient au monoïde, non au groupe affine.\\
\bottomrule
\end{longtable}

\section{Critère d'équivalence généalogique}

Deux voies ne sont pas identifiées parce qu'elles publient la même courbe. On exige un isomorphisme typé qui préserve terme à terme :
\[
 \bigl(
 (x_0,H_0,\Led_0);
 (\widehat e_j,p(\widehat e_j),a(\widehat e_j),P_j,A_j,
 h_{e_j},H_j,\Led_j)_{j=1}^{N}
 \bigr)
\]
et qui commute avec les troncatures et raffinements admis. La classification résultante est plus fine que la taxonomie conventionnelle car elle incorpore l'histoire qui précède l'apparition de \(\R^2\).

\chapter{Réalisations ultérieures et contrôles de type}

\section{Mode électromagnétique circulaire}

Soit
\[
 \zeta=\frac{z-c_{\mathrm{vac}}t}{\lambda},
\]
et
\[
 \mathbf E_\sigma
 =E_0(\mathbf e_1\cos2\pi\zeta
 +\sigma\mathbf e_2\sin2\pi\zeta),
 \qquad
 \mathbf B_\sigma
 =c_{\mathrm{vac}}^{-1}\widehat{\mathbf k}\times\mathbf E_\sigma.
\]
Alors
\[
 \mathbf E_\sigma\perp\mathbf B_\sigma\perp\widehat{\mathbf k},
 \quad |E|=c_{\mathrm{vac}}|B|,
 \quad\omega=c_{\mathrm{vac}}k.
\]
À temps fixé, l'extrémité normalisée du champ dans le fibré phase--position décrit
\[
 (\cos2\pi\zeta,\sin2\pi\zeta,\zeta).
\]
Cette affirmation ne dit pas que toute ligne spatiale de champ est une hélice : elle décrit la trajectoire de l'extrémité du vecteur polarisation avec l'événement marqué.

\section{Compatibilité loxodromique}

Pour
\[
 \Phi(r,\varphi,z,t)
 =\ell\varphi+\nu\log(r/r_0)+k_zz-\omega t,
\]
la transformation
\[
 (r,\varphi)\longmapsto
 ((2+\sqrt3)r,\varphi+\pi/2)
\]
préserve la phase modulo \(2\pi\) exactement lorsque
\[
 \nu\log(2+\sqrt3)+\ell\frac{\pi}{2}\in2\pi\Z.
\]
La dilatation s'interprète comme un changement d'échelle entre solutions ; non comme une production spontanée d'énergie.

\section{Centre électronique et lecteur ultérieur}

L'électron occupe l'unique point fixe \((5,5)\) de l'atlas matériel. Cette centralité n'est pas dérivée de nouveau dans le présent travail. La compatibilité opératorielle antérieure comprend
\[
 C^{(3)}M_{\ph}=-\frac12M_{\ph},
 \qquad
 1-\left(\frac12\right)^2=\frac34=\frac{90}{120},
\]
et
\[
 \|[P_{\Sigma,3},P_{\Pi,3}]\|=\frac{\sqrt3}{4}.
\]
La théorie des correspondances spirales ajoute un lecteur de voies ; elle ne remplace pas la Loi générale des masses et n'utilise pas l'électron pour sélectionner rétrospectivement \(p\), \(\Lambda\) ou \(C\).

\section{Matrice fermée de contrôles négatifs}

\begin{longtable}{p{.38\textwidth}p{.50\textwidth}}
\caption{Réfutateurs et erreurs de type exclus par la construction.}
\label{espiral:tab:falsadores-ap}\\
\toprule
Défaillance&Conséquence\\
\midrule
\endfirsthead
\toprule
Défaillance&Conséquence\\
\midrule
\endhead
La courbe classique sélectionne \(p,\Lambda,C\).&Circularité : le lecteur cesse d'être constructif--génératif.\\
Seule l'holonomie finale est conservée.&Perte des préfixes, de la retenue et du bord ; la naturalité n'est pas attestée.\\
L'invariance sous \(\Gamma_9\) est exigée.&L'avancée de mémoire est effacée ; la relation correcte est l'équivariance.\\
On confond \(U_6\), \(w_6\) et \(\F_3^6\).&Confusion d'objets de cardinalités \(468\), \(243\) et \(729\).\\
Le retour de phase est identifié au retour d'état.&Contradiction avec \(B_{K+9}\neq B_K\) en général.\\
On identifie \(\tau\), \(p\) et \(\kappa_{\mathrm{F}}\).&Confusion du régime, du caractère radial et de la courbure visible.\\
Tout rayon variable est qualifié d'hélice générale.&Affirmation différentielle sans vérification de la condition de Lancret.\\
L'axe \(\varphi\) est identifié à l'électron.&Mélange de la fibre des constantes et de la réalisation matérielle.\\
Pell est interprété comme une amplification énergétique.&Violation de la lecture d'échelle et de la conservation physique.\\
L'hélicité électromagnétique est identifiée au spin \(1/2\).&Confusion entre représentations distinctes.\\
Une cellule ponctuelle est imposée à toute espèce.&Contradiction avec le caractère de fibre ou de multisection des espèces non centrales.\\
Les cinq constructions du continu sont séparées.&Rupture de la généalogie unique antérieure à toute publication géométrique.\\
\bottomrule
\end{longtable}

\chapter{Glossaire mathématique}

\begin{longtable}{p{.22\textwidth}p{.66\textwidth}}
\caption{Symboles et termes directeurs.}
\label{espiral:tab:glosario}\\
\toprule
Symbole ou terme&Définition\\
\midrule
\endfirsthead
\toprule
Symbole ou terme&Définition\\
\midrule
\endhead
\(\Sigma,\Pi\)&Évaluations additive et multiplicative sur le support commun APP.\\
\(\tau\)&Régime local d'orientation TRIT ; ce n'est ni le degré radial ni la torsion de Frenet.\\
\(M_{\ph}\)&Sélecteur tritique de phase opératoire.\\
\(\Gamma_9\)&Holonomie nonadique de la dynamique enrichie.\\
\(H_N\)&Holonomie affine accumulée jusqu'au préfixe \(N\).\\
\(\Lambda_N,\mathcal C_N\)&Partie linéaire et cocycle translationnel de \(H_N\).\\
\(p\)&Cocycle entier du caractère radial.\\
\(a\)&Caractère orienté local qui sépare les généalogies de même degré local.\\
\(L_p,\exp_p\)&Coordonnée puissance--logarithmique et son inverse positive.\\
\(\boxplus_p\)&Loi déformée qui publie l'interaction somme--produit.\\
\(\Led_N\)&Registre enrichi de feuillet, retenue, bord, voie et mémoire.\\
\(\mathcal P_{\rad}^{\enr}\)&Lecteur radial enrichi antérieur à la courbe.\\
\(\mathcal P_{\arq}^{\mathrm{esp}}\)&Évaluation ultérieure dans une courbe de \(\R^2\).\\
\(\kappa_{\mathrm{F}},\mathcal T_{\mathrm{F}}\)&Courbure et torsion de Frenet de la courbe visible.\\
Suspension hélicoïdale radiale&Surface ou feuillet à rayon variable qui conserve la phase, la profondeur et le parcours.\\
Publication non fidèle&Évaluation qui peut identifier des histoires distinctes.\\
Constructivo--generativo&Méthode qui construit l'objet et sa généalogie avant de reconnaître son apparence conventionnelle.\\
\bottomrule
\end{longtable}
  \chapter{Dictionnaire des domaines, opérateurs et représentations}
\chapter*{Notation, types et désambiguïsation des symboles}
\addtocontents{toc}{\protect\enlargethispage{2\baselineskip}}
\addcontentsline{toc}{chapter}{Notation, types et désambiguïsation des symboles}

\begin{longtable}{@{}p{.22\textwidth}p{.70\textwidth}@{}}
\toprule
Symbole & Objet \\ \midrule
\endhead
\(\mathcal A_9\)&Tore arithmétique nonadique de chemins avec deux évaluations
internes.\\
\(\mathrm{TRIT}\)&Coordonnée ternaire orientée du transport ; la sortie
visible n'épuise pas le relèvement enrichi avec retenue.\\
\(\mathrm{TPK}\)&Noyau topologique de phase avec état holonomique intégral, transitions, parcours, prolongements et mémoire.\\
\(M_{\rm ph}\)&Sélecteur ternaire de phase opérative, distinct du transport
cardinal et de l'échantillonnage stroboscopique.\\
\(\tau_3\)&Bloc primitif de phase \((+1,-1,0)\).\\
\(\mathcal N,\mathcal E,\mathcal S,\mathcal O\)&Générateurs du transport
cardinal.\\
\(\operatorname{Obs}_3\)&Échantillonnage stroboscopique de chaque troisième pas.\\
\(\mathcal R_{12}\)&Système temporel dodécaphasique orienté interne au TPK.\\
\(\Theta_{108}=\mathbb Z/108\mathbb Z\)&Calendrier observable
\(12\times9\) ; ce n'est pas l'état TPK complet.\\
\(C_{12}\)&Rotation des douze secteurs temporels.\\
\(\Gamma_{108}\)&Orbite orientée de transport de \(108\) pas.\\
\(\mathcal G_9\)&Élagage unitaire dans la relation d'extension et de
frontière ; il ne constitue pas à lui seul la monodromie.\\
\(\mathfrak M_9\)&Monodromie nonadique avec mémoire : la phase revient et
l'état change.\\
\(\pi_{\rm term}\)&Projection de l'état holonomique terminal sur les trois blocs visibles sélectionnés.\\
\(R_{\rm sgn}\)&Lecture de la coordonnée entière signée du même état holonomique terminal.\\
\(U_{12}^{\rm cat}\)&Concaténation visible \(U_6\Vert U_6\).\\
\(U_{12}^{\rm sgn}\)&Coordonnée signée de l'état holonomique intégral, réversible avec \(K\) au moyen de \(\mathcal H_{12}\) ; elle ne se déduit ni de \(U_{12}^{\rm cat}\), ni de la projection calendaire, ni des trois blocs terminaux isolés.\\
\(\mathcal H_{12}\)&Transformation dodécaphasique réversible entre
\(U_{12}^{\rm sgn}\) et \(K\).\\
\(\mathcal E_{\rm dod}\)&Extracteur intégral
\((D_3,D_4,Q):\mathbb Z^{12}\to\mathbb Z^{25}\) qui reconstruit \(K\).\\
\(\overline{\mathcal C}_9^{\,3}\)&Complétion nonadique cubique.\\
\(A,C^*\)&Coordonnées angulaire et contre-angulaire intrinsèques.\\
\(q_\pm\)&Canaux orientés du lecteur bicouche.\\
\(D_{\rm raw}^{90,120}\)&Différence diagnostique brute de séries dans une
carte fixée ; ce n'est pas un extracteur canonique.\\
\(\iota_{6,8}\)&Inclusion
\(\mathcal F_6\hookrightarrow\mathcal F_8\) d'incidences hexade--octade.\\
\(K_{6\to8}\)&Fonctionnelle angulaire des deux sous-queues de Lambert, postérieure à
\(q_\pm\), aux séries, au pôle, à la positivité et à la minimalité.\\
\(\langle90,120\rangle\)&Semi-groupe spectral qui gouverne la hiérarchie des
tours.\\
\(\varepsilon_{\rm res}^{\circ}\)&Résidu de l'identité angulaire exacte ;
ce n'est pas une extraction finie des objets 90/120 précédents.\\
\(\Delta_4\)&Torsion minimale ou invariant torsionnel global antérieur à
l'espèce ; ce n'est pas la torsion de Cartan.\\
\(\mathsf T_{\alpha,\infty}\)&Prolongement distant de la retenue de la voie
entière de \(\alpha\), typé séparément du générateur orbital
\(T_1^{\rm orb}=L_1\), du compteur neutre et du calendrier.\\
\(N_1^{\rm neu}\)&Nombre de positions neutres dans la première fenêtre de
l'exemple de transport ; il vaut $3$ et n'est pas un opérateur.\\
\(T_1^{\rm orb}:=L_1\)&Second générateur dans la jauge orbitale locale ; dans sa
section de référence, il est abrégé \(T_1\).\\
\(\varepsilon_{\rm clk}\)&Quantum énergétique de l'horloge,
\(\hbar_{\rm int}2\pi/(108t_0)\) ; ce n'est pas une base énergétique canonique sans
déclarer le cadre temporel. Le symbole \(\varepsilon_0\) est réservé à
la permittivité du vide dans l'interface constitutive SI.\\
\(\mathrm{CLK}1\text{--}\mathrm{CLK}6\)&Étiquettes des six équations de
l'horloge finie, du hamiltonien à Schrödinger.\\
\(\mathcal K_{27}\)&Retour composé de vingt-sept mises à jour.\\
\(X_\infty^{\rm cal}\)&Domaine enrichi commun des évaluations
globales.\\
\bottomrule
\end{longtable}

\section*{Cardinaux répétés avec des sources différentes}

L'égalité d'un cardinal ne transporte pas automatiquement le domaine,
l'opération, la mémoire ni la force probante. Le tableau suivant fixe les
collisions qui affectent le plus la lecture.

\begin{longtable}{@{}p{.12\textwidth}p{.27\textwidth}p{.51\textwidth}@{}}
\toprule
Indice & Objeto tipado & Fonction et séparation \\ \midrule
\endhead
\(108\)&\(\Theta_{108}\)&Calendrier \(12\times9\) ; il ferme l'horloge
observable.\\
\(108\)&\(\Pi_{108}^{\rm cal}\)&Projection calendaire d'un curseur avec \(324\)
germes ; elle oublie le second curseur, l'orientation, la retenue, la frontière, le registre
chronologique enrichi et la
survie.\\
\(108\)&\(w_{108}\)&Mot de \(108=18\times6\) trits ; ce n'est pas le calendrier.\\
\(104,976\)&\(X_{\rm TPK}^{\rm seed}\)&Moteur à deux curseurs,
\(104,976=324^2\) ; il conserve un domaine distinct du calendrier.\\
\(1080\)&\(w_{1080}\)&Longueur d'un mot.\\
\(1080\)&\(\Gamma_0(1080)\)&Niveau d'un sous-groupe modulaire ; il exige sa
réalisation propre.\\
\(729\)&\(\mathbb F_3^6\)&Espace ambiant des mots ternaires.\\
\(729\)&\(9^3\)&Cellule cubique/intérieur EMRH.\\
\(729\)&\(\operatorname{Ext}(X_K)\)&Nombre de sous-cylindres examinés à chaque
étape de prolongement.\\
\(729\)&bulk EMRH&Strate intérieure de la complétion en base mille.\\
\(27\)&\(\mathcal K_{27}\)&Opérateur de retour composé ; ce n'est pas le nom d'une
théorie autonome.\\
\(27\)&strate cubique&Contribution de deux coordonnées de frontière dans
\(999=729+243+27\).\\
\(27\)&module ou classificateur&Chaque occurrence exige son application avant
d'être identifiée à une autre.\\
\bottomrule
\end{longtable}

\section*{Six sources distinctes de 90/120}

\begin{longtable}{@{}
>{\raggedright\arraybackslash}p{.21\textwidth}
>{\raggedright\arraybackslash}p{.42\textwidth}
>{\raggedright\arraybackslash}p{.27\textwidth}@{}}
\toprule
Objet & Domaine et opération & Sortie et interdiction de type \\ \midrule
\endhead
\(\mathcal E_{\rm dod}\)&\(K\in\mathbb Z^{12}\) ; différences par pas
\(3,4\) et somme \(Q\)&Coordonnées transversales dodécaphasiques ; ni incidence
ni série angulaire.\\
\(D_{\rm raw}^{90,120}\)&Une carte \(q\) ; différence brute
\(S_{90}(q)-S_{120}(q)\)&Diagnostic d'échelle, distinct de la fonctionnelle de
Barbero--Immirzi et des opérateurs d'incidence.\\
\(\iota_{6,8}\)&\(H\subset O\) ; inclusion de
\(H\times\binom{H}{2}\) dans \(O\times\binom{H}{2}\)&Cardinaux \(90,120\) et
défaut \(-1/12\) ; ni extracteur ni série.\\
\(K_{6\to8}\)&Canaux \(q_\pm\) et sous-queues de Lambert ; combinaison
\(12\Delta S_{90}-\Delta S_{120}\)&Fonctionnelle angulaire relative à ses
prémisses ; elle ne se transporte pas comme un bloc de masse.\\
\(\mathfrak S=\langle90,120\rangle\)&Hauteurs des moments
\(s_n=q_-^n-q_+^n\)&Semi-groupe global de tours ; il ne s'arrête pas à huit hauteurs.\\
\bottomrule
\end{longtable}

\section*{Quatre réalisations de la réversibilité bidirectionnelle}

\begin{longtable}{@{}
  >{\raggedright\arraybackslash}p{.26\textwidth}
  >{\raggedright\arraybackslash}p{.26\textwidth}
  >{\raggedright\arraybackslash}p{.38\textwidth}@{}}
\toprule
Opération & Domaine & Action \\ \midrule
\endhead
\(P_\pm=(I\pm R)/2\)&Algèbre scindée&Projette les deux secteurs
spectraux.\\
Inversion de direction&Route et calendrier composé&Inverse l'orientation du
transport sans effacer son mot.\\
\(\operatorname{Ext}^{\pm}\)&Arbre des prolongements&Compare les valences
futures et passées ; leur bilan peut être lu au moyen de
\(\log(\#\operatorname{Ext}^+/\#\operatorname{Ext}^-)\).\\
\(C_M(\tau)=-\operatorname{rev}(\tau)\)&Mot dodécaphasique&Conjugue l'ordre
et le signe ; ce n'est pas l'involution cellulaire \(\rho_c\).\\
\bottomrule
\end{longtable}

La fonction commune de ces quatre réalisations est de conserver une orientation
réversible. Elle n'autorise pas à identifier retour de phase, retour de route,
retour de mot et retour d'état.
  
\backmatter
\begingroup
\sloppy
\small
\setlength{\parskip}{0pt}
\setlength{\emergencystretch}{3em}
\bibliographystyle{alphaurl}
\begin{thebibliography}{HvdHKN76}

\bibitem[ALY17]{AbeLamYamada2017}
Toshiyuki Abe, Ching~Hung Lam, and Hiromichi Yamada.
\newblock A remark on {$\mathbb Z_p$}-orbifold constructions of the moonshine vertex operator algebra.
\newblock {\em arXiv preprint arXiv:1705.09022}, 2017.
\newblock \href {https://arxiv.org/abs/1705.09022} {\path{arXiv:1705.09022}}.

\bibitem[AW68]{rm:ArakiWoods1968}
H.~Araki and E.~J. Woods.
\newblock A classification of factors.
\newblock {\em Publ. Res. Inst. Math. Sci. Ser. A}, 4:51--130, 1968.
\newblock URL: \url{https://ems.press/content/serial-article-files/41616}.

\bibitem[Bek73]{rm:Bekenstein1973}
J.~D. Bekenstein.
\newblock Black holes and entropy.
\newblock {\em Physical Review D}, 7:2333--2346, 1973.

\bibitem[BG95a]{Barbero}
J.~Fernando Barbero~G.
\newblock Real {Ashtekar} variables for {Lorentzian} signature space-times.
\newblock {\em Physical Review D}, 51:5507--5510, 1995.

\bibitem[BG95b]{Barbero1995}
J.~Fernando Barbero~G.
\newblock Real {Ashtekar} variables for {Lorentzian} signature space-times.
\newblock {\em Physical Review D}, 51:5507--5510, 1995.

\bibitem[Bor92]{Borcherds1992}
Richard~E. Borcherds.
\newblock Monstrous moonshine and monstrous lie superalgebras.
\newblock {\em Inventiones Mathematicae}, 109:405--444, 1992.

\bibitem[{Bur}19]{BIPMSI}
{Bureau International des Poids et Mesures}.
\newblock {\em The International System of Units ({SI} Brochure)}.
\newblock BIPM, S\`evres, 9 edition, 2019.
\newblock Updated online edition.
\newblock \href {https://doi.org/10.59161/AUEZ1291}
  {\path{doi:10.59161/AUEZ1291}}.

\bibitem[Car14]{Caratheodory1914}
Constantin Carath{\'e}odory.
\newblock {"U}ber das lineare ma{\ss} von punktmengen---eine verallgemeinerung
  des l{"a}ngenbegriffs.
\newblock {\em Nachrichten von der Gesellschaft der Wissenschaften zu
  G{"o}ttingen, Mathematisch-Physikalische Klasse}, pages 404--426, 1914.

\bibitem[Car17]{Carnahan2017}
Scott Carnahan.
\newblock 51 constructions of the moonshine module.
\newblock {\em arXiv preprint arXiv:1707.02954}, 2017.
\newblock \href {https://arxiv.org/abs/1707.02954} {\path{arXiv:1707.02954}}.

\bibitem[CLS16]{ChenLamShimakura2016}
Hsian-Yang Chen, Ching~Hung Lam, and Hiroki Shimakura.
\newblock A {$\mathbb Z_3$}-orbifold construction of the moonshine vertex
  operator algebra and some maximal 3-local subgroups of the monster.
\newblock {\em arXiv preprint arXiv:1606.05961}, 2016.
\newblock \href {https://arxiv.org/abs/1606.05961} {\path{arXiv:1606.05961}}.

\bibitem[CN79]{ConwayNorton1979}
John~H. Conway and Simon~P. Norton.
\newblock Monstrous moonshine.
\newblock {\em Bulletin of the London Mathematical Society}, 11:308--339, 1979.

\bibitem[{COD}24]{CODATA2022}
{CODATA Task Group on Fundamental Constants}.
\newblock The 2022 {CODATA} recommended values of the fundamental physical
  constants, 2024.
\newblock URL: \url{https://physics.nist.gov/constants}.

\bibitem[CS99]{ConwaySloane1999}
John~H. Conway and Neil J.~A. Sloane.
\newblock {\em Sphere Packings, Lattices and Groups}.
\newblock Springer, 3 edition, 1999.

\bibitem[Fei78]{rm:Feigenbaum1978}
Mitchell~J. Feigenbaum.
\newblock Quantitative universality for a class of nonlinear transformations.
\newblock {\em Journal of Statistical Physics}, 19(1):25--52, 1978.
\newblock \href {https://doi.org/10.1007/BF01020332}
  {\path{doi:10.1007/BF01020332}}.

\bibitem[FLM88]{FLM1988}
Igor Frenkel, James Lepowsky, and Arne Meurman.
\newblock {\em Vertex Operator Algebras and the Monster}.
\newblock Academic Press, 1988.

\bibitem[GH77]{rm:GibbonsHawking1977}
G.~W. Gibbons and S.~W. Hawking.
\newblock Cosmological event horizons, thermodynamics, and particle creation.
\newblock {\em Physical Review D}, 15:2738--2751, 1977.

\bibitem[Haw75]{rm:Hawking1975}
S.~W. Hawking.
\newblock Particle creation by black holes.
\newblock {\em Communications in Mathematical Physics}, 43:199--220, 1975.

\bibitem[HJ13]{HornJohnson2013}
Roger~A. Horn and Charles~R. Johnson.
\newblock {\em Matrix Analysis}.
\newblock Cambridge University Press, 2 edition, 2013.

\bibitem[Hol96]{Holst1996}
Sören Holst.
\newblock Barbero's hamiltonian derived from a generalized hilbert--palatini
  action.
\newblock {\em Physical Review D}, 53:5966--5969, 1996.

\bibitem[HvdHKN76]{Hehl1976}
Friedrich~W. Hehl, Paul von~der Heyde, G.~David Kerlick, and James~M. Nester.
\newblock General relativity with spin and torsion: Foundations and prospects.
\newblock {\em Reviews of Modern Physics}, 48:393--416, 1976.

\bibitem[IK02]{IosifescuKraaikamp}
M.~Iosifescu and C.~Kraaikamp.
\newblock {\em Metrical Theory of Continued Fractions}.
\newblock Kluwer, 2002.

\bibitem[Imm97a]{Immirzi}
Giorgio Immirzi.
\newblock Real and complex connections for canonical gravity.
\newblock {\em Classical and Quantum Gravity}, 14:L177--L181, 1997.

\bibitem[Imm97b]{Immirzi1997}
Giorgio Immirzi.
\newblock Real and complex connections for canonical gravity.
\newblock {\em Classical and Quantum Gravity}, 14:L177--L181, 1997.

\bibitem[Khi64]{Khinchin}
A.~Ya. Khinchin.
\newblock {\em Continued Fractions}.
\newblock University of Chicago Press, 1964.

\bibitem[Kib61]{Kibble1961}
T.~W.~B. Kibble.
\newblock Lorentz invariance and the gravitational field.
\newblock {\em Journal of Mathematical Physics}, 2:212--221, 1961.

\bibitem[K{"o}36]{Koenig1936}
D{\'e}nes K{"o}nig.
\newblock {\em Theorie der endlichen und unendlichen Graphen}.
\newblock Akademische Verlagsgesellschaft, Leipzig, 1936.

\bibitem[Kol33]{Kolmogorov1933}
Andrei~N. Kolmogorov.
\newblock {\em Grundbegriffe der Wahrscheinlichkeitsrechnung}, volume 2, Heft 3
  of {\em Ergebnisse der Mathematik und ihrer Grenzgebiete}.
\newblock Julius Springer, Berlin, 1933.

\bibitem[Lan82]{rm:Lanford1982}
III Lanford, O.~E.
\newblock A computer-assisted proof of the {Feigenbaum} conjectures.
\newblock {\em Bulletin of the American Mathematical Society}, 6:427--434,
  1982.

\bibitem[Man05]{Manivel2005}
Laurent Manivel.
\newblock Configurations of lines and models of lie algebras.
\newblock arXiv:math/0507118, 2005.
\newblock \href {https://arxiv.org/abs/math/0507118}
  {\path{arXiv:math/0507118}}.

\bibitem[MS64]{rm:MisnerSharp1964}
C.~W. Misner and D.~H. Sharp.
\newblock Relativistic equations for adiabatic, spherically symmetric
  gravitational collapse.
\newblock {\em Physical Review}, 136:B571--B576, 1964.

\bibitem[MS77]{MacWilliamsSloane1977}
Florence~Jessie MacWilliams and Neil J.~A. Sloane.
\newblock {\em The Theory of Error-Correcting Codes}.
\newblock North-Holland, 1977.

\bibitem[Pla01]{Planck}
M.~Planck.
\newblock Ueber das gesetz der energieverteilung im normalspectrum.
\newblock {\em Annalen der Physik}, 309:553--563, 1901.

\bibitem[Sci64]{Sciama1964}
Dennis~W. Sciama.
\newblock The physical structure of general relativity.
\newblock {\em Reviews of Modern Physics}, 36:463--469, 1964.

\bibitem[Ser77]{Serre1977Representations}
Jean-Pierre Serre.
\newblock {\em Linear Representations of Finite Groups}.
\newblock Springer, 1977.

\bibitem[Tak03]{rm:Takesaki}
M.~Takesaki.
\newblock {\em Theory of Operator Algebras II}.
\newblock Springer, 2003.

\bibitem[TMNT21]{Tiesinga2021CODATA}
Eite Tiesinga, Peter~J. Mohr, David~B. Newell, and Barry~N. Taylor.
\newblock {CODATA} recommended values of the fundamental physical constants:
  2018.
\newblock {\em Journal of Physical and Chemical Reference Data}, 50(3):033105,
  2021.
\newblock \href {https://doi.org/10.1063/5.0064853}
  {\path{doi:10.1063/5.0064853}}.

\bibitem[Wie93]{Wien}
W.~Wien.
\newblock Eine neue beziehung der strahlung schwarzer k{\"o}rper zum zweiten
  hauptsatz der w{\"a}rmetheorie.
\newblock {\em Sitzungsberichte der K{\"o}niglich Preussischen Akademie der
  Wissenschaften}, pages 55--62, 1893.

\end{thebibliography}
 
\endgroup
\begingroup
\setlength{\columnsep}{2.2em}\fontsize{9}{10.4}\selectfont
\renewcommand{\indexspace}{\par\vspace{3pt}}
\printindex
\endgroup
\end{document}
